% Options for packages loaded elsewhere
\PassOptionsToPackage{unicode}{hyperref}
\PassOptionsToPackage{hyphens}{url}
\PassOptionsToPackage{space}{xeCJK}
\documentclass[
  12pt,
  fontset=none,oneside,openany]{ctexbook}
\usepackage{xcolor}
\usepackage[a4paper,top=15mm,bottom=15mm,left=15mm,right=15mm,footskip=9mm]{geometry}
\usepackage{amsmath,amssymb}
\setcounter{secnumdepth}{-\maxdimen} % remove section numbering
\usepackage{iftex}
\ifPDFTeX
  \usepackage[T1]{fontenc}
  \usepackage[utf8]{inputenc}
  \usepackage{textcomp} % provide euro and other symbols
\else % if luatex or xetex
  \usepackage{unicode-math} % this also loads fontspec
  \defaultfontfeatures{Scale=MatchLowercase}
  \defaultfontfeatures[\rmfamily]{Ligatures=TeX,Scale=1}
\fi
\usepackage{lmodern}
\ifPDFTeX\else
  % xetex/luatex font selection
  \setmainfont[]{TeX Gyre Termes}
  \setsansfont[]{Microsoft YaHei}
  \setmonofont[]{Consolas}
  \setmathfont[]{Latin Modern Math}
  \ifXeTeX
    \usepackage{xeCJK}
    \setCJKmainfont[]{Microsoft YaHei}
    \setCJKsansfont[]{Microsoft YaHei}
    \setCJKmonofont[]{Microsoft YaHei}
  \fi
  \ifLuaTeX
    \usepackage[]{luatexja-fontspec}
    \setmainjfont[]{Microsoft YaHei}
  \fi
\fi
% Use upquote if available, for straight quotes in verbatim environments
\IfFileExists{upquote.sty}{\usepackage{upquote}}{}
\IfFileExists{microtype.sty}{% use microtype if available
  \usepackage[]{microtype}
  \UseMicrotypeSet[protrusion]{basicmath} % disable protrusion for tt fonts
}{}
\usepackage{setspace}
\makeatletter
\@ifundefined{KOMAClassName}{% if non-KOMA class
  \IfFileExists{parskip.sty}{%
    \usepackage{parskip}
  }{% else
    \setlength{\parindent}{0pt}
    \setlength{\parskip}{6pt plus 2pt minus 1pt}}
}{% if KOMA class
  \KOMAoptions{parskip=half}}
\makeatother
\usepackage{longtable,booktabs,array}
\usepackage{caption}
\captionsetup[table]{skip=6pt}
\newcounter{none} % for unnumbered tables
\usepackage{calc} % for calculating minipage widths
% Correct order of tables after \paragraph or \subparagraph
\usepackage{etoolbox}
\makeatletter
\patchcmd\longtable{\par}{\if@noskipsec\mbox{}\fi\par}{}{}
\makeatother
\setlength{\emergencystretch}{3em} % prevent overfull lines
\providecommand{\tightlist}{%
  \setlength{\itemsep}{0pt}\setlength{\parskip}{0pt}}
\let\NovelUnicodeUnderbrace\underbrace
\let\NovelUnicodeOverbrace\overbrace
\usepackage{mathtools}
\let\underbrace\NovelUnicodeUnderbrace
\let\overbrace\NovelUnicodeOverbrace
\usepackage{adjustbox}
\usepackage{fvextra}
\usepackage{bookmark}
\usepackage{xurl}
\usepackage{newunicodechar}
\newunicodechar{₀}{\ensuremath{{}_0}}
\newunicodechar{₁}{\ensuremath{{}_1}}
\newunicodechar{₂}{\ensuremath{{}_2}}
\newunicodechar{₃}{\ensuremath{{}_3}}
\newunicodechar{ₙ}{\ensuremath{{}_n}}
\newunicodechar{⁰}{\ensuremath{{}^0}}
\newunicodechar{↦}{\ensuremath{\mapsto}}
\newunicodechar{∩}{\ensuremath{\cap}}
\newfontfamily\NovelSymbolFont{DejaVuSans.ttf}
\newunicodechar{�}{{\NovelSymbolFont �}}
\setCJKmainfont{Microsoft YaHei}[Scale=1,BoldFont={Microsoft YaHei Bold}]
\setCJKsansfont{Microsoft YaHei}[Scale=1,BoldFont={Microsoft YaHei Bold}]
\setCJKmonofont{Microsoft YaHei}[Scale=1]
\setCJKfamilyfont{zhhei}{Microsoft YaHei Bold}[Scale=1]
\setCJKfamilyfont{zhsong}{Microsoft YaHei}[Scale=1]
\setCJKfamilyfont{zhkai}{Microsoft YaHei}[Scale=1]
\ctexset{chapter={format=\Large\bfseries\raggedright,beforeskip=8pt,afterskip=16pt}}
\setcounter{secnumdepth}{-1}
\setcounter{tocdepth}{0}
\hypersetup{unicode=true,hidelinks,bookmarksdepth=0,bookmarksopen=true,pdftitle={ggg斗罗同人文}}
\AtBeginDocument{\pdfbookmark[0]{目录}{noveltoc}}
\setlength{\parindent}{2em}
\setlength{\parskip}{0.15\baselineskip}
\setlength{\emergencystretch}{5em}
\setlength{\abovedisplayskip}{5pt}
\setlength{\belowdisplayskip}{5pt}
\setlength{\abovedisplayshortskip}{4pt}
\setlength{\belowdisplayshortskip}{4pt}
\allowdisplaybreaks[4]
\raggedbottom
\pagestyle{plain}
\RecustomVerbatimEnvironment{verbatim}{Verbatim}{breaklines=true,breakanywhere=true,fontsize=\small}
\newsavebox{\NovelMathBox}
\newcommand{\NovelDisplay}[2]{%
  \par\begingroup
  \sbox{\NovelMathBox}{$\displaystyle #2$}%
  \ifdim\wd\NovelMathBox>\linewidth
    \typeout{NOVELWIDE:#1:\the\wd\NovelMathBox:\the\linewidth}%
  \fi
  \[\adjustbox{max width=\linewidth}{\usebox{\NovelMathBox}}\]
  \endgroup
}
\newcommand{\NovelInline}[2]{%
  \begingroup
  \sbox{\NovelMathBox}{$#2$}%
  \ifdim\wd\NovelMathBox>0.55\linewidth
    \typeout{NOVELPROMOTE:#1}%
    \NovelDisplay{#1}{#2}%
  \else
    \ensuremath{#2}%
  \fi
  \endgroup
}
\usepackage{bookmark}
\IfFileExists{xurl.sty}{\usepackage{xurl}}{} % add URL line breaks if available
\urlstyle{same}
% fallback for those not using the hyperref driver hyperxmp:
\makeatletter
\@ifundefined{xmpquote}{\newcommand{\xmpquote}[1]{#1}}{}
\makeatother
\hypersetup{
  hidelinks,
  pdfcreator={LaTeX via pandoc}}

\author{}
\date{}

\begin{document}
\frontmatter

{
\setcounter{tocdepth}{0}
\tableofcontents
}
\setstretch{1.25}
\mainmatter
\chapter{第一章：开端}\label{ux7b2cux4e00ux7ae0ux5f00ux7aef}

唐三出生在附近一个不起眼的小村庄。村里的人都懂点基础的大数理论，提到从前一位强者时，还会说起TREE(3)这个名字。唐三不明白它到底有多大，只知道老人说到这里，都会将声音放低一些。

唐三六岁那年秋天，隔壁的杰克爷爷一手牵着他，一手牵着五岁的妹妹小舞，走了一整天的山路，把兄妹俩送到了学院门口。他爹是村里的铁匠，守着炉子走不开，他们出门的时候，只站在院门口看了一会儿，什么也没说。

那年开春，唐三刚在三河镇的扽魂分殿觉醒了扽魂。一把黑沉沉的锤子，测的人说叫昊天锤，二流，往上练也就那样。

杰克爷爷七十多了，腿脚还算利索。小舞走到半路就走不动了，他把她背起来；过了一会儿唐三也走不动了，他又让小舞下来走一段，换唐三上背。两个孩子换着背，一路上他嘴没闲着，讲故事。

杰克爷爷不识字，可肚子里有故事。早些年镇上来过一个说书先生，在茶棚里讲了一整个冬天，他一天不落地去听，听完就再没忘过。

他讲国王与棋盘上的麦粒，讲数沙子的老先生，讲那些念不完的数名。小舞趴在他背上，听到有趣处便笑，唐三却总追问，故事里那个人究竟怎样算下去。

杰克爷爷能记住故事，细处却答不上来。

``那得问你们学院里的先生。''他说。

走过一道山梁，小舞从他背上溜下来，在路边抓了一把沙。唐三等她将沙漏尽，才拉着她继续往前走。

太阳偏西的时候，他们看见了学院的围墙。

学院门口立着一块石碑，比杰克爷爷还高出两个头。碑上一行一行刻着字，最底下的字刻得又大又深，越往上越小，最上面几行小得看不清。

``杰克爷爷，这上面刻的什么？''小舞仰着头问。

杰克爷爷不识字，摇了摇头。

门口扫地的老门房直起腰来：``这是增长率碑。进这个门的，头一天都得看一遍。''

老门房用扫帚柄点了点最底下的刻痕，说这些是学院用来标示修炼层次的路标。唐三认得下面几个数字，再往上便越来越陌生。

扫帚柄在中间停住，那一行像个小写的w。

``这个叫\allowbreak{}\NovelInline{I00001}{\omega}\allowbreak{}。考核石的刻度，只刻到这里。''

``上面还有呢？''小舞问。

老门房眯着眼往碑顶看了半天。\allowbreak{}\NovelInline{I00002}{\omega}\allowbreak{} 往上还有好多行：\allowbreak{}\NovelInline{I00003}{\omega+1}\allowbreak{}，\allowbreak{}\NovelInline{I00004}{\omega\times 2}\allowbreak{}，\allowbreak{}\NovelInline{I00005}{\omega^2}\allowbreak{}，\allowbreak{}\NovelInline{I00006}{\omega^\omega}\allowbreak{}，再往上是一个像 ε 的字母，再往上的字刻得极小，有的像是一串括号，有的像是一个大写的 Y 后面跟着几个数。最顶上那一行，被风吹雨打得只剩几个模糊的字母。

``上面的，我也看不懂。''他说，``听说一层比一层高，高到有的人一辈子只能看见，摸不着。封号斗罗那样的人物，也才站在那个 Y 字附近。''

他把扫帚放下，冲兄妹俩摆摆手：``进去吧，考核要开始了。''

唐三跟着杰克爷爷往里走，走了几步，又回头看了一眼那块碑。

所有新生都要接受考核，考过才能留下。学院不仅看孩子能否凭直觉凝聚力量，也要在此后的课程里，检验他们能否说清并稳定使用自己的构造。

负责考核的老师乔娜森，实力已经能熟练运用高阶迭代构造大数，看似不高，却已经是很多人一辈子都达不到的高度，她能轻松构造出迭代次方级别的数，远超普通村民的能力范围。

考核用的是一块半人高的考核石。手按上去，把魂力往里一送，石面上就会亮起一道光，光往上爬到哪一行，魂力就是哪一级。石面上的刻度跟门口那块碑一模一样，从 \allowbreak{}\NovelInline{I00007}{0}\allowbreak{} 往上排，最顶上刻着一个 \allowbreak{}\NovelInline{I00008}{\omega}\allowbreak{}，光爬到那儿就到头了，叫满级。

前一个学员将光推到三那一行，围在旁边的老师已露出喜色。轮到下一个时，乔娜森还低头在册子上补刚才的记录。

``这一届学生，恐怕没人能达到增长率 \allowbreak{}\NovelInline{I00009}{4}\allowbreak{} 的级别了。''老师乔娜森感慨道。

``那可不一定，增长率 \allowbreak{}\NovelInline{I00010}{4}\allowbreak{} 的我都见过，不就是隔壁学院的院长吗？他能构造出更高级别的迭代大数，远超我们的想象。''一旁的玉小刚笑着说道。

排在唐三前面的是小舞。她个子矮，踮着脚才把手按到考核石上。光一路往上爬，爬过 \allowbreak{}\NovelInline{I00011}{3}\allowbreak{}，爬过 \allowbreak{}\NovelInline{I00012}{4}\allowbreak{}，在 \allowbreak{}\NovelInline{I00013}{5}\allowbreak{} 那一行停住了。

台下嗡地一声。

``唐三，该你了！''考核老师开口喊道，``施展你的功力让我们看看。''

唐三走上去，把手按了上去。光也是一路往上爬，爬过 \allowbreak{}\NovelInline{I00014}{3}\allowbreak{}，爬过 \allowbreak{}\NovelInline{I00015}{4}\allowbreak{}，最后也停在了 \allowbreak{}\NovelInline{I00016}{5}\allowbreak{}，跟妹妹一样高。

这一下，全场都震惊了！玉小刚瞪大了眼睛，激动地说：``兄妹俩，一个六岁，一个五岁，都推到五！唐三的扽魂还是二流！正常情况下，这个年纪的孩子能推到二就算不错了！''

玉小刚立刻上前，语气恳切：``孩子，可愿意拜我为师？我能教你更高级的大数构造方法，让你变得更强。''

唐三有些懵懂，挠了挠头：``这\ldots\ldots 我感觉我这也很普通啊，就是顺着思路想而已。''

话音刚落，学院院长埃里克就匆匆赶来，他看着唐三，满脸欣慰：``学院很久没有这样的天才学员了，玉小刚是我院最好的教授，擅长高阶大数构造，你跟着他，定能学有所成，构造出更庞大的数。''

唐三连忙躬身：``承蒙学院赏识，愿求教老师。''

考核还没完。

唐三退到一边，后面排着的一个胖乎乎的男孩被叫了上去。他手里攥着半根香肠，上台之前还咬了一口，嚼着就往前走。

``奥斯卡！''乔娜森皱眉，``考核的时候不许吃东西！''

``老师，不吃我使不上劲。''奥斯卡含含糊糊地说。他把最后一口香肠咽下去，手往考核石上一拍。

光呼地往上窜。

窜过 \allowbreak{}\NovelInline{I00017}{3}\allowbreak{}，窜过 \allowbreak{}\NovelInline{I00018}{4}\allowbreak{}、\allowbreak{}\NovelInline{I00019}{5}\allowbreak{}，窜过碑上那一长串数都数不过来的刻度，一口气窜到最顶上那个 \allowbreak{}\NovelInline{I00020}{\omega}\allowbreak{}，停住了。

满级。

台下没有人笑了。乔娜森手里的笔掉在了地上。

奥斯卡回头冲台下咧嘴一笑：``老师，我能再吃一根吗？''

下一个上台的男孩个子不高，肩膀很宽，脸绷得紧紧的。他把手按上去以后，光没有往上窜，而是一格一格地往上走：\allowbreak{}\NovelInline{I00021}{0}\allowbreak{}，\allowbreak{}\NovelInline{I00022}{1}\allowbreak{}，\allowbreak{}\NovelInline{I00023}{2}\allowbreak{}，\allowbreak{}\NovelInline{I00024}{3}\allowbreak{}\ldots\ldots 每走一格都停一下，停稳了，才走下一格。

``太慢了。''有老师小声嘀咕。

可那道光走过的每一格都纹丝不动，稳得像一根钉进墙里的钉子，一点都不晃。走了足有一炷香的工夫，它走到了最顶上：\allowbreak{}\NovelInline{I00025}{\omega}\allowbreak{}。

又一个满级。

乔娜森低头看了看手里的册子：``牢大，扽魂冰红茶\ldots\ldots 品质极高。''她抬起头，又看了他一眼，``品质这么高，怎么推得这么慢？''

``推快了会错。''牢大说。

最后一个被叫到的是个女孩。

``蔡须坤。''

没有人应。

叫到第二遍，那女孩还站在原地，看着台上，好像没听懂这是在叫她。牢大站在她旁边，伸手轻轻碰了一下她的胳膊：``叫你呢。''

她这才走上去。

她上台以后没有马上动，只是站着，看着那块考核石，看了很久，像是在看一样从没见过的东西。然后她伸出手，按了上去。

光一下子就到了顶。

没有往上爬的过程，就是一下子。\allowbreak{}\NovelInline{I00026}{\omega}\allowbreak{} 那一行亮得台下的人都眯起了眼，亮了一瞬，然后整块考核石一下子灭了，灭得干干净净，连一点余光都没剩。

乔娜森愣住了，看了看考核石，又看了看册子。蔡须坤那一栏只记着``品质极高，名称待核''。她能感到那股魂力的非凡，却认不出扽魂的来历，也不敢替它随意定名。

``这\ldots\ldots 也是满级？''

玉小刚盯着那块考核石看了好一会儿，才说：``是满级。考核石到 \allowbreak{}\NovelInline{I00027}{\omega}\allowbreak{} 就到头了。''

他停了一下。

``它刚才那一下，像是还想往上走，走不上去。''

埃里克院长摆摆手：``让她过。''

蔡须坤走下台，走回牢大旁边。牢大小声问她：``你刚才看什么呢？''

``我在看，''她说，``你们是怎么考试的。''

一天之内，三个满级。乔娜森把册子翻过来又翻过去，翻了好几遍，像是怕自己写错了。

散场后，乔娜森没有急着收起册子。她将考核石边缘的铜环重新扣紧，叫来一个负责杂务的高年级学生，先用寻常的魂力试了一遍。石面重新亮起细光，停在靠下的位置，既没有忽然蹿升，也没有熄灭。

``石头没坏？''埃里克问。

``现在看着没坏。''乔娜森蹲下检查接缝，``刚才究竟怎么回事，还得再查。''

她把蔡须坤那一栏翻出来，在旁边添了``待复核''三个小字。有人凑过来，想问那女孩的扽魂到底是什么，乔娜森将册子合上，没有顺着猜。

``品质很高，这一点我看得出来。叫什么，我不知道。''

``那要不要再测一回？''

乔娜森往台下看了一眼。蔡须坤已经站回牢大身旁，牢大正弯腰替她捡起掉到地上的入学凭条。

``先让孩子们安顿。她刚才也没做错什么。''

唐三站得不远，听见了最后一句。他原以为，老师们既然站在台上，就该知道每个人身上的一切。可乔娜森把笔夹回册子时，眉头仍没有松开。

另一边，奥斯卡正蹲着捡那半截掉在地上的香肠。他吹了两口，想往嘴里放，被牢大一把拦住。

``落灰了。''

``洗洗还能吃。''

``那就先洗。''

刚才被围着看的三个满级学员，此刻挤在一只水桶旁边，谁也没有往人群里摆架子。唐三看着他们，攥了一下午的手慢慢松开，掌心已经全是汗。

唐三站在台下，看着那三个人，又低头看了看自己的手。手心里，昊天锤的虚影黑沉沉的，没什么光。

傍晚，老门房听说这一届出了三个满级，扫帚在地上停了很久。

``我在门口扫了四十年地，头一回。''

玉小刚却将三人的册子与唐三兄妹放在一起。他只告诉他们，考核看见的底子，还要靠以后每一次实际运转留下；测到顶，并不等于那块碑已经学完。

埃里克笑着对唐三点头：``好，你先去宿舍，认识一下你的新室友吧，以后你们可以相互切磋，共同进步。''

去宿舍以前，唐三又跑了一趟院门。

杰克爷爷还坐在墙根的石头上。方才送兄妹俩进来时，他便被领去旁边等着，没见着后来的考核，只听见院子里一阵比一阵响的惊呼。见唐三跑来，他忙撑着膝盖站起身。

``能留下了？''

``能。''唐三说，``小舞也能。''

老人一连说了几声好，将搭在肩上的旧布包取下来。里面除了兄妹俩的换洗衣裳，还剩半块已经压扁的干饼。他拿出来，掰开看看，又犹豫着包了回去。

``学院里有饭吃，这个就不留给你们了。''

唐三伸手接过：``您带着，回去还有那么远。''

杰克爷爷笑了笑，把布包的带子替他缩短一截。老人手指粗，那个结打了两次才打紧。唐三站着不动，忽然问，爹是不是知道自己能考上。

``他叫我送你们来。''

``可他都没问考什么。''

``等过年回去，你自己问。''

小舞这时也跑到了门边。她方才嘴里还说要去挑最好的床，看见爷爷要走，却拽住了那根空下来的包带。杰克爷爷摸摸她的脑袋，等她自己松手，才把包推到唐三怀里。

``听先生的话。有事和哥哥说，别一个人跑回村里。''

小舞点了头，又立刻补了一句：``是让他有事和我说。''

老门房在旁边笑了一声。杰克爷爷也笑，伸手指了指渐渐暗下去的山路，说再不走，就要摸黑翻那道梁了。

兄妹俩看着他走过路口。小舞挥了半天手，直到那背影被树遮住，才低头揉了揉鼻子。唐三没有问她怎么了，只将布包换到另一边，腾出靠近她的那只手。

``走吧，先去找床。''

这个学院，一切以实力为准，但这一届似乎有些不一样。

唐三推开宿舍的门，屋里已经有两个人了。奥斯卡盘腿坐在靠窗的床上，腿上摊着一油纸包香肠；牢大在靠门那张床边叠被子，被角折得方方正正，一道褶都没有。

还有第三个人。

小舞站在屋子正中间，两手叉腰。她比唐三早考完，一出考场就跑到这儿来了。

``我是唐三的妹妹，小舞，你们可以叫我乌鲁鲁。''她冲那两个人扬了扬下巴，``这个宿舍，从今日起我来管理！''

奥斯卡嘴里的香肠停住了。牢大叠被子的手也停住了。

``谁让你叫乌鲁鲁的？''唐三说。

``我自己。''

``你住隔壁女生宿舍。''

``那也归我管。''小舞说着，捏了捏手中的燃烧弹，眼神一下子警惕起来，``我不是好欺负的人，我也能构造出远超普通迭代次方的大数。''

``我去你的！''

唐三瞬间大怒，从腰间抽出暗星。

``要不要碰一碰？''

唐三心里有点不服气。他年纪还小，连匕首都用不熟练，可兄妹俩从小打到大，谁也不服谁，两个能构造出高阶大数的人往那儿一站，气场还是十分强烈。本着哥让着妹的想法，唐三没有和小舞争下去。

``乌鲁鲁，今天就到这里吧。''牢大放下被子说。

奥斯卡连忙把油纸包递过去：``吃香肠，吃香肠，一人一根。''

小舞看了看香肠，又看了看唐三，哼了一声，接过一根，收起了气势。

``你们兄妹俩真厉害。''奥斯卡咬着香肠感慨，``扽魂被评成二流，也能硬推到五，真厉害。''

``你靠香肠推到了满级。''牢大说。

``那也是本事。''奥斯卡说。

就这样，四个人算是认识了。

床铺分好，还得去井边提水。小舞才说完这片地方都归自己管，提起空桶便走，过了一会儿却又退回来，问水井在哪。

奥斯卡指向窗外，指了两次，方向还不一样。牢大将被角最后按平，拿起另一只桶：``一起去。''

走到廊下，他们碰见蔡须坤。她捏着入学凭条，正望着墙上的分宿牌，一张脸上没有考核时被全场盯住的紧张，倒是被那几排房号难住了。

小舞伸头看了一眼：``你跟我住那边。''

``是这一间？''

``隔壁。我还没去呢，刚才先管了他们。''

奥斯卡听得直翻白眼。牢大没说话，把蔡须坤的包提过去，将它和小舞的铺盖放到同一张长凳上，免得两人回来又找不到。

井边已经排了几个人。小舞踮着脚想从前面看，牢大伸出手臂拦了一下，她回头瞪他，他只朝队尾偏偏下巴。

``我又没插队。''小舞说着，还是退了回来。

轮到他们时，唐三先把桶放下去。井绳浸过水，比他想的沉，提到一半便磨得手掌发热。他不肯松，手却往下滑了一寸，桶里的水撞在井壁上，响了一声。

牢大从旁边握住绳子，没有整根夺过去，只替他分了些重量。唐三得以换稳手，才一点点将那桶水提上来。

``我自己也能提。''他说。

``知道。''牢大答。

这句话让唐三后面准备好的解释全没了。他愣了一下，蹲下把桶沿的水抹掉，转身去找扁担。奥斯卡已经借来一根，先挂一桶，再去挂另一桶，才发现两头高低差了许多，担子差点把人带歪。

小舞笑得蹲在地上。蔡须坤扶住另一头，等奥斯卡重新摆好肩，才慢慢松手。

回去时，满桶水还是洒出不少。五个人一路提醒，一路争谁方才走得太快，到了宿舍门前，奥斯卡肩上的衣裳已经湿透了。

``下回用魂力提。''小舞说。

``下回先把路让开。''奥斯卡指着她，``刚才拦着笑的就是你。''

牢大将桶搬进门，蔡须坤扶着门等最后一个人进去。唐三放下扁担，手掌里两道红痕还很清楚，心里却不再像刚散场时那样沉。夜饭的钟响起来，他这才想起，自己还不知道饭堂在哪。

小舞从小就是这个脾气。娘走得早，她连娘长什么样都不记得；爹话少，一天说不了几句，她就学会了自己大声说话，先把人吓住再说。唐三知道她外面凶，里头其实怕生。不过课程很快开始，兄妹俩也暂时放下了斗嘴，都想尽快掌握更高级的大数构造方法。

开学头几天，课堂从最基础的练习开始。唐三与小舞做得很快，余下的时间便望着窗外。玉小刚看见了，没有因为两人考核时的成绩，便让他们跳过核对。

当天傍晚，几个人在宿舍门口借光抄完练习。奥斯卡写得快，早早将本子扣住，站起来看唐三的进度。唐三遮了一下，不想让他看见自己划掉重写的那半页。

``我不是来比的。''奥斯卡说，``你刚才那个停顿，我也有过。''

他把自己的本子翻开，前面果然也涂着一块。牢大凑过来，伸手要指出其中一处，奥斯卡却先将纸推回唐三面前，让他自己再找一遍。

蔡须坤坐在廊柱另一侧，迟迟没落笔。小舞以为她不会写，想帮她抄，她却摇头，换了一张纸，只将老师布置的那一小段留下。

等唐三做完，她的纸上也没有多出多少。她拿着那张纸看了又看，最后郑重地折进本子里。小舞不明白这有什么可郑重的，唐三却记得她在考核石前站了那么久，便没有催。

熄灯前，牢大把靠过道的包都往床底收了收，奥斯卡终于将最后半根香肠包好。唐三躺下时还能听见隔壁小舞挪床的声音；过了一会儿，那边也安静了。

第二天的课上，蔡须坤坐在最后一排，一句话没说，只是一直在看。她不怎么看黑板，看的是底下的人：谁先举手，谁偷偷在桌子底下掰手指头，谁算错了脸红。下课的时候，她走到玉小刚跟前，问了一个问题。

``老师，你们人\ldots\ldots 大家都是从一、二、三开始学的吗？''

玉小刚愣了一下：``都是。''

``每一个人都是？''

``每一个人都是。''

蔡须坤点了点头，像是记下了一件很重要的事，转身走了。

往后的普通课程，仍从各自能独立完成的练习开始。

小舞最先耐不住那些重复练习。她做得快，便觉得每天留在课堂上，远没有去后山与哥哥打一场有趣。

``乌鲁鲁！我们该去上课了！''

唐三喊着小舞，可小舞却摆了摆手：``不去不去，天天教的东西太无聊了，都是些基础乘方，谁不会啊。我早就想学习更高级的大数构造方法了，这些基础内容对我来说没用。''

小舞带唐三逃课的事，很快就传遍了学院。玉小刚知道兄妹俩家里没娘，不敢贸然批评小舞，怕伤了她的自尊心，便打算从唐三入手，让唐三劝说小舞。其实小舞也不是完全不爱上课，她特别想去上射击课，她很想摸到枪，但是学院没有。由于没有感兴趣的课程，小舞自然不想去教室。唐三的家庭条件不好，课余时间就找些零工，补贴家用。

就这样，一年多过去，唐三的大数构造能力有了质的飞跃，已经能熟练构造出更高级别的迭代大数，要知道，他这种级别的大数构造能力，对很多人来说，一辈子都突破不了，初级学院毕业的学生，大多也只能停留在基础迭代次方的水平。

一天，玉小刚找到唐三：``小三，老师问你个事，你妹妹小舞，怎么老是不来上课？我知道你们家里的情况，不方便多问，你能不能帮老师劝劝她？你们两个都是百年难遇的天才，在大数构造上很有天赋，可不能浪费天赋啊，只有系统学习，才能掌握更高级的方法，构造出更大的数。''

唐三答应了下来，找到小舞后，他去问：``鲁鲁，你怎么不去上课？是不是有什么难处？课堂上老师会教新的构造技巧，说不定能帮你构造出更大的数。''

小舞撇了撇嘴：``没别的原因，就是太无聊了。都二年级了，还在教乘方，我现在已经能构造出比基础迭代次方更大的数，很快就能突破到更高的层次，就是一直找不到突破的方法，但我能感觉到，快了。''

唐三劝道：``这样下去不是办法，就算无聊，也得去上课，说不定老师讲的某个细节，就能帮你突破瓶颈，掌握更高级的迭代技巧。''其实唐三自己也卡在当前的层次，对''迭代的迭代''理解，始终差了一点，这也是他无法构造出更大数的关键。他的话说完，人也走了。

小舞想了想，觉得唐三说得有道理，便答应重新去上课，只是她的性格依旧孤僻，不怎么和其他同学来往，眼里只有大数构造和提升实力。

可平静的日子没过多久，麻烦就来了。小舞性格怪异，平时总是很消沉、严肃，很多同学都不喜欢她，甚至背地里说她是''冤魂投胎''。一开始，小舞还忍着，可后来，这些同学变本加厉，再也忍不了的小舞，对唐三说：``必须给他们点教训！让他们知道，能构造出更大数的人，不是好欺负的！''

这天，几个同学又来欺负小舞，小舞瞬间爆发，将自己构造的大数能量倾泻而出------那些只能完成基础构造的同学，根本不堪一击，瞬间被能量震倒在地。

``什么？这么厉害！''带头闹事的同学瞬间慌了神，其中一个能构造出基础次方数的同学，被小舞一下子打倒在地，而带头的老大，也只能构造出稍微高阶一点的次方数，根本不敢再嚣张。老大见状，转身就想跑，一回头，牢大不知什么时候已经站在门口，把路堵上了，奥斯卡在他身后探出半个脑袋。老大愣了一下，就被唐三一把抓住。唐三构造的大数级别略微超过小舞，比其他人肯定是强上很多。唐三不想把事情闹大，冷声道：``下次再敢惹事，我就不会这么饶你了。''

说着，一抬手就把老大甩飞了出去。之后，小舞更加努力练功，在之前的基础上不断优化大数构造方法，很快就追上了唐三的水平，距离更高级别的大数构造，只剩下一步之遥。

玉小刚问清双方的经过，处理了那几个闹事的学生，又让唐三与小舞把逃下的课补齐。两人嘴上答应，心里却仍惦记着更高的那些刻痕。

第二天，玉小刚让五个孩子带上纸笔，坐到教室前排。

``门口那块碑，你们都惦记过。''他说，``下面分两堂专门课。先从能算清的箭头讲起，再认一认层级的路标。''

唐三低头磨好笔尖。奥斯卡把香肠收进包里，牢大坐得笔直，小舞挪了挪凳子，蔡须坤则把一张干净的纸铺到了桌上。

\chapter{第二章：高德纳箭头}\label{ux7b2cux4e8cux7ae0ux9ad8ux5fb7ux7eb3ux7badux5934}

专门课开始时，玉小刚没有让五个孩子立刻运转魂力。

他将一粒麦子放到桌上，又在旁边摆了一张画着方格的纸。小舞原本还惦记着前几日那场争执，见老师没有再问，便往前挪了挪凳子。牢大将纸铺平，奥斯卡把香肠包塞到桌底，蔡须坤则望着那粒麦子，等老师往下说。

``来学院的路上，我爷爷讲过棋盘上的麦粒。''唐三说，``可我问怎么算，他就让我来问先生。''

玉小刚点头：``那就借你听过的故事开个头。第一格一粒，往后每格翻倍，先别急着想六十四格，算到第五格。''

\NovelDisplay{M00001R01}{%
1,2,4,8,16,\ldots%
}

``第五格十六粒。''小舞答得很快。

``前五格加在一起呢？''

她顿了一下。奥斯卡在纸上添了几笔，写出三十一。牢大没有往下写，先将``第五格''与``前五格合计''分成两栏。

\NovelDisplay{M00002R01}{%
\text{第}k\text{格的粒数}=2^{k-1}\qquad(k\geq1)%
}

\NovelDisplay{M00003R01}{%
1+2+4+\cdots+2^{k-1}=2^k-1%
}

\NovelDisplay{M00004R01}{%
1+2+4+\cdots+2^{63}=2^{64}-1%
}

``同一个故事，问的东西不同，得到的式子也不同。先弄清在算哪一项，再谈大小。''

唐三记起爷爷在山路上讲故事时，自己总追着问后面还有多少。现在只用一条式子，便能把当时没问完的部分写下。他低头将两个问题都抄全，没把那个减一省掉。

玉小刚又写了两个名字。

\NovelDisplay{M00005R01}{%
\text{古戈尔}=10^{100}%
}

\NovelDisplay{M00006R01}{%
\text{古戈尔普勒克斯}=10^{10^{100}}%
}

``一个是一后面一百个零，另一个是一后面古戈尔个零。即便没有办法把全部数字挤进这间教室，它们仍然是确定的、有限的自然数。''

``那就一直在后面添零？''奥斯卡问。

``当然能继续写大。可你们学构造，要看清生成它的办法。名字叫得再响，也不能代替规则。今日要讲的高德纳箭头，便是一种简洁的写法。''

蔡须坤在纸上写下``数''和``规则''，隔开一小段。玉小刚又让她在旁边留出``层级''的位置，说这一栏等下一堂课再补。

老师让奥斯卡拿出几根香肠。奥斯卡迟疑片刻，只肯拿出三根，又怕小舞伸手，赶紧说这几根还要带回去。

``不吃你的。''玉小刚拿三枚木片代替，``一组里有三枚，摆两组，就是把三个加两回。''

\NovelDisplay{M00007R01}{%
3\times2=3+3=6%
}

\NovelDisplay{M00008R01}{%
3^2=3\times3=9%
}

\NovelDisplay{M00009R01}{%
3^3=3\times3\times3=27%
}

``乘法重复加同样的数，乘方重复乘同样的数。写法都能说成重复，重复的东西却不一样。''

小舞将第二行念成``三个二''，说完自己便停住了。牢大把木片分开，留出两组各三个；又在旁边写了两个相乘的三，让她对着看。

玉小刚没有催，等她重新说清，才添上箭头。

\NovelDisplay{M00010R01}{%
a\uparrow b=a^b%
}

\NovelDisplay{M00011R01}{%
2\uparrow3=2^3=8%
}

\NovelDisplay{M00012R01}{%
3\uparrow2=3^2=9%
}

\NovelDisplay{M00013R01}{%
3\uparrow4=3^4=81%
}

``单箭头就是乘方。左边是底数，右边在这一层表示指数。交换左右两边，通常会改变结果，不能因为数字一样便说它们相同。''

唐三看着八和九，点了点头。他以前见到很长的式子，总先盯着后半截，现在便把底数也圈了起来。

``下面把乘方本身再重复。''玉小刚说。

\NovelDisplay{M00014R01}{%
2\uparrow\uparrow1=2%
}

\NovelDisplay{M00015R01}{%
2\uparrow\uparrow2=2^2=4%
}

\NovelDisplay{M00016R01}{%
2\uparrow\uparrow3=2^{(2^2)}=16%
}

\NovelDisplay{M00017R01}{%
2\uparrow\uparrow4=2^{(2^{(2^2)})}=2^{16}=65536%
}

\NovelDisplay{M00018R01}{%
2\uparrow\uparrow5=2^{65536}%
}

玉小刚在每个式子里的二下方点了小点。小舞数到第四行时，拿笔指着最下方的那个二，问它算不算一层。

``算。这里说塔高四，就是有四个二。不要只数上面的指数，也不要把括号数量当塔高。''

\NovelDisplay{M00019R01}{%
a\uparrow\uparrow b=\underbrace{a^{a^{\cdot^{\cdot^{a}}}}}_{b\text{个}a}\qquad(b\geq1)%
}

``而且，塔从右边算起。''

老师在纸上另写了三个三。奥斯卡算得快，写下了一万九千六百八十三。唐三却皱着眉，没有接着往下抄。

``你先算的是哪边？''他问。

奥斯卡把笔尖移回左边，才发觉自己先把前两个三算成了二十七。

\NovelDisplay{M00020R01}{%
3\uparrow\uparrow3=3^{(3^3)}=3^{27}=7625597484987%
}

\NovelDisplay{M00021R01}{%
\bigl(3^3\bigr)^3=27^3=19683%
}

``第二行的括号已经换了。''玉小刚说，``底数仍是三，三也仍有三个，可它不再是第一行那座塔。写成一条指数链而不加括号时，我们约定从右往左结合。''

奥斯卡望着自己那一行，没再说只是少看了个括号。他将错式留着，在旁边重新算了一遍。

蔡须坤一直没有算那串十进制数字。她先用笔尖点最右边的三，再点它下方的三，最后才点到最左边。

``先把上面那一截算完，下面才知道要乘多少回。''她说。

``是。''玉小刚让她把这句话写在括号旁边。

小舞望着第五层的二，问既然没写出十进制结果，是不是就没算完。

``要看问题要求。若问它是什么，写出精确的指数式已经把它确定了；若问末位数字、大小关系或别的性质，再找相应办法。不能把没铺满纸的精确式，当作一个模糊的`很大'。''

牢大将几个二的例子重新排好，最下面留了一行空处。他想往后接第六层，抬头看见老师已经开始写三箭头，才先把笔收住。

``先看底数二。别一开始便挑最大的。''

\NovelDisplay{M00022R01}{%
2\uparrow\uparrow\uparrow1=2%
}

\NovelDisplay{M00023R01}{%
2\uparrow\uparrow\uparrow2=2\uparrow\uparrow2=4%
}

\NovelDisplay{M00024R01}{%
2\uparrow\uparrow\uparrow3=2\uparrow\uparrow4=65536%
}

\NovelDisplay{M00025R01}{%
2\uparrow\uparrow\uparrow4=2\uparrow\uparrow65536%
}

奥斯卡起初还跟着心算，到最后一行，手便停了。

``这回六万五千五百三十六不是指数，是塔高？''

``是双箭头右边的输入。照双箭头的定义，它表示那座塔里有多少个二。''

小舞低头数前三行的箭头，差点又把第三行写成一座只有三个二的塔。牢大将她的纸推回上一节，让她先比一比双箭头的三和三箭头的三。

\NovelDisplay{M00026R01}{%
2\uparrow\uparrow3=16%
}

\NovelDisplay{M00027R01}{%
2\uparrow\uparrow\uparrow3=65536%
}

``三没有变，箭头数变了，所以三在不同规则里做的事也变了。''唐三说。

玉小刚点头，随后换回底数三。

\NovelDisplay{M00028R01}{%
3\uparrow\uparrow\uparrow1=3%
}

\NovelDisplay{M00029R01}{%
3\uparrow\uparrow\uparrow2=3\uparrow\uparrow3%
}

\NovelDisplay{M00030R01}{%
3\uparrow\uparrow\uparrow3=3\uparrow\uparrow(3\uparrow\uparrow3)%
}

\NovelDisplay{M00031R01}{%
3\uparrow\uparrow\uparrow3=3\uparrow\uparrow7625597484987%
}

``先把里面那一次双箭头算出，结果再交给外面的双箭头。这里不是把原来那座塔简单加高三层，是把原来的结果整个拿来作新高度。''

蔡须坤在内外两处各画了一小道弧。奥斯卡凑近看，才发现她不是把式子拆成两个可以随便调换的位置，而是在标出先后。

``里面还没做完，外面就等着。''她说。

牢大听到这里，拿起笔，将自己上一页省掉的一对括号补了回去。

``四箭头用三箭头，三箭头用双箭头。这样下去，每一层都另外背一段吗？''牢大问。

玉小刚先擦去一小块地方，写下一列较容易看清的数。

\NovelDisplay{M00032R01}{%
a_0=a_1=1,\qquad a_{n+2}=a_{n+1}+a_n%
}

\NovelDisplay{M00033R01}{%
1,1,2,3,5,8,\ldots%
}

``这列数不用逐个取名，因为起点和接法已经交代清楚。递归也是这个道理：说明从哪里开始，再说明如何回到已经定义的情形。''

``只说一直重复，还不够？''小舞问。

``你让牢大一直重复，却不说第一步拿什么，他从哪里下手？''

小舞看了牢大一眼，自己先笑了。牢大没有笑，正在把``起点''两个字写到纸上。

玉小刚规定，本节底数\allowbreak{}\NovelInline{I00028}{a}\allowbreak{}为不小于二的整数，箭头数\allowbreak{}\NovelInline{I00029}{m}\allowbreak{}为正整数，右侧输入\allowbreak{}\NovelInline{I00030}{b}\allowbreak{}为非负整数；\allowbreak{}\NovelInline{I00031}{\uparrow^m}\allowbreak{}表示\allowbreak{}\NovelInline{I00032}{m}\allowbreak{}个箭头，不是把一支箭头拿去乘方。

\NovelDisplay{M00034R01}{%
a\uparrow^1 b=a^b%
}

\NovelDisplay{M00035R01}{%
a\uparrow^m0=1%
}

\NovelDisplay{M00036R01}{%
a\uparrow^{m+1}(b+1)=a\uparrow^m\bigl(a\uparrow^{m+1}b\bigr)%
}

``第三条里，先把右侧输入往回退，得到内部结果，再交给少一支箭头的运算。最后会回到前两条。底数在这一过程中保持原样。''

牢大请老师挑一个能完整回退的例子。玉小刚便没有再换大数，只选了已经算过的双箭头。

\NovelDisplay{M00037R01}{%
2\uparrow\uparrow0=1%
}

\NovelDisplay{M00038R01}{%
2\uparrow\uparrow1=2\uparrow1=2%
}

\NovelDisplay{M00039R01}{%
2\uparrow\uparrow2=2\uparrow(2\uparrow\uparrow1)=2^2=4%
}

\NovelDisplay{M00040R01}{%
2\uparrow\uparrow3=2\uparrow(2\uparrow\uparrow2)=2^4=16%
}

小舞这次从最下行开始往回找，找到零时，伸手点住了那个一。

``所以不是把空塔当成零。''

``本章定义取一。起点不能凭眼熟随意改。''

唐三又沿三箭头试了一小段，发现必须分清眼下到底在哪一层，而不能见到右边的数字减一，便把外面的箭头也一起改掉。

老师让他将刚才出错的位置留着。几个人围过来各看一遍，蔡须坤指向最里面的括号，牢大则在外侧标明原来的箭头数。唐三才重新把那一次回退说全。

四箭头仍照相同规则。

\NovelDisplay{M00041R01}{%
2\uparrow^4 1=2%
}

\NovelDisplay{M00042R01}{%
2\uparrow^4 2=2\uparrow^3 2=4%
}

\NovelDisplay{M00043R01}{%
2\uparrow^4 3=2\uparrow^3 4=2\uparrow\uparrow65536%
}

``箭头多，结果便一定严格变大？''奥斯卡问。

玉小刚让他先看中间一行。底数二、右边也是二时，换了这几种箭头，结果竟然相同。若只背``多一支就一定更大''，连眼前的例子也解释不了。

``比较之前，要把底数、箭头数、右侧输入都看全。某些小输入会碰巧相等；讨论增长，还要看输入变化后的整条规则。''

黑板上的数字暂时没擦。玉小刚让他们各改一个位置，看看自己到底改了什么。

\NovelDisplay{M00044R01}{%
2\uparrow^2 3=16%
}

\NovelDisplay{M00045R01}{%
3\uparrow^2 3=7625597484987%
}

\NovelDisplay{M00046R01}{%
2\uparrow^3 3=65536%
}

``第二行只换底数，第三行只换箭头数。''唐三说，``不能把三处位置都笼统叫成层数。''

小舞又将第一行右边的三改成四，发现这次竟与第三行碰到了一样的结果。

\NovelDisplay{M00047R01}{%
2\uparrow\uparrow4=2\uparrow\uparrow\uparrow3=65536%
}

``所以两种规则其实一样？''

老师没直接回答，让她将两种规则的输入都改成二：第一种仍在输入后添一，再求双箭头；第二种直接用三箭头。

\NovelDisplay{M00048R01}{%
2\uparrow\uparrow(2+1)=16%
}

\NovelDisplay{M00049R01}{%
2\uparrow\uparrow\uparrow2=4%
}

小舞看了一会儿，将自己刚才的话划去。同一输入上碰巧相等，换个输入便不一样，不能凭一个值断定两条函数规则完全相同。

奥斯卡说自己先前切磋时，也常见到两下招式声势差不多，便以为里面接法一样。玉小刚让他先把纸上的判断做好，真正用进魂技，还要另看条件，不能只凭眼睛觉得像。

蔡须坤将那三处分别标作底数、箭头数、右侧输入。牢大看过，又让她指出递归时哪里保持原样、哪里在变。她这次没有顺手越到下一层，而是一处处指完，等他核清才收笔。

课间，奥斯卡从包里拿出另一页笔记，问上面的阿克曼函数能不能与箭头对照。玉小刚没将纸收走，只叫他先把采用的定义写完整。

\NovelDisplay{M00050R01}{%
A(0,n)=n+1%
}

\NovelDisplay{M00051R01}{%
A(m+1,0)=A(m,1)%
}

\NovelDisplay{M00052R01}{%
A(m+1,n+1)=A\bigl(m,A(m+1,n)\bigr)%
}

``这里两个参数都取非负整数。不同书若换了起点，具体数值也可能跟着变；今天按这三条计算。''

牢大先试零行。小舞想直接跳到最后一行，玉小刚却把纸挡住，让她先陪牢大走一遍。

\NovelDisplay{M00053R01}{%
A(1,0)=A(0,1)=2%
}

\NovelDisplay{M00054R01}{%
A(1,1)=A(0,A(1,0))=A(0,2)=3%
}

\NovelDisplay{M00055R01}{%
A(1,n)=n+2%
}

``里面那次的结果，还是送到外面去。''蔡须坤说。

她记得三箭头里那两道弧，在这页上也标了同样的先后。奥斯卡沿着她画的位置，将第二行展开给小舞看。

\NovelDisplay{M00056R01}{%
A(2,0)=A(1,1)=3%
}

\NovelDisplay{M00057R01}{%
A(2,1)=A(1,A(2,0))=A(1,3)=5%
}

\NovelDisplay{M00058R01}{%
A(2,2)=A(1,A(2,1))=A(1,5)=7%
}

\NovelDisplay{M00059R01}{%
A(2,n)=2n+3%
}

``多走一层，原来那一整行便被拿来组织新的调用。''玉小刚说。

\NovelDisplay{M00060R01}{%
A(3,0)=A(2,1)=5%
}

\NovelDisplay{M00061R01}{%
A(3,1)=A(2,A(3,0))=A(2,5)=13%
}

\NovelDisplay{M00062R01}{%
A(3,n)=2^{n+3}-3%
}

\NovelDisplay{M00063R01}{%
A(4,n)=2\uparrow\uparrow(n+3)-3%
}

\NovelDisplay{M00064R01}{%
A(4,0)=2\uparrow\uparrow3-3=13%
}

\NovelDisplay{M00065R01}{%
A(4,1)=2\uparrow\uparrow4-3=65533%
}

奥斯卡在双箭头后面圈住\allowbreak{}\NovelInline{I00033}{n+3}\allowbreak{}。他原先只记了``阿克曼会到指数塔''，若真拿原来的输入直接作塔高，眼前这些值便全对不上。

``名称帮你认方向，精确比较还是回定义。''玉小刚说，``换一种记法，若是在表示同一个结构，不能只凭名字不同，就当成凭空多了一份力量。''

唐三将两页纸并排，第一次觉得这些陌生记号未必是在争谁更神秘。有些地方只是原来的接法换了一身衣裳，需要耐心把边上的偏移也一并核对。

午后，小舞有些坐不住，指向黑板上的三箭头，问把它练熟以后，是不是只剩再加一支箭头。

``可以逐层增加，也可以让层数本身随输入变化。''

老师让五个人各写一行。牢大负责一箭头，奥斯卡负责二箭头，蔡须坤写三箭头，唐三和小舞接着往下写。每一行先固定箭头数，左右两端却都留下输入的位置。

\NovelDisplay{M00066R01}{%
h_m(n)=n\uparrow^m n\qquad(m\geq1,\ n\geq2)%
}

\NovelDisplay{M00067R01}{%
h_1(3)=3\uparrow3%
}

\NovelDisplay{M00068R01}{%
h_2(3)=3\uparrow\uparrow3%
}

\NovelDisplay{M00069R01}{%
h_3(3)=3\uparrow\uparrow\uparrow3%
}

``现在不永远守着某一行。输入是多少，就取第几行在这个输入处的值。''

\NovelDisplay{M00070R01}{%
h(n)=h_n(n)=n\uparrow^n n%
}

\NovelDisplay{M00071R01}{%
h(2)=2\uparrow^2 2=4%
}

\NovelDisplay{M00072R01}{%
h(3)=3\uparrow^3 3%
}

\NovelDisplay{M00073R01}{%
h(4)=4\uparrow^4 4%
}

小舞拿着笔，先点第三行，再往右挪到三的位置，接着又移向第四行。她这才明白，先前增加右侧输入，与这里连选用哪一行也一起改变，是两件事。

``这是对角化的一种例子。''玉小刚说，``对于眼前这列箭头构造，任意先固定的箭头数，在输入充分增大以后，都会被这种取法越过去。不是把无穷多支箭头塞进某一次计算。''

``每一次仍有确定的箭头个数。''牢大说。

``是。也不能拿一列随便乱排的函数，抄个对角就宣称它一定更快。结论要依靠你选的这一列怎样增长。''

唐三在纸上留下两种写法，一种固定行号，一种让行号随输入动。他觉得自己已经看懂，可玉小刚让他合起本子，换一个输入独立说一遍，他仍在中间停住了。

老师没有替他接上，只让他将停住的位置记下来。看懂黑板上的一遍，与自己每一次都能走清，果然还隔着一段路。

临近下课，奥斯卡终于问起考核时听过的那个名字。

玉小刚没有继续把黑板上的数字写长，而是只留出三行。

\NovelDisplay{M00074R01}{%
g_1=3\uparrow^4 3%
}

\NovelDisplay{M00075R01}{%
g_{k+1}=3\uparrow^{g_k}3\qquad(k\geq1)%
}

\NovelDisplay{M00076R01}{%
G=g_{64}%
}

``第一项用四支箭头。以后每一项，把上一项的整个结果拿来决定下一项有多少支箭头，做到第六十四项，便得到葛立恒数。''

小舞盯着下标：``六十四不是六十四支箭头。''

``对，是这一串构造里的第六十四项。''

牢大又指向第二行：``上一项放到了箭头数的位置，左右两端还是三。''

``对。构造规则说得很短，所确定的数却很大。阿克曼的某个值、葛立恒数、一个不断接收输入的函数，不能因为都很大便混写成同一件东西。''

唐三看着那三行，又低头看自己的手：``既然我已经能把定义说出来，是不是也算能构造它了？''

玉小刚没有让他立刻送出魂力，而是将本子合上，请他说明第二项该怎样从第一项得到。唐三能答，问到其中一层箭头如何回退，却又要翻前面的纸。

``你已经知道这套规则指什么，这是进步。接着要能独立推演，再把该走的部分接进自己的运转。像你会画一座高楼，也认得梁柱安在哪里，真要把它盖起来，仍得从地基做。''

牢大听见地基，点了点头。奥斯卡将写错的那页拿回来，没有撕掉；小舞也暂时收起想立刻去后山试一回的念头，在不会独立展开的地方留了一个小记号。

老师让他们在本子末尾留下速查。小舞原本想照着黑板多抄几遍箭头，后来改成先写自己最容易弄错的那个括号。

{\def\LTcaptype{none} % do not increment counter
\begin{longtable}[]{@{}ll@{}}
\toprule\noalign{}
记号或问题 & 先看哪里 \\
\midrule\noalign{}
\endhead
\bottomrule\noalign{}
\endlastfoot
单箭头 & 就是乘方，左右输入不能随意交换 \\
双箭头 & 塔从右算，高度数底数的个数 \\
三箭头 & 内层结果用来作外层双箭头的输入 \\
更多箭头 & 核对零输入的起点，再按统一递归回退 \\
换用阿克曼记法 & 起点、参数和偏移项要一起核对 \\
对角化 & 固定一行，与让行号随输入变化分开 \\
葛立恒数 & 第64项，上一项决定下一项的箭头数 \\
\end{longtable}
}

奥斯卡收好纸，才发现从早上拿出来的香肠一根也没少。他把其中一根分给牢大，另一根掰开递到后面。蔡须坤先将纸上的内外两道弧补好，才伸手接过。

走出教室，小舞跑去看增长率碑。她看了一会儿，指着那个小小的\allowbreak{}\NovelInline{I00034}{\omega}\allowbreak{}，问箭头一直加下去，是不是就能把它算出来。

玉小刚经过门边，听见了，便让她先把这个问题留着。

``下一堂，我们不添箭头。先看那些路标究竟在记什么。''

唐三翻到纸上还空着的第三栏，给``层级''两个字描了一遍，夹回本子里。牢大走在前面，等奥斯卡把散掉的纸角压好，五个人才一同往饭堂去。

\chapter{第三章：队伍尽头的路标}\label{ux7b2cux4e09ux7ae0ux961fux4f0dux5c3dux5934ux7684ux8defux6807}

下一堂专门课，玉小刚把五个人带到了走廊。

小舞怀里还抱着上回的笔记，想着老师是不是要让他们沿增长率碑往上指。结果玉小刚只借来一张长凳，又叫他们搬几张小凳，按顺序摆开。

``今天讲序数。''

``先搬凳子？''奥斯卡问。

``先看先后。''

牢大把最歪的一张扶正。蔡须坤走到凳子旁，伸手摸了摸最前面那张的靠背，没有坐下。唐三翻到纸上那一栏空着的``层级''，听见玉小刚让他们先把笔放一放。

玉小刚从奥斯卡那里借来三枚铜钱。奥斯卡攥着不肯全松手，直到老师答应下课归还，才将它们放到长凳上。

老师拿粉笔作了三个记号，让铜钱从左到右排好，再请小舞将它们倒过来。

``铜钱换了吗？''

``没换。''

``记号谁在前面，换了吗？''

``换了。''

``所以，`有哪些东西'和`东西按什么顺序排'，可以分别来问。谈集合，列举时先写谁，通常不改变里面有哪些元素；谈一个已经指定的顺序，就要把先后认真留下。''

唐三拿笔记下。玉小刚又提醒，换记号与换结构也要分清。若只看三个位置依次排开的样子，它仍是一条三项的队伍。

``我们既不只盯着东西的名字，也不随意忘记先后。序数描述的，是良序的顺序类型。''

小舞皱了皱眉。玉小刚便将凳子上的三枚铜钱换成三块木片，沿原来的位置各放一块。

``原来第一个对应现在第一个，第二个对应第二个，前后关系全保住。说两种排列具有相同顺序类型，便是要求有这种对应，且两边都不遗漏。''

奥斯卡先伸手拿回铜钱，才腾出另一只手，指指前面的凳子：``不管坐谁，前面一个总在后面一个之前。''

``是。今天我们会遇到的不同，往往就在谁前面还有什么，后面又接着什么。''

五张凳子摆好以后，小舞坐上最后一张。玉小刚让她只在第二张和最后一张里挑，又让牢大指出这两张中哪张在先。

牢大没有走到全队最前面，只点了第二张。

``说一条顺序是良序，不能只看整个队伍有开头。它还要求：从中任意挑出一个非空部分，按原来的先后，都能找到那部分最前面的一个。''

``只挑一张呢？''蔡须坤问。

``那一张就是它自己这一部分的最前面。''

老师写下通常顺序的自然数。

\NovelDisplay{M00077R01}{%
0<1<2<3<\cdots%
}

``任选一些自然数，只要不是一个也不选，总能找出其中最小的。可能不包括零，也可能不包括前面许多项，但选出来的那部分仍有自己的开头。''

\NovelDisplay{M00078R01}{%
7<8<9<\cdots%
}

``这一部分从七开始。''唐三说。

``再只挑双数？''

\NovelDisplay{M00079R01}{%
0<2<4<6<\cdots%
}

小舞答零。玉小刚又问，若把所有自然数改按从大到小的方向排，是否仍有同样性质。

奥斯卡想说从最大的开始，话到嘴边便停住了。自然数没有最大的；他无论挑出多少，后面总还能再加一，而按这个倒过来的方向，新得到的那个便要排得更早。

``整个非空部分都找不到最前面的。''牢大说。

``所以不是良序。先说清排列，再检查它是不是良序，不能见到一条队伍便随意贴序数。''

走廊里只有五张小凳。老师说它们只是示意，后面要用规则描述无限的排列，并非真的摆完无限多张。

玉小刚在纸上画了一条长队。开头几人写着零、一、二、三，后面的方向只留一串省略号。

``每个人都在某个有限的位置。你指出任意一个，他后面都还有下一个。这整条排列的顺序类型，记作\allowbreak{}\NovelInline{I00035}{\omega}\allowbreak{}。''

\NovelDisplay{M00080R01}{%
0,1,2,3,\ldots%
}

\NovelDisplay{M00081R01}{%
\omega=\text{自然数按通常先后排列的顺序类型}%
}

``那最后一个人站在哪？''小舞问。

``这里没有最后一个。''

``可你又说整条。''

``我们描述的是由规则确定的完整排列，不是在讲一个人怎样一步步走到某个尽头。说每个人都有下一个，与说全队存在最后一个人，是两件事。''

牢大伸手点住纸上的第三个位置。老师问他后面有没有人，他说有；再点一个更远的位置，答案仍一样。

``任意给定的有限位置都能找到，却没有`最大的有限位置'，再由它加一得到\allowbreak{}\NovelInline{I00036}{\omega}\allowbreak{}。''

\NovelDisplay{M00082R01}{%
0<1<2<\cdots<n<\cdots<\omega%
}

``\allowbreak{}\NovelInline{I00037}{\omega}\allowbreak{}是第一个无限序数。它大于每个有限序数，却没有一个与它紧挨着的最大有限序数。''

唐三想起妹妹前一堂课问过的箭头，便问若塔高写得足够大，能否把这段空隙填上。

玉小刚让他任选有限塔高与有限的箭头数。按上一堂的规则，结果仍是有限自然数，再添一也仍有限。

``所以，\allowbreak{}\NovelInline{I00038}{\omega}\allowbreak{}不是一个特别难算、只是还没写出全部数字的自然数。它在这里记录的排列结构，与你上一堂计算的有限数值有区别。''

蔡须坤将上回留着的第三栏补上。她没有把前两栏划掉，而是各指了一遍，确认它们说的并非同一件事。

小舞听见``没有最后一个''，便把最远处的那张小凳又往后拖了一截。

``要是我说，这张就放在所有人后面呢？''

``那便是另一个排列。''

玉小刚给新添的那位画了一面小旗。前面的自然数队伍保持原样，小旗在每一个原有位置之后。

\NovelDisplay{M00083R01}{%
0,1,2,3,\ldots,\star%
}

\NovelDisplay{M00084R01}{%
\omega+1%
}

``这里有最后一位，原先的\allowbreak{}\NovelInline{I00039}{\omega}\allowbreak{}型队伍却没有。因此两边无法作完整的保序对应，不是同一种顺序类型。''

奥斯卡提议让旗手先往前挪，给后面的人重新编号。牢大望着纸，伸手挡住他的笔。

``挪到前面，顺序就换了。''

``对。给位置换一个号码，与把末尾的人搬到队首，不能混为一谈。''

老师另画一列，这次旗手原本便在所有自然数位置之前。

\NovelDisplay{M00085R01}{%
\star,0,1,2,3,\ldots%
}

\NovelDisplay{M00086R01}{%
1+\omega=\omega%
}

``旗手对应新队伍的零号，原来的零号对应一号，原来的一号对应二号。谁在谁前面，一点也没颠倒。两边所有位置又都能照顾到，所以它与原来的自然数队伍是同一种顺序类型。''

小舞终于不再拖凳子。她蹲在纸边，指指两幅图：一幅永远能找见下一个，一幅多出一个没有后继的末尾。

\NovelDisplay{M00087R01}{%
1+\omega=\omega<\omega+1%
}

\NovelDisplay{M00088R01}{%
\omega+2=(\omega+1)+1%
}

``序数加法的做法，是保留前一段，再把后一段整段接到后面。有限时你习惯交换加数，到这里却不能照搬。''

老师让五个人各画一幅图。奥斯卡总想把旗子画得更显眼，蔡须坤却只在末尾轻轻点了一点，说明重要的是它在所有原有位置之后，换成旗子、凳子还是木片都一样。

唐三在两幅图下各写了一句：第一种没有末尾，第二种有。等老师再问，他终于不用回头找那面旗了。

``现在把一张小凳换成一对。''玉小刚说。

奥斯卡同牢大并排摆好两块木片，小舞再摆下一对，蔡须坤往后接。纸上便出现成对的记号，一对之后还有下一对。

\NovelDisplay{M00089R01}{%
\underbrace{\bullet\ \bullet}_{\text{第0对}}\quad
\underbrace{\bullet\ \bullet}_{\text{第1对}}\quad
\underbrace{\bullet\ \bullet}_{\text{第2对}}\quad\cdots%
}

``每对两项，共有一条\allowbreak{}\NovelInline{I00040}{\omega}\allowbreak{}型的对序列。照次序展开，仍是一个接一个，所有位置前面都只有有限项，也没有最后一项。''

\NovelDisplay{M00090R01}{%
2\times\omega=\omega%
}

``序数乘法\allowbreak{}\NovelInline{I00041}{\alpha\times\beta}\allowbreak{}，是按\allowbreak{}\NovelInline{I00042}{\beta}\allowbreak{}的先后接起\allowbreak{}\NovelInline{I00043}{\alpha}\allowbreak{}型的段。左边看段内，右边看段与段。''

牢大没有急着念公式，而是将自己摆的那一对圈好，标明这是被复制的小段。

``那把整条自然数长队当一段，接两段。''唐三说。

玉小刚让他用两种不同的记号区分。第一条全是空心圈，第二条全是实心点；两条各自都没有最后一项，第二条的第一项却在第一条所有位置之后。

\NovelDisplay{M00091R01}{%
\underbrace{\circ\ \circ\ \circ\ \cdots}_{\text{第一段}}\quad
\underbrace{\bullet\ \bullet\ \bullet\ \cdots}_{\text{第二段}}%
}

\NovelDisplay{M00092R01}{%
\omega\times2=\omega+\omega%
}

``第二段的开头，与\allowbreak{}\NovelInline{I00044}{\omega+1}\allowbreak{}里那一位有什么不同？''蔡须坤问。

``在\allowbreak{}\NovelInline{I00045}{\omega+1}\allowbreak{}里，那一位已是末尾；这里第二段开头后面还有整条长队，所以也不能混同。''

小舞将第二段遮住，只露出第一项，随后一点点把手移开。原来一个小小的乘二，竟然接来了一整段，没有谁因为``都是在后面添''便可以省掉区别。

\NovelDisplay{M00093R01}{%
\omega+1<\omega\times2<\omega\times3<\cdots%
}

``交换乘号两边，会把每段两项换成两段长队，必须先分清段内与段间。''

玉小刚将纸转了一个方向，画了几排书架。

``想象每一排的书按自然数次序排下去，每一排都是\allowbreak{}\NovelInline{I00046}{\omega}\allowbreak{}型。先比排号，排号相同，再比这一排里的位置。第一排全在第二排之前，第二排全在第三排之前。''

老师只画开头几本，剩下用省略号留下。小舞伸手指向第二排，说自己可以先把那里第一本搬来，夹进第一排两本之间。

``搬当然能搬，但搬过以后，就改了我们刚刚指定的顺序。现在先照原来的规则看。''

\NovelDisplay{M00094R01}{%
\omega\times\omega=\omega^2%
}

\NovelDisplay{M00095R01}{%
\text{第}j\text{排的第}i\text{个位置：}\quad\omega\times j+i%
}

\NovelDisplay{M00096R01}{%
i,j\text{均为有限自然数}%
}

``第一排前面没有完整长排，第二排前面有一排，第三排前面有两排。式子把排号和排内位置分开。''

\NovelDisplay{M00097R01}{%
\omega\times2+5<\omega\times3%
}

``第三排里的第几个，与下一排刚开头，也不能仅凭尾巴那个五来比。''牢大说。

玉小刚点头。牢大把自己的图推给奥斯卡，两人一个报排号，一个指出位置，比刚才只盯着符号时快了许多。

``如果再将这样的整座书架当一层，层与层也按自然数次序排列，先比层号，再比排号，最后比排内的位置呢？''

\NovelDisplay{M00098R01}{%
\omega^2\times\omega=\omega^3%
}

\NovelDisplay{M00099R01}{%
\omega^3\times\omega=\omega^4%
}

``每添一层分组，便多一层先后。这是描述排列的规则，不是让人真用无限时间去搬书。''

蔡须坤将``先比层、再比排、再比位置''写在边上。唐三看了，又把先前那条一对一对的队伍翻回来，终于没有把它与这座书架画成同一条直线。

小舞问一直增加层数以后，是否还能留下一个统一的路标。

\NovelDisplay{M00100R01}{%
\omega<\omega^2<\omega^3<\cdots%
}

\NovelDisplay{M00101R01}{%
\omega^\omega=\sup\{\omega^n:n<\omega\}%
}

``\allowbreak{}\NovelInline{I00047}{\sup}\allowbreak{}读作上确界，是不低于列中每一项的最小上界。\allowbreak{}\NovelInline{I00048}{\omega^\omega}\allowbreak{}超过这里每一个固定有限指数的位置，却不是挑一个特别大的有限层数来顶替。''

窗外下起细雨，奥斯卡起身将靠外侧的纸挪进来。图上的书架没有画完，规则却已经说完了，他这次没有再为省略号后面是不是漏了一页而着急。

回到教室，玉小刚没有收走他们的图。他在书架旁添了三座完整的大架，又在后面添两条长排，最后放五个单独的位置。

\NovelDisplay{M00102R01}{%
\omega^2\times3+\omega\times2+5%
}

``从前往后：三段\allowbreak{}\NovelInline{I00049}{\omega^2}\allowbreak{}型排列，接两段\allowbreak{}\NovelInline{I00050}{\omega}\allowbreak{}型排列，再接五项。不能把最后的五挪到最前面，还说只是换了个写法。''

``那跟四座大架呢？''奥斯卡问。

\NovelDisplay{M00103R01}{%
\omega^2\times3+\omega\times2+5<\omega^2\times4%
}

``前三段一样。左边余下的是两排再添五项，还没有达到另一整座\allowbreak{}\NovelInline{I00051}{\omega^2}\allowbreak{}型排列。先比最高层，再看同层有几段，最后才看尾部。''

牢大按顺序划去相同部分，留下真正不同的地方。小舞也不再只看尾部最大的有限数字。

``用递减的ω指数、合并同层的段数来写，便是在学习康托尔标准形。先认清它怎样组织。''

唐三问加法是不是只要把两段尾巴拼上。玉小刚让他回到之前的队首添人。

\NovelDisplay{M00104R01}{%
5+\omega=\omega%
}

\NovelDisplay{M00105R01}{%
\omega+\omega^2=\omega^2%
}

\NovelDisplay{M00106R01}{%
\omega^2+\omega>\omega^2%
}

``接在后面的段，可能使前面的某些尾部在顺序类型里被吸收。要看整条排列，不能拿普通多项式的交换法则照搬。''

\NovelDisplay{M00107R01}{%
\bigl(\omega^2+\omega+5\bigr)+\omega=\omega^2+\omega\times2%
}

小舞先盯住不见了的五，觉得像少算。牢大将那五项画在新接来的长队前，给它们一起重新编号，才让她看出，这一小段与后面那条没有末项的长队接起，又变成一条\allowbreak{}\NovelInline{I00052}{\omega}\allowbreak{}型的排列。

原来不是把五个位置凭空擦掉，而是比较顺序类型时，它们并没有使后面多出一个新的末尾。

午后，玉小刚让他们分别指出\allowbreak{}\NovelInline{I00053}{\omega+3}\allowbreak{}与\allowbreak{}\NovelInline{I00054}{\omega\times2}\allowbreak{}前面紧挨着的那个序数。

小舞很快答出前一个。后一个，她从第二段开头往回找，想找出第一段的最后一位，最后只得停笔。

\NovelDisplay{M00108R01}{%
\omega+3=(\omega+2)+1%
}

``前一个有紧前项，是后继序数。第二个却不能写成某个序数加一，因为第一段没有最后一位。像这样非零、又不是后继的序数，叫极限序数。零另外列作起点。''

\NovelDisplay{M00109R01}{%
\omega,\quad\omega\times2,\quad\omega^2,\quad\omega^\omega%
}

``若要在功法里用极限下标控制计算，便得交代怎样选较低的位置。我们给本章要用的极限选一条严格递增、上确界等于它的基本列，第\allowbreak{}\NovelInline{I00055}{n}\allowbreak{}项记作\allowbreak{}\NovelInline{I00056}{\lambda[n]}\allowbreak{}。''

老师先在书架图上指出排号，再写下本章采用的约定。

\NovelDisplay{M00110R01}{%
\omega[n]=n%
}

\NovelDisplay{M00111R01}{%
\bigl(\omega\times k\bigr)[n]=\omega\times(k-1)+n\qquad(1\leq k<\omega)%
}

\NovelDisplay{M00112R01}{%
\omega^2[n]=\omega\times n%
}

\NovelDisplay{M00113R01}{%
(\omega^\omega)[n]=\omega^n%
}

``第一条选有限位置，第二条在最后一段内往后，第三条增加整排的段数，第四条增加有限层数。每项都低于原极限，整列的上确界却正是它。''

奥斯卡将几条圈起来，想记成看到ω就换\allowbreak{}\NovelInline{I00057}{n}\allowbreak{}。蔡须坤先指了第二条，让他看清前面仍留下的整段。

``这里有的换了，有的没换。''

``因此不能只凭符号长得像就乱替换。''玉小刚说，``要先看整个下标是什么结构，再按它已经选定的基本列取项。不同基本列约定还可能影响后面定义出的具体函数，本章统一用眼前这些。''

老师让他们取第三项。唐三先看图，牢大先写式子，再互相核对。

\NovelDisplay{M00114R01}{%
\omega[3]=3%
}

\NovelDisplay{M00115R01}{%
\bigl(\omega\times3\bigr)[3]=\omega\times2+3%
}

\NovelDisplay{M00116R01}{%
\omega^2[3]=\omega\times3%
}

\NovelDisplay{M00117R01}{%
(\omega^\omega)[3]=\omega^3%
}

小舞问能否接着再降一次。玉小刚说可以依定义逐步做，但要分别看当前下标；用进函数以后，输入也要依函数自己的规则处理。

``前一堂你们已经见过，输入在变，选哪一层也可以变。''玉小刚说，``现在把有顺序的路标放到函数下标处。要使它真能计算，还得另外交代规则，不能只写个高下标就算会了。''

他先约定重复调用的写法。

\NovelDisplay{M00118R01}{%
f^{\langle0\rangle}(x)=x%
}

\NovelDisplay{M00119R01}{%
f^{\langle k+1\rangle}(x)=f\bigl(f^{\langle k\rangle}(x)\bigr)%
}

``尖括号记调用次数，不能当乘方。后一次用前一次的结果，不能总拿最初输入重算。''

\NovelDisplay{M00120R01}{%
f_0(n)=n+1%
}

\NovelDisplay{M00121R01}{%
f_{\alpha+1}(n)=f_\alpha^{\langle n\rangle}(n)%
}

\NovelDisplay{M00122R01}{%
f_\lambda(n)=f_{\lambda[n]}(n)\qquad(\lambda\text{为极限序数})%
}

``这套叫FGH，快速增长层级。本章取正整数输入，后继层统一从\allowbreak{}\NovelInline{I00058}{n}\allowbreak{}出发重复\allowbreak{}\NovelInline{I00059}{n}\allowbreak{}次，极限层用刚才选定的基本列。别的材料若换了次数或基本列，具体值要另核。''

牢大先将第一层完整做过。每次只加一，做\allowbreak{}\NovelInline{I00060}{n}\allowbreak{}次，从\allowbreak{}\NovelInline{I00061}{n}\allowbreak{}走到两倍的\allowbreak{}\NovelInline{I00062}{n}\allowbreak{}。第二层便轮到把倍增这件事再重复。

\NovelDisplay{M00123R01}{%
f_1(n)=2n%
}

\NovelDisplay{M00124R01}{%
f_2(n)=n\times2^n%
}

\NovelDisplay{M00125R01}{%
f_1(3)=6%
}

\NovelDisplay{M00126R01}{%
3\xrightarrow{f_1}6\xrightarrow{f_1}12\xrightarrow{f_1}24%
}

``最后一行的意思是，这三次调用最后得到二十四。''玉小刚指着纸，``把沿途的输入变化留出来，你们便不会把三次都算成三变六。''

\NovelDisplay{M00127R01}{%
f_2(3)=24%
}

\NovelDisplay{M00128R01}{%
f_3(2)=f_2\bigl(f_2(2)\bigr)=f_2(8)=2048%
}

\NovelDisplay{M00129R01}{%
f_3(3)=f_2\bigl(f_2(f_2(3))\bigr)%
}

小舞想接着算三处的值，老师让她先留精确的嵌套式，确认总共调用几次。

``碑上有时只用代表式帮助人认大致层次。这里精确计算，必须依完整定义。你们上一堂的运算层级编号，也不能直接当成这套FGH的下标编号。''

接着，他让唐三指出第一处极限该怎样取项。

\NovelDisplay{M00130R01}{%
f_\omega(n)=f_n(n)%
}

\NovelDisplay{M00131R01}{%
f_\omega(3)=f_3(3)%
}

``这次就取第三层在三处的值。换成四，选第四层。与上一堂对角化的思路接上了，但每次仍在做一个有限输入上的计算。''

\NovelDisplay{M00132R01}{%
f_{\omega+1}(3)=f_\omega\bigl(f_\omega(f_\omega(3))\bigr)%
}

牢大点住最里面的三，说第一遍算过以后，下一遍的输入已经变了。玉小刚点头，让他把这句话念给小舞听。

``它不是简单跳去\allowbreak{}\NovelInline{I00063}{f_4(3)}\allowbreak{}。''唐三说，``后继下标要重新按重复规则走。''

\NovelDisplay{M00133R01}{%
f_{\omega^2}(3)=f_{\omega\times3}(3)%
}

\NovelDisplay{M00134R01}{%
f_{\omega\times3}(3)=f_{\omega\times2+3}(3)%
}

``这两步则都在极限处取基本列，输入仍是三。下标降低的办法，与函数调用后输入怎样变，也得分清。''

唐三把``下标''和``输入''各圈一处，终于明白为什么自己单看一串三，总以为是在做同一种替换。那几步表面相近，回去找的规则却不同。

黑板最后一角，玉小刚画起ω指数塔。

\NovelDisplay{M00135R01}{%
\omega,\quad\omega^\omega,\quad\omega^{\omega^\omega},\quad\ldots%
}

\NovelDisplay{M00136R01}{%
\varepsilon_0=\sup\{\omega,\omega^\omega,\omega^{\omega^\omega},\ldots\}%
}

\NovelDisplay{M00137R01}{%
\omega^{\varepsilon_0}=\varepsilon_0%
}

``这里每个有限高度的塔都仍低于\allowbreak{}\NovelInline{I00064}{\varepsilon_0}\allowbreak{}。它们的上确界是第一个满足这条不动点等式的序数。不是把塔高挑得特别大，便已经拿到极限本身。''

小舞想起原先那面添在队尾的小旗，再看这些层层变高的塔，终于不再问哪一座是最后一座。没有最后一项的序列，也可以有明确的上确界。

更远的\allowbreak{}\NovelInline{I00065}{\varepsilon_\omega}\allowbreak{}等位置以后再认。多头蛇、TREE、PTO与SGH各有自己的问题与规则，老师只提过名字，没有接着展开。

唐三试着换成多下标。玉小刚留下对应，提醒这只是沿用同一套规则的简写，不能凭多一个零就说力量长了一层。

\NovelDisplay{M00138R01}{%
f_{1,0}(n):=f_n(n)=f_\omega(n)%
}

\NovelDisplay{M00139R01}{%
f_{1,m}(n):=f_{\omega+m}(n)%
}

\NovelDisplay{M00140R01}{%
f_{2,0}(n):=f_{\omega\times2}(n)%
}

\NovelDisplay{M00141R01}{%
f_{1,0,0}(n):=f_{\omega^2}(n)%
}

唐三又将更一般的那条留作自己的设想，记为\allowbreak{}\NovelInline{I00066}{F}\allowbreak{}。

\NovelDisplay{M00142R01}{%
F(n):=f_{\omega^n}(n)=f_{\omega^\omega}(n)%
}

``最后一个等号用了已选定的基本列。知道定义以后，还得逐段核实，不能把认得的路标全当成已经练成。''

玉小刚还让他在页边留空：这里先记对应，遇到尚未讲到的极限，必须把它的基本列补清，才能继续实际运算。今天的几条例子，不是一句可以代替所有取项规则的口令。

五个人轮着把今天的例子讲回去。奥斯卡讲到两段长队，仍指错了一次乘号两边；牢大把成对的小段摆回桌上，让他自己改。蔡须坤没有像上次那样先凭感觉走远，而是跟着他们，将每一处``整段''与``段内位置''说清。

{\def\LTcaptype{none} % do not increment counter
\begin{longtable}[]{@{}ll@{}}
\toprule\noalign{}
例子 & 它帮助分清什么 \\
\midrule\noalign{}
\endhead
\bottomrule\noalign{}
\endlastfoot
任取一些凳子找最前一张 & 良序要求每个非空部分都有最前项 \\
自然数长队 & \allowbreak{}\NovelInline{I00067}{\omega}\allowbreak{}没有最后一个有限位置 \\
队首添一位、队尾添一位 & \allowbreak{}\NovelInline{I00068}{1+\omega}\allowbreak{}与\allowbreak{}\NovelInline{I00069}{\omega+1}\allowbreak{}不同 \\
一对一对排、整队接两次 & \allowbreak{}\NovelInline{I00070}{2\times\omega}\allowbreak{}与\allowbreak{}\NovelInline{I00071}{\omega\times2}\allowbreak{}不同 \\
按排号、排内位置取书 & \allowbreak{}\NovelInline{I00072}{\omega^2}\allowbreak{}的分组顺序 \\
逐层分组与基本列 & 极限没有紧前项，接近它要依具体约定 \\
在路标上定义函数 & 下标结构、调用次数和输入分别核对 \\
\end{longtable}
}

下课后，他们把借来的凳子一张张搬回去。奥斯卡确认三枚铜钱都在，揣进包里；小舞还想把最后那张挪到队首，被牢大看了一眼，便笑着放回原位。

经过门口的碑，唐三先看见\allowbreak{}\NovelInline{I00073}{\omega}\allowbreak{}，再往上认出\allowbreak{}\NovelInline{I00074}{\omega+1}\allowbreak{}和\allowbreak{}\NovelInline{I00075}{\omega^2}\allowbreak{}。那些刻痕终于不像一串故意写得陌生的名字。他能说出其中几处为什么排在那里，也知道自己暂时还不能把它们全部用出来。

蔡须坤走到他旁边，问以后是不是还要回来。

``当然。''唐三合上本子，``有些地方，还没走明白。''

小舞在前面喊吃饭，奥斯卡已经跟了过去。牢大等两人收好纸，才跨过门槛。走廊上的凳子都归了位，五个人往后的日子，也终于有了几处能共同指认的路标。

\chapter{第四章：炉火与青石山谷}\label{ux7b2cux56dbux7ae0ux7089ux706bux4e0eux9752ux77f3ux5c71ux8c37}

箭头与序数的两堂专门课结束后，兄妹俩仍要将课堂上的路标一步步接进实际练习。小舞想试试新近接稳的一段，放学以后，便拉唐三去了后山。

``来，我们切磋一下，白天打那几个菜鸟，根本不过瘾，正好试试用新学的箭头运算，提升大数能量的爆发威力。''

小舞说着，递给唐三一把匕首和一套厚重的护甲，``戴上，别被我打伤了。''

唐三看着手中的护甲，眼睛一亮：``这不是 Buchholz 宗门的招牌------Buchholz 突击防弹衣吗？你怎么会有这个？听说这个宗门的装备，能很好地承载大数能量，保护修士不被自身构造的巨数能量反噬。''

小舞摆了摆手：``你不用管这个，Buchholz 和 Madore 是现在最主流的两个宗门，之间的争斗也很多，我们现在还不用管。头盔、护腿、靴子也戴上，别大意，我可不会手下留情。''

两人穿戴齐整，唐三先凝聚起手的力量，向小舞逼近。他把接续铺得很大，落到第二处时却慢了一瞬；小舞一直没有硬接，绕开锤势，从他留出的空处掠了过去。

``看招！''小舞突然施展功力，抛出一个东西。``香肠？这是什么？''唐三正疑惑，那''香肠''就落在了他脚下，瞬间炸开------居然是燃烧弹！燃烧弹爆发的能量，也是基于大数构造的，威力远超普通的火焰。

唐三浑身起火，连忙灭火，心里满是不解：``我构造的大数级别比你高，怎么会被你的燃烧弹伤到？''

``你只顾着往后续，没看我在哪。''小舞说，``再说，你只有四点灵气，我有八点。招式能写多高，身体能撑多久，出手落得准不准，都得一起看。''

``服了服了，到此为止吧！''唐三无奈地说道。

小舞笑着把护甲递给唐三：``这护甲你收着吧，不用还给我了。它能帮你承载更大的大数能量，以后你构造更高级别的数，也能避免被能量反噬。''

唐三连忙推辞：``这怎么行，这么珍贵的装备，还是给你。''

小舞摆了摆手：``你也帮了我，解决了一些闹事玩意，就当是我送你的礼物。以后还要修炼，突破，构造出远超过葛立恒数那样的超级巨数呢。''

切磋之后，兄妹俩躺在后山的草坡上歇气。

``哥，''小舞忽然问，``你还记得娘吗？''

唐三想了想：``不记得了。就记得一个调子，有人哼过。词记不住。''

``爹从来不说娘。''

``嗯。''

``你问过吗？''

``问过一回。''唐三说，``他在炉子边上坐了一晚上，第二天照样打铁，一个字都没说。''

第二天上课，唐三仍想着草坡上小舞提起娘时的神情。

爹一个人守着炉子，兄妹俩读书用的纸和笔，却要一回回托人带来。唐三低头看自己的草稿，翻过用到最后一角的那张纸，心里慢慢有了一个打算。

散课后，小舞一脸无语地过来，递给他一块晒干的果糕：``唐三，玉老师讲的双箭头运算，我还有点没懂，你回头再给我讲讲好不好？''

唐三：你这是在威胁我还是求教呢？行了行了，别这样行不。

小舞：没事，我这人性格不好，你们都知道。

唐三：行吧，我这倒是还有点书籍，你自己去看便是。

等小舞回去休息，唐三独自走到院子里，望着远处的山林，开始认真思索赚钱的法子。他自幼跟着父亲长大，没学过什么经商的门道，但从会走路起就在铁匠炉边上拉风箱，倒是练就了一身打铁的好手艺。唐家祖上三代都是铁匠，父亲话少，可一手打铁的功夫却是顶尖的，唐三耳濡目染，不仅能熟练打造普通的农具、厨具，还能凭着过人的细心和精准的力道，打造出一些造型精巧、耐用耐磨的小物件，甚至能打磨出锋利的短剑坯子。

``打铁，或许是眼下最稳妥的法子。''唐三喃喃自语，眼底渐渐有了光亮。他知道，镇上的铁匠铺不多，大多手艺粗糙，打造的物件要么不耐用，要么造型笨拙。而他不一样，他不仅手艺扎实，还能借着研习大数时培养的逻辑思维，优化打铁的工序------比如计算火候的把控时间、锤炼的次数，甚至能根据不同材质的特性，调整击打力度，让打造出的物件更趁手、更耐用。

他细细盘算起来：首先得找一处简陋的铁匠铺，不用太大，能放下一个熔炉、一副铁砧就好，镇上废弃的老铺子应该有，花点小钱就能租下来；然后是原材料，废铁可以去镇上的废品站收，价格便宜，回来打磨提纯后就能使用；打造的物件也要有讲究，先从最实用的农具、厨具做起，比如锄头、镰刀、铁锅、菜刀，这些是镇上人刚需的，容易卖出去，等积累了一些本钱，再打造一些精巧的小饰品、短剑坯子，卖给镇上的年轻人或者行路人，利润会更高。

唐三又想起父亲，打算等手头宽裕些，替家里的炉子做一把趁手的铁钩。玉老师平日照看他们，他也想亲手打件小物件，找个合适的时候送去。

但他也清楚，打铁是个苦差事，不分昼夜，耗费体力，还要忍受熔炉的高温和铁屑的划伤。可一想到能靠自己的双手撑起日子，能让小舞过得舒心些，能减轻家里的负担，唐三就浑身充满了力气。他暗下决心，每天上完课、和小舞研习完大数运算后，就去铁匠铺打铁，哪怕熬夜到深夜，也一定要坚持下来。

当晚，唐三趁着夜色，悄悄去镇上打探了废弃铁匠铺的消息，又去废品站问了废铁的价格，一一记在心里。回到宿舍，奥斯卡的呼噜打得震天响。他躺在床上，望着黑乎乎的屋顶，轻轻叹了口气，在心里默念：``爹，我会好好努力，靠自己的手艺赚钱，再也不让你和小舞受苦了。''

第二天一早，唐三便租下了那间废弃的铁匠铺，用自己攒下的为数不多的零花钱，买了些简单的工具和废铁。

他先将铁匠铺角落的灰尘清扫干净，把铁砧擦得锃亮，又将熔炉清理妥当，添上晒干的木炭，拿起火折子轻轻一引，橘红色的火苗便顺着木炭蔓延开来，渐渐窜起半尺高，将小小的铁匠铺烘得暖意融融，也熏得他脸颊泛起淡红。

待熔炉温度升至灼热，他便将捡来的废铁切成规整的小块，用长柄铁钳夹起，缓缓送入炉心，看着黝黑的铁块在烈火中渐渐泛红、变软，直至被烧得通体赤红，泛着亮眼的光泽------他凭着多年的经验，不用温度计，仅凭铁块的色泽，便精准判断出火候刚好，这便是他常年耳濡目染练就的本事。

紧接着，他左手紧握长柄铁钳，指节因用力而泛白，虎口死死扣住钳柄，稳稳夹住烧红的铁块，手腕微微翻转，迅速将铁块从熔炉中抽出，手臂顺势前伸，精准地将铁块按在铁砧中央，没有丝毫偏移。

右手猛地抄起身旁沉甸甸的铁锤，手臂高高扬起，肘部微屈，肌肉线条绷得紧实如拉满的弓弦，指尖牢牢扣住锤柄，借着身体下沉的力道，让铁锤带着呼啸的风声重重落下，``哐当------''一声脆响，震得铁砧微微发麻，也震得他肩头轻轻一颤。赤红的铁屑随着撞击飞溅而出，他下意识地偏头避开，左手手腕轻转，微微调整铁块的角度，确保每一次捶打都落在需要塑形的位置。

几轮下来，铁坯渐渐展平。唐三听着落锤的声音，调整手腕，哪里还厚，便多留几锤；薄处先避开，等下一次复烧再修。

打造一把菜刀的坯子，他早已在心中盘算好锤炼的次数与力度，就像研习 Knuth 箭头运算般严谨。他微微弓身，重心下沉，左手铁钳始终稳稳固定着铁块，右手铁锤起落间节奏均匀，每捶打十下，便松开钳柄，手腕翻转，将铁块重新送入熔炉复烧。待铁块再次烧得通体赤红，他又迅速抽出，铁钳精准归位，铁锤再度扬起，动作一气呵成，没有半点停顿。

汗水顺着他的额角滑落，浸湿了额前的碎发，顺着脸颊滴落在滚烫的铁砧上，瞬间蒸发成一缕白烟，他却浑然不觉，眼神始终专注地落在铁块上，眉头微蹙，左手手指轻轻摩挲钳柄，根据铁块的塑形情况，细微调整夹持的力度与角度，右手挥锤的速度与力道也随之微调，时而轻敲细琢，时而重锤猛攻，原本不规则的废铁块，在他精准的动作下，渐渐被捶打得平整光滑，边缘也愈发规整。

待坯子初具雏形，他将铁钳固定好，腾出左手扶住錾子，右手换上小锤，沿需要修整的位置轻敲。铿锵声渐渐低下来，炉火里还留着一点橘红。

他俯下身，沿刀坯边缘一点点去除毛刺。遇到不平的地方，便将坯子重新夹稳，再用錾子修整。炉火映着他的侧脸，锤声已经比先前轻了许多。

夕阳西下时，第一把菜刀的坯子终于打造完成，他将坯子放入冷水之中，``滋啦------''一声，白烟瞬间升腾，水汽弥漫，待水汽散去，一把通体泛着冷光、边缘锋利的菜刀坯子便呈现在眼前，手感沉甸甸的，纹路规整，没有一丝瑕疵。唐三拿起菜刀坯子，轻轻摩挲着光滑的表面，眼底露出一丝浅浅的笑意，这是他靠自己的手艺打造出的第一件物件，也是他撑起日子的第一步。

闲暇时，小舞也来帮忙，递工具、拉风箱，偶尔将课上没接顺的练习拿来。唐三等手里的活告一段落，才放下锤，与她并坐着看。两人说清一处，小舞便在纸角做个小记号，等下一回有空，再各自重算。

小舞是他妹妹，吵归吵，赚了钱总得分她一份。唐三想了想，就让她来铺子里帮忙。

有时也有人请他打造铁制发簪。他换上小锤与细錾，先在冷下来的样坯上比好纹路，再一点点修整。小舞原先只看打刀，第一次见他做这样小的东西，便也搬来凳子，坐在旁边帮着挑选适合的铁条。

小舞蹲在一旁，替他看着需要更换的工具。等铁坯放稳，唐三才接过她递来的湿布，擦去脸上的铁灰，又拿起下一件活。

傍晚，小舞将做好的发簪排在窗边。唐三给最后一件收好口，往外一看，她正凑着光检查纹路；等他过去，便将自己觉得不平的那处指给他看。

铁匠铺的生意日渐红火，可新的难题终究还是横亘在了唐三面前。

他与小舞定下三七分账的规矩后，兄妹俩从小到大的那些小摩擦便时常冒出来，经常因营收核算的事情吵起来。唐三嘴上不说，心里总觉得这丫头是来白拿钱的，一来二去，吵了不止一回。

由于要钻研箭头与超运算的奥义，兄妹俩白天根本无暇打理铁匠铺。铺子里的生意大多集中在日间，村民们赶集采购农具、厨具，旅人挑选随身铁器，每每都因无人值守白白错失客源，不少熟客几次扑空后，便转而投奔了镇上其余的铁匠铺子。

看着日渐减少的客源与账本上停滞不前的营收，唐三夜里独坐铁匠铺，望着炉火沉吟许久，终于向小舞提出了招工的想法。

``铺子缺人照看，我们白日都在学院，根本顾不过来，招一个靠谱的学徒值守日间，打理杂活、接待客人，生意才能稳住。''唐三语气笃定，这是他权衡许久的最优解法。

可小舞当即摇了摇头，满脸不赞同：``没必要的，多招一个人就要多开一份工钱，本来利润就不算极高，刨去人工、物料成本，根本赚不到多少，纯属得不偿失。我们放学后辛苦一点，熬夜赶工、抽空值守就够了。''

``熬夜终究不是长久之计，长期熬身体扛不住，而且日间客源流失的损失，远比人工工钱多。''唐三眉头微蹙。

``那也比白白多花一份钱要强！你就是太想当然了！''小舞性子执拗，寸步不让。

两人各持己见，互不相让，过往积攒的小矛盾彻底爆发，争执声在寂静的铁匠铺里回荡，谁也没能说服谁。这场争吵最终没有任何结果，招工的事情就此搁置，成了两人之间一道无声的隔阂。

风波落幕，喧嚣散去，唐三独自站在炉火旁，看着跳动的火苗，心底的思绪愈发繁杂。他转头看向一旁沉默不语的小舞，怒火中烧------

``你说话啊，你做了多少活！''

``铁匠铺经营成这样，你是猪吗！''

吵了很多回合之后，小舞离开了。

兄妹俩从小吵到大，吵完了第二天照样一块儿吃饭。可这一回，小舞好几天没来铺子。

看着妹妹走远的背影，唐三骂了一句活该，活是我做的，还想分走我的钱，凭什么？？

骂完了，他坐回炉子边上，把账本翻开，在最后一页另起一栏，把小舞那三成一文不少地记了上去。

\begin{center}\rule{0.5\linewidth}{0.5pt}\end{center}

第三天傍晚，铁匠铺的门被人推开了。

进来的是奥斯卡，后头跟着牢大和蔡须坤。

``听说你们兄妹俩吵翻了？''奥斯卡一进门就问，手里还拎着一串香肠，``铺子白天没人看？''

唐三没吭声。

``我白天第三节课没课。''奥斯卡说，``我来看铺子。我在门口烤香肠，香肠一烤，人就来了，人来了顺便就买把菜刀。''

``那第一节、第二节呢？''

``我没课。''牢大说。他走到风箱跟前，看了看，握住把手拉了一下。风箱呼地一声，炉膛里的火往上一窜，稳稳地烧着，不大不小。``我拉风箱。''

唐三看着他拉了几下，每一下的力道都一模一样，炉火的颜色一点都不跳。

``你拉得比我爹还稳。''唐三说。

蔡须坤没有去碰什么东西。她站在铁砧旁边，看着砧子上那把没打完的菜刀坯子，看了很久。

``为什么要自己打刀？''她问。

唐三愣了一下：``不打刀，拿什么切菜？''

``切菜是为了吃。''蔡须坤说，``吃是为了长大。长大了，再打刀。''她像是在把一件事情一节一节捋清楚，``一代一代，都这样？''

``都这样。''

蔡须坤点了点头，又去看那把刀坯子。

那天夜里，小舞回来了。她站在铺子门口，看见里头牢大在拉风箱，奥斯卡在门口烤香肠，蔡须坤在灯底下帮唐三数铜板，一个一个数，数得很认真。

她在门口站了一会儿，走进去，从奥斯卡手里拿过一根香肠，在唐三旁边坐下了。

``账本给我。''她说，``你算得不对。''

唐三把账本推过去，没说话。

小舞翻到最后一页，看见了单记的那一栏，手停了一下，也没说话。

从那以后，铁匠铺白天有人看了。

过了一阵子，唐三开始能把两堂专门课上留下的一些设想真正接进自己的练习，偏偏在一次对角取项时，又反复卡住。

他知道定义写在哪一页，照着读也没有漏字，可合上本子，便分不清自己刚才究竟把哪一个位置跟输入一起改了。

他独自试了几晚，纸上看起来总能接上，实际运转却在同一处停住。小舞陪他核过，也只看出他每一遍停得都差不多。

万般无奈之下，唐三听闻学院高年级的王圣学长，学识渊博、眼界远超同辈，尤其擅长高阶递归与对角化构造，是学院中公认的天才。于是他整理好疑惑，主动登门请教。

王圣出身富裕世家，性情温和、惜才爱学，没有半点高傲架子。见唐三谦逊好学、求知若渴，且基础功底扎实、思维缜密，对大数有着极高的天赋与悟性，便耐心为他拆解对角化的核心逻辑。

王圣没有重新讲一遍整套定义，只请唐三将一次最短的练习独立说出。说到第三步，唐三伸手想翻纸，王圣便让他停在这里，把前面真正用过的那条找回来。

两人从午后坐到傍晚，王圣问得很细，却始终没有把整道结果替他写完。

同一个傍晚，唐三还在王圣那里核笔记，牢大、奥斯卡与蔡须坤则去了学院后的青石山谷。那里要翻过两道山梁，是他们平日练功的地方。唐三和小舞跟着去过几次，赶上铺子有活，便不能总往山里跑。

那天傍晚，牢大最后一个到。

他替唐三拉过半日风箱，袖口还有铁灰，到了谷里，却仍先将平日那段基础练了一遍。奥斯卡坐在石头上等，手里的香肠吃完了两根，忍不住问他，为什么已经会的东西，还要每天这样过。

``怕手生。''牢大说。

奥斯卡捡起枯枝，把自己新整理的练习写到石板上。他沿着入门课里的阿克曼记法，找到了几处比旧笔记简洁的表达，正想拿来给同伴试。

``再往上，也许能省不少力。''他说，``总沿旧箭头一步步搭，太慢了。''

牢大没有马上接话。

他先把奥斯卡写的那几行看过，蹲下来，在旁边留出自己习惯的顺序。两人从前也常交换笔记，这一次却都没有顺着对方的话往下走。

``慢，不等于没用。''牢大说。

``我没说没用。我说换个记法试试。''

``你刚才说旧的太慢。''

奥斯卡握着枯枝，在原处划了一下：``这里确实能写得短。你看一眼再说。''

牢大看了。他知道纸上的写法并没有凭空跳过底层步骤，可心里仍不舒服。那些步骤，他与蔡须坤反复核过许多回，有过错处，也有过终于不再出错的清晨。现在被一句``太慢''带过去，好像那些日子根本没必要。

``我练了三年的基础。''他终于说，``一天一天练到稳，后面才能接。你不能换了一个名字，就把前面的都算成多余。''

奥斯卡的笑意收了。

``那我想试新的，也不能先算成不稳。''

火堆里的柴爆了一声。蔡须坤抬起头，想说什么，牢大已经指向石板上一处没写明的接续。

``你这里呢？''

``我还没试完。''

``没试完，怎么知道一定更好？''

奥斯卡站起来，把枯枝扔回石头边。

``你就是怕出错。只要不试新的，当然永远稳。''

这句话说出口，他便知道重了。

牢大抬头看他，脸上没什么表情，撑着石板的手却一点点收紧。蔡须坤先前也见过那只手：在铺子里，能将风箱拉得丝毫不跳；在深夜，又会为了一个偏差，将整段从头重做。

她伸手按住了牢大要拿起的笔。

``先别写了。''

牢大没有看她。奥斯卡也转过身，过了一会儿，才又坐下。

谷里一下安静下来，连火堆都显得小了。

蔡须坤把两人的纸摊开，想从同一个输入处找共同的接续。她隐约感觉到，写法虽然不同，几处气息相接的位置却很近。她闭上眼，顺着那一点相近的感觉去找，远处几只夜鸟竟落到了更近的石头上，没有叫。

她睁开眼时，牢大正望着那些鸟。

``怎么了？''她问。

牢大摇头，把目光移回纸上。

蔡须坤将入门课留下的对照页取来，放在两人的笔记中间。她没有凭那点感觉宣布谁算对，只请两人各自看一遍课上已经说清的对应，再核眼前这一次实际少了哪一项。

奥斯卡最先低下头。

他沿自己的新写法读过，发现有一处起点被他省掉了。牢大用旧写法接得完整，到了下一段，却也因为不肯换角度，多绕了好几次。

两个人都看见了，没有谁先将对方的错处圈出来。

``这里能对上。''牢大说。

奥斯卡点头，把刚才扔下的枯枝捡回来，补好那一处起点。

``我说它短，是想让你看这段。''他声音低了些，``没有想说你练的东西白费了。''

牢大望着火，好一会儿才开口：``我也没真把你后面写的看完。''

两人重新拿起纸，这一次没有争谁更接近天地本源。奥斯卡说明自己想省掉哪一道重复，牢大指出哪些条件必须留着；蔡须坤坐在两人中间，把他们用不同说法指向同一处的位置记下。

那晚的试练仍出了错。

奥斯卡把内层推得太快，牢大接回时稍慢了一瞬，石板边上的火便猛地窜高。蔡须坤伸手收窄两人外散的气息，等牢大稳住，才让奥斯卡将后面那一段慢慢撤回。

三个人退开，看着一角被熏黑的纸。

奥斯卡先把剩下的香肠拿出来，递给牢大。牢大接过，没有说话，掰下一半给蔡须坤。

``明天再试。''奥斯卡说。

``先把今天没接稳的地方记下来。''

``行。''

蔡须坤将纸转到背面，把两句话都记在了上面。夜鸟从石头上飞起，往林子里去了，她看了片刻，才低头将最后一行写完。

回学院时，三人走得很慢。谷里的火已经熄灭，牢大仍提着灯，等奥斯卡将最初那个漏掉的起点说清，才跨过山道上的最后一道沟。

第二次去王圣那里时，唐三已经能不翻原纸，把那一段接清。他把上一次误认的位置留在页边，自己讲给王圣听，确认终于不是只认得黑板上的名字。

几番交流下来，王圣愈发认可唐三的天赋与心性。他见过无数研习大数的学子，却极少有人能如唐三一般，基础扎实、勤于钻研、悟性卓绝，且能沉下心深耕细节、直面疑难，未来必定能在大数之学上大有作为。

知晓唐三家境清贫，一边打铁谋生、一边潜心求学，时常因生计奔波分心研习，王圣心生惜才之意。他主动提出资助唐三，为他提供研习经费，免去他的生计之忧。

这份帮扶免去了他一部分书费与纸笔的负担，铁匠铺仍要照管，却不必每次为赶齐学费硬撑到后半夜。唐三收下资助时，认真记了账，王圣没有催他将来怎样报答，只让他有新进展便拿来看看。

那一年，唐三满了八岁。

他和小舞将自己能独立接稳的多下标练习重新整理，把入门课里那条F函数也留作继续推演的方向。有些段落已经能实际用出，有些仍只知道定义，两人便将它们分开放，免得为了赶进度，把尚未练成的也记成自己的本事。

可就在两人准备向更高境界冲刺的时候，年关到了。

放假前几天，唐三将铺子里的活一件件排好，准备把未收的尾款结清，再同小舞回村。奥斯卡来帮忙，才看见桌底还压着几件未磨好的农具。

``你今晚不睡了？''

``这几把明天有人来取。''

牢大拿起一把看过，没去碰旁边的铁锤，先问哪一处能由自己来磨。唐三指给他看，他便搬条矮凳坐下。奥斯卡把门边的香肠摊收进来，替他们腾出地方；蔡须坤和小舞并坐在灯下，一人念旧账，一人将付清的那笔勾掉。

``这张上面说，开春才来取。''蔡须坤举起一张纸。

唐三接过，果然是自己看错了日期。他把那件活搁回去，又将另一张拿来核对。才做了一会儿，桌上便少了一半急等着完成的东西。

奥斯卡用胳膊碰了碰牢大：``有时候是他自己不肯停。''

``你也一样。''牢大说。

奥斯卡刚要反驳，就看见蔡须坤抬头瞥过来，只好把话咽回去。山谷里那晚试到纸被熏黑，他起先也说马上能成，结果回学院时险些错过关门。

小舞看他们三人交换眼神，想追问谷里出了什么事。奥斯卡忙让她先把钱数完，蔡须坤却认认真真答了：``我们接错了。''

``你们也会？''

``会。''牢大将磨石浸进水里，``还没改好。''

唐三抬起头。他最近每次去找王圣，回来都要将笔记重新整理一次，却总怕被别人看出又卡住了。现在听牢大说得这样平常，他反倒问起了是哪一处。

``等回来再给你看。''奥斯卡说，``你这几天先别添活了。''

最后一件农具磨完，已经过了饭点。小舞将留着的干粮热了一遍，五个人就着炉边的余温吃。唐三给他们倒水，看见牢大袖子上仍沾着灰，便说等年后多赚些钱，给几人分一份。

``我也不是专为这个来的。''牢大说。

奥斯卡却伸手敲敲账本：``他不要，我要。先记着，等开春多买几张好纸。你这纸一翻身，就不知道原来的字在哪边了。''

唐三笑了一下，在本子最后认真添上。蔡须坤看见了，伸手指指自己，又指指桌上剩下的那几张纸。

``也给你留。''唐三说。

她这才点头，把账本合上。

关门回学院时，唐三又核了一遍留下的单子。第二天，兄妹俩就要回村了，牢大收起钥匙，答应替他们看好门窗。小舞将包里的纸按顺序叠好，最上面放着那页新添的序数图。

\chapter{第五章：父亲的信与同伴的手}\label{ux7b2cux4e94ux7ae0ux7236ux4eb2ux7684ux4fe1ux4e0eux540cux4f34ux7684ux624b}

兄妹俩跟往年一样回了家。杰克爷爷在村口的老槐树底下等着，一手牵一个，把他们领回了家。爹还是那个样子，话少，白天打铁，晚上坐在炉子边上。年过完了，开春了，眼看着再过几天就要回学院。

走前那天晚上，杰克爷爷来家里坐。爹破天荒没去炉子边，也在桌边坐着，手里捧着一碗茶。

小舞缠着杰克爷爷讲故事。

``还讲那个给大数起名字的故事。''小舞说。

杰克爷爷笑着说起那个替巨数起名的孩子。唐三在课上见过古戈尔与古戈尔普勒克斯，小舞却仍爱听取名字那一段，尤其爱听``写到手累''为止的提议。

``后来先生没答应。''杰克爷爷说，``你写得累，别人还不累；总得先说定究竟怎样写，不能谁胳膊有劲，就由谁来定这个数。''

屋里安静了一下。

爹一直没说话，这时候把茶碗放下了。

``写到累的，不算数。''他说。

一家人都看着他。他没再说什么，起身回了炉子边上。

那是唐三记得的，爹走之前说的最后一句话。

那天夜里，唐三睡得很沉。

第二天早上醒来，炉子是冷的。

炉子边上那面墙，原本挂着一排锤子，大大小小，都是唐家传了三代的家伙什。现在只剩一排铁钉。唐三在屋里屋外找了一圈，没找着人，最后在自己枕头底下，摸到了一封信。

信封是黄的。他拆开，里头只有一张纸，字不多：

\begin{quote}
三儿：

爹要出一趟远门，走得急，来不及跟你说。

不用找我。

照看好妹妹，听杰克爷爷的话。你八岁了，是大孩子了。书接着念，杰克爷爷会送你回学院。

等你长大了，要是还想知道爹去了哪儿，
\end{quote}

后面还有一行，被划掉了，划得很用力，一个字也认不出来。末尾没有落款，只盖着一个朱红的圆印，印里一道弯弯的纹，绕了大半圈，到右边没有合上。

唐三拿着信站在院子里，站了很久。

杰克爷爷是被小舞的哭声引过来的。他看了那封信，把两个孩子领回自己家，烧了一锅热粥。

唐三捧着粥碗，声音有些沙哑：``杰克爷爷，我爹走之前，跟您说什么了没有？''

杰克爷爷低头搅着锅里的粥，搅了很久。

``没说什么。''他说，``小三，别怕。你爹是出远门，办完事就回来。有爷爷在。''

唐三没再问。他把信折好，顺着折痕一道一道折回去，贴身揣在胸口那一侧。

学院那边，牢大还不知道唐家出了事。

兄妹俩离开后，铁匠铺暂时不生炉。牢大隔两日过去看看门窗，奥斯卡便顺路陪他，嘴上总说怕他一个人把剩下的活全做了，自己却把忘在门口的半袋木炭也搬了进去。

蔡须坤最后到。她带来几张纸，都是三人上回没接完的练习，角上有一小块被火熏黑的痕迹。奥斯卡往柜台上一铺，才发现旁边还留着唐三未带走的废稿。

``他的别弄混。''牢大说。

蔡须坤将那叠纸拢好，拿铁镇压住，才坐到他们中间。门外很冷，铺里也没有炉火，三人商量了一会儿，最后还是回学院去练。

冬天的青石路上结着薄霜。奥斯卡走到一处背阴地，脚底猛地一滑，纸飞出去半张。牢大伸手拉他，自己也被带得踉跄，最后两人一起扶住了旁边的树。

蔡须坤赶过去，将散掉的纸一张张捡回。她看着奥斯卡发白的指节，问要不要歇一会儿。

``没事。''奥斯卡直起腰，``刚才不算练功。''

``所以才没看脚下。''牢大说。

奥斯卡张了张嘴，这次没顶回去。他揉着被扯疼的胳膊，跟在蔡须坤后面，等下一个转弯时，先伸手摸了摸石上的霜。

那天下午，三人没有急着去接新的一段。牢大先把自己常停的地方让给奥斯卡试，奥斯卡又从头听蔡须坤说了一遍。她仍有几句讲得断断续续，另两人便等着，没有抢先替她补上。

练完以后，奥斯卡忽然说，唐三回来时，也许可以请他从另一头看一眼。

``他不一定比我们快。''牢大说。

``又没叫他跟我们比。''

牢大将纸折起来，想了想，点头。蔡须坤在纸角画了一个很小的记号，将他们停下的位置留下，等着开学后再给唐三看。

她并不知道，唐三此刻也在反复折着一张纸，只是那张纸的落笔人，已经不在他身边。

半个月以后，杰克爷爷又背着小舞、牵着唐三，走了一整天的山路，送兄妹俩回学院。

这一回路上，唐三一句话都不说。杰克爷爷也不劝，走了一段，自己讲开了。

路上，杰克爷爷给他们讲庙里搬金片的人。

小舞听到要一代接一代挪，问，那为什么还挪。

``前一个人做过的，后一个人就接得上。''杰克爷爷说。

唐三低着头，没有接话。走到溪边歇脚，杰克爷爷掰了半块馍给他，见他慢慢吃起来，又讲起大力士与多头蛇。

``砍了还会长，那怎么有个完？''唐三终于问。

``故事里说有它的规矩，照着才能走下去。''杰克爷爷看着他，``我只会讲，讲不清的，你去问先生。三儿，别什么都一个人闷着。''

唐三捧着那半块馍，过了好一会儿才点头。

杰克爷爷站起来，把小舞重新背上。唐三跟在旁边，手伸进衣襟，摸了摸那封折得整齐的信。

到了学院门口，杰克爷爷蹲下来，给唐三把衣领理好。

``三儿，好好念书。''他说，``你爹要是回来了，我让他来看你。''

回学院的路上，小舞一直攥着唐三的袖子。进了学院的门，她忽然开口：``哥，以后我罩着你。''

唐三看了她一眼：``就你？''

``我认真的！''小舞急道，``我构造大数的技巧比你厉害，爹不在，我罩着你。''

唐三看着她红红的眼眶，没再跟她斗嘴。

``行。''他说，``你罩着我，我照看你。爹信里说了，照看好妹妹。''

小舞把脸扭到一边去了。

那天晚上，宿舍里谁都没多问。

奥斯卡把一整包香肠塞进唐三怀里，说是家里新寄来的，吃不完。牢大什么也没说，第二天一早就去了铁匠铺，把炉子生好，一直拉风箱拉到唐三下课过去。

蔡须坤是最后一个来的。她在唐三旁边坐下，看了他很久。

``爹是什么样的？''她问。

唐三想了想：``话少。打铁。''

``会走吗？''

``\ldots\ldots 会。''

蔡须坤点了点头，像是又记下了一件事。

``我没有爹。''她说。

``那你从哪儿来的？''奥斯卡在旁边插嘴。

``林子里。''

奥斯卡噗地笑出声来：``你是猴子变的啊？''

蔡须坤没笑。她看着奥斯卡，看了一会儿，也没解释。

牢大在风箱边上，手底下顿了一下，又接着拉了下去。

当晚铺子里还有一把镰刀没有磨完。唐三将信塞回衣襟，拿起来试了试，才动两下，手便停了。他望着墙上挂工具的铁钉，一时没有分清自己究竟站在镇上，还是已经回到那间空下来的家里。

炉口呼地亮了一下。牢大将风箱的力道缓缓减下去，没有问他为何发呆，只让他把刀坯放稳。

``今天不做了？''唐三说。

``客人明天下午来。''小舞从后面翻出单子，``来得及。''

奥斯卡收好门外的木凳，回来正要说什么，小舞先朝他摆了摆手。他看了唐三一眼，便将余下的香肠重新包好，放在唐三回宿舍时能顺手拿到的地方。

蔡须坤帮着擦桌子。擦到那叠废稿时，她认出自己冬天收过的那几张，便问唐三，等哪天有空，要不要去青石山谷坐一会儿。

``我们还留着一处没接完。''

``你也没接完？''

``嗯。''她把纸放回去，``可以等你。''

唐三没有马上答应。等牢大将最后一点炉火压好，他才将那叠纸抽出，收进自己的包里。

几个人一同走回学院。小舞先前一直攥着他的袖子，这一回松开去赶奥斯卡，抢他怀里最后那根香肠。牢大落后一步，等唐三锁好铺门，才提着灯往前走。

转过街角时，唐三回头看了一眼。铺子里已经黑了，门却锁得好好的，明天下午还有客人要来取那把镰刀。他摸了摸包里的纸，脚下终于跟快了一点。

往后的日子，唐三和小舞照旧一有空就去后山。

他们在入门课上认过的那些路标，如今已经有一部分能实际走过。可纸上推得连贯，出手时仍可能迟半拍。唐三有时忙着续下一段，忘了看小舞的落点；小舞也会图快，先把力发出去，后面再想怎样收。

一回两回还看不出，切磋得多了，小舞送唐三的那件Buchholz护甲便补丁摞着补丁。最先磨坏的是肩带，后来侧面也裂了一道口，远看已经不大像原来的样子。

小舞捏着新换的带子看了半晌，给它起名叫``吊带''。唐三没笑，翻出账本，算明白再照这样打下去，铺子下月赚的钱又得花在修甲上。

王圣资助的纸笔和书费不能总拿来填这个窟窿。兄妹俩便定下规矩，日常练习收小些，真要试更强的接续，就先约好落点与收手的办法。规矩写在包里，真正打起来时，却仍有忘记的时候。

兄妹俩忙着修甲的那几日，牢大三人的练习也停得不大顺利。

他们在青石山谷里选了块平地，先把火堆移远，再用几根枯枝标出站位。牢大用力向来稳，照自己的步子走完并不费事，可奥斯卡从旁边一接，他便总想等对方全做完，再将自己的下一步送出去。

这一下等得太久，奥斯卡早先留下的位置也跟着变了。第一回还能勉强补救，第二回再试，三人便一齐停了手。

``你不用每次都等我点头。''奥斯卡说。

``不等你，怎么知道你准备好了？''

``我刚才喊了。''

``你喊的是纸又飞了。''

蔡须坤将被风卷走的那页纸放回石上，认真地点头：``我也听见这个。''

奥斯卡张嘴笑了，刚才那点烦躁也散了。他自己走到牢大原来站的地方，将两人方才接过的动作慢慢试了一遍，才发觉背对着看，确实分不清哪一声是要开始，哪一声是让大家先停。

三人重新约了几句话，先把空手换位走熟。轮到蔡须坤时，她却在两人开口以前，已经向前挪了半步。

``还没到。''牢大提醒。

``我感觉你下一步要在这里。''

``我知道。''他看着她的脚，``可是我还没走。''

蔡须坤低下头，退回枯枝后面。她方才感觉到的接续没有错，甚至知道前面还能怎样变，可牢大肩上的力尚未收住，她若这时催他跟上，他便不得不临时换掉已经起好的那一下。

奥斯卡没有接着开玩笑，蹲下来把被踩歪的枯枝放正，又用鞋底抹去了一条多余的线。

``照我们说好的走。''他说，``哪怕你先知道，也等一等。''

蔡须坤将手收回身侧。下一回，她一直望着牢大的肩，等听见他说可以，才将脚落到前面。牢大得以按原来的方向收回手，奥斯卡也终于没有在旁边临时补上一截。

只这几步，他们便走了十几回。后来奥斯卡肚子响起来，小舞上次留下的干饼早就分完，他从包里摸出香肠，才发现只剩两根。

``回去吧。''牢大说。

``还有一次，方才已经很近了。''

牢大没有和他争，将旁边的纸递过去。奥斯卡低头看，才发现自己后两遍连站位都记反了。他揉了揉眼，终于把笔放下。

蔡须坤先将火灭干净。等她回来时，牢大已经背起包，奥斯卡正把两根香肠各掰成几段，想让三个人都分得一样长。她站在旁边等了一会儿，接过最短的那段，几口便吃完了。

下山时，奥斯卡还惦记着最后那一次，牢大便将明早的时间告诉他。说到一半，他自己也停了停，把原本要去铺子的一段空了出来。

``先吃饭。''蔡须坤说。

另两人回头看她，她将包往肩上挪好，又重复一遍。奥斯卡终于笑起来，答应明早先在饭堂等，不会再拿两根香肠撑着练到天黑。

那天晚上，后山传出一阵金属撞击声。

巡查的女老师赶来时，小舞正从树干旁翻回来，唐三背后的土坡缺了一角，护甲上全是灰。

``停下！''

两人听见，连忙各退一步。最后一股余力撞在石边，碎土滚了下来。

``大晚上不休息，在这里做什么？''

``切磋。''小舞答。

``这里是学院后山，底下就是宿舍，你们想把学校拆了？''

唐三把武器收好，试着解释他们已经约过范围。老师看了一眼那处塌掉的土角，显然不信。

``跟我到教务处去。''

小舞看了看唐三，唐三也正看她。两人都怕再受一通训斥，下意识便说了不去。

``老师，我们真是在练魂技。''唐三说。

``还没毕业，就拿这个词来糊弄我？''

小舞本来已经收住，听到这里又不服气。她说了句``您看''，便弹身而起，想像往日与唐三试招时那样，让老师看清这一下到底怎样接。

老师没有料到她真会动，下意识抬手一挡。小舞已经借着那点阻力翻过去，腰腿一送，老师脚下失衡，被甩出了几步，重重坐在地上。

唐三脸上的笑还没起来，便看见老师捂着肩膀，顿时知道糟了。

两人赶紧过去扶人。老师缓过那口气，脸色很难看，再没问这算不算魂技。她也不敢继续计较，只让他们立即收手，自己扶着肩膀离开了。

第二天，玉小刚把兄妹俩留在训练场。

``昨天证明了什么？''

小舞不说话。唐三低头看自己护甲上的缺口，也没有答。

玉小刚指向场边立好的木桩：``先打这一个。''

唐三抬锤，木桩被打断，断木飞出了线外。

``再来。''

第二根也断了。这一次唐三收得早了些，锤势却在中途一滞，肩膀被回震得发麻。

小舞站在旁边，起先觉得这没什么可练，等到自己也试，才发现要把动作停在指定位置，比直接把对手甩出去难得多。

``会构造，是一件事；能在别人站在身边时把它收稳，又是一件事。''玉小刚说，``你们昨天要练的，正是这一点。''

几张入门笔记放在场边，谁接续不明白，就自己取回去核。场上的要求始终具体：先看清同伴的位置，按约定出手，到了边界便收。

那一上午，两人打坏的木桩比前几日都多，却没有再把断木送到线外。

下午，牢大来了。

他听过昨天的事，没有笑小舞，只将一块挡板摆在训练场另一侧，让他们从这里看边界。奥斯卡把水与香肠搁在阴凉处，蔡须坤则站到牢大身旁，等他将第一段气息收稳。

唐三见过两人配合，却没有这样近地看过。

牢大先起势，能量在肘部聚起，原本散开的细小波动一点点收拢。蔡须坤没有抢着往里添力，指尖只停在那道气息接回的位置。等牢大点头，她才接上自己的那一段。

木桩向后一震，牢大便收了手。

蔡须坤跟着撤回，脚下只挪了半步。挡板上新添了一道痕，旁边立着的水碗却没有倒。

小舞走过去，看那道痕。

``你们提前说好的？''

``说过。''牢大答，``也试坏过几回。''

奥斯卡想起山谷里熏黑的纸，忍不住笑，牢大看他一眼，他便将水递过来，没有再提。

玉小刚让唐三和小舞分别说明刚才看见了什么。唐三最初只说两股力量接在一起，蔡须坤却指了指牢大的手肘，让他再看开始收回的那一瞬。

第二次，唐三看见了。

``不是一直加上去。你先让他知道，什么时候能收。''

蔡须坤点头。

她仍不擅长将每一种感觉说得像玉小刚那样清楚，便请牢大再慢些做一次，自己只把刚才那一处重现出来。小舞跟着移了一步，问她若换了落点，该怎样让同伴提前知道。

几个人便在场上重新试。

玉小刚没有让他们都挤在同一处。他把五个人拆开，先由唐三和牢大搭一组，小舞站到侧边看，奥斯卡与蔡须坤则去另一侧摆木桩。等一轮做完，再换位置。

唐三起初松了一口气。牢大不像小舞那样爱突然变招，只要照着原先试过的方向送出锤势，似乎便不会再出错。

可真轮到自己，他听见一声``收''，手里的锤子却还在往前。牢大已经撤开，他便赶着去补那点空处，脚底一下踩出了线。

小舞站在旁边喊停，唐三仍想将这一段接完。直到看见蔡须坤伸手拦住搬木桩的奥斯卡，他才猛地醒过来，硬将力量收回。

``差一点。''他说。

小舞蹲到线边，指着被他踩断的细枝：``你差的这一点，刚好就是我站的地方。''

唐三看了她一眼，没再争。他将枝条重新摆好，回到起点。这一次牢大说收，他便先收锤，哪怕后面明明还能续上，也没有再追过去。

换小舞与蔡须坤时，情形又不同。蔡须坤很快能跟上她，小舞便忍不住越走越快，等到转向，才发现对方已经站到了自己原本准备落脚的地方。

她猛地一扭，避开人，自己却在地上滚了一圈。

奥斯卡跑过去拉她。小舞坐在地上，第一句便问蔡须坤，为什么不往旁边让。

蔡须坤一怔，指了指那根枯枝：``说好在这里等。''

小舞看看地上，又看看自己方才滚过的痕迹，耳朵一点点红起来。是她临时觉得另一边更方便，没开口便换了方向。

``我下次先说。''

蔡须坤点头，又补了一句：``我刚才也看见了，可以提醒你。''

两人站回原处。这一回小舞只走到要转身的地方便停下，先伸手指清自己想去哪里。蔡须坤让开一步，她才慢慢翻过去，脚终于稳稳落在空处。

轮到奥斯卡，众人都以为他要先吃。他却把油纸包压在水壶下，拿起场边的挡板，问牢大能不能换自己接一次。

牢大看了看他的手，先让他试最轻的那一下。奥斯卡不大服气，仍将板子抵得很稳；等力量真正落过来，肩头一沉，人才发觉自己把手架得太高，腿也没有站好。

挡板滑了半寸，牢大当即收手。奥斯卡退到线外，甩了甩发麻的手，没有再要求加重。他自己将板子放低，又请唐三站到旁边，帮着看看背后有没有挡住去路。

等几个人都轮过一遍，木桩上的新痕反而没有先前深。小舞累得坐到阴凉处，拿脚拨了拨那包一直没拆的香肠，说今天练得还不如以前热闹。

``你没再摔第二回。''奥斯卡说。

``你板子也没掉第二回。''

两人互看一眼，都笑了。牢大端来水，让奥斯卡先歇歇手，蔡须坤则将倒在地上的枯枝一根根收拢。她留意到唐三最后那次起势以前，已经先向侧边看了一眼，不用再由旁人出声提醒。

玉小刚走过来，没有多夸，只让他们各自说一件今天仍没做好、下回还要接着练的事。

这一次，五个人都没有急着去翻场边的笔记。

那些写在魂技记录里的高阶序数名称仍在：牢大的肘击记着\allowbreak{}\NovelInline{I00076}{\varepsilon_\omega}\allowbreak{}，两人与奥斯卡的挂载配合，合击册上甚至留有\allowbreak{}\NovelInline{I00077}{\varepsilon_{\varepsilon_{\omega+\omega}}}\allowbreak{}的记录。玉小刚没有让他们拿这些名字重新排一遍强弱，只提醒两兄妹：魂技借力有具体条件，合作要在当下能控制的范围内练；看过一次合击，也不等于已经掌握伙伴所有的功法。

到傍晚，唐三能在抬锤前先看一眼小舞，小舞也不再等锤势到了才临时躲。两人仍会吵：一个嫌另一个退得太早，另一个便指出他刚才根本没有照约定收手。可争完以后，他们会把问题留在场上，再试一次。

离开时，小舞走在前面，回头让唐三把护甲拿到铺子去。

``那根带子我来换。''

唐三将护甲卷好，跟了上去。过去两年，他们从学院门口那块看不懂的碑走到这里，仍有许多名字只知道写法，真正能用的东西，却已经一点点落到了手上。

回铺子后，小舞将护甲摊在桌上，才发现要换的还不止一根带子。侧边原有的裂口已经磨得更长，缝在里面的软垫也翻了出来。她拿针穿了好几次，最后烦得把线搁下，盯着那件甲不说话。

``先拆旧的。''唐三将手头的钳子放好，``你隔着旧线，穿不过去。''

``你早说。''

``你刚才没让我说。''

奥斯卡把脸转到门外，没让两人看见自己在笑。牢大过来扶住甲片，等小舞将那段旧线挑开，再一点点替她拨正。蔡须坤拿着新带子，在灯下比了比长短，问多出来的要不要剪掉。

唐三伸手接过，没有立刻剪。他把甲重新套上，抬肩试了一下，小舞从后面拉紧带子，才看见他一动，那一点原本多出来的余量便没了。

``留着。''她说。

蔡须坤将剪子放远些，坐回灯下看她穿线。门外有人来取铁锅，奥斯卡便去搭话。那人问价，他没有自己往下减，转头等小舞说清，再照着账本找钱。

牢大见唐三又要去生炉，先看了看他肩膀：``今天还打？''

``只给他把锅沿磨一下，不用生。''

牢大这才松开风箱的把手，帮着将地上散掉的工具归回原处。蔡须坤看见唐三从包里拿出的纸角压在铁钳底下，连忙抽出来，轻轻抚平，认出正是她年前留给他的那几张。

``这里我看过。''唐三说。

他翻到后面，指给牢大。那一处他也没能完全接下去，只在旁边留了几句自己做过的尝试，没有像往日那样，非要等做完才能拿给别人看。

奥斯卡送完客人回来，俯身一看，先指出自己试过同一个方向。唐三正要把那张抽走，他却按住了纸角。

``留着。省得我们下回忘了，又从这里错一遍。''

小舞终于缝完，咬断线，将护甲扔给唐三。他抬胳膊试过，肩带没有再松，便把那几张纸重新叠好，递给蔡须坤。

她接过去，这次没有独自收进包里，而是放到牢大与奥斯卡中间，等两人一起看完。

那晚，青石山谷里的火又亮起来。

奥斯卡白天看过唐三的练习，吃过饭便把几张旧笔记带来。牢大坐在火旁，将今天没收稳的一处重新记下；蔡须坤把纸接过，指尖停在上面，却没有立即说话。

``怎么了？''牢大问。

``我有时能先感觉到它该往哪里走。''她说，``真要讲清楚，又找不到下一句。''

奥斯卡没有像从前那样，替她把后面全推下去。他将笔放到石头上，让她先把自己确实知道的部分留下。

蔡须坤写了几处，划去一处，最后将纸递给牢大。

牢大看过，指出其中一处自己也接过。奥斯卡便从那里往后问，三个人说到共同不懂的位置，停下来，留给明天。

``我以前总觉得，把每次起手都练稳，就不会有不知道怎么接的时候。''牢大说。

他望着手里那张仍有空处的纸，慢慢将它放平。

``现在知道，还得问。''

奥斯卡把最后一根香肠分开，自己只留一小截：``那就轮着问。反正我也有接错的时候。''

蔡须坤接过来，笑了一下。她很少这样笑，奥斯卡看见，又故意说起自己白天差点弄倒水碗的事，想让她再笑一次。

炉火渐小，三人也不再说话。奥斯卡靠着石头，原本还要核最后一页，手里的纸却先垂下来，慢慢睡着了。

蔡须坤将那张纸接住，放到离火远一点的地方。

过了很久，她忽然问：``要是有一天，我要去一个很高的地方，高到你们够不着呢？''

牢大睁开眼，看了她一会儿。

``那我就在底下给你垫着。''他说，``基底稳，上面才站得住。''

蔡须坤把膝盖抱起来，下巴搁在膝上，看着火堆里最后那一点红，没有再说话。

回学院的路上，牢大提着灯，奥斯卡抱着纸，走两步便打一个哈欠。蔡须坤落在他们身后半步，听着前面两个人说话，过了山梁，才慢慢跟上。

后来几次去山谷，唐三和小舞也赶上了。兄妹俩不能回回都来，铺子有活便得早走；牢大三人却会将当日停下的那张纸单独留下，不叫没来的两人凭空猜前面发生过什么。

小舞最初总嫌他们停得早。等到自己有一次看准了前面的落点，却怎么也说不出后面该如何接，她才肯将笔搁下，先去把已经能做的那一段练稳。

奥斯卡看见了，想说两句，话到嘴边，又想到自己那次忘了起点的笔记。他从包里抽出一张干净纸，分给她半张，自己也重新坐下。

有时争执仍会起来。牢大一急，说话就短；奥斯卡会觉得他又在否定自己的新想法；蔡须坤则常要等两人都停了，才能说出刚才那一处究竟哪里让她不舒服。

唐三来过以后，学会先请他们各自重做一遍。谁的部分没接稳，便由谁慢慢补上，不必把先前所有争论都重新吵一次。小舞坐在石上听，遇到自己看得明白的地方，也会跳下来试。

他们并没有因此一下走得多快。几张纸用完，又添几张；练坏的挡板换了一回，护甲的带子也补了新的线。等到入门课的本子边角都磨圆，唐三翻开末尾那页速查，才发现有些从前只能照抄的名字，已经能和自己亲手做过的事对上。

他将本子拿给牢大看，牢大却没先看那一页，倒是指着最前面提水那晚留下的半页涂改，问他还记不记得。

唐三看了半晌，笑出来。奥斯卡凑近一看，也认出了自己的字迹，说那时候大家还以为考核结束，就能立刻去摸门口碑上更高的刻痕。

小舞催他们收东西。蔡须坤把纸按原来的顺序拢好，走到山道前，等唐三合上本子。牢大提起灯，奥斯卡将空下来的油纸折好，五个人才一前一后往山下去。

初级学院的日子将尽。唐三和小舞还不知道父亲去了哪里，伙伴们也没有走完碑上那些更高的路标。可等到前往中级学院的日子，他们已经约好，同路走一程。

\chapter{第六章：密不可分的伙伴}\label{ux7b2cux516dux7ae0ux5bc6ux4e0dux53efux5206ux7684ux4f19ux4f34}

初级学院的毕业典礼落幕，唐三和小舞背着简单的行囊，踏上了前往中级学院的路途。沿途青山连绵，林间风卷着草木的清香，两人一边赶路，一边切磋着大数构造的技巧，偶尔停下来交流高德纳箭头运算的心得，倒也不觉得枯燥。小舞指尖萦绕着淡淡的灵气，随口报出一串高阶迭代大数，唐三则快速回应，用更精妙的构造方法将其优化，两人的默契，早已在日复一日的相处中刻入骨髓。

\begin{center}\rule{0.5\linewidth}{0.5pt}\end{center}

``老哥，你看前面！''小舞突然停下脚步，指着不远处的林间小道，语气里带着几分好奇和警惕。唐三顺着她指的方向望去，只见两道陌生的身影正站在路边，似乎在争执什么------男子身形挺拔，周身萦绕着冰蓝色的灵气，手中隐约有冰红茶的虚影浮现，眉头紧锁；女子一袭黑衣，手中握着一枚泛着乌黑光泽的篮球，篮球表面萦绕着微弱的篮球序力，正对着男子轻声说着什么，那枚篮球显然是她的扽魂，蕴含着强悍的序力波动。两人周身的灵气波动都不弱，显然也是前往中级学院的修士。

唐三放缓脚步，眼神警惕地打量着两人，低声对小舞说道：``老妹，小心点，他们的序力层级不低，不知道是什么人。''小舞轻轻点头，指尖悄悄凝聚起粉色灵气，转头对唐三小声回应：``老哥，我知道了，你也别大意。''就在这时，那名冰蓝色灵气的男子似乎察觉到了他们的目光，转头看了过来，眉头依旧紧锁，语气带着几分疏离：``你们是谁？也是往中级学院去的？''

唐三上前一步，语气平和却带着几分警惕，拱手道：``在下唐三，这是我老妹小舞，我们确实是前往中级学院报到的。方才见二位似乎在争执，不知发生了什么事？''黑衣女子闻言，无奈地笑了笑，指了指前方不远处一座破旧的木屋：``我叫蔡须坤，他是牢大。我们刚才路过这里，听到木屋地下室里有微弱的呼救声，我想进去看看，牢大却说怕有陷阱，耽误赶路。''牢大补充道，语气缓和了些许：``我和蔡须坤也是要去中级学院，她性子急，见不得有人被困，我是担心贸然进去会有危险。''唐三笑着点头：``能理解，救人固然重要，稳妥些也没错。''小舞也凑上前来，笑着说道：``我和老哥也觉得该去看看，多一个人多一份力，我们一起进去，也好相互有个照应。''

唐三眼神一凝，凝神细听，果然听到木屋方向传来断断续续的呜咽声，微弱却清晰，不似常人的呼救，反倒带着几分稚嫩。``不对劲，这声音像是小孩子的，而且不止一个。''唐三沉声道，``中级学院报名还有三天时间，耽误半天没关系，我们去看看，说不定有人遇到了危险。''

牢大闻言，看了唐三和小舞一眼，又看了看破旧的木屋，沉默片刻后点了点头：``也好，我这冰红茶扽魂能凝聚防御屏障，就算有陷阱也能应付。蔡须坤，你用扽魂篮球探查周围的动静，篮球的序力波动能感知到隐藏的气息，防止有埋伏；唐三、小舞，你们两个负责正面接应，我们暂时分工合作，先看看里面的情况。''几人快速敲定计划，朝着破旧木屋悄悄靠近，唐三边走边轻声对牢大和蔡须坤说道：``我老妹的扽魂是兔子，速度快、爆发力强，近战很厉害；我的话，擅长构造大数能量，能远程支援。''蔡须坤点点头，回应道：``我扽魂是篮球，能操控篮球发动攻击，还能靠篮球序力探查环境；牢大的冰红茶扽魂，防御和近战都很不错，等会儿他来扛伤害，我们配合支援。''牢大也补充道：``放心，有我在，能护住你们。''彼此间的警惕渐渐消散，多了几分默契。

\begin{center}\rule{0.5\linewidth}{0.5pt}\end{center}

木屋看起来废弃已久，门板破旧不堪，布满了灰尘和裂痕，推门进去，一股腐朽的气息扑面而来。地下室的入口藏在木屋角落的杂物堆后面，被一块厚重的石板挡住，石板上刻着粗糙的纹路，隐约能看到上面残留的黑暗系灵气波动------显然，这里确实有异常。

``是狼人的气息。''蔡须坤皱了皱眉，手中扽魂篮球轻轻转动，篮球序力顺着篮球纹路扩散开来，朝着地下室周围探查，``篮球序力能感受到，地下室里有几只狼人，还有七个微弱的生命气息，应该就是呼救的人。''牢大上前一步，双手凝聚起冰蓝色的灵气，冰红茶扽魂全力发动------

\begin{quote}
\textbf{``冰红茶·防御屏障！''}
\end{quote}

一道冰蓝色的屏障瞬间笼罩住四人。``我来移开石板，你们做好战斗准备。''

话音落下，牢大双臂发力，冰蓝色的灵气灌注在石板上，石板缓缓被抬起，露出一个漆黑的地下室入口，呜咽声瞬间变得清晰起来，还夹杂着狼人的低吼。小舞身形一闪，兔子扽魂全力发动，纵身跃到入口旁，指尖凝聚起粉色灵气，随时准备发动攻击；唐三则运转灵气，双手构造出简单的高阶大数能量，高德纳箭头运算快速运转，做好了随时支援的准备。

``吼------！''地下室里传来几声凶狠的低吼，三只体型庞大的狼人从黑暗中冲了出来，周身萦绕着浑浊的灵气，獠牙外露，眼神凶狠，显然是这里的看守。``果然是狼人！''牢大低喝一声，冰蓝色灵气暴涨------

\begin{quote}
\textbf{``冰红茶·肘击！''}
\end{quote}

将冰红茶的能量凝聚在肘部，朝着最前面的狼人狠狠撞去，冲击力裹挟着高阶迭代大数的能量，瞬间将狼人撞得踉跄后退。

蔡须坤见状，厉声催动扽魂篮球：

\begin{quote}
\textbf{``篮球回旋！''}
\end{quote}

手中篮球裹挟着浓郁的篮球序力，飞速旋转起来，化作一道暗黑旋风，如同灌篮般朝着另外两只狼人席卷而去，篮球序力瞬间瓦解了狼人的灵气防御，狼人发出痛苦的嘶吼，动作变得迟缓起来，这一招结合了篮球的冲击力与篮球序力的瓦解力，威力远超普通攻击。

唐三和小舞也立刻加入战斗。小舞纵身跃起，对着唐三喊了一声：``老哥，我去牵制左边那只！''话音未落，她双腿发力，如同敏捷的兔子，朝着一只狼人狠狠蹬去，动作和当初切磋时一样干脆利落，却少了几分打闹，多了几分凌厉，狼人被蹬得重重摔在地上；唐三见状，快速构造出高阶大数能量，对着小舞喊道：``老妹小心，我来帮你！''指尖凝聚的能量精准地朝着狼人身上的薄弱部位射去，高阶能量瞬间击穿了狼人的防御，狼人哀嚎着倒在地上，失去了战斗力。牢大看着两人的配合，忍不住喊道：``唐三，小舞，你们兄妹俩配合得真默契！''蔡须坤也点头附和：``是啊，刚才那一下衔接得太及时了。''

短短几分钟，三只狼人就被四人联手击溃，瘫倒在地，再也无力反抗。蔡须坤操控扽魂篮球，让篮球顺着地下室通道滚动探查，确认没有其他埋伏后，对着众人摆了摆手：``安全了，里面没有其他狼人了。''唐三率先走下地下室，地下室里漆黑一片，他运转灵气，指尖泛起淡淡的光芒，照亮了周围的环境。

\begin{center}\rule{0.5\linewidth}{0.5pt}\end{center}

只见地下室的角落里，七个身形矮小的小矮人被铁链锁住，身上布满了灰尘和伤痕，脸上满是恐惧，看到有人下来，纷纷缩了缩身体，眼神里充满了警惕。``别害怕，我们是来救你们的。''小舞轻声说道，语气温柔，慢慢走上前，``我们不会伤害你们的。''

牢大上前，冰蓝色的灵气凝聚在指尖，轻轻一挥，锁住小矮人的铁链瞬间被冻裂，``咔嚓''几声，铁链纷纷断裂。小矮人们缓缓抬起头，看着四人，眼中的警惕渐渐褪去，取而代之的是感激。``谢谢你们，谢谢你们救了我们。''最小的一个小矮人小声说道，声音还有些颤抖。

唐三蹲下身，轻声问道：``你们怎么会被狼人囚禁在这里？''其中一个看起来年纪稍大的小矮人叹了口气，缓缓说道：``我们是附近山林里的居民，几天前被这些狼人掳来，他们说要把我们当成诱饵，吸引其他修士过来，没想到你们会发现我们。''

蔡须坤看着小矮人们身上的伤痕，手中扽魂篮球泛起微弱的篮球序力，轻轻靠近小矮人的伤口，篮球序力虽然擅长瓦解，却也能缓慢修复伤势：``别担心，我们送你们回去，这些狼人已经被我们击溃，不会再伤害你们了。''奥斯卡虽然不在，但小舞和唐三身上带着一些修炼用的灵气结晶，小舞拿起几枚递给小矮人们，笑着说道：``快拿着补充体力，以后可别再乱跑啦。''唐三也补充道：``要是再遇到危险，就大声呼救，说不定我们还能听到。''蔡须坤看着兄妹俩温柔的样子，轻声说道：``你们兄妹俩心肠真好，难怪配合得这么好。''牢大也笑道：``唐三沉稳，小舞活泼，倒是互补。''小舞吐了吐舌头，挽住唐三的胳膊：``那是，我老哥最疼我了。''唐三无奈地笑了笑，轻轻揉了揉小舞的头发：``你啊，别调皮。''

\begin{center}\rule{0.5\linewidth}{0.5pt}\end{center}

四人搀扶着七个小矮人，慢慢走出地下室，来到木屋外的阳光下。小矮人们看着眼前的四人，脸上满是感激，纷纷鞠躬道谢：``谢谢你们，善良的修士们，如果不是你们，我们恐怕还会被狼人囚禁着。''牢大摆了摆手，笑着说道：``举手之劳而已，你们赶紧回家吧，以后注意安全，不要再单独出来了。''

小矮人们再次道谢后，朝着山林深处走去，直到身影消失在林间小道上，四人才转身看向彼此。经过刚才的并肩作战和相处，几人间的陌生感彻底消散，小舞笑着看向牢大和蔡须坤：``今天真是太谢谢你们了，要不是你们，我们可能也没法这么快救出小矮人们。对了，你们也是去中级学院吗？以后说不定就是同学啦。''蔡须坤笑着点头：``是啊，我们也是去中级学院报到，以后就是同学了。''牢大也哈哈大笑：``可不是嘛，刚才一起并肩作战，也算是战友了。唐三，小舞，你们兄妹俩实力真不错，以后有机会，还能切磋切磋。''

唐三笑着回应：``好啊，求之不得。我老妹最喜欢切磋了，到时候你们可别手下留情。''小舞立刻说道：``老哥，你别揭我老底！不过话说回来，蔡须坤，你的篮球扽魂真厉害，刚才那招回旋太帅了；牢大，你的冰红茶防御也超棒，刚才多亏你护着我们。''蔡须坤被夸得有些不好意思，轻声说道：``你们也很厉害，唐三的大数能量精准度太高了，小舞的速度也快得惊人。既然都是前往中级学院，不如一起同行？路上也好有个照应，我们也能多熟悉熟悉。''

唐三和小舞对视一眼，点了点头。唐三笑着说道：``好啊，有二位同行，确实能安心不少。我和老妹从小一起长大，她性子活泼，有时候有点冲动，以后要是有做得不对的地方，还请二位多担待。''小舞嗔怪地看了唐三一眼：``老哥，你才冲动呢！''蔡须坤笑着说道：``放心吧，我们不会介意的。我和牢大也是从小一起修炼，配合久了，也有不少小毛病。以后我们四个一起，相互照应、相互学习就好。''牢大也点头：``没错，以后就是伙伴了，不用这么客气。''

四人收拾好心情，继续踏上前往中级学院的路途。小舞走在中间，一边挽着唐三的胳膊，一边和蔡须坤聊着修炼的趣事，唐三和牢大则跟在旁边，交流着大数构造和扽魂运用的技巧，阳光透过树叶的缝隙，洒在四人身上，空气中弥漫着淡淡的灵气波动和欢声笑语。这场意外的相遇，不仅让四人联手解救了七个小矮人，更让他们结识了新的伙伴，从陌生到熟悉，彼此间多了一份羁绊。他们对即将到来的中级学院生活，也多了几分期待------毕竟，有兄妹相伴，有伙伴同行，无论遇到什么困难，都能携手并肩，共同面对。而他们不知道的是，中级学院里，还有更多的挑战和惊喜，在等待着他们。

中等学院的校门立在的山阶尽头，比起初级学院的质朴，这里的每一块石砖都刻着淡淡的序数纹路，风一吹，纹路便会泛起微光，那是 FGH 功法与天地灵气交融的痕迹。唐三背着磨得发亮的 Buchholz 防弹衣，指尖摩挲着衣料上的裂痕------那是和小舞三年切磋留下的印记，身后的小舞晃着手里的燃烧弹，奥斯卡则掂着一兜香肠，三人的身影落在山阶上，被余晖揉成温暖的光斑。

刚踏入学院，一股远比初级学院浓郁的灵气便扑面而来，唐三能清晰感受到，这灵气中裹挟着层层叠叠的序数波动，有的浅淡如 ω，有的厚重如 ω\^{}ω\^{}ω，甚至还有几道隐晦的气息，竟触碰到了 ε₀ 的边缘。``这地方的修士，果然和初级学院不是一个量级。''小舞压低声音，指尖捏紧了燃烧弹，眼底却藏着跃跃欲试的光。

玉小刚早已在学院的演武场等候，身旁站着两个身影，一个身形挺拔，手臂肌肉线条凌厉，正是牢大；另一个身形清瘦，指尖绕着一缕淡金色的灵气，自然是蔡须坤。牢大看到唐三，咧嘴一笑，抬手便是一记轻缓的肘击，空气里瞬间泛起一道 ε\_ω 的序数涟漪。``早听说初级学院出了个天才，能在六岁摸到 ω 的门槛，今日一见，果然名不虚传。''

蔡须坤则微微颔首，指尖的金色灵气轻轻飘向奥斯卡，在他的香肠袋上绕了一圈，那些香肠竟瞬间泛起淡淡的银光：``奥斯卡的香肠挂载属性，我已用稍加改良，如今能临时提升挂载序数的稳定性，你们往后切磋，倒能多一层保障。''

奥斯卡眼睛一亮，立刻掏出一根香肠咬了一口，只觉一股温润的灵气顺着喉咙涌入丹田，丹田内的 ω\^{}ω\^{}ω 序数竟微微震颤，似是被激活一般：``这也太神奇了！以前我的香肠只能加些基础灵气，现在居然能触动高阶序数波动了！''

\begin{center}\rule{0.5\linewidth}{0.5pt}\end{center}

玉小刚摆了摆手，目光扫过唐三、小舞和奥斯卡，沉声道：``中等学院的核心，不是单纯提升序数层级，而是将序数功法与自身魂技深度融合。初级学院你们掌握了 FGH 的基础，能构造出 ω\^{}ω\^{}ω 的大数，但那只是纸上谈兵，真正的强者，能让序数成为魂技的骨，灵气成为魂技的血。''

他抬手一指演武场中央的石桩，那石桩上刻着密密麻麻的序数符号，从 ω 到 ε₀，再到浅层 Veblen 序数，层层递进：``今日的第一课，便是\textbf{魂技序数化}。唐三，你用蓝银草；小舞，你用瞬移；奥斯卡，你用香肠挂载，三人配合，击破这根 φ(1,0) 也就是 ε\_0 级的石桩。'' \textgreater{} \textgreater{} 映射过程中蕴含的不动点结构可以启示我们构造出更大的序数。我们总可以找到一个递增的映射 α → f(α)，然后将其不断地作用在某个序数之上，这将增大这个序数的取值。 \textgreater{} \textgreater{} 在进行了任意有限次映射之后，我们将得到这个映射的一个极限点。如果这个极限点不是不动点，那么可以直接从极限点出发继续映射。而如果这个极限点是不动点，则可以通过取后继的方式跳出这个不动点，然后继续进行映射。这样不断地重复下去，我们将抵达下一个极限点，从而可以进行类似的操作。 \textgreater{} \textgreater{} 我们在后文中考虑的映射通常都是具有不动点的映射，因此可以不断地重复迭代到不动点、然后取后继跳出的过程，这样可以不断地得到越来越强大的不动点。等到这一不断地取不动点的过程也难以为继的时候，我们就可以从先前的迭代过程中生成一个更加强大的映射 α → f(α)，然后继续枚举这一新映射的不动点。这里所考虑的映射 f(α) 越强，它所能够得到不动点也就越大。恰当地组织并且标记各映射过程及其不动点的层次，我们将能够得到一个更加强大的序数记号，它可以带领我们在序数世界中前进得更远。我们在后文中将对此进行详细的介绍。

\begin{center}\rule{0.5\linewidth}{0.5pt}\end{center}

唐三闻言，指尖轻抬，数根蓝银草从掌心窜出，蓝银草的藤条上，瞬间浮现出 ω\^{}ω\^{}ω\^{}ω 的序数纹路------这是他目前能稳定操控的最高序数。他心念一动，蓝银草如游龙般缠向石桩，藤条上的序数纹路亮起，试图嵌入石桩的防御层。可石桩上的 φ(1,0)（即 ε\_0）纹路只是轻轻一颤，便将蓝银草的力量弹开，唐三的丹田微微一震，竟生出一丝滞涩感。

``序数层级不够，硬拼是没用的。''牢大的声音从旁传来，他缓步上前，手臂抬起，ε\_ω 的序数涟漪在掌心凝聚，``序数功法的精髓，是层级嵌套与借力，而非单一层级的蛮力冲击。小舞，用你的瞬移术锁定石桩的 φ(1,0) 纹路节点，唐三，将你的蓝银草序数拆解为 ω\^{}ω、ω\^{}ω\^{}ω、ω\^{}ω\^{}ω\^{}ω 三层，层层嵌套，奥斯卡，把你的香肠挂载在蓝银草的每一层序数上，蔡须坤，用序数绑定术将牢大的 ε\_ω 与蓝银草的嵌套序数联结。''

几人立刻行动，小舞的身影化作一道残影，瞬移间指尖点向石桩的几处纹路节点，那些节点瞬间被她的 ω\^{}(ω\^{}(ω²)) 序数标记；唐三深吸一口气，将丹田内的 ω\^{}ω\^{}ω 序数强行拆解，三层高阶 ω 幂次的序数纹路在蓝银草上层层铺开，如三道蓝色的光环；奥斯卡掏出三根香肠，分别贴在蓝银草的三层序数上，香肠的银光与蓝银草的蓝光交融；蔡须坤的指尖金光大盛，一道纤细的金色纽带凭空出现，一头联结着牢大掌心的 ε\_ω（即 φ(1,ω)），一头缠上蓝银草的嵌套序数，序数波动在纽带中飞速流转，竟将 ε\_ω 的力量，层层渗透进蓝银草的每一层纹路里。

\textbf{``动手！''} 玉小刚一声令下。

唐三心念剧震，三层嵌套的蓝银草猛地收紧，缠向石桩，此时的蓝银草，每一层都裹挟着 ε\_ω 的序数力量，再加上奥斯卡香肠的挂载增益，序数波动如海啸般冲击着石桩的防御层。小舞则不断瞬移，在石桩的节点上不断标记，瓦解着石桩的 φ(1,0) 纹路。

只听''咔嚓''一声脆响，那根 φ(1,0) 级的石桩，竟从中间裂开一道缝隙，缝隙中，石桩的序数纹路寸寸碎裂，最终轰然倒塌。

唐三松了口气，蓝银草缓缓收回，丹田内的序数虽有些紊乱，但他能清晰感受到，自己对 ω\^{}ω\^{}ω 的操控，竟比之前熟练了数倍，甚至对 φ(1,0) 的层级，也有了一丝模糊的感知。``原来如此，序数不是单独的数字，而是可以和伙伴的力量相互联结的纽带。''他喃喃道。

小舞擦了擦额角的薄汗，咧嘴一笑，眼底的警惕淡了几分：``以前我总觉得，靠自己的力量就够了，现在才知道，和伙伴一起，能做到单独一个人永远做不到的事。''

牢大走上前，拍了拍唐三的肩膀，肘间的序数涟漪轻轻扫过唐三的丹田，帮他抚平了紊乱的序数：``这就是中等学院的修炼之道，序数为基，伙伴为翼。我们这些修炼序数功法的人，看似在比拼大数的构造，实则在比拼谁能找到最契合的伙伴，将彼此的序数力量发挥到极致。''

蔡须坤点了点头，补充道：``就像我和牢大，我的序数绑定术看似普通，但若没有牢大的 ε\_ω 肘击，便毫无用武之地；而牢大的肘击，若没有我的绑定术，也无法将力量精准传递，更别说和其他人的力量联结。''

\begin{center}\rule{0.5\linewidth}{0.5pt}\end{center}

玉小刚看着眼前的五人，眼中满是欣慰：``你们今日的配合，虽仍有生涩，却已摸到了中等学院修炼的核心。接下来的日子，你们五人将组成一个小队，一起修炼，一起完成学院的任务------猎杀高阶魂兽，拆解它们的魂核序数，提升自身的序数底蕴。''

他抬手一挥，一道卷轴飘到几人面前，卷轴上写着几个字：\textbf{猎杀千年人面魔蛛，提取其魂核中的 φ(ω,0) 序数}（强于 ζ₀ 的高阶递归序数）。

``千年人面魔蛛的魂核，蕴含着 φ(ω,0)（先别管这是啥玩意，后面会介绍，知道比 ε\_0 和 ζ\_0 大就行）级的序数波动，其防御更是达到了 φ(4,0)（也比 ε\_0 和 ζ\_0 大，但是比上面那个小）级别，单打独斗，即便是牢大，也未必能讨到好处。''玉小刚的语气郑重，``这是你们小队的第一个任务，也是对你们配合的第一次考验。记住，序数的力量再强，也抵不过伙伴间的同心协力。''

\begin{center}\rule{0.5\linewidth}{0.5pt}\end{center}

唐三看着卷轴上的字，又看了看身旁的小舞、奥斯卡、牢大和蔡须坤，五人的目光交汇，眼底都燃起了斗志。他能感受到，一种名为''伙伴''的羁绊，正随着序数的波动，在五人之间悄然凝聚。

小舞捏了捏燃烧弹，笑道：``千年人面魔蛛是吧？正好试试我们的配合有多厉害！''奥斯卡拍了拍香肠袋：``放心，我的香肠肯定能给大家挂好增益，保证不掉链子！''牢大活动了一下手臂，ε\_ω 的序数涟漪在周身环绕：``走，让那只蜘蛛见识见识，什么叫高阶序数小队的力量！''

在猎杀千年人面魔蛛之前，五人的身影并肩走出演武场，夕阳将他们的身影交织在一起，空气中的序数波动相互交融，形成一道独特的光环。唐三的蓝银草、小舞的瞬移、奥斯卡的香肠、牢大的肘击、蔡须坤的序数绑定，五种截然不同的力量，因序数而联结，因伙伴而紧密。

他们不知道，这场猎杀千年人面魔蛛的任务，只是中等学院修炼的开始，更不知道，在前方的道路上，还有更强大的对手、更高阶的递归序数、更艰难的考验在等待着他们。但他们心中清楚，只要五人同心，以序数为基，以伙伴为翼，便没有跨不过的坎，没有构造不出的大数，更没有闯不出的传奇。

云雾缭绕的山林间，五人的身影渐渐消失，只留下层层叠叠的高阶序数涟漪，在空气中轻轻荡漾，预示着一场属于他们的序数冒险，正悄然拉开新的序幕。

\chapter{第七章：ε₀和ζ₀}\label{ux7b2cux4e03ux7ae0ux3b5ux2080ux548cux3b6ux2080}

山林边缘的空地上,赵无极身形魁梧如铁塔,周身萦绕着淡淡的φ(3,0)序数波动,比牢大的εω还要厚重几分。他看着眼前的五人小队,声音洪亮如钟,打破了山林的静谧：``在猎杀千年人面魔蛛之前,我会帮你们做战前演练,更重要的是,我要给你们讲清楚ε₀、ζ₀这些关键序数的性质------你们要对付的人面魔蛛,魂核序数是φ(ω,0),而它的防御达到φ(4,0),想要取胜,必须先吃透这些序数的核心,知道它们的强弱层级和突破路径。''

唐三上前一步,拱手问道：``赵老师,我们在初级、中级学院都听过ε₀、ζ₀这些序数,知道它们是高阶递归序数,但始终没弄明白它们的具体性质,还有您说的'基本列',到底是什么意思？''

赵无极点了点头,抬手在空中凝聚出一道淡白色的序数纹路,正是ε₀的符号,缓缓开口：``问得好！首先我们从ε₀说起,它是第一个无法用ω的加法、乘法、幂次构造出来的递归序数,也是你们从基础ω幂次进阶到高阶序数的关键门槛,它的特殊性远超普通高阶序数,我详细给你们说清楚。''

``第一,ε₀是算术序数的绝对上限。''赵无极指尖一动,空中浮现出算术序数的序列,``你们之前掌握的ω、ω²、ω\textasciicircum{}\allowbreak{}ω、ω\textasciicircum{}\allowbreak{}ω\textasciicircum{}\allowbreak{}ω,\allowbreak{}哪怕是你们现在突破后的层级,\allowbreak{}本质上都是’算术序数’——也就是能通过ω的有限次加法、乘法、幂次迭代构造出来的序数。而ε₀,\allowbreak{}无论你用多少次ω的基础运算叠加、迭代,\allowbreak{}都无法构造出它,\allowbreak{}它是第一个跳出’算术运算’范畴的递归序数,\allowbreak{}是从’有限迭代’到’无穷迭代’的分界线。''

``第二,ε₀是最小的'ε-序数'。''他顿了顿,继续讲解,``所谓ε-序数,\allowbreak{}就是满足’ε=\allowbreak{}ω\textasciicircum{}\allowbreak{}ε’这个等式的序数,\allowbreak{}而ε₀是所有满足这个等式的序数中,\allowbreak{}最小的一个。简单说,\allowbreak{}它是ω的无穷幂次迭代（ω\textasciicircum{}\allowbreak{}(ω\textasciicircum{}\allowbreak{}(ω\textasciicircum{}\allowbreak{}(…)\allowbreak{})\allowbreak{})\allowbreak{}）收敛的结果,\allowbreak{}这种’自指式’的无穷迭代,\allowbreak{}是它区别于所有算术序数的核心,\allowbreak{}也是它难以突破的关键——你们之前的迭代,\allowbreak{}都是有限次的,\allowbreak{}而ε₀需要领悟’无穷迭代’的精髓,\allowbreak{}才能真正触及。''

``第三,ε₀的递归性与可构造性。''赵无极的语气愈发郑重,``虽然ε₀无法用算术运算构造,但它是'递归序数',也就是说,我们能通过有限的规则,构造出它的基本列（也就是单调递增、收敛到它的序数序列）,这也是你们能修炼、掌控它的前提。如果是不可递归的序数,哪怕知道它的定义,也无法通过修炼去突破、去操控。而且,ε₀的基本列有明确的构造规则,这也是它比更高阶序数更容易入门的原因。''

``第四,ε₀是高阶序数的'敲门砖'。''他看向五人,``你们想要突破到ζ₀、φ函数序数,必须先突破ε₀------ζ₀是ε序列的无穷迭代,而ε序列的起点就是ε₀；φ(ω,0)等高阶φ函数序数,更是建立在ε₀、ζ₀的基础上。没有ε₀的基础,后续的高阶序数修炼就是空中楼阁,哪怕强行提升,也无法稳定掌控力量。''

唐三眼神发亮,下意识点头：``赵老师,我明白了,ε₀的核心特殊性,就是'跳出算术运算、最小ε-序数、可递归构造、高阶入门门槛'这四点,难怪我们之前始终无法触及它的精髓,原来是没领悟'无穷迭代'和'自指等式'的核心。''

``没错,领悟得很快。''赵无极赞许点头,``既然说到了ε₀,我就顺便给你们讲一讲ε₁------大家要注意,ε1的构造核心的是'跳出ε₀这个不动点',这和你们之前理解的迭代逻辑不一样,我结合你们能听懂的修炼逻辑,一步步讲,结合你们之前掌握的ω、ε₀知识,保证你们能吃透。''

``首先,\allowbreak{}我们要明确一个关键：ε₀是映射α↦ω\textasciicircum{}\allowbreak{}α的第一个不动点,\allowbreak{}什么是不动点？就是满足’ω\textasciicircum{}\allowbreak{}α=\allowbreak{}α’这个等式的序数,\allowbreak{}ε₀就是最小的满足这个等式的序数,\allowbreak{}也就是ω\textasciicircum{}\allowbreak{}(ω\textasciicircum{}\allowbreak{}(ω\textasciicircum{}\allowbreak{}(…)\allowbreak{})\allowbreak{})\allowbreak{}=\allowbreak{}ε₀。''赵无极放慢语速,确保每个人都能跟上,``但序数世界没有尽头,\allowbreak{}ε₀绝非终点,\allowbreak{}想要继续前进,\allowbreak{}第一步就是跳出这个不动点——怎么跳？很简单,\allowbreak{}给ε₀取后继,\allowbreak{}得到ε₀+\allowbreak{}1,\allowbreak{}这样就成功跳出了映射α↦ω\textasciicircum{}\allowbreak{}α的不动点,\allowbreak{}因为ω\textasciicircum{}\allowbreak{}(ε₀+\allowbreak{}1)\allowbreak{}≠ε₀+\allowbreak{}1,\allowbreak{}它不再满足不动点的等式。''

小舞：为什么ε₀不写成ω↑↑ω

赵无极：问得好,这就要涉及到序数和自然数的区别了,首先

ω\textasciicircum{}\allowbreak{}ω=\allowbreak{}sup\{ω\textasciicircum{}\allowbreak{}n,\allowbreak{}n<ω\},\allowbreak{}你想想,\allowbreak{}如果有ω↑↑(ω+\allowbreak{}1)\allowbreak{},\allowbreak{}会发生什么？

唐三那就是ω↑(ω↑↑ω)\allowbreak{}=\allowbreak{}ω\textasciicircum{}\allowbreak{}ε₀=\allowbreak{}ε₀,\allowbreak{}等下,\allowbreak{}没有跳出不动点,\allowbreak{}而是被吸收了,\allowbreak{}

赵无极：对啊,所以我才强调ε₀是算术序数的绝对上限,这样,就必须使用后继跳出不动点,才能突破这一限制。

唐三立刻提笔（指尖凝聚灵气,在空中写下序列）,轻声补充：``赵老师,我是不是可以这样理解：ε₀是这个映射的'平衡点',取后继之后,ε₀+1就偏离了这个平衡点,就能继续通过映射迭代得到更大的序数？''

``完全正确！''赵无极点头,\allowbreak{}指尖在空中浮现出迭代序列,\allowbreak{}“跳出不动点之后,\allowbreak{}我们继续对ε₀+\allowbreak{}1反复应用映射α↦ω\textasciicircum{}\allowbreak{}α,\allowbreak{}就能得到一系列越来越大的序数,\allowbreak{}具体来说就是：

第一步,ε₀+1经过映射,得到ω\^{}(ε₀+1)；

第二步,再将ω\^{}(ε₀+1)代入映射,得到ω\^{}(ω\^{}(ε₀+1))；

第三步,继续代入,得到ω\^{}(ω\^{}(ω\^{}(ε₀+1)))；

以此类推,不断重复这个映射过程,我们就得到了一个单调递增的序数序列：

\{ε₀,\allowbreak{}ω\textasciicircum{}\allowbreak{}(ε₀+\allowbreak{}1)\allowbreak{},\allowbreak{}ω\textasciicircum{}\allowbreak{}(ω\textasciicircum{}\allowbreak{}(ε₀+\allowbreak{}1)\allowbreak{})\allowbreak{},\allowbreak{}ω\textasciicircum{}\allowbreak{}(ω\textasciicircum{}\allowbreak{}(ω\textasciicircum{}\allowbreak{}(ε₀+\allowbreak{}1)\allowbreak{})\allowbreak{})\allowbreak{},\allowbreak{}…\}。”

小舞皱了皱眉,指尖的燃烧弹轻轻转动,追问道：``那ε1就是这个序列的极限吗？和之前ε₀的无穷迭代极限有什么区别？''

``问得好,核心区别就在'跳出不动点'。''赵无极语气郑重,``ε₀是映射α↦ω\textasciicircum{}\allowbreak{}α的第一个不动点,\allowbreak{}而ε1,\allowbreak{}就是这个’跳出ε₀后、反复迭代映射’得到的序列的上确界（也就是极限）,\allowbreak{}用严格的序数表示就是：ε1=\allowbreak{}sup\{ε₀,\allowbreak{}ω\textasciicircum{}\allowbreak{}(ε₀+\allowbreak{}1)\allowbreak{},\allowbreak{}ω\textasciicircum{}\allowbreak{}(ω\textasciicircum{}\allowbreak{}(ε₀+\allowbreak{}1)\allowbreak{})\allowbreak{},\allowbreak{}ω\textasciicircum{}\allowbreak{}(ω\textasciicircum{}\allowbreak{}(ω\textasciicircum{}\allowbreak{}(ε₀+\allowbreak{}1)\allowbreak{})\allowbreak{})\allowbreak{},\allowbreak{}…\}。''

他顿了顿,补充道：``而且这里有个关键性质,\allowbreak{}当我们把ε₁代入映射α↦ω\textasciicircum{}\allowbreak{}α时,\allowbreak{}会发现ω\textasciicircum{}\allowbreak{}ε1=\allowbreak{}sup\{ω\textasciicircum{}\allowbreak{}(ε₀+\allowbreak{}1)\allowbreak{},\allowbreak{}ω\textasciicircum{}\allowbreak{}(ω\textasciicircum{}\allowbreak{}(ε₀+\allowbreak{}1)\allowbreak{})\allowbreak{},\allowbreak{}…\}=\allowbreak{}ε1,\allowbreak{}所以ε1也是映射α↦ω\textasciicircum{}\allowbreak{}α的不动点,\allowbreak{}是紧随ε₀之后的第二个不动点,\allowbreak{}我们可以记为ε1=\allowbreak{}2ndα↦ω\textasciicircum{}\allowbreak{}αfp.（fp.就是不动点fixed point的缩写）。''

``那ε1的两种基本列,是不是也和这个'跳出不动点、反复映射'的逻辑有关？''蔡须坤指尖的金色灵气轻轻流转,适时问道。

``正是如此,\allowbreak{}ε1的两种基本列,\allowbreak{}都源于映射α↦ω\textasciicircum{}\allowbreak{}α的不动点构造逻辑,\allowbreak{}但两种基本列的迭代基础不同,\allowbreak{}而且有一个关键区别——第二种基本列比第一种略微强一点,\allowbreak{}我结合你们能听懂的修炼逻辑,\allowbreak{}详细讲给你们,\allowbreak{}这对你们后续突破ε₀、逼近ε₁至关重要。''赵无极抬手,在空中划分出两个区域,标注''ε₁第一种基本列（ω幂次迭代型）``和''ε₁第二种基本列（ε₀指数塔型）``,逐一讲解：

“第一种基本列,\allowbreak{}是ω幂次迭代型,\allowbreak{}核心是’从ε₀+\allowbreak{}1开始,\allowbreak{}反复迭代映射α↦ω\textasciicircum{}\allowbreak{}α’,\allowbreak{}也是我们最开始构造ε₁的核心路径,\allowbreak{}递归定义如下：

ε₁{[}0{]}=ε₀+1（起点,跳出不动点的第一个序数）；

ε₁{[}1{]}=ω\^{}(ε₁{[}0{]})=ω\^{}(ε₀+1)（第一次映射）；

ε₁[2]=\allowbreak{}ω\textasciicircum{}\allowbreak{}(ε₁[1])\allowbreak{}=\allowbreak{}ω\textasciicircum{}\allowbreak{}(ω\textasciicircum{}\allowbreak{}(ε₀+\allowbreak{}1)\allowbreak{})\allowbreak{}（第二次映射）；

ε₁[3]=\allowbreak{}ω\textasciicircum{}\allowbreak{}(ε₁[2])\allowbreak{}=\allowbreak{}ω\textasciicircum{}\allowbreak{}(ω\textasciicircum{}\allowbreak{}(ω\textasciicircum{}\allowbreak{}(ε₀+\allowbreak{}1)\allowbreak{})\allowbreak{})\allowbreak{}（第三次映射）；

以此类推,\allowbreak{}ε₁[n+\allowbreak{}1]=\allowbreak{}ω\textasciicircum{}\allowbreak{}(ε₁[n])\allowbreak{},\allowbreak{}当n趋向于无穷大时,\allowbreak{}这个序列就收敛到ε₁。这种基本列的优势是,\allowbreak{}每一步都能通过映射计算,\allowbreak{}适合你们日常修炼时,\allowbreak{}逐步积累力量、逼近ε₁。”

牢大挠了挠头,瓮声瓮气地问道：``赵老师,那第二种基本列呢？和第一种有什么不一样,实战中能用吗？还有您说的第二种比第一种强一点,具体强在哪里？''

``第二种基本列,是ε₀指数塔型,迭代基础不再是ω的幂次,而是ε₀自身的指数迭代,核心是'以ε₀为基础,构造ε₀的指数塔序列,取其极限',它和第一种基本列一样,最终收敛到ε₁,但层级略高、力量略强。''赵无极缓缓说道,``具体来说,这种基本列是直接构造ε₀的指数塔序列,作为ε₁的基本列,即：

ε₁[n]=\allowbreak{}ε₀\textasciicircum{}\allowbreak{}ε₀\textasciicircum{}\allowbreak{}…\textasciicircum{}\allowbreak{}ε₀（共n层ε₀指数迭代,\allowbreak{}也就是ε₀的n层指数塔）；

简单说,\allowbreak{}n=\allowbreak{}1时,\allowbreak{}就是ε₀（1层指数塔）；n=\allowbreak{}2时,\allowbreak{}就是ε₀\textasciicircum{}\allowbreak{}ε₀（2层指数塔）；n=\allowbreak{}3时,\allowbreak{}就是ε₀\textasciicircum{}\allowbreak{}ε₀\textasciicircum{}\allowbreak{}ε₀（3层指数塔）；n=\allowbreak{}4时,\allowbreak{}就是ε₀\textasciicircum{}\allowbreak{}ε₀\textasciicircum{}\allowbreak{}ε₀\textasciicircum{}\allowbreak{}ε₀,\allowbreak{}以此类推,\allowbreak{}n越大,\allowbreak{}指数塔层数越多,\allowbreak{}越接近ε₁。用严格的序数表示,\allowbreak{}就是ε₁=\allowbreak{}sup\{ε₀,\allowbreak{}ε₀\textasciicircum{}\allowbreak{}ε₀,\allowbreak{}ε₀\textasciicircum{}\allowbreak{}ε₀\textasciicircum{}\allowbreak{}ε₀,\allowbreak{}ε₀\textasciicircum{}\allowbreak{}ε₀\textasciicircum{}\allowbreak{}ε₀\textasciicircum{}\allowbreak{}ε₀,\allowbreak{}…\}。”

赵无极顿了顿,补充强弱对比的关键细节：``至于为什么第二种比第一种略强,\allowbreak{}核心在于迭代基础的层级差异——第一种基本列的迭代基础是’ω的幂次’,\allowbreak{}哪怕是ω\textasciicircum{}\allowbreak{}(ε₀+\allowbreak{}1)\allowbreak{},\allowbreak{}本质还是以ω为底层；而第二种基本列的迭代基础是’ε₀自身’,\allowbreak{}ε₀本身就比ω高阶,\allowbreak{}它的指数迭代层级自然更高。比如,\allowbreak{}同样是n=\allowbreak{}3的层级,\allowbreak{}第一种是ω\textasciicircum{}\allowbreak{}(ω\textasciicircum{}\allowbreak{}(ω\textasciicircum{}\allowbreak{}(ε₀+\allowbreak{}1)\allowbreak{})\allowbreak{})\allowbreak{},\allowbreak{}第二种是ε₀\textasciicircum{}\allowbreak{}ε₀\textasciicircum{}\allowbreak{}ε₀,\allowbreak{}后者的序数层级的要略高于前者,\allowbreak{}反映在实力上,\allowbreak{}就是第二种基本列能调动的力量,\allowbreak{}比第一种更凝练、更强横。''

奥斯卡问：为什么这两种基本列的上确界都是ε1,

赵无极：首先,\allowbreak{}ω\textasciicircum{}\allowbreak{}εα=\allowbreak{}εα,\allowbreak{}α是任意序数,\allowbreak{}ω\textasciicircum{}\allowbreak{}(εα+\allowbreak{}1)\allowbreak{}=\allowbreak{}εα*\allowbreak{}ω,\allowbreak{}ω\textasciicircum{}\allowbreak{}ω\textasciicircum{}\allowbreak{}(εα+\allowbreak{}1）=\allowbreak{}ω\textasciicircum{}\allowbreak{}(εα*\allowbreak{}ω)\allowbreak{}=\allowbreak{}(ω\textasciicircum{}\allowbreak{}εα)\allowbreak{}\textasciicircum{}\allowbreak{}ω=\allowbreak{}εα\textasciicircum{}\allowbreak{}ω,\allowbreak{}ω\textasciicircum{}\allowbreak{}εα\textasciicircum{}\allowbreak{}ω=\allowbreak{}ω\textasciicircum{}\allowbreak{}(εα·εα\textasciicircum{}\allowbreak{}ω)\allowbreak{}=\allowbreak{}(ω\textasciicircum{}\allowbreak{}εα)\allowbreak{}\textasciicircum{}\allowbreak{}εα\textasciicircum{}\allowbreak{}ω=\allowbreak{}εα\textasciicircum{}\allowbreak{}εα\textasciicircum{}\allowbreak{}ω注意到εα\textasciicircum{}\allowbreak{}ω=\allowbreak{}εα\textasciicircum{}\allowbreak{}(1+\allowbreak{}ω)\allowbreak{}=\allowbreak{}εα·εα\textasciicircum{}\allowbreak{}ω

奥斯卡说：居然是序数吸收率,左乘和右乘不一样的,一般右边乘法才会让序数变大,这么说前者n层,大于后者n-2层又小于后者n-1层,是吗？

赵无极L''那实战和修炼中,我们该怎么选择？``小舞立刻追问道,指尖的燃烧弹泛起淡淡的ε₀(1)波动,似是在推演两种基本列的力量差异。

``很简单,根据场景选择。''赵无极笑着解答,``第一种基本列（ω幂次迭代型）,\allowbreak{}每一步都能通过映射α↦ω\textasciicircum{}\allowbreak{}α计算,\allowbreak{}迭代难度低、稳定性强,\allowbreak{}适合你们日常修炼时,\allowbreak{}逐步积累力量、夯实基础,\allowbreak{}稳步逼近ε₁；第二种基本列（ε₀指数塔型）,\allowbreak{}虽然力量更强,\allowbreak{}但迭代难度更高,\allowbreak{}需要对ε₀的性质有更透彻的领悟,\allowbreak{}适合实战中遇到强敌（比如接近ε₁层级的对手）时,\allowbreak{}快速调动更强的力量,\allowbreak{}实现破局。''

奥斯卡咬着香肠,认真总结：``我明白了,两种基本列最终都收敛到ε₁,但第一种是ω幂次迭代,难度低、适合修炼；第二种是ε₀指数塔迭代,力量略强、适合实战,核心区别就是迭代基础和力量强度,对吗？''

``非常准确！''赵无极赞许点头,``掌握了ε₁的构造和基本列,你们突破ε₀后,就能稳步逼近ε₁。而说到ε₁,就必须讲清楚从ε₀到ζ₀的完整构造------ζ₀是比ε₁、ε₂甚至ε\_ω都更高阶的递归序数,也是你们突破ε序列后的重要目标,更是你们对付千年人面魔蛛（魂核φ(ω,0)）的关键基础,我结合ε序列的构造,详细给你们讲从ε₀到ζ₀的完整路径。''

唐三眼神发亮,立刻说道：``请赵老师讲解,我们一定认真听,这对我们完成任务太重要了。''

赵无极抬手,在空中浮现出ε序列和ζ₀的构造图谱,缓缓讲解：``从ε₀到ζ₀,核心是'不断跳出不动点、迭代映射、取极限'的循环,一共分为三步,每一步都离不开我们刚才讲的ε₀、ε₁的构造逻辑,你们仔细听：''

“第一步,\allowbreak{}构造完整的εₙ序列（n为自然数）。我们已经知道,\allowbreak{}ε₀是映射α↦ω\textasciicircum{}\allowbreak{}α的第一个不动点,\allowbreak{}ε₁是第二个不动点,\allowbreak{}按照同样的逻辑,\allowbreak{}我们可以构造出ε₂、ε₃、…、εₙ,\allowbreak{}具体方法完全一致：

1.构造ε₂：先跳出ε₁这个不动点,\allowbreak{}取ε₁+\allowbreak{}1；再对ε₁+\allowbreak{}1反复应用映射α↦ω\textasciicircum{}\allowbreak{}α,\allowbreak{}得到序列\{ε₁,\allowbreak{}ω\textasciicircum{}\allowbreak{}(ε₁+\allowbreak{}1)\allowbreak{},\allowbreak{}ω\textasciicircum{}\allowbreak{}(ω\textasciicircum{}\allowbreak{}(ε₁+\allowbreak{}1)\allowbreak{})\allowbreak{},\allowbreak{}…\}；取这个序列的上确界,\allowbreak{}就是ε₂,\allowbreak{}它是映射α↦ω\textasciicircum{}\allowbreak{}α的第三个不动点,\allowbreak{}记为ε₂=\allowbreak{}3rdα↦ω\textasciicircum{}\allowbreak{}αfp.；

2.构造ε₃：跳出ε₂,取ε₂+1,反复映射、取极限,得到ε₃（第四个不动点）；

3.以此类推,不断重复'跳出不动点→反复映射→取极限'的过程,我们就能得到所有的εₙ（n为自然数）,即ε₀、ε₁、ε₂、ε₃、\ldots、εₙ、\ldots。''

``第二步,构造ε\_ω（下标为ω的ε序数）。当我们得到了所有自然数下标的εₙ后,这些εₙ构成了一个单调递增的序数序列,我们取这个序列的上确界,就得到了ε\_ω,即ε\_ω=sup\{ε₀,ε₁,ε₂,ε₃,\ldots\}。''赵无极补充道,``这里要注意,\allowbreak{}ε\_ω不是某个具体的εₙ,\allowbreak{}而是所有εₙ的极限,\allowbreak{}它也是映射α↦ω\textasciicircum{}\allowbreak{}α的不动点,\allowbreak{}是第ω个不动点,\allowbreak{}记为ε\_ω=\allowbreak{}ωthα↦ω\textasciicircum{}\allowbreak{}αfp.。''

小舞追问道：``赵老师,那ε\_ω之后,还能继续构造更高的ε序数吗？这和ζ₀又有什么关系？''

``能,而且可以无限构造,而这所有构造的ε序数,进行下标的迭代就是ζ₀。''赵无极语气郑重,继续讲解第三步,``第三步,构造ζ₀------ζ₀是映射α↦ε\_α的第一个不动点,他放慢语速,拆解细节：''我们可以这样理解,当我们不断迭代下标构造更高的ε序数,就是ζ₀。用严格的序数表示就是：

ζ₀=sup\{0,ε₀,ε\_ε₀,ε\_ε\_ε₀,\ldots\}。''

``而且ζ₀有一个关键性质,\allowbreak{}它不仅是映射α↦ω\textasciicircum{}\allowbreak{}α的不动点,\allowbreak{}更是映射α↦ε\_α的第一个不动点——也就是说,\allowbreak{}ε\_ζ₀=\allowbreak{}sup\{ε₀,\allowbreak{}ε₁,\allowbreak{}ε₂,\allowbreak{}…\}=\allowbreak{}ζ₀,\allowbreak{}这也是ζ₀区别于所有ε序数的核心特征,\allowbreak{}我们把ζ₀称为Cantor序数（CO）。''

蔡须坤指尖的金色灵气轻轻缠绕,补充道：``赵老师,我是不是可以这样总结：从ε₀到ζ₀,就是'构造εₙ→构造ε\_ω→构造更高阶ε序数→迭代下标取上确界'的过程,核心始终是'跳出不动点、反复映射、取极限',对吗？''

``完全正确,领悟得非常透彻！''赵无极赞许地点了点头,``简单来说,ε₀是ε序列的起点,ε₁、ε₂\ldots εₙ是有限下标的ε序数,ε\_ω及以上是无限下标的ε序数,而ζ₀,就是所有这些ε序数的'天花板',是ε序列的极限。你们现在的目标,是先突破ε₀,再逐步逼近ε₁,等你们能稳定掌控ε₁的力量,再朝着ε\_ω、ζ₀迈进------而掌握这些,你们才能真正理解φ(ω,0)的层级,才能顺利猎杀千年人面魔蛛。''

唐三默默推演着从ε₀到ζ₀的构造逻辑,\allowbreak{}丹田内的ω\textasciicircum{}\allowbreak{}(ω\textasciicircum{}\allowbreak{}(ω\textasciicircum{}\allowbreak{}ω)\allowbreak{})\allowbreak{}序数纹路微微震颤,\allowbreak{}似是有所领悟,\allowbreak{}他轻声说道：``原来如此,ε₀到ζ₀,就是不断突破不动点、不断迭代升级的过程,我们之前对高阶序数的理解,还是太浅显了。''

``什么是算术序数？''奥斯卡咬着一根香肠,含糊地问道,``是不是就是我们能通过加减乘除、次方迭代构造出来的数？''

``没错。''赵无极笑着点头,\allowbreak{}指尖一动,\allowbreak{}空中又浮现出一串序数：ω、ω+\allowbreak{}1、ω×2、ω²、ω\textasciicircum{}\allowbreak{}ω、ω\textasciicircum{}\allowbreak{}ω\textasciicircum{}\allowbreak{}ω,\allowbreak{}``这些都是算术序数,\allowbreak{}它们的构造都离不开ω的基础运算,\allowbreak{}哪怕是ω\textasciicircum{}\allowbreak{}ω\textasciicircum{}\allowbreak{}ω,\allowbreak{}本质上还是ω的幂次迭代。但ε₀不一样,\allowbreak{}它的定义是’最小的满足ε=\allowbreak{}ω\textasciicircum{}\allowbreak{}ε的序数’,\allowbreak{}简单说,\allowbreak{}就是ω的幂次,\allowbreak{}迭代到ω次,\allowbreak{}最终收敛的那个序数——ω\textasciicircum{}\allowbreak{}(ω\textasciicircum{}\allowbreak{}(ω\textasciicircum{}\allowbreak{}(…)\allowbreak{})\allowbreak{})\allowbreak{},\allowbreak{}这个的结果,\allowbreak{}就是ε₀也就是φ(1,\allowbreak{}0)\allowbreak{}。''

小舞皱了皱眉,指尖的燃烧弹微微转动：``我还是有点懵,既然是无穷迭代,那我们怎么判断它的大小？还有您说的基本列,又是什么？''

``基本列是理解高阶序数的关键,也是你们突破境界、操控序数力量的核心方法。''赵无极的语气变得严肃,``所谓基本列,就是给一个极限序数,构造一个单调递增、最终收敛到这个极限序数的序数序列。ε₀作为极限序数,有两种最基础、最常用的基本列,我结合你们之前掌握的知识,一步步讲给你们听。''

继续说道：``基本列,是从'幂次迭代'的角度构造的,核心是逐步提升ω的幂次层级,最终逼近ε₀。具体来说,就是：

ε₀{[}0{]}=ω（最基础的起点）

ε₀{[}1{]}=ω\^{}ω（第一次迭代,ω的ω次方）

ε₀{[}2{]}=ω\^{}(ω\^{}ω)（第二次迭代,ω的ω\^{}ω次方）

ε₀{[}3{]}=ω\^{}(ω\^{}(ω\^{}ω))（第三次迭代）

以此类推,ε₀{[}n{]}=ω的n层幂次迭代,当n到达ω时,这个序列就收敛到ε₀。''

蔡须坤指尖的金色灵气轻轻流转,补充道：``赵老师,我是不是可以这样理解：第一种基本列适合长期修炼,用来稳步提升序数层级,夯实基础；第二种基本列适合实战,用来快速调动力量,灵活应对敌人？''

``完全正确！''赵无极赞许地点了点头,``蔡须坤你悟性很高,正是这个道理。第一种基本列注重'积累',你们平时修炼FGH功法,其实就是在沿着第一种基本列的路径,逐步逼近ε₀、突破ε₀；而第二种基本列注重'调动',实战中,你们可以通过拆分序数层级,快速凝聚力量,避免蛮力硬拼。''

唐三沉吟片刻,问道：``赵老师,那ε₀之后,就是ζ₀了吧？它和ε₀有什么关系？性质上又有什么不同？''

``问得非常关键,ζ₀就是ε₀之后的下一个重要高阶递归序数,也是你们突破ε₀之后的主要修炼目标。''赵无极抬手,将空中的ε₀符号替换成ζ₀,``ζ₀的定义和ε₀类似,\allowbreak{}但层级更高——ε₀是’最小的ε=\allowbreak{}ω\textasciicircum{}\allowbreak{}ε’,\allowbreak{}而ζ₀是’最小的ζ=\allowbreak{}ε\_ζ’,\allowbreak{}简单说,\allowbreak{}就是把ε₀的构造方法再迭代一次。''

他放慢语速,详细讲解：``我们可以这样理解,\allowbreak{}ε₀是ω的幂次的ω次迭代,\allowbreak{}而ζ₀是ε下标序列的ω迭代。首先,\allowbreak{}我们定义ε₀=\allowbreak{}最小的ε=\allowbreak{}ω\textasciicircum{}\allowbreak{}ε,\allowbreak{}ε₁=\allowbreak{}最小的ε>ε₀且ε=\allowbreak{}ω\textasciicircum{}\allowbreak{}ε,\allowbreak{}ε₂=\allowbreak{}最小的ε序数>ε₁且εα=\allowbreak{}ω\textasciicircum{}\allowbreak{}εα,\allowbreak{}以此类推,\allowbreak{}εₙ就是第n个满足ε=\allowbreak{}ω\textasciicircum{}\allowbreak{}ε的序数。而ζ₀,\allowbreak{}就是这个εα序列的极限,\allowbreak{}也就是ε₀,\allowbreak{}ε\_ε₀,\allowbreak{}ε\_ε\_ε₀,\allowbreak{}…这个序列收敛到的那个序数,\allowbreak{}它比任何一个εα,\allowbreak{}α＜ζ₀都要大。''

``那ζ₀的性质,和ε₀有什么区别呢？''小舞追问道,眼底的疑惑渐渐消散,多了几分专注。

``两者的核心区别,在于'迭代层级'和'力量强度'。''赵无极语气郑重,``第一,ε₀是'一阶无穷迭代',核心是ω的幂次无穷迭代；而ζ₀是'二阶无穷迭代',核心是ε序列的无穷迭代,层级比ε₀高一个维度。第二,在实力对应上,ε₀是中级学院入门的序数,而ζ₀是中级学院学习的第一个高阶的序数,达到ζ₀,就能稳压百年曼陀罗蛇,甚至能与500年曼陀罗蛇不相上下。''

他顿了顿,继续补充：``还有一个关键性质------ε₀和ζ₀,都是'递归序数',也就是说,它们都能通过有限的规则,构造出对应的基本列,这也是你们能够修炼、掌控它们的前提。如果是不可递归的序数,哪怕知道它的定义,也无法通过修炼去突破、去掌控。另外,ε₀的基本列是基于ω的幂次,而ζ₀的基本列,是基于ε序列的迭代,比如ζ₀(n)=εₙ,当n到达ω时,就收敛到ζ₀。''

奥斯卡听得眼睛发亮,连忙问道：``赵老师,\allowbreak{}那我们现在的实力,\allowbreak{}距离ε₀、ζ₀还有多远？唐三现在是ω\textasciicircum{}\allowbreak{}ω\textasciicircum{}\allowbreak{}ω,\allowbreak{}小舞是ω\textasciicircum{}\allowbreak{}((ω³)\allowbreak{}*\allowbreak{}2)\allowbreak{},\allowbreak{}我和牢大、蔡须坤,\allowbreak{}又该怎么朝着这些序数突破？''

赵无极扫过五人,逐一分析：``唐三,\allowbreak{}你刚才听我讲解ε₀的第一种基本列时,\allowbreak{}眼神已有领悟,\allowbreak{}你的ω\textasciicircum{}\allowbreak{}ω\textasciicircum{}\allowbreak{}ω本就接近ε₀第一种基本列的第三项,\allowbreak{}再加上你对幂次迭代的领悟加深,\allowbreak{}此刻应该已经突破到ω\textasciicircum{}\allowbreak{}(ω\textasciicircum{}\allowbreak{}(ω\textasciicircum{}\allowbreak{}ω)\allowbreak{})\allowbreak{},\allowbreak{}也就是ε₀(3)\allowbreak{},\allowbreak{}距离ε₀的门槛只有一步之遥,\allowbreak{}只需再领悟无穷迭代的精髓,\allowbreak{}便能正式突破到ε₀；小舞,\allowbreak{}你对序数嵌套的理解远超常人,\allowbreak{}刚才我讲解ζ₀的ε序列迭代时,\allowbreak{}你已悄然将自身ω\textasciicircum{}\allowbreak{}((ω³)\allowbreak{}*\allowbreak{}2)\allowbreak{}的嵌套幂次,\allowbreak{}转化为初步的ε序列雏形,\allowbreak{}此刻序数层级已提升到ε₀[2],\allowbreak{}也就是ω\textasciicircum{}\allowbreak{}ω\textasciicircum{}\allowbreak{}ω之上、ε₀之下的高阶层级,\allowbreak{}比之前稳固了数倍,\allowbreak{}只要继续打磨,\allowbreak{}也能快速逼近ε₀。''

``至于你们三个,''他看向牢大、蔡须坤和奥斯卡,``牢大你现在是ε\_ω,已经突破了ε₀,ε\_ω是ε序列中比ε₀高的序数,距离ζ₀还有一段距离,需要继续迭代ε序列,逐步逼近ζ₀；蔡须坤你的序数绑定术,核心是联结他人的序数,你自身的序数层级虽然不高,但只要能熟练联结牢大的ε\_ω,甚至未来联结ζ₀级别的力量,你的作用就无可替代；奥斯卡你的香肠挂载,现在能提升高阶序数的稳定性,接下来要尝试挂载ε₀、ζ₀级别的序数波动,让香肠的增益效果更上一层楼。''

牢大握紧拳头,肘间的ε\_ω序数涟漪微微暴涨：``赵老师,那我们该怎么演练？是不是要模拟和千年人面魔蛛的对战,试着运用这些序数的基本列？''

``没错。''赵无极点了点头,身形一闪,来到空地中央,周身浮现出φ(4,0)的序数波动,与千年人面魔蛛的防御层级一致,``我来模拟人面魔蛛的防御,你们五人配合,运用我刚才讲的ε₀、ζ₀的性质,还有序数嵌套、借力的方法,试着突破我的防御。记住,人面魔蛛的魂核是φ(ω,0),比ζ₀还要高,你们不仅要突破防御,还要学会拆解高阶序数的波动,这样才能提取它的魂核。''

唐三眼神一凝,\allowbreak{}指尖轻抬,\allowbreak{}数根蓝银草从掌心窜出,\allowbreak{}藤条上浮现出ω\textasciicircum{}\allowbreak{}(ω\textasciicircum{}\allowbreak{}(ω\textasciicircum{}\allowbreak{}ω)\allowbreak{})\allowbreak{}的序数纹路——这是他领悟ε₀基本列后,\allowbreak{}刚刚突破的新层级,\allowbreak{}比之前的ω\textasciicircum{}\allowbreak{}ω\textasciicircum{}\allowbreak{}ω强横了不少。他心念一动,\allowbreak{}默默推演ε₀的无穷迭代精髓,\allowbreak{}试图将幂次迭代推向更高层次：``小舞,\allowbreak{}你用瞬移锁定赵老师的防御节点,\allowbreak{}标记出φ(4,\allowbreak{}0)\allowbreak{}的薄弱处；奥斯卡,\allowbreak{}给我和牢大挂载香肠增益；蔡须坤,\allowbreak{}联结我和牢大的序数,\allowbreak{}将我的ω\textasciicircum{}\allowbreak{}(ω\textasciicircum{}\allowbreak{}(ω\textasciicircum{}\allowbreak{}ω)\allowbreak{})\allowbreak{}和他的ε\_ω嵌套起来！''

``明白！''众人齐声应道,\allowbreak{}瞬间行动起来。小舞的身影化作残影,\allowbreak{}瞬移间在赵无极周身标记出多个节点,\allowbreak{}每个节点都被她的ε₀(1)\allowbreak{}序数点亮——突破后的序数波动更加凝练,\allowbreak{}标记的精准度也大幅提升；奥斯卡掏出三根香肠,\allowbreak{}分别贴在唐三、牢大和蔡须坤身上,\allowbreak{}银光闪烁间,\allowbreak{}三人的序数波动瞬间稳定了几分；蔡须坤的金色纽带凭空出现,\allowbreak{}一头联结唐三的蓝银草,\allowbreak{}一头联结牢大的肘部,\allowbreak{}将ω\textasciicircum{}\allowbreak{}(ω\textasciicircum{}\allowbreak{}(ω\textasciicircum{}\allowbreak{}ω)\allowbreak{})\allowbreak{}和ε\_ω的力量紧紧绑定,\allowbreak{}层层嵌套；牢大则凝聚全身力量,\allowbreak{}肘间的ε\_ω序数涟漪愈发厚重,\allowbreak{}随时准备发起冲击。

赵无极看着五人默契的配合,嘴角扬起一抹浅笑,周身的φ(4,0)防御纹路缓缓亮起：``很好,就是这样！记住,运用ε₀的基本列,将力量逐步递增,不要急于求成,借力打力,才能突破高阶防御！动手吧！''

``动手！''唐三一声令下,\allowbreak{}嵌套着ε\_ω力量的蓝银草猛地缠向赵无极,\allowbreak{}藤条上的ω\textasciicircum{}\allowbreak{}(ω\textasciicircum{}\allowbreak{}(ω\textasciicircum{}\allowbreak{}ω)\allowbreak{})\allowbreak{}序数纹路不断闪烁,\allowbreak{}隐隐有向ε₀逼近的趋势；牢大则借着蔡须坤的绑定术,\allowbreak{}身形一闪,\allowbreak{}一记肘击狠狠砸向被小舞标记的薄弱节点；小舞不断瞬移,\allowbreak{}ε₀(2)\allowbreak{}的序数波动精准锁定防御纹路的缝隙,\allowbreak{}进一步瓦解防御；奥斯卡则持续补充香肠,\allowbreak{}维持众人的序数稳定性；蔡须坤则专注于操控金色纽带,\allowbreak{}确保力量传递不中断。

``砰！''一声巨响,牢大的肘击砸在防御纹路上,结合唐三蓝银草的缠绕和小舞的标记瓦解,赵无极周身的φ(4,0)纹路瞬间泛起一阵涟漪,出现了一道细微的缝隙。

``不错！''赵无极赞许地说道,``已经能打破我的防御雏形了,但还不够！千年人面魔蛛的防御比我模拟的还要坚固,而且它会主动攻击,你们还要学会在攻击的同时,防御它的φ(ω,0)级序数冲击。接下来,我会加大防御强度,同时模拟它的攻击,你们继续演练！''

五人相视一眼,眼中都燃起了斗志。唐三调整呼吸,借着演练的契机,继续推演ε₀的无穷迭代,蓝银草上的序数纹路愈发接近ε₀；小舞则熟练运用ε₀(2)的力量,标记防御节点的速度更快、更精准,同时也在感悟ε序列的迭代精髓,朝着更高层级迈进；牢大、蔡须坤和奥斯卡也各司其职,配合愈发默契。

阳光透过树叶的缝隙,洒在五人的身影上,序数波动相互交织,形成一道耀眼的光环。他们知道,这场战前演练,不仅是对配合的考验,更是对ε₀、ζ₀等高阶序数的理解与运用的考验。只有吃透这些序数的性质,熟练掌握基本列的运用,才能在猎杀千年人面魔蛛的任务中取胜,才能在中等学院的修炼之路上,迈出更坚实的一步。

而赵无极看着眼前全力以赴的五人,心中暗暗点头：这五个孩子,天赋异禀,又能默契配合,只要稍加引导,吃透这些高阶序数的威力,未来必定能成为魂师大陆上的强者,甚至能触及更高阶的φ函数序数,书写属于他们的传奇。 以下是修改后的完整章节，在赵无极到来之后加入了牢大与唐三的争吵情节：

\begin{center}\rule{0.5\linewidth}{0.5pt}\end{center}

\chapter{第八章：不乐的狂乱}\label{ux7b2cux516bux7ae0ux4e0dux4e50ux7684ux72c2ux4e71}

山林边缘的演练刚结束，赵无极拍了拍手上的灰尘，目光扫过五个气喘吁吁却眼神明亮的少年，满意地点了点头：``今天的演练到此为止，你们对 \allowbreak{}\NovelInline{I00078}{\varepsilon_0}\allowbreak{} 基本列的运用已经初窥门径，接下来几天自行巩固。现在，收拾东西，我们出发前往星斗大森林外围的补给镇。''

五人迅速收拾好行装，沿着林间小道向镇子方向行进。唐三走在最前面，蓝银草的感知如无形的蛛网铺展开来，\allowbreak{}\NovelInline{I00079}{\omega^{\omega^{\omega^\omega}}}\allowbreak{} 的序数波动被他刻意压制到极低，只在感知范围内维持着精密的探测网络。小舞蹦蹦跳跳地跟在旁边，指尖偶尔闪过一丝 \allowbreak{}\NovelInline{I00080}{\varepsilon_0[3]}\allowbreak{} 的微光，像是在把玩一颗看不见的弹珠。牢大走在队伍中间，肘间的 \allowbreak{}\NovelInline{I00081}{\varepsilon_\omega}\allowbreak{} 涟漪被他收敛得干干净净，但每一步踏下去，脚下的落叶都会无声地陷落半分------那是高阶序数自然溢出的重力效应。蔡须坤和奥斯卡一左一右，金色纽带若有若无地联结着彼此，香肠的增益波动也在刻意压低，像一层薄薄的保护膜笼罩着整个小队。

赵无极没有同行。他在出发前便交代过，补给镇距此不过半日路程，沿途又是他走过无数遍的熟路，正好借这段路程考验五人的自主警戒与协同应变能力。``我会在镇子上等你们，''他拍了拍唐三的肩膀，``路上小心，有事用蓝银草的感知频段联系我。''说完便纵身掠入林梢，眨眼间消失在了密林深处。

然而，就在距离补给镇还有不到三里的一处岔路口，唐三的脚步突然停住了。

他的眉头猛地皱起，脸色在瞬间变得凝重。蓝银草的感知捕捉到了一组极其狂暴、且带着明显压迫感的序数波动。那波动的层级\ldots\ldots 让他后背的汗毛根根竖起。

不是 \allowbreak{}\NovelInline{I00082}{\omega}\allowbreak{} 幂次的叠加，不是 \allowbreak{}\NovelInline{I00083}{\omega^\omega}\allowbreak{} 的迭代，甚至不是 \allowbreak{}\NovelInline{I00084}{\varepsilon_0}\allowbreak{} 的门槛------

是 \allowbreak{}\NovelInline{I00085}{\varepsilon_1}\allowbreak{}。

唐三的瞳孔骤然收缩。\allowbreak{}\NovelInline{I00086}{\varepsilon_1}\allowbreak{}，\allowbreak{}\NovelInline{I00087}{\varepsilon}\allowbreak{} 序列中的第二个不动点，满足 \allowbreak{}\NovelInline{I00088}{\omega^{\varepsilon_1} = \varepsilon_1}\allowbreak{} 的恐怖序数。它比唐三自己苦修至今的 \allowbreak{}\NovelInline{I00089}{\varepsilon_0}\allowbreak{} 整整高出了一个完整的 \allowbreak{}\NovelInline{I00090}{\varepsilon}\allowbreak{} 层级，比小舞刚刚触及的 \allowbreak{}\NovelInline{I00091}{\varepsilon_0[3]}\allowbreak{} 基本列更是高了不知多少个维度。在这个层级面前，他们五人中除了牢大，没有任何一人具备正面对抗的资格。

更可怕的是，这股 \allowbreak{}\NovelInline{I00092}{\varepsilon_1}\allowbreak{} 的波动运用方式极其粗暴------不是精密的迭代构造，而是纯粹的蛮力碾压，像是一头被强行催肥的巨兽，浑身上下每一寸肌肤都在向外喷涌着未经驯化的序数洪流。边缘处不断有未收敛的序数碎片向外逸散，形成一种令人窒息的紊乱场，连周围的空气都在微微扭曲。

``前方三百米，岔路口左侧。''唐三低声开口，声音平静却带着一丝罕见的凝重，``有人在用序数力量压制别人。波动层级\ldots\ldots{}\allowbreak{}\NovelInline{I00093}{\varepsilon_1}\allowbreak{}。''

``什么？！''小舞的脸色瞬间变了，指尖那缕 \allowbreak{}\NovelInline{I00094}{\varepsilon_0[3]}\allowbreak{} 的微光不自觉地黯淡了几分。她是五人中对 \allowbreak{}\NovelInline{I00095}{\varepsilon_0}\allowbreak{} 基本列领悟最深的人之一，正因如此，她比任何人都清楚 \allowbreak{}\NovelInline{I00096}{\varepsilon_1}\allowbreak{} 意味着什么------那是一座她穷尽目前所有修为都无法仰望的高山。

牢大的表情也沉了下来。肘间的 \allowbreak{}\NovelInline{I00097}{\varepsilon_\omega}\allowbreak{} 涟漪微微波动------\allowbreak{}\NovelInline{I00098}{\varepsilon_\omega}\allowbreak{} 确实高于 \allowbreak{}\NovelInline{I00099}{\varepsilon_1}\allowbreak{}，在层级上能够形成压制。但他同样清楚，\allowbreak{}\NovelInline{I00100}{\varepsilon_\omega}\allowbreak{} 对 \allowbreak{}\NovelInline{I00101}{\varepsilon_1}\allowbreak{} 的优势，远没有 \allowbreak{}\NovelInline{I00102}{\varepsilon_\omega}\allowbreak{} 对 \allowbreak{}\NovelInline{I00103}{\omega}\allowbreak{} 幂次那种碾压式的差距。\allowbreak{}\NovelInline{I00104}{\varepsilon_1}\allowbreak{} 毕竟是 \allowbreak{}\NovelInline{I00105}{\varepsilon}\allowbreak{} 序列中的第二个不动点，其蛮力输出的上限极其恐怖，如果对方不计后果地疯狂爆发，即便是 \allowbreak{}\NovelInline{I00106}{\varepsilon_\omega}\allowbreak{} 也不能掉以轻心。

唐三迅速用蓝银草的感知频段向赵无极发送了简短的坐标与敌情信息，但信号发出后，反馈回来的只有一片模糊的干扰波纹------那片区域的序数场太过紊乱，\allowbreak{}\NovelInline{I00107}{\varepsilon_1}\allowbreak{} 的蛮力余波像一堵厚重的墙，将一切精密信号都碾成了碎片。赵无极显然还未进入感知范围。

``先靠近观察，不要贸然出手。''唐三压低身形，以战术手势指挥小队呈扇形向目标位置逼近。

穿过一片低矮的灌木丛，眼前的景象让所有人的目光都冷了下来。

岔路口的空地上，一个身材精壮、面容阴鸷的中年男人正单手掐着一个年轻女子的脖子，将她整个人提离了地面。那女子穿着朴素的布裙，脸色已经憋得通红，双手拼命抓着男人的手腕，却根本无法撼动分毫。男人的另一只手上，一团浑浊的 \allowbreak{}\NovelInline{I00108}{\varepsilon_1}\allowbreak{} 序数光芒正疯狂闪烁，像是一条粗糙却极其厚重的锁链，死死缠绕在女子的脖颈和四肢上，将她体内的序数波动完全封锁------那是一种极其低劣却极其霸道的''序数禁锢''手法，没有任何基本列的精密构造，只是将 \allowbreak{}\NovelInline{I00109}{\varepsilon_1}\allowbreak{} 的蛮力直接灌注到对方体内，用纯粹的力量压碎对方的抵抗。

``臭娘们，给脸不要脸！''男人------正是这一带臭名昭著的恶徒不乐------狞笑着凑近女子的脸，声音里满是下流的得意，``老子虽然没能把 \allowbreak{}\NovelInline{I00110}{\varepsilon_1}\allowbreak{} 打磨到精密构造的程度，但论蛮力，这一带谁扛得住？你那个废物未婚夫已经被我一肘子废了序数根基，你现在就算叫破喉咙也没用！乖乖把身上的序数结晶交出来，再陪老子乐呵乐呵，说不定老子还能留你一条命！''

女子的眼中满是绝望，嘴唇已经发紫，却仍然咬着牙不肯发出求饶的声音。她的脖颈上，那团浑浊的序数锁链正在不断收紧，皮肤上已经渗出了细密的血珠。

唐三的目光沉了下来。他在极短的时间内完成了战术判断：对方修为 \allowbreak{}\NovelInline{I00111}{\varepsilon_1}\allowbreak{}，远超自己和小舞的 \allowbreak{}\NovelInline{I00112}{\varepsilon_0}\allowbreak{} 层级。如果正面硬碰，他和小舞的序数构造在 \allowbreak{}\NovelInline{I00113}{\varepsilon_1}\allowbreak{} 的蛮力面前，就像薄纸面对铁锤------不是技巧能弥补的差距，而是层级上的绝对碾压。唯一的优势在于牢大的 \allowbreak{}\NovelInline{I00114}{\varepsilon_\omega}\allowbreak{}，但 \allowbreak{}\NovelInline{I00115}{\varepsilon_\omega}\allowbreak{} 对 \allowbreak{}\NovelInline{I00116}{\varepsilon_1}\allowbreak{} 的压制并非碾压式的，如果对方用人质做盾牌、以自毁式爆发来拼命，牢大也不敢全力出手。

``牢大正面牵制，''唐三通过蓝银草的感知频段向四人传递指令，声音压得极低，``我和小舞从侧翼寻找机会解救人质，蔡须坤和奥斯卡远程策应。记住，不要和对方硬碰------他的层级比我们都高，只能用配合来弥补差距。''

五人同时动了。

牢大率先发难，他一步踏出，肘间的 \allowbreak{}\NovelInline{I00117}{\varepsilon_\omega}\allowbreak{} 涟漪不再收敛，如同实质般的波纹向四周扩散开来。他从正面逼近，每一步都在地面留下深陷的脚印，\allowbreak{}\NovelInline{I00118}{\varepsilon_\omega}\allowbreak{} 的序数波动如同一座无形的大山，朝着不乐碾压而去。

不乐猛地转头，脸上的狞笑僵了一瞬------\allowbreak{}\NovelInline{I00119}{\varepsilon_\omega}\allowbreak{} 的层级确实在他之上，他能清晰地感觉到自己体内的 \allowbreak{}\NovelInline{I00120}{\varepsilon_1}\allowbreak{} 蛮力在这股波动面前被压制得隐隐发颤。但他常年在刀口舔血的直觉告诉他：对方没有立刻出手，说明有顾忌。

他的目光扫过身后被掐住脖子的林婉，嘴角浮起一丝狰狞的笑意。

他将林婉猛地拽到身前，浑浊的 \allowbreak{}\NovelInline{I00121}{\varepsilon_1}\allowbreak{} 序数锁链骤然收紧，林婉发出一声痛苦的闷哼。``别过来！''不乐嘶吼道，声音里带着困兽般的疯狂，``你的 \allowbreak{}\NovelInline{I00122}{\varepsilon_\omega}\allowbreak{} 确实比我强，但你敢动手吗？老子一用力，这女人的喉咙就碎了！你那 \allowbreak{}\NovelInline{I00123}{\varepsilon_\omega}\allowbreak{} 的余波，比老子的蛮力还猛十倍------你出手的瞬间，她就已经死了！''

牢大的脚步硬生生停住了。

他的脸色铁青，额头上青筋暴起。不乐说得没错------\allowbreak{}\NovelInline{I00124}{\varepsilon_\omega}\allowbreak{} 的力量太过强大，如果他全力出手，哪怕只是余波，也足以在击碎不乐的同时，将林婉脆弱的身体一并碾成齑粉。他不敢赌。

就在这一瞬间的僵持中，唐三和小舞从侧翼动了。

唐三的蓝银草无声无息地从地面延伸出去，\allowbreak{}\NovelInline{I00125}{\omega^{\omega^{\omega^\omega}}}\allowbreak{} 的精密序数纹路隐藏在藤条之中，朝着不乐掐住林婉手腕的方向悄然逼近。小舞则从另一侧绕了过去，指尖亮起 \allowbreak{}\NovelInline{I00126}{\varepsilon_0[3]}\allowbreak{} 的柔和光芒，准备用基本列的精密构造一丝一缕地剥离不乐施加在林婉身上的序数禁锢。

两人的计划很明确：牢大正面牵制吸引注意力，他们从侧翼悄悄解除人质危机，等林婉脱离危险后，牢大再全力出手。

计划很完美。

但他们忽略了一件事------\allowbreak{}\NovelInline{I00127}{\varepsilon_1}\allowbreak{} 的蛮力虽然粗糙，却有一个极其恐怖的特性：它的感知范围远超精密构造。

不乐的 \allowbreak{}\NovelInline{I00128}{\varepsilon_1}\allowbreak{} 序数虽然没有任何精密构造可言，但那股蛮力像一滩四处蔓延的洪水，无孔不入地渗透着周围的一切。唐三的蓝银草刚刚延伸到距不乐三米的范围，就被那股浑浊的 \allowbreak{}\NovelInline{I00129}{\varepsilon_1}\allowbreak{} 洪流捕捉到了。

``想从侧面搞鬼？''不乐狞笑一声，眼中闪过一丝狠厉。

他没有选择防守，而是做出了一个所有人都没有预料到的举动------他主动将体内的 \allowbreak{}\NovelInline{I00130}{\varepsilon_1}\allowbreak{} 蛮力再度暴涨，那团浑浊的光芒瞬间膨胀到了极限，像一颗即将爆炸的序数炸弹，朝着四面八方疯狂扩散。

``都给我死！''

紊乱的 \allowbreak{}\NovelInline{I00131}{\varepsilon_1}\allowbreak{} 冲击波轰然炸开。

那是一种完全不计后果的自毁式爆发------不乐将自己体内近半的序数根基一次性倾泻而出，化作一道毁灭性的蛮力风暴。这股力量没有任何技巧、没有任何构造，只有纯粹的、蛮横的、铺天盖地的力量。

但就是这种蛮力，在绝对的层级优势面前，足以碾碎一切花哨的技巧。

唐三首当其冲。他的蓝银草刚刚构筑起 \allowbreak{}\NovelInline{I00132}{\varepsilon_0}\allowbreak{} 基本列的防御纹路，\allowbreak{}\NovelInline{I00133}{\varepsilon_1}\allowbreak{} 的冲击波便如海啸般扑面而来。\allowbreak{}\NovelInline{I00134}{\varepsilon_0}\allowbreak{} 对 \allowbreak{}\NovelInline{I00135}{\varepsilon_1}\allowbreak{}------整整差了一个 \allowbreak{}\NovelInline{I00136}{\varepsilon}\allowbreak{} 层级。那精心编织的藤盾在接触蛮力的瞬间便轰然碎裂，就像一道纸糊的堤坝面对山洪，连一息都没能撑住。

唐三被震退数步，双臂发麻，胸口一阵气血翻涌，嘴角溢出一缕鲜血。

小舞的处境更加危险。她正在全力运转 \allowbreak{}\NovelInline{I00137}{\varepsilon_0[3]}\allowbreak{} 的精密构造去剥离禁锢，根本没有防备这种自毁式的蛮力爆发。\allowbreak{}\NovelInline{I00138}{\varepsilon_1}\allowbreak{} 的冲击波撞上她的瞬间，\allowbreak{}\NovelInline{I00139}{\varepsilon_0[3]}\allowbreak{} 的构造被蛮力硬生生冲散，她体内的序数根基剧烈震荡，像是被人用铁锤砸了一下。她的身体像断了线的风筝一样被掀飞出去，重重地摔在五米外的灌木丛中，嘴角溢出一缕鲜血，脸色瞬间变得惨白。

``小舞！''唐三的瞳孔猛地收缩，蓝银草本能地朝小舞的方向延伸过去。

但这一瞬间的分神，让他的感知网络出现了致命的缺口。

不乐抓住了这个转瞬即逝的机会。他猛地松开林婉------不是因为善心，而是因为林婉已经被蛮力震得昏死过去，对他而言已经失去了人质的价值。他将全部力量集中在右拳上，那团浑浊的 \allowbreak{}\NovelInline{I00140}{\varepsilon_1}\allowbreak{} 光芒凝聚成一道粗糙却极其厚重的光刃，朝着唐三的方向狠狠劈去！

唐三仓促间以蓝银草构筑第二道防御，但体内序数运转因刚才的震荡出现了短暂的滞涩，藤盾的构造速度比平时慢了整整一拍。

``唐三！''牢大吼了一声，肘间的 \allowbreak{}\NovelInline{I00141}{\varepsilon_\omega}\allowbreak{} 涟漪全力爆发，朝着不乐的方向碾压而去。

\allowbreak{}\NovelInline{I00142}{\varepsilon_\omega}\allowbreak{} 的波动终于出手了。那道无形的力量如同一柄巨锤，重重地砸在不乐的光刃上。\allowbreak{}\NovelInline{I00143}{\varepsilon_\omega}\allowbreak{} 对 \allowbreak{}\NovelInline{I00144}{\varepsilon_1}\allowbreak{} 的层级压制在这一刻体现得淋漓尽致------光刃在接触 \allowbreak{}\NovelInline{I00145}{\varepsilon_\omega}\allowbreak{} 的瞬间便出现了细密的裂纹，蛮力被一层一层地剥离、瓦解。

但不乐的蛮力实在太过恐怖。

\allowbreak{}\NovelInline{I00146}{\varepsilon_1}\allowbreak{} 的蛮力输出上限远超常人的想象------那不是一层，而是无数层。\allowbreak{}\NovelInline{I00147}{\varepsilon_\omega}\allowbreak{} 碾碎了第一层，第二层便涌了上来；碾碎了第二层，第三层、第四层、第五层接踵而至。不乐像一头被逼到绝路的疯兽，将自己毕生修炼的序数根基毫无保留地倾泻而出，一层接一层，绵绵不绝。

牢大的 \allowbreak{}\NovelInline{I00148}{\varepsilon_\omega}\allowbreak{} 屏障在疯狂的蛮力冲击下开始剧烈震颤。不是因为他打不过不乐，而是因为不乐根本不在乎自己的死活------他把自己的身体当成了武器，用序数根基的崩溃为代价，换取短暂的、超越自身极限的蛮力爆发。

这种自毁式的疯狂，让 \allowbreak{}\NovelInline{I00149}{\varepsilon_\omega}\allowbreak{} 的压制效率大打折扣。牢大不敢全力出手------\allowbreak{}\NovelInline{I00150}{\varepsilon_\omega}\allowbreak{} 的余波太强了，如果他不加控制地碾压过去，不乐固然会被碾成齑粉，但方圆十米内的唐三、小舞、还有昏倒在地的林婉，都会被 \allowbreak{}\NovelInline{I00151}{\varepsilon_\omega}\allowbreak{} 的余波一并吞噬。

他只能一边压制不乐的蛮力，一边小心翼翼地控制余波的范围，这让他的战力大打折扣。

就在这进退两难之际，不乐做出了最后的疯狂之举。

他将体内残存的全部序数根基凝聚成一团，\allowbreak{}\NovelInline{I00152}{\varepsilon_1}\allowbreak{} 的蛮力已经不再是光芒，而是化作了一道肉眼可见的浑浊波纹，像一颗即将坍缩的黑洞，朝着离他最近的唐三冲去。那团浑浊的力量已经不成形状，所过之处，草木枯萎，大地龟裂，空气都在尖啸。

这是不乐的最后一击------也是他毕生修为的绝唱。

唐三咬紧牙关，蓝银草在身前疯狂编织，但刚才的震荡让他体内的序数运转出现了严重的滞涩，藤盾的构造速度比平时慢了整整一拍。他能感觉到那道浑浊的 \allowbreak{}\NovelInline{I00153}{\varepsilon_1}\allowbreak{} 蛮力正在以不可阻挡之势逼近，\allowbreak{}\NovelInline{I00154}{\varepsilon_0}\allowbreak{} 的防御在这股力量面前脆弱得像一层薄纱------

就在这千钧一发之际------

一道浑厚的序数波动从天而降。

那波动没有任何花哨的光芒，没有刺耳的轰鸣，只是一股沉稳到极致的力量，像一座大山无声无息地压了下来。但就是这股沉默的力量，在接触不乐那团疯狂光芒的瞬间，将其像捏碎一颗鸡蛋一样，轻描淡写地碾成了齑粉。

不乐的身体猛地僵住了。

他瞪大了眼睛，嘴巴张得老大，却发不出任何声音。他感觉到自己体内那团狂暴的 \allowbreak{}\NovelInline{I00155}{\varepsilon_1}\allowbreak{} 蛮力，在这股降临的波动面前，就像一只蚂蚁面对汪洋大海------不是被击败，而是被从根本上否定。他的 \allowbreak{}\NovelInline{I00156}{\varepsilon_1}\allowbreak{} 在这股力量面前，连最低级的 \allowbreak{}\NovelInline{I00157}{\omega}\allowbreak{} 都维持不住，像冰雪遇到烈日一样，无声无息地消融殆尽。

赵无极的身影从半空中缓缓落下，双脚踩在地面上的时候，连灰尘都没有扬起半分。他的表情平静得可怕，目光扫过狼狈不堪的五人------唐三嘴角带血、单膝跪地，小舞昏迷在灌木丛中，牢大浑身紧绷、肘间的 \allowbreak{}\NovelInline{I00158}{\varepsilon_\omega}\allowbreak{} 还在微微颤抖------最后落在瘫软在地、浑身抽搐的不乐身上，声音低沉得像闷雷：

``我说过，有事用蓝银草的感知频段联系我。信号被干扰了，就不知道派人来报信？''

他的语气不算严厉，但五个人都低下了头，脸上满是惭愧。

赵无极没有再说什么。他抬手一挥，一道 \allowbreak{}\NovelInline{I00159}{\varphi(3,0)}\allowbreak{} 的序数波动化作无形的枷锁，将不乐彻底捆缚。这枷锁的构造远比不乐的序数禁锢精密万倍，不仅封锁了他的行动，还彻底切断了他与外界序数波动的联系------从今往后，他连 \allowbreak{}\NovelInline{I00160}{\omega}\allowbreak{} 的最低层级都无法触及，真正成了一个废人。

``带走。''赵无极淡淡地说道。

\chapter{第九章：争吵，兄长的责任}\label{ux7b2cux4e5dux7ae0ux4e89ux5435ux5144ux957fux7684ux8d23ux4efb}

赵无极转身走向林婉，蹲下身检查她的伤势。五个人站在原地，一时间谁都没有说话。空气中弥漫着劫后余生的庆幸，和另一种更加沉重的东西------失败的味道。

沉默只持续了几个呼吸。

牢大忽然动了。

他转过身，目光直直地钉在唐三身上。那双一贯沉稳内敛的眼睛里，此刻翻涌着一种极其复杂的情绪------有后怕，有愤怒，有自责，但更多的，是一种压抑到极点后终于爆发的质问。

``唐三。''

他的声音很低，低到只有唐三能听清。但正是这种低沉，比任何怒吼都更有压迫感。

唐三抬起头，嘴角的血迹还没擦干净，目光与牢大撞在一起。

``你制定的战术。''牢大一字一顿，每个字都像是从牙缝里挤出来的，``正面牵制、侧翼偷袭、远程策应------听起来很完美，是不是？''

唐三沉默了一瞬，沉声道：``我知道哪里出了问题------''

``你当然知道。''牢大打断了他，声音陡然拔高了半分，``你从一开始就知道对方的修为是 \allowbreak{}\NovelInline{I00161}{\varepsilon_1}\allowbreak{}，比你的 \allowbreak{}\NovelInline{I00162}{\varepsilon_0}\allowbreak{} 高出了整整一个 \allowbreak{}\NovelInline{I00163}{\varepsilon}\allowbreak{} 层级。你知道，你仍然让小舞从侧翼切入。你知道 \allowbreak{}\NovelInline{I00164}{\varepsilon_0}\allowbreak{} 的构造挡不住 \allowbreak{}\NovelInline{I00165}{\varepsilon_1}\allowbreak{} 的蛮力，你仍然让她冲了上去。''

他的胸口剧烈起伏着，肘间的 \allowbreak{}\NovelInline{I00166}{\varepsilon_\omega}\allowbreak{} 涟漪随着情绪的波动不受控制地向外扩散，脚下的落叶被震得粉碎。

``你看到小舞被掀飞的时候，你在做什么？你在分神！你的感知网络在那一瞬间出现了缺口，直接导致不乐抓住了机会朝你出手！如果不是我及时挡下那一击，你现在已经------''

他没有说完，但那个未说出口的词，比任何语言都更沉重。

唐三的脸色白了几分。不是因为愤怒，而是因为牢大说的每一个字都是事实。他无法反驳。

但他还是开了口。

``牢大，''唐三的声音沙哑，却异常平静，``你说得对，战术是我制定的，小舞受伤是我的判断失误。但有一件事，我想问你。''

牢大眉头一皱。

``不乐用人质做盾牌的时候，你为什么不出手？''

这句话像一根针，精准地扎在了牢大最不愿触碰的地方。

牢大的脸色瞬间沉了下来：``你什么意思？''

``我的意思是，''唐三直视着他的眼睛，声音不大，却字字清晰，``你的 \allowbreak{}\NovelInline{I00167}{\varepsilon_\omega}\allowbreak{} 对 \allowbreak{}\NovelInline{I00168}{\varepsilon_1}\allowbreak{} 有层级压制。如果你在当时全力出手，不乐根本没有机会发动自毁式爆发，小舞不会受伤，我也不会被反击。''

``你说我不敢全力出手？''牢大的声音陡然冷了下来，眼中的质问变成了锐利的锋芒，``唐三，你动动脑子想一想------不乐手里掐着人质，\allowbreak{}\NovelInline{I00169}{\varepsilon_\omega}\allowbreak{} 的余波范围有多大你比谁都清楚。我全力出手的瞬间，不乐固然会被碾碎，但方圆十米内的林婉、你、还有后来冲上去的小舞，全部都会被 \allowbreak{}\NovelInline{I00170}{\varepsilon_\omega}\allowbreak{} 的余波波及。\allowbreak{}\NovelInline{I00171}{\varepsilon_\omega}\allowbreak{} 对 \allowbreak{}\NovelInline{I00172}{\varepsilon_1}\allowbreak{} 的余波是什么级别，你自己算过没有？''

``我算过。''唐三的声音依然平静，但眼底深处有什么东西在燃烧，``所以我没有要求你全力出手。我让你正面牵制，吸引他的注意力，给我和小舞争取侧翼切入的时间。这个战术的核心不是让你一个人解决战斗，而是------''

``而是让我当靶子。''牢大冷冷地接过了话头，嘴角浮起一丝苦涩的弧度，``让我站在最前面，用 \allowbreak{}\NovelInline{I00173}{\varepsilon_\omega}\allowbreak{} 的威压牵制住一个 \allowbreak{}\NovelInline{I00174}{\varepsilon_1}\allowbreak{} 的疯狗，好让你和小舞从侧面偷偷摸摸地搞小动作。结果呢？你的小动作被发现了，小舞受伤了，你被反击了，最后还是得我出手收拾残局------不，最后连我都收拾不了，还得等赵老师来救场。''

他的声音越来越大，到最后几乎是吼出来的：``唐三，你是队长还是我是队长？这场战斗从头到尾都是你在指挥，你做了所有的判断，下了所有的命令。你判断对方修为 \allowbreak{}\NovelInline{I00175}{\varepsilon_1}\allowbreak{}，却仍然选择正面硬碰；你知道我和 \allowbreak{}\NovelInline{I00176}{\varepsilon_1}\allowbreak{} 之间没有碾压级的差距，却仍然把牵制的重任压在我身上；你知道小舞的 \allowbreak{}\NovelInline{I00177}{\varepsilon_0[3]}\allowbreak{} 挡不住 \allowbreak{}\NovelInline{I00178}{\varepsilon_1}\allowbreak{} 的蛮力，你仍然让她冲了上去！''

``你知道小舞的 \allowbreak{}\NovelInline{I00179}{\varepsilon_0[3]}\allowbreak{} 挡不住 \allowbreak{}\NovelInline{I00180}{\varepsilon_1}\allowbreak{} 的蛮力，你仍然让她冲了上去！''牢大的声音在山林间回荡，连远处的飞鸟都被惊得四散而去。

他向前迈了一步，与唐三之间的距离缩短到不到一臂。那双一向沉稳的眼睛此刻布满了血丝，像是一头被逼到绝路的困兽。

``小舞也算是你的妹妹。''

他的声音忽然低了下来，低到只有两个人能听见。但正是这种低沉，比之前的怒吼更加沉重，更加锋利。

``你就是这么保护她的吗？''

唐三的瞳孔猛地一缩。

``你一个兄长，就是这么当的？''

这句话像一柄钝刀，不是砍，而是慢慢地、一点一点地锯进了唐三的胸口。没有鲜血淋漓的痛快，只有绵长的、闷钝的、几乎令人窒息的疼。

他的嘴唇动了动，却没有发出声音。

牢大看着他，眼中的怒火渐渐被另一种东西取代------那是一种深沉的、无法言说的失望。

``她叫你一声哥。''牢大的声音沙哑，像是砂纸磨过生锈的铁，``她信任你，你说什么她就做什么。你让她从侧翼切入，她连犹豫都没犹豫一下就冲了上去。你知道为什么吗？因为她相信你的判断，她相信你不会让她陷入危险。''

他顿了顿，声音更低了：``但你辜负了这份信任。''

唐三的脸色瞬间变得惨白。

不是被激怒的白，是被戳中了最柔软、最隐秘的痛处之后的白。

他一直以为，自己对小舞的保护已经足够。从武魂觉醒那天起，他就发誓要守护身边的人------小舞是他最重要的妹妹，是他在这世上最不愿看到受到伤害的人。他修炼蓝银草，钻研序数构造，苦练战术判断，一切的一切，都是为了有一天能够拥有足够的力量，护她周全。

但今天，他亲手把她送上了战场。

不是不乐把她推上去的，是他------唐三------亲手把她推上去的。

他制定了战术，他下达了指令，他让小舞从侧翼切入一个 \allowbreak{}\NovelInline{I00181}{\varepsilon_1}\allowbreak{} 强者的攻击范围。他知道 \allowbreak{}\NovelInline{I00182}{\varepsilon_0[3]}\allowbreak{} 挡不住 \allowbreak{}\NovelInline{I00183}{\varepsilon_1}\allowbreak{} 的蛮力，他知道这个战术的成功率极低，但他仍然做了这个决定。因为他觉得，如果配合得当，如果时机精准，小舞有机会在不被发现的���况下解除人质禁锢。

但''如果''没有发生。

小舞被掀飞了，嘴角流血，序数根基震荡。她倒在灌木丛中的那一刻，脸色惨白得像一张纸。

而他------她的兄长------连第一时间跑到她身边都做不到，因为他自己也在被不乐的反击追着打。

``你说得对。''

唐三终于开口了。声音很轻，轻到几乎被风吹散。

牢大微微一怔，他做好了唐三反驳的准备，却没有做好迎接这句话的准备。

``你说得对。''唐三重复了一遍，这一次声音稍微大了一些，但仍然沙哑得厉害，``我不配。''

他低下头，看着自己还在微微发抖的双手。那双手刚才拼命编织蓝银草的防御，却挡不住 \allowbreak{}\NovelInline{I00184}{\varepsilon_1}\allowbreak{} 的一击；那双手刚才试图伸向小舞的方向，却被不乐的反击逼了回来。

``从她叫我哥的那一天起，我就告诉自己，不管发生什么，我都不会让她受伤。''他的声音断断续续，像是一根绷到极限的弦，``但今天我亲手把她送上了一个她根本不该去的战场。我知道 \allowbreak{}\NovelInline{I00185}{\varepsilon_0[3]}\allowbreak{} 挡不住 \allowbreak{}\NovelInline{I00186}{\varepsilon_1}\allowbreak{}，我还是让她去了。''

他闭上眼睛，睫毛微微颤抖：``如果她今天不是被掀飞，而是\ldots\ldots{}''

他没有说下去。那个''而是''后面的可能性，像一把悬在头顶的刀，让他不敢想，却又无法不想。

牢大看着他，眼中的怒火慢慢熄灭了。

取而代之的，是一种同样沉重的东西------自责。

因为他忽然意识到，自己刚才那番话，说得太重了。

``唐三\ldots\ldots{}''他开口，声音里的锋芒已经褪去了大半，只剩下疲惫和懊悔，``我刚才说的话------''

``你没错。''唐三抬起头，打断了牢大。

他的眼眶微微发红，但没有泪水。那双眼睛里翻涌着太多情绪------愧疚、自责、痛苦------但在最深处，有一簇微弱却坚定的火焰，正在慢慢燃起。

``你说的那些话，每一句都是对的。''唐三深吸一口气，声音逐渐平稳下来，``我是她的兄长，我让她受了伤，这是我的错。不管我的初衷是什么，不管我的战术逻辑是什么，结果就是------她受伤了，而我没有保护好她。这一点，我无法辩解。''

他顿了顿，目光变得沉静而坚定：``但你也有你的问题，牢大。''

牢大微微一愣。

``你的 \allowbreak{}\NovelInline{I00187}{\varepsilon_\omega}\allowbreak{} 是我们的底牌，但如果你永远不敢全力打出来，这张底牌就只是一张纸。''唐三的声音平静，却字字清晰，``你的基底稳态确实是你最大的优势，但如果它变成了你不敢出手的理由，那它同时也是你最大的枷锁。''

牢大的脸色变了变，嘴唇微微翕动，想要反驳，但最终没有说话。

``我不是在推卸责任。''唐三看着他，目光坦荡，``我犯的错，我认。但你犯的错，你也得认。我们两个都有问题，这不是谁对谁错的事。''

牢大沉默了很久。

风穿过林间，吹动他额前的碎发。他低头看着自己的双手------那双手做过千万次零误差的基础迭代，打磨出了整个队伍最稳固的基底。但此刻，这双手在微微发抖。不是因为恐惧，而是因为唐三的话像一面镜子，照出了他一直以来不愿面对的东西。

``你说得对。''他终于开口，声音沙哑，却异常坦然，``我确实怕。怕全力爆发的时候，基底出现偏差。怕我花了三年时间搭建的完美架构，在一次冒险中毁于一旦。''

他抬起头，目光与唐三对视：``但今天的事让我明白了一件事------如果完美意味着不敢出手，那这种完美不要也罢。''

唐三看着他，嘴角微微上扬。那是一个极淡的笑容，却包含了太多东西------歉意、认可、以及一种只有并肩作战的人才能理解的默契。

``我不是要逼你冒险。''唐三轻声说道，``我只是希望，下一次遇到同样的情况，你能够相信自己的基底。你花了三年时间打磨的东西，不会那么轻易就碎掉。''

牢大深吸一口气，缓缓点了点头。

就在这时，灌木丛中传来一阵窸窣的响动。小舞扶着树干，摇摇晃晃地站了起来。她的脸色还很苍白，但步伐已经稳了许多。\allowbreak{}\NovelInline{I00188}{\varepsilon_0[3]}\allowbreak{} 的光芒在她指尖重新亮了起来，虽然微弱，却无比坚定。

``你们吵完了？''她的声音还带着一丝虚弱，但语气里已经有了几分平时那种轻快的劲儿，``吵完了就过来帮我看看林婉，她好像快醒了。''

气氛在这一瞬间松弛了下来。

蔡须坤长舒了一口气，拍了拍胸口：``可算结束了，你们俩再吵下去，我都要替你们累了。''

奥斯卡笑着摇了摇头，从储物魂导器中又取出两瓶恢复药剂，一瓶递给唐三，一瓶递给牢大：``先把伤处理了。赵老师那边还不知道什么时候回来，我们得自己撑住。''

唐三接过药剂，一饮而尽。温热的药液顺着喉咙滑下，体内紊乱的序数运转渐渐平稳下来。他走到林婉身边，蹲下身检查了一下她的状况------脖颈上的淤青已经消退了不少，呼吸平稳，意识正在慢慢恢复。

``你\ldots\ldots 你们是谁\ldots\ldots{}''林婉迷迷糊糊地睁开眼，看到的是一张张年轻的面孔。她的声音沙哑，带着劫后余生的惊恐。

``我们是路过的魂师。''唐三温和地说道，蓝银草轻轻托起她的肩膀，让她靠得更舒服一些，``那个伤害你的人已经被制服了，你安全了。''

林婉的目光越过唐三，看到了被 \allowbreak{}\NovelInline{I00189}{\varphi(3,0)}\allowbreak{} 枷锁捆得结结实实、瘫在地上动弹不得的不乐。她的身体猛地一颤，眼中涌出泪水，却拼命咬着嘴唇不让自己哭出声来。

``他\ldots\ldots 他不会再伤害我了吧？''她的声音颤抖。

``不会了。''牢大走到她面前，声音低沉却异常坚定，``我保证。''

他说出这三个字的时候，目光不自觉地扫过自己的双手。那双手还在微微发抖，但抖动的幅度已经比刚才小了很多。

唐三说得对------他花了三年时间打磨的基底，不会那么轻易就碎掉。

\begin{center}\rule{0.5\linewidth}{0.5pt}\end{center}

又过了大约一炷香的时间，赵无极终于回来了。

他的身影从林梢落下，目光扫过五个人------唐三嘴角的血迹已经擦干净了，小舞靠在树上闭目调息，牢大站在原地沉默不语，蔡须坤和奥斯卡一左一右守着林婉。所有人的身上都带着或轻或重的伤，但精神还算稳定。

然后他的目光落在地上瘫软的不乐身上，眉头微微一挑。

``\allowbreak{}\NovelInline{I00190}{\varepsilon_1}\allowbreak{}？''他蹲下身，伸手探了探不乐体内残存的序数波动，嘴角浮起一丝不屑的冷笑，``就这种货色，也值得你们五个联手？''

五个人都低下了头，没有人接话。

赵无极站起身，目光从五个人脸上一一扫过。他不需要问发生了什么------蓝银草的感知频段虽然被干扰了，但赵无极自己的感知能力远超 \allowbreak{}\NovelInline{I00191}{\varepsilon_1}\allowbreak{} 的干扰范围。他在赶回来的路上，就已经通过残留的序数痕迹还原了大半个战斗过程。

``你们今天犯的错误，我全部看到了。''他的声音不急不缓，却像一记记重锤，敲在每个人的心上。

``唐三，你的战术判断有两处失误。第一，你在明知对方修为 \allowbreak{}\NovelInline{I00192}{\varepsilon_1}\allowbreak{} 的情况下，仍然选择让小舞从侧翼切入。\allowbreak{}\NovelInline{I00193}{\varepsilon_0[3]}\allowbreak{} 和 \allowbreak{}\NovelInline{I00194}{\varepsilon_1}\allowbreak{} 之间差了一个完整的 \allowbreak{}\NovelInline{I00195}{\varepsilon}\allowbreak{} 层级，你的侧翼偷袭从一开始就不可能成功。第二，你在感知网络出现缺口的时候，没有第一时间收缩防御，而是试图维持全覆盖。一个人的感知再精密，也不可能同时覆盖五个战场。''

唐三沉声应道：``是，我的错。''

赵无极点了点头，目光转向牢大：``牢大，你的问题不在战术层面，而在心理层面。你的 \allowbreak{}\NovelInline{I00196}{\varepsilon_\omega}\allowbreak{} 对 \allowbreak{}\NovelInline{I00197}{\varepsilon_1}\allowbreak{} 有层级优势，但你不敢全力出手。你的基底稳态是你最大的武器，但你把它变成了不敢冒险的借口。今天如果我不在，你打算怎么办？就这么一直僵持下去，直到不乐力竭而亡？还是等到他的蛮力把你的屏障磨穿？''

牢大沉默了片刻，沉声道：``我不知道。''

``不知道就对了。''赵无极的语气意外地平和，``因为你从来没有面对过这种局面------你的力量明明比对手强，却因为种种顾虑而无法施展。这种'有力使不出'的感觉，比'没有力量'更折磨人。但这就是实战，实战不会给你完美的条件，不会等你准备好了再开始。''

他的目光最后扫过所有人：``但有一件事你们做得不错------唐三在第一时间向我发送了信号。虽然信号被干扰了，但如果没有那个信号，我还要再晚半刻钟才能赶到。半刻钟，足够不乐再杀两个人了。''

他顿了顿，声音变得深沉：``还有一件事------你们吵了。''

五个人的表情同时僵了一下。

``我回来的路上，感知到了你们的情绪波动。''赵无极的嘴角浮起一丝极淡的笑意，``牢大质问唐三，唐三反驳牢大，两个人差点打起来。''

牢大的脸色微微发红，唐三则低下了头。

``但我不打算批评你们。''赵无极的话让所有人都愣住了，``相反，我觉得这是好事。''

``好事？''蔡须坤忍不住插嘴。

``当然。''赵无极负手而立，目光深邃，``你们今天经历了一场惨败，被一个 \allowbreak{}\NovelInline{I00198}{\varepsilon_1}\allowbreak{} 的蛮夫打得手忙脚乱。如果你们事后嘻嘻哈哈、互相安慰、说什么'没关系下次努力'，那才是真正的问题。牢大敢质问唐三，说明他在乎这场战斗的胜负；唐三敢反驳牢大，说明他有自己的判断和立场。两个人都敢把心里的话说出来，而不是憋在心里装没事------这才是一个团队该有的样子。''

他顿了顿，声音变得郑重：``但我要提醒你们一件事------争吵可以有，但必须在战斗结束之后再吵。战斗进行中的时候，指挥权归谁，就听谁的。今天唐三是指定的战术指挥，牢大你有不同意见可以提，但一旦唐三做了决定，你就必须执行。哪怕那个决定在你看来是错的，你也必须先执行，再争论。因为一个执行错误决定的团队，永远比一个在战场上争论不休的团队更有胜算。但是你也记住，你是队长，修为最高的人''

牢大和唐三对视一眼，同时点了点头。

``好了。''赵无极拍了拍手，语气恢复了日常的平淡，``收拾一下，出发去补给镇。小舞和林婉的伤需要处理，不乐要交给镇上的治安官。至于你们今天的失误------''

他回头看了五个人一眼，嘴角扬起一抹意味深长的弧度：

``到了镇子上，我会给你们加训。''

五个人的脸色同时变了。

小舞捂着还在隐隐作痛的胸口，哀嚎了一声：``赵老师，我都受伤了\ldots\ldots{}''

``受伤了更要练。''赵无极头也不回地朝前走去，声音悠悠飘来，``下次遇到 \allowbreak{}\NovelInline{I00199}{\varepsilon_1}\allowbreak{} 的对手，总不能还等我来救场吧？''

五个人面面相觑，不约而同地苦笑了一下。

阳光透过树叶的缝隙，洒在林间的小道上。五人的身影护送着林婉，朝着补给镇的方向缓缓走去。不乐早已被捆得结结实实扔在了牢大的肩上，连挣扎的力气都没有了。

牢大走在队伍中间，沉默不语。他的右手不自觉地攥了攥拳，又松开。肘间的 \allowbreak{}\NovelInline{I00200}{\varepsilon_\omega}\allowbreak{} 涟漪安安静静地流淌着，比之前平稳了许多。

唐三走在最前面，蓝银草的感知重新铺展开来，这一次，他的探测范围比之前小了许多------不再试图覆盖整个战场，而是集中在前方五十米的范围内，精密而克制。

小舞靠在蔡须坤身边，一瘸一拐地走着，指尖的 \allowbreak{}\NovelInline{I00201}{\varepsilon_0[3]}\allowbreak{} 光芒比之前亮了一些。她偶尔抬头看一眼走在前面的唐三和旁边的牢大，嘴角浮起一丝淡淡的笑意。

吵过了，说开了，伤口还在疼。

但他们都知道，下一次再遇到同样的情况，他们不会再犯同样的错误。

不是因为变强了，而是因为------他们终于学会了，如何在不够强的时候，做出正确的选择。

赵无极走在队伍最后，看着眼前这一幕，嘴角扬起一抹不易察觉的弧度。

今天的战斗，是一场惨败。但这场败仗里长出来的东西，比任何胜利都更加珍贵。

\allowbreak{}\NovelInline{I00202}{\varepsilon_0}\allowbreak{} 和 \allowbreak{}\NovelInline{I00203}{\varepsilon_1}\allowbreak{} 之间的鸿沟，不是技巧能填平的。想要跨越这道鸿沟，唯有变强------不是一个人的变强，而是整个团队的变强。

而变强的第一步，从来不是修炼，而是认清自己的不足。

今天，他们认清了。

\chapter{第十章：补给镇的锤炼}\label{ux7b2cux5341ux7ae0ux8865ux7ed9ux9547ux7684ux9524ux70bc}

补给镇不大，拢共也就三条街、两口井、一座歪歪斜斜的钟塔。镇子依着星斗大森林的外围而建，街上走的十个人里有八个是魂师，剩下两个是卖魂师吃喝拉撒的商贩。空气里永远弥漫着一股混杂了药草、烤肉和序数波动的古怪气味。

赵无极在镇子东头租了一间废弃的仓库，四面石墙、屋顶塌了半边，地面坑坑洼洼，勉强能站开五个人。他把不乐交给治安官之后，带着五个少年径直来到这里，一脚踹开锈迹斑斑的铁门，回头看了他们一眼。

``进来。''

五个人鱼贯而入。仓库里光线昏暗，只有屋顶的破洞漏下几缕阳光，在地面上投下不规则的光斑。空气中漂浮着细密的灰尘，被从破洞中吹进来的风搅得团团转。

赵无极站在仓库中央，负手而立，目光从五个人脸上一一扫过。

``从今天开始，每天辰时到这里，酉时收工。''他的声音平淡，却带着不容置疑的威严，``训练内容只有一个------配合。''

``配合？''小舞下意识地摸了摸还在隐隐作痛的胸口，``赵老师，我们不是应该先提升个人修为吗？如果我的 \allowbreak{}\NovelInline{I00204}{\varepsilon_0[3]}\allowbreak{} 能突破到 \allowbreak{}\NovelInline{I00205}{\varepsilon_0[2]}\allowbreak{}，下次遇到 \allowbreak{}\NovelInline{I00206}{\varepsilon_1}\allowbreak{} 的对手------''

``下次遇到 \allowbreak{}\NovelInline{I00207}{\varepsilon_1}\allowbreak{} 的对手，你就算突破到 \allowbreak{}\NovelInline{I00208}{\varepsilon_0[\omega]}\allowbreak{} 也打不过。''赵无极毫不客气地打断了她，``\allowbreak{}\NovelInline{I00209}{\varepsilon_0}\allowbreak{} 和 \allowbreak{}\NovelInline{I00210}{\varepsilon_1}\allowbreak{} 之间的差距，不是基本列的下一项能弥补的。那是整个 \allowbreak{}\NovelInline{I00211}{\varepsilon_0}\allowbreak{} 层级的全部迭代加在一起，才刚刚够到 \allowbreak{}\NovelInline{I00212}{\varepsilon_1}\allowbreak{} 的门槛。你想靠个人修为跨越这道鸿沟，至少要修三年。''

小舞的脸色白了一下，嘴巴张了张，最终没有反驳。

``但配合不一样。''赵无极竖起一根手指，``配合能在短时间内弥补修为差距。你们昨天的问题不是个人实力不够------牢大的 \allowbreak{}\NovelInline{I00213}{\varepsilon_\omega}\allowbreak{} 足以碾压 \allowbreak{}\NovelInline{I00214}{\varepsilon_1}\allowbreak{}，唐三的感知精度也远超同级。问题在于，你们五个人各打各的，像五根散落的筷子，一折就断。但如果把五根筷子捆在一起------''

他随手从地上捡起五根枯枝，攥在掌心，轻轻一折，纹丝不动。

``------就折不断。''

他将枯枝扔到地上，拍了拍手上的灰尘：``今天的训练，从最基础的开始。''

``唐三，把你的蓝银草感知网络铺开。''

唐三依言释放蓝银草，淡绿色的藤蔓无声无息地从地面蔓延开来，\allowbreak{}\NovelInline{I00215}{\omega^{\omega^{\omega^\omega}}}\allowbreak{} 的序数纹路在藤条表面若隐若现，编织成一张精密的感知网络，覆盖了整个仓库。

``现在，把感知权限开放给其他人。''

唐三微微一愣：``全部开放？''

``全部。''赵无极的语气不容商量。

唐三犹豫了一瞬。蓝银草的感知网络是他最核心的能力，将权限完全开放给他人，意味着他的所有感知数据、所有探测结果、甚至他自身的序数波动频率，都会毫无保留地暴露在队友面前。这在实战中是极其危险的事情------如果队友中有任何一个人被敌人控制或策反，他的感知网络就会变成刺向自己心脏的利刃。

但他只犹豫了一瞬，便点了点头。

蓝银草的藤蔓微微颤动，感知网络的边界向外扩展了一圈，将小舞、牢大、蔡须坤和奥斯卡全部纳入了覆盖范围。与此同时，一道极其细微的序数波动从唐三的眉心扩散开来，像是一把钥匙，打开了感知网络的共享通道。

``感受到了吗？''赵无极问道。

小舞闭上眼睛，片刻后点了点头：``我能感觉到唐三的感知范围\ldots\ldots 好远，比我自己的感知范围大了至少三倍。''

``不只是范围。''牢大沉声说道，他的眉头微微皱起，似乎在适应一种全新的感知方式，``我能感觉到唐三感知到的每一个细节------地面裂缝的深度、墙壁上序数残留的浓度、甚至空气流动的轨迹。这些信息\ldots\ldots 比我自己的感知精度高了一个层级。''

``这就是感知共享的意义。''赵无极说道，``唐三的感知精度是五人中最高的，但他一个人的精力有限，不可能同时处理所有信息。如果你们五个人都能接入他的感知网络，就等于每个人都拥有了唐三级别的感知能力。反过来，唐三也能通过你们的感知补充盲区------比如小舞对 \allowbreak{}\NovelInline{I00216}{\varepsilon_0}\allowbreak{} 基本列波动的敏感度，牢大对序数重力场的感知范围，这些都是唐三自己不具备的。''

他顿了顿，目光变得锐利：``但感知共享有一个前提------信任。你们必须完全信任彼此，才能毫无保留地开放自己的感知。昨天那场战斗，唐三的感知网络出现缺口，不是因为他的能力不够，而是因为他不敢把感知权限分出去。他习惯了一个人掌控所有信息，一个人做所有判断。这种独狼式的作战方式，在面对弱敌时没有问题，但面对 \allowbreak{}\NovelInline{I00217}{\varepsilon_1}\allowbreak{} 级别的对手，一个人的感知永远不够用。''

唐三低下头，沉默不语。赵无极说的每一个字，都精准地戳中了他的要害。

``现在，所有人闭上眼睛，用感知网络去'看'彼此。''赵无极的声音放缓了几分，``不要急着分析信息，先学会感受------感受队友的存在，感受他们的序数波动节奏，感受他们的呼吸和心跳。配合的第一步，不是战术，而是默契。默契不是练出来的，是感知出来的。''

五个人闭上了眼睛。

仓库里安静了下来，只有风从破洞中吹过的呜咽声。

起初，感知共享带来的信息洪流让每个人都感到不适。唐三的感知数据太过庞大，太过精密，像是一条奔腾的河流突然灌进了五条细小的溪流，每个人的脑海都被海量的信息淹没------地面的每一道裂纹、空气中的每一缕流动、墙壁上残留的每一丝序数波动，全部涌入感知，密密麻麻，令人头晕目眩。

小舞第一个皱起了眉头，她的 \allowbreak{}\NovelInline{I00218}{\varepsilon_0[3]}\allowbreak{} 构造在信息洪流中剧烈震荡，差点崩溃。牢大的眉心渗出了细密的汗珠，他的基底稳态在努力维持着自身的序数平衡，同时又要消化来自唐三的额外感知数据，两股力量在他体内相互拉扯，像是要把他的意识撕成两半。蔡须坤和奥斯卡的情况稍好一些，但也好不到哪里去------金色纽带的联结能力在感知共享的冲击下变得极不稳定，时断时续。

唐三自己也不好受。他是感知网络的中心节点，所有人的反馈数据都汇聚到他一个人身上。五个人、五种不同的感知方式、五种不同的序数波动频率，全部叠加在一起，像是一场永不停歇的风暴，在他的意识深处呼啸翻滚。

``稳住。''赵无极的声音像一根定海神针，穿透了信息风暴的喧嚣，``不要试图分析所有信息，先学会过滤。唐三，你是中心节点，你来决定哪些信息需要共享，哪些信息可以屏蔽。其他人，不要被动接收，主动去筛选------只接收和你相关的信息，其余的全部过滤掉。''

唐三深吸一口气，开始调整感知网络的信息分配。他将地面震动的数据优先传递给牢大------因为牢大的基底稳态对重力场变化最敏感；将序数波动的频率数据优先传递给小舞------因为她的 \allowbreak{}\NovelInline{I00219}{\varepsilon_0[3]}\allowbreak{} 构造对波动频率的变化最为敏锐；将远程距离的感知数据优先传递给蔡须坤和奥斯卡------因为他们的金色纽带擅长长距离联结。

渐渐地，信息洪流从混沌变得有序。每个人只接收自己最擅长处理的那部分数据，其余的交给更合适的队友去消化。就像一条河流分成了五条支流，每条支流都流向最适合它的田地。

半个时辰后，五个人几乎同时睁开了眼睛。

``感觉到了？''赵无极问道。

小舞点了点头，眼中闪烁着兴奋的光芒：``我能感觉到牢大的位置，精确到厘米。我能感觉到唐三的感知范围边界在哪里，甚至能感觉到奥斯卡的香肠增益波动还有多久会衰减。''

``这就是配合的基础。''赵无极说道，``不是战术，不是指令，而是感知。当你能感觉到队友的存在、状态和需求时，你就不需要等指令再行动------你会在本能的驱动下，做出最正确的选择。''

第二项训练，赵无极在地上用序数光芒画了一个巨大的递归树。

``你们昨天的战斗，最大的问题不是个人实力不够，而是没有人知道什么时候该接棒。''他指着递归树的每一个节点，``牢大是先手，他的 \allowbreak{}\NovelInline{I00220}{\varepsilon_\omega}\allowbreak{} 负责正面压制。唐三和小舞是侧翼，负责解救人质和补位。蔡须坤和奥斯卡是远程，负责增益和干扰。五个人的分工很明确，但你们在实战中的衔接一塌糊涂------牢大停手的时候唐三没跟上，唐三被反击的时候小舞没有补位，小舞受伤的时候没有人第一时间接应。''

他拍了拍手，五道序数光芒从地面升起，在仓库中央凝聚成一个模糊的人形靶标。

``这是我用 \allowbreak{}\NovelInline{I00221}{\varphi(2,0)}\allowbreak{} 的序数构造凝聚的靶标，防御强度相当于 \allowbreak{}\NovelInline{I00222}{\varepsilon_0}\allowbreak{} 的三倍。你们五个人轮流攻击，每个人出手三次，三次之内必须击破靶标的防御层。规则只有一个------每个人出手的时机，必须衔接在上一个人出手的间隙中。不能同时出手，不能等上一个人完全结束再出手，必须在上一个人出手的瞬间，下一个人的攻击就已经在路上了。''

``这不就是\ldots\ldots 接力？''奥斯卡挠了挠头。

``对，就是接力。''赵无极说道，``但不是跑步接力，是序数接力。你们每个人出手的时候，下一个人必须感知到上一个人的攻击节奏、序数波动频率和攻击方向，然后在这个基础上叠加自己的攻击。不是简单的轮流攻击，而是层层叠加、环环相扣，让五个人的一次出手，变成一次完整的、不间断的连续攻击。''

他顿了顿，目光变得锐利：``如果你们能在三次出手之内击破靶标，说明你们的衔接已经及格了。如果做不到------''

他没有说完，但所有人都明白了''做不到''的后果。

``我先来。''牢大率先站了出来。

他踏前一步，肘间的 \allowbreak{}\NovelInline{I00223}{\varepsilon_\omega}\allowbreak{} 涟漪凝聚成一道沉闷的冲击波，重重地砸在靶标上。靶标的防御层在 \allowbreak{}\NovelInline{I00224}{\varepsilon_\omega}\allowbreak{} 的冲击下剧烈震颤，表面出现了细密的裂纹，但并没有碎裂------\allowbreak{}\NovelInline{I00225}{\varphi(2,0)}\allowbreak{} 的防御强度确实不是 \allowbreak{}\NovelInline{I00226}{\varepsilon_\omega}\allowbreak{} 一次出手就能击破的。

但牢大的出手，为下一个人创造了窗口。

唐三的蓝银草在牢大出手的瞬间便已启动。他没有等牢大的攻击完全结束，而是在冲击波命中靶标的前一刹那，将蓝银草的藤条顺着冲击波的轨迹延伸了出去。\allowbreak{}\NovelInline{I00227}{\varepsilon_0}\allowbreak{} 基本列的精密纹路沿着裂纹渗入靶标内部，像是一把精密的手术刀，沿着牢大撕开的伤口进一步切割。

两道攻击几乎无缝衔接------牢大的 \allowbreak{}\NovelInline{I00228}{\varepsilon_\omega}\allowbreak{} 撕开了外层防御，唐三的 \allowbreak{}\NovelInline{I00229}{\varepsilon_0}\allowbreak{} 渗入了内部结构。靶标的裂纹迅速扩大，防御层开始松动。

小舞第三个出手。她的 \allowbreak{}\NovelInline{I00230}{\varepsilon_0[3]}\allowbreak{} 构造精准地找到了唐三的蓝银草撕开的缝隙，将基本列的迭代力量灌注进去。\allowbreak{}\NovelInline{I00231}{\varepsilon_0[3]}\allowbreak{} 的波动频率与唐三的 \allowbreak{}\NovelInline{I00232}{\varepsilon_0}\allowbreak{} 几乎完全一致，两道力量在靶标内部产生了共振，裂纹以肉眼可见的速度扩张。

蔡须坤的金色纽带第四个到达。他没有直接攻击靶标，而是将金色纽带缠绕在小舞的 \allowbreak{}\NovelInline{I00233}{\varepsilon_0[3]}\allowbreak{} 构造上，为她的攻击叠加了一层增幅。金色纽带的增益波动与小舞的序数构造产生了奇妙的化学反应------\allowbreak{}\NovelInline{I00234}{\varepsilon_0[3]}\allowbreak{} 的威力被放大了将近一倍，靶标的防御层在共振与增幅的双重作用下，开始出现大面积的崩解。

最后是奥斯卡。

他的出手最简单，也最不起眼------只是一瓶普通的恢复药剂，被他用力掷向了靶标的核心。

但就是这看似平平无奇的一掷，成了压垮骆驼的最后一根稻草。

药剂瓶在接触靶标核心的瞬间碎裂，药液渗入靶标内部已经被前四人撕开的层层裂纹中。药液本身没有任何攻击力，但它有一个极其特殊的性质------它能溶解序数构造之间的联结。靶标的 \allowbreak{}\NovelInline{I00235}{\varphi(2,0)}\allowbreak{} 防御层之所以坚固，是因为每一层序数构造之间都有精密的联结纽带。药液渗入之后，这些联结被逐一溶解，防御层的结构像多米诺骨牌一样，从内到外，层层崩解。

``砰''的一声闷响，靶标炸裂成漫天的序数碎片，在空气中化为点点光芒，缓缓消散。

五个人愣了一瞬，然后同时露出了笑容。

``三次出手。''赵无极的声音从旁边传来，带着一丝不易察觉的赞许，``不，准确地说，是一次出手。五个人的攻击衔接得太紧密了，紧密到靶标的防御系统根本没来得及区分攻击来源------它把你们五个人的攻击当成了同一次攻击的五次叠加。''

他走到五个人面前，目光中多了一份认真：``这就是配合的意义。不是五个人各打各的，而是五个人的力量融为一体。牢大的 \allowbreak{}\NovelInline{I00236}{\varepsilon_\omega}\allowbreak{} 开路，唐三的 \allowbreak{}\NovelInline{I00237}{\varepsilon_0}\allowbreak{} 渗透，小舞的 \allowbreak{}\NovelInline{I00238}{\varepsilon_0[3]}\allowbreak{} 共振，蔡须坤的金色纽带增幅，奥斯卡的药液瓦解------五个人的能力各不相同，但衔接在一起，就形成了一套完整的攻击链条。这套链条的威力，远超任何一个人的单独出手。''

他顿了顿，声音变得深沉：``昨天你们面对不乐的时候，如果能有今天这种衔接------牢大出手的同时唐三就准备好渗透，唐三被反击的时候小舞就已经在补位，小舞受伤的时候蔡须坤和奥斯卡的增益就已经到了------结局会完全不同。''

五个人沉默了。他们想起了昨天的战斗，想起了那些错过的时机、断裂的衔接、和因为犹豫而浪费的每一个瞬间。

``再来一次。''牢大率先说道，他的眼中没有了昨天的迷茫和恐惧，取而代之的是一种沉稳的坚定。

``再来。''唐三点了点头。

``再来！''小舞、蔡须坤和奥斯卡异口同声。

赵无极看着这五张年轻而坚毅的面孔，嘴角扬起了一抹极淡的笑意。

他抬手一挥，又一道序数光芒凝聚成形。这一次，靶标的防御强度从 \allowbreak{}\NovelInline{I00239}{\varphi(2,0)}\allowbreak{} 提升到了 \allowbreak{}\NovelInline{I00240}{\varphi(2,1)}\allowbreak{}------比之前整整高了一个层级。

``刚才那个是热身。''赵无极的声音平淡，``这个才是正菜。给你们十次出手的机会，击破它。''

五个人对视一眼，同时踏前一步。

接下来的三天，仓库里每天都是同样的场景。

\chapter{第十一章：φ函数的威力}\label{ux7b2cux5341ux4e00ux7ae0ux3c6ux51fdux6570ux7684ux5a01ux529b}

林间空地的序数涟漪尚未散尽，唐三周身的淡金色光晕骤然暴涨，丹田内那道趋近于 ε₀ 的幂次纹路轰然成型，彻底挣脱了有限迭代的桎梏------\textbf{ε₀ 境界，成了。}

身旁的小舞身影轻颤，指尖燃烧弹上的 ε₀(4) 波动层层收敛，转而化作更凝练的 ε₀ 基准光晕，她脚下的青草被高阶序数气息压得伏倒，同样稳稳踏入了 ε₀ 门槛。几日的闭关打磨，两人不仅吃透了 ε₀ 的无穷迭代精髓，更是摸透了 ε 序列的构造逻辑，修为增速远超众人预料。

赵无极负手而立，周身 φ(3,0) 的序数波动缓缓收敛，看着脱胎换骨的两人，浑厚的嗓音响彻空地：

\begin{quote}
``不错，短短数日便突破 ε₀，悟性远超预期。但你们要记着，ε₀、ζ₀ 不过是高阶序数的起点，真正能抗衡千年人面魔蛛 φ(ω,0) 魂核的力量，藏在 Veblen 函数之中，尤其是咱们接下来要讲的二元 φ 函数------这是你们跨越 ε 序列、触碰超阶序数的关键。''
\end{quote}

牢大挠了挠头，肘间突然泛起浓郁的深黄色光晕，一道形似巨锤、萦绕着 ε\_ω 波动的扽魂虚影缓缓浮现，虚影之上布满细密的序数纹路，轻轻晃动便带起阵阵气流：

\begin{quote}
``赵老师，φ 函数？之前您提过人面魔蛛的防御是 φ(4,0)，这和咱们学的 ε、ζ 序数有啥关系？还有我这扽魂，能不能借助 φ 函数的力量变得更强？''
\end{quote}

``关系极大，你的扽魂也能借 φ 函数之力突破。''赵无极抬手凌空勾画，淡青色的序数符号在空中缓缓成型，``咱们此前学的 ω 幂次、ε 序数、ζ 序数，看似杂乱无章，实则都能被 φ 函数统一归纳。先从最简单的一元 Veblen 函数说起，这是二元函数的根基，记法为 φ(α)，本质就是 ω 的 α 次方，也就是：

\NovelDisplay{M00143R01}{%
\varphi(\alpha) = \omega^\alpha%
}

他指尖轻点，符号不断推演：``按照这个定义，基础性质一目了然：φ(0) = ω⁰ = 1；若是后继序数，φ(α+1) = φ(α) · ω；若是极限序数 α，φ(α){[}n{]} = φ(α{[}n{]})。说白了，一元 φ 函数就是咱们之前修炼的 ω 幂次迭代的标准化写法，是整个 φ 体系的底层逻辑，你的扽魂只要将 ε\_ω 的波动与 φ 函数迭代结合，就能进一步凝练力量。''

唐三凝神记录，蓝银草灵气在空中勾勒出推导纹路：

\begin{quote}
``原来如此，那 ε₀ 就是 φ(1,0)？ζ₀ 就是 φ(2,0)？''
\end{quote}

``一点就通。''赵无极赞许颔首，空中符号瞬间切换为二元结构，``这就是 φ(n,α) 二元 Veblen 函数的核心作用------它专门枚举各类映射的不动点，把咱们之前学的 ε、ζ、η 序数全部归类：

\begin{itemize}
\tightlist
\item
  \textbf{φ(1,α) = ε\_α}，对应 ε 序列的所有不动点，φ(1,0) 就是 ε₀，φ(1,1) 就是 ε₁，φ(1,ω) 就是 ε\_ω；
\item
  \textbf{φ(2,α) = ζ\_α}，对应 ζ 序列的所有不动点，φ(2,0) 就是 ζ₀，也就是咱们之前说的 Cantor 序数；
\item
  \textbf{φ(3,α) = η\_α}，以此类推，每提升第一个下标，就代表迭代层级跃升一个维度。''
\end{itemize}

话音刚落，奥斯卡便举着香肠凑了上来，脸上满是困惑：

\begin{quote}
``赵老师，我有点没懂，不动点到底是什么？比如 φ(1,0) = ε₀，为什么说它是映射的不动点？还有，实战里我只需要稳住波动就好，用不用吃透这些不动点的原理啊？''
\end{quote}

蔡须坤也跟着点头，指尖轻轻一旋，一颗银白色、萦绕着序数气息的篮球凭空出现，篮球表面流转着与他气息同源的纹路，轻轻弹跳间便带动周围气流波动：

\begin{quote}
``我也有疑问，您说第一个下标每加 1，迭代维度就升一级，那 φ(3,0) 比 φ(2,0) 强多少？我这篮球武魂，联动不同层级 φ 函数力量时，能不能通过运球节奏校准气息，避免迭代紊乱？''
\end{quote}

赵无极笑着抬手，空中浮现出清晰的映射示意图，耐心解答：

\begin{quote}
``你们问的都是关键问题，先解答奥斯卡的疑惑。所谓不动点，简单说就是'映射后不变的点'，比如映射 β → ω\^{}β，当 β = ε₀ 时，ω\^{}ε₀ = ε₀，这就是不动点------ε₀ 无论怎么进行 ω 幂次迭代，结果还是它本身，这也是它能突破有限迭代的核心原因。至于你说的实战需求，虽然你主要负责稳定波动，但必须懂原理，不然遇到 φ 函数波动紊乱时，你不知道问题出在迭代层级，就无法精准稳住气息。''
\end{quote}

接着他转向蔡须坤，目光落在他手中的篮球上：

\begin{quote}
``关于迭代维度的差距，φ(3,0) = η₀，是映射 β → φ(2,β) = ζ\_β 的不动点，相当于 ζ 序列的'ε₀'，力量层级比 ζ₀ 高出一个完整的迭代维度------简单说，ζ₀ 的力量是'ε 序列迭代'，而 η₀ 是'ζ 序列迭代'，差距就像 ε₀ 和 ω 的差距。至于你的篮球，完全可以借助运球节奏校准气息，运球的快慢、轻重对应 φ 函数的迭代节奏，只要保持运球节奏与 ζ₀、ε₀ 等迭代同频，就能有效避免紊乱，这也是你联动气息的独特优势。''
\end{quote}

牢大挠了挠头，操控着肘间的扽魂虚影轻轻晃动，虚影碰撞地面发出沉闷的声响：

\begin{quote}
``赵老师，那 φ(ω,0) 是所有有限下标 φ 函数的极限，是不是意味着它比任何 φ(n,0) 都强？咱们现在练的 φ(4,0)，和它差距多大？我这扽魂能扛住 φ(4,0) 的反震吗？''
\end{quote}

``问得好。''赵无极神色一凝，``φ(ω,0) 确实是所有 φ(n,0) 的极限，它的迭代维度是'ω 层叠加'，而 φ(4,0) 只是第四层有限迭代，两者差距悬殊------就像 4 和 ω 的差距。但咱们不用急，先吃透 φ(4,0) 的防御拆解技巧，掌握不同层级 φ 函数的联动，后续再逐步突破，才能稳妥应对人面魔蛛的 φ(ω,0) 魂核。尤其要记住，ζ₀ 作为 φ(2,0)，是衔接 ε 序列和更高层级 φ 函数的关键，吃透 ζ₀ 的迭代逻辑，才能顺利过渡到 φ(3,0)、φ(4,0)，这也是咱们模拟作战的核心重点。你的扽魂力量雄厚，只要稳住 ε\_ω 的波动，配合 ζ₀ 的迭代基准，完全能扛住 φ(4,0) 的反震，但切忌急于发力。''

待众人彻底解惑，赵无极放缓语速，拆解二元 φ 函数的递归基本列------这是修炼操控 φ 函数力量的核心，更是实战破防的关键：

\textbf{φ(1,α) 是映射 β → φ(β) = ω\^{}β 的第 1+α 个不动点，基本列分三种情况：}

\textbf{α = 0 时：}

\NovelDisplay{M00144R01}{%
\varphi(1,0)[0] = 0, \quad \varphi(1,0)[n+1] = \varphi(\varphi(1,0)[n])%
}

也就是 ε₀ 的基本列：\{0, 1, ω, ω\^{}ω, ω\^{}ω\^{}ω, \ldots\}，和咱们此前推导的完全一致；

\textbf{α 为后继序数时：} 若 α = β+1（β 为任意序数），则

\NovelDisplay{M00145R01}{%
\varphi(1,\alpha)[n] = \varphi(1,\beta)+1[n] = \varphi(1,\varphi(1,\beta)+1)[n-1]%
}

简单说，就是以 φ(1,β) 的基本列作为迭代基础，跳出不动点之后，继续进行 ω 幂次映射，比如 φ(1,1) = ε₁，它的基本列就是：

\NovelDisplay{M00146R01}{%
\{\varphi(1,1)[0],\ \varphi(1,1)[1],\ \varphi(1,1)[2],\ \dots\} = \{\varepsilon_0+1,\ \omega^{\varepsilon_0+1},\ \omega^{\omega^{\varepsilon_0+1}},\ \dots\}%
}

\textbf{α 为极限序数时：} 若 α 是极限序数（比如 ω），则

\NovelDisplay{M00147R01}{%
\varphi(1,\alpha)[n] = \varphi(1,\alpha[n])%
}

其中 α{[}n{]} 是 α 的基本列。以 φ(1,ω) = ε\_ω 为例，ω 的基本列是 \{0, 1, 2, \ldots, n, \ldots\}，因此

\NovelDisplay{M00148R01}{%
\varphi(1,\omega)[n] = \varphi(1,n) = \varepsilon_n%
}

其基本列就是 \{ε₀, ε₁, ε₂, \ldots, ε\_n, \ldots\}。

讲解完毕，赵无极抬手一挥，周身序数气息暴涨，淡青色的光晕在空地中央凝聚，逐渐形成一个半透明的虚拟怪物------正是模拟千年人面魔蛛的形态，其周身萦绕着淡紫色的 φ(4,0) 防御屏障，屏障上布满了细密的序数纹路，散发着压迫性的气息。更关键的是，屏障的底层纹路清晰可见 ζ₀ 的迭代轨迹，显然是基于 ζ₀（φ(2,0)）叠加两层迭代构建而成，与赵无极此前强调的''ζ₀ 是过渡关键''完美呼应。

\begin{quote}
``好了，理论讲解结束，现在进入模拟作战。''赵无极的声音变得严肃，``我已将虚拟魔蛛的防御设定为 φ(4,0)，和咱们之前预估的一致，它的防御屏障是基于 φ(4,0) 构建的，本质是在 ζ₀（φ(2,0)）的基础上，叠加两层不动点迭代，拆解难度远高于 ε₀、ζ₀ 层级。你们要记住，想要突破 φ(4,0)，必须先吃透 ζ₀ 的迭代逻辑，从 ζ₀ 入手拆解底层屏障，这就是咱们今天演练的核心------做好 ε 序列到 ζ 序列的过渡，再逐步向更高层级突破。''
\end{quote}

他随即分配作战任务：

\begin{quote}
``唐三，你突破了 ε₀，擅长精准操控序数力量，负责主攻------用你最熟练的 ω 幂次迭代，配合 φ(1,0) 的不动点特性，先冲击屏障底层的 ζ₀ 纹路，记住，要遵循 \textbf{φ(2,0){[}n{]} = φ(1,0){[}n{]} = φ(1,φ(1,0)){[}n-1{]}} 的逻辑，先吃透 ζ₀ 的迭代节奏，再逐步提升层级，不可急于求成；

小舞，你的 ε₀ 气息凝练，速度快，负责侧翼牵制，用指尖的燃烧弹干扰屏障的序数流动，重点盯着 ζ₀ 纹路的波动，寻找迭代间隙；

牢大，你催动扽魂，借助其雄厚力量，将 ε\_ω 的波动注入扽魂之中，负责正面承压，用扽魂顶住屏障的反震，同时辅助唐三稳住 ζ₀ 的迭代基准，为他创造进攻窗口，切记不可让扽魂力量失控；

蔡须坤，你操控篮球，通过运球节奏校准自身气息，再以篮球为媒介，按照'嵌套对齐'原则，将唐三的 φ(1,0)、牢大扽魂的 ε\_ω 气息与 ζ₀ 的迭代节奏对齐，运球快慢对应迭代快慢，避免迭代紊乱，同时及时调整众人的气息节奏；

奥斯卡，你全程稳住全场波动，重点校准 ζ₀ 迭代相关的波动偏差，一旦发现唐三拆解 ζ₀ 纹路时出现气息卡顿，或牢大扽魂、蔡须坤篮球的气息出现偏差，立刻介入校准。''
\end{quote}

``明白！''五人齐声应答，迅速进入作战位置。唐三周身淡金色气息亮起，蓝银草灵气缠绕指尖；小舞身形紧绷，指尖燃烧弹蓄势待发；牢大低喝一声，肘间的扽魂虚影暴涨，化作一柄半丈高的巨锤，深黄色的 ε\_ω 波动如纹路般布满锤身，他双手紧握扽魂锤柄，稳稳伫立在前方，锤尖微微下沉，将地面压出细微凹陷，周身气流被扽魂的厚重力量牵引，形成一圈淡淡的气浪，时刻准备承接屏障反震；蔡须坤手中篮球轻轻弹跳，银白色气息与篮球表面的序数纹路完全相融，他双脚开立，运球节奏平稳而有力，篮球每一次落地、弹起，都精准对应着 ζ₀ 的迭代频率，带动周围气息形成有序的波动轨迹，为后续联动做好铺垫；奥斯卡手持香肠，淡蓝色气息弥漫全场，目光紧盯着三人的气息波动。五种气息相互呼应，在空地中形成一道有序的气场。

模拟作战正式开始。唐三率先发难，指尖凝聚起淡金色的序数力量，按照 \textbf{φ(2,0){[}n{]} = φ(1,n){[}n{]}} 的逻辑，先从 ε₀ 迭代过渡到 ζ₀ 层级，试图冲击屏障底层的 ζ₀ 纹路。可就在他的力量即将触碰到 ζ₀ 纹路时，指尖的气息突然出现卡顿，淡金色光束瞬间涣散，不仅没能冲击到屏障，反而被屏障反弹的 ζ₀ 气息震得后退两步，嘴角溢出一丝淡淡的灵气波动。

\begin{quote}
``不对！你太急于求成了！''赵无极的声音及时响起，``ζ₀ 作为 φ(2,0)，是 ε 序列到更高层级的过渡，你没有先稳固 ζ₀ 的迭代基准，直接用 ε₀ 的力量硬冲，自然会被 ζ₀ 的不动点特性反弹。记住，φ(2,0) = ζ₀ 是映射 β → φ(1,β) = ε\_β 的不动点，你要先让自身气息与 ζ₀ 的迭代节奏同频，按照 \textbf{ζ₀{[}n{]} = φ(1,n){[}n{]}} 的基本列，逐步渗透，而不是硬冲！''
\end{quote}

唐三立刻调整，深吸一口气，指尖气息重新凝聚，放缓节奏，先稳固自身 ε₀ 的迭代，再按照基本列逐步过渡到 ζ₀，这一次，淡金色气息稳稳贴合屏障底层的 ζ₀ 纹路，屏障表面的 ζ₀ 纹路出现轻微紊乱，唐三趁机发力，一道凝练的光束砸向纹路薄弱处，淡紫色的屏障终于出现了一道细微的裂痕。

就在众人以为局势好转时，意外突然发生------屏障的反震力量骤然增强，地面都随之微微震颤，牢大急于稳住局势，下意识低吼一声，双手发力将狠狠肘击地面，将自身 φ(1,ω) 的高阶力量强行注入扽魂之中，他的虚影瞬间暴涨半尺，深黄色波动变得狂暴，带着破空之声猛地砸向屏障。这股失控的力量瞬间打破了 ζ₀ 的迭代平衡，导致 ζ₀ 纹路剧烈扭曲、迭代节奏彻底紊乱。屏障瞬间爆发出一股强劲的反弹力量，如巨浪般狠狠肘击，``砰''的一声闷响，牢大连连后退数步，双脚在地面拖出两道深深的沟壑，扽魂虚影剧烈晃动，手上序数纹路变得黯淡，气息大幅紊乱，嘴角溢出一丝灵气；同时反弹力量波及到正在主攻的唐三，唐三的指尖气息瞬间溃散，胸口一阵发闷；蔡须坤手中的篮球也因气息紊乱，运球节奏彻底失控，篮球弹跳幅度忽大忽小，甚至偏离了原本的轨迹，砸向地面后弹起，银白色的联动气息瞬间出现断层，篮球表面的纹路也随之黯淡。

\begin{quote}
``住手！''赵无极沉声喝止，立刻抬手削弱了虚拟魔蛛的屏障力量，``牢大，你犯了一个致命错误！ζ₀ 是过渡层级，气息联动必须遵循'同频对齐'，你强行将高阶力量注入扽魂，让扽魂的厚重力量打破了 ζ₀ 的不动点平衡，不仅自身扽魂不稳、力量溃散，还拖累了全队！蔡须坤，你没能及时通过调整运球节奏校准紊乱的气息，篮球的联动媒介作用没有发挥出来------你本该通过加快或减慢运球速度，贴合 ζ₀ 的紊乱节奏，重新稳住联动气息，可你却任由篮球失控；奥斯卡，你专注于唐三的气息，忽略了 ζ₀ 纹路波动的偏差，也没能及时校准牢大扽魂和蔡须坤篮球的气息！''
\end{quote}

众人停下动作，神色都有些凝重。牢大收起扽魂虚影，满脸愧疚：

\begin{quote}
``对不起赵老师，我太急了，忘了您说的 ζ₀ 是过渡关键，也没控制好扽魂的力量，强行注入高阶气息才出了错。''
\end{quote}

蔡须坤接住失控弹跳的篮球，有些懊恼：

\begin{quote}
``我也有问题，气息紊乱时没能及时调整运球节奏，没能稳住联动的气息。''
\end{quote}

奥斯卡也纷纷低头，认同自己的失误。唐三揉了揉胸口，神色坚定：

\begin{quote}
``赵老师，我明白了，拆解 φ(4,0)，关键是先稳住 ζ₀ 的迭代节奏，不能急于突破，也不能让扽魂、篮球的力量失控。''
\end{quote}

赵无极点点头，语气缓和了几分：

\begin{quote}
``知道问题就好，这就是模拟作战的意义------提前暴露问题，才能在实战中避免致命失误。记住，ζ₀（φ(2,0)）是你们突破 φ(4,0)、甚至应对 φ(ω,0) 的核心过渡，吃透它的迭代逻辑，才能实现 ε 序列到更高层级 φ 函数的平稳衔接。牢大的扽魂力量雄厚，是正面承压的关键，但必须控制力量，与 ζ₀ 迭代同频；蔡须坤的篮球是联动气息的关键，要始终保持运球节奏与迭代节奏一致，才能发挥校准作用。刚才的挫折，就是因为你们没有真正重视 ζ₀ 的过渡作用，也没有用好各自的武魂优势。''
\end{quote}

他重新调整虚拟魔蛛的屏障，着重强化了 ζ₀ 纹路的稳定性：

\begin{quote}
``现在，重新开始。唐三，先专注于 ζ₀ 的迭代渗透，一步一步拆解底层纹路；牢大，催动扽魂，控制力量输出，只辅助稳住 ζ₀ 基准，让扽魂的波动与 ζ₀ 同频，不可强行叠加高阶力量；蔡须坤，操控篮球，保持平稳的运球节奏，以篮球为媒介，重点联动 ζ₀ 相关的气息，及时捕捉紊乱信号，用运球节奏校准；奥斯卡，全程紧盯 ζ₀ 纹路的波动，以及牢大扽魂、蔡须坤篮球的气息，一旦出现偏差，立刻校准。''
\end{quote}

五人重新调整状态，再次进入作战位置。这一次，众人都收敛了急躁的心，唐三缓缓发力，按照 ζ₀ 的基本列，逐步渗透屏障底层纹路；牢大低喝一声，扽魂虚影再次浮现，他双手轻握锤柄，刻意收敛力量，让扽魂锤身的深黄色波动变得平缓，与 ζ₀ 的迭代节奏保持同频，每一次屏障反震传来，他都通过微调扽魂的角度，将反震力量卸到地面，锤尖始终轻轻贴合地面，稳住 ζ₀ 的迭代基准，不再盲目发力；蔡须坤手中的篮球重新恢复平稳弹跳，他双眼紧盯着屏障上的 ζ₀ 纹路，运球节奏随着 ζ₀ 的迭代频率灵活调整，纹路波动加快时，他便加快运球速度，纹路波动放缓时，他便放缓节奏，银白色气息通过篮球的每一次弹跳，精准传递到唐三和牢大身上，及时调整气息衔接，将紊乱的苗头扼杀在摇篮里；奥斯卡紧盯波动，随时准备校准。小舞则抓住 ζ₀ 纹路波动的间隙，指尖燃烧弹精准砸向薄弱处，一步步扩大裂痕。

``咔嚓------''一声脆响，屏障底层的 ζ₀ 纹路彻底被拆解，失去了 ζ₀ 的支撑，φ(4,0) 的防御屏障瞬间弱化了大半。唐三眼神一凝，趁机将 ζ₀ 的力量与 φ(1,0) 叠加，借助蔡须坤篮球联动的合力，一道蕴含着双重迭代力量的金色光束，狠狠砸向屏障的裂痕；与此同时，牢大眼中精光一闪，双手发力，将扽魂锤高高举起，深黄色的 ε\_ω 波动与 ζ₀ 的残余力量相融，锤身纹路重新亮起，带着平稳而厚重的力量，狠狠砸在屏障之上------这一次，扽魂的力量精准可控，没有丝毫紊乱，与唐三的光束形成双重夹击。双重力量叠加之下，淡紫色的 φ(4,0) 防御屏障彻底碎裂，虚拟魔蛛的形态也随之消散。

五人同时松了口气，周身的气息缓缓收敛，牢大收起扽魂虚影，蔡须坤接住篮球，脸上都露出了疲惫却凝重的神情------刚才的挫折，让他们真切感受到了 ζ₀ 过渡的重要性，也明白了如何用好各自的武魂优势，更清楚实战中任何一个失误，都可能导致满盘皆输。唐三看向赵无极，语气郑重：

\begin{quote}
``赵老师，我们明白了，ζ₀ 就是突破高阶 φ 函数的关键，后续我们会重点打磨 ζ₀ 的迭代操控，也会练好扽魂、篮球的配合，避免再出现失误。''
\end{quote}

赵无极满意地点点头：

\begin{quote}
``很好，你们能吸取教训，就是最大的进步。刚才的挫折，只是小试牛刀，真正的千年人面魔蛛，不仅有 φ(ω,0) 的魂核力量，还能主动操控 ζ₀、η₀ 等序数力量反击，远比模拟作战复杂。接下来，我们重点强化 ζ₀ 的迭代操控、气息联动，以及牢大扽魂的力量控制------要练会让扽魂既能发挥厚重承压的优势，又能精准贴合 ζ₀ 迭代节奏，不盲目发力；还要强化蔡须坤篮球的节奏校准能力，让篮球不仅能联动气息，更能在气息紊乱时，通过灵活调整运球节奏，快速稳住局势，务必在实战前，彻底吃透 ζ₀ 的过渡作用，发挥出各自武魂的优势，避免再出现今天这样的失误！''
\end{quote}

话音落下，赵无极再次抬手，虚拟魔蛛的形态重新凝聚，这一次，屏障的 ζ₀ 纹路更加复杂，还加入了轻微的反击机制，显然是为了针对性强化众人对 ζ₀ 的掌控力，以及牢大扽魂、蔡须坤篮球的配合度。新一轮的模拟作战，正式拉开序幕，而这一次，众人的眼神中，多了几分沉稳与坚定------他们知道，今天的挫折，都是为了后续实战中，能稳稳接住人面魔蛛的每一次攻击，也能更好地发挥自身武魂的力量。

\chapter{第十二章：实战初捷，猎杀千年凤尾鸡冠蛇}\label{ux7b2cux5341ux4e8cux7ae0ux5b9eux6218ux521dux6377ux730eux6740ux5343ux5e74ux51e4ux5c3eux9e21ux51a0ux86c7}

唐三又经过几个星期修炼，终于掌握 \allowbreak{}\NovelInline{I00241}{\zeta_0}\allowbreak{} 以内的序数，\allowbreak{}\NovelInline{I00242}{\varepsilon}\allowbreak{} 序数每加一，就会去指数塔极限 \allowbreak{}\NovelInline{I00243}{\varepsilon_\omega}\allowbreak{} 到达 \allowbreak{}\NovelInline{I00244}{\varepsilon_{\varepsilon_0}}\allowbreak{} 的路，正如 \allowbreak{}\NovelInline{I00245}{\omega}\allowbreak{} 到达 \allowbreak{}\NovelInline{I00246}{\varepsilon_0}\allowbreak{}，只不过，每一次后继变成了取指数塔极限。\allowbreak{}\NovelInline{I00247}{f_{\zeta_0}(3) = f_{\varepsilon_{\varepsilon_0}}(3)}\allowbreak{}，依据是 \allowbreak{}\NovelInline{I00248}{\zeta_0 = \sup\{1, \varepsilon_0, \varepsilon_{\varepsilon_0}, \ldots\}}\allowbreak{}，之后变成

\NovelDisplay{M00149R01}{%
f_{\varepsilon_{\omega^{\omega^\omega}}}(3) = f_{\varepsilon_{\omega^{(\omega^2)\cdot 3 + \omega\cdot 2 + 3}}}(3) = f_{\varepsilon_{(\omega^{(\omega^2)\cdot 3 + \omega\cdot 2 + 2})\cdot 2 + (\omega^{\omega^2\cdot 3 + \omega\cdot 2 + 1})\cdot 2 + (\omega^{\omega^2\cdot 3 + \omega + 3})}}(3)%
}

掌握规律后，小舞也能展开 \allowbreak{}\NovelInline{I00249}{\zeta_0}\allowbreak{} 序数，牢大更是突破 \allowbreak{}\NovelInline{I00250}{\zeta_0}\allowbreak{} 的境界，到达了 \allowbreak{}\NovelInline{I00251}{\varphi(3,0)}\allowbreak{}，不仅仅这样，蔡须坤也到达 \allowbreak{}\NovelInline{I00252}{\zeta_\omega}\allowbreak{} 的境界，不过在这之前，奥斯卡总算是突破了 \allowbreak{}\NovelInline{I00253}{\varepsilon_0}\allowbreak{}，到达了 \allowbreak{}\NovelInline{I00254}{\varepsilon_{\varepsilon_0}}\allowbreak{} 的境界，很快就要突破 \allowbreak{}\NovelInline{I00255}{\zeta_0}\allowbreak{}，为了获得灵气，也为了真正检验序数修炼的成果，他们决定前往星斗大森林实战------那里不仅有丰富的灵气来源，更是千年人面魔蛛的栖息地，也是他们必须攻克的实战难关，只是谁也没想到，途中会先遭遇另一只强悍的千年魂兽。

出发前夜，赵无极将五人召集在林间空地，周身 \allowbreak{}\NovelInline{I00256}{\varphi(3,0)}\allowbreak{} 的序数波动比往日更加凝重，打破了夜晚的静谧：``星斗大森林不比演练场，里面的魂兽序数层级杂乱，从基础的 \allowbreak{}\NovelInline{I00257}{\omega}\allowbreak{} 层级到高阶的 \allowbreak{}\NovelInline{I00258}{\varphi}\allowbreak{} 函数层级都有，尤其是千年人面魔蛛，魂核稳定在 \allowbreak{}\NovelInline{I00259}{\varphi(\omega,0)}\allowbreak{}，比牢大的 \allowbreak{}\NovelInline{I00260}{\varphi(3,0)}\allowbreak{} 还要高出两个迭代维度，它的防御依旧是 \allowbreak{}\NovelInline{I00261}{\varphi(4,0)}\allowbreak{}，但实战中会主动运用 \allowbreak{}\NovelInline{I00262}{\zeta_0}\allowbreak{}、\allowbreak{}\NovelInline{I00263}{\eta_0}\allowbreak{} 的序数力量反击，甚至能拆解你们的气息联动，这是你们最大的威胁。

除此之外，森林中还有其他千年魂兽，比如千年凤尾鸡冠蛇，其序数层级在 \allowbreak{}\NovelInline{I00264}{\zeta_0}\allowbreak{} 到 \allowbreak{}\NovelInline{I00265}{\varphi(3,0)}\allowbreak{} 之间，防御约为 \allowbreak{}\NovelInline{I00266}{\varphi(2,0) = \zeta_0}\allowbreak{}，攻击性极强，速度极快，擅长用尾部的凤尾释放 \allowbreak{}\NovelInline{I00267}{\zeta_0}\allowbreak{} 级别的冲击波，更可怕的是它的速度突袭，快到能撕裂空气、扭曲低阶序数波动，你们若遇到，可将其作为猎杀魔蛛前的终极热身。''

他抬手在空中凝聚出星斗大森林的简易图谱，标注出魔蛛和凤尾鸡冠蛇常出没的区域，继续叮嘱：``记住，实战和演练最大的区别，是序数力量的动态博弈------无论是凤尾鸡冠蛇还是人面魔蛛，都会利用自身序数优势反击，甚至吸收你们溢出的 \allowbreak{}\NovelInline{I00268}{\varepsilon}\allowbreak{}、\allowbreak{}\NovelInline{I00269}{\zeta}\allowbreak{} 序数力量。你们必须牢记三个核心：

第一，牢大的 \allowbreak{}\NovelInline{I00270}{\varphi(3,0)}\allowbreak{} 是正面承压的核心，你的冰红茶扽魂不仅能增幅肘击力量，还能叠加团队属性，实战中要始终保持 \allowbreak{}\NovelInline{I00271}{\varphi(3,0)}\allowbreak{} 的稳定迭代，用扽魂的厚重气息顶住魂兽的防御反震，同时释放冰红茶的属性加成，让唐三、小舞的 \allowbreak{}\NovelInline{I00272}{\zeta_0}\allowbreak{} 力量更凝练，蔡须坤的联动更稳定；

第二，蔡须坤的扽魂篮球，要利用 \allowbreak{}\NovelInline{I00273}{\zeta_\omega}\allowbreak{} 的迭代优势，通过运球节奏校准全队气息，尤其是唐三拆解防御时，一旦出现 \allowbreak{}\NovelInline{I00274}{\zeta_0}\allowbreak{}、\allowbreak{}\NovelInline{I00275}{\varphi(1,0)}\allowbreak{} 的迭代紊乱，你要立刻调整运球频率，用篮球的序数波动衔接断层，同时可以借助篮球的弹跳传递联动力量，辅助唐三突破防御；

第三，奥斯卡，你的香肠扽魂不再只是稳定波动，突破 \allowbreak{}\NovelInline{I00276}{\varepsilon_{\varepsilon_0}}\allowbreak{} 后，香肠不仅能挂载高阶序数增益，还能发射出去造成 \allowbreak{}\NovelInline{I00277}{\varphi(1,0)}\allowbreak{} 级别的序数冲击，实战中既要及时给众人补充增益，也要找准时机，用香肠冲击魂兽的防御薄弱处，配合小舞的侧翼牵制。''

赵无极顿了顿，目光扫过五人，语气愈发郑重：``还有序数运用的关键------唐三，你掌握了 \allowbreak{}\NovelInline{I00278}{\zeta_0}\allowbreak{} 以内的序数，拆解防御时，要遵循'从 \allowbreak{}\NovelInline{I00279}{\zeta_0}\allowbreak{} 底层入手，逐步迭代至 \allowbreak{}\NovelInline{I00280}{\varphi(2,0)}\allowbreak{}、\allowbreak{}\NovelInline{I00281}{\varphi(3,0)}\allowbreak{}'的逻辑，利用 \allowbreak{}\NovelInline{I00282}{f_{\zeta_0}}\allowbreak{} 的迭代规律，将自身力量与牢大的 \allowbreak{}\NovelInline{I00283}{\varphi(3,0)}\allowbreak{}、蔡须坤的 \allowbreak{}\NovelInline{I00284}{\zeta_\omega}\allowbreak{} 联动，避免单独发力被魂兽吸收；

小舞，你的 \allowbreak{}\NovelInline{I00285}{\zeta_0}\allowbreak{} 力量凝练，速度快，重点用燃烧弹干扰魂兽的序数流动，尤其是它们调动高阶力量时，用 \allowbreak{}\NovelInline{I00286}{\varepsilon_{\varepsilon_0}}\allowbreak{} 级别的波动打断其迭代节奏。对付凤尾鸡冠蛇，重点拆解其 \allowbreak{}\NovelInline{I00287}{\zeta_0}\allowbreak{} 防御，还要提防它的速度突袭和尾部冲击波------它的突袭快到能规避你们的气息感知，务必时刻警惕；对付人面魔蛛，则要攻克 \allowbreak{}\NovelInline{I00288}{\varphi(4,0)}\allowbreak{} 防御，记住，两者的核心差异的是迭代维度，不可用同一套战术。''

五人齐声应道，眼底满是坚定。牢大握紧拳头，肘间泛起浓郁的冰蓝色光晕，冰红茶扽魂虚影缓缓浮现------不同于演练时的浑厚，此刻的扽魂被 \allowbreak{}\NovelInline{I00289}{\varphi(3,0)}\allowbreak{} 的序数纹路包裹，液体状的冰红茶虚影中夹杂着细密的 \allowbreak{}\NovelInline{I00290}{\eta_0}\allowbreak{} 波动，轻轻晃动便散发着刺骨的寒气，``赵老师放心，我的冰红茶扽魂不仅能扛住 \allowbreak{}\NovelInline{I00291}{\varphi(4,0)}\allowbreak{} 的反震，对付凤尾鸡冠蛇的 \allowbreak{}\NovelInline{I00292}{\zeta_0}\allowbreak{} 防御更是不在话下，肘击时还能叠加 \allowbreak{}\NovelInline{I00293}{\varphi(3,0)}\allowbreak{} 的力量，同时给大家加上冰属性增益，提升力量凝练度，就算它速度再快，我也能及时筑起防御。''

蔡须坤指尖转动，扽魂篮球跃然掌心，篮球表面流转着 \allowbreak{}\NovelInline{I00294}{\zeta_\omega}\allowbreak{} 的序数纹路，银白色气息与篮球完美相融，``我会全程保持运球节奏，校准大家的气息，无论是对付凤尾鸡冠蛇还是人面魔蛛，都能及时衔接紊乱的波动，尤其是它突袭时，我会用篮球波动干扰它的速度节奏，辅助唐三突破防御。''奥斯卡则掏出一根香肠，香肠表面萦绕着 \allowbreak{}\NovelInline{I00295}{\varepsilon_{\varepsilon_0}}\allowbreak{} 的淡蓝色波动，轻轻一抛，香肠便发射出一道细小的光束，击中地面泛起一阵涟漪，``我的香肠既能增益，也能攻击，对付凤尾鸡冠蛇时，刚好能发挥 \allowbreak{}\NovelInline{I00296}{\varphi(1,0)}\allowbreak{} 级别的冲击优势，就算它突袭过来，也能及时用香肠攻击逼退它，为后面打魔蛛热身。''

次日清晨，五人踏入星斗大森林。林间古木参天，灵气浓郁却夹杂着杂乱的序数波动，低阶魂兽的 \allowbreak{}\NovelInline{I00297}{\omega}\allowbreak{}、\allowbreak{}\NovelInline{I00298}{\omega^2}\allowbreak{} 层级气息四处弥漫，偶尔还能感受到高阶魂兽的 \allowbreak{}\NovelInline{I00299}{\varepsilon}\allowbreak{}、\allowbreak{}\NovelInline{I00300}{\zeta}\allowbreak{} 序数波动，让人不敢有丝毫懈怠。

唐三走在最前方，蓝银草轻轻舒展，藤条上的 \allowbreak{}\NovelInline{I00301}{\zeta_0}\allowbreak{} 序数纹路不断闪烁，感知着周围的气息变化；

小舞紧随其后，指尖燃烧弹蓄势待发，\allowbreak{}\NovelInline{I00302}{\varepsilon_{\varepsilon_0}}\allowbreak{} 的波动时刻警惕着四周；牢大走在中间，冰红茶扽魂虚影悬浮在肘间，\allowbreak{}\NovelInline{I00303}{\varphi(3,0)}\allowbreak{} 的厚重气息扩散开来，为众人筑起一道防御屏障；蔡须坤手持篮球，运球节奏平稳，\allowbreak{}\NovelInline{I00304}{\zeta_\omega}\allowbreak{} 的波动与众人的气息保持同频，时刻准备校准；奥斯卡则走在最后，手中握着数根香肠，目光紧盯着前方，随时准备释放增益或攻击。

前行约莫半个时辰，林间的序数波动愈发杂乱，唐三的蓝银草开始轻微震颤，藤条上的 \allowbreak{}\NovelInline{I00305}{\zeta_0}\allowbreak{} 纹路泛起淡淡的警示光晕：``大家小心，周围有低阶魂兽聚集，序数波动在 \allowbreak{}\NovelInline{I00306}{\omega^2}\allowbreak{} 到 \allowbreak{}\NovelInline{I00307}{\varepsilon_0}\allowbreak{} 之间，正好趁机练手，熟悉实战中序数力量的运用，也能收集少量灵气。''话音刚落，三道黑影从两侧的草丛中窜出，正是三只百年幽冥狼------它们身形矫健，体表萦绕着 \allowbreak{}\NovelInline{I00308}{\omega^2}\allowbreak{} 级别的序数波动，双眼泛着凶光，朝着五人扑来。

``正好热身！''牢大低喝一声，肘间的冰红茶扽魂虚影微微亮起，\allowbreak{}\NovelInline{I00309}{\varphi(3,0)}\allowbreak{} 的厚重气息收敛至 \allowbreak{}\NovelInline{I00310}{\omega^2}\allowbreak{} 层级，既不浪费力量，又能精准压制幽冥狼，``奥斯卡，给我加个增益，小舞你侧翼牵制，蔡须坤校准节奏！''奥斯卡立刻抛出一根香肠，淡蓝色的 \allowbreak{}\NovelInline{I00311}{\varepsilon_{\varepsilon_0}}\allowbreak{} 波动落在牢大身上，瞬间提升了他的力量凝练度：``增益生效，放心打！''

第一只幽冥狼率先扑来，\allowbreak{}\NovelInline{I00312}{\omega^2}\allowbreak{} 级别的冲击带着腥风砸向牢大。牢大侧身避开，肘间冰红茶扽魂顺势发力，一记轻缓却有力的肘击砸在幽冥狼的脊背，冰蓝色的 \allowbreak{}\NovelInline{I00313}{\omega^2}\allowbreak{} 波动瞬间炸开，``砰''的一声，幽冥狼被砸得踉跄后退，体表的序数波动紊乱。``蔡须坤，联动我的力量，传递给唐三！''蔡须坤立刻调整运球节奏，篮球弹跳频率贴合牢大的波动，银白色的 \allowbreak{}\NovelInline{I00314}{\zeta_\omega}\allowbreak{} 波动将牢大的冰属性力量传递给唐三，``节奏对齐，力量已传递！''

唐三眼神一凝，蓝银草藤条窜出，\allowbreak{}\NovelInline{I00315}{\zeta_0}\allowbreak{} 力量压制至 \allowbreak{}\NovelInline{I00316}{\omega^\omega}\allowbreak{} 层级，精准缠绕住幽冥狼的四肢，按照 \allowbreak{}\NovelInline{I00317}{\varepsilon_0}\allowbreak{} 基本列的迭代逻辑，逐步收紧力量：``小舞，动手！''小舞身形一闪，瞬移至幽冥狼身前，指尖燃烧弹凝聚 \allowbreak{}\NovelInline{I00318}{\varepsilon_0(1)}\allowbreak{} 级别的波动，狠狠砸在幽冥狼的头部，``嘭''的一声，幽冥狼的 \allowbreak{}\NovelInline{I00319}{\omega^2}\allowbreak{} 波动彻底溃散，倒在地上化作少量灵气。

与此同时，另外两只幽冥狼同时发起攻击，一只扑向奥斯卡，一只扑向蔡须坤。奥斯卡反应极快，掏出两根香肠，一根挂载增益给自己，一根瞬间发射，淡蓝色的 \allowbreak{}\NovelInline{I00320}{\varepsilon_{\varepsilon_0}}\allowbreak{} 波动击中扑来的幽冥狼，打断其攻击节奏：``想偷袭我？没那么容易！''蔡须坤则运球侧身，篮球带着 \allowbreak{}\NovelInline{I00321}{\zeta_\omega}\allowbreak{} 的低阶波动，狠狠砸向幽冥狼的腿部，既不暴露高阶力量，又能限制其行动，``唐三，左边交给你！''

唐三立刻调动蓝银草，藤条分叉，一根缠绕住幽冥狼的腿部，一根凝聚 \allowbreak{}\NovelInline{I00322}{\omega^\omega}\allowbreak{} 力量，轻轻一拉，便将幽冥狼绊倒。牢大趁机上前，肘间冰红茶扽魂发力，\allowbreak{}\NovelInline{I00323}{\varphi(3,0)}\allowbreak{} 力量短暂释放，一记肘击砸在幽冥狼的核心，瞬间击溃其序数波动：``搞定一只！''最后一只幽冥狼见同伴接连被击败，试图逃窜，小舞瞬移追上，指尖燃烧弹叠加 \allowbreak{}\NovelInline{I00324}{\varepsilon_0(2)}\allowbreak{} 力量，精准击中其尾部，幽冥狼踉跄倒地，被唐三的蓝银草彻底束缚，最终失去力量，化作灵气。

五人短暂调息，吸收着幽冥狼残留的灵气。牢大收起冰红茶扽魂，笑道：``这些小怪虽然弱，但实战中调整力量层级很关键，还好有冰红茶的属性加成，力量收放更自如。''蔡须坤擦了擦篮球上的灰尘，说道：``刚才运球校准节奏，正好练了低阶波动的联动，为后面打魂兽热身。''奥斯卡晃了晃手中的香肠，得意道：``我的香肠攻击也很给力，轻松打断小怪的攻击！''

唐三点头，蓝银草再次舒展，藤条上的 \allowbreak{}\NovelInline{I00325}{\zeta_0}\allowbreak{} 纹路恢复警示状态：``大家收敛气息，刚才的战斗可能惊动了周围的高阶魂兽，我们继续前行，注意警惕。''众人立刻调整状态，恢复之前的阵型，继续向星斗大森林深处迈进。又前行约莫一炷香的时间，唐三的蓝银草突然剧烈震颤，藤条上的 \allowbreak{}\NovelInline{I00326}{\zeta_0}\allowbreak{} 纹路瞬间变得黯淡，语气凝重：``小心，前方有高阶序数波动，是 \allowbreak{}\NovelInline{I00327}{\varphi(2,0) = \zeta_0}\allowbreak{}，不是人面魔蛛，应该是赵无极老师提到的千年凤尾鸡冠蛇！''

话音刚落，一阵尖锐的嘶鸣响彻林间，几乎没有预兆，一道七彩残影便如离弦之箭般撕裂空气，带着刺耳的破空声从树梢窜出------速度快到扭曲了周围的灵气，连唐三的蓝银草感知都出现了刹那的延迟，它稳稳落在空地中央，周身的 \allowbreak{}\NovelInline{I00328}{\zeta_0}\allowbreak{} 波动瞬间炸开，正是千年凤尾鸡冠蛇。它身形修长，约莫丈许长，通体覆盖着流光溢彩的七彩鳞片，阳光洒下时折射出耀眼的光晕，头顶的鸡冠鲜红如血，微微凸起，散发着浓郁的 \allowbreak{}\NovelInline{I00329}{\zeta_0}\allowbreak{} 级序数波动，尾端的凤尾展开如扇形，每一根尾羽都带着细密的序数纹路，双眼泛着锐利的金芒，周身灵气波动比幽冥狼浓郁数倍，压迫感远超刚才的小怪，刚一落地，地面便被其 \allowbreak{}\NovelInline{I00330}{\zeta_0}\allowbreak{} 波动震出细密的裂纹，而那撕裂空气的余劲，依旧在林间回荡。

``太好了，正好用它练手，熟悉 \allowbreak{}\NovelInline{I00331}{\zeta_0}\allowbreak{} 防御的拆解技巧，也摸清千年魂兽的攻击套路，尤其是它的速度突袭，为后面打人面魔蛛做准备！''唐三眼神一亮，立刻调整战术，声音沉稳有力，``按计划配合，牢大正面承压，务必顶住它的冲击和尾部冲击波，重点提防它的速度突袭，一旦它发动突袭，立刻展开 \allowbreak{}\NovelInline{I00332}{\varphi(3,0)}\allowbreak{} 防御；蔡须坤校准节奏，重点衔接我和牢大的力量，避免气息紊乱，它突袭时，用篮球波动干扰它的速度迭代；小舞你利用速度优势，侧翼牵制，专打它的七寸和鸡冠薄弱处，同时规避它的突袭；奥斯卡负责全程增益，同时用香肠攻击干扰它的防御迭代，若它突袭奥斯卡，小舞立刻支援！''

``动手！''唐三一声令下，率先发难。蓝银草如锋利的箭矢般窜出，数十根藤条交错缠绕，藤条上的 \allowbreak{}\NovelInline{I00333}{\zeta_0}\allowbreak{} 纹路通体发亮，按照 \allowbreak{}\NovelInline{I00334}{f_{\zeta_0}}\allowbreak{} 的迭代规律，从 \allowbreak{}\NovelInline{I00335}{\zeta_0}\allowbreak{} 底层纹路开始，层层递进，缓缓向凤尾鸡冠蛇的防御屏障逼近，``蔡须坤，校准迭代节奏，配合我渗透屏障，它的防御是 \allowbreak{}\NovelInline{I00336}{\zeta_0}\allowbreak{} 层级，迭代速度极快，而且速度惊人，必须跟上它的节奏，防止它突然突袭！''

蔡须坤立刻回应，运球节奏瞬间加快，篮球每一次落地、弹起都发出清脆的声响，精准对应着凤尾鸡冠蛇防御屏障的 \allowbreak{}\NovelInline{I00337}{\zeta_0}\allowbreak{} 迭代频率，银白色的 \allowbreak{}\NovelInline{I00338}{\zeta_\omega}\allowbreak{} 波动顺着篮球的弹跳，源源不断地传递到唐三的蓝银草上，完美填补了力量衔接的断层，``收到！节奏已校准，与它的 \allowbreak{}\NovelInline{I00339}{\zeta_0}\allowbreak{} 迭代同频，我会时刻警惕它的突袭，一旦有异动，立刻干扰！''与此同时，小舞身形一闪，化作一道残影瞬移到凤尾鸡冠蛇侧面，指尖燃烧弹凝聚起浓郁的 \allowbreak{}\NovelInline{I00340}{\varepsilon_{\varepsilon_0}}\allowbreak{} 波动，毫不犹豫地砸向屏障的薄弱处------蛇身与尾部衔接的位置，``唐三，我打乱它的 \allowbreak{}\NovelInline{I00341}{\zeta_0}\allowbreak{} 迭代了，快趁机突破！''

``嘶------''凤尾鸡冠蛇被小舞的攻击激怒，发出一声尖锐的嘶鸣，尾部猛地一挥，一道凝聚着纯粹 \allowbreak{}\NovelInline{I00342}{\zeta_0}\allowbreak{} 力量的金色冲击波，带着刺耳的破空声，朝着小舞狠狠砸去。就在冲击波即将击中小舞的瞬间，凤尾鸡冠蛇身形骤然一动，七彩残影再次撕裂空气，竟是借着冲击波的掩护，发动了速度突袭，目标直指防御最薄弱的奥斯卡！

它的速度比刚才窜出时更快，周身的 \allowbreak{}\NovelInline{I00343}{\zeta_0}\allowbreak{} 波动被压缩成一道细密的光带，几乎能规避所有低阶感知，转瞬之间便冲到了奥斯卡身前，蛇口微张，带着 \allowbreak{}\NovelInline{I00344}{\zeta_0}\allowbreak{} 波动的獠牙闪着寒光。

``小心！''牢大低喝一声，身形如惊雷般一闪，\allowbreak{}\NovelInline{I00345}{\varphi(3,0)}\allowbreak{} 的力量瞬间爆发，冰红茶扽魂虚影暴涨，堪堪挡在奥斯卡身前，``砰''的一声闷响，冲击波狠狠撞在屏障上，瞬间溃散，而凤尾鸡冠蛇的头部也狠狠撞在屏障上，冰蓝色的 \allowbreak{}\NovelInline{I00346}{\varphi(3,0)}\allowbreak{} 波动与七彩的 \allowbreak{}\NovelInline{I00347}{\zeta_0}\allowbreak{} 波动剧烈碰撞，发出刺耳的滋滋声，牢大脚下的地面被震出深深的凹陷，身体微微后退两步，却依旧稳稳顶住了这一击。

``我的冰红茶增益全力生效，大家全力进攻，它的冲击波虽然强，但更要提防它借冲击波掩护的速度突袭，太快了！''淡蓝色的冰属性光晕瞬间笼罩着五人，唐三的 \allowbreak{}\NovelInline{I00348}{\zeta_0}\allowbreak{} 力量愈发凝练，小舞的速度也提升了几分，蔡须坤的联动也更加稳定。

凤尾鸡冠蛇见突袭未果，眼中凶光更盛，体表的 \allowbreak{}\NovelInline{I00349}{\zeta_0}\allowbreak{} 防御屏障瞬间亮起刺眼的金光，开始以极快的速度迭代力量，屏障的厚度不断增加，同时它的身体微微弓起，头部的鸡冠泛起红光，尾部凤尾轻轻颤动，显然是在积蓄力量，准备发动更迅猛的速度突袭。``它在强化 \allowbreak{}\NovelInline{I00350}{\zeta_0}\allowbreak{} 防御，同时在蓄力突袭，速度会更快，奥斯卡，立刻发射香肠攻击，打断它的迭代节奏和蓄力！

``唐三目光锐利，瞬间察觉其意图，迅速调整自身力量，将 \allowbreak{}\NovelInline{I00351}{\zeta_0}\allowbreak{} 力量与 \allowbreak{}\NovelInline{I00352}{\varphi(1,0)}\allowbreak{} 叠加，按照 \allowbreak{}\NovelInline{I00353}{\zeta_0 = \sup\{1, \varepsilon_0, \varepsilon_{\varepsilon_0}, \ldots\}}\allowbreak{} 的规律，快速迭代力量，让蓝银草的渗透速度更快，''牢大，辅助我压制它的力量，用冰红茶的寒气减缓它的迭代速度和突袭准备；小舞，绕到它身后，牵制它的动作，防止它再次突袭奥斯卡！''

奥斯卡不敢怠慢，立刻掏出三根香肠，指尖发力，香肠如三颗炮弹般射向凤尾鸡冠蛇的头部、鸡冠和尾部，淡蓝色的 \allowbreak{}\NovelInline{I00354}{\varepsilon_{\varepsilon_0}}\allowbreak{} 波动瞬间炸开，狠狠击中蛇身，``嘭嘭嘭''三声闷响，凤尾鸡冠蛇发出一声痛苦的嘶鸣，强化防御的动作瞬间停滞，体表的 \allowbreak{}\NovelInline{I00355}{\zeta_0}\allowbreak{} 屏障波动也出现了明显的紊乱，鳞片上的光泽也黯淡了几分，蓄力突袭的动作也被打断。

``有效！我继续攻击，专打它的薄弱处，打断它的突袭准备！''奥斯卡乘胜追击，不断掏出香肠发射，每一根香肠都精准击中蛇身的薄弱处------鸡冠下方、七寸位置和尾羽根部，每一次击中都能造成 \allowbreak{}\NovelInline{I00356}{\varphi(1,0)}\allowbreak{} 级别的冲击，进一步瓦解着 \allowbreak{}\NovelInline{I00357}{\zeta_0}\allowbreak{} 防御屏障，让其迭代节奏愈发紊乱，更无法顺利蓄力发动突袭。

蔡须坤抓住这个关键间隙，运球节奏陡然加快，篮球带着强劲的 \allowbreak{}\NovelInline{I00358}{\zeta_\omega}\allowbreak{} 力量，狠狠砸向凤尾鸡冠蛇防御屏障的 \allowbreak{}\NovelInline{I00359}{\zeta_0}\allowbreak{} 底层纹路，``唐三，我帮你打破 \allowbreak{}\NovelInline{I00360}{\zeta_0}\allowbreak{} 底层支撑，你趁机渗透！另外我会用篮球波动在周围布下牵制网，就算它发动突袭，也能减缓它的速度！''篮球撞击屏障的瞬间，银白色的波动剧烈炸开，屏障上的 \allowbreak{}\NovelInline{I00361}{\zeta_0}\allowbreak{} 纹路瞬间扭曲、断裂，迭代节奏彻底被打乱，原本厚重的屏障也变得薄弱了几分，同时，篮球弹出的银白色波动在空地周围形成一道细密的波动网，死死锁定着凤尾鸡冠蛇的身形，限制其速度爆发。

唐三眼神一凝，抓住这个千载难逢的间隙，将自身 \allowbreak{}\NovelInline{I00362}{\zeta_0}\allowbreak{} 力量与牢大的 \allowbreak{}\NovelInline{I00363}{\varphi(3,0)}\allowbreak{} 力量通过蔡须坤的篮球紧紧联动，两道力量在蓝银草上叠加，形成一道蕴含着双重高阶序数力量的金色光束，狠狠砸向屏障的裂痕处，``就是现在！突破它的防御！''

``喝！''牢大低吼一声，将 \allowbreak{}\NovelInline{I00364}{\varphi(3,0)}\allowbreak{} 的力量全部爆发，肘间的冰红茶扽魂瞬间暴涨，冰蓝色的光晕几乎笼罩了半个空地，一记势大力沉的肘击，带着刺骨的寒气，狠狠砸在屏障的裂痕上，冰蓝色的波动与金色光束相互叠加，双重力量瞬间撕裂了 \allowbreak{}\NovelInline{I00365}{\zeta_0}\allowbreak{} 的防御屏障，``咔嚓''一声脆响，屏障彻底碎裂，化作漫天光点消散。

凤尾鸡冠蛇发出一声凄厉的嘶鸣，体表的屏障碎裂后，体内的 \allowbreak{}\NovelInline{I00366}{\zeta_0}\allowbreak{} 序数波动彻底暴露在外，却依旧狂暴，它见状，不再蓄力，而是瞬间发动终极速度突袭------七彩身形化作一道模糊的光带，无视蔡须坤布下的波动网，硬生生撕裂一道缺口，朝着唐三猛冲而去，它知道唐三是团队的核心，只要击溃唐三，就能瓦解五人的配合，其速度之快，让唐三的蓝银草都来不及完全反应，藤条刚要缠绕，它便已冲到近前，尾部凤尾带着 \allowbreak{}\NovelInline{I00367}{\zeta_0}\allowbreak{} 波动，狠狠抽向唐三的胸口。

``唐三小心！''小舞尖叫一声，身形瞬移到唐三身前，指尖燃烧弹全力爆发，\allowbreak{}\NovelInline{I00368}{\varepsilon_{\varepsilon_0}}\allowbreak{} 的波动狠狠砸向凤尾鸡冠蛇的头部，试图逼退它；蔡须坤也立刻调整运球节奏，篮球带着 \allowbreak{}\NovelInline{I00369}{\zeta_\omega}\allowbreak{} 的全力波动，狠狠砸向凤尾鸡冠蛇的腹部，干扰其突袭轨迹；牢大更是瞬间冲了过来，肘间冰红茶扽魂再次发力，挡在小舞和唐三身前。``小舞，牵制它的动作，别让它靠近奥斯卡！蔡须坤，校准气息，强化波动网，彻底限制它的速度！

``唐三一边快速后退，一边催动蓝银草，数十根藤条紧紧缠绕住凤尾鸡冠蛇的身体，尤其是它的四肢和尾部，\allowbreak{}\NovelInline{I00370}{\zeta_0}\allowbreak{} 的力量不断收紧，试图限制它的行动，同时藤条上的序数纹路不断闪烁，持续压制着它的力量迭代和速度爆发。小舞身形灵活，不断瞬移，避开凤尾鸡冠蛇的冲击波，同时指尖燃烧弹不断砸向它的七寸处，\allowbreak{}\NovelInline{I00371}{\varepsilon_{\varepsilon_0}}\allowbreak{} 的波动一次次干扰着它的序数迭代，让它的力量难以完全爆发，速度也变得迟缓了几分；蔡须坤则将运球节奏调整到极致，篮球不断弹跳，银白色的波动网变得愈发细密，死死缠绕住凤尾鸡冠蛇的身形，每一次篮球撞击，都能让它的速度出现刹那的停滞。

牢大则继续正面承压，迎着凤尾鸡冠蛇的冲击，肘间的冰红茶扽魂一次次砸在蛇身的鳞片上，每一击都带着 \allowbreak{}\NovelInline{I00372}{\varphi(3,0)}\allowbreak{} 的力量和冰属性的刺骨寒气，砸得鳞片飞溅，蛇身不断颤抖，原本流光溢彩的七彩鳞片，此刻已经变得黯淡无光，多处出现破损，速度也越来越迟缓------刚才的终极突袭消耗了它大量力量，再加上众人的牵制，它的速度已大不如前。

``奥斯卡，给我和唐三补充增益，我要发动全力一击，彻底解决它，不让它再发动突袭！''牢大喊道，声音中带着一丝疲惫，却依旧充满力量。奥斯卡立刻发射两根增益香肠，分别落在唐三和牢大身上，淡蓝色的 \allowbreak{}\NovelInline{I00373}{\varepsilon_{\varepsilon_0}}\allowbreak{} 波动瞬间提升了两人的力量层级，唐三的 \allowbreak{}\NovelInline{I00374}{\zeta_0}\allowbreak{} 力量逼近 \allowbreak{}\NovelInline{I00375}{\varepsilon_{\varepsilon_0}}\allowbreak{}，牢大的 \allowbreak{}\NovelInline{I00376}{\varphi(3,0)}\allowbreak{} 力量也愈发浑厚，冰红茶扽魂的光晕也变得更加浓郁。

蔡须坤将运球节奏调整到与凤尾鸡冠蛇体内的序数迭代频率相反，篮球带着强劲的 \allowbreak{}\NovelInline{I00377}{\zeta_\omega}\allowbreak{} 力量，一次次砸向它的七寸和魂核位置，``打断它的力量迭代，让它无法再发动速度突袭！''篮球撞击的瞬间，银白色的波动炸开，凤尾鸡冠蛇的序数波动剧烈紊乱，\allowbreak{}\NovelInline{I00378}{\zeta_0}\allowbreak{} 级别的力量难以稳定释放，身体也开始不受控制地抽搐，速度彻底被限制。唐三抓住这个机会，催动全身力量，蓝银草紧紧缠绕住凤尾鸡冠蛇的七寸，\allowbreak{}\NovelInline{I00379}{\zeta_0}\allowbreak{} 与 \allowbreak{}\NovelInline{I00380}{\varepsilon_{\varepsilon_0}}\allowbreak{} 的力量叠加，死死锁住它的要害，``牢大，就是现在，全力肘击七寸，彻底击溃它的魂核！''

牢大应声而动，身形一闪，来到凤尾鸡冠蛇的七寸前方，肘间的冰红茶扽魂凝聚了全身的 \allowbreak{}\NovelInline{I00381}{\varphi(3,0)}\allowbreak{} 力量，冰蓝色的光晕耀眼夺目，一记势大力沉的肘击，带着毁天灭地的力量，狠狠砸在凤尾鸡冠蛇的七寸处。``砰------''一声巨响，震得整个林间都在颤抖，凤尾鸡冠蛇发出最后一声凄厉的嘶鸣，体内的 \allowbreak{}\NovelInline{I00382}{\zeta_0}\allowbreak{} 序数波动彻底溃散，魂核也被这一击击碎，身体软软地倒在地上，不再动弹，随后缓缓化作浓郁的灵气，弥漫在整个空地中，七彩鳞片也逐渐失去光泽，散落一地。

五人同时松了口气，周身的气息缓缓收敛，纷纷瘫坐在地上，脸上都露出了疲惫却欣喜的神情，额头上布满了汗水，身上的衣物也有些破损，显然这场战斗消耗了他们大量的力量------尤其是应对凤尾鸡冠蛇的速度突袭，几乎全程紧绷，不敢有丝毫松懈。牢大收起冰红茶扽魂，肘间的 \allowbreak{}\NovelInline{I00383}{\varphi(3,0)}\allowbreak{} 纹路依旧清晰，他擦了擦额头的汗水，笑道：``这凤尾鸡冠蛇果然名不虚传，\allowbreak{}\NovelInline{I00384}{\zeta_0}\allowbreak{} 防御又快又强，还擅长速度突袭，刚才那几次突袭，差点就防不住了，还好我们配合默契，尤其是冰红茶的加成，让我能及时筑起防御，刚好练手，为后面打魔蛛铺垫好了。''

蔡须坤接住篮球，篮球表面的 \allowbreak{}\NovelInline{I00385}{\zeta_\omega}\allowbreak{} 纹路微微黯淡，他揉了揉发酸的手臂，说道：``刚才校准节奏，既要跟上它的 \allowbreak{}\NovelInline{I00386}{\zeta_0}\allowbreak{} 迭代频率，又要在它强化防御、发动突袭时及时调整，尤其是布下波动网牵制它的速度，刚好熟悉了 \allowbreak{}\NovelInline{I00387}{\zeta_0}\allowbreak{} 防御的拆解逻辑，也练熟了应对高速突袭的技巧，后面对付人面魔蛛的 \allowbreak{}\NovelInline{I00388}{\varphi(4,0)}\allowbreak{} 防御，也能更顺手。''奥斯卡则晃了晃手中剩余的香肠，笑着说：``我的香肠攻击也派上了大用场，拆解 \allowbreak{}\NovelInline{I00389}{\zeta_0}\allowbreak{} 防御时，刚好能打断它的迭代，每次击中它的薄弱处，都能看到它的波动紊乱，尤其是在它蓄力突袭时，及时用香肠攻击打断，不然真的会被它偷袭到，下次争取突破 \allowbreak{}\NovelInline{I00390}{\zeta_0}\allowbreak{}，能发射更强的攻击对付魔蛛。''

小舞靠在唐三身边，指尖的燃烧弹缓缓收敛，脸上带着一丝疲惫，却依旧笑容明媚：``是啊，实战比演练更复杂，这只凤尾鸡冠蛇速度太快了，刚才它突袭唐三的时候，我都快吓死了，还好有牢大挡在前面，而且我们的配合也越来越默契了，后面打魔蛛，就能更有底气了。''

唐三看着弥漫在空地上的浓郁灵气，眼底满是领悟，他缓缓站起身，蓝银草轻轻舒展，吸收着周围的灵气，说道：``刚才拆解 \allowbreak{}\NovelInline{I00391}{\zeta_0}\allowbreak{} 防御时，我更熟练了 \allowbreak{}\NovelInline{I00392}{f_{\zeta_0}}\allowbreak{} 的迭代规律，也明白了不同魂兽的防御特性------凤尾鸡冠蛇的 \allowbreak{}\NovelInline{I00393}{\zeta_0}\allowbreak{} 防御侧重迭代速度，攻击擅长速度突袭和冲击波，尤其是它的突袭，能借外力掩护、压缩自身波动，很难防备，而人面魔蛛的 \allowbreak{}\NovelInline{I00394}{\varphi(4,0)}\allowbreak{} 防御侧重层级叠加，还能吸收我们的力量，后面对付魔蛛，我们要调整战术，重点拆解其 \allowbreak{}\NovelInline{I00395}{\varphi(2,0)}\allowbreak{} 的底层支撑，同时提防它的力量吸收能力，也要做好应对高速攻击的准备。''

``大家先调息，吸收这只千年凤尾鸡冠蛇的灵气，巩固自身境界，恢复消耗的力量。''唐三继续说道，``它的灵气非常浓郁，刚好能帮我们提升实力，为后面面对人面魔蛛做好充分准备。''众人纷纷点头，不再说话，闭上眼睛，开始吸收周围的灵气。浓郁的灵气如潮水般涌向五人，与他们体内的序数波动相互交织，形成一道道耀眼的光晕，林间的灵气不断汇聚，围绕在五人身边，唐三体内的 \allowbreak{}\NovelInline{I00396}{\zeta_0}\allowbreak{} 力量愈发凝练，朝着 \allowbreak{}\NovelInline{I00397}{\varepsilon_{\varepsilon_0}}\allowbreak{} 稳步迈进；牢大的 \allowbreak{}\NovelInline{I00398}{\varphi(3,0)}\allowbreak{} 境界更加稳固，冰红茶扽魂的力量也更加强劲；蔡须坤的 \allowbreak{}\NovelInline{I00399}{\zeta_\omega}\allowbreak{} 联动能力愈发娴熟，应对高速突袭的技巧也更加熟练；奥斯卡的 \allowbreak{}\NovelInline{I00400}{\varepsilon_{\varepsilon_0}}\allowbreak{} 力量也得到了提升，距离突破 \allowbreak{}\NovelInline{I00401}{\zeta_0}\allowbreak{} 又近了一步；小舞的 \allowbreak{}\NovelInline{I00402}{\zeta_0}\allowbreak{} 力量也更加凝练，速度也有了小幅提升，应对突袭的反应也更快了。

阳光透过树叶的缝隙，洒在五人身上，林间一片静谧，只有灵气流动的细微声响，以及五人平稳的呼吸声。他们知道，猎杀千年凤尾鸡冠蛇，只是他们星斗大森林实战的第一步，也是猎杀人面魔蛛的热身，真正的挑战还在后面，但这场战斗带来的成长，尤其是应对高速突袭的经验，让他们更加坚定了信心，也为后续的战斗打下了坚实的基础。待他们吸收完灵气，恢复全部力量，便会继续前行，直面那只更强的对手------千年人面魔蛛。

\chapter{第十三章：难以对付的人面魔蛛}\label{ux7b2cux5341ux4e09ux7ae0ux96beux4ee5ux5bf9ux4ed8ux7684ux4ebaux9762ux9b54ux86db}

吸收完千年凤尾鸡冠蛇的浓郁灵气，五人的气息迎来质的飞跃，周身的序数纹路熠熠生辉，各自的扽技也随之凝练成型。唐三缓缓睁开双眼，眼底莹光暴涨，\allowbreak{}\NovelInline{I00403}{\zeta_2}\allowbreak{}的力量彻底稳固，蓝银草藤条上流转着深邃的金色纹路，感知范围大幅扩展，领悟扽技{[}\allowbreak{}\NovelInline{I00404}{\zeta}\allowbreak{}阶藤锁·迭代囚笼{]}------可催动\allowbreak{}\NovelInline{I00405}{\zeta_2}\allowbreak{}序数力量，让蓝银草形成层层迭代的藤笼，既能束缚敌人，又能叠加\allowbreak{}\NovelInline{I00406}{\varphi(1,0)}\allowbreak{}波动进行持续压制；

小舞身形轻盈跃起，周身萦绕着淡淡的粉色光晕，体内\allowbreak{}\NovelInline{I00407}{\zeta_1}\allowbreak{}的力量凝练如丝，完美契合柔骨兔的灵动特质，领悟两大专属扽技：{[}柔骨·瞬影缠{]}，借柔骨兔的极致敏捷瞬移近身，用柔韧的肢体缠绕敌人，牵制其动作并紊乱其序数波动，贴合柔骨兔的缠斗本能；{[}柔骨·裂空踢{]}，将\allowbreak{}\NovelInline{I00408}{\zeta_1}\allowbreak{}力量凝聚于腿部，踢出凌厉的柔劲腿风，既能破防又能精准攻击敌人薄弱处，兼具柔媚与强悍；

牢大活动着手臂，周身\allowbreak{}\NovelInline{I00409}{\varphi(3,\omega)}\allowbreak{}的力量愈发凝实，成功突破至该层级，手中可凝聚自身力量化作专属武器曼巴锤，同时领悟扽技{[}冰红茶·\allowbreak{}\NovelInline{I00410}{\varphi}\allowbreak{}阶增益{]}，饮用冰红茶后可短暂增幅自身力量，强化曼巴锤威力与肘击力量，同时为队友提供冰属性序数增益；

蔡须坤指尖转动篮球，周身\allowbreak{}\NovelInline{I00411}{\zeta_{\varepsilon_0}}\allowbreak{}的迭代节奏愈发流畅，领悟两大扽技：{[}篮球·\allowbreak{}\NovelInline{I00412}{\zeta}\allowbreak{}阶校准{]}，可通过运球节奏精准校准全队气息，甚至能短暂衔接高阶序数力量，填补波动断层；{[}music·声波震慑{]}，催动\allowbreak{}\NovelInline{I00413}{\zeta_{\varepsilon_0}}\allowbreak{}力量注入篮球，通过篮球弹跳发出高频声波，形成无形的声波冲击，对敌方进行震慑，干扰其序数迭代，甚至能短暂让低阶波动的敌人陷入僵直；

奥斯卡咬了一口补充灵气的香肠，体内力量彻底突破至\allowbreak{}\NovelInline{I00414}{\zeta_0}\allowbreak{}，领悟扽技{[}香肠·\allowbreak{}\NovelInline{I00415}{\zeta}\allowbreak{}阶双益{]}，可制作两种高阶香肠------增益香肠能提升队友序数层级，攻击香肠可释放\allowbreak{}\NovelInline{I00416}{\zeta_0}\allowbreak{}级冲击波，干扰敌人防御迭代。

``灵气吸收完毕，我们的实力都有了大幅提升。''唐三收起蓝银草，语气沉稳中带着一丝笃定，``我已突破至\allowbreak{}\NovelInline{I00417}{\zeta_2}\allowbreak{}，领悟扽技{[}\allowbreak{}\NovelInline{I00418}{\zeta}\allowbreak{}阶藤锁·迭代囚笼{]}；小舞达到\allowbreak{}\NovelInline{I00419}{\zeta_1}\allowbreak{}，领悟了柔骨兔专属扽技{[}柔骨·瞬影缠{]}和{[}柔骨·裂空踢{]}，近身缠斗和牵制能力大幅提升；

蔡须坤突破到\allowbreak{}\NovelInline{I00420}{\zeta_{\varepsilon_0}}\allowbreak{}，除了{[}篮球·\allowbreak{}\NovelInline{I00421}{\zeta}\allowbreak{}阶校准{]}，还领悟了{[}music·声波震慑{]}，可通过声波干扰敌人；牢大更是突破至\allowbreak{}\NovelInline{I00422}{\varphi(3,\omega)}\allowbreak{}，能凝聚自身力量化作曼巴锤，还能饮用冰红茶增幅力量，奥斯卡也成功抵达\allowbreak{}\NovelInline{I00423}{\zeta_0}\allowbreak{}，接下来寻找千年人面魔蛛，我们更有底气，但它的\allowbreak{}\NovelInline{I00424}{\varphi(\omega,0)}\allowbreak{}魂核和\allowbreak{}\NovelInline{I00425}{\varphi(4,0)}\allowbreak{}防御依旧强悍，且能吸收我们溢出的序数力量，大家务必配合默契，善用各自扽技，不可大意。''

众人齐声应和，重新调整阵型，周身气息联动愈发紧密：唐三走在最前方，蓝银草舒展，{[}\allowbreak{}\NovelInline{I00426}{\zeta}\allowbreak{}阶藤锁·迭代囚笼{]}随时可发动；小舞紧随其后，身形轻盈如兔，粉色光晕萦绕周身，{[}柔骨·瞬影缠{]}和{[}柔骨·裂空踢{]}随时准备近身缠斗；

牢大走在中间，手中已隐隐有淡金色力量涌动，随时可凝聚曼巴锤，身旁还悬浮着一瓶冰红茶，{[}冰红茶·\allowbreak{}\NovelInline{I00427}{\varphi}\allowbreak{}阶增益{]}随时可激活；蔡须坤手持篮球，运球节奏平稳，{[}篮球·\allowbreak{}\NovelInline{I00428}{\zeta}\allowbreak{}阶校准{]}时刻校准全队气息，{[}music·声波震慑{]}也做好了干扰准备；奥斯卡走在最后，手中握着数根刚制作的\allowbreak{}\NovelInline{I00429}{\zeta_0}\allowbreak{}级香肠，{[}香肠·\allowbreak{}\NovelInline{I00430}{\zeta}\allowbreak{}阶双益{]}随时可提供支援。

星斗大森林深处，古木愈发参天，枝叶交错，遮天蔽日，只有零星的阳光透过枝叶缝隙洒落，形成斑驳的光影。林间的灵气愈发浓郁，却也夹杂着一股阴冷刺骨的压抑感，低阶魂兽的气息几乎消失不见，取而代之的是一种远超\allowbreak{}\NovelInline{I00431}{\zeta_{\varepsilon_0}}\allowbreak{}、逼近\allowbreak{}\NovelInline{I00432}{\varphi(\omega,0)}\allowbreak{}的恐怖序数波动，隐晦而狂暴，如同一张无形的网，笼罩着整片区域，让五人都忍不住心头一紧，呼吸也变得谨慎起来。

唐三的蓝银草突然剧烈震颤，藤条上的\allowbreak{}\NovelInline{I00433}{\zeta_2}\allowbreak{}纹路瞬间黯淡，甚至出现了细微的裂痕，一股阴冷的压迫感从前方的灌木丛中传来，让他浑身的序数波动都出现了轻微的紊乱：``找到了！前方就是千年人面魔蛛，它的魂核波动达到了\allowbreak{}\NovelInline{I00434}{\varphi(\omega,0)}\allowbreak{}，防御正是我们要应对的\allowbreak{}\NovelInline{I00435}{\varphi(4,0)}\allowbreak{}，比预想中还要强悍，大家做好战斗准备，善用扽技！''

话音未落，一道漆黑的身影突然从灌木丛中窜出，速度快得惊人，周身萦绕着浓郁的紫黑色光晕，那光晕中蕴含着\allowbreak{}\NovelInline{I00436}{\varphi(\omega,0)}\allowbreak{}的超限序数波动，阴冷而狂暴，所过之处，灵气瞬间被吞噬，地面留下一道漆黑的痕迹。它身形庞大，通体覆盖着坚硬的黑色外甲，外甲上布满了细密的\allowbreak{}\NovelInline{I00437}{\varphi(4,0)}\allowbreak{}序数纹路，如同坚不可摧的防御屏障，头部生着一张扭曲的人脸，双眼泛着猩红的光芒，八只粗壮的足上布满尖刺，每一根尖刺都萦绕着\allowbreak{}\NovelInline{I00438}{\zeta_0}\allowbreak{}级别的波动，尾部拖着一根细长的毒刺，散发着致命的寒气------正是千年人面魔蛛！

千年人面魔蛛落地的瞬间，八足狠狠蹬地，地面瞬间被震出细密的裂纹，紫黑色的\allowbreak{}\NovelInline{I00439}{\varphi(4,0)}\allowbreak{}防御屏障瞬间展开，将自身包裹其中，同时它发出一声尖锐的嘶鸣，声音刺耳，带着强烈的攻击性，周身的紫黑色波动愈发狂暴，开始主动吸收周围的灵气，甚至试图吞噬五人身上溢出的序数力量。

``动手！按计划配合，发动扽技！''唐三一声令下，率先发难，蓝银草如锋利的箭矢般窜出，数十根藤条交错缠绕，催动扽技{[}\allowbreak{}\NovelInline{I00440}{\zeta}\allowbreak{}阶藤锁·迭代囚笼{]}，藤条上的\allowbreak{}\NovelInline{I00441}{\zeta_2}\allowbreak{}纹路通体发亮，按照\allowbreak{}\NovelInline{I00442}{f_{\zeta_0}}\allowbreak{}的迭代规律，从\allowbreak{}\NovelInline{I00443}{\varphi(4,0)}\allowbreak{}防御的\allowbreak{}\NovelInline{I00444}{\zeta_0}\allowbreak{}底层纹路开始，层层递进，形成迭代藤笼包裹住人面魔蛛，``蔡须坤，发动{[}篮球·\allowbreak{}\NovelInline{I00445}{\zeta}\allowbreak{}阶校准{]}，校准迭代节奏，配合我渗透屏障，它能吸收我们的力量，务必保持气息联动，不可单独发力！''

蔡须坤立刻回应，运球节奏瞬间加快，催动扽技{[}篮球·\allowbreak{}\NovelInline{I00446}{\zeta}\allowbreak{}阶校准{]}，篮球每一次落地、弹起都发出清脆的声响，银白色的\allowbreak{}\NovelInline{I00447}{\zeta_{\varepsilon_0}}\allowbreak{}波动顺着篮球的弹跳，源源不断地传递到唐三的蓝银草上，完美填补了力量衔接的断层，同时校准牢大、小舞的气息，``收到！{[}篮球·\allowbreak{}\NovelInline{I00448}{\zeta}\allowbreak{}阶校准{]}已发动，与它的\allowbreak{}\NovelInline{I00449}{\varphi(4,0)}\allowbreak{}迭代同频，时刻警惕它的力量吸收！''

小舞身形一闪，如柔骨兔般灵动瞬移，催动扽技{[}柔骨·瞬影缠{]}，粉色光晕暴涨，肢体变得愈发柔韧，瞬间缠上人面魔蛛的一只残足，借助柔劲死死牵制其动作，同时指尖凝聚\allowbreak{}\NovelInline{I00450}{\zeta_1}\allowbreak{}力量，发动{[}柔骨·裂空踢{]}，凌厉的腿风狠狠踢向屏障的薄弱处------头部与人身衔接的位置，``唐三，我用{[}柔骨·瞬影缠{]}牵制它，{[}柔骨·裂空踢{]}打乱它的\allowbreak{}\NovelInline{I00451}{\varphi(4,0)}\allowbreak{}迭代了，快趁机突破！''腿风与火焰冲击同时落在屏障上，瞬间炸开一道细小的火星，屏障波动出现短暂紊乱。

``嘶------''人面魔蛛被小舞的攻击激怒，发出一声尖锐的嘶鸣，尾部毒刺猛地一挥，一道凝聚着纯碎\allowbreak{}\NovelInline{I00452}{\zeta_0}\allowbreak{}力量的紫黑色冲击波，带着刺耳的破空声，朝着小舞狠狠砸去。与此同时，它的八足同时发力，身形骤然一动，朝着防御最薄弱的奥斯卡猛冲而去，尖刺上的\allowbreak{}\NovelInline{I00453}{\zeta_0}\allowbreak{}波动愈发浓郁，显然是想先击溃辅助，瓦解五人的配合。

``小心！''牢大低喝一声，身形如惊雷般一闪，先抓起身旁的冰红茶一饮而尽，催动扽技{[}冰红茶·\allowbreak{}\NovelInline{I00454}{\varphi}\allowbreak{}阶增益{]}，\allowbreak{}\NovelInline{I00455}{\varphi(3,\omega)}\allowbreak{}的力量瞬间暴涨，同时手中凝聚自身力量，化作一柄通体金黄、布满序数纹路的曼巴锤，稳稳挡在奥斯卡身前，``砰''的一声闷响，冲击波狠狠撞在曼巴锤上，瞬间溃散，而人面魔蛛的尖刺也狠狠撞在锤身，冰蓝色的增益波动与紫黑色的\allowbreak{}\NovelInline{I00456}{\zeta_0}\allowbreak{}波动剧烈碰撞，发出刺耳的滋滋声，牢大脚下的地面被震出深深的凹陷，身体微微后退两步，却依旧稳稳顶住了这一击，``我的{[}冰红茶·\allowbreak{}\NovelInline{I00457}{\varphi}\allowbreak{}阶增益{]}全力生效，力量已增幅，曼巴锤可硬抗它的攻击，增益也给到大家，大家全力进攻，它的防御太强，而且能吸收我们的力量，不要给它可乘之机！''

淡蓝色的冰属性光晕瞬间笼罩着五人，唐三的\allowbreak{}\NovelInline{I00458}{\zeta_2}\allowbreak{}力量愈发凝练，小舞的速度也提升了几分，蔡须坤的联动也更加稳定。奥斯卡立刻催动扽技{[}香肠·\allowbreak{}\NovelInline{I00459}{\zeta}\allowbreak{}阶双益{]}，掏出三根增益香肠和两根攻击香肠，指尖发力，攻击香肠如两颗炮弹般射向人面魔蛛的头部、腹部，淡蓝色的\allowbreak{}\NovelInline{I00460}{\zeta_0}\allowbreak{}波动瞬间炸开，狠狠击中其外甲，``嘭嘭''两声闷响，人面魔蛛发出一声痛苦的嘶鸣，\allowbreak{}\NovelInline{I00461}{\varphi(4,0)}\allowbreak{}防御屏障波动出现了明显的紊乱，但仅仅片刻，它便吸收了周围的灵气，屏障再次稳固，甚至比之前更加厚重------它的\allowbreak{}\NovelInline{I00462}{\varphi(\omega,0)}\allowbreak{}魂核，能快速吸收灵气和溢出的序数力量，修复自身防御。

``不好，它的修复速度太快了，而且能吸收我们的攻击力量！''唐三脸色一变，他发现自己的蓝银草藤笼刚渗透屏障一丝，力量便被人面魔蛛吸收，藤条上的\allowbreak{}\NovelInline{I00463}{\zeta_2}\allowbreak{}纹路也变得黯淡，``蔡须坤，加快{[}篮球·\allowbreak{}\NovelInline{I00464}{\zeta}\allowbreak{}阶校准{]}节奏，同时发动{[}music·声波震慑{]}，用声波干扰它的力量吸收和修复；小舞，继续用{[}柔骨·瞬影缠{]}和{[}柔骨·裂空踢{]}牵制，打断它的修复节奏；奥斯卡，加大{[}香肠·\allowbreak{}\NovelInline{I00465}{\zeta}\allowbreak{}阶双益{]}输出，用攻击香肠干扰它的力量吸收，增益香肠持续给我们补力！''

众人立刻调整战术，蔡须坤将{[}篮球·\allowbreak{}\NovelInline{I00466}{\zeta}\allowbreak{}阶校准{]}节奏提升到极致，同时催动扽技{[}music·声波震慑{]}，篮球弹跳频率加快，发出高频声波，无形的冲击朝着人面魔蛛席卷而去，``唐三，{[}music·声波震慑{]}已发动，正在干扰它的迭代，力量已联动，集中突破！''声波落在人面魔蛛身上，它发出一声烦躁的嘶鸣，\allowbreak{}\NovelInline{I00467}{\varphi(4,0)}\allowbreak{}屏障的迭代节奏出现明显紊乱，篮球撞击屏障的瞬间，银白色的波动剧烈炸开，屏障上的\allowbreak{}\NovelInline{I00468}{\varphi(4,0)}\allowbreak{}纹路瞬间扭曲、断裂。唐三眼神一凝，抓住这个间隙，将自身\allowbreak{}\NovelInline{I00469}{\zeta_2}\allowbreak{}力量与牢大的\allowbreak{}\NovelInline{I00470}{\varphi(3,\omega)}\allowbreak{}力量通过蔡须坤的篮球紧紧联动，催动扽技{[}\allowbreak{}\NovelInline{I00471}{\zeta}\allowbreak{}阶藤锁·迭代囚笼{]}，藤条瞬间收紧，死死缠绕住屏障裂痕处，两道力量在藤条上叠加，形成一道蕴含着双重高阶序数力量的金色光束，狠狠砸向屏障的裂痕处。

牢大低吼一声，再次灌下一口冰红茶，将{[}冰红茶·\allowbreak{}\NovelInline{I00472}{\varphi}\allowbreak{}阶增益{]}的效果催至极致，\allowbreak{}\NovelInline{I00473}{\varphi(3,\omega)}\allowbreak{}的力量全部爆发，手中曼巴锤瞬间暴涨，通体金黄的锤身萦绕着冰蓝色的增益光晕，一记势大力沉的肘击结合曼巴锤重击，带着刺骨的寒气和强悍的力量，狠狠砸在屏障的裂痕上，冰蓝色的增益波动与金色光束相互叠加，双重力量瞬间撕裂了\allowbreak{}\NovelInline{I00474}{\varphi(4,0)}\allowbreak{}的防御屏障，``咔嚓''一声脆响，屏障出现了一道细小的裂痕。

可就在众人以为局势好转时，人面魔蛛突然发出一声狂暴的嘶鸣，\allowbreak{}\NovelInline{I00475}{\varphi(\omega,0)}\allowbreak{}魂核全力爆发，紫黑色的波动瞬间暴涨，不仅快速修复了屏障的裂痕，还主动吸收了五人溢出的序数力量，周身的波动愈发强悍，八足同时发力，朝着五人发起了全方位的攻击，尖刺上的\allowbreak{}\NovelInline{I00476}{\zeta_0}\allowbreak{}波动化作一道道紫黑色的利刃，密集地射向五人。

``快防御！''唐三大喊一声，催动扽技{[}\allowbreak{}\NovelInline{I00477}{\zeta}\allowbreak{}阶藤锁·迭代囚笼{]}编织成防御网，试图阻挡利刃攻击，可蓝银草的\allowbreak{}\NovelInline{I00478}{\zeta_2}\allowbreak{}力量在\allowbreak{}\NovelInline{I00479}{\varphi(\omega,0)}\allowbreak{}的波动压制下，根本无法完全抵御，几根藤条瞬间被利刃斩断，紫黑色的波动顺着藤条传递到唐三身上，让他浑身一震，嘴角溢出一丝灵气。小舞凭借柔骨兔的敏捷瞬移躲避，同时发动{[}柔骨·瞬影缠{]}缠住一根袭来的尖刺，借力翻身，却还是被另一道利刃擦中手臂，\allowbreak{}\NovelInline{I00480}{\zeta_1}\allowbreak{}波动瞬间紊乱，两大柔骨扽技暂时无法发动；蔡须坤的篮球被利刃击中，{[}篮球·\allowbreak{}\NovelInline{I00481}{\zeta}\allowbreak{}阶校准{]}节奏彻底打乱，{[}music·声波震慑{]}也被迫中断，\allowbreak{}\NovelInline{I00482}{\zeta_{\varepsilon_0}}\allowbreak{}联动中断；牢大的{[}冰红茶·\allowbreak{}\NovelInline{I00483}{\varphi}\allowbreak{}阶承压{]}屏障被多道利刃击中，出现了细密的裂痕，\allowbreak{}\NovelInline{I00484}{\varphi(3,\omega)}\allowbreak{}波动也开始紊乱；奥斯卡被牢大护在身后，却也被波动余劲震得连连后退，手中的香肠险些掉落，{[}香肠·\allowbreak{}\NovelInline{I00485}{\zeta}\allowbreak{}阶双益{]}的输出也出现断层。

``这样下去不行，我们的力量根本无法突破它的\allowbreak{}\NovelInline{I00486}{\varphi(4,0)}\allowbreak{}防御，反而会被它不断吸收力量，迟早会被耗竭！''牢大咬紧牙关，又饮下一口冰红茶维持{[}冰红茶·\allowbreak{}\NovelInline{I00487}{\varphi}\allowbreak{}阶增益{]}效果，手中曼巴锤的光晕渐渐黯淡，显然已经消耗了大量力量，``它的\allowbreak{}\NovelInline{I00488}{\varphi(\omega,0)}\allowbreak{}魂核太强悍了，就算我有冰红茶增幅、握着曼巴锤，也难以抗衡，我们的层级还是不够！''

就在五人陷入绝境，人面魔蛛准备发动致命一击时，一道浑厚的气息突然从林间传来，\allowbreak{}\NovelInline{I00489}{\psi(\psi a(\Omega_{a+2}))}\allowbreak{}的高阶序数波动瞬间笼罩全场，压制住了人面魔蛛的\allowbreak{}\NovelInline{I00490}{\varphi(\omega,0)}\allowbreak{}力量------赵无极及时赶到！他负手而立，周身淡青色的光晕暴涨，\allowbreak{}\NovelInline{I00491}{\psi(\psi a(\Omega_{a+2}))}\allowbreak{}的波动浑厚而凝练，远远超过人面魔蛛的实力，``孽畜，敢伤我的学生！''

话音未落，赵无极抬手凌空一挥，淡青色的序数力量凝聚成一道巨大的掌印，带着\allowbreak{}\NovelInline{I00492}{\psi(\psi a(\Omega_{a+2}))}\allowbreak{}的强悍威力，狠狠砸向人面魔蛛的\allowbreak{}\NovelInline{I00493}{\varphi(4,0)}\allowbreak{}防御屏障。``咔嚓------轰隆！''一声巨响，原本坚不可摧的\allowbreak{}\NovelInline{I00494}{\varphi(4,0)}\allowbreak{}屏障，在赵无极的全力一击下，瞬间碎裂，紫黑色的波动四散开来，人面魔蛛发出一声凄厉的嘶鸣，被掌印余劲震得连连后退，八足有两只被震断，\allowbreak{}\NovelInline{I00495}{\varphi(\omega,0)}\allowbreak{}魂核也出现了紊乱，无法再快速吸收灵气和力量。

``老师！''五人齐声呼喊，脸上露出了希望的光芒。赵无极目光扫过五人，语气凝重：``它的\allowbreak{}\NovelInline{I00496}{\varphi(\omega,0)}\allowbreak{}魂核只是暂时紊乱，防御屏障被我破开，接下来，就靠你们自己了------记住，它的魂核是核心，只要击溃魂核，就能彻底击败它，我不会再出手，这是你们的实战考验！''说完，他便后退几步，收敛气息，静静注视着五人，随时准备在危急时刻兜底。

五人立刻调整状态，虽然浑身疲惫，力量也有消耗，但看到人面魔蛛防御破损、魂核紊乱，心中重新燃起了斗志。唐三擦去嘴角的灵气，重新催动扽技{[}\allowbreak{}\NovelInline{I00497}{\zeta}\allowbreak{}阶藤锁·迭代囚笼{]}，蓝银草藤条上的\allowbreak{}\NovelInline{I00498}{\zeta_2}\allowbreak{}纹路再次亮起，这一次，他不再急于渗透，而是将力量集中在藤条尖端，瞄准人面魔蛛紊乱的\allowbreak{}\NovelInline{I00499}{\varphi(\omega,0)}\allowbreak{}魂核，``大家配合我，集中力量攻击它的魂核，不给它修复的机会！牢大，你试试结合奥斯卡的增益香肠，看看能不能短暂突破层级！''

蔡须坤捡起篮球，拼尽全力催动扽技{[}篮球·\allowbreak{}\NovelInline{I00500}{\zeta}\allowbreak{}阶校准{]}，\allowbreak{}\NovelInline{I00501}{\zeta_{\varepsilon_0}}\allowbreak{}的波动重新与众人气息同频，将唐三的\allowbreak{}\NovelInline{I00502}{\zeta_2}\allowbreak{}、牢大的\allowbreak{}\NovelInline{I00503}{\varphi(3,\omega)}\allowbreak{}力量紧紧联动，``节奏已校准，力量全部传递给你，唐三！我会全程稳住联动，不让气息断层！''

小舞深吸一口气，稳住紊乱的\allowbreak{}\NovelInline{I00504}{\zeta_1}\allowbreak{}波动，再次催动扽技{[}柔骨·瞬影缠{]}，身形如柔骨兔般灵活穿梭，避开人面魔蛛的尖刺攻击，瞬间缠上它的尾部，死死牵制其动作，同时发动{[}柔骨·裂空踢{]}，腿部凝聚全部\allowbreak{}\NovelInline{I00505}{\zeta_1}\allowbreak{}力量，狠狠踢向人面魔蛛的魂核位置，``唐三，我用{[}柔骨·瞬影缠{]}牵制它，{[}柔骨·裂空踢{]}打乱它的魂核波动，快动手！''

奥斯卡立刻反应过来，催动扽技{[}香肠·\allowbreak{}\NovelInline{I00506}{\zeta}\allowbreak{}阶双益{]}，掏出最后两根高阶增益香肠，全力掷向牢大，``牢大，接住！这两根增益香肠能短暂提升你的序数层级，结合你的{[}冰红茶·\allowbreak{}\NovelInline{I00507}{\varphi}\allowbreak{}阶承压{]}，说不定能突破到\allowbreak{}\NovelInline{I00508}{\varphi(\omega,0)}\allowbreak{}！''

牢大眼中闪过一丝决绝，一把接住增益香肠，毫不犹豫地吞下，同时将剩余的冰红茶一饮而尽，香肠的\allowbreak{}\NovelInline{I00509}{\zeta_0}\allowbreak{}增益力量与冰红茶的{[}\allowbreak{}\NovelInline{I00510}{\varphi}\allowbreak{}阶增益{]}效果叠加，与他体内的\allowbreak{}\NovelInline{I00511}{\varphi(3,\omega)}\allowbreak{}力量、自身本源力量相互融合，周身的金色与冰蓝色光晕同时暴涨，气息疯狂攀升，从\allowbreak{}\NovelInline{I00512}{\varphi(3,\omega)}\allowbreak{}一路突破，短暂抵达\allowbreak{}\NovelInline{I00513}{\varphi(\omega,0)}\allowbreak{}的层级！``喝！''牢大发出一声狂暴的低吼，手中的曼巴锤彻底凝实，化作一柄丈高的金色巨锤，锤身布满高阶序数纹路，萦绕着冰蓝色的增益光晕，\allowbreak{}\NovelInline{I00514}{\varphi(\omega,0)}\allowbreak{}的波动浑厚而狂暴，``唐三，我来牵制它，用曼巴锤压制它，你趁机发动致命一击！''

人面魔蛛被众人牵制，魂核紊乱，无法修复防御，也无法吸收力量，看到牢大突然突破到\allowbreak{}\NovelInline{I00515}{\varphi(\omega,0)}\allowbreak{}，眼中闪过一丝恐惧，发出狂暴的嘶鸣，挥舞着剩余的六足，疯狂发起攻击，试图突围。可五人的配合愈发默契，

蔡须坤重新催动{[}篮球·\allowbreak{}\NovelInline{I00516}{\zeta}\allowbreak{}阶校准{]}和{[}music·声波震慑{]}，高频声波持续干扰人面魔蛛的魂核波动，让它动作变得迟缓，同时稳定联动众人力量；牢大握着丈高的曼巴锤，借着冰红茶和香肠的双重增益、以及短暂突破的\allowbreak{}\NovelInline{I00517}{\varphi(\omega,0)}\allowbreak{}力量，狠狠砸向人面魔蛛的身体，每一击都势大力沉，带着刺骨的寒气和强悍的冲击力，死死牵制它的动作，砸得人面魔蛛连连后退，外甲不断破损、开裂；

小舞则靠着{[}柔骨·瞬影缠{]}和{[}柔骨·裂空踢{]}，灵活近身缠斗，不断踢击人面魔蛛的薄弱处，持续干扰魂核波动；奥斯卡虽然耗尽了香肠，却依旧死死盯着人面魔蛛的动作，随时提醒众人躲避攻击；唐三则专注于瞄准魂核，将自身\allowbreak{}\NovelInline{I00518}{\zeta_2}\allowbreak{}力量、蔡须坤的\allowbreak{}\NovelInline{I00519}{\zeta_{\varepsilon_0}}\allowbreak{}力量，以及牢大短暂爆发的\allowbreak{}\NovelInline{I00520}{\varphi(\omega,0)}\allowbreak{}余劲，通过蓝银草全部凝聚在尖端，积蓄力量，准备发动致命一击。

``就是现在！''唐三眼神一凝，催动扽技{[}\allowbreak{}\NovelInline{I00521}{\zeta}\allowbreak{}阶藤锁·迭代囚笼{]}将人面魔蛛死死束缚，藤条上的\allowbreak{}\NovelInline{I00522}{\zeta_2}\allowbreak{}纹路与\allowbreak{}\NovelInline{I00523}{\varphi(1,0)}\allowbreak{}波动叠加，形成一道坚固的囚笼，同时将所有凝聚的力量全部爆发，化作一道蕴含三重高阶序数力量的金色光束，带着毁天灭地的力量，狠狠砸向人面魔蛛的\allowbreak{}\NovelInline{I00524}{\varphi(\omega,0)}\allowbreak{}魂核。

与此同时，牢大的\allowbreak{}\NovelInline{I00525}{\varphi(\omega,0)}\allowbreak{}短暂力量即将耗尽，冰红茶和香肠的增益效果也濒临消失，他拼尽最后一丝力气，挥动曼巴锤，结合增幅后的肘击，一记势大力沉的重击，砸在人面魔蛛的头部，金色的曼巴锤力量与冰蓝色的残余增益波动，一同与金色光束相互呼应，双重重击狠狠落在魂核上。

``嘶------''人面魔蛛发出最后一声凄厉的嘶鸣，魂核被双重力量击中，\allowbreak{}\NovelInline{I00526}{\varphi(\omega,0)}\allowbreak{}的波动瞬间溃散，身体剧烈抽搐，八足不再动弹，周身的紫黑色光晕也渐渐消散，庞大的身躯重重倒在地上，彻底失去了生机，随后缓缓化作浓郁的灵气，弥漫在整个空地中，黑色的外甲和尖刺也逐渐失去光泽，散落一地。

战斗结束的瞬间，牢大体内的增益力量和\allowbreak{}\NovelInline{I00527}{\varphi(\omega,0)}\allowbreak{}的气息瞬间消退，回落至\allowbreak{}\NovelInline{I00528}{\varphi(3,\omega)}\allowbreak{}，手中的曼巴锤也渐渐化作力量消散，他浑身一软，瘫倒在地，大口喘着气，浑身肌肉酸痛，显然短暂突破\allowbreak{}\NovelInline{I00529}{\varphi(\omega,0)}\allowbreak{}、加上持续使用曼巴锤和冰红茶增益，消耗了他大量的本源力量。

五人也纷纷瘫坐在地上，脸上满是疲惫却欣喜的神情，额头上布满了汗水，身上的衣物也有些破损，这场战斗异常艰难------从最初被人面魔蛛压制、力量被吸收，到赵无极出手破开防御，再到牢大结合奥斯卡的香肠和冰红茶短暂突破\allowbreak{}\NovelInline{I00530}{\varphi(\omega,0)}\allowbreak{}、用曼巴锤牵制，众人齐心协力，最终击溃魂核，每一步都拼尽了全力。

``我们\ldots\ldots 我们赢了！''蔡须坤揉了揉发酸的手臂，收起篮球，脸上满是激动，``太不容易了，若不是老师出手破开它的\allowbreak{}\NovelInline{I00531}{\varphi(4,0)}\allowbreak{}防御，若不是牢大结合奥斯卡的香肠短暂突破到\allowbreak{}\NovelInline{I00532}{\varphi(\omega,0)}\allowbreak{}，我们根本无法击败它，\allowbreak{}\NovelInline{I00533}{\varphi(\omega,0)}\allowbreak{}的魂核，果然名不虚传。我的{[}篮球·\allowbreak{}\NovelInline{I00534}{\zeta}\allowbreak{}阶校准{]}最后差点跟不上节奏，还好{[}music·声波震慑{]}起到了干扰作用，不然我们很难牵制住它。''

牢大靠在树干上，缓了许久才开口，声音虚弱却带着自豪：``刚才短暂突破到\allowbreak{}\NovelInline{I00535}{\varphi(\omega,0)}\allowbreak{}，那种力量太强悍了，但也快把我的本源力量耗光了，还好没有白费。靠着冰红茶的增益，我的曼巴锤和肘击力量才能大幅提升，再加上奥斯卡的增益香肠，才能勉强突破，这场战斗，少了谁都不行。''

奥斯卡摊了摊手，脸上满是庆幸和自豪：``我的{[}香肠·\allowbreak{}\NovelInline{I00536}{\zeta}\allowbreak{}阶双益{]}总算派上了大用场，虽然耗尽了所有香肠，但能帮牢大突破到\allowbreak{}\NovelInline{I00537}{\varphi(\omega,0)}\allowbreak{}，击败人面魔蛛，一切都值得！现在我已经达到\allowbreak{}\NovelInline{I00538}{\zeta_0}\allowbreak{}，接下来要尽快打磨扽技，制作更强的香肠。''

小舞靠在唐三身边，手臂上的伤口还在隐隐作痛，却依旧笑容明媚：``是啊，这场战斗虽然艰难，但我们做到了！我的{[}柔骨·瞬影缠{]}和{[}柔骨·裂空踢{]}总算发挥了作用，能死死牵制它、打乱它的节奏，下次我要提升\allowbreak{}\NovelInline{I00539}{\zeta_1}\allowbreak{}力量，让柔劲更足，近身缠斗的威力也能更强。''

唐三缓缓站起身，舒展了一下身体，体内的\allowbreak{}\NovelInline{I00540}{\zeta_2}\allowbreak{}波动逐渐恢复平稳，他看着倒地的人面魔蛛，眼底满是领悟：``这场战斗，我们每个人都拼尽了全力，也充分发挥了各自的扽技优势。我的{[}\allowbreak{}\NovelInline{I00541}{\zeta}\allowbreak{}阶藤锁·迭代囚笼{]}既能束缚又能压制，蔡须坤的{[}篮球·\allowbreak{}\NovelInline{I00542}{\zeta}\allowbreak{}阶校准{]}是联动核心，{[}music·声波震慑{]}更是关键的干扰手段；小舞的{[}柔骨·瞬影缠{]}和{[}柔骨·裂空踢{]}，完美发挥了柔骨兔的近身缠斗和牵制优势；牢大的{[}冰红茶·\allowbreak{}\NovelInline{I00543}{\varphi}\allowbreak{}阶增益{]}能增幅力量，曼巴锤则是核心攻击武器，负责承压和压制，奥斯卡的{[}香肠·\allowbreak{}\NovelInline{I00544}{\zeta}\allowbreak{}阶双益{]}提供支援，正是这份配合，再加上牢大的短暂突破，我们才能击败强悍的人面魔蛛。''

赵无极走上前，看着瘫坐在地上的五人，脸上露出了满意的笑容：``很好，你们没有让我失望。虽然我出手破开了它的\allowbreak{}\NovelInline{I00545}{\varphi(4,0)}\allowbreak{}防御，但后续的战斗，都是你们自己完成的，尤其是牢大，能结合奥斯卡的香肠和冰红茶的增益，短暂突破到\allowbreak{}\NovelInline{I00546}{\varphi(\omega,0)}\allowbreak{}，用曼巴锤死死压制敌人，这份应变能力和毅力，值得肯定；小舞的柔骨兔扽技运用得极为娴熟，近身牵制恰到好处；蔡须坤的{[}music·声波震慑{]}更是起到了关键干扰作用，你们每个人的扽技都运用得很到位，团队配合也愈发默契。千年人面魔蛛的\allowbreak{}\NovelInline{I00547}{\varphi(\omega,0)}\allowbreak{}魂核，确实强悍，你们能在力量被吸收、陷入绝境的情况下，调整战术、拼尽全力，已经非常出色了。''

他顿了顿，继续说道：``这场实战，不仅检验了你们的修炼成果和扽技运用能力，也让你们认识到了自身的不足------就算牢大能短暂突破到\allowbreak{}\NovelInline{I00548}{\varphi(\omega,0)}\allowbreak{}，也无法长时间维持，你们的序数层级，距离\allowbreak{}\NovelInline{I00549}{\varphi(\omega,0)}\allowbreak{}还有很大差距，应对这种高阶魂兽，依旧有些吃力。接下来，你们要吸收这人面魔蛛的灵气，巩固自身境界，打磨各自的扽技，继续加强团队配合，争取早日突破到更高的序数层级，真正具备独当一面的实力。唯有不断提升自身的序数层级，打磨扽技，加强团队配合，才能在充满危险的星斗大森林中站稳脚跟，才能应对未来更多更强的挑战。''

\chapter{第十四章：新成员的到来和的突破}\label{ux7b2cux5341ux56dbux7ae0ux65b0ux6210ux5458ux7684ux5230ux6765ux548cux7684ux7a81ux7834}

人面魔珠打败后，队员们正在吸收灵气，在人面魔蛛的强大灵气滋养下，队员们回复了体力，不过这天地灵气还很让人神清气爽的，更难能可贵的是，人面魔蛛居然爆了馃槔这种天地宝物，要知道，馃槔很罕见，一般只有十万年魂兽才会有馃槔这种东西，佩戴者可以与扽魂产生共鸣，从而大幅提升实力和悟性。吸收完成馃槔后，大家实力得到了提升，一方面，人面魔蛛的灵气滋养，另一方面，就是馃槔的作用，让唐三突破\allowbreak{}\NovelInline{I00550}{\zeta_2}\allowbreak{}，直接到达\allowbreak{}\NovelInline{I00551}{\varphi(3,0)}\allowbreak{}，而牢大到达\allowbreak{}\NovelInline{I00552}{\varphi(4,0)}\allowbreak{}的水平，蔡须坤也到达\allowbreak{}\NovelInline{I00553}{\varphi(3,1)}\allowbreak{}水平，奥斯卡突破到了\allowbreak{}\NovelInline{I00554}{\zeta_0}\allowbreak{}，直接到了\allowbreak{}\NovelInline{I00555}{\zeta_{\zeta_0}}\allowbreak{}

回去的路上，他们在谈论这次冒险的惊心动魄，在中级学院，门口，院长弗兰德正在等带他们，'院长大人，这次实战应道6人，实到6人，平安归来。'赵无极说，弗兰德看到众人平安，说：欢迎你们的回来，现在学院有了新成员，名字叫卢本韦，女性，是全帝国最大的养猪厂厂长卢关的女儿。卢管接任养猪场之后一改往日作风，抓住一些地方南通多的特点，深度改革了养猪场的经营模式，最终抓住了商机，一举成为当地首富，他的女儿卢本韦也来到这里学习，说起来，卢本韦的扽魂就是母猪，技能就是母猪冲撞，利用母猪发疯的力量击倒敌人。

弗兰德：现在你们小队又多了一份力量，卢本韦，来介绍一下自己''

卢本韦：我是你们的新成员，卢本伟，她个子算高一点的，说话冷冷的

`我来自成都，应该听过我的产业吧，'

唐三当然知道，卢氏家族垄断了整个帝国的产业，全国养猪场一半以上都在他们家族手上，家族内部，更是有两大强者护法，没人会招惹如此的势力。

卢本韦：``既然来了，我希望快速融入大家，也想学习一些高深的知识，我卢某还委托老师同学们的帮助和支持了''。老师们，我虽然来自大家庭，但是我可不是来玩的，之后向众人展示了她\allowbreak{}\NovelInline{I00556}{\varphi(2,1)}\allowbreak{}，也就是\allowbreak{}\NovelInline{I00557}{\zeta_1}\allowbreak{}级别的实力。

还有我！我叫诺当，扽魂磁暴环，击中会让人晕眩，降低力量，他淡淡说到，他的磁爆环可很厉害，有着\allowbreak{}\NovelInline{I00558}{\varphi(2,\omega)}\allowbreak{}也就是\allowbreak{}\NovelInline{I00559}{\zeta_\omega}\allowbreak{}的力量。

经过几天的认识，唐三他们熟悉了他的伙伴

弗兰德召集全体队员齐聚课堂，此次授课不再泛讲\allowbreak{}\NovelInline{I00560}{\Gamma_0}\allowbreak{}，而是结合众人实战中接触的\allowbreak{}\NovelInline{I00561}{\varphi(\omega,0)}\allowbreak{}力量，细致拆解从\allowbreak{}\NovelInline{I00562}{\varphi(\omega,0)}\allowbreak{}到\allowbreak{}\NovelInline{I00563}{\varphi(\omega,1)}\allowbreak{}，再进阶至\allowbreak{}\NovelInline{I00564}{\varphi(\omega+1,0)}\allowbreak{}的完整逻辑，全程以对话授课形式展开，鼓励学员结合实战感悟提问，把二元Veblen函数的基本列规则讲透。

弗兰德站在讲台中央，指尖轻点虚空，浮现出清晰的序数迭代脉络，先是定格在众人熟悉的\allowbreak{}\NovelInline{I00565}{\varphi(\omega,0)}\allowbreak{}上，缓缓开口：``咱们上次实战，牢大短暂触及\allowbreak{}\NovelInline{I00566}{\varphi(\omega,0)}\allowbreak{}层级，也就是超Cantor序数，想必大家对这个层级的力量还有印象。先回顾定义：\allowbreak{}\NovelInline{I00567}{\varphi(\omega,0)}\allowbreak{}是所有有限n对应的\allowbreak{}\NovelInline{I00568}{\varphi(n,0)}\allowbreak{}的极限，它的基本列规则非常明确，\allowbreak{}\NovelInline{I00569}{\varphi(\omega,0)[n]=\varphi(n,0)}\allowbreak{}，也就是取n为0、1、2\ldots\ldots 时，对应\allowbreak{}\NovelInline{I00570}{\varphi(0,0)}\allowbreak{}、\allowbreak{}\NovelInline{I00571}{\varphi(1,0)}\allowbreak{}、\allowbreak{}\NovelInline{I00572}{\varphi(2,0)}\allowbreak{}\ldots\ldots 直至上确界，这是所有有限层级Veblen函数的封顶，也是咱们进阶的起点。''

他顿了顿，目光扫过众人：``\allowbreak{}\NovelInline{I00573}{\varphi(\omega,0)}\allowbreak{}并非极限，只是左变量为极限序数\allowbreak{}\NovelInline{I00574}{\omega}\allowbreak{}、右变量为0的特殊情况。今天咱们就顺着二元Veblen函数的规则，一步步走到\allowbreak{}\NovelInline{I00575}{\varphi(\omega,1)}\allowbreak{}，再冲击\allowbreak{}\NovelInline{I00576}{\varphi(\omega+1,0)}\allowbreak{}，这两个层级的迭代逻辑，完全贴合咱们之前学的不动点、基本列规则，谁有疑惑随时打断。''

那在fgh里面怎么展开\allowbreak{}\NovelInline{I00577}{\varphi(\omega,1)}\allowbreak{}？牢大问：

弗兰德：\allowbreak{}\NovelInline{I00578}{f_{\varphi(\omega,1)}(3) = f_{\varphi(3,\varphi(\omega,0)+1)}(3) = f_{\varphi(2,\varphi(2,\varphi(2,\varphi(\omega,0)+1))}(3) =}\allowbreak{}

\allowbreak{}\NovelInline{I00579}{= f_{\varphi(2,\varphi(2,\varphi(1,\varphi(1,\varphi(1,\varphi(\omega,0)+1))))}(3)}\allowbreak{}

\allowbreak{}\NovelInline{I00580}{= f_{\varphi(2,\varphi(2,\varphi(1,\varphi(1,\varphi(\omega,0)^{\varphi(\omega,0)^{\varphi(\omega,0)}}))))}(3)}\allowbreak{}

\allowbreak{}\NovelInline{I00581}{= f_{\varphi(2,\varphi(2,\varphi(1,\varphi(1,\varphi(\omega,0)^{\varphi(\omega,0)^{\varphi(3,0)}}))))}(3)}\allowbreak{}

\allowbreak{}\NovelInline{I00582}{= f_{\varphi(2,\varphi(2,\varphi(1,\varphi(1,\varphi(\omega,0)^{\varphi(\omega,0)^{\varphi(2,\varphi(2,\varphi(2,0)}}))))}(3)}\allowbreak{}

\allowbreak{}\NovelInline{I00583}{= f_{\varphi(2,\varphi(2,\varphi(1,\varphi(1,\varphi(\omega,0)^{\varphi(\omega,0)^{\varphi(2,\varphi(2,\varphi(1,\varphi(1,\varphi(1,0)))))}}))))}(3)}\allowbreak{}

\allowbreak{}\NovelInline{I00584}{= f_{\varphi(2,\varphi(2,\varphi(1,\varphi(1,\varphi(\omega,0)^{\varphi(\omega,0)^{\varphi(2,\varphi(2,\varphi(1,\varphi(1,\omega^{\omega^\omega}))))}}))))}(3)}\allowbreak{}

\allowbreak{}\NovelInline{I00585}{= f_{\varphi(2,\varphi(2,\varphi(1,\varphi(1,\varphi(\omega,0)^{\varphi(\omega,0)^{\varphi(2,\varphi(2,\varphi(1,\varphi(1,\omega^{(\omega^2\cdot 2 + \omega\cdot 2 + 3)}))))}}))))}(3)}\allowbreak{}

牢大：谢谢老师，这么说我还回顾了以前的功法，但是我有个疑惑：在展开3次的时候，\allowbreak{}\NovelInline{I00586}{\varphi(1,\varphi(\omega,0)+1)}\allowbreak{}为什么等于\allowbreak{}\NovelInline{I00587}{\varphi(\omega,0)^{\varphi(\omega,0)^{\varphi(\omega,0)}}}\allowbreak{}

弗兰德：我问你\allowbreak{}\NovelInline{I00588}{\varphi(\omega,0)}\allowbreak{}是不是\allowbreak{}\NovelInline{I00589}{\varepsilon}\allowbreak{}序数？

牢大：\allowbreak{}\NovelInline{I00590}{\varepsilon_{\varphi(\omega,0)} = \varphi(\omega,0)}\allowbreak{}，当然是

弗兰德：那么它的下一个\allowbreak{}\NovelInline{I00591}{\varepsilon}\allowbreak{}序数你觉得怎么求

牢大：懂了，指数塔的极限对吧

弗兰德：非常正确！

弗兰德指尖滑动，将\allowbreak{}\NovelInline{I00592}{\varphi(\omega,0)}\allowbreak{}的右变量从0替换为1，脉络上浮现出\allowbreak{}\NovelInline{I00593}{\varphi(\omega,1)}\allowbreak{}的标识：``首先明确核心定义：\allowbreak{}\NovelInline{I00594}{\varphi(\omega,1)}\allowbreak{}是全体映射\allowbreak{}\NovelInline{I00595}{\alpha \mapsto \varphi(n,\alpha)}\allowbreak{}（n取任意正整数）的第二个共同不动点，而\allowbreak{}\NovelInline{I00596}{\varphi(\omega,0)}\allowbreak{}是第一个共同不动点，这是二者最本质的区别。''

牢大立刻举手，结合实战感悟发问：``院长，我之前短暂停留在\allowbreak{}\NovelInline{I00597}{\varphi(\omega,0)}\allowbreak{}，能感觉到自身序数波动不受任何有限\allowbreak{}\NovelInline{I00598}{\varphi(n,\alpha)}\allowbreak{}影响，也就是您说的'对所有n都是不动点'，那既然已经是不动点了，该怎么跳出这个层级，到达\allowbreak{}\NovelInline{I00599}{\varphi(\omega,1)}\allowbreak{}？直接迭代不行吗？''

``问得精准，这正是关键。''弗兰德赞许点头，指尖勾勒出跳出不动点的逻辑，``不动点的特性就是\allowbreak{}\NovelInline{I00600}{\varphi(n,\varphi(\omega,0)) = \varphi(\omega,0)}\allowbreak{}，直接迭代永远停留在原地，所以必须先取后继打破不动点状态，再做极限收敛。我把推导过程和基本列写清楚，你们对照着理解：''

\allowbreak{}\NovelInline{I00601}{\varphi(\omega,1)}\allowbreak{}核心规则

\begin{enumerate}
\def\labelenumi{\arabic{enumi}.}
\item
  生成逻辑：先对\allowbreak{}\NovelInline{I00602}{\varphi(\omega,0)}\allowbreak{}取后继，得到\allowbreak{}\NovelInline{I00603}{\varphi(\omega,0)+1}\allowbreak{}，以此跳出不动点；再将任意有限n对应的映射\allowbreak{}\NovelInline{I00604}{\varphi(n,\alpha)}\allowbreak{}作用于这个后继，得到\allowbreak{}\NovelInline{I00605}{\varphi(n,\varphi(\omega,0)+1)}\allowbreak{}；取上确界，即为\allowbreak{}\NovelInline{I00606}{\varphi(\omega,1)}\allowbreak{}。
\item
  完整极限式：\allowbreak{}\NovelInline{I00607}{\varphi(\omega,1) = \sup\{\varphi(0,\varphi(\omega,0)+1), \varphi(1,\varphi(\omega,0)+1), \varphi(2,\varphi(\omega,0)+1), \ldots\}}\allowbreak{}
\item
  递归基本列：\allowbreak{}\NovelInline{I00608}{\varphi(\omega,1)[n] = \varphi(n,\varphi(\omega,0)+1)}\allowbreak{}（n为自然数，依次取0、1、2\ldots\ldots 对应生成基本列项）
\end{enumerate}

唐三皱了皱眉，紧接着发问：``院长，按照这个规则，\allowbreak{}\NovelInline{I00609}{\sup\{\varphi(n,\varphi(\omega,0)+1)\}}\allowbreak{}和\allowbreak{}\NovelInline{I00610}{\varphi(\omega,\varphi(\omega,0)+1)}\allowbreak{}是不是同一个序数？我看资料里特意标注了不等号，这两者的差距在哪？''

``这个疑问抓得很细，是很多人容易混淆的误区。''弗兰德指尖分开两条迭代脉络，一条指向\allowbreak{}\NovelInline{I00611}{\varphi(\omega,1)}\allowbreak{}，一条指向更高阶的\allowbreak{}\NovelInline{I00612}{\varphi(\omega,\varphi(\omega,0)+1)}\allowbreak{}，``前者是对有限n取极限，左变量始终局限在有限层级，最终收敛到第二个共同不动点；后者是直接把\allowbreak{}\NovelInline{I00613}{\varphi(\omega,0)}\allowbreak{}的右变量替换为\allowbreak{}\NovelInline{I00614}{\varphi(\omega,0)+1}\allowbreak{}，左变量保持\allowbreak{}\NovelInline{I00615}{\omega}\allowbreak{}不变，迭代层级远超\allowbreak{}\NovelInline{I00616}{\varphi(\omega,1)}\allowbreak{}，二者绝不能等同。简单来说，\allowbreak{}\NovelInline{I00617}{\varphi(\omega,1)}\allowbreak{}是稳步进阶，\allowbreak{}\NovelInline{I00618}{\varphi(\omega,\varphi(\omega,0)+1)}\allowbreak{}是越级跳跃，咱们现阶段只吃透前者即可。''

弗兰德：你要明白\allowbreak{}\NovelInline{I00619}{\varphi(\omega,0) = \sup\{\varphi(n,0), n < \omega\}}\allowbreak{}，很多人把\allowbreak{}\NovelInline{I00620}{\varphi(\omega,\varphi(\omega,0)+1)}\allowbreak{}和\allowbreak{}\NovelInline{I00621}{\varphi(\omega,1)}\allowbreak{}搞混了，举个例子，\allowbreak{}\NovelInline{I00622}{\sup\{\varphi(n,\varphi(\omega,0)+1)\}}\allowbreak{}也就是\allowbreak{}\NovelInline{I00623}{\varepsilon_1}\allowbreak{}，而\allowbreak{}\NovelInline{I00624}{\varphi(\omega,\varphi(\omega,0)+1)}\allowbreak{}是\allowbreak{}\NovelInline{I00625}{\varepsilon_{(\varepsilon_0+1)}}\allowbreak{}

卢本伟清冷开口，补充提问：``那如果继续往右变量迭代，\allowbreak{}\NovelInline{I00626}{\varphi(\omega,2)}\allowbreak{}、\allowbreak{}\NovelInline{I00627}{\varphi(\omega,3)}\allowbreak{}的规则是不是和\allowbreak{}\NovelInline{I00628}{\varphi(\omega,1)}\allowbreak{}一致？只是把起始点换成前一个序数？''

``没错，完全通用。''弗兰德肯定答复，``\allowbreak{}\NovelInline{I00629}{\varphi(\omega,2)}\allowbreak{}是第三个共同不动点，基本列为\allowbreak{}\NovelInline{I00630}{\varphi(\omega,2)[n] = \varphi(n,\varphi(\omega,1)+1)}\allowbreak{}；以此类推，任意\allowbreak{}\NovelInline{I00631}{\varphi(\omega,\alpha+1)}\allowbreak{}都遵循这套规则，极限序数\allowbreak{}\NovelInline{I00632}{\alpha}\allowbreak{}对应的\allowbreak{}\NovelInline{I00633}{\varphi(\omega,\alpha)}\allowbreak{}，基本列则是\allowbreak{}\NovelInline{I00634}{\varphi(\omega,\alpha)[n] = \varphi(\omega,\alpha[n])}\allowbreak{}，这就是二元Veblen函数左变量为\allowbreak{}\NovelInline{I00635}{\omega}\allowbreak{}时的统一逻辑。''

讲透\allowbreak{}\NovelInline{I00636}{\varphi(\omega,1)}\allowbreak{}后，弗兰德将脉络延伸至更高阶的\allowbreak{}\NovelInline{I00637}{\varphi(\omega+1,0)}\allowbreak{}，语气愈发凝重：``突破\allowbreak{}\NovelInline{I00638}{\varphi(\omega,\alpha)}\allowbreak{}的全系列后，咱们就要冲击左变量为后继序数\allowbreak{}\NovelInline{I00639}{\omega+1}\allowbreak{}的初始态------\allowbreak{}\NovelInline{I00640}{\varphi(\omega+1,0)}\allowbreak{}，这个序数是\allowbreak{}\NovelInline{I00641}{\alpha \mapsto \varphi(\omega,\alpha)}\allowbreak{}映射的第一个不动点，也是超越全体\allowbreak{}\NovelInline{I00642}{\varphi(\omega,\alpha)}\allowbreak{}的关键门槛，它的生成逻辑，和咱们之前学的\allowbreak{}\NovelInline{I00643}{\varphi(2,0)}\allowbreak{}、\allowbreak{}\NovelInline{I00644}{\varphi(3,0)}\allowbreak{}完全同源。''

诺当当即提出疑问：``院长，之前您讲过\allowbreak{}\NovelInline{I00645}{\varphi(n+1,0)}\allowbreak{}是映射\allowbreak{}\NovelInline{I00646}{\alpha \mapsto \varphi(n,\alpha)}\allowbreak{}的第一个不动点，基本列是从零开始反复代入映射取极限，那\allowbreak{}\NovelInline{I00647}{\varphi(\omega+1,0)}\allowbreak{}是不是也是这个逻辑？只是把n换成了\allowbreak{}\NovelInline{I00648}{\omega}\allowbreak{}？''

``丝毫不差，诺当的总结很到位，这就是Veblen函数的共性。''弗兰德指尖勾勒出\allowbreak{}\NovelInline{I00649}{\varphi(\omega+1,0)}\allowbreak{}的迭代链条，把基本列和定义清晰呈现，``不管左变量是有限数还是\allowbreak{}\NovelInline{I00650}{\omega}\allowbreak{}，只要是\allowbreak{}\NovelInline{I00651}{\alpha+1}\allowbreak{}的形式，\allowbreak{}\NovelInline{I00652}{\varphi(\alpha+1,0)}\allowbreak{}的生成规则完全一致，咱们对照资料看细节：''

\allowbreak{}\NovelInline{I00653}{\varphi(\omega+1,0)}\allowbreak{}核心规则

\begin{enumerate}
\def\labelenumi{\arabic{enumi}.}
\item
  本质定义：\allowbreak{}\NovelInline{I00654}{\varphi(\omega+1,0) = 1\text{st}\,\alpha \mapsto \varphi(\omega,\alpha)\,\text{fp.}}\allowbreak{}（映射\allowbreak{}\NovelInline{I00655}{\alpha \mapsto \varphi(\omega,\alpha)}\allowbreak{}的第一个不动点）
\item
  完整极限式：\allowbreak{}\NovelInline{I00656}{\varphi(\omega+1,0) = \sup\{0, \varphi(\omega,0), \varphi(\omega,\varphi(\omega,0)), \varphi(\omega,\varphi(\omega,\varphi(\omega,0))), \ldots\}}\allowbreak{}
\item
  递归基本列：\allowbreak{}\NovelInline{I00657}{\varphi(\omega+1,0)[0] = 0}\allowbreak{}，\allowbreak{}\NovelInline{I00658}{\varphi(\omega+1,0)[n+1] = \varphi(\omega,\varphi(\omega+1,0)[n])}\allowbreak{}
\end{enumerate}

蔡须坤握着篮球，思索片刻发问：``院长，我之前用篮球校准序数波动时发现，不动点的迭代都是'闭环'，那\allowbreak{}\NovelInline{I00659}{\varphi(\omega+1,0)}\allowbreak{}作为\allowbreak{}\NovelInline{I00660}{\varphi(\omega,\alpha)}\allowbreak{}的不动点，是不是意味着所有\allowbreak{}\NovelInline{I00661}{\varphi(\omega,\alpha)}\allowbreak{}都被它覆盖？到达这个层级后，是不是就能随意驾驭\allowbreak{}\NovelInline{I00662}{\varphi(\omega,0)}\allowbreak{}、\allowbreak{}\NovelInline{I00663}{\varphi(\omega,1)}\allowbreak{}的力量？''

``正是如此。''弗兰德点头肯定，``\allowbreak{}\NovelInline{I00664}{\varphi(\omega+1,0)}\allowbreak{}已经跳出了左变量为\allowbreak{}\NovelInline{I00665}{\omega}\allowbreak{}的框架，不管右变量怎么迭代，\allowbreak{}\NovelInline{I00666}{\varphi(\omega,\alpha)}\allowbreak{}的所有层级都在它的覆盖范围内。就像你们之前对抗人面魔蛛，\allowbreak{}\NovelInline{I00667}{\varphi(\omega,0)}\allowbreak{}的力量已经极强，但在\allowbreak{}\NovelInline{I00668}{\varphi(\omega+1,0)}\allowbreak{}面前，只是迭代链条里的一个基础项。而且后续的\allowbreak{}\NovelInline{I00669}{\varphi(\omega+1,\alpha)}\allowbreak{}、\allowbreak{}\NovelInline{I00670}{\varphi(\omega+2,\alpha)}\allowbreak{}，都是以此为基础，遵循'后一个是前一个映射的不动点'规则，和\allowbreak{}\NovelInline{I00671}{\varphi(2,\alpha)}\allowbreak{}、\allowbreak{}\NovelInline{I00672}{\varphi(3,\alpha)}\allowbreak{}的进阶逻辑完全一致。''

奥斯卡也举手问道：``那这三个层级的基本列，有没有统一的记忆逻辑？我们实战运用时，能更快对应上迭代节奏。''

弗兰德整理脉络，给出总结：``我给你们梳理成三句口诀，贴合定义好记好用：

\allowbreak{}\NovelInline{I00673}{\varphi(\omega,0)}\allowbreak{}基本列：左\allowbreak{}\NovelInline{I00674}{\omega}\allowbreak{}右0，取有限n，\allowbreak{}\NovelInline{I00675}{\varphi(\omega,0)[n] = \varphi(n,0)}\allowbreak{}；

\allowbreak{}\NovelInline{I00676}{\varphi(\omega,1)}\allowbreak{}基本列：先跳后继，再取n限，\allowbreak{}\NovelInline{I00677}{\varphi(\omega,1)[n] = \varphi(n,\varphi(\omega,0)+1)}\allowbreak{}；

\allowbreak{}\NovelInline{I00678}{\varphi(\omega+1,0)}\allowbreak{}基本列：从零代入，反复迭代，\allowbreak{}\NovelInline{I00679}{\varphi(\omega+1,0)[n+1] = \varphi(\omega,\text{前一项})}\allowbreak{}。

这三个层级，是从有限Veblen函数迈向\allowbreak{}\NovelInline{I00680}{\Gamma_0}\allowbreak{}的必经之路，\allowbreak{}\NovelInline{I00681}{\Gamma_0}\allowbreak{}正是映射\allowbreak{}\NovelInline{I00682}{\alpha \mapsto \varphi(\alpha,0)}\allowbreak{}的第一个不动点，吃透这一段迭代，后续理解\allowbreak{}\NovelInline{I00683}{\Gamma_0}\allowbreak{}就水到渠成。''

弗兰德指尖轻挥，虚空脉络再度延展，从\allowbreak{}\NovelInline{I00684}{\varphi(\omega+1,0)}\allowbreak{}一路延伸出无数分支，\allowbreak{}\NovelInline{I00685}{\varphi(\omega+2,0)}\allowbreak{}、\allowbreak{}\NovelInline{I00686}{\varphi(\omega+3,0)}\allowbreak{}、\allowbreak{}\NovelInline{I00687}{\varphi(\omega\cdot 2,0)}\allowbreak{}、\allowbreak{}\NovelInline{I00688}{\varphi(\omega^2,0)}\allowbreak{}、\allowbreak{}\NovelInline{I00689}{\varphi(\omega^\omega,0)}\allowbreak{}\ldots\ldots 层层叠叠的序数符号铺满半空，看得众人目不暇接，不少人脸上露出困惑之色。

``院长，这些序数看着密密麻麻，迭代规律和\allowbreak{}\NovelInline{I00690}{\varphi(\omega+1,0)}\allowbreak{}一模一样吗？''唐三率先开口，眉头微蹙，``左变量不断拔高，从\allowbreak{}\NovelInline{I00691}{\omega+1}\allowbreak{}到\allowbreak{}\NovelInline{I00692}{\omega\cdot 2}\allowbreak{}，再到\allowbreak{}\NovelInline{I00693}{\omega}\allowbreak{}的更高阶幂次，难道都是沿用'后继序数取不动点、极限序数取基本列'的规则？''

``唐三问到了点子上。''弗兰德颔首，指尖定格在\allowbreak{}\NovelInline{I00694}{\varphi(\omega\cdot 2,0)}\allowbreak{}上，``没错，这套逻辑是通用的。\allowbreak{}\NovelInline{I00695}{\varphi(\omega+2,0)}\allowbreak{}是\allowbreak{}\NovelInline{I00696}{\alpha \mapsto \varphi(\omega+1,\alpha)}\allowbreak{}的第一个不动点，\allowbreak{}\NovelInline{I00697}{\varphi(\omega\cdot 2,0)}\allowbreak{}是全体\allowbreak{}\NovelInline{I00698}{\varphi(\omega+n,0)}\allowbreak{}的极限，基本列就是\allowbreak{}\NovelInline{I00699}{\varphi(\omega\cdot 2,0)[n] = \varphi(\omega+n,0)}\allowbreak{}；再往上，\allowbreak{}\NovelInline{I00700}{\varphi(\omega^2,0)}\allowbreak{}是\allowbreak{}\NovelInline{I00701}{\varphi(\omega\cdot n,0)}\allowbreak{}的极限，\allowbreak{}\NovelInline{I00702}{\varphi(\omega^\omega,0)}\allowbreak{}是\allowbreak{}\NovelInline{I00703}{\varphi(\omega^n,0)}\allowbreak{}的极限，哪怕迭代到\allowbreak{}\NovelInline{I00704}{\varphi(\varphi(\omega,0),0)}\allowbreak{}，规则依旧不变------左变量每提升一阶，就是对前一阶映射取不动点。''

他顿了顿，细化讲解：``比如\allowbreak{}\NovelInline{I00705}{\varphi(\omega+2,0)}\allowbreak{}，就是映射\allowbreak{}\NovelInline{I00706}{\alpha \mapsto \varphi(\omega+1,\alpha)}\allowbreak{}的第一个不动点，基本列和\allowbreak{}\NovelInline{I00707}{\varphi(\omega+1,0)}\allowbreak{}完全一致：\allowbreak{}\NovelInline{I00708}{\varphi(\omega+2,0)[0] = 0}\allowbreak{}，\allowbreak{}\NovelInline{I00709}{\varphi(\omega+2,0)[n+1] = \varphi(\omega+1,\varphi(\omega+2,0)[n])}\allowbreak{}；再往上\allowbreak{}\NovelInline{I00710}{\varphi(\omega+3,0)}\allowbreak{}，就是\allowbreak{}\NovelInline{I00711}{\alpha \mapsto \varphi(\omega+2,\alpha)}\allowbreak{}的不动点，以此类推，左变量每加1，就是对前一层级映射取不动点。''

蔡须坤握着篮球，指尖无意识转动，疑惑开口：``那左变量变成极限序数呢？比如\allowbreak{}\NovelInline{I00712}{\omega\cdot 2}\allowbreak{}、\allowbreak{}\NovelInline{I00713}{\omega^\omega}\allowbreak{}，这类不是后继序数的情况，基本列该怎么取？之前对抗人面魔蛛时，我校准波动碰到极限序数，总会出现断层。''

``这便是二元Veblen函数对极限序数的通用规则。''弗兰德指尖点向虚空的\allowbreak{}\NovelInline{I00714}{\varphi(\omega\cdot 2,0)}\allowbreak{}标识，``只要左变量是极限序数\allowbreak{}\NovelInline{I00715}{\alpha}\allowbreak{}，\allowbreak{}\NovelInline{I00716}{\varphi(\alpha,0)[n] = \varphi(\alpha[n],0)}\allowbreak{}。拿\allowbreak{}\NovelInline{I00717}{\varphi(\omega\cdot 2,0)}\allowbreak{}举例，它的左变量\allowbreak{}\NovelInline{I00718}{\alpha = \omega\cdot 2}\allowbreak{}，基本列\allowbreak{}\NovelInline{I00719}{\alpha[n] = \omega+n}\allowbreak{}，所以\allowbreak{}\NovelInline{I00720}{\varphi(\omega\cdot 2,0)[n] = \varphi(\omega+n,0)}\allowbreak{}，也就是取\allowbreak{}\NovelInline{I00721}{\varphi(\omega+1,0)}\allowbreak{}、\allowbreak{}\NovelInline{I00722}{\varphi(\omega+2,0)}\allowbreak{}\ldots\ldots 的极限；同理，\allowbreak{}\NovelInline{I00723}{\varphi(\omega^2,0)[n] = \varphi(\omega\cdot n,0)}\allowbreak{}，\allowbreak{}\NovelInline{I00724}{\varphi(\omega^\omega,0)[n] = \varphi(\omega^n,0)}\allowbreak{}，哪怕是\allowbreak{}\NovelInline{I00725}{\varphi(\varphi(\omega,0),0)}\allowbreak{}，基本列也是\allowbreak{}\NovelInline{I00726}{\varphi(\varphi(n,0),0)}\allowbreak{}，始终遵循'左极限取基本列项，再套Veblen函数'的逻辑。''

牢大挠了挠头，满脸疑惑地起身：``老师，我越听越迷糊，咱们一直在给左变量'升级'，从有限数到\allowbreak{}\NovelInline{I00727}{\omega}\allowbreak{}，再到\allowbreak{}\NovelInline{I00728}{\omega+1}\allowbreak{}、\allowbreak{}\NovelInline{I00729}{\omega\cdot 2}\allowbreak{}，什么时候是个头？难道要无限迭代下去？而且这些高阶序数，和咱们要突破的\allowbreak{}\NovelInline{I00730}{\Gamma_0}\allowbreak{}到底有什么直接联系？''

卢本韦也淡淡开口，提出众人的共性困惑：``按照之前的规律，\allowbreak{}\NovelInline{I00731}{\varphi(\alpha,0)}\allowbreak{}的\allowbreak{}\NovelInline{I00732}{\alpha}\allowbreak{}越大，层级越高，可\allowbreak{}\NovelInline{I00733}{\alpha}\allowbreak{}本身也能无限迭代，那\allowbreak{}\NovelInline{I00734}{\Gamma_0}\allowbreak{}难道是这些\allowbreak{}\NovelInline{I00735}{\varphi(\alpha,0)}\allowbreak{}的终点？可序数没有最大只有更大，怎么会存在这样一个'终点不动点'？''

诺当紧随其后补充：``还有，咱们之前学的都是左变量固定、迭代右变量，或是左变量小幅进阶，现在左变量本身都成了Veblen函数的结果，比如\allowbreak{}\NovelInline{I00736}{\varphi(\varphi(\omega,0),0)}\allowbreak{}，这种嵌套结构，基本列该怎么写？实战中又该怎么催动力量？''

一时间，课堂上疑问四起，众人盯着虚空中繁复的序数脉络，只觉得晦涩难懂，毕竟此前接触的最高层级不过\allowbreak{}\NovelInline{I00737}{\varphi(\omega,0)}\allowbreak{}，骤然面对如此高阶的迭代逻辑，难免难以消化。

弗兰德不慌不忙，抬手压下众人的议论声，指尖在脉络最顶端画出一个醒目的符号------\allowbreak{}\NovelInline{I00738}{\Gamma_0}\allowbreak{}，周身泛起淡淡的金色光晕，将这个Feferman-Schütte序数的核心奥义凸显出来：``大家的困惑，恰恰是突破\allowbreak{}\NovelInline{I00739}{\Gamma_0}\allowbreak{}的关键。你们记住，\allowbreak{}\NovelInline{I00740}{\Gamma_0}\allowbreak{}不是普通的高阶序数，它是映射\allowbreak{}\NovelInline{I00741}{\alpha \mapsto \varphi(\alpha,0)}\allowbreak{}的第一个不动点，这是定义核心，也是破解所有疑惑的钥匙。''

见众人依旧茫然，弗兰德换了更通俗的说法：``咱们之前所有的进阶，都是在固定\allowbreak{}\NovelInline{I00742}{\varphi(\alpha,0)}\allowbreak{}里的\allowbreak{}\NovelInline{I00743}{\alpha}\allowbreak{}，然后迭代右变量；或是慢慢提升\allowbreak{}\NovelInline{I00744}{\alpha}\allowbreak{}，再取新的不动点。而\allowbreak{}\NovelInline{I00745}{\Gamma_0}\allowbreak{}的特殊之处在于，\allowbreak{}\NovelInline{I00746}{\Gamma_0 = \varphi(\Gamma_0,0)}\allowbreak{}，也就是说，把\allowbreak{}\NovelInline{I00747}{\alpha}\allowbreak{}换成\allowbreak{}\NovelInline{I00748}{\Gamma_0}\allowbreak{}本身，\allowbreak{}\NovelInline{I00749}{\varphi(\alpha,0)}\allowbreak{}的结果还是\allowbreak{}\NovelInline{I00750}{\Gamma_0}\allowbreak{}，这就是'自反不动点'，也是二元Veblen函数的极限。''

奥斯卡立刻举手追问：``院长，那\allowbreak{}\NovelInline{I00751}{\Gamma_0}\allowbreak{}的基本列该怎么写？和之前的\allowbreak{}\NovelInline{I00752}{\varphi(\omega+1,0)}\allowbreak{}一样从零迭代吗？既然它等于自身的Veblen函数值，岂不是永远跳不出这个闭环？''

``问得好，这就是\allowbreak{}\NovelInline{I00753}{\Gamma_0}\allowbreak{}突破的核心法门。''弗兰德指尖轻点\allowbreak{}\NovelInline{I00754}{\Gamma_0}\allowbreak{}，勾勒出它的极限式与基本列，``\allowbreak{}\NovelInline{I00755}{\Gamma_0}\allowbreak{}是全体嵌套\allowbreak{}\NovelInline{I00756}{\varphi(\alpha,0)}\allowbreak{}的上确界，它的完整极限式是：\allowbreak{}\NovelInline{I00757}{\Gamma_0 = \sup\{0, \varphi(0,0), \varphi(\varphi(0,0),0), \varphi(\varphi(\varphi(0,0),0),0), \ldots\}}\allowbreak{}。''

他顿了顿，逐条拆解规则，让众人更容易理解：

\begin{enumerate}
\def\labelenumi{\arabic{enumi}.}
\item
  生成逻辑：从0开始，反复将当前结果代入\allowbreak{}\NovelInline{I00758}{\alpha \mapsto \varphi(\alpha,0)}\allowbreak{}的映射，不断生成\allowbreak{}\NovelInline{I00759}{\varphi(0,0)}\allowbreak{}、\allowbreak{}\NovelInline{I00760}{\varphi(\varphi(0,0),0)}\allowbreak{}、\allowbreak{}\NovelInline{I00761}{\varphi(\varphi(\varphi(0,0),0),0)}\allowbreak{}\ldots\ldots 无限迭代下去，最终收敛的极限就是\allowbreak{}\NovelInline{I00762}{\Gamma_0}\allowbreak{}。这意味着，不管之前的\allowbreak{}\NovelInline{I00763}{\varphi(\alpha,0)}\allowbreak{}迭代到多高，只要还在有限嵌套的范畴内，就始终在\allowbreak{}\NovelInline{I00764}{\Gamma_0}\allowbreak{}的基本列里。
\item
  递归基本列：\allowbreak{}\NovelInline{I00765}{\Gamma_0[0] = 0}\allowbreak{}，\allowbreak{}\NovelInline{I00766}{\Gamma_0[n+1] = \varphi(\Gamma_0[n],0)}\allowbreak{}。和\allowbreak{}\NovelInline{I00767}{\varphi(\omega+1,0)}\allowbreak{}的迭代逻辑同源，但区别在于，\allowbreak{}\NovelInline{I00768}{\varphi(\omega+1,0)}\allowbreak{}是固定左变量为\allowbreak{}\NovelInline{I00769}{\omega+1}\allowbreak{}、迭代右变量，而\allowbreak{}\NovelInline{I00770}{\Gamma_0}\allowbreak{}是左变量直接跟着迭代结果走，形成自我嵌套。
\item
  层级意义：\allowbreak{}\NovelInline{I00771}{\Gamma_0}\allowbreak{}彻底超越了此前所有\allowbreak{}\NovelInline{I00772}{\varphi(\alpha,\beta)}\allowbreak{}的范畴，不管是\allowbreak{}\NovelInline{I00773}{\varphi(\omega,0)}\allowbreak{}、\allowbreak{}\NovelInline{I00774}{\varphi(\omega+1,0)}\allowbreak{}，还是\allowbreak{}\NovelInline{I00775}{\varphi(\omega^\omega,0)}\allowbreak{}、\allowbreak{}\NovelInline{I00776}{\varphi(\varphi(\omega,0),0)}\allowbreak{}，全都是\allowbreak{}\NovelInline{I00777}{\Gamma_0}\allowbreak{}基本列里的某一项，都被\allowbreak{}\NovelInline{I00778}{\Gamma_0}\allowbreak{}稳稳压制。这也是为什么\allowbreak{}\NovelInline{I00779}{\Gamma_0}\allowbreak{}被称为Feferman-Schütte序数，是弱预言论的极限，更是你们现阶段能触及的顶级扽力门槛。
\end{enumerate}

蔡须坤握着篮球的手微微收紧，恍然大悟却又带着新的疑问：``院长，我明白了！之前咱们是'推着\allowbreak{}\NovelInline{I00780}{\alpha}\allowbreak{}往前走'，而\allowbreak{}\NovelInline{I00781}{\Gamma_0}\allowbreak{}是'自己推着自己走'。那实战中想要突破至\allowbreak{}\NovelInline{I00782}{\Gamma_0}\allowbreak{}，是不是要先把自身序数迭代到\allowbreak{}\NovelInline{I00783}{\varphi(\alpha,0)}\allowbreak{}的极限，再让自身波动形成自反闭环，让\allowbreak{}\NovelInline{I00784}{\alpha}\allowbreak{}等于自身层级？''

``完全正确！''弗兰德眼中闪过赞许，``想要突破\allowbreak{}\NovelInline{I00785}{\Gamma_0}\allowbreak{}，不能一味追求左变量拔高，而是要感悟自反不动点的真谛。你们吸收了馃槔的力量，又经过人面魔蛛一战的洗礼，自身序数根基已经足够扎实，只要能吃透'\allowbreak{}\NovelInline{I00786}{\alpha \mapsto \varphi(\alpha,0)}\allowbreak{}不动点'的逻辑，让自身扽力不再依赖外部迭代，而是形成自我循环，就能触摸\allowbreak{}\NovelInline{I00787}{\Gamma_0}\allowbreak{}的门槛。''

牢大一拍大腿，此前的困惑烟消云散：``原来如此！之前总想着往上冲层级，却忽略了不动点的自反性。\allowbreak{}\NovelInline{I00788}{\Gamma_0}\allowbreak{}不是靠硬冲来的，是靠自身波动和\allowbreak{}\NovelInline{I00789}{\varphi(\alpha,0)}\allowbreak{}映射完全契合，达到'我即是映射，映射即是我'的状态！''

卢本韦、诺当、唐三、奥斯卡四人也纷纷点头，眼底闪过明悟，虚空中的\allowbreak{}\NovelInline{I00790}{\Gamma_0}\allowbreak{}符号仿佛与众人的气息产生了微弱共鸣，周身的序数纹路微微发烫，显然是触及了突破的契机。

弗兰德看着众人的状态，沉声收尾：``从\allowbreak{}\NovelInline{I00791}{\varphi(\omega+1,0)}\allowbreak{}到\allowbreak{}\NovelInline{I00792}{\Gamma_0}\allowbreak{}，本质是从'被动迭代'到'自主闭环'的蜕变。\allowbreak{}\NovelInline{I00793}{\varphi(\omega+1,0)}\allowbreak{}只是左变量进阶的一个节点，而\allowbreak{}\NovelInline{I00794}{\Gamma_0}\allowbreak{}是整个二元Veblen函数的终极闭环。接下来你们要做的，就是结合今日所学，在修炼中反复感悟\allowbreak{}\NovelInline{I00795}{\Gamma_0}\allowbreak{}的基本列与不动点特性，借助馃槔的共鸣之力，彻底打通层级壁垒，完成\allowbreak{}\NovelInline{I00796}{\Gamma_0}\allowbreak{}的突破。''

\chapter{第十五章：多元φ函数}\label{ux7b2cux5341ux4e94ux7ae0ux591aux5143ux3c6ux51fdux6570}

课堂散去，众人并未停歇，纷纷寻了学院的静修室闭关。馃槔的共鸣之力萦绕周身，将课堂所学的理论转化为实打实的战力突破。唐三、牢大等人各自盘膝而坐，指尖勾勒 fgh 迭代公式，结合自身层级展开实操，每一次迭代都牵动着周身的序数波动，朝着 \allowbreak{}\NovelInline{I00797}{\Gamma_0}\allowbreak{} 的门槛稳步迈进。

翌日清晨，静修室的光晕缓缓散去，唐三与牢大周身萦绕着淡淡的 \allowbreak{}\NovelInline{I00798}{\Gamma_0}\allowbreak{} 光晕，眼神澄澈而深邃，已然完成终极突破。其余四人也气息大涨，小队整体实力实现质的飞跃。弗兰德伫立在远处，看着众人的蜕变，嘴角微微上扬------这支小队，已然触摸到了序数力量的全新境界，未来的征程，必将更加强悍。

\begin{center}\rule{0.5\linewidth}{0.5pt}\end{center}

率先进入状态的是牢大，他本就肉身强悍、悟性直白，最擅长将晦涩理论转化为蛮力突破。此刻他紧盯 \allowbreak{}\NovelInline{I00799}{\Gamma_0}\allowbreak{} 的基本列，结合自身 \allowbreak{}\NovelInline{I00800}{\varphi(4,0)}\allowbreak{} 的根基，在脑海中推演 fgh 展开式，摒弃杂念，只专注于自反不动点的闭环。

另一侧的唐三心思缜密，擅长精准把控每一处迭代细节，他以自身 \allowbreak{}\NovelInline{I00801}{\varphi(3,0)}\allowbreak{} 为根基，结合蓝银草的细腻感知，一点点校准序数波动，将 fgh 展开到极致。

他先复盘 \allowbreak{}\NovelInline{I00802}{\varphi(\omega,1)}\allowbreak{} 的 fgh 展开，夯实基础：\allowbreak{}\NovelInline{I00803}{f_{\varphi(\omega,1)} = f_{\varphi(2,\varphi(\omega,0)+1)}}\allowbreak{}，再逐步推进到 \allowbreak{}\NovelInline{I00804}{\varphi(\omega+1,0)}\allowbreak{} 的迭代：\allowbreak{}\NovelInline{I00805}{f_{\varphi(\omega+1,0)} = f_{\varphi(\omega,\varphi(\omega,0))}}\allowbreak{}，最终聚焦 \allowbreak{}\NovelInline{I00806}{\Gamma_0}\allowbreak{} 的自反特性。唐三深知，突破 \allowbreak{}\NovelInline{I00807}{\Gamma_0}\allowbreak{} 的关键是''自我映射''，他缓缓将自身序数代入 \allowbreak{}\NovelInline{I00808}{\alpha \mapsto \varphi(\alpha,0)}\allowbreak{}，让 \allowbreak{}\NovelInline{I00809}{\alpha}\allowbreak{} 等于自身当前层级，反复迭代：\allowbreak{}\NovelInline{I00810}{\Gamma_0[n+1] = \varphi(\Gamma_0[n],0)}\allowbreak{}，每一次代入都让馃槔的共鸣更强一分。

他以自身当前层级 \allowbreak{}\NovelInline{I00811}{\varphi(4,0)}\allowbreak{} 为起点，先迭代 fgh 基础项：

\NovelDisplay{M00150R01}{%
f_{\varphi(4,0)}(3) = f_{\varphi(3,\varphi(3,\varphi(3,0)))}(3)%
}

顺着弗兰德讲授的规则，逐步将左变量迭代至极限序数 \allowbreak{}\NovelInline{I00812}{\omega}\allowbreak{}，再延伸至 \allowbreak{}\NovelInline{I00813}{\varphi(\omega+1,0)}\allowbreak{}，最终触碰 \allowbreak{}\NovelInline{I00814}{\Gamma_0}\allowbreak{} 的自反逻辑。

奥斯卡、蔡须坤、卢本韦、诺当四人也各有突破：奥斯卡依托自身辅助系扽魂，将 \allowbreak{}\NovelInline{I00815}{\zeta(\zeta_0)}\allowbreak{} 的层级与 \allowbreak{}\NovelInline{I00816}{\Gamma_0}\allowbreak{} 基本列结合，fgh 展开式虽不如唐三细腻，却胜在稳定；蔡须坤以篮球为媒介，校准每一次迭代节奏，把 \allowbreak{}\NovelInline{I00817}{\varphi(3,1)}\allowbreak{} 的力量融入自反闭环，突破之势迅猛；卢本韦与诺当则凭借新融入的小队氛围，结合自身 \allowbreak{}\NovelInline{I00818}{\varphi(2,1)}\allowbreak{}、\allowbreak{}\NovelInline{I00819}{\varphi(2,\omega)}\allowbreak{} 根基，稳步感悟不动点真谛，虽未完全突破 \allowbreak{}\NovelInline{I00820}{\Gamma_0}\allowbreak{}，却也彻底吃透 \allowbreak{}\NovelInline{I00821}{\varphi(\omega+1,0)}\allowbreak{} 的层级，为后续突破筑牢根基。

经过 20 多天的吸收，牢大直接突破到 \allowbreak{}\NovelInline{I00822}{\varphi(\varepsilon_0,0)}\allowbreak{} 的境界，唐三、小舞分别达到了 \allowbreak{}\NovelInline{I00823}{\varphi(\omega \cdot 2,0)}\allowbreak{} 以及 \allowbreak{}\NovelInline{I00824}{\varphi(\omega+3,0)}\allowbreak{}，蔡须坤达到了 \allowbreak{}\NovelInline{I00825}{\varphi(\omega^\omega,0)}\allowbreak{}，卢本韦、诺当都到了 \allowbreak{}\NovelInline{I00826}{\varphi(\omega+1,0)}\allowbreak{}，奥斯卡也进入了 \allowbreak{}\NovelInline{I00827}{\varphi(\omega,\omega)}\allowbreak{} 的境界了。

\begin{center}\rule{0.5\linewidth}{0.5pt}\end{center}

之后弗兰德继续授课。

弗兰德看着众人眼中对高阶序数的渴求，缓缓将虚空里的二元 Veblen 脉络延展成更繁复的多维结构，三元、四元乃至任意多元的 \allowbreak{}\NovelInline{I00828}{\varphi}\allowbreak{} 函数符号依次浮现，语气沉稳而郑重：

``你们已经吃透了二元 \allowbreak{}\NovelInline{I00829}{\varphi}\allowbreak{} 函数的核心，从 \allowbreak{}\NovelInline{I00830}{\varphi(\omega,0)}\allowbreak{} 走到 \allowbreak{}\NovelInline{I00831}{\Gamma_0}\allowbreak{}，这只是 Veblen 函数体系的起点。真正的力量藏在多元 Veblen 函数里------它能让我们突破二元的局限，构造出远超 \allowbreak{}\NovelInline{I00832}{\Gamma_0}\allowbreak{} 的序数，甚至触及小 Veblen 序数（SVO）的门槛。今天，我们就从三元 \allowbreak{}\NovelInline{I00833}{\varphi}\allowbreak{} 函数开始，严谨拆解多元 \allowbreak{}\NovelInline{I00834}{\varphi}\allowbreak{} 函数的定义、基本列与不动点逻辑，结合实例让你们吃透它。''

弗兰德指尖轻点，虚空里的二元 \allowbreak{}\NovelInline{I00835}{\varphi(\alpha,\beta)}\allowbreak{} 旁，浮现出三元 \allowbreak{}\NovelInline{I00836}{\varphi(1,0,\gamma)}\allowbreak{} 的标识，\allowbreak{}\NovelInline{I00837}{\gamma}\allowbreak{} 代指一个序数。他先回顾二元的根基，再自然过渡到三元：

``我们先回忆二元 \allowbreak{}\NovelInline{I00838}{\varphi}\allowbreak{} 函数的核心：\allowbreak{}\NovelInline{I00839}{\varphi(\alpha+1,\beta)}\allowbreak{} 是映射 \allowbreak{}\NovelInline{I00840}{\gamma \mapsto \varphi(\alpha,\gamma)}\allowbreak{} 的第 \allowbreak{}\NovelInline{I00841}{1+\beta}\allowbreak{} 个不动点；若 \allowbreak{}\NovelInline{I00842}{\alpha}\allowbreak{} 是极限序数，\allowbreak{}\NovelInline{I00843}{\varphi(\alpha,\beta)}\allowbreak{} 是所有映射 \allowbreak{}\NovelInline{I00844}{\gamma \mapsto \varphi(\alpha[n],\gamma)}\allowbreak{} 的第 \allowbreak{}\NovelInline{I00845}{1+\beta}\allowbreak{} 个共同不动点。而三元 \allowbreak{}\NovelInline{I00846}{\varphi}\allowbreak{} 函数，就是对二元 \allowbreak{}\NovelInline{I00847}{\varphi}\allowbreak{} 函数的'左变量映射'再取不动点------这是多元函数的通用逻辑：每多一个变量，就是对前一层映射再取不动点。''

他顿了顿，指着 \allowbreak{}\NovelInline{I00848}{\varphi(1,0,\gamma)}\allowbreak{} 的定义，结合实例拆解：

``首先，我们定义映射 \allowbreak{}\NovelInline{I00849}{\alpha \mapsto \varphi(\alpha,0)}\allowbreak{}，这个映射的第 \allowbreak{}\NovelInline{I00850}{1+\gamma}\allowbreak{} 个不动点，就是三元 \allowbreak{}\NovelInline{I00851}{\varphi}\allowbreak{} 函数 \allowbreak{}\NovelInline{I00852}{\varphi(1,0,\gamma)}\allowbreak{}，也记作 \allowbreak{}\NovelInline{I00853}{\Gamma_\gamma}\allowbreak{}。我给你们举几个关键实例：''

它是映射 \allowbreak{}\NovelInline{I00854}{\alpha \mapsto \varphi(\alpha,0)}\allowbreak{} 的第一个不动点，也是我们之前突破的 \allowbreak{}\NovelInline{I00855}{\Gamma_0}\allowbreak{}：

\begin{itemize}
\tightlist
\item
  \textbf{基本列}：\allowbreak{}\NovelInline{I00856}{\Gamma_0[0] = 0}\allowbreak{}，\allowbreak{}\NovelInline{I00857}{\Gamma_0[n+1] = \varphi(\Gamma_0[n],0)}\allowbreak{}
\item
  \textbf{极限式}：\allowbreak{}\NovelInline{I00858}{\Gamma_0 = \sup\{0, \varphi(0,0), \varphi(\varphi(0,0),0), \varphi(\varphi(\varphi(0,0),0),0), \dots\}}\allowbreak{}
\end{itemize}

它是映射 \allowbreak{}\NovelInline{I00859}{\alpha \mapsto \varphi(\alpha,0)}\allowbreak{} 的第二个不动点，生成逻辑和 \allowbreak{}\NovelInline{I00860}{\varphi(\omega,1)}\allowbreak{} 完全同源：

\begin{itemize}
\tightlist
\item
  \textbf{定义}：\allowbreak{}\NovelInline{I00861}{\varphi(1,0,1) = \sup\{\varphi(n, \varphi(1,0,0)+1) \mid n < \omega\}}\allowbreak{}
\item
  \textbf{基本列}：\allowbreak{}\NovelInline{I00862}{\varphi(1,0,1)[n] = \varphi(n, \Gamma_0+1)}\allowbreak{}
\item
  \textbf{直观理解}：先跳出 \allowbreak{}\NovelInline{I00863}{\Gamma_0}\allowbreak{} 的不动点状态（取 \allowbreak{}\NovelInline{I00864}{\Gamma_0+1}\allowbreak{}），再对每个有限 \allowbreak{}\NovelInline{I00865}{n}\allowbreak{} 取 \allowbreak{}\NovelInline{I00866}{\varphi(n, \Gamma_0+1)}\allowbreak{} 的极限，得到第二个共同不动点。
\end{itemize}

它是映射 \allowbreak{}\NovelInline{I00867}{\alpha \mapsto \varphi(\alpha,0)}\allowbreak{} 的第三个不动点：

\begin{itemize}
\tightlist
\item
  \textbf{基本列}：\allowbreak{}\NovelInline{I00868}{\varphi(1,0,2)[n] = \varphi(n, \varphi(1,0,1)+1)}\allowbreak{}
\item
  以此类推，\allowbreak{}\NovelInline{I00869}{\varphi(1,0,\gamma)}\allowbreak{} 就是这个映射的第 \allowbreak{}\NovelInline{I00870}{1+\gamma}\allowbreak{} 个不动点，\allowbreak{}\NovelInline{I00871}{\gamma}\allowbreak{} 为极限序数时，\allowbreak{}\NovelInline{I00872}{\varphi(1,0,\gamma)[n] = \varphi(1,0,\gamma[n])}\allowbreak{}。
\end{itemize}

唐三立刻举手，结合之前的 \allowbreak{}\NovelInline{I00873}{\Gamma_0}\allowbreak{} 感悟发问：``院长，那 \allowbreak{}\NovelInline{I00874}{\varphi(1,0,1)}\allowbreak{} 和 \allowbreak{}\NovelInline{I00875}{\Gamma_0}\allowbreak{} 的区别是什么？是不是和 \allowbreak{}\NovelInline{I00876}{\varphi(\omega,1)}\allowbreak{} 与 \allowbreak{}\NovelInline{I00877}{\varphi(\omega,0)}\allowbreak{} 的关系一样，只是从'第一个不动点'进阶到'第二个不动点'？''

``完全正确！''弗兰德赞许点头，``\allowbreak{}\NovelInline{I00878}{\Gamma_0}\allowbreak{} 是 \allowbreak{}\NovelInline{I00879}{\alpha \mapsto \varphi(\alpha,0)}\allowbreak{} 的第一个不动点，\allowbreak{}\NovelInline{I00880}{\varphi(1,0,1)}\allowbreak{} 是第二个，\allowbreak{}\NovelInline{I00881}{\varphi(1,0,2)}\allowbreak{} 是第三个\ldots\ldots 以此类推，\allowbreak{}\NovelInline{I00882}{\varphi(1,0,\gamma)}\allowbreak{} 就是这个映射的第 \allowbreak{}\NovelInline{I00883}{1+\gamma}\allowbreak{} 个不动点，它的迭代规则和二元 \allowbreak{}\NovelInline{I00884}{\varphi(\omega,\gamma)}\allowbreak{} 一模一样，只是把'左变量为 \allowbreak{}\NovelInline{I00885}{\omega}\allowbreak{} 的映射'换成了'左变量为任意 \allowbreak{}\NovelInline{I00886}{\alpha}\allowbreak{} 的映射'。''

卢本韦补充提问：``那如果继续在第三个变量上迭代，比如 \allowbreak{}\NovelInline{I00887}{\varphi(1,1,\gamma)}\allowbreak{}，它的定义是什么？是不是和二元 \allowbreak{}\NovelInline{I00888}{\varphi(1,\gamma)}\allowbreak{} 的逻辑一样，是某个新映射的不动点？''

``问得很关键，这是三元函数的核心拓展，我再给你们举两个关键实例。''弗兰德指尖滑动，勾勒出 \allowbreak{}\NovelInline{I00889}{\varphi(1,1,\gamma)}\allowbreak{} 的定义：

它是新映射 \allowbreak{}\NovelInline{I00890}{\alpha \mapsto \varphi(1,0,\alpha)}\allowbreak{} 的第一个不动点：

\begin{itemize}
\tightlist
\item
  \textbf{定义}：\allowbreak{}\NovelInline{I00891}{\varphi(1,1,0) = \text{1st } \delta \mapsto \varphi(1,0,\delta) \text{ fp.}}\allowbreak{}
\item
  \textbf{基本列}：\allowbreak{}\NovelInline{I00892}{\varphi(1,1,0)[0] = 0}\allowbreak{}，\allowbreak{}\NovelInline{I00893}{\varphi(1,1,0)[n+1] = \varphi(1,0, \varphi(1,1,0)[n])}\allowbreak{}
\item
  \textbf{层级}：它远超 \allowbreak{}\NovelInline{I00894}{\Gamma_1}\allowbreak{}（\allowbreak{}\NovelInline{I00895}{\varphi(1,0,1)}\allowbreak{}），是对 \allowbreak{}\NovelInline{I00896}{\Gamma_\gamma}\allowbreak{} 家族再取不动点的产物。
\end{itemize}

它是映射 \allowbreak{}\NovelInline{I00897}{\alpha \mapsto \varphi(1,0,\alpha)}\allowbreak{} 的第二个不动点：

\begin{itemize}
\tightlist
\item
  \textbf{基本列}：\allowbreak{}\NovelInline{I00898}{\varphi(1,1,1)[n] = \varphi(n, \varphi(1,1,0)+1)}\allowbreak{}
\item
  同理，\allowbreak{}\NovelInline{I00899}{\varphi(1,1,\gamma)}\allowbreak{} 是这个映射的第 \allowbreak{}\NovelInline{I00900}{1+\gamma}\allowbreak{} 个不动点；若 \allowbreak{}\NovelInline{I00901}{\gamma}\allowbreak{} 为极限序数，\allowbreak{}\NovelInline{I00902}{\varphi(1,1,\gamma)[n] = \varphi(1,1,\gamma[n])}\allowbreak{}。
\end{itemize}

它是所有 \allowbreak{}\NovelInline{I00903}{\varphi(1,n,0)}\allowbreak{} 的极限，左变量 \allowbreak{}\NovelInline{I00904}{\beta = \omega}\allowbreak{} 是极限序数：

\begin{itemize}
\tightlist
\item
  \textbf{基本列}：\allowbreak{}\NovelInline{I00905}{\varphi(1,\omega,0)[n] = \varphi(1,n,0)}\allowbreak{}
\item
  \textbf{直观理解}：\allowbreak{}\NovelInline{I00906}{\varphi(1,0,0)}\allowbreak{}、\allowbreak{}\NovelInline{I00907}{\varphi(1,1,0)}\allowbreak{}、\allowbreak{}\NovelInline{I00908}{\varphi(1,2,0)}\allowbreak{}、\allowbreak{}\NovelInline{I00909}{\dots}\allowbreak{} 的上确界，是三元函数中左变量为极限序数的典型例子。
\end{itemize}

他顿了顿，总结三元 \allowbreak{}\NovelInline{I00910}{\varphi}\allowbreak{} 函数的通用规则：

``对于三元 \allowbreak{}\NovelInline{I00911}{\varphi}\allowbreak{} 函数 \allowbreak{}\NovelInline{I00912}{\varphi(1,\beta,\gamma)}\allowbreak{}：

\begin{itemize}
\tightlist
\item
  若 \allowbreak{}\NovelInline{I00913}{\beta}\allowbreak{} 是后继序数 \allowbreak{}\NovelInline{I00914}{\beta'+1}\allowbreak{}，\allowbreak{}\NovelInline{I00915}{\varphi(1,\beta'+1,\gamma)}\allowbreak{} 是映射 \allowbreak{}\NovelInline{I00916}{\alpha \mapsto \varphi(1,\beta',\alpha)}\allowbreak{} 的第 \allowbreak{}\NovelInline{I00917}{1+\gamma}\allowbreak{} 个不动点；
\item
  若 \allowbreak{}\NovelInline{I00918}{\beta}\allowbreak{} 是极限序数，\allowbreak{}\NovelInline{I00919}{\varphi(1,\beta,\gamma)}\allowbreak{} 是所有映射 \allowbreak{}\NovelInline{I00920}{\alpha \mapsto \varphi(1,\beta[n],\alpha)}\allowbreak{} 的第 \allowbreak{}\NovelInline{I00921}{1+\gamma}\allowbreak{} 个共同不动点；
\item
  若 \allowbreak{}\NovelInline{I00922}{\gamma}\allowbreak{} 是极限序数，\allowbreak{}\NovelInline{I00923}{\varphi(1,\beta,\gamma)[n] = \varphi(1,\beta,\gamma[n])}\allowbreak{}。
\end{itemize}

这和二元 \allowbreak{}\NovelInline{I00924}{\varphi(\alpha,\beta)}\allowbreak{} 的规则完全一致，只是把'二元 \allowbreak{}\NovelInline{I00925}{\varphi(\alpha,\beta)}\allowbreak{}'换成了'三元 \allowbreak{}\NovelInline{I00926}{\varphi(1,\beta,\gamma)}\allowbreak{}'，本质是不动点的嵌套。''

弗兰德指尖再挥，虚空里的三元 \allowbreak{}\NovelInline{I00927}{\varphi}\allowbreak{} 函数延伸出四元、五元乃至更多变量的结构。他指着最通用的多元 \allowbreak{}\NovelInline{I00928}{\varphi}\allowbreak{} 函数形式 \allowbreak{}\NovelInline{I00929}{\varphi(\#, \alpha, Z, \beta)}\allowbreak{} 讲解：

``多元 Veblen 函数的本质是序列式结构：我们用 \allowbreak{}\NovelInline{I00930}{\#}\allowbreak{} 表示任意有限个前置变量组成的序列，用 \allowbreak{}\NovelInline{I00931}{Z}\allowbreak{} 表示任意有限个 \allowbreak{}\NovelInline{I00932}{0}\allowbreak{} 组成的序列，最终的多元函数可以统一记为 \allowbreak{}\NovelInline{I00933}{\varphi(\#, \alpha, Z, \beta)}\allowbreak{}------不管有多少个变量，都能套进这个框架里。我用实例让你们看懂 \allowbreak{}\NovelInline{I00934}{\#}\allowbreak{} 和 \allowbreak{}\NovelInline{I00935}{Z}\allowbreak{} 的作用。''

他先给出最核心的不动点定义，严格贴合资料：

``对于 \allowbreak{}\NovelInline{I00936}{\varphi(\#, \alpha+1, Z, \beta)}\allowbreak{}：它是映射 \allowbreak{}\NovelInline{I00937}{\gamma \mapsto \varphi(\#, \alpha, \gamma, Z)}\allowbreak{} 的第 \allowbreak{}\NovelInline{I00938}{1+\beta}\allowbreak{} 个不动点；若 \allowbreak{}\NovelInline{I00939}{\alpha}\allowbreak{} 是极限序数，\allowbreak{}\NovelInline{I00940}{\varphi(\#, \alpha, Z, \beta)}\allowbreak{} 是所有映射 \allowbreak{}\NovelInline{I00941}{\gamma \mapsto \varphi(\#, \alpha[n], \gamma, Z)}\allowbreak{} 的第 \allowbreak{}\NovelInline{I00942}{1+\beta}\allowbreak{} 个共同不动点；若 \allowbreak{}\NovelInline{I00943}{\beta}\allowbreak{} 是极限序数，\allowbreak{}\NovelInline{I00944}{\varphi(\#, \alpha, Z, \beta)}\allowbreak{} 是所有 \allowbreak{}\NovelInline{I00945}{\varphi(\#, \alpha, Z, \beta[n])}\allowbreak{} 的极限，即 \allowbreak{}\NovelInline{I00946}{\varphi(\#, \alpha, Z, \beta)[n] = \varphi(\#, \alpha, Z, \beta[n])}\allowbreak{}。''

牢大挠了挠头，满脸困惑：``院长，这个 \allowbreak{}\NovelInline{I00947}{\#}\allowbreak{} 和 \allowbreak{}\NovelInline{I00948}{Z}\allowbreak{} 到底是什么意思？比如四元 \allowbreak{}\NovelInline{I00949}{\varphi(1,0,0,\gamma)}\allowbreak{}，它对应的 \allowbreak{}\NovelInline{I00950}{\#}\allowbreak{}、\allowbreak{}\NovelInline{I00951}{\alpha}\allowbreak{}、\allowbreak{}\NovelInline{I00952}{Z}\allowbreak{}、\allowbreak{}\NovelInline{I00953}{\beta}\allowbreak{} 分别是什么？''

``这个问题问到了多元函数的'结构本质'，用实例拆解给你们看。''弗兰德笑着举例：

\begin{itemize}
\item
  \textbf{结构拆分}：\allowbreak{}\NovelInline{I00954}{\# = (1)}\allowbreak{}，\allowbreak{}\NovelInline{I00955}{\alpha = 0}\allowbreak{}，\allowbreak{}\NovelInline{I00956}{Z = (0)}\allowbreak{}，\allowbreak{}\NovelInline{I00957}{\beta = \gamma}\allowbreak{}
\item
  \textbf{定义}：它是映射 \allowbreak{}\NovelInline{I00958}{\delta \mapsto \varphi(1,\delta,0,0)}\allowbreak{} 的第 \allowbreak{}\NovelInline{I00959}{1+\gamma}\allowbreak{} 个不动点
\item
  \textbf{特殊实例}：\allowbreak{}\NovelInline{I00960}{\varphi(1,0,0,0)}\allowbreak{}（四元第一个不动点）

  \begin{itemize}
  \tightlist
  \item
    基本列：\allowbreak{}\NovelInline{I00961}{\varphi(1,0,0,0)[0] = 0}\allowbreak{}，\allowbreak{}\NovelInline{I00962}{\varphi(1,0,0,0)[n+1] = \varphi(1, \varphi(1,0,0,0)[n], 0, 0)}\allowbreak{}
  \item
    层级：远超三元 \allowbreak{}\NovelInline{I00963}{\varphi(1,0,0)}\allowbreak{}，是对三元不动点再取不动点的产物。
  \end{itemize}
\item
  \textbf{结构拆分}：\allowbreak{}\NovelInline{I00964}{\# = \varnothing}\allowbreak{}，\allowbreak{}\NovelInline{I00965}{\alpha = 2}\allowbreak{}，\allowbreak{}\NovelInline{I00966}{Z = (0,0)}\allowbreak{}，\allowbreak{}\NovelInline{I00967}{\beta = 0}\allowbreak{}
\item
  \textbf{定义}：它是映射 \allowbreak{}\NovelInline{I00968}{\alpha \mapsto \varphi(1,\alpha,0,0)}\allowbreak{} 的第一个不动点
\item
  \textbf{基本列}：\allowbreak{}\NovelInline{I00969}{\varphi(2,0,0,0)[0] = 0}\allowbreak{}，\allowbreak{}\NovelInline{I00970}{\varphi(2,0,0,0)[n+1] = \varphi(1, \varphi(2,0,0,0)[n], 0, 0)}\allowbreak{}
\item
  \textbf{直观理解}：和二元 \allowbreak{}\NovelInline{I00971}{\varphi(2,0)}\allowbreak{}、三元 \allowbreak{}\NovelInline{I00972}{\varphi(2,0,0)}\allowbreak{} 逻辑同源，是对前一层映射取不动点。
\item
  \textbf{结构拆分}：\allowbreak{}\NovelInline{I00973}{\# = (1)}\allowbreak{}，\allowbreak{}\NovelInline{I00974}{\alpha = 0}\allowbreak{}，\allowbreak{}\NovelInline{I00975}{Z = (0,0,0)}\allowbreak{}，\allowbreak{}\NovelInline{I00976}{\beta = 0}\allowbreak{}
\item
  \textbf{定义}：它是映射 \allowbreak{}\NovelInline{I00977}{\delta \mapsto \varphi(1,\delta,0,0,0)}\allowbreak{} 的第一个不动点
\item
  \textbf{基本列}：\allowbreak{}\NovelInline{I00978}{\varphi(1,0,0,0,0)[0] = 0}\allowbreak{}，\allowbreak{}\NovelInline{I00979}{\varphi(1,0,0,0,0)[n+1] = \varphi(1, \varphi(1,0,0,0,0)[n], 0, 0, 0)}\allowbreak{}
\item
  \textbf{规律}：每多一个变量，就是对前一层不动点再取一次不动点，层级呈阶跃式提升。
\end{itemize}

诺当紧随其后提问：``那如果 \allowbreak{}\NovelInline{I00980}{\#}\allowbreak{} 是空序列，比如 \allowbreak{}\NovelInline{I00981}{\varphi(\alpha, Z, \beta)}\allowbreak{}，是不是就退化成三元函数？如果 \allowbreak{}\NovelInline{I00982}{Z}\allowbreak{} 也是空序列，是不是就回到二元 \allowbreak{}\NovelInline{I00983}{\varphi(\alpha,\beta)}\allowbreak{}？''

``完全正确！''弗兰德肯定答复，``这就是多元函数的'退化规则'，我用实例总结：

\begin{itemize}
\tightlist
\item
  当 \allowbreak{}\NovelInline{I00984}{\# = \varnothing}\allowbreak{}、\allowbreak{}\NovelInline{I00985}{Z = \varnothing}\allowbreak{} 时，\allowbreak{}\NovelInline{I00986}{\varphi(\alpha,\beta)}\allowbreak{} 就是我们最熟悉的二元 Veblen 函数（如 \allowbreak{}\NovelInline{I00987}{\varphi(1,0) = \varepsilon_0}\allowbreak{}、\allowbreak{}\NovelInline{I00988}{\varphi(2,0) = \zeta_0}\allowbreak{}）；
\item
  当 \allowbreak{}\NovelInline{I00989}{\# = (1)}\allowbreak{}、\allowbreak{}\NovelInline{I00990}{Z = \varnothing}\allowbreak{} 时，\allowbreak{}\NovelInline{I00991}{\varphi(1,\alpha,\beta)}\allowbreak{} 就是三元 Veblen 函数（如 \allowbreak{}\NovelInline{I00992}{\varphi(1,0,0) = \Gamma_0}\allowbreak{}、\allowbreak{}\NovelInline{I00993}{\varphi(1,1,0)}\allowbreak{}）；
\item
  当 \allowbreak{}\NovelInline{I00994}{\# = (1)}\allowbreak{}、\allowbreak{}\NovelInline{I00995}{Z = (0)}\allowbreak{} 时，\allowbreak{}\NovelInline{I00996}{\varphi(1,\alpha,0,\beta)}\allowbreak{} 就是四元 Veblen 函数（如 \allowbreak{}\NovelInline{I00997}{\varphi(1,0,0,0)}\allowbreak{}）；
\item
  每多一个前置变量到 \allowbreak{}\NovelInline{I00998}{\#}\allowbreak{} 里，就是对前一层映射再取不动点，变量数可以无限增加。
\end{itemize}

他补充道：``资料里还提到一个特殊情况：当 \allowbreak{}\NovelInline{I00999}{\# = \varnothing}\allowbreak{}、\allowbreak{}\NovelInline{I01000}{\alpha = 1}\allowbreak{}、\allowbreak{}\NovelInline{I01001}{Z = \varnothing}\allowbreak{} 时，\allowbreak{}\NovelInline{I01002}{\varphi(1,\beta)}\allowbreak{} 就是 \allowbreak{}\NovelInline{I01003}{\varepsilon_\beta}\allowbreak{}，这是一元 Veblen 函数的特例------\allowbreak{}\NovelInline{I01004}{\varphi(\beta) = \omega^\beta}\allowbreak{}，而 \allowbreak{}\NovelInline{I01005}{\varphi(1,\beta) = \varepsilon_\beta}\allowbreak{}，正好对应 \allowbreak{}\NovelInline{I01006}{\varepsilon}\allowbreak{} 序数的定义，比如 \allowbreak{}\NovelInline{I01007}{\varphi(1,0) = \varepsilon_0}\allowbreak{}、\allowbreak{}\NovelInline{I01008}{\varphi(1,1) = \varepsilon_1}\allowbreak{}。''

弗兰德看着众人对复杂结构的迷茫，将多元函数的基本列整理成和二元一致的口诀，结合实例方便大家记忆：

``不管有多少个变量，多元 \allowbreak{}\NovelInline{I01009}{\varphi}\allowbreak{} 函数的基本列都遵循三条核心规则，和二元函数完全同源，我给你们拆解成口诀 + 实例：''

若最左边的非零变量是后继序数 \allowbreak{}\NovelInline{I01010}{\alpha+1}\allowbreak{}，那么：

\begin{itemize}
\tightlist
\item
  \allowbreak{}\NovelInline{I01011}{\varphi(\#, \alpha+1, Z, \beta)[0] = \varphi(\#, \alpha+1, Z, \beta) + 1}\allowbreak{}（跳出不动点）
\item
  \allowbreak{}\NovelInline{I01012}{\varphi(\#, \alpha+1, Z, \beta)[n+1] = \varphi(\#, \alpha, \varphi(\#, \alpha+1, Z, \beta)[n], Z)}\allowbreak{}
\end{itemize}

\textbf{实例}：\allowbreak{}\NovelInline{I01013}{\varphi(\omega+1,0,0)}\allowbreak{}（左变量 \allowbreak{}\NovelInline{I01014}{\omega+1}\allowbreak{} 是后继序数）

\begin{itemize}
\tightlist
\item
  \allowbreak{}\NovelInline{I01015}{\varphi(\omega+1,0,0)[0] = \varphi(\omega+1,0,0) + 1}\allowbreak{}
\item
  \allowbreak{}\NovelInline{I01016}{\varphi(\omega+1,0,0)[n+1] = \varphi(\omega, \varphi(\omega+1,0,0)[n], 0, 0)}\allowbreak{}
\end{itemize}

这和二元 \allowbreak{}\NovelInline{I01017}{\varphi(\omega+1,0)[n+1] = \varphi(\omega, \varphi(\omega+1,0)[n])}\allowbreak{} 的逻辑完全一样，只是多了 \allowbreak{}\NovelInline{I01018}{\#}\allowbreak{} 和 \allowbreak{}\NovelInline{I01019}{Z}\allowbreak{} 的占位。

若最左边的非零变量是极限序数 \allowbreak{}\NovelInline{I01020}{\alpha}\allowbreak{}，那么：

\begin{itemize}
\tightlist
\item
  \allowbreak{}\NovelInline{I01021}{\varphi(\#, \alpha, Z, \beta)[n] = \varphi(\#, \alpha[n], Z, \beta)}\allowbreak{}
\end{itemize}

\textbf{实例1}：\allowbreak{}\NovelInline{I01022}{\varphi(\omega,0,0)}\allowbreak{}（左变量 \allowbreak{}\NovelInline{I01023}{\omega}\allowbreak{} 是极限序数）

\begin{itemize}
\tightlist
\item
  \allowbreak{}\NovelInline{I01024}{\varphi(\omega,0,0)[n] = \varphi(n,0,0)}\allowbreak{}
\end{itemize}

\textbf{实例2}：\allowbreak{}\NovelInline{I01025}{\varphi(1,\omega,0)}\allowbreak{}（左变量第二个位置的 \allowbreak{}\NovelInline{I01026}{\omega}\allowbreak{} 是极限序数）

\begin{itemize}
\tightlist
\item
  \allowbreak{}\NovelInline{I01027}{\varphi(1,\omega,0)[n] = \varphi(1,n,0)}\allowbreak{}
\end{itemize}

这和二元 \allowbreak{}\NovelInline{I01028}{\varphi(\omega,0)[n] = \varphi(n,0)}\allowbreak{} 的规则一致。

若最右边的变量 \allowbreak{}\NovelInline{I01029}{\beta}\allowbreak{} 是极限序数，那么：

\begin{itemize}
\tightlist
\item
  \allowbreak{}\NovelInline{I01030}{\varphi(\#, \alpha, Z, \beta)[n] = \varphi(\#, \alpha, Z, \beta[n])}\allowbreak{}
\end{itemize}

\textbf{实例}：\allowbreak{}\NovelInline{I01031}{\varphi(1,0,\omega)}\allowbreak{}（右变量 \allowbreak{}\NovelInline{I01032}{\omega}\allowbreak{} 是极限序数）

\begin{itemize}
\tightlist
\item
  \allowbreak{}\NovelInline{I01033}{\varphi(1,0,\omega)[n] = \varphi(1,0,n)}\allowbreak{}
\end{itemize}

这和二元 \allowbreak{}\NovelInline{I01034}{\varphi(\omega,\gamma)[n] = \varphi(\omega,\gamma[n])}\allowbreak{} 的规则一致。

\begin{center}\rule{0.5\linewidth}{0.5pt}\end{center}

蔡须坤握着篮球，指尖转动，结合实战感悟发问：``院长，那 \allowbreak{}\NovelInline{I01035}{\varphi(1,0,1)}\allowbreak{} 也就是 \allowbreak{}\NovelInline{I01036}{\Gamma_1}\allowbreak{} 的基本列该怎么写？是不是和 \allowbreak{}\NovelInline{I01037}{\Gamma_0}\allowbreak{} 一样，从零开始反复代入？''

``没错，我们来完整写一遍 \allowbreak{}\NovelInline{I01038}{\varphi(1,0,1)}\allowbreak{} 的基本列：''弗兰德指尖勾勒：

\NovelDisplay{M00151R01}{%
\varphi(1,0,1)[0] = 1%
}
\NovelDisplay{M00151R02}{%
\varphi(1,0,1)[1] = \varphi(1,0,0) + 1%
}
\NovelDisplay{M00151R03}{%
\varphi(1,0,1)[2] = \varphi(\varphi(1,0,1)[1], 0) = \varphi(\varphi(1,0,0)+1, 0)%
}
\NovelDisplay{M00151R04}{%
\varphi(1,0,1)[n] = \varphi(\varphi(1,0,1)[n-1], 0)%
}

奥斯卡举手补充：``那如果是四元 \allowbreak{}\NovelInline{I01039}{\varphi(1,0,0,0)}\allowbreak{}，它的基本列是不是也是这个逻辑？只是嵌套层数更多？''

``完全正确。''弗兰德点头，``\allowbreak{}\NovelInline{I01040}{\varphi(1,0,0,0)}\allowbreak{} 是映射 \allowbreak{}\NovelInline{I01041}{\delta \mapsto \varphi(\delta,0,0)}\allowbreak{} 的第一个不动点，基本列是：

\NovelDisplay{M00152R01}{%
\varphi(1,0,0,0)[0] = 1%
}
\NovelDisplay{M00152R02}{%
\varphi(1,0,0,0)[n+1] = \varphi(\varphi(1,0,0,0)[n], 0, 0)%
}

它的极限就是四元函数的第一个不动点，层级远超三元 \allowbreak{}\NovelInline{I01042}{\varphi(1,0,1)}\allowbreak{}，更别说 \allowbreak{}\NovelInline{I01043}{\Gamma_0}\allowbreak{} 了。''

弗兰德指尖在虚空最顶端画出一个新的符号------SVO，周身泛起比 \allowbreak{}\NovelInline{I01044}{\Gamma_0}\allowbreak{} 更耀眼的光晕，他沉声讲解：

``当我们不断增加变量数，从二元 \allowbreak{}\NovelInline{I01045}{\varphi(\alpha,\beta)}\allowbreak{} 到三元 \allowbreak{}\NovelInline{I01046}{\varphi(1,0,0)}\allowbreak{}，再到四元 \allowbreak{}\NovelInline{I01047}{\varphi(1,0,0,0)}\allowbreak{}、五元 \allowbreak{}\NovelInline{I01048}{\varphi(1,0,0,0,0)}\allowbreak{}\ldots\ldots 无限迭代下去，这些序数的上确界，就是小 Veblen 序数（Small Veblen Ordinal，SVO），它的定义是：

\NovelDisplay{M00153R01}{%
\text{SVO} = \sup\{\varphi(1), \varphi(1,0), \varphi(1,0,0), \varphi(1,0,0,0), \dots\}%
}

每多一个 \allowbreak{}\NovelInline{I01049}{0}\allowbreak{}，就是多一层不动点嵌套，SVO 就是所有这些有限元不动点的极限。''

卢本韦微微挑眉，清冷开口：``那 SVO 是不是就是所有有限变量 Veblen 函数能构造出的最大序数？超过它就需要超限序数变量了？''

``正是如此。''弗兰德肯定答复，``SVO 是有限变量 Veblen 函数的极限，任何有限个变量的 \allowbreak{}\NovelInline{I01050}{\varphi}\allowbreak{} 函数，不管嵌套多少层，都无法构造出超过 SVO 的序数。想要突破 SVO，就必须引入'无限变量'的 Veblen 函数，那已经超出了你们现阶段的修炼范畴，我们先把有限变量的逻辑吃透。''

牢大一拍大腿，恍然大悟：``原来如此！从 \allowbreak{}\NovelInline{I01051}{\Gamma_0}\allowbreak{} 到 SVO，就是不断给 Veblen 函数加变量、嵌套不动点的过程，每多一个变量，层级就跃升一个大台阶！比如 \allowbreak{}\NovelInline{I01052}{\Gamma_0}\allowbreak{} 是二元到三元的跨越，\allowbreak{}\NovelInline{I01053}{\varphi(1,0,0,0)}\allowbreak{} 是三元到四元的跨越，SVO 就是所有有限变量的终点！''

``没错。''弗兰德笑着总结，``多元 Veblen 函数的核心奥义，就是不动点的无限嵌套：

\begin{itemize}
\tightlist
\item
  \textbf{二元 \allowbreak{}\NovelInline{I01054}{\varphi}\allowbreak{} 函数}：对 \allowbreak{}\NovelInline{I01055}{\alpha \mapsto \varphi(\alpha,0)}\allowbreak{} 取不动点，得到 \allowbreak{}\NovelInline{I01056}{\Gamma_0}\allowbreak{}；
\item
  \textbf{三元 \allowbreak{}\NovelInline{I01057}{\varphi}\allowbreak{} 函数}：对 \allowbreak{}\NovelInline{I01058}{\alpha \mapsto \varphi(\alpha,0,0)}\allowbreak{} 取不动点，得到 \allowbreak{}\NovelInline{I01059}{\varphi(1,0,0,0)}\allowbreak{}；
\item
  \textbf{四元 \allowbreak{}\NovelInline{I01060}{\varphi}\allowbreak{} 函数}：对 \allowbreak{}\NovelInline{I01061}{\alpha \mapsto \varphi(\alpha,0,0,0)}\allowbreak{} 取不动点，得到 \allowbreak{}\NovelInline{I01062}{\varphi(1,0,0,0,0)}\allowbreak{}；
\item
  \ldots\ldots{}
\item
  \textbf{无限迭代后}，得到 SVO。
\end{itemize}

这就是从 \allowbreak{}\NovelInline{I01063}{\Gamma_0}\allowbreak{} 到 SVO 的完整进阶路径，也是你们接下来突破 \allowbreak{}\NovelInline{I01064}{\Gamma_0}\allowbreak{} 后，要攀登的下一座高峰。''

弗兰德看着众人若有所思的样子，主动抛出几个常见误区，结合实例帮大家厘清逻辑：

``绝对不是！我用实例对比给你们看：

\begin{itemize}
\tightlist
\item
  \allowbreak{}\NovelInline{I01065}{\varphi(\omega,0,0)}\allowbreak{}：三元函数，左变量为 \allowbreak{}\NovelInline{I01066}{\omega}\allowbreak{}（极限序数），基本列是 \allowbreak{}\NovelInline{I01067}{\varphi(\omega,0,0)[n] = \varphi(n,0,0)}\allowbreak{}，只是三元函数里的一个普通项；
\item
  \allowbreak{}\NovelInline{I01068}{\varphi(1,0,0)}\allowbreak{}：三元函数的第一个不动点，是所有 \allowbreak{}\NovelInline{I01069}{\varphi(n,0,0)}\allowbreak{} 的极限，层级高于 \allowbreak{}\NovelInline{I01070}{\varphi(\omega,0,0)}\allowbreak{}。就像 \allowbreak{}\NovelInline{I01071}{\varphi(\omega,0)}\allowbreak{} 是二元函数的极限，而 \allowbreak{}\NovelInline{I01072}{\varphi(1,0,0) = \Gamma_0}\allowbreak{} 是二元函数的不动点，二者天差地别。''
\end{itemize}

``不行！我用实例对比 \allowbreak{}\NovelInline{I01073}{\varphi(1,0,1)}\allowbreak{} 和 \allowbreak{}\NovelInline{I01074}{\varphi(1,1,0)}\allowbreak{}：

\begin{itemize}
\tightlist
\item
  \allowbreak{}\NovelInline{I01075}{\varphi(1,0,1)}\allowbreak{}：映射 \allowbreak{}\NovelInline{I01076}{\alpha \mapsto \varphi(\alpha,0)}\allowbreak{} 的第二个不动点，属于 \allowbreak{}\NovelInline{I01077}{\Gamma_\gamma}\allowbreak{} 家族，层级是 \allowbreak{}\NovelInline{I01078}{\Gamma_1}\allowbreak{}；
\item
  \allowbreak{}\NovelInline{I01079}{\varphi(1,1,0)}\allowbreak{}：映射 \allowbreak{}\NovelInline{I01080}{\alpha \mapsto \varphi(1,0,\alpha)}\allowbreak{} 的第一个不动点，是对 \allowbreak{}\NovelInline{I01081}{\Gamma_\gamma}\allowbreak{} 家族再取不动点的产物，层级远超 \allowbreak{}\NovelInline{I01082}{\Gamma_1}\allowbreak{}。多元 Veblen 函数的变量是有序的，最右边的变量是'某映射的第 \allowbreak{}\NovelInline{I01083}{1+\beta}\allowbreak{} 个不动点'，左边的变量是'对前一层映射取不动点'，顺序不能乱。''
\end{itemize}

``当然不是！序数没有最大，只有更大。SVO 只是有限变量 Veblen 函数的极限，引入无限变量后，我们还能构造出远超 SVO 的序数，比如大 Veblen 序数（LVO），甚至是巴赫曼-霍华德序数（BHO）。但对你们现阶段来说，SVO 已经是足够遥远的目标，先把多元函数的基础打牢。''

弗兰德最后看着众人，语气郑重：

``今天我们讲透了多元 Veblen 函数的定义、基本列与核心逻辑，结合实例拆解了 \allowbreak{}\NovelInline{I01084}{\varphi(1,0,1)}\allowbreak{}、\allowbreak{}\NovelInline{I01085}{\varphi(1,1,0)}\allowbreak{}、\allowbreak{}\NovelInline{I01086}{\varphi(1,0,0,0)}\allowbreak{} 等关键节点，它和你们之前突破 \allowbreak{}\NovelInline{I01087}{\Gamma_0}\allowbreak{} 的路径是完全衔接的：

\begin{enumerate}
\def\labelenumi{\arabic{enumi}.}
\tightlist
\item
  \allowbreak{}\NovelInline{I01088}{\Gamma_0}\allowbreak{} 是二元 \allowbreak{}\NovelInline{I01089}{\varphi}\allowbreak{} 函数的自反不动点（\allowbreak{}\NovelInline{I01090}{\Gamma_0 = \varphi(\Gamma_0,0)}\allowbreak{}）；
\item
  三元 \allowbreak{}\NovelInline{I01091}{\varphi(1,0,0)}\allowbreak{} 是对二元不动点再取不动点，是 \allowbreak{}\NovelInline{I01092}{\Gamma_0}\allowbreak{} 的直接进阶；
\item
  每多一个变量，就是对前一层不动点再取不动点，最终通向 SVO。
\end{enumerate}

你们吸收了馃槔的力量，又经过人面魔蛛的实战洗礼，现在要做的，就是把多元函数的不动点逻辑，融入到 \allowbreak{}\NovelInline{I01093}{\Gamma_0}\allowbreak{} 的突破感悟里------当你们能自如操控三元 \allowbreak{}\NovelInline{I01094}{\varphi(1,0,\gamma)}\allowbreak{} 的基本列、分清 \allowbreak{}\NovelInline{I01095}{\varphi(1,0,1)}\allowbreak{} 和 \allowbreak{}\NovelInline{I01096}{\varphi(1,1,0)}\allowbreak{} 的层级差异时，就是彻底稳固 \allowbreak{}\NovelInline{I01097}{\Gamma_0}\allowbreak{} 层级、迈向 SVO 的开始。''

话音落下，虚空里的多元 \allowbreak{}\NovelInline{I01098}{\varphi}\allowbreak{} 函数脉络与众人身上的序数纹路产生强烈共鸣，唐三、牢大、卢本韦等人周身泛起淡淡的金光，显然是触及了更高阶序数的突破契机。

\chapter{第十六章：特训，切磋和打磨}\label{ux7b2cux5341ux516dux7ae0ux7279ux8badux5207ux78cbux548cux6253ux78e8}

讲堂上的多元 \allowbreak{}\NovelInline{I01099}{\varphi}\allowbreak{} 函数脉络早已消散，可那层层嵌套的不动点逻辑、有序变量的严苛规则，依旧像一层密不透风的壁垒，死死困住众人的修为进阶。数月光阴弹指而过，学院后山的静修室里，序数光晕忽明忽暗，却始终无法冲破瓶颈------唐三卡在 \allowbreak{}\NovelInline{I01100}{\varphi(1,1,0)}\allowbreak{} 的门槛前，任凭如何推演基本列，始终摸不透映射 \allowbreak{}\NovelInline{I01101}{\alpha \mapsto \varphi(1,0,\alpha)}\allowbreak{} 的不动点真谛；牢大空有一身蛮力，对着 \allowbreak{}\NovelInline{I01102}{\varphi(2,0,0)}\allowbreak{} 的递归链条反复冲撞，却总在迭代中途魂力紊乱，功亏一篑；蔡须坤以篮球校准序数波动，却始终跨不过 \allowbreak{}\NovelInline{I01103}{\varphi(1,\omega,0)}\allowbreak{} 的极限收敛关；卢本韦、诺当、小舞三人，更是在 \allowbreak{}\NovelInline{I01104}{\Gamma_0}\allowbreak{} 进阶后的 \allowbreak{}\NovelInline{I01105}{\varphi(1,0,0)}\allowbreak{} 层级停滞不前，连最基础的三元不动点闭环无法稳固；就连心思细腻的奥斯卡，也被 \allowbreak{}\NovelInline{I01106}{\varphi(1,0,\varphi(1,0))}\allowbreak{} 的双层嵌套困住，FGH 展开到一半便会脉络断裂，前功尽弃。

\begin{center}\rule{0.5\linewidth}{0.5pt}\end{center}

弗兰德伫立在瞭望台，望着静修室里此起彼伏的黯淡光晕，眉头紧锁。身旁的玉小刚手捧序数典籍，指尖划过泛黄的书页，声音沉稳：``多元函数是序数修炼的第一道天堑，纯靠自行参悟太慢，更何况他们刚突破 \allowbreak{}\NovelInline{I01107}{\Gamma_0}\allowbreak{}，根基尚未完全稳固，强行啃硬骨头只会消磨心性。必须逼他们一把，用实战淬炼悟性，把理论揉碎了灌进骨子里。''

赵无极攥紧拳头，骨节噼啪作响，浑厚的嗓音震得空气微颤：``老夫早就手痒了！普通修炼没用，得来点狠的，让他们在生死边缘吃透不动点的规则！''

三巨头对视一眼，当即敲定主意------这场特训，没有退路，只有通关。

\begin{center}\rule{0.5\linewidth}{0.5pt}\end{center}

翌日清晨，天刚蒙蒙亮，刺耳的集结号便响彻学院。唐三等人来不及休整，周身还带着未散的倦意，便被强行带到后山开辟的序数炼狱场。此地被三位老师布下高阶禁制，空气中浮动着锋利的序数碎片，每一次呼吸都牵扯魂力脉络，稍不留神就会被碎片割伤魂脉、修为倒退。

弗兰德立于禁制中央，周身金色光晕凝成实质，语气没有半分情面：``从今天起，没有课堂，没有答疑，只有闯关。通关标准只有一个：让自身魂力与对应序数的基本列完全共鸣，形成不可逆的不动点闭环。做不到，就别想踏出炼狱场一步！''

玉小刚上前一步，目光扫过七人，声音沉稳却不容置疑。他没有急于下达任务，而是先逐一剖析了每个人的战斗特质------这不仅是分工，更是他们能否在实战中存活下来的根基。

``牢大，你的攻击力是七人中最强的，这一点毋庸置疑。\allowbreak{}\NovelInline{I01108}{\varphi(2,0,0)}\allowbreak{} 级别的递归蛮力，连赵老师都要认真应对。但你的问题同样明显------你太依赖纯粹的输出，缺少增幅手段，单凭自身蛮力冲击高阶不动点，效率极低。奥斯卡，你的辅助增幅能力是牢大最需要的搭档，你的攻击力虽然不强，但你的增幅能让牢大的输出翻数倍，你们是天然的核心组合。''

奥斯卡点了点头，目光中闪过一丝复杂------他清楚自己的攻击力在七人中垫底，但增幅辅助的能力却是不可替代的。

玉小刚继续道：``蔡须坤，你的篮球序数校准天赋异禀，擅长单点突破，能在极短时间内找到对手防御的薄弱点并一击贯穿。但你的短板是持续输出不足，一旦单点突破未能奏效，后续攻击力会断崖式下降。小舞，你恰好相反------近身作战能力极强，身法灵活、爆发迅猛，贴身缠斗几乎无人能及，但你的防御力是七人中最弱的，一旦被对手抓住破绽近身反击，你几乎没有还手之力。''

小舞咬了咬嘴唇，没有反驳。她知道自己的短板------上次被不乐的 \allowbreak{}\NovelInline{I01109}{\varepsilon_1}\allowbreak{} 蛮力一击掀飞，至今心有余悸。

``唐三、诺当，你们两人的控制能力是团队的全局支柱。唐三的蓝银草感知网络可以覆盖整个战场，诺当的序数干扰场能封锁对手的行动路径，你们联手可以控制全场节奏。但------''玉小刚的语气陡然加重，``你们的控制能力再强，也只是'工具'。工具需要有人来使用，需要有人告诉你们该控制哪里、封锁什么方向、在什么时机释放。这个人，是牢大。''

他看向牢大，目光郑重：``牢大，你不只是队伍里攻击力最强的人，你还是战术核心。你的战场直觉和全局判断力，是指挥控制系发挥最大效能的关键。唐三和诺当的控制能力，必须由你来指定战术方向，才能真正发挥全局作用。没有你的指挥，他们的控制就是无的放矢。''

牢大微微一怔，他从未想过自己还承担着这样的角色。他习惯了冲锋在前、蛮力碾压，却忽略了自己在战场上的直觉判断力。

``最后是卢本韦。''玉小刚的目光转向那个沉默寡言的少年，``你的群体作战能力极强，大范围序数铺展、多目标同时压制，是你的拿手好戏。但你的单点突破能力是七人中最弱的，面对高防御的单体目标，你的攻击几乎无法造成有效伤害。所以你的定位是清场和牵制，而非攻坚。''

卢本韦默默点了点头，他的表情平静，似乎对自己的定位早已心知肚明。

玉小刚说完，退后一步，将舞台交给了赵无极。

赵无极跨步踏入禁制中央，浑厚土属性魂力席卷全场，直接碾碎虚浮的序数碎片推演，开启最残酷的实战死训模式。他不借禁制外力，亲自充当终极陪练，周身压制性威压扑面而来，每一寸肌肉都透着沙场狠厉：``都给老夫打起精神！序数道理不是空想的，是在挨打、格挡、反击里练出来的！今天老夫既当考官，也当教头，教你们把不动点迭代融进实战！''

他出手快准狠却留一线生机，每一击都精准戳向魂力迭代薄弱点，逼迫众人在生死瞬间运转对应层级的 \allowbreak{}\NovelInline{I01110}{\varphi}\allowbreak{} 函数力量。炼狱场瞬间化作肉搏战场，尘土飞扬、魂技碰撞声震耳欲聋，赵无极一边以强横攻势施压，一边拆解招式里的序数逻辑，把技巧揉进每一次进攻，让众人在抗压中领悟、在配合中破局。

\textbf{特训铁律：禁止使用外力强行拔高修为，只许靠自身感悟推演迭代；不许躲懒、不许独逃、不许弃队友，一人倒下全队搀扶，一人失误全员共担，每日仅供给最低限度灵气，以战养悟、熬出真章。}

\begin{center}\rule{0.5\linewidth}{0.5pt}\end{center}

特训开启便无喘息之机，赵无极足尖点地，土属性魂力凝实如铅，全程游走在炼狱场中，将多元 \allowbreak{}\NovelInline{I01111}{\varphi}\allowbreak{} 函数的递归规则、不动点逻辑尽数化作拳风、掌风与魂技冲击，不给众人丝毫调息间隙。众人在剧痛中咬牙支撑，每一次格挡、反击、配合，都必须同步完成 FGH 展开与序数迭代，稍有差错便会被强横力道震飞，魂力脉络刺痛难忍，但没有一人示弱退避。

赵无极身形骤闪，率先扑向唐三，右掌横劈而出，掌风裹挟着 \allowbreak{}\NovelInline{I01112}{\varphi(1,1,0)}\allowbreak{} 的三阶迭代节律，直逼唐三胸前魂脉，意图强行打断他的序数运转。唐三不敢大意，立刻催动蓝银草武魂，细密的草藤层层交织成防御网，同时眉心紫极魔瞳全开，口中默念序数递归，实时展开 FGH 运算：

\NovelDisplay{M00154R01}{%
f_{\varphi(1,1,0)}(3) = f_{\varphi(1,0,\varphi(1,0,\varphi(1,0,0)))}(3)%
}

可他刚迭代到第二层，掌风便冲撞而来，蓝银草崩断数根，魂力脉络一阵刺痛，FGH 展开直接中断。

``心神跟着迭代走，不是魂力跟着想法跑！''赵无极厉声提点，第二掌放缓半分，掌风节律刻意外露，``盯着 \allowbreak{}\NovelInline{I01113}{\varphi(1,0,\alpha)}\allowbreak{} 的映射路径，把蓝银草缠绕节奏和 FGH 展开同步！''

唐三深吸一口气，摒弃杂念，再次催动蓝银草，这一次不再硬抗，而是让草藤顺着序数映射轨迹缠绕，同时完成完整 FGH 展开：

\NovelDisplay{M00155R01}{%
f_{\varphi(1,1,0)}(3) = f_{\varphi(1,0,\varphi(1,0,\varphi(1,0,0)))}(3)%
}

掌风撞上草藤的瞬间，冲击力被层层卸去，\allowbreak{}\NovelInline{I01114}{\varphi(1,1,0)}\allowbreak{} 的脉络彻底稳住，防御网坚不可摧。

与此同时，诺当在另一侧悄然展开序数干扰场，淡紫色的波动如蛛网般铺展，试图封锁赵无极的移动路径。他的控制精度极高，干扰场的边界恰好卡在赵无极的落脚点前方半寸------如果赵无极继续前冲，就会一头撞进干扰场的核心区域，行动速度将被大幅拖慢。

但赵无极只是轻蔑一笑，右脚踏地的瞬间，土属性魂力如地龙翻身，直接将干扰场从根部震碎。

``诺当，你的控制太'等'了！''赵无极厉声喝道，``你在等我踩进去？我不会给你这个机会！控制不是守株待兔，是逼对手不得不走进你的陷阱！你和唐三的配合需要牢大的战术指引------牢大，告诉他们该封锁哪里！''

牢大此刻正被赵无极的另一道掌风逼得连连后退，但他的大脑已经在飞速运转。他的战场直觉在这一刻展现得淋漓尽致------尽管自身岌岌可危，他仍然在混乱中捕捉到了赵无极的移动规律。

``唐三！蓝银草封锁左后方三步，不要铺网，集中成线！诺当，干扰场收缩到赵老师右脚落点，等他起跳的瞬间释放！''

两道指令几乎同时出口，精准得像手术刀。唐三的蓝银草瞬间改变形态，从铺展的防御网收缩成一道细如发丝的草线，精准地缠绕在赵无极左后方的地面上；诺当的干扰场也随之收缩，从大面积铺展变为一个拳头大小的紫色光球，悬停在赵无极右脚即将落下的位置。

赵无极的攻势微微一滞------不是因为被击中，而是因为他的移动节奏被打乱了。他不得不调整落脚点，避开了诺当的干扰场，而这一调整，恰好让他的身体撞上了唐三的蓝银草线。草线虽细，却恰好卡在他腰间的魂力运转节点上，让他的第二掌慢了半拍。

就是这半拍的空隙，给了牢大喘息的机会。

``奥斯卡！''牢大吼了一声。

不需要更多的言语。奥斯卡早已准备好了------他的双手泛起柔和的金色光芒，一道增幅光环从掌心扩散而出，精准地笼罩在牢大身上。那不是蛮力的灌输，而是一种精妙的序数增幅------将牢大体内 \allowbreak{}\NovelInline{I01115}{\varphi(2,0,0)}\allowbreak{} 的递归蛮力，通过增幅光环的催化，让每一层递归迭代的效率提升数倍。

牢大的攻击力本就冠绝七人，再加上奥斯卡的增幅催化，他的拳风瞬间暴涨到了一个恐怖的程度。

\NovelDisplay{M00156R01}{%
f_{\varphi(2,0,0)}(3) = f_{\varphi(1,\varphi(1,\varphi(1,0,0),0),0)}(3)%
}

FGH 展开在增幅光环的催化下以前所未有的速度完成嵌套，拳风裹挟着 \allowbreak{}\NovelInline{I01116}{\varphi(2,0,0)}\allowbreak{} 的递归蛮力，轰然砸向赵无极的侧方。

赵无极眼中闪过一丝赞许，但手上的动作毫不含糊------他以掌接拳，轰然对撞，气浪将周围的序数碎片尽数震碎。牢大被震退三步，嘴角渗出血丝，但赵无极的攻势也被这记重拳打断了一瞬。

``好！''赵无极大喝一声，``牢大的攻击力确实恐怖，但如果没有奥斯卡的增幅，你刚才那一拳的 FGH 展开至少要慢两倍，根本打不断我的攻势。记住，蛮力要嵌进递归链，不是往外泄！''

他旋身变向，拳风陡然凌厉了数倍，直扑蔡须坤。

蔡须坤早有准备。他的篮球在掌心飞速旋转，序数脉络沿球面流淌，\allowbreak{}\NovelInline{I01117}{\varphi(1,\omega,0)}\allowbreak{} 的极限收敛规律在他指尖凝聚成一个极小的光点------那是他将全部序数力量压缩到一点后形成的单点突破。

\NovelDisplay{M00157R01}{%
\varphi(1,\omega,0)[n] = \varphi(1,n,0)%
}

\NovelDisplay{M00158R01}{%
f_{\varphi(1,\omega,0)}(3) = \sup\{f_{\varphi(1,0,0)}(3), f_{\varphi(1,1,0)}(3), f_{\varphi(1,2,0)}(3)\}%
}

他没有选择硬接赵无极的拳风，而是在赵无极的攻势即将触及身体的瞬间，侧身闪过，同时将手中的篮球如子弹般射出------那道凝聚了全部力量的光点，精准地命中了赵无极右肩的魂力运转节点。

``啪''的一声脆响，赵无极的右肩微微一麻，拳风偏了半寸。

``单点突破，干净利落！''赵无极赞了一声，但紧接着便以左掌补上，掌风比右拳更加凶猛，``但你的后续呢？突破之后没有第二击，等于白打！''

蔡须坤咬紧牙关，试图再次凝聚第二发单点突破，但体内的序数力量在第一次爆发后已经消耗了大半，光点在掌心闪烁了几下便黯淡下去。他的单点突破确实犀利，但持续输出能力不足------这是他的致命短板。

就在赵无极的左掌即将拍中蔡须坤的瞬间，一道身影从侧面闪电般切入。

小舞。

她的近身作战能力在这一刻展现得淋漓尽致------身形如鬼魅般闪现在赵无极身侧，双腿裹挟着 \allowbreak{}\NovelInline{I01118}{\varphi(1,0,0)}\allowbreak{} 的序数光芒，一记鞭腿直扫赵无极的腰间。她的速度太快了，快到赵无极都不得不侧身闪避，左掌的攻势被迫中断。

``漂亮！''赵无极眼中精光一闪，但紧接着，他的反击便到了------一记肘击直捣小舞面门。

小舞的防御力弱点在这一刻暴露无遗。她勉强以双臂格挡，但赵无极的肘击力量远超她的承受极限，双臂瞬间发麻，整个人被震退数步，踉跄着险些跌倒。

``小舞，你的近身爆发无人能及，但你的防御太脆了！''赵无极厉声喝道，``近身作战不是让你一个人冲进去拼命，是让你在队友创造的空隙中精准切入！牢大，给她指路！''

牢大抹掉嘴角的血迹，目光如炬。他在极短的时间内完成了战术判断------赵无极的右肩被蔡须坤的单点突破命中，魂力运转出现了短暂的滞涩，这意味着他的右侧防御有一个不到半息的窗口期。

``小舞，从右侧切入，不要正面！唐三，蓝银草缠住他的左脚，限制他的移动！诺当，干扰场放在他头顶，压他的起跳高度！卢本韦，大范围铺展牵制他的后退路线！''

四道指令同时下达，精准到可怕。

唐三的蓝银草如毒蛇般缠上赵无极的左脚脚踝；诺当的干扰场从头顶压下，将赵无极的起跳空间压缩到最低；卢本韦的大范围序数铺展如一张无形的网，封锁了赵无极所有后退的路径；而小舞，则从右侧那个不到半息的窗口中闪电切入，一记膝撞精准命中赵无极右肩的旧伤处。

赵无极闷哼一声，被迫后退两步。

但紧接着，他的反击如暴风骤雨般降临------一掌拍向防御最薄弱的小舞。小舞脸色大变，她刚完成近身攻击，正处于旧力已尽、新力未生的空档，根本来不及闪避。

``奥斯卡！''牢大再次怒吼。

奥斯卡的增幅光环瞬间转移到小舞身上，一道金色的防御薄膜在她体表凝聚------那不是攻击增幅，而是防御增幅。奥斯卡的攻击力虽然不强，但他的辅助能力是全方位的，增幅光环不仅能催化攻击力，也能在危急时刻转化为防御屏障。

赵无极的掌风撞上金色薄膜，冲击力被卸去了七成，小舞虽然仍被震退，但没有受伤。

``好配合！''赵无极终于露出了满意的笑容，但他手上的攻势没有丝毫减弱，``但还不够！牢大，你的战术指挥已经很不错了，但你的反应还是慢了半拍------小舞被反击的时候，你不该等奥斯卡来救，你应该提前预判到赵老师的反击路线，提前让奥斯卡准备好防御增幅！''

牢大咬紧牙关，他知道赵无极说得对。他的战场直觉虽然敏锐，但在高压环境下，他的战术反应仍然有延迟。这种延迟在对付弱敌时可以忽略不计，但面对赵无极这种级别的对手，半拍的延迟就可能造成致命的后果。

``再来！''牢大吼了一声，肘间的 \allowbreak{}\NovelInline{I01119}{\varphi(2,0,0)}\allowbreak{} 涟漪再次暴涨。

奥斯卡的增幅光环稳稳笼罩在他身上，金色光芒与 \allowbreak{}\NovelInline{I01120}{\varphi(2,0,0)}\allowbreak{} 的递归蛮力交融，爆发出远超单体的恐怖攻击力。

第一天收训时，七个人横七竖八地瘫在炼狱场的碎石地上，浑身上下没有一块好肉。唐三的蓝银草蔫成了枯草色，软塌塌地垂在地面上，连自动收回武魂的力气都没有了。小舞仰面朝天躺着，双臂呈大字型摊开，身上的淤青密密麻麻，像一幅抽象画。蔡须坤的篮球滚在一旁，球面上的序数纹路都磨花了。卢本韦侧卧着，一动不动，像一尊被遗弃的石像。诺当靠着墙壁，双眼无神地盯着天空，嘴唇微微翕动，似乎在默念什么。奥斯卡趴在地上，脸埋在胳膊里，闷闷地发出呜咽声。

牢大是七个人里唯一还坐着的。他的后背靠在一块断裂的石柱上，双腿伸直，浑身上下也青一块紫一块，但至少还能维持坐姿。他的目光越过众人，落在炼狱场另一端------赵无极正坐在一把不知从哪搬来的太师椅上，翘着二郎腿，手里端着一杯热茶，慢悠悠地吹着茶沫，神态安详得像是来郊游的。

``赵老师。''小舞率先开口了，声音虚弱得像蚊蚋，``我怀疑你是来报复我们的。''

``哦？''赵无极吹了一口茶，不紧不慢。

``上次不乐那件事，你明明可以提前赶到的。''小舞翻了个身，疼得龇牙咧嘴，``你故意晚到了半刻钟，就是为了让我们挨一顿打，好有理由搞这个什么特训。对不对？''

``你想多了。''赵无极喝了口茶，``老夫是真的没收到信号。''

``那你现在为什么笑得那么开心？''蔡须坤冷冷地插了一句，声音里带着罕见的锋利。他坐起身，揉了揉被震得通红的手掌，目光不善地盯着赵无极手中的茶杯，``我们被你用 \allowbreak{}\NovelInline{I01121}{\varphi(2,0,0)}\allowbreak{} 级别的掌风拍了整整六个时辰，浑身骨头都快散架了，你坐在那里喝茶看戏，还笑得出来。''

``那不是看戏，那叫观察。''赵无极放下茶杯，一本正经地纠正，``老夫在观察你们的配合漏洞，好制定明天的训练计划。''

``明天的训练计划？''奥斯卡从地上抬起头，脸上还沾着泥土，眼中满是绝望，``赵老师，明天还有？''

``当然有。''赵无极理所当然地点了点头，``今天只是热身，明天才是正餐。''

``热身？''小舞的声音陡然拔高了八度，差点从地上弹起来，但一动就牵动了浑身的伤，又''嘶''地一声倒了回去，``你用 \allowbreak{}\NovelInline{I01122}{\varphi(2,0,0)}\allowbreak{} 的掌风打我们叫热身？你怎么不用 \allowbreak{}\NovelInline{I01123}{\varphi(3,0,0)}\allowbreak{} 直接把我们送走算了？''

``那多没意思。''赵无极笑了笑，``打死了谁给我练？''

诺当终于开口了。他从头到尾一直靠着墙壁沉默不语，此刻缓缓转过头，目光空洞地看着赵无极，声音平淡得像在念悼词：

``赵老师，我有一个问题。''

``说。''

``您今年贵庚？''

赵无极微微一愣：``老夫七十有三，怎么了？''

``七十有三的人了，''诺当的声音依旧平淡，却字字诛心，``精力还这么旺盛，每天天不亮就把我们从被窝里拽出来，打到日落西山，中间只给一顿饭、两壶水。您不累吗？''

``老夫当然不累。''赵无极翘起二郎腿晃了晃，``老夫修为高，恢复快。''

``所以我们累。''诺当说完这四个字，便不再开口，重新闭上了眼睛。那副''我已经放弃了''的姿态，比任何控诉都有力。

卢本韦始终没有说话。他只是默默地从地上捡起一块碎石，在手指间翻来覆去地捏着，目光落在赵无极身上，眼神幽怨得像一口深井。他不需要开口------那个眼神本身就是一篇控诉。

``行了行了，你们别怨我。''赵无极被这几道目光盯得有些不自在，放下茶杯，干咳了两声，``老夫这也是为你们好。你们想想，上次遇到不乐，你们被一个 \allowbreak{}\NovelInline{I01124}{\varepsilon_1}\allowbreak{} 的蛮夫打得手忙脚乱，最后还是老夫出手收拾的残局。如果下次遇到比不乐更强的对手，老夫不在呢？''

``那我们就跑。''小舞脱口而出。

``跑？''赵无极嗤笑一声，``你跑得掉吗？''

小舞张了张嘴，没有反驳。她跑得确实不快------至少没有赵无极快。

``赵老师，''唐三终于开口了，他的声音虽然虚弱，但条理依然清晰，``您的训练理念我理解，以战养悟、在实战中淬炼配合。但有一个问题------您给的压力太大了。今天是第一天，您就直接用 \allowbreak{}\NovelInline{I01125}{\varphi(2,0,0)}\allowbreak{} 级别的攻势打我们，我们的根基还没有完全稳固，这种强度的冲击，不但不能帮我们突破瓶颈，反而可能导致魂力脉络受损。''

``哦？''赵无极挑了挑眉，``那你觉得，应该用什么强度？''

``至少先从 \allowbreak{}\NovelInline{I01126}{\varphi(1,0,1)}\allowbreak{} 开始，循序渐进。''唐三认真地说道。

``循序渐进？''赵无极摇了摇头，``唐三，你犯了一个错误------你以为修炼是爬楼梯，一步一步往上走。但实战不是。实战不会等你准备好了再开始，敌人不会因为你还没突破就手下留情。上次不乐可没有等你把 \allowbreak{}\NovelInline{I01127}{\varepsilon_0}\allowbreak{} 的基本列推演完毕再出手。''

唐三沉默了。他知道赵无极说的有道理，但道理归道理，浑身上下每一块骨头都在抗议。

``赵老师，''奥斯卡弱弱地举起手，像一只在课堂上不敢大声说话的学生，``我有个小建议\ldots\ldots{}''

``说。''

``您能不能\ldots\ldots 在打我们的时候，顺便给点治疗？''他的声音越来越小，``我是辅助系，我的增幅光环可以给你们加攻击、加防御，但加不了血。您打完了不管治，我们每天都是带着一身新伤叠旧伤上场的。这样下去，别说突破了，能不能活过这周都是问题。''

``你小子。''赵无极笑骂了一声，从椅子上站起来，走到众人面前。他从怀里掏出一个瓷瓶，随手扔在地上，``这是老夫配的金疮药，每天收训之后自己抹。别指望老夫给你们治伤------老夫是教头，不是大夫。''

``教头也不带这么教的。''蔡须坤小声嘟囔了一句，``别的教头是教学生怎么打球，您是直接把球往学生脸上砸。''

``球砸脸上才能练出反应速度。''赵无极面不改色。

``那您怎么不砸自己脸上练练？''小舞接了一句。

空气瞬间安静了。

赵无极低头看着小舞，小舞仰面看着赵无极，两个人大眼瞪小眼了三秒钟。

``\ldots\ldots 你刚才说什么？''赵无极的声音平静得可怕。

``我说------''小舞的求生欲在这一刻达到了巅峰，``我说赵老师您英明神武、德高望重、训练有方、我们心服口服！''

赵无极冷哼一声，重新坐回太师椅上，端起茶杯。

``明天辰时，炼狱场集合。迟到一刻钟，加练两个时辰。''

七个人同时发出了一声哀嚎。

牢大始终没有参与这场埋怨。他只是安静地坐在那里，看着队友们你一言我一语地吐槽赵无极，嘴角微微上扬。他知道，这种埋怨不是真的怨恨------他们心里都清楚，赵无极的严苛是为了他们好。但允许他们发发牢骚、吐吐槽，本身也是团队凝聚的一种方式。

等众人的抱怨声渐渐平息，牢大才缓缓开口，声音低沉而沉稳：

``都省点力气吧。明天还要挨打。''

``\ldots\ldots{}''

六个人同时沉默了。

然后，不知道是谁先笑了一声，紧接着所有人都笑了。笑声在炼狱场中回荡，带着疲惫、带着无奈、也带着一种只有共同经历过苦难的人才能理解的默契。

赵无极坐在太师椅上，听着这阵笑声，嘴角也扬起了一抹不易察觉的弧度。

他端起茶杯，轻轻抿了一口。

茶已经凉了。

\chapter{第十七章：扽魂融合技的威力}\label{ux7b2cux5341ux4e03ux7ae0ux627dux9b42ux878dux5408ux6280ux7684ux5a01ux529b}

炼狱场中，七道身影交错纵横，魂技碰撞声不绝于耳。牢大在奥斯卡的增幅下化作最锋利的矛，每一拳都裹挟着 \allowbreak{}\NovelInline{I01128}{\varphi(2,0,0)}\allowbreak{} 的递归蛮力，打得赵无极不得不认真对待；蔡须坤在牢大的战术指引下，精准地找到赵无极防御的薄弱点，一记记单点突破如针般刺入，虽然无法造成致命伤害，却不断蚕食着赵无极的魂力运转效率；

小舞在唐三和诺当的控制掩护下，一次次从刁钻的角度近身切入，虽然每次都被赵无极的反击震退，但她的近身作战经验在一次次失败中飞速积累；唐三和诺当在牢大的指挥下，控制能力发挥到了极致，蓝银草的封锁与干扰场的压制相互配合，将赵无极的移动空间压缩到了最小；卢本韦则在大范围铺展中扮演着清道夫的角色，他的群体压制虽然单点攻击力薄弱，却有效地牵制了赵无极的走位，让赵无极无法集中力量逐个击破。

而奥斯卡，这个攻击力最弱的辅助，却是整个团队不可或缺的粘合剂。他的增幅光环在七人之间流转，时而催化牢大的攻击力，时而加固小舞的防御，时而为蔡须坤的单点突破提供额外的序数能量，时而又为唐三和诺当的控制场提供持续的魂力补充。他的攻击力虽然不强，但他的辅助能力让整个团队的战力提升了不止一个层级。

赵无极在七人的围攻下，第一次感受到了压力。

不是修为上的压力------他的修为远超七人中的任何一人，哪怕七人联手，在绝对力量上也无法与他抗衡。但配合上的压力是真实的------七个人各司其职、各展所长，在牢大的战术指挥下形成了一个近乎完美的攻防体系。他的每一次进攻都会被控制系限制走位，每一次反击都会被辅助系提前化解，每一次试图逐个击破都会被群体压制牵制。

``不错。''赵无极在连续接了牢大三板斧之后，终于停下了攻势，站在炼狱场中央，大口喘着粗气，嘴角却挂着一抹欣慰的笑容，``这才像一个团队。''

他看向牢大，目光中多了一份郑重：``你的攻击力是矛，奥斯卡的增幅是矛柄，蔡须坤的单点突破是针，小舞的近身缠斗是匕首，唐三和诺当的控制是网，卢本韦的群体压制是墙。矛、柄、针、匕首、网、墙------六样东西拼在一起，才是一副完整的铠甲。缺了任何一样，都不完整。''

他顿了顿，声音变得深沉：``但铠甲的核心是什么？是穿铠甲的人。牢大，你就是那个穿铠甲的人。你的战术指挥，就是铠甲的灵魂。没有你的指挥，矛不知道往哪里刺，网不知道往哪里撒，针不知道往哪里扎。你的攻击力固然重要，但你的战场直觉和战术判断力，才是这个团队真正的核心。''

牢大沉默了片刻，缓缓点了点头。他终于明白了玉小刚今天开场那番话的深意------他不只是队伍里攻击力最强的人，他还是整个团队的战术核心。这个认知让他感到沉甸甸的压力，但也让他第一次清晰地看到了自己的价值所在。

``休息半刻钟。''赵无极摆了摆手，一屁股坐在地上，``然后继续。刚才的配合勉强及格，但衔接还有太多漏洞------唐三的控制和牢大的指令之间有一拍的延迟，小舞的近身切入时机还不够精准，蔡须坤的单点突破之后没有及时撤回，差点被我的反击拍中。这些都是要在接下来的训练中反复打磨的东西。''

炼狱场中，七道身影交错纵横，魂技碰撞声不绝于耳。 牢大在奥斯卡的增幅下化作最锋利的矛，每一拳都裹挟着 φ(2,0,0) 的递归蛮力，打得赵无极不得不认真对待；蔡须坤在牢大的战术指引下，精准地找到赵无极防御的薄弱点，一记记单点突破如针般刺入，虽然无法造成致命伤害，却不断蚕食着赵无极的魂力运转效率；

小舞在唐三和诺当的控制掩护下，一次次从牢大指定的刁钻角度近身切入，虽然每次都被赵无极的反击震退，但她的近身作战经验在一次次失败中飞速积累；唐三和诺当在牢大的指挥下，控制能力发挥到了极致，蓝银草的封锁与干扰场的压制相互配合，将赵无极的移动空间压缩到了最小；

卢本韦则在大范围铺展中扮演着清道夫的角色，他的群体压制虽然单点攻击力薄弱，却有效地牵制了赵无极的走位，让赵无极无法集中力量逐个击破。

而奥斯卡，这个攻击力最弱的辅助，却是整个团队不可或缺的粘合剂。他的增幅光环在牢大的调度下精准流转，时而催化牢大的攻击力，时而加固小舞的防御，时而为蔡须坤的单点突破提供额外的序数能量，时而又为唐三和诺当的控制场提供持续的魂力补充。他的攻击力虽然不强，但他的辅助能力在牢大的战术框架中，让整个团队的战力提升了不止一个层级。

赵无极在七人的围攻下，第一次感受到了真正的压力。不是修为上的压力------他的修为远超七人中的任何一人，哪怕七人联手，在绝对力量上也无法与他抗衡。但配合上的压力是真实的------七个人各司其职、各展所长，在牢大的战术指挥下形成了一个近乎完美的攻防体系。

他的每一次进攻都会被控制系限制走位，每一次反击都会被辅助系提前化解，每一次试图逐个击破都会被群体压制牵制。``不错。''赵无极在连续接了牢大三板斧之后，终于停下了攻势，站在炼狱场中央，大口喘着粗气，嘴角却挂着一抹欣慰的笑容，``这才像一个团队。''

他看向牢大，目光中多了一份郑重：``牢大，你这个队长，当得不错。'' 牢大抹掉嘴角的血迹，沉声说道：``还不够。第二轮，我会把衔接压缩到两次呼吸之内。'' 赵无极哈哈大笑，浑厚的笑声在炼狱场中回荡：``好！那老夫就等着看，你的第二次呼吸，有多快！''

``逐层夯实，别贪快！先把 \allowbreak{}\NovelInline{I01129}{n=0}\allowbreak{}、\allowbreak{}\NovelInline{I01130}{1}\allowbreak{}、\allowbreak{}\NovelInline{I01131}{2}\allowbreak{} 的 FGH 展开算透，再收拢极限！''赵无极足尖轻点落地的篮球，魂力注入其中，帮他校准每层迭代节奏，``\allowbreak{}\NovelInline{I01132}{\varphi(1,0,0)}\allowbreak{} 对应 \allowbreak{}\NovelInline{I01133}{\Gamma_0}\allowbreak{}，FGH 展开 \allowbreak{}\NovelInline{I01134}{f_{\Gamma_0}(3)}\allowbreak{}；\allowbreak{}\NovelInline{I01135}{\varphi(1,1,0)}\allowbreak{} 对应三层迭代，\allowbreak{}\NovelInline{I01136}{f_{\varphi(1,1,0)}(3)}\allowbreak{}，一层一层叠！''

蔡须坤咬牙稳住身形，放慢篮球转速，逐层级嵌入序数力量，同步完成各阶 FGH 运算，篮球转动愈发平稳，极限光晕层层收拢，不再散乱。

与此同时，赵无极反手攻向卢本韦、诺当、小舞三人，专攻他们 \allowbreak{}\NovelInline{I01137}{\Gamma_0}\allowbreak{} 根基薄弱处，逼迫三人抱团配合。

卢本韦率先催动母猪扽魂，施展母猪冲撞，贴合 \allowbreak{}\NovelInline{I01138}{\varphi(1,0,0)}\allowbreak{} 的不动点闭环发力，同步展开 FGH：

\NovelDisplay{M00159R01}{%
f_{\varphi(1,0,0)}(3) = f_{\Gamma_0}(3)%
}

以蛮横韧性正面牵制赵无极；诺当紧随其后甩出磁暴环，环身绽放魂力，磁暴之力吸附序数波动，FGH 展开配合干扰：

\NovelDisplay{M00160R01}{%
f_{\varphi(\varphi(2,0),\omega)}(3) = f_{\varphi(φ(2,0),3)}(3)%
}

打乱赵无极的攻势节奏；小舞施展柔骨魅兔魂技，身形灵动穿梭，顺着不动点脉络缠向赵无极腿脚，以柔劲辅助稳住全队阵型。

三人各司其职，魂技与序数展开无缝衔接，卢本韦主攻牵制、诺当控场干扰、小舞柔劲补位，即便浑身酸痛、魂力消耗巨大，也始终配合默契，没有出现丝毫破绽。

奥斯卡守在阵型最中心，顶着全场压力，专攻 \allowbreak{}\NovelInline{I01139}{\varphi(1,0,\varphi(1,0))}\allowbreak{} 双层嵌套序数，一边完成高难度 FGH 展开，一边释放辅助魂技兜底全队。他屏气凝神，先拆解内层序数，展开基础 FGH：

\NovelDisplay{M00161R01}{%
f_{\varphi(1,0,0)}(3) = f_{\Gamma_0}(3)%
}

再向外层延伸，攻克双层嵌套：

\NovelDisplay{M00162R01}{%
f_{\varphi(1,0,\varphi(1,0,0))}(3) = f_{\varphi(1,0,\varphi(\omega,\varphi(\omega,\varphi(\omega,0))))}(3)%
}

赵无极故意分出一道魂力冲击，直逼奥斯卡脉络，逼他在高压下完成完整展开：``别分心，双层嵌套要先锁内层、再推外层，FGH 一步都不能错！''奥斯卡头晕目眩，魂力几近枯竭，却咬牙坚持，在唐三、牢大的掩护下，终于完成双层 FGH 全展开，同时抛出增益魂剂，为队友补充魂力、梳理紊乱脉络，成为全队最稳固的后盾。

\begin{center}\rule{0.5\linewidth}{0.5pt}\end{center}

本来，面对赵无极的攻击和防御，众人肯定是无法承受的，即使是赵无极收手，也很难说是抵挡，即使是牢大突破 \allowbreak{}\NovelInline{I01140}{\varphi(\omega,0,0)}\allowbreak{} 也不能抵挡半分。

有一次普通的训练，卢本韦被赵无极掀翻，重重摔在地上，奥斯卡见状赶紧发射大香肠支援，这样的场景他已经经历过很多次了。特训的时候，这样的过程每天发生很多次，但是今天有所不同------奥斯卡错误地发射了攻击性香肠，眼看这一操作会导致队友受伤，卢本韦察觉这两个大香肠后，只得发射母猪应对。赵无极见他们失误，本来想批评这一举动，没想到，这两根香肠恰好进入了母猪的两个鼻孔里面，顿时卢本韦的母猪大哼一声。这母猪鼻子本就顺滑，奥斯卡两根香肠进入后，母猪下意识地缩紧了鼻子，两根香肠让母猪无法摄入氧气，只能产生无氧呼吸：

\NovelDisplay{M00163R01}{%
\mathrm{C_6H_{12}O_6} \rightarrow 2\mathrm{C_2H_5OH} + 2\mathrm{CO_2}%
}

反应在母猪体内爆发，瞬间，母猪的鼻孔产生大量气体，极短时间内产生大量气体导致鼻子内部气压升高，最终香肠发射出去了。好巧不巧，香肠正对着赵无极的方向打去，赵无极看两人招式乱了，也恐怕发动不了什么太大的攻击，结果让他意想不到的是，这香肠以极快的速度发射，居然直接穿透了他的防御。

这一击的威力远超 SVO 的水平，甚至隐隐有 \allowbreak{}\NovelInline{I01141}{\psi_\Omega(\Omega_2)}\allowbreak{} 的水平。要知道 \allowbreak{}\NovelInline{I01142}{\psi_\Omega(\Omega_2)}\allowbreak{} 是一个巨大的序数（后面会讲），是 \allowbreak{}\NovelInline{I01143}{\varphi}\allowbreak{} 函数可望不可及的高度。虽说由于攻击力仅仅达到了 \allowbreak{}\NovelInline{I01144}{\psi_\Omega(\varepsilon_{\Omega+1})}\allowbreak{} 多一点的水平（也很大，比上一个小点），但是也是出乎赵无极的预料了。他没有防住，一下子击破他的防御，这一根香肠击中头部，他连连后退。按理说他这样子的实力应该不会影响的，他疏忽了，不过他却暗暗高兴，因为他们的战斗力现在变强了很多。

看不懂的可以类比幽冥白虎。

\begin{center}\rule{0.5\linewidth}{0.5pt}\end{center}

他们将 \allowbreak{}\NovelInline{I01145}{\varphi}\allowbreak{} 函数的最后一位的规则由不动点进位改为''最后一位+1 对应于序数值''，得到 \allowbreak{}\NovelInline{I01146}{\phi}\allowbreak{} 函数：

\NovelDisplay{M00164R01}{%
\phi(S, \alpha+1) = \phi(S, \alpha) + 1%
}

这样，原来的 \allowbreak{}\NovelInline{I01147}{\phi(2,0,0) = \varphi(1,1)}\allowbreak{}，他们很厉害，直接让两个函数追平，也就是说，在一个足够大的序数后，这两种函数表示的序数一样。换言之：

设 \allowbreak{}\NovelInline{I01148}{\psi_A}\allowbreak{}、\allowbreak{}\NovelInline{I01149}{\psi_B}\allowbreak{} 是两个将序数映射为序数的函数，若存在序数 \allowbreak{}\NovelInline{I01150}{\alpha}\allowbreak{} 使得

\NovelDisplay{M00165R01}{%
\psi_A(\alpha) \leq \psi_B(\alpha) < \psi_A(\alpha+1)%
}

则称 \allowbreak{}\NovelInline{I01151}{\psi_A}\allowbreak{} 与 \allowbreak{}\NovelInline{I01152}{\psi_B}\allowbreak{} 在 \allowbreak{}\NovelInline{I01153}{\alpha}\allowbreak{} 处追平，\allowbreak{}\NovelInline{I01154}{\alpha}\allowbreak{} 称为\textbf{追平点}。

我们可以考虑如下的两个序数函数：

\NovelDisplay{M00166R01}{%
\psi_A(\alpha) = \alpha, \quad \psi_B(\alpha) = 2 \cdot \alpha%
}

直观地说，这两个序数函数的增长过程可以被视为两个人赛跑，这两个人以相同的节奏跑动，其中第一个人每步跑 1 米，第二个人每步跑 2 米。在最开始，两个人自然是站在同一起跑线上的，即有：

\NovelDisplay{M00167R01}{%
\psi_A(0) = \psi_B(0) = 0%
}

随着两人跑的步数越来越多，他们的距离似乎会越拉越大，至少在有限步之内是这样的。但是令人惊奇的是，在两个人跑了 \allowbreak{}\NovelInline{I01155}{\omega}\allowbreak{} 步之后，他们竟然奇迹般地相遇了：

\NovelDisplay{M00168R01}{%
\psi_A(\omega) = \psi_B(\omega) = \omega%
}

这就是说，序数 \allowbreak{}\NovelInline{I01156}{\omega}\allowbreak{} 抹平了二人之间的差距------两个人让弱的追平强的。记录第一个追平点是 \allowbreak{}\NovelInline{I01157}{\varphi_1(0)}\allowbreak{}，第二个追平点是 \allowbreak{}\NovelInline{I01158}{\varphi_1(1)}\allowbreak{}，让 \allowbreak{}\NovelInline{I01159}{\phi}\allowbreak{} 追平 \allowbreak{}\NovelInline{I01160}{\varphi_1}\allowbreak{} 得到 \allowbreak{}\NovelInline{I01161}{\varphi_2(0)}\allowbreak{}，以此类推得到 \allowbreak{}\NovelInline{I01162}{\varphi_\omega(0)}\allowbreak{}。这香肠承载的力量，正是这 \allowbreak{}\NovelInline{I01163}{\varphi_\omega(0)}\allowbreak{} 得到了突破，以及团队的力量。

那记裹挟着 \allowbreak{}\NovelInline{I01164}{\varphi_\omega(0)}\allowbreak{} 磅礴力量的香肠破空而出，狠狠砸中赵无极的刹那，整个炼狱场都陷入死寂。众人僵在原地，看着赵无极踉跄后退的身影，手心冷汗直冒，方才慌乱间的失误操作，竟误打误撞撞破了序数修炼的全新壁垒，连他们自己都不敢置信。

\begin{center}\rule{0.5\linewidth}{0.5pt}\end{center}

赵无极缓缓直起身，抬手抹去唇角淡淡的血迹，非但没有怒意，反倒仰天大笑，浑厚的笑声震得禁制微微颤动：``好！好得很！老夫没白逼你们这一场！''他望着满脸错愕的众人，眼神里的冷峻尽数化作欣慰，``你们以为这只是巧合？若非这四十九天的炼狱打磨，若非你们把不动点迭代、FGH 展开刻进本能，这般绝境逢生的追平契机，就算送上门来也抓不住！''

众人怔怔站着，浑身的伤痛仿佛在这一刻尽数消散，心底积压的怨气与不甘，也随着这声朗笑烟消云散。

回想这数十日的煎熬，他们曾恨极了赵无极的狠辣无情，恨极了没有喘息的死战、永不停歇的推演，恨极了每一次魂力耗尽后的剧痛、每一次迭代失败的挫败。唐三无数次对着 \allowbreak{}\NovelInline{I01165}{\varphi(1,1,0)}\allowbreak{} 的脉络彻夜苦思，蓝银草被序数碎片割得支离破碎；牢大一遍遍冲撞递归链条，经脉反复受损又强行愈合；蔡须坤的双手被篮球磨得血肉模糊，只为校准那一丝极限收敛波动；卢本韦、诺当、小舞三人抱团硬扛威压，魂力透支昏厥是家常便饭；奥斯卡更是顶着巨大压力，在双层嵌套的迷雾里反复摸索，数次脉络断裂险些伤及根基。

他们曾无数次盼着这场炼狱早点结束，盼着重回安逸的课堂，可此刻望着彼此衣衫破烂、伤痕累累却眼神发亮的模样，心底只剩下浓烈的感激。

感激这场没有退路的死训，逼他们摒弃了浮躁空想，把晦涩的序数理论揉碎了融进骨子里；感激三位老师的严苛磨砺，让他们在生死边缘逼出了潜藏的悟性，打破了固有的思维桎梏；更感激身边不离不弃的同伴，在一次次跌倒后搀扶，在一次次失误中共担，让孤军奋战变成了众志成城。

\begin{center}\rule{0.5\linewidth}{0.5pt}\end{center}

唐三掌心的蓝银草轻轻颤动，周身 \allowbreak{}\NovelInline{I01166}{\varphi(1,1,0)}\allowbreak{} 的光晕愈发凝练，原本滞涩的三元不动点脉络彻底通畅，FGH 迭代行云流水，甚至能隐约触碰到更深层级的序数门槛；牢大周身蛮力与 \allowbreak{}\NovelInline{I01167}{\varphi(2,0,0)}\allowbreak{} 递归链完美契合，每一次吐纳都带着沉稳的序数波动，再也不会出现发力紊乱的窘境；蔡须坤手中的篮球缓缓转动，\allowbreak{}\NovelInline{I01168}{\varphi(1,\omega,0)}\allowbreak{} 的极限收敛光晕层层收拢，彻底跨过了此前难以逾越的关卡；卢本韦、诺当、小舞三人周身 \allowbreak{}\NovelInline{I01169}{\Gamma_0}\allowbreak{} 根基坚如磐石，\allowbreak{}\NovelInline{I01170}{\varphi(1,0,0)}\allowbreak{} 的不动点闭环稳固无比，配合之间毫无破绽；奥斯卡眼底闪过透亮的光芒，双层嵌套序数的 FGH 展开信手拈来，甚至能举一反三，摸索出更多辅助推演的技巧。

而最珍贵的收获，远不止修为的突破。他们读懂了序数修炼从不是闭门造车，而是心性与默契的双重淬炼；明白了绝境之中的坚守，远胜无数次安逸的参悟；更懂得了团队羁绊迸发的力量，足以突破看似不可逾越的天堑，让弱势的函数追平强势，让不可能变成可能。

弗兰德与玉小刚相视一笑，眼底满是赞许。赵无极拍了拍身上的尘土，语气依旧硬朗，却少了几分凌厉，多了几分期许：``别得意太早，这只是起点，高阶序数的路还长着呢。''

众人没有应声，只是不约而同地看向身边的同伴，抬手紧紧相握。周身的序数光晕交织缠绕，不再是此前的忽明忽暗，而是汇成一股沉稳厚重的光流，那是突破瓶颈的底气，是生死与共的羁绊，更是破茧重生的坚毅。

曾经痛恨的炼狱，如今成了蜕变的熔炉；曾经避之不及的严苛，如今成了进阶的阶梯。他们望着远处连绵的后山，眼神坚定，再也没有了此前的迷茫与困顿。这场炼狱特训，磨平了心性的浮躁，筑牢了修为的根基，更让他们明白，真正的强者，从不是顺境中高歌，而是绝境中破局，在痛苦中成长，在羁绊中前行。

\begin{center}\rule{0.5\linewidth}{0.5pt}\end{center}

炼狱场的高阶禁制缓缓散去，后山的清风裹挟着草木气息扑面而来，吹散了连日来的血腥与疲惫。四十九天的炼狱特训，彻底改写了这支小队的修为格局，也磨出了刻入骨髓的韧性与默契。

唐三周身序数光晕流转，已然突破至 \allowbreak{}\NovelInline{I01171}{\varphi(1,0,0,0)}\allowbreak{} 层级，三元不动点的桎梏彻底打破，四元序数脉络初现雏形，蓝银草武魂与序数递归完美共鸣，哪怕不动用魂力，也能隐约感知周遭的序数波动；牢大与蔡须坤双双踏足 SVO 境界，蛮力与校准技巧尽数升华，拳风、球影之间皆暗含高阶序数节律，战力较此前翻了数倍；诺当、小舞等人也顺利突破 \allowbreak{}\NovelInline{I01172}{\varphi(2,0,0)}\allowbreak{}，\allowbreak{}\NovelInline{I01173}{\Gamma_0}\allowbreak{} 根基稳固无匹，不动点闭环坚不可摧；奥斯卡的双层嵌套 FGH 推演早已炉火纯青，甚至能触碰到多层序数的辅助法门，成为小队最可靠的后盾。

曾经让他们寸步难行的序数壁垒，如今已成脚下基石；曾经痛恨至极的残酷特训，此刻尽数化作破茧成蝶的底气。众人站在炼狱场出口，回望这片洒满汗水与血水的土地，眼底没有怨怼，只有释然与敬重------若非这场破釜沉舟的磨砺，他们依旧是困在静修室里苦苦参悟的井底之蛙，永远触不到高阶序数的门槛。

弗兰德、玉小刚与赵无极三位老师伫立在不远处，看着脱胎换骨的众人，脸上满是欣慰。玉小刚合上手中的序数典籍，轻声开口：``你们的成长，超出了我们的预料。但序数之路永无止境，切不可因此骄躁。''

``少废话，以后修炼别想偷懒！''赵无极依旧是一副硬朗模样，却抬手扔给众人几瓶修复魂剂，语气不自觉柔和了几分，``好好休整三日，后续的修炼，只会更难。''

众人躬身行礼，心底的感激溢于言表。这三位老师看似严苛，却处处留一线生机，每一次呵斥、每一击施压，都是为了逼他们突破自我，这般良苦用心，他们早已彻悟。

\chapter{第十八章：扽魂殿的邀请}\label{ux7b2cux5341ux516bux7ae0ux627dux9b42ux6bbfux7684ux9080ux8bf7}

休整的三日里，小队并未松懈，而是趁热打铁巩固修为，把特训期间领悟的序数逻辑、FGH迭代技巧反复打磨，让实战经验彻底转化为本能。唐三时常静坐参悟四元序数的映射规律，牢大对着山石演练递归蛮力，蔡须坤借着篮球校准极限收敛，卢本韦、诺当、小舞三人磨合配合细节，奥斯卡则潜心钻研辅助型序数魂技，整个小队的凝聚力愈发强悍。

就在众人休整完毕，准备开启新一轮修炼时，一道金光破空而来，径直落在学院广场中央，化作一枚镌刻着繁复序数纹路的金色请柬，请柬周身散发着威严而中正的气息，一看便知出自高阶势力之手。

唐三上前拾起请柬，指尖触碰的瞬间，一股温和的魂力涌入体内，请柬上的文字缓缓浮现：烫金的字体苍劲有力，落款处赫然是------\textbf{扽魂殿}。

小队众人瞬间围拢过来，眼神中满是激动与郑重。

扽魂殿，是整个帝国最顶尖的正义组织，坐拥无数强者，执掌扽魂觉醒、等级评定、秩序维护大权，与无恶不作的帅气集团、美丽集团常年抗衡，守护着整片大陆的安定。但凡心怀正义的扽魂师，无不以加入扽魂殿为荣。

唐三攥紧请柬，思绪飘回年少之时：他自幼孤苦，正是扽魂殿的使者无偿为他觉醒扽魂，指引他踏上序数修炼之路，这份恩情，他始终铭记于心。他抬头看向同伴，声音带着难掩的热忱：

\begin{quote}
``扽魂殿邀请我们前往总部，参与精英扽魂师登记，这是我们的机遇，也是践行初心的机会。''
\end{quote}

``守护世界，惩治恶人！''小舞眼睛发亮，柔骨魅兔武魂轻轻颤动，满是期待。

牢大攥紧拳头，浑身蛮力蠢蠢欲动：``早就想会会那些作恶多端的家伙，扽魂殿，正好是我们的舞台！''

蔡须坤转动手中篮球，序数光晕在球面流转：``既能精进修为，又能坚守正义，此行非去不可。''

卢本韦、诺当、奥斯卡等人也纷纷点头，眼神坚定。他们历经炼狱磨砺，早已不是只顾修为的散兵游勇，而是心怀信念、彼此羁绊的小队，渴望用自己的力量守护一方安定。

请柬上标注了集合时限与地点，扽魂殿总部坐落于帝国核心腹地，汇聚了整片大陆的顶尖强者，更是对抗帅气、美丽两大邪恶集团的前沿阵地。众人深知，此行既是机遇，也暗藏凶险，高阶序数的较量、正邪势力的交锋，远比学院特训更加残酷。

但他们毫无惧色。

历经炼狱生死，他们早已练就钢铁心性；同伴并肩而立，他们便拥有破局一切的底气。

众人迅速收拾行装，告别三位老师，踏上了前往扽魂殿总部的路途。阳光洒在他们身上，将身影拉得修长，周身的序数光晕交织成盾，既是修为的象征，也是守护正义的誓言。

一行人沿着官道行至半途，日头渐渐西斜，Oscar擦了擦额头的薄汗，忍不住开口：

\begin{quote}
``从这里到扽魂殿总部，还要穿过这片荒林，路程不短，大家注意戒备。''
\end{quote}

唐三颔首，目光扫过四周暗沉的林木：

\begin{quote}
``这片林子太过僻静，少有人烟，容易藏着隐患，都打起精神。''
\end{quote}

话音刚落，众人转而踏入荒僻古林腹地，此地远离城镇，古木参天蔽日，粗壮的枝桠交错纠缠，将天光切割成细碎的光斑，地上积着厚厚的腐叶，踩上去发出细碎的闷响，林间弥漫着阴冷潮湿的雾气，连空气都透着压抑的凝滞。

``不对劲，这雾气里有邪性魂力，还混着一股陌生的古老序数波动。''蔡须坤指尖轻转篮球，球面光晕突然剧烈震颤，竟泛起一层极淡的银灰色纹路，绝非寻常邪修所有。忽然间，林深处炸开一股紊乱的黑色序数波动，气息阴邪暴戾，如同毒蛇吐信般蔓延开来，周遭的草木瞬间被这股邪力侵蚀，叶片发黑卷曲，连枝头的鸟鸣都戛然而止，死寂得可怕。更诡异的是，雾气深处隐约闪过一道极快的银影，转瞬即逝，紫极魔瞳全开的唐三也只捕捉到一抹残影，完全无法辨识来路。

``是帅气集团的人，但这股波动不纯，背后还有别的势力！''唐三眉心紫极魔瞳骤然睁开，紫芒刺破浓雾，脸色微沉，``阵型戒备！小舞左翼、牢大正面，诺当、卢本韦守右翼！''

众人瞬间身形错落形成攻防阵，只见几名身着玄黑劲装、面带煞气的邪异修士，脖颈处都纹着一枚从未见过的暗金色纹路图腾，绝非普通喽啰。他们正呈合围之势围猎一名落单的低阶扽魂师，黑色魂力缠如毒蛇，口中还念念有词：

\begin{quote}
``快取他的序数根基，大人要的东西不能耽误！''
\end{quote}

那名扽魂师瘫倒在地，周身光晕黯淡，已是强弩之末。

``住手！''小舞身形率先窜出，脚下腐叶飞溅，厉声呵斥道，``光天化日之下截杀同脉，不知廉耻！''她柔骨魅兔魂技全速运转，顺着 φ(2,0,0) 的不动点脉络缠向为首黑衣人，动作干脆利落。牢大紧随其后，催动 SVO 境界的蛮力，爆喝一声：

\begin{quote}
``滚远点！''
\end{quote}

拳风裹挟着规整的递归节律，砸向侧面两名小喽啰，拳风过境，连身旁的粗木都震得簌簌发抖。

蔡须坤手腕翻转，篮球带着 φ(1,ω,0) 的收敛光晕飞旋而出：

\begin{quote}
``别让他们跑了！''
\end{quote}

球面流转的光晕刚触碰到黑衣人，竟被那暗金色图腾微微弹开，诡异至极；卢本韦、诺当左右合围，牢牢挡住邪修退路；奥斯卡则退守后方，一边快速展开 FGH 辅助运算，一边高声提醒：

\begin{quote}
``大家稳住魂力脉络，别被邪力和那陌生纹路干扰！''
\end{quote}

f\_φ(1,0,φ(1,0))(3) 的双层嵌套光晕铺开，柔和的魂力涌入受伤的低阶扽魂师体内，稳住其濒临断裂的序数脉络。缠斗间，一名邪修突然捏碎一枚漆黑令牌，发出尖锐的破空声，却未召来援兵，反倒让林间那股古老波动彻底消散。不过半柱香功夫，几名黑衣人便被尽数制服，黑色序数波动缓缓消散，林间的阴冷雾气也渐渐散去，重新透出几分天光，唯独那抹银灰色残影与暗金图腾，成了未解的谜团。

获救的扽魂师挣扎着起身，对着众人拱手道谢，声音虚弱却带着惊惧：

\begin{quote}
``多谢诸位出手相救，我是扽魂殿信使，奉命传递紧急情报！''
\end{quote}

唐三上前搀扶，沉声追问：

\begin{quote}
``他们脖颈的暗金图腾、林间的古老波动，究竟是什么？''
\end{quote}

信使脸色煞白，压低声音道：

\begin{quote}
``不清楚\ldots 只听说两大集团最近勾结了神秘隐世势力，专门截杀持有高阶序数资质的魂师，总部那边也在查这股神秘力量，早已戒备森严！''
\end{quote}

众人对视一眼，心底的紧迫感更甚，牢大攥紧拳头：

\begin{quote}
``不管是什么来头，到了扽魂殿定要查个水落石出！''
\end{quote}

经此小插曲，小队众人不再是单纯奔赴正义殿堂，更带着对未知神秘势力的戒备，脚下步伐愈发急促。

前方，是帝国最顶尖的扽魂殿堂，是强者云集的舞台，更是正邪对决的战场。而这支从炼狱里走出的小队，即将带着满身锋芒，在这片全新的天地里，书写属于他们的传奇，迎战前所未有的挑战。

\begin{center}\rule{0.5\linewidth}{0.5pt}\end{center}

踏入扽魂殿正殿的刹那，所有人不自觉屏住呼吸，连呼吸都放轻了几分。

玄玉巨殿直插云霄，亿万道序数铭文自殿顶垂落，如星河倒悬，每一道纹路都烙印着万载岁月的威严。日光穿堂而过，铭文同时亮起金光，汇成一道横贯天际的神圣光桥，凡光所及，邪祟不侵，威压厚重如天地倾覆，却又中正平和，只压妄心，不欺善类。殿内无风自动，只有古老而庄严的气息在流淌，每一寸空气都铭刻着秩序、正义与守护，让人心生敬畏，不敢有半分亵渎。

大殿两侧，扽魂卫持杖肃立，身姿如岳，气息沉凝如万古磐石。玄色战甲覆身，肩刻守序徽记，目光如寒刃，不怒自威，周身散发出的气息沉稳得令人心悸，却始终守礼如仪，半步不越。这里没有张扬的威压，却有一种不容侵犯的凛然正气------扽魂殿立世之本，从来不是强权，而是守序；从来不是威慑，而是诛邪。

正殿之中，早已站满来自整片大陆的各路精英，鸦雀无声，落针可闻。

有人身披厚重铠甲，铠甲上刻着斩邪留痕的旧伤，一看便是常年在边境与恶势力厮杀的守护者；

有人身着素色长袍，气质温润，却眼神如炬，显然是隐世不出、专研正道序数的长者；

还有来自各大城邦的代表队、学院的精英师长、镇守一方的城邦守护者，他们来自不同地域、不同宗门，却在此刻保持着同样的肃穆垂立，连呼吸都刻意放轻，唯恐惊扰了殿中威严。

人群之中，白发苍苍的老魂师佝偻着身躯，却依旧挺直脊梁，目光虔诚地望着定序珠，不敢有半分歪斜；

刚成年的年轻修士紧握双拳，难掩激动，却依旧严守礼仪，垂首肃立，不敢左顾右盼；

更有不少低阶扽魂师，曾受过扽魂殿庇护，此刻能站在这座殿堂里，眼中满是感恩与敬畏，连大气都不敢喘。

整座大殿寂静如渊，唯有定序珠的微光缓缓流转，无形的威严笼罩全场，不威自怒，不言而信。

正殿中央，``守序诛邪''四个鎏金大字高悬，字迹苍劲如刀劈斧凿，金光内敛，却透着贯穿古今的意志，一眼望去，便让人心中邪念尽消、敬畏心生。大字之下，百丈高的序数神台巍然矗立，台心定序珠莹白温润，散发着能抚平一切紊乱、稳固一切脉络的神圣波动，乃是天地秩序所化，万邪辟易，万众敬仰。它不只是镇殿至宝，更是扽魂殿公正无私、不偏不倚、至高无上的威严象征。

少顷，三道身影缓步踏上神台。

脚步未落，殿内气息已然一凝。

为首者正是扽魂殿梅天狸长老，白袍金边，面容威严而慈和，周身没有半分凌厉气息，却仿佛与整个大殿的道义、秩序、威严融为一体。她目光扫过全场，温和却有千钧之力，所过之处，所有人不自觉垂首躬身，心悦诚服。

其他两位长老紧随其后，气息沉稳，神色肃穆，不言自威，仅仅是站立在神台之上，便让殿内威严再增三分。

大长老缓缓抬手，浑厚而中正的声音透过魂力传遍大殿每一个角落，不高亢、不张扬，却字字如钟鸣，叩击人心，带着天地秩序般的威严：

\begin{quote}
``今日，精英扽魂师登记大典，启。''
\end{quote}

一语落下，定序珠光芒大放，纯白圣光普照全殿，驱散一切阴邪、紊乱与躁动。殿内所有魂师同时躬身行礼，不是对长老，而是对扽魂殿至高无上的威严与正道致敬。

梅天狸长老声音继续回荡，庄严而清晰，带着不容置疑的威严：

\begin{quote}
``诸位，今日齐聚于此，不为论高低，不为较强弱。登记、定序、定级，皆不是为了名位，不是为了荣光，更不是为了权势。''
\end{quote}

他顿了顿，目光如炬，照亮殿中每一寸角落，声威浩荡：

\begin{quote}
``扽魂殿立世万载，一脉相承，唯守八字------守序，诛邪，护民，安道。''
\end{quote}

八字出口，定序珠光芒暴涨，序数铭文齐齐共鸣，金光冲天而起，整座大殿都在微微震颤，仿佛天地都在回应这亘古不变的誓言，威严直冲九霄。

\begin{quote}
``序，是天地之序，是修炼之序，是人心之序。我扽魂殿所守，不是强权之序，不是私利之序，而是天地正道之序，弱者可安、善者可存之序。我们定序数，不是划分尊卑，而是标定责任；不是彰显强弱，而是明确担当。''
\end{quote}

梅天狸的声音越发庄重，带着穿透岁月的厚重与威严：

\begin{quote}
``世间总有邪祟横行，有恶势力祸乱四方，有野心者篡改序数、残害同道、掠夺根基。他们无视秩序，践踏善意，鱼肉弱小，以强凌弱。而我扽魂殿，便是为此而立------邪祟一出，必诛；恶行一现，必惩；弱者一难，必救；正道一危，必赴。''
\end{quote}

殿内依旧寂静，可每一个字都落在众人心中，燃起一团滚烫的火焰，也更显扽魂殿威严如天。

\begin{quote}
``我们不拜强权，不奉私利，不趋炎附势。我们只站在正义一边，只站在弱小一边，只站在秩序一边。无论对方势力多大、根基多深、手段多诡，只要触犯道义、祸乱大陆，扽魂殿便会出手。不问来路，不问背景，不问强弱------唯善是护，唯恶是诛。''
\end{quote}

左边的一位长老上前一步，声音沉稳有力，补充道，声威震彻大殿：

\begin{quote}
``今日登记在册者，从此便身负扽魂殿印记。这印记不是荣耀，不是身份，不是特权。它是一道约束：不可恃强凌弱，不可滥伤无辜，不可篡改序数，不可为非作歹。它更是一道誓言：见恶必除，见危必救，见乱必平，见邪必斩。''
\end{quote}

右边的长老抬手一指殿外，声音肃然，威严尽显：

\begin{quote}
``大陆之上，仍有阴影潜行。帅气集团、美丽集团勾结隐世邪道，残害魂师、掠夺根基、扰乱序数，所行之处，民不安、师不宁、秩序崩坏。扽魂殿与之抗衡，不是为胜负，不是为霸权，而是为------让善良得以立足，让修炼得以纯粹，让弱小得以安稳，让正道得以长存。''
\end{quote}

梅天狸长老再次开口，声音升至最高点，却依旧中正庄严，威严如天规降临：

\begin{quote}
``今日登记，是入正道之门，是担守序之责。往后岁月，无论风雨险阻、无论强敌当前、无论诱惑万千，勿忘今日之敬畏，勿忘扽魂之初心：我等持序而立，为正义而生，为万民而战，为正道而死！守序，不折腰！诛邪，不退避！''
\end{quote}

\begin{quote}
``守序，不折腰！诛邪，不退避！''
\end{quote}

殿内所有扽魂师、所有精英、所有卫者，同时躬身齐喝。

声浪如雷，震动大殿，直冲云霄。定序珠光芒万丈，序数铭文漫天闪耀，天地共鸣，万邪俯首，扽魂殿的威严在这一刻达到顶峰。

阳光穿过鎏金大字，洒在众人身上，也洒在''守序诛邪''四个大字之上。

这一刻，没有等级高低，没有修为强弱，只有一群心怀正道的人，在至高殿堂之上，共同立下最庄严的承诺。

\begin{center}\rule{0.5\linewidth}{0.5pt}\end{center}

登记仪式正式开始。

没有喧嚣，没有浮华，只有肃穆、庄重与神圣。

来自大陆四方的魂师依次上前，每一位都在定序珠前驻足一瞬，躬身行礼，神色恭敬。有人闭目凝神，重温初心；有人微微躬身，向正道致意；有人目光坚定，暗下决心。没有人在意次序先后，所有人都在这一刻，感受着扽魂殿至高无上的威严，再次确认自己为何而修、为何而战。

扽魂殿的威严，从来不在殿堂之高，不在铭文之繁，不在卫者之强。而在那八个字里，千世万代，永不更改，威震八方，万邪臣服。

\chapter{第十九章：初赛惊鸿}\label{ux7b2cux5341ux4e5dux7ae0ux521dux8d5bux60caux9e3f}

扽魂殿精英登记大典落幕，高级扽魂师精英大赛的号角便紧随而至。这场赛事汇聚了整片大陆最顶尖的年轻扽魂师，既有Madore宗、Buchholz宗这样传承悠久的老牌宗门精锐，也有扽魂殿直属培养的强队，每一支参赛队伍都带着实打实的战绩与底气，唯有果糕七怪战队，因出身普通学院、此前从未涉足大型赛事，被所有参赛者轻飘飘地归为''陪跑队伍''。

大赛初赛采用单循环积分制，共设十组赛场，每组八支队伍，两两对决，胜者积3分，平局积1分，负者积0分，每组前四晋级复赛。果糕七怪战队被分在第七组，同组对手不仅有边陲城邦的烈风队，更有Madore宗、Buchholz宗两支老牌宗门小队------这两支队伍传承百年，弟子的序数修为普遍偏高，序链运用娴熟，是第七组的夺冠热门，也是果糕七怪战队晋级路上的最大阻碍。

初赛首日，第七组赛场外围人声鼎沸，围观者大多抱着''看老牌宗门碾压新人''的心态，目光多聚焦在Madore宗、Buchholz宗小队身上，对果糕七怪战队的首场对决几乎无人问津。果糕七怪战队的对手是来自边陲城邦的烈风队，这支队伍常年在边境与邪修厮杀，实战经验丰富，周身魂力运转沉稳，序链节律清晰，一开场便率先发难。

裁判一声令下，烈风队队长手持序刃，周身 \allowbreak{}\NovelInline{I01174}{\varphi(1,1,0)}\allowbreak{} 的光晕暴涨，带着凌厉的风属性魂力，直扑果糕七怪战队阵型，另外两名队员紧随其后，形成夹击之势。``小舞左翼牵制，诺当控场，卢本韦正面格挡！''唐三语速快得几乎咬字不清，眉心紫极魔瞳骤然迸出两道凌厉紫芒，\allowbreak{}\NovelInline{I01175}{\varphi(1,0,0,1)}\allowbreak{} 的四元序数光晕如潮水般铺展开来，蓝银草瞬间窜出，根根草藤泛着莹白序链光泽，如附骨之疽般缠绕向烈风队队长的序刃，同时指尖飞速掐诀，完成完整FGH展开：

\NovelDisplay{M00169R01}{%
f_{\varphi(1,0,0,1)}(3) = f_{\varphi(1,0,0,\varphi(1,0,0,\varphi(1,0,0)))}(3) = f_{\varphi(1,0,0,\varphi(1,0,0,\varphi(\omega,\varphi(\omega,\varphi(\omega,0)))))}(3)%
}

草藤层层交织，试图锁住对方的序链运转。

可烈风队早有防备，队长突然弃刃，掌心凝聚出一团漆黑如墨的暗金色魂力，那魂力翻涌间带着刺鼻的腥气，竟是沾染了邪修浊气的邪异序技------这种序技虽不纯正，却能精准撕裂正常序数脉络，刚一浮现，周遭空气便被搅得扭曲，赛场地面裂开细密的纹路。``不好！是邪异序技！''唐三心头巨震，嘶吼着提醒队友，可话音未落，蓝银草的序数脉络便被暗金色邪力狠狠侵蚀，翠绿的草藤瞬间变得焦黑干枯，发出''滋滋''的灼烧声，缠绕之力瞬间溃散，断裂的草藤带着火星飞溅。

烈风队队长眼中闪过狠厉，身形如鬼魅般欺近，掌心暗金色魂力暴涨，带着撕裂空气的锐响，狠狠拍向唐三胸口。就在这千钧一发之际，牢大已然爆喝出声，吼声震得赛场看台嗡嗡作响，她猛地掏出一瓶冰红茶，仰头一饮而尽，瓶身碎裂的瞬间，冰爽的能量顺着喉咙涌入体内，周身 \allowbreak{}\NovelInline{I01176}{\varphi(1@\omega)}\allowbreak{} 的高阶序数光晕骤然亮起，SVO境界的蛮力与序数力量、冰红茶的增幅完美融合，掌心递归节律飞速运转，完成完整FGH展开：

\NovelDisplay{M00170R01}{%
f_{\varphi(1@\omega)}(3) = f_{\varphi(1@3)}(3)%
}

严格遵循 \allowbreak{}\NovelInline{I01177}{f_{\varphi(1@\omega)}(n) = f_{\varphi(1@n)}(n)}\allowbreak{} 的规则，每一层迭代精准衔接，一拳砸出，拳风裹挟着狂暴的序链力量与冰爽劲气，狠狠撞向烈风队队长侧翼，``嘭''的一声巨响，气浪席卷全场，尘土飞扬，烈风队队长被震得连连后退，嘴角渗出鲜血。

可牢大急于救主，序链迭代节奏出现细微偏差，FGH展开出现卡顿，\allowbreak{}\NovelInline{I01178}{f_{\varphi(1@\omega)}(3)}\allowbreak{} 未能完全收敛，狂暴的魂力瞬间反噬自身，她闷哼一声，脚步踉跄着后退数步，胸口剧烈起伏，周身 \allowbreak{}\NovelInline{I01179}{\varphi(1@\omega)}\allowbreak{} 的光晕变得忽明忽暗。``牢大，稳住！''蔡须坤厉声提醒，周身同样泛起 \allowbreak{}\NovelInline{I01180}{\varphi(1@\omega)}\allowbreak{} 的序数光晕，手中篮球飞速旋转，球面凝聚起尖锐的收敛光晕，精准击中烈风队另一名队员的魂力节点，``咔嚓''一声脆响，对方序链瞬间崩断，惨叫着倒飞出去。此时诺当身形一闪，双手催动磁暴环，淡蓝色的磁暴光晕瞬间扩散，精准笼罩住烈风队剩余队员，``嗡''的一声低鸣，磁暴之力顺着对方序链钻进经脉，烈风队队员瞬间陷入眩晕，身形踉跄，序链运转彻底紊乱。

赛场外围瞬间炸开了锅，围观者的惊呼声此起彼伏，原本轻蔑的目光尽数变成震惊：``这新人战队居然能接住烈风队的暗招？''``那 \allowbreak{}\NovelInline{I01181}{\varphi(1@\omega)}\allowbreak{} 的运用，比有些老牌宗门还娴熟！''卢本韦抓住时机，周身土黄色光晕暴涨，催动母猪扽魂全力，身形如奔雷般冲向眩晕的烈风队队员，一头狠狠撞向对方小腹，``嘭''的一声，将对方撞得倒飞出去，重重摔在地上。她抬手示意队友，高声嘶吼：``目标失衡，让队友一鼓作气打败敌方！''蔡须坤没有停歇，手腕翻转，篮球带着收敛光晕再次飞射而出，精准缠绕住烈风队剩余队员的序链，唐三则趁机催动四元序数，蓝银草重新焕发生机，根根草藤泛着莹白光泽，死死锁住烈风队队长的魂力脉络，不让其有喘息之机。最终，烈风队队长无力支撑，抬手示意认输，果糕七怪战队拿下首场胜利。

首场胜利并未让果糕七怪战队松懈，她们深知，真正的考验还在后面------Madore宗、Buchholz宗两支老牌宗门小队，才是她们晋级路上的最大障碍。这两支队伍传承百年，弟子的序数修为普遍偏高，Madore宗擅长序链压制，其队长更是达到SVO高阶，周身 \allowbreak{}\NovelInline{I01182}{\varphi(1,\omega,0)}\allowbreak{} 的极限收敛光晕凝练，FGH展开娴熟，能精准干扰对手的序链运转；Buchholz宗则擅长防御反击，序链闭环严密，凭借 \allowbreak{}\NovelInline{I01183}{\varphi(1,0,0)}\allowbreak{} 的韧性，能死死守住阵型，等待反击时机。事实上SVO不是别的，就是 \allowbreak{}\NovelInline{I01184}{\varphi(1@\omega)}\allowbreak{}，\allowbreak{}\NovelInline{I01185}{\varphi(1@0) = \varphi(0)}\allowbreak{}，\allowbreak{}\NovelInline{I01186}{\varphi(1@1) = \varphi(1,0)}\allowbreak{}，\allowbreak{}\NovelInline{I01187}{\varphi(1@2) = \varphi(1,0,0)}\allowbreak{}，实际上，可以理解为 \allowbreak{}\NovelInline{I01188}{1@n}\allowbreak{} 是1在n号位置。

果糕七怪战队的第二场对决，便遭遇了Madore宗小队。Madore宗弟子身着统一的深紫色宗门服饰，衣摆绣着古朴的序链纹路，周身序数光晕凝实如玉，甫一出场，脚下的赛场地面便被序链力量压出细密的裂痕，强烈的压制力如无形的巨石，沉甸甸地压在全场观众心头。

队长站在阵型前方，\allowbreak{}\NovelInline{I01189}{\varphi(1,\omega,0)}\allowbreak{} 的极限收敛光晕如实质般铺展开来，紫金色的光芒刺得人睁不开眼，他抬眼扫过果糕七怪战队，语气里的傲慢几乎要溢出来：``新人战队，识相点直接认输，免得被我的序链撕碎，连体面都留不下。''话音未落，Madore宗两名弟子身形如离弦之箭，序链运转间发出''嗡嗡''的低沉震颤，FGH展开流畅得毫无破绽，魂力裹挟着凌厉的序链锐芒，如两道闪电般直扑果糕七怪战队的侧翼，所过之处，空气都被撕裂，留下两道扭曲的气痕。

``卢本韦，守住左翼！蔡须坤，校准收敛节律！''唐三的声音冷静得没有一丝波澜，眉心紫极魔瞳骤然迸出两道紫金色锐芒，穿透层层序链光晕，精准捕捉到Madore宗序链的致命破绽------其FGH展开虽娴熟，却在 \allowbreak{}\NovelInline{I01190}{\varphi(1,\omega,0)}\allowbreak{} 的迭代衔接处，存在一丝微不可查的断层，那是高阶序数运转时最容易被突破的弱点。话音未落，蓝银草如暴雨般窜出，根根草藤泛着莹白的序链光泽，如钢鞭般死死缠绕住Madore宗弟子的序链，草藤与序链碰撞的瞬间，发出''滋滋''的刺耳声响，火星四溅。

奥斯卡周身柔光暴涨，多层嵌套FGH展开瞬间拉满，\allowbreak{}\NovelInline{I01191}{f_{\varphi(1,0,\varphi(1,0))}(3)}\allowbreak{} 的柔和光晕如潮水般笼罩全队，不仅快速补充众人魂力，更在每个人周身形成一层序链防护，强化序链韧性。牢大借着辅助光晕的加持，再次掏出冰红茶一饮而尽，周身 \allowbreak{}\NovelInline{I01192}{\varphi(1@\omega)}\allowbreak{} 的光晕骤然变得炽烈，肌肉贲张间，骨骼发出''咔咔''的脆响，掌心递归节律飞速运转，\allowbreak{}\NovelInline{I01193}{f_{\varphi(1@\omega)}(3) = f_{\varphi(1@3)}(3)}\allowbreak{} 的FGH展开完美闭环，蛮力顺着序链层层叠加，一拳砸出，拳风裹挟着狂暴的序链力量，带着''轰隆''的巨响，狠狠撞向Madore宗队长的序链节点。

这场对决打得天昏地暗，激烈程度远超首场，序数波动碰撞的轰鸣声震耳欲聋，几乎要掀翻赛场的顶棚。Madore宗的序链压制如潮水般层层递进，紫金色的序链光泽与果糕七怪战队的莹白、炽烈光晕交织碰撞，赛场之上尘土漫天，碎石飞溅，原本平整的赛场地面被硬生生砸出一个个深坑。

Madore宗队长见状，眼中闪过狠厉，猛地催动全部魂力，\allowbreak{}\NovelInline{I01194}{\varphi(1,\omega^\omega,0)}\allowbreak{} 的极限收敛光晕暴涨数倍，序链迭代速度瞬间加快，试图用高阶序数强行撕裂果糕七怪战队的防御，他掌心凝聚出一团紫金色的序链光球，光球运转间，FGH展开的节律如惊雷般炸响，直扑唐三而去。就在这千钧一发之际，诺当再次催动磁暴环，淡蓝色磁暴光晕瞬间扩大，精准命中紫金色序链光球，同时笼罩住Madore宗队长，磁暴之力瞬间干扰其序链运转，让他陷入短暂眩晕，光球运转节奏大乱。

蔡须坤趁机身形一闪，手中篮球带着 \allowbreak{}\NovelInline{I01195}{\varphi(1@\omega)}\allowbreak{} 的收敛光晕，如闪电般精准击中序链光球的核心节点，``嘭''的一声巨响，光球瞬间炸裂，狂暴的序链力量四散飞溅，蔡须坤被气浪狠狠掀飞，重重摔在地上，嘴角涌出一大口鲜血，周身 \allowbreak{}\NovelInline{I01196}{\varphi(1@\omega)}\allowbreak{} 的光晕瞬间黯淡了几分。

另一侧，卢本韦催动母猪扽魂的全力，周身土黄色光晕厚重如盾，死死挡住Madore宗弟子的侧翼攻击，序链碰撞的冲击力让她浑身震颤，手臂青筋暴起。见对方两名弟子攻势不减，卢本韦眼神一厉，身形猛地前冲，借着蛮力一头撞向对方，``嘭嘭''两声，将两名Madore宗弟子撞得身形失衡、连连后退，她趁机高声嘶吼：``目标失衡，让队友一鼓作气打败敌方！''吼声未落，牢大借着冰红茶的残余增幅，再次催动 \allowbreak{}\NovelInline{I01197}{\varphi(1@\omega)}\allowbreak{} 的高阶力量，身形如铁塔般冲上前，一拳接一拳砸向眩晕未消的Madore宗队长，每一拳都裹挟着狂暴的序链力量，赛场地面被砸得剧烈震颤。

``蔡须坤！''``卢本韦！''众人齐声嘶吼，眼底满是悲愤，斗志被彻底点燃。唐三抓住Madore宗队长眩晕未消、序链紊乱的破绽，眉心紫极魔瞳光芒更盛，四元序数全力运转，蓝银草疯狂生长，根根草藤穿透漫天尘土，如灵蛇般死死锁住对方队长的序链，序链力量如利刃般钻进对方经脉，让其魂力运转瞬间停滞。牢大趁着对方失神之际，再次催动 \allowbreak{}\NovelInline{I01198}{\varphi(1@\omega)}\allowbreak{} 的高阶力量，FGH展开完美无缺，肘击Madore宗队长的胸口，``咔嚓''一声脆响，对方胸骨凹陷，喷出一大口紫金色的魂力精血，身体如断线的风筝般倒飞出去，重重撞在赛场外围的屏障上，屏障瞬间裂开蛛网般的裂痕。

他挣扎着撑起身体，没有丝毫怨怼，反而抬眼望向果糕七怪战队，眼底的傲慢彻底褪去，取而代之的是浓烈的敬佩，嘴角溢出鲜血，却缓缓抬手，对着果糕七怪战队深深一揖------这是老牌宗门对强者的认可，是技不如人的坦荡，更是对对手实力的由衷折服。剩余的Madore宗弟子见状，虽满心不甘，却也纷纷停下反抗，对着果糕七怪战队躬身示意后，才搀扶着队长缓缓退下，没有丝毫耍赖推诿。

果糕七怪战队终于拿下这场艰难的胜利，蔡须坤挣扎着站起身，卢本韦也拄着手臂勉强站稳，众人周身或多或少都带着伤势，却个个目光坚定，周身的序链光晕虽有些黯淡，却依旧带着不屈的锋芒。而此时，Buchholz宗小队正在不远处的赛场作战，她们全程紧盯着这边的战况，看到Madore宗惨败却依旧保持体面，更看到果糕七怪战队的强悍，神色愈发凝重，眼神中满是警惕------这个不起眼的新人战队，已然成为她们晋级路上最可怕的威胁。

两场连胜过后，果糕七怪战队彻底摆脱了''陪跑队伍''的标签，名气渐渐传开。围观者中，有扽魂殿的长老驻足观看，对着身边的人赞叹：``Madore宗、Buchholz宗的序链运用已是顶尖，却没想到被一群新人逼到这般地步，这果糕七怪战队，果然有过人之处。''

初赛第三场，果糕七怪战队遭遇了Buchholz宗小队。Buchholz宗弟子周身序链光晕厚重如石，防御闭环严密，序链迭代沉稳，FGH展开毫无破绽，一开场便摆出死守阵型，试图消耗果糕七怪战队的魂力。唐三改变战术，让卢本韦继续正面牵制，诺当释放磁暴环干扰对方序链节律，蔡须坤则用篮球精准打击对方防御薄弱节点，牢大则伺机寻找反击时机。

奥斯卡的辅助始终在线，多层嵌套FGH展开不断，为队友补充魂力，强化序链力量。激战半个时辰后，Buchholz宗的防御终于出现破绽------其序链迭代出现卡顿，诺当抓住时机，催动磁暴环全力爆发，淡蓝色磁暴光晕笼罩住Buchholz宗全队，让对方陷入短暂眩晕，序链闭环瞬间松动。卢本韦趁机催动母猪扽魂，身形奔冲而出，一头撞向Buchholz宗队长，将其撞得身形失衡，她高声嘶吼：``目标失衡！''唐三抓住机会，催动四元序数全力一击，蓝银草死死锁住对方队长的序链，牢大则掏出冰红茶一饮而尽，借着增幅一拳砸破对方的防御闭环，蔡须坤同步发动收敛攻击，彻底打乱对方的序链运转。最终，Buchholz宗小队无力支撑，抬手认输。

初赛落幕时，果糕七怪战队以五战全胜、零失误的战绩，稳居第七组第一，顺利晋级复赛。Madore宗、Buchholz宗虽未能阻止果糕七怪战队的步伐，却也让她们看清了自身的不足------相较于老牌宗门的序链底蕴，她们的FGH展开熟练度仍有差距，且在高阶序数的持久作战上，还需进一步打磨。

赛场出口，夕阳洒在果糕七怪战队众人身上，序数光晕与余晖交相辉映。唐三望着远处Madore宗、Buchholz宗小队离去的方向，语气坚定：``复赛的对手会更强，我们还要继续打磨序链，不能有丝毫松懈。''众人纷纷点头，眼底满是锋芒------她们从默默无闻的新人，一步步击败老牌宗门精锐，用实力证明，炼狱特训的磨砺、生死与共的羁绊，远比天赋更能成就强大。

而不远处的看台上，Madore宗、Buchholz宗的队长并肩而立，望着果糕七怪战队的身影，神色凝重却不见怨毒。Madore宗队长捂着胸口，指尖还残留着紫金色的魂力痕迹，语气里没有不甘，只有实打实的敬佩：``我输得口服心服，这果糕七怪战队，每一个人的序链运用都远超同龄修士，尤其是她们的羁绊与韧性，是我们Madore宗弟子所欠缺的。''他抬眼望向赛场中央，眼底燃起炽热的斗志，攥紧了拳头，声音沉而坚定：``今日之败，我记下了。

复赛也好，决赛也罢，我会带着Madore宗的弟子，日夜打磨序链，精进FGH展开，定要在决战之上，与她们堂堂正正再决高下，用实力赢回属于Madore宗的荣耀！``Buchholz宗队长微微颔首，目光紧紧锁定唐三等人：''她们的序链运用虽还有瑕疵，但羁绊极强，实力更是不容小觑。复赛再遇，我们必须全力以赴，既要守住Buchholz宗的颜面，也不能让Madore宗的蛰伏白费。''

初赛的喧嚣渐渐散去，复赛的号角即将吹响。果糕七怪战队的传奇，才刚刚拉开序幕，而Madore宗、Buchholz宗的蛰伏，也让这场精英大赛，多了几分未知与期待。

\chapter{第二十章：非递归序数与Ω的坍缩}\label{ux7b2cux4e8cux5341ux7ae0ux975eux9012ux5f52ux5e8fux6570ux4e0eux3c9ux7684ux574dux7f29}

初赛结束后，玉小刚见众人对战时仍显吃力------即便全员已突破 \allowbreak{}\NovelInline{I01199}{\varphi(2,0,0)}\allowbreak{}，却依旧无法对老牌宗门形成压制。为此，玉小刚决定向众人讲解非递归序数。

讲堂内，气氛凝重。玉小刚站在黑板前，粉笔在黑板上敲击出清脆的声响。

``你们目前掌握的序数，都有清晰的基本列，对吗？''

诺当点头答道：``是的。比如 \allowbreak{}\NovelInline{I01200}{\varphi(1,2,1) = \sup\{\varphi(1,2,0)+1, \varphi(1,2,\varphi(1,2,0)+1), \dots\}}\allowbreak{}。''

``事实上，所有递归序数，都是通过序数的后继、乘法、幂运算不断作用得到的。''玉小刚的声音沉稳有力，``比如你们修炼中接触的 \allowbreak{}\NovelInline{I01201}{\varepsilon_0 = \varphi(1,0)}\allowbreak{}，是 \allowbreak{}\NovelInline{I01202}{\omega^{\omega^{\dots}}}\allowbreak{} 的极限；\allowbreak{}\NovelInline{I01203}{\varepsilon_1}\allowbreak{} 是 \allowbreak{}\NovelInline{I01204}{\varepsilon_0^{\varepsilon_0^{\dots}}}\allowbreak{} 的极限。哪怕是你们刚突破的 \allowbreak{}\NovelInline{I01205}{\varphi(2,0,0)}\allowbreak{}，本质上还是这几种运算层层叠加的结果。''

他顿了顿，加重语气问道：``那么你们想一想，是否存在无法用这种方式构造的序数？我给你们举个具体的例子，就像你们修炼时，用魂力催动魂技、叠加魂环效果，能不断提升战力，但总有一个`天花板'，无论怎么叠加基础魂技，都达不到神级战力------\textbf{非递归序数，就是递归运算的`战力天花板'}。''

他进一步解释：``咱们先明确一个最基础的非递归序数例子，就是丘奇-克莱尼序数，记作 \allowbreak{}\NovelInline{I01206}{\omega_1^{CK}}\allowbreak{}。它是第一个非递归序数，也是所有递归序数的上确界，就像 \allowbreak{}\NovelInline{I01207}{\omega}\allowbreak{} 是所有自然数的上确界一样------无论你们把递归运算做多少次，哪怕是叠无穷多层指数塔，得到的序数都小于 \allowbreak{}\NovelInline{I01208}{\omega_1^{CK}}\allowbreak{}。它无法通过任何小于它自身次数的递归运算得到。''

``再通俗点说，你们之前练的 \allowbreak{}\NovelInline{I01209}{\varphi(1,0)}\allowbreak{}、\allowbreak{}\NovelInline{I01210}{\varphi(2,0)}\allowbreak{}、\allowbreak{}\NovelInline{I01211}{\varphi(2,0,0)}\allowbreak{}，就像你们的魂帝、魂斗罗境界，靠层层修炼；而 \allowbreak{}\NovelInline{I01212}{\omega_1^{CK}}\allowbreak{} 这样的非递归序数，不是靠`叠加基础修炼'就能触及的，必须跳出递归的框架------咱们今天说的 \allowbreak{}\NovelInline{I01213}{\Omega}\allowbreak{}，也是咱们后续修炼要突破的核心门槛。''

奥斯卡有些难以置信：``不会吧？我感觉这样运算下去序数只会越来越大，无论多大的序数都能表示。''

牢大却若有所思地摇了摇头：``未必。\allowbreak{}\NovelInline{I01214}{\omega}\allowbreak{} 大于一切自然数，自然数再大也超不过 \allowbreak{}\NovelInline{I01215}{\omega}\allowbreak{}。或许真的存在这样的序数，就像玉老师说的。''

弗兰德此时也开口了，他目光如炬：``牢大说得对。这便是非递归序数，记作 \allowbreak{}\NovelInline{I01216}{\Omega}\allowbreak{}。\allowbreak{}\NovelInline{I01217}{\Omega}\allowbreak{} 无法通过小于 \allowbreak{}\NovelInline{I01218}{\Omega}\allowbreak{} 次的递归运算得到，就像 \allowbreak{}\NovelInline{I01219}{\omega_1^{CK}}\allowbreak{}，它囊括了所有递归序数，却永远无法被递归运算构造出来。''

``什么是递归运算？''唐三敏锐地抓住了关键。

``简单说，就是以递归方式作用于序数。''弗兰德解释道，``比如将 \allowbreak{}\NovelInline{I01220}{\zeta_0+1}\allowbreak{} 代入 \allowbreak{}\NovelInline{I01221}{\alpha \mapsto \varepsilon_\alpha}\allowbreak{} 的运算，最终会得到 \allowbreak{}\NovelInline{I01222}{\varepsilon_{\zeta_0+1}}\allowbreak{}，这就像你们递归使用魂技，用前一次的魂技效果催动下一次，本质还是在`递归框架'内。''

小舞眼中闪过一丝明悟：``我懂了，就是递归层面的映射。''

``你们暂时不必深究，但要记住：\allowbreak{}\NovelInline{I01223}{\Omega}\allowbreak{} 是一切递归序数无法逾越的高峰。''弗兰德总结道，``那有人的修为能达到 \allowbreak{}\NovelInline{I01224}{\Omega}\allowbreak{} 吗？我也不清楚。古老传说中，Madore与Buchholz两兄弟创立MOCF与BOCF，最终突破平凡递归境界，抵达神级水准。他们所创的Madore宗与Buchholz宗，也是大陆上首屈一指的宗门。''

``那我们学习哪一种OCF？''蔡徐坤问道。

``MOCF。下面我来讲解。''弗兰德转身在黑板上写下巨大的 \allowbreak{}\NovelInline{I01225}{\psi_\Omega}\allowbreak{}。

``首先认识 \allowbreak{}\NovelInline{I01226}{\psi_\Omega}\allowbreak{} 函数，其作用是对序数进行坍缩------这也是OCF被称为序数坍缩函数（Ordinal Collapsing Function）的缘由。Madore所创OCF规则如下：\allowbreak{}\NovelInline{I01227}{\psi_\Omega(n) = \varepsilon_n, n < \zeta_0}\allowbreak{}。''

蔡徐坤反应极快：``那 \allowbreak{}\NovelInline{I01228}{\psi_\Omega(0) = \varphi(1,0)}\allowbreak{}。''

``正确。并且 \allowbreak{}\NovelInline{I01229}{\psi_\Omega(\psi_\Omega(0)) = \varphi(1, \varphi(1,0))}\allowbreak{}。''

唐三紧接着问：``那 \allowbreak{}\NovelInline{I01230}{\psi_\Omega(\Omega)}\allowbreak{} 是什么？''

玉小刚接话道：``\allowbreak{}\NovelInline{I01231}{\psi_\Omega(\Omega) = \sup\{\psi_\Omega(0), \psi_\Omega(\psi_\Omega(0)), \psi_\Omega(\psi_\Omega(\psi_\Omega(0))), \dots\}}\allowbreak{}。''

牢大目光一闪：``这不就是 \allowbreak{}\NovelInline{I01232}{\varphi(2,0)}\allowbreak{} 吗？''

``没错。''弗兰德点头，``\allowbreak{}\NovelInline{I01233}{\psi_\Omega(\Omega)}\allowbreak{} 是映射 \allowbreak{}\NovelInline{I01234}{\alpha \mapsto \psi_\Omega(\alpha)}\allowbreak{} 的不动点，即满足 \allowbreak{}\NovelInline{I01235}{\alpha < \psi_\Omega(\Omega)}\allowbreak{}，且是该映射的最小不动点。''

卢本伟追问：``那 \allowbreak{}\NovelInline{I01236}{\psi_\Omega(\Omega+1)}\allowbreak{} 如何定义？''

``\allowbreak{}\NovelInline{I01237}{\psi_\Omega}\allowbreak{} 内层加1，取指数塔极限，也就是 \allowbreak{}\NovelInline{I01238}{\varphi(1, \psi_\Omega(\Omega)+1)}\allowbreak{}。''

``也就是 \allowbreak{}\NovelInline{I01239}{\varepsilon_{\zeta_0+1}}\allowbreak{}？''诺当确认道。

``对。''

``那 \allowbreak{}\NovelInline{I01240}{\psi_\Omega(\Omega \cdot 2) = \varepsilon_{\zeta_0 \cdot 2}}\allowbreak{} 吗？''

玉小刚立刻纠正：``不对。\allowbreak{}\NovelInline{I01241}{\Omega}\allowbreak{} 在此处的作用是总结不动点。\allowbreak{}\NovelInline{I01242}{\psi_\Omega(\Omega \cdot 2)}\allowbreak{} 是映射 \allowbreak{}\NovelInline{I01243}{\alpha \mapsto \psi_\Omega(\Omega+\alpha)}\allowbreak{} 的最小不动点。''

牢大迅速在脑海中构建模型：``\allowbreak{}\NovelInline{I01244}{\psi_\Omega(\Omega \cdot 2) = \sup\{\psi_\Omega(\Omega), \psi_\Omega(\Omega+\psi_\Omega(\Omega)), \psi_\Omega(\Omega+\psi_\Omega(\Omega+\psi_\Omega(\Omega))), \dots\}}\allowbreak{}，对吗？''

``对！你们觉得它对应哪个序数？''

蔡徐坤脱口而出：``\allowbreak{}\NovelInline{I01245}{\zeta_1}\allowbreak{}。''

``完全正确。明白 \allowbreak{}\NovelInline{I01246}{\psi_\Omega(\Omega+\psi_\Omega(\Omega))}\allowbreak{} 为何不会卡在不动点，而不是 \allowbreak{}\NovelInline{I01247}{\varepsilon_{\zeta_0}}\allowbreak{} 吗？''

小舞条理清晰地分析道：``因为 \allowbreak{}\NovelInline{I01248}{\psi_\Omega(\Omega+1) = \varepsilon_{\zeta_0+1}}\allowbreak{}，于是 \allowbreak{}\NovelInline{I01249}{\psi_\Omega(\Omega+\psi_\Omega(\Omega)) = \varepsilon_{\zeta_0 \cdot 2}}\allowbreak{}，从而跳出不动点。\allowbreak{}\NovelInline{I01250}{\psi_\Omega(\Omega+\psi_\Omega(\Omega+\psi_\Omega(\Omega)))) = \varphi(1, \varphi(2,0)+\varphi(1, \varphi(2,0) \cdot 2))}\allowbreak{}。根据序数吸收率：若 \allowbreak{}\NovelInline{I01251}{\beta}\allowbreak{} 为极限序数，则 \allowbreak{}\NovelInline{I01252}{\alpha+\beta = \alpha+\beta[n]}\allowbreak{}。于是 \NovelDisplay{I01253}{\varphi(2,0)+\varphi(1, \varphi(2,0) \cdot 2) = \sup\{\varphi(2,0)+\varphi(1, \varphi(2,0) \cdot \varphi(2,0)[n]) | n<\omega\}}。而 \allowbreak{}\NovelInline{I01254}{\varphi(2,0)[n] = \varphi(1, \varphi(2,0))[n-1]}\allowbreak{}，\allowbreak{}\NovelInline{I01255}{\varphi(2,0)[0] = \varphi(1,0)}\allowbreak{}。又 \allowbreak{}\NovelInline{I01256}{\varphi(1, \varphi(2,0) \cdot \varphi(2,0)[n]) + \varphi(2,0)}\allowbreak{} 介于 \allowbreak{}\NovelInline{I01257}{\varphi(1, \varphi(2,0) \cdot \varphi(2,0)[n])}\allowbreak{} 与 \allowbreak{}\NovelInline{I01258}{\varphi(1, \varphi(2,0) \cdot \varphi(2,0)[n+1])}\allowbreak{} 之间。由夹逼性，两者极限相同，故 \allowbreak{}\NovelInline{I01259}{\varphi(1, \varphi(2,0)+\varphi(1, \varphi(2,0) \cdot 2)) = \varphi(1, \varphi(1, \varphi(2,0) \cdot 2))}\allowbreak{}，最终收敛到 \allowbreak{}\NovelInline{I01260}{\zeta_1}\allowbreak{}。''

玉小刚赞许道：``解释得非常好。那 \allowbreak{}\NovelInline{I01261}{\psi_\Omega(\Omega \cdot 3)}\allowbreak{} 是什么？''

奥斯卡回答：``自然是映射 \allowbreak{}\NovelInline{I01262}{\alpha \mapsto \psi_\Omega(\Omega \cdot 2 + \alpha)}\allowbreak{} 的不动点。''

``不错，它就是 \allowbreak{}\NovelInline{I01263}{\zeta_2}\allowbreak{}。''

诺当总结道：``依此类推，\allowbreak{}\NovelInline{I01264}{\psi_\Omega(\Omega \cdot n) = \zeta_{n-1} = \varphi(2, n-1)}\allowbreak{}，\allowbreak{}\NovelInline{I01265}{n=\omega}\allowbreak{} 时亦然。''

弗兰德补充：``对 \allowbreak{}\NovelInline{I01266}{\varphi(3,0)}\allowbreak{} 以下均成立：\allowbreak{}\NovelInline{I01267}{\psi_\Omega(\Omega \cdot \omega) = \zeta_\omega}\allowbreak{}，\allowbreak{}\NovelInline{I01268}{\psi_\Omega(\Omega \cdot (\omega+1)) = \zeta_{\omega+1}}\allowbreak{}。''

唐三举一反三：``\allowbreak{}\NovelInline{I01269}{\psi_\Omega(\Omega^2) = \psi_\Omega(\Omega \cdot \Omega)}\allowbreak{}，是映射 \allowbreak{}\NovelInline{I01270}{\alpha \mapsto \psi_\Omega(\Omega \cdot \alpha)}\allowbreak{} 的不动点，对吗？''

玉小刚点头：``看来你们已初步理解OCF中 \allowbreak{}\NovelInline{I01271}{\Omega}\allowbreak{} 的作用------总结不动点。集合论层面的严格定义，等你们能熟练展开OCF后再讲。''

小舞猜测：``\allowbreak{}\NovelInline{I01272}{\psi_\Omega(\Omega^2) = \varphi(3,0)}\allowbreak{}，我猜得对吗？''

``对，正是如此。\allowbreak{}\NovelInline{I01273}{\psi_\Omega(\Omega \cdot \psi_\Omega(\Omega)) = \zeta_{\zeta_0}}\allowbreak{}，\allowbreak{}\NovelInline{I01274}{\psi_\Omega(\Omega \cdot \psi_\Omega(\Omega+1)) = \zeta_{\varepsilon_{\zeta_0+1}}}\allowbreak{}。''

诺当继续推导：``\allowbreak{}\NovelInline{I01275}{\psi_\Omega(\Omega^2+\Omega)}\allowbreak{} 是映射 \allowbreak{}\NovelInline{I01276}{\alpha \mapsto \psi_\Omega(\Omega^2+\alpha)}\allowbreak{} 的不动点，那就是 \allowbreak{}\NovelInline{I01277}{\varphi(2, \varphi(3,0)+1)}\allowbreak{}，对吧？''

弗兰德引导道：``再想想 \allowbreak{}\NovelInline{I01278}{\psi_\Omega(\Omega^2 \cdot 2)}\allowbreak{}。''

牢大自信地说道：``把式子补到以 \allowbreak{}\NovelInline{I01279}{\Omega}\allowbreak{} 结尾，写成 \allowbreak{}\NovelInline{I01280}{\psi_\Omega(\Omega^2 + \Omega \cdot \Omega)}\allowbreak{}，再取映射 \allowbreak{}\NovelInline{I01281}{\alpha \mapsto \psi_\Omega(\Omega^2 + \Omega \cdot \alpha)}\allowbreak{} 的不动点，也就是 \allowbreak{}\NovelInline{I01282}{\varphi(3,1)}\allowbreak{}。''

蔡徐坤发现了规律：``我发现规律了：\allowbreak{}\NovelInline{I01283}{\psi_\Omega()}\allowbreak{} 内部乘2，就是下一个同性质序数。\allowbreak{}\NovelInline{I01284}{\varphi(3,1)}\allowbreak{} 就是下一个具有 \allowbreak{}\NovelInline{I01285}{\varphi(3,0)}\allowbreak{} 性质的序数。''

玉小刚肯定道：``牢大说得很对，确实如此。''

随着讲解的深入，众人的思维愈发敏捷。

唐三：``那 \allowbreak{}\NovelInline{I01286}{\varphi(4,0) = \psi_\Omega(\Omega^3)}\allowbreak{}。'' 弗兰德：``是的。进而 \allowbreak{}\NovelInline{I01287}{\psi_\Omega(\Omega^\omega) = \varphi(\omega, 0)}\allowbreak{}。''

奥斯卡提出了一个关键问题：``老师，我有个问题不知该不该问。'' 弗兰德：``直说。'' 奥斯卡：``\allowbreak{}\NovelInline{I01288}{\psi_\Omega(\Omega+\varphi(3,1))}\allowbreak{} 有定义吗？\allowbreak{}\NovelInline{I01289}{\psi_\Omega(\varphi(3,0))}\allowbreak{} 又是什么？''

弗兰德严肃地解释：``现阶段可视为非标准序列。\allowbreak{}\NovelInline{I01290}{\psi_\Omega(\Omega+\varphi(3,1)) = \psi_\Omega(\Omega \cdot 2) = \varphi(3,0)}\allowbreak{}。非标准序列中，超出部分一律视作 \allowbreak{}\NovelInline{I01291}{\Omega}\allowbreak{}：只要 \allowbreak{}\NovelInline{I01292}{\varphi(3,1) > \psi_\Omega(\Omega \cdot 2) = \varphi(3,0)}\allowbreak{}，就把 \allowbreak{}\NovelInline{I01293}{\varphi(3,1)}\allowbreak{} 看作 \allowbreak{}\NovelInline{I01294}{\Omega}\allowbreak{}，结果恒为 \allowbreak{}\NovelInline{I01295}{\varphi(3,0)}\allowbreak{}。哪怕是 \allowbreak{}\NovelInline{I01296}{\varphi(3,0)}\allowbreak{}、\allowbreak{}\NovelInline{I01297}{\varphi(3,0,0)}\allowbreak{}，结果都一样，永远超不过 \allowbreak{}\NovelInline{I01298}{\psi_\Omega(\Omega \cdot 2)}\allowbreak{}。\allowbreak{}\NovelInline{I01299}{\psi_\Omega(\varphi(3,0)) = \psi_\Omega(\varphi(2,0)) = \psi_\Omega(\Omega) = \varphi(2,0)}\allowbreak{}。明白吗？以后再按定义严格推导非标准式。''

玉小刚继续出题测试：``\allowbreak{}\NovelInline{I01300}{\psi_\Omega((\Omega^\omega) \cdot 2) = \varphi(\omega, 1)}\allowbreak{}，没问题吧？''

小舞迅速作答：``那 \allowbreak{}\NovelInline{I01301}{\psi_\Omega(\Omega^{\omega+1}) = \psi_\Omega(\Omega^\omega + \Omega)}\allowbreak{}，即映射 \allowbreak{}\NovelInline{I01302}{\alpha \mapsto \psi_\Omega(\Omega^\omega+\alpha)}\allowbreak{} 的不动点。''

弗兰德点头：``正确。猜猜 \allowbreak{}\NovelInline{I01303}{\psi_\Omega(\Omega^\Omega)}\allowbreak{} 是多少。''

唐三脱口而出：``\allowbreak{}\NovelInline{I01304}{\Gamma_0}\allowbreak{}，也就是 \allowbreak{}\NovelInline{I01305}{\varphi(1,0,0)}\allowbreak{}。''

弗兰德欣慰道：``看来你们已领悟OCF的基本模式。''

牢大紧接着说：``那 \allowbreak{}\NovelInline{I01306}{\psi_\Omega((\Omega^\Omega) \cdot 2)}\allowbreak{} 就是 \allowbreak{}\NovelInline{I01307}{\varphi(1,0,1)}\allowbreak{}。因为 \allowbreak{}\NovelInline{I01308}{\psi_\Omega(\Omega^\Omega + \Omega^n) = \varphi(n, \varphi(1,0,0)+1)}\allowbreak{}，逻辑很清晰。''

玉小刚进一步推导：``\allowbreak{}\NovelInline{I01309}{\psi_\Omega((\Omega^\Omega) \cdot \Omega) = \psi_\Omega(\Omega^{\Omega+1})}\allowbreak{}，是映射 \allowbreak{}\NovelInline{I01310}{\alpha \mapsto \psi_\Omega((\Omega^\Omega) \cdot \alpha)}\allowbreak{} 的不动点，即 \allowbreak{}\NovelInline{I01311}{\varphi(1,1,0)}\allowbreak{}。''

奥斯卡感叹：``快接近我目前的境界了。很难想象 \allowbreak{}\NovelInline{I01312}{\varphi(1, \Omega+1)}\allowbreak{} 这种 \allowbreak{}\NovelInline{I01313}{\Omega}\allowbreak{} 指数塔极限有多庞大。''

弗兰德鼓励道：``一步一步来。接下来：\allowbreak{}\NovelInline{I01314}{\varphi(1,1,1) = \psi_\Omega((\Omega^{\Omega+1}) \cdot 2)}\allowbreak{}，\allowbreak{}\NovelInline{I01315}{\varphi(1,2,0) = \psi_\Omega(\Omega^{\Omega+2})}\allowbreak{}，能理解吗？''

众人纷纷点头。

弗兰德加快了语速：``我加快速度。\allowbreak{}\NovelInline{I01316}{\psi_\Omega(\Omega^{\Omega \cdot 2}) =}\allowbreak{} 映射 \allowbreak{}\NovelInline{I01317}{\alpha \mapsto \psi_\Omega(\Omega^{\Omega+\alpha})}\allowbreak{} 的不动点，也就是 \allowbreak{}\NovelInline{I01318}{\varphi(2,0,0)}\allowbreak{}。因为 \allowbreak{}\NovelInline{I01319}{\psi_\Omega(\Omega^{\Omega+n}) = \varphi(1,n,0)}\allowbreak{}，\allowbreak{}\NovelInline{I01320}{\alpha \mapsto \varphi(1,\alpha,0)}\allowbreak{} 的不动点自然是 \allowbreak{}\NovelInline{I01321}{\varphi(2,0,0)}\allowbreak{}。接下来 \allowbreak{}\NovelInline{I01322}{\psi_\Omega(\Omega^{\Omega \cdot 3})}\allowbreak{}，你们应该能明确了，就是 \allowbreak{}\NovelInline{I01323}{\varphi(3,0,0)}\allowbreak{}。''

唐三深吸一口气：``进展好快，一下子就到多元 \allowbreak{}\NovelInline{I01324}{\varphi}\allowbreak{} 了。我感觉对序数的理解还需要再夯实。''

``下面开始例题。''弗兰德在黑板上写下一行复杂的公式，``\allowbreak{}\NovelInline{I01325}{\psi_\Omega(\Omega \cdot \psi(\Omega \cdot \psi(0)))}\allowbreak{} 是什么？''

小舞计算道：``\allowbreak{}\NovelInline{I01326}{= \varphi(2, \varphi(2, \varphi(1,0)))}\allowbreak{}。''

``对。那 \allowbreak{}\NovelInline{I01327}{\psi_\Omega(\Omega \cdot 3 + 1)}\allowbreak{}？''

牢大抢答道：``我知道，是 \allowbreak{}\NovelInline{I01328}{\varphi(1, \varphi(2,2)+1)}\allowbreak{}。因为 \allowbreak{}\NovelInline{I01329}{\psi_\Omega(\Omega \cdot 3) = \varphi(2,2)}\allowbreak{}，\allowbreak{}\NovelInline{I01330}{\varphi(1, \varphi(2,2)+1)}\allowbreak{} 就是 \allowbreak{}\NovelInline{I01331}{\varphi(2,2)}\allowbreak{} 的指数塔极限。''

``正确。来个难点的：\allowbreak{}\NovelInline{I01332}{\psi_\Omega(\Omega^\Omega + \Omega^{\psi_\Omega(\Omega^\omega)})}\allowbreak{}。''

蔡徐坤试探着问：``是 \allowbreak{}\NovelInline{I01333}{\varphi(\omega, \varphi(1,0,0)+1)}\allowbreak{} 吗？''

卢本伟摇头反驳：``依我看不是。应该是 \allowbreak{}\NovelInline{I01334}{\varphi(\varphi(\omega,0), \varphi(1,0,0)+1)}\allowbreak{}。\allowbreak{}\NovelInline{I01335}{\psi_\Omega(\Omega^\omega) = \varphi(\omega,0)}\allowbreak{}，\allowbreak{}\NovelInline{I01336}{\Omega^\cdot}\allowbreak{} 的作用是把上层数字放到 \allowbreak{}\NovelInline{I01337}{\varphi}\allowbreak{} 函数第二个位置。''

弗兰德判定：``卢本伟正确。那 \allowbreak{}\NovelInline{I01338}{\psi(\Omega^\omega + \Omega^2)}\allowbreak{} 呢？''

奥斯卡回答：``是 \allowbreak{}\NovelInline{I01339}{\varphi(3, \varphi(\omega,0)+1)}\allowbreak{}。大部分知识我已经掌握了。''

玉小刚趁热打铁：``下面展开：\allowbreak{}\NovelInline{I01340}{\psi_\Omega(\Omega^\Omega + \Omega^2 \cdot \psi_\Omega((\Omega^\Omega) \cdot 2))}\allowbreak{}。基本功扎实，修为才能稳固。''

牢大皱眉：``这怎么展开？又不以 \allowbreak{}\NovelInline{I01341}{\Omega}\allowbreak{} 结尾。而且 \allowbreak{}\NovelInline{I01342}{\psi_\Omega((\Omega^\Omega) \cdot 2) < \Omega}\allowbreak{}。''

唐三沉思片刻：``一定有办法。先写成 \allowbreak{}\NovelInline{I01343}{\psi_\Omega(\Omega^\Omega + \Omega^2 \cdot \psi_\Omega(\Omega^\Omega + \Omega^\Omega))}\allowbreak{} 这种形式？''

小舞提出思路：``先不管括号层数，先找到以 \allowbreak{}\NovelInline{I01344}{\Omega}\allowbreak{} 结尾的结构。\allowbreak{}\NovelInline{I01345}{\psi_\Omega(\Omega^\Omega + \Omega^2 \cdot \psi_\Omega(\Omega^\Omega + \Omega^\Omega))}\allowbreak{} 是映射 \allowbreak{}\NovelInline{I01346}{\alpha \mapsto \psi_\Omega(\Omega^\Omega + \Omega \cdot 2 \cdot \psi_\Omega(\Omega^\Omega +\Omega^\alpha)}\allowbreak{} 的不动点。''小舞顺着思路补充道，眼神中闪烁着兴奋的光芒。

牢大猛地一拍大腿，眼中精光爆射：``我懂了！\NovelDisplay{I01347A}{\psi_\Omega(\Omega^\Omega + \Omega^2 \cdot \psi_\Omega(\Omega^\Omega + \Omega^\Omega))}
\NovelDisplay{I01347B}{=\sup\{\psi_\Omega(\Omega^\Omega + \Omega^2 \cdot \psi_\Omega(\Omega^\Omega + \Omega)), \psi_\Omega(\Omega^\Omega + \Omega^2 \cdot \psi_\Omega(\Omega^\Omega + \Omega^{\psi_\Omega(\Omega^\Omega + \Omega^2 \cdot \psi_\Omega(\dots))})), \dots\}}''

他深吸一口气，语速极快地继续推导：``用 \allowbreak{}\NovelInline{I01348}{\varphi}\allowbreak{} 函数表示，首先 \allowbreak{}\NovelInline{I01349}{\psi_\Omega(\Omega^\Omega + \Omega^2 \cdot \psi_\Omega((\Omega^\Omega) \cdot 2))}\allowbreak{} 介于 \allowbreak{}\NovelInline{I01350}{\psi_\Omega(\Omega^\Omega)}\allowbreak{} 与 \allowbreak{}\NovelInline{I01351}{\psi_\Omega((\Omega^\Omega) \cdot 2)}\allowbreak{} 之间，也就是在 \allowbreak{}\NovelInline{I01352}{\varphi(1,0,0)}\allowbreak{} 与 \allowbreak{}\NovelInline{I01353}{\varphi(1,0,1)}\allowbreak{} 之间。\allowbreak{}\NovelInline{I01354}{\Omega^2}\allowbreak{} 对应 \allowbreak{}\NovelInline{I01355}{\varphi(3,0)}\allowbreak{}，\allowbreak{}\NovelInline{I01356}{\psi_\Omega((\Omega^\Omega) \cdot 2)}\allowbreak{} 对应 \allowbreak{}\NovelInline{I01357}{\varphi(1,0,1)}\allowbreak{}。写出来应该是\ldots\ldots''

蔡徐坤此时也反应过来了，抢着说道：``应该是 \allowbreak{}\NovelInline{I01358}{\varphi(\varphi(4, \varphi(1,0,0)+1), 0)}\allowbreak{}\ldots\ldots 等一下，我再理一理。''

玉小刚突然抬手打断，目光如炬扫视众人：``看吧，这是你们第一个致命错误：没判断是否标准序列就直接展开。''

众人正讨论得热火朝天，此刻都像被掐住脖子的鸭子，瞬间沉默下来。讲堂内只剩下窗外风吹树叶的沙沙声。

玉小刚走到黑板前，用粉笔重重地圈出那个式子：``我问你们，刚才的 \allowbreak{}\NovelInline{I01359}{\psi_\Omega(\Omega^\Omega + \Omega^2 \cdot \psi_\Omega((\Omega^\Omega) \cdot 2))}\allowbreak{} 是不是标准序列？\allowbreak{}\NovelInline{I01360}{\psi_\Omega((\Omega^\Omega) \cdot 2)}\allowbreak{} 的值是多少？''

牢大额头渗出一层细密的汗珠，大脑飞速运转，片刻后沉声道：``\allowbreak{}\NovelInline{I01361}{\psi_\Omega((\Omega^\Omega) \cdot 2) > \psi_\Omega(\Omega^\Omega + \Omega^3) \ge \psi_\Omega(\Omega^\Omega + \Omega^2 \cdot \psi_\Omega((\Omega^\Omega) \cdot 2))}\allowbreak{}。所以\ldots\ldots 它等价于 \allowbreak{}\NovelInline{I01362}{\psi_\Omega(\Omega^\Omega + \Omega^3)}\allowbreak{}。''

``没错！''弗兰德大喝一声，声如洪钟，``战场上没人给你纠错机会。一旦展开错误，魂力运行轨迹就会偏差，轻则走火入魔，重则被敌人抓住弱点一击毙命！那是致命的！''

牢大深吸一口气，强行压下心头的慌乱，重新梳理逻辑：``所以就是 \allowbreak{}\NovelInline{I01363}{\psi_\Omega(\Omega^\Omega + \Omega^2 \cdot \Omega)}\allowbreak{} 的展开，即 \allowbreak{}\NovelInline{I01364}{\varphi(4, \varphi(1,0,0)+1)}\allowbreak{}。也就是： \NovelDisplay{M00171R01}{%
\psi_\Omega(\Omega^\Omega + \Omega^3) = \sup\{\psi_\Omega(\Omega^\Omega + \Omega^2), \psi_\Omega(\Omega^\Omega + \Omega^2 \cdot \psi_\Omega(\Omega^\Omega + \Omega^2)), \dots\}%
} 用 \allowbreak{}\NovelInline{I01365}{\varphi}\allowbreak{} 表示： \NovelDisplay{M00172R01}{%
\varphi(4, \varphi(1,0,0)+1) = \sup\{\varphi(3, \varphi(1,0,0)+1), \varphi(3, \varphi(3, \varphi(1,0,0)+1)), \dots\}%
}

弗兰德微微点头，神色稍缓：``现在明白了？基础扎实才能提升。继续多元 \allowbreak{}\NovelInline{I01366}{\varphi}\allowbreak{} 与 OCF 的讲解。既然已知 \allowbreak{}\NovelInline{I01367}{\psi_\Omega(\Omega^{\Omega \cdot 3})}\allowbreak{}，那 \allowbreak{}\NovelInline{I01368}{\psi_\Omega(\Omega^{\Omega^2})}\allowbreak{} 也就不远了。我直接给出：它就是 \allowbreak{}\NovelInline{I01369}{\varphi(1,0,0,0)}\allowbreak{}。我们把 \allowbreak{}\NovelInline{I01370}{\varphi(1,1)}\allowbreak{} 记作 \allowbreak{}\NovelInline{I01371}{\varphi(1@1, 1@0)}\allowbreak{}，表示 1 在 1 号位，1 在 0 号位。\allowbreak{}\NovelInline{I01372}{\varphi(1,0,0,0) = \varphi(1@3)}\allowbreak{}，懂吗？''

牢大眼神一亮，迅速接话：``\allowbreak{}\NovelInline{I01373}{(1@4, \omega @3, \varphi(2,0)@0)}\allowbreak{} 就是 \allowbreak{}\NovelInline{I01374}{\varphi(1, \omega, 0, 0, \zeta_0)}\allowbreak{}？''

``没错。@ 表示位置序号。''弗兰德赞许道，``好，\allowbreak{}\NovelInline{I01375}{\psi_\Omega(\Omega^{\Omega^2}) = \varphi(1,0,0,0)}\allowbreak{} 没问题。那 \allowbreak{}\NovelInline{I01376}{\psi_\Omega((\Omega^{\Omega^2}) \cdot 2) = \varphi(1,0,0,1)}\allowbreak{} 也没问题吧？''

唐三点头：``可以理解。''

``继续。''弗兰德语速加快，``\allowbreak{}\NovelInline{I01377}{\psi_\Omega(\Omega^{\Omega^2+1})}\allowbreak{} 是什么？''

小舞脱口而出：``映射 \allowbreak{}\NovelInline{I01378}{\alpha \mapsto \psi_\Omega((\Omega^{\Omega^2}) \cdot \alpha)}\allowbreak{} 的不动点，把 \allowbreak{}\NovelInline{I01379}{\Omega}\allowbreak{} 往下移就行。看来是 \allowbreak{}\NovelInline{I01380}{\varphi(1,0,1,0)}\allowbreak{}。''

``化成 @ 记法呢？''

``\allowbreak{}\NovelInline{I01381}{\varphi(1@3, 1@1)}\allowbreak{}。''

``正确。下一个：\allowbreak{}\NovelInline{I01382}{\psi_\Omega(\Omega^{\Omega^2+\Omega})}\allowbreak{} 谁来答？''

牢大自信满满：``我知道。是映射 \allowbreak{}\NovelInline{I01383}{\alpha \mapsto \psi_\Omega(\Omega^{\Omega^2+\alpha})}\allowbreak{} 的不动点，即 \allowbreak{}\NovelInline{I01384}{\varphi(1@3, 1@2)}\allowbreak{}，也就是 \allowbreak{}\NovelInline{I01385}{\varphi(1,1,0,0)}\allowbreak{}。''

玉小刚满意地点头：``很好。那 \allowbreak{}\NovelInline{I01386}{\psi_\Omega(\Omega^{\Omega^2+\Omega \cdot n})}\allowbreak{} 也显而易见了，你们解析！''

蔡徐坤推了推眼镜：``自然是 \allowbreak{}\NovelInline{I01387}{\varphi(1,n,0,0)}\allowbreak{}。乘 \allowbreak{}\NovelInline{I01388}{n}\allowbreak{} 就像 \allowbreak{}\NovelInline{I01389}{\psi_\Omega(\Omega \cdot n) = \varphi(2,n)}\allowbreak{} 一样，OCF 与多元 \allowbreak{}\NovelInline{I01390}{\varphi}\allowbreak{} 完全契合。''

``很好。给出 \allowbreak{}\NovelInline{I01391}{\psi_\Omega(\Omega^{\Omega^2 \cdot 2})}\allowbreak{} 的基本列，并说明理由。''

卢本伟举手道：``我知道。\allowbreak{}\NovelInline{I01392}{\psi_\Omega(\Omega^{\Omega^2 \cdot 2}) = \psi_\Omega(\Omega^{\Omega^2 + \Omega \cdot \Omega})}\allowbreak{}。根据前面结论： \NovelDisplay{M00173R01}{%
\psi_\Omega(\Omega^{\Omega^2 \cdot 2}) = \sup\{\psi_\Omega(\Omega^{\Omega^2+\Omega}), \psi_\Omega(\Omega^{\Omega^2+\Omega \cdot \psi_\Omega(\Omega^{\Omega^2+\Omega})}), \dots\}%
} 写成 \allowbreak{}\NovelInline{I01393}{\varphi}\allowbreak{} 函数就是： \NovelDisplay{M00174R01}{%
\varphi(2,0,0,0) = \sup\{\varphi(1,1,0,0), \varphi(1,\varphi(1,1,0,0),0,0), \dots\}%
}

弗兰德终于露出了笑容：``这次对了。再看 \allowbreak{}\NovelInline{I01394}{\varphi(\omega,0,0,0) = \varphi(\omega @3)}\allowbreak{}，对应 \allowbreak{}\NovelInline{I01395}{\psi_\Omega}\allowbreak{} 如何表示？''

奥斯卡回答：``\allowbreak{}\NovelInline{I01396}{\psi_\Omega(\Omega^{\Omega^2 \cdot \omega})}\allowbreak{}？''

``是。那么 \allowbreak{}\NovelInline{I01397}{\psi_\Omega(\Omega^{\Omega^3}) = \varphi(1@4)}\allowbreak{}，懂了吗？''

唐三：``懂。那 SVO 境界，即 \allowbreak{}\NovelInline{I01398}{\varphi(1@\omega)}\allowbreak{}，不就是 \allowbreak{}\NovelInline{I01399}{\psi_\Omega(\Omega^{\Omega^\omega})}\allowbreak{} 吗？''

弗兰德神色一肃：``正是。但现在要思考：\allowbreak{}\NovelInline{I01400}{\varphi(1@5, \omega @3, \omega^\omega @2, \varepsilon_1@0)}\allowbreak{} 用 \allowbreak{}\NovelInline{I01401}{\psi_\Omega}\allowbreak{} 如何表达？''

牢大略一思索：``先有 \allowbreak{}\NovelInline{I01402}{\psi_\Omega(\Omega^{\Omega^4})}\allowbreak{}，再处理 \allowbreak{}\NovelInline{I01403}{\omega @3}\allowbreak{}。''

蔡徐坤试探着问：``\allowbreak{}\NovelInline{I01404}{\psi_\Omega(\Omega^{\Omega^4 + (\Omega^2) \cdot \omega + \Omega \cdot \omega^\omega} \cdot \varepsilon_1)}\allowbreak{}？''

唐三摇头纠正：``不对，指数运算优先级要注意。是 \allowbreak{}\NovelInline{I01405}{\psi_\Omega(\Omega^{\Omega^4 + \Omega^2 \cdot \omega + \Omega \cdot \omega^\omega} \cdot \varepsilon_1)}\allowbreak{}。''

牢大点头确认：``对，就是这个。''

弗兰德环视众人，目光灼灼：``唐三正确。以后还要多加练习，务必熟练掌握。距离决赛只剩一个月，能否成功，就看这一个月的修炼进度了。''

\begin{center}\rule{0.5\linewidth}{0.5pt}\end{center}

理论课结束，真正的考验才刚开始。你觉得接下来是安排一场模拟战检验成果，还是让玉小刚讲讲更深层的``非标准序列''陷阱？

\chapter{第二十一章：序元Veblen函数}\label{ux7b2cux4e8cux5341ux4e00ux7ae0ux5e8fux5143veblenux51fdux6570}

OCF的讲解结束后，果糕七怪战队的队员们虽已掌握 \allowbreak{}\NovelInline{I01406}{\psi_\Omega}\allowbreak{} 函数与多元 \allowbreak{}\NovelInline{I01407}{\varphi}\allowbreak{} 函数的对应关系，但在实战中仍存在高阶序数展开卡顿、不动点判断失误的问题------尤其是面对Madore宗、Buchholz宗可能掌握的更高阶序数技巧，现有修为仍有不足。

次日清晨，修炼场的晨光尚未完全铺展，玉小刚便带着厚厚的典籍，与弗兰德并肩站在众人面前，神色严肃。

``昨日我们讲解了OCF与多元Veblen函数的关联，你们也初步掌握了 \allowbreak{}\NovelInline{I01408}{\varphi(1 @3)}\allowbreak{}、\allowbreak{}\NovelInline{I01409}{\varphi(1 @4)}\allowbreak{} 这类简单的序元记号，''玉小刚抬手，指尖凝聚出一道莹白的序链，在空气中勾勒出 \allowbreak{}\NovelInline{I01410}{\varphi(1 @\omega)}\allowbreak{} 的符号，``但多元Veblen函数的位置仅限于自然数，就像你们修炼时只能突破有限层次的境界，无法触及更高阶的战力天花板。今天，我们重点讲解一组高阶序元函数------从 \allowbreak{}\NovelInline{I01411}{\varphi(1 @\omega)}\allowbreak{} 到LVO，先结合实例吃透每个序数的构造，再总结通用规则，这是你们突破SVO境界、应对复赛强敌的关键。''

奥斯卡揉了揉发胀的额头，率先开口提问：``玉老师，我们先从最基础的 \allowbreak{}\NovelInline{I01412}{\varphi(1 @\omega)}\allowbreak{} 开始吧？之前您说它就是SVO，具体怎么理解它的构造和作用呢？''

弗兰德上前一步，语气放缓，结合实例拆解：

\begin{quote}
很好的问题，我们先从 \allowbreak{}\NovelInline{I01413}{\varphi(1 @\omega)}\allowbreak{}（小Veblen序数SVO）讲起，这是我们今天所有高阶序数的基础。大家还记得多元Veblen函数的序列吗？\allowbreak{}\NovelInline{I01414}{\varphi(1 @0)=\varphi(1)}\allowbreak{}、\allowbreak{}\NovelInline{I01415}{\varphi(1 @1)=\varphi(1,0)}\allowbreak{}、\allowbreak{}\NovelInline{I01416}{\varphi(1 @2)=\varphi(1,0,0)}\allowbreak{}，以此类推，\allowbreak{}\NovelInline{I01417}{\varphi(1 @n)}\allowbreak{} 就是''1后面跟 \allowbreak{}\NovelInline{I01418}{n}\allowbreak{} 个0''的多元函数。而 \allowbreak{}\NovelInline{I01419}{\varphi(1 @\omega)}\allowbreak{}，就是这个无穷序列的上确界，简单说就是把 \allowbreak{}\NovelInline{I01420}{n}\allowbreak{} 换成极限序数 \allowbreak{}\NovelInline{I01421}{\omega}\allowbreak{}，取所有 \allowbreak{}\NovelInline{I01422}{\varphi(1 @n)}\allowbreak{} 的极限。
\end{quote}

他抬手在空中补充了一组序列，指尖序链流转：

\begin{quote}
你们看，\allowbreak{}\NovelInline{I01423}{\varphi(1 @0)=\varphi(1)}\allowbreak{}，\allowbreak{}\NovelInline{I01424}{\varphi(1 @1)=\varphi(1,0)}\allowbreak{}，\allowbreak{}\NovelInline{I01425}{\varphi(1 @2)=\varphi(1,0,0)}\allowbreak{}，\ldots\ldots，当 \allowbreak{}\NovelInline{I01426}{n}\allowbreak{} 不断增大，\allowbreak{}\NovelInline{I01427}{\varphi(1 @n)}\allowbreak{} 的序数也不断提升，而 \allowbreak{}\NovelInline{I01428}{\varphi(1 @\omega)}\allowbreak{} 就是这串无穷序列的''终点''，即：

\NovelDisplay{M00175R01}{%
\text{SVO} = \sup\{\varphi(1 @0),\, \varphi(1 @1),\, \varphi(1 @2),\, \dots\}%
}

另外，它还有一个核心性质------是所有映射 \allowbreak{}\NovelInline{I01429}{\delta \to \varphi(\delta @n)}\allowbreak{} 的第一个共同不动点，也就是说，不管 \allowbreak{}\NovelInline{I01430}{n}\allowbreak{} 取哪个自然数，\allowbreak{}\NovelInline{I01431}{\varphi(\text{SVO} @n) = \text{SVO}}\allowbreak{}，这也是它能成为''战力天花板突破口''的关键。
\end{quote}

卢本韦皱了皱眉，追问：``那 \allowbreak{}\NovelInline{I01432}{\varphi(1 @\omega,\, 1 @0)}\allowbreak{} 呢？它和SVO有什么关系？实战中我们怎么展开它的基本列？''

玉小刚点头，赞许地看了卢本韦一眼，结合实例讲解：

\begin{quote}
问得很关键，\allowbreak{}\NovelInline{I01433}{\varphi(1 @\omega,\, 1 @0)}\allowbreak{} 是在SVO基础上的下一个高阶序数，我们可以把它理解为''第二个共同不动点''------既然SVO是所有映射 \allowbreak{}\NovelInline{I01434}{\delta \to \varphi(\delta @n)}\allowbreak{} 的第一个共同不动点，那 \allowbreak{}\NovelInline{I01435}{\varphi(1 @\omega,\, 1 @0)}\allowbreak{} 就是第二个。
\end{quote}

他在地面上写下基本列公式，耐心拆解：

\begin{quote}
它的构造很简单，先给SVO取后继，得到 \allowbreak{}\NovelInline{I01436}{\text{SVO}+1}\allowbreak{}，然后让这个后继分别作用于每个 \allowbreak{}\NovelInline{I01437}{\varphi(\delta @n)}\allowbreak{}，再取所有结果的上确界，也就是：

\NovelDisplay{M00176R01}{%
\varphi(1 @\omega,\, 1 @0) = \sup\{\varphi((\text{SVO}+1) @0),\, \varphi((\text{SVO}+1) @1),\, \varphi((\text{SVO}+1) @2),\, \dots\}%
}

对应的基本列就是 \allowbreak{}\NovelInline{I01438}{\varphi(1 @\omega,\, 1 @0)[n] = \varphi(\text{SVO}+1 @n)}\allowbreak{}，比如 \allowbreak{}\NovelInline{I01439}{n=0}\allowbreak{} 时，就是 \allowbreak{}\NovelInline{I01440}{\varphi((\text{SVO}+1) @0) = \varphi(\text{SVO}+1)}\allowbreak{}；\allowbreak{}\NovelInline{I01441}{n=1}\allowbreak{} 时，就是 \allowbreak{}\NovelInline{I01442}{\varphi((\text{SVO}+1) @1) = \varphi(\text{SVO}+1,\, 0)}\allowbreak{}，以此类推，实战中只要按这个规律展开，就不会出错。
\end{quote}

小舞举手，眼神明亮：``老师，那 \allowbreak{}\NovelInline{I01443}{\varphi(1 @\omega,\, 1 @1)}\allowbreak{} 是不是第三个共同不动点？它的构造和 \allowbreak{}\NovelInline{I01444}{\varphi(1 @\omega,\, 1 @0)}\allowbreak{} 类似吗？''

``不完全是共同不动点，而是另一种迭代层次。''弗兰德笑着解答，结合规则实例：

\begin{quote}
\allowbreak{}\NovelInline{I01445}{\varphi(1 @\omega,\, 1 @1)}\allowbreak{} 的核心是''映射迭代升级''------它不再是针对所有 \allowbreak{}\NovelInline{I01446}{\delta \to \varphi(\delta @n)}\allowbreak{}，而是针对映射 \allowbreak{}\NovelInline{I01447}{\delta \to \varphi(1 @\omega,\, \delta @0)}\allowbreak{} 的第一个不动点。简单说，\allowbreak{}\NovelInline{I01448}{\varphi(1 @\omega,\, 1 @1)}\allowbreak{} 就是满足：

\NovelDisplay{M00177R01}{%
\varphi(1 @\omega,\, \varphi(1 @\omega,\, 1 @1) @0) = \varphi(1 @\omega,\, 1 @1)%
}

的第一个序数。
\end{quote}

他补充道：

\begin{quote}
举个通俗的例子，\allowbreak{}\NovelInline{I01449}{\varphi(1 @\omega,\, 1 @0)}\allowbreak{} 是''第二个共同不动点''，而 \allowbreak{}\NovelInline{I01450}{\varphi(1 @\omega,\, 1 @1)}\allowbreak{} 是''以 \allowbreak{}\NovelInline{I01451}{\varphi(1 @\omega,\, \delta @0)}\allowbreak{} 为映射的第一个不动点''，相当于在SVO的基础上，又提升了一个迭代层级。它的基本列的构造，和我们之前讲的''\allowbreak{}\NovelInline{I01452}{\beta}\allowbreak{} 为后继序数''的规则一致，后续我们讲规则时再详细拆解。
\end{quote}

唐三指尖微动，在地面上写下 \allowbreak{}\NovelInline{I01453}{\varphi(2 @\omega)}\allowbreak{}，轻声追问：``弗兰德老师，\allowbreak{}\NovelInline{I01454}{\varphi(2 @\omega)}\allowbreak{} 又是什么？它比SVO和 \allowbreak{}\NovelInline{I01455}{\varphi(1 @\omega,\, 1 @0)}\allowbreak{} 更高阶吗？''

``没错，\allowbreak{}\NovelInline{I01456}{\varphi(2 @\omega)}\allowbreak{} 是比前两者更高阶的序数，它是另一个'上确界+共同不动点'的组合。''玉小刚上前，详细讲解实例：

\begin{quote}
我们先看一组序列：\allowbreak{}\NovelInline{I01457}{\varphi(1 @\omega,\, 1 @0)}\allowbreak{}、\allowbreak{}\NovelInline{I01458}{\varphi(1 @\omega,\, 1 @1)}\allowbreak{}、\allowbreak{}\NovelInline{I01459}{\varphi(1 @\omega,\, 1 @2)}\allowbreak{}、\ldots\ldots，这组序列的上确界，就是 \allowbreak{}\NovelInline{I01460}{\varphi(2 @\omega)}\allowbreak{}，即：

\NovelDisplay{M00178R01}{%
\varphi(2 @\omega) = \sup\{\varphi(1 @\omega,\, 1 @0),\, \varphi(1 @\omega,\, 1 @1),\, \varphi(1 @\omega,\, 1 @2),\, \dots\}%
}
\end{quote}

他进一步补充核心性质：

\begin{quote}
同时，\allowbreak{}\NovelInline{I01461}{\varphi(2 @\omega)}\allowbreak{} 也是所有映射 \allowbreak{}\NovelInline{I01462}{\delta \to \varphi(1 @\omega,\, \delta @n)}\allowbreak{} 的第一个共同不动点------也就是说，对于任意自然数 \allowbreak{}\NovelInline{I01463}{n}\allowbreak{}，\allowbreak{}\NovelInline{I01464}{\varphi(1 @\omega,\, \varphi(2 @\omega) @n) = \varphi(2 @\omega)}\allowbreak{}。实战中，我们可以这样理解：\allowbreak{}\NovelInline{I01465}{\varphi(1 @\omega)}\allowbreak{} 是''第一层天花板''，\allowbreak{}\NovelInline{I01466}{\varphi(1 @\omega,\, 1 @n)}\allowbreak{} 是''天花板下的进阶''，而 \allowbreak{}\NovelInline{I01467}{\varphi(2 @\omega)}\allowbreak{} 是''第二层天花板''。
\end{quote}

蔡徐坤拿出随身携带的笔记，快速记录，同时提问：``那 \allowbreak{}\NovelInline{I01468}{\varphi(\omega @\omega)}\allowbreak{} 呢？它是不是把 \allowbreak{}\NovelInline{I01469}{\varphi(n @\omega)}\allowbreak{} 的序列取上确界得到的？''

``完全正确，蔡徐坤的思路很清晰。''弗兰德赞许道，结合实例拆解：

\begin{quote}
我们先梳理 \allowbreak{}\NovelInline{I01470}{\varphi(n @\omega)}\allowbreak{} 的序列：\allowbreak{}\NovelInline{I01471}{\varphi(1 @\omega) = \text{SVO}}\allowbreak{}，\allowbreak{}\NovelInline{I01472}{\varphi(2 @\omega)}\allowbreak{} 是第二层天花板，\allowbreak{}\NovelInline{I01473}{\varphi(3 @\omega)}\allowbreak{} 是所有映射 \allowbreak{}\NovelInline{I01474}{\delta \to \varphi(2 @\omega,\, \delta @n)}\allowbreak{} 的第一个共同不动点，以此类推，\allowbreak{}\NovelInline{I01475}{\varphi(n @\omega)}\allowbreak{} 就是不断迭代''上一层天花板''的结果。而 \allowbreak{}\NovelInline{I01476}{\varphi(\omega @\omega)}\allowbreak{}，就是这组无穷序列的上确界，即：

\NovelDisplay{M00179R01}{%
\varphi(\omega @\omega) = \sup\{1,\, \varphi(1 @\omega),\, \varphi(2 @\omega),\, \dots\}%
}
\end{quote}

他抬手写下 \allowbreak{}\NovelInline{I01477}{\varphi(\omega @\omega,\, 1 @0)}\allowbreak{} 的基本列，补充道：

\begin{quote}
比如 \allowbreak{}\NovelInline{I01478}{\varphi(\omega @\omega,\, 1 @0)}\allowbreak{}，它是所有映射 \allowbreak{}\NovelInline{I01479}{\delta \to \varphi(n @\omega,\, \delta @n)}\allowbreak{} 的第二个共同不动点，基本列是：

\NovelDisplay{M00180R01}{%
\varphi(\omega @\omega,\, 1 @0)[n] = \varphi(n @\omega,\, (\varphi(\omega @\omega)+1) @n)%
}

简单说，就是给 \allowbreak{}\NovelInline{I01480}{\varphi(\omega @\omega)}\allowbreak{} 取后继，然后让这个后继分别作用于 \allowbreak{}\NovelInline{I01481}{\varphi(n @\omega,\, \delta @n)}\allowbreak{}，再取极限，和 \allowbreak{}\NovelInline{I01482}{\varphi(1 @\omega,\, 1 @0)}\allowbreak{} 的构造逻辑一致，只是迭代层次更高。
\end{quote}

牢大皱起眉，沉吟道：``老师，那 \allowbreak{}\NovelInline{I01483}{\varphi(1 @\omega+1)}\allowbreak{} 呢？\allowbreak{}\NovelInline{I01484}{\omega+1}\allowbreak{} 是后继序数，和 \allowbreak{}\NovelInline{I01485}{\varphi(1 @\omega)}\allowbreak{}（\allowbreak{}\NovelInline{I01486}{\omega}\allowbreak{} 是极限序数）的构造有什么区别？''

玉小刚点头，针对性讲解：

\begin{quote}
\allowbreak{}\NovelInline{I01487}{\varphi(1 @\omega+1)}\allowbreak{} 的关键的是''\allowbreak{}\NovelInline{I01488}{\beta}\allowbreak{} 为后继序数''------这里的 \allowbreak{}\NovelInline{I01489}{\beta = \omega+1}\allowbreak{}，是极限序数 \allowbreak{}\NovelInline{I01490}{\omega}\allowbreak{} 的后继。它的核心性质是：\allowbreak{}\NovelInline{I01491}{\varphi(1 @\omega+1,\, \gamma @0)}\allowbreak{} 是映射 \allowbreak{}\NovelInline{I01492}{\delta \to \varphi(\delta @\omega)}\allowbreak{} 的第 \allowbreak{}\NovelInline{I01493}{1+\gamma}\allowbreak{} 个不动点。
\end{quote}

他举了一个具体例子：

\begin{quote}
比如 \allowbreak{}\NovelInline{I01494}{\gamma=0}\allowbreak{} 时，\allowbreak{}\NovelInline{I01495}{\varphi(1 @\omega+1,\, 0 @0)}\allowbreak{} 就是映射 \allowbreak{}\NovelInline{I01496}{\delta \to \varphi(\delta @\omega)}\allowbreak{} 的第一个不动点，也就是说：

\NovelDisplay{M00181R01}{%
\varphi(\varphi(1 @\omega+1,\, 0 @0) @\omega) = \varphi(1 @\omega+1,\, 0 @0)%
}

它的基本列构造，遵循''\allowbreak{}\NovelInline{I01497}{\beta}\allowbreak{} 为后继序数''的规则，后续我们会统一梳理，大家现在重点记住：\allowbreak{}\NovelInline{I01498}{\varphi(1 @\omega+1)}\allowbreak{} 的迭代层级，比 \allowbreak{}\NovelInline{I01499}{\varphi(1 @\omega)}\allowbreak{} 更高，因为它的''参数 \allowbreak{}\NovelInline{I01500}{\beta}\allowbreak{}``从 \allowbreak{}\NovelInline{I01501}{\omega}\allowbreak{}（极限）变成了 \allowbreak{}\NovelInline{I01502}{\omega+1}\allowbreak{}（后继）。
\end{quote}

牢大听得认真，忍不住提问：``那如果 \allowbreak{}\NovelInline{I01503}{\beta}\allowbreak{} 是任意极限序数 \allowbreak{}\NovelInline{I01504}{\alpha}\allowbreak{}，比如 \allowbreak{}\NovelInline{I01505}{\alpha = \omega \times 2}\allowbreak{}、\allowbreak{}\NovelInline{I01506}{\alpha = \omega^2}\allowbreak{}，\allowbreak{}\NovelInline{I01507}{\varphi(1 @\alpha)}\allowbreak{} 该怎么理解？''

``问得非常好，这就是 \allowbreak{}\NovelInline{I01508}{\varphi(1 @\alpha)}\allowbreak{}（\allowbreak{}\NovelInline{I01509}{\alpha}\allowbreak{} 为极限序数）的核心------把 \allowbreak{}\NovelInline{I01510}{\beta}\allowbreak{} 推广到任意极限序数，迭代层次进一步提升。''弗兰德详细讲解，结合实例：

\begin{quote}
比如 \allowbreak{}\NovelInline{I01511}{\alpha = \omega \times 2}\allowbreak{}（\allowbreak{}\NovelInline{I01512}{\omega}\allowbreak{} 乘以2，也是一个极限序数），\allowbreak{}\NovelInline{I01513}{\varphi(1 @\omega \times 2,\, \gamma @0)}\allowbreak{} 就是所有映射 \allowbreak{}\NovelInline{I01514}{\delta \to \varphi(\delta @(\omega+n))}\allowbreak{} 的第 \allowbreak{}\NovelInline{I01515}{1+\gamma}\allowbreak{} 个共同不动点。
\end{quote}

他补充道：

\begin{quote}
再比如 \allowbreak{}\NovelInline{I01516}{\alpha = \omega^2}\allowbreak{}（\allowbreak{}\NovelInline{I01517}{\omega}\allowbreak{} 的平方），\allowbreak{}\NovelInline{I01518}{\varphi(1 @\omega^2,\, \gamma @0)}\allowbreak{} 就是所有映射 \allowbreak{}\NovelInline{I01519}{\delta \to \varphi(\delta @(\omega \times n))}\allowbreak{} 的第 \allowbreak{}\NovelInline{I01520}{1+\gamma}\allowbreak{} 个共同不动点。简单说，\allowbreak{}\NovelInline{I01521}{\alpha}\allowbreak{} 是哪个极限序数，就针对''\allowbreak{}\NovelInline{I01522}{\delta @\alpha}\allowbreak{} 的基本列''取共同不动点，构造逻辑和 \allowbreak{}\NovelInline{I01523}{\varphi(1 @\omega)}\allowbreak{} 一致，只是范围更广、层级更高。
\end{quote}

唐三皱起眉，仔细思考：``老师，那LVO（大Veblen序数）就是这些 \allowbreak{}\NovelInline{I01524}{\varphi(1 @\alpha)}\allowbreak{} 序列的极限吗？它的核心性质是什么？''

``没错，LVO是我们今天讲解的最高阶序数，也是序元Veblen函数的极限。''玉小刚语气严肃，结合实例拆解：

\begin{quote}
我们先看一组迭代序列：\allowbreak{}\NovelInline{I01525}{0}\allowbreak{}、\allowbreak{}\NovelInline{I01526}{\varphi(1 @0)}\allowbreak{}、\allowbreak{}\NovelInline{I01527}{\varphi(1 @\varphi(1 @0))}\allowbreak{}、\allowbreak{}\NovelInline{I01528}{\varphi(1 @\varphi(1 @\varphi(1 @0)))}\allowbreak{}、\ldots\ldots，这组序列是不断迭代 \allowbreak{}\NovelInline{I01529}{\varphi(1 @\alpha)}\allowbreak{} 得到的------先取 \allowbreak{}\NovelInline{I01530}{\varphi(1 @0)}\allowbreak{}，再把 \allowbreak{}\NovelInline{I01531}{\alpha}\allowbreak{} 换成 \allowbreak{}\NovelInline{I01532}{\varphi(1 @0)}\allowbreak{} 得到 \allowbreak{}\NovelInline{I01533}{\varphi(1 @\varphi(1 @0))}\allowbreak{}，再把 \allowbreak{}\NovelInline{I01534}{\alpha}\allowbreak{} 换成这个结果，以此类推，而LVO就是这组无穷序列的上确界，即：

\NovelDisplay{M00182R01}{%
\text{LVO} = \sup\{0,\, \varphi(1 @0),\, \varphi(1 @\varphi(1 @0)),\, \dots\}%
}
\end{quote}

他强调核心性质：

\begin{quote}
LVO的关键特征，是它是映射 \allowbreak{}\NovelInline{I01535}{\alpha \to \varphi(1 @\alpha)}\allowbreak{} 的第一个不动点，也就是说：

\NovelDisplay{M00183R01}{%
\text{LVO} = \varphi(1 @\text{LVO})%
}

这是什么意思呢？就是当 \allowbreak{}\NovelInline{I01536}{\alpha = \text{LVO}}\allowbreak{} 时，\allowbreak{}\NovelInline{I01537}{\varphi(1 @\alpha)}\allowbreak{} 等于 \allowbreak{}\NovelInline{I01538}{\alpha}\allowbreak{} 本身，这是之前所有序数都不具备的性质------比如 \allowbreak{}\NovelInline{I01539}{\varphi(1 @\omega) \neq \omega}\allowbreak{}，\allowbreak{}\NovelInline{I01540}{\varphi(1 @\varphi(1 @\omega)) \neq \varphi(1 @\omega)}\allowbreak{}，而只有LVO，能满足这个不动点条件，这也是它''大Veblen序数''名字的由来。
\end{quote}

众人纷纷点头，快速记录下各个序数的实例和核心性质，奥斯卡忍不住追问：``老师，例子我们大概懂了，现在能不能梳理一下通用规则？这样我们实战中遇到类似序数，就能自己展开了。''

``当然可以，例子是基础，规则是核心。''弗兰德走到修炼场中央，指尖序链流转，将核心规则逐一展现在空气中，``我们分两类规则来梳理------\textbf{不动点规则}（决定序数的构造）和\textbf{基本列规则}（决定序数的展开），所有我们刚才讲的例子，都能对应到这些规则里。''

本质上是根据 \allowbreak{}\NovelInline{I01541}{\alpha}\allowbreak{}、\allowbreak{}\NovelInline{I01542}{\beta}\allowbreak{} 的类型（后继/极限序数），判断 \allowbreak{}\NovelInline{I01543}{\varphi(\#,\, \alpha @\beta,\, \gamma @0)}\allowbreak{} 的构造方式，其中 \allowbreak{}\NovelInline{I01544}{\#}\allowbreak{} 是任意不变序列，\allowbreak{}\NovelInline{I01545}{\gamma @0}\allowbreak{} 是末项（负责不动点计数），\allowbreak{}\NovelInline{I01546}{\alpha @\beta}\allowbreak{} 是除末项外最左边第一个非零项（负责迭代层次）：

\allowbreak{}\NovelInline{I01547}{\varphi(\#,\, \alpha+1 @\beta+1,\, \gamma @0)}\allowbreak{}，是映射 \allowbreak{}\NovelInline{I01548}{\delta \to \varphi(\#,\, \alpha @\beta+1,\, \delta @\beta)}\allowbreak{} 的第 \allowbreak{}\NovelInline{I01549}{1+\gamma}\allowbreak{} 个不动点。

\begin{quote}
\textbf{示例}：\allowbreak{}\NovelInline{I01550}{\varphi(1 @1)}\allowbreak{}（\allowbreak{}\NovelInline{I01551}{\alpha=0+1}\allowbreak{}，\allowbreak{}\NovelInline{I01552}{\beta=0+1}\allowbreak{}），就是映射 \allowbreak{}\NovelInline{I01553}{\delta \to \varphi(0 @1,\, \delta @0)}\allowbreak{} 的第 \allowbreak{}\NovelInline{I01554}{1+0}\allowbreak{} 个不动点。
\end{quote}

\allowbreak{}\NovelInline{I01555}{\varphi(\#,\, \alpha @\beta+1,\, \gamma @0)}\allowbreak{}，是所有映射 \allowbreak{}\NovelInline{I01556}{\delta \to \varphi(\#,\, \alpha[n] @\beta+1,\, \delta @\beta)}\allowbreak{} 的第 \allowbreak{}\NovelInline{I01557}{1+\gamma}\allowbreak{} 个共同不动点。

\begin{quote}
\textbf{示例}：\allowbreak{}\NovelInline{I01558}{\varphi(\omega @1)}\allowbreak{}（\allowbreak{}\NovelInline{I01559}{\alpha=\omega}\allowbreak{} 极限，\allowbreak{}\NovelInline{I01560}{\beta=0+1}\allowbreak{}），就是所有 \allowbreak{}\NovelInline{I01561}{\delta \to \varphi(\omega[n] @1,\, \delta @0)}\allowbreak{} 的共同不动点 以下是从''规则2''之后的示例处接续，直到章节末尾的完整内容：
\end{quote}

\begin{center}\rule{0.5\linewidth}{0.5pt}\end{center}

\begin{quote}
\textbf{示例}：\allowbreak{}\NovelInline{I01562}{\varphi(\omega @1)}\allowbreak{}（\allowbreak{}\NovelInline{I01563}{\alpha=\omega}\allowbreak{} 极限，\allowbreak{}\NovelInline{I01564}{\beta=0+1}\allowbreak{}），就是所有 \allowbreak{}\NovelInline{I01565}{\delta \to \varphi(\omega[n] @1,\, \delta @0)}\allowbreak{} 的共同不动点。
\end{quote}

\allowbreak{}\NovelInline{I01566}{\varphi(\#,\, \alpha+1 @\beta,\, \gamma @0)}\allowbreak{}，是所有映射 \allowbreak{}\NovelInline{I01567}{\delta \to \varphi(\#,\, \alpha @\beta,\, \delta @\beta[n])}\allowbreak{} 的第 \allowbreak{}\NovelInline{I01568}{1+\gamma}\allowbreak{} 个共同不动点。

\begin{quote}
\textbf{示例}：\allowbreak{}\NovelInline{I01569}{\varphi(1 @\omega)}\allowbreak{}（\allowbreak{}\NovelInline{I01570}{\alpha=0+1}\allowbreak{}，\allowbreak{}\NovelInline{I01571}{\beta=\omega}\allowbreak{} 极限），就是所有 \allowbreak{}\NovelInline{I01572}{\delta \to \varphi(0 @\omega,\, \delta @\omega[n])}\allowbreak{} 的共同不动点。
\end{quote}

\allowbreak{}\NovelInline{I01573}{\varphi(\#,\, \alpha @\beta,\, \gamma @0)}\allowbreak{}，是所有映射 \allowbreak{}\NovelInline{I01574}{\delta \to \varphi(\#,\, \alpha[n] @\beta,\, \delta @\beta[n])}\allowbreak{} 的第 \allowbreak{}\NovelInline{I01575}{1+\gamma}\allowbreak{} 个共同不动点。

\begin{quote}
\textbf{示例}：\allowbreak{}\NovelInline{I01576}{\varphi(\omega @\omega)}\allowbreak{}（\allowbreak{}\NovelInline{I01577}{\alpha=\omega}\allowbreak{} 极限，\allowbreak{}\NovelInline{I01578}{\beta=\omega}\allowbreak{} 极限），就是所有 \allowbreak{}\NovelInline{I01579}{\delta \to \varphi(\omega[n] @\omega,\, \delta @\omega[n])}\allowbreak{} 的共同不动点。
\end{quote}

玉小刚补充道：

\begin{quote}
这4条规则可以统一成一句话：\allowbreak{}\NovelInline{I01580}{\varphi(\#,\, \alpha @\beta,\, \gamma @0)}\allowbreak{}，是第 \allowbreak{}\NovelInline{I01581}{1+\gamma}\allowbreak{} 个满足''对所有 \allowbreak{}\NovelInline{I01582}{\delta < \alpha}\allowbreak{}、\allowbreak{}\NovelInline{I01583}{\sigma < \beta}\allowbreak{}，都是映射 \allowbreak{}\NovelInline{I01584}{\rho \to \varphi(\#,\, \delta @\beta,\, \rho @\sigma)}\allowbreak{} 的不动点''的序数。大家记住，判断序数构造，先看 \allowbreak{}\NovelInline{I01585}{\alpha}\allowbreak{} 和 \allowbreak{}\NovelInline{I01586}{\beta}\allowbreak{} 是后继还是极限，再对应到这4条规则，就不会出错。
\end{quote}

``接下来是序元Veblen函数的基本列规则，共13条，我们结合刚才的例子，挑核心的、实战中常用的几条重点讲解，剩下的可以结合参考资料巩固：''

对任意序数 \allowbreak{}\NovelInline{I01587}{\alpha \geq \beta}\allowbreak{}：

\NovelDisplay{M00184R01}{%
\alpha + \beta[n] = \alpha + \beta[n]%
}

\begin{quote}
\textbf{示例}：\allowbreak{}\NovelInline{I01588}{\varphi(1 @\omega) + \omega[n] = \varphi(1 @\omega) + n}\allowbreak{}，因为 \allowbreak{}\NovelInline{I01589}{\omega[n] = n}\allowbreak{}，这条和多元函数一致，很简单。
\end{quote}

\NovelDisplay{M00185R01}{%
\varphi(0 @0) = 1%
}

\begin{quote}
注意和 \allowbreak{}\NovelInline{I01590}{\varphi(1 @0) = \varphi(1)}\allowbreak{} 区分，这是序元函数的最基础序数，不要混淆。
\end{quote}

\NovelDisplay{M00186R01}{%
\varphi(\#,\, \alpha+1 @\beta+1)[0] = 0%
}

\NovelDisplay{M00187R01}{%
\varphi(\#,\, \alpha+1 @\beta+1)[n+1] = \varphi(\#,\, \alpha @\beta+1,\, \varphi(\#,\, \alpha+1 @\beta+1)[n] @\beta)%
}

\begin{quote}
\textbf{示例}：\allowbreak{}\NovelInline{I01591}{\varphi(1 @1)[0] = 0}\allowbreak{}，\allowbreak{}\NovelInline{I01592}{\varphi(1 @1)[1] = \varphi(0 @1,\, 0 @0) = 1}\allowbreak{}，就是这条规则。
\end{quote}

\NovelDisplay{M00188R01}{%
\varphi(\#,\, \alpha+1 @\beta+1,\, \gamma+1 @0)[0] = \varphi(\#,\, \alpha+1 @\beta+1,\, \gamma @0) + 1%
}

\NovelDisplay{M00189R01}{%
\varphi(\#,\, \alpha+1 @\beta+1,\, \gamma+1 @0)[n+1] = \varphi(\#,\, \alpha @\beta+1,\, \varphi(\#,\, \alpha+1 @\beta+1,\, \gamma+1 @0)[n] @\beta)%
}

\begin{quote}
\textbf{示例}：\allowbreak{}\NovelInline{I01593}{\varphi(1 @1,\, 1 @0)[0] = \varphi(1 @1,\, 0 @0) + 1}\allowbreak{}，就是这条。
\end{quote}

\NovelDisplay{M00190R01}{%
\varphi(\#,\, \alpha @\beta,\, \gamma @0)[n] = \varphi(\#,\, \alpha @\beta,\, \gamma[n] @0)%
}

\begin{quote}
\textbf{示例}：\allowbreak{}\NovelInline{I01594}{\varphi(1 @1,\, \omega @0)[n] = \varphi(1 @1,\, n @0)}\allowbreak{}，因为 \allowbreak{}\NovelInline{I01595}{\omega[n] = n}\allowbreak{}。
\end{quote}

\NovelDisplay{M00191R01}{%
\varphi(\#,\, \alpha @\beta+1)[n] = \varphi(\#,\, \alpha[n] @\beta+1)%
}

\begin{quote}
\textbf{示例}：\allowbreak{}\NovelInline{I01596}{\varphi(\omega @1)[n] = \varphi(n @1)}\allowbreak{}，因为 \allowbreak{}\NovelInline{I01597}{\omega[n] = n}\allowbreak{}。
\end{quote}

\NovelDisplay{M00192R01}{%
\varphi(\#,\, \alpha+1 @\beta)[n] = \varphi(\#,\, \alpha @\beta,\, 1 @\beta[n])%
}

\begin{quote}
\textbf{示例}：\allowbreak{}\NovelInline{I01598}{\varphi(1 @\omega)[n] = \varphi(0 @\omega,\, 1 @n) = \varphi(1 @n)}\allowbreak{}，这就是SVO的基本列，我们之前重点讲过。
\end{quote}

\NovelDisplay{M00193R01}{%
\varphi(\#,\, \alpha @\beta)[n] = \varphi(\#,\, \alpha[n] @\beta)%
}

\begin{quote}
\textbf{示例}：\allowbreak{}\NovelInline{I01599}{\varphi(\omega @\omega)[n] = \varphi(n @\omega)}\allowbreak{}，因为 \allowbreak{}\NovelInline{I01600}{\omega[n] = n}\allowbreak{}。
\end{quote}

弗兰德总结道：

\begin{quote}
这8条是实战中最常用的基本列规则，剩下的5条，本质上是这8条的延伸，结合 \allowbreak{}\NovelInline{I01601}{\alpha}\allowbreak{}、\allowbreak{}\NovelInline{I01602}{\beta}\allowbreak{}、\allowbreak{}\NovelInline{I01603}{\gamma}\allowbreak{} 的类型组合即可。大家记住一个核心：基本列的展开，关键看 \allowbreak{}\NovelInline{I01604}{\alpha}\allowbreak{}（非零项的序数）、\allowbreak{}\NovelInline{I01605}{\beta}\allowbreak{}（参数序数）、\allowbreak{}\NovelInline{I01606}{\gamma}\allowbreak{}（末项序数）是后继还是极限，再对应到相应规则，就能精准展开。
\end{quote}

顿了顿，弗兰德话锋一转，语气变得更为严谨：``在继续深入练习之前，有一个关键概念必须吃透------序元Veblen函数的\textbf{标准形式}。我们之前讲的所有序数展开、不动点判断，都必须建立在'序数是标准形式'的基础上，实战中如果遇到非标准形式的序链，盲目展开只会出错，甚至被对方抓住漏洞反击。下面我就给大家详细讲解定义，结合我们刚才练过的例子，大家务必牢记。''

他抬手凝聚序链，将标准形式的三条定义清晰展现在空气中，逐条讲解：

\begin{quote}
\textbf{定义} 序元Veblen函数的标准形式，共有三条核心规定，每条都和我们实战息息相关，听仔细了。
\end{quote}

\begin{quote}
\allowbreak{}\NovelInline{I01607}{0}\allowbreak{} 是标准形式。这是最基础、最直观的一条，大家可以把 \allowbreak{}\NovelInline{I01608}{0}\allowbreak{} 理解为序链的''起点''，所有高阶序数的展开，最终都能追溯到 \allowbreak{}\NovelInline{I01609}{0}\allowbreak{}，比如我们之前讲的 \allowbreak{}\NovelInline{I01610}{\varphi(1 @1)[0] = 0}\allowbreak{}，这里的 \allowbreak{}\NovelInline{I01611}{0}\allowbreak{} 就是标准形式，也是序链展开的初始项，实战中不会出错。
\end{quote}

\begin{quote}
一个序数是由若干个 \allowbreak{}\NovelInline{I01612}{\varphi}\allowbreak{} 函数的加和构成的，如果它们都是标准形式，而且前者不小于后者，那么这个序数是标准形式。

举个通俗的例子，比如 \allowbreak{}\NovelInline{I01613}{\varphi(2 @\omega) + \varphi(1 @\omega)}\allowbreak{}，这里 \allowbreak{}\NovelInline{I01614}{\varphi(2 @\omega)}\allowbreak{} 和 \allowbreak{}\NovelInline{I01615}{\varphi(1 @\omega)}\allowbreak{} 都是我们之前讲过的标准形式，而且 \allowbreak{}\NovelInline{I01616}{\varphi(2 @\omega) \geq \varphi(1 @\omega)}\allowbreak{}，所以这个加和序数就是标准形式；但如果是 \allowbreak{}\NovelInline{I01617}{\varphi(1 @\omega) + \varphi(2 @\omega)}\allowbreak{}，前者小于后者，就不是标准形式，实战中遇到这种序链，说明对方序链构造有问题，我们可以直接针对性破解。
\end{quote}

\begin{quote}
如果 \allowbreak{}\NovelInline{I01618}{\alpha_1, \beta_1, \dots, \alpha_n, \beta_n}\allowbreak{} 都是标准形式，而且 \allowbreak{}\NovelInline{I01619}{\varphi(\alpha_n @\beta_n, \dots, \alpha_1 @\beta_1)}\allowbreak{} 大于每一个 \allowbreak{}\NovelInline{I01620}{\alpha_i}\allowbreak{} 和 \allowbreak{}\NovelInline{I01621}{\beta_i}\allowbreak{}，那么这个 \allowbreak{}\NovelInline{I01622}{\varphi}\allowbreak{} 函数就是标准形式。

这一条是实战中判断高阶序链是否合法的核心，大家结合我们练过的例子来理解------比如 \allowbreak{}\NovelInline{I01623}{\varphi(1 @\omega,\, 1 @0)}\allowbreak{}，这里 \allowbreak{}\NovelInline{I01624}{\alpha_1 = 1}\allowbreak{}、\allowbreak{}\NovelInline{I01625}{\beta_1 = 0}\allowbreak{}，都是标准形式，而且 \allowbreak{}\NovelInline{I01626}{\varphi(1 @\omega,\, 1 @0) > 1}\allowbreak{}、\allowbreak{}\NovelInline{I01627}{\varphi(1 @\omega,\, 1 @0) > 0}\allowbreak{}，所以它是标准形式；再比如 \allowbreak{}\NovelInline{I01628}{\varphi(\omega @\omega)}\allowbreak{}，\allowbreak{}\NovelInline{I01629}{\alpha_n = \omega}\allowbreak{}、\allowbreak{}\NovelInline{I01630}{\beta_n = \omega}\allowbreak{}，都是标准形式，且 \allowbreak{}\NovelInline{I01631}{\varphi(\omega @\omega) > \omega}\allowbreak{}，因此也是标准形式。
\end{quote}

弗兰德刻意停顿，观察队员们的反应，见大家都在认真记录，继续补充：

\begin{quote}
大家要注意，这条规则的关键是''\allowbreak{}\NovelInline{I01632}{\varphi}\allowbreak{} 函数值大于所有的 \allowbreak{}\NovelInline{I01633}{\alpha_i}\allowbreak{} 和 \allowbreak{}\NovelInline{I01634}{\beta_i}\allowbreak{}``，如果不满足这个条件，就不是标准形式。比如假设存在 \allowbreak{}\NovelInline{I01635}{\varphi(1 @\omega,\, 2 @0)}\allowbreak{}，其中 \allowbreak{}\NovelInline{I01636}{\alpha_1 = 2}\allowbreak{}、\allowbreak{}\NovelInline{I01637}{\beta_1 = 0}\allowbreak{}，若 \allowbreak{}\NovelInline{I01638}{\varphi(1 @\omega,\, 2 @0) < 2}\allowbreak{}，那它就不是标准形式，实战中这种序链无法稳定存在，我们只需释放低阶序链就能破解。
\end{quote}

弗兰德语气加重，做出总结：

\begin{quote}
标准形式的作用就是''筛选合法序链''。实战中，我们首先要判断对方释放的序链是否为标准形式，若是，再按之前讲的基本列规则展开破解；若不是，直接抓住其构造漏洞反击即可。掌握了标准形式，我们才能避免''无效展开''，提升实战效率。
\end{quote}

他抬手收起序链，衔接之前的内容：``清楚了标准形式的定义，接下来我们就结合这个定义，分 \allowbreak{}\NovelInline{I01639}{\alpha}\allowbreak{}、\allowbreak{}\NovelInline{I01640}{\beta}\allowbreak{} 分别为后继序数和极限序数的四种情况，进一步细化基本列展开规则，把每一种情况的展开逻辑吃透，确保实战中万无一失。''

``我们来做两个练习，巩固一下例子和规则。''玉小刚抬手，在空中写下第一个练习：

\begin{quote}
\textbf{练习1}：\allowbreak{}\NovelInline{I01641}{\varphi(1 @\omega,\, 1 @0)}\allowbreak{} 的基本列是什么？谁来回答？
\end{quote}

小舞率先举手，语气流畅：

\begin{quote}
\allowbreak{}\NovelInline{I01642}{\varphi(1 @\omega,\, 1 @0)}\allowbreak{} 中，\allowbreak{}\NovelInline{I01643}{\alpha = 0+1}\allowbreak{}（后继），\allowbreak{}\NovelInline{I01644}{\beta = \omega}\allowbreak{}（极限），\allowbreak{}\NovelInline{I01645}{\gamma = 0+1}\allowbreak{}（后继），对应第7条规则的延伸------

\NovelDisplay{M00194R01}{%
\varphi(\#,\, \alpha+1 @\beta,\, \gamma+1 @0)[n] = \varphi(\#,\, \alpha @\beta,\, (\varphi(\#,\, \alpha+1 @\beta,\, \gamma @0)+1) @\beta[n])%
}

这里 \allowbreak{}\NovelInline{I01646}{\alpha=0}\allowbreak{}，\allowbreak{}\NovelInline{I01647}{\beta=\omega}\allowbreak{}，\allowbreak{}\NovelInline{I01648}{\gamma=0}\allowbreak{}，所以基本列是：

\NovelDisplay{M00195R01}{%
\varphi(1 @\omega,\, 1 @0)[n] = \varphi(0 @\omega,\, (\varphi(1 @\omega,\, 0 @0)+1) @n) = \varphi(\text{SVO}+1 @n)%
}

和老师之前讲的例子一致！
\end{quote}

``完全正确。''弗兰德赞许道，再出一个难点的：

\begin{quote}
\textbf{练习2}：\allowbreak{}\NovelInline{I01649}{\varphi(\omega @\omega,\, 1 @0)}\allowbreak{} 的基本列是什么？
\end{quote}

牢大沉吟片刻，缓缓开口：

\begin{quote}
\allowbreak{}\NovelInline{I01650}{\varphi(\omega @\omega,\, 1 @0)}\allowbreak{} 中，\allowbreak{}\NovelInline{I01651}{\alpha = \omega}\allowbreak{}（极限），\allowbreak{}\NovelInline{I01652}{\beta = \omega}\allowbreak{}（极限），\allowbreak{}\NovelInline{I01653}{\gamma = 0+1}\allowbreak{}（后继），对应第8条规则的延伸------

\NovelDisplay{M00196R01}{%
\varphi(\#,\, \alpha @\beta,\, \gamma+1 @0)[n] = \varphi(\#,\, \alpha[n] @\beta,\, (\varphi(\#,\, \alpha @\beta,\, \gamma @0)+1) @\beta[n])%
}

这里 \allowbreak{}\NovelInline{I01654}{\alpha[n] = n}\allowbreak{}，\allowbreak{}\NovelInline{I01655}{\beta[n] = n}\allowbreak{}，\allowbreak{}\NovelInline{I01656}{\gamma = 0}\allowbreak{}，所以基本列是：

\NovelDisplay{M00197R01}{%
\varphi(\omega @\omega,\, 1 @0)[n] = \varphi(n @\omega,\, \varphi(\omega @\omega)+1 @n)%
}

和 \allowbreak{}\NovelInline{I01657}{\varphi(1 @\omega,\, 1 @0)}\allowbreak{} 的逻辑一致，只是 \allowbreak{}\NovelInline{I01658}{\alpha}\allowbreak{} 和 \allowbreak{}\NovelInline{I01659}{\beta}\allowbreak{} 都是极限序数。
\end{quote}

``非常好，牢大的理解很透彻。''玉小刚点头，随即话锋一转，``光懂理论不够，实战中序链展开瞬息万变，我们现在就补充几个实战展开例子，大家现场演练，熟悉操作细节。''

\chapter{第二十二章：序元veblen函数和OCF的实战融合}\label{ux7b2cux4e8cux5341ux4e8cux7ae0ux5e8fux5143veblenux51fdux6570ux548cocfux7684ux5b9eux6218ux878dux5408}

玉小刚的话音刚落，修炼场上的气氛再次凝重起来。他走到场地中央，左手凝聚出一团莹白色的序元 \allowbreak{}\NovelInline{I01660}{\varphi}\allowbreak{} 函数光团，右手则托起一团暗金色的 \allowbreak{}\NovelInline{I01661}{\psi_\Omega}\allowbreak{} 函数光团。两股力量在空气中交织，发出低沉的嗡鸣。

``光懂理论不够，实战中序链展开瞬息万变。''玉小刚目光扫过众人，``我们现在要解决的核心问题，是如何将\textbf{序元Veblen函数}与\textbf{序数坍缩函数（OCF）}进行无缝对接。在之前的训练中，你们虽然掌握了各自的展开规则，但在高压对抗下，往往因为无法快速在两种体系间切换，导致反击慢半拍。今天，我们就把这两种体系彻底打通。''

弗兰德接过话头，语气严肃：``首先，为了方便实战交流，我们统一一个记号：在不引起歧义的情况下，\allowbreak{}\NovelInline{I01662}{\psi_\Omega}\allowbreak{} 可以简写为 \allowbreak{}\NovelInline{I01663}{\psi}\allowbreak{}。大家必须牢记以下这组核心对应关系，这是连接两个世界的桥梁。''

他抬手在空中写下第一组公式，光芒映照在众人脸上：

\begin{quote}
\NovelDisplay{M00198R01}{%
\psi(\Omega^{\Omega^\omega} \cdot 2) = \varphi(1 @\omega, 1 @0)%
} \NovelDisplay{M00199R01}{%
\psi(\Omega^{\Omega^\omega + 1}) = \varphi(1 @\omega, 1 @1)%
} \NovelDisplay{M00200R01}{%
\psi(\Omega^{\Omega^\omega + \Omega}) = \varphi(1 @\omega, 1 @2)%
} \NovelDisplay{M00201R01}{%
\psi(\Omega^{\Omega^\omega + \Omega^2}) = \varphi(1 @\omega, 1 @3)%
} \NovelDisplay{M00202R01}{%
\psi(\Omega^{\Omega^\omega \cdot 2}) = \varphi(2 @\omega)%
}
\end{quote}

\begin{quote}
\NovelDisplay{M00203R01}{%
\psi(\Omega^{\Omega^\omega \cdot \psi(\Omega)}) = \varphi(\varphi(1 @1) @\omega)%
} \NovelDisplay{M00204R01}{%
\psi(\Omega^{\Omega^\omega \cdot \psi(\Omega^2)}) = \varphi(\varphi(1 @2) @\omega)%
} \NovelDisplay{M00205R01}{%
\psi(\Omega^{\Omega^\omega \cdot \psi(\Omega^\omega)}) = \varphi(\varphi(1 @\omega) @\omega)%
} \NovelDisplay{M00206R01}{%
\psi(\Omega^{\Omega^\omega \cdot \psi(\Omega^{\Omega^\omega \cdot \psi(\Omega^{\Omega^\omega})})}) = \varphi(\varphi(\varphi(1 @\omega) @\omega) @\omega)%
}
\end{quote}

``直到这里，''弗兰德写下最后一个长公式，``我们看到了三重嵌套的威力。这种序数在实战中几乎不可直接计算，必须依赖 \allowbreak{}\NovelInline{I01664}{\psi}\allowbreak{} 函数的坍缩特性来快速定位其量级。''

\begin{quote}
\NovelDisplay{M00207R01}{%
\psi(\Omega^{\Omega^{\omega+1}}) = \psi(\Omega^{\Omega^\omega \cdot \Omega}) = \varphi(1 @(\omega + 1))%
}
\end{quote}

``接下来有：'' \textgreater{} \NovelDisplay{M00208R01}{%
\psi(\Omega^{\Omega^{\omega+1}} \cdot 2) = \varphi(1 @(\omega + 1), 1 @0)%
} \textgreater{} \NovelDisplay{M00209R01}{%
\psi(\Omega^{\Omega^{\omega+1} + \Omega^\omega}) = \varphi(1 @(\omega + 1), 1 @\omega)%
}

``这意味着，当我们跨越了 \allowbreak{}\NovelInline{I01665}{\omega}\allowbreak{} 的某个界限后，整个体系会进入一个新的循环，但此时的基准线已经大幅提升。''

``大家注意观察，''弗兰德指着 \allowbreak{}\NovelInline{I01666}{\varphi(1 @\omega, 1 @0)}\allowbreak{} 到 \allowbreak{}\NovelInline{I01667}{\varphi(2 @\omega)}\allowbreak{} 的变化，``这里展示了如何利用 \allowbreak{}\NovelInline{I01668}{\Omega}\allowbreak{} 的不同幂次来驱动 \allowbreak{}\NovelInline{I01669}{\varphi}\allowbreak{} 函数第二个参数的增长。\allowbreak{}\NovelInline{I01670}{\Omega^\omega \cdot 2}\allowbreak{} 是一个关键节点，它标志着从`单参数迭代'向`双参数协同'的质变。'' \textgreater{} \NovelDisplay{M00210R01}{%
\psi(\Omega^{\Omega^{\omega+1}} \cdot 2) = \varphi(2 @(\omega + 1))%
} \textgreater{} \NovelDisplay{M00211R01}{%
\psi(\Omega^{\Omega^{\omega + 2}}) = \varphi(1 @(\omega + 2))%
} \textgreater{} \NovelDisplay{M00212R01}{%
\psi(\Omega^{\Omega^{\omega \cdot 2}}) = \varphi(1 @(\omega \cdot 2))%
}

``注意看指数部分的变化，''弗兰德指着公式解释道，``\allowbreak{}\NovelInline{I01671}{\Omega}\allowbreak{} 的指数结构直接对应了 \allowbreak{}\NovelInline{I01672}{\varphi}\allowbreak{} 函数中 \texttt{@} 符号后的参数层级。\allowbreak{}\NovelInline{I01673}{\Omega^\omega}\allowbreak{} 对应的是 \allowbreak{}\NovelInline{I01674}{\omega}\allowbreak{} 层级的迭代，而系数的变化则对应了参数的加法或乘法迭代。''

\begin{quote}
\NovelDisplay{M00213R01}{%
\psi(\Omega^{\Omega^{\psi(\Omega)}}) = \varphi(1 @(\varphi(2 @1)))%
} \NovelDisplay{M00214R01}{%
\psi(\Omega^{\Omega^{\psi(\Omega^2)}}) = \varphi(1 @(\varphi(3 @2)))%
} \NovelDisplay{M00215R01}{%
\psi(\Omega^{\Omega^{\psi(\Omega^\omega)}}) = \varphi(1 @(\varphi(1 @\omega)))%
} \NovelDisplay{M00216R01}{%
\psi(\Omega^{\Omega^{\psi(\Omega^{\Omega^\omega})}}) = \varphi(1 @(1 @\varphi(1 @\omega)))%
}
\end{quote}

``这一组是 nested（嵌套）结构，''玉小刚补充道，``当 \allowbreak{}\NovelInline{I01675}{\psi}\allowbreak{} 函数内部出现递归调用时，对应到 \allowbreak{}\NovelInline{I01676}{\varphi}\allowbreak{} 函数就是参数的层层嵌套。比如 \allowbreak{}\NovelInline{I01677}{\psi(\Omega^{\Omega^{\psi(\Omega^\omega)}})}\allowbreak{}，最内层的 \allowbreak{}\NovelInline{I01678}{\psi(\Omega^\omega)}\allowbreak{} 对应 \allowbreak{}\NovelInline{I01679}{\varphi(1 @\omega)}\allowbreak{}，也就是SVO；外层的结构则将其作为参数再次代入。这种结构在防御高阶精神攻击时尤为重要。''

\begin{quote}
\NovelDisplay{M00217R01}{%
\psi(\Omega^{\Omega^\Omega}) = \varphi(1 @(1, 0)) = \text{LVO}%
}
\end{quote}

``至此，我们达到了多元Veblen函数的极限------大Veblen序数（LVO）。''弗兰德的声音提高了几分，``这是目前你们能触及的理论天花板。记住这个等式，它是检验你们修为是否圆满的标尺。''

``理论讲完了，现在进入实战融合阶段。''玉小刚一挥手，修炼场中央升起四座模拟对战台，``我们将分为两组，进行`双体系即时翻译'训练。要求是：当对手释放 \allowbreak{}\NovelInline{I01680}{\varphi}\allowbreak{} 序链时，你们必须在0.5秒内用 \allowbreak{}\NovelInline{I01681}{\psi}\allowbreak{} 函数构建防御；反之亦然。''

\textbf{场景模拟}：Madore宗弟子释放 \allowbreak{}\NovelInline{I01682}{\varphi(1 @(\omega + 1))}\allowbreak{} 攻击。

唐三眼神一凝，大脑飞速运转：``\allowbreak{}\NovelInline{I01683}{\varphi(1 @(\omega + 1))}\allowbreak{}，根据对应表，这是 \allowbreak{}\NovelInline{I01684}{\psi(\Omega^{\Omega^{\omega+1}})}\allowbreak{} 的量级！小舞，准备柔骨锁，目标锁定 \allowbreak{}\NovelInline{I01685}{\Omega^{\Omega^{\omega+1}}}\allowbreak{} 的坍缩点！''

小舞身形一闪，粉色光芒瞬间化作一道精密的锁链，精准地缠绕在虚拟敌人的序力节点上：``收到！\allowbreak{}\NovelInline{I01686}{\psi}\allowbreak{} 函数已激活，防御屏障强度匹配 \allowbreak{}\NovelInline{I01687}{\Omega^{\Omega^{\omega+1}}}\allowbreak{}，对方序链被成功中和！''

\textbf{玉小刚点评}：``反应速度极快。唐三对 \allowbreak{}\NovelInline{I01688}{\varphi(1 @(\omega + 1))}\allowbreak{} 到 \allowbreak{}\NovelInline{I01689}{\psi}\allowbreak{} 的映射已经形成了肌肉记忆。注意，如果是 \allowbreak{}\NovelInline{I01690}{\varphi(1 @(\omega + 1), 1 @0)}\allowbreak{}，对应的就是 \allowbreak{}\NovelInline{I01691}{\psi(\Omega^{\Omega^{\omega+1}} \cdot 2)}\allowbreak{}，防御强度需要翻倍，下次注意这个系数变化。''

\textbf{场景模拟}：Buchholz宗弟子释放 \allowbreak{}\NovelInline{I01692}{\varphi(2 @\omega)}\allowbreak{} 压制场域。

奥斯卡大喝一声：``诺当，辅助序力输出！目标序数 \allowbreak{}\NovelInline{I01693}{\varphi(2 @\omega)}\allowbreak{}，对应 \allowbreak{}\NovelInline{I01694}{\psi(\Omega^{\Omega^\omega \cdot 2})}\allowbreak{}！我需要你把这个量级的序力放大三倍！''

诺当手中光芒大作，周身泛起淡蓝色序力光环：``辅助解析完成，正在将 \allowbreak{}\NovelInline{I01695}{\Omega^{\Omega^\omega \cdot 2}}\allowbreak{} 的坍缩能转化为增幅光环！''

随着光环笼罩，原本沉重的压制感瞬间减轻。奥斯卡趁机抛出一根镜像肠，上面流转着 \allowbreak{}\NovelInline{I01696}{\psi(\Omega^{\Omega^\omega \cdot 2})}\allowbreak{} 的金色纹路，瞬间复制了对方的压制力并反弹回去。

\textbf{弗兰德点评}：``配合默契。奥斯卡没有试图直接用 \allowbreak{}\NovelInline{I01697}{\varphi}\allowbreak{} 函数去硬抗，而是利用 \allowbreak{}\NovelInline{I01698}{\psi}\allowbreak{} 函数的`坍缩'特性，将庞大的 \allowbreak{}\NovelInline{I01699}{\varphi(2 @\omega)}\allowbreak{} 压缩成可操作的能量块，这是非常高明的战术。''

\textbf{场景模拟}：面对 \allowbreak{}\NovelInline{I01700}{\varphi(\varphi(1 @\omega) @\omega)}\allowbreak{} 这种三重嵌套的高阶防御。

牢大武魂附体，周身金光暴涨，作为队长展现出强大的统率力：``须坤，这是 \allowbreak{}\NovelInline{I01701}{\psi(\Omega^{\Omega^\omega \cdot \psi(\Omega^\omega)})}\allowbreak{}！它的防御核心在于中间的 \allowbreak{}\NovelInline{I01702}{\psi(\Omega^\omega)}\allowbreak{}，也就是SVO！只要击碎这个内核，外层就会崩塌！''

蔡须坤身形化作一道黑色闪电，动作灵动而迅猛：``明白！突袭突刺，目标 \allowbreak{}\NovelInline{I01703}{\psi(\Omega^\omega)}\allowbreak{} 奇点！''

只听``轰''的一声巨响，虚拟防御盾片碎裂。牢大紧随其后，烈光波轰出，直接将剩余序力轰散。

\textbf{玉小刚点评}：``这就是`擒贼先擒王'在序数战中的应用。面对复杂的嵌套结构，不要试图全面展开，利用 \allowbreak{}\NovelInline{I01704}{\psi}\allowbreak{} 函数找到最内层的关键不动点（如SVO），集中火力突破，效率最高。''

\textbf{场景模拟}：独自面对 \allowbreak{}\NovelInline{I01705}{\varphi(1 @\omega, 1 @0)}\allowbreak{} 的持续压制。

卢本韦深吸一口气，额头渗出汗水，但他眼神坚定：``\allowbreak{}\NovelInline{I01706}{\varphi(1 @\omega, 1 @0)}\allowbreak{} 对应 \allowbreak{}\NovelInline{I01707}{\psi(\Omega^{\Omega^\omega} \cdot 2)}\allowbreak{}\ldots{} 稳住，不能乱！

他双手结印，强行构建出 \allowbreak{}\NovelInline{I01708}{\psi}\allowbreak{} 函数的防御模型，虽然序力波动剧烈，但始终没有崩溃。

\textbf{弗兰德点评}：``卢本韦的表现值得肯定。在单人面对双参数高阶序数时，能准确调用 \allowbreak{}\NovelInline{I01709}{\psi}\allowbreak{} 函数进行抗压，说明基础已经打牢。接下来要练习的是如何在抗压的同时寻找反击机会。''

夕阳西下，修炼场上的光芒逐渐收敛。队员们虽然大汗淋漓，但眼神中却充满了自信。

``今天的训练非常成功。''弗兰德看着众人，欣慰地点头，``你们已经开始学会像呼吸一样自然地使用 \allowbreak{}\NovelInline{I01710}{\psi}\allowbreak{} 和 \allowbreak{}\NovelInline{I01711}{\varphi}\allowbreak{} 两种语言。记住，\allowbreak{}\NovelInline{I01712}{\psi(\Omega^{\Omega^\Omega}) = \text{LVO}}\allowbreak{} 只是开始，在那之后还有 \allowbreak{}\NovelInline{I01713}{\varphi(1 @\omega, 1 @0)}\allowbreak{} 乃至更广阔的天地。''

玉小刚合上典籍，语重心长地说：``接下来的日子，我们要把这些对应关系刻进骨子里。复赛在即，Madore宗和Buchholz宗的高手绝不会给我们思考的时间。只有做到`见招拆招，瞬发瞬解'，我们才能在这场精英大赛中走得更远。''

众人齐声应诺，声音响彻修炼场。果糕七怪战队的序数修为，在这一刻，真正迈入了一个全新的境界。

\chapter{第二十三章：半决赛，鏖战Madore七人组}\label{ux7b2cux4e8cux5341ux4e09ux7ae0ux534aux51b3ux8d5bux93d6ux6218madoreux4e03ux4ebaux7ec4}

序数精英大赛半决赛赛场，穹顶的高阶序数光纹如沸腾星河般剧烈震颤，连空气都因两股极致序数力量的对峙而微微扭曲。当Madore宗七名精英弟子缓步登场时，全场倒吸一口凉气------这支全员精通坍缩序数的战队，是夺冠头号热门，每一人的序数境界都足以碾压常规对手，七人并肩而立时，威压如泰山压顶，压得赛场地面的序数纹路微微凹陷，连观众席的呐喊都瞬间噤声。

\begin{itemize}
\tightlist
\item
  \textbf{队长·赛伊德}：银发如霜，墨眸深寒，周身暗紫色坍缩序链如活物般翻涌，指尖每一次轻动，都牵引着虚空微微震颤，执掌 \allowbreak{}\NovelInline{I01714}{\psi(\Omega^{\Omega^{\omega^2}})}\allowbreak{} 巅峰序数力量，举手投足间皆有空间坍缩的威能；
\item
  \textbf{副队长·雷斯}：魁梧如铁塔，银灰刚猛序力在周身奔涌咆哮，拳面萦绕螺旋状的坍缩序纹，掌控对标 \allowbreak{}\NovelInline{I01715}{\varphi(2 @\omega)}\allowbreak{} 的 \allowbreak{}\NovelInline{I01716}{\psi(\Omega^{(\Omega^\omega) \cdot 2})}\allowbreak{} 强攻之力，每一寸肌肉都透着毁天灭地的爆发力；
\item
  \textbf{Levi}：队内核心强攻手，境界直达 \allowbreak{}\NovelInline{I01717}{\psi(\Omega^{\Omega^3})}\allowbreak{} 高阶坍缩序数，指尖一动便引动天地雷威；
\item
  \textbf{Tonelli}：筑固态序数壁垒，能扛住高阶序数攻击；
\item
  \textbf{Kael}：施瞬发高速斩击，速度快至留下残影；
\item
  \textbf{Maren}：放序力干扰，能打乱敌方序力流转；
\item
  \textbf{Voss}：叠坍缩增幅，可瞬间拔高队友序数层级。
\end{itemize}

七道序力彼此共鸣、层层嵌套，织成密不透风的绝杀网，连空气都被压得凝滞。

\begin{itemize}
\tightlist
\item
  \textbf{牢大}：指尖飞速勾勒序元纹路，脑海中疯狂推演序数对应逻辑，心底早已绷紧了弦------他清楚，Madore宗的每一人都不是善茬，稍有不慎便会万劫不复；
\item
  \textbf{小舞}：粉色序力蓄势待发，眼神如鹰隼般紧盯敌方动向，指尖微微蜷缩，随时准备发起致命突袭；
\item
  \textbf{奥斯卡}：指尖金色辅助序纹微微闪烁，暗自警惕自身成为敌方突破口，心底反复告诫自己必须稳住全队序力；
\item
  \textbf{蔡徐坤、唐三、诺当、卢本韦}：屏息凝神，序力悄然流转。
\end{itemize}

这场七人对七人的巅峰序数对决，从开局就弥漫着生死一线的窒息感。

``不堪一击，全员出手，直接碾压！''赛伊德眼神骤冷，嘴角勾起一抹傲慢的弧度，根本不给果糕七怪任何调整与思考的机会，率先催动全身坍缩序力，发起绝杀攻势。

只见他双臂缓缓抬起，双手结出繁复的Madore宗专属印诀，周身暗紫色序链疯狂汇聚、缠绕，\allowbreak{}\NovelInline{I01718}{\psi(\Omega^{\Omega^{\omega^2}})}\allowbreak{} 坍缩纹路在半空具象为深渊漩涡，紫黑色烈焰裹挟着细碎的空间碎片倾泻而下，所过之处，空气被灼烧成扭曲的虚影，连序数光纹都被高温电离出焦糊的气息。这便是他的终极杀招\textbf{死亡火场}------火焰不仅能焚尽肉身，更能直接吞噬序数结构，从根源瓦解敌方防御，让其序力彻底溃散。

同一瞬，雷斯双脚狠狠跺向赛场地面，整块青石地面瞬间崩裂出蛛网般的裂痕，全身银灰 \allowbreak{}\NovelInline{I01719}{\psi(\Omega^{(\Omega^\omega) \cdot 2})}\allowbreak{} 序力尽数灌注双拳，金色强攻序纹缠绕拳面，肌肉贲张，猛然向前拍出：

\begin{quote}
\textbf{``Madore宗强攻绝技------大力金刚掌！''}
\end{quote}

遮天蔽日的金色巨掌轰然压下，掌风震出层层虚空涟漪，每一寸掌面都透着碾碎一切序数防御的威势，精准封死了果糕七怪所有闪避与突围的路线。

Madore剩余五人同步联动，序力彼此共鸣，发动协同攻击：

\begin{itemize}
\tightlist
\item
  \textbf{Levi}：指尖幽绿雷芒暴涨，依托 \allowbreak{}\NovelInline{I01720}{\psi(\Omega^{\Omega^3})}\allowbreak{} 高阶坍缩序力，指尖凝聚数十道幽绿雷刃，雷刃旋转飞舞，裹挟着刺耳的破空声，朝着果糕七怪周身射去，雷刃所过之处，空气炸裂，留下一道道焦黑的痕迹；
\item
  \textbf{Tonelli}：周身土金序数壁垒瞬间扩张，如铜墙铁壁般笼罩赛场左侧，彻底锁死果糕七怪的退路，壁垒表面流转着厚重的坍缩序纹，连高阶攻击都难以撼动；
\item
  \textbf{Kael}：身形化作数道残影，高速穿梭于雷刃与掌风之间，瞬发高速斩击四面合围，每一道残影都携着 \allowbreak{}\NovelInline{I01721}{\psi(\Omega^n)}\allowbreak{} 层级的序力，斩向七人的破绽之处；
\item
  \textbf{Maren}：指尖轻弹，无数淡灰色序力干扰丝飘出，缠上果糕七怪的序力流转路径，打乱其序力节奏；
\item
  \textbf{Voss}：抬手将自身坍缩序力汇入赛伊德与雷斯的技能之中，瞬间拔高二者技能威力，让死亡火场与大力金刚掌的威势再增三分。
\end{itemize}

七道序数攻击交织成密不透风的绝杀网，将果糕七怪死死笼罩，威压之强让七人周身序力都开始微微紊乱。

``快筑防！严格对位！''唐三仓促嘶吼，可情急之下的致命失误瞬间浮现------他竟未用高阶 \allowbreak{}\NovelInline{I01722}{\varphi}\allowbreak{} 序数对位，仅以基础 \allowbreak{}\NovelInline{I01723}{\varphi(1 @\omega)}\allowbreak{} 序链仓促凝盾，全然忘了层级差距带来的致命碾压。

蓝银色蓝银序力藤蔓仓促蔓延，交织成简易的防御囚笼，可刚一触碰到赛伊德 \allowbreak{}\NovelInline{I01724}{\psi(\Omega^{\Omega^{\omega^2}})}\allowbreak{} 的死亡火场，淡金色 \allowbreak{}\NovelInline{I01725}{\varphi}\allowbreak{} 纹就被暗紫色坍缩纹疯狂吞噬、撕裂，蓝银藤蔓以肉眼可见的速度枯萎、碎裂，连唐三的指尖都被震得发麻，心底咯噔一沉，深知这波防御根本不堪一击。

小舞见状，眼神一厉，纵身跃起，粉色序力凝聚成数枚尖锐的追踪导弹，弹头刚浮现出粉色序纹，就朝着雷斯的大力金刚掌射去，试图拆解掌风。可导弹刚触碰到掌面的银灰坍缩纹，就被掌劲震成齑粉，连近身的机会都没有，小舞心头一紧，暗叫不好。

蔡徐坤仓促掷出回旋篮球，篮球表面的基础序力瞬间被坍缩纹瓦解，直飞而出，竟直直撞向Levi的幽绿雷刃，被雷刃轻易劈成两半，碎片四散飞溅，险些波及自身。

\textbf{瞬息之间，果糕七怪的防御全线崩溃。}

七人被两股超强高阶序数力量的余波狠狠震飞，嘴角溢出鲜血，重重砸在赛场角落的地面上，骨骼作响，序力紊乱不堪，连起身都显得艰难。

唐三撑着颤抖的手臂起身，紧盯着逼近的紫黑火海与金色巨掌，心底一片冰凉------他从未想过，自己精心推演的序数对应，竟会因一时慌乱出现如此致命的偏差，全队竟会瞬间陷入必死绝境。他咬着牙，指尖再次勾勒序纹，却发现序力流转滞涩，根本无法凝聚有效防御。

小舞咬牙撑起身，粉色序力微弱流转，满心焦灼却无力回天，看着越来越近的火海，心底掠过一丝绝望，却依旧死死盯着敌方，不肯放弃。

奥斯卡蜷缩在赛场角落，指尖金色辅助序纹黯淡，看着周身紊乱的序力，又望向逼近的攻击，心底满是自责------若不是自己成为敌方重点针对的目标，全队也不会陷入如此困境，他攥紧拳头，指甲嵌进掌心，疼得眼前发黑，却依旧强撑着序力，准备随时做出最后辅助。

牢大攥紧双拳，指节泛白，冰蓝序力在周身微弱流转，却根本无法抵挡高阶坍缩序力，他看着赛伊德那副傲慢的嘴脸，心底满是不甘与愤怒------难道就要这样惨败于此？他不甘心！

诺当快速调整序力，粉色电磁序力在掌心微微跳动，脑中飞速思索破局之法，可面对七人的协同攻击，他竟一时想不出任何应对之策，额头渗出冷汗，心底满是焦灼。

``就这点序元掌控力，也配与Madore宗为敌？''赛伊德冷笑一声，声音冰冷刺骨，死亡火场的威势再度暴涨，紫黑色火焰如毒蛇般缠绕，步步紧逼；雷斯双拳再次蓄力，金色掌印重新凝聚，掌风之中透着终结一切的威压，眼看就要彻底碾碎果糕七怪的最后一丝抵抗。

``全员冷静！摒弃慌乱！严格对位------''

唐三强压体内翻涌的气血与喉间的腥甜，嘶吼着传递精准的序数对应逻辑，每一个字都清晰传入队友耳中，同时用蓝银草轻轻缠绕住身边的队友，传递着''绝不放弃''的信念：

\begin{quote}
\begin{itemize}
\tightlist
\item
  \textbf{赛伊德} \allowbreak{}\NovelInline{I01726}{\psi(\Omega^{\Omega^{\omega^2}})}\allowbreak{} → 用 \allowbreak{}\NovelInline{I01727}{\varphi(1 @\omega^2)}\allowbreak{} 精准对位拆解！
\item
  \textbf{雷斯} \allowbreak{}\NovelInline{I01728}{\psi(\Omega^{(\Omega^\omega) \cdot 2})}\allowbreak{} → 用 \allowbreak{}\NovelInline{I01729}{\varphi(2 @\omega)}\allowbreak{} 正面抗衡！
\item
  \textbf{Levi} \allowbreak{}\NovelInline{I01730}{\psi(\Omega^{\Omega^3})}\allowbreak{} → 用 \allowbreak{}\NovelInline{I01731}{\varphi(\omega @4)}\allowbreak{} 牵制！
\item
  \textbf{Tonelli、Kael、Maren、Voss} → 由诺当、小舞、蔡徐坤、牢大分头牵制！
\end{itemize}

\textbf{层级不匹配，必被碾压！}
\end{quote}

这一瞬，果糕七怪的心神瞬间安定下来------唐三的冷静分析，是他们唯一的生机。

奥斯卡瞬间会意，双手飞速拨动金色辅助序纹，周身金色辅助序纹暴涨，璀璨的金环瞬间笼罩全队，开启全域增幅：

\begin{quote}
\textbf{``序力增幅·速启·稳序！全员序力流转速度提升50\%，稳固自身序数构造，屏蔽敌方序力干扰！''}
\end{quote}

金环之中，细密的金色序纹流转，果糕七怪紊乱的序力瞬间变得平稳流畅，滞涩的高阶序数展开也不再卡顿，唐三的 \allowbreak{}\NovelInline{I01732}{\varphi(1 @\omega^2)}\allowbreak{} 序纹、蔡徐坤的 \allowbreak{}\NovelInline{I01733}{\varphi(2 @\omega)}\allowbreak{} 序纹、诺当的 \allowbreak{}\NovelInline{I01734}{\varphi(3 @\omega)}\allowbreak{} 序纹，都开始清晰浮现，为校准对位、释放技能争取到了宝贵的喘息之机。奥斯卡心底松了口气，却依旧不敢松懈------他知道，这只是转机，想要彻底扭转战局，还需全队齐心协力。

千钧一发之际，牢大猛地从怀中掏出冰蓝色的冰红茶，仰头一饮而尽，冰凉的液体顺着喉咙滑落的瞬间，全身冰蓝序力轰然爆发，周身冰雾弥漫，指尖扩散开一圈圈冰蓝色的魅惑序纹：

\begin{quote}
\textbf{``魅惑·寒魄迷神！''}
\end{quote}

柔和却极具穿透力的冰蓝波纹横扫Madore七人，瞬间干扰了他们的序数神经与动作节奏。雷斯蓄力的动作骤然迟缓，拳面的金色坍缩纹出现剧烈波动；赛伊德操控死亡火场的节奏被打乱，紫黑火焰的吞吐变得紊乱；Levi、Kael等人的反应速度、攻击速度大幅下降，原本密不透风的绝杀攻击网，瞬间出现了致命的破绽！牢大稳住呼吸，紧盯着敌方的破绽，心底清楚，这是全队唯一的转机，必须死死抓住。

``诺当，磁爆炸弹扰乱Tonelli与Voss的序力联动！小舞牵制Kael的高速斩击！蔡徐坤盯紧Levi的雷刃！牢大稳住魅惑效果！''唐三高声下令，同时双手按地，全身序元 \allowbreak{}\NovelInline{I01735}{\varphi}\allowbreak{} 之力尽数爆发，蓝银色蓝银序力裹挟着 \allowbreak{}\NovelInline{I01736}{\varphi(1 @\omega^2)}\allowbreak{} 的高阶金色序纹，疯狂蔓延、层层交织：

\begin{quote}
\textbf{``蓝银囚笼·\allowbreak{}\NovelInline{I01737}{\varphi(1 @\omega^2)}\allowbreak{} 专属序盾！''}
\end{quote}

巨型蓝银序盾瞬间矗立在果糕七怪身前，盾面布满规整的 \allowbreak{}\NovelInline{I01738}{\varphi(1 @\omega^2)}\allowbreak{} 高阶不动点序纹，淡金色的序纹与赛伊德袭来的暗紫色 \allowbreak{}\NovelInline{I01739}{\psi(\Omega^{\Omega^{\omega^2}})}\allowbreak{} 坍缩序纹狠狠碰撞在一起！两种高阶序数纹路瞬间疯狂对冲，发出刺耳的摩擦声，淡金色 \allowbreak{}\NovelInline{I01740}{\varphi}\allowbreak{} 纹牢牢抵住紫黑色坍缩纹，以同等的序数层级，稳稳中和、瓦解着死亡火场的吞噬之力，紫黑火焰再也无法前进一步。唐三紧盯着序盾表面的纹路变化，不敢有丝毫松懈，心底反复核对序数层级，生怕再次出现偏差，同时用蓝银草加固盾面，维持序盾稳定。

\begin{center}\rule{0.5\linewidth}{0.5pt}\end{center}

同一瞬，蔡徐坤旋身蓄力，左脚狠狠踏地稳住身形，全身序元 \allowbreak{}\NovelInline{I01741}{\varphi}\allowbreak{} 之力尽数灌注手中的橙色篮球，篮球表面飞速缠绕上 \allowbreak{}\NovelInline{I01742}{\varphi(2 @\omega)}\allowbreak{} 的螺旋金色序环，篮球旋转间，发出嗡鸣的序力波动。他猛然转身、全力掷出：

\begin{quote}
\textbf{``回旋篮球·\allowbreak{}\NovelInline{I01743}{\varphi(2 @\omega)}\allowbreak{} 序斩！''}
\end{quote}

篮球携着金色弧线划破长空，精准击中雷斯掌芯的 \allowbreak{}\NovelInline{I01744}{\psi(\Omega^{(\Omega^\omega) \cdot 2})}\allowbreak{} 序数核心，\allowbreak{}\NovelInline{I01745}{\varphi(2 @\omega)}\allowbreak{} 的金色序纹如锋利的刀刃，层层切割着 \allowbreak{}\NovelInline{I01746}{\psi}\allowbreak{} 坍缩掌纹，金色巨掌瞬间崩碎成无数序力碎片，碎片四散飞溅，却被蔡徐坤周身的序力屏障拦下。蔡徐坤掌心微麻，却眼神坚定，知道自己必须扛住正面强攻，为全队争取调整时间。

\begin{center}\rule{0.5\linewidth}{0.5pt}\end{center}

与此同时，诺当双手合于胸前，粉色电磁序力瞬间转化为蓝紫电弧光球，厉声暴喝：

\begin{quote}
\textbf{``磁爆炸弹·序链崩解！''}
\end{quote}

数十颗蓝紫色磁爆炸弹呼啸而出，精准落在Tonelli的土金壁垒与Voss的坍缩序纹间隙之中，下一秒轰然炸裂！狂暴的电磁风暴瞬间席卷开来，蓝紫色电弧疯狂跳跃、撕扯，狠狠干扰着Tonelli的固态壁垒构造与Voss的坍缩增幅序纹，土金壁垒表面出现细密的裂痕，Voss的序力流转瞬间紊乱，再也无法拔高雷斯与赛伊德的技能威力。诺当紧盯着爆炸后的效果，心底松了口气，却依旧快速补充序力，准备应对敌方的反扑。

\begin{center}\rule{0.5\linewidth}{0.5pt}\end{center}

小舞纵身跃起，粉色序力凝聚成数道\textbf{追踪导弹·序缠版}，导弹表面缠绕着淡粉色序纹，精准缠上Kael的高速残影，干扰其斩击轨迹，让Kael的斩击偏离目标，无法形成有效合围。小舞眼神锐利，紧盯着Kael的动向，心底清楚，必须牵制住这位高速强攻手，否则全队侧翼将彻底暴露。

\begin{center}\rule{0.5\linewidth}{0.5pt}\end{center}

牢大则持续催动\textbf{魅惑·寒魄迷神}，冰蓝魅惑序纹不断扩散，进一步削弱Madore七人的反应速度与序力掌控力，Levi的雷刃轨迹微微偏差，Tonelli的壁垒加固动作变得迟缓，赛伊德的坍缩序链也出现了短暂的紊乱。牢大咬着牙，维持着魅惑效果，冰蓝序力在周身不断消耗，却依旧不肯停歇------他知道，自己的魅惑是全队牵制敌方的关键，绝不能掉链子。

\begin{center}\rule{0.5\linewidth}{0.5pt}\end{center}

战局暂稳，观众席瞬间爆发出震耳欲聋的欢呼声，呐喊声浪直冲云霄！

可就在众人刚松一口气的瞬间，Madore剩余五人突然做出联动姿态，五人齐齐抬手，将自身全部序力顺着坍缩纹路，尽数汇入赛伊德与雷斯体内！

原本就强横的 \allowbreak{}\NovelInline{I01747}{\psi(\Omega^{\Omega^{\omega^2}})}\allowbreak{} 与 \allowbreak{}\NovelInline{I01748}{\psi(\Omega^{(\Omega^\omega) \cdot 2})}\allowbreak{} 坍缩序数，瞬间暴涨数倍，暗紫色、银灰色的坍缩序纹变得更加浓郁、更加狂暴，\textbf{直接超出了固定序数层级的阈值！}

唐三的 \allowbreak{}\NovelInline{I01749}{\varphi(1 @\omega^2)}\allowbreak{} 蓝银序盾、蔡徐坤的 \allowbreak{}\NovelInline{I01750}{\varphi(2 @\omega)}\allowbreak{} 回旋篮球，在两股暴涨后的坍缩序数威压下，如同遭遇飓风的薄纸，瞬间崩裂成漫天金色与蓝银色的序力碎片！

``不好！''唐三心底咯噔一沉，刚稳住的战局瞬间再度陷入危机，Madore七人组的威压如泰山压顶，让全队序力都开始微微紊乱。

``全员串联序力！用 \allowbreak{}\NovelInline{I01751}{\varphi(1 @\omega^3)}\allowbreak{} 高阶共振防御阵！''唐三仓促间嘶吼，指尖飞速勾勒出新的序元纹路，周身蓝银色序力裹挟着更强的波动，朝着队友们蔓延而去。

奥斯卡心领神会，双手飞速结印，周身金色辅助序纹暴涨，璀璨的金环再度扩张，笼罩住整个赛场角落：

\begin{quote}
\textbf{``序力共振·高阶防御！''}
\end{quote}

\begin{itemize}
\tightlist
\item
  \textbf{牢大}猛地将自身冰蓝序力全部催动，冰蓝魅惑序纹瞬间转化为厚重的冰盾序力，朝着Madore七人组的方向汇聚，形成一道坚不可摧的冰蓝色防御屏障；
\item
  \textbf{诺当}双手合十，粉色电磁序力转化为无数道蓝紫色电磁光束，交织成密集的干扰网，试图削弱坍缩序数的威压；
\item
  \textbf{小舞}纵身跃起，粉色序力凝聚成数道追踪导弹，朝着Madore七人组的技能落点射去，试图拆解其攻势；
\item
  \textbf{蔡徐坤}快速调整自身序力，将 \allowbreak{}\NovelInline{I01752}{\varphi(2 @\omega)}\allowbreak{} 回旋篮球与自身序力融合，化作一道金色的序力光矛，朝着赛伊德的死亡火场方向刺去，试图分散其注意力。
\end{itemize}

可这一切，在Madore七人组的高阶坍缩序数面前，都显得如此脆弱！

赛伊德嘴角勾起一抹傲慢的冷笑，双臂缓缓抬起，周身暗紫色坍缩序链如同活物般翻涌，\allowbreak{}\NovelInline{I01753}{\psi(\Omega^{\Omega^{\omega^2}})}\allowbreak{} 坍缩纹路在半空具象化为更加恐怖的火海：

\begin{quote}
\textbf{``死亡火场·高阶坍缩版！''}
\end{quote}

紫黑色火焰如同吞噬一切的巨兽，朝着果糕七怪席卷而来，所过之处，空气被灼烧成扭曲的虚影，连序数光纹都被高温电离出焦糊的气息，牢大的冰蓝色防御屏障瞬间融化，小舞的追踪导弹被轻易吞噬，蔡徐坤的金色序力光矛也被火焰包裹，变得黯淡无光。

雷斯双拳蓄力，金色掌印重新凝聚，掌风之中透着终结一切的威压：

\begin{quote}
\textbf{``大力金刚掌·高阶坍缩版！''}
\end{quote}

金色巨掌再次朝着果糕七怪压下，掌面的坍缩序纹变得更加浓郁、更加狂暴，直接超出了果糕七怪固定序数层级的阈值！

``坚持住！我们不能放弃！''唐三嘶吼着，周身蓝银色序力再度爆发，试图凝聚新的防御体系，可序力刚一接触到坍缩序数，就被瞬间吞噬、瓦解，连他的指尖都被震得发麻，心底咯噔一沉，深知这波攻击根本不堪抵挡。

就在这千钧一发之际，奥斯卡突然想起自身储存的 \allowbreak{}\NovelInline{I01754}{\varphi(3 @\omega^3)}\allowbreak{} 终极共振辅助序纹------那是他之前偷偷推演、从未使用的底牌，能够短暂融合全队序力，形成一道超越固定序数层级的防御体系！

``全员配合！催动 \allowbreak{}\NovelInline{I01755}{\varphi(3 @\omega^3)}\allowbreak{} 终极共振辅助序纹！''奥斯卡高声嘶吼，双手飞速拨动辅助序纹，周身金色辅助序纹变得更加璀璨，一道淡金色的序力屏障瞬间笼罩住果糕七怪周身，开始快速吸收、转化全队序力，朝着更高阶的序数层级凝聚。

唐三见状，眼中闪过一丝希望，立刻高声下令：

\begin{quote}
\begin{itemize}
\tightlist
\item
  \textbf{牢大}！你的冰蓝序力融入共振屏障，强化防御稳定性！
\item
  \textbf{小舞}！你的粉色序力补充序力流转速度！
\item
  \textbf{蔡徐坤}！你的金色序力攻击敌方技能破绽！
\item
  \textbf{诺当}！你的电磁序力干扰敌方序力联动！
\end{itemize}
\end{quote}

牢大、小舞、蔡徐坤、诺当纷纷应声，周身序力尽数汇入淡金色共振屏障之中，屏障光芒暴涨，原本黯淡的 \allowbreak{}\NovelInline{I01756}{\varphi(3 @\omega^3)}\allowbreak{} 序纹变得更加清晰、更加耀眼，一道超越固定序数层级的防御体系，终于在绝境之中凝聚成型！

\textbf{``终极共振防御·正式启动！''}

唐三高声喝出，淡金色共振屏障与Madore七人组的高阶坍缩序数技能狠狠碰撞在一起！两种远超常规序数层级的力量相互冲击、相互抵消，发出震耳欲聋的轰鸣声响，赛场地面的序数纹路被震得微微凹陷，连空气都因两股极致序数力量的对峙而微微扭曲。

Madore七人组见状，眼神骤冷，赛伊德再度催动坍缩序数，紫黑色火焰的威势再度暴涨：

\begin{quote}
\textbf{``死亡火场·终极坍缩！''}
\end{quote}

雷斯双拳紧握，金色掌印重新凝聚，掌风之中透着终结一切的威压：

\begin{quote}
\textbf{``大力金刚掌·终极坍缩！''}
\end{quote}

Levi、Kael等人也纷纷催动自身高阶坍缩序力，朝着果糕七怪周身射去，无数道坍缩序力交织成密不透风的绝杀网，将果糕七怪彻底笼罩！

果糕七怪的共振屏障开始微微震颤，淡金色的序纹变得愈发黯淡，牢大的冰蓝序力、小舞的粉色序力、蔡徐坤的金色序力、诺当的电磁序力，都因持续抵挡高阶坍缩序数而开始快速消耗，心底咯噔一沉，深知这样下去，迟早会被彻底击破！

``不能放弃！绝对不能放弃！''唐三咬紧牙关，指尖飞速勾勒出新的序元纹路，周身蓝银色序力再度爆发，``我们还有最后一张底牌------\allowbreak{}\NovelInline{I01757}{\psi(\Omega^{(\Omega^\omega) \cdot 3})}\allowbreak{} 序数协同攻击阵！''

奥斯卡立刻会意，双手飞速拨动辅助序纹，周身金色辅助序纹尽数汇入唐三体内：

\begin{quote}
\textbf{``序力协同·终极共振！''}
\end{quote}

牢大、小舞、蔡徐坤、诺当也纷纷将自身剩余序力全部汇入共振屏障之中，淡金色的共振屏障光芒再度暴涨，与唐三的蓝银色序力融合，化作一道更加恐怖的金色序力光柱，朝着Madore七人组的方向射去！

序力光柱表面，密密麻麻刻着果糕七怪的序数对应逻辑，每一道纹路都透着超越固定序数层级的威压，所过之处，空气被炸裂，序数光纹被灼烧成焦糊的气息，连Madore七人组的高阶坍缩序数技能，都开始出现微微动摇！

Madore七人组见状，脸色骤变，赛伊德急忙催动坍缩序数，试图抵挡这道金色序力光柱，可序力光柱刚一接触到紫黑色火焰，就被瞬间吞噬、瓦解，连火焰都开始出现细密的裂痕；雷斯双拳蓄力，金色掌印重新凝聚，却根本无法阻挡序力光柱的前进势头；Levi、Kael等人也纷纷催动自身坍缩序力，可依旧无法阻止序力光柱朝着自身逼近。

``完了！彻底完了！''观众席中，不少观众看着这道恐怖的金色序力光柱，心底满是绝望，他们知道，果糕七怪即将迎来最终的胜利！

可就在这时，Madore七人组突然做出联动姿态，五人齐齐抬手，将自身全部剩余序力顺着坍缩纹路，尽数汇入赛伊德与雷斯体内！

他们的力量暴涨到 \allowbreak{}\NovelInline{I01758}{\psi(\Omega^{\Omega^{\omega^\omega}})}\allowbreak{} 与 \allowbreak{}\NovelInline{I01759}{\psi(\Omega^{\Omega^{\omega \cdot 2}})}\allowbreak{} 坍缩序数，暗紫色、银灰色的坍缩序纹变得更加浓郁、更加狂暴，直接超出了 \allowbreak{}\NovelInline{I01760}{\varphi(3 @\omega^3)}\allowbreak{} 终极共振辅助序纹的层级阈值！

金色序力光柱刚一接触到暴涨后的坍缩序数，就被瞬间包裹、吞噬，淡金色的 \allowbreak{}\NovelInline{I01761}{\varphi}\allowbreak{} 纹被暗紫色的 \allowbreak{}\NovelInline{I01762}{\psi}\allowbreak{} 坍缩纹疯狂侵蚀，序力光柱表面的纹路开始快速模糊、碎裂，最终化作漫天序力碎片，消散在空气之中！

唐三、牢大、小舞、蔡徐坤、诺当等人被震得连连后退，嘴角溢出鲜血，周身序力紊乱不堪，可他们依旧死死盯着Madore七人组，眼神之中透着不屈的光芒------他们已经尽力了，可终究还是无法抵挡层级差距如此巨大的高阶坍缩序数！

``哈哈哈！就这点序元掌控力，也配与Madore宗为敌？''赛伊德冷笑一声，声音冰冷刺骨，死亡火场的威势再度暴涨，紫黑色火焰如同毒蛇般缠绕，步步紧逼；雷斯双拳再次蓄力，金色掌印重新凝聚，掌风之中透着终结一切的威压，眼看就要彻底碾碎果糕七怪的最后一丝抵抗。

``我们不会放弃！绝对不会放弃！''唐三猛地站起身，周身蓝银色序力再度爆发，指尖飞速勾勒出新的序元纹路，``Madore宗的坍缩序数虽强，但我们的序数信念更强！''

牢大、小舞、蔡徐坤、诺当等人也纷纷起身，周身序力开始快速恢复，眼神之中透着坚定的信念：

\begin{quote}
\textbf{``绝不放弃！''}
\end{quote}

果糕七怪的序数对应逻辑，开始以肉眼可见的速度变得清晰、变得坚定，原本紊乱的序力瞬间变得平稳流畅，滞涩的高阶序数展开也不再卡顿，一道超越 \allowbreak{}\NovelInline{I01763}{\varphi(1 @\varphi(1,0))}\allowbreak{} 终极共振辅助序纹、直指 \allowbreak{}\NovelInline{I01764}{\psi(\Omega^{\Omega^{\omega \cdot 2}})}\allowbreak{} 层级的序数协同攻击阵，终于在绝境之中凝聚成型！

\textbf{``序数协同攻击阵·终极信念版！''}

唐三高声喝出，果糕七怪的序数对应逻辑交织成密不透风的绝杀网，朝着Madore七人组的方向射去！这道攻击阵，不仅蕴含着全队超越固定序数层级的序力，更蕴含着他们永不放弃、坚守初心的序数信念！

Madore七人组见状，脸色终于出现了一丝慌乱，可他们别无选择，只能催动自身全部坍缩序数，试图抵挡这道攻击阵！

\textbf{``砰------！''}

两道远超常规序数层级的力量狠狠碰撞在一起！发出震耳欲聋的轰鸣声响，赛场地面的序数纹路被震得彻底碎裂，连空气都被炸裂成无数碎片，序数光纹被灼烧成焦糊的气息，漫天序力碎片四处飞溅，遮天蔽日！

不知过了多久，漫天序力碎片终于缓缓落下，赛场之中恢复了短暂的平静。

果糕七怪瘫坐在赛场角落，满身汗水与鲜血，周身序力紊乱不堪，可他们的眼神却无比坚定，嘴角勾起一抹胜利的笑容------\textbf{他们成功了！} 他们凭借着永不放弃的序数信念，打破了Madore七人组层级差距巨大的高阶坍缩序数，赢得了这场艰难的序数对决！

Madore七人组则面色惨白，瘫坐在不远处，周身坍缩序力消耗殆尽，眼神之中满是不甘与难以置信------他们明明拥有层级优势，明明拥有更强的坍缩序数，却最终还是输给了果糕七怪的序数信念！

观众席瞬间爆发出震耳欲聋的欢呼声，呐喊声浪直冲云霄！所有人都站起身，朝着果糕七怪疯狂呐喊，为他们永不放弃、坚守初心的序数精神喝彩！

唐三撑着颤抖的手臂起身，看着身边同样疲惫却笑容坚定的队友，心底满是温暖------他知道，这场序数对决的胜利，不仅仅是层级优势的胜利，更是序数信念的胜利！正是因为他们永不放弃、坚守初心，才最终打破了极限，赢得了比赛！

小舞走到唐三身边：

\begin{quote}
``老哥，我们赢了！我们真的赢了！''
\end{quote}

他声音沙哑却坚定：

\begin{quote}
``行了，老妹，我们还有更大的目标，接下来，我们要继续努力，解锁更多高阶序数技能，不辜负大家的期待！''
\end{quote}

牢大、蔡徐坤、诺当、奥斯卡等人也纷纷起身，围聚在一起，脸上满是笑容与疲惫，他们知道，这场序数对决的胜利，将成为他们序数修炼之路中，最珍贵、最难忘的一段回忆！

\chapter{第二十四章：暗影密谋，前路疑云}\label{ux7b2cux4e8cux5341ux56dbux7ae0ux6697ux5f71ux5bc6ux8c0bux524dux8defux7591ux4e91}

半决赛的硝烟尚未散尽，序数赛场的穹顶光纹仍在微微震颤，残留着两股高阶序数碰撞后的余波。果糕七怪满身汗渍与血痕，却难掩眼底的光芒，他们简单休整片刻，便踏着夕阳的余晖，朝着决赛场地出发------那是一座矗立在赛场中心的高阶序数殿堂，也是他们距离冠军荣耀最近的地方。

牢大走在最前方，指尖依旧习惯性地勾勒着序元纹路，脑海中反复推演着决赛可能遇到的对手，以及应对的序数对应逻辑。经过与Madore宗的死战，他愈发谨慎，周身蓝银色序力若有若无地流转，时刻警惕着周围的动静。

唐三紧随其后，蓝色序力轻轻萦绕在周身，时不时侧头看向牢大，眼底满是劫后余生的庆幸与对未来的期待。

奥斯卡走在队伍中间，指尖金色辅助序纹缓缓闪烁，默默为众人补充着消耗的序力，嘴角还带着一丝疲惫的笑意。

小舞、蔡须坤、诺当并肩而行，三人低声交谈着半决赛的得失，语气中既有胜利的喜悦，也有对决赛的敬畏。

与此同时，赛场之外的暗黑深渊之中，一座高耸入云的黑色塔楼正散发着诡异的光晕------那便是\textbf{帅气教}的巢穴，\textbf{帅气塔}。

塔楼通体由被掠夺的序数残骸堆砌而成，每一块砖石都封印着一位被击败的序数修炼者的序力精华，塔身表面流转着密密麻麻的银色掠夺序纹，如同无数条饥饿的毒蛇，贪婪地吞噬着周围一切游离的序数能量。塔楼周围百丈之内，寸草不生，连空气都被扭曲成暗紫色的漩涡，任何靠近的生物都会被瞬间抽干序力，化为一具枯骨。

塔楼顶端的议事大殿内，烛火摇曳，映得一道高大的身影愈发阴沉，他便是帅气教教主------\textbf{胡赞}。

胡赞身着玄色长袍，衣摆上绣着扭曲的银色序纹，那是帅气教专属的\textbf{掠夺序纹}，每一道纹路都透着贪婪与邪恶。他端坐于黑曜石王座之上，周身萦绕着浓郁的黑色序力，那序力中夹杂着无数细碎的、被掠夺来的帅气能量，隐隐透着令人心悸的压迫感。

他指尖轻叩王座扶手，眼神阴鸷地盯着身前的水镜，水镜之中，正是果糕七怪赶路的身影，以及远处扽魂殿决战团队的踪迹。

\begin{quote}
``果糕七怪\ldots\ldots 扽魂殿\ldots\ldots{}''胡赞的声音沙哑而冰冷，带着一丝不屑与贪婪，``一群不知天高地厚的小家伙，还有那群伪善的扽魂殿杂碎，竟敢阻碍我帅气教的大业，抽取世间帅气，成就无上扽力，谁也拦不住！''
\end{quote}

水镜之中，扽魂殿决战团队正稳步前行，他们周身散发着纯净的扽力，是帅气教最痛恨的存在------扽魂殿一直致力于打压邪恶扽师，摧毁帅气塔，阻止胡赞的掠夺之举，多年来与帅气教势同水火，不死不休。

胡赞站起身，走到水镜前，干枯的手指轻轻抚过镜面上果糕七怪的倒影，嘴角勾起一抹阴冷的笑容。他转身走向大殿深处，那里悬挂着一幅巨大的古卷------\textbf{《掠夺真经》}，上面记载着帅气教创立千年以来的核心秘术。

\begin{quote}
``千年了\ldots\ldots{}''胡赞低声呢喃，眼神中燃烧着疯狂的火焰，``千年前，帅气教的先祖便已参透天地间的终极奥秘------\textbf{世间万物皆有序力，而序力的本质，便是帅气}。山川河流、日月星辰、生灵的喜怒哀乐，无一不蕴含着纯粹的帅气能量。谁若能掌控所有帅气，便能掌控整个序数世界！''
\end{quote}

他猛然转身，目光如炬，对着空旷的大殿高声宣告，仿佛在向天地宣誓：

\begin{quote}
``我胡赞，誓要完成先祖未竟之业！我要将世间所有帅气尽数掠夺，汇入帅气塔中，铸就\textbf{无上扽力}，突破序数的极限，成为这天地间唯一的\textbf{帅气之神}！''
\end{quote}

大殿两侧的墙壁上，镶嵌着数百颗漆黑的水晶，每一颗水晶中都封印着一个被帅气教击败的序数修炼者的序力精华。这些水晶如同活物般微微跳动，散发着微弱的光芒，那是被囚禁的序力在绝望中挣扎的痕迹。

胡赞走到一颗水晶前，指尖轻触，水晶中便传来一声微弱的惨叫------那是一位曾经名震一方的序数大师，三年前被帅气教击败后，序力被尽数掠夺，灵魂被封印在水晶之中，永世不得超生。

\begin{quote}
``这就是违抗帅气教的下场。''胡赞冷冷说道，``每一个被掠夺的序数修炼者，都会成为帅气塔的养料，他们的序力、他们的智慧、他们的生命，都将化为我的力量。果糕七怪\ldots\ldots 扽魂殿\ldots\ldots 他们身上的序力如此纯粹、如此强大，若是能尽数掠夺，我的无上扽力必将更进一步！''
\end{quote}

胡赞重新回到水镜前，目光锁定在牢大身上，眼神中闪过一丝贪婪：

\begin{quote}
``这个小家伙\ldots\ldots 他的序元 \allowbreak{}\NovelInline{I01765}{\varphi}\allowbreak{} 函数掌控力远超同辈，蓝银序力中蕴含着极其纯粹的序数信念，若是能将他掠夺，我的掠夺序纹必将进化到更高的层次。还有那个辅助系的小子，金色辅助序纹竟能融合全队序力，这种能力简直是天生的'序力容器'，若能将他的序力据为己有，我的帅气塔便能再增十倍威能！''
\end{quote}

他转头看向水镜中扽魂殿的方向，眼神变得更加阴冷：

\begin{quote}
``至于扽魂殿那群伪善的家伙\ldots\ldots 他们身上的纯净扽力，是帅气教最大的威胁，也是最好的养料。纯净扽力与掠夺序纹本就是同源异流，若能将其融合，我便能打破正邪序力的壁垒，创造出前所未有的\textbf{帅气序力}，届时，天地之间，再无人能与我抗衡！''
\end{quote}

胡赞重新坐回王座，指尖凝聚起一缕黑色序力，在虚空中勾勒出一幅复杂的序数阵列------那是帅气教暗中布下的\textbf{掠夺大阵}，覆盖了整个决赛赛场。

\begin{quote}
``这场序数精英大赛，从预选赛开始，便在我的掌控之中。Madore宗、Buchholz宗\ldots\ldots 这些所谓的强队，不过是我安插的棋子，用来试探果糕七怪的实力底线。如今半决赛已结束，果糕七怪的真实实力已经暴露，他们的序数对应逻辑、序力流转方式、团队配合节奏，全部被我尽收眼底。''
\end{quote}

他冷笑着，指尖在虚空中轻轻一划，一幅决赛赛场的立体投影浮现在大殿中央：

\begin{quote}
``决赛赛场上，我早已布下天罗地网。当果糕七怪与扽魂殿决战团队在赛场上拼得两败俱伤之时，掠夺大阵便会启动，将所有人的序力尽数吸入帅气塔中。届时，不仅果糕七怪和扽魂殿将彻底覆灭，连赛场上的数万名观众，也都将成为帅气教的养料！''
\end{quote}

他的声音越来越疯狂，周身黑色序力暴涨，大殿中的烛火全部熄灭，只剩下水镜中映出的幽冷光芒：

\begin{quote}
``这就是帅气教的终极计划------\textbf{以大赛为饵，以决赛为局，掠夺世间一切帅气，成就无上扽力}！谁也拦不住！谁也逃不掉！''
\end{quote}

``教主。''殿外传来两道恭敬却带着狠戾的声音，两道身影快步走入大殿，单膝跪地，周身同样萦绕着黑色掠夺序纹，``请教主吩咐。''

这两人便是胡赞的两个儿子------\textbf{胡勇}与\textbf{胡善}。

\begin{itemize}
\tightlist
\item
  \textbf{胡勇}：身形魁梧，面容凶悍，周身序力狂暴，擅长近距离掠夺帅气与序力；
\item
  \textbf{胡善}：身形瘦削，眼神阴鸷，心思缜密，擅长隐蔽暗杀与序力干扰。
\end{itemize}

兄弟二人配合默契，是胡赞最得力的手下，更是帅气教中最凶悍的暗杀者------二人的序数实力早已突破常规阈值，高达 \allowbreak{}\NovelInline{I01766}{Y(1,3,4,2,8,5,10)}\allowbreak{}，这是一个极其巨大的序数层级，强横到令人发指，即便是声名远扬的弗兰德，在他们二人面前，也如同蝼蚁般渺小，根本不堪一击。

胡赞抬眼，目光落在二人身上，语气冰冷而威严：

\begin{quote}
``胡勇、胡善，你们二人即刻出发，暗中跟随果糕七怪与扽魂殿决战团队，寻得最佳时机，了断他们的踪迹。''
\end{quote}

``是！''二人齐声应道，语气中没有丝毫犹豫。

``记住，''胡赞前倾身体，指尖凝聚起一缕黑色序力，注入二人眉心，``务必隐蔽行事，不可暴露帅气教的踪迹。切不可急于求成，需隐蔽蛰伏，寻得最合适的时机再行事。扽魂殿团队必定会在决赛现场附近待命，届时一并处置，取走他们身上的东西，带回帅气塔，助我突破更高阶的扽力境界！''

``属下明白！''胡勇与胡善眼中闪过一丝贪婪的光芒，眉心的黑色序力微微闪烁，周身的掠夺气息愈发浓郁。他们深知，果糕七怪与扽魂殿成员身上的东西，若是能尽数带回，必将让他们的实力大增，也能得到教主的重赏。

``去吧，莫要让我失望。''胡赞挥了挥手，重新靠回王座，眼神再次投向水镜，嘴角勾起一抹阴狠的笑容，``这场序数精英大赛，终将成为我帅气教崛起的垫脚石，世间帅气，终将归我所有！''

胡勇与胡善躬身退下，身形一闪，便消失在大殿之中，化作两道黑影，朝着果糕七怪赶路的方向疾驰而去，周身的序力被刻意压制，隐蔽得毫无痕迹。

此时的果糕七怪，正行走在一条通往决赛场地的林间小径上。小径两旁古木参天，枝叶交错，遮挡住了大部分阳光，林间弥漫着淡淡的雾气，空气中夹杂着草木的清香与微弱的序力波动。

``牢大，还有多久才能到决赛场地？''小舞揉了揉有些酸胀的肩膀，语气中带着一丝疲惫。半决赛的死战消耗了他们大量的序力，即便有奥斯卡的辅助补充，依旧难以完全恢复。

牢大停下脚步，指尖凝聚起一缕序力，感知着周围的序数节点，缓缓说道：

\begin{quote}
``不远了，穿过这片序数密林，前面就是决赛殿堂的外围。大家再坚持一下，趁这段时间，尽快恢复序力，决赛的对手，只会比Madore宗更强。''
\end{quote}

众人纷纷点头，放慢脚步，一边赶路，一边运转序力恢复自身。奥斯卡则加快了辅助的节奏，金色序纹笼罩着全队，源源不断地为众人补充着序力，嘴里还念叨着：

\begin{quote}
``放心吧，有我在，保证大家决赛时都能满血状态，绝对不能输！''
\end{quote}

就在这时，牢大的眼神骤然一凝，指尖的序纹瞬间停滞，语气变得警惕起来：

\begin{quote}
``等等，有动静！''
\end{quote}

众人瞬间停下脚步，全身序力瞬间绷紧，警惕地看向四周。小舞粉色序力凝聚，眼神如鹰隼般扫视着林间的每一个角落；牢大攥紧双拳，冰蓝序力悄然流转，随时准备应对突发状况；蔡须坤、诺当也纷纷做好战斗准备，目光紧紧盯着林间的阴影。

只见林间的雾气中，几道黑影一闪而过，速度极快，几乎只留下一道模糊的残影，周身散发着微弱的、诡异的序力波动，那序力波动既不是坍缩序数，也不是他们熟悉的 \allowbreak{}\NovelInline{I01767}{\varphi}\allowbreak{} 序力，带着一丝冰冷的掠夺感，一闪而逝，仿佛从未出现过。

``什么东西？''牢大低声喝问，冰蓝序力朝着黑影消失的方向探去，却只捕捉到一丝残留的微弱序力，根本无法判断对方的身份与实力。

唐三眉头紧锁，指尖快速勾勒序纹，推演着那丝残留的序力波动，语气凝重地说道：

\begin{quote}
``不清楚，那序力很诡异，带着掠夺性，不像是大赛的参赛队伍，也不像是扽魂殿的人。''
\end{quote}

``会不会是路过的散修？''奥斯卡有些疑惑地说道，``这片密林里，偶尔也会有一些低阶扽师或者参赛选手路过，可能是我们太过敏感了。''

小舞摇了摇头，眼神依旧警惕：

\begin{quote}
``不像，散修的序力不会这么诡异，而且速度这么快，显然是有备而来，不像是偶然路过。''
\end{quote}

蔡须坤缓缓说道：

\begin{quote}
``不管是什么人，我们都要小心。不过现在我们要赶去决赛，不宜节外生枝，只要他们没有主动攻击我们，暂时先不要理会，尽快赶到决赛场地再说。''
\end{quote}

牢大沉默片刻，点了点头。他能感觉到，那几道黑影的实力不弱，但似乎并没有立刻动手的意思，而且他们现在确实需要尽快赶到决赛场地，恢复序力，备战决赛。若是在这里与不明身份的人纠缠，只会消耗更多的序力，得不偿失。

\begin{quote}
``大家提高警惕，继续赶路，奥斯卡，加强辅助序力，一旦有异常，立刻通知大家。''唐三沉声吩咐道，周身蓝银色序力再次流转，时刻感知着周围的动静，心底却隐隐有些不安------那丝诡异的序力波动，总让他觉得，事情并没有那么简单。
\end{quote}

众人应声，再次出发，脚步比之前更加谨慎。林间的雾气似乎更浓了，那几道黑影再也没有出现，仿佛只是众人的错觉。可只有唐三知道，那不是错觉，那丝残留的序力波动，如同一个隐患，潜藏在他们前行的道路上。

而在他们身后不远处的密林阴影中，胡勇与胡善并肩而立，眼神阴鸷地盯着果糕七怪的背影，周身的序力被压制到极致。

``大哥，要不要现在动手？''胡善低声问道，指尖凝聚起一缕黑色序力，眼神中透着狠戾。

胡勇摇了摇头，眼神凝重却带着绝对的自信，缓缓说道：

\begin{quote}
``不急，教主早有吩咐，让我们隐蔽蛰伏，寻得最合适的时机再行事。他们现在虽序力未恢复，但警惕性极高，牢大的推演能力又强，我们二人虽有 \allowbreak{}\NovelInline{I01768}{Y(1,3,4,2,8,5,10)}\allowbreak{} 的序数实力，远超他们想象，即便弗兰德在此也只能俯首称臣，但贸然行事难免暴露行踪，坏了教主的大事。耐心等待，等他们卸下防备、露出破绽的那一刻，我们再行动，到时候，他们再无还手之力。''
\end{quote}

胡善点了点头，眼中闪过一丝阴狠的笑意，收回指尖的序力，与胡勇一同再次隐入阴影之中：

\begin{quote}
``大哥说得对，唯有耐心等待，才能万无一失。我们就一路跟随，紧盯他们的一举一动，等到那最合适的一刻，再按教主的吩咐行事，取走他们身上的东西，向教主复命。''
\end{quote}

二人如同两匹蛰伏的恶狼，紧紧跟随着果糕七怪的脚步，眼底藏着不为人知的恶意，耐心等待着时机的到来。

\begin{center}\rule{0.5\linewidth}{0.5pt}\end{center}

果糕七怪依旧在林间赶路，阳光透过枝叶的缝隙，洒下斑驳的光影。他们虽然依旧警惕着周围的动静，但渐渐放松了紧绷的神经------毕竟，他们刚刚赢得了一场艰难的半决赛，满心都是即将到来的决赛，谁也没有察觉，暗处有两道冰冷的目光，始终追随着他们的脚步，一场针对他们与扽魂殿决战团队的隐秘算计，正在悄然酝酿。

\chapter{第二十五章：扩展序元Veblen函数}\label{ux7b2cux4e8cux5341ux4e94ux7ae0ux6269ux5c55ux5e8fux5143veblenux51fdux6570}

果糕七怪穿过序数密林，终于抵达了决赛殿堂外围的备战区。这是一座半开放式的修炼场，穹顶镶嵌着密密麻麻的序数光纹，散发着柔和而深邃的光芒，空气中弥漫着浓郁的序力气息，连呼吸都能感受到序数能量的流转。

``大家先休整片刻，恢复序力。''牢大环顾四周，确认没有异常后，沉声说道。

众人纷纷落座，奥斯卡第一时间催动金色辅助序纹，为全队进行深层序力补充。小舞靠在石柱上，闭目调息；蔡须坤与诺当低声交流着半决赛中的配合细节；卢本韦则蹲在一旁，指尖在地面上反复勾勒着序元纹路，似乎在推演着什么。

唐三走到牢大身边，眉头微蹙，压低声音说道：

\begin{quote}
``牢大，密林里那几道黑影的事，我越想越不对劲。那丝序力波动的层级\ldots\ldots 远超我们目前接触过的任何对手，甚至远超Madore宗的赛伊德。如果那不是散修，而是某个隐藏势力的探子\ldots\ldots{}''
\end{quote}

牢大目光一沉，点了点头：

\begin{quote}
``我也感觉到了。那丝序力中带着一种\ldots\ldots 掠夺性，和我们熟知的 \allowbreak{}\NovelInline{I01769}{\varphi}\allowbreak{} 序力、\allowbreak{}\NovelInline{I01770}{\psi}\allowbreak{} 坍缩序力都不同，像是某种我们从未接触过的体系。决赛在即，我们必须尽快提升实力，否则一旦那个势力出手，我们连还手的机会都没有。''
\end{quote}

话音刚落，备战区的入口处传来一阵沉稳的脚步声。弗兰德与玉小刚并肩走来，弗兰德的周身金色序力比往日更加凝练，眼神中透着一股前所未有的严肃。

\begin{quote}
``你们的警觉性不错，那丝序力波动确实不简单。''弗兰德开门见山，语气沉重，``但现在不是追查的时候。决赛在即，你们目前的序数修为还不够。今天，我要传授你们一套全新的知识------\textbf{扩展序元Veblen函数}。这是你们突破现有战力天花板、应对决赛中可能出现的未知对手的关键。''
\end{quote}

玉小刚补充道：

\begin{quote}
``这套知识极为深奥，涉及序元Veblen函数的极限以及超越极限后的扩展体系。你们必须在决赛前将其完全掌握，否则面对真正的强敌，你们连对方的序数层级都无法解析。''
\end{quote}

众人闻言，纷纷起身，神色肃然。唐三眼中闪过一丝兴奋------他等这一天已经等了很久。

弗兰德走到修炼场中央，抬手凝聚序链，莹白色的光带在空气中缓缓展开，勾勒出一幅宏大的序数图谱。图谱的末端，一个金色的符号熠熠生辉------

\NovelDisplay{M00218R01}{%
\varphi(1 @(1, 0)) = \text{LVO}%
}

\begin{quote}
``这是我们之前讲到的\textbf{大Veblen序数（LVO）}，序元Veblen函数的极限。''弗兰德的声音沉稳而有力，``LVO的诞生，源于我们对序元Veblen函数位置变量的一次'进位'------我们定义了映射 \allowbreak{}\NovelInline{I01771}{\rho \mapsto \varphi(1 @\rho)}\allowbreak{} 的第一个不动点，便得到了LVO。''
\end{quote}

他指向图谱中LVO之后的空白区域，那里一片虚无，什么都没有。

\begin{quote}
``但LVO并不是终点。我们可以对位置变量继续'进位'------既然 \allowbreak{}\NovelInline{I01772}{\alpha @\beta}\allowbreak{} 中的 \allowbreak{}\NovelInline{I01773}{\beta}\allowbreak{} 可以是一个序数，那么它为什么不能是一个\textbf{由一系列序数组成的数组}呢？''
\end{quote}

唐三瞳孔微缩，瞬间捕捉到了关键：

\begin{quote}
``老师，您的意思是\ldots\ldots 把位置变量从单个序数扩展为\textbf{序数数组}？''
\end{quote}

``没错。''弗兰德赞许地点头，``这就是\textbf{扩展序元Veblen函数}的核心思想。位置变量不再是一个序数，而是一个数组，比如 \allowbreak{}\NovelInline{I01774}{(1, 0)}\allowbreak{}、\allowbreak{}\NovelInline{I01775}{(1, 1)}\allowbreak{}、\allowbreak{}\NovelInline{I01776}{(\omega, 0)}\allowbreak{}，甚至可以是嵌套的数组，比如 \allowbreak{}\NovelInline{I01777}{(1 @(1, 0))}\allowbreak{}。各版本扩展序元Veblen函数的规则不尽相同，我们下面介绍梅天理给出的规则。''

弗兰德抬手，空气中浮现出密密麻麻的公式，他将规则分为四大类，逐一讲解。

\begin{quote}
此时Veblen函数可以记为 \allowbreak{}\NovelInline{I01778}{\varphi(\#, \alpha @\beta)}\allowbreak{}，其基本列为： \NovelDisplay{M00219R01}{%
\varphi(\#, \alpha @\beta)[n] = \varphi(\#, \alpha[n] @\beta)%
}
\end{quote}

弗兰德解释道：

\begin{quote}
``这条规则和我们之前学的序元Veblen函数完全一致，没什么特别的。当最低非零位是极限序数时，直接对该极限序数取基本列即可。''
\end{quote}

弗兰德语气加重：

\begin{quote}
``这条规则比较复杂，需要根据 \allowbreak{}\NovelInline{I01779}{\beta}\allowbreak{} 的类型进一步细分。''
\end{quote}

\NovelDisplay{M00220R01}{%
\varphi(\#, (\alpha + 1) @\beta)[n] = \varphi(\#, \alpha @\beta, 1 @\beta[n])%
}

弗兰德举例：

\NovelDisplay{M00221R01}{%
\varphi(1 @\zeta_1, 4 @\varepsilon_5, 1 @\omega^4) = \sup\{\varphi(1 @\zeta_1, 4 @\varepsilon_5, 1 @0), \varphi(1 @\zeta_1, 4 @\varepsilon_5, 1 @\omega^3), \varphi(1 @\zeta_1, 4 @\varepsilon_5, 1 @\omega^3 \cdot 2), \dots\}%
}

\begin{quote}
``注意这里的结构------当 \allowbreak{}\NovelInline{I01780}{\beta}\allowbreak{} 是极限序数时，展开方式和我们之前讲的规则7（\allowbreak{}\NovelInline{I01781}{\beta}\allowbreak{} 为极限、\allowbreak{}\NovelInline{I01782}{\alpha}\allowbreak{} 为后继）是一致的：在 \allowbreak{}\NovelInline{I01783}{\alpha}\allowbreak{} 的位置减一，同时在末尾新增一项 \allowbreak{}\NovelInline{I01784}{1 @\beta[n]}\allowbreak{}。''
\end{quote}

小舞若有所思地点头：

\begin{quote}
``所以 \allowbreak{}\NovelInline{I01785}{\varphi(1 @\omega^4)}\allowbreak{} 的基本列就是不断在末尾追加 \allowbreak{}\NovelInline{I01786}{1 @\omega^3}\allowbreak{}、\allowbreak{}\NovelInline{I01787}{1 @(\omega^3 \cdot 2)}\allowbreak{}\ldots\ldots 直到极限？''
\end{quote}

``完全正确。''弗兰德赞许道。

\NovelDisplay{M00222R01}{%
\varphi(\#, (\alpha + 1) @(\beta + 1))[0] = \varphi(\#, \alpha @(\beta + 1))%
}

\NovelDisplay{M00223R01}{%
\varphi(\#, (\alpha + 1) @(\beta + 1))[n + 1] = \varphi(\#, \alpha @\beta, X[n] @\beta)%
}

其中 \allowbreak{}\NovelInline{I01788}{X = \varphi(\#, (\alpha + 1) @(\beta + 1))}\allowbreak{}，即函数自身。

弗兰德举例说明：

\NovelDisplay{M00224R01}{%
\varphi(1, 0, 0) = \sup\{1, \varphi(1, 0), \varphi(\varphi(1, 0), 0), \dots\}%
}

\NovelDisplay{M00225R01}{%
\varphi(1 @\omega^2, 2 @(\omega + 1)) = \sup\{\varphi(1 @\omega^2, 1 @(\omega + 1)),%
} \NovelDisplay{M00226R01}{%
\varphi(1 @\omega^2, 1 @(\omega + 1), \varphi(1 @\omega^2, 1 @(\omega + 1)) @\omega),%
} \NovelDisplay{M00227R01}{%
\varphi(1 @\omega^2, 1 @(\omega + 1), \varphi(1 @\omega^2, 1 @(\omega + 1), \varphi(1 @\omega^2, 1 @(\omega + 1)) @\omega) @\omega), \dots\}%
}

玉小刚补充道：

\begin{quote}
``这条规则的核心是\textbf{自指}------\allowbreak{}\NovelInline{I01789}{X[n]}\allowbreak{} 中引用了函数自身的前一项，形成递归嵌套。\allowbreak{}\NovelInline{I01790}{\varphi(1, 0, 0)}\allowbreak{} 就是最经典的例子：第零项是 \allowbreak{}\NovelInline{I01791}{1}\allowbreak{}，第一项是 \allowbreak{}\NovelInline{I01792}{\varphi(1, 0)}\allowbreak{}，第二项是 \allowbreak{}\NovelInline{I01793}{\varphi(\varphi(1, 0), 0)}\allowbreak{}\ldots\ldots 每一项都是对前一项的再次嵌套。这种自指结构在实战中极为强大，因为它的序力会随着嵌套层数的增加而指数级暴涨。''
\end{quote}

此时Veblen函数可以记为 \allowbreak{}\NovelInline{I01794}{\varphi(\#, \alpha_1 @\beta_1, (\alpha + 1) @0)}\allowbreak{}，其基本列的第零项为：

\NovelDisplay{M00228R01}{%
\varphi(\#, \alpha_1 @\beta_1, (\alpha + 1) @0)[0] = \varphi(\#, \alpha_1 @\beta_1, \alpha @0) + 1%
}

弗兰德语气凝重：

\begin{quote}
``这一情形最为复杂，需要根据 \allowbreak{}\NovelInline{I01795}{\alpha_1}\allowbreak{} 和 \allowbreak{}\NovelInline{I01796}{\beta_1}\allowbreak{} 的类型进一步细分。''
\end{quote}

\textbf{子情形2.3.1：\allowbreak{}\NovelInline{I01797}{\alpha_1}\allowbreak{} 为极限序数}

\NovelDisplay{M00229R01}{%
\varphi(\#, \alpha_1 @\beta_1, (\alpha + 1) @0)[n] = \varphi(\#, \alpha_1[n] @\beta_1, (\varphi(\#, \alpha_1 @\beta_1, \alpha @0) + 1) @0)%
}

举例：

\NovelDisplay{M00230R01}{%
\varphi(\omega @\omega, 1 @0) = \sup\{\varphi(\omega @\omega) + 1, \varphi(1 @\omega, \varphi(\omega @\omega) + 1 @0), \varphi(2 @\omega, \varphi(\omega @\omega) + 1 @0), \dots\}%
}

弗兰德特别强调：

\begin{quote}
``大家注意 \allowbreak{}\NovelInline{I01798}{\varphi(\omega @\omega, 1 @0)}\allowbreak{} 的展开------这正是我们在第十七章实战演练中重点练习过的序链！现在你们可以看到它的全貌了：第零项是 \allowbreak{}\NovelInline{I01799}{\varphi(\omega @\omega) + 1}\allowbreak{}，之后每一项都将 \allowbreak{}\NovelInline{I01800}{\omega @\omega}\allowbreak{} 中的 \allowbreak{}\NovelInline{I01801}{\omega}\allowbreak{} 替换为 \allowbreak{}\NovelInline{I01802}{n}\allowbreak{}，同时将 \allowbreak{}\NovelInline{I01803}{\varphi(\omega @\omega) + 1}\allowbreak{} 作为新的参数嵌入。这就是扩展序元Veblen函数的威力------它能将我们之前学过的所有序数纳入统一的框架中。''
\end{quote}

唐三闻言，眼中闪过一丝恍然：

\begin{quote}
``所以我们在半决赛中用 \allowbreak{}\NovelInline{I01804}{\varphi(1 @\omega^2)}\allowbreak{} 对位赛伊德的 \allowbreak{}\NovelInline{I01805}{\psi(\Omega^{\Omega^{\omega^2}})}\allowbreak{}，其实只是触及了扩展序元Veblen函数的皮毛\ldots\ldots 真正的完全体，远比我们想象的强大。''
\end{quote}

``没错。''弗兰德点头，``继续往下看。''

\begin{center}\rule{0.5\linewidth}{0.5pt}\end{center}

\textbf{子情形2.3.2.1：\allowbreak{}\NovelInline{I01806}{\alpha_1}\allowbreak{} 为后继序数，\allowbreak{}\NovelInline{I01807}{\beta_1}\allowbreak{} 为极限序数}

\NovelDisplay{M00231R01}{%
\varphi(\#, (\alpha_1 + 1) @\beta_1, (\alpha + 1) @0)[n]%
} \NovelDisplay{M00232R01}{%
= \varphi(\#, \alpha_1 @\beta_1, 1 @\beta_1[n], (\varphi(\#, (\alpha_1 + 1) @\beta_1, \alpha @0) + 1) @0)%
}

举例：

\NovelDisplay{M00233R01}{%
\varphi(2 @\varepsilon_0, 1 @0) = \sup\{\varphi(2 @\varepsilon_0) + 1, \varphi(1 @\varepsilon_0, 1 @\omega, \varphi(2 @\varepsilon_0) + 1 @0), \varphi(1 @\varepsilon_0, 1 @\omega^2, \varphi(2 @\varepsilon_0) + 1 @0), \dots\}%
}

\begin{center}\rule{0.5\linewidth}{0.5pt}\end{center}

\textbf{子情形2.3.2.2：\allowbreak{}\NovelInline{I01808}{\alpha_1}\allowbreak{} 为后继序数，\allowbreak{}\NovelInline{I01809}{\beta_1}\allowbreak{} 为后继序数}

\NovelDisplay{M00234R01}{%
\varphi(\#, (\alpha_1 + 1) @(\beta_1 + 1), (\alpha + 1) @0)[n + 1]%
} \NovelDisplay{M00235R01}{%
= \varphi(\#, \alpha_1 @(\beta_1 + 1), \varphi(\#, (\alpha_1 + 1) @(\beta_1 + 1), (\alpha + 1) @0)[n] @\beta_1)%
}

举例：

\NovelDisplay{M00236R01}{%
\varphi(2, 0, 1) = \sup\{\varphi(2, 0, 0) + 1, \varphi(1, \varphi(2, 0, 0) + 1, 0), \varphi(1, \varphi(1, \varphi(2, 0, 0) + 1, 0), 0), \dots\}%
}

\NovelDisplay{M00237R01}{%
\varphi(3, 0, 0, 3) = \sup\{\varphi(3, 0, 0, 2) + 1, \varphi(2, \varphi(3, 0, 0, 2) + 1, 0, 0), \dots\}%
}

\NovelDisplay{M00238R01}{%
\varphi(\alpha @0) = \omega^\alpha%
}

弗兰德轻描淡写地带过：

\begin{quote}
``这条很简单，当 \allowbreak{}\NovelInline{I01810}{\varphi}\allowbreak{} 函数内部只有一个参数且 \allowbreak{}\NovelInline{I01811}{\beta = 0}\allowbreak{} 时，退化为普通的 \allowbreak{}\NovelInline{I01812}{\omega}\allowbreak{} 幂次。这是最基础的规则，大家应该已经烂熟于心了。''
\end{quote}

\NovelDisplay{M00239R01}{%
(\varphi(\#_1) + \varphi(\#_2) + \dots + \varphi(\#_k))[n] = \varphi(\#_1) + \varphi(\#_2) + \dots + \varphi(\#_k)[n]%
}

\begin{quote}
``多个 \allowbreak{}\NovelInline{I01813}{\varphi}\allowbreak{} 函数求和时，只有最后一项取基本列，其余项保持不变。这和我们之前讲的标准形式定义是一致的。''
\end{quote}

弗兰德话锋一转，语气变得更加严肃：

\begin{quote}
``以上规则覆盖了位置变量为单个序数的情况。但当位置变量本身也是一个\textbf{数组}时，比如 \allowbreak{}\NovelInline{I01814}{\beta_1 = (1, 0)}\allowbreak{} 或更复杂的嵌套数组，我们需要递归地应用上述规则。''
\end{quote}

\NovelDisplay{M00240R01}{%
\varphi(\#, (\alpha + 1) @(1, 0))[0] = \varphi(\#, \alpha @(1, 0))%
}

\NovelDisplay{M00241R01}{%
\varphi(\#, (\alpha + 1) @(1, 0))[n] = \varphi(\#, \alpha @(1, 0), 1 @\varphi(\#, \alpha @(1, 0)))%
}

举例：

\NovelDisplay{M00242R01}{%
\varphi(2 @(1, 0)) = \sup\{\varphi(1 @(1, 0)),, \varphi(1 @(1, 0), 1 @\varphi(1, 0)),,%
} \NovelDisplay{M00243R01}{%
\varphi(1 @(1, 0),  1 @\varphi(1 @\varphi(1, 0)))),, \dots\}%
}

弗兰德指着公式中的 \allowbreak{}\NovelInline{I01815}{\varphi(1 @(1, 0))}\allowbreak{}，目光灼灼：

\begin{quote}
``这就是LVO！\allowbreak{}\NovelInline{I01816}{\varphi(1 @(1, 0))}\allowbreak{} 就是LVO本身。所以 \allowbreak{}\NovelInline{I01817}{\varphi(2 @(1, 0))}\allowbreak{} 的基本列中，每一项都在前一项的基础上，将LVO作为新参数嵌入，形成层层嵌套。这种结构的序力增长是\textbf{超指数级}的，远超你们目前掌握的任何序链。''
\end{quote}

弗兰德继续写道：

\NovelDisplay{M00244R01}{%
\varphi(1 @(\omega, 0)) = \sup\{\omega, \varphi(1 @(1, 0)), \varphi(1 @(2, 0)), \varphi(1 @(3, 0)), \dots\}%
}

\NovelDisplay{M00245R01}{%
\varphi(1 @(1 @(1, 0))) = \sup\{\varphi(1 @\omega), \varphi(1 @(1 @\varphi(1 @\omega))),%
} \NovelDisplay{M00246R01}{%
\varphi(1 @(1 @\varphi(1 @(1 @\varphi(1 @\omega))))), \dots\}%
}

\begin{quote}
``看到了吗？\allowbreak{}\NovelInline{I01818}{\varphi(1 @(1 @(1, 0)))}\allowbreak{} 的展开中，出现了三层嵌套------最内层是 \allowbreak{}\NovelInline{I01819}{\varphi(1 @\omega)}\allowbreak{}（即SVO），中间层将其代入 \allowbreak{}\NovelInline{I01820}{1 @(\cdot)}\allowbreak{}，最外层再次代入。这种深度嵌套的结构，在实战中可以形成\textbf{多层防御屏障}，每一层都比前一层强横数倍。''
\end{quote}

弗兰德的语气变得无比庄重，他在空气中写下了一个金色的符号：

\NovelDisplay{M00247R01}{%
\text{BHO} = \sup\{\varphi(1, 0), \varphi(1 @(1, 0)), \varphi(1 @(1 @(1, 0))), \dots\}%
}

\begin{quote}
``这就是\textbf{Bachmann-Howard序数（BHO）}，扩展序元Veblen函数的极限。从LVO出发，我们继续对位置变量进行'进位'------\allowbreak{}\NovelInline{I01821}{\varphi(1 @(1, 0))}\allowbreak{} 是LVO，\allowbreak{}\NovelInline{I01822}{\varphi(1 @(1 @(1, 0)))}\allowbreak{} 是在LVO的基础上再次进位，如此反复迭代，取所有迭代的极限，便得到了BHO。''
\end{quote}

他停顿片刻，目光扫过众人：

\begin{quote}
``BHO的意义在于，它标志着序元Veblen函数体系的\textbf{终极边界}。在此之后，Veblen函数的扩展过于复杂，强度也没有预想中的那样高。因此，我们对Veblen函数的扩展就到此为止，在后文中将采用其他的模式在序数世界之中继续前进。''
\end{quote}

玉小刚接过话头，在空气中写下了一组全新的对应公式：

\begin{quote}
``BHO并不是Veblen函数的终点。在BHO之后，\allowbreak{}\NovelInline{I01823}{\psi}\allowbreak{} 函数继续向前推进，与扩展序元Veblen函数形成了全新的对应关系。你们必须牢记以下对应：''
\end{quote}

\NovelDisplay{M00248R01}{%
\psi(\Omega^{\Omega^\Omega} \cdot 2) = \varphi(1 @(1, 0), 1 @0)%
} \NovelDisplay{M00249R01}{%
\psi(\Omega^{\Omega^\Omega + 1}) = \varphi(1 @(1, 0), 1 @1)%
} \NovelDisplay{M00250R01}{%
\psi(\Omega^{\Omega^\Omega + \Omega^\omega}) = \varphi(1 @(1, 0), 1 @\omega)%
} \NovelDisplay{M00251R01}{%
\psi(\Omega^{\Omega^\Omega} \cdot 2) = \varphi(2 @(1, 0))%
} \NovelDisplay{M00252R01}{%
\psi(\Omega^{\Omega^{\Omega + 1}}) = \varphi(1 @(1, 1))%
} \NovelDisplay{M00253R01}{%
\psi(\Omega^{\Omega^{\Omega + 1}} \cdot 2) = \varphi(2 @(1, 1))%
} \NovelDisplay{M00254R01}{%
\psi(\Omega^{\Omega^{\Omega + 2}}) = \varphi(1 @(1, 2))%
} \NovelDisplay{M00255R01}{%
\psi(\Omega^{\Omega^{\Omega + \omega}}) = \varphi(1 @(1, \omega))%
} \NovelDisplay{M00256R01}{%
\psi(\Omega^{\Omega^{\Omega + \psi(\Omega^\Omega)}}) = \varphi(1 @(1, \varphi(1 @(1, 0))))%
} \NovelDisplay{M00257R01}{%
\psi(\Omega^{\Omega^{\Omega \cdot 2}}) = \varphi(1 @(2, 0))%
} \NovelDisplay{M00258R01}{%
\psi(\Omega^{\Omega^{\Omega \cdot 2 + 1}}) = \varphi(1 @(2, 1))%
} \NovelDisplay{M00259R01}{%
\psi(\Omega^{\Omega^{\Omega \cdot 3}}) = \varphi(1 @(3, 0))%
} \NovelDisplay{M00260R01}{%
\psi(\Omega^{\Omega^{\Omega \cdot \omega}}) = \varphi(1 @(\omega, 0))%
} \NovelDisplay{M00261R01}{%
\psi(\Omega^{\Omega^{\Omega \cdot \psi(\Omega^\Omega)}}) = \varphi(1 @(\varphi(1 @(\omega, 0)), 0))%
}

玉小刚逐一解释：

\begin{quote}
``注意看，这些对应关系中，\allowbreak{}\NovelInline{I01824}{\psi}\allowbreak{} 函数中 \allowbreak{}\NovelInline{I01825}{\Omega}\allowbreak{} 的指数结构发生了质变------从之前的 \allowbreak{}\NovelInline{I01826}{\Omega^\omega}\allowbreak{} 层级跃升到了 \allowbreak{}\NovelInline{I01827}{\Omega^\Omega}\allowbreak{} 层级。这意味着，\allowbreak{}\NovelInline{I01828}{\psi}\allowbreak{} 函数已经突破了序元Veblen函数的极限，进入了\textbf{扩展序元Veblen函数}的\#\# 五、超越Veblen：\allowbreak{}\NovelInline{I01829}{\psi}\allowbreak{} 函数与扩展序元Veblen函数的对应（续）
\end{quote}

玉小刚逐一解释：

\begin{quote}
``注意看，这些对应关系中，\allowbreak{}\NovelInline{I01830}{\psi}\allowbreak{} 函数中 \allowbreak{}\NovelInline{I01831}{\Omega}\allowbreak{} 的指数结构发生了质变------从之前的 \allowbreak{}\NovelInline{I01832}{\Omega^\omega}\allowbreak{} 层级跃升到了 \allowbreak{}\NovelInline{I01833}{\Omega^\Omega}\allowbreak{} 层级。这意味着，\allowbreak{}\NovelInline{I01834}{\psi}\allowbreak{} 函数已经突破了序元Veblen函数的极限，进入了\textbf{扩展序元Veblen函数}的领域。''
\end{quote}

他指向第一组对应：

\begin{quote}
``\allowbreak{}\NovelInline{I01835}{\psi(\Omega^{\Omega^\Omega} \cdot 2) = \varphi(1 @(1, 0), 1 @0)}\allowbreak{}------这里的 \allowbreak{}\NovelInline{I01836}{\Omega^{\Omega^\Omega}}\allowbreak{} 正是LVO在 \allowbreak{}\NovelInline{I01837}{\psi}\allowbreak{} 函数中的对应，而乘以 \allowbreak{}\NovelInline{I01838}{2}\allowbreak{} 则标志着从LVO向 \allowbreak{}\NovelInline{I01839}{\varphi(1 @(1, 0), 1 @0)}\allowbreak{} 的跨越。这是超越LVO后的第一个序数，也是扩展序元Veblen函数的起点。''
\end{quote}

弗兰德补充道：

\begin{quote}
``大家特别注意这组对应------\allowbreak{}\NovelInline{I01840}{\psi(\Omega^{\Omega^{\Omega + \psi(\Omega^{\Omega^\Omega})}}) = \varphi(1 @(1, \varphi(1 @(1, 0))))}\allowbreak{}。这里出现了\textbf{三重嵌套}：\allowbreak{}\NovelInline{I01841}{\psi(\Omega^{\Omega^\Omega})}\allowbreak{} 对应 \allowbreak{}\NovelInline{I01842}{\varphi(1 @(1, 0))}\allowbreak{} 即LVO，而外层的 \allowbreak{}\NovelInline{I01843}{\Omega^{\Omega^{\Omega + \psi(\Omega^{\Omega^\Omega})}}}\allowbreak{} 则将LVO作为 \allowbreak{}\NovelInline{I01844}{\Omega}\allowbreak{} 指数的一部分，最终对应到 \allowbreak{}\NovelInline{I01845}{\varphi(1 @(1, \text{LVO}))}\allowbreak{}。这种嵌套结构在实战中极为恐怖------它的序力层级已经不是简单的叠加，而是\textbf{指数级的嵌套增长}。''
\end{quote}

\begin{quote}
``这正是我急着传授你们这套知识的原因。你们在半决赛中展现的实力，在真正的顶级强者面前，不过是蝼蚁。但如果你们能熟练掌握扩展序元Veblen函数的展开与对位，至少能在关键时刻保住性命。''
\end{quote}

\begin{center}\rule{0.5\linewidth}{0.5pt}\end{center}

\chapter{第二十六章：实战推演：用扩展序元Veblen}\label{ux7b2cux4e8cux5341ux516dux7ae0ux5b9eux6218ux63a8ux6f14ux7528ux6269ux5c55ux5e8fux5143veblen}

``理论讲完了，现在进入实战推演。''弗兰德一挥手，修炼场中央升起一座虚拟对战台，暗紫色的序链在台上翻涌，模拟出一个极其强大的对手。

\begin{quote}
``假设对手释放 \allowbreak{}\NovelInline{I01846}{\varphi(1 @(1, 0), 1 @0)}\allowbreak{} 级别的攻击------这是超越LVO后的第一阶序链，对应 \allowbreak{}\NovelInline{I01847}{\psi(\Omega^{\Omega^\Omega} \cdot 2)}\allowbreak{}。你们需要在3秒内展开它的基本列，找到破绽。谁来？''
\end{quote}

蔡须坤挺身而出，指尖飞速勾勒序纹。她深吸一口气，脑海中飞速回忆刚才学到的规则：

\begin{quote}
``\allowbreak{}\NovelInline{I01848}{\varphi(1 @(1, 0), 1 @0)}\allowbreak{} 中，最低非零位是 \allowbreak{}\NovelInline{I01849}{1 @0}\allowbreak{}，其中 \allowbreak{}\NovelInline{I01850}{\alpha = 0}\allowbreak{}（后继序数，\allowbreak{}\NovelInline{I01851}{\alpha = 0 + 0}\allowbreak{}，不对\ldots\ldots{}\allowbreak{}\NovelInline{I01852}{1 = 0 + 1}\allowbreak{}，所以 \allowbreak{}\NovelInline{I01853}{\alpha = 0}\allowbreak{}，\allowbreak{}\NovelInline{I01854}{\alpha + 1 = 1}\allowbreak{}），\allowbreak{}\NovelInline{I01855}{\beta = 0}\allowbreak{}。这属于情形2.3------\allowbreak{}\NovelInline{I01856}{\beta}\allowbreak{} 为0的情况。''
\end{quote}

她顿了顿，继续推演：

\begin{quote}
``前面还有 \allowbreak{}\NovelInline{I01857}{\varphi(1 @(1, 0))}\allowbreak{} 作为前置参数，即 \allowbreak{}\NovelInline{I01858}{\alpha_1 = 1}\allowbreak{}（后继序数），\allowbreak{}\NovelInline{I01859}{\beta_1 = (1, 0)}\allowbreak{}（数组）。这属于子情形2.4.3.1------\allowbreak{}\NovelInline{I01860}{\beta_1}\allowbreak{} 为 \allowbreak{}\NovelInline{I01861}{(1, 0)}\allowbreak{} 的情况。''
\end{quote}

弗兰德眼中闪过一丝赞许：

\begin{quote}
``继续展开。''
\end{quote}

蔡须坤额头渗出汗水，但指尖的序纹越来越清晰：

他深吸一口气，逐一代入：

\begin{quote}
``第零项：\allowbreak{}\NovelInline{I01862}{\varphi(1 @(1, 0), 1 @0)[0] = \varphi(1 @(1, 0)) + 1 = \text{LVO} + 1}\allowbreak{}''
\end{quote}

\begin{quote}
``第一项：\allowbreak{}\NovelInline{I01863}{\varphi(1 @(1, 0), 1 @0)[1] = \varphi(1 @(\text{LVO} + 1))}\allowbreak{}''
\end{quote}

\begin{quote}
``第二项：\allowbreak{}\NovelInline{I01864}{\varphi(1 @(1, 0), 1 @0)[2] = \varphi(1 @\varphi(1 @(\text{LVO} + 1)))}\allowbreak{}''
\end{quote}

\begin{quote}
``以此类推，每一项都将前一项的结果作为新参数嵌入，形成层层嵌套！''
\end{quote}

弗兰德猛然拍手：

\begin{quote}
``完全正确！唐三的推演非常精准。大家看------''他指向虚拟对战台上的暗紫色序链，``对手的攻击序链 \allowbreak{}\NovelInline{I01865}{\varphi(1 @(1, 0), 1 @0)}\allowbreak{} 的基本列就是这样的嵌套结构。第零项是 \allowbreak{}\NovelInline{I01866}{\text{LVO} + 1}\allowbreak{}，这是攻击的'种子'；第一项将 \allowbreak{}\NovelInline{I01867}{\text{LVO} + 1}\allowbreak{} 作为参数嵌入，攻击强度暴涨；第二项再次嵌套，强度再翻数倍。如果我们不能在它完成第三项嵌套之前击破种子，后续的攻击将不可阻挡！''
\end{quote}

牢大沉声问道：

\begin{quote}
``那破绽在哪里？''
\end{quote}

``破绽就在第零项。''玉小刚接过话头，``\allowbreak{}\NovelInline{I01868}{\text{LVO} + 1}\allowbreak{} 是整个嵌套结构的根基，它对应 \allowbreak{}\NovelInline{I01869}{\psi}\allowbreak{} 函数中的 \allowbreak{}\NovelInline{I01870}{\psi(\Omega^{\Omega^\Omega}) + 1}\allowbreak{}。如果我们用 \allowbreak{}\NovelInline{I01871}{\psi(\Omega^{\Omega^\Omega} + 1)}\allowbreak{} 级别的序链进行精准打击，就能在第零项尚未成长为第一项之前，将其彻底瓦解。''

弗兰德切换虚拟场景，暗紫色序链变得更加浓郁：

\begin{quote}
``第二轮，假设对手释放 \allowbreak{}\NovelInline{I01872}{\varphi(2 @(1, 0))}\allowbreak{} 级别的攻击。谁来展开？''
\end{quote}

小舞举手，粉色序力在指尖流转，眼神专注：

\begin{quote}
``我来试试。\allowbreak{}\NovelInline{I01873}{\varphi(2 @(1, 0))}\allowbreak{} 中，\allowbreak{}\NovelInline{I01874}{\alpha = 1}\allowbreak{}（后继序数，\allowbreak{}\NovelInline{I01875}{\alpha + 1 = 2}\allowbreak{}），\allowbreak{}\NovelInline{I01876}{\beta = (1, 0)}\allowbreak{}（数组）。这属于情形2.4.1------\allowbreak{}\NovelInline{I01877}{\beta_1}\allowbreak{} 为 \allowbreak{}\NovelInline{I01878}{(1, 0)}\allowbreak{} 的情况。''
\end{quote}

她深吸一口气，缓缓展开：

\begin{quote}
``根据规则2.4.1：\allowbreak{}\NovelInline{I01879}{\varphi(\#, (\alpha + 1) @(1, 0))[0] = \varphi(\#, \alpha @(1, 0))}\allowbreak{}''
\end{quote}

\begin{quote}
``\allowbreak{}\NovelInline{I01880}{\varphi(\#, (\alpha + 1) @(1, 0))[n] = \varphi(\#, \alpha @(1, 0), 1 @\varphi(\#, \alpha @(1, 0)))}\allowbreak{}''
\end{quote}

\begin{quote}
``代入 \allowbreak{}\NovelInline{I01881}{\alpha = 1}\allowbreak{}：第零项是 \allowbreak{}\NovelInline{I01882}{\varphi(1 @(1, 0)) = \text{LVO}}\allowbreak{}''
\end{quote}

\begin{quote}
``第一项是 \allowbreak{}\NovelInline{I01883}{\varphi(1 @(1, 0), 1 @0)}\allowbreak{}''
\end{quote}

\begin{quote}
``第二项是 \allowbreak{}\NovelInline{I01884}{\varphi(1 @(1, 0), 1 @\varphi(1@0)))}\allowbreak{}''
\end{quote}

她顿了顿，眼神越来越亮：

\begin{quote}
``每一项都将前一项的结果作为参数，嵌入到 \allowbreak{}\NovelInline{I01885}{1 @(\cdot)}\allowbreak{} 的位置，形成嵌套。而且这个嵌套是在LVO的基础上进行的，所以序力增长的起点就远超常规序链！''
\end{quote}

弗兰德赞许点头：

\begin{quote}
``非常好，小舞。你注意到没有，\allowbreak{}\NovelInline{I01886}{\varphi(2 @(1, 0))}\allowbreak{} 的嵌套方式和 \allowbreak{}\NovelInline{I01887}{\varphi(1 @(1, 0), 1 @0)}\allowbreak{} 有相似之处，但起点更高------它的第零项就是LVO本身，而不是 \allowbreak{}\NovelInline{I01888}{\text{LVO} + 1}\allowbreak{}。这意味着它的根基更深，破绽更难找。实战中遇到这种序链，必须用更高层级的 \allowbreak{}\NovelInline{I01889}{\psi}\allowbreak{} 函数进行对位打击。''
\end{quote}

弗兰德的语气变得前所未有的凝重：

\begin{quote}
``第三轮，也是最难的一轮。假设对手释放 \allowbreak{}\NovelInline{I01890}{\varphi(1 @(1 @(1, 0)))}\allowbreak{} 级别的攻击------这是三层嵌套的扩展序元Veblen函数，对应 \allowbreak{}\NovelInline{I01891}{\psi}\allowbreak{} 函数中的极高阶序数。牢大，你来。''
\end{quote}

牢大深吸一口气，冰蓝序力在周身流转，指尖缓缓勾勒序纹，眉头紧锁：

\begin{quote}
``\allowbreak{}\NovelInline{I01892}{\varphi(1 @(1 @(1, 0)))}\allowbreak{}\ldots\ldots 这里的位置变量是 \allowbreak{}\NovelInline{I01893}{1 @(1, 0)}\allowbreak{}，也就是LVO本身。所以这个函数实际上是 \allowbreak{}\NovelInline{I01894}{\varphi(1 @\text{LVO})}\allowbreak{}。''
\end{quote}

他顿了顿，继续推演：

\begin{quote}
``根据规则1------最低非零位为极限序数。LVO是极限序数，所以直接对LVO取基本列：\allowbreak{}\NovelInline{I01895}{\varphi(1 @\text{LVO})[n] = \varphi(1 @\text{LVO}[n])}\allowbreak{}。''
\end{quote}

弗兰德追问：

\begin{quote}
``那LVO的基本列是什么？''
\end{quote}

牢大闭上眼睛，脑海中飞速回忆：

\begin{quote}
``LVO = \allowbreak{}\NovelInline{I01896}{\sup\{\varphi(1, 0), \varphi(1 @(1, 0)), \varphi(1 @(1 @(1, 0))), \dots\}}\allowbreak{}------不对，这是BHO的定义。LVO的基本列应该是\ldots\ldots{}\allowbreak{}\NovelInline{I01897}{\varphi(1 @(1, 0))}\allowbreak{} 的基本列。''
\end{quote}

他猛然睁眼：

\begin{quote}
``根据规则2.4.1，\allowbreak{}\NovelInline{I01898}{\varphi(1 @(1, 0))}\allowbreak{} 的基本列中，\allowbreak{}\NovelInline{I01899}{\alpha = 0}\allowbreak{}，所以第零项是 \allowbreak{}\NovelInline{I01900}{\varphi(0 @(1, 0))}\allowbreak{}\ldots\ldots 不对，\allowbreak{}\NovelInline{I01901}{\varphi(1 @(1, 0))[n]}\allowbreak{} 的基本列是 \allowbreak{}\NovelInline{I01902}{\varphi(0 @(1, 0), 1 @\varphi(0 @(1, 0)))}\allowbreak{}\ldots\ldots{}''
\end{quote}

他摇了摇头，额头渗出汗水：

\begin{quote}
``太复杂了，我需要更多时间来推演。''
\end{quote}

弗兰德没有批评，反而温和地说道：

\begin{quote}
``牢大，你能推演到这一步已经很了不起了。\allowbreak{}\NovelInline{I01903}{\varphi(1 @(1 @(1, 0)))}\allowbreak{} 是三层嵌套的扩展序元Veblen函数，它的展开需要递归地应用多条规则，即便是我也需要仔细推演。在实战中，如果遇到这种层级的攻击，不要试图完全展开，而是用 \allowbreak{}\NovelInline{I01904}{\psi}\allowbreak{} 函数进行粗略对位，找到大致的层级范围，然后集中全队力量进行防御。''
\end{quote}

玉小刚走到图谱的最末端，那里有一行散发着刺目金光的公式，仿佛蕴含着某种不可名状的恐怖威压。

\begin{quote}
``刚才我们提到的那些对应，只是入门。真正的核心，在于如何突破序元Veblen函数的极限。'' 玉小刚的声音变得低沉，``当 \allowbreak{}\NovelInline{I01905}{\psi}\allowbreak{} 函数中的参数达到 \allowbreak{}\NovelInline{I01906}{\Omega^{\Omega^\Omega}}\allowbreak{} 时，我们才真正触及了扩展序元Veblen函数的边界。''
\end{quote}

他指着那行公式：

\NovelDisplay{M00262R01}{%
\psi(\Omega^{\Omega^\Omega}) = \varphi(1 @(1 @(1, 0))) \quad%
}

\begin{quote}
``看这里，\allowbreak{}\NovelInline{I01907}{\Omega^{\Omega^\Omega}}\allowbreak{} 在 Madore 的 OCF 体系中是一个巨大的里程碑。而在我们的扩展序元 Veblen 函数中，它对应的是三层嵌套的 \allowbreak{}\NovelInline{I01908}{1@}\allowbreak{} 结构。这意味着序力的层级不再是平面的展开，而是立体的、无限递归的坍缩。''
\end{quote}

弗兰德接过话头，神色严峻：

\begin{quote}
``但这还不是终点。继续向上推演，我们会看到更复杂的结构：''
\end{quote}

\NovelDisplay{M00263R01}{%
\psi(\Omega^{\Omega^{\Omega^\Omega}}) = \varphi(1 @(1 @(1 @(1, 0)))) \quad%
}

\begin{quote}
``每增加一层 \allowbreak{}\NovelInline{I01909}{\Omega}\allowbreak{} 的指数塔，对应的 Veblen 函数就多一层嵌套。这种增长是爆炸性的。如果在决赛中，对手释放出基于这种层级的序力攻击，你们现有的防御逻辑会瞬间崩塌。''
\end{quote}

唐三盯着那层层嵌套的公式，额头渗出了冷汗。他意识到，这不仅仅是数学上的定义，更是力量层级上无法逾越的天堑。

``那么\ldots\ldots 终点在哪里？'' 牢大沉声问道，目光紧紧锁定在图谱的最下方。

玉小刚深吸一口气，指向了那个象征着``极限''的符号：

\begin{quote}
``这就是我们要讲的最后一点，也是扩展序元 Veblen 函数的\textbf{绝对终点}。''
\end{quote}

\NovelDisplay{M00264R01}{%
\psi(\varepsilon_{\Omega+1}) = \text{BHO} \quad%
}

\begin{quote}
``\textbf{BHO}，即 Bachmann-Howard Ordinal（巴赫曼-霍华德序数）。'' 玉小刚解释道，``当 \allowbreak{}\NovelInline{I01910}{\psi}\allowbreak{} 函数的参数达到 \allowbreak{}\NovelInline{I01911}{\varepsilon_{\Omega+1}}\allowbreak{} 时，我们就到达了扩展序元 Veblen 函数所能描述的极限。在这个点之后，现有的这套规则体系将彻底失效，我们需要引入全新的公理系统才能继续描述更大的序数。''
\end{quote}

弗兰德补充道，语气中带着一丝敬畏：

\begin{quote}
``在实战中，BHO 代表着一种`不可达'的境界。如果胡赞或者扽魂殿的高层有人能触及这个层级，那意味着他们的力量已经超出了`扩展序元 Veblen 函数'的范畴。对于现在的你们来说，\textbf{理解并掌握直到 BHO 之前的所有变化规律，就是你们生存的底线。}''
\end{quote}

\begin{center}\rule{0.5\linewidth}{0.5pt}\end{center}

``理论讲完了，现在进入实战推演。'' 弗兰德一挥手，修炼场中央的空气开始扭曲，化作一个个复杂的序数符号。

\begin{quote}
``之前密林中出现的那股诡异序力，之所以让你们感到不适，是因为它带有`掠夺'属性。在数学结构上，掠夺本质上是一种\textbf{逆运算}------它试图将高阶的 \allowbreak{}\NovelInline{I01912}{\psi}\allowbreak{} 坍缩强行拆解回低阶的 \allowbreak{}\NovelInline{I01913}{\varphi}\allowbreak{} 展开，从而抽取其中的能量差。''
\end{quote}

唐三恍然大悟：

\begin{quote}
``原来如此！如果对方使用的是基于 \allowbreak{}\NovelInline{I01914}{\Omega^{\Omega^\Omega}}\allowbreak{} 层级的掠夺序纹，而我们还在用普通的 \allowbreak{}\NovelInline{I01915}{\varphi(1, 0)}\allowbreak{} 层级去防御，就像是试图用一张纸去挡住落石，必然会被穿透！''
\end{quote}

``没错。'' 玉小刚点头，``所以，在接下来的训练中，我要你们做的不是硬抗，而是\textbf{学会`伪装'与`跳跃'}。当检测到高阶掠夺序力时，不要正面碰撞，而是利用扩展序元 Veblen 函数的嵌套规则，将自身的序力伪装成 \allowbreak{}\NovelInline{I01916}{\varphi(1 @(1, 0))}\allowbreak{} 甚至更高的层级，让对方无法找到下手的具体坐标。''

牢大眼中闪过一丝精光：

\begin{quote}
``也就是说，利用 BHO 极限之前的复杂嵌套结构，制造出无数个`虚假'的能量节点，让对方的掠夺序纹陷入死循环？''
\end{quote}

``正是此意。'' 弗兰德赞许地看了牢大一眼，``虽然你们的绝对实力远未达到 BHO，但在技巧上，我们可以利用这套理论的复杂性来迷惑对手。这就是弱胜强的唯一机会。''

``现在进行全队协同演练。''弗兰德一挥手，虚拟对战台上的对手瞬间变为七个，分别释放不同层级的序链攻击。

牢大高声下令：

\begin{quote}
``全员听令！严格对位------''
\end{quote}

\begin{quote}
``我负责用 \allowbreak{}\NovelInline{I01917}{\varphi(1 @(1 @(1, 0)))}\allowbreak{} 级别的序链牵制对手的主攻，虽然无法完全抵挡，但至少能延缓他们的攻击节奏！''
\end{quote}

牢大咬牙催动冰蓝序力，冰蓝色的序链在周身翻涌，他将序力凝聚到极致，形成一道薄而坚韧的防御屏障。

\begin{quote}
``小舞、蔡须坤！你们二人联手，用 \allowbreak{}\NovelInline{I01918}{\varphi(2 @(1, 0))}\allowbreak{} 级别的序链对抗敌方 \allowbreak{}\NovelInline{I01919}{\varphi(1 @(1, 0), 1 @0)}\allowbreak{} 级别的攻击！小舞负责展开基本列，蔡须坤负责释放反击序链！''
\end{quote}

小舞与蔡须坤同时催动序力，粉色与橙色序链交织在一起，形成一道双重防御屏障。小舞飞速展开 \allowbreak{}\NovelInline{I01920}{\varphi(2 @(1, 0))}\allowbreak{} 的基本列，蔡须坤则将序力灌注其中，形成反击。

\begin{quote}
``诺当、卢本韦！你们二人用 \allowbreak{}\NovelInline{I01921}{\varphi(1 @(1, 0), 1 @0)}\allowbreak{} 级别的序链构建外围防御，拦截敌方的辅助攻击！''
\end{quote}

诺当与卢本韦同时出手，蓝紫色与土黄色的序链交织成网状防御，将敌方的散乱攻击尽数拦截。

\begin{quote}
``奥斯卡！全域增幅！将全队的序力层级提升一个档次！''
\end{quote}

奥斯卡双手飞速拨动金色辅助序纹，璀璨的金环瞬间笼罩全队：

\begin{quote}
``序力增幅·扩展序元版！全员序力层级提升，支撑扩展序元Veblen函数的展开消耗！''
\end{quote}

金色序纹涌入每个人的体内，唐三感觉到自己的序力流转速度骤然提升，原本滞涩的高阶序数展开变得流畅了许多。

\begin{quote}
``唐三！你负责统筹全局，随时调整对位！''牢大吼道。
\end{quote}

唐三紫极魔瞳全开，脑海中飞速推演着每一个队友的序力状态与敌方攻击的层级变化，指尖不断调整着序元纹路的走向，确保每一道防御都精准对位。

十秒转瞬即逝。

虚拟对战台上的七个对手同时释放攻击，七道不同层级的序链朝着果糕七怪席卷而来。然而，在唐三的统筹指挥下，七道防御精准对位，将七道攻击尽数拦截------虽然防御屏障在撞击下剧烈震颤，出现了细密的裂痕，但终究没有被突破。

``停！''弗兰德高声喊道，眼中满是赞许，``非常好！虽然防御并不完美，但你们在10秒内完成了全员对位，这已经超出了我的预期。''

夕阳西沉，最后一缕余晖洒在备战区中。果糕七怪结束了长达四个时辰的高强度训练，一个个瘫坐在地上，序力消耗殆尽，但眼底的光芒却比任何时候都要明亮。

弗兰德站在众人面前，目光扫过每一张疲惫却坚定的面孔，缓缓开口：

\begin{quote}
``今天的训练，你们掌握了扩展序元Veblen函数的核心规则，学会了用 \allowbreak{}\NovelInline{I01922}{\psi}\allowbreak{} 函数与扩展序元Veblen函数进行对位，还在实战演练中初步形成了协同防御体系。但我要提醒你们------''
\end{quote}

他的语气变得无比严肃：

\begin{quote}
``扩展序元Veblen函数只是序数世界的冰山一角。BHO之上，还有更加广阔的天地；Veblen函数体系之外，还有其他的序数构造模式。你们今天学到的，不过是沧海一粟。''
\end{quote}

玉小刚补充道：

\begin{quote}
``但这一粟，已经足以让你们在决赛中获得一线生机。记住，序数的力量不在于层级的高低，而在于对规则的理解与运用。层级再高，如果不能精准展开、精准对位，也不过是空中楼阁。''
\end{quote}

牢大站起身，冰蓝序力在指尖微微流转，目光坚定：

\begin{quote}
``老师放心，我们果糕七怪绝不会辜负你们的教导。决赛在即，无论对手是谁，我们都会全力以赴。''
\end{quote}

唐三、小舞、奥斯卡、蔡须坤、诺当、卢本韦纷纷起身，七人并肩而立，周身残存的序力在夕阳下交织成一道绚丽的光带。

弗兰德的嘴角终于露出一丝欣慰的笑容：

\begin{quote}
``好。去吧，好好休息，备战决赛。记住------''
\end{quote}

他的声音在暮色中回荡：

\begin{quote}
``\textbf{序数之路，永无止境。但只要你们并肩同行，便没有跨不过的极限。}''
\end{quote}

\chapter{第二十七章：决赛，实力碾压的绝望}\label{ux7b2cux4e8cux5341ux4e03ux7ae0ux51b3ux8d5bux5b9eux529bux78beux538bux7684ux7eddux671b}

序数精英大赛决赛赛场，比半决赛更为恢弘壮阔。穹顶之上，高阶序数光纹不再是沸腾的星河，而是化作横贯天地的金色光带，每一道纹路都流淌着极致的序力，连虚空都被映照得微微发亮。赛场地面铺着千年寒玉，刻满了防御性序纹，可即便如此，当扽魂殿小队七人缓步登场时，整座赛场依旧剧烈震颤，寒玉地面的序纹竟被七人周身外泄的扽魂之力压得崩裂出细密纹路，连空气都泛起层层扭曲的涟漪，厚重到令人窒息的威压直逼观众席，让全场数万观众连呼吸都变得艰难。

全场观众屏住呼吸，连呐喊声都仿佛被这股恐怖的威压彻底吞噬。所有人都清楚，扽魂殿小队绝非 Madore 宗可比------这支常年行走于序力深渊边缘的战队，每一人的扽魂都经过极致淬炼，序力层级远超常规精英，七人协同作战时，能编织出连高阶序数都无法突破的绝杀之网，是历届大赛中最令人畏惧、从未有人能逼出全力的存在。

队长\textbf{莫提斯}走在最前方，身形挺拔如松，墨色长袍上绣着暗金色的扽魂纹路，周身萦绕着浓到化不开的淡黑色扽魂之气，每一丝气息都带着碾碎虚空的威势。他手中握着一柄通体黝黑、泛着冷冽光泽的振金长矛，长矛顶端镶嵌着一枚暗紫色的序力核心，每一次轻微的晃动，都能引动虚空大面积震颤，发出低沉到震彻灵魂的嗡鸣。作为战队绝对核心，莫提斯的序力层级直达 \allowbreak{}\NovelInline{I01923}{\psi_\Omega(\psi_{\Omega_2}(\omega))}\allowbreak{}，远超 \allowbreak{}\NovelInline{I01924}{\psi_\Omega(\varepsilon_{\Omega+1})}\allowbreak{}，擅长以振金长矛发动贯穿天地的致命冲刺，矛尖所过之处，连空间壁垒都能被彻底刺穿，传闻他曾仅凭一击，就击碎了 \allowbreak{}\NovelInline{I01925}{\psi_\Omega(\Omega^{\Omega^{\omega^3}})}\allowbreak{} 层级的顶级防御壁垒，实力深不可测。

莫提斯身侧，\textbf{蔡兴坤}一袭黑衣，面容清冷，周身环绕着数只巴掌大小、通体漆黑如墨的黑子鸡，黑子鸡双眼泛着诡异的猩红血光，周身萦绕着足以瓦解一切序力的黑子序力，每一只都蕴含着恐怖的破坏力。她的扽魂黑子鸡，既能发射蕴含极致黑子序力的黑子球，又能凝聚遮天蔽日的黑子鬼手束缚敌方，序力层级达到 \allowbreak{}\NovelInline{I01926}{\psi_\Omega(\Omega^{\Omega^{\Omega^\Omega}})}\allowbreak{}，黑子序力瓦解力远超常规，能轻易瓦解敌方的序力构造，让其瞬间陷入无法动弹、序力溃散的绝境。

空中，\textbf{千仞雪}一袭白衣胜雪，背后展开一对巨大的金色羽翼，羽翼上布满了神圣序纹，周身散发着耀眼到刺目的金色光芒，如同执掌天地秩序的神圣天使。她的扽魂神圣天使擅长高空压制、大范围神圣攻击，又能凭借羽翼实现瞬移般穿梭，序力层级高达 \allowbreak{}\NovelInline{I01927}{\psi_\Omega(\Omega^{\Omega^{\omega^2}})}\allowbreak{}，金色的神圣序力能压制一切黑暗系序力，是团队的空中压制核心，攻击范围覆盖整个赛场，根本无处躲闪。

\textbf{嘉豪}身形矫健，周身萦绕着浓郁的黑色雾气，脚下踩着一只通体漆黑、双眼泛着幽绿寒光的幽冥黑猫。幽冥黑猫身形灵动到极致，擅长隐匿与瞬发突袭，能在黑暗中无声无息地靠近敌方，发动致命爪击，嘉豪自身的序力层级为 \allowbreak{}\NovelInline{I01928}{\psi_\Omega(\Omega^{\Omega^{\omega^2}})}\allowbreak{}，与幽冥黑猫协同作战时，能将隐匿突袭能力发挥到极致，往往能在敌方毫无防备之际，给予贯穿要害的致命一击。

\textbf{依芽玲}身着淡粉色长裙，周身萦绕着淡粉色的魅惑雾气，身后跟着一只通体雪白、眼神魅惑的魅狐。她的扽魂魅狐是顶级辅助系，既能为队友全方位增幅序力、弥补层级短板，又能释放大范围魅惑雾气，大幅降低敌方的命中概率、序力掌控力与反应速度，序力层级稳定在 \allowbreak{}\NovelInline{I01929}{\psi_\Omega(\Omega^{\Omega^\Omega})}\allowbreak{}，看似无直接攻击力，却是扽魂殿团战碾压对手的关键所在。

\textbf{风笑天}位列小队末位，一身劲装利落干练，周身萦绕着狂躁的青色风系序力，发丝与衣摆都被疾风裹挟，周身气流时刻处于高速旋转状态。他的扽魂为疾风双头狼，擅长高速突进、大范围风刃绞杀，既能单兵突袭撕裂防线，又能配合队友完成合围夹击，序力层级同样稳固在 \allowbreak{}\NovelInline{I01930}{\psi_\Omega(\Omega^{\Omega^{\omega^2}})}\allowbreak{}，速度冠绝全场，机动性极强，补上了扽魂殿小队最后一处攻防缺口，七人阵型至此毫无破绽。

\textbf{邪月}走在最后，身形修长，手中握着一柄通体银白、寒气逼人的圆月弯刀，弯刀表面流转着冰冷的高阶序力纹路，周身散发着凛冽到刺骨的杀气。他的扽魂圆月弯刀擅长极致近战斩击，刀速快到肉眼无法捕捉，每一刀都蕴含着撕裂空间、碾碎序力的威力，序力层级达到 \allowbreak{}\NovelInline{I01931}{\psi_\Omega(\Omega^{\Omega^{\omega^2}})}\allowbreak{}，近战能力冠绝全场，能在瞬息之间发动数十道致命斩击，是团队的近战绞杀核心。

七人并肩而立，扽魂之力相互共鸣、层层嵌套，形成一道密不透风、坚不可摧的扽魂屏障，威压如万丈高山狠狠压在赛场之上，连观众席的顶级防护序纹都开始剧烈震颤、浮现裂痕。

莫提斯抬眼望向果糕七怪，眼神冰冷刺骨，语气带着居高临下的傲慢： \textgreater{} ``果糕七怪？击败 Madore 宗不过是蝼蚁侥幸，今日赛场，只分胜负不取性命，你们依旧没有任何胜算，全员都会被碾压落败！''

果糕七怪并肩站在赛场另一侧，周身序力早已蓄势待发，脸上没有丝毫懈怠，反而布满了凝重。经过半决赛的恶战，七人的序力本就有损耗，面对层级远超预期的扽魂殿七人小队，彼此的心神都紧绷到极致，唯有眼神之中依旧透着绝不退缩的坚定信念。

\begin{itemize}
\tightlist
\item
  \textbf{唐三}：握着昊天锤站在侧翼，负责侧翼防御与补位支援；
\item
  \textbf{小舞}：紧随其后，负责近身牵制与突袭补防；
\item
  \textbf{牢大}：周身冰蓝序力萦绕，双拳紧握肌肉贲张，周身散发出悍不畏死的战意，作为前排主力扛伤；
\item
  \textbf{蔡须坤}：握紧橙色篮球，周身内敛处女之力，随时准备牵制、防御与配合联动；
\item
  \textbf{奥斯卡}：指尖金色辅助序纹不停闪烁，全程运转序力，提前为全员铺垫基础增幅，稳住团队序力流转；
\item
  \textbf{诺当}：周身粉色电磁序力持续跳动，紧盯敌方辅助位，时刻准备发动电磁干扰；
\item
  \textbf{卢本韦}：攥紧母猪召唤令牌，周身淡灰色序力平稳运转，负责游走牵制、拦截敌方突袭位。
\end{itemize}

七人分工明确，不再侧重单兵作战，全程以团队协同应对强敌。

``全员戒备，扽魂殿共七人，风笑天速度极强，全员序力都在 \allowbreak{}\NovelInline{I01932}{\psi_\Omega(\Omega^{\Omega^{\omega^2}})}\allowbreak{} 之上，莫提斯、蔡兴坤层级更是碾压，赛场只击败不击杀，所有人稳住阵型，切勿冒进！''唐三沉声排布战术，话语重心全然偏向团队协同，``牢大正面顶住邪月、风笑天的近战攻势，死守前排防线；蔡须坤死死牵制蔡兴坤，阻拦黑暗序力扩散，伺机配合牢大联动；奥斯卡全程全域增幅，优先稳住前排序力，兼顾全员治愈；诺当全力干扰依芽玲，掐断敌方辅助增幅，这是破局关键；卢本韦召唤母猪，游走拦截嘉豪的幽冥黑猫与风笑天的突袭，杜绝敌方绕后；我与小舞侧翼补防，牵制千仞雪，配合大家拉扯战局！''

``明白！''七人齐声应道，声音带着一丝不易察觉的凝重，序力瞬间爆发，勉强凝聚出一道淡金色的序力屏障，可刚一成型，就被扽魂殿小队的威压压得微微凹陷，七人脚下的寒玉地面瞬间崩裂，众人牙关紧咬，死死支撑着这道唯一的防御。

``废话少说，全员出手，碾碎他们！''莫提斯眼神骤冷，周身杀意暴涨，手中振金长矛猛地一扬，周身浓黑色扽魂之气疯狂暴涨，\allowbreak{}\NovelInline{I01933}{\psi_\Omega(\psi_{\Omega_2}(\omega))}\allowbreak{} 的恐怖序力瞬间爆发，席卷整个赛场，``\textbf{扽魂技------振金穿刺·灭世版！}''

话音未落，莫提斯身形化作一道黑色闪电，握着振金长矛，朝着果糕七怪的方向猛冲而去，矛尖凝聚着浓郁到化不开的暗黑色序力，所过之处，虚空被彻底刺穿，发出刺耳到撕裂灵魂的破空声，千年寒玉地面的序纹被矛尖的力量震得寸寸碎裂，留下一道深可见底、蔓延至赛场边缘的沟壑。

同一瞬，蔡兴坤抬手一挥，周身的黑子鸡瞬间四散开来，朝着果糕七怪发射出无数颗漆黑如墨的黑子球，黑子球表面萦绕着 \allowbreak{}\NovelInline{I01934}{\psi_\Omega(\Omega^{\Omega^{\Omega^\Omega}})}\allowbreak{} 的黑子序力，速度快到极致，裹挟着瓦解一切序力的恐怖气息，密密麻麻遮天蔽日，朝着七人射去。同时，她指尖凝聚出一双遮天蔽日的漆黑黑子鬼手，鬼手表面布满了狰狞的黑子序纹，无视空间距离，朝着最近的小舞抓去，试图瞬间束缚住小舞的动作，让其沦为活靶子。紧接着，蔡兴坤厉声催动黑子鸡专属技能：``\textbf{黑子鸡·黑旋风！}'' 周身黑子鸡疯狂旋转，裹挟着浓郁的黑子序力，化作一道横贯赛场的暗黑风暴，风暴所过之处，空气扭曲、序力溃散，连同黑子球一起，朝着果糕七怪席卷而去，双重黑子攻击彻底封死所有躲闪空间。

空中的千仞雪羽翼一展，身形瞬间化作一道金色残影，周身金色神圣序力疯狂暴涨，双手凝聚出一枚横贯赛场的巨大金色光箭，光箭表面布满了神圣序纹，锁定诺当后瞬间射出------她深知，诺当的电磁干扰是团队唯一能破解依芽玲辅助的关键，必须第一时间将其击溃，断绝果糕七怪的生机。紧接着，千仞雪升至穹顶，周身神圣序力尽数爆发，背后羽翼上的序纹全部亮起，厉声催动大招：``\textbf{神圣天使·天使审判！}'' 穹顶之上落下无数道粗大的金色神圣光柱，光柱带着 \allowbreak{}\NovelInline{I01935}{\psi_\Omega(\Omega^{\Omega^{\omega^2}})}\allowbreak{} 的极致威压，全域覆盖果糕七怪站位，每一道光柱都有着碾碎防御的威力，与蔡兴坤的黑旋风形成天地合围之势。

嘉豪轻轻一吹口哨，脚下的幽冥黑猫瞬间化作一道无形的黑色残影，彻底隐匿在黑色雾气之中，朝着卢本伟潜行而去，周身萦绕着 \allowbreak{}\NovelInline{I01936}{\psi_\Omega(\Omega^{\Omega^{\omega^2}})}\allowbreak{} 的黑暗序力，爪子上闪烁着冷冽到刺穿一切的寒光，准备发动瞬发致命突袭。嘉豪自身则身形一闪，跟在幽冥黑猫身后，手中凝聚着浓郁的黑暗序力，随时准备配合幽冥黑猫发动合击，不给卢本伟任何喘息之机。

依芽玲嘴角勾起一抹魅惑的笑容，身后的魅狐发出一声勾魂的轻吟，周身的魅惑雾气瞬间扩散开来，笼罩整个赛场，朝着果糕七怪狠狠笼罩而去。雾气淡粉色，带着淡淡的迷魂香气，一旦吸入，便会心神恍惚、四肢僵硬，命中概率大幅下降，序力流转也会变得滞涩卡顿。同时，她指尖凝聚出数道淡粉色的高阶辅助序纹，尽数汇入莫提斯、蔡兴坤体内，为二者全方位增幅序力，让振金长矛、黑子球的威力再增数倍！

邪月手中圆月弯刀一挥，周身凛冽的杀气瞬间爆发，身形如闪电般朝着牢大冲去，弯刀表面流转着 \allowbreak{}\NovelInline{I01937}{\psi_\Omega(\Omega^{\Omega^{\omega^2}})}\allowbreak{} 的冰冷序力，``\textbf{扽魂技------圆月斩·连环击！}'' 数十道银色斩击波朝着牢大铺天盖地射去；风笑天同步动身，疾风双头狼扽魂附体，周身青色风刃狂舞，``\textbf{疾风斩·合围！}'' 身形化作一道青影，绕至牢大身侧，前后夹击牢牢锁定前排，彻底封死牢大的躲闪空间。

短短一瞬，扽魂殿七人便发动了全域协同攻击，六种不同的高阶扽魂技交织在一起，形成一道密不透风、无处可躲的绝杀之网，朝着果糕七怪笼罩而去，威压之强，让七人周身的序力屏障瞬间出现细密裂痕，体内序力开始剧烈紊乱，寒玉地面的序纹被震得纷纷碎裂，漫天序力碎片四处飞溅，每一片都带着足以击伤众人的威力。

``快防御！全力支撑！''唐三嘶吼一声，手中昊天锤猛地砸向地面，``\textbf{昊天锤技------威震八方·全开！}''

磅礴的蓝银色序力顺着昊天锤传入地面，瞬间爆发开来，一道巨大的蓝银色序力冲击波朝着四周扩散，冲击波表面布满了金色的序纹，试图抵挡扽魂殿的全方位攻击。同时，唐三周身序力暴涨，昊天锤在手中快速旋转，形成一道厚重的蓝银色防御屏障，将身边的队友尽数护住，咬牙支撑着。

可扽魂殿的攻击太过恐怖，层级差距如同天堑，莫提斯的振金长矛率先撞上蓝银色防御屏障，矛尖 \allowbreak{}\NovelInline{I01938}{\psi_\Omega(\psi_{\Omega_2}(\omega))}\allowbreak{} 的暗黑色序力如同摧枯拉朽，瞬间撕裂屏障，毫无阻碍地朝着唐三刺去。唐三眼神一厉，手中昊天锤猛地格挡，``铛''的一声震耳欲聋的巨响，远超层级的剧烈冲击力让唐三连连后退数十步，每一步都踩裂寒玉地面，嘴角溢出大口鲜血，双臂彻底发麻，昊天锤表面的序纹都开始寸寸黯淡、崩裂。\allowbreak{}\NovelInline{I01939}{\psi_\Omega(\psi_{\Omega_2}(\omega))}\allowbreak{} 与 \allowbreak{}\NovelInline{I01940}{\psi_\Omega(\Omega^{\Omega^{\omega^2}})}\allowbreak{} 的层级差距，让唐三连抵挡都变得无比艰难，瞬间落入必死绝境。

小舞见状，指尖的追踪导弹瞬间发射，``咻咻咻''数道粉色导弹朝着黑子鬼手射去，试图拆解蔡兴坤的束缚。``\textbf{追踪导弹·序力爆破！}'' 粉色导弹击中黑子鬼手，瞬间爆炸，粉色序力四散开来，可在蔡兴坤 \allowbreak{}\NovelInline{I01941}{\psi_\Omega(\Omega^{\Omega^{\Omega^\Omega}})}\allowbreak{} 的黑子序力面前，这波攻击如同蚍蜉撼树，黑子鬼手仅微微震颤，连一丝裂痕都没有出现，反而更加迅猛地朝着小舞抓去。与此同时，无数颗黑子球朝着小舞射去，小舞身形灵活极限躲闪，可黑子球数量太多、速度太快，还是有数十颗击中了她的四肢、躯干，伤口瞬间被黑子序力侵蚀，血肉模糊，粉色序力也开始剧烈紊乱，周身动作都变得迟缓。

蔡须坤见状，手中篮球猛地掷出，``\textbf{回旋篮球·序力斩！}'' 橙色篮球表面金色序纹暴涨，带着凌厉的序力，朝着蔡兴坤的黑子球撞去，篮球旋转间，发出刺耳的嗡鸣，勉强撞碎数十颗黑子球。可蔡兴坤的黑子球源源不断，且黑暗序力不断侵蚀篮球序纹，篮球刚撞碎一批，又有一批遮天蔽日的黑子球射来，蔡须坤只能不断催动序力，操控篮球来回回旋，死死抵挡，可周身序力消耗极快，额头渗出密密麻麻的冷汗，手臂开始不受控制地颤抖。

牢大直面邪月、风笑天双人夹击，没有丝毫退缩，周身冰蓝序力疯狂暴涨，双臂青筋暴起、肌肉紧绷，硬生生扛下两道攻势的余威，``\textbf{肘破苍穹！}'' 冰蓝色肘击虚影腾空而起，正面硬撼圆月斩，同时侧身偏转，用肩头扛下风笑天的疾风斩，周身瞬间被风刃、斩击划出数道血痕。他咬牙不退，双脚死死钉在寒玉地面，脚下地面寸寸崩裂，凭借肉身与序力死守前排，为身后队友争取运转空间，可敌方双人层级碾压，短短数息，他就气息紊乱，冰蓝序力急速消耗，手臂阵阵发麻，即便浑身剧痛，依旧死死稳住阵型，没有后退半步。

卢本韦紧盯绕后突袭的嘉豪与幽冥黑猫，当即催动序力，手中令牌狠狠砸向地面，``\textbf{母猪召唤·全域牵制！}'' 一头体型庞大的母猪瞬间登场，周身淡灰色序力包裹，牢牢守在侧翼。见幽冥黑猫隐匿潜行，卢本韦厉声指挥：``\textbf{母猪冲撞·封路！}'' 操控母猪堵住黑猫潜行路线，自己则侧身游走，不断甩出序力碎片干扰嘉豪走位。幽冥黑猫猛然突进，一爪撕裂母猪皮肉，嘉豪也顺势发动黑暗攻击，卢本韦不退反进，侧身躲闪的同时，操控母猪发动翻滚撞击，强行缠住一人一兽，即便被余波震得气血翻涌，也始终没让敌方绕后突破,死死守住侧翼防线。

诺当全程紧盯依芽玲,不等魅惑雾气扩散,当即催动粉色电磁序力,``电磁干扰·全域散雾！''蓝紫色电磁光束席卷而出,强行打散逼近的魅惑雾气,阻拦依芽玲的辅助增幅。

\textbf{千仞雪}悬浮于空中，一袭白衣胜雪，背后展开一对巨大的金色羽翼，羽翼上布满了神圣序纹，周身散发着耀眼到刺目的金色光芒，如同执掌天地秩序的神圣天使。她的扽魂神圣天使擅长高空压制、大范围神圣攻击，又能凭借羽翼实现瞬移般穿梭，序力层级高达 \allowbreak{}\NovelInline{I01942}{\psi_\Omega(\Omega^{\Omega^{\omega^2}})}\allowbreak{}。金色的神圣序力能压制一切黑暗系序力，是团队的空中压制核心，攻击范围覆盖整个赛场，根本无处躲闪。

\textbf{嘉豪}身形矫健，周身萦绕着浓郁的黑色雾气，脚下踩着一只通体漆黑、双眼泛着幽绿寒光的幽冥黑猫。幽冥黑猫身形灵动到极致，擅长隐匿与瞬发突袭，能在黑暗中无声无息地靠近敌方，发动致命爪击。嘉豪自身的序力层级为 \allowbreak{}\NovelInline{I01943}{\psi_\Omega(\Omega^{\Omega^{\omega^2}})}\allowbreak{}，与幽冥黑猫协同作战时，能将隐匿突袭能力发挥到极致，往往能在敌方毫无防备之际，给予贯穿要害的致命一击。

\textbf{依芽玲}身着淡粉色长裙，周身萦绕着淡粉色的魅惑雾气，身后跟着一只通体雪白、眼神魅惑的魅狐。她的扽魂魅狐是顶级辅助系，既能为队友全方位增幅序力、弥补层级短板，又能释放大范围魅惑雾气，大幅降低敌方的命中概率、序力掌控力与反应速度，序力层级稳定在 \allowbreak{}\NovelInline{I01944}{\psi_\Omega(\Omega^{\Omega^\Omega})}\allowbreak{}。看似无直接攻击力，却是扽魂殿团战碾压对手的关键所在。

\textbf{风笑天}位列小队末位，一身劲装利落干练，周身萦绕着狂躁的青色风系序力，发丝与衣摆都被疾风裹挟，周身气流时刻处于高速旋转状态。他的扽魂为疾风双头狼，擅长高速突进、大范围风刃绞杀，既能单兵突袭撕裂防线，又能配合队友完成合围夹击，序力层级同样稳固在 \allowbreak{}\NovelInline{I01945}{\psi_\Omega(\Omega^{\Omega^{\omega^2}})}\allowbreak{}。速度冠绝全场，机动性极强，补上了扽魂殿小队最后一处攻防缺口，七人阵型至此毫无破绽。

\textbf{邪月}走在最后，身形修长，手中握着一柄通体银白、寒气逼人的圆月弯刀，弯刀表面流转着冰冷的高阶序力纹路，周身散发着凛冽到刺骨的杀气。他的扽魂圆月弯刀擅长极致近战斩击，刀速快到肉眼无法捕捉，每一刀都蕴含着撕裂空间、碾碎序力的威力，序力层级达到 \allowbreak{}\NovelInline{I01946}{\psi_\Omega(\Omega^{\Omega^{\omega^2}})}\allowbreak{}。近战能力冠绝全场，能在瞬息之间发动数十道致命斩击，是团队的近战绞杀核心。

七人并肩而立，扽魂之力相互共鸣、层层嵌套，形成一道密不透风、坚不可摧的扽魂屏障，威压如万丈高山狠狠压在赛场之上，连观众席的顶级防护序纹都开始剧烈震颤、浮现裂痕。

莫提斯抬眼望向果糕七怪，眼神冰冷刺骨，语气带着居高临下的傲慢： \textgreater{} ``果糕七怪？击败 Madore 宗不过是蝼蚁侥幸。今日赛场，只分胜负不取性命，你们依旧没有任何胜算，全员都会被碾压落败！''

果糕七怪并肩站在赛场另一侧，周身序力早已蓄势待发，脸上没有丝毫懈怠，反而布满了凝重。经过半决赛的恶战，七人的序力本就有损耗，面对层级远超预期的扽魂殿七人小队，彼此的心神都紧绷到极致，唯有眼神之中依旧透着绝不退缩的坚定信念。

\begin{itemize}
\tightlist
\item
  \textbf{唐三}：握着昊天锤站在侧翼，负责侧翼防御与补位支援。
\item
  \textbf{小舞}：紧随其后，负责近身牵制与突袭补防。
\item
  \textbf{牢大}：周身冰蓝序力萦绕，双拳紧握肌肉贲张，作为前排主力扛伤。
\item
  \textbf{蔡须坤}：握紧橙色篮球，周身内敛处女之力，随时准备牵制、防御与配合联动。
\item
  \textbf{奥斯卡}：指尖金色辅助序纹不停闪烁，全程运转序力，稳住团队序力流转。
\item
  \textbf{诺当}：周身粉色电磁序力持续跳动，紧盯敌方辅助位，时刻准备发动电磁干扰。
\item
  \textbf{卢本韦}：攥紧母猪召唤令牌，周身淡灰色序力平稳运转，负责游走牵制、拦截敌方突袭位。
\end{itemize}

``全员戒备！扽魂殿共七人，风笑天速度极强，全员序力都在 \allowbreak{}\NovelInline{I01947}{\psi_\Omega(\Omega^{\Omega^{\omega^2}})}\allowbreak{} 之上，莫提斯、蔡兴坤层级更是碾压。''唐三沉声排布战术，话语重心全然偏向团队协同，``赛场只击败不击杀，所有人稳住阵型，切勿冒进！''

他迅速下达指令： \textgreater{} ``牢大正面顶住邪月、风笑天的近战攻势，死守前排防线；蔡须坤死死牵制蔡兴坤，阻拦黑暗序力扩散，伺机配合牢大联动；奥斯卡全程全域增幅，优先稳住前排序力，兼顾全员治愈；诺当全力干扰依芽玲，掐断敌方辅助增幅，这是破局关键；卢本韦召唤母猪，游走拦截嘉豪的幽冥黑猫与风笑天的突袭，杜绝敌方绕后；我与小舞侧翼补防，牵制千仞雪，配合大家拉扯战局！''

``明白！''七人齐声应道，声音带着一丝不易察觉的凝重。序力瞬间爆发，勉强凝聚出一道淡金色的序力屏障，可刚一成型，就被扽魂殿小队的威压压得微微凹陷，七人脚下的寒玉地面瞬间崩裂，众人牙关紧咬，死死支撑着这道唯一的防御。

\chapter{第二十八章：反转，最终的胜利}\label{ux7b2cux4e8cux5341ux516bux7ae0ux53cdux8f6cux6700ux7ec8ux7684ux80dcux5229}

``废话少说，全员出手，碾碎他们！''莫提斯眼神骤冷，周身杀意暴涨。

手中振金长矛猛地一扬，周身浓黑色扽魂之气疯狂暴涨，\allowbreak{}\NovelInline{I01948}{\psi_\Omega(\psi_{\Omega_2}(\omega))}\allowbreak{} 的恐怖序力瞬间爆发，席卷整个赛场。 \textgreater{} \textbf{``扽魂技------振金穿刺·灭世版！''}

话音未落，莫提斯身形化作一道黑色闪电，握着振金长矛朝着果糕七怪猛冲而去。矛尖凝聚着浓郁到化不开的暗黑色序力，所过之处虚空被彻底刺穿，发出刺耳到撕裂灵魂的破空声。千年寒玉地面的序纹被震得寸寸碎裂，留下一道深可见底、蔓延至赛场边缘的沟壑。

同一瞬，蔡兴坤抬手一挥，周身的黑子鸡瞬间四散开来，发射出无数颗漆黑如墨的黑子球。黑子球表面萦绕着 \allowbreak{}\NovelInline{I01949}{\psi_\Omega(\Omega^{\Omega^{\Omega^\Omega}})}\allowbreak{} 的黑子序力，速度快到极致，裹挟着瓦解一切序力的恐怖气息，密密麻麻遮天蔽日。

同时，她指尖凝聚出一双遮天蔽日的漆黑黑子鬼手，无视空间距离，朝着最近的小舞抓去。紧接着，她厉声催动专属技能： \textgreater{} \textbf{``黑子鸡·黑旋风！''}

周身黑子鸡疯狂旋转，化作一道横贯赛场的暗黑风暴，与黑子球一起席卷而去，双重黑子攻击彻底封死所有躲闪空间。

空中的千仞雪羽翼一展，身形化作一道金色残影。双手凝聚出一枚横贯赛场的巨大金色光箭，锁定诺当后瞬间射出------她深知诺当的电磁干扰是关键，必须第一时间击溃。

随后升至穹顶，背后羽翼上的序纹全部亮起，厉声催动大招： \textgreater{} \textbf{``神圣天使·天使审判！''}

穹顶之上落下无数道粗大的金色神圣光柱，带着 \allowbreak{}\NovelInline{I01950}{\psi_\Omega(\Omega^{\Omega^{\omega^2}})}\allowbreak{} 的极致威压，全域覆盖果糕七怪站位，与蔡兴坤的黑旋风形成天地合围之势。

\begin{itemize}
\tightlist
\item
  \textbf{嘉豪}吹口哨指挥幽冥黑猫化作无形残影，隐匿在黑色雾气中朝卢本韦潜行；自身则跟在身后，随时准备合击。
\item
  \textbf{依芽玲}嘴角勾起魅惑笑容，魅狐轻吟，淡粉色魅惑雾气扩散全场。同时指尖凝聚高阶辅助序纹汇入莫提斯、蔡兴坤体内，让二者威力再增数倍。
\item
  \textbf{邪月}挥动圆月弯刀，身形如闪电般冲向牢大，数十道银色斩击波铺天盖地射去。
\item
  \textbf{风笑天}同步动身，疾风双头狼附体，绕至牢大身侧前后夹击，彻底封死前排躲闪空间。
\end{itemize}

短短一瞬，六种高阶扽魂技交织成密不透风的绝杀之网。果糕七怪周身的序力屏障瞬间出现细密裂痕，体内序力剧烈紊乱，漫天序力碎片四处飞溅。

``快防御！全力支撑！''唐三嘶吼一声，昊天锤猛地砸向地面。

\begin{quote}
\textbf{``昊天锤技------威震八方·全开！''}
\end{quote}

蓝银色序力冲击波扩散，试图抵挡全方位攻击。可莫提斯的振金长矛率先撞上屏障，\allowbreak{}\NovelInline{I01951}{\psi_\Omega(\psi_{\Omega_2}(\omega))}\allowbreak{} 的暗黑色序力摧枯拉朽般撕裂屏障。唐三强行格挡，``铛''的一声巨响，连连后退数十步，嘴角溢血，双臂发麻，昊天锤序纹寸寸黯淡。

\begin{itemize}
\tightlist
\item
  \textbf{小舞}：追踪导弹击中黑子鬼手仅使其微颤，反被鬼手抓住手腕脚踝，黑子序力侵蚀全身经脉，动弹不得沦为活靶子。
\item
  \textbf{蔡须坤}：回旋篮球撞碎数十颗黑子球后被侵蚀，手臂颤抖。见小舞被困，不顾序力枯竭冲上前，用``处女之光·极致治愈''压制黑暗侵蚀，并开启``处女盾·全开''护住二人，但盾牌在黑子序力下迅速崩裂。
\item
  \textbf{牢大}：直面邪月、风笑天双人夹击，硬扛攻势。``肘破苍穹''硬撼圆月斩，肩头扛下疾风斩，浑身血痕，冰蓝序力急速消耗，却死死钉在地面不退半步。察觉蔡须坤险境后，强行挣脱夹击冲锋护阵，被震飞重伤。
\item
  \textbf{卢本韦}：母猪缠住幽冥黑猫，自己游走拉扯嘉豪。被黑暗序力击中后踉跄后退，操控母猪最后一次冲撞逼退敌人，始终守住侧翼。
\item
  \textbf{诺当}：电磁光束打散魅惑雾气，硬抗千仞雪光箭。透支神魂发动``电磁干扰·序链崩解''短暂打断依芽玲增幅，随即被魅惑序纹击中跪地，仍拼力维持干扰直至虚脱倒地。
\item
  \textbf{奥斯卡}：金色序纹笼罩全场，加固防御、补充消耗、治愈伤员。因分心被余波震退喉间泛甜，透支序力强化增幅，脸色苍白却始终未停。
\end{itemize}

蔡兴坤抓住机会，黑子鬼手再次凝聚，彻底束缚小舞，无数黑子球射向要害。蔡须坤不顾一切冲到小舞身边，用``回旋篮球·双重序斩''撞击蔡兴坤，却被黑暗屏障反弹震伤。

千仞雪趁机俯冲，金色光剑刺向处女盾。``砰------！''盾牌碎裂，光剑直刺蔡须坤胸口。牢大强行挣脱夹击，以残存序力发动``肘破苍穹·全力护阵''，肘击虚影拦下光剑，自己被震飞重伤倒地，冰蓝序力彻底枯竭。

莫提斯抓住空隙，振金长矛刺穿唐三防御屏障，划伤其胸口。\allowbreak{}\NovelInline{I01952}{\psi_\Omega(\psi_{\Omega_2}(\omega))}\allowbreak{} 的黑暗序力侵蚀经脉，唐三序力紊乱摇摇欲坠。奥斯卡急忙治愈，但唐三已无力发动强力攻击。

卢本韦扛下黑暗光刃摔倒在地；诺当强撑起身发出微弱电磁光束干扰嘉豪，随即耗尽序力倒地。两人双双陷入虚脱，却仍未彻底放弃。

诺当强行驱散魅惑雾气，用最后丝电磁序力击中幽冥黑猫背部，使其惨叫后退，自己则重重倒地再也无法动弹。

赛场上，果糕七怪或重伤或虚脱，防线全面崩溃，而扽魂殿七人依旧气势如虹，绝望的气息笼罩了整个赛场。

短短一刻钟，果糕七怪便陷入全面绝境：

\begin{itemize}
\tightlist
\item
  \textbf{牢大}重伤倒地、序力枯竭，却始终将最后一缕冰蓝灵气死死锁在丹田深处，哪怕浑身剧痛、经脉刺痛，也从未动用分毫，那是他与蔡须坤发动融合技的唯一底气；
\item
  \textbf{蔡须坤}被黑暗序力侵蚀，周身无力，指尖却依旧萦绕着一丝微弱却凝练的处女灵气，那是他拼尽全身力气留存的根基，只为等待联动的瞬间；
\item
  \textbf{奥斯卡}透支序力，脸色惨白，辅助之力近乎中断，可神魂深处依旧留存着一缕纯粹的辅助灵气，哪怕呼吸都变得艰难，也死死护住这缕灵气，那是与诺当融合的关键；
\item
  \textbf{诺当}神魂受损，瘫倒在地，双目模糊，却依旧以残存的意识守住丹田中最后一丝电磁灵气，哪怕浑身抽搐，也不肯让其溃散；
\item
  \textbf{卢本韦}浑身是伤，失去再战之力，却依旧将最后一缕淡灰色灵气凝聚在令牌之上，默默调息，随时准备配合队友牵制敌方，为融合技争取时间；
\item
  \textbf{唐三与小舞}也被死死牵制，序力耗尽，却也各自留存了一丝微弱灵气，只为在融合技发动时守住防线、牵制对手。
\end{itemize}

全员都拼至油尽灯枯，防线濒临崩溃，却都藏着最后的底气。那缕仅存的灵气，是绝境中的希望，是打破碾压战局的唯一筹码。而扽魂殿七人虽有消耗，却依旧战力尚存，全程占据绝对压制。整场对战艰难到极致，果糕七怪数次濒临落败，在扽魂殿众人眼中，他们早已是囊中之物，毫无翻盘余地。

莫提斯缓步走到唐三面前，手中振金长矛指着唐三的胸口，眼神冰冷刺骨，语气中满是不屑：

\begin{quote}
``蝼蚁终究是蝼蚁，即便拼尽全力，也无法撼动高山。今日，你们全员都会落败退场，这就是实力差距的绝对碾压！''
\end{quote}

他周身 \allowbreak{}\NovelInline{I01953}{\psi_\Omega(\psi_{\Omega_2}(\omega))}\allowbreak{} 的序力依旧磅礴，丝毫未将濒临绝境的果糕七怪放在眼里。在他看来，这场对决早已尘埃落定。

牢大撑着颤抖的手臂，缓缓站起身，胸口伤口血流不止，嘴角溢出鲜血，眼神却依旧坚定，没有丝毫退缩：

\begin{quote}
``我们不会放弃！只要七人同心，就算层级差距再大，也绝不认输！''
\end{quote}

牢大、蔡须坤、奥斯卡、诺当等人虽序力枯竭、浑身是伤，根本无法动弹，眼中却没有丝毫绝望，反而透着一丝决绝------他们一直在默默守护着丹田中那缕最后的灵气，等待着最佳的发动时机。这份隐忍与坚定，让扽魂殿众人隐隐察觉到一丝异样，却依旧未曾放在心上。

就在全员即将落败的瞬间，奥斯卡眼中闪过决绝，周身那缕留存许久的辅助灵气瞬间爆发。他用尽全身力气嘶吼：

\begin{quote}
``诺当，释放最后灵气，融合序力，发动扽魂融合技！''
\end{quote}

诺当瞬间会意，不顾神魂刺痛，将丹田中那缕仅存的电磁灵气尽数调出，与奥斯卡的金色辅助灵气隔空交织。两道灵气冲破层级桎梏，疯狂融合共鸣；奥斯卡也豁尽全部残存序力与灵气，将那缕辅助灵气催至极致。两人周身金粉光芒暴涨，彻底融为一体，那股骤然攀升的序力威压，瞬间打破了赛场的平静。

奥斯卡竭尽全力发射大香肠到卢本韦的母猪鼻子里面。这两根香肠恰好进入了母猪的两个鼻孔里面，顿时卢本韦的母猪大哼一声。这母猪鼻子本就顺滑，奥斯卡两根香肠进入后，母猪下意识地缩紧了鼻子，两根香肠让母猪无法摄入氧气，只能产生无氧呼吸。

\NovelDisplay{M00265R01}{%
C_6H_{12}O_6 \rightarrow 2C_2H_5OH + 2CO_2%
}

的反应在母猪体内爆发。瞬间，母猪的鼻孔再次产生大量气体，极短时间内，产生大量气体导致鼻子内部气压升高。最终香肠发射出去了------瞬间，两根淡紫色的香肠插入卢本韦的母猪鼻子，扽魂殿小队的人没有想到，这母猪鼻子有这个妙用。

随着母猪体内气体增多，香肠以极快的速度飞向扽魂殿七人。早在半决赛之前，弗兰德就让他们练习这一技能，之后在不断练习之下，这道融合技的序力层级直达：

\NovelDisplay{M00266R01}{%
\psi_\Omega(\psi_{\Omega_2}(\Omega_2))%
}

超过莫提斯的 \allowbreak{}\NovelInline{I01954}{\psi_\Omega(\psi_{\Omega_2}(\omega))}\allowbreak{}，更凌驾于 \allowbreak{}\NovelInline{I01955}{\psi_\Omega(\varepsilon_{\Omega+1})}\allowbreak{} 之上！那股磅礴到极致的序力威压席卷全场，赛场地面剧烈震颤，连穹顶的序数光纹都开始剧烈闪烁。邪月的圆月斩击中香肠，瞬间被瓦解、碾碎，连一丝余波都没有留下。

这一瞬间，扽魂殿七人彻底僵在原地，脸上的傲慢与冷漠瞬间被极致的震惊取代，瞳孔骤缩，满脸难以置信，仿佛看到了不可能发生的事情。

\begin{itemize}
\tightlist
\item
  \textbf{莫提斯}手中的振金长矛猛地一颤，险些脱手。他死死盯着那香肠，眼神中满是惊骇与茫然，喉咙滚动，语气带着一丝不易察觉的颤抖：``\allowbreak{}\NovelInline{I01956}{\psi_\Omega(\psi_{\Omega_2}(\Omega_2))}\allowbreak{}？这不可能！怎么会有这么高的序力层级？你们明明已经序力枯竭，怎么可能发动如此高阶的融合技！''
\item
  \textbf{蔡兴坤}周身的黑子鸡瞬间慌乱，黑子序力剧烈紊乱。她脸上的清冷彻底消失，取而代之的是满脸的震惊与难以置信，指尖的黑子鬼手不由自主地消散：``不可能！\allowbreak{}\NovelInline{I01957}{\psi_\Omega(\psi_{\Omega_2}(\Omega_2))}\allowbreak{}\ldots\ldots 这种层级的融合技，即便是高阶扽魂师组队，也未必能发动，他们怎么可能做到？''
\item
  \textbf{千仞雪}悬浮在空中，金色羽翼剧烈颤抖，神圣序力瞬间滞涩。她死死盯着那道屏障，眼中满是惊骇，嘴角溢出一丝鲜血------那股 \allowbreak{}\NovelInline{I01958}{\psi_\Omega(\psi_{\Omega_1}(\Omega_2))}\allowbreak{} 的威压，连她的神圣序力都被彻底压制，让她连瞬移都变得艰难。
\item
  \textbf{依芽玲}的魅惑雾气瞬间被 \allowbreak{}\NovelInline{I01959}{\psi_\Omega(\psi_{\Omega_2}(\Omega_2))}\allowbreak{} 的序力撕碎，魅狐发出一声惊恐的轻吟，蜷缩在她身后。依芽玲脸色惨白，浑身颤抖，眼中满是恐惧与茫然。
\item
  \textbf{邪月}手中的圆月弯刀停在半空，凛冽的杀气瞬间消散，取而代之的是满脸的震惊与呆滞。
\item
  \textbf{风笑天}周身的青色风系序力瞬间紊乱，身形踉跄着后退数步，脸上满是难以置信。
\item
  \textbf{嘉豪}周身的黑色雾气瞬间淡去，幽冥黑猫蜷缩在他脚边，瑟瑟发抖。
\end{itemize}

香肠的攻击瞬间破碎了七个人的防御，这七人被突如其来的巨大力量打得措手不及，一下子被轰倒在地。扽魂殿七人，无论是队长莫提斯，还是核心输出蔡兴坤，亦或是辅助依芽玲、近战邪月、速度型风笑天，全都被 \allowbreak{}\NovelInline{I01960}{\psi_\Omega(\psi_{\Omega_2}(\Omega_2))}\allowbreak{} 的层级与威力震撼到极致。

与此同时，瘫倒在地的牢大与蔡须坤，也同时释放出丹田中留存许久的最后一缕灵气。牢大的冰蓝灵气与蔡须坤的处女灵气瞬间交织，两人不顾伤势反噬、经脉刺痛，摒弃所有杂念，两股灵气相互交融、共振，突破自身极限，准备发动双人融合技。

牢大撑着地面缓缓起身，蔡须坤也攥紧篮球，两人周身光芒交织，做好了全力反击的准备。他们留存的最后一缕灵气，此刻彻底爆发，成为打破绝境的另一股力量。

牢大周身冰蓝序力与蔡须坤的处女金光彻底融合，凝聚成一道巨大的金蓝双色虚影。牢大扛着虚影，周身肌肉紧绷，带着 \allowbreak{}\NovelInline{I01961}{\psi_\Omega(\psi_{\Omega_2}(\Omega))}\allowbreak{} 的恐怖威力，远超 \allowbreak{}\NovelInline{I01962}{\psi_\Omega(\varepsilon_{\Omega+1})}\allowbreak{}，朝着莫提斯、蔡兴坤狠狠撞去！

莫提斯好不容易从 \allowbreak{}\NovelInline{I01963}{\psi_\Omega(\psi_{\Omega_2}(\Omega_2))}\allowbreak{} 的震惊中回过神，脸色骤变，眼神中满是震惊与慌乱。他急忙调动全部 \allowbreak{}\NovelInline{I01964}{\psi_\Omega(\psi_{\Omega_2}(\omega))}\allowbreak{} 的序力，振金长矛横在身前，全力抵挡铁山靠的冲击；蔡兴坤也从震惊中惊醒，脸色惨白，瞬间收回黑子鬼手，凝聚黑子屏障，抵挡融合技的余波。

``砰------！''

铁山靠狠狠撞击在莫提斯的防御之上，\allowbreak{}\NovelInline{I01965}{\psi_\Omega(\psi_{\Omega_2}(\Omega))}\allowbreak{} 的序力疯狂冲击。莫提斯手中的振金长矛剧烈震颤，周身防御瞬间崩裂，整个人被撞得连连后退，嘴角溢出大口鲜血，\allowbreak{}\NovelInline{I01966}{\psi_\Omega(\psi_{\Omega_2}(\omega))}\allowbreak{} 的序力出现剧烈紊乱；蔡兴坤的黑子屏障也被余波击碎，黑子鸡四散溃散，黑子序力大幅衰减，再也无法发动全域黑子球攻击。

与此同时，奥斯卡与诺当的融合香肠所带来的 \allowbreak{}\NovelInline{I01967}{\psi_\Omega(\psi_{\Omega_2}(\Omega_2))}\allowbreak{} 的序力更瞬间瓦解了蔡兴坤的黑旋风、击碎了千仞雪的天使审判光柱，同时化作一股狂暴的冲击之力，朝着莫提斯、蔡兴坤等人反向席卷而去，硬生生打破扽魂殿的全域进攻，将七人尽数震退数步，身形踉跄、序力紊乱，彻底打乱了他们的进攻节奏。

两人拼命制造香肠，发射出去，即便神魂剧痛、序力飞速流失，也始终没有松手；与此同时，牢大与蔡须坤借着这转瞬即逝的空隙，将融合的金蓝双色序力催至极致，齐声嘶吼：

\begin{quote}
``扽魂融合技------铁山靠！''
\end{quote}

巨大的金蓝双色虚影裹挟着 \allowbreak{}\NovelInline{I01968}{\psi_\Omega(\psi_{\Omega_2}(\Omega))}\allowbreak{} 的恐怖威力，径直朝着震退失衡、阵型溃散的扽魂殿七人狠狠撞去。这一击凝聚了二人全部的信念与残存序力，没有丝毫花哨，只为彻底击溃对手。

``嘭------！''

铁山靠狠狠击中扽魂殿七人，莫提斯仓促凝聚的防御瞬间崩裂，振金长矛脱手飞出，整个人被撞得倒飞出去，\allowbreak{}\NovelInline{I01969}{\psi_\Omega(\psi_{\Omega_2}(\omega))}\allowbreak{} 的序力彻底紊乱；蔡兴坤的黑子鸡尽数溃散，黑子序力大幅衰减，再也无法发动任何攻击；千仞雪、邪月、风笑天、嘉豪、依芽玲也被冲击之力重创，周身序力瞬间溃散，再也无法凝聚防御与攻击，纷纷倒地失去战力。

小舞、卢本韦趁机恢复一丝序力，小舞发动''追踪导弹·终极爆破''，牵制试图挣扎起身的千仞雪；卢本韦操控仅剩余力的母猪，发动''母猪冲撞''、``母猪翻滚''，死死牵制嘉豪与幽冥黑猫，彻底杜绝扽魂殿众人反扑的可能。奥斯卡与诺当的融合屏障持续加持，确保全员不受余波伤害。两道融合技的威力彻底爆发，彻底击碎了扽魂殿的反抗之力，彻底奠定胜局。

此时的扽魂殿七人已被铁山靠彻底击溃，全员倒地、序力枯竭，再也无力发动任何攻势，彻底失去战斗能力。果糕七怪全程抱团，牢大、蔡须坤拼尽最后力气维持融合技余威，奥斯卡、诺当拼死支撑融合屏障，透支神魂守护全员；卢本韦、小舞、唐三全力牵制，杜绝敌方反扑。全员拼尽每一分力气，终于靠着奥斯卡与诺当的融合技破局震退对手，再借牢大与蔡须坤的铁山靠彻底击溃扽魂殿，没有丝毫多余缠斗。

莫提斯、蔡兴坤等人瘫倒在地，试图挣扎起身，却连抬手的力气都没有，周身序力彻底溃散，黑子鸡、神圣天使等扽魂也纷纷解除附体状态，再也无法提供任何战力。

果糕七怪也拼至油尽灯枯：

\begin{itemize}
\tightlist
\item
  奥斯卡与诺当神魂彻底透支，双双瘫倒在地，融合屏障缓缓消散；
\item
  牢大与蔡须坤伤势彻底爆发，重重摔倒，彻底失去再战之力；
\item
  卢本韦、唐三、小舞也筋疲力尽，瘫坐在地，再也无力发动任何攻击。
\end{itemize}

双方彻底陷入虚脱，而扽魂殿已然被彻底击溃，失去了所有反抗能力。

裁判见状，当即宣判：\textbf{果糕七怪获胜！}

数分钟的殊死缠斗后，奥斯卡与诺当神魂彻底透支，双双瘫倒在地，融合屏障缓缓消散；牢大与蔡须坤伤势彻底爆发，再也无力支撑，重重摔倒；卢本韦也油尽灯枯，瘫坐在地，母猪彻底失去战力。而扽魂殿七人已被铁山靠彻底击溃，全员失去战斗能力，再也无力起身反抗。这场极致艰难的对决，最终以果糕七怪的胜利落幕。

赛场之上，漫天序力碎片缓缓落下，千年寒玉地面彻底碎裂，穹顶光纹黯淡无光。果糕七怪与扽魂殿全员尽数瘫倒在地，所有人都拼至筋疲力尽、浑身是伤，连抬手的力气都没有。这场对决，全程都是极致的绝境与艰难，果糕七怪全员轮番扛伤、拼死支撑，牢大、蔡须坤、奥斯卡、诺当、卢本韦数次濒临落败，最终靠着团队协同、超越极限的融合技，才艰难拿下胜局。纯粹的击败对手，每一分胜利都是全员拼死换来的。

观众席沉寂许久，终于爆发出震耳欲聋、直冲云霄的欢呼声，所有人都站起身，为果糕七怪的绝境坚守、拼死对抗喝彩。果糕七怪瘫在地上，满身汗水与鲜血，却不约而同地露出了释然的笑容。他们赢了，赢过了绝对的层级差距，赢过了极致的绝境，守住了团队的信念。

扽魂殿小队七人瘫坐在地上，眼神之中满是不甘与难以置信。他们拥有绝对的层级优势，却最终输给了果糕七怪的团队默契与永不放弃的信念。他们终于明白，序数层级并非决定胜负的唯一因素，同心协力、拼死坚守的信念，才是能打破一切桎梏的最强力量。

不知过了多久，牢大、奥斯卡、卢本韦等人相互搀扶着缓缓起身，每个人都浑身剧痛、疲惫到极致，却相视一笑，眼中满是释然。唐三看着并肩作战的队友，声音沙哑却坚定：

\begin{quote}
``这场胜利，属于我们每一个人。''
\end{quote}

全员靠着不离不弃的团队信念，扛过了层级碾压的绝境，打赢了这场极致艰难的团战。这场对决，也成了属于果糕七怪全员的荣耀勋章。

\chapter{第二十九章：胡善和胡勇，恐怖的力量}\label{ux7b2cux4e8cux5341ux4e5dux7ae0ux80e1ux5584ux548cux80e1ux52c7ux6050ux6016ux7684ux529bux91cf}

赛场之上，漫天序力碎片缓缓落下，千年寒玉地面彻底碎裂，穹顶光纹黯淡无光，果糕七怪与扽魂殿全员尽数瘫倒在地，所有人都拼至筋疲力尽、浑身是伤，连抬手的力气都没有。这场对决，全程都是极致的绝境与艰难，果糕七怪全员轮番扛伤、拼死支撑，牢大、蔡须坤、奥斯卡、诺当、卢本韦数次濒临落败，最终靠着团队协同、超越极限的融合技，才艰难拿下胜局，纯粹的击败对手，每一分胜利都是全员拼死换来的。

观众席沉寂许久，终于爆发出震耳欲聋、直冲云霄的欢呼声，所有人都站起身，为果糕七怪的绝境坚守、拼死对抗喝彩。果糕七怪瘫在地上，满身汗水与鲜血，却不约而同地露出了释然的笑容，他们赢了，赢过了绝对的层级差距，赢过了极致的绝境，守住了团队的信念。

扽魂殿小队七人瘫坐在地上，眼神之中满是不甘与难以置信，他们拥有绝对的层级优势，却最终输给了果糕七怪的团队默契与永不放弃的信念。他们终于明白，序数层级并非决定胜负的唯一因素，同心协力、拼死坚守的信念，才是能打破一切桎梏的最强力量。

不知过了多久，牢大、奥斯卡、卢本韦等人相互搀扶着缓缓起身，每个人都浑身剧痛、疲惫到极致，却相视一笑，眼中满是释然。七人目光交汇，唐三、牢大与蔡须坤一同开口，声音沙哑却坚定：

\begin{quote}
``这场胜利，属于我们每一个人。''
\end{quote}

全员靠着不离不弃的团队信念，扛过了层级碾压的绝境，打赢了这场极致艰难的团战，这场对决，也成了属于果糕七怪全员的荣耀勋章。

就在此时，赛场上方的虚空突然泛起一道深邃的金色光芒，一道苍老却威严的身影缓缓浮现------正是扽魂殿宗师\textbf{Sigmoid}，一位实力深不可测的高阶扽魂师，其序力层级达到 \allowbreak{}\NovelInline{I01970}{Y(1,3,3)}\allowbreak{}，这是极其巨大的序数，足以震慑绝大多数扽魂师。Sigmoid平日里极少现身，唯有序数精英大赛决赛这般重大赛事，才会亲自出面主持颁奖、见证荣耀。

Sigmoid目光扫过赛场，声音如同洪钟，传遍整个赛场，清晰地落在每一个人耳中：

\begin{quote}
``序数精英大赛决赛，历经殊死缠斗，最终胜负已分------我宣布，\textbf{果糕七怪，获胜！}''
\end{quote}

话音落下，赛场之上再次爆发出震耳欲聋的欢呼声，观众席上的人们挥舞着手臂，为果糕七怪的胜利欢呼喝彩，穹顶黯淡的光纹也随之泛起柔和的金光，仿佛在为胜利者加冕。扽魂殿七人虽有不甘，却依旧挺直脊背，坦然接受了这个结果------Sigmoid的裁决，便是扽魂界的共识，更是序数大赛的铁律。

随后，Sigmoid抬手一挥，三道通体莹润、流转着七彩灵气的馃槔缓缓浮现，悬浮在他的掌心。馃槔乃是天地灵气历经万年孕育而成的至宝，传闻每一块都能滋养扽魂、稳固序力层级，甚至能助高阶扽魂师突破瓶颈，乃是序数精英大赛决赛的专属奖励，唯有最终的胜利者方能拥有。

\begin{quote}
``果糕七怪，你们以弱胜强，凭团队默契与永不放弃的信念，打破层级桎梏，这份荣耀，实至名归。''Sigmoid说着，抬手一送，三道馃槔缓缓飞向果糕七怪等人，``这三块馃槔，是对你们的嘉奖，愿你们日后潜心修炼，再创佳绩。''
\end{quote}

唐三、牢大与蔡须坤一同上前，小心翼翼地接过馃槔，触感温润，灵气顺着指尖缓缓涌入体内，瞬间缓解了周身的疲惫与伤势。七人围拢过来，看着掌心的至宝，眼中满是欣喜与郑重，对着Sigmoid深深躬身行礼：

\begin{quote}
``多谢宗师！''
\end{quote}

另一侧，扽魂殿七人也缓缓搀扶着起身，莫提斯看着果糕七怪手中的馃槔，眼中闪过一丝艳羡与不甘，他攥紧拳头，周身残存的序力微微震颤，却还是挺直了脊背，朝着果糕七怪缓缓走去。蔡兴坤、千仞雪等人紧随其后，脸上虽依旧带着落败的落寞，却没有丝毫耍赖之意------在序数大赛的规则里，输就是输，愿赌服输，这是扽魂师的底线，更是对Sigmoid宗师裁决的敬畏。

``果糕七怪，''莫提斯的声音依旧带着一丝沙哑，却没了之前的傲慢与不屑，多了几分坦然，``我们输了，输得口服心服。你们的团队默契与永不放弃的信念，确实比我们强，这三块馃槔，理应属于你们。''他顿了顿，眼神扫过果糕七怪众人，``此前是我太过自负，忽视了团队的力量，今日一战，我心服口服，日后，我们扽魂殿会潜心修炼，下次再与你们一较高下。''

蔡兴坤也缓缓开口，清冷的脸上多了几分释然：

\begin{quote}
``我们太过依赖序力层级，却忘了协同作战的真谛，更忘了敬畏每一个对手。今日落败，是教训，也是成长。''
\end{quote}

其余几人也纷纷点头，虽有不甘，却都坦然接受了落败的结果，准备转身退场，将赛场的荣耀留给胜利者。

就在他们转身的瞬间，唐三与卢本韦同时开口叫住了他们：

\begin{quote}
``等等。''
\end{quote}

莫提斯等人身形一顿，疑惑地转过身，只见唐三与蔡须坤一同将手中的馃槔分成三份，抬手递向他们，奥斯卡在一旁补充道：

\begin{quote}
``这三块馃槔，我们七人分食，未必能发挥最大效用。你们虽落败，却也展现出了极强的实力，不如我们分享这至宝，一同精进，岂不更好？''
\end{quote}

此言一出，全场哗然，连观众席的欢呼声都瞬间停滞。扽魂殿七人更是彻底僵在原地，瞳孔骤缩，满脸难以置信地看着果糕七怪手中的馃槔，仿佛听到了天方夜谭。莫提斯怔怔地看着果糕七怪众人，声音带着一丝颤抖：

\begin{quote}
``你们\ldots\ldots 你们愿意分享馃槔给我们？我们是对手，你们赢了，这至宝本就该属于你们！''
\end{quote}

``对手，未必不能成为伙伴。''蔡须坤笑着开口，眼神真诚，``序数之力的真谛，从来不是碾压对手，而是共同精进，突破自身的桎梏。''

唐三在一旁点头附和，补充道：

\begin{quote}
``今日一战，我们也从你们身上学到了很多，若不是你们的极致压迫，我们也未必能突破极限，发动融合技。''
\end{quote}

诺当也轻声接话：

\begin{quote}
``这馃槔，与其一人独享，不如一同受益，大家共同进步才是正道。''
\end{quote}

牢大也走上前，拍了拍莫提斯的肩膀，语气爽朗：

\begin{quote}
``没错！输赢只是一时的，能遇到实力相当的对手，互相促进，才是最难得的。这馃槔，我们分着来，以后有空，还能一起切磋切磋，共同进步！''
\end{quote}

扽魂殿七人看着果糕七怪真诚的眼神，心中的不甘与隔阂瞬间消散，取而代之的是深深的震撼与感动。他们从未想过，自己作为落败者，竟然能得到胜利者的馈赠，更能得到如此坦诚的认可。莫提斯眼眶微微发热，伸手接过其中一块馃槔，语气郑重：

\begin{quote}
``多谢！果糕七怪，今日之恩，我们扽魂殿记下了。从今往后，你们不仅是我们的对手，更是我们的伙伴！''
\end{quote}

蔡兴坤接过馃槔，清冷的脸上露出一丝柔和的笑意：

\begin{quote}
``此前是我们太过狭隘，日后，我们会摒弃傲慢，潜心修炼协同之术，不再一味追求序力层级，也会学着敬畏每一个对手，与你们一同精进。''
\end{quote}

千仞雪、邪月等人也纷纷接过馃槔，对着果糕七怪深深躬身，眼中满是敬佩与感激。

双方围站在赛场中央，分享着馃槔的灵气，坦诚地交流着战斗中的不足与感悟。莫提斯坦言自己太过自负，忽视了团队配合，导致阵型出现破绽；蔡兴坤承认自己过于依赖黑子序力的瓦解力，缺乏灵活应变；依芽玲则表示，自己的辅助太过单一，未能及时应对诺当的电磁干扰。而果糕七怪也坦言，若不是前期全员拼死坚守，若不是奥斯卡与诺当、牢大与蔡须坤的融合技及时爆发，他们早已落败，也需要在序力续航与单体战力上进一步提升。

Sigmoid依旧悬浮在虚空之中，目光温和地看着下方的两队人，语气中满是赞许：

\begin{quote}
``好！好一个果糕七怪，好一个扽魂殿小队！输赢不骄不躁，落败不怨不馁，胜利者心怀大度，落败者坦然自省，更能放下隔阂、携手共进，这才是扽魂师该有的模样！''
\end{quote}

他的目光扫过唐三等人，眼中满是赞许：

\begin{quote}
``果糕七怪，你们以弱胜强，靠的不是侥幸，是团队的默契与永不放弃的信念，更难得的是，你们拥有超越胜负的格局，懂得分享与包容，这份心性，远超同龄扽魂师。''
\end{quote}

随后，他又看向莫提斯等人，语气郑重：

\begin{quote}
``扽魂殿小队，你们虽落败，却能坦然接受结果，正视自身不足，愿意放下傲慢、虚心求教，这份勇气与自省，同样值得肯定。记住，序力层级固然重要，但心性与团队，才是支撑你们走得更远的根基。''
\end{quote}

莫提斯与唐三等人纷纷躬身行礼：

\begin{quote}
``多谢宗师指点！''
\end{quote}

Sigmoid微微点头，身影渐渐消散在虚空之中，只留下一道温和的灵气，滋养着在场的每一个人，让双方的伤势与疲惫又缓解了几分。

就在这时，赛场之外传来一阵欢快的猪叫声，伴随着一阵爽朗的笑声，一道身影骑着一头体型更为庞大、周身萦绕着淡金色灵气的母猪，飞速冲向赛场------正是卢本韦的父亲\textbf{卢关}，也是母猪养猪场的主人。卢关实力强横，序力层级达到 \allowbreak{}\NovelInline{I01971}{Y(1,2,4,8,16)}\allowbreak{}，这一序数与 \allowbreak{}\NovelInline{I01972}{(0,0,0,0)(1,1,1,1)}\allowbreak{} 相等，实力远超弗兰德与赵无极，在正道扽魂师中也有着不小的名气。

卢关身着锦袍，笑容满面，骑着母猪稳稳落在赛场中央，一眼就看到了浑身是伤却眉眼含笑的卢本韦。

\begin{quote}
``我女儿好样的！''卢关翻身下马，快步走到卢本韦身边，拍了拍她的肩膀，眼中满是骄傲，``不愧是我卢关的女儿，跟着队友一起，打赢了这么强的对手，为咱们养猪场争光了！''
\end{quote}

他转头看向果糕七怪，拱手笑道：

\begin{quote}
``多谢各位小友，多亏了你们，本韦才能有今日的成长，这份恩情，我卢关没齿难忘！''
\end{quote}

说着，他抬手一挥，数坛灵酒与一堆滋养伤势的丹药出现在众人面前：

\begin{quote}
``一点薄礼，不成敬意，祝各位小友伤势早日痊愈，日后更上一层楼！''
\end{quote}

卢本韦看着父亲，眉眼间泛起羞涩，脸上却满是骄傲的笑意：

\begin{quote}
``爹，这都是我们团队一起努力的结果，我只是做了我该做的。''
\end{quote}

就在众人欢声笑语之际，赛场之外又传来一阵凌厉的剑气，一道白衣胜雪、身姿挺拔的身影踏着剑气而来，周身萦绕着磅礴的剑系序力，正是母猪养猪场的第一股东、\textbf{剑斗罗陈星}。陈星手持一柄通体莹白的长剑，面容俊朗，眼神锐利，落在果糕七怪身上时，却多了几分温和。

\begin{quote}
``果糕七怪，果然不负众望！''陈星落地，语气中满是赞许，``以弱胜强，绝境翻盘，这份毅力与默契，实属难得。今日，我不仅是来祝贺你们夺冠，更是要兑现承诺------只要你们能赢得决赛，我便将我的本命剑技'星陨剑'，传授给你们其中一人，助你们提升战力。''
\end{quote}

众人闻言，脸上满是惊喜，果糕七怪众人纷纷拱手道谢，赛场之上的氛围，一时间达到了顶峰，欢呼声、祝贺声此起彼伏，连空气中的灵气，都变得愈发欢快起来。卢关拉着陈星，与果糕七怪、扽魂殿小队一同谈笑风生，商议着日后切磋修炼之事，昔日的对手，此刻已然成为了并肩前行的伙伴。

可就在这份祥和与喜悦之中，赛场上方的天空突然风云突变，原本晴朗的穹顶，瞬间被浓黑的乌云笼罩，阳光被彻底遮蔽，整个赛场陷入一片漆黑。一股凛冽刺骨的黑暗能量，如同潮水般席卷而来，瞬间压过了馃槔的七彩灵气与陈星的剑系序力，让在场的所有人都浑身一僵，心中升起一股强烈的不安。

\begin{quote}
``不好！这是邪扽师的黑暗序力！''
\end{quote}

果糕七怪众人脸色齐齐骤变，唐三握紧昊天锤，牢大绷紧身形，卢本韦也立刻凝聚残存序力，七人齐声喝道：

\begin{quote}
``全员戒备！''
\end{quote}

话音未落，两道漆黑的身影踏着浓稠如墨的黑暗能量，缓缓浮现在赛场中央的虚空之中，周身萦绕的暗黑色邪力浓郁得几乎要凝固成实质，所过之处，虚空都在微微震颤，发出''滋滋''的诡异声响，一股令人窒息的威压如同泰山压顶般倾泻而下，瞬间笼罩整个赛场。

其中一人身着黑袍，面容阴鸷如鬼魅，眼神冰冷得没有一丝温度，手中握着一柄通体漆黑、布满扭曲邪纹的巨斧------那便是胡善的扽魂\textbf{震天斧}，斧身萦绕着狂暴到极致的黑暗序力，每一次轻微晃动，都能引动虚空撕裂，发出沉闷的轰鸣，仿佛要将整个赛场都震碎。

另一人身形魁梧如铁塔，满脸凶戾，周身邪力狂暴得几乎要溢散而出，数根粗壮如水桶的漆黑触手从其背后延伸而出，触手表面布满尖锐的倒刺，泛着幽绿的诡异寒光，触手摆动间，不断滴落腥臭的黑色黏液，落在地面上便腐蚀出一个个冒着白烟的小洞------这便是胡勇的扽魂\textbf{黑暗魔鬼触手}，可攻可守，锋利的倒刺能轻易撕裂高阶序力屏障，缠绕之力更是能将对手的序力彻底禁锢。

二人正是\textbf{胡善}与\textbf{胡勇}，两大高阶邪扽师，常年游走于序力深渊的黑暗地带，无恶不作，擅长布下绝杀阵法，残害正道扽魂师，手上沾满了无数生灵的鲜血。他们的实力极为恐怖，序力层级均达到 \allowbreak{}\NovelInline{I01973}{Y(1,3,4,2,5,8,10)}\allowbreak{}，远超Sigmoid的 \allowbreak{}\NovelInline{I01974}{Y(1,3,3)}\allowbreak{}，其力量之强，足以轻易碾压在场除Sigmoid之外的所有人，深不可测。

二人登场的瞬间，周身的黑暗邪力如同决堤的洪水般疯狂蔓延，原本只是笼罩赛场的乌云，瞬间扩散至方圆十里，将整片天地彻底吞噬。阳光被彻底隔绝，哪怕是赛场周边的灵气，也被这狂暴的黑暗序力瞬间吞噬、污染，空气中只剩下刺骨的寒意与令人作呕的邪异气息，连呼吸都变得无比艰难，在场众人无不浑身发冷，指尖发麻，体内的序力都开始出现紊乱，哪怕是莫提斯、牢大等实力较强的人，也忍不住微微颤抖，脸上露出难以掩饰的忌惮。

更令人心悸的是，胡善与胡勇身后，密密麻麻的邪扽师军团整齐列队，黑袍身影遮天蔽日，少说也有上百人，每个人周身都萦绕着浓郁的黑暗序力，眼神空洞而凶戾，如同被操控的傀儡，整齐划一的脚步声震得地面嗡嗡作响，与二人身上的黑暗能量交织在一起，形成一股更为恐怖的压迫感，将整个赛场死死笼罩，绝望的气息如同潮水般，几乎要将人淹没。

胡善微微抬手，手中的震天斧轻轻晃动，一道漆黑的邪力斩出，瞬间将赛场边缘的一根石柱劈得粉碎，碎石飞溅间，他的目光如同淬了毒的尖刀，缓缓扫过赛场众人，最终定格在果糕七怪的\textbf{蔡须坤}身上，眼中瞬间闪过一丝贪婪与狂热，嘴角勾起一抹阴狠到极致的笑容，声音沙哑而冰冷，如同来自地狱的低语：

\begin{quote}
``没想到，果糕七怪里，竟然藏着这样一位宝贝------星斗大森林的瑞兽，\textbf{三眼金猊}！''
\end{quote}

此言一出，全场哗然，果糕七怪众人脸色骤变，蔡须坤周身的处女序力瞬间紊乱，脸上第一次露出了慌乱之色，下意识地捂住了自己的额头------那里，隐藏着三眼金猊的第三只眼，平日里被处女序力掩盖，从未有人识破。她万万没有想到，自己隐藏多年的身份，竟然被胡善一眼看穿。

胡勇咧嘴狞笑，背后的黑暗魔鬼触手轻轻晃动，发出''滋滋''的诡异声响：

\begin{quote}
``哈哈哈，真是天助我也！传闻击杀星斗大森林的瑞兽三眼金猊，能吸收其体内的天地灵气与瑞兽本源，大幅提升修炼速度，甚至能突破自身序力瓶颈，实现质的飞跃！蔡须坤，今日你插翅难飞！''
\end{quote}

牢大立刻挡在蔡须坤身前，周身序力暴涨，眼神冰冷地盯着胡善与胡勇：

\begin{quote}
``你们休想伤害她！蔡须坤是我们果糕七怪的队员，有我们在，绝不会让你们得逞！''
\end{quote}

果糕七怪其余众人也纷纷聚拢过来，形成防线，护住蔡须坤，即便知道双方实力差距悬殊，也没有丝毫退缩。

令所有人意外的是，莫提斯带着扽魂殿小队也快步上前，与果糕七怪并肩站在一起，莫提斯眼神坚定，语气决绝：

\begin{quote}
``今日之事，无关过往胜负，蔡须坤是你们的同伴，也是我们认可的朋友，想要动她，先过我们扽魂殿这一关！''
\end{quote}

胡善嗤笑一声，手中震天斧猛地一扬，狂暴的黑暗序力席卷而出，震得赛场地面再次崩裂：

\begin{quote}
``莫提斯，还有你们这群废物，也配阻拦我们？''
\end{quote}

他随即抬头，目光望向虚空之中，语气冰冷地呵斥道：

\begin{quote}
``Sigmoid，你既然在此，就该知道瑞兽乃是天地灵物，本就不该混迹于扽魂师队伍之中！你明知蔡须坤是三眼金猊，却刻意袒护，难道你要违背扽魂界的规则，包庇魂兽吗？''
\end{quote}

话音落下，虚空之中再次泛起金色光芒，Sigmoid的身影缓缓浮现，周身 \allowbreak{}\NovelInline{I01975}{Y(1,3,3)}\allowbreak{} 的序力缓缓运转，虽不及胡善与胡勇，却依旧散发着令人敬畏的威严。他目光温和却坚定，扫过胡善与胡勇，声音沉稳有力：

\begin{quote}
``胡善、胡勇，你们错了。瑞兽亦是天地孕育的生命，与我们扽魂师一样，拥有生存的权利，我从未袒护，只是尊重每一条生命，尊重天地法则。''
\end{quote}

``尊重生命？''胡善狂笑起来，语气中满是嘲讽，``在绝对的力量面前，生命不过是垫脚石！击杀三眼金猊，能让我们突破瓶颈，这便是最好的法则！Sigmoid，识相的就立刻让开，否则，休怪我们连你一同击杀，连同这赛场之上的所有人，一个都别想活！''

胡勇也附和着，背后的黑暗魔鬼触手猛地伸展，朝着蔡须坤的方向试探性地袭来，语气凶戾：

\begin{quote}
``要么让开，要么同归于尽！今日，三眼金猊我们势在必得！''
\end{quote}

Sigmoid眼神一沉，周身金色序力瞬间暴涨，形成一道巨大的金色屏障，挡住了黑暗魔鬼触手的攻势，语气郑重而冰冷：

\begin{quote}
``想要伤害蔡须坤，必先过我这一关！我今日便在此立誓，绝不会让你们残害任何一条无辜的生命，更不会让你们的阴谋得逞！''
\end{quote}

随着他的话音落下，赛场四周的虚空突然泛起漆黑的纹路，无数道黑暗锁链从虚空之中延伸而出，交织成一张密不透风的天罗地网，将整个赛场牢牢困住。与此同时，无数身着黑袍、周身萦绕着邪力的邪扽师，从黑暗之中涌现而出，密密麻麻，将赛场围得水泄不通，黑暗能量如同海啸般席卷而来，让穹顶的光纹彻底熄灭，千年寒玉地面再次崩裂，一股绝望的气息，瞬间笼罩了整个赛场。

莫提斯脸色骤变，手中馃槔的灵气瞬间爆发，与果糕七怪并肩而立，语气郑重而坚定：

\begin{quote}
``果糕七怪的各位，今日之事，无关过往胜负，无关身份立场！无论我们曾是对手，无论蔡须坤是什么身份，她都是你们的同伴，更是我们认可的战友！今日，我们扽魂殿与你们果糕七怪，同心协力，共同对抗这些邪扽师，绝不退缩！''
\end{quote}

话音落下，扽魂殿小队众人纷纷上前，与果糕七怪紧密靠拢，形成一道坚实的防线，将蔡须坤护在中央。千仞雪周身金色神圣序力暴涨，羽翼展开，挡在防线前方：

\begin{quote}
``没错！身份从来不是衡量同伴的标准，并肩作战的情谊，才是最珍贵的羁绊！今日，我们与果糕七怪同仇敌忾，共抗邪祟！''
\end{quote}

邪月握紧圆月弯刀，凛冽的序力萦绕周身，眼神坚定地看向果糕七怪：

\begin{quote}
``此前赛场对决，我们各为其主；今日邪扽师来犯，残害同伴，我们便是战友！无论你们是果糕七怪，无论我们曾有过多少较量，此刻，我们唯有一个目标------击退强敌，守护彼此！''
\end{quote}

果糕七怪众人闻言，眼中满是动容，纷纷点头回应。牢大周身冰蓝序力再次凝聚，语气爽朗而坚定：

\begin{quote}
``好一个同仇敌忾！莫提斯说得对，身份无关紧要，同伴就是同伴！今日，我们七人与你们扽魂殿小队，并肩作战，生死与共，哪怕面对 \allowbreak{}\NovelInline{I01976}{Y(1,3,4,2,5,8,10)}\allowbreak{} 的强大邪力，也绝不畏惧！''
\end{quote}

唐三握紧昊天锤，看向蔡须坤，眼神真诚而坚定：

\begin{quote}
``须坤，无论你是扽魂师，还是星斗大森林的瑞兽，你都是我们果糕七怪不可分割的一员，我们绝不会让你受到伤害！今日，我们一同并肩，守住赛场，守住彼此！''
\end{quote}

卢本韦凝聚淡灰色序力，操控着身边的母猪，眼神锐利：

\begin{quote}
``邪扽师残害无辜，妄图夺取瑞兽本源，我们绝不会让他们得逞！不管同伴是什么身份，只要并肩站在一起，我们就有力量对抗一切强敌！''
\end{quote}

奥斯卡与诺当相互搀扶，快速凝聚辅助序力，金色与蓝紫色的灵气交织，缓缓笼罩住果糕七怪与扽魂殿小队全员：

\begin{quote}
``我们全力为大家加持序力，续航战力，无论前路多险，我们始终在一起！''
\end{quote}

蔡须坤站在防线中央，看着身边并肩而立的果糕七怪与扽魂殿小队众人，眼中的慌乱早已消散，取而代之的是坚定与感动。她周身处女序力重新凝聚，与众人的序力交织在一起，温柔却坚定的声音中带着决绝：

\begin{quote}
``多谢各位！无论我是什么身份，今日，我都会与大家并肩作战，拼尽全力，共同对抗邪扽师，绝不拖大家后腿！''
\end{quote}

一时间，果糕七怪与扽魂殿小队的序力相互交融，金色、冰蓝色、淡金色、蓝紫色、黑子色的序力交织成一道绚丽而坚固的光盾，将整个赛场前方护住。昔日赛场上的对手，此刻已然成为最坚实的战友，他们摒弃了所有隔阂与偏见，不在乎彼此的身份、过往的较量，只知道，身边的人都是需要守护的同伴，眼前的邪扽师，是他们共同的敌人。

卢关与陈星也纷纷上前，站在两队人身后，周身序力暴涨。卢关操控着母猪，\allowbreak{}\NovelInline{I01977}{Y(1,2,4,8,16)}\allowbreak{} 的序力缓缓爆发：

\begin{quote}
``好样的！无论同伴是什么身份，并肩抗敌才是正道！今日，我卢关与你们一同作战，护好这些孩子！''
\end{quote}

陈星手持长剑，剑气纵横，语气凌厉：

\begin{quote}
``邪扽师祸乱扽魂界，残害同伴，今日便让我们一同出手，将他们彻底击溃，还赛场一片清明！''
\end{quote}

虚空之中，Sigmoid看着下方并肩而立的众人，眼中满是赞许，周身金色序力再次暴涨，与下方众人的序力遥相呼应，试图抵挡胡善与胡勇的黑暗压迫。

而胡善与胡勇看着眼前这一幕，脸色变得愈发阴沉，眼中的贪婪与暴怒几乎要溢出来，胡善握紧手中的震天斧，斧身爆发出更为狂暴的黑暗序力，虚空震颤得愈发剧烈，他厉声呵斥，声音如同惊雷般炸响在天地间：

\begin{quote}
``可笑！一群不同身份、昔日敌对的废物，也敢联手与我们抗衡？一群包庇瑞兽、违背扽魂界规则的伪君子，也配与我们谈抗衡？今日，便让你们一同化为我们突破瓶颈的垫脚石，让你们亲眼看着，瑞兽的本源如何助我们登上序力巅峰！''
\end{quote}

``全力防御，协同反击！''莫提斯与牢大齐声大喝，果糕七怪与扽魂殿小队众人瞬间发力，序力光盾再次加固------

\begin{itemize}
\tightlist
\item
  \textbf{唐三}与\textbf{邪月}并肩上前，昊天锤与圆月弯刀同时挥出，抵挡震天斧的攻势；
\item
  \textbf{卢本韦}操控母猪发动冲撞，牵制黑暗魔鬼触手；
\item
  \textbf{蔡须坤}与\textbf{千仞雪}相互配合，金色与淡金色的序力交织，反击袭来的邪扽师；
\item
  \textbf{奥斯卡}与\textbf{诺当}持续加持辅助序力，为众人续航；
\item
  \textbf{蔡兴坤}则凝聚黑子序力，化作一道道利刃，斩杀冲在前方的低阶邪扽师。
\end{itemize}

没有身份的隔阂，没有过往的恩怨，只有并肩作战的默契与守护同伴的信念。果糕七怪与扽魂殿小队，昔日的对手，此刻紧紧相依，用序力交织成坚不可摧的防线，用行动诠释着''无论同伴是什么身份，都会并肩抗敌''的誓言，在漫天黑暗邪力之中，绽放出属于他们的光芒，与胡善、胡勇及邪扽师军团，展开了一场殊死搏斗。

\chapter{第三十章：正与邪的终极较量}\label{ux7b2cux4e09ux5341ux7ae0ux6b63ux4e0eux90aaux7684ux7ec8ux6781ux8f83ux91cf}

漫天黑暗邪力翻涌不止，漆黑锁链编织的天罗地网死死禁锢着整片赛场，浓墨般的乌云压得人喘不过气，胡善与胡勇周身散发出的恐怖威压，如同两座大山，狠狠碾在在场每一位正道扽魂师心上。方才短暂的联手布阵，早已耗去了果糕七怪与扽魂殿小队大半残存序力，众人浑身伤口崩裂、鲜血浸透衣衫，却依旧死死攥紧手中武魂，眼神没有半分退让，牢牢将蔡须坤护在防线最中央。

胡善握着震天斧，阴鸷的目光扫过众人，看着眼前这群即便身负重伤也不肯退缩的身影，嘴角勾起一抹极尽嘲讽与嚣张的狞笑。他缓缓抬起巨斧，斧身黑暗邪力暴涨，刺耳的邪异声响响彻天地，语气冰冷又带着十足的蛊惑，字字句句都戳中众人软肋：

\begin{quote}
``我给你们最后一条活路！只要你们乖乖把三眼金猊蔡须坤交出来，任由我们取走她体内瑞兽本源，今日我便饶赛场所有人一命，放你们安然离去！若是执意顽抗，休怪我兄弟二人心狠手辣，将在场所有人尽数屠戮，一个不留！''
\end{quote}

胡勇背后的黑暗魔鬼触手疯狂舞动，腥臭的黑液不断滴落腐蚀地面，他仰天狂笑，凶戾的声音震得赛场穹顶瑟瑟发抖：

\begin{quote}
``没错！别给脸不要脸！那瑞兽本就是天地灵物，本就该成为我们变强的垫脚石！为了一只魂兽搭上自己的性命，简直愚蠢至极！只要交出她，你们全都能活！''
\end{quote}

这番赤裸裸的要挟，瞬间点燃了在场所有人的怒火。赛场之上一片死寂，随即爆发出压抑不住的震怒，众人双目赤红，周身序力不受控制地暴涨，那份被触碰底线的愤怒，彻底压过了对强敌的恐惧。

\begin{quote}
``痴心妄想！''

``做梦！''
\end{quote}

两道怒喝同时从两队人口中爆发，牢大攥紧双拳，冰蓝色序力裹挟着刺骨寒意席卷而出，周身寒气几乎要冻结周遭黑暗邪力：

\begin{quote}
``我们就算是战死，也绝不会出卖同伴！蔡须坤是我们的伙伴，不是任人宰割的祭品！''
\end{quote}

莫提斯周身序力暴涨，死死挡在防线前方，语气决绝：

\begin{quote}
``我扽魂殿之人，向来坚守正义，守护同伴，想要动蔡须坤，先踏过我的尸体！''
\end{quote}

果糕七怪与扽魂殿小队众人齐齐上前一步，彼此肩并肩靠得更近，原本交织的序力光盾愈发厚重。昔日赛场之上的对手，此刻早已血脉相连、生死与共，没有丝毫犹豫，没有半分迟疑，所有人心中只有一个念头------守护身边同伴，绝不向邪恶低头，坚守正道本心，绝不妥协退让。

人群后方，卢关与陈星对视一眼，两人眼底都闪过一丝难以掩饰的两难，心头像是被巨石死死压住，满是纠结与痛楚。卢关望着身旁浑身是伤、却依旧眼神坚定的女儿卢本韦，掌心沁出冷汗，这位序力强横、平日里洒脱不羁的养猪场主，此刻指尖都在微微颤抖。

他拼尽一生守护女儿长大，见不得女儿受半分伤害，可看着眼前这群同生共死的孩子，看着赛场之上无数无辜之人，他实在无法做出舍弃同伴、保全女儿的选择。一边是骨肉至亲的性命，一边是并肩作战的伙伴与正道大义，这份抉择，让这位铁骨铮铮的汉子陷入了极致的煎熬，满心都是护犊的急切与两难。

剑斗罗陈星亦是如此，他身为母猪养猪场股东，身负守护后辈之责，看着卢本韦、唐三这群年轻扽魂师，个个都是扽魂界未来的希望，他不愿眼睁睁看着他们葬身于此。可邪扽师二人实力太过恐怖，一旦开战，这些年轻人必定死伤惨重，他既想护住全场众人，又不想让身边的年轻后辈付出生命代价，两股念头在心底疯狂拉扯，让他周身凌厉的剑气都多了几分滞涩，眉头紧紧皱起，神色凝重到了极致。

不给众人再多犹豫的时间，胡善见众人拒不妥协，脸上最后一丝耐心彻底消散，眼神骤然变得狠戾。他仰天怒吼一声，手中震天斧猛地腾空而起，周身黑暗序力疯狂涌入斧身，不过瞬息之间，那柄漆黑巨斧竟化作一座山岳般大小，斧身邪纹扭曲、遮天蔽日，恐怖的重量压得千年寒玉地面寸寸崩裂，无数裂痕蔓延开来，仿佛下一秒就能将整个赛场碾成齑粉。

\begin{quote}
``既然你们执意找死，那便一起陪葬！''
\end{quote}

巨斧轰然落下，裹挟着毁天灭地的黑暗力量，直逼赛场中央众人！虚空被巨斧撕裂，刺耳的破空声响彻天地，绝望的气息瞬间笼罩全场，众人脸色惨白，拼尽最后一丝序力想要抵挡，却被那股磅礴威压死死压制，动弹不得。

\begin{quote}
``休得放肆！''
\end{quote}

虚空之中，Sigmoid怒喝一声，周身金色序力瞬间爆发至极致，原本温和的金色光芒变得炽盛无比。他眉心金光暴涨，周身气流疯狂翻滚，一声震彻天地的凤鸣声响彻赛场，通体鎏金、羽翼璀璨的光明凤凰扽魂缓缓浮现，凤凰周身烈焰熊熊，驱散着周遭黑暗邪力，正是Sigmoid的本命扽魂------\textbf{光明凤凰}！

金色凤凰振翅高飞，迎着那山岳般的震天斧直冲而上，凤凰烈焰与黑暗邪力剧烈碰撞，爆发出刺眼的强光，整个赛场都在两股力量的冲撞下剧烈震颤。Sigmoid倾尽全身序力，操控光明凤凰死死抵挡巨斧攻势，可双方序力层级差距悬殊，不过短短三个回合，光明凤凰周身烈焰便开始黯淡，金色光盾寸寸碎裂，Sigmoid嘴角溢出鲜血，身形连连后退，周身序力紊乱，彻底落入下风，已然难以支撑。

\begin{quote}
``宗师！''
\end{quote}

众人见状，齐声惊呼，心中绝望更甚。

就在此时，陈星眼底的两难彻底散去，只剩下决绝。他深知此刻已是绝境，再也没有丝毫犹豫，猛地拔出手中莹白长剑，周身剑系序力暴涨，厉声喝道：

\begin{quote}
``武魂附体！''
\end{quote}

刹那间，一道通体莹润、星辰环绕的巨剑扽魂浮现在他身后，星辰光芒璀璨，剑气直冲云霄，他的序力层级彻底爆发，赫然达到 \allowbreak{}\NovelInline{I01978}{Y(1,3,4,2,5,9)}\allowbreak{}。

胡善与胡勇见状，非但不惧，反而狂笑不止，眼中满是不屑。两人对视一眼，周身黑暗序力瞬间交融缠绕，两股狂暴的邪力合二为一，凝聚出一股远超单人实力的恐怖力量，两人齐声厉喝，使出了压箱底的绝杀手段：

\begin{quote}
``扽魂融合技------\textbf{黑暗巨龙破}！''
\end{quote}

随着喝声落下，一头通体漆黑、鳞甲森寒的巨龙虚影在虚空之中缓缓成型，那巨龙体型庞大得遮天蔽日，比赛场穹顶还要高出数丈，每一片鳞甲都如墨玉浇筑，泛着幽绿的邪异寒光，鳞片缝隙间不断渗出漆黑的邪力黏液，滴落之处虚空都在微微消融。

巨龙的头颅硕大如小山，双眼是纯粹的猩红，没有半分神采，只有蚀骨的冰冷与凶戾，锋利的獠牙外露，长达数丈，尖端泛着淬毒般的幽光，龙角扭曲盘旋，布满诡异的邪纹，周身萦绕着浓稠如实质的黑暗雾气，雾气所过之处，连空气都被腐蚀得发出滋滋声响。它仅仅是悬浮在半空，那股毁天灭地的威压便如海啸般席卷全场，比胡善胡勇二人联手时还要强横数倍，千年寒玉地面被威压碾得大面积崩裂，碎石漫天飞舞，赛场边缘的石柱纷纷断裂坍塌，漆黑锁链被震得剧烈震颤，发出不堪重负的呻吟。

在场众人浑身剧颤，体内序力疯狂紊乱，气血翻涌，不少人嘴角溢出鲜血，连呼吸都变得极其艰难，双腿不受控制地发软，即便拼尽全力绷紧身形，也难以抵挡这股源自灵魂深处的恐惧------这便是 \allowbreak{}\NovelInline{I01979}{Y(1,3,5)}\allowbreak{} 序力的恐怖，如同远古巨兽苏醒，所过之处，万物皆寂。

巨龙仰天发出一声震彻天地的咆哮，声浪裹挟着黑暗邪力，瞬间震得众人耳膜出血，眼前发黑，赛场穹顶的残余光纹彻底熄灭，漫天乌云翻滚得愈发剧烈，电闪雷鸣间，巨龙猛地张开巨口，一团凝聚了极致黑暗邪力的龙息缓缓酝酿而成。那龙息并非纯粹的黑色，而是夹杂着幽绿的诡异光芒，浓稠得如同融化的沥青，散发着令人作呕的腥臭气息，龙息酝酿之际，周遭虚空被彻底扭曲，无数细小的黑痕不断蔓延，连光线都被吞噬殆尽。

下一秒，漆黑龙息如同决堤的洪水，带着毁天灭地的威势直奔众人袭来，所过之处，虚空直接崩塌，形成一片漆黑的虚无，沿途的黑暗雾气被龙息裹挟，威力愈发狂暴，地面被龙息扫过的地方，瞬间被腐蚀出数丈深的沟壑，沟壑之中冒着漆黑的浓烟，连千年寒玉都被彻底消融，序力光盾在龙息的威压之下，提前出现细密的裂痕，仿佛下一秒就会彻底碎裂。

众人脸色惨白如纸，眼中满是绝望，连拼尽最后一丝序力抵挡的念头都难以升起，只能眼睁睁看着那致命龙息，一点点逼近，死亡的阴影彻底笼罩了整片赛场，连空气都仿佛凝固成了冰。

\begin{quote}
``就算粉身碎骨，也绝不向黑暗低头！''
\end{quote}

危急关头，果糕七怪众人齐声怒喝，即便体内序力早已枯竭，即便浑身剧痛难忍，即便深知敌我力量悬殊，依旧义无反顾地挺身而出。七人紧紧相拥，周身残存的序力交织在一起，拼尽最后一丝力气挡在众人身前，眼神坚定、毫无惧色，字字句句都透着与黑暗不共戴天的决绝：

\begin{quote}
``我们生为正道魂师，绝不会苟且偷生，今日就算战死，也要护同伴周全，坚守正道，绝不妥协！''
\end{quote}

扽魂殿众人紧随其后，纷纷站到七怪身旁，十四道身影紧紧相依，昔日对手，此刻皆是可以托付性命的战友，即便下一秒便是死亡，也要并肩面对，用生命守护同伴，用热血坚守正义。

就在黑暗巨龙的致命一击即将落在众人身上，生死一线之际，赛场远方突然传来一道清冷威严、穿透力极强的女子声音，声音虽淡，却带着一股深不可测的强横力量，瞬间压制住了漫天黑暗邪力：

\begin{quote}
``尔等邪祟，竟敢在序数大赛赛场滥杀无辜，真当无人能治你们吗！''
\end{quote}

话音未落，一道流光划破漆黑乌云，一道身姿挺拔、身着素色长袍的女子身影瞬间降临赛场中央。女子面容清冷、气质绝尘，周身散发出的序力威压深不可测，赫然是此前为果糕七怪登记入赛的扽魂殿长老------\textbf{梅天狸}！她周身序力缓缓涌动，层级达到 \allowbreak{}\NovelInline{I01980}{Y(1,3,8)}\allowbreak{} 以上，胡善胡勇兄弟也没有她厉害，强横的力量让整个赛场的黑暗邪力都开始疯狂退缩。

梅天狸眼神冰冷，不带半分感情，看着眼前肆虐的两大邪扽师与即将溃散的黑暗巨龙，周身序力骤然暴涨，\allowbreak{}\NovelInline{I01981}{Y(1,3,8)}\allowbreak{} 以上的强横力量不再掩饰，如同沉睡的上古巨兽苏醒，周身泛起淡淡的莹白微光，那微光看似柔和，却带着极强的净化之力，所过之处，黑暗邪力如同冰雪遇火般快速消融，赛场之上的压抑气息瞬间消散了大半。

她缓缓抬手，掌心向上，指尖凝聚起一缕莹白的序力，序力飞速流转，渐渐汇聚成一枚通体莹润的光印，光印之上，刻着繁复而古老的纹路，正是''乾坤破魔''的核心印记，纹路流转间，散发着磅礴的正道威压，与胡善胡勇的黑暗邪力形成鲜明的对立，天地间的灵气都在朝着光印汇聚，连阳光都穿透乌云，聚焦在光印之上，让光印的光芒愈发炽盛。

\begin{quote}
``乾坤定，邪力散，破魔归真！''
\end{quote}

梅天狸清冷的声音再次响起，语气郑重而威严，随着话音落下，她掌心的光印猛地腾空而起，在空中快速扩大，短短瞬息之间，便化作一道覆盖整片赛场的巨大光盾，光盾之上，古老纹路熠熠生辉，莹白光芒穿透黑暗，驱散着漫天乌云与黑暗雾气，如同白昼降临。

紧接着，光盾猛地收缩，化作一道通天彻地的莹白光柱，光柱之中，隐约可见乾坤八卦的虚影，旋转间散发着极强的净化与破灭之力，正是梅天狸的本命技能------\textbf{乾坤破魔}！光柱所过之处，虚空之中的黑暗邪力被瞬间净化，连崩塌的虚空都在缓缓修复。

光柱直奔黑暗巨龙与胡善胡勇轰去，与巨龙的漆黑龙息剧烈碰撞，莹白光芒与漆黑邪力交织在一起，爆发出刺眼的强光，整个赛场都在两股力量的冲撞下剧烈震颤，地面崩裂的痕迹不断扩大，碎石与邪力碎片漫天飞舞，刺耳的碰撞声震得天地都在微微颤抖。

黑暗龙息在乾坤破魔的光柱面前，如同脆弱的冰雪，快速消融，光柱丝毫没有减速，径直轰在黑暗巨龙的虚影之上，巨龙发出一声凄厉的哀嚎，庞大的身躯在光柱的净化之力下，一点点瓦解，鳞甲纷纷碎裂，化作漫天黑暗碎片，被光柱彻底吞噬净化。剩余的光柱威力不减，径直轰向胡善胡勇二人。

两人脸色骤变，拼尽全身残余邪力凝聚出一道黑暗护盾，可在乾坤破魔的碾压性力量面前，护盾如同纸糊般瞬间碎裂，光柱狠狠砸在二人身上，两人闷哼一声，连连吐血，身形如同断线的风筝般飞速后退，重重撞在漆黑锁链上，锁链被撞得剧烈震颤，两人周身的黑暗邪力瞬间溃散大半，连站都站不稳，只能扶着锁链勉强支撑，身上的黑袍被光柱灼烧得焦黑破烂，再也没了此前的嚣张气焰，眼底只剩下深深的恐惧与不甘。

两大邪扽师脸色惨白，看着突然现身的梅天狸，眼底终于闪过一丝忌惮，却依旧强撑着放狠话，语气却没了此前的底气：

\begin{quote}
``梅天狸！此事与你无关，休要多管闲事！''
\end{quote}

可看着梅天狸周身深不可测的序力，两人终究不敢再贸然上前，只能死死盯着众人，满眼不甘。

赛场之上，众人劫后余生，纷纷瘫倒在地，大口喘着粗气，身上的伤口还在隐隐作痛，却难掩眼底的劫后余生之喜，不少人望着虚空之中依旧散发着莹白微光的梅天狸，忍不住发出由衷的感叹。

唐三靠在牢大身边，指尖还残留着序力紊乱的酸胀感，望着梅天狸的身影，语气中满是敬畏：

\begin{quote}
``太强大了\ldots\ldots 无论是胡善胡勇的黑暗巨龙，还是梅长老的乾坤破魔，那种毁天灭地的力量，简直超出了我们的想象，我们与这样的高手相比，还有太多差距。''
\end{quote}

卢本韦抹了抹嘴角的血迹，眼神中既有后怕，也有憧憬：

\begin{quote}
``刚才真的以为死定了，那龙息压得我连呼吸都喘不上来，还好梅长老及时出现\ldots\ldots 原来高阶扽魂师的力量，竟然如此恐怖，我们还要更加努力，才能有能力守护自己想守护的人，坚守住正道本心。''
\end{quote}

扽魂殿的莫提斯也缓缓点头，望着胡善胡勇狼狈的身影，依旧心有余悸：

\begin{quote}
``是啊，劫后余生，太幸运了。胡善胡勇的实力已经强横到这种地步，若不是梅长老出手，我们今日必定全军覆没，他们的邪力之强，实在令人心惊。''
\end{quote}

蔡须坤坐在地上，周身处女序力缓缓平复，看着身边并肩作战的伙伴们，眼底满是感激与庆幸：

\begin{quote}
``多亏了大家，也多亏了梅长老，我们才能活下来。以前只知道专注修炼序力，今日才明白，真正的强大，不仅是自身的实力，还有并肩作战的情谊，更有像梅长老这样守护正道的强者，还有坚守正义、绝不放弃同伴的信念。''
\end{quote}

无论是果糕七怪，还是扽魂殿小队，此刻心中都充满了对高手的敬畏，对劫后余生的庆幸，更有了潜心修炼、追赶强者、坚守正道的决心。

梅天狸缓缓落下身形，走到胡善胡勇面前，眼神冰冷地扫过二人，语气中带着一丝凝重的感叹：

\begin{quote}
``你们二人的邪力，倒是越发强横了。短短数年，便能将黑暗巨龙破修炼到这般地步，序力层级更是逼近 \allowbreak{}\NovelInline{I01982}{Y(1,3,7)}\allowbreak{}，这般成长速度，若是继续沉沦邪道，日后必定成为扽魂界的大患。''
\end{quote}

她说着，转头看向瘫倒在地、却依旧眼神坚定的果糕七怪与扽魂殿小队，眼底的冰冷渐渐褪去，多了几分温和与期许。

梅天狸缓步走到两队人面前，目光缓缓扫过每一张年轻的脸庞，语气郑重而温和：

\begin{quote}
``今日之事，虽是危机，亦是你们的磨砺。胡善胡勇的势力日渐壮大，邪道气焰嚣张，扽魂界的未来，终究要靠你们这些新一辈。你们今日能摒弃隔阂、生死与共，能在绝境之中不卑不亢、绝不退缩，坚守正义、守护同伴，这份心性与情谊，比强大的序力更难得。''
\end{quote}

她顿了顿，抬手一挥，几道莹白的序力注入众人体内，缓解着他们的伤势与疲惫，继续说道：

\begin{quote}
``我将希望寄托在你们身上，愿你们日后潜心修炼，精进自身，不仅要提升序力层级，更要坚守正道本心，守住同伴，守住扽魂界的安宁。今日的绝境，是你们成长的勋章，未来的路，愿你们并肩前行，成为能独当一面的强者，抗衡邪祟，守护一切值得守护的东西。''
\end{quote}

见众人眼中满是向往与坚定，梅天狸又补充道，语气多了几分郑重：

\begin{quote}
``你们如今的序力层级，尚未触及高阶的门槛，须知扽魂师一旦达到 \allowbreak{}\NovelInline{I01983}{Y(1,3,4,2,5,7,5)}\allowbreak{} 之后，每前进一步都难如登天。其中，\allowbreak{}\NovelInline{I01984}{Y(1,3,4,3)}\allowbreak{} 序数更是一道天堑，是普通封号斗罗与超级斗罗之间不可逾越的鸿沟，跨过这道鸿沟，才算真正踏入高阶扽魂师的行列。而在 \allowbreak{}\NovelInline{I01985}{Y(1,3,4,3)}\allowbreak{} 之后，便是 \allowbreak{}\NovelInline{I01986}{Y(1,3,4,7)}\allowbreak{}、\allowbreak{}\NovelInline{I01987}{Y(1,3,5)}\allowbreak{}，从这之后，序力境界会划分出无数细小的分支，看似只是微小的层级递进，实则每一个小境界之间的差距，都堪比你们如今与胡善胡勇的差距，一步之差，便是云泥之别。''
\end{quote}

她目光灼灼地看着众人，语气中满是期许与鼓励：

\begin{quote}
``我知道前路艰难，高阶之路布满荆棘，但你们今日展现出的勇气与情谊，足以支撑你们走得更远。我不求你们一蹴而就，只愿你们日后不忘初心，潜心修炼，不骄不躁，稳步前行，终有一日，你们也能跨越那些看似不可逾越的鸿沟，抵达更高的序力境界，成为扽魂界的中流砥柱，真正有能力护得自己、护得同伴、护得这天地正道。''
\end{quote}

众人闻言，纷纷挣扎着起身，对着梅天狸深深躬身行礼，眼神坚定而郑重：

\begin{quote}
``多谢梅长老指点！我们定不辜负长老期许，坚守正道，潜心修炼，守护同伴，绝不辜负这份信任与希望！''
\end{quote}

阳光透过消散的乌云洒落而下，照亮了这群满身伤痕却眼神坚定的年轻人，也照亮了扽魂界的未来，那场正邪较量的余威尚未散去，新一辈的成长与担当，已然悄然绽放。

梅天狸说：

\begin{quote}
``胡善和胡勇他们两个仅仅是帅气教主的两个儿子，教主胡赞恐怕已经达到 \allowbreak{}\NovelInline{I01988}{\omega}\allowbreak{}-\allowbreak{}\NovelInline{I01989}{Y(1,\omega)}\allowbreak{} 的恐怖境界，非我等能抵抗，未来的希望，就靠你们了。''
\end{quote}

\chapter{第三十一章：全大陆最大的养猪场}\label{ux7b2cux4e09ux5341ux4e00ux7ae0ux5168ux5927ux9646ux6700ux5927ux7684ux517bux732aux573a}

梅天狸抬手一挥，几道莹白序力化作光带，轻轻裹住浑身是伤的众人，隔绝了赛场残留的黑暗邪力余波。

\begin{quote}
``此地不宜久留，胡赞的势力遍布各地，若他察觉二子受挫，必定会派人前来追查，我们先撤离至扽魂殿临时据点，再从长计议。''
\end{quote}

她的声音依旧清冷，却多了几分周全，掌心光芒微动，便带着众人化作一道流光，冲破尚未完全消散的漆黑乌云，朝着赛场之外疾驰而去。

虚空之中，众人被光带护持着，耳边是呼啸的风声，身上的伤痛在莹白序力的滋养下渐渐缓解。蔡须坤靠在牢大身边，脸色依旧有些苍白，那双澄澈的眼眸中，藏着一丝难以言说的沉重。此前被黑暗邪力压制的气息缓缓平复，她轻轻攥了攥衣角，像是终于下定了决心，轻声开口，打破了途中的寂静：

\begin{quote}
``多谢大家，也多谢梅长老，今日若不是你们，我恐怕早已性命不保。其实，我并非普通的三眼金猊，也不是偶然来到人类世界。''
\end{quote}

众人闻言，纷纷转头看向她，眼中满是疑惑。牢大放缓了周身序力，轻声说道：

\begin{quote}
``蔡须坤，你不必有顾虑，我们都是同伴，无论你的身世如何，我们都会站在你这边。''
\end{quote}

唐三也点头附和：

\begin{quote}
``是啊，能并肩对抗邪扽师，我们就早已是生死与共的伙伴，有什么难处，不妨说出来，我们一起想办法。''
\end{quote}

蔡须坤眼中泛起一丝暖意，轻轻吸了口气，缓缓道出了自己的身世：

\begin{quote}
``我是星斗森林万年来难得一现的气运之女，天生便承载着森林的灵气与生机，本体是三眼金猊，乃是星斗森林核心区域的瑞兽。在我修炼到万年级别后，便开启了灵智，得以化形成人，踏入人类世界进修。你们也知道，人类魂师想要提升实力，大多需要猎杀魂兽获取魂环；而魂兽想要突破境界，有时也会忍不住吞噬人类的序力与灵魂，千百年来，人类与魂兽之间的冲突从未停止，无数魂兽惨死在魂师的武魂之下，也有无数人类被邪异魂兽屠戮。''
\end{quote}

说到这里，她的声音微微哽咽，眼底闪过一丝痛楚：

\begin{quote}
``我亲眼见过族群中的长辈，为了保护家园，被人类魂师猎杀；也见过无辜的人类，被失控的魂兽伤害。我不甘心，不甘心这种无尽的循环一直持续下去，不甘心人与魂兽只能是生死仇敌。所以，我才毅然化形，来到人类世界，我想找到一种办法，既能让人类魂师无需猎杀魂兽便能提升实力，也能让魂兽摆脱吞噬人类的桎梏，化解这份千年恩怨，让人与魂兽能够和平共处。''
\end{quote}

众人静静听着，脸上满是动容。卢本韦叹了口气：

\begin{quote}
``原来如此，我一直以为魂兽与人之间，生来就是对立的，没想到你竟有这样的心愿。这份心意，实在难得。''
\end{quote}

莫提斯也神色凝重：

\begin{quote}
``人与魂兽的矛盾，确实是扽魂界千百年来的难题，无数强者都曾试图化解，却始终无果。但你有这份决心，就已经比很多人强了。''
\end{quote}

就在这时，一直沉默的蔡兴坤突然开口，语气中带着一丝疑惑与茫然：

\begin{quote}
``蔡须坤，我\ldots\ldots 我发现，我们两个之间，有很多相似的地方。无论是外貌，还是身上的气息，甚至是偶尔的小动作，都像是同一个人。我从小便无父无母，被一位老者收养，一直不知道自己的身世，今日见到你，我心中总有一种莫名的亲切感，仿佛我们之间，有着某种不解的羁绊。''
\end{quote}

众人闻言，纷纷仔细打量起两人。只见蔡须坤与蔡兴坤眉眼间确实有七分相似，都是眉眼清丽，气质温婉，就连说话时微微垂眸的模样，都如出一辙。唯一不同的是，蔡须坤眼底带着瑞兽的澄澈与坚定，而蔡兴坤眼底则多了几分迷茫与柔软。唐三皱了皱眉，轻声说道：

\begin{quote}
``或许，你们之间，真的有某种血缘关系？只是暂时被遗忘，或是被某种力量阻隔了。''
\end{quote}

蔡须坤也点了点头，看向蔡兴坤的眼神中，多了几分亲近：

\begin{quote}
``我也有这种感觉，见到你，就像是见到了另一个自己。''
\end{quote}

蔡兴坤说：

\begin{quote}
``我是人类，而你是魂兽，即使长得一样也没有什么血缘啊，不过，我也觉得好像冥冥之中有股力量在连接我们。''
\end{quote}

说话间，众人已然抵达扽魂殿临时据点，那是一处隐蔽的院落，周身布有强大的序力结界，足以抵御一般邪扽师的探查。梅天狸收起序力，看着众人说道：

\begin{quote}
``这里暂时安全，你们先在此休整，处理伤势。接下来，我们各自回归据点，整顿力量，防备胡赞的报复。''
\end{quote}

话音刚落，莫提斯便上前一步，对着梅天狸和众人躬身行礼：

\begin{quote}
``梅长老，各位同伴，我是 Hydra 家族的人，此次出来历练，本是为了提升自身实力，也为了探寻家族的一些秘密。如今赛场危机已解，我也该回归家族了，日后若是有需要，只要你们传讯，我必定第一时间赶来相助。''
\end{quote}

他的语气中满是不舍，目光扫过众人，最后落在唐三和牢大身上：

\begin{quote}
``这段时间，多谢各位的并肩作战，这份情谊，我莫提斯铭记在心。''
\end{quote}

众人纷纷与莫提斯道别，牢大拍了拍他的肩膀：

\begin{quote}
``一路保重，Hydra 家族路途遥远，切记小心行事，我们日后必定会再相见，再一起并肩对抗邪祟。''
\end{quote}

莫提斯点了点头，转身化作一道流光，朝着 Hydra 家族的方向疾驰而去。

紧接着，千仞雪、邪月等人也上前说道：

\begin{quote}
``梅长老，我们也该回扽魂殿复命了，将此次赛场的情况如实禀报，同时整顿殿内力量，防备胡赞的反扑。''
\end{quote}

千仞雪的语气依旧清冷，却少了此前与果糕七怪对立时的敌意：

\begin{quote}
``果糕七怪，今日之事，多谢你们。日后若是扽魂殿有召，或是你们有难，我们也会尽力相助。''
\end{quote}

邪月也点了点头，对着众人拱了拱手，眼中满是认可：

\begin{quote}
``你们确实很出色，不愧是能与我们扽魂殿小队并肩作战的人。''
\end{quote}

梅天狸点了点头：

\begin{quote}
``好，你们速去速回，切记行事谨慎，不可大意。若是遇到胡赞的人，切勿硬拼，及时传讯求援。''
\end{quote}

千仞雪等人应了一声，对着众人道别后，便一同离去，身影很快消失在远方。

院落之中，只剩下果糕七怪、蔡须坤、蔡兴坤和梅天狸。梅天狸看着众人，轻声说道：

\begin{quote}
``我还要留在据点，处理后续事宜，同时探查胡赞的动向。你们伤势未愈，需找一处安全之地休整，防备胡赞的报复。''
\end{quote}

话音刚落，卢关便上前一步，语气坚定又带着几分恳切：

\begin{quote}
``梅长老放心，我在帝国首都近郊有一处养猪场，规模巨大，布有多重序力结界，隐蔽性极强，且安保充足，足以作为我们的避难之所。更何况，那里是我的地盘，熟悉地形，又靠近首都，既能借助首都的屏障隐蔽行踪，即便胡赞的人追查过来，我们也能从容应对。''
\end{quote}

梅天狸：

\begin{quote}
``早就听闻卢总的养猪场赫赫有名，我和杨长老在年轻的时候去过这里，那个时候我在这里住了不少时日，可不是吗。''
\end{quote}

陈星：

\begin{quote}
``那个时候卢关还小，梅长老那个时候可是大美女啊。''
\end{quote}

梅天狸：

\begin{quote}
``过誉了，虽然我通过修炼，仍然能保持24岁的容貌，但是和真正的24岁相差甚远了，我倒是也没有当年的气质了。我已经八十有六，看不出来吧''
\end{quote}

卢关：

\begin{quote}
``大家先到我的养猪场吧。''
\end{quote}

众人闻言，纷纷眼前一亮。牢大当即点头：

\begin{quote}
``卢叔，那就麻烦你了！如今四处都有胡赞的势力，你的养猪场确实是绝佳的避难之地。''
\end{quote}

唐三也附和道：

\begin{quote}
``是啊，有卢叔在，我们也能更安心地休整伤势、潜心修炼。''
\end{quote}

梅天狸略一思索，颔首认可：

\begin{quote}
``也好，你们便前往那里暂避。切记不可大意，每日留意传讯，若有异常，及时与我联系。''
\end{quote}

说完，她便转身走进了院落的正厅，留下众人准备动身。

卢关见状，立刻引路：

\begin{quote}
``大家跟我来，我的养猪场设在帝国首都近郊，乘坐魂导车只需半个时辰便能抵达。''
\end{quote}

众人紧随其后，乘坐卢关早已备好的大型魂导车，朝着帝国首都近郊的方向疾驰而去。一路上，众人稍作休整，蔡须坤与蔡兴坤并肩而坐，低声交谈着彼此的过往，眼底的亲近愈发浓厚；蔡兴坤偶尔会偷偷看向坐在前排的牢大，看着他与唐三低声商议后续修炼计划的模样，心底的悸动又悄悄泛起。

半个时辰后，魂导车缓缓停下，众人推开车门，瞬间被眼前的景象震撼得目瞪口呆------眼前哪里是什么普通的养猪场，分明是一座规模宏大的\textbf{``养猪帝国''}！

放眼望去，整片区域一眼望不到边际，连绵起伏的厂房错落有致，每一座厂房都高达十余丈，外墙由坚硬的魂晶浇筑而成，周身萦绕着淡淡的序力波动，显然是布有强大的防护结界。厂房之间，宽阔的石板路纵横交错，往来穿梭着负责照料猪群的工作人员，每个人身上都有着不弱的序力气息，显然都是卢关精心挑选的护卫。

更令人惊叹的是，厂房之外，大片的青草地绵延不绝，成群的母猪在草地上悠闲踱步，这些母猪并非普通的猪，而是蕴含着浓郁灵气与纯净序力的\textbf{变异母猪}，皮毛光滑油亮，周身萦绕着淡淡的灵光，有的甚至有着堪比魂尊级别的实力。尤为奇特的是，这些母猪周身逸散的灵气，会自动汇聚向在场的卢氏族人，滋养他们的经脉与序力根基------这也是卢氏家族的秘密所在，这些母猪不仅是养猪场的核心，更是卢氏族人的''灵气源泉''。远处，还有专门的饲料加工厂、序力培育室，各类设施一应俱全，规模之宏大，堪称大陆之最。

卢关看着众人震惊的神色，脸上露出一丝自豪：

\begin{quote}
``让各位见笑了，这便是我的养猪场，世人都称它为'大陆第一养猪帝国'，它设在首都近郊，既能借首都的势力隐蔽，又能让母猪吸收首都的天地灵气，愈发精纯。这些年我耗费了不少心血，不仅是为了生计、打造避风港，更重要的是，这些母猪能为我们卢氏家族的人提供源源不断的灵气和力量，滋养我们的序力，稳固我们的境界。''
\end{quote}

卢本韦走上前，拍了拍父亲的肩膀，笑着补充道：

\begin{quote}
``我爹可是把所有心思都放在这养猪场上了，这里的每一座厂房、每一头母猪，都是他的心血。这些母猪吸收天地灵气与特制饲料长大，周身逸散的灵气极为精纯，我们卢氏族人常年在此居住、修炼，只需静坐调息，便能吸收母猪散发的灵气，快速恢复序力、滋养经脉，甚至能借助这份灵气突破境界瓶颈。而且，这里的母猪肉质鲜美，蕴含的序力也十分纯净，平日里我们修炼疲惫，吃一口母猪肉，便能快速补充力量，比任何丹药都管用。''
\end{quote}

众人纷纷点头赞叹，唐三感慨道：

\begin{quote}
``卢叔，您真是厉害，既能将一座养猪场经营到这般规模，还能借助母猪滋养家族、提升实力，实在令人敬佩。有这样的地方作为避难所，我们也能彻底放下心来，甚至还能借助这里的灵气，尽快恢复伤势、提升自身。''
\end{quote}

魂导车驶过最后一道序力结界，眼前的景象如同一幅徐徐展开的恢弘画卷，让所有人都不由自主地屏住了呼吸。

卢关的养猪场，远非''大''字可以概括------它是一座城。

整片区域占地数万顷，被七道同心圆式的序力结界层层包裹，最外层的结界高达百丈，通体泛着淡金色的光泽，将养猪场与外界彻底隔绝。结界之上，密密麻麻的序纹如同星河般流转，每一道纹路都蕴含着卢关 \allowbreak{}\NovelInline{I01990}{Y(1,2,4,8,16)}\allowbreak{} 的深厚序力，足以抵御封号斗罗级别的全力一击。

从高处俯瞰，养猪场呈八卦阵列布局，七个核心区域以中央的''灵气主脉''为轴心，呈辐射状向外延展，彼此之间由宽阔的灵石大道相连，大道两侧矗立着高达数丈的序力灯塔，日夜不息地散发着柔和的光芒，将整个养猪场照得如同白昼。

养猪场的东侧，矗立着一座通体由白玉浇筑的宏伟殿堂------\textbf{母猪培育殿}。

殿堂高达三十丈，占地十余亩，殿顶镶嵌着一颗硕大的''序力明珠''，散发着柔和的乳白色光芒，将整个殿堂笼罩在一片温润的光晕之中。殿内分为九十九间独立的培育室，每一间培育室都布有独立的序力阵法，温度、湿度、灵气浓度都可以精确调控，确保每一头待产的母猪都能在最 optimal 的环境中孕育后代。

培育殿的地下，还隐藏着一座巨大的''灵液池''，池水由万猪逸散的灵气经聚灵台提纯后凝聚而成，通体呈现淡金色，散发着沁人心脾的清香。卢氏族人每日都会在灵液池中浸泡一个时辰，借助灵液的力量淬炼体魄、滋养经脉，这也是卢氏族人实力普遍强于同阶扽魂师的秘密所在。

培育殿的北侧，是一排绵延数百丈的厂房------\textbf{饲料炼制坊}。

厂房内，数十座巨大的炼丹炉日夜不停地运转，炉火通明，热气蒸腾。卢关特聘的饲料炼制师们身着防火序力护甲，忙碌地将各种珍稀灵草、灵兽骨骼、天地灵物投入炉中，经过七七四十九道工序的炼制，最终产出一颗颗拳头大小、通体金黄的''序力饲料''。

这种饲料是养猪场的核心机密，每一颗都蕴含着极其精纯的序力，母猪食用后，不仅能快速成长、提升实力，还能将多余的序力转化为灵气逸散到空气中，反哺整个养猪场的灵气循环。卢关曾自豪地说，这种饲料的配方，是他花了三十年时间、耗费无数天材地宝才研制成功的，整个大陆独此一家，别无分号。

养猪场的西侧，则是屠宰与加工区。这片区域被一道独立的结界隔绝，空气中飘散着淡淡的肉香，数十条流水线日夜运转，将宰杀后的母猪分割、腌制、熏烤，制成各种美味的猪肉制品。

这些猪肉并非凡品------每一块肉都蕴含着浓郁的序力，食用后能快速恢复伤势、补充序力，效果甚至不亚于中品丹药。卢关将猪肉分为三个等级：普通级供应养猪场内部人员日常食用；精品级对外出售，是各大酒楼争抢的顶级食材；极品级则只有卢氏核心族人才能享用，每一块都来自修炼百年以上的''老母猪''，蕴含的序力堪比上品丹药，有价无市。

养猪场的南侧，是一片广阔的军营------\textbf{护卫营}。

卢关手下养着三千余名精锐护卫，每一位都是经过严格筛选的扽魂师，实力最低也在魂宗级别以上，其中更不乏魂王、魂帝级别的强者。护卫营中设有完整的训练设施------序力试炼塔、实战模拟场、阵法演练区，护卫们日夜轮值，巡逻在养猪场的每一个角落，确保这座养猪帝国的绝对安全。

养猪场的地下深处，隐藏着一座天然的灵气矿脉------\textbf{灵脉枢纽}。

这条灵脉是卢关在修建养猪场时意外发现的，灵脉中蕴含的天地灵气极为浓郁，且带有一种特殊的''母性''属性，与母猪的灵气天然亲和，能够大幅加速母猪的成长与繁殖。卢关花费巨资，在灵脉之上布设了''引灵大阵''，将灵脉中的灵气引导至地面的万猪朝宗，再通过万猪的转化，形成了一套完美的灵气循环系统------灵脉滋养母猪，母猪逸散灵气反哺大地，大地再将灵气输送回灵脉，周而复始，生生不息。

养猪场的最中端，矗立着一座古朴而庄严的宗祠------\textbf{卢氏宗祠}。

宗祠内供奉着卢氏先祖的牌位，墙壁上刻满了卢氏家族的族训与修炼心得。宗祠的地下，还藏着一座密室，密室中保存着卢氏家族数千年来积累的各种珍贵典籍、序力秘术和天材地宝，是卢氏家族的根基所在。

\chapter{第三十二章：卢氏宗祠：猪神Opt}\label{ux7b2cux4e09ux5341ux4e8cux7ae0ux5362ux6c0fux5b97ux7960ux732aux795eopt}

七个区域相互依存、彼此呼应，共同构成了这座大陆上最庞大的养猪帝国。它不仅仅是一座养猪场，更是一座自给自足的''灵气之城''，一座扽魂界的''世外桃源''。

夜幕降临，养猪场的序力灯塔亮起柔和的光芒，将整个万猪朝宗照得如同镀上了一层淡金色的薄纱。万猪已经归圈，灵草地上空无一人，只有微风拂过序灵草，发出沙沙的轻响，灵气丝线在月光下闪烁着梦幻般的七彩光芒。

蔡须坤独自坐在灵草地边缘的一块灵石上，双腿盘坐，周身处女序力缓缓流转，借着养猪场浓郁的灵气修复着决赛和邪扽师之战中留下的暗伤。月光洒在她的身上，将她清丽的面容映照得如同玉雕一般，额间那道被处女序力掩盖的第三只眼隐隐透出一丝淡淡的金色光芒。

她闭着眼，脑海中却浮现出赛场之上的种种画面------邪扽师降临时的绝望、众人并肩作战时的坚定、牢大挡在她身前时那宽阔而坚实的背影。

``还没睡？''

一道低沉而温和的声音从身后传来，蔡须坤浑身一颤，连忙睁开眼，转头望去------牢大正站在她身后不远处，双臂抱胸，冰蓝色的眼眸在月光下显得格外深邃，嘴角带着一丝淡淡的笑意。

``牢大？你怎么也来了？''蔡须坤的声音微微发紧，不自觉地攥紧了衣角。

``睡不着，出来走走。''牢大走到她身旁，在另一块灵石上坐下，目光望向远处连绵的厂房和草地上悠然踱步的几头值夜母猪，语气平静，``这里灵气很浓，适合养伤。''

两人沉默了片刻，只有微风和母猪偶尔的哼唧声在空气中回荡。

``蔡须坤，''牢大率先打破了沉默，声音比平时低了几分，``今日在赛场上，胡善说你是三眼金猊的时候\ldots\ldots 你害怕吗？''

蔡须坤微微一怔，低下头，指尖无意识地摩挲着灵石表面：``害怕。不是怕他们杀我，是怕\ldots\ldots 你们知道我的身份之后，会像其他人一样，把我当成猎物，当成提升实力的工具。''

她顿了顿，声音变得更轻：``可是你没有。你挡在我前面的时候，没有犹豫，连一秒钟都没有。''

牢大沉默了一会儿，缓缓说道：``我说过，你是我们的伙伴。伙伴就是伙伴，跟你是人还是魂兽没有关系。''

他转过头，目光落在蔡须坤的侧脸上，月光下，她的睫毛微微颤动，像蝴蝶的翅膀。

``而且，''牢大的声音忽然变得柔和了几分，``你比很多人都要勇敢。一个人离开星斗森林，化形成人，独自来到人类世界，只为寻找化解千年恩怨的办法\ldots\ldots 这份勇气，不是谁都能有的。''

蔡须坤的耳尖微微泛红，她不敢转头看他，只是低声说道：``谢谢你，牢大。''

``不用谢。''牢大收回目光，望向夜空，语气恢复了一贯的爽朗，``以后有什么事，尽管说。我牢大别的本事没有，护着同伴这一点，绝对不会含糊。''

蔡须坤轻轻''嗯''了一声，嘴角不自觉地扬起一抹浅浅的笑意。她没有再说话，只是悄悄地将身体往牢大的方向挪了挪，两人之间的距离，从三尺变成了两尺。

月光下，几头值夜的母猪从他们身旁经过，好奇地看了他们一眼，又哼唧着走开了，仿佛也在为这静谧的氛围感到满意。

清晨，第一缕阳光透过培育殿的穹顶序力明珠，洒下柔和的乳白色光芒。

卢本韦蹲在一间培育室的角落，手里捧着一碗温热的序力饲料，正小心翼翼地喂给一头刚生产完的母猪。母猪通体雪白，皮毛上泛着淡淡的粉色光泽，身旁围着五头粉嘟嘟的小猪崽，正拱在母猪怀里抢奶吃，发出''嘤嘤''的细碎声响。

``慢点吃，别抢。''卢本韦轻声哄着，伸手将一头挤到最外面的小猪崽轻轻拨回母猪怀里，指尖沾了一点饲料，递到小猪崽嘴边，小猪崽立刻伸出粉嫩的小舌头舔了舔，满足地哼了一声。

她的嘴角挂着温柔的笑意，发丝被晨风吹得有些凌乱，几缕碎发贴在脸颊上，衬得她的面容格外柔和。

奥斯卡靠在培育室的门框上，手里拿着一根香肠，却没有吃，只是静静地看着卢本韦的背影。晨光从她身后照过来，在她周身镀上一层淡淡的金边，空气中弥漫着母猪和猪崽的气息，混着序灵草的清香，竟然让他觉得格外安心。

``你看什么呢？''卢本韦忽然回头，正好对上奥斯卡的目光，脸颊瞬间红了。

``没\ldots\ldots 没看什么。''奥斯卡连忙移开视线，假装若无其事地咬了一口香肠，``我就是来看看母猪和小猪崽的，顺便\ldots\ldots 顺便看看你。''

最后一句话说得很轻，几乎被猪崽的哼唧声盖住了。

卢本韦的耳尖红得快要滴血，她连忙低下头，假装整理饲料碗，嘴里含糊地''嗯''了一声。

奥斯卡走到她身旁，蹲下来，和她一起看着那窝小猪崽。两人之间的距离很近，近到卢本韦能闻到他身上淡淡的香肠味和序力的清新气息。

``你看那头最小的，''奥斯卡伸手指了指角落里一头特别小的猪崽，它正努力地拱着母猪的肚子，却怎么也抢不过哥哥姐姐们，``它最弱，但是最努力。''

卢本韦看着那头小猪崽，忍不住笑了：``它像你。''

``像我？''奥斯卡一愣。

``嗯，''卢本韦歪了歪头，认真地说道，``看起来吊儿郎当的，其实比谁都努力，比谁都上心。''

奥斯卡的心猛地跳了一下，他转头看向卢本韦，发现她正认真地看着自己，那双清澈的眼眸里，倒映着晨光和他自己的影子。

``卢本韦，''奥斯卡的声音忽然变得认真起来，没有了平时的嬉皮笑脸，``决赛的时候，我把香肠塞给你的时候，你有没有觉得\ldots\ldots 我有点烦？''

卢本韦愣了一下，随即摇头，嘴角的笑意更深了：``不烦。''

``那以后，我还给你塞香肠，你也不许嫌烦。''奥斯卡的语气带着几分霸道，眼底却藏着温柔。

卢本韦低下头，看着怀里的小猪崽，轻轻''嗯''了一声，声音细若蚊蚋：``\ldots\ldots 好。''

角落里，那头最小的猪崽终于抢到了奶，满足地哼了一声，蜷缩在母猪怀里，沉沉睡去。

培育室外，阳光正好，万猪朝宗的灵草地上，序灵草在晨风中轻轻摇曳，灵气丝线如同金色的丝线，编织出一幅温暖而美好的画卷。

接下来的几日，众人在养猪场中安心休养，日子过得平静而充实。

每日清晨，众人会在万猪朝宗的灵草地上进行晨练，借助浓郁到液化的灵气加速序力恢复。母猪们在草地上悠然踱步，偶尔有几头调皮的母猪会凑到众人身边，用鼻子拱一拱他们的手心，讨要序灵草吃。牢大冰蓝色的序力天生带着一丝寒意，母猪们似乎格外喜欢他，每次他晨练时，总有三五头母猪围在他身边，蹭着他的裤腿不肯走，惹得众人哈哈大笑。

蔡须坤每次看到这一幕，都会忍不住弯起嘴角。她会悄悄走到牢大身边，假装也在晨练，实则借着母猪的掩护，偷偷看他几眼。有一次，一头母猪忽然拱了蔡须坤一下，将她推到了牢大身上，蔡须坤猝不及防，整个人撞进了牢大怀里，两人的脸瞬间都红了。牢大下意识地伸手扶住她的腰，触感柔软，他的手指微微一僵，随即若无其事地松开，低声说了句''小心''。蔡须坤''嗯''了一声，快步走开，耳尖红得通透，身后那两头''肇事''的母猪则悠闲地嚼着序灵草，仿佛什么都没发生过。

午后，众人会在饲料炼制坊帮忙，学习序力饲料的炼制方法。奥斯卡凭借着对食材的敏锐感知，很快掌握了饲料炼制的火候控制技巧，炼制出的饲料品质甚至超过了一些老师傅。卢本韦则负责饲料的分配和母猪的喂养，她对每一头母猪的习性都了如指掌，哪头母猪喜欢吃什么口味的饲料，哪头母猪最近心情不好需要多哄哄，她都记得清清楚楚。奥斯卡常常以''送饲料''为由，跑到卢本韦负责的片区，一边帮忙搬饲料，一边和她聊天，两人配合默契，效率比其他区域高出一倍。

傍晚时分，众人会聚在卢氏宗祠前的广场上，围坐在一起，品尝卢关特意准备的极品母猪肉宴。母猪肉经过特殊的序力烹饪，肉质鲜嫩多汁，入口即化，浓郁的序力在唇齿间绽放，每一口都如同在品尝天地灵气。卢关还会拿出珍藏多年的灵酒，与众人把酒言欢，席间笑声不断，气氛温馨而融洽。

酒过三巡，卢关会讲起他年轻时闯荡扽魂界的经历，讲他如何发现灵脉、如何培育出第一头变异母猪、如何一步步将养猪场发展成大陆第一养猪帝国。他的故事生动有趣，常常逗得众人捧腹大笑。

``你们别小看养猪，''卢关拍着桌子，满脸自豪，``这世上的道理，都藏在养猪里。母猪要养好，得懂它的脾气，给它最好的环境，让它吃得饱、睡得好，它自然会给你最好的回报。做人也是一样，心要正，路要宽，善待身边的人，日子自然会越过越好。''

众人纷纷点头，唐三若有所思地说道：``卢叔说得对，养猪和修炼，其实道理是相通的。''

酒宴散后，众人三三两两地散去。牢大会送蔡须坤回住处，两人走在灵石大道上，月光将两人的影子拉得很长。奥斯卡则会找借口和卢本韦一起去巡查母猪，两人并肩走在灵草地上，头顶是满天繁星，脚下是柔软的序灵草，空气中弥漫着母猪和灵草的气息，一切都刚刚好。

\begin{center}\rule{0.5\linewidth}{0.5pt}\end{center}

这座养猪帝国，不仅是一座避难所，更是一个温暖的港湾。在这里，没有邪扽师的威胁，没有生死的考验，只有母猪的哼唧声、灵草的清香、和少年少女心底悄悄萌芽的情愫。

而这份宁静与温暖，正是他们拼尽全力守护的意义所在。

养猪场的最中央，灵草地尽头，矗立着一座远比其余六座核心区域更加古老、更加庄严的建筑------\textbf{卢氏宗祠}。

与周围魂晶浇筑的新式厂房不同，宗祠通体由一种罕见的''太古玄石''砌成，石面呈深沉的墨黑色，却隐隐透出暗金色的光泽，仿佛石头内部封印着一轮永不熄灭的太阳。宗祠高达百丈，占地数十亩，却只有一进院落、一座正殿，没有多余的殿宇楼阁，简洁到了极致，却又庄严到了极致。

正殿的屋顶以一个巨大的\textbf{猪首图腾}为脊饰------那猪首双目圆睁，獠牙外露，却并不凶恶，反而透着一种浑厚而温厚的神性，仿佛天地间第一头母猪降世时的模样。猪首图腾通体由纯金铸造，历经万年风雨依旧熠熠生辉，每当阳光照在其上，便会折射出一道暗金色的光柱，直冲云霄，方圆百里都能看到。

宗祠大门高二十丈，宽十丈，门楣上刻着四个苍劲有力的古篆大字------\textbf{``万猪之源''}，字迹深邃入石三分，每一笔都蕴含着极其恐怖的序力。据传这四个字乃是猪神Opt亲手所书，万年过去，字迹中蕴含的序力依旧没有消散分毫，反而随着岁月沉淀愈发浑厚，任何心怀不轨之人靠近大门十丈之内，便会被字迹中蕴含的序力震得气血翻涌、寸步难行。

大门两侧，各矗立着一尊高达十丈的石雕母猪，石雕通体由整块太古玄石雕刻而成，母猪神态安详，双目微闭，腹部浑圆饱满，仿佛正在孕育天地间的万物生灵。石雕母猪的周身刻满了密密麻麻的序纹，那些序纹并非后世扽魂师所刻，而是万年之前便已存在的''天然序纹''------据传，这是猪神Opt在开辟养猪场时，以无上序力将天地法则直接刻入石头之中，形成了这两尊守护石雕。

推开沉重的玄石大门，正殿之内，一股浓郁到近乎液化的灵气扑面而来，那灵气与万猪朝宗的灵气截然不同------万猪朝宗的灵气温润而柔和，如同春风拂面；而宗祠内的灵气则厚重而庄严，如同站在万丈高山之巅，俯瞰天地万物，让人不由自主地心生敬畏。

正殿中央，矗立着一座\textbf{巨大无比的人形神像}------\textbf{猪神Opt}。

神像通体由一种从未见过的''序力神玉''雕琢而成，玉质温润如脂，通体呈现淡金色，在阳光下散发着柔和而神圣的光芒。神像高达\textbf{八十丈}，几乎与正殿的穹顶齐平，站在神像脚下仰望，只觉得自身渺小如蝼蚁，一股源自灵魂深处的敬畏感油然而生。

猪神Opt的神像是一位身披长袍的老者形象，面容慈祥而威严，长须垂胸，双目微闭，嘴角带着一抹淡然的微笑，那笑容中蕴含着无尽的包容与温厚，如同天地之父，孕育万物而不求回报。他的左手托着一头小猪崽，小猪崽通体金光，憨态可掬，与Opt慈祥的面容相映成趣；右手握着一柄短杖，短杖顶端雕刻着一枚精致的猪首，与屋顶的猪首图腾遥相呼应。

Opt身披的长袍上，刻满了密密麻麻的序纹------那些序纹并非装饰，而是真正的序数法则。从最基础的 \allowbreak{}\NovelInline{I01991}{\omega}\allowbreak{}、\allowbreak{}\NovelInline{I01992}{\varepsilon_0}\allowbreak{}，到极其复杂的 Veblen 函数、\allowbreak{}\NovelInline{I01993}{\psi}\allowbreak{} 函数，再到连弗兰德都未曾见过的 \allowbreak{}\NovelInline{I01994}{ω-Y(1,ω)\text{，}PTO(ZFC+I0)}\allowbreak{} 以上的高阶序数，仿佛猪神Opt将自己对序数的全部理解，都刻入了这身长袍之中。长袍的下摆自然垂落，衣褶之间隐约可见流动的金色光芒，如同一条条序力河流在石像体内奔涌不息。

神像的周身，环绕着九道金色的序力光环，每一道光环都缓缓旋转，散发着不同属性的序力波动------有的炽热如火，有的冰冷如霜，有的狂暴如雷，有的柔和如水，九道光环交织在一起，形成了一个完美的序力循环，源源不断地向整个宗祠、乃至整个养猪场输送着精纯至极的灵气。

神像基座高达十丈，由整块太古玄石雕成，基座正面刻着一行古老的铭文，字迹苍劲有力，正是猪神Opt的亲笔：

\begin{quote}
\textbf{``天地生万物，万物归于序。序之力，不在高低，而在循环。母猪孕万物，万物养母猪，此乃天道。''}
\end{quote}

这短短三十二个字，却蕴含着极其深奥的序数哲理。唐三第一次看到这段铭文时，紫极魔瞳不受控制地自动运转，他盯着那三十二个字看了整整一个时辰，脑海中轰然炸开------

``这不是普通的铭文\ldots\ldots 这是序数法则！''唐三震惊地说道，``猪神Opt在万年前就已经领悟了序数循环的真谛------序力不是单向的攀升，而是如同母猪孕育万物、万物反哺母猪的循环。这和我们之前学的序数理论完全不同，这是一种更高层次的理解！''

养猪场有一个延续了万年的传统------\textbf{每逢星期日，全场上万名员工都会齐聚宗祠，祭拜猪神Opt}。

这一传统源自卢氏第一祖。万年前，第一祖得猪神神念，开卢氏一脉，便定下规矩：每周日辰时，养猪场所有员工，无论身份高低、实力强弱，都必须放下手中一切事务，齐聚宗祠，向猪神Opt行祭拜之礼。这一规矩代代相传，万年不曾中断，即便是在养猪场最艰难的时期，也从未有人敢违抗。

每逢星期日清晨，天还未亮，养猪场便已忙碌起来。

负责清扫的员工会将宗祠内外的每一块石砖、每一根廊柱都擦拭得一尘不染；负责布置的员工会在宗祠前的广场上摆满供桌，供桌上摆放着精心准备的祭品------最新鲜的序灵草、最醇厚的灵酒、最精美的序力饲料，以及一头用整块灵玉雕刻的小猪崽。

卢关作为养猪场的主人，每次祭拜都会亲自换上祖传的玄色祭袍，祭袍上绣着金色的猪首图腾，庄重而肃穆。他会提前一个时辰来到宗祠，检查每一处细节，确保祭拜大典万无一失。

``祭拜猪神，容不得半点马虎。''卢关常常这样叮嘱众人，``Opt老祖宗保佑我们养猪场万年兴盛，我们理应以最虔诚的心意回报。''

\chapter{第三十三章：猪神传说}\label{ux7b2cux4e09ux5341ux4e09ux7ae0ux732aux795eux4f20ux8bf4}

辰时初刻，养猪场的钟声准时敲响，浑厚的钟声回荡在整个养猪场上空，方圆数十里都能听到。

钟声落下的瞬间，养猪场的所有工作全部停止。饲料炼制坊的炉火暂时封存，母猪培育殿的序力阵法暂时关闭，护卫营的巡逻暂时暂停，就连万猪朝宗的灵草地上，那些正在悠然踱步的变异母猪们仿佛也感受到了某种召唤，纷纷停下脚步，朝着宗祠的方向发出低沉的哼鸣，如同在回应远古的呼唤。

上万名员工从养猪场的四面八方汇聚而来，沿着宽阔的灵石大道，浩浩荡荡地走向宗祠。队伍的最前方，是卢关带领的卢氏核心族人，他们身着统一的玄色祭袍，步伐整齐划一，神情肃穆而庄重；队伍中间，是饲料炼制师、培育师、护卫队等各部员工，他们身着各自部门的制服，手持祭品，鱼贯而行；队伍最后方，是养猪场的杂役、仆从、学徒，他们虽然身份卑微，但脸上的虔诚与敬畏，却与前方的高层毫无二致。

当队伍抵达宗祠前的广场时，广场已经被布置得庄严肃穆。广场中央摆着一座巨大的香炉，香炉中燃烧着由九十九种灵草混合制成的''敬神香''，烟气袅袅升腾，散发着沁人心脾的清香，那香气中蕴含着微弱的序力，吸入之后，浑身经脉都会微微发热，精神为之一振。

卢关站在香炉前方，面向宗祠正殿，双手抱拳，深深鞠躬。上万名员工紧随其后，齐刷刷地弯腰鞠躬，动作整齐划一，没有丝毫杂音，只有衣袍摩擦的细微声响。

三鞠躬毕，卢关朗声诵读祭文，祭文的内容万年未变，正是猪神Opt留下的那段铭文：

\begin{quote}
``天地生万物，万物归于序。序之力，不在高低，而在循环。母猪孕万物，万物养母猪，此乃天道。''
\end{quote}

上万名员工齐声跟诵，上千道声音汇聚在一起，回荡在宗祠上空，回荡在万猪朝宗的灵草地上，回荡在整个养猪场的每一个角落。那声音浑厚而庄严，如同远古的钟鸣，震得空气都在微微颤动。

跟诵完毕，便是最神圣的环节------\textbf{猪神赐福}。

卢关独自走进正殿，来到八十丈高的猪神Opt巨像脚下，双手按在基座的序纹之上，催动自身 \allowbreak{}\NovelInline{I01995}{Y(1,2,4,8,16)}\allowbreak{} 的序力。刹那间，神像周身九道金色光环猛然亮起，一道金色的光柱从神像掌心射出，穿过正殿的穹顶，直冲云霄，在养猪场上空化作一朵巨大的金色莲花，莲花缓缓旋转，洒下无数金色的光点。

光点如同细雨般飘落，落在每一个员工的身上。那些光点蕴含着猪神Opt残留的神力，能够温养经脉、修复暗伤、提升序力，虽然效果微弱，但日积月累，养猪场员工的整体实力远超同阶扽魂师，这便是万年祭拜的真正意义所在。

光点落在万猪朝宗的母猪们身上，母猪们纷纷发出满足的哼鸣，周身灵气愈发浓郁，有几头天赋异禀的母猪甚至在光点的滋养下当场突破，周身绽放出耀眼的灵光，引来众人一阵欢呼。

赐福结束后，卢关从正殿走出，面向众人，朗声说道：

\begin{quote}
``猪神庇佑，我养猪场万年兴盛。愿诸君勤勉修炼，善待母猪，坚守正道，不负Opt老祖宗之期望！''
\end{quote}

\begin{quote}
``谨遵家主之命！''上万名员工齐声应道，声震云霄。
\end{quote}

祭拜大典至此结束，众人各自散去，回到岗位上继续工作。但每个人的脸上都带着满足与安宁的笑容，仿佛刚刚经历了一场灵魂的洗礼。

宗祠两侧的墙壁上，刻满了壁画，记录着猪神Opt的传奇一生。

壁画以极其精湛的技艺雕刻而成，每一幅画面都栩栩如生，仿佛将万年前的历史直接定格在了石壁之上。

第一幅壁画描绘的是Opt以无上神力，在天地之间开辟出一片广袤的灵地，将自身的序力注入大地之中，形成了第一条灵脉。灵脉之上，Opt以自身精血培育出了第一批变异母猪，这些母猪继承了Opt的神性，能够吸收天地灵气并转化为纯净的序力，反哺大地。

第二幅壁画中，Opt站在灵脉之上，左手托着一头小猪崽，右手的短杖点向大地，金色的序力从杖尖涌出，化作无数金色的丝线，渗入大地深处，形成了养猪场的雏形------万年前，Opt亲手种下的第一株序灵草，亲手培育的第一头变异母猪，亲手布下的第一道序力结界。

第三幅壁画描绘的是Opt在养猪场中修炼万年，最终突破序力的极限，成就''猪神''之位的场景。壁画中，Opt的身躯化作一道金色的光柱，直冲云霄，周身环绕着无数道序力光环，每一道光环都代表着一个序数层级的突破。最终，Opt的身影消失在光柱之中，只留下一座巨大的神像，和一段永恒的传说。

第四幅壁画描绘的是卢氏先祖------一个普通的牧猪人------在万猪朝宗的灵草地上，意外获得了猪神Opt残留的一缕神念。神念化为一道金色的光芒，融入牧猪人的体内，从此，卢氏族人便拥有了与母猪心灵相通的天赋，能够借助母猪的灵气修炼序力，代代相传，绵延至今。

众人抵达养猪场的第一个星期日，便赶上了祭拜大典。

卢关特意嘱咐卢本韦，带着果糕小队和蔡兴坤一同参加。``你们是养猪场的贵客，也是本韦的同伴，理应一同祭拜猪神，接受赐福。''

清晨，众人换上卢关准备的玄色祭袍，跟随卢氏族人一同走向宗祠。

蔡须坤走在队伍中，望着前方那座越来越近的黑色建筑，心中涌起一股莫名的亲切感。她是魂兽化形之人，对天地间的神性有着天然的感知，她能感受到宗祠中蕴含的那股古老而温柔的力量，那力量与星斗森林深处的某种气息有着微妙的相似------都是属于天地万物最原始、最纯粹的力量。

当众人走进正殿，抬头仰望那座八十丈高的猪神Opt巨像时，所有人都被震撼得说不出话来。

``好大\ldots\ldots{}''卢本韦喃喃自语，虽然她从小在养猪场长大，每次祭拜都会被神像的宏伟所震撼，但今天，有同伴在身边一同仰望，那份震撼似乎又多了几分。

奥斯卡站在卢本韦身旁，目光却没有落在神像上，而是偷偷看着卢本韦仰望神像时的侧脸，心底悄悄泛起一丝温柔。

牢大站在蔡须坤身旁，两人之间的距离比平时近了一些。蔡须坤仰望着神像，眼眶微微泛红------她是魂兽，她比任何人都能感受到猪神Opt神像中蕴含的那股温柔而强大的力量。

``原来\ldots\ldots 魂兽也可以成神。''蔡须坤喃喃自语，``猪神Opt虽然不是魂兽，但它以母猪为伴，以孕育万物为道，它理解魂兽，也尊重魂兽\ldots\ldots 那我，是不是也有可能，找到属于自己的路？''

牢大侧头看了她一眼，没有说话，只是默默地往她身边靠了靠，用肩膀轻轻碰了碰她的肩膀。

蔡须坤微微一怔，转头看向牢大，正好对上他温和的目光，心底涌起一股暖流，嘴角不自觉地弯了起来。

当金色的光点从穹顶洒落，落在每个人身上时，蔡须坤闭上眼睛，感受着那股温暖的力量。她仿佛听到了一声遥远的叹息，那叹息来自万年之前，来自猪神Opt的最后一缕神念------

\begin{quote}
``万物皆有序，万物皆可成神。''
\end{quote}

蔡须坤的眼角滑落一滴泪水，那不是悲伤的泪，而是释然的泪。她知道，自己的路，就在前方。

祭拜结束后，众人并未立刻散去。

卢关带着卢本韦、果糕小队和蔡兴坤，绕到正殿后方一处僻静的偏殿。偏殿内没有神像，只有一张古朴的石桌，石桌上摆着一壶刚沏好的灵茶，茶香清冽，与宗祠中那股厚重的灵气交织在一起，令人心神安宁。

卢关亲自为众人斟茶，随后在石桌旁坐下，神色比方才在祭台上更加肃穆。

``今日带你们来祭拜，不只是为了让你们感受猪神赐福。''卢关缓缓说道，``而是有一件事，我想让你们知道。''

众人纷纷看向他。

卢关的目光落在卢本韦身上，又扫过果糕小队的每一个人，最后停留在蔡须坤身上，语气低沉而郑重：

``你们方才在正殿中，看到了Opt老祖宗神像长袍上的序纹，也看到了基座上的铭文。但那些，只是猪神留给世人的表象。''

他顿了顿，端起茶杯，轻轻抿了一口，才继续说道：

``猪神Opt之所以能够成神，并非因为他修炼到了某个递归序数的极限，而是因为他最终领悟了宇宙法则。他的修为，早已不能用递归序数来表示。''

这句话如同一颗石子投入平静的湖面，在众人心中激起层层波澜。

唐三第一个反应过来，紫极魔瞳微微闪烁，沉声问道：``不能用递归序数表示？难道\ldots\ldots 猪神Opt的境界，已经超越了 PTO(ZFC+I0)？''

卢关点了点头，目光中带着深深的敬畏：

``不仅超越了。在猪神Opt尚未成神之前，他的修为就已经达到了 PTO(ZFC+I0) 的境界。那是当今世间任何修炼方法------无论是魂导科技、传统冥想，还是各大势力的秘传功法------都无法企及的巅峰。''

此言一出，偏殿内一片寂静。

卢本韦虽然从小听父亲提起过猪神Opt的传说，但如此直白地听到 PTO(ZFC+I0) 被定义为``未成神时的境界''，依旧让她心头剧震。

奥斯卡忍不住问道：``那\ldots\ldots 成神之后的猪神，究竟是什么境界？''

卢关摇了摇头，目光望向正殿的方向，仿佛能透过墙壁看到那座八十丈高的神像：

``没有人知道。因为猪神Opt成神之后，便不再属于任何一个可以被定义、被计算的序数体系。他领悟了宇宙法则的终极奥义------循环与孕育。他不再是`更高'，而是`圆满'。''

他看向蔡须坤，语气忽然变得温和：

``这也是为什么，魂兽也可以成神。猪神Opt告诉我们，成神的关键不在于血脉，不在于种族，而在于是否领悟了`循环'的真谛。母猪孕育万物，万物反哺母猪，这便是天道。任何生灵，只要顺应这份天道，便有机会触摸到神的门槛。''

蔡须坤浑身一震，眼眶再次泛红。她终于明白，方才在赐福中感受到的那股温暖力量，为何会让她如此动容------因为猪神Opt的道，从来不是高高在上的，而是包容万物的。

卢关最后看向卢本韦，目光中带着期许：

``本韦，你是卢氏后人，也是养猪场的未来。今日我将这些话告诉你，是希望你能明白，卢氏一脉传承万年的，不仅仅是养猪的技术和序力修炼的方法，更是猪神Opt留下的这份`循环之道'。善待母猪，敬畏天地，坚守正道------这不仅是家训，更是通往更高境界的唯一路径。''

卢本韦郑重地点了点头：``父亲，我明白了。''

偏殿内再次陷入沉默，但每个人的心中都已掀起了惊涛骇浪。

猪神Opt，一个在万年前便已站在 PTO(ZFC+I0) 巅峰的存在，最终却选择了一条与所有修炼者截然不同的道路------不是继续向上攀登递归的阶梯，而是回归循环，回归孕育，回归万物本源。

这，便是万猪之源的真正含义。

宗祠正殿的后面，有一道隐蔽的暗门，暗门之后是一条向下延伸的石阶，通往宗祠地下的密室。

密室不大，只有方圆三丈，四壁由太古玄石砌成，密室的正中央，悬浮着一颗拳头大小的金色光球，光球缓缓旋转，散发着温暖而柔和的光芒，整个密室的灵气都源自这颗光球。

这便是猪神Opt留下的最后一缕神念------``猪神令''。

卢关曾告诉卢本韦，猪神令是卢氏家族的至宝，每一代卢氏厂长在接任之时，都要进入密室，接受猪神令的传承。传承之后，厂长的序力会与养猪场的灵脉彻底融合，能够调动整个养猪场的灵气为己用，实力暴涨数倍。但猪神令每年只能传承一次，且只有卢氏血脉之人才能承受传承的力量，外人若强行接触，便会被猪神令蕴含的恐怖序力瞬间碾碎。

密室墙壁上，还刻着一幅特殊的序纹图案------那是一头母猪的轮廓，母猪体内刻满了密密麻麻的序数公式，从最基础的 omega 到极其复杂的 ω-Y(1,ω),乃至PTO(ZFC+I0)，仿佛猪神Opt将自己对序数的全部理解，都刻入了这头母猪的体内。

卢本韦在参观密室时，盯着那幅序纹图案看了许久，忽然浑身一震：

``我明白了\ldots\ldots 猪神Opt在万年前就已经领悟了序数的终极奥义------序数的力量不在于层级的高低，而在于循环与孕育。母猪孕育万物，万物反哺母猪，这就是序数的终极法则。我们之前学的Veblen函数、psi 函数，都是在追求'更高'，但猪神Opt告诉我们，`更高'的尽头，是'回归'。''

这番话让所有人都陷入了沉思。或许，序数之路的终点，并不是一个无穷大的数字，而是一个完美的循环------如同养猪场的灵气循环，如同母猪孕育万物的天道轮回。

\chapter{第三十四章：进阶的BOCF}\label{ux7b2cux4e09ux5341ux56dbux7ae0ux8fdbux9636ux7684bocf}

在养猪场安顿好之后，众人惊叹于这养猪场规模真的极其宏大，一眼望不到边的大小，卢本韦说：你们要知道，我的养猪场可不是普通养猪场，这养猪场多达几十万头母猪，母猪都是有灵气的，卢关在养猪场，可以短暂发挥出Y(1,3)的巨大力量。

要说那蔡须坤和蔡兴坤，这两人同人喜欢上了牢大，但是她们也是密切的战友。所以，就算两人略有抵触，但是一切以队伍整体为重。

卢关说：``现在有必要和大家讲一讲更高层的BOCF。如果我猜得没错，LVO之前的那些结构，你们已经掌握了不少。''

牢大点头：``经过果糕八怪的切磋和努力，我已经基本掌握了BHO以内的序数。唐三他们，也摸到了LVO以内的路。''

卢关没有急着往后讲，而是在众人面前写下了这一阶段的约定。

\NovelDisplay{M00267R01}{%
\Omega=\omega_1^{\mathrm{CK}},\qquad\psi=\psi_\Omega%
}

``这里的 \allowbreak{}\NovelInline{I01996}{\Omega}\allowbreak{} 是Church--Kleene序数，本身可数，却已经超出了递归序数的范围。后面的 \allowbreak{}\NovelInline{I01997}{\Omega_n}\allowbreak{} 也按这套可数的非递归层级理解。我们借助这些更高的层次，构造 \allowbreak{}\NovelInline{I01998}{\Omega}\allowbreak{} 以下的递归序数。''

他又补上了三个起点。

\NovelDisplay{M00268R01}{%
\psi(0)=1%
}

\NovelDisplay{M00269R01}{%
\psi(\Omega)=\varepsilon_0%
}

\NovelDisplay{M00270R01}{%
\psi_{\Omega_2}(0)=\Omega%
}

``BOCF里，零参数不是自动跳到下一个 \allowbreak{}\NovelInline{I01999}{\varepsilon}\allowbreak{} 数。先把这几行记牢。''

卢关转向牢大：``你说掌握了BHO，那么它的表达式与基本列，也要认得出来。''

\NovelDisplay{M00271R01}{%
\mathrm{BHO}=\psi(\Omega_2)%
}

\NovelDisplay{M00272R01}{%
\psi_{\Omega_2}(\Omega_2)=\varepsilon_{\Omega+1}%
}

\NovelDisplay{M00273R01}{%
\psi\bigl(\varepsilon_{\Omega+1}\bigr)=\psi(\Omega_2)%
}

``最后一行是这个关口的另一种表达。实际推演时，我们用标准形式 \allowbreak{}\NovelInline{I02000}{\psi(\Omega_2)}\allowbreak{}。以下方括号表示基本列取项，所有下标从零开始。''

\NovelDisplay{M00274R01}{%
\bigl(\psi(\Omega_2)\bigr)[0]=1%
}

\NovelDisplay{M00275R01}{%
\bigl(\psi(\Omega_2)\bigr)[1]=\psi(\Omega)%
}

\NovelDisplay{M00276R01}{%
\bigl(\psi(\Omega_2)\bigr)[2]=\psi(\Omega^2)%
}

\NovelDisplay{M00277R01}{%
\bigl(\psi(\Omega_2)\bigr)[3]=\psi(\Omega^\Omega)%
}

``现在展开 \allowbreak{}\NovelInline{I02001}{f_{\psi(\Omega_2)}(3)}\allowbreak{}。对极限下标，采用 \allowbreak{}\NovelInline{I02002}{f_\lambda(3)=f_{\lambda[3]}(3)}\allowbreak{}。同一个输入三，每一步都取当前下标基本列中下标为三的项。''

牢大接过笔，先写出了第一步。

\NovelDisplay{M00278R01}{%
f_{\psi(\Omega_2)}(3)=f_{\psi(\Omega^{\Omega})}(3)%
}

唐三跟着往下算。

\NovelDisplay{M00279R01}{%
=f_{\psi(\Omega^{\psi(\Omega^{\psi(\Omega^{\omega})})})}(3)%
}

\NovelDisplay{M00280R01}{%
=f_{\psi(\Omega^{\psi(\Omega^{\psi(\Omega^{4})})})}(3)%
}

小舞看着越来越深的括号，忍不住把纸转向自己。她先沿着最里面的指数检查了一遍，才接着写下去。

\NovelDisplay{M00281R01}{%
=f_{\psi(\Omega^{\psi(\Omega^{\psi(\Omega^{3}\times\psi(\Omega^{3}\times\psi(\Omega^{3}\times\omega)))})})}(3)%
}

\NovelDisplay{M00282R01}{%
=f_{\psi(\Omega^{\psi(\Omega^{\psi(\Omega^{3}\times\psi(\Omega^{3}\times\psi(\Omega^{3}\times3)))})})}(3)%
}

\NovelDisplay{M00283R01}{%
=f_{\psi(\Omega^{\psi(\Omega^{\psi(\Omega^{3}\times\psi(\Omega^{3}\times\psi(\Omega^{3}\times2+\Omega^{2}\times\psi(\Omega^{3}\times2+\Omega^{2}\times\psi(\Omega^{3}\times2+\Omega^{2}\times\omega)))))})})}(3)%
}

``这里是函数值相等，因为我们用了极限下标的定义。''卢关说，``里面的序数下标却在下降。别把两件事混在一起，更不能把取出的基本列项直接说成原序数。''

蔡兴坤看得有些发怔：``这样看来，记住一个关口的名字，还远远不够。''

``你得知道它怎样一步步退下来。''卢关指了指最外层的 \allowbreak{}\NovelInline{I02003}{\psi}\allowbreak{}，``往上修炼与往下展开，正好要互相印证。''

他又写下两个更高的关口。

\NovelDisplay{M00284R01}{%
\psi_{\Omega_2}(\Omega_2\times2)=\varepsilon_{\Omega+2}%
}

\NovelDisplay{M00285R01}{%
\psi_{\Omega_2}(\Omega_2^2)=\zeta_0\text{ aft }\Omega%
}

``它们也提示我们，可以继续研究下面两种更高的标准参数。不过，内层函数的取值，和最外层的折叠结果要分别核对，不能任意删掉嵌套层。''

\NovelDisplay{M00286R01}{%
\psi(\Omega_2\times2)%
}

\NovelDisplay{M00287R01}{%
\psi(\Omega_2^2)%
}

唐三问：``这是不是说明，参数每增加一，就进入下一个 \allowbreak{}\NovelInline{I02004}{\varepsilon}\allowbreak{} 数？''

``先看清楚这里增加的是什么。''卢关在 \allowbreak{}\NovelInline{I02005}{\Omega_2\times2}\allowbreak{} 下方画了一道线，``从 \allowbreak{}\NovelInline{I02006}{\Omega_2}\allowbreak{} 到 \allowbreak{}\NovelInline{I02007}{\Omega_2\times2}\allowbreak{}，增加的是一整个 \allowbreak{}\NovelInline{I02008}{\Omega_2}\allowbreak{}。BOCF里，参数加一的行为不同。''

\NovelDisplay{M00288R01}{%
\psi(\Omega+1)=\varepsilon_0\times\omega%
}

\NovelDisplay{M00289R01}{%
\psi(\Omega\times2)=\varepsilon_1%
}

\NovelDisplay{M00290R01}{%
\psi_{\Omega_2}(\Omega_2+1)=\varepsilon_{\Omega+1}\times\omega%
}

``下一项 \allowbreak{}\NovelInline{I02009}{\varepsilon}\allowbreak{} 数可以由指数塔的极限得到，但要写清底数和极限；不能把塔里一长串省略号当成任意参数加一的规则。''

小舞问：``那么，\allowbreak{}\NovelInline{I02010}{\Omega}\allowbreak{} 之上还有更高的非递归层级吗？''

``有。\allowbreak{}\NovelInline{I02011}{\Omega_2}\allowbreak{} 就是接下来要使用的一层。''

陈星这时走了过来：``在继续之前，我先给你们解释一下刚才那个 \allowbreak{}\NovelInline{I02012}{\text{aft}}\allowbreak{}。''

他写下一个众人已经熟悉的式子。

\NovelDisplay{M00291R01}{%
\zeta_0=\varepsilon_{\zeta_0}%
}

``若说`在 \allowbreak{}\NovelInline{I02013}{\zeta_0}\allowbreak{} 之后的第一个 \allowbreak{}\NovelInline{I02014}{\varepsilon}\allowbreak{} 数'，我们在这段讨论中简写为：''

\NovelDisplay{M00292R01}{%
\varepsilon_0\text{ aft }\zeta_0=\varepsilon_{\zeta_0+1}%
}

``第二个，则是：''

\NovelDisplay{M00293R01}{%
2\text{ th }\bigl(\varepsilon_0\text{ aft }\zeta_0\bigr)=\varepsilon_{\zeta_0+2}%
}

卢本韦想了想：``如果说第二个 \allowbreak{}\NovelInline{I02015}{\zeta}\allowbreak{} 数，就是在 \allowbreak{}\NovelInline{I02016}{\zeta_0}\allowbreak{} 之后再找一个同样性质的序数？''

\NovelDisplay{M00294R01}{%
2\text{ th }\zeta_0=\zeta_1=\zeta_0\text{ aft }\zeta_0%
}

卢关说：``是的，女儿。这个说法能帮助你辨认关口。不过，必须先讲清楚在枚举什么，不能把它当成对所有序数都适用的平移公式。''

陈星接着问：``那我们刚才讨论的，LVO之后的下一个 \allowbreak{}\NovelInline{I02017}{\Gamma}\allowbreak{} 型关口呢？''

蔡兴坤答道：

\NovelDisplay{M00295R01}{%
\Gamma_0\text{ aft }\mathrm{LVO}=\psi(\Omega^{\Omega^\Omega}+\Omega^\Omega)%
}

``这里的 \allowbreak{}\NovelInline{I02018}{\text{aft}}\allowbreak{} 仍是针对指定关口的说法。''陈星说，``有了它，你们能先理解次序；真正展开，还是要回到BOCF的参数和基本列。''

牢大回过头，重新看了看最初的零参数：``也就是说，\allowbreak{}\NovelInline{I02019}{\psi_{\Omega_2}(0)}\allowbreak{} 是 \allowbreak{}\NovelInline{I02020}{\Omega}\allowbreak{}；要到 \allowbreak{}\NovelInline{I02021}{\varepsilon_{\Omega+1}}\allowbreak{}，参数得走到 \allowbreak{}\NovelInline{I02022}{\Omega_2}\allowbreak{}。''

``对。你也可以把它与最外层作一个对照。''

\NovelDisplay{M00296R01}{%
\psi(0)=1,\qquad\psi_{\Omega_2}(0)=\Omega%
}

\NovelDisplay{M00297R01}{%
\psi(\Omega)=\varepsilon_0,\qquad\psi_{\Omega_2}(\Omega_2)=\varepsilon_{\Omega+1}%
}

``但这不等于对任何公式，都只要把某个 \allowbreak{}\NovelInline{I02023}{\Omega}\allowbreak{} 换成 \allowbreak{}\NovelInline{I02024}{\Omega_2}\allowbreak{} 就行。每一层都有自己的作用范围。''

\NovelDisplay{M00298R01}{%
\Omega\leq\psi_{\Omega_2}(\alpha)<\Omega_2%
}

``左端包含等号。零参数恰好取到下界，不能漏掉。''

唐三问：``如果里面遇到的是 \allowbreak{}\NovelInline{I02025}{\Omega}\allowbreak{}，和遇到 \allowbreak{}\NovelInline{I02026}{\Omega_2}\allowbreak{}，该怎样区别？''

卢关在纸上写出两道已经化为BOCF标准形的题。

\NovelDisplay{M00299R01}{%
\psi(\Omega_2\times\Omega)%
}

\NovelDisplay{M00300R01}{%
\psi(\Omega_2^{2})%
}

``这就是你们理解进阶BOCF的关键。第一道在末尾遇到 \allowbreak{}\NovelInline{I02027}{\Omega}\allowbreak{} 层的索引，回填时用最外层的 \allowbreak{}\NovelInline{I02028}{\psi}\allowbreak{}；第二道遇到 \allowbreak{}\NovelInline{I02029}{\Omega_2}\allowbreak{} 层的索引，则要用 \allowbreak{}\NovelInline{I02030}{\psi_{\Omega_2}}\allowbreak{}。''

第一道的前三个非零下标项是：

\NovelDisplay{M00301R01}{%
\bigl(\psi(\Omega_2\times\Omega)\bigr)[1]=\psi(\Omega_2\times\omega)%
}

\NovelDisplay{M00302R01}{%
\bigl(\psi(\Omega_2\times\Omega)\bigr)[2]=\psi(\Omega_2\times\psi(\Omega_2\times\omega))%
}

\NovelDisplay{M00303R01}{%
\bigl(\psi(\Omega_2\times\Omega)\bigr)[3]=\psi(\Omega_2\times\psi(\Omega_2\times\psi(\Omega_2\times\omega)))%
}

``前一项的整体结果，被填回后面的 \allowbreak{}\NovelInline{I02031}{\Omega_2}\allowbreak{} 系数右侧。''牢大说，``所以重复出现的是最外面的 \allowbreak{}\NovelInline{I02032}{\psi}\allowbreak{}。''

第二道则是：

\NovelDisplay{M00304R01}{%
\bigl(\psi(\Omega_2^{2})\bigr)[1]=\psi(\Omega_2\times\Omega)%
}

\NovelDisplay{M00305R01}{%
\bigl(\psi(\Omega_2^{2})\bigr)[2]=\psi(\Omega_2\times\psi_{\Omega_2}(\Omega_2\times\Omega))%
}

\NovelDisplay{M00306R01}{%
\bigl(\psi(\Omega_2^{2})\bigr)[3]=\psi(\Omega_2\times\psi_{\Omega_2}(\Omega_2\times\psi_{\Omega_2}(\Omega_2\times\Omega)))%
}

奥斯卡皱着眉：``看起来都在复制，可这里多出来的，是 \allowbreak{}\NovelInline{I02033}{\psi_{\Omega_2}}\allowbreak{}。''

``你终于看见区别了。''陈星说，``前一道用 \allowbreak{}\NovelInline{I02034}{\Omega}\allowbreak{} 以下的结果回填；后一道则先生成 \allowbreak{}\NovelInline{I02035}{\Omega}\allowbreak{} 到 \allowbreak{}\NovelInline{I02036}{\Omega_2}\allowbreak{} 这一层的项。只记住`复制'，却不记住在哪一层复制，就会把两道题混成一道。''

蔡须坤又把几个具体值写在旁边，让自己的理解有个落脚点。

\NovelDisplay{M00307R01}{%
\psi_{\Omega_2}(\Omega_2\times\Omega)=\varepsilon_{\Omega\times2}%
}

\NovelDisplay{M00308R01}{%
\psi_{\Omega_2}\bigl(\Omega_2\times\psi_{\Omega_2}(\Omega_2)\bigr)=\varepsilon_{\varepsilon_{\Omega+1}}%
}

``这些都是具体结构的对照。''陈星提醒，``不能从中推出任意一个含 \allowbreak{}\NovelInline{I02037}{\Omega}\allowbreak{} 的函数都能整体平移。要先保证参数与折叠层都符合规则。''

唐三问：``\allowbreak{}\NovelInline{I02038}{\Omega_2}\allowbreak{} 的层次，我还是有些模糊。''

``把它看成比 \allowbreak{}\NovelInline{I02039}{\Omega}\allowbreak{} 更高的容许层级。''陈星说，``它不是 \allowbreak{}\NovelInline{I02040}{\Omega}\allowbreak{} 后面紧挨着的那个序数，也不是说 \allowbreak{}\NovelInline{I02041}{\Omega}\allowbreak{} 之后只有一个非递归序数。我们在较高层借助新的折叠，处理前一层没有提供的结构。''

牢大点头：``像研究 \allowbreak{}\NovelInline{I02042}{\varepsilon}\allowbreak{} 数时，也得区分`下一个序数'与`下一个满足同样闭包性质的序数'。''

``这个区别记住就好。现在再看，折叠函数嵌套以后，有时能化成更短的表达式。''

\NovelDisplay{M00309R01}{%
\psi\bigl(\psi_{\Omega_2}(0)\bigr)=\psi(\Omega)%
}

\NovelDisplay{M00310R01}{%
\psi_{\Omega_2}(\Omega)=\Omega^2%
}

\NovelDisplay{M00311R01}{%
\psi\bigl(\psi_{\Omega_2}(\Omega)\bigr)=\psi(\Omega^2)%
}

\NovelDisplay{M00312R01}{%
\psi\bigl(\psi_{\Omega_2}(\Omega_2)\bigr)=\psi(\Omega_2)%
}

小舞把最后两行比了又比：``能省去一层，必须看这一层的参数。不能见到两层 \allowbreak{}\NovelInline{I02043}{\psi}\allowbreak{}，就一律删掉里面那层。''

``正是如此。''陈星说，``也不能看见某一对式子相等，就认为往括号里任意加一个数以后，它们还遵守同一条化简。先检查整个表达式的标准形，再做下一步展开。''

他写出了BHO之上一个更高关口的后继变化。

\NovelDisplay{M00313R01}{%
\psi(\Omega_2^2+1)=\psi(\Omega_2^2)\times\omega%
}

\NovelDisplay{M00314R01}{%
\psi(\Omega_2^2+\Omega)=\varepsilon_{\psi(\Omega_2^2)+1}%
}

``加一，和加 \allowbreak{}\NovelInline{I02044}{\Omega}\allowbreak{}，依然不同。要进入下一个 \allowbreak{}\NovelInline{I02045}{\varepsilon}\allowbreak{} 数，不能省去中间那层极限。''

小舞接着问：``再往后，还是可以研究不动点？''

\NovelDisplay{M00315R01}{%
\psi(\Omega_2^2+\Omega^2)%
}

\NovelDisplay{M00316R01}{%
\alpha\longmapsto\psi(\Omega_2^2+\Omega\times\alpha)%
}

``对。它对应从合适起点迭代这条映射得到的不动点。''陈星说，``注意参数中的 \allowbreak{}\NovelInline{I02046}{\Omega}\allowbreak{} 乘数。到了更大的固定前缀之后，这条局部规律还可以继续发挥作用。''

卢关接过话：``比如下面这个节点，就是后面研究两倍结构的基础。''

\NovelDisplay{M00317R01}{%
\psi(\Omega_2^{2}+\psi_{\Omega_2}(\Omega_2))%
}

``你们要考虑的两倍结构，则是：''

\NovelDisplay{M00318R01}{%
\psi(\Omega_2^{2}\times2)%
}

唐三刚想把末尾的 \allowbreak{}\NovelInline{I02047}{\Omega_2}\allowbreak{} 换成一个最外层 \allowbreak{}\NovelInline{I02048}{\psi}\allowbreak{}，陈星已经按住了他的笔。

``末尾这一层需要的是 \allowbreak{}\NovelInline{I02049}{\psi_{\Omega_2}}\allowbreak{}。你若换成 \allowbreak{}\NovelInline{I02050}{\psi}\allowbreak{}，就把回填范围压到另一层去了。顺着基本列看，不会弄混。''

牢大写道：

\NovelDisplay{M00319R01}{%
\bigl(\psi(\Omega_2^{2}\times2)\bigr)[1]=\psi(\Omega_2^{2}+\Omega_2\times\Omega)%
}

\NovelDisplay{M00320R01}{%
\bigl(\psi(\Omega_2^{2}\times2)\bigr)[2]=\psi(\Omega_2^{2}+\Omega_2\times\psi_{\Omega_2}(\Omega_2^{2}+\Omega_2\times\Omega))%
}

\NovelDisplay{M00321R01}{%
\bigl(\psi(\Omega_2^{2}\times2)\bigr)[3]=\psi(\Omega_2^{2}+\Omega_2\times\psi_{\Omega_2}(\Omega_2^{2}+\Omega_2\times\psi_{\Omega_2}(\Omega_2^{2}+\Omega_2\times\Omega)))%
}

``开头固定的那一份不动，后面按规则回填。''

陈星点头：``这次没有把原序数本身放到右边来冒充基本列项。每一个取出的项，都必须比原序数小。''

蔡兴坤又指着一个中间节点问：``那这一道，也是在固定前缀后面做迭代？''

\NovelDisplay{M00322R01}{%
\psi(\Omega_2^{2}+\psi_{\Omega_2}(\Omega_2\times\Omega))%
}

\NovelDisplay{M00323R01}{%
\alpha\longmapsto\psi\bigl(\Omega_2^2+\psi_{\Omega_2}(\Omega_2\times\alpha)\bigr)%
}

``对，要保留里面的 \allowbreak{}\NovelInline{I02051}{\Omega_2}\allowbreak{} 乘数，并从这道题相应的起点开始。''陈星说，``换了前缀，回填时携带的结构也要一起带上，不能只靠记忆里一个`不动点'的名称。''

他又问：``如果再把乘法的层次提高一层呢？''

\NovelDisplay{M00324R01}{%
\psi(\Omega_2^{3})%
}

小舞这次没有直接在指数上堆复制。她按乘法的末项写出了三项。

\NovelDisplay{M00325R01}{%
\bigl(\psi(\Omega_2^{3})\bigr)[1]=\psi(\Omega_2^{2}\times\Omega)%
}

\NovelDisplay{M00326R01}{%
\bigl(\psi(\Omega_2^{3})\bigr)[2]=\psi(\Omega_2^{2}\times\psi_{\Omega_2}(\Omega_2^{2}\times\Omega))%
}

\NovelDisplay{M00327R01}{%
\bigl(\psi(\Omega_2^{3})\bigr)[3]=\psi(\Omega_2^{2}\times\psi_{\Omega_2}(\Omega_2^{2}\times\psi_{\Omega_2}(\Omega_2^{2}\times\Omega)))%
}

``保持的是 \allowbreak{}\NovelInline{I02052}{\Omega_2^2}\allowbreak{} 这一份，再让它右边的因子按相应层级增长。''

``很好。乘法的共尾展开，不能随手换成乘方。''

现在，就可以继续引入更高的自然数下标了。

\NovelDisplay{M00328R01}{%
\psi_{\Omega_3}(0)=\Omega_2%
}

\NovelDisplay{M00329R01}{%
\psi_{\Omega_3}(\Omega_3)=\varepsilon_{\Omega_2+1}%
}

\NovelDisplay{M00330R01}{%
\psi\bigl(\varepsilon_{\Omega_2+1}\bigr)=\psi(\Omega_3)%
}

``\allowbreak{}\NovelInline{I02053}{\Omega_3}\allowbreak{} 是接下来的第三层。你们已经知道怎样区分零参数与较大参数，就可以把这种分层思路继续往上用。''

含有自然数下标的各层，取上确界得到 \allowbreak{}\NovelInline{I02054}{\Omega_\omega}\allowbreak{}。这里的 \allowbreak{}\NovelInline{I02055}{\Omega_\omega}\allowbreak{} 本身仍然可数。

\NovelDisplay{M00331R01}{%
\Omega_\omega=\sup_{1\leq n<\omega}\Omega_n%
}

\NovelDisplay{M00332R01}{%
\mathrm{BO}=\psi(\Omega_\omega)=\sup_{1\leq n<\omega}\psi(\Omega_n^2)%
}

这个巨大的关口，叫作Buchholz Ordinal，简称BO。

\begin{quote}
万年以前，Buchholz创立BOCF修炼方法，让魂师有了巨大突破，人类第一次有了和魂兽匹敌的资本。他创立的Buchholz宗至今仍然是大陆上数一数二的宗门。作为第一任宗主，他开创盛世，集万千信仰之力，最终飞升成神。同时代的Madore也发展了另一支折叠修炼方法。两人因为规则上的分歧而决裂，却都是划时代的人物。两个宗门因他们而建立，因他们而兴起，传承直到今天。
\end{quote}

牢大抬头问：``那这么说，在您之上，还有很多境界？''

陈星：``我停滞于 \allowbreak{}\NovelInline{I02056}{Y(1,3,4,2,5,9)}\allowbreak{} 很多年了。如果不是母猪的滋养，这辈子别想突破 \allowbreak{}\NovelInline{I02057}{Y(1,3,4,2,5,8,10,5)}\allowbreak{}。不管怎么样，我还要继续努力。''

\begin{center}\rule{0.5\linewidth}{0.5pt}\end{center}

片刻过后，牢大决定再思考一会，今天学习的知识太多了，OCF的奇妙让他有了初步的认识。但是，这母猪厂，不，应该叫做母猪城。这养猪场设有结界，一般人也无法进入。今天好不容易进来一回，那是要参观的。

母猪的大小决定了力量，不仅如此，养猪场的几万员工还可以骑上母猪，冲锋陷阵。厂内员工也在修炼，他们也是勤劳刻苦，这其中很多人都有BO以上的实力。然而，这个养猪场来源已久，并且厂内员工都在星期五这一天去猪猪供奉殿，也就是整个养猪厂的核心。

过了几天，小队继续学习OCF。

\chapter{第三十五章 初等序列与养猪场的秘密}\label{ux7b2cux4e09ux5341ux4e94ux7ae0-ux521dux7b49ux5e8fux5217ux4e0eux517bux732aux573aux7684ux79d8ux5bc6}

陈星说：``你们已经了解了基本的BOCF规则，但理论听懂了，还得靠练习巩固。先看这一题。''

\NovelDisplay{M00333R01}{%
\psi(\Omega_5^{2}+\psi_{\Omega_3}(\Omega_4^{2}+\psi_{\Omega_4}(\Omega_4\times\Omega_3)))%
}

牢大先把范围说了出来：``\allowbreak{}\NovelInline{I02058}{\psi_{\Omega_n}}\allowbreak{} 在我们使用的相应层级里，输出从 \allowbreak{}\NovelInline{I02059}{\Omega_{n-1}}\allowbreak{} 起、严格小于 \allowbreak{}\NovelInline{I02060}{\Omega_n}\allowbreak{} 的项。下界要包含等号。''

陈星：``对。遇到最里面的 \allowbreak{}\NovelInline{I02061}{\Omega_3}\allowbreak{}，这次要用 \allowbreak{}\NovelInline{I02062}{\psi_{\Omega_3}}\allowbreak{} 生成较低的索引。先取基本列下标一。''

\NovelDisplay{M00334R01}{%
\psi(\Omega_5^{2}+\psi_{\Omega_3}(\Omega_4^{2}+\psi_{\Omega_4}(\Omega_4\times\Omega_2)))%
}

蔡须坤接着写下下标二的项：

\NovelDisplay{M00335R01}{%
\psi(\Omega_5^{2}+\psi_{\Omega_3}(\Omega_4^{2}+\psi_{\Omega_4}(\Omega_4\times\psi_{\Omega_3}(\Omega_4^{2}+\psi_{\Omega_4}(\Omega_4\times\Omega_2)))))%
}

``回填进去的，正是上一轮在关键层生成的结构。外面的 \allowbreak{}\NovelInline{I02063}{\Omega_5^2}\allowbreak{} 保持不动。''

陈星点了点头，又让她检查了一遍括号：``这是下标二的项，别再说成与原式直接相等。完整的关系是：''

\NovelDisplay{M00336R01}{%
\psi(\Omega_5^{2}+\psi_{\Omega_3}(\Omega_4^{2}+\psi_{\Omega_4}(\Omega_4\times\Omega_3)))=\sup_{n<\omega}\bigl(\psi(\Omega_5^{2}+\psi_{\Omega_3}(\Omega_4^{2}+\psi_{\Omega_4}(\Omega_4\times\Omega_3)))\bigr)[n]%
}

小舞问：``如果末尾是 \allowbreak{}\NovelInline{I02064}{\Omega_3}\allowbreak{}，但前面只看见 \allowbreak{}\NovelInline{I02065}{\psi_{\Omega_2}}\allowbreak{}，该怎么办？比如这一题。''

\NovelDisplay{M00337R01}{%
\psi(\Omega_5^{2}+\psi_{\Omega_2}(\Omega_4^{2}+\psi_{\Omega_4}(\Omega_4^{2})+\Omega_3^{2}))%
}

``先别急着把下标改掉。''陈星说，``回忆两倍结构。那一题里，虽然最外层是 \allowbreak{}\NovelInline{I02066}{\psi}\allowbreak{}，回填时仍然得用 \allowbreak{}\NovelInline{I02067}{\psi_{\Omega_2}}\allowbreak{}。''

卢本韦把前一章的第二个取项重新写了出来。

\NovelDisplay{M00338R01}{%
\bigl(\psi(\Omega_2^{2}\times2)\bigr)[2]=\psi(\Omega_2^{2}+\Omega_2\times\psi_{\Omega_2}(\Omega_2^{2}+\Omega_2\times\Omega))%
}

``所以，这一题也要补上末尾所需的那一层？''

``对。但不是随意加个壳，而是在递归共尾结构要求的位置，使用相应的折叠函数。你们可以先比较两个范围。''

\NovelDisplay{M00339R01}{%
\Omega_4^2+\psi_{\Omega_4}(\Omega_4^2)>\Omega_3%
}

\NovelDisplay{M00340R01}{%
\Omega\leq\psi_{\Omega_2}\bigl(\Omega_4^2+\psi_{\Omega_4}(\Omega_4^2)\bigr)<\Omega_2%
}

``\allowbreak{}\NovelInline{I02068}{\psi_{\Omega_2}}\allowbreak{} 的输出落在 \allowbreak{}\NovelInline{I02069}{\Omega_2}\allowbreak{} 以下。如果当前需要的是 \allowbreak{}\NovelInline{I02070}{\Omega_3}\allowbreak{} 这一层的回填索引，就不能只拿这个较低层的输出照搬。需要引入 \allowbreak{}\NovelInline{I02071}{\psi_{\Omega_3}}\allowbreak{}，并且把这一题的前缀一起带上。''

奥斯卡指着最里面的平方：``现在是 \allowbreak{}\NovelInline{I02072}{\Omega_3^2}\allowbreak{}，回填以后，前面还应当留下一个 \allowbreak{}\NovelInline{I02073}{\Omega_3}\allowbreak{} 因子？''

``是。它的乘法结构也不能弄丢。''

牢大重新写下下标为一的项：

\NovelDisplay{M00341R01}{%
\psi(\Omega_5^{2}+\psi_{\Omega_2}(\Omega_4^{2}+\psi_{\Omega_4}(\Omega_4^{2})+\Omega_3\times\Omega_2))%
}

接着，第二轮把刚才生成的那一层带了回来。

\NovelDisplay{M00342R01}{%
\psi(\Omega_5^{2}+\psi_{\Omega_2}(\Omega_4^{2}+\psi_{\Omega_4}(\Omega_4^{2})+\Omega_3\times\psi_{\Omega_3}(\Omega_4^{2}+\psi_{\Omega_4}(\Omega_4^{2})+\Omega_3\times\Omega_2)))%
}

诺当沿着每一个右括号检查过去，才慢慢点头：``这次用的是 \allowbreak{}\NovelInline{I02074}{\psi_{\Omega_3}}\allowbreak{}。虽然它被写在 \allowbreak{}\NovelInline{I02075}{\psi_{\Omega_2}}\allowbreak{} 里面，承担的仍是 \allowbreak{}\NovelInline{I02076}{\Omega_3}\allowbreak{} 那一层的任务。''

陈星说：``你们可以这样记，但还得把后半句也记全。''

\begin{quote}
\textbf{末尾看共尾，回填找对应层；层级找对后，前缀、系数与起点也不能丢。}
\end{quote}

``不能看到末尾一个 \allowbreak{}\NovelInline{I02077}{\Omega_n}\allowbreak{}，就在任意位置塞一个 \allowbreak{}\NovelInline{I02078}{\psi_{\Omega_n}}\allowbreak{}。如果这个层次已经能按更低层的规则传递，还得先顺着规则找到真正的关键位置。''

他又给出了一道更深的题。

\NovelDisplay{M00343R01}{%
\psi(\Omega_6\times\Omega_5+\Omega_6\times\psi_{\Omega_2}(\Omega_4^{2}\times\psi_{\Omega_4}(\Omega_4\times\Omega_3^{\Omega_3})))%
}

``来，找出关键位置。''

蔡兴坤看了好一会儿：``外面有很大的系数，但最里面真正要处理的，是 \allowbreak{}\NovelInline{I02079}{\Omega_3^{\Omega_3}}\allowbreak{} 的指数。它需要 \allowbreak{}\NovelInline{I02080}{\Omega_3}\allowbreak{} 层的回填；前面的 \allowbreak{}\NovelInline{I02081}{\psi_{\Omega_2}}\allowbreak{} 把这层结构继续向外折叠。''

``好。先取最短的一轮，让大家看清变动的位置。''

\NovelDisplay{M00344R01}{%
\psi(\Omega_6\times\Omega_5+\Omega_6\times\psi_{\Omega_2}(\Omega_4^{2}\times\psi_{\Omega_4}(\Omega_4\times\Omega_3^{\Omega_2})))%
}

蔡须坤接着多回填了一轮。

\NovelDisplay{M00345R01}{%
\psi(\Omega_6\times\Omega_5+\Omega_6\times\psi_{\Omega_2}(\Omega_4^{2}\times\psi_{\Omega_4}(\Omega_4\times\Omega_3^{\psi_{\Omega_3}(\Omega_4^{2}\times\psi_{\Omega_4}(\Omega_4\times\Omega_3^{\Omega_2}))})))%
}

陈星：``这一步就把 \allowbreak{}\NovelInline{I02082}{\psi_{\Omega_3}}\allowbreak{} 明确写出来了。它的参数里，保留了当前关键层携带的乘数与前缀。不要把整道题外面的所有系数一股脑复制进去。''

两人把前后两项并排核对。一项看着长了很多，表示的序数却没有变大；它仍然是原序数以下的基本列项。

陈星看着她们，终于露出一点笑意：``虽然你们离我的 \allowbreak{}\NovelInline{I02083}{Y(1,3,4,2,5,9)}\allowbreak{} 还有相当距离，不过，\allowbreak{}\NovelInline{I02084}{\psi(\Omega_\omega)}\allowbreak{} 已经能够达到 \allowbreak{}\NovelInline{I02085}{Y(1,2,4,8)}\allowbreak{}。路是有的，只是还要一步一步走。''

过了一段时间，众人在BOCF修炼中遇到困难。牢大突破到了 \allowbreak{}\NovelInline{I02086}{\psi(\Omega_2^2\times2)}\allowbreak{} 的境界，其他人大多停留在 \allowbreak{}\NovelInline{I02087}{\psi(\Omega_2^2)}\allowbreak{} 上面稍微一点的位置。究其原因，他们还是难以把不同``层''的作用分清，经常在回填时选错位置。卢本韦也很苦恼。她的父亲卢关终于为她和伙伴们指明了另一条练习的路------学习\textbf{初等序列（Primitive Sequence System）}。

它结构简单，能让他们先把``从末项找关键位置、划分好部坏部、按规则复制''这些动作练熟。魂殿大供奉当年也是从这些基础出发，用 \allowbreak{}\NovelInline{I02088}{\omega}\allowbreak{}-\allowbreak{}\NovelInline{I02089}{Y}\allowbreak{} 序列走到了前所未有的 \allowbreak{}\NovelInline{I02090}{\omega}\allowbreak{}-\allowbreak{}\NovelInline{I02091}{Y(1,\omega)}\allowbreak{} 境界，达到极限斗罗层次。

``你们要明白，帅气集团依靠吞噬帅气，恐怕已经超过了大供奉的实力。''卢关看着他们，``新一代的希望，是你们。''

陈星：现在来介绍初等序列。下面是初等序列的Python代码：

\begin{verbatim}
def Primative_Sequence_System(nums):
    # 首先定义Primative Sequence System函数

    # 空序列判断
    if not nums:
        return "序列不合法"
    # 空序列是没有意义的，或者说，它代表序数0，所以第一个元素必须是1

    # 第一个元素必须是 1
    if nums[0] != 1:
        return "序列不合法"
    # 序列(1)实际上就是序数1，这样我们规定了初始值

    # 合法性校验：后一个数 <= 前一个数 + 1
    for i in range(len(nums) - 1):
        if nums[i + 1] > nums[i] + 1:
            return "序列不合法"
    # 现在，序列的合法性需要检验，你们也看到了，这个序列的前一项最多是后一项＋1，不然序列不合法

    # 查找 key_index
    last_number = nums[-1]
    # 好，现在找关键索引
    key_index = -1
    for i in range(len(nums) - 2, -1, -1):
        if nums[i] < last_number:
            # 首先，创建变量key_number，在倒数第一个元素开始寻找，而且是倒着找，
            # 找到序列左边，第一个小于last_number的数字，也就是序列第i个元素，
            # 之后把关键索引赋值i，并停止查找
            key_index = i
            break

    key_number = nums[key_index]
    good_part = nums[:key_index]
    bad_part = nums[key_index:-1]
    # 这之后，分别用得到的列表赋值这些变量
    # key_number就是那个数字，good_part（好部）是关键数字之前的列表，
    # 坏部是包括关键数字到倒数第二个数字的列表。

    return {
        "last_number": last_number,
        "key_number": key_number,
        "good_part": good_part,
        "bad_part": bad_part
    }


def expand_sequence(nums, n):
    """
    展开序列函数
    :param nums: 原始序列
    :param n: 展开次数（bad_part 复制 n 次）
    :return: 展开后的新序列 / 不合法提示
    """
    # 第一步：拆分得到 good / bad / last
    split_res = Primative_Sequence_System(nums)
    # 如果不合法，直接返回提示
    # 一现在直接调用上面的函数，把序列分割
    if isinstance(split_res, str):
        return split_res
    # 调用isinstance来判断是不是字符串

    # 提取部分
    good = split_res["good_part"]
    bad = split_res["bad_part"]
    # 展开规则：good + bad * n
    new_sequence = good + bad * n
    return new_sequence
    # 之后展开就是好部 + 坏部 * n


# 测试调用
if __name__ == "__main__":
    input_str = input("请输入空格分隔的数字序列：")
    num_list = list(map(int, input_str.split()))
    expand_times = int(input("请输入展开次数 n："))
    result = expand_sequence(num_list, expand_times)
    print("展开后的序列：", result)
    # 最后的调用函数也很简单，就是输入序列，和展开
\end{verbatim}

牢大：那么序列 \texttt{1\ 1} 是什么？

陈星：2，最后面添加1，表示序数+1。

唐三：\texttt{1\ 2} 是什么？

陈星：首先，一个序列表示一个序数，序列 \allowbreak{}\NovelInline{I02092}{X = \sup\{\text{expand\_sequence}(X, n), n < \omega\}}\allowbreak{}。\texttt{1\ 2} 展开是 \texttt{1\ 1\ 1\ 1......} 懂吗？所以 \texttt{1\ 2} 就是 \allowbreak{}\NovelInline{I02093}{\omega}\allowbreak{}。

小舞：序列加一个1表示序数+1，那么，\texttt{1\ 2\ 1} 是 \allowbreak{}\NovelInline{I02094}{\omega+1}\allowbreak{} 吗？

陈星：是的，我目前修炼的 \allowbreak{}\NovelInline{I02095}{Y}\allowbreak{} 序列，就是以它为基础的。行了，现在你们还是要巩固基础，不过下一个序数 \allowbreak{}\NovelInline{I02096}{\omega \cdot 2}\allowbreak{} 你们要想想怎么来。

牢大：是 \texttt{1\ 2\ 2} 吗？

陈星：不是，\texttt{1\ 2\ 2} 比 \allowbreak{}\NovelInline{I02097}{\omega \cdot 2}\allowbreak{} 大，现在你们打开PyCharm，运行这个代码。

蔡兴坤打开了PyCharm，得到如下结果：

\begin{verbatim}
请输入空格分隔的数字序列：1 2 2
请输入展开次数 n：3
展开后的序列： [1, 2, 1, 2, 1, 2]
进程已结束，退出代码为 0
\end{verbatim}

诺当：就是说 \allowbreak{}\NovelInline{I02098}{[1,2,2] = \sup\{[1,2], [1,2,1,2], [1,2,1,2,1,2], \ldots\}}\allowbreak{}。

奥斯卡：这是为什么，这和OCF有什么区别？

众人看着远去的陈星，顿时陷入迷茫。

众人接下来自己思考，序列 \texttt{1\ 2\ 1\ 2} 怎么展开。首先，key\_number 是第三项的1，那么 good\_part 就是前面的 \texttt{1\ 2}，然而 bad\_part 不包括 last\_number，所以 bad\_part 是 \texttt{1}。之后，删去 last\_number，之后直接复制 bad\_part，那就是：\texttt{1\ 2\ 1\ 1\ 1......} 这步正好是 \allowbreak{}\NovelInline{I02099}{\omega \cdot 2}\allowbreak{} 吗？

那么接下来，\texttt{1\ 2\ 1\ 2\ 1\ 2} 的规律就很明确了，展开正是 \texttt{1\ 2\ 1\ 2\ 1\ 1\ 1.......} 嗯，这是 \allowbreak{}\NovelInline{I02100}{\omega \cdot 3}\allowbreak{}。可以想象，\texttt{1\ 2\ 2} 就是 \texttt{1\ 2\ 1\ 2\ 1\ 2\ 1\ 2........} 就是 \allowbreak{}\NovelInline{I02101}{\omega^2}\allowbreak{} 啊。

这样一来 \texttt{1\ 2\ 2\ 1\ 2\ 2} 当然是 \allowbreak{}\NovelInline{I02102}{(\omega^2) \cdot 2}\allowbreak{}，\texttt{1\ 2\ 2\ 2} 就是 \allowbreak{}\NovelInline{I02103}{\omega^3}\allowbreak{}。找到规律后，\texttt{1\ 2\ 3} 序列的坏部是2，好部是1，自然，它就是 \allowbreak{}\NovelInline{I02104}{\sup\{[1,2], [1,2,2], [1,2,2,2], \ldots\}}\allowbreak{} 也就是 \allowbreak{}\NovelInline{I02105}{\omega^\omega}\allowbreak{}。

接下来，众人发现，如果序列 \allowbreak{}\NovelInline{I02106}{A}\allowbreak{}、\allowbreak{}\NovelInline{I02107}{B}\allowbreak{} 都是1开头，那么他们拼起来的序列就是 \allowbreak{}\NovelInline{I02108}{A, B}\allowbreak{} 表示序数的相加。比如 \texttt{1\ 2\ 2} 和 \texttt{1\ 2}，一个是 \allowbreak{}\NovelInline{I02109}{\omega^2}\allowbreak{}，一个是 \allowbreak{}\NovelInline{I02110}{\omega}\allowbreak{}，而 \texttt{1\ 2\ 2\ 1\ 2} 是 \allowbreak{}\NovelInline{I02111}{\omega^2 + \omega}\allowbreak{}。

之后，\texttt{1\ 2\ 2\ 1\ 2\ 2} 的规律和之前是一样的。

牢大：这样一来，\texttt{1\ 2\ 3\ 2} 展开是 \allowbreak{}\NovelInline{I02112}{\sup\{[1,2,3], [1,2,3,1,2,3], \ldots\}}\allowbreak{}，按照这个规律，\texttt{1\ 2\ 3\ 2} 是 \allowbreak{}\NovelInline{I02113}{(\omega^\omega) \cdot \omega}\allowbreak{}，之后，\texttt{1\ 2\ 3\ 2\ 2} 是 \allowbreak{}\NovelInline{I02114}{\omega^{\omega+2}}\allowbreak{}。

奥斯卡：这样一来，\texttt{1\ 2\ 3\ 2\ 3} 自然就是 \allowbreak{}\NovelInline{I02115}{\omega^{\omega \cdot 2}}\allowbreak{}，因为 \texttt{1\ 2\ 3\ 2\ 3} 展开就是 \texttt{1\ 2\ 3\ 2\ 2\ 2\ 2.....}

之后诺当快速分析： - \texttt{1\ 2\ 3\ 2\ 3\ 2\ 3} = \allowbreak{}\NovelInline{I02116}{\omega^{\omega \cdot 3}}\allowbreak{} - \texttt{1\ 2\ 3\ 3} 是啥呢？它的坏部是 \texttt{2\ 3}，好部是 \texttt{1}，展开是 \texttt{1\ 2\ 3\ 2\ 3\ 2\ 3\ 2\ 3.......} 之后 \texttt{1\ 2\ 3\ 3} 是 \allowbreak{}\NovelInline{I02117}{\omega^{\omega^2}}\allowbreak{}

就在这个时候，陈星过来了，说：你们好好学习初等序列，之后，等你们到了我这个境界，就会知道它有多重要。

众人说：多谢，不过，我们也不能长久住下去，不过，中级学院的日子已经过去。

就在众人准备收拾行李走的时候，卢关说：现在你们走，还不合适。

为什么？牢大说。

卢关：你们也看到了，帅气集团越来越嚣张了，如果你们散开了，恐怕还是会遇到危险，就算梅长老和我有千米感知，也没办法救你们。比赛的时候，我和梅长老就在旁边观看，这才免得被那两个兄弟弄死。帅气集团现在强大的很，实力快速增长。并且，我们不知道帅气教主的基地，也就是帅气堡的位置。

你们安心修炼几天，蔡须坤，你是三眼金猊，星斗大森林的气运之子。如果我猜想没错，你正在寻找人类和魂兽共生的一些办法。

蔡兴坤：不错，但是这要怎么解决，魂兽家族大多厌恶人类。我身上使命无比沉重，不然这世界恐怕还会在战争中。

卢关：接下来，你们应该去往星斗大森林，老夫养母猪30余年，厂内百万母猪早已和我有所共鸣，而且，我发明了一种新的方式，可以解决困扰大陆的生育问题，并且魂兽可以以魂灵形态寄托在人体。

啊？众人都惊呆了。但是，好戏才刚刚开始。

\chapter{第三十六章：魂灵？一种新的存在方式}\label{ux7b2cux4e09ux5341ux516dux7ae0ux9b42ux7075ux4e00ux79cdux65b0ux7684ux5b58ux5728ux65b9ux5f0f}

卢关话音落下，整座房间瞬间陷入一片哗然。唐三、小舞、牢大一行人满脸惊愕，万万没想到这位深耕养猪三十余年、手握百万灵气母猪的强者，不仅掌控着顶尖的OCF序数修炼法门与初等序列体系，竟还研究出了能改写斗罗大陆格局的全新生灵形态，甚至能同时破解人兽矛盾与大陆生灵的生育难题。

卢关抬手轻压，示意众人安静下来，目光透过木屋的窗棂，望向远方无边无际的巨型养猪场。漫山遍野的灵气母猪悠然踱步，周身萦绕着温润醇厚的生命灵气，无数细碎的法则波纹交织飘荡，那是数十万头灵韵母猪长年共鸣形成的本源气场，温和又磅礴。

``我知晓你们心中满是疑惑，便从头细细道来。''卢关缓缓开口，苍老的嗓音沉稳有力，缓缓揭开魂灵的奥秘，``你们切莫以为魂灵是我凭空创造的产物，也莫要觉得这片养猪场只是我们卢家三代人的根基------它的历史，早已跨越万年。万年前，这片土地上便诞生过一位神级强者，他才是魂灵之路的第一位开拓者，也是这片养猪场的初代缔造者，而我们卢家，便是他一脉单传的后人，耗费近百年光阴接续他的遗志，才最终完善了这大道。我今日的成果，既是站在父辈与祖辈的根基之上，更离不开那位万年前神级强者留下的珍贵遗产与严苛嘱托。''

众人皆是一怔，目光中多了几分敬畏与诧异，没人想到这般颠覆大陆的全新体系，竟能追溯到万年前，更没想到眼前这片看似普通的养猪场，竟藏着如此悠久的秘辛，承载着一位神级强者的未竟理想，更有着一脉单传的严苛传承。

``万年前，这片土地还未这般规模，却已是大陆上最纯净的生命灵气汇聚之地，那位神级强者，他有个称号 \textbf{Opt}，他便是我们卢家的先祖。彼时的他，已触及神境门槛，掌控着最顶尖的生命法则，修为远超如今大陆上的所有强者，却始终忧心人兽厮杀的浩劫与大陆生灵本源的衰败。万年前的大陆，魂兽族群远比今日庞大，人类与魂兽的厮杀更为惨烈，无数族群在战火中覆灭，天地间的戾气几乎要侵蚀生命本源，生灵繁衍日渐艰难，这一切，都被先祖看在眼里，记在心里。''

``先祖最早提出'生灵共生'的构想，也是第一位真正着手研究魂灵的人------那时，没有任何修炼体系可借鉴，没有序数法则，没有OCF坍缩理论，先祖仅凭神级修为的感悟，从零开始探索魂灵的奥秘。他缔造这片养猪场，饲养家畜，并非为了生计，而是看中了家畜温顺的品性与强大的生命韧性，将其作为生命灵气的载体，摸索生灵本源联结的规律，为魂灵的研究奠定基础。他想要打造一种无需猎杀、无需掠夺，能让人类与魂兽共生共荣的灵体，彻底终结人兽死局，这便是魂灵最初的雏形。''

卢关顿了顿，语气中多了几分敬佩与惋惜，继续说道：``先祖的研究，远比我们想象中艰难。他以神级精神力渗透魂兽识海，尝试剥离纯粹的本源灵力与意识碎片，搭建最初的灵气共鸣结界，试图将魂兽本源与家畜灵气融合，培育出稳定的共生灵体。可即便强如神级强者，也难以突破时代的局限------彼时天地间的戾气过重，魂兽的杀戮本能早已深入本源，即便剥离意识碎片，也难以彻底净化；更重要的是，没有初等序列的规整之力，没有非递归序数的压制之法，剥离的意识碎片混乱驳杂，根本无法形成稳定的灵体，要么瞬间溃散，要么被戾气反噬，引发狂暴。先祖耗费了整整数百年光阴，最终也未能成功培育出真正的魂灵。''

``最让先祖忧心的，并非研究失败，而是他早已看透魂灵研究的隐患------这种能改写大陆格局的力量，若是落入恶势力手中，必然会被扭曲利用，强行掠夺魂兽本源、炼制邪异灵体，届时会酿成滔天大祸，让大陆陷入更深的浩劫。因此，先祖在飞升神界前，立下严苛的祖训：\textbf{魂灵研究一脉单传，仅传卢家心性纯良、坚守正道之人，绝不外泄半分研究成果}，直至有人能真正完善魂灵体系，找到可控的制衡之法，再考虑将其公之于众，护大陆安宁。''

``他将自己所有的研究心得、实验数据、失败经验，全部记录在手札之中，藏于养猪场深处的密室，布下强大的神级结界，守护这份遗产，再将开启密室的方法与祖训一同传给唯一的后人，便在这片他亲手缔造的养猪场中飞升神界。万年来，卢家子孙世代守护这片养猪场，恪守祖训，一脉单传，默默守护着先祖的遗志与研究遗产，却始终无人能触及魂灵研究的核心，直到我的祖父卢承业。''

``百年前，祖父继承了先祖的手札与祖训，被其中的研究理念与先祖的家国情怀所震撼，于是下定决心，接续先祖与历代卢家先人的心愿，重新开启魂灵研究。祖父虽是精通生命法则的修士，修为却远不及万年前的先祖，也没有神级修为的加持，只能从最基础的地方入手，依托先祖留下的养猪场灵气底蕴，一步步摸索。他没有选择高阶魂兽，而是前往星斗大森林的边缘，寻找性情温顺、本源纯净的低阶魂兽，小心翼翼地尝试与它们建立信任，复刻先祖当年的实验，同时结合自己的感悟，弥补先祖研究中的不足。''

``他的第一步，是尝试剥离魂兽浅层的意识碎片与本源灵力。那时没有任何法则工具可以借助，祖父只能凭借自身的精神力，一点点渗透到魂兽的识海之中，小心翼翼地剥离那些不携带杀戮戾气的浅层意识------这是最危险的一步，稍有不慎，便会被魂兽的本能戾气反噬，轻则精神受损，重则识海崩塌。有一次，祖父尝试剥离一头八万年的风狼魂兽的意识碎片，因风狼的戾气比预想中浓郁，他被戾气冲击，昏迷了整整三天三夜，醒来后嘴角还带着血迹，却只是休息了一天，便又重新投入到研究中。''

``剥离出意识碎片与本源灵力后，祖父便尝试将它们与母猪的灵气结合，想要打造出一种'半灵体'------这便是魂灵的最初形态。他将剥离的意识碎片与本源灵力，小心翼翼地注入到母猪的灵气共鸣场中，试图借助母猪的温和灵气，压制魂兽意识中的暴戾，让其形成稳定的灵体。可现实远比预想中艰难，那些剥离的意识碎片极其脆弱，一旦脱离魂兽本体，便会快速溃散；即便注入母猪的灵气场中，也难以与母猪的灵气融合，要么相互排斥，要么被母猪的灵气同化，失去魂兽的本源特性。''

``祖父没有放弃，他不断调整实验方案，改变灵气注入的速度与剂量，尝试不同种类的魂兽与母猪的搭配，甚至摸索出了简单的灵气调和之法，用母猪的生命灵气稀释魂兽意识中的戾气。他还搭建了简陋的灵气结界，将实验用的母猪与魂兽意识碎片封闭其中，维持稳定的灵气环境，减少外界干扰。那些年，他失败了无数次，实验室里常常弥漫着灵气溃散的气息，案头的记录本上，写满了失败的实验数据与修改后的方案，有些页面甚至被泪水浸湿------那是他为实验失败的惋惜，也是为无法实现'共生'理想的无奈。''

``为了寻找突破，祖父走遍了大陆的名山大川，拜访了当时所有精通生命法则的修士，可那时的人们都沉迷于修炼，认为他的'共生构想'是天方夜谭，要么拒绝接见，要么敷衍了事。唯有一位隐居的老修士，认可他的理念，给了他一句点拨------\textbf{`生灵之本，在于秩序，无秩序则本源混乱'}。这句话，祖父记了一辈子，也为他的研究指明了方向，可惜他始终无法找到'规整意识、梳理本源'的方法。''

``祖父一生都在探寻人与生灵和谐共存的办法，他穷尽毕生心血，尝试了无数种实验方案，却始终没能造出真正的魂灵。他培育的半灵体，最多只能维持几个时辰，便会彻底溃散，既无法与人类建立共生羁绊，也无法自主成长，甚至偶尔还会因意识混乱，引发小规模的灵气暴走。到了晚年，祖父的身体日渐衰弱，精神力也大不如前，却依旧没有停下研究的脚步，他每天都会坐在猪圈旁，看着那些被他滋养的灵气母猪，一遍遍推演实验方案，直到生命的最后一刻。''

``他临终前，将所有的研究手稿、实验数据，连同万年前先祖留下的手札、开启密室的方法与祖训，一同交给了我的父亲，叮嘱他恪守一脉单传的祖训，一定要继续走下去，完成先祖与他未竟的'共生'理想，守护好这份研究遗产，绝不能让其落入恶人之手。祖父虽然没能取得重大成果，却在先祖研究的基础上，完善了生命灵气运转的细节，总结出更安全的魂兽意识剥离方法，为我们后续的研究避开了无数弯路。而万年前先祖留下的，不仅是手札与养猪场的灵气底蕴，更是那份跨越万年的执着与守护之心------他当年搭建的灵气共鸣阵基，至今仍在这片土地下运转，滋养着一代又一代的灵气家畜，也为我们后续的研究，提供了最核心的基础。''

说到此处，卢关的眼中泛起一丝泪光，那是对祖父一生执着的敬佩，也是对三代人传承的感慨。众人静静聆听，心中满是震撼，没人想到，如今成熟完善的魂灵体系，背后是一位老者一生的孤独探索与无数次失败的坚守。

``我的父亲，继承了祖父毕生的遗志与所有研究手稿，一辈子扎根这片养猪场，继续深耕探索。他扩大养殖规模，一代代选育驯化，培育出第一批真正拥有灵气的母猪，搭建起最初的群体生命共鸣阵法。比起祖父，他更进一步，摸索出了温和萃取魂兽多余本源灵力的办法，避开了强行掠夺带来的致命反噬，还划分了不同魂兽的本源层级，总结出无数生灵调和的经验。''

``但瓶颈依旧无法突破。''卢关神色惋惜，``父亲终生未曾接触 \allowbreak{}\NovelInline{I02118}{Y}\allowbreak{} 序列这样的递归序数法则。魂兽的意识混乱驳杂，与生俱来的杀戮天性难以压制，他炼制出的半成品灵体残缺不堪，无法稳固形态，更不能与人类魂师达成长久平衡的羁绊。一旦脱离灵气滋养，便会快速消散，甚至触发魂兽本源暴走。终其一生，他也没能完成魂灵共生的终极突破，却完善了灵力配比、本源剥离、群体结界的核心框架，补全了祖辈遗留的无数漏洞，给我留下了一套完整且可行的研究基础。''

\begin{quote}
万年前神级强者开疆拓土、坚守守护之心，卢家世代一脉单传、接续深耕，一代代步步试探、层层铺垫，虽无颠覆性的重大成果，却硬生生在先祖开辟的道路上，不断前行，既守护着研究遗产，也为最终的突破筑牢了根基。
\end{quote}

``我便是承接万年前先祖的遗志与祖训，以及父辈与祖辈的毕生心血，在此基础上，独自深耕三十余年。依托先祖留下的养猪场灵气底蕴、万年阵基，再结合日后参悟的BOCF序数法则、高阶OCF坍缩运算，以 \allowbreak{}\NovelInline{I02119}{\Omega}\allowbreak{} 系列非递归序数压制生灵野性，用初等序列梳理混乱的灵体意识，才彻底冲破了先祖、祖辈与父辈无法跨越的壁垒，解决了灵体溃散、意识暴走、血脉排斥、生育壁垒等所有难题，也找到了魂灵的制衡之法，终于有底气打破一脉单传的祖训，将先祖未完成的构想，化为了完整、成熟、可普及的魂灵体系，既完成先祖遗愿，也护大陆安宁。''

``斗罗大陆万年以来，始终困在一道无解的死局之中。''卢关收回思绪，神色重新变得郑重，开始完整详解魂灵的核心设定，``人类魂师想要变强，便要猎杀魂兽夺取魂环、掠夺魂骨，杀戮积攒无尽戾气；魂兽饱受屠戮之苦，本能仇视人类，星斗大森林常年戒备森严，人兽战火一触即发。除此之外，大陆生灵早已出现本源衰败的隐患，魂兽繁衍速度逐年锐减，高阶魂师血脉固化、生灵生育壁垒层层叠加，这两大顽疾，困扰了大陆数十万年。''

``而我借助整片养猪城的结界之力，让百万头觉醒灵气的母猪形成稳固的生命共鸣网络。常年浸润在序数法则、OCF坍缩之力与初等序列的秩序规则之中，结合祖辈流传的生命本源理论，最终彻底完善了魂灵这一全新的存在形式。''

所谓\textbf{魂灵}，是剥离魂兽本体狂暴野性与杀戮戾气，萃取其纯粹生命灵力、法则碎片与温和意识本源，再经由灵气驯化、OCF高阶序数法则平衡压制后，诞生的无形灵体生命。它没有血肉实体，不被肉身桎梏束缚，以精神力与天地生命灵气为核心载体，游离于实体生灵与抽象法则之间，是一种兼具意识、成长性与共生性的特殊存在。

魂灵的诞生条件，绝非随意便可炼制：

\begin{itemize}
\tightlist
\item
  \textbf{其一}：需要庞大且纯粹的生命灵气作为基底，自家养猪场百万头灵气母猪交织的共鸣气场，便是最完美的孕育温床；
\item
  \textbf{其二}：必须依靠初等序列的规整之力，梳理魂兽混乱的本源意识，避免灵体暴走；
\item
  \textbf{其三}：还要借助非递归序数的坍缩规则，压制魂兽刻在血脉中的杀戮本能，从根源上杜绝暴戾反噬。
\end{itemize}

{\def\LTcaptype{none} % do not increment counter
\begin{longtable}[]{@{}lll@{}}
\toprule\noalign{}
对比项 & 传统魂环/魂骨 & 魂灵 \\
\midrule\noalign{}
\endhead
\bottomrule\noalign{}
\endlastfoot
\textbf{获取方式} & 猎杀魂兽，强行掠夺本源 & 自愿共生，魂兽主动分割本源 \\
\textbf{魂兽状态} & 死亡 & 肉身完整无损，寿命不折损 \\
\textbf{成长性} & 年限固定，无法成长 & 可随魂师OCF境界同步进化 \\
\textbf{副作用} & 血脉排斥，煞气侵蚀心神 & 双向滋养，无排斥反应 \\
\textbf{因果} & 结下无解的因果孽债 & 斩断杀戮因果 \\
\end{longtable}
}

魂灵走的是\textbf{自愿共生}的道路，魂兽无需死亡，仅需主动分割一缕本源意识与灵力化作灵体，肉身完整无损、寿命不会折损，反而能借助人类魂师的修炼道路感悟高阶序数法则，突破种族天赋的天花板；人类魂师无需杀戮，以魂灵替代魂环供给魂技，彻底斩断杀戮因果，更不存在年限桎梏，魂灵可跟随魂师的OCF境界、序列修为同步进化成长。

一旁的陈星适时补充，结合自身 \allowbreak{}\NovelInline{I02120}{Y}\allowbreak{} 序列的修炼境界，完善了魂灵的法则适配性：``魂灵的强弱，严格绑定序数层级与序列认知。比如你们如今大多停留在 \allowbreak{}\NovelInline{I02121}{\psi_\Omega(\Omega_2)}\allowbreak{} 上下的境界，便只能承载对应层级的普通魂灵；若是强行融合蕴含 \allowbreak{}\NovelInline{I02122}{\Omega_3}\allowbreak{} 法则的神兽魂灵，会直接导致识海崩塌、OCF展开紊乱，修炼根基彻底损毁。我停滞 \allowbreak{}\NovelInline{I02123}{Y(1,3,4,2,5,9)}\allowbreak{} 多年，能容纳更强的本源魂灵，便是高阶序列带来的底气。''

魂灵成型后，会自主寄宿在魂师的精神识海与血脉脉络之中，形成\textbf{双向滋养的平衡羁绊}：

\begin{itemize}
\tightlist
\item
  \textbf{魂师 → 魂灵}：以自身精神力、修炼的法则灵气滋养魂灵，稳固灵体形态；
\item
  \textbf{魂灵 → 魂师}：反哺生命本源，弱化魂师突破序数壁垒时的法则撕裂反噬，衍生专属法则魂技，补齐修炼短板。
\end{itemize}

二者一荣俱荣、一损俱损，羁绊牢牢绑定，从根本上杜绝了魂灵叛变的可能。

{\def\LTcaptype{none} % do not increment counter
\begin{longtable}[]{@{}
  >{\raggedright\arraybackslash}p{(\linewidth - 4\tabcolsep) * \real{0.3333}}
  >{\raggedright\arraybackslash}p{(\linewidth - 4\tabcolsep) * \real{0.3333}}
  >{\raggedright\arraybackslash}p{(\linewidth - 4\tabcolsep) * \real{0.3333}}@{}}
\toprule\noalign{}
\begin{minipage}[b]{\linewidth}\raggedright
类别
\end{minipage} & \begin{minipage}[b]{\linewidth}\raggedright
来源
\end{minipage} & \begin{minipage}[b]{\linewidth}\raggedright
特性
\end{minipage} \\
\midrule\noalign{}
\endhead
\bottomrule\noalign{}
\endlastfoot
\textbf{普通魂灵} & 寻常魂兽凝练 & 基础生命灵力充沛，适配低阶魂师稳固修为 \\
\textbf{元素魂灵} & 掌控水火风雷等自然元素的魂兽 & 自带专属元素法则，大幅强化属性战力 \\
\textbf{神兽魂灵} & 星斗大森林上古神兽、远古异种 & 蕴藏高阶不动点序数碎片，成长性顶尖，最契合气运之子 \\
\textbf{集体本源魂灵} & 群居生灵共鸣凝聚 & 不依附单一个体，可编织大范围生命法则网络，破解生育难题的核心关键 \\
\end{longtable}
}

小舞听得满心触动，轻声发问：``卢前辈，这种集体魂灵，真的能解决生灵生育的困境吗？''

卢本韦向前踏出一步，接过话语，眼中满是对父亲与祖辈研究的敬佩：``没错，祖辈早已发现群居生灵的生命联结，父亲完善了共鸣阵法，而祖父与父亲遗留的研究数据，让父亲看清了血脉壁垒的本质。集体本源魂灵相互交织，会编织出覆盖整片大陆的生命法则屏障，修复生灵常年因杀戮、法则破损留下的本源损伤，弱化血脉隔阂与生育壁垒。无论是繁衍艰难的魂兽族群，还是后期血脉固化的人类魂师，都会被温和的生命灵气持续滋养，长久下去，大陆生灵的繁衍难题自然会迎刃而解。''

一直沉默的蔡须坤此刻浑身微微震颤，三眼金猊的气运血脉本能被深深触动。她身负平衡人兽两族的天命，长久以来都因魂兽对人类的偏见而束手无策，此刻终于看到了希望，急忙开口追问：``卢前辈，魂灵共生，真的能彻底化解人与魂兽的仇恨吗？''

``这正是我推行魂灵的最终目的，也是卢家三代人共同的心愿。''卢关目光坚定，``万年仇恨的根源是掠夺与杀戮，如今魂兽无需惧怕猎杀，反而和人类共生突破种族上限；人类不必依靠屠戮变强，以平和的方式获取力量。当共生成为常态，隔阂会慢慢消融，人兽共存，将不再是空谈。''

奥斯卡挠了挠头，满脸疑惑：``既然魂灵这般完美，就没有任何隐患吗？若是被恶人利用，强行炼制邪异魂灵该如何是好？''

``隐患从来都存在。''卢关的神色骤然凝重，话语间多了几分肃杀，``这也是我迟迟没有对外公开魂灵研究，也不让你们贸然离开养猪城的原因，更是万年前先祖立下一脉单传祖训的核心用意。\textbf{帅气集团}横行大陆，以吞噬帅气本源、掠夺法则力量为生，行事残暴不择手段。一旦他们知晓魂灵的存在，必然会抓捕大量魂兽，强行剥离本源、扭曲炼制邪魂灵，用来强化自身诡异的特殊修为，甚至妄图篡改生命法则，以邪异的方式掌控生灵生死。先辈们当年也曾遭遇邪道势力觊觎，祖父的不少研究记录险些被毁，这也是我们恪守祖训、不敢贸然推广研究的缘由，如今我完善了魂灵体系与制衡之法，才敢让你们知晓这份跨越万年的传承与守护。''

唐三闻言眉头紧锁，瞬间明白当下的处境。众人如今刚刚吃透进阶OCF与初等序列的基础，修为尚且稚嫩，若是分散行动，以帅气集团的恐怖实力，根本无力抗衡。

``所以接下来，你们的目的地，便是星斗大森林。''卢关缓缓道出后续安排，``星斗是大陆魂兽的根源之地，蕴藏无数上古本源，是孕育顶尖神兽魂灵的核心秘境。蔡须坤，你身为三眼金猊，身负天地气运，是联结人兽两族最好的纽带。你们一行人前往星斗，一边打磨自身修为，吃透 \allowbreak{}\NovelInline{I02124}{\Omega_3}\allowbreak{} 序数展开规则、稳固PRSS序列基础，一边联络各大魂兽族群，传播魂灵共生的理念。这也是延续我卢家三代人的夙愿，让先辈未完成的理想，真正照耀整片大陆。''

众人豁然开朗，原本迷茫的修行之路瞬间清晰。牢大握紧双拳，心中暗下决心，一定要尽快突破 \allowbreak{}\NovelInline{I02125}{\psi_\Omega(\Omega_2 \cdot 2)}\allowbreak{} 之上的境界，吃透更高阶的OCF展开，守护这份来之不易的新生希望；唐三低头沉思，不断推演魂灵与序数坍缩函数的结合方式，规划着后续的修炼方向；蔡兴坤、蔡须坤姐妹相视一眼，终于找到了完成自身使命的道路。

OCF序数攀登无穷法则，三代人百年深耕铸就魂灵共生之道，三条全新的道路相互交织，让果糕七怪的眼界彻底打开，明白了这片大陆的力量远不止魂环与武魂那般简单。

夕阳落下，养猪城的结界泛起淡淡的金光，万千灵气母猪发出温顺的低鸣。养猪场深处的密室之中，一排排泛黄的古籍手札静静存放，最上方的一本，封皮刻着古朴的纹路，字迹早已模糊，却能隐约辨认出 \textbf{``Opt''} 三字，那是万年前神级强者的毕生心血

\chapter{第三十七章：传递，到达EBO的BOCF}\label{ux7b2cux4e09ux5341ux4e03ux7ae0ux4f20ux9012ux5230ux8fbeeboux7684bocf}

之前，众人经过分析 \texttt{1\ 2\ 3\ 3} 是 \allowbreak{}\NovelInline{I02126}{\omega^{\omega^2}}\allowbreak{}，那么 \texttt{1\ 2\ 3\ 3\ 2} 是 \allowbreak{}\NovelInline{I02127}{\omega^{\omega^2+1}}\allowbreak{}，自然 \texttt{1\ 2\ 3\ 3\ 2\ 3} 是 \allowbreak{}\NovelInline{I02128}{\omega^{\omega^2+\omega}}\allowbreak{}，之后，\texttt{1\ 2\ 3\ 3\ 2\ 3\ 2\ 3} 是 \allowbreak{}\NovelInline{I02129}{\omega^{\omega^2+\omega \cdot 2}}\allowbreak{}。

而对于 \texttt{1\ 2\ 3\ 3\ 2\ 3\ 3}，那展开是 \texttt{1\ 2\ 3\ 3\ 2\ 3\ 2\ 3\ 2\ 3\ 2\ 3......} \texttt{2\ 3} 是坏部，这么说，\texttt{1\ 2\ 3\ 3\ 2\ 3\ 3} 是 \allowbreak{}\NovelInline{I02130}{\omega^{(\omega^2) \cdot 2}}\allowbreak{} 也就很自然了。

经过分析，\texttt{1\ 2\ 3\ 3\ 3} 展开是 \texttt{1\ 2\ 3\ 3\ 2\ 3\ 3\ 2\ 3\ 3......} 众人几乎不怎么费力气，就分析到了 \allowbreak{}\NovelInline{I02131}{\omega^{\omega^3}}\allowbreak{}，之后，\texttt{1\ 2\ 3\ 4} 是再自然不过的事情。但是，这样一来，初等序列就只有 \allowbreak{}\NovelInline{I02132}{\varepsilon_0}\allowbreak{} 了，他们不明白，为什么要学这个，这功法太弱了。唐三很困惑，为什么要学习这东西，有什么用呢？

陈星这个时候走了过来。他看过众人留下的推演，在旁边重新写下了三个式子。

\NovelDisplay{M00346R01}{%
\psi(0)=1%
}

\NovelDisplay{M00347R01}{%
\psi(\Omega)=\varepsilon_0%
}

\NovelDisplay{M00348R01}{%
\psi_{\Omega_2}(0)=\Omega%
}

``后面的练习统一使用BOCF，也就是Buchholz序数折叠函数。''他说，``这里的 \allowbreak{}\NovelInline{I02133}{\psi}\allowbreak{} 都是 \allowbreak{}\NovelInline{I02134}{\psi_\Omega}\allowbreak{} 的简写。先把起点记牢，同一个符号在不同体系里未必表示同一个序数。''

他又在旁边补上了这一阶段采用的约定。

\NovelDisplay{M00349R01}{%
\Omega=\omega_1^{\mathrm{CK}}%
}

``这里的 \allowbreak{}\NovelInline{I02135}{\Omega}\allowbreak{} 是Church--Kleene序数，本身可数，却已经超出了递归序数的范围。后面的 \allowbreak{}\NovelInline{I02136}{\Omega_n}\allowbreak{} 和 \allowbreak{}\NovelInline{I02137}{\Omega_\omega}\allowbreak{} 也都按可数序数来理解。我们借助更高的非递归层级，构造 \allowbreak{}\NovelInline{I02138}{\Omega}\allowbreak{} 以下更大的递归序数。''

唐三盯着第一行：``所以，光记住一个折叠函数的符号，还不够？''

``对。参数的起点和展开的规则，都得记清。比如一个很容易混淆的地方：''

\NovelDisplay{M00350R01}{%
\psi(\Omega+1)=\varepsilon_0\times\omega%
}

\NovelDisplay{M00351R01}{%
\psi(\Omega\times2)=\varepsilon_1%
}

``BOCF里，参数加一，和走到下一个 \allowbreak{}\NovelInline{I02139}{\varepsilon}\allowbreak{} 数，是两件不同的事。''

陈星接着给出了一道复杂的练习：

\NovelDisplay{M00352R01}{%
\psi(\Omega_3^{2}\times\psi_{\Omega_3}(\Omega_3\times\Omega_2^{\Omega^{2}}))%
}

``连续展开五次。每次从基本列里选一项，写清楚选的是哪一项就行。选出的项应当比展开前小，不能把展开的方向弄反。''

小舞先抬起手：``这里最深处是 \allowbreak{}\NovelInline{I02140}{\Omega^2}\allowbreak{}，但不能只看见一个大下标，就直接套上 \allowbreak{}\NovelInline{I02141}{\psi_{\Omega_2}}\allowbreak{}。要先看它的递归共尾结构怎么传到外层。''

她先说明，\allowbreak{}\NovelInline{I02142}{\operatorname{cf}_{\mathrm{rec}}}\allowbreak{} 记录的是本章展开规则中的递归共尾类型，随后写出了两个较短的式子。

\NovelDisplay{M00353R01}{%
\operatorname{cf}_{\mathrm{rec}}(\Omega^2)=\Omega%
}

\NovelDisplay{M00354R01}{%
\operatorname{cf}_{\mathrm{rec}}\bigl(\psi_{\Omega_3}(\Omega_3\times\Omega_2^{\Omega^2})\bigr)=\Omega%
}

``这一层的 \allowbreak{}\NovelInline{I02143}{\Omega}\allowbreak{} 索引结构可以穿过 \allowbreak{}\NovelInline{I02144}{\psi_{\Omega_3}}\allowbreak{} 传出来。最外面的 \allowbreak{}\NovelInline{I02145}{\psi_\Omega}\allowbreak{} 再按规则回填，得到由自然数索引的基本列。''

陈星点了点头：``先算最短的一轮回填。''

\NovelDisplay{M00355R01}{%
\psi(\Omega_3^{2}\times\psi_{\Omega_3}(\Omega_3\times\Omega_2^{\Omega\times\omega}))%
}

``再多回填一轮，就得到我们这次实际选择的项。下面所有方括号都表示基本列取项；这里沿用从零开始编号的约定，\allowbreak{}\NovelInline{I02146}{[2]}\allowbreak{} 指下标为二的项。''

小舞写下了第一次展开：

\NovelDisplay{M00356R01}{%
\psi(\Omega_3^{2}\times\psi_{\Omega_3}(\Omega_3\times\Omega_2^{\Omega\times\psi(\Omega_3^{2}\times\psi_{\Omega_3}(\Omega_3\times\Omega_2^{\Omega\times\omega}))}))%
}

唐三抄到中间时，把纸上的原式又往指数里套了一层。

牢大伸手按住他的笔：``先等等。你现在复制的，是外面整道题，还是这一轮规则真正要求回填的部分？''

``不是越完整越好吗？''

``你多放进去一层，结果可能反而比原来的指数还大。基本列的项若比原序数更大，这就已经错了。''

陈星把最外层的系数与内层的指数分别指给他们看：``折叠函数不是随便嵌套的壳。要顺着末项的共尾结构找回填位置，同时检查新参数能否成为标准项。不能因为出现了同一个 \allowbreak{}\NovelInline{I02147}{\psi}\allowbreak{}，就把整串式子原样复制。''

牢大没有急着继续堆长式，而是指着第一次展开中较小的那个内层指数：``现在轮到它下降了。为了看清位置，接下来的四次都取基本列下标一。''

他写出了第二次展开：

\NovelDisplay{M00357R01}{%
\psi(\Omega_3^{2}\times\psi_{\Omega_3}(\Omega_3\times\Omega_2^{\Omega\times\psi(\Omega_3^{2}\times\psi_{\Omega_3}(\Omega_3\times\Omega_2^{\Omega}))}))%
}

卢本韦接过笔：``也就是说，这一步变的是内层的 \allowbreak{}\NovelInline{I02148}{\Omega\times\omega}\allowbreak{}。前面已经搭好的外层结构，要留着。''

第三次：

\NovelDisplay{M00358R01}{%
\psi(\Omega_3^{2}\times\psi_{\Omega_3}(\Omega_3\times\Omega_2^{\Omega\times\psi(\Omega_3^{2}\times\psi_{\Omega_3}(\Omega_3\times\Omega_2^{\omega}))}))%
}

``到了这里，最里面的 \allowbreak{}\NovelInline{I02149}{\Omega}\allowbreak{} 又降成了 \allowbreak{}\NovelInline{I02150}{\omega}\allowbreak{}。''卢本韦甩了甩有些发酸的手，``长得吓人，但每次真正动的地方并没有那么多。''

陈星：``很好。小舞，第四次。''

\NovelDisplay{M00359R01}{%
\psi(\Omega_3^{2}\times\psi_{\Omega_3}(\Omega_3\times\Omega_2^{\Omega\times\psi(\Omega_3^{2}\times\psi_{\Omega_3}(\Omega_3\times\Omega_2))}))%
}

奥斯卡凑近看了一遍，才写下第五次：

\NovelDisplay{M00360R01}{%
\psi(\Omega_3^{2}\times\psi_{\Omega_3}(\Omega_3\times\Omega_2^{\Omega\times\psi(\Omega_3^{2}\times\psi_{\Omega_3}(\Omega_3\times\Omega))}))%
}

``五次做完了。''他长出一口气，``这回没有凭感觉多加一层。''

``还可以继续。''牢大把纸往自己面前拉了拉，又接着往下算了三步。

\NovelDisplay{M00361R01}{%
\psi(\Omega_3^{2}\times\psi_{\Omega_3}(\Omega_3\times\Omega_2^{\Omega\times\psi(\Omega_3^{2}\times\psi_{\Omega_3}(\Omega_3\times\omega))}))%
}

\NovelDisplay{M00362R01}{%
\psi(\Omega_3^{2}\times\psi_{\Omega_3}(\Omega_3\times\Omega_2^{\Omega\times\psi(\Omega_3^{2}\times\psi_{\Omega_3}(\Omega_3))}))%
}

\NovelDisplay{M00363R01}{%
\psi(\Omega_3^{2}\times\psi_{\Omega_3}(\Omega_3\times\Omega_2^{\Omega\times\psi(\Omega_3^{2}\times\Omega_2^{\Omega_2})}))%
}

蔡须坤注意到最后一行的变化：``这一次不只是把里面某个数变小。\allowbreak{}\NovelInline{I02151}{\psi_{\Omega_3}(\Omega_3)}\allowbreak{} 本身也展开了。''

``对。''陈星说，``在BOCF中，它与零参数的值不同。''

\NovelDisplay{M00364R01}{%
\psi_{\Omega_3}(0)=\Omega_2%
}

\NovelDisplay{M00365R01}{%
\psi_{\Omega_3}(\Omega_3)=\varepsilon_{\Omega_2+1}%
}

\NovelDisplay{M00366R01}{%
\bigl(\psi_{\Omega_3}(\Omega_3)\bigr)[1]=\Omega_2^2%
}

``最后这个等号，说的是单独展开这一层时的基本列取项，不能把右边直接说成原序数。刚才整道题的回填还受外层规则影响，不能将每一部分都机械地换成各自相同下标的项。你们之前那种把展开步骤全部用普通等号串起来的写法，也要改。''

众人重新比过每一步。式子在展开时有时会明显变长，表示的序数却始终下降；复制出来的结构，是为了描述更复杂的下降过程。

蔡须坤又问：``那如果在较低层的折叠里面，不断塞进已经折叠过的东西，就能一直往上走吗？''

陈星摇头，写下了一个较短的反例。

\NovelDisplay{M00367R01}{%
\psi(\psi(\Omega))=\psi(\Omega)%
}

``右边就是 \allowbreak{}\NovelInline{I02152}{\varepsilon_0}\allowbreak{}。把它再次塞进外面的 \allowbreak{}\NovelInline{I02153}{\psi}\allowbreak{}，并不会自动超过它。按这一套BOCF的标准项规则，左边需要化简，不能只数括号来判断强弱。''

唐三看了一会儿：``所以，关键不在复制了多少次，而在复制之后，新参数有没有真正越过原来的结构。''

``正是如此。遇到更深的例子，同样先检查末项的递归共尾类型，再检查它能穿过哪些折叠层。外层的函数下标很高，不代表它就一定是当前需要发动的那一层。''

蔡须坤又试着处理了一道涉及更高下标的题。

\NovelDisplay{M00368R01}{%
\psi(\Omega_6\times\Omega_5+\Omega_6\times\psi_{\Omega_2}(\Omega_4^{2}\times\psi_{\Omega_4}(\Omega_4\times\Omega_3^{\Omega_3})))%
}

她没有把整条式子胡乱复制，而是先沿最里面的 \allowbreak{}\NovelInline{I02154}{\Omega_3}\allowbreak{} 找到需要回填的位置。取基本列下标一，可以得到：

\NovelDisplay{M00369R01}{%
\psi(\Omega_6\times\Omega_5+\Omega_6\times\psi_{\Omega_2}(\Omega_4^{2}\times\psi_{\Omega_4}(\Omega_4\times\Omega_3^{\Omega_2})))%
}

``这一步的回填用到了 \allowbreak{}\NovelInline{I02155}{\psi_{\Omega_3}(0)=\Omega_2}\allowbreak{}。''她说，``前面那些更大的系数还在原位。每一层都检查过，才知道究竟该动哪一处。''

众人感觉这真是一场酣畅淋漓的展开。然而，序数BO仍然是他们眼前最大的壁垒。自从初等序列分析到了极限，他们总觉得 \texttt{1\ 2\ 3\ 4……} 不该是终点；倘若还可以接着写下 \texttt{1\ 3}，会发生什么？可那已经越出了他们眼下这套初等序列的规则。

苦修之中，他们也渐渐发现，BOCF的展开和PRSS有相似的地方。

小舞把一段初等序列的展开写在纸上。这次她特意用了箭头，提醒自己这是展开过程的示意。

\NovelDisplay{M00370R01}{%
[1\;2\;2\;2]\longmapsto[1\;2\;2\;1\;2\;2\;\cdots]%
}

``都是找到关键位置，再按照规则复制一段结构。''

陈星说：``可以这样理解它们的共同之处。不过，不能因此断言任意一个加法、乘法、乘方表达式，都能用同一条复制公式处理。末项不同，递归共尾类型不同，规则就可能不同。拿一个确定的BOCF例子来比较，反而更清楚。''

\NovelDisplay{M00371R01}{%
\psi(\Omega\times2)=\sup_{n<\omega}\bigl(\psi(\Omega\times2)\bigr)[n]%
}

牢大把笔移到上确界的右边：``写出基本列的符号还不够，得看每一项究竟怎样生成。''

``对。''陈星说，``不能把任意一列套进上确界，就称它为基本列。真正算这一步，得依照BOCF的回填规则。''

他写下了其中的前三个非零下标的项。

\NovelDisplay{M00372R01}{%
\bigl(\psi(\Omega\times2)\bigr)[1]=\psi(\Omega+1)%
}

\NovelDisplay{M00373R01}{%
\bigl(\psi(\Omega\times2)\bigr)[2]=\psi\bigl(\Omega+\psi(\Omega+1)\bigr)%
}

\NovelDisplay{M00374R01}{%
\bigl(\psi(\Omega\times2)\bigr)[3]=\psi\bigl(\Omega+\psi(\Omega+\psi(\Omega+1))\bigr)%
}

``每次把上一步得到的结构带回来，却仍然是原序数以下的项。这样的复制、回填、下降，是你们理解Hydra结构时真正要抓住的地方。''

为了构造增长更快的大数函数，人们引入序数，并通过FGH把合适的序数记号用于定义自然数函数。为了构造更大的可数序数记号，又引入 \allowbreak{}\NovelInline{I02156}{\Omega_n}\allowbreak{} 这些更高层的形式符号，用折叠函数描述它们在较低层留下的结构。

陈星提醒他们，折叠并不是把一个大序数原封不动地搬到小序数里。它有自己的闭包条件和标准项规则。

\begin{itemize}
\tightlist
\item
  \allowbreak{}\NovelInline{I02157}{\psi_{\Omega_3}}\allowbreak{} 的值落在 \allowbreak{}\NovelInline{I02158}{\Omega_2\leq\alpha<\Omega_3}\allowbreak{} 这一层；
\item
  \allowbreak{}\NovelInline{I02159}{\psi_{\Omega_2}}\allowbreak{} 的值落在 \allowbreak{}\NovelInline{I02160}{\Omega\leq\alpha<\Omega_2}\allowbreak{} 这一层；
\item
  \allowbreak{}\NovelInline{I02161}{\psi=\psi_\Omega}\allowbreak{} 用来构造 \allowbreak{}\NovelInline{I02162}{\Omega}\allowbreak{} 以下的序数记号。
\end{itemize}

``尤其是前两层，零参数的值就在下界上。不能漏掉这个起点，说它们一上来就是某个更大的 \allowbreak{}\NovelInline{I02163}{\varepsilon}\allowbreak{} 数。''

如果把FGH也类比为一种``折叠''，就像是把无限层次的记号落实为自然数的增长过程。不过，FGH不是这些 \allowbreak{}\NovelInline{I02164}{\psi}\allowbreak{} 函数中的一员；这种类比能帮助理解，却不能代替定义。

陈星在纸上写出了他们正在追赶的界限。

\NovelDisplay{M00375R01}{%
\mathrm{BO}=\sup_{1\leq n<\omega}\psi(\Omega_n^2)%
}

\NovelDisplay{M00376R01}{%
\mathrm{BO}=\psi(\Omega_\omega)%
}

``这里从各有限层的平方取起。\allowbreak{}\NovelInline{I02165}{\Omega_n^2}\allowbreak{} 随着层数增加，仍然共尾于 \allowbreak{}\NovelInline{I02166}{\Omega_\omega}\allowbreak{}。你们也可以选另一条合适的共尾列，但要确认它确实达到了同一个上确界。''

``BO是一个新的关口，走到这里，前面有限层数的折叠还可以继续往上推广。今天就点拨到这里，后续的努力，要靠你们自己。''

过了三个月，牢大率先突破BO的实力，之后是蔡须坤，再之后是诺当他们。但是，对于BO以上的结构，他们还不太清楚。于是陈星继续科普BO之后的BOCF。

陈星：``先看后继。BO已经是 \allowbreak{}\NovelInline{I02167}{\varepsilon}\allowbreak{} 函数的不动点，所以它之后的下一个 \allowbreak{}\NovelInline{I02168}{\varepsilon}\allowbreak{} 数，可以写成这样：''

\NovelDisplay{M00377R01}{%
\psi(\Omega_\omega+\Omega)=\varepsilon_{\mathrm{BO}+1}%
}

``这里增加的是 \allowbreak{}\NovelInline{I02169}{\Omega}\allowbreak{}。若只增加一，就只是：''

\NovelDisplay{M00378R01}{%
\psi(\Omega_\omega+1)=\mathrm{BO}\times\omega%
}

唐三点头：``增加一个单位，和跨过下一个 \allowbreak{}\NovelInline{I02170}{\varepsilon}\allowbreak{} 数，果然不是一回事。''

陈星：``再回忆一个有限层的例子。''

\NovelDisplay{M00379R01}{%
\psi(\Omega_2^2+\Omega^2)%
}

``它对应从适当起点迭代下面这个映射而得到的不动点。''

\NovelDisplay{M00380R01}{%
\alpha\longmapsto\psi(\Omega_2^2+\Omega\times\alpha)%
}

唐三：``那么在BO之后，也有类似的结构？''

\NovelDisplay{M00381R01}{%
\psi(\Omega_\omega+\Omega^2)%
}

\NovelDisplay{M00382R01}{%
\alpha\longmapsto\psi(\Omega_\omega+\Omega\times\alpha)%
}

陈星：``对。但不要省掉参数里的乘数 \allowbreak{}\NovelInline{I02171}{\Omega}\allowbreak{}。少了这个乘数，迭代的起点与过程都会改变。''

他接着将两个式子并排写下。

\NovelDisplay{M00383R01}{%
\psi(\Omega_2^2\times2)%
}

\NovelDisplay{M00384R01}{%
\psi(\Omega_\omega\times2)%
}

``到有限层的 \allowbreak{}\NovelInline{I02172}{\Omega_2}\allowbreak{}，可以利用相应的折叠层处理。到了 \allowbreak{}\NovelInline{I02173}{\Omega_\omega}\allowbreak{}，要先看它自身的共尾列。''

\NovelDisplay{M00385R01}{%
\Omega_\omega=\sup_{1\leq n<\omega}\Omega_n%
}

``这里的 \allowbreak{}\NovelInline{I02174}{\Omega_\omega}\allowbreak{} 仍然可数。''陈星说，``不过，它已经大于 \allowbreak{}\NovelInline{I02175}{\Omega=\omega_1^{\mathrm{CK}}}\allowbreak{}，不能因此称它为递归序数。眼下要抓住的是这条展开路径：沿着 \allowbreak{}\NovelInline{I02176}{\Omega_1,\Omega_2,\Omega_3,\ldots}\allowbreak{}，就能逐层逼近 \allowbreak{}\NovelInline{I02177}{\Omega_\omega}\allowbreak{}。''

奥斯卡问：``之前说的递归共尾度，是不是就在说基本列的长度？''

``更准确地说，是在本章允许的展开规则下，需要哪一层的索引来构造基本列。''陈星说，``这里讲的是递归共尾度。我们暂时不展开更深的定义，先比较两种情形：\allowbreak{}\NovelInline{I02178}{\Omega_\omega}\allowbreak{} 可以沿着已有的 \allowbreak{}\NovelInline{I02179}{\Omega_n}\allowbreak{} 逐层逼近；后继层 \allowbreak{}\NovelInline{I02180}{\Omega_{\omega+1}}\allowbreak{} 则作为新的容许层级处理，不能照搬前面那条列。''

诺当把两种情形并排记下。这几处公式使用的都是本章约定的递归共尾记法。

\NovelDisplay{M00386R01}{%
\operatorname{cf}_{\mathrm{rec}}(\Omega_\omega)=\omega%
}

\NovelDisplay{M00387R01}{%
\operatorname{cf}_{\mathrm{rec}}(\Omega_{\omega+1})=\Omega_{\omega+1}%
}

陈星：``第一行表示，按自然数下标依次取 \allowbreak{}\NovelInline{I02181}{\Omega_n}\allowbreak{}，就有现成的逼近路径。第二行表示，在这套递归展开规则中，\allowbreak{}\NovelInline{I02182}{\Omega_{\omega+1}}\allowbreak{} 本身承担新一层的索引作用。要处理这一层，就得使用相应的 \allowbreak{}\NovelInline{I02183}{\psi_{\Omega_{\omega+1}}}\allowbreak{}。别把这两种展开方式混在一起。''

唐三：``那么，原来那个两倍的式子，可以先写成上确界？''

\NovelDisplay{M00388R01}{%
\psi(\Omega_\omega\times2)=\sup_{1\leq n<\omega}\psi(\Omega_\omega+\Omega_n^2)%
}

陈星：``可以。这里与BO处一样，选平方项仍是一条共尾列。''

小舞刚要把上确界的左端再抄进右边，牢大已经把手指伸了过来。

``刚才才检查过的毛病，别又犯一次。''

她笑着收回笔，改用一个确定的有限层展示基本列回填。

\NovelDisplay{M00389R01}{%
\bigl(\psi(\Omega_\omega+\Omega_2^2)\bigr)[1]=\psi(\Omega_\omega+\Omega_2\times\Omega)%
}

\NovelDisplay{M00390R01}{%
\bigl(\psi(\Omega_\omega+\Omega_2^2)\bigr)[2]=\psi\bigl(\Omega_\omega+\Omega_2\times\psi_{\Omega_2}(\Omega_\omega+\Omega_2\times\Omega)\bigr)%
}

陈星：``注意，这两行都是比原式小的项。要讨论平方，也得看清它的乘法结构，不能把共尾项换到指数上去。''

\NovelDisplay{M00391R01}{%
\psi(\Omega_\omega^2)=\sup_{1\leq n<\omega}\psi(\Omega_\omega\times\Omega_n)%
}

``如果误写成 \allowbreak{}\NovelInline{I02184}{\Omega_\omega^{\Omega_n}}\allowbreak{}，那就换成另一道题了。''

他接着问：``那么，\allowbreak{}\NovelInline{I02185}{\varepsilon_{\Omega_\omega+1}}\allowbreak{} 又该怎么写？''

牢大想起了前面的零参数：``\allowbreak{}\NovelInline{I02186}{\psi_{\Omega_{\omega+1}}}\allowbreak{} 这一层会出现，但在BOCF中，它的零参数只是 \allowbreak{}\NovelInline{I02187}{\Omega_\omega}\allowbreak{}。''

\NovelDisplay{M00392R01}{%
\psi_{\Omega_{\omega+1}}(0)=\Omega_\omega%
}

\NovelDisplay{M00393R01}{%
\psi_{\Omega_{\omega+1}}(\Omega_{\omega+1})=\varepsilon_{\Omega_\omega+1}%
}

陈星：``正是如此。因此，这个较高的关口，可以写成：''

\NovelDisplay{M00394R01}{%
\psi\bigl(\varepsilon_{\Omega_\omega+1}\bigr)=\psi\bigl(\psi_{\Omega_{\omega+1}}(\Omega_{\omega+1})\bigr)%
}

\NovelDisplay{M00395R01}{%
\psi\bigl(\varepsilon_{\Omega_\omega+1}\bigr)=\psi(\Omega_{\omega+1})%
}

``内层虽然可以用 \allowbreak{}\NovelInline{I02188}{\varepsilon}\allowbreak{} 数表示，最外层仍要按BOCF化成标准项，不能仅凭括号深浅比较。''

奥斯卡：``下一个 \allowbreak{}\NovelInline{I02189}{\varepsilon}\allowbreak{} 数呢？''

\NovelDisplay{M00396R01}{%
\psi_{\Omega_{\omega+1}}(\Omega_{\omega+1}\times2)=\varepsilon_{\Omega_\omega+2}%
}

\NovelDisplay{M00397R01}{%
\psi\bigl(\psi_{\Omega_{\omega+1}}(\Omega_{\omega+1}\times2)\bigr)=\psi(\Omega_{\omega+1}\times2)%
}

陈星：``这里是 \allowbreak{}\NovelInline{I02190}{\varepsilon}\allowbreak{} 下标再加一，得到 \allowbreak{}\NovelInline{I02191}{\Omega_\omega+2}\allowbreak{}。别把前一个 \allowbreak{}\NovelInline{I02192}{\varepsilon}\allowbreak{} 数整个塞回它的下标。''

他又给出了两层折叠之间的对应。

\NovelDisplay{M00398R01}{%
\psi\bigl(\psi_{\Omega_{\omega+1}}(\Omega_{\omega+1}^2)\bigr)=\psi(\Omega_{\omega+1}^2)%
}

``再往后，比如这个：''

\NovelDisplay{M00399R01}{%
\psi(\Omega_{\omega+1}^2+\Omega_\omega)%
}

卢本韦：``末项是 \allowbreak{}\NovelInline{I02193}{\Omega_\omega}\allowbreak{}，沿着 \allowbreak{}\NovelInline{I02194}{\Omega_n}\allowbreak{} 取共尾项，前面的平方不动。''

\NovelDisplay{M00400R01}{%
\psi(\Omega_{\omega+1}^2+\Omega_\omega)=\sup_{1\leq n<\omega}\psi(\Omega_{\omega+1}^2+\Omega_n)%
}

蔡须坤：``下标本身也可以继续走。比如从 \allowbreak{}\NovelInline{I02195}{\omega+n}\allowbreak{} 走向 \allowbreak{}\NovelInline{I02196}{\omega\times2}\allowbreak{}？''

\NovelDisplay{M00401R01}{%
\sup_{1\leq n<\omega}\psi(\Omega_{\omega+n}^2)=\psi(\Omega_{\omega\times2})%
}

陈星：``可以。到后继下标，就有相应的折叠上界。例如下面这一层的两个起点，必须区分。''

\NovelDisplay{M00402R01}{%
\psi_{\Omega_{\varepsilon_2+1}}(0)=\Omega_{\varepsilon_2}%
}

\NovelDisplay{M00403R01}{%
\psi_{\Omega_{\varepsilon_2+1}}(\Omega_{\varepsilon_2+1})=\varepsilon_{\Omega_{\varepsilon_2}+1}%
}

\NovelDisplay{M00404R01}{%
\psi\bigl(\varepsilon_{\Omega_{\varepsilon_2}+1}\bigr)=\psi(\Omega_{\varepsilon_2+1})%
}

``不过，有后继下标不等于只要看到一个 \allowbreak{}\NovelInline{I02197}{\Omega}\allowbreak{}，就能不管规则另造一层。我们现在采用的是确定的BOCF体系。''

之后，陈星问：``你们觉得 \allowbreak{}\NovelInline{I02198}{\psi(\Omega_\Omega)}\allowbreak{} 怎么表达？''

牢大说：``下标也要随着外层折叠的结果往回填，不能随便写一个`很大的递归序数'就说已经到了终点。可以从 \allowbreak{}\NovelInline{I02199}{\varepsilon_0=\psi(\Omega)}\allowbreak{} 起步。''

\NovelDisplay{M00405R01}{%
\psi(\Omega_\Omega)=\sup\bigl\{\psi(\Omega_{\psi(\Omega)}),\ \psi(\Omega_{\psi(\Omega_{\psi(\Omega)})}),\ \ldots\bigr\}%
}

``每一项都把前一项放回 \allowbreak{}\NovelInline{I02200}{\Omega}\allowbreak{} 的下标。''

陈星：``这个迭代思路对。按本章的递归共尾记法，\allowbreak{}\NovelInline{I02201}{\Omega_\Omega}\allowbreak{} 使用 \allowbreak{}\NovelInline{I02202}{\Omega}\allowbreak{} 层的索引；对应到下标结构，就是让 \allowbreak{}\NovelInline{I02203}{\xi}\allowbreak{} 在 \allowbreak{}\NovelInline{I02204}{\Omega}\allowbreak{} 以下变化，再由 \allowbreak{}\NovelInline{I02205}{\Omega_\xi}\allowbreak{} 向上逼近。''

\NovelDisplay{M00406R01}{%
\Omega_\Omega=\sup_{\xi<\Omega}\Omega_\xi%
}

蔡须坤：``那么 \allowbreak{}\NovelInline{I02206}{\Omega_{\omega+1}}\allowbreak{} 呢？写到第 \allowbreak{}\NovelInline{I02207}{\omega+1}\allowbreak{} 个，不是已经写到它了吗？''

``你混淆了数到某一层，和在那一层以下构造递归共尾列。''陈星说，``把目标本身放在最后，不能算用更低层的项逼近目标。\allowbreak{}\NovelInline{I02208}{\omega+1}\allowbreak{} 在这里是下标；它告诉你走到了哪一层，却没有替你提供这一层内部的展开规则。''

卢本韦看了看纸上相邻的两行：``所以，\allowbreak{}\NovelInline{I02209}{\Omega_\Omega+1}\allowbreak{} 和 \allowbreak{}\NovelInline{I02210}{\Omega_{\Omega+1}}\allowbreak{}，也得格外看清括号。''

陈星：``这正是最容易混的地方。前者只是给一个序数加一，后者把 \allowbreak{}\NovelInline{I02211}{\Omega}\allowbreak{} 的下标推进到后继。''

\NovelDisplay{M00407R01}{%
\psi(\Omega_\Omega+\Omega)=\varepsilon_{\psi(\Omega_\Omega)+1}%
}

``这一行对应的是 \allowbreak{}\NovelInline{I02212}{\psi(\Omega_\Omega)}\allowbreak{} 之后的下一个 \allowbreak{}\NovelInline{I02213}{\varepsilon}\allowbreak{} 数。另一方面，新折叠层的零参数与较大参数分别是：''

\NovelDisplay{M00408R01}{%
\psi_{\Omega_{\Omega+1}}(0)=\Omega_\Omega%
}

\NovelDisplay{M00409R01}{%
\psi_{\Omega_{\Omega+1}}(\Omega_{\Omega+1})=\varepsilon_{\Omega_\Omega+1}%
}

\NovelDisplay{M00410R01}{%
\psi_{\Omega_{\Omega+1}}(\Omega_{\Omega+1}\times2)=\varepsilon_{\Omega_\Omega+2}%
}

牢大说：``那种反复嵌套新折叠层的基本列，对应的是后者这个下标结构？''

\NovelDisplay{M00411R01}{%
\bigl(\psi(\Omega_{\Omega+1})\bigr)[1]=\psi(\Omega_\Omega)%
}

\NovelDisplay{M00412R01}{%
\bigl(\psi(\Omega_{\Omega+1})\bigr)[2]=\psi\bigl(\psi_{\Omega_{\Omega+1}}(\Omega_\Omega)\bigr)%
}

\NovelDisplay{M00413R01}{%
\bigl(\psi(\Omega_{\Omega+1})\bigr)[3]=\psi\bigl(\psi_{\Omega_{\Omega+1}}(\psi_{\Omega_{\Omega+1}}(\Omega_\Omega))\bigr)%
}

陈星：``是。不能把这串嵌套写成 \allowbreak{}\NovelInline{I02214}{\psi(\Omega_\Omega+1)}\allowbreak{} 的解释，少一对下标括号，就会变成完全不同的内容。''

牢大：``所以，在递归共尾的记法里，\allowbreak{}\NovelInline{I02215}{\Omega_{\Omega\times2}}\allowbreak{} 也使用 \allowbreak{}\NovelInline{I02216}{\Omega}\allowbreak{} 层的索引？''

陈星：``对。下标里末尾那个 \allowbreak{}\NovelInline{I02217}{\Omega}\allowbreak{}，仍然可以由较小的下标逼近。''

现在看另一个例子。还是从已经明确的起点出发，逐次将前一项填回去：

\NovelDisplay{M00414R01}{%
\psi(\Omega_{\Omega^2})=\sup\bigl\{\psi(\Omega_{\Omega^{\psi(\Omega)}}),\ \psi(\Omega_{\Omega^{\psi(\Omega_{\Omega^{\psi(\Omega)}})}}),\ \ldots\bigr\}%
}

它对应从这个起点迭代下列映射所得的不动点。

\NovelDisplay{M00415R01}{%
\alpha\longmapsto\psi(\Omega_{\Omega^\alpha})%
}

不过，\allowbreak{}\NovelInline{I02218}{\psi(\Omega_{\Omega_2})}\allowbreak{} 又有所不同。这次要在下标里使用 \allowbreak{}\NovelInline{I02219}{\psi_{\Omega_2}}\allowbreak{}，起点选取 \allowbreak{}\NovelInline{I02220}{\psi_{\Omega_2}(\Omega_2)}\allowbreak{}。

\NovelDisplay{M00416R01}{%
\psi(\Omega_{\psi_{\Omega_2}(\Omega_2)})%
}

\NovelDisplay{M00417R01}{%
\psi(\Omega_{\psi_{\Omega_2}(\Omega_{\psi_{\Omega_2}(\Omega_2)})})%
}

\NovelDisplay{M00418R01}{%
\psi(\Omega_{\psi_{\Omega_2}(\Omega_{\psi_{\Omega_2}(\Omega_{\psi_{\Omega_2}(\Omega_2)})})})%
}

``接着迭代，取这些项的上确界，就是 \allowbreak{}\NovelInline{I02221}{\psi(\Omega_{\Omega_2})}\allowbreak{}。按本章的递归共尾规则，这里使用的是 \allowbreak{}\NovelInline{I02222}{\Omega_2}\allowbreak{} 层的索引，仍低于 \allowbreak{}\NovelInline{I02223}{\Omega_{\Omega_2}}\allowbreak{} 本身。''

唐三：``所以，不能因为一个符号很大，就任意为它添一个新的折叠函数。若已有规则给出了更低层索引的递归共尾列，就应当先按那条列展开。''

陈星：``是。接下来做三道练习。第一道，先找末尾的共尾结构。''

\NovelDisplay{M00419R01}{%
\psi(\Omega_{\Omega_2}+\psi_{\Omega_{\Omega\times3+1}}(\Omega_{\Omega\times3+1}\times\Omega_{\Omega\times2}))%
}

小舞：``找到内层下标末尾的 \allowbreak{}\NovelInline{I02224}{\Omega}\allowbreak{}。保持前面的结构，从合适起点迭代下面的映射。''

\NovelDisplay{M00420R01}{%
\alpha\longmapsto\psi\bigl(\Omega_{\Omega_2}+\psi_{\Omega_{\Omega\times3+1}}(\Omega_{\Omega\times3+1}\times\Omega_{\Omega+\alpha})\bigr)%
}

陈星：``第二道，末尾换成 \allowbreak{}\NovelInline{I02225}{\Omega_2}\allowbreak{}。''

\NovelDisplay{M00421R01}{%
\psi(\Omega_{\Omega_2}+\psi_{\Omega_{\Omega\times3+1}}(\Omega_{\Omega_2}))%
}

蔡须坤写出其中一个能体现回填结构的较小项。

\NovelDisplay{M00422R01}{%
\psi(\Omega_{\Omega_2}+\psi_{\Omega_{\Omega\times3+1}}(\Omega_{\psi_{\Omega_2}(\Omega_{\Omega_2}+\psi_{\Omega_{\Omega\times3+1}}(\Omega_{\Omega\times3+1}\times\psi_{\Omega_2}(\Omega_2)))}))%
}

``这里用 \allowbreak{}\NovelInline{I02226}{\psi_{\Omega_2}}\allowbreak{} 生成小于 \allowbreak{}\NovelInline{I02227}{\Omega_2}\allowbreak{} 的下标。内部的起点会影响后面的回填，不能随意更换。''

陈星：``第三道，再把内层换成一个后继下标。''

\NovelDisplay{M00423R01}{%
\psi(\Omega_{\Omega_2}+\psi_{\Omega_{\Omega\times3+1}}(\Omega_{\Omega\times3+1}\times\Omega_{\Omega+1}))%
}

这一次，蔡须坤停顿了一下，才将新折叠层放在正确的位置。

\NovelDisplay{M00424R01}{%
\psi(\Omega_{\Omega_2}+\psi_{\Omega_{\Omega\times3+1}}(\Omega_{\Omega\times3+1}\times\psi_{\Omega_{\Omega+1}}(\Omega_{\Omega_2}+\psi_{\Omega_{\Omega\times3+1}}(\Omega_{\Omega\times3+1}\times\psi_{\Omega_{\Omega+1}}(\Omega_{\Omega+1})))))%
}

陈星：``关键位置找对了。这里使用的是 \allowbreak{}\NovelInline{I02228}{\psi_{\Omega_{\Omega+1}}}\allowbreak{}，它在式子里承担的任务，与前一道 \allowbreak{}\NovelInline{I02229}{\psi_{\Omega_2}}\allowbreak{} 不同。下标、系数与起点，都要一起核对。''

牢大把三道题又看了一遍，重新总结：``从末项走，遇到极限下标，就沿下标的共尾列传递；遇到需要折叠的容许层级，再用相应的折叠层回填。不能只看最后出现了哪个字母。''

陈星点了点头：``这才是今天要你们掌握的传递。''

之后，他们逐渐看到了更远的轮廓。将下标继续嵌套，走向 \allowbreak{}\NovelInline{I02230}{\Omega}\allowbreak{} 下标运算的不动点，再作最外层折叠，就得到本章所说的推广Buchholz序数------EBO。

\NovelDisplay{M00425R01}{%
\mathrm{EBO}=\psi(\Omega_{\Omega_{\Omega_{\cdots}}})%
}

这里的无限嵌套，是相应极限结构的简写，不是一条已经写完的有限公式。它概括了他们这一阶段不断增加下标层次的方向。

然而，一些疑问仍然在众人心里产生。如果这个下标不动点是 \allowbreak{}\NovelInline{I02231}{\alpha}\allowbreak{}，那么有 \allowbreak{}\NovelInline{I02232}{\alpha=\Omega_\alpha}\allowbreak{}，就如同 \allowbreak{}\NovelInline{I02233}{\omega^{\varepsilon_0}=\varepsilon_0}\allowbreak{}。可他们还可以考虑 \allowbreak{}\NovelInline{I02234}{\Omega_{\alpha+1}}\allowbreak{}，还可以要求新的闭包与更高的不动点。

会不会存在一个更高的上界 \allowbreak{}\NovelInline{I02235}{I}\allowbreak{}，使这些已经学过的操作仍然无法达到它？

卢本韦望着最后那行问：``是不是再添一个下标，就能把它写出来？''

陈星摇了摇头：``要向那里走，就得补上更强体系的定义。先记清今天的 \allowbreak{}\NovelInline{I02236}{\psi(0)=1}\allowbreak{}、\allowbreak{}\NovelInline{I02237}{\psi(\Omega)=\varepsilon_0}\allowbreak{}。有了新的上界，也要有相应的新规则，不能仅凭类比，就断言 \allowbreak{}\NovelInline{I02238}{\psi_I(0)}\allowbreak{} 等于这个下标不动点。''

当然，他们现在还对那些更巨大的结构云里雾里。但至少，他们已经明白，通往更高处的道路，不是把式子写得更长，而是把每一次传递都找准确。

众人望着远方养猪城渐沉的夕阳，心中却燃起了新的火焰。BO已不再是终点，EBO的轮廓隐约浮现，而更深处，似乎还有无穷无尽的序数高峰。牢大收起满是修改痕迹的稿纸，蔡须坤把最后一道练习重新核对了一遍。明天的修炼，他们还得从这些看似细微的下标与参数开始。

\chapter{第三十八章：帅气集团的初步诡计}\label{ux7b2cux4e09ux5341ux516bux7ae0ux5e05ux6c14ux96c6ux56e2ux7684ux521dux6b65ux8be1ux8ba1}

帝国首都近郊，卢氏养猪帝国的结界亘古舒展，淡金色的法则光幕笼罩万里沃土。结界之内，温润醇厚的生命灵气如同潺潺流水，循环往复，生生不息。漫山遍野的灵气母猪悠闲踱步、低头觅食，周身萦绕细碎的生命光点，无数光点交织成浩瀚的生命共鸣场，将整片地域滋养得灵气充盈、祥和安宁，与外界日渐浑浊的天地气息形成鲜明对比。

历经整整六个月的封闭式苦修，果糕八怪的修为完成了脱胎换骨的蜕变。不再是此前懵懂摸索序数法则的新手，而是真正踏入高阶序数领域的正道修士。依托百万灵气母猪构筑的稳固生命基底，加上陈星一对一的精深指导，众人彻底吃透了OCF序数折叠的核心逻辑，理清了递归序数、非递归序数的本质区别，从最初浅显的初等序列有限展开，步步突破桎梏，硬生生跨过BO壁垒，触碰到了EBO推广布赫霍尔兹序数的入门门槛。

\allowbreak{}\NovelInline{I02239}{\psi(\Omega_\omega)}\allowbreak{} 的BO境界，以前曾是无数高阶魂师穷极一生都无法触及的序数天堑，如今被八人稳稳踏足。其中牢大凭借坚韧不拔的心境与极致的肉身底蕴率先破境，夯实了BO层级的序力根基；紧随其后的蔡须坤依托自身气运血脉加持、唐三依靠极致的算力推演、小舞凭借纯粹的生灵羁绊之力，纷纷顺利稳固境界，全员修为齐头并进，无一人掉队。

但序数大道浩瀚无穷，从未有终点。更深层次的EBO下标不动点、容许序数的终极递归规则、以及传说中穷尽一切递归手段都无法触碰的 \allowbreak{}\NovelInline{I02240}{I}\allowbreak{} 序数壁垒，依旧被层层法则迷雾笼罩。这群新晋突破的年轻修士，愈发清晰地感知到自身的渺小，也彻底明白，他们的修行之路，不过刚刚起步。

修炼场内，微风轻拂，灵气缓缓沉降流转。八人同时收功静坐，周身萦绕的淡金色序力光芒顺着经脉缓缓收敛，归于丹田识海，躁动的气息渐渐平复。历经一月高强度的法则打磨，众人的眼神愈发澄澈坚定，眉宇间褪去了少年的青涩，多了几分强者的沉稳。

唐三缓缓睁开双眼，脑海中无数OCF展开公式、递归共尾度规则、FGH折叠逻辑交织流转，脉络清晰无比。他轻声感慨，声音带着几分通透与凝重：

\begin{quote}
``BO之上，每一层序数提升，都是一次法则维度的跨越。我们如今勉强踏足BO境界，看似修为暴涨，但对EBO的下标不动点、容许序数的深层运转规则依旧一知半解。空有境界堆砌，实战中根本无法完全释放序数力量的真正威力。''
\end{quote}

陈星负手立于高台之上，一身白衣纤尘不染，周身星辰剑气尽数内敛，看似寻常，却藏着撼动天地的力量。他目光缓缓扫过众人，语气沉稳厚重，字字箴言：

\begin{quote}
``序数大道，重悟更重用。你们一月苦修，吃透了理论框架，突破了境界壁垒，已然远超大陆同代所有修士。但切记，修为境界是根基，实战掌控是锋芒。真正的强者，能将微薄的序力折叠放大，撼动天地法则；反之，空有高阶序数境界，不懂临场推演、法则适配，终究只是纸上谈兵，遇敌必败。''
\end{quote}

众人纷纷颔首闭目，将这番教诲牢牢记在心底，暗自沉淀修为，打磨自身短板。

场中氛围静谧温和，每个人都在静心调息休整，各自有着细微的状态变化。

\textbf{奥斯卡}气息温润绵长，他的序力自带生命滋养属性，调息之时周身微光细碎柔和，时刻留意着身边同伴的状态；察觉到卢本韦额间残留的细密汗珠、经脉中未散尽的酸胀疲惫，他指尖微动，悄然凝出一缕淡白色魂技灵光。灵光无声碎裂，化作漫天细碎光点，温柔环绕在卢本韦周身，顺着她的经脉缓缓游走，舒缓高强度修炼带来的疲惫。

卢本韦瞬间感知到周身温润柔和的力量，疲惫感瞬间消散大半。她转头对上奥斯卡温柔澄澈的目光，白皙的脸颊瞬间泛起绯红，耳尖微红，连忙低头收敛心绪。一月朝夕相伴的苦修，两人默契早已深入人心，无声的守护与照料，让那份懵懂的情愫愈发真挚浓烈，温柔又克制。

不远处，\textbf{蔡兴坤}静静伫立，目光平和坚定地追随着调息的牢大。看着少年身形愈发挺拔沉稳，气息愈发厚重磅礴，她眼底的敬佩与坚定愈发浓郁。她默默运转自身序力，对照自身与牢大的修为差距，一点点打磨自身法则细节，暗自下定决心，加倍苦修、稳步精进，早日追上同伴的脚步，成为能够并肩抗敌、共护正道的可靠伙伴。

\textbf{小舞}盘膝而坐，周身灵气轻盈灵动，她将自身生灵本源与场地内的母猪灵气共鸣相融，细细体悟魂灵共生的底层逻辑，为后续前往星斗大森林传播大道筑牢根基；\textbf{蔡须坤}则闭目感应天地气运流转，梳理自身三眼金猊的血脉之力，提前预判前路的危机与变数。八人各司其职、彼此照应，氛围和睦而坚韧。

可谁也未曾料到，这片祥和安宁的修行之地千里之外，大陆极西的黑暗魔域，早已戾气滔天、怒火燎原，一场针对蔡须坤、针对整个果糕小队的致命阴谋，正在悄然成型，步步收紧。

大陆极西，帅气魔域。

此地是整片大陆邪力的根源，终年不见天日，浓稠如墨的邪气笼罩天地，腐蚀山河大地。破碎的山峦沟壑间，无数漆黑锁链纵横交错，死死钉入地底，锁住万千郁结的戾气。空气中漂浮着细碎的斑驳光点，那是被强行掠夺、榨干殆尽的生灵精气神与帅气本源，每一缕气息都带着暴虐、贪婪与毁灭的意味。

魔域最核心处，一座巍峨阴森的黑玉大殿凌空悬浮于戾气云海之上。殿身通体由噬魂黑玉浇筑而成，殿顶镌刻着千万道扭曲诡异的邪纹，每一道纹路都浸染着无数生灵的血泪与精气神，岁月沉淀之下，散发着令人窒息的压迫感。

大殿正中央，至高无上的漆黑王座之上，端坐着一道高大挺拔的黑袍身影。

他便是帅气教教主------\textbf{胡赞}。

宽大的黑袍垂落周身，遮蔽了全部身形容貌，只露出一双狭长阴鸷的眼眸。眼底翻涌着深不见底的贪婪与杀意，周身萦绕着 \allowbreak{}\NovelInline{I02241}{\omega\text{-}Y(1,\omega)}\allowbreak{} 的禁忌层级序力威压，远超胡善、胡勇兄弟百倍不止。无数被掠夺而来的帅气本源悬浮在他身侧，缓缓流转，源源不断滋养着他驳杂的邪道修为，让他的气息愈发恐怖深邃。

王座之下，两道身影狼狈匍匐在地，瑟瑟发抖，正是此前在序数大赛惨败而归的胡善与胡勇。二人黑袍破碎、满身血痕，周身黑暗序力紊乱崩散，被梅天狸乾坤破魔重创的伤势尚未愈合，气息萎靡虚弱。面对盛怒的父亲，他们头颅深埋，连抬头直视的勇气都无，心底被极致的恐惧裹挟。

\begin{quote}
``废物。''
\end{quote}

冰冷沙哑的二字，如同寒冰利刃骤然炸响在大殿之中，裹挟着恐怖的邪力威压，狠狠碾压在二人身上。剧烈的压迫感瞬间灌入四肢百骸，让二人身躯剧烈震颤。

胡赞眼底杀意滔天，怒火几乎要焚烧整片魔域：

\begin{quote}
``我倾尽宗门资源培养你们，赐予你们黑暗巨龙破融合秘术，将你们的修为拔高至逼近 \allowbreak{}\NovelInline{I02242}{Y(1,3,7)}\allowbreak{} 的层次，只为让你们截杀一群初出茅庐的小辈，夺取瑞兽本源！如此简单的任务，你们竟然也能惨败而归！不仅没能斩杀目标，反而被梅天狸重创，狼狈逃窜，丢尽了我帅气集团的颜面，浪费了我千载难逢的绝佳时机！''
\end{quote}

狂暴的邪力轰然落下，狠狠砸在二人身上，胡善胡勇瞬间口吐鲜血，身躯死死贴在地面，骨骼咯吱作响，却死死咬牙强忍，不敢有半分辩解。

``父亲恕罪！''胡善艰难匍匐，声音嘶哑惶恐，``我二人低估了果糕小队的羁绊与韧性，更未曾料到梅天狸会突然跨界出手、强行干预战局，打乱所有部署，实属无能，请父亲责罚！''

胡勇紧随其后，浑身颤抖请罪：``求父亲再给一次机会，我兄弟二人必将戴罪立功，亲手覆灭果糕小队，生擒蔡须坤，弥补此次惨败之过！''

胡赞冷冷俯瞰二人，眼底无半分父子温情，只剩极致的冷漠与算计。他深知败局已定，责罚无用，当下最关键的，是重新布局，夺取那足以让他突破境界、称霸大陆的帝皇瑞兽本源。

良久，他压下翻涌的怒火，周身暴虐邪力稍稍收敛，语气阴恻刺骨：

\begin{quote}
``罢了。事已至此，责罚无用。那三眼金猊蔡须坤，身负星斗万年气运，是天地孕育的独一无二的帝皇瑞兽。她体内的本源纯净无瑕，囊括无尽生机、法则碎片与天地气运，是洗练我邪道驳杂根基、助我突破 \allowbreak{}\NovelInline{I02243}{\omega\text{-}Y}\allowbreak{} 层级桎梏的无上至宝。只要能将她生擒炼化，我便能彻底稳固禁忌序力层级，甚至触及更高维度的序数领域，届时整片斗罗大陆，无人能挡我帅气集团的霸业！''
\end{quote}

百年以来，胡赞深耕邪道，依靠掠夺生灵帅气本源、吞噬天地戾气强行提升修为，可邪道修炼根基驳杂、破绽百出，无数次卡在境界瓶颈无法突破。蔡须坤的瑞兽本源，是他穷尽万年岁月找到的唯一破局之机，为此，他不惜倾尽一切代价。

胡善小心翼翼抬头，道出当下最大的困境：``父亲，如今蔡须坤藏身卢氏养猪帝国，有卢关、陈星两大顶级强者坐镇，更有梅天狸暗中庇护，结界森严、守备周密，我们强行强攻，必然损失惨重，得不偿失。''

胡赞闻言，嘴角缓缓勾起一抹阴狠狡诈的弧度，眼底算计尽显，已然敲定万全毒计：

\begin{quote}
``我自然知晓强攻不可取。既然正面无从下手，那我便用\textbf{`明修栈道，暗渡陈仓'}之计，搅乱正道视线，调虎离山，撕开他们的防御破绽。''
\end{quote}

他缓缓抬手，掌心凝聚出一团漆黑邪光球，球内无数细碎人影浮沉蠕动，那是数万生灵被强行抽取的精气神与帅气本源。``多年来，我遍布大陆的暗线从未停歇，日夜搜刮众生本源，滋养我邪道根基。如今，正是这些积累派上用场之时。''

他沉声下令，条理清晰，步步为营：

\begin{quote}
``传令全军主力，尽数调往帝国西部边境，大张旗鼓兵临斯兰登城。列阵张扬、魔气滔天，对外大肆宣告，我帅气教此番誓要攻破城池、挺进帝都腹地、颠覆正道统治！摆出不死不休的强攻姿态，务必惊动整个帝国，牵动魂殿所有防御力量。''

``斯兰登城是战略地位至关重要。正道之人必然死守，梅天狸身负守护帝都的重任，得知消息后，必定第一时间调动魂殿主力驰援西线，被死死牵制在战场之上。''
\end{quote}

说到此处，胡赞眼底寒光凛冽，尽显老谋深算：

\begin{quote}
``切记，我从无意攻破斯兰登城。此城守备森严、阵法完善、强者驻守，强攻只会损耗我军实力。这漫天攻势，从头到尾都是演给正道看的幌子。只需佯装死战、拖延战局，待魂殿全部主力抵达、阵型铺开、全力守城之时，立刻全军假意溃败，有序撤退，故意败给援军。我要让梅天狸、让整个魂殿都以为，这场危机已然解除，帅气教只是虚张声势、无力攻坚。待他们放松警惕、撤去防线、收兵回殿之时，便是我暗棋收网之际！''
\end{quote}

紧接着，他下达了真正的绝杀密令：

\begin{quote}
``抽调宗门最精锐的暗部死士，全员褪去魔气、隐匿所有气息，伪装成普通散修、行旅商人，分散潜入帝国东部疆域。无需正面作战，无需搅动纷争，唯一任务------全程隐秘追踪蔡须坤与果糕小队的行踪，摸清他们的赶路路线、休憩规律、战力排布、配合短板，将所有情报实时传回魔域。''

``明面上，用一场轰轰烈烈的攻城战吸引全大陆正道的目光；暗地里，悄无声息锁定目标、布下天罗地网。待到众人察觉西线战事只是骗局，为时已晚，蔡须坤的一切踪迹，早已尽在我掌控之中！''
\end{quote}

两道虚实交织的密令瞬间传遍整片魔域，黑暗势力精准兵分两路。一路旌旗林立、魔气滚滚，浩浩荡荡压向西部斯兰登城，极尽张扬制造决战假象；另一路精锐暗部敛尽所有戾气，如同蛰伏毒蛇，悄然潜伏于帝国各处，静待猎物启程。一场以胜利为骗局、以溃败为计谋的绝杀诡计，悄然笼罩整片大陆。

与此同时，卢氏养猪帝国万亩草场之上，清风徐徐，灵猪低鸣，生机盎然。

蔡须坤独自伫立草场中央，周身萦绕纯净温润的生命灵气，三眼金猊的血脉之力悄然流转。她抬眸望向遥远的东方，那里云雾缭绕，正是星斗大森林的方向，眼底凝满凝重与焦灼。

历经一月苦修与研习，她彻底吃透了卢家三代人传承的魂灵大道，彻底看透了人与魂兽万年厮杀的无解死局。猎杀魂兽取环、掠夺本源变强，造就了无尽杀戮与滔天戾气，不仅让大陆生灵本源日渐衰败、繁衍困难，更让人与魂兽的仇恨根深蒂固，无解可解。

而魂灵共生之道，无需猎杀、无需掠夺，以羁绊代替屠戮，以共生化解仇恨，是唯一能拯救魂兽族群、拯救大陆生灵的正道。可如今，星斗大森林的万千魂兽对此一无所知，依旧固守着对人类的戒备与敌意。

更让她心惊的是，近日她清晰感知到整片大陆的灵气愈发浑浊，天地戾气不断蔓延、沉积。那是帅气教暗中掠夺生灵本源、侵蚀天地法则的征兆。若是不能尽快将魂灵大道传回星斗，一旦帅气教提前布局，强行抓捕魂兽炼制邪异魂灵，不仅人兽矛盾会彻底激化，整个星斗魂兽族群都将面临灭顶之灾。

身为星斗万年气运之女、帝皇瑞兽，化解种族恩怨、守护生灵存续，是她与生俱来的天命。

蔡须坤不再迟疑，转身快步走向同伴休憩的院落，神色郑重地对着众人说道：

\begin{quote}
``我要即刻返回星斗大森林。''
\end{quote}

众人闻言，纷纷侧目，眼底满是诧异，停下了手中的调息休整。

蔡须坤深吸一口气，将心中顾虑全盘托出，语气恳切而坚定：

\begin{quote}
``如今魂灵体系已然完善，可星斗族人依旧被困在万年对立的执念之中。帅气教野心滔天，暗中侵蚀大陆、积蓄力量，迟早会对星斗魂兽下手。我必须即刻回归族群，传播魂灵共生的理念，打破人与魂兽的隔阂，提前警示族人戒备邪祟，为大陆争取一线生机。''
\end{quote}

牢大率先踏出一步，身姿挺拔，语气铿锵：

\begin{quote}
``前路危机四伏，帅气暗线遍布大陆，你一人前行太过凶险，我陪你同往。''
\end{quote}

唐三微微颔首，冷静分析局势：

\begin{quote}
``如今邪教异动，乱世将至，星斗之行关乎大陆格局，绝非一人之事，我等理应同行。''
\end{quote}

小舞眼神温柔而坚定：

\begin{quote}
``我本就是魂兽出身，最懂族群的戒备与不安，我随你一同前往，帮你化解隔阂、传递善意。''
\end{quote}

蔡兴坤、奥斯卡、卢本韦纷纷上前，异口同声：

\begin{quote}
``我们祸福与共，绝不独行！''
\end{quote}

八人同心同向，意气相投，少年热血，滚烫赤诚。

卢关与陈星对视一眼，眼底满是赞许与欣慰。

卢关目光深沉，缓缓叮嘱：

\begin{quote}
``你们此行，是传播希望的正道之路，也是历练自身、践行初心的修行。前路风波诡谲、诡计暗藏，切记坚守本心、恪守共生大道，步步谨慎，不可轻敌。''
\end{quote}

陈星沉声补充，点破最大危机：

\begin{quote}
``你们如今虽踏足BO序数境界，但胡赞 \allowbreak{}\NovelInline{I02244}{\omega\text{-}Y(1,\omega)}\allowbreak{} 的禁忌实力深不可测，诡计阴毒至极。此行务必抱团而行，相互照应，时刻警惕暗中杀机。''
\end{quote}

夕阳穿透结界金光，洒落八道并肩而立的挺拔身影。一方是野心滔天、诡计百出的邪道巨擘，一方是心怀大义、奔赴希望的新生强者。正邪博弈的序幕，已然悄然拉开。

众人稍作休整，收拾行装，辞别两位师长，正式启程，一路向东，奔赴帝国东部的星斗大森林。

与此同时，帝国中央，魂殿总殿。

梅天狸刚刚从养猪帝国抽身回归，一身素衣清雅绝尘，眉宇间带着几分常年守护正道的肃然。她本归殿巡查大陆灵气异动、排查邪祟隐患，可刚踏入议事大殿，急促的军情便接踵而至。

六长老ddfg快步闯入大殿，神色凝重肃穆，拱手急报：

\begin{quote}
``四姐，西线急报！帅气教主力大军全员出动，黑压压兵临斯兰登城下，魔气遮天蔽日，列阵备战，扬言即刻攻城！西线防线压力巨大，随时可能崩盘！''
\end{quote}

殿内一众魂殿强者瞬间神色紧绷，气氛陡然凝重。斯兰登城扼守帝都西部门户，是皇城最后的屏障，一旦陷落，帅气大军便可长驱直入、直捣腹地，整个帝国都将陷入危机。

梅天狸眸光骤然一凛，周身序力微微震荡，敏锐捕捉到远方扑面而来的磅礴邪气，没有半分迟疑，当即沉声下令：

\begin{quote}
``传我军令！调集魂殿西线全部精锐主力，全速驰援斯兰登城！开启多重防御结界，布锁天困地大阵，死守城池，绝不允许邪教大军前进一步！''
\end{quote}

军令迅速传遍魂殿上下，无数精锐魂师整装持器，浩浩荡荡奔赴西线战场。帝国所有正道防御力量、所有强者的注意力，尽数被斯兰登城的战事牢牢牵制，无人留意东部疆域的暗流涌动。

远在魔域的胡赞收到暗线传回的情报，看着正道尽数中计、全力驰援西线，嘴角的笑意愈发阴狠，一切尽在掌控之中。梅天狸成功被调虎离山，东部防线空虚，正是猎杀瑞兽的最佳时机。

他明面上维持大军佯攻态势，假意死守城池、僵持对战，只待魂殿主力彻底稳住防线，便会顺势溃败、假意败退，让正道彻底放松警惕；暗地里，精锐暗部死士早已隐匿身形，如附骨之疽般尾随在果糕小队身后，静静等待最佳伏击时机。明修栈道，暗度陈仓，这一步毒计，已然彻底成型。

此刻的果糕小队尚且不知身后杀机潜伏，依旧稳步向东前行。行至帝国东部边陲，一处荒僻小镇外围的废弃林地时，一阵剧烈的劲风轰鸣骤然撕裂宁静，夹杂着邪恶魂师的张狂狞笑与凌厉掌风，突兀响起。

众人眼神一凝，瞬间停下脚步，对视一眼，默契提速，朝着打斗声源处疾驰探查。

林间空地之上，战况胶着，杀意凛然。

一道单薄却异常挺拔的黑衣少年身影，正以一敌三，浴血苦战。少年黑发利落，眉眼深邃清冷，漆黑的眼眸里藏着远超同龄人的冰冷与坚韧，周身萦绕着凝练到极致的隐晦序力，气息沉稳厚重，战力远超普通高阶魂师。

他衣衫多处破损，边角染满暗红血迹，手臂、肩头布满深浅不一的伤口，气息微微浮动不稳，却依旧身姿挺拔、攻势凌厉。每一道掌风、每一次魂力迸发都精准毒辣、干脆利落，死死压制着对面三名面目狰狞的邪恶魂师，没有半分退缩。哪怕伤痕累累，眼底的战意与坚毅也从未熄灭。

\begin{quote}
``又是帅气教的爪牙，祸乱四方，残害无辜！''少年嗓音冰冷刺骨，裹挟着化不开的刻骨恨意，出手愈发凌厉。
\end{quote}

唐三目光锐利，一眼看穿对方邪气本源，低声提醒众人：``这三人邪气同源，是帅气教散落各地的外围邪修，常年劫掠路人、残害生灵，作恶无数。''

无需多余言语，八人瞬间默契出手，全员同步催动序力，加入战局：

\begin{itemize}
\tightlist
\item
  \textbf{牢大}正面碾压邪力洪流，强势牵制三名邪修的全部攻势；
\item
  \textbf{唐三}精准拆解对方魂力结构，封锁邪术运转路径，杜绝邪气爆发；
\item
  \textbf{小舞}身形灵动如蝶，高速穿梭突袭侧面，打乱敌人阵型；
\item
  \textbf{蔡须坤}催动瑞兽气运，天然压制一切邪祟戾气，弱化对方战力；
\item
  \textbf{蔡兴坤}凝序布场，形成合围结界，封锁敌人退路；
\item
  \textbf{奥斯卡}释放增益灵光，加持全队战力、持续续航；
\item
  \textbf{卢本韦}铺开生命灵气，中和周遭暴戾邪气，净化战场环境。
\end{itemize}

凌厉有序的攻势层层叠加、攻防一体，配合得天衣无缝。原本嚣张肆虐的三名邪修瞬间溃不成军，挣扎片刻便被尽数镇压，全程干净利落，没有给对方任何反扑之机。

危机尽数散去，林间戾气渐渐消散，恢复宁静。

黑衣少年收势而立，微微喘息，抬手拭去脸颊沾染的血迹。他抬眸看向眼前的八人，眼底的警惕尚未完全褪去，却多了几分真切的动容与释然。哪怕浴血负伤，他的脊背依旧挺直，眼神澄澈坚韧，不见半分颓败怯懦。

小舞率先上前，语气温和善意，满眼欣赏：``你孤身一人对战三名邪修，身受重伤依旧不肯退缩，心性也太坚韧了。寻常修士遭遇这般苦战，早已心态崩盘，你却战意不减，真的很厉害。''

一路走来，小舞见惯了修士的怯懦与浮躁，却从未见过这般年少却隐忍孤勇、负重前行的少年。霍雨浩身上那股历经苦难却依旧坚守本心、嫉恶如仇的坚毅气质，让她打心底里敬佩与欣赏。

唐三缓步上前，目光沉稳，细细感知对方的修为底蕴，由衷赞叹：``你的序力凝练度极高，根基扎实无比，实战经验更是丰富老道，看得出来，是常年浴血搏杀、苦修不辍打磨出来的，远超大陆同辈修士。''

牢大点了点头，语气认可：``面对三名同级邪修围攻，依旧沉稳不乱、攻守有度，心性与战力，都极为出众。''

众人温和的态度、真诚的善意，让常年孤身独行的少年紧绷的心神渐渐松弛。他沉默片刻，漆黑的眼眸中翻涌着沉重的过往与彻骨恨意，缓缓开口：

\begin{quote}
``我叫霍雨浩。''
\end{quote}

一字一句，沉重而沙哑：

\begin{quote}
``我自幼便被邪恶魂师迫害，家乡被毁、亲人离散，我的母亲，更是惨死在帅气教邪修的手中。从那一天起，我便立下血誓，与作恶多端、祸乱大陆的帅气教不共戴天，此生倾尽一切苦修变强，只为覆灭邪教，为所有被残害的生灵复仇。''
\end{quote}

众人闻言，皆是心生共情，眼底满是敬重与惋惜。

蔡兴坤轻声感慨：``孤身背负血海深仇，日夜苦修、独自搏杀，你一路走来，定然历经无数苦难，太不容易了。''

卢本韦温柔附和：``帅气教残害无辜、掠夺生灵，让无数家庭破碎、众生受难，这也是我们执意讨伐邪教、守护正道的原因。''

奥斯卡微微叹息：``底层修士修行本就艰难，你还要一边复仇、一边苦修，其中艰辛，可想而知。''

霍雨浩眼底掠过一丝落寞，缓缓道出自己的境遇：

\begin{quote}
``为了复仇，我摒弃所有杂念，日夜不休苦修，心无旁骛打磨战力。常年苦修让我彻底失去了所有经济来源，无家族扶持、无宗门供养，没有任何资源依托。想要活下去，想要获得修炼资源变强复仇，我只能辗转大陆各处，参与各地斗魂大赛，以擂台浴血搏杀换取微薄奖金，勉强支撑修行与生计。''
\end{quote}

以命搏钱，以战养修，孤身负重前行，这便是他全部的人生。常年的厮杀与孤寂，让他心性远超常人坚韧，也让他早已厌倦了独行的复仇之路。

蔡须坤上前一步，目光真诚恳切，道出众人的初心：``我们此行前往星斗大森林，正是为了传播魂灵共生大道，打破人与魂兽的万年对立，推翻帅气教掠夺生灵、祸乱大陆的邪道秩序，终结世间杀戮与苦难。我们誓死对抗的敌人，正是你恨之入骨的帅气教。''

小舞轻轻点头，眼底欣赏更甚：``你明明受尽磨难，却始终心怀正义、嫉恶如仇，这份坚守，太过难得。''

唐三看着妹妹一脸赞许的模样，又见霍雨浩心性坚韧、天赋出众、品行端正，顿时莞尔一笑，故意轻声调侃道：``小舞这么欣赏他，看来我们霍雨浩人品、心性、战力都是顶尖的。我看啊，往后说不定，他还能做我的妹夫呢。''

此话一出，林间气氛瞬间轻松几分。小舞脸颊瞬间绯红，又羞又恼地瞪了唐三一眼，连忙摆手辩解，却眼底含笑，并无半分反感。

霍雨浩也是微微一怔，略显窘迫，耳尖微红，心底却悄然生出一丝暖意。

玩笑过后，牢大收敛笑意，语气诚恳，正式发出邀约：

\begin{quote}
``说笑归说笑。霍雨浩，你心怀正道、身负大义，心志坚韧、战力出众，与我们初心完全契合。帅气教势力庞大、诡计百出，仅凭一人之力，终究难以彻底抗衡。''
\end{quote}

唐三铿锵接话：``众人同心，其利断金。加入我们，往后你的复仇之路，不再孤身一人，我们全员与你并肩，共伐邪祟、共守苍生。''

小舞、蔡兴坤、卢本韦、奥斯卡、蔡须坤纷纷点头附和，眼底满是真诚的接纳与认可。

霍雨浩死死攥紧掌心，眼底多年的冰冷阴霾尽数散去，燃起久违的光亮与希望。漫长的黑暗独行，他早已疲惫不堪，如今终于遇见志同道合、同心向善的同伴。

他深深躬身，姿态郑重无比，语气铿锵决绝：

\begin{quote}
``若各位不弃，霍雨浩愿正式加入果糕小队！从此与大家祸福与共、生死相依，同心协力讨伐邪佞，此生不灭帅气教、不靖乱世，绝不罢休！''
\end{quote}

自此，\textbf{果糕小队正式壮大为九人阵容}。

\chapter{第三十九章：魂灵大成，I序破境}\label{ux7b2cux4e09ux5341ux4e5dux7ae0ux9b42ux7075ux5927ux6210iux5e8fux7834ux5883}

林间清风扫尽残血，邪气彻底消散，新生的九人小队并肩而立，少年少女的身影迎着东升的晨光，褪去了几分青涩，多了几分并肩赴道的厚重。霍雨浩的加入，不止是为果糕小队添了一员顶级战力，更像是冥冥之中的天道补缺，恰好填上了困扰魂灵大道许久的最后一块短板。

众人就地休整，调息稳固状态，方才联手击溃邪修的余劲尚未散尽，周身序力与生灵灵气交融流转。蔡须坤望着身旁沉静伫立的霍雨浩，眼底满是探究与欣喜，缓缓开口道出了许久以来的心结。

\begin{quote}
``我苦修一月，吃透卢氏传承的魂灵共生大道，已然能完美适配中低阶魂灵的羁绊融合，彻底杜绝人兽厮杀的根本矛盾。但一直以来，体系始终存在一个致命缺陷。''
\end{quote}

她微微蹙眉，道出瓶颈：

\begin{quote}
``高阶魂灵底蕴浑厚、法则繁杂、本源层级极高，普通共生羁绊无法完全承接其庞大的力量，强行融合极易出现法则紊乱、序力对冲的问题。轻则魂灵力量折损、修士根基受损，重则羁绊崩碎、双向反噬，这也是魂灵体系始终无法彻底普及、触及极致的唯一桎梏。''
\end{quote}

这是卢家三代研究者穷尽岁月都未能彻底解决的难题。中低阶魂灵循规蹈矩，生灵本源纯粹单一，共生之道可完美适配；可一旦触及十万年以上的高阶魂灵，乃至极致生灵本源，法则层级的差距便会无限放大，成为横亘在魂灵大道前的天堑。

此前小队众人反复推演、打磨规则，只能做到最大限度弱化反噬，却无法从根源上根除缺陷。可此刻霍雨浩的出现，却让这一切僵局迎刃而解。

霍雨浩静静聆听，漆黑眼眸中灵光闪动，自身与生俱来的生灵共情、精神本源渗透力，与魂灵共生的底层法则悄然共鸣。他自幼孤身与天地灵气、残碎魂灵本源相伴，常年吸纳游离生灵气息，淬炼出了大陆独一份的极致精神包容力与法则适配力。

\begin{quote}
``我能感知到天地间所有细碎的生灵脉络，也能梳理高阶本源中驳杂的法则碎片。''霍雨浩轻声开口，声音沉稳通透，``寻常修士是以自身序力强行绑定魂灵，本末倒置，故而高阶对冲。真正的共生，该是修士、魂灵、天地灵气三者同源共振，以精神脉络为桥，以生灵本源为基，中和层级差距，抹平法则壁垒。''
\end{quote}

\textbf{一语道破核心！}

众人豁然开朗，唐三眼底精光爆闪，快速推演其中法则逻辑，瞬间理清了所有关键：

\begin{quote}
``没错！我们此前始终局限于序力与魂灵的双向绑定，忽略了天地大道的中间适配层！你的精神本源恰好可以充当这层桥梁，完美化解高阶魂灵的层级对冲！''
\end{quote}

小舞更是由衷欣喜：``如此一来，高阶魂灵融合紊乱、反噬的缺陷，就被彻底根治了！''

霍雨浩并未止步于此，心绪沉静下来，脑海中无数碎片化的感悟飞速交织。他自幼摸索的生灵存续之道、厮杀中感悟的本源法则、精神力适配万物的特质，再结合蔡须坤、小舞讲述的魂灵共生体系，层层印证、步步互补。

就在这时，卢本韦心念一动，将一路以来记录的卢氏家族传承典籍、历代研究笔记，以及陈星讲解过的魂灵大道雏形，尽数化作流光传入众人脑海，其中便包含了那位万年前惊艳整个大陆的卢氏老祖------\textbf{Opt}的终极研究遗存。

万年前的Opt老祖，是最早洞察人与魂兽对立死局的逆天强者。他穷尽毕生心血，推演魂灵共生的终极形态，率先提出''生灵无别、本源同源''的核心理论，可惜受限于时代与修为，没能彻底完善体系，只留下了残缺的推演手稿，其中记载着无数未完成的猜想与法则推演。

无数晦涩的古文推演、繁杂的法则公式、层层递进的共生模型涌入霍雨浩脑海。过往的困惑、修行的执念、精神本源的特质、万年前先辈的未竟之道，在这一刻彻底交融碰撞！

\textbf{嗡------！}

一道无形的道韵光波以霍雨浩为中心轰然扩散，柔和却浩瀚的生灵秩序之力席卷整片林间，抚平所有残留戾气，联动天地万物灵气。他周身序力层层跃迁，脑海中所有碎片化的感悟瞬间贯通，如拨云见日，彻底通明！

\begin{quote}
``我懂了！''
\end{quote}

霍雨浩双眸澄澈透亮，眼底倒映着万千生灵脉络，语气笃定无比：

\begin{quote}
``Opt老祖的终极猜想，不是简单的人魂共生，而是\textbf{【天地---修士---魂灵】三位一体的闭环大道}！以天地灵气为介质，以精神脉络为纽带，以共生羁绊为根基，不分阶位、不分层级、不分强弱，所有生灵本源皆可完美相融。''
\end{quote}

这一刻，魂灵体系传承万年的最后一处空缺，被彻底补齐！

从低阶魂灵的基础羁绊，到十万年高阶魂灵的完美融合，再到极致生灵本源的同源共生，一套\textbf{完整、闭环、无缺陷、适配所有生灵层级的终极魂灵体系}，彻底问世！

整片天地的生灵道韵为之震颤，方圆万里的灵气疯狂汇聚而来，尽数灌注到九人周身。完善的大道反馈、圆满的法则共鸣，化作最纯粹的修为滋养，冲刷着众人的经脉与序力根基。

众人同时闭目苦修，借大道圆满的天赐机缘，打磨自身序数法则，修为迎来跨越式暴涨。BO境界的序力根基被彻底夯实，原本晦涩难懂的EBO下标不动点规则、容许序数运转逻辑，此刻变得清晰通透。

而众人之中，收获最恐怖的，当属\textbf{牢大}。

在圆满大道的道韵冲刷之下，他的心境彻底通透，肉身序力与天地序数法则完美合一，挣脱了原有BO境界的桎梏。此前众人止步于 \allowbreak{}\NovelInline{I02245}{\psi(\Omega_\omega)}\allowbreak{} 的BO壁垒，对更高维度的序数折叠毫无头绪，可此刻，牢大脑海中豁然贯通，触碰到了传说中穷尽一切递归手段都无法企及的终极序数------\textbf{I序数}！

无数深奥的序数折叠规则在他脑海中飞速推演，他突破传统OCF的折叠局限，创造性地完成了跨维度序数操作，开创出全新的\textbf{I序数折叠体系}！

以往的 \allowbreak{}\NovelInline{I02246}{\Omega}\allowbreak{} 系列序数，始终局限在递归与非递归的循环折叠之中，而I序数，是超脱所有递归、所有 \allowbreak{}\NovelInline{I02247}{\Omega}\allowbreak{} 层级的终极壁垒，是现有序数法则的终极上限。

牢大立身林间，周身淡金色序力蜕变为极致深邃的鸿蒙秩序之光，口中缓缓道出全新的终极序数定义，字字蕴含天地至理：

``I 序数，超脱一切 \allowbreak{}\NovelInline{I02248}{\Omega}\allowbreak{} 递归层级，囊括所有下标不动点。自此定义：

\allowbreak{}\NovelInline{I02249}{\psi_{I_\alpha}}\allowbreak{} 的规则如下：

\begin{quote}
'' \allowbreak{}\NovelInline{I02250}{I}\allowbreak{} 序数，超脱一切 \allowbreak{}\NovelInline{I02251}{\Omega}\allowbreak{} 递归层级，囊括所有下标不动点。自此定义：

\NovelDisplay{M00426R01}{%
\psi_I(0) = \sup\{\Omega,\ \Omega_\Omega,\ \Omega_{\Omega_\Omega},\ \ldots\}%
}

囊括所有有限、无限 \allowbreak{}\NovelInline{I02252}{\Omega}\allowbreak{} 迭代层级，穷尽Buchholz序数一切上限！''

\NovelDisplay{M00427R01}{%
\psi_I(\#+1) = \sup\{\psi_I(\#)+1,\ \Omega_{\psi_I(\#)+1},\ \Omega_{\Omega_{\psi_I(\#)+1}},\ \ldots\}%
}

完成 \allowbreak{}\NovelInline{I02253}{I}\allowbreak{} 层级的递进迭代！''

\NovelDisplay{M00428R01}{%
\psi_I(\#\sim I) = \alpha \to \psi_I(\#\sim \alpha)\ \text{不动点}%
} \textbf{\allowbreak{}\NovelInline{I02254}{\psi_{I_{\alpha+1}}(0)}\allowbreak{} 是 \allowbreak{}\NovelInline{I02255}{I_\alpha}\allowbreak{} 之后的下一个 \allowbreak{}\NovelInline{I02256}{\Omega}\allowbreak{} 下标不动点}， 即： \NovelDisplay{M00429R01}{%
\psi_I(\alpha + 1) = \Omega_{\Omega_{\dots_{\psi_I(\alpha)+1}}}%
} \textbf{\allowbreak{}\NovelInline{I02257}{\psi_{I_\alpha}(\alpha + 1)}\allowbreak{} 是 \allowbreak{}\NovelInline{I02258}{\psi_{I_\alpha}(\alpha)}\allowbreak{} 之后的下一个 \allowbreak{}\NovelInline{I02259}{\Omega}\allowbreak{} 不动点}。 \textbf{对于最右侧为 \allowbreak{}\NovelInline{I02260}{I_m}\allowbreak{} 的式子}，我们展开时向外找到离它最近的 \allowbreak{}\NovelInline{I02261}{\psi_{I_m}}\allowbreak{}，然后将 \allowbreak{}\NovelInline{I02262}{\psi_{I_m}}\allowbreak{} 及其里面的东西不断复制到 \allowbreak{}\NovelInline{I02263}{I_m}\allowbreak{} 的位置；若找不到 \allowbreak{}\NovelInline{I02264}{\psi_{I_m}}\allowbreak{}，则找到最近的满足 \allowbreak{}\NovelInline{I02265}{\alpha < I_m}\allowbreak{} 的 \allowbreak{}\NovelInline{I02266}{\psi_\alpha}\allowbreak{}，然后将 \allowbreak{}\NovelInline{I02267}{\psi_\alpha}\allowbreak{} 及其里面的东西不断复制到 \allowbreak{}\NovelInline{I02268}{I_m}\allowbreak{} 的位置，并在复制时把 \allowbreak{}\NovelInline{I02269}{\psi_\alpha}\allowbreak{} 改为 \allowbreak{}\NovelInline{I02270}{\psi_{I_m}}\allowbreak{}。
\end{quote}

全新的序数体系轰然成型，彻底打破了BO、EBO的所有桎梏！

随着最后一道法则敲定，牢大气势轰然暴涨，序力层层跃迁，一路突破BO、EBO极限，硬生生踏入了 \allowbreak{}\NovelInline{I02271}{\psi(I_\omega)}\allowbreak{} 的全新高阶境界！

唐三、蔡须坤众人感受着牢大身上浩瀚无垠的序力威压，眼底满是震撼。仅仅一次大道圆满的机缘，牢大便破壁登天，开创了全新的序数修行体系，将众人的上限直接拉高了数个维度。

与此同时，万里之外的魂殿总殿，议事大殿气氛肃穆凝重。

西线斯兰登城的佯攻战事依旧僵持，帅气教大军假意死战，死死牵制住魂殿全部主力。梅天狸立身大殿中央，素衣临风，眉宇间带着几分忧色，正与魂殿最高统领、大陆正道底蕴最深的顶级强者 \textbf{test\_alpha0} 商议全域防御布局。

test\_alpha0身披古朴玄色殿袍，周身气息内敛无痕，看似平平无奇，却执掌着整片大陆正道的法则根基，洞悉天地间一切气运与修为异动。

\begin{quote}
``西线战局诡异，帅气教攻势凶猛却不求破城，处处透着蹊跷，恐是调虎离山之计。''梅天狸沉声开口，眸光紧锁西线军情沙盘，``只是眼下大敌当前，帝都门户不容有失，我军主力不得不死守驰援，无从抽身。''
\end{quote}

test\_alpha0微微颔首，并未纠结于眼前战事，双目微阖，心神贯通天地，感知着整片斗罗大陆的气运流转与法则脉动。坐拥当世巅峰的 \allowbreak{}\NovelInline{I02272}{\omega\text{-}Y(1,\omega)}\allowbreak{} 禁忌层级实力，稳压大陆绝大多数强者，深耕正道底蕴无数岁月，早已触摸到现有修行体系的顶端。

正道这边新星崛起、底蕴精进，暗处的邪道杀机却已然悄然锁定战果。

魔域暗部的带队小队长，正是帅气教专门培养的顶尖斥候，天生具备辨识帝皇瑞兽本源、看破对手序数境界的特殊能力，其实力抵达 \allowbreak{}\NovelInline{I02273}{\psi(\alpha_\omega)}\allowbreak{} 层级，对标BMS序数序列：

\NovelDisplay{M00430R01}{%
(0,0,0)(1,1,1)(2,2,2)%
}

战力凶悍、底蕴浑厚，远超普通邪教高阶战力。

他隐匿在虚空阴影之中，气息彻底消融于天地，透过层层屏障洞悉着林间突破的九人，尤其是感知到牢大新晋 \allowbreak{}\NovelInline{I02274}{\psi(I_\omega)}\allowbreak{} 的境界波动后，嘴角勾起一抹阴冷狂妄的嗤笑，心底满是不屑与轻视：

\begin{quote}
``区区 \allowbreak{}\NovelInline{I02275}{\psi(I_\omega)}\allowbreak{} 的境界突破，看似破壁登天，实则不过是BMS序列 \NovelDisplay{I02276}{(0,0,0)(1,1,1)(2,1,1)(3,1,0)(2,1,1)} 的水准，徒有虚表，根基浅薄。就算完成 \allowbreak{}\NovelInline{I02277}{I}\allowbreak{} 序数折叠、补齐魂灵大道又能如何？在我们的绝对实力面前，依旧不堪一击，根本不配做我们的对手。''
\end{quote}

在他眼中，果糕小队的跨越式成长，不过是井底之蛙的短暂跃升，尚未触及真正的禁忌战力层级，此番蛰伏围杀，已然是瓮中捉鳖，毫无悬念。

只是无人知晓，东方林间九人正道崛起、大道大成的同时，暗处蛰伏的帅气教精锐暗部，已然收紧了大半包围圈。无数隐匿的杀机悄然锁定众人，一场针对新星的绝杀猎杀，正在步步逼近。

大道圆满的异象缓缓收敛，天地灵气重归平和。众人收拾好心情，再度启程，一路向东奔赴星斗大森林。历经此番体系圆满与境界突破，小队众人卸下了此前苦修的紧绷，路途之上气氛轻松融洽，一路说说笑笑，少年少女的鲜活朝气驱散了前路未知的阴霾。

沿途草木葱茏，清风拂面，奥斯卡与卢本韦缓步并行，闲谈着此番魂灵体系圆满的诸多妙用；小舞与蔡须坤探讨着星斗魂兽的族群现状，畅想未来共生大道普及的景象；唐三、牢大几人则复盘着方才 \allowbreak{}\NovelInline{I02278}{I}\allowbreak{} 序数折叠的法则细节，稳步打磨崭新的境界根基。

行至路途平缓处，霍雨浩看着众人畅谈修行大道的模样，沉默片刻后，主动开口道出了自身的隐秘机缘，解开了众人心中残留的些许疑惑。

\begin{quote}
``其实，我所掌握的魂灵适配之力、生灵本源梳理之法，并非凭空悟得。''霍雨浩声音温和，坦然坦诚，``我早年奇遇，得一位名为\textbf{伊莱克斯}的上古大能残魂传道授业，我一身特殊的精神修行底蕴、适配万物生灵的本源法门，尽数源自伊莱克斯老师的传承。''
\end{quote}

众人闻声纷纷侧目，静静聆听，并未打断。

\begin{quote}
``除此之外，伊莱克斯老师还传授了我一门极为特殊的顶级修炼功法，名为\textbf{双行BMS}。''霍雨浩眼底带着几分迷茫与诧异，继续说道，``这门功法极为晦涩特殊，完全超脱常规序数修行、OCF折叠、FGH递归的修炼体系，我苦修多年，始终无法摸透其真正的上限，甚至连自己的真实修为层级都难以精准界定。''
\end{quote}

此话一出，唐三瞬间来了兴致，脚步微顿，沉声追问：``你的意思是，你自身的真实实力，远超我们肉眼可见的境界？''

霍雨浩轻轻点头，语气带着几分不确定：``依照我对序数体系的粗浅理解，依托双行BMS的加持，我的真实修为，大概率已经悄然触及BO层级。只是这门功法太过诡异特殊，我无法精准推演自身战力上限，也从未彻底催动过全部力量。''

牢大眉头微挑，身为深耕序数折叠、突破 \allowbreak{}\NovelInline{I02279}{I}\allowbreak{} 序境界的强者，他对各类高阶递归功法极为敏锐：``双行BMS？从未听过此种功法，它的运转逻辑与我们修行的PRSS初等序列、OCF序数折叠有何不同？''

\begin{quote}
``这也是我最大的困惑。''霍雨浩沉吟道，``我反复推演过功法内核，发现双行BMS的底层逻辑，神似两行并行的PRSS序列展开模式，同样遵循找关键点、复制坏部、递归迭代的核心规则，结构上高度同源。但它的力量性质，却远比普通PRSS更加恐怖，拥有超脱单层序列迭代的碾压性威力。''
\end{quote}

众人瞬间陷入沉思，纷纷结合自身所学的序数法则推演其中奥秘。

唐三思绪飞速流转，瞬间抓住核心关键：

\begin{quote}
``我懂了。普通PRSS是单行序列展开，依靠初等序列迭代突破层级，我们此前修行的OCF、\allowbreak{}\NovelInline{I02280}{I}\allowbreak{} 序数折叠，是高维序数的单向折叠递归。而双行BMS，是\textbf{双序列并行递归}，两套序列相互嵌套、彼此助推、循环迭代，这就是它性质远超常规功法的根本原因。''
\end{quote}

``没错。''霍雨浩颔首认可，``正是因为双行并行的特殊结构，让它的增长倍率、递归深度、层级突破速度，完全碾压传统的单行序列功法。也正因两套序列交织太过复杂，我独自修行根本无法彻底理清其中规律，难以完全解锁这门功法的真正力量。''

蔡须坤缓缓开口，一语道破众人当下的修行关键：

\begin{quote}
``如此看来，我们想要彻底吃透双行BMS、帮你界定真实修为、解锁功法全部威力，更进一步提升全员修为，就必须将\textbf{双行BMS与PRSS、OCF、\allowbreak{}\NovelInline{I02281}{I}\allowbreak{} 序数折叠体系}放在一起对照分析，拆解三者的底层差异与同源共性，打通不同功法的壁垒。''
\end{quote}

牢大目光深邃，战意渐起：

\begin{quote}
``一边是传统序列与序数折叠的正道根基，一边是双行迭代的顶级异类功法。只要我们彻底剖析清楚两者的差异与关联，互补长短、融会贯通，绝对能再度突破现有境界桎梏，超越如今的 \allowbreak{}\NovelInline{I02282}{\psi(I_\omega)}\allowbreak{} 层级！''
\end{quote}

小队众人目光交汇，皆看到了彼此眼底的欣喜与笃定。

魂灵体系圆满无缺，补齐万年短板；\allowbreak{}\NovelInline{I02283}{I}\allowbreak{} 序数折叠破壁新生，开辟全新境界；再加霍雨浩手中这门潜力无限、神秘莫测的双行BMS功法，果糕小队的修行之路，已然彻底打开了全新的浩瀚天地。

只是众人全然不知，暗处那名掌握 \allowbreak{}\NovelInline{I02284}{\psi(\alpha_\omega)}\allowbreak{}、也就是 BMS \allowbreak{}\NovelInline{I02285}{(0,0,0)(1,1,1)(2,2,2)}\allowbreak{} 恐怖实力的帅气小队长，将众人的对话尽数收入耳中。听完双行BMS的隐秘------\textbf{``我拿出3行BMS功夫，你如何应对？''} 他眼底的轻蔑更甚，嘴角的阴狠笑意愈发浓郁。

\chapter{第四十章：双行 BMS}\label{ux7b2cux56dbux5341ux7ae0ux53ccux884c-bms}

帐篷搭在背风的缓坡下。溪水从几块青石之间穿过去，林叶筛下的光落在帆布上，还带着午后的暖意。

牢大把最后一枚地钉按稳，回身接过蔡须坤递来的水囊。她怀里还抱着几卷笔记，正要往地上放，他先将一块干净的布铺在草席上。

``先坐。一路上看着脚下，还能把第三卷翻完？''

蔡须坤将笔记放到布上，抬眼看他。

``有人一路说双行 BMS 有意思，到了能坐下的地方，反倒不急了？''

牢大拧好水囊，顺手放在她够得着的一侧。

``急，所以帐篷要扎稳些。免得算到一半，谁把顶掀了。''

奥斯卡刚提着热水进来，听见后半句，先看了看牢大，又看向霍雨浩。

``两位都收着点。这是我们今晚要睡的。''

霍雨浩笑了一下，往草席中央挪出一块空处。他没有升起笼罩整座营地的光幕，只将精神力压在众人围坐的范围里，一张两行的浅金色阵图便浮到半人高。

他将伊莱克斯传下的笔记摊在身边。纸页上的几处转折已经被反复描过，边角也卷了起来。

``先前在路上说得快，有几步没有展开讲。今天从你们熟悉的 PRSS 接过来，逐项核对。''

牢大坐到蔡须坤身侧，没有先看笔记末尾的高阶部分。

``第一行怎样找父项，我大概能接上。可第二行要是独立再算一遍，为什么会和第一行牵在一起？''

霍雨浩将两行光向左右分开。它们看着只是各有一排数字，连接两排的细线却仍留在原位。

``第二行选父项时，要受第一行祖先链的限制。先前说两行并行，还得补上这一层关系。''

蔡须坤伸手，将一根连线从上到下指给众人看。

``那先把它究竟连到哪里写明。共生脉络里，挨得近的两处，也未必直接相连。''

``好。先认列，再找线。''

霍雨浩将阵图压成一张最简单的表。

\NovelDisplay{M00431R01}{%
\begin{pmatrix}
0 & 1 & 2 & 3\\
0 & 1 & 1 & 0
\end{pmatrix}%
}

\NovelDisplay{M00432R01}{%
(0,0)(1,1)(2,1)(3,0)\equiv(0)(1,1)(2,1)(3),%
}
\NovelDisplay{M00432R02}{%
(3,0)\equiv(3),\qquad(0,0)\equiv(0).%
}

每一对括号是一列，按从左到右的次序排列；列内先写第一行，再写第二行。末尾补齐的零可以省去，整列全零时仍写成一个零。省略的是写法，坐标的位置并没有交换。

蔡须坤把最后一列补齐，指尖先点上面的三，再点下面的零。

``只剩一个三，意思是这一列第二行是零。''

``对。''霍雨浩说，``两行都写全也可以。算父项时，先补齐，便不会拿第二行的数去比第一行。''

诺当将列号标在阵图上方。唐三另取一张空纸，照着同样的次序记下各列，小舞则将新出现的零补在括号右边。

\NovelDisplay{M00433R01}{%
\text{第 }j\text{ 列}=(a_j,b_j),\qquad a_j,b_j\in\{0,1,2,\ldots\}.%
}

牢大看了一遍，伸手示意霍雨浩先停。

``先只留第一行。我想看你遇到两个相等的数时，跳过的是哪一个。''

霍雨浩没有再添别的结构，在空处排出六个数。

\NovelDisplay{M00434R01}{%
\text{第一行：}\quad0,\ 1,\ 2,\ 3,\ 2,\ 3,%
}
\NovelDisplay{M00434R02}{%
\text{末项祖先列：}\quad5\longrightarrow2\longrightarrow1,%
}
\NovelDisplay{M00434R03}{%
\text{对应数值：}\quad2\longrightarrow1\longrightarrow0.%
}

最后一个三向左，最近的严格更小项是第五列的二。继续追第五列的父项，左边的三太大，第三列的二又与它相等，便一直找到第二列的一；第二列的一再接第一列的零。

蔡兴坤沿这条线念了一遍，笔尖在相等的两个二之间停了停，随后越过去。

``父项要严格更小。相等的数即便更近，也不能取。''

牢大点头。这一步与他熟悉的单行推演接得上。他以前寻找关键位置时，也会先向左收拢可用的连接；只是到了双行，末列的关键项可能落在下一行，便不能一直盯着上面。

霍雨浩将正式规则留在旁边。

\NovelDisplay{M00435R01}{%
p_1(j)=\max\{q:1\le q<j,\ a_q<a_j\},%
}
\NovelDisplay{M00435R02}{%
p_2(j):\quad\text{沿 }p_1(j),\ p_1(p_1(j)),\ldots\text{ 依次找},%
}
\NovelDisplay{M00435R03}{%
\text{取第一个满足 }b_q<b_j\text{ 的列 }q;\quad\text{无候选则父项不存在。}%
}

``第一行从整行左侧找；第二行沿本列第一行的祖先列，从近到远找。''他用两根手指分开示意，``具体碰到哪一个，待会儿放进例子里看。''

卢本韦已经拿起笔，等着阵图增加难度，霍雨浩却把刚才的四列推回中央。

``先算末列第二行为零的情形。这个最容易和单行的旧经验接上。''

末列下面是零，上面是三。于是从下往上找到的最后非零项，就是第一行的三；它的父项是第三列的二，坏根便落在第三列。

牢大亲自划出边界。

最左两列是好部，留在原处；第三列是坏部，既包含坏根，也到末列之前为止。最右边的三单独拿掉，不能混进要复制的那一段。

霍雨浩把最后非零行号记为 h，坏根列号记为 r，并将坏根标橙、原坏部其余部分标红、待删末列标紫，好部仍用正文色。展开第一项时新增的部分用蓝色，展开第二项时再新增的部分用绿色。

\NovelDisplay{M00436R01}{%
h=1,\quad r=3,\quad\text{坏根}=\textcolor[HTML]{B45309}{(2,1)},%
}
\NovelDisplay{M00436R02}{%
\text{好部}=(0)(1,1),\quad\text{坏部}=\textcolor[HTML]{B45309}{(2,1)},%
}
\NovelDisplay{M00436R03}{%
(\delta_1,\delta_2)=(0,0).%
}

蔡须坤在末列与坏根之间比了比。

``上面三减二，还差一。这里为什么是零阶差？''

``双行的截断规则。''霍雨浩指向末列下方那个零，``第二行的阶差恒为零；当末列第二行也是零，第一行的阶差同样取零。不能把普通的整列相减，直接当成最终阶差。''

她将原本准备标在第一行旁的一点光收回，重新落到坏部的边界上。

坏部只有一列，阶差为零，复制以后各坐标保留原值。展开第一项，在保留下来的原坏部后添一列；展开第二项，再多添一列相同的内容。

\NovelDisplay{M00437R01}{%
\bigl((0)(1,1)\textcolor[HTML]{B45309}{(2,1)}\textcolor[HTML]{7E22CE}{(3)}\bigr)[1]%
}
\NovelDisplay{M00437R02}{%
=(0)(1,1)\textcolor[HTML]{B45309}{(2,1)}\textcolor[HTML]{1D4ED8}{(2,1)}%
}
\NovelDisplay{M00437R03}{%
\bigl((0)(1,1)\textcolor[HTML]{B45309}{(2,1)}\textcolor[HTML]{7E22CE}{(3)}\bigr)[2]%
}
\NovelDisplay{M00437R04}{%
=(0)(1,1)\textcolor[HTML]{B45309}{(2,1)}\textcolor[HTML]{1D4ED8}{(2,1)}\textcolor[HTML]{15803D}{(2,1)}%
}

牢大将第一项的结果推到一边，正要沿这张阵图继续，蔡须坤按住了最初的输入纸。

``你换了原式。''

他低头看见纸上的末列已经不一样，手便停住了。

``我把顺次取项，当成接着对结果再施一次了。''

霍雨浩将原式放回两人中间。

``这里说第一项、第二项，都针对同一个原式。两项展开的结果可以并排看，计算时仍从原式的父项、坏根和阶差出发。''

蔡须坤松开纸，留给牢大重新接阵。他这一次先看最左端的输入，再把多出来的那一列送到末尾。

\NovelDisplay{M00438R01}{%
\bigl((0)(1,1)(2,1)(3)\bigr)[k]=(0)(1,1)\underbrace{(2,1)\cdots(2,1)}_{k+1\text{ 列}},\quad k\ge0.%
}

编号取多少，便在原坏部之后再添多少段。好部与原坏部都没有被替换掉，因此结果中相同的那一列，比新添的段数多一个。

诺当数完列，终于将先前握着的笔落下。

``阶差为零，只是数不增加，整段内容仍然继续复制。''

霍雨浩点了点头。

牢大将那张式子递回蔡须坤，食指压住她刚才提醒的位置。

``这一处给我留着。我看完新式子，容易顺手拿它当下一次输入。''

``已经圈上了。''

她把纸翻给他看，最初的四列上方画着一个清楚的框。他低头笑了笑，将自己的记录也放在她旁边。

霍雨浩等两人核完，没有立刻换成长式，而是收起中间两列，只留下两列最短的双行结构。

这回末列的第二行是一个一。

它上面的第一行只有最左零列可作祖先，第二行也在同一列找到严格更小的零。坏根落在第一列，好部为空，坏部就是这一列零。

\NovelDisplay{M00439R01}{%
p_1(2)=p_2(2)=1,\quad h=2,\quad r=1,%
}
\NovelDisplay{M00439R02}{%
(\delta_1,\delta_2)=(1-0,0)=(1,0).%
}

第二行的阶差仍为零，第一行这次却能增加一。

小舞先接第一项。零列保留，新添的一列由零成为一；蔡兴坤接着把第二项所需的后一列补上，得到二。这里的第二行始终是零，省略以后，看上去便成了单行数列。

\NovelDisplay{M00440R01}{%
\bigl((0)(1,1)\bigr)[0]=(0).%
}
\NovelDisplay{M00440R02}{%
\bigl((0)(1,1)\bigr)[1]=(0)(1).%
}
\NovelDisplay{M00440R03}{%
\bigl((0)(1,1)\bigr)[2]=(0)(1)(2).%
}
\NovelDisplay{M00440R04}{%
\bigl((0)(1,1)\bigr)[3]=(0)(1)(2)(3).%
}

牢大只看了一眼，便认出了这些数列。

``PRSS 的那条路，从这里接回来了。''

他将对应的序数写到纸上。霍雨浩逐项核对，没有让阵图靠直觉越过中间的层次。

\NovelDisplay{M00441R01}{%
(0)(1)=\omega,%
}
\NovelDisplay{M00441R02}{%
(0)(1)(2)=\omega^{\omega},%
}
\NovelDisplay{M00441R03}{%
(0)(1)(2)(3)=\omega^{\omega^{\omega}}.%
}

有限高度的幂塔，一层一层往上走，它们的上确界是第一个满足以自身为指数仍保持不变的 epsilon 数。

\NovelDisplay{M00442R01}{%
(0)(1,1)=\sup\{1,\omega,\omega^\omega,\omega^{\omega^\omega},\ldots\}=\varepsilon_0.%
}

蔡须坤没有将两行光阵拉得更大，只让刚才算出的几个有限项依次亮起。原式在最上方，每次取出的有限项则排在下面，各有明确的编号。

``能不断取更高的有限项，并不意味着现在已经接出一条无穷长的阵。''

``每一次用的仍是有限的那一项。''霍雨浩说，``写下省略号，是说明还可以继续。''

牢大缓缓收拢掌边的序力。他本来习惯在感到通路顺畅时，再多送入一层力量；这一次，他只让已经核过的第一项走完，便停在原位。

蔡须坤看见他收手，指尖也跟着松开。

``这一段比方才轻。''

``因为只接了算清楚的部分。''他偏过头，``你倒是盯得很紧。''

``刚才谁拿错了输入？''

牢大将那张被她圈过的纸往怀边收了一点，没有反驳。奥斯卡在旁边笑出声，顺手把温水递了一圈。

帐外有鸟振翅掠过，影子在帆布上轻轻一晃。霍雨浩抬头看了一眼，等众人放好杯子，才将接下来要算的五列摆开。

\NovelDisplay{M00443R01}{%
(0)(1,1)(2,2)(3,1)(4,2)%
}

这一回第一行从零到四依次上升，末列四的祖先列也便依次向左排列。第二行末项是二，沿刚才留下的第一行祖先链寻找，最近的严格更小项是第四列的一。

\NovelDisplay{M00444R01}{%
\text{末列第一行祖先列：}\quad4\to3\to2\to1,%
}
\NovelDisplay{M00444R02}{%
\text{候选列的第二行：}\quad1,\ 2,\ 1,\ 0;\qquad p_2(5)=4.%
}

坏根落在第四列。好部有前三列，坏部只有第四列；末列删去以后，要从原来的那列出发，新添按阶差调整过的内容。

\NovelDisplay{M00445R01}{%
\text{好部}=(0)(1,1)(2,2),\quad\text{坏部}=\textcolor[HTML]{B45309}{(3,1)},%
}
\NovelDisplay{M00445R02}{%
(\delta_1,\delta_2)=(4-3,0)=(1,0).%
}

卢本韦先写出普通差值，第二行也差一。蔡兴坤用笔尾在那一项旁边点了点，她便将最终阶差第二项改为零。

``下面只用来定坏根，这一次不随阶差增加。''

霍雨浩让她将这句话说完，又补了一句。

``下面的原值仍复制过去。不能把不增加，写成把这一行删掉。''

牢大将第四列留在蔡须坤手边，自己接起右边的新增部分。第一项只添一列，上面由三到四，下面的一保留；第二项保留前一段，再接上面为五、下面仍为一的一列。

\NovelDisplay{M00446R01}{%
\bigl((0)(1,1)(2,2)\textcolor[HTML]{B45309}{(3,1)}\textcolor[HTML]{7E22CE}{(4,2)}\bigr)[1]%
}
\NovelDisplay{M00446R02}{%
=(0)(1,1)(2,2)\textcolor[HTML]{B45309}{(3,1)}\textcolor[HTML]{1D4ED8}{(4,1)}%
}
\NovelDisplay{M00446R03}{%
\bigl((0)(1,1)(2,2)\textcolor[HTML]{B45309}{(3,1)}\textcolor[HTML]{7E22CE}{(4,2)}\bigr)[2]%
}
\NovelDisplay{M00446R04}{%
=(0)(1,1)(2,2)\textcolor[HTML]{B45309}{(3,1)}\textcolor[HTML]{1D4ED8}{(4,1)}%
}
\NovelDisplay{M00446R05}{%
\quad\textcolor[HTML]{15803D}{(5,1)}%
}

蓝色的一段与绿色的一段分得清楚。霍雨浩将坏根所在的原列标成橙色，原坏部其余部分用红色；这里坏部只有坏根一列，便只留橙色。待删末列标紫，好部仍用正文色。

``两段都回到橙色的原列计算。''他说，``倍数不同，原列却是同一列。''

蔡须坤沿着下面那排一看过去，没有再给它们添光。牢大原本托着新增部分的手，也稳稳停在既定的位置。

他正要让阵图再走一遍，霍雨浩忽然收住最右列，把第一行的四换成了二。

第二行仍是二，前面四列没有改变。

\NovelDisplay{M00447R01}{%
(0)(1,1)(2,2)(3,1)(2,2)%
}

``先只找父项。''他说，``别沿用刚才的坏根。''

牢大的目光先落在第二行。第四列的一离得最近，也小于二，他的手刚向那处伸去，蔡须坤便将上方的连线抬了起来。

``先看第一行。''

末列现在是二。它向左遇到三，太大；再遇到另一个二，相等，都越过去，父项才落在第二列的一。沿这一行继续追，下一处是第一列的零。

第四列与第三列并不在这条祖先链里。

\NovelDisplay{M00448R01}{%
p_1(5)=2,\qquad\text{第一行祖先列：}\quad2\to1,%
}
\NovelDisplay{M00448R02}{%
p_2(5)=2,\qquad\text{第三、第四列不在本次候选中。}%
}

牢大收住伸出去的手，沿蔡须坤亮出的两处位置重新看第二行。

第二列的一严格小于二，可以作父项。坏根便从刚才的第四列移到了第二列。

``原先那个一，数值没错，位置不合条件。''

``对。''蔡须坤将第四列下面的一留在原处，``它在这一行还在，却没有资格参与末列这一回的父项选择。''

霍雨浩把精神共享放得更细，让每个人都能看见同样的列号。两处数值相同，来自哪一列，却不能混为一谈。

``你们感觉哪一处牵动，可以从那里开始查。''他说，``父项最后落在哪里，要按刚才这条链核对。''

牢大点头，掌心的序力没有再压向原先那处接点。

``只改一个数，先前的坏根和分段都得重算。不能因为刚才已经走顺，就把旧线直接接过来。''

蔡须坤看了他一眼，将第二列与后面两列一起托起。

这次好部只剩最前面的一列零。坏部从第二列开始，一直延伸到第四列；阶差仍是一与零，但要复制的内容已有三列。

\NovelDisplay{M00449R01}{%
\text{好部}=(0),\quad\text{坏部}=(1,1)(2,2)(3,1),%
}
\NovelDisplay{M00449R02}{%
r=2,\quad(\delta_1,\delta_2)=(2-1,0)=(1,0).%
}

牢大负责保持这三列的相对位置，蔡须坤依次检查下面的原值，霍雨浩记录新添部分与原列的对应。

\NovelDisplay{M00450R01}{%
\textcolor[HTML]{B45309}{(1,1)}\longmapsto\textcolor[HTML]{1D4ED8}{(1+q,1)},%
}
\NovelDisplay{M00450R02}{%
\textcolor[HTML]{B91C1C}{(2,2)}\longmapsto\textcolor[HTML]{1D4ED8}{(2+q,2)},%
}
\NovelDisplay{M00450R03}{%
\textcolor[HTML]{B91C1C}{(3,1)}\longmapsto\textcolor[HTML]{1D4ED8}{(3+q,1)},\quad1\le q\le k.%
}

这里 k 是这次要展开的项号，q 表示其中第几项新添的那一段。第一项取一倍阶差，第二项再新添的部分取两倍阶差；完整的第二项里，前面按一倍阶差生成的那一段也仍保留。

三列不能各自挑一个方便的编号。算第一项，就在这一整段里统一用同一个倍数，再按原来的次序接回去。

唐三稳住前面的零列，诺当按顺序报数。小舞与蔡兴坤在旁边核对边界，卢本韦则把先前那一列坏部的旧答案收好，给新结果腾出地方。

\NovelDisplay{M00451R01}{%
\bigl((0)\textcolor[HTML]{B45309}{(1,1)}\textcolor[HTML]{B91C1C}{(2,2)}\textcolor[HTML]{B91C1C}{(3,1)}\textcolor[HTML]{7E22CE}{(2,2)}\bigr)[1]%
}
\NovelDisplay{M00451R02}{%
=(0)\textcolor[HTML]{B45309}{(1,1)}\textcolor[HTML]{B91C1C}{(2,2)}\textcolor[HTML]{B91C1C}{(3,1)}\textcolor[HTML]{1D4ED8}{(2,1)}%
}
\NovelDisplay{M00451R03}{%
\quad\textcolor[HTML]{1D4ED8}{(3,2)}\textcolor[HTML]{1D4ED8}{(4,1)}%
}

牢大一直等到最后一列落定，才缓缓松开托住它的手。

``原来五列，现在七列。''

蔡须坤却没有随他去看最右端。她将原末列与新添部分的首列并到比较的位置，先看第一行，再看第二行。

第一行都是二，第二行却分别是二与一。

\NovelDisplay{M00452R01}{%
(0)(1,1)(2,2)(3,1)(2,1)(3,2)(4,1)%
}
\NovelDisplay{M00452R02}{%
\prec(0)(1,1)(2,2)(3,1)(2,2),\qquad1<2.%
}

``结果变长，仍是向原式下面取了一项。''她说，``得看共同前缀后首先不同的位置。''

霍雨浩将更右边的两列暂时压暗，留下那处最先出现的差别。

``已经在这里分出顺序，后面再多添几列，也不会把前面这一处改掉。只有前面的各列全相同，较短的一边又恰好是另一边的前缀时，才能直接用长短比较。''

牢大把阵图往左收了一点，给她留出看两列的空处。他先前一直托着增长的末端，此时才把手移到那个更早出现差别的位置。

``阵看着更宽，不代表它越过了原式。''

他没有再添力，将比较的结果直接记到自己的纸上。

霍雨浩将三张主例按先后放好，原式与各自的坏根都留在纸边。

零阶差的一张，靠重复原坏部生成新项；一列坏部、非零阶差的一张，添出的每段只占一列；最后一张的阶差虽然也是一与零，坏部却有三列，每项新添的内容必须整段计算。

``现在把共同的步骤收起来。''他说。

牢大先按住原式，给末列单独留出空位。

``先找末列最后一个非零项。若下面为零，就看上面；下面非零，便从下面找父项。父项所在列是坏根。''

蔡须坤接着划出两道边界。

``好部到坏根之前，坏部从坏根开始、到末列之前。末列删去，原好部与原坏部都保留。''

霍雨浩先在页边补了一类起始情形：末列全零时，原式对应后继，删去这一列得到前驱。下面的阶差通式只针对末列非零的情况，不能把全零末列也送去找坏根。

随后，他才将两种阶差写在下面。

\NovelDisplay{M00453R01}{%
b_m=0:\quad r=p_1(m),\quad(\delta_1,\delta_2)=(0,0),%
}
\NovelDisplay{M00453R02}{%
b_m>0:\quad r=p_2(m),\quad(\delta_1,\delta_2)=(a_m-a_r,0).%
}

``r 是坏根列号，m 是原末列号。''他将它们与前面的列号对齐，``第二行的阶差恒为零。末列第二行为零时，第一行阶差也为零；不为零时，第一行才取末列减坏根。''

唐三将第一张的坏根代入，蔡兴坤核对最后一张。两者都落回刚才的结果。

牢大指向一般式里要复制的那一段。

``只有两行时，坏部第一行都能按这个差增加？''

``这里最后非零项若在第二行，坏根就在末列第一行的祖先链上。''霍雨浩沿那条链指回去，``坏根之后、末列之前，第一行若有一项不大于坏根第一行的值，就会截断末列向前追到这个坏根的链。这与我们已经找到的祖先关系不合。''

他把几列中间项逐一亮起。

``所以这一段第一行都接回这个根。双行只可能增加第一行，第二行的阶差为零，保留原值即可。更多行时，还要逐行查提升条件，那是后面的功课。''

蔡须坤顺着原坏部重新走完，没有把最后那句当成可以略过父链的理由。坏根若找错，前面的保证便不成立。

最后的新添坐标与列数，只需留下两行。

\NovelDisplay{M00454R01}{%
(a_j,b_j)\longmapsto(a_j+q\delta_1,b_j),\quad r\le j<m,\quad1\le q\le k,%
}
\NovelDisplay{M00454R02}{%
\text{第 }k\text{ 项总列数}=m-1+k(m-r),\qquad k=0,1,2,\ldots.%
}

q 对应第几项新添的那一段，k 对应这次要展开到的项号。完整第 k 项的新添内容，要依次包含按一倍、两倍，一直到 k 倍阶差得到的各段。每段都从同一段原坏部出发，内部列的顺序保持不变。

诺当将列数公式用在刚才的七列结果上。原式五列，坏根列号二，每项新添三列；第一项删去原末列再加三列，正好七列。

``列数可以帮忙查有没有少接一列。''他说，``父项与坐标，还是要另外核。''

``每个编号都是有限的。''牢大又看向阵图最右端，``选第几项，就把这一项算完整。''

霍雨浩应了一声，将还没用到的空处收回。浅金色阵图缩到草席中央，众人都能伸手指到具体的位置。

蔡须坤伸展了一下握笔的手腕。牢大将杯子推近，她却先把他的笔记翻开，在``分段''后面补上了``以新坏根重算''几个字。

``你的字写得比我的还小。''她说。

``地方都让给阵图了。''

他将纸向她那边移了移，等她补完，才接过笔。两人合看一张图，谁要核对哪一处，便把纸略转过去，彼此不用再伸手去够。

霍雨浩看着几张已经写齐的纸，轻轻吐出一口气。

``规则讲到这里，再对照你们熟悉的坍缩记号。''

他在新页上写下 BOCF 的基本约定。外层未另标下标的 ψ，以 Ω 为坍缩上界；更高层的下标则直接写出对应的 Ω 层级。

\NovelDisplay{M00455R01}{%
\psi(0)=1,\qquad\psi(\Omega)=\varepsilon_0,%
}
\NovelDisplay{M00455R02}{%
\psi_{\Omega_2}(0)=\Omega,\qquad\psi_{\Omega_3}(0)=\Omega_2.%
}

``同一个函数名，得先看采用哪套定义。''他指着第一行，``这一页里，ψ 在零处是一，在 Ω 处才是刚才的 epsilon 零。后面的对应都照这一套来。''

蔡须坤在页角添上``BOCF''，又把下标写法圈清楚。她以前读不同传承时，也遇见过相似的符号；若省掉约定，单凭字面拼起来，前后就可能落在不同层次。

牢大没有催霍雨浩把整张对照谱全投出来，而是把最短的双行结构放到纸边。

``先从它后面添零、添一开始。这样能看清同样的尾数，为什么会起不同作用。''

霍雨浩选了六条相邻关系，逐一排开。

\NovelDisplay{M00456R01}{%
(0)(1,1)(0)=\psi(\Omega)+1.%
}
\NovelDisplay{M00456R02}{%
(0)(1,1)(0)(1,1)=\psi(\Omega)\times2.%
}
\NovelDisplay{M00456R03}{%
(0)(1,1)(1)=\psi(\Omega)\times\omega.%
}
\NovelDisplay{M00456R04}{%
(0)(1,1)(1)(1)=\psi(\Omega)\times\omega^{2}.%
}
\NovelDisplay{M00456R05}{%
(0)(1,1)(1)(2,1)=\psi(\Omega)^{2}.%
}
\NovelDisplay{M00456R06}{%
(0)(1,1)(1)(2,1)(2)=\psi(\Omega)^{\omega}.%
}

最末添一列零，得到后继。将完整的基础单元重复一次，形成两个同类单元的有限相加；末列换成第一行为一、第二行为零，则坏部回到整段基础单元，零阶差使这段内容不断重复，对应乘到 ω。

牢大从表中挑出这一条，没有先照着右端念答案，而是自己重新分部。

末列的一向左找到第一列的零作父项，好部为空，坏部是前面那两列。阶差为零，展开第一项便将这两列再接一遍。

\NovelDisplay{M00457R01}{%
\bigl(\textcolor[HTML]{B45309}{(0)}\textcolor[HTML]{B91C1C}{(1,1)}\textcolor[HTML]{7E22CE}{(1)}\bigr)[1]%
}
\NovelDisplay{M00457R02}{%
=\textcolor[HTML]{B45309}{(0)}\textcolor[HTML]{B91C1C}{(1,1)}\textcolor[HTML]{1D4ED8}{(0)}\textcolor[HTML]{1D4ED8}{(1,1)}%
}

蔡须坤将重复单元之间的边界逐一标好。

``这里先有完整的一段，又添完整的一段。末尾那个一不是给序数加一，和刚才添零的后继不同。''

``而且要看这个一落在什么前缀之后。''霍雨浩说，``不能把任意一个末尾的一，都读成固定的一种运算。''

她把这句话写在表下，只在``这几个结构中''旁边画了一道短线。原式的父项、坏根与坏部，都还留在前页，随时能翻回去核对。

牢大接着挑出平方那一条。他让蔡须坤保持前缀，自己沿末列第二行寻找父项。坏部回到中间那列第一行为一、第二行为零的内容，第一行阶差是一。

展开第二项时，新添的两列分别在原列基础上按一倍、两倍阶差生成。

\NovelDisplay{M00458R01}{%
\bigl((0)(1,1)\textcolor[HTML]{B45309}{(1)}\textcolor[HTML]{7E22CE}{(2,1)}\bigr)[2]%
}
\NovelDisplay{M00458R02}{%
=(0)(1,1)\textcolor[HTML]{B45309}{(1)}\textcolor[HTML]{1D4ED8}{(2)}\textcolor[HTML]{15803D}{(3)}%
}

``看见一个二，不能直接说平方。''牢大看着新添的那一段，``先把整个式子算出，再核对它怎样逼近原来的结构。''

``这里的平方，是整个原式对应的结果。''蔡须坤说，``新增段里的二、三，各有它们在同一套列关系中的位置。''

霍雨浩让两人将各自负责的部分合在一起。前缀未动，后面的第一行逐步上升，第二行的零也没有被另添出一层。

帐内的光落在蔡须坤掌背上。她收手时，牢大的手还停在阵图边缘，两人手指轻轻碰了一下，他便替她稳住正在翻动的那页纸。

``再往上看？''他问。

她点头，把纸压在两人中间。

接下来几条，开始把原有的结构放入更深的坍缩参数。霍雨浩选出其中五条，足以看清有限重复、ω 倍与幂次之间的差别。

\NovelDisplay{M00459R01}{%
(0)(1,1)(1,1)=\psi(\Omega\times2).%
}
\NovelDisplay{M00459R02}{%
(0)(1,1)(2)=\psi(\Omega\times\omega).%
}
\NovelDisplay{M00459R03}{%
(0)(1,1)(2,1)=\psi(\Omega^{2}).%
}
\NovelDisplay{M00459R04}{%
(0)(1,1)(2,1)(2,1)=\psi(\Omega^{3}).%
}
\NovelDisplay{M00459R05}{%
(0)(1,1)(2,1)(3)=\psi(\Omega^{\omega}).%
}

蔡兴坤看见前面两个相同的双行列，想将它们直接当成两段基础单元相加，牢大却把算有限相加时用的那张纸推到了旁边。

``那一条中间有零列。这里没有。''

两张式子并排，分段的位置立刻显了出来。相同的列形，在不同前缀与父链下，不能只按出现次数套用同一个算术。

霍雨浩随即把相邻同列的那一条展开第一项。

\NovelDisplay{M00460R01}{%
\bigl(\textcolor[HTML]{B45309}{(0)}\textcolor[HTML]{B91C1C}{(1,1)}\textcolor[HTML]{7E22CE}{(1,1)}\bigr)[1]%
}
\NovelDisplay{M00460R02}{%
=\textcolor[HTML]{B45309}{(0)}\textcolor[HTML]{B91C1C}{(1,1)}\textcolor[HTML]{1D4ED8}{(1)}\textcolor[HTML]{1D4ED8}{(2,1)}%
}

新添的一段有两列。第一列下方是零，后一列下方是一；两列共同参与后面的结构，不能只拿最后一个一，解释整个原式。

蔡须坤顺着新段看完，又回到对照表第一行。

``这里变化进入了 ψ 的参数。前面的后继，却是在整个 ψ 值之后再加一。''

霍雨浩点头，将参数内外的两个位置分别点亮。

``区分括号内外，再看加法、乘法和幂放在哪里。对照谱里的式子很长，真正要核的是这些位置。''

牢大正要往后翻，蔡须坤却按住了页角。

``前面从二次幂、三次幂走到 ω 次幂，我还能接上。可下一页就出现 Ω 的下标二，中间总该有路。''

``有。先将刚才零阶差的例子接回它的序数。''霍雨浩把前页的坏部重新点亮，``末列是单独的三，删掉后，留下一个二与一。每多展开一项，就再添同样的一列。''

他让牢大数坏部，让蔡须坤数 ψ 参数中的幂次。

\NovelDisplay{M00461R01}{%
(0)(1,1)\underbrace{(2,1)\cdots(2,1)}_{n\text{ 列}}=\psi(\Omega^{n+1}),\quad n\ge1,%
}
\NovelDisplay{M00461R02}{%
\psi(\Omega^\omega)=\sup_{n\ge1}\psi(\Omega^{n+1}).%
}

这里 n 数的是末尾二与一的总列数，原来就有的那列也包括在内。因此一列对应二次幂，两列对应三次幂，后面的幂次始终比列数多一。有限次重复没有终点，原式才由这些已经核过的展开项逼近到 ω 次幂。

``这行上确界，是沿这个 BMS 的基本序列取的。''霍雨浩补了一句，``不能拿着任意一串 ψ 记号，就不检查展开而将极限移进括号。''

牢大将方才准备翻过去的纸重新铺平。蔡须坤松开手，两人的笔落在同一行前后。

``那么在 ω 的上面，先改变指数内部？''她问。

``对，先区分加一、乘到 ω，再看指数里出现新的坍缩值。下面几条保留完整前缀，才能辨认改变发生在哪里。''

\NovelDisplay{M00462R01}{%
(0)(1,1)(2,1)(3)(2,1)=\psi(\Omega^{\omega+1}).%
}
\NovelDisplay{M00462R02}{%
(0)(1,1)(2,1)(3)(2,1)(3)=\psi(\Omega^{\omega\times2}).%
}
\NovelDisplay{M00462R03}{%
(0)(1,1)(2,1)(3)(4,1)=\psi(\Omega^{\psi(\Omega)}).%
}
\NovelDisplay{M00462R04}{%
(0)(1,1)(2,1)(3)(4,1)(5,1)=\psi(\Omega^{\psi(\Omega^{2})}).%
}
\NovelDisplay{M00462R05}{%
(0)(1,1)(2,1)(3,1)=\psi(\Omega^{\Omega}).%
}

前两行沿着刚才的结构继续组织，指数从 ω 变成 ω 加一，再变成 ω 乘二。第三、第四行则让 ψ 的值进入指数：先是 ψ 在 Ω 处的值，随后是 ψ 在 Ω 平方处的值。最外层的 ψ 仍然包住整段 Ω 的幂，里面另出现的 ψ 也必须连同自己的括号一起读。

蔡须坤用指尖隔开内外两层括号。

``第三行里，指数本身就是前面学过的 epsilon 零。不能把里面的 ψ 擦掉，直接说成指数已经到了 Ω。''

霍雨浩点头，将第三行与最后一行之间的空白留宽些。

``这里还容得下许多中间项。我们挑有代表性的结构，说明指数怎样继续抬升；表中相邻的两行，并不表示中间没有其他序数。''

牢大将最后一行的末列拉近。之前的三下面补上了一，但它选父项的方式也随之改变。他先沿第一行走过第三、第二、第一列，再检查这些列的第二行。

\NovelDisplay{M00463R01}{%
\text{末列第一行祖先列：}\quad3\to2\to1;\quad\text{第二行：}1,1,0,%
}
\NovelDisplay{M00463R02}{%
p_2(4)=r=1,\qquad(\delta_1,\delta_2)=(3-0,0)=(3,0),%
}
\NovelDisplay{M00463R03}{%
\text{好部为空；坏部}=\textcolor[HTML]{B45309}{(0)}\textcolor[HTML]{B91C1C}{(1,1)(2,1)}.%
}

第三列、第二列下方都是一，不能作末项第二行那个一的父项；走到第一列，下方的零才满足严格小于。坏根因此退到最前面的零列，坏部包含三列，第一行阶差也变成三。

``所以不能还照刚才那样，只重复一列二与一。''牢大把旧的坏部标记收回，``这次要将零、一与一、二与一整段接上，而且第一行都加三。''

蔡须坤将新增段接到他指的位置。霍雨浩把原来的坏根染成赭色，旧坏部其余两列染红，末列染紫，新增的三列则一起染蓝。

\NovelDisplay{M00464R01}{%
\bigl(\textcolor[HTML]{B45309}{(0)}\textcolor[HTML]{B91C1C}{(1,1)}\textcolor[HTML]{B91C1C}{(2,1)}\textcolor[HTML]{7E22CE}{(3,1)}\bigr)[1]%
}
\NovelDisplay{M00464R02}{%
=\textcolor[HTML]{B45309}{(0)}\textcolor[HTML]{B91C1C}{(1,1)}\textcolor[HTML]{B91C1C}{(2,1)}\textcolor[HTML]{1D4ED8}{(3)}\textcolor[HTML]{1D4ED8}{(4,1)}%
}
\NovelDisplay{M00464R03}{%
\quad\textcolor[HTML]{1D4ED8}{(5,1)}%
}
\NovelDisplay{M00464R04}{%
=\psi(\Omega^{\psi(\Omega^{2})}).%
}

蓝色新增段从单独的三起头，接着是四与一、五与一。末列下方原来的一已经随紫色末列删去；新增段下方的零、一、一，分别来自旧坏部的三列。

蔡须坤没有急着看最后的对应值，先将三个增量逐个核完。

``零加三是一段的开头，一加三、二加三接在后面。三列都属于展开第一项新增的同一段。''

``而这个完整展开项，对应的是外层 Ω 的指数里再放一个 ψ。''霍雨浩指向末行，``它还小于原式的 ψ 在 Ω 的 Ω 次幂处的值。继续取项时，嵌套会向更深处生长。''

牢大看了看前面只添一列的展开，又看向新添的三列。

``同样叫展开第一项，新增多少列，由坏部长度决定。阶差决定怎么抬升，这两件事得一起算。''

霍雨浩让光线继续向后铺开，接出下一组有代表性的台阶。

\NovelDisplay{M00465R01}{%
(0)(1,1)(2,1)(3,1)(3,1)=\psi(\Omega^{\Omega^{2}}).%
}
\NovelDisplay{M00465R02}{%
(0)(1,1)(2,1)(3,1)(4)=\psi(\Omega^{\Omega^{\omega}}).%
}
\NovelDisplay{M00465R03}{%
(0)(1,1)(2,1)(3,1)(4,1)=\psi(\Omega^{\Omega^{\Omega}}).%
}
\NovelDisplay{M00465R04}{%
(0)(1,1)(2,1)(3,1)(4,1)(5,1)=\psi(\Omega^{\Omega^{\Omega^{\Omega}}}).%
}

第一行将外层 Ω 的指数推进到 Ω 的平方；第二行将那一层推进到 Ω 的 ω 次幂。第三、第四行继续增加幂塔的高度。幂塔从右向左结合，右上角那层完整的幂，才是左边一层的指数。

蔡须坤将最外层的 ψ 留在原处，只用笔尖顺着指数由里向外走了一遍。

``这一组还是把更复杂的内容放在指数里。到了下一组，才要解释 Ω 的下标二。''

``现在回头算那个最短的三列式子。''霍雨浩说。

他保留首列零与中间的一与一，将末列换成二与二。牢大顺着第一行选到第二列；末项第二行是二，第二列下方的一也已经严格小于它，于是父项就在这里停住。

\NovelDisplay{M00466R01}{%
p_1(3)=p_2(3)=r=2,%
}
\NovelDisplay{M00466R02}{%
\text{好部}=(0),\quad\text{坏部}=\textcolor[HTML]{B45309}{(1,1)},%
}
\NovelDisplay{M00466R03}{%
(\delta_1,\delta_2)=(2-1,0)=(1,0).%
}

坏部只有中间这一列。删掉末列后，首列零仍留在好部，中间的一与一也留在原处。每新增一段，都以原坏部为基准，将第一行加上相应倍数的阶差。

\NovelDisplay{M00467R01}{%
\textcolor[HTML]{B45309}{(1,1)}\longmapsto(1+q,1),\qquad1\le q\le k.%
}

``新增第一段是二与一，第二段是三与一，第三段是四与一。''蔡须坤边说边将它们排在原来的一与一后面。

牢大这次没有将第几项混成反复展开的次数。他让每一行都从同一个原式出发，再让霍雨浩把对应值写在右侧。

\NovelDisplay{M00468R01}{%
\bigl((0)(1,1)(2,2)\bigr)[0]=(0)(1,1)=\psi(\Omega).%
}
\NovelDisplay{M00468R02}{%
\bigl((0)(1,1)(2,2)\bigr)[1]=(0)(1,1)(2,1)=\psi(\Omega^{2}).%
}
\NovelDisplay{M00468R03}{%
\bigl((0)(1,1)(2,2)\bigr)[2]=(0)(1,1)(2,1)(3,1)=\psi(\Omega^{\Omega}).%
}
\NovelDisplay{M00468R04}{%
\bigl((0)(1,1)(2,2)\bigr)[3]=(0)(1,1)(2,1)(3,1)(4,1)=\psi(\Omega^{\Omega^{\Omega}}).%
}

第零项只删除末列，所以回到 ψ 在 Ω 处的值；第一项再添一列二与一，得到 Ω 平方的那一层；第二项保留原来的两列，再依次添二与一、三与一，来到 Ω 的 Ω 次幂；第三项又添出四与一，对应更高一层的幂塔。

每一项都比原式低，却能沿着第一行连续上升、第二行保持为一的结构，走到任意有限高度。先前那几组很长的式子，此时被同一条展开序列串了起来。

\NovelDisplay{M00469R01}{%
\psi(\Omega_2)=\sup\bigl\{\psi(\Omega),\ \psi(\Omega^2),\ \psi(\Omega^\Omega),%
}
\NovelDisplay{M00469R02}{%
\qquad\psi(\Omega^{\Omega^\Omega}),\ \psi(\Omega^{\Omega^{\Omega^\Omega}}),\ldots\bigr\}.%
}

``这里取上确界的是整个 ψ 值。''蔡须坤把指尖停在最外面的括号，``右边不是说 Ω 的下标二，等于这些 Ω 幂塔本身的极限。''

``对。下标二标明 BOCF 参数中的另一层；对应到这个 BMS，它的基本序列恰好逐项给出这些已经核过的坍缩值。''霍雨浩将两边的最外层同时点亮，``从展开看懂它怎样逼近，再写最后的等号，才不会把参数本身和坍缩后的值混在一起。''

牢大终于翻过这一页。他没有再催着看高处，而是将坏根退到第一列的例子，与坏根停在第二列的例子并排夹好。

``前一个坏部有三列，后一个只有一列。式子短了，所能统摄的层级却更高。以后碰见相似形状，还是先找父项。''

霍雨浩在旁边留下最后四条对照，接着说明更高的两行对角结构怎样进入不同的 Ω 层级。

\NovelDisplay{M00470R01}{%
(0)(1,1)(2,2)=\psi(\Omega_2).%
}
\NovelDisplay{M00470R02}{%
(0)(1,1)(2,2)(3,2)=\psi(\Omega_2^{2}).%
}
\NovelDisplay{M00470R03}{%
(0)(1,1)(2,2)(3,3)=\psi(\Omega_3).%
}
\NovelDisplay{M00470R04}{%
(0)(1,1)(2,2)(3,3)(4,4)=\psi(\Omega_4).%
}

第二行的数增大后，仍要按同样的祖先限制选父项；它会参与更高坍缩层的组织，却不能简单读成``看到第二行是几，就把 ψ 的下标写成几''。

``要看整段。''蔡须坤指向表中共有的前缀，``少一列、换一个父项，后面的组织就可能不同。''

牢大把最后一条与前面比较了一遍。矩阵仍然只有两行，列中出现的数却能继续变大。行数固定，没有把它限制成只能写几个固定层级。

他重新看向霍雨浩。

``只要长得像这样，都能当成标准记号？''

``还得查来源。''

霍雨浩将必要的局部检查写在纸上。

\NovelDisplay{M00471R01}{%
a_1=b_1=0,\qquad0\le b_j\le a_j,%
}
\NovelDisplay{M00471R02}{%
a_j>0\Longrightarrow p_1(j)\text{ 存在且 }a_j\le a_{p_1(j)}+1,%
}
\NovelDisplay{M00471R03}{%
b_j>0\Longrightarrow p_2(j)\text{ 存在且 }b_j\le b_{p_2(j)}+1.%
}

首列全零，同列下项不大于上项；非零项必须有父项，并且不能比父项大出一以上。仅仅写出两排非负整数，还没有完成这些检查。

他留下一个最短的错误例子。

\NovelDisplay{M00472R01}{%
(0)(1,1)(3,1):\quad p_1(3)=2,\qquad3>1+1.%
}

蔡须坤找到那项三的父项，牢大则沿同列上下核了一遍。列内上下关系虽然没问题，第一行从父项的一直接到了三，却过不了父项约束。

``能找到父项，与数值满足约束，还要分开查。''她将两处判断各留一个记号。

霍雨浩把错误例子收在纸边，没有继续堆同类题。

``通过局部检查之后，具体表达式怎样由已知入口生成，也应当留下依据。标准性不能靠看起来短不短来决定。展开后变长的式子，照样可以属于标准表达式。''

牢大摸了摸笔记最末的空白页。

``双行不断往高处取，它的入口在哪里？''

霍雨浩将对照谱翻到两行结构的尽头。那里不是一张已经铺了无穷多列的阵图，而是一列零与一列三个一。

\NovelDisplay{M00473R01}{%
(0)(1,1,1)=\psi(\Omega_{\omega}).%
}

这个式子已经有第三行。它标出前面有限双行标准结构所能不断趋近的边界，而非那些结构中的最大一项。表中从这里往后，才开始继续组织三行的更高层次。

霍雨浩先将双行对角结构的规律补在它下面。

\NovelDisplay{M00474R01}{%
(0)(1,1)(2,2)\cdots(n,n)=\psi(\Omega_n),\qquad n\ge2.%
}

这里 n 始终是有限整数。对角线向右增长，便依次到达不同的有限下标层级；把这些层级继续取上去，才趋向刚才写出的边界。不能在一个普通有限输入里，直接放进一列无穷长的坐标。

蔡须坤将三行入口停在阵图上方，暂时没有把它接进大家掌中的练习。

``这一张留作边界。我们正在练的，还是下面已经核过的有限项。''

``是。''霍雨浩说，``三行的详细规则以后另算。今天可以先沿传承中已知的取项，从这个入口走回刚才做过的一张双行主例，给它留一条明确的生成路径。''

他让诺当记箭头编号，让唐三核对每一步的末列；牢大与蔡须坤照着结果接阵，其他人依次查新增的列数。

\NovelDisplay{M00475R01}{%
(0)(1,1,1)%
}
\NovelDisplay{M00475R02}{%
\xrightarrow{[3]}(0)(1,1)(2,2)(3,3)%
}
\NovelDisplay{M00475R03}{%
\xrightarrow{[1]}(0)(1,1)(2,2)(3,2)%
}
\NovelDisplay{M00475R04}{%
\xrightarrow{[1]}(0)(1,1)(2,2)(3,1)(4,2)%
}

第一步从三行入口取第三项，得到双行的四列；后面两步再各自取第一项，终于落到刚才末列是四与二的那个主例。

霍雨浩把三个箭头依次核完。每一步都有明确的输入、取项编号和输出；现在这张主例既有具体展开算法，也有从边界生成它的路径。

牢大将这条路径折在原题后面，手指还按着那张明显更长的结果。

``长短不能决定标准性。由哪一步产生，才是我们这里留下的依据。''

蔡须坤把自己的记录也合拢，将两次输入的父链并排留好，又圈出牢大差点沿用旧坏根的位置。

``下次再遇到第一行忽然下降，就先重查祖先链。''

``你倒是记得我错在哪里。''

``你让我留着的。''

她抬眼望向牢大。他低低笑了一声，将两人的纸叠好，又把她刚才落在草席边的笔递回掌中。

霍雨浩正收拢光幕，听见两人的话，也笑了。

``这张先别收。最后一次，换你们接，我只看结果。''

三行入口与较高的对照谱淡去，只留最后那张需要重新筛选祖先列的双行阵图。

牢大先按住原式的位置，蔡须坤沿第一行祖先链找到第二行父项。她报出坏根以后，他才将好部、坏部和待删末列各自分开；新添部分由两人共同接起，先核第一行的阶差，再对照下面保留的原值。

霍雨浩将精神共享收窄，留给其他队员各查一处。唐三确认前缀，诺当数列，小舞核顺序，蔡兴坤与卢本韦各自代入编号，奥斯卡则把两张演算纸对齐，省得有人把第二项新添的一段抄到第一项下面。

阵图走完一遍，前后的光没有相互拖动。

牢大没有立刻加大序力。他等蔡须坤把末列旁的两个数又看一遍，才与她一起收手。

她掌心留下些许温热，抬头时，正看见他还守着那一小段新接的线。方才领悟更高折叠时那股迫人的气息已经敛下，他低着头，仍在认真确认最后一列有没有接错。

蔡须坤伸手，轻轻碰了一下他的手背。

``可以了。''

牢大翻过手，短短握住她的指尖，又松开。

``再核一遍，总比带着错处走远好。''

帐外的溪水声仍在。奥斯卡将杯子放回原处，唐三靠着帐柱开始复盘，小舞拿着刚核过的纸，与蔡兴坤低声交换下一步的编号。

牢大与蔡须坤仍并肩坐着，各自将那段有限回路映入识海。霍雨浩见众人已经知道从哪里查起，也没有再往光幕上添新的整页公式。

他只把伊莱克斯笔记中另一个三行表达式写到后页，留在众人能够看见的位置。

\NovelDisplay{M00476R01}{%
(0)(1,1,1)(2,1)(1,1,1)%
}

``三行里还有这样的结构。''霍雨浩的声音放轻了些，``初步展开可以继续算，可这里的序数怎样组织、哪一步发生了变化，我还没分析清楚。先把今天的双行规则练稳，后面再一起往下查。''

\chapter{第四十一章：死斗，帅气小队的威力}\label{ux7b2cux56dbux5341ux4e00ux7ae0ux6b7bux6597ux5e05ux6c14ux5c0fux961fux7684ux5a01ux529b}

林间清风温润，帐内道韵绵长。

果糕九人小队尽数沉心悟道，识海之中双行BMS序列层层铺展，ε₀、ζ₀、Γ₀等序数脉络交替流转，每个人都在疯狂吸纳全新的递归大道，打磨自身序力根基。霍雨浩最后的话音余韵未落，关于三行BMS的混沌未知还萦绕在众人心头，整片山谷的祥和氛围骤然被一股刺骨的死寂撕碎。

没有风声异动，没有魂力先兆，纯粹、冰冷、碾压一切的高阶序力，瞬间锁死了整片山谷的每一寸空间。

``嗡------！''

漆黑的虚空阴影从天地缝隙中疯狂涌出，原本隐匿蛰伏的杀机彻底爆发。数十道黑袍身影撕裂空气，错落落地，周身萦绕着浓郁的邪道序力，每一人的气息都凶悍磅礴，远超寻常邪教精锐。为首的帅气小队小队长负手而立，黑袍猎猎，眼底满是戏谑与阴冷的轻蔑。

他周身浮动的序数纹路清晰浮现，BMS序列：0 111 222，远超果糕小队众人此刻的双行BMS根基，对应的ψ(a\_ω)层级力量轰然铺开，形成一道密不透风的压制力场。

``钻研双行BMS？推演序数天梯？''小队长嗤笑出声，声音冰冷刺骨，``一群井底之蛙，刚摸到大数修行的门槛，便以为踏足大道之巅？今日，就让你们见识一下，真正的高阶序列之力，究竟何等恐怖！''

猝不及防的绝杀偷袭，骤然降临！

帐内悟道的九人瞬间惊醒，所有人神识紧绷，周身序力瞬间暴涨，本能般撑起防御屏障。唐三眼神骤凝，身形瞬间挡在众人最前方，掌心序力激荡，OCF层级的防御壁垒刹那成型；牢大新晋ψ(I\_ω)的序数之力全面铺开，鸿蒙色的秩序之光笼罩周身，抵消着外界的邪道压制；而队列中央的霍雨浩双目骤然澄澈透亮，漆黑瞳孔深处，万千精神光点炸裂铺开，周身三道极致魂灵虚影同步显化，三大顶级魂灵之力与双行BMS功法、生灵共生大道瞬间交融，层层铺开稳固全队的共生羁绊，强行对冲来袭的混沌法则，杜绝众人序力脉络被强行扰乱。

悬浮在霍雨浩周身的三道魂灵虚影，气息各异、层次分明，皆是伴随他一路修行、适配生灵大道与序数体系的顶尖存在。最通透柔和的是天梦冰蚕魂灵，通体莹白剔透，周身萦绕淡淡寒雾与细碎精神光点，身躯圆润绵长，自带万古沉寂的空灵气息，作为霍雨浩最早的本命魂灵，它承载着极致精神本源，是所有精神系序力、生灵共情之力的根基。紧随其侧的是冰碧帝皇蝎魂灵，通体冰蓝璀璨，蝎身棱角凌厉，鳞片泛着冰封万物的冷冽光泽，周身流转极致极寒序力，自带绝对冰封的法则特性，肉身防御与攻坚能力冠绝同级。最后一道虚影缥缈灵动，是雪帝魂灵，身形似少女凌空而立，白衣胜雪、气韵超然，周身飘落细碎冰晶，执掌极致冰雪领域与天道极寒奥义，是霍雨浩领域类魂技的核心载体。

三道魂灵各司其职、完美联动，瞬间催动全系魂技，适配当下的绝境防御。霍雨浩眉心灵眸武魂光芒大放，率先激活天梦冰蚕第一魂技·精神探测共享，无边精神脉络瞬间铺展整片山谷，将敌方所有人的站位、攻击轨迹、BMS序列运转节奏、序力波动破绽尽数精准捕捉，毫无遗漏同步共享给八名队友。不同于普通精神探查，依托双行BMS功法加持，此番探测直接穿透敌方混沌遮蔽，解析出对手0 111层级序列的迭代漏洞，为众人拆解攻势、规避杀招提供精准指引。

紧接着，他瞬发精神干扰魂技，无形精神浪潮席卷四方，针对性冲击帅气小队众人的识海，扰乱其序列推演节奏与序力运转轨迹。寻常精神干扰只能撼动心神，可此刻结合生灵同源大道，这道技能能够直接扭曲对手高阶BMS序列的展开逻辑，让敌方0 111层级的递归迭代出现短暂卡顿，压制其混沌拆解特性，稍稍削弱对方的碾压攻势。

危机感极致攀升之际，霍雨浩眸光一凛，冰蓝光芒暴涨，冰碧帝皇蝎护体魂技·冰皇护体瞬间全开。一层细密坚韧的冰晶铠甲覆盖自身与全队众人周身，冰晶纹路暗藏双行BMS序列结构，完美适配正道序数法则，专门抵御邪道混沌序力的侵蚀拆解。冰晶铠甲每一寸都流转极寒序力，既能物理缓冲绝杀攻势，又能冻结来袭的混沌能量，防止众人序力脉络持续紊乱、根基受损。

为彻底稳住濒临破碎的阵型闭环，霍雨浩抬手凌空一握，雪帝领域魂技·帝寒天·雪舞耀阳轰然展开。无边冰雪领域以他为中心扩散，凛冽寒风与纯净冰晶充斥整片防御区域，形成独立的极寒法则力场。领域之内，邪道序力的压制效果大幅削弱，敌方序列迭代速度被迫放缓，而己方众人的序力恢复速度、法则稳固程度大幅提升，雪帝的极寒奥义与蔡须坤的魂灵共生大道完美契合，进一步加固了【天地-修士-魂灵】三位一体的闭环，硬生生稳住了即将崩塌的阵型。

可差距，一目了然。

帅气小队此番来袭的精锐，无一弱者，全员突破双行BMS极限，半数人手掌握0 111 221以上的高阶序列，为首小队长更是稳稳拿捏0 111 222的恐怖战力。这套层级的序列力量，早已脱离了普通OCF折叠与双行递归的范畴，每一次出手都带着混沌无序的碾压特性，完美克制果糕小队尚未完全稳固的新生大道。

``动手，不留活口！''

小队长一声冷喝，大手骤然下压。

数道漆黑的序力利刃破空袭来，带着扭曲天地的混沌气息，狠狠劈向众人的防御屏障。那些利刃之上，BMS高阶序列疯狂流转，0 111的三行雏形力量初显锋芒，每一道攻击的递归深度、力量倍率，都远超果糕小队的理解范畴。

``全员结阵！稳固魂灵共生羁绊！''蔡须坤厉声大喝，瞬间调动圆满的魂灵大道，周身生灵灵气汇聚，与霍雨浩的精神脉络、众人的序力融为一体，形成三位一体的防御闭环。

小舞身形一闪，柔骨羁绊全力展开，化作万千柔韧光带，缠绕加固整个阵型；奥斯卡瞬间催动食物系序力，源源不断的增幅buff倾泻而出，补足众人战力损耗；卢本韦、诺当二人守住侧翼，运转研习的序数展开规则，层层拆解对手的攻击序列。

轰然巨响响彻山谷！

正邪序力剧烈碰撞，漫天能量风暴疯狂肆虐，地面草木尽数碾碎，山石崩裂纷飞。果糕小队凭借完美的大道闭环、默契至极的配合，硬生生扛住了第一轮绝杀突袭。

但压制，已然无处不在。

帅气小队的高阶序列攻击，自带混沌拆解特性，专门克制双行BMS与I序数折叠的正道体系。每一次碰撞，众人周身的序力都会剧烈震颤，稳固的序数脉络被强行打乱，刚刚成型的I序境界根基不断遭受冲击。霍雨浩全程催动魂灵之力维系防御，天梦冰蚕持续消耗精神本源，不间断修正众人紊乱的识海与序列脉络；冰碧帝皇蝎的冰晶铠甲不断被混沌之力击碎、又不断重组，极寒序力持续透支；雪帝领域的范围不断收缩，法则之力逐渐衰弱，三道魂灵虚影愈发黯淡，他的精神力与序力已然濒临透支。

``噗！''

``不要硬接！对方的序列迭代速度，比我们双行体系快数倍！''牢大沉声嘶吼，眼底满是凝重。他的I序数折叠已然是当下正道顶尖突破，可面对三行BMS雏形的碾压，依旧显得捉襟见肘，每一次抬手反击，都要耗费数倍气力，肉身与序力根基承受着巨大负荷。霍雨浩见状，咬牙催动残力，瞬发灵魂冲击，凝练极致的精神光束直刺一名敌方精锐识海，强行打断其序列展开，为牢大的反击创造转瞬即逝的契机。

战斗瞬间进入白热化，却从一开始就是单方面的压制。 以下是衔接在''所有攻击都精准锁定阵型核心------补齐魂灵大道、掌控共生桥梁的蔡须坤。``之后的续写段落：

\begin{center}\rule{0.5\linewidth}{0.5pt}\end{center}

他们无比清楚，蔡须坤是整套终极魂灵体系的关键，是九人阵型的枢纽，只要将她斩杀，果糕小队的三位一体闭环瞬间崩塌，全员战力都会大打折扣。

三道裹挟着混沌序力的杀招，避开所有人的阻拦，刁钻无比地直刺蔡须坤要害，气息阴冷致命，封死了她所有闪避空间。

``休想伤她！''

唐三怒喝一声，身形瞬闪，挡在蔡须坤身前，掌心万千序力壁垒层层叠加，将Ω\_ω层级防御催动到极致。轰鸣巨响中，他硬生生扛下两道绝杀攻击，手臂青筋暴起，虎口炸裂出血，整个人被震得连连后退，脚下地面踏出数道深深的裂痕。

几乎同一时间，小舞、奥斯卡、诺当三人同时挺身而出，舍弃自身攻势，全力护住阵型中路。小舞柔骨缠身，以肉身硬抗剩余一道杀招，娇嫩的脊背瞬间被混沌之力撕裂出一道狰狞伤口，鲜血浸染衣衫；奥斯卡透支本源，强行催动高阶增幅屏障，面色瞬间惨白如纸；诺当咬牙推演反向序列，以自身识海负荷为代价，弱化对手攻击威力。

趁着对手攻势暂缓的间隙，牢大抓住战机，周身鸿蒙序力暴涨，全新的I序数体系全力爆发，ψ(I\_ω)的境界威压轰然席卷，一拳轰出，序数折叠之力层层推进，瞬间击溃两名帅气小队外围精锐的序力屏障，将其肉身直接震碎。

短短数息交锋，果糕小队拼死反击，成功斩杀两名敌方高手，击碎三道绝杀攻势，守住了阵型核心。可敌方后续追兵源源不断，数名身披黑袍的精锐踏着混沌序力极速突进，死死咬住小队突围的路线，眼看就要冲破侧翼防线，将众人截杀在山谷之中。

绝境之中，蔡须坤骤然停下脚步。

她浑身浴血，魂灵共生大道的闭环已濒临崩碎，识海中每一寸脉络都在剧烈震颤。但她的眼神却在这一刻变得无比清明，仿佛所有的恐惧、痛苦、疲惫都在一瞬间被焚烧殆尽，只剩下一团纯粹的、不计代价的决绝。

``须坤！不要------！''牢大瞬间读懂了她的意图，嘶声暴喝，伸手要去拉她。

可蔡须坤已经动了。

她双手合十，十指交错结印，周身残存的所有魂灵本源在这一刻被毫无保留地引爆。天地之间，无数生灵的虚影从虚空中浮现------飞禽走兽、草木虫鱼、上古遗种、远古神兽，万千生灵的气息如百川归海般涌入她的体内，与她的魂灵大道融为一体。

那是她从未在人前展露过的底牌，是她以自身性命为引、以魂灵共生大道为炉，将【天地-修士-魂灵】三位一体闭环彻底逆转、坍缩、压缩成一点之后，才能催动的绝杀一击。

``魂灵归一------万灵寂灭弹！''

蔡须坤厉喝出声，双手猛然前推。

一颗巨大的光球在她掌心凝聚成形。那不是普通的光球------它的表面流转着层层叠叠的序数纹路，双行BMS序列在其内部疯狂迭代，I序数折叠的力量将其递归深度推向了一个前所未有的极限。光球每膨胀一寸，周围的虚空便扭曲一分，地面龟裂下陷，空气被压缩成肉眼可见的涟漪，整片山谷都在这一击面前瑟瑟发抖。

\textbf{ψ(ψ\_I(0))}------这个序数就是普通OCF的极限，在这一刻被蔡须坤以燃烧性命为代价，硬生生堆叠而出。

那已经不再是双行BMS能够描述的范畴，而是触及了I序数自指的递归深渊。光球绽放的光芒刺穿云层，将半边天际映成惨白，呼啸的风压将方圆数里的草木尽数碾成齑粉。

果糕小队众人被这股力量震得齐齐后退，每个人眼中都浮现出难以置信的神色------他们从未想过，蔡须坤竟能催动如此恐怖的力量。

``须坤\ldots\ldots{}''小舞眼眶泛红，她清楚地看到，蔡须坤的肌肤正在寸寸龟裂，鲜血从每一道裂缝中渗出，她的生命力正随着这颗光球的成型而飞速流逝。

也是她最强的一击。

蔡须坤嘴角溢出血丝，却露出了一个释然的笑容。她回头望了牢大一眼，那一眼里没有恐惧，没有遗憾，只有无尽的温柔。

``走。''她用尽最后的力气，轻声说出这个字。

巨大的光球脱手而出，化作一道毁天灭地的光柱，裹挟着ψ(ψ\_I(0))的恐怖威压，直直轰向帅气小队小队长！

那一瞬间，天地失色。

光柱所过之处，虚空被撕裂出漆黑的裂痕，大地被犁出深不见底的沟壑，连空间本身都在这一击面前发出痛苦的哀鸣。这是果糕小队自组建以来最强的一击，是蔡须坤以命相搏、以魂灵大道为代价换来的终极绝杀------

然后，小队长伸出了一只手。

没有序列纹路的暴涨，没有全力运转的架势，他甚至没有改变站立的姿态。他只是平平淡淡地伸出了右手，五指微张，像是要接住一片飘落的树叶。

光柱撞击在他的掌心。

没有爆炸。没有轰鸣。没有能量对冲的滔天巨浪。

ψ(ψ\_I(0))的恐怖力量，那道凝聚了蔡须坤全部生命、燃烧了魂灵共生大道一切根基的绝杀一击------在接触到他掌心的瞬间，就像一滴水落入了无底深渊。

无声无息地，消散了。

光球表面疯狂流转的序数纹路一寸寸黯淡、熄灭，递归迭代戛然而止，I序数折叠的递归深渊被一只无形的手从根源上抹平。那股足以撕裂虚空的力量，连小队长掌心的皮肤都未能撼动分毫，便如冰雪遇骄阳般消融殆尽，化为漫天飘散的光点，缓缓坠落。

小队长收回手，轻轻弹了弹指尖残留的光尘，动作随意得像是在拂去一粒灰尘。

``ψ(ψ\_I(0))？''他语气平淡，不过是0 111 211 31 2甚至带着一丝无聊，``在0 111 222面前，你连让我认真拆解的资格都没有。三行序列对双行体系的碾压，不是靠燃烧性命就能弥补的。这是维度的差距。''

蔡须坤的身体猛地一僵，瞳孔中最后的光芒缓缓黯淡。她张了张嘴，似乎想说什么，却只吐出一口鲜血，整个人如断线的风筝般向后倒去。

``须坤！！''

牢大疯了一般冲上前，将她接入怀中。

她靠在牢大怀中，嘴角还挂着那抹释然的笑，目光涣散地望着天空，嘴唇微微翕动。

``\ldots\ldots 没用的牢大\ldots\ldots 他的层级太高了\ldots\ldots{}''

声音细若蚊蚋，却像一把钝刀，一下一下割在在场每个人的心上。

她拼尽全力的一击，甚至没有让对方后退半步。

帅气小队长负手而立，居高临下地俯视着这群浑身浴血的少年，目光中没有愤怒，没有赞赏，甚至没有轻蔑------只有漠然。那种面对蝼蚁奋力挣扎时的漠然。

``我说过，''他缓缓开口，声音不大，却清晰地传入每一个人的耳中，``你们的大道圆满、序数突破，在三行BMS的力量面前，不过是孩童戏法。''

他顿了顿，目光扫过牢大怀中重伤的蔡须坤，语气依旧平淡如水：

``这一击，算是你们今天\ldots\ldots 唯一让我记住的东西。''

绝境生死一刻，卢本韦骤然沉声大喝：``奥斯卡，借你本命香肠！联动合击！''

``收到！极致增幅·本命爆裂香肠！''

奥斯卡眸光一凝，透支仅剩的本源序力，掌心瞬间凝聚出一根通体鎏金、缠绕双层BMS序列纹路的巨型增幅香肠。这并非普通辅助魂技，而是他打磨许久、适配序数大道的顶级本命造物，内含极致压缩的生灵序力与爆裂势能，是他压箱底的攻坚杀招。

卢本韦身形瞬动，精准接手鎏金香肠，同时转头看向正在正面抵挡攻势、周身萦绕鸿蒙序力的牢大。无需多言，牢大瞬间洞悉其意，濒临透支的I序数之力毫无保留倾泻而出，ψ(I\_ω)的高阶序数折叠本源层层叠加、灌注而入！

I序终极序数的递归特性、双行BMS的并行迭代逻辑、再加上奥斯卡的极致增幅本源，三者瞬间拧成一股前所未有的秩序洪流！卢本韦握持鎏金香肠，周身序数纹路疯狂闪烁，他摒弃常规序列拆解，强行将牢大灌注的I序折叠力量、奥斯卡的本源爆裂势能压缩收敛，完成绝杀前置蓄力。

``融合技------香肠暴击！''

两道吼声同步炸响，划破山谷硝烟！奥斯卡指尖魂力暴涨，彻底解锁本命香肠的终极投射特性，不再是近身增幅、原地爆发，而是化作一道凝练到极致的鎏金流光，裹挟着层层叠叠的双行BMS迭代序列，直接射入卢本伟发射的母猪鼻子里面，之后母猪大哼一声，两个鼻孔被巨大的香肠填满，之后母猪把香肠直接喷射出去，这一次，不仅仅是两个人的融合技能，是三个人的，牢大用力肘击猪屁股，将母猪直接击飞，母猪受到刺激，鼻孔直接张开，将大香肠发射出去，之后，香喷射出去，由于反作用力弹了回来。

喷射的大香肠不再局限于普通增幅爆裂，而是牢牢承接牢大透支灌输的ψ(I\_ω)序数根基，完成层级破壁跃升！无数递归序列在冲击波内部重构迭代，冲破原有双行BMS与I序的层级桎梏，硬生生堆叠出ψ(K\_ω)的恐怖爆发力量！

K序数超脱常规Ω迭代与I序折叠的递归上限，是凌驾于ψ(I\_ω)之上的全新临界层级，此刻经由三人完美联动彻底引爆。

璀璨鎏金洪流席卷前路，序力碾压波纹瞬间扩散整片山谷，冲在最前方的数名帅气小队精锐甚至来不及催动0 111 211层级的序列防御，体内混沌BMS序列瞬间被强行拆解、崩碎、湮灭。

连惨叫声都未曾响起，数名高阶精锐的身躯连同其运转一生的序数根基、混沌道韵，尽数被ψ(K\_ω)的极致秩序之力碾为虚无，前路追杀的包围圈瞬间被撕开一道宽大的缺口！

滔天烟尘翻滚四溢，短暂的强势爆发，让紧绷的死局骤然松动。

可这份惊艳绝伦的绝杀突破，仅仅持续了一瞬。

``区区伪临界K层级，也敢班门弄斧？''

冰冷淡漠的嗤笑轰然压过爆鸣，山谷上空的邪气骤然凝固、暴涨！原本伫立后方的帅气小队小队长，周身黑袍无风狂舞，原本0 111 221的序列纹路瞬间暴涨迭代，彻底解锁全部战力，深邃漆黑的序数光芒笼罩天地，完整的BMS：0 111 222恐怖层级轰然现世！

相较于众人临时爆发的残缺ψ(K\_ω)，0 111 222的三行完整序列，是维度与递归层级的全方位碾压。这是彻底成型的高阶混沌序列，远超双行体系、临时K层级爆发的上限，整片山谷的正道序力瞬间被死死压制，刚刚铺开的鎏金冲击波以肉眼可见的速度萎缩、崩解、消散。

恐怖的压制力场轰然落下，笼罩果糕小队全员。

噗------噗------噗！

唐三、牢大、卢本韦、奥斯卡四人同时气血翻涌，大口呕出鲜血。刚刚透支本源打出绝杀的几人体力、序力、识海全部濒临枯竭，牢大的I序数脉络剧烈震颤，险些崩碎；卢本韦与奥斯卡的融合技本源彻底耗尽，浑身魂力一空，面色惨白如纸，脱力身形踉跄欲坠。

霍雨浩强忍识海剧痛，瞬间收拢所有魂灵力量，将雪帝领域、精神屏障全部收缩贴身，死死护住全队残存的生机，低声急喝：``留力！全部保留本源，不要反击！准备突围！''

众人瞬间心神会意，强行压下体内翻涌的气血与剧痛，放弃所有后续攻势，不再消耗一丝一毫多余力量。所有人默契封存仅剩的序力、精神本源与魂灵底蕴，将全部体力集中于身法突围，不再与0 111 222层级的恐怖力量正面对抗。

小舞死死扶住脱力的奥斯卡，蔡须坤强撑受伤之躯，催动最后一丝魂灵共生之力串联全队，稳固众人濒临溃散的身形与脉络。

``走！''

唐三咬牙爆喝，携着牢大一同顶在最前，以残破的防御壁垒硬扛队长溢出的混沌威压，掩护全队顺着刚刚撕开的缺口，带着一身重伤与残存气力，向着星斗大森林深处全力遁逃。

战力层级的绝对差距，绝非默契与意志可以弥补。

为首小队长冷眼看着战死的两名手下，脸上没有丝毫波澜，只有愈发浓郁的杀意：``负隅顽抗，徒劳无功。你们的大道圆满、序数突破，在三行BMS的力量面前，不过是孩童戏法！''

话音落下，他不再留手，周身0 111 222的高阶序列全力运转，整片山谷的邪道气息瞬间暴涨，一股碾压全场的恐怖力量轰然落下。

狂风肆虐，邪气滔天。

果糕小队众人齐齐被震退，人人带伤，气血动荡不止。唐三的防御壁垒布满裂痕，随时可能崩碎；牢大的I序数之力持续消耗，后继乏力；霍雨浩浑身汗透，识海阵阵刺痛，三道魂灵虚影近乎透明，精神脉络饱受混沌气息持续压制，生灵共情之力大幅衰减，维系全队的共生羁绊已然濒临断裂；蔡须坤的魂灵共生大道剧烈震颤，闭环体系濒临破碎。

绝境之中，霍雨浩强行压下识海剧痛，燃烧少量精神本源，催动高阶组合魂技，以精神群体虚弱叠加雪帝领域之力，笼罩所有敌方人手。双重削弱效果同步生效，既压制对手的序力增幅，又迟缓其序列展开速度，为全队突围争取关键窗口。

再缠斗下去，全员必死无疑。

``撤！立刻突围！''唐三当机立断，沉声嘶吼，``所有人聚拢，护住须坤，向东突围，直奔星斗大森林深处！''

没有丝毫犹豫，众人瞬间收拢阵型，以唐三、牢大为锋，霍雨浩居中调度，死死护住受伤的蔡须坤，摒弃所有反击，全力朝着山谷东侧突围。

帅气小队众人紧随追击，高阶序力铺天盖地，一道道绝杀技能不断轰击在众人身后，每一次冲击都让小队阵型剧烈晃动，众人伤口撕裂，血水不断滴落，染红了沿途的青草山石。

惨烈的逃亡死斗，正式拉开序幕。

与此同时，万里之外的魂殿总殿议事大殿。

此前众人的所有注意力，都被西线连绵的战事牵制。帅气教大军佯攻斯兰登城，攻势凶猛、死战不退，死死牵制住魂殿全部主力，营造出全线压境、即将破城的假象，让所有人都误以为这是敌方的主战场。

一直静默伫立的ddfg瞬间睁眼，眸光锐利如锋，沉声暴喝：``骗局！全盘都是调虎离山之计！''

梅天狸身躯一震，瞬间洞悉所有阴谋！

西线的凶猛攻势，从来都不是为了攻城掠地，只是为了牵制魂殿主力，拖延正道驰援的脚步！帅气教真正的目的，从来不是西线战场，而是悄然崛起、补齐万年大道、诞生全新序数体系的果糕小队！

``好毒辣的算计！''梅天狸声音冰冷，眼底满是怒意与焦灼，``他们忌惮正道新星崛起，忌惮双行、三行BMS与I序数体系彻底成型，不惜倾尽西线兵力演戏，只为围杀一众少年！''

果糕小队众人衣衫染血、气力透支，每一步前行都耗尽全身力气。身后的追杀从未停歇，高阶序列的碾压之力如影随形，死亡的阴影牢牢笼罩在众人头顶。霍雨浩全程透支魂灵本源，天梦冰蚕持续解析追兵攻击轨迹，不间断共享预判信息；冰碧帝皇蝎残力护住众人后路，冻结部分追击的绝杀攻势；雪帝领域收缩成贴身屏障，抵挡散落的混沌余波，以近乎枯竭的力量，死死护住全队生机。

但没有一人退缩，没有一人放弃。

牢大始终挡在最前，以肉身硬抗绝大多数攻势；他不断透支序力，用I序数折叠之力开辟突围通道；霍雨浩倾尽精神本源，梳理众人紊乱的法则，维系最后的共生羁绊；其余众人彼此扶持、拼死抵挡，带着满身伤痕，向着生机所在的星斗深处，奋力狂奔。

生死一线之间，远方天际，四道璀璨正道霞光，正破开云层，急速赶来。

星斗大森林外围，血染荒草，风卷残烟。

果糕小队一行人带伤狂奔，衣衫尽数被血水浸透，每个人的序力都濒临枯竭。唐三顶着残破的防御壁垒在前开路，霍雨浩魂灵虚影黯淡无光，精神屏障堪堪护住全队法则脉络，奥斯卡与卢本伟透支本源过后气息奄奄，小舞、诺当死死护住中间的蔡须坤，一行人拼尽残余气力，朝着星斗深处的密林生机狂奔。

可身后的追杀碾压，从未断绝。

再这样下去，全员都会被追上围剿，无人能够幸存。

狂奔途中，牢大骤然驻足。

他周身原本璀璨的鸿蒙秩序之光早已黯淡大半，ψ(I\_ω)的序数根基布满裂痕，浑身经脉断裂无数，血水顺着四肢不断滴落，沉重的呼吸在风声中格外清晰。但他脊背依旧挺拔，如同擎天立柱，死死挡在所有人的身后。

所有人脚步一顿，回头望去，眼底皆是骤然的心悸与不安。

牢大目光越过众人，牢牢落在被众人护在中心、身受重伤、气息虚弱的蔡须坤身上，眼底所有的凌厉杀伐尽数化作温柔与决绝。自小队成立以来，蔡须坤便是魂灵大道的基石，是全队的软肋，也是他唯一想要拼尽一切守护的人。

如今绝境横亘，追兵势不可挡，唯有断后，方能为生路

\chapter{第四十二章：牢大坠机，曼巴殉道}\label{ux7b2cux56dbux5341ux4e8cux7ae0ux7262ux5927ux5760ux673aux66fcux5df4ux6b89ux9053}

\begin{center}\rule{0.5\linewidth}{0.5pt}\end{center}

星斗大森林外围，荒草染血，朔风卷泣。

天穹低垂如铅，压得人喘不过气。残阳最后一缕余光被乌云吞噬，整片密林陷入昏黑，唯有遍地残骸上未干的血迹，在暮色中泛着暗红的微光。空气中弥漫着焦灼的灵力余味与泥土翻涌的腥气，风穿过枯枝的呜咽声，像是大地在为即将发生的一切提前哀鸣。

刚刚挣脱香肠暴击绝杀余波的果糕小队，全员带伤，气力枯竭。唐三顶着龟裂的序力壁垒艰难开路，每一步落脚都在地面碾出深深的裂痕，壁垒上的裂纹如蛛网般不断蔓延，随时可能彻底崩碎；霍雨浩周身魂灵虚影黯淡欲灭，精神屏障勉强维系全队脉络，他的太阳穴青筋暴起，鼻血早已干涸成两道暗红的痕迹；奥斯卡与卢本伟透支本源，气息奄奄，每呼吸一次都伴随着胸腔内骨骼错位的闷响；小舞与诺当贴身护住重伤的蔡须坤，她的裙摆被鲜血浸透，却依旧死死将蔡须坤护在怀中，不让一丝杀机靠近。

一行人拼尽最后余力，向着星斗密林深处仓皇遁逃。

可身后的追杀阴霾，从未散去半分。

帅气小队裹挟着 \textbf{0 111 222} 三行完整BMS混沌序列的恐怖威压，踏碎长空紧追不舍。漆黑的混沌邪气凝作万千杀刃，死死锁定众人气息，高阶序列的维度碾压之力如万丈悬山镇压而下，让每一名队员的递归脉络剧烈震颤，气血逆流，逃亡的步伐愈发沉重艰难。

所有人都心知肚明------以当下全员重伤、战力尽失的状态，一旦被追兵合围，果糕小队必将全员覆灭。新生的正道序数大道、传承的魂灵共生体系，都将在此彻底断绝。

死亡的阴影，已经贴上了每一个人的后颈。

\begin{center}\rule{0.5\linewidth}{0.5pt}\end{center}

千钧一发之际，一道挺拔的身影骤然驻足，硬生生停在了逃亡队伍的最后方。

\textbf{是牢大。}

他停下的那一刻，像是一柄长枪扎入大地，纹丝不动。

他周身原本璀璨的鸿蒙秩序神光黯淡大半，此前强行催动 ψ(M\_ω) 爆发的后遗症彻底爆发------浑身经脉寸寸崩裂，皮肉外翻，血水顺着四肢不断滴落，在脚下汇成一汪暗红的水洼。原本稳固的 ψ(I\_ω) 序数道基布满蛛网裂痕，濒临溃散，每一次心跳都伴随着道基碎裂的细微声响，像是瓷器在烈火中一寸寸龟裂。

极致的虚弱与剧痛蚕食着他的肉身，可他的脊背依旧挺拔如松，宛若擎天之柱，死死挡在所有队友身前，挡在漫天杀机之前。

风卷起他残破的衣袍，猎猎作响，像一面千疮百孔却不肯倒下的战旗。

众人脚步骤停，回头望去，眼底尽数涌上惊惧、不舍与沉痛。

没有人说话。所有人都知道他要做什么，却没有一个人愿意先开口。

\begin{center}\rule{0.5\linewidth}{0.5pt}\end{center}

牢大的目光越过所有人，最终牢牢落在气息虚弱、面色苍白的蔡须坤身上。

那双在杀伐战场上凛冽凌厉、令敌人闻风丧胆的眸光，瞬间柔化，只剩满心的牵挂与决绝------像是暴雪夜里最后一盏烛火，微弱，却倾尽所有地在燃烧。

自小队组建以来，蔡须坤以柔弱之躯撑起全队魂灵共生大道，默默滋养众人序力脉络，分担所有厮杀苦痛。她是他在刀光剑影中回头时唯一能看见的温柔，是他满身血污时唯一不舍得弄脏的手，是他乱世征途里唯一的救赎，也是他拼尽一切想要守护的毕生执念。

前路生路渺茫，后有绝世杀机。

唯有断后，可护全员周全。

他比任何人都清楚这一点。

\begin{quote}
``所有人听我命令，即刻全速撤退！''
\end{quote}

牢大的声音沙哑低沉，像是砂纸碾过生锈的铁片，却带着小队队长不容置喙的绝对威严，穿透呼啸风声，砸在每人心头。

\begin{quote}
``直奔星斗核心腹地，保存正道火种，护住蔡须坤------不许回头，不许停留！''
\end{quote}

最后四个字，他咬得极重。

唐三咬牙怒吼：``牢大！要走一起走！对方 0 111 222 序列战力碾压一切，你孤身断后必死无疑！''

``我是队长。服从命令。''

牢大轻轻摇头，语气坚定得没有一丝松动。他没有再看唐三，而是转头望向泪眼婆娑的蔡须坤。

那一刻，漫天杀机、遍体鳞伤、油尽灯枯------统统从他脸上褪去。他看着她，像是看着此生最后的落日，声音放得极柔，柔得不像是一个即将赴死之人，倒像是临行前在门口替她掖好被角的寻常丈夫------

\begin{quote}
``放心，我有脱身之法，一定能追上你们。''\,''
\end{quote}

这是谎言。

所有人都心知肚明。

面对成型的三行混沌序列强者合围，孤身一人，绝境断后，本就是九死一生的死局。可局势已然没有半分退路，小队全员重伤枯竭，无人再战，唯有牺牲一人，方能为所有人搏出一线生机。

蔡须坤泪水滚落，虚弱的声音哽咽颤抖：``我不走，我要和你一起\ldots\ldots{}''

``听话。''

牢大眸光温柔却决绝。他抬起满是血污的手，似乎想要替她擦去泪水，手伸到一半，又缓缓放下------他怕自己手上的血，弄脏了她的脸。

\begin{quote}
``活下去，带着我们的大道，好好活下去。''
\end{quote}

霍雨浩攥紧双拳，指甲刺入掌心，鲜血顺着指缝滴落。识海清晰感知到牢大视死如归的决然，他强忍悲痛沉声开口，号令全队撤退。

众人望着那道孤身伫立、逆对杀机的背影，满心悲壮不舍，终究只能咬牙转身，倾尽最后气力，向着密林深处极速遁去。

蔡须坤被小舞拖着前行，她不断回头，视线模糊得看不清那道身影，只知道那道背影越来越小、越来越远，最终被密林的黑暗彻底吞没。

她不知道的是，那道背影始终没有动过，一直面朝她的方向，直到她的身影彻底消失在林间。

\begin{center}\rule{0.5\linewidth}{0.5pt}\end{center}

风声萧瑟，转瞬之间，林间天地，只剩牢大一人。

他缓缓吐出一口带血的气，肩膀微微塌了一下------那是只有独处时才会流露的疲态。但仅仅一瞬，他的脊背便重新挺直，像一柄钉入大地的长枪。

\textbf{身前是覆压天地的绝世死敌，身后是挚爱与队友的唯一生路。}

他缓缓抬起双臂，摆出迎战的架势。残破的衣袖垂落，露出布满裂痕的小臂，血珠顺着指尖一滴一滴坠入泥土。

他在等。

等她们走远一些，再远一些。

远到他死了，她们也听不见。

\begin{center}\rule{0.5\linewidth}{0.5pt}\end{center}

轰隆------！

漆黑邪气轰然落地，大地震颤，云雾翻涌。帅气小队队长携副队长、四名精锐全员合围而至，层层混沌序列力场封锁整片天地，彻底封死所有闪避退路，绝杀之局，彻底成型。

黑袍翻飞间，帅气小队队长周身 \textbf{0 111 222} 完整三行BMS序列熠熠生辉，漆黑的维度威压碾压四方，眼神冰冷戏谑，带着对正道修士的极致轻蔑：

\begin{quote}
``为了情爱与队友，甘愿自蹈死地？正道的执念，果然愚蠢可笑。既然你主动赴死，今日我便成全你，让你的大道与性命，一同湮灭于此！''
\end{quote}

四名精锐分列四方，\textbf{0 111 21} 层级的混沌序列交织成锁杀灭阵，空间层层折叠挤压，封死一切闪避角度；副队长伫立侧方，\textbf{0 111 221} 的高阶序列蓄势待发，凛冽杀机死死锁定牢大，只待总攻号令。

无边绝境，彻底笼罩。

牢大巍然伫立，孤身对千钧杀机，无半分退避之意。肉身崩裂剧痛难忍，序力濒临枯竭，道基摇摇欲坠，寻常修士早已心脉崩断、弃械臣服，可他的战意，却在极致死局中愈发炽烈------

像是暴风雪中最后一簇篝火，风越大，烧得越烈。

\textbf{置之死地，而后生！}

濒临神魂湮灭的刹那，他的心神彻底空明通透。过往苦修的I序数折叠、Ω层级递归迭代、双行BMS并行逻辑、M基数临界奥义，无数晦涩深奥的天地法则，在他识海之中疯狂交织、碰撞、贯通、升华。

像是混沌中劈开的第一道闪电。

I序桎梏破碎，M序临界突破，一尊凌驾于二者之上的全新高阶序数基数，轰然成型！

超脱寻常递归、碾压双行体系的K序终极序数，绝境破壁，现世降临！

鸿蒙色的秩序神光骤然自他残破的身躯中炸裂暴涨，崩碎的道基被强行重塑，枯竭的序力迎来短暂的极致复苏，苍茫霸道的秩序道韵席卷四方，硬生生将周遭的混沌邪气层层震退。

那一刻，他周身绽放的光芒，照亮了整片昏暗的星斗密林------像是黑夜中骤然升起的第二颗太阳，炽烈、滚烫、不计后果地燃烧着。

``哦？绝境突破，倒是有些天赋。''队长眼底掠过一丝讶异，随即尽数化为冰冷蔑视，``可惜，区区临时成型的K序雏形，根基虚浮，底蕴浅薄，也敢抗衡完整的 0-111-222 三行序列？螳臂当车，不自量力！''

\begin{quote}
``全员出手，速斩此子！''
\end{quote}

杀令落下，天地皆寂。

四名精锐同时催动混沌奥义，漆黑的序列杀招漫天袭来，层层叠叠的递归链条绞杀而至，欲将这位绝境突破的正道天才瞬间碾碎。

牢大眼底赤红，战意焚天，明知战力悬殊、结局已定，依旧逆杀而上！

他身躯骤然旋拧，腰背筋骨爆发出阵阵脆响，全身凝练的K序神力尽数汇聚于肘锋，空间秩序被强行撬动、改写！

\begin{quote}
\textbf{``肘翻天地！''}
\end{quote}

一声低喝震彻山谷！凌厉肘势横空碾压，硬生生斩断迎面袭来的混沌序列，暴力拆解敌方递归公式。一名率先扑杀的精锐猝不及防，护体序列瞬间崩碎，体内气血逆流，整条手臂骨骼炸裂，整个人如遭神山撞击，惨叫着横飞十余丈，落地重伤晕厥，彻底失去战力。

余下三人见状，怒极反扑，左右夹击、后方偷袭，漆黑序列利爪撕裂长空，夺命针雨封锁八方，招招直指神魂要害。

牢大K序领域全开，极致的秩序感知洞悉所有杀机轨迹。他踏残影、避绝杀，在漫天夺命攻势中辗转腾挪，衣袍被气劲撕碎，皮肉再添无数血痕，却始终护住核心道基，未露半分破绽。

闪避刹那，他沉身扎根大地，双拳凝聚毕生底蕴与K序终极神力，轰然砸落地面！

\begin{quote}
\textbf{``曼巴重锤！''}
\end{quote}

轰鸣巨响震彻四野，大地崩裂百丈裂痕，磅礴浩瀚的秩序洪流顺着地脉喷涌炸开，化作环形巨浪清空全场混沌之力。两名夹击精锐的护体屏障脆如薄纸，身躯被巨浪狠狠掀飞，经脉尽碎、道基受损，坠地之后抽搐不止，再无起身之力。

最后一名后方精锐趁机偷袭，漆黑序列长矛裹挟破空之势直刺后心，欲破其功、断其道基。牢大不闪不避，以残破肉身硬抗绝杀，长矛刺穿皮肉、渗血入骨，剧痛贯穿全身，可他道基稳固分毫未损。

他低头看了一眼贯穿后背的长矛，嘴角竟微微上扬------像是嘲笑对手的无能，又像是早已将疼痛视作日常。

趁敌方招式滞涩、力道穷尽的瞬间，牢大反手一肘悍然撞出！

咔嚓！

骨裂巨响响彻山谷，那名精锐肋骨寸断，脏腑震裂，惨叫倒飞，彻底溃败。

\textbf{瞬息之间，四名合围高阶精锐，尽数被牢大孤身击溃！}

山谷烟尘弥漫，邪气散尽，满地狼藉。

牢大独自伫立战场中央，浑身浴血，伤痕累累，后背贯穿伤口汩汩淌血，刚刚成型的K序神光已然黯淡斑驳。强行破壁的 ψ(K\_ω) 本就是无根之木、无源之水，经此极致鏖战，序力彻底透支，道基裂痕飞速蔓延，濒临彻底崩塌。

他大口喘息，口鼻溢出血沫，视线恍惚眩晕，体力、魂力、神魂本源尽数抵达极限。

可他不能倒。

他咬着牙，用长矛撑住身体，不让膝盖触地------因为他知道，一旦跪下，就再也站不起来了。

而真正的死局，才刚刚降临。

``杂鱼尽数落败，该轮到我们了。''

冰冷的话音落下，帅气小队队长、副队长双双压上，残余所有战力尽数集结，无解的合围杀阵彻底成型。\textbf{0-111-222} 与 \textbf{0-111-221} 的高阶混沌序列层层叠加，漆黑的毁灭之力笼罩整片天地，维度碾压、法则拆解、层级压制三重杀机并行，这是帅气小队倾尽全员底蕴的 \textbf{终极合击}，只为一瞬抹杀牢大。

漫天混沌洪流倾覆而下，碾压一切秩序，破灭所有生机！

牢大望着那覆压天地的绝杀攻势，心中已然知晓结局。临时突破的K序雏形，根本无法抗衡成型的三行混沌序列合击，此战无解，此命必绝。

可他心中没有半分悔恨，唯有滚烫至死的守护执念。

脑海中翻涌着与蔡须坤并肩同行的朝夕------她温柔的眉眼、隐忍的倔强、默默的付出，还有自己未曾兑现的安稳诺言。他多想活着，多想护她余生安稳，多想陪她看遍山河月明，可乱世杀伐，生死无常，他别无选择。

他是果糕小队的队长，是全队的壁垒；他是蔡须坤的守护者，是她此生唯一的偏爱。

\textbf{身可死，道可灭，唯羁绊与爱意，永不消亡！}

为了队友的生路，为了挚爱的平安，为了正道的火种，他愿燃尽一切，以身殉道！

\begin{quote}
\textbf{``燃烧全部本源生命力------''}
\end{quote}

他的声音已经嘶哑到几乎听不清，却依旧咬字清晰，一字一顿，像是在念诵最后的祷文。

\begin{quote}
\textbf{``奥义·坠机。''}
\end{quote}

心火燃尽，道念归真，这一刻，万千杂念尽数沉淀，只剩他心底最纯粹、最滚烫的独白------

\begin{quote}
\emph{我知道，我今天必死无疑。}

\emph{临时突破的K序终究是镜花水月，抵不住 0-111-222 的滔天大势，抵不住这整支小队的倾巢合击。我撑到现在，早已油尽灯枯，肉身崩裂、道基残缺、神魂欲碎，我的每一寸筋骨都在嘶吼着溃败，我的修行之路，到此为止了。}

\emph{可我一点都不后悔。}

\emph{身为果糕小队的队长，我自当立于万人之前，挡于绝境之后。全队的希望不能断，正道的火种不能灭，那些一路陪我浴血厮杀、并肩悟道的伙伴，绝不能葬送在此地。我是队长，这份牺牲，我不扛，无人可扛。}

\emph{最对不起的，是须坤。}

\emph{我曾在心底无数次许诺，乱世平定，大道安稳，我便卸去一身杀伐，护你岁岁平安，陪你看遍山河落日，让你不必再背负伤痛、不必再直面生死。你温柔坚韧，默默撑起全队的魂灵共生大道，把温柔留给所有人，把苦痛藏给自己。你是我黑暗乱世里唯一的救赎，是我穷尽一生想要守护的温柔执念。}

\emph{我多想活下去，多想兑现诺言，多想余生与你并肩，岁岁年年，安稳相守。}

\emph{可我不能。}

\emph{我若退，全队皆亡，你也难逃一死。我唯一能护你的方式，便是以我之死，换你余生安稳；以我道灭，铺你逃亡生路。}

\emph{爱意无声，殉道无言。}

\emph{今日，我燃尽生命，崩碎道基，散尽神魂，不求盛名，不求留名，只求我的牺牲能拦住漫天杀机，护我的队友全身而退，护我挚爱平安远去。}

\emph{须坤，忘了沙场纷争，忘了乱世疾苦，好好活下去。}

\emph{若有来生，我不求天骄大道，不求绝世战力，只求安稳遇你，平凡相守，护你一生无忧，不负遇见，不负深情。}

\textbf{此生以身殉道，以命护卿，足矣。}

\textbf{无憾，亦无悔。}
\end{quote}

撕心裂肺的怒吼响彻天地，牢大义无反顾引爆自身所有生命精血、毕生修为、K序道基与神魂本源！

熊熊生命之火自他血脉神魂中冲天而起，黯淡的秩序神光骤然极致暴涨，鸿蒙色的终极秩序洪流横贯天地，跨越层级桎梏，凝聚成他此生最强、也是最后的绝杀一击！

这是赌上性命的绝唱，是绝境少年最后的反抗，是守护一切的终极执念！

那一刻，他的身躯在鸿蒙神光中变得近乎透明，血肉、骨骼、经脉、神魂，一切有形之物都在燃烧，化作纯粹的力量达到了,0 111 22的威力。他像一颗坠向大地的流星，拖着长长的光尾，义无反顾地撞向那覆压天地的黑暗。

\textbf{轰隆------------！！！}

秩序洪流与混沌合击轰然对撞，黑白两极极致力量疯狂炸裂，山峦崩摧、云雾撕碎、大地塌陷，整片星斗外围剧烈震颤，暴乱的能量风暴席卷万里空域！

霸道的K序燃命之力疯狂拆解敌方混沌序列，硬生生撕裂全员合击的核心脉络！首当其冲的副队长猝不及防，护体高阶序列瞬间崩碎，恐怖的秩序余力轰然轰击肩头，骨骼炸裂、血肉翻飞，一条手臂直接被废，凄厉的惨叫响彻山谷，重伤倒飞出局。

\textbf{仅凭一己之力，牢大便重创敌方核心战力，撕碎无解合击，创下绝境逆天的杀伐奇迹！}

可奇迹，终究无力逆转死局。

奥义迸发的瞬间，牢大的生命本源彻底燃尽，神魂灯火骤然黯淡，K序道基寸寸湮灭，毕生修为尽数归零。燃烧生命力的恐怖反噬席卷全身，本就残破的肉身瞬间千疮百孔，彻底失去所有战力。

他挡下了合击的正面冲击，重创了副队长，却再也无力抵御残余的混沌绝杀之力。

残存的 \textbf{0 111 222} 层级混沌威压裹挟毁灭一切的力量，狠狠碾压在毫无防备、彻底力竭的牢大身上。

噗------！

滚烫精血狂喷而出，牢大挺拔的身躯剧烈震颤，周身所有秩序光芒彻底熄灭。

极致的绝望席卷神魂，可他最后的心念，依旧是远方逃亡的队友，是他拼尽性命守护的蔡须坤。

他的嘴唇微微翕动，没有发出声音，但识海中最后一缕残存的神念，拼尽全力凝聚成一句话------

\begin{quote}
\emph{我此生无缘与你相守，但我以命为你铺路，护你平安余生。}

\textbf{足矣。}

\textbf{无憾，无悔。}
\end{quote}

下一秒，漫天混沌毁灭之力彻底吞没他的身躯。

骨骼崩碎，肉身湮灭，神魂脉络寸寸消散，刚刚破壁成型的 ψ(K\_ω) 高阶大道彻底崩塌归零。

那道挺拔如松的身影，那双向她望去时永远温柔的眼眸，那双伸到一半又放下的满是血污的手------

尽数化为齑粉，散入星斗密林的风中。

就在他神魂彻底消散的最后一瞬，燃尽生命与序数大道的终极余波轰然炸开！

万里长空剧烈震颤，星斗大森林灵气暴乱、草木倾覆、山峦移位，无数天地法则、序数序列紊乱暴走，形成万年难遇的天骄陨落异象。

\textbf{大道悲鸣，天地同悲。}

远在密林深处仓皇奔逃的果糕小队众人，几乎在同一瞬间，感受到了一股撕裂神魂的剧烈悸动。

霍雨浩脚步猛地一顿，识海深处传来一声尖锐到极致的悲鸣------那是他与牢大神魂共鸣的羁绊脉络，在一瞬间寸寸崩断。每一根脉络断裂的痛楚都清晰无比，像是有人用烧红的铁针一根一根地刺穿他的神魂，可比起肉身之痛万倍的，是那股从灵魂深处涌上来的、无法忽视的讯息------

那道与他并肩作战、以队长之威号令全队的强大神魂，消失了。

不是沉睡，不是昏迷，不是暂时隐匿。

是彻底消散，如同烛火被风吹灭，如同墨迹被水化开，干干净净，不留一丝痕迹。

霍雨浩双膝一软，重重跪倒在地，双手死死抱住头颅，喉咙里发出一声压抑到极致的嘶吼。他的指甲刺穿头皮，鲜血顺着脸颊淌下，可他却浑然不觉，只是反复呢喃着一个字------

``不\ldots\ldots 不\ldots\ldots 不\ldots\ldots{}''

唐三身形剧烈摇晃，头顶龟裂的序力壁垒轰然碎裂，他单膝跪地，一口鲜血喷出三尺远。他感知到了------那道始终走在队伍最前方、替他挡下无数次致命攻击的坚实背影，那股永远沉稳厚重、如山岳般不可撼动的秩序气息，在这一刻，彻底归于虚无。

他猛地回头望向身后那片被漆黑浓云笼罩的密林方向，眼眶赤红，嘴唇剧烈颤抖。

``队长\ldots\ldots{}''蔡须坤说

声音嘶哑破碎，像是从胸腔深处硬生生撕扯出来的。

奥斯卡瘫坐在地，浑身脱力，连抬起一根手指的力气都没有，可他的眼泪却止不住地涌出，无声地滚落，砸在满是血污的衣襟上。卢本伟仰面朝天，喉结剧烈滚动，死死咬住下唇，咬到鲜血淋漓也不肯发出一点声音------他知道，如果此刻崩溃，队长的牺牲就白费了。

小舞浑身颤抖，泪水早已模糊了视线，她下意识收紧了怀抱------

``须坤！须坤！''小舞慌乱地低头看去，只见蔡须坤双目圆睁，嘴唇青紫，气息微弱到几乎感知不到。她的胸口剧烈起伏，每一次呼吸都像是在吞咽碎玻璃，鲜血从嘴角不断溢出，可她的双手却死死攥着胸口一枚温润的玉坠------那是牢大临行前塞入她掌心的。

她感知到了。

在那道鸿蒙神光炸裂天际的瞬间，在那股撕裂天地的终极余波席卷而来的刹那，她体内与牢大神魂相连的魂灵共生脉络，在一瞬间全部断裂。每一根脉络断裂时，都像是在她心口剜去一块肉，痛到她几乎昏厥。

可她不肯昏过去。

因为她清晰地感知到了牢大最后的神念------那道穿越万里密林、穿透层层杀机、精准落入她灵魂深处的温柔讯息。

\emph{``好好活下去。''}

只有五个字。

没有告别，没有遗言，没有不舍，只有他一贯的、不容置疑的温柔命令。

蔡须坤的泪水无声地滑落，淌过苍白的脸颊，滴落在掌心的玉坠上。她的嘴唇微微翕动，像是想要回应什么，却发不出一丝声音。她的手指死死攥紧玉坠，指节泛白，骨节咯咯作响，像是攥着这世上最后一根救命稻草。

她感知到了他的死。

感知到了他燃尽一切的那一刻，感知到了他最后的温柔，感知到了他消散前最后的心念------不是恐惧，不是不甘，而是她的名字。

须坤。

须坤。

须坤。

他的最后一缕神念，从头到尾，只喊了她的名字。

``啊------！！！''

一声撕心裂肺的悲鸣终于从蔡须坤胸腔中迸发而出，沙哑破碎。她将玉坠死死按在心口，蜷缩起身躯，浑身剧烈颤抖，泪水与鲜血混在一起，浸透了身下的泥土。

小舞紧紧抱住她，将她的头按在自己怀中，自己的眼泪却怎么也止不住，一滴一滴落在蔡须坤的发间。

``须坤\ldots\ldots 须坤\ldots\ldots{}''

她不知道该说什么。任何安慰在这一刻都显得苍白可笑。她能做的，只有抱着她，像牢大临走前嘱托的那样------护住她。

霍雨浩缓缓站起身，双腿颤抖得几乎站不稳，他抬手抹去脸上的血与泪，目光望向身后那片漆黑的密林，眼底翻涌着滔天的悲痛与刻骨铭心的恨意。

``走。''

她的声音沙哑到几乎听不清，却带着一种前所未有的坚定。

她拖着重伤身体，再次而又坚定的说：``队长用命换来的路，我们不能浪费。''

唐三沉默片刻，缓缓站起，擦去嘴角的血迹，目光沉重却坚毅：``走。''

奥斯卡咬着牙撑地起身，卢本伟一言不发地站起，小舞扶着蔡须坤缓缓站起------蔡须坤的双腿软得几乎无法支撑身体，全靠小舞搀扶，可她的手，始终死死攥着那枚玉坠，一刻也不曾松开。

没有人再回头。

不是不想，是不敢。

他们怕一回头，就会看见那道本该站在身后的挺拔身影，然后发现那只是幻觉。

他们怕一回头，就会彻底崩溃，再也迈不动脚步。

他们只能往前走，带着牢大的遗志，带着他的牺牲，带着他用命换来的生路，一步一步，走向密林深处。

风过林梢，像是有人在轻声叹息。

\begin{center}\rule{0.5\linewidth}{0.5pt}\end{center}

星斗大森林外林要道。

两道正道霞光撕裂长空，梅天狸长老与ddfg长老全速驰援而来，满心焦灼，日夜兼程，只为救下身陷死局的果糕小队一众天骄。

可就在二人即将踏入核心战场的刹那，一声震天狼嚎炸响四野，滔天蛮荒兽威骤然封锁前路！

\textbf{二十万年顶级凶兽------灭世狼王} 踏步而出，漆黑鬃毛倒竖，猩红双眸饱含戾气，庞大的身躯遮天蔽日，浑厚的兽域之力死死锁死整片要道。

人类与魂兽的千年矛盾根深蒂固。在灭世狼王眼中，所有踏入星斗的人类强者皆是侵略者，绝不允许二人入场厮杀、惊扰森林生灵！

\begin{quote}
``人类强者，止步于此！''
\end{quote}

狼王戾气滔天，二十万年蛮荒兽力全面爆发，无解兽域封锁天地，拼死阻拦二人驰援之路。

两位长老心急如焚，明知战场瞬息万变，却无可奈何，只能全力出手，与灭世狼王展开死战。顶级正道法则与蛮荒兽力疯狂对冲，山林崩裂、灵气暴乱，战况极致惨烈。

灭世狼王凭借星斗地界加持，拼死缠斗、死战不退，不惜损耗本源也要拖延时间。两位长老虽实力碾压，却无法瞬间破局脱身，只能被动鏖战，眼睁睁看着时间飞速流逝。

数息之后，二人终于击碎兽域、重创狼王，冲破阻拦极速奔赴战场。

可入目之处，只剩漫天暴乱的能量余波、崩塌的山峦与满地狼藉。

那道撼动天地的大道陨落异象，已然彻底落幕。

梅天狸身躯剧震，心神骤凉，眼底涌上无尽的沉痛与无力。他缓缓环顾四周------碎裂的大地、被秩序之力烧焦的草木、空气中尚未散尽的鸿蒙神韵残息，以及\ldots\ldots 那四名横七竖八、重伤昏迷的帅气小队精锐。

他蹲下身，指尖轻触地面上一滩尚未凝固的血迹，指腹传来微弱的秩序波动------那是牢大的气息。

是那个孩子留下的，最后的痕迹。

梅天狸闭上眼，良久，一声苍老的叹息自喉间溢出。

ddfg伫立风中，望着空荡荡的血染山谷，声音沙哑哽咽：

\begin{quote}
``晚了\ldots\ldots 我们终究，来晚了一步。''
\end{quote}

他们冲破重重阻碍，击退二十万年凶兽，终究没能护住那位以身殉道、守护羁绊的少年。

梅天狸缓缓站起身，目光扫过满目疮痍的战场，最后落在山谷中央------那里什么都没有，没有遗体，没有遗物，连一丝完整的神魂残片都没有留下。那个少年将自己燃烧得干干净净，像是从未存在过。

可大地上那道深达百丈的掌印还在。

那是他最后站立的地方。

梅天狸走到掌印旁，沉默良久，缓缓单膝跪地，将右拳抵在左胸，行了一个正道最高规格的英灵祭礼。

ddfg紧随其后，同样单膝跪地，神色肃穆。

``果糕小队队长，牢大------''

梅天狸苍老的声音在空旷的山谷中回荡，带着无尽的沉痛与敬意。

``以身殉道，以命护道。正道序数一脉，铭记此恩。''

``英灵不朽，浩气长存。''

山风呼啸而过，卷起漫天落叶与尘埃，像是天地在无声地回应。

那道深达百丈的掌印中，隐隐泛起一丝微弱到几乎不可察觉的鸿蒙光芒，转瞬即逝------

像是回应，又像是告别。

\begin{center}\rule{0.5\linewidth}{0.5pt}\end{center}

风过荒山，大道哀鸣。

残阳终于彻底沉入地平线，夜幕降临，星斗大森林上空繁星漫天，冷冽的星光洒落在满目疮痍的山谷中，照亮了那道孤零零的掌印。

那位以血肉为盾、以神魂为薪，护队友、护挚爱、护正道的少年，永远留在了这片血染的星斗山林。

他没有留下遗体，没有留下遗言，甚至没有留下一缕残魂。

他留下的，只有队友活着的呼吸，只有挚爱掌心温热的玉坠，只有正道序数一脉延续的火种------

和那道深深刻入大地的掌印，证明他曾在这里，最后站立过。

化作乱世之中，最悲壮、最滚烫的正道绝唱。

\chapter{第四十三章：一线残魂，绝境希望}\label{ux7b2cux56dbux5341ux4e09ux7ae0ux4e00ux7ebfux6b8bux9b42ux7eddux5883ux5e0cux671b}

\begin{center}\rule{0.5\linewidth}{0.5pt}\end{center}

残阳如血，浸染星斗荒谷。

方才那场撼动天地的序力绝杀缓缓落幕，黑白交织的毁灭风暴逐层溃散，肆虐四野的能量暴乱渐渐平息。可整片山谷早已不复往日生机，大地崩裂出纵横千里的沟壑，焦黑土石堆叠满地，断裂的古木与残破的草木浸透暗红血痕，凌厉的杀伐之气与混沌余威萦绕不散。短短数息，这座灵气充沛的林间谷地，彻底沦为一片死寂萧瑟的修罗残局。

方才为所有人劈开生路的，是果糕小队的队长------\textbf{牢大}。

在整支帅气小队三行BMS混沌合击的无解围剿之下，绝境逼命之际，牢大彻底破壁，引爆毕生修为，催动终极奥义「坠机」。以燃烧全部生命本源、崩碎千年道基、倾覆周身所有序数脉络为代价，催生出霸道绝伦的ψ(Π\_ω)秩序洪流。

这道挣脱体系桎梏的极致力量，硬生生撕裂了层层叠叠的混沌杀局，正面重创敌方核心副队长，废其本源、碎其战力，彻底击溃了邪道小队的碾压攻势。漫天无解杀机尽数被一人拦下，濒临覆灭的果糕小队得以脱身逃生。

他赢下了战局，护住了所有人，却唯独输掉了自己。

狂风卷过焦土，山谷中央空空荡荡，再无那道永远挡在阵前的挺拔身影。

毁灭性的混沌之力席卷全场，将他的肉身筋骨、血脉经脉彻底碾灭，毕生苦修的I序数、K序道基层层崩塌碎裂，贯穿神魂与躯体的所有秩序脉络寸寸湮灭。天地无情，未曾为这位殉道少年留下半分残骨、一缕余血，唯有一方漆黑焦土，默默镌刻下他以命护道的悲壮终章。

昔日撑起整支小队的擎天壁垒，终究化作清风流云，消散于天地之间。

\begin{center}\rule{0.5\linewidth}{0.5pt}\end{center}

战场之上，压抑的死寂令人窒息。

幸存的帅气小队精锐伫立残墟，人人带伤、心神震颤，眼底藏着难以掩饰的后怕与忌惮。空气中久久不散的K序燃命余威，带着一股置之死地、玉石俱焚的决绝执念，哪怕主人已然陨落，依旧震慑着这群久经杀伐、心性冷血的邪道修士。他们从未料到，一个已然油尽灯枯、临时绝境突破的正道少年，竟能爆发出撼动整支精锐小队的逆天战力，将一场必胜围剿打成惨败残局。

小队副队长的境遇最为凄惨狼狈。

他整条右臂被秩序洪流彻底轰碎，筋骨寸断、脏腑爆裂，周身流转的0 111 221高阶序列纹路尽数黯淡崩裂、紊乱暴走。狂暴的混沌之力持续侵蚀其神魂本源，让他气血逆流、口鼻不断溢出血沫，一身纵横多年的顶尖战力彻底报废。此刻的他瘫倒在焦土之上，动弹不得，生机不断流逝，彻底沦为弃子。

其余四名精锐亦是护体序列斑驳破损、气血亏虚严重，军心彻底涣散。一场看似碾压的绝杀战局，被一名少年以命翻盘，极致的恐惧笼罩众人，无人再敢久留。

\begin{center}\rule{0.5\linewidth}{0.5pt}\end{center}

就在整片战场陷入死寂僵持的瞬间，万里长空骤然霞光炸裂！

两道横贯云海的璀璨正道金光撕裂云层，裹挟着浩瀚厚重的顶级修为威压，瞬息抵达星斗荒谷上空。纯正磅礴的正道法则倾泻而下，瞬间凝滞了天地间躁动的混沌邪气，层级碾压的恐怖气场笼罩四野，让整片山谷的空气都为之凝固。

是魂殿驰援的顶级长老，\textbf{梅天狸}与\textbf{ddfg}！

高空之上，帅气小队黑袍队长眸光骤然一缩，面色瞬间沉如寒渊。他深耕高阶序列修行多年，眼界毒辣，一眼便洞悉两位长老的修为深不可测，绝非自己这支残损小队能够抗衡。原本清扫战场、追杀残余、斩草除根的心思瞬间烟消云散，心底只剩极致的警惕与退意。

大势已去，滞留此地，唯有全军覆没。

他行事冷酷果决，从无半分犹豫，当即转身，欲破空遁逃。

\begin{quote}
``队长！救我！带我走！''
\end{quote}

瘫倒在地的副队长察觉到撤离的动向，强忍撕心裂肺的肉身与神魂剧痛，拼尽最后一丝气力嘶吼求救。他追随队长征战数载，浴血拼杀、屡立战功，为小队闯下无数威名与资源，从未想过自己重伤废残、失去战力后，迎来的会是无情舍弃。眼底浓烈的绝望与哀求，在死寂的山谷中显得格外凄厉刺耳。

可回应他的，只有彻骨的寒凉与无情。

黑袍队长脚步未顿，自始至终未曾回头一眼，淡漠冰冷的声音穿透风声，不带半分同僚情义：

\begin{quote}
``你本源尽损，序列崩废，再无利用价值。''

``战场唯利，无用之人，不必拖累全队。''

``自生自灭。''
\end{quote}

短短三字，斩断数年浴血情谊，道尽帅气教残酷冰冷的生存法则。在这群邪道修士眼中，从来只有战力高低、利用价值，无关情义、无关陪伴，一旦沦为累赘，便会被毫不犹豫地彻底舍弃。

话音落定，黑袍队长袖袍一挥，周身混沌序列全力运转，化作一道漆黑残影撕裂长空，瞬息绝尘而去、遁逃无踪。余下几名轻伤精锐彻底军心崩盘，不敢多做片刻停留，纷纷四散奔逃，顷刻间撤离一空。

喧嚣散尽，荒谷重归死寂。

血染的焦土之上，只剩那名重伤濒死的副队长孤零零僵卧原地。他望着空空荡荡的长空，数年征战的热血、执念、忠心尽数落空，眼底最后一丝光亮彻底熄灭，只剩无边无际的冰冷绝望，静静等待死亡降临。

\begin{center}\rule{0.5\linewidth}{0.5pt}\end{center}

数息之间，两道璀璨正道霞光轰然落地。

梅天狸与ddfg身形稳稳伫立在战场核心，浩荡纯正的正道威压席卷四野，瞬间镇压、涤荡了山谷内残存的所有混沌邪气，驱散了漫天阴冷杀机，让躁动的天地缓缓平复，并迅速治疗了小队的伤势。

二人目光快速扫遍全场，入目皆是崩裂的山峦、倾覆的草木、暴乱未消的能量余波，以及奄奄一息的敌方副队长。可那道逆势抗天、以身护友的少年身影，却遍寻无迹。

空气中，那道独属于ψ(K\_ω)的霸道秩序道韵，已然微弱到近乎断绝，濒临彻底消散。

\begin{quote}
``来晚了。''
\end{quote}

\begin{center}\rule{0.5\linewidth}{0.5pt}\end{center}

与此同时，星斗大森林深处，浓林蔽日、阴风凄凄，簌簌落叶纷飞，宛若天地同悲。

侥幸突围的果糕小队，没有半分劫后余生的庆幸，铺天盖地的悲恸与愧疚，早已将所有人彻底吞没。

一行人浑身带伤、气血透支，踉跄着驻足幽深林地，无人敢于挪动半步。方才山谷之中，牢大孤身断后、燃命殉道的画面，一遍遍在众人脑海中循环回放，刻骨铭心、挥之不去。每一寸他们踏过的生路，都是少年用性命铺就；每一次得以存续的呼吸，都承载着他倾尽所有的守护。

林间死寂无声，唯有众人粗重颤抖的喘息，以及一颗颗滚烫泪水砸落枯叶的细碎声响。压抑的悲伤笼罩全队，让人几近窒息。

过往数年朝夕相伴、并肩浴血的点点滴滴，在此刻尽数翻涌心头，化作最锋利的刀刃，狠狠割裂每个人的心神。

\textbf{唐三}闭上双眼，眼底通红酸涩，满心皆是无力与自责。初入星斗森林的懵懂青涩，一次次绝境遇险的惶恐无助，一次次历练崩盘的心态崩塌，每一次都是牢大挺身而出，稳压强敌、布局破局，沉稳安抚全队，带着众人一次次绝境翻盘。他永远是全队最坚实的靠山、最可靠的后盾，永远把危险挡在身前，把生机留给队友。可这一次，他独自扛下了无解的死局，把所有生路尽数赠予众人，自己却尸骨无存、消散天地。

\textbf{霍雨浩}心神震颤，识海空空荡荡，久违的惶恐与无助席卷全身。自小队组建以来，牢大霸道沉稳的秩序道韵，始终是全队最稳固的定心丸，能镇压一切紊乱杀机。二人无数次并肩作战、攻防相辅，牢大总能精准预判险境，替他补齐所有破绽，为他兜底护行。可如今，那道熟悉安心的道韵彻底断绝，紧密相连的共生羁绊轰然断裂，从此再无人为他遮风挡雨，再无人护全队周全。往日热血并肩的时光，此刻尽数化作剜心刺骨的痛楚。

\textbf{小舞}双肩不住颤抖，泪眼朦胧，脑海里满是少年温柔可靠的模样。平日里的嬉笑打闹、绝境中的伸手相护、负伤后的轻声宽慰、历练时的细心包容，桩桩件件、历历在目。牢大素来温柔沉稳，护着队内每一个人的成长，无论前路多险、敌人多强，他永远伫立阵前，替众人扛下所有风雨。她从未想过，那个永远从容笃定、无所不能的队长，会以如此惨烈的方式，永远离开大家。

\textbf{奥斯卡、卢本伟与诺当}三人低头伫立，满心愧疚无以复加。一路走来，牢大始终包容他们的短板与稚嫩，悉心带队成长，将历练机缘、修炼资源尽数偏向队友，自己却永远冲锋在最险前线。无数次生死对局，他们依仗着牢大的庇护安稳输出、稳步进阶，可在最后的终极死局里，他们拼尽全力也无法替他分担分毫，只能眼睁睁看着他燃尽道基、湮灭身形，徒留满心悔恨，扎根心底、无处释怀。

全队之中，最痛苦、最煎熬的，莫过于\textbf{蔡须坤}。

她浑身冰凉，四肢百骸皆是刺骨寒意，泪水早已浸透衣襟，身躯摇摇欲坠、几近不支。队友们看见的是牢大以身殉道的大义凛然，唯有她，看得见这份牺牲背后极致的温柔与偏爱。无数个日夜的默默守护，危难之时下意识的护她周全，临行前那句安抚众人的温柔谎言，绝境里至死未改的牵挂执念。

她能清晰感知到，自己魂灵共生的血脉脉络深处，一缕细碎微弱、摇摇欲坠的神魂正在顽强摇曳。那是牢大留在世间的最后痕迹，是他拼尽一切、舍弃所有，只为换她平安的证明。

她今日侥幸存活的每一分生机，皆是他用性命换来。这份厚重的恩情与羁绊，压得她神魂俱痛、几近崩溃。

\begin{center}\rule{0.5\linewidth}{0.5pt}\end{center}

无尽悲恸压抑周身，窒息的绝望笼罩众人。良久，蔡须坤咬紧牙关，强行压下眼底汹涌的泪水与酸涩，抬首望向荒谷的方向，声音沙哑颤抖，带着全队人最后的、近乎偏执的希冀，隔空发问：

\begin{quote}
``ddfg长老、梅天狸长老\ldots\ldots 牢大他，真的彻底陨落了吗？有没有\ldots\ldots 哪怕一丝复活的可能？''
\end{quote}

这一句话，道出了所有人心底最后的期盼。一众少年少女纷纷抬眸，通红的眼眸里死死攥着一缕渺茫微光，屏息等待答案。他们无论如何也无法接受，那个护了他们一路、撑起整支小队的少年，就此彻底消散、永归虚无。

\begin{center}\rule{0.5\linewidth}{0.5pt}\end{center}

荒谷之上，风尘萧瑟。

两位长老听闻林间颤抖的问询，相视一眼，神色复杂凝重，眼底满是唏嘘与沉重。

ddfg嗓音沙哑干涩，眼底翻涌着难以掩饰的沉痛与惋惜。他指尖轻触地面残留的能量痕迹，神识瞬间回溯整场战局，将牢大燃尽本源、崩碎道基、以身殉道、独拦全军的悲壮画面，看得分毫毕现。

\begin{quote}
``他以肉身作盾、魂力为薪，燃尽毕生一切，硬生生拖住整支邪道精锐，为果糕小队换来了全员生路。''
\end{quote}

梅天狸眉头紧锁，神色凝重如铁，周身灵气尽数铺展，细致探查山谷每一寸角落，不肯放过分毫痕迹，语气沉重无比：

\begin{quote}
``肉身彻底湮灭，道基归零崩塌，周身序数脉络寸寸崩碎。肉身、道基、魂力三重尽毁，这般伤势，寻常修士早已神魂俱灭、彻底消散，再无轮回存续的可能。''
\end{quote}

二人心底沉至谷底，却依旧不肯彻底放弃。牢大是新生代序数大道的破壁者，是甘愿舍己护人的少年英雄，是正道未来不可或缺的希望，若就此彻底陨落，必将是整片大陆正道的莫大损失。

就在二人遍寻无果、满心绝望之际，ddfg的神识骤然剧烈一震！

他眸光骤然锁定星斗密林深处、果糕小队仓皇逃离的方向，一缕缥缈至极、微弱欲熄的神魂气息，正于萧瑟寒风中顽强摇曳、死死存续，未曾彻底消散于天地！

\begin{quote}
``还有残魂！''
\end{quote}

ddfg沉声低喝，沙哑的语气中炸开一丝难以置信的惊喜，让死寂绝望的战局，骤然透出一缕绝境微光。

梅天狸心神巨震，即刻凝神溯源，顺着那缕微弱的神魂气息层层探查，瞬息洞悉全部真相，眼底瞬间涌上动容与恍然：

\begin{quote}
``是蔡须坤！是她的帝皇瑞兽血脉秘术！''

``是独一无二的瑞兽护魂大道，保住了他最后一缕本源神魂！''
\end{quote}

方才牢大燃尽一切、身形湮灭的刹那，灭世般的混沌之力席卷天地，万物皆被侵蚀粉碎。逃亡途中的蔡须坤，清晰感知到挚爱生机断绝、道基崩塌、神魂溃散的极致危机，彻骨的悲痛与刻骨羁绊瞬间引爆了她深埋血脉的本命秘术。

无形温润的神圣瑞光跨越空间阻隔，精准笼罩住牢大即将彻底溃散的本源神魂。她以二人朝夕相伴、生死与共的羁绊为锚，以自身本命瑞气为屏障，硬生生在混沌之力的极致冲刷下，死死锁住了他最后一缕本心残魂。

这是绝境之中唯一的奇迹，是生死羁绊催生的逆天生机。

肉身可碾，道基可毁，大道可崩，唯独他滚烫的守护执念、纯粹的正道本心，得以留存天地，未归虚无。

\begin{quote}
``只是这缕残魂状态极差。''梅天狸敛去动容，神色重新覆上浓重凝重，``无肉身承载，无道基维系，只能暂时寄存在蔡须坤的魂灵共生脉络之中，靠着她源源不断输出的瑞兽灵气勉强存续。''
\end{quote}

她轻轻摇头，眼底满是忧虑：

\begin{quote}
``残魂脆弱如风中残烛，毫无自保之力，但凡有一丝外力侵扰、灵气断绝，便会彻底飘散，再无半分挽回之机。''
\end{quote}

\begin{center}\rule{0.5\linewidth}{0.5pt}\end{center}

话音落下，荒谷之内再度陷入死寂，气氛压抑得令人窒息。

ddfg望着密林深处那缕摇曳欲熄的魂息，轻声长叹，语气复杂难言：

\begin{quote}
``他留住了一线魂息，便有重生的火种。可这份希望，太过微弱，近乎虚无缥缈。''
\end{quote}

片刻后，ddfg缓缓开口，沉稳的声音穿透林间风声，带着难以掩饰的疲惫：

\begin{quote}
``诸位，不必绝望，他\ldots\ldots 并未彻底消散。''
\end{quote}

一语落地，林间死寂轰然打破！

果糕小队众人浑身巨震，濒临破碎的心神骤然被托住，通红的眼眸中瞬间炸起璀璨微光。那一缕濒临熄灭的希望星火，骤然刺破了笼罩全队的无边黑暗。

梅天狸抬手轻引，将那缕寄宿在蔡须坤血脉中、微弱欲熄的神魂气息，清晰展露在众人感知之中，轻声道：

\begin{quote}
``是蔡须坤的帝皇瑞兽血脉秘术，以生死羁绊为锚，绝境护魂，保住了他最后一缕本心残魂。这是天意，也是你们二人的羁绊奇迹，让他得以残魂存续、未归虚无。''
\end{quote}

刹那间，热泪翻涌而出，所有人紧绷的心弦骤然松弛，悲恸之中生出无尽期许。

只要残魂尚在，就还有重逢的可能。只要魂息未灭，他们就还有再见队长的机会。

\begin{center}\rule{0.5\linewidth}{0.5pt}\end{center}

可这份来之不易的欣喜与希冀，仅仅存续片刻，便被ddfg沉重冰冷的话语，狠狠击碎殆尽。

\begin{quote}
``残魂虽在，可重生之路，难如登天，近乎无解。''
\end{quote}

ddfg字字沉重，清晰道出逆天重生的苛刻条件，将众人刚刚升起的希望，狠狠打入万丈深渊：

\begin{quote}
``他肉身湮灭、道基尽碎、序数脉络全盘崩塌。想要重塑肉身、重修道基、修复破碎的序数大道，必须同时集齐\textbf{乾坤造化丹}与\textbf{天地灵水}两样旷世至宝，缺一不可。''

``乾坤造化丹可逆转生死、重构血肉，是重塑肉身的唯一神物，万年难寻、藏于绝域；天地灵水源自天地本源，可滋养残魂、修复大道，世间存世寥寥、无价无市。''

``二者兼得，以你我当下的力量与资源，毫无可能。''

乾坤造化丹万年难遇，隐匿于世间最顶级的绝域秘境，可遇而不可求；天地灵水更是绝世稀缺，放眼整片大陆存世寥寥，千金不换、权势难夺。二者缺一不可，同时集齐更是近乎天方夜谭。以当下魂殿的底蕴、众人的修为与资源，想要寻得两样至宝，无异于痴人说梦。
\end{quote}

\begin{center}\rule{0.5\linewidth}{0.5pt}\end{center}

梅天狸静静立在萧瑟风里，望着那缕在血脉中苟延残喘的微弱魂息，良久，才轻轻吐出一声叹息。

这一声极轻，风一吹便要散，却压尽了千言万语的沉重。她活过漫长岁月，见过天骄辈出、见过大道无常，见过无数生死逆转、绝境逢生。可今日，她清清楚楚看着------少年以命殉道、以心护友，人间最赤诚的情义，硬生生换来的，只是一缕将熄残魂、一场无解天意。

她太清楚这两样至宝的分量。不是苦修可补，不是机缘可觅，是整片天地都吝啬给予的一线生机。世人皆道残魂未灭便是奇迹，可只有她心知肚明：这不是救命的火种，是上天留给这群孩子、最残忍的温柔酷刑。

奇迹给了他们念想，天道断了他们归途。

这一声轻叹，是悲悯，是惋惜，更是对这场徒劳牺牲、无望重逢的无声定论。温柔，却彻骨冰冷，彻底压灭了林间最后一丝侥幸。

\begin{center}\rule{0.5\linewidth}{0.5pt}\end{center}

果糕小队众人静静伫立，无人争辩、无人反驳，死寂瞬间吞没整片林地。

方才骤然燃起的希望，在这一刻彻底沉淀、冷却，化作彻骨的冰凉与通透的绝望。他们不再是懵懂的孩童，历经无数历练、踏遍诸多秘境，早已深知天地至宝的稀缺与难得，更懂修行界的规则与壁垒。

他们终于彻底明白，长老口中的''难如登天''从不是客套的告诫，而是冰冷刺骨的事实。

不是努力不够、不是机缘未到，而是这条重生之路，从根源上就堵死了所有可能。

\textbf{唐三}喉间哽咽，心底最后一丝侥幸轰然碎裂。他终于认清，牢大留给他们的从来不是重生的希望，而是一份遥遥无期、永远无法兑现的念想。那缕残魂存续的意义，仿佛只是为了让他们清醒地承受这份求而不得的煎熬。

\textbf{霍雨浩}识海阵阵抽痛，彻底读懂了长老叹息里的沉重。残魂未灭，却形同禁锢，只能依附血脉苟延残喘，无法苏醒、无法归位，更无法重塑身形。这份活着的残影，远比彻底的陨落更磨人，时时刻刻提醒着众人那场悲壮的牺牲与无解的遗憾。

\textbf{小舞}眼眶通红，泪水无声滚落，心底酸涩泛滥。原来奇迹之后，是更残忍的绝境。上天仁慈地留住了他一缕残魂，却又极致残忍地断绝了所有重逢的可能。

\textbf{奥斯卡、卢本伟、诺当}三人垂首伫立，满心的期盼尽数归零。他们终于知晓，自己哪怕穷尽余生、踏遍万里山河、倾尽所有修为资源，也绝对凑不齐这两样旷世至宝，根本无力偿还少年以命换命的深情守护。

最痛的依旧是\textbf{蔡须坤}。

她清晰感知着脉络中那缕微弱摇曳的魂息，指尖微微颤抖，心底一片荒芜。是她留住了他最后的残魂，可也亲手留住了这份无尽的煎熬与无望。她护住了他的本心，却没有半分能力渡他归来。

她终于彻底懂得，自己日夜牵挂的重逢，终究只是一场触不可及的幻梦。

\begin{center}\rule{0.5\linewidth}{0.5pt}\end{center}

冰冷的真相随风漫入林间，字字诛心，彻底碾碎了众人心中刚刚燃起的微光。

方才眼底滚烫的希冀瞬间消散，取而代之的是更深、更沉、更让人绝望的无力。

原来不是彻底陨落的绝路，而是一条看得见、却永远走不通的生路。

咫尺希望，万丈天堑。这般可望而不可即的生机，比起彻底的死寂陨落，更让人窒息、更让人崩溃。

林间寒风呜咽，木叶萧萧纷飞，天地悲意更胜从前。

果糕小队众人伫立幽深密林，眼底光亮彻底黯淡，满心期许尽数落空。刚刚燃起的生机火种被无情掐灭，所有人再度被无边无际的沉痛、愧疚与绝望，彻底吞没

\chapter{第四十四章：残魂枷锁，前路危局}\label{ux7b2cux56dbux5341ux56dbux7ae0ux6b8bux9b42ux67b7ux9501ux524dux8defux5371ux5c40}

林间悲风未歇，彻骨的绝望依旧死死缠绕着果糕小队每一个人。

众人伫立幽深密林，眼底最后一点微光彻底熄灭，方才奇迹复生的欣喜，早已被逆天无解的重生条件碾得粉碎。看得见的生机，摸不着的归途，这份残忍的侥幸，比彻底的陨落更令人崩溃。所有人沉默伫立，泪水无声滑落，胸腔里塞满了愧疚、悲痛与无力，整片山林都浸泡在沉沉哀恸之中。

荒谷战场之上，萧瑟风声猎猎作响。

ddfg长老收敛眼底唏嘘悲悯，抬眸望向谷地中央那名奄奄一息的帅气小队副队长，周身骤然升腾起凛冽刺骨的正道威压。方才面对少年殉道、天人永隔的柔软尽数褪去，取而代之的是执掌正道秩序、审判邪祟的冰冷肃杀。

整片荒谷的气流瞬间凝固，狂暴余波尽数平息，唯有厚重的法则镇压之力笼罩四方，让人喘不过气。

那名被舍弃的副队长依旧瘫倒焦土之上，断臂残躯血流不止，本源破碎、序列崩废，早已失去所有战力，只剩最后一丝残喘苟延。他望着步步走来的正道长老，眼底瞬间被极致的恐惧填满，浑身剧烈颤抖。

他追随帅气小队数年，杀伐无数、造孽深重，深谙正道修士的准则------除恶务尽，绝不姑息。更何况，方才那场死局围剿，他是亲手施压、催动绝杀的核心元凶之一，亲手将那位赤诚少年逼至燃命殉道的绝境。

\begin{quote}
``尔等仗混沌邪序，屠戮正道修士，践踏生灵道义，以强凌弱、围杀天骄，造下无边杀孽。''
\end{quote}

ddfg声音沉如惊雷，字字铿锵，裹挟着天地正道的审判之力，响彻整片荒谷：

\begin{quote}
``同队战友，重伤便弃、无情无义，邪道凉薄，卑劣至极。乱世杀伐，自有天道轮回，今日，我便以正道之名，清算你所有罪孽。''
\end{quote}

话音落下，不待对方求饶辩解，一抹纯净凛然的正道秩序神光自指尖迸发而出。

没有惊天动地的轰鸣，只有绝对层级的法则碾压。柔和却霸道的神光瞬穿虚空，精准落在那名副队长残破的身躯之上。

凄厉短促的惨叫骤然响起，又瞬息湮灭。

他周身残存的混沌邪气瞬间被净化殆尽，紊乱暴走的序列脉络寸寸崩解，满身罪孽与杀伐根基被彻底拔除。身躯与残余神魂在正道神光的涤荡下化作点点飞灰，随风消散，世间再无半分痕迹。

一招之下，罪徒伏诛，罪孽清零。

随后ddfg眸光扫遍四野，神识铺展千里，排查战场所有残留邪气与隐匿余孽。方才四散奔逃的几名轻伤精锐，早已吓破胆、亡命逃窜，可他们身上沾染的正道血气与混沌气息，早已成为无法抹去的追踪烙印。

\begin{quote}
``邪孽余毒，不容存世。''
\end{quote}

他冷喝一声，抬手结印，天地间游离的正道法则瞬间汇聚，数道秩序流光破空追出，横跨山林千里，精准追上逃窜的残敌。

天际远端接连响起数声闷响，所有帅气小队残余精锐，尽数被法则之力湮灭肃清，无一漏网。

至此，整场围剿果糕小队的邪道势力，彻底除名，片甲不存。

荒谷彻底重归寂静，邪气尽散、杀孽清零，唯独那场少年殉道的悲壮，永远镌刻在这片土地之上，久久不散。

\begin{center}\rule{0.5\linewidth}{0.5pt}\end{center}

密林深处，悲伤依旧蔓延，却悄然生出一丝隐忍的韧劲。

良久，萧瑟风声里，\textbf{唐三}缓缓抬起头，眼底通红，泪水早已流干，嗓音沙哑得近乎破碎：

\begin{quote}
``牢大以命换我们全员存活，不是让我们在这里沉溺悲痛、止步不前的。''

``他护住了我们，护住了正道火种，我们不能让他的牺牲，白白作废。''
\end{quote}

\textbf{霍雨浩}攥紧颤抖的双拳，识海之中牢大昔日护队、并肩作战的画面一遍遍闪过，心底的悲痛渐渐沉淀为沉甸甸的责任：

\begin{quote}
``队长不在了，但果糕小队不能散，正道的路，我们必须接着走下去。''
\end{quote}

\textbf{小舞、奥斯卡、卢本伟、诺当}几人纷纷抬首，满目沉痛却眼神坚定。他们失去了最可靠的队长、最坚实的后盾，可身上却背负着牢大倾尽性命换来的生机与希望。

所有人的目光，不约而同、尽数汇聚向人群中央的\textbf{蔡须坤}。

她依旧面色苍白、身躯虚弱，泪痕斑驳的脸颊难掩极致的憔悴，脉络中那缕摇摇欲坠的残魂依旧微弱飘摇，时刻牵动着她的心神。可所有人都清楚，她是留住牢大最后执念的人，是小队魂灵共生大道的根基，更是此刻全队唯一的精神锚点。

\begin{quote}
``须坤，由你来当新的队长。''唐三沉声开口，语气郑重无比，``这是牢大拼尽全力护住的队伍，也唯有你，能扛起这份责任，守住我们的道。''
\end{quote}

众人纷纷颔首，无人有半分异议。

\begin{quote}
``请须坤带队！''
\end{quote}

整齐肃穆的呼声穿透林间悲风，褪去了慌乱与脆弱，多了几分沉凝的力量。

蔡须坤身躯微颤，缓缓抬眸。眼底没有接任的欣喜，只有彻骨的寒凉与无尽的哀思。她清晰感知着血脉中那缕微弱的魂息，感知着牢大残留的守护执念，心口像是被巨石死死压住，喘不过气。

她多想沉溺悲痛，多想停在此地，可她不能。

是牢大以命铺路，换她安稳、换全队存续。她唯有扛起这份重担，带着他的执念、带着全队的正道火种继续前行，才不负他以身殉道的深情与牺牲。

唐三望着她憔悴却坚韧的模样，轻声补道，语气恳切：

\begin{quote}
``须坤，我们都知道你最难，你守着他最后的残魂，也守着他最后的牵挂。如今队长不在，只有你能稳住全队、守住我们的道。''
\end{quote}

霍雨浩郑重开口，目光笃定：

\begin{quote}
``没错，我们所有人都会听你调度、随你进退。无论前路多难，我们并肩扛住，绝不溃散。''
\end{quote}

小舞红着眼眶上前半步，声音哽咽却坚定：

\begin{quote}
``须坤，别怕，我们都陪着你，跟着你好好走下去，不辜负牢大的付出。''
\end{quote}

众人纷纷附和，声声恳切，掷地有声。

听罢众人的信任与托付，蔡须坤眼底翻涌的泪意终是被她硬生生压下。她挺直单薄却挺拔的脊背，抬手轻轻按在心口那缕飘摇的残魂之上，声音带着未散的沙哑泪音，却字字铿锵、无比坚定：

\begin{quote}
``我接下队长之位。''

``从今往后，果糕小队由我带队。我会守住所有人，守住我们与牢大一同坚守的正道，带着他的执念好好活下去，走完我们未走完的路。''

``我在此立誓，只要我蔡须坤一日尚存，小队绝不散、正道绝不灭！''
\end{quote}

简单几句誓言，承载了千钧重量，藏着无人知晓的煎熬与孤勇。

悲痛未消，责任已落肩。昔日被牢大拼命守护的温柔少女，在这场惨烈的生死离别中，褪去稚嫩、扛起责任，硬生生站成了全队新的壁垒。

整顿心绪，压下翻涌的酸涩，蔡须坤抬眸望向星斗大森林雾气沉沉的核心腹地，沉声号令，语气沉稳有度，尽显新任队长的果敢气度：

\begin{quote}
``全员休整片刻，妥善包扎伤势，调整状态。休整完毕，即刻向星斗核心行进。''

``不负牺牲，不负本心，不负牢大，不负我们所有人的大道。''
\end{quote}

众人应声颔首，尽数收敛悲色，压下心底绝望，目光坚定，默默起身整理伤势，循着幽深密林的方向，准备再度启程。

\begin{center}\rule{0.5\linewidth}{0.5pt}\end{center}

可就在众人即将迈步前行之际，一道清冷沉稳的声音骤然从林间上空落下，拦住了众人的脚步。

是\textbf{梅天狸}。

她与ddfg一同踏落林间，望着前方迷雾重重的星斗核心，神色凝重肃穆，眼底满是深深的忧虑，彻底打断了众人的行进节奏。

\begin{quote}
``你们暂且止步。''

``我知晓你们想砥砺前行、延续正道，可如今贸然踏入星斗核心，太过凶险，前路危机远超你们想象。''
\end{quote}

众人闻声驻足，纷纷抬眸望向长老，眼底满是疑惑。经历此番死战，他们早已身心俱疲、无惧凶险，为何此刻前路依旧被强行阻拦？

梅天狸望着茫茫林海，缓缓道出其中根源，一语道破最隐秘、最致命的深层危机：

\begin{quote}
``人类与魂兽，积怨千年、矛盾根深蒂固，从未真正和解。''

``你们人类修士踏入星斗，向来只为历练、猎取魂环、争夺资源，在所有魂兽眼中，人类皆是入侵者、掠夺者。如今你们身负正道名号深入核心，在魂兽看来，并非悟道修行，而是彻头彻尾的欺骗与伪装。''
\end{quote}

话音稍顿，她眼底忧虑更盛，道出了当下最棘手的祸根：

\begin{quote}
``方才你们深陷死战之时，二十万年修为的灭世狼王强行锁死要道、拼死阻拦驰援，我一时暴怒，出手重创了它的本源根基。''

``此战动静极大，本源创伤震荡整片星斗，所有高阶魂兽皆有所感。如今整片森林的魂兽族群，已然对人类修士布满戒备与敌意。''

``我重创其王族凶兽本源，这份恩怨，早已牵连整片星斗魂兽族群。此刻你们踏入核心，无论本心善恶、目的纯粹与否，都不会有任何魂兽愿意相信。''

``等待你们的，不会是机缘历练，只会是整片星斗魂兽的集体敌视、无尽截杀。''
\end{quote}

一番话落地，林间再度陷入死寂。

刚刚接过重担、勉强燃起前行信念的果糕小队，心头再度被一层厚重的阴霾笼罩。

后方是无解的重生绝境，前路是整片魂兽族群的敌视死局。

失去了最坚实的队长壁垒，背负着无解的执念与遗憾，如今又深陷人类与魂兽的千年恩怨漩涡之中。

风雨前路，步步荆棘，层层绝境。

\begin{center}\rule{0.5\linewidth}{0.5pt}\end{center}

整片林间的压抑阴霾里，众人皆陷入前路无望的沉寂，唯有\textbf{霍雨浩}眉头紧蹙，精神识海飞速运转，梳理着眼下死局的唯一破局点。片刻后，他抬眸开口，声音冷静沉稳，刺破满场低迷：

\begin{quote}
``长老，若想踏入星斗核心、安稳探寻机缘，唯一的办法，便是化解恩怨、消解魂兽的敌意。''
\end{quote}

所有人闻声转头，目光尽数落在他的身上。

霍雨浩目光坚定，字字清晰道出核心关键：

\begin{quote}
``人类与魂兽的积怨根深蒂固，寻常说辞、刻意示好皆无用。如今祸根的源头，便是被重创本源的灭世狼王。我们唯有主动救治它、弥补此次过错，才有一丝机会打破偏见，换取整片星斗魂兽族群的谅解与信任。除此之外，再无第二条路可走。''
\end{quote}

这番话一语点醒众人。

众人方才只看到前路的无尽敌视，却未曾想过，恩怨可解、死局可破。绝望的心底，再度浮出一丝微弱却切实的转机。

梅天狸眸光微动，神色稍缓，眼底的凝重却未曾褪去。她深知霍雨浩所言句句属实，这是唯一的破局之法，可其中难度依旧不容小觑。她沉吟片刻，郑重颔首：

\begin{quote}
``你所言极是，这是唯一能缓和人兽矛盾、破除眼下危局的途径。事不宜迟，我们即刻折返星斗森林边缘，寻灭世狼王。''
\end{quote}

事已至此，别无选择。

蔡须坤压下心中繁杂心绪，以新任队长的身份沉声号令，全队即刻调转方向，跟随两位长老，循着来时的路途，折返星斗大森林外围。

\begin{center}\rule{0.5\linewidth}{0.5pt}\end{center}

一路疾行，林间气氛肃穆压抑。众人身心俱疲，身后是牢大无解的复活绝境，身前是亟待化解的人兽恩怨，每一步前行，都走得沉重无比。

不多时，一行人抵达森林边缘的空旷林地。

此前被重创本源、强行逼退的\textbf{灭世狼王}，正盘踞在荒草地中央。庞大的兽躯匍匐在地，漆黑鬃毛凌乱枯萎，周身原本磅礴浩瀚的二十万年兽域之力彻底溃散，浑身兽纹黯淡无光，口鼻不断溢出暗红血丝，呼吸微弱紊乱。

而它周身，十余头高阶狼形魂兽警惕环绕，獠牙外露、低吼不断，兽瞳死死锁定一行人，周身戾气暴涨，随时准备拼死护主、发起突袭。

整片外围林地，再度被紧绷的对峙气氛笼罩。

这群狼族魂兽清晰感知到首领的本源重创，深知眼前人类强者实力恐怖，心中积攒了极致的怨恨与畏惧，没有半分松懈，死死阻拦众人靠近狼王半步。

\begin{quote}
``诸位无需戒备，我等今日前来，并非厮杀，只为弥补过错、救治狼王。''
\end{quote}

梅天狸抬手示意全队止步，收敛周身所有威压，褪去所有锋芒，姿态平和坦诚，主动释放出毫无恶意的正道气息，试图安抚躁动的狼族魂兽。

可族群的戒备根深蒂固，过往的厮杀掠夺、今日的重创之仇，绝非一句安抚便能消解。狼族低吼不止，敌意未减分毫，对峙的紧绷气氛丝毫没有松动。

眼见言语无用，梅天狸不再迟疑。为了全队前路，为了彻底化解人兽死结，她主动迈步上前，无视周遭狼族的威慑低吼，掌心徐徐升腾起温润纯净的生命灵光。

柔和的翠绿圣光流淌萦绕，不含半分杀伐与掠夺之力，纯粹只为滋养生灵、修复伤势，缓缓笼罩住奄奄一息的灭世狼王。

圣光入体的瞬间，灭世狼王庞大的身躯猛地一颤，低沉的痛哼从喉间溢出。

梅天狸眉心微蹙，神色彻底沉了下来。透过圣光的本源探查，她已然洞悉了最残酷的真相，语气满是惋惜无奈：

\begin{quote}
``抱歉，我尽力了。''

``你的本源根基已彻底碎裂，大道脉络全盘崩坏，并非寻常外伤，根本无法修复。''

``二十万年苦修底蕴一朝散尽，修为已然彻底报废，寿元大幅折损，根本撑不过下一次天地天劫。哪怕有顶级生命之力滋养，也只能勉强维系残命，再无复原可能。''
\end{quote}

冰冷的真相落下，狼族的低吼瞬间停滞，整片林地彻底陷入死寂。

灭世狼王艰难抬起沉重的头颅，猩红黯淡的兽瞳牢牢注视着身前的梅天狸，眼底满是浓浓的疑惑与不解。

它很清楚，眼前这位人类长老实力深不可测，方才一念之间便能重创它二十万年本源，足以碾压整片狼族族群。以她的实力，完全可以趁它重伤濒危、无力反抗之际，抽干它残存的所有本源能量、夺取它的十万年兽核，尽数化作自身修为底蕴，无需耗费心力为自己疗伤。

明明二者方才是生死对立的仇敌，明明她可以顺势掠夺、获益良多，可她最终选择放下敌意，耗费自身灵力，徒劳为自己医治残伤。

这份反常的温柔与善意，彻底颠覆了灭世狼王对人类修士的固有认知。

自古以来，人类踏入星斗，皆是掠夺、厮杀、索取，从无这般明知无解、依旧倾力施救的举动。

它眼底的暴戾与戒备缓缓褪去，取而代之的是深深的茫然与迟疑。周身紧绷的兽躯彻底松弛，死死对峙的戾气悄然消散，萦绕整片林地的窒息紧张气氛，终于缓缓化解、烟消云散。

\chapter{第四十五章：瑞兽降世，狼王归灵}\label{ux7b2cux56dbux5341ux4e94ux7ae0ux745eux517dux964dux4e16ux72fcux738bux5f52ux7075}

\begin{center}\rule{0.5\linewidth}{0.5pt}\end{center}

紧绷的对峙气氛，依旧笼罩着这片外围林地。

灭世狼王眼底的暴戾虽散，可周遭盘旋的狼族高阶魂兽，依旧獠牙森寒、低吼不止。它们世代生于星斗、死于星斗，刻在血脉深处的记忆里，从来只有人类的掠夺、偷袭与背信弃义。

梅天狸的救治暂时稳住了狼王的伤情，却修复不了破碎的根基。狼群看见她耗损灵力，依旧没有放松戒备：它们的首领，此前正是伤在她的手下。

善意可以化解仇恨，却难以抹平千年积怨。

梅天狸看着依旧紧绷对峙的狼族，轻轻蹙眉，微微摇头，眼底满是无奈：

\begin{quote}
``伤是我造成的，它们不会立刻信我。''梅天狸低声道，``先稳住狼王，别再逼它们后退了。''
\end{quote}

ddfg伫立一旁，神色沉凝：

\begin{quote}
``留在这里僵持，对伤者也无益。''ddfg看了看护住首领的狼群，``我们先退开些，给它们一点余地。''
\end{quote}

果糕小队众人静静伫立。狼族的领地挡在通往核心的路上，绕行并非毫无可能，却会耽搁更多时日。比赶路更令他们在意的，是不能将眼前这份伤留在身后，便当作什么都没有发生。

\begin{center}\rule{0.5\linewidth}{0.5pt}\end{center}

就在众人沉默之际，一道清晰而坚定的女声响起。

\begin{quote}
``让我来吧。''
\end{quote}

\textbf{蔡须坤}缓步走出队列，立于众人身前。

她身姿单薄，脚步还有些虚浮，却仍走到了队伍前面。眼角未干的泪痕没有被掩住，心口那缕微弱魂息更让她不敢分神。牢大留下的牵挂支撑着她，可接过队长的位置，并不意味着她已经不再悲伤。

唐三微微一怔：

\begin{quote}
``须坤，你要做什么？''
\end{quote}

\begin{quote}
``它们刚刚受了伤，不能只凭我们一句话就放下戒备。''蔡须坤看了一眼梅天狸，轻声道，``我能让它们认出我的本源，先停下攻击。至于信不信我们，还得看我们接下来做什么。''
\end{quote}

话音落下的刹那，蔡须坤不再压制自身血脉。

嗡------！

温润的金色瑞光自她体内涌起，穿过近旁的林间雾霭。她没有任由气息向远处扩散，只将它缓缓送到狼群面前，让那股属于帝皇瑞兽的本源气息清楚地显露出来。

细碎的金光在她身侧浮沉，古老而温和的气息缓缓铺开。那是星斗魂兽熟悉的祥瑞气运，与寻常人类的魂力迥然不同。它能滋养生灵，却不会替任何一头魂兽作出选择。

蔡须坤身后，一道鎏金的帝皇瑞兽虚影缓缓凝现，与她的呼吸相应。她仍保持着人形，显露的是血脉本源的投影。

\textbf{瑞兽本源显现！}

整片林地瞬间寂静无声。

原本低吼不止的狼群渐渐收住了逼近的脚步。几头年长的狼兽辨认出金光中的气息，缓缓垂下头；护着幼崽的母狼仍留在原处，没有立即靠近。敬意压住了它们的攻击冲动，方才发生的重创却并未从记忆中抹去。

就连奄奄一息、肉身根基重创的灭世狼王，也艰难抬起头。它黯淡的兽瞳中掠过震动，随后收敛残存的威势，向蔡须坤微微低首。它认出了瑞兽气息，却仍将一缕目光留在梅天狸身上。

\textbf{帝皇瑞兽！}

承载祥瑞气运、能为魂兽修行带来裨益的帝皇瑞兽。

狼族的旧传中，瑞兽出现往往意味着附近魂兽能得到气运滋养，因而受到诸多族群珍重。这样的地位来自它对生灵的裨益，并不意味着它的修为高过所有凶兽，更不意味着它有权替受伤者免除仇怨。狼群认出了她，终于愿意暂缓对峙，听一听她要说什么。

\begin{center}\rule{0.5\linewidth}{0.5pt}\end{center}

不等众兽回神，蔡须坤抬手轻展掌心。

金色瑞气缓缓洒落，先停在狼群身前。见几头伤狼没有再退，蔡须坤才将气息送近，替它们舒缓受损脉络。几头幼狼的呼吸渐渐平稳，狼王的伤势却没有因此逆转；那处被打碎的肉身根基，仍在不断流失生机。

\begin{quote}
``我是帝皇瑞兽，也愿意为今天的话作保。''
\end{quote}

蔡须坤收住外放的气息，声音清楚地传到近旁狼兽耳中。

\begin{quote}
``我知道，你们受的伤不是假的，也不会要求你们因为认出我，就当作什么都没发生。''

``梅天狸长老出手重创狼王，是为了赶去救我们。她愿意承担自己那一击的责任，我们也不会把你们眼前的伤，推给别的人。''
\end{quote}

她直视前方的狼群，声音没有拔高。

\begin{quote}
``我们这次来，不为猎取魂环，也不会趁伤夺骨、强迫任何魂兽献祭。救治之后，你们仍可拒绝让路。若必须自保，我们会设法脱身，不把你们的族群当作报复的对象。''

``我们想去核心，为牢大仅存的魂息寻找生机。若有需要你们相助的地方，也会先问过你们。''
\end{quote}

她说完便停了下来，没有让瑞光在狼群身上留下任何强制烙印。

梅天狸走到她身侧，俯身将又一缕疗伤灵力送入狼王伤处。

\begin{quote}
``那一击是我出的。即便你不原谅，我也不会收回已经答应的救治。''
\end{quote}

灭世狼王注视她许久，才向护在身前的群狼发出一声低鸣。最前面的几头狼退开半步，给伤者与施救者之间留出一条窄路。

此刻的狼群，是敬重瑞兽，也愿意暂时观察这些人类；这与彻底信任，仍有距离。

几头成年狼仍守在幼崽身旁，眼底留着审慎。过去那些猎杀与失信，它们都没有忘记；眼下停战，是狼王与族群共同作出的选择。

即便蔡须坤已经作保，狼群仍在默默观望------它们要看，这支小队说过的话，能不能落到实处。

蔡须坤洞悉狼族心底最后的犹疑。

她清楚，本源气息能让狼群听她开口，真正能让它们放下戒心的，却是此后的行动。

无需她开口调度，经历牢大浴血殉道淬炼过的果糕小队，早已刻入骨髓般默契，全员不约而同运转力量，展开一场极致纯粹、毫无私心的团队协作。

狼群安静地看着，渐渐允许更多队员走近伤者。

\textbf{唐三}放出蓝银草，清理倒伏的细枝，为伤狼垫出干燥的卧处。草茎绕过伤口，只负责固定包扎用的布条；遇到不肯让他靠近的狼兽，他便停手等在一旁，没有强行伸过去。

\textbf{霍雨浩}收窄精神探测，留意近旁林木与风声。他将发现的陌生气息告知狼王，由熟悉领地的狼兽判断是否需要戒备，既没有擅自清剿周边魂兽，也没有为了替谁站岗，继续透支尚未恢复的精神力。

\textbf{小舞}蹲在一头受伤母狼附近，用清水浸湿布条，先让它闻过，才慢慢擦去外伤旁的泥污。母狼将受惊的幼崽拢到腹下，舔了舔它的耳朵，小舞便留出位置，不再惊扰它们。

\textbf{奥斯卡}分出恢复香肠，先向狼王说明只能帮助恢复体力，治不了根基重伤；\textbf{卢本韦、诺当}则与狼族守卫商量，在它们指定的边界轮流警戒。\textbf{蔡兴坤}收好散落的药布，将清水送到众人手边，又把一件外衣披在蔡须坤肩上。

整支小队攻防有度、补给有序、彼此兜底、同心共生，没有一人争抢机缘、没有一人觊觎兽核、没有一人借机施压。

两位长老的实力足以压住眼前的狼群，小队成员却都还带着伤，谁也没有将停战视作理所当然。他们放慢动作，等狼兽自己决定是否接受帮助。

狼族众兽静静伫立，全程看在眼里、感在魂中。

它们见过无数人类队伍：得利则进、无利则退，遇兽则猎、遇弱则欺。

可今日，它们亲眼看见------人类小队可以并肩守护、可以无私救赎、可以负重存善、可以心怀苍生。

瑞兽的气息让它们愿意听，眼前这一件件小事，则让这次停战渐渐有了信任的分量。

几头伤势较轻的狼兽主动上前，低头嗅了嗅队员手中的药布。有一头幼狼碰了碰小舞的袖角，母狼这次没有将它立刻叼回去。并非所有狼都已放下防备，但双方终于可以在同一片林地中安稳停留。

这份信任先在眼前的狼群与小队之间生根。至于森林更深处、从未见过他们的其他族群，还需要另外争取。

\begin{center}\rule{0.5\linewidth}{0.5pt}\end{center}

对峙渐缓后，\textbf{霍雨浩}再次来到灭世狼王身边。他先与梅天狸确认伤势，又借天梦冰蚕的协助辨认狼王的魂息，迟迟没有立即开口。

蔡须坤看出他的犹豫：

\begin{quote}
``你想到了别的办法？''
\end{quote}

霍雨浩点头，目光仍落在狼王身上。

\begin{quote}
``有一条路，但必须先把风险说清楚。''
\end{quote}

\begin{quote}
``狼王前辈。''
\end{quote}

霍雨浩语气诚恳，不卑不亢：

\begin{quote}
``长老救不回的，是你已经崩坏的肉身根基。可你的主意识还清醒，魂息也没有完全散掉。我们先前摸索过魂灵共生之法，或许能趁这部分尚在，将它保存下来。''

``这不是治好原来的身体，而是换一种承载魂息的方式。一旦开始，能不能稳住契约，要靠双方配合，我也不能保证一定成功。''
\end{quote}

灭世狼王微微抬眸，眼底带着一丝疑惑与希冀。

\begin{quote}
``若你愿意，可以与须坤尝试缔结魂灵契约。她的瑞气有助于养护伤魂，你也能保留自己的意识。但今后的恢复需要时间，你们会互相影响，她若遭遇重创，你也会受到牵连。''

``你可以拒绝。即便答应，也不等于把全部力量交出来。契约的边界，必须由你们共同确认。''
\end{quote}

此话一出，全场皆惊。

狼王望了一眼守在身侧的族群。没有人催促它，也没有人伸手索取魂环。对它而言，这条路仍有风险，却比等待肉身衰败，多了一种能够继续守望族群的可能。

灭世狼王静静凝视着身前圣洁庄严的帝皇瑞兽虚影，心中思绪翻涌。

它在这片领地生活了漫长岁月，一直将守住族群看得比自己的生死更重。重伤让它愤怒，也让它明白眼下身体的极限。面前的梅天狸没有逃避责任，其他人也没有趁伤逼迫；它愿意考虑这个提议，却仍有自己的要求。

\begin{quote}
``我若离开肉身，不能再替它们守在这里。你们不许借我的名字驱使狼群。日后若有机会，我要回来看看。''
\end{quote}

蔡须坤认真答应后，灭世狼王才缓缓低头，将一缕清晰的意念送到她面前。那缕意念没有求饶，只有它反复确认后的决定。

它愿意尝试魂灵共生，保存尚未散去的意识，与蔡须坤相伴前行。它仍记得梅天狸造成的伤，也承认她此刻正尽力补救；接受这条生路，并不意味着往事从未发生。

\begin{quote}
``我记住了。''蔡须坤轻轻颔首，``今后我们一起养好伤。你的族群，不归我支配。''
\end{quote}

霍雨浩借三道魂灵相助，先在狼王清醒的意念与蔡须坤的精神回应之间搭起联系。梅天狸维持狼王最后的肉身生机，ddfg守住外围，隔绝杂乱气息。蔡须坤缓缓送出瑞气，没有去拉扯狼王尚未放开的那部分魂息。

第一次牵引便令她胸口发闷。霍雨浩立刻让她停下，等狼王主动收束意念，才重新接续。双方一点点确认共生的联系，直到那道魂影终于能离开伤躯，停在她身侧。

狼王的旧躯渐渐失去生息，魂影也远比生前淡薄。它低低呼唤一声，群狼听见熟悉的声音，才慢慢围了上来。

蔡须坤下意识按住心口。狼王保住了意识，牢大的那缕残魂却仍沉寂着。

\begin{quote}
``伤势不同。''霍雨浩看见她的动作，低声道，``狼王缔约前，意识还完整。牢大连维持自我的魂息都只剩一点，我们不能拿他的残魂来试。''
\end{quote}

她点了点头，先将那份急切压住，重新护好心口的微光。

\begin{center}\rule{0.5\linewidth}{0.5pt}\end{center}

嗡------！

契约初步稳住后，狼王保留下来的魂力沿共生脉络缓缓流动。霍雨浩不断提醒蔡须坤放慢，绝大部分魂力仍由狼王收束，用来维系自身；只有双方能够承受的一小部分，渐渐与她原有的序力相合。

蔡须坤原本枯竭的序力终于开始回转。此前修行中已有所理解、却迟迟无法稳住的那道关隘，在反复调整共生联系时逐渐清晰。她没有停下梳理，直到新增的力量能顺着原有脉络运行，才试着向前迈出一步。

她身周的金光亮起又收住，紊乱的序力一段段接续，终于跨过眼前的瓶颈。

\NovelDisplay{M00477R01}{%
\psi(I_\omega)%
}

她触及了这个此前已有所参悟的层级。新成的脉络还很生涩，需要继续稳固，远不到能任意挥霍力量的时候。

序力在经脉中缓缓周转，比起方才有了支撑。梅天狸检查过她的伤势，叮嘱她暂时不要强行出手。狼王的魂影也只在外停留片刻，就回到共生脉络中休养；它保下了意识，离恢复生前的状态仍很远。

果糕小队众人看着焕然一新、气场沉稳的新任队长，眼底满是欣慰与敬佩。

梅天狸与ddfg相视一眼，皆是暗自点头，眼底满是赞许。

蜕变落幕，神光敛去。

蔡须坤缓缓睁眼，先确认狼王的魂息已经稳住，随后又将意识轻轻探向心口。牢大的残魂没有因为她的突破而苏醒，仍需她小心养护。那份无力没有消失，她能做的，只是将这一步走稳，再继续找下一条路。

她抬眸望向远方云海，声音清亮坚定，响彻林间，既是自语，也是对全队、对逝去之人的郑重承诺：

\begin{quote}
``牢大留下的魂息还在。''

``我们要继续找让他回来的办法，也要让每个人平安走出这里。今天你们帮我撑住了这一步，往后，我们一起走。''
\end{quote}

字字铿锵，掷地有声。

悲恸沉淀为力量，遗憾化作前行的铠甲。

队长的位置仍让她感到沉重。可这次她抬起头时，身旁已经站着愿意和她一同分担的人。

\begin{center}\rule{0.5\linewidth}{0.5pt}\end{center}

蔡须坤转身，目光扫过唐三、霍雨浩、小舞、奥斯卡、卢本韦、诺当与蔡兴坤，又望向两位长老。

\begin{quote}
``先把伤势稳住，跟狼群交代好，再往核心走。我们去拜见帝天，为牢大问一问。''
\end{quote}

风声再起，林海翻涌。

等狼群将狼王的旧躯安置妥当，午后的光已从枝叶间落下。狼群没有再挡路，几头熟悉林间路径的成年狼送他们走出领地，才在交界处停住。蔡须坤回身道别，随后与伙伴们一起，向星斗更深处前行。

\begin{center}\rule{0.5\linewidth}{0.5pt}\end{center}

辞别狼族领地，果糕小队一行人的身影渐渐沉入星斗大森林的幽深腹地。

越往核心深处行进，周遭的天地灵气便愈发浓郁醇厚，近乎凝成实质，在林间缓缓流淌。可与之相伴的，是层层叠叠、愈发厚重的苍茫死寂，没有半分外围林地的鲜活喧闹。沿途草木苍劲古老，参天古木遮天蔽日，将天穹彻底遮蔽，林间光线昏暗幽深，每一缕风掠过林海，都裹挟着远古荒古的厚重气息，压得人心神沉静，不敢肆意妄动。

走出狼族领地后，偶遇的魂兽有的退入林间，有的停在远处审视他们。瑞兽气息能减少一部分敌意，却不能让所有族群自动放行。一行人避开巢穴，没有追赶退走的魂兽，也没有将狼王的名声当作通行令。

沿途没有再发生正面冲突，仍不能放松警惕。霍雨浩只在需要时展开探测，两位长老则交替走在队伍前后，遇见不宜靠近的气息，便带着众人绕行。

\begin{center}\rule{0.5\linewidth}{0.5pt}\end{center}

夜幕笼罩林海，古木遮去星月，只有零星灵草的微光点缀前路。连续数个时辰的赶路，加上尚未痊愈的伤势，让队员们的脚步渐渐沉重。梅天狸看了看队伍，没有再催他们前进。

所有人的脚步都悄然放缓，呼吸粗重，额角布满细密冷汗，四肢百骸皆是酸胀疲惫。连日的死战、离别、和解、突破层层波折叠加，早已耗尽众人大半心神，极致的倦意如同潮水般反复席卷而来。

梅天狸见状，轻声开口体恤众人：

\begin{quote}
``原地休整片刻，连夜赶路损耗过大，养好精神，我们再继续前行。''
\end{quote}

众人在背风处安顿下来，分食恢复体力的食物。ddfg先巡过近旁的林地，梅天狸则坐在能照看全队的位置，两人商定轮流守夜。确认有人警戒后，伤重的队员才陆续靠着古木睡下，林间渐渐只剩呼吸声与枝叶轻响。

\begin{center}\rule{0.5\linewidth}{0.5pt}\end{center}

蔡须坤靠在苍老的树干上，将瑞气缓缓收拢。她听见梅天狸低声让她安心休息，才稍稍松开一直绷紧的肩背。疲惫、思念与白日勉强压住的悲痛一齐涌上，她护着心口那缕魂息，缓缓睡了过去。

坠入梦境的一瞬，所有压迫、冰冷、沉重、死寂被彻底剥离。

这里没有幽深密林的昏暗，没有前路未知的惶恐，也没有生死离别的刺骨寒意。整片天地温柔得近乎不真实，暖融融的柔光铺遍大地，流云缓慢舒展，晚风轻软似絮，连空气都带着昔日并肩休憩时的草木清甜。时间仿佛被彻底凝滞，万物静默，唯余安稳。

唯有一片澄澈温暖的晴空，晚风轻柔，草木安然，是他们无数次并肩历练、安稳休憩的寻常林间。

一道熟悉到刻入神魂、入骨难忘的身影，静静立在柔光深处。身姿挺拔，肩背宽阔，眉眼坦荡温柔，不染半分杀伐戾气，更无战场殉道的惨烈苍白。此刻的牢大，完好无损、安然无恙，是蔡须坤记忆里最温暖、最可靠、永远挡在全队身前的模样，真实得让她心口骤然发酸。

他没有战死沙场的惨烈伤痕，没有燃命殉道的虚弱苍白，一如从前无数次护在队伍前方的模样，温柔又坚定。

蔡须坤瞳孔剧烈震颤，心口骤然一紧，压在心底的委屈与思念瞬间决堤。她脚步发虚，下意识想要奔上前拥抱那道身影，可梦境的时空朦胧虚幻，脚下似踏在绵软流云之上，每一步都轻飘飘的，始终无法靠近。她不敢眨眼，生怕这唯一的重逢幻影转瞬破碎。

牢大望着她，眼底带着温柔的笑意，声音平和厚重，和从前无数次安抚全队、指引前路时别无二致：

``须坤，我看到了，你现在是新的队长了。'' ``你做得很好，你没有让队伍散掉，也没有让我的意志断掉。''

他缓缓开口，字字落在蔡须坤心底，清晰无比：

``你救下了狼王，也走过了眼前这一道关。可新得的力量还要慢慢练，别一有了起色，就又把自己逼到最后。'' ``我从前护着你们，往后，你也要学会让他们护着你。队长不能替所有人走完路。''

他顿了顿，目光微微低垂，声音忽然轻了下去，像是怕惊碎什么易碎的东西：

``须坤，我知道你很难。'' ``一闭上眼睛，就会想起我消散的样子，对不对？'' ``你总觉得，是你不够强，才让我替你挡下了那一击。你总觉得，是你留住了我的残魂，却救不回我的命。你总觉得，你活着，是欠了我的。''

蔡须坤浑身剧烈一颤，眼泪夺眶而出。

每一个字，都像烧红的针，精准刺入她心底最隐秘、最溃烂、从来不敢触碰的伤口。

她张了张嘴，想要否认，想要说''不是的''，可喉咙像被一只无形的手死死扼住，一个字也说不出来。因为她知道，他说的每一个字，都是真的。

牢大看着她崩溃的模样，眼底的笑意没有消失，却多了一层极深极深的疼惜。他缓缓抬起手，像是想要像从前那样揉一揉她的发顶，可指尖伸到半途，又停住了------他知道，他碰不到她。

手悬在半空，停了三息。

然后，他缓缓收回，垂在身侧。

``须坤，听我说。''

他的声音依旧温柔，却带上了一丝近乎恳求的郑重：

``最后那一刻，我只想让你们走出去。''

蔡须坤猛地抬头，泪眼模糊地望着他。

``你把能做的都做了，连我最后的一点魂息也护住了。须坤，别再把所有没能挽回的事，都算成自己的错。''

他的声音很轻，像许多次危急过后坐在她身旁，等她慢慢平复呼吸时那样。

``那是我自己的选择。我想护住你们，也想看见你们往后好好活着。''

蔡须坤的膝盖一软，跪倒在虚无的梦境之中。双手死死捂住嘴，不让自己哭出声来，可滚烫的泪水从指缝间疯狂涌出，砸落在金色的草地上，化作点点碎光。

她想伸手抱住他，像过去许多次受伤后那样，抓住他的衣袖确认他还在。可她的手穿过那道身影，只抓到一把温热的流光。

牢大低头看着她，眼底终于浮上一层薄薄的水光。

他蹲下身，与她平视。明明触碰不到，他却像是在隔着万里山河，轻轻捧住了她的脸。

``须坤，你听好了。'' ``你留住了我的残魂，不是欠我的。你已经护住了最后的一点可能。'' ``可你不能只为了这一点可能，就把自己熬坏。还在你身边的人，也都在担心你。''

他的声音忽然哽咽了一瞬，又迅速压了下去，重新扬起那个她最熟悉的、永远沉稳笃定的笑容：

``你答应我一件事。'' ``以后难过的时候，不要一个人扛。唐三在，霍雨浩在，小舞在，所有人都在。你护着他们，他们也护着你。别像我一样，什么都往自己身上揽。''

他站起身，后退一步。

身后的柔光开始缓缓收拢。她想挽留，身旁的草木却已一点点失去轮廓，梦境的边缘渐渐崩塌。

``须坤，该醒了。''

蔡须坤猛地抬头，瞳孔骤缩，所有理智在这一刻彻底崩碎。

``醒来以后，先看看身边的人。''

他笑着，眼底却分明有泪光一闪而过。

``别只顾着追一道影子。''

``你还有自己的路要走。''

``往前走的时候，也记得让他们陪着你。''

话音落尽，他的身躯开始变得通透、轻盈，周身溢出点点纯白微光，如同风中残烛、晨雾流光。他没有回头，只是静静朝着天光尽头走去，背影温柔却决绝，一步步淡化、虚化、消融在朦胧柔光之中。

蔡须坤疯了一样扑上去，双手拼命去抓他消散的衣角，指尖穿过一缕又一缕流光，什么都抓不住。她跪在地上，一遍又一遍地伸手、抓握、扑空，像一个溺水的人拼命想要抓住最后一根浮木，却只能眼睁睁看着它从指缝间滑落。

``牢大------！'' ``你别走！求你了------！'' ``我不要当队长，我不要当英雄，我只要你活着------！''

她的嘶吼声撕裂了整片梦境的温柔，可回应她的，只有越来越淡的光点，和越来越远的背影。

最后一点微光消散的瞬间，她恍惚看见他侧过脸，嘴唇微动，无声地说了一句话。

她没有听清。

但她读懂了口型。

------``好好活。''

``牢大------！！''

一声嘶哑到变调的悲鸣撕裂夜空，蔡须坤骤然惊醒，冷汗浸透衣衫，心口剧烈抽痛，呼吸急促滚烫。眼底的温热柔光彻底消散，残留的掌心余温瞬间冷却，只剩现实林海刺骨的夜风与无边幽暗。

她下意识地攥紧胸口衣襟，指尖死死按在心口那缕残魂之上，像是怕它也跟着梦境一同消散。指节泛白，浑身止不住地颤抖。

夜色依旧深沉。梅天狸在不远处抬起头，先确认周围没有异动，才低声唤她的名字。蔡兴坤被惊醒，摸索着递来水囊；霍雨浩也撑起身，问她心口是不是又疼了。

蔡须坤没有立刻回答，先去确认那缕魂息。微光仍在原来的地方，安安静静地受着瑞气养护，没有随梦中的身影消散。

她不知道这场梦有多少只是自己的思念。那几句话却留了下来，连同伙伴递到掌心的温水，一起将她从方才的惊惶中慢慢拉回。

梦醒，人空，余悲难歇。

无数并肩作战的画面、无数被他守护的瞬间、那场惨烈的殉道离别，一幕幕翻涌心头，尖锐的酸涩与悲痛席卷全身，几乎将她彻底淹没。失去队长的遗憾、背负执念的沉重、孤身扛队的孤勇，万般情绪交织缠绕，压得她心口发痛。

她微微低头，指尖微微颤抖，眼底泪光闪烁，无尽的悲伤几乎要冲破防线。

她想像白日那样迅速压住泪水，刚抬起头，却见蔡兴坤已经坐到她身旁。

她想起梦中牢大最后的笑容，想起他让自己看看身边的人。霍雨浩没有催问，蔡兴坤也没有叫她坚强，只是陪着她，等她把那口哽在喉中的气慢慢顺下去。

\begin{quote}
``我梦见牢大了。''她终于说。
\end{quote}

声音一出，泪水又落下来。她没有松开护着残魂的手，也没有再将脸转向没人看见的地方。

蔡兴坤靠近了一点。霍雨浩低声应了一句，起身将快要滑落的外衣重新搭好。两人都留在她身边，直到她终于能够平稳呼吸。

逝者已矣，执念长存，意志不灭。

她仍要寻找牢大归来的办法，仍要护好这缕来之不易的魂息。狼王的契约不能让牢大立即重生，新得的序力也不能替她免去悲伤；但前路不再需要她一个人扛。

她轻轻握了握蔡兴坤的手，终于肯靠回树干，再休息片刻。

夜色将尽，天光微亮。

\chapter{第四十六章：龙域觐见，问情天机}\label{ux7b2cux56dbux5341ux516dux7ae0ux9f99ux57dfux89d0ux89c1ux95eeux60c5ux5929ux673a}

\begin{center}\rule{0.5\linewidth}{0.5pt}\end{center}

蔡须坤深吸一口气，压下所有心绪，目光重新望向星斗核心的幽暗迷雾，身姿挺拔如松，静待全员休整完毕，再度踏上前路。

可越是深入腹地，众人心中的凝重便愈发浓烈。

\textbf{霍雨浩}留意着近旁动静，只在路况难辨时送出一缕精神探测。此刻他忽然停下，低声提醒众人：

\begin{quote}
``大家小心，前方气场不对劲。这片区域的天地规则极度凝练，威压层层叠加，绝非普通十万年魂兽所能掌控。''
\end{quote}

霍雨浩的精神力远比寻常同阶魂师敏锐，此时却还没有从连日的透支中完全恢复。前方迷雾里，一股沉重力量截住了他的探测。他立刻收回那缕精神力，额角仍渗出冷汗，不敢再贸然向内试探。

\begin{quote}
``是王者气息。''
\end{quote}

已然化作魂灵、蛰伏在蔡须坤体内的\textbf{灭世狼王}，此刻传出一道带着极致敬畏的意念音，回荡在蔡须坤的心神之间，语气满是臣服与惶恐。

\begin{quote}
``星斗大森林真正的主宰，万兽共主的代言人------\textbf{金眼黑龙王帝天}，就坐镇在这片核心秘境深处。''
\end{quote}

灭世狼王拥有二十万年修为，对星斗核心的气息并不陌生。它知道面前是谁，也清楚双方的差距；仅剩魂灵之身的自己，更不宜在这里强行释放威压。

\begin{center}\rule{0.5\linewidth}{0.5pt}\end{center}

\textbf{帝天}，统御整片星斗万千魂兽的无上至尊，是扎根星斗本源、横贯万古岁月的魂兽共主代言人。

他有着\textbf{八十九万年}的深厚修为，是星斗大森林最强大的守护者之一。各族对他的敬畏，既来自黑龙的威势，也来自他在漫长岁月中的守护。可众人并无资格仅凭远处的一缕气息，就断言他与魂殿大长老谁高谁低。森林中的族群也各有领地与生存法则，并非时时刻刻都受他一念驱使。

狼王曾在自己的领地内统领一族，如今感受到这股气息，才更清楚眼前的差距。它只将辨认出的情况告知蔡须坤，没有贸然替帝天作出任何承诺。

\begin{center}\rule{0.5\linewidth}{0.5pt}\end{center}

随着众人前行，林间的风渐渐弱了。来自核心的龙威拦住道路，没有朝外围无止境扩散，却足够令靠近的人停下脚步。霍雨浩示意大家止住探测，唐三也将蓝银草收回近旁，避免被当作试探禁区的举动。

即便隔着数里密林，这股威压仍令伤后的队员呼吸发紧。奥斯卡扶住身旁的诺当，蔡兴坤则看了看蔡须坤，先将脚步放慢。众人没有硬闯，在两位长老的示意下暂时退到一处空地。

\textbf{小舞}下意识攥紧手心，轻声道：

\begin{quote}
``好恐怖的威压\ldots\ldots 这才是星斗大森林真正的底蕴吗？''
\end{quote}

\textbf{唐三}握紧了收回的蓝银草，额角微微见汗，低声道：

\begin{quote}
``先别继续往前了。我们带着伤，不能在这里强撑。''
\end{quote}

他将身旁的枝条轻轻拨开，给后面的队员留出落脚的位置，没有再向林中试探。

\textbf{蔡须坤}走到空地前缘，掌心缓缓透出金色瑞光。她没有用它去冲撞龙威，只试着将自己的气息送过去，让核心深处的守护者辨认。

这股气息对帝天而言并不陌生。她曾从他身边离开，如今却带着人类同伴、一位受伤的新魂灵，以及一道残存的人类魂息回来。她不知道他会怎样看待这些变化，只能先把自己的身份表明。

她望向林中，将手示意众人停在身后。

\begin{quote}
``他认得我。''蔡须坤轻声道，``先让我通报，不要再往里面探。我们是来求助的。''
\end{quote}

\begin{center}\rule{0.5\linewidth}{0.5pt}\end{center}

就在全队停步等候之际，梅天狸从后面走到队前，开口唤住蔡须坤。

是\textbf{梅天狸}。

她抬眸望向深处迷雾，神色郑重。帝天的力量远在小队之上，连她也不愿贸然试探。蔡须坤与他相识，可以争取开口的机会，却不能代所有同行者保证对方一定放行。

尤其狼王的伤由她造成，此事不能只靠蔡须坤从中担保。若帝天问起，她必须亲自交代。

\begin{quote}
``须坤，诸位队员。''
\end{quote}

梅天狸将护持队员的灵力收在身侧，没有向林中释放威压。她看向蔡须坤，又扫过其余队员，语气郑重。

\begin{quote}
``这里已经接近核心，不能擅闯。''

``狼王的伤因我而起，待会儿我会当面说明。若对方不许继续前进，我们先退回来，再想办法。''
\end{quote}

她目光扫过众人，字字恳切，郑重做出决断：

\begin{quote}
``我与你们一同拜见他。''

``我能做的是护好你们、把经过说清楚，不是保证任何风险都不存在。你们也不必为了求一句答复，强行承受不必要的威压。''
\end{quote}

一旁的\textbf{ddfg}也微微颔首，沉声道：

\begin{quote}
``我在后面照应。''ddfg道，``先等回应，有了许可再走。''
\end{quote}

有两位长老留在身侧，众人稍稍定下心来。他们仍停在原地，没有将这份照看当作闯入核心的把握。

历经生死离别、磨合成长，又得长老全力守护，果糕小队众人心中暖意涌动，原本紧绷的心弦也稍稍松弛。

蔡须坤转头望向二人，郑重颔首：

\begin{quote}
``多谢两位长老。''
\end{quote}

蔡须坤等众人站稳，才朝林中开口。

\begin{quote}
``大家先等在这里。''蔡须坤向林中开口，``帝天，我是蔡须坤。我带了几位同伴来，想请你看一看一缕受伤的魂息，也想把狼王的事向你说清楚。''
\end{quote}

瑞光在空地前缘停住。数息后，前方厚重的龙威缓缓分开一条通路，林间传来低沉的回应。

\begin{quote}
``让她过来。同行者收敛魂力，不得惊扰核心。''
\end{quote}

蔡须坤这才回头示意，带众人沿那条通路缓缓前行。

\begin{center}\rule{0.5\linewidth}{0.5pt}\end{center}

星斗核心腹地，万古死寂沉沉笼罩四方。

越接近核心，越能清楚感受到林中力量的厚重。那股力量没有主动碾向他们，却令伤后的众人难以轻松。梅天狸与ddfg放慢脚步，帮队员稳住气息；蔡须坤也没有催促，等落在后面的人跟上，才继续向前。

梅天狸始终护在队伍身侧，用灵力帮伤重的队员调匀呼吸。即便对方已收敛威压，这片核心区域常年沉积的气息仍然厚重。她不敢疏忽，留意着每个人的状态。

ddfg走在后方，看了一眼核心深处，低声道：

\begin{quote}
``仅从这里，根本看不出他力量的深浅。''
\end{quote}

梅天狸缓缓颔首，语气里带着发自心底的敬畏：

\begin{quote}
``我也摸不透。''梅天狸低声道，``只凭这点感知，判断不了他的真实战力。等见到他，不要拿外面的传闻来试探。''
\end{quote}

梅天狸见过不少强者，此刻却不愿为帝天的力量仓促下结论。她将注意放回队员身上，确认每个人都还能跟上。

\begin{center}\rule{0.5\linewidth}{0.5pt}\end{center}

全队屏息凝神，步步谨慎前行，整片秘境死寂得令人心生压抑。就在这时，秘境最深处的幽暗虚空缓缓蠕动，打破了长久的沉寂。

一双横贯丈余、鎏金璀璨的龙之竖瞳，自无尽黑暗中缓缓睁开。

金色眸光穿过雾气，先停在蔡须坤身上，随后才移向她身后的两位长老与其余队员。周围龙气轻轻翻涌，众人同时停下脚步，没有继续靠近。

帝天守护星斗漫长岁月，见过许多为魂环与魂骨而来的人类，因此对陌生队伍始终保有戒备。但蔡须坤是他认识的瑞兽，他没有将她与同行者一并当作入侵者驱逐。

他愿意听，并不意味着已经相信。尤其蔡须坤身侧多出的那缕狼王魂息，令他的目光停留得更久。

漆黑龙气在湖畔收拢，显出一具庞大的龙影轮廓。帝天没有出手，目光却始终留意着狼王魂息与蔡须坤之间的联系。

\begin{quote}
``你回来了。''

``先前离开的时候，你身边没有这道魂息。狼王的身体在哪里？这些人，又为什么跟着你来到这里？''
\end{quote}

低沉的声音回荡在林间，听不出欢迎，也没有立即问罪。小队众人屏息站定，知道接下来的解释不能含糊。

梅天狸与ddfg护在队伍两侧，没有展开攻势。蔡须坤向前半步，将一直收敛的瑞兽气息缓缓展开，让帝天看清她此刻的状况。

\begin{center}\rule{0.5\linewidth}{0.5pt}\end{center}

唯独\textbf{蔡须坤}立身队前，身姿挺拔，无惧无退。

她看见帝天的目光，知道他最在意的是什么。狼王成为魂灵，是他们共同确认过的选择，不能只凭她一句自愿便带过。

心念既定，蔡须坤不再刻意压制自身血脉与气运。

嗡------！

金色瑞气在她身周流转，没有冲向天空，也没有试图与龙威争持。帝天顺着那股气息，辨认出她仍未复原的伤势，以及被护在本源深处的另一点魂光。

蔡须坤将这道询问传入共生联系，等狼王回应。一缕墨色魂影随后缓缓浮出，停在她身侧。

灭世狼王的魂影还很淡，身周没有生前完整的威势。它自行走到蔡须坤身侧，抬头望向帝天。只是维持这次显形，便让轮廓轻轻晃了一下。

它向帝天低首致意，没有躲到蔡须坤身后，随后主动送出一缕清晰的意念。那是它自己的声音，不受蔡须坤代答。

\begin{center}\rule{0.5\linewidth}{0.5pt}\end{center}

帝天的竖瞳微微收缩。让他意外的不是蔡须坤的人形，也不是早已认识的瑞兽血脉，而是这头原本固守领地的凶兽，竟能以自己的意识回应，与她维持一道稳定的共生联系。

他也留意到霍雨浩身边的魂灵气息。对魂灵之法，他并非全无听闻；但旁人的传闻，与亲眼看见狼王在肉身重创后仍能保持自我，终究不同。

他没有因此就断言人兽恩怨可以尽解，而是先辨认那道契约有没有强迫、抽取与束缚的痕迹。

帝天再次开口，这一次，问题直接落在新成的契约上。

\begin{quote}
``我知道你是谁，也没有忘记你离开时的模样。''帝天看向蔡须坤，``我问的是这道联系。你们有没有逼迫狼王，借救治之名取走它的力量？''
\end{quote}

蔡须坤没有抢着回答，只望向狼王。帝天也将目光移到它身上，等它自己开口。

\begin{quote}
``若有不能由你自己作主的地方，现在说。''
\end{quote}

他转向立在蔡须坤身侧的狼王，继续问道。

\begin{quote}
``还有，你的伤究竟从何而来？我不接受把一切推给人兽旧怨的说辞。''
\end{quote}

\begin{center}\rule{0.5\linewidth}{0.5pt}\end{center}

面对万古至尊的恳切发问，灭世狼王不敢有半分怠慢，上前微微俯首行礼，以魂灵独有的神魂共鸣道音，清晰响彻整片秘境。

它先说明自己为阻拦人类强者穿过领地，与急于救援的梅天狸交手，被她重创根基。梅天狸没有回避，站出来承认自己的出手与当时的判断，也承认未能控制好力道。

狼王随后讲起众人折返救治、蔡须坤显露本源，以及霍雨浩提出魂灵契约的经过。它没有将自己说成毫无疑虑，也没有替梅天狸抹去责任，只将自己最终同意的理由交代清楚。

\begin{quote}
``我若不答应，他们也不会取我的魂环。契约中没有驱使狼群的条件。我还能自己说话，也能留住力量养伤。这是我愿意试的路。''
\end{quote}

至于帅气教对魂兽的猎杀，则是另一件事。小队讲述了刚刚遭遇的围杀，霍雨浩补上自己见到的经过，将二者分开说清，没有用那场袭击掩盖狼王受伤的原因。

狼王的声音时断时续。帝天没有催促，等它将能够确认的经过讲完，才又开口询问。

帝天听完，又向狼王确认了几处契约细节。见它能独立收回魂息，才缓缓敛下紧绷的龙气。

\begin{quote}
``你们救下了它，这一点我承认。造成的伤，也不能因为事后救治就一笔勾销。''
\end{quote}

梅天狸应了一声，没有为自己争辩。

帝天收回审视契约的目光，谈起众人所说的那场围杀。

\begin{quote}
``原来如此。''

``近些年屠戮我星斗子民最多、祸乱世间最深的祸源，正是你们所言的帅气教。''帝天语气冰冷，藏着积压许久的万古怒意，``其教众贪婪暴戾、杀伐无度，四处猎杀魂兽、掠夺兽核与本源精气，无数无辜子民惨死其手。教主胡赞杀伐滔天，双手沾满星斗生灵鲜血，罪孽滔天，罄竹难书。''
\end{quote}

帝天一直留意着森林外围的猎杀，却不能只凭模糊传闻就掌握所有据点。核心另有需要守护的事，贸然带走大批强者，也会留下空隙。这样的顾虑使他暂时留在这里，却不是束缚他永远不能离开的天地禁令。

\begin{quote}
``我会查清这些猎杀从何处来，也会让各族加强警戒。''帝天的声音冷了下来，``留在核心，不等于允许他们不断试探。真到了必须出手的时候，我自会出去。''
\end{quote}

话音一转，他的目光温柔落回蔡须坤身上，眼底满是悲悯与惋惜。

\begin{quote}
``先让我看一看你护着的那缕魂息。''

蔡须坤闻言，缓缓放开最外层的一线瑞气。帝天没有将残魂拉出，只隔着保护仔细辨认，片刻便示意她重新收好。
\end{quote}

\begin{quote}
``伤得很重。狼王是在意识完整时转为魂灵，你的同伴只剩一点本源魂息，二者不能用同一办法处置。若要让他真正回来，养魂之外，还须解决肉身与道基的问题。''

``就我知道的药物线索，\textbf{乾坤造化丹}值得继续寻找。但你们不能把它当作服下便能立刻复活人的保证。''
\end{quote}

蔡须坤指尖微微收紧。她记得ddfg此前说过，乾坤造化丹用于重构血肉，天地灵水用于滋养残魂、修复大道。帝天提起的药名，终于与那条渺茫的办法接上了。

\begin{quote}
``长老还说，需要天地灵水。您知道这两样东西的去处吗？''
\end{quote}

帝天没有立刻给出肯定的答复。

\begin{quote}
``关于乾坤造化丹，我听过的线索指向\textbf{乾坤问情谷}。至于天地灵水，我没有能交给你们的存物，也没有可靠的出处。两样东西如何配合你同伴的伤，还要找到材料之后再判断。''

``问情谷与神明之力有关，传闻会以情念试人。我没有亲自深入，知道的只是旧闻，不能替你们保证谷中一定还有丹药，更不能保证你们能取到。''
\end{quote}

帝天顿了顿，又道：

\begin{quote}
``我现今的序力在 \allowbreak{}\NovelInline{I02286}{\omega\text{-}Y(1,ω)}\allowbreak{} 上面，仍不会仅凭修为便断言能破开神明遗留的规则。你们更不能将这一点旧闻，当作可以立刻入谷的把握。''
\end{quote}

林中静了下来。众人听明白了：他们得到的是一条可以查证的方向，还不是拿在手中的救人之法。

全队众人默然无言，心头沉甸甸的压着一块巨石。前路本就荆棘遍布、危机四伏，如今又添上一处连万古龙王都忌惮万分的绝世险地，重重磨难接踵而至。

\begin{center}\rule{0.5\linewidth}{0.5pt}\end{center}

短暂的沉寂过后，蔡须坤先低头确认牢大的魂息仍然平稳，才重新抬起眼。她想立刻出发，却也明白身后所有人都还带着伤，不能再凭一腔急切将伙伴推入未知险地。

她仍会继续找，却要先把线索查实，再与同伴商量如何走。狼王才刚稳住契约，牢大的残魂更受不得反复折腾，眼下谁都不能再轻率冒险。

心绪落定，蔡须坤再度抬首，望向秘境深处那道至高无上的龙影，心中依旧藏着无数疑问，迫切想要探寻更多线索。

乾坤问情谷凶险莫测，连帝天这等万古至尊都没有十足把握，她必须尽可能摸清所有隐秘，为全队的秘境之行铺路，为牢大的复生搏出一线生机。

\begin{quote}
``龙王前辈。''
\end{quote}

蔡须坤出声请教，语气恭敬诚恳，带着满心探寻：

\begin{quote}
``晚辈还想确认：关于问情谷的位置、丹药的来历，或是怎样进入的记录，前辈可有可靠的卷册或见闻？我们想先查清楚，再决定如何前往。''
\end{quote}

这是她当下最迫切的诉求，只要能多一丝线索，前路的凶险便能少一分，全队便多一分胜算。

然而，面对她的追问，威震星斗、洞悉万古法理的金眼黑龙王帝天，却罕见地陷入了长久的沉默。

幽暗秘境之中，龙威缓缓收敛。帝天将自己知道的旧闻又想了一遍，终于摇了摇头。

\begin{quote}
``我没有可以交给你们的入谷记录。''

``方才说的是传闻中的危险与机缘。确切的入口、现存丹药和内部试炼，我都不能作准。你们应当另找亲历者留下的记载，相互印证。''
\end{quote}

帝天缓缓道出真相，语气坦荡无虚：

\begin{quote}
``我在星斗活得久，不等于知道大陆上一切秘境。更不能为了让你们安心，就把没有验证过的话说成事实。''
\end{quote}

此话落下，众人一时没有接话。帝天能为他们分辨残魂的现状，却不能替他们查清问情谷。接下来的路，仍需要找到更可靠的依据。

霍雨浩默默记下了帝天能够确认的部分，又将其余未明的地方留空。至少这些旧闻并非毫无用处；接下来找到古籍，便有了可以逐一对照的问题。

蔡须坤的目光暗了一下，随即郑重道谢。至少她已经知道，不能只凭这点传闻便带全队闯谷。

众人沉思之际，一道清柔却沉稳的女声打破了沉寂。

\begin{quote}
``等等，或许\ldots\ldots 我知道一点线索。''
\end{quote}

\textbf{卢本韦}缓步走出队伍，眉头微蹙。她还没有想起完整的内容，只记得家中族谱旁的一条旧目。方才听到问情谷的名字，那点记忆才渐渐接上。

众人闻声皆是一怔，纷纷转头望向她。

卢本韦抬眸，目光坚定，沉声开口：

\begin{quote}
``我家藏经阁里收着先祖留下的旧籍。我没读过问情谷的正文，只在家族旧目中见过它的名字。''

``猪神Opt老祖似乎留过相关手记。正文如今还在不在、有没有取丹之法，我现在不敢保证。可我们能回去问长辈，也许比在这里猜更有用。''
\end{quote}

原本一无所知的绝境前路，骤然浮出一条尘封千年的隐秘线索。

卢本韦继续正色补充，将古籍秘闻娓娓道来：

\begin{quote}
``那部分典藏平日不对外开放，许多卷册也已残缺。先回养猪场，请长辈核对目录、查一查实物，才知道能不能用。''

``就算找到，也得逐条辨认。老祖当年的经历，不一定就是我们今天能照着走的路。''
\end{quote}

这番话没有让问情谷变得容易，却给众人指明了一件眼下能做的事：先回养猪场，把那卷可能存在的手记找出来。

没有人再催着立即入谷。寻找材料与寻找可靠记载，是两件都不能省去的事；而他们终于有了一处可以先去核查的地方。

\begin{center}\rule{0.5\linewidth}{0.5pt}\end{center}

帝天听见Opt的名字，目光微动，却没有替卢氏家中的藏书作证。他只将话说得更明白了一些。

\begin{quote}
``若有亲历者的手记，值得去找。但残卷没有写到的地方，不能由你们自己补成确定的答案。''
\end{quote}

卢本韦点了点头，将自己记得的旧目与可能的存放处又认真想了一遍。

蔡须坤黯淡的眸光瞬间重燃璀璨光芒，紧绷的心弦终于稍稍松弛，目光灼灼地看向卢本韦：

\begin{quote}
``你记得的是旧目，正文还得回去找，对吗？''
\end{quote}

\begin{quote}
``对。''卢本韦道，``我记得有过这条记录，其他的不乱保证。回去找到多少，我们就按多少来商量。''
\end{quote}

\begin{center}\rule{0.5\linewidth}{0.5pt}\end{center}

前路依旧凶险未知，但此刻的果糕小队，已然不再是盲目闯荡。

他们还没有拿到古籍，更没有找到那两样材料。可返回养猪场、向卢氏长辈求证的方向已经清楚。蔡须坤向帝天道谢，又看了看身边的伙伴，终于没有再独自急着往前赶。

先稳住伤势，护好魂息，再去寻找那卷可能仍在的手记。

\chapter{第四十七章：序境初成，古经藏阁}\label{ux7b2cux56dbux5341ux4e03ux7ae0ux5e8fux5883ux521dux6210ux53e4ux7ecfux85cfux9601}

\begin{center}\rule{0.5\linewidth}{0.5pt}\end{center}

星斗核心秘境的苍茫龙威缓缓退散，厚重凝滞的天地气韵渐渐归于平和。

果糕小队众人伫立原地，心头的沉重并未因为卢本韦的一条线索而尽散。不过，他们终于有了一件眼下可以去做的事：返回养猪场，核查家族旧目中是否仍有Opt老祖关于问情谷的手记。古籍尚未拿到，其中能留下多少有用文字，谁也不能先作保证。

\textbf{蔡须坤}深深吐纳一口气，压下心底翻涌的复杂心绪，抬眸望向身旁一众伙伴，声音沉稳笃定，敲定前路方向。

\begin{quote}
``帝天前辈已经把知道的线索告诉我们了，我们先回养猪场，看看旧目对应的手记还在不在。''

``先养好伤，再查问情谷和两样材料的去处。有了可靠的记录，我们一起商量下一步。''
\end{quote}

号令落下，全队众人应声颔首，无人有半分迟疑。

\begin{center}\rule{0.5\linewidth}{0.5pt}\end{center}

牢大战死后，继续同行的是八名队员：唐三、霍雨浩、小舞、奥斯卡、诺当、蔡须坤、蔡兴坤与卢本韦。梅天狸和ddfg仍在队伍两侧护送，没有算入这八人的编制之中。

可所有人心底都清楚，这支队伍早已不复从前。

昔日坐镇阵前、护佑全员的\textbf{牢大}，已经无法再以血肉之躯挡在他们面前。蔡须坤护着的那点残魂仍在，却还不能清醒回应。众人下意识留出的那一处位置，时刻提醒着他们失去过什么，也提醒着他们，为何还要继续寻找归来的办法。

无人提及，却无人忘却。

\begin{center}\rule{0.5\linewidth}{0.5pt}\end{center}

众人向帝天告辞后，收拾心绪，循着获准通行的路径，缓缓走出星斗核心区域。

返程不再承受核心近处的龙威，众人却也没有因此松懈。两位长老交替探路，霍雨浩只在必要时短暂送出精神探测，避开沿途巢穴。行进与休息间留出的空隙，才被他们用来调息。

狼王的魂息仍在蔡须坤体内休养。契约稳住之后，它保留了大部分力量滋养自身，并没有向整支队伍无条件输送二十万年的修为。蔡须坤也依着霍雨浩的提醒，一点点梳理新增的共生联系，将刚触及的 \allowbreak{}\NovelInline{I02287}{\psi(I_\omega)}\allowbreak{} 层级慢慢稳住。

其余人各有需要补上的地方。\textbf{唐三}回顾突围时防线崩开的先后，尝试缩短蓝银草收放的间隙；\textbf{霍雨浩}分段调息，先养护透支的精神力，再复盘自己没能及时辨清的攻击。\textbf{小舞}每次落脚都比平日缓些，仍在适应伤后的筋骨，遇到地面湿滑处，便与诺当相互扶一把。

\textbf{奥斯卡}按众人的恢复情况分配食物，自己也没有因为赶路而省下休息。\textbf{卢本韦}一边梳理实战中的失误，一边回想家族旧目，怕漏掉与问情谷有关的一个字；\textbf{蔡兴坤}则将众人在交手时看见的线索记下，休整时拿出来，与霍雨浩逐条对照。

一开始，几个人都想把修行赶得更快。梅天狸看出他们的急切，让他们先停下强行运转的序力。

\begin{quote}
``复盘可以慢慢来。伤势没有稳住，不能靠再压一层力量，去盖住原来的问题。''
\end{quote}

她与ddfg交替走在前后，休息时也会留意众人的调息。遇到有人因用力过急牵动旧伤，两位长老便先帮忙稳住，再等气息平复后讲明问题出在哪一段。

几日下来，有人稳住了原有修为，有人在旧瓶颈上有了新的体会，进度并不相同。卢本韦记不清的一处交手细节，由霍雨浩补上；奥斯卡发现恢复食物消耗过快，也提出来请众人一起调整。真正令他们稍感踏实的，是伤势逐步稳定，队伍终于能在夜间轮换休息，不必一味透支着赶路。

\begin{center}\rule{0.5\linewidth}{0.5pt}\end{center}

数日后，众人离开星斗林海，又走过返程的道路，终于在清晨望见了养猪场熟悉的轮廓。

卢本韦猛地停下脚步，瞳孔骤然放大，眼底瞬间涌上一层滚烫的亮光。

天际尽头，晨雾尚未散尽，那片浩瀚无垠的养猪场如同蛰伏在云海之中的远古巨兽，磅礴壮阔的轮廓在金色朝阳的映照下缓缓浮现。连绵起伏的巨型畜栏从地平线一端铺展至另一端，看不到边际，乳白色的灵雾在栏顶缓缓流淌，折射出七彩瑞光，宛如仙境。

高耸的饲料塔错落矗立，塔身铭刻着卢氏家族的族徽。远处，溪水汇入场中的灵泉池，水雾蒸腾间，成群灵猪在护工照看下缓缓踱步。卢本韦望见熟悉的栏舍与巡逻路线，脚下终于不再像林间赶路时那样紧绷。

``到了\ldots\ldots 到了！'' 卢本韦的声音骤然拔高，带着压抑不住的颤抖与狂喜，她猛地转过身，双手紧紧攥住身旁蔡须坤的手臂，指尖因为太过用力而微微发白，眼眶瞬间泛红，泪水在眼底疯狂打转。

``你们看！那就是我家！那就是卢氏养猪场！''

她激动得语无伦次，一只手松开蔡须坤，拼命朝天际那片恢弘壮阔的建筑群指去，手臂挥得又高又急，仿佛生怕旁人看不见似的。

``看到最中间那座塔了吗？塔后面就是祖祠。''她朝远处指去，``上面的金猪徽记，也是族中先辈留下的。''

她的声音越来越快，越来越亮，像是要把积攒了数日的紧张、疲惫、恐惧，在这一刻尽数宣泄出来。 ``还有那边那片碧蓝色的灵泉池！池水能加速灵猪的血脉觉醒，小时候我天天趴在池边看小猪仔泡澡，毛茸茸的挤成一团，可可爱了！''

说到这里，她忽然吸了吸鼻子，声音带上了一丝哽咽，笑容却愈发灿烂： ``我好想它们\ldots\ldots 好想家\ldots\ldots{}''

她用力抹了一把眼角的泪花，深吸一口气，挺直腰背，朝着养猪场的方向大步迈出，脚步轻快得像要飞起来，裙摆被风扬起，猎猎作响。

``快走吧，我们先回祖宅，找我爹问问。''

她一边说一边回头，脸上挂着又哭又笑的灿烂表情，眼底是劫后余生、归家心切的滚烫光芒。

众人望着她兴奋雀跃的背影，连日积压在心底的沉重与阴霾，竟也不由自主地消散了几分。

这里有守卫与防护阵，也有熟悉的长辈和可以休整的地方。它给了众人暂时落脚的安稳，却不能让大陆上的纷争凭空消失。

踏入养猪场地界，熟悉的灵气与有人照看的道路，让众人一直绷着的肩背渐渐放松。伤势与疲惫仍在，至少这一刻，他们可以安心坐下休息。

卢本韦带着众人径直走入家族腹地，一路轻车熟路，很快便寻到了家族掌权人------她的父亲，\textbf{卢关}。

众人与卢关早已相识，曾在赛场并肩抗敌，也在这里受过他的照看。今日远远望见女儿带着伙伴归来，他立即放下手中的茶盏，起身迎了过去。

卢关先看向女儿，继而数过同行的队员。数到最后，他脚步微微一顿，目光又落回队列前方------那里没有他记得的牢大。

他没有立刻追问，只让人送来清水与干净的坐垫。孩子们带着怎样的消息，他还不知道；但衣衫上的旧血与尚未平稳的呼吸，已经说明这趟远行并不顺利。

卢本韦快步走到父亲面前。方才归家的欣喜还没有散尽，一见他欲言又止的神情，喉头却先哽住了。

卢关在众人面前停下，目光从女儿移到其余几人身上，声音仍是他们熟悉的温和。

\begin{quote}
``你们回来了。''
\end{quote}

熟悉的声音让众人心里一松。蔡须坤应了一声，抬起头，已经看见卢关向队伍后方望去的目光。

卢关留意到众人尚未恢复的伤势，又向两位长老点头致意，侧身让开庭院入口。

\begin{quote}
``一路凶险奔波，辛苦了，快些落座调息，休整片刻再说不迟。''
\end{quote}

待众人落座，卢关才低声问起一路的经过。他没有省略那个最难开口的问题，却也留出时间，等孩子们慢慢回答。

\begin{center}\rule{0.5\linewidth}{0.5pt}\end{center}

卢本韦看着父亲这般亲和模样，心中暖意涌动，没有半分隐瞒，将此番星斗之行的所有经历娓娓道来。

她先讲帅气小队的围杀，以及牢大燃尽本源为众人争取生路的经过。说到残魂被蔡须坤护住时，她停了一会儿，才继续讲梅天狸为赶去救援与狼王交手、随后折返救治的事。

狼族暂时停战、狼王自愿尝试魂灵契约，都是具体发生在那一处领地里的选择。卢本韦没有将其说成整片森林与人类从此和解，也没有替长老掩去重伤狼王的责任。

最后，她将帝天能够确认的情况与只听过的旧闻分开说明：乾坤造化丹的线索指向问情谷，天地灵水仍无可靠出处。家族旧目上可能有老祖手记的记录，所以他们想回来查证。

全程听完所有始末，卢关眉宇间的平和缓缓褪去，取而代之的是浓重的凝重与后怕。即便早已对世间邪祟有所认知，可听闻众人一次次深陷死局、浴血搏杀，他依旧心头震颤，愈发心疼这群年少坚韧的孩子。

他沉默良久，望着眼前这群尚且年少、却历经无数生死绝境的少年少女，语气满是唏嘘与郑重。

\begin{quote}
``你们这一行，太过冒险了。''
\end{quote}

卢关说完，将视线停在蔡须坤护着心口的手上，许久没有再开口。帅气教的手段他并不陌生，却没有想到孩子们这次会付出这样的代价。

\begin{quote}
``帅气教的手段与势力，都不能轻看。''卢关沉声道，``你们见到的只是其中一支队伍，背后还有多少人、多少据点，都要继续查。''

``只靠零散队伍，遇到这样的围杀，很难及时互相救援。要遏止帅气教，还得让帝国、宗门与扽魂殿互通消息，查明他们的据点和动向。''
\end{quote}

卢关目光坚定，道出破局关键：

\begin{quote}
``你们把这次亲眼见到的事写下来，我会与两位长老商量如何送交。查证与应对是我们这些长辈也该承担的事，不能全压在你们身上。''
\end{quote}

一旁静立聆听的\textbf{梅天狸}，闻言深深颔首，眼底满是认同与愤然。

\begin{quote}
``此教恶行滔天，祸乱苍生，屠戮生灵、践踏正道，属实罪无可赦，天地共诛。''
\end{quote}

梅天狸看向眼前这群带着伤的孩子，想起他们在林中问起牢大时的声音，神色慢慢沉了下来。

\begin{quote}
``这次你们能互相照应着回来，已经很不容易。''梅天狸看向众人，声音低了些，``我和ddfg没有及时赶到，也有我们要反省的地方。你们年纪还轻，不必用再闯一次险地，去证明自己比谁更能扛。''
\end{quote}

感慨过后，卢关收敛心绪，不再纠结于过往凶险与未来大势，目光落向卢本韦，又缓缓扫过在场众人，眼底的亲切与赞许愈发浓重，语气郑重起来。

\begin{quote}
``藏经阁有些关系传承的旧卷，平日只许指定族人查阅。你们先前看过的是修行与魂灵研究的部分，问情谷手记所在的旧藏，还没有核过实物。''

``这次我带你们进去找。能找到多少，先看见了再说。''
\end{quote}

卢关已经知道众人为何而来，也知道女儿记住的目录确有出处。他愿意让他们进入旧藏查阅，却仍要先确认卷册如今的情况，免得再给出一个落不了地的承诺。

话音落下，卢关转身抬手，做出一个请行的手势，语气温和却笃定：

\begin{quote}
``随我来。''
\end{quote}

\begin{center}\rule{0.5\linewidth}{0.5pt}\end{center}

一行人紧随其后，穿过层层古朴庭院、清幽廊道，踏入了养猪场核心区域。

此地与世隔绝，气韵古朴厚重，周遭古木参天、石碑林立，斑驳的岁月痕迹遍布四方，处处皆是万年传承的沧桑底蕴。空气中浮动着淡淡的书香与道韵，无数古老的规则气息萦绕不散，沉淀着卢氏家族跨越万古的修行秘辛。

一座恢弘古朴的双层阁楼静静伫立在禁地中心，飞檐翘角刻满上古符文，木梁石壁浸润着万年岁月的气韵，庄严肃穆，不容亵渎。

这里，便是卢氏家族的根基所在------\textbf{万年藏经阁}。

\begin{quote}
``祖上留下的许多卷册，都收在这里。历经搬迁与战乱，也有不少书已经散失，旧目录不能全当作现存书目。''
\end{quote}

卢关驻足阁楼门前，望着眼前肃穆古楼，又侧首看向身旁一众少年，语气敬畏而郑重，更添几分恳切。

\begin{quote}
``有关问情谷的手记，我记得旧目里确实有过。只是那部分旧藏久未整理，能否找到、还有多少完整文字，现在都要重新查。''
\end{quote}

阁楼木门厚重古朴，缓缓推开的瞬间，一股浩瀚苍茫的万古道韵扑面而来，无数尘封万年的文字光影流转浮动，仿佛跨越万古岁月，静待来人揭晓终极秘辛。

门已经打开。众人先将手上的水迹擦净，跟在卢关身后，轻轻踏过门槛。

\begin{center}\rule{0.5\linewidth}{0.5pt}\end{center}

众人紧随卢关踏入藏经阁楼内，满目皆是层层叠叠的古朴书卷。黝黑沉木打造的书架顶天立地，木纹里浸透千年墨香与道韵，无数古籍错落陈列，竹纸、兽皮、绢帛各类材质交织，岁月沉淀的厚重气息扑面而来，刚一踏入，便让人下意识敛息凝神，心生肃穆敬畏。

可细细端详，满目沧桑残破更让人心头沉沉。

万年光阴流转，沧海数次桑田更迭，大地震荡、气运衰败、战火浩劫、岁月风化，无数上古孤本、老祖手记惨遭湮灭。许多书架上空空荡荡，只余下常年尘封的浅淡印痕，证明曾经有典籍在此陈列；残存的书卷大多难逃损毁厄运，竹纸书页脆如枯叶，轻轻触碰便有细碎纸屑簌簌脱落，不少页面虫蛀斑驳、破洞连片，字迹被岁月磨得模糊残缺，通篇晦涩难辨。

有的古籍整册褪色发黑，墨色彻底消融，只余下空白泛黄的纸页，曾经记载的通天大道、秘境秘辛，尽数化作虚无；有的书卷仅剩寥寥残片，边角卷曲朽烂，拼接不出完整内容；更有不少上古秘典早已彻底失传，只在家族典籍目录中留下一个名号，徒留后人嗟叹。

偌大一座万年藏经阁，看似底蕴浩瀚，实则十不存一，万古传承的秘辛，大半都已葬于岁月长河之中。

卢关抬手轻轻拂去一卷古籍封面的薄尘，目光掠过满目残卷，眼底满是无奈与惋惜，沉声开口道：

\begin{quote}
``卢氏先祖世代守藏，殚精竭虑护持古籍，奈何天道轮转、岁月无情，人力终究难逆万古沧桑。历经数次天地浩劫与世代更迭，家族古籍目录错乱、典藏移位，无数上古至秘尽数失传，就连Opt老祖留存的乾坤问情谷秘典，我也多年未曾寻得踪迹。''
\end{quote}

此前他只记得族谱残页留有记载，知晓老祖曾留下一卷秘境手记，却因藏经阁万册典籍杂乱尘封、不少藏书移位混叠，始终无法确定具体藏处。今日为帮众人探寻救世机缘，卢关不再依托固有记忆，决意逐层逐架仔细翻找。

他抬步穿梭在林立的书架之间，指尖逐一抚过每一卷残破古籍，目光快速扫过泛黄封面与残缺页脚。阁楼顶层、暗格夹层、密藏书匣，但凡存放古老残卷的角落，他皆一一细细排查。许多书架积满千年厚尘，轻轻一动便尘雾弥漫，朽碎的纸屑簌簌飘落，枯燥又费力的搜寻过程，让空气都多了几分沉闷。

众人安静伫立一旁，无人打扰，静静看着卢关躬身翻找。偌大藏经阁书海浩渺，残缺典籍数不胜数，想要在这万千残卷中，寻出一卷不知所踪、年代万古的神级手记，无异于大海捞针。

\begin{center}\rule{0.5\linewidth}{0.5pt}\end{center}

几番逐层筛查、翻遍大半藏书之后，卢关终于在阁楼最深处、一处被尘封已久的隐秘暗格之中，停下了动作。此处远离中心祭台，隐蔽无比，常年无人触碰，暗格柜门布满斑驳锈迹，被数层老旧书卷堆叠掩盖，也正因如此，才侥幸躲过无数岁月浩劫与人为翻动，未曾遗失。

他挪开堆叠的残卷，推开暗格柜门，内里静置着一方古朴琉璃玉匣。匣侧还留着一道褪色的编号，与旧目中一项并不显眼的记载相合。

\begin{quote}
``找到了。''
\end{quote}

卢关长舒一口气，将玉匣移到案上。旧库几度搬迁，这件藏品又被归在``杂录''一栏，早已偏离原来的架位。他以前只知道手记的名目，直到这次沿着旧编号核查，才终于将实物对上。

匣身仍有淡淡莹光，是先辈留下、后来又经族人补护的封存纹路。它减缓了纸皮的腐坏，也挡不住入匣前已经形成的断裂和缺损。

卢关依家传方法解开匣上的封存纹路。内里并排放着两份旧卷：一份较薄，外签隐约可辨``问情''二字；另一份卷得更厚，题着``星斗纪事''，暂时还没有展开。

卢关先取出较薄的问情谷手记，将卷边残存的笔迹与匣内签记相互对照，才小心展开。兽皮经年封存，仍有韧性，几处断口却早已失去对应的残片；防护能留住剩下的部分，无法让已经遗失的文字重新出现。

大片大片的字迹沿着纹路彻底湮灭，化作空白斑驳的皮面，再也无从寻觅。唯有零星碎片化的区域，还残留着老祖落笔的雄浑墨痕，墨色深沉如暗星，裹挟着一丝不朽的神性微光，与普通古籍的凡俗笔墨截然不同。

\begin{center}\rule{0.5\linewidth}{0.5pt}\end{center}

\textbf{卢本韦}屏气凝神，双手小心翼翼接过这卷绝世残篇，触感微凉厚重，指尖触碰之处，隐约有微弱的神泽流转，涤荡心神。她将兽皮古卷缓缓平铺展开，众人立刻围拢上前，目光紧紧锁定卷面，不敢错过分毫细节。

这份问情谷手记的绝大部分正文已经缺失。尚可连读的，主要是开篇一段行路记载、中段几句机缘提示，以及末尾的一行警示。匣里另有星斗旧录，却不能拿另一卷中的内容补成这份手记失去的部分。

她凝眸细辨，逐字轻声念出残存的模糊字迹：

\begin{quote}
``乾坤问情谷者，西北入落日森林八十余里，树木渐疏，遂雾起，进其谷。''
\end{quote}

这一句字迹尚可辨认，下一行却沿着断口失去了后文。卢本韦将它抄下，正要再念，霍雨浩轻声问：

\begin{quote}
``西北，是从哪里往西北？这里的八十余里，要从哪个位置起算？''
\end{quote}

卢关仔细看过断口，摇了摇头。

\begin{quote}
``起点可能就在缺掉的文字里。落日森林是个方向，这段行程还要和现在的舆图核对，不能照着残句就认定找到了谷口。''
\end{quote}

无人知晓前路雾障是寻常天地异象，还是困杀万古修士的绝杀迷阵；无人知晓谷口有无镇守禁制、有无先天凶煞；更无人知晓踏入迷雾的刹那，究竟会直面机缘，还是坠入死局。

跨越数行空白卷面，才有断断续续的残字勉强可辨，拼凑出残缺的秘辛：

\begin{quote}
``内有神明之力，亦正亦邪\ldots\ldots 情念为锁，真心为钥\ldots\ldots 可获乾坤造化丹，塑神基，补天道\ldots\ldots{}''
\end{quote}

短短数句残缺不全，字句之间断层无数。神明之力究竟源自何物？是上古神祇遗留的道统，还是秘境自生的阴阳煞力？亦正亦邪的力量如何判定、如何制衡？情念试炼的规则是什么、关卡有几重、失败的代价为何？造化丹除了逆天改命、修补道基，是否还有更为恐怖的反噬弊端？所有关键隐秘，尽数被万古岁月彻底抹去。

\begin{center}\rule{0.5\linewidth}{0.5pt}\end{center}

关于机缘的那一段，到这里便断了。再往下，只有几处无法连读的墨痕，没有现成的试炼次序与取丹方法。

唯独古卷最底端，整片残破卷面的最后一角，一行字迹异常深重、笔画铿锵决绝，毫无风化残缺，完整留存万古岁月，赫然烙印在兽皮之上：

\begin{quote}
\textbf{``非我等修为者，勿入！''}
\end{quote}

这行字比其余墨痕更深，侥幸留在没有断裂的一角。卢关辨认过后，示意众人先停下猜测。

\begin{quote}
``这是老祖留给后来者的告诫。至少说明在他看来，修为不足的人不宜贸然进入。至于他有没有受伤、谷中是什么力量，仅凭这一句，还看不出来。''
\end{quote}

警示下方还有一道断口，说明卷尾曾有别的内容。如今残片不在，无法判断写的是亲历、避险之法还是别的嘱托。若日后能找到别卷或旧抄本，才有机会再作核对。

\begin{center}\rule{0.5\linewidth}{0.5pt}\end{center}

阁楼之内寂静无声，唯有众人细微的呼吸声轻轻回荡。

卢本韦将手停在古卷边缘，没有再来回摩挲脆弱的裂痕。族中向来敬重Opt老祖，她也习惯了长辈们讲述先祖的功绩。如今亲眼看见他留下的告诫，才发现那些故事之外，还有许多他们不知道的经历。

老祖到底记录过多少，只能从残存的装订与断口中判断。眼下看不到破关经验，可能是那些文字遗失了，也可能本来就没有详写，不能直接断言他当年只敢留下警告。

蔡须坤看了看那行警示，又望向中段``可获乾坤造化丹''的残字。两处文字都必须留下：取丹的可能与不宜轻入的告诫，不能只挑自己愿意相信的一处。

老祖当年入谷，究竟看见了什么、遭遇了什么、错失了什么、忌惮了什么？

万古谜团，尽数尘封，无解无答。

残缺的古卷，寥寥的线索，看似为众人指明了救世之路，实则掀开了一层深不见底的无尽黑暗。

前路可能有造化，却还隔着无法确认的入口、试炼与代价。众人得到了一点可核查的方向，离真正救人仍很远。

未知的凶险，如同漫天阴霾，沉沉笼罩在果糕小队所有人的心头。

\chapter{第四十八章：古史秘辛，谷藏天威}\label{ux7b2cux56dbux5341ux516bux7ae0ux53e4ux53f2ux79d8ux8f9bux8c37ux85cfux5929ux5a01}

藏经古阁的死寂沉沉延续良久，满室残卷无声伫立，上古岁月遗留的斑驳裂痕，仿佛无声叩问着人心。

众人的目光仍停在问情谷手记上。``非我等修为者，勿入''几个字，让方才急于动身的念头慢慢冷却。老祖的告诫必须认真对待，但它并没有替他们排出秘境与天下强者的高低。

残札缺页太多，谷中试炼的内容与力量来源仍未说明。众人将能读通的句子反复看过，也没有找到可以安心入谷的办法。

\begin{center}\rule{0.5\linewidth}{0.5pt}\end{center}

卢关等众人抄完残文，才将问情谷手记重新垫好，取出匣内另一份卷得较厚的旧录。

\begin{quote}
``这是《星斗纪事》，与方才那份问情谷手记分别成卷。你们提到帝天，我想里面的几段旧事，也许有助于理解他守着什么。''
\end{quote}

这份旧录有几处首尾相接的段落尚在，记录比问情谷残札连贯，却没有谷内试炼的内容。卢关先将这一点说清，才展开卷面。

卢关辨过笔迹与旧签，按能读通的段落缓缓诵出，断开的地方便停住，没有自行补成先祖原话。

\textbf{《星斗纪事》残文}

\begin{quote}
昔入星斗，见龙王帝天，修为八十八万载，统御兽众，护一方林海。

时人兽相争，猎杀与报复不休，伤者日众。兽群亦有生息之困，强弱相夺，战火不止。

吾与帝天交手，破其攻势，止其出林兵戈。后与之约：不以滥杀掠人为修行之途，护群兽，慎启争端。

瑞兽承天地祥瑞，本源可泽万兽。吾留护源之法，助稳其栖息之地，托帝天护持。

然旧恨未消，护源不能尽弭人兽之争。宜各守其约，徐图共存。
\end{quote}

能连读的这几段，在卷面上各有对应的墨迹。再往后的几处缺页，卢关便略过，只将新辨出的卷次记在旁边，以备日后继续核对。

\begin{center}\rule{0.5\linewidth}{0.5pt}\end{center}

卢本韦看看旧录，又望了望蔡须坤。

\begin{quote}
``万年前是八十八万年，帝天如今是八十九万年。这里记的是当时的修为？''
\end{quote}

卢关点头。

\begin{quote}
``是。不能把我们现在知道的年限，挪到老祖当年的交手上。旧录说老祖破了他的攻势、止住战事，没有说把他的道基毁尽。''
\end{quote}

唐三将两卷记录分开，低声问：

\begin{quote}
``老祖能胜过当年的帝天，又对问情谷留了这样的警示，是不是说明谷里的力量更强？''
\end{quote}

\begin{quote}
``能说明那地方不能轻看，还排不出这样的高低。''卢关将问情谷残札推近一点，``留下告诫的人，未必是在那里败了。谷中也可能有与交手不同的规则，不能只拿修为去撞。我们连试炼内容都没找到，先别替老祖把结论说满。''
\end{quote}

霍雨浩望着``情念为锁，真心为钥''几个残字，慢慢接上了他的意思。

\begin{quote}
``同一句提示，可能对不同的人有不同要求。就算战力强过守谷者，也未必等于懂得怎么通过。''
\end{quote}

\begin{quote}
``至少不能把两件事混成一件。''梅天狸道，``先把亲眼能确认的记录留下。我们需要的是办法，不是一张凭猜测排出来的强弱表。''
\end{quote}

卢关这才回到《星斗纪事》，将原文与自己的理解分开讲述。

\begin{quote}
``先祖当年确实压住过帝天的一次出林行动，后来才谈成护持与止戈的约定。他能胜这一战，不等于旧录里那头龙从此没有自己的意志，更不等于所有族群都因此归服。''
\end{quote}

\begin{quote}
``人类为魂环与资源猎杀，魂兽因争地与报复伤人，各有纠葛。旧录里说的是止住当时的战事，从未说全体魂兽必须吃人才能修炼。猎物、领地、灵气与族群旧怨，也不能全用一种办法解决。''
\end{quote}

蔡须坤听到这里，目光重新落在``瑞兽''二字上。

\begin{quote}
``护源之法，是说老祖造出了瑞兽本源？''
\end{quote}

\begin{quote}
``不是这个意思。''卢关指着原文，``这里先说瑞兽承天地祥瑞，随后才说留法护持。先祖帮助稳定的是它栖息、受养的环境。瑞兽的本源并非从一件法器中造出来，也不是可以随意分给别人的修为。''
\end{quote}

旧录没有写那时瑞兽的名字与具体经历。蔡须坤没有将这几行直接接成自己的身世，只将涉及护源的原句认真记下。它能解释帝天为什么重视瑞兽，却还不能回答关于她的一切。

\begin{quote}
``那他如今还留在核心\ldots\ldots{}''诺当迟疑了一下。
\end{quote}

\begin{quote}
``有守护的责任，不等于永远不能离开。''梅天狸道，``他亲口说过，必要时会出去。旧录里的约定是慎启争端，不能把它读成任何情况下都不准行动。''
\end{quote}

卢关合上旧录中已经读过的部分。

\begin{quote}
``留护源之法，能让附近族群多一份生息的条件，却不会把从前的仇恨一并抹去。你们与狼群谈成停战、救下狼王，意义就在这些具体事情上，不必将它说成已经替整片大陆完成了和解。''
\end{quote}

\begin{center}\rule{0.5\linewidth}{0.5pt}\end{center}

藏经阁内安静了一会儿。唐三将方才推到一半的强弱比较划去，换成几项尚待查证的问题：入谷起点、试炼规则、丹药是否仍在，以及失败之后能否退出。

小舞看着那张纸，轻声道：

\begin{quote}
``这么说，我们知道的危险，和我们自己想出来的危险，也得分开。''
\end{quote}

\begin{quote}
``对。''霍雨浩应道，``不知道，不代表可以放心；也不能因为害怕，就把所有猜测都当成事实。''
\end{quote}

卢本韦望着第一次完整展开在自己面前的两份旧卷，心里仍有震动。她熟悉的是家族流传的故事，今天才得到许可细看这些旧藏。故事里的``先祖无所不能''，到了纸上，原来也有没有写明、不能替他作答的地方。

``族中常讲先祖的功绩，许多事也能在旧录里找到痕迹。可我们后人知道得终究不全。''

卢关轻叹一声，语气带着无尽唏嘘与怅然。

``后来的族谱记着他成神、离开此界的事，之后就再没有可供我们直接询问的消息。如今能依靠的，是这些文字与先辈保存下来的实物。''

``至于问情谷的通行方法，现存手记中没有完整答案。缺掉的那些部分究竟记过什么，我们也不能替他断言。''

``若想再往前找，还得看看别处有没有抄本、访录或亲历者的线索。我们族里的这两卷，今天能给出的就是这些。''

\begin{center}\rule{0.5\linewidth}{0.5pt}\end{center}

偌大藏经阁，墨香依旧，道韵浮沉，此刻却只剩无尽的茫然与压抑。

线索比出发时多了一点：落日森林的地名、残缺的行程、情念相关的提示，以及可能取得造化丹的记载。入口尚未确认，取丹方法仍空着，天地灵水的下落也没有因此得到解答。

众人将这几条分别抄下，没有将``真心为钥''改写成真心足够便一定能过关。若要踏入迷雾，仍需先查地点，再设法获得更多关于规则与退路的消息。

诺当紧握双拳，眼底锋芒收敛，只剩凝重：

\begin{quote}
``我们现在带着伤，入口都没认准。这样就往谷里闯，连该防什么都不知道。''
\end{quote}

小舞轻轻颔首，眉眼间满是审慎：

\begin{quote}
``丹药是值得找的线索，可牢大还需要养魂和重塑道基。不能只盯着一个药名，就把别的条件都忘了。''
\end{quote}

众人心绪沉沉。返程中一点点恢复的力量，到了这些尚未解开的空白前，并没有变成足以放心入谷的把握。

蔡须坤原本已经按住纸角，想再问一次能否即刻启程，听见小舞的话，才慢慢松开了手。她想救牢大，却不能将``尽快''只理解为尽快闯入。

卢关看着一众少年凝重的模样，并未出言宽慰。他深知，唯有认清凶险，方能敬畏大道，方能在绝境之中寻得一线生机。

\begin{quote}
``这两卷现存的内容，今天能辨认的就这些。先抄清、核对，往后若有新消息，还可以拿回来比照。''
\end{quote}

他缓缓开口，声音沉稳落地。

蔡须坤看着自己抄下的两样材料，缓缓开口：

\begin{quote}
``我们先不闯谷。入口、取丹之法、天地灵水，都还要继续找。牢大的魂息由我护着，其他的事，我们和长辈们一起想办法。''
\end{quote}

这句话说得很慢。决定暂缓，并没有让她轻松，可她终于将最急切的念头摆到同伴面前，而没有独自替所有人决定。

两份旧卷并排放在案上，一份留下了旧战与护源的经过，一份只给出残缺的问情谷线索。它们各自能证明什么，众人终于分清了。

\textbf{先养伤，查清线索，再决定如何入谷。}

\begin{center}\rule{0.5\linewidth}{0.5pt}\end{center}

沉闷的气氛笼罩藏经古阁许久，压得一众少年心绪沉甸甸的，无人率先开口打破死寂。卢关望着眼前这群少年少女坚毅又青涩的面庞，眼底的凝重缓缓褪去，取而代之的是一片温和的期许，众人紧绷的心绪也稍稍缓和。

他目光落向身侧的卢本韦，语气放缓，带着几分欣慰与笃定，朗声对众人说道：

``你们这般年纪，已经能接触到很深的序数结构，往后还有不少路可走。''卢关看着众人，``陈星长老当年也有逐步打牢根基的时候。别拿一时进度，去抵消眼下该做的准备。''

一旁伫立的梅天狸闻言，微微含笑点头，眼底满是唏嘘与赞叹，接过话茬感慨道：

``有些概念能推到，不代表身体和魂息已经跟得上。''她道，``把一道难题解明白，再将力量稳稳用出来，中间也要练。别因为想追上谁，就略过这一步。''

众人闻言心头微松，些许底气悄然回升，先前被上古禁忌压制的颓势稍稍消解。

\begin{center}\rule{0.5\linewidth}{0.5pt}\end{center}

卢关抬手轻轻拍了拍卢本韦的肩头，目光真挚而郑重，字字落地有声：

``闺女，我想先把下一任执掌人的责任交给你学着担。''他说，``你还要与伙伴修行，场里的事务仍由我主持。等你回来，把该学的事务也学会，我们再正式交接。''

``这次我会将猪神令交给你暂时携行。它是先祖的遗物，也是我们这些后人要守的责任。先学会认识它、照看它，别急着将里面的传承都引进自己体内。''

话音落下，卢关转身迈步走出藏经古阁。众人心中微动，纷纷紧随其后，踏着古朴青石古道，一路行至养猪场核心的专属祠堂。

众人走过高塔后的祖祠主殿，在卢关带领下转入一间古朴的传承偏殿。这里梁柱不饰繁纹，只有历代守藏者添置的灯盏与供案。中央供着猪神Opt老祖的神位，香火安静地燃着。

卢关率先净手焚香，向神位躬身行礼。众人随之致意，感念先祖留下的庇护与研究，也记着今日所读的那些仍未解开的空白。

\begin{center}\rule{0.5\linewidth}{0.5pt}\end{center}

祭拜礼毕，卢关领众人到偏殿后的传承密室，先将连着地脉的纹路稳住，才从供座上取出猪神令。他托在掌心，示意卢本韦先看清令上的印记。

那是一枚拇指大小的小猪令牌，由远古精铁锻成，通体呈深邃哑光的铁色。边缘有历代持令者抚过的细微痕迹，纹路却仍清晰。小猪体态浑圆，四肢蜷伏，腹间刻着极小的``Opt''字样。它看着朴素，握近时却能感觉到一缕缓慢而温厚的力量。

卢关先让女儿辨认外层的守护纹路，再指给她看更深处与地脉相连的部分。令中留着Opt老祖的一点本源力量，也有历代持令者补入的印记，足以作为辨认传承与引导守护之用，却不是将历代先辈的全部修为都装在其中。

\begin{quote}
``完整承接传承，要在密室里与地脉相合，一年至多一次，还要看你能承受多少。今天先认印，不作灌顶。离开养猪场，你也不能隔着千里就把整片地脉的力量随意搬走。''
\end{quote}

``此乃\textbf{猪神令}，是先祖留下的信物，也见证历代执掌人的交接。''卢关缓缓道，``持令不是只拿着一件物事。族中认定的传承、对场务的承担、与守护阵法的联系，都要学。''

``你与它的气息相合，若能用好外层的守护纹路，可以帮你定住一时紊乱的心神，或在遇到阴煞侵扰时多撑片刻。每次引动都要消耗力量，不能任意挥霍。''

``至于问情谷，现有记载没有说此令能在那里通行。能不能护住你们，更不能拿来先作担保。往后你还要随伙伴们修行，带着它慢慢熟悉。遇到不懂的，记下来，回来再与我一起看。''

\begin{center}\rule{0.5\linewidth}{0.5pt}\end{center}

卢本韦依着父亲的指点，将手指轻轻搭上令牌。等那缕灵力与印记接稳，卢关才将它放入她掌心。

古朴的小猪神令轻轻震颤，腹间的Opt印记浮起一线金光。那道光顺着铁纹慢慢流动，映亮了卢本韦的手指，也让她清楚感受到自己送入的灵力正在往哪里走。

密室门外，守护纹路随之泛起一层微光，很快便停在廊下，没有继续向远处扩散。回应来自令牌与近旁阵法原有的联系，没有凭空垂落无尽灵力。

等那缕回应稳住，卢本韦试着让灵力多走一段。她急着多送一点，指尖顿时发麻，光纹也跟着乱了。卢关立刻按住她的手腕，让她将力量收细。

众人望着令牌与阵纹的回应，没有上前打断。

等纹路重新稳住，卢关才松开手，眼里有了笑意。

\begin{quote}
``记住刚才的感觉。它能回应你，说明可以继续学，不是已经把里面的东西都学会了。''
\end{quote}

卢本韦低头记下那道最外层纹路的走向，将令牌仔细收好。她接下的是一份开始学习承担的责任；场里的日常事务，仍由父亲与熟悉场务的族人照管。

\begin{center}\rule{0.5\linewidth}{0.5pt}\end{center}

回到偏殿后，令牌上的光已经收敛。梅天狸看过卢本韦仍有些发麻的手指，适时开口，语气认真。

``令牌能回应，是件好事。你方才多送了一点力量就险些失控，往后要把这一层慢慢练熟。''

她看向其余队员，继续道：

\begin{quote}
``你们在那场围杀里露出的短处，不是拿到一件信物就能补上的。精神力耗尽以后，谁来辨认下一道攻击；前面的人撑不住，后面怎样接替------这些都要重新练。问情谷的规则还没查清，越是这样，越得先有稳住自己、照看同伴的本事。''
\end{quote}

众人闻言纷纷收敛心中的振奋，神色再度沉凝，默默点头认可。经历方才上古秘辛的震撼，他们早已认清自身实力的渺小。

梅天狸等众人安静下来，才提出接下来的安排。

``先随我回扽魂殿，安置伤势，护好牢大的残魂。我会请Asheep长老一道，带你们修习\textbf{反射序数大法}。ddfg负责沿途照应，也会将这次遭遇的情况送回殿中。''

``这里说的反射，是从序数结构中辨认可以下探、对应的规律，不是把敌人的招式原样反弹回去。你们先前已经学过不少序数构造，接下来要用另一套方法去校验它们，再学着把推演与自身的力量接稳。''

``这会从基础讲起，不表示你们原来的境界要退回去，也不表示听完一门功法就一定能闯过问情谷。先让自己有更多判断与自保的办法，才好谈后面的事。''

众人终于有了眼下能落实的安排。修行、养魂与查证线索，都需要时间，却可以一道往前推，不必停在``立刻闯谷''与``彻底放弃''之间反复拉扯。

卢关点头应下。

\begin{quote}
``我会将这两卷可辨的部分另抄一份，原件留在阁中。你们带着抄本去，也请殿里的长老看看。落日森林旧路与天地灵水的消息，我继续查，找到可靠的再传过去。''
\end{quote}

他看向蔡须坤。

\begin{quote}
``牢大的事，我们记着。你先把自己和那缕魂息护稳，不必为了今天还没有答案，就逼自己明天一定闯进去。''
\end{quote}

蔡须坤应了一声，将抄好的材料名单收进怀中。众人又回到藏经阁，将临行前要带的抄本仔细整理。卢本韦拿起干净的纸，坐到案边，与霍雨浩逐字核对残文；唐三和蔡兴坤帮忙标明断口，遇到不认识的字，便留白请卢关确认。

他们还没有找到让牢大归来的办法，也没有因此停下。去扽魂殿之前，先将手里这一点真实的线索带稳。

\chapter{第四十九章 扽魂开讲，反射初基}\label{ux7b2cux56dbux5341ux4e5dux7ae0-ux627dux9b42ux5f00ux8bb2ux53cdux5c04ux521dux57fa}

养猪场的传承交接尘埃落定，第二日天光大亮，果糕小队八人便在梅天狸的引领下辞别卢关，径直奔赴大陆正道腹地的扽魂殿。

扽魂殿屹立于苍莽群山之巅，不倚宗门、不附帝国，是世间少数专研序数大道的中立修行圣地。殿宇依山层叠直入云霄，石壁镌刻着无数古老序数符文，终年被清濛序雾笼罩，寻常修士难窥门径。asheep 长老便是此殿执掌者之一，与梅天狸相交数十年，二人皆是当世序数大道的顶尖学者，一手''反射序数大法''独步天下。

踏入殿门，周遭喧嚣尽数隔绝。正殿空旷肃穆，正前方一座白玉高台，两侧分列八座素色蒲团，天地序力在此自发成流、循节律运转，天然便是参悟序道的绝佳道场。众人依序落座屏息凝神，心知此番修行关乎救世机缘，半分松懈不得。

梅天狸缓步走上高台，素手轻扬，一道莹白光幕在半空展开，其上浮动着细密的序数符号，正是反射模式的基础符文阵列。一旁的 asheep 长老负手而立，身形挺拔如松，目光沉凝，静待理论讲罢再传授实战凝练之法。

\begin{center}\rule{0.5\linewidth}{0.5pt}\end{center}

``修行至序境，你们便该知晓：所谓序数大道，本质是对天地规则集合的推演与折叠。寻常修士只知堆砌灵力、硬破壁垒，终其一生也难窥 I 境门径；而反射序数大法，便是教你们以规则驭规则，通过'反射'操作层层提纯道则集合，逐级攀登序数层级，最终以自身道基映照诸天序力。''

``今日我从最根基的反射模式讲起。这套体系并非单一记号，而是一整套完备的集合运算系统，所有推演皆有严谨定则，容不得半分臆测。切记：\textbf{本体系所有讨论，皆限定于可数序数范畴之内}------便是你们当下修行所能触及的全部天地序力。''

梅天狸指尖轻点光幕，第一组定义清晰浮现。

``反射模式的一切推演，皆始于 onto 操作。它以一个序数集合为作用对象，输出一个全新的序数集合。我们将'第 n 层谓词反射作用于集合 X'记作 Π\_n onto X，为书写简便，简写为 n-X。''

``若被作用的集合是全体可数序数集 Ord，可将其省略，直接记作 n。''

光幕上定格出严格定义：

\begin{quote}
\textbf{定义 1（onto 简写规则）}：对任意自然数 n 与序数集合 X， \NovelDisplay{M00478R01}{%
n\text{-}X = \Pi_n \text{ onto } X%
} 若 \allowbreak{}\NovelInline{I02288}{X = \text{Ord}}\allowbreak{}（全体可数序数类），则简记为 \allowbreak{}\NovelInline{I02289}{n}\allowbreak{}，即 \allowbreak{}\NovelInline{I02290}{n = \Pi_n \text{ onto } \text{Ord}}\allowbreak{}。
\end{quote}

``举两个最基础的实例：

\begin{itemize}
\tightlist
\item
  \allowbreak{}\NovelInline{I02291}{\Pi_1 \text{ onto }}\allowbreak{} 全体序数，得到全体极限序数的集合，对应记号 \allowbreak{}\NovelInline{I02292}{1}\allowbreak{}；
\item
  \allowbreak{}\NovelInline{I02293}{\Pi_2 \text{ onto }}\allowbreak{} 全体序数，得到全体容许序数的集合，对应记号 \allowbreak{}\NovelInline{I02294}{2}\allowbreak{}。
\end{itemize}

前者是你们早已踏足的领域，后者便是你们如今正试图叩开的 I 境根基。至于 \allowbreak{}\NovelInline{I02295}{\Pi_3}\allowbreak{} 及以上，待夯实前两层根基后再谈不迟。''

``第二，交集运算。集合的交集 \allowbreak{}\NovelInline{I02296}{\cap}\allowbreak{}，即同时属于两个集合的元素构成的新集合。在反射模式中，为简化书写，我们将 \allowbreak{}\NovelInline{I02297}{\cap}\allowbreak{} 简写为空格。''

``这里有两条铁则，务必刻进识海：

\begin{itemize}
\tightlist
\item
  \textbf{其一}，交集与 onto 运算优先级完全等同；
\item
  \textbf{其二}，所有运算一律从右向左结合。
\end{itemize}

差一个次序，最终集合便天差地别；日后凝练道则、叠加反射，但凡错了一分，轻则修为滞涩，重则道基崩损。''

光幕同步浮现规则与示例：

\begin{quote}
\textbf{运算定则}：\allowbreak{}\NovelInline{I02298}{\cap}\allowbreak{} 与 onto 优先级相同，表达式从右向左依次运算。

\textbf{例 1}：\allowbreak{}\NovelInline{I02299}{2\ 1\text{-}2 = 2 \cap (1\text{-}2) = \Pi_2\text{反射集} \cap (\Pi_1 \text{ onto } \Pi_2\text{反射集})}\allowbreak{}

\textbf{例 2}：\allowbreak{}\NovelInline{I02300}{2\text{-}3\ 1\text{-}3\ 2\text{-}3 = 2\text{-}(3 \cap (1\text{-}(3 \cap (2\text{-}3))))}\allowbreak{} 即 \allowbreak{}\NovelInline{I02301}{\Pi_2 \text{ onto } [\Pi_3 \cap (\Pi_1 \text{ onto } (\Pi_3 \cap (\Pi_2 \text{ onto } \Pi_3)))]}\allowbreak{}
\end{quote}

台下蔡须坤指尖在膝头默默推演示例，眸光沉凝。唐三蹙眉片刻，侧首向霍雨浩低声确认：``最右侧的 2-3 是第一层内核，往左依次裹上 1-、3、2-，对吗？''霍雨浩精神力微动，在识海中拆解完运算层级，轻轻颔首。

``第三，元素选取操作。反射模式研究的所有集合，元素皆是序数，天然可按从小到大排序。''

集合 \allowbreak{}\NovelInline{I02302}{X}\allowbreak{} 的最小元素，记作 \allowbreak{}\NovelInline{I02303}{\min X}\allowbreak{}；第二小、第三小分别记作 \allowbreak{}\NovelInline{I02304}{2\text{nd } X}\allowbreak{}、\allowbreak{}\NovelInline{I02305}{3\text{rd } X}\allowbreak{}；第 n 小则为 \allowbreak{}\NovelInline{I02306}{\text{nth } X}\allowbreak{}。若无歧义，可将 \allowbreak{}\NovelInline{I02307}{\min}\allowbreak{} 省略，直接用集合名指代其最小元；若易与序数本身混淆，则在外层加括号，以 \allowbreak{}\NovelInline{I02308}{(X)}\allowbreak{} 代表 \allowbreak{}\NovelInline{I02309}{\min X}\allowbreak{}。

\begin{quote}
\textbf{定义 2（最小元约定）}：对任意非空序数集合 \allowbreak{}\NovelInline{I02310}{X}\allowbreak{}，\allowbreak{}\NovelInline{I02311}{\min X}\allowbreak{} 为 \allowbreak{}\NovelInline{I02312}{X}\allowbreak{} 中最小序数；无歧义时 \allowbreak{}\NovelInline{I02313}{X}\allowbreak{} 可代指 \allowbreak{}\NovelInline{I02314}{\min X}\allowbreak{}，有歧义时记作 \allowbreak{}\NovelInline{I02315}{(X)}\allowbreak{}。
\end{quote}

``再讲 aft 操作。\allowbreak{}\NovelInline{I02316}{A \text{ aft } B}\allowbreak{}，意为集合 \allowbreak{}\NovelInline{I02317}{A}\allowbreak{} 中所有大于序数 \allowbreak{}\NovelInline{I02318}{B}\allowbreak{} 的元素。与 \allowbreak{}\NovelInline{I02319}{\min}\allowbreak{} 结合，\allowbreak{}\NovelInline{I02320}{\min(A \text{ aft } B)}\allowbreak{} 便是'集合 \allowbreak{}\NovelInline{I02321}{A}\allowbreak{} 中大于 \allowbreak{}\NovelInline{I02322}{B}\allowbreak{} 的最小元素'，日常推演中可直接简写为 \allowbreak{}\NovelInline{I02323}{A \text{ aft } B}\allowbreak{}。''

\begin{quote}
\textbf{定义 3（aft 操作）}：对任意序数集合 \allowbreak{}\NovelInline{I02324}{A}\allowbreak{} 与序数 \allowbreak{}\NovelInline{I02325}{\beta}\allowbreak{}， \NovelDisplay{M00479R01}{%
A \text{ aft } \beta = \{\alpha \in A \mid \alpha > \beta\}%
} \allowbreak{}\NovelInline{I02326}{\min(A \text{ aft } \beta)}\allowbreak{} 可简记为 \allowbreak{}\NovelInline{I02327}{A \text{ aft } \beta}\allowbreak{}。
\end{quote}

一直沉默的 asheep 长老此时开口，声线低沉厚重：``这套操作看似琐碎，却是日后定位修为节点、测算突破契机的根基。实战之中，对手的神通威能对应某一层序数，你若能精准算出其层级之上的下一个反射节点，便能以弱胜强、后发先至------这便是反射大法'以序御力'的核心。''

讲完基础符号，梅天狸语气稍缓，光幕上的符号开始层层递进：``现在我们从最底层的 \allowbreak{}\NovelInline{I02328}{\Pi_1}\allowbreak{} 反射开始，完整推演一遍层级递进，让你们直观感受反射的力量。''

``\allowbreak{}\NovelInline{I02329}{\Pi_1 \text{ onto }}\allowbreak{} 一个集合，核心效果是取出该集合的所有极限点。所谓极限点，便是集合中前 \allowbreak{}\NovelInline{I02330}{\lambda}\allowbreak{} 个元素的上确界，其中 \allowbreak{}\NovelInline{I02331}{\lambda}\allowbreak{} 取遍所有小于该点的极限序数。通俗说，极限点就是集合里无法通过后继运算抵达、只能靠无穷累加逼近的节点。''

\begin{quote}
\textbf{定理 1（\allowbreak{}\NovelInline{I02332}{\Pi_1}\allowbreak{} 反射的语义）}：对任意序数集合 \allowbreak{}\NovelInline{I02333}{X}\allowbreak{}， \NovelDisplay{M00480R01}{%
\Pi_1 \text{ onto } X = \{\alpha \mid \alpha \text{ 是 } X \text{ 的极限点}\}%
} 即 \allowbreak{}\NovelInline{I02334}{X}\allowbreak{} 中所有''前极限序数个元素的上确界''构成的集合。
\end{quote}

``我们先推演最基础的 \allowbreak{}\NovelInline{I02335}{1 = \Pi_1 \text{ onto } \text{Ord}}\allowbreak{}。全体序数的前 \allowbreak{}\NovelInline{I02336}{\omega}\allowbreak{} 个元素是 \allowbreak{}\NovelInline{I02337}{0, 1, 2, \ldots}\allowbreak{}，其上确界为 \allowbreak{}\NovelInline{I02338}{\omega}\allowbreak{}，故 \allowbreak{}\NovelInline{I02339}{\omega \in 1}\allowbreak{}；前 \allowbreak{}\NovelInline{I02340}{\omega \cdot 2}\allowbreak{} 个元素的上确界为 \allowbreak{}\NovelInline{I02341}{\omega \cdot 2}\allowbreak{}，同样属于 \allowbreak{}\NovelInline{I02342}{1}\allowbreak{}。以此类推，所有极限序数皆在其中。因此：

\NovelDisplay{M00481R01}{%
1 = \{\omega, \omega \cdot 2, \omega \cdot 3, \ldots, \omega^2, \ldots, \omega^\omega, \ldots, \omega^{\omega^\omega}, \ldots\}%
}

也就是全体极限序数的集合。''

话音落下，众人只觉识海中轰然一响。过往无数次懵懂的境界突破，此刻都有了清晰的数学对应：小舞肉身凝练的一次次极限蜕变，奥斯卡灵力提纯的一层层瓶颈，本质都是在踏过自身道则集合里的一个个极限点。

``在此基础上再做一次 \allowbreak{}\NovelInline{I02343}{\Pi_1}\allowbreak{} 反射，便是 \allowbreak{}\NovelInline{I02344}{1\text{-}1 = \Pi_1 \text{ onto } 1}\allowbreak{}。取 \allowbreak{}\NovelInline{I02345}{1}\allowbreak{} 的前 \allowbreak{}\NovelInline{I02346}{\omega}\allowbreak{} 个元素 \allowbreak{}\NovelInline{I02347}{\omega, \omega \cdot 2, \omega \cdot 3, \ldots}\allowbreak{}，其上确界为 \allowbreak{}\NovelInline{I02348}{\omega^2}\allowbreak{}，故 \allowbreak{}\NovelInline{I02349}{\omega^2}\allowbreak{} 是 \allowbreak{}\NovelInline{I02350}{1\text{-}1}\allowbreak{} 的最小元。同理，\allowbreak{}\NovelInline{I02351}{1\text{-}1}\allowbreak{} 便是全体 \allowbreak{}\NovelInline{I02352}{\omega^2}\allowbreak{} 的倍数构成的集合：

\NovelDisplay{M00482R01}{%
1\text{-}1 = \{\omega^2, \omega^2 \cdot 2, \ldots, \omega^3, \ldots, \omega^\omega, \ldots, \omega^{\omega^\omega}, \ldots\}%
}

继续迭代，\allowbreak{}\NovelInline{I02353}{1\text{-}1\text{-}1}\allowbreak{} 对应 \allowbreak{}\NovelInline{I02354}{\omega^3}\allowbreak{} 的倍数，\allowbreak{}\NovelInline{I02355}{1\text{-}1\text{-}1\text{-}1}\allowbreak{} 对应 \allowbreak{}\NovelInline{I02356}{\omega^4}\allowbreak{} 的倍数。有限次 1- 操作，规律一目了然。''

``有限次迭代好懂，那超限次数呢？比如 \allowbreak{}\NovelInline{I02357}{(1\text{-})^\omega}\allowbreak{}，重复 \allowbreak{}\NovelInline{I02358}{\omega}\allowbreak{} 次 1- 操作，该如何定义？''

梅天狸语气郑重：``我们定下通则：当上标为极限序数 \allowbreak{}\NovelInline{I02359}{\lambda}\allowbreak{} 时，\allowbreak{}\NovelInline{I02360}{(1\text{-})^\lambda}\allowbreak{} 等于所有 \allowbreak{}\NovelInline{I02361}{\beta < \lambda}\allowbreak{} 时 \allowbreak{}\NovelInline{I02362}{(1\text{-})^\beta}\allowbreak{} 的交集。''

\begin{quote}
\textbf{定义 4（超限 1- 迭代）}：对任意极限序数 \allowbreak{}\NovelInline{I02363}{\lambda}\allowbreak{}， \NovelDisplay{M00483R01}{%
(1\text{-})^\lambda = \bigcap_{\beta < \lambda} (1\text{-})^\beta%
}
\end{quote}

``据此推演，\allowbreak{}\NovelInline{I02364}{(1\text{-})^\omega = \bigcap_{n < \omega} (1\text{-})^n}\allowbreak{}。\allowbreak{}\NovelInline{I02365}{(1\text{-})^1}\allowbreak{} 是 \allowbreak{}\NovelInline{I02366}{\omega^1}\allowbreak{} 倍数，\allowbreak{}\NovelInline{I02367}{(1\text{-})^2}\allowbreak{} 是 \allowbreak{}\NovelInline{I02368}{\omega^2}\allowbreak{} 倍数\ldots\ldots 所有有限次幂的交集，便是 \allowbreak{}\NovelInline{I02369}{\omega^\omega}\allowbreak{} 的倍数集合。因此：

\NovelDisplay{M00484R01}{%
(1\text{-})^\omega = \{\omega^\omega, \omega^\omega \cdot 2, \ldots, \omega^{\omega+1}, \ldots, \omega^{\omega \cdot 2}, \ldots\}%
}

更一般地，对任意序数 \allowbreak{}\NovelInline{I02370}{\alpha}\allowbreak{}，有恒等式：

\NovelDisplay{M00485R01}{%
\min (1\text{-})^\alpha = \omega^\alpha%
}

也就是说，每一层 1- 反射的最小元，恰好对应 \allowbreak{}\NovelInline{I02371}{\omega}\allowbreak{} 的 \allowbreak{}\NovelInline{I02372}{\alpha}\allowbreak{} 次幂。''

讲到此处，卢本韦眸光一亮。家族古籍中零散记载的''\allowbreak{}\NovelInline{I02373}{\omega}\allowbreak{} 幂阶道韵''，此刻终于与严谨的反射定义严丝合缝地对应上了，心中诸多模糊疑惑豁然开朗。

``只靠序数上标，终究还在 \allowbreak{}\NovelInline{I02374}{\omega}\allowbreak{} 幂的框架里打转。想要再进一步，就要引入不动点进位模式。''

梅天狸指尖划过光幕，一行新的定义浮现：

\begin{quote}
\textbf{定义 5（一阶不动点集）}： \NovelDisplay{M00486R01}{%
(1\text{-})^{1,0} = \{\alpha \mid \alpha = \min (1\text{-})^\alpha\}%
}
\end{quote}

``结合方才的恒等式 \allowbreak{}\NovelInline{I02375}{\min (1\text{-})^\alpha = \omega^\alpha}\allowbreak{}，这个定义等价于 \allowbreak{}\NovelInline{I02376}{\alpha = \omega^\alpha}\allowbreak{}。满足这个等式的序数，便是你们熟知的 \allowbreak{}\NovelInline{I02377}{\varepsilon}\allowbreak{} 序数。所以------

\NovelDisplay{M00487R01}{%
(1\text{-})^{1,0} = \{\varepsilon_0, \varepsilon_1, \ldots, \varepsilon_\omega, \ldots, \zeta_0, \ldots\}%
}

全体 \allowbreak{}\NovelInline{I02378}{\varepsilon}\allowbreak{} 序数，尽在这一个集合之中。''

``在此之上再做一次 1- 反射，得到 \allowbreak{}\NovelInline{I02379}{(1\text{-})^{1,1} = 1\text{-}(1\text{-})^{1,0}}\allowbreak{}，便是下标为极限序数的 \allowbreak{}\NovelInline{I02380}{\varepsilon}\allowbreak{} 序数构成的集合，如 \allowbreak{}\NovelInline{I02381}{\varepsilon_\omega}\allowbreak{}、\allowbreak{}\NovelInline{I02382}{\varepsilon_{\omega^2}}\allowbreak{} 等等。我们沿用弱 \allowbreak{}\NovelInline{I02383}{\varphi}\allowbreak{} 模式的进位规则：\allowbreak{}\NovelInline{I02384}{1\text{-}(1\text{-})^{\#,\alpha} = (1\text{-})^{\#,\alpha+1}}\allowbreak{}；遇上标为极限序数时，依旧取交集。由此可延伸出 \allowbreak{}\NovelInline{I02385}{(1\text{-})^{1,\omega}}\allowbreak{}、\allowbreak{}\NovelInline{I02386}{(1\text{-})^{1,\varepsilon_0}}\allowbreak{}，直至新的不动点层级：

\NovelDisplay{M00488R01}{%
(1\text{-})^{2,0} = \{\alpha \mid \alpha = \min (1\text{-})^{1,\alpha}\}%
}

这便是全体 \allowbreak{}\NovelInline{I02387}{\zeta}\allowbreak{} 序数的集合。再往上，\allowbreak{}\NovelInline{I02388}{(1\text{-})^{1,0,0}}\allowbreak{} 对应 \allowbreak{}\NovelInline{I02389}{\Gamma}\allowbreak{} 序数，乃至 @ 记号的超限嵌套，理论上足以覆盖所有递归序数。''

梅天狸话音稍顿，望着台下神色各异的众人，缓声道：``你们不必急于吃透所有高阶嵌套。今日所讲的 \allowbreak{}\NovelInline{I02390}{\Pi_1}\allowbreak{} 反射体系，是反射模式的地基。唯有把极限点、超限迭代、不动点这三层逻辑刻进道基，日后才能承接 \allowbreak{}\NovelInline{I02391}{\Pi_2}\allowbreak{} 容许序数、乃至递归不可达的层级。''

理论讲罢，asheep 长老上前一步，袖袍一挥，八道淡金色序力丝线便分别没入众人体内。众人只觉丹田道基微微发烫，自身的灵力流转竟自发地按照 1- 反射的节律开始提纯。

``梅长老讲的是理，我讲的是用。''asheep 长老声线平稳，``反射序数大法的修行，便是将这套集合运算内化于自身道基：

\begin{itemize}
\tightlist
\item
  每运转一次 1- 反射，你们便提纯一次自身道则的极限点，灵力凝练度跃迁一个 \allowbreak{}\NovelInline{I02392}{\omega}\allowbreak{} 层级；
\item
  凝出 \allowbreak{}\NovelInline{I02393}{(1\text{-})^\omega}\allowbreak{}，便是将道基打磨至 \allowbreak{}\NovelInline{I02394}{\omega^\omega}\allowbreak{} 层级的精纯，灵力生生不息、循环不绝；
\item
  待到证得 \allowbreak{}\NovelInline{I02395}{(1\text{-})^{1,0}}\allowbreak{} 不动点，你们的道基便会达成自洽------自身便是 \allowbreak{}\NovelInline{I02396}{\varepsilon}\allowbreak{} 层级的序力源泉，无需外求便可自生序力。''
\end{itemize}

``你们当下卡在 I 境门槛，道基尚且停留在 \allowbreak{}\NovelInline{I02397}{\Gamma}\allowbreak{} 序数上下，距 \allowbreak{}\NovelInline{I02398}{\varepsilon}\allowbreak{} 圆满尚有距离。接下来半月，你们便在此殿闭关，先将 \allowbreak{}\NovelInline{I02399}{\Pi_1}\allowbreak{} 反射体系修至 \allowbreak{}\NovelInline{I02400}{(1\text{-})^{1,0}}\allowbreak{} 圆满，凝出自身第一重不动点道基。''

话音落下，殿内一时寂静，众人皆在闭目消化所学，将理论与自身修为相互印证。

蔡须坤闭目凝神，瑞兽本源在体内缓缓流转，将方才的定义、定理一一映照在识海。他能清晰感知到，自身道则集合原本杂乱的极限点，此刻正循着反射规则重新排布，壁垒松动的速度比往日快了数倍。少顷他睁开眼，沉声问道：``梅长老，依此推演，\allowbreak{}\NovelInline{I02401}{\Pi_2}\allowbreak{} 反射便是比 \allowbreak{}\NovelInline{I02402}{\Pi_1}\allowbreak{} 更高一层的提纯，所得的容许序数无法用更低层的递归运算抵达，对吗？''

梅天狸微微颔首，眼中掠过赞许：``不错。\allowbreak{}\NovelInline{I02403}{\Pi_1}\allowbreak{} 取极限点，尚可用递归的基本列逼近；可 \allowbreak{}\NovelInline{I02404}{\Pi_2}\allowbreak{} 所得的容许序数，根本不存在长度短于自身的递归基本列------这便是 \allowbreak{}\NovelInline{I02405}{\omega}\allowbreak{}、I 这类序数的本质：非递归，不可由下而上抵达，只能由上而下折叠。这也是你们接下来要攻克的核心。''

另一侧，霍雨浩已在精神识海中构建出层层嵌套的集合模型，从 \allowbreak{}\NovelInline{I02406}{1}\allowbreak{} 到 \allowbreak{}\NovelInline{I02407}{1\text{-}1}\allowbreak{}，再到 \allowbreak{}\NovelInline{I02408}{(1\text{-})^\omega}\allowbreak{}、\allowbreak{}\NovelInline{I02409}{(1\text{-})^{1,0}}\allowbreak{}，层级分明。他隐隐感觉到，若能将精神力按反射层级折叠，其探测与洞察能力必将再上一个台阶。唐三指尖萦绕着蓝银草的嫩绿色光芒，草木生长的节律天然契合极限点的递进，他已然开始尝试用 1- 反射梳理蓝银本源，只觉原本驳杂的草木道韵飞速凝练，隐隐有了 \allowbreak{}\NovelInline{I02410}{\varepsilon}\allowbreak{} 层级的雏形。

日头西斜，第一日的讲学暂告一段落。白玉高台上的光幕渐渐散去，殿内的序力雾气却愈发浓郁。八人端坐蒲团之上，周身皆萦绕着深浅不一的序力光晕，道基在反射规则的打磨下悄然蜕变。

梅天狸与 asheep 长老并肩立于殿门处，望着殿内潜心修行的少年们，神色欣慰。

``这群孩子悟性惊人，尤其是蔡须坤与卢本韦，一个道基浑厚、一个家学渊源，怕是不出十日便能凝出 \allowbreak{}\NovelInline{I02411}{(1\text{-})^{1,0}}\allowbreak{} 不动点。''asheep 长老缓缓开口。

梅天狸轻轻摇头，目光望向远方天际，语气带着几分凝重：``这还远远不够。乾坤问情谷的凶险远超想象，仅凭 \allowbreak{}\NovelInline{I02412}{\Pi_1}\allowbreak{} 圆满根本不够看。待他们夯实 \allowbreak{}\NovelInline{I02413}{\Pi_1}\allowbreak{} 根基，便要立刻切入 \allowbreak{}\NovelInline{I02414}{\Pi_2}\allowbreak{} 容许序数，讲解递归不可达与容许点进位\ldots\ldots{}''

\begin{center}\rule{0.5\linewidth}{0.5pt}\end{center}

日头西斜，殿内序雾缓缓流淌，梅天狸讲完 \allowbreak{}\NovelInline{I02415}{\Pi_1}\allowbreak{} 反射的不动点进位体系，话音稍作停顿，留给众人消化推演的空隙。

台下八人闭目凝神，各自将方才整套 \allowbreak{}\NovelInline{I02416}{\Pi_1}\allowbreak{} 反射的定义、迭代、不动点在识海中一遍遍复盘。霍雨浩的精神力最是精微，率先在层层嵌套的集合模型里，捕捉到了一股异常熟悉的结构感。他眉心微蹙，精神识海中两套体系并排铺展开------一边是今日所学的 \allowbreak{}\NovelInline{I02417}{(1\text{-})}\allowbreak{} 上标进位，一边是早年修习 \allowbreak{}\NovelInline{I02418}{\varphi}\allowbreak{} 函数时烂熟于心的不动点层级，两相映照，竟严丝合缝。

``不对\ldots\ldots{}''他低声自语，声音虽轻，却在寂静的正殿里格外清晰，``这 \allowbreak{}\NovelInline{I02419}{\Pi_1}\allowbreak{} 反射的不动点模式，怎么和我们早年学过的 \allowbreak{}\NovelInline{I02420}{\varphi}\allowbreak{} 函数，几乎是一一对应的？''

此言一出，众人纷纷睁眼，眸中皆闪过恍然之色。唐三指尖顿住，蓝银草的微光在指节间敛去，沉声接过话头：``雨浩不说我还未深想，如今对照来看，确实如此。\allowbreak{}\NovelInline{I02421}{\min(1\text{-})^\alpha}\allowbreak{} 对应 \allowbreak{}\NovelInline{I02422}{\omega^\alpha}\allowbreak{}，恰是 \allowbreak{}\NovelInline{I02423}{\varphi}\allowbreak{} 函数的第零层根基；\allowbreak{}\NovelInline{I02424}{(1\text{-})^{(1,0)}}\allowbreak{} 对应全体 \allowbreak{}\NovelInline{I02425}{\varepsilon}\allowbreak{} 序数，正是 \allowbreak{}\NovelInline{I02426}{\varphi(1,0)}\allowbreak{}；\allowbreak{}\NovelInline{I02427}{(1\text{-})^{(2,0)}}\allowbreak{} 对应 \allowbreak{}\NovelInline{I02428}{\zeta}\allowbreak{} 序数，便是 \allowbreak{}\NovelInline{I02429}{\varphi(2,0)}\allowbreak{}；再往上 \allowbreak{}\NovelInline{I02430}{(1\text{-})^{(1,0,0)}}\allowbreak{} 对应 \allowbreak{}\NovelInline{I02431}{\Gamma}\allowbreak{} 序数，也就是 \allowbreak{}\NovelInline{I02432}{\varphi(1,0,0)}\allowbreak{}\ldots\ldots 连进位规则都一模一样。''

蔡须坤缓缓颔首，瑞兽本源映照出的道则纹路在识海中两两重合，他语气沉定，补全了其中的对应逻辑：``不止是最终的不动点集合，连进位规则都完全契合弱 \allowbreak{}\NovelInline{I02433}{\varphi}\allowbreak{} 模式。每做一次 1- 反射，相当于在最末位加一；遇上极限序数取交集，对应极限位置的基本列逼近；新的不动点层级向左进位，和 \allowbreak{}\NovelInline{I02434}{\varphi}\allowbreak{} 函数的多元参数进位毫无二致。说白了，这套 \allowbreak{}\NovelInline{I02435}{\Pi_1}\allowbreak{} 反射体系，本质就是用反射的语言，重新包装了一遍弱 \allowbreak{}\NovelInline{I02436}{\varphi}\allowbreak{} 函数的构造。''

``可既然结构完全对应，为何还要单独成一门反射大法？''小舞蹙眉不解，``若是如此，我们直接修 \allowbreak{}\NovelInline{I02437}{\varphi}\allowbreak{} 函数不也一样？''

卢本韦轻轻摇头，指尖摩挲着袖间猪神令的纹路，家族古籍里的记载与今日所学相互印证，她一语点破了其中的关键：``结构对应，不代表强度对等。你们仔细想，\allowbreak{}\NovelInline{I02438}{\Pi_1}\allowbreak{} 反射再怎么嵌套、怎么叠不动点，自始至终都跳不出递归序数的范畴------它能构造出的最高序数，也不过是递归序数的极限，和 \allowbreak{}\NovelInline{I02439}{\varphi}\allowbreak{} 函数的天花板别无二致。''

``可我们之前接触过的 OCF，靠 \allowbreak{}\NovelInline{I02440}{\Omega}\allowbreak{} 折叠能直接跨过递归序数的壁垒，触及非递归的容许层级。''霍雨浩瞬间反应过来，眸中光亮明灭，``也就是说，\allowbreak{}\NovelInline{I02441}{\Pi_1}\allowbreak{} 反射构造序数的能力，其实远不如带大序数折叠的 OCF？它看似体系庞大、层层迭代，实则始终困在递归的框架里，连 \allowbreak{}\NovelInline{I02442}{\Omega}\allowbreak{} 这一道门槛都跨不过去？''

``正是如此。''

高台上的梅天狸闻言，眼中掠过一丝赞许，缓缓开口印证了众人的推断。

``你们能自行勘破这一层，足见悟性。\allowbreak{}\NovelInline{I02443}{\Pi_1}\allowbreak{} 反射取极限点、叠不动点，本质上都是在可数递归序数的范围内做内部提纯，无论叠多少层 \allowbreak{}\NovelInline{I02444}{(1\text{-})}\allowbreak{}、多少位不动点，终究逃不出递归运算的边界。它的构造强度，确实与 \allowbreak{}\NovelInline{I02445}{\varphi}\allowbreak{} 函数相当，远不及以 \allowbreak{}\NovelInline{I02446}{\Omega}\allowbreak{} 为折叠子的 OCF。''

asheep 长老负手立于旁侧，低沉的声音补充道：``既如此，为何还要从 \allowbreak{}\NovelInline{I02447}{\Pi_1}\allowbreak{} 反射学起？因为反射是骨架，\allowbreak{}\NovelInline{I02448}{\varphi}\allowbreak{} 函数与 OCF 都只是具体的构造工具。学会了反射的底层逻辑，你们才能看懂 OCF 折叠的本质，才能理解 \allowbreak{}\NovelInline{I02449}{\Pi_2}\allowbreak{} 容许序数为何是更高一层的天地。''

``\allowbreak{}\NovelInline{I02450}{\Pi_1}\allowbreak{} 反射是序道的'内功心法'，教你们理解极限点、迭代、不动点的底层规律；\allowbreak{}\NovelInline{I02451}{\varphi}\allowbreak{} 函数是这套规律的一种写法，OCF 是更强的写法。没有 \allowbreak{}\NovelInline{I02452}{\Pi_1}\allowbreak{} 反射打底，直接去啃 OCF，只会知其然不知其所以然，遇上高阶逆乱秩序便会根基不稳。''

众人闻言尽皆默然。

此前他们只觉得反射序数大法玄奥高深，是远超寻常功法的绝世传承，心中未尝没有几分震撼与敬畏。可此刻对照 \allowbreak{}\NovelInline{I02453}{\varphi}\allowbreak{} 函数一一拆解开来才发现，自己惊为天人的 \allowbreak{}\NovelInline{I02454}{\Pi_1}\allowbreak{} 反射体系，居然只是和早已学过的 \allowbreak{}\NovelInline{I02455}{\varphi}\allowbreak{} 函数同阶的基础构造，连 OCF 的强度门槛都尚未摸到。

震撼褪去，取而代之的是更深的清醒。

蔡须坤轻叹一声，语气里再无半分此前突破境界的自得：``原以为摸到了反射大道的门径，如今看来，\allowbreak{}\NovelInline{I02456}{\Pi_1}\allowbreak{} 反射不过是把我们早已熟知的东西，换了一套更底层的语言重新讲了一遍。真正的高阶序道，还在 \allowbreak{}\NovelInline{I02457}{\Pi_2}\allowbreak{} 容许序数之后。''

``难怪长老说这只是地基。''唐三眸光深沉，望向高台上的光幕，``地基再牢，也还在平地之上。唯有跨过 \allowbreak{}\NovelInline{I02458}{\Pi_2}\allowbreak{} 这道坎，触及容许序数、递归不可达的层级，才算真正踏入高阶序道的大门，才能拥有抗衡乾坤问情谷禁忌之力的资本。''

诺当紧握双拳，眼底战意非但未减，反而愈发炽烈：``也好。越早认清根基浅薄，越能沉下心苦修。先把 \allowbreak{}\NovelInline{I02459}{\Pi_1}\allowbreak{} 反射彻底刻进道基，把底层逻辑打牢，再攻 \allowbreak{}\NovelInline{I02460}{\Pi_2}\allowbreak{} 容许序数便是。''

殿外暮色渐沉，殿内烛火次第亮起。八人周身的序力光晕愈发沉稳，再无半分初见反射体系时的浮躁与惊艳。他们终于明白，反射序数大法的真正价值，不在于 \allowbreak{}\NovelInline{I02461}{\Pi_1}\allowbreak{} 这一层能构造多强的序数，而在于它搭建了一套可无限向上延伸的规则框架------\allowbreak{}\NovelInline{I02462}{\Pi_1}\allowbreak{} 是第一层，\allowbreak{}\NovelInline{I02463}{\Pi_2}\allowbreak{} 是第二层，往后还有 \allowbreak{}\NovelInline{I02464}{\Pi_3}\allowbreak{}、\allowbreak{}\NovelInline{I02465}{\Pi_4}\allowbreak{}\ldots\ldots 无穷无尽的高阶反射，才是这条大道真正的浩瀚之处。

梅天狸望着台下少年们神色的变化，与 asheep 长老对视一眼，皆微微颔首。

``能看破'\allowbreak{}\NovelInline{I02466}{\Pi_1}\allowbreak{} 与 \allowbreak{}\NovelInline{I02467}{\varphi}\allowbreak{} 同阶、强度不及 OCF'这一层，你们才算真正入了反射的门。''她声音平缓，带着几分期许，``今日 \allowbreak{}\NovelInline{I02468}{\Pi_1}\allowbreak{} 的理论便讲到这里。半月之内，将其内化于道基，凝出 \allowbreak{}\NovelInline{I02469}{(1\text{-})^{1,0}}\allowbreak{} 不动点。之后，我们便正式开启 \allowbreak{}\NovelInline{I02470}{\Pi_2}\allowbreak{} 容许序数的修行------到那时，你们才会见识到，反射体系真正的力量。''

殿内烛火摇曳，序力流转不息。

反射序数的大门已然开启，而更浩瀚的 \allowbreak{}\NovelInline{I02471}{\Pi_2}\allowbreak{} 层级、更凶险的秘境前路，正静静等待着他们。

\chapter{第五十章：容许点进位}\label{ux7b2cux4e94ux5341ux7ae0ux5bb9ux8bb8ux70b9ux8fdbux4f4d}

第二天，扽魂殿中灵纹缓缓流转，虚空中那张反射式对照表静静悬浮，果糕小队众人望着层层递进的序数层级，心神久久未能平复。他们本以为吃透一阶反射、初识二阶集合，便已触碰到序道根基，此刻才惊觉，容许序数的背后，藏着远比想象中更深邃的秩序奥秘。

梅天狸见众人神色，便知他们尚未完全融会贯通，抬手轻拂，对照表隐去，声线平和地开启了更深一层的讲解：``方才我们只说了结论 ------ \allowbreak{}\NovelInline{I02472}{\Pi_2}\allowbreak{} onto 全体序数，得到全体容许序数的集合，简记为 \allowbreak{}\NovelInline{I02473}{2}\allowbreak{}。可 \allowbreak{}\NovelInline{I02474}{\Pi_2}\allowbreak{} onto 任意集合的作用，暂时是难以说明的，我们便从全体序数出发，把这套二阶反射的底层逻辑彻底讲透。''

``首先，我们重新明确容许序数的定义。'' asheep 长老接过话头，清冷的声音在殿内回荡，字字精准，``所谓容许序数，可以大致地理解为无法通过比它小的序数进行递归运算得到的序数。其精准的定义依靠集合论，本系列不做深入展开。与之定义类似的，还有一个`非递归序数'的概念。由于它们在 \allowbreak{}\NovelInline{I02475}{\lambda\alpha.\Omega_{\alpha+1}\text{-stb.}}\allowbreak{} 以后才会产生区别，并且本文不区分这两个概念。''

霍雨浩眉头微蹙，拱手发问：``长老，`递归运算得到'这个说法还是太过抽象。修行之中，我们该如何直观判断一个序数是不是容许序数？总不能全靠公理断言吧？''

``问得极好，这正是入门最关键的判据。'' 梅天狸微微颔首，指尖划出一道序数序列，``换用另一个更直观的概念：\textbf{不存在长度小于它的、递归表达的基本列的序数}。也就是说，如果我们找到了某个序数的一条长度小于它本身的基本列，就能立刻断言它不是容许序数。这为我们排除很多容许序数的备选提供了一条可行的方案。''

唐三眸光一动，立刻抓住了核心：``也就是说，能被低级基本列逼近的，都不是容许序数；唯有无法被低级序列穷尽的，才是高阶容许秩序？''

``正是如此。'' asheep 长老道，``最小的容许序数，是 Church-Kleene 序数，记为 \allowbreak{}\NovelInline{I02476}{\omega_1^{CK}}\allowbreak{}，在我们的体系中写作 \allowbreak{}\NovelInline{I02477}{\Omega}\allowbreak{} ------ 没错，就是在 OCF 中出现的那个 \allowbreak{}\NovelInline{I02478}{\Omega}\allowbreak{}。过去常有人将 \allowbreak{}\NovelInline{I02479}{\Omega}\allowbreak{} 定义为最小的不可数序数 \allowbreak{}\NovelInline{I02480}{\omega_1}\allowbreak{}，但这其实是过时的旧观念。要实现折叠递归序数的功能，\allowbreak{}\NovelInline{I02481}{\omega_1^{CK}}\allowbreak{} 已经足够了。''

小舞轻声追问：``那第二个容许序数呢？是 \allowbreak{}\NovelInline{I02482}{\Omega}\allowbreak{} 的两倍、平方，还是更高阶的不动点？''

``都不是。'' 梅天狸摇头，``第二个容许序数是 \allowbreak{}\NovelInline{I02483}{\omega_2^{CK}}\allowbreak{}，记作 \allowbreak{}\NovelInline{I02484}{\Omega_2}\allowbreak{}。它无法通过包括 \allowbreak{}\NovelInline{I02485}{\Omega}\allowbreak{} 在内的比它小的序数通过递归运算得到，这意味着 \allowbreak{}\NovelInline{I02486}{\Omega_2}\allowbreak{} 不是 \allowbreak{}\NovelInline{I02487}{\Omega \times 2}\allowbreak{}、\allowbreak{}\NovelInline{I02488}{\Omega^2}\allowbreak{}、\allowbreak{}\NovelInline{I02489}{\varepsilon_{\Omega+1}}\allowbreak{}、\allowbreak{}\NovelInline{I02490}{\varphi(\Omega,1)}\allowbreak{} 之类的东西，而是远在它们之上的存在。''

``再往后，还有 \allowbreak{}\NovelInline{I02491}{\Omega_3, \Omega_4, \dots, \Omega_{\omega+1}, \dots, \Omega_{\Omega+1}, \dots}\allowbreak{} 它们一同组成了容许序数的集合，也就是集合 \allowbreak{}\NovelInline{I02492}{2}\allowbreak{}。''

诺当敏锐地捕捉到了表述里的细节，立刻问道：``长老，您列举时跳过了 \allowbreak{}\NovelInline{I02493}{\Omega_\omega}\allowbreak{} 和 \allowbreak{}\NovelInline{I02494}{\Omega_\Omega}\allowbreak{}，单独列出了 \allowbreak{}\NovelInline{I02495}{\Omega_{\omega+1}}\allowbreak{} 和 \allowbreak{}\NovelInline{I02496}{\Omega_{\Omega+1}}\allowbreak{}。难道 \allowbreak{}\NovelInline{I02497}{\Omega_\omega}\allowbreak{}、\allowbreak{}\NovelInline{I02498}{\Omega_\Omega}\allowbreak{} 都不是容许序数？''

``观察很准。'' asheep 长老指尖轻点，虚空中浮现出三条基本列： * ``\allowbreak{}\NovelInline{I02499}{\Omega_\omega}\allowbreak{} 存在一条长度为 \allowbreak{}\NovelInline{I02500}{\omega}\allowbreak{} 的基本列 \allowbreak{}\NovelInline{I02501}{\{\Omega, \Omega_2, \Omega_3, \dots\}}\allowbreak{}，所以它是非容许的； * \allowbreak{}\NovelInline{I02502}{\Omega_\Omega}\allowbreak{} 存在一条长度为 \allowbreak{}\NovelInline{I02503}{\Omega}\allowbreak{} 的基本列 \allowbreak{}\NovelInline{I02504}{\{\Omega, \Omega_2, \dots, \Omega_\omega, \dots\}}\allowbreak{}，所以它是非容许的； * 更进一步地，\allowbreak{}\NovelInline{I02505}{\text{OFP} = \psi_I(0)}\allowbreak{}，也就是满足 \allowbreak{}\NovelInline{I02506}{\alpha = \Omega_\alpha}\allowbreak{} 的最小序数，也存在一条长度为 \allowbreak{}\NovelInline{I02507}{\omega}\allowbreak{} 的基本列 \allowbreak{}\NovelInline{I02508}{\{\Omega, \Omega_\Omega, \Omega_{\Omega_\Omega}, \dots\}}\allowbreak{}，它同样是非容许的。''

``事实上，对于大多数的 \allowbreak{}\NovelInline{I02509}{\Omega_\alpha}\allowbreak{}，其中 \allowbreak{}\NovelInline{I02510}{\alpha}\allowbreak{} 为极限序数时，\allowbreak{}\NovelInline{I02511}{\alpha}\allowbreak{} 的基本列总是会成为突破口，让我们找到 \allowbreak{}\NovelInline{I02512}{\Omega_\alpha}\allowbreak{} 的一条长度小于 \allowbreak{}\NovelInline{I02513}{\Omega_\alpha}\allowbreak{} 的基本列。''

卢本韦抚着掌心微微发烫的猪神令，沉吟道：``那是不是意味着，不存在极限序数 \allowbreak{}\NovelInline{I02514}{\alpha}\allowbreak{}，使得 \allowbreak{}\NovelInline{I02515}{\Omega_\alpha}\allowbreak{} 为容许序数？''

话音落下，殿内短暂一静。梅天狸与 asheep 长老对视一眼，都从对方眼中看到了赞许。

``这个问题，正是二阶反射走向深层的核心。'' 梅天狸缓缓道，``不过在回答它之前，我们先来看一个更基础的操作：对集合 \allowbreak{}\NovelInline{I02516}{2}\allowbreak{} 执行 \allowbreak{}\NovelInline{I02517}{1\text{ onto}}\allowbreak{}，会得到什么。''

``\allowbreak{}\NovelInline{I02518}{1-2}\allowbreak{}，即取出集合 \allowbreak{}\NovelInline{I02519}{2 = \{\Omega, \Omega_2, \dots, \Omega_{\omega+1}, \dots, \Omega_{\Omega+1}, \dots\}}\allowbreak{} 中所有的\textbf{极限点}。''

``\allowbreak{}\NovelInline{I02520}{2}\allowbreak{} 中的前 \allowbreak{}\NovelInline{I02521}{\omega}\allowbreak{} 个序数的上确界是 \allowbreak{}\NovelInline{I02522}{\sup \{\Omega, \Omega_2, \dots\} = \Omega_\omega}\allowbreak{}，所以 \allowbreak{}\NovelInline{I02523}{\Omega_\omega \in 1-2}\allowbreak{}； \allowbreak{}\NovelInline{I02524}{2}\allowbreak{} 中的前 \allowbreak{}\NovelInline{I02525}{\omega \times 2}\allowbreak{} 个序数的上确界是 \allowbreak{}\NovelInline{I02526}{\Omega_{\omega \times 2}}\allowbreak{}，所以 \allowbreak{}\NovelInline{I02527}{\Omega_{\omega \times 2} \in 1-2}\allowbreak{}；\ldots\ldots{} 如此一来，\allowbreak{}\NovelInline{I02528}{1-2}\allowbreak{} 就是`\allowbreak{}\NovelInline{I02529}{\Omega_\alpha}\allowbreak{}，其中 \allowbreak{}\NovelInline{I02530}{\alpha}\allowbreak{} 为极限序数'的集合。''

她顿了顿，加重了语气：``需要注意的是，根据我们刚才的讨论，这样的序数大多数都是非容许的。这意味着 \allowbreak{}\NovelInline{I02531}{1-2}\allowbreak{} 这个集合中的大部分元素甚至不属于 \allowbreak{}\NovelInline{I02532}{2}\allowbreak{}。我们对 \allowbreak{}\NovelInline{I02533}{2}\allowbreak{} 这个集合进行了 \allowbreak{}\NovelInline{I02534}{1\text{ onto}}\allowbreak{} 的操作，得到了一个性质更差的集合 ------ 这种事在反射里，是相当常见的。''

诺当愕然：``反射操作\ldots\ldots{} 竟然不是越强叠加强，反而可能让集合性质变差？''

``反射是结构重塑，不是单向强化。'' asheep 长老平静道，``高阶反射作用于低阶集合，可能提纯；低阶反射作用于高阶集合，往往只会筛出极限点，而丢失原有的容许性质。这是序道修行中极易走偏的误区，你们务必牢记。''

``继续递推，\allowbreak{}\NovelInline{I02535}{1-1-2}\allowbreak{}，也就是对 \allowbreak{}\NovelInline{I02536}{1-2}\allowbreak{} 再取一次极限点，可以得到`\allowbreak{}\NovelInline{I02537}{\Omega_\alpha}\allowbreak{}，其中 \allowbreak{}\NovelInline{I02538}{\alpha}\allowbreak{} 为 \allowbreak{}\NovelInline{I02539}{\omega^2}\allowbreak{} 的倍数'的集合。同样地，我们可以通过取交集得到 \allowbreak{}\NovelInline{I02540}{(1-)^\omega 2}\allowbreak{}、\allowbreak{}\NovelInline{I02541}{(1-)^{\omega^2} 2}\allowbreak{}，乃至 \allowbreak{}\NovelInline{I02542}{(1-)^\Omega 2}\allowbreak{}。''

``由于 \allowbreak{}\NovelInline{I02543}{\Omega = \min 2}\allowbreak{}，所以 \allowbreak{}\NovelInline{I02544}{(1-)^\Omega 2}\allowbreak{} 又可以写作 \allowbreak{}\NovelInline{I02545}{(1-)^{(2)} 2}\allowbreak{}。这里在上标上省略了 \allowbreak{}\NovelInline{I02546}{\min}\allowbreak{}，并且在外面加了一层括号用以和序数 \allowbreak{}\NovelInline{I02547}{2}\allowbreak{} 区分。''

唐三顺着推导思路继续追问：``那如果我们把不动点模式也套用到二阶集合上，也就是 \allowbreak{}\NovelInline{I02548}{(1-)^{1,0} 2}\allowbreak{}，会得到什么？应当不再是普通的 \allowbreak{}\NovelInline{I02549}{\varepsilon}\allowbreak{} 序数集合了吧？''

``很好，已经摸到门道了。'' 梅天狸笑道，``类似 \allowbreak{}\NovelInline{I02550}{(1-)^{1,0}}\allowbreak{} 的定义，\allowbreak{}\NovelInline{I02551}{(1-)^{1,0} 2 = \{\alpha \mid \alpha = \min (1-)^\alpha 2\}}\allowbreak{}。不过，现在它代表的可就不是`\allowbreak{}\NovelInline{I02552}{\Omega_\alpha}\allowbreak{}，其中 \allowbreak{}\NovelInline{I02553}{\alpha}\allowbreak{} 是 \allowbreak{}\NovelInline{I02554}{\varepsilon}\allowbreak{} 序数'的集合了，而是映射 \allowbreak{}\NovelInline{I02555}{\alpha \mapsto \Omega_{\omega^\alpha}}\allowbreak{} 的\textbf{不动点}。''

``不难发现，这样的不动点必然也是 \allowbreak{}\NovelInline{I02556}{\alpha \mapsto \omega^\alpha}\allowbreak{} 的不动点，也就是说，\allowbreak{}\NovelInline{I02557}{(1-)^{1,0} 2}\allowbreak{} 的每一个元素 \allowbreak{}\NovelInline{I02558}{\alpha}\allowbreak{} 都满足 \allowbreak{}\NovelInline{I02559}{\alpha = \Omega_\alpha}\allowbreak{} ------ 它们都是 OFP。''

``因此，\allowbreak{}\NovelInline{I02560}{(1-)^{1,0} 2}\allowbreak{} 是全体 OFP 的集合：\allowbreak{}\NovelInline{I02561}{\{\psi_I(0), \psi_I(1), \dots, \psi_I(\Omega), \dots, \psi_I(I), \dots\}}\allowbreak{}，也可以利用 \allowbreak{}\NovelInline{I02562}{\Phi}\allowbreak{} 函数写成 \allowbreak{}\NovelInline{I02563}{\{\Phi(1,0), \Phi(1,1), \dots, \Phi(1,\Omega), \dots, \Phi(2,0), \dots\}}\allowbreak{}。\allowbreak{}\NovelInline{I02564}{\psi_I}\allowbreak{} 函数之于 \allowbreak{}\NovelInline{I02565}{\Phi}\allowbreak{} 函数，正如同 \allowbreak{}\NovelInline{I02566}{\psi}\allowbreak{} 函数之于 \allowbreak{}\NovelInline{I02567}{\varphi}\allowbreak{} 函数。''

``\allowbreak{}\NovelInline{I02568}{(1-)^{1,0} 2}\allowbreak{} 过后，同样也有 \allowbreak{}\NovelInline{I02569}{(1-)^{2,0} 2}\allowbreak{}、\allowbreak{}\NovelInline{I02570}{(1-)^{1,0,0} 2 \dots}\allowbreak{} 它们可以用多元 \allowbreak{}\NovelInline{I02571}{\Phi}\allowbreak{} 函数表示。不过，不管 \allowbreak{}\NovelInline{I02572}{(1-)^{a,b,c,\dots} 2}\allowbreak{} 的上标的层级多深，只要它还是不动点的形式，它就是非容许的。''

讲到此处，话题终于绕回了方才卢本韦抛出的那个核心问题。

asheep 长老往前踏出一步，周身序道气息骤然凝重，殿内灵纹都随之微微震颤。

``现在，我们回到最初的问题：是否存在极限序数 \allowbreak{}\NovelInline{I02573}{\alpha}\allowbreak{}，使得 \allowbreak{}\NovelInline{I02574}{\Omega_\alpha}\allowbreak{} 为容许序数？''

``我们不妨先推导，若这样的 \allowbreak{}\NovelInline{I02575}{\alpha}\allowbreak{} 存在，它需要满足哪些条件： * 第一，\allowbreak{}\NovelInline{I02576}{\alpha}\allowbreak{} 是 OFP。因为 \allowbreak{}\NovelInline{I02577}{\Omega_\alpha}\allowbreak{} 通过`遍历下标'的形式可以找到一条长度至多为 \allowbreak{}\NovelInline{I02578}{\alpha}\allowbreak{} 的基本列，要让 \allowbreak{}\NovelInline{I02579}{\Omega_\alpha}\allowbreak{} 的基本列长度等于它自身，唯有 \allowbreak{}\NovelInline{I02580}{\alpha = \Omega_\alpha}\allowbreak{}； * 第二，\allowbreak{}\NovelInline{I02581}{\alpha}\allowbreak{} 是容许序数，即 \allowbreak{}\NovelInline{I02582}{\alpha \in 2}\allowbreak{}。由`\allowbreak{}\NovelInline{I02583}{\alpha = \Omega_\alpha}\allowbreak{}'和`\allowbreak{}\NovelInline{I02584}{\Omega_\alpha}\allowbreak{} 为容许序数'可直接推出； * 第三，\allowbreak{}\NovelInline{I02585}{\alpha \in 1-2}\allowbreak{}。由`\allowbreak{}\NovelInline{I02586}{\alpha = \Omega_\alpha}\allowbreak{}'和`\allowbreak{}\NovelInline{I02587}{\alpha}\allowbreak{} 为极限序数'可直接推出。''

``总结一下，满足上述条件的 \allowbreak{}\NovelInline{I02588}{\alpha}\allowbreak{} 属于 \allowbreak{}\NovelInline{I02589}{2}\allowbreak{} 和 \allowbreak{}\NovelInline{I02590}{1-2}\allowbreak{} 的交集，记为 \allowbreak{}\NovelInline{I02591}{\alpha \in 2 \cap 1-2}\allowbreak{}。在我们的简写体系中，由于 \allowbreak{}\NovelInline{I02592}{\cap}\allowbreak{} 可以用空格表示，所以它记作：\allowbreak{}\NovelInline{I02593}{2 \ 1-2}\allowbreak{}。''

\begin{quote}
\emph{``Anything else? No that's all.''}
\end{quote}

小舞眨了眨眼，轻声问：``那\ldots\ldots{} \allowbreak{}\NovelInline{I02594}{2}\allowbreak{} 和 \allowbreak{}\NovelInline{I02595}{1-2}\allowbreak{} 的交集，真的存在元素吗？我们推导了半天条件，可万一它是空集呢？''

``这正是最关键的地方。'' 梅天狸郑重道，``实质性的问题其实并未得到解决 ------ 仅靠定义推导，我们无法证明它非空。所幸，根据 ZFC 公理体系，我们可以直接证明 \allowbreak{}\NovelInline{I02596}{2}\allowbreak{} 和 \allowbreak{}\NovelInline{I02597}{1-2}\allowbreak{} 的交集存在元素，我们将其命名为\textbf{递归不可达序数}。''

``最小的递归不可达序数记作 \allowbreak{}\NovelInline{I02598}{I}\allowbreak{}，次小的是 \allowbreak{}\NovelInline{I02599}{I_2 \dots}\allowbreak{} 就和 \allowbreak{}\NovelInline{I02600}{\Omega, \Omega_2}\allowbreak{} 差不多。''

``总之，\allowbreak{}\NovelInline{I02601}{2 \ 1-2 = \{I, I_2, \dots, I_{\omega+1}, \dots, I_{I+1}, \dots, I(1,0), \dots\}}\allowbreak{}。注意到这里也没有列出 \allowbreak{}\NovelInline{I02602}{I_\omega}\allowbreak{}，\allowbreak{}\NovelInline{I02603}{I_\omega}\allowbreak{} 并非容许序数，道理和 \allowbreak{}\NovelInline{I02604}{\Omega_\omega}\allowbreak{} 完全一致。''

``接下来，我们还可以对 \allowbreak{}\NovelInline{I02605}{2 \ 1-2}\allowbreak{} 进行若干次 \allowbreak{}\NovelInline{I02606}{1-}\allowbreak{}，例如：\allowbreak{}\NovelInline{I02607}{1-2 \ 1-2 = \{I_\omega, I_{\omega \times 2}, \dots, I_I, \dots, \text{IFP}, \dots\}}\allowbreak{}。推导过程和 \allowbreak{}\NovelInline{I02608}{1-2}\allowbreak{} 是一致的，注意 \allowbreak{}\NovelInline{I02609}{1-2 \ 1-2}\allowbreak{} 表示的实际上是 \allowbreak{}\NovelInline{I02610}{1\text{ onto }(2 \ 1-2)}\allowbreak{}。\allowbreak{}\NovelInline{I02611}{(1-)^{1,0} 2 \ 1-2}\allowbreak{} 自然是全体 IFP，也就是满足 \allowbreak{}\NovelInline{I02612}{\alpha = I_\alpha}\allowbreak{} 的序数构成的集合。''

卢本韦若有所思：``\allowbreak{}\NovelInline{I02613}{\Omega}\allowbreak{} 对应 \allowbreak{}\NovelInline{I02614}{\omega^\alpha}\allowbreak{} 的不动点，\allowbreak{}\NovelInline{I02615}{I}\allowbreak{} 对应 \allowbreak{}\NovelInline{I02616}{\Omega_\alpha}\allowbreak{} 的不动点\ldots\ldots{} 好像有一个统一的规律？''

``没错，这就是\textbf{容许点}的概念。'' 梅天狸道，``在更早的文章里，我曾将`容许点'刻画为`足以放进下标输出各种不动点'的序数。但在接触过反射之后，我们可以有一个更好的定义： \textbf{序数 \allowbreak{}\NovelInline{I02617}{\Theta}\allowbreak{} 为函数 \allowbreak{}\NovelInline{I02618}{f(\alpha)}\allowbreak{} 的容许点，等同于 \allowbreak{}\NovelInline{I02619}{\Theta}\allowbreak{} 是 \allowbreak{}\NovelInline{I02620}{f(\alpha)}\allowbreak{} 的不动点且 \allowbreak{}\NovelInline{I02621}{\Theta}\allowbreak{} 为容许序数。}''

``于是我们可以看到： * \allowbreak{}\NovelInline{I02622}{\Omega}\allowbreak{} 是函数 \allowbreak{}\NovelInline{I02623}{f(\alpha) = \omega^\alpha}\allowbreak{} 的容许点，\allowbreak{}\NovelInline{I02624}{\psi_\Omega}\allowbreak{}（即 \allowbreak{}\NovelInline{I02625}{\psi}\allowbreak{}）输出 \allowbreak{}\NovelInline{I02626}{f(\alpha) = \omega^\alpha}\allowbreak{} 的各种层级的不动点； * \allowbreak{}\NovelInline{I02627}{I}\allowbreak{} 是函数 \allowbreak{}\NovelInline{I02628}{f(\alpha) = \Omega_\alpha}\allowbreak{} 的容许点，\allowbreak{}\NovelInline{I02629}{\psi_I}\allowbreak{} 输出 \allowbreak{}\NovelInline{I02630}{f(\alpha) = \Omega_\alpha}\allowbreak{} 的各种层级的不动点。''

``更一般地，若函数 \allowbreak{}\NovelInline{I02631}{f(\alpha)}\allowbreak{} 的值域是集合 \allowbreak{}\NovelInline{I02632}{A}\allowbreak{}，那么 \allowbreak{}\NovelInline{I02633}{f(\alpha)}\allowbreak{} 的容许点组成的集合是 \allowbreak{}\NovelInline{I02634}{2 \ 1-A}\allowbreak{}。''

``在更以后的地方，提到的某函数的`马洛点'\,`K 点'等`X 点'概念，也都指代`某序数是 X 序数且是该函数的不动点'。''

唐三沉吟片刻，顺着容许点的逻辑继续推演：``既然 \allowbreak{}\NovelInline{I02635}{\Omega}\allowbreak{} 是 \allowbreak{}\NovelInline{I02636}{\omega^\alpha}\allowbreak{} 的容许点，\allowbreak{}\NovelInline{I02637}{I}\allowbreak{} 是 \allowbreak{}\NovelInline{I02638}{\Omega_\alpha}\allowbreak{} 的容许点，那函数 \allowbreak{}\NovelInline{I02639}{f(\alpha) = I_\alpha}\allowbreak{} 的容许点又是什么？难道是比 \allowbreak{}\NovelInline{I02640}{I}\allowbreak{} 更高一阶的 \allowbreak{}\NovelInline{I02641}{M}\allowbreak{} 之类的大序数？''

``很多人都会有这个猜测，受 IMK 并称的习惯影响，容易直接跳到 \allowbreak{}\NovelInline{I02642}{M}\allowbreak{}。'' asheep 长老微微摇头，``但实际上，函数 \allowbreak{}\NovelInline{I02643}{f(\alpha) = I_\alpha}\allowbreak{} 的最小的容许点仅仅是 \allowbreak{}\NovelInline{I02644}{I(1,0)}\allowbreak{}。''

``这意味着 \allowbreak{}\NovelInline{I02645}{I(1,0)}\allowbreak{} 可以作为下标输出 \allowbreak{}\NovelInline{I02646}{f(\alpha) = I_\alpha}\allowbreak{} 的不动点，例如 \allowbreak{}\NovelInline{I02647}{\psi_{I(1,0)}(0)}\allowbreak{} 是最小的 IFP。受《大数入门》的影响，部分读者会认为 \allowbreak{}\NovelInline{I02648}{I(1,0)}\allowbreak{} 与 IFP 等同，实则并非如此。''

``根据前面提到的规则 ------ 若函数 \allowbreak{}\NovelInline{I02649}{f(\alpha)}\allowbreak{} 的值域是集合 \allowbreak{}\NovelInline{I02650}{A}\allowbreak{}，那么 \allowbreak{}\NovelInline{I02651}{f(\alpha)}\allowbreak{} 的容许点组成的集合是 \allowbreak{}\NovelInline{I02652}{2 \ 1-A}\allowbreak{} ------ 不难知道，函数 \allowbreak{}\NovelInline{I02653}{f(\alpha) = I_\alpha}\allowbreak{} 的全体容许点的集合是 \allowbreak{}\NovelInline{I02654}{2 \ 1-2 \ 1-2}\allowbreak{}。''

``\allowbreak{}\NovelInline{I02655}{f(\alpha) = I_\alpha}\allowbreak{} 的次小的容许点是 \allowbreak{}\NovelInline{I02656}{I(1,1)}\allowbreak{}。这里，二元 \allowbreak{}\NovelInline{I02657}{I}\allowbreak{} 的最后一位实质上相当于 \allowbreak{}\NovelInline{I02658}{\Omega}\allowbreak{} 或者 \allowbreak{}\NovelInline{I02659}{I}\allowbreak{} 的下标。于是：\allowbreak{}\NovelInline{I02660}{(2 \ 1-)^2 2 = 2 \ 1-2 \ 1-2 = \{I(1,0), I(1,1), \dots, I(1, \omega+1), \dots\}}\allowbreak{}。\allowbreak{}\NovelInline{I02661}{I(1,\omega)}\allowbreak{} 同样是非容许的。''

``\allowbreak{}\NovelInline{I02662}{(1-)^{1,0} 2 \ 1-2 \ 1-2}\allowbreak{} 是函数 \allowbreak{}\NovelInline{I02663}{f(\alpha) = I(1, \alpha)}\allowbreak{} 的不动点，记作 \allowbreak{}\NovelInline{I02664}{I(1,a)\text{FP}}\allowbreak{}。最小的 \allowbreak{}\NovelInline{I02665}{I(1,a)\text{FP}}\allowbreak{} 也并非 \allowbreak{}\NovelInline{I02666}{I(2,0)}\allowbreak{}，而是 \allowbreak{}\NovelInline{I02667}{\psi_{I(2,0)}(0)}\allowbreak{}。''

``\allowbreak{}\NovelInline{I02668}{I(2,0)}\allowbreak{} 才是函数 \allowbreak{}\NovelInline{I02669}{f(\alpha) = I(1, \alpha)}\allowbreak{} 的容许点。多元 \allowbreak{}\NovelInline{I02670}{I}\allowbreak{} 函数的这种进位方式，称为\textbf{容许点进位}。''

讲到这里，梅天狸抬手一挥，虚空中重新浮现出完整的反射式总结列表，比先前的版本更加详尽：

\NovelDisplay{M00489R01}{%
2 = \{\Omega, \Omega_2, \dots, \Omega_{\omega+1}, \dots, I, \dots\}%
}
\NovelDisplay{M00489R02}{%
1-2 = \{\Omega_\omega, \Omega_{\omega \times 2}, \dots, \Omega_{\omega^2}, \dots, \psi_I(0), \dots, I, \dots\}%
}
\NovelDisplay{M00489R03}{%
(1-)^{1,0} 2 = \{\psi_I(0), \psi_I(1), \dots, \psi_I(\omega), \dots, \psi_I(I), \dots, I, \dots\}%
}
\NovelDisplay{M00489R04}{%
2 \ 1-2 = \{I, I_2, \dots, I_{\omega+1}, \dots, I(1,0), \dots\}%
}
\NovelDisplay{M00489R05}{%
(1-)^{1,0} 2 \ 1-2 = \{\psi_{I(1,0)}(0), \psi_{I(1,0)}(1), \dots, \psi_{I(1,0)}(\omega), \dots, I(1,0), \dots\}%
}
\NovelDisplay{M00489R06}{%
(2 \ 1-)^2 2 = \{I(1,0), I(1,1), \dots, I(1, \omega+1), \dots, I(2,0), \dots\}%
}
\NovelDisplay{M00489R07}{%
(1-)^{1,0} (2 \ 1-)^2 2 = \{\psi_{I(2,0)}(0), \psi_{I(2,0)}(1), \dots, \psi_{I(2,0)}(\omega), \dots, I(2,0), \dots\}%
}
\NovelDisplay{M00489R08}{%
(2 \ 1-)^3 2 = \{I(2,0), I(2,1), \dots, I(2, \omega+1), \dots, I(3,0), \dots\}%
}

``像这样继续往前，自然会遇到 \allowbreak{}\NovelInline{I02671}{(2 \ 1-)^\omega 2}\allowbreak{}，我们期望它是 \allowbreak{}\NovelInline{I02672}{\{I(\omega,0), I(\omega,1), \dots\}}\allowbreak{}。'' 梅天狸的声音忽然放缓，带着一丝深意，``但如果我们仍用定义 \allowbreak{}\NovelInline{I02673}{(1-)^\omega}\allowbreak{} 时的交集来定义它，会怎么样呢？''

asheep 长老接过话头，清冷的声音里藏着未尽之意：``这一步，已经触及了反射体系深层的分歧。下一期，才会真正开始涉及真伪的争端。''

\chapter{第五十一章：真伪之争}\label{ux7b2cux4e94ux5341ux4e00ux7ae0ux771fux4f2aux4e4bux4e89}

殿内灵纹渐渐归于平静，漫天序数符号缓缓隐去。果糕小队众人只觉心神震荡，从一阶极限反射到二阶容许序数，再到递归不可达、容许点进位，短短一次传道，便为他们推开了一扇远超认知的序道大门。

他们此刻才真正明白，为何 Opt 老祖能以序道镇压八十八万年龙王、平定万古格局 ------ 这套反射序数体系，本身就是对天地秩序的层层解构与重塑，根基越深，能执掌的规则便越强。

``今日所传，便是二阶反射的全部筑基内容。'' 梅天狸环视众人，``从今日起，你们便在扽魂殿闭关，从 \allowbreak{}\NovelInline{I02674}{1-2}\allowbreak{} 开始逐层感悟，先吃透一阶反射对二阶集合的作用，再逐步接触 \allowbreak{}\NovelInline{I02675}{2 \ 1-2}\allowbreak{} 与递归不可达的道韵。''

``切记，序道修行稳扎稳打，不可贪功冒进。容许序数、递归不可达序数，皆是天地本源秩序，稍有不慎便会道基紊乱、走火入魔。''

殿内灵纹的最后一缕微光缓缓敛入石壁，扽魂殿重归一片沉凝的寂静。

果糕小队众人依旧盘膝端坐，久久未能从方才的传道震撼中抽身。许久之后，诺当才缓缓松开紧握的双拳，喉结微动，吐出一句满是复杂的慨叹，一语道破了所有人心中的震荡：

``想我前几日突破至 \allowbreak{}\NovelInline{I02676}{I}\allowbreak{} 序数境界时，还道自己已然触摸到了序道高阶门槛，修为远超同辈，足以应对前路凶险\ldots\ldots{} 今日听二位长老讲完完整的反射体系才知晓，原来我引以为傲的境界，竟也仅仅是反射序数大道里最粗浅的入门知识。''

此言一出，众人皆默然点头，心绪翻涌。小舞轻轻颔首，眉眼间满是恍然与唏嘘：``此前我只觉得递归不可达序数已是遥不可及的至高境界，谁能想到，最小的递归不可达序数 \allowbreak{}\NovelInline{I02677}{I}\allowbreak{}，不过是 \allowbreak{}\NovelInline{I02678}{2 \ 1-2}\allowbreak{} 集合里最微小的第一个元素。往后还有 \allowbreak{}\NovelInline{I02679}{I(1,0)}\allowbreak{}、容许点进位、多元 \allowbreak{}\NovelInline{I02680}{I}\allowbreak{} 函数\ldots\ldots{} 无穷无尽的层级环环相扣，我们连二阶反射的真正门径都才刚摸到。''

霍雨浩眉头微蹙，指尖无意识地在虚空中勾勒着方才记下的反射式，沉声道：``我们先前对序道的认知，实在太过浅薄。从前只盯着修为境界的高低、灵力的厚薄，却从未窥见这套反射体系的全貌。\allowbreak{}\NovelInline{I02681}{\Pi_1}\allowbreak{} 极限反射是筑基打底，\allowbreak{}\NovelInline{I02682}{\Pi_2}\allowbreak{} 容许序数才算踏入二阶门径，而递归不可达，不过是二阶反射交集运算里诞生的第一类高阶不动点。再往上还有 \allowbreak{}\NovelInline{I02683}{\Pi_3}\allowbreak{}、\allowbreak{}\NovelInline{I02684}{\Pi_4}\allowbreak{} 乃至更高阶的反射奥秘，我们更是闻所未闻。''

唐三眸光深沉，心中早已将前后脉络串联起来，语气郑重：``也难怪乾坤问情谷的禁忌之力如此恐怖，连 Opt 老祖都要留下严诫。我们这点 \allowbreak{}\NovelInline{I02685}{I}\allowbreak{} 序数层级的修为，在完整的反射序数体系面前，不过是尘埃一般的存在。谷中潜藏的力量能颠覆天地秩序，正是因为它触及了我们远未抵达的高阶序则。''

卢本韦指尖摩挲着衣襟内温热的猪神令，远古精铁肌理中沉眠的老祖道韵微微起伏，与殿内残留的序道气息隐隐共鸣。她轻声开口，语气里带着对先祖传承的更深敬畏：``先祖 Opt 老祖当年以序道定乾坤，只身镇压八十八万年龙王、调和人兽两族秩序，想来早已将这套反射体系修至极致，执掌了远超 \allowbreak{}\NovelInline{I02686}{\Pi_2}\allowbreak{}、\allowbreak{}\NovelInline{I02687}{\Pi_3}\allowbreak{} 的高阶序则。我们今日习得的全部内容，不过是老祖万古传承中最根基的皮毛。''

众人闻言皆是心头一震。他们此前突破序境、踏足 \allowbreak{}\NovelInline{I02688}{I}\allowbreak{} 层级时，心中未尝没有几分少年天骄的自得，只觉自己已走在同辈前列，距离顶尖强者不再遥远。可今日一堂传道讲完，整套反射序数的浩瀚体系徐徐铺开，他们才惊觉自己昔日眼界之狭小 ------ 自己拼尽全力突破的境界，在这条大道上，竟连入门的第一层台阶都尚未完全站稳。没有沮丧，没有颓唐，短暂的震撼过后，众人心底反而燃起了更盛的求索之意。先前突破境界带来的浮躁尽数褪去，取而代之的是对大道的敬畏，以及沉心苦修的笃定。

他们终于明白，为何梅天狸长老与 asheep 长老执意要他们在此闭关夯实根基。反射大道层层递进、环环相扣，根基差之毫厘，前路便谬以千里。若不把 \allowbreak{}\NovelInline{I02689}{\Pi_1}\allowbreak{}、\allowbreak{}\NovelInline{I02690}{\Pi_2}\allowbreak{} 这两层基础彻底吃透，贸然踏入乾坤问情谷，在更高阶的逆乱秩序面前，只会落得道基崩碎的下场。

``苦修吧。'' 诺当深吸一口气，眼底锋芒重聚，却比往日多了几分沉稳，``从前是我坐井观天。今日方知反射大道浩瀚，我们唯有从最基础的层级一步步打磨，方能真正掌控这份力量。''

``不错。'' 唐三微微颔首，缓缓闭上双目，``从今日起，沉心推演 \allowbreak{}\NovelInline{I02691}{1-2}\allowbreak{}、\allowbreak{}\NovelInline{I02692}{2 \ 1-2}\allowbreak{}，吃透二阶反射的每一条规则，直至融会贯通。''

殿外天光流转，殿内气息归宁。果糕小队众人收敛所有杂念，闭目沉神，灵台之上序数流转，正式开启了二阶反射体系的沉潜苦修。

扽魂殿的石门缓缓闭合，将外界喧嚣尽数隔绝。门内，属于新一代序道修士的进阶之路，才刚刚正式启程。

\begin{center}\rule{0.5\linewidth}{0.5pt}\end{center}

扽魂殿的石门缓缓闭合，将少年们的苦修与外界隔绝。梅天狸与 asheep 长老没有停留，转身沿殿后密道一路下行，直抵扽魂殿最核心的禁地 ------ 良序圣殿。

这座圣殿藏于山体腹心，不接天光，却通体被柔和的金白色序光笼罩。殿壁之上，无数细密序纹交织成舒展的良序羽翼形态，层层叠叠蔓延至穹顶，每一道纹路都流淌着沉凝如山的本源序力。此处是扽魂殿立道以来的根基所在，也是大长老 \texttt{test\_alpha0} 的闭关之所，非殿中核心执掌者不得踏入半步。

殿中空旷，唯有正中一座白玉云床。\texttt{test\_alpha0} 端坐其上，一身素白道袍，须发皆白却丝毫不显衰颓，眉目间依稀镌刻着当年纵横天下的锐利锋芒。他周身羽翼状的序纹缓缓沉浮，气息深不见底，正是早已登临 \allowbreak{}\NovelInline{I02693}{\omega\text{-Y}(1,\omega)}\allowbreak{} 境界的当世序道巨擘。

``来了。'' 他未睁眼，声线苍老却平正，在空旷圣殿中泛起淡淡回音。

梅天狸与 asheep 上前躬身行礼。``大长老。''

``\allowbreak{}\NovelInline{I02694}{\Pi_2}\allowbreak{} 反射讲完了？那群孩子悟性如何？''

梅天狸抬头，语气带着难掩的赞许：``回大长老，\allowbreak{}\NovelInline{I02695}{\Pi_1}\allowbreak{} 体系一日便吃透，今日 \allowbreak{}\NovelInline{I02696}{\Pi_2}\allowbreak{} 容许序数、递归不可达与容许点进位的全套推演，八人全程跟上，无一人掉队。尤其是对 \allowbreak{}\NovelInline{I02697}{I}\allowbreak{} 序数的理解，并非死记结论，而是从`容许交集极限点'的本源处便已通透。''

asheep 在旁补了一句：``蔡须坤道基浑厚，卢本韦有家学渊源，唐三、霍雨浩悟性最是惊艳，常能顺着推演自行触达下一层逻辑。''

\texttt{test\_alpha0} 听罢，缓缓睁开双眼。眸中金白色序光一闪而逝，随即化作几分讶异，随即便是一声带着追忆的长叹。

``了不起\ldots\ldots{} 真是了不起。'' 他声音放轻了些，仿佛透过眼前的虚空，看见了自己年少时鲜衣怒马的模样，``想老夫当年，正是意气风发的年纪。二十九岁，便修到了俗世封号斗罗的修为，序道上稳稳踏过 \allowbreak{}\NovelInline{I02698}{\text{Y}(1,3)}\allowbreak{} 门槛，在同辈之中一骑绝尘。那时候宗门大比连战百场未尝一败，多少前辈断言我是扽魂殿千年不遇的奇才。''

``及至六十六岁，老夫踏破层层桎梏，摸到了 \allowbreak{}\NovelInline{I02699}{\omega\text{-Y}(1,\omega)}\allowbreak{} 的境界边界，一举登临当世序道巅峰之列。那时候纵横南北，遍访天下名家，同辈之中难寻抗手，只觉天地辽阔，序道巅峰不过咫尺。''

他说到此处，微微摇头，眉宇间掠过一丝岁月沉淀的怅然：``如今活了三百二十余年，见惯了宗门起落、天才凋零，本以为再难见到能比肩当年的后辈。没想到这群孩子，理解 \allowbreak{}\NovelInline{I02700}{I}\allowbreak{} 序数的速度，竟比老夫当年还早了一年有余。''

``千古奇才，莫过于此。''

殿中寂静片刻，唯有羽翼序纹缓缓流淌的微光。梅天狸望着云床上的老者，心中亦是感慨万千。她深知大长老当年二十九岁封道、六十六岁登峰的战绩，在序道界流传百年，一直是后辈修士仰望的传奇。纵是这般惊才绝艳的人物，如今也对果糕小队的天资给出了如此高的评价。

她轻声宽慰：``大长老道基深厚，寿元绵长，定能亲眼见到他们追上您的脚步。''

\texttt{test\_alpha0} 却只是淡淡一笑，抬手抚过膝头的道袍，语气平静：``大道有定数，老夫的路早已走到头了。这副身子骨，撑不了多少年，怕是等不到他们成长到我这般境界了。''

他话锋一转，语气郑重起来：``不过既然他们悟性足够，就不必按部就班磨日子。\allowbreak{}\NovelInline{I02701}{\Pi_2}\allowbreak{} 根基扎稳之后，立刻切入 \allowbreak{}\NovelInline{I02702}{\Pi_3}\allowbreak{} 反射与马洛基数的内容。''

``弟子明白。'' 梅天狸与 asheep 齐齐躬身领命。

圣殿之中序光静静流淌，老者的目光穿过厚重的山壁，仿佛落在了远处正殿里闭关的八个少年身上。当年他意气风发，以一己之力撑起扽魂殿的序道传承；如今薪火相传，新一代的天骄正在破土而出，即将接过这副沉甸甸的担子。

扽魂殿正殿静如止水。

八名少年少女盘膝端坐，周身流转的序道灵纹规整而稳固。经过一夜闭关沉淀，众人早已吃透二阶反射的核心根基：容许序数的判定准则、递归不可达序数 \allowbreak{}\NovelInline{I02703}{I}\allowbreak{} 的由来、二元与多元 \allowbreak{}\NovelInline{I02704}{I}\allowbreak{} 函数的进位逻辑、容许点的本质定义，尽数融入道基，再无半分晦涩滞涩。

殿门隔绝俗世喧嚣，却挡不住殿内传道道韵流转。梅天狸与 asheep 长老自后山良序圣殿踏序而出，一袭素衣映着漫天序数微光，缓步落至高台。虚空中，完整版二阶反射对照表徐徐展开，层层符号递进排布，唯独最顶端的极限迭代项 \allowbreak{}\NovelInline{I02705}{(2\ 1\text{-})^{\omega}\ 2}\allowbreak{} 悬而未决，留白处暗藏序道最关键的真伪分歧，也是困住无数序道修士的核心误区。

梅天狸眸光扫过众人，见八人气息凝练、心神澄澈，已然具备深究极限序数迭代的根基，当即开口传道，声线清泠庄重：

\begin{quote}
``昨日收尾，我们停在 \allowbreak{}\NovelInline{I02706}{(2\ 1\text{-})^{\omega}\ 2}\allowbreak{} 的边界之上。有限次 \allowbreak{}\NovelInline{I02707}{(2\ 1\text{-})}\allowbreak{} 迭代秩序恒定，可一旦触及 \allowbreak{}\NovelInline{I02708}{\omega}\allowbreak{} 极限无穷，体系便会分裂为两套截然不同的规则，衍生出 \textbf{psd 伪定义}与 \textbf{real 真定义}的千古之争。今日，我们先破真伪之辨，再穷尽 \allowbreak{}\NovelInline{I02709}{I}\allowbreak{} 函数的容许进位极限，最终触碰 \allowbreak{}\NovelInline{I02710}{\Pi_2}\allowbreak{} 体系不可逾越的天堑------\textbf{递归马洛序数}。''
\end{quote}

\begin{center}\rule{0.5\linewidth}{0.5pt}\end{center}

话音落，蔡须坤率先睁眼拱手，眸中满是审慎的求索：

\begin{quote}
``长老，弟子推演之时发现，有限 \allowbreak{}\NovelInline{I02711}{(2\ 1\text{-})^n\ 2}\allowbreak{} 皆对应固定的 \allowbreak{}\NovelInline{I02712}{I}\allowbreak{} 序数集合，性质统一、边界清晰。为何唯独极限迭代的 \allowbreak{}\NovelInline{I02713}{(2\ 1\text{-})^{\omega}\ 2}\allowbreak{} 会出现两套定义？二者的核心差异，究竟落在何处？''
\end{quote}

``此问直击本源。'' asheep 长老上前半步，清冷声线带着严苛的数理秩序，抬手引动虚空序纹，两套泾渭分明的序列同时浮现，一者宽泛包容、一者纯粹凝练：

\begin{quote}
``分歧根源，在于极限序数的迭代规则不同。有限迭代收敛唯一，\allowbreak{}\NovelInline{I02714}{\omega}\allowbreak{} 极限迭代却有两种处理方式，由此分出 \textbf{psd} 与 \textbf{real} 两套体系，二者形式相似、性质天差地别，也是序道修行最易混淆、最易走火的关键。''
\end{quote}

\begin{quote}
``先讲 \textbf{psd} \allowbreak{}\NovelInline{I02715}{(2\ 1\text{-})^{\omega}\ 2}\allowbreak{}，为本系列专栏默认标准定义。''
\end{quote}

asheep 指尖划过浅层序纹，无数二元 \allowbreak{}\NovelInline{I02716}{I}\allowbreak{} 序数尽数浮现：

\begin{quote}
``它遵循反射闭包规则，不取交集收敛，而是收纳所有 \allowbreak{}\NovelInline{I02717}{I(\omega, \cdot)}\allowbreak{} 型序数，是全体二元极限 \allowbreak{}\NovelInline{I02718}{I}\allowbreak{} 序数的集合。包含 \allowbreak{}\NovelInline{I02719}{I(\omega, 0)}\allowbreak{}、\allowbreak{}\NovelInline{I02720}{I(\omega, 1)}\allowbreak{}、\allowbreak{}\NovelInline{I02721}{I(\omega, \omega)}\allowbreak{}、\allowbreak{}\NovelInline{I02722}{I(\omega, \varepsilon_0)}\allowbreak{} 等一切衍生形态，不分容许与非容许，只论层级归属。''
\end{quote}

霍雨浩眉头微展，瞬间吃透关键细节：

\begin{quote}
``弟子明白了！此前所学，\allowbreak{}\NovelInline{I02723}{\alpha}\allowbreak{} 为零或后继序数时，\allowbreak{}\NovelInline{I02724}{I(\omega, \alpha)}\allowbreak{} 皆存在长度为 \allowbreak{}\NovelInline{I02725}{\omega}\allowbreak{} 的基本列，属于非容许序数。psd 定义包容万象，将这些不稳定的非容许序数尽数纳入，因此集合层级驳杂、秩序松散。''
\end{quote}

``分毫不错。'' 梅天狸颔首接过话头，转而点亮深层纯白序纹，展出截然相反的真性规则：

\begin{quote}
``而 \textbf{real} \allowbreak{}\NovelInline{I02726}{(2\ 1\text{-})^{\omega}\ 2}\allowbreak{}，是严苛的真性序数集合，取所有有限迭代的交集收敛，定义为：

\NovelDisplay{M00490R01}{%
\text{real.}(2\ 1\text{-})^{\omega}\ 2 = \bigcap_{n<\omega} (2\ 1\text{-})^{n}\ 2%
}
\end{quote}

``交集筛选的核心作用，是去芜存菁、剔除一切不稳定秩序。'' 她加重语气，字字剖析本质：

\begin{quote}
``能进入 real 集合的序数，必须同时满足\textbf{双重铁律}： 1. 是 \allowbreak{}\NovelInline{I02727}{\Omega_\alpha}\allowbreak{}、\allowbreak{}\NovelInline{I02728}{I_\alpha}\allowbreak{}、\allowbreak{}\NovelInline{I02729}{I(1, \alpha)}\allowbreak{} 等所有前置序函数的公共不动点； 2. 本身是纯粹的\textbf{容许序数}，无任何低级基本列可以逼近。''
\end{quote}

小舞眸光清亮，轻声总结：

\begin{quote}
``也就是说，psd 是包罗万象的序域集合，容纳所有极限迭代产物，包含大量非容许序数；real 是提纯到极致的序道真核，只留存稳固的高阶容许秩序？''
\end{quote}

``正是如此。'' asheep 微微点头，道出二者的递进转化关系：

\begin{quote}
``psd \allowbreak{}\NovelInline{I02730}{(2\ 1\text{-})^{\omega}\ 2}\allowbreak{} 囊括的 \allowbreak{}\NovelInline{I02731}{I(\omega, \alpha)}\allowbreak{} 大多非容许，秩序残缺；而 real \allowbreak{}\NovelInline{I02732}{(2\ 1\text{-})^{\omega}\ 2}\allowbreak{}，正是函数 \allowbreak{}\NovelInline{I02733}{f(\beta) = I(\omega, \beta)}\allowbreak{} 的最小容许点。伪定义经一次 \allowbreak{}\NovelInline{I02734}{2\ 1\text{-}}\allowbreak{} 反射提纯，即可蜕变为真定义，固定公式为：

\NovelDisplay{M00491R01}{%
\text{real.}(2\ 1\text{-})^{\omega}\ 2 = (2\ 1\text{-})(2\ 1\text{-})^{\omega}\ 2 = (2\ 1\text{-})^{\omega+1}\ 2%
}
\end{quote}

卢本韦摩挲掌心温热的猪神令，先祖传承的道韵隐隐共鸣：

\begin{quote}
``原来序道迭代，便是这般去伪存真的过程。先广纳万序立框架，再提纯本源定真意，层层收敛、步步圆满，方能逼近大道核心。''
\end{quote}

众人此刻尽数豁然开朗，连日积攒的疑惑一扫而空。

他们终于彻底明白，为何 \allowbreak{}\NovelInline{I02735}{I(\omega, \alpha)}\allowbreak{} 与此前的 \allowbreak{}\NovelInline{I02736}{I(n, \alpha)}\allowbreak{} 性质截然相反：有限参数的 \allowbreak{}\NovelInline{I02737}{I(n, \alpha)}\allowbreak{} 多为稳固容许序数，而极限参数的 \allowbreak{}\NovelInline{I02738}{I(\omega, \alpha)}\allowbreak{} 自带 \allowbreak{}\NovelInline{I02739}{\omega}\allowbreak{} 长度基本列，尽数非容许。也正因这份非容许特性，psd 与 real 的真伪割裂才会必然诞生。

梅天狸抬手一挥，虚空序纹层层攀升，顺着容许点进位规则继续推演：

\begin{quote}
``吃透真伪之分，我们便可顺延迭代脉络，理清高阶反射式与多元 \allowbreak{}\NovelInline{I02740}{I}\allowbreak{} 函数的对应关系。\allowbreak{}\NovelInline{I02741}{(2\ 1\text{-})^{\omega+1}\ 2}\allowbreak{} 对应 \allowbreak{}\NovelInline{I02742}{I(\omega+1, \cdot)}\allowbreak{} 容许序数，\allowbreak{}\NovelInline{I02743}{(2\ 1\text{-})^{\omega+2}\ 2}\allowbreak{} 对应 \allowbreak{}\NovelInline{I02744}{I(\omega+2, \cdot)}\allowbreak{} 容许序数，以此类推。''

``当迭代上标抵达不动点层级，便诞生 \allowbreak{}\NovelInline{I02745}{(2\ 1\text{-})^{(1,0)}\ 2}\allowbreak{}，对应 \allowbreak{}\NovelInline{I02746}{\alpha = I(\alpha, 0)}\allowbreak{} 的全体不动点，最小元为 \allowbreak{}\NovelInline{I02747}{\psi_I(1,0,0)(0)}\allowbreak{}；再递进至 \allowbreak{}\NovelInline{I02748}{(2\ 1\text{-})^{(1,1)}\ 2}\allowbreak{}，便是 \allowbreak{}\NovelInline{I02749}{I(1, 0, \cdot)}\allowbreak{} 型容许序数\ldots\ldots 直至迭代触及多元 \allowbreak{}\NovelInline{I02750}{I}\allowbreak{} 函数的极限，抵达 \allowbreak{}\NovelInline{I02751}{(2\ 1\text{-})^{(1 @ \omega)}\ 2}\allowbreak{}、\allowbreak{}\NovelInline{I02752}{(2\ 1\text{-})^{(1 @ (1,0))}\ 2}\allowbreak{} 的深层层级。''
\end{quote}

诺当凝神推演片刻，骤然蹙眉：

\begin{quote}
``长老，这套 \allowbreak{}\NovelInline{I02753}{(2\ 1\text{-})}\allowbreak{} 迭代与容许点进位可以无限叠加、无限深化，难道真的没有尽头？我们可以依靠这套规则，无限攀升序道层级？''
\end{quote}

此话一出，殿内序道气息骤然凝重。

\begin{center}\rule{0.5\linewidth}{0.5pt}\end{center}

\chapter{第五十二章：递归马洛序数}\label{ux7b2cux4e94ux5341ux4e8cux7ae0ux9012ux5f52ux9a6cux6d1bux5e8fux6570}

asheep 长老缓缓摇头，清冷声线道出终极桎梏：

\begin{quote}
``无尽迭代，终究是闭环轮回。正如一阶 \allowbreak{}\NovelInline{I02754}{1\text{-}}\allowbreak{} 操作永远无法触及容许序数 \allowbreak{}\NovelInline{I02755}{\Omega}\allowbreak{}，穷尽所有 \allowbreak{}\NovelInline{I02756}{(2\ 1\text{-})}\allowbreak{} 反射、所有多元 \allowbreak{}\NovelInline{I02757}{I}\allowbreak{} 函数进位、所有真伪极限迭代，也永远无法突破一层天堑。''
\end{quote}

虚空之巅，骤然浮现一道漆黑无垠的序界壁垒，横亘在所有 \allowbreak{}\NovelInline{I02758}{I}\allowbreak{} 序数、所有二阶反射式的顶端，封死了一切攀升路径。壁垒之上，一枚极简的符号缓缓亮起，镇压万序、亘古不动。

\begin{quote}
``这便是二阶反射的终极边界，\allowbreak{}\NovelInline{I02759}{\Pi_2}\allowbreak{} 体系的封顶秩序------\textbf{递归马洛序数}。''
\end{quote}

梅天狸望着那枚 \allowbreak{}\NovelInline{I02760}{M}\allowbreak{} 符号，郑重传道：

\begin{quote}
``记其本源：最小的递归马洛序数记为 \allowbreak{}\NovelInline{I02761}{M}\allowbreak{}，是集合 \allowbreak{}\NovelInline{I02762}{2\text{-}2}\allowbreak{} 的最小元。若 \allowbreak{}\NovelInline{I02763}{2\ 1\text{-}2}\allowbreak{} 是容许序数与极限点的交集、诞生递归不可达序数 \allowbreak{}\NovelInline{I02764}{I}\allowbreak{}，那 \allowbreak{}\NovelInline{I02765}{2\text{-}2}\allowbreak{} 便是超脱所有 \allowbreak{}\NovelInline{I02766}{2\ 1\text{-}}\allowbreak{} 体系的至高高阶交集，是 ZFC 公理体系下天然存在、不可递归、不可迭代的本源序数。''
\end{quote}

蔡须坤心神巨震，瞬间厘清三代序数的层级鸿沟：

\begin{quote}
``弟子彻底通透！\allowbreak{}\NovelInline{I02767}{\Omega}\allowbreak{} 是一阶容许根基，\allowbreak{}\NovelInline{I02768}{I}\allowbreak{} 是二阶迭代衍生的高阶容许序数，二者皆可依靠反射规则构造推演；唯独 \allowbreak{}\NovelInline{I02769}{M}\allowbreak{}，跳出了所有二阶反射闭环，是无法通过任何 \allowbreak{}\NovelInline{I02770}{I}\allowbreak{} 函数、任何 \allowbreak{}\NovelInline{I02771}{(2\ 1\text{-})}\allowbreak{} 操作生成的域外至高秩序！''
\end{quote}

``没错。'' asheep 点头补充，顺势展开完整的马洛序数序列：

\begin{quote}
``递归马洛序数并非孤例，有无穷无尽层级，依次记为 \allowbreak{}\NovelInline{I02772}{M}\allowbreak{}、\allowbreak{}\NovelInline{I02773}{M_2}\allowbreak{}、\allowbreak{}\NovelInline{I02774}{M_3 \dots}\allowbreak{}，构成完整集合 \allowbreak{}\NovelInline{I02775}{2\text{-}2 = \{M, M_2, M_3, \dots, M_{\omega+1}, \dots\}}\allowbreak{}。但诸位切记一条铁律：\textbf{所有真正的递归马洛序数，必然是容许序数}。''

``诸如 \allowbreak{}\NovelInline{I02776}{M_\omega}\allowbreak{}、\allowbreak{}\NovelInline{I02777}{M_\Omega}\allowbreak{} 这类极限下标马洛数，自带低级基本列逼近，尽数属于非容许序数，被彻底剔除出马洛真序之列。这份纯粹稳固，是马洛序数碾压所有二阶迭代序数的根本所在。''
\end{quote}

为让众人彻底贯通新旧体系，梅天狸直接展开马洛折叠的核心机理，虚空浮现无数对应等式：

\begin{quote}
``所有你们苦修至今的多元 \allowbreak{}\NovelInline{I02778}{I}\allowbreak{} 序数、所有二阶反射层级，皆是马洛序数的折叠产物。我们专属的马洛折叠函数 \allowbreak{}\NovelInline{I02779}{\psi_M}\allowbreak{}，对应普通 OCF 的 \allowbreak{}\NovelInline{I02780}{\psi}\allowbreak{} 函数，正如多元 \allowbreak{}\NovelInline{I02781}{I}\allowbreak{} 对应多元 \allowbreak{}\NovelInline{I02782}{\varphi}\allowbreak{} 函数一般。''

\NovelDisplay{M00492R01}{%
\psi_M(M^\omega) = I(\omega, 0)%
} \NovelDisplay{M00493R01}{%
\psi_M(M^M) = I(1, 0, 0)%
} \NovelDisplay{M00494R01}{%
\psi_M(M^{M^\omega}) = I(1 @ \omega)%
} \NovelDisplay{M00495R01}{%
\psi_M(M^{M^M}) = I(1 @ (1, 0))%
}

``\ldots\ldots{}''
\end{quote}

一连串精准的序道等式悬浮虚空，字字印证大道真相。

众人望着眼前的对应关系，心神尽数震颤。

他们此前奉为高阶大道的多元 \allowbreak{}\NovelInline{I02783}{I}\allowbreak{} 序数、穷尽无数迭代的反射层级、耗费数日苦修的容许点进位，原来从来不是序道高点，仅仅是马洛本源秩序舒展折叠后的细微一隅。所谓二阶反射的浩瀚体系，终究只是 \allowbreak{}\NovelInline{I02784}{M}\allowbreak{} 序数铺开的疆域之内的衍生规则。

\begin{quote}
``\allowbreak{}\NovelInline{I02785}{\psi_M}\allowbreak{} 并非粗浅的定式公式，而是一套拥有完整闭环逻辑、可逐层枚举、无限推演的活序规则，总计\textbf{四条核心法则}，囊括后继、极限、容许点、兜底四大场景，也是你们三日闭关的核心修行内容。'' asheep 长老眸光肃穆，摒弃此前简略概述，将精准严苛的序道规则尽数显化于虚空，逐条拆解，适配修士推演修行：

``这套规则可完美类比常规 OCF 体系中 \allowbreak{}\NovelInline{I02786}{\psi}\allowbreak{} 与多元 \allowbreak{}\NovelInline{I02787}{\varphi}\allowbreak{} 的对应关系，\allowbreak{}\NovelInline{I02788}{\psi_M}\allowbreak{} 对应 \allowbreak{}\NovelInline{I02789}{\psi}\allowbreak{}、马洛序数 \allowbreak{}\NovelInline{I02790}{M}\allowbreak{} 对应 \allowbreak{}\NovelInline{I02791}{\Omega}\allowbreak{}、多元 \allowbreak{}\NovelInline{I02792}{I}\allowbreak{} 函数对应多元 \allowbreak{}\NovelInline{I02793}{\varphi}\allowbreak{}，整套推演逻辑完全互通，只需吃透类比规律，便可贯通两大序道体系。后续更深层的复杂延展，专栏专项篇章会详细举证，今日我们立足实操，落地每一条规则、每一步枚举。''
\end{quote}

虚空序纹重组，四条工整严谨的 \allowbreak{}\NovelInline{I02794}{\psi_M}\allowbreak{} 核心规则悬浮于众人眼前，字字可证、步步可推，无半分模糊晦涩：

\begin{quote}
\textbf{规则一·初始基值：} \allowbreak{}\NovelInline{I02795}{\psi_M(0) = \Omega}\allowbreak{}，是整套函数的推演根基，所有序阶攀升皆由此始发。

\textbf{规则二·后继递推：} 若 \allowbreak{}\NovelInline{I02796}{\psi_M(\alpha+1) > \psi_M(\alpha)}\allowbreak{}，则 \allowbreak{}\NovelInline{I02797}{\psi_M(\alpha+1)}\allowbreak{} 为 \allowbreak{}\NovelInline{I02798}{\psi_M(\alpha)}\allowbreak{} 的二阶后继，记作 \allowbreak{}\NovelInline{I02799}{2\ \text{aft}\ \psi_M(\alpha)}\allowbreak{}。此规则自带约束，可规避序道迭代卡顿、数值卡死的乱象，杜绝类似常规 OCF 中 \allowbreak{}\NovelInline{I02800}{\psi(\zeta_0+1)}\allowbreak{} 的失效问题。

\textbf{规则三·容许定点：} 凡 \allowbreak{}\NovelInline{I02801}{\psi_M}\allowbreak{} 算式内部以 \allowbreak{}\NovelInline{I02802}{M}\allowbreak{} 收尾，无论内部嵌套加、乘、幂任意运算，\allowbreak{}\NovelInline{I02803}{\psi_M(\#\sim M)}\allowbreak{} 皆为函数 \allowbreak{}\NovelInline{I02804}{f(\alpha) = \psi_M(\#\sim\alpha)}\allowbreak{} 的最小容许点。其''收尾判定''逻辑，与序道体系中 \allowbreak{}\NovelInline{I02805}{\omega}\allowbreak{} 收尾定点规则同源，是生成高阶稳固序数的核心。

\textbf{规则四·极限兜底：} 所有不满足前两条规则的极限情形，统一遵循基本列递推，即 \allowbreak{}\NovelInline{I02806}{\psi_M(\alpha)[n] = \psi_M(\alpha[n])}\allowbreak{}。这条规则适配一切超长竖式、超限极限，哪怕是基本列长度超越 \allowbreak{}\NovelInline{I02807}{\omega}\allowbreak{} 的 \allowbreak{}\NovelInline{I02808}{\psi_M(\Omega)}\allowbreak{}、\allowbreak{}\NovelInline{I02809}{\psi_M(M \times I)}\allowbreak{} 等复杂算式，皆可稳妥推演。
\end{quote}

规则公示完毕，梅天狸抬手引动微光，将虚空序数划分为两大阶段：

\begin{quote}
``\allowbreak{}\NovelInline{I02810}{\psi_M}\allowbreak{} 的推演分为\textbf{两大境界}，OFP 之前为基础 Omega 迭代阶段，OFP 至 IFP 为多元 \allowbreak{}\NovelInline{I02811}{I}\allowbreak{} 高阶迭代阶段，循序渐进，方能不乱序基。''
\end{quote}

八人凝神静气，依序闭目推演，开始亲手落地第一条基础规则。初入推演，众人心中皆存一处共性误区------他们下意识以为，只要是 \allowbreak{}\NovelInline{I02812}{\psi_M}\allowbreak{} 迭代生成的序数，便天然是容许序数。

蔡须坤最先勘透根基，指尖凝出淡金序纹，轻声推演：

\begin{quote}
``依规则一，\allowbreak{}\NovelInline{I02813}{\psi_M(0) = \Omega}\allowbreak{}。再以规则二后继递推，\allowbreak{}\NovelInline{I02814}{\psi_M(1)}\allowbreak{} 便是 \allowbreak{}\NovelInline{I02815}{\Omega}\allowbreak{} 的二阶后继，对应 \allowbreak{}\NovelInline{I02816}{\Omega_2}\allowbreak{}；\allowbreak{}\NovelInline{I02817}{\psi_M(2) = \Omega_3}\allowbreak{}、\allowbreak{}\NovelInline{I02818}{\psi_M(3) = \Omega_4}\allowbreak{}，有限后继序数，尽数逐次递推，秩序规整无误。''
\end{quote}

霍雨浩接续推演极限序数，眸中灵光闪烁：

\begin{quote}
``极限 \allowbreak{}\NovelInline{I02819}{\omega}\allowbreak{} 无后继，适配规则四兜底条款。\allowbreak{}\NovelInline{I02820}{\psi_M(\omega)[n] = \psi_M(\omega[n])}\allowbreak{}，逐层展开便是 \allowbreak{}\NovelInline{I02821}{\Omega_{n+1}}\allowbreak{}，收敛终值为 \allowbreak{}\NovelInline{I02822}{\Omega_\omega}\allowbreak{}。''
\end{quote}

众人顺着二人的推演脉络层层延伸，纷乱的序道符号在掌心有序排布，自成章法：

\NovelDisplay{M00496R01}{%
\psi_M(\omega+1) = \Omega_{\omega+1}, \quad \psi_M(\omega \times 2) = \Omega_{\omega \times 2}, \quad \psi_M(\omega^2) = \Omega_{\omega^2}, \quad \psi_M(\varepsilon_0) = \Omega_{\varepsilon_0}%
}

乃至嵌套自迭代的高阶算式，也尽数被众人勘破：

\NovelDisplay{M00497R01}{%
\psi_M(\psi_M(0)) = \Omega_\Omega, \quad \text{二次嵌套便是 } \Omega_{\Omega_\Omega}%
}

推演至此，小舞忍不住开口道出心中认知：

\begin{quote}
``长老，弟子观 OFP 之前所有 \allowbreak{}\NovelInline{I02823}{\psi_M(\alpha) = \Omega_{1+\alpha}}\allowbreak{}，层级规整、层层递增，无一断裂，想来这些序数尽数都是容许序数吧？''
\end{quote}

此言一出，恰好道破众人心中共通的错悟。八人皆默认：只要序道迭代平稳收敛、无崩坏乱象，便是稳固容许序数。

asheep 长老适时出声，清冷道音点破误区，一语敲碎众人的浅层认知：

\begin{quote}
``\textbf{错。} 容许与否，从来不看迭代是否规整，只看是否存在更低阶基本列逼近、是否为函数最小不动容许点。OFP 之前的 \allowbreak{}\NovelInline{I02824}{\psi_M(\alpha)}\allowbreak{} 虽排布有序、层层递推，但绝大多数只是普通极限序数、后继序数，并非天然容许。你们此刻的误区，也是无数修士卡在马洛 OCF 门外的根本------\textbf{以序貌定序性，而非以规则定本质}。''
\end{quote}

梅天狸顺势补全定论，帮众人夯实根基：

\begin{quote}
``你们只需牢记铁律：在触及 \allowbreak{}\NovelInline{I02825}{M}\allowbreak{}、生成真正的马洛容许点之前，OFP 之前所有 \allowbreak{}\NovelInline{I02826}{\psi_M(\alpha)}\allowbreak{} 统一满足 \allowbreak{}\NovelInline{I02827}{\psi_M(\alpha) = \Omega_{1+\alpha}}\allowbreak{}。这一阶段只铺陈层级、不诞生真容许，唯有突破 OFP、接入 \allowbreak{}\NovelInline{I02828}{M}\allowbreak{} 的折叠规则，序道才会从'普通迭代'跃变为'容许定点迭代'。''
\end{quote}

众人闻言纷纷敛神自省，将这份误区死死刻入道心，彻底摒弃''有序即为容许''的粗浅认知，开始以四条核心规则为唯一准绳，重新审视每一步推演。

``突破 OFP，入局马洛真序。'' 梅天狸话音落下，虚空 \allowbreak{}\NovelInline{I02829}{M}\allowbreak{} 符号骤然亮起，威压席卷整座殿宇，所有基础 Omega 序纹尽数收敛，全新的高阶迭代体系缓缓铺开。

诺当潜心推演片刻，忽然心神巨震，脱口而出：

\begin{quote}
``\allowbreak{}\NovelInline{I02830}{\psi_M(M)}\allowbreak{}！此式不再适配基础递推，依规则三判定，算式以 \allowbreak{}\NovelInline{I02831}{M}\allowbreak{} 收尾，是 \allowbreak{}\NovelInline{I02832}{f(\alpha) = \psi_M(\alpha)}\allowbreak{} 的最小容许点！弟子此前误判，本以为此值是 OFP，如今彻悟------它绝非 OFP，而是实打实的\textbf{递归不可达序数 \allowbreak{}\NovelInline{I02833}{I}\allowbreak{}}！至于 OFP，以我们当下的序阶尚且无法表示，需待后续深修方可触及。''
\end{quote}

这一步突破，彻底打通了 \allowbreak{}\NovelInline{I02834}{\psi_M}\allowbreak{} 与多元 \allowbreak{}\NovelInline{I02835}{I}\allowbreak{} 函数的本源关联，众人此前苦修的 \allowbreak{}\NovelInline{I02836}{I}\allowbreak{} 序数、容许点、二阶迭代，瞬间尽数溯源，与马洛折叠规则完美契合。经历方才纠错，众人此刻推演愈发严谨，步步循规、不敢妄断。

卢本韦推演后继延伸，严格对照规则二：

\begin{quote}
``依后继递推规则，\allowbreak{}\NovelInline{I02837}{\psi_M(M+1) = \Omega_{I+1}}\allowbreak{}、\allowbreak{}\NovelInline{I02838}{\psi_M(M+\omega) = \Omega_{I+\omega}}\allowbreak{}、\allowbreak{}\NovelInline{I02839}{\psi_M(M+\Omega) = \Omega_{I+\Omega}}\allowbreak{}。我已归纳通式，凡 \allowbreak{}\NovelInline{I02840}{M+\alpha}\allowbreak{} 结构，统一遵循 \allowbreak{}\NovelInline{I02841}{\psi_M(M+\alpha) = \Omega_{I+\alpha}}\allowbreak{} 的规律！''
\end{quote}

霍雨浩紧接着推演嵌套迭代，精准区分规则边界，再度纠正自身旧错：

\begin{quote}
``此前我曾混淆规则三与规则四的适用场景！如今看清：\allowbreak{}\NovelInline{I02842}{\psi_M(M+\psi_M(M))}\allowbreak{} 一式，唯有内层 \allowbreak{}\NovelInline{I02843}{\psi_M(M)}\allowbreak{} 可依规则三定点化为 \allowbreak{}\NovelInline{I02844}{I}\allowbreak{}，外层嵌套迭代必须交由规则四兜底展开，最终得 \allowbreak{}\NovelInline{I02845}{\psi_M(M+I) = \Omega_{I \times 2}}\allowbreak{}。同理，\allowbreak{}\NovelInline{I02846}{\psi_M(M+\psi_M(M+1)) = \Omega_{\Omega_{I+1}}}\allowbreak{}，规则边界绝不能混淆！''
\end{quote}

众人接连突破，很快触及全新的序阶壁垒------\allowbreak{}\NovelInline{I02847}{M \times 2}\allowbreak{} 层级，此处恰好又是一处极易误判的容许节点。

蔡须坤凝神沉心，先摒除心中预判，纯以规则推演：

\begin{quote}
``\allowbreak{}\NovelInline{I02848}{\psi_M(M \times 2)}\allowbreak{}，算式以 \allowbreak{}\NovelInline{I02849}{M}\allowbreak{} 收尾，适配规则三，为函数 \allowbreak{}\NovelInline{I02850}{f(\alpha) = \psi_M(M+\alpha) = \Omega_{I+\alpha}}\allowbreak{} 的最小容许点。我此前险些再入误区，误以为 \allowbreak{}\NovelInline{I02851}{I}\allowbreak{} 可自我迭代收敛、成为此函数不动点，实则细细验证便知：\allowbreak{}\NovelInline{I02852}{I}\allowbreak{} 根本不是 \allowbreak{}\NovelInline{I02853}{f(\alpha) = \Omega_{I+\alpha}}\allowbreak{} 的不动点，无法自洽收敛！因此全新高阶容许序数诞生，\allowbreak{}\NovelInline{I02854}{\psi_M(M \times 2) = I_2}\allowbreak{}！''
\end{quote}

双长老微微颔首，赞许众人知错即改、循规推演的精进状态。

以此为基，众人顺势铺开整套倍数迭代体系，全程严守规则、无半分臆断：

\NovelDisplay{M00498R01}{%
\psi_M(M \times 2 + 1) = \Omega_{I_2 + 1}%
}

通式可归纳为：

\NovelDisplay{M00499R01}{%
\psi_M(M \times 2 + \alpha) = \Omega_{I_2 + \alpha}%
}

顺延迭代：

\NovelDisplay{M00500R01}{%
\psi_M(M \times 3) = I_3%
}

随后众人递归延伸所有超限倍数情形：

\NovelDisplay{M00501R01}{%
\psi_M(M \times \omega) = I_\omega, \quad \psi_M(M \times \Omega) = I_\Omega, \quad \psi_M(M \times \psi_M(M)) = I_I%
}

至此众人彻底总结出 IFP 之前的\textbf{绝对铁律}：在 IFP 壁垒之前，所有 \allowbreak{}\NovelInline{I02855}{\psi_M(M \times \alpha)}\allowbreak{} 恒等于 \allowbreak{}\NovelInline{I02856}{I_\alpha}\allowbreak{}。

序道推演继续攀升，踏入幂次层级，多元 \allowbreak{}\NovelInline{I02857}{I}\allowbreak{} 函数的完整形态终于尽数现世，每一步推演都对应清晰规则，众人的容许判定认知已然彻底成熟。

\allowbreak{}\NovelInline{I02858}{\psi_M(M^2)}\allowbreak{} 为 \allowbreak{}\NovelInline{I02859}{M \times \alpha}\allowbreak{} 迭代体系的最小容许点，对应双元 \allowbreak{}\NovelInline{I02860}{I}\allowbreak{} 函数根基 \allowbreak{}\NovelInline{I02861}{I(1, 0)}\allowbreak{}。

后续幂次延伸有序铺开：

\NovelDisplay{M00502R01}{%
\psi_M(M^2 + 1) = \Omega_{I(1,0)+1}%
} \NovelDisplay{M00503R01}{%
\psi_M(M^2 + M) = I_{I(1,0)+1}%
} \NovelDisplay{M00504R01}{%
\psi_M(M^2 \times 2) = I(1, 1)%
} \NovelDisplay{M00505R01}{%
\psi_M(M^2 \times 3) = I(1, 2)%
} \NovelDisplay{M00506R01}{%
\psi_M(M^3) = I(2, 0)%
} \NovelDisplay{M00507R01}{%
\psi_M(M^4) = I(3, 0)%
}

当幂次抵达超限 \allowbreak{}\NovelInline{I02862}{\omega}\allowbreak{}：

\NovelDisplay{M00508R01}{%
\psi_M(M^\omega) = I(\omega, 0)%
}

小美顺势归纳出通用大道，再无此前的认知疏漏：

\begin{quote}
``如今我已彻底明晰！只要满足 \allowbreak{}\NovelInline{I02863}{\omega \leq \alpha < I(a, 0)\text{ FP}}\allowbreak{}，必然有 \allowbreak{}\NovelInline{I02864}{\psi_M(M^\alpha) = I(\alpha, 0)}\allowbreak{}，这是幂次迭代的通用铁律，无一例外。''
\end{quote}

推演抵达 \allowbreak{}\NovelInline{I02865}{M^M}\allowbreak{} 层级，三阶 \allowbreak{}\NovelInline{I02866}{I}\allowbreak{} 函数轰然现世，彻底印证马洛折叠的至高包容性，众人已然能自主区分定点容许与普通极限的差异。

\allowbreak{}\NovelInline{I02867}{\psi_M(M^M)}\allowbreak{} 为 \allowbreak{}\NovelInline{I02868}{I(\alpha, 0)}\allowbreak{} 体系的不动容许点，是纯正稳固的马洛真序，生成三元基底 \allowbreak{}\NovelInline{I02869}{I(1, 0, 0)}\allowbreak{}。

众人逐层落地高阶嵌套迭代，每一式皆精准无误：

\NovelDisplay{M00509R01}{%
\psi_M(M^M + 1) = \Omega_{I(1,0,0)+1}%
} \NovelDisplay{M00510R01}{%
\psi_M(M^M + M) = I_{I(1,0,0)+1}%
} \NovelDisplay{M00511R01}{%
\psi_M(M^M + M^2) = I(1, I(1,0,0)+1)%
} \NovelDisplay{M00512R01}{%
\psi_M(M^M + M^\omega) = I(\omega, I(1,0,0)+1)%
} \NovelDisplay{M00513R01}{%
\psi_M(M^M + M^{\psi_M(M^M)}) = I(I(1,0,0), 1)%
} \NovelDisplay{M00514R01}{%
\psi_M(M^M + M^{\psi_M(M^M)+1}) = I(I(1,0,0)+1, 0)%
} \NovelDisplay{M00515R01}{%
\psi_M(M^M \times 2) = I(1, 0, 1)%
}

卢本韦推演至此，彻底贯通类比大道，感慨出声：

\begin{quote}
``弟子终于完全明白！此体系与常规 OCF 完全同源，只需将 \allowbreak{}\NovelInline{I02870}{M}\allowbreak{} 替换为 \allowbreak{}\NovelInline{I02871}{\Omega}\allowbreak{}、\allowbreak{}\NovelInline{I02872}{I}\allowbreak{} 替换为 \allowbreak{}\NovelInline{I02873}{\varphi}\allowbreak{}，所有推演规则、迭代逻辑、序数对应全然一致。此前的误区，皆是不肯循规则实证、仅凭表象臆断所致！''
\end{quote}

众人加快推演速度，踏入超高阶嵌套幂次层级，四元、五元乃至无限维 \allowbreak{}\NovelInline{I02874}{I}\allowbreak{} 函数接连诞生：

\NovelDisplay{M00516R01}{%
\psi_M(M^{M+1}) = I(1, 1, 0)%
} \NovelDisplay{M00517R01}{%
\psi_M(M^{M \times 2}) = I(2, 0, 0)%
} \NovelDisplay{M00518R01}{%
\psi_M(M^{M^2}) = I(1, 0, 0, 0)%
} \NovelDisplay{M00519R01}{%
\psi_M(M^{M^3}) = I(1, 0, 0, 0, 0)%
} \NovelDisplay{M00520R01}{%
\psi_M(M^{M^4 + M^3 + M^2 \times 4 + M \times 5 + 1} \times 5) = I(1, 1, 4, 5, 1, 4)%
}

超限塔极限随之现世：

\NovelDisplay{M00521R01}{%
\psi_M(M^{M^\omega}) = I(1 @ \omega)%
} \NovelDisplay{M00522R01}{%
\psi_M(M^{M^M}) = I(1 @ (1, 0))%
}

最终，纯 \allowbreak{}\NovelInline{I02875}{\psi_M}\allowbreak{} 与单 \allowbreak{}\NovelInline{I02876}{M}\allowbreak{} 序数的推演抵达终极壁垒：

\NovelDisplay{M00523R01}{%
\psi_M(\varepsilon_{M+1}) = \psi_M(\psi_{\Omega_{M+1}}(0))%
}

这是单马洛序数迭代所能触及的绝对巅峰，再无延展空间。

``单 \allowbreak{}\NovelInline{I02877}{M}\allowbreak{} 序数的推演，仅为根基。'' asheep 长老见状，再展更高维度的序道体系：

\begin{quote}
``真正的高阶马洛序道，还有带下标的 \allowbreak{}\NovelInline{I02878}{M_\nu}\allowbreak{} 拓展体系，衍生全新的 \allowbreak{}\NovelInline{I02879}{\psi_{M_\nu}}\allowbreak{} 函数，规则框架与基础 \allowbreak{}\NovelInline{I02880}{\psi_M}\allowbreak{} 同源，仅初始基值迭代升级。''
\end{quote}

虚空再显拓展规则，逐条公示，与基础规则一一对应：

\begin{quote}
\textbf{规则一·初始基值：} \allowbreak{}\NovelInline{I02881}{\psi_{M_\nu}(0) = \Omega_{M_\nu + 1}}\allowbreak{}

\textbf{规则二·后继递推：} 当 \allowbreak{}\NovelInline{I02882}{\psi_{M_\nu}(\alpha+1) > \psi_{M_\nu}(\alpha)}\allowbreak{} 时，\allowbreak{}\NovelInline{I02883}{\psi_{M_\nu}(\alpha+1) = 2\ \text{aft}\ \psi_{M_\nu}(\alpha)}\allowbreak{}

\textbf{规则三·容许定点：} \allowbreak{}\NovelInline{I02884}{\psi_{M_\nu}}\allowbreak{} 算式内部以 \allowbreak{}\NovelInline{I02885}{M_\nu}\allowbreak{} 收尾时，为 \allowbreak{}\NovelInline{I02886}{f(\alpha) = \psi_{M_\nu}(\# \sim \alpha)}\allowbreak{} 的最小容许点

\textbf{规则四·极限兜底：} 其余所有极限情形，遵循 \allowbreak{}\NovelInline{I02887}{\psi_{M_\nu}(\alpha)[n] = \psi_{M_\nu}(\alpha[n])}\allowbreak{} 递推
\end{quote}

众人顺势推演下标马洛体系的全新序数，依旧循规而行，彻底杜绝臆断误区：

\NovelDisplay{M00524R01}{%
\psi_{M_2}(0) = \Omega_{M+1}, \quad \psi_{M_2}(M) = \Omega_{M \times 2}%
}

此二式完全依托规则四兜底展开。

继而深化高阶迭代：

\NovelDisplay{M00525R01}{%
\psi_{M_2}(M_2) = I_{M+1}%
} \NovelDisplay{M00526R01}{%
\psi_{M_2}(M_2^2) = I(1, M+1)%
} \NovelDisplay{M00527R01}{%
\psi_{M_2}(M_2^\omega) = I(\omega, M+1)%
} \NovelDisplay{M00528R01}{%
\psi_{M_2}(M_2^{M_2}) = I(1, 0, M+1)%
} \NovelDisplay{M00529R01}{%
\psi_{M_2}(M_2^{M_2^{M_2}}) = I(1 @ (1,0), M+1 @ 0)%
}

更高阶基底顺势诞生：

\NovelDisplay{M00530R01}{%
\psi_{M_3}(0) = \Omega_{M_2 + 1}%
}

层层下标、层层拔高，序道维度无限拓宽。

全程推演落幕，八人心中的错悟尽数拔除。他们彻底铭记修行关键：序道真伪、容许与否，从来不靠观感判定，唯以四条核心规则为唯一准绳。OFP 前无真容许、\allowbreak{}\NovelInline{I02888}{M}\allowbreak{} 迭代方生真定点、极限下标多为非容许，这些曾经晦涩的定理，如今尽数化作切身推演的道果。

梅天狸收尾总结，道尽二阶圆满的真谛：

\begin{quote}
``在 \allowbreak{}\NovelInline{I02889}{M}\allowbreak{} 出现之前，\allowbreak{}\NovelInline{I02890}{\psi_M(\alpha) = \Omega_{1+\alpha}}\allowbreak{}，尽数对应容许序数下标；抵达 \allowbreak{}\NovelInline{I02891}{M}\allowbreak{} 层级后，便开始折叠出全套多元 \allowbreak{}\NovelInline{I02892}{I}\allowbreak{} 体系，层层递进、无限深化，直至抵达纯 \allowbreak{}\NovelInline{I02893}{\psi_M}\allowbreak{} 可推演的极限 \allowbreak{}\NovelInline{I02894}{\varepsilon_{M+1}}\allowbreak{}。''

``psd 与 real 的真伪之辨，让你们勘破极限迭代的虚实；容许点进位，让你们穷尽二阶迭代的极致；而递归马洛序数，让你们窥见 \allowbreak{}\NovelInline{I02895}{\Pi_2}\allowbreak{} 体系的真正边界。''
\end{quote}

\begin{quote}
``\allowbreak{}\NovelInline{I02896}{\psi_M}\allowbreak{} 负责折叠生成容许序数，而完整的马洛 OCF 大道，是将这些容许序数再次折叠、退层、嵌套，最终生成纯粹的递归序数。'' 梅天狸适时补全终极闭环，点出序道终极奥义：

``为规避序阶混乱、虚实倒置的乱象，完整体系会以 \allowbreak{}\NovelInline{I02897}{M}\allowbreak{} 系列马洛序数为典范基底，而非 \allowbreak{}\NovelInline{I02898}{\Omega}\allowbreak{}、\allowbreak{}\NovelInline{I02899}{I}\allowbreak{}，这套退层展开、嵌套计数的深层规则，便是你们后续需要深耕的终极内容。''
\end{quote}

众人此刻已然彻底通透，全程亲手推演、知错纠错、逐层验证，将四条核心规则、两大迭代阶段、单 \allowbreak{}\NovelInline{I02900}{M}\allowbreak{} 与下标 \allowbreak{}\NovelInline{I02901}{M}\allowbreak{} 体系的所有逻辑，尽数烙印道基。

他们终于真切明白，psd 与 real 的真伪之辨是序道的虚实根基，容许点进位是迭代攀升的阶梯，而 \allowbreak{}\NovelInline{I02902}{\psi_M}\allowbreak{} 马洛折叠，便是统合所有二阶序道、衍生万千 \allowbreak{}\NovelInline{I02903}{I}\allowbreak{} 序数、锚定 \allowbreak{}\NovelInline{I02904}{\Pi_2}\allowbreak{} 体系边界的至高大道。

\chapter{第五十三章：折叠真意，不动点之殇}\label{ux7b2cux4e94ux5341ux4e09ux7ae0ux6298ux53e0ux771fux610fux4e0dux52a8ux70b9ux4e4bux6b87}

扽魂殿内，随着梅天狸长老将 \allowbreak{}\NovelInline{I02905}{\psi_M}\allowbreak{} 的四条核心法则与两大迭代阶段全盘托出，原本凝滞的空气仿佛被某种宏大的秩序重新梳理。八名少年少女并未急于睁眼，而是沉浸在方才那场惊心动魄的推演余韵中。

卢本韦掌心的猪神令微微震颤，他深吸一口气，打破了沉默：``诸位，双长老所传之法，虽名为`马洛折叠'，实则是一场关于`不动点'的极致操弄。此前我们只知 \allowbreak{}\NovelInline{I02906}{I}\allowbreak{} 是容许点，却不知在 \allowbreak{}\NovelInline{I02907}{\psi_M}\allowbreak{} 的视界下，万物皆可被折叠为 \allowbreak{}\NovelInline{I02908}{M}\allowbreak{} 的函数。''

``正是。''霍雨浩接过话头，眸中灵光流转，仿佛看穿了虚空中那些繁复符号的本质，``规则三所言`以 \allowbreak{}\NovelInline{I02909}{M}\allowbreak{} 收尾即为最小容许点'，这不仅是生成的规则，更是\textbf{退层}的钥匙。当我们面对一个以 \allowbreak{}\NovelInline{I02910}{M}\allowbreak{} 结尾的庞大算式时，它不再是一个静止的数值，而是一个正在坍缩、正在寻找自身不动点的动态过程。''

asheep 长老微微颔首，清冷的声音再次响起，带着不容置疑的威严：``既然你们已窥见门径，便由我引出 \allowbreak{}\NovelInline{I02911}{\psi_M}\allowbreak{} 体系中最晦涩、也最核心的\textbf{退层展开机制}。此乃连接`表象序数'与`本源马洛'的唯一桥梁。''

虚空中的序纹骤然变幻，不再是简单的等式对应，而是化作了一条条深邃的逻辑链条。

``你们需铭记，\allowbreak{}\NovelInline{I02912}{\psi_M}\allowbreak{} 的本质，是将高阶的马洛序数 \allowbreak{}\NovelInline{I02913}{M}\allowbreak{} `翻译'回低阶的 \allowbreak{}\NovelInline{I02914}{\Omega}\allowbreak{} 语言。'' asheep 长老抬手，虚空中浮现出第一条基础公理的变体：

\begin{quote}
\textbf{规则一（基值锚定）：} \NovelDisplay{M00531R01}{%
\psi_M(0) = \varepsilon_0%
} ``这是万法之始。但在进入高阶之前，你们必须理解\textbf{后继规则}的严苛性。''
\end{quote}

asheep 指尖轻点，一道流光划过：

\begin{quote}
\textbf{规则二（后继跃迁）：} 当 \allowbreak{}\NovelInline{I02915}{\psi(\alpha+1) > \psi(\alpha)}\allowbreak{} 时，有： \NovelDisplay{M00532R01}{%
\psi(\alpha+1) = (1\text{-})^{1,0} \text{ aft } \psi(\alpha)%
} ``注意，这里的 \allowbreak{}\NovelInline{I02916}{(1\text{-})^{1,0}}\allowbreak{} 并非简单的加法，而是二阶反射操作符的具象化。它意味着每一次后继，都是在当前序数之上，强行叠加了一层`不可达'的性质。''
\end{quote}

蔡须坤眉头微蹙，显然在消化这其中的深意：``长老，若遇内部结构极其复杂，尤其是嵌套了多层 \allowbreak{}\NovelInline{I02917}{\psi_M}\allowbreak{} 的情况，该如何剥离？''

``问到了关键。''梅天狸长老适时介入，她的声音柔和却直击要害，``这便是\textbf{规则三}与\textbf{规则四}的分野------\textbf{不动点的判定与展开}。''

她挥手间，虚空中显现出一段极为复杂的嵌套结构，正是困扰无数修士的``黑暗森林''：

\begin{quote}
\textbf{规则三（\allowbreak{}\NovelInline{I02918}{M}\allowbreak{} 收尾的退层）：} 当 \allowbreak{}\NovelInline{I02919}{\psi}\allowbreak{} 函数的内部以 \allowbreak{}\NovelInline{I02920}{M}\allowbreak{} 收尾时，需从内到外找到最近的、包含这个 \allowbreak{}\NovelInline{I02921}{M}\allowbreak{} 的最内层下标不超过 \allowbreak{}\NovelInline{I02922}{M}\allowbreak{} 的函数，记该函数为 \allowbreak{}\NovelInline{I02923}{\psi_t}\allowbreak{}。 此时，整个表达式可表示为 \allowbreak{}\NovelInline{I02924}{\psi(h_0(\psi_t(h(M))))}\allowbreak{}，其中 \allowbreak{}\NovelInline{I02925}{h_0}\allowbreak{} 和 \allowbreak{}\NovelInline{I02926}{h}\allowbreak{} 为递归函数。 \textbf{核心操作：} 将表达式``退层''为 \allowbreak{}\NovelInline{I02927}{\psi(h_0(\psi_t(h(\psi_M(h(M))))))}\allowbreak{}，然后将 \allowbreak{}\NovelInline{I02928}{\psi_M(h(M))}\allowbreak{} 作为整体放入下标，展开为无穷嵌套： \NovelDisplay{M00533R01}{%
\psi(h_0(\psi_t(h(\psi_{\psi_M(h(M))}(h(\psi_{\psi_M(h(M))}(h(\dots))))))))%
} ``此处，''梅天狸加重语气，``\allowbreak{}\NovelInline{I02929}{\psi_{\psi_M(h(M))}}\allowbreak{} 即为函数计数 \allowbreak{}\NovelInline{I02930}{\alpha \to \psi_M(h(\alpha))}\allowbreak{} 的\textbf{不动点}。你们要将这个 \allowbreak{}\NovelInline{I02931}{M}\allowbreak{} 看作是一个不断自我引用的深渊，每一次展开，都是向深渊更深处的一次凝视。''
\end{quote}

小舞听得入神，喃喃自语：``以自身为下标，无限递归\ldots\ldots 这便是马洛序数`大基数'特性的体现吗？''

``不错。''asheep 长老冷冷补充，``而当遇到后继序数收尾时，逻辑则截然不同。''

\begin{quote}
\textbf{规则四（后继序数的处理）：} 当 \allowbreak{}\NovelInline{I02932}{\psi}\allowbreak{} 函数内部以后继序数收尾，且外层为 \allowbreak{}\NovelInline{I02933}{\psi_M}\allowbreak{} 时，将这个 \allowbreak{}\NovelInline{I02934}{\psi_M}\allowbreak{} 及其内部当做普通的 \allowbreak{}\NovelInline{I02935}{\Omega}\allowbreak{} 序数处理。 即： \NovelDisplay{M00534R01}{%
\psi(h_0(\psi_t(h(\psi_M(\alpha+1))))) = \psi(h_0(\psi_t(h(\psi_{\psi_M(\alpha+1)}(h(\psi_{\psi_M(\alpha+1)}(\dots)))))))%
} ``这里 \allowbreak{}\NovelInline{I02936}{\psi_t}\allowbreak{} 是最近的包含 \allowbreak{}\NovelInline{I02937}{\psi_M(\alpha+1)}\allowbreak{} 的函数。切记，此时它不再享有 \allowbreak{}\NovelInline{I02938}{M}\allowbreak{} 的`特权'，而是退化为普通的 \allowbreak{}\NovelInline{I02939}{\Omega}\allowbreak{} 序列进行展开。''
\end{quote}

``理论终归抽象，唯有实战方能证道。''梅天狸长老目光如炬，``现在，我将给出三个层级的实例，你们需在脑海中同步推演。''

``首先是 \allowbreak{}\NovelInline{I02940}{\psi(\psi_M(\omega) + \psi_{\psi_M(1)}(\psi_M(3)))}\allowbreak{} 的展开。''

霍雨浩反应极快，脱口而出：``此式以后继序数 \allowbreak{}\NovelInline{I02941}{(3)}\allowbreak{} 收尾，且外层是 \allowbreak{}\NovelInline{I02942}{\psi_M}\allowbreak{}。依\textbf{规则四}，应将 \allowbreak{}\NovelInline{I02943}{\psi_M(3)}\allowbreak{} 视作整体，向外寻找最近的 \allowbreak{}\NovelInline{I02944}{\psi_t}\allowbreak{}，即 \allowbreak{}\NovelInline{I02945}{\psi_{\psi_M(1)}}\allowbreak{}。随后将其退化为 \allowbreak{}\NovelInline{I02946}{\Omega}\allowbreak{} 序列！''

``正解。''梅天狸赞许点头，虚空显现出最终形态： \NovelDisplay{M00535R01}{%
\psi(\psi_M(\omega) + \psi_{\psi_M(3)}(\psi_{\psi_M(3)}(\psi_{\psi_M(3)}(\dots))))%
} ``看似繁复，实则不过是 \allowbreak{}\NovelInline{I02947}{\psi(\Omega_\omega + \psi_1(\Omega_4))}\allowbreak{} 罢了。这就是\textbf{规则四}的威力------将高阶伪装瞬间拆解为低阶的累加。''

``接下来，\allowbreak{}\NovelInline{I02948}{\psi(M^2)}\allowbreak{}。''asheep 长老的声音如同寒冰撞击，``注意，这是乘法结构，需运用\textbf{规则三}。''

卢本韦沉声道：``\allowbreak{}\NovelInline{I02949}{\psi(M^2) = \psi(M \times M)}\allowbreak{}。首先将其退层为 \allowbreak{}\NovelInline{I02950}{\psi(M \times \psi_M(M^2))}\allowbreak{}。此时，\allowbreak{}\NovelInline{I02951}{\psi_M(M^2)}\allowbreak{} 成为新的核心，作为整体进行展开\ldots\ldots''

随着他的推演，虚空中浮现出令人眩晕的嵌套： \NovelDisplay{M00536R01}{%
\psi(M \times \psi_{\psi_M(M^2)}(M \times \psi_{\psi_M(M^2)}(M \times \dots)))%
} ``Quiet easily done.'' 卢本韦收势，嘴角勾起一抹自信的笑意，``只要抓住了`以 \allowbreak{}\NovelInline{I02952}{M}\allowbreak{} 收尾'这一牛鼻子，乘积不过是加法的迭代延伸。''

殿内的气氛陡然凝重。asheep 长老抛出了最后的试炼：``\allowbreak{}\NovelInline{I02953}{\psi(\psi_{\psi_M(M)}(1))}\allowbreak{}。''

``这是\ldots\ldots{}\textbf{规则六}的领域！''蔡须坤瞳孔微缩，``最为黑暗复杂的一条规则。''

他闭上眼，神识疯狂运转：``此式以 \allowbreak{}\NovelInline{I02954}{\psi_{\psi_M(M)}(1)}\allowbreak{} 收尾。根据\textbf{规则三}的定义，\allowbreak{}\NovelInline{I02955}{\psi_{\psi_M(M)}(1)}\allowbreak{} 本身就是函数 \allowbreak{}\NovelInline{I02956}{\alpha \to \psi_M(\alpha)}\allowbreak{} 的\textbf{次小不动点}（位于 \allowbreak{}\NovelInline{I02957}{\psi_{\psi_M(M)}(0)}\allowbreak{} 之后）。''

``我们要做的，不是简单的替换，而是\textbf{重构}。''蔡须坤大喝一声，指尖在虚空中划出一道完美的闭环： \NovelDisplay{M00537R01}{%
\psi(\psi_{\psi_M(M)}(0) + \psi_M(\psi_{\psi_M(M)}(0) + \psi_M(\psi_{\psi_M(M)}(0) + \dots)))%
} ``更一般地，对于位于 \allowbreak{}\NovelInline{I02958}{\gamma}\allowbreak{} 之后的最小不动点，我们将其展开为 \allowbreak{}\NovelInline{I02959}{\gamma + \psi_M(h(\gamma + \psi_M(h(\dots))))}\allowbreak{} 的形式。这不再是线性的增长，而是指数级的爆发！''

随着最后一个实例的解析完毕，扽魂殿内陷入了长久的死寂。八人皆被那无穷无尽的递归之美所震撼。

梅天狸长老望着众人，缓缓道出了最终的结论：``通过引入 \allowbreak{}\NovelInline{I02960}{M}\allowbreak{} 的行为，我们能得到的极限表达式是 \allowbreak{}\NovelInline{I02961}{\psi(M^{M^{M^{\cdot^{\cdot}}}})}\allowbreak{}，这在序数分析中被称为 \textbf{SRO}。这一序数在 BMS 中被表示为 \NovelDisplay{I02962}{(0,0,0)(1,1,1)(2,1,1)(3,1,1)(3,1,0)(4,2,0)}。''

她顿了顿，目光深邃：``接下来，我们可以进一步引入带下标的 \allowbreak{}\NovelInline{I02963}{M_2, \dots, M_\omega, \dots}\allowbreak{} 等序数，进一步增强 Mahlo OCF。不过，引入了这些序数以后，OCF 的规则会进一步复杂，本文暂不做介绍了。''

``但你们需知晓，''asheep 长老补充道，声音仿佛来自远古，``今日所学，不过是 \allowbreak{}\NovelInline{I02964}{\Pi_2}\allowbreak{} 体系冰山一角。真正的马洛序数，是通往更高维度的阶梯。当你以为触碰到天花板时，那仅仅是下一层地板的开始。''

众人齐齐躬身，声音汇聚成一股洪流，震荡着整座大殿： ``弟子谨遵教诲！真伪已辨，折叠已成，愿以此基，攀登无尽序道！''

殿门缓缓开启，外界的晨光洒入，与殿内未散的序道微光交织在一起。八名少年少女踏出大殿，他们的步伐依旧轻盈，但每一步落下，仿佛都踩在了宇宙最深层的律动之上。

晨光熹微，扽魂殿外的广场上，八名少年少女并未散去。方才殿内的推演虽已落幕，但那种思维被强行拉伸至极限的酸胀感依旧残存。

卢本韦望着天边初升的旭日，手中的猪神令此刻已不再震颤，而是散发出一种沉稳的幽光。他转头看向霍雨浩，眼中满是未尽的思索：``雨浩，方才长老所言 SRO 之境，虽已是凡人思维的极限，但我总觉得，那仅仅是`一元'马洛函数的终点。若 \allowbreak{}\NovelInline{I02965}{M}\allowbreak{} 不再是孤独的行者，而是有了同伴，有了层级，那序数的天空又会是何等景象？''

霍雨浩微微颔首，目光穿透云层，仿佛看到了更高维度的架构：``师兄所言极是。\allowbreak{}\NovelInline{I02966}{M}\allowbreak{} 代表的是`大'，是`不可达'。但当这种`不可达'本身也需要被度量、被分类时，我们就需要引入\textbf{多元 \allowbreak{}\NovelInline{I02967}{M}\allowbreak{} 函数}。这不仅仅是数量的增加，更是性质的质变。''

就在这时，asheep 长老的身影再次浮现于广场上空。这一次，他的周身不再环绕着单一的黑色序纹，而是交织着红蓝双色的复杂光带，宛如两条相互缠绕的巨龙。

``看来，你们已不满足于 \allowbreak{}\NovelInline{I02968}{\Pi_2}\allowbreak{} 体系的温床。'' asheep 的声音带着一种金属般的质感，``既如此，便随我踏入\textbf{多元 \allowbreak{}\NovelInline{I02969}{M}\allowbreak{} 函数}与\textbf{带分号进制}的领域。在这里，`容许点'将不再是静止的锚点，而是流动的江河。''

``在上一课中，我们枚举了递归马洛序数 \allowbreak{}\NovelInline{I02970}{2-2}\allowbreak{}。'' asheep 长老抬手，虚空中浮现出一串金色的符号序列： \NovelDisplay{M00538R01}{%
2-2 = \{M, M_2, M_3, \dots, M_{\omega+1}, \dots\}%
} ``这是基础。但现在，我要你们对这套集合进行\textbf{\allowbreak{}\NovelInline{I02971}{1-}\allowbreak{} 反射}。想象一面镜子，照出的不是影像，而是性质的升华。''

随着长老的话语，那串金色符号骤然分裂、重组，化作了一片银色的海洋： \NovelDisplay{M00539R01}{%
1-2-2 = \{M_\omega, M_{\omega \times 2}, \dots, M_{\omega^2}, \dots, M_M, \dots\}%
}

蔡须坤惊呼出声：``这是\ldots\ldots 全体以极限序数为下标的 \allowbreak{}\NovelInline{I02972}{M_\alpha}\allowbreak{}！原本离散的 \allowbreak{}\NovelInline{I02973}{M_n}\allowbreak{} 被压缩成了连续的流！''

``正是。'' asheep 点头，``继续反射，便是 \allowbreak{}\NovelInline{I02974}{1-1-2-2}\allowbreak{}，那是 \allowbreak{}\NovelInline{I02975}{\Omega}\allowbreak{} 级别的跨越。而当我们进行 \allowbreak{}\NovelInline{I02976}{(1-)^{(2)} 2-2}\allowbreak{} 操作时，我们触及了 \allowbreak{}\NovelInline{I02977}{I}\allowbreak{}（不可达基数）的领域。''

梅天狸长老此时也飘然而至，她补充道：``但这还不够。我们要寻找的是函数 \allowbreak{}\NovelInline{I02978}{f(\alpha) = M_\alpha}\allowbreak{} 的不动点。根据定义，满足 \allowbreak{}\NovelInline{I02979}{\alpha = \min (1-)^\alpha 2-2}\allowbreak{} 的集合，即为 \textbf{MFP}。相信顺着前几篇文章看到这里的读者不会再有什么疑问了。''

她目光扫过众人，继续说道：``那么，函数 \allowbreak{}\NovelInline{I02980}{f(\alpha) = M_\alpha}\allowbreak{} 的容许点又如何呢？根据前文所讲的容许点的性质，可以知道，这个函数的全体容许点构成了集合 \allowbreak{}\NovelInline{I02981}{2 \ 1-2-2}\allowbreak{}。我们设定多元 \allowbreak{}\NovelInline{I02982}{M}\allowbreak{} 函数仍然为容许点进制，将最小的 \allowbreak{}\NovelInline{I02983}{2 \ 1-2-2}\allowbreak{} 记作 \allowbreak{}\NovelInline{I02984}{M(1,0)}\allowbreak{}，其余的容许点按大小以此类推。这样一来，就有：'' \NovelDisplay{M00540R01}{%
2 \ 1-2-2 = \{M(1,0), M(1,1), \dots, M(1,\omega+1), \dots\}%
}

``而上文的 MFP 也可以用 \allowbreak{}\NovelInline{I02985}{\psi_{M(1,0)}}\allowbreak{} 函数表示了：'' \NovelDisplay{M00541R01}{%
(1-)^{1,0} 2-2 = \{\psi_{M(1,0)}(0), \psi_{M(1,0)}(1), \dots, M(1,0), \dots\}%
}

``接着，我们可以效仿多元 \allowbreak{}\NovelInline{I02986}{I}\allowbreak{}，给出多元 \allowbreak{}\NovelInline{I02987}{M}\allowbreak{} 函数的反射式。'' asheep 长老的声音愈发洪亮，``因为我们的多元 \allowbreak{}\NovelInline{I02988}{M}\allowbreak{} 仍使用容许点进制，所以其内核与多元 \allowbreak{}\NovelInline{I02989}{I}\allowbreak{} 是完全一致的。''

虚空中，复杂的公式如瀑布般倾泻而下： \NovelDisplay{M00542R01}{%
2 \ 1-2 \ 1-2-2 = \{M(2,0), M(2,1), \dots, M(2,\omega+1), \dots, M(3,0), \dots\}%
} \NovelDisplay{M00543R01}{%
(2 \ 1-)^{\omega} 2-2 = \{M(\omega,0), M(\omega,1), \dots, M(\omega,\omega), \dots, M(\omega+1,0), \dots\}%
}

``这里的 \allowbreak{}\NovelInline{I02990}{(2 \ 1-)^{\omega} 2-2}\allowbreak{} 也是 psd. 定义的。'' asheep 特意强调了这一点。

紧接着，更深层的嵌套结构显现： \NovelDisplay{M00544R01}{%
(2 \ 1-)^{1,0} 1-2-2 = \{\psi_{M(1,0,0)}(0), \psi_{M(1,0,0)}(1), \dots, \psi_{M(1,0,0)}(\omega), \dots, M(1,0,0), \dots\}%
} \NovelDisplay{M00545R01}{%
(2 \ 1-)^{1,1} 1-2-2 = \{M(1,0,0), M(1,0,1), \dots, M(1,0,\omega+1), \dots, M(1,1,0), \dots\}%
} \NovelDisplay{M00546R01}{%
(2 \ 1-)^{1,0,0} 1-2-2 = \{\psi_{M(1,0,0,0)}(0), \psi_{M(1,0,0,0)}(1), \dots, \psi_{M(1,0,0,0)}(\omega), \dots, M(1,0,0,0), \dots\}%
}

霍雨浩看着这些繁复的符号，喃喃自语：``\allowbreak{}\NovelInline{I02991}{(2 \ 1-)^{1@ \omega} 1-2-2}\allowbreak{}\ldots\ldots 这是 \NovelDisplay{I02992}{\{M(1@\omega), M(1@\omega, 1@0), \dots\}} 的集合。这种 \allowbreak{}\NovelInline{I02993}{1@ \omega}\allowbreak{} 的结构，简直是将维度的折叠玩弄到了极致。''

``接下来，会有一个序数能折叠 \allowbreak{}\NovelInline{I02994}{(2 \ 1-)^{a,b,c,\dots} 2-2}\allowbreak{}。'' asheep 长老抛出了重磅炸弹，``也就是，任意深度的 \allowbreak{}\NovelInline{I02995}{(2 \ 1-)}\allowbreak{} 操作 onto 到 \allowbreak{}\NovelInline{I02996}{2-2}\allowbreak{} 这个集合上，都得不到它。这个序数就是 \allowbreak{}\NovelInline{I02997}{2-2 \ 1-2-2}\allowbreak{}，即 \allowbreak{}\NovelInline{I02998}{2 \text{ onto } 2 \ 1-2-2}\allowbreak{}。''

``还可以更深入地了解一下 \allowbreak{}\NovelInline{I02999}{2-2 \ 1-2-2}\allowbreak{} 这个大序数的性质。暂时把它记作 \allowbreak{}\NovelInline{I03000}{X}\allowbreak{} 吧。'' asheep 长老解释道。

``首先，\allowbreak{}\NovelInline{I03001}{X}\allowbreak{} 要能折叠 \allowbreak{}\NovelInline{I03002}{(2 \ 1-)}\allowbreak{} 操作，所以它至少要是递归马洛序数，即 \allowbreak{}\NovelInline{I03003}{X \in 2-2}\allowbreak{}；'' ``然后，需要考虑一下这样的 \allowbreak{}\NovelInline{I03004}{X}\allowbreak{} 相对于普通的递归马洛序数有哪些不同之处。\allowbreak{}\NovelInline{I03005}{X}\allowbreak{} 能折叠出函数 \allowbreak{}\NovelInline{I03006}{f(\alpha) = M_\alpha}\allowbreak{} 的容许点，而这些容许点有什么特性呢？回到容许点的定义：''

\begin{quote}
\textbf{序数 \allowbreak{}\NovelInline{I03007}{\Theta}\allowbreak{} 为函数 \allowbreak{}\NovelInline{I03008}{f(\alpha)}\allowbreak{} 的容许点，等同于 \allowbreak{}\NovelInline{I03009}{\Theta}\allowbreak{} 是 \allowbreak{}\NovelInline{I03010}{f(\alpha)}\allowbreak{} 的不动点且 \allowbreak{}\NovelInline{I03011}{\Theta}\allowbreak{} 为容许序数。}
\end{quote}

``可以知道，这些容许点都是 \allowbreak{}\NovelInline{I03012}{f(\alpha) = M_\alpha}\allowbreak{} 的不动点，也即 \allowbreak{}\NovelInline{I03013}{\text{MFP} = (1-)^{1,0} 2-2}\allowbreak{}。可以证明，折叠它们的 \allowbreak{}\NovelInline{I03014}{X}\allowbreak{} 也要是 MFP。''

asheep 长老顿了顿，总结道：``所有的 \allowbreak{}\NovelInline{I03015}{(1-)^{1,0} 2-2}\allowbreak{} 自然也是 \allowbreak{}\NovelInline{I03016}{1-2-2}\allowbreak{}。那么，就得到了 \allowbreak{}\NovelInline{I03017}{2-2 \ 1-2-2}\allowbreak{} 同时是 \allowbreak{}\NovelInline{I03018}{2-2}\allowbreak{} 和 \allowbreak{}\NovelInline{I03019}{1-2-2}\allowbreak{} 的结论。更进一步的证明显示，\allowbreak{}\NovelInline{I03020}{2-2 \ 1-2-2}\allowbreak{} 与 \allowbreak{}\NovelInline{I03021}{2-2 \cap 1-2-2}\allowbreak{} 等价------这可不是一句无意义的废话。我们对 \allowbreak{}\NovelInline{I03022}{2-2 \ 1-2-2}\allowbreak{} 的定义实际上是 \allowbreak{}\NovelInline{I03023}{2-(2 \ 1-2-2)}\allowbreak{}，这里证明了交集和最左侧的 onto 运算次序可以交换。''

``可以对 \allowbreak{}\NovelInline{I03024}{2-2 \ 1-2-2}\allowbreak{} 继续进行 \allowbreak{}\NovelInline{I03025}{1-}\allowbreak{} 和 \allowbreak{}\NovelInline{I03026}{2 \ 1-}\allowbreak{} 操作，得到更大的序数。可以使用多元 \allowbreak{}\NovelInline{I03027}{M(1;0)}\allowbreak{} 来表示它们。''

梅天狸长老接管了讲解，她指着虚空中新出现的符号： \NovelDisplay{M00547R01}{%
1-2-2 \ 1-2-2 = \{M(1;\omega), M(1;\omega \times 2), \dots M(1;\omega^2), \dots\}%
} \NovelDisplay{M00548R01}{%
(1-)^{1,0} 2-2 \ 1-2-2 = \{\psi_{M(1;1,0)}(0), \psi_{M(1;1,0)}(1), \dots M(1;1,0), \dots\}%
}

``\allowbreak{}\NovelInline{I03028}{M(1;1,0)}\allowbreak{} 就是函数 \allowbreak{}\NovelInline{I03029}{f(\alpha) = M(1;\alpha)}\allowbreak{} 的容许点，它之于 \allowbreak{}\NovelInline{I03030}{M(1;n)}\allowbreak{}，就好比 \allowbreak{}\NovelInline{I03031}{I(1,0)}\allowbreak{} 之于 \allowbreak{}\NovelInline{I03032}{I_n}\allowbreak{}，\allowbreak{}\NovelInline{I03033}{M(1,0)}\allowbreak{} 之于 \allowbreak{}\NovelInline{I03034}{M_n}\allowbreak{}。''

她继续列举： \NovelDisplay{M00549R01}{%
2 \ 1-2-2 \ 1-2-2 = \{M(1;1,0), M(1;1,1), \dots, M(1;1,\omega+1), \dots\}%
} \NovelDisplay{M00550R01}{%
2 \ 1-2-2 \ 1-2-2 \ 1-2-2 = \{M(1;2,0), M(1;2,1), \dots, M(1;2,\omega+1), \dots\}%
} \NovelDisplay{M00551R01}{%
(2 \ 1-)^{\omega} 2-2 \ 1-2-2 = \{M(1;\omega,0), M(1;\omega,1), \dots, M(1;\omega,\omega), \dots\}%
}

``同样地，我们会有 \allowbreak{}\NovelInline{I03035}{(2 \ 1-)^{1,0,0} 1-2-2 \ 1-2-2}\allowbreak{}，\allowbreak{}\NovelInline{I03036}{(2 \ 1-)^{1@ \omega} 1-2-2 \ 1-2-2}\allowbreak{}，以及折叠它们的 \allowbreak{}\NovelInline{I03037}{2-2 \ 1-2-2 \ 1-2-2}\allowbreak{}。''

``这便是带分号的多元 \allowbreak{}\NovelInline{I03038}{M}\allowbreak{} 函数最低位的分号的进位规则了。它是马洛点进位。序数 \allowbreak{}\NovelInline{I03039}{\Theta}\allowbreak{} 为函数 \allowbreak{}\NovelInline{I03040}{f(\alpha)}\allowbreak{} 的马洛点，指的是 \allowbreak{}\NovelInline{I03041}{\Theta}\allowbreak{} 是 \allowbreak{}\NovelInline{I03042}{f(\alpha)}\allowbreak{} 的不动点且 \allowbreak{}\NovelInline{I03043}{\Theta}\allowbreak{} 为递归马洛序数。''

asheep 长老补充道：``\allowbreak{}\NovelInline{I03044}{M}\allowbreak{} 是函数 \allowbreak{}\NovelInline{I03045}{f(\alpha) = \omega^\alpha}\allowbreak{} 的马洛点，\allowbreak{}\NovelInline{I03046}{M(1;0)}\allowbreak{} 是函数 \allowbreak{}\NovelInline{I03047}{f(\alpha) = M_\alpha}\allowbreak{} 的马洛点，\allowbreak{}\NovelInline{I03048}{M(2;0)}\allowbreak{} 是函数 \allowbreak{}\NovelInline{I03049}{f(\alpha) = M(1;\alpha)}\allowbreak{} 的马洛点。''

自然地，还可以继续得到： \NovelDisplay{M00552R01}{%
2 \ 1-2-2 \ 1-2-2 \ 1-2-2 = \{M(2;1,0), M(2;1,1), \dots, M(2;1,\omega+1), \dots\}%
} \NovelDisplay{M00553R01}{%
2-2 \ 1-2-2 \ 1-2-2 \ 1-2-2 = \{M(3;0), M(3;1), \dots, M(3;\omega+1), \dots\}%
} \NovelDisplay{M00554R01}{%
(2-2 \ 1-)^{\omega} 2-2 = \{M(\omega;0), M(\omega;1), \dots, M(\omega;\omega), \dots\}%
}

``注意 psd. 和 real. 定义。这里的 \allowbreak{}\NovelInline{I03050}{(2-2 \ 1-)^{\omega} 2-2}\allowbreak{} 是 psd. 定义的，所以这个集合里面能含有 \allowbreak{}\NovelInline{I03051}{M(\omega;\omega)}\allowbreak{}。''

``我们将最小的 \allowbreak{}\NovelInline{I03052}{2-2 \ 1-2-2}\allowbreak{} 记作 \allowbreak{}\NovelInline{I03053}{M(1;0)}\allowbreak{}，次小的记为 \allowbreak{}\NovelInline{I03054}{M(1;1)}\allowbreak{}，以此类推，就有：'' \NovelDisplay{M00555R01}{%
2-2 \ 1-2-2 = \{M(1;0), M(1;1), \dots, M(1;\omega+1), \dots, M(1,1,0), \dots\}%
}

梅天狸长老最后总结道：``\allowbreak{}\NovelInline{I03055}{\psi_{M(1;0)}}\allowbreak{} 函数可以输出多元 \allowbreak{}\NovelInline{I03056}{M}\allowbreak{}，其规则与 \allowbreak{}\NovelInline{I03057}{\psi_M}\allowbreak{} 函数输出多元 \allowbreak{}\NovelInline{I03058}{I}\allowbreak{} 类似。有：''

\NovelDisplay{M00556R01}{%
\psi_{M(1;0)}(0) = M(1,0)%
} \NovelDisplay{M00557R01}{%
\psi_{M(1;0)}(1) = M(1,1)%
} \NovelDisplay{M00558R01}{%
\psi_{M(1;0)}(M(1;0)) = M(2,0)%
} \NovelDisplay{M00559R01}{%
\psi_{M(1;0)}(M(1;0)^\omega) = M(\omega,0)%
} \NovelDisplay{M00560R01}{%
\psi_{M(1;0)}(M(1;0)^{M(1;0)}) = M(1,0,0)%
} \NovelDisplay{M00561R01}{%
\psi_{M(1;0)}(M(1;0)^{M(1;0)^{M(1;0)}}) = M(1@(1,0))%
}

随着最后一个公式落下，广场上的空气仿佛凝固了。这些符号不再是枯燥的数学表达，它们化作了实质的压力，笼罩在八人头顶。

卢本韦深吸一口气，擦去额头的汗水：``原来如此\ldots\ldots 从 \allowbreak{}\NovelInline{I03059}{M}\allowbreak{} 到 \allowbreak{}\NovelInline{I03060}{M(1;0)}\allowbreak{}，再到 \allowbreak{}\NovelInline{I03061}{M(1,0,0)}\allowbreak{} 乃至 \allowbreak{}\NovelInline{I03062}{M(1@(1,0))}\allowbreak{}，这不仅仅是数字的增长，这是`描述能力'的指数级爆炸。每一个分号，每一个逗号，甚至那个 @ 符号，都是通往新世界的钥匙。''

霍雨浩眼中闪烁着狂热的光芒：``这就是多元 \allowbreak{}\NovelInline{I03063}{M}\allowbreak{} 函数的真意！它告诉我们，所谓的`大'，是没有止境的。只要我们愿意定义新的规则，愿意寻找新的不动点，序数的阶梯就能无限延伸！''

asheep 长老看着众人，微微颔首：``今日所学，名为`多元'，实为`无限'。你们手中的笔， désormais 不再只是记录工具，而是创世的神器。去吧，将这些规则融入你们的灵魂，下一次推演，我要看到属于你们自己的`新序数'！''

八人齐齐躬身，声音汇聚成一股洪流，震荡着整座大殿： ``弟子谨遵教诲！真伪已辨，折叠已成，愿以此基，攀登无尽序道！''

\chapter{第五十四章：不可转换之壁}\label{ux7b2cux4e94ux5341ux56dbux7ae0ux4e0dux53efux8f6cux6362ux4e4bux58c1}

广场上的晨风似乎停滞了一瞬。随着 asheep 长老指尖那红蓝光带的剧烈碰撞，一股比之前更为沉重、更为晦涩的威压笼罩了全场。这不再是单纯的``大''，而是一种结构上的``质变''。

``你们以为 \allowbreak{}\NovelInline{I03064}{M}\allowbreak{} 的多元扩展便是终点？'' asheep 长老冷笑一声，目光扫过众人，``\allowbreak{}\NovelInline{I03065}{M}\allowbreak{} 虽强，却仍可被`转换'，仍可被更低阶的逻辑所捕捉。但在序数的深渊中，存在着一种\textbf{不可被转换}的绝对壁垒。今日，我便带你们见识这堵墙------\textbf{不可转换序数}，以及通往 \allowbreak{}\NovelInline{I03066}{\Pi_3}\allowbreak{} 反射的荆棘之路。''

霍雨浩眉头紧锁，盯着虚空中那些不断自我吞噬的符号：``长老，资料中提及，当 \allowbreak{}\NovelInline{I03067}{2--}\allowbreak{} 操作的层级不断加深，甚至引入 \allowbreak{}\NovelInline{I03068}{(2-2 \cap 1-)^{1@\omega}}\allowbreak{} 这种折叠操作时，我们触及了一个新的极限。这个极限被称为 \textbf{\allowbreak{}\NovelInline{I03069}{2-2-2}\allowbreak{}}。''

``不错。'' asheep 长老点头，``\allowbreak{}\NovelInline{I03070}{2-2-2}\allowbreak{} 是一个大序数折叠的产物，它超越了 \allowbreak{}\NovelInline{I03071}{M}\allowbreak{} 的范畴。我们将构成这个集合的最小元素记为 \textbf{\allowbreak{}\NovelInline{I03072}{N}\allowbreak{}}（Non-projectible，不可转换）。''

他大手一挥，虚空中显现出一串全新的序列，其复杂程度远超之前的 \allowbreak{}\NovelInline{I03073}{M}\allowbreak{} 序列： \NovelDisplay{M00562R01}{%
2-2-2 = \{N, N_2, \dots, N_{\omega+1}, \dots\}%
}

``注意看，'' asheep 长老的声音变得低沉，``\allowbreak{}\NovelInline{I03074}{N}\allowbreak{} 的性质与 \allowbreak{}\NovelInline{I03075}{M}\allowbreak{} 截然不同。\allowbreak{}\NovelInline{I03076}{M}\allowbreak{} 是马洛点，而 \allowbreak{}\NovelInline{I03077}{N}\allowbreak{} 是\textbf{不可转换点}。这意味着你无法通过任何低阶的投影操作将 \allowbreak{}\NovelInline{I03078}{N}\allowbreak{} 映射回更小的序数集合中而不丢失信息。它是自成一体的孤岛。''

卢本韦迅速在脑海中构建模型：``既然有了 \allowbreak{}\NovelInline{I03079}{N}\allowbreak{}，那么针对 \allowbreak{}\NovelInline{I03080}{N}\allowbreak{} 的操作逻辑应当与 \allowbreak{}\NovelInline{I03081}{M}\allowbreak{} 同构？''

``正是！\textbf{多元 \allowbreak{}\NovelInline{I03082}{N}\allowbreak{} 函数}与\textbf{带分号 \allowbreak{}\NovelInline{I03083}{N}\allowbreak{}} 的规则与 \allowbreak{}\NovelInline{I03084}{M}\allowbreak{} 函数完全一致。'' asheep 长老赞许道，``例如：'' \NovelDisplay{M00563R01}{%
(1-)^{1,0} 2-2-2 = \{\psi_{N(1,0)}(0), \psi_{N(1,0)}(1), \dots, N(1,0), \dots\}%
} ``这是函数 \allowbreak{}\NovelInline{I03085}{f(\alpha) = N_\alpha}\allowbreak{} 的不动点，我们称之为 \textbf{NFP}。同理，我们有：'' \NovelDisplay{M00564R01}{%
2 \cap 1-2-2-2 = \{N(1,0), N(1,1), \dots, N(1,\omega+1), \dots\}%
} ``这是函数 \allowbreak{}\NovelInline{I03086}{f(\alpha) = N_\alpha}\allowbreak{} 的容许点。再进一步，加深马洛点的层次，可以得到 \allowbreak{}\NovelInline{I03087}{(2-2 1-)^\omega 2-2-2}\allowbreak{} 之类的东西。折叠它们的是 2-2-2∩1-2-2-2。也可以证明，这个集合实际上等同于 \allowbreak{}\NovelInline{I03088}{2-2-2}\allowbreak{} 和 \allowbreak{}\NovelInline{I03089}{1-2-2-2}\allowbreak{} 的交集。''

蔡须坤插话道：``所以，2-2-2∩1-2-2-2 就是函数 \allowbreak{}\NovelInline{I03090}{f(\alpha) = N_\alpha}\allowbreak{} 的 \textbf{\allowbreak{}\NovelInline{I03091}{N}\allowbreak{} 点}？''

``完全正确。'' asheep 长老眼中闪过一丝精光，``你们已经掌握了 \allowbreak{}\NovelInline{I03092}{N}\allowbreak{} 的呼吸节奏。''

``然而，\allowbreak{}\NovelInline{I03093}{N}\allowbreak{} 并非终局。'' asheep 长老话锋一转，虚空中的符号再次扭曲，``在 \allowbreak{}\NovelInline{I03094}{2-2-2 \cap 1-2-2-2}\allowbreak{} 之后，又可以加深 \allowbreak{}\NovelInline{I03095}{2-2-2 \cap 1-}\allowbreak{} 操作的层级，并引入一个更大的序数折叠这一操作。这样得到的就是 \textbf{\allowbreak{}\NovelInline{I03096}{2-2-2-2}\allowbreak{}} 了。''

他顿了顿，语气中带着一丝无奈：``这个序数过于`不整'，以至于它并没有一个专有的名字和符号，我这里按照个人习惯将其记作 \textbf{\allowbreak{}\NovelInline{I03097}{Q}\allowbreak{}}。'' \NovelDisplay{M00565R01}{%
2-2-2-2 = \{Q, Q_2, \dots, Q_{\omega+1}, \dots\}%
} ``\allowbreak{}\NovelInline{I03098}{2-2-2-2}\allowbreak{} 折叠 \allowbreak{}\NovelInline{I03099}{(2-2-2 1-)^{a,b,c,\dots}}\allowbreak{}。到这里，足以勾勒出一个对 \textbf{\allowbreak{}\NovelInline{I03100}{2 \text{ onto}}\allowbreak{}} 的印象了。''

asheep 长老双手合十，仿佛在演示某种挤压的动作：``我们可以这样理解 \allowbreak{}\NovelInline{I03101}{2 \text{ onto}}\allowbreak{} 集合 \allowbreak{}\NovelInline{I03102}{X}\allowbreak{} 的行为：\textbf{\allowbreak{}\NovelInline{I03103}{2 \text{ onto } X}\allowbreak{} 是对 \allowbreak{}\NovelInline{I03104}{(X \cap 1-)}\allowbreak{} 的折叠。}''

小舞若有所思：``也就是说，集合 \allowbreak{}\NovelInline{I03105}{X}\allowbreak{} 越强，\allowbreak{}\NovelInline{I03106}{2 \text{ onto}}\allowbreak{} 所折叠的操作也越强？反之，如果 \allowbreak{}\NovelInline{I03107}{X}\allowbreak{} 本身性质不够好，也会`连累' \allowbreak{}\NovelInline{I03108}{2 \text{ onto}}\allowbreak{} 的强度一起变弱？''

``悟性极佳。'' asheep 长老点头，``关于很弱的 \allowbreak{}\NovelInline{I03109}{2 \text{ onto}}\allowbreak{} 大概会在本系列的某一篇番外讲到。但现在，我们要看的是通往 \allowbreak{}\NovelInline{I03110}{2-2-2-2-2}\allowbreak{} 的路。函数 \allowbreak{}\NovelInline{I03111}{f(\alpha) = Q_\alpha}\allowbreak{} 有不动点、容许点、马洛点、\allowbreak{}\NovelInline{I03112}{N}\allowbreak{} 点、\allowbreak{}\NovelInline{I03113}{Q}\allowbreak{} 点。不断加深 \allowbreak{}\NovelInline{I03114}{Q}\allowbreak{} 点的层级，折叠`取 \allowbreak{}\NovelInline{I03115}{Q}\allowbreak{} 点'操作的就是 \allowbreak{}\NovelInline{I03116}{2-2-2-2-2}\allowbreak{} 了。''

``接下来，我们要面对的是真正的深水区------\allowbreak{}\NovelInline{I03117}{(2-)}\allowbreak{} 的反射以及 \allowbreak{}\NovelInline{I03118}{\Pi_3}\allowbreak{} 的一些枚举。'' asheep 长老的声音仿佛来自远古的回响。

``在上一篇文章的末尾，我们来到了 \allowbreak{}\NovelInline{I03119}{psd. (2-)^\omega}\allowbreak{} 以及 \allowbreak{}\NovelInline{I03120}{real. (2-)^\omega = (2-)^{\omega+1}}\allowbreak{}。现在，可以对它们进行 \allowbreak{}\NovelInline{I03121}{2-}\allowbreak{} 反射以继续前进了。''

虚空中浮现出层层叠叠的指数塔： \NovelDisplay{M00566R01}{%
(2-)^{\omega+2} \text{ 折叠任意深度的 } ((2-)^{\omega+1} 1-)^{a,b,c,\dots} (2-)^{\omega+1}%
} \NovelDisplay{M00567R01}{%
(2-)^{\omega+3} \text{ 折叠任意深度的 } ((2-)^{\omega+2} 1-)^{a,b,c,\dots} (2-)^{\omega+2}%
} \NovelDisplay{M00568R01}{%
\dots%
} \NovelDisplay{M00569R01}{%
\min (2-)^{\omega \times 2} = \sup\{\min (2-)^{\omega+1}, \min (2-)^{\omega+2}, \min (2-)^{\omega+3}, \dots\}%
} \NovelDisplay{M00570R01}{%
(2-)^{\omega \times 2 + 1} = real. (2-)^{\omega \times 2} \text{ 折叠 “取 } (2-)^{\omega \times 2} \text{ 点” 的操作}%
}

卢本韦看着那令人眼花缭乱的公式，沉声道：``上标上的序数可以一直增加，每次经过极限序数的时候都会碰到一点关于 \allowbreak{}\NovelInline{I03122}{psd.}\allowbreak{} 和 \allowbreak{}\NovelInline{I03123}{real.}\allowbreak{} 的破事。''

``没错。'' asheep 长老叹了口气，``不过，对于不想搭理这些破事的读者，记住 \textbf{\allowbreak{}\NovelInline{I03124}{psd.}\allowbreak{} 规则}和 \textbf{``\allowbreak{}\NovelInline{I03125}{real.}\allowbreak{} 相当于 \allowbreak{}\NovelInline{I03126}{+1}\allowbreak{}''} 就已经足够了。''

他继续推演：``接下来，还会有 \allowbreak{}\NovelInline{I03127}{(2-)^\Omega = (2-)^{(2)}}\allowbreak{}，\allowbreak{}\NovelInline{I03128}{(2-)^M = (2-)^{(2-2)}}\allowbreak{}，以及 \allowbreak{}\NovelInline{I03129}{(2-)^{(2-)^\omega}}\allowbreak{} 和 \allowbreak{}\NovelInline{I03130}{(2-)^{(2-)^{(2-)^\omega}}}\allowbreak{} \ldots\ldots 这样做的极限就是 \textbf{\allowbreak{}\NovelInline{I03131}{(2-)^{1,0}}\allowbreak{}} 了。''

``注意，'' asheep 长老特意强调，``到上一句话的用词是`极限'而不是`不动点'。这样用词是因为这里也有 \allowbreak{}\NovelInline{I03132}{psd. vs. real.}\allowbreak{} 的定义问题。对于 \allowbreak{}\NovelInline{I03133}{(2-)^{1,0}}\allowbreak{}，有如下定义：'' \NovelDisplay{M00571R01}{%
\alpha \in psd. (2-)^{1,0} \Leftrightarrow \forall \beta < \alpha, \exists \gamma \in (2-)^\beta, \gamma < \alpha%
} \NovelDisplay{M00572R01}{%
\alpha \in real. (2-)^{1,0} \Leftrightarrow \forall \beta < \alpha, \alpha \in (2-)^\beta%
}

霍雨浩喃喃自语：``简而言之，\allowbreak{}\NovelInline{I03134}{psd. (2-)^{1,0}}\allowbreak{} 只要求对于所有 \allowbreak{}\NovelInline{I03135}{\beta < \alpha}\allowbreak{}，在 \allowbreak{}\NovelInline{I03136}{\alpha}\allowbreak{} 的下方都存在一个 \allowbreak{}\NovelInline{I03137}{(2-)^\beta}\allowbreak{} 的序数，对 \allowbreak{}\NovelInline{I03138}{\alpha}\allowbreak{} 自身的性质未作出要求；而 \allowbreak{}\NovelInline{I03139}{real. (2-)^{1,0}}\allowbreak{} 要求 \allowbreak{}\NovelInline{I03140}{\alpha}\allowbreak{} 自身就是 \allowbreak{}\NovelInline{I03141}{(2-)^\beta}\allowbreak{} 序数。''

``\allowbreak{}\NovelInline{I03142}{psd. (2-)^{1,0}}\allowbreak{} 显然不是容许序数。'' asheep 长老补充道，``可以用类似 OFP 的方式为它找到一条长度为 \allowbreak{}\NovelInline{I03143}{\omega}\allowbreak{} 的递归基本列。不想搭理破事的读者依旧可以把 \allowbreak{}\NovelInline{I03144}{real. (2-)^{1,0}}\allowbreak{} 理解为 \allowbreak{}\NovelInline{I03145}{(2-)^{1,1}}\allowbreak{}，它折叠`取 \allowbreak{}\NovelInline{I03146}{(2-)^{1,0}}\allowbreak{} 点'的操作。''

\chapter{第五十五章：递归弱紧致序数的深渊凝视}\label{ux7b2cux4e94ux5341ux4e94ux7ae0ux9012ux5f52ux5f31ux7d27ux81f4ux5e8fux6570ux7684ux6df1ux6e0aux51ddux89c6}

``最后，是 \allowbreak{}\NovelInline{I03147}{2-(2-)^{1,1} = (2-)^{1,2}, (2-)^{1,\omega}, (2-)^{1,\Omega}, \dots}\allowbreak{} 以及这样做的极限 \allowbreak{}\NovelInline{I03148}{(2-)^{2,0}}\allowbreak{}。上标的规则没什么变化。进一步加深上标，就会得到 \allowbreak{}\NovelInline{I03149}{(2-)^{1,0,0}, (2-)^{1@\omega}, (2-)^{1@(1,0)}, \dots}\allowbreak{} 之类的东西。''

asheep 长老深吸一口气，抛出了最后的重磅概念：``是时候，把这些恼人的上标也送进 OCF 里折叠了。折叠 \allowbreak{}\NovelInline{I03150}{2-}\allowbreak{} 操作的序数称为\textbf{递归弱紧致序数}，写作 \textbf{\allowbreak{}\NovelInline{I03151}{K}\allowbreak{}}。全体递归弱紧致序数构成的集合也是 \allowbreak{}\NovelInline{I03152}{\Pi_3}\allowbreak{} 反射作用于全体序数得到的集合，在反射模式中记作 \textbf{3}。''

\NovelDisplay{M00573R01}{%
3 \text{ 折叠 } (2-)^{a,b,c,\dots}%
}

``这就是 \allowbreak{}\NovelInline{I03153}{\Pi_3}\allowbreak{} 的领域。'' asheep 长老的身影开始逐渐淡化，只留下最后一句告诫，``从这里开始，每一步都如履薄冰。但唯有跨过这道坎，你们才能真正触摸到序数分析的王座。''

广场上的空气仿佛凝固成了实质。随着 asheep 长老抛出``递归弱紧致序数 \allowbreak{}\NovelInline{I03154}{K}\allowbreak{}''这一概念，原本在众人脑海中已经构建好的 \allowbreak{}\NovelInline{I03155}{\Pi_2}\allowbreak{} 大厦瞬间显得如此渺小。

``\allowbreak{}\NovelInline{I03156}{K}\allowbreak{}\ldots\ldots''霍雨浩低声重复着这个符号，感到一股来自灵魂深处的战栗，``如果说 \allowbreak{}\NovelInline{I03157}{M}\allowbreak{} 是马洛，\allowbreak{}\NovelInline{I03158}{N}\allowbreak{} 是不可转换，那么 \allowbreak{}\NovelInline{I03159}{K}\allowbreak{} 就是这一切的集大成者，是 \allowbreak{}\NovelInline{I03160}{\Pi_3}\allowbreak{} 反射的具象化。''

asheep 长老的身影在晨光中若隐若现，他的声音不再仅仅是威严，更带有一种引导者面对未知时的狂热：``正是。所有的递归弱紧致序数，都同时身兼数职：它们是容许序数、递归不可达序数、递归马洛序数、递归不可转换序数\ldots\ldots 它是所有`大'性质的交汇点。''

他抬手一指，虚空中那原本代表 \allowbreak{}\NovelInline{I03161}{M}\allowbreak{} 和 \allowbreak{}\NovelInline{I03162}{N}\allowbreak{} 的光带瞬间崩解，重组成一个深邃无比的黑色漩涡，漩涡中心赫然写着一个巨大的数字：\textbf{3}。

``我们将全体递归弱紧致序数构成的集合记为 \textbf{3}。'' asheep 长老解释道，``在这个集合中，我们可以像之前枚举 \allowbreak{}\NovelInline{I03163}{M}\allowbreak{} 和 \allowbreak{}\NovelInline{I03164}{N}\allowbreak{} 一样，继续我们的征程。''

\NovelDisplay{M00574R01}{%
3 = \{K, K_2, K_3, \dots, K_{\omega+1}, \dots\}%
}

蔡须坤盯着那个序列，快速说道：``这意味着我们可以对集合 3 施加 \allowbreak{}\NovelInline{I03165}{\Pi_1}\allowbreak{}、\allowbreak{}\NovelInline{I03166}{\Pi_2}\allowbreak{} 的反射操作，从而得到更丰富的结构！''

``反应很快。'' asheep 点头，``看这里：'' \NovelDisplay{M00575R01}{%
1-3 = \{K_\omega, K_{\omega \times 2}, \dots, K_\Omega, \dots, K_K, \dots\}%
} ``这是 \allowbreak{}\NovelInline{I03167}{K}\allowbreak{} 的极限点序列。而当我们引入 \allowbreak{}\NovelInline{I03168}{(1-)^{1,0}}\allowbreak{} 操作时：'' \NovelDisplay{M00576R01}{%
(1-)^{1,0} 3 = \{\psi_{K(1,0)}(0), \psi_{K(1,0)}(1), \dots, \psi_{K(1,0)}(\omega), \dots, K(1,0), \dots\}%
} ``这便是所谓的 \textbf{KFP}（\allowbreak{}\NovelInline{I03169}{K}\allowbreak{}-Fixed Point）构成的集合。虽然 KFP 是个冷门概念，但它是通往更高阶的必经之路。''

紧接着，长老的手指在虚空中划出更复杂的轨迹： \NovelDisplay{M00577R01}{%
2 1-3 = \{K(1,0), K(1,1), \dots, K(1,\omega+1), \dots\}%
} ``这是函数 \allowbreak{}\NovelInline{I03170}{f(\alpha) = K_\alpha}\allowbreak{} 的容许点。把它们放进 OCF 的下标能输出丰富的 KFP。''

``再进一步，'' asheep 长老的声音提高了几分，``\allowbreak{}\NovelInline{I03171}{2 1-2 1-3}\allowbreak{} 折叠任意深度的 \allowbreak{}\NovelInline{I03172}{(2 1-)^{a,b,c,\dots} 3}\allowbreak{}。它们也是函数 \allowbreak{}\NovelInline{I03173}{f(\alpha) = K_\alpha}\allowbreak{} 的马洛点。''

``而 \allowbreak{}\NovelInline{I03174}{2-2 1-3}\allowbreak{} 折叠任意深度的 \allowbreak{}\NovelInline{I03175}{(2-2 1-)^{a,b,c,\dots} 3}\allowbreak{}。它们也是函数 \allowbreak{}\NovelInline{I03176}{f(\alpha) = K_\alpha}\allowbreak{} 的 \allowbreak{}\NovelInline{I03177}{N}\allowbreak{} 点。''

卢本韦此时插话道：``那么 \allowbreak{}\NovelInline{I03178}{(2-)^{1,0} 1-3}\allowbreak{} 呢？''

``问得好。'' asheep 赞许道，``\allowbreak{}\NovelInline{I03179}{(2-)^{1,0} 1-3}\allowbreak{} 是满足`对于任意 \allowbreak{}\NovelInline{I03180}{\beta < \alpha}\allowbreak{}，都存在 \allowbreak{}\NovelInline{I03181}{\gamma \in (2-)^\beta 1-3}\allowbreak{} 使得 \allowbreak{}\NovelInline{I03182}{\gamma < \alpha}\allowbreak{}'的全体序数 \allowbreak{}\NovelInline{I03183}{\alpha}\allowbreak{} 构成的集合。它也是集合 \allowbreak{}\NovelInline{I03184}{(2-)^{1,0}}\allowbreak{} 与集合 \allowbreak{}\NovelInline{I03185}{1-3}\allowbreak{} 的交集。''

``现在，我们要开始一段奇妙的旅程。'' asheep 长老目光深邃，``\textbf{\allowbreak{}\NovelInline{I03186}{3 \cap 1-}\allowbreak{}} 是集合 3 与集合 \allowbreak{}\NovelInline{I03187}{1-3}\allowbreak{} 的交集，折叠任意深度的 \allowbreak{}\NovelInline{I03188}{(2-)^{a,b,c,\dots} 1-3}\allowbreak{}。与 \allowbreak{}\NovelInline{I03189}{2 1-2}\allowbreak{} 类似，它也可以被理解为全体`足以使得 \allowbreak{}\NovelInline{I03190}{K_\alpha}\allowbreak{} 是递归弱紧致序数的极限序数 \allowbreak{}\NovelInline{I03191}{\alpha}\allowbreak{}'所构成的集合，以及函数 \allowbreak{}\NovelInline{I03192}{f(\alpha) = K_\alpha}\allowbreak{} 的 \allowbreak{}\NovelInline{I03193}{K}\allowbreak{} 点。''

他顿了顿，语气变得庄重：``从 3 到 \textbf{\allowbreak{}\NovelInline{I03194}{3 \cap 1-}\allowbreak{}}，我们相当于把从 Ord 到 3 的路又走了一遍。''

``再走第二遍的话，我们就会路过：'' \NovelDisplay{M00578R01}{%
2 1-3 \cap 1-, \quad 2-2 1-3 \cap 1-, \quad (2-)^{1,0} 1-3 \cap 1-, \dots%
} ``最终到达：'' \NovelDisplay{M00579R01}{%
3 \cap 1-3 \cap 1- = (3 \cap 1-)^2 3%
}

``再走第三遍，到达 \allowbreak{}\NovelInline{I03195}{(3 \cap 1-)^3 3}\allowbreak{}。'' ``一直走，离 \allowbreak{}\NovelInline{I03196}{(3 \cap 1-)^\omega 3}\allowbreak{} 越来越近。'' ``花吃奶的劲突破超限序数，到达 \allowbreak{}\NovelInline{I03197}{(3 \cap 1-)^{1,0} 3}\allowbreak{}。''

``而 \allowbreak{}\NovelInline{I03198}{2-3}\allowbreak{} 可以折叠任意深度的 \allowbreak{}\NovelInline{I03199}{(3 \cap 1-)^{a,b,c,\dots} 3}\allowbreak{}。正如我们上一篇所说的，\allowbreak{}\NovelInline{I03200}{2-X}\allowbreak{} 折叠 \allowbreak{}\NovelInline{I03201}{(X \cap 1-)}\allowbreak{} 的操作。\allowbreak{}\NovelInline{I03202}{X}\allowbreak{} 本身强得可怕的时候，\allowbreak{}\NovelInline{I03203}{2-X}\allowbreak{} 就会更加强得可怕。''

``以上，还只是 \allowbreak{}\NovelInline{I03204}{\Pi_3}\allowbreak{} 反射能形成的丰富结构中最靠前的一部分。''

``继续构造一个函数 \allowbreak{}\NovelInline{I03205}{f(\alpha)}\allowbreak{} 用来表示`第 \allowbreak{}\NovelInline{I03206}{\alpha}\allowbreak{} 个 \allowbreak{}\NovelInline{I03207}{2-3}\allowbreak{}'，然后继续前进。'' asheep 长老的手势愈发繁复，``中间又会经过 \allowbreak{}\NovelInline{I03208}{2 1-2-3}\allowbreak{}, \allowbreak{}\NovelInline{I03209}{2-2 1-2-3}\allowbreak{}, \allowbreak{}\NovelInline{I03210}{(2-)^{1,0} 1-2-3}\allowbreak{}，而 \textbf{\allowbreak{}\NovelInline{I03211}{3 \cap 1-2-3}\allowbreak{}} 是这个 \allowbreak{}\NovelInline{I03212}{f(\alpha)}\allowbreak{} 的 \allowbreak{}\NovelInline{I03213}{K}\allowbreak{} 点。''

``然后，是 \allowbreak{}\NovelInline{I03214}{(3 \cap 1-)^2 2-3}\allowbreak{}。'' ``\allowbreak{}\NovelInline{I03215}{(3 \cap 1-)^\omega 2-3}\allowbreak{}。'' ``\allowbreak{}\NovelInline{I03216}{(3 \cap 1-)^{1,0} 2-3}\allowbreak{}。''

``继续枚举吧。''

``\allowbreak{}\NovelInline{I03217}{2 1-3 \cap 2-3}\allowbreak{}。这是`计数 \allowbreak{}\NovelInline{I03218}{3 \cap 2-}\allowbreak{} 的函数'的容许点。'' ``\allowbreak{}\NovelInline{I03219}{2-2 1-3 \cap 2-3}\allowbreak{}。这是`计数 \allowbreak{}\NovelInline{I03220}{3 \cap 2-}\allowbreak{} 的函数'的马洛点。'' ``\allowbreak{}\NovelInline{I03221}{3 \cap 1-3 \cap 2-3}\allowbreak{}。这是`计数 \allowbreak{}\NovelInline{I03222}{3 \cap 2-}\allowbreak{} 的函数'的 \allowbreak{}\NovelInline{I03223}{K}\allowbreak{} 点。'' ``\allowbreak{}\NovelInline{I03224}{3 \cap 1-3 \cap 1-3 \cap 2-3}\allowbreak{}。'' ``\allowbreak{}\NovelInline{I03225}{(3 \cap 1-)^\omega 3 \cap 2-3}\allowbreak{}。''

``这里是一个易错点。'' asheep 长老特意停顿，``\allowbreak{}\NovelInline{I03226}{2-3 \cap 1-3 \cap 2-3}\allowbreak{} 折叠任意深度的 \allowbreak{}\NovelInline{I03227}{(3 \cap 1-)^{a,b,c,\dots} 2-3}\allowbreak{}。根据前面提到的`\allowbreak{}\NovelInline{I03228}{2-3}\allowbreak{} 折叠 \allowbreak{}\NovelInline{I03229}{(3 \cap 1-)^{a,b,c,\dots}}\allowbreak{}'可以知道。折叠 \allowbreak{}\NovelInline{I03230}{(3 \cap 1-)^{a,b,c,\dots} 2-3}\allowbreak{} 的应该是`计数 \allowbreak{}\NovelInline{I03231}{3 \cap 2-}\allowbreak{} 的函数'的 \allowbreak{}\NovelInline{I03232}{2-3}\allowbreak{} 点，也就是 \textbf{\allowbreak{}\NovelInline{I03233}{2-3 \cap 1-3 \cap 2-3}\allowbreak{}}，而并非 \allowbreak{}\NovelInline{I03234}{2-3 \cap 2-3}\allowbreak{}。''

``\allowbreak{}\NovelInline{I03235}{2-3 \cap 1-2-3 \cap 1-3 \cap 2-3 = (2-3 \cap 1-)^2 3 \cap 2-3}\allowbreak{}。'' ``\allowbreak{}\NovelInline{I03236}{2-2-3 \cap 1-3 \cap 2-3}\allowbreak{} 折叠任意深度的 \allowbreak{}\NovelInline{I03237}{(2-3 \cap 1-)^{a,b,c,\dots} 3 \cap 2-3}\allowbreak{}。'' ``\allowbreak{}\NovelInline{I03238}{2-2-2-3 \cap 1-3 \cap 2-3}\allowbreak{}。'' ``\allowbreak{}\NovelInline{I03239}{(2-)^\omega 3 \cap 1-3 \cap 2-3}\allowbreak{}。''

``然后，有 \allowbreak{}\NovelInline{I03240}{3 \cap 2-3 \cap 1-3 \cap 2-3}\allowbreak{} 折叠任意深度的 \allowbreak{}\NovelInline{I03241}{(2-)^{a,b,c,\dots} 3 \cap 1-3 \cap 2-3}\allowbreak{}。它是`计数 \allowbreak{}\NovelInline{I03242}{3 \cap 2-}\allowbreak{} 的函数'的 \allowbreak{}\NovelInline{I03243}{3 \cap 2-}\allowbreak{} 点。这是我们第一次在`XX点'的概念中涉及到交集。\allowbreak{}\NovelInline{I03244}{2-3 \cap 2-3}\allowbreak{} 要折叠的正是这个`取 \allowbreak{}\NovelInline{I03245}{3 \cap 2-}\allowbreak{} 点'的概念。''

``再取一次 \allowbreak{}\NovelInline{I03246}{3 \cap 2-}\allowbreak{} 点看看。'' ``\allowbreak{}\NovelInline{I03247}{3 \cap 1-3 \cap 2-3 \cap 1-3 \cap 2-3}\allowbreak{}。''

``什么折叠任意深度的 \allowbreak{}\NovelInline{I03248}{(3 \cap 1-)^{a,b,c,\dots} 2-3}\allowbreak{} 呢？注意到从 \allowbreak{}\NovelInline{I03249}{2-3}\allowbreak{} 到这里的过程相当于从 Ord 到 \allowbreak{}\NovelInline{I03250}{2-3}\allowbreak{}，所以它应该是 \allowbreak{}\NovelInline{I03251}{2-3 \cap 1-2-3}\allowbreak{}，也就是 \allowbreak{}\NovelInline{I03252}{f(\alpha)}\allowbreak{} 的 \allowbreak{}\NovelInline{I03253}{2-3}\allowbreak{} 点。''

``利用`取 \allowbreak{}\NovelInline{I03254}{2-3}\allowbreak{} 点'的操作，又可以有：'' \NovelDisplay{M00580R01}{%
2-3 \cap 1-2-3 \cap 1-2-3 = (2-3 \cap 1-)^2 2-3%
} \NovelDisplay{M00581R01}{%
(2-3 \cap 1-)^\omega 2-3%
} \NovelDisplay{M00582R01}{%
(2-3 \cap 1-)^{1,0} 2-3%
}

``在这之后，\allowbreak{}\NovelInline{I03255}{2-2-3}\allowbreak{} 能折叠任意深度的 \allowbreak{}\NovelInline{I03256}{(2-3 \cap 1-)^{a,b,c,\dots} 2-3}\allowbreak{}，也就是`取 \allowbreak{}\NovelInline{I03257}{2-3}\allowbreak{} 点'这么一个操作。''

``\allowbreak{}\NovelInline{I03258}{2-2-2-3}\allowbreak{} 折叠任意深度的 \allowbreak{}\NovelInline{I03259}{(2-2-3 \cap 1-)^{a,b,c,\dots} 2-2-3}\allowbreak{}。这一部分和前面的 \allowbreak{}\NovelInline{I03260}{2-2-2}\allowbreak{} 没什么不同，只是底层由`\allowbreak{}\NovelInline{I03261}{1-}\allowbreak{} 反射'变为了`\allowbreak{}\NovelInline{I03262}{3 \cap 1-}\allowbreak{} 操作'。''

``\allowbreak{}\NovelInline{I03263}{(2-)^\omega 3}\allowbreak{}。'' ``\allowbreak{}\NovelInline{I03264}{(2-)^{1,0} 3}\allowbreak{}。在它下面的每一个 \allowbreak{}\NovelInline{I03265}{\beta}\allowbreak{}，都存在一个 \allowbreak{}\NovelInline{I03266}{(2-)^\beta 3}\allowbreak{} 序数小于它。''

``接下来考虑一个能折叠任意深度的 \allowbreak{}\NovelInline{I03267}{(2-)^{a,b,c,\dots} 3}\allowbreak{} 的序数。这个序数能折叠 \allowbreak{}\NovelInline{I03268}{2-}\allowbreak{} 的操作，所以是一个 \allowbreak{}\NovelInline{I03269}{\Pi_3}\allowbreak{} 序数；同时，它又至少是 \allowbreak{}\NovelInline{I03270}{2-3}\allowbreak{}（任意的 \allowbreak{}\NovelInline{I03271}{(2-)^\beta 3}\allowbreak{}，其中 \allowbreak{}\NovelInline{I03272}{\beta > 0}\allowbreak{}，都是 \allowbreak{}\NovelInline{I03273}{2-3}\allowbreak{}）。因此，这个序数至少是 3 和 \allowbreak{}\NovelInline{I03274}{2-3}\allowbreak{} 的交集。依旧可以证明，这个条件已经足够了。也就是：''

\NovelDisplay{M00583R01}{%
3 \cap 2-3 \text{ 折叠 } (2-)^{a,b,c,\dots} 3%
}

``\allowbreak{}\NovelInline{I03275}{2-3 \cap 1-3 \cap 2-3 \cap 1-3 \cap 2-3}\allowbreak{} 折叠任意深度的 \allowbreak{}\NovelInline{I03276}{(3 \cap 1-)^{a,b,c,\dots} 3 \cap 2-3 \cap 1-3 \cap 2-3}\allowbreak{}。''

``\allowbreak{}\NovelInline{I03277}{2-2-3 \cap 1-3 \cap 2-3 \cap 1-3 \cap 2-3}\allowbreak{} 折叠任意深度的 \allowbreak{}\NovelInline{I03278}{(2-3 \cap 1-)^{a,b,c,\dots} 3 \cap 2-3 \cap 1-3 \cap 2-3}\allowbreak{}。''

``\allowbreak{}\NovelInline{I03279}{3 \cap 2-3 \cap 1-3 \cap 2-3 \cap 1-3 \cap 2-3 = (3 \cap 2-3 \cap 1-)^2 3 \cap 2-3}\allowbreak{} 折叠任意深度的 \allowbreak{}\NovelInline{I03280}{(2-)^{a,b,c,\dots} 3 \cap 1-3 \cap 2-3 \cap 1-3 \cap 2-3}\allowbreak{}。''

``再才有 \allowbreak{}\NovelInline{I03281}{2-3 \cap 2-3}\allowbreak{} 折叠任意深度的 \allowbreak{}\NovelInline{I03282}{(3 \cap 2-3 \cap 1-)^{a,b,c,\dots} 3 \cap 2-3}\allowbreak{}。''

``\allowbreak{}\NovelInline{I03283}{2-3 \cap 2-3 \cap 1-2-3 \cap 2-3}\allowbreak{} 是`计数 \allowbreak{}\NovelInline{I03284}{2-3 \cap 2-3}\allowbreak{} 的函数'的 \allowbreak{}\NovelInline{I03285}{2-3 \cap 2-3}\allowbreak{} 点。''

``\allowbreak{}\NovelInline{I03286}{2-2-3 \cap 2-3}\allowbreak{} 折叠任意深度的 \allowbreak{}\NovelInline{I03287}{(2-3 \cap 2-3 \cap 1-)^{a,b,c,\dots} 2-3 \cap 2-3}\allowbreak{}。''

``\allowbreak{}\NovelInline{I03288}{3 \cap 2-3 \cap 2-3}\allowbreak{} 折叠任意深度的 \allowbreak{}\NovelInline{I03289}{(2-)^{a,b,c,\dots} 3 \cap 2-3}\allowbreak{}。''

``\allowbreak{}\NovelInline{I03290}{2-3 \cap 2-3 \cap 2-3}\allowbreak{} 折叠任意深度的 \allowbreak{}\NovelInline{I03291}{(3 \cap 2-3 \cap 2-3 \cap 1-)^{a,b,c,\dots} 3 \cap 2-3 \cap 2-3}\allowbreak{}。简直疯了。''

asheep 长老深吸一口气，目光扫过众人：``稍微理一下这一段的逻辑：我们对一个表达式不断地进行 \allowbreak{}\NovelInline{I03292}{2-}\allowbreak{} 的操作，折叠 \allowbreak{}\NovelInline{I03293}{2-}\allowbreak{} 操作的是最前面有个`3'（注意空格）的表达式；然后对这个有`3'的表达式再不断进行 \allowbreak{}\NovelInline{I03294}{2-}\allowbreak{} 操作，再折叠它\ldots\ldots 最终形成一条 \allowbreak{}\NovelInline{I03295}{\dots 3 \cap 2-3 \cap 2-3 \cap 2-3}\allowbreak{} 的链条。每一次最基础的 \allowbreak{}\NovelInline{I03296}{2-}\allowbreak{} 操作，都要折叠一次`取自己点'的操作。好一个折叠再折叠。''

``延伸这个链条到超限序数，就能看到 \allowbreak{}\NovelInline{I03297}{(3 \cap 2-)^\omega 3}\allowbreak{}。'' ``还有 \allowbreak{}\NovelInline{I03298}{(3 \cap 2-)^{1,0} 3}\allowbreak{}。''

``现在，终于可以打开 3 onto 的大门了。''

广场上的晨光已经彻底铺展开来，将八名少年少女的影子拉得极长。方才那番关于 \allowbreak{}\NovelInline{I03299}{3 2-3}\allowbreak{} 折叠 \allowbreak{}\NovelInline{I03300}{(2-)^{a,b,c,\dots} 3}\allowbreak{} 的推演仍在众人脑海中回荡，而 asheep 长老的身影已重新凝聚于虚空之中，周身的光带由红蓝双色渐变为一道深邃的紫金色。

``你们以为打开 \allowbreak{}\NovelInline{I03301}{3 \text{ onto}}\allowbreak{} 的大门便是终局？'' asheep 长老的声音带着一种近乎戏谑的冷意，``那不过是序数深渊中又一层薄冰罢了。今日，我要带你们看看这扇门后的世界------\textbf{\allowbreak{}\NovelInline{I03302}{\Pi_3 \text{ onto}}\allowbreak{} 的行为}，以及由此一路畅通至 \textbf{\allowbreak{}\NovelInline{I03303}{\Pi_\omega}\allowbreak{}} 的壮丽征途。''

霍雨浩深吸一口气，率先开口：``长老，根据之前的经验，\allowbreak{}\NovelInline{I03304}{2 \text{ onto } X}\allowbreak{} 是对 \allowbreak{}\NovelInline{I03305}{(X \cap 1-)}\allowbreak{} 的折叠。那么 \allowbreak{}\NovelInline{I03306}{3 \text{ onto}}\allowbreak{} 是否也遵循类似的规律？''

``正是。'' asheep 长老点头，虚空中浮现出两行金色的公式：

\NovelDisplay{M00584R01}{%
3 \text{ onto } X \text{ 是对 } (X \cap 2-) \text{ 的折叠。}%
}

蔡须坤立刻反应过来：``也就是说，\allowbreak{}\NovelInline{I03307}{3 \text{ onto}}\allowbreak{} 与 \allowbreak{}\NovelInline{I03308}{2 \text{ onto}}\allowbreak{} 的行为如出一辙，只是折叠的对象从 \allowbreak{}\NovelInline{I03309}{X \cap 1-}\allowbreak{} 升级为了 \allowbreak{}\NovelInline{I03310}{X \cap 2-}\allowbreak{}！''

``不止如此。'' asheep 长老的目光变得深邃，``我们可以将此规律推广至任意 \allowbreak{}\NovelInline{I03311}{k \ge 2}\allowbreak{}：''

\NovelDisplay{M00585R01}{%
k \text{ onto } X \text{ 是对 } (X \cap (k-1)-) \text{ 的折叠。}%
}

小舞喃喃道：``\allowbreak{}\NovelInline{I03312}{k}\allowbreak{} 越大，折叠的操作越强\ldots\ldots 这意味着，理论上我们已经畅通无阻，直通 \allowbreak{}\NovelInline{I03313}{\Pi_\omega}\allowbreak{} 了？''

``理论上，是的。'' asheep 长老嘴角微微上扬，``但理论归理论，一步步枚举到它，可没那么轻松。来吧。''

\begin{center}\rule{0.5\linewidth}{0.5pt}\end{center}

\chapter{第五十六章：反射序数的终点}\label{ux7b2cux4e94ux5341ux516dux7ae0ux53cdux5c04ux5e8fux6570ux7684ux7ec8ux70b9}

``首先，有：''

\NovelDisplay{M00586R01}{%
3-3 \text{ 折叠任意深度的 } (3\ 2-)^{a,b,c,\dots}\ 3%
}

asheep 长老的声音沉稳而有力，``通过上面对 \allowbreak{}\NovelInline{I03314}{3-}\allowbreak{} 的刻画不难察觉这一点。''

紧接着，一连串令人目眩的符号如暴雨般倾泻而下：

\NovelDisplay{M00587R01}{%
3\ 1-3-3 \text{。计数上一行那玩意的函数的 } K \text{ 点。}%
}

\NovelDisplay{M00588R01}{%
(3\ 1-)^2\ 3-3 \text{。计数上一行那玩意的函数的 } K \text{ 点。}%
}

\NovelDisplay{M00589R01}{%
2-3\ 1-3-3 \text{ 折叠任意深度的 } (3\ 1-)^{a,b,c,\dots}\ 3-3 \text{。}%
}

\NovelDisplay{M00590R01}{%
2-2-3\ 1-3-3 \text{ 折叠任意深度的 } (2-3\ 1-)^{a,b,c,\dots}\ 3-3 \text{。}%
}

\NovelDisplay{M00591R01}{%
3\ 2-3\ 1-3-3 \text{ 折叠任意深度的 } (2-)^{a,b,c,\dots}\ 3\ 1-3-3 \text{。}%
}

卢本韦的额头开始渗出细密的汗珠，但他仍紧紧跟上节奏，大声念出下一行：

\NovelDisplay{M00592R01}{%
3\ 1-3\ 2-3\ 1-3-3 \text{ 是计数上一行那玩意的函数的 } K \text{ 点。}%
}

\NovelDisplay{M00593R01}{%
2-3\ 1-3\ 2-3\ 1-3-3 \text{ 折叠任意深度的 } (3\ 1-)^{a,b,c,\dots}\ 3\ 2-3\ 1-3-3 \text{。}%
}

\NovelDisplay{M00594R01}{%
2-3\ 2-3\ 1-3-3 \text{ 折叠任意深度的 } (3\ 2-3\ 1-)^{a,b,c,\dots}\ 3-3 \text{。}%
}

\NovelDisplay{M00595R01}{%
3-3\ 1-3-3 \text{ 折叠任意深度的 } (3\ 2-)^{a,b,c,\dots}\ 3\ 1-3-3 \text{。}%
}

\NovelDisplay{M00596R01}{%
3\ 1-3-3\ 1-3-3 \text{ 折叠任意深度的 } (2-)^{a,b,c,\dots}\ 1-3-3\ 1-3-3 \text{。}%
}

\NovelDisplay{M00597R01}{%
2-3\ 1-3-3\ 1-3-3 \text{ 折叠任意深度的 } (3\ 1-)^{a,b,c,\dots}\ 3-3\ 1-3-3 \text{。}%
}

\NovelDisplay{M00598R01}{%
3\ 2-3\ 1-3-3\ 1-3-3 \text{ 折叠任意深度的 } (2-)^{a,b,c,\dots}\ 3\ 1-3-3\ 1-3-3 \text{。}%
}

\NovelDisplay{M00599R01}{%
3-3\ 1-3-3\ 1-3-3 = (3-3\ 1-)^2\ 3-3 \text{ 折叠任意深度的 } (3\ 2-)^{a,b,c,\dots}\ 1-3-3\ 1-3-3 \text{。}%
}

霍雨浩感到自己的思维正在被一股无形的力量拉伸至极限，但他咬紧牙关，继续跟随长老的推演。

``然后------'' asheep 长老的声音陡然拔高，

\NovelDisplay{M00600R01}{%
2-3-3 \text{ 折叠任意深度的 } (3-3\ 1-)^{a,b,c,\dots}\ 3-3 \text{。}%
}

\NovelDisplay{M00601R01}{%
2-2-3-3 \text{ 折叠任意深度的 } (2-3-3\ 1-)^{a,b,c,\dots}\ 2-3-3 \text{。}%
}

\NovelDisplay{M00602R01}{%
3\ 2-3-3 \text{ 折叠任意深度的 } (2-)^{a,b,c,\dots}\ 3-3 \text{。}%
}

\NovelDisplay{M00603R01}{%
2-3\ 2-3-3 \text{ 折叠任意深度的 } (3\ 2-3-3\ 1-)^{a,b,c,\dots}\ 3\ 2-3-3 \text{。}%
}

\NovelDisplay{M00604R01}{%
3-3\ 2-3-3 \text{ 折叠任意深度的 } (3\ 2-)^{a,b,c,\dots}\ 3-3 \text{。}%
}

\NovelDisplay{M00605R01}{%
3-3-3 \text{ 折叠任意深度的 } (3-3\ 2-)^{a,b,c,\dots}\ 3-3 \text{。}%
}

梅天狸长老此时飘然而至，她看着众人几近崩溃的神情，淡淡补充道：``别忘了，这只是 \allowbreak{}\NovelInline{I03315}{\Pi_3}\allowbreak{} 反射能形成的丰富结构中的一小段。继续------''

\NovelDisplay{M00606R01}{%
2-3-3-3 \text{ 折叠任意深度的 } (3-3-3\ 1-)^{a,b,c,\dots}\ 3-3-3 \text{。}%
}

\NovelDisplay{M00607R01}{%
3\ 2-3-3-3 \text{ 折叠任意深度的 } (2-)^{a,b,c,\dots}\ 3-3-3 \text{。}%
}

\NovelDisplay{M00608R01}{%
3-3\ 2-3-3-3 \text{ 折叠任意深度的 } (3\ 2-)^{a,b,c,\dots}\ 3-3-3 \text{。}%
}

\NovelDisplay{M00609R01}{%
3-3-3\ 2-3-3-3 \text{ 折叠任意深度的 } (3-3\ 2-)^{a,b,c,\dots}\ 3-3-3 \text{。}%
}

\NovelDisplay{M00610R01}{%
3-3-3-3 \text{ 折叠任意深度的 } (3-3-3\ 2-)^{a,b,c,\dots}\ 3-3-3 \text{。}%
}

asheep 长老抬手，虚空中那无穷无尽的符号链条终于汇聚成了一个巨大的极限符号：

\NovelDisplay{M00611R01}{%
(3-)^\omega = \sup\{3,\ 3-3,\ 3-3-3,\ 3-3-3-3,\ \dots\}%
}

``以及------''

\NovelDisplay{M00612R01}{%
(3-)^{1,0} \text{ 是满足"对于所有 } \beta < \alpha \text{，在 } \alpha \text{ 的下方都存在一个 } (3-)^\beta \text{"的序数 } \alpha \text{ 构成的集合。}%
}

``现在，'' asheep 长老的声音如同雷鸣般在广场上空炸响，``是时候迎接新的纪元的到来了。''

虚空中的紫金色光带猛然炸裂，化作一片璀璨的星河。星河的中心，一个巨大的数字缓缓浮现：\textbf{4}。

\NovelDisplay{M00613R01}{%
4 \text{ 折叠任意深度的 } (3-)^{a,b,c,\dots}%
}

``我们将 \allowbreak{}\NovelInline{I03316}{\Pi_4}\allowbreak{} 序数记作 \textbf{\allowbreak{}\NovelInline{I03317}{\kappa}\allowbreak{}}。'' asheep 长老庄严宣告。

\NovelDisplay{M00614R01}{%
4 = \{\kappa,\ \kappa_2,\ \kappa_3,\ \dots,\ \kappa_{\omega+1},\ \dots\}%
}

蔡须坤看着那个序列，嘴角忍不住抽搐了一下：``\allowbreak{}\NovelInline{I03318}{(1-)^{1,0}\ 4}\allowbreak{} 是 \allowbreak{}\NovelInline{I03319}{\kappa \text{FP}}\allowbreak{} 的集合\ldots\ldots 这名字，真是幽默。''

asheep 长老难得地露出一丝笑意：``幽默归幽默，但它的力量不容小觑。继续枚举------''

\NovelDisplay{M00615R01}{%
2\ 1-4 \text{ 是 } f(\alpha) = \kappa_\alpha \text{ 的容许点。}%
}

\NovelDisplay{M00616R01}{%
3\ 1-4 \text{ 是 } f(\alpha) = \kappa_\alpha \text{ 的 } K \text{ 点。}%
}

\NovelDisplay{M00617R01}{%
4\ 1-4 \text{ 是 } f(\alpha) = \kappa_\alpha \text{ 的 } \kappa \text{ 点。}%
}

\NovelDisplay{M00618R01}{%
2-4 \text{ 折叠任意深度的 } (4\ 1-)^{a,b,c,\dots}\ 4 \text{。}%
}

\NovelDisplay{M00619R01}{%
3\ 2-4 \text{ 折叠任意深度的 } (2-)^{a,b,c,\dots}\ 4 \text{。}%
}

霍雨浩此时已经完全进入了状态，他接过了长老的话头，声音坚定而清晰：

\NovelDisplay{M00620R01}{%
2-3\ 2-4 \text{ 折叠任意深度的 } (3\ 2-4\ 1-)^{a,b,c,\dots}\ 3\ 2-4 \text{。}%
}

\NovelDisplay{M00621R01}{%
3-3\ 2-4 \text{ 折叠任意深度的 } (3\ 2-)^{a,b,c,\dots}\ 3\ 2-4 \text{。}%
}

\NovelDisplay{M00622R01}{%
4\ 2-4 \text{ 折叠任意深度的 } (3-)^{a,b,c,\dots}\ 2-4 \text{。}%
}

\NovelDisplay{M00623R01}{%
2-4\ 2-4 \text{ 折叠任意深度的 } (4\ 2-4\ 1-)^{a,b,c,\dots}\ 4\ 2-4 \text{。}%
}

\NovelDisplay{M00624R01}{%
3-4 \text{ 折叠任意深度的 } (4\ 2-)^{a,b,c,\dots}\ 4 \text{。}%
}

\NovelDisplay{M00625R01}{%
2-3-4 \text{ 折叠任意深度的 } (3-4\ 1-)^{a,b,c,\dots}\ 3-4 \text{。}%
}

\NovelDisplay{M00626R01}{%
2-2-3-4 \text{ 折叠任意深度的 } (2-3-4\ 1-)^{a,b,c,\dots}\ 2-3-4 \text{。}%
}

\NovelDisplay{M00627R01}{%
3\ 2-3-4 \text{ 折叠任意深度的 } (2-)^{a,b,c,\dots}\ 3-4 \text{。}%
}

\NovelDisplay{M00628R01}{%
3-3\ 2-3-4 \text{ 折叠任意深度的 } (3\ 2-)^{a,b,c,\dots}\ 3-4 \text{。}%
}

\NovelDisplay{M00629R01}{%
4\ 2-3-4 \text{ 折叠任意深度的 } (3-)^{a,b,c,\dots}\ 2-3-4 \text{。}%
}

卢本韦紧随其后，声音越来越快：

\NovelDisplay{M00630R01}{%
2-4\ 2-3-4 \text{ 折叠任意深度的 } (4\ 2-3-4\ 1-)^{a,b,c,\dots}\ 4\ 2-3-4 \text{。}%
}

\NovelDisplay{M00631R01}{%
3-4\ 2-3-4 \text{ 折叠任意深度的 } (4\ 2-)^{a,b,c,\dots}\ 3-4 \text{。}%
}

\NovelDisplay{M00632R01}{%
3-3-4 \text{ 折叠任意深度的 } (3-4\ 2-)^{a,b,c,\dots}\ 3-4 \text{。}%
}

\NovelDisplay{M00633R01}{%
4\ 3-4 \text{ 折叠任意深度的 } (3-)^{a,b,c,\dots}\ 4 \text{。}%
}

\NovelDisplay{M00634R01}{%
2-4\ 3-4 \text{ 折叠任意深度的 } (4\ 3-4\ 1-)^{a,b,c,\dots}\ 4\ 3-4 \text{。}%
}

\NovelDisplay{M00635R01}{%
3-4\ 2-4\ 3-4 \text{ 折叠任意深度的 } (4\ 2-)^{a,b,c,\dots}\ 4\ 3-4 \text{。}%
}

\NovelDisplay{M00636R01}{%
3-4\ 3-4 \text{ 折叠任意深度的 } (4\ 3-4\ 2-)^{a,b,c,\dots}\ 4\ 3-4 \text{。}%
}

\NovelDisplay{M00637R01}{%
4-4 \text{ 折叠任意深度的 } (4\ 3-)^{a,b,c,\dots}\ 4 \text{。}%
}

\NovelDisplay{M00638R01}{%
4-4-4 \text{ 折叠任意深度的 } (4-4\ 3-)^{a,b,c,\dots}\ 4-4 \text{。}%
}

asheep 长老的身影开始膨胀，仿佛要与整个天空融为一体。他的声音从四面八方传来：

\NovelDisplay{M00639R01}{%
5 \text{ 折叠任意深度的 } (4-)^{a,b,c,\dots}%
}

\NovelDisplay{M00640R01}{%
6 \text{ 折叠任意深度的 } (5-)^{a,b,c,\dots}%
}

\NovelDisplay{M00641R01}{%
\dots%
}

``而 \textbf{\allowbreak{}\NovelInline{I03320}{\Pi_\omega}\allowbreak{} 反射}------'' asheep 长老的声音在最高点骤然停顿，然后以无比庄严的语调落下最后的话语，``当然，也是 \allowbreak{}\NovelInline{I03321}{psd.}\allowbreak{} 的。你问 \allowbreak{}\NovelInline{I03322}{real.\ \Pi_\omega}\allowbreak{}？到稳定阶段再去搞它------\textbf{在终点等着我们。结束了。}''

广场上陷入了长久的沉默。八名少年少女呆立原地，他们刚刚见证了从 \allowbreak{}\NovelInline{I03323}{\Pi_3 \text{ onto}}\allowbreak{} 一路狂飙至 \allowbreak{}\NovelInline{I03324}{\Pi_\omega}\allowbreak{} 的壮丽旅程。那无穷无尽的反射式链条，如同一条蜿蜒的巨龙，从脚下延伸至天际的尽头。

就在众人还沉浸在震撼之中时，梅天狸长老悄然出现在广场的另一侧。她的声音轻柔却清晰，如同一泓清泉注入众人干涸的脑海：

``你们是否想过，是否存在一套简单的规律，以便判断某一个反射式折叠了什么操作？''

众人纷纷转头，目光中满是渴望。

``答案是肯定的。'' 梅天狸长老微微一笑，``并且这一规律奇妙地与 \textbf{PrSS}（Protocol of Sequence System）存在联系。''

她抬手，虚空中浮现出四个步骤，每一个步骤都以金色的光芒书写：

\textbf{第一步：替换。}

``把反射式中间的空格和 \allowbreak{}\NovelInline{I03325}{-}\allowbreak{} 都替换为逗号，然后倒序排列。''

她以 \allowbreak{}\NovelInline{I03326}{2-4\ 2-3-4}\allowbreak{} 为例： \NovelDisplay{M00642R01}{%
2-4\ 2-3-4 \xrightarrow{\text{替换}} 2, 4, 2, 3, 4 \xrightarrow{\text{倒序}} 4, 3, 2, 4, 2%
}

\textbf{第二步：补项。}

``若第一项为 \allowbreak{}\NovelInline{I03327}{k}\allowbreak{}，在前面补上 \allowbreak{}\NovelInline{I03328}{1, 2, \dots, k-1}\allowbreak{}；若式子中后一项减前一项的差 \allowbreak{}\NovelInline{I03329}{> 1}\allowbreak{}，也在中间补上连续自然数，使得任意一项都不超过前一项 \allowbreak{}\NovelInline{I03330}{+1}\allowbreak{}。这样，序列就变为了一个标准的 PrSS 序列。补上的项需要做好标记。''

上一步得到的 \allowbreak{}\NovelInline{I03331}{4, 3, 2, 4, 2}\allowbreak{} 在这里补项为： \NovelDisplay{M00643R01}{%
\color{red}{1, 2, 3}, 4, 3, 2, \color{red}{3}, 4, 2%
}

霍雨浩仔细观察着那些红色的标记项，喃喃道：``补上的项是为了让序列满足 PrSS 的合法性条件\ldots\ldots''

\textbf{第三步：展开。}

``按照 PrSS 规则展开序列。复制坏部的时候需要连标记一并复制，但坏根除外。''

上一步的结果在这里展开为： \NovelDisplay{M00644R01}{%
\color{red}{1, 2, 3}, 4, 3, 2, \color{red}{3}, 4, (\color{red}{1, 2, 3}, 4, 3, 2, \color{red}{3}, 4), (\color{red}{1, 2, 3}, 4, 3, 2, \color{red}{3}, 4), (\color{red}{1, 2, 3}, 4, 3, 2, \color{red}{3}, 4), \dots%
}

\textbf{第四步：还原。}

``去掉被标记的项，倒序回来，将逗号还原为空格和 \allowbreak{}\NovelInline{I03332}{-}\allowbreak{}。此时的结果就是原来的反射式要折叠的结构。''

上一步的展开结果还原为： \NovelDisplay{M00645R01}{%
\dots (4\ 2-3-4\ 1-)(4\ 2-3-4\ 1-)(4\ 2-3-4\ 1-)4\ 2-3-4%
}

即： \NovelDisplay{M00646R01}{%
(4\ 2-3-4\ 1-)^{a,b,c,\dots}\ 4\ 2-3-4%
}

梅天狸长老目光扫过众人，缓缓说道：``对照一下上面的枚举结果------''

\begin{quote}
\allowbreak{}\NovelInline{I03333}{2-4\ 2-3-4}\allowbreak{} 折叠任意深度的 \allowbreak{}\NovelInline{I03334}{(4\ 2-3-4\ 1-)^{a,b,c,\dots}\ 4\ 2-3-4}\allowbreak{}。
\end{quote}

``完全一致。''

广场上响起了此起彼伏的惊叹声。这套四步法则------替换、补项、展开、还原------就像一把万能钥匙，将那些看似混沌无序的反射式链条统统纳入了可计算的框架之中。

卢本韦长舒一口气，感慨道：``从替换到还原，四步之间环环相扣。PrSS 的坏部复制规则竟然能精确地捕捉到反射式中`折叠任意深度'的语义\ldots\ldots 这简直是数学与组合之间最优美的桥梁。''

asheep 长老的身影最后一次浮现在众人面前，他的声音已经变得极为遥远，仿佛来自宇宙的尽头：

``你们今日所学的，是从 \allowbreak{}\NovelInline{I03335}{\Pi_3 \text{ onto}}\allowbreak{} 到 \allowbreak{}\NovelInline{I03336}{\Pi_\omega}\allowbreak{} 的完整路径，以及 PrSS 与反射式之间的深层对应。记住这套法则，它是你们未来在序数深渊中前行的指南针。''

``而 \allowbreak{}\NovelInline{I03337}{\Pi_\omega}\allowbreak{} 之后\ldots\ldots'' 长老的声音渐渐消散在晨风中，``那便是稳定阶段的领域了。到了那里，\allowbreak{}\NovelInline{I03338}{real.}\allowbreak{} 与 \allowbreak{}\NovelInline{I03339}{psd.}\allowbreak{} 的区分将不再是细枝末节，而是决定生死的关键。好生准备吧。''

八名少年少女齐齐躬身，声音汇聚成一股洪流：

``弟子谨遵教诲！''

晨光洒满广场，将他们的身影镀上了一层金色的光辉。他们知道，从今日起，自己已不再是序数世界的旁观者，而是真正的行者。前方，\allowbreak{}\NovelInline{I03340}{\Pi_\omega}\allowbreak{} 的终点线隐约可见，而那条线的背后，是更加浩瀚无垠的未知。

\chapter{第五十七章：队长的柔情}\label{ux7b2cux4e94ux5341ux4e03ux7ae0ux961fux957fux7684ux67d4ux60c5}

往后数月，扽魂殿晨钟暮鼓不歇，八人日夜苦修，沉浸在反射序数的浩瀚道则之中。众人愈发惊觉，这套看似和初等序列、基础迭代一脉相承的序道体系，实则藏着天地鸿沟。初等序列的推演有迹可循、层层渐进，可高阶反射序数的全域归一、\allowbreak{}\NovelInline{I03341}{\Pi_\omega}\allowbreak{} 终极统合，早已超脱常规迭代逻辑，是维度与规则的终极跃迁。

\begin{center}\rule{0.5\linewidth}{0.5pt}\end{center}

蔡须坤端坐蒲团，周身序纹平稳流转，作为队长，她率先完成全域复盘，指尖轻抬，虚空自动铺开全新的序集推演线条，无需长老指引，自主接续上一轮的 \allowbreak{}\NovelInline{I03342}{\Pi_3}\allowbreak{} 核心架构。

``既然 \allowbreak{}\NovelInline{I03343}{3\ 2-3}\allowbreak{} 可完整折叠任意深度的 \allowbreak{}\NovelInline{I03344}{(2-)^{a,b,c,\dots} 3}\allowbreak{} 序列，那我们便顺着道途，自主补全后续所有层级推演。''

她话音落下，第一道自主推演的序集缓缓浮现，规整而严谨：

\NovelDisplay{M00647R01}{%
2\ 1-3\ 2-3%
}

``此为计数 \allowbreak{}\NovelInline{I03345}{3\ 2-3}\allowbreak{} 序数函数的全体容许点。''蔡须坤轻声解读，语气笃定无错，``容许点定阶、不动点固锚、\allowbreak{}\NovelInline{I03346}{K}\allowbreak{} 点破界，这是底层不变的铁律，适配所有高阶反射体系。''

紧随其后，霍雨浩心神联动，接续推演，第二重序集次第绽放：

\NovelDisplay{M00648R01}{%
2-2\ 1-3\ 2-3%
}

``马洛点层级迭代，折叠所有 \allowbreak{}\NovelInline{I03347}{2\ 1-3\ 2-3}\allowbreak{} 的容许点迭代操作，进一步收紧 \allowbreak{}\NovelInline{I03348}{3\ 2-3}\allowbreak{} 体系的秩序边界，让全域迭代更稳固。''

卢本韦天生对极致嵌套的序道敏感，瞬间踏出更高一层推演，虚空符号骤变：

\NovelDisplay{M00649R01}{%
3\ 1-3\ 2-3%
}

``\allowbreak{}\NovelInline{I03349}{K}\allowbreak{} 点成型！''她声音带着一丝振奋，``这是计数 \allowbreak{}\NovelInline{I03350}{3\ 2-3}\allowbreak{} 函数的递归弱紧致点，彻底超脱普通容许、马洛层级，让 \allowbreak{}\NovelInline{I03351}{3\ 2-3}\allowbreak{} 具备自我折叠、自我升华的本源属性。''

小队八人默契联动，无需言语，各自执掌一层推演维度，层层递进、环环相扣，无数繁杂的高阶序集在扽魂殿虚空接连铺展，无人出错、无人卡顿。

\NovelDisplay{M00650R01}{%
3\ 1-3\ 1-3\ 2-3, \quad (3\ 1-)^\omega 3\ 2-3, \dots%
}

超限嵌套层层铺开，从有限迭代直达极限超限。

\begin{center}\rule{0.5\linewidth}{0.5pt}\end{center}

推演至关键节点，蔡须坤抬手止住众人节奏，点出整片体系最核心的易错壁垒，声音清晰响彻殿宇：

\begin{quote}
\textbf{``诸位谨记一处核心误区！折叠任意深度 \allowbreak{}\NovelInline{I03352}{(3\ 1-)^{a,b,c,\dots} 2-3}\allowbreak{} 序列的序集，绝非 \allowbreak{}\NovelInline{I03353}{2-3\ 2-3}\allowbreak{}，而是 \allowbreak{}\NovelInline{I03354}{2-3\ 1-3\ 2-3}\allowbreak{}！''}
\end{quote}

``遵循 \allowbreak{}\NovelInline{I03355}{2-3}\allowbreak{} 折叠所有 \allowbreak{}\NovelInline{I03356}{(3\ 1-)}\allowbreak{} 迭代的底层规则，此序集是计数 \allowbreak{}\NovelInline{I03357}{3\ 2-3}\allowbreak{} 函数的 \allowbreak{}\NovelInline{I03358}{2-3}\allowbreak{} 点，是体系专属的折叠锚点，混淆二者，整个高阶嵌套体系会直接崩塌失真。''

众人纷纷颔首，将这一关键禁忌深烙道心，随即继续深度推演。

新一轮嵌套迭代尽数现世：

\NovelDisplay{M00651R01}{%
2-3\ 1-2-3\ 1-3\ 2-3 = (2-3\ 1-)^2 3\ 2-3%
}

二次迭代折叠彻底成型。

\NovelDisplay{M00652R01}{%
2-2-3\ 1-3\ 2-3%
}

统合所有 \allowbreak{}\NovelInline{I03359}{(2-3\ 1-)}\allowbreak{} 深层迭代。

\NovelDisplay{M00653R01}{%
2-2-2-3\ 1-3\ 2-3, \quad (2-)^\omega 3\ 1-3\ 2-3%
}

有限至超限的迭代链条完美闭环。

\begin{center}\rule{0.5\linewidth}{0.5pt}\end{center}

小舞凝视虚空新生的序集，瞬间洞悉全新大道特性，轻声感慨：``终于明白交集序数的真意了！\allowbreak{}\NovelInline{I03360}{3\ 2-3\ 1-3\ 2-3}\allowbreak{} 折叠所有 \allowbreak{}\NovelInline{I03361}{(2-)^{a,b,c,\dots} 3\ 1-3\ 2-3}\allowbreak{} 序列，它是计数 \allowbreak{}\NovelInline{I03362}{3\ 2-3}\allowbreak{} 函数的 \allowbreak{}\NovelInline{I03363}{3\ 2-3}\allowbreak{} 点，是我们首次真正意义上，在'XX点'体系中完成集合交集的高阶运用。''

``没错。''蔡须坤接续推演，眸光愈发深邃，``也正因如此，\allowbreak{}\NovelInline{I03364}{2-3\ 2-3}\allowbreak{} 才应运而生，专门折叠'求取 \allowbreak{}\NovelInline{I03365}{3\ 2-3}\allowbreak{} 点'的高阶操作，收纳这一层全新的迭代规则。''

推演并未止步，众人顺着这套折叠再生、迭代再生的闭环逻辑，持续深挖：

\begin{itemize}
\tightlist
\item
  二次求取 \allowbreak{}\NovelInline{I03366}{3\ 2-3}\allowbreak{} 点，诞生 \allowbreak{}\NovelInline{I03367}{3\ 1-3\ 2-3\ 1-3\ 2-3}\allowbreak{}；
\item
  对应折叠序集 \allowbreak{}\NovelInline{I03368}{2-3\ 1-3\ 2-3\ 1-3\ 2-3}\allowbreak{}、\allowbreak{}\NovelInline{I03369}{2-2-3\ 1-3\ 2-3\ 1-3\ 2-3}\allowbreak{} 次第成型；
\item
  最终固化出 \allowbreak{}\NovelInline{I03370}{(3\ 2-3\ 1-)^2 3\ 2-3}\allowbreak{} 的二次超限迭代结构。
\end{itemize}

层层折叠、层层再生，看似繁杂繁复，实则每一步都严丝合缝、循规蹈矩。八人此刻彻底吃透了这套近乎偏执的迭代逻辑：不断对核心表达式施加 \allowbreak{}\NovelInline{I03371}{2-}\allowbreak{} 反射，以带 \allowbreak{}\NovelInline{I03372}{3}\allowbreak{} 前缀的高阶序集完成折叠，再对折叠后的新体系继续迭代，生生造出无限延伸的超长序链------

\NovelDisplay{M00654R01}{%
\dots 3\ 2-3\ 2-3\ 2-3%
}

``每一次基础 \allowbreak{}\NovelInline{I03373}{2-}\allowbreak{} 反射，都伴随一次自我点位的折叠收纳。''梅天狸静静旁观，微微颔首，眼中满是赞许，``折叠再折叠，迭代再迭代，这便是 \allowbreak{}\NovelInline{I03374}{\Pi_3}\allowbreak{} 超反射最霸道的攀升之路。''

\begin{center}\rule{0.5\linewidth}{0.5pt}\end{center}

众人顺势将这条有限长的嵌套链条，延伸至超限维度，自主推演出新的终极极限序集：

\NovelDisplay{M00655R01}{%
(3\ 2-)^\omega 3%
}

随后再度升维，固化出更高阶的封顶结构：

\NovelDisplay{M00656R01}{%
(3\ 2-)^{1,0} 3%
}

至此，\allowbreak{}\NovelInline{I03375}{3 \text{ onto}}\allowbreak{} 的终极大门，被八人彻底亲手推开。

\begin{center}\rule{0.5\linewidth}{0.5pt}\end{center}

Asheep 长老适时开口，考验众人的融会贯通之力：``既已勘破 \allowbreak{}\NovelInline{I03376}{3 \text{ onto}}\allowbreak{}，那你们来说说，高阶 \allowbreak{}\NovelInline{I03377}{\text{onto}}\allowbreak{} 运算的通用本质是什么？''

这一次，无需队长率先开口，全员心神共鸣，异口同声，道破通用铁律：

\begin{quote}
\textbf{``2 onto \allowbreak{}\NovelInline{I03378}{X}\allowbreak{}，折叠 \allowbreak{}\NovelInline{I03379}{X \cap 1-}\allowbreak{} 所有操作；''}

\textbf{``3 onto \allowbreak{}\NovelInline{I03380}{X}\allowbreak{}，折叠 \allowbreak{}\NovelInline{I03381}{X \cap 2-}\allowbreak{} 所有操作；''}

\textbf{``推广至全域通用：\allowbreak{}\NovelInline{I03382}{k}\allowbreak{} onto \allowbreak{}\NovelInline{I03383}{X}\allowbreak{}，即为对 \allowbreak{}\NovelInline{I03384}{X \cap (k-1)-}\allowbreak{} 全体迭代操作的终极折叠！''}
\end{quote}

一字不差、分毫不差。

两大长老眸光骤亮，果糕小队已然不是单纯记忆公式，而是真正掌握了序道攀升的底层源代码，自此可自主推演所有高阶 \allowbreak{}\NovelInline{I03385}{\Pi}\allowbreak{} 层级，直通 \allowbreak{}\NovelInline{I03386}{\Pi_\omega}\allowbreak{} 终极领域。

\begin{center}\rule{0.5\linewidth}{0.5pt}\end{center}

得到通用公理加持，八人开启全速自主枚举，从 \allowbreak{}\NovelInline{I03387}{\Pi_3}\allowbreak{} 圆满突破，一路向上，无人停顿、无人滞涩。

\textbf{首先是三阶反射的终极闭环：}

\NovelDisplay{M00657R01}{%
3-3 \text{ 折叠全部 } (3\ 2-)^{a,b,c,\dots} 3 \text{ 迭代}%
}

\NovelDisplay{M00658R01}{%
3\ 1-3-3, \quad (3\ 1-)^2 3-3 \text{ 层层递进，固化 } K \text{ 点层级}%
}

\NovelDisplay{M00659R01}{%
2-3\ 1-3-3, \quad 2-2-3\ 1-3-3 \text{ 完成逐层折叠}%
}

\NovelDisplay{M00660R01}{%
3\ 2-3\ 1-3-3 \text{ 收纳二阶全域迭代}%
}

后续 \allowbreak{}\NovelInline{I03388}{1-}\allowbreak{} 嵌套、\allowbreak{}\NovelInline{I03389}{2-}\allowbreak{} 嵌套、\allowbreak{}\NovelInline{I03390}{3-}\allowbreak{} 嵌套层层叠加，迭代出无数细分序集，最终成型 \allowbreak{}\NovelInline{I03391}{(3-3\ 1-)^2 3-3}\allowbreak{} 的二次超限结构。

小队顺势推演三阶反射的终极形态：

\NovelDisplay{M00661R01}{%
2-3-3, \quad 2-2-3-3, \quad 3\ 2-3-3 \text{ 逐级锁死底层秩序}%
}

\NovelDisplay{M00662R01}{%
2-3\ 2-3-3, \quad 3-3\ 2-3-3 \text{ 完成跨层折叠}%
}

\NovelDisplay{M00663R01}{%
3-3-3, \quad 3-3-3-3 \text{ 持续堆叠，直至抵达三阶极限：}%
}

\NovelDisplay{M00664R01}{%
(3-)^\omega = \sup\{3,\ 3-3,\ 3-3-3,\ \dots\}%
}

最后固化三阶终极封顶序集 \allowbreak{}\NovelInline{I03392}{(3-)^{1,0}}\allowbreak{}，圆满整个三阶反射体系。

\begin{center}\rule{0.5\linewidth}{0.5pt}\end{center}

\textbf{三阶圆满，四阶顺势开启。}

蔡须坤指尖一点，全新的 \allowbreak{}\NovelInline{I03393}{\Pi_4}\allowbreak{} 序数符号 \allowbreak{}\NovelInline{I03394}{\kappa}\allowbreak{} 现世，定义出四阶反射基底：

\NovelDisplay{M00665R01}{%
4 = \{\kappa,\ \kappa_2,\ \kappa_3,\ \dots,\ \kappa_{\omega+1},\ \dots\}%
}

众人默契接续，逐层推演四阶全套体系：

\begin{itemize}
\tightlist
\item
  \allowbreak{}\NovelInline{I03395}{(1-)^{1,0}\ 4}\allowbreak{} 对应 \allowbreak{}\NovelInline{I03396}{\kappa \text{FP}}\allowbreak{} 不动点全集；
\item
  \allowbreak{}\NovelInline{I03397}{2\ 1-4}\allowbreak{}、\allowbreak{}\NovelInline{I03398}{3\ 1-4}\allowbreak{}、\allowbreak{}\NovelInline{I03399}{4\ 1-4}\allowbreak{} 依次对应 \allowbreak{}\NovelInline{I03400}{\kappa}\allowbreak{} 序列的容许点、\allowbreak{}\NovelInline{I03401}{K}\allowbreak{} 点、\allowbreak{}\NovelInline{I03402}{\kappa}\allowbreak{} 点；
\item
  \allowbreak{}\NovelInline{I03403}{2-4}\allowbreak{}、\allowbreak{}\NovelInline{I03404}{3\ 2-4}\allowbreak{}、\allowbreak{}\NovelInline{I03405}{2-3\ 2-4}\allowbreak{}、\allowbreak{}\NovelInline{I03406}{3-3\ 2-4}\allowbreak{}、\allowbreak{}\NovelInline{I03407}{4\ 2-4}\allowbreak{} 逐层折叠对应层级迭代；
\item
  \allowbreak{}\NovelInline{I03408}{2-4\ 2-4}\allowbreak{}、\allowbreak{}\NovelInline{I03409}{3-4}\allowbreak{}、\allowbreak{}\NovelInline{I03410}{2-3-4}\allowbreak{}、\allowbreak{}\NovelInline{I03411}{2-2-3-4}\allowbreak{}\ldots\ldots 无数细分序集井然落地；
\item
  最终以 \allowbreak{}\NovelInline{I03412}{4-4}\allowbreak{}、\allowbreak{}\NovelInline{I03413}{4-4-4}\allowbreak{} 封顶，圆满 \allowbreak{}\NovelInline{I03414}{\Pi_4}\allowbreak{} 完整体系。
\end{itemize}

\begin{center}\rule{0.5\linewidth}{0.5pt}\end{center}

道途已通，攀升再无桎梏。

五阶、六阶\ldots\ldots 有限 \allowbreak{}\NovelInline{I03415}{\Pi}\allowbreak{} 层级层层突破，迭代规律完全通用，\allowbreak{}\NovelInline{I03416}{\text{onto}}\allowbreak{} 运算铁律全域适配。小队八人如同复制升华，以统一的底层逻辑，极速推演所有高阶有限反射体系，一路畅通无阻，最终望见所有有限 \allowbreak{}\NovelInline{I03417}{\Pi}\allowbreak{} 层级的终极归宿------

\NovelDisplay{M00666R01}{%
psd.\ \Pi_\omega \text{ 全域反射}%
}

漫天无尽序链缓缓收敛，最终归一。

\begin{center}\rule{0.5\linewidth}{0.5pt}\end{center}

推演至终点，众人并未止步，而是主动回过头来，以全套高阶序集反向验证 PrSS 序列规律，完成最后的闭环悟道。

蔡须坤随意选取复杂反射式 \allowbreak{}\NovelInline{I03418}{2-4\ 2-3-4}\allowbreak{}，当众复刻四步推演法则，为众人巩固终极捷径：

\begin{quote}
**``第一步，替换倒置------** \allowbreak{}\NovelInline{I03419}{2-4\ 2-3-4}\allowbreak{} 转为序列 \allowbreak{}\NovelInline{I03420}{4, 3, 2, 4, 2}\allowbreak{}；''

**``第二步，补全断层------** 前置补 \allowbreak{}\NovelInline{I03421}{1, 2, 3}\allowbreak{}，中间补 \allowbreak{}\NovelInline{I03422}{3}\allowbreak{}，得到标准 PrSS 序列 \allowbreak{}\NovelInline{I03423}{1, 2, 3, 4, 3, 2, 3, 4, 2}\allowbreak{}；''

**``第三步，迭代展开------** 复制坏部、保留标记，无限延伸序列结构；''

**``第四步，逆向还原------** 去标记、倒置、复现符号，最终锁定折叠结构为 \allowbreak{}\NovelInline{I03424}{(4\ 2-3-4\ 1-)^{a,b,c,\dots} 4\ 2-3-4}\allowbreak{}。''
\end{quote}

结果与序道本质分毫不差，完美印证规律。

\begin{center}\rule{0.5\linewidth}{0.5pt}\end{center}

霍雨浩缓缓开口，道出全员最终的道心彻悟：``原来所有繁杂的反射枚举、嵌套折叠、层级攀升，从来不是杂乱的符号堆砌。PrSS 是表象捷径，\allowbreak{}\NovelInline{I03425}{k}\allowbreak{} 阶 \allowbreak{}\NovelInline{I03426}{\text{onto}}\allowbreak{} 是底层公理，点位迭代是攀升骨架，三者合一，便是整个序数反射体系的终极真理。''

``二阶筑基，三阶破壁，四阶拓界，全域归宗。''蔡须坤起身而立，周身序道威压浑然一体，``我们从前是依序悟道，如今已然可以自主造序、自主推演、自主封顶。''

八人同时起身，八道凝练至极的序道道韵交织相融，贯通整座扽魂殿。

他们彻底吃透了从多元 \allowbreak{}\NovelInline{I03427}{M}\allowbreak{} 函数、\allowbreak{}\NovelInline{I03428}{M\text{FP}}\allowbreak{} 迭代、\allowbreak{}\NovelInline{I03429}{\Pi_3}\allowbreak{} 超反射、\allowbreak{}\NovelInline{I03430}{k}\allowbreak{} 阶 \allowbreak{}\NovelInline{I03431}{\text{onto}}\allowbreak{} 通用公理，直至 \allowbreak{}\NovelInline{I03432}{\Pi_\omega}\allowbreak{} 全域反射的全部内容，更掌握了 PrSS 这套可以解析一切反射式的终极密钥。

\textbf{大道万千，折叠为一。}

此刻的果糕小队，才算真正握牢了反射序数的无上威力，静待下一步，全域超限归一，登临有限序道之巅。

\begin{center}\rule{0.5\linewidth}{0.5pt}\end{center}

可道途登顶的最后一步，从来最为艰难。

队长蔡须坤，卡在了最关键的临门一步。

她的道心最稳、推演最细，从二阶筑基到三阶破壁，每一层序道规则都打磨得极致圆满，可越是通透，越能感知 \allowbreak{}\NovelInline{I03433}{\Pi_\omega}\allowbreak{} 全域反射的恐怖壁垒。那不是单一序数的强度桎梏，而是统合所有有限反射、收纳万千维度规则、归一全部迭代折叠的终极大道枷锁。

无数个日夜，她端坐序阵中央，一遍遍推演 PrSS 终极规律，一遍遍复刻 \allowbreak{}\NovelInline{I03434}{k}\allowbreak{} 阶 \allowbreak{}\NovelInline{I03435}{\text{onto}}\allowbreak{} 通用公理，从细碎序集到全域链条，反复复盘、层层拆解，可每当即将触碰 \allowbreak{}\NovelInline{I03436}{\Pi_\omega}\allowbreak{} 的边界，便会被一股无形的力量弹回，道心震颤、序纹紊乱，数次推演尽数崩塌。

瓶颈如天堑，横亘身前，寸步难行。

``为何\ldots\ldots 规则尽数通透，推演毫无错漏，却始终无法归一全域？''

蔡须坤缓缓睁眼，眼底布满疲惫，周身流转的序光黯淡无光。数月苦修无果，无数次突破失败，哪怕她心性坚韧，此刻也难免生出无力之感。她执掌全队道途，背负众人登顶序道的期许，可自己却卡在终极门槛，无法带领小队踏足巅峰。

殿内晚风轻拂，序道余韵微凉，孤寂与焦灼缠上心头，压得她几乎喘不过气。

\begin{center}\rule{0.5\linewidth}{0.5pt}\end{center}

就在她道心浮动、几近溃散的刹那，一缕极淡、却无比熟悉的温暖意识，悄然从她灵魂最深处苏醒。

没有天地异象，没有序光涌动，唯有一抹温柔且坚定的灵魂气息，缓缓包裹住她躁动的道心，熟悉到让她身躯骤然一僵。

一道温润的神念，直接响彻她的识海，带着久别重逢的宽慰，亦带着一如既往的笃定：

\begin{quote}
\emph{``小坤，道途登顶，从不是一蹴而就。通透规则是根基，勘破本心是天梯，不必急，不必慌。''}
\end{quote}

是牢大。

是早已陨落、舍身护佑众人离去的前任队长。

岁月流转，生死相隔，已是数月未见。自那场绝境追杀、那场惨烈献祭之后，牢大的神魂气息便彻底沉寂，消散于天地之间，只留下一缕残魂，隐秘寄居在蔡须坤的灵魂深处，无声守护，从未现世。

蔡须坤僵坐原地，心脏骤然紧缩，眼眶瞬间泛红。无数压抑的思念、过往的绝境、惨烈的别离，在这一刻尽数翻涌而上。

``牢大\ldots\ldots？''

她在心底轻声呢喃，声音带着难以抑制的颤抖。

那道温柔的神念再度响起，带着浅浅的笑意，温柔抚平她所有的焦躁与自我怀疑：

\begin{quote}
\emph{``我一直在。看着你们步步精进、踏破层层桎梏，看着你扛起队长之责，护住全队、稳步前行，做得很好。''}

\emph{``瓶颈不是桎梏，是道心最后的淬炼。你看懂了万千序道规则，却唯独没看懂，反射序数的终极，从来不是推演符号，而是舍己归一、承载万序的本心。你怕突破失败、怕辜负众人、怕前路无途，这份执念，困住了你的 \allowbreak{}\NovelInline{I03437}{\Pi_\omega}\allowbreak{} 大道。''}
\end{quote}

寥寥数语，如醍醐灌顶，瞬间击穿蔡须坤心中所有的迷茫。

积压数月的压力、隐忍已久的悲痛、生死相隔的遗憾，在这一刻彻底决堤。滚烫的泪水夺眶而出，顺着脸颊滑落，滴落在身前序阵之上，漾开细碎的光纹。

数月苦修的疲惫、绝境求生的阴影、扛起重担的孤勇，尽数化作泪水宣泄而出。她以为自己早已习惯别离、扛起责任，却在听到这道熟悉的声音时，瞬间变回了当年那个被前辈护在身后的少年。

``我好想你\ldots\ldots{}''她低声哽咽，字字酸涩。

灵魂深处的暖意愈发浓厚，温柔包容着她所有的脆弱，无声慰藉，静静陪伴。没有多余的言语，却胜过万千开导，让她纷乱的道心渐渐平复、归于澄澈。

\begin{center}\rule{0.5\linewidth}{0.5pt}\end{center}

待心绪缓缓平复，蔡须坤抬手拭去泪水，眼底的迷茫尽数褪去，取而代之的是极致的通透与坚定。

她终于彻底明白 PrSS 规律的终极内核，读懂了 \allowbreak{}\NovelInline{I03438}{k}\allowbreak{} 阶 \allowbreak{}\NovelInline{I03439}{\text{onto}}\allowbreak{} 折叠的至高真谛。

所谓 \allowbreak{}\NovelInline{I03440}{\Pi_\omega}\allowbreak{} 全域反射，是收纳所有低维秩序、承载万千迭代规则、以一己之身归一万序的终极境界。\textbf{折叠万物者，必先承载万物；统合万序者，必先舍弃自我。}

这一刻，她不再执着于''突破境界''，不再执念于''登顶巅峰''，所有的私心、执念、桎梏尽数消散。

道心空明，万序归宗。

\textbf{轰------！}

无形道韵骤然炸开，整座扽魂殿的序链尽数震颤，漫天繁杂的有限 \allowbreak{}\NovelInline{I03441}{\Pi}\allowbreak{} 层级序列层层跃升、飞速归一。无数嵌套、折叠、迭代的序道规则，尽数涌入蔡须坤的身躯。

\allowbreak{}\NovelInline{I03442}{psd.\ \Pi_3}\allowbreak{}、\allowbreak{}\NovelInline{I03443}{psd.\ \Pi_4}\allowbreak{}、\allowbreak{}\NovelInline{I03444}{psd.\ \Pi_5}\allowbreak{}\ldots\ldots 所有有限序数层级彻底消融，万千序道归一原点，一道横贯天地的终极序界轰然成型！

\textbf{\allowbreak{}\NovelInline{I03445}{\Pi_\omega}\allowbreak{} 全域反射，成！}

极致浩瀚、稳固纯粹的全域序道威压席卷整座殿宇，凌驾所有低维反射、所有迭代嵌套，是当前有限序道体系的至高尽头。

\begin{center}\rule{0.5\linewidth}{0.5pt}\end{center}

突破刹那，过往被追杀的绝境画面，不受控制地涌入蔡须坤脑海。

当年强敌环伺、死局无解，全队濒临覆灭，是牢大挺身而出，孤身挡下所有攻势，爆发出远超自身境界、近乎逆天的恐怖力量，硬生生撕开生路，将所有人护离绝境。彼时的她，只知前辈拼死相救、燃尽一切，却始终不懂，那般碾压强敌、破碎天地的力量，究竟源自何等境界。

直至此刻登临 \allowbreak{}\NovelInline{I03446}{\Pi_\omega}\allowbreak{}，她才后知后觉，心神巨震。

\textbf{原来当年牢大拼死爆发的力量，赫然正是 \allowbreak{}\NovelInline{I03447}{\Pi_\omega}\allowbreak{} 全域反射的终极之力！}

可她比任何人都清楚，\allowbreak{}\NovelInline{I03448}{\Pi_\omega}\allowbreak{} 境界的力量何等霸道狂暴，其生命燃烧的反噬更是恐怖至极。

规则铁律历历在目：强行催动全域反射之力，无需片刻，仅仅一分钟，便可彻底枯竭所有生命本源、神魂根基，无一丝残存、无半点退路。

蔡须坤身形微颤，眼眶再度湿润。

她终于彻底读懂了那场舍身献祭的真相。

牢大当年明明知晓刚刚突破 \allowbreak{}\NovelInline{I03449}{\text{I}}\allowbreak{} 序境界的他，催动 \allowbreak{}\NovelInline{I03450}{\Pi_\omega}\allowbreak{} 之力便是必死之局，燃烧生命、修为尽毁灭，却依旧毫不犹豫、毅然爆发。

他没有给自己留半分退路，没有留一丝生机，以彻底陨落为代价，以一己全域大道，换全队安然无恙、平安脱身。

何为队长？何为守护？何为舍身取义？

这一刻，蔡须坤彻彻底底，刻入骨髓，融入道心。

\begin{center}\rule{0.5\linewidth}{0.5pt}\end{center}

灵魂深处，那缕温暖的神魂气息轻轻震荡，似是欣慰，似是释然，随后缓缓沉寂，归于静谧，继续默默守护在她灵魂深处，无声相伴，岁岁无忧。

蔡须坤伫立序阵中央，周身 \allowbreak{}\NovelInline{I03451}{\Pi_\omega}\allowbreak{} 全域序光流转不息，澄澈眼底藏着热泪，亦藏着坚定。

前人燃尽己身，铺就前路万丈。

往后，她便接过重担，携牢大遗志，护全队登顶，踏遍序道巅峰，不负牺牲，不负初心，不负这场以命换来的万丈前程。

\begin{center}\rule{0.5\linewidth}{0.5pt}\end{center}

数月未见，牢大那道残魂神念依旧温润如初，稳稳熨帖住蔡须坤濒临溃散的道心。

一句提点，破开万般桎梏。

蔡须坤凝滞许久的序道瞬间轰鸣复苏，原本无数次推演崩塌、死死困住她的 \allowbreak{}\NovelInline{I03452}{\Pi_\omega}\allowbreak{} 天堑壁垒，应声碎裂、彻底消融。万千散落的有限 \allowbreak{}\NovelInline{I03453}{\Pi}\allowbreak{} 层级序链、无尽折叠迭代的反射规则、贯通全域的 PrSS 大道脉络，尽数朝着她的身躯汇聚归一。

璀璨至极的纯白序光自她周身炸裂而出，横扫整座扽魂殿，凌驾所有低阶反射序数的威压徐徐绽放，浩瀚、稳固且极致圆满。

\textbf{\allowbreak{}\NovelInline{I03454}{\Pi_\omega}\allowbreak{} 全域反射境界，彻底突破！}

道心通透的刹那，当年绝境追杀、牢大燃命护队的破碎画面汹涌灌入脑海，\allowbreak{}\NovelInline{I03455}{\Pi_\omega}\allowbreak{} 境界的终极奥义与昔日那场悲壮的献祭完美重合。她终于彻彻底底看清，那一日牢大毫无保留爆发的逆天之力，正是这归一万序、燃尽本源的全域反射真谛，是以一分钟生命枯竭、神魂俱灭的必死结局，为全队换来的一线生机。

\begin{center}\rule{0.5\linewidth}{0.5pt}\end{center}

万般情绪翻涌交织，突破的狂喜、勘破真相的震撼、缅怀逝者的悲痛、负重前行的酸涩，瞬间击溃了蔡须坤所有的坚韧伪装。

方才还端坐如松、沉稳执掌全队的队长，此刻伫立序阵中央，任由滚烫的泪水无声滑落，一滴接着一滴，砸在冰凉的殿宇地面，漾开细碎的序纹涟漪。

全场骤然陷入一片死寂。

所有人望着序光璀璨中默默垂泪的蔡须坤，尽数默然低头，心底五味杂陈。

在全队所有人的印象里，蔡须坤永远是最沉稳可靠的队长。她修为冠绝全队，心性坚韧果决，带队闯绝境、破桎梏、渡难关，无论遭遇多大的危机、多严苛的瓶颈，她都始终冷静自持、岿然不动，素来眉眼清冷、神色坚定，从未有过半分脆弱。身为少女，她却早早扛起全队的重担，藏起所有柔软与私情，硬生生撑起了整支果糕小队，众人从未见过她落下一滴眼泪。

可此刻，这位从不轻易落泪的队长，立于全域序道之巅，褪去了所有坚硬的伪装。滚烫的泪水无声滚落，没有嘶吼，没有哽咽，唯有无声的悲戚，却比任何痛哭都更让人心疼。

这不是突破的悲苦，是时隔数月，终于彻悟牢大当年舍身真相、直面那场生死别离的绵长恸念。昔日护佑众人的前辈燃尽一切、不留退路的决绝，此刻尽数化作沉甸甸的念想，压在了她的心头。

\begin{center}\rule{0.5\linewidth}{0.5pt}\end{center}

八人齐齐垂首，无人言语，无人敢打破这份肃穆的沉静。殿内只剩流转不息的 \allowbreak{}\NovelInline{I03456}{\Pi_\omega}\allowbreak{} 道韵，衬得这份缅怀与酸涩，愈发浓重。

八人之中，卢本韦向来心思敏感、共情最甚，看着素来坚韧强大的队长独自承受着突破的重压与生死的遗憾，鼻头骤然一酸。

方才为队友突破圆满大道的欣喜荡然无存，温热的泪水毫无预兆地涌上眼眶，顺着脸颊簌簌落下，转瞬便湿了衣襟。他没有放声痛哭，只是死死抿着唇，肩膀微微颤抖，沉默地看着蔡须坤，眼底满是心疼与酸涩。

一旁的诺当最先察觉卢本韦的失态，见平日里开朗跳脱、最爱嬉闹的少女此刻默默垂泪，神色落寞又酸涩，心头顿时一软。他迈步上前，轻轻拍了拍卢本韦的后背，动作温柔又稳妥，低声安抚道：``别难过，队长不是悲戚，是释然，也是新生。牢大前辈的牺牲从不是遗憾，是护我们走到如今的前路。''

唐三亦缓步走近，眸光温和沉静，抬手轻轻抚平卢本韦微颤的肩头，声音沉稳有力，带着让人安心的力量：``道途登顶必经悲欢，队长扛起了前辈的遗志，也扛起了我们所有人的未来。泪水是感念，更是铭记，不必强忍。''

两人一柔一稳，默默陪伴在侧，一点点抚平卢本韦心底的酸涩与动容。

\begin{center}\rule{0.5\linewidth}{0.5pt}\end{center}

不远处，奥斯卡静静伫立，看着落泪的卢本韦，看着满殿的大道余韵与悲欢交织的画面，心底生出前所未有的坚定。他望着少女单薄的背影，在心中无声默念，字字铿锵、句句赤诚：

\begin{quote}
\emph{往后漫漫道途，我要永远保护好卢本韦，护她岁岁无忧，不受风霜、不遇危厄，倾尽所有，护她周全。}
\end{quote}

\begin{center}\rule{0.5\linewidth}{0.5pt}\end{center}

另一侧，霍雨浩凝视着沐浴序光、含泪挺立的蔡须坤，又转头看向身侧眉眼温柔的小舞，心中百感交集。他缓缓伸出手臂，轻轻将小舞拥入怀中，怀抱温暖而坚定。

待心绪沉淀，他轻声开口，语气郑重无比，许下横贯前路的诺言：

\begin{quote}
\emph{``小舞，也是给全队所有人。以后无论前路多险、困境多难，但凡遇到分毫危险，我都会拼尽全力保住所有人。哪怕我并非队长，亦与全队同进退、共生死，不离不弃，至死不渝。''}
\end{quote}

小舞轻轻靠在他的胸膛，静静点头，眼底满是信任与笃定。

\begin{center}\rule{0.5\linewidth}{0.5pt}\end{center}

殿内序光依旧璀璨，\allowbreak{}\NovelInline{I03457}{\Pi_\omega}\allowbreak{} 的全域道韵萦绕不绝。

泪水是感念过往的赤诚，坚守是奔赴未来的初心。一场极致的境界突破，终究让全员道心愈发凝练，情义愈发深重。前人燃灯铺路，后人执灯前行，果糕小队的序道征途，自此不负牺牲，不负初心，不负彼此。

\chapter{第五十八章：多逗号系统和反射模式的对应}\label{ux7b2cux4e94ux5341ux516bux7ae0ux591aux9017ux53f7ux7cfbux7edfux548cux53cdux5c04ux6a21ux5f0fux7684ux5bf9ux5e94}

一众魂师研习反射法则之后，心中依旧满是困惑，诸多晦涩的序数规则始终难以吃透。见状，扽魂殿圣女千仞雪起身解惑。她不仅是殿内天赋卓绝的圣女，更是ddfg长老的亲孙女，修行之路与众人截然不同，修的是一套极为特殊的多逗号系统。

扽魂殿多数长老皆是test\_alpha0的徒子徒孙，曾执掌圣女之位的梅天狸长老也位列其中。ddfg长老与梅天狸长老自幼相伴长大，情谊深厚。千仞雪的年岁比蔡须坤还要稍长几岁，众人早年便在魂师大赛中有过一面之缘，此番故地重逢，又见同为修行同道，亲切感油然而生。此番与果糕小队众人相聚探讨修为，更是格外投缘。

千仞雪眉眼淡然，缓缓道出自己的修行困惑与感悟：

\begin{quote}
``我自幼跟随爷爷修习多逗号系统，可始终难以将这套体系与反射序数完美对应。这套规则听着简单，实则门槛极高，当迭代序数 \allowbreak{}\NovelInline{I03458}{\varphi(1@(1@(1@\cdots(1,0)\cdots)))}\allowbreak{} 攀升至 \allowbreak{}\NovelInline{I03459}{\psi(\varepsilon_{\Omega+1})}\allowbreak{} 层级后，便会衍生出双逗号规则，同时完成单逗号折叠。最简单的例子便是 \allowbreak{}\NovelInline{I03460}{\varphi(1,,0)=\Omega}\allowbreak{}。''
\end{quote}

蔡须坤闻言立刻追问：``雪姐，如此说来，双逗号的序数层级，是凌驾于所有单逗号序数之上的吗？''

``正是如此。''千仞雪颔首确认，``但 \allowbreak{}\NovelInline{I03461}{\varphi(1,\varphi(1,,0)+1)}\allowbreak{} 便已然对应 \allowbreak{}\NovelInline{I03462}{\varepsilon_{\Omega+1}}\allowbreak{}。我反复推演得出，\allowbreak{}\NovelInline{I03463}{\varphi(1,,1)}\allowbreak{} 即是 \allowbreak{}\NovelInline{I03464}{\Omega_2}\allowbreak{}，只因它的层级超越了所有依托 \allowbreak{}\NovelInline{I03465}{\varphi(1,,0)}\allowbreak{} 构建的单逗号序数，这个推演结果应当无误。''

唐三眸光一亮，瞬间抓住关键：``照这个规则推导，\allowbreak{}\NovelInline{I03466}{\psi(\varphi(1,,\omega))}\allowbreak{} 便是BO层级，对吗？''

``我亦是这般判定。''千仞雪应声作答，``不过 \allowbreak{}\NovelInline{I03467}{\varphi(2,,0)}\allowbreak{} 的层级，远比众人想象的更为浩瀚。''

蔡兴坤细细梳理层级递进规律，开口补充道：``我看在 \allowbreak{}\NovelInline{I03468}{\varphi(2,,0)}\allowbreak{} 之前，尚且存在 \allowbreak{}\NovelInline{I03469}{\varphi(1,,(1,0))}\allowbreak{} 这一层级。''

诺当瞬间顿悟，脱口而出：``原来如此，这不就是OFP规则吗！那 \allowbreak{}\NovelInline{I03470}{\varphi(1,,(1,1))}\allowbreak{} 对应的是什么层级？''

千仞雪耐心解答：``对应 \allowbreak{}\NovelInline{I03471}{\Omega_{\psi_I(0)+1}}\allowbreak{}。这种逐级进位的推演方式，便是我修习的弱φ模式。''

奥斯卡疑惑发问：``雪姐，序数的强度变化，不会影响这套基础推演规则吧？''

``不会。''千仞雪轻轻摇头，举例佐证，``就比如 \allowbreak{}\NovelInline{I03472}{\varphi(1,,(2,0))}\allowbreak{} 仅仅对应 \allowbreak{}\NovelInline{I03473}{\psi_I(1)}\allowbreak{}。我知晓你们一众修习的是MOCF体系，对吗？''

蔡须坤点头应声：``没错，我们修行的正是MOCF。''

``为了方便你们融会贯通，我将弱φ模式与BOCF的核心规则拆解开来，与大家共享。''千仞雪缓缓道出两套体系的核心奥义，``BOCF体系中，核心规则为 \allowbreak{}\NovelInline{I03474}{\psi_B(\Omega)=\varepsilon_0}\allowbreak{}、\allowbreak{}\NovelInline{I03475}{\psi_B(\#+1)=\psi(\#)\cdot\omega}\allowbreak{}，其余基础规则不变。其中 \allowbreak{}\NovelInline{I03476}{\psi_B(\Omega^2)=\psi_M(\Omega)}\allowbreak{}，更进一步的 \allowbreak{}\NovelInline{I03477}{\psi_B(\Omega^2\cdot n)=\psi_M(\Omega\cdot n)=\zeta_n}\allowbreak{}。''

``而在弱φ模式中，我们可以简化规则辅助理解，令 \allowbreak{}\NovelInline{I03478}{\varphi(1,0)+1=\varphi(1,1)}\allowbreak{}。依照这套规则推演，\allowbreak{}\NovelInline{I03479}{\psi_I(I)}\allowbreak{} 便已然对应 \allowbreak{}\NovelInline{I03480}{\varphi(1,,(1,0,0))}\allowbreak{}。''

蔡须坤顺着规则推演，眼前豁然开朗：``那若是令 \allowbreak{}\NovelInline{I03481}{\Omega_{\psi_I(0)+1}=\varphi(1,,\varphi(1,,(1,0))+1)}\allowbreak{}，便可直接推出 \allowbreak{}\NovelInline{I03482}{\varphi(1,1)=\psi_I(1)}\allowbreak{}，这个对应方式完全可行！''

``这个推演思路十分精妙，完全契合体系规则。''千仞雪赞许道。

唐三神色凝重，感知到了更深层的层级壁垒：``如此看来，\allowbreak{}\NovelInline{I03483}{\varphi(2,,0)}\allowbreak{} 能够折叠所有基于单逗号构建的 \allowbreak{}\NovelInline{I03484}{\varphi(1,,)}\allowbreak{} 序数，这已然触及了I序数的本源！按照 \allowbreak{}\NovelInline{I03485}{(1\text{-})^{(1,0)}\,2}\allowbreak{} 的规则定义，\allowbreak{}\NovelInline{I03486}{\varphi(1,,(1,0))}\allowbreak{} 便是其核心载体，也就是说 \allowbreak{}\NovelInline{I03487}{\Omega_{I+1}=\varphi(1,,\varphi(2,,0)+1)}\allowbreak{}。''

蔡兴坤接续推演，层层递进：``那 \allowbreak{}\NovelInline{I03488}{\varphi(1,,(1,\varphi(2,,0)+1))}\allowbreak{}，便是囊括所有层级的 \allowbreak{}\NovelInline{I03489}{\Omega_{\cdots\cdots\Omega_{I+1}}}\allowbreak{}，也就是 \allowbreak{}\NovelInline{I03490}{\psi_{I_2}(0)}\allowbreak{}！''

``这层级太过浩瀚可怖。''千仞雪微微惊叹，``照此推演，\allowbreak{}\NovelInline{I03491}{\varphi(1,,(1,\varphi(2,,0)\cdot 2))}\allowbreak{} 恐怕对应 \allowbreak{}\NovelInline{I03492}{\psi_{I_2}(I)}\allowbreak{}，而 \allowbreak{}\NovelInline{I03493}{\varphi(1,,(2,\varphi(2,,0)+1))}\allowbreak{} 便是 \allowbreak{}\NovelInline{I03494}{\psi_{I_2}(I_2)}\allowbreak{}。''

小舞听得认真，缓缓梳理出规律：``难怪如此，那 \allowbreak{}\NovelInline{I03495}{\varphi(2,,1)}\allowbreak{} 岂不是就是 \allowbreak{}\NovelInline{I03496}{I_2}\allowbreak{}？它能够折叠所有依托 \allowbreak{}\NovelInline{I03497}{\varphi(2,,0)}\allowbreak{} 构建的单逗号序数，这么说来 \allowbreak{}\NovelInline{I03498}{\varphi(1,,(1,0,\varphi(2,,0)+1))}\allowbreak{}，对应的便是 \allowbreak{}\NovelInline{I03499}{\psi_{I_2}(I_2^{I_2})}\allowbreak{}。''

千仞雪点头确认，总结出核心规律：``没错，\allowbreak{}\NovelInline{I03500}{\varphi(2,,1)}\allowbreak{} 正是 \allowbreak{}\NovelInline{I03501}{I_2}\allowbreak{}。我发现这套体系有一个通律：序数末尾的数值，永远对应第n阶层级。以此类推，\allowbreak{}\NovelInline{I03502}{\varphi(2,,(1,0))}\allowbreak{} 是所有满足 \allowbreak{}\NovelInline{I03503}{\alpha\to\varphi(2,,\alpha)}\allowbreak{} 的不动点，也就是IFP层级。不过这仅仅对应反射规则中 \allowbreak{}\NovelInline{I03504}{(1\text{-})^{(1,0)}\,2}\allowbreak{} 的1-2阶，尚未触及 \allowbreak{}\NovelInline{I03505}{2\;1\text{-}2\;1\text{-}2}\allowbreak{} 的更高层级，而 \allowbreak{}\NovelInline{I03506}{\varphi(3,,0)}\allowbreak{} 已然抵达 \allowbreak{}\NovelInline{I03507}{I(1,0)}\allowbreak{} 层级。''

诺当立刻跟进推演：``若这个规律成立，那 \allowbreak{}\NovelInline{I03508}{\varphi(2,,\varphi(3,,0)+1)}\allowbreak{} 便是 \allowbreak{}\NovelInline{I03509}{I_{I(1,0)}+1}\allowbreak{}！要知道 \allowbreak{}\NovelInline{I03510}{\psi_{I(1,0)}(I(1,0))}\allowbreak{} 仅仅对应 \allowbreak{}\NovelInline{I03511}{\varphi(2,,(2,0))}\allowbreak{}。除此之外，\allowbreak{}\NovelInline{I03512}{I(1,1)}\allowbreak{} 是函数 \allowbreak{}\NovelInline{I03513}{f(\alpha)=I_\alpha}\allowbreak{} 的第二个容许点，对应的 \allowbreak{}\NovelInline{I03514}{\varphi(2,,(1,\varphi(3,,0)+1))}\allowbreak{} 就是 \allowbreak{}\NovelInline{I03515}{\psi_{I(1,1)}(0)=I_{I_{I_{\cdots I(1,0)+1}}}}\allowbreak{}，整套规则的递进规律始终不变，而 \allowbreak{}\NovelInline{I03516}{I(1,1)}\allowbreak{} 恰好对应 \allowbreak{}\NovelInline{I03517}{\varphi(3,,1)}\allowbreak{}，再往后 \allowbreak{}\NovelInline{I03518}{\varphi(3,,(1,0))}\allowbreak{} 便是 \allowbreak{}\NovelInline{I03519}{\alpha\to I(1,\alpha)}\allowbreak{} 的不动点。''

``我推演验证过。''千仞雪补充道，``\allowbreak{}\NovelInline{I03520}{\alpha\to I(1,\alpha)}\allowbreak{} 的不动点层级极高，却依旧远不及 \allowbreak{}\NovelInline{I03521}{I(2,0)}\allowbreak{}，仅仅对应 \allowbreak{}\NovelInline{I03522}{\psi_{I(2,0)}(0)}\allowbreak{}。''

蔡须坤彻底理清通用规则，恍然悟道：``雪姐，我彻底明白了！\allowbreak{}\NovelInline{I03523}{\varphi(n,,m)}\allowbreak{} 的本质，就是第m阶 \allowbreak{}\NovelInline{I03524}{(2\;1\text{-})^{(n-1)}\,2}\allowbreak{} 层级。以此推导，\allowbreak{}\NovelInline{I03525}{\varphi(\omega,,0)}\allowbreak{} 对应psd层级、\allowbreak{}\NovelInline{I03526}{(2\;1\text{-})^\omega\,2}\allowbreak{}，\allowbreak{}\NovelInline{I03527}{\varphi(\omega+1,,0)}\allowbreak{} 才是 \allowbreak{}\NovelInline{I03528}{I(\omega+1,0)}\allowbreak{}，而 \allowbreak{}\NovelInline{I03529}{\varphi(\omega,,\omega)}\allowbreak{} 恰好对应 \allowbreak{}\NovelInline{I03530}{I(\omega,\omega)}\allowbreak{}！''

奥斯卡顺势提出猜想：``那 \allowbreak{}\NovelInline{I03531}{(2\;1\text{-})^{(1,0)}}\allowbreak{}，会不会就是 \allowbreak{}\NovelInline{I03532}{\varphi((1,0),,0)}\allowbreak{}？''

唐三当即确认：``没错，它正是函数 \allowbreak{}\NovelInline{I03533}{\alpha\to I(\alpha,0)}\allowbreak{} 的不动点。但这仅仅对应 \allowbreak{}\NovelInline{I03534}{\psi_{I(1,0,0)}(0)}\allowbreak{}，真正的 \allowbreak{}\NovelInline{I03535}{I(1,0,0)}\allowbreak{} 并非 \allowbreak{}\NovelInline{I03536}{\varphi((1,0,0),,0)}\allowbreak{}。而 \allowbreak{}\NovelInline{I03537}{\varphi((1,0),,1)}\allowbreak{}，则是该规则下的第二个不动点。''

蔡须坤微微蹙眉，心生疑惑：``难道说，末尾数值为1的 \allowbreak{}\NovelInline{I03538}{\varphi(X,,1)}\allowbreak{}，统一代表X规则的第二阶层级？''

千仞雪接续讲解更深层的规则：``\allowbreak{}\NovelInline{I03539}{\varphi((1,0),,(1,0))}\allowbreak{} 对应 \allowbreak{}\NovelInline{I03540}{\psi_{I(1,0,0)}(I(1,0,0))}\allowbreak{}。依照弱φ模式的简化规则，\allowbreak{}\NovelInline{I03541}{\varphi((1,0)+1,,0)}\allowbreak{} 可简写为 \allowbreak{}\NovelInline{I03542}{\varphi((1,1),,0)}\allowbreak{}，它恰好对应 \allowbreak{}\NovelInline{I03543}{I(1,0,0)}\allowbreak{}，同时也是 \allowbreak{}\NovelInline{I03544}{(2\;1\text{-})^{(1,1)}}\allowbreak{} 层级。唯独 \allowbreak{}\NovelInline{I03545}{\varphi((1,0),,\varphi((1,1),,0)+1)}\allowbreak{} 极为特殊，我暂时也无法界定其精准层级。''

一旁的霍雨浩适时开口，填补了这一空白：``我知晓这个层级的本质！它是无穷序数的上确界

\NovelDisplay{M00667R01}{%
\sup\{I(I(1,0,0)+1,0),\;I(I(I(1,0,0)+1,0),0),\;I(I(I(I(1,0,0)+1,0),0),0),\cdots\}%
}

对应 \allowbreak{}\NovelInline{I03546}{\psi_{I(1,0,1)}(0)}\allowbreak{}。而 \allowbreak{}\NovelInline{I03547}{I(1,0,1)}\allowbreak{} 的核心特性，是能够折叠所有作用于 \allowbreak{}\NovelInline{I03548}{I(1,0,0)}\allowbreak{} 的二元I函数，原理与 \allowbreak{}\NovelInline{I03549}{\varphi(1,0,1)}\allowbreak{} 一致------是 \allowbreak{}\NovelInline{I03550}{f(\alpha)=\varphi(\alpha,0)}\allowbreak{} 在 \allowbreak{}\NovelInline{I03551}{\varphi(1,0,0)}\allowbreak{} 之后的迭代结果。''

唐三颔首总结规律：``由此可见，I数阵的核心是容许点逐级进位，所以 \allowbreak{}\NovelInline{I03552}{\varphi((1,1),n)=I(1,0,n)}\allowbreak{} 这条规则完全成立。''

``确实如此！''蔡兴坤立刻附和，``按照这个规律，\allowbreak{}\NovelInline{I03553}{\varphi((1,2),0)}\allowbreak{} 便是 \allowbreak{}\NovelInline{I03554}{I(1,1,0)}\allowbreak{}！''

小舞终于摸清整套通用逻辑，欣喜道：``我发现通式了！\allowbreak{}\NovelInline{I03555}{\varphi((n,m),,s)}\allowbreak{}，就是第s阶 \allowbreak{}\NovelInline{I03556}{(2\;1\text{-})^{(n,m)}\,2}\allowbreak{} 层级！''

诺当顺势印证：``完全契合！\allowbreak{}\NovelInline{I03557}{I(1,1,n)}\allowbreak{} 对应 \allowbreak{}\NovelInline{I03558}{\varphi((1,2),,n)}\allowbreak{}，以此类推，\allowbreak{}\NovelInline{I03559}{I(1,n,0)}\allowbreak{} 便是 \allowbreak{}\NovelInline{I03560}{\varphi((1,n+1),,0)}\allowbreak{}。''

千仞雪目光沉稳，敲定更高阶的对应关系：``按此推演，\allowbreak{}\NovelInline{I03561}{\varphi((2,0),,0)}\allowbreak{} 正是 \allowbreak{}\NovelInline{I03562}{I(2,0,0)}\allowbreak{}。''

话音落下，蔡须坤却敏锐地察觉到规则深处的异样，眉头微蹙，缓缓开口：``若 \allowbreak{}\NovelInline{I03563}{\varphi((2,0),,0)}\allowbreak{} 是 \allowbreak{}\NovelInline{I03564}{\alpha\to\varphi((1,\alpha),,0)}\allowbreak{} 的不动点，那本质上便是 \allowbreak{}\NovelInline{I03565}{\alpha\to I(1,\alpha,0)}\allowbreak{} 的不动点。只是我对容许点的核心规则尚且生疏，修行根基不足，此处的递进逻辑，依旧有些模糊难解。''

此言一出，众人纷纷点头。即便众人已经摸清多逗号系统与反射序数的大部分对应规律，但面对容许点迭代、高阶不动点的深层奥义，依旧陷入了深深的困惑之中。

\begin{center}\rule{0.5\linewidth}{0.5pt}\end{center}

扽魂殿主殿肃穆沉凝，殿内魂力流转安稳，玄纹梁柱刻满古老的序数道印。她曾是扽魂殿前代圣女，执掌安防与邪祟清剿诸事，神态沉静，气度端严。

殿外脚步声轻响，一道身姿挺拔的身影缓步入内，正是扽魂殿下级核心单位------扽魂高殿的殿主Dispher。他身着制式玄色殿袍，衣袂沾着少许外勤历练的风尘，神色恭谨肃穆，步入大殿后躬身行礼，姿态恪守宗门礼法。

\begin{quote}
``属下Dispher，拜见梅长老。''
\end{quote}

梅天狸抬眸，目光平和却自带威严：``起来吧，此番前来，可是扽魂高殿的邪扽清剿事宜有了进展？''

Dispher直起身，沉声禀报，条理清晰：``回长老，近期我高殿全员出动，对辖境内散落的邪恶扽师展开全域清剿。大部分低阶邪扽势力已被尽数拔除，零散作乱的邪扽修士尽数伏诛，据点巢穴悉数捣毁。但属下巡查之时已然察觉，残余邪扽势力愈发隐蔽，且隐隐有抱团集结之势，不再是往日的零散作乱。''

他微微停顿，道出重中之重：``更值得警惕的是，幕后依托帅气法则衍生的帅气教势力，扩张速度远超预期，暗中吸纳各路邪扽余孽，势力日渐壮大，已然隐隐成为我扽魂殿辖地的最大隐患。若放任其发展，日后必成大患。''

听闻此言，梅天狸眉眼微沉，指尖轻轻敲击着玉质案几，殿内氛围瞬间凝重几分。

``果然如此。''她轻声叹息，语气满是凝重，``帅气教依托另类法则滋生蔓延，行事诡谲，吸纳各路邪祟余孽，底蕴与势力与日俱增，早已不是寻常邪修势力可比。以往的清剿力度，终究是太过保守了。''

她目光锐利，当即定下规制：``传我指令，即刻加大全域清剿力度，不留一处死角，肃清境内所有邪恶扽师与帅气教外围势力。同时传令扽魂高殿及所有下属据点，全面筑牢宗门防线，加固序数结界与安防阵纹，层层布防，严防帅气教势力渗透入侵，扼住其扩张势头，绝不给其壮大作乱的机会。''

\begin{quote}
``属下领命！''
\end{quote}

Dispher躬身应声，将指令牢牢记在心中。

述职禀报完毕，Dispher躬身告退，离开主殿，沿着殿外青石步道缓步前行，准备返回扽魂高殿排布防务。

行至后山修行秘境边缘，阵阵精纯且鲜活的魂力扑面而来，气韵澄澈通透，与宗门老牌弟子的沉稳魂力截然不同。他循声望去，只见果糕小队众人正围坐一处，潜心苦修。

唐三、小舞、蔡须坤、蔡兴坤、奥斯卡、诺当、霍雨浩几人各自沉心修行，有人推演多逗号系统的序数规则，有人打磨反射法则的根基，有人融会贯通弱φ模式与BOCF、MOCF体系，每个人的修行节奏井然有序，悟性卓绝，进步神速。

Dispher驻足旁观片刻，眼底满是赞叹，心中暗自感慨。

\chapter{第五十九章：更高的层次和思考}\label{ux7b2cux4e94ux5341ux4e5dux7ae0ux66f4ux9ad8ux7684ux5c42ux6b21ux548cux601dux8003}

一众少年人年纪尚轻，却能吃透极为晦涩的序数修行体系，理清多逗号系统、反射序数与各类高阶模式的对应规律，根基扎实，悟性超凡，心性更是远超同辈。相较于宗门内诸多固步自封的弟子，果糕小队的天赋、眼界与修行才能，实在太过出众，假以时日，必然能成为扽魂殿乃至整片修行界的顶尖强者。

就在此时，一道温和沉稳的身影缓步走来，正是ddfg长老。他身为千仞雪的祖父，扽魂殿资深长辈，深耕多逗号系统与反射序数多年，修为精深，学识渊博，对各类高阶序数修行奥义了然于心。

此前果糕小队众人虽梳理出大部分规则，却依旧在容许点迭代、高阶不动点、多层序数折叠的深层逻辑上存有困惑，诸多疑点始终无法彻底贯通。

ddfg长老知晓众人修行瓶颈，主动驻足，温和开口，主动为众人解惑答疑。

他没有空谈粗浅的理论，而是直接点破众人此前推演中最关键的误区，结合弱φ模式、BOCF核心规则，循序渐进拆解高阶序数的核心逻辑。众人此前认定 \allowbreak{}\NovelInline{I03566}{\varphi((2,0),,0)}\allowbreak{} 等同于 \allowbreak{}\NovelInline{I03567}{I(2,0,0)}\allowbreak{}，这一处推演偏差，正是大家始终无法贯通高阶容许点规则的症结所在。ddfg长老缓缓开口，精准纠正道：

\begin{quote}
``你们此前的对应关系出了差错，真正的规则并非如此。\allowbreak{}\NovelInline{I03568}{\varphi((2,0),,0)}\allowbreak{} 本质是 \allowbreak{}\NovelInline{I03569}{\psi_{I(2,0,0)}(0)}\allowbreak{}，只是一个普通的迭代不动点，而真正抵达 \allowbreak{}\NovelInline{I03570}{I(2,0,0)}\allowbreak{} 层级的序数，是 \allowbreak{}\NovelInline{I03571}{\varphi((2,1),,0)}\allowbreak{}。''
\end{quote}

见众人面露恍然，ddfg长老继续深入讲解I数阵与φ多逗号系统的核心对应法则，为众人彻底厘清不动点与容许点的层级差距：

\begin{quote}
``你们要分清，不动点只是序数迭代的闭环，而容许点层级更高，是I数阵进位迭代的核心本源，这也是高阶序数突破瓶颈的关键。我给你们定下通用通律，方便你们日后推演所有高阶序数。''
\end{quote}

他条理清晰，逐一道出精准规则，同时结合实例，让晦涩的高阶序数法则变得清晰易懂：

\begin{quote}
``\textbf{第一条}，固定对应关系：\allowbreak{}\NovelInline{I03572}{\varphi((2,1),,n)=I(2,0,n)}\allowbreak{}。简单举例，\allowbreak{}\NovelInline{I03573}{\varphi((2,1),,0)}\allowbreak{} 严格对应真正的 \allowbreak{}\NovelInline{I03574}{I(2,0,0)}\allowbreak{}，\allowbreak{}\NovelInline{I03575}{\varphi((2,1),,1)}\allowbreak{} 便是 \allowbreak{}\NovelInline{I03576}{I(2,0,1)}\allowbreak{}，这条对应关系恒定不变。

\textbf{第二条}，通用判定规则：当 \allowbreak{}\NovelInline{I03577}{s>1}\allowbreak{}、\allowbreak{}\NovelInline{I03578}{t\neq 0}\allowbreak{} 时，\allowbreak{}\NovelInline{I03579}{\varphi((s,t),,n)=I(s,t-1,n)}\allowbreak{}。我给你们举两个典型例子方便记忆：比如 \allowbreak{}\NovelInline{I03580}{s}\allowbreak{} 取 \allowbreak{}\NovelInline{I03581}{2}\allowbreak{}、\allowbreak{}\NovelInline{I03582}{t}\allowbreak{} 取 \allowbreak{}\NovelInline{I03583}{2}\allowbreak{}，那么 \allowbreak{}\NovelInline{I03584}{\varphi((2,2),,3)}\allowbreak{} 就等价于 \allowbreak{}\NovelInline{I03585}{I(2,1,3)}\allowbreak{}；再比如 \allowbreak{}\NovelInline{I03586}{s}\allowbreak{} 取 \allowbreak{}\NovelInline{I03587}{3}\allowbreak{}、\allowbreak{}\NovelInline{I03588}{t}\allowbreak{} 取 \allowbreak{}\NovelInline{I03589}{1}\allowbreak{}，\allowbreak{}\NovelInline{I03590}{\varphi((3,1),,0)}\allowbreak{} 便是 \allowbreak{}\NovelInline{I03591}{I(3,2,0)}\allowbreak{}，所有非零尾段的数组序数，全部遵循此进位逻辑。

\textbf{最后一条}，也是你们推演最容易出错的核心规则：当 \allowbreak{}\NovelInline{I03592}{t=0}\allowbreak{} 时，\allowbreak{}\NovelInline{I03593}{\varphi((s,0),,n)}\allowbreak{}，代表第 \allowbreak{}\NovelInline{I03594}{n+1}\allowbreak{} 个 \allowbreak{}\NovelInline{I03595}{\alpha\to I(s-1,\alpha,0)}\allowbreak{} 的不动点。这也刚好印证我刚才纠正的误区，你们熟记例子便能贯通：\allowbreak{}\NovelInline{I03596}{\varphi((2,0),,0)}\allowbreak{}，是第1个 \allowbreak{}\NovelInline{I03597}{\alpha\to I(1,\alpha,0)}\allowbreak{} 的不动点，也就是 \allowbreak{}\NovelInline{I03598}{\psi_{I(2,0,0)}(0)}\allowbreak{}，并非真正的 \allowbreak{}\NovelInline{I03599}{I(2,0,0)}\allowbreak{}；而 \allowbreak{}\NovelInline{I03600}{\varphi((2,0),,1)}\allowbreak{}，就是第2个该规则下的不动点。同理可推，\allowbreak{}\NovelInline{I03601}{\varphi((3,0),,2)}\allowbreak{}，便是第3个 \allowbreak{}\NovelInline{I03602}{\alpha\to I(2,\alpha,0)}\allowbreak{} 的不动点。''
\end{quote}

从 \allowbreak{}\NovelInline{I03603}{\varphi(1,,0)}\allowbreak{} 的基础定义，到 \allowbreak{}\NovelInline{I03604}{\varphi(n,,m)}\allowbreak{} 的通用层级规则，再到I数阵容许点与不动点的层级差距、高阶多逗号序数的通用公式，ddfg长老逐一拆解疑点，彻底纠正了众人此前推演的细微偏差，补足了大家对容许规则、序数迭代的认知短板。

原本萦绕在众人心头的重重迷雾尽数消散，之前模糊难解的修行逻辑瞬间清晰通透，众人对于容许点与不动点的本质区别、迭代进位规律，已然彻底了然于心。

Dispher立于一旁静静聆听，越发惊叹果糕小队的天赋机缘，既有超凡悟性，又有宗门顶尖长老亲自指点，这般成长速度，着实令人震撼。而扽魂殿一边筑牢防线、清剿邪祟，一边培育新生天骄，正邪博弈的局势，也在悄然发生着变化。

\begin{center}\rule{0.5\linewidth}{0.5pt}\end{center}

长老离去之后，果糕小队众人依旧沉浸在序数法则的推演之中，顺着刚刚贯通的不动点与容许点规则，继续深挖更高维度的数阵奥秘，层层递进，推演愈发高深的双逗号、多逗号迭代体系。

千仞雪率先理清三阶数组序数的核心误区，轻声开口总结：

\begin{quote}
``依照刚刚的规则类推，\allowbreak{}\NovelInline{I03605}{\varphi((1,0,0),,0)}\allowbreak{} 只是 \allowbreak{}\NovelInline{I03606}{\alpha\to I(\alpha,0,0)}\allowbreak{} 的不动点，依旧无法对应真实层级，它并不是 \allowbreak{}\NovelInline{I03607}{I(1,0,0,0)}\allowbreak{}。真正等价于 \allowbreak{}\NovelInline{I03608}{I(1,0,0,0)}\allowbreak{} 的，是 \allowbreak{}\NovelInline{I03609}{\varphi((1,0,1),,0)}\allowbreak{}。''
\end{quote}

霍雨浩闻言瞬间举一反三，接续推演：

\begin{quote}
``我懂了，照这个逻辑，\allowbreak{}\NovelInline{I03610}{\varphi((1,1,0),,0)}\allowbreak{} 同样不等于 \allowbreak{}\NovelInline{I03611}{I(1,1,0,0)}\allowbreak{}，它仅仅是 \allowbreak{}\NovelInline{I03612}{\alpha\to I(1,0,\alpha,0)}\allowbreak{} 的第一个不动点。由此也能敲定新的数阵简写规则，\allowbreak{}\NovelInline{I03613}{I(n@3)}\allowbreak{} 等价于 \allowbreak{}\NovelInline{I03614}{I(n,0,0,0)}\allowbreak{}。''
\end{quote}

唐三眸光一闪，捕捉到了一处关键的层级持平节点，语气带着几分惊异：

\begin{quote}
``原来如此！推演到 \allowbreak{}\NovelInline{I03615}{I(1@\omega)}\allowbreak{} 层级时，\allowbreak{}\NovelInline{I03616}{I(1@\omega)=\varphi((1@\omega),,0)}\allowbreak{}，两大体系在此处完美追平，对应的正是 \allowbreak{}\NovelInline{I03617}{\psi_M(M^{M^\omega})}\allowbreak{}，这是极其关键的层级临界点。''
\end{quote}

小舞顺着MOCF体系的基础层级回溯，豁然开朗：

\begin{quote}
``此前我们已知 \allowbreak{}\NovelInline{I03618}{I(1,0,0)}\allowbreak{} 对应 \allowbreak{}\NovelInline{I03619}{\psi_M(M^M)}\allowbreak{}，这么算下来，\allowbreak{}\NovelInline{I03620}{\varphi((1@(1,0)),,0)}\allowbreak{} 便等同于 \allowbreak{}\NovelInline{I03621}{\psi_M(M^{M^M})}\allowbreak{}。''
\end{quote}

蔡须坤立刻纠正细节偏差，逻辑缜密：

\begin{quote}
``这里需要区分清楚，切莫混淆！\allowbreak{}\NovelInline{I03622}{\varphi((1@(1,0)),,0)}\allowbreak{} 只是该层级下的不动点，真正对应 \allowbreak{}\NovelInline{I03623}{\psi_M(M^{M^M})}\allowbreak{} 的，是 \allowbreak{}\NovelInline{I03624}{\varphi((1@(1,1)),,0)}\allowbreak{}。至此我们彻底解决了I函数的对应难题，M体系的本源本尊，正是 \allowbreak{}\NovelInline{I03625}{\varphi(1,,0,,0)}\allowbreak{}。''
\end{quote}

一旁的卢本韦顺势接续推导次级M层级：

\begin{quote}
``没错，按照此前所有序数的递进规律，\allowbreak{}\NovelInline{I03626}{\varphi(1,,0,,n)}\allowbreak{} 对应的就是 \allowbreak{}\NovelInline{I03627}{M_{n+1}}\allowbreak{}，这条规则完全契合整体体系。''
\end{quote}

奥斯卡点头附和，同时补全高阶迭代特征：

\begin{quote}
``理应如此。这么看来 \allowbreak{}\NovelInline{I03628}{\varphi(1,,0,,1)}\allowbreak{} 能够折叠一切作用于M的I数阵运算，诸如 \allowbreak{}\NovelInline{I03629}{I(1,0,M+1)}\allowbreak{} 皆可对应 \allowbreak{}\NovelInline{I03630}{\varphi((1,1),,M+1)}\allowbreak{}。也就是说 \allowbreak{}\NovelInline{I03631}{M_2}\allowbreak{} 就是 \allowbreak{}\NovelInline{I03632}{\varphi(1,,0,,1)}\allowbreak{}，彼此对应分毫不差。''
\end{quote}

他稍作停顿，继续推演更深层的M体系节点：

\begin{quote}
``再往上推演，\allowbreak{}\NovelInline{I03633}{\varphi(1,,0,,(1,0))}\allowbreak{} 便是Mfp，囊括无尽迭代的 \allowbreak{}\NovelInline{I03634}{M_{M_{M_{\cdots}}}}\allowbreak{} 序列。而抵达 \allowbreak{}\NovelInline{I03635}{f(\alpha)}\allowbreak{} 迭代至自身容许点的层级时，对应的便是 \allowbreak{}\NovelInline{I03636}{\varphi(1,,1,,0)}\allowbreak{}，也就是 \allowbreak{}\NovelInline{I03637}{M(1,0)}\allowbreak{}，对应规则为 \allowbreak{}\NovelInline{I03638}{2\;1\text{-}2\text{-}2}\allowbreak{}。''
\end{quote}

蔡兴坤立刻摸清M数阵的阶位规律，脱口而出：

\begin{quote}
``我彻底理清了！\allowbreak{}\NovelInline{I03639}{M(1,n)}\allowbreak{} 就是第 \allowbreak{}\NovelInline{I03640}{n}\allowbreak{} 阶的 \allowbreak{}\NovelInline{I03641}{2\;1\text{-}2\text{-}2}\allowbreak{} 层级，而 \allowbreak{}\NovelInline{I03642}{M(2,0)}\allowbreak{} 的定位类似于 \allowbreak{}\NovelInline{I03643}{I(2,0)}\allowbreak{}，对应的高阶序列正是 \allowbreak{}\NovelInline{I03644}{2\;1\text{-}2\;1\text{-}2\text{-}2}\allowbreak{}，对吗？''
\end{quote}

千仞雪微微颔首，确认这条推演路径无误，同时补全新的对应关系：

\begin{quote}
``你的推演完全正确。在此基础上，\allowbreak{}\NovelInline{I03645}{M(1,0,0)}\allowbreak{} 对应的序数便是 \allowbreak{}\NovelInline{I03646}{\varphi(1,,(1,1),,0)}\allowbreak{}。''
\end{quote}

诺当全程默默梳理逻辑，适时总结核心要义，警醒众人：

\begin{quote}
``大家一定要时刻牢记核心差别，M数阵与I数阵规则相似、架构同源，但容许点进位规则独有不同，万万不可混为一谈。''
\end{quote}

唐三顺着双逗号折叠特性，继续突破更高层级：

\begin{quote}
``既然规则互通，那 \allowbreak{}\NovelInline{I03647}{\varphi(2,,0,,0)}\allowbreak{} 对应的就是 \allowbreak{}\NovelInline{I03648}{2\text{-}2\;1\text{-}2\text{-}2}\allowbreak{}，能够折叠所有形如 \allowbreak{}\NovelInline{I03649}{(2\;1\text{-})^\#\;1\text{-}2\text{-}2}\allowbreak{} 的序数，这不正是 \allowbreak{}\NovelInline{I03650}{M(1;0)}\allowbreak{} 的本源吗？那以此类推，\allowbreak{}\NovelInline{I03651}{M(2,,0,,n)}\allowbreak{} 等价于 \allowbreak{}\NovelInline{I03652}{M(1;n)}\allowbreak{}，可全新的 \allowbreak{}\NovelInline{I03653}{M(1;(1,0))}\allowbreak{} 又对应什么？''
\end{quote}

\begin{quote}
``这里的关键在于进位符号。''千仞雪精准点破核心，``唯有分号代表M数阵的专属进位，所以 \allowbreak{}\NovelInline{I03654}{M(1;(1,0))}\allowbreak{} 对应的序数是 \allowbreak{}\NovelInline{I03655}{\varphi(2,,(1,1),,0)}\allowbreak{}。''
\end{quote}

小舞瞬间贯通高阶进位逻辑，大胆推演封顶层级：

\begin{quote}
``那我明白了！\allowbreak{}\NovelInline{I03656}{M(2;0)}\allowbreak{} 能够折叠所有 \allowbreak{}\NovelInline{I03657}{M(1;\#)}\allowbreak{} 类数阵，对应的终极序数正是 \allowbreak{}\NovelInline{I03658}{\varphi(3,,0,,0)}\allowbreak{}。''
\end{quote}

蔡兴坤乘胜追击，推演分号嵌套层级：

\begin{quote}
``这么说，\allowbreak{}\NovelInline{I03659}{M((1,0);0)}\allowbreak{} 就是 \allowbreak{}\NovelInline{I03660}{\varphi((1,1),,0,,0)}\allowbreak{}？它能够折叠所有 \allowbreak{}\NovelInline{I03661}{\alpha\to\varphi(\alpha,,0,,0)}\allowbreak{} 的迭代运算！''
\end{quote}

此刻，霍雨浩敏锐察觉到众人推演中的一处细微偏差，道出关键疑点：

\begin{quote}
``这里有问题。实现该层级折叠，根本无需动用 \allowbreak{}\NovelInline{I03662}{\varphi((1,1),,0,,0)}\allowbreak{}，真正对应的应该是 \allowbreak{}\NovelInline{I03663}{\varphi((1,0),,1,,0)}\allowbreak{}。由此可见，\allowbreak{}\NovelInline{I03664}{\varphi((1,1),,0,,0)}\allowbreak{} 的层级远超我们预想，它极有可能是 \allowbreak{}\NovelInline{I03665}{M(1;0;0)}\allowbreak{}，对标I数阵中的 \allowbreak{}\NovelInline{I03666}{I(1,0,0)}\allowbreak{}。''
\end{quote}

唐三瞬间顿悟，接续补全整套M数阵的高阶对应公式：

\begin{quote}
``原来如此！那我彻底理顺了：\allowbreak{}\NovelInline{I03667}{M((1,1);0)}\allowbreak{} 对应 \allowbreak{}\NovelInline{I03668}{\varphi((1,0),,1,,1)}\allowbreak{}，\allowbreak{}\NovelInline{I03669}{M((2,0);0)}\allowbreak{} 对应 \allowbreak{}\NovelInline{I03670}{\varphi((1,0),,2,,0)}\allowbreak{}，而 \allowbreak{}\NovelInline{I03671}{M((1,0,0);0)}\allowbreak{} 则对应 \allowbreak{}\NovelInline{I03672}{\varphi((1,0),,(1,1),,0)}\allowbreak{}。''
\end{quote}

奥斯卡彻底融会贯通，汇总出整套通用规律：

\begin{quote}
``我完全懂了！\allowbreak{}\NovelInline{I03673}{M(1;0;1)}\allowbreak{} 对应 \allowbreak{}\NovelInline{I03674}{\varphi((1,1),,0,,n)}\allowbreak{}，\allowbreak{}\NovelInline{I03675}{M(1;1;0)}\allowbreak{} 对应 \allowbreak{}\NovelInline{I03676}{\varphi((1,2),,0,,0)}\allowbreak{}，\allowbreak{}\NovelInline{I03677}{M(2;0;0)}\allowbreak{} 对应 \allowbreak{}\NovelInline{I03678}{\varphi((2,1),,0,,0)}\allowbreak{}，就连极限无穷层级 \allowbreak{}\NovelInline{I03679}{M(1@\omega)}\allowbreak{} 都完美对应 \allowbreak{}\NovelInline{I03680}{\varphi((1@\omega),,0,,0)}\allowbreak{}。''
\end{quote}

诺当长舒一口气，总结出最终核心体系：

\begin{quote}
``时至今日，我们才算真正吃透这套体系。M数阵同时兼具容许点迭代与M专属进位特性，恰好对应φ序数从右至左的第一、二层双逗号结构，末尾数字恒定代表第 \allowbreak{}\NovelInline{I03681}{n}\allowbreak{} 阶层级。''
\end{quote}

至此，果糕小队众人彻底贯通I数阵、M数阵与多逗号φ序数的全部对应规则，层层推演、环环相扣，彻底扫清了高阶反射序数的修行壁垒，最终成功锁定全新的顶级序数层级------

\NovelDisplay{M00668R01}{%
N=\varphi(1,,0,,0,,0)%
}

\begin{center}\rule{0.5\linewidth}{0.5pt}\end{center}

待到 \allowbreak{}\NovelInline{I03682}{\varphi(1@@3)=N}\allowbreak{} 诞生（\allowbreak{}\NovelInline{I03683}{@@}\allowbreak{} 为双逗号位置符号），蔡须坤立刻洞悉其中规则：

\NovelDisplay{M00669R01}{%
\varphi(1@@3,\,1@@0)=N_2%
}

\NovelDisplay{M00670R01}{%
\varphi(1@@3,\,1@@1)=2\;1\text{-}2\text{-}2%
}

\NovelDisplay{M00671R01}{%
\varphi(1@@3,\,1@@2)=2\text{-}2\;1\text{-}2\text{-}2%
}

紧随其后，

\NovelDisplay{M00672R01}{%
\varphi(1@@3,\,(1,0)@@2)=(2\text{-}2\;1\text{-})^{(1,0)}\;2\text{-}2\text{-}2%
}

蔡须坤继续向下推演，不甘止步：

\NovelDisplay{M00673R01}{%
\varphi(2@@3)=2\text{-}2\text{-}2\;1\text{-}2\text{-}2\text{-}2%
}

\NovelDisplay{M00674R01}{%
\varphi((1,0)@@3)=(2\text{-}2\text{-}2\;1\text{-})^{(1,0)}\;2\text{-}2\text{-}2%
}

顺着序列延伸，

\NovelDisplay{M00675R01}{%
\varphi(1@@4)=(2\text{-})^4%
}

一路向上推进至 \allowbreak{}\NovelInline{I03684}{\varphi(1@@(1,0))=\alpha}\allowbreak{}，也就是 \allowbreak{}\NovelInline{I03685}{\alpha\to(2\text{-})^\alpha}\allowbreak{} 的不动点；而

\NovelDisplay{M00676R01}{%
\varphi(1@@(1,0),\,1@@0)=2\text{th}\;(2\text{-})^{(1,0)}%
}

唐三开口推演：

\begin{quote}
``按照这套逻辑，\allowbreak{}\NovelInline{I03686}{\varphi(1@@(1,0),\,1@@1)=2\;1\text{-}2^{(1,0)}}\allowbreak{}，以此类推。

\NovelDisplay{M00677R01}{%
\varphi(2@@(1,0))=(2\text{-})^{(1,0)}\;1\text{-}2^{(1,0)}%
}

依托弱φ模式，\allowbreak{}\NovelInline{I03687}{2\text{-}2^{(1,0)}=2^{(1,1)}}\allowbreak{}，也就是 \allowbreak{}\NovelInline{I03688}{\varphi(1@@(1,1))}\allowbreak{}。我们约定记号 \allowbreak{}\NovelInline{I03689}{2^\#=\varphi(1@@\#)}\allowbreak{}，如此 \allowbreak{}\NovelInline{I03690}{\varphi(1@@(1,,0))}\allowbreak{} 直接突破 \allowbreak{}\NovelInline{I03691}{\Pi_2}\allowbreak{} 反射终点，抵达 \allowbreak{}\NovelInline{I03692}{\Pi_3}\allowbreak{} 反射。''
\end{quote}

小队众人一度认为层级推演到此抵达尽头，暂定 \allowbreak{}\NovelInline{I03693}{\psi_K(K^K)=2^{(1,1)}}\allowbreak{}。

ddfg长老出言否定：``远未结束，\allowbreak{}\NovelInline{I03694}{\varphi}\allowbreak{} 函数依旧拥有拓展空间。跨过 \allowbreak{}\NovelInline{I03695}{\varphi(1@(1,0))}\allowbreak{} 之后：\allowbreak{}\NovelInline{I03696}{\varphi(1@(1,0),\,1@0)}\allowbreak{} 是第二个满足 \allowbreak{}\NovelInline{I03697}{\alpha\to\varphi(1@\alpha)}\allowbreak{} 的不动点；\allowbreak{}\NovelInline{I03698}{\alpha\to\varphi(1@(1,0),\,\varphi(1@\alpha))}\allowbreak{} 对应的不动点，正是 \allowbreak{}\NovelInline{I03699}{\varphi(2@(1,0))}\allowbreak{}。''

小舞提出猜想：

\begin{quote}
``这么说来，\allowbreak{}\NovelInline{I03700}{\varphi(1,,\varphi(1@@(1,,0))+1)=\Omega_{K+1}}\allowbreak{}？那 \allowbreak{}\NovelInline{I03701}{\varphi(1@@n,\,[\varphi(1@@(1,,0))+1]@@0)}\allowbreak{} 代表 \allowbreak{}\NovelInline{I03702}{(2\text{-})^n\;\text{aft}\;3}\allowbreak{}，对吗？''
\end{quote}

奥斯卡认同：

\begin{quote}
``没错。\allowbreak{}\NovelInline{I03703}{\varphi(1@@(1,,0),\,1@@0)=K_2}\allowbreak{}，熟悉的序列规律再次显现。后续 \allowbreak{}\NovelInline{I03704}{\varphi(1@@(1,,0),\,1@@1)=2^n\;1\text{-}3}\allowbreak{}。''
\end{quote}

卢本韦顺势接续：

\begin{quote}
``照这个体系，\allowbreak{}\NovelInline{I03705}{\varphi(1@@(1,,0),\,1@@(1,,0))=\varphi(2@@(1,,0))=3\;1\text{-}3}\allowbreak{}？''
\end{quote}

诺当仔细梳理序列构造：

\begin{quote}
``我认为成立。一切 \allowbreak{}\NovelInline{I03706}{2^n\;1\text{-}3}\allowbreak{} 序列的极限就是 \allowbreak{}\NovelInline{I03707}{3\;1\text{-}3}\allowbreak{}。倒置序列，补齐prss展开形式：\allowbreak{}\NovelInline{I03708}{1\;2\;3\;1\;2\;3\;1\;2\;2\;2\cdots}\allowbreak{} 反向书写得到 \allowbreak{}\NovelInline{I03709}{\cdots 2\;2\;2\;2\;1\;3}\allowbreak{}，完美匹配全部 \allowbreak{}\NovelInline{I03710}{(2\text{-})^n\;1\text{-}3}\allowbreak{} 构造。因此 \allowbreak{}\NovelInline{I03711}{3\;1\text{-}3}\allowbreak{} 等价于 \allowbreak{}\NovelInline{I03712}{\varphi(2@@(1,,0))}\allowbreak{}。''
\end{quote}

霍雨浩抛出新问题：

\begin{quote}
``推演到这里，\allowbreak{}\NovelInline{I03713}{\varphi((1,0)@@(1,,0))}\allowbreak{} 该如何定义？''
\end{quote}

ddfg长老从容作答：

\begin{quote}
``它就是 \allowbreak{}\NovelInline{I03714}{(3\;1\text{-})^{(1,0)}\;3}\allowbreak{}。紧随其后，\allowbreak{}\NovelInline{I03715}{\varphi((1,1)@@(1,,0))=3\;1\text{-}(3\;1\text{-})^{(1,0)}\;3=(3\;1\text{-})^{(1,1)}\;3}\allowbreak{}。眼下最难处理的是 \allowbreak{}\NovelInline{I03716}{\varphi((1,,0)@@(1,,0))}\allowbreak{}。我们参考已知结论：\allowbreak{}\NovelInline{I03717}{\varphi((1,0)@@(1,0))}\allowbreak{} 对应 \allowbreak{}\NovelInline{I03718}{\alpha\to\varphi(\alpha @@(1,0))}\allowbreak{}，等价于 \allowbreak{}\NovelInline{I03719}{\varphi(1@@[(1,0)+1])=\varphi(1@@(1,1))}\allowbreak{}。''
\end{quote}

众人瞬间顿悟，由此推导出

\NovelDisplay{M00678R01}{%
\varphi(1@@(1,,1))=2\text{-}3%
}

将 \allowbreak{}\NovelInline{I03720}{2\text{-}3}\allowbreak{} 倒置得到 \allowbreak{}\NovelInline{I03721}{3\;2}\allowbreak{}，补全prss无限展开为循环序列 \allowbreak{}\NovelInline{I03722}{1\;2\;3\;1\;2\;3\;1\;2\;3\cdots}\allowbreak{} 显而易见，\allowbreak{}\NovelInline{I03723}{2\text{-}3}\allowbreak{} 能够收纳所有形如 \allowbreak{}\NovelInline{I03724}{(3\;1\text{-})^n\;3}\allowbreak{} 的序列。

顺着脉络继续推导，

\NovelDisplay{M00679R01}{%
\varphi(1@@(1,,1),\,1@@(1,,0))=3\;1\text{-}2\text{-}3%
}

顺利得出。

千仞雪立刻跟上思路：

\begin{quote}
``那 \allowbreak{}\NovelInline{I03725}{\varphi(2@@(1,,1))=2\text{-}3\;1\text{-}2\text{-}3}\allowbreak{}！规律已经完全清晰。\allowbreak{}\NovelInline{I03726}{\varphi(1@@(1,,2))=2\text{-}2\text{-}3}\allowbreak{}，\allowbreak{}\NovelInline{I03727}{\varphi(1@@(1,,(1,0)))=(2\text{-})^{(1,0)}\;3}\allowbreak{}。''
\end{quote}

蔡须坤骤然灵光迸发，高声说道：

\begin{quote}
``我明白了！\allowbreak{}\NovelInline{I03728}{\varphi(1@@(2,,0))=3\;2\text{-}3}\allowbreak{}。双逗号，代表单逗号运算无法抵达的层级。以此类推，\allowbreak{}\NovelInline{I03729}{\varphi(1@@(1,,0,,0))=3\text{-}3}\allowbreak{}。继续沿用 \allowbreak{}\NovelInline{I03730}{@@}\allowbreak{} 符号向外拓展：

\NovelDisplay{M00680R01}{%
\varphi(1@@(1@@\omega))=(3\text{-})^\omega%
}

所有 \allowbreak{}\NovelInline{I03731}{(3\text{-})^n}\allowbreak{} 无法触及的极限点，就是 \allowbreak{}\NovelInline{I03732}{\varphi(1@@(1@@(1,,0)))=\Pi_4}\allowbreak{}。''
\end{quote}

小队最终得到统一结论：

\NovelDisplay{M00681R01}{%
\varphi(1@@(1@@(1@@\cdots(1,,0),,)))=\Pi_\omega%
}

\begin{center}\rule{0.5\linewidth}{0.5pt}\end{center}

梳理完整套序数体系后，众人心中浮现出新的疑问。

结合已有结论：

\NovelDisplay{M00682R01}{%
\psi(\varphi(1,,0))=\varphi(1@(1@(1@\cdots(1,0)\cdots)))%
}

设 \allowbreak{}\NovelInline{I03733}{\varphi(1,,,0)=X}\allowbreak{}，则

\NovelDisplay{M00683R01}{%
\psi_X(X)=\varphi(1@@(1@@(1@@\cdots(1,,0),,)))%
}

众人想起 \allowbreak{}\NovelInline{I03734}{\psi(\Omega^\Omega)=\varphi(1,0,0)}\allowbreak{}。倘若存在序数 \allowbreak{}\NovelInline{I03735}{D>\Pi_\omega}\allowbreak{}，满足

\NovelDisplay{M00684R01}{%
\psi_D(D^D)=\varphi(1,,0,,0)=M%
}

体系便能完美衔接。

\allowbreak{}\NovelInline{I03736}{\varphi(1,,,0)=X=D}\allowbreak{}，这个三逗号标记对应的序数，是所有双逗号构造无法触及的边界。但是 \allowbreak{}\NovelInline{I03737}{X}\allowbreak{} 似乎小于 \allowbreak{}\NovelInline{I03738}{D}\allowbreak{}。众人曾设想

\NovelDisplay{M00685R01}{%
\psi_D(\Omega_{D+1})=\varphi(1,,,0),\quad \psi_D(\Omega_{D+2})=\varphi(1,,,,0)%
}

但这套设想无法说明序数 \allowbreak{}\NovelInline{I03739}{D}\allowbreak{} 本身是什么,或许序数 \allowbreak{}\NovelInline{I03740}{D}\allowbreak{} 的层级，远远凌驾于 \allowbreak{}\NovelInline{I03741}{\varphi(1,,,0)}\allowbreak{} 之上。

小队又想，如果有一个序数\allowbreak{}\NovelInline{I03742}{a>\Pi_\omega}\allowbreak{}，让\allowbreak{}\NovelInline{I03743}{\psi_a(\Omega_{a+n})=\varphi(1,,...n+1\text{个逗号},0)}\allowbreak{},如此，便让\allowbreak{}\NovelInline{I03744}{\psi_a(\Omega_{a+1})=\varphi(1,,0)}\allowbreak{},这样一来，\allowbreak{}\NovelInline{I03745}{D=\Omega_{a+1}}\allowbreak{}就成立了，虽然他们还有一些疑惑，但是没准有一天，a,D的谜题就解决了。

\chapter{第六十章：空台试炼，以战证道}\label{ux7b2cux516dux5341ux7ae0ux7a7aux53f0ux8bd5ux70bcux4ee5ux6218ux8bc1ux9053}

时光荏苒，倏忽数月。

扽魂殿后山的修行秘境始终灵气充盈，道韵绵长，日复一日地静心苦修，彻底淬炼了果糕小队众人的根基。数月以来，众人摒弃杂念，沉心钻研高阶序数法则，将多逗号系统、弱φ模式、I数阵与M数阵的所有对应规则彻底融会贯通。

从基础的单逗号、双逗号序数迭代，到高阶多逗号层级嵌套，再到 \allowbreak{}\NovelInline{I03746}{\Pi_\omega}\allowbreak{} 终极边界的推演奥秘，曾经晦涩难懂、困住无数高阶魂师的反射序数奥义，已然化作众人骨子里的修行底蕴。那些层层递进的层级规律、不动点与容许点的核心差异、数阵进位的专属奥义，尽数被他们吃透摸清，再无半分模糊滞涩。

修为的蜕变肉眼可见。

原本众人的魂力、精神力尚且局限于常规魂师体系，如今依托顶尖序数法则的加持，自身魂力本质完成了质的跃迁。每一缕魂力都镌刻着古老的序数道印，每一次魂力流转都遵循高阶反射规则，攻防、续航、感知、屏障，所有修行维度尽数暴涨。

唐三的蓝银领域融入序数折叠之力，可层层嵌套禁锢对手；小舞的瞬移身法契合数阵迭代节奏，虚实难辨，轨迹无迹可循；蔡须坤、蔡兴坤两位女修深耕体系推演，魂力运转精准无瑕，每一次输出都极尽法则极致；奥斯卡的香肠武魂叠加序数增幅，辅助效果突破原有极限；诺当、霍雨浩更是将精神力与序数奥义完美融合，感知覆盖千里，预判洞悉先机。

整支小队，已然彻底褪去昔日青涩，褪去了赛场切磋的稚嫩，坐拥碾压同阶的恐怖底蕴，在理论修行的道路上，已然走到了整片修行界的前沿。

可理论圆满，终究差了实战淬炼的锋芒。

这一日，众人如常围坐秘境之中，复盘 \allowbreak{}\NovelInline{I03747}{\Pi_\omega}\allowbreak{} 层级的终极规律，深究序数D的未解之谜时，一道清冷端庄的身影缓步踏碎云海，降临秘境之中。

正是梅天狸长老。

她一身素白殿袍，衣袂流云，眉眼平和却自带经年沉淀的威严，周身萦绕着厚重的序数道韵与杀伐气息。作为扽魂殿前代圣女，执掌邪祟清剿与宗门安防多年，她见惯了修行界的生死搏杀，眼光毒辣至极，一眼便看穿了众人当下的修行桎梏。

众人见状，纷纷起身行礼，神色恭敬。数月相处解惑，梅天狸与ddfg长老的悉心指点，早已让众人对两位长辈心生敬畏------没想到，梅天狸居然也突破了，实力达到了 \allowbreak{}\NovelInline{I03748}{Y(1,3,9)}\allowbreak{} 的恐怖层次。

梅天狸目光扫过众人，缓缓开口，声线沉稳，一语点破关键：

\begin{quote}
``我观你们数月苦修，已然彻底走完反射序数的修行道路，理论根基、法则领悟、魂力底蕴，放眼同阶近乎无敌，哪怕对比宗门老牌核心弟子，也毫不逊色。''
\end{quote}

话音微顿，她话锋一转，带着几分严苛：

\begin{quote}
``但你们有一个致命短板，也是所有潜心悟道修士最容易陷入的误区------理论已满，实战尚虚。''

``你们精通万千序数规则，能推演无尽高阶层级，可你们的力量，依旧停留在推演与苦修之中，未曾真正融入生死战场。''

``法则悟得再透，不能用于搏杀、破敌、守道，便只是纸上谈兵。真正的强者，既要通晓大道至理，更要敢在厮杀中证道，在攻防中淬炼本心。''
\end{quote}

众人闻言，皆是默然颔首。

他们心中清楚，数月以来皆是静坐推演、静心苦修，从未经历真正的高强度攻防搏杀。一身浩瀚磅礴的序数之力，确实缺少实战打磨，无法做到收发由心、极致运用，面对诡谲战局与凌厉攻势，难免会应变不足。

\begin{quote}
``今日，我便为你们补上这最后一块短板。''
\end{quote}

梅天狸袖袍轻挥，刹那间云海翻涌，天地灵气剧烈震荡。一股浩瀚无垠的魂力直冲云霄，撕开层层云层！

轰隆隆------

天际巨响轰鸣，无尽云雾汇聚翻涌，无数玄奥古老的序数道纹自虚空浮现，纵横交错，熠熠生辉。道纹链接天地，凝聚成一方无边无际的巨型悬空平台，自万米高空缓缓成型。

平台通体由云晶与道纹构筑，平整辽阔，一望无际，四周云海环绕，罡风凛冽，高空罡气呼啸流转，隔绝了大地灵气，自成一方独立战场空间。平台之上刻满层层叠叠的序数结界，可锁死战场空间，承受超阶魂力碰撞，任凭惊天动地的攻防爆发，也不会波及下方宗门疆域。

\begin{quote}
``随我来。''
\end{quote}

梅天狸轻身一跃，足尖踏空，稳稳落于高空平台中央，身姿卓然，气度凛然。

唐三众人不敢迟疑，各自催动魂力，踏空而起，紧随其后，尽数落在这座悬浮云端的巨型试炼台之上。

立足高台，凛冽罡风拂面，高空稀薄的灵气、肆虐的空间气流，瞬间让众人精神一振，周身毛孔尽数舒张。

梅天狸转过身，目光落于众人身上，沉声颁布试炼规则，字字铿锵，掷地有声：

\begin{quote}
``此地名为【序数云台】，是我扽魂殿专属的高阶试炼战场，结界稳固，空间封闭，可承受极致战力碰撞。''

``今日试炼，规则唯一且严苛------接下来，我会亲自对你们发起进攻，你们所有人联手，在一炷香的时间让我后退半步，便是试炼合格。''
\end{quote}

话音落下，众人神色一凛，瞬间摆正姿态，心中不敢有半分懈怠。

梅天狸乃是扽魂殿老牌长老，与ddfg长老自幼相伴，深耕序数法则修行无数岁月，修为深不可测，远非刚刚悟道的众人所能比拟。以一人之力对抗整支顶尖小队，还要支撑一炷香时长，这份试炼难度，远超众人以往任何一次历练。

就在众人凝神备战之时，虚空之中再度泛起灵气波动，几道熟悉的身影踏云而来，稳稳落于云台一侧，静静伫立观战。

为首之人，正是玉小刚！

紧随其旁的，还有弗兰德、赵无极等众人昔日的恩师。几人身姿沉稳，目光澄澈，带着期许与肃穆，静静注视着场中的果糕小队。

一年之前，众人尚在学院修行，受几位恩师悉心教导，打下坚实魂师根基。如今众人踏入扽魂殿，登顶序数修行大道，已然成长为一方天骄，再度与恩师重逢于高空试炼战场，心中百感交集。

\begin{quote}
``诸位老友，劳烦远道而来，为这群孩子坐镇试炼。''梅天狸对着几人微微颔首，语气谦和。
\end{quote}

玉小刚微微抬手，目光落在唐三众人身上，满是欣慰与郑重：

\begin{quote}
``梅长老客气了。这群孩子是我亲手带出的弟子，如今得以修行顶尖大道，有幸参与高阶试炼，我等自然要来见证一番。''
\end{quote}

他上前一步，声音洪亮，响彻整座云台，传入众人耳中：

\begin{quote}
``唐三，小舞，诸位！昔日我教你们魂师根基、武魂搭配、战术配合，今日这场试炼，不再是简单的魂力切磋、套路比拼！''

``梅长老的攻势，毫无保留，极尽杀伐之道！我与诸位师长今日在此坐镇，便是要让这场试炼，化作你们真正的生死战场！''

``摒弃切磋的侥幸，放下比试的温和！在这里，没有点到为止，只有全力攻防、绝境求生！你们所学的所有序数法则、魂力奥义、战术体系，全部要在这场实战中落地生根！''
\end{quote}

弗兰德适时开口，语气凌厉：

\begin{quote}
``记住！真正的修行战场，从不会给你们推演法则、梳理规律的时间！敌人不会等你们理顺层级、规划战术，只会极尽所能，破防绝杀！''
\end{quote}

赵无极沉声道：

\begin{quote}
``你们理论修行已然圆满，如今要练的，是瞬间法则运用、临场战术应变、绝境协同制衡！撑过这一炷香，你们的战力将会完成真正的蜕变！''
\end{quote}

几位恩师的话语，如同警钟轰鸣，彻底击碎众人心中最后的松懈。

众人神色尽数凝重，原本松弛的心神瞬间紧绷到极致。他们终于彻底明白，这场试炼绝非普通考核，而是一场复刻生死厮杀的极致淬炼。

梅天狸目光锐利，周身魂力悄然升腾，层层序数道纹在体表流转绽放，恐怖的威压席卷整座云台，空气瞬间凝固，凛冽的杀伐气息扑面而来。

\begin{quote}
``准备好了吗？''
\end{quote}

她轻声开口，语气平淡，却带着不容置疑的威严。

蔡须坤深吸一口气，眸光坚定，率先踏步而出，沉声喝道：

\begin{quote}
``全员列阵，备战！''
\end{quote}

话音落下，小舞、蔡须坤、蔡兴坤、奥斯卡、诺当、霍雨浩瞬间站位，众人默契十足，顷刻间组成攻防战阵。每个人周身魂力暴涨，序数法则熠熠生辉，层层结界、道道攻势瞬间蓄势待发。

云端风急，道纹轰鸣。

玉小刚一行人凝神注视，目光灼灼。

梅天狸身姿挺立，威压渐盛，属于高阶序数修行者的恐怖战力，彻底苏醒。

一炷香生死试炼，真正的实战大道淬炼，自此，正式开启！

\begin{center}\rule{0.5\linewidth}{0.5pt}\end{center}

云端罡风陡然狂乱，整座序数云台的道纹都随着梅天狸的魂力起伏轻轻震颤。唐三虽率先喝出备战指令，但场上真正的核心指挥却始终是果糕小队的队长蔡须坤。历经无数次体系推演与团战磨合，这位沉稳聪慧的女队长临场调度早已远超众人，冷静有度，滴水不漏。

\begin{quote}
``全员听令，标准联防阵型，梯度布防，守住正面威压！''
\end{quote}

蔡须坤声音沉稳有力，穿透呼啸罡风，精准落进每个人耳中。话音落下的刹那，原本松散站位的众人瞬间移位、错落排布，攻防阵型瞬息成型，默契度已然臻至顶尖小队水准。

可下一瞬，一股无边无际的恐怖威压骤然倾覆而下。

众人愕然抬头，只见立于云台中央的梅天狸始终静立不动，身姿挺拔如初，指尖未抬、袖袍未动，甚至连周身流转的序数道纹都未曾爆发半分攻势。她就这般凭空伫立在半空，双脚未踏台地，悬空而立，静谧得仿佛一尊俯瞰凡尘的大道神像。

没有招式，没有术法，没有魂力冲击。

唯有纯粹、厚重、碾压一切的高阶魂力威压，如亿万钧山岳轰然压落，死死笼罩整座云台。

咔嚓------

无形的气压撕裂空气，台面上细密的云晶纹路微微塌陷，众人周身的气流尽数凝固。每个人的胸膛都像是被巨手死死按住，呼吸骤然滞涩，胸腔发闷，头皮发麻，四肢百骸都传来沉重的压迫感，仿佛周身空间都被高阶序数法则彻底锁死。

这便是真正的境界差距。无需出手，仅凭底蕴威压，便足以压制一群吃透反射序数的顶尖天骄。

\begin{quote}
``稳住！别被威压打乱心神！催动全部武魂与魂灵，法则全开，固守阵型！''蔡须坤沉声厉喝，清丽容颜上凝着一抹坚毅，强行压下胸腔的闷窒感，率先爆发全部战力。
\end{quote}

嗡------！

璀璨的魂力光晕自她周身炸裂开来，两道截然不同的魂体力量同步苏醒。通体鎏金、纹路规整的篮球扽魂悬浮身前，球面流转着层层序数迭代光纹，自带极致的平衡制衡之力，可稳固阵型、分摊伤害；紧随其后，苍劲暴戾的狼王魂灵破体而出，青灰色狼瞳凛冽刺骨，漫天狼影萦绕周身，杀伐气息瞬间冲散部分凝滞气压，双重魂体相辅相成，攻防一体，让这位少女队长的气场陡然凌厉数倍。

蔡兴坤紧随姐姐步伐，清冷娇俏的面容尽是肃杀，瞬间催动本命武魂。漆黑凝练的黑子鸡扽魂展翅腾空，墨色魂力羽翼席卷四方，羽翼掠过之处，紊乱的高空罡风尽数被撕裂、中和，化作无形魂力反哺阵型，为众人撑开一方相对舒缓的攻防区域，稳稳补足团队防御短板。

一侧的卢本韦不敢怠慢，少女心神极致凝练，武魂全开！温润厚重的母猪扽魂通体泛着柔和的土黄色光晕，落地生根般的厚重气息瞬间铺开，层层厚实的魂力屏障叠加成型，极致的稳固防御牢牢托住整个小队的阵型基底，硬生生扛住了大半倾泻而下的威压，为众人争取到喘息之机。

\begin{quote}
``极致增幅，全员续航拉满！''
\end{quote}

奥斯卡低喝一声，香肠扽魂全力催动，一道道流光溢彩的增幅香肠飞速凝聚成型，精准落向每一位队友手中。力量增幅、速度增幅、精神稳固、魂力续航，多重buff瞬间覆盖全员，抵消威压带来的实力压制，让众人的战力始终维持在巅峰状态。

唐三眸光沉凝，不再保留半分实力。右手黑芒暴涨，沉重霸道的昊天锤轰然现世，锤身铭刻的古老纹路与序数道纹交相辉映，厚重的破阵之力席卷周身，随时准备轰碎一切来袭攻势；同时蓝银领域悄然铺开，细密的蓝银藤蔓扎根虚空，缠绕交织，形成密不透风的辅助禁锢网，攻防预备完毕。

小舞身形一闪，身姿轻盈蜕变，柔光亮起，柔骨兔真身开启。身形变得愈发灵动缥缈，周身力道柔韧十足，兼具极致速度与近身爆发，随时可借助瞬移身法突袭、补防、救援，成为阵型中最灵活的机动点。

诺当双目澄澈，精神力极致迸发，一圈圈明暗交错的磁暴环自周身扩散开来，狂暴的磁暴魂力纵横交错，形成大范围的力场制衡。无形磁力扭曲周遭空间，微微卸去梅天狸的碾压式威压，同时封锁云台空间，杜绝一切突袭可能。

最后一瞬，霍雨浩眼眸骤亮，纯白光芒刺破凝滞空气，顶级灵眸扽魂全开，浩瀚精神力铺天盖地席卷整座战场，千里战局尽收眼底，预判着梅天狸一切可能的出手轨迹。与此同时，多重魂灵齐齐现身，冰火两仪眼、雪帝、冰碧帝皇蝎等魂灵气息交融，极致的冰控与精神震慑之力层层叠加，稳固团队防线。

刹那之间，果糕小队全员战力彻底解锁，八大武魂、多重魂灵的力量交织融汇，序数法则层层联动，攻防阵型固若金汤。

云端战场之上，少年天骄们尽数凝神备战，气息凝于一点，直面那伫立虚空、不动如山的恐怖对手。

云台侧方，玉小刚、弗兰德几人屏息凝望，眼底满是惊叹。这群孩子不仅理论修行冠绝同辈，武魂底蕴、团战默契、临战心态，早已远超他们的想象。

半空之中，梅天狸静静俯瞰下方严阵以待的众人，淡漠的眼眸中终于掠过一丝微不可察的赞许，沉寂已久的周身魂力，开始缓缓涌动。

\begin{center}\rule{0.5\linewidth}{0.5pt}\end{center}

不等梅天狸率先出手，阵型前方的蔡须坤、蔡兴坤两位少女姐妹已然默契达成共鸣。

\begin{quote}
``合击技------铁山靠！''
\end{quote}

两道清冽喝声同时炸响，姐妹二人身形骤然狂飙突进。蔡须坤的篮球扽魂光影暴涨，狼王魂灵的凶煞之力尽数灌注少女纤巧却极具爆发力的肉身，体魄瞬间绷紧，柔韧线条贲张，带着制衡一切的厚重冲势碾压而出；蔡兴坤的黑子鸡扽魂双翼收拢，全部魂力凝练于纤细肩背，舍弃远程牵制，凝聚极致爆发力。

姐妹二人肩背合一，魂力、魂灵、序数法则彻底交融归一，整套合击技将数月苦修的高阶序数之力尽数引爆！

轰！！

一声沉闷炸裂的巨响震荡整片序数云台，强横的冲击波以二人为中心席卷四方，台面云晶道纹层层震颤，狂风倒卷，气流爆裂。这是二人倾尽底蕴的全力一击，融汇多逗号系统的迭代之力，层层叠加、层层增幅，足以碾碎无数高阶魂师的防御，是他们现阶段最强的肉身合击杀招。

劲风滚滚，威势滔天，两道合一的磅礴劲气直直撞向悬空静立的梅天狸！

然而下一秒，令人窒息的一幕降临。

面对足以崩山裂地的合击猛攻，梅天狸身形纹丝不动，悬空伫立的身姿稳如万古山岳，脸上无波无澜。她甚至没有挪动半步，仅仅是随意抬起白皙纤细的右手，轻描淡写的横挡在身前。

没有惊天动地的碰撞，没有剧烈的魂力对冲。

咔嚓------噗！

姐妹二人倾尽一切的铁山靠攻势，狠狠撞在梅天狸单薄的手掌之前。极致的爆发力瞬间炸开，却像是一头狠狠撞在了亘古不灭的万丈山岳之上，所有冲击、所有力道、所有序数叠加的法则之力，在接触的瞬间尽数消解、归零！

汹涌的攻势如同泥牛入海，连一丝涟漪都无法掀起。

蔡须坤、蔡兴坤只觉得一股浩瀚无边、无从抗衡的反震之力骤然传回肉身，少女纤细的身躯巨震，浑身气血剧烈翻腾，骨骼阵阵发麻，向前冲击的身形瞬间被死死锁死，再难寸进分毫。她们拼尽全力催动的绝杀合击，在梅天狸面前，脆弱得不堪一击。

二人慌忙收力倒飞而出，裙袂翻飞，落地后踉跄数步才勉强站稳，清丽的脸庞上满是震撼与难以置信。

\begin{quote}
``太强了\ldots\ldots{}''蔡须坤低声喃喃，清丽眉眼间满是动容，``长老的魂力根本看不到尽头，不像山岳有顶、沧海有涯，反倒如同无边大海，浩瀚无垠、包容万法，我们的全力一击，连她的魂力表层都无法撼动。''
\end{quote}

一旁的蔡兴坤轻点螓首，心有余悸。刚刚那一瞬间，她清晰感受到了层级上的绝对碾压，那是法则、底蕴、境界全方位的天堑差距，根本无法靠战术、合击、爆发力弥补。

\begin{center}\rule{0.5\linewidth}{0.5pt}\end{center}

就在众人心神震颤、尚且未从合击被碾压的落差中缓过神之际，梅天狸淡漠的眸光骤然彻冷，眼底所有温润尽数褪去，取而代之的是横跨万古的死寂与苍茫。她周身原本温顺蛰伏的序数道纹瞬间暴走、倒卷、冲天而起，亿万道深邃鎏金碧海霞光撕裂云层、横贯九天，将万米高空彻底浸染成厚重、静谧的深海玄色，连天地间的日月灵光都被强行压制、黯淡无光！

嗡------------------！！！

这不是魂力轰鸣，是高阶反射序数大道的本源道音震荡。整座序数云台瞬间被漆黑的空间裂隙包裹，坚不可摧的云晶台面以肉眼可见的速度龟裂、塌陷、消融，台面铭刻的万千防御道纹、结界符文、序数禁制成片崩碎、化作飞灰，连扽魂殿专为高阶试炼打造的锁空阵法都在极速瓦解。方圆万里高空灵气彻底逆流、坍缩、湮灭，罡风归零、流云静止、天地死寂，整片空间的修行规则被强行改写，进入专属梅天狸的序数法则领域！

梅天狸窈窕的人形躯体在漫天霞光中层层虚化、褪尽凡质，数年克制、万千收敛的大道底蕴彻底解禁。一股源自太古、诞生于序数本源、执掌清邪镇道的无上威压轰然倾覆天地，它不暴戾、不狂躁，却带着维度层面的绝对碾压，如同凡尘蝼蚁直面苍穹万界，从灵魂根源锁死所有人的反抗意志，让每一寸血肉、每一缕精神力都本能战栗、臣服。

轰隆隆------------------------！！！

\chapter{第六十一章：四位一体，108母猪大阵}\label{ux7b2cux516dux5341ux4e00ux7ae0ux56dbux4f4dux4e00ux4f53108ux6bcdux732aux5927ux9635}

万丈云层轰然崩碎、化为虚无，虚空大面积塌陷褶皱，漆黑的空间裂痕纵横交错蔓延至视野尽头，整片高空天地濒临崩解，恐怖的空间乱流被无形道力强行禁锢、无法肆虐，唯有极致的法则压迫笼罩四方！

一尊真正意义上的大道真身，自虚空裂隙与碧海道霞中缓缓苏醒、现世！

足足三百余米的庞大海狸真身横亘苍穹，体量远超寻常山岳，身躯巍峨厚重，遮蔽整片高空天光，将整座序数云台彻底吞入沉沉道域阴影之中。它通体覆盖致密如神金的深海玄青皮毛，每一根绒毛都镌刻着Y序列本源道纹，早已超脱 \allowbreak{}\NovelInline{I03749}{\Pi_\omega}\allowbreak{} 反射序数的极限桎梏。梅天狸的真实修为早已踏足至高的Y序列强度，这是彻底凌驾于反射序数体系之上的顶级修行层级，所谓多逗号迭代、I/M数阵、\allowbreak{}\NovelInline{I03750}{\Pi_\omega}\allowbreak{} 终极边界，尽数属于低阶反射序数规则，在Y序列的本源大道面前，如同萤火比之皓月，低级且渺小。粗壮堪比山岳的四肢扎根云海，踏碎虚空乱流，厚重的躯干自带万古不动的镇压势，每一次轻微的呼吸都牵动Y序列天地法则迭代。一双竖瞳深邃如寂灭深海，冷寂、苍茫、无悲无喜，俯瞰下方众人之时，如同高阶大道俯瞰低维规则，层级之差早已不是量级差距，而是维度上的绝对鸿沟！

这绝非普通魂师的武魂真身、绝非寻常魂灵显化！这是梅天狸深耕大道本源，突破反射序数天花板，踏入Y序列高阶领域后，以自身千年道基凝练的序列道体真身。常规修士穷尽一生钻研的反射序数体系，只是修行路上的基础法理，而Y序列是凌驾其上、可绝对支配、归零所有低阶序数规则的至高体系。常规魂师的战力靠魂力、靠魂环、靠武魂，而它的每一寸躯体、每一缕气息，都是Y序列固化的至高法则，对果糕小队现阶段的反射序数修为，是彻彻底底、无解的降维打击！

吼------------------------------------！！！

一声横贯九天十地、震彻万古的道吼轰然炸开！

这不是兽鸣，是序数大道的本源咆哮！肉眼可见的深蓝色巨型法则音浪层层叠叠、摧枯拉朽席卷全场，所过之处虚空直接塌陷湮灭、灵气彻底归零、残余结界尽数粉碎。原本濒临破碎的云台结界瞬间被压得极致凹陷，漫天防御道纹彻底泯灭，整座悬空战场摇摇欲坠，若非扽魂殿基底阵法死死支撑，早已伴随这声道吼彻底崩毁！

极致的法则威压瞬间锁死整片空间，形成密闭的镇压道域！

果糕小队全员瞬间神魂巨震、肉身僵死！

众人脑海中熟记于心的多逗号系统、弱φ模式、I/M数阵规则瞬间紊乱崩碎，精神力被强行压制滞涩，魂力运转卡顿滞停，浑身气血疯狂逆流翻涌。所有人叠加到极致的武魂屏障、磁暴力场、蓝银领域、精神壁垒，在道域威压下如同薄纸轻纱，瞬间层层黯淡、龟裂、溃散！哪怕全员尽数催动毕生修为、叠加所有序数防御，依旧被这股无解的大道之力死死镇压，连站直身躯都需要倾尽全部意志力！

\begin{quote}
``稳住！全员固守阵型！''
\end{quote}

蔡须坤脸色剧变，厉声嘶吼。

全员瞬间催动全部防御手段！卢本韦母猪扽魂屏障全力撑开，奥斯卡增幅buff拉满，诺当磁暴力场层层阻隔，唐三蓝银领域编织天罗地网，霍雨浩精神屏障锁死心神，蔡氏二女死守前路，小舞机动补位稳固阵型！

即便众人早已全开武魂魂灵，融汇所有序数修行底蕴，依旧被这一声咆哮震得节节败退。

噔！噔！噔！

每个人的双腿剧烈发颤，全身骨骼密布承压细纹，皮肉被道域压得紧绷刺痛，裙袂衣袍被狂暴的法则气流撕扯得猎猎破碎，嘴角不断溢出细密血痕。她们每后退一步，台面云晶都会被踩出细碎裂痕，所谓的圆满序数修为、顶尖小队战力，在这尊太古道体面前，脆弱得不堪一击，差距悬殊到让人心底生出极致的绝望与无力。

云台侧边，玉小刚、弗兰德、赵无极几人彻底浑身僵硬、瞳孔骤缩，心神震颤到极致。几人征战半生、阅尽天下强者，见过魂帝、魂圣、封号斗罗的惊天对决，却从未见过这种脱离常规战力体系的碾压！无需动手、无需攻杀，仅凭真身道体与本源道域，便彻底压制七位吃透高阶序数法则的顶尖天骄。这已经不是修为差距，而是修行维度、大道层级的天堑鸿沟，让旁观者彻底失语、心生敬畏！

\begin{center}\rule{0.5\linewidth}{0.5pt}\end{center}

狂风肆虐，道纹崩碎，整座序数云台的压制力依旧层层叠加、持续攀升。

巍峨三百米海狸太古真身盘踞长空，镇压万方，周身流转的Π层级序数神纹自成绝对防御道壁，万法不侵、诸攻无效。而在这尊神山般的巨兽脊背正中，梅天狸一袭素白殿袍纤尘不染、衣袂轻垂，身姿挺拔从容，淡漠伫立在无尽道韵中央，仿佛执掌这片序数天地的神祇，俯瞰芸芸天骄，无喜无悲。

任凭下方果糕小队全员倾尽底牌、轮番猛攻、法则尽出，漫天绚烂霸道的魂力攻势、序数杀招、合击秘术层层叠叠轰涌而上，一旦触及海狸真身外围的道纹屏障，便瞬间被法则归零、序数消解，连最表层的皮毛都无法撼动，连一丝涟漪、一点微光都无法留存，彻底沦为无用之功。

唐三全力抡动昊天锤，裹挟层层序数折叠之力，锤影破空轰鸣，重重砸落，只溅起点点细碎的法则微光，连一丝印痕都无法留下；小舞借助瞬移身法近身突袭，柔骨极致力道尽数爆发，却被无形的大道屏障轻轻弹开；蔡须坤、蔡兴坤二女再度联动发力，篮球扽魂制衡之力、黑子鸡魂力翼刃轮番轰击，所有叠加的高阶序数招式尽数落空，如同投石入海，掀不起半点风浪。

其余众人的磁场禁锢、增幅秘术、领域封锁齐齐全开，无数绚烂霸道的魂力招式铺满整片云台，可落在海狸真身厚重的法则防御上，尽数被无声消解、层层归零。

这一刻的碾压直白而残酷：众人苦修数月、日夜推演吃透的反射序数奥义、多逗号迭代体系、I/M数阵规律，只是低阶凡人修行的顶端法理；而梅天狸的海狸真身，是Y序列至高本源的固化规则。低维的反射法则，在高阶Y序列大道面前天然失效，凡人悟道再精，终究触不到真正的大道层级，二者隔着不可逾越的修行维度天堑！

久攻不破，阵型承压越来越重，倒退的步伐愈发艰难，所有人心底都泛起浓浓的无力感。

\begin{quote}
``全员稳住阵型！我开启全域精神探测共享，捕捉对方法则破绽！''
\end{quote}

危急关头，霍雨浩眸光骤然凛冽，灵眸扽魂全力暴走，浩瀚无垠的精神力瞬间铺展笼罩整座高空战场。极致细密的精神探测穿透层层狂暴风压与紊乱灵气，精准扫描海狸真身的每一寸法则纹路、每一处气息流转节点，同时将捕捉到的所有动态、破绽轨迹实时共享给每一位队友，帮众人精准锁定攻击落点、规避威压死角。

有了精神共享的精准加持，众人压力骤减，攻防节奏瞬间规整，不再盲目强攻。

\begin{quote}
``老韦，正面强攻！牵制真身壁垒！''蔡须坤清冷喝声穿透罡风。
\end{quote}

卢本韦应声颔首，少女眸光坚毅，不再保留半分魂力，本命母猪扽魂全力催动！

嗡的一声厚重震颤，土黄色魂力暴涨铺开，一头体型敦实、防御极致的巨型母猪武魂凌空显化，身躯饱满厚实，带着落地生根的厚重霸势，直面前方恐怖巨兽轰然冲击而出！

这头母猪武魂足有十几米之高，已然是卢本韦现阶段所能催动的极限体型，承载着极致的厚重冲击与破防之力，是她最顶尖的牵制招式。

可在横跨数百米、如山岳横亘苍穹的海狸真身面前，这十几米大小的母猪武魂渺小得可怜，如同蝼蚁撼巨峰、细石撞沧海，悬殊的体型差距触目惊心。

渺小的土黄色武魂带着倾尽一切的力道狠狠撞向海狸真身的躯干，下一秒，无声无息的法则屏障微微震颤，仅仅一缕外泄的道韵之力，便将母猪武魂的冲击尽数消解。

嘭------！

一声闷响，土黄色魂力瞬间溃散，母猪武魂应声黯淡、层层虚化，被直接弹飞回阵。卢本韦身躯巨震，气血翻腾，硬生生逼退三步，堪堪稳住身形，眼底满是震撼。

全力一击，连对方的皮毛都未曾触及。

高台之上，梅天狸静立巨兽脊背，神色依旧淡然无波，仿佛方才那铺天盖地的全员猛攻、倾力牵制，不过是清风拂衣，不值一提。

\begin{center}\rule{0.5\linewidth}{0.5pt}\end{center}

绝境之下，众人再无保留。

卢本韦咬牙凝力，声线带着极致的决绝，高声喝道：

\begin{quote}
``奥斯卡，联动！终极合击------108母猪大阵！''
\end{quote}

\begin{quote}
``收到！全力联动！''
\end{quote}

奥斯卡面色惨白，早已被道域威压压得气血翻涌，此刻却骤然绷紧心神，燃烧自身魂力本源，彻底解禁香肠扽魂的全部底牌！

刹那间，整片序数云台的土黄色魂力疯狂暴涨、铺天盖地，密密麻麻的母猪武魂自虚空源源不断显化成型，一头、十头、百头\ldots\ldots 转瞬之间，整整一百零八头体型敦实、承载厚重序数之力的巨型母猪凌空排布，以极精准的阵道轨迹环绕全场，凝聚成一圈硕大无比、规整森严的圆形大阵。

一百零八头母猪首尾相接、环锁天地，阵眼流转层层叠叠的反射序数灵光，将所有残余的低阶序数规则尽数收拢、叠加、熔炼。最诡异且震撼的是，每一头母猪的鼻孔里面，都笔直插着一根粗壮凝练、流光溢彩的大香肠，荤素双魂完美契合，魂力纹路彼此交融、互补共生，形成独一无二的攻防合击阵基。

\begin{quote}
``全员输功！灌注双阵核心！''蔡须坤当即沉声下令。
\end{quote}

话音落下，唐三、小舞、蔡兴坤、霍雨浩、诺当五人毫无保留，彻底撤去自身防御，周身魂力尽数剥离、倾泻而出，化作五道雄浑的魂力洪流，尽数汇入卢本韦与奥斯卡体内。

一瞬间，两人身躯被璀璨魂力包裹，全身经脉胀痛欲裂，浑身汗水淋漓，脸色愈发苍白。这是榨干自身修为本源的极限献祭，每一缕魂力都被压榨到极致，没有半分留存。

奥斯卡浑身剧颤，双目布满血丝，已然透支到极致。他倾尽体内最后一丝魂力本源，咬紧牙关猛然催发，苍穹之上，一根横贯长空大香肠轰然凝聚成型！

咻------！

香肠迅速插入母猪鼻子里面，母猪大哼一声，这母猪鼻子本就顺滑，奥斯卡两根香肠进入后，母猪下意识地缩紧了鼻子，两根香肠让母猪无法摄入氧气，只能产生无氧呼吸，

\NovelDisplay{M00686R01}{%
\text{C}_6\text{H}_{12}\text{O}_6 \longrightarrow 2\,\text{C}_2\text{H}_5\text{OH} + 2\,\text{CO}_2%
}

的反应在母猪体内爆发，瞬间，母猪的鼻孔产生大量气体，极短时间内，产生大量气体导致鼻子内部气压升高，最终香肠发射出去了，这一次，香肠经过第一个母猪的鼻子里面之后，诺当立刻放出磁暴环击中母猪头部，瞬间母猪就高潮了。之后在剧烈的快感下把香肠喷出去。获得巨大的速度，由于母猪是兴奋状态，香肠喷出去的速度很快。然而，喷出去香肠之后，有到了另一个母猪鼻子里面，经过108头母猪的阵型，最后一个母猪承载着最强的力量。

一百零八头母猪同时迸发通体序数灵光，层层叠叠的叠加推力源源不断灌注核心香肠，一重又一重的阵力不断叠加、爆发、推送。一百零八道厚重的土系序数之力、一百零八道增幅香肠的精纯灵力、一百零八层阵道迭代规则，尽数汇聚于一点！嗡------！！！诺当和卢本韦的配合恰到好处，母猪鼻子里面的气压达到临界值的时候，母猪高潮，下意识缩紧鼻子，之后喷出去。每一次配合都恰到好处

香肠获得了最大加速

之后，蔡须坤用出了本源的力量，处女之力，当处女之光直接作用于香肠上面的时候，将香肠变成了诡异的蓝色

这一记倾尽全队之力的终极合击，硬生生冲破了无数修士终生无法逾越的 \allowbreak{}\NovelInline{I03751}{0}\allowbreak{}、\allowbreak{}\NovelInline{I03752}{111}\allowbreak{}、\allowbreak{}\NovelInline{I03753}{222}\allowbreak{} 序数壁垒，打破了反射序数的层级桎梏！不止如此，它的力量阈值疯狂攀升，直接到达了众多人可望不可及的 \allowbreak{}\NovelInline{I03754}{0\;111\;222\;322\;422}\allowbreak{}，隐隐触摸到了传说中 \allowbreak{}\NovelInline{I03755}{0\;1111}\allowbreak{} 的高阶序数门槛，已然半只脚踏入超越常规反射序数的全新领域，是果糕小队此刻能打出的最强一击！

蓝色辉光横贯长空，气场滔天，直面山岳般的海狸真身轰然射杀而去！

云海之巅，梅天狸伫立巨兽脊背，淡漠的眼眸终于掠过一丝微弱的讶异。

她未曾动用任何杀招，仅仅轻轻摆动海狸真身那粗壮绵长、覆满Y序列神纹的巨型尾巴！

漆黑深海纹路缠绕的巨型狸尾破空横扫，带着Y序列至高法则的镇压之力，裹挟漫天道韵，径直迎上那道突破序数壁垒的蓝色香肠攻势！

下一瞬，天地倾覆，万象炸裂！

轰------------------------------------！！！

极致的力量碰撞炸开无尽璀璨的光潮，蓝、金、玄黑三色法则光晕疯狂迸发、席卷九天！恐怖的能量冲击波瞬间撕裂剩余的云台结界，高空云层被尽数清空，整片虚空剧烈震颤、塌陷、翻涌！

这等层级的对撞早已超越常规魂师对战的极限，哪怕是远在十里之外的宗门民居、山林旷野，所有生灵都能清晰感受到脚下大地的剧烈震动，一股磅礴浩荡、窒息压抑的能量风暴席卷四野，家家户户门窗震颤、灯火摇曳，寻常修士无不骇然抬头，凝望高空那片倾覆天地的璀璨光海！

序数云台之上，强光刺眼，劲风灭世，毁灭性的能量浪潮以碰撞点为中心疯狂席卷四方，果糕小队众人尽数被狂暴冲击波掀飞，重重摔落在云晶台面之上，浑身经脉刺痛、魂力透支枯竭，人人气血翻涌、衣衫破碎，连抬手的力气都没有，只能怔怔凝望那一场低维极致之力与高阶Y序列大道的终极碰撞。

\begin{center}\rule{0.5\linewidth}{0.5pt}\end{center}

云台侧边，观战的弗兰德与赵无极目睹这惊天动地的一击，皆是神色复杂，满心感慨，久久无法回神。

弗兰德长叹一声，眼底满是惊叹与动容，沉声感慨道：

\begin{quote}
``回想当初在学院特训之时，我们为打磨他们的抗压能力，仅仅是随意释放出自身的魂力威压，便足以将这群孩子逼至绝境，险象环生，连立足攻防都无比艰难。''
\end{quote}

赵无极重重点头，粗犷的面容上满是震撼，语气满是唏嘘：

\begin{quote}
``短短数月光阴，翻天覆地啊。如今他们施展出的这记108母猪大阵，层级之高、威力之狂暴、序数突破之恐怖，已然完全超越了我们的修为上限。换做是我们二人正面硬接，根本无力抵挡，足以被这一招彻底碾压。''
\end{quote}

\begin{quote}
``从被我们随意威压逼入绝境，到如今打出能够碾压我们的绝杀招式，这群孩子的修炼速度、成长跨度，简直堪称妖孽，世间罕见，当真不负天才之名！''弗兰德由衷赞叹，目光死死望着场中满身狼狈却依旧傲骨未折的众人，满心欣慰。
\end{quote}

二人的感慨发自肺腑，也道尽了果糕小队的逆天成长，可唯有玉小刚眸光深沉，透过这惊天一击，彻底看透了梅天狸深不可测的实力层级，心中暗自推演、剖析着双方的大道差距。

全场所有人，包括突破极限的果糕小队，尽数深耕于反射序数体系，无论多逗号迭代、I/M数阵嵌套，还是 \allowbreak{}\NovelInline{I03756}{0}\allowbreak{} \allowbreak{}\NovelInline{I03757}{111}\allowbreak{} \allowbreak{}\NovelInline{I03758}{222}\allowbreak{} 乃至 \allowbreak{}\NovelInline{I03759}{0\;1111}\allowbreak{} 的序数壁垒，统统属于低维反射规则的迭代范畴，哪怕推演至极致，也依旧被锁死在低阶大道框架之内。

而梅天狸的 \allowbreak{}\NovelInline{I03760}{Y(1,3,9)}\allowbreak{} 序列修为，是完全超脱反射序数体系的高阶维度大道，属于层级碾压式的规则跃迁，而非简单的魂力数值差距。常规修士的战力提升，是在原有规则内不断叠加、迭代、突破上限，可Y序列强者，拥有改写、归零、压制一切低阶序数法则的绝对权柄。

更令人惊悚的是，方才迎击终极香肠的狸尾横扫，仅仅是她真身无意识的道韵外泄，只动用了 \allowbreak{}\NovelInline{I03761}{Y(1,3)}\allowbreak{} 的基础力量，尚且未触及自身 \allowbreak{}\NovelInline{I03762}{Y(1,3,9)}\allowbreak{} 的完整道体底蕴、未施展任何序列专属杀招、未催动半点Y序列结界领域。

简单来说，果糕小队是倾尽低维大道的全部极限，燃烧本源、献祭魂力、全员联动，打出了反射序数的天花板之力；而梅天狸，只是被动溢出了高阶序列的一丝基础力量。

一方是拼尽全力的极致绝杀，一方是随手为之的本能防御，二者的差距，早已不是修为深浅、战力高低，而是低维规则与高维大道的本质隔绝。若是梅天狸全力催动 \allowbreak{}\NovelInline{I03763}{Y(1,3,8)}\allowbreak{} 完整道力，整片反射序数构筑的战力体系，会被瞬间全盘瓦解、彻底归零，根本不存在抗衡的可能。

这一记绝杀的力量层级早已被精准定义。众人拼死突破的 \allowbreak{}\NovelInline{I03764}{0}\allowbreak{} \allowbreak{}\NovelInline{I03765}{111}\allowbreak{} \allowbreak{}\NovelInline{I03766}{222}\allowbreak{} 序数壁垒，对应序列模板为 \allowbreak{}\NovelInline{I03767}{Y(1,2,4,8,15)}\allowbreak{}，而他们隐约触摸到的 \allowbreak{}\NovelInline{I03768}{0\;1111}\allowbreak{} 高阶门槛，正是 \allowbreak{}\NovelInline{I03769}{Y(1,2,4,8,16)}\allowbreak{} 的极限阈值。这已是目前他们所能企及的绝对巅峰，是这群少年少女倾尽苦修底蕴、团队协同换来的终极极致。

可他们面对的梅天狸，修为根基早已定格在碾压一切低阶序列的至高层级------\allowbreak{}\NovelInline{I03770}{Y(1,3,9)}\allowbreak{}。

\begin{center}\rule{0.5\linewidth}{0.5pt}\end{center}

惊天动地的能量爆鸣缓缓消散，漫天璀璨光潮逐步褪去，动荡的虚空慢慢归于平稳。

让所有人难以置信的一幕悄然发生。

三百余米的巍峨海狸真身依旧伫立云海，通体玄青神纹流转不息，周身Y序列道域稳固如初，没有一丝破损、半分黯淡。方才足以突破反射序数极限的绝杀一击，倾尽百猪阵力、全队本源的蓝色辉光，最终只换来一个微乎其微的结果------

镇压万方的太古巨兽身躯，在极致对冲的反震力下，仅仅微微向后退却半步。

仅此而已。

半寸的位移，便是这无解低维绝杀所能撼动至高Y序列大道的全部极限，渺小、微弱，却又足以让全场旁观者心神巨震。玉小刚几人瞳孔骤缩，彻底失语，他们终于真切看懂了序列层级的恐怖差距：低阶序列的极致爆发，在高阶Y序列的本源法则面前，终究只是徒劳挣扎。

随着威压缓缓收敛，漫天碧海霞光尽数回笼，笼罩天地的道域屏障悄然消散。海狸巍峨的真身开始层层虚化、褪尽巨影，磅礴的太古道韵缓缓内敛，最终重新凝聚为纤细挺拔的白衣人影，梅天狸身姿卓然，稳稳立在云台中央，神色依旧淡漠从容，不染半点尘埃。

她抬眸望向瘫落台面、魂力耗尽、满身狼狈却眼神坚毅的果糕小队，白皙指尖轻轻一抬。嗡------

一道温润纯粹的纯白光束自天际垂落，精准笼罩整支小队，轻柔包裹每一个人的身躯。极致温和的治愈道韵流淌四肢百骸，方才对战留下的经脉损伤、气血逆流、骨骼承压细纹尽数修复，透支枯竭的魂力被缓缓补足，撕裂的衣袍、体表的伤痕瞬间消融殆尽。

短短数息时间，众人满身伤势彻底痊愈，状态重回饱满，唯有精神深处的疲惫，无声诉说着方才绝境死战的消耗。

梅天狸缓步上前，清冷的声音响彻整座序数云台，字字清晰，振聋发聩：

\begin{quote}
``你们的配合、战术布局、绝境爆发力，还有突破自身极限的执念，远超我的预料，确实让我倍感惊喜。''
\end{quote}

话音微顿，她目光锐利，语气添了几分凝重与教诲：

\begin{quote}
``但你们必须牢牢记住今日的对战感悟------在绝对的大道实力面前，一切战术、配合、低阶极致的爆发，终究都是虚无。''

``Y序列的层级压制，是维度的隔绝，是法则的宿命。就算你们到达 \allowbreak{}\NovelInline{I03771}{Y(1,2,4,8,16)}\allowbreak{} 极致之力，连我随手 \allowbreak{}\NovelInline{I03772}{Y(1,3)}\allowbreak{} 的道韵防御都无法真正突破，这便是修行最残酷的真相。''

``术可补天，道可镇世。精妙的战术能让你们同阶无敌，唯有不断深耕大道、突破序列天花板、提升自身绝对修为，才能让你们真正立于不败之地。''

``今日试炼合格，望你们铭记此战的差距与震撼，往后潜心苦修，挣脱反射序数的桎梏，向着真正的高阶序列大道奋勇攀登。''
\end{quote}

\chapter{第六十二章：登云朝圣，良序圣殿}\label{ux7b2cux516dux5341ux4e8cux7ae0ux767bux4e91ux671dux5723ux826fux5e8fux5723ux6bbf}

序数云台的余波缓缓散尽。

漫天溃散的法则光晕一点点归于沉寂，震颤不休的云晶台面彻底平稳，崩裂的虚空裂隙缓缓愈合。方才那场横跨维度的极致碰撞，终究彻底落幕，只余下空气中残留的淡淡道韵，无声印证着方才那绝望又璀璨的一战。

果糕小队众人立于高台之上，周身伤势尽数被梅天狸的治愈光束修复，体魄完好如初，可心神深处的震撼与敬畏，却久久无法平息。

方才他们倾尽全队底蕴、燃烧本源魂力催发出的一百零八母猪大阵，突破层层反射序数壁垒，触摸到Y序列下位门槛，已然是同辈修士可望而不可及的巅峰造诣。可在梅天狸随手一记 \allowbreak{}\NovelInline{I03773}{Y(1,3)}\allowbreak{} 层级的狸尾横扫面前，依旧形同蝼蚁撼山，仅仅逼得对方半步后退。

这种维度层面的绝对差距，远比任何重创溃败，更让人心生敬畏。

弗兰德、赵无极与玉小刚三人伫立侧边，望着身姿挺拔、气度愈发沉稳的一众少年少女，眼中满是欣慰与感慨。短短数月，这群孩子已然彻底挣脱了普通魂师的修行桎梏，踏入了序数大道的全新领域，成长速度之快，堪称逆天。

\begin{center}\rule{0.5\linewidth}{0.5pt}\end{center}

试炼尘埃落定，云台之上紧绷的战斗气氛彻底消散，难得迎来片刻安稳。

玉小刚、弗兰德、赵无极三人走上前来，围在果糕小队众人身侧，眼底满是欣慰与感慨。昔日稚嫩懵懂、连基础魂力掌控都略显生涩的孩童，如今已然成长为能撬动序数壁垒、撼动高阶大道的顶尖天骄，脱胎换骨的蜕变，让三位恩师心中百感交集。

玉小刚望着眼前身姿挺拔、气度沉稳的众人，缓缓开口，语气温和又厚重：

\begin{quote}
``你们方才的战斗，我们全程看在眼里。从被威压碾压的绝境，到全员联动突破序数极限，你们的成长、韧性与默契，早已远超我当初的预期。''
\end{quote}

他顿了顿，目光扫过每一个人：

\begin{quote}
``我从前教你们的，是魂师根基、武魂配合、战术博弈，是世俗魂道的方寸之术。而如今你们踏入扽魂殿，钻研序数大道，接触Y序列高阶法则，已然走上了一条我们毕生都无法触及的至高修行路。''
\end{quote}

弗兰德笑着感慨，语气满是唏嘘：

\begin{quote}
``想当年在史莱克特训，随便一点魂力威压就能把你们逼得节节败退、险象环生。谁能想到短短一年多，你们便能施展出108母猪大阵这种突破维度极限的招式，连我们二人的修为都能彻底碾压。长江后浪推前浪，你们的天赋与毅力，当真逆天。''
\end{quote}

赵无极重重点头，粗犷的面容少了几分平日的严肃，多了几分柔和：

\begin{quote}
``但你们切记，今日梅长老展现的层级压制，才是修行真正的天花板。你们突破的 \allowbreak{}\NovelInline{I03774}{0\;111\;22}\allowbreak{} 序数壁垒，在真正的高阶Y序列面前，依旧渺小不堪。越是见识过高深大道，越要懂得谦卑，不可因一时突破而自满。''
\end{quote}

蔡须坤微微颔首，身姿恭谨，轻声回应：

\begin{quote}
``老师教诲，我们铭记于心。今日一战，我们才真正明白，我们所学的反射序数，不过是大道皮毛，在绝对实力面前，所有技巧与极限突破都有穷尽。''
\end{quote}

唐三亦是正色开口：

\begin{quote}
``我们不会沉溺于当下的成绩，今日的差距，便是我们日后苦修的动力。''
\end{quote}

小舞、霍雨浩、卢本伟、蔡兴坤等人纷纷附和，眼底褪去了年少浮躁，多了沉稳与敬畏。

玉小刚见状，心中愈发宽慰，缓缓抬手：

\begin{quote}
``去吧，你们已然通过试炼，拥有接触扽魂殿核心的资格。跟随梅长老前行，好好拜见大长老，潜心修习更高深的大道，未来的你们，注定远超我们这些故人。''
\end{quote}

众人齐齐躬身行礼，郑重辞别三位恩师。

\begin{center}\rule{0.5\linewidth}{0.5pt}\end{center}

也就在此时，梅天狸眸光平和，扫过面前气息沉稳、眼底依旧有锋芒的众人，淡淡开口：

\begin{quote}
``你们方才的表现，已然完全通过了本次生死试炼。团队默契、绝境韧性、法则突破能力，皆已达标，足以正式接触扽魂殿的核心层级。''
\end{quote}

话音落下，她袖袍轻扬，身姿缓缓腾空，立于云台高空，望向云海之上更辽阔的九天穹顶。

\begin{quote}
``你们已然窥见了Y序列的门槛，触摸到了高阶大道的皮毛。今日，我便带你们前往我扽魂殿至高圣地，拜见我的授业恩师，也是执掌整座扽魂殿底蕴、统御序数大道的大长老------test\_alpha0。''
\end{quote}

此言一出，众人心神齐齐一震，下意识屏住呼吸。

梅天狸已是踏足 \allowbreak{}\NovelInline{I03775}{Y(1,3,9)}\allowbreak{} 高阶序列的顶尖强者，稳压所有反射序数修士，深不可测。而能成为她的师父、身为扽魂殿大长老的存在，修为层级必然早已超脱众人认知，是真正屹立在修行顶端的至强者。

\begin{quote}
``随我上来。''
\end{quote}

梅天狸轻声一语，身形率先踏破层层云海，沿着序数云台的延伸轨迹，向着更高的九天之上缓步前行。

众人不敢迟疑，齐齐敛去周身杂念，整理身形，紧随其后，踏着辽阔无垠的云端平台，步步向上登临。

\begin{center}\rule{0.5\linewidth}{0.5pt}\end{center}

众人应声启程，脚下的云端古道缓缓延展。此前进行生死决战的序数云台，本就高悬万仞高空、凌驾世俗云海之上，已是常人难以企及的云端绝境，可如今众人前行的道路，却是笔直向上、层层攀升，彻底超脱了决战云台的高度。

这是一片真正的九天高层空域，远在决战云台之上，二者如同凡尘与九霄的悬殊差距。脚下的云晶栈道一路向上绵延，直通天境尽头，越往上走，天地灵气愈发纯粹凝练，不再是寻常气态灵气，而是化作肉眼可见的乳白色灵雾，层层叠叠，萦绕周身。

四周的空间壁垒愈发稳固厚重，空气中流转的不再是普通反射序数道纹，而是一丝丝极其细微、至高无上的Y序列本源气息，静默无声，却自带镇压心神的磅礴威压。

一路走来，风罡绝迹，流云静止。

下方的山川大地、宗门楼宇尽数化作微末小点，彻底被云海隔绝。整片天地只剩下纯白云道、澄澈穹顶，以及无处不在的苍茫道韵，空灵、神圣、肃穆，让人不由自主心生朝圣之心。

蔡须坤、蔡兴坤、卢本韦三位少女步履轻盈却端正，眼底满是敬畏；唐三、小舞、霍雨浩几人亦是神色庄重，不敢有半分懈怠。每向上一步，他们都能清晰感受到自身魂力、序数法则被天地道韵缓缓淬炼、打磨、净化。

这里的每一缕空气，都远超秘境苦修的极致环境，是真正的大道圣地。

不知登临数千丈云道，彻底脱离了决战云台所在的高空层域，下方那片刚刚经历惊天对决的序数试炼台，已然缩成云海之下一枚微不足道的微小亮斑，彻底沦为下层凡尘。而在九天最顶层的云海尽头，一座镇压万古、凌驾诸天的神圣宏伟圣殿，终于完整展露在众人眼前------\textbf{良序圣殿}。

它绝非寻常宗门殿宇所能比拟，整座殿堂以太古神晶为骨、诸天序数理纹为肤，通体剔透纯白，无一丝杂色，亿万道细密繁复的神圣道纹如星河脉络般交织流淌，自殿基、殿柱直至万丈穹顶生生不息，每一道纹路都是最本源的高阶序列法则，普照八方、涤荡乾坤。圣殿悬空而立，不倚山岳、不踏云海，独立于九天穹顶之上，凌驾所有云端空域，是真正意义上的「天上圣境，道统之巅」。

无尽柔和却至尊至极的圣光自殿体深处溢流而出，铺展成千万道秩序光带，环裹整座圣殿，层层叠叠、轮转不息。圣殿穹顶直接接驳九天星河，漫天星辰垂落点点星辉，与殿内道纹交相辉映，自成一方圆满大道天地。这里是整片扽魂殿的道统源头，是所有序数法则的诞生之地，更是大长老test\_alpha0静坐悟道三百年的至高修炼圣地，层级与威严，冠绝整个宗门。

整座圣殿自带天然无上威仪，无需阵法禁锢、无需护法镇守，万法不侵、万邪退避、万道俯首。但凡低阶法则、浅薄魂力、躁动心神，一旦踏入圣殿千里领域，便会被瞬间净化抚平。没有凌厉的杀伐威压，却有着岁月沉淀、大道独尊的苍茫厚重，是高等维度对低维生灵的天然俯瞰，让所有生灵发自内心生出匍匐朝圣的敬畏，不敢生出半分亵渎之心。天地间的风、云、光、序，尽数在此殿之前沉寂俯首，整片九天空域，唯剩圣殿独尊。

通体由纯白神晶构筑，万千神圣纹路交织成片，无尽柔光自殿体深处流淌而出，普照八方。圣殿穹顶链接九天星河，殿身悬浮于万层云海之上，无基无台，自成天地，周身环绕层层神圣秩序光带，正是扽魂殿最至高无上的核心圣地------良序圣殿。

这里，便是大长老test\_alpha0的专属修炼之地，也是整个扽魂殿序数大道的源头所在。

圣殿四周没有任何防御阵法，没有任何值守强者，却自带万法不侵、诸邪退散的至高气场。寻常修士只需靠近百丈范围，周身紊乱法则、浅薄魂力便会被瞬间净化，心境浮躁尽数抚平，唯有极致的神圣与威严笼罩身心。

梅天狸驻足圣殿千丈之外，回身看向身后列队肃立的果糕小队，语气郑重几分：

\begin{quote}
``良序圣殿承托整片扽魂殿的大道道基，汇聚万千序数本源，是我殿道统核心。大长老身居此处悟道万载，执掌顶级Y序列法则，你们待会入殿觐见，谨守本心，恭肃行礼即可。''
\end{quote}

\begin{quote}
``是！''
\end{quote}

众人齐声应答，声音整齐沉稳，无半分慌乱。

梅天狸微微颔首，随即转身，衣袂轻扬，率先朝着那座悬浮九天的神圣圣殿缓步走去。

云海铺道，圣光引路。

果糕小队众人紧随其后，一步步踏向这座承载至高大道的良序圣殿，准备觐见这位传说中的扽魂殿至高大长老。

越是靠近圣殿，神圣肃穆的道韵便愈发厚重，周身流动的Y序列气息愈发精纯逼人，哪怕是众人早已稳固的心境，也不由得微微紧绷，每一步落脚都无比端正恭敬。

踏入圣殿殿门的刹那，万丈柔光瞬间笼罩周身，殿内并无想象中的奢华雕饰，唯有一片澄澈无垠的纯白天地，穹顶流转着无尽序数神纹，天地道韵自成循环，静谧得仿佛隔绝了世间一切喧嚣。

圣殿正中央，一道身影静静盘坐于虚空莲台之上。

那是一位白发垂肩的老者。满头银发如雪如霜，丝丝缕缕垂落肩头，没有半分杂乱，面容苍老却不颓败，眉眼间镌刻着岁月沉淀的沧桑，沟壑浅浅的纹路里藏着阅尽万古的通透。明明只是静静端坐，却自带镇压九天的无上威严，无需释放半点威压，便让整片圣殿的所有道纹尽数俯首沉寂。

他便是扽魂殿大长老------\textbf{test\_alpha0}。

世人皆知，大长老已然活过三百余载岁月。三百年光阴，足以看遍世间兴衰、大道更迭，见证无数修士起起落落、无数法则迭代新生。大陆修行界的风起云涌、序数大道的起源演变、各路天骄的崛起陨落，尽数收纳于他的眼底心中。

整个扽魂殿，上至诸位掌权长老，下至核心执事、亲传弟子，半数以上皆是出自他的道统传承，遍布整座宗门，徒子徒孙无数，是扽魂殿名副其实的道统源头、定海神针。

而最令人敬畏的，是他深不可测的绝顶修为。

梅天狸已是踏足 \allowbreak{}\NovelInline{I03776}{Y(1,3,9)}\allowbreak{} 层级的顶尖强者，足以碾压一切反射序数修士，可在这位白发老者面前，依旧只是晚辈。test\_alpha0深耕高阶Y序列大道三百年，早已超脱常规序列层级的桎梏，道行之深、法则掌控之圆满，远远超出所有人的认知边界，是整片大陆为数不多真正触摸到序列大道本源的至强者。

在他面前，所谓的多逗号迭代、I/M数阵、反射序数极致壁垒，尽数是孩童般的粗浅法理；即便是Y序列高阶力量，也只是他万千道统中的一隅分支。

梅天狸驻足莲台之下，微微躬身行礼，姿态恭敬至极：

\begin{quote}
``师尊，弟子天狸，带领本次试炼合格的弟子，前来觐见。''
\end{quote}

果糕小队众人见状，齐齐躬身俯首，身姿端正，神色肃穆，无人敢有半分轻慢。

虚空莲台上，白发老者缓缓抬眸。

那双历经三百年岁月的眼眸，澄澈、深邃、苍茫，仿佛藏着整片天地的大道流转，淡淡扫来的目光，没有凌厉杀意，却让众人心神瞬间澄澈，所有浮躁尽数消融，只剩满心敬畏。

目光缓缓扫过果糕小队每一人，看着这群少男少女眼底澄澈的韧劲、未被磨平的锋芒，以及方才跨越序数极限、绝境联动的绝佳底蕴，test\_alpha0苍老的眉眼间，难得掠过一丝极淡的赞许与感慨。

他声线低沉厚重，带着跨越三百年岁月的沙哑沧桑，缓缓开口，声音不高，却响彻整座静谧神圣的良序圣殿：

\begin{quote}
``你们这一代人，潜力远超老夫预料。''

``区区数月，便从粗浅的反射序数体系中挣脱桎梏，硬生生摸到Y序列的大道门槛，甚至能凭借凡人极致的低维之力，逼得天狸身拥 \allowbreak{}\NovelInline{I03777}{Y(1,3,9)}\allowbreak{} 道体的强者退后半步。这般悟性、绝境破限的韧性、浑然一体的团队默契，纵览老夫三百年修行所见，亦是百年难遇的顶尖天骄。''
\end{quote}

老者微微颔首，目光悠远，透过圣殿星河穹顶，仿佛望向了百年前那片战火纷飞的天地，语气渐渐沉凝，裹挟着无尽遗憾：

\begin{quote}
``你们天赋卓绝、未来不可限量，或许，你们能完成老夫百年前，未能完成的夙愿。''
\end{quote}

众人闻言心头一凛，纷纷屏息凝神，静静聆听大长老的过往。

\begin{quote}
``距今整整百年之前，世间有一尊祸乱诸天的绝世魔头------帅气教教主，胡赞。''
\end{quote}

test\_alpha0字字沉重，每一句都裹挟着跨越百年的刻骨怅然：

\begin{quote}
``百年之前，世间祸乱丛生，一尊绝世魔头横行诸天，正是帅气教教主------胡赞。彼时老夫已然勘破光明Y序列本源，执掌世间最纯粹的光明正道之力。光明克噬暗、正道镇邪祟，天克地制，我对胡赞的黑暗邪道，是与生俱来、层级碾压的绝对克制。''

``那一战，鏖战七日七夜，昏天黑地、日月倾覆，大陆空域崩碎千里，法则秩序尽数紊乱，正邪碰撞的余波震裂山河、殃及亿兆生灵。战局从头到尾，皆牢牢攥在老夫手中，我以光明序列层层碾压、步步锁死，将胡赞的黑暗邪术尽数拆解、压制、击溃。他节节败退、修为溃散、底牌尽空，彼时的他，已然油尽灯枯、无路可逃，老夫只需最后一指，便可斩除这世间百年祸根。''
\end{quote}

说到此处，老者澄澈的眼底骤然翻涌百年未散的戾气与滔天不甘，苍老的嗓音紧绷沙哑，满是无尽悔恨：

\begin{quote}
``可我万万没有料到，此獠歹毒冷血、丧尽天良！绝境之中，他不惜触犯天地禁忌，催动禁灭黑暗邪术，硬生生蛮横抽取千万苍生与生俱来的本源帅气气韵，尽数吞噬炼化为己身邪力！''

``一夜之间，万民失色、生灵涂炭，整片大地生机枯竭、民不聊生。而胡赞凭借数万百姓的帅气之力，黑暗邪力暴增十倍不止，硬生生打破光明序列的克制枷锁、突破自身修为上限！原本必死之局，被他以苍生性命强行逆转！老夫手握绝对正道优势，占尽天时地利道统，却眼睁睁看着对手浴邪翻盘，最终力竭道崩、含恨败退。那是我三百年修行生涯，最屈辱、最不甘、最刻骨铭心的惨败。''

``百年悠悠而过，胡赞盘踞世间，以帅气教祸乱四方、荼毒苍生，黑暗邪道蔓延天下，无数百姓饱受掠夺之苦。老夫苦修三百年，阅尽大道兴衰，却始终弥补不了当年的遗憾，斩不掉这尊世间魔头。这份无力与缺憾，萦绕百年，终生难消。''
\end{quote}

test\_alpha0缓缓敛去眼底沧桑戾气，抬眸望向眼前朝气炽盛、傲骨铮铮的果糕小队，苍茫的目光中，从暮年的颓憾死寂，骤然燃起一抹炽热、极致滚烫的期许，新旧光影剧烈碰撞：

\begin{quote}
``老夫如今岁至暮年，大道定型、修为封顶，寿数将尽，再无突破之力，毕生遗憾，至死难偿。我做不到的，你们可以！''

``你们是初生朝阳、无尽可期，拥有老夫早已逝去的锐气、潜力与无限可能。你们能绝境破限、逆伐层级，能以凡人之躯撼动高阶大道，这份天赋与韧性，远超当年的我。老夫百年前没能守住的苍生太平、没能斩断的黑暗祸根，今日尽数托付于你们。''

``潜心苦修，挣脱低阶序数桎梏，登临Y序列至高巅峰！终有一日，愿你们踏碎黑暗、击溃胡赞、覆灭邪道，以少年锋芒，了结老夫百年憾事，救万民于水火，还天下朗朗乾坤、永世太平！''
\end{quote}

庄严神圣的圣殿之内，老者的期许字字铿锵、振聋发聩，化作沉甸甸的大道传承与使命，深深烙印在每一位少年少女心底。

\begin{center}\rule{0.5\linewidth}{0.5pt}\end{center}

test\_alpha0静静俯瞰着众人，眼底的沧桑遗憾尽数褪去，只剩一片肃穆郑重，声音沉如天钟，回荡在整座纯白圣殿之间：

\begin{quote}
``今日，老夫给你们一场机缘，亦是一份世间重托。''

``你们天资卓绝、心性坚韧，身负绝境破限的逆天潜力，远超历代同辈天骄。老夫在此立下规制------如若你们能在十年之内，全员尽数突破，人人踏足 \allowbreak{}\NovelInline{I03778}{Y(1,3)}\allowbreak{} 序列境界，挣脱低阶反射序数的所有桎梏，真正踏入高阶大道之列。''

``那老夫便破格册封，让你们全员位列扽魂殿长老之位，执掌殿内道统权柄，承继高阶序列修行底蕴，与我、与天狸一同镇守扽魂大道。''
\end{quote}

他微微前倾身形，目光锐利而恳切，带着三百年大道传承的厚重，沉声发问：

\begin{quote}
``如今，老夫问你们------你们可愿意潜心苦修、砥砺前行，十年破境、登临序列，自此守卫世间正道，扛起这份庇护苍生、镇灭邪祟的沉重重任吗？''
\end{quote}

话音落定，整座良序圣殿一片寂然，唯有星河道纹缓缓流转，无声等候一众少年的答复。

没有半分迟疑，没有半分犹豫。

蔡须坤率先上前一步，身姿挺拔如松，眉宇间褪去所有青涩，只剩凛然正气，音色清亮铿锵，响彻整座纯白圣殿：

\begin{quote}
``我等愿意！承蒙大长老栽培与厚望，我辈定当日夜苦修、砥砺本心，挣脱序数桎梏，踏碎大道壁垒！''
\end{quote}

紧随其后，唐三沉步踏出，眼神坚定如铁，字字掷地有声：

\begin{quote}
``十年为限，全员破境，登临 \allowbreak{}\NovelInline{I03779}{Y(1,3)}\allowbreak{} 序列！此生恪守正道，不畏强权、不惧黑暗，以己身大道庇护苍生！''
\end{quote}

卢本韦、奥斯卡并肩而立，神色肃穆，高声应和：

\begin{quote}
``愿承大道重任！苦修不辍，破限精进，练成无上序列之力，终结百年黑暗祸乱！''
\end{quote}

小舞、蔡兴坤、霍雨浩、诺当齐齐躬身立誓，少年声响层层叠加，整齐洪亮，震彻殿宇，撼动漫天序神道纹：

\begin{quote}
\textbf{``我等愿担世间重任！''}

\textbf{``十年苦修，全员踏足 \allowbreak{}\NovelInline{I03780}{Y(1,3)}\allowbreak{}！''}

\textbf{``镇守正道，诛灭邪祟！''}

\textbf{``击溃胡赞，还天下太平！''}
\end{quote}

全员齐声立誓，初心灼灼，傲骨铮铮。少年热血与大道敬畏交融在一起，化作纯粹且坚定的道心，回荡在九天圣境之中。原本静静流转的圣殿神纹，仿佛受到少年誓言的感召，骤然熠熠生辉，万千圣光垂落，温柔包裹住果糕小队众人，为他们烙印下不可磨灭的正道道誓。

他们清楚知晓，这份承诺绝非虚名浮华，而是沉甸甸的使命------是大长老百年未圆的遗憾，是天下苍生亟待救赎的期盼，更是他们从此踏上高阶大道、以身为刃镇护世间的毕生道途。

虚空莲台之上，test\_alpha0静坐百年、古井无波的身形，第一次微微松动。

那双阅尽三百年沧桑、常年覆着苍茫死寂的眼眸深处，原本沉寂无光的瞳孔，骤然亮起一星极柔、极滚烫的光点。那是积压百年的不甘、屈辱、遗憾尽数消融后的释然，是苍生夙愿终有传承的动容。常年紧绷、肃穆冷硬的眉眼纹路缓缓舒展，岁月刻下的沟壑仿佛都温柔淡化，一抹极浅、极珍贵的浅笑，缓缓漾在苍老的唇角，不张扬，却足以撼动整座圣殿的道心。这是他战败百年、独守黑暗残局百年以来，第一次发自肺腑的宽慰与欣喜。

他雪白垂落的长眉微微轻颤，极细微的动作，却是至强者心绪波动的极致体现。原本静静叠于膝上、执掌三百年大道、稳如磐石的枯瘦手掌，缓缓抬起，动作缓慢、庄重、沉稳，带着抚平百年乱世沧桑的厚重，轻轻向外虚按。指尖流转的纯净光明序列柔光，温柔舒展，不再有半分镇压强横，只剩长者的庇佑与期许。

伴随着他的心绪动容，整座良序圣殿共鸣震颤。漫天星河道纹轮转生辉，万千圣光簌簌垂落，温柔笼罩下方立誓的少年们，天地道韵自发俯首，印证这一场正道传承。

老者胸腔微震，苍老沙哑的嗓音褪去了往日的沉冷孤寂，染上明显的温热震颤，连道三声，字字恳切：

\begin{quote}
``好\ldots\ldots 好\ldots\ldots 好！''
\end{quote}

三声赞叹，一声比一声厚重，一声比一声滚烫，藏着百年孤守的落幕，也藏着新生大道的开篇。眼底盘踞百年的阴霾与戾气，在此刻被少年们滚烫的赤诚道心彻底驱散、一空如洗。

他轻轻吐息，声如暮钟悠远，满是释怀与笃定：

\begin{quote}
``百年憾事，终有传人。苍生太平，终有望矣。''
\end{quote}

一丝极淡、极温柔的笑意缓缓攀上他布满岁月沟壑的眉眼，那是尘封百年、历经战败屈辱、看尽苍生流离后，第一次发自心底的动容与宽慰。三百年风霜铸就的威严冷肃尽数消融，取而代之的是长者的慈爱、道祖的期许，以及夙愿有望的滚烫暖意。

\begin{quote}
``百年憾事，终有传人。苍生太平，终有望矣。''
\end{quote}

\begin{center}\rule{0.5\linewidth}{0.5pt}\end{center}

殿内圣光缓缓收敛，漫天序神道纹归于平和，良序圣殿肃穆的氛围褪去几分沉郁，只剩绵长悠远的大道气韵萦绕其间。

梅天狸见师尊心绪舒展，当即缓步上前，身姿恭谨行礼，出声请示道：

\begin{quote}
``师尊，此番试炼，果糕小队全员圆满通关，已然彻底吃透反射序数体系，将低维法则打磨至极致圆满，再无精进余地。''
\end{quote}

她目光落向身后台姿端正、道心稳固的众人，继续恭敬问道：

\begin{quote}
``弟子斗胆请教，他们修尽反射序数壁垒，前路已无低阶精进空间，下一步该去往何处修行，方能稳步突破、登临更高阶的Y序列大道？''
\end{quote}

虚空莲台之上，test\_alpha0闻言微微抬眸，苍老的眼眸透着洞悉万物的深邃，略一沉吟，缓缓开口指点，声音厚重悠远，响彻整座圣殿：

\begin{quote}
``尔等既已圆满炼化所有反射序数法理，低维天地再无滋养尔等道体的机缘，固守此处，只会停滞不前。''

``老夫为你们指一条无上机缘之路。''
\end{quote}

老者抬手指向诸天东南方向，眸光穿透层层云海与万里山河，精准落点大陆极境：

\begin{quote}
``自此地出发，一路向东南直行，直至抵达整片修行大陆的疆域边缘，再向东横渡沧海八百里海域，便可抵达海外秘境------\textbf{蓬莱岛}。''

``蓬莱岛内，屹立着一座先天道胎凝成的无上神山------\textbf{坤仑山}。它并非世俗凡山，亦非普通秘境灵峰，而是整片诸天维度极少数诞生于混沌初开的序列道山，自成一方独立高阶大道领域，彻底隔绝大陆低维法则，是专为高阶序列悟道而生的绝世圣地。

大陆之内，无论是世俗魂道、反射序数壁垒，还是寻常Y序列法理，皆属于「有限低维大道」，法则固化、上限锁死；可坤仑山不同，此地萦绕的是\textbf{无限维度的终极序数理则}，整座山体每一寸岩土、每一缕云雾、每一缕清风，都有着无上序列本源。

山上无四季更迭、无昼夜轮转，常年悬垂诸天万道霞光，星河垂落为瀑，道纹凝地为脉。山间云海并非凡俗水汽，而是凝练到极致的序列道韵原液，翻涌流转间层层叠叠、包罗万象，低阶修士一旦踏入，周身所有粗浅法则会瞬间被净化、重塑、升华，道基瑕疵尽数涤荡。

山体贯穿数层虚空秘境，上接域外星河，下通混沌道根，山中灵脉永续不竭、大道生生不息。随处可见天然凝成的道台、石座、悟道莲台，皆是千万年道韵沉淀而成，静坐其上，便可被动吸纳终极序列法理，体悟维度跃迁的至高真谛。

整座神山寂静、孤高、超然世外，无凡尘喧嚣，无纷争杂念，唯有万古流转的大道轰鸣回荡。它比云端之上的良序圣殿更为超脱，圣殿是扽魂殿的道统核心，而坤仑山，是诸天终极序列的悟道本源秘境，层级完全凌驾大陆所有修行之地。''
\end{quote}

test\_alpha0语气愈发肃穆郑重，缓缓道出这座秘境的镇守底蕴与至强者修为：

\begin{quote}
``如此无上圣地，自然有顶级大能镇守。坤仑山由大供奉------\textbf{最菜萌新}亲自坐镇。他与老夫是天道认证的正统神之使者，执掌超脱常规修行体系的终极大道，早已登临 \allowbreak{}\NovelInline{I03781}{\omega\text{-}Y(1,\omega)}\allowbreak{} 的无上终极序列境界。

此等修为，早已跳出有限序数、常规Y序列的桎梏，触摸到维度无限叠加的大道尽头。所谓反射序数极限、\allowbreak{}\NovelInline{I03782}{Y(1,3,9)}\allowbreak{} 层级，在 \allowbreak{}\NovelInline{I03783}{\omega\text{-}Y}\allowbreak{} 终极序列面前，不过是孩童初学的粗浅法理。大供奉道行浩瀚难测，洞悉序列终极真谛，手握无尽大道权柄，是整片诸天最顶尖的悟道者之一。''

``正因有他坐镇，坤仑山的终极道韵得以永续不散，秘境机缘纯正且充沛。你们果糕小队已然修尽低维反射序数，卡在高阶Y序列的入门瓶颈，前往坤仑山悟道，便是从「有限法理」踏入「无限大道」的绝佳契机。你们可近距离观摩终极序列运转、体悟维度层级差距、打磨自身道基，破除修行桎梏，极速积淀底蕴，为十年全员稳稳突破 \allowbreak{}\NovelInline{I03784}{Y(1,3)}\allowbreak{}、执掌正道重任，铸就最稳固、最纯粹的大道根基。''
\end{quote}

一番指点，彻底为果糕小队拨开了前路迷雾。众人心中豁然开朗，同时愈发敬畏天地大道的浩瀚无垠，眼底齐齐亮起炙热而纯粹的向往之色。

此前他们困守大陆低维法则，纵使倾尽极致、突破序数壁垒，也始终被桎梏在有限大道之中，处处受限、难窥真章。如今听闻坤仑山超然世外的终极道韵、\allowbreak{}\NovelInline{I03785}{\omega}\allowbreak{} 级序列的无上真谛，还有顶尖神之使者坐镇传道，众人心中积压已久的修行渴望彻底被点燃。

原来真正的大道，从不是固化的迭代推演，而是无限维度的层层跃迁。

无人言语，众人却各自在心底暗自立誓。方才许下的十年破境、镇守正道、抚平百年祸乱的誓言，不再是空洞的承诺，而是有了清晰可寻的前路与实实在在的机缘。他们收敛心中躁动，褪去所有侥幸与浮躁，道心被打磨得愈发澄澈坚固。

卢本韦诚恳地询问大长老：

\begin{quote}
``我们现在就去吗？''
\end{quote}

test\_alpha0缓缓摇头，语气温和而沉稳：

\begin{quote}
``等你们修为更加扎实之后，自然可以前往。万事有序，不可急躁。''

``等你们到达 \allowbreak{}\NovelInline{I03786}{0\;111\;222}\allowbreak{} 的境界之后，小狸（梅天狸）会亲自带领你们前往。''
\end{quote}

一颗颗滚烫坚定的道心，在神圣圣殿的道韵滋养下悄然蜕变、扎根生长，为接下来的跨海修行、大道破境，埋下了最坚实的伏笔。

\chapter{第六十三章：帝国密谋，暗夜刺杀}\label{ux7b2cux516dux5341ux4e09ux7ae0ux5e1dux56fdux5bc6ux8c0bux6697ux591cux523aux6740}

可九天之下，大陆腹地的黑暗暗流，早已汹涌滔天，吞噬万里河山。

大陆极西，帅气教百年总坛，漆黑邪殿终年黑云覆顶、邪气锁空，无半分天光正道。自星斗大森林一战落幕，这座沉寂已久的黑暗王座之上，始终萦绕着暴戾到极致的森冷怒意。

帅气教教主胡赞端坐王座，黑袍翻涌，周身流转的黑暗气韵剧烈躁动，无数被蛊惑万民的帅气本源翻腾不息，却压不住他心底滔天的愠怒与忌惮。殿内阴风呼啸，无数邪道执事、麾下长老匍匐在地，噤若寒蝉，无人敢发一言，整片邪殿被死寂又狂暴的威压彻底笼罩。

星斗一役，是他布局已久、志在必得的绝杀之局。

为彻底掐灭果糕小队这颗正道新星，根除未来最大隐患，他先后布设两层杀局，步步紧逼、层层绝杀，只为将这群潜力逆天的少年天骄扼杀于未成之时。

\begin{center}\rule{0.5\linewidth}{0.5pt}\end{center}

\textbf{第一次刺杀}，恰逢大陆青年魂师大赛盛况空前，四方天骄齐聚赛场。胡赞派出自己两位嫡子------沉稳阴狠、精于暗杀阵法的长子胡善，暴戾悍猛、专修黑暗肉身的次子胡勇。二人伪装散修混入赛场，蛰伏隐忍，伺机待发，趁果糕小队登台对战、心神专注擂台之际，骤然爆发全部邪力，撕裂赛场防护阵纹，欲借赛事混乱暗下杀手、一击斩草除根。

奈何正道有护，天骄有佑。就在杀招临身的刹那，梅天狸瞬身降临，高阶Y序列道韵凌空铺开，仅凭一己之力便镇压兄弟二人全部邪功，瓦解暗杀阵法，击溃联手杀招。胡善、胡勇惊惧交加，深知战力悬殊，只能狼狈遁走，第一次精心谋划的刺杀，无功而返。

\begin{center}\rule{0.5\linewidth}{0.5pt}\end{center}

一次失败，不足以平息胡赞的杀心，反而让他愈发忌惮果糕小队的逆天潜力。他深知这群少年绝境破限、愈战愈强，放任成长，日后必是颠覆自己黑暗霸业的最大阻碍。

为此，他不惜动用教中底蕴，派出常年镇守总坛、身经百战的顶级精锐猎杀小队，由资深副队长统领，全员精通黑暗隐匿、序列绞杀、禁术自爆，是专为猎杀正道天骄打造的死士精锐，开启\textbf{第二次跨境追杀}。

这支精锐小队隐匿踪迹、跨域追踪，最终在星斗大森林外围布下九重黑暗绝杀大阵，以完整 \allowbreak{}\NovelInline{I03787}{0\;111\;222}\allowbreak{} 三行BMS混沌序列锁死天地，布下无解死局，誓要全歼果糕小队。

战局初期，果糕小队全员苦战、节节败退，重伤力竭、阵法濒临崩碎，已然走到全员覆灭的绝境。可谁也未曾料到，绝境之中，小队队长牢大毅然舍身断后，以一己之躯拦下漫天杀机，更在死局之中破壁突破，踏出全新 \allowbreak{}\NovelInline{I03788}{\psi(K_\omega)}\allowbreak{} 终极序数，逆势击溃四名高阶精锐。

最终，牢大燃尽生命、崩碎道基，以奥义·坠机殉道赴死，不仅硬生生拖住绝杀战局、为队友搏出一线生机，更重创敌方核心副队长，硬生生将胡赞的第二次绝杀布局彻底撕碎、击碎。

虽最终灭杀牢大，可帅气教损兵折将、精锐重创，两次倾尽心力的刺杀追杀，尽数未能覆灭果糕小队，甚至被对方以悲壮殉道的方式，破局脱身、留存正道火种。

\begin{center}\rule{0.5\linewidth}{0.5pt}\end{center}

\begin{quote}
``两次绝杀，两次落空\ldots\ldots{}''
\end{quote}

胡赞缓缓抬手，指尖漆黑邪力滋滋暴涨，眼底戾气翻涌、寒芒刺骨，声音沙哑冰冷，裹挟着百年难遇的震怒：

\begin{quote}
``一群初出茅庐的少年，折我子嗣、损我精锐、破我杀局！区区正道雏苗，竟三番五次坏我大业！''
\end{quote}

他真正忌惮的，从来不是当下战力尚且稚嫩的果糕小队，而是他们近乎无解的绝境潜力、悍不畏死的正道道心，以及那一次次置之死地而后生的逆天韧性。百年前他憾败于test\_alpha0的光明正道，侥幸以万民气韵翻盘续命，深知正道天骄一旦羽翼丰满、登临高阶序列，便是他黑暗霸业的末日。

牢大的绝境破壁、以身殉道，更是让他心生寒意。未成天骄尚且有如此逆天战力与决绝执念，若是放任剩余队员稳步修行、十年破境、铸就阴阳双核无瑕道体，未来整片大陆的黑暗格局，必将彻底倾覆。

数次深呼吸后，胡赞强行压下心底翻涌的杀意与焦躁，暴戾的神色渐渐褪去，取而代之的是深沉阴翳与滔天野心。

强攻刺杀，已然行不通。扽魂殿长老守护森严、天骄气运深厚，再贸然出手，只会惊动良序圣殿的至强者，招致灭顶之灾。

既然杀不尽、灭不绝，那便换一种玩法------\textbf{以疆域养邪气，以万民筑帝业}。

\begin{center}\rule{0.5\linewidth}{0.5pt}\end{center}

百年蛰伏，胡赞从未停下积蓄力量的脚步。当年凭借数万百姓帅气气韵翻盘之后，他便洞悉终极大道：个人战力终有上限，掌控天下万民、执掌万里疆域、掠夺世间气韵，方能铸就不败黑暗基业。

这些年来，他刻意避开扽魂殿视线，暗中布局、步步蚕食，以帅气教邪术渗透四方行省，腐化官员、收服宗门、颠覆地方秩序，悄无声息掌控了大陆西境整整六个行省。

整片西境沃土，早已脱离俗世正道管控，律法失效、正邪颠倒，彻底沦为帅气教的私人疆域，成为滋养黑暗邪道的温床。

不止于此，为快速扩张势力、积累气韵底蕴，胡赞下令麾下邪教大军全员出动，以雷霆之势疯狂吞并周边城镇、边陲村落，所过之处，尽数被黑暗序数理纹笼罩，天地气韵被强行篡改。

教中修士游走街巷、深入民间，编织虚假盛世谎言，扭曲正邪大道，大肆蛊惑淳朴万民。他们宣扬正道腐朽、天道不公，唯有帅气教能庇佑众生、赐人荣光，日复一日侵染百姓心神。

寻常民众不识高阶大道、难辨正邪真伪，久而久之，尽数被黑暗邪力洗脑同化，从抵触抗拒变为盲从信奉，自愿供奉气韵、追随邪道，沦为胡赞修炼黑暗功法、构筑霸业的源源养分。

万民归心，疆域在手，邪气日盛，帝业初成。

胡赞抬眸俯瞰下方整片黑暗疆域，万千民众的信仰气韵隔空汇聚、涌入其身，滋养着他愈发恐怖的黑暗修为。他眼底燃起滔天帝王野心，冰冷霸道的声音响彻整座邪殿：

\begin{quote}
``数次刺杀无果，无妨。区区残剩少年，无足为惧，待本座大业功成，弹指可灭！''

``百年前，test\_alpha0以光明正道压我一头，害我毕生留憾。今日，我便以万民为基、以疆域为土，建帅气帝国，立黑暗正统！''

``颠覆正道秩序，横扫扽魂殿道统，一统整片诸天大陆！从今往后，黑暗为尊，帅气为道，世间再无光明正道！''
\end{quote}

霸道话音落下，整座邪殿邪气暴涨、黑云翻滚，无数黑暗序数理纹疯狂跳动，整片西境的混沌气息遥相呼应、震天动地。

\begin{center}\rule{0.5\linewidth}{0.5pt}\end{center}

一边是九天之上，良序圣殿传道授业、正道火种传承不息，少年天骄立誓十年破境、守护苍生；一边是大地之下，黑暗帝国悄然成型、邪气蔓延万里，祸乱世间、荼毒万民。

正邪棋局彻底对立，乱世格局已然成型。

果糕小队背负队友殉道的血海深仇、大长老的百年遗憾、天下万民的太平期许，即将远赴海外蓬莱坤仑山，踏上行棋悟道、破境变强的逆天征途。

而胡赞蛰伏蓄力、构筑帝朝、蓄谋颠覆天下的黑暗阴谋，也已然拉开序幕。

一场席卷整片大陆、关乎正邪存亡、苍生命运的旷世大战，正在黑暗与光明的对峙之中，悄然酝酿，步步逼近。

\begin{center}\rule{0.5\linewidth}{0.5pt}\end{center}

邪殿王座之上，胡赞眸光阴鸷，野心不止于割据西境、积蓄气韵。他深知想要真正建立帅气帝国、名正言顺一统大陆，必先撼动俗世正统皇权，掌控天下权柄。

当今天斗帝国，雄踞大陆中土疆域，幅员辽阔、万民富庶，是整片大陆最正统的俗世皇朝，民心所向、气运鼎盛，也是阻碍黑暗帝国崛起的最大世俗壁垒。只要掌控天斗皇权，便可借皇朝正统收拢天下民心、吸纳天地气韵，彻底抹平正道的世俗根基。

雪夜大帝执掌天斗帝国多年，勤政爱民、民心稳固，皇朝气运绵长，寻常手段难以撼动。可其常年身染旧疾、体弱气虚，朝堂之中暗流涌动、派系林立，防务松弛、内患丛生，正是颠覆夺权的最佳时机。

深思熟虑之后，胡赞当即定下惊天毒计------\textbf{朝堂政变，弑帝夺权}！

他当即传令，召两名嫡系嫡孙入殿领命。

两道黑衣身影躬身踏入邪殿，身姿挺拔、气息阴诡，周身萦绕着精纯内敛的黑暗序列气息，正是胡赞亲手培育、悉心教导的两名天才后人------长孙\textbf{胡继}、次孙\textbf{胡峰}。

二人自幼浸染帅气教黑暗道统，苦修混沌序列邪术，心性隐忍狠绝、手段阴毒缜密，远超胡善、胡勇两兄弟，是胡赞暗中培养、用于颠覆俗世皇朝的隐秘利刃，常年隐匿暗处、不涉明面杀局，只为等待夺权篡位的终极时机。

\begin{quote}
``孙儿听命。''
\end{quote}

二人单膝跪地，声线冰冷沉稳，静待号令。

胡赞俯身垂眸，漆黑眼底翻涌着帝王野心与森冷杀机，声音低沉而霸道，字字皆是诛心毒策：

\begin{quote}
``命你二人即刻隐匿行踪，乔装潜入天斗皇城。''

``蛰伏皇宫内外，伺机刺杀雪夜大帝！''

``大帝一死，天斗皇朝群龙无首，朝堂派系互相倾轧、四方疆域必然大乱。你们趁机联络我教潜伏在天斗境内的暗子势力，搅动朝局、煽动兵变，借乱世之势发动政变，彻底掌控天斗皇权！''
\end{quote}

胡继眸光一厉，沉声道：

\begin{quote}
``祖父放心，我二人定不负命。暗中蛰伏、一击必杀，乱其朝纲、夺其帝位，助祖父奠定万世帝业！''
\end{quote}

胡峰亦是颔首应声，周身黑暗邪力微微震颤，杀机暗藏：

\begin{quote}
``天斗皇室气运虽盛，却挡不住我教混沌序列之力。此次潜入，必成弑帝政变大业，为帅气帝国铺平天下一统之路！''
\end{quote}

二人领命之后，不再多言，身形一晃，化作两道极淡的黑影，消融于漫天黑云之中，气息彻底隐匿，无声无息奔赴中土天斗皇城。

\begin{center}\rule{0.5\linewidth}{0.5pt}\end{center}

两日之后，繁华鼎盛、门禁森严的天斗皇城之内，多了两名容貌温润、气质儒雅的宫外幕僚。

无人知晓，这两名看似文弱、谈吐得体的陌生幕僚，正是改换皮囊、隐去杀机的胡继与胡峰。

为完美潜入深宫、避开层层筛查，二人动用了帅气教传承百年的禁忌秘法------\textbf{幻面藏序术}。配合独门伪装魂导器，将周身所有混沌序列邪气、黑暗魂力彻底剥离、封印、掩盖。体表气息被强行替换为普通文人的淡薄灵气，骨相、容貌、气质尽数重塑，连周身血脉波动、神魂频率都被伪造成寻常魂师模样。

这套伪装魂导器乃是帅气教专为深宫渗透、朝堂刺杀打造的秘宝，可瞒魂力探查、可避阵法侦测、可躲神识扫略，哪怕是天斗皇室专属的高阶护宫检测阵，也无法窥破分毫破绽。

皇城宫门守卫轮番核验身份、探查魂力、扫描邪祟，层层关卡严密周密，可落在二人身上，皆只探测出干净平和的寻常魂师气息，无半分邪异、无半点杀机。整套伪装天衣无缝，骗过了所有皇城守卫、宫廷执事与巡防长老。

成功潜入皇城腹地，胡继、胡峰并未急于直面刺杀雪夜大帝，而是遵循毒计------\textbf{先斩羽翼、再弑君王}。

天斗太子仁厚储君，深得朝臣拥护，二皇子手握部分皇城兵权、掌控禁军调度，二人皆是雪夜大帝左膀右臂，也是天斗皇朝最坚固的两道屏障。只要除掉两位皇子，皇城兵权溃散、储位悬空、朝堂群龙无首，雪夜大帝便彻底沦为孤家寡人。

二人分工明确、配合无间。胡继心思缜密、擅长隐匿绝杀与毒术，负责潜伏皇宫内苑，伺机投毒；胡峰杀伐干脆、身法迅捷，专司暗处突袭、瞬杀截命。

入夜三更，皇城深宫灯火渐寂，禁军巡防节奏渐缓，正是深宫戒备最松懈之时。

两道儒雅身影借着园林暗影、宫墙死角，无声穿梭于重重宫阙之间，伪装魂导器持续运转，全程遮蔽气息，宛若两道无形幽魂，游走在森严皇宫之中，无一人察觉异常。

太子东宫之内，太子正伏案批阅奏折，心神松弛、毫无防备。胡峰蛰伏殿外廊柱阴影之中，指尖凝出一缕微不可查的黑暗序列细刃，借着夜风穿窗而入，无声无息掠过虚空，精准洞穿太子护体魂力，一击封喉。

全程无声无息，无波动、无巨响、无魂力外泄，堂堂天斗太子，未及反应分毫，便颓然伏案，当场殒命。

几乎同一时刻，胡继潜入二皇子练武偏殿。二皇子常年习武、体魄强健、警惕性远超常人，可在完美伪装与高阶黑暗秘法面前，所有戒备皆形同虚设。胡继掌心凝出混沌暗劲，贴身瞬袭，黑暗序列之力瞬间封禁其周身经脉、凝固其魂力流转，趁着二皇子错愕瞬间，利落绝杀，一击毙命。

前后不过数息时间，天斗皇朝两大支柱，尽数悄无声息陨落于深宫之中。

\begin{center}\rule{0.5\linewidth}{0.5pt}\end{center}

斩杀两位皇子、斩断皇朝羽翼后，二人并未抽身撤离，而是收敛所有杀伐痕迹，抹去殿内邪气残留，将现场伪装成突发暴毙的假象，天衣无缝，毫无破绽。

随后，胡继取出一枚通体透明、毫无异色的无声毒丹。此毒无臭无味、入水即融、无形无迹，不沾染邪气、不触发阵法警报、不被魂力探查识破，是帅气教专门针对帝王体魄炼制的绝世奇毒。

他借着内侍换水上茶的间隙，身形微晃，借旁人视线死角，指尖微弹，将毒丹悄然融入雪夜大帝御用的御用玉盏茶汤之中。

毒丹遇水即化，彻底消融，茶汤澄澈依旧，不见丝毫异样。剧毒静静蛰伏在茶汤之内，只待雪夜大帝饮下，便可慢慢侵蚀心脉、腐朽气运、瓦解根基，最终无声无息崩逝，落得久病不治、寿元耗尽的假象。

做完这一切，胡继与胡峰身形悄然退入深宫暗影，依旧保持儒雅幕僚模样，气息平和、神色淡然，混入宫外人流之中。

整座天斗皇城，无人知晓深宫已遭惊天祸乱，无人知晓储君与皇子已然惨死，更无人知晓当朝大帝已然身中无解剧毒。

致命杀机，深藏九重宫阙之内；皇朝倾覆，已然近在咫尺。

自此，黑暗棋局再落重子。

外有六省疆域割据、万民被蛊惑同化，内有绝世刺客潜伏皇城、图谋弑帝政变。胡赞步步为营、里外夹击，以俗世皇权为突破口，以天下万民为养料，疯狂堆砌黑暗霸业，妄图彻底颠覆千年正道秩序，让黑暗帅气之道，主宰整片诸天大陆。

\begin{center}\rule{0.5\linewidth}{0.5pt}\end{center}

就在众人心神沉淀、默默展望前路之时，莲台之上的test\_alpha0再度开口，嗓音苍茫厚重，道出一段世间极少有人知晓的高阶修行秘辛，专门为他们铺垫未来突破Y序列的核心根基。

\begin{quote}
``老夫再传你们一桩顶级修行真谛，是你们未来突破封号斗罗、踏入 \allowbreak{}\NovelInline{I03789}{Y(1,3)}\allowbreak{} 境界的关键根本，也是高阶序列修士脱胎换骨的核心所在。''

``诸天修行体系之中，但凡踏足极限斗罗层次的强者，体内必然会诞生\textbf{双魂核}，其一为\textbf{递归魂核}，其二为\textbf{非递归魂核}。双核并存，方能承载序数大道的循环推演与维度跃迁，方能撑起高阶Y序列的法则体量。''

``世间绝大多数修士，受凡俗道基、浅薄法理桎梏，突破之时，注定天生失衡。递归魂核为主核，体量雄浑、道韵厚重，承载一身修行根基；而非递归魂核，也就是世人所称的\textbf{同位共振第二魂核}，天生孱弱、体量悬殊，远小于主核，只能为辅佐附庸，永远无法做到真正的道衡圆满。绝大多数修行者的双魂核差距极大，递归主核包揽九成以上的魂力运转与道基承载，非递归魂核形同虚设，根本无法发挥真正的双核之力。''

``也正因这份先天失衡，寻常封号斗罗魂力续航薄弱、法则运转滞涩，境界壁垒重重，战力上限被死死锁死，同阶差距难以逾越，终生难有破格战力。世人皆知双魂核是极限斗罗的标配，却无人知晓其中的序列真谛。''

``\textbf{递归魂核}，为秩序主核，承接自身全部修行底蕴，承载序数轮回、法则沉淀、魂力根基，是修士一身修为的本源载体，稳定、厚重、守正，执掌大道根基的存续。''

``\textbf{非递归魂核}，也就是同位共振第二魂核，为变异辅核，不参与基础魂力储存，专职承接域外道韵、维度共振、破限蜕变，负责拆解敌法、叠加buff、突破境界枷锁、衍生全新序列术法，是修士的\textbf{破道之核、破格之核}。''

``普通修士双核悬殊，主核独大、辅核萎靡，相当于空有稳固根基，却无破限之力，辅核的共振、蜕变、破法能力彻底作废，只能沦为摆设。看似双核齐备，实则依旧是单核修行，续航有尽头、法则有僵化、战力有封顶，永远卡在常规封号斗罗的桎梏之中，终生无法触及高阶序列的破格之力。''

``这种天生的失衡，是万千修士的修行桎梏，看似无碍，实则锁死了上限，终生无法触摸终极序列真谛，哪怕踏入Y序列，也永远停留在下品层级，难以精进、难以破限。''
\end{quote}

test\_alpha0眸光锐利，字字郑重，为众人定下无上修行准则：

\begin{quote}
``但你们不同。你们身负逆天破限之姿，承老夫正道嘱托、怀苍生重任，老夫不许你们落凡俗窠臼。''

``老夫要求你们，在初入封号斗罗、突破 \allowbreak{}\NovelInline{I03790}{Y(1,3)}\allowbreak{} 境界的那一刻，务必打磨\textbf{双核同源、炼至均等圆满}。让非递归魂核与递归魂核体量等同、道韵同源、力量持平，铸就独一无二的\textbf{阴阳互补双魂核道体}。''

``一核轮回推演、镇守根基，一核共振破限、衍生新法，阴阳相生、双衡归一。唯有这般完美均等的阴阳互补双魂核，才能打破天地修行桎梏，形成闭环循环的灵气道域。两核体量对等、道韵互通、力量互济，主核守正、辅核出奇，秩序与共振完美交融，形成生生不息的大道闭环。''

``这便是你们破格战力的核心根源：常规修士魂力越用越枯、道力越打越衰，而完美阴阳双核可以内核互生、域外引灵、自我循环修复，源源不断滋生精纯魂力、永续补给大道灵气，彻底杜绝魂力枯竭、道力透支的弊端，持久战、爆发作、守御战皆立于不败之地。''

``更关键的是，失衡双核只能承载低阶序数法理，而均等阴阳双魂核，可以完美承接Y序列高阶维度的跃迁力量，放大所有序列术法的威力，破除同阶战力枷锁。也正因这无瑕道体的破格加持，你们初入封号斗罗、踏入 \allowbreak{}\NovelInline{I03791}{Y(1,3)}\allowbreak{} 境界的瞬间，便拥有寻常超级斗罗的恐怖战力，做到境界初成、战力超阶，跨阶碾压、同阶无敌。
\end{quote}

\begin{quote}
``这是老夫专为你们这批破限天骄留下的无上道基法门，避开了诸天修士千万年的先天修行缺陷，让你们从踏入高阶序列的第一步，就铸就无瑕道体，为日后抗衡胡赞黑暗邪道、登临诸天大道巅峰，铸就万古不败根基！''
\end{quote}

这番秘传心法，远超寻常魂师认知，彻底刷新了众人对封号斗罗与序列修行的认知。众人尽数凝神铭记，将阴阳双核均等圆满的至高要求深深烙印道心，暗自笃定，此番坤仑山修行，必打磨无瑕道基，不负大长老悉心传道。

\begin{center}\rule{0.5\linewidth}{0.5pt}\end{center}

\chapter{第六十四章：真正的n行Bashicu矩阵}\label{ux7b2cux516dux5341ux56dbux7ae0ux771fux6b63ux7684nux884cbashicuux77e9ux9635}

双魂核的运转图还悬在桌边。

卢本韦将两道尚未成形的回路看了一遍，没有急着往里添魂力。先前 test\_alpha0 说起阴阳相济时，她已经在心里试着接过几次；两边各自运转并不难，难的是接到一起以后，怎样让原先稳定的部分仍旧稳定。

她指了指图上交会的地方。

``若后面还要加新的层次，这里也得跟着改？''

``有的要改，有的要保留。''大长老从莲台上下来，走到长桌前，``今天把这件事算清楚。''

他没有立即给那幅双魂核图添线，而是另铺一张纸，写下小队已经熟悉的双行 BMS。

小舞将纸转向众人，指尖依次经过父项、坏根与末列。唐三接着划出好部和坏部，蔡兴坤在旁边写出阶差。几人的动作接得很顺，霍雨浩连行号也没有重复标，只在纸边留下了这一回取项的编号。

test\_alpha0 看过他们的结果，把新纸压在旧纸右边。

``双行的规矩，你们已经用熟了。今天不从头背。我们从它多出一行的地方开始。''

test\_alpha0 在纸边留下几种颜色的标记，给后面的分部图和各项展开结果作对照。

\NovelDisplay{M00687R01}{%
\textcolor[HTML]{B45309}{\text{橙：坏根}}\quad \textcolor[HTML]{B91C1C}{\text{红：原坏部其余部分}}\quad \textcolor[HTML]{7E22CE}{\text{紫：删去的末列}}%
}
\NovelDisplay{M00687R02}{%
\textcolor[HTML]{1D4ED8}{\text{蓝：展开第一项时新增的部分}}\quad \textcolor[HTML]{15803D}{\text{绿：展开第二项时再新增的部分}}%
}
\NovelDisplay{M00687R03}{%
\textcolor[HTML]{6B7280}{\text{灰：因提升条件未通过而保留原值的项}}%
}

好部仍用正文色。坏根本就在坏部里，单用橙色标出它的位置；新增部分的通式也用蓝色。展开第一项时，新添的一段标蓝；展开第二项时保留这段蓝色，再添的一段标绿。灰色只标本次可提升范围内因父链条件未通过而保留原值的项，单因本行阶差为零而不改变的数，不据此涂灰。

蔡须坤将椅子往前挪了些。

``多一行，就得多找一次父项？''

问完这句，蔡须坤将注意力沉入识海片刻。牢大的肉身早已不在，保留下来的残魂仍寄于她的魂灵共生脉络，凝成一道淡淡的人影。

她把刚记下的算式也映了进去。

``一起看。''她在心里轻声说。

那道人影微微颔首。蔡须坤重新握好笔，听见大长老接着道：

``是。先看最短的一张。''

纸面亮起两列光。

\NovelDisplay{M00688R01}{%
\textcolor[HTML]{B45309}{(0)}\textcolor[HTML]{7E22CE}{(1,1,1)}%
}

最左边的零列补齐以后，三行都是零。右边一列则是三个一。没有旁的列挡在中间，第二列每一行的父项都落在第一列。

蔡兴坤先将三条线连好，随后把笔尖停在第三行。

``末列最后一个非零项在这里。坏根是第一列，好部为空，坏部只有最前面这一列零。''

她没有把末列也框进坏部。诺当顺手在两列之间标了边界。

\NovelDisplay{M00689R01}{%
p_1(2)=p_2(2)=p_3(2)=1,%
}
\NovelDisplay{M00689R02}{%
h=3,\quad r=1,\quad\text{好部}=\varnothing,\quad\textcolor[HTML]{B45309}{\text{坏部}=(0)}.%
}

卢本韦沿用双行时的算法，先把末列与坏根作差，再将最后非零行及其下方的阶差截为零。

\NovelDisplay{M00690R01}{%
\textcolor[HTML]{7E22CE}{(1,1,1)}-\textcolor[HTML]{B45309}{(0,0,0)}=(1,1,1),%
}
\NovelDisplay{M00690R02}{%
(\delta_1,\delta_2,\delta_3)=(1,1,0).%
}

``所以新增的第三行，决定了这一回截到哪里。''她说，``第一行、第二行各能增加一，第三行不增加。''

小舞先展开第一项：原坏部那一列零仍保留，右侧新添的一列是 \allowbreak{}\NovelInline{I03792}{(1,1)}\allowbreak{}。展开第二项时，这一列继续保留，再在后面添上 \allowbreak{}\NovelInline{I03793}{(2,2)}\allowbreak{}。

\NovelDisplay{M00691R01}{%
\bigl(\textcolor[HTML]{B45309}{(0)}\textcolor[HTML]{7E22CE}{(1,1,1)}\bigr)[0]=\textcolor[HTML]{B45309}{(0)}%
}
\NovelDisplay{M00691R02}{%
\bigl(\textcolor[HTML]{B45309}{(0)}\textcolor[HTML]{7E22CE}{(1,1,1)}\bigr)[1]=\textcolor[HTML]{B45309}{(0)}\textcolor[HTML]{1D4ED8}{(1,1)}%
}
\NovelDisplay{M00691R03}{%
\bigl(\textcolor[HTML]{B45309}{(0)}\textcolor[HTML]{7E22CE}{(1,1,1)}\bigr)[2]=\textcolor[HTML]{B45309}{(0)}\textcolor[HTML]{1D4ED8}{(1,1)}\textcolor[HTML]{15803D}{(2,2)}%
}

霍雨浩给三行结果分别标上``展开第零项''\,``展开第一项''\,``展开第二项''。这里仍沿用他们先前的约定：\allowbreak{}\NovelInline{I03794}{[0]}\allowbreak{} 删去末列，不新添内容；\allowbreak{}\NovelInline{I03795}{[1]}\allowbreak{} 复制并调整坏部，添上一段；\allowbreak{}\NovelInline{I03796}{[2]}\allowbreak{} 则在删去末列后保留原好部、原坏部，再接上两段依次调整过的坏部。

这些编号都针对同一个原式。展开第二项，不是把展开第一项所得的结果再展开一次；保留下来的原坏部，也不算在新添的部分之中。

蔡须坤托住展开第二项时再新添的那一列，看向仍压在左边的双行例子。

``三行的这一项，展开之后落成了双行的东西？''

``这一例是。''小舞把两张纸靠近，``但多出来的第三行没有白写。它决定了我们从哪里向下取项。''

她在旧纸上补了同样编号的一步。

\NovelDisplay{M00692R01}{%
\bigl(\textcolor[HTML]{B45309}{(0)}\textcolor[HTML]{7E22CE}{(1,1)}\bigr)[2]=\textcolor[HTML]{B45309}{(0)}\textcolor[HTML]{1D4ED8}{(1)}\textcolor[HTML]{15803D}{(2)},%
}
\NovelDisplay{M00692R02}{%
\bigl(\textcolor[HTML]{B45309}{(0)}\textcolor[HTML]{7E22CE}{(1,1,1)}\bigr)[2]=\textcolor[HTML]{B45309}{(0)}\textcolor[HTML]{1D4ED8}{(1,1)}\textcolor[HTML]{15803D}{(2,2)}.%
}

同样展开第二项，双行的最短结构给出一行的数列，三行的最短结构则给出双行的数列。新增的一行改变了展开的层次，具体的一次展开，仍然只留下有限行、有限列。

蔡须坤望着那一列已经删去的 \allowbreak{}\NovelInline{I03797}{(1,1,1)}\allowbreak{}。

``我们得到其中一项，不是把原来的第三行擦掉，再说两张表相等。''

``对。''test\_alpha0 将旧末列的光收回，``箭头上写什么，要和你实际做的事相符。''

梅天狸替霍雨浩压住翻起的纸角。殿外的风从窗边经过，桌上的两张小纸都没有挪动。

大长老让众人将这一步记在页首，只留了三条短注：末列找最后非零行，沿这一行找坏根，阶差到这一行之前为止。

``先记住它在这一例里怎么起作用。''他说，``下一张不增加行数，只把列排得复杂一点。''

第二张纸上仍是三行。

\NovelDisplay{M00693R01}{%
\begin{pmatrix}
0 & 1 & 2 & 3 & 2\\
0 & 1 & 1 & 1 & 2\\
0 & 1 & 1 & 0 & 0
\end{pmatrix}%
}

\NovelDisplay{M00694R01}{%
(0)(1,1,1)(2,1,1)(3,1)(2,2)%
}

蔡须坤先把省略的零补在自己本子上。\allowbreak{}\NovelInline{I03798}{(3,1)}\allowbreak{} 是 \allowbreak{}\NovelInline{I03799}{(3,1,0)}\allowbreak{}，末列 \allowbreak{}\NovelInline{I03800}{(2,2)}\allowbreak{} 是 \allowbreak{}\NovelInline{I03801}{(2,2,0)}\allowbreak{}；列内坐标的前后位置都保留，只省略末尾补齐用的零。

她写完便把笔让给诺当。诺当在上方标出一至五的列号，小舞沿末列第一行的二向左寻找父项。

第四列是三，比二大；第三列也是二，与二相等。两处都越过，父项落在第二列的一。再往前追，第二列的一的父项是第一列的零。

\NovelDisplay{M00695R01}{%
p_1(5)=2,\quad p_1(2)=1,\quad p_1(1)\text{ 不存在},%
}
\NovelDisplay{M00695R02}{%
\text{末列第一行的祖先列：}\quad\textcolor[HTML]{B45309}{2}\longrightarrow1.%
}

小舞将这条父链画好，手指向下一行移去。

``第二行只能沿刚才留下的祖先列找。第三列、第四列第二行虽都是一，这一回都不在候选里。''

``那么最靠右的那个一，也不能因为位置近就选它。''卢本韦说。

她在第四列下方轻轻点了一下，没有落笔。小舞沿第二列、第一列依次看去，在第二列找到严格小于二的一，便将末列第二行的父项也连到那里。

霍雨浩将父项记作 \allowbreak{}\NovelInline{I03802}{p_i(j)}\allowbreak{}，表示第 \allowbreak{}\NovelInline{I03803}{j}\allowbreak{} 列第 \allowbreak{}\NovelInline{I03804}{i}\allowbreak{} 行的父项所在列号。前一下标对应行，括号内对应列，行列都从一编号。

\NovelDisplay{M00696R01}{%
p_1(j)=\max\{q:1\le q<j,\ a_{1,q}<a_{1,j}\},%
}
\NovelDisplay{M00696R02}{%
\text{第一行的候选集合为空时，父项不存在。}%
}
\NovelDisplay{M00696R03}{%
i>1:\quad\text{沿第 }i-1\text{ 行的祖先链，依次检查}%
}
\NovelDisplay{M00696R04}{%
p_{i-1}(j),\ p_{i-1}(p_{i-1}(j)),\ \ldots,%
}
\NovelDisplay{M00696R05}{%
\text{取第一个满足 }a_{i,q}<a_{i,j}\text{ 的列 }q\text{ 为 }p_i(j),%
}
\NovelDisplay{M00696R06}{%
\text{找不到时，父项不存在；该列也没有下一行的候选链。}%
}

``第三行就沿第二行自己的祖先链继续找。''霍雨浩念到这里，翻了一下上一页，``每往下一行，上一行已经选出的父链便是新的限制。''

``别只核对末列。''test\_alpha0 将笔尾移到第三列，``你们用它再走一次。''

第三列第一行的二，父项在第二列；这一行的祖先列于是为第二列、第一列。到了第二行，第三列的一与第二列的一相等，不能取相等项作父项，便继续追到第一列的零。

接下来找第三行父项，用的是第二行已经选出的祖先链。那条链只剩第一列，第二列不能重新放回来。

\NovelDisplay{M00697R01}{%
p_1(3)=2:\quad2\to1,%
}
\NovelDisplay{M00697R02}{%
a_{2,3}=1,\quad a_{2,2}=1,\quad a_{2,1}=0,%
}
\NovelDisplay{M00697R03}{%
p_2(3)=1:\quad1,\qquad p_3(3)=1.%
}

诺当看着自己的连线，将第三行候选旁边原本空着的位置划掉。

``往下算时，可用的列会被上面的关系约束。不能到了新一行，又从整行左边重新搜。''

蔡须坤将三行连线分开托在掌中。第一行留两处，第二行改成一处，第三行便沿那一处继续。她先前觉得``多找一次父项''只是多做一遍，现在才看清，每一遍的起点都由上一遍决定。

``这就是双行往上增加行数时，最容易漏掉的地方。''test\_alpha0 说，``选完以后，把实际选出的链留下。''

蔡兴坤接过笔，开始定坏根。

她先看末列，而没有看整张矩阵最下面出现过多少个非零数。末列第三行是零，最后非零项在第二行；这一项的父项在第二列，因此坏根是第二列 \allowbreak{}\NovelInline{I03805}{(1,1,1)}\allowbreak{}。

\NovelDisplay{M00698R01}{%
p_2(5)=2,\quad\textcolor[HTML]{B45309}{\text{坏根}=(1,1,1)},%
}
\NovelDisplay{M00698R02}{%
\text{好部}=(0),\qquad\textcolor[HTML]{B91C1C}{\text{坏部}=}\textcolor[HTML]{B45309}{(1,1,1)}\textcolor[HTML]{B91C1C}{(2,1,1)(3,1)},%
}
\NovelDisplay{M00698R03}{%
\textcolor[HTML]{7E22CE}{\text{删去的末列}=(2,2)}.%
}

唐三按边界拆开阵盘。好部只有第一列零，坏部是第二至第四列，最右边的 \allowbreak{}\NovelInline{I03806}{(2,2)}\allowbreak{} 单独取下。

他让三列坏部保持原来的相对位置，又把好部按回桌面最左侧。

``这一段保留，复制它，也不会把它从原位搬走。''

蔡兴坤在旁边算阶差。她先写了末列减坏根的各行差值，随后在第二行上画了一道截线。

\NovelDisplay{M00699R01}{%
\text{最后非零行号 }h=2,\quad\text{坏根列号 }r=2,%
}
\NovelDisplay{M00699R02}{%
\delta_1=\textcolor[HTML]{7E22CE}{2}-\textcolor[HTML]{B45309}{1}=1,\quad\delta_2=0,\quad\delta_3=0,%
}
\NovelDisplay{M00699R03}{%
(\delta_1,\delta_2,\delta_3)=(1,0,0).%
}

诺当按刚才查出的分部给原矩阵染色，坏根一列留作橙色，坏部其余列标红，待删的末列标紫。好部在最左边，仍是原来的正文色。

\NovelDisplay{M00700R01}{%
\begin{pmatrix}
0 & \textcolor[HTML]{B45309}{1} & \textcolor[HTML]{B91C1C}{2} & \textcolor[HTML]{B91C1C}{3} & \textcolor[HTML]{7E22CE}{2}\\
0 & \textcolor[HTML]{B45309}{1} & \textcolor[HTML]{B91C1C}{1} & \textcolor[HTML]{B91C1C}{1} & \textcolor[HTML]{7E22CE}{2}\\
0 & \textcolor[HTML]{B45309}{1} & \textcolor[HTML]{B91C1C}{1} & \textcolor[HTML]{B91C1C}{0} & \textcolor[HTML]{7E22CE}{0}
\end{pmatrix}%
}

卢本韦看见末列第二行的二与坏根第二行的一之间也差一，却被截掉，指尖在那处停了停。

``第一张截在第三行，这张截在第二行。用的是末列最后非零项所在的行，不是整张纸一共有几行。''

``正是。''test\_alpha0 说。

他在阵盘下面又添了一条全零行，让纸面暂时变成四行。坏根的位置没变，最后非零行仍是第二行，前两行的判断也没有改变。

小舞将那条全零行折到纸后。

``补零是把坐标写齐，不是另添了一层非零结构。''

梅天狸将蔡兴坤最先写的差值留在旁边，没有让她擦去。差值与最终阶差并排放着，截断发生在哪里，一眼就能看见。

接下来，众人逐列核对坏部的提升位置。这一例只有第一行的阶差非零，三列第一行都沿父链接回坏根，故都能增加；其余行的相应阶差为零，数值保留。

\NovelDisplay{M00701R01}{%
\textcolor[HTML]{B45309}{(1,1,1)}\longmapsto \textcolor[HTML]{1D4ED8}{(1+k,1,1)}%
}
\NovelDisplay{M00701R02}{%
\textcolor[HTML]{B91C1C}{(2,1,1)}\longmapsto \textcolor[HTML]{1D4ED8}{(2+k,1,1)}%
}
\NovelDisplay{M00701R03}{%
\textcolor[HTML]{B91C1C}{(3,1)}\longmapsto \textcolor[HTML]{1D4ED8}{(3+k,1)}%
}
\NovelDisplay{M00701R04}{%
k=1,2,\ldots\quad\text{为展开项编号，右侧仅列本项新添的一段。}%
}

霍雨浩将``原列''和``展开第 \allowbreak{}\NovelInline{I03807}{k}\allowbreak{} 项时新添部分的对应列''写在表的两边。右边的式子都从原列出发：展开第一项时，新添的一段按一倍阶差计算；展开第二项时，再新添的一段按两倍阶差计算。

``两段新增内容的差，正好又是一倍阶差。''他说，``记录时仍从原值出发，写成原值加 \allowbreak{}\NovelInline{I03808}{k}\allowbreak{} 倍，便能直接核对展开第 \allowbreak{}\NovelInline{I03809}{k}\allowbreak{} 项时新添的部分。''

奥斯卡展开第一项，将新添的三列接到唐三留下的原坏部之后。诺当依次读出列号，小舞对照各列坐标，卢本韦则查新增部分的首尾。

\NovelDisplay{M00702R01}{%
\bigl((0)\textcolor[HTML]{B45309}{(1,1,1)}\textcolor[HTML]{B91C1C}{(2,1,1)}\textcolor[HTML]{B91C1C}{(3,1)}\textcolor[HTML]{7E22CE}{(2,2)}\bigr)[1]%
}
\NovelDisplay{M00702R02}{%
=(0)\textcolor[HTML]{B45309}{(1,1,1)}\textcolor[HTML]{B91C1C}{(2,1,1)}\textcolor[HTML]{B91C1C}{(3,1)}\textcolor[HTML]{1D4ED8}{(2,1,1)}%
}
\NovelDisplay{M00702R03}{%
\quad\textcolor[HTML]{1D4ED8}{(3,1,1)}\textcolor[HTML]{1D4ED8}{(4,1)}%
}

阵盘在右侧长出一段，原来的好部与坏部仍安稳留在前面。

诺当数了一遍，原来五列，删一列再加三列，结果是七列。他在新添部分的末端标上边界，停住了笔。

``变长了，位置却往下取了一步。''

唐三将原末列与展开第一项时新添部分的首列并排放在比较的位置。两者第一行都是二，第二行分别是二与一，首先出现差别的坐标已经决定顺序。

\NovelDisplay{M00703R01}{%
\text{首先不同的坐标：}\quad1<2,%
}
\NovelDisplay{M00703R02}{%
\bigl((0)(1,1,1)(2,1,1)(3,1)(2,2)\bigr)[1]%
}
\NovelDisplay{M00703R03}{%
\prec(0)(1,1,1)(2,1,1)(3,1)(2,2),%
}
\NovelDisplay{M00703R04}{%
\text{比较使用相同前缀之后首先出现区别的位置。}%
}

``展开第二项、第三项时，展开第一项所添的那一段仍保留在前面。''卢本韦顺着纸面看过去，``后面追加得更长，也不会越过这里最早出现的差别。''

test\_alpha0 点了点头，将阵盘收低。

``你们现在已经能完整走一遍。下一张只增加一个问题：同一行的阶差非零，坏部里是否每一项都跟着增加？''

蔡兴坤的笔已经伸到空纸上，听见后半句，又将笔尖抬起。

``得先查每一项的父链。''

``那就让这一回的阶差，真的有两行非零。''

第三张阵盘落到桌面，三行，八列。

\NovelDisplay{M00704R01}{%
\begin{pmatrix}
0 & 1 & 2 & 3 & 4 & 5 & 3 & 2\\
0 & 1 & 2 & 3 & 1 & 2 & 2 & 2\\
0 & 1 & 2 & 1 & 1 & 2 & 2 & 2
\end{pmatrix}%
}

\NovelDisplay{M00705R01}{%
(0)(1,1,1)(2,2,2)(3,3,1)(4,1,1)%
}
\NovelDisplay{M00705R02}{%
(5,2,2)(3,2,2)(2,2,2)%
}

卢本韦先看末列。三个二沿各行的父链都接回第二列的一个一，最后非零行在第三行，坏根仍是第二列。

\NovelDisplay{M00706R01}{%
p_1(8)=p_2(8)=p_3(8)=2,\qquad h=3,\quad r=2,%
}
\NovelDisplay{M00706R02}{%
\text{好部}=(0),%
}
\NovelDisplay{M00706R03}{%
\text{坏部为原第二至第七列，末列为原第八列。}%
}
\NovelDisplay{M00706R04}{%
(\delta_1,\delta_2,\delta_3)=(1,1,0).%
}

唐三留住第一列好部。原第二至第七列合成坏部，第八列末列取下。蔡兴坤算完阶差，在前两行下面各标一个一，第三行标零。

与上一张不同，这一次第二行是否参加提升，会直接改变各项展开时新增部分的数。

``先不急着拼完整结果。''test\_alpha0 将原第五列单独推近，``查这一列。''

原第五列是 \allowbreak{}\NovelInline{I03810}{(4,1,1)}\allowbreak{}。蔡须坤先沿第一行的四往前追，父项依次经过第四、第三、第二、第一列。霍雨浩在这些列号下面抄出对应的第二行数值。

\NovelDisplay{M00707R01}{%
\text{原第五列第一行祖先列：}\quad\textcolor[HTML]{B91C1C}{4}\to\textcolor[HTML]{B91C1C}{3}\to\textcolor[HTML]{B45309}{2}\to1,%
}
\NovelDisplay{M00707R02}{%
\text{相应的第二行数值：}\quad\textcolor[HTML]{B91C1C}{3},\ \textcolor[HTML]{B91C1C}{2},\ \textcolor[HTML]{B45309}{1},\ 0,%
}
\NovelDisplay{M00707R03}{%
a_{2,5}=1,\qquad p_2(5)=1<r=2.%
}

她的手指停在那串数值上。

``原第五列第二行是一。三太大，二太大，下一个一相等，都不能取。父项落在第一列的零。''

卢本韦看向唐三手边。

第一列在好部里，坏根是第二列。

她将原本准备写在第二行旁边的加一符号收回，留下原值一。第一行仍能增加，所以展开第一项时，这一列应成为 \allowbreak{}\NovelInline{I03811}{(5,1,1)}\allowbreak{}。

\NovelDisplay{M00708R01}{%
\textcolor[HTML]{B91C1C}{(4,\textcolor[HTML]{6B7280}{1},1)}\longmapsto\textcolor[HTML]{1D4ED8}{(4+k,\textcolor[HTML]{6B7280}{1},1)},%
}
\NovelDisplay{M00708R02}{%
\delta_2=1,\quad\text{但原第五列第二行保留 }1.%
}

``并没有删掉第二行。''小舞补了一句，``它仍复制原来的数，只是不参加这一次阶差提升。''

test\_alpha0 让他们把正确结果先留在纸上，随后才在外面的演示阵盘中，将这一项试着多推了一格。

蔡须坤掌边的细光立即向左偏了一下。唐三稳住好部，卢本韦接着托住新添的部分，两人的手之间出现了一点拉扯。

大长老将那点偏斜限制在桌面上，既没有让它回传到众人体内，也没有替他们立即消去。

``你们已经算过父项。现在看见的，是我依照这条关系搭出的练习回路。哪里接错，便去核对哪里。''

蔡须坤沿第二行把连线重新点亮。那根线越过坏根，回到第一列。

卢本韦撤下硬添的一格，让这一项带着原来的一复制过去。细光重新贴回桌面，唐三掌中的好部也恢复平稳。

``我刚才若只看它在坏部里，就会让整行都加。''她看着自己没有落下的那个符号，``父项落在好部，正好把它分开了。''

``这个判断有行号，也有列号，能复查。''大长老说，``你们感到哪里不顺，可以从那里开始找；最后还得把路走出来。''

蔡须坤应了一声，将感知到的那处牵扯标在纸边，父链则完整留在正文里。两种记录没有互相替代。

诺当接着推来原第六列 \allowbreak{}\NovelInline{I03812}{(5,2,2)}\allowbreak{}。

这一列第一行的五，父项是第五列的四。沿第一行祖先链查第二行，最近的严格更小项，是第五列的一；因此第二行父项也在第五列。

``第五列在坏部里面。''诺当把位置标出来，又看了一眼刚刚保留原值的那一个一，``但它自己第二行不能提升。''

卢本韦没有立即把第六列的二改成三。她沿第二行的父链继续追：第六列先到第五列，第五列再到第一列，中间没有经过第二列坏根。

\NovelDisplay{M00709R01}{%
p_1(6)=p_2(6)=p_3(6)=5,%
}
\NovelDisplay{M00709R02}{%
p_2(5)=1<r=2,%
}
\NovelDisplay{M00709R03}{%
\text{原第六列第二行祖先列：}\quad5\longrightarrow1,%
}
\NovelDisplay{M00709R04}{%
2\notin\{5,1\}:\quad\text{这一行未接回坏根。}%
}

``这一项的父项虽然在坏部里，父项那一行却已经停住。''她说，``所以它也不能只凭`父项在坏部'就加一。''

蔡兴坤将两列的第二行连在一起。展开第一项时，原第五、第六列在新增部分中对应的第二行分别仍是一、二。它们都保留原值，但各自保留自己的数。

\NovelDisplay{M00710R01}{%
\textcolor[HTML]{B91C1C}{(5,\textcolor[HTML]{6B7280}{2},2)}\longmapsto\textcolor[HTML]{1D4ED8}{(5+k,\textcolor[HTML]{6B7280}{2},2)},%
}
\NovelDisplay{M00710R02}{%
k=1:\ \textcolor[HTML]{1D4ED8}{(6,\textcolor[HTML]{6B7280}{2},2)},\qquad k=2:\ \textcolor[HTML]{15803D}{(7,\textcolor[HTML]{6B7280}{2},2)}.%
}

蔡须坤的目光在第五列的一与坏根那一行的一之间停了片刻。识海里，牢大的声音轻轻响起：

``数值相同，也不代表是同一个父项。''

她低低应了一声，将两个一所在的列号分别标出，再把纸转向霍雨浩。

霍雨浩把两处列号重新写重了一些。第五列的一与第二列的一相等，位置却不同，第二行的父链也不相同；寻找祖先，核对的是实际连接到哪一列。

他先前想把规则缩成``父项在坏部即可''，此时在旁边补了一段：还要查这个父项在同一行是否参加提升，沿链继续追到坏根。

test\_alpha0 等他补完，才让诺当把原第七列移过来。

第七列是 \allowbreak{}\NovelInline{I03813}{(3,2,2)}\allowbreak{}。它第一行先接第三列，第二行则沿候选链接回第二列坏根。前面两个位置没有提升，并不妨碍它凭自己的父链参加提升。

\NovelDisplay{M00711R01}{%
p_1(7)=3,\quad p_2(7)=2=r,%
}
\NovelDisplay{M00711R02}{%
(3,2,2)\longmapsto\textcolor[HTML]{1D4ED8}{(3+k,2+k,2)}.%
}

小舞将展开第一项时原第七列对应的结果 \allowbreak{}\NovelInline{I03814}{(4,3,2)}\allowbreak{} 放到 \allowbreak{}\NovelInline{I03815}{(6,2,2)}\allowbreak{} 后面。

``不能从第五列画一条横线，把右边整排都锁住。每一列按各自的父链算。''

奥斯卡将横放的尺子移开，笑着给她腾出位置。

``那我这把尺子今天只能用来排齐纸。''

小舞把几张纸推正：``这个也要紧，别把第六列抄到第七列下面。''

霍雨浩抬头看了一眼，果然将自己刚记下的列号挪回原位。几人笑了片刻，原先绷紧的手腕才渐渐松开。

六列坏部逐一核完，能够提升到的行数分别为二、二、二、一、一、二。第三行的阶差为零，无论对应位置的父项如何，数值都按本次规则保留；真正需要区分的，是前两行中哪些项允许加上非零阶差。

蔡须坤将原第五、第六列第二行的数涂成灰色。其余分部的颜色照旧，只把这两个未通过提升条件的位置单独显出来。

\NovelDisplay{M00712R01}{%
\begin{pmatrix}
0 & \textcolor[HTML]{B45309}{1} & \textcolor[HTML]{B91C1C}{2} & \textcolor[HTML]{B91C1C}{3} & \textcolor[HTML]{B91C1C}{4} & \textcolor[HTML]{B91C1C}{5} & \textcolor[HTML]{B91C1C}{3} & \textcolor[HTML]{7E22CE}{2}\\
0 & \textcolor[HTML]{B45309}{1} & \textcolor[HTML]{B91C1C}{2} & \textcolor[HTML]{B91C1C}{3} & \textcolor[HTML]{6B7280}{1} & \textcolor[HTML]{6B7280}{2} & \textcolor[HTML]{B91C1C}{2} & \textcolor[HTML]{7E22CE}{2}\\
0 & \textcolor[HTML]{B45309}{1} & \textcolor[HTML]{B91C1C}{2} & \textcolor[HTML]{B91C1C}{1} & \textcolor[HTML]{B91C1C}{1} & \textcolor[HTML]{B91C1C}{2} & \textcolor[HTML]{B91C1C}{2} & \textcolor[HTML]{7E22CE}{2}
\end{pmatrix}%
}

第五列的一、第六列的二虽都保留原值，灰色没有把它们合成同一个数。第三行因阶差为零而保留，仍按原分部着色。

\NovelDisplay{M00713R01}{%
\textcolor[HTML]{B45309}{(1,1,1)}\longmapsto \textcolor[HTML]{1D4ED8}{(1+k,1+k,1)}%
}
\NovelDisplay{M00713R02}{%
\textcolor[HTML]{B91C1C}{(2,2,2)}\longmapsto \textcolor[HTML]{1D4ED8}{(2+k,2+k,2)}%
}
\NovelDisplay{M00713R03}{%
\textcolor[HTML]{B91C1C}{(3,3,1)}\longmapsto \textcolor[HTML]{1D4ED8}{(3+k,3+k,1)}%
}
\NovelDisplay{M00713R04}{%
\textcolor[HTML]{B91C1C}{(4,\textcolor[HTML]{6B7280}{1},1)}\longmapsto \textcolor[HTML]{1D4ED8}{(4+k,\textcolor[HTML]{6B7280}{1},1)}%
}
\NovelDisplay{M00713R05}{%
\textcolor[HTML]{B91C1C}{(5,\textcolor[HTML]{6B7280}{2},2)}\longmapsto \textcolor[HTML]{1D4ED8}{(5+k,\textcolor[HTML]{6B7280}{2},2)}%
}
\NovelDisplay{M00713R06}{%
\textcolor[HTML]{B91C1C}{(3,2,2)}\longmapsto \textcolor[HTML]{1D4ED8}{(3+k,2+k,2)}%
}
\NovelDisplay{M00713R07}{%
k=1,2,\ldots\quad\text{为展开项编号，右侧仅列本项新添的一段。}%
}

蔡兴坤负责代入展开第一项所用的编号，小舞与诺当接列，唐三保持原前缀。卢本韦逐一读出第五、第六、第七列对应的结果，让霍雨浩最后核一遍。

\NovelDisplay{M00714R01}{%
\bigl((0)\textcolor[HTML]{B45309}{(1,1,1)}\textcolor[HTML]{B91C1C}{(2,2,2)}\textcolor[HTML]{B91C1C}{(3,3,1)}\textcolor[HTML]{B91C1C}{(4,1,1)}\textcolor[HTML]{B91C1C}{(5,2,2)}\textcolor[HTML]{B91C1C}{(3,2,2)}\textcolor[HTML]{7E22CE}{(2,2,2)}\bigr)[1]%
}
\NovelDisplay{M00714R02}{%
=(0)\textcolor[HTML]{B45309}{(1,1,1)}\textcolor[HTML]{B91C1C}{(2,2,2)}\textcolor[HTML]{B91C1C}{(3,3,1)}\textcolor[HTML]{B91C1C}{(4,\textcolor[HTML]{6B7280}{1},1)}%
}
\NovelDisplay{M00714R03}{%
\quad\textcolor[HTML]{B91C1C}{(5,\textcolor[HTML]{6B7280}{2},2)}\textcolor[HTML]{B91C1C}{(3,2,2)}\textcolor[HTML]{1D4ED8}{(2,2,1)}\textcolor[HTML]{1D4ED8}{(3,3,2)}\textcolor[HTML]{1D4ED8}{(4,4,1)}%
}
\NovelDisplay{M00714R04}{%
\quad\textcolor[HTML]{1D4ED8}{(5,\textcolor[HTML]{6B7280}{1},1)}\textcolor[HTML]{1D4ED8}{(6,\textcolor[HTML]{6B7280}{2},2)}\textcolor[HTML]{1D4ED8}{(4,3,2)}%
}

这一次，众人没有再从左至右抄第二张相同的大表。展开第二项时再新增的一段，已经可以直接由通式代入二算出。霍雨浩只抽查了原第五、原第六两列：分别成为 \allowbreak{}\NovelInline{I03816}{(6,1,1)}\allowbreak{} 与 \allowbreak{}\NovelInline{I03817}{(7,2,2)}\allowbreak{}。

``第一行继续增加，第二行一直带着各自原值。''他在这两处落下小勾，``不能把展开第一项时不提升，理解为只有第一项保留原值。''

卢本韦将笔放下，原以为这一张已经做完。test\_alpha0 却没有收阵盘，只伸指把原第五列第二行的一改成了二。

其余位置没有改，末列也没有改。

``这回先别重抄整张表。哪些关系必须重新查？''

蔡兴坤立即按住第五列和第六列之间的线。

``这一列自己的父项，还有后面把它当父项的那些位置。''

她沿第五列第一行的祖先链，重新比较第二行的三、二、一、零。原值变成二以后，三仍太大，二仍相等，第二列的一已经严格更小，父项便落回坏根。

\NovelDisplay{M00715R01}{%
\text{原第五列 }(4,1,1)\longrightarrow(4,2,1),%
}
\NovelDisplay{M00715R02}{%
a_{2,5}=2:\quad p_2(5)=2=r,%
}
\NovelDisplay{M00715R03}{%
(4,2,1)\longmapsto\textcolor[HTML]{1D4ED8}{(4+k,2+k,1)},%
}
\NovelDisplay{M00715R04}{%
k=1:\quad\textcolor[HTML]{1D4ED8}{(5,3,1)}.%
}

紧接着，诺当查第六列。它第二行仍是二，刚改过的第五列也是二，不能再取第五列作父项；继续沿候选向前，最后同样落到第二列的一。

两列第二行都接回了坏根，便都参加阶差提升。

\NovelDisplay{M00716R01}{%
\text{第五列改成 }(4,2,1)\text{ 后：}\quad p_2(6)=2=r,%
}
\NovelDisplay{M00716R02}{%
(5,2,2)\longmapsto\textcolor[HTML]{1D4ED8}{(5+k,2+k,2)},%
}
\NovelDisplay{M00716R03}{%
k=1:\quad\textcolor[HTML]{1D4ED8}{(6,3,2)}.%
}

卢本韦将改动前后各展开第一项的局部结果并排留下：原先是 \allowbreak{}\NovelInline{I03818}{(5,1,1)}\allowbreak{}、\allowbreak{}\NovelInline{I03819}{(6,2,2)}\allowbreak{}，改动后是 \allowbreak{}\NovelInline{I03820}{(5,3,1)}\allowbreak{}、\allowbreak{}\NovelInline{I03821}{(6,3,2)}\allowbreak{}。

前者第一处的变化包含两件事：输入本身从一改成二，随后在新展开中再加一；后者输入仍是二，却因父项重新选择而开始参加提升。

小舞用笔圈住原列，让两次变化的来源始终摆在结果旁边。

``不能只盯着答案差了多少。先改了什么，哪条父链因此改变，再由新父链决定怎么复制。''

test\_alpha0 将新旧两张小纸交给卢本韦。

``你开始问，后面加了层次，前面的接法是否也要改。先把具体关系查出来，答案才有落处。''

卢本韦把两张纸夹进双魂核运转图的旁页。她没有立即把这一列的数照搬进魂核图里，只在那个交会点边上写下：先查连接，再定调整。

梅天狸给众人重新添了水。窗外的云影已经移过一段石阶，霍雨浩本子上的墨也干了一半。

大长老等他们歇了片刻，才铺开下一张纸。

``三行里，父链、坏根、阶差与提升条件都已经见过。现在把这些放到四行，交给你们自己做。''

\NovelDisplay{M00717R01}{%
\begin{pmatrix}
0 & 1 & 2 & 3 & 4 & 5 & 6 & 7\\
0 & 1 & 2 & 3 & 2 & 1 & 2 & 3\\
0 & 1 & 2 & 1 & 0 & 1 & 2 & 1\\
0 & 1 & 1 & 1 & 0 & 1 & 1 & 1
\end{pmatrix}%
}

\NovelDisplay{M00718R01}{%
(0)(1,1,1,1)(2,2,2,1)(3,3,1,1)(4,2)%
}
\NovelDisplay{M00718R02}{%
(5,1,1,1)(6,2,2,1)(7,3,1,1)%
}

没有人急着按整列相减。

诺当先把八列排齐，唐三留出分割好部、坏部的位置；蔡兴坤等父项确定以后再算阶差。小舞与卢本韦各拿一张空纸，分别追末列的父链，霍雨浩在两人之间核对列号。

末列第一行是七，向左的父链依次经过六、五、四、三、二、一、零，七个祖先列都在其中。

第二行末项是三。沿第一行的祖先链比较，最近的严格更小项是第七列的二；沿第二行继续追，第七列的二接第六列的一，第六列的一接第一列的零。

第二行的祖先链于是缩短成第七、第六、第一列。

\NovelDisplay{M00719R01}{%
\text{第一行末项祖先列：}\quad\textcolor[HTML]{B91C1C}{7}\to\textcolor[HTML]{B91C1C}{6}\to\textcolor[HTML]{B91C1C}{5}\to\textcolor[HTML]{B91C1C}{4}\to\textcolor[HTML]{B91C1C}{3}\to\textcolor[HTML]{B91C1C}{2}\to\textcolor[HTML]{B45309}{1},%
}
\NovelDisplay{M00719R02}{%
\text{第二行末项祖先列：}\quad\textcolor[HTML]{B91C1C}{7}\to\textcolor[HTML]{B91C1C}{6}\to\textcolor[HTML]{B45309}{1}.%
}

蔡须坤将这三列从光阵里轻轻托起，第三行原末项一仍在右边。

``现在给第三行挑的，就是这三处。''

第七列第三行是二，太大；第六列第三行是一，相等。直到第一列的零，才找到严格更小项。第三行选完以后，自己的祖先链只剩第一列。

小舞在第四行拉出一条最短的线。

第四行末项也是一，沿第三行留下的链，父项同样落在第一列的零。

\NovelDisplay{M00720R01}{%
\text{第三行末项祖先列：}\quad\textcolor[HTML]{B45309}{1},%
}
\NovelDisplay{M00720R02}{%
\text{第四行末项祖先列：}\quad\textcolor[HTML]{B45309}{1}.%
}

卢本韦与小舞交换草纸。两人连出的父项一致，卢本韦却把自己第三行旁边已经划去的第六列保留下来，没有重新画进第四行。

``我起初想让第四行继续用第二行筛下来的三列。''她用笔尾点着那道删线，``第三行已经选过，就得以第三行的实际祖先链为准。''

小舞把自己纸上同一处指给她。

``我也在这里停了一下。多一行，便多经过这一层限制。''

霍雨浩核完两张纸，才请蔡兴坤落下阶差。她先确认最后非零行是第四行，父项在第一列；唐三随即让好部的位置空着，将原第一至第七列连成坏部。

\NovelDisplay{M00721R01}{%
h=4,\quad p_4(8)=1,\quad r=1,\quad\text{好部}=\varnothing,%
}
\NovelDisplay{M00721R02}{%
\text{坏部为原第一至第七列；原末列不参加复制。}%
}
\NovelDisplay{M00721R03}{%
(\delta_1,\delta_2,\delta_3,\delta_4)=(7,3,1,0).%
}

最后一列 \allowbreak{}\NovelInline{I03822}{(7,3,1,1)}\allowbreak{} 删除。坏根本身是零列，所以前三行差值分别为七、三、一；第四行到了截断处，阶差置零。

唐三将最前面那列零复制到展开第一项时应新添的位置，先把它与后面几列分开。

``根列要先作为起点，和第一张一样。后面的列再沿父链检查，不能把根本身也排除。''

test\_alpha0 从众人身后走到桌角，将坏根与刚才被删除的末列连了一下。

``把这个起点写明。''大长老说，``定下坏根以后，坏根在最后非零行之前参与提升；坏部其余各列，再沿父链逐行检查是否接到它。''

他没有给零列编出一串父项，而是在根的位置单独放了一枚小小的光点。

``根的处理先写明，后面的检查才有起点。不能拿后面各列的条件，把根本身排除。''

\NovelDisplay{M00722R01}{%
\textcolor[HTML]{B45309}{(0)}\longmapsto\textcolor[HTML]{1D4ED8}{(7k,3k,k)},\qquad k\ge1,%
}
\NovelDisplay{M00722R02}{%
k=1:\ \textcolor[HTML]{1D4ED8}{(7,3,1)},\qquad k=2:\ \textcolor[HTML]{15803D}{(14,6,2)}.%
}

展开第一项时，零列在新增部分中成为 \allowbreak{}\NovelInline{I03823}{(7,3,1)}\allowbreak{}；展开第二项时，再新增的对应列成为 \allowbreak{}\NovelInline{I03824}{(14,6,2)}\allowbreak{}。补齐后最下面一项仍是零，按记法省略。

蔡须坤把两项各自新添的根列并排留好，再往后查各列。

她最先停在原第五列 \allowbreak{}\NovelInline{I03825}{(4,2)}\allowbreak{}。补齐以后，第三行与第四行都是零，第三行没有严格更小的非负项可作父项；这一列允许增加到第二行，第三行不参加提升。

唐三则在原第四列 \allowbreak{}\NovelInline{I03826}{(3,3,1,1)}\allowbreak{} 下方找到另一条线。

这一列第三行的一确有父项，直接落在第一列坏根。它虽然只是一个一，仍能增加第三行的阶差。

\NovelDisplay{M00723R01}{%
p_3(5)\text{ 不存在：}\quad(4,2)\longmapsto\textcolor[HTML]{1D4ED8}{(4+7k,2+3k)},%
}
\NovelDisplay{M00723R02}{%
p_3(4)=1=r:\quad(3,3,1,1)\longmapsto\textcolor[HTML]{1D4ED8}{(3+7k,3+3k,1+k,1)}.%
}

蔡须坤将第三行的两个位置分开点亮，一个保留零，一个由一成为二。

``这一项留着零，是父项条件没有通过。旁边的一能增加，是它的父链接到了根。不能看见某一行里有零，就把整行都当作不能动。''

她说完，又把自己掌边的光压低了些，好让卢本韦看见那两根各自的线。

霍雨浩将原第五列第三行的零标成灰色，原第四列第三行的一仍留红色。坏根是橙色的第一列零，最后待删的一列仍标紫色。

\NovelDisplay{M00724R01}{%
\begin{pmatrix}
\textcolor[HTML]{B45309}{0} & \textcolor[HTML]{B91C1C}{1} & \textcolor[HTML]{B91C1C}{2} & \textcolor[HTML]{B91C1C}{3} & \textcolor[HTML]{B91C1C}{4} & \textcolor[HTML]{B91C1C}{5} & \textcolor[HTML]{B91C1C}{6} & \textcolor[HTML]{7E22CE}{7}\\
\textcolor[HTML]{B45309}{0} & \textcolor[HTML]{B91C1C}{1} & \textcolor[HTML]{B91C1C}{2} & \textcolor[HTML]{B91C1C}{3} & \textcolor[HTML]{B91C1C}{2} & \textcolor[HTML]{B91C1C}{1} & \textcolor[HTML]{B91C1C}{2} & \textcolor[HTML]{7E22CE}{3}\\
\textcolor[HTML]{B45309}{0} & \textcolor[HTML]{B91C1C}{1} & \textcolor[HTML]{B91C1C}{2} & \textcolor[HTML]{B91C1C}{1} & \textcolor[HTML]{6B7280}{0} & \textcolor[HTML]{B91C1C}{1} & \textcolor[HTML]{B91C1C}{2} & \textcolor[HTML]{7E22CE}{1}\\
\textcolor[HTML]{B45309}{0} & \textcolor[HTML]{B91C1C}{1} & \textcolor[HTML]{B91C1C}{1} & \textcolor[HTML]{B91C1C}{1} & \textcolor[HTML]{B91C1C}{0} & \textcolor[HTML]{B91C1C}{1} & \textcolor[HTML]{B91C1C}{1} & \textcolor[HTML]{7E22CE}{1}
\end{pmatrix}%
}

两处都在第三行，灰色的零不参加提升，红色的一则能在各项展开时加上相应倍数的本行阶差。第四行的阶差本来就是零，便不因数值保持原样而另标灰色。

卢本韦核完后，展开第一项，将它们接进新增的部分，动作比先前三行时更稳。唐三托住新增的根列，奥斯卡负责把右侧的纸往外展开，桌面不够宽，便将空墨碟移到了窗边。

七列坏部逐一完成检查，可提升到的行数为三、三、三、三、二、三、三。

\NovelDisplay{M00725R01}{%
\textcolor[HTML]{B45309}{(0)}\longmapsto \textcolor[HTML]{1D4ED8}{(7k,3k,k)}%
}
\NovelDisplay{M00725R02}{%
\textcolor[HTML]{B91C1C}{(1,1,1,1)}\longmapsto \textcolor[HTML]{1D4ED8}{(1+7k,1+3k,1+k,1)}%
}
\NovelDisplay{M00725R03}{%
\textcolor[HTML]{B91C1C}{(2,2,2,1)}\longmapsto \textcolor[HTML]{1D4ED8}{(2+7k,2+3k,2+k,1)}%
}
\NovelDisplay{M00725R04}{%
\textcolor[HTML]{B91C1C}{(3,3,1,1)}\longmapsto \textcolor[HTML]{1D4ED8}{(3+7k,3+3k,1+k,1)}%
}
\NovelDisplay{M00725R05}{%
\textcolor[HTML]{B91C1C}{(4,2)}\longmapsto \textcolor[HTML]{1D4ED8}{(4+7k,2+3k)}%
}
\NovelDisplay{M00725R06}{%
\textcolor[HTML]{B91C1C}{(5,1,1,1)}\longmapsto \textcolor[HTML]{1D4ED8}{(5+7k,1+3k,1+k,1)}%
}
\NovelDisplay{M00725R07}{%
\textcolor[HTML]{B91C1C}{(6,2,2,1)}\longmapsto \textcolor[HTML]{1D4ED8}{(6+7k,2+3k,2+k,1)}%
}
\NovelDisplay{M00725R08}{%
k=1,2,\ldots\quad\text{为展开项编号，右侧仅列本项新添的一段。}%
}

霍雨浩一边写通式，一边让蔡兴坤报出展开第一项时的两处新增结果：原第四列对应 \allowbreak{}\NovelInline{I03827}{(10,6,2,1)}\allowbreak{}，原第五列对应 \allowbreak{}\NovelInline{I03828}{(11,5)}\allowbreak{}。展开第二项时，再新增的对应两列分别是 \allowbreak{}\NovelInline{I03829}{(17,9,3,1)}\allowbreak{} 与 \allowbreak{}\NovelInline{I03830}{(18,8)}\allowbreak{}。

蔡兴坤报完，又将展开第二项时原第五列对应的新添一列补齐给众人看：最下面两个零仍是零。

``这里从原值出发，加两倍阶差。不能因为展开所得的结果已经接得很长，就把本来不参加的两行也补上。''

小舞展开第一项，将新添的七列按顺序接好。唐三松开原坏部之前，先让诺当核完所有边界：好部为空，原坏部仍保留，末列已删，新增七列从本项新添的根列开始。

\NovelDisplay{M00726R01}{%
\bigl(\textcolor[HTML]{B45309}{(0)}\textcolor[HTML]{B91C1C}{(1,1,1,1)}\textcolor[HTML]{B91C1C}{(2,2,2,1)}\textcolor[HTML]{B91C1C}{(3,3,1,1)}\textcolor[HTML]{B91C1C}{(4,2)}\textcolor[HTML]{B91C1C}{(5,1,1,1)}\textcolor[HTML]{B91C1C}{(6,2,2,1)}\textcolor[HTML]{7E22CE}{(7,3,1,1)}\bigr)[1]%
}
\NovelDisplay{M00726R02}{%
=\textcolor[HTML]{B45309}{(0)}\textcolor[HTML]{B91C1C}{(1,1,1,1)}\textcolor[HTML]{B91C1C}{(2,2,2,1)}\textcolor[HTML]{B91C1C}{(3,3,1,1)}\textcolor[HTML]{B91C1C}{(4,2)}%
}
\NovelDisplay{M00726R03}{%
\quad\textcolor[HTML]{B91C1C}{(5,1,1,1)}\textcolor[HTML]{B91C1C}{(6,2,2,1)}\textcolor[HTML]{1D4ED8}{(7,3,1)}\textcolor[HTML]{1D4ED8}{(8,4,2,1)}\textcolor[HTML]{1D4ED8}{(9,5,3,1)}%
}
\NovelDisplay{M00726R04}{%
\quad\textcolor[HTML]{1D4ED8}{(10,6,2,1)}\textcolor[HTML]{1D4ED8}{(11,5)}\textcolor[HTML]{1D4ED8}{(12,4,2,1)}\textcolor[HTML]{1D4ED8}{(13,5,3,1)}%
}

光阵向右铺开，展开第一项时新添的七列清清楚楚地排在原坏部后面。最前面的零列没有移动，新添部分中的根列却已不再是零。

诺当俯身比较两处不同的根列，拿笔沿原列与新增部分之间的连线划了一遍。

``复制的是它在这一回结构中的位置关系，再按规则改数。不是找一列长得一样的东西放过去。''

test\_alpha0 点了点头，示意众人将阵盘降到纸面。随着众人松手，光没有散，依照刚才核好的连接平稳运转了一遍。

大长老将四张主例按顺序摆开，没有再铺第五张更高、更宽的矩阵。

``从三行换到四行，你们方才增加的是追链与核对的工作。若再增加有限行数，规则怎么写？''

霍雨浩拉过最后一张空纸。蔡兴坤将第一张最短三行结构放在左上角，卢本韦把第三张的两条不同父链压在右边；一张纸记录根怎样确定，另一张纸提醒他们哪些项不能整行照加。

``先定记号。''霍雨浩说。

他把矩阵写成有限的 \allowbreak{}\NovelInline{I03831}{n}\allowbreak{} 行、\allowbreak{}\NovelInline{I03832}{m}\allowbreak{} 列，第 \allowbreak{}\NovelInline{I03833}{j}\allowbreak{} 列第 \allowbreak{}\NovelInline{I03834}{i}\allowbreak{} 行的元素记作 \allowbreak{}\NovelInline{I03835}{a_{i,j}}\allowbreak{}。每个元素都是非负整数，不同长度的列在比较和找父项时补齐末尾零。

\NovelDisplay{M00727R01}{%
a_{i,j}\in\{0,1,2,\ldots\},\quad1\le i\le n,\quad1\le j\le m.%
}

第一行父项向左取最近的严格更小项。其余各行沿上一行已经确定的祖先链，按从近到远的次序，取本行第一个严格更小项。若找不到，这一项便没有父项，也不能凭空给下一行添候选。

小舞将前面的父项公式翻到这一页旁边，没有再抄一遍。蔡须坤沿原第四个例子的四行末项，核对这句话有没有漏掉``上一行已经确定''几个字。

接着，霍雨浩将空矩阵与末列全零的情况写在页边。它们是已经掌握的起始情形，不必为了多行另换一套规则。

\NovelDisplay{M00728R01}{%
\text{空矩阵}\longleftrightarrow0,\quad(0)\longleftrightarrow1,\quad(0)(0)\longleftrightarrow2,%
}
\NovelDisplay{M00728R02}{%
\text{末列全零：对应序数}=\text{删末列后的对应序数}+1,%
}
\NovelDisplay{M00728R03}{%
\bigl((0)(1)\bigr)[k]=\underbrace{(0)(0)\cdots(0)}_{k+1\text{ 列}}\quad(k\ge0).%
}

``末列非零时，才往下写最后非零行与坏根。''他说。

\NovelDisplay{M00729R01}{%
1\le n,m<\omega,\qquad h=\max\{i:a_{i,m}>0\},%
}
\NovelDisplay{M00729R02}{%
r=p_h(m),\qquad\text{坏根是原第 }r\text{ 列},%
}
\NovelDisplay{M00729R03}{%
\delta_i=\begin{cases}a_{i,m}-a_{i,r},&1\le i<h,\\0,&h\le i\le n.\end{cases}%
}

这里的 \allowbreak{}\NovelInline{I03836}{h}\allowbreak{} 是末列最后非零行号，\allowbreak{}\NovelInline{I03837}{r}\allowbreak{} 是这一行父项所在的列号。最后非零行及其下方的阶差全部为零，上方各行才取末列减坏根。

蔡兴坤将三张三行例子并排核对。最短那张与第三张的 \allowbreak{}\NovelInline{I03838}{h}\allowbreak{} 都是三，第二张的 \allowbreak{}\NovelInline{I03839}{h}\allowbreak{} 却是二；纸面同为三行，阶差的截断位置仍有区别。

``若末列非零，却找不到它在这一行的父项呢？''奥斯卡问。

``先检查输入与前面的父链。''test\_alpha0 指着 \allowbreak{}\NovelInline{I03840}{r=p_h(m)}\allowbreak{}，``这里的通式以相关父项存在为前提。不能因为急着取坏根，随意补一个列号。''

唐三随后在纸上划出各部分边界。

\NovelDisplay{M00730R01}{%
\text{好部：原第 }1,\ldots,r-1\text{ 列； }r=1\text{ 时为空},%
}
\NovelDisplay{M00730R02}{%
\text{坏部：原第 }r,\ldots,m-1\text{ 列；末列单独删去},%
}
\NovelDisplay{M00730R03}{%
\text{删末列后，原好部与原坏部均保留。}%
}

他把坏根一列圈在坏部里，又将末列圈到外面。原好部、原坏部均保留，各项展开时新添的部分接在其后。

``定下整体阶差，还没决定每个位置都加。''卢本韦把第三张里那两列推到公式边，``最后要把这一项的条件写出来。''

霍雨浩为每个坏部列记一个数 \allowbreak{}\NovelInline{I03841}{t_j}\allowbreak{}，表示这一列从第一行开始，允许提升到哪一行。它记录的是一段从上向下连续通过的行，不表示列中最大的数，也不表示这一列有几个非零项。

坏根先设 \allowbreak{}\NovelInline{I03842}{t_r=h-1}\allowbreak{}。其余列从左往右计算，每列从第一行往下核对。由于父项总在左边，查到某个父项时，它允许提升到的行数已经算过。

\NovelDisplay{M00731R01}{%
t_r=h-1,%
}
\NovelDisplay{M00731R02}{%
j=r+1,\ldots,m-1:\quad\text{逐行检查 }1\le i<h,%
}
\NovelDisplay{M00731R03}{%
\text{本行失败}\iff p_i(j)\text{ 不存在，或存在但有}%
}
\NovelDisplay{M00731R04}{%
p_i(j)<r\quad\text{或}\quad t_{p_i(j)}<i,%
}
\NovelDisplay{M00731R05}{%
t_j=\min\bigl(\{i-1:1\le i<h,\ \text{第 }i\text{ 行失败}\}\cup\{h-1\}\bigr).%
}

诺当将失败的三种情形分别指给众人：父项不存在；父项落在坏根之前；父项虽然在坏部中，却没有在这一行通过提升条件。

``第三张的第五列，是第二种。''蔡须坤说。

``第六列，是第三种。''卢本韦接着说，``它父项在第五列，位置仍在坏部，第二行却接不上参加提升的根。''

小舞指向四行例子的第五列。

``这里第三行没父项，是第一种。''

霍雨浩将对应位置写到公式下方，三种失败原因各有一个已经算过的落点。

到了第一处失败的行，这一列的可提升范围便在上一行结束；若一直没有失败，则允许提升到第 \allowbreak{}\NovelInline{I03843}{h-1}\allowbreak{} 行。后面的行即便有别的数，也不能跳过中间已经失败的一行，另起一段继续加。

``这样检查，沿同一行父链继续向前，最终要接回参加提升的坏根。''test\_alpha0 说，``父项在左边，左边的结论已经有了，便可以一列一列算下去。''

他让霍雨浩用第三张原第七列再核一次。它虽然排在两个不能提升第二行的位置之后，父项却直接接回坏根，因此 \allowbreak{}\NovelInline{I03844}{t_7=2}\allowbreak{}，没有被前两列的失败横向截住。

``是按自己的父项继承条件。''霍雨浩把最后一句写全，``不是按紧挨着的前一列继承。''

最后，各项展开时新添部分中的数可以统一写出。

\NovelDisplay{M00732R01}{%
a_{i,j}^{(k)}=\begin{cases}a_{i,j}+k\delta_i,&i\le t_j,\\a_{i,j},&i>t_j,\end{cases}%
}
\NovelDisplay{M00732R02}{%
r\le j<m,\quad1\le i\le n,\quad k=1,2,\ldots.%
}

展开第 \allowbreak{}\NovelInline{I03845}{k}\allowbreak{} 项时，最后新添的一段从原坏部的对应列算起：允许提升的行，加 \allowbreak{}\NovelInline{I03846}{k}\allowbreak{} 倍本行阶差；其余行保留原值。完整的第 \allowbreak{}\NovelInline{I03847}{k}\allowbreak{} 项，还包括前面按一倍至 \allowbreak{}\NovelInline{I03848}{k-1}\allowbreak{} 倍阶差计算得到的各段新增内容。这里的``保留''，包括保留零，也包括保留原有的非零数。

蔡兴坤将 \allowbreak{}\NovelInline{I03849}{(4,1,1)}\allowbreak{} 对应的那一行代入检查，卢本韦则用四行例子的 \allowbreak{}\NovelInline{I03850}{(3,3,1,1)}\allowbreak{} 检查第三行是否会正确增加。两处结果都与先前相同。

诺当数完列数，在最下方添了一行。

\NovelDisplay{M00733R01}{%
\text{展开每一项时新添的列数}=m-r,%
}
\NovelDisplay{M00733R02}{%
\text{第 }k\text{ 项总列数}=m-1+k(m-r)\quad(k\ge0).%
}

原矩阵 \allowbreak{}\NovelInline{I03851}{m}\allowbreak{} 列，先删一列；每展开一项，就按相应倍数的阶差新添一段 \allowbreak{}\NovelInline{I03852}{m-r}\allowbreak{} 列，所以展开第 \allowbreak{}\NovelInline{I03853}{k}\allowbreak{} 项时，共新添 \allowbreak{}\NovelInline{I03854}{k(m-r)}\allowbreak{} 列。这个数可以检查有没有漏列、重复列，却不能代替逐项找父项，也不能代替比较序数大小。

``每回的 \allowbreak{}\NovelInline{I03855}{k}\allowbreak{} 都是有限的。''诺当说。

``对。''test\_alpha0 将最右端还没画出的地方留空，``写省略号，是说明这一族取项可以继续。你们眼下写出的、用来运转的，仍是某个有限项。''

蔡须坤望了一眼桌边的双魂核图。她先前习惯把光阵越铺越长的样子叫作``无限展开''，如今落笔时却在旁边补上了取项编号。

她知道自己具体算了哪一项，也知道尚未算出的部分并未因为一个省略号就全部完成。

test\_alpha0 将第一张最短结构移到一般式下面。

``它也可以检验行数的变化。''

末列写有限个一，首列全零。无论取多少个有限行数，各行父项都回到首列，最后非零行之前的阶差均为一，最后非零行的阶差为零。

\NovelDisplay{M00734R01}{%
n\ge2,\quad(0)(\underbrace{1,\ldots,1}_{n\text{ 个}}),%
}
\NovelDisplay{M00734R02}{%
h=n,\quad r=1,\quad(\delta_1,\ldots,\delta_n)=(\underbrace{1,\ldots,1}_{n-1\text{ 个}},0),%
}
\NovelDisplay{M00734R03}{%
\text{展开第 }q\text{ 项时新添的部分}=(\underbrace{q,\ldots,q}_{n-1\text{ 个}}),\qquad1\le q\le k,%
}
\NovelDisplay{M00734R04}{%
\bigl((0)(\underbrace{1,\ldots,1}_{n\text{ 个}})\bigr)[k]\quad(k\ge2)%
}
\NovelDisplay{M00734R05}{%
=(0)(\underbrace{1,\ldots,1}_{n-1\text{ 个}})(\underbrace{2,\ldots,2}_{n-1\text{ 个}})\cdots%
}
\NovelDisplay{M00734R06}{%
\quad(\underbrace{k,\ldots,k}_{n-1\text{ 个}}),%
}
\NovelDisplay{M00734R07}{%
k=0:\ (0),\qquad k=1:\ (0)(\underbrace{1,\ldots,1}_{n-1\text{ 个}}).%
}

小舞将 \allowbreak{}\NovelInline{I03856}{n=3}\allowbreak{} 代入，得到开头已经算过的三行例子；再代入四，只需多保留一行参与增加的坐标。

``这里再增加一行，坏根仍在首列，只是截断处往下移一行。''她将通式与页首对齐，``各项展开时新添的部分，便多保留一个随编号增加的坐标。''

霍雨浩在代入 \allowbreak{}\NovelInline{I03857}{n=3}\allowbreak{} 的那一行旁落下小勾，将通式夹在四张纸的后面。若忘了哪一项怎样处理，还能顺着行号查回原例。

他正要把一般式收进夹页，test\_alpha0 又推来一张较窄的纸。

``计算会做了，还要留下表达式的来源。''

窄纸上写着七列目标。

\NovelDisplay{M00735R01}{%
(0)(1,1,1)(2,1)(1,1)(2,2,1)(3,1)(4,2)%
}

唐三看见最前面的零列，蔡兴坤先检查同列上下的数，卢本韦则沿各非零项找父项。他们把已经熟悉的局部检查做完，霍雨浩在旁边留下条件。

\NovelDisplay{M00736R01}{%
a_{i,1}=0\quad(1\le i\le n),%
}
\NovelDisplay{M00736R02}{%
a_{i+1,j}\le a_{i,j}\quad(1\le i<n,\ 1\le j\le m),%
}
\NovelDisplay{M00736R03}{%
a_{i,j}>0\ \Longrightarrow\ p_i(j)\text{ 存在，且}%
}
\NovelDisplay{M00736R04}{%
a_{i,j}\le a_{i,p_i(j)}+1.%
}

``它们是标准表达式应当满足的检查条件。''test\_alpha0 说，``但若要确认眼前这一张，光说`看着都过了'还没有把生成过程交代出来。''

霍雨浩看着目标，在纸边写下``从哪里来''。

大长老指向他写的字。

``今天给它一条实际能复算的路径。每个箭头都按同一套取项规则写。''

他先说明边界处的记法：从 BMS 边界取第 \allowbreak{}\NovelInline{I03858}{n}\allowbreak{} 项，得到一列零与一列 \allowbreak{}\NovelInline{I03859}{n+1}\allowbreak{} 个一。这里每个实际取出的矩阵都只有有限行，不把一列无穷多个一当作普通有限输入。

\NovelDisplay{M00737R01}{%
\text{BMS 边界}[n]=(0)(\underbrace{1,\ldots,1}_{n+1\text{ 个}}),%
}
\NovelDisplay{M00737R02}{%
n=0,1,2,\ldots,\qquad\text{本次从 }n=3\text{ 出发：}%
}
\NovelDisplay{M00737R03}{%
\text{BMS 边界}[3]=(0)(1,1,1,1).%
}

这一次从四行的入口出发。蔡兴坤认出它与页首的最短结构属于同一族，先展开第二项，得到零列以及 \allowbreak{}\NovelInline{I03860}{(1,1,1)}\allowbreak{}、\allowbreak{}\NovelInline{I03861}{(2,2,2)}\allowbreak{} 两列。后面的每一步，则重新核对当前末列。

\NovelDisplay{M00738R01}{%
(0)(1,1,1,1)%
}
\NovelDisplay{M00738R02}{%
\xrightarrow{[2]}(0)(1,1,1)(2,2,2)%
}
\NovelDisplay{M00738R03}{%
\xrightarrow{[1]}(0)(1,1,1)(2,2,1)%
}
\NovelDisplay{M00738R04}{%
\xrightarrow{[1]}(0)(1,1,1)(2,2)(3,3,1)%
}
\NovelDisplay{M00738R05}{%
\xrightarrow{[0]}(0)(1,1,1)(2,2)%
}

走到 \allowbreak{}\NovelInline{I03862}{(2,2,1)}\allowbreak{} 时，卢本韦先算出展开第一项会新添两列。她让小舞暂缓收光，霍雨浩记下 \allowbreak{}\NovelInline{I03863}{[1]}\allowbreak{} 以后，再将两列完整接好。

结果的尾部是 \allowbreak{}\NovelInline{I03864}{(2,2)(3,3,1)}\allowbreak{}。接着取一次 \allowbreak{}\NovelInline{I03865}{[0]}\allowbreak{}，删去最后一列，才只留下所需的 \allowbreak{}\NovelInline{I03866}{(2,2)}\allowbreak{}。

``不能看见目标要少一个坐标，就把原末列最后的一直接擦去。''卢本韦将两个编号连起来，``那样没有说明中间用的哪条规则。''

小舞照着她的记录收回末列，另一手仍压着前面的 \allowbreak{}\NovelInline{I03867}{(2,2)}\allowbreak{}。

``生成以后再截去，纸上便要留下这两步。''

后面的路径开始出现更长的坏部。唐三负责分段，诺当报展开第一项时应新添多少列；小舞按列接出，卢本韦等新增部分完整落定以后才记录截断。

\NovelDisplay{M00739R01}{%
(0)(1,1,1)(2,2)%
}
\NovelDisplay{M00739R02}{%
\xrightarrow{[1]}(0)(1,1,1)(2,1,1)%
}
\NovelDisplay{M00739R03}{%
\xrightarrow{[1]}(0)(1,1,1)(2,1)(3,2,1)%
}
\NovelDisplay{M00739R04}{%
\xrightarrow{[1]}(0)(1,1,1)(2,1)(3,2)%
}
\NovelDisplay{M00739R05}{%
\xrightarrow{[1]}(0)(1,1,1)(2,1)(3,1)%
}
\NovelDisplay{M00739R06}{%
\xrightarrow{[1]}(0)(1,1,1)(2,1)(3)(4,1,1)%
}
\NovelDisplay{M00739R07}{%
\quad(5,1)%
}
\NovelDisplay{M00739R08}{%
\xrightarrow{[0]}(0)(1,1,1)(2,1)(3)(4,1,1)%
}
\NovelDisplay{M00739R09}{%
\xrightarrow{[0]}(0)(1,1,1)(2,1)(3)%
}

取到末尾 \allowbreak{}\NovelInline{I03868}{(3,1)}\allowbreak{} 的那一步，展开第一项会新添三列。所需的 \allowbreak{}\NovelInline{I03869}{(3)}\allowbreak{} 出现在新增部分的开头，其后还有两列。于是展开以后，须再取两次 \allowbreak{}\NovelInline{I03870}{[0]}\allowbreak{}，逐列删去尾端，不能把应当新添的三列写成只生成了一列。

霍雨浩把这两次零项写在同一段左侧。诺当沿列数再核一遍，确认每删一次，纸面上的末列确实改变一次。

奥斯卡端着水走到桌边，见几人还在数，便先把杯子放远了些。蔡兴坤抬头喝了一口，重新对准下一步的末列。

后面重复的 \allowbreak{}\NovelInline{I03871}{(2,1)}\allowbreak{} 展开第一项时，也会新添三列。它们照样先完整写出，再逐列截去。众人这一次没有停下来争论记法，只依照相同约定继续向下。

\NovelDisplay{M00740R01}{%
(0)(1,1,1)(2,1)(3)%
}
\NovelDisplay{M00740R02}{%
\xrightarrow{[1]}(0)(1,1,1)(2,1)(2,1)%
}
\NovelDisplay{M00740R03}{%
\xrightarrow{[1]}(0)(1,1,1)(2,1)(2)(3,1,1)%
}
\NovelDisplay{M00740R04}{%
\quad(4,1)%
}
\NovelDisplay{M00740R05}{%
\xrightarrow{[0]}(0)(1,1,1)(2,1)(2)(3,1,1)%
}
\NovelDisplay{M00740R06}{%
\xrightarrow{[0]}(0)(1,1,1)(2,1)(2)%
}
\NovelDisplay{M00740R07}{%
\xrightarrow{[1]}(0)(1,1,1)(2,1)(1,1,1)(2,1)%
}
\NovelDisplay{M00740R08}{%
\xrightarrow{[0]}(0)(1,1,1)(2,1)(1,1,1)%
}
\NovelDisplay{M00740R09}{%
\xrightarrow{[1]}(0)(1,1,1)(2,1)(1,1)(2,2,1)%
}
\NovelDisplay{M00740R10}{%
\quad(3,2)%
}
\NovelDisplay{M00740R11}{%
\xrightarrow{[1]}(0)(1,1,1)(2,1)(1,1)(2,2,1)%
}
\NovelDisplay{M00740R12}{%
\quad(3,1)(4,2,1)%
}
\NovelDisplay{M00740R13}{%
\xrightarrow{[1]}(0)(1,1,1)(2,1)(1,1)(2,2,1)%
}
\NovelDisplay{M00740R14}{%
\quad(3,1)(4,2)%
}

七列目标最终落到纸上。霍雨浩将最后一个编号补好，没有再把结果另外重抄一遍。

生成路径中的每一步都从当前矩阵出发，取一个明确的有限编号，得到下一张矩阵；需要截去末列的地方，也各自留下 \allowbreak{}\NovelInline{I03872}{[0]}\allowbreak{}。

小舞从最后一行往上核，诺当从第一行往下核，两人的手指在中间相遇，便交换了各自已经查过的部分。

卢本韦看向三条局部检查，又看向这一整条路径。

``前面的检查帮我们发现明显问题；这一张的标准性，有这条从边界取项下来的路径作依据。''

``是。''test\_alpha0 将笔放在目标旁边，``给出哪个结论，就把支撑它的依据留在哪里。''

霍雨浩没有把``会展开''\,``通过局部检查''\,``已给出生成路径''三个说法写成同一句。它们在纸上各有对应位置，他日后翻看，便能知道哪件事已经完成，哪件事还需要补证。

梅天狸将几张纸按今天的顺序叠好。第一页是最短的三行结构，第二页定坏根、算截断；接着是同一行不同位置的提升条件，随后是四行综合与一般规则，最后才是完整生成路径。

卢本韦没有立刻松开最后一张纸。

``大长老，我开始问的是两边接到一起以后，添新层次要不要改旧连接。''她将双魂核图重新拉到面前，``今天这些算式，还不能替我把两颗魂核做出来。''

test\_alpha0 笑了一下，示意她接着说。

``但我知道要先查什么了。先把原先各部分怎样连接弄清楚，再分清哪些位置能够随这一步调整，哪些要保留。父项变了，还要追查依赖它的后续位置。''

她指向第三张改动前后的两列。

``不能因为都排在一段里，就要求处处同步增加。也不能只盯着新添的那一段，漏掉它的关系已经接回前面。''

``这一步想清楚，接下来练才有依据。''大长老说，``魂核的体量、属性与供给怎样相济，还要另行打磨。今天的有限阵盘先在外面练稳，别把尚未查明的回路直接引进体内。''

他让众人将第三张阵盘复原成改动以前的样子，隐去父项连线，只保留数字与列号。

连线暗去以后，先前需要核对的位置仍清楚地留在众人本子上。

唐三稳住好部，诺当对齐列号，小舞接原坏部。蔡须坤与卢本韦各查一处第二行，蔡兴坤算出阶差以后等在旁边，直到霍雨浩点头，才展开第一项，将算好的新增部分送出。

卢本韦托到原第六列对应的位置时停了半息。

她没有因为两边的光都已亮起就硬往前推，而是先沿第五列第二行的父链接回第一列，确认此处保留二，再接到后面能提升的第七列。

阵盘平稳铺开，唐三掌中的前缀没有被拖动。

唐三慢慢收回托住阵盘的手，光仍按核好的路线走完一遍。奥斯卡这才将刚才留在远处的水递过来，卢本韦接杯子时，指尖还带着一点细光。

``这回停的地方有原因了。''她低头看了一眼自己的记录。

蔡须坤将她那张纸转正，在原第六列下方补了一道通向第五列的小线。

``下次忘了，从这里继续追。''

识海中的牢大也看完了那条父链，轻声道：``这回讲得很清楚。''

蔡须坤按着纸角的手指微微停了一下。她抿起嘴角，只在心里回了一句：``你的那份，我也记着。''

淡淡的魂影应了一声。她把这一页抚平，继续看向桌上的阵盘。

蔡兴坤掌心也留下一道细细的 BMS 序数道纹。她没有立即收进体内，而是依照自己刚写好的有限回路，让纹路在掌外重新亮了一遍。

test\_alpha0 托住她的手，等她将根、截断与提升条件都核过，才示意她缓缓收拢。

这道纹路所能承接的，仍是今日核定的有限结构。蔡兴坤将自己的演算纸折在它旁边，没有把掌心的一次稳定，当成今后每一次都无需检查。

小舞替大长老收起笔，看见他还在检查那条生成路径，便问：``您也要把这些父项一项一项看过去？''

``要。看得快一些，仍是在看。''

小舞笑了笑，将最后一张纸压平，露出差点被折痕挡住的取项编号。

霍雨浩等他核完，才将长纸折好。原先画在纸边的双魂核问题没有被挤走，旁边多了四个各有用处的例子，也多了几条能指出具体位置的回答。

窗外云影缓缓移过石阶。梅天狸合上墨碟，诺当将余下的空纸收进夹页，众人正要起身，殿外忽然传来一阵急促的脚步声。

test\_alpha0 搁下笔，望向殿门。

\chapter{第六十五章：碎紫来报，暗流涌动}\label{ux7b2cux516dux5341ux4e94ux7ae0ux788eux7d2bux6765ux62a5ux6697ux6d41ux6d8cux52a8}

黑板上那条标准 BMS 的定义还没擦。

找父列，找坏根，算阶差，往后抄，抄的时候哪一行跟着往上加、哪一行不动，一条一条写得清清楚楚。八个人坐在底下，盯着那几行字，谁也没出声。大长老刚才把最后一条讲完，粉笔往板槽里一放，殿里就静下来了。这条定义他们啃了小半个月，今天才算从头到尾看全了。

霍雨浩的本子摊在膝盖上，笔尖停在半空，停了很久才落下去，在那几行字底下画了一道横线。

奥斯卡嘴里叼着半根香肠，忘了嚼。

殿外这时候传来了脚步声。

脚步很急，一步赶着一步，石阶上咚咚地响。来人身上的魂力压不住，隔着殿门就涌进来一截，修为不低，可乱得厉害。

``禀报大长老！''

殿门让人一把推开了。

进来的是个穿紫袍的老人，须发花白，袍角上全是土。他三步并作两步走到殿中央，单膝跪下，膝盖砸在地砖上，咚的一声。

蔡须坤认得他。天斗帝国的护国斗罗，帝国大将军，碎紫斗罗。大赛那年颁奖，就是他把奖牌挂到她们脖子上的。那时候他满面红光，嗓门大得半个会场都听得见。

今天他跪在那儿，两只手撑着膝盖，额头上的汗一直往下淌。

``大长老，''他说，``天斗城出事了。''

test\_alpha0 坐在上首，没动。

``起来说。''

碎紫斗罗站起来，喘了两口气，才往下说。

``陛下中了毒。''他说，``昏过去三天了，一直没醒。太医院的人轮着看，看不出是什么毒，药灌下去，一点用都没有。那毒怪得很，不在血里，也不在经脉里，像是\ldots\ldots{}''

他停了一下，像是在找词。

``像是什么东西，在一点一点地改陛下身上的魂力。''他说，``太医院的老院正说，他行医五十年，没见过这样的。他说这不像毒，倒像是序列在往里头渗。''

殿里八个人同时抬起了头。

test\_alpha0 的目光动了一下，没有打断他。

``还有太子。''碎紫斗罗的声音压低了，``这阵子太子殿下不对劲。先是说城外闹匪，把禁军三大营调走了两营。后来又关起门来，一个一个地召见朝里的大臣，召见了谁，说了什么，末将一概不知。到了昨天\ldots\ldots{}''

他咬了一下牙。

``昨天太子调了禁军，把皇宫里里外外围了三层。说是护着陛下。''碎紫斗罗抬起头，``末将说不上来哪儿不对，可就是不对。这几天末将心里一直有个地方在疼，疼得睡不着，可要说是哪儿疼，又指不出来。''

``像是身上有一道魂力在隐隐作痛，''test\_alpha0 说，``可你找遍了，找不到伤口在哪儿。''

碎紫斗罗猛地抬起头来。

``大长老，您怎么\ldots\ldots{}''

test\_alpha0 没回答他。他从上首站起来，背着手，慢慢走到殿门口，往天斗城那个方向看了一会儿。

天斗城在山底下。从这儿看下去，城墙是灰的，屋脊是黑的，皇宫那一片琉璃瓦在太阳底下亮着，看着跟平常没什么两样。

他看了很久。

``真霜。''他说。

\begin{center}\rule{0.5\linewidth}{0.5pt}\end{center}

话音刚落，殿里就冷了。

冷是从殿后头来的。先是脚底下，地砖上起了一层白，一寸一寸往前爬；然后是手上，奥斯卡叼着的那半根香肠，香肠皮上结了一层细细的霜。小舞哈了一口气，那口气一出嘴就成了雾。

殿后的侧门里走出来一个人。

他一百多岁了，头发白里透着一点蓝，像冬天河面上结的那层冰，在脑后挽着，插着一根木簪。身上是一件洗得发白的旧袍子，袖口磨出了毛边。腰背挺得笔直，比站在殿中央的碎紫斗罗还直。他走得不快，每一步落下去，脚底下的地砖就结出一圈霜纹，霜纹往四边散开，散成一朵一朵的冰花。

扽魂殿的长老，真霜斗罗。

殿里的人都知道他。他在藏卷楼管了几十年的卷宗，一年到头不下楼，殿里年轻一辈见过他的不多。可殿里的人也都知道，他的境界停在 Y(1,3,6)，修的是极致之冰。冰这一个字，大陆上能修到他这一步的封号斗罗，十根指头数得过来。

他走到殿中央，朝上首抱了抱拳。

``徒儿在。''

test\_alpha0 没有回身，还是看着山底下的天斗城。

``带果糕小队，去天斗城。''他说，``把叛乱平了。陛下中毒的事，顺便查清楚。''

``是。''

``太子要是真反了，''test\_alpha0 说，``不必留手。''

真霜斗罗的眼皮抬了一下。

``徒儿领命。''

八个人坐在底下，你看看我，我看看你。事情来得太突然，可谁也没说不去。刚才黑板上那几行定义还在脑子里转，转着转着，好像把八个人的魂力都串在了一根线上，一个人动一下，别的人都觉得出来。

奥斯卡把冻硬了的半根香肠揣回了袋子里。

``走了。''蔡须坤站起来说。

八个人跟着站了起来。

``等等。''

test\_alpha0 转过身来了。

他走到殿后头那面墙跟前，在墙上按了一下。墙上开了一个暗格。他从暗格里取出来一样东西，捏在手里，走回来。

那东西拳头大小，圆的，通体漆黑，表面刻满了纹路，一圈套着一圈，密得看不清起头和收尾。纹路里有光在走，走得很慢，像是里头关着什么活物，在一下一下地喘气。

碎紫斗罗看见那东西，脸色一下子变了。

``八级定装魂导炮弹？''

test\_alpha0 没理他。他走到蔡须坤跟前，把那东西放进了她手里。

炮弹很沉，沉得蔡须坤的手往下坠了一坠。贴着掌心的那一面是温的。

炮弹出自二长老 Smile Lee 之手。八级的，殿里也没有几颗。

test\_alpha0 看着她手里那一颗，想起的是另一颗。九级的，Smile 头一回做成，做了七年。一百六十八年前，那一颗从天州城楼上打了出去，打的是二十里外星斗大森林边上那只翡翠天鹅。帝天一剑劈开了一大半，剩下那点余火，燎了碧姬一只翅膀。炮弹出膛的时候，整座城楼往下一沉，瓦片掉了一地。

这些，他一句都没提。

``这里头封着 0-Y 序列的功法。''test\_alpha0 说，``炸开以后，够把一个封号斗罗打成重伤。''

他看着八个人，一个一个地看过去。

``可你们给我记住。''他说，``用了它，用的人要虚好几天，功力掉一大截。所以不到万不得已，不许动它。''

八个人谁也没出声。八级定装魂导炮弹这几个字，他们在书上见过，在学院的讲台底下听老师说过，没想到有一天会揣在自己身上。炸一回，能把一场仗翻过来；炸完了，扔炮弹的人要躺好几天。那几天里要是再撞上硬茬子，就只能干等着挨打。

``当成最后一张底牌。''test\_alpha0 说，``局面收拾不住了，别的法子都用尽了，才能用。到时候炮弹一炸，里头的 0-Y 序列会顺着爆开的那一下灌进你们身子里，把你们一下子顶到头。可顶完了以后那几天怎么过，得你们自己想法子。''

``是。''蔡须坤说。

她把炮弹收进了贴身的储物魂导器里。收进去以后，她隔着衣服按了按那个地方。

奥斯卡凑过来，小声问：``沉不沉？''

``沉。''蔡须坤说。

``那我那四十根香肠，你帮我也揣一半。''

蔡须坤看了他一眼。奥斯卡缩了缩脖子，退回去了。

真霜斗罗转过身，朝上首深深一拜。

``大长老放心。''

test\_alpha0 点了一下头。

八个人跟在真霜斗罗和碎紫斗罗后头出了殿门。test\_alpha0 站在殿门口，看着他们一个一个走下石阶，走到最后一个，是卢本韦，两只手插在兜里。

他一直看到他们拐过山道，看不见了，才回殿里去。

\begin{center}\rule{0.5\linewidth}{0.5pt}\end{center}

下了山就是官道。

真霜斗罗走在最前头。他身上散出来一层薄薄的冰雾，不冷，罩在十个人身上，路边的风沙吹过来，撞上那层雾就落了地。碎紫斗罗走在他边上，八个人跟在后头。

走出去十来里，真霜斗罗开了口。

``陛下倒下去之前，''他问碎紫斗罗，``吃过什么，见过谁？''

``那天晚上，陛下跟太子一块儿用的膳。''碎紫斗罗说，``就他们父子两个，内侍都让退出去了。用完膳太子走了，陛下进了寝宫，第二天早上就没醒过来。''

``太子那天晚上，脸色怎么样？''

碎紫斗罗想了想。``好得很。''他说，``那阵子朝里事多，陛下愁得头发白了一片，太子倒是一天比一天精神。末将那时候还觉得，东宫到底是年轻。''

真霜斗罗没再问。

再往后，一路上没人说话。

官道两边是收过的麦地，麦茬子一排一排地立着，地头上堆着麦秸垛。往天斗城去的路上，本来这时候该是一车一车拉粮的，今天一辆都没碰见。倒是往外走的人不少，挑着担子的，推着独轮车的，车上堆着铺盖卷跟锅碗，一家老小跟在后头，闷着头赶路。

碎紫斗罗拦下一个推车的汉子，问他上哪儿去。

``出城躲躲。''汉子说，眼睛不敢往他身上的紫袍看，``城里头说是要搜人，挨家挨户地搜。俺家婆娘害怕。''

``搜什么人？''

``说是身上带着什么\ldots\ldots 序什么的。''汉子挠了挠头，``俺也听不明白。反正官爷们见着谁不顺眼，就往外拖。''

碎紫斗罗松了手，让他走了。他站在路边，看着那辆独轮车吱扭吱扭地走远，脸色沉得像要下雨。

卢本韦走在最后头，手揣在袖子里，一直在摸那块猪神令。摸着摸着，她小声跟边上的蔡兴坤说：``我爸说过，城里要是乱了，头一个跑的，就是种地的跟推车的。他们鼻子最灵。''

蔡兴坤看了她一眼，没接话。黑子鸡蹲在她肩膀上，脖子一伸一伸的，往天斗城那边看。

蔡须坤走在八个人最前头，一只手隔着衣服按着那颗炮弹。她在想碎紫斗罗那句话，像是序列在往里头渗。还有大长老站在殿门口看天斗城的那个样子。

陛下中毒，要真是序列在往里渗，那就不光是谁下了一包药的事了。太子造不造反，倒成了小事。

她的魂力在身子里动了一下。这一个月学 BMS，学得魂力都有了脾气，碰上不对劲的东西，它自己就先紧起来。就像推一条表达式推到半截，坏根一冒头，整条式子都跟着一沉。

``快点。''真霜斗罗在前头说，``天黑之前，要进城。''

他说完，脚底下一点，人就窜出去了。冰雾拖在他身后，在官道上拉出一道长长的白痕。碎紫斗罗跟着窜了出去。八个人咬着牙，一个跟着一个，把魂力催到脚底下，追着那道白痕往天斗城赶。

\begin{center}\rule{0.5\linewidth}{0.5pt}\end{center}

到天斗城外头的时候，太阳已经挨着西边的山尖了。

城墙底下，真霜斗罗和碎紫斗罗并肩站住。八个人在他们身后散开，各自站好位置，魂力都提到了七成。

``不对。''碎紫斗罗压着嗓子说，``上回我出城的时候，城门上四队人轮班。现在\ldots\ldots{}''

他没往下说，可谁都看见了。城墙上每隔三丈就站着一个甲士，手里攥着长戟。箭楼上站满了弓弩手，弩臂在夕阳底下泛着黑光。护城河里头又加了三道拒马，河滩上撒满了铁蒺藜。

这不像护着陛下，像是在守一座围城。

真霜斗罗没说话，接着往前走。他身上那层冰雾比刚才厚了一点，脚底下不紧不慢，一步是一步。

离城门还有百来丈的时候，城门洞里有人喝了一声。

``站住！''

一个披着重甲的将领从城门洞里大步走出来，身后跟着二十几个魂师，身上的魂力云都翻起来了，最低的也是魂王。将领手里横着一把长刀，眼睛在他们身上一个一个地扫过去，扫到真霜斗罗身上，停住了。

``陛下龙体抱恙，卧床不起。朝中已经下了旨，天斗城即刻戒严。''他一字一顿地说，``从现在起，没有太子殿下亲批的手令，谁都不许进城。''

他顿了一下。

``你们，也不许。''

碎紫斗罗的脸一下子涨红了。他往前跨了一步。

``本座是护国斗罗碎紫，奉旨出城办事，现在回城。你拦我？''

将领的脸上一点变化都没有。他把长刀往上抬了一寸。

``戒严令下了，不分高低。''他说，``碎紫斗罗要是不服，可以硬闯。末将拦不住，可末将的人会放箭。''

他话音刚落，城墙上几十架弩机一齐转了过来，弩口黑洞洞地对着城下。

八个人里有人的魂力云已经从身上冒出来了，是诺当。

真霜斗罗抬了一下手。

``退后。''

声音不大。八个人都往后退了半步。诺当身上那团魂力云收了回去。

真霜斗罗往前走，走到离那个将领十来丈的地方，站住了。他身上的魂力一点都没往外放，就那么站着，看着那个将领。

``你说戒严。''他开口了。声音不高，城墙上每一个人都听得清清楚楚，``我问你，戒严令是谁下的？''

将领皱了一下眉，像是没想到他问这个。

``自然是太子殿下\ldots\ldots{}''

``太子殿下拿什么名义下的？''

``自然是护着陛下\ldots\ldots{}''

``陛下昏了三天，太医院一点法子都没有。''真霜斗罗说，``一个昏过去的皇帝，怎么批的戒严令？''

将领的嘴唇动了两下，没说出话来。

``戒严令要是太子一个人下的，陛下没有朱批，''真霜斗罗说，``那这东西就叫兵变，不叫戒严。''

城墙上的甲士你看我，我看你。攥着弩机的手，一只一只都紧了。

那将领的额头上冒出了一层细汗。他咬着牙，嗓门提了起来。

``你血口喷人！太子殿下奉旨监国，一切都合\ldots\ldots{}''

``合什么？''

真霜斗罗这三个字说得很轻。

说完，城门口就冷了。

\begin{center}\rule{0.5\linewidth}{0.5pt}\end{center}

冷是从他脚底下起来的。

一圈白从他脚边往外扑，扑到哪儿，哪儿就冻上。地上的水洼一眨眼冻成了冰，路边的草冻成了一根一根绿色的冰棍，风里的土粒子冻成了细细的冰渣，簌簌地往下掉。那圈白扑过了城门洞，扑上了城墙，扑上了箭楼。

城墙上的弩机不转了。弓弦冻硬了，弩臂上的机括咔咔地响，就是扳不动。甲士们的盔甲上结了一层白霜，有人牙齿打架，打得咯咯响，有人腿一软，扶着墙垛子才站住。

那将领手里的长刀，当的一声掉在地上。他的手指头冻木了，攥不住了。

真霜斗罗背着手站着。他身后起了一面冰，薄薄的，亮亮的，映着他的背影，也映着西边快落下去的太阳。

他连扽魂都没放。

``我再问你一回。''他说，``戒严令，合不合规矩？''

那将领的嘴唇冻得发紫，浑身哆嗦，咬着牙，一个字都不说。他身后那二十几个魂师，一个个缩着脖子，谁也不敢出声。

碎紫斗罗走上来，看着那将领，看了一会儿，脸上的红慢慢退下去了。他跟着太子好些年了，太子是什么人他知道，做事果断，可从来不莽撞。今天这一出，后头怕是还有别的东西。

``真霜前辈，''他压低了嗓子，``这事别在城门口闹大。太子要真反了，咱们打进去就是。要是这里头有什么误会，在城门底下动了手，往后就收拾不住了。''

真霜斗罗的眼睛动了一下。他身上的冷收回去三分。

城门口那层白，一点一点地退了。城墙上的弩机又能转了，甲士们这才觉出手脚是自己的，一个个松了一口气，可再也没人敢把弩口对着城下了。

那将领一屁股坐在了地上，大口大口地喘气。他抬头看真霜斗罗，眼睛里全是怕。

``给你一炷香。''真霜斗罗说，``一炷香以后，太子的亲笔手令，太医院拿出来的陛下亲笔旨意，两样里拿不出一样\ldots\ldots{}''

他转过身去，背对着城门，看着八个人。

``我们就自己进去拿。''

蔡须坤站在八个人最前头，手心里出了一层汗。她身子里的魂力又紧了一下，比刚才在官道上紧得还厉害。这座城里头，有东西已经歪了，歪得很远。

她隔着衣服按了按那颗炮弹。炮弹贴着她的肋骨，在里头轻轻地震了一下。

\begin{center}\rule{0.5\linewidth}{0.5pt}\end{center}

一炷香没到，真霜斗罗就不等了。

他看了一眼坐在地上的那个将领。那人张着嘴，喘着气，一个字也没说出来。他身后那些甲士，一个个缩着脖子，刚才那股子横劲儿，跟着手指头上的冻一块儿化没了。刚才那一下他们都尝到了，封号斗罗是什么分量。弩机冻住了，弓弦冻住了，连身子里的魂力都冻得转不动，人就跟封在一块冰里一样。

那还只是真霜斗罗身上漏出来的那么一点。

真霜斗罗看着那些人，脸上没有火气，也没有瞧不起。

他活了一百多岁了。调虎离山，封锁消息，打着护着谁的名头把谁关起来，这一套他年轻的时候见过不下十回。回回都差不多，换一拨人，换一座城，再唱一遍。

他甚至觉得有点乏味。

``走。''

他没再看那将领一眼，抬脚往城门里走。走得不快，脚底下一步一个霜印，在青石板上排成一溜。

碎紫斗罗跟上去。八个人排成楔子，一个跟着一个进了城门洞。那将领张嘴想喊，抬头看见真霜斗罗的背影，伸出去的那只手抖了两下，又垂下来了。

``让\ldots\ldots 让他们进去。''他跟身边的副将说。

副将的脸色很难看。``将军，太子殿下的命令\ldots\ldots{}''

``你拦得住？''将领吼了一声，嗓子都劈了，``那是封号斗罗！你当你是谁？''

副将不吭声了。

\begin{center}\rule{0.5\linewidth}{0.5pt}\end{center}

进了城，比城外还静。

天斗城的正街又宽又直，两边铺子一家挨着一家。平日这个时辰，正是点灯的时候，酒楼里划拳，茶馆里说书，街上的人挤得走不动道。今天街上空荡荡的，铺子大多早早上了门板，有几扇窗户里透出一点光，刚透出来，就让人从里头拿布遮上了。

城门里头那面墙上贴着一张告示，纸是新的，糨糊还没干透。碎紫斗罗走过去，扯下来看了一眼，脸色就沉了。

告示上写着：雪星亲王勾结奸人，于陛下饮食中下毒，事败潜逃。太子奉旨监国，全城戒严。有知亲王下落者，报官重赏；藏匿者，同罪。

底下盖着东宫的大印。

``雪星亲王下毒？''碎紫斗罗把告示攥成了一团，``他是陛下的亲弟弟。''

没人接他的话。

再往里走，街两边那些关着的门板后头，有人在小声说话。他们走过去，说话声就停；走过去了，又响起来。断断续续的，飘进耳朵里几句。

``\ldots\ldots 说是太子还活着，好好的，在宫里主事\ldots\ldots{}''

``\ldots\ldots 是雪星亲王下的毒，亲眼有人看见的\ldots\ldots{}''

``\ldots\ldots 二皇子也跟着他跑了，昨夜里，从西门\ldots\ldots{}''

``嘘，别说了。''

门板后头又没声了。

霍雨浩一路听着，一路在心里记。这几句话，他在三条街上听了三遍，三遍都是这么说的，一个字都不差。街坊传闲话，传到第二家就走样了，传到第三家就面目全非。这几句话传了三条街，还是原样。

他把这件事记在了本子上，记完了，在底下画了一道线。

真霜斗罗脚底下没停，一道神念从他身上散开去，铺满了方圆好几里。

碎紫斗罗走在他边上，眼睛扫着两边关得严严实实的门窗，眉头越皱越紧。

``左边。''霍雨浩忽然开口了。

他的精神力早就铺出去了，这时候往一个方向收了收。

``左边那条巷子口，有一队兵在挨家挨户地搜。''他压着嗓子说，``五个人一组，每组里有一个魂师，手里拿着一个仪器，是魂导探测仪。''

唐三的蓝银草从他脚边钻进了石板缝里，顺着缝往那条巷子底下爬过去。过了一会儿，他点了点头。

``他们走得太齐了。''他说，``脚步一个点儿，跟一个人走的一样。那个探测仪，军中没有这个样式，是专门做的。''

``专门探序列的。''蔡须坤说。她身子里那股紧，从进了城门起就没松过，``他们找的人，身上有序列的东西。''

小舞皱着眉，看着那队兵刚搜过的几户人家。``那几户的门槛上，有东西留下来。很淡。''她说，``不像扽魂留的。''

真霜斗罗没开口。可他的神念比谁铺得都远。霍雨浩说的那条巷子里，那几户搜过的人家里，他也闻到了那一丝东西。

序列的气味。

淡得很，细得跟一根蛛丝似的。不是他这样的人，根本闻不出来。可它确实在那儿，在那几户人家里，在那些关严了的门窗后头。

``他们在找人。''他开口了。

碎紫斗罗一愣。``找谁？''

真霜斗罗往皇宫那边看过去。皇宫那一片灯火通明，禁军的火把把半边天都照红了。

``找一个知道得太多的人。''他说。

他话音还没落，奥斯卡的脚步就停了。

他的鼻子动了两下。

奥斯卡的扽魂是香肠，他这个人鼻子也灵，灵得出了名。学院食堂今天做什么菜，他在宿舍里闻一闻就知道；谁兜里揣了块糖，他隔着三张桌子都闻得出来。这时候风从右边的巷子里吹过来，他闻到了一点味道。

血腥气，很重。底下压着苦味，是药。再底下，还有一丝很淡的魂力的气味，淡得快散了，可那是高阶魂师身上才有的。

``有情况。''他抬手拦住了身边的人，嗓子压得很低，``右边巷子里，有血。很重，是内伤出的血，五脏六腑让什么撞过了，血从里头往外渗。还有人在跑，跑得快，可脚底下不稳，左腿伤了，落脚比右脚轻三成。''

几个人都愣了一下。

蔡兴坤挑了一下眉。``隔这么远，你闻得出来？''

``闻得出来。''奥斯卡点头，``那股魂力的味儿，是皇家的。我在宫里待过一阵子，错不了。''

真霜斗罗转过来看他。

``哪边？''

``右边第三条巷子，往里四百来步。''奥斯卡说，``在往西跑。跑得越来越快，可伤拖着他，撑不了多久。''

碎紫斗罗的脸色变了。

他一句话都没说，身子一晃，就窜了出去。紫光在他身上一闪，人已经到了那条巷子口，拐进去不见了。

``跟上。''真霜斗罗说。

八个人跟着冲进了巷子。奥斯卡跑在中间，一边跑一边耸着鼻子，嘴里一直没停。

``往左拐\ldots\ldots 直走五十步\ldots\ldots 又拐了，进了前头那座宅子的后门\ldots\ldots{}''

巷子很窄，两边是高墙，墙头上长着草。脚步声在巷子里撞来撞去，听着像有好几十个人在跑。

\begin{center}\rule{0.5\linewidth}{0.5pt}\end{center}

碎紫斗罗先到的。

他落在一座宅子后头的窄巷里，一眼就看见了一个人影。那人正扒着墙往上翻，左胳膊垂在身边，使不上劲，右腿也是瘸的，翻了两回没翻上去。身上的锦袍撕得一条一条的，露出底下贴身的软甲，软甲上全是刀口和烧焦的印子。

碎紫斗罗站在那儿，没动。

他认出了这个人。

几乎就在同时，巷子那头响起了一片脚步声，甲叶子哗哗地响。一队禁军冲进了巷子，领头的军官手里举着一个探测仪，仪器上的指针乱跳，嗡嗡地叫。

``在那儿！别让他跑了！''

碎紫斗罗一步跨过去，站在了那个翻墙的人和那队禁军中间。

他背着手，身上的紫光漫出来，把整个巷口堵得严严实实。冲在最前头的十二个禁军一齐刹住了脚，像是一头撞上了一堵看不见的墙。

``什么人拦路？''那军官喝了一声，长刀横起来，身上的魂力云腾了起来，``我等奉太子殿下的命，捉拿要犯。让开，不然\ldots\ldots{}''

``不然怎么样？''

碎紫斗罗的声音不大，在窄巷里撞了一圈，嗡嗡地响。他慢慢抬起右手，指尖上绕着一缕紫光。就这么一抬手，那队禁军身上的魂力全都滞了一下。

``本座碎紫斗罗，帝国大将军。''

那军官的脸一下子白了，刀差点脱手。封号斗罗，还是大将军。他身后那十二个人，齐刷刷往后退了一步，身上的魂力云都淡了下去。

``你们奉谁的命，在城里满街地搜一位亲王？''碎紫斗罗盯着他，``太子殿下？一个连自己死活都做不了主的太子，也配调禁军满城追一位亲王？''

那军官的脸铁青，嘴唇动了好几回，一个字也说不出来。

``收起刀，带上你的人，滚。''碎紫斗罗说，``再让本座看见你们在城里闹事，本座不讲情面。''

那军官咬着牙，看了碎紫斗罗两眼，到底没敢动手。他铁青着脸，一挥手，带着那十二个人退出了巷子。脚步声越来越远，拐了个弯，听不见了。

巷子里静了。那个探测仪的嗡嗡声跟着那队人走远了，走到巷子尽头还在叫，叫得一声比一声尖，指针一直朝着这边。

碎紫斗罗听着那声音远了，才转过身来。

那个翻墙的人已经从墙上滑了下来，背靠着墙根坐着，大口大口地喘气。他的头发散下来，遮住了半张脸。

后头的脚步声到了。真霜斗罗从巷子口走进来，八个人跟在他身后，奥斯卡跑在最前头，鼻子还在耸。

碎紫斗罗站在那个人跟前，低头看着他，眉头拧成了一个疙瘩。

那人抬起头来，把挡在脸上的头发拨开了。

碎紫斗罗往后退了半步。

``王\ldots\ldots 王爷？！''

\begin{center}\rule{0.5\linewidth}{0.5pt}\end{center}

\chapter{第六十六章：进宫救帝}\label{ux7b2cux516dux5341ux516dux7ae0ux8fdbux5babux6551ux5e1d}

靠在墙根底下的，是雪星亲王。

陛下的亲弟弟，天斗帝国皇位的第二顺位。宫里的人管他叫二皇叔，碎紫斗罗这些年见他，都喊王爷。

这时候他半坐半躺地靠着墙，嘴角挂着一道干了的血，头发散乱，遮住了半张脸。锦袍上全是泥，一只袖子扯掉了半截。可他抬起头来的那一下，那双眼睛碎紫斗罗一眼就认出来了。眼睛里有血丝，有火，就是没有怕。

真霜斗罗和八个人这时候也到了。

奥斯卡第一个开口，嗓门压不住，带着点得意：``我就说那股味儿错不了，是皇家的人。''

蔡须坤没接他的话。她蹲下来，隔着一臂远，看着雪星亲王身上那几处伤。她身子里的魂力动了一下，往那几处伤上探了探。

``他身上的伤，一半是刀伤。''她说，``另一半是反噬。他往谁身上渡过气，渡到一半，让那边的东西顶回来了，伤了经脉。''

霍雨浩的精神力从亲王身上过了一遍，点了点头。

``左胳膊的经脉断了好几截，右腿骨头裂了。''他说，``身子里的魂力是乱的，有三处经脉让什么东西啃过。先受的内伤，跑的时候又让人追着砍，伤上加伤。''

``先别急着下结论。''卢本韦开口了。她的手在袖子里，指头摸着一块东西，是她家传的猪神令，平时揣在身上，有事没事拿手指头摸一摸上头的纹路，``听他自己说。''

雪星亲王抬起头来，脸上全是土和血。他看了看碎紫斗罗，又看了看真霜斗罗，再看看蹲在跟前的蔡须坤和她身后那七个半大孩子。他的眼神一个一个地扫过去，扫得很慢，带着戒备，也带着疑心。

``你们\ldots\ldots{}''他的嗓子哑得像砂纸，``是来拿我的？''

``来查清楚的。''真霜斗罗说，``你说。''

雪星亲王笑了一下。那一下笑得很难看，扯动了嘴角的伤，又渗出血来。他用那只还能动的胳膊撑着墙，把身子坐直了一点。

``今天夜里，''他说，``有人要杀我。''

``从头说。''

雪星亲王吸了一口气，胸口起伏了两下，才往下说。

``子时前后，我府里来了一个人，说陛下不行了，让我赶紧进宫。''他说，``那人穿着内侍的衣裳，脸生得很，我府上的人说他是从后门进来的，走也是从后门走的。我那时候急，没顾上多想。我进了宫，到了寝宫。太子倒在龙榻边上，七窍都在流血，已经没气了。''

巷子里没人出声。

``我想救他。''雪星亲王说，``我把魂力往他身上渡，渡进去，就没了。泥牛入海一样，一点回响都没有。我再渡，还是没了。就跟他身子里有一张嘴，张着，我渡多少，它吃多少。''

``然后禁军就进来了。''

他停了一下。

``寝宫里就我一个人，站在太子身边，手还按在他胸口上。''他说，``他们一进来，就认定是我。我说什么都没用。我只能跑。''

``等一下。''唐三开口了，``你说太子死了？''

``我亲眼看见的。''雪星亲王说，``七窍流血，一口气都没有了。那毒厉害得很，我的魂力一碰他，就让什么东西抽走了。那不像是毒，倒像是\ldots\ldots{}''

``序列侵蚀。''真霜斗罗说。

雪星亲王猛地抬起头来。``你\ldots\ldots 你知道？''

真霜斗罗没回答。他慢慢蹲下来，跟亲王平视，一双眼睛像两面冰镜子，看着亲王的脸，看了很久。

``你身上一样药都没带。''他说，``所以下毒的不是你。''

碎紫斗罗一愣。``前辈，身上没带药，那不正说明他\ldots\ldots{}''

``真下毒的人，会不给自己留解药？''真霜斗罗扭头看了他一眼，``你觉得一个亲王，会蠢到杀了人不带解药，不留退路，站在死人边上等着让人抓？''

碎紫斗罗张了张嘴，没话说了。

诺当在一边开了口。``还有他身上的伤。''他说，``队长和雨浩都看过了，他是渡气渡伤的。他要是下毒的人，渡什么气？他这是在跟一样他不认得的东西顶着，没在演。''

小舞也点了点头。``他跑的路也不对。''她说，``真是他下的毒，头一件事就该往城外跑。他一直在城里兜圈子，像在找什么。''

雪星亲王转过头去看小舞，看了好一会儿。他的眼眶红了一下，又压下去了。

``我在找药。''他说，``太子中的那种毒，我在一本老书上见过。要解，得三样药，其中一样，城里有一家药铺有。我跑出来以后，一直在找那家铺子。''

``哪家？''卢本韦问。

``城西，杏林堂。''雪星亲王说，``掌柜的是我几十年的老朋友，他手上有一味药，叫玄冰玉露。我能想到的，只有这一样，兴许压得住那东西。''

\begin{center}\rule{0.5\linewidth}{0.5pt}\end{center}

真霜斗罗站起来，往城西那边看了一眼。

``太子死了。陛下昏着。王爷让人栽成了凶手。禁军把全城都封了。''他一样一样地数，``这几件事，指的是同一个方向。''

他转过身来，看着八个人。

``在背后动手的，太子算不上，王爷也算不上。''他说，``动手的，是替太子写告示的那个人。''

蔡兴坤皱着眉头。``长老，城门里那张告示，底下有太子的亲笔画押\ldots\ldots{}''

``假的。''真霜斗罗说，语气平平的，像是在说一碗菜咸了，``那几个字我看了一眼。一笔下去，笔锋是稳的，写字的人心跳也是稳的，喘气也是匀的。一个七窍流血、快断气的人，写不出那样的字。''

``那是别人替他写的。''

他的目光沉下去一点。

``能仿太子的字，能调禁军封城，能在陛下中毒的同时把这一整套都摆好。''他说，``这不是一个人干得了的。''

碎紫斗罗的脸色变了。``前辈的意思是\ldots\ldots{}''

``背后有人在指点。''真霜斗罗说，``而且那个人，对序列懂得很深。陛下的毒，太子的死，满城搜身上带序列的人，这三件事是一根线上串着的三颗珠子。串这根线的那只手\ldots\ldots{}''

他没往下说。

夜风从巷子口灌进来，吹得巷子里一扇破门帘子扑扑地响。雪星亲王靠着墙，浑身在抖，分不清是伤疼，还是冷。

远处，皇宫那一片还是灯火通明，火把的红光映在夜空上，一明一暗的。

真霜斗罗背着手站着，往皇宫那边看。他脚底下的霜无声无息地往外漫，漫了一圈，把巷子里的青石板都盖白了。

``碎紫。''

``在。''

``送王爷回府，多派人手守着。''真霜斗罗说，``送到了，你再来找我。他们八个，跟我进宫。''

``是。''

真霜斗罗走过去，伸出两根手指，在雪星亲王的左胳膊上点了一下，又在右腿上点了一下。点过的地方结了一层薄冰，冰底下的血不往外渗了。亲王咬着牙哼了一声，额头上的汗一下子冒了出来。

``断了的经脉我接不上。''真霜斗罗说，``冻住了，撑到天亮不出事。''

碎紫斗罗把手指头放进嘴里，打了一声呼哨。巷子口的房顶上跳下来两个穿便装的汉子，是他的亲兵，一直远远地跟着。两个人一左一右，把雪星亲王架了起来。

``王爷，得罪了。''

雪星亲王让他们架着，走了两步，回过头来，看着真霜斗罗。

``真霜长老。''他说，``陛下那边\ldots\ldots{}''

``我去看。''真霜斗罗说。

亲王点了点头，让人架着走了。碎紫斗罗跟在后头，走到巷子口，回头看了他们一眼，紫光一闪，人不见了。

真霜斗罗抬脚往巷子外头走，霜跟着他走。

八个人跟在他身后，排成了一个楔子。蔡须坤走在最前头，她身子里的瑞兽本源转得很稳，稳稳地铺开一层，护着身后的人。蔡兴坤走在她边上，眼睛不停地往两边扫。唐三和霍雨浩一左一右，蓝银草从唐三脚边钻进石缝，一路往前爬，霍雨浩的精神力铺出去，两个人一个摸地底下，一个摸地面上，把四周百来丈罩得严严实实。小舞和诺当护着两边后头。奥斯卡走在当中，随时准备往谁嘴里塞香肠。卢本韦走在最后，手还在袖子里，摸着那块猪神令。

八个人走得很齐，一声不吭，走进了天斗城的夜色里。

皇宫那片灯火，在前头等着他们。

\begin{center}\rule{0.5\linewidth}{0.5pt}\end{center}

城门口，那个将领还坐在地上。

真霜斗罗那一行人的背影早就看不见了，长街尽头黑黢黢的，冰霜的寒气一点一点地散了。城墙上的甲士们搓着手，跺着脚，谁也不说话。

那将领一直坐着，一动不动，像城墙根底下一尊泥胎。额头上的冷汗顺着下巴往下滴，在青石板上洇出一小片湿印子。

等到最后那一丝寒气也散了，那股压得人喘不过气的封号斗罗的分量也一点都觉不出来了，他才慢慢地、很慢地吐出了一口气。

那口气很长。

``\ldots\ldots 走了。''他说。

他身边的副将凑了过来，声音还在打颤：``将\ldots\ldots 将军\ldots\ldots{}''

``闭嘴。''

将领的声音一下子冷了。

他站起来，拍了拍甲上的土，不慌不忙。刚才瘫在地上那副样子，跟着这一拍，全拍没了。他往四下看了一圈：长街上没人，真霜斗罗那一行人的气息一点都没剩下，窗户后头那些偷看的百姓，也早就缩回去了。

他的右手伸进了甲胄里头的暗袋。

暗袋里有一样东西，冰凉的，拇指盖那么大，是个小圆盘。漆黑，摸上去不像铁，也不像玉，表面光溜溜的，一道纹都没有。圆盘外头裹着一层隔魂力的东西，探测仪扫过去，只当是一片甲叶。

他的指头在圆盘边上按了一下。

咔。

很轻的一声，轻得只有他自己听得见。圆盘当中亮起了一点光，灰白色的，冷的，一闪就没了。

一道波动从圆盘里出去了。看不见，摸不着，没有魂力的味道，穿过城墙，穿过天斗城一层一层的屋脊，直直地往皇宫那边去了。

将领把圆盘塞回暗袋，理了理甲，脸上又换回了那副惶惶恐恐的样子。他转过身，跟副将说话。

``传下去，今夜的事，谁也不许往外说。那几个人进城的事，不用报上去。''

副将愣了一下。``可是将军，太子殿下那边\ldots\ldots{}''

``太子殿下自有安排。''将领的嘴角往上动了一下，那一下在火把底下看着发冷，``咱们守好城门就是。别的事\ldots\ldots{}''

他没往下说，往皇宫那边看了一眼。

\begin{center}\rule{0.5\linewidth}{0.5pt}\end{center}

真霜斗罗带着八个人穿过长街。越往皇宫走，街上越静，静得只剩下他们自己的脚步声，还有远处巡夜的甲叶子哗哗响。

街两边的铺子全上了门板。门板上让人拿白灰画了记号，有的画一个圈，有的画一个叉。画叉的那几家，门板让人踹开过，又拿木条胡乱钉上了，门缝里黑洞洞的，没有灯。

一家药铺门口，招牌让人摘下来扔在了地上，摔成了两截。台阶上撒了一地的药材，当归，黄芪，还有些叫不上名字的根根草草，让人踩烂了，混在泥里。

唐三蹲下来，捡起一截干巴巴的根子，在指头上捻了捻。

``搜过了。''他说，``搜得很急，柜子全拉开了，抽屉倒在地上。''

``搜什么？''小舞问。

``杏林堂是城西的。''霍雨浩说，``这一家不是。可要找一味药，挨家搜，最省事。''

几个人互相看了一眼，没再说话。

有人比王爷先想到了那味玄冰玉露。

真霜斗罗没停步。他走过那家药铺门口的时候，往门里看了一眼，脚底下的霜往门槛里头漫了一下，又收回来了。

走到一半，后头一阵风，碎紫斗罗赶上来了。

``王爷送到了。''他说，``府里留了我二十个亲兵。''

真霜斗罗点了一下头。

皇宫外头还有一道护城河，河上架着五座白石桥，当中那一座只有皇帝的车驾走得。往常夜里桥头各站着两个禁军，提着灯笼。今天五座桥头都空着，灯笼挂在桥栏上，风一吹，一晃一晃的，没人管。

真霜斗罗带着他们走了当中那一座。脚底下的霜一路铺过去，把白石桥面铺得更白了。

碎紫斗罗走在他身侧，落后半步。他在禁军里待了几十年，从百夫长一路干到大将军，皇宫哪道门归哪个营守，几时换班，巡夜的一队几个人，他闭着眼睛都背得出来。

皇宫正门到了。

正门外头是一片广场，青石铺地，方方正正，平日里百官上朝，就在这儿下轿下马。广场两边各立着一排石灯，灯里的火都点着，照得广场亮堂堂的，连石缝里的青苔都看得见。

奥斯卡走到广场边上，鼻子又动了两下。

``龙涎香。''他小声说，``宫门里头飘出来的。新点的，烧得很旺。''

``宫里点龙涎香，有什么稀奇。''蔡兴坤说。

``香是对的。''奥斯卡说，``可烧得太旺了。像是谁怕身上别的味儿盖不住。''

没人接他的话。

正门外头站着两排甲士，手里拿着长戟，火把照着他们的脸，一明一暗的。甲是禁军的甲，戟是禁军的戟，站法也是禁军的站法。

碎紫斗罗的脚步顿了一下。

他只看了两眼，就看出了三处不对。

头一处是脸。他在禁军几十年，宫门上当过差的将领，他起码都见过一面。眼前这二十几个人，没有一张脸是他见过的。

第二处是站相。禁军天天操练巡夜，兵油子气重，站得再齐，眼神里也带着几分懈怠。这二十几个人，一个个眼睛沉着，气息收着，站在那儿像二十几把没出鞘的刀。这种站相，是见过血的老兵才有的。

第三处，是修为。

碎紫斗罗不动声色地从他们身上一个一个地扫过去。最弱的那个，也稳稳在 BO 境界。站在最前头的两个领队，气息又沉又厚，往外压的那股劲儿，已经快摸到封号斗罗的门槛了。

二十几个 BO 往上的人，守着一座宫门。

碎紫斗罗袖子里的手指头攥紧了。他侧过头，嗓子压得很低：``前辈，宫门的人换了。不是禁军的人。''

真霜斗罗没回头，点了一下头。他早就看出来了。他看着那两排甲士，脚底下的霜往前又漫了几分。

``接着走。''

\begin{center}\rule{0.5\linewidth}{0.5pt}\end{center}

八个人也觉出不对来了。

蔡须坤走在八个人最前头。她的瑞兽本源这时候转得比平时快，感知也比平时灵。她的目光从那两排甲士身上扫过去，眉头皱了一下。那些人身上的魂力压着，可底子厚得很，禁军里没有这样的兵。

她把这件事记下了，没出声。

就在队伍又往前走了几步的时候，她停住了。

后脖颈上一阵凉，从脖子一直凉到腰。

刚才那一下，她的瑞兽本源抓住了一样东西。不在前头，也不在两边，在后头。

一道光。

很快，很隐蔽。从他们来的那个方向射过来，从队伍头顶上掠过去，直直地往皇宫里去了。没有魂力的味道，没有扽魂的气息，什么魂技都不像。就像一根看不见的丝线，在夜空里一闪就没了。

她猛地回头。

长街空荡荡的，两边的房子黑黢黢的。没有魂力，没有气息，什么都没有。

``怎么了？''蔡兴坤在她边上小声问。

蔡须坤没马上回答。她盯着来时的方向看了三口气，确认后头没人跟着，才转回来，脸沉着。

``刚才有一道光。''她的声音压得很低，只有蔡兴坤和前头的唐三、霍雨浩听得见，``从咱们后头来的，从头顶上过去，进宫了。很快，没有魂力，不是魂技。''

唐三的蓝银草往后头铺了出去。过了一会儿，他摇了摇头。``什么都没有。没留下魂力，也没留下气味。''

霍雨浩的精神力往后头罩了几百丈，也摇头。``我也没探到。干净得过头了。''

``就是因为干净，才不对。''蔡须坤说，``能把痕迹抹得这么干净的光，寻常传信的法子做不到。那东西，咱们没见过。''

她没把后半句说出来。她想说的是，序列。

卢本韦走在最后头，这时候往前凑了半步。

``从城门那边来的？''

``差不多是那个方向。''蔡须坤说。

卢本韦没再问。她回头往城门那边看了一眼，长街黑黢黢的，什么也看不见。她想起那个瘫坐在地上的将领，想起他抬头看真霜长老的那一眼。那一眼里全是怕，怕得太足了，足得像是装出来的。

几个人互相看了一眼，都明白了。有人在盯着他们。盯着他们的那个人，手里的东西他们摸不着，也追不上。

``都听见了？''真霜斗罗的声音从前头传过来，平平的，像是早就知道，``不用回头。该来的，总会来。''

他脚底下一点都没慢，霜在他脚边漫开，一步一步往宫门走去。

八个人收了心，重新排好队，跟了上去。蔡须坤还走在最前头，可她分了一缕感知，一直往后头探着。那道光扎在她心里，拔不出来。

\begin{center}\rule{0.5\linewidth}{0.5pt}\end{center}

宫门越来越近。

为首的一个甲士往前跨了一步，长戟一横，拦住了路。他的声音又平又硬，一点起伏都没有。

``来者何人？宫禁重地，没有许可，谁都不许进。''

真霜斗罗站住了。他看了那甲士一眼，又看了看他身后那二十几个人。

他没报名号，也没掏令牌，只是把右手抬了起来。

冷从他身上漫出去。宫门前的石阶上，一层白霜从他脚下往前铺，铺到那甲士的脚尖跟前。

``老夫真霜。''他说，``让开。''

火把在夜风里呼呼地响。

那甲士没让。

他横着长戟，戟尖离真霜斗罗的胸口只有三尺，眼睛平平地看着前头，不怕，也不犹豫。他身后那二十几个人一齐往前踏了半步，长戟一齐举起来，魂力从他们身上一齐涌出来，在宫门前结成了一堵看不见的墙。

``陛下龙体欠安，在寝宫静养。''那甲士说，像是在背一道早就背熟了的诏书，``太子殿下有令：谁也不许进宫觐见。违令者，格杀勿论。''

宫门前一下子静了。

碎紫斗罗的脸沉了下去。他往前跨了一步，紫光从他身上涌出来，大将军的威压往前一推。

``太子殿下有令？本座是帝国大将军，进宫面圣，用得着太子准许？''

``大将军恕罪。''那甲士的声音还是那样，不软，也不硬，``太子殿下有口谕：陛下病重这些日子，一切军政都由太子代行。进宫的人，要有太子的手令。没有手令，就算是封号斗罗，也不能进。''

``好大的胆子。''碎紫斗罗气得笑了一声。他身上的紫光猛地一涨，威压像山洪一样朝那甲士压了过去，``一个看门的，敢拦本座？''

那甲士的身子晃了一下，脚底下却跟生了根一样，一步都没退。他身后那二十几个人一齐催动魂力，二十几道魂力拧成一张网，硬是把碎紫斗罗那一压给接住了。

碎紫斗罗的瞳孔缩了一下。

霍雨浩站在八个人中间，眼睛一直盯着那二十几个人的胸口。他的精神力贴着他们的甲胄过了一遍，眉头慢慢皱起来了。

``他们喘气是一个点儿。''他小声说，``二十几个人，吸一口，呼一口，一点不差。''

``练出来的？''诺当问。

``练不出来。''霍雨浩说，``心跳也是一个点儿。''

诺当没再问。

一群 BO 境界的守卫，挨了封号斗罗一压，脸色都不变，硬接了下来。寻常的精锐，做不到。

``大将军。''那甲士开口了，声音还是硬邦邦的，``末将只是奉命行事。太子殿下有令在先，末将不敢违抗。大将军要是执意要闯\ldots\ldots{}''

他把长戟往下一压，戟尖对准了碎紫斗罗的喉咙。

``末将只有一死。''

戟尖离碎紫斗罗的喉咙不到一尺。碎紫斗罗站着没动，紫光在他身上一涨一落，涨了三回，落了三回。他在禁军里待了几十年，拿戟尖对着他喉咙的，这是头一个。

\begin{center}\rule{0.5\linewidth}{0.5pt}\end{center}

八个人的脸色同时变了。

蔡须坤往前走了一步。她身子里的瑞兽本源转得飞快，她看着那甲士，又看着他身后那二十几个人，脑子里像让什么劈了一下，刚才那一个一个零零碎碎的地方，一下子连成了一串。

宫门口这些生面孔，个个是 BO 境界。领头的挨了碎紫斗罗一压，一步都没退，这份定力，寻常的魂尊没有。他嘴里那道太子口谕，一个字一个字的，说得那么顺，像是早就备好的词。还有\ldots\ldots{}

太子要是真活着，他怎么不出来？

她的心往下一沉。

她又想起进城时门板后头那几句闲话。太子还活着。是雪星亲王下的毒。二皇子也跟着跑了。三条街，三遍，一个字都不差。那是有人在往外放，街坊之间传不成这样。

现在，这些人又打着太子的名号守着宫门，不让真霜长老和碎紫将军进去。

太子要真活着，真在宫里主事，他凭什么不让两个封号斗罗进宫见陛下？

除非他不敢。

除非坐在那个位子上的，根本就不是太子。

蔡须坤后背一阵发凉。她转过头，跟身后的唐三、霍雨浩对上了眼睛。两个人都看着她，眼神里是一样的意思。

太子已经死了。眼前这一出，是一场戏。

``长老。''蔡须坤快步走到真霜斗罗身边，瑞兽本源在她周身铺开一层，把声音拢住，只有真霜斗罗和碎紫斗罗听得见，``这些人不是禁军。看修为，看站相，看他们怎么配合，像是谁私底下养的兵。他们嘴里那道太子口谕，多半是假的。''

真霜斗罗点了一下头，脸上没什么变化。

``我觉得，太子殿下多半已经\ldots\ldots{}''蔡须坤的声音更低了，``不然，他们没有道理拦咱们。''

``本座也是这么想的。''碎紫斗罗在另一边说，紫光在他身上一明一暗，``这些人修为不过魂尊，可配合得太熟了，不是三两天练得出来的。禁军里头，本座从来没见过这样的编制。''

三个人互相看了一眼，心里都有了数。

真霜斗罗抬起眼睛，往宫门里头看。他脚底下的霜已经铺满了整片石阶，宫门前冷得滴水成冰。

``太子口谕？''他的声音平平的，``老夫倒要看看，这位太子殿下，藏在哪儿。''

\begin{center}\rule{0.5\linewidth}{0.5pt}\end{center}

他话音刚落，宫门里头响起了脚步声。

那脚步声不紧不慢，一步一步，每一步都落在差不多的地方，又稳又从容，像是一个闲人在自家园子里散步。脚步声越来越近，跟着脚步声一块儿飘出来的，还有一股龙涎香的味道。

宫门两边，几乎同时响起了一片甲叶子撞在一起的声音。

八个人扭头去看，瞳孔一齐缩了一下。

宫门两边的宫墙上，不知道什么时候站满了人。他们穿着深灰色的重甲，手里拿着长戟，从墙头上、廊柱后头、飞檐底下一个接一个地出来，像潮水一样涌到宫门前的广场上。十个人一排，二十个人一列，排得整整齐齐，脚底下一点声音都没有。一眨眼的工夫，宫门前就立起了两堵黑压压的人墙。

一百\ldots\ldots 两百\ldots\ldots{}

整整两百个甲士。

蔡须坤的瑞兽本源猛地一颤。她的感知从那两百个人身上扫过去，心往下一坠。

每一个都是 BO 境界。

两百个 BO 境界的人，摆在一座宫门前头。这不是一支守宫门的兵。

可让蔡须坤后背发紧的，是站在这两百个人最前头的那一个。

宫门正中的阴影里，走出来一个人。

他穿着一身黑甲，跟别人的甲不一样，肩甲上刻着一头张着嘴咆哮的兽，胸甲正中嵌着一枚暗金色的徽章，徽章上的纹路又细又密，看着像一张阵图。他脸上罩着一副覆面盔，只露出一双眼睛。

蔡须坤看着那双眼睛。

那双眼睛里什么都没有。没有杀气，没有战意，连活人该有的那点神气都没有。只有一样东西，冷的，齐整的，一格一格的，像一台机括在里头咬着齿轮转。

那黑甲将领往前踏了一步。

就这一步，真霜斗罗的眉头猛地一皱。

他脚底下铺出去的那片霜，往回缩了三分。

没人打它，也没人压它。是它自己缩回去的，像一头野兽碰上了比自己大得多的东西，还没想明白，身子就先往后退了。

真霜斗罗活了一百多年，自己的冰头一回这样。

碎紫斗罗也变了脸色。那黑甲将领踏出那一步的时候，他觉出了一样东西。那东西深得看不见底，大得望不到边，可又收得严严实实，一丝都没往外漏。要不是他对魂力的感应比常人灵，根本觉不出来。

``这个人的修为\ldots\ldots{}''碎紫斗罗的声音压得极低，只有身边的真霜斗罗听得见，``看不到底。''

真霜斗罗没接话。他看着那副覆面盔，看着盔后头那双一格一格转着的眼睛，脚底下的霜又往外铺了一寸，铺到那黑甲将领的影子边上，停住了。

他心里有了一个名字。

他没说出来。

\begin{center}\rule{0.5\linewidth}{0.5pt}\end{center}

\chapter{第六十七章：假面太子}\label{ux7b2cux516dux5341ux4e03ux7ae0ux5047ux9762ux592aux5b50}

黑甲将领站定了，抬起了右手。

他抬得很慢，慢得像在水里打拳，每一节指头怎么弯，胳膊上每一块肉怎么动，都看得清清楚楚。可就是这么慢慢地一抬，宫门前所有人身上的魂力都跟着顿了一下。

嗡。

一声闷响从地底下传上来，低得人胸口发麻。宫门前的青石地面抖了一抖，从黑甲将领脚底下裂开了一道一道的细纹，像一张蜘蛛网，眨眼间铺满了整片广场。

他身后那两百个甲士，一齐催动了魂力。

两百道光柱从两百个人身上冲起来，冲到半空里，撞在一块儿，嗡嗡地响，震得人耳朵眼发疼。那股魂力往四周推开，推得广场边上那两排石灯里的火苗一齐往外倒，又把地上的裂纹碾得更细了。

可那两百道光柱冲上去以后，没往天上散。它们拐了个弯，一道一道地，全往黑甲将领身上落。

两百道光，落进一个人身上。

宫门上那几排铜钉嗡嗡地响起来，一颗挨着一颗地抖。广场边上石灯里的火苗，本来一齐往外倒，这时候又一齐拐了回来，火尖全朝着黑甲将领那边，像一排让人牵着鼻子的牲口。

黑甲将领身上一点光都没亮，一丝魂力都没往外冒。他就站在那儿，周身像是多了一层看不见、摸不着的东西，那东西撑在他身子四周，撑得空气都变了样。

那两百个甲士的脸，在光柱里一张一张地变了。

本来是二百张不一样的脸，有方的有圆的，有老的有少的。光柱一起，那些脸上的神气就一层一层地往下褪，褪得发灰，褪得发木。眼珠子不转了，嘴角不动了，到最后，两百双眼睛跟黑甲将领那一双一模一样，冷冷的，齐齐的，一格一格地转。

小舞抓住了霍雨浩的袖子。

八个人的脸色一齐变了。

蔡须坤身子里的瑞兽本源乱了。它在她胸口里一下一下地撞，撞得她每一根经脉都绷紧了，像有人在她耳朵边上敲锣。她的脸一下子白了。

她看出来了。那两百个人的魂力，一点不剩地全灌进了那个人身上。灌进去以后，在里头让人揉成一团，压紧，再一点一点地重新拼起来。拼出来的那个东西，比两百个人加在一块儿还大。

唐三的蓝银草从石缝里探过去，碰了一下那些甲士的脚跟，又缩回来了。他的脸色变了。

``他们在结阵。''他说，``魂力没往一块儿堆，是在\ldots\ldots 对齐。每个人身上的魂力都让人重新排了一遍，排进一个序列里头去了。''

霍雨浩往后退了一步，额头上冒出了一层汗。

``我的精神力，''他说，``让人弹回来了。''

``弹回来了？''诺当扭头看他，``你的精神力，连封号斗罗的精神屏障都能摸进去一点。''

``我知道。''霍雨浩的嗓子发干，``可在那个人跟前，我那点精神力，跟一粒沙子往山上撞一样。''

他话没说完，鼻子底下流出来一道血。他拿袖子抹了一下，抹完了又流。小舞从怀里摸出一块帕子，塞到他手里，手还抓着他的袖子没放。

八个人里，卢本韦的反应最怪。

她没说话，就站在那儿，盯着那个黑甲将领看。她的手还在袖子里，指头一下一下地摸着那块猪神令上的纹路。她看那个人的眼神，不像是怕，倒像是在认，认一样她在家里的老书上见过、可从来没亲眼见过的东西。

``0 1111。''她小声说了一句。

声音很轻，只有站在她边上的蔡兴坤听见了。

蔡兴坤猛地转过头来。``你说什么？''

卢本韦没答，只是轻轻摇了一下头，眼睛一直没离开那个黑甲将领。她的嘴唇在动，像是在心里默念着什么。

过了一会儿，她才开口，声音还是很低。

``我家里有本老书，曾祖手抄的。''她说，``后头几页写的是序数修行。0 1111 那一页上画了一个人，身上一圈一圈的，跟他一个样子。书上说，走到这一步的人，身上不冒光，冒光的是他身边的人。''

``书上还说什么了？''

``说这一步走上去容易，走下来难。''卢本韦说，``走上去的人，身上的东西越攒越多，攒到后来，自己就压不住了。曾祖在那一页边上写了一行小字。''

``写的什么？''

``慎之。''卢本韦说。

蔡兴坤看了一眼那两百道往黑甲将领身上落的光柱，没再说话。

\begin{center}\rule{0.5\linewidth}{0.5pt}\end{center}

真霜斗罗的目光从宫门深处收回来，落在了那个黑甲将领身上。

他看了一会儿，开口了。

``序数修行。''他说，``两百个 BO 境界的人，合在一块儿，到头了就是这个样子。''

碎紫斗罗扭头看他。

``这个人，寻常魂师那条路上练不出来。''真霜斗罗的声音沉下去了，``他的力气，不从扽魂里来，也不从魂骨里来，是从序数里来的。他直接跟天地间那套序数的规矩搭上了，把那套规矩化进了自己身上。序数修行这一路，管他这个境界叫 0 1111，写成 Y 序列，是 Y(1,2,4,8,16)。''

``所以本座的魂力刚才往回缩\ldots\ldots{}''碎紫斗罗喃喃地说。

``你的魂力没缩。''真霜斗罗摇头，``是他借着身后这两百个人，把自己顶到了 Y(1,3,4,3) 上下。那个层级，天生就压在魂力那一套上头。你修到头的魂力，碰上序数，就跟小孩手里的玩意儿一样。''

他停了一下，往那两百个甲士身上扫了一眼，声音更低了。

``可老夫在意的，是这个阵。这个人倒在其次。''

碎紫斗罗一愣。``阵？''

``两百个魂尊，再多也是魂尊。''真霜斗罗说，``放在平时，别说两百个，一千个、一万个魂尊绑在一块儿，也碰不着封号斗罗一根头发。可现在\ldots\ldots{}''

他看着那两百个甲士，眼睛里有一点疑惑。

``他们让这个人领着，拿序列对齐的法子把魂力合到一块儿，打出来的分量，快到 Y(1,3,4,3) 了。''他说，``单拎出来哪一个，都比不上你碎紫的 Y(1,3,4,2,5,8,10,5)。合起来，能跟封号斗罗掰一掰手腕。''

碎紫斗罗的瞳孔缩了一下。

这句话有多重，他比谁都清楚。他自己就是 Y(1,3,4,2,5,8,10,5) 的封号斗罗，他知道这个层次意味着什么。他从十几岁当兵，一刀一枪地往上爬，爬了一百多年，才把这一串数一个一个地填满。每填一个数，身上就要添几道疤。眼前这两百个魂尊，让一个 Y(1,2,4,8,16) 的人领着，竟然能打出跟他平起平坐的分量。

魂尊跟封号斗罗之间，隔着的是萤火跟月亮。一百个魂尊绑在一块儿，也碰不着一个封号斗罗。这是魂师这一行里谁都知道的死理。

今天这条死理，让人当着他的面掰断了。

``不对。''碎紫斗罗的嗓子发干，``魂尊就是魂尊，魂力有多少，有多纯，都是有顶的。怎么叠，怎么合，也跨不过去这么大一道坎。除非\ldots\ldots{}''

``除非有人在后头动了手脚。''真霜斗罗接过他的话，``那两百个魂尊，魂力本身一点没涨，涨的是用法。那个人\ldots\ldots{}''

他看着黑甲将领，眼睛里多了一层戒备。

``他有法子，把两百个人的魂力，从魂力换成序数。魂力有顶，序数没有。一换过去，往上叠就不归魂师这一套管了，归序数那一套管。那一套的规矩，是往上翻着长。''

``所以两百个魂尊的魂力\ldots\ldots{}''碎紫斗罗的声音在抖。

``对齐以后，就不往上加了，一层一层往上乘。''真霜斗罗说，``所以两百个魂尊，带上一个 Y(1,2,4,8,16)，就敢站在两个封号斗罗跟前。他们没变强，他们的力气让人改写了。''

碎紫斗罗沉默了一会儿，慢慢吐出一口气。

``所以他们敢拦咱们。''他说，``扯着太子的旗子是一回事，人家手里是真有这个东西。''

他看了一眼那两百个人，又看了一眼真霜斗罗。

``前辈，这个阵，破得了吗？''

``破得了。''真霜斗罗说。

他说得很平，跟说今天天冷一样。

真霜斗罗没接这句话。他还看着那个黑甲将领，可眼底下浮出了一点说不清的东西。

不对。

他在心里想。

照他对序数修行的了解，Y(1,2,4,8,16) 的人，是能领着低一截的魂师对齐序列。可那要极细的功夫，代价也大，领头的人自己要挨反噬。眼前这个黑甲将领，气息是深，可真霜斗罗隐隐觉得，他的气息不稳。

他像是在压着什么。

还有那两百个魂尊。对齐序列，对人的身子、对人的神识要求都极高，寻常魂尊的经脉跟神识，根本扛不住这样的冲刷。他们扛住了，就说明他们的身子，早让什么东西改过了。

真霜斗罗的目光沉了一沉。他心里有了一个猜想，没说出口。

他又看了一眼那些甲士的胸甲。每一副胸甲正中，都刻着一个小小的徽记，跟黑甲将领胸前那一枚是一个样子，只是小了好几圈。徽记的纹路他在藏卷楼的卷宗里见过，画在一张从西边送回来的验尸图上，画在死人的手腕上。

\begin{center}\rule{0.5\linewidth}{0.5pt}\end{center}

就在这时候，宫门的阴影里，慢慢走出来一个人。

火把的光照在他身上，给他镀了一层暖黄。

他二十五六岁，个子高挑，脸生得端正好看，眉眼之间带着从小养出来的贵气。身上是明黄的太子常服，头上戴着紫金冠，腰间挂着一块龙纹玉佩。他走路不快不慢，每一步都走得端端正正，一抬手一投足，都是在上头待惯了的人才有的样子。

脚步声停了。刚才飘出来的那股龙涎香，就是从他身上来的。

他嘴角带着笑，温温和和的。他的目光从真霜斗罗身上扫过去，又扫过碎紫斗罗，最后落在八个人身上，在他们脸上停了一下。

``真霜前辈，碎紫将军。''

他的声音温润，带着几分关切，几分歉意，像是在接一位远道来的贵客。

``深夜劳动二位，孤没能出迎，是孤的不是。父皇龙体欠安，孤不得不替父皇看着这一摊事。宫里戒严，多有得罪，还请前辈不要见怪。''

他拱了拱手，礼数一点不差，挑不出一处毛病。

拱完了手，他又朝八个人笑了一下。

``这几位小友，孤记得。''他说，``大赛那年，孤亲手递过赏。''

他的目光在八个人脸上挨个过了一遍，过到蔡须坤脸上，停了一下。

``孤还记得，那时候是九个。''他说。

八个人没一个出声。诺当的拳头在身子边上攥紧了，骨节咯咯地响。蔡须坤站着，脸上一点变化都没有。

真太子递赏那天，记得的是她个子小。他不会记得那天台上站了几个人，更不会两年以后还拿这个数出来说。

记着这个数的人，是数过他们的人。两年前星斗大森林外头那条路上，有人数过一回。

太子。

或者说，一个自称是太子的人。

\begin{center}\rule{0.5\linewidth}{0.5pt}\end{center}

宫门前静得能听见火把燃烧的噼啪声。

蔡须坤看着那个太子，瑞兽本源在她身子里一阵一阵地往上顶，顶得她心口发闷。她的目光钉在那张脸上，一点细处都不放过。

她见过太子。大赛那年在宫里领赏，太子就站在陛下身边，替陛下把赏赐一样一样地递下来。递到她跟前的时候，他低头看了她一眼，说了一句''这么小的个子，怎么打赢的''，说完自己先笑了，笑得眉毛往上一挑，一点架子都没有。后来殿里有事进宫，她又见过两回，两回他都在跟侍卫比箭，射偏了就骂一句，骂完了接着射。

五官是像的，身量也差不多。可眼神不对。

真太子的眼睛里有少年人的那种劲儿，冲，张扬，看人的时候眉毛总是往上挑着。眼前这个人的眼神，太稳了，太从容了，稳得不像一个二十五六岁的人。倒像一个活了很久很久的人，躲在这张脸后头，往外看。

还有那股龙涎香。

她的鼻翼动了一下，侧头跟奥斯卡对了一眼。奥斯卡几乎看不出地摇了一下头，鼻子动得比刚才还快。

他早就闻出来了。这股龙涎香太新了，新得像是刚点上的。一个天天在宫里理事的太子，身上那股香应该是浸透了的，沉的，厚的，不会这么浮在外头。

让奥斯卡更在意的，是香底下压着的另一股味儿。很淡，淡得快透明了，像是药，可他叫不出是什么药。没有草木的清气，没有石头的涩味，倒像是冬天的雪化了以后，地上那股空空的、什么都不剩的气息。他说不上来那是什么，可他的鼻子告诉他，活人身上不该有这个味儿。

霍雨浩的精神力也摸回来了。他的脸色变了一点，精神力悄悄传到身边唐三的耳朵里：``他的精神力不对。面上看，就是个寻常的高阶魂师。可底下压着一层很厚的东西，像一扇门，门后头藏着什么。''

唐三的蓝银草也摸了一遍，他的目光沉了沉，精神力回过去：``他的魂力也不对。面上是魂斗罗顶头的样子，可太匀了，匀得一点起伏都没有。真魂斗罗的魂力，不会这么平。像是有人一点一点修出来的。''

八个人各摸各的，摸到的是同一个结果。

这个人，不是太子。

可没人开口。

那个太子笑吟吟地站在宫门前，身后是两百个 BO 境界的甲士，跟前是那个 Y(1,2,4,8,16) 的黑甲将领，像一尊铁铸的像，一动不动。火把照着每个人的脸，一明一暗。宫门关着，长戟横着，整片宫门前像是一只摆好了的笼子，把他们罩在里头。

真霜斗罗背着手站着，看着那个太子，脸上不恼，也不轻慢。

``太子殿下。''他开口了，声音平平的，像在说今晚的月亮，``老夫有些日子没见殿下了。殿下的气色，比上回见的时候好了不少。''

太子笑了一下，答得很从容：``劳前辈惦记。父皇病着，孤日夜操心，本该憔悴几分。只是这几日服了些灵药调养，气色好了一点，让前辈见笑了。''

``灵药？''真霜斗罗的眼睛动了一下，``什么灵药，能让殿下在陛下病重的时候，还有这么好的气色？''

太子脸上的笑没变，眼神却凝了一下。

``前辈要是好奇，''他说，``等父皇身子好了，孤禀明父皇，给前辈也备一份。''

``那就先谢过殿下。''真霜斗罗说，``城门里那张告示，殿下画的押，老夫也看了。殿下的字，比去年送到藏卷楼的那份奏本，写得稳多了。''

``孤这些日子静心练字。''太子说。

``练得好。''真霜斗罗说，``一个笔画都不抖。''

太子笑着，没接这句话。

可就在真霜斗罗说''一个笔画都不抖''的时候，太子的眼底下闪过了一点东西。灰白的，冷的，一闪就没了，快得像是火把晃了一下。

蔡兴坤肩上的黑子鸡，就在那一下，把脖子缩了回去，缩得只剩一个脑袋顶。它的爪子扣紧了蔡兴坤的肩膀，扣得她疼。她没出声，只是抬手，把它往自己领口里拢了拢。

碎紫斗罗这时候往前走了一步。

``殿下。''他说，``末将想先去寝宫看看陛下。''

``父皇刚睡下。''太子笑着说，``太医说了，不宜惊动。将军一片忠心，孤替父皇记下了。明日一早，孤亲自陪将军进去。''

``那殿下的手令呢？''碎紫斗罗说，``末将进不得寝宫，进宫总进得。宫门上的人说，要殿下的手令。''

太子从袖子里摸出一卷黄绫，托在手上，没往前递。

``手令在这儿。''他说，``不过今夜，孤不批。''

他笑着把黄绫又收回了袖子里。

碎紫斗罗的手攥紧了，骨节咯咯地响。

真霜斗罗没再接话。

他就那么看着太子，看得很静。他的眼睛像两面冰镜子，对面那个人身上一层一层的东西，一面一面地照过去。

他记得这位太子。十来年前，先帝还在的时候，有一年开春，宫里把十几岁的太子送到扽魂殿的藏卷楼读了三个月的书。那孩子坐不住，看半个时辰就要下楼去院子里扔石子，扔完了再上来，脸红扑扑的，满头的汗。有一回他扔石子砸碎了七层的一扇窗，自己跑上来认错，站在长案跟前，脖子梗着，说窗户钱从他的月例里扣。

眼前这个人站在宫门口，一身明黄，一脸温和，站得端端正正。

真霜斗罗没从他身上看见那个扔石子的孩子。

\begin{center}\rule{0.5\linewidth}{0.5pt}\end{center}

就在两边这么僵着的时候，蔡须坤又觉出了一样东西。

很细，不在这个太子身上，也不在宫门口这些甲士身上。在宫里头，在一重一重的宫墙后头，在灯火最亮的那一片，寝宫那边。

她往那边看了一眼。寝宫在宫城的最深处，隔着好几重宫墙，只看得见几道飞檐，飞檐底下的灯笼一盏一盏亮着，亮得发白。别处的灯火都在晃，就那一片不晃，像是那儿连风都没有。

一丝气味。

极淡，极细，像蛛丝一样，在夜风里飘过来，飘过去。她的瑞兽本源碰上它，往回缩了一下。它不像哪一种扽魂，哪一种魂技，她认得的味道里没有这一种。它更像一样脏东西，看不见，摸不着，可它在那儿，一点一点地往这座皇宫的每一个角落里渗。

真霜斗罗也觉出来了。

他闭了一下眼，把感知推到了头，去摸那一丝气味。过了一会儿，他睁开眼，眼底下有东西沉了一沉。

序列之毒。

这种毒不伤肉，不伤经脉，它从最底下的规矩下手，悄没声地改人身上的魂力怎么走。中了这种毒的人，面上一点事没有，身子里的魂力早就让人改了道。像一条河，在源头上让人悄悄挖了一锹，下游的水照样在流，可已经不是原来那条河里的水了。

他转头往四周看。宫墙，屋脊，角楼。那一丝气味哪儿都有，像一张看不见的网，把整座皇宫兜在里头。一间屋，一条廊，地上的砖，房上的梁，喝的水，喘的气，没有一处不是浸透了的。

他的手指头攥紧了，指甲掐进了掌心。

整座皇宫，都脏了。

这意味着什么，他一下子就想明白了。这座皇宫里的每一个人，昏着的陛下，中了毒的太子，宫女内侍，禁军侍卫，身子里早就让这种毒啃过了。他们的魂力，他们的念头，说不定连他们想干什么，都已经让人拿捏在手里了。

能下这么大一盘手的人\ldots\ldots{}

真霜斗罗的脑子里，冒出了一个名字。

序数修行这一行里，这个名字谁都听过，可亲眼见过他的人不多。

胡帅。

序列之毒就是他弄出来的。这个人的毒跟谁的都不一样，不碰肉身，不碰经脉，直接在序列那一层上改人身上的东西。

更让真霜斗罗心往下沉的，是另一件事。

胡帅，是胡赞的长子。

胡赞这个名字在序数修行这一行里有多重，在场的人里没有谁比真霜斗罗清楚。大陆最西边有一座堡，叫帅气堡，外头的人只知道那儿是做大买卖的帅气集团的老窝，藏卷楼的卷宗里却记着另一个名字：帅气教。胡赞就坐在那座堡里，坐了不知道多少年。藏卷楼里跟帅气教有关的卷宗，真霜斗罗管了几十年，一年比一年厚。胡赞本人深得看不见底，真霜斗罗活了一百多年，也只是听人说过，从没见过他出手。胡帅是胡赞的大儿子，从小跟着他爹修，他有多深，只会比外头传的更深。

胡帅要是插手了，那就不止是插手。看这满宫的序列之毒，这就是他亲手布的。一个 Y(1,2,4,8,16) 的黑甲将领，两百个让序列对齐改过身子的魂尊，一座从根上泡在毒里的皇宫。这几样东西，指向的是同一个人。

胡帅在后头。这一整盘，是他摆的。

真霜斗罗的心往下一沉。

他回过头，看着身后的八个人。蔡须坤，唐三，霍雨浩，小舞，诺当，奥斯卡，蔡兴坤，卢本韦，一个个脸色都很沉。他们觉出不对来了，可还没想到这件事有多大。

不能在这儿耗。

他一下子就拿定了主意。他跟碎紫斗罗两个人联手，兴许能跟那个黑甲将领和他身后的阵打个平手。可胡帅既然亲自下了场，宫里头就一定还藏着比眼前这些更硬的东西。在这儿耗下去，宫门进不去，还要一脚踩进人家挖好的坑里。

何况这八个孩子，修为最高的也才魂斗罗。真打起来，那阵上随便溢出来一点余波，就够他们受的。

得尽快冲出去。

真霜斗罗的眼神一下子利了。他侧过头，声音压得只有碎紫斗罗和八个人听得见：

``胡帅插手了。这个人是胡赞的长子，深得看不见底，眼前这个假太子跟那个黑甲将领，加起来也量不出他的斤两。''

碎紫斗罗的脸色一下子变了。

``胡赞\ldots\ldots{}''他的嘴唇动了一下，``帅气堡那个胡赞？''

``就是他。''

碎紫斗罗没再出声。他在禁军里待了一辈子，帅气堡的名字他听过不知道多少回。西边的城报上来的案子，十件里有好几件牵着那座堡，可回回到了朝里，就没了下文。他攥着的拳头松开了，又攥紧了。

宫门口，那个太子还笑吟吟地站着。他看着真霜斗罗侧过去的脸，笑着开了口。

``前辈在跟几位小友商量什么？''他说，``天冷，进宫喝杯热茶，慢慢商量也不迟。''

真霜斗罗没理他。

``不能再拖了。''真霜斗罗说。他脚底下的霜一下子往外扑开，扑得又快又急，``准备冲出去。''

``往外冲？''碎紫斗罗问。

``往里冲。''真霜斗罗说，``陛下在里头。''

碎紫斗罗看了他一眼，点了一下头，身上的紫光一下子亮了。

八个人听见了。

唐三脚边的蓝银草从石缝里一根一根地钻出来，粗了一圈，贴着地面往四边铺开。诺当身上漫出一层灰色的魂力，两只手垂在身边，指头尖上有细细的电光在跳。奥斯卡把手伸进袋子，摸出一把恢复香肠，挨个往身边人手里塞，塞到蔡须坤那儿，她摇了摇头，他硬塞进了她的手心。蔡兴坤肩上那只黑子鸡把一身的毛都奓了起来，脖子伸得笔直。小舞松开了霍雨浩的袖子，身子往下一沉，脚后跟离了地。霍雨浩把那块帕子团成一团揣进怀里，鼻子底下的血还没干。

卢本韦把两只手从袖子里拿了出来。

她右手攥着那块猪神令，攥得发烫，指头肚上让纹路硌出了一道一道的印子。她低着头，嘴唇动了一下，说了两个字，声音小得只有边上的蔡兴坤听见了。

``慎之。''

蔡兴坤扭头看了她一眼。卢本韦没抬头，把猪神令又往手心里按了按，按完了，揣回了袖子里，两只手空着，垂在身边。

蔡须坤站在最前头，隔着衣服按了一下那颗炮弹。炮弹贴着她的肋骨，一下一下地在震。

\begin{center}\rule{0.5\linewidth}{0.5pt}\end{center}

宫城外头，西华门外那条长街上，停着一辆马车。

马车很大，黑漆的车厢，车辕上挂着两盏灯笼，灯笼上写着一个''雪''字。车前车后站着几十个人，穿着王府的号衣，一个个垂着手，不说话。长街两头都让人堵上了，堵得严严实实，一个百姓都过不去。

车帘掀开了一角。

车里坐着一个人，五十来岁，身子发福，脸是圆的，下巴上留着三缕长须，修得一丝不乱。身上是亲王的紫金蟒袍，腰里束着玉带，头上戴着九旒冕。

他往宫门那个方向看了一眼。从这儿看不见宫门，只看得见宫墙上头那一片冲天的光，还有光底下一层薄薄的白，白得发冷，正一圈一圈地往外扑。

``冰。''他说，``真霜来了。''

车外有人低声应了一句什么。

他把车帘放下了。车里黑下来，只有他袖口上一道金边，在黑里头亮着。

``让城外的人进来。''他说，``三千个，一个都别少。''

车外的人应了一声，脚步声跑远了。

车里又静下来。他从袖子里摸出一枚棋子，黑的，捏在指头缝里，转过来，转过去。

``皇帝在里头。''他像是在跟车里另一个人说话，可车里就他一个，``峰儿会把他拿出来。活的。''

他把棋子往掌心里一扣。

``真霜挡得住峰儿。三千个人，再加上老夫，他挡不住。''

马车外头，灯笼上那个雪字让风吹得一晃一晃。长街两头，堵路的人一个都没动。

\begin{center}\rule{0.5\linewidth}{0.5pt}\end{center}

\chapter{第六十八章：冰霜怒涛}\label{ux7b2cux516dux5341ux516bux7ae0ux51b0ux971cux6012ux6d9b}

``准备冲出去''这几个字刚落地，真霜斗罗身上的气息就变了。

变得很快，一口气的工夫。

他身上原来那层平平和和的东西，像一片冰面，冰底下压着什么，谁也看不见。这时候冰面碎了，一块一块地裂开，裂缝底下那样东西露了出来。

碎紫斗罗猛地扭过头去。

他觉出来了。真霜斗罗身上的魂力，一口气翻了好几番。他在往外放，放得一点不留。像一口压了很多年的井，井盖让人一把掀了，底下的水一下子冲了上来。

Y(1,3,6)。

这才是真霜斗罗原本的样子。

碎紫斗罗的瞳孔缩紧了。他也是封号斗罗，同一层次的人，身上的分量他一摸就摸得出来。可这时候他摸着真霜斗罗身上那股东西，再摸摸自己身上的 Y(1,3,4,2,5,8,10,5)，觉得两样东西不在一个地方。差在底子上，多一点少一点倒在其次。

河跟海不一样。

``碎紫。''真霜斗罗没回头，声音平得发冷，``带他们退后三十丈。接下来的事，用不着他们。''

碎紫斗罗张了张嘴，想说什么。对上真霜斗罗侧过来的那一眼，话又咽回去了。

那双眼睛已经不是刚才那个和和气气的老人的眼睛了。冷，硬，碰不得。眼珠子底下有两点冰蓝的光，亮着，四周的热气一靠近就没了。

``\ldots\ldots 是。''碎紫斗罗转过身，朝八个人一挥手，``走！退后三十丈！''

蔡须坤还想说什么，唐三一把拉住了她的胳膊。

``真霜长老要动真格的了。''唐三摇了一下头，``咱们在这儿，只会碍事。''

八个人往后撤，撤得很快，一眨眼就退出去三十丈。

\begin{center}\rule{0.5\linewidth}{0.5pt}\end{center}

宫门前，只剩下真霜斗罗一个人。

他对面是那两百个甲士，最前头站着那个黑甲将领。黑甲将领后头，宫门的阴影边上，站着那个太子。

太子脸上的笑，头一回裂了一道缝。他的眼睛凝了一下，嘴角的弧度没那么松了，眼底下那一点灰白的光，闪得比刚才快了。

``前辈\ldots\ldots{}''他刚开口。

``闭嘴。''

真霜斗罗只说了两个字。

就这两个字，宫门前的空气冻住了。

真冻住了。空气里的水汽，在他话音落下的那一下，一齐结成了冰。冰碴子细得跟粉一样，密密麻麻地悬在半空里，把整片宫门前变成了一只大水晶球。火把的火苗让冰包住了，火光在冰里一折一折地散开，散成一道一道细碎的彩光，一点声音都没有。

然后，领域开了。

一点预兆都没有，一点蓄力都没有。真霜斗罗只是往前迈了一步。那一步落地，他脚底下荡开了一圈冰蓝的光，像往水里扔了一块石头，一圈一圈地往外漾。

光漾到哪儿，哪儿就冷。

冷得一下子到了底。刚才还是夜里那点凉意，一口气的工夫，就成了万年冰窟里的冷。广场的青石地上结了一层冰，一寸厚，冰底下的石板冻得咔咔地裂。宫门两边的火把，一齐冻成了冰蓝的冰柱子，火苗封在冰里头，还保持着烧的样子。

冰霜领域。

真霜斗罗压箱底的东西，极致之冰走到头，就是这个。

在这片领域里，他就是冰的主人。一片雪，一缕寒气，一寸冰面，都是他手指头的延伸。冷到什么地步他说了算，冰往哪儿长他说了算。不归他的热，进来一点，就让它吞一点，吞到一点不剩。

可这片领域真正吓人的地方，在压。

领域铺开的那一下，宫门前每一个人都觉出了一样东西涌过来。那东西比风沉，比魂力撞人还往下一层。它在冻，冻的是人身上一切在动的力气。

那两百个甲士头一个挨上。

最外圈那几十个，脸色一齐变了。他们身子里的魂力，碰上这股冷，一下子就稠了，像河里灌进了泥浆，流得慢了好几倍。有人想催魂力顶住，魂力刚出经脉，就让外头的冷冻成了冰渣子，又倒灌回去，震得他们一口血喷出来，喷在冰面上，转眼也冻成了红的冰花。

可反应最大的，是站在最前头的黑甲将领。

他的身子猛地一僵。

那双一直像死水一样的眼睛里，头一回有了动静。他吃了一惊，那一惊是从骨头缝里往外冒的。

他觉出来了。领域漫进他身子里的那一下，他身上那一套拿序数搭起来的东西，像碰上了天敌，一下子往里缩。他从来没有过这样的感觉。像是有一只看不见的手，伸进他的经脉里，一寸一寸地，往外抽他的力气。

Y(1,2,4,8,16)。

这是他的境界。序数的规矩在他身子里搭起了一座又大又精细的架子，每一节都像齿轮一样咬得严丝合缝，源源不断地往外送力气，送出来的力气，魂师那一套够不着。

可在冰霜领域里，那些齿轮开始冻住了。

先是最末那一节，十六。像一盏灯让风吹灭了，那一点光在冷里头一下子就没了。

然后是八，是四，是二，一节跟着一节灭了。

黑甲将领的脸一下子白得跟纸一样。他身上的力气像退潮的水，退得飞快。每灭一节，他就往下掉一层。

他想往回够，够 Y(1,2,4,8,15)。

够不着。

他的精神力拼命地往那些冻住的节上撞，想把它们重新撞亮。可冷已经钻进了他力气的根上，那些节像冰封住的泉眼，他怎么催，都催不出一滴水来。

连 Y(1,2,4,8,15) 都使不出来了。

刚才还能跟封号斗罗掰手腕的一个 Y(1,2,4,8,16)，在真霜斗罗的领域里，连低一格的力气都调不动。

``这\ldots\ldots{}''黑甲将领的声音头一回在抖，``这是什么领域\ldots\ldots{}''

他想不明白。序数修行的力气，从天地的规矩里来，魂力那一套管不着它，这是他一直拿来压人的本钱。可眼前这个老人的领域，能把序数的规矩一块儿冻住。

这片领域最吓人的地方，就在这儿。

它冻的不光是冷热，它从力气最底下那一层下手，凡是不归冰管的力气，在里头都转不动。魂师的魂力是这样，序数修行的那一套规矩也是这样。只要不是冰，进来了，就压，就削，一直削到没有。

对那个黑甲将领来说，这一压最要命。他的力气全靠序数的规矩转着，规矩要一节一节咬得准，才使得出来。冷一来，齿轮冻住了，规矩就转不动了，规矩转不动，力气就没了出处。

太子的脸色，也终于变了。

他眼底那一点灰白的光，疯了一样地闪，像风里的一截蜡烛。他往后退了半步，嘴唇动了动，一点声音都没出来。

\begin{center}\rule{0.5\linewidth}{0.5pt}\end{center}

三十丈外，碎紫斗罗站着，一动不动。

他的嘴半张着，眼睛里映着宫门前那一片冰蓝的光，还有光当中那个背着手的老人。

整个过程他都看见了。

那个黑甲将领，刚才还气势压人，能跟封号斗罗掰手腕，进了领域以后，就像让人抽了脊梁骨，力气一节一节地往下掉，连 Y(1,2,4,8,15) 都使不出来了。

碎紫斗罗的喉结上下滚了一下。

他自己就是封号斗罗。Y(1,3,4,2,5,8,10,5)，放在整个天斗帝国，也是站在顶上的那几个人里的一个。他知道封号斗罗的力气是什么分量，也知道要压住一个 Y(1,2,4,8,16) 的序数修行者，得费多大的劲。

换了他自己，对上同一个人，赢是能赢的，可少说也得缠斗几十招，说不定还要挂点彩。

真霜斗罗连手都没出。

他就是把领域铺开了。

光是领域往下一压，就把那个黑甲将领从 Y(1,2,4,8,16) 压到了连 Y(1,2,4,8,15) 都够不着的地步。

同是封号斗罗。

碎紫斗罗在这一刻想明白了一件事。封号斗罗跟封号斗罗，站的不是一层楼。Y(1,3,6) 跟他的 Y(1,3,4,2,5,8,10,5)，看着都是一串数，中间隔着的东西，比数字上那一点差多得多。

真霜斗罗的 Y(1,3,6)，是把什么都往一处收。不分心，不散开，一身的魂力压得密得吓人，压成了一把冰做的刀。他自己的 Y(1,3,4,2,5,8,10,5)，串长，铺得开，管得宽，可论密，跟真霜斗罗比，差得远。

一把薄刀，跟一面厚盾。

薄刀只有一道刃，可那道刃什么都切得开。厚盾结实，可盾到底只是挡。

碰上真能切开一切的东西，铺得宽反倒成了累赘。

``真霜前辈\ldots\ldots{}''碎紫斗罗小声说，他自己都没听出来自己声音里带着敬，``这就是 Y(1,3,6)\ldots\ldots{}''

他到今天才明白，序数修行这一行里，为什么提起真霜斗罗，人人都要低一低头。他的魂力算不上最多，魂技也花不到哪儿去。可是他把冰这一个字，修到了头。

\begin{center}\rule{0.5\linewidth}{0.5pt}\end{center}

宫门前，领域还在往外铺。

黑甲将领单膝跪在了冰面上，两只手撑着冰，指节都白了。他身上还剩一层很薄的暗金色的光，那是序数的规矩在他身子里挣最后一下。那层光一明一灭，随时要熄。

他慢慢抬起头，隔着满天的冰碴子，看着真霜斗罗。

``Y(1,3,6)\ldots\ldots{}''他的声音哑得像铁片子在石头上刮，``能把序数冻住\ldots\ldots 名不虚传。''

真霜斗罗背着手站着，周身一层冰蓝的光，一点声音都没有。他看了黑甲将领一眼，又看了看他身后那两百个已经倒的倒、跪的跪的甲士，嘴角往上动了一下，那一下很冷。

``胡帅的狗。''他说，声音不大，冷得往骨头里钻，``回去告诉你主子。''

他慢慢抬起右手，指尖上亮起了一点冰蓝的光，在满天的冰碴子里，像一颗掉下来的星星。

``天斗帝国的皇宫，轮不到他撒野。''

话落了，宫门前一片死静。

黑甲将领跪在冰上，身上那层暗金的光明明灭灭。那两百个甲士有的跪着，有的趴着，魂力乱成一团，一时半会儿谁也站不起来。

太子站在宫门的阴影边上，脸色灰了。他眼底那一点灰白的光暗到了底，像一根快烧完的灯芯。

他往后退了一步，又退了一步。宫门的阴影把他吞进去了。

真霜斗罗没追，也没回头。

``碎紫。''他说，``带人进去。''

碎紫斗罗一下子回过神来，重重地点了一下头。``是！''

他转身朝八个人大步走过去，身上的紫光一涨，刚才那点愣怔一扫就没了，又是那个统领三军的大将军。

``走。''

\begin{center}\rule{0.5\linewidth}{0.5pt}\end{center}

八个人跟在碎紫斗罗身后，从跪了一地的甲士中间穿过去，进了皇宫正门。

一跨过门槛，蔡须坤身子里的瑞兽本源就猛地一颤。

序列之毒的气味，比外头浓了十倍都不止。

宫门外那点毒，细得跟蛛丝一样，有一下没一下的。宫门里头，毒像雾一样，到处都是，往哪儿钻的都有。空气里，地砖上，廊柱上，飞檐上，每一寸地方都浸透了。蔡须坤甚至觉得出来，那毒在往她身子里钻，像无数根细得看不见的针，从她身上每一个毛孔往里扎。

``屏住气，拿魂力护住身子。''蔡须坤压着嗓子说，瑞兽本源在她身子里飞快地转，把钻进来的毒一丝一丝地往外逼，``这儿的毒太浓了，别碰。''

八个人一齐催动魂力，身上罩了一层薄薄的护罩。毒撞在护罩上，发出极细的嘶嘶声，像醋泼在铁上。护罩肉眼看得见地在往薄里去，可一时半会儿还撑得住。

奥斯卡一边走，一边从袋子里摸出几根香肠，塞给走在边上的小舞和诺当。

``嚼着，''他说，``嚼着护罩长得快。''

诺当把香肠塞进嘴里，嚼了两口，身上那层灰色的护罩果然厚了一点。

碎紫斗罗走在最前头。

他在这座皇宫里进进出出几十年，哪道门通哪条廊，哪个拐角后头有台阶，闭着眼都摸得着。他领着八个人穿过前殿，绕过御花园，从内侍值房的屋檐底下掠过去，一路往寝宫那边赶，脚底下一步都没停过。

路上碰见的人不少。

前殿的廊子底下，有个宫女拿着一把扫帚在扫地。地上什么都没有，青砖干干净净的，她还在扫，一下，一下，扫的是同一块砖。扫帚的竹枝磨秃了，在砖上划出沙沙的响。

御花园边上的亭子里，一个老内侍站在石桌跟前倒茶。茶壶嘴对着茶杯，杯子早满了，茶水漫出来，顺着石桌往下淌，淌了一地，在青砖缝里结了一层薄冰。他手上的壶还歪着，壶里早空了，他还是那么歪着。

值房门口站着四个禁军，刀拄在地上，脸朝着同一个方向，眼睛睁着，一眨不眨。九个人从他们跟前走过去，四个人连眼珠子都没转一下。

谁也没拦他们，谁也没出声。九个人从这些人身边走过去，像一阵风从院子里穿过去，那些人照旧扫地，照旧倒茶，照旧站着。

小舞走过那个扫地的宫女身边，脚底下慢了一步。

那个宫女看着十五六岁，脸圆圆的，鼻尖上有几粒雀斑。脸白得没有血色，嘴唇发青。她低着头，一下一下地扫，嘴里在念叨什么，声音极小。

小舞弯下腰去听。

``\ldots\ldots 卯时扫前殿\ldots\ldots 辰时扫前殿\ldots\ldots 卯时扫前殿\ldots\ldots{}''

就这一句，翻来覆去。

小舞直起身来，伸手在她肩膀上轻轻拍了一下。``姐姐？''

宫女没抬头，扫帚也没停。

``别碰她。''霍雨浩在后头说，声音压得很低，``她的精神力，让人拿线缝死了。就留了一个口子，口子里只剩下扫地这一件事。你叫她，她听不见。''

小舞把手收回来，在衣服上蹭了一下，又蹭了一下。

``都让人拿住了。''碎紫斗罗头也不回，声音沉着，拳头上一层紫光凝起来又散开，``序列之毒啃到了脑子里。这些人，眼下已经由不得他们自己了。''

唐三的目光从那个宫女脸上扫过去，又从亭子里那个倒茶的老内侍脸上扫过去。他的眉头拧着。

``他们还有救吗？''

``不知道。''碎紫斗罗的嗓子发干，``要是毒啃得太深\ldots\ldots{}''

他没往下说。

奥斯卡回头看了那个老内侍一眼，从袋子里摸出一根香肠，想了想，又塞回去了。

``我的香肠管魂力，''他小声说，``管不了这个。''

没人接他的话。

碎紫斗罗一路没回头。可走过御花园那座假山的时候，他的脚底下慢了半步，往假山后头看了一眼。

假山后头有一口井，井沿上搭着一只木桶，桶里的水结了冰。

``上个月，''他说，声音很低，像是说给自己听，``陛下还在这口井边上跟本座下棋。他输了，非说是井里的蛤蟆叫得他分了心。''

没人吭声。

碎紫斗罗把头扭回去，加快了步子。

卢本韦走在最后。她经过那四个禁军的时候，往他们脖子上看了一眼。四个人的脖子侧面，都有一道灰白的细线，从耳朵底下往衣领里头钻。她把那块猪神令从袖子里摸出来，攥在手心里，指头按着上头的纹路，按得很紧。

\begin{center}\rule{0.5\linewidth}{0.5pt}\end{center}

寝宫的门关着。

门前站着两排禁军，一边十个，刀都出了鞘，竖在胸前。他们的眼睛空着，脸是灰的。皮肤上爬着一层灰白的纹路，像蜘蛛网，从脖子一直爬到手背上，有几个人的纹路已经爬到了脸上，从下巴往颧骨上走。

九个人走到台阶底下，二十个禁军一齐转过头来。

二十双眼睛，转得一样快，停得也一样齐。

蔡须坤的脚底下顿了一下。

这些人比外头那些啃得深。外头的人只是丢了魂，这些人身子里还剩着一点东西，那点东西让人拧过了，拧成了一个死扣：守着这扇门，谁来都不让进。

``让开。''碎紫斗罗说。

二十把刀一齐往前一递。

碎紫斗罗没再说第二句。他身上紫光一涨，右手往前一推，一掌拍了出去。掌风从台阶上卷过去，卷得两排人一齐往后飞。二十个禁军撞在廊柱上，撞在门框上，滑下来，躺了一地，刀叮叮当当地落在台阶上。

诺当跑过去，蹲下来，挨个摸了摸脖子。

``都有气。''他说，``昏过去了。''

``得罪了。''碎紫斗罗冲地上那一片人说了一句，抬手推开了寝宫的门。

\begin{center}\rule{0.5\linewidth}{0.5pt}\end{center}

门一开，里头的气味扑了出来。

龙涎香，药汤子，还有一样别的，搅在一块儿，又闷又稠，堵得人嗓子眼发紧。屋里的窗户全关着，帘子也全放下来了，只点着几盏灯，灯火是青的，照得满屋子发暗。

正当中摆着一张大龙榻。

榻上躺着一个人，锦被盖到下巴，只露出一张脸。脸瘦得颧骨都凸出来了，白得跟窗户纸一样，眼窝陷下去两个坑，嘴唇上起了一层干皮。

雪夜大帝。

蔡须坤在大赛那年见过他。那时候他坐在高台上，嗓门洪亮，笑起来整个看台都听得见，给她递赏的时候还拍了拍她的肩膀。

眼前这个人，她差一点没认出来。

她身子里的瑞兽本源猛地一震，震得她心口发麻。她快步走到榻前，两只手悬在大帝的胸口上头，没往下按，瑞兽本源从她掌心里涌出去，隔着锦被，往大帝身子里探。

探了一会儿，她的脸色越来越难看。

``已经到骨头里了。''她说，嗓子在抖，``这毒，寻常药石碰不着它。它在改陛下身子里的规矩。陛下经脉里的魂力，快有一半让它改成了别的东西。照这么下去\ldots\ldots{}''

``多久？''碎紫斗罗问。他站在榻尾，两只手攥着拳，指节都白了。

``最多三天。''蔡须坤咬了一下嘴唇，``三天以内，陛下的魂力会让它改干净。改干净了，陛下就成了一个空壳子。到那时候，躺在这儿的就是一个傀儡，人家让他往哪儿走，他就往哪儿走。''

屋里没人说话。

灯火跳了一下。帘子后头有一只飞蛾在撞窗纸，扑扑地响。

蔡兴坤站在门边，手按在门框上，没往里走。她肩膀上的黑子鸡缩成一团，脑袋埋进翅膀底下，一动不动。

``让我来。''

门口有人说话，声音平平的。

真霜斗罗什么时候进来的，屋里没一个人听见脚步声。他站在门槛里头，身上还挂着领域散下来的一点冷气，白里泛蓝的发梢上结着一层细霜。他一进来，屋里那股闷稠的气味就往四边退，退到墙根底下，贴着帘子缩成一团。灯火的青色淡了一点，屋里冷了好几度。

蔡兴坤往边上让了一步。她肩膀上的黑子鸡把脑袋从翅膀底下伸出来，看了真霜斗罗一眼，又缩回去了。

``外头呢？''碎紫斗罗问。

``跪着。''真霜斗罗说，``那个太子跑了，那个领头的也跑了。剩下的两百个，一时半会儿起不来。''

他走到龙榻前，低头看着雪夜大帝的脸。

看了很久。

他第一回见这个人，这个人还是个十来岁的少年，跟在老皇帝后头来扽魂殿，站在良序圣殿的台阶底下，仰着头看殿顶上的冰棱子，看得脖子都酸了。后来少年当了皇帝，每年开春都要来殿里坐一坐，跟他下一盘棋，输了就耍赖，非说那一步是手滑。

真霜斗罗伸出右手。

五个指头张开，手心朝下，悬在大帝的胸口上头三寸。一缕冰蓝的魂力从他手心里淌出来，细得跟一根线一样，往大帝身子里渗。

Y(1,3,6) 的魂力。

那缕魂力一进去，蔡须坤的瑞兽本源就猛地一跳。她还悬着手，这一下，她看见了。

大帝的经脉里，盘着一条一条灰白的东西，细长，软，像蛇。它们缠在经脉壁上，一口一口地啃大帝身子里剩下的那点魂力，啃一口，那一截就变灰一截。真霜斗罗那缕冰蓝的魂力一进去，正撞上它们。

冷气走到哪儿，哪儿的灰蛇就缩。

它们缩得很快，一碰上冷，身子就僵了，硬了，表面结出一层白霜。有几条想往回咬，张嘴去啃那缕冰蓝的魂力，牙还没落下去，嘴就冻上了。

真霜斗罗的冰只认一个字，冻。不归冰管的东西，碰上了就冻。那些灰蛇想改规矩，规矩还没改动，先让冻住了。

灰蛇一条一条地退，从大帝的心脉上退下来，从几条大经脉上退下来，一直退到经脉的边边角角。真霜斗罗的魂力追过去，在它们外头结了一层薄薄的冰壳子，把它们封在里头。

封住了。可也只是封住了，一条都没拔出来。

真霜斗罗的眉头皱了一下。

他觉出来了。这毒的根，不全在大帝身上。它是从外头来的。屋里的气，桌上的水，墙上的漆，帘子上的穗子，一样一样地往大帝身子里渗。他在里头封住一条，外头又渗进来一丝。冰冻得住河，冻不住上游。

他把手收回来。指尖上那点冰蓝的光一点一点灭了。

他的脸比进门的时候白了一层。拿 Y(1,3,6) 的魂力去顶序列之毒，一寸一寸地封，他也费了不小的力气。

``先压住了。''他说，``撑不长。''

``毒是从这座宫里来的。''他看了一眼四周的墙，``根扎在这一整座皇宫里头。陛下待在这儿一天，毒就往他身子里渗一天。老夫的冰能拖着它，拔不掉它。''

碎紫斗罗的眉头拧成一个疙瘩。``那怎么办？''

真霜斗罗没马上答。他的目光从屋里每一个人脸上扫过去，扫过碎紫斗罗，扫过蔡须坤，最后在八个人身上停住了。

``把陛下挪出去。''他说，``出了宫墙，毒就淡了。离了根，老夫再拿冰慢慢地往外拔，配上药，陛下的身子自己能缓过来。''

``挪？''碎紫斗罗一愣，``陛下这个样子，经得起颠？''

``不用走远。''真霜斗罗说，``出了宫墙就行。宫里头的毒最浓，过了宫墙，一下子就淡了。''

他停了一下，声音沉下去。

``可得快。老夫这层冰，最多撑三天。三天以内，陛下必须出这座宫。要不然\ldots\ldots{}''

他没往下说。屋里的人都听懂了。

碎紫斗罗吸了一口气，身上的紫光往上一涨。``那就现在走。''

他转过身，看着八个人。

``你们八个，护着陛下。''他说，``本座在前头开路，真霜前辈压在后头。宫里的禁军让毒拿住了，里头有硬手。这一路，少不了有人拦。''

蔡须坤点了点头，瑞兽本源在她身子里稳稳地转了一圈。``明白。''

唐三点了点头。霍雨浩，小舞，诺当，奥斯卡，蔡兴坤，卢本韦，一个跟着一个点了头。

\begin{center}\rule{0.5\linewidth}{0.5pt}\end{center}

要挪人，得先有个抬的东西。

唐三往屋里看了一圈，蹲下来，蓝银草从他手心里钻出去，顺着地砖缝爬到墙角，从那儿拖出两根挑帘子用的长竹竿。草叶子一根缠一根，在两根竹竿当中织出一张网，织得又密又匀，像一张软床。

``担架。''唐三说。

小舞把龙榻上的一床薄被抽下来，铺在蓝银草织的网上。诺当跟蔡兴坤一个托头，一个托腿，把大帝从榻上挪过去。大帝轻得吓人，诺当托着他的肩膀，手底下全是骨头。

真霜斗罗伸出一个指头，在大帝的眉心点了一下。一点冰蓝的光渗进去，大帝的脸上结了一层极薄的霜，眉毛上白了。

``这层霜挡着外头的毒。''他说，``路上别让人碰他的脸。''

奥斯卡从袋子里摸出一根香肠，在手里掂了掂，又放回去了。他转头去看蔡须坤：``队长，陛下能吃香肠么？''

``他嚼不动。''蔡须坤说。

``那我切碎了。''

``他咽不下去。''

奥斯卡张了张嘴，把袋子口扎紧了，没再说话。

卢本韦从地上捡起一把禁军落下的刀，在手里转了一圈，又扔了。她把猪神令揣回袖子里，走到担架边上，跟诺当换了个手，自己托住了大帝的肩膀。

``我力气大。''她说。

霍雨浩站在门口，精神力往外铺出去，铺过前殿，铺过御花园，一直铺到宫门那边。他闭着眼，过了一会儿，睁开了。

``一路上，''他说，``都还是刚才那样，扫地的扫地，站着的站着。没人在等咱们。''

碎紫斗罗点了一下头，没说什么，眉头却拧得更紧了。

\begin{center}\rule{0.5\linewidth}{0.5pt}\end{center}

真霜斗罗最后看了一眼空了的龙榻。

锦被掀在一边，枕头上压出一个浅坑，坑里落着几根白头发。榻边的小几上搁着一只药碗，碗底还剩一口黑的药汤，汤面上结了一层皮。

他把那只碗拿起来，凑到鼻子底下闻了一下，又放下了。

然后他慢慢吐出一口气。

``走吧。''他说。

他的声音平平的，跟一百年前他头一回走进这座皇宫的时候一样。

``天斗帝国的天子，不该死在这种阴沟里。''

\begin{center}\rule{0.5\linewidth}{0.5pt}\end{center}

\chapter{第六十九章：深渊}\label{ux7b2cux516dux5341ux4e5dux7ae0ux6df1ux6e0a}

担架在当中。

卢本韦托着头，诺当托着脚，蓝银草织的网兜着雪夜大帝，一颠都不颠。唐三走在担架左边，手里攥着一把蓝银草的草梢，草梢连着担架的网，网要是哪儿松了，他手指头一动就能补上。蔡兴坤走在右边，黑子鸡站在她肩膀上，脖子伸得老长。蔡须坤在担架前头，瑞兽本源铺开一层，把大帝跟他们几个罩在里头。霍雨浩跟小舞一前一后，奥斯卡夹在当中。

碎紫斗罗在最前头开路。真霜斗罗压在最后头，离担架十来步，一步一步地走，脚底下的霜跟着他走。

他们从寝宫出来，穿过内宫的甬道，绕过太液池，顺着御花园西边的长廊往宫门那边赶。

宫墙上的火把还亮着，让领域散出来的冷气一激，火苗子一跳一跳的。橘红的火光跟地上的白霜搅在一块儿，把甬道照得一截亮一截暗。

一路上没人拦。

前殿廊子底下那个宫女还在扫那块砖。亭子里那个老内侍还歪着那只空壶。值房门口那四个禁军还拄着刀站着，脸朝着同一个方向。

快到二门的时候，从门洞里出来了七八个禁军，刀出了鞘，直着眼往这边走。碎紫斗罗没停步，一掌推过去，七八个人一齐飞出去，撞在门洞的墙上，滑下来，不动了。

就这么一回。

再往前，又没人了。

蔡须坤的步子慢下来半步。她身子里的瑞兽本源一直在转，转得越来越急，转得她后脖颈子发紧。

``太顺了。''她压着嗓子说，``从寝宫到这儿，像是有人把路扫干净了，等着咱们走。''

唐三点了一下头。他的蓝银草顺着地砖缝往前爬，爬出去几十丈，一路摸过去，什么都没摸着。

``咱们进来的时候，宫门口拦着两百个人，一个黑甲的，一个太子。''他说，``现在咱们要把人往外带，倒一个拦的都没有了。''

``那两百个呢？''小舞问，``真霜长老说他们跪着，起不来。''

``起不来，也该有人看着他们。''唐三说，``可我的草摸过去，宫门那边的地上，一个跪着的都没有了。''

霍雨浩一直没说话。他的精神力早就铺出去了，铺过二门，铺过前殿，往宫门外头摸。

摸到宫门那儿，他的脚底下一顿，脸色变了。

``前头，''他说，``宫门外头，有人。很多人。''

碎紫斗罗一下子站住了。

``多少？''

霍雨浩的额头上冒出一层汗。他把精神力又往外推了推，推得太阳穴一跳一跳地疼。

``三千往上。''他说，``还有魂斗罗。二十个，也许更多。他们站得很齐，在等着。''

三千人。二十个魂斗罗。

碎紫斗罗的瞳孔缩了一下。他在城里管着兵，城里有多少兵，在谁手里，他一清二楚。三千人，比禁军在宫门外头的整个营还多。这么些人要进城，要在宫门外头摆开，他的人不可能一点消息都没传回来。

除非这些人本来就不在城里。

他想起来了。他们进城的时候，城门是开着的。城门那个将领，手里攥着一个圆盘子，往宫里打了一道光。

``走快些。''真霜斗罗在后头说。

他的声音还是平的，可脚底下的霜往前扑了一大片，一直扑到担架底下。

碎紫斗罗一咬牙，身上紫光一涨，领着人往宫门那边冲过去。

\begin{center}\rule{0.5\linewidth}{0.5pt}\end{center}

皇宫正门。

九个人冲出宫门，脚底下一齐钉住了。

宫门外的广场上，火把一根挨着一根，插满了。

黑压压的人从宫门前一直排到长街的尽头，把整片广场围得严严实实。长戟竖着，刀出着鞘，火光照在刀刃上，白花花的一片。三千个人身上溢出来的魂力搅在一块儿，压在广场上空，压得人喘气都短了一截。

刚才跪在这儿的那两百个甲士，一个都不见了。冰面让人砸碎了，碎冰扫到了广场两边，堆成两道矮墙。

大军的正前头，站着一个人。

五十来岁，身子富态，脸盘圆润，下巴上留着三绺长须，修得一根不乱。身上是一件紫金的蟒袍，亲王的规制，腰上束着玉带，头上戴着九旒的冠，冠上的玉珠子垂下来，在火光里一晃一晃的。

他背着手站着，看着宫门里冲出来的这几个人，脸上带着笑。

雪星亲王。

碎紫斗罗的嘴张开了，半天没合上。

一个时辰以前，他在城西那条巷子里，亲手把这个人从墙根底下扶起来。那时候这个人满脸是血，左胳膊吊着，右腿一瘸一拐，让他的两个亲兵架着送回了王府。

``王爷？''他的声音一下子拔高了，紫光从他身上一冲，冲得脚底下的碎冰都蹦起来，``你\ldots\ldots 你怎么在这儿？！''

雪星亲王笑了一下。笑得温温和和的，像长辈看着一个不懂事的小辈。

``碎紫将军。''他说，``深更半夜，带着人闯宫，把陛下从寝宫里抬出来。本王倒要问问你，这是要反么？''

碎紫斗罗的拳头攥紧了，指节咔咔地响。

``你\ldots\ldots{}''

``他是假的。''

真霜斗罗从宫门里走出来，走到碎紫斗罗身边，站定了。

他的目光从广场上扫过去。三千甲士，二十几个魂斗罗，最前头那个穿着蟒袍的亲王。他脸上没什么变化，眼底下有一点东西落定了，像一个早就猜到的数，终于算出来了。

``真的那一个，''他说，``左胳膊的经脉断了好几截，右腿的骨头裂了。一个时辰以前，雨浩亲手摸过的。''

霍雨浩在后头点了一下头。

``这一个，两条腿站得稳稳当当，''真霜斗罗说，``左手背在后头，攥着右手的腕子，攥得一点都不抖。''

碎紫斗罗的眼睛一下子转到那个亲王背在后头的手上。

``还有。''真霜斗罗说，``真的雪星亲王，修为就到魂王。一个魂王，站在两个封号斗罗跟前，喘气不会这么匀。''

那个亲王听着，脸上的笑一点都没变。

碎紫斗罗的脑子里，一样一样的东西往回翻。

进城的时候，城门大开着。告示上盖着东宫的印，说下毒的是雪星亲王。满街的人传的话，一个字都不差：太子还活着，是雪星亲王下的毒，二皇子也跟着他跑了。城门那个将领手里的圆盘子，往宫里打的那一道光。宫门前那两百个甲士，那个黑甲的，那个太子，拦得那么死，一放又放得那么干净。

还有那个真亲王。浑身是伤，靠在墙根底下，偏偏让他们撞上了。

``是个套。''碎紫斗罗的嗓子发干，``从城门那儿就是。他们放出风来，说亲王造反，就为了让本座带人闯宫\ldots\ldots{}''

``闯宫只是头一步。''真霜斗罗说，``你再往下想。''

碎紫斗罗愣了一下。

``他们要的，是咱们把陛下抬出来。''真霜斗罗看着那个亲王，声音很低，``陛下在宫里，他们是挟天子。陛下让碎紫大将军从寝宫里抬出来，抬到宫门口，让亲王领着三千人截住了，那就是亲王救驾，大将军劫驾。明儿一早，满城的告示就换了。''

碎紫斗罗的后背一下子凉透了。

``宫门口那两百个人，那个黑甲的，是拿来拖的。''真霜斗罗接着说，``拖住咱们，拖到这三千人从城外头进来，摆开。进宫的路他们留着，出宫的路他们堵着。咱们进得去，抬得出人来，到了这儿，就走不了了。''

``可他们没算到前辈的领域\ldots\ldots{}''

``算到了。''真霜斗罗说，``老夫的领域一开，那两百个人跪下了，他们照样能走。你看这地上，冰都让人扫干净了。人家把账算到了老夫头上。''

那个亲王一直没插嘴。他背着手站着，笑眯眯地听着两个人说话，像在戏园子里听一出跟自己没关系的戏。

真霜斗罗把目光落回他身上。

``你是谁？''

亲王没答。

他把背在后头的右手拿到前头来，抬起来，往下轻轻一按。

三千个甲士同时催动了魂力。

三千道光柱从广场上冲起来，冲到半空里，横着竖着搅成一张网，把宫门前的天都罩住了。那二十几个魂斗罗往前迈了一步，一齐迈的，脚底下的青石板一齐裂了。二十几道魂斗罗的魂力压下来，压在宫门前每一寸地方上。

担架上，大帝脸上那层霜抖了一下，裂了一道细缝。

蔡须坤的脸一下子白了。她身子里的瑞兽本源猛地往外一涌，扑在担架上头，死死压住那层霜。

``长老！''她喊，``陛下脸上的霜撑不住！''

真霜斗罗没回头。

他的眼睛还看着那个亲王。看着那只抬起来又按下去的右手。

那个按手的样子，太稳了，太随便了。三千个人，二十几个魂斗罗，他按一下手就全动了，他自己身上连一丝魂力都没往外冒。

真霜斗罗心里那个名字，从宫门口闻见序列之毒的时候就冒出来了。他一直压着，没往外说。这时候它自己顶了上来。

他闭上了眼睛。

他把精神力全收拢了，收成一根线，细得跟冰锥的尖子一样，悄没声地往那个亲王身子里扎过去。

一息。两息。三息。

他睁开了眼。

那张一百多年没怎么变过颜色的脸，这时候白了一层。

Y(1,3,7,12)。

他自己是 Y(1,3,6)。碎紫是 Y(1,3,4,2,5,8,10,5)。进城以来，他一直当自己是这座城里最硬的那一个。

眼前这个穿着亲王皮的人，是 Y(1,3,7,12)。

Y(1,3,6) 往上，是 Y(1,3,6,\ldots) 那一长串，一层一层往上叠，叠到头，才摸到 Y(1,3,7)。Y(1,3,7,12) 还要再往上去好几层。他站的地方，比真霜斗罗的头顶还高出一大截。

真霜斗罗这才想明白，刚才为什么一直没看穿。

那层皮做得好，那是一回事。要紧的是对面比他高。高出这么多，他的感知探上去，就像拿一根竹竿去量天，量到竹竿的头，天还在上头。他只看得见对面肯让他看见的那一层。

胡帅。

弄出序列之毒的人，胡赞的长子。

真霜斗罗慢慢吐出一口气。他侧过头，声音压得只有碎紫斗罗一个人听得见。

``碎紫。''

``前辈。''

``这个人是胡帅。''真霜斗罗说，``Y(1,3,7,12)。''

碎紫斗罗的瞳孔一下子缩成了针尖。

``Y(1,3,7,12)\ldots\ldots{}''他把这几个字咬在牙缝里，``那比前辈的 Y(1,3,6)\ldots\ldots{}''

``高一大截。''真霜斗罗说，``真动起手来，老夫不是他的对手。''

碎紫斗罗的脸白了。

他刚才在三十丈外看着真霜斗罗铺开领域，一个 Y(1,2,4,8,16) 的人，连手都没让他出就跪下了。在他心里，那就是人能走到的头了。

现在这个人亲口说，打不过。

``那\ldots\ldots{}''碎紫斗罗的声音发涩，``怎么办？''

真霜斗罗没马上答。他的目光往后扫了一下，扫过八个孩子，扫过担架上的大帝，又扫过广场上那一片刀枪。

最后落回胡帅身上。

胡帅还是那么站着，蟒袍让夜风吹得往后鼓，冠上的玉珠子在火光里晃。他一点都不急，像站在自家园子里，等客人把茶喝完。

他用不着急。

``碎紫。''真霜斗罗说，``带他们，带着陛下，从左边冲出去。老夫拖住他。''

``前辈\ldots\ldots{}''

``别废话。''真霜斗罗身上的冰蓝光一下子亮起来，比刚才在宫门口亮得多，``老夫能撑三十息。三十息以内，你们冲出去。''

碎紫斗罗的嘴唇抖了两下，最后重重地点了一下头。

``走！''

他转过身，紫光一冲，领着八个人跟担架往广场左边扑过去。左边那一片的甲士站得稀一些，有一道矮墙，墙后头是一条窄巷子。

真霜斗罗一个人留在了宫门前。

他对着三千甲士，二十几个魂斗罗，还有那个穿着亲王皮的 Y(1,3,7,12)，慢慢抬起了两只手。

冰蓝的光从他身子里涌出来，涌得一点不留。宫门前的冷一下子往下坠，坠得比刚才在这儿压那个黑甲将领的时候还深。半空里的水汽一齐冻住，冻成满天的冰碴子，在火光里亮成一片。

冰霜领域，全开了。

胡帅脸上的笑，到这时候才收起来。

他隔着满天的冰碴子看着真霜斗罗，看了一会儿，眼底下闪过一点东西，像是赏识。

``真霜斗罗。''他开口了，声音一下子变了，亲王那副温吞吞的嗓子没了，换成了一把又低又厚的嗓子，沉得像庙里的钟，``Y(1,3,6)，也敢挡在 Y(1,3,7,12) 跟前。''

他抬起右手，手心里亮起了一团暗金的光。

``你值得老夫认真一回。''

\begin{center}\rule{0.5\linewidth}{0.5pt}\end{center}

碎紫斗罗领着人往左边冲，冲出去二十来丈，奥斯卡站住了。

``小奥！''蔡须坤在前头回过头来，``你干什么？走！''

奥斯卡没答。他的手伸进腰上的袋子里，在一堆香肠里摸，摸了两下，摸出来一根。

那根香肠比他平常给人吃的粗一圈，外皮上爬着一层细细的金色纹路，纹路一明一暗，像在喘气。香味一出来，连担架边上的黑子鸡都把脖子伸过来了。

``你们先走。''奥斯卡说。他的声音平得很，一点都不像他平常那个嬉皮笑脸的样子，``三十息不够。真霜长老拼了命，也拖不了那么久。''

蔡须坤看着他，嘴张了一下。

``小奥\ldots\ldots{}''

``超级大香肠。''奥斯卡说，``金纹的，我拿魂力一遍一遍地熬了三天。吃下去，一时半会儿能把人的魂力顶到一百五十成。''

他又从袋子里摸出一根，一模一样的，两根并在手里。

``金纹的，就这两根。''他说，``给两位前辈。''

唐三接过去一根，低头看了看。香肠在他手心里是热的，金纹一跳一跳地亮，里头那股魂力冲得他手腕子发麻。

``一百五十成\ldots\ldots{}''唐三抬起头，``你的魂力，够同时给两个封号斗罗顶上去？''

奥斯卡咧了一下嘴，笑得比哭还难看。

``够不够，''他说，``都得够。''

他把那根香肠从唐三手里拿回来，两根攥在一块儿，身上的魂力云一层压一层地翻了起来。魂力从他身子里往外倒，倒得他的脸一下子没了血色。金色的光从他手心里冒出来，拧成两股，一股缠一根香肠。

``去！''

他一扬手，两根香肠飞了出去。

两道金光从夜空里划过去，一道往宫门前那片冰碴子里飞，一道往前头碎紫斗罗的后背上飞。

\begin{center}\rule{0.5\linewidth}{0.5pt}\end{center}

真霜斗罗听见了后头的风声。

他没回头，右手往后一探，一把攥住了那根裹着金光的香肠。

香肠一碰着手心，热乎乎的魂力就从手心里往经脉里钻。他的魂力本来是冷的，冷得跟冬天的河一样，冻了一百多年，流得又稳又慢。这股热一进去，河面上的冰一下子化开了，底下的水哗地往前冲，冲得比原来快了好几倍。

一百五十成。

真霜斗罗的眼睛亮了一下。

他的 Y(1,3,6)，让这股热顶着，往上硬生生地撑开了一层。

Y(1,3,6,12)。

他身上的冰蓝光一下子浓了好几倍，冰霜领域往外一扩，扩出去一倍远。广场上那些甲士的脚底下，冰一层一层地往上长，长到了膝盖。

前头，碎紫斗罗也接住了那一根，看都没看，一口塞进嘴里，嚼了两下咽下去。金光从他喉咙底下往全身走，他身上那一串 Y(1,3,4,2,5,8,10,5) 一阵乱颤，第六节的 8 让那股劲儿顶着，往上跳成了 9。

Y(1,3,4,2,5,9)。

他身上的紫光翻起来，翻得像一条龙。

他站住了，回头看了一眼。

宫门前，真霜斗罗一个人站在冰碴子里，对面是那个穿蟒袍的人。

碎紫斗罗把脚底下的方向一拧。

``唐三。''他说，``你们护着陛下，往巷子里走。别回头。''

唐三张嘴想说什么，碎紫斗罗已经不在那儿了。

真霜斗罗听见身边落下了一个人。

碎紫斗罗站在他左边，身上的紫光跟他身上的冰蓝光碰在一块儿，嗡嗡地响。帝国大将军的脸绷得很紧，嘴角却还挂着一点苦笑。

``前辈。''他说，``碎紫不才，陪前辈打这一架。''

真霜斗罗侧过头，看了他一眼。

``你疯了。''他说，``Y(1,3,7,12)，你我两个加起来，也未必\ldots\ldots{}''

``未必打不赢。''碎紫斗罗把他的话截住了，右拳上紫光一点一点往里收，``打不赢，拖他三十息，总拖得住。''

真霜斗罗没说话。

过了一会儿，他点了一下头。

``那就一起。''

\begin{center}\rule{0.5\linewidth}{0.5pt}\end{center}

两个人一齐动了。

真霜斗罗在前，右手往前一推。Y(1,3,6,12) 的魂力从他手心里推出去，推成一整面冰蓝的光墙，光墙往前碾，碾过去的地方，空气冻成了冰，冰又碎成了刀片子，亿万片冰刀子裹在光墙里，一齐往胡帅身上压。

碎紫斗罗跟在后头，两只拳头同时砸出去。Y(1,3,4,2,5,9) 的魂力化成一只紫色的大鸟，两只翅膀张开几十丈，从光墙的缝里钻过去，直扑胡帅的脸。

胡帅脸上一点笑都没了。

他手心里那团暗金的光一下子炸开了。Y(1,3,7,12) 的魂力从他身子里往外涌，涌成一面暗金的大盾，竖在他身前。

三样东西撞在了一块儿。

撞上的那一下，一点声音都没有。

声音太大了，大过了耳朵能接住的那个数，耳朵里只剩下一片空。

然后才响。

轰的一声，从三个人当中炸开。宫门前的青石地往下一塌，石板一块一块地掀起来，在半空里碎成粉。气浪往四面推出去，推倒了一片又一片的甲士，三千人里，离得近的那几百个让它掀起来，像秋天的叶子一样往后飘。那二十几个魂斗罗也站不住，一个个吐着血往后倒飞。

等烟尘往下落，撞的那个地方，多了一个大坑。

坑有几十米宽。坑壁上的石头一半化了，一半冻着，冻住的那一半上还挂着烧红的渣子。坑沿还在往里塌，碎石子哗哗地往下滚。

气浪一直推出去几百丈才散。

左边的巷子口，八个人让气浪推得东倒西歪。蔡须坤扑在担架上，拿瑞兽本源把大帝整个罩住了。奥斯卡一条腿跪在地上，两只手撑着地，嘴唇白得发青。同时给两个封号斗罗顶魂力，他身上的魂力差不多让抽干了。

奥斯卡抬起头，往坑那边看了一眼，嗓子里挤出一句：``值了。''

\begin{center}\rule{0.5\linewidth}{0.5pt}\end{center}

坑底下，尘土一点一点地散开。

三个人站在里头。

胡帅站在坑正当中。蟒袍碎了一大半，露出底下贴身穿的一件暗金的内甲。内甲上全是细细的裂纹，暗金的光从裂纹里往外渗，像一只快要破的蛋。他嘴角挂着一道血，脸色还是稳的。

他对面，真霜斗罗跟碎紫斗罗并排站着。

真霜斗罗的长袍让气浪撕开了好几道口子，露出两条又白又瘦的胳膊。他嘴角也挂着血，眼睛里那两点冰蓝的光一点都没暗。

碎紫斗罗最惨。身上的紫甲碎了大半，左胳膊吊在身边，往外拐着，撞的时候折了。他的右手还攥着拳，拳头上紫光一明一灭，眼睛瞪着胡帅，像一头让人逼到墙角的狼。

三个人对着，谁都没先动，谁都在喘。

``平了。''胡帅抬手擦了一下嘴角的血，眼睛里有一点意外，``Y(1,3,6,12)，再添一个 Y(1,3,4,2,5,9)，再添上那一手增幅\ldots\ldots 倒是老夫没算到。''

真霜斗罗没答话。

他心里清楚，刚才那一下是平了，可这个平是借来的。那根香肠顶上来的劲儿有时候，劲儿一退，他就从 Y(1,3,6,12) 掉回 Y(1,3,6)。到那时候，再对上 Y(1,3,7,12)，他一招都接不住。

得快。

他侧过头，跟碎紫斗罗对了一眼。

碎紫斗罗看懂了。

``前辈。''他把右拳上的紫光又攥紧了一点，``再来一回。''

真霜斗罗点了一下头。

冰蓝的光跟紫光又在坑底下亮起来。

巷子口，八个人看着坑那边，蔡须坤刚把担架往巷子里推了两步。

就在这时候，宫门里头，有样东西醒了。

先是冷。真霜斗罗的冰是一种冷，那种冷是干净的，冷到底了，是透明的。这一种冷带着潮气，从宫门深处一点一点地漫出来，贴着地皮走，像深井底下的水往上返。

然后是黑。

宫门里那一片灯火，一盏接着一盏地灭了。灭得很快，从最里头的寝宫往外灭，灭过二门，灭过前殿，一直灭到宫门口。灭掉的那一片，黑得看不见东西，连宫墙的影子都没了。

那一片黑从宫门里漫出来，漫上了广场。

蔡须坤身子里的瑞兽本源尖叫了一声。

那一声在她骨头里响，又尖又细，像一只小兽让火燎了毛。她的脸一下子白得没了血色，后背上的汗毛一根一根地竖起来。

``那是什么\ldots\ldots{}''唐三的嗓子发干。他的蓝银草往宫门那边探过去，草梢刚碰着那片黑，就蔫了，一截一截地发黑，往回缩。他把草收回来，手心里全是汗。

霍雨浩的精神力撞了上去。

撞上的那一下，他觉得自己一头撞在了一堵铁墙上，墙是黑的，冰凉的，往外渗着水。精神力让墙弹回来，弹得他往后一仰，鼻子底下一热，两道血淌了下来。

``不对\ldots\ldots{}''他拿袖子抹了一把鼻血，眼睛还瞪着宫门，``这东西，比我在星斗大森林最里头碰上过的哪一头魂兽都凶\ldots\ldots{}''

黑一路往外漫。火把碰上它，噗地灭了。地上的冰碰上它，化成水，水又冒着白烟干了。脚底下的青石板让它舔过去，发出嘶嘶的响，像泼了一层酸。

八个人一步一步地往巷子里退。

那片黑的最里头，有一个人慢慢走了出来。

明黄的太子常服。紫金冠。腰上那块龙纹玉佩。

太子。

蔡须坤的瞳孔一下子缩紧了。

还是那张脸，二十五六岁，端正好看。可那个温温和和、礼数一点不差的人没了。这个人走路的样子是歪的，肩膀一高一低，脑袋往前探着，像一头刚从窝里钻出来的东西，在闻空气里的味儿。

他的眼睛全成了灰白色。

没有黑眼珠，没有眼白，两只眼眶里只剩下一片灰白，空的，像两颗泡了很久的鱼眼珠子。灰白的光从眼眶里往外溢，溢成两点鬼火，在黑地里飘。

黑从他身子里往外冒，冒出来就转，在他身边转成一个大旋涡。旋涡里头，隐隐约约有一张一张的脸，扭着，张着嘴，没有声音。

``胡峰。''

坑底下，真霜斗罗开口了。

他认得这股黑。这股黑底下藏着的那张脸，他也认出来了。

胡峰。胡帅的亲侄子，胡赞次子胡继的儿子。打小让胡帅带在身边，胡家最里头那一套黑的序列，传到了他身上。胡帅阴，躲在后头使劲儿。胡峰暴，什么都摆在明面上，下手又狠又急，是胡帅手里最快的一把刀。

叔侄两个。一个穿上亲王的皮，在宫门外头堵着。一个穿上太子的皮，在宫里头藏着。

真霜斗罗这才把这一整盘看全了。他们要的是天斗皇室。皇帝攥在谁手里，这座城就跟谁姓。

``你们觉着，''胡峰开口了。太子那副温润的嗓子没了，换成了一把又低又哑的嗓子，像铁在石头上拖，``把孤从宫门口吓回去，这事就完了？''

他的灰白眼睛从八个人身上扫过去，一个一个地扫，扫得很慢。扫到巷子里那副担架上，停住了。

他的嘴角往上一挑。

``果糕小队\ldots\ldots{}''他把这四个字含在嘴里，嚼了一遍，像在尝一口好菜，``你们踏进天斗城那一刻，就已经在我这盘棋上了。''

他往前走了一步，那片黑跟着他往前漫了一丈。

``我要的东西，在你们后头。''他说，眼睛还钉在那副担架上，``把担架上那个人留下。剩下的，可以走。''

``他要陛下。''唐三压着嗓子说。

``天斗的皇帝，活的，比死的值钱。''胡峰说，``死了的，还得再找一个来坐那把椅子。''

巷子口一下子静了。

蔡须坤没动。她身子里那只小兽还在尖叫，叫得她牙根发酸。她两只手还按在担架上，按得指头都白了。

唐三往前跨了一步，挡在了担架前头。

接着是霍雨浩，鼻子底下的血还没擦干净。小舞从左边站过来，诺当从右边站过来。卢本韦把担架那头交给蔡兴坤，自己走到最前头，袖子里那块猪神令攥在手心里，攥得发烫。奥斯卡还跪在地上，他撑着地，晃了两下，也站起来了。

谁都没说话。

胡峰看着他们，看了一会儿，笑出了声。

``好。''他说，``我叔父拖住那两个老东西。至于你们\ldots\ldots{}''

他的灰白眼睛一下子缩紧了。黑从他身子里猛地往外一涨，涨得那个旋涡转得飞快，里头那些脸一齐张大了嘴。

``我来把你们，彻底击碎。''

黑像一堵墙一样塌了下来，塌成无数条又粗又长的黑触手，朝着巷子口的八个人扑了过去。

\begin{center}\rule{0.5\linewidth}{0.5pt}\end{center}

\chapter{第七十章：帅气本源}\label{ux7b2cux4e03ux5341ux7ae0ux5e05ux6c14ux672cux6e90}

黑触手朝巷子口扑过去的那一下，真霜斗罗扭过头，看的却是坑当中的胡帅。

他的瞳孔一下子缩紧了。

刚才三个人撞那一下，他只当胡帅身上的东西是序数的另一种走法。可胡峰身上那股黑一冒出来，他就觉出来了。胡峰身上的黑，跟胡帅身上那层暗金的光，是一个窝里出来的。

再往深里摸，胡峰那点只是一根枝杈。根长在胡帅身上。

``帅气本源。''真霜斗罗说，声音沉下去了，``原来是这个。你修的那一套，根子上就跟序数两样。''

碎紫斗罗扭过头。``前辈，什么叫帅气本源？''

``比序数还老的一路。''真霜斗罗说，``它不靠天地的规矩转，拿自己身上的帅气当根，自己转自己的。帅气越足，力气越大。胡帅这一身帅气\ldots\ldots{}''

他的目光又落回胡帅身上。

``已经黑了。''

胡帅身上那点帅气，早就让什么东西泡过了，泡得发黑，发苦，拧成了另一副样子。它像一口没底的井，什么东西靠近，它就往里吸。

真霜斗罗身上的冰，正让它一丝一丝地往过吸。

他的领域刚才还压得住那个黑甲将领，冻得住序数的规矩，这时候碰上胡帅，冻不住了。冰蓝的魂力像一条河碰上了漩涡，不由自主地往胡帅那边歪，歪过去，散了。站在胡帅跟前，像站在一头大兽的嘴边上，它也不咬，就那么张着嘴，一口一口地吸。

``见鬼。''碎紫斗罗也觉出来了。他身上的紫光一阵乱颤，像让人掐住了脖子，``本座的魂力\ldots\ldots 在往外漏！''

``他在吃。''真霜斗罗的声音一下子硬了，``不能再拖了。开真身！''

\begin{center}\rule{0.5\linewidth}{0.5pt}\end{center}

扽魂真身。

封号斗罗压箱底的东西。

真霜斗罗身上腾起了一大团冰蓝的魂力云，一层压一层，翻着，越翻越大，翻到后来拧成了一根冰蓝的光柱，直冲到夜空里去。

光柱里头，走出来一头熊。

一头冰熊，身子有一百五十多米长。浑身披着一层半透明的冰甲，冰甲一片压一片，每一片上都刻满了细密的冰蓝纹路。两只前掌各有一座小山那么大。它喘一口气，鼻子里喷出来的就是一阵冰碴子风。两只眼睛是两点冰蓝的亮光，看到哪儿，哪儿的水汽就冻成冰棱子，噼里啪啦往地上掉。

远古冰熊王。

真霜斗罗站在冰熊王的头顶上，长袍让冷风吹得往后扯。熊是他，他是熊。熊抬一抬爪子，他手指头就跟着动一动。

Y(1,3,6,12) 的魂力灌进了冰熊王的身子里。冰熊王仰起头吼了一声，吼声里裹着冷，方圆百丈的地面上一下子又结了一层冰，冰底下的石头冻得往上拱。

碎紫斗罗的真身，是另一个样子。

一团紫色的魂力云在他头顶上翻，越翻越急，翻成了一道紫色的闪电，把他整个人裹在里头。

他在闪电里往小里缩。

闪电一散，碎紫斗罗不见了。原地站着一个三尺来高的小人，通身都是紫色的光，光稠得往下淌，像一团压实了的水银。

极速模式。

碎紫斗罗的真身不往大里长，往快里走。在这个样子底下，他快得眼睛跟不上，光都追不上他。

``前辈，本座去缠住他！''

小人里传出来一声喊。

喊声还没落，原地就空了。地上留着一道紫色的影子，一圈让风刮起来的碎冰。

\begin{center}\rule{0.5\linewidth}{0.5pt}\end{center}

胡帅站在坑底下，看着跟前一头一百五十米的冰熊，一道快得看不见的紫光，嘴角往上挑了一下。

``真身？''他说，``好。老夫也认真一回。''

他抬起两只手。

身上剩下的那半件蟒袍，让他身子里冒出来的东西一冲，化成灰飞走了。底下那件暗金的内甲咔咔地往开裂，裂缝里往外冒的光，换了颜色。暗金褪下去，冒出来的是一种干干净净、亮得扎眼的光。

帅气之光。

``帅气真身。''

胡帅的声音一下子空了，远了，像从天上掉下来的。

他身上开始发光。先是手心里一团，拳头那么大。一口气的工夫，长成了磨盘那么大。再一口气，几丈宽了。光的颜色从暗金往白里走，白到了头，边上又泛出一圈金。

胡帅在那团光里慢慢往上浮。头上的冠化成了一圈光，冠上垂着的那些玉珠子一颗颗飞起来，变成一个个小光球，绕着他转。他的脸在光里看不清了，只剩下一个亮着的轮廓。

Y(1,3,7,12) 的帅气本源，在真身里放开了。刚才那往里吸的黑劲儿收起来了，换成了往外打的东西。

``去。''

胡帅只说了一个字。

他手心里那团光往前一喷，喷出来一道光，十几丈粗，白得刺眼，直直地朝冰熊王射了过去。

光走过的地方，天上让它烧出一道黑缝子。半空里的冰碴子还没碰着光，就化成了白汽。连光自己都在往弯里拐，像是扛不住自己身上那么大的劲儿。

冰熊王吼了一声，一百五十米的身子往前一扑，抬起前掌想挡。

光太快了。

前掌还没抬到一半，光就打在了它的右腿上。

``吼！''

冰熊王疼得一声长嚎。那道白光像烧红的铁签子扎进了一块黄油，几丈厚的冰甲一下子就穿了。冰甲上那些冰蓝的纹路一闪，一碎，一灭，冰层在光里化成了水，水又化成了一大团白汽，往天上冒。

右腿一软，冰熊王往右边歪了下去。它的膝盖砸在地上，地面跟着一阵乱颤，裂开一道几十丈长的口子。

真霜斗罗在熊头上晃了一下。

腿底下一阵钻心的疼，从熊的右腿一直传到他自己的膝盖上。他的瞳孔缩成了一个针尖。

光还在烧。熊腿要是挪不开，整条腿都得烧穿。熊一倒，站在熊头上的他就整个露在胡帅跟前。

来不及想了。

``冰暴突击！''

真霜斗罗两只手往前一推，身上 Y(1,3,6,12) 的魂力一点没剩，全灌进了冰熊王的身子里。

冰熊王的两只眼睛一下子亮得刺眼。

它没站起来。它拿右膝盖顶着地，一百五十米的身子往前一弹，整个弹了出去。

轰的一声，它起跳的地方塌下去一个上百米的大坑。泥，石头，冰，一齐让它掀上了几百丈高的天。冰熊王像一座倒下来的冰山，可这座山快得吓人，往前冲出去，一眨眼就是几十丈。

白光在它身后划了一道长长的焦痕。右腿上烧穿的冰甲在冲的时候一片一片往下掉，露出底下冰蓝色的肉，肉上全是血。它顾不上，四只爪子拼了命地往地上刨，一爪子一个坑，几口气的工夫就冲出去几百丈，冲出了那道光。

光扑了个空，射到天上去了，把夜空烧出一道好几里长的焦印子。

冰熊王在几百丈外刹住了。四只爪子插进地里，往前犁出四道几十丈长的深沟，才停下来。它的右腿在抖，伤口上冒着白汽，冰蓝色的血一滴一滴往下落，落在地上嘶嘶地响。

真霜斗罗站在熊头上，脸色很难看。

冰暴突击把那一下躲过去了，可熊腿伤得重。往后冰熊王跑不快了，也使不出全力了。

他没退。

他退不了。八个孩子还在胡峰手底下，碎紫还在缠着胡帅。他这一退，什么都散了。

他蹲下来，一只手按在熊的脑袋上，冰蓝的魂力从手心里淌出去，顺着熊的身子往下走，走到那条伤腿上，在伤口上结了一层薄冰，把血先封住了。

``老伙计。''他说，``还撑得住么？''

熊没马上应。它把脑袋往下低了低，拿鼻子在他脚边拱了一下，拱得他身子一晃。这头熊跟了他一百多年，他头一回进扽魂殿那年，它还没他腰高。

冰熊王低低地吼了一声，四只爪子往地上一撑，又站起来了。右腿有点跛，可站得稳稳当当。它仰起头又吼了一声，比刚才那一声还响。

真霜斗罗心里有数。Y(1,3,6,12) 对 Y(1,3,7,12)，中间隔着一大层。照常理，这一层是翻不过去的。

可他用不着赢。

他只要拖。

``碎紫！''他喊了一声，声音穿过满场的轰响，``放开了打！一点别留！''

远处那个紫色的小人一闪。

碎紫斗罗听懂了。打不赢胡帅，那就死死缠住他，让他腾不出手去帮胡峰。缠一下少一分血，也得缠。

``明白！''

小人身上的紫光往里一收，又往外一炸。极速模式让他催到了头，三尺高的小身子上一下子爬满了细细的裂纹。那是快过了头，身子扛不住了。他顾不上，化成一道紫线，又朝胡帅扑过去。

真霜斗罗也动了。冰熊王拖着伤腿，从另一边往胡帅身上压过去，两只前掌高高举起，掌心里各攒起一团冰蓝的光，五丈多宽，冷得能冻住一座山。

一个在前头绕，一个从上头砸。

胡帅的眉头动了一下。

他觉出来了。这两个人跟刚才不一样了，刚才是试，现在是拼。碎紫的那个小人身上裂纹越来越多，每打一拳，都在拿自己的命往里填。真霜站在一头瘸了腿的熊头上，把 Y(1,3,6,12) 一丝一丝地往外榨。

``拼了？''胡帅说，声音平平的，眼底下那点暗金闪得快了些，``好。老夫看看你们能撑多久。''

他身上的光一炸，炸成了一片一片的光刃，又薄又快，一半朝冰熊王飞，一半朝碎紫斗罗飞。

三样东西又撞在了一块儿。

光把整片宫门前照得跟白天一样。气浪横着扫过去，把几百丈以内本来就塌了的房子又削平了一遍。

碎紫斗罗让气浪掀飞了，小人身上的裂纹又多了几道。他咬着牙在半空里一拧，站稳了，又是一道紫线扑回去。

冰熊王的左肩让一片光刃劈中了，冰甲碎了一大片，露出底下冰蓝色的肉。它吼着，一步不退，一座山似的堵在胡帅跟前。

胡帅的真身是强，Y(1,3,7,12) 也是实打实地压着他们两个。可这两个人配合得太熟了。碎紫快，绕着他转，打乱他出手的节奏，他那道光老是锁不住人。真霜慢，可那头熊太大了，往前一堵，他往哪边挪都有一堵冰墙。他想甩开这两个人去帮胡峰，甩不开。

他一连出了几十手，一手都没落实。

胡帅像一只让两条狼围住的老虎。哪一条狼单拎出来，都不够他一口吃的。两条狼一前一后，他就是走不了。

``好一个真霜，好一个碎紫。''胡帅的声音里头一回有了点分量，``拿两个封号斗罗的命，来拖老夫\ldots\ldots 你们倒舍得。''

可他还是拿不下来。

真霜斗罗站在熊头上，嘴角挂着一道血。魂力往外榨得太狠，身子里在往回反。他拿袖子擦了一下，看着胡帅。

``胡帅。''他说，``你的对手，是我们两个。''

\begin{center}\rule{0.5\linewidth}{0.5pt}\end{center}

巷子口，八个人正跟那些黑触手缠成一团。

担架让他们放到了巷子里头，唐三的蓝银草从墙缝里钻出来，把担架四个角钉在了地上，草叶子织成一张网，在大帝身上又罩了一层。

蔡须坤站在最前头。瑞兽本源从她身子里往外涌，涌成一面金色的光墙，把八个人跟巷子口一块儿罩住。唐三的蓝银草顺着光墙往外长，长成一张大网，一根黑触手扑过来，网就缠上去，缠住了往回拽。霍雨浩站在光墙后头，鼻子底下的血还没干，精神力一刻不停地往外摸，哪根触手要从哪边扑过来，他提前半口气就喊出来。

``左上！''

``底下，地底下有一根！''

``右边三根，一起来！''

可触手太多了，也太硬了。

胡峰身上的黑像没有底，一个劲儿地往外冒。蓝银草缠断一根，断口那儿就长出两根来，断两根，长出四根。光墙让它们一下一下地抽，抽得一阵一阵地晃。

``这么耗不行！''蔡须坤咬着牙，``咱们的魂力在往下掉，他那边一点都没少！''

话刚出口，胡峰的身子猛地一颤。

他身上又冒出来一样东西。

这一回冒出来的，跟刚才那股黑不一样。蔡须坤认得这个味儿。刚才胡帅在坑底下放出来的，就是这个。

帅气。

可它又跟胡帅的不全一样。

蔡须坤的瑞兽本源一阵乱跳。她看见了，胡峰身子里那套序列，在一节一节地往外亮。

0。

头一个数亮起来，胡峰身上的气一沉。

111。

三个一，一齐亮了。他身边那团黑一下子翻起来，跟那股新冒出来的帅气搅到一块儿，搅成了一样说不出颜色的东西，又黑，又亮。

222。

三个二跟着亮了。胡峰身上的气一下子涨上去，涨得人喘不过气来。他那两只灰白的眼睛里，竖着裂开了两道金色的缝，像两扇门让人从里头推开了一条缝。门后头，黑跟帅气搅在一块儿，看不见底。

0 111 222。

胡峰真正的东西，到这时候才拿出来。

胡帅身上的帅气，是让黑泡坏了的。胡峰身上的，是他自己把黑跟帅气捏到一块儿去的，捏成了一样新东西。

``胡家压箱底的，就是这个\ldots\ldots{}''霍雨浩的声音在抖。他的精神力刚碰上那股东西，就让它顶了回来，鼻子底下的血又涌出来了。

胡峰慢慢抬起两只手。

又黑又亮的东西在他两只手当中攒起来，攒成了一座山的影子。半透明的，通体黑，山上流着一道一道暗金色的纹路，像一座拿黑铸出来的山。它浮在胡峰头顶上，慢慢地转，转一圈，四周的空气就往里挤一分。地上的石板一块一块地碎。巷子口八个人只觉得身上压了一座看不见的山，膝盖一个劲儿地往下弯。

胡峰的嘴角往上一挑。

``泰山压顶。''

他两只手往下一按。

那座黑山像一块从天上掉下来的大石头，朝巷子口的八个人砸了下来。

``往我这儿聚！''蔡须坤嗓子都喊劈了，瑞兽本源在她身子里烧起来，金光往外喷。

几个人刚要把魂力往她那面光墙上堆。

``须坤！''

边上有人喊了她一声。

蔡兴坤。

她什么时候站到蔡须坤身边的，谁也没看见。她浑身裹着一层土黄色的魂力，厚得像披了一层泥甲。两只脚往地里一扎，脚底下的青石一下子化开，化成一层几尺厚的石头，顺着她的腿往上爬，爬成了一身石甲，跟她长在了一块儿。

她肩膀上的黑子鸡早就跳下来了，蹲在她脚后跟边上，一身毛奓着。

``融合技。''蔡兴坤说，土黄色的魂力从她身上往外淌，跟蔡须坤的金光碰在一起，缠在一起，``铁山靠！''

蔡须坤一下子就明白了。

她没犹豫。那面金色的光墙一下子换了样子，从挡在头顶上的一个罩子，变成了八个人身前的一堵墙，墙根扎进了地里。土黄的石头跟金色的瑞兽本源在墙里头搅匀了，搅成一层，又硬，又韧。

铁山靠。

蔡须坤跟蔡兴坤的融合技。瑞兽本源管韧，管补，哪儿碎了哪儿就长回来。地底下的力气管硬。两样合在一块儿，名字里那个靠字，说的就是这个：拿自己当山，拿地当根，上头砸下来什么，都顺着往地里卸。

轰。

黑山砸在了墙上。

巷子口的地一下子往下沉。方圆几十丈的地整块往下陷了好几尺，裂缝从墙根底下往四边爬。墙面让黑山压得往里凹，土黄的石头一层一层地碎，金光一明一灭。

泰山压顶的劲儿不光在砸。那股又黑又亮的东西从墙缝里往里渗，往每个人的脑子里钻，要把人心里那点硬气碾平了。蔡须坤嘴角流出了血。蔡兴坤两只脚陷进了地里，陷到了膝盖，膝盖以下的石头让她压成了粉。

墙一直没塌。

它在抖，满是裂纹，可就是没塌。蔡须坤跟蔡兴坤的魂力在墙里头拼了命地转，碎掉的石头让瑞兽本源一补，又长回来，旧的一层碎了，新的一层顶上去。黑山往下碾，墙往上补，两边就这么耗着。

一息，两息，三息\ldots\ldots{}

耗到第十息，黑山先撑不住了。

0 111 222 再凶，泰山压顶这一下也费力气。黑山上的暗金纹路暗了，山也一圈一圈地往小里缩。

黑山散了。墙也跟着轰的一声碎了，碎成一地的石头渣子跟金色的光点。

胡峰站在原地，没接着往下打。

那些金色的光点飘到他跟前，有一点落在他手背上，化了。他低头看了一眼手背，鼻子动了一下。

那光里有点味儿。他说不上来是什么，可他身子里那点黑帅气碰上它，自己往前探了一下，像一条饿了好几天的狗，闻见了肉。

他抬起头，往那堵碎墙后头看了一眼，眉头皱起来了。

八个人从石头渣子里站起来。

身上都挂了彩。蔡须坤嘴角带着血，蔡兴坤两条腿在打颤，唐三的袖子让碎石头划开了几道口子，霍雨浩鼻子底下的血还在往下滴。

一个都没倒。

蔡须坤拿手背抹了一下嘴角的血，身子里的瑞兽本源慢慢稳下来。

``挡完了。''她说，``该咱们了。''

\begin{center}\rule{0.5\linewidth}{0.5pt}\end{center}

卢本韦头一个动。

她两只手往地上一拍，魂力从手心里灌进地里。脚底下的地一阵乱颤，一道粉红的光从地缝里冲上来。

``出来！''

地底下拱出来一头母猪。

一头大母猪，粉红的皮，三丈高，五丈长，四条腿跟四根柱子一样，两颗獠牙往外翻着。两只眼睛血红，鼻孔里喷出两股粉红色的热气。它四个蹄子落地，方圆几丈的石板一齐碎了。

卢本韦的扽魂，母猪。

名字听着土，这头猪可一点都不好惹。皮糙肉厚，跟万年魂兽有得一比，撞一下就是千钧的力气。八个人里，往前顶的活儿，一向是它的。

母猪落了地，仰起脖子嚎了一嗓子，四个蹄子在地上一刨，朝胡峰冲了过去。

奥斯卡也动了。

他咧着嘴笑，手伸进袋子里，一把抓出来一大捆香肠。这一捆跟刚才那两根金纹的不一样，每一根外皮上都爬着赤红的纹路，捏在手里烫手，嘶嘶地往外冒热气。

爆裂香肠。

``本韦，往左！''

卢本韦一听就懂。她一把揪住母猪脖子上的鬃毛往左一拽，母猪冲着冲着猛地往左一拐，四个蹄子在地上刨出一道深深的弧。这一拐，正好把胡峰左边的路堵死了。胡峰要躲，只能往右躲。

右边，是奥斯卡那一捆香肠要落的地方。

``接着！''

奥斯卡一扬手，三十根香肠一齐飞了出去。红光从夜空里划过去，像下了一场流星雨，铺天盖地地往胡峰右边砸。每一根香肠在半空里都在胀，越胀越大，红纹越来越亮，嘶嘶地叫。

碰上什么，就炸。一根香肠里攒着的劲儿，抵得上一个魂尊拼了命的一击。奥斯卡一口气扔出去三十根。

母猪堵路，香肠收人。

这一手他们两个练过不知多少回。卢本韦管把人往哪边赶，奥斯卡管在那边等着。一个赶，一个打，一点缝都没留。

三十道红光一齐落到胡峰身边。

轰，轰，轰，轰\ldots\ldots{}

一连串的炸响在胡峰身边炸开，红火冲上了天。气浪把方圆几十丈的黑触手炸成了一截一截的碎渣子。母猪站在火边上，一动不动，粉红的皮上连一道白印子都没留下。

火慢慢往下落。

烟里头站着一个人。

胡峰身上那件明黄的太子常服烧没了一半，剩下的一半挂在身上，边上还冒着火星子。脸上蹭了几道黑，左边的头发燎掉了一绺。他抬手在肩膀上拍了两下，把火星子拍灭了，拍得不紧不慢，像是在掸灰。

三十根香肠，就换来这么点东西。

``啧。''奥斯卡咂了一下嘴，``我熬了一个月的。''

胡峰没看他。他的灰白眼睛转过去，落在了那头母猪身上，上上下下打量了一遍。

``一头猪。''他说，``胡家的人，让一头猪拦了道。''

母猪鼻子里喷出两股热气，前蹄在地上刨了两下。

卢本韦趴在猪背上，手里揪着鬃毛，看着胡峰，说：``它不光会拦道。''

胡峰笑了一下，笑得嘴角往一边扯。

巷子里头，小舞跟诺当一左一右守着担架。担架让蓝银草钉在地上，大帝脸上那层霜方才让黑山一压，又裂了一道缝，蔡须坤的瑞兽本源分出来一缕，一直罩在那上头。小舞蹲在担架边上，两只手攥成拳头，眼睛钉着巷子口，脚尖一下一下地点地，随时要往外蹿。诺当站着没动，指头尖上的电光越跳越密，噼噼啪啪地响，在他手背上爬来爬去。

``还不到时候。''诺当说，不知道是说给小舞听，还是说给自己听。

唐三站在巷子口，蓝银草缠在他手腕子上，他一直盯着胡峰身子里那套序列。

0 111 222。

他在心里过了一遍。这一串后头，还有空着的地方。胡峰把它亮到 222 就停了，没往下亮。

是亮不了，还是不肯亮，唐三看不出来。

霍雨浩往前走了一步，站到了唐三边上。

他鼻子底下的血还在淌。方才撞那堵黑墙，又顶了一回泰山压顶，他的精神力已经去了一大半，脑袋里一阵一阵地发空，像一口让人舀干了的缸。

小舞在巷子里看见了，站起来喊了一声：``雨浩！''

``够。''霍雨浩没回头。

他把袖子里那块帕子掏出来，在鼻子底下狠狠按了一下，又塞回去。

然后他的眼睛变了颜色。

两只眼珠子成了深深的冰蓝。他的精神力一下子全放开了，从他身上往外涌，涌得又冷又宽，在他身前攒成了一道看不清的光门。

光门里头，有两个影子慢慢浮了出来。

头一个，是一只蝎子。

通体冰蓝，个头有一座小山那么大。身上披着半透明的冰蓝甲壳，两只大钳子像两把冰做的大刀，一开一合，带起一阵刮骨头的冷风。尾巴高高地翘着，尾钩尖上挂着一滴冰蓝的水珠，晶亮晶亮的。那一滴是极寒的毒，滴进一片湖里，一片湖都能冻实了。

冰帝。

第二个，是一个披着白斗篷的姑娘。

脸很清冷，好看得让人不敢多看。一头银白的长头发披下来，发梢上亮着幽幽的蓝光。两只眼睛是冰蓝色的，眼底下封着一大片雪地。她光着脚，踩在半空里，一步迈出去，脚底下的空气就结成一片冰蓝的雪花。

雪帝。

霍雨浩的两个魂灵，是伊莱克斯给他的。

两个影子从光门里一跃出来，一左一右，落在胡峰两边。

就在这时候，蔡须坤闭上了眼睛。

她站在原地没动，瑞兽本源在她身子里转得飞快。她脑子里那片地方，叫精神识海的，让人从里头撞开了。

识海最深处，传出来一声狼嚎。

那一声又低又沉，不像这世上该有的声音。它从很远很远的地方来，从人还没有的年月里来。巷子口每一个人听见了，骨头里都跟着打了一个颤。坑底下的胡帅听见了，手上慢了半拍。绕着他转的碎紫斗罗听见了，那道紫线也顿了一顿。

蔡须坤的识海深处，一头大兽睁开了眼。

灭世狼王。

十九万年的凶兽。

它从蔡须坤的识海里走出来，像一道灰色的闪电，把眼前的夜撕开了一道口子，落在了巷子口。

一身灰黑的毛，每一根都硬得跟钢针一样，毛跟毛的缝里流着暗红的光，像石头缝里淌着的岩浆。四只爪子是四团黑，黑得看不见底，踩下去，地上就留下一个看不见底的爪印。两只眼睛是两团烧着的暗红火苗。

十九万年的修为，一点没藏，全放了出来。

狼王在当中，冰帝在左，雪帝在右。三头东西站成一个品字，把胡峰围在了当中。三股威压从三边一齐往下压，压得胡峰身边剩下的那点黑一个劲儿地往里缩。

冰帝先动的。

它两只大钳子猛地一合，钳尖上射出一道冰蓝的寒光，贴着地皮，直奔胡峰的两条腿。寒光扫过去的地方，石板一下子冻住，白霜顺着地面往四边爬。

胡峰脚底下一蹬，往后躲。

他的后路，雪帝已经站在那儿了。

她没动手。她就抬起一只手，手心朝外，轻轻一推。

胡峰身后的半空里，一下子落下来一堵雪墙。雪是一片一片的，轻得像鹅毛，可落到一块儿，就冻成了一整块，冻得严丝合缝，比石头还硬。

胡峰的后背撞在了雪墙上。

墙没碎。他让墙顶得往前一栽，脚底下一个趔趄，脸上头一回没了笑。

他站住了，左右看了一眼。左边是那只山一样的蝎子，钳子还张着。右边是那个光着脚的姑娘，站在半空里，低头看着他，一点表情都没有。

前头，那头灰黑的狼慢慢低下了头。

它喉咙里滚出一串很低的响动，咕噜咕噜的，像远处在打雷。两团暗红的火苗眯成了两条缝，前爪从地上抬了起来。

爪子上那四团黑，一点一点地，往长里伸。

\begin{center}\rule{0.5\linewidth}{0.5pt}\end{center}

\chapter{第七十一章：震天斧}\label{ux7b2cux4e03ux5341ux4e00ux7ae0ux9707ux5929ux65a7}

雪帝那堵雪墙还竖在胡峰身后。墙面上一层一层地往外冒冰蓝的纹路，冷气从纹路里渗出来，渗到胡峰的后背上，他后脖颈子上的汗毛结了一层霜。

前头是灭世狼王，左边是冰帝，后头是雪帝。

胡峰让三个魂灵夹在了当中。他身上那套 0 111 222 让三股威压一挤，往里缩，身上那层暗金的光一明一灭，明的时候越来越短。

狼王没给他喘气的工夫。

它没蓄力，没往后退两步助跑，连吼都没吼。

它就抬起了右前爪。

往下一挥。

裂空。

一点声音都没有，一点光都没有。

就一道爪印。

那道爪印从狼王的爪尖上拉出去，横着拉过几百丈，拉过的地方，半空里留下了一道黑缝。缝的边是碎的，像一块布让人从当中撕开了，撕开的口子后头，什么都没有，黑漆漆一片。

爪印走过去，胡峰身边那些黑触手一根一根地断了，断得一点声音都没有，像拿热刀子划黄油。

胡峰的瞳孔缩成了针尖。他身子里的 0 111 222 拼了命地转，想在身前竖起一层东西来。

来不及。

裂空走过的那条道上，地方早让它占死了。胡峰往哪边躲，那道爪印都在那儿等着他。

嗤。

爪印从胡峰的左肩膀上划进去，一直划进胸口。

暗红的血喷出来，在夜里拉出一条血线。胡峰让那股劲儿掀得往后飞。伤口边上，暗金的纹路疯了一样地闪，想把裂开的肉往一块儿拽。

拽不回来。裂空撕开的不光是肉。

胡峰的伤口里传出一声尖叫，又细又利，像一匹绸子让人从当中扯开。那一声是从他身子里头出来的，嘴都没张。他的脸一下子白得跟纸一样，灰白的眼睛睁得溜圆，眼角一跳一跳的。

``这\ldots\ldots{}''他的嗓子哑了，嘴角往外冒血，``十九万年\ldots\ldots 不可能\ldots\ldots 一个人，怎么可能跟十九万年的凶兽结契\ldots\ldots{}''

狼王没理他。

它站在原地，两团暗红的火苗低着，看着往后飞出去的胡峰，像看一只扑腾的耗子。

狼王身后，蔡须坤的脸白了一层。

裂空那一下，是从她身上抽的。狼王出一爪子，她身子里的魂力就让它拽走一大截，拽得她心口一空，腿肚子发软。瑞兽本源在她身子里转着，转得发烫，一边往外给狼王递魂力，一边还要分出一缕去巷子里罩着担架。

她咬着后槽牙站着，没让膝盖弯下去。

狼王回头看了她一眼。就一眼，两团暗红的火苗在她脸上停了停，又转回去了。

\begin{center}\rule{0.5\linewidth}{0.5pt}\end{center}

诺当一直守在巷子里的担架边上，一句话都没说过。

灰色的魂力在他身上转着，转得不声不响。他的两只眼睛很静，像两汪没有风的水。从胡峰从宫门里走出来的那一刻起，他就在等。

裂空落在胡峰身上的那一下，他动了。

他两只手在身前画了一个圈。

灰色的魂力从他指头尖上流出来，在他身前攒成了一个环，三丈多宽。环是拿无数道细细的灰纹路编起来的，纹路跟纹路当中跳着电弧，嗡嗡地响。

磁暴环。

诺当练得最熟的那一手。

环一成形，就朝胡峰飞了过去。它一边飞，一边转，越转越快，转的时候往外甩电弧，一道一道又粗又长，灰色的，像一条条灰蟒在半空里甩尾巴。

电弧扫到哪儿，哪儿剩下的黑就碎了。那些还没断的黑触手让电弧一抽，跟让雷劈了的枯枝一样，焦了，黑了，断了，化成灰。

嗡！

环飞到胡峰跟前，猛地一张，从三丈张到了三十丈，把胡峰整个罩在了当中。

几百道电弧从环的里圈一齐往里射，射成一只笼子，把胡峰锁在里头。电弧打在他身上，噼啪噼啪地炸，他身上那层暗金的纹路让电一打，一闪，一碎，一灭。

胡峰在环里吼了一声。

他的身子让几百道电弧抽得一抽一抽的，0 111 222 让环压着，又往里缩了一截，身上的光暗得像快烧完的蜡烛头。

裂空撕开了他的肉，磁暴环锁住了他的劲儿。

胡峰在环里挣扎着抬起头。他的灰白眼睛隔着电弧，往巷子口那几个人身上看。

母猪堵在正面，鼻孔里还冒着热气。奥斯卡站在侧面，手又伸进了袋子里。狼王在头顶上，冰帝在左边，雪帝在后头，磁暴环在身边转。

他往哪边看，哪边都有东西。

刚才在他眼里，这几个人跟一窝蚂蚁差不多，一脚就能碾平。这一窝蚂蚁这时候拿着他看不懂的东西，把他一步一步逼到了墙角。

``不可能\ldots\ldots{}''他的嗓子发哑，嘴里跟胸口上一齐往外冒血，``你们就是一群魂师\ldots\ldots 怎么可能\ldots\ldots{}''

蔡须坤站在狼王身边，瑞兽本源在她身子里稳稳地转。她隔着电弧，看着胡峰。

``胡峰。''她的声音不大，在电弧的噼啪声里听得清清楚楚，``你叔父这盘棋，下到这儿就够了。''

\begin{center}\rule{0.5\linewidth}{0.5pt}\end{center}

几百丈外，大坑那边，真霜斗罗站在冰熊王的头顶上，从头到尾都看见了。

他看着那头灰黑的狼，又看了看站在狼边上的那个丫头。

``十九万年的灭世狼王\ldots\ldots{}''他低低地说了一句，``这丫头身上，到底还揣着多少东西。''

他跟这个丫头在殿里待了这么些日子，教过她，骂过她，看着她在良序圣殿里一个数一个数地往下抠。他一直当她是个底子好、肯下苦功的孩子。他不知道她脑子里还蹲着这么一头东西。

冰熊王在他脚底下低低地哼了一声。右腿在抖，伤口上那层冰让热气一熏，化了一层，又渗出血来。真霜斗罗拍了拍熊脖子，把冰又补上，转过头去看胡帅。

胡帅还浮在半空里，一身白光，光里一张看不清的脸。

他也在看巷子口。

他看见胡峰让狼撕了一爪子，让电锁在了环里。看见那八个孩子一个都没倒。看见真霜跟碎紫两个人绕着他，咬着牙不松口。

然后他笑了。

这一回的笑，跟方才在宫门前那个笑不一样。方才是笑一个客人茶喝得慢。这一回，像一个下了几十年棋的人，棋盘上忽然冒出来一个他没算过的子儿，他非但不慌，反倒把身子往前探了探，要看个仔细。

``有意思。''胡帅的声音从白光里传出来，又低又远，``老夫布了这么多年的局，头一回碰上这样的岔子。''

他的目光从八个孩子身上扫过去，从真霜跟碎紫身上扫过去，最后停在那头狼身上。

``十九万年\ldots\ldots{}''他嘴角往上一挑，``好。很好。''

他手心里的光又攒起来了，比刚才还亮。

\begin{center}\rule{0.5\linewidth}{0.5pt}\end{center}

磁暴环里，电弧还在一下一下地抽胡峰。他左肩上那道爪印还在往外渗血，血顺着胳膊往下滴，滴到地上，冒着黑烟。

八个人刚松了半口气。

胡峰笑了。

笑声从电弧当中钻出来，又哑又低，一声连着一声，越笑越响。他的灰白眼睛在电光里一明一暗，血顺着下巴往下淌，他的嘴越咧越大。

``你们\ldots\ldots 是真强。''他从牙缝里往外挤字，一个字带一口血沫子，``我叔父看走眼了。你们几个，是正菜。''

他抬起了右手，手心朝上。

那一下抬得很慢，慢得每个人都看得清清楚楚。可没人上去拦。每个人的手指头都在那一下僵住了，骨头里有个东西在往后拽。

他手心里冒出来一团魂力，暗金色的，在他手里攒，一边攒，一边往外拉，拉出了一个形状。

蔡兴坤的瞳孔一下子缩成了针尖。

她的脸一下子白了。身上那层土黄色的魂力乱颤，颤得跟害了冷一样。那是从血里头往上冒的东西，比怕还早一步。她的嘴唇在抖，嗓子里挤出来一句话，干巴巴的：

``那个样子\ldots\ldots 不可能\ldots\ldots 那是\ldots\ldots{}''

她一把抓住了蔡须坤的胳膊，攥得指节都白了。

``须坤！退！快往后退！''

蔡须坤猛地扭过头。她从来没见过蔡兴坤这个样子。蔡兴坤一向稳得跟一座山一样，刚才泰山压顶压下来，她两只脚陷进地里陷到膝盖，脸色都没变过。这时候她的眼珠子一直往胡峰手心里瞟，瞟一下，攥着蔡须坤胳膊的手就紧一下。那团光攒出来的样子，她认得。在这儿，它不该出现。

``那是什么\ldots\ldots{}''

蔡须坤的话还没说完，胡峰手心里那团光定住了。

定成了一样东西。

一把斧头。

通体暗金。斧柄一丈来长，上头刻满了帅气的纹路。斧刃两尺多宽，弯得像半个月亮，刃口上淌着一层薄薄的金光，像化开的金子顺着刃口往下流。斧头两边各嵌着一颗暗红的石头，石头里头有一个小小的旋涡，一直在转。

蔡兴坤的脸上一点血色都没了。

震天斧。

她认得这把斧子。

胡家正根上的人才有的扽魂，一代传一代，斧身里压着黑帅气本源最里头的那一点。拿着震天斧的胡家人，同一个层次里，几乎没人打得过。这把斧子砍下去，能把人一身的魂力压成一根针那么细，再一口气放出去，砍什么都是一条直道。

``震天斧\ldots\ldots{}''蔡兴坤的声音在抖，``他把震天斧整个化实了\ldots\ldots 他才多大\ldots\ldots{}''

``蔡兴坤。''胡峰的声音一下子冷了，清清楚楚的。他的灰白眼睛里映着斧头上的金光，``你们蔡家的事，你倒知道得不少。''

他的手指头顺着斧柄上的纹路慢慢往下摸。摸到哪一道，哪一道就亮起来，亮得刺眼。斧子在他手里嗡嗡地响，那响声闷，沉，一下一下的，像很老很老的战鼓，每一下都敲在人的心口上。

然后，他身上的劲儿往上翻了。

胡峰身子里那一套序列又往下亮了。

322。

三个数一齐亮起来。胡峰身上的气一下子拔高了一大截。他的身子在电弧里往外胀，暗金的纹路从皮底下往外钻，一直爬满了全身。灰白眼睛里那两道竖着的金缝，一下子张开了，张成了两团烧着的金火苗。

还没完。

42。

最后两个数砸下来，像两下重锤砸在了半空里。四周的空气吱嘎响了一声，像一块玻璃，碎之前先吭一声。

0 111 222 322 42。

拿着震天斧的胡峰，全部的东西都在这儿了。

磁暴环碎了。

几百道电弧碰上他身上那股劲儿，跟蛛丝一样，一扯就断。环在他身边炸成了一片光点子，让夜风一吹，散了。

诺当往后退了一步，脸色发白。他的魂力跟那个环是连着的，环一碎，他胸口里像让人捶了一拳。

胡峰从碎掉的环里走了出来。斧头扛在肩上，暗金的纹路在他身上流来流去，像活的。刚才那个让人逼到墙角、满嘴是血的人不见了。这个人走路，一步是一步，踩得结结实实。

``这会儿，''他说，声音低，稳，肩上的斧刃在月亮底下发着冷光，``轮到我了。''

\begin{center}\rule{0.5\linewidth}{0.5pt}\end{center}

八个人的脸色一齐变了。

0 111 222 322 42 压下来，像一座看不见的山，压得每个人的骨头都咯吱咯吱地响。刚才那一下泰山压顶，身上的酸疼还没散，这一座比那一座还沉。

灭世狼王的两团火苗往里一收。它也觉出来了。冰帝把两只钳子收到了胸前，雪帝落回了地面，三个魂灵一齐往后撤了半步，摆出了守的架子。

蔡兴坤松开了蔡须坤的胳膊。她身上的土黄魂力转得飞快，脸还是白的，眼睛一直钉在那把斧子上。

``化实了\ldots\ldots{}''她嘴里念叨着，``再添上 322 跟 42\ldots\ldots 这会儿的胡峰，跟刚才那个，两样了。''

``兴坤姐。''蔡须坤压着嗓子问，``这斧子什么来头？''

蔡兴坤的喉咙动了一下，声音发涩。

``胡家正根上最硬的扽魂。''她说，``斧子里压着黑帅气本源最里头那一点。拿着它的人，能把一身的魂力都从斧子里过一遍，压紧了再放出去。打出来的东西，寻常魂技没法比。''

她停了一下。

``我只在书上见过它的画。''她说，``画底下就写了一行字：见斧，退。''

胡峰没再给他们说话的工夫。

他把斧子从肩膀上举起来，举过头顶。暗金的斧刃在夜里划出一道金弧，斧刃指着的地方，空气扭起来，像让火烤着。

``去死。''

两只手握着斧柄，往下一劈。

轰！

斧刃上飞出一道金光，十几丈宽，上百丈长，弯弯的，像一轮金色的月牙，横着朝八个人扫过来。

月牙扫过的地上，犁出了一道好几丈深的沟，沟两边的石板让余劲儿一震，全碎成了粉。半空里那些黑帅气让月牙一带，拧成一条一条暗金的蛟龙，冲在月牙前头开道，碰上什么，碾什么，吞什么。

``散开！''蔡须坤嗓子都喊劈了。

八个人往八个方向扑出去。

可月牙太快了，扫得也太宽了。

卢本韦跟奥斯卡没躲干净。

月牙的边擦过了卢本韦的左肩膀，把她整个人掀了起来。她在半空里翻了好几个跟头，重重地砸在几十丈外的一堆碎砖里。左肩上的衣裳豁开了一道大口子，底下的肉也豁开了，深得能看见骨头，血一下子涌出来。

奥斯卡让气浪掀翻了，在地上滚出去好几丈才停住，嘴角流出了血。

狼王、冰帝、雪帝三个挡住了月牙正中那一截，也让那股劲儿推得一步一步往后退。狼王的四只爪子在地上犁出四道深沟。冰帝的大钳子上裂开了一道一道细纹。雪帝那件白斗篷撕掉了一大半，挂在肩膀上，让风一吹，飘成了一条一条。

巷子里，小舞扑在担架上，拿后背挡住了一阵飞过来的碎砖。砖头砸在她背上，咚咚地响，她一声没吭。诺当蹲在担架另一头，两只手撑着地，脸色还是白的，磁暴环碎的那一下，他胸口里那一拳的劲儿还没缓过来。

唐三的蓝银草还钉着担架的四个角，草叶子让气浪刮断了几根，他手指头一勾，断了的地方又长出新的来。

``陛下呢？''唐三喊。

``还喘着。''小舞从担架上抬起头，脸上蹭了一道灰，``霜没碎。''

霍雨浩从一堆碎石里爬出来，拿袖子擦了一把鼻血。

``一下\ldots\ldots{}''他嗓子发干，``就一下，阵就散了\ldots\ldots 这股劲儿，比刚才高出去不止一层\ldots\ldots{}''

胡峰扛着斧子，往这边走过来。

一步一步，走得不快。他一脚踩下去，脚底下的地就裂开，塌下去一块。0 111 222 322 42 的纹路在他身上流着，斧刃在月亮底下发着金光，光照到谁脸上，谁的脸就白一层。

``还有谁？''他说，声音平平的，``你们压箱底的东西，我都接住了。现在，该我了。''

他走到二十来丈外，站住了，斧子往地上一拄。灰白的眼睛落到了蔡须坤身上，落了好一会儿。

他吸了吸鼻子。

``我说怎么一直闻着不对。''他说，``那堵墙上的金光，这头十九万年的狼，还有你身上这点味儿\ldots\ldots{}''

他的嘴角慢慢咧开了。

``我爹手底下的人在外头找了两年的东西，自己跑到天斗城来了。''

他把斧子从地上拔起来，扛回肩膀上。

``皇帝我先不要了。''他说，``把瑞兽交出来。剩下的人，还能走。担架上那个，你们也抬走。''

巷子口一下子静了。

没有一个人回头看她。

唐三往前挪了半步，站到了蔡须坤边上。霍雨浩擦了一把鼻血，也站了过去。巷子里，小舞从担架上抬起头，看了胡峰一眼，又低下头，两只手护着大帝的脑袋。

蔡须坤站在狼王边上，嘴角那道血已经干了。她往地上啐了一口，啐出来的是红的。

``做梦。''她说。

胡峰点了点头，像是早就知道她会这么答。

\begin{center}\rule{0.5\linewidth}{0.5pt}\end{center}

碎砖堆里，卢本韦咬着牙撑起了身子。

左肩膀在往外冒血，一阵一阵地冒，把半边衣裳都浸透了，黏在身上。她拿右手按着伤口按了一下，按不住，索性不按了。

她抬起头，隔着满天的烟尘，看着往这边走过来的胡峰。

她的嘴角往上一挑。

``奥斯卡。''

另一边，奥斯卡从地上爬了起来，嘴角还挂着血。他一听卢本韦这一声，眼睛一下子亮了。

他听得出来。这一声是在叫他干活。

``晓得了。''他咧开嘴笑了，笑得有点疯，手已经伸进了袋子里，摸出来一根香肠。

这一根跟刚才那两根金纹的又不一样。外皮上的纹路是红的跟金的拧在一块儿的，一圈一圈地缠着。香肠两头各箍着一个环，环是拿魂力铸出来的，亮着，扣得严丝合缝，看着就是专为一样什么东西做的。

超级大香肠。赤金纹的这一批，是奥斯卡专门给卢本韦那头母猪熬的。

``老规矩！''奥斯卡一扬手，把香肠朝卢本韦扔过去，赤金的光从烟尘里划过去，``我数三下！''

卢本韦右手一伸，接住了。她的魂力已经往地底下灌进去了。

``出来！''

粉红的光从地缝里冲上来，母猪又从地底下拱了出来。三丈高的一头大猪，四个蹄子落地，方圆几丈的地面又碎了一层。它鼻孔里喷着两股粉红的热气，两只血红的眼睛钉住了远处的胡峰。

卢本韦一翻身，爬上了猪背。左手揪住猪脖子上的鬃毛，左肩上的伤一扯，血又冒出来一片，她的脸白了一下，手没松。右手攥着那根超级大香肠，对准了母猪的鼻孔。

``对准了。''她说，声音稳得一点都不像一个正在淌血的人。

香肠头往母猪鼻孔里一塞，塞进去了。两头那两个环一碰着猪鼻子，自己就收紧了，扣得死死的，把香肠跟猪鼻子扣成了一体。赤金的纹路在猪鼻孔里亮起来，跟母猪身子里的魂力对上了，一跳一跳地应。

``一！''奥斯卡在那边喊，两只手已经抬起来了。

母猪觉出鼻子里塞进了东西，脑袋一甩就要往外甩。卢本韦一把按住它的鬃毛，魂力从手心里灌进猪脖子，硬把它那一甩给压住了。

``乖，忍一忍。''她趴下去，贴着猪耳朵说，声音软了一点，底下又压着一点狠，``一会儿就让你痛快。''

``二！''

奥斯卡两只手在身前画出了一个阵。阵画得很细，一道套一道，当中是一个漏斗一样的口子，口子朝着母猪的鼻孔。这个阵是他专门为这一手画的，画了不下几百遍，闭着眼都画得出来。

香肠炮。

道理简单得很。超级大香肠塞进母猪鼻孔，拿母猪身子里那一大缸魂力当火药，再拿奥斯卡这个阵把火药的劲儿全拧到一个方向上去。母猪憋到头，一口气冲出来，香肠就让那口气从鼻孔里顶出去，快得没边，打成一颗炮弹。

简单，粗，管用。

``三！走你！''

奥斯卡的阵一下子全亮了。赤金的魂力从阵里喷出来，灌进母猪的鼻孔，跟超级大香肠里的魂力撞在了一块儿，嗡嗡地共振。

母猪的眼珠子一下子瞪圆了。

鼻子里痒。

痒得钻心。香肠在它鼻孔里胀，越胀越大，越胀越烫，红的金的纹路亮得刺眼，一鼻孔的热往里顶，像鼻子里头塞了一座要开口子的火山。

母猪的身子弓起来了。

四个蹄子往地里一插，插得没了蹄腕子。脊背往上拱，拱成一张拉满了的弓。胸口一鼓一鼓地往外撑，肺里的气越吸越多，越压越紧，鼻子里那股劲儿一口气的工夫就顶到了头。

然后。

``阿\ldots\ldots 嚏！''

母猪打了一个喷嚏。

这个喷嚏打出来，跟炸了差不多。

超级大香肠在母猪鼻孔里攒的那点劲儿，一下子全从鼻孔里喷了出来。赤金的光从猪鼻子里冲出去，冲成一根三丈多粗的光柱。光柱最前头，那根超级大香肠已经成了一颗赤金色的炮弹，外皮让风擦得着了火，烧成了一团白光。

轰！

炮弹出去的那一下，母猪让它往后一坐，三丈高的身子往后滑出去几十丈，四个蹄子在地上犁出四道沟。卢本韦死死揪着鬃毛，让风吹得睁不开眼，眯着一条缝，一直盯着那颗炮弹。

炮弹快得眼睛跟不上，只看得见一道赤金色的线，从烟尘里穿过去，身后拖着一条长长的尾巴。尾巴里的空气让它烧空了，尖着嗓子叫。

胡峰的瞳孔一下子缩了。

他觉出来了。这一下跟刚才八个人打过来的东西都两样。刚才那些是炸，是咬，是锁。这一下，整团劲儿都压在一个尖上，冲着一个点往里钻。

他手里的震天斧自己嗡地响了一声，斧刃上的金光一阵一阵地跳。

``有点意思。''胡峰的嘴角往上一挑，灰白的眼睛里映着那颗越来越近的赤金炮弹。

他没躲。

他把斧子举起来了。

斧柄上的纹路一下子全亮了，0 111 222 322 42 一点不剩地灌进了斧身。斧刃上的金光亮得刺眼，斧子里压着的那点黑帅气本源让他整个叫醒了，劲儿从刃口上往外冒，比刚才劈那道月牙的时候还足。

``来。''

胡峰两只手握着斧柄，迎着那颗炮弹，劈了下去。

赤金的炮弹跟暗金的斧刃，在广场当中撞上了。

撞上的那一下，又是一片空。

然后才响。

那一声响得人魂儿都在晃。气浪往四面推出去，方圆几百丈的碎砖烂瓦让它一推，全推成了粉，地面整个往下塌了一丈。八个人让它掀了起来，狼王、冰帝、雪帝一齐吼了一声，让它推出去几百丈。

撞的那个地方，赤金跟暗金拧在一块儿，拧成了一个大光球。光球往外胀，越胀越大，面上爬满了细细的裂纹，像一个要破壳的蛋。

光球炸了。

赤金的碎片跟暗金的碎片一齐往四面八方飞，一片碎片落到地上，就炸出一个几丈宽的坑。宫门前那一片广场，这一下子全成了焦土，坑连着坑。

烟尘起来了，好一会儿都没落。

大坑那边，碎紫斗罗听见了这一声。

他那道紫线在半空里一拐，就要往广场这边扑。他人还没拐过来，一道白光已经横在了他前头。胡帅连头都没回，一只手往边上一挡，那道光就把他的路截死了。

紫色的小人撞在白光上，身上的裂纹又多了一道，让光推着倒飞回去。

``你的对手，''胡帅说，``是老夫。''

碎紫斗罗在半空里拧住了身子，咬着牙，又扑了回去。冰熊王拖着伤腿，从另一边堵上来。三个人又搅在了一块儿。

\begin{center}\rule{0.5\linewidth}{0.5pt}\end{center}

烟慢慢散开。

胡峰站在广场当中，斧子拄在地上。

头上的紫金冠碎了，明黄的太子常服又烂了一大块，露出底下暗金的内甲。他嘴角挂着一道血，脸上还是那副不急不忙的样子。

那颗赤金的炮弹，让他劈开了。

震天斧从正当中下去，把香肠炮的炮弹劈成了两半。两半从他身子两边飞过去，砸进远处的废墟里，炸出两个大坑，可已经没了往里钻的那股劲儿。

一斧子，两半。干干净净。

几十丈外，母猪四个蹄子还插在地里，粉红的皮上多了一道烧焦的黑印子。卢本韦趴在猪背上，左肩的伤口让这一震又裂开了，血顺着胳膊往下淌，半边衣襟全红了。

奥斯卡一条腿跪在地上，嘴角挂着血，脸白得跟纸一样。

``劈\ldots\ldots 劈开了？''他嗓子发干，``那可是超级大香肠\ldots\ldots 一整根的劲儿都压在一个尖上\ldots\ldots{}''

卢本韦趴在猪背上，没出声，眼睛钉着胡峰手里那把斧子。

斧子上，就多了一道新划痕。细细的一道，在暗金的斧面上亮着。

``震天斧\ldots\ldots{}''她哑着嗓子说，``胡家最硬的扽魂，就这个样子\ldots\ldots{}''

胡峰把斧子提起来。斧刃在月光底下泛着冷光，0 111 222 322 42 在斧身上流着，一点都没见少。刚才那一斧子，在他那儿好像没费多大劲儿。

``香肠炮？''他的嘴角往上挑了挑，``有点门道。一身的劲儿全拧到一个尖上，这一手，算我见过最机灵的几样里头的一样。''

他停了一下，灰白的眼睛里冷了一层。

``可力气差得太多，机灵就没用。''

他又把斧子举起来了。刃口上那团金光又攒起来，这一回，金光里头多了一丝红，红得发暗。那是斧子里的黑帅气本源，让他又往深里叫醒了一层。

八个人的脸色一齐变了。

霍雨浩往前站了一步，脑袋里一阵一阵地发空。冰帝的大钳子上那几道裂纹还没合上，雪帝站在它边上，低着头，一只手捏着那件撕烂的斗篷，斗篷的碎边上一点一点地结霜，霜连着霜，慢慢把口子补上。补得很慢。

狼王蹲在蔡须坤身前，两团火苗眯着，喉咙里一直在滚那串低低的响动。

谁都看得出来，再挨这么一斧子，挡不挡得住，难说。

卢本韦在猪背上撑起了身子。她拿右手背擦了一下嘴角的血，擦完了，嘴角往上一挑，挑得有点倔。

``奥斯卡。''她的声音发虚，可一个字是一个字。

``嗯。''奥斯卡也撑着地站了起来。他手里不知道什么时候又多了一根赤金纹的超级大香肠，两头的环亮着，``还有七根。''

``够了。''卢本韦拍了拍母猪的脖子。母猪低低地哼了一声，鼻孔里还在冒烟，冒的是方才那一炮剩下的烟。

``再来一回。''

\begin{center}\rule{0.5\linewidth}{0.5pt}\end{center}

\chapter{第七十二章：香肠篮球锤}\label{ux7b2cux4e03ux5341ux4e8cux7ae0ux9999ux80a0ux7beeux7403ux9524}

大坑那边。

真霜斗罗站在冰熊王的头顶上，看着半空里的胡帅。

胡帅还浮在那团白光里。Y(1,3,7,12) 的劲儿从他身上往外铺，铺开几百丈，铺到哪儿，哪儿就归他管。刚才碎紫那个小人让他一巴掌拍飞了一回，冰熊王那一下冰暴突击，也只是勉强从他那道光底下逃出来。硬碰硬，两个封号斗罗绑在一块儿，也不是他的对手。

真霜斗罗活了一百多年，硬碰硬本来就不是他最拿手的。

``老伙计。''他低声说。

他身上的魂力换了个走法。刚才是往外泼，这时候是往外渗，细得跟一根线一样，悄没声地往四边走。

冰熊王低低地哼了一声，它懂。它抬起爪子，落下去，落得很轻，像踩在一层薄冰上，生怕把冰踩裂了。每落一下，爪子底下就留下一层冰膜，薄得几乎看不见。

冰膜在往外爬。

从熊爪子底下往四面八方爬，爬得像蛛网，一点声音都没有。它不像冰霜领域那样一下子铺满一大片，冷得人打哆嗦。它贴着地，顺着碎石跟碎石当中的缝，顺着地上裂开的那些口子，一点一点地往里钻，往外延，一根搭着一根，织在一块儿。

胡帅这时候的心思，全在碎紫身上。

那个紫色的小人跟鬼一样，绕着他转，一会儿在左边，一会儿在头顶上，每一拳都是 Y(1,3,4,2,5,9) 的全部劲儿。哪一拳他都接得住，可碎紫太快了，快得他那道光锁不住人，他一刻都松不得。

他没往脚底下看。

直到真霜斗罗开口。

``封。''

就一个字。

冰熊王两只前掌往当中一合。

方圆几百丈的冰膜，一齐竖了起来。

它们一齐往里收。地上每一个角落里那些薄得看不见的冰膜，一下子全站起来，站成了一面一面的墙，从四面八方往胡帅那儿合。

墙合得不快，可哪儿都有。它们是从胡帅脚底下每一寸地里、每一道缝里一齐冒出来的，像一只正在收口的笼子。

胡帅的瞳孔一缩。

他这才觉出脚底下不对。

晚了。冰墙从地里冒出来那一下，就已经长到了他的腰。他手心里的白光一炸，一道光劈过去，把跟前那面墙劈碎了。

碎了一面，四面八方又补上来几面。

``冰牢？''胡帅的声音冷了一点。

他身上的白光一阵一阵地往外炸，靠近的冰墙一面接一面地让他劈碎。可冰墙太多了，铺得太宽了。他劈碎一面，补上来三面。劈碎三面，补上来九面。真霜斗罗把 Y(1,3,6,12) 一点不剩地灌进了这座冰牢，冰墙长出来的快慢，跟胡帅劈的快慢，几乎咬平了。

可真霜斗罗要的，不光是拿墙把他围住。

``碎紫，看着。''真霜斗罗说。

他身上的魂力又变了一回。冰熊王那边不再往外长新墙了，那些已经碎了的墙，让他拿魂力一拢，拢在了半空里。

胡帅劈碎的那些冰墙，一块都没落地。它们碎成了拳头大小的冰片，一片一片浮在冰牢里头。每一片都让真霜斗罗拿魂力磨过了，磨得跟镜子一样亮，斜着，歪着，一片对一片，角度一点不差。

胡帅又攒起一道白光，朝一面冰墙劈过去。

光劈碎了墙。

碎片没散。

那些镜子一样的冰片让真霜斗罗一拨，重新排了一排，迎着那道白光，把光弹了回去。

光撞在另一片冰上，又弹了一回。

又撞上第三片，又弹。第四片，第五片，第十片。

那道白光在冰牢里头的冰片之间来回乱撞，撞一回拐一个弯，像一头关在一只万花筒里的野兽，在无数面镜子当中东一头西一头地撞，找不着门。

胡帅的脸色到这时候才变了。

他的光出不去了。

他打出去的每一道光，都让冰片弹回来，在冰牢里头绕圈子，绕成一个圈，绕不出去。一道一道白光在笼子里游，像一笼子白蛇。

``反着照的冰牢\ldots\ldots{}''胡帅的声音从冰牢里传出来，压着火，``好手段。''

他把光一收。

冰牢里那些绕着圈的白光一下子全没了，那些冰片没了魂力撑着，一片一片往下掉。

真霜斗罗没让他喘这口气。

冰熊王四只爪子往地上一拍，冰膜又从地里涌上来，在胡帅身边又竖起一层一层的墙。新的冰片在墙跟墙当中重新排好，等着下一道光。

胡帅劈碎一层，真霜斗罗就补一层。

胡帅劈碎两层，真霜斗罗就补三层。

冰牢越垒越多，墙越来越厚，冰片越来越密。真霜斗罗把 Y(1,3,6,12) 往外倒，倒得跟不要命一样。一百五十米的冰熊王站在那儿，跟一座停不下来的冰窑一样，一层一层地往外出冰，把冰牢垒得一眼望不到头。

胡帅让困在了冰牢最里头。

他的光打出去就弹回来，他的感知让一层一层的冰墙隔着，他的眼睛让无数片冰镜子照得七零八落，看哪儿都是自己。他有一身 Y(1,3,7,12) 的劲儿，这时候像一只撞进了蛛网的老虎，越挣，网缠得越紧。

``碎紫！''真霜斗罗的声音隔着冰墙传出来，``老夫撑不了多久！''

远处，那个紫色的小人一闪。

碎紫斗罗听见了，冰牢一层一层垒起来，他也都看见了。

他知道，这是他的机会。

``超载。''他低低地说了一声。

紫色的小人身上一阵乱颤。

超载，是碎紫斗罗压在最底下、轻易不碰的那一手。极速模式是把快推到头。超载是拿魂力往极速模式里烧，一边烧一边换快，烧完了就没了。往后好几天，他得躺在床上，连拳头都攥不紧。

这时候他顾不上了。

小人身上爬满了裂纹，裂纹里往外冒紫光，刺得人睁不开眼。那光是魂力在烧，像一根蜡烛烧到最后一截，火苗子一下子蹿得老高。

碎紫斗罗快得没边了。

他从原地消失的那一下，连影子都没留。没有风声，四周的空气一动都没动。他穿过去了，空气才反应过来，在他身后嘭地合上。

冰牢里的胡帅，一下子觉出了不对。

他那点感知从冰墙碎开的缝里漏出去，碰上了一丁点东西，细得几乎摸不着。有个东西在往他这边来，快得他想不明白。

``什么\ldots\ldots{}''

他猛地一转身，手心里的白光攒起来。

身后没人。

他转身之前，碎紫就已经换了地方。紫色的小人在冰牢的缝里钻来钻去，像一道在迷宫里乱窜的闪电，闪一下换一个方向，再闪一下又换一个。

胡帅的光扫过来扫过去，一道都没扫中。冰牢挡着他的眼，碎紫的快搅着他的感知，他一身的劲儿使出来，打不着跟前那只蚊子。

``该死！''

胡帅的声音里，头一回冒出了真火。

他一巴掌劈碎了跟前的冰墙，从冰牢里冲了出来。

碎紫斗罗不在那儿。

他冲出来的那一下，紫色的小人已经绕到了他的侧后头，一拳砸在了他的腰眼上。

那一拳劲儿不大，可太快了，快得他身上那层光罩还没来得及往那儿聚。

砰。

胡帅的身子晃了一下。腰上那件暗金的内甲上，多了一个浅浅的拳头印子。

不疼。

可他胡帅，Y(1,3,7,12)，让一个 Y(1,3,4,2,5,9) 打中了。

他的脸一下子沉下来了。

``好。很好。''他的声音又低又冷，身上的白光一下子涨了一倍，``你们要玩，老夫陪你们玩到底。''

真霜斗罗在熊头上蹲了下来。

他的两条腿在打颤。一百多年了，他头一回觉得膝盖底下发软。Y(1,3,6,12) 是借来的，借来的东西有借期。他身子里那股热，比刚接住香肠的时候淡了一截，冷正一点一点地从骨头缝里往回爬。

冰熊王的右腿也在抖。伤口上的冰补了三回，化了三回。

他在心里数着。

奥斯卡那根香肠，还能撑多久，他说不准。冰牢还能垒几层，他也说不准。他就知道一件事，那边的孩子得快。

他抬起手，又往冰牢上补了一层墙。

\begin{center}\rule{0.5\linewidth}{0.5pt}\end{center}

广场这边。

卢本韦趴在猪背上，左肩的血还在往下滴，滴在粉红的猪皮上，一滴一个红点子。奥斯卡一条腿跪着，手里攥着那根新的赤金香肠。两个人隔着十来丈对了一眼。

香肠炮让人劈开了。

炮弹的劲儿够了，可那把斧子太硬了。0 111 222 322 42，再添上胡家最硬的扽魂，正面撞上去，撞不开。

``得换个打法。''卢本韦咬着牙，撑着猪脖子坐直了，眼睛把场上的人挨个扫了一遍。

远处，真霜长老的冰牢把胡帅困住了，碎紫将军在冰牢外头一圈一圈地绕。可两个人都撑不长。冰牢早晚让胡帅劈穿，碎紫将军那一身紫光烧完了，也就停了。

得赶在冰牢塌之前，把胡峰撂倒。

胡峰扛着斧子，站在几十丈外，看着他们。0 111 222 322 42 在他身上流着，斧刃在月光底下亮得扎眼。他不急，等着他们自己凑上来。

``正面打不过。''卢本韦说，嗓子又急又低，``他一个人，把咱们打过去的东西全扛下来了。得把所有人的劲儿捏到一块儿，砸在一个点上。''

``一个点\ldots\ldots{}''奥斯卡嘴里念了一遍，低头看了一眼手里的香肠。

他的眼睛一下子亮了。

他把手伸进袋子里，在里头翻，翻出来一样东西，圆滚滚的。

一颗篮球。

球皮上刻着一圈一圈的魂力纹路，磨得有点旧了。这是他随身带着的，平常拿来练香肠篮球，就是一边拍球一边往筐里扔香肠，扔准了就能吃。

``香肠插在篮球上！''奥斯卡说得飞快，``篮球有分量，香肠有爆劲儿，让一个人抡起来\ldots\ldots 那就成锤子了！''

卢本韦的眼睛一亮。

炮弹是走直线的，一斧子迎上去，从当中就劈开了。锤子两样。锤子有分量，往下坠，抡起来还带一道弧。震天斧劈得开迎面飞来的炮弹，从头顶上砸下来的一把大锤，它劈不劈得开，还两说。

``所有人的魂力，往一个人身上送。''卢本韦说，``送给队长。''

蔡须坤站在狼王边上，听见了，扭过头来。

``我？''

``你身子里有瑞兽。''卢本韦说，``七个人的魂力七个样，搁谁身上都得打架。就你那个本源，揉得匀。''

蔡须坤没再问。她抬手在狼王脖子上按了一下。狼王看了她一眼，往后退了几步，蹲在了巷子口，两团火苗一直盯着胡峰。冰帝跟雪帝也退到了它两边，一左一右，守着底下那副担架。

\begin{center}\rule{0.5\linewidth}{0.5pt}\end{center}

事情一下子就定了。

奥斯卡把那根赤金香肠往篮球的气嘴里一插。香肠两头的环碰上球皮上的纹路，自己就扣紧了。赤金的纹路从香肠上往外爬，爬满了整颗球，把香肠跟球长成了一样东西。

一颗爬满赤金纹路的香肠篮球，在奥斯卡手里冒着热气。

``兴坤！''卢本韦扭头喊，``带我们上去！''

蔡兴坤一句话没问。

她肩膀上的黑子鸡早就跳到了地上，这时候它咯的叫了一声，一扑棱，跳回了她的后背上。

蔡兴坤两只脚往地上一跺，土黄色的魂力在她身子里转起来，转得飞快。

黑子鸡在她背上一抖，身子往两边一摊。

一对翅膀从她后背上张开了。漆黑的，一边一丈多，一张就是三丈多宽，每一根羽毛都泛着铁一样的暗光。

黑子鸡，蔡兴坤的扽魂。飞的那一路，少见，能驮的分量比同一层次的飞行魂兽多出去好几倍。

翅膀一扇，地上的碎砖让风刮得往四边滚。蔡兴坤的脚离了地，往上一窜。她两只手往下一抄，一手揪住卢本韦的后领子，一手揪住奥斯卡的后领子，把两个人一块儿提了起来。翅膀驮着三个人，扇得又稳又有劲儿，几下就上了几十丈高。

母猪在底下仰着脑袋，哼了一声。

``去！''卢本韦在半空里冲它喊，``守着小舞！''

母猪甩了甩耳朵，四个蹄子一刨，跑到巷子口，往担架前头一横，横成了一堵粉红的墙。

唐三那边也动了。

他手心里的蓝银草从地底下往上钻，钻出来几十根，一根比一根粗，像一条条绿色的大蟒。几根缠住了诺当，几根缠住了霍雨浩，几根缠住了蔡须坤。草身子往下一缩，缩得紧紧的，跟拉满的弹弓一样，再往上一弹。

三个人让蓝银草一甩，甩上了半空。

唐三自己踩着一根草，草身一弹，也上去了。

小舞没上。

她站在担架边上，唐三扔下来一根蓝银草，草梢缠在她手腕子上，另一头一直连到半空里去。

``我守着陛下。''她仰着脑袋喊，``劲儿顺着草送上去！''

前后不到三息。

胡峰猛地抬起头。灰白的眼睛里映着半空里那几个人影。他头一个念头就是举斧子，往上劈一道。

斧子刚举起来，一道紫光闪到了他脸跟前。

碎紫斗罗。

紫色的小人从冰牢那边绕过来，就为了这一下，一拳砸在胡峰的脸上。

那一拳没多大劲儿，可把胡峰的斧子砸偏了半寸。

一道金色的月牙从八个人身边擦过去，在夜空里划了一道弧，一个人都没碰着。

紫色的小人一闪就没了，又扑回冰牢那边去了。

``谢了！''卢本韦在半空里扯着嗓子喊。

胡峰抬手摸了一下脸。

鼻梁上让那一拳砸出了一道血口子，不深，血顺着鼻翼往下淌。他把手指头上的血看了看，舔了一下，抬头又往半空里看。

他没再劈第二斧子。他把斧子扛回肩膀上，站在原地，看着半空里那几个人往一块儿凑，像看一群鸟往一棵树上落。

``凑吧。''他低声说，``凑齐了，省得我一个一个去找。''

半空里，几个人往一块儿凑。蔡兴坤的黑子鸡翅膀扇得稳稳的，一下一下，把几个人的分量都扛住了。唐三的蓝银草在几个人当中来回地织，一根缠一根，织成了一张大网，把每个人的腰都拴在了网上，连底下小舞手腕上那一根也织进去了。

霍雨浩挂在网上，往底下看了一眼。

从这么高往下看，整片宫门前一眼就看全了。大坑那边一座一座的冰墙，垒得跟一座冰城一样，冰城里一团白光在撞。广场上坑连着坑，到处冒着烟。再往外，长街那头，刚才让气浪掀翻的那三千甲士，正一排一排地往起站，站起来就往后退，退到了几条街口，又一排一排地站齐了。

他们没往这边冲。

二十几个魂斗罗站在他们前头，一动不动，也往这边看着。

``他们在等。''霍雨浩说。

``等什么？''卢本韦问。

霍雨浩没答上来。

``所有人，往队长那儿送！''唐三喊。

蓝银草的草身里，一下子亮起了一道一道的光。每一根草都成了一条道，道的那一头，全通到蔡须坤身上。

七个人的魂力顺着草往蔡须坤身子里淌。

各人的魂力各是各的味儿。霍雨浩那一缕冷得她后槽牙发酸，奥斯卡送过来的那一缕，居然还带着点香肠味儿。小舞那一缕最细，从地上一路爬上来，一直没断。

这些魂力进了蔡须坤的身子，先是乱的，你撞我，我撞你。瑞兽本源在当中一转，把它们一样一样地揉到一块儿，揉匀了，压紧了，压成了一道金色的劲儿，在她经脉里哗哗地淌。

蔡须坤的身子一震。

她的两只眼睛变成了金色。这一层金比瑞兽本源平常那点金深，深得发暗，像很老很老的金子。金色从她眼睛里往外漫，漫到她全身，把她整个人裹在了一层金光里头。

她身子里那一套序列，跟着亮了。

0。

111。

222。

0 111 222。

七个人的劲儿压在她一个人身上，把她顶到了胡峰刚才没拿斧子时候的那个地方。

她伸出了右手。

奥斯卡把那颗香肠篮球朝她扔了过去。

蔡须坤一把接住，五个指头一扣。球皮上的赤金纹路一下子往她手心里钻，跟她的瑞兽本源搅在了一块儿。香肠篮球在她手里胀了起来，从一颗篮球胀成了一扇磨盘那么大，赤金的纹路越来越亮，烫得她手心里滋滋地响。

香肠篮球锤。

蔡须坤一只手攥着那颗磨盘大的锤头，从半空里往下看着胡峰。

金光在她身上流着，瑞兽本源在她身子里一圈一圈地转，转得又沉又稳。0 111 222 顺着她的胳膊往下走，从指头尖灌进锤头。锤头上的赤金纹路一下子全亮了，亮得刺眼。

胡峰仰着头，往上看。

他看见了她手里那把锤，看见了锤头上流来流去的赤金纹路。他也看见了她身上那层金光，跟他刚才拿斧子之前，差不了多少。

``有点意思。''他嘴角一挑，把斧子慢慢举了起来，``一窝蚂蚁，也敢\ldots\ldots{}''

蔡须坤一锤子砸了下来。

没有花样，没有拐弯。

从上往下，直直地砸。

就这么一砸，七个人的魂力全在里头，0 111 222 全在里头。锤头从半空里往下坠，坠得像一块金色的陨石，直奔胡峰的脑袋顶。

锤子底下的空气让它挤得没处去，尖着嗓子叫。锤头上的赤金纹路越坠越亮，亮得像一团烧着的火。

胡峰嘴角那点笑，一下子僵住了。

他觉出来了。这一锤子跟刚才哪一下都两样。刚才八个人打过来的，东一下，西一下，散的。这一锤子，八个人的劲儿全捏成了一团，压在一个头上。

``什\ldots\ldots{}''

他本能地把斧子横了过来。

轰！

锤子砸在了斧子上。

撞的那一下，一圈白色的气浪从撞的地方往外推，推出去几十丈。地上的碎石、尘土、烂木头，让它一推，全飞上了半空。

胡峰的两只脚往地里一沉，沉下去半尺。

他的瞳孔一下子缩成了针尖。

沉。

锤子上压下来的那股劲儿，沉得他没想到。那股劲儿一点都不往四边散，攒成一团，死沉死沉地往下坠。像有人搬了一座山，搁在了他的斧子上，又顺着斧柄压到他胳膊上，压到他肩膀上。

震天斧在他手里嗡地叫了一声，叫得发颤。斧柄上那些暗金的纹路一闪一闪的，像一根绷到头的琴弦。

``这股劲儿\ldots\ldots{}''胡峰从牙缝里往外挤字，``一个魂师\ldots\ldots 哪来这么大的劲儿\ldots\ldots{}''

蔡须坤没答他。

她两只金色的眼睛往下看着他，一眨不眨。锤子压在斧子上，一点一点往下按。0 111 222 在锤头里转着，七个人的魂力全往锤子跟斧子碰着的那一点上灌。

胡峰咬紧了牙。

他两条胳膊在抖。0 111 222 322 42 在他身子里转得飞快，要把锤子顶回去。顶不动。锤子太沉了，沉得那把斧子都在叫唤。

他想起了进宫以前，叔父在马车里跟他说的那几句话。

那时候叔父手里捏着一枚棋子，在指头缝里转来转去。``皇帝要活的。''叔父说，``真霜交给老夫。扽魂殿跟来的那几个孩子，碰上了，碾过去就是，别缠。''

他当时笑了一声。几个孩子，还用得着叔父特意交代？

刚才闻出那个丫头身上的东西，他才想起他爹说过的一句话：瑞兽那一路，最会把别人的力气往自己身上揉。

他也没往心里去。一个魂斗罗都不到的丫头，能揉多少？

这时候那把锤子压在他斧子上，他知道能揉多少了。

``给我\ldots\ldots 碎！''

胡峰吼了一声。

震天斧上猛地炸出一道金光，刺得人睁不开眼。0 111 222 322 42 一点不剩地从斧刃上倒出去，倒成一道斧芒，要把锤子连着拿锤子的人一块儿劈成两半。

又撞上了。

金光跟赤金光在撞的地方拧成一团，拧成一个大光球。光球一下一下地往外胀，面上爬满了细纹，咯吱咯吱地响。

大坑那边，冰牢里头。

胡帅一巴掌劈开了跟前三层冰墙，从缝里往广场这边看了一眼。

他看见了半空里那团金光，看见了胡峰两只脚陷在地里，斧子横在头顶上，让一把锤子往下压。

他往那边迈了一步。

一步还没落地，脚底下又竖起来一面墙，墙上斜着十几片冰镜子，把他刚攒起来的那点光又弹回了他自己脸上。紫色的小人从他左边一闪，一拳砸在他肋巴骨上。

``老夫说过，''真霜斗罗站在熊头上，嘴角的血已经流到了下巴，``你的对手，是我们两个。''

胡帅回过头，看了他一眼。那一眼里头，头一回没了笑。

蔡须坤的胳膊一震，虎口裂开了，血顺着锤柄往下淌。

她没退。

她退不了。她身后是七个人，七个人的魂力都在这一锤子里。她这一退，大家掏出来的东西全白掏了。

``大家\ldots\ldots{}''她从牙缝里挤出来，``再给我一点！''

半空里，几个人一齐咬紧了牙。

唐三的蓝银草一下子粗了一圈，草身里的光亮得发白，他把自己往草里灌，灌得脸都青了。

霍雨浩脑袋里那口缸早见了底。他拿精神力在蔡须坤的识海边上搭了一座桥，从桥上往过送的，已经是最后那一层底子。他鼻子底下的血一直没停过，滴在蓝银草上，顺着草往下淌。

诺当伸出一只手，手心里还剩着磁暴环碎掉的那几丝电，他把电拧成一根细线，搭在锤头上。

卢本韦伏在蔡兴坤的胳膊底下，冲着底下喊了一声。地上那头母猪仰起脖子嚎了一嗓子，身子里的粉红魂力顺着契约往卢本韦身上涌，又从卢本韦身上顺着草往上送。她左肩的伤口一下子又裂开了，她没管。

奥斯卡把身子里最后一点魂力也榨了出来，榨成一根赤金的细线，直接接在了篮球的气嘴上。榨完了，他的脑袋往下一耷拉，挂在蓝银草的网上，还剩一口气，嘴里念叨着：``值\ldots\ldots{}''

蔡兴坤那对黑翅膀扇得越来越急，扇一下，土黄色的魂力就从她背上往外冒一阵，顺着草往蔡须坤那边走。

地上，小舞跪在担架边上，攥着手腕上那根草。她把身子里能掏的都掏了，掏得眼前一阵一阵发黑。

她低头看了一眼担架。

大帝脸上那层霜底下，眼皮动了一下。

小舞以为自己眼花了，凑过去盯着看。过了两口气，那只眼皮又动了一下，眼珠子在眼皮底下转了半圈。担架边上，大帝搭在被子外头的一根手指头，往里勾了勾。

``陛下？''小舞小声叫。

没人应。可那根手指头没再松开，一直勾着。

八个人的劲儿，一齐往上涨了一截。

锤头上的赤金纹路，一下子全成了金的。

一点红都没剩下，纯纯的金。

锤子又沉了。

这一回沉，是因为八个人的劲儿全拧到了一个方向上。往下。

那颗香肠篮球锤，这时候就是一座往下掉的山。

胡峰的脸一下子没了血色。

斧子上压下来的那股劲儿，一下子涨到了他想都没想过的地方。他把斧子转了个角度想往边上卸，卸不掉。分量就是分量，压下来了，什么巧劲儿都使不上。

真像一座山压在了他斧子上。

他的两条胳膊咯咯地响，骨头在叫。两只脚又往地里陷了半尺，膝盖以下的石头全碎成了粉。

``不可能\ldots\ldots{}''他从牙缝里往外挤，``这股分量\ldots\ldots 不可能\ldots\ldots{}''

蔡须坤从半空里往下看着他。两只金色的眼睛很静，静得像一口深井，井里什么都没有，就一个念头。

``胡峰。''她的声音从半空里落下来，不大，穿过满场的轰响，落进他耳朵里，``下去。''

锤子在八个人的劲儿底下，往下走。

一寸。又一寸。

蔡须坤的两条胳膊都在抖。虎口上的血早顺着锤柄淌到了手腕子上，又从手腕子滴下去，滴到胡峰的斧刃上，嗞的一声冒了一缕烟。

她觉得自己的骨头在响。锤子压着斧子，斧子顶着锤子，两头的劲儿全挤在她两条胳膊里，挤得她肩膀里头一阵一阵地发麻。

底下巷子口，那头灰黑的狼仰起脖子，长长地嚎了一声。

那一声顺着她的识海钻进来，在她脑子里转了一圈。她发麻的肩膀里，一下子又多了一点劲儿。

震天斧拖着长音哀叫了一声。斧柄上那些暗金的纹路，一根跟着一根地断，一根跟着一根地灭。胡峰的两只脚已经整个陷进了地里，膝盖以下全埋在碎石头底下。他的脸白得跟纸一样，灰白的眼珠子在往两边转，转得没个落处。

可胡峰到底是胡峰。

他的嘴角往上一扯，扯出一个疯得吓人的笑。

``好\ldots\ldots 好一把香肠篮球锤\ldots\ldots{}''他的嗓子哑得像破锣，灰白眼睛里那两团金火苗一下子蹿高了，``你们要比劲儿，那就比到头！''

震天斧在他手里，吼了一声。

\begin{center}\rule{0.5\linewidth}{0.5pt}\end{center}

\chapter{第七十三章：八级定装魂导炮弹}\label{ux7b2cux4e03ux5341ux4e09ux7ae0ux516bux7ea7ux5b9aux88c5ux9b42ux5bfcux70aeux5f39}

锤子跟斧子顶在一块儿，顶了二十息。

蔡须坤的两只眼睛还是金的。八个人的魂力顺着蓝银草一刻不停地往锤子里灌，那股纯金的分量压着震天斧，一寸一寸往下按。胡峰的两只脚早整个陷进了地里，膝盖以下的石板全成了粉。

胡峰没倒。

他的两条胳膊抖得厉害，骨头咯咯地响，脸白得跟纸一样。可他的嘴角扯着，扯出一个笑。那个笑不苦，倒有点像赌桌上的人，把最后一把筹码往前推之前，先笑一下。

``你们\ldots\ldots 是真强。''他从牙缝里往外挤，一个字一口血沫子，``可你们弄错了一件事。''

他眼睛里那两团金火苗，一下子蹿高了。

``你们当我就这么点东西？''

\begin{center}\rule{0.5\linewidth}{0.5pt}\end{center}

蔡须坤觉出来了。

锤子底下的劲儿，一下子变了。

刚才是一点一点往上顶。这时候，像山塌了一样往上涌。胡峰身上的劲儿，一口气的工夫翻了一倍。震天斧上的暗金光从暗往亮里走，亮得刺眼，斧柄上那些断了的纹路一根一根全接上了，接上了还往外冒光，一条一条，像流着金水的沟。

``不好！''蔡须坤的瞳孔一缩，``他在攒劲儿！''

晚了。

胡峰两只手猛地攥紧了斧柄。他的身子又往下沉了半尺，膝盖都碰着地了。就在他眼看要让锤子压趴下的那一下，比刚才凶得多的劲儿从他身子里炸了出来。

这一下他什么都不管了。身子里剩下的那点魂力，不管是哪一层的，一口气全拧到一块儿，拧得死紧，然后往外放。

``斧破山河。''

他的声音一下子远了，像从很远很远的古战场上传过来的鼓声。他两只手把斧子往上一举，斧刃上的光攒成了一道从没见过的大光刃，二十多丈宽，上百丈长，像一轮刚从地边上冒出来的金太阳，朝着香肠篮球锤劈了上去。

两样东西在夜空里撞上了。

又是那一片空。

然后才响。

那一声响得人魂儿都往外飘。气浪往四边推，方圆几百丈的地面整块往下塌了好几尺。八个人让气浪掀了起来。底下巷子口，狼王、冰帝、雪帝一齐吼了一声，让它推出去几百丈。

撞的地方，赤金的光跟暗金的光拧成一团，拧成一个大光球。光球往外胀，面上全是裂纹。

光球炸了。

赤金的碎片跟暗金的碎片往四面八方飞。香肠篮球锤让斧破山河一震，锤头上的赤金纹路一大片一大片地灭了。

蔡须坤的胳膊一麻。

她觉出来了。锤子里那七个人的劲儿，在往下掉，掉得飞快。

``怎么\ldots\ldots{}''

她自己的劲儿还在。见了底的，是大家的魂力。刚才为了让锤子一直往下压，七个人早把自己榨到了头。胡峰这一下斧破山河，全部的劲儿捏在一个尖上往外炸，一下子就把七个人合起来的那点东西捅穿了。

蓝银草在她身边一根一根地蔫下去，草身里的光一节一节地灭。霍雨浩搭在她识海边上的那座桥，裂了，那边送过来的东西断断续续，一阵有，一阵没有。诺当搭在锤头上那根电线，啪的一声断了。底下巷子口，母猪让气浪掀了个四脚朝天，卢本韦身上那条从母猪那儿过来的魂力，一下子掐断了。

锤子上那股往下坠的金色分量，一眨眼就少了七成。

胡峰觉出来了。

他的嘴咧开了，咧得很大。

``轻了！''

震天斧上猛地又炸出一道金光。0 111 222 322 42 剩下的那点劲儿一下子全倒了出来，倒成一道横着扫的斧芒，把香肠篮球锤连着拿锤子的蔡须坤，一块儿往天上掀了出去。

轰！

蔡须坤让斧芒一冲，跟一只断了线的风筝一样，往天上倒飞。锤子从她手里脱了出去，赤金的光在半空里划了一道弧，越飞越远。

半空里，挂在黑子鸡翅膀跟蓝银草上的那几个人，也让斧芒的余波扫着了。

卢本韦，奥斯卡，唐三，霍雨浩，诺当，蔡兴坤。

蔡兴坤背上那对黑翅膀让这一扫震得乱抖，羽毛一把一把地往下掉。她拼了命地扇，想稳住，可那股劲儿太凶了，黑子鸡驮着这么多人，本来就到了头，这一下，驮不住了。

``抓住\ldots\ldots{}''蔡兴坤的声音让风一刮，碎成了几截。

晚了。

几个人一齐从网上掀了出去，像让弹弓弹出去的石子，往几个方向飞，在半空里翻了好几个跟头，各自摔在几十丈外。

卢本韦砸进了一片烂瓦堆里，左肩上那道伤口整个崩开了，血往外涌。奥斯卡磕在一块大石头上，嘴里吐出一口血。唐三身边的蓝银草全枯了，他在地上滚出去好几丈才停住。霍雨浩脑袋里嗡的一声，鼻血跟开了闸一样往下淌。诺当的后背在地上犁出一道长沟。蔡兴坤背上的翅膀掉了一大半毛，勉强张开，歪歪斜斜地滑下来，一头栽在地上，黑子鸡从她背上滚下来，变回了一只鸡，扑腾着翅膀，咯咯地叫。

一眨眼，八个人全散了。

胡峰站在广场当中，斧子拄在地上，大口大口地喘。

斧破山河这一下，把他身上的魂力也抽得差不多了。0 111 222 322 42 在往下缩，身上的暗金纹路一明一灭。可他脸上挂着笑，赢了的人才有的那种笑。

``完了。''他的嗓子哑，可声音里透着得意，``你们那点合起来的东西，让我\ldots\ldots{}''

他的话没说完。

他看见了蔡须坤。

蔡须坤还在往下掉。

斧芒把她掀上去有几十丈高，掀到头，人就往下坠。风从她耳朵边上呼呼地刮，刮得她眼睛都睁不开。她身上那层金光让斧芒撕掉了一大半，剩下的一点贴在皮上，一闪一闪的。

坠到一半，她伸手一抓。

抓住了一块石头。

那是刚才光球炸开的时候崩上天的一块大石板，让余劲儿托着，还在半空里打着转，还没落下去。她一只脚踩上去，身子一拧，站住了。

石板晃了两晃，托着她，慢慢往下沉。

她低头往下看。胡峰站在广场当中，仰着脸看她。

她的衣裳撕烂了一大半，嘴角挂着血，两只眼睛还是金的，可金得发暗，像一盏油快熬干了的灯。

她的右手，攥着一样东西。

拳头那么大，圆的，通体漆黑。面上刻满了纹路，一圈套一圈，密得看不清哪儿起头哪儿收尾。纹路里有光在走，走得很慢，一下一下的，像里头关着个活物，在喘气。

八级定装魂导炮弹。

出城那天，大长老从殿后头的暗格里取出来，搁在她手心里的。他当时就说了几句话：里头封着 0-Y 序列的功法，炸开了够把一个封号斗罗打成重伤。用了它，用的人要虚好几天。不到万不得已，不许动。

一路上她隔着衣服按过它好几回。在城门口按过，在巷子里按过，在宫门前也按过。每按一回，就又把手拿开了。

她站在石板上，往底下看了一眼。

卢本韦躺在烂瓦堆里，一只手捂着肩膀，血从指头缝里往外冒。奥斯卡趴在大石头边上，脑袋耷拉着，看不出还醒没醒。唐三撑着胳膊想站起来，撑了两回，又趴下了。霍雨浩坐在地上，仰着脸，鼻子底下两道血。

巷子口，小舞还跪在担架边上，一只手攥着那根早就枯了的草。

不到万不得已。

她把这几个字在嘴里过了一遍。

这时候，她把它从贴身的储物魂导器里拿了出来。

还不止这样。

刚才斧破山河跟锤子撞的那一下，撞散的魂力像下了一场暴雨，往四边乱溅。炮弹揣在她怀里，贴着她的肋骨，那些溅出来的魂力有一点让炮弹面上的纹路吸了进去。这时候那些纹路里走的光，比出城那天快多了，快得一跳一跳的。

蔡须坤把炮弹攥紧了。

她身子里头，有一样东西醒了。

这样东西的颜色，跟她平常用的那点金不一样，跟刚才七个人合起来的赤金也不一样。它更老，更干净。

是绿的。

一点翠绿的光从她胸口里冒了出来。那绿像开春以后头一根从土里拱出来的芽，嫩，可是一个劲儿地往上长，挡都挡不住。绿光从她胸口往外走，顺着右胳膊往下淌，淌进了她手心里那颗黑炮弹。

帝皇瑞兽的处女之力。

这是瑞兽那一脉身上压得最深、最干净，也最不许碰的东西。它不归哪一种属性，也不归哪一条规矩。它就是活着本身。瑞兽是魂兽里头的王，它那一身本事，说到底，是能管住活着这件事：能生出来，能给出去，能往大里催，也能一把掐灭。

处女之力，是瑞兽这一身活气，一回都还没给过别人以前，在身子里攒了几万年的那一份。

一放出去，就收不回来了。

蔡须坤放了。

那一缕绿从她手心灌进炮弹里，跟炮弹里本来的东西撞在一块儿。炮弹里封着的那点 0-Y 序列，本来是冷的，硬的，一节扣一节。绿一进去，那些节一下子活了，像冬天冻住的河让春水一冲，冰面上一道一道地裂开。

炮弹面上那八圈纹路，一圈接一圈全亮了。漆黑的炮弹皮底下透出了绿光，一开始是一点一点的，后来连成了片，整颗炮弹绿得发亮。

它在她手心里往大里胀。拳头那么大，胀成了人头那么大，又胀成了一个西瓜那么大。面上那些纹路让它撑得吱嘎吱嘎地响，纹路边上开始往外裂细缝。

它没炸。

蔡须坤拿那一份处女之力，硬生生地把它摁着。

她得把它送到胡峰跟前。

扔，扔不准，胡峰躲得开。她也没有能把它打出去的东西。

她得自己拿着它，一直摁到胡峰身上。

\begin{center}\rule{0.5\linewidth}{0.5pt}\end{center}

胡峰看见了那颗绿炮弹。

他的瞳孔一下子缩紧了。

眼睛里那两团金火苗一阵乱颤，颤得根本由不得他。那是他身子在替他怕，比他脑子快了一步。他不知道那颗炮弹到底有多凶，可他浑身每一根骨头都在跟他喊同一个字。

跑。

``叔父！''

胡峰张嘴就喊了出来。刚才那副从容，那副一字一顿的腔调，全没了。声音又急又尖，劈了。

``叔父！帮我\ldots\ldots{}''

他的声音穿过满场的轰响，穿过冰牢那一堵堵碎墙，一直传到了大坑那边。

胡帅听见了。

他正一巴掌劈碎跟前第四层冰墙。这一声钻进他耳朵里，他的手停在了半空里。

他这个侄儿，打小跟在他身边，挨过打，挨过饿，在死人堆里趴过三天三夜，从来没这么喊过他。

他扭过头。

他看见了蔡须坤手里那颗绿炮弹。看见了炮弹面上那几圈快要撑裂的纹路，看见了炮弹里头那团又烫又活、一个劲儿往外拱的东西。

他的瞳孔一下子缩成了针尖。

他这辈子见过的魂导器，数都数不过来。炮弹，见过上百种。能毁城、能削山的东西，他也见过。可那颗炮弹里攒着的那团劲儿，密得他心里一沉。绿的那一份跟里头的 0-Y 搅在了一块儿，炸开会是什么样子，他算不出来。他只算得出来一件事。

胡峰扛不住。

胡峰就在它跟前。

``该死！''

胡帅的声音里一点从容都没了。他身上那团白光猛地一炸，炸得比刚才哪一回都亮，Y(1,3,7,12) 一点不留地往外倒。冰牢他不管了，碎紫他不管了，他这一晚上布的那些局，他也顾不上了。

他就剩一个念头。

把胡峰拽出来。

他整个人化成了一道白光，往广场那边射了过去。快得碎紫那一身超载都没跟上。紫色的小人一拳砸过去，砸进了他留下的影子里，攥住了一把风。

``不好！''碎紫斗罗在半空里拧过身来。

他看清了胡帅往哪儿去。

往胡峰那儿。

往那颗绿炮弹那儿。

\begin{center}\rule{0.5\linewidth}{0.5pt}\end{center}

蔡须坤看见了那道白光。

它从大坑那边射过来，像一颗从天上砸下来的星星，快得她躲都没处躲。她心里清楚，炮弹还没摁到胡峰身上，这道白光就会先到，把她连人带炮弹一块儿截住。

她的嘴角往上挑了一下。

``等的就是你。''

她一扬手，把炮弹扔了出去。

炮弹飞的方向，对着的是胡帅的脸。

那颗绿炮弹从她手里飞出去，化成一道绿光，从满天的烟尘里穿过去，直奔胡帅的面门。

胡帅的瞳孔一缩。

他没躲。

他躲不了。他一躲，炮弹就从他身边穿过去，落到他身后的胡峰身上。胡峰这时候身上的魂力快空了，挨上这一下，活不成。

他只能挡。

``破！''

胡帅吼了一声，身上的白光在身前一拢，拢成了一面好几尺厚的光盾。光盾上刻满了帅气纹路，一道一道地转，转得飞快。他把 Y(1,3,7,12) 全灌进了这面盾里。

绿炮弹撞在了白盾上。

撞上的那一下，什么都绿了。

绿光把宫门前的一切全吞了。那团又烫又活的东西在撞的那一点上往外冲，冲得光盾往里凹，盾上的纹路一根接一根地断，一根接一根地灭。胡帅的魂力还在往盾里灌，断了的纹路又一根根接回去，接回去又断。两边就这么咬着，咬了两口气。

然后炮弹炸了。

那一声，比今晚哪一声都响。

绿光跟白光在炸开的地方拧在一块儿，拧成一个一百多丈宽的大光球，光球面上绿的白的纹路搅在一起，像一颗正在生出来的星星。

光球又炸了。

绿的碎片跟白的碎片往四面八方飞。气浪往外推，方圆几百丈的地面整块往下塌了好几丈。剩下的那点废墟全成了粉，碎石成了土，连半空里飘着的那些魂力，都让这一下冲得干干净净。

胡帅让气浪从正当中推了出去，推出去几百丈。

他那一身帅气真身，一下子全散了。白光像让大风吹灭的一排蜡烛，一片一片地灭。蟒袍早没了，底下那件暗金的内甲上全是裂纹，血从裂纹里往外渗，把他整个人染成了暗红的。

他没死。

Y(1,3,7,12) 的人，一颗八级的炮弹要不了他的命。炸开的那一下，他身上本能地又裹了最后一层光，把最要命的那一截挡在了外头。

可胡峰就在炸开的地方边上。

炮弹没打着他，可余劲儿结结实实地撞在了他身上。斧破山河已经把他抽得差不多了，0 111 222 322 42 让这一撞，一下子往里缩，缩得身上的暗金纹路只剩一点火星子。

他往后飞了出去，像一只断了线的风筝。

一口血从他嘴里喷出来，在夜里拉出一条红线。

他的脸一下子白透了，眼睛里那两团金火苗暗得快没了。震天斧从他手里脱出去，在半空里翻了好几个跟头，一头扎进了远处的地里，斧柄还在嗡嗡地抖。

他一头砸进了一百多丈外的废墟里，砸出一个几丈宽的坑。

巷子口，绿光扑过来的那一下，母猪一头拱到担架前头，拿后背把担架整个挡住了。狼王往母猪身前一站，冰帝把两只大钳子合起来，竖在狼王边上，雪帝张开那件破斗篷，在它们头顶上落下了一层雪。

绿光撞在雪上，雪化了。撞在钳子上，钳子上的裂纹又多了几道。撞在狼王身上，狼王四只爪子往地里插得更深了，一步没退。

小舞趴在担架上，两只手护着大帝的脑袋，后背让碎石子砸得一下一下地响。

等那阵光过去，她抬起头。担架上，大帝脸上那层霜碎了一角，露出来半边眉毛。眉毛底下，那只眼皮又跳了两下。

\begin{center}\rule{0.5\linewidth}{0.5pt}\end{center}

真霜斗罗看见了。

他站在冰熊王的头顶上，看着胡峰从炸开的地方往后飞出去。

他的嘴角往下一压。

就是这时候。

``老伙计。''他说，声音又低又快。他身上的魂力换了个去处，冰牢不要了，剩下的那点 Y(1,3,6,12)，一点不留地全灌进了冰熊王的身子里。

``死亡翻滚。''

冰熊王两只眼睛猛地亮成了两团冰蓝的火。

它一百五十米的身子转了起来。它拿右前掌往地上一戳，戳成一根轴，整个身子绕着这根轴转，越转越快，转成了一根冰蓝色的大钻头。

四只爪子在转的时候伸成了四道冰蓝的光刃。每一道光刃上，都压着 Y(1,3,6,12) 的魂力，压着一百五十米冰熊的全身分量。转得越快，劲儿甩得越狠，光刃在半空里划过去，划出一道一道的黑缝。

冰熊王压箱底的那一手。

这根一百五十米的冰钻头在地上碾出一道好几丈深的沟，朝着胡峰飞出去的方向碾了过去，快得声音都追不上。沟两边的地让它一碾，碎了，冻了，成了粉。白霜从沟里往两边扑，扑出去几十丈，扑到哪儿哪儿冻得咔咔响。

碎紫斗罗也动了。

刚才那一炸，他也挨着了。紫色的小人身上全是裂纹，裂纹比好的地方还多，身子里的魂力剩不了几口。可他的两只眼睛亮得吓人。

``超载\ldots\ldots 再来一回。''

他把身子里最后那一点魂力全塞进了极速模式里。紫色的小人亮了最后一下，亮得刺眼。

``虚空破碎。''

碎紫斗罗不见了。

他快到了头，快得把自己挤进了虚空里头。他在虚空里穿，每闪一下就跨出去几十丈，在实打实的这边，留下一道一道紫色的裂痕。

那些裂痕是空间本身让他撕开的口子，一时半会儿合不上。裂痕边上闪着紫光，像夜里一道一道紫色的闪电，在胡峰往后飞的那条道上交在一起，拢在一起，织成了一张大网。

碎紫斗罗在胡峰头顶上冒了出来。

那张网一下子收紧了。

轰！

冰熊王的死亡翻滚从正面碾过来。碎紫斗罗的虚空破碎从上头压下来。

两样东西在胡峰的落点上一齐炸开。

冰蓝的光刃跟紫色的裂痕缠成一团，缠成一个大旋涡。旋涡正当中，胡峰让两边一夹，像一片秋风里的叶子，让碾，让撕，让甩。

死亡翻滚的光刃在他身上碾出了几十道口子，口子深得能看见骨头，血一冒出来就让冷冻住了，冻成暗红的冰碴子。虚空破碎的裂痕在他身上又划出了几十道细口子，一道比一道深，深到骨头缝里。

胡峰让这两下子打得飞出去好几百米。

他在半空里划了一道弧，血在身后拖了一条长长的红线。紫金冠早碎了，明黄的太子常服成了一条一条的布片子，底下的身子上全是伤，全是青紫。那件暗金的内甲让死亡翻滚一碾，碎得一片一片的，碎片在半空里飘，像一群蝴蝶。

他眼睛里那两团金火苗，灭了。

0 111 222 322 42 在他身上一节一节地散了，身上的暗金纹路一根都不剩。他身上那股吓人的气，一下子掉到了底，像一颗从天上掉下来的星星，在夜里拖着最后一道光。

轰。

胡峰砸进了几百米外的一片废墟里，砸出了一个好几丈深的人形坑。碎石跟尘土在他身边扬起来，把他埋住了。

坑里没动静了。

\begin{center}\rule{0.5\linewidth}{0.5pt}\end{center}

战场一下子静了。

冰熊王在几百丈外停下了翻滚，一百五十米的身子让惯性又带着往前滑了几十丈，才停住。四只爪子插进地里，身子一起一伏地喘。刚才那一下死亡翻滚，把它的力气也抽空了。真霜斗罗站在它头顶上，脸白得跟纸一样。冰牢垒了那么久，又来了这一下，他身子里的魂力已经榨到了底。

碎紫斗罗在胡峰落的那个坑上头冒了出来。紫色的小人一冒出来，就碎了，碎成一片紫色的光点子，让夜风一吹，散了。碎紫斗罗本人从光点里掉下来，一条腿跪在了地上，大口大口地喘。脸白得没有一点血色，嘴角挂着血，身上的魂力一滴都不剩了。

可他还活着。

他抬起头，隔着几百米，往那个坑里看。

``完了\ldots\ldots 没有？''他的嗓子哑得厉害。

真霜斗罗没答。

他看见了。坑沿上，有一只手从碎石头底下慢慢伸了出来。

一只全是血、全是青紫的手。手指头在抖。

可它还在动。

那只手在碎石头里慢慢攥紧了，指甲掐进了手心，血从指头缝里渗出来。

然后，碎石头的缝里，露出来一只灰白的眼睛。

那只眼睛里的金火苗已经灭了，就剩一片空空的灰白。可那片灰白最底下，有个东西正在醒，那东西比金火苗还吓人。它里头一点人的意思都没有了。

\begin{center}\rule{0.5\linewidth}{0.5pt}\end{center}

胡帅浮在半空里。

身上那点白光一明一灭，跟坏了的灯笼一样。帅气真身早散了，蟒袍没了，暗金的内甲上全是裂口，血从裂口里往外渗，渗得他半边身子都是暗红的。

他低下头，往胡峰那个坑里看了一眼。

又抬起头，往四边看了一圈。

真霜站在那头瘸了腿的冰熊头上，冷冷地看着他。碎紫一条腿跪在坑边上，魂力空了，脸白得像纸，可眼睛还瞪着他。半空里，那个丫头踩着一块石板往下落，两只眼睛的金色快褪干净了，可她还站着。巷子口那边，那头灰狼，那只冰蝎子，那个光脚的姑娘，三个慢慢往一块儿聚，三双眼睛全朝着他。

胡帅的脸沉下去了。

他心里那本账，一下子就算完了。他这会儿这个样子，同时对付一个真霜、一个碎紫，再加一窝不知道还藏着什么的孩子，打不下来。胡峰已经不能动了。再在这儿耗，一点好处都捞不着。

``撤。''

他就说了一个字。

他身上的白光猛地一炸，炸成一根光柱。他整个人在光柱里往下一扎，扎进了胡峰那个坑里，又从坑里冲出来，胳肢窝底下夹着一个血糊糊的人。

``想跑？''真霜斗罗一催冰熊王，一百五十米的身子朝那根光柱压过去。

胡帅太快了。

光柱在夜空里一拐，化成一道白线，往皇城那边射了过去，快得声音都跟不上。冰熊王身子大，腿又瘸了，论快，差得远。几口气的工夫，那道白线就钻进了皇城的轮廓里，不见了。

坑里只剩下了一个人形的印子，印子里一滩血。

远处地上，震天斧还插在那儿。斧柄抖了一抖，化成了一团暗金的光，往皇城那边飞过去了。

``追不上了\ldots\ldots{}''碎紫斗罗撑着地，想站起来，两条腿一个劲儿地打颤，又跪了回去。

真霜斗罗站在熊头上，看着皇城那边。

他没追。

追，追不上，这是一回事。他心里还压着另一件事。

胡帅往皇城里跑，不光是为了逃命。

城外那三千个人，刚才退到了几条街口，一排一排地站齐了，就是不往这边冲。二十几个魂斗罗站在最前头，也一直不动手。他们是在等。

等胡帅回去。

真霜斗罗在天斗城里住了一百多年，这座城有几道墙，几道门，墙底下埋着什么，他一清二楚。皇城那几道墙，是老皇帝那辈子请人修的，墙根底下埋着一圈阵眼，本来是拿来守城的。

要是让人改过了呢。

``所有人！''真霜斗罗的声音一下子拔高了，又急又硬，``往陛下那儿靠！马上！''

八个人从烂瓦堆里，从碎石头里，从地上，一个一个地往起爬。

``胡帅在皇城里布了阵。''真霜斗罗说，``他跑回去，就是要把阵起起来！''

\begin{center}\rule{0.5\linewidth}{0.5pt}\end{center}

蔡须坤从石板上跳下来，落地的时候腿一软，差一点跪下去。

她扶着一截断墙站住了。身子里头空落落的。那一份处女之力放出去了，炮弹炸了，0-Y 序列顺着炸开的那一下往她身子里灌了一阵，又散了。散完了，就剩下一个空壳子，手脚又沉又软，像在水里泡了一整夜。

大长老说过，用的人要虚好几天。

她现在知道是什么滋味了。

唐三一瘸一拐地走过来，扶住了她的胳膊。

``还能走？''

``能。''蔡须坤说。

奥斯卡让诺当拽着胳膊从大石头边上拖起来，脑袋还晕着，眼睛眯成一条缝，往蔡须坤这边看。

``队长。''他说，舌头有点大，``你刚才扔的那个，是大长老给的那个？''

``嗯。''

``那你\ldots\ldots{}''

``虚几天。''蔡须坤说，``死不了。''

奥斯卡张了张嘴，没再问。他把手伸进袋子里摸了摸，摸出来一根最普通的恢复香肠，掰成两截，一截塞进她手里，一截塞进自己嘴里。

蔡须坤看了看手里那半截香肠，咬了一口。

她扭头往巷子口看。母猪已经翻过身来了，四个蹄子撑着地，站在担架前头，身上全是灰。小舞跪在担架边上，两只手护着大帝的脑袋，冲这边拼命地招手。

``陛下醒了！''小舞喊，``眼睛睁开了！''

几个人互相搀着，往巷子口跑。

真霜斗罗的话没说错。

\begin{center}\rule{0.5\linewidth}{0.5pt}\end{center}

\chapter{第七十四章：养猪场第一股东phyrion}\label{ux7b2cux4e03ux5341ux56dbux7ae0ux517bux732aux573aux7b2cux4e00ux80a1ux4e1cphyrion}

几个人跑到巷子口的时候，大帝已经让小舞扶着坐了起来。

他靠在母猪的肚子上，身上裹着那床薄被，脸还是白的，颧骨底下两个坑。眼睛睁着，睁得很费劲，眼皮一会儿往下耷拉一下，又硬撑开。他看着围过来的这几个满身是血的孩子，嘴唇动了动。

``\ldots\ldots 真霜呢？''

``在后头。''蔡须坤蹲下来，``陛下，您别说话。''

大帝没听她的。他把一只手从被子里伸出来，往宫门那边指了指。

``城楼。''他说，``上城楼。''

真霜斗罗正好赶到，冰熊王一瘸一拐地跟在他身后。他听见了这两个字，点了一下头。

``陛下说得对。''他说，``站高了，看得见阵眼。''

宫门两边各有一座城楼，城楼底下的台阶有一百来级。诺当跟蔡兴坤抬着担架，唐三的蓝银草在担架底下垫着，一级一级往上走。碎紫斗罗让奥斯卡架着一条胳膊，走一级喘一口。真霜斗罗没让冰熊王上来，那么大的身子，一脚就能把城楼踩塌。他从熊头上跳下来，自己走。

走到一半，皇城正当中，冲起来一根光柱。

金色的，粗得吓人，从皇城最当中那片宫殿里拔地而起，一直冲到天上。半边天让它染成了暗金色。

八个人一齐站住了。

光柱冲到云彩里，猛地一炸。

炸开的光没散。它往四边铺，铺成了一整张天幕，从头顶上罩下来，把整座皇城扣在了底下。天幕上密密麻麻刻满了暗金色的帅气纹路，纹路当中流着黑帅气本源的那股味儿，跟胡峰身上那股一模一样。

阵起来了。

天幕底下，从皇城的四面八方，涌出来一群人。

二十几个，全穿着暗金的甲，脸上扣着铁面具。看不见脸，只看得见面具上两个窟窿，窟窿里往外冒着光。霍雨浩的精神力往那边一探，脸色就变了。

``魂斗罗。''他说，``最弱的也是 Y(1,2,4,8,16)。最强的那个，Y(1,2,4,8,16,30)。''

二十几个人身上冒着同一种味儿，跟胡峰那股黑帅气是一个窝里的。他们按着一个说不出名堂的阵势站开，站成一个大圆圈，魂力在一个跟一个之间流，流成了一个环。

紧跟着，是那三千甲兵。

刚才退到街口的那三千人，这时候从各条街上涌进来，在二十几个魂斗罗外头排成了一个大方阵。一个人的修为不高，可三千个人的魂力让阵一拧，拧成了一条河，跟里头那个环接上了。环里的劲儿一下子涨了一大截。

八个人抬着担架，爬上了城楼。

大帝让他们扶着，靠在城楼的垛口上，往底下看。

他的脸一点一点地沉下去，沉得发青。

蔡须坤站在他边上，两只眼睛里那点金色早褪光了，只剩下自己原来的颜色。她看着那张从天上扣下来的暗金天幕，嘴唇发白。

``阵。''她说，``他把整座皇城弄成了一个阵。''

真霜斗罗扶着垛口，一级一级地喘匀了气，往底下看。

二十几个魂斗罗，三千甲兵，再让这个阵一拧，合起来的魂力，已经比一个封号斗罗到头的分量还多。他自己呢，刚才跟胡帅那一场，Y(1,3,6,12) 早退回去了，身子里的魂力剩不到四成。冰熊王的右腿还在淌血，站在城楼底下，一个劲儿地舔伤口。碎紫靠在垛口边上，坐着都费劲。

``真霜。''大帝开口了，声音沉，一个字一个字地往外挤，``这个阵，你破得了么？''

真霜斗罗没马上答。

他看着底下那张天幕，看着二十几个魂斗罗的魂力在环里转圈，看着三千甲兵的魂力一道一道地往环里汇，汇进去，环就厚一层。

要是他身子是满的，有 Y(1,3,6,12) 撑着，兴许能硬砸开一个口子。

这时候\ldots\ldots{}

``\ldots\ldots 破不了。''

他说得很平，很干脆。

城楼上一下子没声了。

封号斗罗说破不了。那就是说，他们出不去了。

底下，暗金的天幕在皇城上头慢慢地转。转一圈，就厚一层。魂斗罗们的魂力在阵里一圈一圈地绕，绕得一刻不停，像一架不知道累的水车。

三千甲兵齐声吼了一嗓子。

暗金的魂力从三千个人身上一齐冒出来，冒成三千道金光，射进了天幕里。天幕让这三千道光一灌，越来越亮，越来越厚，像一口倒扣下来的金锅，把城楼上的人罩得死死的。

``它在往里收。''霍雨浩说。他的精神力贴着天幕在摸，摸一圈，脸就白一层，``四边都在往里收\ldots\ldots 最多半个时辰，就收到咱们这儿了。到那时候\ldots\ldots{}''

他没往下说。谁都知道到那时候会是什么样子。一口锅往里收，收到最后，锅里的东西让它挤成一团。

``有人能来么？''蔡须坤扭头问大帝，``陛下，天斗城还有别的封号斗罗么？''

大帝摇了摇头，脸色灰得跟城砖一样。

``他起阵之前，四门早就让人封死了。''他说，``城里能动的封号斗罗，就剩真霜一个。''

城楼上没人说话了。

真霜斗罗闭上了眼睛。城楼底下，冰熊王低低地哼了一声，声音发闷，右腿一抖一抖的。

头顶上那口锅，还在一点一点地往下压。

就在这时候，皇城东边，传过来一声嚎。

那一声不像人，也不像寻常的魂兽。粗，野，又蛮又横，嚎得理直气壮，像一头很老很老的大东西在跟谁宣告：这块地，归我了。

城楼上所有人一齐扭过头，往东边看。

皇城东边那道城墙，让大阵罩着，本来该是一点缝都没有的。这时候那道几丈厚的暗金城墙在抖，抖得厉害。墙那一边，有一个很大很大的粉色影子，在往这边挪。影子太大了，大得把半边天都遮住了，月亮让它挡在了后头。

轰！

城墙炸了。

几丈厚的城墙碰上那股蛮劲儿，跟纸糊的一样，一下子碎成了漫天的石头跟土。烟尘里冲出来一个粉色的大东西，像一座会走路的粉色山头，冲着大阵这边碾了过来。

一头母猪。

一头一百米高的母猪。

它比真霜斗罗的冰熊王还大出一倍多。四条腿跟四根柱子一样，一只蹄子就有好几丈宽，落一下地，整座皇城都跟着颤一下。一身粉红的皮上爬着一层细细的金纹路，像拿化开的金子浇在上头。两颗獠牙往外翻着，一颗就有十几丈长，牙尖上闪着铁一样的冷光。两只眼睛血红血红的。

这头母猪的背上，站着一个人。

个子高大，穿着一身深灰的西装，剪裁得板板正正。领口系着一条暗红的领带，袖口露出一对金色的袖扣。头发梳得一丝不乱，脸上棱是棱角是角，眉毛底下压着在上头待惯了的人才有的那种稳当。左手腕子上戴着一块魂导表，右手随随便便地插在裤兜里。

他站在一头百米母猪的背上，跟站在自家办公室里差不多。

城楼上，八个人看着那个人，脸色一齐变了。

蔡须坤的嘴张开了，半天才挤出两个字：``卢叔\ldots\ldots{}''

``是我爸！''

卢本韦一嗓子喊了出来，喊得劈了。她本来让蔡兴坤扶着靠在垛口上，这一下自己站直了，左肩上的伤好像一下子不疼了。

唐三瞪着眼，嘴张了又合，合了又张：``养猪场的卢关？他怎么来了？''

``第三回了。''霍雨浩说。他看着那个穿西装的人，声音很低，``这是咱们第三回见他。''

第三回。

头一回见卢关，是梅天狸长老把他们从死地里捞出来以后。那时候八个人一身的伤，卢本韦说，去我家吧。她家就是养猪场。他们在养猪场的总部见到了她爸，一个坐在大办公室里、穿着笔挺西装、翘着二郎腿、手里端着一杯现磨咖啡的男人。他笑眯眯地听完了他们的来意，起身带着他们把整个养猪场转了一圈。那一回他们头一次见到了那头百米母猪，站在猪圈外头仰着脖子看了半天。卢关这个人到底有多深，那一回他们没看出来，只觉得看不透。

后来又遭了一回大难，几个人拖着一身伤，又回了养猪场。卢关站在大门口等他们，西装还是那么挺，脸上还是那副笑。他挨个拍了拍他们的肩膀，多余的话一句都没说，就一句：回来了就好。

这一回，他骑着那头母猪，撞塌了皇城的城墙。

养猪场离天斗城一百来里地。路不算远，可要把一头百米母猪、两百号人从养猪场拉到这儿，再赶在大阵合拢之前撞进来，碰巧路过，路过不成这样。

要么有人提前叫了他，要么他自己算着日子，一直在等。

蔡须坤看着那张不慌不忙的脸，嗓子发紧。

``卢叔叔\ldots\ldots{}''她小声说了一句，``谢谢。''

一百米外，猪背上那个人好像听见了。他偏过头，往城楼这边看了一眼，冲城楼点了一下头。

点得很轻，很随便，像老熟人在街上碰见了，打个招呼。

八个人都看懂了那一下是什么意思。客气什么。

母猪后头，烟尘里，一队一队的人跟着涌了进来。两百来号，穿着一样的灰布褂子，手里提着铁锨、扛着粪叉，打头那个腰上还拴着一只喂猪的瓢。一个个精壮，走起路来脚底下带风。养猪场的人。

队伍最前头，走着另一个人。

他没骑猪，也没穿西装。身上是一件深蓝的长袍，洗得有点发白了。脸清瘦，冷，眼睛半闭着。他走在两百来号壮汉前头，身子骨看着比谁都单薄，可他一个人往那儿一走，身后那两百号人的动静，全让他压住了。

像一潭水，看不见底。

卢本韦看着那个人，嘴唇哆嗦了一下。

``\ldots\ldots 爷爷。''

phyrion。卢关的叔父，卢本韦管他叫爷爷。

Y(1,3,10)。

phyrion 的扽魂，也是母猪。

可他那头母猪跟卢关那头两样。卢关那头大，横，一身的蛮劲儿往外冒。phyrion 把自己那头母猪身上多余的东西一样一样都削掉了，削了一辈子，个头、花样、架子，统统削没了，最后就剩下一样。

往前撞。

phyrion 的眼睛慢慢睁开了。

一双粉红的眼睛。卢关那头母猪的眼是血红的，凶。他这双眼是粉的，润，像开春以后枝头上的桃花。可那点粉往里看，看到最底下，城楼上每一个魂师的后脖颈子都一凉。

他抬起头，看了一眼头顶上那张暗金天幕，又往下看了一眼二十几个魂斗罗跟三千甲兵摆的阵，嘴角往上动了一下。

``这就是胡帅的阵？''他说。声音平平的，像在说养猪场新进的一批饲料，``还行。''

他把他的母猪叫了出来。

卢关那头母猪一出来，地都在颤。phyrion 这一头，就那么安安静静地，在他脚边上冒了出来。

一头小猪。

跟谁家圈里养的家猪差不多大，通身雪白，一根杂毛都没有。皮光溜溜的，像一块白玉。身上没有金纹，没有魂力往外冒，连一点气都闻不着。它就那么站在 phyrion 腿边上，拿一双粉红的圆眼睛东瞅西瞅，瞅瞅天上那张金锅，又瞅瞅底下那些戴铁面具的人。

城楼上，真霜斗罗扶着垛口的手一下子攥紧了。碎紫斗罗本来半闭着眼，这时候眼睛一下子睁开了。

能摸到封号斗罗那一层的人，这时候都觉出来了。

那头小白猪身子里，压着东西。

压得很紧。一身的劲儿全往里收，收了又收，收了再收，收得越来越小，越来越密，最后收成了一个点。一点都不往外漏。外头看着什么都没有，可那个点里头攒着的分量，真霜斗罗光是隔着几百丈往那边一摸，后背上的汗就下来了。

``那\ldots\ldots 那是什么？''奥斯卡的嘴张得老大，手里攥着的那根超级大香肠掉到了脚面上，他都没觉出来，``一头\ldots\ldots 小猪？''

``别小看它。''真霜斗罗的声音很低，``那头小猪身子里的东西，比卢关那头百米的还吓人。''

``啥？''卢本韦瞪圆了眼，``比我爸那头还\ldots\ldots 那它要是冲起来\ldots\ldots{}''

她的话没说完。

phyrion 弯下腰，在小白猪的背上轻轻拍了一下。

``去吧。''

小白猪歪了歪脑袋，像是在琢磨这句话。

然后它迈开了腿。

头一步。蹄子落在青石板上，嗒的一声，脆生生的。

就这一声嗒，整座皇城的地面都颤了一下。

第二步，它跑起来了。

看着一点都不快，小短腿一倒一倒的，跟村口追着鸡跑的小猪没两样。可它每迈一步，离那个阵就近出去一大截。五十丈，二十丈，十丈。

阵里那二十几个魂斗罗也看见它了。

``什么玩意儿？''站在最前头那个魂斗罗冷笑了一声，铁面具底下的声音瓮声瓮气的，``一头猪？派一头猪来破阵？''

他抬手一挥。旁边几个魂斗罗一齐出手，暗金的魂力拧成了几条又粗又长的铁链子，哗啦啦地朝小白猪缠过去。

小白猪没躲。

它连看都没看那几条链子一眼。

它就是接着跑。

链子碰上它身子的那一下，嗞的一声，像冷水泼在烧红的铁板上，一下子全化了，化成满天的魂力渣子。小白猪身上一根毛都没乱。

那几个魂斗罗的脸色，隔着铁面具都看得出变了。

``什\ldots\ldots{}''

小白猪没停。

它跑到了阵的边上。

一头撞了上去。

没有魂技，没有光，连蓄力都没有。一头猪，低下脑袋，往前一拱。

就这么一下。

撞上的那一下，四下里一点声音都没有了。声音也让它撞碎了。

暗金的天幕上，小白猪拱的那个地方，凹进去了一个小坑。

小坑在往大里长。拳头大，磨盘大，几丈宽。

轰！

天幕碎了。

从小白猪拱的那一点开始，往四面八方裂，像一块玻璃让一颗石子砸中了，裂纹一眨眼就爬满了整张天。暗金的碎片哗啦啦地往下掉，掉到半空里就化成了金色的光点子，一点一点地飘，一点一点地散。

胡帅拿二十几个魂斗罗、三千甲兵撑起来的这个阵，碎了。

一下。

城楼上，八个人全愣住了。

唐三的嘴张了又合，合了又张，半天一个字都没挤出来。

``我的精神力\ldots\ldots{}''霍雨浩两只眼睛直勾勾的，``摸不到那头猪的底。往里探，探多深都是满的\ldots\ldots{}''

``那是猪？''卢本韦的声音在抖，``那真是一头猪？我从小在猪圈里长大的，我咋没见过这样的猪？''

蔡须坤站在垛口边上，看着小白猪拱出来的那道口子。那道口子从阵边上一直拉到阵的另一头，笔直笔直的。她嘴唇发白，过了好一会儿，才说了一句：

``\ldots\ldots 养猪场的水，到底有多深。''

\begin{center}\rule{0.5\linewidth}{0.5pt}\end{center}

还没完。

小白猪撞碎了天幕，脚底下一点没停，一头扎进了魂斗罗跟甲兵摆的阵里头。

百米母猪跟在后头，从那道口子里冲了进去。卢关站在猪背上，西装的下摆让风吹得往后扯。他往底下看了一眼，右手从裤兜里抽出来，随手挥了一下。

``清场。''

就两个字。

百米母猪仰着脖子嚎了一嗓子，四个蹄子刨着地，一座粉色的山一样碾进了阵里。两百来号养猪场的人扯着嗓子一齐吼，举着铁锨粪叉跟着冲了进去。

小白猪还在撞。

它的个头还是那么小，可身子里那个点，这时候放开了。它像一颗粉白的流星，在阵里穿过来，穿过去，每穿一回，就带走一大片暗金的光。

十六个魂斗罗，从 Y(1,2,4,8,16) 到 Y(1,2,4,8,16,30)，让它撞上了，当场就没了。

暗金的甲碎了，铁面具瘪了，身子里的魂力让它一拱，搅成了一锅粥，散了。有的连一声都没来得及叫，人就化成了一片光点子，让夜风吹走了。

剩下那几个魂斗罗，跟阵里那几个魂圣，让小白猪的余劲儿扫着了，一个个飞出去，甲碎了，嘴里往外喷血，脸白得跟纸一样。有个魂圣让余劲儿正扫在腰上，浑身的骨头断了一大半，瘫在地上，再也没爬起来。

三千甲兵更惨。

小白猪那点余劲儿，对魂斗罗都这么狠，落到修为差得远的甲兵身上，就跟镰刀割麦子一样，一排一排地往下倒。暗金的甲碎了，人也碎了。

十来口气的工夫，三千甲兵，死的死，伤的伤，倒下去一大半。

剩下的那些，崩了。

他们扔了刀，扔了戟，扒了身上的甲，扭头就跑。一个跑了，后头就跟着跑了一大片，满街都是往外逃的人，你踩我，我踩你。

小白猪在阵的另一头停了下来。

它歪着脑袋，拿那双粉红的圆眼睛看了看四下里乱跑的甲兵，又看了看地上那些再也不会动的魂斗罗，鼻子里哼哼了两声。

那两声哼哼，听着像一头猪在嫌今天的食儿不好吃。

卢关骑着百米母猪走了过来。一身西装，脚底下是一头粉色的大山，怎么看怎么不搭。他从猪背上一跳，皮鞋落在青石板上，嗒的一响。他扯了扯领带，往满地的狼藉上扫了一眼，最后看向 phyrion。

``叔父。''他说，口气里带着一点请示，``要不要追？或者喊两嗓子，让剩下那些降了？''

phyrion 往四下乱跑的那些甲兵身上扫了一眼，又往皇城深处看了一眼，那是胡帅逃进去的方向。

``不用追。''他说，``跑不远。''

他停了一下，又补了一句。

``能拿力气解决的事，别跟这种人费口舌。''

卢关愣了一下，跟着就笑了，点点头。他转过身，冲着那两百来号人一挥手。

``打扫干净。活的留下，死的不用管。''

``是！''两百来号人一齐应了一声，散开了，动作麻利得很，一看就是平常在养猪场里练出来的。

\begin{center}\rule{0.5\linewidth}{0.5pt}\end{center}

城楼上，卢本韦已经往台阶底下冲了。

她一级跨两级，跑得跌跌撞撞，左肩上那道伤口一扯，又裂开了，血顺着胳膊往下淌，把半边衣裳又浸了一遍。她顾不上，从城楼上一路冲下来，冲过满地的碎石头，往那头百米母猪那边跑。

``爸！''

那一声喊出来，抖得厉害。这一嗓子她在外头憋了好些日子了，到她爸跟前，才敢喊出来。

卢关听见了。

他嘴角那点总裁的笑，一下子僵住了。

他低下头看过去。一个浑身是血的丫头正跌跌撞撞地朝他跑过来。左肩在淌血，衣裳烂了一大半，脸上全是灰跟血，头发糊在脑门上，要多狼狈有多狼狈。

可那双眼睛，跟他一模一样，倔，亮，隔着一层泪看着他。

卢关脸上那副样子，一下子全垮了。

什么总裁，什么在上头待惯了的稳当，一下子全碎成了渣。

他往前迎过去，皮鞋踩在碎石头上，一个趔趄，差点没站稳。他三步并成两步，在离她还有三步远的地方，一把把她揽了过来。

``你这丫头\ldots\ldots{}''他的嗓子一下子哑了，胳膊收紧，把她整个人箍在怀里。那条暗红的领带蹭上了她的血，他看都不看，``怎么弄成这样？伤哪儿了？让爸看看\ldots\ldots{}''

``爸，我没事\ldots\ldots{}''卢本韦把脸埋在他胸口上，声音闷闷的，带着鼻音，``就左肩挨了一下，不疼\ldots\ldots{}''

``不疼？流这么多血还不疼？''卢关又急又心疼，一只手按在她左肩上，魂力从手心里淌出去，一点一点地往伤口里探。探了一圈，骨头没断，他才长长地吐了一口气。胳膊还是没松。

卢本韦在他怀里扑哧笑了一声，笑着笑着，眼泪就下来了。

``爸，你真厉害。''她闷着声音说，``那头大猪一冒头，我就认出来了。''

卢关拿手揉了揉她的头发，揉得乱七八糟。

``你更厉害。''他说，声音低，里头压着一点得意，``爸在养猪场都听说了。八个娃娃，硬扛胡家那个小子，还把他打得爬不起来。''

卢本韦的眼泪掉得更凶了，她咬着嘴唇，没让自己哭出声。

``本韦。''

身后有人叫她，声音平平的。

卢本韦从她爸怀里抬起头，泪眼模糊地往那边看。

phyrion 站在几步外头，深蓝的长袍让夜风吹得轻轻地飘。他脸上那股冷，这时候一点都看不见了，刚才那一身叫封号斗罗都后背发凉的东西，全收了回去。这会儿他看着，就是个瘦瘦的、上了年纪的老头。

他手里端着一碗汤。

是养猪场的人拿随身带的锅灶现熬的排骨汤，汤面上漂着几片葱花，热气往上冒，在夜风里一歪一歪的，香得人鼻子发酸。

``\ldots\ldots 爷爷。''卢本韦的声音一下子全软了，眼泪一颗接一颗地往下掉。

phyrion 走过来，把汤碗搁进她手里。

``先喝汤。''他说，声音平平的，跟刚才那句能拿力气解决的事别费口舌，像两个人说的，``肩膀上的伤不重，血流得多了。喝完了让你爸给你治。''

卢本韦两只手捧着碗。碗是烫的，烫得她冰凉的手指头一下子缓过来了。她低头喝了一口，汤顺着嗓子往下走，走到肚子里，整个人都热了。

phyrion 伸出手。

那只手又瘦又有劲儿，指节上全是老茧。他把她脑门上糊着的那几绺头发拨到耳朵后头去，拨得很轻。这只手刚才拍了一下猪背，那头猪就撞碎了一座城的阵。

``瘦了。''他说，声音很轻，``在外头吃不好吧。回去让灶上多给你做点。你爱吃的红烧猪蹄，明天就做。''

卢本韦喝着汤，眼泪又下来了。这一回是笑着掉的。

``爷爷，你那头猪真厉害。''她一边抽搭一边说，``比我爸那头大的还厉害。''

phyrion 的嘴角往上动了一下。

``它啊。''他低头看了一眼脚边。那头小白猪正拿鼻子拱地上的砖缝，拱得哼哼唧唧的，``平常在养猪场闲着，就爱撞东西玩。撞着撞着，撞出点名堂来了。''

小白猪好像听懂了，抬起头，拿那双粉红的圆眼睛看了看卢本韦，哼哼了两声，慢悠悠地走过去，拿脑袋在她小腿上蹭了蹭。

温温的，热热的。

卢本韦低头看着脚边这头刚把一座大阵拱碎了的小白猪，这时候跟一条小狗似的蹭她的腿，一下子没忍住，笑出了声。

``爷爷，我能摸摸它不？''

``摸吧。''phyrion 说，``它不咬人。''

卢本韦蹲下来，伸出手，小心翼翼地在小白猪脑袋上摸了一下。小白猪舒服得眯起了眼，哼了一声，整个身子往地上一趴，摊成一团白，像一只大猫。

卢关站在边上看着，脸上挂着笑。他把那条沾了血的领带又扯了扯，往满地的碎砖烂瓦上扫了一眼，又落回女儿身上。

``叔父。''他低声说，``本韦跟着这帮孩子，吃了不少苦。''

phyrion 没回头。他就站着，看着蹲在地上摸猪的卢本韦，那双粉红的眼睛里映着她小小的一个影子。

``本韦是卢家的孩子。''他说，``卢家的孩子，吃不着苦。''

他停了停。

``只会长大。''

\begin{center}\rule{0.5\linewidth}{0.5pt}\end{center}

城楼上。

真霜斗罗扶着垛口，一直盯着 phyrion 的后背。

他脸上的颜色变了好几回，最后定在一种说不上来的样子上，像一个老头子回头看自家后院，看见一棵自己从没注意过的树，已经长得比房子还高了。

Y(1,3,10)。

他在嘴里把这几个字又念了一遍。

Y(1,3,10)，Y(1,3,6,12)。摆在纸上，前头都是一个一，一个三，后头的数也没差多少。外行人看了，兴许觉得也就高出去一截。

真霜斗罗比谁都清楚这一截有多长。

Y(1,3,6) 再往上，是 Y(1,3,7)，Y(1,3,8)，Y(1,3,9)。每往上挪一个数，当中都夹着一整条他走不完的路。他自己在 Y(1,3,6) 上站了几十年，往 Y(1,3,7) 那边看，看得见，够不着。phyrion 站的那个地方，在 Y(1,3,9) 那一长串的尽头还要再往上。

他抬头看天，phyrion 站在天上头。

再看看他自己。刚才跟胡帅那一场打下来，魂力去了一大半，冰熊王的腿还瘸着，借来的 Y(1,3,6,12) 早退回去了。这时候真动起手来，他能使出来的，连 Y(1,3,5) 都够呛。

Y(1,3,5) 对 Y(1,3,10)。

真霜斗罗闭上眼，长长地吐了一口气。

能拿力气解决的事，别跟这种人费口舌。

他这才咂摸出这句话的味儿来。说这话的人一点都没摆架子。站在 Y(1,3,10) 上往下看，他们刚才拼了命、差点把命搭进去的那个阵，就是一头猪拱一下的事。他说的是实话。

那个人站得太高了。

高得真霜斗罗活了一百多年，头一回觉得自己矮。

``前辈\ldots\ldots{}''碎紫斗罗靠在垛口边上，撑着想坐直，嗓子哑得快听不见了，``那一位\ldots\ldots 就是养猪场的大股东？''

``嗯。''

``他到底\ldots\ldots 什么境界？''

真霜斗罗没马上答。他看着底下那个瘦瘦的深蓝色的背影，看了一会儿。

``Y(1,3,10)。''他说，声音很轻，可每个字都沉，``比老夫，高了整整两个大层次。''

碎紫斗罗的脸色一下子变了。

一个封号斗罗，当着人说另一个人比自己高两个大层次。这话有多重，碎紫斗罗不用想就知道。

\begin{center}\rule{0.5\linewidth}{0.5pt}\end{center}

剩下几个人站在城楼上，往底下看。

卢本韦蹲在地上，一只手捧着汤碗，一只手摸着小白猪的脑袋。她爸站在边上笑。她爷爷站在另一边，不说话。

蔡须坤的目光在卢关身上停了很久。

西装笔挺，不慌不忙，嘴角挂着那点笑。跟头一回在养猪场办公室里见他的时候一样，跟那回他站在大门口说回来了就好的时候也一样。

三回了。三回都是在他们快撑不住的时候。

她以为自己已经把这个人看明白了一些。这一回又发现，还差得远。

``这就是\ldots\ldots{}''唐三的嗓子发涩，``顶上那一层的人么\ldots\ldots{}''

没人接他的话。

几个人心里都在过同一件事。他们拼了命才摸到的那块天花板，在 phyrion 脚底下，就是一块地砖。

可卢关站在 phyrion 边上，一点都不缩。不低头，也不往前抢，就那么站着。

奥斯卡凑到蔡须坤身边，压着嗓子说：``队长，往后\ldots\ldots 咱们是不是得抱紧卢本韦的大腿？''

蔡须坤没理他。

她的目光越过 phyrion 的后背，往皇城深处看。那是胡帅钻进去的地方。

``走。''她说。声音已经稳下来了，``胡帅还没逮着。这一仗，还没打完。''

底下，卢关好像听见了。他把卢本韦从地上拉起来，拿拇指替她把脸上的泪痕擦了擦。又整了整西装的领子，脸上那点当爹的软一收，又成了那个总裁。他抬起头，冲城楼上点了点头。

``八位小友。''他说，声音不高，稳，``谢谢你们照看本韦。这份情，我卢关记下了。''

他转过身，往皇城深处走。

``叔父，咱们进去转一圈，看看还剩下什么。''

phyrion 跟在他身后，深蓝的长袍在夜风里飘。小白猪一溜小跑地跟在 phyrion 脚边，粉红的圆眼睛东张西望。

卢本韦跟在最后。她爸拿魂力先把她左肩的伤口封住了，她手里还捧着那碗没喝完的排骨汤，走了几步，回过头，冲城楼上咧开嘴笑了。牙上还沾着汤。

``等我回来，请你们去养猪场吃红烧猪蹄！''

蔡须坤看着她的背影，嘴角往上动了一下。

``好。''

\begin{center}\rule{0.5\linewidth}{0.5pt}\end{center}

\chapter{第七十五章：Y(1,3,10)}\label{ux7b2cux4e03ux5341ux4e94ux7ae0y1310}

皇城上头的烟，还没散干净。

大阵碎下来的那些暗金光点，还有一些在半空里飘，一点一点地往下落，落到半截就灭了，像一场下不到地上的雪。远处废墟里，冰熊王身上冒出来的冷气一阵一阵地往上走，走不高就散了。

碎紫斗罗让人从城楼上扶了下来，盘着腿坐在一堆碎瓦上，脸白得跟纸一样，闭着眼在往回攒魂力。攒了半天，攒回来的那点，刚够他把腰挺直。

真霜斗罗站在冰熊王的头顶上，看着皇城深处。卢关他们三个往里走了，走了快半个时辰，还没出来。

他一直没说话。

冰熊王趴在废墟里，一百五十米的身子蜷着，那条伤腿一抖一抖的。它时不时低低地吼一声，像是腿疼，又像是不服气。刚才城楼上那一句破不了，它在底下也听见了。

过了好一会儿，真霜斗罗叹了一口气。

``这个家伙\ldots\ldots{}''

碎紫斗罗睁开眼，隔着满天飘着的光点子，往他看的那个方向看。

``前辈说的是 phyrion？''

真霜斗罗没答。他看着皇城深处，那里黑漆漆的，什么都看不见。

``他比老夫小不了几岁。''过了一阵，他说。

碎紫斗罗没接话。

``可他这一身本事\ldots\ldots{}''真霜斗罗的声音沉下去了，``老夫不服都不行。''

到了他们这一层，魂力多少，魂力云厚不厚，早就不顶用了。顶用的，是身上那一串 Y 序列，能往前走到哪儿。

``Y(1,3,10) 是什么意思，''真霜斗罗说，``你该知道。''

碎紫斗罗沉默了一会儿，点了点头。

``知道。''

真霜斗罗伸出一只手。

冰蓝的魂力在他手心里聚起来。四下里一下子冷了，半空里的水汽在他指头尖上凝成了细细的冰碴子。冰碴子让魂力一拨，一颗一颗地排开，排成了一行字，悬在半空里。

Y(1,3,10)。

他手指头一动，那一行字往后头长。一颗冰碴子接一颗冰碴子地凝出来，排上去，在夜空里慢慢铺开。

Y(1,3,9,27,81,243,729,2187,6561,19683,59049\ldots\ldots)

一串冰蓝的数，挂在半空里，一个比一个大，一直往后头排，排到看不清了，还在排。每一个数都幽幽地发着冷光。

碎紫斗罗看着这一串数，脸色越来越白。

``九后头是二十七，二十七后头是八十一。''真霜斗罗说，``一直这么乘下去，乘一辈子，乘到头，才挨着那个十。''

他的声音又低又慢，像在给碎紫讲，又像在跟自己说。

``打 Y(1,3,7) 起，往上挪一步，就是一道沟。老夫在 Y(1,3,6) 上站了几十年，往 Y(1,3,7) 那边看，看得清清楚楚，一步都迈不过去。''

``这些年，倒在这条路上的人，老夫见得多了。''

他的手慢慢攥紧了。那一串冰数在他指头缝里碎了，碎成冰粉，从指缝里簌簌地往下掉。

``天分高得吓人的，宗门把家底全往身上堆的，老夫都见过。到头来，照样走不到。有人花上几十年，上百年，就为了让后头那个数往上挪一格。''

碎紫斗罗听着，一句没插。

他知道真霜没往大里说。到了他们这一层，修炼这件事，早就变成了一步一步地往自己那条路的前头挪。Y(1,3,7) 往后，每一步都在掂量一个人的天分，悟性，心气，还有命长不长。

真霜斗罗往远处看了一眼。

那几个孩子正在废墟里互相搀着往起站。蔡须坤坐在一截断柱子上，脸白得发青，眼睛却还亮着。唐三蹲在她边上，手心里那点蓝银草细细的，一根一根地往外冒，冒出来就去缠几个人身上的伤口。霍雨浩盘着腿坐着，精神力一点一点地往外铺，鼻子底下的血总算止住了。奥斯卡蹲在地上，把袋子里剩下的香肠一根一根地往外掏，摆了一地，嘴里念念有词地数。诺当闭着眼，靠在墙根底下，攒魂力。

几个人都狼狈得不像样子，可都还站在这片地上。

他们身上的魂力还在动。刚才那一仗打下来，有些东西正在他们身子里一点一点地沉下去，化开，长进去。

``这几个娃娃\ldots\ldots{}''真霜斗罗的嘴角往上动了一下，``往后这块地上，能指望的，兴许就是他们。''

他停了一下。

``长得太快了。快得老夫都有点吃不消。尤其是蔡家那个丫头。刚碰上胡峰那会儿，让人家压着打。今天，借着大伙儿的劲儿，能跟他面对面地顶。''

他摇了摇头。

``老夫年轻的时候，没这么快。''

``你把我们看得太高了。''

身后有人说话，声音平平的。

几个人一齐回过头。

phyrion 不知道什么时候已经站在了冰熊王边上。他身上那件深蓝的长袍刮破了几处，衣角上全是灰跟石头末子。卢关跟在他后头，怀里抱着卢本韦。卢本韦睡着了，脑袋歪在她爸胳膊弯里，嘴角上还挂着一点排骨汤的印子，手里攥着那只空碗不撒手。小白猪跟在卢关脚后跟，一溜小跑，跑两步，哼一声。

phyrion 的脸还是那么清瘦，那么静，看不出累。可碎紫斗罗留了个心眼。他身上那股看不见底的东西，比刚才撞阵之前，弱了那么一丝。Y(1,3,10) 的人，拱那一下，也不是白拱的。

可那双眼睛一点没变。粉红的，润，看不出一点波澜。

真霜斗罗看着他，笑了一下。

``回来了。''

phyrion 点了点头。他的目光从冰熊王那条伤腿上扫过去，又从碎紫斗罗那张白脸上扫过去，最后落在真霜斗罗脸上。

``熊腿伤得不重，养半个月就好。''他说，``碎紫将军魂力是空了，根子没动，三天能缓回来七成。''

他停了一下。

``倒是你。城楼上那句破不了，说得挺干脆。''

真霜斗罗苦笑了一声。``破不了就是破不了。老夫活这么大岁数，不图那张脸。''

phyrion 的嘴角往上动了动，没再往下说。

``里头怎么样？''真霜斗罗问。

``跑干净了。''phyrion 说，``胡帅叔侄俩，从北边一个水门钻出去的。水门底下的铁栅栏让人锯断了，锯口上的锈都起了，有些日子了。''

``早就留好了退路。''真霜斗罗说。

``嗯。''

几个人都没说话。

过了一会儿，真霜斗罗又开口了。

``phyrion，老夫跟你说句实在话。你这一身本事，老夫是服的。''

phyrion 听完，摇了摇头。

``没什么好服的。''

真霜斗罗的眉毛挑了一下。``没什么好服的？''

phyrion 抬起头，往天上看。刚才那些飘着的光点子已经灭干净了，天又露了出来，星星一颗一颗地挂着，冷冷地往下看。

``你刚才说 Y(1,3,10)。''他的声音轻下来了，``你知道我从 Y(1,3,9) 走到 Y(1,3,10)，走了多少年？''

没人答。

``四十年。''

四下里一下子静了。风停了，冰熊王不哼了，碎紫斗罗连气都忘了喘。

真霜斗罗的瞳孔缩了一下。``四十年？''

phyrion 点头。``整整四十年。我一百二十八岁那年，才算跨过那道坎。''

真霜斗罗没吭声。

四十年。从九到十，就那一个数。

他想起了他师兄。

他师兄年轻时候，是一百年里头都难碰见一个的那种人。他从 Y(1,3,7) 往 Y(1,3,8) 走，走了一百二十年。走到了，转过头又去冲 Y(1,3,9)，死在了半道上。

phyrion 从九走到十，四十年。

这里头是什么分量，真霜斗罗比谁都清楚。到了 Y(1,3,9) 往上，一步一步都是拿十年当一个数来算的。有的人把半辈子耗进去，也就是在 Y(1,3,9) 那道门槛跟前来回打转。phyrion 拿四十年，走完了许多人几辈子都走不完的路。

可这四十年落在 phyrion 自己身上，是什么滋味，就是另一回事了。

phyrion 低下头，看了看自己的两只手。

两只手早就不年轻了。指节粗大，皮上全是细细的褶子，指肚上结着厚厚的老茧。四十年，一层一层磨出来的。

``所以啊。''他轻轻笑了一声，笑里带着点自嘲，``我倒觉得，我是那个没多少指望的人了。''

真霜斗罗皱起了眉头。``别这么说。''

phyrion 摇了摇头。

``一百三十一了。''他说，``到这个岁数，自己这副身子还能走多远，心里有数。Y(1,3,10) 往上，还有路。那条路，我怕是没多少工夫走了。''

他抬起头，往废墟那边看。

蔡须坤还坐在那截断柱子上。卢本韦让她爸放了下来，靠在蔡须坤肩膀上接着睡，蔡须坤一只手扶着她，一只手拿着一卷布条，一圈一圈地给她缠左肩。唐三跟霍雨浩凑在一块儿，低着头小声说话。奥斯卡终于把香肠数完了，脸一下子垮下来，嘴里嘟囔：``不够了不够了，回去得再熬\ldots\ldots{}''诺当靠在墙根，闭着眼，嘴角挂着一点笑。

几个人累得都快站不住了，可眼睛里都还有光。

phyrion 看着他们，眼神一点一点地软下来。

``我这把老骨头，没什么好指望的。''他说，``指望在他们身上。''

真霜斗罗顺着他的目光看过去。

``他们年轻。''phyrion 慢慢地说，``往后还有几十年、上百年去试。错个十回八回都不要紧。在 Y(1,3,7) 跟前蹲上几十年，再从头来，也不要紧。他们有工夫。''

他的声音越说越轻，轻得像夜风里最后一片叶子。

``我没有了。''

真霜斗罗半天没出声。

冰熊王在他脚底下低低地哼了一声，像是在应这句话。碎紫斗罗低着头，魂力空了，连叹口气的劲儿都没有。

过了好一会儿，真霜斗罗笑了。

``所以你才老盯着他们看。''

phyrion 也不推。

``嗯。''他说，``往后站在前头的，该是他们。我们这几个老东西，站不了几年了。''

他的目光落到了蔡须坤身上。那个丫头一身的伤，还坐得笔直，低着头给人缠布条，缠得一丝不苟。

夜风从废墟上刮过去，卷起一阵灰。

真霜斗罗看着眼前这个比自己小不了几岁的人。今晚之前，他一直把 phyrion 当成一个本事比自己大的同辈，一个往后少不了要并肩打仗的人。这时候他头一回换了个眼光看他。

一个走了很长的路、走到这会儿开始回头看后来人的老头。养猪的。拿四十年走完了别人几辈子走不完的路，还嫌自己没工夫。

``phyrion。''真霜斗罗说。

phyrion 转过头。

真霜斗罗看着他，冰蓝的眼睛里有一点火，那点火他年轻的时候常有，后来很多年没见过了。

``四十年从九走到十，''他说，``不等于你这条路就到头了。''

phyrion 笑了一下。

``兴许吧。''他说，``不管怎么样，趁我还走得动，替他们多挡几回风。''

\begin{center}\rule{0.5\linewidth}{0.5pt}\end{center}

他说完，转过身，朝废墟那边走过去。

八个人看见他过来，一个个想往起站。phyrion 摆了摆手，让他们坐着。

``大长老 test\_alpha0 殿下，''他说，``托我捎几句话给你们。''

蔡须坤一愣。

``大长老？''

``你们出城那天，他就往养猪场递了信。''phyrion 说，``信上说，天斗城要出事，让我们备着。我们在养猪场等了三天。今天夜里，城里那一下光一冲上天，我跟小关就出门了。''

八个人互相看了一眼。

原来是这样。

卢关能赶在大阵合拢之前撞进来，大长老在他们出门那天就把这一步算好了。

``他让我跟你们说，''phyrion 接着说，声音一下子正经起来，``蓬莱岛的事，不急。你们还得再多练练，再去。''

``你们今晚这一仗，是打下来了，可根子还没扎稳。魂力往上蹿得太快，身子跟不上。这会儿就去蓬莱岛那种地方练，练不出好来，只会把自己练坏。''

蔡须坤沉默了一会儿，点了点头。

她知道这是实话。今天夜里这一仗，他们是拿命垫着才撑下来的。蓬莱岛那边要是比今晚还狠，就凭他们现在这副样子，去了也是白去。

``先养伤。''phyrion 说，口气里没有商量的余地，``把根子扎牢了，再去。''

``是。''蔡须坤说。

phyrion 没马上走。他看了看靠在蔡须坤肩上睡着的卢本韦，又看了看蔡须坤。

``他还捎了一句。''phyrion 说，``他说，魂力跑在根子前头，这个滋味，他尝过。''

``大长老也有过？''

``天州。''phyrion 说。

蔡须坤没听懂。

``我年轻的时候往西边送过一趟猪，路过天州。''phyrion 说，``城外有个大坑，方圆两三里，坑壁烧成了黑的，亮得能照见人影。坑底立着一座石像，比城楼还高，两只手拄着剑，脸朝着星斗大森林。''

``是大长老？''

``是他。坑沿上修了一座纪念堂，四面墙上刻满了名字。''phyrion 说，``殿里的史册上，那一仗只记了三行。那个坑，几十丈深。''

他没再往下说，蔡须坤也没问。卢本韦在她肩膀上动了动，嘴里咕哝了一句什么，又睡沉了。

\begin{center}\rule{0.5\linewidth}{0.5pt}\end{center}

``咳\ldots\ldots 咳咳\ldots\ldots{}''

城楼那边传过来一阵咳嗽，咳得撕心裂肺。

所有人一齐扭过头去。

大帝还靠在城楼的垛口上，刚才谁都没顾上他。这时候他身子在晃，脸灰得跟城砖一样。刚才他扶着垛口看了那么久的阵，一直硬撑着，撑到这会儿，撑不住了。

更不对劲的是他的嘴唇。嘴唇在变色，一点一点地发青，发紫，眼看着往深里走。

``陛下！''小舞就在他边上，一把扶住了他。

大帝两只手撑着垛口，指节都白了。脑门上冒出一层冷汗，喘气又急又乱。他每往外呼一口气，嘴里就带出来一缕黑雾，黑雾在夜风里散开，腥得人想吐。

``毒\ldots\ldots{}''他的嗓子哑得快听不见，``是胡帅\ldots\ldots 那个阵\ldots\ldots{}''

phyrion 的脸色一下子变了。

他人在原地一闪，化成一道粉色的光，再出现的时候，已经站在了城楼上，站在大帝身边。快得真霜斗罗都没看清他是怎么上去的。

``别动。''他说。一只手按在了大帝的胸口上。

粉色的魂力从他手心里淌出来，温，厚，像开春以后化开的一汪暖水，慢慢往大帝身子里渗。那汪暖水在大帝的经脉里走，走到哪儿，哪儿那些往外蔓延的黑东西就往回缩，让它一点一点地压住，堵死，往外赶。

大帝的身子在他手底下抖了几下。嘴唇上那层青紫一点一点褪下去了。他张嘴吐出一口黑气，黑气在半空里变成一缕一缕的黑烟，腥臭腥臭的，让夜风一吹，散了。

``蚀骨散魂毒。''phyrion 说，声音冷下来了，粉红的眼睛里闪过一道寒光，``胡帅埋在阵里的后手。阵一起，这毒就顺着阵里的魂力走，专找阵里身子最虚、底子早让人啃过的那一个钻。慢慢啃人的经脉，啃魂力的根子。不管它，三天以内，一身魂力就废了。''

大帝让他的魂力压着，脸上一点一点有了血色。他吃力地抬起头，看着眼前这个穿深蓝长袍的瘦老头，嘴唇哆嗦着。

``多谢\ldots\ldots phyrion 前辈\ldots\ldots{}''

phyrion 没吭声。他手上的魂力还在往里灌，把那些最顽的毒一点一点地往经脉的梢上逼，逼到头，拿一层薄薄的粉色魂力裹住，封死。

``先封住了。''他把手收回来，``只是封住。要拔干净，少说七天。这七天里，陛下一点魂力都不能动。''

大帝点了点头，让小舞扶着，慢慢坐了下来。他闭着眼喘了一会儿，再睁开的时候，脸上总算又有了点当皇帝的样子。

他的眼睛里全是血丝，可亮得吓人。他看着城楼底下那几个孩子，那几张满是伤、还硬撑着的脸，嘴唇抖了一下。

底下几个人互相搀着，也上了城楼。

大帝看着他们，一个一个地看过去。

``果糕小队。''他开口了，嗓子哑，一个字一个字说得很重，``今天夜里，要是没有你们，天斗帝国\ldots\ldots{}''

他没说完。谁都知道他后头那半句是什么。

``朕欠你们一条命。''

蔡须坤刚要开口推，大帝抬起手，把她拦住了。

``别推。''他说，``朕这辈子，不欠人家的。''

他吸了一口气，目光从八个人脸上挨个扫过去，最后停在蔡须坤脸上。

``天斗皇室传了一千多年。''他说，``历代的皇帝，修炼之外，都会把自己练出来的门道记下来，一辈一辈攒着，攒成了一部书。''

``这部书里记的东西，外头哪儿都找不着。''

他的声音一下子郑重起来。

``它记的，是投影序数怎么练。''

八个人一齐愣住了。

``投影序数？''霍雨浩的瞳孔一缩，精神力不由自主地往外一冲。

``往简单里说，''大帝慢慢地说，``人练到一个坎上，过不去了，正面是一堵墙。投影序数，是教你从墙的侧面绕，绕出一条别的路来。''

真霜斗罗站在城楼底下，听见这句话，瞳孔猛地一缩。

``从侧面绕\ldots\ldots{}''他低声念叨，冰蓝的眼睛里亮了一下，``怪不得\ldots\ldots 怪不得打 Y(1,3,7) 起，一步一道沟。寻常练 Y 序列，走的都是正面，拿脑袋往墙上撞。投影序数，走的是另一条道。''

phyrion 站在大帝身边，粉红的眼睛里也动了一下。他没出声，就那么听着。

大帝点了点头，接着往下说。

``可投影序数最要紧的地方，还不在过坎上。''

他的目光又从八个人脸上扫过去，声音更沉了。

``它真正要紧的，是能破开 Y(1,3)。''

城楼上一下子没声了。

Y(1,3)。

多少魂师练了一辈子，就停在 Y(1,2,4,8,16) 那一串上，Y(1,3) 那道门槛，连摸都摸不着。天分够的，下的苦功也够的，到了那儿，照样过不去。Y(1,2,4,8,16) 往后那一串，加多长都还在门外头。要跨过去，得把自己从根上换一遍。

``投影序数，是破 Y(1,3) 的钥匙。''大帝说，``也是练阴阳互补双魂核的法子。''

阴阳互补双魂核。

这几个字落进八个人耳朵里，像打了一个雷。

``双魂核？''唐三手心里的蓝银草一下子蹿出来一大截，``一个人\ldots\ldots 能有两个魂核？''

大帝摇了摇头。``阴阳互补双魂核，是把自己那一串 Y 序列从当中劈开，一半阴，一半阳，两半对着顶，又谁也离不了谁，在身子里头转成一个圈，自己转自己的。''

``这部书，是朕能拿出来的，最金贵的东西。''

八个人你看看我，我看看你，谁都说不出话来。

蔡须坤吸了一口气。她单膝跪了下去，右手按在胸口，行了一个最郑重的礼。

``果糕小队，谢陛下赐书。''

唐三跟着跪下了。霍雨浩，奥斯卡，诺当，蔡兴坤，小舞，一个接一个跪了下去。卢本韦还靠在蔡须坤肩膀上睡着，蔡须坤一跪，她身子一歪，醒了，迷迷糊糊地看了一圈，也跟着跪下了，手里那只空碗还攥着。

大帝看着他们，嘴角往上动了动。

``起来吧。''他的声音软了一些，``书不急。等你们伤养好了，朕让人送过去。''

\begin{center}\rule{0.5\linewidth}{0.5pt}\end{center}

就在这时候，phyrion 的身子一僵。

他那双粉红的眼睛一下子缩紧了，那张一直都平平静静的脸上，头一回浮起了一层沉。他身上的魂力猛地往外一放，粉色的光从他身子里涌出去，像一张看不见的网，一下子罩住了整座皇城。

他在找。

Y(1,3,10) 的魂力像涨潮一样，往皇城的每一个角落里漫。每一座殿，每一条廊，每一块砖底下，他一处都没漏。

刚才他跟卢关在皇城里转了一圈，转的是地面上。这一回，他往地底下摸。

摸到了。

皇城东南角，有一座早就废了的偏殿。偏殿底下，藏着一间地下室。地下室的口子上，压着一道封印，封印上刻的纹路，跟胡帅身上那股黑帅气是一个窝里出来的。

封印已经碎了。

phyrion 的感知从碎掉的封印里钻了进去，往地下室里一探。

他的脸色一下子全变了。

``\ldots\ldots 大皇子。''他说。声音一下子冷了，沉了。

``还有二皇子。''

所有人一齐看向他。

``东南角，地底下有一间密室。''phyrion 说，声音压得很低，粉红的眼睛里烧着一团冷火，``大皇子\ldots\ldots 已经没了。''

``二皇子还剩一口气。''

城楼上像是让人一下子抽干了气。

大帝的脸，一下子白透了。

``什么？！''

他的声音又尖又抖，整个人从地上猛地往起一挣，魂力不管不顾地往外冲，冲了一半，就是一阵撕心裂肺的咳嗽。一口黑雾从他嘴里喷出来，刚让 phyrion 封住的毒，又在他身子里翻起来了。

``陛下！''小舞跟诺当一左一右，死死扶住了他。

``朕的\ldots\ldots 皇儿\ldots\ldots{}''大帝的声音一下子碎了。刚才那点当皇帝的样子，一点都没剩下，就剩一个当爹的，两只手死死攥着小舞的胳膊，攥得她骨头都疼。

phyrion 一个字都没再说。

他身子一闪，化成一道粉光，没了。

下一口气的工夫，他已经站在了皇城东南角那座废偏殿里。

偏殿的房梁塌了一半，地上全是落下来的瓦片跟鸟粪。正当中的地砖让人撬开了，底下露出一个黑洞洞的口子，口子边上散着那道封印的碎片，一片一片的，还在冒着一点黑气。

phyrion 一脚踩上去，把碎片踩成了粉，身子一矮，钻进了口子里。

底下又潮又暗，血腥味直往鼻子里冲。

他手心里亮起一团粉光，把整间密室照亮了。

密室不大，四面是石墙，墙上长着一层青苔。墙角堆着几只空水罐，一只倒了，罐口朝外，里头干的。地上有两道拖出来的印子，从门口一直拖到正当中，印子边上是干了的血，黑红黑红的。

正当中，两个人倒在血里。

一个仰面躺着，眼睛还睁着，脸扭着。身上是一件明黄的袍子，袍子上绣着天斗皇室的金龙，让血泡透了，成了暗红色。手里攥着半截断剑，剑身上沾着一层黑帅气，跟胡帅身上那股是一个味儿。

phyrion 认得这张脸。这是大皇子，也就是太子。

真亲王那天夜里在寝宫看见的，就是他。phyrion 蹲下来看了一眼，太子的眼角、鼻孔、耳朵眼、嘴角，七处都挂着干了的血道子，黑的。地上那两道拖出来的印子，一道就是他的。寝宫里那具尸首，让人从龙榻边上拖走，挪到了这儿。挪他的人，转过身就穿上了他的衣裳，学着他的样子，站到了宫门口。

太子胸口上还有一个窟窿，碗口大，边上焦黑。窟窿里一滴血都没有。人死了以后才开的。

他手里那半截剑，是谁塞进去的，phyrion 不知道。

太子边上不到三步远，墙根底下，蜷着另一个人。

二皇子。

浑身是血，蜷成一团，脸灰得发青，嘴唇紫得发黑，跟刚才大帝毒发的时候一个样。身上大大小小全是口子，最深的一道从左肩膀一直拉到右腰上，差一点就把人劈成了两半。

可他还在喘。

很轻，很浅，隔好一会儿才喘一下。身子里那点魂力还在转，转得跟一盏快熬干了的油灯一样，风一吹就要灭。

城里传了一夜的话，说二皇子跟着雪星亲王跑了。他哪儿都没跑。他一直在这儿，在他哥边上，在这间地底下的屋子里，一个人熬着。

phyrion 站在密室当中，那团粉光照着这两个人。

他的脸，一点一点地冷下来。

他这辈子很少动气。这时候他也没动气。动气是热的。他身上那点东西是冷的，冷得一点热乎气都不剩，什么都冻住了，最后就剩下一个念头，又干又硬，像一根冻在冰里的钉子。

``\ldots\ldots 胡帅。''

他说得很轻，轻得几乎听不见。

整间地下室，跟着抖了一下。墙上的青苔簌簌地往下掉。

\begin{center}\rule{0.5\linewidth}{0.5pt}\end{center}

他弯下腰，把二皇子从墙根底下抱了起来。抱得很稳，一点没颠着。二皇子的脑袋靠在他胳膊上，嘴唇动了一下，眼睛没睁开。

``别说话。''phyrion 说，``爷爷带你出去。''

二皇子不是他的孙子。他也没改口。

他抱着人，又看了一眼地上的太子。

他把身上那件深蓝的长袍脱下来，盖在了太子脸上。

\begin{center}\rule{0.5\linewidth}{0.5pt}\end{center}

城楼上，大帝让小舞跟诺当扶着，一直望着东南角。

他看见一道粉光从东南角那边升起来，往这边飘。飘得很慢，跟来的时候那一道两样。

光落在城楼底下，phyrion 从光里走出来，上身就剩一件白布褂子。怀里抱着一个人，浑身是血。

大帝往前扑了一步，让小舞死死拽住了。

``活着。''phyrion 抬起头，冲城楼上说，``二皇子还活着。''

他停了一下。

``太子\ldots\ldots 我拿袍子盖上了。待会儿让人去抬。''

大帝站在垛口后头，两只手抓着垛口，抓得指甲都劈了。他没出声，嘴张着，半天，一个字都没挤出来。

碎紫斗罗撑着墙站了起来，一步一步挪到 phyrion 跟前，低头看着他怀里那个人。看了好一会儿，他伸出手，在二皇子额头上碰了一下，又缩了回来。

``二殿下小时候，''他哑着嗓子说，``是本座教他骑的马。''

没人接话。

蔡须坤从城楼上下来，蹲在 phyrion 跟前，两只手悬在二皇子胸口上头。瑞兽本源早让那颗炮弹抽空了，她咬着牙往外挤，挤出来一缕，细得跟头发丝一样，往二皇子身子里送。

``够了。''phyrion 说，``你自己都站不住了。''

蔡须坤没听，又挤了一缕。phyrion 看了她一眼，没再拦。

天快亮了。

东边天上泛起了一层灰白。皇城里静得出奇。烟散了，光点子灭干净了，满天的星星一颗一颗地往下沉，沉到天边，就看不见了。

城楼底下的废墟还在，东边那道撞塌了的城墙还豁着一个大口子。宫门前那个大坑里，冰熊王蜷着身子，睡着了，伤腿还在一下一下地抽。

一夜过去了。

谁都知道，天亮了以后，好些事跟昨天不一样了。

\begin{center}\rule{0.5\linewidth}{0.5pt}\end{center}

\chapter{第七十六章：守护所爱之人}\label{ux7b2cux4e03ux5341ux516dux7ae0ux5b88ux62a4ux6240ux7231ux4e4bux4eba}

那一整天，八个人几乎都在睡。

天亮以后，宫里的人一拨一拨地出来收拾。抬人的，扫瓦的，往大坑里填土的，乱糟糟的一直没停。八个人让养猪场的人领着，在宫门前那片空地上搭的几顶帐篷里倒头就睡，睡到太阳偏了西，才一个一个地醒过来。

醒过来，天又黑了。

养猪场的人在帐篷外头拿木板跟油布围了一圈，挡风，也把外头的动静挡在了外头。围子当中生了一盆炭火，火不旺，红红的，安安静静地烧。火光照在每个人脸上，脸上的伤，脸上没擦干净的血，都让它照得暖了一层。

卢关派人送了吃的跟药过来。排骨汤，烤面饼，几碟咸菜，还有一大包拿油纸裹得严严实实的药膏。送东西的那个小伙子把药膏往地上一搁，说：``老爷子特意交代的，给那个肩膀上挨了一下的丫头多抹点。''

卢本韦脸一红，没吭声。

奥斯卡管分饭。

他把最大的那一碗排骨汤端给了卢本韦，又把自己那张烤饼掰下来一半，塞进她手里。

``你多吃点。''他嘟囔着，``瘦成这样，像什么话。''

``我爷爷也这么说。''卢本韦说。

``那说明你爷爷有眼光。''

他转过身，拿起那包药膏，挨个往别人手里分。分到蔡兴坤那儿，蔡兴坤伸手来接，他往前一递，两个人的手指头尖碰了一下。

两个人一齐往回缩了缩。

蔡兴坤低下头，把药膏攥在手里，没说话。黑子鸡蹲在她脚边上，歪着脑袋看了奥斯卡一眼，咯了一声。

奥斯卡挠了挠后脑勺，眼睛不知道往哪儿放，最后还是飘到了卢本韦那边。

卢本韦两只手捧着汤碗，一小口一小口地喝。左肩让她爸拿魂力封住了，外头又缠了一层布，脸还是白的。火光照在她脸上，把那些灰跟血都照软了，露出来一张还带着点孩子气的脸，腮帮子让汤烫得一鼓一鼓的。

奥斯卡看着她，心口那儿软了一下。

蔡兴坤他是有点意思的。这一点他自己清楚。她长得好看，做事利索，说话不多，可说一句是一句。队里分吃的，他不止一回给她多留一份。打起来，他也不止一回下意识地往她前头挡。

可他心里最想着的，是卢本韦。

头一回见面，那丫头一巴掌拍在他肩膀上，拍得他差点趴下，说以后咱俩一块儿打。他揉着肩膀，心想这是哪来的愣头青。

后来每回要出门打仗，她都偷偷把大家的东西翻一遍，谁的魂导器松了，谁的绑腿没扎紧，她比谁都上心，嘴上从来不说。

再后来，好几回最要命的时候，他回头一看，她都在。一回都没扔下他自己跑过。

蔡兴坤让他心里跳一下，脸上热一下，想往跟前凑，又不太敢。

卢本韦给他的，是另一样东西。是踏实。不管出了什么事，一回头，她在。

他想护着她。跟她爹是谁、她爷爷是谁、她家里有多大一片养猪场，都扯不上关系。她就是卢本韦，他就想护着她。

可打起仗来，他从没因为这点心思偏过她。该挡的，他替所有人挡。该冲的，他冲在前头。他是队里管补给的，金纹的香肠给了两个前辈，赤金的给了母猪，恢复的那一把，见人就塞。他那颗篮球，也是给全队开道用的。

大伙儿的事在前头，自己的心思在后头。这是他给自己定的规矩，一回都没破过。

``看什么看。''卢本韦从汤碗上抬起头，``汤洒你鞋上了。''

奥斯卡低头一看，手里那碗给诺当的汤歪了，汤顺着碗沿往下淌，淌了他一鞋面。

``哎哟。''

诺当伸手把碗接过去，看了他一眼，又看了卢本韦一眼，什么都没说，低头喝汤。

围子的门帘一掀，卢关弯着腰钻了进来。

他把西装脱了，搭在胳膊上，衬衫袖子挽到胳膊肘，领带也不知道扔哪儿去了。他往火盆边上一蹲，先伸手在卢本韦脑门上摸了摸，又捏起她的左胳膊，轻轻转了转。

``还疼不疼？''

``不疼。''

``嘴硬。''卢关说。他从兜里掏出一个小瓷瓶，倒出两粒黑药丸，塞进她嘴里，``嚼了，别吐。''

卢本韦嚼了两下，脸皱成一团，硬咽了下去。

卢关又扭头去看蔡须坤。

``丫头，你那颗炮弹，''他说，``大长老信里提过一句。说要是你们用上了，用的人得虚好几天，让我们带点补气的东西。''他把另一个瓷瓶搁在她脚边，``一天一粒，别多吃。''

``谢谢卢叔叔。''

``谢什么。''卢关站起来，拍了拍膝盖上的灰，``老爷子在西边偏殿守着二皇子，一天没出来了。我给他送饭去。你们早点睡。''

他掀开门帘，又回过头，看了卢本韦一眼，什么都没说，出去了。

火堆另一边，霍雨浩坐着，一直没怎么动。

他的精神力从醒过来就铺出去了，铺得很薄，像一张软软的网，把这一圈帐篷、这一盆火、火边上这几个人全兜在了里头。

所有人都松下来了，他还在看着。这是他的习惯，也是他给自己揽的差事。队里得有一双眼睛一直睁着，他愿意当这双眼睛。

小舞坐在他边上。

她靠着他的肩膀，一头银白的长头发散下来，搭在他胳膊弯里，像一只团着的猫。昨天夜里她一直守在担架边上，碎砖头一阵一阵地往她背上砸，她拿后背替大帝挡着，挡了好几回。这会儿她后背上青了一大片，衣裳底下缠着布，一动就疼。

她一声都没喊过。

霍雨浩侧过脸看了她一眼。

火光照在她半边脸上。那双粉色的大眼睛半眯着，睫毛长长的，在脸上投下一小片影子。嘴唇抿着，也不知道是在忍疼，还是什么都没想，就那么靠着他，暖着。

霍雨浩心里软得一塌糊涂。

他想起了昨天夜里，巷子口，胡峰刚把那座黑山砸下来以后。

那时候他脑袋里那口缸已经快见底了。撞那堵黑墙撞了一回，顶泰山压顶又顶了一回，再往外掏，掏出来的就是底子。底子掏空了，人要傻好些日子，往后精神力能不能长回来，都两说。

他还是把冰帝跟雪帝放出来了。

他记得那一下自己根本没想。脑子还没转过来，手已经伸出去了，身子替他先应了。

他那时候就知道一件事。胡峰往前一步，巷子里头就是担架，担架边上就是小舞。他身后还站着队长，站着唐三，站着一整个队的人。

两个魂灵出来那一下，他脑袋里的识海晃了三晃，鼻血哗地就下来了。

他一点都不后悔。

小舞是他心尖上的人。这个队，是他的家。为了家里这些人，他把底子掏出来，掏多少都认。

``雨浩\ldots\ldots{}''小舞的声音很轻，带着累透了以后那股懒劲儿，``你脑袋还疼不疼？''

``不疼了。''霍雨浩笑了一下。

``骗人。''

小舞伸出一只手，轻轻按在他的太阳穴上。她的手指头凉凉的，按得很轻。

``你识海还在晃。''她说，``我摸得出来。''

霍雨浩没说话。他把她的手从自己太阳穴上拿下来，攥在手心里。

``睡吧。''他说，声音压得又低又软，``我守着。''

小舞看着他。那双粉色的眼睛里映着火，也映着他。她没再说什么，往他肩膀上又靠紧了一点，眼睛慢慢合上了。

\begin{center}\rule{0.5\linewidth}{0.5pt}\end{center}

蔡须坤坐在火堆最外头。

她面前那碗排骨汤，一口没动。烤饼搁在碗沿上，整整齐齐的，也没动。她两只眼睛看着火，看得很远，像是在看火，又像是隔着火，在看一个很远很远的地方。

炮弹炸完以后，她身子一直是空的。手脚又沉又软，坐着都得拿手撑着膝盖。大长老说的那几天，这才是头一天。

她的目光从火上挪开，往边上扫。

奥斯卡蹲在卢本韦边上，笨手笨脚地给她弄肩膀上那圈布条。弄得很轻，很小心，像是生怕把她弄疼了。卢本韦一边喝汤一边含含糊糊地说：``行了行了，你别弄了，越弄越紧。''嘴上嫌他，嘴角是翘着的。

蔡须坤的嘴角跟着动了一下。

她又往另一边看。

霍雨浩攥着小舞的手，低着头看她。小舞已经睡着了，靠在他肩膀上，喘气喘得匀匀的。霍雨浩还睁着眼，精神力还铺着，替所有人看着。

蔡须坤看着这些，一样一样地看过去。

他们各自都有了想护着的人。

她低下头，看着自己的两只手。

两只手上全是伤。旧的叠着新的，刀口压着烧痕。最旧的那一道，在右手手背上，一道白印子，弯弯的，已经平了。

这只手，牵过牢大的手。

两年前那个开春，星斗大森林外头那条路上，十一个人围上来。牢大让大家带着她往后退。退的时候，她回头看了一眼，看见牢大站在路当中，冲她这边摆了摆手。

后来就没有后来了。

牢大。她的爱人。

那个说往后打架他站最前头的少年。她拿什么主意，他从来不吭声，就站过来，站到她边上去。最后那一回，他一个人站在那条路上，把自己一层一层地烧，烧了一整个时辰，帅气教追上来的那十一个人，一个都没让过去。

他什么话都没留下。连一块骨头都没留下。

蔡须坤闭上了眼睛。

她是队长。她不能让这点难受把自己吃了。队里的人得靠她，往后的路得靠她走。

可一到夜里人都睡了，她一闭眼，就是牢大那张脸。

那张总带着点笑的脸。那双一直往她这边看的眼睛。

``队长。''

有人叫她。

她睁开眼。唐三站在她跟前，手里端着一碗排骨汤，往她跟前递，递得有点笨。

``你一口都没吃。''他说，声音压得很低，``不管怎么样，得吃点。''

蔡须坤看着那碗汤，没动。看了一会儿，她伸出手接了过来。

``谢谢。''

她低头喝了一口。汤是温的，不烫了，香味从嘴里一直漫到鼻子里，顺着嗓子往下走，走到肚子里，空落落的那块地方，填进去了一点。

唐三在她边上坐下了。

``队长。''他的声音压得更低了，``我知道你在想什么。''

蔡须坤没吭声。

``牢大\ldots\ldots{}''唐三说到这两个字，嗓子发涩，顿了一下，``他要是看见你这样，不会高兴的。''

``我知道。''

蔡须坤的声音很平，平得有点冷。可唐三看见了，她攥着碗的那几根手指头，白了。

``我知道。''她又说了一遍。两只眼睛里映着火，也蒙着一层薄薄的水，``可我知道，跟我不难受，是两回事。''

唐三不说话了。

他不知道该说什么。宽心的话说出来太轻，他也说不出口。他能做的，就是在这儿坐着，陪着她。

火盆里的炭噼啪响了一声。

蔡须坤把那碗汤一口一口喝完了，空碗搁在一边。她吸了一口气，抬起头，眼睛里那层水还在，可目光已经硬起来了。

``明天一早，进宫。''她说，声音又是队长的声音了，``陛下说的那部书，领了。往后怎么走，领完再说。''

``是。''唐三点了点头，站了起来。

他走了两步，又停住了，回过头来看她。

``队长。''

``嗯？''

``牢大\ldots\ldots{}''他说，``他要是看见了今天夜里，肯定挺得意。''

蔡须坤没答。

她低下头。眼睛里那层薄薄的水，到底聚成了一颗，一点声音都没有地滑下来，落在她右手手背上，落在那道最旧的白印子上。

\begin{center}\rule{0.5\linewidth}{0.5pt}\end{center}

诺当跟蔡兴坤坐在帐篷口那一边，背靠着背。

诺当手里攥着一只小酒葫芦，是养猪场送饭的小伙子偷偷塞给他的。他拔了塞子，闻了闻，又塞上了，没喝。

``不喝？''蔡兴坤问。她没回头，后脑勺抵着他的后脑勺。

``留着。''诺当说，``等胡帅那个老东西让人逮住了，再喝。''

蔡兴坤嗯了一声。

过了一会儿，她说：``你那个环，碎的时候，疼不疼？''

``疼。''诺当说，``像让人当胸捶了一拳。''

``现在呢？''

``还疼。''

蔡兴坤没再问。她把手伸到背后，摸着了诺当的胳膊，在他胳膊上拍了两下，拍完了又把手收回去了。

黑子鸡从她怀里钻出来，蹦到诺当的膝盖上，蹲下了，脑袋往翅膀底下一埋，睡了。

诺当低头看着那只鸡，没赶它。

\begin{center}\rule{0.5\linewidth}{0.5pt}\end{center}

夜深了。

火盆里的火一点一点矮下去，只剩下一层暗红的炭，偶尔噼啪一下。

几个人陆陆续续都睡了。奥斯卡靠在油布围子上打起了呼噜，一声长一声短。卢本韦枕着他那件外衣，睡得正香，嘴角上还沾着一点汤。霍雨浩还睁着眼，小舞枕在他胳膊弯里，喘气喘得又轻又匀。唐三盘着腿坐在角落里，左手手心里冒着一点蓝银草，细细的，绕着他手指头转。他的右手搁在膝盖上，攥着拳。攥了一会儿，他把拳头松开，低头看了一眼手心。手心里什么都没有，只有一层汗。他又把拳头攥上了。诺当跟蔡兴坤还背靠着背，一个歪着头，一个低着头，都闭上了眼。黑子鸡在诺当膝盖上，也睡着了。

蔡须坤还坐在火盆边上。

她看着那一盆快灭的炭。

她想起了 phyrion 昨天夜里在废墟上说的那句话。往后站在前头的，该是他们。

她又低头看了看手背上那道白印子，再抬头看了看四下里睡着的这几张脸。

牢大不在了。

这个队还在。

他们一块儿死里逃生过，一块儿往上长，也一块儿把一个人丢在了那条路上。打到今天，大伙儿还挨在一块儿。这份情分是让血跟火一遍一遍淬出来的。疼也疼过，哭也哭过，最后还是挨着坐在一盆火边上。

蔡须坤撑着膝盖站起来。腿一软，她扶了一下油布围子，站稳了。

她走到卢本韦边上，弯下腰，把滑到一边的那件外衣往上拉了拉，盖住她的肩膀。又走到奥斯卡跟前，替他把歪到一边的领子拽正了。

她一个一个地走过去，走得很慢，动作很轻。给小舞掖了掖衣角，把唐三手边那只快倒的水碗扶正，从诺当膝盖上把那只黑子鸡抱起来，塞回蔡兴坤怀里。

这是她的队。

是牢大拿命护下来的东西。

往后，她接着护。

火盆里最后一块炭啪地响了一声，蹦起来一小撮火星子。

火星子往上飘，飘到半空里，一颗一颗地亮着，转着，慢慢落下来，落到半截就灭了。

远处，皇城的影子在月亮底下黑沉沉的，一声不出。

再往里，宫里西边一座偏殿，窗户还亮着灯。

二皇子躺在榻上，浑身缠满了布，脸还是青的，嘴唇还是紫的，可胸口一起一伏，比昨天夜里匀多了。太医来了三拨，看了一眼，都说不敢下手。最后是 phyrion 拿魂力一点一点地往里渡，渡了一整天，才把那口气稳住。

phyrion 坐在榻边一张小凳子上。身上那件深蓝的长袍已经换了一件新的，是卢关让人从养猪场捎来的，有点大，袖子挽了两道。

小白猪趴在他脚边，粉红的圆眼睛半眯着，隔一会儿哼一声。

他没看榻上的人。他一直看着窗外头，看着北边。

北边是那个水门。铁栅栏底下锯开的那个口子，昨天夜里他亲手摸过。锯口上的锈，少说一个月。

一个月以前，胡帅就把退路留好了。

phyrion 坐着，一动不动。窗外头的夜风刮进来，把灯苗吹得一歪一歪的。

天亮还有好一阵。

\begin{center}\rule{0.5\linewidth}{0.5pt}\end{center}

第二天一早，太阳从皇城东边那道撞塌了的城墙豁口里照进来，照在一片烂瓦上。

八个人刚洗了把脸，啃了两口饼，宫里就来人了。一队穿金甲的禁卫，客客气气地站在油布围子外头，说陛下请几位小友进内廷。

八个人跟着他们进了宫。

一路上，昨天夜里那些扫地的宫女，倒茶的内侍，拄着刀的禁军，都不见了。换了一批新面孔，脸上有血色，眼珠子会转，见了他们，低着头往边上让。

正殿上，静得很。

大帝坐在龙椅上。脸还是白的，带着大病一场还没缓过来的那种白。可 phyrion 一整夜两头跑，抽空又给他压了一回毒，那蚀骨散魂毒让一层粉色的魂力裹得严严实实，一时半会儿翻不起来。他坐得笔直，又是一个皇帝的样子了。

殿下两边站着人。

真霜斗罗站在左边头一个，换了一身干净的旧袍子，背挺得笔直。碎紫斗罗站在他边上，左胳膊吊在胸前，脸还是白的，人站得稳稳当当。右边站着 phyrion，一件挽着袖子的深蓝长袍，粉红的眼睛静静地看着殿当中。

phyrion 边上，站着一个人，八个人看见了，一齐愣了一下。

雪星亲王。

真的那一个。左胳膊吊着，右腿底下拄着一根拐，脸上的伤结了痂。他身后站着碎紫斗罗那两个亲兵，一边一个。他看见八个人进来，冲他们点了点头，点得很重。

``果糕小队。''大帝开口了。声音在殿里回着，不大，稳。

蔡须坤往前迈了一步，领着身后几个人，一齐单膝跪了下去，右手按在胸口。

``臣等，拜见陛下。''

大帝看着底下这几张年轻的脸，看了好一会儿。他慢慢站起身来，目光从他们身上扫过去，最后停在了蔡须坤身上。

``前天夜里，皇城差一点就翻了。''他说，``没有你们几个拿命去顶，天斗这会儿已经是另一副样子了。朕是天斗的主子，受了你们这份护国的恩，不重重地赏，朕往后拿什么去服人？''

他吸了一口气，声音一下子拔高了。

``胡帅从北边水门跑了。''大帝说，``朕已经下了旨，天斗境内，见着他，格杀勿论。太子\ldots\ldots 三天以后出殡。二皇子那边，命保住了。''

他停了一下。

他的目光往殿下两边扫了一圈，在真霜斗罗身上停了停，在碎紫斗罗吊着的胳膊上停了停，最后落在 phyrion 身上。

``三位前辈，朕另有谢。''他说，``真霜，你那头熊的腿，朕让御马监最好的兽医去看。碎紫，你那条胳膊养好之前，禁军的事交给副将，你给朕在家里躺着。''

碎紫斗罗张了张嘴，想说什么，让真霜斗罗拿胳膊肘碰了一下，又闭上了。

``phyrion 前辈。''大帝说，``养猪场往后在天斗境内的买卖，一文税都不收。''

phyrion 点了一下头，没说话。站在殿门口的卢关，脸上那点笑一下子深了。

``这些是朕自家的事。朕要说的，是你们的事。''

``朕已经定了。打今天起，果糕小队八个人，全封帝国子爵。赐地，每年从国库里领俸。''

殿上的人一下子都动了一下。

子爵。在天斗，这是一脚就把他们踩进了帝国最上头那一层里。

奥斯卡跟卢本韦对看了一眼，两个人的眼睛都瞪圆了。蔡须坤的眼睛里也亮了一下，可她还跪着，没急着谢恩。

大帝看出来了，嘴角往上动了一下。

``这只是个头。''他说。

他转过身，从身后的龙案上拿起一本书。

书不厚，封皮旧了，边角磨得起了毛，上头一点花样都没有，就写着四个字，笔画古拙。

《投影序数》。

``朕前天夜里在城楼上许过你们。''大帝说，声音一下子沉了，``这是皇室传下来的书。里头写的是怎么破 Y(1,3)。今天，朕把它交给果糕小队。你们拿回去，好好琢磨。一千多年攒下来的东西，别糟蹋了。''

蔡须坤两只手高高举过头顶。

``果糕小队，谢陛下赐书。''

大帝点点头，走下台阶，把那本书放进了她手里。

书很轻。可蔡须坤接住它的时候，两只胳膊往下一沉。

大帝没回龙椅上去。他就站在台阶底下，离他们几步远，看着他们。

``朕知道，''他说，``前天夜里那一场，你们身上落下的东西，得有工夫慢慢化开。光有工夫还不够，得有地方。''

``所以朕再给你们一样。打今天起，果糕小队，进天斗皇家修炼堂。''

``皇家修炼堂？''

这一回，连唐三都没绷住，脸上的样子一下子变了。

``对。''大帝点头，``皇家修炼堂，是历代天斗皇帝跟皇族宗亲闭关的地方。天斗城里天地灵气最浓的一块地，就在那儿。历代先皇练出来的门道，也都留在那儿的墙上。皇族正根以外的人，就是功劳再大的臣子，能进去的也没几个。''

``你们在里头练六个月。''他说，``六个月以后，你们要是觉得火候到了，phyrion 前辈自会领你们去蓬莱岛。''

phyrion 站在右边，听见自己的名字，没出声，就点了一下头。

八个人这回是真绷不住了。

皇家修炼堂。天斗的魂师，做梦都想进去的地方。外头人都说，在里头练一天，顶得上在外头苦练十天。

奥斯卡的嘴张开了，半天没合上。卢本韦伸手在他胳膊上掐了一把，掐得他嘶了一声，嘴才合上。霍雨浩攥着小舞的手，攥紧了。诺当低着头，嘴角往上翘了一下，又压下去了。

蔡须坤吸了一口气，把心口那股往上顶的东西硬压了下去。

她知道这份赏后头压着什么。前天夜里，这个人没了一个儿子，差一点又没了一个，差一点连整座城都没了。他把能掏的都掏了出来，往他们手里塞。

``果糕小队，领旨。''她说，``一定不让陛下白看我们一回。''

八个人一齐喊了一声，声音在殿里撞来撞去，好半天才散。

大帝点了点头，像是总算放下了一桩事。他转过身，往台阶上走，走到一半，回过头来。

``phyrion。''他说，``劳烦前辈，亲自送他们去皇家修炼堂。''

``嗯。''phyrion 应了一声，转身往殿外走。

八个人又行了一回礼，站起来，跟在他后头。

走到殿门口，碎紫斗罗叫住了蔡须坤。

``丫头。''他说。他的左胳膊还吊着，拿右手在她肩膀上拍了一下，拍得很轻。

``将军\ldots\ldots{}''

``别说谢。''碎紫斗罗摆了摆手，``前天夜里，本座欠你们的。往后慢慢还。''

真霜斗罗站在他边上，没说话，冲她点了一下头。他那双冰镜子一样的眼睛，看她的时候，比从前软了一点。

雪星亲王拄着拐，一瘸一拐地挪过来，在殿门口站住，冲八个人拱了拱手，拱得很低。

``杏林堂那味玄冰玉露，''他说，``本王找着了。二殿下那边，太医说能用。''

``那就好。''蔡须坤说。

\begin{center}\rule{0.5\linewidth}{0.5pt}\end{center}

走出殿门，太阳一下子照在了脸上。

早上的太阳，不烫，亮。八个人一齐眯了一下眼睛。

phyrion 走在前头，深蓝的长袍让风吹得往后飘。小白猪一溜小跑地跟在他脚边，跑两步，回头看一眼，看他们跟没跟上。

卢本韦小跑了几步，追上了 phyrion。

``爷爷，''她说，``红烧猪蹄呢？你说第二天就做的。''

phyrion 头都没回。

``修炼堂的灶上，我打过招呼了。''他说，``中午就有。''

卢本韦一下子乐了，扭过头冲后头喊：``听见没！中午有猪蹄！''

奥斯卡第一个应了一声。霍雨浩跟小舞对看了一眼，都笑了。唐三摇了摇头，嘴角也是往上的。

蔡须坤走在八个人最后头。

她低头看了一眼手里那本《投影序数》，封皮上那四个字让太阳一照，旧得发黄。她又回头看了一眼。宫殿一重一重地压在那儿，屋脊上还落着前天夜里崩上去的碎石头，有一块卡在一只脊兽的嘴里，没人去拿。

从今天起，果糕小队跟从前不一样了。往后要走的路，比前天夜里那一场还难走，也比它宽。

她把书揣进怀里，贴着胸口按了按。

牢大。

她在心里叫了一声。

``你看见了么。''

没人应。风从宫墙底下吹过来，把她的头发吹到了脸上。

往皇家修炼堂去，要从前殿的廊子底下过。

廊子底下有个宫女在扫地。十五六岁，圆脸，鼻尖上几粒雀斑。前天夜里，她在这儿一遍一遍地扫同一块砖，嘴里念叨着卯时扫前殿。

这会儿她扫得好好的，扫一下，往前挪一步，地上的落叶让她归成了一小堆。听见脚步声，她抬起头，看见这一队人，愣了一下，赶紧把扫帚往身后一藏，低下头往边上让。

小舞走过她跟前，冲她笑了一下。

宫女也笑了，脸一下子红了，低着头接着扫。

亭子里那个老内侍也在。他坐在石凳上，手里端着一只茶杯，自己喝。杯子是满的，一滴都没洒。

她把头发拨开，抬起脚，跟上了前头那几个人。她走得不快，一步一步，走得很稳，朝着皇家修炼堂那边去了。

\chapter{第七十七章：2-投影序数的初步介绍}\label{ux7b2cux4e03ux5341ux4e03ux7ae02-ux6295ux5f71ux5e8fux6570ux7684ux521dux6b65ux4ecbux7ecd}

皇家修炼堂的最深处，一片被七层封印笼罩的地下宫殿中，果糕小队八人围坐在一片由魂力晶石铺就的圆形平台之上。

穹顶高达数十米，镶嵌着无数散发着柔和光芒的魂力晶石。空气中弥漫着浓郁到几乎液化的天地灵气，每一次呼吸都能感受到魂力在经脉中微微躁动。

Phyrion站在平台中央。

白色母猪趴在他的脚边，粉红色的圆眼睛半眯着，偶尔发出一声低沉的''哼哼''。

``从今天开始，''Phyrion的声音平静而清晰，在空旷的地下宫殿中回荡，``我要教你们投影序数。不是纸上谈兵的理论，而是真正能帮助你们突破 0 111 22 壁垒的实践之法。''

他抬起手，粉红色的魂力在掌心凝聚，化作了一行古朴的文字，悬浮在半空中------

\textbf{``投影序数讲义''}

``这门学问，''Phyrion说。我作为皇家学会的会长，将这些内容融会贯通之后，今天传授给你们。''

他的目光扫过每一个人。``认真听。这关乎你们能否突破 0 111 22。''

Phyrion的魂力在空中画出了三个巨大的符号------

\NovelDisplay{M00741R01}{%
\Omega \quad I \quad M%
}

``首先，''他说，``我们把传统的非递归序数，比如 \allowbreak{}\NovelInline{I03873}{\Omega}\allowbreak{}、\allowbreak{}\NovelInline{I03874}{I}\allowbreak{}、\allowbreak{}\NovelInline{I03875}{M}\allowbreak{} 等东西称作 \textbf{1-投影序数}。''

奥斯卡挠了挠头：``非递归序数？就是那些不能用更小的序数通过递归运算得到的大序数，也就是基本列长度=自己的序数？''

``没错。''Phyrion点了点头，``\allowbreak{}\NovelInline{I03876}{\Omega}\allowbreak{} 是第一个不可数序数，\allowbreak{}\NovelInline{I03877}{I}\allowbreak{} 是第一个弱不可达序数，\allowbreak{}\NovelInline{I03878}{M}\allowbreak{} 是第一个马洛序数。它们都是 `1-投影序数'------也就是投影体系中的基础砖石。''

他的魂力再次凝聚，在空中出现了一个全新的符号------

\NovelDisplay{M00742R01}{%
a%
}

``然后想象存在一系列巨大的序数，我们称为 \textbf{2-投影序数}，简称 \textbf{2-P}。最小的 2-P 记作 \allowbreak{}\NovelInline{I03879}{a}\allowbreak{}''

``2-投影序数？''唐三皱了皱眉，``比1-投影序数更高一级的存在？''

``对。''Phyrion说，``如果说1-投影序数是'不可被递归运算触达的序数'，那么2-投影序数就是'不可被1-投影运算触达的序数'。它是更高阶的存在。\allowbreak{}\NovelInline{I03880}{a}\allowbreak{} 是最小的2-投影序数。''

他在空中补了一行注释：

\begin{quote}
\textbf{注意：} \allowbreak{}\NovelInline{I03881}{\Omega_{a+1}}\allowbreak{}（即 \allowbreak{}\NovelInline{I03882}{a}\allowbreak{} 的下一个 \allowbreak{}\NovelInline{I03883}{\Omega}\allowbreak{} 序数）、\allowbreak{}\NovelInline{I03884}{I_{a+1}}\allowbreak{}（\allowbreak{}\NovelInline{I03885}{a}\allowbreak{} 的下一个 \allowbreak{}\NovelInline{I03886}{I}\allowbreak{} 序数）也是 1-投影序数。
\end{quote}

``\allowbreak{}\NovelInline{I03887}{\Omega_{a+1}}\allowbreak{} 是 \allowbreak{}\NovelInline{I03888}{a}\allowbreak{} 之后的第一个 \allowbreak{}\NovelInline{I03889}{\Omega}\allowbreak{} 序数，\allowbreak{}\NovelInline{I03890}{I_{a+1}}\allowbreak{} 是 \allowbreak{}\NovelInline{I03891}{a}\allowbreak{} 之后的第一个 \allowbreak{}\NovelInline{I03892}{I}\allowbreak{} 序数。它们相对于 \allowbreak{}\NovelInline{I03893}{a}\allowbreak{} 来说，都是1-投影序数。记住这一点，后面会反复用到。''

Phyrion深吸了一口气，粉红色的魂力在空中凝结成了三行公式，每一行都散发着淡淡的光芒------

``现在，我们定义 \allowbreak{}\NovelInline{I03894}{\psi_a}\allowbreak{} 函数。这是投影序数体系中最核心的工具。''

\NovelDisplay{M00743R01}{%
\psi_a(0) = \Omega%
}

``第一条规则：\allowbreak{}\NovelInline{I03895}{\psi_a}\allowbreak{} 在 0 处的值等于 \allowbreak{}\NovelInline{I03896}{\Omega}\allowbreak{}这是起点。''

\NovelDisplay{M00744R01}{%
\psi_a(X+1) = \psi_a(X) \cdot \omega%
}

``第二条规则：每往后推一步，就把前一个值乘以 \allowbreak{}\NovelInline{I03897}{\omega}\allowbreak{}。这意味着 \allowbreak{}\NovelInline{I03898}{\psi_a(1) = \Omega \cdot \omega}\allowbreak{}，\allowbreak{}\NovelInline{I03899}{\psi_a(2) = \Omega \cdot \omega^2}\allowbreak{}，以此类推。''

\NovelDisplay{M00745R01}{%
\psi_a(X \sim a) = \psi_a(X \sim \psi_a(X \sim \psi_a(X \sim \dots)))%
}

``第三条规则------也是最重要的一条。''Phyrion的声音变得严肃起来，``其中 \allowbreak{}\NovelInline{I03900}{\sim}\allowbreak{} 是 \allowbreak{}\NovelInline{I03901}{+}\allowbreak{} 或 \allowbreak{}\NovelInline{I03902}{\times}\allowbreak{} 或 \allowbreak{}\NovelInline{I03903}{\hat{}}\allowbreak{} 或其他运算。''

``当 \allowbreak{}\NovelInline{I03904}{\psi_a}\allowbreak{} 的自变量涉及到 \allowbreak{}\NovelInline{I03905}{a}\allowbreak{} 本身的时候，就触发了自指递归。比如 \allowbreak{}\NovelInline{I03906}{\psi_a(a)}\allowbreak{}，就等于 \allowbreak{}\NovelInline{I03907}{\psi_a(\psi_a(\psi_a(\dots)))}\allowbreak{}------无限嵌套。\allowbreak{}\NovelInline{I03908}{\psi_a(a^2)=\psi_a(a \times \psi_a(a \times \psi_a(a \times \dots)))}\allowbreak{}。''

``这意味着，\allowbreak{}\NovelInline{I03909}{a}\allowbreak{} 在 \allowbreak{}\NovelInline{I03910}{\psi_a}\allowbreak{} 内部扮演的角色，就像 \allowbreak{}\NovelInline{I03911}{\Omega}\allowbreak{} 在普通的 OCF 中扮演的角色一样------它是折叠的锚点。''

霍雨浩的精神识海微微震动。他能感觉到，这些公式背后蕴含着某种深层的结构------那是一种他从未接触过的、关于序数本身的''投影''逻辑。

Phyrion的魂力在空中画出了一个巨大的条件语句------

``接下来是投影的核心规则。听好了。''

``如果 \allowbreak{}\NovelInline{I03912}{Y}\allowbreak{} 是一个 \textbf{1-投影序数}，则：''

\NovelDisplay{M00746R01}{%
\psi_a(Y) = \sup\{\psi_a(Z) \mid Z < Y\}%
}

``其中 \allowbreak{}\NovelInline{I03913}{Z}\allowbreak{} 是所有小于 \allowbreak{}\NovelInline{I03914}{Y}\allowbreak{} 的序数。''

``换句话说，''Phyrion看着众人的表情，解释道，``当 \allowbreak{}\NovelInline{I03915}{\psi_a}\allowbreak{} 的自变量碰到一个1-投影序数的时候，它不是简单地取前一步的值再乘以 \allowbreak{}\NovelInline{I03916}{\omega}\allowbreak{}，而是取所有小于这个1-投影序数的自变量对应的 \allowbreak{}\NovelInline{I03917}{\psi_a}\allowbreak{} 值的上确界,或者叫做取并集。''

``这就是'投影'这个名字的由来------你把一个巨大的2-投影序数，'投影'到了一个由所有更小的序数构成的集合上，取它们的上确界。''

Phyrion在空中展开了第一个示例------

``因为 \allowbreak{}\NovelInline{I03918}{\Omega_{a+1}}\allowbreak{} 是一个 1-投影序数，所以 \allowbreak{}\NovelInline{I03919}{\psi_a(\Omega_{a+1})}\allowbreak{} 应当大于所有的 \allowbreak{}\NovelInline{I03920}{\psi_a(X)}\allowbreak{}，其中 \allowbreak{}\NovelInline{I03921}{X < \Omega_{a+1}}\allowbreak{}。''

``这意味着 \allowbreak{}\NovelInline{I03922}{\psi_a(\Omega_{a+1})}\allowbreak{} 大于：''

\NovelDisplay{M00747R01}{%
\psi_a(a),\quad \psi_a(\varepsilon_{a+1}),\quad \psi_a(\Gamma_{a+1}),\quad \psi_a(BO_{a+1}),\quad \psi_a(SHO_{a+1}),\quad \dots%
}

``更直观------虽然不太严谨------地说，你需要在里面取尽 \allowbreak{}\NovelInline{I03923}{a}\allowbreak{} 的所有递归运算。所有你能用 \allowbreak{}\NovelInline{I03924}{a}\allowbreak{} 构造出来的序数，不管多大，只要它小于 \allowbreak{}\NovelInline{I03925}{\Omega_{a+1}}\allowbreak{}，对应的 \allowbreak{}\NovelInline{I03926}{\psi_a}\allowbreak{} 值都已经被包含在了这个上确界之中。''

蔡须坤的金色双瞳微微闪烁。她开始理解了------这不是简单的''变大''，而是一种结构性的''折叠''，SHO这个序数何其巨大，已经达到了Y(1,3)的地步了。

Phyrion继续展开第二个示例------

``\allowbreak{}\NovelInline{I03927}{\psi_a(\Omega_{a+1}^3)}\allowbreak{}，它要大于：''

\NovelDisplay{M00748R01}{%
\psi_a(\Omega_{a+1}^2 \cdot a),\quad \psi_a(\Omega_{a+1}^2 \cdot \varepsilon_{a+1}),\quad \psi_a(\Omega_{a+1}^2 \cdot \Gamma_{a+1}),\quad \psi_a(\Omega_{a+1}^2 \cdot BO_{a+1}),\quad \dots%
}

``看到区别了吗？自变量从 \allowbreak{}\NovelInline{I03928}{\Omega_{a+1}}\allowbreak{} 变成了 \allowbreak{}\NovelInline{I03929}{\Omega_{a+1}^3}\allowbreak{}，那么上确界的范围就扩展到了所有小于 \allowbreak{}\NovelInline{I03930}{\Omega_{a+1}^3}\allowbreak{} 的序数。这些序数可以写成 \allowbreak{}\NovelInline{I03931}{\Omega_{a+1}^2 \cdot X}\allowbreak{} 的形式，其中 \allowbreak{}\NovelInline{I03932}{X}\allowbreak{} 取遍所有小于 \allowbreak{}\NovelInline{I03933}{\Omega_{a+1}}\allowbreak{} 的值。''

``再看一个更远的例子。\allowbreak{}\NovelInline{I03934}{\psi_a(\Omega_{a+2}^2)}\allowbreak{} 要大于：''

\NovelDisplay{M00749R01}{%
\psi_a(\Omega_{a+2} \cdot \varepsilon_{\Omega_{a+1}+1}),\quad \psi_a(\Omega_{a+2} \cdot \Gamma_{\Omega_{a+1}+1}),\quad \psi_a(\Omega_{a+2} \cdot BO_{\Omega_{a+1}+1}),\quad \dots%
}

``注意这里的下标------\allowbreak{}\NovelInline{I03935}{\Omega_{a+2}}\allowbreak{} 是 \allowbreak{}\NovelInline{I03936}{a}\allowbreak{} 之后的第二个 \allowbreak{}\NovelInline{I03937}{\Omega}\allowbreak{} 序数，而中间出现的 \allowbreak{}\NovelInline{I03938}{\varepsilon_{\Omega_{a+1}+1}}\allowbreak{}、\allowbreak{}\NovelInline{I03939}{\Gamma_{\Omega_{a+1}+1}}\allowbreak{} 等都是基于 \allowbreak{}\NovelInline{I03940}{\Omega_{a+1}}\allowbreak{} 构造的序数。投影的范围更大了，但规则是一样的。''

``现在，''Phyrion的声音微微一顿，``我要讲一个关键的区别。\allowbreak{}\NovelInline{I03941}{a}\allowbreak{} 和其他非递归序数------比如 \allowbreak{}\NovelInline{I03942}{M}\allowbreak{}------之间的本质差异，就体现在这里。''

他在空中写下了 \allowbreak{}\NovelInline{I03943}{M}\allowbreak{} 的对应规则：

``对于 \allowbreak{}\NovelInline{I03944}{M}\allowbreak{} 来说：''

\NovelDisplay{M00750R01}{%
\psi_M(\Omega_{M+1}) = \psi_M(\psi_{\Omega_{M+1}}(\psi_{\Omega_{M+1}}(\psi_{\Omega_{M+1}}(\dots))))%
}

``它只是找层展开------也就是不断地把 \allowbreak{}\NovelInline{I03945}{\psi_{\Omega_{M+1}}}\allowbreak{} 嵌套在自身之中。而非像 \allowbreak{}\NovelInline{I03946}{\psi_a(\Omega_{a+1})}\allowbreak{} 一样一定要取遍所有小于 \allowbreak{}\NovelInline{I03947}{\Omega_{a+1}}\allowbreak{} 的序数。''

``这是根本性的不同。\allowbreak{}\NovelInline{I03948}{M}\allowbreak{} 的展开是'逐层递归'，而 \allowbreak{}\NovelInline{I03949}{a}\allowbreak{} 的展开是'穷尽取并'。''

这时，白色母猪忽然抬起头，``哼哼''了两声。Phyrion微微一笑：``她说得对。让我补充一下当年 Opt 老祖的观点。''

\begin{quote}
\textbf{opt：} 更进一步地，不仅仅 \allowbreak{}\NovelInline{I03950}{\Omega_{X+1}}\allowbreak{} 是取并，\allowbreak{}\NovelInline{I03951}{\psi_a(I_{a+1})}\allowbreak{}、\allowbreak{}\NovelInline{I03952}{\psi_a(M_{a+1})}\allowbreak{}、甚至 \allowbreak{}\NovelInline{I03953}{\psi_a(a\text{ 容许点})}\allowbreak{} 都是取并。
\end{quote}

``最后，在结束基础部分之前，''Phyrion说，``让我从递归的角度给你们一个补充视角。''

``从递归的角度来说，我们记 \allowbreak{}\NovelInline{I03954}{\Omega_{a+1}}\allowbreak{} 为 \allowbreak{}\NovelInline{I03955}{D}\allowbreak{}：''

\NovelDisplay{M00751R01}{%
D = \Omega_{a+1}%
}

``然后我们有以下对应关系：''

\NovelDisplay{M00752R01}{%
\psi_{\psi_a(\# \sim D)}(0) = \psi_a(\# \sim 0)%
}

\NovelDisplay{M00753R01}{%
\psi_{\psi_a(\# \sim D)}(X+1) = \psi_a(\# \sim (X+1))%
}

\NovelDisplay{M00754R01}{%
\psi_{\psi_a(\# \sim D)}(\psi_a(\# \sim D)) = \psi_a(\# \sim a)%
}

\NovelDisplay{M00755R01}{%
\psi_{\psi_a(\# \sim D)}(\psi_a(\# \sim D)^2) = \psi_a(\# \sim a^2)%
}

``这里的 \allowbreak{}\NovelInline{I03956}{\#}\allowbreak{} 是一个占位符，代表任意的前缀表达式。\allowbreak{}\NovelInline{I03957}{D}\allowbreak{} 在递归视角下充当了一个'步进器'的角色------每增加一个 \allowbreak{}\NovelInline{I03958}{D}\allowbreak{}，就相当于在 \allowbreak{}\NovelInline{I03959}{\psi_a}\allowbreak{} 内部向前走了一步。''

``比如我们可以扽出 \allowbreak{}\NovelInline{I03960}{I = \psi_a(D^2)}\allowbreak{}，这样就可以得出：''

Phyrion在空中画出了一张完整的表格------

{\def\LTcaptype{none} % do not increment counter
\begin{longtable}[]{@{}
  >{\raggedright\arraybackslash}p{(\linewidth - 2\tabcolsep) * \real{0.5000}}
  >{\raggedright\arraybackslash}p{(\linewidth - 2\tabcolsep) * \real{0.5000}}@{}}
\toprule\noalign{}
\begin{minipage}[b]{\linewidth}\raggedright
左边
\end{minipage} & \begin{minipage}[b]{\linewidth}\raggedright
右边
\end{minipage} \\
\midrule\noalign{}
\endhead
\bottomrule\noalign{}
\endlastfoot
\allowbreak{}\NovelInline{I03961}{\psi_{\psi_a(D^2)}(0) = \psi_a(D \cdot 0) = \Omega}\allowbreak{} & \\
\allowbreak{}\NovelInline{I03962}{\psi_{\psi_a(D^2)}(1) = \psi_a(D \cdot 1) = \Omega_2}\allowbreak{} & \\
\allowbreak{}\NovelInline{I03963}{\psi_{\psi_a(D^2)}(2) = \psi_a(D \cdot 2) = \Omega_3}\allowbreak{} & \\
\allowbreak{}\NovelInline{I03964}{\psi_{\psi_a(D^2)}(\omega) = \psi_a(D \cdot \omega) = \Omega_\omega}\allowbreak{} & \\
\allowbreak{}\NovelInline{I03965}{\psi_{\psi_a(D^2)}(\psi_a(D^2)) = \psi_a(D \cdot a) = OFP}\allowbreak{} & \\
\allowbreak{}\NovelInline{I03966}{\psi_{\psi_a(D^2)}(\psi_a(D^2)^2) = \psi_a(D \cdot a^2) =}\allowbreak{} 第不动点个 \allowbreak{}\NovelInline{I03967}{OFP}\allowbreak{} & \\
\end{longtable}
}

``符合我们对 \allowbreak{}\NovelInline{I03968}{\psi_I}\allowbreak{} 的认知。''

唐三的眼睛亮了起来：``所以 \allowbreak{}\NovelInline{I03969}{D^2}\allowbreak{} 对应的是 \allowbreak{}\NovelInline{I03970}{I}\allowbreak{} 的递归结构？\allowbreak{}\NovelInline{I03971}{\psi_a(D^2)}\allowbreak{} 就相当于 \allowbreak{}\NovelInline{I03972}{I}\allowbreak{}，而 \allowbreak{}\NovelInline{I03973}{\psi_{\psi_a(D^2)}}\allowbreak{} 就是 \allowbreak{}\NovelInline{I03974}{\psi_I}\allowbreak{}？''

``没错。''Phyrion赞许地点了点头，``这就是递归视角的威力。它让你能用 \allowbreak{}\NovelInline{I03975}{D = \Omega_{a+1}}\allowbreak{} 的幂次来表示更高层级的序数结构。''

\begin{center}\rule{0.5\linewidth}{0.5pt}\end{center}

Phyrion停顿了一下，让众人消化这些信息。

白色母猪站起身来，缓缓绕着果糕小队走了一圈。她经过每个人身边时，都会释放出一缕温润的魂力，轻轻触碰他们的精神识海------那不是探查，而是一种''帮助理解''的引导。

诺当喃喃道：``这头猪\ldots\ldots 比我会教书。''

白色母猪''哼''了一声，用鼻子拱了拱他的小腿，像是在说''你才知道''。

Phyrion重新站定。

``基础概念讲完了。接下来进入详解部分。''

``首先，我们重新定义起点------''

\NovelDisplay{M00756R01}{%
\psi_a(0) = \Omega%
}

``然后，和 OCF 类似。

``那么 \allowbreak{}\NovelInline{I03976}{\psi_a}\allowbreak{} 里面每 \allowbreak{}\NovelInline{I03977}{+1}\allowbreak{}，和 BOCF 一样，对应的是整个序数*ω。''

``也就是说------''Phyrion在空中逐一展开公式------

\NovelDisplay{M00757R01}{%
\psi_a(1) = \psi_{\Omega_2}(0)%
}

\NovelDisplay{M00758R01}{%
\psi_a(2) = \psi_{\Omega_2}(1)%
}

\NovelDisplay{M00759R01}{%
\psi_a(\omega) = \psi_{\Omega_2}(\omega)%
}

\NovelDisplay{M00760R01}{%
\psi_a(\psi_a(0)) = \psi_{\Omega_2}(\Omega)%
}

\NovelDisplay{M00761R01}{%
\psi_a(\psi_a(1)) = \psi_{\Omega_2}(\psi_{\Omega_2}(0))%
}

``看到规律了吗？左边是 \allowbreak{}\NovelInline{I03978}{\psi_a}\allowbreak{} 的世界，右边是 \allowbreak{}\NovelInline{I03979}{\psi_{\Omega_2}}\allowbreak{} 的世界。\allowbreak{}\NovelInline{I03980}{\psi_a}\allowbreak{} 里面的每一个值，都能在 \allowbreak{}\NovelInline{I03981}{\psi_{\Omega_2}}\allowbreak{} 里面找到对应。''

``然后到达这样做的不动点。''

``在投影里面，\allowbreak{}\NovelInline{I03982}{\psi_a}\allowbreak{} 的不动点我们也用 \allowbreak{}\NovelInline{I03983}{a}\allowbreak{} 折叠，也就是：''

\NovelDisplay{M00762R01}{%
\psi_a(a) = \psi_{\Omega_2}(\Omega_2)%
}

奥斯卡愣住了：``等等------\allowbreak{}\NovelInline{I03984}{\psi_a(a) = \psi_a(\psi_a(\psi_a(\dots)))}\allowbreak{} 呀？''

``没错。''Phyrion说，``这就是第三条规则的体现。当自变量等于 \allowbreak{}\NovelInline{I03985}{a}\allowbreak{} 的时候，\allowbreak{}\NovelInline{I03986}{\psi_a(a)}\allowbreak{} 就是无限嵌套 \allowbreak{}\NovelInline{I03987}{\psi_a(\psi_a(\psi_a(\dots)))}\allowbreak{} 的不动点。而这个不动点，恰好对应 \allowbreak{}\NovelInline{I03988}{\psi_{\Omega_2}(\Omega_2)}\allowbreak{}。''

``\allowbreak{}\NovelInline{I03989}{a}\allowbreak{} 在 \allowbreak{}\NovelInline{I03990}{\psi_a}\allowbreak{} 中的角色，就如同 \allowbreak{}\NovelInline{I03991}{\Omega_2}\allowbreak{} 在 \allowbreak{}\NovelInline{I03992}{\psi_{\Omega_2}}\allowbreak{} 中的角色。''

``那么后面仍然是类似的------''

\NovelDisplay{M00763R01}{%
\psi_a(a+1) = \psi_{\Omega_2}(\Omega_2 + 1)%
}

\NovelDisplay{M00764R01}{%
\psi_a(a \times 2) = \psi_{\Omega_2}(\Omega_2 \times 2)%
}

``在这里，\allowbreak{}\NovelInline{I03993}{\psi_a}\allowbreak{} 里面的 \allowbreak{}\NovelInline{I03994}{a}\allowbreak{} 发挥的作用和 \allowbreak{}\NovelInline{I03995}{\psi_{\Omega_2}}\allowbreak{} 里面的 \allowbreak{}\NovelInline{I03996}{\Omega_2}\allowbreak{} 是完全一样的。''

``一直到------''Phyrion看着众人，``有人能猜到到哪里都一样吗？''

唐三沉思片刻：``一直到 \allowbreak{}\NovelInline{I03997}{\psi_a(\varepsilon_{a+1})}\allowbreak{} 都一样的？''

``对的。''

\NovelDisplay{M00765R01}{%
\psi_a(\varepsilon_{a+1}) = \psi_{\Omega_2}(\psi_{\Omega_3}(\Omega_3))%
}

``这个 \allowbreak{}\NovelInline{I03998}{\varepsilon_{a+1}}\allowbreak{}，按照 MOCF 的规则也可以写成 \allowbreak{}\NovelInline{I03999}{\psi_{\Omega_{a+1}}(\Omega_{a+1})}\allowbreak{}。''

``也就是：''

\NovelDisplay{M00766R01}{%
\psi_a(\psi_{\Omega_{a+1}}(\Omega_{a+1})) = \psi_{\Omega_2}(\psi_{\Omega_3}(\Omega_3))%
}

霍雨浩的瞳孔微微收缩。他开始看到了一种深层的对称结构------左边用 \allowbreak{}\NovelInline{I04000}{\Omega_{a+1}}\allowbreak{} 构造的东西，右边用 \allowbreak{}\NovelInline{I04001}{\Omega_2}\allowbreak{}（也就是对应 \allowbreak{}\NovelInline{I04002}{\psi_{\Omega_2}}\allowbreak{} 的第二个 \allowbreak{}\NovelInline{I04003}{\Omega}\allowbreak{} 序数）构造的东西，完全一一对应。

``哪些细节？''

``继续，加快一点。''Phyrion说。

\NovelDisplay{M00767R01}{%
\psi_a(\psi_{\Omega_{a+1}}(a)) = \psi_{\Omega_2}(\psi_{\Omega_3}(\Omega_2))%
}

\NovelDisplay{M00768R01}{%
\psi_a(\psi_{\Omega_{a+1}}(\psi_{\Omega_{a+1}}(0))) = \psi_{\Omega_2}(\psi_{\Omega_3}(\psi_{\Omega_3}(0)))%
}

``一直到 \allowbreak{}\NovelInline{I04004}{\psi_{\Omega_2}(\Omega_3)}\allowbreak{}。''

Phyrion的声音在这一刻变得格外清晰------

``\textbf{这个地方就开始出现差异了。}''

\NovelDisplay{M00769R01}{%
\psi_a(\psi_{\Omega_{a+1}}(\Omega_{a+1})) = \psi_{\Omega_2}(\Omega_3)%
}

``看到没有？左边有两层------\allowbreak{}\NovelInline{I04005}{\psi_a}\allowbreak{} 嵌套着 \allowbreak{}\NovelInline{I04006}{\psi_{\Omega_{a+1}}}\allowbreak{}，而 \allowbreak{}\NovelInline{I04007}{\psi_{\Omega_{a+1}}}\allowbreak{} 的自变量是 \allowbreak{}\NovelInline{I04008}{\Omega_{a+1}}\allowbreak{} 本身。但右边只剩一层------\allowbreak{}\NovelInline{I04009}{\psi_{\Omega_2}(\Omega_3)}\allowbreak{}。''

``这就是投影的'层级压缩'现象。左边的双层结构，被投影压缩成了右边的单层结构。''

``然后继续------''

\NovelDisplay{M00770R01}{%
\psi_a(\psi_{\Omega_{a+1}}(\Omega_{a+1}) + 1) = \psi_{\Omega_2}(\Omega_3 + 1)%
}

\NovelDisplay{M00771R01}{%
\psi_a(\psi_{\Omega_{a+1}}(\Omega_{a+1}) \times 2) = \psi_{\Omega_2}(\Omega_3 + \psi_{\Omega_3}(\Omega_3))%
}

奥斯卡挠了挠头：``这个\ldots\ldots{}\allowbreak{}\NovelInline{I04010}{\times 2}\allowbreak{} 怎么变成加号了？''

``好问题。''Phyrion说，``这里确实跳步了。让我细扽一下------其实很简单：''

\NovelDisplay{M00772R01}{%
\psi_a(\psi_{\Omega_{a+1}}(\Omega_{a+1}) + a) = \psi_{\Omega_2}(\Omega_3 + \Omega_2)%
}

\NovelDisplay{M00773R01}{%
\psi_a(\psi_{\Omega_{a+1}}(\Omega_{a+1}) + \psi_{\Omega_{a+1}}(0)) = \psi_{\Omega_2}(\Omega_3 + \psi_{\Omega_3}(0))%
}

\NovelDisplay{M00774R01}{%
\psi_a(\psi_{\Omega_{a+1}}(\Omega_{a+1}) + \psi_{\Omega_{a+1}}(\psi_{\Omega_{a+1}}(0))) = \psi_{\Omega_2}(\Omega_3 + \psi_{\Omega_3}(\psi_{\Omega_3}(0)))%
}

``看，左边的加法部分------\allowbreak{}\NovelInline{I04011}{+ a}\allowbreak{}、\allowbreak{}\NovelInline{I04012}{+ \psi_{\Omega_{a+1}}(0)}\allowbreak{}、\allowbreak{}\NovelInline{I04013}{+ \psi_{\Omega_{a+1}}(\psi_{\Omega_{a+1}}(0))}\allowbreak{}------对应右边的 \allowbreak{}\NovelInline{I04014}{+ \Omega_2}\allowbreak{}、\allowbreak{}\NovelInline{I04015}{+ \psi_{\Omega_3}(0)}\allowbreak{}、\allowbreak{}\NovelInline{I04016}{+ \psi_{\Omega_3}(\psi_{\Omega_3}(0))}\allowbreak{}。一一对应。''

``取极限，就有：''

\NovelDisplay{M00775R01}{%
\psi_a(\psi_{\Omega_{a+1}}(\Omega_{a+1}) \times 2) = \psi_{\Omega_2}(\Omega_3 + \psi_{\Omega_3}(\Omega_3))%
}

``这就是那个跳步的结果。''

``继续------''

\NovelDisplay{M00776R01}{%
\psi_a(\psi_{\Omega_{a+1}}(\Omega_{a+1})^2) = \psi_{\Omega_2}(\Omega_3 + \psi_{\Omega_3}(\Omega_3)^2)%
}

``这里不是重点，讲快一点------''

\NovelDisplay{M00777R01}{%
\psi_a(\psi_{\Omega_{a+1}}(\Omega_{a+1} + 1)) = \psi_{\Omega_2}(\Omega_3 + \psi_{\Omega_3}(\Omega_3 + 1))%
}

\NovelDisplay{M00778R01}{%
\psi_a(\psi_{\Omega_{a+1}}(\Omega_{a+1} + a)) = \psi_{\Omega_2}(\Omega_3 + \psi_{\Omega_3}(\Omega_3 + \Omega_2))%
}

\NovelDisplay{M00779R01}{%
\psi_a(\psi_{\Omega_{a+1}}(\Omega_{a+1} + \psi_{\Omega_{a+1}}(0))) = \psi_{\Omega_2}(\Omega_3 + \psi_{\Omega_3}(\Omega_3 + \psi_{\Omega_3}(0)))%
}

``类似地，再取极限，就是：''

\NovelDisplay{M00780R01}{%
\psi_a(\psi_{\Omega_{a+1}}(\Omega_{a+1} \times 2)) = \psi_{\Omega_2}(\Omega_3 \times 2)%
}

Phyrion指着这个公式说：``可以发现左边又比右边多一层了。''

奥斯卡嘟囔道：``我还以为是 \allowbreak{}\NovelInline{I04017}{\psi_{\Omega_2}(\Omega_3 + \Omega_2)}\allowbreak{} 诶\ldots\ldots{}''

``总之这些都是细节，不是重点。\textbf{重点是这里的现象值得关注。}''

``什么现象？''蔡须坤问。

``层级压缩的现象。''Phyrion说，``左边的嵌套结构在特定节点会被压缩成右边的更简洁的形式。但这种压缩不是随意的------它遵循严格的对应规则。''

\begin{center}\rule{0.5\linewidth}{0.5pt}\end{center}

\chapter{第七十八章：小队思考和温情}\label{ux7b2cux4e03ux5341ux516bux7ae0ux5c0fux961fux601dux8003ux548cux6e29ux60c5}

火堆旁的沉默持续了很久。

最终，蔡须坤开口了：``所以，这和我们突破 0 111 22 有什么关系？''

Phyrion看着她，粉红色的瞳孔中闪过一丝赞许。

``0 111 22 的壁垒，就是反射序数的极限，即便你们使出吃奶的力气拓展，也很难到达0 111 22 331

蔡须坤闭上了眼睛。

在那面镜子中，她看到了自己的序列------那些展开的部分，那些她从未注意到的盲区。

她也看到了牢大。

不是牢大的身影，而是牢大留在她中的印记------那些在无数次并肩修炼中，两颗魂核相互映照、相互投影所留下的痕迹。

``原来如此。''她低声说。

泪水无声地滑落，但她的嘴角却微微上翘。

``牢大\ldots\ldots 你一直都在。''

\begin{center}\rule{0.5\linewidth}{0.5pt}\end{center}

夜深了。

Phyrion站在修炼堂的入口处，回头看了一眼那些沉浸在冥想中的年轻人。

白色母猪趴在他的脚边，粉红色的圆眼睛半眯着。

``他们能突破吗？''Phyrion低声问。

白色母猪''哼''了一声，那声音里带着一丝笃定。

Phyrion笑了笑。

``也是。我教猪配种的时候，成功率都有百分之九十八。''

白色母猪用鼻子狠狠拱了他的小腿一下。

``疼疼疼------开玩笑的。''

夜风拂过，带走了最后一丝喧嚣。

灵气如潮水般在空旷的地下宫殿中涌动，魂力晶石散发出柔和的光芒，将整个空间笼罩在一片淡金色的宁静之中。

果糕小队的众人已经结束了今天长达十二个时辰的修炼。Phyrion早已带着白色母猪离开了，留下他们自行消化今天所学的投影序数知识。

但没有人立刻入睡。

\begin{center}\rule{0.5\linewidth}{0.5pt}\end{center}

唐三盘坐在石室的角落，双膝之上摊开着那本《投影序数》，指尖轻轻翻过一页又一页。他的眉头微微蹙起，显然是在某个推导环节遇到了瓶颈。

一杯温热的茶忽然递到了他的面前。

``喝点吧，你都看了三个时辰了。''

霍雨浩站在他的身旁，手里端着一个粗陶茶壶。那茶壶是他从修炼堂的储物间里翻出来的，里面泡的是修炼堂特供的灵茶------用这里浓郁的灵气滋养过的茶叶，入口温润，能舒缓过度运转的精神力。

唐三接过茶杯，微微一笑：``谢了。''

``你那个地方卡住了？''霍雨浩在他身旁坐下，目光扫过摊开的书页。

``\allowbreak{}\NovelInline{I04018}{\psi_a(\psi_{\Omega_{a+1}}(\Omega_{a+1}) \times 2)}\allowbreak{} 到 \allowbreak{}\NovelInline{I04019}{\psi_{\Omega_2}(\Omega_3 + \psi_{\Omega_3}(\Omega_3))}\allowbreak{} 这一步。''唐三指着书页上的一行公式，``Phyrion说这里跳步了，但我总觉得中间的推导不应该只是'细扽一下'那么简单。''

霍雨浩沉思了片刻，他的灵眸在黑暗中微微发光------那是精神力高度集中的表现。

``你试试这样想，''他缓缓说道，``\allowbreak{}\NovelInline{I04020}{\Omega_{a+1}}\allowbreak{} 到 \allowbreak{}\NovelInline{I04021}{\Omega_3}\allowbreak{} 的对应，本质上是投影压缩了一整层递归。所以 \allowbreak{}\NovelInline{I04022}{\psi_{\Omega_{a+1}}(\Omega_{a+1})}\allowbreak{} 里面包含的所有内容，在投影之后都会被压缩到 \allowbreak{}\NovelInline{I04023}{\psi_{\Omega_3}(\Omega_3)}\allowbreak{} 的层级里。\allowbreak{}\NovelInline{I04024}{\times 2}\allowbreak{} 的意思，就是把这整层内容复制一份叠加上去。''

唐三的瞳孔微微收缩，随即露出了恍然大悟的表情。

``对\ldots\ldots 我之前一直把它当成简单的乘法来处理，但实际上它是对整个递归结构的一次完整复制。''

霍雨浩点了点头，嘴角带着一丝笑意：``所以Phyrion说'细扽一下就很简单'，其实是因为他站在 \allowbreak{}\NovelInline{I04025}{\psi_{\Omega_2}}\allowbreak{} 的视角往回看，当然简单。但对于我们这些还在 \allowbreak{}\NovelInline{I04026}{\psi_a}\allowbreak{} 层级挣扎的人来说，这一步至少要转三个弯。''

唐三忍不住笑了：``Phyrion教猪配种的时候有耐心，教我们的时候倒是经常跳步。''

两人相视一笑，石室里的气氛轻松了不少。

\begin{center}\rule{0.5\linewidth}{0.5pt}\end{center}

另一边，奥斯卡正四仰八叉地躺在石板上，双手抱头，嘴里嘟嘟囔囔地念叨着什么。

``\allowbreak{}\NovelInline{I04027}{\psi_a(0) = \Omega}\allowbreak{}\ldots\ldots{}\allowbreak{}\NovelInline{I04028}{\psi_a(1) = \psi_{\Omega_2}(0)}\allowbreak{}\ldots\ldots{}\allowbreak{}\NovelInline{I04029}{\psi_a(a) = \psi_{\Omega_2}(\Omega_2)}\allowbreak{}\ldots\ldots{}''

他的声音越来越小，越来越含糊，眼看就要睡着了。

一条毯子轻轻地盖在了他的身上。

奥斯卡迷迷糊糊地睁开眼睛，看到卢本韦正弯着腰，把一条从储物间里找来的羊毛毯子盖到他身上。

``你这样躺着会着凉的。''卢本韦的声音一如既往地平淡，但动作却很轻，生怕惊到他似的，``修炼堂的灵气虽然浓郁，但地底的温度比外面低不少。''

``谢了老卢。''奥斯卡嘿嘿一笑，把毯子裹了裹，``你说我是不是太笨了？Phyrion今天讲的那些公式，蔡须坤和唐三一听就懂，霍雨浩自己还能推导，我就只能死记硬背。''

卢本韦在他旁边坐了下来，沉默了一会儿。

``你和他们不一样。''她说。

``什么意思？''

``蔡须坤是队长，她的Y序列感知力是天赋。唐三是理论大师，霍雨浩有灵眸加持的精神力。他们的路和你不同，不代表你比他们差。''

卢本韦顿了一下，目光落在奥斯卡手中的笔记本上------那上面密密麻麻地写满了公式，每一行旁边都标注着奥斯卡自己的理解和注释，字迹虽然潦草，但内容极其详尽。

``你今天记了整整三十页笔记。''卢本韦说，``全队最多的就是你。''

奥斯卡愣了一下，低头看了看自己的笔记本，忽然笑了。

``嘿，这么一说好像还真是。''

``笨鸟先飞，也是一种天赋。''卢本韦站起身，拍了拍奥斯卡的肩膀，``早点睡。明天Phyrion还要讲递归视角的补充，那部分更烧脑。''

``知道了知道了。''奥斯卡把毯子拉到了下巴，``老卢你也早点休息啊。''

卢本韦没有回头，只是挥了挥手，但是她的心里，早就波涛汹涌。

\begin{center}\rule{0.5\linewidth}{0.5pt}\end{center}

蔡须坤一个人坐在修炼堂最高处的石台上。

这里是整个地下空间的制高点，可以俯瞰整个修炼堂。魂力晶石的光芒从穹顶洒落，像是夜空中的星辰。

她的手中握着那本《投影序数》，但目光却不在书上。她在看穹顶。

那些闪烁的晶石，让她想起了很久以前的一个夜晚。那时候果糕小队还只是一个籍籍无名的小队，他们在荒野中扎营，头顶也是这样一片星光。

牢大坐在她的身旁，给她讲bms的基础理论。

``须坤，你知道为什么bms要从0开始吗？''``因为万物从虚无中诞生？''

``不。因为0代表着无限的可能性。在0的位置上，你可以成为任何东西。''

蔡须坤的金色双瞳在晶石的光芒中微微闪烁。``队长。''一个声音打断了她的回忆。

她转过头，看到唐三不知什么时候走到了石台的旁边，手里端着两杯灵茶。

``睡不着？''唐三将一杯茶递给她。``嗯。''蔡须坤接过茶杯，轻轻吹了吹，``在想一些事情。''

唐三没有追问。他在她身旁坐下，两人沉默地喝了一会儿茶。``今天的课，你觉得怎么样？''唐三问。

``Phyrion讲得很好。''蔡须坤说，``但我总觉得，他教我们的不仅仅是投影序数本身。''

``你是说\ldots\ldots{}''

``他在教我们一种思维方式。''蔡须坤的目光变得深邃，``投影的本质不是计算，而是视角的转换。

``而我们现在之所以卡在 0 111 22，就是因为我们还站在序列内部看序列。我们需要跳出来，从更高的维度审视自己。''

唐三静静地听着，然后轻轻点了点头。

``你说得对。Phyrion让我们把武魂当作'伙伴'而不是'工具'，其实也是同一个道理------武魂是我们的一面镜子，它能从外部看到我们看不到的东西。''

蔡须坤微微一笑，低头喝了一口茶。

``唐三。''``嗯？''``谢谢你。谢谢大家。''

唐三有些意外地看了她一眼。

蔡须坤的声音很轻，但很认真：``我知道大家修炼都很累。但我看到奥斯卡笔记记得比谁都认真，看到卢本韦默默照顾每一个人，看到你和霍雨浩深夜还在讨论公式，看到大家都在拼命\ldots\ldots{}''

她的声音微微颤了一下。

``牢大以前跟我说过，一个人的力量是有限的。但一个小队的力量，是无限的。''``我现在明白了。''

唐三沉默了片刻，然后缓缓说道：``队长，你不用一个人扛。''

蔡须坤微微一怔。

``牢大不在了，但你还有我们。''唐三的声音平静而坚定，``果糕小队不只是你的责任，也是我们每个人的。你照顾我们的时候，也让我们照顾你。''

蔡须坤的金色双瞳中闪过一丝水光。

她没有说话，只是轻轻地点了点头。

头顶的魂力晶石依旧闪烁着柔和的光芒，像是一片永不熄灭的星空。

\begin{center}\rule{0.5\linewidth}{0.5pt}\end{center}

石室的另一端，霍雨浩正在帮卢本韦按摩肩膀。

``你今天修炼太猛了，肩膀的经脉都僵硬了。''霍雨浩的手法很专业，灵眸的精神力让他能精确地感知到肌肉和经脉的每一个细微变化。

``没事，明天就好了。''卢本韦嘴硬。

``明天Phyrion要讲 \allowbreak{}\NovelInline{I04030}{a_2}\allowbreak{} 和 \allowbreak{}\NovelInline{I04031}{a_3}\allowbreak{} 的内容，那部分需要大量的精神力运算。你现在不把经脉疏通，明天脑子会转不动的。''

卢本韦沉默了一下，然后老老实实地坐好，让霍雨浩继续按。

``雨浩。''``嗯？''

``你那个精神探测\ldots\ldots 能不能帮我看看，我的Y序列现在到底卡在哪个位置？''

霍雨浩的灵眸微微闪烁，精神力如同水流一般渗入卢本韦的经脉之中。

片刻后，他收回了精神力，若有所思地说：

``你的展开态已经基本完成了，'0 111'的部分很稳定。但收敛态'22'还没有形成。你的bms现在就像一条只放出去、但还没收回来的线。''

``那怎么办？''

``Phyrion说的阴阳互补双魂核，一阴一阳，一个递归一个非递归。你现在只有非递归展开，没有递归化。你需要找到那个让序列向内折叠的契机。''

卢本韦皱了皱眉：``怎么找？''

霍雨浩想了想，忽然笑了。

``这个我不确定。但Phyrion说明天白色母猪会帮我们一个一个感知盲区。也许她能帮你找到答案。''

卢本韦无奈地叹了口气：``堂堂封号斗罗级别的强者，要靠一头猪来找答案。''

``她不是普通的猪。''霍雨浩笑着说，``她可是Phyrion的魂灵。''``有区别吗？''

``区别大了。普通的猪只会哼哼。她会哼哼着嫌你笨。''卢本韦忍不住笑出了声。

夜更深了。修炼堂里安静了下来，只剩下灵气流动的细微声响。

奥斯卡裹着卢本韦给他盖的毯子，打着轻微的呼噜。唐三和霍雨浩各自回到石室，继续默默推演着公式。卢本韦的肩膀松快了不少，终于沉沉地睡去了。

而蔡须坤依旧坐在最高处的石台上。

她合上了《投影序数》，将它轻轻放在膝上。

然后她抬起头，望着穹顶的晶石，嘴角微微上扬。

``牢大，''她在心中默默地说，``你说得对。一个人的力量是有限的。''``但我有一群最好的伙伴。''

``我们一定会突破 0 111 22。''``一定。''

魂力晶石的光芒温柔地洒在她的身上，像是某种无声的祝福。在这个寂静的夜晚，皇家修炼堂里的八个人，各自怀揣着梦想与牵挂，沉入了梦乡。

而明天，Phyrion的课还会继续。投影序数的奥秘，还远远没有揭开。

唐三沉默了很久，然后缓缓说道：

``所以投影序数的本质，就是不断地'跳级'。每跳一次，体系就变得更强大。而每一次跳级，都建立在前一次跳级的基础之上。''

``可以这么理解。''Phyrion说。

``那\ldots\ldots{}''霍雨浩犹豫了一下，``这个跳级的过程，有没有尽头？''

Phyrion看了他一眼，粉红色的瞳孔中闪过一丝复杂的光芒。

``在投影体系内部，没有尽头。\allowbreak{}\NovelInline{I04032}{a_n}\allowbreak{} 可以引入任意多个，每一层都建立在前一层之上。但如果我们跳出投影体系本身\ldots\ldots{}''

他没有继续说下去。

白色母猪在这一刻''哼''了一声。那声音低沉而悠长，像是在回应Phyrion未尽的话语。

果糕小队众人感觉到了一股奇异的波动------那不是魂力的波动，而是某种更深层的东西。像是整个修炼堂的灵气都在那一瞬间微微震颤了一下。

然而，最难的部分，是最小的lift展开

phyrion:老夫修为达到Y(1,3,10)你们知道为什么吗，众人不解， 他慢慢说到：我的武魂也是魂灵，你们要知道，当年我的父亲，也是本韦的曾祖，曾经游历大陆

``好了。''Phyrion拍了拍手，将众人的注意力拉了回来，``今天的课就到这里。你们需要消化一下

他转身准备离开，走了两步又停了下来，回头看了一眼白色母猪。

``对了，她今天会留在你们身边。如果你们在推导过程中遇到卡住的地方，可以让她帮你们感知一下盲区。''

白色母猪''哼''了一声，慢悠悠地站了起来，走到果糕小队众人面前，粉红色的圆眼睛一个一个地扫过他

Phyrion站在石台中央，粉红色的魂力在他周身缓缓流转，像是一层薄薄的光纱。白色母猪趴在他的脚边，粉红色的圆眼睛半眯着，像是在听，又像是在审视。

沉默持续了很久。然后Phyrion开口了。

``老夫修为达到 \allowbreak{}\NovelInline{I04033}{Y(1,3,10)}\allowbreak{}。''他的声音不大，但每一个字都清晰地落在了每个人的耳中。

``你们知道为什么吗？''众人面面相觑。

蔡须坤第一个反应过来：``因为你的武魂是魂灵------有意识、有灵性的魂灵。她不是工具，而是伙伴。''

``对，也不全对。''Phyrion缓缓摇了摇头，``武魂是魂灵，这只是条件之一。真正的原因\ldots\ldots{}''

他停顿了一下，目光变得悠远，像是穿透了修炼堂的穹顶，望向了很久很久以前的某个地方。

``要从我的父亲说起。''``卢承业。''``也是本韦的曾祖。''

卢本韦的身体微微一震。她张了张嘴，但最终没有说话，只是静静地听着。

``我的父亲，卢承业。和你们很多人想象的不一样，他不是一个强大的魂师。他的魂力天赋平平，修炼速度缓慢，年轻时候甚至被人嘲笑是'卢瑟福家的废物'。''

``但他有一样东西，是任何人都比不了的。''``耐心。''

``他花了整整三十年的时间，游历大陆。不是为了修炼，不是为了战斗，而是为了研究一个问题------''

Phyrion的声音低沉了下来。``魂灵。''

``在那个年代，魂灵的概念还不存在。魂师与武魂之间的关系，被所有人视为理所当然的主从关系------魂师驱动武魂，武魂执行意志。没有人想过，武魂可以有独立的意识。''

``opt老祖记载：武魂不应该只是工具。如果能让武魂拥有独立的意识，让魂师与武魂之间形成真正的共生关系，那么魂师的力量将会发生质的飞跃。''``他开始研究。''

``他走访了大陆上所有的魂师家族、所有的宗门、所有的隐世高人。他收集了数以千计的武魂样本，记录了数以万计的魂师与武魂之间的互动数据。他试图找到一种方法，让武魂从'被驱使的工具'变成'有意识的伙伴'。''

``但他失败了。''``一次又一次地失败。''Phyrion的声音里带上了一丝苦涩。

``他尝试过用魂力灌注来唤醒武魂的意识------失败了。武魂在魂力灌注下只会变得狂暴，失去控制。''

``他尝试过用精神链接来建立魂师与武魂之间的意识通道------失败了。精神链接建立的一瞬间，魂师的精神识海就会被武魂的本能意识冲垮。''

``他尝试过用药物、用阵法、用古老的仪式------全部失败了。''

``每一次失败，都会让他离答案更近一步。但每一次，那最后一步，他始终迈不过去。''

``直到他遇到了一位老修士。''

\begin{center}\rule{0.5\linewidth}{0.5pt}\end{center}

Phyrion的目光变得格外柔和，像是透过记忆的迷雾，看到了一个遥远的身影。

``那是在大陆东南方的一座小岛上。那座岛没有名字，地图上找不到它的位置。我父亲是在一次航海中偶然发现它的------暴风雨把他的船吹偏了航线，他在海上漂流了七天七夜，最后被海浪冲上了那座岛的沙滩。''

``岛上有一个人。''

``一个老修士。没有人知道他的名字，没有人知道他来自哪里，没有人知道他活了多久。我父亲说，那个老修士看起来没有那么的华贵。

``但我父亲说，他一眼就看出来了------那个老修士的实力，深不可测。''

``不是魂力的深不可测。而是某种更深层的东西。那个老修士隐居于岛内的一座高山之巅，但他的气息却能遍布整座岛屿。我父亲后来回忆说，他踏上那座岛的第一步，就感觉到自己被一股无形的力量笼罩了------那气息无处不在，在空气中，在海水中，在沙砾中，在每一片树叶的脉络里。整座岛屿就像是一个巨大的阵法，而那个老修士，就是阵法的核心。''

``我父亲顺着那股气息，攀上了高山，找到了老修士的住所。''

``他在茅庐前坐下，把自己的研究全部告诉了老修士。''

``老修士听完之后，沉默了很久。''

``然后他说了一句话。''

Phyrion的声音在这一刻变得异常缓慢，像是在一字一句地重现当年那个老修士的语气。

``\textbf{生灵之本，在于序。}''``\textbf{无序，则本源混乱。}''

``我父亲愣住了。他反复咀嚼这句话，越咀嚼越觉得震撼。''

``生灵之本，在于序------所有生灵的本质，都建立在'秩序'之上。魂力是秩序，经脉是秩序，魂核是秩序，甚至连意识本身，也是一种秩序。''

``无序，则本源混乱------如果秩序被打乱，生灵的本源就会崩溃。这就是为什么他的每一次尝试都会失败------他试图用外部的力量去唤醒武魂的意识，但没有先建立起'秩序'。没有秩序的意识，就是混乱。混乱的意识，只会摧毁一切。''

``我父亲在那座岛上住了三年。''

``三年里，老修士教了他很多东西。不是魂力修炼的方法，不是战斗的技巧，而是------如何感知秩序。如何在一个生灵的体内，找到那条贯穿一切的'序线'。如何让意识沿着序线生长，而不是像野草一样四处蔓延。''

``三年后，我父亲离开了那座岛。他带着老修士教给他的一切，继续研究魂灵。''

``他又花了二十年。''

``二十年里，他建立了魂灵理论的基础框架。他提出了'意识序线'的概念------每一个生灵的体内，都存在一条无形的序线，它贯穿精神识海、经脉、魂核，将所有的一切串联成一个有序的整体。''

``他提出了'本源规整'的方法------通过感知序线，将武魂散乱的本能意识沿着序线重新排列，使之从混乱变为有序。''

``但他到死，也没有完成最后一步。''

\begin{center}\rule{0.5\linewidth}{0.5pt}\end{center}

Phyrion闭上了眼睛。

``他找不到'规整意识、梳理本源'的具体方法。他知道原理，他知道方向，但他找不到那条路。''

``他临终前把我叫到床前，把他所有的研究手稿、笔记、数据，全部交给了我和我哥。''

``他说：`孩子，我走了一辈子的路，走到了这里。剩下的路，要靠你自己走了。'\,''

``然后他就走了。''

修炼堂里安静得能听到灵气流动的声音。

卢本韦低下了头。她的手指紧紧攥着衣角，指节发白。那是她的曾祖------一个她从未见过面、却流淌在她血脉中的人。

``我继承了父亲的一切。''Phyrion的声音恢复了平静，但那种平静之下，藏着极深的力量，``他的理论，他的方法，他的失败，他的遗憾。我花了十五年的时间，在他的基础上继续研究。''``然后我找到了。''

``父亲失败的原因，不是他的理论有错。而是他缺少一样东西。''``什么？''唐三问。``一个活的合作者。''

Phyrion低下头，看着脚边的白色母猪。

``父亲的研究对象，都是已经契约的武魂。那些武魂在契约之后，意识已经被魂师的精神力压制了。它们不是没有意识，而是意识被封锁了。父亲试图从外部打破这个封锁，但每一次打破，都会导致意识的混乱------因为封锁本身就是一种秩序，打破封锁就等于打破秩序。''

``而我换了一个思路。''``我不去打破封锁。我去和武魂对话。''

``我在契约之前，就和她建立了联系。不是精神链接，不是魂力灌注，而是------对话。我用我的意识去触碰她的意识，沿着她体内的序线，一点一点地引导她的意识从混乱走向有序。''``这个过程花了整整十年。''

``十年里，她从一个只有本能反应的魂兽，慢慢变成了有独立思考能力的存在。她学会了感知序线，学会了梳理自己的本源，学会了用秩序来约束自己的力量。''

``然后有一天，她开口说了第一句话。''

白色母猪在这一刻''哼''了一声。Phyrion笑了。``她说的是------`你太慢了'。''

众人都忍不住笑了。紧张的气氛稍稍缓解了一些。

``从那天起，我们就是真正的伙伴了。''Phyrion的声音变得柔和，``她帮我感知我感知不到的东西，我帮她梳理她梳理不了的东西。我们的Y序列在共生中不断成长，最终走到了 \allowbreak{}\NovelInline{I04034}{Y(1,3,10)}\allowbreak{}。''

``这就是我能走到这一步的原因。不是因为我的天赋比别人高，而是因为我继承了父亲卢承业的智慧，然后用十年的时间，和一个有意识的魂灵建立了真正的共生关系。''

沉默再次降临。

然后Phyrion的声音忽然变了。

那是一种很微妙的变化------语速没有变，音量没有变，但所有人都感觉到了一股淡淡的苦涩，像是秋天的风里夹杂着一丝寒意。``可惜\ldots\ldots 我那哥哥。''他叹了口气。

那声叹息很轻，但在寂静的修炼堂里，却像是一块石头落入了深潭。

``卢承业一生痴迷于魂灵的研究，对我的哥哥------也就是本韦的祖父------疏于管教。我哥哥从小在父亲的光环下长大，心里既崇拜又怨恨。他渴望得到父亲的认可，但父亲的眼里永远只有那些手稿和公式。''

``后来父亲去世了，哥哥继承了卢家的家业，也继承了父亲的部分手稿。但他没有继承父亲的耐心。''

``他试图走捷径。''``他想用父亲留下的理论框架，强行制造魂灵。不是用十年的时间去和一个武魂慢慢建立联系，而是用阵法、用药物、用精神力的暴力灌注，试图在一夜之间唤醒武魂的意识。''

Phyrion的声音低沉得几乎听不见。

``他失败了。''``不，比失败更糟。''

``他临终前把卢关叫到床前，把我父亲留给他的手稿，连同他自己这些年的研究笔记，一起交给了卢关。

``他说：`不要走我的老路。秩序不能强求，只能引导。'\,''

Phyrion的目光落在了卢本韦的身上。

``所以本韦，你身上流淌的，不只是你曾祖的血脉。还有你祖父用生命换来的教训。''

卢本韦的身体微微颤抖了一下。她张了张嘴，似乎想说什么，但最终只是用力地点了点头。

``哎。''Phyrion又叹了口气，``我哥哥是个好人。只是太急了。''

他没有继续说下去。

白色母猪在这一刻站了起来，慢慢地走到Phyrion的身旁，用脑袋轻轻蹭了蹭他的小腿。那动作温柔得不像是一头封号斗罗级别的魂灵，更像是一个沉默的伙伴，在用自己的方式安慰着自己的搭档。

Phyrion低下头，伸手摸了摸她的脑袋。

``不说了。''他说，``今天讲到这里。''他的目光扫过众人，粉红色的瞳孔中恢复了往日的平静。

``你们要记住------魂灵不是天生的。它是用时间、耐心和信任，一点一点养出来的。你们每个人的武魂，都有可能成为你们的魂灵。但前提是------你们要愿意花时间去了解它，去尊重它，去把它当作一个独立的生命来对待。''

``而不是一个工具。''

``秩序不能强求，只能引导。这是我父亲用一辈子悟出来的道理，也是我哥哥用一辈子验证过的教训。''

他转身向修炼堂的出口走去。白色母猪跟在他的身后，粉红色的圆眼睛最后扫了果糕小队众人一眼。

那一眼里，有太多太多的东西。走了几步，Phyrion忽然停了下来，没有回头。

``本韦。''``在。''卢本韦的声音有些沙哑。

``你曾祖卢承业的手稿，还有你祖父的笔记，我都整理过一部分。改天给你看。''卢本韦的眼眶红了。

``\ldots\ldots 谢谢。''Phyrion没有再说话，带着白色母猪走出了修炼堂。

石门缓缓关闭。果糕小队众人坐在原地，久久没有动弹。

\chapter{第七十九章：引入多个2-投影序数}\label{ux7b2cux4e03ux5341ux4e5dux7ae0ux5f15ux5165ux591aux4e2a2-ux6295ux5f71ux5e8fux6570}

清晨，皇家修炼堂。

魂力晶石的光芒从穹顶洒落，将地下宫殿照得通亮。果糕小队众人围坐在修炼堂中央的石台旁，面前摊开着各自的笔记。昨天的课让他们熬了大半个通宵，但此刻每个人的眼中都闪烁着兴奋的光芒。

Phyrion站在石台正前方，粉红色的瞳孔扫过众人。白色母猪趴在他的脚边，粉红色的圆眼睛半眯着，像是在打盹，但果糕小队所有人都知道，她比任何人都清醒。

``昨天的课，消化得怎么样？''Phyrion问。

唐三率先开口：``投影对应关系的基本框架我理解了。\allowbreak{}\NovelInline{I04035}{\psi_a}\allowbreak{} 里面的结构和 \allowbreak{}\NovelInline{I04036}{\psi_1}\allowbreak{} 里面的结构一一对应，区别只在于尺度。但有几个细节我还没完全想通。''

``继续扽的话，还有：''

\NovelDisplay{M00781R01}{%
\psi_a(\psi_{\Omega_{a+1}}(\Omega_{a+1}^2)) = \psi_{\Omega_2}(\Omega_3^2)%
}

``也就是说，在 \allowbreak{}\NovelInline{I04037}{\psi_{\Omega_2}(\Omega_3)}\allowbreak{} 之后，\allowbreak{}\NovelInline{I04038}{\psi_{\Omega_{a+1}}}\allowbreak{} 里面的东西总能和 \allowbreak{}\NovelInline{I04039}{\psi_{\Omega_2}}\allowbreak{} 里面的东西对应。区别仅仅是 \allowbreak{}\NovelInline{I04040}{\Omega_{a+1}}\allowbreak{} 和 \allowbreak{}\NovelInline{I04041}{\Omega_3}\allowbreak{}。''

``这就是投影对应关系的第一条核心规律。''

``然后我们进入更高的层级。''

\NovelDisplay{M00782R01}{%
\psi_a(\psi_{\Omega_{a+1}}(\psi_{\Omega_{a+2}}(0))) = \psi_{\Omega_2}(\psi_{\Omega_4}(0))%
}

``也好理解吧。取到指数塔了现在。左边引入了 \allowbreak{}\NovelInline{I04042}{\Omega_{a+2}}\allowbreak{}------也就是 \allowbreak{}\NovelInline{I04043}{a}\allowbreak{} 之后的第二个 \allowbreak{}\NovelInline{I04044}{\Omega}\allowbreak{} 序数。右边对应的是 \allowbreak{}\NovelInline{I04045}{\psi_{\Omega_4}}\allowbreak{}。''

\NovelDisplay{M00783R01}{%
\psi_a(\psi_{\Omega_{a+1}}(\Omega_{a+2}(0))) = \psi_{\Omega_2}(\Omega_4)%
}

``像这样：\allowbreak{}\NovelInline{I04046}{\Omega_{a+1}}\allowbreak{} 对应 \allowbreak{}\NovelInline{I04047}{\Omega_3}\allowbreak{}，\allowbreak{}\NovelInline{I04048}{\Omega_{a+2}}\allowbreak{} 对应 \allowbreak{}\NovelInline{I04049}{\Omega_4}\allowbreak{}。''

``这种对应一直延续------''

\NovelDisplay{M00784R01}{%
\psi_a(\psi_{\Omega_{a+1}}(\Omega_{a+\omega})) = \psi_{\Omega_2}(\Omega_\omega)%
}

\NovelDisplay{M00785R01}{%
\psi_a(\psi_{\Omega_{a+1}}(\Omega_{a+a})) = \psi_{\Omega_2}(\Omega_{\Omega_2})%
}

Phyrion在这里特别停顿了一下，强调道：

\begin{quote}
\textbf{注意：} \allowbreak{}\NovelInline{I04050}{a}\allowbreak{} 对应的是 \allowbreak{}\NovelInline{I04051}{\Omega_2}\allowbreak{}，因为右边的最外层是 \allowbreak{}\NovelInline{I04052}{\psi_{\Omega_2}}\allowbreak{}。

\allowbreak{}\NovelInline{I04053}{\psi_{\Omega_2}(\Omega_\Omega)}\allowbreak{} 是 \allowbreak{}\NovelInline{I04054}{\psi_a(\psi_{\Omega_{a+1}}(\Omega_{a+\psi_a(0)}))}\allowbreak{} 喵。不要忘了 \allowbreak{}\NovelInline{I04055}{\psi_a(0) = \Omega}\allowbreak{} 这第一条规则。
\end{quote}

``这里有一个容易搞混的地方。\allowbreak{}\NovelInline{I04056}{a}\allowbreak{} 对应的是 \allowbreak{}\NovelInline{I04057}{\Omega_2}\allowbreak{} 而不是 \allowbreak{}\NovelInline{I04058}{\Omega}\allowbreak{}。为什么？因为右边的最外层是 \allowbreak{}\NovelInline{I04059}{\psi_{\Omega_2}}\allowbreak{}，而 \allowbreak{}\NovelInline{I04060}{\psi_{\Omega_2}}\allowbreak{} 里面的 \allowbreak{}\NovelInline{I04061}{\Omega}\allowbreak{} 对应的是 \allowbreak{}\NovelInline{I04062}{\psi_a}\allowbreak{} 里面的 \allowbreak{}\NovelInline{I04063}{0}\allowbreak{}------\allowbreak{}\NovelInline{I04064}{\psi_a(0) = \Omega}\allowbreak{}。所以 \allowbreak{}\NovelInline{I04065}{a}\allowbreak{} 要比 \allowbreak{}\NovelInline{I04066}{\Omega}\allowbreak{} 高一级，对应 \allowbreak{}\NovelInline{I04067}{\Omega_2}\allowbreak{}。''

``而 \allowbreak{}\NovelInline{I04068}{\Omega_\Omega}\allowbreak{} 在右边出现的时候，左边对应的是 \allowbreak{}\NovelInline{I04069}{\Omega_{a+\psi_a(0)}}\allowbreak{}------注意不是 \allowbreak{}\NovelInline{I04070}{\Omega_{a+\Omega}}\allowbreak{}，而是 \allowbreak{}\NovelInline{I04071}{\Omega_{a+\psi_a(0)}}\allowbreak{}。因为 \allowbreak{}\NovelInline{I04072}{\psi_a(0) = \Omega}\allowbreak{}，所以这实际上就是 \allowbreak{}\NovelInline{I04073}{\Omega_{a+\Omega}}\allowbreak{}，但写法上必须用 \allowbreak{}\NovelInline{I04074}{\psi_a(0)}\allowbreak{} 来表达，以保持投影体系的自洽性。''

唐三深深地吸了一口气。他感觉到了------这套体系的精妙程度远超他的想象。每一个细节都严丝合缝，每一个对应都不是巧合，而是规则本身的必然结果。

``总之，''Phyrion继续道，``我们取 \allowbreak{}\NovelInline{I04075}{\Omega}\allowbreak{} 下标的不动点，就得到了：''

\NovelDisplay{M00786R01}{%
\psi_a(\psi_{\Omega_{a+1}}(\psi_{I_{a+1}}(0))) = \psi_{\Omega_2}(\psi_I(0))%
}

``现在 \allowbreak{}\NovelInline{I04076}{I_{a+1}}\allowbreak{} 对应右边的 \allowbreak{}\NovelInline{I04077}{I}\allowbreak{}。''

\NovelDisplay{M00787R01}{%
\psi_a(\psi_{\Omega_{a+1}}(I_{a+1})) = \psi_{\Omega_2}(I)%
}

\NovelDisplay{M00788R01}{%
\psi_a(\psi_{\Omega_{a+1}}(M_{a+\omega})) = \psi_{\Omega_2}(M_\omega)%
}

\NovelDisplay{M00789R01}{%
\psi_a(\psi_{\Omega_{a+1}}(K_{a+\omega})) = \psi_{\Omega_2}(K_\omega)%
}

``始终一一对应，因为它们规则是一样的。''

Phyrion收回了所有魂力，地下空间重新陷入了安静。只有魂力晶石的柔和光芒照亮着每一个人的面孔。

``\textbf{这个其实就是'投影'------把 \allowbreak{}\NovelInline{I04078}{\psi_{\Omega_2}}\allowbreak{} 里面的东西投影到了 \allowbreak{}\NovelInline{I04079}{\psi_{\Omega_{a+1}}}\allowbreak{} 里面。}''

``左边的世界和右边的世界，结构完全相同。区别只在于尺度------左边用的是 \allowbreak{}\NovelInline{I04080}{a}\allowbreak{} 和 \allowbreak{}\NovelInline{I04081}{\Omega_{a+1}}\allowbreak{} 这些2-投影序数级别的巨大存在，右边用的是 \allowbreak{}\NovelInline{I04082}{1}\allowbreak{} 和 \allowbreak{}\NovelInline{I04083}{\Omega_2}\allowbreak{}、\allowbreak{}\NovelInline{I04084}{\Omega_3}\allowbreak{} 这些普通的大序数。但它们之间的对应关系是完美的、无损的、一一对应的。''

``这就是投影序数的本质。''

``\allowbreak{}\NovelInline{I04085}{\psi_a(\psi_{\Omega_{a+1}}(\Omega_{a+1}) \times 2)}\allowbreak{} 到 \allowbreak{}\NovelInline{I04086}{\psi_1(\Omega_3 + \psi_{\Omega_3}(\Omega_3))}\allowbreak{} 的推导。你说'细扽一下就很简单'，但我扽了一晚上。''

Phyrion笑了笑：``那说明你扽对了方向。今天我们先解决你昨天的疑问，然后进入新的内容。''

他抬起手，粉红色的魂力在虚空中凝聚，开始书写公式。

``在继续讲 \allowbreak{}\NovelInline{I04087}{\psi_a}\allowbreak{} 的展开之前，我们需要先解决一个基础问题。''

Phyrion的声音变得严肃起来。

``昨天我们讲到，\allowbreak{}\NovelInline{I04088}{\psi_a}\allowbreak{} 内部可以放 \allowbreak{}\NovelInline{I04089}{\Omega_{a+1}}\allowbreak{}，也可以放 \allowbreak{}\NovelInline{I04090}{\Omega_{a+2}}\allowbreak{}、\allowbreak{}\NovelInline{I04091}{I_{a+1}}\allowbreak{}、\allowbreak{}\NovelInline{I04092}{M_{a+1}}\allowbreak{}\ldots\ldots 但这样岂不是意味着我们又要引入反射的规则了吗？''

蔡须坤皱了皱眉：``昨天你提到了这个问题，但没有展开。''

``因为现在才是展开的时候。''Phyrion说，``为了用投影本身------而非反射的规则------来表示它们，我们引入 \allowbreak{}\NovelInline{I04093}{a_2}\allowbreak{}。''

他在虚空中写下了一个符号：

\NovelDisplay{M00790R01}{%
a_2%
}

``\textbf{\allowbreak{}\NovelInline{I04094}{a_2}\allowbreak{}}，即第二个 2-投影序数。''

``你们还记得 \allowbreak{}\NovelInline{I04095}{a}\allowbreak{} 是什么吗？\allowbreak{}\NovelInline{I04096}{a}\allowbreak{} 是最小的 2-投影序数。而 \allowbreak{}\NovelInline{I04097}{a_2}\allowbreak{} 是第二个。它们之间的关系，就像 \allowbreak{}\NovelInline{I04098}{\Omega}\allowbreak{} 和 \allowbreak{}\NovelInline{I04099}{\Omega_2}\allowbreak{} 之间的关系一样------\allowbreak{}\NovelInline{I04100}{a_2}\allowbreak{} 是 \allowbreak{}\NovelInline{I04101}{a}\allowbreak{} 之上的下一个同级别的存在。''

``与 \allowbreak{}\NovelInline{I04102}{\psi_a}\allowbreak{} 类似的，我们有 \allowbreak{}\NovelInline{I04103}{\psi_{a_2}}\allowbreak{}：''

Phyrion的魂力在虚空中书写出三条核心规则：

\NovelDisplay{M00791R01}{%
\psi_{a_2}(0) = \Omega_{a+1}%
}

\NovelDisplay{M00792R01}{%
\psi_{a_2}(X+1) = \psi_{a_2}(X) \cdot \omega%
}

\NovelDisplay{M00793R01}{%
\psi_{a_2}(X \sim a_2) = \psi_{a_2}(X \sim \psi_{a_2}(X \sim \psi_{a_2}(X \sim \dots)))%
}

``看到了吗？''Phyrion指着这三条规则，``结构和 \allowbreak{}\NovelInline{I04104}{\psi_a}\allowbreak{} 完全一样。唯一的区别是起始值------\allowbreak{}\NovelInline{I04105}{\psi_a(0) = \Omega}\allowbreak{}，而 \allowbreak{}\NovelInline{I04106}{\psi_{a_2}(0) = \Omega_{a+1}}\allowbreak{}。''

``这意味着 \allowbreak{}\NovelInline{I04107}{\psi_{a_2}(\Omega_{a_2+1})}\allowbreak{} 大于：''

\NovelDisplay{M00794R01}{%
\psi_{a_2}(a_2),\quad \psi_{a_2}(\varepsilon_{a_2+1}),\quad \psi_{a_2}(\Gamma_{a_2+1}),\quad \psi_{a_2}(BO_{a_2+1}),\quad \psi_{a_2}(SHO_{a_2+1}),\quad \dots%
}

卢本韦：\allowbreak{}\NovelInline{I04108}{\psi_{a_2}(\Omega_{a_2+1})=\Omega_{a+2}}\allowbreak{}对吗

phyrion:正是

``也就是说，\allowbreak{}\NovelInline{I04109}{\psi_{a_2}}\allowbreak{} 是从 \allowbreak{}\NovelInline{I04110}{\Omega_{a+1}}\allowbreak{} 开始往上构建的。\allowbreak{}\NovelInline{I04111}{\psi_a}\allowbreak{} 是从 \allowbreak{}\NovelInline{I04112}{\Omega}\allowbreak{} 开始往上构建的。两者平行，但层级不同。''

``同样地，如果 \allowbreak{}\NovelInline{I04113}{Y}\allowbreak{} 是 1-投影序数，则 \allowbreak{}\NovelInline{I04114}{\psi_{a_2}(Y)}\allowbreak{} 大于所有 \allowbreak{}\NovelInline{I04115}{\psi_{a_2}(Z)}\allowbreak{}，其中 \allowbreak{}\NovelInline{I04116}{Z < Y}\allowbreak{}。这和 \allowbreak{}\NovelInline{I04117}{\psi_a}\allowbreak{} 的取并规则完全一致。''

霍雨浩举手提问：``那 \allowbreak{}\NovelInline{I04118}{a_2}\allowbreak{} 能做什么 \allowbreak{}\NovelInline{I04119}{\psi_a}\allowbreak{} 做不到的事？''

``好问题。''Phyrion说，``\allowbreak{}\NovelInline{I04120}{\psi_a}\allowbreak{} 能表示的东西是有限的。当我们需要表示 \allowbreak{}\NovelInline{I04121}{\Omega_{a_2+2}}\allowbreak{}、\allowbreak{}\NovelInline{I04122}{I_{a_2+1}}\allowbreak{} 等等这些 1-投影序数的时候，\allowbreak{}\NovelInline{I04123}{\psi_a}\allowbreak{} 就不够用了。所以我们需要引入 \allowbreak{}\NovelInline{I04124}{\psi_{a_3}}\allowbreak{}。''

他继续在虚空中书写：

\NovelDisplay{M00795R01}{%
\psi_{a_3}(0) = \Omega_{a_2+1}%
}

\NovelDisplay{M00796R01}{%
\psi_{a_3}(X+1) = \psi_{a_3}(X) \cdot \omega%
}

\NovelDisplay{M00797R01}{%
\psi_{a_3}(X \sim a_2) = \psi_{a_3}(X \sim \psi_{a_3}(X \sim \psi_{a_3}(X \sim \dots)))%
}

``注意这里的第三条规则。''Phyrion特别强调，``里面的折叠点是 \allowbreak{}\NovelInline{I04125}{a_2}\allowbreak{} 而不是 \allowbreak{}\NovelInline{I04126}{a_3}\allowbreak{}。这和 \allowbreak{}\NovelInline{I04127}{\psi_{a_2}}\allowbreak{} 的第三条规则里折叠点是 \allowbreak{}\NovelInline{I04128}{a_2}\allowbreak{} 是一致的------因为 \allowbreak{}\NovelInline{I04129}{a_2}\allowbreak{} 是 \allowbreak{}\NovelInline{I04130}{a_3}\allowbreak{} 之下的那个 2-投影序数，就像 \allowbreak{}\NovelInline{I04131}{\psi_a}\allowbreak{} 的折叠点是 \allowbreak{}\NovelInline{I04132}{a}\allowbreak{} 本身一样。''

``等一下，''唐三忽然开口，``\allowbreak{}\NovelInline{I04133}{\psi_{a_2}}\allowbreak{} 的第三条规则里，折叠点也是 \allowbreak{}\NovelInline{I04134}{a_2}\allowbreak{}。但 \allowbreak{}\NovelInline{I04135}{\psi_{a_3}}\allowbreak{} 的第三条规则里，折叠点写的是 \allowbreak{}\NovelInline{I04136}{a_2}\allowbreak{} 而不是 \allowbreak{}\NovelInline{I04137}{a_3}\allowbreak{}。这是为什么？''

Phyrion看了唐三一眼，目光中闪过一丝赞许。

``因为你问到了关键点。''他说，``\allowbreak{}\NovelInline{I04138}{\psi_{a_3}}\allowbreak{} 的起始值是 \allowbreak{}\NovelInline{I04139}{\Omega_{a_2+1}}\allowbreak{}，它是在 \allowbreak{}\NovelInline{I04140}{a_2}\allowbreak{} 的基础上构建的。所以它的折叠点自然是 \allowbreak{}\NovelInline{I04141}{a_2}\allowbreak{}------就像 \allowbreak{}\NovelInline{I04142}{\psi_{a_2}}\allowbreak{} 的起始值是 \allowbreak{}\NovelInline{I04143}{\Omega_{a+1}}\allowbreak{}，它的折叠点是 \allowbreak{}\NovelInline{I04144}{a}\allowbreak{} 一样。每一层 \allowbreak{}\NovelInline{I04145}{\psi}\allowbreak{} 函数的折叠点，都是它之下那一层的 2-投影序数。''

``如果 \allowbreak{}\NovelInline{I04146}{Y}\allowbreak{} 是 1-投影序数，则 \allowbreak{}\NovelInline{I04147}{\psi_{a_3}(Y)}\allowbreak{} 大于所有 \allowbreak{}\NovelInline{I04148}{\psi_{a_3}(Z)}\allowbreak{}，其中 \allowbreak{}\NovelInline{I04149}{Z < Y}\allowbreak{}。规则和之前完全一致。''

Phyrion收回了手，目光扫过众人。

``\textbf{以此类推，我们便可以引入任意的 \allowbreak{}\NovelInline{I04150}{a_n}\allowbreak{}。}''

``每一个 \allowbreak{}\NovelInline{I04151}{a_n}\allowbreak{} 都是一个 2-投影序数。每一个 \allowbreak{}\NovelInline{I04152}{\psi_{a_n}}\allowbreak{} 都从 \allowbreak{}\NovelInline{I04153}{\Omega_{a_{n-1}+1}}\allowbreak{} 开始构建。每一个 \allowbreak{}\NovelInline{I04154}{\psi_{a_n}}\allowbreak{} 的折叠点都是 \allowbreak{}\NovelInline{I04155}{a_{n-1}}\allowbreak{}。整个体系像一座无限延伸的阶梯，每一层都建立在前一层的基础之上。''

奥斯卡小声嘟囔：``这不就是套娃吗\ldots\ldots{}''

``你可以这么理解。''Phyrion居然没有反驳，``但这不是一个简单的套娃。每一层套娃里面的东西，都比外面那一层大得多。就像 \allowbreak{}\NovelInline{I04156}{\psi_a}\allowbreak{} 里面的 \allowbreak{}\NovelInline{I04157}{a}\allowbreak{} 对应 \allowbreak{}\NovelInline{I04158}{\psi_1}\allowbreak{} 里面的 \allowbreak{}\NovelInline{I04159}{\Omega_2}\allowbreak{}------一个巨大的的 \allowbreak{}\NovelInline{I04160}{a}\allowbreak{}，在投影之后变成了一个小小的 \allowbreak{}\NovelInline{I04161}{\Omega_2}\allowbreak{}。''

白色母猪在这一刻睁开了粉红色的圆眼睛，缓缓扫了奥斯卡一眼。

奥斯卡打了个寒颤：``我\ldots\ldots 我就是打个比方。''

``好了，\allowbreak{}\NovelInline{I04162}{a_2}\allowbreak{} 和 \allowbreak{}\NovelInline{I04163}{a_3}\allowbreak{} 的框架搭好了。现在我们回到 \allowbreak{}\NovelInline{I04164}{\psi_a}\allowbreak{} 本身，解决一个昨天遗留的关键问题。''

Phyrion的声音变得格外认真。

``如何才能得到 \allowbreak{}\NovelInline{I04165}{\psi_a(\Omega_{a+1})}\allowbreak{}？''

KNK------不对，是果糕小队的众人面面相觑。这个问题他们确实没有想通。

Phyrion没有直接给出答案。他说：

``在这里我们需要做一个假设------''

他停顿了一下，目光变得深邃。

``\textbf{假设我们可以让 \allowbreak{}\NovelInline{I04166}{\psi_{\Omega_{a+1}}}\allowbreak{} 里面的东西任意大，并且这样得到的式子都有意义。}''

``里面的东西越大，就代表越高级的递归操作。而'递归操作'的极限------你们应该不陌生吧？''

蔡须坤的瞳孔微微收缩：``容许序数。''

``没错。''Phyrion点了点头，``容许序数就是递归操作的极限。它是所有递归操作都无法触及的边界。''

``左边是 \allowbreak{}\NovelInline{I04167}{\psi_a(\Omega_{a+1})}\allowbreak{} 的时候，代表我们把左边的 \allowbreak{}\NovelInline{I04168}{\psi_{\Omega_{a+1}}}\allowbreak{} 里面的东西取到了极限。相对应的，右边的 \allowbreak{}\NovelInline{I04169}{\psi_1}\allowbreak{} 里面的东西也取到了极限。''

他的魂力在虚空中凝聚出一个简洁而震撼的等式：

\NovelDisplay{M00798R01}{%
\boxed{\psi_a(\Omega_{a+1}) = \Omega_2}%
}

``一个容许序数。''

整个修炼堂陷入了短暂的沉默。

霍雨浩率先打破了沉默：``等一下。\allowbreak{}\NovelInline{I04170}{\psi_1}\allowbreak{} 始终输出 \allowbreak{}\NovelInline{I04171}{\Omega_2}\allowbreak{} 以下的序数。那么我们把'\allowbreak{}\NovelInline{I04172}{\psi_1}\allowbreak{} 能输出的东西'扔到一个集合里面去取极限，得到的自然就是 \allowbreak{}\NovelInline{I04173}{\Omega_2}\allowbreak{} 了。''

``对。''Phyrion说，``实际上，我们是在假设：对于任意一个 \allowbreak{}\NovelInline{I04174}{\psi_1(0)}\allowbreak{} 到 \allowbreak{}\NovelInline{I04175}{\Omega_2}\allowbreak{} 之间的序数，我们总能找到一个 \allowbreak{}\NovelInline{I04176}{\psi_{\Omega_{a+1}}(0)}\allowbreak{} 到 \allowbreak{}\NovelInline{I04177}{\Omega_{a+1}}\allowbreak{} 之间的序数与之对应，并且这种对应关系保序。''

他的目光变得异常锐利。

``\textbf{这种对应关系，其实就是投影。}''

唐三深吸了一口气。他终于明白了------投影不仅仅是一种计算工具，它本质上是一种保序的双射。\allowbreak{}\NovelInline{I04178}{\psi_a}\allowbreak{} 之所以能把 \allowbreak{}\NovelInline{I04179}{\psi_{\Omega_{a+1}}}\allowbreak{} 里面的巨大结构压缩成 \allowbreak{}\NovelInline{I04180}{\psi_1}\allowbreak{} 里面的较小结构，是因为两者之间存在完美的对应关系。这种对应关系不是人为构造的，而是规则本身的必然结果。

``所以 \allowbreak{}\NovelInline{I04181}{\psi_a(\Omega_{a+1}) = \Omega_2}\allowbreak{}，不是因为我们'算'出了这个结果，而是因为我们'假设'了投影的完备性，然后取极限得到的。''唐三缓缓说道。

``可以这么理解。''Phyrion说，``但'假设'这个词不够准确。更准确地说，是投影体系本身的自洽性要求这个等式成立。如果 \allowbreak{}\NovelInline{I04182}{\psi_a(\Omega_{a+1})}\allowbreak{} 不等于 \allowbreak{}\NovelInline{I04183}{\Omega_2}\allowbreak{}，那么整个投影体系就会出现矛盾。''

``现在我们承认 \allowbreak{}\NovelInline{I04184}{\psi_a(\Omega_{a+1}) = \Omega_2}\allowbreak{}，并让我们继续前进。''

Phyrion的魂力再次在虚空中书写：

\NovelDisplay{M00799R01}{%
\psi_a(\Omega_{a+1} + 1) = \psi_{\Omega_3}(1)%
}

``还是*ω关系。和前面一样。''

蔡须坤点了点头：``和前面 \allowbreak{}\NovelInline{I04185}{\psi_a(1) = \psi_1(1)}\allowbreak{} 的模式一样。\allowbreak{}\NovelInline{I04186}{\Omega_{a+1}}\allowbreak{} 对应 \allowbreak{}\NovelInline{I04187}{\Omega_2}\allowbreak{}，所以 \allowbreak{}\NovelInline{I04188}{\Omega_{a+1} + 1}\allowbreak{} 对应 \allowbreak{}\NovelInline{I04189}{\Omega_2 + 1}\allowbreak{}，而 \allowbreak{}\NovelInline{I04190}{\psi_1(\Omega_2 + 1) = \psi_{\Omega_3}(1)}\allowbreak{}。''

``对。''Phyrion继续书写，``然后------''

\NovelDisplay{M00800R01}{%
\psi_a(\Omega_{a+1} + \psi_a(\Omega_{a+1} + 1)) = \psi_{\Omega_3}(\psi_{\Omega_3}(1))%
}

\NovelDisplay{M00801R01}{%
\psi_a(\Omega_{a+1} + a) = \psi_{\Omega_3}(\Omega_3)%
}

Phyrion在这里停了一下，特别指出：

``我们发现现在的 \allowbreak{}\NovelInline{I04191}{a}\allowbreak{} 对应到 \allowbreak{}\NovelInline{I04192}{\psi_{\Omega_3}}\allowbreak{} 里面的 \allowbreak{}\NovelInline{I04193}{\Omega_3}\allowbreak{} 了。为什么？因为现在我们可以用比 \allowbreak{}\NovelInline{I04194}{a}\allowbreak{} 小的 \allowbreak{}\NovelInline{I04195}{\psi_a(\Omega_{a+1})}\allowbreak{} 直接把 \allowbreak{}\NovelInline{I04196}{\Omega_2}\allowbreak{} 表示出来。这是前面的 \allowbreak{}\NovelInline{I04197}{\psi_1}\allowbreak{} 阶段我们做不到的。''

``在 \allowbreak{}\NovelInline{I04198}{\psi_1}\allowbreak{} 阶段，\allowbreak{}\NovelInline{I04199}{\Omega_2}\allowbreak{} 是 \allowbreak{}\NovelInline{I04200}{\psi_1}\allowbreak{} 无法触及的上界。但在 \allowbreak{}\NovelInline{I04201}{\psi_a}\allowbreak{} 阶段，\allowbreak{}\NovelInline{I04202}{\psi_a(\Omega_{a+1}) = \Omega_2}\allowbreak{} 已经把 \allowbreak{}\NovelInline{I04203}{\Omega_2}\allowbreak{} 纳入了可表示的范围。所以 \allowbreak{}\NovelInline{I04204}{a}\allowbreak{} 在 \allowbreak{}\NovelInline{I04205}{\psi_a}\allowbreak{} 中的角色，就自动升级成了 \allowbreak{}\NovelInline{I04206}{\Omega_3}\allowbreak{} 在 \allowbreak{}\NovelInline{I04207}{\psi_{\Omega_3}}\allowbreak{} 中的角色。''

``自然地，\allowbreak{}\NovelInline{I04208}{a}\allowbreak{} 的指数塔对应 \allowbreak{}\NovelInline{I04209}{\Omega_3}\allowbreak{} 的指数塔。''

\NovelDisplay{M00802R01}{%
\psi_a(\Omega_{a+1} + \psi_{\Omega_{a+1}}(0)) = \psi_{\Omega_3}(\psi_{\Omega_4}(0))%
}

\NovelDisplay{M00803R01}{%
\psi_a(\Omega_{a+1} + \psi_{\Omega_{a+1}}(\Omega_{a+1})) = \psi_{\Omega_3}(\Omega_4)%
}

``现在是 \allowbreak{}\NovelInline{I04210}{a}\allowbreak{} 对应 \allowbreak{}\NovelInline{I04211}{\Omega_3}\allowbreak{}，\allowbreak{}\NovelInline{I04212}{\Omega_{a+1}}\allowbreak{} 对应 \allowbreak{}\NovelInline{I04213}{\Omega_4}\allowbreak{}，\ldots\ldots{}''

Phyrion的语速加快了一些：

\NovelDisplay{M00804R01}{%
\psi_a(\Omega_{a+1} + \psi_{\Omega_{a+1}}(\Omega_{a+\omega})) = \psi_{\Omega_3}(\Omega_\omega)%
}

\NovelDisplay{M00805R01}{%
\psi_a(\Omega_{a+1} + \psi_{\Omega_{a+1}}(I_{a+\omega})) = \psi_{\Omega_3}(I_\omega)%
}

\NovelDisplay{M00806R01}{%
\psi_a(\Omega_{a+1} + \psi_{\Omega_{a+1}}(K_{a+\omega})) = \psi_{\Omega_3}(K_\omega)%
}

``又是一样的关系。''Phyrion说，``我们再和前面一样假设 \allowbreak{}\NovelInline{I04214}{\psi_{\Omega_{a+1}}}\allowbreak{} 里面的东西可以任意大，再取一次极限，就是------''

\NovelDisplay{M00807R01}{%
\psi_a(\Omega_{a+1} \times 2) = \Omega_3%
}

蔡须坤的眼睛亮了：``\allowbreak{}\NovelInline{I04215}{\psi_a(\Omega_{a+1} \times \omega)}\allowbreak{} 是 \allowbreak{}\NovelInline{I04216}{\Omega_\omega}\allowbreak{} 吧？''

``嗯呐。''Phyrion微微一笑，``后面就完全一样了。''

\NovelDisplay{M00808R01}{%
\psi_a(\Omega_{a+1} \times \omega) = \Omega_\omega%
}

``你看，''Phyrion收回了手，``整个体系的规律是非常优美的。\allowbreak{}\NovelInline{I04217}{\psi_a(\Omega_{a+1})}\allowbreak{} 是一个分水岭------在它之前，\allowbreak{}\NovelInline{I04218}{\psi_a}\allowbreak{} 里面的结构对应的是 \allowbreak{}\NovelInline{I04219}{\psi_1}\allowbreak{} 里面的结构；在它之后，\allowbreak{}\NovelInline{I04220}{\psi_a}\allowbreak{} 里面的结构对应的是 \allowbreak{}\NovelInline{I04221}{\psi_{\Omega_3}}\allowbreak{} 里面的结构。''

``而 \allowbreak{}\NovelInline{I04222}{\psi_a(\Omega_{a+1} \times 2) = \Omega_3}\allowbreak{} 又是一个分水岭------在它之后，\allowbreak{}\NovelInline{I04223}{a}\allowbreak{} 对应的不再是 \allowbreak{}\NovelInline{I04224}{\Omega_2}\allowbreak{}，而是 \allowbreak{}\NovelInline{I04225}{\Omega_3}\allowbreak{}。整个投影的层级又往上跳了一级。

``这就是投影序数的自举能力。每到一个新的容许序数，整个体系就会自动升级，开始对应更高层级的结构。''

\chapter{第八十章：从OFP到I(1,0,0)}\label{ux7b2cux516bux5341ux7ae0ux4eceofpux5230i100}

清晨的皇家修炼堂，灵气如潮水般在穹顶下缓缓流动。魂力晶石散发出淡金色的柔光，将石台照得通亮。

果糕小队众人已经围坐在石台旁，面前摊开着各自的笔记。经过前两天的学习，他们对投影对应关系已经有了基本的认知，但每个人的眼底都带着一种隐隐的紧张------他们知道，今天的内容会比之前更难。

Phyrion站在石台正前方，粉红色的瞳孔扫过众人。白色母猪趴在他的脚边，粉红色的圆眼睛半眯着，看上去漫不经心，但果糕小队所有人都知道，她比任何人都清醒。

``前两节课，我们讲了 \allowbreak{}\NovelInline{I04226}{\psi_a(\Omega_{a+1})}\allowbreak{} 到 \allowbreak{}\NovelInline{I04227}{\psi_a(\Omega_{a+1} \times \omega)}\allowbreak{} 的展开。''Phyrion的声音平静而清晰，``今天我们继续。从 \allowbreak{}\NovelInline{I04228}{\psi_a(\Omega_{a+1} \times a)}\allowbreak{} 开始。''

他抬起手，粉红色的魂力在虚空中凝聚，开始书写公式。

``那么 \allowbreak{}\NovelInline{I04229}{\psi_a(\Omega_{a+1} \times a)}\allowbreak{} 呢？''

Phyrion的目光变得格外认真。

``在写这个式子之前，有一个重要的注意事项。''

他在虚空中写下一行醒目的字：

\begin{quote}
\textbf{注意：} \allowbreak{}\NovelInline{I04230}{\psi_a(\Omega_{a+1} \times \psi_a(0))}\allowbreak{} 已经是 \allowbreak{}\NovelInline{I04231}{\Omega_\Omega}\allowbreak{} 了。
\end{quote}

``还记得第一条规则吗？\allowbreak{}\NovelInline{I04232}{\psi_a(0) = \Omega}\allowbreak{}。所以 \allowbreak{}\NovelInline{I04233}{\psi_a(\Omega_{a+1} \times \psi_a(0))}\allowbreak{} 就是 \allowbreak{}\NovelInline{I04234}{\psi_a(\Omega_{a+1} \times \Omega)}\allowbreak{}，它等于 \allowbreak{}\NovelInline{I04235}{\Omega_\Omega}\allowbreak{}。''

``这意味着什么？意味着在 \allowbreak{}\NovelInline{I04236}{\Omega_{a+1} \times a}\allowbreak{} 到达之前，\allowbreak{}\NovelInline{I04237}{\Omega_{a+1} \times \Omega}\allowbreak{} 就已经把 \allowbreak{}\NovelInline{I04238}{\Omega}\allowbreak{} 的下标推到了 \allowbreak{}\NovelInline{I04239}{\Omega}\allowbreak{} 本身。''

唐三皱起了眉头：``那 \allowbreak{}\NovelInline{I04240}{\psi_a(\Omega_{a+1} \times a)}\allowbreak{} 会是什么？''

Phyrion没有直接回答。他看着众人，等待着。

蔡须坤闭上了眼睛，金色双瞳在眼睑下微微转动。她在感知------用她的Y序列感知力去触碰那个答案的边缘。

片刻后，她睁开眼睛，试探性地说：``\allowbreak{}\NovelInline{I04241}{\psi_I(0)}\allowbreak{}？''

Phyrion的嘴角微微上扬。

``对的。是 \textbf{OFP}。''

他在虚空中写下了那个等式：

\NovelDisplay{M00809R01}{%
\psi_a(\Omega_{a+1} \times a) = \psi_I(I)%
}

``一出现 \allowbreak{}\NovelInline{I04242}{a}\allowbreak{}，就是不动点。''

奥斯卡挠了挠头：``等等，为什么一出现 \allowbreak{}\NovelInline{I04243}{a}\allowbreak{} 就不动点了？''

``因为 \allowbreak{}\NovelInline{I04244}{a}\allowbreak{} 在这个体系里扮演的角色，就是'折叠点'。''Phyrion解释道，``\allowbreak{}\NovelInline{I04245}{\psi_a}\allowbreak{} 里面的 \allowbreak{}\NovelInline{I04246}{a}\allowbreak{}，对应的是投影另一侧的不动点操作。当 \allowbreak{}\NovelInline{I04247}{\Omega_{a+1}}\allowbreak{} 的系数从有限的数字变成 \allowbreak{}\NovelInline{I04248}{a}\allowbreak{} 的时候，就相当于在投影的另一侧执行了一次'取不动点'的操作。''

``而 \allowbreak{}\NovelInline{I04249}{\psi_I(I)}\allowbreak{}，就是 \allowbreak{}\NovelInline{I04250}{I}\allowbreak{} 函数的第一个值------第一个不可达序数。也就是 OFP。''

``你们要记住这个规律：\textbf{在 \allowbreak{}\NovelInline{I04251}{\psi_a}\allowbreak{} 里面，\allowbreak{}\NovelInline{I04252}{a}\allowbreak{} 的出现标志着不动点的诞生。}''

Phyrion的目光变得深邃起来。

``在继续展开之前，我要先讲一个关键的认知。这个认知，决定了你们能不能真正理解投影序数的力量。''

他转过身，面对着石台正中央的那块巨大的魂力黑板，用指尖缓缓写下两个字：

\textbf{\allowbreak{}\NovelInline{I04253}{a}\allowbreak{} 与 \allowbreak{}\NovelInline{I04254}{I}\allowbreak{}}

``你们刚才看到了------\allowbreak{}\NovelInline{I04255}{\psi_a(\Omega_{a+1} \times a) = \psi_I(I)}\allowbreak{}。\allowbreak{}\NovelInline{I04256}{a}\allowbreak{} 在投影的另一侧对应的是 \allowbreak{}\NovelInline{I04257}{I}\allowbreak{}，是不可达序数的入口。''

``但这只是 \allowbreak{}\NovelInline{I04258}{a}\allowbreak{} 的第一重身份。''

他的声音变得低沉而有力。

``在 \allowbreak{}\NovelInline{I04259}{\psi_a(\Omega_{a+1} \times a \times \text{递归运算})}\allowbreak{} 这个阶段------也就是说，\allowbreak{}\NovelInline{I04260}{a}\allowbreak{} 作为 \allowbreak{}\NovelInline{I04261}{\Omega_{a+1}}\allowbreak{} 的系数的乘数出现的这个阶段------\allowbreak{}\NovelInline{I04262}{a}\allowbreak{} 代表的是 \allowbreak{}\NovelInline{I04263}{I}\allowbreak{}。\allowbreak{}\NovelInline{I04264}{a}\allowbreak{} 的每一次递归运算，都对应着 \allowbreak{}\NovelInline{I04265}{I}\allowbreak{} 的层级的跃迁。''

他在虚空中书写：

\NovelDisplay{M00810R01}{%
a \longleftrightarrow I%
}

``\allowbreak{}\NovelInline{I04266}{a}\allowbreak{} 的加法，对应 \allowbreak{}\NovelInline{I04267}{I}\allowbreak{} 的不动点迭代。\allowbreak{}\NovelInline{I04268}{a}\allowbreak{} 的乘法，对应 \allowbreak{}\NovelInline{I04269}{I}\allowbreak{} 的自变量增长。\allowbreak{}\NovelInline{I04270}{a}\allowbreak{} 的幂次，对应 \allowbreak{}\NovelInline{I04271}{I}\allowbreak{} 本身的运算。''

``但我要特别强调一件事。''

Phyrion的目光扫过众人，粉红色的瞳孔中闪烁着一种近乎敬畏的光芒。

``\allowbreak{}\NovelInline{I04272}{a}\allowbreak{} 之所以能代表 \allowbreak{}\NovelInline{I04273}{I}\allowbreak{}，不是因为我们人为地建立了这种对应关系。而是因为 \textbf{\allowbreak{}\NovelInline{I04274}{a}\allowbreak{} 本身就足够大到让任何容许序数望尘莫及}。''

``你们回想一下------\allowbreak{}\NovelInline{I04275}{\Omega}\allowbreak{} 是容许序数，\allowbreak{}\NovelInline{I04276}{\Omega_2}\allowbreak{} 是容许序数，\allowbreak{}\NovelInline{I04277}{\Omega_3}\allowbreak{} 是容许序数\ldots\ldots{}\allowbreak{}\NovelInline{I04278}{I}\allowbreak{} 是不可达序数，也是容许序数。这些容许序数一个比一个大，一个比一个强。但在 \allowbreak{}\NovelInline{I04279}{a}\allowbreak{} 面前，它们全部黯然失色。''

``\allowbreak{}\NovelInline{I04280}{a}\allowbreak{} 是 2-投影序数。它不是一个普通的容许序数------它是所有容许序数的'投影极限'。它大到什么程度？大到它内部的每一次递归运算，都能在投影的另一侧引发整个 \allowbreak{}\NovelInline{I04281}{I}\allowbreak{} 层级的跃迁。''

唐三深吸了一口气。他终于明白了------为什么 Phyrion 说''投影很强''。不是因为投影的技巧有多精妙，而是因为投影的基石------\allowbreak{}\NovelInline{I04282}{a}\allowbreak{}------本身就站在一个所有容许序数都无法企及的高度上。

``所以，''Phyrion继续说道，``当你们看到 \allowbreak{}\NovelInline{I04283}{\psi_a(\Omega_{a+1} \times a)}\allowbreak{} 的时候，不要把它简单地理解为'一个系数变成了 \allowbreak{}\NovelInline{I04284}{a}\allowbreak{}'。你要理解的是------\allowbreak{}\NovelInline{I04285}{a}\allowbreak{} 的出现，意味着整个体系从 \allowbreak{}\NovelInline{I04286}{\Omega}\allowbreak{} 的层级跃迁到了 \allowbreak{}\NovelInline{I04287}{I}\allowbreak{} 的层级。这个跃迁，是由 \allowbreak{}\NovelInline{I04288}{a}\allowbreak{} 本身的巨大性所保证的。''

``现在我们知道了 \allowbreak{}\NovelInline{I04289}{\psi_a(\Omega_{a+1} \times a) = \psi_I(I)}\allowbreak{}，我们继续用BOCF模式，我们继续。''

Phyrion的魂力在虚空中飞速书写：

\NovelDisplay{M00811R01}{%
\psi_a(\Omega_{a+1} \times a + \Omega_{a+1}) = \psi_I(I+1)%
}

``\allowbreak{}\NovelInline{I04290}{\psi_a}\allowbreak{} 里面每加一个 \allowbreak{}\NovelInline{I04291}{\Omega_{a+1}}\allowbreak{}，代表输出的 \allowbreak{}\NovelInline{I04292}{\Omega}\allowbreak{} 下标 \allowbreak{}\NovelInline{I04293}{+1}\allowbreak{}。这个规律从前面的课就已经建立起来了。''

``然后------''

\NovelDisplay{M00812R01}{%
\psi_a(\Omega_{a+1} \times a + \Omega_{a+1} \times \psi_a(0)) = \psi_I(I+\psi_I(0))%
}

``注意这里，\allowbreak{}\NovelInline{I04294}{\psi_a(0) = \Omega}\allowbreak{}，所以 \allowbreak{}\NovelInline{I04295}{\Omega_{a+1} \times \psi_a(0)}\allowbreak{} 就是 \allowbreak{}\NovelInline{I04296}{\Omega_{a+1} \times \Omega}\allowbreak{}。在投影的另一侧，\allowbreak{}\NovelInline{I04297}{\psi_I(I)}\allowbreak{}内部多了\allowbreak{}\NovelInline{I04298}{\psi_I(0)}\allowbreak{}。''

``继续------''

\NovelDisplay{M00813R01}{%
\psi_a(\Omega_{a+1} \times a + \Omega_{a+1} \times \psi_a(\Omega_{a+1} \times a)) = \psi_I(I+\psi_I(I))%
}

``这个式子有点长。但你们仔细看------左边 \allowbreak{}\NovelInline{I04299}{\psi_a}\allowbreak{} 的自变量里面出现了 \allowbreak{}\NovelInline{I04300}{\psi_a(\Omega_{a+1} \times a)}\allowbreak{}，也就是我们刚刚算出来的 \allowbreak{}\NovelInline{I04301}{\psi_I(0)}\allowbreak{}。所以在投影的另一侧，\allowbreak{}\NovelInline{I04302}{\Omega}\allowbreak{} 的下标变成了 \allowbreak{}\NovelInline{I04303}{\psi_I(0) \times 2}\allowbreak{}。''

``总之------''Phyrion加快了语速，``\allowbreak{}\NovelInline{I04304}{\psi_a(\Omega_{a+1} \times a \times 2)}\allowbreak{} 是下一个 OFP，也就是 \allowbreak{}\NovelInline{I04305}{\psi_I(I \times 2)}\allowbreak{}。''

蔡须坤微微点头：``这个我会。\allowbreak{}\NovelInline{I04306}{a \times 2}\allowbreak{} 就是第二个不动点。''

``对。''Phyrion说，``然后------''

\NovelDisplay{M00814R01}{%
\psi_a(\Omega_{a+1} \times a \times \psi_a(0)) = \psi_I(I\times \Omega)%
}

\NovelDisplay{M00815R01}{%
\psi_a(\Omega_{a+1} \times a \times \psi_a(\Omega_{a+1} \times a)) = \psi_I(I\times \psi_I(I))%
}

``\allowbreak{}\NovelInline{I04307}{a}\allowbreak{} 的系数里面出现了 \allowbreak{}\NovelInline{I04308}{\psi_a(0) = \Omega}\allowbreak{}，对应 \allowbreak{}\NovelInline{I04309}{\psi_I}\allowbreak{} 的自变量变成了 \allowbreak{}\NovelInline{I04310}{\Omega}\allowbreak{}。出现了 \allowbreak{}\NovelInline{I04311}{\psi_a(\Omega_{a+1} \times a) = \psi_I(I)}\allowbreak{}，对应 \allowbreak{}\NovelInline{I04312}{\psi_I}\allowbreak{} 的自变量变成了 \allowbreak{}\NovelInline{I04313}{\psi_I(I)}\allowbreak{}。''

``一一对应，始终如一。''

``然后------''

\NovelDisplay{M00816R01}{%
\psi_a(\Omega_{a+1} \times a^2) = \psi_I(I^2)%
}

``看到了吗？\allowbreak{}\NovelInline{I04314}{a^2}\allowbreak{} 对应 \allowbreak{}\NovelInline{I04315}{I}\allowbreak{}。\allowbreak{}\NovelInline{I04316}{a}\allowbreak{} 的自乘，在投影的另一侧变成了 \allowbreak{}\NovelInline{I04317}{I}\allowbreak{}------也就是第一个不可达序数本身。''正如\allowbreak{}\NovelInline{I04318}{ψ(Ω)[mocf]=ψ(Ω)[bocf]}\allowbreak{}

Phyrion在这里停了一下

``继续------''

\NovelDisplay{M00817R01}{%
\psi_a(\Omega_{a+1} \times a^3) = \psi_I(I^3)%
}

\NovelDisplay{M00818R01}{%
\psi_a(\Omega_{a+1} \times a^a) = \psi_I(I^I)%
}

``\allowbreak{}\NovelInline{I04319}{a}\allowbreak{} 的连乘直接对应 \allowbreak{}\NovelInline{I04320}{I}\allowbreak{} 的幂次。\allowbreak{}\NovelInline{I04321}{a}\allowbreak{} 乘了自己 \allowbreak{}\NovelInline{I04322}{a}\allowbreak{} 次，在投影的另一侧就是 \allowbreak{}\NovelInline{I04323}{I^I}\allowbreak{}。''

Phyrion的声音变得格外有力。

``\textbf{所以投影很强。}''

``你们要理解这句话的分量。\allowbreak{}\NovelInline{I04324}{a}\allowbreak{} 的每一次乘法叠加，都在投影的另一侧引发了 \allowbreak{}\NovelInline{I04325}{I}\allowbreak{} 的幂次跃迁。这不是线性的对应------这是结构性的对应。左边的一次运算，在右边对应的是整个层级的跃迁。''

``然后，和之前一样，\allowbreak{}\NovelInline{I04326}{\psi_{\Omega_{a+1}}}\allowbreak{} 会再次出现。''

Phyrion的魂力在虚空中继续书写：

\NovelDisplay{M00819R01}{%
\psi_a(\Omega_{a+1} \times \psi_{\Omega_{a+1}}(0)) = \psi_I(\psi_{\Omega_{I+1}}(0))%
}

``\allowbreak{}\NovelInline{I04327}{\psi_{\Omega_{a+1}}(0)}\allowbreak{} 对应 \allowbreak{}\NovelInline{I04328}{\psi_{\Omega_{I+1}}(0)}\allowbreak{}。注意这里的变化------左边是 \allowbreak{}\NovelInline{I04329}{\Omega_{a+1}}\allowbreak{}，右边对应的是 \allowbreak{}\NovelInline{I04330}{\Omega_{I+1}}\allowbreak{}。因为现在投影的基准已经从 \allowbreak{}\NovelInline{I04331}{\Omega_2}\allowbreak{} 升级到了 \allowbreak{}\NovelInline{I04332}{\psi_I(0) = I}\allowbreak{}，所以 \allowbreak{}\NovelInline{I04333}{\Omega}\allowbreak{} 的下标也相应地从 \allowbreak{}\NovelInline{I04334}{a+1}\allowbreak{} 对应到了 \allowbreak{}\NovelInline{I04335}{I+1}\allowbreak{}。''

``继续------''

\NovelDisplay{M00820R01}{%
\psi_a(\Omega_{a+1} \times \psi_{\Omega_{a+1}}(\Omega_{a+1})) = \psi_I(\Omega_{I+1})%
}

\NovelDisplay{M00821R01}{%
\psi_a(\Omega_{a+1} \times \psi_{\Omega_{a+1}}(I_{a+1})) = \psi_I(I_2)%
}

``\allowbreak{}\NovelInline{I04336}{I_{a+1}}\allowbreak{} 对应 \allowbreak{}\NovelInline{I04337}{I_2}\allowbreak{}。和之前 \allowbreak{}\NovelInline{I04338}{I_{a+1}}\allowbreak{} 对应 \allowbreak{}\NovelInline{I04339}{I}\allowbreak{} 的模式一样，只不过现在基准升级了。''

\NovelDisplay{M00822R01}{%
\psi_a(\Omega_{a+1} \times \psi_{\Omega_{a+1}}(I_{a+\omega})) = \psi_I(I_\omega)%
}

``又对应上了。''Phyrion说，``规律始终不变------左边用什么序数，右边就用对应的序数。区别只在于尺度。''

``好，现在进入一个新的阶段。''

Phyrion的声音微微提高了一些。

``和前面一样，我们来看看 \allowbreak{}\NovelInline{I04340}{\psi_a(\Omega_{a+1}^2)}\allowbreak{} 是什么。''

他看着众人，等待着。

蔡须坤沉思了片刻，然后缓缓说道：``\allowbreak{}\NovelInline{I04341}{I}\allowbreak{}。''

``嗯呐。''Phyrion点了点头，``浅显一点的话，就可以理解成\textbf{把下标取出来的操作}。\allowbreak{}\NovelInline{I04342}{\Omega_{a+1}}\allowbreak{} 的一次方对应 \allowbreak{}\NovelInline{I04343}{\psi_I(I)}\allowbreak{}，也就是 OFP。\allowbreak{}\NovelInline{I04344}{\Omega_{a+1}}\allowbreak{} 的平方，就是把 OFP 的'下标'取出来，得到了 \allowbreak{}\NovelInline{I04345}{I}\allowbreak{} 本身。''

``那么继续------''

\NovelDisplay{M00823R01}{%
\psi_a(\Omega_{a+1}^2 + \Omega_{a+1}) = \Omega_{I+1}%
}

``\allowbreak{}\NovelInline{I04346}{\Omega_{a+1}^2}\allowbreak{} 的基础上加一个 \allowbreak{}\NovelInline{I04347}{\Omega_{a+1}}\allowbreak{}，对应 \allowbreak{}\NovelInline{I04348}{\Omega}\allowbreak{} 的下标变成 \allowbreak{}\NovelInline{I04349}{I+1}\allowbreak{}。''

\NovelDisplay{M00824R01}{%
\psi_a(\Omega_{a+1}^2 + \Omega_{a+1} \times \psi_a(\Omega_{a+1}^2)) = \Omega_{I \times 2}%
}

``\allowbreak{}\NovelInline{I04350}{\psi_a(\Omega_{a+1}^2) = I}\allowbreak{}，所以 \allowbreak{}\NovelInline{I04351}{\Omega_{a+1} \times \psi_a(\Omega_{a+1}^2)}\allowbreak{} 就是 \allowbreak{}\NovelInline{I04352}{\Omega_{a+1} \times I}\allowbreak{}。对应 \allowbreak{}\NovelInline{I04353}{\Omega}\allowbreak{} 的下标变成 \allowbreak{}\NovelInline{I04354}{I \times 2}\allowbreak{}。''

\NovelDisplay{M00825R01}{%
\psi_a(\Omega_{a+1}^2 + \Omega_{a+1} \times a) = \psi_{I_2}(I_2)%
}

``出现了 \allowbreak{}\NovelInline{I04355}{a}\allowbreak{}，又是不动点。\allowbreak{}\NovelInline{I04356}{\psi_{I_2}(I_2)}\allowbreak{}------第二个 \allowbreak{}\NovelInline{I04357}{I}\allowbreak{} 函数的第一个值。''

\NovelDisplay{M00826R01}{%
\psi_a(\Omega_{a+1}^2 \times 2) = I_2%
}

\NovelDisplay{M00827R01}{%
\psi_a(\Omega_{a+1}^2 \times 3) = I_3%
}

``\allowbreak{}\NovelInline{I04358}{\Omega_{a+1}^2}\allowbreak{} 的系数直接对应 \allowbreak{}\NovelInline{I04359}{I}\allowbreak{} 的下标。\allowbreak{}\NovelInline{I04360}{\times 2}\allowbreak{} 就是 \allowbreak{}\NovelInline{I04361}{I_2}\allowbreak{}，\allowbreak{}\NovelInline{I04362}{\times 3}\allowbreak{} 就是 \allowbreak{}\NovelInline{I04363}{I_3}\allowbreak{}。''

\NovelDisplay{M00828R01}{%
\psi_a(\Omega_{a+1}^2 \times \psi_a(0)) = I_\Omega%
}

``\allowbreak{}\NovelInline{I04364}{\psi_a(0) = \Omega}\allowbreak{}，所以是 \allowbreak{}\NovelInline{I04365}{I_\Omega}\allowbreak{}。''

\NovelDisplay{M00829R01}{%
\psi_a(\Omega_{a+1}^2 \times \psi_a(\Omega_{a+1}^2)) = I_I%
}

``\allowbreak{}\NovelInline{I04366}{\psi_a(\Omega_{a+1}^2) = I}\allowbreak{}，所以是 \allowbreak{}\NovelInline{I04367}{I_I}\allowbreak{}。''

Phyrion在这里停了一下，然后写下了一个关键的等式：

\NovelDisplay{M00830R01}{%
\psi_a(\Omega_{a+1}^2 \times a) = \psi_{I(1,0)}(I(1,0))%
}

``看到了吗？\allowbreak{}\NovelInline{I04368}{a}\allowbreak{} 又出现了，又是不动点。但这次的不动点不是 \allowbreak{}\NovelInline{I04369}{\psi_{I_n}(I_n)}\allowbreak{} 这种简单的形式了------它是 \allowbreak{}\NovelInline{I04370}{\psi_{I(1,0)}(I(1,0))}\allowbreak{}。''

``也就是说，\allowbreak{}\NovelInline{I04371}{\Omega_{a+1}^2 \times a}\allowbreak{} 对应的不动点，已经跳到了 \allowbreak{}\NovelInline{I04372}{I(1,0)}\allowbreak{} 的层级。这就是二重逗号 Veblen 函数的入口------\allowbreak{}\NovelInline{I04373}{I(1,0)}\allowbreak{} 是二元 Veblen 函数 \allowbreak{}\NovelInline{I04374}{I-\varphi(1,0)}\allowbreak{} 在大序数尺度上的对应物。''

``投影和二重逗号 Veblen 函数有着深刻的关联。\allowbreak{}\NovelInline{I04375}{I}\allowbreak{} 本身就是一元 I-Veblen 函数的不动点枚举，而 \allowbreak{}\NovelInline{I04376}{I(1,0)}\allowbreak{} 则是二元 I-Veblen 函数的第一个值。投影体系通过 \allowbreak{}\NovelInline{I04377}{a}\allowbreak{} 的递归运算，自然而然地触及了 I-Veblen 函数的层级结构。''

奥斯卡瞪大了眼睛：``这\ldots\ldots 感觉好简单的样子。''

Phyrion微微一笑：``是不难的。我实际上还把里面的细节揉碎了讲了。像他们只讲大概的话，感觉会更简单。''

``然后------''

\NovelDisplay{M00831R01}{%
\psi_a(\Omega_{a+1}^3) = I(1,0)%
}

``中间过程应该不用再写了吧。\allowbreak{}\NovelInline{I04378}{\Omega_{a+1}^2 \times a}\allowbreak{} 的不动点是 \allowbreak{}\NovelInline{I04379}{\psi_{I(1,0)}(0)}\allowbreak{}，那么 \allowbreak{}\NovelInline{I04380}{\Omega_{a+1}^3}\allowbreak{} 就是把 \allowbreak{}\NovelInline{I04381}{I(1,0)}\allowbreak{} 的'下标'取出来，得到 \allowbreak{}\NovelInline{I04382}{I(1,0)}\allowbreak{} 本身。''

``后面又是------''

\NovelDisplay{M00832R01}{%
\psi_a(\Omega_{a+1}^3 + \Omega_{a+1}) = \Omega_{I(1,0)+1}%
}

\NovelDisplay{M00833R01}{%
\psi_a(\Omega_{a+1}^3 + \Omega_{a+1}^2) = I_{I(1,0)+1}%
}

``同样的模式。加 \allowbreak{}\NovelInline{I04383}{\Omega_{a+1}}\allowbreak{} 就是 \allowbreak{}\NovelInline{I04384}{\Omega}\allowbreak{} 下标 \allowbreak{}\NovelInline{I04385}{+1}\allowbreak{}，加 \allowbreak{}\NovelInline{I04386}{\Omega_{a+1}^2}\allowbreak{} 就是跳到 \allowbreak{}\NovelInline{I04387}{I}\allowbreak{} 的下标。''

\NovelDisplay{M00834R01}{%
\psi_a(\Omega_{a+1}^3 \times 2) = I(1,1)%
}

``\allowbreak{}\NovelInline{I04388}{\times 2}\allowbreak{} 对应 \allowbreak{}\NovelInline{I04389}{I(1,1)}\allowbreak{}------二元 Veblen 函数的第二个参数 \allowbreak{}\NovelInline{I04390}{+1}\allowbreak{}。''

\NovelDisplay{M00835R01}{%
\psi_a(\Omega_{a+1}^3 \times a) = \psi_{I(2,0)}(I(2,0))%
}

``又是不动点。这次是 \allowbreak{}\NovelInline{I04391}{\psi_I(2,0)(0)}\allowbreak{}。''

\NovelDisplay{M00836R01}{%
\psi_a(\Omega_{a+1}^4) = I(2,0)%
}

``\allowbreak{}\NovelInline{I04392}{\Omega_{a+1}^4}\allowbreak{} 对应 \allowbreak{}\NovelInline{I04393}{I(2,0)}\allowbreak{}。''

\NovelDisplay{M00837R01}{%
\psi_a(\Omega_{a+1}^\omega) = I(\omega, 0)%
}

``\allowbreak{}\NovelInline{I04394}{\omega}\allowbreak{} 次方对应 \allowbreak{}\NovelInline{I04395}{I(\omega, 0)}\allowbreak{}。''

Phyrion深吸了一口气，然后写下了一个关键的等式：

\NovelDisplay{M00838R01}{%
\psi_a(\Omega_{a+1}^a) = \psi_{I(1,0,0)}(I(1,0,0))%
}

``注意这里。\allowbreak{}\NovelInline{I04396}{a}\allowbreak{} 次方------又是不动点。这次的不动点是 \allowbreak{}\NovelInline{I04397}{\psi_{I(1,0,0)}(I(1,0,0))}\allowbreak{}。''

``也就是说，\allowbreak{}\NovelInline{I04398}{\Omega_{a+1}}\allowbreak{} 的幂次每升高一级，投影另一侧的 Veblen 函数就多一个参数。\allowbreak{}\NovelInline{I04399}{\Omega_{a+1}^2}\allowbreak{} 对应 \allowbreak{}\NovelInline{I04400}{I}\allowbreak{}，\allowbreak{}\NovelInline{I04401}{\Omega_{a+1}^3}\allowbreak{} 对应 \allowbreak{}\NovelInline{I04402}{I(1,0)}\allowbreak{}，\allowbreak{}\NovelInline{I04403}{\Omega_{a+1}^4}\allowbreak{} 对应 \allowbreak{}\NovelInline{I04404}{I(2,0)}\allowbreak{}\ldots\ldots 而 \allowbreak{}\NovelInline{I04405}{\Omega_{a+1}^a}\allowbreak{} 对应的是三元\allowbreak{}\NovelInline{I04406}{I-\phi}\allowbreak{} 函数的第一个不动点。''

Phyrion忽然收起了所有魂力。

整个修炼堂陷入了短暂的黑暗，只有魂力晶石的微光照亮着他的面容。

``现在，我要讲今天最重要的一个认知。''

他的声音变得异常严肃。

``前面我说过，在 \allowbreak{}\NovelInline{I04407}{\psi_a(\Omega_{a+1} \times a \times \text{递归运算})}\allowbreak{} 的阶段，\allowbreak{}\NovelInline{I04408}{a}\allowbreak{} 代表的是 \allowbreak{}\NovelInline{I04409}{I}\allowbreak{}。但现在，我们已经进入了 \allowbreak{}\NovelInline{I04410}{\psi_a(\Omega_{a+1}^a \times \text{递归运算})}\allowbreak{} 的阶段------\allowbreak{}\NovelInline{I04411}{\Omega_{a+1}}\allowbreak{} 的指数里面出现了 \allowbreak{}\NovelInline{I04412}{a}\allowbreak{}。''

``在这个阶段，\allowbreak{}\NovelInline{I04413}{a}\allowbreak{} 的身份发生了本质的变化。''

他在虚空中写下了两个字：

\textbf{\allowbreak{}\NovelInline{I04414}{a}\allowbreak{} 与 \allowbreak{}\NovelInline{I04415}{M}\allowbreak{}}

``在 \allowbreak{}\NovelInline{I04416}{\psi_a(\Omega_{a+1}^a \times \text{递归运算})}\allowbreak{} 这个阶段------也就是说，\allowbreak{}\NovelInline{I04417}{a}\allowbreak{} 作为 \allowbreak{}\NovelInline{I04418}{\Omega_{a+1}}\allowbreak{} 的指数的递归运算出现的这个阶段------\allowbreak{}\NovelInline{I04419}{a}\allowbreak{} 代表的是 \allowbreak{}\NovelInline{I04420}{M}\allowbreak{}。''

``\allowbreak{}\NovelInline{I04421}{M}\allowbreak{}------Mahlo 序数。正则序数的不动点。所有 \allowbreak{}\NovelInline{I04422}{\Omega}\allowbreak{}、所有 \allowbreak{}\NovelInline{I04423}{I}\allowbreak{}、所有 Veblen 函数的层级都无法触及的存在。''

``你们要理解这意味着什么。''

Phyrion的目光变得格外深邃。

``同一个 \allowbreak{}\NovelInline{I04424}{a}\allowbreak{}，在 \allowbreak{}\NovelInline{I04425}{\Omega_{a+1} \times a}\allowbreak{} 的阶段代表 \allowbreak{}\NovelInline{I04426}{I}\allowbreak{}，在 \allowbreak{}\NovelInline{I04427}{\Omega_{a+1}^a}\allowbreak{} 的阶段代表 \allowbreak{}\NovelInline{I04428}{M}\allowbreak{}。这不是矛盾------而是因为 \allowbreak{}\NovelInline{I04429}{a}\allowbreak{} 本身就大到不可思议。它的不同'侧面'，对应着不同层级的序数结构。''

``当 \allowbreak{}\NovelInline{I04430}{a}\allowbreak{} 作为系数出现时，它对应的是 \allowbreak{}\NovelInline{I04431}{I}\allowbreak{} 的层级------因为系数阶段的递归运算是乘法和加法，复杂度相对较低。''

``但当 \allowbreak{}\NovelInline{I04432}{a}\allowbreak{} 作为指数出现时，它对应的是 \allowbreak{}\NovelInline{I04433}{M}\allowbreak{} 的层级------因为指数阶段的递归运算是幂次和指数塔，复杂度发生了质的飞跃。''

``\textbf{这就是 \allowbreak{}\NovelInline{I04434}{a}\allowbreak{} 的恐怖之处。它在不同的递归层级上，能代表完全不同的巨大序数。}''

唐三沉默了很久，然后缓缓说道：``所以投影序数的强大，不仅仅在于投影本身，更在于投影的基石------\allowbreak{}\NovelInline{I04435}{a}\allowbreak{}------本身的巨大性。\allowbreak{}\NovelInline{I04436}{a}\allowbreak{} 就像一座无限高的山峰，从山脚看和从山顶看，看到的风景完全不同。''

``可以这么理解。''Phyrion说，``但更准确地说，\allowbreak{}\NovelInline{I04437}{a}\allowbreak{} 不是一个'固定高度的山峰'。它是一个\textbf{无限的阶梯}。每往上走一步，你看到的'风景'都会发生质的变化。而投影体系，就是沿着这个阶梯不断攀升的过程。''

``到这里，我需要停一下。''

Phyrion的声音变得异常严肃。

``如果你们认为 \allowbreak{}\NovelInline{I04438}{\psi_a(\Omega_{a+1}^{\Omega_{a+1}}) = I(1,0,0)}\allowbreak{} 的话------''

他停顿了一下，目光如刀。

``\textbf{那可就是踩坑了。}''

\chapter{第八十一章：杀上反射极限}\label{ux7b2cux516bux5341ux4e00ux7ae0ux6740ux4e0aux53cdux5c04ux6781ux9650}

``这个坑，是当年 test\_alpha0 也踩过的。''

唐三的瞳孔微微收缩。他本能地感觉到了某种危险------不是物理上的危险，而是思维上的陷阱。

``那正确的展开是什么？''他问。

Phyrion重新释放出魂力，在虚空中一笔一划地书写：

\NovelDisplay{M00839R01}{%
\psi_a(\Omega_{a+1}^a + \Omega_{a+1}) = \Omega_{\psi_{I(1,0,0)}(I(1,0,0))+1}%
}

``\allowbreak{}\NovelInline{I04439}{\Omega_{a+1}^a}\allowbreak{} 的基础上加一个 \allowbreak{}\NovelInline{I04440}{\Omega_{a+1}}\allowbreak{}，对应 \allowbreak{}\NovelInline{I04441}{\Omega}\allowbreak{} 的下标变成 \allowbreak{}\NovelInline{I04442}{\psi_{I(1,0,0)}(I(1,0,0)) + 1}\allowbreak{}。注意------不是 \allowbreak{}\NovelInline{I04443}{I(1,0,0) + 1}\allowbreak{}，而是 \allowbreak{}\NovelInline{I04444}{\psi_{I(1,0,0)}(I(1,0,0) + 1)}\allowbreak{}。''

``这两者有什么区别？''霍雨浩问。

``\allowbreak{}\NovelInline{I04445}{I(1,0,0)}\allowbreak{} 是一个容许序数，而 \allowbreak{}\NovelInline{I04446}{\psi_{I(1,0,0)}(I(1,0,0))<I(1,0,0)}\allowbreak{}

``继续------''

\NovelDisplay{M00840R01}{%
\psi_a(\Omega_{a+1}^a + \Omega_{a+1}^2) = I_{\psi_{I(1,0,0)}(I(1,0,0))+1}%
}

\NovelDisplay{M00841R01}{%
\psi_a(\Omega_{a+1}^a + \Omega_{a+1}^\omega) = I(\omega, \psi_{I(1,0,0)}(I(1,0,0))+1)%
}

``\allowbreak{}\NovelInline{I04447}{\Omega_{a+1}^a + \Omega_{a+1}^\omega}\allowbreak{} 对应 \allowbreak{}\NovelInline{I04448}{I(\omega, \psi_{I(1,0,0)}(I(1,0,0))+1)}\allowbreak{}。看到了吗？二元 Veblen 函数的第二个参数变成了 \allowbreak{}\NovelInline{I04449}{\psi_{I(1,0,0)}(I(1,0,0))+1}\allowbreak{}。''

\NovelDisplay{M00842R01}{%
\psi_a(\Omega_{a+1}^a \times 2) = \psi_{I(1,0,0)}(I(1,0,0)\times 2)%
}

``\allowbreak{}\NovelInline{I04450}{\times 2}\allowbreak{} 对应 \allowbreak{}\NovelInline{I04451}{\psi_{I(1,0,0)}(I(1,0,0)\times 2)}\allowbreak{}------第二个不动点。''

\NovelDisplay{M00843R01}{%
\psi_a(\Omega_{a+1}^a \times a) = \psi_{I(1,0,0)}(I(1,0,0)^2)%
}

``\allowbreak{}\NovelInline{I04452}{\times a}\allowbreak{} 又是不动点。这次的不动点是 \allowbreak{}\NovelInline{I04453}{\psi_{I(1,0,0)}(I(1,0,0)^2)}\allowbreak{}。''

Phyrion写下了一个关键的等式：

\NovelDisplay{M00844R01}{%
\psi_a(\Omega_{a+1}^{a+1}) = I(1,0,0)%
}

``这里需要特别注意。\allowbreak{}\NovelInline{I04454}{\Omega_{a+1}^{a+1}}\allowbreak{} 对应的是 \allowbreak{}\NovelInline{I04455}{I(1,0,0)}\allowbreak{}。''

``而 \allowbreak{}\NovelInline{I04456}{I(1,0,0)}\allowbreak{}，实际上是 \allowbreak{}\NovelInline{I04457}{(2\ 1\text{-})^{(1,1)}\ 2}\allowbreak{} 的一种表达。指数上的 \allowbreak{}\NovelInline{I04458}{a+1}\allowbreak{}，代表的就是这个 \allowbreak{}\NovelInline{I04459}{(1,1)}\allowbreak{}。''

``你们要理解------\allowbreak{}\NovelInline{I04460}{a+1}\allowbreak{} 不是一个简单的数字。在这个体系里，它代表的是一种结构。\allowbreak{}\NovelInline{I04461}{a}\allowbreak{} 是不动点操作，\allowbreak{}\NovelInline{I04462}{+1}\allowbreak{} 是在不动点之上的第一次迭代。所以 \allowbreak{}\NovelInline{I04463}{a+1}\allowbreak{} 代表的是'在不动点之上再走一步'。''

``继续------''

\NovelDisplay{M00845R01}{%
\psi_a(\Omega_{a+1}^{a+2}) = I(1,1,0)%
}

\NovelDisplay{M00846R01}{%
\psi_a(\Omega_{a+1}^{a+\omega}) = I(1,\omega,0)%
}

\NovelDisplay{M00847R01}{%
\psi_a(\Omega_{a+1}^{a \times 2}) = \psi_{I(2,0,0)}(I(2,0,0))%
}

``\allowbreak{}\NovelInline{I04464}{a \times 2}\allowbreak{} 对应 \allowbreak{}\NovelInline{I04465}{\psi_I(2,0,0)(0)}\allowbreak{}------又一个不动点。''

\NovelDisplay{M00848R01}{%
\psi_a(\Omega_{a+1}^{a \times 2 + 1}) = I(2,0,0)%
}

\NovelDisplay{M00849R01}{%
\psi_a(\Omega_{a+1}^{a^2}) = I(1,0,0,0)%
}

``\allowbreak{}\NovelInline{I04466}{a^2}\allowbreak{} 对应四元 Veblen 函数的 \allowbreak{}\NovelInline{I04467}{I(1,0,0,0)}\allowbreak{}。\allowbreak{}\NovelInline{I04468}{a}\allowbreak{} 的连乘每多一次，Veblen 函数的参数就多一个。''

\NovelDisplay{M00850R01}{%
\psi_a(\Omega_{a+1}^{a^ω}) = I(1@\omega)%
}

``\allowbreak{}\NovelInline{I04469}{a}\allowbreak{} 连乘 \allowbreak{}\NovelInline{I04470}{\omega}\allowbreak{} 次，对应 \allowbreak{}\NovelInline{I04471}{I(1@\omega)}\allowbreak{}------这里的 \allowbreak{}\NovelInline{I04472}{@}\allowbreak{} 表示参数列表的延伸。''

\NovelDisplay{M00851R01}{%
\psi_a(\Omega_{a+1}^{a^a+ 1}) = I(1@(1,0))%
}

``\allowbreak{}\NovelInline{I04473}{a^a+1}\allowbreak{}，对应 \allowbreak{}\NovelInline{I04474}{I(1@(1,0))}\allowbreak{}。注意这是容许点进位，一定不能错''

Phyrion停了下来，目光扫过众人。

``实际上，这里的 \allowbreak{}\NovelInline{I04475}{a}\allowbreak{} 对应的是 \allowbreak{}\NovelInline{I04476}{\psi_M}\allowbreak{} 里面的 \allowbreak{}\NovelInline{I04477}{M}\allowbreak{}。''

``也就是说，当 \allowbreak{}\NovelInline{I04478}{\Omega_{a+1}}\allowbreak{} 的指数里面出现 \allowbreak{}\NovelInline{I04479}{a}\allowbreak{} 的复杂运算时，投影的另一侧已经到达了 \allowbreak{}\NovelInline{I04480}{M}\allowbreak{} 的层级。这就是 \allowbreak{}\NovelInline{I04481}{a}\allowbreak{} 的第二重身份------在指数阶段，\allowbreak{}\NovelInline{I04482}{a}\allowbreak{} 不再是 \allowbreak{}\NovelInline{I04483}{I}\allowbreak{}，而是 \allowbreak{}\NovelInline{I04484}{M}\allowbreak{}。''

``总之------''

\NovelDisplay{M00852R01}{%
\psi_a(\Omega_{a+1}^{a^ω}) = \psi_M(M^\omega)%
}

``看到了吗？左边是 \allowbreak{}\NovelInline{I04485}{\Omega_{a+1}}\allowbreak{} 的指数里面 \allowbreak{}\NovelInline{I04486}{a}\allowbreak{} 连乘了 \allowbreak{}\NovelInline{I04487}{\omega}\allowbreak{} 次，右边是 \allowbreak{}\NovelInline{I04488}{\psi_M(M^\omega)}\allowbreak{}。\allowbreak{}\NovelInline{I04489}{a}\allowbreak{} 对应 \allowbreak{}\NovelInline{I04490}{M}\allowbreak{}，\allowbreak{}\NovelInline{I04491}{\Omega_{a+1}}\allowbreak{} 对应 \allowbreak{}\NovelInline{I04492}{M}\allowbreak{} 之上的结构。''

``然后------''

\NovelDisplay{M00853R01}{%
\psi_a(\Omega_{a+1}^{\psi_{\Omega_{a+1}}(M_{a+\omega})}) = \psi_M(M_\omega)%
}

``\allowbreak{}\NovelInline{I04493}{\psi_{\Omega_{a+1}}(M_{a+\omega})}\allowbreak{} 对应 \allowbreak{}\NovelInline{I04494}{\psi_M(M_\omega)}\allowbreak{}。又和之前一样对应上了。''

Phyrion深吸了一口气，然后写下了一个被方框框住的等式：

\NovelDisplay{M00854R01}{%
\boxed{\psi_a(\Omega_{a+1}^{\Omega_{a+1}}) = M}%
}

整个修炼堂陷入了短暂的沉默。

\allowbreak{}\NovelInline{I04495}{M}\allowbreak{}。

那个传说中的大序数------Mahlo 序数。它代表着正则序数的不动点，是所有 \allowbreak{}\NovelInline{I04496}{\Omega}\allowbreak{}、\allowbreak{}\NovelInline{I04497}{I}\allowbreak{}、以及更高层级序数都无法触及的存在。

而现在，它就这样出现在了 \allowbreak{}\NovelInline{I04498}{\psi_a(\Omega_{a+1}^{\Omega_{a+1}})}\allowbreak{} 的位置上。

Phyrion的目光变得深邃。

``从现在开始，我们进入一个新的领域。''

``后面写 OCN 也不会变得好懂了。我就直接写反射了。''

他在虚空中写下了第一个反射记号的等式：

\NovelDisplay{M00855R01}{%
\psi_a(\Omega_{a+1}^{\Omega_{a+1}}) = 2\text{-}2%
}

``\allowbreak{}\NovelInline{I04499}{M}\allowbreak{} 在反射记号里面，写作 \allowbreak{}\NovelInline{I04500}{2\text{-}2}\allowbreak{}。''

``然后继续------''

\NovelDisplay{M00856R01}{%
\psi_a(\Omega_{a+1}^{\Omega_{a+1}} + \Omega_{a+1}) = 2 \text{ aft } 2\text{-}2%
}

``加一个 \allowbreak{}\NovelInline{I04501}{\Omega_{a+1}}\allowbreak{}，对应 \allowbreak{}\NovelInline{I04502}{2 \text{ aft } 2\text{-}2}\allowbreak{}------`在 \allowbreak{}\NovelInline{I04503}{2\text{-}2}\allowbreak{} 之后的第一个 \allowbreak{}\NovelInline{I04504}{2}\allowbreak{}'。''

\NovelDisplay{M00857R01}{%
\psi_a(\Omega_{a+1}^{\Omega_{a+1}} + \Omega_{a+1}^2) = 2\ 1\text{-}2 \text{ aft } 2\text{-}2%
}

``加 \allowbreak{}\NovelInline{I04505}{\Omega_{a+1}^2}\allowbreak{}，对应 \allowbreak{}\NovelInline{I04506}{2\ 1\text{-}2 \text{ aft } 2\text{-}2}\allowbreak{}。''

\NovelDisplay{M00858R01}{%
\psi_a(\Omega_{a+1}^{\Omega_{a+1}} + \Omega_{a+1}^\omega) = (2\ 1\text{-})^\omega\ 2 \text{ aft } 2\text{-}2%
}

``加 \allowbreak{}\NovelInline{I04507}{\Omega_{a+1}^\omega}\allowbreak{}，对应 \allowbreak{}\NovelInline{I04508}{(2\ 1\text{-})^\omega\ 2 \text{ aft } 2\text{-}2}\allowbreak{}。''

Phyrion的书写速度越来越快：

\NovelDisplay{M00859R01}{%
\psi_a(\Omega_{a+1}^{\Omega_{a+1}} + \Omega_{a+1}^a) = (2\ 1\text{-})^{(1,0)}\ 2 \text{ aft } 2\text{-}2%
}

\NovelDisplay{M00860R01}{%
\psi_a(\Omega_{a+1}^{\Omega_{a+1}} + \Omega_{a+1}^{a+1}) = (2\ 1\text{-})^{(1,1)}\ 2 \text{ aft } 2\text{-}2%
}

\NovelDisplay{M00861R01}{%
\psi_a(\Omega_{a+1}^{\Omega_{a+1}} + \Omega_{a+1}^{a \times 2}) = (2\ 1\text{-})^{(2,0)}\ 2 \text{ aft } 2\text{-}2%
}

\NovelDisplay{M00862R01}{%
\psi_a(\Omega_{a+1}^{\Omega_{a+1}} + \Omega_{a+1}^{a^2}) = (2\ 1\text{-})^{(1,0,0)}\ 2 \text{ aft } 2\text{-}2%
}

``看到了吗？反射记号里面的结构，和之前 Veblen 函数的结构完全对应。\allowbreak{}\NovelInline{I04509}{(2\ 1\text{-})}\allowbreak{} 的指数，就是 Veblen 函数的参数。''

``然后------''

\NovelDisplay{M00863R01}{%
\psi_a(\Omega_{a+1}^{\Omega_{a+1}} \times 2) = \text{2nd } 2\text{-}2%
}

``\allowbreak{}\NovelInline{I04510}{\times 2}\allowbreak{} 对应第二个 \allowbreak{}\NovelInline{I04511}{2\text{-}2}\allowbreak{}。''

\NovelDisplay{M00864R01}{%
\psi_a(\Omega_{a+1}^{\Omega_{a+1}} \times \omega) = 1\text{-}2\text{-}2%
}

``\allowbreak{}\NovelInline{I04512}{\times \omega}\allowbreak{} 对应 \allowbreak{}\NovelInline{I04513}{1\text{-}2\text{-}2}\allowbreak{}。''

\NovelDisplay{M00865R01}{%
\psi_a(\Omega_{a+1}^{\Omega_{a+1}} \times a) = (1\text{-})^{(1,0)}\ 2\text{-}2%
}

``\allowbreak{}\NovelInline{I04514}{\times a}\allowbreak{} 对应 \allowbreak{}\NovelInline{I04515}{(1\text{-})^{(1,0)}\ 2\text{-}2}\allowbreak{}。又是不动点。''

Phyrion的语速进一步加快：

\NovelDisplay{M00866R01}{%
\psi_a(\Omega_{a+1}^{\Omega_{a+1}+1}) = 2\ 1\text{-}2\text{-}2%
}

\NovelDisplay{M00867R01}{%
\psi_a(\Omega_{a+1}^{\Omega_{a+1}+2}) = 2\ 1\text{-}2\ 1\text{-}2\text{-}2%
}

\NovelDisplay{M00868R01}{%
\psi_a(\Omega_{a+1}^{\Omega_{a+1}+a}) = (2\ 1\text{-})^{(1,0)}\ 2\text{-}2%
}

\NovelDisplay{M00869R01}{%
\psi_a(\Omega_{a+1}^{\Omega_{a+1}+a^2}) = (2\ 1\text{-})^{(1,0,0)}\ 2\text{-}2%
}

``指数上的 \allowbreak{}\NovelInline{I04516}{\Omega_{a+1}+1}\allowbreak{}、\allowbreak{}\NovelInline{I04517}{\Omega_{a+1}+2}\allowbreak{}、\allowbreak{}\NovelInline{I04518}{\Omega_{a+1}+a}\allowbreak{}，在反射记号里面变成了 \allowbreak{}\NovelInline{I04519}{2\ 1\text{-}2\text{-}2}\allowbreak{}、\allowbreak{}\NovelInline{I04520}{2\ 1\text{-}2\ 1\text{-}2\text{-}2}\allowbreak{}、\allowbreak{}\NovelInline{I04521}{(2\ 1\text{-})^{(1,0)}\ 2\text{-}2}\allowbreak{}。''

``继续------''

\NovelDisplay{M00870R01}{%
\psi_a(\Omega_{a+1}^{\Omega_{a+1} \times 2}) = 2\text{-}2\ 1\text{-}2\text{-}2%
}

\NovelDisplay{M00871R01}{%
\psi_a(\Omega_{a+1}^{\Omega_{a+1} \times 2} \times 2) = \text{2nd } 2\text{-}2\ 1\text{-}2\text{-}2%
}

\NovelDisplay{M00872R01}{%
\psi_a(\Omega_{a+1}^{\Omega_{a+1} \times 2 + 1}) = 2\ 1\text{-}2\text{-}2\ 1\text{-}2\text{-}2%
}

Phyrion在这里停了下来，看着众人。

蔡须坤忽然开口：``其实不用写这么细的吧？''

Phyrion微微一笑：``你觉得你都会了？''

``嗯，我知道规律了。''蔡须坤说，``反射记号里面的结构，就是不断地在已有的结构前面加上新的层。每加一个 \allowbreak{}\NovelInline{I04522}{\Omega_{a+1}}\allowbreak{} 的幂次，就在最前面加一个新的 \allowbreak{}\NovelInline{I04523}{2\ 1\text{-}}\allowbreak{}。''

``那我来考考你。''Phyrion说，``\allowbreak{}\NovelInline{I04524}{\psi_a(\Omega_{a+1}^{\Omega_{a+1}^2})}\allowbreak{} 是什么？''

蔡须坤闭上了眼睛。

片刻后，她睁开眼睛，试探性地说：``\allowbreak{}\NovelInline{I04525}{2\text{-}2\ 1\text{-}2\text{-}2\ 1\text{-}2\text{-}2}\allowbreak{}？''

Phyrion摇了摇头。

``不对。''

蔡须坤愣了一下。

``你们看------''Phyrion在虚空中书写，``你们忽略了中间的步骤。''

\NovelDisplay{M00873R01}{%
\psi_a(\Omega_{a+1}^{\Omega_{a+1} \times 2 + a}) = (2\ 1\text{-})^{(1,0)}\ 2\text{-}2\ 1\text{-}2\text{-}2%
}

``\allowbreak{}\NovelInline{I04526}{\Omega_{a+1} \times 2 + a}\allowbreak{}------在 \allowbreak{}\NovelInline{I04527}{\Omega_{a+1} \times 2}\allowbreak{} 的基础上加了 \allowbreak{}\NovelInline{I04528}{a}\allowbreak{}。这对应的是在反射记号里面取了一次不动点。''

\NovelDisplay{M00874R01}{%
\psi_a(\Omega_{a+1}^{\Omega_{a+1} \times 3}) = 2\text{-}2\ 1\text{-}2\text{-}2\ 1\text{-}2\text{-}2%
}

``\allowbreak{}\NovelInline{I04529}{\Omega_{a+1} \times 3}\allowbreak{} 才是 \allowbreak{}\NovelInline{I04530}{2\text{-}2\ 1\text{-}2\text{-}2\ 1\text{-}2\text{-}2}\allowbreak{}。不是 \allowbreak{}\NovelInline{I04531}{\Omega_{a+1}^2}\allowbreak{}。''

``然后------''

\NovelDisplay{M00875R01}{%
\psi_a(\Omega_{a+1}^{\Omega_{a+1} \times \omega}) = (2\text{-}2\ 1\text{-})^\omega\ 2\text{-}2%
}

\NovelDisplay{M00876R01}{%
\psi_a(\Omega_{a+1}^{\Omega_{a+1} \times a}) = (2\text{-}2\ 1\text{-})^{(1,0)}\ 2\text{-}2%
}

\NovelDisplay{M00877R01}{%
\psi_a(\Omega_{a+1}^{\Omega_{a+1} \times a \times 2}) = (2\text{-}2\ 1\text{-})^{(2,0)}\ 2\text{-}2%
}

\NovelDisplay{M00878R01}{%
\psi_a(\Omega_{a+1}^{\Omega_{a+1} \times a^2}) = (2\text{-}2\ 1\text{-})^{(1,0,0)}\ 2\text{-}2%
}

``看到了吗？\allowbreak{}\NovelInline{I04532}{\Omega_{a+1} \times a}\allowbreak{} 对应的是 \allowbreak{}\NovelInline{I04533}{(2\text{-}2\ 1\text{-})^{(1,0)}\ 2\text{-}2}\allowbreak{}。\allowbreak{}\NovelInline{I04534}{a}\allowbreak{} 的连乘运算对应 \allowbreak{}\NovelInline{I04535}{(2\text{-}2\ 1\text{-})}\allowbreak{} 的指数。''

``然后------''

Phyrion写下了一个关键的等式：

\NovelDisplay{M00879R01}{%
\psi_a(\Omega_{a+1}^{\Omega_{a+1}^2}) = 2\text{-}2\text{-}2%
}

``所以这个是 \allowbreak{}\NovelInline{I04536}{N}\allowbreak{}。''

``三层 \allowbreak{}\NovelInline{I04537}{2\text{-}}\allowbreak{}，代表三层反射。\allowbreak{}\NovelInline{I04538}{2\text{-}2}\allowbreak{} 是 \allowbreak{}\NovelInline{I04539}{M}\allowbreak{}，\allowbreak{}\NovelInline{I04540}{2\text{-}2\text{-}2}\allowbreak{} 是 \allowbreak{}\NovelInline{I04541}{N}\allowbreak{}------比 \allowbreak{}\NovelInline{I04542}{M}\allowbreak{} 更高一级的反射序数。''

``从 \allowbreak{}\NovelInline{I04543}{N}\allowbreak{} 继续------''

Phyrion的魂力在虚空中飞速书写，公式一个接一个地浮现：

\NovelDisplay{M00880R01}{%
\psi_a(\Omega_{a+1}^{\Omega_{a+1}^2} + \Omega_{a+1}) = 2 \text{ aft } 2\text{-}2\text{-}2%
}

\NovelDisplay{M00881R01}{%
\psi_a(\Omega_{a+1}^{\Omega_{a+1}^2} + \Omega_{a+1}^{\Omega_{a+1}}) = 2\text{-}2 \text{ aft } 2\text{-}2\text{-}2%
}

``\allowbreak{}\NovelInline{I04544}{+ \Omega_{a+1}^{\Omega_{a+1}}}\allowbreak{} 对应 \allowbreak{}\NovelInline{I04545}{+ 2\text{-}2}\allowbreak{}。因为 \allowbreak{}\NovelInline{I04546}{\Omega_{a+1}^{\Omega_{a+1}}}\allowbreak{} 本身就是 \allowbreak{}\NovelInline{I04547}{2\text{-}2}\allowbreak{}。''

\NovelDisplay{M00882R01}{%
\psi_a(\Omega_{a+1}^{\Omega_{a+1}^2} \times 2) = \text{2nd } 2\text{-}2\text{-}2%
}

\NovelDisplay{M00883R01}{%
\psi_a(\Omega_{a+1}^{\Omega_{a+1}^2} \times a) = (1\text{-})^{(1,0)}\ 2\text{-}2\text{-}2%
}

``\allowbreak{}\NovelInline{I04548}{\times a}\allowbreak{} 又是不动点。''

\NovelDisplay{M00884R01}{%
\psi_a(\Omega_{a+1}^{\Omega_{a+1}^2+1}) = 2\ 1\text{-}2\text{-}2\text{-}2%
}

\NovelDisplay{M00885R01}{%
\psi_a(\Omega_{a+1}^{\Omega_{a+1}^2+a}) = (2\ 1\text{-})^{(1,0)}\ 2\text{-}2\text{-}2%
}

\NovelDisplay{M00886R01}{%
\psi_a(\Omega_{a+1}^{\Omega_{a+1}^2+\Omega_{a+1}}) = 2\text{-}2\ 1\text{-}2\text{-}2\text{-}2%
}

\NovelDisplay{M00887R01}{%
\psi_a(\Omega_{a+1}^{\Omega_{a+1}^2+\Omega_{a+1} \times 2}) = 2\text{-}2\ 1\text{-}2\text{-}2\ 1\text{-}2\text{-}2\text{-}2%
}

\NovelDisplay{M00888R01}{%
\psi_a(\Omega_{a+1}^{\Omega_{a+1}^2+\Omega_{a+1} \times a}) = (2\text{-}2\ 1\text{-})^{(1,0)}\ 2\text{-}2\text{-}2%
}

``规律和之前完全一样。每加一个 \allowbreak{}\NovelInline{I04549}{\Omega_{a+1}}\allowbreak{} 的幂次，就在反射记号的前面加一个新的结构。''

\NovelDisplay{M00889R01}{%
\psi_a(\Omega_{a+1}^{\Omega_{a+1}^2 \times 2}) = 2\text{-}2\text{-}2\ 1\text{-}2\text{-}2\text{-}2%
}

\NovelDisplay{M00890R01}{%
\psi_a(\Omega_{a+1}^{\Omega_{a+1}^2 \times a}) = (2\text{-}2\text{-}2\ 1\text{-})^{(1,0)}\ 2\text{-}2\text{-}2%
}

\NovelDisplay{M00891R01}{%
\psi_a(\Omega_{a+1}^{\Omega_{a+1}^3}) = 2\text{-}2\text{-}2\text{-}2%
}

``\allowbreak{}\NovelInline{I04550}{\Omega_{a+1}^3}\allowbreak{} 在指数上------对应四层 \allowbreak{}\NovelInline{I04551}{2\text{-}}\allowbreak{}。''

``然后------''

\NovelDisplay{M00892R01}{%
\psi_a(\Omega_{a+1}^{\Omega_{a+1}^\omega}) = (2\text{-})^\omega%
}

\NovelDisplay{M00893R01}{%
\psi_a(\Omega_{a+1}^{\Omega_{a+1}^a}) = (2\text{-})^{(1,0)}%
}

``\allowbreak{}\NovelInline{I04552}{\omega}\allowbreak{} 次方对应 \allowbreak{}\NovelInline{I04553}{(2\text{-})^\omega}\allowbreak{}，\allowbreak{}\NovelInline{I04554}{a}\allowbreak{} 次方对应 \allowbreak{}\NovelInline{I04555}{(2\text{-})^{(1,0)}}\allowbreak{}。''

Phyrion深吸了一口气，然后写下了一个让所有人都屏住呼吸的等式：

\NovelDisplay{M00894R01}{%
\psi_a(\Omega_{a+1}^{\Omega_{a+1}^{\Omega_{a+1}}}) = 3%
}

整个修炼堂陷入了死一般的寂静。

\allowbreak{}\NovelInline{I04556}{3}\allowbreak{}。

不是 \allowbreak{}\NovelInline{I04557}{2\text{-}2\text{-}2\text{-}2}\allowbreak{}，不是 \allowbreak{}\NovelInline{I04558}{(2\text{-})^\omega}\allowbreak{}，不是 \allowbreak{}\NovelInline{I04559}{(2\text{-})^{(1,0)}}\allowbreak{}。

而是 \allowbreak{}\NovelInline{I04560}{3}\allowbreak{}。

一个简简单单的 \allowbreak{}\NovelInline{I04561}{3}\allowbreak{}。

``三层指数塔------\allowbreak{}\NovelInline{I04562}{\Omega_{a+1}^{\Omega_{a+1}^{\Omega_{a+1}}}}\allowbreak{}------对应的是 \allowbreak{}\NovelInline{I04563}{3}\allowbreak{}。''Phyrion的声音在寂静中格外清晰，``反射记号里面的 \allowbreak{}\NovelInline{I04564}{3}\allowbreak{}，代表的是三层反射的极限。''

蔡须坤低声说：``这个和反射对应上了\ldots\ldots{}''

``嗯哼。''Phyrion点了点头，``有一定的相通之处。''

``然后------''

\NovelDisplay{M00895R01}{%
\psi_a(\Omega_{a+1}^{\Omega_{a+1}^{\Omega_{a+1}+1}}) = 2\text{-}3%
}

\NovelDisplay{M00896R01}{%
\psi_a(\Omega_{a+1}^{\Omega_{a+1}^{\Omega_{a+1} \times 2}}) = 3\ 2\text{-}3%
}

\NovelDisplay{M00897R01}{%
\psi_a(\Omega_{a+1}^{\Omega_{a+1}^{\Omega_{a+1}^2}}) = 3\text{-}3%
}

\NovelDisplay{M00898R01}{%
\psi_a(\Omega_{a+1}^{\Omega_{a+1}^{\Omega_{a+1}^a}}) = (3\text{-})^{(1,0)}%
}

``\allowbreak{}\NovelInline{I04565}{3\text{-}3}\allowbreak{}，\allowbreak{}\NovelInline{I04566}{(3\text{-})^{(1,0)}}\allowbreak{}------\allowbreak{}\NovelInline{I04567}{3}\allowbreak{} 的层级也在不断地自我迭代。''

``然后------''

Phyrion写下了最后一组等式：

\NovelDisplay{M00899R01}{%
\psi_a(\Omega_{a+1}^{\Omega_{a+1}^{\Omega_{a+1}^{\Omega_{a+1}}}}) = 4%
}

``四层指数塔对应 \allowbreak{}\NovelInline{I04568}{4}\allowbreak{}。五层指数塔对应 \allowbreak{}\NovelInline{I04569}{5}\allowbreak{}。''

``反射记号里面的数字，就是指数塔的高度。每多一层指数塔，数字就 \allowbreak{}\NovelInline{I04570}{+1}\allowbreak{}。''

Phyrion收回了手，目光扫过众人。

``到目前为止，我们一直在反射记号的有限数字层级里面打转。但现在------''

他的声音变得格外庄重。

``我们要进入稳定层级了。''

他在虚空中写下了一个等式：

\NovelDisplay{M00900R01}{%
\psi_a(\varepsilon_{\Omega_{a+1}+1}) = \Pi_\omega%
}

``注意这个式子。左边不再是 \allowbreak{}\NovelInline{I04571}{\Omega_{a+1}}\allowbreak{} 的幂次塔了。左边出现了 \allowbreak{}\NovelInline{I04572}{\varepsilon_{\Omega_{a+1}+1}}\allowbreak{}------这是一个全新的东西。''

``\allowbreak{}\NovelInline{I04573}{\varepsilon}\allowbreak{} 数------\allowbreak{}\NovelInline{I04574}{\varepsilon}\allowbreak{} 序数------是指数塔的不动点。\allowbreak{}\NovelInline{I04575}{\varepsilon_0}\allowbreak{} 是 \allowbreak{}\NovelInline{I04576}{\omega^{\omega^{\omega^{\dots}}}}\allowbreak{} 的不动点。\allowbreak{}\NovelInline{I04577}{\varepsilon_{\Omega_{a+1}+1}}\allowbreak{} 是以 \allowbreak{}\NovelInline{I04578}{\Omega_{a+1}+1}\allowbreak{} 为下标的 \allowbreak{}\NovelInline{I04579}{\varepsilon}\allowbreak{} 序数。''

``而在投影的另一侧，它对应的是 \allowbreak{}\NovelInline{I04580}{\Pi_\omega}\allowbreak{}。''

``\allowbreak{}\NovelInline{I04581}{\Pi_ω}\allowbreak{}------这是稳定层级的入口。''

唐三的瞳孔微微收缩。他感觉到了------从这一刻开始，整个体系的复杂度将会发生质的飞跃。

\begin{center}\rule{0.5\linewidth}{0.5pt}\end{center}

Phyrion转过身，面对着魂力黑板，写下了今天的最后一行字：

\begin{itemize}
\tightlist
\item
  \textbf{系数阶段}（\allowbreak{}\NovelInline{I04582}{\Omega_{a+1} \times a\text{的递归运算}}\allowbreak{}）：\allowbreak{}\NovelInline{I04583}{a}\allowbreak{} 代表 \allowbreak{}\NovelInline{I04584}{I}\allowbreak{}
\item
  \textbf{指数阶段}（\allowbreak{}\NovelInline{I04585}{\Omega_{a+1}^{a\text{的递归运算}}}\allowbreak{}）：\allowbreak{}\NovelInline{I04586}{a}\allowbreak{} 代表 \allowbreak{}\NovelInline{I04587}{M}\allowbreak{}
\item
  \textbf{指数塔阶段}（\allowbreak{}\NovelInline{I04588}{\Omega_{a+1}^{\Omega_{a+1}^{\dots}}}\allowbreak{}）：\allowbreak{}\NovelInline{I04589}{a}\allowbreak{} 代表反射序数的层级
\end{itemize}

``这就是投影的终极奥义。''

``今天的课，就到这里。''

他转过身，向修炼堂的出口走去。白色母猪跟在他的身后，粉红色的圆眼睛最后扫了果糕小队众人一眼。

走了几步，Phyrion忽然停了下来，没有回头。

``你们回去之后，把今天讲的所有公式重新推导一遍。特别要注意三个地方------''

``理解了这三个等式，后面的内容就会顺畅很多。''

石门缓缓关闭。

果糕小队众人坐在原地，久久没有动弹。

石台上密密麻麻地布满了粉红色的魂力痕迹，那些公式像是一幅巨大的星图，记录着从 OFP 到 Veblen 层级的漫漫征途。

唐三低头看着自己的笔记本------上面已经写满了整整四十页的公式和注释。他在最后一页的空白处，缓缓写下了三行字：

\begin{quote}
\allowbreak{}\NovelInline{I04590}{a}\allowbreak{} 在系数阶段代表 \allowbreak{}\NovelInline{I04591}{I}\allowbreak{}

\allowbreak{}\NovelInline{I04592}{a}\allowbreak{} 在指数阶段代表\allowbreak{}\NovelInline{I04593}{M}\allowbreak{}

\allowbreak{}\NovelInline{I04594}{a}\allowbreak{} 足够大到任何容许序数望尘莫及
\end{quote}

他合上笔记本，深吸了一口气。

窗外，晨光正好。

\chapter{第八十二章：最小的lift提升}\label{ux7b2cux516bux5341ux4e8cux7ae0ux6700ux5c0fux7684liftux63d0ux5347}

phyrion负手而立，目光如炬，缓缓扫过众人，沉声道：``回顾昨日所学------''

他抬手，指尖在虚空中轻轻一划，两道算式浮现而出，灵光流转不定：

\NovelDisplay{M00901R01}{%
\psi_a(\Omega_{a+1} \times \omega) = \Omega_\omega%
} \NovelDisplay{M00902R01}{%
\psi_a(\psi_{\Omega_{a+1}}(\Omega_{a+\omega})) = \psi_{\Omega_2}(\Omega_\omega)%
}

``那么你们认为------''他的声音不疾不徐，却带着不容置疑的分量，``下面这个等式\ldots\ldots{}''

又一道算式在虚空中凝成，金光微颤：

\NovelDisplay{M00903R01}{%
\psi(\psi_a(\Omega_{a+1} \times \omega)) = \psi(\psi_a(\psi_{\Omega_{a+1}}(\Omega_{a+\omega})))%
}

``------成立与否？''

蔡须坤率先抱拳，朗声道：``后者并非标准式，远没有 \allowbreak{}\NovelInline{I04595}{\psi(\Omega_2)}\allowbreak{} 大。''

phyrion微微摇头，嘴角浮起一丝意味深长的笑意：``这是OCF的规则。投影是OCN，二者区别极大。前者是定义严格的函数，后者不过是一种大数记号。在OCN的框架之内，两者可以视作等量齐观。换言之------''

他顿了顿，加重了语气，一字一顿道：``在投影之中，\allowbreak{}\NovelInline{I04596}{\psi}\allowbreak{} 内部的表达式更大，并不意味着它所表示的递归序数也更大。''

此言一出，众人面面相觑，心中皆是一震。这一句话，颠覆了他们长久以来的直觉。

phyrion没有给他们消化的时间，紧接着说道：``不止如此。你们想想，\allowbreak{}\NovelInline{I04597}{\psi(\psi_a(a))}\allowbreak{} 等于 \allowbreak{}\NovelInline{I04598}{\psi(\psi_{\Omega_2}(\Omega_2))}\allowbreak{}。而 \allowbreak{}\NovelInline{I04599}{\psi(\psi_a(\psi_{\Omega_{a+1}}(\Omega_{a+1})))}\allowbreak{}，你们大可以将其视作------''

\NovelDisplay{M00904R01}{%
\psi(\psi_{\Omega_2}(\varepsilon_{\Omega_2+1})) = \psi(\psi_{\Omega_2}(\psi_{\Omega_3}(\Omega_3)))%
}

``再往后推------''

\NovelDisplay{M00905R01}{%
\psi(\psi_{\Omega_2}(\psi_{\Omega_3}(\varepsilon_{\Omega_3+1}))) = \psi(\psi_{\Omega_2}(\psi_{\Omega_3}(\psi_{\Omega_4}(\Omega_4)))) = \psi(\psi_a(\psi_{\Omega_{a+1}}(\psi_{a_2}(a_2))))%
}

唐三忍不住开口，眼中闪过一丝恍然：``前辈，您想说的是\ldots\ldots 什么的展开？''

``正是 \allowbreak{}\NovelInline{I04600}{\psi(\psi_a(\Omega_{a+1} \times \omega))}\allowbreak{} 的展开。''phyrion目中精光一闪，``继续往下------''

\NovelDisplay{M00906R01}{%
\psi(\psi_a(\psi_{\Omega_{a+1}}(\psi_{a_2}(\psi_{\Omega_{a_2+1}}(\psi_{a_3}(a_3)))))) = \psi(\psi_{\Omega_2}(\psi_{\Omega_3}(\psi_{\Omega_4}(\psi_{\Omega_5}(\Omega_5)))))%
}

他抬手一引，一道漫长的上确界序列如长龙般铺展开来，灵光层层叠叠，蔚为壮观：

\NovelDisplay{M00907R01}{%
\psi(\psi_a(\Omega_{a+1} \times \omega)) = \sup\Big[\psi(\psi_a(a)),\ \psi(\psi_a(\psi_{\Omega_{a+1}}(\psi_{a_2}(a_2)))),\ \psi(\psi_a(\psi_{\Omega_{a+1}}(\psi_{a_2}(\psi_{\Omega_{a_2+1}}(\psi_{a_3}(a_3)))))),\ \ldots\Big]%
}

诺当皱眉，若有所思道：``莫非这里面的 \allowbreak{}\NovelInline{I04601}{a_n}\allowbreak{} 代替了 \allowbreak{}\NovelInline{I04602}{\Omega_n}\allowbreak{} 的作用？''

phyrion没有直接回答，而是反问道：``请思考，\allowbreak{}\NovelInline{I04603}{\psi(\Omega_2+1)}\allowbreak{} 用投影表达式如何书写？记住，投影的模式是BOCF------括号里面加一，外面乘 \allowbreak{}\NovelInline{I04604}{\omega}\allowbreak{}。''

卢本韦略一沉吟，答道：``\allowbreak{}\NovelInline{I04605}{\psi(\psi_a(a))}\allowbreak{} 已经是 \allowbreak{}\NovelInline{I04606}{\psi(\Omega_2)}\allowbreak{} 了，那这应该是 \allowbreak{}\NovelInline{I04607}{\psi(\psi_a(a+1))}\allowbreak{}。''

``不错。''phyrion点头，``那么再回答，\allowbreak{}\NovelInline{I04608}{\psi(\Omega_2 + \psi_{\Omega_2}(\Omega_2))}\allowbreak{} 呢？''

蔡兴坤脱口而出：``\allowbreak{}\NovelInline{I04609}{\psi(\psi_a(a) \times 2)}\allowbreak{}。''

唐三紧接着追问：``\allowbreak{}\NovelInline{I04610}{\psi(\Omega_2 + \psi_{\Omega_2}(\Omega_2+1))}\allowbreak{} 等于 \allowbreak{}\NovelInline{I04611}{\psi(\psi_a(a+1))}\allowbreak{}，对吗？''

``没错。''phyrion赞许地看了他一眼，``既然你们已经明白此处 \allowbreak{}\NovelInline{I04612}{a}\allowbreak{} 的含义，那么 \allowbreak{}\NovelInline{I04613}{\psi(\psi_a(a \times 2)) = \psi(\Omega_2 \times 2)}\allowbreak{} 自然也就水到渠成了。下一步------''

\NovelDisplay{M00908R01}{%
\psi(\psi_a(\psi_{\Omega_{a+1}}(\Omega_{a+1}))) = \psi(\Omega_3)%
}

\NovelDisplay{M00909R01}{%
\psi(\Omega_3 + \Omega_2) = \psi(\psi_a(\psi_{\Omega_{a+1}}(\Omega_{a+1}) + a))%
}

\NovelDisplay{M00910R01}{%
\psi(\Omega_3 + \psi_{\Omega_3}(\Omega_3)) = \psi(\psi_a(\psi_{\Omega_{a+1}}(\Omega_{a+1}) \times 2))%
}

\NovelDisplay{M00911R01}{%
\psi(\Omega_3 \times 2) = \psi(\psi_a(\psi_{\Omega_{a+1}}(\Omega_{a+1} \times 2)))%
}

蔡须坤恍然大悟，一拍大腿：``我懂了！没想到 \allowbreak{}\NovelInline{I04614}{a_n}\allowbreak{} 和 \allowbreak{}\NovelInline{I04615}{\Omega_{a_n+1}}\allowbreak{} 都能起到 \allowbreak{}\NovelInline{I04616}{\Omega_n}\allowbreak{} 的折叠作用，所以这才是BO的展开式子。这里的 \allowbreak{}\NovelInline{I04617}{a}\allowbreak{}、\allowbreak{}\NovelInline{I04618}{a_2}\allowbreak{}\ldots\ldots 居然分别对应 \allowbreak{}\NovelInline{I04619}{\Omega_2}\allowbreak{}、\allowbreak{}\NovelInline{I04620}{\Omega_4}\allowbreak{}，当真神奇！''

小舞轻声感叹，目光中满是惊叹：``lift展开好生奇妙，就好像许多相似的东西叠在一起。\allowbreak{}\NovelInline{I04621}{a}\allowbreak{} 不断增大的下标，宛如平移一般，将 \allowbreak{}\NovelInline{I04622}{a}\allowbreak{} 的某种运算一层一层地平移到 \allowbreak{}\NovelInline{I04623}{a_2}\allowbreak{}、\allowbreak{}\NovelInline{I04624}{a_3}\allowbreak{}\ldots\ldots{}''

phyrion微微颔首，目光中似乎掠过一丝追忆：``在更大的尺度上，譬如 \allowbreak{}\NovelInline{I04625}{\psi(\psi_a(\Omega_{a+1} \times a \times \omega))}\allowbreak{}，同样可以用lift提升来展开。具体来说------''

他指尖轻点，虚空中浮现出一串层层嵌套的表达式：

\NovelDisplay{M00912R01}{%
\psi(\psi_a(\Omega_{a+1} \times a+\psi_{\Omega_{a+1}}(\psi_{a_2}(\Omega_{a_2+1} \times a_2+\psi_{\Omega_{a_2+1}}(\psi_{a_3}(\Omega_{a_3+1} \times a_3...))))))%
}

唐三问道：``lift提升究竟有什么用处？''

phyrion目光幽远，似在追忆往事，声音也低沉了几分：``其中妙不可言。当年那位老修士将此法传给我父亲，再由父亲传入我耳中。虽已丢失部分真谛，但你们日后若遇困境，一定要善加利用。''

蔡兴坤若有所思，试探着问道：``如此说来，\allowbreak{}\NovelInline{I04626}{\psi(\psi_a(\Omega_{a+1}^2 \times \omega))}\allowbreak{} 是否等于------''

\NovelDisplay{M00913R01}{%
\psi\Big(\psi_a\big(\Omega_{a+1} \times \psi_{\Omega_{a+1}}\big(\psi_{a_2}(\Omega_{a_2+1} \times \psi_{\Omega_{a_2+1}}(\psi_{a_3}(\cdots)))\big)\big)\Big)%
}

``正是。''phyrion道，``不过在此之前，我们先做一点简要分析。''

他抬手在空中写下一行算式：

\NovelDisplay{M00914R01}{%
\psi(\psi_a(\Omega_{a+1}^2 \times \omega)) = \psi(I_\omega)=ψ(\phi(1,,(ω,0)))%
}

\NovelDisplay{M00915R01}{%
\psi(\psi_a(\Omega_{a+1}^2+\psi_{\Omega_{a+1}}(\psi_{a_2}(\Omega_{a_2+1}^2+\psi_{\Omega_{a_2+1}}(\psi_{a_3}(\Omega_{a_3+1}^2+...))))))%
}

借助这种方法，我们写出来两种不同的lift展开表达式。之后注意了，我们先不要管这两个东西------在小的地方，我们要进行非递归分析。

首先，\allowbreak{}\NovelInline{I04627}{\psi(\psi_a(\Omega_{a+1} \times a+\psi_{\Omega_{a+1}}(\Omega_{a+1})))}\allowbreak{} 是一个关键点。部分人认为这是 \allowbreak{}\NovelInline{I04628}{\psi(I+\psi_I(I+1))=\psi(\phi(1,,(1,1)))}\allowbreak{}。

但是，这是一大误区。实际上 \allowbreak{}\NovelInline{I04629}{\psi(\psi_a(\Omega_{a+1} \times a+a))}\allowbreak{} 就是 \allowbreak{}\NovelInline{I04630}{\psi(\varepsilon_{\psi_a(\Omega_{a+1} \times a)+1})=\psi(I+\psi_I(I+1))}\allowbreak{}，而 \allowbreak{}\NovelInline{I04631}{\psi(\psi_a(\Omega_{a+1} \times a+a\text{的递归运算}))=\psi(I+\psi_I(I+1)\text{的递归运算})}\allowbreak{}。

蔡须坤恍然道：``那岂不是 \allowbreak{}\NovelInline{I04632}{\psi(\psi_a(\Omega_{a+1} \times a+\psi_{\Omega_{a+1}}(\Omega_{a+1})))=\psi(I+\psi_I(I+2))=\psi(\phi(1,,(1,2)))}\allowbreak{}？''

``正是。''phyrion抬手制止了众人急于推进的势头，``现在你们就别想太远。你们要吃透这一段的非递归分析，才能继续。回答这个式子------''

\NovelDisplay{M00916R01}{%
\psi(\psi_a(\Omega_{a+1} \times a+\psi_{\Omega_{a+1}}(\Omega_{a+1}+1)))%
}

唐三思索片刻，答道：``我觉得可以化成 \allowbreak{}\NovelInline{I04633}{\psi(\psi_a(\Omega_{a+1} \times a+\psi_{\Omega_{a+1}}(\Omega_{a+1})\times \omega))}\allowbreak{}。''

小舞摇头，轻声却坚定地说：``应当是：\allowbreak{}\NovelInline{I04634}{\psi(I+\psi_I(I+2)+\psi_{\psi_I(I+1)}(\psi_I(I+2))\times \omega)}\allowbreak{}。''

``很好。''phyrion目光中露出赞许，``现在继续猜测------''

\NovelDisplay{M00917R01}{%
\psi(\psi_a(\Omega_{a+1} \times a+\psi_{\Omega_{a+1}}(\Omega_{a+1}+a)))%
}

蔡兴坤沉声道：``这个 \allowbreak{}\NovelInline{I04635}{a}\allowbreak{} 是不动点，我猜测，恐怕已经是 \allowbreak{}\NovelInline{I04636}{\psi(I+\psi_I(I+2)+\psi_I(I+1))=\psi(φ(1,,(1,2))+φ(1,,(1,1)))}\allowbreak{}。''

``对。''phyrion点头，``你们能懂这个非标准表达式吗？我要向你们介绍一下------虽然这是非标准表达式，但是 \allowbreak{}\NovelInline{I04637}{\psi_I(I+n)=\Omega_{\psi_I(I)+n}}\allowbreak{} 仍然成立。''

他顿了顿，继续道：``好，接下来 \allowbreak{}\NovelInline{I04638}{\psi(\psi_a(\Omega_{a+1} \times a+\psi_{\Omega_{a+1}}(\Omega_{a+1}+a\text{的递归运算})))=\psi(I+\psi_I(I+2)+\psi_I(I+1)\text{的递归运算})}\allowbreak{}。比如说------''

\NovelDisplay{M00918R01}{%
\psi(\psi_a(\Omega_{a+1} \times a+\psi_{\Omega_{a+1}}(\Omega_{a+1}+\psi_{\Omega_{a+1}}(\Omega_{a+1}))))%
}

蔡须坤接道：``这不就是 \allowbreak{}\NovelInline{I04639}{\psi(I+\psi_I(I+2)+\psi_{\psi_I(I+2)}(\psi_I(I+2)))}\allowbreak{} 吗？之后，\allowbreak{}\NovelInline{I04640}{\psi(\psi_a(\Omega_{a+1} \times a+\psi_{\Omega_{a+1}}(\Omega_{a+1}\text{的递归运算})))=\psi(I+\psi_I(I+2)\text{的递归运算})}\allowbreak{}。之后，我们就来到了------''

\NovelDisplay{M00919R01}{%
\psi(\psi_a(\Omega_{a+1} \times a+\psi_{\Omega_{a+1}}(\Omega_{a+2})))=\psi(I+\psi_I(I+3))%
}

``你说的很对，非常正确。''phyrion道，``之后，我们有 \allowbreak{}\NovelInline{I04641}{\psi(\psi_a(\Omega_{a+1} \times a+\psi_{\Omega_{a+1}}(\Omega_{a+\omega})))=\psi(\psi_I(I+\omega))}\allowbreak{}。但是注意了------''

他的声音陡然一沉：

\NovelDisplay{M00920R01}{%
\psi(\psi_a(\Omega_{a+1} \times a+\psi_{\Omega_{a+1}}(\Omega_{a\times 2})))%
}

``------这是一大难点。''

卢本韦眉头紧锁，半晌才道：``嗯，让我想想\ldots\ldots{}\allowbreak{}\NovelInline{I04642}{\psi(\psi_a(\Omega_{a+1} \times a+\psi_{\Omega_{a+1}}(\Omega_{a+\psi_a(\Omega_{a+1} \times a)})))}\allowbreak{} 是 \allowbreak{}\NovelInline{I04643}{\psi(I+\psi_I(I+\psi_I(I)))}\allowbreak{}。''

奥斯卡接过话头，眼中渐渐亮起光芒：``接下来，我来想想------\allowbreak{}\NovelInline{I04644}{\psi(\psi_a(\Omega_{a+1} \times a+\psi_{\Omega_{a+1}}(\Omega_{a+\psi_a(\Omega_{a+1} \times a+a)})))}\allowbreak{} 就是 \allowbreak{}\NovelInline{I04645}{\psi(I+\psi_I(I+\psi_{\psi_I(I+1)}(\psi_I(I+1))))}\allowbreak{}。这里，\allowbreak{}\NovelInline{I04646}{\psi_{\psi_I(I+1)}(\psi_I(I+1))=\varepsilon_{\psi_I(I)+1}}\allowbreak{}。''

诺当顺着奥斯卡的思路，提出：

\NovelDisplay{M00921R01}{%
\psi(\psi_a(\Omega_{a+1} \times a+\psi_{\Omega_{a+1}}(\Omega_{a+\psi_a(\Omega_{a+1} \times a+\psi_{\Omega_{a+1}}(\Omega_{a+\omega}))})))%
}

是 \allowbreak{}\NovelInline{I04647}{\psi(I+\psi_I(I+\psi_{\psi_I(I+1)}(\psi_I(I+\omega))))}\allowbreak{}。

``很好。''phyrion道，``现在，\allowbreak{}\NovelInline{I04648}{\psi(\psi_a(\Omega_{a+1} \times a+\psi_{\Omega_{a+1}}(\Omega_{a\times 2})))=\psi(I+\psi_I(I+\psi_I(I+1)))}\allowbreak{}。之后，\allowbreak{}\NovelInline{I04649}{\psi(\psi_a(\Omega_{a+1} \times a+\psi_{\Omega_{a+1}}(\Omega_{a}\text{的递归运算})))=\psi(I+\psi_I(I+\psi_I(I+1))\text{的递归运算})}\allowbreak{}。''

霍雨浩眼前一亮，接口道：``之后 \allowbreak{}\NovelInline{I04650}{\psi(\psi_a(\Omega_{a+1} \times a+\psi_{\Omega_{a+1}}(\Omega_{\Omega_{a+1}})))}\allowbreak{} 就是 \allowbreak{}\NovelInline{I04651}{\psi(I+\psi_I(I+\psi_I(I+2)))}\allowbreak{}，好比之前我们推导的------''

\NovelDisplay{M00922R01}{%
\psi(\psi_a(\Omega_{a+1} \times a+\psi_{\Omega_{a+1}}(\Omega_{a+1})))=\psi(I+\psi_I(I+2))%
}

``它同时也是 \allowbreak{}\NovelInline{I04652}{\psi(\psi_a(\Omega_{a+1} \times a+\psi_{\Omega_{a+1}}(\psi_{a_2}(\Omega_{a_2+1} \times \psi_{a_2}(0)))))}\allowbreak{}。''

``至此，我们继续推演。''phyrion道，``不难发现------''

\NovelDisplay{M00923R01}{%
\psi(\psi_a(\Omega_{a+1} \times a+\psi_{\Omega_{a+1}}(\Omega_{\Omega_{a+1}+1})))=\psi(\psi_I(I+\psi_I(I+2)+1))%
}

\NovelDisplay{M00924R01}{%
\psi(\psi_a(\Omega_{a+1} \times a+\psi_{\Omega_{a+1}}(\Omega_{\psi_{\Omega_{a+2}}(\Omega_{a+2})})))=\psi(\psi_I(I+\psi_{\psi_I(I+3)}(\psi_I(I+3))))%
}

唐三目光灼灼，语气中带着几分确信：``现在我们几乎可以明确------\allowbreak{}\NovelInline{I04653}{\psi(\psi_a(\Omega_{a+1} \times a+\psi_{\Omega_{a+1}}(\Omega_{\Omega_{a+2}})))=\psi(\psi_I(I+\psi_I(I+3)))}\allowbreak{}。之后------''

\NovelDisplay{M00925R01}{%
\psi(\psi_a(\Omega_{a+1} \times a+\psi_{\Omega_{a+1}}(\psi_{a_2}(\Omega_{a_2+1} \times \psi_{a_2}(\Omega_{a_2+1}\times a)))))%
}

``记得是 \allowbreak{}\NovelInline{I04654}{\psi(I+\psi_I(I+\psi_I(I+\psi_I(I+1))))}\allowbreak{}。''

phyrion终于露出了满意的神色：``现在我们终于明确------''

\NovelDisplay{M00926R01}{%
\psi(\psi_a(\Omega_{a+1} \times a+\psi_{\Omega_{a+1}}(\psi_{a_2}(\Omega_{a_2+1} \times a_2))))=\psi(I\times 2)=ψ(φ(1,,(2,0)))=\psi(\psi_a(\Omega_{a+1} \times a\times 2))%
}

他话锋一转：``接下来，事情就又开始复杂了。''

\NovelDisplay{M00927R01}{%
\psi(\psi_a(\Omega_{a+1} \times a+\psi_{\Omega_{a+1}}(\psi_{a_2}(\Omega_{a_2+1} \times a_2+a))))%
}

``------这是什么？''

小舞轻声道：``我觉得，我们应当考虑更小的------''

\NovelDisplay{M00928R01}{%
\psi(\psi_a(\Omega_{a+1} \times a+\psi_{\Omega_{a+1}}(\psi_{a_2}(\Omega_{a_2+1} \times a_2))+a))%
}

唐三则道：``不如我们思考一下------''

\NovelDisplay{M00929R01}{%
\psi(\psi_a(\Omega_{a+1} \times a+\psi_{\Omega_{a+1}}(\psi_{a_2}(\Omega_{a_2+1} \times a_2)))+\psi(\psi_a(\Omega_{a+1} \times a)))%
}

``直到这里，仍然是 \allowbreak{}\NovelInline{I04655}{\psi(I\times 2+\psi_I(I))}\allowbreak{}。这就是非标准式子的好处------不过，这最后会转化成标准式子。''

phyrion接道：``现在，我们就要注意了------''

\NovelDisplay{M00930R01}{%
\psi(\psi_a(\Omega_{a+1} \times a+\psi_{\Omega_{a+1}}(\psi_{a_2}(\Omega_{a_2+1} \times a_2)))+\psi(\psi_a(\Omega_{a+1} \times a+a)))%
}

``现在只是 \allowbreak{}\NovelInline{I04656}{\psi(I\times 2+\psi_{\psi_I(I+1)}(\psi_I(I+1)))}\allowbreak{}。所以说------''

\NovelDisplay{M00931R01}{%
\psi(\psi_a(\Omega_{a+1} \times a+\psi_{\Omega_{a+1}}(\psi_{a_2}(\Omega_{a_2+1} \times a_2))\times \omega))%
}

``就是 \allowbreak{}\NovelInline{I04657}{\psi(I\times 2+\psi_{\psi_I(I+1)}(\psi_I(I\times 2))\times \omega)}\allowbreak{}。''

蔡须坤忽然道：``小舞妹妹考虑的------''

\NovelDisplay{M00932R01}{%
\psi(\psi_a(\Omega_{a+1} \times a+\psi_{\Omega_{a+1}}(\psi_{a_2}(\Omega_{a_2+1} \times a_2))+a))%
}

``居然是 \allowbreak{}\NovelInline{I04658}{\psi(I\times 2+\psi_I(I+1))}\allowbreak{}！就好比 \allowbreak{}\NovelInline{I04659}{\psi(\psi_a(\Omega_{a+1} \times a+a))=\psi(\psi_I(I+1))}\allowbreak{}。而我注意到------''

\NovelDisplay{M00933R01}{%
\psi(\psi_a(\Omega_{a+1} \times a+\psi_{\Omega_{a+1}}(\psi_{a_2}(\Omega_{a_2+1} \times a_2))+a\text{的递归运算}))%
}

``就是 \allowbreak{}\NovelInline{I04660}{\psi(I\times 2+\psi_I(I+1)\text{的递归运算})}\allowbreak{}。''

phyrion道：``所以，我们自然就可以推导这样的式子------''

\NovelDisplay{M00934R01}{%
\psi(\psi_a(\Omega_{a+1} \times a+\psi_{\Omega_{a+1}}(\psi_{a_2}(\Omega_{a_2+1} \times a_2)\times \omega)))=\psi(I\times 2+\psi_{\psi_I(I+2)}(\psi_I(I\times 2))\times \omega)%
}

\NovelDisplay{M00935R01}{%
\psi(\psi_a(\Omega_{a+1} \times a+\psi_{\Omega_{a+1}}(\psi_{a_2}(\Omega_{a_2+1} \times a_2+\Omega_{a+1}\text{的递归运算}))))=\psi(I\times 2+\psi_I(I+2)\text{的递归运算})%
}

蔡兴坤长叹一声，目中满是震撼：``天啊，太神奇了\ldots\ldots 我感觉我的功力在提升。''

``好。''phyrion道，``接下来，我们要知道 \allowbreak{}\NovelInline{I04661}{\psi(\psi_a(\Omega_{a+1} \times a+\psi_{\Omega_{a+1}}(\psi_{a_2}(\Omega_{a_2+1} \times a_2)\times 2)))}\allowbreak{} 实际上已经是 \allowbreak{}\NovelInline{I04662}{\psi(I\times 2+\psi_I(I\times 2))}\allowbreak{}。现在能推导出------''

\NovelDisplay{M00936R01}{%
\psi(\psi_a(\Omega_{a+1} \times a+\psi_{\Omega_{a+1}}(\psi_{a_2}(\Omega_{a_2+1} \times a_2+\psi_a(\Omega_{a+1} \times a)))))%
}

``------的大小吗？''

卢本韦眼中精光一闪：``我领悟了！是 \allowbreak{}\NovelInline{I04663}{\psi(I\times 2+\psi_I(I\times 2)\times \psi_I(I))}\allowbreak{}。之后------''

\NovelDisplay{M00937R01}{%
\psi(\psi_a(\Omega_{a+1} \times a+\psi_{\Omega_{a+1}}(\psi_{a_2}(\Omega_{a_2+1} \times a_2+a))))%
}

``就是 \allowbreak{}\NovelInline{I04664}{\psi(I\times 2+\psi_I(I\times 2)\times \psi_I(I+1))}\allowbreak{}。''

霍雨浩接着道，声音中带着压抑不住的兴奋：``非常漂亮！通过------''

\NovelDisplay{M00938R01}{%
\psi(\psi_a(\Omega_{a+1} \times a+\psi_{\Omega_{a+1}}(\psi_{a_2}(\Omega_{a_2+1} \times a_2+\psi_{a_2}(\Omega_{a_2+1} \times a_2)))))%
}

``就是 \allowbreak{}\NovelInline{I04665}{\psi(I\times 2+\psi_I(I\times 2)^2)}\allowbreak{}！''

然后，果糕小队一致得到------

\NovelDisplay{M00939R01}{%
\psi(\psi_a(\Omega_{a+1} \times a+\psi_{\Omega_{a+1}}(\psi_{a_2}(\Omega_{a_2+1} \times a_2+a_2))))=\psi(I\times 2+\psi_I(I\times 2+1))%
}

\allowbreak{}\NovelInline{I04666}{a_2}\allowbreak{} 的递归运算自然代表着 \allowbreak{}\NovelInline{I04667}{\psi_I(I\times 2+1)}\allowbreak{} 的递归运算。

\NovelDisplay{M00940R01}{%
\psi(\psi_a(\Omega_{a+1} \times a+\psi_{\Omega_{a+1}}(\psi_{a_2}(\Omega_{a_2+1} \times a_2+\psi_{\Omega_{a_2+1}}(\Omega_{a_2+1})))))=\psi(I\times 2+\psi_I(I\times 2+2))%
}

接下来------

\NovelDisplay{M00941R01}{%
\psi(\psi_a(\Omega_{a+1} \times a+\psi_{\Omega_{a+1}}(\psi_{a_2}(\Omega_{a_2+1} \times a_2+\psi_{\Omega_{a_2+1}}(\psi_{a_3}(\Omega_{a_3+1} \times a_3))))))=\psi(I\times 3)%
}

最终------

\NovelDisplay{M00942R01}{%
\psi(\psi_a(\Omega_{a+1} \times a\times \omega))=\psi(I\times \omega)%
}

phyrion在黑板上写下了一个式子：

\NovelDisplay{M00943R01}{%
\psi(\psi_a(\Omega_{a+1}^2+a))%
}

``现在，我们要考虑一个问题------这个式子，代表了什么序数？''

蔡须坤几乎是脱口而出：``\allowbreak{}\NovelInline{I04668}{\psi(\varepsilon_{I+1})}\allowbreak{}。在多逗号记法里，就是 \allowbreak{}\NovelInline{I04669}{\psi(\varphi(1,,(1@(1@\cdots(1,0)))))}\allowbreak{}。''

她的声音干脆利落，金色的双瞳中闪烁着自信的光芒。帝皇瑞兽的本源处女之力在她体内微微涌动，让她的思维比常人敏锐数倍。

Phyrion微微点头，目光转向唐三：``那么，你觉得 \allowbreak{}\NovelInline{I04670}{\psi(\psi_a(\Omega_{a+1}\times \psi_{\Omega_{a+1}}(\Omega_{a+1})))}\allowbreak{} 是什么？''

唐三沉吟片刻，眉头微蹙：``我觉得\ldots\ldots 还是 \allowbreak{}\NovelInline{I04671}{\psi(\varepsilon_{I+1})}\allowbreak{}，还是上面那个东西。但是------''

他的目光变得锐利起来。

``经过前面的推导，\allowbreak{}\NovelInline{I04672}{\psi(\psi_a(\Omega_{a+1}^2+a\text{的递归运算}))=\psi(\Omega_{I+1}\text{的递归运算})}\allowbreak{}。''

Phyrion的眼中闪过一丝赞许：``你们发现了 \allowbreak{}\NovelInline{I04673}{\psi(\psi_a(\Omega_{a+1}\times a\text{的递归运算}))=\psi(I\text{的递归运算})}\allowbreak{} 吗？''

小队众人齐声答道：``是。这和接下来的内容有关系吗？''

``有关系。而且关系很大。''

Phyrion的声音忽然沉了下来。他转身面对众人，表情变得前所未有的严肃。

``之前的规则现在还是适用的。唯一不同的是------''

他在黑板上写下了两行式子，左边一行，右边一行。

左边：

\NovelDisplay{M00944R01}{%
\psi\Big(\psi_a\big(\Omega_{a+1} \times \psi_{\Omega_{a+1}}\big(\psi_{a_2}(\Omega_{a_2+1} \times \psi_{\Omega_{a_2+1}}(\psi_{a_3}(\cdots)))\big)\big)\Big)%
}

右边：

\NovelDisplay{M00945R01}{%
\psi(\psi_a(\Omega_{a+1}^2+\psi_{\Omega_{a+1}}(\psi_{a_2}(\Omega_{a_2+1}^2+\psi_{\Omega_{a_2+1}}(\psi_{a_3}(\Omega_{a_3+1}^2+...))))))%
}

``看到了吗？''Phyrion用指尖点了点两行式子之间的空隙，``左边这个嵌套结构，在很小的地方就和右边这个式子追平了。''

小舞歪着头看了看，金色的马尾在肩头轻轻晃动：``它们似乎就差了一个 \allowbreak{}\NovelInline{I04674}{\Omega}\allowbreak{} 的差距。我感觉这并不重要，而且这个差距很小。''

``这要看在什么层面了。''Phyrion的语气变得意味深长，``有时候，差距也没有这么小。''

他擦掉了黑板上的式子，重新写下一个：

\NovelDisplay{M00946R01}{%
\psi(\psi_a(\Omega_{a+2}))=\psi(\phi(1,,,0))%
}

``这个，你们见到过吗？''

小队众人的表情瞬间变了。

``这已经是 \allowbreak{}\NovelInline{I04675}{0,111,22}\allowbreak{} 了。''唐三低声说道，眉头紧锁，``在表达式内层不觉得什么，但如果在外面直接放置这个式子\ldots\ldots 那就是一个巨大无比的序数。''

Phyrion点了点头，又在黑板上写下了另一个等式：

\NovelDisplay{M00947R01}{%
\psi\Big(\psi_a\big(\Omega_{a+1} \times \psi_{\Omega_{a+1}}\big(\psi_{a_2}(\Omega_{a_2+1} \times \omega)\big)\big)\Big)=\psi(\psi_a(\Omega_{a+1}^2+\psi_{\Omega_{a+1}}(\psi_{a_2}(\Omega_{a_2+1}\times \omega))))%
}

``这一块，是无足轻重的。你们知道为什么吗？''

众人摇头。

``因为这就是个\textbf{初始值选取问题}，没有什么大不了的。''Phyrion的语气变得轻松了一些，``你们如果想的话，可以将这两个弄成一样的，完全没问题------因为这对于整体没有什么影响。''

诺兰忽然开口了：``那么 \allowbreak{}\NovelInline{I04676}{\psi(\psi_a(\Omega_{a+1}^{\Omega_{a+1}} \times \omega))}\allowbreak{} 也有两种展开表达式？''

``没错。''Phyrion看向卢本韦，``你来写。''

卢本韦站起身来，走到黑板前，写下了两行：

\NovelDisplay{M00948R01}{%
\psi\Big(\psi_a\big(\Omega_{a+1}^{\psi_{\Omega_{a+1}}\big(\psi_{a_2}(\Omega_{a_2+1}^{\psi_{\Omega_{a_2+1}}(\psi_{a_3}(\cdots))})\big)}\big)\Big)%
}

\NovelDisplay{M00949R01}{%
\psi(\psi_a(\Omega_{a+1}^{\Omega_{a+1}}+\psi_{\Omega_{a+1}}(\psi_{a_2}(\Omega_{a_2+1}^{\Omega_{a_2+1}}+\psi_{\Omega_{a_2+1}}(\psi_{a_3}(\Omega_{a_3+1}^{\Omega_{a_3+1}}+...))))))%
}

他退后一步，看着自己写的式子，若有所思。

蔡兴坤接道：``往往迁移至对于好点，后一种差了一个 \allowbreak{}\NovelInline{I04677}{\Omega}\allowbreak{}，但是后一种好理解。''

``很好。''Phyrion的目光在众人脸上扫过，``那你们知道，在更大的地方，一般怎么处理的吗？比如说------''

他在黑板上重重地写下了：

\NovelDisplay{M00950R01}{%
\psi(\psi_a(\Omega_{a+2}\times \omega))%
}

蔡须坤的眼睛亮了起来。帝皇瑞兽的本源之力在她体内涌动，让她的直觉比任何人都要敏锐。她几乎是本能地开口了：

``类比一下------应该是：

\NovelDisplay{M00951R01}{%
\psi(\psi_a(\psi_{\Omega_{a+2}}(\psi_{a_2}(\psi_{\Omega_{a_2+2}}(\psi_{a_3}(\cdots))))))%
}

Phyrion的嘴角微微上扬：``嗯，悟性非常好。''

他的目光转向霍雨浩：``现在你们觉得 \allowbreak{}\NovelInline{I04678}{\psi(\psi_a(\Omega_{a+2}))=\psi(\psi_a(\varepsilon_{\Omega_{a+1}}))}\allowbreak{} 是不是对的？''

霍雨浩想了想，点了点头：``我觉得就应该这样规定。''

``不过------''Phyrion话锋一转，``有时候，你们遇到了 \allowbreak{}\NovelInline{I04679}{\psi(\psi_a(\Omega_{a+2}+\Omega_{a+1}))=\psi(\psi_a(\varepsilon_{\Omega_{a+1}}))}\allowbreak{} 也不必惊讶。这都是一样的。''

他的声音变得低沉而郑重。

``投影内部表达式不决定外部大小------这就是OCN和OCF的本质区别。''

``这就是投影的初始值问题。你们记住------''

他的目光如刀锋般扫过每一个人的脸。

``不管怎么样，\textbf{灵活切换}才是根本。''

漫长的推演终于告一段落。

霍雨浩长长地呼出一口气，整个人像是被抽去了力气一般，靠在了椅背上：``好累啊\ldots\ldots{}''

众人深有同感，纷纷点头。连续数个时辰的高强度运算，让每个人的神识都近乎枯竭，脑海中那些层层嵌套的符号仍在嗡嗡作响，挥之不去。

phyrion看着众人疲惫的面容，冷峻的面容上罕见地浮现出一丝温和的笑意。他收起虚空中密密麻麻的算式，轻声道：

``我给你们做好吃的。你们这修炼，比我当年的强度还高。''

夜色已深。窗外万籁俱寂，唯有远处偶尔传来几声虫鸣。方才那些灵光闪烁的算式早已消散于虚空，只余下满室静谧，和众人心中久久不散的震撼。

\chapter{第八十三章：扽力的奥秘}\label{ux7b2cux516bux5341ux4e09ux7ae0ux627dux529bux7684ux5965ux79d8}

夜色深沉。

修炼堂的魂力晶石已经暗淡了大半，只有角落里几颗还散发着微弱的暖光，像是夜空中最后的几颗星辰。

果糕小队众人横七竖八地靠在石椅上，一个个面色苍白，眼神涣散。连续数个时辰的高强度推演，让每个人的神识都近乎枯竭。那些层层嵌套的 \allowbreak{}\NovelInline{I04680}{\psi}\allowbreak{} 符号仍在脑海中嗡嗡作响，像一群不肯散去的飞蛾。

霍雨浩长长地呼出一口气，整个人像是被抽去了力气一般，软软地靠在了椅背上。

``好累啊\ldots\ldots{}''

众人深有同感，纷纷点头。

唐三揉了揉太阳穴，苦笑道：``我感觉我的精神力已经被榨干了。再算下去，我怕我的识海要裂开。''

蔡须坤趴在石台上，金色的双瞳半眯着，连说话的力气都快没了：``我感觉\ldots\ldots 我连手指头都不想动了\ldots\ldots{}''

卢本韦靠在墙角，双眼无神地望着穹顶，嘴唇微微翕动，像是在无意识地默念着公式。

小舞歪着头靠在奥斯卡的肩膀上，已经快要睡着了。

就在这片疲惫与寂静之中，Phyrion从修炼堂的侧门走了进来。

他的手中端着一个托盘，上面摆满了热气腾腾的食物------一碗碗色泽金黄的浓汤，一盘盘散发着诱人香气的烤肉，还有几碟精致的小菜。白色母猪跟在他的身后，嘴里叼着一个比她脑袋还大的食盒，粉红色的圆眼睛在昏暗中闪闪发亮。

``来。''Phyrion把托盘放在石台中央，声音比上课时柔和了许多，``吃点东西。你们这修炼，比我当年的强度还高。''

食物的香气像是一只无形的手，瞬间把所有人从半昏迷的状态中拽了回来。

奥斯卡第一个弹了起来，以和他魂力完全不匹配的敏捷冲到了托盘前面：``Phyrion老师，您简直是救世主！''

众人纷纷围了上来。

吃饱喝足之后，众人的气色好了许多。但那种深入骨髓的疲惫感并没有完全消散------神识的消耗不是靠吃东西就能恢复的。

Phyrion坐在石台边，看着众人，沉默了片刻。

然后他开口了。

``今天的推演，你们表现得不错。''他说，``但你们有没有注意到一个问题？''

``什么问题？''唐三问。

``你们在推演的后半段，明显感觉到力不从心了。不是智力上的不足，而是------''他停顿了一下，``一种更深层的疲惫。像是有什么东西在被不断地抽取，而你们却无力补充。''

众人面面相觑。

他说得没错。那种疲惫感和平时修炼后的疲劳完全不同。平时修炼后的疲劳是魂力的消耗，休息一段时间就能恢复。但今天的那种疲惫\ldots\ldots 更像是灵魂深处被什么东西掏空了。

``那是什么？''霍雨浩问。

Phyrion伸出一根手指，粉红色的魂力在指尖凝聚，形成了一个微小的光点。

``那就是扽力。''

\begin{center}\rule{0.5\linewidth}{0.5pt}\end{center}

``扽力。''

唐三重复了一遍这个词，眉头微微皱起。他从未在任何典籍中见过这个概念。

``你们一定很好奇，扽魂殿为什么叫做'扽魂殿'。''Phyrion说，``很多人以为这个名字和武魂、魂灵有关。但实际上，这个名字来源于两种力量------''

他在虚空中写下了两个字：

\textbf{扽力} 与 \textbf{魂力}

``魂力，你们已经很熟悉了。它是魂师修炼的基础，决定了你的攻击强度、防御能力、魂技威力。简单来说------''

他的目光扫过众人。

``\textbf{魂力决定爆发的上限。}''

``而扽力------''

Phyrion的声音变得低沉而郑重。

``\textbf{扽力决定你的持续力。}''

``持续力？''奥斯卡有些困惑。

``对。持续力。''Phyrion解释道，``你们今天的推演就是最好的例子。你们的智力足够，你们的魂力也足够，但你们的扽力不足。所以当推演进行到后半段的时候，你们的神识开始枯竭，不是因为魂力不够，而是因为扽力被耗尽了。''

``魂力像是你手中的剑------锋利、强大、一击致命。但扽力像是你握剑的手臂------它决定了你能挥剑多少次，能战斗多久，能在极限状态下坚持多长时间。''

``只有这两者协调发展，才能突破境界。''

``那么，扽力怎么修炼？''蔡须坤问。她的金色双瞳中闪烁着求知的光芒。

Phyrion竖起了两根手指。

``两条路。''

``第一条路------学习记号。''

``你们这几天学的投影记号、反射记号、Veblen记号------这些不仅仅是数学工具。每掌握一种新的记号，每理解一种新的序数结构，你们的扽力就会获得相应的提升。因为记号本身就是对序数本质的一种'刻画'------当你的认知能够触及更深层的序数结构时，你的灵魂就会被拉伸、被扩展，扽力自然随之增长。''

``第二条路------''

Phyrion的目光变得格外认真。

``\textbf{扽西。}''

``扽西，说白了就是\textbf{分析}。''

众人的表情微微一滞。分析？

``你们别小看这两个字。''Phyrion看出了众人的疑惑，解释道，``扽西的分析，不是普通的思考或推理。它是一种用灵魂去'拆解'和'透视'记号结构的深层认知活动。''

``当你对着两种不同的记号进行扽西------也就是分析------的时候，你的灵魂会同时'拉住'这两种记号，深入它们的内部结构，拆解它们的运作机制，透视它们之间的关联与差异。这个过程，就像是用灵魂的手术刀，一层一层地剥开记号的外壳，直达核心。''

Phyrion伸出了双手，左手掌心浮现出一个粉红色的符号------那是投影记号的标志。右手掌心浮现出一个淡金色的符号------那是bms记号的标志。

``当两种记号在你的识海中被同时分析，它们之间就会产生一种张力。这种张力，就是扽力的来源。你用灵魂去'拉住'这两种记号，拆解它们、透视它们、理解它们------在这种高强度的认知活动中，你的扽力就会被不断地锻炼和强化。''

``但是------''

他的表情变得严肃起来。

``两种记号之间的差异越大，扽西------也就是分析------就越困难。''

``如果你分析的是两种相近的记号------比如投影记号和投影记号的变体------那难度很低，但提升也很小。就像分析两道同类型的数学题，毫不费力，但也锻炼不了什么。''

``但如果你分析的是两种差异巨大的记号------比如投影记号和bms记号，或者反射记号和bms记号------那难度就会成倍增加。就像同时分析两个完全不同领域的难题，稍有不慎就会让你的神识过载。''

``可是------''

Phyrion的声音忽然提高了。

``\textbf{两种差别巨大的记号扽西------也就是分析------结束之后，你的扽力就会获得巨大的提升。}''

``提升的幅度，和两种记号之间的差异成正比。差异越大，分析越难，但成功之后的回报也越丰厚。''

``这就是扽魂殿的核心修炼法门。我们之所以叫做'扽魂殿'，就是因为我们通过'扽'------通过分析、拆解、透视不同的记号------来锤炼灵魂的力量。''

众人听得入了神。

小舞忍不住问道：``那Phyrion老师，您的扽力\ldots\ldots 是怎么修炼到现在的程度的？''

Phyrion沉默了片刻。

然后他低下头，看了一眼脚边的白色母猪。

白色母猪''哼''了一声，用脑袋蹭了蹭他的小腿。

``你们都知道，我的武魂是魂灵。''Phyrion缓缓说道，``白色母猪------她不是普通的武魂，她是神级扽魂。''

``神级扽魂？''霍雨浩瞪大了眼睛。

``对。神级扽魂。''Phyrion说，``扽魂的品质，直接决定了魂师扽力的上限。普通扽魂能让你的扽力修炼事半功倍，而神级扽魂------''

他停了一下。

``神级扽魂能让你的扽力修炼速度提升十倍以上。''

``但这还不是全部。''

Phyrion的目光变得深邃起来，像是穿透了时间的迷雾，望向了很久以前的某个时刻。

``我之所以能突破到 \allowbreak{}\NovelInline{I04681}{Y(1,3,10)}\allowbreak{} 的境界，不仅仅是因为她是神级扽魂。更重要的是------''

他伸出手，轻轻抚摸着白色母猪的脑袋。

``------魂灵的融合。''

``当魂师与魂灵之间建立了真正的共生关系，两者之间的扽力就会互相交融、互相增益。我的扽力在融合的过程中飞速提升，她的扽力也在融合的过程中飞速提升。我们两个人的扽力像两条河流汇入同一片大海，彼此激荡，彼此壮大。''

``这就是为什么我之前说，魂灵不是工具，而是伙伴。魂灵的融合不仅仅是战力上的叠加，更是扽力上的质变。''

说到这里，Phyrion忽然话锋一转。

他的目光从白色母猪身上移开，缓缓扫过果糕小队的每一个人。

``不过------''

他的嘴角微微上扬，露出了一个意味深长的笑容。

``老夫作为养猪场的第一股东，我可以感受到------''

他的目光首先落在了卢本韦的身上。

``卢本韦，你的母猪扽魂------''

卢本韦一愣。他低头看了看自己腰间的魂导器，那里封印着他的扽魂------一头通体雪白、体型比Phyrion脚边那只还要大上一圈的母猪。

``------品质甚至比我这只更高。''

这句话像一颗石子投入了平静的湖面，修炼堂里瞬间炸开了锅。

``什么？！''卢本韦瞪大了眼睛，``比您的还高？''

``嗯。''Phyrion点了点头，语气平淡得像是在说今天天气不错，``我能感受到。你的母猪扽魂的血脉纯度极高，远超我当年遇到的大多数神级扽魂。这也是为什么你的修炼速度一直很快------扽魂的品质决定了扽力的起点，而你的起点，比绝大多数人都要高。''

卢本韦张了张嘴，半天没说出话来。他低头看了看自己的母猪扽魂，眼神里多了一种从未有过的复杂情绪------震惊、骄傲，还有一丝隐隐的不安。

Phyrion没有给他太多消化的时间，目光已经转向了蔡须坤。

``蔡须坤------''

蔡须坤微微一怔，下意识地坐直了身体。

``你的扽魂是篮球。''Phyrion说，``但篮球只是你扽魂的外在形态。真正的核心，是你身为帝皇瑞兽的本源处女之力。''

蔡须坤的瞳孔猛地一缩。

``帝皇瑞兽的本源处女之力------''Phyrion的声音变得格外郑重，``那是一种极其纯粹、极其强大的生命本源力量。它和篮球扽魂融合在一起，形成了一种独一无二的复合扽魂------外在是篮球的灵动与精准，内在是帝皇瑞兽的磅礴本源。''

``这种复合扽魂的品质------''

他停了一下，目光中闪过一丝赞叹。

``------同样是神级的。''

蔡须坤低下头，看着自己的双手。金色的光芒在她的指尖若隐若现，那是帝皇瑞兽本源之力的外在表现。她一直知道自己的力量与众不同，但从未想过，这份力量的本质竟然是''帝皇瑞兽的本源处女之力''。

``你的扽力修炼，有着双重加持。''Phyrion继续说道，``篮球扽魂赋予你分析记号时的精准与敏锐，帝皇瑞兽的本源之力则赋予你源源不断的扽力储备。这两者结合在一起，让你的扽力修炼速度远超常人。''

蔡须坤缓缓点了点头，金色的双瞳中闪烁着若有所思的光芒。

Phyrion的目光又转向了蔡兴坤。

``还有你，蔡兴坤。你的黑子鸡------''

蔡兴坤的肩膀上，一只通体漆黑的小鸡正歪着脑袋，用一双圆溜溜的眼睛好奇地打量着四周。它的羽毛在魂力晶石的光芒下泛着幽蓝色的光泽。

``------同样是神级扽魂。''

Phyrion看着这一幕，嘴角的笑意更深了。但紧接着，他的笑容缓缓收敛，目光变得深邃起来。

``不过------''他的声音低了下去，语气中多了一种说不清的凝重。``黑子鸡\ldots\ldots 并不普通。''

蔡兴坤一愣：``不普通？''

Phyrion没有直接回答。他伸出右手，粉红色的魂力在掌心凝聚，形成了一个微型的感知领域。他将这个领域缓缓向黑子鸡推了过去------

就在感知领域触碰到黑子鸡的一瞬间，Phyrion的瞳孔猛地一缩。

他的手掌像是被什么东西弹了一下，感知领域瞬间溃散。

``\ldots\ldots 果然。''Phyrion收回了手，目光变得异常深沉。

``果然什么？''蔡兴坤紧张地问。他下意识地伸手摸了摸头顶上的黑子鸡。黑子鸡歪着脑袋，``咕''了一声，用喙啄了啄他的手指，一副人畜无害的模样。

Phyrion沉默了片刻，然后缓缓说道：

``黑子鸡的体内\ldots\ldots 封印着一股极其庞大的力量。''

修炼堂里的空气瞬间凝固了。

Phyrion的目光紧紧盯着黑子鸡。黑子鸡浑然不觉，正用爪子在蔡兴坤的头发里挠痒痒。

``我无法判断这股力量的本质是什么。''Phyrion的声音变得格外低沉，``但可以确定的是------它远超黑子鸡作为扽魂本身应有的力量上限。换句话说，黑子鸡体内封印的东西，不是它自己的力量，而是某种\ldots\ldots 来自外部的、被刻意封入的力量。''

他停顿了一下，目光从黑子鸡身上移开，重新看向蔡兴坤。

``但你要记住一件事------黑子鸡的封印，和你自身的扽力修炼是息息相关的。当你的扽力提升到某个临界点的时候，封印可能会自行松动。到那时候，你体内的力量会发生什么变化，连我也无法预料。''

``所以------好好修炼。不是为了去解封，而是为了在解封的那一天，你有能力承受住那股力量。''

蔡兴坤缓缓点了点头。他低头看着掌心里的黑子鸡，眼神变得复杂起来。

他轻咳了一声，继续说道：``你们看------卢本韦的母猪，蔡须坤的篮球加帝皇瑞兽本源处女之力，蔡兴坤的黑子鸡。''他缓缓说道，``三个神级扽魂。你们果糕小队，居然坐拥三个神级扽魂。''

他的目光变得认真起来。

``这就是你们修炼速度远超常人的根源。''

``神级扽魂不仅仅是战力上的加成。更重要的是，它对扽力的修炼有着巨大的增幅效果。普通扽魂能让你的扽力修炼事半功倍，而神级扽魂------它能让你在同样的时间内，获得数倍于常人的扽力提升。''

``你们今天的推演，强度确实很高。但你们能坚持到最后，不仅仅是因为你们的智力足够，更是因为你们的神级扽魂在背后默默地支撑着你们的扽力消耗。''

``如果不是神级扽魂，换做普通魂师来承受今天这种强度的推演------''

Phyrion停顿了一下。

``------恐怕中途就会神识崩溃。''

果糕小队众人面面相觑，一时间都不知道该说什么好。

卢本韦低头看了看自己的母猪扽魂，母猪正用粉红色的圆鼻子拱着他的手心，像是在撒娇。

蔡须坤摊开手掌，金色的光芒在掌心缓缓流转，篮球扽魂的轮廓在光芒中若隐若现。她能感觉到，在那份灵动与精准之下，还有一股更加深沉、更加磅礴的力量在涌动------那是帝皇瑞兽的本源处女之力，是她与生俱来的天赋。

蔡兴坤则低头看着掌心里的黑子鸡。黑子鸡正眯着眼睛打盹，小小的身体微微起伏，看起来毫无威胁。但蔡兴坤知道，在这副蠢萌的外表之下，沉睡着连Phyrion都无法窥探的未知力量。

他不由自主地握紧了手掌，将黑子鸡轻轻拢在掌心。

黑子鸡''咕''了一声，没有反抗，只是把脑袋往他的掌心里缩了缩，找了个更舒服的姿势继续睡。

``所以，''Phyrion总结道，``你们拥有得天独厚的优势。神级扽魂给了你们极高的起点，但这并不意味着你们可以懈怠。起点高，只是意味着你们能走得更快。能走多远，还是要看你们自己的扽力修炼。''

``记住------魂力决定爆发的上限，扽力决定你的持续力。只有这两者协调发展，才能突破境界。而你们的神级扽魂，就是你们扽力修炼的最大助力。好好利用这份天赋。''

Phyrion站起身来，走到魂力黑板前。

他用指尖在黑板上画了两个圆------一个黑色，一个白色------它们相互交叠，形成了一个太极的图案。

``你们知道阴阳互补双魂核吗？''

唐三的身体微微一震。他当然知道------那是魂师修炼中最高深的境界之一。

``阴阳互补双魂核，要求魂师在体内同时凝聚两颗属性相反但又相互补充的魂核。这两颗魂核之间的平衡极其微妙------稍有不慎，就会崩溃。''

Phyrion的目光扫过众人。

``而要维持这种平衡，仅仅靠魂力是不够的。你需要极其强大的扽力------强大到能够同时'拉住'两颗属性完全相反的魂核，让它们在你的体内长期共存而不崩溃。''

``这就好比扽西------你用灵魂同时分析两种差异巨大的记号。而阴阳互补双魂核，就是用灵魂同时分析两种差异最大的力量------阴与阳。''

``所以，阴阳互补双魂核对扽力的要求极高。扽力不够的人，根本无法凝聚双魂核------即便勉强凝聚了，也会在短时间内崩溃。''

``这也是为什么，我要求你们在学习记号的同时，必须同步修炼扽力。因为你们将来要走的路，不仅仅是记号的推演，更是扽力的锤炼。''

Phyrion收回了手，转过身来，目光落在了果糕小队众人的身上。

``我听说，你们之前和梅天狸长老切磋过。''

众人的表情瞬间变得微妙起来。那场切磋给他们留下了深刻的印象------梅天狸长老的实力远超他们的想象。

``那一战，''Phyrion的语气平淡，但每个字都重若千钧，``她发挥的力量，恐怕连十分之一都不到。''

唐三的瞳孔微微收缩。

十分之一都不到？

他们在那场切磋中已经拼尽了全力，而梅天狸长老只是轻描淡写地化解了他们所有的攻击。当时他们已经觉得梅天狸长老深不可测了，但现在Phyrion告诉他们------那连十分之一都不到？

``我虽然境界比她高，''Phyrion继续说道，语气中没有丝毫的傲慢，只有一种平静的陈述，``但是她的扽力，强于我。''

这句话在修炼堂里引起了短暂的沉默。

Phyrion------\allowbreak{}\NovelInline{I04682}{Y(1,3,10)}\allowbreak{} 的封号斗罗，扽魂殿的长老之一------亲口承认梅天狸的扽力强于自己。

``她的扽力修炼方式和我不一样。''Phyrion解释道，``我依靠的是魂灵融合带来的扽力飞跃。而她------她走的是纯粹的记号扽西之路。她花了数十年的时间，不断地分析差异越来越大的记号对，将扽力锤炼到了一个极其恐怖的程度。''

``如果说我的扽力是一条宽阔的河流，那她的扽力就是一片深不见底的湖泊。表面上波澜不惊，但底下蕴藏着无穷无尽的力量。''

``所以你们要记住------境界高不代表一切。扽力，往往才是决定一个魂师真正底蕴的关键。''

Phyrion的目光变得悠远起来。

他转过身，望向修炼堂的穹顶，仿佛穿透了石壁，望向了某个遥远的地方。

``而你们见过的大长老------test\_alpha0------''

他的声音变得格外低沉。

``他的实力，深不可测。底蕴之深厚，极其罕见。''

``我父亲卢承业当年游历大陆时，在那座无名小岛上见过一位老修士。我父亲说，那位老修士的气息遍布整座岛屿，实力深不可测。''

Phyrion停顿了一下。

``我父亲见到的那位老修士，已经是传说中的存在了。但test\_alpha0------''他的声音几乎低不可闻。

``恐怕和我父亲见到的那位老修士一样强大。甚至------''

他又停顿了一下，像是在斟酌措辞。

``甚至更甚一筹。''

修炼堂里安静得能听到灵气流动的声音。

``更甚一筹''------这四个字的分量，在座的每一个人都能感受到。那位老修士已经是Phyrion的父亲卢承业一生所见最强之人，而test\_alpha0可能比那位老修士还要强。

``我虽然没有亲眼见过test\_alpha0出手，''Phyrion缓缓说道，``但从他偶尔流露出的气息来判断------他的扽力，已经不是'深厚'两个字能形容的了。那是一种\ldots\ldots{}''

他搜索了很久，才找到了一个词。``------无穷。''

``他的扽力像是无穷无尽一般，无论你怎么感知，都探不到底。我怀疑，他可能已经扽西过差异最大的记号对------那种差异，已经不是投影记号和反射记号之间的差异能比的了。''

``那种差异，是不同体系之间的差异。是投影体系和反射体系、反射体系和bms体系、bms体系和更高层级Y,MN体系之间的差异。''

``如果真是这样的话------''

Phyrion没有继续说下去。

但所有人都明白了他的意思。

如果test\_alpha0真的扽西过不同体系之间的记号对，那他的扽力已经达到了一个凡人无法想象的高度。

Phyrion收回了目光，重新看向众人。

``所以我今天跟你们讲这些，不是为了打击你们的信心，而是为了让你们明确修炼的方向。''

``从今天开始，你们的修炼将分为两条线。''

``第一条线------记号的学习和推演。这是你们已经在做的事情。投影记号、反射记号、Veblen记号------每一种记号的掌握，都会提升你们的认知深度，从而间接提升扽力。''

``第二条线------扽西的实战训练。你们将开始学习如何同时分析两种不同的记号，用灵魂去拆解它们、透视它们，在张力中锤炼扽力。''

``一开始，你们会从差异较小的记号对开始。比如投影记号和它的变体。然后逐步过渡到差异更大的记号对------投影记号和反射记号，反射记号和Veblen记号。''

``每成功扽西一次，你们的扽力就会获得一次显著的提升。而差异越大的记号对，提升的幅度就越恐怖。''

``但我要警告你们------''

Phyrion的语气变得异常严肃。

``扽西是有风险的。如果两种记号之间的差异超过了你灵魂的承受能力，分析就会失败。失败的后果，轻则神识受损，重则灵魂崩溃。''

``所以，永远不要贪功冒进。在你的扽力没有达到相应的水平之前，不要去分析差异过大的记号对。''

``循序渐进，方为正道。''

Phyrion说完了。

他站起身来，拍了拍白色母猪的脑袋。白色母猪''哼''了一声，慢悠悠地走到了角落里，蜷缩成一团，闭上了粉红色的圆眼睛。

``今天就到这里。''Phyrion说，``回去好好休息。明天开始，扽西训练。''

他转身向门口走去。

走了几步，他忽然停了下来。

``对了。''他没有回头，声音淡淡的，``明天扽西训练的第一课------投影记号与反射记号的扽西。这两种记号之间的差异，足够让你们喝一壶的了。''

``做好心理准备。''

石门缓缓打开，又缓缓关闭。

Phyrion的身影消失在了夜色之中。

\begin{center}\rule{0.5\linewidth}{0.5pt}\end{center}

修炼堂里，果糕小队众人面面相觑。

沉默了许久，霍雨浩率先打破了沉默。

``\ldots\ldots 我突然觉得，之前的推演训练，简直就是在度假。''

众人苦笑。

卢本韦低头看了看自己的母猪扽魂，喃喃自语：``比Phyrion老师的还高\ldots\ldots 我这母猪，居然这么厉害？''

他身边的母猪''哼''了一声，用鼻子拱了拱他的下巴，像是在说''那当然''。

蔡须坤摊开手掌，看着掌心流转的金色光芒。篮球扽魂的轮廓在光芒中若隐若现，而在那层灵动之下，是一股更加深沉的力量------帝皇瑞兽的本源处女之力。她第一次如此清晰地感受到了这股力量的存在。

``原来如此\ldots\ldots{}''她低声说道，``难怪每次扽西的时候，我总能感觉到一股额外的力量在支撑着我。那不是篮球的力量，是\ldots\ldots 我本身的力量。''

蔡兴坤没有说话。

蔡兴坤的脑海中回响着Phyrion的话------``黑子鸡体内封印着一股极其庞大的力量''、``来自外部的、被刻意封入的力量''、``在你没有足够的能力驾驭它之前，强行解封只会带来灾难''。

``好好修炼。''Phyrion的声音在他耳边回响，``不是为了去解封，而是为了在解封的那一天，你有能力承受住那股力量。''

蔡兴坤缓缓点了点头。

唐三站起身来，走到窗前。窗外，月光如水，洒在修炼堂前的石板路上。远处的山峦在月色下若隐若现，像是沉睡的巨人。

他想起了Phyrion说的那句话------``魂力决定爆发的上限，扽力决定你的持续力。''

他又想起了test\_alpha0------那个深不可测的大长老。如果test\_alpha0的扽力真的无穷无尽，那他的持续力该有多么恐怖？

``路还很长啊。''唐三低声说道。

身后传来了蔡须坤的声音：``再长也得走。''

唐三回头看去，只见蔡须坤正站在石台旁，金色的双瞳在月光下熠熠生辉。她的嘴角微微上扬，带着一种不服输的倔强。

``Phyrion老师说，差异越大的记号对，扽西之后的提升越恐怖。''蔡须坤说，``那我就分析差异最大的。''

小舞伸了个懒腰，打了个哈欠：``你们慢慢聊，我先去睡觉了。明天还要被折磨呢。''

奥斯卡拍了拍肚子：``我也去睡了。今晚吃的烤肉真香，Phyrion老师的手艺不比食堂差。''

众人纷纷散去。

蔡兴坤最后一个离开。

他走在月光下，影子被拉得很长很长。肩膀上的黑子鸡在睡梦中翻了个身，爪子无意识地收紧了一些。

就在蔡兴坤转身的那一瞬间，黑子鸡的羽毛深处，有什么东西闪烁了一下------极其微弱，极其短暂，像是深海中一条发光的鱼，一闪即逝。

但没有人注意到。

月光下，蔡兴坤的背影渐渐远去。肩膀上的黑子鸡继续呼呼大睡，浑然不知自己体内沉睡着怎样的秘密。

只有夜风知道。

\chapter{第八十四章：T-1Y}\label{ux7b2cux516bux5341ux56dbux7ae0t-1y}

众人将这句话默默记在心中。

Phyrion沉默了片刻，然后深吸了一口气。

``现在，我要传授你们另一种功法。''

他的声音忽然变得格外庄重，像是即将揭开某个尘封已久的秘密。

``它的名字叫做------\textbf{T-1Y}。''

``全称：\textbf{拓展-1Y序列}。''

修炼堂里的空气仿佛凝固了一瞬。

``Y序列无疑是极其强大的。''Phyrion缓缓说道，``但是，掌握-1Y的拓展，才是根基。没有-1Y的拓展，Y序列就像一座没有地基的高塔------看似宏伟，实则脆弱。''

他转身面向黑板，提笔写下了五条规则。每一条规则落笔时，魂力黑板都会微微震颤，像是在回应这些规则的分量。

\textbf{一、初始值是 \allowbreak{}\NovelInline{I04683}{1}\allowbreak{}。}

\textbf{二、\allowbreak{}\NovelInline{I04684}{\#, 1 = \# + 1}\allowbreak{}。}

\textbf{三、\allowbreak{}\NovelInline{I04685}{G, r, B, l = G, r, B, l-1, B, l-1, B, l-1, \dots}\allowbreak{}}

\textbf{四、\allowbreak{}\NovelInline{I04686}{\#, \alpha = \sup\{\#, \alpha[n]\}}\allowbreak{}，如果 \allowbreak{}\NovelInline{I04687}{\alpha}\allowbreak{} 是极限序数。}

\textbf{五、\allowbreak{}\NovelInline{I04688}{(\#)}\allowbreak{} 本身也可以写进表达式，用于简写。比如 \allowbreak{}\NovelInline{I04689}{(1,2)=\omega}\allowbreak{}。}

五条规则写完，Phyrion放下笔，转过身来。

``接下来，你们自主扽西，提升扽力。''

\begin{center}\rule{0.5\linewidth}{0.5pt}\end{center}

果糕小队众人立刻进入了状态。

扽西------分析。用灵魂去拆解、透视这些规则的内在结构。蔡须坤的金色双瞳微微发光，帝皇瑞兽的本源处女之力在她体内缓缓流转，让她的思维变得异常清晰。蔡兴坤肩膀上的黑子鸡歪着脑袋，圆溜溜的眼睛盯着黑板上的规则，似乎在以自己的方式理解着这一切。卢本韦的母猪扽魂趴在地上，粉红色的鼻子有节奏地翕动着，像是在嗅探规则中隐藏的气息。

小队很快发现了规律。

蔡须坤率先开口，声音清亮：``从最简单的开始------''

\NovelDisplay{M00952R01}{%
1 \sim 1%
}

``这是初始值。''

卢本韦接道：

\NovelDisplay{M00953R01}{%
1, 1 \sim 2%
}

\NovelDisplay{M00954R01}{%
1, 1, 1 \sim 3%
}

``规则二在起作用------每加一个 \allowbreak{}\NovelInline{I04690}{1}\allowbreak{}，就加一。''

蔡兴坤忽然插嘴：``那 \allowbreak{}\NovelInline{I04691}{1, 2}\allowbreak{} 呢？''

众人沉默了片刻，然后异口同声：

\NovelDisplay{M00955R01}{%
1, 2 \sim \omega%
}

``因为规则三------\allowbreak{}\NovelInline{I04692}{1, 2}\allowbreak{} 展开成 \allowbreak{}\NovelInline{I04693}{1, 1, 1, 1, \dots}\allowbreak{}，取上确界，就是 \allowbreak{}\NovelInline{I04694}{\omega}\allowbreak{}。''

小舞兴奋地拍了一下手：``那 \allowbreak{}\NovelInline{I04695}{1, 2, 1, 2}\allowbreak{} 就是------''

\NovelDisplay{M00956R01}{%
1, 2, 1, 2 \sim \omega \times 2%
}

``对！''唐三点头，``两个 \allowbreak{}\NovelInline{I04696}{\omega}\allowbreak{} 的叠加。''

蔡须坤的眼睛越来越亮：``那 \allowbreak{}\NovelInline{I04697}{1, 2, 2}\allowbreak{} 呢？按照规则三，\allowbreak{}\NovelInline{I04698}{1, 2, 2}\allowbreak{} 展开成 \allowbreak{}\NovelInline{I04699}{1, 2, 1, 2, 1, 2, \dots}\allowbreak{}------''

\NovelDisplay{M00957R01}{%
1, 2, 2 \sim \omega^2%
}

``漂亮！''Phyrion在角落里淡淡地评价了一句。

蔡兴坤猛地站了起来：``那 \allowbreak{}\NovelInline{I04700}{1, 3}\allowbreak{} 就是------''

\NovelDisplay{M00958R01}{%
1, 3 \sim \omega^\omega%
}

``因为 \allowbreak{}\NovelInline{I04701}{1, 3}\allowbreak{} 展开成 \allowbreak{}\NovelInline{I04702}{1, 2, 2, 2, 2, \dots}\allowbreak{}，而每一个 \allowbreak{}\NovelInline{I04703}{2}\allowbreak{} 都对应一个 \allowbreak{}\NovelInline{I04704}{\omega}\allowbreak{} 的指数叠加------''

``等等。''唐三忽然说道，他的目光变得锐利起来，``这一点很相似于 \allowbreak{}\NovelInline{I04705}{2\text{-}3}\allowbreak{} 展开成 \allowbreak{}\NovelInline{I04706}{(3\cap 1-)^n 3}\allowbreak{}。''

众人一愣。

``我们学过反射的规则------''唐三的声音变得急促起来，``反射的规则就是T-1Y序列的规则！这一点就是 \allowbreak{}\NovelInline{I04707}{1,2,3}\allowbreak{} 的简写。\allowbreak{}\NovelInline{I04708}{1,3}\allowbreak{} 展开成 \allowbreak{}\NovelInline{I04709}{1,2,2,2,2,\dots}\allowbreak{}------而这个对应是非常巧妙的。''

``这个对应的序数是 \allowbreak{}\NovelInline{I04710}{\omega^{\omega+1}}\allowbreak{}。''蔡须坤接过话头，金色的双瞳中闪烁着兴奋的光芒。

``再举一个例子。''唐三越说越快，'' \allowbreak{}\NovelInline{I04711}{2\text{-}2}\allowbreak{} 和 \allowbreak{}\NovelInline{I04712}{1\text{-}2\text{-}2}\allowbreak{} 对应的就是 \allowbreak{}\NovelInline{I04713}{1, 2, 2, 1, 2, 2}\allowbreak{} 展开成 \allowbreak{}\NovelInline{I04714}{1, 2, 2, 1, 2, 1, 2, 1, 2, \dots}\allowbreak{}，对应 \allowbreak{}\NovelInline{I04715}{(2\cap 1-)^n 2\cap 1\text{-}2\text{-}2}\allowbreak{}。''

小队在推演的过程中，不知不觉发现------对于反射的理解，又更进了一层。

那些曾经晦涩难懂的概念，那些在反射体系中看似高不可攀的结构------比方说最难的 \allowbreak{}\NovelInline{I04716}{M,N}\allowbreak{} 这样的东西------总能在T-1Y序列里找到类似物。就像是在迷雾中行走，忽然发现脚下的路和一个早已熟悉的迷宫完全重合。

蔡须坤深吸一口气：``在反射里面，\allowbreak{}\NovelInline{I04717}{\psi(\psi_a(\varepsilon_{\Omega_{a+1}+1}))=\Pi_\omega}\allowbreak{}。''

``而在这里------''

\NovelDisplay{M00959R01}{%
1, \omega \sim \varepsilon_0%
}

``但是这里有着反射的规律。因为反射序数用的是同样的模式，所以出现了相同的结果------这是正常的。''

Phyrion在角落里微微点头，没有说话。

小队继续推导。速度越来越快，像是齿轮咬合之后越来越顺畅的机械。

\NovelDisplay{M00960R01}{%
1, \omega, 2 \sim \varepsilon_0 \times \omega = \omega^{\varepsilon_0+1}%
}

\NovelDisplay{M00961R01}{%
1, \omega, 2, 3 \sim \varepsilon_0 \times \omega^\omega%
}

\NovelDisplay{M00962R01}{%
1, \omega, 2, \omega \sim \varepsilon_0^2%
}

蔡兴坤忽然喊道：``我发现了一个规律------序列加上 \allowbreak{}\NovelInline{I04718}{\omega}\allowbreak{}，就是下一个 \allowbreak{}\NovelInline{I04719}{\varepsilon}\allowbreak{} 序数！''

``不对，先别急。''唐三抬手制止了他，``先看这个------''

\NovelDisplay{M00963R01}{%
1, \omega, 3 \sim \omega^{\omega^{\varepsilon_0+1}}%
}

Phyrion这时候开口了：``你们推导的是对的。有时候，用 \allowbreak{}\NovelInline{I04720}{\omega^{\omega^{\varepsilon_0+1}}}\allowbreak{} 的写法会更好。这取决于你们想要做什么，目标有什么行为。''

他顿了顿，补充道：``\allowbreak{}\NovelInline{I04721}{1,\omega,3}\allowbreak{} 展开成 \allowbreak{}\NovelInline{I04722}{1,\omega,2,\omega,2,\omega,2,\omega,\dots}\allowbreak{}，恰好对应乘上 \allowbreak{}\NovelInline{I04723}{\omega}\allowbreak{}。''

众人恍然。

推导继续。

\NovelDisplay{M00964R01}{%
1, \omega, \omega \sim \varepsilon_1%
}

\NovelDisplay{M00965R01}{%
1, \omega+1 \sim \varepsilon_\omega%
}

霍雨浩忽然笑了：``这里，\allowbreak{}\NovelInline{I04724}{1,\omega+1 = 1,\omega,\omega,\omega,\omega,\dots}\allowbreak{}，这就不难理解了。''

``对。''唐三点头，``\allowbreak{}\NovelInline{I04725}{\omega+1}\allowbreak{} 作为参数，展开之后就是无限个 \allowbreak{}\NovelInline{I04726}{\omega}\allowbreak{} 的叠加，直接跳到 \allowbreak{}\NovelInline{I04727}{\varepsilon_\omega}\allowbreak{}。''

推导继续加速。

\NovelDisplay{M00966R01}{%
1, \omega+1, 2 \sim \omega^{\varepsilon_\omega+1}%
}

\NovelDisplay{M00967R01}{%
1, \omega+1, n \sim \omega^{\omega^{\cdots^{\varepsilon_\omega+1}}}%
}

蔡兴坤猛地拍了一下桌子：``规律确认了！序列加上 \allowbreak{}\NovelInline{I04728}{\omega}\allowbreak{}，就是下一个 \allowbreak{}\NovelInline{I04729}{\varepsilon}\allowbreak{} 序数！''

\NovelDisplay{M00968R01}{%
1, \omega+1, \omega \sim \varepsilon_{\omega+1}%
}

\NovelDisplay{M00969R01}{%
1, \omega+1, \omega, \omega+1 \sim \varepsilon_{\omega \times 2}%
}

唐三接道：``这里 \allowbreak{}\NovelInline{I04730}{1,\omega+1,\omega,\omega+1 \sim \varepsilon_{\omega \times 2} = 1,\omega+1,\omega,\omega,\omega,\dots}\allowbreak{}，所以说这里面的 \allowbreak{}\NovelInline{I04731}{\varepsilon_{\omega \times 2}}\allowbreak{} 就是这样来的。''

小舞忽然眼睛一亮：``我发现了！序列加上 \allowbreak{}\NovelInline{I04732}{\omega+1}\allowbreak{}，\allowbreak{}\NovelInline{I04733}{\varepsilon}\allowbreak{} 的下标就乘上 \allowbreak{}\NovelInline{I04734}{\omega}\allowbreak{}！''

\NovelDisplay{M00970R01}{%
1, \omega+1, \omega+1 \sim \varepsilon_{\omega^2}%
}

\NovelDisplay{M00971R01}{%
1, \omega+1, \omega+1, \omega \sim \varepsilon_{\omega^2+1}%
}

\NovelDisplay{M00972R01}{%
1, \omega+1, \omega+1, \omega, \omega+1 \sim \varepsilon_{\omega^2+\omega}%
}

小舞越说越兴奋，马尾辫在脑后甩来甩去：``看到了吗？每次加上 \allowbreak{}\NovelInline{I04735}{\omega+1}\allowbreak{}，下标就从 \allowbreak{}\NovelInline{I04736}{\omega^2}\allowbreak{} 跳到 \allowbreak{}\NovelInline{I04737}{\omega^2 \times \omega = \omega^3}\allowbreak{}------不对，是从 \allowbreak{}\NovelInline{I04738}{\omega^2+\omega}\allowbreak{} 开始继续叠加\ldots\ldots{}''

``别急。''蔡须坤抬手示意她冷静，``先看更大的。''

\NovelDisplay{M00973R01}{%
1, \omega+2 \sim \varepsilon_{\omega^\omega}%
}

\NovelDisplay{M00974R01}{%
1, \omega+2, \omega \sim \varepsilon_{\omega^\omega+1}%
}

\NovelDisplay{M00975R01}{%
1, \omega+2, \omega+1 \sim \varepsilon_{\omega^{\omega+1}}%
}

\NovelDisplay{M00976R01}{%
1, \omega+2, \omega+1, \omega+2 \sim \varepsilon_{\omega^{\omega \times 2}}%
}

卢本韦点了点头：``我看，到这里，规律非常清楚了。''

他的声音变得沉稳而有力，像是在宣读一份经过反复验证的实验报告。

``接下来------''

\NovelDisplay{M00977R01}{%
1, \omega+3 \sim \varepsilon_{\omega^{\omega^\omega}}%
}

卢本韦的声音在这里顿了一下。

\NovelDisplay{M00978R01}{%
1, \omega \times 2 \sim \varepsilon_{\varepsilon_0}%
}

修炼堂里忽然安静了。

\allowbreak{}\NovelInline{I04739}{\varepsilon_{\varepsilon_0}}\allowbreak{}。

这个序数的出现，像是一块石头投入了平静的湖面。

蔡兴坤的眼神忽然变得恍惚起来。

\allowbreak{}\NovelInline{I04740}{\varepsilon_{\varepsilon_0}}\allowbreak{}\ldots\ldots 当年，这个序数困住了牢大很长时间。牢大如果还在的话\ldots\ldots{}

她的思绪不由自主地飘向了远方。那些和牢大一起修炼的日子，那些在黑板前争论公式的深夜，那些失败后互相鼓励的清晨------一切都像是昨天发生的事。

``蔡兴坤？''小舞轻声唤道。

蔡兴坤回过神来，微微摇了摇头，将那股涌上心头的情绪压了下去。

``\ldots\ldots 没事。''她的声音很轻，但很坚定，``为了修炼，先把这件事放到一边。''

她深吸一口气，将注意力重新拉回到黑板上。金色的双瞳中，悲伤被一种更加坚韧的东西取代了。

推导继续。

\NovelDisplay{M00979R01}{%
1, \omega \times 2, \omega+1 \sim \varepsilon_{\omega^{\varepsilon_0+1}}%
}

``\allowbreak{}\NovelInline{I04741}{\omega+n}\allowbreak{} 对应 \allowbreak{}\NovelInline{I04742}{n}\allowbreak{} 层指数塔。''唐三解说道，``同样是 \allowbreak{}\NovelInline{I04743}{\omega}\allowbreak{} 为底数，不断往上爬升的指数塔。''

\NovelDisplay{M00980R01}{%
1, \omega \times 2, \omega \times 2 \sim \varepsilon_{\varepsilon_1}%
}

\NovelDisplay{M00981R01}{%
1, \omega \times 3 \sim \varepsilon_{\varepsilon_{\varepsilon_0}}%
}

众人的呼吸越来越急促。他们能感觉到，扽力在体内翻涌，像是即将沸腾的水。每一次新的序数被推导出来，他们的灵魂就会被拉伸一分，扽力就会增长一分。

然后------

\NovelDisplay{M00982R01}{%
1, \omega^2 \sim \zeta_0%
}

蔡须坤写下了最后这个等式，手指微微颤抖。

\allowbreak{}\NovelInline{I04744}{\zeta_0}\allowbreak{}。

\allowbreak{}\NovelInline{I04745}{\varepsilon}\allowbreak{} 的不动点。\allowbreak{}\NovelInline{I04746}{\varepsilon_{\varepsilon_{\varepsilon_{\cdots}}}}\allowbreak{}。

小队众人快速分析，规律是非常显然的------但正是这种''显然''，才最令人震撼。从 \allowbreak{}\NovelInline{I04747}{1}\allowbreak{} 到 \allowbreak{}\NovelInline{I04748}{\zeta_0}\allowbreak{}，每一步推导都环环相扣，每一个序数的跃迁都自然而然，像是宇宙本身的法则在引导着他们。

尽管他们很早就掌握了 \allowbreak{}\NovelInline{I04749}{\zeta_0}\allowbreak{}，但是对于接下来的突破，这是大有帮助的。

蔡须坤缓缓放下笔，转过身来看着Phyrion。

``Phyrion老师，''她说，``T-1Y序列的根基，我现在理解了。它不是Y序列的附属品------它是Y序列的地基。所有的序数结构，都可以从T-1Y的规则中生长出来。''

``继续。''

蔡须坤走到黑板前，提笔写下：

\NovelDisplay{M00983R01}{%
1, \omega^2, 2 \sim \omega^{\zeta_0+1}%
}

``从 \allowbreak{}\NovelInline{I04750}{\zeta_0}\allowbreak{} 出发，规律和之前完全一致。''她的声音清亮而笃定，笔尖在黑板上飞速移动。

\NovelDisplay{M00984R01}{%
1, \omega^2, 2, \omega^2 \sim \omega^{\zeta_0 \times 2}%
}

\NovelDisplay{M00985R01}{%
1, \omega^2, 3 \sim \omega^{\omega^{\zeta_0+1}}%
}

``参数 \allowbreak{}\NovelInline{I04751}{2}\allowbreak{} 对应指数塔的第一层，参数 \allowbreak{}\NovelInline{I04752}{3}\allowbreak{} 对应第二层------和之前 \allowbreak{}\NovelInline{I04753}{\omega}\allowbreak{} 作为参数时的行为完全一样。''

卢本韦接过话头：``那 \allowbreak{}\NovelInline{I04754}{\omega}\allowbreak{} 作为参数------''

\NovelDisplay{M00986R01}{%
1, \omega^2, \omega \sim \varepsilon_{\zeta_0+1}%
}

``\allowbreak{}\NovelInline{I04755}{\zeta_0}\allowbreak{} 的下一个 \allowbreak{}\NovelInline{I04756}{\varepsilon}\allowbreak{} 不动点。''

蔡兴坤紧跟着写下：

\NovelDisplay{M00987R01}{%
1, \omega^2, \omega, \omega^2 \sim \varepsilon_{\zeta_0 \times 2}%
}

``这一点，和之前的规律一样，没什么区别。''他说。

小队众人的推导速度越来越快，像是已经找到了节奏的乐手，每个人都能精准地接上下一个人的旋律。

\NovelDisplay{M00988R01}{%
1, \omega^2, \omega+1 \sim \varepsilon_{\omega^{\zeta_0+1}}%
}

诺兰忽然开口了，声音中带着一丝感慨：``感觉距离极限很快了。''

没有人接话。但所有人都知道他说的是什么------T-1Y序列的极限。那个他们尚未触及的、隐藏在规则深处的终极序数。

推导继续。

\NovelDisplay{M00989R01}{%
1, \omega^2, \omega \times 2 \sim \varepsilon_{\varepsilon_{\zeta_0+1}}%
}

小舞忽然拍了一下手：``现在 \allowbreak{}\NovelInline{I04757}{\omega \times 2}\allowbreak{} 代表一层 \allowbreak{}\NovelInline{I04758}{\varepsilon}\allowbreak{} 序数！''

她的眼睛亮晶晶的，像是发现了什么了不起的宝藏。

``对。''唐三点头，``\allowbreak{}\NovelInline{I04759}{\omega \times 2}\allowbreak{} 作为参数，就是在 \allowbreak{}\NovelInline{I04760}{\varepsilon}\allowbreak{} 的层面再做一次递归------和之前 \allowbreak{}\NovelInline{I04761}{\omega+1}\allowbreak{} 在 \allowbreak{}\NovelInline{I04762}{\omega}\allowbreak{} 层面做的事情完全对应。''

蔡须坤的笔尖在黑板上重重地点了一下，写下了一个让所有人都屏住呼吸的等式：

\NovelDisplay{M00990R01}{%
1, \omega^2, \omega^2 \sim \zeta_1%
}

修炼堂里安静了一瞬。

\allowbreak{}\NovelInline{I04763}{\zeta_1}\allowbreak{}。

\allowbreak{}\NovelInline{I04764}{\zeta_0}\allowbreak{} 的下一个不动点。

``两次用 \allowbreak{}\NovelInline{I04765}{\omega^2}\allowbreak{}，代表 \allowbreak{}\NovelInline{I04766}{\varepsilon}\allowbreak{} 下标不动点取两次。''蔡须坤的声音平静，但她握笔的手指微微发白------那是扽力在体内剧烈涌动的表现。

蔡兴坤接道：

\NovelDisplay{M00991R01}{%
1, \omega^2+1 \sim \zeta_\omega%
}

``之后 \allowbreak{}\NovelInline{I04767}{1,\omega^2+1 = 1,\omega^2,\omega^2,\omega^2,\dots}\allowbreak{}''他解释道，``\allowbreak{}\NovelInline{I04768}{\omega^2+1}\allowbreak{} 展开成无限个 \allowbreak{}\NovelInline{I04769}{\omega^2}\allowbreak{} 的叠加，直接跳到 \allowbreak{}\NovelInline{I04770}{\zeta_\omega}\allowbreak{}。''

唐三看着黑板上越来越密集的公式，忽然笑了：'' \allowbreak{}\NovelInline{I04771}{+n}\allowbreak{} 是 \allowbreak{}\NovelInline{I04772}{n}\allowbreak{} 层指数塔的性质不变哎。''

他的语气中带着一丝如释重负的轻松------在无穷序数的世界里，``不变''就是最好的消息。

\NovelDisplay{M00992R01}{%
1, \omega^2+\omega \times 2 \sim \zeta_{\varepsilon_{\varepsilon_0}}%
}

\NovelDisplay{M00993R01}{%
1, \omega^2 \times 2 \sim \zeta_{\zeta_0}%
}

蔡须坤的笔尖停了一下。

\allowbreak{}\NovelInline{I04773}{\zeta_{\zeta_0}}\allowbreak{}------\allowbreak{}\NovelInline{I04774}{\zeta}\allowbreak{} 函数的不动点。她深吸一口气，继续写。

\NovelDisplay{M00994R01}{%
1, \omega^3 \sim \eta_0%
}

\allowbreak{}\NovelInline{I04775}{\eta_0}\allowbreak{}。

一个新的层级。

小舞倒吸了一口凉气：``我们已经到 \allowbreak{}\NovelInline{I04776}{\eta_0}\allowbreak{} 了\ldots\ldots{}''

但推导没有停下。

\NovelDisplay{M00995R01}{%
1, \omega^\omega \sim \varphi(\omega, 0)%
}

蔡兴坤瞪大了眼睛：``天啊，我们推导得好快！''

他的声音里带着一种近乎不敢相信的颤抖。从 \allowbreak{}\NovelInline{I04777}{1}\allowbreak{} 到 \allowbreak{}\NovelInline{I04778}{\varphi(\omega, 0)}\allowbreak{}------这在以前是不可想象的。他们只用了两天的扽西训练，就走完了别人可能需要数年才能走完的路。

但蔡兴坤的心里，除了震撼之外，还有一丝隐隐的酸楚。

如果牢大还在，看到这一幕，该有多高兴啊。

蔡须坤没有停下来。她的金色双瞳中燃烧着一种近乎执拗的光芒，帝皇瑞兽的本源处女之力在她体内翻涌，将她的思维推向了极限。

\NovelDisplay{M00996R01}{%
1, \omega^\alpha \sim \varphi(\alpha, 0)%
}

``这是一个通用的规律。''她的声音变得低沉而有力，``\allowbreak{}\NovelInline{I04779}{\omega}\allowbreak{} 的任意次幂作为参数，都对应 \allowbreak{}\NovelInline{I04780}{\varphi}\allowbreak{} 函数的第一个参数。''

\NovelDisplay{M00997R01}{%
1, \varepsilon_0 \sim \varphi(\varphi(1, 0), 0)%
}

唐三低声念出了这个等式，声音中带着一丝敬畏。

\allowbreak{}\NovelInline{I04781}{\varepsilon_0}\allowbreak{} 作为参数------这已经不是在 \allowbreak{}\NovelInline{I04782}{\omega}\allowbreak{} 的层面上操作了，而是在 \allowbreak{}\NovelInline{I04783}{\varepsilon}\allowbreak{} 的层面上。T-1Y序列的表达能力，远超他们的想象。

然后，蔡须坤写下了最后一行。

她的笔尖在黑板上停顿了片刻，像是在积蓄最后的力量。然后，她重重地写下了：

\NovelDisplay{M00998R01}{%
1, \Omega \sim \varphi(1, 0, 0) = \Gamma_0%
}

\allowbreak{}\NovelInline{I04784}{\Gamma_0}\allowbreak{}。

费弗巴赫-舒滕巴赫序数。

\allowbreak{}\NovelInline{I04785}{\varphi(1,0,0)}\allowbreak{}------第一个使得 \allowbreak{}\NovelInline{I04786}{\alpha = \varphi(\alpha, 0)}\allowbreak{} 成立的序数。

整个修炼堂陷入了长久的沉默。

\begin{center}\rule{0.5\linewidth}{0.5pt}\end{center}

黑板上密密麻麻地写满了公式，从最初的 \allowbreak{}\NovelInline{I04787}{1 \sim 1}\allowbreak{} 到最终的 \allowbreak{}\NovelInline{I04788}{1, \Omega \sim \Gamma_0}\allowbreak{}，每一步推导都清晰可见。那是一条从最朴素的自然数出发，一路攀升至无穷之巅的道路。

蔡须坤放下笔，退后一步，看着满墙的公式。

她的手在微微发抖。不是因为疲惫------虽然确实很累------而是因为扽力。一整天的扽西训练，让她的扽力获得了前所未有的飞跃。帝皇瑞兽的本源处女之力在她体内翻涌，像是即将破茧的蝴蝶，拼命地挣动着翅膀。

``最终推导完毕。''她的声音有些沙哑，``到达了极限表达式。''

Phyrion从角落里走了出来，看着黑板上的公式，沉默了许久。

``不错。''他最终只说了两个字。

但所有人都知道，这两个字的分量有多重。

\begin{center}\rule{0.5\linewidth}{0.5pt}\end{center}

好累啊。

又是很累的一天。

众人横七竖八地靠在石椅上，连说话的力气都快没了。扽力的消耗远比魂力更加折磨人------它不像魂力那样可以通过冥想来快速恢复，扽力的恢复需要时间，需要沉淀，需要灵魂自己去慢慢修补那些被拉伸过的裂痕。

蔡兴坤靠在墙角，黑子鸡蹲在他的膝盖上，歪着脑袋看着他。

蔡须坤坐在石台边，双手撑着台面，低着头，金色的发丝遮住了她的脸。

修炼堂里安静得只能听到灵气流动的声音。

然后，蔡兴坤开口了。

``你还记得吗？''她的声音很轻，像是怕惊碎了什么。

蔡须坤没有抬头，但她知道蔡兴坤在问什么。

``记得。''

``那天\ldots\ldots{}''蔡兴坤的声音开始发颤，``牢大站在我们前面，背对着我们。她的背影好高好高，挡住了所有的攻击。''

蔡须坤的手指在石台上微微收紧。

``她说------`你们往后退。'\,''蔡兴坤的声音越来越低，``\,`我给你们争取时间。'\,''

``然后她\ldots\ldots{}''

她没有说下去。

她说不下去了。

蔡须坤终于抬起了头。金色的双瞳中，泪水无声地滑落。

``她回头看了我们一眼。''蔡须坤的声音也在颤抖，``就一眼。然后她笑了。''

``她笑了。''

``在那种情况下，她居然笑了。''

蔡兴坤闭上了眼睛。泪水从她的眼角滑落，滴在黑子鸡的羽毛上。黑子鸡''咕''了一声，用脑袋蹭了蹭他的手指，像是在安慰她。

``她把自己的所有力量都爆发了出来。''蔡须坤的声音断断续续，``为了给我们争取撤退的时间\ldots\ldots 她一个人挡住了\ldots\ldots{}''

她咬住了嘴唇，不让自己哭出声来。

``她说过，她最大的遗憾，就是没能看到我们走到今天这一步。''

蔡兴坤猛地捂住了脸，肩膀剧烈地颤抖着。

``他要是能看到\ldots\ldots 我们推导到了 \allowbreak{}\NovelInline{I04789}{\Gamma_0}\allowbreak{} \ldots\ldots{}''他的声音从指缝间挤出来，破碎而嘶哑，``他该有多高兴啊\ldots\ldots{}''

蔡须坤站起身来，走到蔡兴坤身边，蹲下身，伸手抱住了她。

``他会看到的。''他的声音很轻，但很坚定，``他一直在看着我们。''

``每一次我们推导出一个新的序数，每一次我们的扽力获得新的突破------他都在看着。''

``他牺牲了自己，不是为了让我们哭泣，而是为了让我们走到更远的地方。''

蔡兴坤的哭声渐渐低了下去。

蔡须坤站起身来，走到黑板前。

她伸出手，轻轻触碰了一下那个最终等式------\allowbreak{}\NovelInline{I04790}{1, \Omega \sim \Gamma_0}\allowbreak{}。指尖传来一阵微弱的温热，像是黑板上残留的扽力在回应她的触碰。

``牢大，''她低声说道，``我们扽完了，极限 \allowbreak{}\NovelInline{I04791}{\Gamma_0}\allowbreak{} 。''

``接下来，我们会走得更远。''

``替你，也替我们自己。''

窗外，夕阳已经沉入了山脊之下。最后一缕余晖透过窗棂，洒在黑板上的公式上，将那些密密麻麻的符号染成了金红色。

Phyrion站在门口，看着这一幕，没有说话。

白色母猪趴在他的脚边，粉红色的圆眼睛半睁半闭，像是在打盹，又像是在沉思。

``走吧。''Phyrion终于开口了，声音比平时柔和了许多，``回去好好休息。''

``明天------''

他停顿了一下。

``明天，我们开始真正的扽西训练。''

众人缓缓起身，拖着疲惫的身体走出了修炼堂。

\chapter{第八十五章：超限反射的阶梯}\label{ux7b2cux516bux5341ux4e94ux7ae0ux8d85ux9650ux53cdux5c04ux7684ux9636ux68af}

清晨的修炼堂里，魂力晶石散发着柔和而稳定的光芒。

果糕小队众人端坐在石椅上，目光齐聚在正前方的魂力黑板前。经过昨日的休整，他们的精神状态恢复了不少，但眼底深处仍残留着对未知领域的敬畏与渴望。

Phyrion站在黑板中央，白色母猪安静地趴在他脚边。他的助手KNK(猫娘)蹲在一旁的高脚椅上，尾巴轻轻摆动，猫耳微微抖动，似乎在期待着什么。

``昨天你们到达了 \allowbreak{}\NovelInline{I04792}{\Gamma_0}\allowbreak{}。''Phyrion的声音平静地响起，``但那只是T-1Y序列的极限。今天，我们要换一个视角------用投影记号重新审视从 \allowbreak{}\NovelInline{I04793}{\Pi_\omega}\allowbreak{} 开始的所有序数结构。''

他提笔在黑板上写下一行字：``左侧是 \allowbreak{}\NovelInline{I04794}{\Pi}\allowbreak{} 记法，右侧是对应的投影表达式。这张表很长，但每一行都不是孤立的------它们是一条完整的阶梯。''

随后，他写下了第一组公式：

\NovelDisplay{M00999R01}{%
\Pi_\omega \sim \psi_a(\varepsilon_{T+1})%
}

``起点，\allowbreak{}\NovelInline{I04795}{Ω_{a+1}=T}\allowbreak{}。''他说，这一点，在T-1Y里面是\allowbreak{}\NovelInline{I04796}{ω}\allowbreak{}，则在这里面是\allowbreak{}\NovelInline{I04797}{T}\allowbreak{},这是一个简写

蔡须坤的金色双瞳微微发亮：``\allowbreak{}\NovelInline{I04798}{\Pi_\omega}\allowbreak{} 对应 \allowbreak{}\NovelInline{I04799}{\psi_a(\varepsilon_{T+1})}\allowbreak{}------这是 \allowbreak{}\NovelInline{I04800}{\varepsilon}\allowbreak{} 函数在 \allowbreak{}\NovelInline{I04801}{T+1}\allowbreak{} 处的取值。''

``对。''Phyrion点头，笔尖不停。

\NovelDisplay{M01000R01}{%
2\text{nd}\ \Pi_\omega \sim \psi_a(\varepsilon_{T+1}\times 2)%
}

\NovelDisplay{M01001R01}{%
\Pi_1\text{-}\Pi_\omega \sim \psi_a(\varepsilon_{T+1}\times\omega)%
}

\NovelDisplay{M01002R01}{%
(\Pi_1\text{-})^{(1,0)}\text{-}\Pi_\omega \sim \psi_a(\varepsilon_{T+1}\times a)%
}

``\allowbreak{}\NovelInline{I04802}{2\text{nd}\ \Pi_\omega}\allowbreak{} 是两倍，\allowbreak{}\NovelInline{I04803}{\Pi_1\text{-}\Pi_\omega}\allowbreak{} 是 \allowbreak{}\NovelInline{I04804}{\omega}\allowbreak{} 倍，\allowbreak{}\NovelInline{I04805}{(\Pi_1\text{-})^{(1,0)}\text{-}\Pi_\omega}\allowbreak{} 是不动点倍。''他解释道。

卢本韦盯着黑板，眉头微蹙：``所以 \allowbreak{}\NovelInline{I04806}{\Pi_1\text{-}}\allowbreak{} 这个前缀，本质上就是在 \allowbreak{}\NovelInline{I04807}{\Pi_\omega}\allowbreak{} 上做\allowbreak{}\NovelInline{I04808}{n\text{th}}\allowbreak{}递归？''

``没错。''Phyrion的目光变得锐利，``接下来，指数塔开始攀升,你们学过T-1Y，就知道这里面的规律了''

\NovelDisplay{M01003R01}{%
\Pi_2\ \Pi_1\text{-}\Pi_\omega \sim \psi_a(T^{\varepsilon_{T+1}+1})%
}

\NovelDisplay{M01004R01}{%
\Pi_2\ \Pi_1\text{-}\Pi_2\ \Pi_1\text{-}\Pi_\omega \sim \psi_a(T^{\varepsilon_{T+1}+2})%
}

\NovelDisplay{M01005R01}{%
(\Pi_2\ \Pi_1\text{-})^\omega\text{-}\Pi_\omega \sim \psi_a(T^{\varepsilon_{T+1}+\omega})%
}

\NovelDisplay{M01006R01}{%
(\Pi_2\ \Pi_1\text{-})^{(1,0)}\text{-}\Pi_\omega \sim \psi_a(T^{\varepsilon_{T+1}+a})%
}

蔡兴坤的眼睛亮了起来：``\allowbreak{}\NovelInline{I04809}{\Pi_2\ \Pi_1\text{-}}\allowbreak{} 对应的是 \allowbreak{}\NovelInline{I04810}{\times T}\allowbreak{}！指数部分从 \allowbreak{}\NovelInline{I04811}{+1}\allowbreak{} 到 \allowbreak{}\NovelInline{I04812}{+2}\allowbreak{}，再到 \allowbreak{}\NovelInline{I04813}{+\omega}\allowbreak{}、\allowbreak{}\NovelInline{I04814}{+a}\allowbreak{}------它们的展开T-1Y序列里的规则完全一致,仅仅是倒过来！''

``很好。''Phyrion点头，``继续。''

\NovelDisplay{M01007R01}{%
\Pi_2\text{-}\Pi_2\ \Pi_1\text{-}\Pi_\omega \sim \psi_a(T^{\varepsilon_{T+1}+T})%
}

\NovelDisplay{M01008R01}{%
\Pi_2\text{-}\Pi_2\text{-}\Pi_2\ \Pi_1\text{-}\Pi_\omega \sim \psi_a(T^{\varepsilon_{T+1}+T^2})%
}

\NovelDisplay{M01009R01}{%
(\Pi_2\text{-})^\omega\ \Pi_1\text{-}\Pi_\omega \sim \psi_a(T^{\varepsilon_{T+1}+T^\omega})%
}

\NovelDisplay{M01010R01}{%
(\Pi_2\text{-})^{1,0}\ \Pi_1\text{-}\Pi_\omega \sim \psi_a(T^{\varepsilon_{T+1}+T^a})%
}

唐三低声念道：``\allowbreak{}\NovelInline{I04815}{\Pi_2\text{-}}\allowbreak{} 前缀让指数部分出现了 \allowbreak{}\NovelInline{I04816}{T}\allowbreak{} 的幂次\ldots\ldots{}\allowbreak{}\NovelInline{I04817}{T}\allowbreak{}、\allowbreak{}\NovelInline{I04818}{T^2}\allowbreak{}、\allowbreak{}\NovelInline{I04819}{T^\omega}\allowbreak{}、\allowbreak{}\NovelInline{I04820}{T^a}\allowbreak{}。这是指数塔的第二层。''

\NovelDisplay{M01011R01}{%
\Pi_3\ \Pi_1\text{-}\Pi_\omega \sim \psi_a(T^{\varepsilon_{T+1}+T^T})%
}

\NovelDisplay{M01012R01}{%
\Pi_4\ \Pi_1\text{-}\Pi_\omega \sim \psi_a(T^{\varepsilon_{T+1}+T^{T^T}})%
}

\NovelDisplay{M01013R01}{%
\Pi_\omega\ \Pi_1\text{-}\Pi_\omega \sim \psi_a(T^{\varepsilon_{T+1}\times 2})%
}

小舞倒吸了一口凉气：``\allowbreak{}\NovelInline{I04821}{\Pi_3}\allowbreak{} 对应 \allowbreak{}\NovelInline{I04822}{T^T}\allowbreak{}，\allowbreak{}\NovelInline{I04823}{\Pi_4}\allowbreak{} 对应 \allowbreak{}\NovelInline{I04824}{T^{T^T}}\allowbreak{}------数字每增加 \allowbreak{}\NovelInline{I04825}{1}\allowbreak{}，指数塔就升高一层！而 \allowbreak{}\NovelInline{I04826}{\Pi_\omega}\allowbreak{} 直接把指数部分的加法变成了乘法！''

``对。''Phyrion的声音沉稳有力，``\allowbreak{}\NovelInline{I04827}{\Pi_\omega}\allowbreak{} 作为前缀时，它的作用是将底层的递归结构提升一个维度。''

\NovelDisplay{M01014R01}{%
(\Pi_\omega\ \Pi_1\text{-})^\omega \sim \psi_a(T^{\varepsilon_{T+1}\times\omega})%
}

\NovelDisplay{M01015R01}{%
(\Pi_\omega\ \Pi_1\text{-})^{(1,0)} \sim \psi_a(T^{\varepsilon_{T+1}\times a})%
}

``\allowbreak{}\NovelInline{I04828}{\omega}\allowbreak{} 次和不动点次递归。''卢本韦快速说道，``模式自相似。''

Phyrion没有停顿，笔尖继续在黑板上飞舞：

\NovelDisplay{M01016R01}{%
\Pi_2\text{-}\Pi_\omega \sim \psi_a(T^{T^{\varepsilon_{T+1}+1}})%
}

\NovelDisplay{M01017R01}{%
\Pi_3\ \Pi_2\text{-}\Pi_\omega \sim \psi_a(T^{T^{\varepsilon_{T+1}+T}})%
}

\NovelDisplay{M01018R01}{%
\Pi_\omega\ \Pi_2\text{-}\Pi_\omega \sim \psi_a(T^{T^{\varepsilon_{T+1}\times 2}})%
}

\NovelDisplay{M01019R01}{%
\Pi_3\text{-}\Pi_\omega \sim \psi_a(T^{T^{T^{\varepsilon_{T+1}+1}}})%
}

蔡须坤深吸一口气：``两层指数塔、三层指数塔\ldots\ldots{}\allowbreak{}\NovelInline{I04829}{\Pi_2\text{-}}\allowbreak{} 是两层，\allowbreak{}\NovelInline{I04830}{\Pi_3\text{-}}\allowbreak{} 是三层。前缀中的数字直接对应指数塔的层数。''

然后，Phyrion写下了一个转折性的等式：

\NovelDisplay{M01020R01}{%
\Pi_\omega\text{-}\Pi_\omega \sim \psi_a(\varepsilon_{T+2})%
}

修炼堂里安静了一瞬。

从 \allowbreak{}\NovelInline{I04831}{T^{T^{T^{\cdots}}}}\allowbreak{} 的指数塔，直接跃迁到了 \allowbreak{}\NovelInline{I04832}{\varepsilon_{T+2}}\allowbreak{}。

``\allowbreak{}\NovelInline{I04833}{\Pi_\omega\text{-}\Pi_\omega}\allowbreak{}。''Phyrion的声音变得低沉，``两个 \allowbreak{}\NovelInline{I04834}{\Pi_\omega}\allowbreak{} 用 \allowbreak{}\NovelInline{I04835}{\text{-}}\allowbreak{} 连接，意味着 \allowbreak{}\NovelInline{I04836}{\varepsilon}\allowbreak{} 的下标增加了 \allowbreak{}\NovelInline{I04837}{1}\allowbreak{}。这和之前T-1Y序列中 \allowbreak{}\NovelInline{I04838}{1,\omega,\omega \sim \varepsilon_1}\allowbreak{} 的模式完全一致。''

\NovelDisplay{M01021R01}{%
\Pi_2\ \Pi_1\text{-}\Pi_\omega\text{-}\Pi_\omega \sim \psi_a(T^{\varepsilon_{T+2}+1})%
}

\NovelDisplay{M01022R01}{%
\Pi_\omega\text{-}\Pi_\omega\ \Pi_1\text{-}\Pi_\omega\text{-}\Pi_\omega \sim \psi_a(T^{\varepsilon_{T+2}\times 2})%
}

\NovelDisplay{M01023R01}{%
\Pi_2\text{-}\Pi_\omega\text{-}\Pi_\omega \sim \psi_a(T^{T^{\varepsilon_{T+2}+1}})%
}

\NovelDisplay{M01024R01}{%
\Pi_\omega\text{-}\Pi_\omega\text{-}\Pi_\omega \sim \psi_a(\varepsilon_{T+3})%
}

``在新的基底 \allowbreak{}\NovelInline{I04839}{\varepsilon_{T+2}}\allowbreak{} 上，所有的展开规则重新开始。''唐三说道，``自相似性贯穿始终。''

\NovelDisplay{M01025R01}{%
(\Pi_\omega\text{-})^\omega \sim \psi_a(\varepsilon_{T+\omega})%
}

\NovelDisplay{M01026R01}{%
(\Pi_\omega\text{-})^{(1,0)} \sim \psi_a(\varepsilon_{T+a})%
}

``\allowbreak{}\NovelInline{I04840}{\omega}\allowbreak{} 层 \allowbreak{}\NovelInline{I04841}{\Pi_\omega\text{-}}\allowbreak{} 对应 \allowbreak{}\NovelInline{I04842}{\varepsilon_{T+\omega}}\allowbreak{}，不动点层对应 \allowbreak{}\NovelInline{I04843}{\varepsilon_{T+a}}\allowbreak{}。''蔡兴坤快速总结。

Phyrion擦掉了黑板的上半部分，只留下最后几行作为参照。然后他写下了新的起点：

\NovelDisplay{M01027R01}{%
\Pi_{\omega+1} \sim \psi_a(\varepsilon_{T\times 2})%
}

\allowbreak{}\NovelInline{I04844}{\Pi_{\omega+1}}\allowbreak{}。

从 \allowbreak{}\NovelInline{I04845}{\Pi_\omega}\allowbreak{} 到 \allowbreak{}\NovelInline{I04846}{\Pi_{\omega+1}}\allowbreak{}------这是一个全新的层级，这意味着，我们到达了\allowbreak{}\NovelInline{I04847}{real.\Pi_ω}\allowbreak{}。

``投影 vs 反射稳定的细节，其实很值得你们扽一扽。''Phyrion忽然说道，语气中带着一丝深意，``从 \allowbreak{}\NovelInline{I04848}{\Pi_\omega}\allowbreak{} 到 \allowbreak{}\NovelInline{I04849}{\Pi_{\omega+1}}\allowbreak{} 的跃迁，在投影体系下是突变的，在反射体系下是渐进的。这种差异，就是两种记号体系的本质区别之一。''

他没有过多解释，而是继续写下去：

\NovelDisplay{M01028R01}{%
\Pi_2\ \Pi_1\text{-}\Pi_{\omega+1} \sim \psi_a(T^{\varepsilon_{T\times 2}+1})%
}

\NovelDisplay{M01029R01}{%
\Pi_{\omega+1}\ \Pi_1\text{-}\Pi_{\omega+1} \sim \psi_a(T^{\varepsilon_{T\times 2}\times 2})%
}

\NovelDisplay{M01030R01}{%
\Pi_2\text{-}\Pi_{\omega+1} \sim \psi_a(T^{T^{\varepsilon_{T\times 2}+1}})%
}

\NovelDisplay{M01031R01}{%
\Pi_\omega\text{-}\Pi_{\omega+1} \sim \psi_a(\varepsilon_{T\times 2+1})%
}

``看到了吗？''Phyrion转身面向众人，``在 \allowbreak{}\NovelInline{I04850}{\Pi_{\omega+1}}\allowbreak{} 的体系中，所有的展开规则和 \allowbreak{}\NovelInline{I04851}{\Pi_\omega}\allowbreak{} 完全一致------只不过基底从 \allowbreak{}\NovelInline{I04852}{\varepsilon_{T+1}}\allowbreak{} 换成了 \allowbreak{}\NovelInline{I04853}{\varepsilon_{T\times 2}}\allowbreak{}。''

\NovelDisplay{M01032R01}{%
\Pi_{\omega+1}\ \Pi_\omega\text{-}\Pi_{\omega+1} \sim \psi_a(\varepsilon_{T\times 3})%
}

\NovelDisplay{M01033R01}{%
(\Pi_{\omega+1}\ \Pi_\omega\text{-})^\omega \sim \psi_a(\varepsilon_{T\times\omega})%
}

\NovelDisplay{M01034R01}{%
(\Pi_{\omega+1}\ \Pi_\omega\text{-})^{(1,0)}\sim \psi_a(\varepsilon_{T\times a})%
}

``\allowbreak{}\NovelInline{I04854}{\Pi_{\omega+1}\ \Pi_\omega\text{-}}\allowbreak{} 这个组合，对应的是 \allowbreak{}\NovelInline{I04855}{\varepsilon}\allowbreak{} 下标的乘法递归。''蔡须坤说道，``\allowbreak{}\NovelInline{I04856}{\omega}\allowbreak{} 次对应 \allowbreak{}\NovelInline{I04857}{T\times\omega}\allowbreak{}，不动点次代表\allowbreak{}\NovelInline{I04858}{T\times a}\allowbreak{}。

\NovelDisplay{M01035R01}{%
\Pi_{\omega+1}\text{-}\Pi_{\omega+1} \sim \psi_a(\varepsilon_{T^2})%
}

\NovelDisplay{M01036R01}{%
(\Pi_{\omega+1}\text{-})^\omega \sim \psi_a(\varepsilon_{T^\omega})%
}

\NovelDisplay{M01037R01}{%
(\Pi_{\omega+1}\text{-})^{(1,0)} \sim \psi_a(\varepsilon_{T^a})%
}

``\allowbreak{}\NovelInline{I04859}{\Pi_{\omega+1}\text{-}\Pi_{\omega+1}}\allowbreak{} 对应 \allowbreak{}\NovelInline{I04860}{\varepsilon_{T^2}}\allowbreak{}------下标从乘法变成了幂运算。''唐三的声音中带着一丝敬畏，``这和 \allowbreak{}\NovelInline{I04861}{\Pi_\omega\text{-}\Pi_\omega}\allowbreak{} 对应 \allowbreak{}\NovelInline{I04862}{\varepsilon_{T+2}}\allowbreak{} 的模式是一致的，只不过操作的层级更高了。''

\NovelDisplay{M01038R01}{%
\Pi_{\omega+2} \sim \psi_a(\varepsilon_{T^T})%
}

\NovelDisplay{M01039R01}{%
\Pi_{\omega+3} \sim \psi_a(\varepsilon_{T^{T^T}})%
}

``\allowbreak{}\NovelInline{I04863}{\Pi_{\omega+n}}\allowbreak{} 对应 \allowbreak{}\NovelInline{I04864}{\varepsilon}\allowbreak{} 下标的 \allowbreak{}\NovelInline{I04865}{n}\allowbreak{} 层指数塔。''卢本韦总结道。

Phyrion再次擦掉黑板，开辟了新的区域。

``现在进入 \allowbreak{}\NovelInline{I04866}{\varepsilon}\allowbreak{} 的递归。''他说。

\NovelDisplay{M01040R01}{%
\Pi_{\omega\times 2} \sim \psi_a(\varepsilon_{\varepsilon_{T+1}})%
}

\NovelDisplay{M01041R01}{%
\Pi_{\omega\times 2+1} \sim \psi_a(\varepsilon_{\varepsilon_{T\times 2}})%
}

\NovelDisplay{M01042R01}{%
\Pi_{\omega\times 2+2} \sim \psi_a(\varepsilon_{\varepsilon_{T^T}})%
}

\NovelDisplay{M01043R01}{%
\Pi_{\omega\times 3} \sim \psi_a(\varepsilon_{\varepsilon_{\varepsilon_{T+1}}})%
}

``\allowbreak{}\NovelInline{I04867}{\Pi_{\omega\times 2}}\allowbreak{} 对应两层 \allowbreak{}\NovelInline{I04868}{\varepsilon}\allowbreak{} 嵌套，\allowbreak{}\NovelInline{I04869}{\Pi_{\omega\times 3}}\allowbreak{} 对应三层。''蔡兴坤说道，``\allowbreak{}\NovelInline{I04870}{\omega\times n}\allowbreak{} 对应 \allowbreak{}\NovelInline{I04871}{n}\allowbreak{} 层 \allowbreak{}\NovelInline{I04872}{\varepsilon}\allowbreak{} 递归。''

然后是 \allowbreak{}\NovelInline{I04873}{\zeta}\allowbreak{} 层级：

\NovelDisplay{M01044R01}{%
\Pi_{\omega^2} \sim \psi_a(\zeta_{T+1})%
}

\NovelDisplay{M01045R01}{%
\Pi_{\omega^2+1} \sim \psi_a(\zeta_{T\times 2})%
}

\NovelDisplay{M01046R01}{%
\Pi_{\omega^2+2} \sim \psi_a(\zeta_{T^T})%
}

\NovelDisplay{M01047R01}{%
\Pi_{\omega^2+\omega} \sim \psi_a(\zeta_{\varepsilon_{T+1}})%
}

\NovelDisplay{M01048R01}{%
\Pi_{\omega^2\times 2} \sim \psi_a(\zeta_{\zeta_{T+1}})%
}

``\allowbreak{}\NovelInline{I04874}{\Pi_{\omega^2}}\allowbreak{} 对应 \allowbreak{}\NovelInline{I04875}{\zeta_{T+1}}\allowbreak{}------\allowbreak{}\NovelInline{I04876}{\varepsilon}\allowbreak{} 的不动点。''蔡须坤的声音清亮而笃定，``\allowbreak{}\NovelInline{I04877}{\omega^2+n}\allowbreak{} 对应 \allowbreak{}\NovelInline{I04878}{\zeta}\allowbreak{} 下标的递增，\allowbreak{}\NovelInline{I04879}{\omega^2+\omega\times 2}\allowbreak{} 对应 \allowbreak{}\NovelInline{I04880}{\zeta}\allowbreak{} 下标出现 \allowbreak{}\NovelInline{I04881}{\varepsilon}\allowbreak{}，\allowbreak{}\NovelInline{I04882}{\omega^2\times 2}\allowbreak{} 对应 \allowbreak{}\NovelInline{I04883}{\zeta}\allowbreak{} 的递归。''

\allowbreak{}\NovelInline{I04884}{\eta}\allowbreak{} 与 \allowbreak{}\NovelInline{I04885}{\varphi}\allowbreak{} 层级紧随其后：

\NovelDisplay{M01049R01}{%
\Pi_{\omega^3} \sim \psi_a(\eta_{T+1})%
}

\NovelDisplay{M01050R01}{%
\Pi_{\omega^4} \sim \psi_a(\varphi(4,T+1))%
}

\NovelDisplay{M01051R01}{%
\Pi_{\omega^\omega} \sim \psi_a(\varphi(\omega,T+1))%
}

\NovelDisplay{M01052R01}{%
\Pi_{\omega^\omega+1} \sim \psi_a(\varphi(\omega,T\times 2))%
}

\NovelDisplay{M01053R01}{%
\Pi_{\omega^\omega+\omega} \sim \psi_a(\varphi(\omega,\varepsilon_{T+1}))%
}

\NovelDisplay{M01054R01}{%
\Pi_{\omega^\omega\times 2} \sim \psi_a(\varphi(\omega,\varphi(\omega,T+1)))%
}

\NovelDisplay{M01055R01}{%
\Pi_{\omega^{\omega+1}} \sim \psi_a(\varphi(\omega+1,T+1))%
}

\NovelDisplay{M01056R01}{%
\Pi_{\varepsilon_0} \sim \psi_a(\varphi(\varepsilon_0,T+1))%
}

KNK的耳朵猛地竖了起来：``\allowbreak{}\NovelInline{I04886}{\Pi_{\omega^\omega}}\allowbreak{} 对应 \allowbreak{}\NovelInline{I04887}{\varphi(\omega,T+1)}\allowbreak{} 喵！Veblen函数正式登场了！''

``\allowbreak{}\NovelInline{I04888}{\varphi}\allowbreak{} 的第一个参数可以是任意序数。''Phyrion的声音变得低沉而有力，``\allowbreak{}\NovelInline{I04889}{\omega}\allowbreak{}、\allowbreak{}\NovelInline{I04890}{\omega+1}\allowbreak{}、\allowbreak{}\NovelInline{I04891}{\varepsilon_0}\allowbreak{}------每一个都忠实地反映在 \allowbreak{}\NovelInline{I04892}{\Pi}\allowbreak{} 的下标中。''

然后是不可数基数的跃迁：

\NovelDisplay{M01057R01}{%
\Pi_\Omega \sim \psi_a(\varphi(\Omega,T+1))%
}

\NovelDisplay{M01058R01}{%
\Pi_I \sim \psi_a(\varphi(\psi_a(T^2),T+1))%
}

\NovelDisplay{M01059R01}{%
\Pi_M \sim \psi_a(\varphi(\psi_a(T^T),T+1))%
}

\NovelDisplay{M01060R01}{%
\Pi_{\Pi_\omega} \sim \psi_a(\varphi(\psi_a(\varepsilon_{T+1}),T+1))%
}

\allowbreak{}\NovelInline{I04893}{\Pi_\Omega}\allowbreak{}、\allowbreak{}\NovelInline{I04894}{\Pi_I}\allowbreak{}、\allowbreak{}\NovelInline{I04895}{\Pi_M}\allowbreak{}------从可数到不可数，从不可达到马洛。

蔡兴坤的肩膀上，黑子鸡忽然``咕''了一声，圆溜溜的眼睛盯着黑板上的公式，似乎闪过了一丝异样的光芒。

\NovelDisplay{M01061R01}{%
\Pi\text{-}\Phi(1,0) \sim \psi_a(\varphi(a,T+1))%
}

``\allowbreak{}\NovelInline{I04896}{a}\allowbreak{}是一种不动点。''Phyrion缓缓说道，``它把前面所有的结构封装在了一个表达式里。''

Phyrion没有停下。他以 \allowbreak{}\NovelInline{I04897}{\Pi\text{-}\Phi(1,0)}\allowbreak{} 为新的基底,这个函数是每取\allowbreak{}\NovelInline{I04898}{\Pi}\allowbreak{}的不动点，之后加上1，继续向上攀登相当于\allowbreak{}\NovelInline{I04899}{f(α)→\Pi_α}\allowbreak{}的φ函数,是下标的不动点进位。

\chapter{第八十六章：拓展反射层级}\label{ux7b2cux516bux5341ux516dux7ae0ux62d3ux5c55ux53cdux5c04ux5c42ux7ea7}

第二天，小队继续修炼，接着推导更高层级的内容

\NovelDisplay{M01062R01}{%
\Pi_1\text{-}\Pi\text{-}\Phi(1,0) \sim \psi_a(\varphi(a,T+1)\times\omega)%
}

\NovelDisplay{M01063R01}{%
\Pi_2\ \Pi_1\text{-}\Pi\text{-}\Phi(1,0) \sim \psi_a(T^{\varphi(a,T+1)+1})%
}

\NovelDisplay{M01064R01}{%
\Pi\text{-}\Phi(1,0)\ \Pi_1\text{-}\Pi\text{-}\Phi(1,0) \sim \psi_a(T^{\varphi(a,T+1)\times 2})%
}

\NovelDisplay{M01065R01}{%
\Pi_2\text{-}\Pi\text{-}\Phi(1,0) \sim \psi_a(T^{T^{\varphi(a,T+1)}})%
}

\NovelDisplay{M01066R01}{%
\Pi_\omega\text{-}\Pi\text{-}\Phi(1,0) \sim \psi_a(\varepsilon_{\varphi(a,T+1)+1})%
}

\NovelDisplay{M01067R01}{%
\Pi_{\omega^2}\text{-}\Pi\text{-}\Phi(1,0) \sim \psi_a(\zeta_{\varphi(a,T+1)+1})%
}

\NovelDisplay{M01068R01}{%
\Pi_{\omega^\omega}\text{-}\Pi\text{-}\Phi(1,0) \sim \psi_a(\varphi(\omega,\varphi(a,T+1)+1))%
}

\NovelDisplay{M01069R01}{%
\Pi_\Omega\text{-}\Pi\text{-}\Phi(1,0) \sim \psi_a(\varphi(\Omega,\varphi(a,T+1)+1))%
}

``所有的规则在 \allowbreak{}\NovelInline{I04900}{\Pi\text{-}\Phi(1,0)}\allowbreak{} 之上重新生效。''唐三低声说道，``乘法、指数塔、\allowbreak{}\NovelInline{I04901}{\varepsilon}\allowbreak{}、\allowbreak{}\NovelInline{I04902}{\zeta}\allowbreak{}、\allowbreak{}\NovelInline{I04903}{\varphi}\allowbreak{}、\allowbreak{}\NovelInline{I04904}{\Omega}\allowbreak{}------全部重现。''

\NovelDisplay{M01070R01}{%
\Pi\text{-}\Phi(1,0)\text{-}\Pi\text{-}\Phi(1,0) \sim \psi_a(\varphi(a,T+2))%
}

\NovelDisplay{M01071R01}{%
(\Pi\text{-}\Phi(1,0)\text{-})^{\omega} \sim \psi_a(\varphi(a,T+\omega))%
}

\NovelDisplay{M01072R01}{%
(\Pi\text{-}\Phi(1,0)\text{-})^{(1,0)} \sim \psi_a(\varphi(a,T+a))%
}

``\allowbreak{}\NovelInline{I04905}{\Pi\text{-}\Phi(1,0)\text{-}\Pi\text{-}\Phi(1,0)}\allowbreak{} 对应 \allowbreak{}\NovelInline{I04906}{\varphi(a,T+2)}\allowbreak{}------\allowbreak{}\NovelInline{I04907}{T}\allowbreak{} 的下标增加了 \allowbreak{}\NovelInline{I04908}{1}\allowbreak{}。''蔡须坤说道，``和 \allowbreak{}\NovelInline{I04909}{\Pi_\omega\text{-}\Pi_\omega}\allowbreak{} 对应 \allowbreak{}\NovelInline{I04910}{\varepsilon_{T+2}}\allowbreak{} 的模式完全一致。''

\NovelDisplay{M01073R01}{%
\Pi_{\Pi\text{-}\Phi(1,0)+1} \sim \psi_a(\varphi(a,T\times 2))%
}

\NovelDisplay{M01074R01}{%
\Pi_{\Pi\text{-}\Phi(1,0)+\omega} \sim \psi_a(\varphi(a,\varepsilon_{T+1}))%
}

\NovelDisplay{M01075R01}{%
\Pi_{\Pi\text{-}\Phi(1,0)\times 2} \sim \psi_a(\varphi(a,\varphi(a,T+1)))%
}

\NovelDisplay{M01076R01}{%
\Pi_{\Pi\text{-}\Phi(1,0)\times\omega} \sim \psi_a(\varphi(a+1,0))%
}

\NovelDisplay{M01077R01}{%
\Pi_{\Pi\text{-}\Phi(1,0)\times\omega^2} \sim \psi_a(\varphi(a+2,0))%
}

\NovelDisplay{M01078R01}{%
\Pi_{\Pi\text{-}\Phi(1,0)\times\Omega} \sim \psi_a(\varphi(a+\Omega,0))%
}

\NovelDisplay{M01079R01}{%
\Pi_{\Pi\text{-}\Phi(1,0)^2} \sim \psi_a(\varphi(a+\psi_a(\varphi(a,T+1)),0))%
}

\NovelDisplay{M01080R01}{%
\Pi_{\Omega_{\Pi\text{-}\Phi(1,0)+1}} \sim \psi_a(\varphi(a+\psi_a(\varphi(a,T+1)+T),0))%
}

\NovelDisplay{M01081R01}{%
\Pi_{\Pi_{\Pi\text{-}\Phi(1,0)+1}} \sim \psi_a(\varphi(a+\psi_a(\varphi(a,T\times 2)),0))%
}

``\allowbreak{}\NovelInline{I04911}{\Pi\text{-}\Phi(1,0)}\allowbreak{} 作为下标时，它的每一种运算都对应着 \allowbreak{}\NovelInline{I04912}{\varphi}\allowbreak{} 第一个参数的复杂化。''卢本韦的声音越来越快，``\allowbreak{}\NovelInline{I04913}{+1}\allowbreak{}、\allowbreak{}\NovelInline{I04914}{+\omega}\allowbreak{}、\allowbreak{}\NovelInline{I04915}{\times 2}\allowbreak{}、\allowbreak{}\NovelInline{I04916}{\times\omega}\allowbreak{}、\allowbreak{}\NovelInline{I04917}{\times\Omega}\allowbreak{}、平方、\allowbreak{}\NovelInline{I04918}{\Omega}\allowbreak{} 下标、\allowbreak{}\NovelInline{I04919}{\Pi}\allowbreak{} 下标------每一种操作都在 \allowbreak{}\NovelInline{I04920}{\varphi}\allowbreak{} 的第一个参数中找到了精确的对应。''

Phyrion的笔尖继续飞舞，速度越来越快。

\NovelDisplay{M01082R01}{%
\Pi\text{-}\Phi(1,1) \sim \psi_a(\varphi(a\times 2,T+1))%
}

\NovelDisplay{M01083R01}{%
\Pi\text{-}\Phi(1,\omega) \sim \psi_a(\varphi(a\times\omega,T+1))%
}

\NovelDisplay{M01084R01}{%
\Pi\text{-}\Phi(1,\Omega) \sim \psi_a(\varphi(a\times\Omega,T+1))%
}

\NovelDisplay{M01085R01}{%
\Pi\text{-}\Phi(1,\Pi\text{-}\Phi(1,0)) \sim \psi_a(\varphi(a\times\psi_a(\varphi(a,T+1)),T+1))%
}

``\allowbreak{}\NovelInline{I04921}{\Pi\text{-}\Phi(1,n)}\allowbreak{} 的第二个参数，对应 \allowbreak{}\NovelInline{I04922}{\varphi}\allowbreak{} 第一个参数的乘法因子。''蔡兴坤说道，``\allowbreak{}\NovelInline{I04923}{1}\allowbreak{} 对应 \allowbreak{}\NovelInline{I04924}{\times 2}\allowbreak{}，\allowbreak{}\NovelInline{I04925}{\omega}\allowbreak{} 对应 \allowbreak{}\NovelInline{I04926}{\times\omega}\allowbreak{}，\allowbreak{}\NovelInline{I04927}{\Omega}\allowbreak{} 对应 \allowbreak{}\NovelInline{I04928}{\times\Omega}\allowbreak{}，\allowbreak{}\NovelInline{I04929}{\Pi\text{-}\Phi(1,0)}\allowbreak{} 对应 \allowbreak{}\NovelInline{I04930}{\times\psi_a(\varphi(a,T+1))}\allowbreak{}。''

\NovelDisplay{M01086R01}{%
\Pi\text{-}\Phi(2,0) \sim \psi_a(\varphi(a^2,T+1))%
}

\NovelDisplay{M01087R01}{%
\Pi\text{-}\Phi(\omega,0) \sim \psi_a(\varphi(a^\omega,T+1))%
}

\NovelDisplay{M01088R01}{%
\Pi\text{-}\Phi(1,0,0) \sim \psi_a(\varphi(a^a,T+1))%
}

\NovelDisplay{M01089R01}{%
\Pi\text{-}\Phi(1@\omega) \sim \psi_a(\varphi(a^{a^\omega},T+1))%
}

\NovelDisplay{M01090R01}{%
\Pi\text{-}\Phi(1@(1,0)) \sim \psi_a(\varphi(a^{a^a},T+1))%
}

\NovelDisplay{M01091R01}{%
\Pi\text{-}\Phi(1@(1,,0)) \sim \psi_a(\varphi(\varepsilon_{a+1},T+1))%
}

``\allowbreak{}\NovelInline{I04931}{\Pi\text{-}\Phi}\allowbreak{} 的第一个参数从 \allowbreak{}\NovelInline{I04932}{2}\allowbreak{} 到 \allowbreak{}\NovelInline{I04933}{\omega}\allowbreak{} 到 \allowbreak{}\NovelInline{I04934}{(1,0,0)}\allowbreak{} 到 \allowbreak{}\NovelInline{I04935}{1@\omega}\allowbreak{} 到 \allowbreak{}\NovelInline{I04936}{1@(1,0)}\allowbreak{}------对应 \allowbreak{}\NovelInline{I04937}{\varphi}\allowbreak{} 第一个参数从 \allowbreak{}\NovelInline{I04938}{a^2}\allowbreak{} 到 \allowbreak{}\NovelInline{I04939}{a^\omega}\allowbreak{} 到 \allowbreak{}\NovelInline{I04940}{a^a}\allowbreak{} 到 \allowbreak{}\NovelInline{I04941}{a^{a^\omega}}\allowbreak{} 到 \allowbreak{}\NovelInline{I04942}{a^{a^a}}\allowbreak{}。''唐三的声音中带着一丝颤抖，``而 \allowbreak{}\NovelInline{I04943}{1@(1,,0)=1@(1@(1@.....))}\allowbreak{} 直接跳到了 \allowbreak{}\NovelInline{I04944}{\varepsilon_{a+1}}\allowbreak{}------这是由于\allowbreak{}\NovelInline{I04945}{f(α)→\Pi_α}\allowbreak{}相对于\allowbreak{}\NovelInline{I04946}{\psi_a(\varphi(a\text{的递归运算},T+1))}\allowbreak{}，正如\allowbreak{}\NovelInline{I04947}{f(α)→ω^α}\allowbreak{}对于\allowbreak{}\NovelInline{I04948}{ψ(Ω\text{的递归运算})}\allowbreak{}。''

然后是 \allowbreak{}\NovelInline{I04949}{\Pi\text{-}I}\allowbreak{} 层级,相当于\allowbreak{}\NovelInline{I04950}{f(α)→\Pi_α}\allowbreak{}的\allowbreak{}\NovelInline{I04951}{I}\allowbreak{}点：

\NovelDisplay{M01092R01}{%
\Pi\text{-}I(1,0) \sim \psi_a(\varphi(T,1))%
}

\NovelDisplay{M01093R01}{%
\Pi\text{-}I(1,0)\text{-}\Pi\text{-}I(1,0) \sim \psi_a(\varphi(T,2))%
}

\NovelDisplay{M01094R01}{%
\Pi_{\Pi\text{-}I(1,0)+1} \sim \psi_a(\varphi(T,T))%
}

\NovelDisplay{M01095R01}{%
\Pi_{\Pi\text{-}I(1,0)\times\omega} \sim \psi_a(\varphi(T+1,0))%
}

\NovelDisplay{M01096R01}{%
\Pi_{\Pi\text{-}I(1,0)^2} \sim \psi_a(\varphi(T\times\psi_a(\varphi(T,1)),0))%
}

\NovelDisplay{M01097R01}{%
\Pi\text{-}\Phi(1,\Pi\text{-}I(1,0)+1) \sim \psi_a(\varphi(T+a,0))%
}

\NovelDisplay{M01098R01}{%
\Pi\text{-}\Phi(1,\Pi\text{-}I(1,0)+2) \sim \psi_a(\varphi(T+a\times 2,0))%
}

\NovelDisplay{M01099R01}{%
\Pi\text{-}\Phi(2,\Pi\text{-}I(1,0)+1) \sim \psi_a(\varphi(T+a^2,0))%
}

\NovelDisplay{M01100R01}{%
\Pi\text{-}\Phi(1@(1,,0),\Pi\text{-}I(1,0)+1) \sim \psi_a(\varphi(T+\varepsilon_{a+1},0))%
}

``\allowbreak{}\NovelInline{I04952}{\Pi\text{-}I(1,0)}\allowbreak{} 对应 \allowbreak{}\NovelInline{I04953}{\varphi(T,1)}\allowbreak{}------\allowbreak{}\NovelInline{I04954}{T}\allowbreak{} 成为了 \allowbreak{}\NovelInline{I04955}{\varphi}\allowbreak{} 的第一个参数。''蔡须坤的声音清亮而坚定，``之后所有的展开规则，都在 \allowbreak{}\NovelInline{I04956}{\varphi(T,\cdot)}\allowbreak{} 的基础上重新生效。''

\NovelDisplay{M01101R01}{%
\Pi\text{-}I(1,1) \sim \psi_a(\varphi(T\times 2,0))%
}

\NovelDisplay{M01102R01}{%
\Pi\text{-}I(1,\omega) \sim \psi_a(\varphi(T\times\omega,0))%
}

\NovelDisplay{M01103R01}{%
\Pi\text{-}I(1,x)\text{-}\Phi(1,0) \sim \psi_a(\varphi(T\times a,0))%
}

\NovelDisplay{M01104R01}{%
\Pi\text{-}I(2,0) \sim \psi_a(\varphi(T^2,0))%
}

\NovelDisplay{M01105R01}{%
\Pi\text{-}I(\omega,0) \sim \psi_a(\varphi(T^\omega,0))%
}

\NovelDisplay{M01106R01}{%
\Pi\text{-}I(x,0)\text{-}\Phi(1,0) \sim \psi_a(\varphi(T^a,0))%
}

\NovelDisplay{M01107R01}{%
\Pi\text{-}I(1,0,0) \sim \psi_a(\varphi(T^{a+1},0))%
}

\NovelDisplay{M01108R01}{%
\Pi\text{-}I(1@(1,,0)) \sim \psi_a(\varphi(T^{\varepsilon_{a+1}},0))%
}

``\allowbreak{}\NovelInline{I04957}{\Pi\text{-}I}\allowbreak{} 的参数变化，精确对应 \allowbreak{}\NovelInline{I04958}{\varphi}\allowbreak{} 第一个参数中 \allowbreak{}\NovelInline{I04959}{T}\allowbreak{} 的运算。''卢本韦快速总结，``乘法、幂运算、\allowbreak{}\NovelInline{I04960}{a}\allowbreak{} 次递归、\allowbreak{}\NovelInline{I04961}{\varepsilon_{a+1}}\allowbreak{}------全部一致。''

Phyrion的笔尖在黑板上重重地写下了接下来的内容：

\NovelDisplay{M01109R01}{%
\Pi\text{-}M(1,0) \sim \psi_a(\varphi(T^T,0))%
}

\NovelDisplay{M01110R01}{%
\Pi\text{-}N(1,0) \sim \psi_a(\varphi(T^{T^2},0))%
}

\NovelDisplay{M01111R01}{%
\Pi\text{-}(\Pi_2\text{-})^\omega(1,0) \sim \psi_a(\varphi(T^{T^\omega},0))%
}

\NovelDisplay{M01112R01}{%
\Pi\text{-}(\Pi_2\text{-})^a(1,0) \sim \psi_a(\varphi(T^{T^a},0))%
}

\NovelDisplay{M01113R01}{%
\Pi\text{-}K(1,0) \sim \psi_a(\varphi(T^{T^T},0))%
}

``\allowbreak{}\NovelInline{I04962}{\Pi\text{-}M}\allowbreak{}、\allowbreak{}\NovelInline{I04963}{\Pi\text{-}N}\allowbreak{}、\allowbreak{}\NovelInline{I04964}{\Pi\text{-}K}\allowbreak{}------这些大基数记号对应的是 \allowbreak{}\NovelInline{I04965}{T}\allowbreak{} 的指数塔。''唐三的声音越来越急促，``\allowbreak{}\NovelInline{I04966}{T^T}\allowbreak{}、\allowbreak{}\NovelInline{I04967}{T^{T^2}}\allowbreak{}、\allowbreak{}\NovelInline{I04968}{T^{T^\omega}}\allowbreak{}、\allowbreak{}\NovelInline{I04969}{T^{T^a}}\allowbreak{}、\allowbreak{}\NovelInline{I04970}{T^{T^T}}\allowbreak{}------每一层指数塔的升高，都对应着一个新的大基数记号。''

\NovelDisplay{M01114R01}{%
\Pi\text{-}(\Pi_\omega)(1,0) \sim \psi_a(\varphi(\varepsilon_{T+1},0))%
}

\NovelDisplay{M01115R01}{%
\Pi\text{-}(\Pi\text{-}\Phi(1,0))(1,0) \sim \psi_a(\varphi(\varphi(a,T+1),0))%
}

\NovelDisplay{M01116R01}{%
\Pi\text{-}(\Pi\text{-}I(1,0))(1,0) \sim \psi_a(\varphi(\varphi(T,1),0))%
}

\NovelDisplay{M01117R01}{%
\Pi\text{-}(\Pi\text{-}(\Pi_\omega)(1,0))(1,0) \sim \psi_a(\varphi(\varphi(\varepsilon_{T+1},0),0))%
}

``\allowbreak{}\NovelInline{I04971}{\Pi\text{-}(\cdots)(1,0)}\allowbreak{} 这个结构，将括号内的整个 \allowbreak{}\NovelInline{I04972}{\Pi}\allowbreak{} 表达式转化为 \allowbreak{}\NovelInline{I04973}{\varphi}\allowbreak{} 的第一个参数。''蔡须坤说道，``这是一种元层面的操作------用 \allowbreak{}\NovelInline{I04974}{\Pi}\allowbreak{} 记法来描述 \allowbreak{}\NovelInline{I04975}{\varphi}\allowbreak{} 函数的参数结构。''

然后，Phyrion停下了笔。

他深吸一口气，转身面向众人。

``接下来------''他的声音变得格外庄重，``是 \allowbreak{}\NovelInline{I04976}{real.\Pi(1,0)}\allowbreak{}。''

他在黑板上重重地写下了：

\NovelDisplay{M01118R01}{%
\text{real.}\Pi(1,0) \sim \psi_a(\Gamma_{T+1})%
}

\NovelDisplay{M01119R01}{%
\Pi(1,0) \sim \psi_a(\Gamma_{T+1})%
}

\allowbreak{}\NovelInline{I04977}{\Gamma_{T+1}}\allowbreak{}。

费弗巴赫-舒滕巴赫序数在 \allowbreak{}\NovelInline{I04978}{T+1}\allowbreak{} 处的取值。

``\allowbreak{}\NovelInline{I04979}{\text{real.}\Pi(1,0)}\allowbreak{} 都对应 \allowbreak{}\NovelInline{I04980}{\Gamma_{T+1}}\allowbreak{}。''Phyrion的声音在寂静的修炼堂中回荡，``这是 \allowbreak{}\NovelInline{I04981}{\Pi}\allowbreak{} 记法中的一个里程碑------从 \allowbreak{}\NovelInline{I04982}{\text{二元}\varphi}\allowbreak{} 函数跨越到了 \allowbreak{}\NovelInline{I04983}{\Gamma}\allowbreak{} 函数。''

\NovelDisplay{M01120R01}{%
\Pi_\omega\text{-}\Pi(1,0) \sim \psi_a(\varepsilon_{\Gamma_{T+1}})%
}

\NovelDisplay{M01121R01}{%
\Pi\text{-}\Phi(1,0)\text{-}\Pi(1,0) \sim \psi_a(\varphi(a,\Gamma_{T+1}))%
}

\NovelDisplay{M01122R01}{%
\Pi\text{-}I(1,0)\text{-}\Pi(1,0) \sim \psi_a(\varphi(T,\Gamma_{T+1}))%
}

\NovelDisplay{M01123R01}{%
\Pi(1,0)\text{-}\Pi(1,0) \sim \psi_a(\Gamma_{T+2})%
}

\NovelDisplay{M01124R01}{%
\Pi_{\Pi(1,0)+1} \sim \psi_a(\Gamma_{T\times 2})%
}

``在 \allowbreak{}\NovelInline{I04984}{\Gamma_{T+1}}\allowbreak{} 的基础上，所有的规则再次重新生效。''蔡兴坤的声音发颤，``\allowbreak{}\NovelInline{I04985}{\varepsilon}\allowbreak{}、\allowbreak{}\NovelInline{I04986}{\varphi(a,\cdot)}\allowbreak{}、\allowbreak{}\NovelInline{I04987}{\varphi(T,\cdot)}\allowbreak{}、\allowbreak{}\NovelInline{I04988}{\Gamma}\allowbreak{} 的递增、\allowbreak{}\NovelInline{I04989}{\Gamma}\allowbreak{} 的乘法递归------全部重现。''

\NovelDisplay{M01125R01}{%
\Pi_{\Pi(1,0)\times\omega} \sim \psi_a(\varphi(1,1,T+1))%
}

\NovelDisplay{M01126R01}{%
\Pi\text{-}\Phi(1,\Pi(1,0)+1) \sim \psi_a(\varphi(1,a,T+1))%
}

\NovelDisplay{M01127R01}{%
\Pi\text{-}I(1,\Pi(1,0)+1) \sim \psi_a(\varphi(1,T,1))%
}

\NovelDisplay{M01128R01}{%
\Pi(1,1) \sim \psi_a(\varphi(2,0,T+1))%
}

\NovelDisplay{M01129R01}{%
\Pi(1,x)\text{-}\Phi(1,0) \sim \psi_a(\varphi(a,0,T+1))%
}

phyrion:这里，\allowbreak{}\NovelInline{I04990}{\Pi(1,x)\text{-}\Phi(1,0)}\allowbreak{}表示\allowbreak{}\NovelInline{I04991}{f(α)→\Pi(1,α)}\allowbreak{}的不动点计数

\NovelDisplay{M01130R01}{%
\Pi(1,x)\text{-}I(1,0) \sim \psi_a(\varphi(T,0,1))%
}

\NovelDisplay{M01131R01}{%
\Pi(2,0) \sim \psi_a(\varphi(1,0,0,T+1))%
}

\NovelDisplay{M01132R01}{%
\Pi(\omega,0) \sim \psi_a(\varphi(1@\omega,T+1))%
}

\NovelDisplay{M01133R01}{%
\Pi(x,0)\text{-}\Phi(1,0) \sim \psi_a(\varphi(1@a,T+1))%
}

\NovelDisplay{M01134R01}{%
\Pi(x,0)\text{-}I(1,0) \sim \psi_a(\varphi(1@T,1))%
}

\NovelDisplay{M01135R01}{%
\Pi(1,0,0) \sim \psi_a(\varphi(1@(1,0),T+1))%
}

\NovelDisplay{M01136R01}{%
\Pi(1,0,0,0) \sim \psi_a(\varphi(1@(1,0,0),T+1))%
}

\NovelDisplay{M01137R01}{%
\Pi(1@\omega) \sim \psi_a(\varphi(1@(1@\omega),T+1))%
}

\NovelDisplay{M01138R01}{%
\Pi(1@x)\text{-}\Phi(1,0) \sim \psi_a(\varphi(1@1@a,T+1))%
}

\NovelDisplay{M01139R01}{%
\Pi(1@x)\text{-}I(1,0) \sim \psi_a(\varphi(1@1@T,1))%
}

\NovelDisplay{M01140R01}{%
\Pi(1@(1,0)) \sim \psi_a(\varphi(1@1@(1,0),T+1))%
}

Phyrion的笔尖几乎没有停顿。一行又一行的公式如瀑布般倾泻而下，密密麻麻地铺满了整面黑板。

``\allowbreak{}\NovelInline{I04992}{\Pi(1,\cdot)}\allowbreak{} 的参数变化，对应 \allowbreak{}\NovelInline{I04993}{\varphi}\allowbreak{} 函数多参数结构的演化。''唐三的声音已经变得嘶哑，``从双参数到三参数，从 \allowbreak{}\NovelInline{I04994}{@}\allowbreak{} 记法到多逗号记法------\allowbreak{}\NovelInline{I04995}{\Pi}\allowbreak{} 记法的表达能力，远远超出了我们的想象。''

最后，Phyrion写下了表格的最后三行：

\NovelDisplay{M01141R01}{%
\Pi(1,,0) \sim \psi_a(\varphi(1@(1,,0),T+1))=\psi_a(\Omega_{a+2}+T)%
}

他放下笔，转过身来。

修炼堂里鸦雀无声。

Phyrion的目光扫过每一个人的脸。

``这张表，从 \allowbreak{}\NovelInline{I04996}{\Pi_\omega}\allowbreak{} 到 \allowbreak{}\NovelInline{I04997}{\Pi(1,,0)}\allowbreak{}，涵盖了投影记号体系中几乎所有重要的序数结构。''他的声音低沉而郑重，``你们现在看到的，不是零散的公式------而是一套完整的、自相似的、层层递进的记号体系。''

``我要问你们------在这一整段推导中，你们感受到了什么？''

沉默了片刻，蔡须坤率先开口了。

``自相似性。''她的声音清亮而笃定，``从 \allowbreak{}\NovelInline{I04998}{\Pi_\omega}\allowbreak{} 到 \allowbreak{}\NovelInline{I04999}{\Pi_{\omega+1}}\allowbreak{}，从 \allowbreak{}\NovelInline{I05000}{\Pi_{\omega+1}}\allowbreak{} 到 \allowbreak{}\NovelInline{I05001}{\Pi\text{-}\Phi(1,0)}\allowbreak{}，从 \allowbreak{}\NovelInline{I05002}{\Pi\text{-}\Phi(1,0)}\allowbreak{} 到 \allowbreak{}\NovelInline{I05003}{\Pi(1,0)}\allowbreak{}------每一次跃迁之后，所有的展开规则都会在新基底上重新生效。乘法、指数塔、\allowbreak{}\NovelInline{I05004}{\varepsilon}\allowbreak{}、\allowbreak{}\NovelInline{I05005}{\zeta}\allowbreak{}、\allowbreak{}\NovelInline{I05006}{\varphi}\allowbreak{}、大基数------每一层结构都是上一层的精确复刻。''

``这种自相似性，不是巧合------它是投影记号体系的内在法则。''

``很好。''Phyrion点头，``还有呢？''

唐三接道：``记号之间的对应关系。\allowbreak{}\NovelInline{I05007}{\Pi}\allowbreak{} 记法和投影记号是同一枚硬币的两面。左边的每一个 \allowbreak{}\NovelInline{I05008}{\Pi}\allowbreak{} 表达式，都在右边找到了精确的投影对应。而这种对应不是任意的------它遵循着严格的、可推导的规则。''

``理解了这种对应，就理解了投影与反射的本质区别。''

``很好。''Phyrion的目光中闪过赞许，``还有呢？''

蔡兴坤沉默了片刻，然后缓缓说道：``扽力。''

``在这一整段推导中，我能感觉到扽力在不断地被拉伸。''他的声音变得低沉，``每跨越一个新的层级，我的灵魂就被扩展一分。尤其是从 \allowbreak{}\NovelInline{I05009}{\Pi\text{-}\Phi(1,0)}\allowbreak{} 到 \allowbreak{}\NovelInline{I05010}{\Pi(1,0)}\allowbreak{}，从 \allowbreak{}\NovelInline{I05011}{\Pi(1,0)}\allowbreak{} 到 \allowbreak{}\NovelInline{I05012}{\Pi(1,,0)}\allowbreak{}------每一次跃迁，扽力的增长都是之前所有增长的总和。''

``这就是扽西的本质。''Phyrion的声音变得郑重，``你们不是在简单地记忆公式------你们是在用灵魂去分析、去拆解、去透视这些序数结构。每一次分析，都是对灵魂的一次锤炼。''

``而差异越大的记号对，扽西之后的提升就越恐怖。''他顿了顿，``从 \allowbreak{}\NovelInline{I05013}{\Pi_\omega}\allowbreak{} 到 \allowbreak{}\NovelInline{I05014}{\Pi(1,,0)}\allowbreak{}，你们扽西的记号对之间的差异，已经远远超过了普通修炼者一生的积累。''

``你们的扽力，在今天获得了质的飞跃。''

众人沉默了。

他们能感受到------Phyrion说得没错。今天的扽西训练，让他们的扽力获得了前所未有的提升。那种灵魂被拉伸、被扩展的感觉，既痛苦又畅快。

蔡须坤低下头，看着自己的双手。金色的光芒在她的指尖流转，比昨日更加明亮、更加纯粹。帝皇瑞兽的本源处女之力在她体内翻涌，像是即将破茧的蝴蝶。

蔡兴坤摸了摸肩膀上的黑子鸡。小黑鸡歪着脑袋看着他，圆溜溜的眼睛里满是无辜。但蔡兴坤知道，在这副蠢萌的外表之下，封印着连Phyrion都无法窥探的未知力量。

``牢大\ldots\ldots''他在心里默默说道，``我们又在变强了。''

Phyrion走到门口，回头看了众人一眼。

``回去好好休息。''他说，``这张表的内容，足够你们消化很长时间了。''

他顿了顿，嘴角微微上扬。

``但记住------这只是 \allowbreak{}\NovelInline{I05015}{\Pi}\allowbreak{} 阶梯的全貌。真正的修炼，才刚刚开始。''

石门关闭。

窗外，午后的阳光斜斜地洒进修炼堂，将黑板上密密麻麻的公式染成了金红色。

从 \allowbreak{}\NovelInline{I05016}{\Pi_\omega}\allowbreak{} 到 \allowbreak{}\NovelInline{I05017}{\Pi(1,,0)}\allowbreak{}------那是一条用数字和符号铺就的阶梯，从最朴素的 \allowbreak{}\NovelInline{I05018}{\varepsilon_{T+1}}\allowbreak{} 出发，一路攀升至 \allowbreak{}\NovelInline{I05019}{\psi_a(\Omega_{a+2}+T)}\allowbreak{} 的巅峰。

而在这条阶梯的尽头，还有更高的山峰在等待着他们。

\chapter{第八十七章：殊途同归的两条路}\label{ux7b2cux516bux5341ux4e03ux7ae0ux6b8aux9014ux540cux5f52ux7684ux4e24ux6761ux8def}

修炼堂已经三天没有亮过灯了。

唐三靠在墙角，盯着天花板上某处斑驳的纹路发呆。那上面有一道细小的裂缝，蜿蜒曲折，像极了他在黑板上写过的某条推导链------但最终通向的，不过是一个 \allowbreak{}\NovelInline{I05020}{+T}\allowbreak{}。

蔡须坤坐在窗边，金色的光芒在她指尖若隐若现，微弱得像是随时会熄灭的烛火。她没有像往常那样主动开口分析，也没有去翻看笔记。那些密密麻麻的公式此刻在她眼中不再是阶梯，而是一座牢笼------他们拼尽全力攀上去，回头一看，脚下不过是沧海一粟。

蔡兴坤蹲在角落里，黑子鸡趴在他肩膀上，一人一鸡都蔫头耷脑。

``\allowbreak{}\NovelInline{I05021}{\psi_a(\Omega_{a+2}+T\times\omega)}\allowbreak{}。''他喃喃地重复着这句话，像是在咀嚼一块怎么也咽不下去的苦药，``我们把整个 \allowbreak{}\NovelInline{I05022}{\Pi}\allowbreak{} 阶梯从底到顶全部拆解完毕，分析了这么多记号对\ldots\ldots 结果就是让下标从 \allowbreak{}\NovelInline{I05023}{\Omega_{a+2}}\allowbreak{} 变成了 \allowbreak{}\NovelInline{I05024}{\Omega_{a+2}+T\times\omega}\allowbreak{}。''

他抬起头，目光空洞。

``\allowbreak{}\NovelInline{I05025}{\psi_a(\Omega_{a+2}\times\omega)}\allowbreak{} 到底有多大？我们连想都不敢想。''

沉默像潮水一样漫过整个房间。

直到第四天清晨，石门的缝隙中透进一道光。

Phyrion走了进来。他没有拿粉笔，也没有走向黑板。他只是站在门口，看着这群消沉的年轻人，沉默了很久。

然后他说了一句话：

``你们觉得 \allowbreak{}\NovelInline{I05026}{+T}\allowbreak{} 很小，是因为你们还站在 \allowbreak{}\NovelInline{I05027}{T}\allowbreak{} 的高度往下看。''

没有人抬头。

``等你们有一天站在 \allowbreak{}\NovelInline{I05028}{\Omega_{a+2}\times\omega}\allowbreak{} 的高度回头看------''他的声音很轻，却清晰地落入每一个人的耳中，``你们会明白，今天这个让你们沮丧的 \allowbreak{}\NovelInline{I05029}{+T}\allowbreak{}，是后来所有跃迁的基石。''

他转身离去，石门在身后缓缓合拢。

唐三的手指微微动了一下。

蔡须坤抬起头，看向窗外。晨光正从地平线上涌出，将远处的山脊镀上一层金边。

蔡兴坤肩膀上的黑子鸡歪了歪脑袋，发出一声细小的鸣叫。

没有人说话。但有什么东西，在这间昏暗了三天的修炼堂里，悄悄地重新亮了起来。

\begin{center}\rule{0.5\linewidth}{0.5pt}\end{center}

这一切的根源，在于 \allowbreak{}\NovelInline{I05030}{\psi_a(\Omega_{a+2})}\allowbreak{} 之大------它足以折叠 \allowbreak{}\NovelInline{I05031}{\psi_a(\Omega_{a+1})}\allowbreak{} 的递归运算。因为 \allowbreak{}\NovelInline{I05032}{\psi_a(\Omega_{a+2})}\allowbreak{} 相当于多逗号的 \allowbreak{}\NovelInline{I05033}{\varphi(1,,,0)}\allowbreak{} 级别，正如 \allowbreak{}\NovelInline{I05034}{\psi_a(\Omega_{a+1})}\allowbreak{} 之大足以折叠 \allowbreak{}\NovelInline{I05035}{\psi_a(a\text{ 的递归运算})}\allowbreak{} 一样，层级之间的鸿沟，每上一阶便是一道天堑。

Phyrion的声音仿佛还在石壁间回荡。他此前留下的那个问题，此刻终于被众人重新拾起------

\NovelDisplay{M01142R01}{%
\psi(\psi_a(\Omega_{a+2}\times \omega))%
}

的 lift 展开。然而这一次，他给出的不再是单一的表达式，而是两条看似平行的路径：

\textbf{路径一：} \NovelDisplay{M01143R01}{%
\psi(\psi_a(\psi_{\Omega_{a+2}}(\psi_{a_2}(\psi_{\Omega_{a_2+2}}(\psi_{a_3}(\cdots))))))%
}

\textbf{路径二：} \NovelDisplay{M01144R01}{%
\psi(\psi_a(\Omega_{a+2}+\psi_{\Omega_{a+2}}(\psi_{a_2}(\Omega_{a_2+2}+\psi_{\Omega_{a_2+2}}(\psi_{a_3}(\cdots))))))%
}

\textbf{路径三：} \NovelDisplay{M01145R01}{%
\psi(\psi_a(\Omega_{a+2}+\psi_{\Omega_{a+1}}(\psi_{a_2}(\Omega_{a_2+2}+\psi_{\Omega_{a_2+1}}(\psi_{a_3}(\cdots))))))%
}

``这两条路------''Phyrion的声音在石壁间回荡，``是不是一样的？''

没有人说话。直觉告诉他们，从结构上看，二者确实相似。但''相似''二字太轻了，他们需要的是一把精确的手术刀，将每一层嵌套剖开，看看骨头里到底长着什么。

\begin{center}\rule{0.5\linewidth}{0.5pt}\end{center}

蔡须坤率先打破了沉默。她站起身，走到黑板前，指尖凝出一缕金光，在石面上刻下第一行字。

``先看起点。''她的声音恢复了往日的沉稳，``路径一的最内层------''

\NovelDisplay{M01146R01}{%
\psi(\psi_a(\psi_{\Omega_{a+2}}(\psi_{a_2}(\psi_{\Omega_{a_2+2}}(0)))))%
}

``\allowbreak{}\NovelInline{I05036}{\psi_{\Omega_{a_2+2}}(0)=\Omega_{a_2+2}}\allowbreak{}，往上折叠一层------\allowbreak{}\NovelInline{I05037}{\psi_{a_2}(\Omega_{a_2+2})=\Omega_{a_2+2}}\allowbreak{}，再往上------\allowbreak{}\NovelInline{I05038}{\psi_{\Omega_{a+2}}(\Omega_{a_2+2})=\Omega_{a+2}}\allowbreak{}。所以起点就是------''

\NovelDisplay{M01147R01}{%
\psi(\psi_a(\Omega_{a+2}))%
}

她顿了顿，目光扫过众人。``而路径二的起点，同样是 \allowbreak{}\NovelInline{I05039}{\psi(\psi_a(\Omega_{a+2}))}\allowbreak{}。起点一样，不过------''

``这仅仅是起点。''她加重了语气，``接下来，才是重点。''

\begin{center}\rule{0.5\linewidth}{0.5pt}\end{center}

小舞接过话头，她的手指在空气中划出一道道光线，将下一步展开投影在半空中：

``路径一的第二步------在 \allowbreak{}\NovelInline{I05040}{\psi_{\Omega_{a_2+2}}}\allowbreak{} 内部填入 \allowbreak{}\NovelInline{I05041}{\Omega_{a_2+2}}\allowbreak{}：

\NovelDisplay{M01148R01}{%
\psi(\psi_a(\psi_{\Omega_{a+2}}(\psi_{a_2}(\psi_{\Omega_{a_2+2}}(\Omega_{a_2+2})))))%
}

\allowbreak{}\NovelInline{I05042}{\psi_{\Omega_{a_2+2}}(\Omega_{a_2+2})=\Omega_{a_2+2}}\allowbreak{}，所以内层不变。但 \allowbreak{}\NovelInline{I05043}{\psi_{a_2}(\Omega_{a_2+2})}\allowbreak{}------``她眯起眼睛，''这是 \allowbreak{}\NovelInline{I05044}{\varepsilon_{\Omega_{a_2+1}}}\allowbreak{}。''

蔡兴坤猛地抬起头。``不对------\allowbreak{}\NovelInline{I05045}{\psi_{a_2}(\Omega_{a_2+2})}\allowbreak{} 应该是 \allowbreak{}\NovelInline{I05046}{\Omega_{a_2+2}}\allowbreak{} 的下一个 \allowbreak{}\NovelInline{I05047}{\psi_{a_2}}\allowbreak{} 不动点\ldots\ldots{}''

``不，''小舞摇了摇头，``你看，\allowbreak{}\NovelInline{I05048}{\psi_{a_2}}\allowbreak{} 的临界点就是 \allowbreak{}\NovelInline{I05049}{\Omega_{a_2+1}}\allowbreak{}，而 \allowbreak{}\NovelInline{I05050}{\Omega_{a_2+2}}\allowbreak{} 远在它之上。所以 \allowbreak{}\NovelInline{I05051}{\psi_{a_2}(\Omega_{a_2+2})=\varepsilon_{\Omega_{a_2+1}}}\allowbreak{}------这是 \allowbreak{}\NovelInline{I05052}{\psi_{a_2}}\allowbreak{} 在 \allowbreak{}\NovelInline{I05053}{\Omega_{a_2+2}}\allowbreak{} 处的折叠结果。''

她继续向上推进：``再往上一层------\allowbreak{}\NovelInline{I05054}{\psi_{\Omega_{a+2}}(\varepsilon_{\Omega_{a_2+1}})}\allowbreak{}。这里 \allowbreak{}\NovelInline{I05055}{\varepsilon_{\Omega_{a_2+1}}}\allowbreak{} 远小于 \allowbreak{}\NovelInline{I05056}{\Omega_{a+2}}\allowbreak{} 的临界结构，所以折叠后------\allowbreak{}\NovelInline{I05057}{\psi_{\Omega_{a+2}}(\varepsilon_{\Omega_{a_2+1}})=\Omega_{a+2}}\allowbreak{}。''

``等等。''蔡须坤皱眉，``你是说它直接坍缩回 \allowbreak{}\NovelInline{I05058}{\Omega_{a+2}}\allowbreak{}？''

``不------''小舞咬了咬嘴唇，重新审视自己的推导，``我犯了个错误。\allowbreak{}\NovelInline{I05059}{\varepsilon_{\Omega_{a_2+1}}}\allowbreak{} 虽然在 \allowbreak{}\NovelInline{I05060}{\Omega_{a+2}}\allowbreak{} 面前很小，但 \allowbreak{}\NovelInline{I05061}{\psi_{\Omega_{a+2}}}\allowbreak{} 的作用不是简单地坍缩，它是在 \allowbreak{}\NovelInline{I05062}{\Omega_{a+2}}\allowbreak{} 的框架内寻找不动点。''

她深吸一口气，重新来过。

``\allowbreak{}\NovelInline{I05063}{\psi_{a_2}(\Omega_{a_2+2})}\allowbreak{} 的结果，应当是 \allowbreak{}\NovelInline{I05064}{\Omega_{a_2+2}}\allowbreak{} 层级内的某个大序数。考虑到 \allowbreak{}\NovelInline{I05065}{a_2}\allowbreak{} 层级的 \allowbreak{}\NovelInline{I05066}{\psi}\allowbreak{} 函数在 \allowbreak{}\NovelInline{I05067}{\Omega_{a_2+2}}\allowbreak{} 处的行为------''

``是 \allowbreak{}\NovelInline{I05068}{\varepsilon_{\Omega_{a_2+1}}}\allowbreak{}。''蔡须坤平静地纠正道，``这一点小舞没有说错。关键是下一步------\allowbreak{}\NovelInline{I05069}{\psi_{\Omega_{a+2}}(\varepsilon_{\Omega_{a_2+1}})}\allowbreak{}。''

沉默了片刻。

``\allowbreak{}\NovelInline{I05070}{\varepsilon_{\Omega_{a_2+1}}}\allowbreak{} 在 \allowbreak{}\NovelInline{I05071}{\Omega_{a+2}}\allowbreak{} 的尺度下，''蔡须坤缓缓说道，``确实太小了。\allowbreak{}\NovelInline{I05072}{\psi_{\Omega_{a+2}}}\allowbreak{} 在这个位置的取值，就是 \allowbreak{}\NovelInline{I05073}{\Omega_{a+2}}\allowbreak{} 本身。''

``所以路径一的第二步------''唐三接过话头，``实际上给出的信息量是 \allowbreak{}\NovelInline{I05074}{\psi_{\Omega_{a+2}}(\Omega_{a+2})=\Omega_{a+3}}\allowbreak{}。''

``对。''蔡须坤点头，``\allowbreak{}\NovelInline{I05075}{\Omega_{a+3}}\allowbreak{}。然后 \allowbreak{}\NovelInline{I05076}{\psi_a(\Omega_{a+3})}\allowbreak{}------这就是 BHO 在 \allowbreak{}\NovelInline{I05077}{\Omega_{a+1}+1}\allowbreak{} 层级的位置。''

\NovelDisplay{M01149R01}{%
\psi(\psi_a(\Omega_{a+3}))=\psi(\psi_a(BHO_{\Omega_{a+1}+1}))%
}

\begin{center}\rule{0.5\linewidth}{0.5pt}\end{center}

``但是有个问题。''卢本韦突然开口，他的声音沙哑，像是很久没有说话了，``你们看路径二在相应位置的值。''

他在空中划出路径二的对应步骤：

\NovelDisplay{M01150R01}{%
\psi(\psi_a(\Omega_{a+2}+\psi_{\Omega_{a+1}}(\psi_{a_2}(\Omega_{a_2+1}\times \omega))))%
}

``\allowbreak{}\NovelInline{I05078}{\psi_{a_2}(\Omega_{a_2+1}\times \omega)}\allowbreak{}------这个值，在 \allowbreak{}\NovelInline{I05079}{\Omega_{a+1}}\allowbreak{} 的框架下，被 \allowbreak{}\NovelInline{I05080}{\psi_{\Omega_{a+1}}}\allowbreak{} 折叠后，仅仅是 \allowbreak{}\NovelInline{I05081}{\Omega_{a+1}\times \omega}\allowbreak{}。''

他看向众人，目光锐利。

``也就是说------路径二在这里给出的，只是 \allowbreak{}\NovelInline{I05082}{\psi(\psi_a(\Omega_{a+2}+\Omega_{a+1}\times \omega))}\allowbreak{}。''

\begin{center}\rule{0.5\linewidth}{0.5pt}\end{center}

霍雨浩的眼睛突然亮了。

``我明白了！''他猛地站起来，差点碰翻身后的椅子，``你们看------\allowbreak{}\NovelInline{I05083}{\psi(\psi_a(\Omega_{a+2}+\psi_{\Omega_{a+2}}(\psi_{a_2}(\Omega_{a_2+1}\times \omega))))}\allowbreak{}，这个表达式，它的结构对应的是什么？''

他抓起粉笔，在黑板上飞快地写下：

\NovelDisplay{M01151R01}{%
\psi(\psi_a(\Omega_{a+2}+BHO_{\Omega_{a+1}+1}))%
}

``这和 \allowbreak{}\NovelInline{I05084}{\psi(\psi_a(\Omega_{a+1}+BHO_{a+1}))}\allowbreak{} 是同构的！而在我们之前的分析中------''

``\allowbreak{}\NovelInline{I05085}{\psi(\psi_a(\Omega_{a+1}+BHO_{a+1}))=\psi(\psi_a(\Omega_{a+1}\times 2+\varepsilon_{a+1}))}\allowbreak{}。''蔡须坤接过他的话，声音微微发颤，``所以\ldots\ldots{}''

``\allowbreak{}\NovelInline{I05086}{\psi(\psi_a(\Omega_{a+2}+BHO_{\Omega_{a+1}+1}))=\psi(\psi_a(\Omega_{a+2}\times 2+\varepsilon_{\Omega_{a+1}}))}\allowbreak{}。''霍雨浩一字一顿地说完，缓缓坐回椅子上。

\begin{center}\rule{0.5\linewidth}{0.5pt}\end{center}

诺当一直没有说话。此刻他缓缓举起手，像是在课堂上提问的学生。

``即使是 \allowbreak{}\NovelInline{I05087}{\psi(\psi_a(BO_{\Omega_{a+1}+1}))}\allowbreak{}------''他的声音很轻，但每个字都清晰可闻，``也就是路径一继续推进------

\NovelDisplay{M01152R01}{%
\psi(\psi_a(\psi_{\Omega_{a+2}}(\psi_{a_2}(\psi_{\Omega_{a_2+2}}(0)\times \omega))))%
}

它的结果，也才仅仅是------''

\NovelDisplay{M01153R01}{%
\psi(\psi_a(\Omega_{a+2}+\psi_{\Omega_{a+2}}(\psi_{a_2}(\Omega_{a_2+1}\times \omega))))%
}

``和路径二在同一步给出的结果，一模一样。''

小舞轻轻叹了口气：``对于路径二更深层的嵌套------\allowbreak{}\NovelInline{I05088}{\psi(\psi_a(\Omega_{a+2}+\psi_{\Omega_{a+1}}(\psi_{a_2}(\Omega_{a_2+2}+\psi_{\Omega_{a_2+1}}(\psi_{a_3}(\cdots))))))}\allowbreak{}------我们甚至都不知道它最终会通向哪里。''

\begin{center}\rule{0.5\linewidth}{0.5pt}\end{center}

蔡兴坤肩膀上的黑子鸡突然拍了拍翅膀，像是被什么激灵了一下。蔡兴坤自己也坐直了身子。

``换个角度想。''他摸了摸下巴，``考虑 \allowbreak{}\NovelInline{I05089}{\psi(\psi_a(\Omega_{a+2}+\psi_{\Omega_{a+1}}(\psi_{a_2}(\Omega_{a_2+1}\times a_2))))}\allowbreak{}------这是个不动点。你们觉得，这个不动点，会不会恰好就是 \allowbreak{}\NovelInline{I05090}{\psi(\psi_a(\Omega_{a+2}+\Omega_{a+1}\times a))}\allowbreak{}？''

蔡须坤立刻摇头。``作为队长，我有必要提醒各位------不能只是猜测。我们应当严格分析。''她转向黑板，写下一个新的表达式：

``对于 \allowbreak{}\NovelInline{I05091}{\psi(\psi_a(\Omega_{a+2}+\psi_{\Omega_{a+1}}(\psi_{a_2}(\Omega_{a_2+1}\times \psi_a(\Omega_{a+2})))))}\allowbreak{}------这个值是什么？''

\begin{center}\rule{0.5\linewidth}{0.5pt}\end{center}

唐三沉默了很久。他的目光落在天花板上那道蜿蜒的裂缝上，像是在沿着某条看不见的推导链向上攀爬。

``我觉得，''他终于开口，声音不大，却异常笃定，``应当是 \allowbreak{}\NovelInline{I05092}{\psi(\psi_a(\Omega_{a+2}+\Omega_{a+1}\times \psi_a(\Omega_{a+2})))}\allowbreak{}。''

``理由呢？''蔡须坤追问。

``\allowbreak{}\NovelInline{I05093}{\psi_{\Omega_{a+1}}}\allowbreak{} 在折叠 \allowbreak{}\NovelInline{I05094}{\psi_{a_2}(\Omega_{a_2+1}\times \psi_a(\Omega_{a+2}))}\allowbreak{} 时，\allowbreak{}\NovelInline{I05095}{\psi_a(\Omega_{a+2})}\allowbreak{} 虽然在绝对意义上极大，但在 \allowbreak{}\NovelInline{I05096}{\Omega_{a+1}}\allowbreak{} 的框架内，它只是一个'系数'。\allowbreak{}\NovelInline{I05097}{\psi_{a_2}}\allowbreak{} 将 \allowbreak{}\NovelInline{I05098}{\Omega_{a_2+1}\times \psi_a(\Omega_{a+2})}\allowbreak{} 折叠后，得到的结果在 \allowbreak{}\NovelInline{I05099}{\psi_{\Omega_{a+1}}}\allowbreak{} 的作用下，最终呈现为 \allowbreak{}\NovelInline{I05100}{\Omega_{a+1}\times \psi_a(\Omega_{a+2})}\allowbreak{}。''

\begin{center}\rule{0.5\linewidth}{0.5pt}\end{center}

霍雨浩紧接着补充：``但是这里有一个关键的思考------\allowbreak{}\NovelInline{I05101}{\psi(\psi_a(\Omega_{a+2}+\Omega_{a+1}\times \psi_a(\Omega_{a+2})))}\allowbreak{} 的递归运算，只需要 \allowbreak{}\NovelInline{I05102}{\psi(\psi_a(\Omega_{a+2}+\Omega_{a+1}\times \psi_a(\Omega_{a+2})+\Omega_{a+1}))}\allowbreak{} 就可以折叠。''

他在黑板上画了一条箭头，指向路径二的对应位置：

``这应当对应于 \allowbreak{}\NovelInline{I05103}{\psi(\psi_a(\Omega_{a+2}+\psi_{\Omega_{a+1}}(\psi_{a_2}(\Omega_{a_2+1}\times \psi_a(\Omega_{a+2})+\Omega_{a_2+1}))))}\allowbreak{}。''

小舞点头，目光灼灼：``是的。因为在最外层，是 \allowbreak{}\NovelInline{I05104}{1}\allowbreak{}-投影序数的指数塔。在 \allowbreak{}\NovelInline{I05105}{\psi_{\Omega_{a+1}}}\allowbreak{} 函数内部，往往代表着下一个 \allowbreak{}\NovelInline{I05106}{\Omega}\allowbreak{} 点------而加上 \allowbreak{}\NovelInline{I05107}{\Omega_{a_2+1}}\allowbreak{}，正是'下一个 \allowbreak{}\NovelInline{I05108}{\Omega}\allowbreak{} 点'的意思。两层嵌套的对应关系，在这里完美吻合。''

\begin{center}\rule{0.5\linewidth}{0.5pt}\end{center}

``接下来就有不一样的地方了。''卢本韦的声音突然变得急促，像是发现了什么关键的岔路口。

``\allowbreak{}\NovelInline{I05109}{\psi(\psi_a(\Omega_{a+2}+\psi_{\Omega_{a+1}}(\psi_{a_2}(\Omega_{a_2+1}\times \psi_a(\Omega_{a+2}+1)))))}\allowbreak{}------这个表达式，虽然等于 \allowbreak{}\NovelInline{I05110}{\psi(\psi_a(\Omega_{a+2}+\Omega_{a+1}\times \psi_a(\Omega_{a+2}+1)))}\allowbreak{}------''

他停顿了一下，像是在确认自己的判断。

``但是在它\textbf{前面}，还隔着一整片我们尚未跨越的区域------\allowbreak{}\NovelInline{I05111}{\psi(\psi_a(\Omega_{a+2}+\Omega_{a+1}\times \psi_a(\Omega_{a+2}+\Omega_{a+1})))}\allowbreak{}。''

奥斯卡皱起眉头：``小韦，你的意思是------\allowbreak{}\NovelInline{I05112}{\psi(\psi_a(\Omega_{a+2}+\psi_{\Omega_{a+1}}(\psi_{a_2}(\Omega_{a_2+1}\times a))))}\allowbreak{} 可能并不是我们之前猜测的那个值，而是 \allowbreak{}\NovelInline{I05113}{\psi(\psi_a(\Omega_{a+2}+\Omega_{a+1}\times \psi_a(\Omega_{a+2}+\Omega_{a+1})))}\allowbreak{}？''

``对。''卢本韦重重地点头，``因为有 \allowbreak{}\NovelInline{I05114}{\psi_{\psi_a(\Omega_{a+2}+\Omega_{a+1})}}\allowbreak{} 函数从中作祟。''

他转身面向黑板，粉笔在石面上划出刺耳的声响：

``这就导致------\allowbreak{}\NovelInline{I05115}{\psi(\psi_a(\Omega_{a+2}+\psi_{\Omega_{a+1}}(\psi_{a_2}(\Omega_{a_2+1}\times \psi_a(\Omega_{a+2}+\psi_{\Omega_{a+1}}(\psi_{a_2}(\Omega_{a_2+1}\times \omega)))))))}\allowbreak{}------''

他深吸一口气，写下最终结果：

``实际上仅仅是 \allowbreak{}\NovelInline{I05116}{\psi(\psi_a(\Omega_{a+2}+\Omega_{a+1}\times \psi_{\psi_a(\Omega_{a+2}+\Omega_{a+1})}(\psi_a(\Omega_{a+2}+\Omega_{a+1}\times \omega))))}\allowbreak{}。''

\begin{center}\rule{0.5\linewidth}{0.5pt}\end{center}

唐三盯着黑板上的表达式，瞳孔微微收缩。

``所以说------''他的声音低沉而缓慢，像是在确认一个令人震惊的事实，``\allowbreak{}\NovelInline{I05117}{\psi(\psi_a(\Omega_{a+2}+\psi_{\Omega_{a+1}}(\cdots)))}\allowbreak{} 根本到不了 \allowbreak{}\NovelInline{I05118}{\psi_a(\Omega_{a+2}+\Omega_{a+1})}\allowbreak{}。''

他转身面对众人。

``因此------\allowbreak{}\NovelInline{I05119}{\psi(\psi_a(\Omega_{a+2}+\psi_{\Omega_{a+1}}(\psi_{a_2}(\Omega_{a_2+1}\times a))))=\psi(\psi_a(\Omega_{a+2}+\Omega_{a+1}\times \psi_a(\Omega_{a+2}+\Omega_{a+1})))}\allowbreak{}。''

奥斯卡猛地拍了一下大腿：``我领悟了！那么------\allowbreak{}\NovelInline{I05120}{\psi(\psi_a(\Omega_{a+2}+\psi_{\Omega_{a+1}}(\psi_{a_2}(\Omega_{a_2+1}\times \Omega_{a+1}))))}\allowbreak{} 将会是------''

\NovelDisplay{M01154R01}{%
\psi(\psi_a(\Omega_{a+2}+\Omega_{a+1}\times \psi_a(\Omega_{a+2}+\Omega_{a+1}\times 2)))%
}

``对！''蔡须坤的眼中闪过一道金光，``而之后------\allowbreak{}\NovelInline{I05121}{\psi(\psi_a(\Omega_{a+2}+\psi_{\Omega_{a+1}}(\psi_{a_2}(\Omega_{a_2+1}\times \psi_{a_2}(\Omega_{a_2+1}\times \omega)))))}\allowbreak{}------才是------''

\NovelDisplay{M01155R01}{%
\psi(\psi_a(\Omega_{a+2}+\Omega_{a+1}\times \psi_a(\Omega_{a+2}+\Omega_{a+1}\times \omega)))%
}

\begin{center}\rule{0.5\linewidth}{0.5pt}\end{center}

整个修炼堂安静了下来。只有粉笔在石面上最后划动的声音。

蔡须坤走到黑板中央，将所有结论汇总，一笔一划地写下：

\NovelDisplay{M01156R01}{%
\psi(\psi_a(\Omega_{a+2}+\psi_{\Omega_{a+1}}(\psi_{a_2}(\Omega_{a_2+1}\times a_2))))=\psi(\psi_a(\Omega_{a+2}+\Omega_{a+1}\times a))%
}

``以此类推------''

\NovelDisplay{M01157R01}{%
\psi(\psi_a(\Omega_{a+2}+\psi_{\Omega_{a+1}}(\psi_{a_2}(\Omega_{a_2+2}))))=\psi(\psi_a(\Omega_{a+2}\times 2))%
}

她放下粉笔，退后一步。

``这一点，在 \allowbreak{}\NovelInline{I05122}{\psi(\psi_a(\Omega_{a+2}+\psi_{\Omega_{a+2}}(\Omega_{a+2})))}\allowbreak{} 中，同样可以做到。''

``但是------''她的声音突然沉了下来，``对于 \allowbreak{}\NovelInline{I05123}{\psi(\psi_a(\Omega_{a+2}+\psi_{\Omega_{a+2}}(\Omega_{a+2}\times \omega)))}\allowbreak{}------它实际上是 \allowbreak{}\NovelInline{I05124}{\psi(\psi_a(\Omega_{a+2}\times 2+\Omega_{a+1}\times \omega))}\allowbreak{}。''

\begin{center}\rule{0.5\linewidth}{0.5pt}\end{center}

唐三缓缓站起身。他走到黑板前，盯着那两行并列的等式，目光从路径一的结论缓缓移向路径二的结论。

然后他轻声说：

``没想到啊。''

\NovelDisplay{M01158R01}{%
\psi(\psi_a(\Omega_{a+2}+\psi_{\Omega_{a+1}}(\psi_{a_2}(\Omega_{a_2+\omega}))))=\psi(\psi_a(\Omega_{a+2}\times 2+\Omega_{a+1}\times \omega))%
}

他转过身，面对所有人。

``也就是说------这两条看似截然相反的路，通向的是同一个地方。''

\begin{center}\rule{0.5\linewidth}{0.5pt}\end{center}

震惊像一道无声的闪电，劈过每个人的心头。

两条路，方向截然相反，却在终点轰然汇合。

``殊途同归\ldots\ldots{}''蔡兴坤喃喃道，肩膀上的黑子鸡也发出一声低鸣，像是在附和。

但这份震惊很快被更深的忧虑所取代。因为他们已经触碰到了拓展反射所能到达的极限。再往下------

就是 \allowbreak{}\NovelInline{I05125}{\psi(\psi_a(\Omega_{a+2}\times \Omega_{a+1}\times \omega))}\allowbreak{} 的领域。

而这里，将开始与 \allowbreak{}\NovelInline{I05126}{\psi(\psi_a(\Omega_{a+2}^2 \times \omega))}\allowbreak{} 产生区别。但区别究竟在哪里？没人能说清。

唐三重新靠回墙角，目光再次落向天花板上那道蜿蜒的裂缝。但这一次，他的眼中不再有迷茫。

``接下来的任务------''蔡须坤的声音在修炼堂中回荡，清晰而坚定，``找到 \allowbreak{}\NovelInline{I05127}{\psi(\psi_a(\Omega_{a+2}\times \Omega_{a+1}\times \omega))}\allowbreak{} 和 \allowbreak{}\NovelInline{I05128}{\psi(\psi_a(\Omega_{a+2}^2 \times \omega))}\allowbreak{} 各自的 lift 提升。''

``以及------''她指向黑板上那两条已经汇合的路径的尽头，``\allowbreak{}\NovelInline{I05129}{\psi(\psi_a(\Omega_{a+2}\times\psi_{\Omega_{a+2}}(\psi_{a_2}(\Omega_{a_2+2}\times \psi_{\Omega_{a_2+2}}(\psi_{a_3}(\cdots))))))}\allowbreak{}''

``和------''

``\allowbreak{}\NovelInline{I05130}{\psi(\psi_a(\Omega_{a+2}\times\psi_{\Omega_{a+1}}(\psi_{a_2}(\Omega_{a_2+2}\times \psi_{\Omega_{a_2+1}}(\psi_{a_3}(\cdots))))))}\allowbreak{}''

\chapter{第八十八章：六个月的尽头}\label{ux7b2cux516bux5341ux516bux7ae0ux516dux4e2aux6708ux7684ux5c3dux5934}

皇家修炼堂的石门在清晨第一缕光里推开。

堂内那面黑板还留着昨夜没擦完的推演，粉笔灰积在板槽和地面上。\allowbreak{}\NovelInline{I05131}{\psi(\psi_a(\Omega_{a+2}+\psi_{\Omega_{a+1}}(\psi_{a_2}(\cdots))))}\allowbreak{}，这样一行东西曾经让他们整宿睡不着，现在看过去，只是一行字。

六个月的闭关，到今天为止。

大半年前，大长老把他们八个叫进扽魂殿的修炼堂，亲手在黑板上写下了第一行 BMS。

扽魂殿那间修炼堂里，他教过的人不少。Smile Lee 在那儿坐过，梅天狸在那儿坐过，两百来年前，卢守衡也在那儿坐过------他是卢承业的祖父，后来死于魂灵反噬。这支队伍他早就留意了，大赛上的名次，他看在眼里。

功法讲到快收尾的时候，天斗城出了宫变。碎紫斗罗一路跑上山来报的信。大长老没有亲自出殿。他让真霜斗罗带着他们八个进城，临走从暗格里取出一颗八级定装魂导炮弹，搁在了队长手里。那一夜打下来，八个人一个都没少。雪夜大帝封了他们子爵，赐了一部《投影序数》，又准他们进天斗皇家修炼堂闭关六个月。

他们进修炼堂那天，大长老下了一趟山。他在皇家修炼堂的黑板上把最后一段讲完，把石门一关，自己回了良序圣殿。

走之前，他只留下一句话：二十五岁之前，突破封号斗罗。到了，就册封长老。

扽魂殿的长老，大多是九十五级往上才封的。

蔡须坤第一个走出来。她在门槛外站住，回头看了那面黑板一眼，没说话，伸手把石门带上了。

小舞跟着出来，伸了个懒腰。``我感觉像做了很长一个梦。醒过来发现梦里学的东西全都还在。''

``半年前谁能想到，''奥斯卡靠在门框上，``咱们这群人能站到 0 111 22 以上。''

诺当没搭腔，只是握了握拳。半年前他的气息还是浮的，现在沉下去了，沉得听不见。

霍雨浩推了推眼镜，目光在每个人身上过了一遍，最后停在蔡须坤那里。``队长走得最远。0 111 221。这不是把壁垒捅穿，是换了个地方站。''

``不只是我。''蔡须坤摇头，``卢本韦跟我差不了多少。''

所有人转过去看她。

卢本韦在阶下的树影里，手插在兜里。``别看我，''她说，``该走的路走完了而已。''

蔡兴坤肩上的黑子鸡扑棱了一下翅膀，短促地叫了一声。她抬手把它的毛顺下去，没说话。

唐三最后一个出来。他没看黑板，抬头看了会儿天。

十六岁。他们在十六岁出头这年，把 0 111 22 走穿了。

\begin{center}\rule{0.5\linewidth}{0.5pt}\end{center}

天斗城以北百里，卢氏养猪场。

场子里的人都知道后山那块青石不许搬。至于为什么，问遍全场只有一个人答得上来，而那个人从来不说。

phyrion 盘坐在青石上，膝头摊着一卷纸。

账面上他占着这片场子的第一股，可他一年到头不管账，也很少露面。场里年轻的伙计只当这位股东脾气怪。老一辈的人知道得多一点------知道他姓卢，知道他是老厂长带回来的，也知道整个场子里再没有第二个人比他年纪大。

纸上的符号不属于任何已有的体系，是他自己造的，用来写 1Y 序列。

最后一个符号落下。

他睁开眼。风从后山过来，把纸角掀起又按下，他伸手压住，慢慢折好收进怀里。

1Y 序列的良序化，证完了。半年，一个人，没跟任何人提过。

最后一笔落下去的时候，他身上的魂力往上走了一级。九十八。

他没起身庆祝，只是伸手在身下那块石头上按了按。

这块石头他曾祖坐过。那个人死在一百九十一年前，比他出生早了六十年。他这辈子没见过那张脸，只见过族谱上那个名字，和名字底下那一行小字：

卒于魂灵。

然后他朝东南方向看了很久。那边是海。

``蓬莱岛。''他说，``也该去看看了。''

\begin{center}\rule{0.5\linewidth}{0.5pt}\end{center}

扽魂殿最高处，良序圣殿。十二根白玉柱撑着穹顶，四面没有墙，神像的目光越过整座天斗城。

大长老 test\_alpha0 已经三天没有移动过了。

石台上摊着一卷玉简，刻满了他自己造的符号。iblp 的良序化证明，最后六个月，一笔一笔刻到这里。

他写下最后一行，手停住了。

玉简末尾是那个层次的名字：PTO(Z₃)。

极限斗罗之上，多少人连它的门槛在哪儿都摸不着。

他用了大半辈子的那些东西------ω-Y 序列、iblp、一层一层往下叠的递归结构------底下究竟是怎么搭起来的，他年轻时没有回头看过。那时候他只管往上走，走得比谁都快，快到把根基甩在了后头。

iblp 是他自己削出来的。他见过 Laver 表，握不住，就照着它削，削掉十成里的九成九，剩下一个自己握得住的影子。影子已经大得离谱。

一百六十八年前的那一夜之后，他开始往回补。

七十三年后，他先把 ω-Y 序列的良序化证完了，落在 PTO(KP+ω₁)。那已经是旁人一辈子望不到边的地方，可对他来说，只是垫脚的那一层。

真正难的在上头。iblp 比 ω-Y 高出不知多少，他又推了九十五年。

补了一百六十八年，补到今天早上这最后一行。

他放下刻刀，很长时间没有动。

不是累。

是六个月前坐下来的时候，他只想把 iblp 的良序化收个尾。可写到最后那几行，有一样搁了一百六十八年的东西，在他脑子里忽然动了一下。

一样他一直以为自己这辈子都看不懂的东西。

他盯着那卷玉简。

``原来是这样。''他对着空殿说了一句。

他站起来，走到圣殿边缘。脚下是扽魂殿层叠的殿宇，再远处是天斗城的屋脊，晨光铺了满城。

那群孩子今天出关。

七个月前，真霜呈上来的战报里，有一句他一直记着：胡家的人当街索要瑞兽。此前殿里也收过零星消息，说那支队伍的队长身上有魂兽气息。他让人沿胡家的线索往下查，却没有凭一个叫法给那孩子定下来历。瑞兽、契约、借来的气运，都有可能；今天重新浮起的那段旧纹路，或许能让他分清。

\begin{center}\rule{0.5\linewidth}{0.5pt}\end{center}

当天下午，八个人回了一趟扽魂殿。

良序圣殿外。

卢本韦背着一个不大的行囊站在树下。脚步声从身后过来，她转过身。

``大长老。''

test\_alpha0 走到她跟前，上下打量了一遍。``比我算的要快。你们这批人。''

卢本韦没接话。

``到了蓬莱岛，你们会见到大供奉，最菜萌新。''test\_alpha0 说，``见到他的时候，报你曾祖的名号。''

卢本韦的眉毛动了一下。``卢承业。''

``对。''test\_alpha0 点头，``一百三十六年前，他找过我。在这座圣殿上，跪在石台前面，跟我提了一个构想------魂灵。''

风穿过白玉柱，发出一点低响。

``那时候他很年轻。比你现在大不了几岁，衣服磨得起毛，鞋是破的，眼睛亮得吓人。''test\_alpha0 说，``他说人跟魂兽之间不该只有猎杀这一种关系。他想让魂兽的灵魂跟魂师共存，不是吞，是合。''

``我当时觉得他疯了。''

``因为构想本身？''

``不是。''test\_alpha0 摇头，``因为我比他先走过那条路，而且走死了。''

他背过身去，看着远处。

``在他来找我之前三十二年，我也想过一个人进星斗大森林。也是去谈共存的。''

卢本韦猛地抬起头。

``您也\ldots\ldots{}''

``我也。''test\_alpha0 说，``魂灵那两个字，我比你曾祖早得多就知道了。''

他停了一下。

``再往前二十三年，我去过你们家一场葬礼。''

卢本韦的动作停住了。

``一百九十一年前。那年我一百二十六岁，坐上大长老这个位子，已经三十七年了。''

``死的是你曾祖的祖父，卢守衡。''

``我知道那件事。''卢本韦说。

test\_alpha0 看着她。``你们家怎么说的？''

``族谱上写了四个字。卒于魂灵。''

``就四个字？''

``就四个字。''卢本韦说，``我问过我爹，我爹说他也只知道这四个字。''

风穿过柱子。

``那我告诉你后面的。''test\_alpha0 说，``我在场。''

``他找的是一头柔骨兔。''

卢本韦的手指停住了。

``八万年。''test\_alpha0 说。

她猛地抬起头。

``八万年。''test\_alpha0 又说了一遍，``你听清了。不是八千，是八万。''

``那头兔子是自愿的。这点我要先讲清楚------它在你们场子后山住了四十多年，是它自己点的头。柔骨兔那一族性子软，八万年也还是软的，它从头到尾没有伤过人。''

``问题不在它愿不愿意。''

``问题在八万年。''

test\_alpha0 的手在石栏上停了一下。

``一个人类魂师，撑死了能扛多少？他那时候是八十七级。八万年的魂力灌进一具人的身子，那不叫融合，那叫往一个碗里倒一条河。''

``Opt 留下的那套东西，只教到'怎么合'。没有'怎么解'，也没说过合之前该先称一称两边的分量。''

``头四天是好的。''

``那四天他能听见兔子在想什么，兔子也能听见他。场子里的人说，那几天他一直在笑。''

``第五天，兔子的魂力失控了。''

``不是它要伤他。是那股东西太大，它自己也按不住。八万年的魂力在一具人的身体里翻过来，先烧的是经脉，再往里烧。''

``反噬。''

``烧了一夜。''test\_alpha0 说，``天亮前他就没了。那头兔子在院子里守到最后，等他断气，它自己也散了。''

殿上很长时间没有人说话。

``葬礼上我站在最后一排。''

``抱着那本笔记站在最前头的，是卢守衡的儿子卢景川------你曾祖卢承业的父亲。他那年十三岁。不哭，就是站着。''

``景川成年后接掌了养猪场，是卢承业之前的那任厂长。场子和这本笔记，后来都传到了承业手里。''

test\_alpha0 转过身来。

``你背上那本，就是他当年抱着的那本。''

卢本韦的手不自觉地摸了一下行囊带子。

过了一会儿她才开口。``小舞------唐三的妹妹------她的扽魂就是柔骨兔。''

``我知道。''test\_alpha0 说，``今天早上我看见她的时候，也想起了那一夜。''

``那之后我想了二十三年。''test\_alpha0 说，``我想，融合这条路是死过人的。那是不是路子本身不对？两个种族之间的事，是不是根本不该塞进一个人的身体里解决？''

``该在星斗大森林里解决。''

``所以一百六十八年前那年开春，我打算一个人进星斗大森林。''

``我那时候已经能展开 ω-Y 序列了，会用，可说不清它为什么成立。iblp 也一样，那一层在 PTO(Z₃)，我用得了，也有几分把握，离出神入化还远。''

``路我画好了。从天州往西进星斗大森林，一直走到核心区，去找熊君。那些超过十万年的魂兽，智慧不比人差，听得懂我说什么。我想着，坐下来，总能说上几句。''

他停了一下。

``我没走成。''

``星斗大森林先出事了。''

卢本韦看着他。

``熊君有个儿子。''test\_alpha0 说，``二十万年。星斗大森林里排得上号的那一批，熊君之下就是它。''

``那年开春，它死在了星斗大森林的核心区。''

``二十万年的魂兽\ldots\ldots 怎么会\ldots\ldots{}''

``不是围杀。''

test\_alpha0 的语气没有起伏。

``围杀要人多，要拿命去填。这一趟一条命都没填。''

``进星斗大森林的一共十三个人。十二个封号斗罗，在外围站着。''

``十二个封号斗罗围杀二十万年\ldots\ldots{}''

``他们没动手。''

卢本韦愣住了。

``那十二个从头到尾没动过手。''test\_alpha0 说，``他们是去看的。''

``看什么？''

``看中间那个人怎么杀。''

殿上静下来。

``这些是我后来查出来的。那块地我去看过。星斗大森林的核心区是禁猎的，他们是私自进去的，天没亮进去，天黑以前就出来了。''

``地上的印子我数过。外围十二个，站得规规矩矩，像在观礼。''

``中间只有一个落点。''

``一个？''

``一个。''test\_alpha0 说，``从头到尾，那个人站在那儿没挪过窝。''

卢本韦的手指凉了。

``它反抗了吗？''

``抗了。''test\_alpha0 说，``那块地四周有四棵三人合抱的树，全断了。断口在离地二十丈。''

``那是熊在挣。''

``它够不着他。''

``那个人一步都没有退，也一步都没有进。''

``草是焦的，可树没倒，土没翻。那不是打出来的。''

``是那头熊站在他面前，一点一点被抽干的。''

卢本韦很久没有出声。

``能做到这件事的人，这世上一只手数得过来。''test\_alpha0 说。

他停了一下。

``我是其中一个。''

``另一个叫胡赞。''

``帅气教的教主。''

``那天站在中间的就是他。''

殿上很静。

``帅气教跟扽魂殿，那时候已经对着干了？''卢本韦问。

``对着干了几十年。''test\_alpha0 说，``可你别把那时候的帅气教想成今天这样。''

``那时候的帅气教，还远没有今天这么大。堡是有了，墙根底下也已经站着骷髅，可他们在外头还没做起买卖，没有'集团'那张皮。在外头干得最多的事，就是进星斗大森林，杀魂兽，把杀来的东西往自己身上堆。''

``他们信这个。年份越高的越管用，杀得越狠，长得越快。''

``扽魂殿禁猎，他们就私自进；扽魂殿抓人，他们就换个地方接着杀。一年到头就是这点事，来来回回，我们抓过，罚过，也打过。''

``他为什么亲自去？''

``为他儿子。''

test\_alpha0 转过身来。

``二十万年的帅气，抽出来没浪费，当场就灌进了他儿子身上。那十二个封号斗罗是去观礼的，也是去做见证的，见证教主的儿子在那一天拿到了什么。''

``他儿子叫胡继，那年二十九岁，八十几级。''

``三个月以后，他升了封号斗罗。''

``那不是修出来的。''test\_alpha0 说，``那是他爹从一头二十万年的魂兽身上，给他刮下来的。''

卢本韦的手攥紧了。

``帅气抽干了，他们还没走。''

``他们没取魂骨。取魂骨要剖，剖了，这东西就废了。''test\_alpha0 说，``他们当场把它的魂剥了下来。帅气抽干，魂剥干净，剩下那具壳子，就是傀儡。''

``刚炼成的傀儡，腿脚还不灵。他们拿铁链锁着，拖出了星斗大森林。''

``熊君赶到的时候，那片空地上只剩一道血印子。十几丈长，一条直线，一直拖到星斗大森林外头。''

``东边那片洼地里，还躺着它那一族两百多头。大的小的都有，有几头还没长开。帅气一头一头抽过，魂一头一头剥过，壳子扔在那儿，堆成了山。''

卢本韦的脸色白了。

``拖出星斗大森林的那一具，一路往西运，运了四十天，运到了大陆最西边。''test\_alpha0 说，``帅气堡。''

``那是帅气教的总坛。再往西就没有路了，只有海。''

``他们在外面做买卖的时候叫帅气集团，在堡里的时候叫帅气教。一个东西，两张皮。''

``那具傀儡到今天还在。一百六十八年，还在他们手里。他们拿它打过三场仗，三场都赢了。每打一场，就把它摆出来给人看一次。''

殿上很久没有声音。

卢本韦站在那儿，半天没动。

``天州那一仗，''她说，``殿里的史册上只有三行。兽潮犯境，总殿两位长老殉职于天州，大长老亲至，退敌。''

``三行够写了。''test\_alpha0 说，``写史册的人，只看见了最后那一天。''

他转过身去，望着西边。殿外的天很高，西边天尽头什么都没有。

``前头那几天，你听我从头讲。''他说，``那几天的事，有一半是我后来才知道的。星斗大森林里怎么起的兵，总殿那两位长老怎么去的天州，我一样都没看见。看见的人，后来一个一个讲给我听。''

他接着往下讲。讲得很慢，一句一句，一百六十八年前的那个春天，就这么一点一点回到了眼前。

\begin{center}\rule{0.5\linewidth}{0.5pt}\end{center}

消息是第二天传进星斗大森林最深的地方的。

送信的是一条碧鳞蛇，三万年，平日在东边洼地一带巡林。它游进核心区的时候，背上的鳞掉了一半，身上沾的血，不知道是谁的。

熊君正趴在半山腰的洞口晒太阳。碧鳞蛇在它跟前盘下来，话才说了一半，它已经站起来了。

它是自己走去洼地的，没让谁跟着。

先到的是那片空地。四棵三人合抱的树断在离地二十丈的地方，断口还是新的，渗着树浆。草焦了一圈，焦得很齐，圈子当中干干净净，什么都没有。只有一道血印子，十几丈长，一条直线，从空地中间一直拖出去，拖向星斗大森林外头。

熊君低下头，顺着那道印子闻过去，闻到星斗大森林的边上，停住了。再往外，是人的地方。

它掉头去了东边的洼地。

洼地里堆着它那一族两百多头。大的压着小的，小的蜷在大的肚子底下。有几头还没长开，毛是软的。一头一头都干瘪了，像晒过的皮囊，魂没了，帅气也没了，只剩一副壳子。

熊君站在洼地边上，看了很久。

它那个儿子出生的时候，星斗大森林下了一整夜的雨。小东西生下来没有巴掌大，眼睛睁不开，拱在母熊的肚皮底下找奶吃。它守在旁边，守了一夜。后来长到十万年，长到二十万年，个头追上了它的一半，还是改不了那个毛病，困了就往它身上靠。

这些年它一直在等。等那小东西再长一长，长到能替它守着这一族。

它没有吼。它走下去，用鼻子把最上头那一头小的拱开，底下还压着一头更小的。它又拱，拱到最后，前爪踩进了一摊黑水里，就那么站着，一直站到天黑。

\begin{center}\rule{0.5\linewidth}{0.5pt}\end{center}

星斗大森林最深的地方有一片黑水潭。潭边一块青石，帝天坐在上头。

他是人的样子。个子很高，一身黑，膝上横着一把黑剑。

那天夜里，潭边来了十大凶兽里的七个。有的蹲在潭边，有的盘在树上，只露出一双眼睛。没有谁先开口。

熊君是最后到的。它从林子里走出来，在潭边趴下，把下巴搁在前爪上。

``我儿子叫人抽干了。''它说，``洼地里还有两百多口。''

赤王身上的火苗往上蹿了一蹿。``谁干的？''

``一个人。''熊君说，``站着，没挪窝。外头还有十二个，站着看。''

潭边静了。能一个人站着把二十万年的熊君之子抽干，这样的人，在场的谁都没见过。

``一千年前那纸约，''万妖王慢慢开了口，``人不进大森林深处，魂兽不出大森林。约是谁签的？''

``扽魂殿。''帝天说。

他一直没动，这时候才抬起眼皮。

``人类魂师的头，就是扽魂殿。约是它签的，它管不住自己的人，这笔账，就记在它头上。''

熊君把头抬了起来。

``我要去。''

潭边一棵树上，碧姬收拢了翅膀。她是一只翡翠天鹅，平日在潭边替受了伤的魂兽治伤，很少开口。

``扽魂殿有个大长老。''碧姬说，``不到九十岁就坐上了那个位子。听说人那边，没有谁敢跟他动手。''

赤王身上的火苗又蹿了一蹿，没吭声。

``人。''帝天说，``到了九十九级，也就那样。''

他说这话的时候，没有一点要压人的意思，像在说一件谁都知道的事。五百年前，扽魂殿有一个九十九级的极限斗罗独自进了星斗大森林，走到这片黑水潭边上，再没出去。那一回，帝天只出了一爪。

潭边再没有谁提那个大长老。

帝天没有马上答熊君。

潭水黑得看不见底。他坐在青石上，一只手按着膝上的剑。他胸口底下压着一团金光，压了一万年，平日里安安静静的，那几天一直不怎么安稳。

过了很久，他站了起来。

``那就去。''他说，``到了城下，先跟扽魂殿要人。不交，就踏平。''

树上的碧姬把翅膀张开了一半，又收了回去。

``我跟着去。''她说，``打起来，总得有个治伤的。''

那一夜，帝天做了一件一万年都没做过的事。他把整片星斗大森林的精锐，一头不留，全放了出来。

\begin{center}\rule{0.5\linewidth}{0.5pt}\end{center}

兽潮是第三天天亮出的林。

最先觉出来的是鸟。星斗大森林东边那一大片林子里的鸟，天还没亮就全飞了，黑压压一片往东边去，飞过天州城头的时候，把半边天都遮暗了。城里的人站在街上抬头看，谁也说不出是怎么回事。

然后是地。天州城里，桌上的碗自己往桌边上挪。

离星斗大森林最近的那座小城，城墙两丈高，城里住着三千来口人，大半是猎户和采药的。守城的魂师一共十一个，最高的四十几级。

兽潮到的时候是晌午。

城门楼上的魂师头领姓何，四十三级，在这座城里守了二十年，打过狼，打过野猪，最大的一回打退过一头三千年的铁背熊。他站在城楼上往西看，看见麦田那头起了一道黑线。黑线越来越粗，越来越近，到后来，麦田没了，只剩一片往这边涌的脊背。

他回过头，冲城里喊了一句：往东跑，别回头。

喊完，他把城门楼上那面旗拔下来，绑在自己背上，拎着刀下了城墙。

城墙没撑到太阳偏西。往东去的那条土路上，挤满了人，推车的，牵牛的，背着包袱的。有人跑着跑着回头看了一眼，看见城门楼上那面旗不在了。

第二座小城大一些，城墙也厚一些，撑到了半夜。城当中有座钟楼，从天黑起一直在敲，敲得又急又乱，隔着几十里地都听得见。后半夜，钟声停了。钟声停下没多久，城里就起了火，火光隔着几十里地也看得见。

从星斗大森林里出来的，一共二十万。兽潮最后头，碧姬张着翅膀跟着，一层翡翠色的光从她身上铺出去，落在受了伤的魂兽身上。十大凶兽来了五个：帝天，熊君，万妖王，赤王，还有跟在后方救治的碧姬。

万妖王和赤王各带走一支，沿西北、西南两侧的丘陵向天州逼近。两支合起来近十万，隔着山梁，与正面的兽群遥相呼应。它们还没有绕到城东，天州通往天斗城的官道，暂时仍在守军手里。

剩下的约十万头，由熊君带着，直往天州西门压过去。帝天走在后面，碧姬跟在他附近。麦田一片接一片倒下，远远望去，像有一层黑水漫过了田埂。

\begin{center}\rule{0.5\linewidth}{0.5pt}\end{center}

\chapter{第八十九章：兽潮}\label{ux7b2cux516bux5341ux4e5dux7ae0ux517dux6f6e}

逃出来的人是夜里到天州的。

西门开了一条缝，放进来几百个，男女老少，有的光着脚，有的身上带着伤，有个汉子背上背着他娘，进了城门，把人放下，才发现他娘已经没气了。

天州扽魂分殿的殿主，连夜在西城门楼上见了这些人。

殿主五十来岁，八十一级，在天州待了二十年，一夜里涌进来几百个逃难的，他头一回见。

``多少？''殿主问。

逃来的人里有个年轻猎户，嘴唇抖着，答不上来。问了三遍，他才说出一句：``数不清。''

过了一会儿，他又说了一句。

``它们不吃人。''他说，``也不抢东西。踩过去就走，一路往东。''

``领头的是什么？''

``一头熊。''猎户说，``站起来，比城楼还高。''

殿主皱了一下眉，没往下问。

他走到垛口边上，往西看。西边天上有一片红，是第二座小城还在烧。红光底下，看不清有多少东西在动。

州长也上了城楼，站在他身后，手里攥着一卷城防图。``殿主，要不要先跟天斗那边要人？''

``要。''殿主说，``折子你今夜就递，医疗兵、城防营，能要的都要。我这边往总殿递信。''

他当夜写了一封信。信不长：西境魂兽出林，两城遇袭，兽群数不清，为首者似一巨熊，分殿力不能支，请总殿速援。

魂导鸽是半夜放出去的，天亮前到了天斗城。

\begin{center}\rule{0.5\linewidth}{0.5pt}\end{center}

总殿值夜的执事把信收下，看了一眼，没敢压，连夜送进了长老院。

那一夜在院里的，是周长老和叶长老。

两位长老一个姓周，一个姓叶，都是九十六级的封号斗罗，在总殿当了几十年长老。西边几州出了事，十回里有七八回，是他们两个去的。

周长老把信看完，递给叶长老。

``一头熊。''叶长老说，``站起来比城楼还高。''

``十万年上下。''周长老说。

``又是哪伙猎魂的进了星斗大森林，惹急了一头十万年的。''叶长老把信往桌上一拍，``两座城都破了。''

星斗大森林边上，人跟魂兽起冲突，这些年不是一回两回了。领头的多半是一头年份高的，顶天了十万年。总殿往西边派人，回回都是一个打法：找着领头的，宰了，剩下的自己就散了。

``等不到天亮了。''叶长老说，``拖一个时辰，就多死一片人。''

周长老没有马上应。他比叶长老大十来岁，头发白了一半，说话慢，做事也慢。他把信拿回来，又看了一遍。

``十万年的，''他最后说，``咱们两个，够了。''

两个人给二长老千古东风留了一张条子：西境魂兽出林，两城遇袭，似为十万年魂兽作乱。周、叶二长老即刻驰援天州，事后详报。

条子是天亮前搁到千古东风案上的。早上千古东风看见了，提笔在后头批了一个字：准。

往年这样的事，总殿都是这么批的。

条子没往良序圣殿送。大长老那几天在推演，殿里的规矩是，没有要紧的事，不去扰他。

\begin{center}\rule{0.5\linewidth}{0.5pt}\end{center}

天刚亮，城外先传过来一声``交人''。州长正让分殿派人问清，两个身影已经落在了西城门楼上。

州长迎上来，才说了半句``城外的魂兽要找人''，叶长老已经走到了垛口边。

从城墙根底下往西，一直到天边，麦田里全是魂兽。最前头伏着一头巨熊，暗金色的毛，肩背高过了城墙。它身旁的兽群还没有往前扑，正朝城头张望。

``就是它。''叶长老说，``杀了领头的，其余的自然就散。''

``二位长老，它们这次\ldots\ldots''

``两座城都没了，还等它们把这一座也拆了？''

叶长老没让州长说完。周长老往城外扫了一眼，把腰间的令牌交给一名执事，拔出了剑。

``开门。我们出去，你们守住城墙。''

西城门开了一条缝。两位长老一前一后走了出去，身后没带一个人。

州长站在城楼上看着。他五十多岁，在天州当了九年州长，头一回看见城外的麦田变了颜色。他还想让人把分殿昨夜记下的问话送过去，两个人已经掠过了护城河。

叶长老先动手。

最外头那群风狼尚未扑起，一片青光已经压进兽群。一掌下去，五六头风狼飞出麦垄，落地就没再起来。周长老紧跟着落下，一剑横扫，几头铁甲犀连着断角栽进了田里。

没有问话，也没有停手。

两人径直往那头熊面前杀。叶长老一掌推开一头护着幼兽的母狼，下一掌已经蓄在身前；周长老的剑光从他身侧过去，把正往后退的兽群截断了。

兽群里响起一声短促的哀叫。

那头巨熊抬起了头。

它先看见地上那几具新倒下的尸首，又看见叶长老抬起的手。前些日子，洼地里那摊黑水的气味，仿佛又冲进了鼻腔。

``还要杀多少！''

这一句是人话。吼声落下，熊君已经站直了身子。

五十万年的气息骤然压开，叶长老掌中的青光猛地一滞。直到这时，他才看清那只探到头顶的熊掌------每一根暗金色的利爪，都比他整个人还长。

``退！''周长老喊。

叶长老转身的同时，护体魂力已经提到了顶。第一道爪芒落下来，那层青光从中间裂开，整个人被拍进了田埂，连抬起第二只手的工夫都没有。

周长老的剑随即斩在熊君腕上。

剑刃切开了外层的毛，却没能再往里走。熊君翻过前爪，五道暗金色的光从剑身上掠过去。剑断了，周长老胸前的护体魂力也跟着散开。

州长在城楼上只来得及看见两片青光先后熄灭。前一刻还在往兽群深处冲的两个人，此刻一个倒在断裂的田埂里，一个摔在护城河外，再没站起来。

执事手里的令牌掉到了砖地上，滚了两滚。

州长张着嘴，喉咙里一点声音也没有。过了好一会儿，他才回过身，冲城楼底下吼了一句：关门！上闩！三根全上！

城门轰的一声合上了。三根铁木门闩，一根一根落进槽里。

熊君踩着翻开的麦垄往前走。它身后，刚才退开的魂兽重新聚拢，低伏着身子，跟到了护城河边。

二十里外，帝天仍在林边。他看着城头升起的扽魂殿旗号，手里的黑剑没有抬起来。

熊君仰头，冲着天州城吼了一声。

吼的是熊族的话，城里没有一个人听得懂。可城墙上的人，手里的刀都握紧了。

\begin{center}\rule{0.5\linewidth}{0.5pt}\end{center}

那一天，天州打了一整天。

西城墙最先挨的。风狼一群一群往上扑，铁甲犀拿角顶城门。城头上的魂导炮响一阵，停一阵。弓手射完了箭，开始往下扔砖。

城楼上有一门魂导炮。它炸翻的魂兽，倒在麦田里，过不了多大工夫，身上亮起一层绿光，又爬了起来。炮手是个十七八岁的后生，炮管打红了，他拿湿麻布往上捂，烫得直甩手，甩完了接着打。打到晌午，他耳朵聋了，旁边的人冲他喊话，他只看得见嘴在动。

晌午过后，西城门的门闩断了一根。

后半晌，又断了一根。三人合抱的铁木，最后只剩一根还横在门上，门板中间裂开了一道缝，缝里能看见外头魂兽的眼睛。

城门里头，铁匠把自家的铁砧一车一车推过来，堆在门后头。门每挨一下，那堆铁就往后滑一尺，几十个人扑上去，拿肩膀顶回来。

西边那一段城墙上，已经爬上来几十头，跟守城的人撕在一块儿。西城起了火，街上全是往东跑的老百姓。

州长前一天夜里就往天斗城递了求援的折子。天斗帝国的医疗兵正往这边赶，最快也得天黑才到。

城防营的一个校尉趴在城墙上，吊着左胳膊，拿右手捏着炭条在一张纸上写字。写完一行，他把纸叠起来，绑在一只魂导传信鸽的腿上，往东一扬。那是城防营里最好的一只，最新的样式，半个时辰能飞一千里。

这一回，他没写给总殿的执事。他写的是：呈大长老。

头一封信是分殿夜里发的，总殿来了两位长老，两位长老死在了城外。他只知道，能把这座城救下来的，整个扽魂殿只有一个人。

鸽子飞过东城的时候，底下的官道上，天斗帝国的医营马车正一车挨着一车往天州赶。领头的军医抬头看了一眼那只鸽子，冲车夫喊了一句：再快些。

\begin{center}\rule{0.5\linewidth}{0.5pt}\end{center}

良序圣殿里很静。

test\_alpha0 坐在石台前头，面前摊着一卷推演，ω-Y 序列推到一半，他停在一处，看了半个时辰没往下写。

石台的角上压着一张地图，是他前些日子自己画的。一条墨线从天斗城出来，往西，过天州，进星斗大森林，弯弯绕绕，一直画到核心区。线的尽头画了一个小圈，圈旁边写着两个字：熊君。

他原打算过了这个月就动身。

一只魂导传信鸽从殿门外头扑进来，翅膀擦过柱子，掉下来几根毛，落在他的石台上。

他把纸解下来，展开。

纸上就一行字，字是歪的，几处炭迹被手指蹭开了。

兽潮至。两长老殁。城门将破。

纸上没写为什么。

他看完那一行字，手里那张纸就没了，连灰都没剩。

两位长老，他都认得。周长老往西边跑过不下十回，从来没有失过手。叶长老年轻些，下手更重。九十六级的封号斗罗，两个一块儿，到天州半天就没了。

能在半天里杀掉两个九十六级的，星斗大森林里数得出来，一只手就够了。那样的东西出了星斗大森林，这件事就不会是哪伙猎魂的人惹出来的小事。

他低头看了一眼石台角上那张地图。墨线从天州穿过去，天州在线的半道上。他要去找的那一头，先找到了天州。

他站起来的时候，千古东风正从殿外快步走进来，手里攥着一本册子，脸是白的。他看见石台上那只鸽子，脚底下顿了一顿。

``头一封信，昨天夜里到的。''千古东风说，``周、叶二位说是十万年魂兽作乱，要去天州驰援，我批了一个准。我以为\ldots\ldots{}''

``我知道。''test\_alpha0 说。

他没有多问。

``我跟你一块儿去。''千古东风说。

``殿里得有人坐镇。''test\_alpha0 说，``你留下。''

千古东风张了张嘴，没再说什么，把那本册子攥得更紧了。

他转过头。

``Smile。''

殿柱后头走出来一个人，九十八级巅峰斗罗，九级魂导师。右手戴着一枚亲手炼制的银灰色储物魂导戒，备用核心、修复器件和定装魂导炮弹，分存其中。那年他六十六岁，还不是二长老。

``跟我走。''

Smile Lee 应了一声，拇指在魂导戒上轻轻按了一下。里面那颗九级定装魂导炮弹刚刚完工，封禁阵纹才合拢不久。他还没来得及试射，天州告急的消息就到了。

七年的工夫，总算能派上用场。他要让师父亲眼看看，这颗炮弹能为天州打出多大的缺口。

test\_alpha0 走出殿门，站在台阶上，往西看了一眼。天斗城的天是晴的，西边天尽头压着一线黑，细得像一根头发。

那根头发底下，是天州。

他从台阶上一步跨了出去。

从良序圣殿下去，落到天州城里，用了不到一炷香。

\begin{center}\rule{0.5\linewidth}{0.5pt}\end{center}

他往下落的时候，头顶上的云跟着动了。

天州正上方那一块乌云，越转越快，转成了一个漩涡，甩开的地方透出蓝来。等他落到西城门楼上，天州头顶已经是一整个蓝色的漩涡，照得城楼上一张张脸都是蓝的。

漩涡外头，天边还是黑的，乌云压在城外那一大片魂兽头上，黑得像锅底。

那片蓝，是他自己的魂力，天上的云跟着它走，绞成了那个样子。

他一落下，先看见一头风狼正扑向断了胳膊的守兵。

袖子向外一展，一股蓝色的魂力托住了狼的腰，把它从人身前移开。紧接着，贴在西城墙上的几十头，正往上爬的、挤在墙根底下的，都被这股力量推过护城河，落回麦田。它们挣扎着爬起来，想再往前，河边却像多了一道看不见的坎。

他没有追击，只把魂力留在城墙与兽群之间。垛口上空了，守城的人还举着刀，刀底下已经什么都没有了。

街上跑的人停住了脚。城墙上扔砖的弓手，手里还举着砖，都抬起了头。

写信的那个校尉靠在垛口底下，吊着左胳膊，仰着脸看头顶那团蓝。从鸽子出手，到这个人落在城楼上，还不到一个时辰。

他往东边看了一眼，又看回城外。那些刚才还要越过垛口的魂兽，全被隔到了河外；守兵从兽爪底下爬起来，扶着墙站住，还不知道该往哪里挥刀。

校尉的右手指腹上仍沾着炭灰。写那封信的时候，他左臂痛得抬不起来，只能把纸压在砖缝里，一笔一笔写下``城门将破''。他不敢写还能撑多久，怕鸽子尚未飞到总殿，这句话就已经过了时。

现在，他把手掌按在仍然立着的城墙上，指头不自觉地收紧了。

州长就在城楼上，离他不到两丈。一身官袍烧掉了半边，手里攥着半截断刀。他看看头顶那个漩涡，又看看落下来的这个人，愣了一会儿，眼泪先下来了。

他转过身，冲着城墙上那些还站着的人，把那半截断刀举了起来。

``救星来了！天州必胜！''

城墙上的人跟着喊，嗓子都是哑的，一段一段传下去，一直传到了城东。

喊完，州长把断刀往地上一放，冲着 test\_alpha0 单膝跪了下去。

test\_alpha0 把他扶了起来。那句''救星''，他没接。

Smile Lee 随后落在城楼上，目光扫过残损的炮位，又落到城墙上早已熄灭的防御阵纹。

``城里交给你。''test\_alpha0 说，``别让人踩死人。城墙上的东西，能修的都修起来。''

``是。''

Smile Lee 下了城墙。

城里比城墙上还乱。西城的火顺着风往东烧，烧到哪条街，哪条街的人就往外涌。东门那边挤成了一团，几千人往一扇门里涌，有人摔倒了就再没站起来。Smile Lee 站到东门口一辆翻倒的粮车上，把扩音的魂导器扣在嘴边，一条街一条街地点人：哪几个魂师守哪条街口，老的小的先走，青壮的留下来搬砖。乱成一锅粥的人群，一点一点有了个样子。

点完了人，他回到城墙上。西墙的控制阵盘裂成了三瓣，几处汇流口已经烧黑，原本通向城东主阵室的阵纹，只剩近处几段还在忽明忽暗。

Smile Lee 从魂导戒里取出备用控制核心，先接过了这几段的控制。他把断开的阵盘固定在一起，沿裂口续上导流纹，再将一枚银灰色的核心压进中央凹槽。九十八级的魂力顺着他的指尖灌进去，城墙内侧埋着的阵线，一条接一条重新亮了起来。

西城墙上，一层淡金色的应急光幕先立了起来。

``这只能先护住墙头。''他把一张标好位置的城防图递给校尉，``带人去接南北两侧的阵眼。城东的汇流节点烧穿了，等我过去处理。四面的魂力接通，护城大阵才能真正开启。''

校尉喊来几名还能走动的魂师，接过他分下的修复器件，照着图奔下了城楼。

他又去看城楼上那门魂导炮。炮管已经红得发暗，那个十七八岁的炮手还蹲在边上拿湿麻布捂。Smile Lee 让他让开，从魂导戒中取出一只备用散热环，换下炮尾烧损的那只，拍了拍炮管。

``打吧。''他说，``这回不烫手了。''

炮手愣了半天，才听清这一句。

修完了炮，他没歇。魂导戒上的空间阵纹微微一亮，一排排拳头大小的魂导炮弹便整齐落在城楼内侧。

他从城防营要了十几个魂师，一人领了一批，从城墙上缒下去，顺着护城河外头那道土埂往麦田里埋。

一颗跟一颗隔着三丈，排成三道，从西城门口一直排到城墙拐角。每埋好一颗，他就拿指头在地上划一道，把这一颗跟上一颗连起来。

``踩着就炸？''一个正在埋弹的年轻魂师问。

``一颗响了，这一道上的全跟着响。''Smile Lee 说。

埋到第三道的时候，兽群里有几头风狼盯上了他们，压低了身子往这边窜。Smile Lee 没抬头，手掌往地上一按。

最外头那一道土埂底下，一串闷响从西往东滚过去，泥土跟火光一齐掀上了天。那几头风狼让掀回了麦田里，后头的魂兽在麦田边上站住了，一时没往前挪。

城墙上有人吹了一声口哨。

Smile Lee 拍去手上的土，身形一提，已回到城墙上。

\begin{center}\rule{0.5\linewidth}{0.5pt}\end{center}

而在 Smile Lee 入城疏散百姓、抢修防御的时候，西城门楼上，test\_alpha0 已经在与兽潮对话。

二十里外，星斗大森林的边上，帝天原先还坐在那块石头上。

天州头顶那团蓝色的漩涡转起来的时候，他抬起了头。

这股魂力，跟五百年前死在黑水潭边上的那个九十九级不大一样。它是一层一层往上叠的，叠得很高。他只知道人类那边又来了一个，这一个比早上那两个硬。

他胸口底下那团金光，轻轻动了一下。

帝天站了起来，把膝上的黑剑提在手里。

他刚站直，龙吟便从林边滚向天州。隔着二十里，震得城楼上的瓦片直跳。

龙吟后头，跟着一个人的声音。

``卑鄙的人类！屠杀族人之恨，今日必须报仇！交出那个人。否则，踏平天州！''

城外那一大片魂兽，听见这一声，全往前压了一步。

是帝天。城楼上落下来的这个人，他不认得，也懒得去认。来的是谁，在他那儿都一样。

test\_alpha0 没有立刻放出领域。

``你要的是哪个人？''他问，``在什么地方，什么时候动的手？''

``杀我族人的凶手！''帝天的声音隔着二十里传过来，``人类魂师闯进森林，杀了熊君那个二十万年的儿子，又屠了洼地里两百多头同族。你们扽魂殿签的约，你们就得把人交出来！''

test\_alpha0 听完，回头问分殿殿主。殿主脸色发白，摇了摇头。近来猎魂队进出森林的名册都在分殿，可没有谁报过深处的洼地，更没有一份记着二十万年魂兽的卷宗。

帅气教是偷着进去的。城楼上的人，谁也没有接到过那件事的消息。

``我刚到，不能凭一句话认下凶手是谁。''test\_alpha0 转回身，``先停攻。天州停炮，你们退过护城河。我同你们去看出事的地方，查清是谁动的手。若是人类魂师干的，扽魂殿不包庇。''

``退？''帝天说，``人没交到本座面前，一步都不退。''

``那就先让百姓从东门走。查凶手用不着拿他们抵命。''

``我的族人死的时候，可有人让它们先走？''

test\_alpha0 看了一眼城下。兽群被隔在河外，仍在往前挤。他扶着垛口，放慢了声音。

``凶手是谁，眼下还没有查清。我可以留在天州，替总殿担保。你们给我查证的时间，也让城里的人有条活路。''

最前头，那头巨熊低下头看他。隔着一道城墙，他闻得到它嘴里的血腥气。

``你们这种人，''熊君说，``我前些日子刚见过一个。''

\begin{center}\rule{0.5\linewidth}{0.5pt}\end{center}

殿上，test\_alpha0 停了下来。

``那句话，我当时没听懂。''他说，``我以为它说的是'人类'。''

``它说的是'你们这种人'。它指的是能一个人站在那儿不挪窝、把一头二十万年的魂兽抽干的那种人。那个层次，它前些日子刚在自己儿子身上见过一次。''

卢本韦的呼吸停住了。

``那天我站在城楼上，身上带着展开一半的 ω-Y 序列。在它眼里，我跟胡赞是同一样东西。''

``可那一趟您根本不在场。''

``我不在场。它儿子死的时候，我在这座殿上，扽魂殿连报都没收到。''test\_alpha0 说，``它不知道我叫什么，也不知道我是哪一边的。它只知道前些日子有一个人站在那儿没挪窝，把它儿子抽干了。这会儿城楼上又站着一个人，身上是一样的东西。''

``你要它怎么分？''

卢本韦说不出话。

\begin{center}\rule{0.5\linewidth}{0.5pt}\end{center}

熊君又说了一句。

``人类，你也配谈和平？''

它低头扫过脚下的麦田。周、叶两位长老杀死的魂兽，还倒在它身后；护城河外，那片暗下去的血迹，仍混在翻开的泥土里。

``早上来的那两个，也是替你们担保的？''

test\_alpha0 顺着它的目光看过去，停了一息。

``他们先动手的事，我会查。可城里那些人，没有杀你儿子。''

二十里外，帝天的声音又滚了过来。

``本座不等。交出凶手，或者同这座城一起死。''

``一定要继续打？''

帝天没有再答。熊君抬起前掌，往地上重重一顿。兽群里响起一片低吼，最前排的铁甲犀抵住了河边那道无形的阻隔，往前顶出了半步。

test\_alpha0 把手从垛口上收了回来。

``既然各位不肯罢休，那我只能由一人代表扽魂殿，来平息这场事件。''

声音不大，却清清楚楚地传遍了城外，连二十里外星斗大森林边上的树梢，都跟着轻轻晃了一下。

话音落下，他把良序领域放了出去。

一层看不见的东西从城楼上往外铺。铺过护城河，铺过埋着炮弹的那道土埂，一直铺到二十里外星斗大森林的边上。城外的风停了。麦穗不晃了，烧着的草棚上冒出来的那几缕烟，直直地停在半空，不往上走。

天州头顶，那团蓝色的漩涡转动了一下。

最外缘的一圈蓝，剥下来许多细碎的光。蓝白色，比雪沫子还细，直接从漩涡里飘落，越过城楼的檐角，一片一片落向城外。没有风托着，它们却没有直坠下去，而是沿着良序领域铺开的方向，缓缓散向更远的麦田。

守军仰头看着。那些光从他们头顶经过，到了城墙外才低下来，落到魂兽的鼻尖前，悬在离地数尺的地方，仍按着天上漩涡的方向，慢慢打转。

那是 iblp 的碎片。真正的核心还压在 test\_alpha0 体内，只有泄出的一点映进了头顶的蓝色漩涡，又从漩涡里一片片落下来。麦田上方渐渐铺满了细光，越靠近城墙，落得越密。

一头风狼不知道厉害，抬起鼻子，去嗅离它最近的那一点。

鼻尖碰上去，那头狼就没了。没有血，也没叫出声，原地连一根毛都没剩，地上只留下四个爪印。

边上的魂兽一头一头往后缩，伏低了身子，脑袋贴着地，眼珠子跟着那些亮点转。一头铁甲犀退得急了，后腿蹭着了一点，当场湮灭，从后腿起，一直没到角尖。

麦田里，再没有一头敢动。

漩涡边缘又剥下几点细光，沿着领域伸向远处的纹路飘出去，越过整片兽群，一直落到了二十里外。帝天提着黑龙剑站在林边，其中一点落在了剑刃上。

剑身猛地一抖，嗡的一声，抖得帝天整条胳膊都跟着一震。他低头去看。剑刃上那一点蓝亮了一亮，灭了。落过的地方，乌黑的剑身上褪掉了指甲盖大的一块，露出底下一片灰白。

帝天把剑握紧了，没出声。

过了好一会儿，兽潮深处才炸起来一声咆哮，震天动地，算是给城楼上那句话的回答。

兽潮最后头，亮起来一层光。

翡翠色的，从星斗大森林边上铺过来，一丝一丝，细得跟蛛丝一样，落到河外那些负伤的魂兽身上。先前挨了炮火的，刚从城墙上退下去的，断腿往回接，裂开的甲往一块儿合。麦田里横七竖八躺着的那一片，一头接一头，又爬起来了。

Smile Lee 就蹲在城楼那门魂导炮边上。他抬起头，顺着那些光丝往回看，一直看到二十里外。光丝的那一头，星斗大森林边上的半空里，浮着一只翡翠色的大鸟，两只翅膀张开，有十几丈宽。

翡翠天鹅。十大凶兽里的碧姬。

他在心里算了一笔账。守军拼命打伤一批，那边转眼又给治好；前头那些才被推回去的魂兽，也重新围到了河边。只要那只鸟还张着翅膀，这场攻城就能一轮接一轮地打下去。

得先剪了后头那根线。

他拇指在魂导戒上一抹，一只长条魂导匣落进掌中。匣面上三道封禁阵纹依次熄灭，匣盖无声滑开。里头躺着一颗炮弹，比小臂长，壳子是暗银色的，上头的纹路密得看不出头尾。

九级定装魂导炮弹。他做了七年，才刚刚做成，就这一颗。

来天州的时候，他就想过，要用它在这场兽潮里大展身手。眼下，二十里外那一丝丝翡翠色的光，给了他一个不能再等的理由。

他从戒中取出备用的强化炮膛，换下城楼这门炮的旧膛。掌心贴住炮身，魂力沿阵纹走了一遍，将各处接合牢牢锁住。炮架随他的魂力平稳转向西边，瞄具里的标线一点点收拢，停在碧姬张开的双翼之间。他又校过一次落点，才把那颗九级炮弹送进炮膛。

test\_alpha0 在垛口边上回过头，看了他一眼。

Smile Lee 也看着他。

test\_alpha0 没说话，把头转了回去。

Smile Lee 将魂力送入发射阵纹。

炮弹出膛的那一下，整座城楼往下一沉，瓦片哗啦啦往下掉。一道白线从城楼上拉出去，贴着兽潮的头顶，一直往西，二十里，眨两下眼就到了。

星斗大森林边上，帝天还提着剑站在那块石头跟前。

他没挪脚，把黑龙剑往上一撩。

一道黑光从剑刃上翻出去，迎着那道白线劈了上去。白线让它从当中劈开，炸成了两团，一团往北，一团往南，落在空地上，掀起两座土山。

炸开的余火还是往里卷了一截，扑在了碧姬身上。

那只翡翠色的大鸟在半空里晃了一下，往下坠了十几丈，又硬撑着停住了。她左边那只翅膀上，翎毛着了一大片，冒着烟。连在前头魂兽身上的那些光丝，一下子断了好些，麦田里那几百头，又有不少软软地趴了回去。

碧姬没掉下来。她把左翅往身上收了收，往后退了半里地，身上的光又一丝一丝地亮起来，比刚才稀了，也暗了。

城楼上，Smile Lee 拍了一下滚烫的炮尾。``让人挡了一大半。''他说，``可她再扶人，得慢了。''

他看着那片稀下去的翡翠色光丝，又看了一眼炮口升起的白烟，嘴角不自觉地动了一下。头一回实战，隔着帝天那一剑，仍能打出这样的结果，七年的工夫没有白费。

帝天回头看了一眼碧姬。她伤翼上的火还没灭，几片烧焦的翎毛从半空里落下来。

再转过头时，他的目光越过前头的兽潮，钉在了城楼上那名魂导师身上。那个人还把手搭在炮尾上，正端详自己的成果。

握剑的手一点点收紧，黑龙剑上那条龙的鳞片，亮起了一线暗金。炮火伤到了碧姬，城楼上那点得意又落进他眼里。他胸口压着的怒意，终于全翻了上来。

test\_alpha0 看着二十里外，没接这句话。

\begin{center}\rule{0.5\linewidth}{0.5pt}\end{center}

二十里外，叫人喘不过气的气息升了起来。

帝天的领域铺开了。星斗大森林边上的树一棵棵朝他那边倾过去，天边的乌云也翻滚着往他头顶聚。

兽潮最前头，熊君回过头，看了一眼，往边上让开了一步。

``人类！''

帝天的声音从二十里外滚过来，震得城头上的守军纷纷捂住了耳朵。

``敢在本座面前炮击碧姬。今天，本座就让你们付出代价！''

话音未落，他手持黑龙剑冲天而起，化作一道黑光，直扑天州城。二十里，转眼就到。剑锋所指，是城楼上那门还冒着烟的魂导炮，和炮边上的 Smile Lee。

黑龙剑上，漆黑的剑气翻滚着拖出几十丈长，剑气底下的麦田被犁出一道又深又长的沟。剑还没到，城头上那几门魂导炮的炮管，先被剑风削掉了半截。

在帝天看来，人类修到头，也不过是九十九级的极限斗罗。他的实力，早就远远超出了极限斗罗那一层。除了当年的 Opt，他从没把哪一个人类真正放在眼里。

Smile Lee 抬起头，那道黑光已经压到了城楼跟前，剑风刮得他睁不开眼。

test\_alpha0 看清了那道黑光对着的是谁。

他这辈子收的徒弟不多，Smile 是跟他最久的一个。那孩子十几岁进殿，在他眼皮底下一道阵纹一道阵纹地学。如今已是九十八级巅峰斗罗、九级魂导师，刚才打出去的那颗炮弹，也是他亲手造出来的。

一个连人带炮要往他徒弟身上劈的东西，正从二十里外冲过来。

他胸口里那团火，一下子顶上了天灵盖。

方才他只让一点 iblp 随着蓝色漩涡飘落，核心始终没有动。这一回，他将更深的一层放了出来，接进头顶的漩涡。

天州头顶那团蓝色的漩涡猛地一涨。一圈套一圈往外扩，扩出去好几倍，一直扩到城外麦田的上空，把压在兽潮头上的乌云硬生生推开了一大片。漩涡当中那一点蓝，亮得扎眼，城墙上的守军抬手去挡，指头缝里漏进来的全是蓝光。麦田里的魂兽一头一头伏低了身子，连熊君都往后挪了半步。

城头上，几个守军把遮在眼前的手放低了一点，想看清城外，却又被蓝光刺得眯起眼。一名老兵指着麦田，嘴唇动了好几下，才挤出一句话。

``那头熊\ldots\ldots 退了。''

旁边的人没接话。方才还把他们压得喘不过气的兽群，此刻一头头伏着，连吼声都低了下去。有人攥着刀站在垛口边，半晌忘了把手放下。

州长扶住垛口，仰头望着那片扩出城墙的蓝。他刚才还能喊出``天州必胜''，这会儿却张着嘴，喉咙里一点声音也没有。目光从退去的乌云移到 test\_alpha0 身上，又移回天上，扶着石砖的手越攥越紧。

一座城的守军和几门旧魂导炮，苦苦撑了一整天。Smile Lee 才把西墙的防御接起来，眼前这个人放出的魂力，竟已经压过了城外那片黑天。

Smile Lee 也怔住了。

他的感知追着蓝光向外铺去。最外面的几层，他还能辨出 ω-Y 序列的结构；再往里，熟悉的递归纹路便一层层隐没了。更深处有什么正在展开，他看得见，却认不出来。

他下意识地想再辨一层，目光却停住了。师父的气息仍在往上攀，转眼就越过了他能够判断的地方。

``师父\ldots\ldots''

后半句话没能出口。他跟在 test\_alpha0 身边这么多年，竟从没见过这副景象。

test\_alpha0 一步跨到他前头，从城楼上迎了出去，迎着那道黑光，把右手往前一伸，掌心朝上。

iblp，Infinite Basic Laver Pattern。这是他照着 Laver 表削出来的东西。良序之神头一回借他 I3 的时候，他看见了 Laver 表，可握不住，只能照着它一刀一刀地削，削到自己握得住为止。削完了，剩下的连 Laver 表一个零头都不到。

在这世上，魂力、肉身、灵魂、修为、意志，说到底都是一层一层的基本图案织出来的。普通魂师的力量打在表层，ω-Y 序列打在序数结构上，iblp 打在这一切的下面。它不在哪一层上动手脚，它从最底下，把一样东西的纹样拆开。

他动用它的时候，对面那个人还在其次，他真正怕的，是它本身。那东西大得离谱，手稍微抖一下，他自己也要被卷进去。

可他身后是一座城，城楼上还站着他的徒弟。

头顶的漩涡垂下一线幽蓝，落进他摊开的掌心，旋转、折叠，往里收。ω-Y 序列的递归结构做骨，iblp 做核，越收越紧，越收越亮，收到最后，凝成了一柄剑。剑是透明的，隔着它能看见远处星斗大森林的轮廓，可剑影落下去的地方，什么东西都像被抹淡了一层。

良序圣剑。

这是 test\_alpha0 照着良序之神的样子凝出来的剑。三百年里，他只凝过三次。天州城外，是头一次。

帝天看见了那柄透明的剑，黑龙剑却没有半点迟疑。

原本拖在剑后的漆黑剑气骤然收紧，压成一道窄而凝实的黑弧。黑弧前沿，一线空间被切开，正对着 test\_alpha0 的胸口；弧后的余势仍直指城楼，连他身后那门炮也一并罩了进去。

test\_alpha0 抬起右腕，良序圣剑由下向上斜挑。

一线幽蓝迎着黑弧升起。

两道剑气先在半空咬住，黑弧向里压，幽蓝却沿着它最凝实的中线切了进去。碰撞处亮得城头的人睁不开眼，随后，黑色剑气从两侧翻开，碎成一片片薄刃，落向城外麦田。

帝天已经随剑穿过了那些碎刃。他没有退，黑龙剑压着尚未散尽的黑光，劈到了良序圣剑上。

两柄剑在城门外百丈的半空中，正面撞在了一起。

轰！

一圈肉眼看得见的气浪从两剑相交的地方炸开。西墙那层应急光幕猛地一晃，护住了正面的守军，城楼两侧尚未接通的地方，仍有碎砖哗啦啦落下来。有人被震得跌坐在地，抱着头往垛口底下缩。

就是这头一剑，帝天的手臂麻了一下。

``回城东去。''test\_alpha0 没回头，``把护城大阵开启。''

Smile Lee 将手边的器件收入魂导戒，掠过城楼，直奔城东的主阵室。修复南北阵眼的魂师已经到了指定的位置，传讯器里，报出了一声接一声的``已接通''。

他拆下烧穿的汇流环，把自己的备用核心接了上去。掌心压住主控阵盘，九十八级的魂力灌进全城阵线，城东那道熄灭许久的阵纹，终于贯通到了城西。

``起！''

八处阵眼依次亮起，淡金色的光从四面城墙向上合拢，在蓝色漩涡下方，撑起了另一层穹顶。

天州的护城大阵，由他重新开启了。

西门下沿仍缺着一角。那里的阵基被铁甲犀撞裂，埋在断木和砖石底下，城门还在挨撞，修阵的人靠不过去。Smile Lee 把这个位置单独标红，先将其余几处的魂力稳住。

帝天没有立刻收剑，反而沉肩压腕，把力道又往下送了一重。

黑龙剑嗡嗡震颤，良序圣剑却稳稳横在那里。帝天看得清楚，握着它的那只右手，连指节都没有动一下。

他的眼神沉了下去。

下一瞬，黑龙剑猛地抽回，随即化作一片漆黑的剑幕。龙吟声骤然拔高，乌云随剑势倾压下来，城外的麦田陷进阴影，连伏在地上的魂兽都缩起了脖颈。

test\_alpha0 抬腕迎上。

两剑相击的声响密得分不出先后。黑蓝两色的火星沿剑刃迸开，落进麦田，烧出一串焦黑的坑。帝天每一剑都借着上一剑的回震蓄力，攻势越来越重，良序圣剑却始终挡在他面前，没有让开半寸。

帝天忽然撤了力。

黑龙剑贴着透明的剑身滑开，一道剑芒从两人交错的衣袍间穿出去，直奔城楼。炮口还在冒烟，Smile Lee 刚离开，那座炮架便被剑气照得漆黑。

test\_alpha0 横移半步，剑锋一转，截住了那道剑芒。

剑气在他身前崩散，余波却仍越过肩头，扫向城墙。城东主阵室内，Smile Lee 掌下的阵盘亮起一道红线。他并起两指，将南北两侧的魂力调往西墙，淡金色的光幕随之加厚。黑色的余波撞得光幕深深凹陷，又被分散到相邻的阵眼，碎成了烟雾。烟雾散去，城楼檐角那只脊兽仍蹲在原处。

帝天已经逼到他头顶。

他把黑龙剑横过来，剑脊上的暗金龙纹一片片亮起，聚在剑尖。

剑气翻卷而出，转眼长成一条数十丈的黑龙。龙首压下，张开的巨口将 test\_alpha0 连人带剑罩在里面。龙吟先一步撞向城墙，震得护城光幕一阵波动，城头的守军捂着耳朵伏了下去。

巨龙合拢双颚，四周的黑暗也随之往内挤。它要把那柄剑、握剑的人，还有他周围能够转身的地方，一起吞进腹中。

test\_alpha0 没有后退，右手向上一送。

剑尖前亮起一点幽蓝，直穿进黑龙上颚。

黑龙的头颅猛地一顿。龙身仍在往下压，鳞片沿着脊背一层层竖起，帝天灌进去的魂力从龙尾奔向龙首，要把那一点蓝强压回去。

良序圣剑没有弯。

蓝光沿着剑尖向外延开，先拆开了两只龙角，接着是上颚，再往后，连接龙颈的鳞纹也一片片散了。新灌进去的黑色魂力刚接上，就失去了落脚的地方，反卷成一圈黑雾。

龙吟断在半空。

几十丈的龙躯仍维持着俯冲的姿势，身子却从头到尾散开，露出后面帝天绷紧的脸。

良序圣剑的剑尖，只向上移了数寸。

帝天翻身压下，剑势由纵斩转为横扫，再从肋下反挑而起。黑龙剑几乎没有停顿，两人之间尽是交错的剑光。test\_alpha0 随着他的变招转动剑锋，每次拦截的幅度都很小，却总能截在力道最盛的地方。

帝天的黑暗领域随剑势向外再涨一重。麦田上刚露出的天光又被压暗，地上的影子一条条抬起，向着 test\_alpha0 所在的地方合拢。

良序领域的蓝色纹路就在那些影子底下亮起来。

两层领域在半空交错，挤得地面翻开一道道泥浪。黑暗想把整个空间拖向帝天，蓝纹却一处接着一处钉住，始终没有让出通往城墙的那一面。帝天每往前送一分魂力，剑身上便多一分沉重的回震。

乌云被他催得越来越低，到了城墙上方，却再也压不进去。天上的蓝色漩涡仍在转动，下方由 Smile Lee 维持的淡金光幕稳稳撑着，云层沿穹顶向两侧翻涌。城外暗得看不清麦垄，城内的旗杆仍稳稳立着。

又一次两刃相撞，帝天腕骨猛地一麻。

他立即收剑，强提魂力，将那股麻意压了下去。可剑再递出去时，手臂已经比先前沉了几分。黑龙剑上的龙吟短了一截，原本随心而动的剑气，竟在剑锋前迟滞了一瞬。

帝天盯住了那柄透明的剑。

方才每一次交锋，都有一丝异样的力量顺着剑身传回来。他起初还能震开，接连强攻之后，那股东西却在手臂里积了下来，越催魂力，越觉滞涩。

他抬眼看向对面。

test\_alpha0 的左手仍背在身后。剑风不断扯动他的袍角，右手中的良序圣剑却平稳得没有一丝颤动。

对面那个人的气息，一点都没乱。

良序圣剑上传过来的东西，比魂力深得多。那是宇宙本源里的良序之力。它不跟黑龙剑硬碰，只顺着剑身往里渗，一直渗进帝天握剑的手臂。被它碰过的地方，力道都被重新排得整整齐齐，再也使不出一点乱劲。

帝天出剑越来越慢，剑气也短了下去。他的领域被那股力量压着往回收，收到最后，只剩周身三丈。

帝天忽然抽回黑龙剑，左掌在剑脊上重重一按。

暗金色的龙纹没有再化成剑气。一缕漆黑的火从剑身上升起来，沿着龙纹烧遍整柄剑，又顺着他的手臂漫到肩头。

``万龙焱！''

火焰从他身后轰然铺开。

一条条火龙从黑焱里抬起头，张口、舒爪，密密麻麻地挤满了半空。最前一批直扑 test\_alpha0，后面的则绕向两侧，要从剑锋顾不到的地方合拢过来。黑火还没碰到麦田，底下的麦穗已经卷曲，贴着地面烧成了一片灰。

迎着这片龙焱，test\_alpha0 将良序圣剑横在身前，右腕一转。

``序元剑风。''

透明的剑身上分出了一柄同样的长剑，紧接着，又是第二柄、第三柄。越来越多的分身从本体中脱出，在他身周首尾相接，排成一圈圈倾斜的剑轮。

剑轮开始旋转。

最下方的一圈卷起麦田上的尘土，上面的剑影随之接续，一层叠着一层，连成了一道高速转动的龙卷。每一柄分身都带着凝实的锋刃，幽蓝色的轨迹在风里交织，始终循着固定的次序穿行。

第一条火龙撞了上来。

龙首刚探入剑风，就被迎面转来的剑锋切开。后面的黑焱翻卷着补上，第二圈剑影已经接到，一绞，将尚未成形的龙躯也扯散了。

更多火龙压下来，从上方、从两侧，接连扑进龙卷。黑火一度将蓝色剑影全盖住，里面却仍传出密密的剑鸣。每响一声，就有一截火龙从风里甩出来，还没飞远，又被回旋的剑锋切成了细碎的火星。

帝天握紧黑龙剑，将火势再往前推。

test\_alpha0 的剑尖也随之向前一送。

龙卷离开了他身前，迎着万龙焱倒卷过去。最前面几条火龙被推得翻转身躯，后面的还在往外涌，两股黑焱撞在一起，火势顿时乱了一瞬。

帝天立即撤去前方火龙，把剩下的黑焱收在剑前，凝成一道厚重的火幕。

剑风撞上火幕，轰地炸开一圈蓝黑相间的气浪。帝天横剑抵住，双臂猛地一沉。外层的剑影被他震散，内层却仍沿着原来的次序转来，一圈接着一圈，逼得他只能将力量压在正面。

最后一圈剑影掠过剑脊。

帝天刚把它挡开，右腕还在回震，万龙焱也才从剑锋上散去，来不及重新接成下一道攻势。

test\_alpha0 抓住了这一下。

他已穿过散开的火星，到了帝天身前。四周尚未散尽的剑影一柄接一柄归回本体，良序圣剑上的蓝光随之收紧。

他双手握剑，高高举起，重重劈下。

帝天横过黑龙剑去架。两柄剑撞上的那一刻，他整个人往下一沉，从半空里直坠下去十几丈，两只脚踩进麦田，一直陷到小腿。黑龙剑被压得弯了一弯，剑刃离他自己的脸只剩一寸。他两条胳膊上的青筋全鼓了起来，咬着牙往上顶，顶了好一会儿，才把那柄剑顶开半尺，趁着这半尺的空当往后一撤，一口气退出去三里，落在城外麦田的正当中。

城东主阵室里，Smile Lee 正把第二枚备用核心送进辅助槽。那一剑劈下来的时候，他透过敞开的窗抬头看了一眼。隔着大半座城，那柄剑比城楼还高，他看得清清楚楚。

他手里的核心停了半息，左手仍压在主控阵盘上。

师父出手的样子，他见过不止一回。可这一剑，他是头一回见。帝天那样的存在，接这一剑接得那么吃力，两只脚都陷进了地里。

他忽然觉得，师父手里的力量，叫人害怕。换作是任何一个普通的极限斗罗，挨这一剑，恐怕都劈成两半了。

阵盘上亮起一线过载的红光。他猛地回过神，把那块核心按进槽底，接过西墙传回来的压力。这时他才发现，自己手心里全是汗。

\begin{center}\rule{0.5\linewidth}{0.5pt}\end{center}

\chapter{第九十章：天州}\label{ux7b2cux4e5dux5341ux7ae0ux5929ux5dde}

帝天让那一剑劈退了三里。

兽潮最前头，熊君一直看着。它看着帝天两只脚陷在麦田里，看着那柄黑龙剑叫人压弯了，看着半空里那个人把透明的长剑往身侧一垂，停在西城门楼前方，剑尖斜指着城外。

那个人身上，是它前些日子在儿子身上见过的那样东西。

西北的麦田里，泥土忽然裂开了。

几条粗大的藤蔓贴地疾行，沿着蓝白色碎光之间的空隙钻过护城河，到了 test\_alpha0 下方，猛地向上窜起。一条绕向脚踝，另外几条从左右合拢，要把他连人带剑缠在里面。

丘陵上的万妖王，将扎进泥土的根须绷紧了。

藤蔓刚碰到袍角，test\_alpha0 周身的魂力便向外一震。

贴近他的那一圈藤蔓瞬间炸碎，断口沿着绷直的藤身向外崩开，一直碎到护城河边。半空里落下一片青黑色的碎屑，万妖王闷哼一声，探进麦田的根须齐齐抖了一下。

火光几乎在同一刻扑到了。

赤王等的就是藤蔓收拢的这一瞬。它从南侧丘陵跃下，贴着麦田掠来，身上的火焰拖成一道长线。等它看清那些藤蔓已经碎了，探出的前爪离 test\_alpha0 只剩数尺。

test\_alpha0 将良序圣剑垂在右侧，左手握拳，迎着它一拳打出。

拳头穿过火焰，正中赤王胸前。

一声闷响，赤王的身子在半空里猛地弓了起来。前扑的势头当场倒转，身上的火也被压得向后翻卷。它一路倒飞，掠过数里麦田，重重撞进西南那片丘陵，把半面土坡砸得塌了下去。

碎石和断树翻滚着落下来，赤王撑着前爪从泥里挣出半个身子，张嘴喷出一口血。火焰在獠牙之间窜了两次，都没能吐出来。

test\_alpha0 收回左手，右手中的剑仍稳稳垂着。

万妖王没有再催动藤蔓。赤王伏在远处塌下的土坡旁，也没能立刻站起。两侧原本正往前压的兽群，随着这两次失手停在了领域边缘，没有再向城后合拢。

熊君没再等帝天。它抬起前掌，在地上重重按了一下，冲身后吼了两个字。

``攻城！''

城外那一整片，动起来了。

麦田上头还悬着那些蓝白色的碎亮点。冲在最前头的那一排撞上去，当场没了一片，连声都没出。后头的没再躲，踩着前头空出来的地方，接着往前涌。

兽潮涌到护城河外那道土埂，Smile Lee 埋下的另外两道魂导炮弹一齐炸了。火光从西城门口一直烧到城墙拐角，冲在最前头的几百头魂兽让掀翻在地，后头的踩着它们的身子，接着往前涌。

西城门响了一声。最后那根门闩，断了。

三人合抱的铁木从当中崩开，两扇城门向里倒去，砸起漫天尘土。门边那座早已撞裂的阵基，也跟着塌进了瓦砾。西门下沿缺着的那一角，骤然张开。

兽群盯上了这个缺口。最前面上千头风狼压低身子，朝门洞扑来，獠牙上还挂着前头那两座小城里带出来的血。风狼后面是铁甲犀，鼻端那根独角像攻城锤一样往前顶。

就在城门倒下的那一刻，一道人影从半空落了下来。

他落得不快，衣袍连一道褶子都没有，就那么站在了城门外三丈的地方，正好挡住那道豁口。手里那柄透明的剑垂在身侧，剑尖离地一寸。

冲在最前面的一头风狼收不住脚，张嘴就朝他扑过来。离他还有一丈远，那头风狼像撞上了一堵看不见的墙，整个身子在半空里顿了一顿，重重摔在地上，四肢抽搐，再没爬起来。

后面的风狼群硬生生刹住了脚，伏低身子，喉咙里呜呜地低吼，没有一头敢再往前。

test\_alpha0 没有回头，也没再开口。该说的话，刚才在城楼上已经说过了。他抬起眼，望着城外那一望无际的兽潮。

兽潮深处，又是一声震天动地的咆哮。

\begin{center}\rule{0.5\linewidth}{0.5pt}\end{center}

数以万计的魂兽同时向前冲，大地在它们脚下剧烈地颤抖，天州城的城墙跟着一阵阵呻吟。

test\_alpha0 把剑换到左手，抬起右手，五指张开。

一圈幽蓝色的光，以他为中心，朝四面八方荡开。光所过之处，天色一下子暗了下来。城外的大地上浮出密密麻麻的纹路，一层套着一层，一直铺到看不见的地方。那是 ω-Y 序列铺开的递归结构，最上头压着的，是 iblp 的威压。

冲在最前面的几百头魂兽，像被一只看不见的大手从天上按了下来，一头接一头趴倒在地。几头十万年级别的魂兽嘶吼着，修为在这片空间里被一层一层往下压，压到最后，连站起来的力气都没有了。

铺天盖地的冲锋，在城门外两百丈的地方，戛然而止。

城头上死一般地静。不知是谁先喊了一声，紧接着，守军爆发出一阵欢呼。

欢呼声还没落下，城外就响起了一阵叫人牙酸的骨头摩擦声。

兽潮最前方，熊君一点一点站了起来。

它站直了足有百丈，比天州城的城墙高出一大截。两只前掌探出一丈多长的利爪，嘴角淌着血，四条腿在威压底下不停地打颤，可它还是挺直了腰。

五十万年修为，大陆第一熊。

``就算你压住了所有人，''熊君的声音像闷雷一样从城外滚过来，``你也杀不了我！''

test\_alpha0 没有反驳。威压镇得住十万年，镇不死五十万年，这一点他比谁都清楚。

他把良序圣剑换回了右手。

熊君看着那柄剑，浑身的毛根根倒竖。五十万年的本能在它身体里疯狂地叫，叫它跑。

它没有跑。

它仰起头，发出一声长啸。这一声，冲着的是它身后。

长啸声里，原本被威压死死按在地上的七头十万年魂兽，一头接一头地站了起来。打头的是一头肩生骨刺的赤红巨猿，后面跟着一条几十丈长的碧鳞巨蟒。它们浑身都在发抖，七窍里往外渗血，可还是一步一步挪到了熊君身边，跟它站成了一排。

它们本来可以一直趴着的。威压只是压着它们，test\_alpha0 不会动它们分毫。

是它们自己站起来的。

城头上的欢呼声，一下子断了。

七头十万年魂兽同时放出了各自的领域，一层叠着一层，全压在熊君身后。熊君的气息猛地暴涨，它把两只前爪高高举起，暗金色的魂力顺着爪尖往上攀，烧得四周的空气都扭曲了。

撕天爪。

扽魂殿的卷宗里，这一招是记着名字的。一百年间，死在这一爪下面的封号斗罗，有名有姓的就有十几个。有几具尸首是 test\_alpha0 亲自去收的，从当中断开，断口平得像镜子。那些卷宗上写的都是同一句话：违禁入林，遇袭身亡。

一百年里，熊君从没走出过星斗大森林一步。

这一回，它为了儿子，走出来了。

``吼！''

熊君挟着七重领域，一爪撕下。

天空像是被这一爪划开了一道口子。暗金色的爪芒带着毁天灭地的气势直扑下来，扑向 test\_alpha0，也扑向他身后那道洞开的城门。

test\_alpha0 握紧良序圣剑，把 iblp 压在剑尖上。压着，不敢全放。

然后，往前一刺。

``神圣突刺。''

剑尖和爪芒，在城门外三十丈的地方，正面撞在了一起。

一道裂痕，出现在 test\_alpha0 脚下的序数空间上。撕天爪的力量硬是在 ω-Y 序列的递归结构上撕开了一道口子，口子往两边飞快地张开，一直张到他脚底下。三百多年里，能在他的序数空间上撕开口子的魂兽，只有熊君一个。

那一刻，他是佩服它的。

可撕天爪也只走到了这一步。

iblp 从最底下那层开始拆。爪芒还在往前推，织出它的那层图案却已经开始崩散。肉、骨、魂力，五十万年一天一天熬出来的修为，在剑尖面前一层一层失了依托。

剑没有停。

良序圣剑穿过熊君的前爪，穿过它的胸膛，从它后背透了出去。

熊君离 test\_alpha0 只剩两丈。它低下头，看着胸口那柄透明的剑，又看着自己那只前爪从边缘开始，一点一点淡下去。不断，也不碎，就那么淡下去，像一幅画被人从最底下那一笔开始往回抹。

它抬起头，看了 test\_alpha0 最后一眼。

百丈高的身躯轰然往前倒下，砸在城门外的大地上。整座天州城跟着震了三震。它的头落在 test\_alpha0 脚边，离那扇破门，还有三十丈。

test\_alpha0 想让剑停在熊君这一层。剑没有停。那时候的 iblp，他用得了，却还没用到家，那一刺已经压着力了，余波还是收不住。站在熊君身后的七头十万年魂兽，就那么站在原地，从里头散开了，连一声叫都没来得及发出。

再往外，是被威压按在地上的四十九头万年魂兽。余波扫过去，一样化成了虚无。

五十七。

后来，test\_alpha0 一头一头地数过。

\begin{center}\rule{0.5\linewidth}{0.5pt}\end{center}

城东主阵室里，Smile Lee 掌下那道标红的阵纹，终于不再剧烈跳动。

西门外的兽群被压住了。

``就是现在。''他对着传讯器说，``把备用阵座送到门洞内侧，不用搬原来的。按我给的位置落下。''

几名城防魂师抬着备用阵座，从街角冲了出来。他们越过倒下的门板，把阵座卡进门内两侧的砖缝，依次接上导流纹。外头那头百丈高的熊就倒在三十丈外，几个人不敢多看，接好一处便拍一下传讯器。

最后一声传回来，Smile Lee 立即合上了主控阵盘边缘的银环。

淡金色的光从门洞两侧接到了一起。西门下沿最后那个缺口，终于封住了。

他又将辅助核心全部接入，把各处阵眼的负载重新分匀。光幕从城西连过城北、城南，一直落到东门外。东门仍留着一道能控制开合的通道，撤退的人排着队，一批一批从那里往外走。

城墙上的守军扒着垛口往外看。隔着一层流转的淡金光，城外的麦田、倒下的巨熊，还有三里外重新站直的黑衣人，都看得清清楚楚。

Smile Lee 望了一眼西边，没离开主阵室。他把右掌压实，盯住了阵盘上代表城西的那道纹路。

帝天低头看了一眼握剑的那只手。虎口裂开了一道口子，血顺着剑柄往下淌。

再抬起头的时候，他眼里那点对人类的轻慢，已经一丝都不剩了。

他往城门那边看了一眼。熊君倒在那儿，头朝着那扇破门，它身后那七个，一个都不剩。

万龙焱没能烧穿那道剑风，黑龙剑也没能挡住随后压下的一击。方才抵挡剑影时留下的麻意，还盘在他右腕里。再拿剑与火焰硬拼下去，只会继续落进对方掌握的节奏。

他把黑龙剑换到左手，慢慢抬起了右手。

先变的是眼睛。瞳仁一点一点收窄，竖成了一条线，颜色也跟着变，变成了熔金一样的金色。星斗大森林里的魂兽叫他金眼黑龙王，叫的就是这双眼睛。

然后是那只手。皮肉底下翻出一片一片漆黑的龙鳞，每一片鳞的边上都镶着一道暗金色的细纹。指甲往外长，弯成五根锋利的钩，鳞片一路往上爬，爬到了肘弯。转眼之间，他右臂的尽头，已经是一只狰狞的龙爪。

龙神爪。

龙威是跟着这只爪子一块儿出来的。城外还活着的魂兽，离得近的离得远的，一头接一头趴了下去，脑袋贴着地，尾巴夹得死紧。城墙上的守军膝盖发软，有人扶着垛口，才没跪下去。

帝天把那只龙爪举过了头顶。

他身后的天上，跟着浮起了一只爪子的影子。形状跟他手上那只一模一样，大小却没法比，五根爪尖张开，从城外麦田的上空，一直罩到了天州城头。爪影底下的光，一下子就没了。天州头顶那团蓝色的漩涡被一点一点挤小，乌云压过了城外的麦田。

爪子还没落下，半空里先多出了五道黑线。那是爪尖划过的地方，空间叫它撕开了口子，口子里头什么都没有，黑得看不见底。龙爪底下的田地，在爪芒落下之前，就已经往下陷了三尺。

这一爪要是落下来，城门，门后的大半座天州城，连同站在城门前的 test\_alpha0，都在爪下。城里，Smile Lee 还守在城东，满城的百姓还没撤完。

test\_alpha0 认得这一爪。

五百年前，扽魂殿一位九十九级的极限斗罗独自进了星斗大森林，从此再没回来。殿里找了很多年，最后才确认，他是死在龙神爪底下的。九十九级的极限斗罗，一爪。

刚才对熊君那一刺，他是压着出的。压着，还是多搭进去了五十六条命。

这一回，他压不住，也不能压。

\begin{center}\rule{0.5\linewidth}{0.5pt}\end{center}

test\_alpha0 把 iblp 彻底放开，一层不留。

紧接着，他做了一件三百多年来从没做过的事。良序之神交到他手上的那样东西，Laver 表本身，他从里头掐出了一丝，融进了剑里。极少的一丝，少到连他自己都怀疑它在不在。

良序圣剑的光芒，在这一刻暴涨到了极致，亮得他自己都看不清剑身。

天州头顶，那团蓝色的漩涡忽然向内收紧。

最外面一圈先卷了回来，接着是第二圈、第三圈。一层层魂力挤向中心，原先还能透出云影的浅蓝，渐渐压成了深蓝。越往中心，颜色越深，厚得近乎凝成了实物。漩涡还在缩，每缩一分，里面的能量便密上一分，城头上的人已经分不清那些旋转的纹路，只看见一片深蓝悬在头顶。

黑云从四周压来，撞到那片蓝的边缘，反倒被逼得翻卷出去。

漩涡当中，慢慢显出了一道笔直的轮廓。

先是剑脊，随后，两侧的蓝光向内折叠，压成剑身。最下方的一点光越收越尖，从云底垂出。仍在旋转的魂力顺着剑身向上汇去，将剩下的轮廓一寸寸补齐。

整团漩涡，最终凝成了一柄巨剑。

剑尖悬在城外麦田的上空，剑身一直没入高处的云层。深蓝色的剑脊里，一线白光贯穿上下，剑锋附近的空气被压得扭曲，远处的黑云隔着它望过去，像被分在了两片不同的天空里。

把修阵的人送到各处阵眼后，写信的校尉已经回到了西城门楼。他原本正用右肩顶住一个松开的炮架，这时忽然停住，仰起了头。

那柄剑的影子越过炮架，越过垛口，一直铺到了身后的街上。他抬起右手遮眼，蓝光从指缝间漏下来，照得掌上的炭灰也清清楚楚。

校尉先看了看云里的剑，又去找城外握剑的人。两柄剑，一大一小，映着同样的蓝光。他扶住垛口，半晌没能把气吐出来。

放飞传信鸽的时候，他想的是请总殿再来一位长老，替天州多挡一阵，让街上的伤员能挪到东门。如今那个人站在城外，头顶的整片天空，都成了他的剑。

test\_alpha0 向前跨出一步，身形越过护城河，掠到了城外三里的半空。

帝天正将龙神爪往下压。那只遮住半边天的爪影迎面落来，五道黑线向着城池延伸，所过之处，麦田连着底下的泥土一齐塌陷。

test\_alpha0 双手握住良序圣剑，举过头顶。

云层里的巨剑，也随着他的动作扬了起来。

``神圣审判！''

他把手中的剑往下一劈。

天上那柄巨剑随之斩落。

深蓝色的剑锋切进黑云，从高处一路压下，剑脊里的白光骤然铺开。云层从当中断成两半，裂口沿着剑势向西延伸，一直拉到星斗大森林上方。巨剑下方，尚未触到剑锋的麦田已经伏倒，泥土被迎面推来的力量刮去一层。

城东，Smile Lee 从阵盘上抬起头。隔着淡金色的光幕，他看见那柄由漩涡凝成的巨剑劈开了云层。城墙、炮楼，连同他脚下的主阵室，都映进了剑身落下的深蓝里。

传讯器里有人惊呼，声音刚起，又被他压了下去。

``各处守住阵眼。都别松手。''

他把全部心神收回阵盘。西墙传回的压力已经开始往上攀。

巨剑的剑锋与龙神爪，在城外三里的地方，轰然相撞。

那一下，天地之间什么声音都没有了。

\begin{center}\rule{0.5\linewidth}{0.5pt}\end{center}

百余里外，极山宗的最高峰上，hypcos 正盘膝坐在峰顶。

四周没有树，只有一圈被山风磨得发白的岩石。薄云贴着脚下的山腰流过，他的呼吸很长，体内的魂力循着修炼的路径，一层层收向深处。

忽然，那股魂力停了一瞬。

hypcos 睁开眼睛，转头望向天边。

一阵巨大的能量波动，越过重重山岭，从百余里外传了过来。到了这里，已经没有震耳的轰鸣，可其中那股沉重的压迫感，仍让他下意识地收紧了掌心。

他站起身，走到峰顶的岩石边。

遥远的天边，云分成了蓝黑两种颜色。黑云层层翻卷，几乎连成了一片，蓝光却从中间透出来，将厚重的云层撑开。两种颜色交界的地方，忽然亮起一道白里透蓝的光，隔着这样远的距离，仍清楚得让人移不开眼。

那道光一闪，黑云便向两侧翻开。紧跟着传来的第二阵波动，比先前更强。

hypcos 的衣袖被山风卷起。他没有动，望着天边的神色却一点点凝重起来。

那阵波动里，有一股气息是他认得的。

魂力一层扣着一层，向外铺展时，各处的衔接严整得没有一丝偏差。撞进其中的杂乱力量，被一缕缕分开，再沿着同一个方向压了出去。隔了百余里，那种熟悉的良序气息仍然清楚。

``test\_alpha0\ldots\ldots''

他低低念出了这个名字。

不会错。蓝色云层里的那道力量，是良序圣殿的大长老在释放。

可认出是谁，并没有让他松一口气。外面那几层魂力，他能辨认；藏在深处、随着那道蓝白剑光一齐落下的东西，却让他刚探出去的感知猛地收了回来。

他从没见过 test\_alpha0 把力量放到这个地步。

他不自觉地扣住身旁的岩石，指下簌簌落了几粒碎屑。他低头看了一眼，又抬起头，盯住那片正在被劈开的黑云。

究竟是什么对手，逼得他动用这样的一剑？

那是天州所在的方向。可相隔百余里，群山遮住了城池，hypcos 看不见地上的兽群，也看不见城头正在苦撑的人。他认出了大长老，却还不知道另一股力量属于谁，更不知道这场交锋因何而起。

蓝黑交错的云仍在缓缓翻动。

他没有再坐回去。

\begin{center}\rule{0.5\linewidth}{0.5pt}\end{center}

城外三里，巨剑压住了龙神爪。

五根爪尖猛地向内扣紧，暗金色的纹路从爪腕一直亮到尖端，竟将深蓝剑锋夹住了一瞬。剑与爪之间，一片片黑色裂口接连张开，刚张开，又被两边挤来的力量压合。

帝天右臂上的鳞片全立了起来。龙威沿着爪影向前奔涌，他把最后一重力量也推了上去，要将那柄巨剑从中间撕开。

test\_alpha0 握剑的双手仍然很稳。

他将剑锋再往下压了半寸。

云层间的巨剑随之下沉，剑脊深处那一线蓝白光透过了暗金纹路。龙神爪表面仍在向外喷涌力量，撑住它的纹样却开始一处接着一处失去衔接。帝天想把这些地方重新接上，剑锋已经从爪影中央斩了进去。

两股力量在这一刻彻底错开。爪影被截在城池之外，神圣审判的剑气顺着裂口，继续向帝天所在的方向压去。半空里那五道黑线，过了好一会儿，才一点一点合上。

test\_alpha0 在这里挡住了帝天这一击，碰撞的余波却仍向四面八方挤开。贴着地面卷起的气浪裹着碎石，越过三里麦田，一路扑向天州，重重撞上了 Smile Lee 开启的护城大阵。

西边的光幕整片向内凹下去。

城东主阵室里，八道阵纹同时亮红。Smile Lee 两掌按住阵盘，将西墙承受的冲击一股一股引向南北，再经城东的辅助核心分散出去。第一枚核心的外壳裂开，第二枚紧接着接上，淡金色的光在城墙上急急地流，像有一整圈火沿着城池跑过去。

门洞内侧新接的阵座抖得最厉害。一名城防魂师扑上去，双手抱住底座，另两个人把魂力接到他背上。三个人的鞋底在砖地上拖出半尺，阵座终于没再滑。

Smile Lee 手背上的筋一根根绷起。他把最后一处松开的阵纹压回原位，光幕终于止住了内陷，沿着城墙一点一点撑回来。

碎石落在光幕外，噼啪响了好一阵。城墙上的守军伏在垛口底下，等那阵响动过去，才敢抬起头。这一次，城内没有再落下砖石。

校尉扶着一个被震倒的炮手，从垛口底下慢慢站起来。刚才金色光幕陷下来，他几乎以为整座城都要跟着塌了；再抬眼，城墙还在，街上的伤员还在往东挪，那只罩住天州的龙爪，却已经被挡在了三里之外。

他望着城外，眼睛被蓝白光刺得发酸，连呼吸都忘了接。直到炮手抓住他的右臂，他才猛地回过神。

``扶伤员进去。''他说，声音哑得不像自己，``别堵在这儿，让出路来。''

说完，他又忍不住往那个人站着的地方看了一眼。

烟尘慢慢落下来。城外三里，两股力量撞上的那块地方，塌下去一个大坑。坑口方圆两三里，边沿齐整得像拿刀削过，坑壁上的土烧成了黑亮的一层，往下少说有几十丈深。

帝天已经不在坑边了。

他被这一剑推到了离城十里的麦田里。

天上那只爪影先散了。跟着是他右臂上那只龙爪，从正当中被劈成了两半，从指尖一直裂到手腕。剑气劈开龙爪，去势未尽，顺着那道裂口往里走，刺进了他的身体。

帝天低下头，看了一眼自己的胸口。漆黑的血顺着伤口往下淌，浸透了那一身黑衣。

神圣审判的剑气，绝大部分叫他挡下来了。龙神爪接了一截，他自己的身子又接了一截。从他身边漏过去的，只是剩下的那一点。

就这一点，从他身后那铺满了整片麦田、一直延伸到星斗大森林边缘的十万兽潮里，从头扫到了尾。

剑气碰上魂兽，先没的是修为。iblp 和那一丝 Laver table 的力量压下来，一头魂兽几万年攒下的魂力，一圈一圈往外散。跟着是身子，从皮毛到骨头，一寸一寸地瓦解，碎下来的东西还没落地，就没了。

剑气最底下还压着一层，是宇宙本源的良序之力。它不碰皮肉，专往根上走，把这些魂兽修炼的根基一并毁了。

剑气一路推到星斗大森林边上，才停下来。它推过去的地方，十万魂兽大军，一头都没剩下。

碧姬一直在兽潮最后头，翅膀张着，身上那层翡翠色的光一丝一丝连到前头的魂兽身上，哪一头伤了，光就往哪一头身上多走一点。十万头，十万丝。

剑气扫过去的时候，那十万丝光是一下子断的。

断下来的劲，顺着光，一下子全回到了她身上。碧姬从半空里掉了下来，摔在星斗大森林边上的草里。翅膀上的翎毛掉了一大片，她张嘴吐了一口血，血是绿的。她想把翅膀收起来，左边那一只早先让城楼上那颗九级炮弹燎过，收到一半，就耷拉了下去，再也抬不起来。

那还只是帝天没挡住的那一点余波。

云层间的巨剑渐渐散去。尚未消尽的蓝光沿着剑身原先所在的位置卷回来，在天州上空重新聚成一团漩涡。它比先前淡了许多，仍在缓缓转动，城外那些悬着的细光，也跟着它转。

这一回，test\_alpha0 没有去数。

数不过来。

他的左肩忽然一凉。低头一看，三道深可见骨的口子，从肩膀一直拉到胸前。龙神爪被劈开之前，爪尖终究还是够着了他。伤不重，半个月就合上了，可那三道疤，一直留到了今天。

三百多年，帝天是头一个让他见了血的魂兽。他也是头一个劈开帝天龙爪、把剑刺进帝天身体的人。

从那一刻起，两边之间的这笔账，就再也算不清了。

\begin{center}\rule{0.5\linewidth}{0.5pt}\end{center}

test\_alpha0 收住剑势，从城外三里返回，落在西城门前。

十里之外，帝天抬起头，与他隔空对视。

帝天还站着。他是拿黑龙剑撑着站的，剑尖扎进土里半尺。胸口那道口子还在往外冒血，他每喘一口气，肩膀就往下塌一点。右臂上那只裂成两半的龙爪垂在身侧，抬不起来了。

他扑不动了。

test\_alpha0 看了他一会儿，把良序领域收了回来。

满城外悬着的那些 iblp 碎片，一点一点往上浮。细光沿着来时的纹路升回天上，没进那团蓝色的漩涡。最后一片也收尽，漩涡才向中心缓缓合拢，露出后面的云。城外的风又刮起来了，坑里冒出来的烟，这才歪歪斜斜地往天上去。两侧还活着的那些魂兽，身上一松，把脑袋从地上抬了起来。

test\_alpha0 站在城门前，开了口。声音不大，可他知道，十里，帝天听得见。

``帝天，你若下次再敢来犯，老夫必将亲手斩杀你。''

帝天张口，吐出一大口血。

过了好一会儿，他的声音才从十里外传过来，带着几分艰难。

``人类，今日之事，本座记下了。''

说完，他仰起头，喉咙里滚出一声悠长的龙吟。龙吟传遍了整片战场。西北丘陵上，万妖王将残断的藤蔓收回土里；另一侧，赤王撑起身子，从撞塌的土坡旁一瘸一拐地退回兽群，胸前的火光已经暗了。它们带来的两支兽群始终没能合围，此刻与战场边缘侥幸活下来的魂兽一道，纷纷掉头，朝着星斗大森林退去。

帝天是最后一个走的。左手提着黑龙剑，右手那只被劈开的龙爪垂在身侧，一步一步，走回了星斗大森林深处。

后来 test\_alpha0 听说，帝天回到森林最深的地方，整整躺了五十年。爪子上的伤几年就长好了，真正要命的是刺进身体里的那一剑。剑上那一丝 Laver 表的力量，在他身子里头一层一层地拆，他用了整整五十年，才一点一点化干净。

天黑以后，天斗帝国的医疗兵赶到了。马车一辆接一辆从东门进来，伤员从城墙根底下一个一个抬上车，城东那几条街上，点了一夜的灯。

Smile Lee 等兽群退进林子，才将护城大阵降到维持的强度。他在主阵室旁坐了一会儿，手指还在不受控制地发抖。换下烧裂的核心，把各处阵眼交给城防魂师轮值，他又拿着器件上了城墙。

那一夜，哪段光幕暗了，他就去哪段补。天亮以前，砸坏的防御魂导器一座一座修了回来，西门外的光仍连成一片。第二天，他带人把熊君抬回了星斗大森林边上，放下，人就回来了。

没有人取它的魂骨。

城外那个坑，后来一直没有填。

过了些年，天州人在坑底立起了一座 test\_alpha0 的石像。石像比天州的城楼还高，双手拄着剑，脸朝着星斗大森林的方向。坑沿上又修了一座纪念堂，堂里的墙上刻满了名字，那天死在城里的守军和百姓，还有总殿来的那两位长老，一个挨着一个。

\begin{center}\rule{0.5\linewidth}{0.5pt}\end{center}

殿上，test\_alpha0 的手还按在石栏上。

``那天晚上让我睡不着的，不是熊君。''

``是帝天转身时候的那个样子。''

``他没怕我。他在记账。''

卢本韦背上起了一层凉。

``他记到今天，一百六十八年了。''

卢本韦沉默了一会儿。

``大长老。''

``嗯。''

``您那天为什么不杀了他？''

test\_alpha0 转过头来。

``熊君死了，正面的十万兽潮没了，帝天被您推回去十里，那只龙爪劈成了两半，身上还扎着您那一剑，两翼也没再围上来。''卢本韦的声音很平，``那是一百六十八年里最好的一次机会。您那天要是动了手，就没有今天这笔账。''

test\_alpha0 看了她很久。

``你很像你曾祖。''他说，``一百三十六年前，他坐在你现在站的地方，问过我同一句话。''

``您怎么答的？''

``我没答。''

风从柱子之间过去。

``我能杀。''

test\_alpha0 说这三个字的时候，像在说今天天气不错。

``但是他不能死。''

卢本韦看着他，等着后面的话。

``帝天是黑龙一族的族长。''test\_alpha0 说，``可他管的，不止黑龙一族。星斗大森林里，各族争领地、争水源，翻起旧年的血债，最后都得有人压住，有人调停。''

``龙族和其他族群之间，也有不知道积了多少年的恩怨。他站在那里，有的敬他，有的怕他，也有的恨他。可真到了要把一族赶尽杀绝的时候，他说停，还能让双方停下来。''

卢本韦回想起他方才讲过的兽潮。

那些一同向天州冲去的魂兽，原来也彼此有仇。

``他一死，谁来管？''test\_alpha0 问。

卢本韦没有立刻回答。

``旧仇先翻出来，接着就是抢地盘、争谁来做主。今天这一族报仇，明天另一族报复，打到最后，连本来不想打的，也会被卷进去。''

他的手指在石栏上慢慢收紧。

``没有能让各族都认下的人接住这个位置，魂兽就会陷进无穷无尽的族群内斗。它们不会因为人类杀了帝天，就从此安分下来。''

``可他回去以后，躺了五十年。''卢本韦说。

``重伤，和死了，是两回事。''

test\_alpha0 抬眼看她。

``他还活着，旧规矩就还有人撑着。今天趁乱灭了别族，来日帝天起身，这笔账还得算。可他真死了，那些压着没动的人，就要重新打算了。''

``森林里先乱，败下来的族群便往外迁，抢新的地方活命。边上的村镇要不要防？被挤走的弱小魂兽，又去哪里？''

卢本韦垂下眼，过了片刻，才重新开口。

``所以，您把他打退，却没有追。''

``我要让兽潮退回去，也得想清楚，它们退回去以后会怎样。''

test\_alpha0 停了一下。

``还有，我手里的力量。''

卢本韦看向他按在栏上的手。

``天州那一剑，到底有多大的威力？''

``不知道。''

test\_alpha0 答得很快。

``就算只融进去少量的 Laver table，也可能已经到了 PTO(Z\_ω) 的力量。甚至更强。''

``那一年我一百四十九岁。iblp 我用得了，也有几分把握，可离出神入化还远。iblp 那一层在 PTO(Z₃)，收拾五十万年，够了。要收住它，不够。再往里掺一丝真东西，剑落下去到了哪一层，我当时说不准，今天也只敢说个大概。''

``熊君那一剑，就是这么落下去的。我想落在一头熊身上。落下去的，是一整层。''

``帝天那一剑，也是这么落下去的。''

``我的力量，我自己都无法控制。''test\_alpha0 说，``越往深处借，剑落下去以后能收住多少，我就越没有把握。''

卢本韦看着他。

``我那时候手里能动的，不止我自己那一层。''test\_alpha0 说，``还有 PTO(ZFC+I3)。''

``那是什么层次？''

``你现在这样问，说明你还想象不到。''test\_alpha0 没有笑，``换个说法。我自己那一层，在 I3 面前，是一粒灰。''

``Laver table 的全部力量，一层不留地展开，没有同等层次的力量护着，一样东西，直接就从'存在过'里面抹掉了。''

``可要是对面也有足够强的力量顶着，就得先破开它。神力借得到，不等于谁都能一剑抹掉。''

卢本韦皱起眉。

``这世上有哪一门功法能修到那儿？''

test\_alpha0 看了她一眼。那一眼里有点别的东西，像是终于有人问到了这一句。

``没有。''

``一门都没有。人这一族的修法，往死里推，推到头是 ω-Y(1,ω)。再往上，路是断的。没有功法，没有前人，什么都没有。''

``我能比那走得远一截，是因为 iblp。那是我自己立的，前头没有人走过。有了它，我才握得住比 ω-Y(1,ω) 更强的力量。''

``可就算多走了这一截，离 I3 也还远得看不见。''

``那您怎么够得着？''

test\_alpha0 没有马上回答。他抬手指了指头顶。

穹顶之上，良序之神像垂着眼睛，金色的光从那双眼睛里落下来，铺了满殿。神座正面镌着初代良序神王的名字：Zermelo。

``我是良序之神的大供奉。''

殿上静下来。

``那一层不是我的力量，卢本韦。那是祂的。''

``我有资格开口向祂借，祂也一直借给我。上一代大长老把大长老和大供奉这两个位子，一块儿交到我手上。从那天起，两百二十八年了，祂从来没有拒绝过一次。''

``可借来的东西，终究不是自己的。''

他把手放下。

``祂头一回借我 I3，就让我看见了一件事：在 ZFC+I3 底下，Laver 表是良序的。不管从哪一张表往下走，总走得到底。''

``我看见了，可我握不住。那股力量从指缝里往外漏，漏一点，我身上就塌一块。''

``所以才有了 iblp。我照着 Laver 表削，一直削到自己握得住为止。削完了，剩下的连它一个零头都不到。''

``就这一个零头，再掺上一丝真东西，在天州城外抹掉了十万条命，叫帝天躺了五十年。''

殿上很静。

``我连自己削出来的 iblp 都没收住。托着 I3 下去会怎样，你自己想。''

``那现在呢？''

``现在，我对自己修出的那一层，比当年明白得多。''

test\_alpha0 看着自己的手。

``那一夜之后我开始往回走。七十三年以后，先把 ω-Y 序列的良序化证完了。那是垫脚的一层。''

``iblp 在它上头，高出不知多少。我又推了九十五年。''

``今天早上刻完的那张证明，就是 iblp 的。到今天我才真正懂了，那一剑是怎么落下去的。''

``如今我自己站到了 Z\_ω。祂借给我的，也从 I3 换成了 I1。''

``iblp 的道理证清了，借来的神力却还在更高处。知道自己能借到 I1，不等于我已经能把它每一分都收放自如。''

他停了一下。

``可天州城外那些命，是在我还没懂的时候抹掉的。''

卢本韦想了一会儿。

``I3 到底有多大？''

``我也不知道。''test\_alpha0 说，``我没跟你客气，我是真不知道。''

``我能借到它，能把它展开，能指着一个方向让它落下去。可它落下去以后铺开多宽，我按不住。''

他想了想，换了个说法。

``就像你手里有一把刀，要去割面前一尺远的一根线。''

``刀是好刀。可这把刀有多长，你不知道。你只知道它比你的胳膊长，比这座殿长，比你眼睛能看见的一切都长。''

``你挥下去，线断了。''

``线后面那些东西呢？''

卢本韦说不出话。

``熊君那一回，我已经见过一次了。''test\_alpha0 说，``我只想抹掉一头五十万年的。''

``跟着没的是七头十万年的，外围又散了四十九头万年上下的。''

``那七头是自己爬起来站到它后面去的，这我认。''

``外围那四十九头，是跟着来攻城的，这我也认。可剑刺出去那一刻，它们已经被威压按在地上，连头都抬不起来了。''

``它们是被余波扫到的。''

殿上静了一下。

``五十七条命，只有一条是我想要的。''

``那还是我削过、只敢放出一小段的力量，我都没收住。I3 整个铺开是多大，我说不上来。大到我连'多大'这句话都问不出口。''

他抬起眼睛。

``真要为了斩尽帝天，把那股力量继续展开，剑落下去，会停在哪里，我不知道。''

test\_alpha0 望向殿外。

``也许只死他一个，也许余波越过兽潮，扫进森林深处。那些没有来攻城的族群，那些还在巢里的幼兽，会不会一起被抹掉，我没有把握。''

``他的死，会让各族失去调停的人；我的剑，还可能毁掉它们赖以活下去的地方。这两个后果，都不能不算。''

卢本韦的手慢慢从行囊带子上松开。

她方才想着的，是帝天倒下以后，天州便再也不必防着他。可大长老说到这里，她才发觉，自己的念头一直停在那座城外，没有再往远处想。

``如果魂兽真的灭绝了呢？''她低声问。

``人类也会面临危机。''

test\_alpha0 答得很慢。

``魂兽在山林里捕食、繁衍，也在维持那里原有的平衡。哪一族多了，哪一族少了，影响的都不止它们自己。真把魂兽全抹掉，留下来的山林、水土，还会是原来的样子吗？''

``人住的城，看着离森林很远。可城里人喝的水、种粮的土地、赖以生活的山川，都和外面连着。森林里的平衡坏了，后果迟早要落到人身上。''

他转过头，看着卢本韦。

``我不能护住一座天州，却给后面的人留下一个活不下去的世间。''

殿上静了很久。

``天州那十万条命，是我杀的。''他说，``我必须记着。可已经杀了那么多，更不能拿这个当理由，觉得余下的也可以一并抹掉。''

卢本韦望着他，终于听明白了那句''我能杀''后面的分量。

他能将帝天斩在城外，却没有办法保证，斩下那一剑以后，两族还能怎样活下去。

``祂肯借我力量，落剑的人却是我。''test\_alpha0 说，``后果也得由我来担。我不能等到剑落完了，再说一句自己不知道。''

他在栏边坐了下来，这是卢本韦第一次见他坐。

``还有一件事，你该知道。''

卢本韦抬起头。

``猪神 Opt，当年也站到过帝天面前。他的力量比我更强，对那股力量的掌握，也远在我之上。可他最后留下了帝天，还托付给他一件要紧的事。那件事，到今天也没有做完。''

``所以您要跟我讲他？''

test\_alpha0 点了点头。

``你姓卢，背着家里的魂灵笔记，来问我为什么不杀帝天。我就得让你知道，你的先祖当年有能力结束一场仗，为什么还要替人和魂兽，去找另一条路。''

``要讲清这件事，得从一万年前讲起。''

``一万年前，人跟魂兽打过一场。那一场你们的史书上不写。''

``不写的原因是，那一场没有赢家。''

卢本韦愣了一下。

``一万年前，魂兽占上风，可人抗得住。''test\_alpha0 说，``那时候没有魂导器，没有魂师学院，连魂力怎么修都是几百年一点一点摸出来的。可人会结阵，会设伏，肯拿十条命换一头万年的。''

``所以那不是一场仗。那是一千年里没停过的争端。这边屠一个村，那边剥一窝崽，来来回回，谁也吃不掉谁。''

他看了卢本韦一眼。

``你刚才听我讲熊君的儿子。一万年前是这样，一百六十八年前是这样，今天星斗大森林东边那片洼地还是这样。''

``这件事，一万年里一天都没停过。''

``最后那十一年才真正打成了一场。人类一路退到山脉以南。退到那儿，两边都打得动，可两边都看出来了：再打下去，谁也活不成。''

``就在那个当口，出了一个人。''

``封号诛仙\ldots\ldots{}''

``Opt。''卢本韦接了下去。

test\_alpha0 停住了。

``他本来是人。''卢本韦说，``那片场子是他开的，城北一百里，今天还在。族谱第一页就是他，第二页往后空了九千多年，才接上有名有姓的。''

``'猪神'这个叫法是魂兽先叫起来的。它们不知道他姓什么叫什么，只知道那个人是从一片猪场里走出来的。''

``后来两族共举，信仰化神。''

``后山那块青石是他坐过的，场里的规矩是不许搬。''

她顿了一下。

``我姓卢。这些我三岁就会背了。''

test\_alpha0 点了点头。他脸上没有一丝意外。

``我知道你会背。''

卢本韦看着他。

``一百三十六年前，你曾祖跪在这座殿上，也背过这一段。''test\_alpha0 说，``你刚才那几句，跟他当年说的，一个字都不差。''

``我让你先说，是想知道这一百三十六年里，你们家这段话变过没有。''

``变了吗？''

``一个字都没变。''

他站起身。

``这就是我要跟你讲的事。''

``一段话传一万年，一个字都不变，那它就成了供在堂上的牌位。''

``你们家的说法里，一万年前那一仗，是怎么结束的？''

卢本韦愣了一下。

``\ldots\ldots 说是他出面，谈成了和。''

test\_alpha0 没有立刻回答。

``两边罢手，人退到山脉以南，魂兽守着星斗大森林。''卢本韦一字一句地说，像在背一段听了十六年的东西，``从那以后各活各的。''

``我曾祖笔记的第一页就是这么写的。''

``你曾祖那本笔记，是从卢守衡那里传下来的。景川保住它，再交给了承业。至于守衡记下的旧事，也是从上一辈听来的。''test\_alpha0 说，``一万年，一辈一辈传下来，传到最后，剩下的是一句话。''

``哪一句？''

``\,`谈成了和。'\,''

他转过身来。

``卢本韦，你们家那句话没有错。''

``错在它只有四个字。''

\begin{center}\rule{0.5\linewidth}{0.5pt}\end{center}

\chapter{第九十一章：一万年前那一炷香}\label{ux7b2cux4e5dux5341ux4e00ux7ae0ux4e00ux4e07ux5e74ux524dux90a3ux4e00ux70b7ux9999}

风从十二根白玉柱之间穿过去。

卢本韦看着他。

``他一个人上了战场------这一句你们家是有的。他找的是谁，你们家没有。''

``他找的是帝天。''

卢本韦抬起头。

``一万年前的帝天，比今天还强。他是黑龙一族的族长，也是各大魂兽族群都肯认下的统领。他点头，才有人能把停战的约定带回各族，压住那些还想继续打的。''

``Opt 找他，也是因为这个。要把仗停下来，得有一个能代表各族、又能让约定落到实处的人。''

``Opt 站到他面前，展开了 PTO(ZFC+I0)。''

test\_alpha0 念这个名字的时候，语气变了。

卢本韦没有立刻听懂。``那比 I1\ldots\ldots{}''

``你刚才听我讲过 I1，现在要再往上一层。''test\_alpha0 说。

``我从良序之神那里借来的是 I1。那已经是这座殿三百年来能借到的顶。''

``I0 在 I1 上头。''

``高多少？''

``不是高多少的问题。''test\_alpha0 说，``是 I1 到 I0 中间那一段，没有梯子。我借不到，扽魂殿借不到，我这辈子见过的任何一个人都借不到。''

``那他怎么会有？''

``因为他不是借的。''

卢本韦点了点头。``因为他是神。''

``对。''test\_alpha0 说，``这个你们家记着。''

``记着。八个字。两族共举，信仰化神。''

test\_alpha0 看了她很久。

``那这八个字底下是什么，你们家写了吗？''

卢本韦张了张嘴，没有出声。

他走到殿边。

``在站到帝天面前之前，他已经在大陆上走了很多年。''

``哪儿有人跟魂兽打起来，他就去哪儿。他不站边------人杀红了眼，他按住人；魂兽要屠村，他按住魂兽。他的实力足够按住任何一边，所以两边都动不了他，也都恨不起来。''

``走了多久？''

``帝天摊给我的那一段里没有数。''test\_alpha0 说，``只知道很久。久到星斗大森林里的魂兽开始认得他，开始在他走过的路口等他。''

他转过身来看着卢本韦。

``'两族共举'这四个字，写下来一笔就完了。''

``落到地上是什么样子，你想过没有？''

``是他一个人，在两边中间，站了几百年。''

``这边杀红了眼要他让开，那边死了崽子要他偿命。他谁也不帮，谁也不杀，就站在那儿，一年一年地站。''

``站到最后，魂兽那一族认他------不是怕他，是认他。''

``整个魂兽族群一起认下一个人类，这件事一万年里只有过这一回。''

殿上很静。

``认的兽多了，年头久了，就有了香火。有了香火，就有了信仰之力。''

``八个字里的'化神'两个字，是这么来的。''

test\_alpha0 停了一下。

``所以他身上那一层不是借的，也不是修出来的。''

``是两族一起把他举上去的。''

卢本韦很长时间没有出声。

然后她低声说了一句。

``又是几个字。''

test\_alpha0 没有接话。

``你现在再回头看'魂灵'那两个字。''过了一会儿他才说，``他不是坐在屋里空想出来的。''

``他站在中间站了几百年，站出了一个活的样子------人跟魂兽是能共存的，因为他自己就是那个证据。''

``他把那个样子写下来，起名叫魂灵。''

风从柱子之间过去。

``一万年前有过一个活的证明。''

``一万年后，那条路还是没人走通。''

test\_alpha0 停了很久。

``这一百六十八年我想明白过一件事------当年我拿 I3 对着帝天，算不出压下去会塌多少，连敢不敢动手都拿不准。''

``他拿 I0 对着帝天，压得规规矩矩，一根毛都没伤着。那一层是他自己的，他吃得透。''

``差的不只是层次。''

他接着往下说。

``他用它压住了帝天。''

``压了多久？''

``一炷香。''test\_alpha0 说，``不多不少，一炷香。帝天动不了，可也没伤着。他从一开始就没想杀帝天。''

卢本韦怔住了。

``那我们家说的'谈成了和'\ldots\ldots{}''

``是真的。''test\_alpha0 说，``只不过他谈的方式，是先把对面按在地上，按到他连一根指头都抬不起来，然后再开口。''

``那一炷香里他说了什么，我不知道。帝天没有把那一段摊给我看。''

``我只知道他说完以后，做了一件事。''

``什么事？''

test\_alpha0 沉默了一会儿。

``他以神力，塑了一样东西出来。''

``塑？''

``塑。''test\_alpha0 说，``不是找来的，不是抢来的，也不是从谁身上剥下来的。是他一手一手做出来的。''

殿上很静。

``那样东西，叫帝皇瑞兽的本源。''

卢本韦的嘴唇动了动，没有声音。

``那之前也有强大的魂兽，帝天就是其中一个。可让整个族群都有机会往上走，而不只是少数凶兽独自熬过去------这样的路，那时候还没有。''

``是他给的。''

test\_alpha0 的手指在石栏上停了一下。

``他把那样东西交到帝天手上，然后转身走了。人退到山脉以南，魂兽守着星斗大森林。''

``你们家那四个字说的就是这个。''

``然后他做了第三件事。''test\_alpha0 说，``这一件，你们家是知道的。''

``他提出了魂灵。''

卢本韦点了点头。``这个我知道。是老祖提的。我们家那本笔记，第一页就是这两个字。''

``你知道，你曾祖也知道。''test\_alpha0 说，``所以他那年跪在这座殿上，说的不是他自己想出来的东西。他是重提。''

``人跟魂兽不该只有猎杀这一种关系，魂兽的灵魂可以跟魂师共存，不是吞，是合------这些话，Opt 在一万年前说过一遍。''

``说给谁听？''

``说给一个刚被他按在地上、又被他亲手递了一条活路的族群听。''

风从柱间穿过去，很长时间没有别的声音。

``所以帝天肯把本源接下来，不是因为他信人。''test\_alpha0 说，``是因为那样东西是那个人凭空做出来交到他手上的。他这一族往后所有的指望，都在那里头。他不接也得接。''

卢本韦想了一会儿。

``我有一件事想不通。''

``说。''

``他给了魂兽一条往上走的路。''她说，``那为什么一万年下来，走到前面的是人？''

test\_alpha0 看了她很久。

``因为他给的是种子，不是果子。''

``那样东西要长。一万年，它一直在帝天身上长，还没长成。整整一万年，魂兽那边守着一个还没长成的指望。''

``人这一万年干了什么？''

``建学院，造魂导器，把怎么修魂力、怎么养扽魂、怎么带徒弟，一代一代写成本子，写完了再改，改完了再传。''

``一边守着一样东西等它长，一边把自己会的东西一遍一遍教下去。''

``一万年下来，就是今天这个样子。''

殿上安静了一会儿。

test\_alpha0 转过身看着她。

``这件事你记住，卢本韦。不是因为它跟魂兽有关。''

``是因为你们家守那四个字，也守了一万年。''

卢本韦没有说话。

``那魂灵为什么没成？''

``因为他留下的那套东西是残的。''

卢本韦抬起头。

``方向他有了，完整的路子还没有。''test\_alpha0 说，``最初怎么合，他留下过办法。可合完了谁说了算？两边的力量差得太远怎么办？魂兽不肯呢？肯了以后反悔呢？这些，他没有留下完整的答案。''

``他指了方向，也留下了最初的尝试，却没把整条路走通。''

``一万年里，零星有人试过，可始终没补成一套能让旁人照着走的办法。''

他看着卢本韦。

``直到你曾祖。''

殿上静了一下。

``结果你看得见。这一万年，人越走越高，魂兽越掉越低。Opt 留下的魂灵构想，始终没成为两族都能沿着走的路。''

``一万年以后，你曾祖跪在这座殿上，把同样的话又说了一遍。''

他停住了。

``而听的人是我。''

test\_alpha0 的声音第一次有点不稳。

``卢本韦，我当年很清楚自己在做什么。我知道跪在那儿的是谁的后人，我知道那套东西最早是谁说的，我甚至知道你家后山那块石头是谁坐过的。''

``我还是告诉他：走不通。''

``为什么？''

test\_alpha0 沉默了一会儿。

``你还记得我讲过的那场葬礼吗。''

卢本韦点头。

``五十五年以后，那个抱着笔记站在最前头的孩子，卢景川，已经老了。他的儿子卢承业二十二岁，跪在这座殿上跟我说魂灵。''

``手里拿的，还是先祖留下、由他祖父试过的那套东西。关键的几处，一条都没补上。''

test\_alpha0 的声音低下去。

``我那天满脑子都是一百九十一年前那一夜。一个八十七级的魂师，一头八万年的兔子，烧了一整晚。''

``我告诉他走不通，不是因为我不信这条路。''

``是因为那本笔记上写着怎么合，没写怎么解，也没写该先称一称两边有多重。他拿着一模一样的东西，要去走一模一样的路。''

``我怕他死在同一件事上。''

殿上又静了一阵。

``那时候，本源还在帝天身上。他一死，本源就散。''test\_alpha0 说，``散了以后，魂兽不会立刻绝种。它们会一代比一代低地掉下去。三百年出不了一头十万年的，两千年出不了一头万年的。再往后，星斗大森林里剩的就只是野兽，随便哪个魂师都能猎。''

``人类连仗都不用打。等就行了。''

卢本韦站在那里，很久没说话。

``您怎么知道的？''她终于问，``这种事，魂兽不会告诉人。''

test\_alpha0 看着远处。

``那十里。''

卢本韦愣住了。

``我说我知道他在算什么。''test\_alpha0 说，``我没说完。他算完之后，转身之前，给我看了一样东西。''

``我到今天也说不清那是什么方式。不是开口，不是魂力。更像是他把那一万年整个摊在我面前，让我自己看。''

``我看完了。''

``然后他才转身走的。''

他停了一下。

``他赌我看得懂。''

卢本韦的手指慢慢松开了。

``所以他不是忍。''

``他是忍。''test\_alpha0 说，``他也是在赌。那天他要是扑上来，输了，整个族群跟着陪葬。他赢不起。所以他把自己的命摊开给我看，让我来决定要不要动手。''

``那是一头魂兽，跟一个刚杀了他十万同族的人类，做的第一笔交易。''

test\_alpha0 站起来了。

``到此为止的事，我揣了一百六十八年。''

``可他给我看的那段纹路，我一直没看懂。''

卢本韦抬起头。

``我一直以为是帝天没让我看全。''

test\_alpha0 摇了摇头。

``不是他藏了。是我那时候还不懂自己手里的东西。''

``他摊给我的那一段，讲的是一样东西怎么长起来的------什么叫序，什么叫一层一层往上垒，垒到最后为什么能立得住。''

``我这些年推的就是这个。iblp 的良序化，说到底也是这一件事。''

``今天早上刻完最后一行，我闭上眼睛，一百六十八年前那段纹路自己浮上来了。''

``这一回，我才看懂那份本源怎样在里面一点点长起来。可当年的记忆只到当年，没有后来的事。''

``下午你们回殿，我又看见了她。她身上那层气息，跟那段旧纹路接在了一处。''

``到这时候，我才确认------真霜报上来的那个'瑞兽'，不是另一头偶然生出的祥瑞。就是 Opt 留给帝天的那份本源。''

他停了一下。

``帝天守的不是一件器物，是它们一族的祖宗。而祖宗不是死物。''

风把玉简的边页掀了一下，又落回去。

``那本源在帝天身上压了一万年，这一万年它一直在长。''

``一万年，够它自己把那条路重新走一遍。''

卢本韦的声音有点干。``您是说\ldots\ldots{}''

``它化形了。''

``它化形了。''

卢本韦往后退了半步。

``什么时候？''

``十六年前。今天见到她，我才认出化形前后的分界：她身上已经走完的那一段，接着十六年的生长纹路。跟她如今的年纪，正好对得上。''

``原来那段就停了？''

``不是停。是走完了。''test\_alpha0 说，``它离开帝天了。它现在是个人，在这个世上走着，吃饭睡觉，有名有姓------多半连自己是什么都不知道。''

``那帝天呢？''

``他守了一万年的东西，十六年前从他手里走掉了。''test\_alpha0 说，``所以那笔账上现在多了一条。他这一百六十八年不只是在记我。''

殿上安静了一会儿。

test\_alpha0 转过身来，看着卢本韦。他斟酌了一下措辞，像是准备把一样很重的东西轻轻放下。

``卢本韦，有些话我现在告诉你，对你没好处，对她也------''

``她知道。''

test\_alpha0 停住了。

``你说什么？''

``她知道自己是什么。''卢本韦说，``我们也知道。''

卢本韦这辈子没在这张脸上见过那样的表情。不是吃惊。是一个算了三百年的人，忽然发现自己整盘棋里有一步根本没算进去。

``从什么时候？''

``大赛那年。决赛打完那天夜里。''

test\_alpha0 缓缓坐回了栏边。

``决赛最后一场，对面那队的控场一掌拍在队长后心上。''卢本韦说，``队长跪下去，两只手撑在台子上。就那一下，她手背上透出一层金纹，一闪就没了。台底下没人看见。我们八个都看见了。''

``那时候我们已经一块儿练了好几年。一锅里吃饭，一排铺上睡觉，谁夜里翻几回身，大家都知道。她身上多出来一样东西，瞒不住。''

``那天夜里，她自己讲了。''

``她说她不是人。''

``她说她是从星斗大森林最深的地方走出来的。在一样东西里头长了一万年，长够了，才化成人形。''

``她说她走出来那天，心里只装着一件事：魂兽这一族一万年越掉越低，她得替它们找一条往回走的路。''

test\_alpha0 的手指动了一下。

``她说她进学院，进这支队，一开始就是为了这个。她想看明白人是怎么一年一年变强的，看明白了，好回去替魂兽用。''

``她说她瞒了我们这么些年。''

卢本韦停了一下。

``然后她说，还有一样东西，比她的名字还早就在她身上。''

``她说那是一段很长、很旧的东西，不是她的，是当初把她捏出来的那个人留在里头的。那个人的想法，那个人的念头，那个人一辈子没做成的事。''

``最底下压着两个字。''

殿上静了下来。

``她说那两个字是'魂灵'。''

test\_alpha0 看着卢本韦。

``那两个字出来的时候，我手里的碗掉了。''卢本韦说，``那是我们家那本笔记第一页上的字。''

``她说她一开始只想着魂兽，后来在这支队里待久了，那两个字就一天比一天沉。她说她现在不想只替一边找路了。''

``她想要两边都活。''

``她说她一直不敢讲，怕我们把她当魂兽，也怕我们知道她当初是揣着什么进来的。''

``她说她这样的东西，迟早有人要来抢。跟着她走，要死人。''

test\_alpha0 没有出声。

``然后她问我们要不要换个队长。''

风从柱子之间过去。

``奥斯卡当时在啃鸡腿。''卢本韦说，``他把骨头放下，问了一句：`那你明天早上还煮不煮粥？'\,''

``她说煮。''

``他说那就没事了。''

test\_alpha0 的手在膝盖上慢慢收紧。

``小舞抱了她一下。霍雨浩问了她三个问题，都是关于那两个字底下还压着什么，问完低头记进本子里，后来那本子记满了三本。唐三一句话没说，第二天天一亮，照常站到了她身后。''

``诺当把自己那块魂骨给了她。''

``牢大说，往后打架，他站最前头。''

卢本韦的声音停了一下，又接下去。

``宫变那一夜，您没出殿。''卢本韦说，``您把真霜长老叫出来，让他带我们去。临走给了队长一颗炮弹，说以防万一。''

``这些您都知道。''

test\_alpha0 点了一下头。

``战报里有的事，跟我们当时听见、看见的，还不完全一样。''

``宫门口那个太子，是假的。真太子那天夜里就死了。穿着他的衣裳站在宫门口的，是胡峰。胡赞次子胡继的儿子，胡帅的亲侄子。''

``城外头领着三千人堵我们的那个雪星亲王，也是假的。是胡帅。胡赞的长子。''

``Y(1,3,7,12)。比真霜长老高。''

test\_alpha0 的眉头动了一下。

``真霜长老放了冰熊王，跟碎紫将军两个人，在宫门前把胡帅缠了一整夜。''

``我们八个在边上，对的是胡峰。''

``他一开始要的是陛下。''

``在巷子口，他让我们把担架上的人留下，剩下的可以走。他们叔侄那一夜，是冲着天斗皇室去的。''

``打到一半，队长的本源露出来了，那头狼也出来了。他闻出来了。''

``闻出来以后，他喊了一句话。''

``把他的原话再讲一遍。''

``\,`皇帝我先不要了。把瑞兽交出来，剩下的人，还能走。'\,''

殿上静下来。

``我们没有一个人回头看她。''卢本韦说，``那两个字是什么意思，我们早就知道了。''

``他喊的时候，手里已经拎着震天斧了。0 111 222 322 42。''

卢本韦的声音很平。

``我们打不过他。奥斯卡的香肠喂空了，诺当的环让他碎了，我左肩上挨了一斧子。队长把那头狼放出来，堵在巷子口，一步没退。''

``最后是队长把您给的那颗炮弹掏出来，掺上她自己身子里的东西，扔了出去。她扔的时候，对着的是胡帅。胡帅要救他侄儿，只能拿身子去挡。''

``炸完了，胡帅夹着胡峰跑了。跑之前在皇城里起了一个阵，二十几个魂斗罗，三千甲兵，要把我们连陛下一块儿压死。''

``真霜长老说，破不了。''

test\_alpha0 看着她。

``后来呢？''

``后来东边的城墙塌了。''

``我爸骑着那头一百米的猪撞进来的，后头跟着场子里两百来号人。他前头走着一个人，穿一件深蓝的袍子，脚边跟着一头白猪，家猪那么大。''

``那头白猪往阵上一拱，阵就碎了。''

``二十几个魂斗罗，没了十六个。三千甲兵，跑的跑，倒的倒。''

``真霜长老说，那人那时候是九十七级，Y(1,3,10)。比胡帅高，比他自己高了两个大层次。''

test\_alpha0 很久没有出声。

``phyrion。''他说，``信是我递的。我只让他备着。''

``他没光备着。''卢本韦说，``我从小管他叫爷爷，那天夜里才头一回看见他出手。打完了，他端了一碗排骨汤给我。''

``后来的事，您都知道。陛下封了子爵，给了书，我们进了皇家修炼堂。六个月里，队里再没有人提过那一夜的事。''

test\_alpha0 很久没有说话。

后来卢本韦回想，觉得他接下来问的那句，是她这辈子听过最像个老人说的话。

``你们\ldots\ldots 就这样？''

``就这样。''

``没有人退出。''

``没有。''

``没有人躺在那儿睡不着，想着自己身边睡的是一头一万年的魂兽，想着外头还有一整个帅气教要来抢她。''

卢本韦想了想。

``我们睡不着是因为推不出 bms。''

test\_alpha0 笑了。

那是卢本韦第一次看见这位大长老笑。笑得很短，而且不好看，像是很多年没用过那块肌肉了。

``一百六十八年。''他说。

``什么？''

``我花了一百六十八年，想的全是这两个种族能不能坐下来。''

``你们这群人花了一个晚上。''

``还是拿一碗粥解决的。''

他站起来了。

然后他停住。

``\,`把瑞兽交出来'，这句话，真霜回来就报过我。''

卢本韦抬起头。``那您那时候\ldots\ldots{}''

``听说过。此前也有传闻，说你们队长身上带着魂兽的气运。可传闻是一回事，她究竟是什么，是另一回事。''

``您没查？''

``查过。真霜能确认的，是胡家的人抢她，她身边还有一头契约来的狼。他没看见那份本源怎样运转，也不知道 Opt 当年交给帝天的是什么。''

test\_alpha0 的目光落到皇家修炼堂的方向。

``你们在堂里闭关这半年，堂外一直有殿里的人护着，胡家那条线也一直在查。我没有拿那两个字去逼她交代。''

``今天我读懂了旧纹路，又亲眼见到她，才把这两件事对上。''

他转过身来。卢本韦第一次在这张脸上看见那种冷。

``胡家七个月前就已经盯住她了。他们究竟只认出一头瑞兽，还是连这份本源也认出来了，现在得重新查。''

``那一夜胡峰本来不知道。''卢本韦说，``他是冲着皇室去的，打到一半才闻出来。''

``胡峰此前不知道，不等于胡家别的人没盯着。''test\_alpha0 说，``真霜的战报里也记了，你们两年前在星斗大森林外头遭截，动手的跟胡继那条线有牵连。胡家认到了哪一层，不能只看胡峰那一夜说的话。''

``也许，他们查得比我们更早。''

他走到殿边，看着底下那座城。

``帅气教。大陆最西边那座堡，一百六十八年没歇过。''

``你想过没有，他们要一个活的帝皇瑞兽，是要拿她做什么？''

卢本韦的手凉了下去。

``帅气教那套东西，我跟你讲过。杀魂兽，把杀来的往自己身上堆，年份越高的越管用。''test\_alpha0 说，``一头二十万年的熊，够胡赞把一个八十几级的儿子送上封号斗罗。''

``那帝皇瑞兽身上有什么，你想。''

``魂兽这一族往上走的那条路，是从她身上开出来的。她身上压着的不是年份，是命数------一整个族群能长到多高，都系在那上面。''

``命运之力。''

殿上静了下来。

``谁把那样东西抽到自己身上，谁就不用再一头一头地杀了。''

``所以他们不会一刀杀了她。''test\_alpha0 说，``一刀下去，那东西就散了。''

``他们要的是活着抽。''

他停了一下。

``就像一百六十八年前，那头熊站在圈子正中间，一点一点被抽干那样。''

殿上很久没有人说话。

``那我剩下的话就好说了。你们知道她是什么，可有两件事你们不知道。''

``第一，她身上那段'别人的一辈子'，是你祖宗一手塑出来的。她跟你们家的账，一万年前就开始记了。''

卢本韦的呼吸慢了下来。

``第二------''

test\_alpha0 看向官道的方向。

``托付的那个还活着。他守了一万年的东西，十六年前从他手里走了。他还在星斗大森林里，他在等。''

``等什么？''

``我不知道。''test\_alpha0 说，``我只知道，他把那份本源托在身上一万年。我到今天，才把她和当年看见的东西认到一起。他究竟在等什么，还得往下看。''

卢本韦等他把那句话说完，才开口。

``大长老。''

``嗯。''

``我们见过他。''

test\_alpha0 转过身来。

``你说什么？''

``两年前。''卢本韦说，``梅长老带我们进的星斗大森林。''

殿上静了下来。

``梅天狸？''

``嗯。''

test\_alpha0 的脸色变了。

``她带你们进去干什么？''

卢本韦没有立刻回答。

``大长老，这件事得从更早说起。''

``帅气教找上我们，不止一回。''

test\_alpha0 转过头看着她。

``头一回，是全大陆高级魂师大赛。''

``那一届的冠军，是你们。''test\_alpha0 说，``这个我知道。扽魂殿那一队拿了第二。''

``梅长老那天是来给扽魂殿小队庆功的。''卢本韦说，``酒都备下了。''

test\_alpha0 的嘴角动了一下，没有说话。

``决赛那一场她从头看到尾。我们那时候个个比对面低好几级，是一招一招硬磨下来的。她后来跟我们说，那一场她越看越坐不住------不是气自家输了，是没见过那样打的孩子。''

``决赛打完三天以后颁奖。颁完奖那天夜里，我们在回去的路上叫人截了。''

``几个人，一句话不说，上来就奔队长去。''

``我们知道是冲她来的。三天前她才跟我们讲过，迟早有人要来抢。我们八个没商量，就把她围在了中间。牢大站在最前头。''

``梅长老跟在后头。她说她本来是想看看这几个孩子是哪家的。''

``她出手没过三招，那几个人就退了。''

``她认得那几个人的路数，帅气教的。扽魂殿跟帅气教是死仇，她不能看着帅气教的人在她眼皮底下动手。''

``她挡到队长前面的那一下，从队长身上感觉到了一样东西。她回过头，看见我们八个围成一圈，把队长护在当中。''

``她什么都没问。''

``从那一夜起，她就知道队长是什么了，也知道我们都知道。''

test\_alpha0 的手在石栏上停住。

``原来她两年多以前，就已经知道了。''

``两年多。''卢本韦说，``她知道队长带着魂兽本源，也知道我们没有被蒙着。她没问过我们一个字，也没把那一夜的事告诉您。''

``她只是从那以后，开始带我们的队。''

殿上静了很久。

``第二回是开春。''

``我们队原来有九个人。''

殿上安静下来。

``第九个叫牢大。''

``那年开春，我们在星斗大森林外围历练。还是梅长老带的队。''

test\_alpha0 抬起头。

``进星斗大森林第四天，东边出了动静。有人报上来，说邻边一个村子叫魂兽群围了。''

``梅长老赶过去了。''

殿上静了一下。

``她走了不到一个时辰，人就来了。''

``十一个。领头的是胡继手底下一个堂主。''

test\_alpha0 的手指在石栏上停住。

``那个村子\ldots\ldots{}''

``没有魂兽群。''卢本韦说，``事后查过。那三天里那个村子什么都没发生过。''

``报信的人也再没找到。''

``是算好的。''test\_alpha0 说。

``是算好的。''

风从柱子之间过去。

``我们那时候最高的三十几级。对面十一个人，一个魂帝，其余都是魂王。那几个人追得太紧，没人留下拦一拦，我们退不出去。''

``梅长老那时已经在往回赶，可赶到我们这里，还差一刻钟。''

``得有人把这一刻钟撑过去。''

卢本韦停了一下。

``牢大留下了。那条退路窄，他挡住的是追在最前头的几个人。''

``他让我们带着队长往后退，退到星斗大森林边上，等梅长老。他说他顶得住。''

她的声音很平。

``他三十几级。他知道，顶住了，自己也活不下来。''

test\_alpha0 没有出声。

``他把自己烧了。''

殿上静了下来。

``魂力烧完烧魂骨，魂骨烧完烧身子。他就站在那条路上，一层一层往里烧。一刻钟，追在最前头的那几个人没能从他身边过去。''

``梅长老赶回来的时候，那几个强攻过来的人也受了伤，正跟着同伴撤走。''

``那条路上只剩下一圈焦印子。''

``连一块骨头都没留下。''

test\_alpha0 闭上了眼睛。

风从十二根白玉柱之间穿过去。

``这件事跟帝天有什么关系？''

``那天在场的不只有我们。''卢本韦说，``星斗大森林里有一头狼。''

``灭世狼王。''

test\_alpha0 睁开眼睛。

``十九万年。''卢本韦说。

test\_alpha0 的眉头动了一下。``差一万年。''

``差一万年。''

``它过不去。''

``它自己算出来过不去。''卢本韦说，``二十万年那道天劫，它准备了几千年。它跟队长说，它把能想的法子都想完了，最后算明白------扛不住。''

``所以它不修了。它躺在星斗大森林边上，等最后那口气。''

``它看完了全程。''

``它看见十一名人类追杀一群少年，也看见其中几个人被一个三十几级的少年拦在路上。它看见了那个少年为什么不跑。''

殿上静得可怕。

``三天以后，它来找队长。''

``它自己走出来的，走了三十多里，走到我们扎营的地方。走到的时候它已经站不住了。''

``它跟队长说了一句话。''

``说什么？''

卢本韦看着大长老。

``它说：`我把最后这点东西给你。不是给魂兽的，是给那个孩子的。'\,''

test\_alpha0 的手在石栏上停住了。

``然后呢？''

``然后它成了队长的魂灵。''

殿上的风好像停了一下。

``魂灵。''test\_alpha0 把这两个字重复了一遍。

``真霜回来跟我报过。''他说，``说那丫头身边跟着一头狼，十九万年，是契约来的。''

``真霜长老当它是契约。碎紫将军也是。''卢本韦说，``那一夜谁都顾不上细看。他们没往那两个字上想。''

``十九万年的灭世狼王，自愿的。''卢本韦说，``没有称两边的分量，没有'怎么合'，也没有'怎么解'。''

``它走进去的时候，队长哭了。''

``它说：`别哭。一万年了，总得有一个成的。'\,''

test\_alpha0 很长时间没有说话。

``它是怎么知道那句话的？''过了很久他才问，``一万年这三个字，一头狼从哪儿听来的？''

``我们也问过。''卢本韦说，``它没答。''

``它只说了一句：它去过星斗大森林里最深的地方。''

殿上安静下来。

``灭世狼王死之前，在星斗大森林里待过两天。''卢本韦说，``它去见了帝天。''

``它把那件事从头到尾讲给帝天听。讲完了才出来找的队长。''

test\_alpha0 站在那里，一动不动。

``所以两年前我们进星斗大森林的时候，帝天早就知道了。''卢本韦说，``知道队长是谁，知道牢大是怎么死的，也知道灭世狼王最后去了哪儿。''

殿上很久没有声音。

``还有一件事。''卢本韦说。

test\_alpha0 转过头来。

``雨浩身上也有。''

``一只冰蝎子，一个穿白斗篷的姑娘。那一夜他也放出来了，真霜长老一样当成了契约。''

``是伊莱克斯给他的。''卢本韦说，``队长是头一个，雨浩是第二个。''

test\_alpha0 没有说话。

过了很久，他把那三个字低低地念了一遍。

``第二个。''

\begin{center}\rule{0.5\linewidth}{0.5pt}\end{center}

\chapter{第九十二章：总得有一个成的}\label{ux7b2cux4e5dux5341ux4e8cux7ae0ux603bux5f97ux6709ux4e00ux4e2aux6210ux7684}

``接着说。''test\_alpha0 说，``你们走岔了。''

``我们走岔了。''

卢本韦的声音很平，像在复述一件早就放下的事。

``走到第几天我记不清了。星斗大森林越往里越黑，天光从树顶上漏下来，一丝一丝的。后来前面忽然开阔了------一片空地，四周的树全断了。''

``断口在离地二十丈。''

test\_alpha0 的手在石栏上握紧了。

``是那块地方。''

``我们到的时候，他已经站在那儿了。''

``人的样子，很高，一身黑。手里提着一把剑，剑也是黑的，剑身上盘着一条龙，鳞一片一片都刻出来了。''

``我们没想跑。不是不敢，是想不起来跑这件事。他站在那儿，我们就动不了。''

``梅长老挡在最前面。九十七级，八十四岁上的。''

卢本韦停了一下。

``她挡了不到一息。''

风从柱子之间过去。

``可他没动手。他看着梅长老，看了很久。握剑的那只手，往回收了一点。''

``然后他开口了。''

``他说什么？''

卢本韦想了想，一个字一个字地说。

``他说：`请你们快速离开。'\,''

test\_alpha0 站在那里，很长时间没有动。

``就这一句？''

``就这一句。''卢本韦说，``他说的是'请'。''

``说完他转身走了。''

``梅长老带我们出来，路上问过我们一件事------他为什么放我们走。''

``她想不明白。我们当时也想不明白。''

test\_alpha0 慢慢地说：``她回来，没跟我提过这件事。''

``她提了。''卢本韦说。

test\_alpha0 转过头。

``她禀过您。她说队伍进星斗大森林历练走岔了路，遇上一头大凶兽，报了扽魂殿的名头，唬过去了。''

殿上安静得可怕。

``她没提那头凶兽是帝天。''test\_alpha0 低声说，``队里折了一个孩子，我知道。可牢大是为护住谁留下的，那头狼又为什么主动找上你们，她没有讲。''

``她也没把帝天那句话带给我。''

``从大赛那一夜到今天，两年多，她每回来禀报，讲的都是修炼、历练。连真霜报来瑞兽的消息以后，她也没有补上这段。她是长期带队的人，能告诉我的，本该比一份战报多。''

他停了很久。

``她跟了我六十年。九十七级是我一手看着她走上来的。''

``她为什么不说？''

卢本韦看着他，没有马上回答。

``出星斗大森林那天夜里，梅长老坐在火堆边上，一个人坐了很久。''

``后来她跟我们说了一句话。''

``她说：`这件事我不能跟师父讲。'\,''

test\_alpha0 的手在石栏上停住了。

``我们问她为什么。''

``她说，师父这一百多年，听见'魂灵'两个字就摇头。她说她要是把队长是什么、队长身上又多了个什么、雨浩身上又多了两个什么，一样一样禀上去，师父第一件事一定是把队长从队里带走。''

``带到一个他觉得安全的地方，盯着，护着，谁也碰不着。''

殿上静了下来。

``她说那样队长就活得下来了。''

卢本韦停了一下。

``\,`可她就不是队长了。'\,''

test\_alpha0 很长时间没有说话。

``她不是怕您罚她。''卢本韦说。

``她是怕您太想护着我们。''

``他为什么放你们走？''过了很久他才问。

``当时我们不知道。''卢本韦说，``今天知道了。''

她看着大长老。

``我们队里有一样他守了一万年的东西。''

test\_alpha0 闭上了眼睛。

``还有一个。''卢本韦说，``梅长老后来跟我们讲过一个细节。她说那头凶兽听见'扽魂殿'三个字的时候，停了一下。''

``不是听见扽魂殿。''test\_alpha0 说。

殿上很静。

``是听见您。''

test\_alpha0 没有否认。

``他认得你们队长身上那样东西，所以他不能动手。动了，那一万年就白守了。''他说，``可他停那一下，说那个'请'字，是因为梅天狸。''

``梅长老？''

test\_alpha0 睁开眼睛。

``她跟了我六十年。她的功法，是我一笔一笔教出来的，她身上有我的东西。旁人闻不出来，他闻得出来。''

``天州那一回，他那只龙爪叫我劈成了两半，那一剑扎进了他身子。他回到星斗大森林最深的地方，躺了整整五十年，才把剑上那一丝东西化干净。''

``他在梅天狸身上闻见那一点我，头一个想起来的，就是那五十年。''

卢本韦没有出声。

``他不想再打一回。''test\_alpha0 说，``他手里那把剑，叫黑龙剑，是他这一百多年里最趁手的东西。可他比谁都清楚，黑龙剑跟良序圣剑，差得远。''

他停了一下。

``他还清楚另一件事。天州那一剑，用的是 iblp，Laver 表本身只掺进去一丝。光那一丝，他化了五十年。良序之神交到我手上的那样东西，整个的，我一回都没往外拿过。''

``所以他在算另一笔账。''

``一百六十八年前他隔着十里算过我一次，算的是再扑上来能不能赢。两年前他又算了一次，算的不是这个。''

``他算的是：如果他动了你们，我会不会来。''

卢本韦抬起头。

``他算出来的结果是，我会。''

test\_alpha0 的声音很平。

``我拿那样东西对着他，算不出会塌多少，所以一百六十八年前我收了手。''

``可要是他杀了梅天狸，再杀八个十四岁的孩子------''

他停了一下。

``那我就不算了。''

殿上安静下来。

``整片星斗大森林塌不塌，我不会再去算第二遍。他很清楚这一点。''

test\_alpha0 转回身，看着官道的方向。

``他比我自己还清楚我是个什么人。''

卢本韦沉默了很久。

``大长老，还有一件事我一直没想明白。''

``说。''

``他转身之前，多看了一眼。''

test\_alpha0 抬起头。

``不是看队长。''卢本韦说，``我当时以为他在看队长。今天想想，不是。''

``那是看什么？''

``是看我们站的样子。''

殿上静了下来。

``那天我们七个人，站成一个半圆，把队长挡在中间。''

``我们当时知道她是什么，可不知道那头兽是谁，也不知道那块地上死过谁的儿子。''卢本韦说，``我们只知道前面有个东西很大很可怕，我们队长在后头。''

``所以我们就那么站着。''

test\_alpha0 很久没有出声。

然后他转过身，一步一步走回石台边，扶着那卷刻了六个月的玉简站住了。

``他看了多久？''

``不长。几息。''

``然后呢？''

``然后他说了那句话，转身走了。''

test\_alpha0 低头看着玉简。

``卢本韦。''

``在。''

``一万年前，Opt 站在两族中间站了几百年，站出一个活的样子，把那个样子写下来，起名叫魂灵。''

``一万年里，没有第二个人站到过那儿。''

他停了一下。

``两年前的春天，一个三十几级的孩子，为了一个他知道是魂兽、却仍然认作伙伴的队长，一个人留在了星斗大森林里。''

``一头十九万年的狼在不远处看完了全程。它自己再走一步就是天劫，它知道自己过不去，于是走了三十多里路，把最后一点东西给了那个孩子没能护完的人。''

``这两件事，帝天都听说了。是那头狼死之前，亲口讲给他听的。''

``然后过了没多久，你们走岔了路，走到他面前，剩下的七个站成一个半圆，把那个人挡在中间。''

``你们知道的不多，懂的更少，就那么站着。''

``前一件是他听来的。这一件是他亲眼看见的。''

殿上很长时间没有人说话。

``他守了一万年的东西，那一年在他眼前活过来了。''test\_alpha0 说，``不是因为本源化了形------是那件事真的成了。''

``人为魂兽死，魂兽为人舍命。一万年没有成过的东西，他在一个死了的少年和一头快死的狼身上，一次看了两回。''

``所以他收了手。''

``他不是放过你们。''

``他是看见了。''

风从十二根白玉柱之间穿过去。

test\_alpha0 抬起头。

``\,`请你们快速离开。'\,''

``他那样的，一万年里，没跟人类说过一个'请'字。''

``这个字，一半是说给梅天狸身上那一点我听的。天州的账他没忘，可他不想再打了。''

``另一半是说给你们的。他看见了，他愿意再等一等。''

``还有一件，不算我告诉你的，算我提醒你。''

``你们家那位股东，今天早上动身了。''

卢本韦抬起头。``您怎么知道？''

``我不知道。我是猜的。''test\_alpha0 说，``他读过你们家那本笔记，也亲眼见过她在宫变那夜放出的力量。有些地方，他也许比我先认出来。''

``那一夜，我只让他备着。他自己领着那头小白猪进城，先把你们救了下来。他到底看懂了多少，得等你见到他，亲口去问。''

``phyrion 那个人，这辈子出手的次数，一只手数得过来。你曾祖那本笔记他从头读到尾，补进去的那些他也读过。他轻易不动，不是因为不敢。''

``他在等一个条件。一百三十一年，那个条件没出现过。''

``宫变那一回他出手，是为了让你们活下来。那还不是他在等的东西。''

他看着卢本韦。

``现在出现了。''

卢本韦的手攥紧了。

``他说过的。''她说，``陛下当着他的面说的，六个月以后，他领我们去蓬莱岛。''

``他是说过。''test\_alpha0 说，``话还算数。他就是先走了一步。''

卢本韦的呼吸慢了下来。

``他会去哪儿？''

test\_alpha0 看向东南方。

``跟你们一个方向。''

``那这些要不要告诉她？''

``你们自己商量。''test\_alpha0 说，``这半年是你们一块儿熬的，我没资格替你们定。''

``还有一句。到了蓬莱岛，这些事去问最菜萌新。''

``他比我懂。''test\_alpha0 说，``一万年前站在 Opt 边上的，是坤神。那套东西，是他们两个一起留下来的。最菜萌新是坤神的大供奉，坤神那边一代一代传下来的东西，都在他手里。我知道的这些，都是从他那儿听来的。两百多年前，他挑了我一个人讲。''

``为什么挑您？''

``我问过他。''test\_alpha0 说，``他说，因为我是扽魂殿的大长老，而这件事早晚要有一个扽魂殿的人知道。''

他顿了顿。

``你曾祖在那座岛上待过三年。他知道的恐怕比我还多，只是一辈子没跟我提过。''

他往圣殿的方向走了两步，停住。

``卢本韦。''

``在。''

``我方才有一句话说错了。''

``哪一句？''

``我说，天州城外那些命，是在我还没懂的时候抹掉的。''test\_alpha0 说，``我原本以为，我那时候没懂的，只是 iblp。''

他没有再往下说，转身朝殿内走去。

卢本韦站在原地，没有动。

test\_alpha0 在石台前停住了。

``还有一件我没说。''他说，``你曾祖那天问完'您为什么不杀他'，我没答他。我也没告诉他，那时候我已经替帝天瞒了三十二年。''

卢本韦抬起头看他。

``到今天是一百六十八年。''

``嗯。''

``您嘴上说人跟魂兽坐不下来。''卢本韦说，``可您跟帝天，隔着十里，已经坐下来过一次了。''

test\_alpha0 没有反驳。

风从白玉柱之间过去，很长时间没有人再开口。

``那之后，星斗大森林安静了整整五十年。天州城外那片麦田，五十年里没有一头魂兽敢靠近。''

他停了一会儿。

``我原本是要进星斗大森林去谈的。最后是在天州城门口打的。打完了，左肩上那三道口子还在往外渗血。''

``打完了我没有回天斗城。我一路往西走。''

卢本韦抬起头。

``不是去讨说法的。''test\_alpha0 说，``是去拿人的。''

``星斗大森林的核心区是禁猎的。皇室一道令，扽魂殿一道令，压了三百年。他们私自进去，杀了一头二十万年的，连带它那一族两百多头，把尸首拖走炼了东西。''

``按律，那十三个一个都跑不掉，动手的那个该绑到天斗城的刑台上。''

``那为什么没有？''

test\_alpha0 沉默了一会儿。

``因为动手的那个就是胡赞本人。''

``而我得走到帅气堡门口才拿得到他。''

殿上静了一下。

``那地方你没去过，我给你讲讲。''

``最西边那一片，早年是没有人住的。风大，地是盐碱，种什么死什么。他把堡建在那儿，建了一百多年。''

``堡是黑的。不是刷的黑漆，是那块地方的石头本身就是黑的。''

``墙根底下，一圈一圈全是骷髅。''

卢本韦的手指紧了一下。

``不是堆在那儿吓人的。''test\_alpha0 说，``每一具都是炼过的，都站着，都朝着一个方向。风一过，那一片会响。''

``我从堡门外一里地开始数，数到三千就没再数下去。''

``后来我才明白，那不是墙，也不是装饰。''

``那是他的家底。''

殿上很久没有声音。

``他就站在门口等我。''test\_alpha0 说，``他早知道我会去。他在那儿等了三天。''

``他用的是帅气教的禁术。抽人的帅气，填进自己身上。那些年他到底抽过多少人，我查不清。''

``他本人就不弱于我。神力借下来，也得先破开他那一身东西。我从来没有把握说我一定赢他。''

``我要拿他，就得先打赢他。''

``打不赢？''

``换个地方，也得打过才知道。''test\_alpha0 说，``可那是他的堡，他的地，他脚底下还踩着一百多年攒的东西。''

他停了一下。

``我算不出那些骷髅有多少是能动的，算不出他一声令下会起来多少，也算不出那一整片黑石头底下还埋着什么。''

``那天在天州城外，我算不出把 I3 铺到帝天身上会塌多少。站在帅气堡门口，我还是算不出。''

``三百年里我怕过两回。一回是在天州城外，怕我自己手里那样东西。''

``另一回就是那天。''

``我在他门口站了半个时辰。''

``他一直看着我，没说话，也没动。他知道我在算什么。''

``最后我转身走了。''

``走出十几步，他在后面开口了。''

``他说什么？''

test\_alpha0 沉默了一下。

``他说：`大长老，天州那一趟，辛苦了。'\,''

卢本韦猛地抬起头。

``他知道熊君认错了人。''test\_alpha0 说，``他从头到尾都知道。星斗大森林里的账，记到扽魂殿头上，兽潮冲着天斗城去，打成什么样，他都只赚不赔。''

殿上静了很久。

``这一百六十八年，扽魂殿跟帅气教动过四次手。''test\_alpha0 说，``四次都打到一半停了。''

``为什么停？''

``每一次都是同一个原因------谁也吃不掉谁，接着打下去，死的是城里的人。''

他转过身来看着卢本韦。

``你听出来了没有。''

卢本韦想了一会儿。

``一万年前人跟魂兽，也是谁也吃不掉谁。''

``对。''

test\_alpha0 点了点头。

``这世上最难的事从来不是赢。是有人不肯罢手，而你又打不死他。''

卢本韦沉默了一会儿。

``大长老。''

``嗯。''

``\,`你们这种人，我前些日子刚见过一个'，这句话您是什么时候听懂的？''

test\_alpha0 没有立刻回答。

``就在帅气堡门口。''他说，``胡赞站在我面前，我看着他的时候，忽然就懂了。''

``那时候熊君已经死了十一天。''

殿上很长时间没有人说话。

``这一百六十八年，我想过无数回。''test\_alpha0 说，``要是我落到天州城楼上之前，就知道前些日子站在星斗大森林里的是谁，我会不会换一种说法。''

``我会不会先告诉它：那个人不是我，那个人是我要去抓的。''

``它未必肯信。''

``可我连试的机会都没有。''

他停了很久。

``熊君为它儿子死了。我为一座城，杀了十万。''

``动手的那个一天牢都没坐。''

test\_alpha0 停了一下。

``卢本韦，有件事你要记牢。''

``胡赞没有死。''

卢本韦猛地抬起头。

``他今年三百零三岁，比我小十几岁。''test\_alpha0 说，``一百六十八年前他能一个人站着抽干一头二十万年的，这一百六十八年他没闲着。''

``那他现在\ldots\ldots{}''

``他还坐在帅气堡里，教主还是他。这些年他很少露面，外头的事，多半是胡继在跑。''

``胡继今年一百九十七岁，身上背着那二十万年。''

test\_alpha0 看着她。

``宫变那一夜，扮成太子、在巷子口喊'把瑞兽交出来'的胡峰，是胡继的儿子，胡赞的孙子。一年前才升的封号斗罗。''

``帅气教找了你们三回。头一回派的是几个探路的，第二回派的是一个堂主。第三回，是宫变那一夜。胡家的长子跟孙子本来是冲着天斗皇室去的，撞上了你们，胡峰才闻出她来。''

``一回比一回近。''

``可胡赞本人，一回都没有出面。''

``他要是出面了，那一夜就算 phyrion 那头猪撞过去，也未必撞得开。''

``帅气堡今天还在那儿，堡底下那座炉子今天还烧着。''

``天州那一趟之后，我在帅气堡门口站过半个时辰。从堡门往里看，能看见正殿。那时候殿上供的是一块匾，一个'帅'字。''

``前几年，殿里派进去的人带回来一句话：匾撤了，换上去的是一尊像。''

``谁的像？''

``他自己的。''test\_alpha0 说，``堡里上上下下，早晚各一炷香。''

卢本韦想起了那八个字。两族共举，信仰化神。

``Opt 是两族一块儿认的他。''test\_alpha0 说，``胡赞用不着。人这一边，他收香火，也用禁术抽帅气；魂兽那一边，他直接往身上堆。''

他的声音很平。

``我拿它没办法。''

``往后三十二年，我没再提过共存两个字。''

``然后你曾祖来了。他说完'魂灵'两个字，我把星斗大森林那一趟从头到尾讲给他听，跟今天讲给你听的差不多。讲完我告诉他：这条路走不通，你祖父卢守衡已经拿命试过了。''

``那他怎么说？''卢本韦问。

test\_alpha0 的嘴角动了一下。

``他什么也没反驳。只说了一句------`也许不是所有魂兽都愿意和解，但只要有一个肯伸手，这条路就值得走。'\,''

``我当时以为他说的是那头兔子。''

test\_alpha0 摇了摇头。

``过了很多年我才想明白。伸手的那个，不是魂兽。''

``然后他转身去了蓬莱岛。三年。''

``后来他修到了 97 级超级斗罗。可那三年他真正做成的，跟修为没关系。''

test\_alpha0 说到这里停了一下。

``他把他祖父死在上面的那一段补上了。''

``怎么解。''

``怎么解。''test\_alpha0 说，``还有结之前怎么称两边的分量、差多少年以内才敢结、撑不住的时候先松哪一边、出了事算谁的。一百九十一年前那一夜没有答案的，他一条一条写了下来。''

``八万年配八十七级，这种事在他那本笔记里是头一条禁令。''

``写在哪儿？''

``你们家那本笔记里。''

卢本韦的手指在行囊带子上动了一下。

``你背着的那本。''

她没有否认。

``那还差什么？''她问。

test\_alpha0 看了她很久。

``一百九十一年前那一回，他们以为差的是一头肯点头的魂兽。那头兔子点过头了，还是死了两个。''

``你曾祖以为差的是'怎么解'。他把它补上了，可直到你们之前，那些办法还没能立成一套可传开的体系。''

他转过身去。

``我想了很久。我觉得魂灵这条路，卡的地方从一开始就不在'怎么合'。''

``是那两样东西本来就是两样。合得再好，中间那道缝一直在，缝在那儿，早晚要裂。''

他停了一下。

``除非有一个，本来就不用合。''

卢本韦抬起头看他。

test\_alpha0 没有再说下去。

他转回来看着卢本韦。

``他走了以后，我在这座殿上又看了一百三十六年的人间。战争、背叛、仇杀，我见得够多。''

``可我也见过魂师替自己的魂兽挡下致命的一击。见过魂兽在主人死了以后守在坟前，一直守到自己老死。''

``我活了三百多年。''test\_alpha0 说，``前面那一半，我以为自己把这个世界看得够清楚了。后面这一半，我一直在想同一件事------当年那个疯子，可能是对的。''

卢本韦看着他，没有立刻回答。

test\_alpha0 沉默了很久。

``从天州城门口回来那天起，我就不信这条路了。''

他说得很慢。

``不是不想信。是我亲眼看着它塌了两回。一回是你曾祖的祖父，一回是我自己。''

``从那以后，凡是有人来跟我提魂灵，我只有一句话回过去------走不通。''

``这一百六十八年，我用这三个字劝回去过多少人，我自己都记不清了。''

卢本韦抬起头。

``其中有你曾祖。''

殿上静了下来。

``他跪在这儿的那天，我看着他，想起的是他父亲卢景川------那个十三岁抱着笔记站在坟前的孩子。如今，儿子又要去走一条烧死过祖父的路。''

``我听见了他说的话，却没有认真想过，他能不能把缺的那一段补上。''

test\_alpha0 的声音低了下去。

``我这一辈子劝人，从来都觉得自己是在救命。''

``我今天才想，那一百六十八年里，我拦下来的那些人当中，会不会有一个本来是走得通的。''

风从十二根白玉柱之间穿过去，很长时间没有别的声音。

``现在看来------''

他停了很久。

``我可能真的错了。''

殿上安静得能听见风。

过了一会儿，test\_alpha0 才重新开口。

``压力也是动力。''他说完这句话，自己停了一下，``这话对你说好像轻了。''

``你们家走这条路，走了一万年。Opt 说过一遍，你曾祖说过一遍。两回都没能把路走通。''

``到了蓬莱岛，见到最菜萌新，把它说第三遍。''

卢本韦重重点了一下头。``我不会让他们失望的。''

test\_alpha0 的目光越过卢本韦，落到不远处的蔡兴坤身上。黑子鸡站在她肩上，歪着脑袋打量这位大长老，那眼神不太像一只鸡该有的。

``蔡兴坤。''他忽然开口。

蔡兴坤愣了一下，下意识抬手护住肩上的鸡。``我？''

``你那只鸡，还有你们那个篮球。''test\_alpha0 盯着黑子鸡看了几秒，``两百二十多年前，我坐上这个位置没几年，在另一个地方见过同一种气息。当时我以为自己看错了。''

她肩上的鸡忽然不叫了。

``什么地方？''

``蓬莱岛。''test\_alpha0 说，``去了自己看。这种事我说了不算数，得你自己撞上。''

他从袖子里摸出一封信。信封上没写字，封口压着一方小小的朱印。他把信递给卢本韦。

``见了最菜萌新，当面给他。''

卢本韦接过来，揣进了怀里。

卢本韦向他行了一礼，转身去寻队友。脚步声沿着殿阶远去，Smile Lee 才从偏殿走过来，站到师父身旁。

test\_alpha0 望着殿外，过了一会儿，才发现徒弟还有话要说。

``师父，天州那一仗\ldots\ldots''

Smile Lee 停了一下，目光落到他左肩。隔着完好的衣袍，当然看不见伤口，可当年龙神爪留下的那三道疤，他一直记得。

``我当时是不是太冲动了？''

殿外的风吹进来，掀动了石台上一页纸。

``那颗九级炮弹，当年才刚做出来，我一心想着能派上用场。''Smile Lee 说，``看到碧姬在后面治伤，我就想，正好能试一试。后来帝天扑上来，您又受了伤。这些年，我有时还会想------那时候，我是不是太想让这门炮立功了？''

test\_alpha0 转过身，看了他一会儿。

``为守护苍生，你没有错，不必自责。''

Smile Lee 没有马上接话。

``你装弹以前，帝天就已经说过要踏平天州。前头那两座城，也早就破了。''test\_alpha0 说，``碧姬不断替兽群治伤，守军打倒一批，又起来一批。你想截断它，让那座城少挨几轮攻势，这个判断没有错。''

他抬手，按了一下徒弟的肩。

``当年头一回把自己造的九级炮弹打出去，心急是有的。炮该怎么用，反击该怎么防，这些可以回头检点。可别把帝天的怒气，都算成你的过失。''

Smile Lee 垂在身侧的手，慢慢松开了一点。

test\_alpha0 又看向西边。

``如果帝天去了帅气堡，下场只会更惨。''

``真正挑起兽潮的人，却一直在那里。''Smile Lee 低声说。

``是。胡赞到今天，还坐在堡里。''

test\_alpha0 把手从他肩上放下来。

``帅气教已经成为两大帝国最大的威胁了。可帅气堡那片地盘，既不属于帝国，也不属于扽魂殿。''

``一道禁令送过去，得有人能让它生效。他们守着自己的堡，养自己的势力。帝国的军令进不去，总殿派人去，也不能指望拿着令牌就把人带出来。''

Smile Lee 沉默了片刻。

``一百六十八年了。''他说。

``所以这笔账，还没有结。''test\_alpha0 说。

殿下传来几声说话声。卢本韦已经赶上了同伴，八个人正商量着下山的路。Smile Lee 朝那边望了一眼，又收回目光。

test\_alpha0 也听见了。

``天州那一仗，你护住了身后的人。往后的事，照样要把人护住。''

Smile Lee 点了点头。

``是，师父。''

\begin{center}\rule{0.5\linewidth}{0.5pt}\end{center}

黄昏。

天斗城门口，夕阳把八个人的影子拉出去很长，一直铺到城外的官道上。

蔡须坤回头看了一眼天斗城的高墙，然后转过身，把每个人挨个看了一遍。

小舞、奥斯卡、霍雨浩、诺当、蔡兴坤、卢本韦、唐三。

``走吧。''

没有人再犹豫。

他们迈出城门。身后的天斗城在夕阳里一点点矮下去，前面是他们从没踏进过的地方。

风从东边过来，带着海的味道。

\begin{center}\rule{0.5\linewidth}{0.5pt}\end{center}

城墙上。

test\_alpha0 站在垛口后面，看着那八个影子在官道上越走越小。

走在最前面的那个忽然停了一下，回过头，朝城墙这边看了一眼。

隔着三里地，什么也看不清。

但她知道他在那儿。她一直知道。

test\_alpha0 站直了身子，对着那个方向，慢慢拱了拱手。

一百六十八年前，他在天州城外杀了她十万同族。

三里地之外，那个影子停了很久。

然后她转过身，追上了前面的队伍。

一万年前，一个从猪场里走出来的人，把整片星斗大森林的王按在地上，然后亲手做了一样东西交到他手里，说，我们可以一起活。

那份神力，是两族多年积攒的信仰给他的。

可他提出的魂灵构想，没有真正走通。他的后人隔了近一万年又接着往下走，还是没能把它立成一套完整的体系。

那本源在帝天身上压了一万年，自己走完了那条路，落了地，有了名字，今年十六岁，带着七个人往海边走，一路上还在操心谁的行囊捆紧了没有。

她知道自己是什么。

那七个人也知道。

他们照样一起熬了六个月。

天快黑透的时候，test\_alpha0 转身往圣殿走。

一百三十六年前，有个年轻人跪在那座殿上问过他一句话，他没答。

现在他答得出来了。

只是那个人早就不在了，而答案是一群十六岁的孩子替他写完的。

\begin{center}\rule{0.5\linewidth}{0.5pt}\end{center}

很多年以后，魂灵体系真正立起来的那一天，最前面署了三个名字。

后两个，一个姓卢，一个戴眼镜。

头一个，是从一万年前走回来的。

她出星斗大森林的时候，只想替一边找路。

她把名字写上去的时候，写的是两边。

\begin{center}\rule{0.5\linewidth}{0.5pt}\end{center}

\chapter{第九十三章：那封信}\label{ux7b2cux4e5dux5341ux4e09ux7ae0ux90a3ux5c01ux4fe1}

出天斗城的第四天，官道分了岔。

往东的那条宽，铺过碎石，车辙压得很深，一路通到海边。往南的那条窄，两边是齐腰的野草，再往里走就进了山。

唐三在岔口停住了。

这三天他一直走在队长身后，话本来就少，今天一早起来更少，走着走着脚步就慢下来，慢到奥斯卡差点踩上他的脚后跟。

``怎么了？''奥斯卡问。

唐三没答，只是看着往南那条路。

小舞从后面走上来，站到他旁边，也看了一会儿。然后她轻轻''啊''了一声。

``哥，是这儿吧。''

``往南两天。''唐三说。

蔡须坤走过来，没问是哪儿，只问：``要绕多久？''

``来回四天。''

蔡须坤回头看了一眼其他人。没有人说话。诺当已经把行囊换到了另一边肩膀上，换到了朝南的那一侧。

``那就绕。''她说。

\begin{center}\rule{0.5\linewidth}{0.5pt}\end{center}

卢本韦走在最后。

从出城那天起她就一直走在最后。前面七个人的背影，她一个一个都看得见：队长、唐三、小舞、奥斯卡、霍雨浩、诺当，还有肩上蹲着一只黑鸡的蔡兴坤。

出城前在殿上听到的那些话，她一个字都还没说出去。

大长老说，你们自己商量。她想了四天，还没想好这个''商量''该从哪一句开头。从帝天开始，还是从一万年前那一炷香开始，还是从队长身上那样东西是谁一手塑出来的开始。

哪一句都太重。她背着那本笔记，笔记也重。

霍雨浩有一回落到后面来，跟她并排走了一段。他什么也没问，只是走。走了一里多地，他才开口：``你要是想先跟一个人说，可以跟我说。''

卢本韦看了他一眼。

``我没问你是什么事。''霍雨浩推了推眼镜，``我只是看你四天没翻过那本笔记了。''

说完他又走回前面去了。

卢本韦低头看了看怀里的笔记。封皮磨得起了毛，边上用麻线重新锁过一回，是她十岁那年自己锁的。

她把笔记往怀里又按了按，没翻开。

\begin{center}\rule{0.5\linewidth}{0.5pt}\end{center}

往南的路越走越窄，第五天就只剩一条踩出来的土道，贴着山脚绕。

唐三走得比前几天快，走到拐弯的地方就停一停，看看山，看看沟，再往前走。有两回他停下来，指着路边一块石头或者一棵歪脖子树跟小舞说话，小舞就凑过去看，看完摇头，说不记得了。

``你那时候才七岁，又不常出村。''唐三说。

``那你记得？''

``记得一点。''

奥斯卡在后面听着，嘴里嚼着一根草茎，嚼了半天忍不住问：``三哥，你爹以前是干什么的？''

唐三想了想。

``打铁。''

``铁匠啊。''

``我爷爷是铁匠，我太爷爷也是。''唐三说，``到我爹是第三代。村里的锄头、镰刀、犁铧、门环，十户里有九户是我们家打的。''

``收钱吗？''

``收得少。手头紧的人家送点菜，送两个鸡蛋，他也收。''

霍雨浩在旁边插了一句：``你学过吗？''

``学过拉风箱。''唐三说，``他教到我六岁，我就去学院了。''

``那家伙什呢？''

唐三没答上来。他走了好一会儿才说：``他走的时候，锤子一把都没留。''

霍雨浩没再问。他看了卢本韦一眼，卢本韦也正看着他。两个人都没说话。

那天夜里他们歇在山坳里。唐三守前半夜，一个人坐在火边，把什么东西从怀里掏出来看了一会儿，又放回去。卢本韦躺在另一边，眼睛闭着，没睡着。

\begin{center}\rule{0.5\linewidth}{0.5pt}\end{center}

第六天晌午，他们看见了那个村子。

村子不大，窝在两座山中间，几十户人家，屋顶都是灰的。村口有一棵老槐树，树底下一口井，井沿被磨得发亮。

唐三在槐树底下站了很久。

``八年了。''小舞说。

``八年。''

他六岁那年秋天，他爹守着炉子走不开，托了隔壁的老杰克。老杰克一手牵着他，一手牵着五岁的小舞，从这棵树底下走出去，走了一整天的山路，把兄妹俩送到了学院门口。往后两年，每到年关兄妹俩回来一趟。八岁那年开春，他爹走了，他们就再没回来过。

村里有人认出了他。

是个挑水的中年女人。她走到井边放下桶，盯着唐三看了半天，又看看他旁边的小舞，桶也忘了提。

``你是\ldots\ldots 唐家那个小子？''

``是我。''唐三说，``王婶。''

王婶是他娘的亲姐姐。他娘走得早，村里人都叫她王婶，他从小跟着叫，叫惯了，一直没改口。

王婶的眼圈一下子红了。她看了小舞一眼，又看了一眼，说：``跟你娘一个样。''她拉着两个人问了好些话，长高了，瘦了，在外头吃得好不好。唐三一句一句地答，答到后来，他问了一句：

``杰克爷爷呢？''

王婶的话停住了。

她低下头，把桶提起来，又放下。

``去年冬天走的。''她说，``腊月里。那年雪大，他扫完你家院子里的雪回来，第二天就没起来。''

唐三没有动。

``八十七了，走得不遭罪。村里人帮着办的，葬在后山，挨着他老伴。''

王婶停了一下，又说：``这八年，他每个月都去扫一回你家那院子。门上的锁他也上油。我们说人都走了还扫什么，他说那两个孩子会回来的，回来总得有个干净地方落脚。''

小舞捂住了嘴。

唐三站在井边看着那口井，过了很久才说：``谢谢王婶。''

声音很平，跟平时一样。只有站在他旁边的小舞看见，他垂在身侧的那只手一直在抖。

\begin{center}\rule{0.5\linewidth}{0.5pt}\end{center}

老屋在村子最东头。

门上的锁果然是亮的，一点锈都没有。唐三从怀里摸出一把钥匙，那把钥匙他也带了八年。插进去一拧，就开了。

院子干干净净。东边墙角是一座铁匠炉，炉膛冷了八年，里头的灰被人掏过，砧子上盖着一块油布。炉子边上的墙钉着一排铁钉，原先是挂锤子的，现在一把都没有，只剩钉子。墙根下的柴码得整整齐齐，是老人的码法，一层压一层，缝都对着缝。屋里的桌椅还在原来的地方，桌上盖着一块布，布上落着一层薄灰，是从腊月到现在的灰。再往前的，都被人擦掉了。

队伍在院子里停下来。没有人往屋里走，都站在门口，等唐三先进去。

唐三在屋子中间站了很久。

里屋的门框上有一排刻痕，一道一道横着刻，旁边刻着小小的数：一、二、三、四、五、六、七、八。刻到八就没有了。旁边还有一排矮一些的，只刻到七。

小舞走过去，把手按在最上面那道七上。那道刻痕只到她胸口。

``我那时候这么点儿。''她说。

唐三站在她身后，看着那道八。他伸手摸了一下，刀口很齐，一刀下去，没有补第二刀。

然后他蹲下来，从床底下拖出一只小木箱。箱子里是他小时候的东西：一把木头削的小刀，几块磨圆了的石子，一双早就穿不下的布鞋。布鞋底子补过好几层，针脚很粗。

``杰克爷爷补的。''小舞在他身后说。

``嗯。''

\begin{center}\rule{0.5\linewidth}{0.5pt}\end{center}

后山的坟很新，土还是黄的。碑是一块青石，字是村里的教书先生写的，写得有点歪：杰克之墓。

唐三在坟前跪了下来。

小舞跪在他旁边。其他人站在后面，没有人说话。奥斯卡从行囊里摸出三根香，是出城那天在天斗城买的，他当时也说不上来买它做什么。他把香点着，递给唐三。

唐三接过去，插在坟前。风从山上下来，香头的火一明一暗。

``六岁那年，他送我们去学院。''唐三忽然开口，``走了一整天。''

``半路上小舞走不动了，他背她。后来我也走不动了，他让小舞下来走一段，又背我。他那时候就七十多了，两个换着背，背了二十里山路，一句累都没说。到了学院门口，他从怀里掏出两个馍，一人一个，还是热的。他一路揣在怀里捂着。''

``八岁那年开春，我爹走了。过了半个月，他又送我们回学院，还是那条路。到了门口，他说，三儿，好好念书。你爹要是回来了，我让他来看你。''

唐三停了一下。

``我爹没回来。''

``他也没再来过。学院离这儿太远，他腿脚不好。''

``我想着等我出息了，回来看他。''

他没有再往下说。

小舞趴在坟上哭出了声。唐三伸手一下一下拍着她的背，自己一滴眼泪都没掉。

蔡须坤走过去，在唐三另一边蹲下来。她什么也没说，只是把手放在他肩膀上。

霍雨浩一直站在后面，没往前凑。等小舞哭得喘不上气了，他才走过去，蹲在她另一边，把水囊拧开递过去。小舞接过去喝了一口，呛着了，咳了好几下。他伸手把粘在她脸上的头发拨到耳朵后面，又把水囊接回来，拧上盖子。

唐三把拍着小舞后背的那只手收了回来，往旁边让了让。

当年在学院，头一个跟小舞说''霍雨浩那人话少，心细，你多看两眼''的，就是他。

那天下午，他们在坟前一直待到太阳落山。

\begin{center}\rule{0.5\linewidth}{0.5pt}\end{center}

天擦黑的时候，有人来敲院门。

是个瘦高的老头，穿一件洗得发白的长衫，手里提着一盏灯笼。他在门口站定了，先把灯笼举高照了照唐三的脸，才放下来。

``是长大了。''他说，``眉眼没怎么变。''

``先生。''唐三说。

是村里的教书先生，碑上的字就是他写的。他进了院子，也不坐，从袖子里摸出一个折成方块的纸包，递给唐三。

``杰克托我写的。''他说，``他不识字，说一句我写一句。写完了他让我念了三遍，改了两个地方。那是前年开春的事，他那时候腿就不大好了。''

``他说你要是回来了，就交给你。你要是不回来，就跟他一块儿埋了。''

唐三接过纸包，拆开。纸是先生平时写字用的毛边纸，字是先生的字，一笔一画很端正。

\begin{quote}
三儿：

爷爷不识字，这是托先生写的。

你爹走那天夜里来找过我。他留了一袋钱，让我接着供你们兄妹念书。我问他上哪去，他不说。他就说了一句，别让孩子等他。

我没听他的。我跟你说他会回来。你别怪爷爷。爷爷是怕你一个人在外头，心里没个盼头。

前年秋天，有个外乡人来村里打听你们家，问你爹，也问你们兄妹俩，问得很细。我说不认识。那人住了一宿就走了。你在外头要当心。

钱没用完，剩下的埋在院子里枣树底下，往东两步。

杰克
\end{quote}

唐三看完了，又从头看了一遍。

先生在旁边说：``那个外乡人我也见过。三十来岁，说话客气，不是来买东西的，也不像走亲戚的。他在村里转了一天，最后在你家门口站了很久。第二天一早就走了。''

``往哪边走的？''

``往北。''

唐三把信折起来，向先生深深鞠了一躬。

先生摆摆手，提着灯笼走了。走到院门口又回过头：``杰克走之前那几天，嘴里念叨的都是你们俩。他说三儿那孩子心重，什么事都往肚子里咽。''

院门关上了。

唐三在院子里站了一会儿，走到墙角那棵枣树底下，往东量了两步，蹲下来用手刨。土不深，刨了没几下就碰到一个陶罐。罐口用油纸封着，揭开，里面是一个布袋。

布袋是旧的，颜色已经分不清了。他解开来，里面是碎银子，不多，底下垫着一张纸。纸上什么都没写，只按着一个手印，是杰克的。

``这是他每一回从里头拿钱，就在账上按一个？''奥斯卡小声问。

``不是。''小舞说，她的声音哑了，``杰克爷爷不会记账。这是告诉哥，这罐子他动过，他没藏着。''

唐三把布袋系好，放进怀里，跟那封信搁在一处。

\begin{center}\rule{0.5\linewidth}{0.5pt}\end{center}

夜里，他们歇在老屋。

屋子小，八个人挤一挤也就挤下了。奥斯卡和诺当在院子里生了火，蔡须坤煮了一锅粥。粥熬得很稠，她往唐三碗里多舀了一勺，唐三没推。

吃完了，大家各自找地方躺下。蔡兴坤的黑子鸡蹲在房梁上，一只眼睛闭着，一只眼睛睁着。

卢本韦坐在门槛上，把那本笔记摊在膝盖上，半天没翻过一页。

身后有脚步声。唐三在她旁边坐下了。

两个人都没说话。院子里的火快灭了，只剩一层红炭。

过了很久，唐三从怀里摸出一样东西。

一封信。

信封是黄的，边角磨得发白。折痕的地方已经快断了，被人用细线从里头缝过，缝得很仔细。

``我爹留的。''唐三说，``他走之前那天夜里，放在我枕头底下。''

``我带了八年。''

他把信递过来。卢本韦没有马上接。

``你确定给我看？''

``嗯。''唐三说，``就是想让你看。''

她接过来，借着炭火的光展开。

纸很薄，字不多。

\begin{quote}
三儿：

爹要出一趟远门，走得急，来不及跟你说。

不用找我。

照看好妹妹，听杰克爷爷的话。你八岁了，是大孩子了。书接着念，杰克爷爷会送你回学院。

等你长大了，要是还想知道爹去了哪儿，
\end{quote}

后面还有一行，被划掉了。划得很用力，一道压着一道，墨都透到了背面，一个字也认不出来。

信的末尾没有落款，只有一个印。

印是圆的，比指甲盖大一点，朱红色。里头一道弯弯的纹，从左边起笔，绕了大半圈，到右边没有合上，留着一个口子。

卢本韦盯着那个印看了很久。

``你娘呢？''她问。

``信里一个字都没提。''唐三说，``我记事的时候，她就已经不在了。王婶是她亲姐姐，跟我说过，她叫王冬，进城送货，走夜路摔进了山沟。我只记得一个声音，哼过一首歌。调子我还记得，词记不住了。''

``小舞呢？''

``她比我还小，什么都不记得。''

卢本韦想了想，又问：``你爹的扽魂是什么？''

``不知道。''唐三说，``我从没见他放出来过。村里人也没见过。''

``你的呢，是随他？''

``说不准。''唐三先摊开左手，掌心里钻出一根蓝银草，细细的，绿莹莹的，绕着他手指头转了一圈，``这个你们都见过。''

他又摊开右手，掌心里浮起一把锤子的虚影，黑沉沉的，没什么光，一会儿又散了。

``昊天锤。''他说，``双生的。觉醒那天两样一块儿出来。测的人拿草看了一眼，说是废扽魂，又看了看锤子，说二流，往上练也就那样。我爹也在。他一句话没说，回家在炉子边上坐了一晚上。倒是杰克爷爷高兴得不行，说二流也是锤，比镰刀强。''

卢本韦看着他两只手。

``宫变那一夜，你一锤都没抡。''

``用不上。''唐三说，``那一夜要拴人，要把八个人的魂力往一处送，锤子干不了这个，得用草。''

他把两只手都收了回去。

``你是神级，队长和兴坤也是，牢大也是。''他说，``我拿一把二流的锤子，跟着你们走到今天。有时候我也想，我爹要是在，会不会知道这把锤子还能怎么用。''

卢本韦没有马上接话。她想了想，才说：``你老拿我们几个比。''

``队里是我和小舞垫底。''

``队里是。''卢本韦说，``大赛那年你十四，对面那队主攻十七，魂力比你高两级，你三锤把他砸下了台。小舞那场，对面三个人没一个近得了她的身。整个天斗城跟你们同岁的，我没见过第二对。''

唐三笑了一下，很短。

``宫变那一夜，你拿草把我们拴在一块儿，拴到最后一根草都枯了，手也没松。''卢本韦说，``那个活儿，换了谁都干不来。''

唐三没接这句话。

炭火噼啪响了一声。

``这八年，''唐三说，``我没跟任何人提过要找他们。''

``为什么？''

唐三想了想。

``他叫我不用找。''他说，``我八岁的时候，以为不用找的意思就是他会回来。''

``后来知道不是了。可我一直觉得，我要是去找了，就是没听他的话。''

卢本韦没有出声。

``今天在坟前，我想，杰克爷爷等了八年，等到死也没等到他。''唐三说，``我要是再不找，就轮到我了。''

他转过头看着卢本韦。

``卢本韦，你能不能帮我查一查？''

``查我爹，还有我娘。''

卢本韦低头看着那封信。

``为什么来找我？''她问，``队长、雨浩，都比我会查东西。''

唐三指了指信末尾那个印。

``你们家养猪场。''他说，``你们家的猪卖了一万年，卖遍了整个大陆。各地的商帮，没有一家不跟你们家打交道。''

``商帮认印。''

卢本韦怔了一下，然后慢慢点了点头。

``每个地方的商帮都有一本印谱。''她说，``哪家商号在当地做过买卖，立过契，押过货，印都要拓一份留在帮里，免得日后赖账。一本印谱用几十年，旧的封存，新的接着拓。''

``卢家的买卖走到哪儿，哪儿的商帮就得给卢家几分面子。我要翻他们的印谱，没有哪一家会说不。''

``只要这个印在哪儿正经盖过一回，那儿的印谱里多半就有底。''

``要是没有呢？''

卢本韦看着他。

``要是没有，就说明盖这个印的人，从来没跟正经买卖打过交道。''她说，``那也是一条线。''

唐三点了点头。

卢本韦又看了一会儿那个印，把它的样子记下了：圆圈，弯纹，右边那个没合上的口子。

``明天让雨浩照着描一份，他手稳。''她说，``岔口往东有个三河镇，卢家往这一带收猪，走的就是那儿的商帮。我拿着拓样去找他们的管事。本地的印谱里要是没有，就让他们往周边几个州的商帮传。商帮之间有自己传信的路子，半个月就能有回音。''

``谢谢。''唐三说。

``不用。''卢本韦把信折好，顺着原来的折痕一道一道折回去，递还给他，``你收好。原件别离身。''

唐三接过信，放回怀里，贴着胸口那一侧。

他没有马上走。

``卢本韦。''

``嗯。''

``你这几天也有事。''唐三说，``你一直走在最后。''

卢本韦没有说话。

``我不问。''唐三说，``你想说的时候，我在。''

他站起来，回屋里去了。

\begin{center}\rule{0.5\linewidth}{0.5pt}\end{center}

卢本韦一个人在门槛上又坐了很久。

炭火灭了，院子里只剩月光。屋里传来很轻的鼾声，是诺当。房梁上的黑子鸡换了一只眼睛睁着。

她想着那个印，想着那道没合上的口子。然后她把笔记翻开，从后往前，一页一页地翻，翻过曾祖补上去的那几十页。

翻到其中一页，她的手停了一下。

那一页上只有一句话，是曾祖从蓬莱岛回来以后写的。她小时候读不懂，后来读懂了，就不怎么敢翻那一页。那句话说：查一件事，先查它为什么被人埋起来。

她接着往前翻，翻过更早的、字迹已经褪成淡褐色的那些页，一直翻到第一页。第一页上还是那两个字。

她看了一会儿，把笔记合上了。

她想起那道被划掉的一行，想起那个往北走了的外乡人，想起一个打了半辈子铁的男人，在一个春天的夜里把墙上的锤子全摘走了，只留下一封不让人找的信。

这支队伍里，有人把一封信揣了八年，一个字没跟人提；有人把自己是什么揣了好些年，直到被人在门外喊破才说出来。现在她自己也揣着一样东西了。

唐三今晚把信拿出来，是因为杰克爷爷等到死也没等到。

那她呢，她要等到什么时候。

屋里有人翻了个身。她抬起头，看见队长坐了起来，隔着门槛往外看了她一眼，没说话，又躺了回去。

卢本韦把笔记抱在怀里。

``明天。''她对自己说，``明天路上，先跟队长说。''

\begin{center}\rule{0.5\linewidth}{0.5pt}\end{center}

第二天一早，蔡须坤又熬了一锅粥。

八个人围着院子里那张矮桌坐下，桌子小，碗挨着碗。小舞挨着霍雨浩坐，把自己碗里的咸菜拨了一半到他碗里。唐三喝了两口，把碗放下了。

``有件事，我跟大家说一下。''

没有人再动筷子。

唐三把那封信从怀里拿出来，摊在桌子中间，又把杰克那张毛边纸摊在旁边。

``我爹留的，我带了八年。昨天夜里给卢本韦看过了，我请她帮我查。''他说，``今天也给你们看。以后路上要是碰见什么跟这个有关的，你们也替我留个心。''

奥斯卡第一个凑过去，看完了，挠了挠头：``\,`不用找我'\ldots\ldots 三哥，你爹这话说得也太干了。''

``他说话一直这样。''

诺当看了很久，一个字没说，把信往唐三那边推回去一点，推得很轻，怕碰坏了折痕。

蔡兴坤肩上的黑子鸡伸着脖子去啄那个印，被她一巴掌按了回去。

霍雨浩没有看前面，他一直盯着最后那道被划掉的地方。他把信拿起来，对着太阳举高，眯着眼看了半天。

``七个字。''他说。

``什么？''

``划掉的这一行，底下原本是七个字。''霍雨浩说，``划的时候笔是横着走的，碰到字的地方墨会堆一下，堆出来的疙瘩数得出来。第三个字和第六个字的笔画多，划到那里墨透得最深。''

唐三看着他。

``字是什么我认不出来。''霍雨浩把信放下，``可你爹写完这七个字，想了很久才划掉。前面几行墨是一口气下来的，这一行起笔之前，笔尖在纸上停过，洇了一个点。''

院子里静了一会儿。

``他本来是想告诉你的。''小舞小声说。

唐三没接话。他把信拿回来，一道一道折好。

霍雨浩把眼镜摘下来，拿衣角擦。镜片本来不脏，他擦了很久。

小舞先放下了筷子。

``三哥，你至少还有封信。''霍雨浩说。

桌上的人都看向他。

``我七岁那年，家里夜里来了几个人。穿黑衣裳，袖口绣一道金边。''霍雨浩说，``我爹去开的门。他们没进屋，就站在门口，念了一个名字，又念了一个。是我爹和我娘的。''

``我娘把我塞进灶台底下，拿柴火挡上。我从柴缝里往外看，看见他们把我爹我娘带走了。我爹出门的时候回了一下头，看的是灶台。''

``第三天，有人在镇子外头的河滩上找到了他们。''

他把眼镜戴回去。

``身上一道伤都没有。镇上的人都说我娘是镇上最好看的，我爹走在街上，卖菜的都要多看他两眼。河滩上那两个人，我站在跟前，认了半天才认出来。眉眼一处没少，脸还是那张脸，可就是什么都没了。像一盏灯，灯芯还在，火让人掐了。''

``后来我才知道，这叫抽帅气。帅气教早些年只猎魂兽，后来魂兽不好猎了，就开始抓人。魂兽身上抽魂力，人身上抽帅气，抽干了就扔。''

奥斯卡手里的筷子放下了。

``宫变那一夜，胡峰在巷子口念'果糕小队'那四个字，念得慢慢的，像在嘴里嚼。''霍雨浩说，``我一听那个腔调，手就攥上了。跟那天夜里在我家门口念名字的，是一个腔调。''

``我那三本笔记，前面记的都是魂灵。第三本后半本，我另起了几页，记的是帅气教。他们在哪儿出没，用什么法子，抽一个人要多久。''

诺当抬起头看着他。

``我没跟你们说过。''霍雨浩说，``我怕说了，你们往后看我的时候，就只看得见这一件事。''

桌子底下，小舞的手伸过去，把他攥着的那只手掰开，握住了。

奥斯卡看了看她。

小舞没躲他的眼神。``他跟我说过。''她说，``去年冬天，在学院后山。他说完了，我们俩在那儿坐到天黑。''

``那现在呢？''唐三问。

``现在你把信拿出来了。''霍雨浩说。

他推了推眼镜，低头喝粥。

喝完了，他从行囊里摸出那第三本笔记，翻过后半本记帅气教的那几页，找了一张空页。他把唐三的信要过去压平，照着信末那个印一笔一笔描下来，描完了，在旁边写：

唐三之父。八年前春离家，不知所踪。在村中以打铁为业，离家时炉边锤具尽数带走。扽魂从未示人。母，不详。

信末一行被划去，原七字，第三、第六字笔画多。

前年秋，有外乡人来村打听唐家，三十来岁，去向北。

他写字很快，字又小又密。写完了，他把笔记转过去给唐三看了一眼，才合上。

蔡须坤一直没开口。等霍雨浩把笔记收起来，她才说：``到了蓬莱岛，这两件事都算上。''

``队长------''

``你昨天夜里说，她想说的时候你在。''蔡须坤看了卢本韦一眼，又看回唐三，``这话不光你会说。''

唐三低下头，把碗里剩下的粥喝完了。

卢本韦坐在桌子另一头，手放在膝盖上的笔记上，没抬头。她听见队长那句话里转过来的那一下，知道那不全是说给唐三听的。

\begin{center}\rule{0.5\linewidth}{0.5pt}\end{center}

动身的时候，唐三把门锁上了。

他锁得很仔细，锁完又拉了两下，确认拉不开。然后他从行囊里摸出一小瓶油，也不知道是什么时候去村里要的，往锁眼里滴了两滴。

小舞站在旁边看着，眼泪又下来了。

霍雨浩从后面走上来，没说话，把她肩上的行囊卸下来，挎到自己肩上。小舞抬手抹了一把脸，拽住了他的袖子。

``走吧。''唐三说。

他们从村口那棵老槐树底下走过去。王婶站在井边，冲他们挥手。

唐三回头看了一眼。

那间屋子在村子最东头，门关着，锁是亮的。

然后他转过身，跟上了队伍。

\begin{center}\rule{0.5\linewidth}{0.5pt}\end{center}

往北走了两天，回到那个岔口。

岔口往东二里有个镇子，叫三河镇，是这一带山里人赶集的地方。镇子不大，一条正街，正街当中让出一片青石铺的空场。空场北面是扽魂殿的分殿，三进的院子，青砖黛瓦，门前立着两根旗杆，门楣上的徽记擦得锃亮。这附近几十个村子，孩子六岁觉醒扽魂要到这儿来；不觉醒的日子，落户、分田、立契、断官司，也都在这儿办。

街口还有一座商帮会馆，门脸也气派，挂着一排灯笼。

卢本韦要去商帮会馆，队伍就在镇上歇了半天。

唐三没跟着去会馆。他在分殿门口站了一会儿，走了进去。

分殿的前厅很宽，一排柜台后面坐着四五个执事。柜台前排着队，有来落户的，有来立田契的，还有两家为了一道田埂吵到这儿来，一个执事正拿着地图给他们比划。

前厅东边隔着一道屏风，屏风后头是觉醒扽魂的地方。唐三六岁那年就是在那儿觉醒的，他爹站在屏风外头等。

唐三排了一会儿，轮到最边上那个柜台。柜台后面的执事五十来岁，穿着分殿的青袍，袖口绣着徽记。他抬头看见唐三胸口别着的魂师牌子，搁下了手里的笔。

``登记？''

``查一个人的登记。''唐三报了村子的名字，``姓唐，打铁的。我是他儿子。''

执事看了看他的牌子，又看了看他的脸，没多问，起身往后头的档房去了。去了一盏茶的工夫，出来的时候抱着一本厚册子，蓝布封皮，四角包着铜，端端正正放到柜台上。

``魂师的名册单立一本。山里的魂师少，一本能用二十年。''他一边翻一边说，``你们那个村子，这些年统共就登过两三个魂师，好找。''

他翻到一页，停住了，把册子掉了个头推过来，手指点着其中一行。

唐三低头看。

行首写着一个名字：唐舞麟。

唐三盯着那三个字看了很久。村里人叫他爹，叫的是''唐铁匠''，杰克爷爷叫他''舞麟''，他一直以为是''五林''。这三个字怎么写，他今天头一回看见。

舞。

他想起小舞的名字。小时候他问过爹，妹妹为什么叫小舞。爹正在炉子边上拉风箱，没回头，只说了一句，随我。

他那时候不懂，只当是随口一句。

那一行底下连着好几行，是同一个名字，一年一行，每行后面写着年份、魂力和经手人的签押。最早那一行是十八年前，魂力四十一级。往下一年一行，四十一，四十二，四十二，四十三。再往下几行，一直是四十三。

最后一行是八年前的春天。四十三级。

再往下，就没有了。下一行是别的村子、别的名字。

唐三的手指按在那行字上，很久没动。

``这人后来没来过？''他问。

``册子上没有，那就是没来过。''执事说，``按规矩，魂师一年来登一回，领当年的补贴。三年不登，册子上就不再往下续了，也不销，就这么挂着。''

他往那几行字后面瞅了一眼，``咦''了一声。

``补贴他一回都没领过。''执事指了指每行末尾空着的那一格，``四十多级的魂宗，一个月的魂币够山里一户人家吃半年。年年来登，年年不领。''

``那他来登什么？''唐三问。

执事想了想，摇头：``这我就说不上来了。按说不登也不犯什么法。''

唐三把那几行字从头到尾又看了一遍。

他去学院以前，家里吃的是自家种的菜，打铁换来的鸡蛋，冬天不够了，就去杰克爷爷家借半袋粮。他爹每年开春会出一趟远门，推着独轮车，说是去赶集，一走七八天，回来的时候车是空的，手里什么都不提。

他那时候以为，赶集就是推着车出去转一圈。

``能抄一份吗？''他问。

执事从柜台底下摸出一张纸，一支笔。唐三把那几行字一笔一笔抄下来，年份、魂力，一格一格对齐了。抄完他又核了一遍，把纸折好，揣进怀里，跟那封信放在一起。

``谢谢。''

``没什么。''执事把册子合上，抱起来往档房走，走到门口又回过头，``小伙子，你爹是四十三级的魂宗，在这山里头，算是顶了天的人物了。''

``他在村里打了一辈子锄头。''唐三说。

执事愣了一下，没再说话，进了档房。

\begin{center}\rule{0.5\linewidth}{0.5pt}\end{center}

\chapter{第九十四章：没有词的歌}\label{ux7b2cux4e5dux5341ux56dbux7ae0ux6ca1ux6709ux8bcdux7684ux6b4c}

唐三从分殿出来的时候，卢本韦正站在街口的商帮会馆门口。会馆的管事是个穿绸衫的胖子，手里捧着霍雨浩早上描的那张拓样，陪在她旁边，一脸的笑。会馆门口搭着一个茶摊，其他人都坐在那儿歇脚。

本地的印谱，管事已经陪她翻过了。三十年的拓样，一页一页对下来，没有这个印。

她看见唐三的脸色，抬手让管事等一等。

唐三走过去，把那张纸递给她。卢本韦看完，转手递给了茶摊上的霍雨浩。

霍雨浩从头看到尾，摸出第三本笔记，翻到早上那一页，在底下接着写：

父名唐舞麟，魂宗。三河镇扽魂分殿有登记，十八年前起，年年登，魂力四十一至四十三级，补贴分文未领。末次登记八年前春，此后不再登。

写完，他把登记那几行又誊了一份，递回给卢本韦。卢本韦把它跟拓样叠在一起，从袖子里摸出一枚指节大的小印，蘸了印泥，在两张纸的骑缝上按了一下。

``拓样往周边几个州传。''她对管事说，``名字也一块儿查。唐舞麟，这十几年在哪家商号买过铁料，卖过铁器，都报上来。''

管事接过去，低头看见骑缝上那个印，腰一下子弯了下去，连声应着，双手捧着进了会馆。

``四十三级。''她说，``一个魂宗，在山里藏了十年，每年下山登一回记，一个魂币不拿。''

``他是想让册子上有他。''唐三说。

卢本韦抬头看了他一眼。

``册子上有他，就说明他那年还活着，还在这儿。''唐三说，``他不领钱，是不想让人知道他缺什么。他来登记，是想让某个会来翻这本册子的人知道，他还在。''

``那个人是谁？''

``不知道。''唐三说，``可八年前春天，他不登了。''

他看着街上，过了一会儿，又开口。

``上回见你爹，他四十八岁，就快摸到封号斗罗了。''

卢本韦没否认。``还差一点。''

``你干爷爷九十八级，你曾祖九十七级。''唐三说，``四十三级搁在你们家，说出去都没人当回事。''

``我们家也就是熬得久。''卢本韦说，``要说天才，大长老殿下才是。二十三岁封号，三十七岁就是极限斗罗了。我爹四十八岁还差一点，在他跟前也不够看。''

茶摊那边忽然没了声音。

奥斯卡端着碗，碗沿停在嘴边。诺当剥了一半的花生掉回了盘子里。蔡兴坤肩上的黑子鸡也不动了，歪着脑袋。

``多少？''奥斯卡问。

``二十三封号，三十七极限。''卢本韦说。

奥斯卡把碗放下了。``闭关那天他跟我们说，二十五岁之前都得给他到封号斗罗，到了就封长老。我当时还琢磨，这老头是不是没算过数。''

``殿里的长老，都是九十五级往上才封的。''诺当说。

他难得开口，一桌人都看了他一眼。

``他算过。''霍雨浩推了推眼镜，``他是照着自己算的。''

``那他还多给了我们两年。''小舞小声说。

``是留给我们摔跟头的。''蔡须坤说。

没有人笑。茶摊老板提着壶过来添水，看这一桌人都不说话，添完就走了。

唐三站在会馆门口，一直没吭声。二十三岁，离他现在不过七年。

``再说，你爹用不着跟谁比。''卢本韦说。

``我没跟谁比。''唐三说，``我就是想不通。''

``一个四十三级的魂宗，在大陆上算不上什么。没人会为了他那点魂力追到山里来，他身上也没什么值得别人抢的。''

``算不上什么，也用不着打铁。''卢本韦说，``四十三级的魂宗，进天斗城找一家贵族府上做供奉，一年的俸禄够在城里置一座宅子。不愿意给人看门，就去学院教低年级，也是个体面的先生。''

``他哪条路都走得通，偏偏守着山里一座炉子，给人打锄头。''

唐三低头看着脚下的石板。``我小时候以为，我们家祖祖辈辈都是打铁的，他也就只会打铁。''

``他打铁，种地，一年下一回山，连补贴都不领。''他说，``放着城里的日子不过，在山里一窝十年，最后又非走不可。为什么？''

卢本韦没答。

``走的时候为什么要把锤子全带走，为什么不让我找，那一行写好了为什么又划掉。''唐三说，``还有那个外乡人，前年秋天来找的是他，还是我们。''

``八年了，我一件都没想通。''

他说完，把手伸进怀里，隔着衣裳按了按那封信，又把手拿出来。

街上有人赶着驴车过去，车轱辘压着石板，一路响过去。

卢本韦回头看了一眼会馆。

``拓样明天一早就往外传。''她说，``半个月。''

唐三点点头。

那天下午，八个人从三河镇出来，上了往东的官道。路宽了，碎石铺得平平整整，车辙一道压着一道，一直通到海边。

\begin{center}\rule{0.5\linewidth}{0.5pt}\end{center}

从三河镇往东走到第五天，商帮的人追上来了。

是个十七八岁的小伙计，骑一匹枣红马，马跑得浑身是汗。他在官道上把队伍截住，下了马，先认人，认准了卢本韦，才从怀里掏出一个油纸包，双手递过去。

``卢小姐，三河镇的拓样往北传，头一站到的落枫。落枫城南的会馆里，有个老账房认得这个印。''

卢本韦拆开油纸包。里面是一张短笺，字写得很急：拓样已核，本地印谱无。南会馆孙姓老账房，自称见过此印多年，欲见唐姓之人。

唐三站在她旁边，看完了，没说话。

``落枫在哪儿？''奥斯卡问。

``北边。''小伙计往身后指了指，``从这儿往西北走，岔上去落枫的官道，快的话六天。''

所有人都看着蔡须坤。

``去落枫。''蔡须坤说，``蓬莱岛晚些去，岛跑不了。''

\begin{center}\rule{0.5\linewidth}{0.5pt}\end{center}

往落枫去的这六天，卢本韦把殿上听来的那些话说了。第二天夜里先跟队长说的，两个人在火堆边上坐到后半夜。隔了两天，又当着所有人说了一遍。说完以后，队伍里安静了好几天。蔡须坤守夜，从那以后都守后半夜。

唐三那件事，谁都没有再提。

唐三自己也不提。他每天夜里等别人都躺下了，一个人走到离营地远一点的地方，把昊天锤放出来，一锤一锤地往下砸。砸的是空处，什么也不碰，可每一锤落下去，脚底下的土都跟着震一下。

诺当有一回起夜，站在树后头看了半天，回来跟奥斯卡说，三哥那锤子砸得跟打铁一样，一下是一下，不带虚的。

奥斯卡翻了个身，说，他家本来就是打铁的。

小舞也知道。她不去看，只是每天早上起来，把唐三的水囊灌满，放在他行囊边上。霍雨浩陪她去河边打水，两个人一路走一路说话，说的什么，别人听不见。

第六天傍晚，他们进了落枫城。

\begin{center}\rule{0.5\linewidth}{0.5pt}\end{center}

落枫城是北陵州南边最大的一座城。城墙是黄土夯的，外头包了一层砖，城门口两排枫树，这个时节叶子刚开始红。城里的街比天斗城窄，铺面挨着铺面，一直开到城根底下。

正街当中有一片空场，空场北面是落枫的中殿。五进的院子，门前两根旗杆比三河镇那两根高出一截，徽记是鎏金的，太阳一照，老远就晃眼。

唐三走过空场的时候，多看了一眼。

空场上人不少，挑担的、推车的、牵孩子的，可走到中殿门前那一段，人都往两边让，贴着对面的铺面走过去，好像那块地上铺着什么扎脚的东西。没有人往那扇门里看。

霍雨浩也注意到了。他看了看那扇门，又看了看绕着走的人，没说话。

城南是小买卖人扎堆的地方，米铺、布庄、杂货铺，一家挨一家。商帮的南会馆就在这一片，门脸不大，两层小楼，门口挂着一块黑底金字的匾。

卢本韦进去不到一盏茶的工夫，会馆的管事就迎出来了。五十来岁，一身绸袍，走在她半步后头，腰一直没直起来过。

``卢小姐。''管事说，``一路辛苦。孙老先生在后院等着，从拓样送来那天起，他天天问。''

``他在会馆多少年了？''卢本韦问。

``七年。''管事说，``早年在城西市开过铺子，后来铺子关了，投到我们南会馆来记账。账记得细，人也老实，就是不爱说话。''

``那天拓样送到，账房里几个年轻的围着看，他路过瞅了一眼，就站住了。他说这个印他见过，见过好些年。他说要是来的人姓唐，他有话跟他说。''

\begin{center}\rule{0.5\linewidth}{0.5pt}\end{center}

孙老先生的屋子在会馆后院最里头，一排平房的最后一间。

屋子很小。一张床，一张桌，一把椅子，墙上钉着一排木架子，架上码满了账本，书脊朝外，一本本都用毛笔写了年份。窗户对着后墙，只透进来一条光。屋里一股药味。

八个人挤不进去。蔡须坤让其他人在院子里等，只带了唐三和小舞。卢本韦站在门口，霍雨浩站在她旁边，门没关。

老人坐在床沿上，两只手搭在膝盖上，手背上全是斑。他看见唐三进来，眼睛眯起来，盯着看了很久。

``你是他儿子。''老人说。

``是。''唐三说。

老人又看小舞。看了一会儿，他的嘴动了动，没出声。

``这是我妹妹。''唐三说。

``我知道。''老人说，``像你娘。''

小舞的手一下子攥紧了。

老人让他们坐。屋里只有一把椅子，唐三没坐，蹲在老人面前。小舞站在他身后，霍雨浩从门口走进来，站到小舞旁边。

``你爹打的铁，''老人说，``落枫城里没有不认得的。''

\begin{center}\rule{0.5\linewidth}{0.5pt}\end{center}

``那是十七八年前的事了。我在落枫城西市开一间杂货铺，锅碗瓢盆，针头线脑，什么都卖。''

``有一年开春，铺子门口来了两口子，推着一辆独轮车，车上盖着草帘子。男的不说话，女的跟我讲价。帘子一掀开，是一车铁器，菜刀，剪子，锄头，镰刀，还有几把柴刀。''

``我拿起一把菜刀，在磨刀石上走了两下，就知道这东西不一般。钢口匀，刃口薄，敲一下，声音清得跟磬一样。我干了半辈子杂货，没见过山里人能打出这种东西。''

``刀背上都有一个印。''

老人伸出一根手指，在膝盖上画了一个圈，画到右边停住，没合上。

``就是这个。''

唐三没动。

``我问他们，这么好的手艺，怎么不去城里开个铺子，挂个字号。女的笑了，说他不爱跟人打交道，打完了拿出来换点油盐就行。我说那你们登个商号，让商帮拓个印，往后价钱好说。她说不登，登了就有人来找。''

``我当时没听懂。''

``后来他们每年开春来一趟，秋收前再来一趟。东西都放我这儿，我按市价给钱，不压价，也不上铺子的账。我那时候想，这两口子是有来路的人，不想让人知道他们在哪儿，我不问就是了。''

老人停了一下，咳了两声。霍雨浩从桌上端过水碗递给他，他喝了一口，接着说。

``你娘叫王冬，话多。你爹搬东西，她就坐在我铺子门口的条凳上等，一边等一边哼歌。那调子我听了好几年，到现在还记得。''

``她跟我说，家里有个小子，刚会走路，又怀上一个。''老人看着小舞，``就是你了。''

小舞捂住了嘴。

\begin{center}\rule{0.5\linewidth}{0.5pt}\end{center}

``出事那年，你一岁多。''老人对唐三说，``你妹妹还在襁褓里。''

``落枫城里有一家姓------''

他停住了，看了看门口。卢本韦会意，把门掩上了一半。

``有一家，在落枫城里说一不二。家里出过两个魂圣，跟城主府是姻亲，城里一半的铺面是他们家的。那年秋天，他家的管家找到我铺子里来，说府上要打一批兵器，听说有个山里的铁匠手艺好，让我把人叫来。''

``我不敢不叫。''

``你爹你娘来了。那天你爹还是推着那辆独轮车，车上盖着草帘子，你娘走在他旁边，怀里抱着你妹妹。''

``那管家坐在我铺子里，腿翘着，茶是我沏的，他一口没喝。他拿眼睛上上下下把你爹看了一遍，又看你娘，然后从怀里拿出一张契，写得明明白白：往后你爹打的铁，只许卖给他们府上，打什么，打多少，府上说了算。工钱给得很足，足到山里人一辈子见不着的数。''

``你爹没说话，就看着那张契。''

``你娘拿起来看了一遍，当着管家的面，撕了。''

老人说到这儿，嘴角动了一下，像是想笑，又没笑出来。

``她说，我们家的铁，是打给种地的人使的。''

``西市那么多人看着。那管家脸都青了，站起来，把撕成两半的契从地上捡起来，折好，揣进袖子里，一句话没说就走了。''

``你爹一直没开口。等人走远了，他才跟你娘说了一句，你不该撕。''

``你娘说，不撕，你就签了。''

``你爹没再说话。''

屋里没人出声。

``那年冬天，你娘一个人进城来送货。''老人说，``那一趟她没出城。''

``过了半个月，落枫城外的山沟里找着了人。城主府说是走夜路，失足摔下了崖。''

``我去看过。''

他没有往下说。他两只手在膝盖上搓了搓，搓了很久。

唐三蹲在那儿，一动不动。小舞站在他身后，眼泪顺着下巴往下滴，一滴一滴落在唐三的肩膀上。霍雨浩的手搭在她背上。

\begin{center}\rule{0.5\linewidth}{0.5pt}\end{center}

``你爹来收的尸。''老人说。

``他在我铺子里坐了一夜，一句话没说。天亮的时候，他问我，那府里有几个供奉，都是什么修为。''

``我跟他说，你别去。你家里还有两个孩子。''

``他没去。他把你娘背回山里去了。''

``打那以后，他还是每年开春来一趟。东西照样放我这儿，可人不在铺子里坐了。放下东西，他就往城北去，在那府外头的茶棚里坐一天，天黑了再走。一年一回，年年如此。''

``我知道他在看什么。那府里的供奉，哪年添了一个，哪年走了一个，哪个是魂圣，哪个是魂帝。他看了六年。''

老人扶着床沿站起来，走到墙边的木架子前头，在最底下一层摸了半天，抽出一本薄薄的小册子。册子是自己拿线订的，封皮上一个字都没写。

``铺子的账上没有他们。''他说，``可我自己记了一本。''

他把册子递给唐三。

唐三翻开。头几页是一行一行的小字，写得很密：某年开春，唐家，菜刀二十，剪子十，锄头八。某年秋，唐家，柴刀十二，镰刀二十。每一行后面写着给了多少钱。

翻到后面，字变了。

某年开春，唐，锄头六，往城北。

某年开春，唐，镰刀十，往城北。

一年一行，六行。每一行后面都跟着那三个字：往城北。

最后一行只有四个字：唐，空手来。

唐三看着那一行，很久没有翻页。后面没有了，后面全是空白。

唐三的肩膀动了一下。

``八年前开春，他最后来了一趟。那回没带铁器，车也没推，空着手来的。他在我铺子里坐了一盏茶的工夫，问我那府里今年又添了人没有。我说添了一个，是个魂圣。''

``他点点头，站起来，跟我说，孙掌柜，这些年多谢了。''

``然后就走了。往城北去的。''

老人抬起头看着唐三。

``打那以后，落枫城里再没人见过他。''

``那府里什么动静也没有。没有打斗，没有死人，没有抓人。就好像他从来没来过一样。''

``孩子，我说句不好听的。你爹四十来级，那府里光魂圣就好几个。他要是去寻仇，多半斗不过。人没了，连尸首都没有，那府里又一点动静都没传出来\ldots\ldots{}''

他没说完。

``恐怕是让人暗地里处置了。''

\begin{center}\rule{0.5\linewidth}{0.5pt}\end{center}

屋子里很久没有声音。

后墙那条光挪了一点，落到了架子上的账本上。

唐三开口的时候，嗓子是哑的。

``扽魂殿呢？''

老人看着他。

``落枫城有扽魂殿。''唐三说，``出了人命，城主府说是摔的，扽魂殿的人没来看过？我爹一个四十三级的魂宗，在殿里登了十年的册子，说没就没了，也没人问一句？''

老人叹了口气。

``落枫城归北陵州管，城里那座是中殿。''他说，``周边几十个镇的分殿都归它管。你们三河镇那座分殿，往上报，头一站就是落枫中殿，再往上才是北陵大殿。''

``落枫中殿的殿主，姓什么，你知道吗？''

唐三摇头。

``他娶的是那家的三小姐。''老人说，``逢年过节，那府里摆酒，上席坐的就是他。''

``孩子，恐怕殿主都和他们是一伙的。你说，他帮谁呢？''

唐三蹲在那儿，右手慢慢握成了拳。掌心里有一点黑沉沉的光闪了一下，又灭了。

门外，卢本韦的手按在门框上，指节发白。

\begin{center}\rule{0.5\linewidth}{0.5pt}\end{center}

老人又咳了一阵，喝了几口水，缓过来了。

``他走后第二年，我那铺子就开不下去了。''他说，``有人来查税，查了三回。我知道是怎么回事。他那几年年年在我铺子里落脚，那府里的人不会不知道。我把铺子盘了，投到这儿来，一待就是七年。''

``我这把年纪，不怕什么了。可我那时候怕，我一家老小都在落枫城。''

``这些话，我憋了十几年，没跟人说过。''

他看着唐三。

``前天看见那个印，我就想，总算有人来问了。''

他从枕头底下摸出一张折好的纸，没递给唐三，看了看，递给了蔡须坤。

``那家的名字，还有殿主的名字，都在上头。''他说，``给你，不给他。''

蔡须坤接过去，没打开，揣进了怀里。

``他是你们队长？''老人问唐三。

``是。''

``那你听她的。''老人说，``你爹当年要是有人拦着\ldots\ldots{}''

他摆了摆手，没说下去。

\begin{center}\rule{0.5\linewidth}{0.5pt}\end{center}

临走的时候，小舞在门口停住了。

``孙爷爷。''她说，``我娘哼的那个歌，您还记得吗？''

老人愣了一下。

他靠在床头，闭上眼睛，想了一会儿。然后他张开嘴，哼了起来。

他嗓子老了，调子有几处是劈的，哼得断断续续。头两句一出来，唐三就转过了身。

他站在门口，听着。

第三句，老人哼岔了一个音，停住了，想不起来往下怎么走。

唐三接了下去。

他哼得很轻，一个字都没有，只有调子。那调子他记了十四年，从来没哼出声过。接到第四句，老人又跟上了，一老一少，在那间小屋子里把一首没有词的歌哼完了。

小舞靠在霍雨浩身上，哭得站不住。

老人睁开眼睛，看着唐三。

``就是这个。''他说，``她每回坐在我门口，哼的就是这个。''

老人看着他们，又说了一句。

``那个印，我问过你娘是什么意思。''

唐三抬起头。

``她说，是个没画完的圈。''老人说，``我问她为什么不画完。她笑，说等他哪天画完了，再告诉我。''

``我等到今天，也没等到。''

他朝唐三伸出手。唐三愣了一下，才明白，把那本拿线订的小册子还给了他。

老人从枕头底下摸出一支笔，蘸了蘸桌上砚台里干了一半的墨，翻到最后那一行''唐，空手来''，在底下又添了一行。他手抖得厉害，一行字写了很久。

唐三站在旁边看着。那一行写的是：

某年秋，唐家的孩子来了。

老人把笔搁下，把册子合上，又递回给唐三。

``拿着吧。''他说，``我记完了。''

\begin{center}\rule{0.5\linewidth}{0.5pt}\end{center}

从后院出来，其他人都在院子里等着。

没有人问里面说了什么。门没关严，该听见的都听见了。

奥斯卡蹲在墙根底下，手里攥着一根不知道从哪儿捡来的木棍，攥得咔咔响。诺当靠墙站着，一只手按在墙上，墙皮让他按下来一块。蔡兴坤肩上的黑子鸡浑身的毛都奓起来了，一声不吭。

唐三走到院子当中，站住了。

``我去城北。''他说。

没有人说话。

``唐三。''蔡须坤说。

``他们杀了我娘。''唐三说，``我爹去了，八年没回来。那个殿主坐在他们家上席喝酒。''

``我知道。''

``那你让我------''

``我说了，我知道。''蔡须坤看着他，``你现在去，你能干什么？''

唐三没答。

卢本韦站在廊下，这时候开口了。

``落枫城那种人家，我从小见得多。''她说，``年年有人拎着礼到我们家场子里来，求着做生意。在我们家门口，他们笑得比谁都客气。回到自己那块地上，他们就是天。''

``他们不怕一个铁匠。可他们怕比他们大的。''

蔡须坤看了她一眼，转回来对着唐三。

``那府里好几个魂圣，中殿的殿主是他们家女婿，城主府是他们家姻亲。你一个人，拿着一把二流的锤子，走到那府门口，然后呢？''

``然后你就是你爹。''

唐三的嘴唇抖了一下。

``我也想去。''奥斯卡站起来，把木棍扔了，``队长，这事就这么算了？''

``没说算了。''

蔡须坤把那张纸从怀里拿出来，还是没有打开。

``一万年了。''她说，``人这边的，魂兽那边的，当家的我见得多了。没有哪一个是干干净净的。扽魂殿那么大，分总殿底下是各州的大殿，大殿底下是城里的中殿，再往下是镇上的分殿，几百上千座。出一个跟人穿一条裤子的殿主，不稀奇。''

``没有完美的执政者。''

``可这件事，也不能就这么放着。''

她看着唐三。

``我们先找梅长老。''她说，``这张纸，要交到她手上。扽魂殿自己的殿主，让扽魂殿的长老去查。她查不下来，还有大长老。''

``之后怎么处理，让她来。''

唐三看着她，看了很久。

``那个殿主，也是扽魂殿的人。''他说。

``梅长老也是。''蔡须坤说，``你那一手锤法，头一回砸准，是她站在旁边一锤一锤给你数出来的。咱们几个，哪一个没跟她学过东西？''

``扽魂殿里有那个殿主，也有她。''

唐三低下头，没再说话。

院子里静了一会儿。

霍雨浩开口了：``蓬莱岛呢？''

``先不去了。''蔡须坤说，``我们在落枫住下。''

几个人都看向她。

``梅长老要查，总得有人在落枫。那府在哪儿，门朝哪边开，城主府跟中殿平时怎么走动，这些东西，信里写不清楚，她来了也得有人从头摸。''蔡须坤说，``我们人已经在这儿了。也许帮得上她。''

``可谁也不许一个人往城北去。''她看着唐三，``你也一样。''

唐三没应声，过了一会儿，点了一下头。

``信怎么走？''霍雨浩问。

``不能走殿里的驿路。''卢本韦说，``殿里的信，出了落枫头一站就得过中殿的手。走商帮的鸽子，北陵州的会馆收了，再让他们去找梅长老。''

蔡须坤点点头，转身跟管事要了纸笔，就在院子里的石桌上写。

她写得很慢。从三河镇的登记册写起，写孙老先生，写那张契，写唐家娘子，写那府里的供奉，写中殿的殿主坐在上席。每写完一段，她就停一停，把前面的看一遍。写到最后，她把那张名单誊了一份，夹进信里。

霍雨浩把自己第三本笔记里那几页拆下来，也叠好了放进去，那上面有描下来的印，有登记册上的年份，有孙老先生小册子里那几行''往城北''。

卢本韦把信封好，从袖子里摸出那枚小印，在封口的蜡上按了一下，交给管事。

``最快的鸽子。''她说，``今天就放。''

管事双手接过去，小跑着进了后院。

霍雨浩一直没说话。这时候他走到唐三旁边，站定了。

``抽帅气的，抓了人往河滩上一扔。''他说，``打死了人，说是自己摔的。三哥，我比你早知道这些，早了八九年。''

``我那时候也想一个人去。''

唐三看了他一眼。

``后来我没去，是因为你们。''霍雨浩说，``你现在也有我们。''

诺当从墙边走过来，一只大手按在唐三肩膀上，按了一下，什么也没说，又走开了。

唐三站在院子当中，低着头。过了很久，他把右手慢慢松开了。

``好。''他说。

\begin{center}\rule{0.5\linewidth}{0.5pt}\end{center}

那天夜里，他们住在会馆的客房里。

客房在二楼，窗户朝北。月亮很亮，城里的屋顶一片挨一片，灰瓦上落着一层白。再往北，是城北那片大宅子，高墙一道套一道，墙头上挂着灯笼。

唐三睡不着，坐在窗户边上。

小舞推门进来，在他旁边坐下了。过了一会儿，霍雨浩也进来了，靠着门框站着。

三个人谁都没说话。

过了很久，唐三又哼起了那个调子。很轻，一个字都没有。

小舞跟着哼。她记不全，哼到一半就接不上了，唐三就放慢一点，等她跟上来。

霍雨浩站在门口听着，没有出声。

窗外打更的敲过了二更。那首歌哼了一遍，又哼了一遍。

隔壁屋里，蔡须坤坐在床沿上，把那张纸打开了。

她看了很久，又折好，重新揣回怀里。

有人敲门。是卢本韦。

她进来，在桌边坐下，两个人隔着一盏油灯。隔壁的调子透过墙传过来，很轻。

``白天那句话。''卢本韦说，``没有完美的执政者。你是说给唐三听的，还是说给你自己听的？''

蔡须坤想了一会儿。

``都有。''她说。

``一万年前那一回，你们家那位老祖宗跟帝天坐下来谈。两边身后站着的，手上都不干净。人这边杀过魂兽，魂兽那边吃过人。可谈成了。''

``要是非得等一个干干净净的人来谈，那一回就谈不成，两边到今天还在杀。''

卢本韦没出声。

``我不是替那个殿主说话。''蔡须坤说，``该查的，一个都跑不了。我就是想，这回能不能比那回干净一点。''

卢本韦看着她，过了一会儿，点了点头。

她又坐了一会儿，起身走了。走到门口，她回过头：``鸽子傍晚就放出去了。北陵的会馆找到梅长老，快的话五六天。''

``那就等五六天。''蔡须坤说。

然后她起身，推开窗户，看着城北。城北那片高墙底下，灯笼一盏一盏亮着，哪一盏是那家的，她看不出来。

她就那么站着，一直站到了后半夜。

\begin{center}\rule{0.5\linewidth}{0.5pt}\end{center}

第二天一早，唐三去了城北。

他没一个人去。蔡须坤跟着，霍雨浩也跟着。三个人走得不快，跟逛街一样，从城南穿过正街，绕开中殿门前那片空场，往北走。

城北的街宽，街两边是高墙，墙头冒出来的树都修得齐齐整整。走到一条横街口，霍雨浩停了一下，朝左边偏了偏头。

那座府占了大半条街。门口一对石狮子，朱漆大门，门楣上一块匾。门口站着两个门子，腰里挂着刀。

府对面的街角，有一个茶棚。

茶棚很旧，四根竹竿撑着一块油布，底下摆着三张桌子。守棚子的是个老婆婆，正拿蒲扇扇炉子。

唐三在茶棚里站了一会儿，挑了最靠墙角的那张桌子坐下。那个位置正对着府门，隔着一条街，门里进出什么人，看得清清楚楚。

老婆婆提着壶过来，给三个人倒茶。倒到唐三那碗，她停了一下，眯着眼睛看他。

``你这孩子，''她说，``长得像早些年常来的一个人。''

唐三端着茶碗，没动。

``每年开春来一回。''老婆婆说，``一个人，就坐你这个位置，一壶茶从早上喝到天黑，话也不说，就看着对面。我问过他一回看什么，他说看人。''

``后来呢？''蔡须坤问。

``后来就不来了。''老婆婆想了想，``最后一回，是七八年前开春。那天他没坐到天黑，坐到半下午，茶钱搁在桌上，起身往那边去了。''

她拿蒲扇往府墙那头指了指。府墙拐过去，是一条窄巷。

``往那巷子里去的。''

她说完，提着壶又回炉子边上去了。

唐三看着那条巷子。巷口有一棵老枫树，叶子红了一半，风一吹，落下来几片，落在巷子里的青石板上。

他端着那碗茶，一直没喝。

蔡须坤坐在他对面，也没催他。霍雨浩把笔记摊在膝盖上，照着府门、门子、巷口那棵树，一笔一笔画了一张图。

太阳从东边挪到了头顶，又往西偏了一点。

唐三把茶碗放下了。

``走吧。''他说，``明天再来。''

从那天起，他们在落枫住了下来。

\begin{center}\rule{0.5\linewidth}{0.5pt}\end{center}

\chapter{第九十五章：长老令}\label{ux7b2cux4e5dux5341ux4e94ux7ae0ux957fux8001ux4ee4}

在落枫住下的头五天，唐三每天早上都去城北那个茶棚。

蔡须坤排了班，每天换两个人陪着他。头一天是她自己和霍雨浩，第二天是诺当和奥斯卡，第三天是小舞和卢本韦。三个人进了茶棚，要一壶最便宜的茶，坐到过午，再换个方向绕回城南。

老婆婆不问他们是谁。她只是每天把墙角那张桌子擦干净，茶端上来的时候，唐三那碗总比别人的满一点。

霍雨浩那张图越画越细。

府门什么时辰开，什么时辰关，每天进出几趟车，车上盖不盖布，他都记下来了。府里有两个人出门从不走正门，走西边的角门，一个瘦高，一个矮胖，走路脚底下都不带声。

``魂圣。''蔡须坤看了一眼就说，``两个都是。''

第三天，小舞在正街上看见一面墙。

墙在中殿斜对面，是一家关了门的布庄。墙上贴满了纸，一层压一层，底下的早就发黄了，边角卷起来，风一吹哗哗响。

小舞凑过去看。全是寻人的。

寻夫，寻妻，寻儿子，寻闺女。有画像的，有写生辰八字的，有写''必有重谢''的。最底下那一层的纸已经脆了，字迹洇成一团，只看得出一个年份，是十一年前。

``怎么贴在这儿？''小舞问旁边一个卖烧饼的。

卖烧饼的看了她一眼，又看了一眼对面中殿门口那两根旗杆。

``报了也没用。''他压低了声音，``去中殿报人丢了，头一回还给你登个册子，第二回人家就不让进门了。贴在这儿，好歹路过的人能看一眼。''

``丢的人多吗？''

卖烧饼的没答。他把一张烧饼包好塞给小舞，没收钱，推着车走了。

那天夜里，小舞把这件事跟霍雨浩说了。霍雨浩听完，第二天一早就去了那面墙，从头到尾抄了一遍。

一共一百三十七张。

他抄完回来，把笔记摊在桌上，一张一张对年份，对完了，抬起头。

``十二年前开始多起来的。''他说，``之前一年三四张，之后一年十几张。''

``丢的人里，男的女的都有，十几岁到四十来岁。''他停了一下，``寻人纸上写了长相的，差不多都写了一句，长得周正。''

屋里没人说话。

霍雨浩把笔记合上，手在封皮上按了很久。

\begin{center}\rule{0.5\linewidth}{0.5pt}\end{center}

第六天夜里，南会馆的后门被人敲了三下。

管事去开的门，开了门就愣住了。门外站着两个人，都穿着灰布袍子，风帽压得很低。前头那个把风帽往后一掀，是个年轻女人，看着不过二十出头，眉眼清清淡淡，一头黑发挽得一丝不乱。

``我找蔡须坤。''她说。

蔡须坤是从楼上跑下来的，鞋都没穿好。她在楼梯口站住了，叫了一声：``梅长老。''

梅天狸看着她，看了一会儿，走过去，伸手把她耳朵边上一绺乱发拨了回去。

``瘦了。''她说。

后面那个人也掀了风帽。六十来岁，方脸，下巴上一圈短胡子，眉毛很浓。他朝楼上下来的几个人扫了一眼，点了点头，没说话。

``北陵大殿的严殿主。''梅天狸说，``我从北陵过来，他跟着一块儿来的。''

严殿主朝她欠了欠身，才跟着往楼上走，一直落后她半步。

严殿主八十六级。四十年前，他还是北陵大殿一个管卷宗的执事，梅长老巡察北陵，他在一旁磨了三天墨。那时候她就是这副模样。四十年过去，他自己胡子白了一半，她一点没变。

在北陵州，严殿主是数得着的人物。可在梅天狸跟前，他只能站着。她是总殿的长老，九十七级，又是上一代的圣女。一州大殿的殿主进了总殿，连长老院的门槛都迈不进去。

\begin{center}\rule{0.5\linewidth}{0.5pt}\end{center}

会馆的管事把楼上最里头那间屋子腾出来，关了门，又亲自守在楼梯口。

八个人加上梅天狸和严殿主，挤了一屋子。油灯挑得很亮。

梅天狸先看的是唐三。

她看了很久，才开口：``你娘的事，信我看了。你爹的事，信我也看了。''

唐三站着，没出声。

``殿里欠你一句话。''梅天狸说，``这句话，等我把事办完了再说。现在说，轻了。''

她转过头，对严殿主点了一下。

严殿主从怀里拿出一个油布包，打开，里面是一叠账页，边角都用朱笔圈过。

``半年前，北陵大殿核各中殿的年报。''他说话很慢，一个字一个字往外吐，``落枫中殿报上来的案数，跟底下分殿年底直接报到大殿的对不上。分殿报上来有一百件，到了落枫中殿往上转，只剩七十来件。少的那些，差不多都是人口走失。''

``我派人来查过一回。落枫中殿的账做得干净，人也客气，查了半个月，什么都没查出来。''

``我那时候就该亲自来。''

他把账页往桌子中间推了推。

霍雨浩看了一眼，把自己的笔记翻开，放在账页旁边。

严殿主低头看了看那一百三十七张寻人纸的抄本，又看了看自己的账页，眉头越皱越紧。

``这是哪儿来的？''

``中殿对面那面墙。''霍雨浩说，``一张一张抄的。''

严殿主没再问。他拿手指点着两边的年份，一行一行往下对，对了很久。

``对得上。''他说。

梅天狸看向蔡须坤：``你信里说那府里有两个魂圣。''

``两个。一个瘦高，一个矮胖，走西边角门。''

``府墙拐角那条巷子，''霍雨浩说，``西边第三段墙，砖跟两边的不一样。两边是老砖，发黑，那一段是后砌的，颜色浅，砖缝里的灰也新。我量过，一丈二尺长。''

梅天狸看着他。

``你们在这儿住了六天，''她说，``就干了这些？''

``还喝了六天茶。''奥斯卡说。

梅天狸没笑。她把桌上的东西一样一样看过去，看完了，把手放在那张图上。

``明天。''她说，``先去那条巷子，再去中殿。你们都跟着。''

``都去？''严殿主抬起头。

``都去。''梅天狸说，``事是他们查出来的。唐家的孩子，有权看着。''

\begin{center}\rule{0.5\linewidth}{0.5pt}\end{center}

第二天一早，梅天狸先去了巷子。

她没穿长老的袍子，还是那身灰布衣裳，走在最前头。看着跟队里这帮孩子也差不了几岁，像个带着师弟师妹出来办事的师姐。严殿主换了一身寻常的绸衫，跟在她后头。

巷子窄，两个人并排都嫌挤。巷口那棵老枫树落了一地的叶子，踩上去沙沙响。

梅天狸走到霍雨浩说的那一段墙前面，站住了。

她伸出手，掌心贴在墙上，贴了很久。

没有人出声。唐三站在她身后，盯着那只手。

``这一段墙底下，''梅天狸说，``原先有过一场架。''

她把手收回来。

``砖是新的，地基是旧的。地基里还压着一点魂力的残迹，八年了，散得差不多了，只剩一层皮。''她看着唐三，``两股。一股是那种人的，阴的，黏的，跟抽帅气那帮人一个路数。另一股\ldots\ldots{}''

她停了一下。

``另一股很沉。''她说，``硬，直，不拐弯。像一锤砸下去。''

唐三的手指动了一下。

``谁赢谁输，看不出来。''梅天狸说，``只看得出来，那一锤砸得很重，把地基都砸裂了一道缝。后来砌墙的人把缝填上了，可填不回去。''

唐三往前走了一步，蹲下来，把手按在墙根底下那道缝上。

缝很细，填着灰，不仔细看看不出来。他的手指顺着那道缝摸过去，摸到头，又摸回来。

他在那儿蹲了很久。小舞走过去，蹲在他旁边，把手也按了上去。

\begin{center}\rule{0.5\linewidth}{0.5pt}\end{center}

落枫中殿的正门，辰时开。

梅天狸让严殿主在空场边上的茶摊坐着，自己也坐下了，要了一壶茶。

``你们先进去。''她对蔡须坤说，``就说要调唐家娘子的卷宗。看他怎么说。''

蔡须坤点点头，带着其他七个人穿过空场。

空场上的人照旧绕着走。八个半大孩子直直地冲着中殿大门过去，好几个挑担子的都停下来看。

门口的执事拦住了他们。问明白了来意，执事的脸色变了变，让他们在门厅等着，自己往里跑。

等了一炷香的工夫，里头出来一个人。

五十来岁，白面长须，一身绛紫的殿主袍子，腰里系着玉带，手上一枚扳指，走路袍角一点不乱。后头跟着四个执事。

他在门厅当中站住了，没往前走，隔着一丈远把八个人看了一遍，从左看到右，又从右看到左。看到唐三的时候，他的目光停了一下。

``调卷宗？''他开口了，声音不高，``谁让你们来的？''

``我们自己来的。''蔡须坤说，``十四年前，城外山沟里死了一个女人，城主府说是失足。我们想看看中殿的卷宗。''

殿主笑了一下。

``中殿的卷宗，是给殿里的人看的。''他说，``你们是哪个殿的？分殿的执事？中殿的执事？还是哪位长老的门下？''

没人答。

``都不是。''他点点头，``一群学院里的孩子，跑出来历练，走到落枫，听了几句街上的闲话，就想进中殿翻卷宗。''

``你们知不知道，中殿的卷宗库，连城主都得递帖子才能进？''

奥斯卡往前走了半步，被诺当一把拽了回来。

殿主的目光又落到了唐三身上。

``你姓什么？''

``唐。''

殿主脸上那点笑，一下子淡了。

他没再问。他转过身，对身后的执事摆了摆手：``送客。往后这几个人再来，不用通报。''

``卷宗的事------''蔡须坤说。

``卷宗？''殿主回过头，看着她，``十四年前的案子，早就结了。失足坠崖，无可疑处。城主府验过，中殿核过，白纸黑字，还有什么好看的？''

``小姑娘，我劝你一句。落枫城不大，路窄，天黑了不好走。带着你这帮同学，早点回学院去。''

两个执事上前一步，伸手来推。

门口传来一个声音。

``白纸黑字，底下押的是谁的印？''

\begin{center}\rule{0.5\linewidth}{0.5pt}\end{center}

梅天狸从门外走进来。

还是那身灰布衣裳，头发挽得一丝不乱。她一步一步走进门厅，走得不快，脚步声在石板地上一下一下响。

严殿主跟在她身后，已经换上了大殿殿主的袍子。玄色，绣金，胸口一个''北陵''。

落枫中殿的殿主看见严殿主，脸色变了一变，拱了拱手：``严殿主怎么来了？也不提前知会一声，我好出城去接。''

他又看了一眼梅天狸。一个二十来岁的年轻女人，灰布衣裳，身上没佩徽记。他只当是严殿主身边带的哪个执事，目光停了一瞬就挪开了。

``严殿主身边这位是\ldots\ldots 新提的执事？''

严殿主没答。他站在梅天狸身后，落后她半步，一动不动。

梅天狸没理他。她从袖子里摸出一块令牌，往前一递。

令牌只有巴掌大，乌木底子，正面一个篆体的''长''字，四周一圈金纹，背面刻着扽魂殿的徽记，徽记底下一行小字。

殿主低头看了一眼那行小字。

他脸上的血色一下子退干净了。

``长\ldots\ldots 长老\ldots\ldots{}''

他又抬头看那张脸。总殿的梅长老，上一代的圣女，八十四岁上破的九十七级，这件事整个扽魂殿没有不知道的。他怎么也没法把那个名字，跟眼前这张二十来岁的脸对到一块儿。

``总殿长老，梅天狸。''梅天狸说，``持长老令，查落枫中殿。''

殿主往后退了一步，撞在身后执事的身上。他张了张嘴，想说什么，嗓子里只挤出来一点气声。

然后他忽然挺直了腰。

``梅长老。''他说，声音还稳，``下官落枫中殿殿主姚某，在殿里二十年，没有出过一件错案。长老要查，下官自然开库。可长老今天带着一群来历不明的孩子闯门，又不走大殿的公文------''

梅天狸抬起了眼睛。

门厅里的空气忽然沉了下去。

没有风，没有声音，也没有光。就是沉。像是整个门厅被人从上面按住了，按进了很深很深的水底下。

站在殿主身后的四个执事，膝盖一软，接二连三跪了下去。

唐三站在门厅边上，只觉得身上一重，像背上压了一袋米。他看了看两边，队里其他人也都只是晃了一下，站住了。蔡兴坤肩上那只黑子鸡，连毛都没动一下。

那股魂力是冲着一个人去的。

姚殿主脚底下的石板裂了。

先是一道缝，从他左脚底下往外爬，爬到一尺远，分成了两道，又分成四道。他整个人往下一沉，膝盖弯了，又拼命撑起来，两只手死死按着膝盖，按得指节发白。

霍雨浩站在唐三旁边，嘴唇动了动，说得很轻。

``Y(1,3,9)。''

唐三转过头看他。

``那股魂力展开的样子。''霍雨浩说，眼睛一直盯着梅天狸，``一层压一层，每一层往外铺的时候都照着前一层的样子再铺一遍。一、三、九，三层。我在书上见过这个写法，没想到能亲眼看见。''

他停了一下。

``我数不到底。''

唐三没接话。他想起大赛那年。

那时候牢大还在，队里九个人。梅长老说要看看他们的成色，让九个人一块儿上，不许留手。他们也真没留手，九个扽魂一起放出来，魂力拧成一股，是他们那时候能打出来的最狠的一下。

梅长老站在原地，只抬了一下手。

那一下撞上去，像撞进了一团棉花里，连个响都没有，就散了。九个人被自己的力道弹回来，摔了一地。奥斯卡的香肠撒了半条街。

事后霍雨浩问她那是什么。她说，Y(1,3)。

``那回是两层。''唐三低声说。

``一、三。''霍雨浩说，``今天是三层。''

两层，就接住了九个人拼命的一下。三层，全压在一个人身上。

姚殿主六十多级，在落枫城里坐了十五年的殿主，城主见了他都要起身。他这辈子没见过这个东西，也没想过有人拿这个东西对着他。

姚殿主撑不住了。

他的膝盖砸在裂开的石板上，砸出一声闷响。他想用手撑着地站起来，手按下去，胳膊抖得像筛糠，撑了两回，都没撑起来。

最后他整个人瘫坐在地上，后背靠着一根柱子，两条腿软得像是没有骨头。那身绛紫的袍子铺在碎石板上，玉带歪到了一边。

梅天狸一步一步走到他跟前，低头看着他。

``不走大殿的公文。''她说，``严殿主就站在这儿。你要公文，让他当场给你写。''

姚殿主的嘴唇抖着。

``还要吗？''

他拼命摇头。

梅天狸把那股魂力收了。

门厅里一下子轻了。跪着的执事们大口喘着气。姚殿主瘫在柱子底下，还是站不起来，脸上全是汗，顺着长须往下滴。

``唐家娘子的卷宗。''梅天狸说，``底下押的是谁的印？''

姚殿主没答。

``我再问一遍。''

``\ldots\ldots 是我的。''

\begin{center}\rule{0.5\linewidth}{0.5pt}\end{center}

卷宗是严殿主亲自去库里调出来的。

薄薄一册，封皮上写着年份，里头只有三页。第一页是城主府的验尸文书，写着''失足坠崖，颅骨碎裂，无他伤''。第二页是中殿的核验，写着''与城主府所验相符，无可疑处，结案''。第三页什么都没有，只在右下角押着一方朱印。

梅天狸把卷宗递给唐三。

唐三接过来，一页一页看完了。看到那个''无他伤''，他停了很久。

孙老先生说过，他去看过。孙老先生说到那儿就没往下说。

唐三把卷宗合上，递回给梅天狸，一个字也没说。

姚殿主被两个北陵来的执事架着，坐在一张椅子上。他已经能坐了，可两只手还在抖，茶碗端起来，洒了一半。

梅天狸在他对面坐下。

``从头说。''她说。

\begin{center}\rule{0.5\linewidth}{0.5pt}\end{center}

姚殿主说了一个多时辰。

十五年前，他刚升落枫中殿殿主的第二年，城北韩府的家主请他吃酒。酒桌上多了两个生面孔，一个瘦高，一个矮胖，说是府上新请的供奉。

那顿酒吃完，他袖子里多了一张银票。

``头两年，只是些小事。''他说，``韩府的人在城里打了人，伤了人，报到中殿来，我往下压一压。''

``后来\ldots\ldots{}''

他停住了，看了一眼严殿主，又看了一眼梅天狸。

``后来，大概十二年前，城里开始丢人。''

``那两个供奉，是帅气教的。''姚殿主说，``两个都是魂圣，比我高出一大截。我头一回知道的时候，想抽身已经晚了。''

``他们在城里挑人，挑长得好的，夜里弄走，弄到韩府后院的地窖里。抽完了，人就没用了，往城外一扔，扔到山沟里，扔到河里。''

霍雨浩的手按在桌沿上。桌沿是木头的，被他按出了几道白印。

小舞的手从桌子底下伸过去，握住了他。

``丢了人，家里人来中殿报案。''姚殿主说，``我让执事登个册子，登完了往库里一放，不往上报。报一回，登一回。报两回，就不让进门了。''

``一百三十七张。''霍雨浩说。

姚殿主愣了一下。

``中殿对面那面墙上，贴了一百三十七张寻人纸。''霍雨浩说，声音很平，``你每天从那面墙底下过。''

姚殿主低下了头。

``唐家娘子呢？''梅天狸问。

``那是另一回事。''姚殿主说，``她撕了韩府的契，韩家主咽不下这口气。那两个供奉说，这种人，留着也是个麻烦，顺手就处置了。城主府那边是韩家的姻亲，验尸文书是现成的，送到中殿来，我押了印。''

``就这么简单？''蔡须坤说。

姚殿主没敢看她。

``唐舞麟。''梅天狸说。

姚殿主的身子缩了一下。

``八年前开春，那人夜里从西边巷子摸进去的。''他说，``那天夜里的事，我不在场。第二天一早，韩家主派人来叫我。我去的时候，巷子里那段墙已经塌了，地上全是碎砖。矮的那个供奉躺在后院，胸口塌下去一块，躺了三个月才下地。''

唐三抬起了头。

``人呢？''梅天狸问，``唐舞麟呢？''

``瘦的那个说，处置了。''姚殿主说，``我问他尸首呢，他说不用我管。他让我把那段墙重新砌起来，让城主府把那一片的更夫换了。我照办了。''

``你见着尸首了？''

``\ldots\ldots 没有。''

屋子里很静。

唐三坐在那儿，两只手放在膝盖上，慢慢握紧了。

``前年秋天，''梅天狸说，``有人去南边山里打听唐家的两个孩子。''

姚殿主闭上了眼睛。

``是我派的人。''他说，``帅气教的人让我派的。他们说，唐舞麟那一锤，砸伤了他们的人，这种人家，要看看还有没有后。要是有，就记下来，等以后用得着。''

小舞的手一下子攥紧了霍雨浩的。

蔡须坤站了起来。

她没说话，就那么站着，看着姚殿主。她站起来的那一刻，窗户纸哗地响了一声，外头没有风。

梅天狸伸出手，按在她手腕上。

蔡须坤看了梅天狸一眼，慢慢坐回去了。

``帅气教那边，谁跟你接头？''梅天狸问。

``每年秋天从西边来一个人，收东西。''姚殿主说，``抽出来的那些，装在黑陶罐子里，一罐一罐封好，从韩府后门运走。来收的人每年不一样，都戴着面具。''

``罐子怎么运？''卢本韦问。

``混在韩家药材行的货里。''姚殿主说，``挂药材行的旗，走西边的商道。''

卢本韦转头看了霍雨浩一眼。``韩家药材行每年秋天往西发一趟大货，商帮的册子上都有。发到哪儿，谁接的货，一查就知道。''

梅天狸点了一下头，示意她记下。

``去年开春那回换了人。来的是个年轻的，二十来岁，没戴面具，那两个供奉见了他都得弯腰，一口一个小少爷。有一回那个瘦的叫漏了嘴，叫了一声峰少爷，叫完自己脸都白了。''

霍雨浩的眼睛眯了起来。

``他说什么了？''

``他说，落枫这边的买卖做得不错。''姚殿主说，``他还说，让落枫的人怕中殿，比让他们怕帅气教好用。人怕帅气教，会去找扽魂殿。人要是连扽魂殿都怕了，就没地方去了。''

``他说这话的时候，一直在笑。''

屋子里没有人出声。

霍雨浩慢慢地说了两个字。

``胡峰。''

姚殿主茫然地抬起头。那人从头到尾没报过名字。

``胡继的儿子，胡赞的孙子。''霍雨浩说。

他的手按在怀里那本笔记上，按得很紧。宫变那一夜，扮成太子、在巷子口念''果糕小队''的，就是这个人。那个腔调，他到今天都记得。

蔡须坤看着他，没出声。去年开春，胡峰来过落枫。

\begin{center}\rule{0.5\linewidth}{0.5pt}\end{center}

那天夜里，严殿主带着北陵来的人封了韩府。

梅天狸亲自去的。唐三他们没能跟去，蔡须坤把所有人都按在南会馆里，自己守在楼梯口，一夜没合眼。

后半夜，城北那边亮了一下。

很短，一闪就灭了。灭了以后，整座落枫城的狗都叫了起来，叫了很久才停。

天亮的时候，梅天狸回来了。

她身上那件灰布衣裳，袖口破了一道口子。她进门的时候，蔡须坤迎上去，她摆摆手，说没事。

``瘦的那个跑了。''她说，``他身上带着帅气教的遁符，我追到城外，没追上。矮的那个没跑成。''

``地窖里还有四个人。''

她停了一下。

``活着。''

\begin{center}\rule{0.5\linewidth}{0.5pt}\end{center}

韩府封了，梅天狸没走。

``一座中殿压了十二年的案子，不可能只经姚某人一个人的手。''她说，``底下的执事，城主府的人，关卡上收税的，一环扣一环。少一环，这条路就走不通。''

她把韩府的账本摊在南会馆的桌上。账本是从地窖里一个铁箱子里翻出来的，一共九本，记得比商帮的账还细：哪年哪月，从哪儿弄来的人，抽了几罐，给了谁多少银子。

``这是帅气教的规矩。''梅天狸说，``他们自己人也不信自己人，什么都要记。''

接下来半个月，八个人跟着她，一件一件往下查。

霍雨浩拿一百三十七张寻人纸跟韩府的账本对。账本上不写名字，只写日子，写''男，十九''\,``女，二十三''。他一张一张对日子，对上一张，就在抄本上画一个圈。对到第九天，画了一百一十二个圈。

剩下二十五张对不上。梅天狸看了，说这二十五个，要么是真走丢了，要么是别家干的，让严殿主带人另查。

卢本韦去了商帮。韩家药材行十二年里往西发了十二趟大货，每一趟出城都要过西门的关卡，关卡上验货的是城主府的税吏。十二趟，税单上写的都是''药材，无夹带''，签的是同一个人的名字。

小舞和奥斯卡去了城外的村子。好几户丢了孩子的人家说，当年去中殿报案，是城西巡检司的人拦在半路上，说中殿忙，让他们回去等消息。他们等了十年。

诺当跟着严殿主去了落枫下面的几个镇子。有两个镇的分殿执事，收了中殿递下来的条子，把人口走失的案子改成了''逃役''，往上报的时候，就不算丢人了。

唐三跟着梅天狸去了城主府。城主是韩家的姻亲，一问三不知。问到第三天，他手下的主簿先扛不住了，把十四年前那份验尸文书的底稿，从自己家的夹墙里拿了出来。

底稿上写的是：颅骨碎裂，颈侧有指痕。

誊到正本上，只剩了前半句。

唐三看着那张底稿，看了很久，把它还给了梅天狸。

蔡须坤哪儿都去。她不怎么说话，就跟在梅天狸后头，看她怎么问人，怎么翻账，怎么在一张纸上找出三个对不上的日子。有一回梅天狸回头看她，问她看出什么了。

``您问人的时候，先问他家里几口人。''蔡须坤说。

``为什么？''

``他答了这个，后面就不好撒谎了。''

梅天狸看了她一会儿，没说对不对，转过身接着翻账。

半个月查下来，名单从十一个人，变成了二十六个。

\begin{center}\rule{0.5\linewidth}{0.5pt}\end{center}

处决那天，落枫中殿门口的空场上搭了一座台子。

韩府家主，当年拿着契去西市的那个管家，没跑掉的那个供奉。落枫中殿的姚殿主，经手压案的五个执事。城主府验尸的仵作，改文书的主簿，西门关卡上签了十二年税单的税吏，城西巡检司拦人的巡检和他手下三个差役。落枫下面两个镇子改案卷的分殿执事。还有韩家药材行里跟着押货的两个掌柜，地窖里看门的三个家丁。

城主革了职，押去北陵受审。

二十六个人跪在台上，把台子跪满了。

严殿主念的罪状。念了一个多时辰，一百一十二个名字，一个一个念过去。念到一半，台底下有人开始哭，哭声从前头传到后头，传满了整片空场。

唐三站在人群最后面，听着。

念到''唐家娘子''四个字的时候，他身边的小舞把脸埋进了霍雨浩的肩膀。

念完了，严殿主把罪状卷起来，退到一边。

梅天狸没让刽子手上台。

她自己走上去，从台子这头走到那头。她走得很慢，每走到一个人跟前，就停一下。她一停，那人就低下头去，再也没抬起来。

从头到尾，没有一点声音。

她走到最后一个人跟前，停了一下，然后转过身，站在台子当中。

严殿主站在台边上，看着那一排低下去的头，长长吐了一口气。

``正义终于到了。''他说。

台底下最前面一排，站着一个中年女人，怀里抱着一张发黄的纸。纸是从中殿对面那面墙上揭下来的，边角都脆了，上头画着一个十几岁的姑娘。

她没哭。她看着台上，嘴唇动了动，声音很小，小得旁边的人都没听清。

``迟到的正义，也能算正义吗？''

梅天狸听见了。

她从台上走下来，穿过人群，走到那个女人跟前。人群往两边让开了一条道。

那女人看着她，往后退了半步，没再退。

``当然不算。''梅天狸说。

女人愣住了。

``你闺女丢了十一年。''梅天狸看了一眼她怀里那张纸，``这十一年，你每天从那面墙底下过，看着这张纸一年比一年黄。今天台上这些人，换不回来。''

``四十年前，我来北陵巡察过一回。那时候落枫的中殿还干净。后来落枫报上去的册子，一年比一年干净，殿里年年给它记个上等。我就再没来过。''

``这四十年里，我要是再来一回，哪怕一回\ldots\ldots{}''

她没有说下去。

``这件事，有我一份。''

女人抱着那张纸，站在那儿，眼泪终于下来了。

梅天狸没有走。她站在女人面前，等她哭完。

哭完了，女人把那张纸叠好，揣进怀里，朝梅天狸点了一下头，转身走了。

台底下很久没有声音。然后，人群最前面，一个拄着拐的老头往前挪了一步，又挪了一步，挪到了中殿门口。

他站在那两根旗杆底下，站了一会儿，抬脚迈进了门槛。

\begin{center}\rule{0.5\linewidth}{0.5pt}\end{center}

\chapter{第九十六章：藏卷楼}\label{ux7b2cux4e5dux5341ux516dux7ae0ux85cfux5377ux697c}

那天傍晚，梅天狸在南会馆楼上那间屋子里，把八个人叫齐了。

``韩府的账本，你们都看过了。''她说，``帅气教在落枫埋了十几年的人，图的不光是抽几个人的帅气。''

``你们听见姚某人那句话了。让人怕中殿，比让人怕帅气教好用。''

``说这句话的是胡峰，胡继的儿子，胡赞的孙子。胡赞把自己的孙子派到一座中殿来，图的不光是几罐帅气。''

``它要的，是让底下的老百姓觉得，扽魂殿跟那些人是一伙的。落枫城里这些年，老百姓见了中殿门口那对旗杆，是绕着走的。''

她停了一下。

``今天台底下那个女人问我那句话，我一直在想。''

``殿里有位真霜长老，今年一百三十六岁，管了殿里六十年的卷宗。前些年我去看他，他跟我说过一句话。他说，这五十年，各地报上来的案子，贪的、压的、欺上瞒下的，一年比一年多。五十年前一年几件，如今一年几十件。''

``我那时候以为，殿大了，人多了，难免的。''

``现在我不这么想了。''

``报上来的多，是看得见的。看不见的，在落枫这样的地方。''

``落枫的册子为什么干净，你们在这儿这些天也看见了。丢了人，只敢往墙上贴一张纸，不敢进中殿的门。孙老先生跟你们说那些旧事，说到一半，先要看一眼门口。''

``帅气教在这儿，管的不光是中殿的账，还管老百姓嘴里说什么。案子先压在嘴里，再压在中殿，有的地方，连分殿也一块儿压。一层一层压上去，到了大殿，就只剩一本干净册子。''

她走到窗边，往西看了一眼。

``北陵往西，顺着商道一直走，走到头就是帅气教的地界。韩家的药材走这条路，走了十二年。''

``帅气教要是从几十年前就开始往殿里埋人，一座中殿一座中殿地埋，埋得不声不响\ldots\ldots{}''

她没说完。

``落枫恐怕只是离它近的一处。''

屋子里没有人出声。

``那怎么办？''奥斯卡问。

``一处一处查。''梅天狸说，``回去我先报给二长老，再去见真霜长老，把他那六十年的卷宗从头翻一遍。''

``二长老？''奥斯卡问。

``Smile Lee。''梅天狸说，``十级魂导师。''

霍雨浩抬起头：``书上说，魂导师到九级就到顶了。''

``在他之前是。''梅天狸说。

屋子里静了一下。

``这件事交到他手里，比交到我手里合适。''梅天狸说，``这是殿里的事。你们该去蓬莱岛，还去蓬莱岛。''

她转过头，看着唐三。

``你爹的事，我没查完。那个瘦的，我会接着追。追到了，我亲自问他。''

唐三点了点头。

``姚某人说处置了。''他说，``可他没见着尸首。''

``没有。''

``那两个供奉，一个胸口塌了，躺了三个月。''唐三说，``那一锤，是我爹砸的。''

梅天狸看着他。

``没人见过他死。''唐三说，``在见着之前，我当他还活着。''

梅天狸没有说他对，也没有说他不对。她走过去，在唐三肩膀上拍了两下。

``你娘那件事，''她说，``殿里欠你一句对不住。现在事办完了，我替殿里说。''

``对不住。''

唐三站着，过了很久，朝她弯下腰，行了一礼。

\begin{center}\rule{0.5\linewidth}{0.5pt}\end{center}

第二天一早，梅天狸和严殿主往北去了。

队伍在南会馆门口送他们。梅天狸走到街口，又折了回来。她从袖子里摸出一块拇指大的玉符，青白色，拴着一根红绳，挂到了蔡须坤脖子上。

``下回有事，别写信了。''她说，``捏碎了就行。''

蔡须坤低头看了看那块玉符，笑了一下，点点头。

梅天狸转身走了。

唐三没跟着回会馆。他一个人去了城北，又在那个茶棚里坐下了。

老婆婆给他倒了茶。今天对面府门口没有门子了，大门上贴着北陵大殿的封条，门口那对石狮子底下，落了一地的枫叶。

唐三坐了一会儿，站起来，走进了那条窄巷。

他在那段新砌的墙跟前站住，蹲下来，又把手按在墙根那道缝上。

``爹。''他说。

声音很轻，巷子里没有人。

他蹲了很久，站起来，拍了拍手上的灰，走出巷口。

小舞和霍雨浩在巷口等着他。小舞手里提着一包烧饼，是正街上那个卖烧饼的硬塞给她的，说今天不要钱。

``哥，吃不吃？''

唐三接过一个，咬了一口。

``走吧。''他说，``去蓬莱岛。''

\begin{center}\rule{0.5\linewidth}{0.5pt}\end{center}

梅天狸回到天斗城，是离开落枫的第十四天。

她没回自己的住处，也没去长老院换衣裳。进了扽魂殿的大门，她穿过前面三进院子，一直往后走，走到藏卷楼底下，才停住脚。

藏卷楼七层，砖是青的，窗是窄的，一年到头关着。守楼的执事认得她，赶紧开门。

``真霜长老在几层？''

``七层。''执事说，``这几天没下来过。''

梅天狸上了七层。

真霜长老坐在一张长案后头，案上摊着十几本卷宗，都翻开着。他一百三十六岁了，头发白里带着一点蓝，腰还挺得笔直，手边一盏油灯，灯芯剪得很短。

他听见脚步声，没抬头。

``北陵的事，我听说了。''他说，``你的信比你先到。''

``我来借您的卷宗。''梅天狸说。

``哪一年的？''

``五十年的。''

真霜长老这才抬起头来，看了她一会儿。

``我等这句话，等了十几年了。''他说。

他站起来，走到身后那一排木架子前头，从最底下一层开始，一格一格往上指。

``这一层是五十年前的，一年一册，那时候薄。往上一层，厚一点。再往上，又厚一点。''他的手指一路往上走，走到最上头那一层，停住了，``这是去年的。''

最底下那一册，一指来厚。最上头那一册，跟一块砖差不多。

``我跟你说过，一年比一年多。''真霜长老说，``你那时候说殿大了，人多了，难免。''

梅天狸没说话。

``我那时候也这么想。''真霜长老说，``可我管了六十年的卷，有一件事我一直没想通。''

``什么事？''

``这些案子，东边多，西边少。''他说，``越往西，越少。最西边那几州，一年报上来三五件，有的年份一件都没有。五十年了，一直是这样。''

他把手放下，回过身来。

``西边那几州，人不比东边少，商道比东边还热闹。哪一州的老百姓能干净成这样？我往长老院递过两回条子，批回来的都是一句：西边民风淳朴，治理有方。''

梅天狸没有马上接话。

``落枫的册子，在北陵也是年年最干净的。''她说。

``西边那几州，''真霜长老说，``都是落枫。''

屋里静了一会儿。油灯的灯芯噼啪响了一声。

``这回要是下去查，''他说，``算我一个。''

梅天狸愣了一下：``您一百三十六了。''

``腿脚是慢了。''真霜长老说，``看账不用腿。''

他的境界停在 Y(1,3,6)，停了快二十年。可殿里管卷的人都知道，论翻账、对账、从一堆纸里揪出一个对不上的数，扽魂殿上下没有人比他快。

\begin{center}\rule{0.5\linewidth}{0.5pt}\end{center}

当天夜里，大长老的殿上点了灯。

殿上只有三个人。test\_alpha0 站在石栏边上，背对着门。梅天狸站在殿当中。靠东边那根柱子底下，放着一张矮凳，凳子上坐着一个人，面前摊着一块油布，油布上摆满了零件。

那人五十来岁的模样，穿一件洗得发白的短褂，袖子挽到胳膊肘，手上沾着油。他一边听，一边拿一把细镊子，把一根头发丝粗的银线往一块铜片上的小孔里穿，穿了好几回，才穿过去。

扽魂殿的二长老，Smile Lee。

他今年二百三十四岁，九十九级的极限斗罗，大陆上唯一的十级魂导师，也是扽魂殿魂导一脉的总负责人。锻造一行里的人提起他，不叫名字，只叫一声神匠。

他早年拜在大长老门下。论辈分，梅天狸进门的时候，得管他叫一声师兄；论年纪，他跟一百九十一年前死在魂灵上的那位卢家前辈只差五岁，算是同一辈的人。那位前辈走的时候四十八岁，八十七级，儿子才十三。

梅天狸从三河镇的登记册说起，说到孙老先生，说到那张契，说到落枫中殿门口绕着走的人，说到那一百三十七张寻人纸，说到姚某人瘫在柱子底下，说到胡峰，说到韩家药材行十二年往西发了十二趟货，说到台底下那个女人。

她说得很平，一件一件，没有落下什么。

说到最后，她停了一下。

``四十年前，北陵那一趟是我去的。''她说，``我走了以后，再没去过。''

test\_alpha0 转过身来。

``殿里往各州巡察，是照着报上来的案数分的。''他说，``哪一州案子多，就往哪一州多派人。北陵往西那几州，册子上年年最干净，年年评上等。四十年，巡察一趟一趟都派去了东边。''

``帅气教要的就是这个。它把西边的册子做干净了，殿里的眼睛，自己就转到东边去了。''

``出事的多在分殿、中殿。可分殿、中殿头上有大殿，大殿头上有长老院；朝廷那边，城主府头上还有皇室的检查团。一层一层往上，每一层都该看见，每一层都没看见。这一漏，漏的是上头每一层，轮不到你一个人背。''

梅天狸抬起头。

``我平日不管这些。''test\_alpha0 说，``殿里的事有你们分着担，我只在出大事的时候露面。可五十年了，这件事一点一点攒成了大事，我坐在这座殿上，没看出来。''

``所以你那一份，我替你背一半。另一半你自己背着，背着它去干活。别背着它睡觉。''

梅天狸看着师父，过了一会儿，低下头，应了一声。

``那几个孩子呢？''test\_alpha0 问。

``都好。''梅天狸说，``姓霍的那个，把中殿对面墙上的寻人纸一张一张抄下来，一百三十七张，跟韩府的账本对了九天。卢家那个丫头去商帮翻税单，翻出来十二年的货一张都没验过。唐家兄妹跟着我跑了半个月，一句苦没叫过。''

``蔡须坤呢？''

``她一直跟在我后头，看我怎么问人。''梅天狸说，``有一回我问她看出什么了，她说，我问人的时候，先问人家家里几口人。''

test\_alpha0 愣了一下，然后笑了。

那是梅天狸这天夜里头一回看见师父笑。笑得很短，一下就收住了。

``一万年，''他说，``总算有人教她点正经东西了。''

矮凳那边，Smile Lee 把镊子放下了。

``真霜的卷，你拿到了？''

``拿到了。五十年的，明天一早送过来。''

``越往西越少。''Smile Lee 说，``真霜说的？''

``他说的。''

Smile Lee 点了点头，低头又拿起那块铜片，对着灯看了看。

``唐家那个印，你描下来了？''

梅天狸从袖子里拿出那张拓样，递过去。

Smile Lee 接过来，看了很久。一个圈，从左边起笔，绕了大半圈，到右边没合上。

``锻造这一行，打得好的铁匠，都有自己的印。''他说，``这个，我回头看看。''

他把拓样折好，揣进了短褂的口袋里，又拿起镊子。

``严家那个小子是怎么查出落枫不对的？''

``对账。''梅天狸说，``分殿年底直接报到大殿的案数，跟中殿往上转的对不上。''

``那就还用这个法子。''Smile Lee 说，``每一州的分殿，都有自己往上报的底子。中殿能改它往上转的，改不了分殿手里留的。一州一州地对，对不上的，就是有人伸了手。''

``西边光对账不够。''梅天狸说，``分殿要是也伸了手，底子一样是干净的。老百姓不敢报的，压根上不了账。落枫那一百三十七件，是霍家那孩子从中殿对面墙上的寻人纸里抄出来的。''

Smile Lee 手里的镊子停了一下。

``那就连墙一块儿对。''

``一共多少座中殿？''test\_alpha0 问。

``天斗这边，三百一十二座。''Smile Lee 说，``分殿四千出头。''

``对得完吗？''

``人对，对不完。''Smile Lee 把手里那块铜片举起来，``所以我做这个。''

铜片只有巴掌大，中间嵌着一块指甲盖大的晶石，四周一圈细得看不清的银纹。

``你在落枫那条巷子里摸到的那股魂力，阴的，黏的。''他对梅天狸说，``我这几天照着你信里写的，调了十几回。帅气教那套东西，抽了人的帅气往自己身上堆，堆久了，魂力里就带着那股味道。人走了，东西搬了，味道散得慢。''

``这块片子，碰上那股味道，会发黑。''

他把铜片放在油布上。

``我打算做两百块。''

``两百块片子，各顾各的，一州一州往回报，太慢。''他接着说，``还得有个母盘。两百块片子，哪一块发黑，母盘上都看得见。''

``母盘要几级？''test\_alpha0 问。

``九级。''

``用什么料？''

Smile Lee 手上停了一下。

``幽邃玄钢。''他说，``库里那一块。''

test\_alpha0 看了他一眼，没再问。

``两百块要多久？''

``一个月。''Smile Lee 说，``一个月以后，人下去。''

``派谁？''test\_alpha0 问。

``长老，还有各州大殿的殿主。''Smile Lee 说，``两个人一组，一个长老带一个殿主。不穿袍子，不挂徽记，不走公文，扮成走商的、赶考的、投亲的，一座中殿一座中殿地暗访。''

``殿主不查自己那一州。北陵的去查南边，南边的去查东边。自己管的地方，底下人都认得他的脸，也难免有几分情面。''

``落枫那座中殿的殿主坐了十五年，北陵派人查过一回，客客气气地回来了。明着去，人家就明着请你吃酒，账早做干净了等你。''

``暗着去，它来不及做。''

test\_alpha0 看着他：``你呢？''

``片子和母盘做完，交出去，我就回魂导院关门。''Smile Lee 说，``我一出天斗城，西边那一路上的人就都知道了。我不动，比我动管用。''

``再说，''他低下头，又去穿第二根银线，``母盘只管看哪儿黑了。黑了以后要拿人，拿的是魂圣，是魂斗罗，还有身上揣着遁符的。落枫那个瘦的，梅丫头都没追上。殿里的魂导器，还差得远。''

他说这话的时候，脸上一点笑都没有。

银线穿到一半，他的手停了一下。

一百六十八年前，他也这么蹲着，蹲在天州城东的城墙上，拆开一座叫砸裂了的防御魂导器，往里头换核心。那天天州城墙上的防御魂导器，砸裂了一大半。城门外那一剑劈下来的时候，他抬了一下头，手里的核心停在了半空。后来那块核心他按了两回才按到底，按完了，手心里全是汗。

他看了看自己的手。手上沾着油，没有汗。

他把银线捻直，接着往孔里穿。

test\_alpha0 看了他一会儿，没再说什么。

\begin{center}\rule{0.5\linewidth}{0.5pt}\end{center}

test\_alpha0 在石栏边站了很久。

``城主府的仵作，城主府的主簿，西门关卡的税吏，城西巡检司的巡检。''他说，``天狸，你在落枫台上处决的那二十六个人里，有几个是殿里的？''

梅天狸算了一下。

``姚某人，五个执事，两个镇上的分殿执事。八个。''

``剩下十八个呢？''

``韩家的人，帅气教的人，还有\ldots\ldots{}''梅天狸顿了一下，``城主府和关卡上的，是朝廷的人。''

``朝廷的人，你处决了。''

``他们认了罪。''

``我知道。''test\_alpha0 说，``我没说你处决错了。我是说，他们是朝廷的人。''

``殿里跟朝廷是两家。殿里管魂师，朝廷管地方，各管各的，有事一块儿办。你是殿里的长老，处决几个朝廷的吏，朝廷不会说什么。''

``可不说，不等于没记着。''

殿上静了一下。

``殿里的人下去，查殿里的人，谁也拦不住。''test\_alpha0 说，``可落枫那条路，一环扣一环。殿里的人压案子，城主府的人改文书，关卡上的人放货，巡检司的人拦人。少了哪一环都走不通。''

``殿里那几环，我们自己砍。朝廷那几环，得朝廷自己点头。''

``不然殿里的人走到哪一州，城主府面上客客气气，底下把账一烧，把人一藏，他们什么也摸不着。''

Smile Lee 头都没抬：``那就一块儿查。''

``一块儿查，得先一块儿坐下来。''test\_alpha0 说，``明天我去一趟皇宫。''

梅天狸看了师父一眼。

师父是良序之神的大供奉，平日不出这座殿，殿里的庶务都由各位长老和供奉分着管，皇宫那边，每年只在年节里去一回，坐一坐，喝一盏茶就回来。年节还没到，师父开口说要进宫，就是把这件事当成了大事。

Smile Lee 手里的镊子也停了一下，又接着往下穿。

``进宫的时候，跟陛下要一个人。''他说。

``Dispher？''test\_alpha0 说。

``是他。''Smile Lee 说，``殿里的人撒下去，朝廷那几环也得一块儿拔。朝廷那头，得有一个他这样的人。母盘交给他，他看得懂。''

test\_alpha0 没再问。那孩子他是看着长大的：在魂导院跟着 Smile 学了十几年，出师以后进了御史台，一查就是二十年的账，一个铜板都没漏过。九级魂导师，朝廷那边就他一个。

``殿里这头呢？''梅天狸问。

``真霜长老说，他要去。''

test\_alpha0 看了看梅天狸，又看了看 Smile Lee，点了点头。

``一个管卷的，一个查账的。''他说，``够了。''

``还有一件。''他看着梅天狸，``帅气教跟朝廷、跟殿里，是有一纸约的。''

梅天狸的眉头动了一下。

``那纸约签的时候，你还没入门。''test\_alpha0 说，``明天你跟我一块儿去，听一听。''

\begin{center}\rule{0.5\linewidth}{0.5pt}\end{center}

天斗皇宫的议事殿，大半年前宫变那一夜，让宫门前传过来的气浪震塌过半边屋顶。宫门前那个大坑早填平了，皇城东边那道让猪撞塌的城墙，也重新砌了起来。

屋顶已经修好了，梁是新换的，漆味还没散。殿里的地砖有几块颜色不一样，是后补的。

test\_alpha0 进殿的时候，皇帝从上首站了起来，下了两级台阶来迎。

天斗帝国跟扽魂殿，是两家。帝国管地方，管百姓，管赋税；扽魂殿管魂师。两家各立门户，大事上一块儿商量。可大陆上但凡能打的，十个里有九个在扽魂殿的册子上。帝国的皇帝见了扽魂殿的大长老，是要起身的。

大长老三百多岁了。上首那把椅子，他看着换过五个人坐。

他不管朝政，殿里的庶务也不管。平日里，天下的事是皇帝和宰相在办，殿里的事是各位长老在办，他坐在自己那座殿上，一年只出来一回，年节里进宫坐一坐。可国号改过，年号换过，先帝一个一个送走了，他一直在。新皇帝登基，头一件事是去他殿上走一趟；他点了头，上首那把椅子才算坐稳。

真到了要拿主意的时候，满朝的人，不管坐在哪一边，最后都会转过头去看他。

上首那把椅子旁边，另摆了一把，跟皇帝的并排，一样高。test\_alpha0 没去坐，只是朝皇帝点了一下头。

皇帝回到上首坐下。四十来岁，面色有些发黄，是宫变那回中的蚀骨散魂毒落下的病根，到现在还没养回来。他左手边坐着宰相，右手边是一位白须的老亲王，皇帝的叔祖，宫里论辈分最高的人。再往下，刑部尚书，户部尚书，禁军统领，一溜排开。

最末一把椅子上，坐着御史台的御史 Dispher。一身青袍，膝盖上摊着一本薄册子，手里捏着一支笔，从进殿起就在记，一句话没说。朝廷里的九级魂导师，只有他一个。

扽魂殿这边来了五个人。test\_alpha0，Smile Lee，梅天狸，还有从藏卷楼里下来的真霜长老。他一头白里泛蓝的头发，腰挺得笔直，一进殿就在 test\_alpha0 左手边那把椅子上坐下了，怀里抱着一摞卷宗。

第五个是个瘦高的中年人，一身青灰的袍子，腰里挂着一块铜牌，牌子上刻着一个''察''字。他进殿以后一直没坐，抱着胳膊站在 test\_alpha0 身后，脸上没什么表情。

扽魂殿检查团的总领事，川风斗罗。

test\_alpha0 没有绕弯子。他把落枫的事从头讲了一遍，讲得比梅天狸那天夜里还简短。讲完了，他让梅天狸把那二十六个人的名单呈上去。

名单在殿上传了一圈。传到户部尚书手里，他看了一眼西门关卡那个税吏的名字，脸色变了变。传到刑部尚书手里，他在''城主革职，押北陵受审''那一行上停了很久。

``落枫城主，是本部三年前考过的。''刑部尚书说，``考评是上等。''

``他的考评是谁写的？''Smile Lee 问。

刑部尚书没接话。

户部尚书清了清嗓子：``西门关卡那个税吏，户部这边的簿子上，十二年没出过差错。每年的税银都是足数上缴的。''

``足数。''川风斗罗在后头说了一句，``一车药材的税，一文不少。车上装的是什么，他没看过。''

户部尚书的脸红了一下，没再说话。

皇帝把名单拿过来，从头看到尾，放下了。

``大长老的意思是？''

``查。''test\_alpha0 说，``殿里的人，殿里派长老和大殿殿主下去查。朝廷的人，请陛下派人查。两边的人撒下去，得有一个人统着。查到谁，是殿里的归殿里，是朝廷的归朝廷，谁也不护着谁。''

``出事的多在分殿、中殿。''他接着说，``可分殿、中殿的年报，大殿看了几十年，没看出来。各州的城主府，皇室检查团三年走一趟，走了几十年，也没查出来。''

``底下伸手的要查，上头闭眼的也要查。北陵大殿要查，皇室检查团也要查。''

皇帝的眉头动了一下。坐在下首的一个中年官员，身上穿着皇室检查团的绯袍，低下了头。

``一州一州查？''

``一州一州查。''

皇帝的目光往下首扫了一圈，停在最末那把椅子上。

``Dispher。''

御史站了起来，册子还攥在手里。

``这件事，朝廷这头，你统着。''

``臣领旨。''Dispher 说。

``殿里这头，''test\_alpha0 说，``真霜统着。''

真霜长老抱着卷宗站起来，朝皇帝点了一下头。

皇帝也朝他点了点头。宫变那一夜，是这位老人带着八个孩子进的宫，骑着一头冰熊，在宫门前把胡帅缠了一整夜。

Dispher 转过身，朝他深深行了一礼。他在御史台查账的头几年，常往扽魂殿的藏卷楼跑，楼上七层那张长案，他借着坐过不少回。

``两头要是意见不合呢？''皇帝问。

``商量。''真霜长老说。

``商量不下呢？''

``魂师的事，听殿里的。地方上的事，听朝廷的。''test\_alpha0 说，``两边都沾的\ldots\ldots{}''

``先听我的，''真霜长老说，``再跟陛下报。''

殿上静了一下。宰相张了张嘴，看了看皇帝，又合上了。

皇帝想了想，点头：``就这么办。''

Dispher 坐回去，在册子上记了一行。

殿上静了一会儿。

宰相开口了。他六十多岁，说话慢条斯理。

``大长老，落枫一事，骇人听闻，该查，该办。''他说，``只是有一句话，老臣不得不问。''

``请讲。''

``帅气教与朝廷、与扽魂殿，是有约在先的。''宰相说，``五十五年前，就在这座殿上，先帝和大长老，还有帅气教那边，三家共签。帅气教只取魂兽，不犯人间。五十五年来，大体相安。''

``落枫一事，会不会是韩家私自勾结帅气教的几个人，胡作非为？若是如此，查办韩家、查办落枫便是。若一州一州地查，查到哪里都要牵扯帅气教，这纸约\ldots\ldots{}''

他没说完。

``这纸约，还算不算数？''

\begin{center}\rule{0.5\linewidth}{0.5pt}\end{center}

\chapter{第九十七章：从签订那天起}\label{ux7b2cux4e5dux5341ux4e03ux7ae0ux4eceux7b7eux8ba2ux90a3ux5929ux8d77}

川风斗罗笑了一声。

声音不大，可殿上太静，所有人都听见了。

宰相转过头看他。

``川风。''test\_alpha0 说。

``大长老，让我说两句。''川风斗罗说，``这两句话，我憋了二十年了。''

test\_alpha0 没拦他。

川风斗罗往前走了两步，走到殿当中，站定了。

``宰相大人说，大体相安。''他说，``我当了二十年检查团的总领事，这二十年里，检查团往西边派过六十一回人。''

``六十一回，四十三回报上来的东西，都跟落枫一个样。人丢了，案子压了，账对不上。有的是城主府压的，有的是殿里压的，有的两边一块儿压。''

``每一回报上去，上头批下来的都是同一句话：有约在先，不宜深究。''

他看着宰相。

``有约在先。''

``那纸约，我看过原件。''川风斗罗说，``帅气教只取魂兽，不犯人间。写得明明白白，三方的印押在底下，一个都不少。''

``可大人想过没有，帅气教为什么肯签？''

宰相没答。

``因为那一年，大长老跟胡赞在帅气堡门口打了第四回，打到一半停了。''川风斗罗说，``谁也吃不掉谁。再打下去，死的是城里的人。''

``明面上谁也奈何不了谁，才签的这纸约。''

``它签，是因为它打不进来。我们签，是因为我们打不进去。两边心里都清楚，这纸约不过是把刀插回鞘里，刀还在手上。''

``打了四回，胡赞也看明白了。正面打，它讨不到好。刀不能明着拔，那就换个法子用。''

``只有我们自己，把它当了真。''

他停了一下。

``这么多年来，帅气教答应我们只针对魂兽掠夺帅气，就是个彻头彻尾的骗局。''

``从签订的时候起，就没有遵守过。''

殿上没有一点声音。

``真霜长老管了六十年的卷。''川风斗罗说，``他那一架子卷宗，最底下一册，一指厚。最上头一册，一块砖厚。从哪一年开始往上涨的？''

他转过头，看着梅天狸。

梅天狸说：``五十五年前。''

川风斗罗转回来，看着宰相。

``就是签约那一年。''

宰相的手在膝盖上握了一下，又松开了。

禁军统领一直坐在最下首，这时候开口了。他四十出头，脸上一道疤，从眉骨一直划到耳根。

``宫变那一夜，''他说，``末将的人中了毒，站在二门上，刀拄着地，跟一排木头桩子一样。帅气教的人从末将眼皮子底下进进出出，末将一根指头都动不了。后来才听那几个孩子说，那一夜有人在外头喊'把瑞兽交出来'，喊的是帅气教的人。''

``毒解了以后，末将递了折子上来，说帅气教的人闯宫，该当如何。批回来的，也是四个字。''

他没说是哪四个字。殿上的人都知道。

宰相低下了头。

``它签了约，转过头就往殿里埋人，往城主府里埋人，往关卡上埋人。''川风斗罗说，``它不用打进来。它让我们自己的人，替它开门。''

``落枫压了一百三十七个人。胡家那个小孙子站在韩府的院子里说，让落枫的人怕中殿，比让他们怕帅气教好用。''

``大人，这句话您听明白了吗？''

``它要的从来就不是几罐帅气。它要的是有一天，天斗城里的老百姓看见扽魂殿的旗杆，看见朝廷的衙门，都绕着走。''

``到了那一天，它不用打，这座城自己就塌了。''

他说完，退回到 test\_alpha0 身后，又抱起了胳膊。

\begin{center}\rule{0.5\linewidth}{0.5pt}\end{center}

殿上很久没有人说话。

老亲王一直闭着眼睛，这时候睁开了。

``五十五年前签约那天，老夫在场。''他说，声音有点哑，``那时候老夫还不到四十。先帝签完了字，跟老夫说了一句话。''

``他说，这纸约，能换二十年太平，就算值了。''

``五十五年了。''

他看着皇帝。

``陛下，二十年太平是先帝换来的。剩下这三十五年，是咱们欠的。''

皇帝沉默了很久。

他转过头，看着梅天狸。

``梅长老，你在落枫的台上，有个百姓问你一句话。大长老方才说了，朕想听你自己说一遍。''

梅天狸站起来。

``她问，迟到的正义，也能算正义吗。''她说，``我说，不算。''

``为什么不算？''

``因为她闺女回不来了。''梅天狸说，``台上砍了多少颗脑袋，都回不来了。''

``那今天查的这些，算什么？''

梅天狸想了想。

``算还账。''她说，``账还完了，也不算正义。可不还，就连账都不算了。''

皇帝看着她，很久没说话。

``大长老。''他转回来，``那纸约，是撕，还是不撕？''

所有人都看向 test\_alpha0。

test\_alpha0 一直站在殿当中，没有坐下。他听完了川风的话，也听完了老亲王的话，脸上的神情一直没怎么变。

``那纸约，是我跟先帝一块儿签的。''他说。

``签之前，我跟先帝说过一句话。我说，帅气教的约，不知道哪一天就会撕，还请陛下留心。''

``先帝说，打了这么久，最后伤的都是百姓，总得给百姓一个交代。再说，他们答应了，只取魂兽的帅气。''

他停了一下。

``先帝说这话的时候，我想起了天州那一回。我原本是要进星斗大森林去谈共存的，还没动身，星斗大森林的兽潮先打到了天州城门口。打完了，城门外头躺着熊君。它身后那十万兽潮，连灰都没剩下。城里死的人，我到今天也没数清。''

``打来打去，最后倒下的，从来不是坐在桌边签字的人。''

``签字那天，胡赞坐在桌子对面。他落笔之前，看着我说了一句话。''

``他说，人类的共同对手是魂兽。你我不过是修炼的功法不一样。你和我再这么打下去，魂兽那边，说不定就缓过来了。''

test\_alpha0 停了一下。

``熊君扑上来的样子，我是亲眼见过的。他这句话，我那时候挑不出错。''

``我就没再说什么，签了。''

``现在回头看，他那句话里，只有一句是真的。''他说，``功法不一样。''

``那纸约，当年看着是挑不出错的。''

``从星斗大森林那一回以后，殿里限着魂兽。大森林外头的人死得太多了，殿里得先护着人。一百年下来，星斗大森林一圈一圈往回缩，魂兽聚的地方，散了大半。''

``帅气教说它只取魂兽，跟殿里走的，看着是一条路。签约那天，殿上没有一个人觉得不对。''

``可星斗大森林已经缩成那样了。它要的那么多，上哪儿去取？''

他停了一下。

``这一句，当年没有人问。我也没问。''

``川风说得对。只有我们自己，把它当了真。''

他转向皇帝。

``约是废纸。''他说，``可撕它，得挑时候。''

``今天在这座殿上撕了，明天帅气教就知道了。它埋在各州的那些人，一夜之间就会缩回去，缩得干干净净，Dispher 的人下去，什么都摸不着。''

``它要是觉得缩不回去，就会狗急跳墙。西边那几州的城，头一个遭殃。''

``所以不撕。''test\_alpha0 说，``约放着，让它以为我们还当真。Dispher 的人下去，不打旗号，暗访。查的是殿里和朝廷自己的人。查实一个，办一个，只当是查贪查渎，不提帅气教三个字。''

``等它埋的人拔得差不多了，那纸约，它自己就会撕。''

``到那一天，''他说，``扽魂殿站在前头。''

皇帝看着他，看了很久。

``刑部，户部。''皇帝说，``各出二十个人，归 Dispher 调遣。各州城主府，见 Dispher 的令，如见朕的旨意。谁关门，谁就不用再开了。''

``皇室检查团这几十年的案卷，一并交出来，让 Dispher 从头翻。''

Dispher 又站起来，朝皇帝行了一礼。

``殿里下去的长老和殿主，由真霜调度。''test\_alpha0 说，``要跟朝廷的人一块儿动的，他们两个商量着来。''

Dispher 转过身，朝 test\_alpha0 和真霜长老各行了一礼，坐下，在册子上又记了一行。

那个穿绯袍的官员站起来领旨，领完了，站在那儿半天没坐下。

刑部尚书和户部尚书站起来，领了旨。

宰相也站了起来。他朝 test\_alpha0 拱了拱手，又朝川风斗罗拱了拱手，什么也没说，坐了回去。

\begin{center}\rule{0.5\linewidth}{0.5pt}\end{center}

从皇宫出来，天已经黑了。

Smile Lee 在宫门口就跟他们分了手，说还有一百九十九块片子没做，头也不回地往扽魂殿那边走了。Dispher 和真霜长老跟梅天狸要了落枫的全部卷宗，说今夜就把头一批下去的人排出来，一块儿走了。川风斗罗跟在他后头，走了几步，又回过头来，朝 test\_alpha0 行了一礼。

``大长老，今天殿上那些话，我说得冲了。''

``不冲。''test\_alpha0 说，``二十年前就该有人这么说。''

川风斗罗站了一会儿，转身走了。

宫门外只剩下师徒两个。

梅天狸跟在师父后头，走了一段，开口了。

``师父，签约那天，先帝说完那句话，您真的就什么都没再说？''

test\_alpha0 走了几步，才说：``说了。''

``说了什么？''

``我跟先帝说，这纸约我签。哪天它撕了，头一个站出来的，也是我。''

梅天狸没再问。

两个人沿着宫墙往东走。宫墙很长，墙头上每隔一段挂着一盏灯，灯底下有禁军在巡。

``那几个孩子，''test\_alpha0 说，``现在到哪儿了？''

``算日子，该到海边了。''梅天狸说，``再过几天，就该上船了。''

test\_alpha0 在宫墙的拐角停下来，往东边看了一眼。东边是黑的，城墙外头就是官道，官道一直通到海边。

``唐家那个孩子，''他说，``你信里写，他说没人见过他爹死，在见着之前，他当他爹还活着。''

``是。''

test\_alpha0 点了点头。

``那一锤能把一个魂圣的胸口砸塌，''他说，``四十三级的魂宗，砸不出来。''

梅天狸愣了一下，停住了脚。

``师父？''

test\_alpha0 没有解释。他在拐角那儿又站了一会儿，转过身，往扽魂殿走去。

``让 Smile 做片子的时候，多做一块。''他说，``给那个孩子留着。''

走到殿门口，他又停了一下。

``明天起，我也关门了。''他说，``外头的事，有 Dispher。天塌下来之前，别来叫我。''

梅天狸应了一声。

\begin{center}\rule{0.5\linewidth}{0.5pt}\end{center}

那天夜里，扽魂殿魂导院地底下那间锻造室的炉火，一直烧到天明。

魂导院在扽魂殿的最北边，一座灰石头砌的方楼，地上三层，地下两层。殿里上上下下的魂导器，从图纸到出炉，最后都要过这座楼，也都要过二长老的手。

锻造室在最底下一层。屋子是圆的，四面墙都是整块的黑石，石面上刻满了隔火的纹路。当中一座炉子，炉膛比人还高，炉口封着三道铁门。

平日里，这间屋子只有二长老一个人进。今天门口挤了一堆魂导院的人，七级的、八级的都有，站在台阶上，谁也不敢往下走。

川风斗罗是被二长老叫下来的。他抱着真霜长老那五十册卷宗，找了个离炉子最远的角落，把卷宗按州分开，一州一摞，从东往西排。

Smile Lee 把短褂脱了，只穿一件单衣，站在炉子跟前。

他脚边放着一块料。

那块料只有两个拳头大，黑的，黑得发沉，灯光照上去，表面浮着一层极淡的青。放在地上的时候，脚底下的石板往下陷了半寸。

幽邃玄钢。

魂导院的库里只有这一块，是四十年前从大陆最北边的冰川底下凿出来的。四十年里，魂导院做过三回九级魂导器的图纸，三回都想动它，三回都被二长老拦下来了。

``这回舍得了？''川风斗罗在角落里问。

Smile Lee 没答。他把那块料夹起来，放进了炉膛。

炉子里的火本来是红的。料一进去，火就暗了下去，暗成一团黑红，像是被那块料吸走了什么。

Smile Lee 把手按在炉壁上。

他没放扽魂，身上也不见一丝魂力外泄。台阶上那些魂导师只看见炉壁上那一圈一圈隔火的纹路，从他手按着的地方开始亮，一道一道往外亮，亮到整座炉子都成了白的。

炉膛里的火从黑红变成橙，从橙变成白，从白变成一种看久了眼睛发疼的青。

那块料还是黑的。

烧了一个时辰，料的边上才泛出一线暗红。

Smile Lee 开了炉门，把它夹出来，放上砧子，拿起锤子。

锤子不大，比寻常铁匠的锤子还小一号，锤头是灰的，看不出是什么料。他第一锤落下去的时候，声音不响，闷闷的一下，像敲在一口很深的井上。

台阶上的人都缩了一下脖子。

第二锤，落在第一锤左边一指的地方。第三锤，落在两锤中间偏上。第四锤，又回到第一锤的位置，力道轻了一半。

川风斗罗看了一会儿，看不出门道，问：``你这锤点是怎么排的？''

``照着一张表排的。''Smile Lee 说，一边说一边落锤，``第几锤落在哪儿，由前面的锤落在哪儿推出来。前面错一锤，后面全错。''

``表是谁定的？''

``我自己推的。''Smile Lee 说，``推了三十年，推到两千零四十八锤，才推圆。''

``师父当年给我讲过一种表，一行推一行，推到后来，数长得比山还高。我拿它来排锤点，排了三十年，才排出一张不打架的。''

他落下一锤，又说：``那张表，是师父从良序之神那儿看来的。他看了两百多年，到今天也说不清，它那股劲儿到底是从哪儿来的。''

川风斗罗没听懂，也没再问。他低下头，接着分卷宗。

又落了十几锤，Smile Lee 忽然开口，像是说给自己听的。

``一百六十八年前，天州那一回，我九十八级，跟着师父落下去，在城墙里头守着城里的秩序，修城防的魂导器。城门外那一剑，我是隔着城墙看见的。''

``那时候我想，等我到了九十九级，总能看清一点。''

他把锤子举起来，停了一下，又落下去。

``如今到了。回头再看那一剑，还是看不见顶。''

锤声一下一下，闷闷的，不快，也不慢。

那块幽邃玄钢在砧子上一点一点变了形。起先是一块疙瘩，后来扁了，后来圆了，边上薄，中间厚，厚的地方鼓起来一圈，像一只倒扣的碗。

每落一锤，碗面上就多一道极细的纹，细得看不见，只有灯光斜着照过去的时候，才能看见一片密密麻麻的银光。

两千零四十八锤落完，天已经过了三更。

Smile Lee 放下锤子，拿起一根针。针是幽邃玄钢打的，比头发丝粗不了多少，是他三十年前拿边角料做的，一直没用过。

他开始往那只碗上刻法阵。

一层刻完，又在上头刻第二层。第二层的纹路压在第一层的纹路缝里，一根不碰一根。第三层压在第二层的缝里。

台阶上一个八级魂导师看着看着，腿一软，坐到了台阶上。

``九层。''他小声跟旁边的人说，``他要刻九层。''

``九级魂导器，不就是九层？''旁边一个年轻的问。

``别人刻九层，是九张图叠在一块儿。''那个八级魂导师说，``他刻的九层，是一张图。拆开哪一层，剩下八层都废。''

刻到第七层，Smile Lee 停了一下，从案子上拿起一块铜片，就是他在大长老殿上穿银线的那种，放在碗边上比了比，又接着刻。

刻到第九层的时候，台阶最上头那扇门缝里，透进来一点白。

天亮了。

Smile Lee 把针放下，把那只碗翻过来，放在长案当中。

碗面上什么都看不见了。九层法阵刻完，纹路反倒全隐进了钢里，只剩下一层极淡的青，跟没动过的时候一模一样。

他把那十二块做好的铜片，一块一块摆在碗的四周，围成一圈。

然后他把手按在碗上。

碗面上亮起了十二个小点，一个铜片一个，排成一圈，淡淡的，是白的。

``两百块片子做完，''他说，``这上头就是两百个点。哪一个点黑了，我就知道是哪一州，哪一座殿。''

``它叫什么？''川风斗罗问。

Smile Lee 想了想。

``照幽。''他说。

川风斗罗走过来，看着那只碗，看了很久。

``从哪儿开始？''

``这你得问他们。''Smile Lee 朝台阶上抬了抬下巴。

台阶上的魂导师们让开了一条道，有几个认出了走在后头那个，小声说了一句：``九级的那个，御史台的。''

前头下来的是真霜长老，扶着墙，一级一级往下走，走得很慢。Dispher 跟在他后头半步，还是那身青袍，腋下夹着那本薄册子，册子里夹满了纸条。两个人一夜没睡，眼睛却都很清醒。

``师父。''Dispher 在长案前站定，叫了一声。

Smile Lee 嗯了一声，头都没抬。

真霜长老看了看那只碗，又看了看碗边那十二块铜片。

``北陵。''他说，``梅长老刚走过一遍，路熟。从北陵往西，一州一州往西查。''

``头一组谁去？''

``梅长老，带南边的一个殿主，''真霜长老说，``去北陵西边那一州。落枫的事刚办完，北陵认得她的人多了，她不在北陵露面。北陵那边，换两个生面孔去。''

川风斗罗挑了一下眉毛：``总殿的长老，上一代的圣女，也换便装？''

``换。''真霜长老说，``她那张脸，西边认得的人不多，可她那身长老袍子，谁都认得。袍子一穿，人还没进城，城里就有人往中殿报信了。''

``便装出去，不为别的，就为了让通风报信的人来不及报。''

``严殿主往东去，查别家的中殿。''

川风斗罗点了点头：``检查团呢？''

``检查团不动。''真霜长老说，``你们一动，满大陆都知道了。你们照旧往西派人，照旧报上来，照旧让上头批'有约在先'。让它看着，什么都没变。''

川风斗罗愣了一下，然后慢慢笑了。

``朝廷那头呢？''他问。

``刑部户部那四十个人，拆开了掺进每一组里，一组两三个。''Dispher 说，``扮伙计，扮车夫，扮账房先生。殿里的人盯中殿，他们盯城主府、关卡、巡检司。''

``城主府那边，我先不发令。令一发，账就烧了。等片子黑了，人拿住了，令再下去。''

真霜长老点了点头。

``查到头呢？''川风斗罗问。

真霜长老和 Dispher 都没答，一起看向 Smile Lee。

Smile Lee 把那只碗用一块油布包起来，放到长案最里头。

``一个月以后，连着两百块片子，一块儿交给你们。母盘放在 Dispher 那儿，他看得懂。''他说，``查到头，就是帅气堡。''

Dispher 点了点头，在册子上记了一笔。

川风斗罗没说话。他回头看了一眼角落里那一摞一摞的卷宗。越往西，摞得越矮。最东边那一摞，比他坐着还高；最西边那几摞，薄得立不住，歪在地上。

检查团往西边派过六十一回人。他知道那几摞底下压着的是什么。

真霜长老、Dispher 和川风斗罗抱着卷宗上了台阶。台阶上的魂导师们也散了，一个一个回头看，谁也不说话。

Smile Lee 没送。他拿起火钳，把炉膛里暗下去的火又拨旺了。还有一百八十八块片子。片子做完了，还有拿人用的东西要做。

锻造室的门，从里面关上了。

\begin{center}\rule{0.5\linewidth}{0.5pt}\end{center}

大长老的殿门，是同一个时辰关上的。

殿里只点着一盏灯。石台上摊着的，还是两个多月前刻完的那卷 iblp，末尾那一行 PTO(Z₃)，在灯底下泛着一点青。旁边压着另一卷，边角已经磨得发白，是九十五年前刻完的 ω-Y。

test\_alpha0 把两卷都卷起来，放回架子上。

石台空了。他从架子最底下抽出一块新玉简，摆在正当中。

良序。

他年轻的时候，有人问过他，这两个字到底是什么意思。他那时候想了很久，给了一句最笨的话：不管从哪儿往下走，走多少步，总有走到底的那一天。一条往下走的路，不会没完没了。

ω-Y 是这样，iblp 也是这样。他花了一百六十八年，把这两条路一步一步走到了底。

可路不止这两条。

他闭上眼睛。

两百多年前，他刚接下大供奉的位子，头一回向良序之神借力。借来的是 I3。托着那一层往下看，他看见了 Laver 表：两格的，四格的，八格的\ldots\ldots 一张比一张大，一张套着一张往上长。每一张表的第一行都是一个周期，周期也一张比一张长，可长得极慢，慢到要翻到一张大得没法想的表，周期才肯从十六长到三十二。

祂让他看清了：在 ZFC+I3 底下，Laver 表是良序的。

看清了，握不住。他只好照着它削，削出了 iblp。

再后来，他又往回添。把一张一张的表一根一根捻起来，像捻纱线一样，织成一整匹。

Laver 表纱，LTY。

魂力、肉身、灵魂、修为，说到底都是 Laver 图案织出来的。iblp 看的是一根线，LTY 看的是整匹布。

这些年，祂借给他的从 I3 换成了 I1。托着 I1 往下看，他又看清了一件事：在 ZFC+I1 底下，LTY 也是良序的。

结论就摆在那儿，是祂给的。可它凭什么良序，Laver 表那股力量到底从哪儿来，他看了几十年，还是看不懂。

借来的结论，用得上，看不懂。ω-Y 和 iblp，他都是一步一步自己走到底才看懂的。LTY 比这两条都长，也比这两条都深，也得这么走。

test\_alpha0 睁开眼，拿起刻刀，在玉简第一行刻下了三个字母。

LTY。

刻完了，他在底下又刻了一行小字：

一条往下走的路，总有尽头。

灯芯短了一截。他没去剪，就着那点光，刻下了第二行。

\begin{center}\rule{0.5\linewidth}{0.5pt}\end{center}

魂导院地底下，Smile Lee 用了二十六天，把两百块片子做完了。

第二十七天一早，真霜长老和 Dispher 一块儿下来取。母盘''照幽''上亮起两百个白点，密密地排成一片，像一张摊开的天斗地图。Dispher 抱着母盘，真霜长老抱着那两百块片子的清单，一前一后上了台阶。

``师父，走了。''Dispher 在台阶上回头说了一句。

Smile Lee 摆了摆手。

门关上以后，Smile Lee 在长案上铺开一张新图纸。

纸上画的是一只环，比手镯大一圈。环上刻着九层法阵，最里头那一层留着一道口子，还没画完。

落枫那个瘦供奉，是捏碎了一张遁符跑的。遁符一碎，人就从原地抽走，连梅丫头都追不上。要拿这样的人，得有一样东西，在他捏碎遁符的那一瞬间，把他身边一丈之内的路全锁死。

他把这只环叫作''锁遁''。

料还没找齐，法阵还差最里头那一层。Smile Lee 盯着那道没画完的口子看了很久，拿起笔，又放下，又拿起来。

炉子里的火，一直没灭。

天斗城的西门外，天刚亮，几个穿粗布衣裳的人各自牵着马，混在赶早集的人群里出了城。有的扮成走商，有的扮成投亲，有一个女人头上包着蓝布头巾，看着二十出头，走路很轻。

谁也没往她身上多看一眼。

\begin{center}\rule{0.5\linewidth}{0.5pt}\end{center}

\chapter{第九十八章：望潮镇}\label{ux7b2cux4e5dux5341ux516bux7ae0ux671bux6f6eux9547}

从落枫出来，往东南走了整整一个月，就闻见了海。

先是风变了，潮乎乎的，带着一股咸腥。再往前走半天，路两边的田没了，换成一片一片晒盐的滩。滩上的人赤着脚，拿木耙子把盐往一堆拢，太阳底下白得晃眼。

奥斯卡站在坡上看了半天，说：``这么多盐，够腌多少香肠。''

没人理他。

翻过最后一道土坡，海就在底下了。

灰蓝灰蓝的一片，一直铺到天边，铺到看不见的地方。浪一层一层往岸上推，推到沙滩上碎成白沫，退回去，又推上来。

小舞站在坡顶，半天没动。

``哥，我头一回见海。''

``我也是。''唐三说。

霍雨浩走到小舞旁边，也没说话，陪她站着。风把小舞的头发吹到他脸上，他伸手拨开，拨完了，手就没放下来，搭在她肩上。

蔡须坤站在最前头，看着东边。

东边什么也没有，只有海。可她看得很专注，像是在看一样隔着很远、却认得的东西。

``队长？''卢本韦走过去。

``没什么。''蔡须坤说，``风里有股味道。''

``什么味道？''

蔡须坤想了想，没答上来，摇了摇头。

蔡兴坤落在最后头。她肩上那只黑子鸡，从闻见海味那一刻起就不太对劲。平时它蹲在她肩上，眯着一只眼睛，半天动一下。今天它一直伸着脖子，脑袋朝着东边，一动不动，连毛都不抖。

``你怎么了？''蔡兴坤小声问它。

黑子鸡没理她。

\begin{center}\rule{0.5\linewidth}{0.5pt}\end{center}

坡底下是个渔镇，叫望潮镇，百来户人家，房子都是石头垒的，屋顶压着一排一排的石块，怕被海风掀了。镇子东头是个小码头，停着二三十条渔船，大的不过能坐十来个人。

进了镇子，头一样扑过来的不是鱼腥味，是鸡叫。

家家户户都养鸡。院墙上蹲着鸡，屋顶的石块上蹲着鸡，码头的缆桩上也蹲着一只，缩着脖子，眯着眼睛看海。石头院子里一群一群地刨食，见了生人也不躲。镇上的鸡比镇上的人还多。

奥斯卡看得直咽口水：``这地方的鸡腿，得便宜成什么样。''

``海边的人，养鸡做什么？''霍雨浩问。

没人答得上来。打鱼的人家，按说该养鸭，鸭子不怕水。

蔡兴坤肩上的黑子鸡，一进镇子，就把脑袋缩回来了。院墙上、屋顶上那些鸡，一只一只都转过头来看它，看了一会儿，又都转回去，朝着东边。

黑子鸡是蔡兴坤的扽魂，神级。别人的扽魂放出来要耗魂力，用完就收；她这只从六岁觉醒那天起就没收回去过，吃饭睡觉都蹲在她肩上，看着跟一只真鸡没什么两样。镇上这么多鸡，它是头一回把脑袋缩回去。

这是大陆的最东南角。再往东，就没有路了。

他们在镇上唯一一家客栈落了脚。客栈老板是个独眼的老头，收了钱，把钥匙往柜台上一拍，多一句话都没有。

夜里，八个人挤在一间屋里，把手头的东西凑了一遍。

蓬莱岛的事，大长老早就交代过。宫变之前，他在扽魂殿的修炼堂给他们讲 BMS，讲到快收尾的那几天，把蓬莱岛也一并说了：从大陆东南角出海，往正东六百里；两百二十多年前，他刚坐上大长老的位子没几年，去过一回；岛上主事的，是大供奉最菜萌新，跟他是一样的身份。最菜萌新认得他，也认得卢承业。

``报我的名号也行。''大长老当时说，``报你曾祖的，比报我的好使。他在岛上住过三年，我只去过一回。''

八个人那时候都在，都记着。

``往哪儿走，说了。''奥斯卡掰着手指头，``到了找谁，报谁的名号，也说了。怎么过去，一个字没说。''

``他说这种事他说了不算，得我们自己撞上。''卢本韦说。

``商帮那边呢？''

``问过了。''卢本韦说，``总会的海图我让人翻了个遍，往东画到三百里就断了，三百里以外，图上是空的。商帮的船从来不往那边走。路太远，浪太大，跑一趟挣的，不够赔一条船。''

霍雨浩翻开笔记，找到一页，念道：``往东三百里，海图止。再三百里，有岛。''

``六百里是真的。''唐三说，``怎么走过去，是另一回事。''

没有人接话。

卢本韦把她那本笔记拿出来，翻到曾祖补的那几十页里。

一百三十六年前，卢承业一个人去了蓬莱岛，待了三年。回来以后，他在这本笔记里补了几十页。补的都是魂灵，怎么合，怎么养，怎么不烧死人。关于那座岛本身，他几乎一个字都没写。

卢本韦一页一页往后翻，翻到那一页，那一句''查一件事，先查它为什么被人埋起来''，停了一下，又接着往后翻。

在那一页的背面，靠着装订线的地方，有一行很小的字。小到她翻了十几年笔记，从来没注意到。

她把油灯挪过来，凑近了看。

字是曾祖的字，写得很急，墨也淡：

出东南角，正东六百里。雾不开，船不靠。等它来接。

六百里，跟大长老说的一里不差。

``等谁来接？''小舞凑过来问。

卢本韦摇了摇头。

屋子里安静下来。窗外的浪一下一下拍着码头，拍得很有耐心。

\begin{center}\rule{0.5\linewidth}{0.5pt}\end{center}

第二天一早，八个人分头去打听。

唐三和奥斯卡去了码头。

码头上的渔民正在收网，一条一条船靠过来，鱼倒在木板上，噼里啪啦地跳。奥斯卡蹲在一条船边上，递过去两根香肠，跟船老大套了半天近乎，才开口问：

``大叔，往东六百里，有个岛吗？''

船老大正在理网的手停了一下。

``蓬莱岛。''他说，``你们找蓬莱岛。''

他没问找去干什么，抬手往东边一指。

``在那边。镇上三岁的孩子都知道在那边。''

``您去过吗？''

``没有。''船老大说，``往东一百里，是咱们打鱼的地方。过了两百里，浪就起来了，一浪比一浪高，高过桅杆。咱们这种船，走不到一半就得翻。''

``那就没人去过？''

``有。老水手里头有去过的。''船老大说，``听他们讲，快到的地方有一堵雾，雾开了，才看得见岛。我爹那一辈，镇上去过的能数出五六个。到我这一辈，一只手数得满，都老了，出不动了。''

他把香肠还给奥斯卡，转身又去理他的网了。

唐三这才注意到，船头拴着一只竹笼，笼里蹲着一只芦花鸡，脑袋朝着东边。再往码头上看，一条一条船，船头都有这么一只笼子。

``大叔，打鱼还带鸡？''

``老辈传下来的规矩。''船老大头都没抬，``出海带只鸡。鸡叫了才开船，鸡不叫，天再好也不走。''

``为什么？''

``敬坤神。''

``坤神？''

船老大这回抬起头，看了唐三一眼，像是看一个不知道太阳打哪边出来的人。

``海上的神。''他说，``你们内陆来的不知道。海边的人家，家家都供。''

``坤神管什么？''

``管鸡。''船老大说，``鸡叫，就是坤神点了头。我爹这么干，我爷爷也这么干。照着干的，都回来了。''

唐三想起昨天进镇的时候，霍雨浩问的那一句。海边的人养鸡做什么。

现在有答案了。

霍雨浩和小舞去了镇子西头。那里住着几个不出海的老渔民，腿脚不好了，天天坐在墙根底下晒太阳，补网，说闲话。

小舞搬了个小板凳坐在他们跟前，听他们说了一上午。

说的东西七零八落。

蓬莱岛在东边，这一样谁都知道，老头们说起来，手都往一个方向指。岛上的人一万年没怎么上过岸，这一样也谁都知道。咱们不过去，他们不过来，一万年就这么过来了。

说到坤神，几个老头说得最齐。

坤神是打蓬莱岛上起来的。坤神的本事都在鸡身上。鸡知道风，知道浪，知道哪天出海回得来，哪天回不来。海边的人养鸡，是跟坤神借一双眼睛。

``坤神有多大？''小舞问。

一个缺了门牙的老头想了想，说：``北边有个猪神，你们听说过吧？''

小舞愣了一下，点点头。

``跟猪神一个辈分。''老头说，``老辈人都这么说。''

霍雨浩手里的笔停了。他抬起头，看了小舞一眼。小舞也在看他。

两个人都没说话。

再往下，就一人一个说法了。有的说蓬莱岛是会走的，今年在东边，明年就往北去了。有的说岛上的人不会老，他爷爷的爷爷见过一个从岛上回来的，走的时候二十岁，回来的时候还是二十岁，可家里人都死光了。有的说岛根本就不在海上，在雾里，雾散了，岛就没了。

霍雨浩一条一条记下来，记了三页，越记眉头皱得越紧。

``都是听来的。''他跟小舞说，``没有一个是自己见过的。''

``有一个。''小舞说。

她朝墙根底下努了努嘴。最边上坐着一个老婆婆，一上午都没开口，就那么补网，脚边卧着三四只老母鸡。别人说什么，她都不接，可说到''会走的岛''的时候，她手里的梭子停了一下。

小舞走过去，蹲下来。

``婆婆，您见过吗？''

老婆婆没抬头，接着补网。补了很久，才说了一句：

``不是岛会走。是它想让谁看见，谁才看得见。''

``那它想让谁看见？''

老婆婆终于抬起头来，看了小舞一眼，又看了看站在后面的霍雨浩。

``反正不是你们俩。''她说。

\begin{center}\rule{0.5\linewidth}{0.5pt}\end{center}

卢本韦去了镇上的商帮分号。分号小得可怜，就一个管事，一个伙计，平日里收收海货。

管事看了卢家的印，又看了看她，没动。十六岁的姑娘，揣着总会的印，跑到大陆最东南角的小分号来翻账，换谁都得掂量掂量。

卢本韦没多说，把扽魂放了出来。

屋当中多了一头母猪。通身白里透粉，个头比寻常的猪大出一圈，四只蹄子稳稳踩在地上。它一出来，分号里那股海货的腥气都淡了，伙计手里的算盘珠子停在半道，拨不下去。

母猪扽魂，大陆上本来就没几头。神级的，只出在卢家。卢本韦出门在外，放一回扽魂，比亮卢家的印还管用。印能仿，扽魂仿不了。

管事从椅子上站了起来，站得太急，带翻了椅子。他没顾上扶，转身进了里屋，把分号里所有的簿子都搬了出来。

簿子上什么都没有。

往东三百里以外，一笔买卖都没有做过。

诺当陪着她翻了一上午，一页一页地翻。翻到最后一本，诺当忽然伸出一根指头，按在一行字上。

那是一笔很小的账。二十多年前，分号卖出去一条小船，买主没留名字，只在簿子上画了一个记号。

记号是一只鸡。

卢本韦盯着那只鸡看了很久。

\begin{center}\rule{0.5\linewidth}{0.5pt}\end{center}

蔡须坤和蔡兴坤哪儿都没去。

她们在海边坐了一整天。

蔡兴坤肩上的黑子鸡，从早上起就一直朝着东边。中午的时候，它从她肩上跳下来，跳到沙滩上，一步一步往海里走。走到浪打湿了它的爪子，它停住了，就站在那儿，朝着东边。

``回来。''蔡兴坤喊它。

它没回来。

蔡须坤看着那只鸡，看了很久。

``它在等。''她说。

``等什么？''

``不知道。''蔡须坤说，``我也有点像在等。''

太阳落山的时候，海面上起了一层薄雾，从东边慢慢漫过来，漫到离岸几里地的地方，停住了，不动了，像一堵墙。

黑子鸡朝着那堵雾墙，叫了一声。

叫声不大，可整片沙滩都安静了。浪还在拍，可好像离得很远。

雾墙里，有什么东西，亮了一下。

\begin{center}\rule{0.5\linewidth}{0.5pt}\end{center}

那天夜里，八个人把一天打听来的东西凑到了一块儿。

霍雨浩念到坤神那一段，念到''跟猪神一个辈分''，停了一下。

屋里的人都看卢本韦。

卢本韦坐在床沿上，那本笔记搁在膝盖上，没翻开。

``我们家那段话，我三岁就会背。''她说，``一万年前，人退到山脉以南，魂兽守着星斗大森林。一片猪场里走出来一个人，两族共举，信仰化神。''

``那段话里，没有第二个神。''

没有人接话。

``大长老跟我说过，我们家那段话没有错，错在太短。''卢本韦低头看着那本笔记，``现在看，又短了一截。''

过了一会儿，霍雨浩开口了。

``大长老说，最菜萌新跟他是一样的身份。''

``大长老是良序之神的大供奉。''唐三说。

霍雨浩低头看了看笔记上那两个字。坤神。

没有人再往下说。

蔡兴坤肩上的黑子鸡蹲在那儿，脑袋朝着东边，一晚上没动过。

敲门声响了。敲了三下，很轻。

客栈的独眼老板提着一壶热水进来，把水壶放在桌角，没走。他在门口站定，用那只好眼睛把八个人挨个看了一遍，最后停在卢本韦身上，腰弯了下去。

``下午分号的伙计上我这儿喝茶，说今天有位姑娘，在分号里放出来一头母猪。''他说，``是您？''

卢本韦点了点头。

老板的腰又弯低了一些。

``那就对上了。卢姑娘，白天多有怠慢。''

他没往屋里走，就站在门口，两只手垂着。

``两个多月前，有位先生来过镇上。一个人来的，穿得像个管账的，袖子上还沾着饲料。他在码头上蹲了一整天，跟那几个老水手说话。说的什么，我没听着，就看见他们一个一个把手往东边指。''

``第二天早上，他没坐船，也没跟人说一声，自己往海里走了。''

``往海里走？''奥斯卡问。

``踩着浪走的。''老板说，``那天浪不小，两边的浪一人多高。他脚底下那一道，平得像铺了一条石板路，他走过去小半个时辰，那一道才合上。镇上的人全跑到码头上看，没有一个敢出声。''

``我在海边六十多年，南来北往的魂师见过不少。能叫海给他让路的，就见过这一个。''

``第三天傍晚，他又踩着浪回来了，鞋都没湿。在我这儿又住了一夜。走之前跟我说，后头会来一拨人，里头有位姓卢的姑娘，扽魂是头母猪。让我等见着了，替他捎几句话。''

``母猪扽魂，我没见过。今天镇上出了一头。''

卢本韦站了起来。

``他让您捎什么话？''

老板想了想，一句一句往外说，说得很慢，像是怕说错一个字。

``头一句，六百里，没错。岛在雾后头，他看见轮廓了。''

``第二句，他看了一眼就回来了，没往跟前去。''

``第三句，''老板顿了一下，``剩下的路，让你们自己走。''

说完，他又弯了弯腰。

``几位这几天的房钱，小店不收了。''

卢本韦要开口，他摆了摆手。

``卢家的人住我这儿，是我这小店的脸面。''

他退出去，把门轻轻带上，下楼去了。

屋子里很久没有人说话。

``干爷爷。''卢本韦说。

``你干爷爷以前去过蓬莱岛？''小舞问。

``没有。''卢本韦说，``干爷爷是曾祖收的干儿子。曾祖那本笔记他从头读到尾，可那座岛，他一回都没上过。我们家去过蓬莱岛的，一直只有曾祖一个。''

她看着窗外那片黑下来的海。

``这一回，他也只看了一眼。''

\begin{center}\rule{0.5\linewidth}{0.5pt}\end{center}

第二天一早，独眼老板亲自把早饭端了上来，一样一样摆好，退出去的时候把门带上了。

八个人又凑到了一块儿。

``曾祖写的是，等它来接。''卢本韦说。

``干爷爷说的是，剩下的路，让咱们自己走。''唐三说。

``大长老说的也是这个，得我们自己撞上。''蔡须坤说，``先把船备下。''

这话没人反驳。九十八级的超级斗罗能踩着浪走，一群十六岁的孩子不能。六百里海路，总得有条船。

卢本韦又去了商帮分号。管事一听是要租船往东，脸都白了。可昨天那头母猪还在他眼前晃，他没敢说不，咬着牙把码头上最结实的一条船指给了她：两根桅杆，能坐二十来个人，船底新刷过桐油，是分号自己的船，平时拉海货用。

``船有了。''卢本韦回来说，``可谁会开？''

八个人互相看了一圈。

唐三会打铁，霍雨浩会记账，奥斯卡会烤香肠，诺当力气大，能拉帆，可帆该往哪边拉，他不知道。卢家的生意做遍了大陆，卢本韦坐过的船都是有人开的。小舞和蔡兴坤连海都是头一回见。

``雇一个。''蔡须坤说。

\begin{center}\rule{0.5\linewidth}{0.5pt}\end{center}

雇人比租船难。

码头上的渔民，一听是去蓬莱岛，没有一个点头的。给多少钱都不去。有个年轻的犹豫了一下，被他娘从后头一把拽走了，拽出去老远还在骂。

问到晌午，问到了码头最北头。

那儿单独泊着一条小船，船身旧得发黑。船头刻着一只鸡，刻得很粗，年头久了，叫海水泡得只剩个轮廓。卢本韦的眼睛在那只鸡上停了一下。

船头还拴着一只鸡笼。笼里蹲着一只芦花鸡，毛都秃了一半，看着比船还老。船上坐着一个老头，头发胡子全白了，皮肤晒得跟船板一个颜色，正低着头补帆。

船老大说的那一只手数得满的老水手里头，还能出海的，就剩他一个。

奥斯卡把来意说了。

老头没抬头，也没停手。

``蓬莱岛，我去过三回。''他说，``头一回跟我爹，后两回是我自己。三回都回来了。''

``上过岛吗？''唐三问。

``上过一回。''老头说，``两回在雾外头等了三天，雾没开，掉头回来了。一回雾开了，船靠上去，我在滩上站了一炷香。''

``岛上什么样？''奥斯卡问。

老头手里的针穿过去，拉紧了，才说：``跟这儿差不多。鸡多。''

他抬起头，看了卢本韦一眼。

``两个多月前，你们家那个管账的，也在我这条船上坐过，问的跟你们差不多。''他说，``我跟他说，雾不开，船不靠。他听完，自己走着去了。''

雾不开，船不靠。曾祖笔记背面那一行字里，也是这六个字。

再问，老头就不说了。

``那这回\ldots\ldots{}''

``这回去不去，不是我说了算。''老头朝船头那只鸡笼抬了抬下巴，``它叫了，我就去。它不叫，你们给我一座金山也没用。''

芦花鸡缩在笼子里，眼睛半睁半闭，一点动静都没有。

蔡须坤就是在这个时候，觉出不对的。

\begin{center}\rule{0.5\linewidth}{0.5pt}\end{center}

起先只是胸口发闷。

她以为是海风吹多了，往后退了一步，想站到背风的地方去。可那股闷没散，反倒往下沉，沉进了血里。

血热了。

不烫，一跳一跳的，跳得很齐。这股热在她血里睡了十六年，她从来没觉出来过。这会儿它醒了，每跳一下，都朝着东边。

东边是海。昨天傍晚，那堵雾墙就是从那边漫过来的。

``队长？''卢本韦第一个发现，伸手扶住她。

几步外，蔡兴坤也扶住了缆桩。她肩上的黑子鸡炸开了一身毛，脖子伸得笔直，朝着东边。

``队长，''蔡兴坤按着胸口，脸白了一下，``我心口热。''

蔡须坤想回她一句，张开嘴，没说出来。

因为血底下还有一样东西，也醒了。

比血更深。深到她这十六年，一次都没碰到过底。它没有形状，没有方向，说不上是冷是热，是一团混混沌沌的东西。血往东涌，它就往回拽；血往上烧，它就往下压。

两股力量在她身子里撞在一起，撞在正当中。

正当中是本源。帝皇瑞兽的本源在她身体最深处，一万年来一直是安静的，像一口井。血和混沌每撞一下，那口井就晃一下。井晃一下，她眼前就黑一下。

她的膝盖一软，跪在了码头的木板上。

``须坤！''

唐三冲过来，和卢本韦一左一右架住她。霍雨浩一把扯开她的袖子，她的手腕上，一道一道金色的纹路正从皮肤底下往外透，透到一半，又被什么灰黑色的东西从底下吞回去，吞回去，又透出来。

``魂力乱了。''霍雨浩说，声音在抖，``两股，在打架。''

诺当张开两只大手，想往她背上按，按下去之前又停住了，不敢碰。

奥斯卡慌得直往怀里掏香肠，掏出来了，又不知道该往哪儿塞。

蔡兴坤攥着缆桩，想过来，腿软了一下，没迈开步子。她的袖口底下也透出几道金纹，很淡，一闪一闪的。她咬着牙站着，没倒。

蔡须坤的嘴角渗出一道血。

她体内那两股力量越撞越狠。血往外冲，混沌往下拽，本源夹在中间，晃得一下比一下厉害。她觉得自己站在一道悬崖边上，脚底下的石头一块一块往下掉。

就在这时候，她身后有人把她托住了。

\begin{center}\rule{0.5\linewidth}{0.5pt}\end{center}

码头上所有人都看见了。

一个半透明的人影，从蔡须坤背后浮了出来。

个子不高，肩膀很宽，脸绷得紧紧的。十四岁的样子，跟两年前一模一样。他单膝跪在她身后，两只手按在她背上，一动不动。

他的手按下去的地方，一层很淡的光铺开来。那光不亮，不往上冲，也不往下压，就那么稳稳地铺着，一层一层地铺，铺得纹丝不动。

血撞过来，撞在那层光上，停住了。

混沌拽过来，拽在那层光上，也停住了。

那口井不晃了。两股力量中间那道悬崖，被他一层一层垫平了。

蔡须坤手腕上的纹路慢慢暗下去，金色的，灰黑色的，都退回了皮肤底下。她的呼吸一点一点匀了。

缆桩那边，蔡兴坤袖口底下那几道淡金的纹，也退了。她松开缆桩，大口喘了几口气。

码头上没有一点声音。

奥斯卡手里的香肠掉在了木板上。

``牢\ldots\ldots{}''他张了张嘴，第二个字没说出来。

诺当往前走了一步，又停住了。唐三架着蔡须坤的那只手，一下子攥紧了。小舞捂住了嘴，霍雨浩摘下眼镜，又戴上，又摘下来。

卢本韦没有动。她看着那个人影，眼眶一点一点红了。

人影没有看他们。他一直看着蔡须坤，看着她的后脑勺，等她的呼吸彻底匀了，才开口。

声音很轻，像是隔着很远的地方传过来的。

``我说过，在底下给你垫着。''

蔡须坤没有回头。

她跪在那儿，两只手撑在木板上，肩膀一抖一抖的，眼泪一颗一颗砸在木板上。

她一直都知道他在。

两年前，星斗大森林外头那条路上，他一层一层往里烧，烧到最后，肉身什么都没剩下。可就在他烧完的那一刻，她身体最深处的本源动了一下，把他最后那一点魂，收了进来。

一收就是两年。

本源护着他，他就在里头，一声不吭，跟着她一块儿修。她破一层，他跟着破一层。她从没跟任何人说过。他只剩一缕魂，往后也只能是一缕魂。说了，他们会想见他，会盼他现一回身。可现一回身要耗他多少，她心里没数。

这两年，她一回都没叫过他。

``你出来干什么。''她终于开口，声音哑得不成样子，``你才修到哪儿，出来一回，得耗多少\ldots\ldots{}''

``推快了会错。''牢大说，``可这回不推，你就掉下去了。''

蔡须坤说不出话来了。

两年前，就是这个人，让他们带着她往后退，自己站在那条路上。九个人的队伍，活下来八个，是他一个人换的。现在他连身子都没有了，还是这样，一声不吭地站到她身后，把她托住。

她哭得喘不上气。

牢大等她哭了一阵，才又开口。

``底下那一团，你先别去碰。''他说，``我在那口井边上坐了两年，天天挨着它。它比你的血老，比这口井也老。''

``那是什么？''蔡须坤哑着嗓子问。

``混沌之力。''牢大说，``叫什么，我知道。是什么，我也说不上来。''

人影的轮廓开始淡了，从边上一点一点往里褪。

他终于抬起头，朝码头上那几个人看了一眼。看到奥斯卡的时候，他停了一下。

``香肠还没吃完？''

奥斯卡的眼泪一下子就下来了。卢本韦走到他边上，没说话，把他的手攥住了。

``没\ldots\ldots 没吃完。''奥斯卡说，``给你留着。''

牢大好像笑了一下，很淡，看不太清。

然后他低下头，最后看了蔡须坤一眼，身影一点一点沉下去，沉回了她的背里，沉回了那口井一样的本源里。

码头上的风又吹起来了。

蔡须坤跪在那儿，很久都没有起身。

蔡兴坤走过来，在她旁边蹲下，没说话，伸手扶住了她的胳膊。黑子鸡那一身炸开的毛，慢慢顺了下去。

\begin{center}\rule{0.5\linewidth}{0.5pt}\end{center}

码头最北头，那条旧船上，补帆的老头手里的针不知道什么时候停了。

他看着这边，看了很久。看的是跪在地上的那两个姑娘，和那只黑鸡。

船头的鸡笼里，那只秃了一半毛的芦花鸡，不知什么时候站了起来，伸长了脖子，朝着东边。

\begin{center}\rule{0.5\linewidth}{0.5pt}\end{center}

\chapter{第九十九章：六百里}\label{ux7b2cux4e5dux5341ux4e5dux7ae0ux516dux767eux91cc}

芦花鸡是第二天天不亮叫的。

头一声哑，第二声亮，第三声拖得很长。三声叫完，望潮镇的鸡跟着叫了起来，从码头叫到镇子西头，一家接一家，把整个镇子都叫醒了。

八个人下楼的时候，老头已经在码头上了。他没上自己那条旧船，抱着鸡笼，站在商帮那条两桅船的跳板前头。

``它叫了。''他说，``走。''

码头上站了不少人。鸡一叫，镇上的人就出来了，站在码头上，不说话，就那么看着。独眼老板也来了，挑着两只筐，一只装着干粮，一只装着腌鱼，搁在跳板边上，冲卢本韦弯了弯腰，没多说什么。

墙根底下补网的那个老婆婆也在，站在人堆最后头。

小舞上跳板的时候，她开了口。

``丫头。''

小舞回过头。

老婆婆没看她。她看着小舞身后，看着蔡须坤，看着蔡兴坤，看着蔡兴坤肩上那只黑鸡。

``我那天说的那俩，''她说，``在你后头。''

说完，她转身走了，拄着一根竹竿，一步一步往镇子里去，没再回头。

蔡须坤走在最后。昨天跪过的地方，木板上还有一小块暗色。她从那儿跨过去，上了船。

\begin{center}\rule{0.5\linewidth}{0.5pt}\end{center}

头一百里，海是平的。

打鱼的船一条一条落在后头，到晌午，最后一条也看不见了。四面都是水，水连着天，分不出哪儿是边。

奥斯卡趴在船帮上吐了两回。卢本韦在他边上站着，一只手揪着他的后领子，怕他一头栽下去。吐完了，她把水壶塞到他手里。

奥斯卡漱了口，擦擦嘴，爬起来给大家烤香肠。头一根递给了老头，第二根递给了卢本韦。

``大爷，您怎么称呼？''

``叫老头就行。''老头接过香肠，``镇上都这么叫。''

霍雨浩坐在船尾，膝盖上摊着笔记。每过一个时辰，他问老头一回：往哪边走，风从哪边来，船走得多快。问完了，记一行。

``记这个做什么？''小舞蹲在他旁边看。

``海图画到三百里就断了。''霍雨浩说，``三百里以后的，咱们自己画。''

小舞看了一会儿，看不懂，把头靠在他胳膊上，看海。

卢本韦坐到老头边上。

``您那条旧船，是从分号买的？''

``二十三年了。''老头说。

``账上没留名字，画了只鸡。''

``我不识字。''老头咬了一口香肠，``画只鸡，管事就认了。镇上谁不认得我家的鸡。''

他朝船头的鸡笼抬了抬下巴。

``这只跟了我十一年。它娘跟过我，它姥姥跟过我爹。往上数，一辈一辈，都是这一窝。''

``您信坤神？''小舞在那边问。

老头想了想。

``信不信的，''他说，``我爹跟着鸡走，回来了。我跟着鸡走，回来了三回。''

``那您爹\ldots\ldots{}''

``我爹最后一回出海，鸡没叫。''老头说，``那天天好，他说天这么好，不等了。''

他没往下说。小舞也没再问。

蔡须坤站在船头。

血里那股热一直在，从出港就在，一跳一跳，跳得很稳。往东走一里，就热一分。底下那一团也在，沉着，没动。她按着船帮站着，没跟任何人说。

蔡兴坤走过来，站在她旁边。

``你也热？''蔡须坤问。

``一点点。''蔡兴坤说。

过了一会儿，蔡兴坤又说：``队长，黑子鸡从我六岁觉醒那天起，睡觉脑袋就朝东。搬过多少回住处，都朝东。我一直以为它认生。''

``现在呢？''

``现在觉得，它是在认路。''

两个人就那么站着，谁也没再开口。黑子鸡蹲在蔡兴坤肩上，脖子伸得笔直，朝着东边。

太阳落进海里以前，把整片海烧成了一层红。

奥斯卡坐在船尾，手里攥着一根香肠。香肠拿油纸包着，包了三层，他攥了一下午，一口没动。

诺当在他旁边坐下。

``不饿？''

``这根不是我的。''奥斯卡说，``给牢大留的。''

``他吃得着？''

``吃不着也留着。''

诺当没说话，把那只大手搁在奥斯卡后脑勺上，按了一下。

卢本韦从他们身后经过，停了一下，把奥斯卡被风吹翻的领子翻好，才往船头走。她在蔡须坤边上站了一会儿。

``两年？''她问。

``两年。''蔡须坤说。

``你一个人扛了两年。''

蔡须坤没接这句。她看着那层红一点一点暗下去，过了很久，才说：``他没让我一个人扛过。''

\begin{center}\rule{0.5\linewidth}{0.5pt}\end{center}

过了两百里，浪起来了。

先是船头一沉一起，沉得深，起得高。到了傍晚，整条船被一个浪托起来，托到半空，再扔下去。船底砸在水面上，砰的一声，砸得人牙根发酸。

天黑下来，起了风。

老头把着舵，两条腿叉开，钉在甲板上一样。

``不会水的下舱！''他喊，``留两个力气大的拉帆！''

诺当没等第二声，已经攥住了帆绳。唐三跟过去，攥住了另一根。

浪一个比一个高。有一个从左边打过来，高过了桅杆，在船头上方悬了一下，才砸下来。甲板上的水没到膝盖，退下去的时候卷走了一只木桶，连声响都没有。

蔡须坤从舱口探出半个身子，手已经抬起来了。血里那股热跟着往上一涌，底下那一团也跟着一翻，她眼前黑了一下。

一只手从后头攥住了她的胳膊。是蔡兴坤。

``队长，下去。''

蔡须坤站了一会儿，把手放下，退回了舱里。

舱里黑，一股桐油味混着海水的腥。霍雨浩把笔记揣在怀里，外头裹了两层油布。船一歪，他整个人撞在舱壁上，两只手还是护着胸口。

半夜，舵断了。

咔的一声，很脆。老头手里的舵把一下子轻了，整条船横过来，被浪推着打转。

``舵销！''老头吼，``舵销断了！''

唐三松开帆绳，扑到船尾。

舵柱和舵叶之间，原本穿着一根拇指粗的铁销，这会儿断成了两截，一截掉进了海里，一截还卡在孔里。没有这根销，舵叶就是一块挂在船尾的木板，吃不上力。

``有备用的没有？''

``没有！''老头说，``谁出海带这个！''

唐三回头看了一眼。舱口挂着一盏马灯，灯底下盘着锚链，一环一环，每一环都有拇指粗。

``诺当！''

诺当把帆绳往桅杆上绕了两圈，扣死，冲过来。

``舵叶你给我扶住，别让它晃。''

诺当趴到船尾，两只手伸下去，攥住了舵叶的边。一个浪打过来，他整个人被拍在船板上，手没松。

唐三扯过锚链，找准一环，把昊天锤放了出来。

锤子黑沉沉的，在马灯底下不反光。没有炉子，也没有火。他把魂力压在锤面上，把那一环链子搁在锚爪上，一锤砸了下去。

一锤，那一环扁一分。两锤，再扁一分。

船在打转，甲板一会儿往左斜，一会儿往右斜。卢本韦从舱口爬出来，把马灯摘下来，举在他手边。风把灯吹得乱晃，她拿身子挡着。

唐三没抬头。链环被砸开，砸直，砸成一根，再一点一点砸圆，砸细。

他想起玉小刚那句话。二流有二流的练法，锤子不在好坏，在落得准不准。

船歪得最厉害的时候，锤子也没落偏过。

第四十七锤，铁销成了。他攥着那根砸得发烫的铁，爬到船尾，对准孔，一锤送了进去。

舵叶一沉，吃上了力。

``好了！''

老头把舵往回一扳，船头慢慢转回来，重新顶住了浪。

天快亮的时候，主帆被风撕开一道口子，一截帆绳缠死在桅杆顶上。小舞把鞋一蹬，顺着桅杆爬了上去。她在桅杆顶上被风甩得晃了三晃，一刀割断了那截绳子，又顺着下来。霍雨浩在底下伸着两只手，一直没放下来。

风是天亮以后停的。

奥斯卡从舱里钻出来，往每个人嘴里塞了一根香肠。香肠是他的扽魂，吃下去，灵气一点一点往回补。诺当两只手掌全磨破了，攥着香肠的地方印出两道血。小舞撕了自己一截袖子，给他一圈一圈缠上。

老头蹲在船尾，把那根新铁销看了半天，拿指头敲了敲。

``这手活，''他说，``比镇上的铁匠强。''

``我爹是打铁的。''唐三说。

老头点点头，没再问。

唐三坐在甲板上，低头看自己的手。虎口震裂了，血已经干了，黑红黑红的一道。他看了一会儿，把手攥上了。

小舞在他旁边坐下来。

``哥，你砸铁那会儿，跟爹一个样。''

唐三没抬头。``你还记得爹打铁什么样？''

``记得。''小舞说，``不说话。一锤一锤的，谁叫他都不应。你刚才也是。''

唐三没吭声。他把昊天锤又放出来，拿袖子把锤面上的海水一点一点擦干了，才收回去。

\begin{center}\rule{0.5\linewidth}{0.5pt}\end{center}

第三天早上，霍雨浩在笔记上画了一道横线。

``三百里。''他说，``海图到这儿就没了。''

横线往下，是一整页空白。他在空白那页的左边点了一个点，写上''望潮镇''，从点上往右画了一条线，画到纸边上，打了个问号。

海变了颜色。先前是灰蓝，过了三百里，越走越深，深到发黑。海上看不见鸟。天是阴的，白天看不见太阳，夜里看不见星星。

老头不看天。他看鸡。

芦花鸡蹲在船头的笼子里，脑袋朝着哪边，他就把船头往哪边扳。黑子鸡蹲在蔡兴坤肩上，脑袋也朝着那边。两只鸡，一只秃了半身毛，一只黑得发亮，朝着同一个方向，一整天没转过头。

第四天夜里，船底下有东西过去了。

先是海面鼓起来一道，从船左边一直鼓到看不见的地方。然后整条船轻轻往上一抬，又落了下去。马灯照下去，水里是一片黑，比船大出不知多少倍，一片一片的鳞，每一片都有磨盘大，慢慢地从船底下往右滑。

诺当的手摸向了腰。老头一把按住了他。

``别动。''老头的声音压得很低，``海里的东西。它不惹你，你别惹它。''

两只鸡都没叫。芦花鸡把脑袋埋进了翅膀底下，黑子鸡缩着脖子，眼睛眯成一条缝。

那片黑滑了一炷香的工夫，才滑完。海面又平了。

奥斯卡憋了半天的一口气，这才吐出来。

``那是什么？''

``不知道。''老头说，``我爹那回也见过。他说，在这片海里，咱们是客，它们是主人家。''

卢本韦看着那片平下去的海面，半天没说话。

``一万年前那一炷香，''她说，``谈的是星斗大森林。人退到山脉以南，魂兽守着大森林。''

``海里呢？''唐三问。

``没人跟海里谈过。''

那天后半夜，霍雨浩没睡。马灯挂在桅杆上，他坐在灯底下，一笔一笔补白天的航程。

小舞披着毯子出来，在他边上坐下。

``你记这么细，出了海也记。''她说，``风往哪边刮都记。''

霍雨浩的笔停了一下。

``我爹娘出事那年，没人记。''他说，``后来我去查，那一年镇上发生过什么，谁来过，谁走了，什么都查不到。就跟那一年没有过一样。''

他把那一行写完，吹了吹墨。

``我不想再有哪一年，是查不到的。''

小舞没说话，把毯子分了一半，搭到他肩上。

霍雨浩还是一个时辰记一行。

四百里。四百五十里。五百里。

蔡须坤血里的热，一天比一天厉害，底下那一团也一天比一天沉。她吃得少，话也少，大半时间坐在船头，闭着眼睛。

背上那一处，一直是温的。

第五天凌晨，别人都睡了。蔡须坤一个人坐在船头，把一只手背到身后，按在那一处上。

``你在吗？''她说。声音小得只有自己听得见。

没有人答。手底下那一处温了一下，比刚才温一点，又慢慢落了回去。

她把手收回来，抱着膝盖，一直坐到天亮。

\begin{center}\rule{0.5\linewidth}{0.5pt}\end{center}

第五天傍晚，雾墙到了。

霍雨浩笔记上，这一行写的是五百六十里。

雾从海面上立起来，一直立到天上，白里发灰，左右两头望不到边。它不飘，不散，不往前漫，也不往后退，就立在那儿。离得近了，能看见雾的边沿，齐得像刀切过。

老头落了帆，抛了锚。

两只鸡都站了起来。芦花鸡在笼子里，黑子鸡在蔡兴坤肩上，脖子伸得笔直，冲着那堵雾，一声不叫。

蔡须坤扶着船帮，血里那股热一下子涌到了顶。旁边的蔡兴坤也按住了心口。两个人对看了一眼，谁都没说什么。

``到了。''老头说，``就在这儿等。''

``等多久？''奥斯卡问。

``头一回跟我爹，等了三天，没开。''老头说，``第二回，等了三天，没开。第三回，第二天早上开的。''

那天夜里，老头坐在船尾抽烟，说起了他那第三回。

``雾是早上开的。开了一道口子，刚好过一条船。我把船摇进去，摇了小半天，看见岛了。''

``上去了？''小舞问。

``上去了。在滩上站了一炷香。''

``有人吗？''

``没人。''老头说，``满滩的鸡，都看着我。看了一炷香，没一只过来。我就知道了，等的不是我。我就回来了。''

他把烟锅在船帮上磕了磕。

第二天，雾没开。

唐三拆了一块舱板，在上头刻了一道箭头，箭头朝东，拿绳子拴好，扔进了雾里。绳子一截一截放出去，放到二十丈，停住。他再一截一截往回拉。

木板拉回来了。箭头朝西。

他又扔了一回。拉回来，箭头还是朝西。

霍雨浩在笔记上记了一笔，没说话。

老头在船尾看着，烟锅里的火一明一暗。

``我爹说过，''他说，``那里头，东西南北不作数。''

第三天，雾还是没开。

第三天夜里，八个人坐在甲板上。马灯挂在桅杆上，照着一圈人，照不到那堵雾。

``水还够九天，干粮够十二天。''霍雨浩说，``再等几天，等得起。''

``那就再等等？''奥斯卡说。

``曾祖在这儿等过。''卢本韦说，``他等来了。''

她顿了一下。

``干爷爷在这儿看了一眼，回去了。他说，剩下的路，让咱们自己走。''

``大长老也没说等。''唐三说。

蔡须坤一直没开口。她坐在最靠船头的地方，面朝着雾，脸在灯底下白得发青。

过了很久，她说：``老人家上过岛，满滩的鸡看着他，没一只过来。等的是谁，咱们在外头猜不出来。进去了，就知道了。''

她站起来。

``不等了。''

老头听了，没拦，也没点头。他又装了一袋烟，抽完了，才开口。

``这条船不进雾。''他说，``我爹说过，进雾的船，他没见回来过一条。''

他站起来，走到船尾，把拴在那儿的一条舢板解了下来。

``挤一挤，坐得下八个。桨四副。''他说，``我在这儿等你们三天。三天没出来，我就回去，跟镇上说，有八个娃娃进了雾。''

舢板落水的时候，他从怀里摸出一个小布袋，塞给蔡兴坤。袋子里是一把谷子。

``给它的。''他朝黑子鸡抬了抬下巴，``进去以后，它要是不肯往前看，别逼它。''

\begin{center}\rule{0.5\linewidth}{0.5pt}\end{center}

舢板进雾的时候，天刚蒙蒙亮。

雾里没有声音。

桨划下去，听不见水响。浪没了，海面平得像一块磨过的铁。前后左右都是白的，坐在船尾看不见船头的人，把手伸出去，看不见自己的指头。

霍雨浩从怀里摸出一块木片，拿刀在上头刻了个''一''字，扔进水里。

``记个号。''

四个人划桨，诺当和唐三一边，卢本韦和奥斯卡一边，划一下，数一下，两边一样使劲。黑子鸡蹲在船头，脑袋朝着哪边，他们就往哪边划。

划了一个时辰，船头咚的一声，碰上了一样东西。

诺当伸手捞起来。一块木片，上头刻着一个''一''。

没有人说话。

霍雨浩又刻了一块，刻''二''，扔下去。这回不看鸡，数着桨，一下一下，笔直地往前划。

又一个时辰。诺当捞上来两块木片，一块刻着''一''，一块刻着''二''。

黑子鸡蹲在船头，脑袋一会儿朝这边，一会儿朝那边，转了一圈，又转一圈，最后缩了回去。

蔡兴坤想起老头那句话，没逼它。她把它抱下来搁在膝盖上，从布袋里抓了一小撮谷子，摊在手心。黑子鸡低头啄了两口，不转了，缩在她怀里，眼睛眯着。

``方向靠不住。''霍雨浩说。他的声音在雾里闷闷的，像隔着一层棉被，``在这雾里，往前划跟往回划，是一回事。''

``那还有什么靠得住？''奥斯卡问。

没人答得上来。

诺当一直没说话。他把桨在手里掂了掂。

``我划。''他说，``你们说往哪儿，我就往哪儿。''

蔡须坤坐在船中间，两只手攥着船帮，指节都白了。进了雾，她身子里那两股力量又顶上了，血往前涌，底下那一团往回拽，比在码头上那一回还狠。她咬着牙，一声没出。

蔡兴坤握住了她的手。

握了一会儿，蔡兴坤忽然抬起头。

``热。''她说。

``什么？''

``我心口热。''蔡兴坤说，``刚才那一圈，热了，又凉了，又热了。往那两块木片那儿划的时候，是凉的。''

霍雨浩一下子坐直了。

``队长也\ldots\ldots{}''

``队长那边太重。''蔡兴坤说，``她一用，底下那个就跟着起来。''

她转过头，看着蔡须坤。

``队长，你别用了。我来。''

蔡须坤看着她，看了很久，点了点头。

霍雨浩把笔记翻到空白的一页，写了一行字，举起来给大家看：

只往热处走。

``划一百下，停一回。''他说，``热了，接着往前。没热，退回来，换个方向再试。''

``要是一直试不出来呢？''奥斯卡问。

``试得出来。''唐三说。

他把桨横在膝盖上。

``大长老讲良序那天说过，不管从哪儿往下走，走多少步，总有走到底的那一天。''他说，``咱们离岛还有几十里。每一回都比上一回热，就是每一回都近了一截。近一截，少一截，少不了没完没了。''

``也绕不了圈子。''霍雨浩说，``绕回原地，热就得跟原来一样。每一回都比上一回热，就回不到原地。''

没有人再问。

蔡兴坤闭上眼睛，一只手按在心口，一只手攥着蔡须坤。

四支桨落下去。霍雨浩数着，一，二，三\ldots\ldots 数到一百，停。

``热了。''

接着划。

``热了。''

``\ldots\ldots 没有。''

退回来，船头往左偏一点，再划。

``热了。''

就这么一百下一百下地往前挪。热了，往前；没热，退回来，换个方向。霍雨浩每一回都在笔记上记一笔，热了画一道，没热画个圈。一页记满了，翻过去，接着记。

有一段，蔡兴坤一连说了七回''没有''。往左偏，没有。往右偏，没有。退回来再往左，还是没有。霍雨浩笔记上一连画了七个圈。

奥斯卡的桨停在半空。

``是不是\ldots\ldots 到头了？''

``到不了头。''霍雨浩说，``没到岛，就还有更热的地方。是咱们偏得不够。''

卢本韦伸过手去，按着奥斯卡的手腕，把那支桨按回了水里。

``划。''

第八回，船头往右扳了大半个弯。

``热了。''

雾里分不出时辰。胳膊从酸到麻，从麻到没知觉。诺当手上缠的那截袖子早就透了，血顺着桨杆往下流，他换了个握法，接着划。小舞把卢本韦替下来，卢本韦歇了一会儿，又把奥斯卡替下来。四支桨，一直没停过。

蔡兴坤脸上的汗一层一层往下淌。她袖口底下那几道淡金的纹又透了出来，一闪一闪，一回比一回亮。

``热了。''她说。嗓子已经哑了。

有一回，她说完''热了''，身子往旁边一歪。小舞扔下桨，从后头挪过去，让她靠在自己身上。

``还行？''

``还行。''蔡兴坤闭着眼睛，``接着划。''

蔡须坤攥着蔡兴坤的手，攥得很紧。她身子里那两股力量还在顶，她一直没松劲，没让底下那一团起来。

推快了会错。她在心里把这句话过了一遍。不推，也不松，就这么顶着。

背上那一处，一直是温的。

霍雨浩的笔记记到第三页的时候，雾薄了。

先是能看见船头了。后来能看见船头前头的水。水面上起了皱，一道一道的，是浪。

然后有了声音。浪拍船帮的声音，风的声音。再后来，从前头很远的地方，传过来一声鸡叫。

黑子鸡在船头，脖子伸得笔直，嘴闭着。叫的是前头。

又一声。又一声。一声接一声，远远近近，像是一整个镇子的鸡，同时醒了。

蔡兴坤睁开眼睛，嘴角动了一下，身子一软，靠在了小舞身上。

舢板冲出了雾。

\begin{center}\rule{0.5\linewidth}{0.5pt}\end{center}

雾墙外头，两桅船上，老头正抽着这一天的第三袋烟。

船头笼子里的芦花鸡忽然站起来，冲着雾叫了一声。

老头把烟锅放下，看着那堵雾，看了很久。雾没开，一丝缝都没开。

他又把烟锅捡了起来，没起锚。

\begin{center}\rule{0.5\linewidth}{0.5pt}\end{center}

天是亮的。太阳刚从海面上出来，照在一座岛上。

岛不大，一眼望得到两头。一座青山从当中拱起来，山坡上一层一层都是田，田埂边上是石头垒的房子，屋顶压着石块，跟望潮镇的一模一样。海滩是白的，白沙上一群一群的鸡在刨食，礁石上站着公鸡，一只比一只叫得响。

滩头立着一座石头牌坊。牌坊顶上蹲着一只石头刻的鸡，刻得很粗，跟老头那条旧船船头上那只一个样，叫海风磨得只剩个轮廓。

牌坊底下没有人。整片滩上，看不见一个人。可山坡上那些石头房子，烟囱里都有烟，一缕一缕往上飘。风从岛上吹下来，带着一股晒过的谷子味。

八个人回过头去。

雾还在身后立着，一道墙，一丝缝都没开。

卢本韦看着那堵雾，想起曾祖笔记背面那一行小字。等它来接。

他们没等。

霍雨浩低头看了看笔记。第三页最后一笔，是一道。他在底下补了一行字：

六百里。

奥斯卡伸手摸了摸怀里。那根拿油纸包了三层的香肠还在，压扁了一点，没湿。

舢板的底蹭上了沙子。

满滩的鸡，一只一只停了下来，抬起头，看着船上的人。看了一会儿，都看向了船中间那两个姑娘。

蔡须坤扶着蔡兴坤，先下了船。

脚踩上沙子的那一下，她血里那股热，静了。底下那一团，也静了。

十六年来，头一回这么静。

\begin{center}\rule{0.5\linewidth}{0.5pt}\end{center}

\chapter{第一百章：灯还亮着}\label{ux7b2cux4e00ux767eux7ae0ux706fux8fd8ux4eaeux7740}

梅天狸出城的第十天夜里，藏卷楼七层的灯还亮着。

这盏灯亮了六十年。天斗城里起早的人都知道，扽魂殿后头那座青砖楼，最高那扇窄窗里头，夜夜有一点光。光底下坐着一个管卷宗的老人，很少下楼。

今天长案上没摊卷宗，摊着的是一只碗。

照幽。碗面上两百个白点，这十天一直在动，一点一点往西挪。东边那几十个挨得紧，越往西越稀，到最西边，零零落落只剩七八个，像撒在沙地上的几粒米。

Dispher 坐在长案对面，那本薄册子摊在膝盖上，一页一页往后翻。

``头一批下去的十几组，走得最远的，昨天过了北陵大殿的地界。''他说，``梅长老那一组在最前头。照这个走法，一州一州往西推，推到最西边，''他翻到最后一页，``得三个月。''

``三个月。''真霜长老说。

``东边的中殿一座一座出事，消息是捂不住的。''Dispher 说，``最西边那几州的人不傻。东边倒一座，他们就烧一本账。等咱们的人推过去，账烧干净了，人也挪走了。''

真霜长老没说话。他伸出一根指头，点在碗沿最西边那几个白点上。

``苍岭州。''他说。

``苍岭州。''Dispher 说，``挨着界河，过了河就是帅气教的地界。这五十年，苍岭州往上报的案子，一共十九件。''

``东边一座中殿，一年就不止十九件。''

``所以得先到。''Dispher 合上册子，``不等东边推过去，咱们直接去最西边。他们等的是从东边一路查过来的人，咱们从他们背后进去。''

真霜长老看了他一会儿。

``咱们？''

``朝廷那头我统着，殿里那头您统着。''Dispher 说，``苍岭州要是跟我想的一样，城主府和殿里是一伙的，那片子一黑，当场就得拿人。拿殿里的人，得您点头；拿朝廷的人，得我点头。两个人都在，才来得及。''

``一个御史，一个管卷的，一块儿离了天斗城，''真霜长老说，``第二天满城都知道了。''

``所以不能让人知道。''

Dispher 从册子里抽出一张纸条，放在长案上。

``御史台那边，我递了折子，说要闭门核检查团这几十年的案卷，三个月不见客。我那间屋子的灯，书童每天夜里会点上，点到天亮。''

他抬头看了一眼窗户。

``您这盏灯亮了六十年。再亮三个月，没人会多看一眼。''

真霜长老也看了一眼那盏灯。灯芯剪得很短，火苗不高。

``楼下的执事跟了我三十年。''他说，``灯，他会点。''

``那就只剩一件事。''Dispher 说。

``什么事？''

``脸。''

Dispher 指了指自己。

``我在御史台二十年，西边那几州的城主府，十个里有七个见过我。您的脸见过的人少，可一百三十六岁，满头白里泛蓝的头发，腰板挺成这样，走到哪儿都扎眼。''

``还有腿。''真霜长老说，``从天斗城到苍岭州，两千来里路。我这两条腿，走不到一半。''

他想了想，从长案底下抽出一册薄薄的卷宗，翻开看了一眼，又合上。

``殿里有个人，''他说，``脸和腿，都归他管。''

\begin{center}\rule{0.5\linewidth}{0.5pt}\end{center}

那人是半夜来的。

楼下的执事守了三十年的门，开门的时候，看见门外站着一个年轻人，以为是城南哪家书院的学生走错了地方。

``公子，夜里不借书。''

年轻人笑了一下，没解释，递过去一张折好的字条。执事看了字条上那两个字，手一抖，把门开大了，退到一边。

年轻人上了七层。

Dispher 站起来的时候，先看见的是头发。

银色的，很长，用一根素白的带子松松束在脑后，一直垂过腰。楼梯口的风吹进来，发梢动了一下，在灯底下泛着一层很淡的光。

他个子不算高，比 Dispher 矮小半个头，身形清瘦，一身月白的长衫，袖口洗得发软。腰里没挂玉，没挂牌，只别着一把合起来的折扇。

脸生得极好看。眉毛细长，眼尾微微往下垂，鼻梁挺，嘴唇薄，下颌收得很干净。线条是柔的，柔得有几分像女子，可站在那儿，又一点不叫人觉得弱。看人的时候，眼睛里带着三分笑，像书院里那种教了好些年书、脾气极好的先生。

看着二十五六岁。

``蝶神。''真霜长老说。

Dispher 愣了一下。

扽魂殿供奉殿里有一位蝶神斗罗，姓赵，名雪谣，九十五级，扽魂是光明女神蝶。这个名字他在卷宗里见过不止一回。可几十年来，殿里的人提起这位供奉，说法从来对不上。有人说是个老太太，有人说是个络腮胡子的大汉，有人说就是城南开药铺的郎中。

``真霜长老。''年轻人朝长案拱了拱手，声音不高，很温和，``好些年没收到您的条子了。''

``三十二年。''真霜长老说。

``三十二年。''年轻人笑了笑，转向 Dispher，``这位是御史台的 Dispher 大人。九级魂导师，Smile 师兄的徒弟。''

``您认得我？''

``你在城南我那间铺子里，抓过三回治咳嗽的药。''年轻人说，``最后一回是去年腊月。你嫌药苦，多要了一包糖。''

Dispher 张了张嘴，没说出话来。

城南那家药铺他记得。坐堂的是个五十来岁的郎中，左脸上一块青胎记，说话带着南边口音，抓药的时候手很稳。

``那个郎中\ldots\ldots{}''

``是我。''

真霜长老在一旁说：``他今年一百一十四了。''

Dispher 又看了一眼那张脸。灯底下，那张脸比他还年轻二十岁。

\begin{center}\rule{0.5\linewidth}{0.5pt}\end{center}

真霜长老把苍岭州的事说了一遍。说得很短，跟他翻账一样，一句是一句。

赵雪谣听完，没先问苍岭州。他走到长案后头，蹲下来，把一只手按在真霜长老的膝盖上。

``您这两条腿，''他说，``三十二年前给您看的时候，还能一口气走到城门口。''

``现在走到楼下，要歇两回。''

赵雪谣没再说什么。他背后亮了一下。

一只蝶从他肩后展开来。翅膀很大，展开比他整个人还宽，薄得透光。底色是极淡的蓝，往边上一点一点转成紫，翅脉是银的，跟他的头发一个颜色。翅膀扇了一下，满屋子落下一层细细的光，像灯底下飘着的灰，可那灰是亮的。

光明女神蝶。

他的境界停在 Y(1,3,5)，比真霜长老低一格。可殿里管治疗的那一脉，从上到下，没有谁的手比他稳。

光落在真霜长老的膝盖上，一点一点往里渗。

一炷香以后，真霜长老扶着长案站起来，在屋里走了一圈。走到窗户底下，又走回来，没扶墙。

``三十年没这么松快过。''他说。

``撑得到苍岭州，撑得回来。''赵雪谣说，``路上每三天，我给您续一回。''

他转过身，看着 Dispher。

``该你了。''

Dispher 还没回过神来，那只蝶已经从他头顶上掠过去了。翅膀扇了一下，落下来的光换了颜色，从亮的变成灰的，细细地扑在他脸上，凉丝丝的，像落了一层霜。

他觉得脸上的皮紧了一下，又松开了。

赵雪谣从袖子里摸出一面小铜镜，递给他。

镜子里是一个四十来岁的账房先生。脸比他圆一点，颧骨低一点，眼角多了两道纹，左边眉毛里藏着一颗小痣。皮肤是常年坐在柜台后头、不怎么见太阳的那种白。

Dispher 抬手摸了摸脸。摸上去是自己的脸，镜子里不是。

``鳞粉。''赵雪谣说，``摸得着的是你的，看得见的是我的。魂力在七十级以下的，看不破。七十级往上，离得近了，盯着你看一炷香，也许能看出点不对。''

``也许？''

``那就别让人盯着你看一炷香。''

轮到真霜长老的时候，赵雪谣绕着他走了一圈。

他拿手指头从真霜长老的鬓角上抹过去，抹过的地方，那点蓝就没了，只剩下白。

``头发上那点蓝，扎眼，得盖掉。白头发留着。老人家往西收药材，满头白发，正好。''他说，``腰得弯一点。''

真霜长老没动。

``一百三十六岁的人，腰挺成这样，''赵雪谣说，``是告诉人家你练过。''

真霜长老站了一会儿，慢慢把腰塌下去两寸。

赵雪谣又扇了一下翅膀。光落下去，老人脸上的棱角软了，眼睛浑了一点，嘴角往下耷拉，看着是一个常年咳嗽、脾气不太好的老东家。

最后是他自己。他抬手往脑后一捋，那一头银发从发根开始暗下去，一寸一寸，暗成了墨色。脸没怎么动，只是眉眼间那点秀气收了收，像一个读过几年书、被家里逼着出来学做买卖的少东家。

``老东家，少东家，一个账房。''他说，``往西收药材的。''

``收药材？''真霜长老说。

``韩家药材行在落枫倒了。''Dispher 说，他已经明白过来了，``西边那条商道上，收药材的买卖空出来了。空出来的地方，总有人想补。''

``咱们去补。''赵雪谣说。

他把那面小铜镜收回袖子里，朝真霜长老欠了欠身。

``祖父，路上委屈您了。''

真霜长老看了他一眼，没理他。

\begin{center}\rule{0.5\linewidth}{0.5pt}\end{center}

他们是第三天四更出的城，走的是东门。

一辆骡车，车上码着十几只药箱，最底下那只，垫着三层油布，装的是照幽。一人身上贴肉揣着一块片子。

骡车出了东门，往东走了二十里，拐上往北的岔路，又走了两天，才折向西。

西门那边，十几天前出去的那十几组人，都是从西门走的。西边那一路上的眼睛，盯的也是西门。

``打个措手不及，''Dispher 坐在车辕上说，``就得先让人家觉得，咱们压根没来。''

一路往西，走了二十三天。

赵雪谣每三天给真霜长老续一回腿。续完了，他就坐在车尾，打开那把折扇，扇面上什么也没画，他也不扇，就那么看着。

Dispher 每天夜里打开最底下那只药箱，看一回照幽。两百个点，有的挪了，有的停了，都还是白的。

过了第十九天，路两边的地开始发黄。树少了，沙多了，风里带着土腥味。

第二十三天傍晚，骡车翻过一道土梁，苍岭城就在底下。

城墙是黄土夯的，比落枫的矮，可厚。城门口站着四个兵，进城的人排成一溜，一个一个地问：打哪儿来，到哪儿去，叫什么，做什么买卖。问完了，旁边一个书吏拿笔记下来。

``收药材的。''Dispher 把路引递过去，``天斗城来的，姓张。''

兵丁看了看路引，又看了看车上的药箱，掀开一只闻了闻，挥手让他们过去了。

进了城，街是干净的，干净得出奇。铺子都开着，人也不少，可说话的声音很小。茶棚里坐满了人，一桌一桌地喝茶，没有人大声说话。棚子的柱子上钉着一块木牌，上头四个字：莫谈闲事。

骡车慢慢往城里走。走到城当中，路边是苍岭中殿。门口两根旗杆，旗子垂着，没有风。

中殿对面是一面长长的白墙。

墙上干干净净。一张纸都没有。

Dispher 让骡车放慢了一点。他看着那面墙，看了很久。

``落枫那面墙上，有一百三十七张。''他小声说。

车里，真霜长老咳嗽了一声。那声咳嗽是赵雪谣教他的，咳得很像。

``落枫好歹还有人敢贴。''老人说。

赵雪谣坐在车尾，手里那把折扇合着。他没往墙上看，看的是街角。街角蹲着一个卖烧饼的，烧饼摊上的炉子没生火，那人的眼睛一直跟着骡车走。

Dispher 低下头，在册子上记了一行字：

苍岭城。墙上无纸。

\begin{center}\rule{0.5\linewidth}{0.5pt}\end{center}

他们在城西一家客栈落了脚。

第二天一早，三个人分了头。

Dispher 去了城东的药市。

药市是一条长街，两边十一家药铺，门口摆着笸箩，甘草、黄芪、当归，一样一样晒着。他扮的是账房，进一家，问一家的价，问完了记在册子上，挑便宜的买一点。

价钱都公道。掌柜的也客气，有问必答，问到哪味药哪个山头出的、哪个月收成好，说得头头是道。

第三家，他买了二斤黄芪，付钱的时候顺口问了一句：

``掌柜的，咱们这市上抽税，是城主府的人来收，还是\ldots\ldots{}''

掌柜的数钱的手停了一下。

``官家的事，小的不懂。''

``那往年落枫韩家来收货，''Dispher 又问，``是不是秋天来？''

``不认得，不认得。''掌柜的把黄芪包好，往他怀里一塞，``客官慢走。''

Dispher 走到门口，隔壁铺子的伙计正在收笸箩。那伙计没抬头，声音压得很低：

``客官，买药就买药。''

那天他走完了十一家。问价，十一家都说得清清楚楚。问到税，十一家都说不懂。有两家，话没问完，掌柜的就去后头找东西，找了半天没出来。

\begin{center}\rule{0.5\linewidth}{0.5pt}\end{center}

真霜长老在城北那个茶棚里坐了一整天。

他要了一壶最便宜的粗茶，坐在最靠里的角落，咳一阵，喝一口，再咳一阵。没有人多看这个老头一眼。

茶棚里说的都是家常。麦子什么价，谁家闺女下个月出门子，东头那口井的水又咸了。说着说着，笑一阵，又接着说。

后半晌，邻桌一个年轻后生喝多了两碗，嗓门大起来：

``听说东边落枫，中殿的殿主叫人当街押上了台子\ldots\ldots{}''

他旁边坐着的中年人一脚踩在他的鞋面上。

后生的话断了。一桌人都不吭声了，有的低头喝茶，有的拿指头敲桌面。过了一会儿，那后生站起来，往桌上丢了几个铜板，走了。

茶棚又热闹起来，接着说麦子，说井水。

真霜长老端着茶碗，看见茶棚的掌柜往街角那边看了一眼。街角那个卖烧饼的，今天还蹲在那儿，炉子还是没生火。

\begin{center}\rule{0.5\linewidth}{0.5pt}\end{center}

赵雪谣在客栈门口摆了一张桌子。

桌上立一块木牌，牌上写着：张记药行，少东家坐堂，义诊三日，不收分文。

头一个时辰，没有一个人过来。走过去的人都往这边瞟一眼，又把眼睛挪开，脚底下没停。

后来来了一个拉车的，腰闪了，扶着墙一步一步挪过来，往凳子上一坐，龇着牙说不出话。赵雪谣让他趴在桌上，手在他腰眼上按了一会儿，指缝里漏出一点光，很淡，太阳底下谁也看不出来。

拉车的站起来，扭了扭腰，愣住了，又扭了扭。

``神了。''他说。

这一声喊出去，人就来了。发烧的孩子，咳嗽的老人，手上生疮的妇人，一个接一个，排到了街口。

赵雪谣一个一个看。看完了，开方子，抓药，嘴里跟人家闲聊，问的都是些不相干的事：吃得下饭吗，夜里睡得着吗，家里几口人。

``家里几口人？''

一个抱着孩子的妇人张嘴就答：``五口\ldots\ldots{}''

她停住了。

``四口。''她说，``四口。''

``原来是五口？''赵雪谣问，声音很温和。

``一直是四口。''妇人低下头，把孩子往怀里搂了搂，``一直是四口。''

赵雪谣没再问，把药包好，递给她。

那一天，他看了四十三个人。问家里几口人，有九个，先说了一个数，又改了口。九个都是往小里改。改口之前，九个都先往街角那边看了一眼。

天快黑的时候，最后来了一个老太太，领着一个七八岁的小孙子。孩子咳了半个月，咳得小脸发白。

赵雪谣给孩子看完，开了药，又多包了一包冰糖。

老太太没走。她攥着那包药站在桌子边上，嘴唇动了好几回，像是有话要说。

赵雪谣等着。

老太太往街角看了一眼。烧饼摊上的那个人，今天头一回站了起来，拍了拍身上的土。

老太太把嘴闭上了。

``少东家是好人。''她说，``好人，就早些走吧。''

她拉着孙子，走了。

\begin{center}\rule{0.5\linewidth}{0.5pt}\end{center}

夜里，三个人在客栈楼上那间屋子里碰了头。

Dispher 把册子摊开：``药市十一家，问到税，都说不懂。''

``茶棚里提了一句落枫。''真霜长老说，``一桌人不吭声了。''

赵雪谣坐在窗边，没点灯。

``我看了四十三个人。''他说，``问家里几口人，九个改了口，都往小里改。''

Dispher 拿笔的手停了一下。

他在册子上记下一行：九户，少了人。

``他们不肯说。''真霜长老说。

``嘴是闭上了。''赵雪谣看着窗外，``数闭不上。''

窗外，街角那个烧饼摊早收了。可对面屋檐底下的黑影里，有一点火星，一明一暗，有人在那儿抽烟。

\begin{center}\rule{0.5\linewidth}{0.5pt}\end{center}

那点火星亮到了三更。

三更一过，楼下的门被人拍响了。拍得很急，一下接一下，整座客栈都醒了。

``巡捕房办案！开门！''

掌柜的披着衣裳跑去开门。门一开，呼啦啦涌进来六个人，青衣皂靴，腰里挂着铁尺和锁链。领头的是个黑脸汉子，三十多岁，胸口挂一块铜牌，刻着一个''捕''字。他身后跟着一个中年男人，哭哭啼啼的，手里攥着一个药包。

``张记药行的人住哪间？''

掌柜的手往楼上一指，指完了，就缩到柜台后头去了。

楼梯咚咚地响。门没等敲，就被一脚踹开了。

屋里点着灯。三个人都没睡，也都没动。真霜长老歪在床头，裹着被子咳嗽。Dispher 坐在桌边，账本摊着，手里还捏着笔。赵雪谣坐在窗下，折扇合着，搁在膝盖上。

``就是他们！''中年男人从捕头身后挤出来，指着赵雪谣，``就是这个少东家！白天给我爹看病，开了这包药。我爹喝下去，天没黑就咽了气！''

他把药包往桌上一摔。纸包散开，滚出几截黑褐色的根子。

赵雪谣看了一眼那几截根子，又看了一眼那个人。

``断肠草。''他说，``张记的方子里，没有这一味。''

``你说没有就没有？''捕头冷笑一声，``人死了，药在这儿，苦主在这儿。''

``今天我看了四十三个人。''赵雪谣说，声音还是很温和，``四十三张脸，我都记得。没有你，也没有你爹。''

中年男人的哭声顿了一下，又接着哭，哭得更响了。

捕头没接这话。他一挥手：``搜！''

两个巡捕扑到墙角那摞药箱跟前，一只一只往下掀，药材撒了一地。掀到第三只，一个巡捕的手往袖子里一缩，再伸出来，从药材底下翻出一个油纸包，举得高高的。

``头儿！这儿还有！''

油纸包打开，又是一包断肠草，比桌上那包还多。

Dispher 看见了那只缩进袖子里的手。他没出声，低头在账本上记了一笔。

``人证物证都齐了。''捕头从腰里抖出锁链，哗啦一响，``贩卖毒药，谋财害命。三个都带走，押进大牢，明天过堂。''

他走到墙角，弯下腰，伸手去掀最底下那只药箱。

那只药箱垫着三层油布。

真霜长老不咳了。

他靠在床头，被子还裹在身上，只是把眼皮抬了起来，看了捕头一眼。

就一眼。

屋里的灯火没晃，窗户纸也没响。可捕头伸出去的那只手，停在了药箱上头半寸的地方，再也落不下去。

他想把手收回来，收不回来。想直起腰，直不起来。整间屋子的分量好像一下子全压到了他一个人背上，压得他骨头缝里咯吱咯吱地响。他身上那点魂力，二十几级的大魂师，攒了二十年的那点东西，在这一眼底下缩成一团，缩到最里头，一动都不敢动。

另外五个巡捕也一样。举油纸包的那个，胳膊还举在半空；抖锁链的那个，锁链还垂在手上，一节一节绷得笔直，一声不响。他们的眼珠子还能转，转来转去，转不出这间屋子。汗顺着鬓角往下淌，淌进了领子里。

那个中年男人张着嘴，哭声卡在喉咙里，像被人掐住了脖子。

这屋里三个人，随便拎出哪一个，都是封号斗罗的分量。真霜长老今天，只用了一眼。

赵雪谣站起来，走到窗边，把半开的窗户轻轻合上了。

对面屋檐底下那点火星，看不见了。

他回过身，看着屋里这一群一动不动的人，笑了笑。

``几位官爷，''他说，``站着累，可也坐不下。那就站着说吧。''

Dispher 放下笔，起身走到举油纸包的那个巡捕跟前，把油纸包从他手里拿下来，又把手伸进那人的袖子，摸出一张叠起来的草纸。草纸上沾着一点黑褐色的末子，跟包里的断肠草一个颜色。

他把草纸摊在桌上，跟那两包药并排放好，坐回去，翻开账本，蘸了蘸墨。

``一个一个来。''他对捕头说，``先说你。谁让你来的？''

捕头的嘴唇抖了抖。

``\ldots\ldots 总捕头。''声音是从牙缝里挤出来的，``总捕头吩咐的。说城里来了三个外乡人，打听不该打听的，按老规矩办。''

``老规矩是什么？''

捕头不说话了。

真霜长老的眼皮又往上抬了一点。

捕头的膝盖咯噔响了一声。``栽个罪名，关进大牢。''他喘了一口气，``关上几天，就\ldots\ldots 就没了。''

``没了是什么意思？''Dispher 问。

``每月初一，牢里来一辆车。车上盖着油布，夜里来，夜里走，往西走。''捕头说，``拉走的人，没有一个回来过。''

屋里静了一会儿。

Dispher 在账本上写了一行：每月初一，往西。

他转过头，看着那个''苦主''。

``你爹呢？''

中年男人的哭声早没了，嘴还张着。``在\ldots\ldots 在家。''他说，``活着。巡捕房给了二两银子，让我来哭一场。''

Dispher 从怀里摸出那块片子，往捕头腰上那块铜牌跟前凑了凑。片子中间那块晶石，边上泛起了一层灰，很淡，像有人对着镜子哈了一口气。

``他身上没有。''赵雪谣说，``是牌子上沾的。这块牌子，过过别人的手。''

``总捕头，城主府，中殿。''真霜长老一样一样地数，``一根线上的。今晚就去，把巡捕房拿了，天亮之前，总捕头跑不了。''

赵雪谣拿折扇在手心里敲了一下：``一晚上的事。''

Dispher 没有马上接话。他看着账本上那一行字，看了很久。

``仅仅我们去审判这些官员，是不合适的。''他说。

真霜长老看了他一眼。

``我在御史台二十年，审过的官不算少。''Dispher 说，``审一个官，要苦主，要证人，要卷宗，要一座公堂，要当着这一方的老百姓审。''

``苍岭的苦主，一个都不敢开口。九户人家，家里少了人，连少了都不敢认。我们三个今晚把巡捕房拿了，替他们把案子审完，审出来的东西是我们的，跟他们没有关系。''

``明天苍岭城里只会多一个说法：天斗来了三个更狠的。''

``那面墙上，还是一张纸都不会有。''

赵雪谣的扇子停在了手心里。

``还有。''Dispher 说，``巡捕房上头是城主府，城主府上头是谁，中殿跟哪一条线连着，每月初一那辆车往西走到哪儿，这些我们一样都还没摸着。今晚动了巡捕房，上头一夜之间就能把线剪干净。''

``落枫那一回，梅长老来得够快了，韩府的账还是烧了一半。''

真霜长老沉默了一会儿。

他低头看着地上那个捕头。

``总捕头，几级。''

``六\ldots\ldots 六十出头。''捕头的牙磕了一下，``六十三。''

``城主。''

``七十一。''

``中殿殿主。''

``七十四。''捕头咽了一口唾沫，``听说开春破了七十五。''

真霜长老点了点头，转向窗下。

``你的领域，能开多大？''

赵雪谣手里的折扇转了半圈。

``百丈之内，绝无问题。''

他拿扇子在膝盖上轻轻敲了一下。

``我有个好主意。''他说，``我直接开领域，定向保护百姓。我的魂力，对付这些七十多级的，还是很够的。''

地上那几个巡捕连喘气都放轻了。

``光底下，没作过恶的人，刀砍上去伤不着，魂技打上去也伤不着。''赵雪谣说，``作过恶的，我不护。''

真霜长老看着他，看了好一会儿。

``你是要让人们自己拿起武器，端了这些玩意。''

赵雪谣笑了笑，没说是，也没说不是。

Dispher 低头看着账本，手指在''九户，少了人''那一行上按了一下。

``好主意。''他说，``这样避免了误判。''

真霜长老转过来看他。

``我们三个在苍岭待了两天。''Dispher 说，``认得的脸，连地上这几位算上，不到十张。哪个巡捕手上沾过人命，哪个只是跟着跑腿，哪个狱卒半夜往牢里偷偷塞过一个馍，我们分不出来。''

``苍岭人分得出来。''

``谁家少了人，是谁敲的门，谁套的锁链，那一家子记了好几年了。夜里睡不着，一遍一遍地记。''

他把账本合上。

``他们认人，比我们准。''

赵雪谣点了点头。

``就算认错了，''他说，``光底下，打错了人，也打不疼。''

``还差一样。''他又说，``领域开着，没人第一个站出来，那也就是一片光。苍岭人低着头走了五十年，光底下，照样低着头走过去。''

Dispher 看着桌上那两包断肠草。

``中殿对面那面墙，''他说，``一张纸都没有。''

``明天早上，让它有纸。''

\begin{center}\rule{0.5\linewidth}{0.5pt}\end{center}

那天夜里，客栈楼上那间屋子的灯，一直亮到四更。

Dispher 把账本翻到新的一页，让捕头一句一句地说。谁吩咐的，哪天，原话是什么，车几时来，往哪边走。捕头说一句，他写一句，写完了念一遍，问一句''对不对''。

写到车，他停了笔。

``这个月的车，哪天来？''

``初一。''捕头说，``还有四天。''

Dispher 把''四天''两个字写了上去。

``苦主''那一张短。姓什么，住哪条巷子，哪天谁给了他二两银子，让他上哪儿哭。写完了，那人把二两银子摸出来，放在纸边上，又往里推了推。

屋里没有印泥。赵雪谣从药箱里翻出一块朱砂，搁在茶碗底，兑了两滴水，拿扇骨慢慢磨开。

``张记药行，''他说，``药材倒是真的。''

捕头把大拇指按进碗里，又按在纸上，看着那个红印子，看了很久。

六个巡捕，一人一个手印，按在捕头那张纸底下，一个都没少。

赵雪谣起身，把窗户推开一条缝。对面屋檐底下那点火星还在，一明一暗。他把两根手指伸出窗缝，轻轻一弹。

几点银灰的粉末飘过街面。

那点火星往下一沉，掉在地上，滚了两滚，灭了。黑影里有人顺着墙根滑坐下去，脑袋歪到了一边。

``睡到天亮。''赵雪谣把窗户合上，``醒了，只当自己昨晚多喝了两碗。''

捕头还跪在地上，看着纸上那一溜红手印。看着看着，他又开了口。

``还有一个。''

Dispher 抬起头。

``天亮以后，午时三刻，白墙底下要斩一个。''捕头说，``城主府的马夫，姓石。半个月前夜里在夜市上纵马，踩死了两个人，一个老的，一个小的。城主亲自审的，总捕头陪审，是铁案。斩字的告示，五更就贴出去。''

``城主府的马夫，''Dispher 说，``骑的是城主府的马。''

``是。''

``那天夜里，谁看见他骑了？''

捕头张了张嘴。

``城主府的人说的。''

Dispher 没往纸上写。他把笔搁在砚台边上，看着窗户纸。窗户纸外头，天还黑着。

``天亮以后，''他说，``那面墙上头一张纸，是他们贴的。''

\begin{center}\rule{0.5\linewidth}{0.5pt}\end{center}

\chapter{第一百零一章：午时三刻}\label{ux7b2cux4e00ux767eux96f6ux4e00ux7ae0ux5348ux65f6ux4e09ux523b}

四更天，客栈后院的井台上落了一层霜。

Dispher 坐在井台边上。赵雪谣站在他跟前，那只蝶只放出来半边翅膀，贴着地皮扇，灰色的光一层一层扑在 Dispher 脸上。

``脸要糙。''赵雪谣说，``牢里的人一年见不着几回太阳，可也不白，是那种捂出来的黄。''

光落了三回。Dispher 摸了摸脸，摸上去还是自己的。铜镜里是个二十出头的后生，塌鼻梁，左脸一溜麻子，嘴唇干得起了皮。

``牢里要是来了七十级往上的，''赵雪谣把镜子收回袖子，``别抬头。''

``一炷香。''Dispher 说，``我记着。''

真霜长老站在廊下，披着那件旧棉袍。

``御史台的大人，''他说，``去当狱卒。''

``狱卒手里有钥匙，有灯。''Dispher 说，``牢里什么样，只有他们看得见。''

捕头在院门口等着。领口里那点鳞粉还在，他一晚上没敢往别处多看一眼。

``刘二这两天害病，''他说，``我就说是总捕头叫补的人。老吕不会多问。''

``老吕是谁？''

``管夜班的。在牢里守了三十年。''捕头顿了一下，``腿瘸。''

两个人一前一后出了客栈，往城西走。

街上一个人都没有，铺子的门板全上着。走过城当中，中殿那两根旗杆在黑里竖着，像两根插在地上的筷子。中殿前头的门窗都是黑的，后院里有一扇窗户亮着，四更天了，还亮着。

捕头走过去的时候，脚底下快了几步，头一直没往那边偏。

\begin{center}\rule{0.5\linewidth}{0.5pt}\end{center}

苍岭大牢在城西，挨着城墙根。门口两盏灯笼，一盏灭了。

老吕坐在台阶顶上。五十多岁，裹一件油黑的羊皮袄，怀里揣着一只酒葫芦。他看了捕头一眼，又看了 Dispher 一眼。

``刘二的班？''

``刘二的班。''捕头说。

老吕嗯了一声，从腰里解下一串钥匙，站起来，拖着一条腿往下走。

台阶是石头的，往下一共三十七级。越走越潮，墙上渗着水，霉味里混着屎尿味。Dispher 数着台阶，数到二十几级，胸口那块片子凉了一下。

捕头送到台阶底下就走了。老吕把钥匙串挂到 Dispher 腰上。

``东头七间，西头五间。一个时辰走一趟，看看有没有死的。''他说，``最里头那间是死牢，午时三刻的。别跟他说话。''

``他犯的什么事？''

老吕没答，回到台阶底下那张小桌边上坐下，拔开酒葫芦喝了一口。

Dispher 提着灯往东走。

栅栏后头是黑的。灯照过去，黑里一张一张的脸往后缩。有男人，有女人，有个十二三岁的孩子，抱着膝盖坐在稻草上，眼珠子跟着灯转。东头七间，二十六个。西头五间，十一个。

走到东头第三间，有人开了口。

``天亮了没有？''

是个中年汉子，胡子乱成一团，两只手攥着栅栏。

``还没。''

``今儿初几？''

``二十七。''

汉子松了手，坐回稻草上，不说话了。

Dispher 走回台阶底下。老吕从羊皮袄里摸出一个油纸包，里头是四个杂面馍，捂得还有点温乎。

``东头第三间，第五间，西头第二间，一间一个。''他说着，把剩下那个掰开，``这半个给那孩子。''

``总捕头要是问起来\ldots\ldots{}''

``他不会问。''老吕又喝了一口，``问了，你就说是我。''

Dispher 把馍送进去。那半个馍递到孩子手里，孩子两口就吞了，把指头一根一根舔干净，舔完了，仰着脸问：

``叔，外头下霜了没有？''

``下了。''

孩子点点头，缩回稻草里去了。

东头第三间那个汉子接过馍，没吃，抬头看着他。

``新来的？''

``嗯。''

``夜市那天，我闺女在我跟前买糖人。''汉子的声音压得很低，``马是从北头冲下来的。枣红马，四个蹄子钉着银掌。马上那个人穿一身白，腰里挂着玉。''

他把馍攥在手心里。

``我在城主府门口跪了一天，说的就是这几句。第二天，我就进来了。''

Dispher 站在栅栏外头，没出声。

``马夫我认得。''汉子说，``石柱子。他骑不起那样的马，也穿不起那样的衣裳。''

最里头那间是石头砌的，门比别处矮，栅栏是铁的。

里头只有一个人，靠墙坐着，脚上一副镣，手上一副铐。二十七八岁，肩膀很宽，两只手摊在膝盖上，掌心横着几道老茧，是常年攥缰绳勒出来的。

灯照过去，他没躲。

``天亮了没有？''

又是这一句。

``还没。''

那人点了点头，把脸转过去，对着墙。

台阶上头响起了脚步声。

老吕手忙脚乱地把酒葫芦往羊皮袄里塞，站起来，那条瘸腿一软，扶住了桌子。

下来三个人。前头一个打灯笼，管家模样。中间那个五十来岁，白面短须，穿一件寻常的灰绸袍子，大拇指上套着玉扳指。后头那个又高又壮，黑脸，腰里挂刀，胸口一块银牌，刻着''总捕''两个字。

Dispher 低下了头。

片子贴在他胸口，凉得像一块冰。凉是从那个黑脸汉子身上过来的，他往前走一步，片子就凉一分。

``城主。''老吕弯着腰，``总捕头。''

城主走过 Dispher 身边，脚步停了停。

``新来的？''

``刘二害病，''老吕抢着说，``补的。''

城主的眼睛在 Dispher 脸上停了一下。

Dispher 盯着自己的鞋尖。鞋尖上沾着一块泥，是台阶上蹭的。他听见自己的心跳，一下，两下，三下。

城主的眼睛挪开了。

``灯。''

Dispher 提着灯，跟在后头，走到最里头。

城主隔着铁栅栏站着。

``石柱。''

墙边那个人没回头。

``你娘，你儿子，城主府养着。''城主说。声音不高，像家里长辈交代事情。``你娘那条腿，开春我请大夫给她看。你儿子，过两年送进学堂。''

``午时上去，一个字别说。''

``说了，你娘活不过初一。''

牢里很静。东头有人翻了个身，稻草响了一声。

过了很久，墙边那人的肩膀动了一下。他点了点头。

城主转身往外走。总捕头跟在后头，走到台阶底下，压着嗓子说：

``三公子那匹马，蹄铁换过了。夜市那两家，老的那家给了十两，没人闹。小的那家\ldots\ldots{}''

``东头第三间。''城主说。

``初一的车，一块儿拉走。''总捕头说，``客栈那三个外乡人，天亮前也押进来。''

``琛儿这两天在府里抄经。''城主踩上头一级台阶，``午时过了，再让他出门。''

灯笼的光一级一级往上走，走没了。

老吕还弯着腰站在桌子边上。站了一会儿，他慢慢坐下去，把酒葫芦掏出来，拔开塞子，又塞上了。

``听见了？''他没看 Dispher。

``听见了。''

``听见了，就当没听见。''老吕说，``我在这底下三十年，三十年就是这么过来的。''

Dispher 提着灯，又往里走了一趟。

东头第五间，一个女人扒着栅栏，等灯过来。

``小兄弟，''她说，``东街那家豆腐坊，还开着吗？''

Dispher 想了想。药市往北拐过去，是有一家豆腐坊，门口支着一口大锅，早上冒热气。

``开着。''

女人点点头，退回黑里，把脸埋进膝盖。过了一会儿，黑里传出很轻的抽气声，一下，一下。

墙上刻满了道道。四道一竖，一道一横，一捆一捆的，从地面刻到人伸手够得着的地方。有的刻得深，有的只浅浅一划，刻到一半就停了。

最里头那间，石柱还对着墙坐着。

灯过来的时候，他开了口。

``斩人那天，每家要出一个人去看。''

Dispher 没接话。

``我家就两口人。''石柱说，``我娘，我儿子。''

他没再往下说。

Dispher 回到台阶底下。老吕把酒葫芦抱在怀里，眼睛闭着，没睡。

``初一那车，''Dispher 问，``走了多少年了？''

老吕没睁眼。

``我来的那年就有。''他说，``头几年，一个月拉一两个。这几年，一回拉七八个。拉到哪儿，我不知道。赶车的脸蒙着，不说话，手上戴着手套，大夏天也戴。''

``三十年，没见回来过一个。''

天快亮的时候，换班的人下来了。Dispher 把钥匙串解下来，搁在桌上，搁在酒葫芦旁边。

``老吕。''他说，``那四个馍，我记着。''

老吕抬头看他，没明白。

\begin{center}\rule{0.5\linewidth}{0.5pt}\end{center}

客栈楼上，真霜长老和赵雪谣都没睡。

Dispher 进门，在桌边坐下，脸还是那张麻脸。他把夜里的事说了一遍，从东头第三间说到城主那三句话，一个字没改。

说完了，屋里没人出声。

``城主看了你多久？''真霜长老问。

``三下心跳。''

``他没认出来。''

``他没往下看。''Dispher 说，``一个城主，看一个狱卒，看三下心跳就够了。''

``石柱子的娘，''赵雪谣开了口，``是不是领着个孙子，孩子咳了半个月。''

Dispher 看着他。

``我给孩子看病那天，问过一句孩子爹在哪儿。''赵雪谣说，``她说，在城主府喂马。''

``她还说，好人，就早些走吧。''

扇子在他手里合着，一动不动。

``午时三刻斩人。''Dispher 说，``苍岭的规矩，斩人那天，每家出一个人去看，不去的记名字。捕头夜里说的。''

``全城的人都在白墙底下。''真霜长老说。

``墙上贴什么，他们看了也不敢认。''Dispher 说，``石柱子一个字不说，他们更不敢认。''

``得让他们看见点别的。''

他伸出一根指头，在桌面上画了一条线，从城西画到城当中。

``东头七间，西头五间，三十七个人。九户少了的人，有几个就在里头。卖糖人的也在里头，他看见过那匹马。''

真霜长老看着桌面上那条看不见的线。

``劫狱。''他说。

``劫狱。''

老人低头看了看自己那双手。手背上的皮松了，青筋一根一根鼓着。

``一百三十六年，''他说，``头一回。''

``囚车午时初出牢门。''Dispher 说，``囚车一走，牢里只剩老吕他们几个。您进去。''

``我在白墙底下。''赵雪谣说，``刀砍下去，伤不着他。''

``我上台子。''Dispher 说。

``中殿殿主呢？''赵雪谣问，``七十五。''

``总捕头身上沾着的东西，八成是从中殿那边过来的。''Dispher 说。

``他要是出来，''真霜长老说，``我在。''

赵雪谣没再问。他在真霜长老跟前蹲下来，一只手按在老人膝盖上。离上回续腿，正好三天。光落进去，老人站起来，在屋里走了两步，没扶墙。

捕头从牢里回来，就一直守在门外。Dispher 出去，跟他交代了一句：

``去跟总捕头报，客栈那三个，押进去了，关在西头。''

捕头的脸一下子白了：``总捕头要是下去看\ldots\ldots{}''

``午时以前，''Dispher 说，``他没工夫下牢里数人头。''

\begin{center}\rule{0.5\linewidth}{0.5pt}\end{center}

天亮以后，白墙上有了第一张纸。

是城主府五更贴的。一个斗大的''斩''字，底下几行小字：斩犯石柱一名，纵马踏死二命，午时三刻，就地正法。

纸贴得很高，伸手够不着。

苍岭的刑场就是白墙底下这块空地。巳时刚过，人就来了。一家一个，从四面八方的巷子里出来，往这儿聚。巡捕拿着册子站在路口，来一个，勾一个。

人越聚越多，挤满了半条街，没有一个人说话。

茶棚的掌柜来了。药市那十一家，来了十一个伙计。街角那个卖烧饼的也来了，揉着脖子，脸上还带着在墙根底下睡了一夜的印子，站在巡捕旁边，眼睛在人堆里转来转去。

那个老太太天不亮就来了，坐在白墙底下，一直坐到巡捕拿铁尺把她赶到人堆最前头。她手里牵着孙子。

孩子上过半年私塾。他仰着脸，看着墙上那张纸，一个字一个字地念：

``斩\ldots\ldots 犯\ldots\ldots 石\ldots\ldots 柱\ldots\ldots{}''

老太太一把捂住了他的嘴。

孩子在她手底下眨着眼睛，不念了。

赵雪谣站在人堆里，离她不远。今天没摆义诊的桌子，折扇合着，垂在手边。

中殿门口两根旗杆底下搭了一座台子。台上一张长桌，桌上一个签筒。城主坐在桌后，换了一身官袍。总捕头站在他身边。

中殿的大门关着。

街角，一匹枣红马慢慢踱过来，停在人堆外头。马上的人二十出头，一身白，腰里挂着玉。他嘴里嚼着什么，看着白墙，看得很有兴致。

城主看见了他，眉头动了一下，没说话。

午时，囚车到了。

囚车从城西过来，走得很慢。走过人堆跟前的时候，石柱抬了一下头。他看见了最前头那个老太太，看见了她手里牵着的孩子。

他把头低了下去，再没抬起来。

两个巡捕把他从车里拖下来，按跪在白墙底下。刽子手提着一把鬼头刀，站到他身后。

总捕头展开一张纸，念罪状：本月十二夜，城东夜市，城主府马夫石柱私骑府中马匹，纵马伤人，踏死二命，人证物证俱在，本人画押认罪。

他念得很快。念完了，把纸一卷。

老太太往前挤了一步，铁尺横过来，把她挡了回去。

石柱没抬头。

日头一点一点往正中挪。

城主看了看天，从签筒里抽出一支令签，往地上一扔。

``斩。''

刽子手含了一口酒，喷在刀上，双手把刀抡圆了，重重砍了下去。

刀砍在石柱后颈上。

当的一声，像砍在一块青石上。

刽子手两条胳膊让刀震得往上一弹，鬼头刀差点脱了手。他低头去看。石柱的后颈上，连一道白印子都没有。

石柱自己也不知道出了什么事。他跪在那儿，后脖颈上只觉得凉了一下，像有人往上头吹了口气。

刽子手的脸白了。他咬着牙，又抡了一刀。

当。

还是那一声。刀口卷了一块，石柱的头发一根都没断。

人堆里有人抬起了头。

赵雪谣的头发，从发根开始一寸一寸褪下去，墨色褪回了银色。他身后，那只蝶展开了翅膀。

这一回翅膀没收着。银色的翅脉一根一根亮起来，蓝的转成紫的，紫的转成白的。光从翅膀上落下去，落在白墙上，落在囚车上，落在每一个站着的人肩膀上，往外铺，铺过了半条街，铺到街口那口井，才停住。

百丈。

城主站了起来。

他身上腾起一团魂力云，云里压着七层颜色，两层黄，两层紫，三层黑，一层压着一层，绕着他翻滚。台子底下的人往后退，退成一片。

``妖人。''城主说，``拿下。''

一道青色的风从他掌心推出去，贴着地面卷进人堆，卷向赵雪谣站的地方。

风是实打实撞上去的。前排七八个人让它撞了个满怀，衣角呼啦啦往后扯，一条头巾吹上了天。风过去了，那七八个人还站在原地，低头看自己的手，看自己的胸口。身上一道口子都没有。

城主那团魂力云里，第四层亮了。一根风柱从台子上拔起来，砸向人堆正当中。风柱砸下去，街面上的土全掀了起来，两边铺子的门板哐哐直响。土落下来，底下的人一个一个还站着，有人咳嗽，有人揉眼睛，没有一个倒下的。

有个孩子从他娘身后探出手，往那阵风里抓了一把，抓了个空。

城主额头上见了汗。

他把魂力云里七层一齐催亮了。第七层黑得发沉，他身后升起一头青色大雕的影子，翅膀一张，遮住了半座台子。雕影俯冲下来，一头扎进人堆。整条街刮起一阵狂风，屋檐上的瓦片哗啦啦往下掉，掉在人头上，掉在人肩膀上。

风停了。

瓦片碎了一地。满街的人，一个没少，一个没伤。

雕影散了。城主站在台子上，喘着粗气，脸色发青。七十一级的魂力，一口气倒了个干净，底下连一个人的皮都没碰破。

Dispher 从人堆里走出来。那张麻脸已经没了，又是那个账房先生。他走到白墙底下，从怀里掏出两张纸，一张一张贴在那个''斩''字旁边。

捕头那一张，按着七个红手印。

他转过身，对着满街的人开了口。声音不大，可光底下很静，每一个字都送得很远。

``今天五更以前，我在苍岭大牢里当了一夜狱卒。''

``城主到死牢门口，对石柱说了三句话。''

``你娘，你儿子，城主府养着。''

``午时上去，一个字别说。''

``说了，你娘活不过初一。''

老太太手里的孩子被攥疼了，叫了一声。

``总捕头跟城主说，三公子那匹马，蹄铁换过了。''

满街的人，一个一个，慢慢把头转了过去，转向街角。

枣红马上那个年轻人不嚼了。

``胡说八道！''城主一掌拍在桌上，``拿下！统统拿下！''

巡捕往前冲。冲在头一个的，抡起铁尺就往 Dispher 头上砸。铁尺砸实了，Dispher 连眼睛都没眨一下。后头几个巡捕的铁尺砸在前排的人身上，砸在肩膀上，砸在背上，砸得结结实实，挨打的人晃都没晃。

巡捕们看着自己手里的铁尺，往后退了一步。

人堆里还是没人动。一个一个看着墙，看着台子，看着街角那匹马，脚底下像生了根。

Dispher 站在墙底下，没再说话。

\begin{center}\rule{0.5\linewidth}{0.5pt}\end{center}

午时初，囚车出牢门的时候，老吕还没走。

他换了班，没回家，一直坐在台阶顶上。囚车从他跟前过去，他扶着墙站起来，看着那辆车拐过街角，又坐了回去。

过了一会儿，街那头走过来一个老头。

满头白发，腰弯着，拄一根柴棍，咳两声，走几步。走到台阶跟前，停住了。

白天当班的两个狱卒从门洞里探出头来：``干什么的？''

老头抬起眼皮，看了他们一眼。

两个狱卒就那么探着头，不动了。

老头没再理他们。他看着老吕，什么也没说，冲他伸出一只手。手背上的皮松了，青筋一根一根鼓着，掌心是空的。

老吕看着那只手，看了很久。

他撑着墙站起来，拖着那条瘸腿走到门洞跟前，从一个狱卒腰上把钥匙串摘下来，走回来，搁进那只手里。

``东头七间，西头五间。''他说，``底下黑，我给您照亮。''

\begin{center}\rule{0.5\linewidth}{0.5pt}\end{center}

城西那边传来了脚步声。

很杂，很慢，一大片，从街那头过来。

走在最前头的是老吕，拖着那条瘸腿，腰里挂着那串钥匙。他身后跟着一溜人，衣裳破烂，脸色发灰，眯着眼睛，一个搀着一个，在大太阳底下走得磕磕绊绊。最后头是真霜长老，拄着那根柴棍，走得也不快。

三十七个人。

人堆的西边先乱了。

一个女人尖叫了一声，喊的是一个名字，喊完就往那边扑，被人挡住，又扑。一个汉子从人堆里往外挤，挤掉了一只鞋。那个十二三岁的孩子从队伍里跑出来，一头撞进一个妇人怀里，妇人跪在了地上。

那天抱着孩子来看病的妇人也在人堆里。家里原来五口，后来是四口。她看着队伍里一个男人，怀里的孩子往下滑，她都没去扶。

卖糖人的汉子穿过人堆，一直走到白墙底下。他站在石柱旁边，回过身，抬起手，指着街角。

``那天夜里，''他说，``骑在马上的，是他。''

枣红马上的年轻人笑了一声，笑得很短。

``一个卖糖人的。''他说。

没有人接他的话。

街上没有一点声音。

老太太弯下腰，从墙根底下捡起了半块砖头。

她把砖头扔了出去。

砖头飞过人头顶，砸在城主的额角上。

城主捂着头往后一仰，撞翻了签筒，令签撒了一地。他的指缝里淌下血来。

满街的人都看见了。鬼头刀砍在石柱脖子上，连个白印子都没留下；七十多级的魂力砸进人堆，没碰破一个人的皮。半块砖头砸在城主头上，血下来了。

挑水的抡起了扁担。

人从四面八方往台子上涌。

刽子手扔了鬼头刀，往后退，退到墙根，蹲了下去。

药市第三家的掌柜挤在最前头，扯着嗓子喊：``税是总捕头的人收的！一年收三回！''喊完了，他自己愣了一下，像是不认得自己的声音。

街角那个卖烧饼的转身往巷子里钻。没钻进去，叫茶棚的掌柜一把揪住了后领子。

``三年了，''茶棚掌柜说，``你在我棚子对面蹲了三年。''

巡捕的刀拔出来了。一刀砍在一个卖菜汉子的肩膀上，砍透了棉袄，棉絮翻出来，底下的皮肉一点事没有。总捕头身上冒起一团魂力云，六层颜色压在一块儿，一刀劈向卖糖人的汉子。刀落在汉子脑门上，像落在一张纸上，轻飘飘地滑开了。

扁担从旁边抡过来，砸在总捕头手腕上，咔的一声，刀掉了。他抱着手腕跪下去，疼得满头是汗。三四个人扑上去，把他按倒在地上。

城主往后退，退到中殿门口，抡起拳头砸门。

门没开。

人涌上台子，把他从门口拖了下来。他的官帽滚到了台阶底下，额角上的血糊了半张脸。

街角，枣红马人立起来。三公子勒着缰绳要往巷子里冲，几只手同时抓住了马笼头。马蹄底下是一副新钉的铁掌，钉帽还亮着。三公子从马背上滚下来，摔在青石板上，半边脸蹭掉了一层皮，那身白衣裳踩满了泥。

巡捕们挨了一顿乱打。有的抱着头蜷在地上，鼻子嘴里都是血。有一个年轻的，让人围着打了十几下，拳头扁担落在他身上，跟落在棉花上一样。他自己也愣着，摸摸脸，摸摸胳膊，一点不疼。

打他的人停了手，你看我，我看你。

``别打了。''一个卖针线的婆子挤过来，``他没干过。他去年开春才进的巡捕房，他娘就住我们那条巷子。他夜里来敲过门，敲完了，又偷偷把门给我们带上了。''

那几个人把手里的扁担放下了。

老吕被一群年轻后生围住了。他身上那件狱卒的衣裳，在人堆里扎眼。有个后生揪住他的领子，一拳打在他脸上。老吕的头偏了一下，又转回来，脸上什么也没有。那后生看看自己的拳头，又举了起来。

``别动他！''

队伍里好几个人一齐喊。

``他给我们塞过馍！''

那个十二三岁的孩子从他娘怀里挣出来，跑过去，挡在老吕跟前，张开两只胳膊。

举起来的拳头，慢慢放下了。

Dispher 在册子上记了一行：狱卒老吕，四个馍。

\begin{center}\rule{0.5\linewidth}{0.5pt}\end{center}

石柱身上的镣铐，是铁匠铺的人拿锤子砸开的。

他站起来，脚底下发软，一步一步挪到他娘跟前。

``娘。''他说，``我一个字都没说。''

老太太抬手打了他一下，打在肩膀上，又打了一下。第三下没打下去。

孙子抱住了他的腿。

老太太从人缝里看见了赵雪谣。头发成了银色的少东家，站在光底下，冲她点了一下头。

``少东家。''她说，``你没走。''

``没走。''赵雪谣说。

Dispher 走上那座台子。城主坐过的那张长桌翻倒了，他把桌子扶起来，摆正，把签筒捡起来，搁在桌角，把册子摊开，笔蘸了墨。

他没坐那把椅子，站在桌边。

``谁家少了人，''他说，``上来说。''

台子底下的人你看我，我看你。

头一个上来的是那个抱孩子的妇人。她男人跟在她后头，脸上的灰还没洗。她说了一个日子，说了一个巡捕的名字，说那天夜里是谁敲的门，谁套的锁链。

墙根底下捆着的那一串里，有一个巡捕把头埋了下去。

那个挨了十几下没伤着的年轻巡捕也上了台子。他把巡捕房这两年夜里出去拿人的日子，一个一个报出来，哪天，哪条巷子，谁带的队。报到后来，他不看台子底下了，只看着 Dispher 手里的笔。

第二个上来的是卖糖人的汉子。第三个是药市第三家的掌柜。

从未时说到申时，台子底下排起了队。Dispher 写满了半本册子。每写完一个，他念一遍，问一句''对不对''，底下的人点了头，他再写下一个。

日头偏西的时候，街上静了下来。

城主、三公子、总捕头，还有十几个巡捕，捆成一串，坐在白墙底下，坐在那个''斩''字下头。没有人再打他们。看着他们的是一群拿扁担、拿锄头的苍岭人，换着班，谁也不走。

赵雪谣沿着墙根，从这头走到那头，一个一个看过去。

城主歪在墙根底下，坐不起来。额角那道口子还在渗血，左边肋条塌下去一块，喘一口气，哼一声。三公子抱着一条腿，那条腿拧着，脚尖朝里。总捕头的手腕肿得像个馒头，嘴里少了两颗牙。

再往后是巡捕。有四个躺着，血糊了满脸，站不起来。有七个坐着，脸肿了，胳膊上几道红印子，自己还能揉。客栈里那个捕头，一只眼青了，一条胳膊脱了臼，耷拉在身子边上。

卖烧饼的那个缩在最边上，脸上五个指印，还红着。

赵雪谣看完了，回到台子底下。

``光底下挨的打，一下是一下。''他说，``落到谁身上有多重，是照着他手上沾的东西来的。沾得多，就重。沾得少，就轻。''

Dispher 把册子翻开，一页一页地对。

一下午记下来的，是谁家哪天少了人，谁敲的门，谁套的锁链。躺着的那四个巡捕，名字在册子上出了二十几回。坐着的那七个，三回五回。城主、三公子、总捕头，差不多每一页都有。

伤得重的，册子上名字多。伤得轻的，名字少。一个都没对错。

真霜长老站在他边上，看着那一页一页的字，点了点头。

Dispher 合上册子，走上台子。

``城主，三公子，总捕头，还有躺着的这四个。''他说，``斩。''

``坐着的七个，革去差事，各杖四十，发去界河修堤，三年。''

``捕头，革去差事，杖六十，修堤五年。卖烧饼的，枷号十日。''

墙根底下，城主撑着地，把头抬了起来。嗓子是哑的。

``本官\ldots\ldots 是朝廷命官\ldots\ldots 苍岭城主\ldots\ldots{}''他喘了一口气，``你一个收药材的账房\ldots\ldots 凭什么\ldots\ldots{}''

Dispher 没答他。他把手伸进账房衣裳的里襟，摸出一块令牌，举了起来。

铜底，巴掌大，正面一个''令''字，背面盘着一条龙。

``天斗御史台，Dispher。''他说，``各州城主府，见此令，如见圣上旨意。''

城主张着嘴，没再出声。

真霜长老把柴棍靠在囚车上。他的腰原本塌着两寸，这时候一点一点直了起来，直到底。赵雪谣抬手一拂，老人脸上那层浑浊散了，棱角回来了，眼睛也清了。

他从怀里摸出一块乌木令牌，正面一个篆体的''长''字，四周一圈金纹。

``扽魂殿，真霜。''他说。

赵雪谣什么也没掏。他身后那只蝶还张着翅膀，银发垂过了腰。

``供奉殿，赵雪谣。''

街上静了好一会儿。

人堆里有个老人先跪了下去，后头跟着跪了一片。

``不用跪。''Dispher 说，``苦主有了，证人有了，卷宗有了，公堂也有了。这一堂是你们自己审的。我们三个，只是把字写上。''

\begin{center}\rule{0.5\linewidth}{0.5pt}\end{center}

申时末，刑场还是那块刑场。

鬼头刀还是那把鬼头刀，刀口卷着一块。刽子手双手捧着刀，看了看 Dispher。

Dispher 弯下腰，从白墙根底下捡起一支令签。签子上沾着土，是午时城主扔下来的那一支。

他把令签扔了出去。

刀落下去。

这一回，没有当的一声。

一个，两个，三个。七个。

白墙底下先是静的。

后来不知道是谁，拍了一下手。一下，又一下。拍手的声音从墙根底下传出去，传过半条街，传到街口那口井。有人一边拍一边喊好，喊得嗓子都劈了。有人拍着拍着蹲了下去，两只手捂着脸，手掌拍得通红。

老太太没拍手。她攥着孙子的手站在最前头，一直看到最后一刀落下去。石柱站在她身后。

赵雪谣把领域收了。光一寸一寸缩回他背后，缩进那对翅膀里。头发还是银的，没再染回去。

中殿的大门从早到晚，一直关着。

Dispher 走到门前，从怀里摸出那块片子。片子中间的晶石，灰了一大片，灰得发暗。

他把片子收回去，什么也没说。

真霜长老走过来，站在他旁边，也看着那扇门。

``七十五。''老人说，``关了一整天门。外头这么大动静，一声都没漏进去。''

``今天夜里，''Dispher 说，``他要么烧账，要么跑。''

``殿里的账，''真霜长老说，``明天我来算。''

赵雪谣收着扇子走过来。

``后门，''他说，``我早上来的时候，落了点东西。''

天擦黑的时候，中殿后院那扇窗户又亮了。

白墙上的纸，也是这时候多起来的。

起先是一张。一个教私塾的先生蹲在墙根，替一个不识字的妇人写她男人的名字，哪年哪月，初一，往西。写完了，妇人拿浆糊贴了上去。

后来是两张，五张，二十张。排着队找先生写字的人，从墙根排到了街口。

卖糖人的汉子也排在队里。轮到他，他说了一个名字，三个字，又说了一个岁数：六岁。先生写完了，他自己把纸贴上去，贴得很低，贴在一个六岁的孩子伸手够得着的地方。

Dispher 站在人后头看着。

他把册子翻开，在''苍岭城。墙上无纸。``那一行底下，又写了一行：

苍岭城。墙上有纸。

\begin{center}\rule{0.5\linewidth}{0.5pt}\end{center}

\chapter{第一百零二章：银发女子}\label{ux7b2cux4e00ux767eux96f6ux4e8cux7ae0ux94f6ux53d1ux5973ux5b50}

蔡须坤站在沙子上，没动。

血里那股热静下去了，底下那一团也静下去了。十六年，身子里头一直有东西在顶，在拽，劲一下子没了，她晃了一下，伸手扶住了船帮。

蔡兴坤脚一沾沙子，膝盖就软了。小舞在后头一把捞住她，两个人一块儿坐倒在沙子上。诺当的十根指头攥着桨杆攥了一整天，松不开，是唐三一根一根替他掰开的。掌心里那截袖子跟肉粘在一块儿，黑红黑红的。

奥斯卡把剩下的香肠一人塞了一根。那根拿油纸包了三层的，他没动，又揣了回去。

滩上有一道湿印子，比舢板停的地方高出十几丈。

``潮在涨。''霍雨浩说，``船漂走了，雾里那条道就白找了。''

六个人把肩膀抵在船帮上。沙子软，一脚陷到脚脖子，往前顶一步，往后溜半步。舢板泡了一整天水，沉得像块石头。

``一，二。''唐三喊。

一尺，又一尺。潮水在后头漫上来，漫过了他们刚才坐的地方。推过那道湿印子，六个人一松劲，全躺倒在沙子上。

满滩的鸡一直看着。没有一只过来，也没有一只走开。

\begin{center}\rule{0.5\linewidth}{0.5pt}\end{center}

从滩上到牌坊，是一溜石阶。阶面长着青苔，每一级都比大陆上的高出一截。

唐三背着蔡兴坤。蔡须坤不让人扶，自己扶着石头，迈一级，喘一口气。

走到半道，石阶塌了三四丈，只剩一面长着青苔的陡坡。小舞把鞋一脱，腰上缠了缆绳往上爬，爬到一半脚底下的石头掉了，整个人吊在石缝上。底下霍雨浩的脸一下子白了。她腰一拧，贴回坡上，接着爬，爬到顶上把绳子拴在石桩上。

剩下的人攥着绳子一个一个上去。诺当的手攥不住，是上头五个人拿绳套拴着他的腰，一尺一尺拽上来的。

再往上，台阶边上摊着一张补了一半的渔网，梭子还别在网眼里。台阶当中搁着一筐鸡蛋。坡上的烟囱都在冒烟，可一路上一个人都没看见。

卢本韦回头看了一眼。雾墙还在海上立着，一丝缝都没开。

``曾祖写的是，等它来接。''她说。

``那咱们这算什么？''奥斯卡问。

``不请自来。''

唐三背着蔡兴坤，迈上最后一级台阶，抬起头，停住了。

牌坊底下站着人。

\begin{center}\rule{0.5\linewidth}{0.5pt}\end{center}

五个人，看着都跟他们差不多大。谁也没听见他们是什么时候来的。

一个黑瘦的少年，赤着脚，手里拎着一柄鱼叉。一个短头发的姑娘，腰里别一把柴刀。一个高个子少年背着手站在后头，身后还躲着一个圆脸的小个子。

当中那一个，是个姑娘。

头发是白的，白里透着一层银，没有束，一直披过了腰。海风从坡下往上吹，那几个少年的衣角直扑腾，她的头发一缕都没乱。

眼睛是紫色的。颜色很深，可清得能看见底，像山涧里的一汪水，水底铺着一层紫晶。

她的脸生得极精致。眉毛细长，弯弯地扫进鬓角，螓首蛾眉，说的大概就是这个样子。皮肤白得发凉，像一块冰里冻着的玉。鼻梁，嘴唇，下颌，一样一样看过去，挑不出一处不妥帖的地方。

她个子很高。唐三他们几个男孩子比她高一点；女孩子里，只有卢本韦跟她差不多，小舞、蔡须坤、蔡兴坤都比她矮。肩背挺直，脖颈修长，腰身和两条腿都长，上下比例匀称得像拿尺子量过。看年纪跟他们差不多，身量却已经是个成年女子了。

一身素白的粗布衣裳，领口交叠着。领口往上一点，锁骨上头，脖子底下，有一枚浅淡的椭圆形胎记，不凑近了看不出来。

她两只手垂在身侧，什么也没做。可她往那儿一站，坡上的风好像就小了。

她的眼睛从他们脸上一个一个看过去，看一个，停一口气。脸上什么都没有。看到诺当的时候，诺当的脚往后挪了半步，他自己都没觉出来。

然后她开了口。声音不高，清清冷冷的。

``你们是谁，为什么到达这个岛屿？''

黑瘦少年手里的鱼叉，叉尖往上抬了一寸。

蔡须坤从唐三身后走了出来，站到最前头。

``天斗城来的，八个人。''她说，``来找最菜萌新大供奉。''

高个子少年嗤了一声。短头发的姑娘把手搭在了柴刀上。

那姑娘没答话。她看着蔡须坤，看了很久，脸上还是什么都没有。

``回去。''她说。

就两个字。

黑瘦少年往前跨了一步，鱼叉一横，堵住了往上去的路。后头那个高个子少年手里多了一根杖，杖身细长，杖头上一点一点亮着，像攥了一把星星。诺当往前一站，挡在蔡须坤身前半个身子，两只缠着布的手张开了。唐三背着蔡兴坤，没往后退。

两拨人隔着三级台阶，谁都没再动。

那姑娘的眼睛在诺当张开的两只手上停了一下。

她身后亮了。

光是银的，从她脚底下升起来，一直升到牌坊顶上。光里浮出一条龙，通体银色，鳞片一片压着一片，亮得晃眼。龙身绕着牌坊盘了半圈，龙头垂下来，悬在她头顶上，一双眼睛也是紫的。

坡上的鸡全趴下了。海风停了一停。

霍雨浩的嘴唇动了动，声音很轻：

``银龙王。''

这三个字他在书里见过，在最老的那几册里，后头跟着四个字：上古异兽。

诺当张着的两只手僵在了半空。唐三背上的蔡兴坤把脸埋进了他的肩膀。这一回压过来的分量，压的是骨头，是骨头里的血。

蔡须坤把蔡兴坤往唐三身后推了推，自己往前一步，两脚分开，膝盖微屈，摆开了架势。身子里那口井静着，底下那一团也睡着，她能用的只有自己这一副身子。她把拳头攥紧了。

一只手轻轻拍在了她肩膀上。

蔡须坤回过头。

卢本韦站在她身后，手搭在她肩上，拍了一下。

``我来。''

她从蔡须坤身边走过去，迈上一级台阶，跟那姑娘面对面站着。两个人差不多高，眼睛平平地对着。

卢本韦把扽魂放了出来。

石阶上多了一头母猪。通身白里透粉，个头比寻常的猪大出一圈，四只蹄子稳稳地踩在青苔上。它晃了晃耳朵，抬起头，看了一眼头顶上那条龙。

就看了一眼。

那条龙往后缩了。

龙头从半空一寸一寸往下沉，银色的鳞一片一片贴紧了身子，贴得严严实实，像冬天里缩着脖子的鸟。绕着牌坊的那半圈龙身往回收，收到那姑娘背后，盘成了一团。那双紫色的龙眼垂了下去，不再看台阶底下的任何一个人。

光暗下去一半。

黑瘦少年手里的鱼叉当的一声掉在石阶上。短头发的姑娘捂住了嘴。高个子少年往前迈了一步，又站住了，脸上白得没有一点血色。

母猪哼了一声，往前拱了一步。

龙又缩了一寸。

\begin{center}\rule{0.5\linewidth}{0.5pt}\end{center}

那姑娘站在原地，脸上还是什么都没有。

可她垂在身侧的那只手，指头蜷起来了。

她把魂力往上催。催到一半，催不上去了。魂力还在，一点没少，可一到身后那条龙那儿，就像一道水撞上了堤，堵在那儿，过不去。龙不肯起来。她再催，龙把头埋得更低了。

压着它的，跟魂力没关系。她感觉得出来。

对面这个姑娘，年纪比她小，脸上那点稚气还没褪干净。可那头猪往那儿一站，就有一样东西从它身上漫过来，比魂力深，比魂力云老，一直漫进她的血里。她的血碰上它，自己就低了头。

她攥着的那一身本事，一样都使不出来。

一颗汗从她额角沁出来，顺着那张白得发凉的脸，慢慢往下淌。

她看着那头母猪，看了很久。

卢本韦看着她，没有往前逼。

``我们来找人，不打架。''卢本韦说，``找着了，把话说完，就走。''

母猪在她脚边卧了下来，打了个响鼻。

牌坊底下，静了很久。

先开口的是那姑娘。

``你们那条两根桅杆的船，''她说，``在我们后头跟了三天。''

唐三愣了一下。``前头有船？''

霍雨浩把笔记翻得哗哗响，从望潮镇翻到雾墙，一页一页，没有一行写着船。

``离得远。''那姑娘说，``远到你们看不见。我们看得见。''

蔡须坤没出声。

那样远的一条船，队里能摸得着的，本来只有她。可出了望潮镇以后，她的血一天比一天热，底下那一团一天比一天沉，她大半时间坐在船头闭着眼睛，跟身子里那两股力量较劲。前头海上有什么，她一点都没觉出来。

``我们停在雾外头，你们也停在雾外头。''那姑娘说，``我们进了雾，你们也进了雾。''

``你们的魂力是新的。还没压实。''

卢本韦没接话。

``十六岁，到这一步。''她看了一眼蔡须坤，``跟我差不多。''

``这世上修得比我快的，一只手数得过来。扽魂殿那位大长老算一个，一万年前的猪神算一个。剩下的，我没见过。''

``我有银龙王，修炼的地方也是最好的。你们有什么？''

她停了一下。

``修得这么快，还一路跟着我们，我只知道是什么人。拿魂兽的命，往自己身上堆。帅气教。''

``帅气教干过的事，我见过。''

她说这句的时候，声音还是清清冷冷的，可垂在身侧的那只手，又攥紧了。

小舞的脸一下子涨红了，往前蹿了一步，被霍雨浩一把拽住了袖子。

蔡须坤开了口。她声音不大，也不急。

``我们没跟着谁。往东六百里，是殿里的长辈指的路，雾里是跟着心口的热一百下一百下试出来的。''她说，``修为是闭了六个月关，功法是殿里的长辈一个字一个字教的，一头魂兽都没碰过。''

``帅气教的人，''卢本韦说，``见了我们，是要动刀子的。''

那姑娘没说话。她又看了一眼卧在石阶上的母猪。母猪眯着眼睛，打了个哈欠。

过了一会儿，她身后那团缩着的银光，慢慢收了回去，收得干干净净。后头那个高个子少年看了她一眼，手里那根杖也散了。

``古月娜。''她说。

``我不是岛上的人。''她说，``从北边来的。''

她朝身后偏了偏头。

``叶阳。薛峰。''

高个子少年把背着的手放了下来，冲他们点了一下头，算是打过招呼。圆脸的小个子从他身后探出半个身子，小声说了一句''薛峰''，又缩了回去。

``刚才，''古月娜说，``是我看错了。''

拿鱼叉的黑瘦少年弯腰把叉子捡起来，攥在手里，没再抬叉尖。

牌坊底下那股绷着的劲，一点一点松了。小舞长长出了一口气。唐三把背上的蔡兴坤往上托了托，诺当那两只张着的手，也慢慢放了下来。

蔡须坤看着古月娜，想了一会儿。

``既然都是为了对付帅气教，''她说，``不如，一起行动。''

古月娜没有马上答。她回头看了一眼叶阳。叶阳没说话，冲她点了一下头。薛峰从他身后探出脑袋，看看古月娜，又看看那头母猪，也跟着点了点头。

``好。''古月娜说。

她的眼睛又落到卢本韦脚边。母猪已经睡着了，鼻子里一下一下地出着气。

``你们队伍里的人，居然这么强。''她说，``母猪扽魂，我也是第一次见到。''

卢本韦弯下腰，在母猪耳朵后头挠了两下。母猪哼了一声，没睁眼。

\begin{center}\rule{0.5\linewidth}{0.5pt}\end{center}

两拨人在牌坊底下的石阶上坐了下来。

奥斯卡手一翻，又变出一把香肠，一人一根。递到薛峰跟前，薛峰看了看古月娜，见她没出声，才伸手接了，咬了一小口，眼睛一下子亮了。

``我们三个也是上岛来历练的。''叶阳说，``船比你们早到了一步。''他朝拿鱼叉的黑瘦少年和短头发的姑娘偏了偏头，``他们俩是岛上的孩子，这几天跟着我们。''

小舞凑过去，看着古月娜。``你多大？''

``十七。''

比蔡须坤大一岁。刚才那一会儿，两边都探过底，修为跟蔡须坤差不多。

霍雨浩把雾里那一整天讲了一遍。一百下一百下地划，热了往前，不热退回来，笔记上画满了道和圈。叶阳听得很认真，听到''只往热处走''，看了蔡兴坤一眼。

唐三讲了舵销。诺当举起那两只缠着布的手，给薛峰看。薛峰盯着看了半天，小声说了句''疼不疼''，诺当说不疼。

古月娜一直没怎么开口。等他们都讲完了，她才问了一句：

``你们这身功法，是谁教的？''

``大长老。''小舞说，脸上带着点得意，``扽魂殿的大长老。我们八个都是他教出来的，修炼堂里一个字一个字讲的。''

古月娜看着她。

``哪个大长老。''

``test\_alpha0。''

古月娜站了起来。

她脸上头一回有了东西。

``他。''她说。声音不高，可是在抖。

``帅气教在西边长了一百多年。一座堡，长成今天这个样子，人埋进了各州的城里，埋进了扽魂殿自己的殿里。''

``他是扽魂殿的大长老。这一百多年，他管过吗？''

``西边丢了多少人，没人敢出一声。老百姓的死活，他坐在殿上，看不见。''

``帅气堡，到今天还在。''

她身后那团银光又亮起来一点，一闪，又压了下去。卢本韦脚边的母猪在睡梦里哼了一声。

``你又知道他什么！''小舞喊了一声。

古月娜没看她。

``放任帅气教长成这样，漠视民生的人，''她说，``他教出来的人，我不跟着走。''

没有人答得上来。

古月娜低头看着坐在石阶上的蔡须坤。

``合作的事，''她说，``算了。''

她转过身，往坡上走。叶阳站起来，看了蔡须坤一眼，像是想说什么，最后没说，跟了上去。薛峰手里还攥着那半根香肠，看看古月娜的背影，又看看那头睡着的母猪，站了一会儿，也走了。

蔡须坤坐在石阶上，没动。她看着古月娜的背影。

那条银龙从牌坊底下升起来的时候，她就认出来了。

身子里静了一整天的那口井，轻轻晃了一下。

她站了起来。

``古月娜。''

古月娜已经走到了半坡，脚步停住了，没回头。

``你也知道天州那一战，我们损失了多少精锐。''蔡须坤说，``我又何尝不知道。''

坡上静了一下。叶阳回过头来，看着她，眉头皱起来。小舞张了张嘴，被卢本韦按住了手。

``我们''这两个字，古月娜听见了。她慢慢转过身来。

蔡须坤没再说话。她把两只手摊开，垂在身侧，闭上了眼睛。

身子最深处那口井，她头一回往上提。

先亮的是手背。几道淡金的纹从腕子底下透出来，一路往上爬，爬过小臂，爬过脖子，爬到了眼角。金光不刺眼，温温的，像冬天里晒透了的一块石头。底下那一团混沌没有动，在这座岛上，它还睡着。

满坡的鸡全抬起了头。卢本韦脚边那头睡着的母猪，睁开了眼睛。

古月娜身后，那团刚收回去的银光自己亮了起来。银龙从她背后探出半个身子，没有张爪，也没有盘起来。它把头低下去，低到了石阶上，冲着蔡须坤的方向，一动不动。

古月娜站在半坡上。

她那张一直没有表情的脸，这时候白得像纸。她看着蔡须坤身上那些金纹，看了很久，嘴唇动了几回，才发出声音来。

``居然是你。''

她想起来了。

那时候她们都在星斗大森林里。大森林最深的地方，帝天身上压着一团金光，压了一万年。那团光有一天开了口，说它要化形，要到人那边去，去找一条魂灵的路。人跟魂兽合在一块儿，谁也不杀谁的那条路。

古月娜那时候就站在边上。

她选的是另一条。

人有一句老话，叫师夷长技以制夷。人靠什么压着魂兽？功法，魂导器，一样一样造出来的东西。那她就去学。学会了，再拿人的这些东西，回过头去制住人。

帝天是她的长辈。帝天跟她说过一句话，她一直记着。

``test\_alpha0 是个大魔王。他不光杀了熊君一族。我带着星斗大森林里的魂兽去讨一个公道，反倒叫他打成了重伤。''

那一年，一个往人那边去找共存，一个往人那边去学怎么制住人。

两条路，从同一片星斗大森林里走出来，走得完全不一样。

古月娜站在半坡上，看着石阶上那个满身金纹的姑娘。

``你去找你的魂灵，''她说，``找到了他门下。''

蔡须坤身上的金纹还亮着。她抬起头，看着半坡上的古月娜。

``我记得。''她说，``那天你也在。''

古月娜从半坡上走下来，一级一级，走到她跟前两级台阶的地方，站住了。叶阳和薛峰跟在后头，没出声。

``刚才我说从北边来，''她说，``没说全。''

``我是从星斗大森林里出来的。''

唐三和霍雨浩对看了一眼。小舞张着嘴，半天没合上。叶阳和薛峰站在她身后，谁也没说话。

``你在帝天身上待了一万年，''古月娜说，``可你不是龙。''

``我不是。''蔡须坤说，``我是猪神留下来的。''

``我跟帝天，都是龙族。''

古月娜看着她身上那些金纹，停了一会儿。

``星斗大森林里是怎么讲天州的，你记得。''

``一千年前，星斗大森林跟人类签过约。人不进大森林深处，魂兽不出大森林，互不相犯。''

``约是人先破的。熊君的儿子叫人抽干了，两百多头熊在洼地里堆成了山。破约的是哪一伙人，星斗大森林里分不清，也不想分。在大森林看来，人类魂师的头，就是扽魂殿。''

``帝天带着星斗大森林的魂兽去讨这笔账。讨到天州城门口，熊君死了，十万同族没了，帝天挨的那一剑，在大森林最深的地方化了五十年。''

``这些年，星斗大森林里叫他伤过的，数都数不过来。再让他这么砍下去，魂兽这一族，迟早断在他手里。''

``星斗大森林里讲到这儿，就讲完了。''

卢本韦开了口。

``大长老讲的，是另一头。''

``兽潮压到城下，总殿两位长老赶去天州，当是寻常一头十万年魂兽作乱，到了城外不问缘由，见着魂兽就杀，激怒了熊君，两个都死在它的爪下。城门裂了，城里的人往东跑。大长老落下去的时候，不知道熊君的儿子是怎么死的。熊君冲他说'你们这种人'，他到后来才听懂。''

``打完那一仗，他没回天斗城，一路往西，到帅气堡门口去拿人。''

``他不知道。''古月娜说，``所以就可以杀十万？''

没有人答得上来。

蔡须坤身上的金纹还亮着。她往上走了一级台阶，跟古月娜隔着一级站着。

``我是帝皇瑞兽的本源。''她说，``天州那一回，我就在帝天身上。''

叶阳的眉头一下子皱紧了。薛峰从他身后探出头来，嘴张着，忘了合上。拿鱼叉的黑瘦少年手一松，叉子又掉在了石阶上。

``星斗大森林里是怎么起的兵，我比谁都清楚。''蔡须坤说，``帝天，熊君，万妖王，赤王，都在兽潮里。后头跟着二十万魂兽大军。''

``天州前头那两座小城，兽潮过去，都破了。万妖王和赤王各带着一支兽群，从两侧丘陵逼近。帝天被一剑逼退以后，万妖王想用藤蔓缠住大长老，藤蔓当场就被魂力震碎。赤王趁势扑上去，挨了一拳，飞出去数里。两支兽群也没能完成合围。''

``大长老落进天州城里的时候，帝天下的令，是踏平这座城。''

``这哪里是来讨公道的。这是来报仇的。''

``多年的和平，就这么破了。''

古月娜站在她上面一级，脸白得像纸，一个字也没说出来。

``两边讲的，我都听过。''蔡须坤说，``星斗大森林里讲的是真的，殿上讲的也是真的。谁站在哪儿，就先看见哪一截。''

她顿了一下。

``那笔账是帅气教欠下的。一百六十八年，记在了扽魂殿头上。''

古月娜站在上面那一级台阶上，很久没说话。

她身后那条银龙，还伏在石阶上，没有起来。

过了一会儿，她开了口，眼睛没离开蔡须坤。

``我想跟你单独说几句。''

卢本韦看了蔡须坤一眼。蔡须坤点了点头。

\begin{center}\rule{0.5\linewidth}{0.5pt}\end{center}

两个人顺着石阶往下走，一直走到滩上。

潮已经退了，白沙上留着一道一道水印子，舢板搁在高处，船底朝着天。满滩的鸡跟着她们挪了几步，又停下了。

蔡须坤身上的金纹慢慢暗下去，收回了皮肤底下。古月娜身后那条银龙，也一点一点淡了，散进了海风里。

两个人在舢板边上站着，看着海上那堵雾。

``兽潮那一次，''古月娜先开了口，``魂兽高层，确实误会了扽魂殿。''

蔡须坤没说话。

``熊君儿子那具壳子，帅气教每打一场仗，就摆出来一回。星斗大森林里的魂兽看过一回，就什么都明白了。''

``可那时候，前头两座城已经破了，熊君已经倒在天州城门外头了。''

她停了很久。

``这场战争，没有罪人，也没有真正的赢家。''

蔡须坤看着那堵雾，过了好一会儿才开口。

``那一剑扎进来的时候，我就在帝天身上。''

古月娜没出声。

``大长老身后，是天州城里的百万百姓。帝天的龙神爪已经抬起来了，那是杀招。换了谁站在那扇城门底下，都只能那样接。''

``可那时候，他自己的力量，他是控制不住的。''

``iblp 他用得了，还没用到家。再掺上那一丝 Laver 表，那一剑的分量，极有可能已经越过了 PTO(Z\_ω)。那一层，他自己到今天才走到。''

``一百六十八年前的他，握着一样比自己大得多的东西。剑落下去，铺多宽，铺多深，由不得他。''

``十万魂兽灰飞烟灭，''她说，``也由不得他。''

海风从滩上吹过去，把古月娜的头发吹起来一缕，又落下了。

``可你在台阶上说的那句话，''蔡须坤说，``我也认。''

``魂兽这一族要是真断了，那也是可怕的结果。一万年前那一炷香，就白烧了。''

``我去人那边，找的是两边都活得下去的那条路。''

古月娜转过头来，看了她很久。

当年那团金光离开帝天，古月娜一直以为，它是去了人那一边。这会儿她才看明白，它走的那条路，两头都连着。

``罪人有。''蔡须坤说，``不在这场仗里。在西边。''

浪一下一下拍上来，又退回去。两个人在舢板边上站着，谁也没再开口，一直站到太阳落进了海里。

\begin{center}\rule{0.5\linewidth}{0.5pt}\end{center}

两个人回到牌坊底下的时候，天已经黑了。

石阶上生了一堆火。奥斯卡蹲在火边烤香肠，诺当坐在他旁边，两只缠着布的手搁在膝盖上。叶阳和薛峰没走，坐在火堆另一头，跟小队隔着一截，谁也没先开口。拿鱼叉的黑瘦少年和短头发的姑娘已经回家去了。

火光照过来，所有人都抬起了头。

古月娜在火边站住了。她看了一圈，看唐三，看小舞，看诺当那两只手，最后看着卢本韦脚边那头又睡着了的母猪。

``白天那些话，''她说，``是我说重了。''

小舞别过脸去，哼了一声，没接话。霍雨浩在旁边碰了碰她的胳膊。

``你们那位大长老，我还是信不过。''古月娜说，``可你们，我信。''

她停了一下。

``合作的事，还算数吗？''

蔡须坤看了一眼卢本韦。卢本韦点了点头。唐三把手里的柴往火里送了送，说：``算数。''

小舞憋了半天，最后说：``算数就算数。下回别拿龙吓人。''

叶阳笑了一下。这是他上岛以来头一回笑。薛峰从他身后挪出来，挪到奥斯卡边上，眼睛盯着火上那几根香肠。奥斯卡把烤得最好的一根递了过去。

``明天一早，''蔡须坤说，``一块儿上山，找大供奉。''

``一块儿。''古月娜说。

她在火边坐了下来，就坐在蔡须坤旁边。

火堆不大，十一个人围着，刚好坐满一圈。

\begin{center}\rule{0.5\linewidth}{0.5pt}\end{center}

\chapter{第一百零三章：最菜萌新}\label{ux7b2cux4e00ux767eux96f6ux4e09ux7ae0ux6700ux83dcux840cux65b0}

鸡是天不亮就叫的。

先是牌坊顶上那只石头鸡旁边的一只芦花公鸡，后来是坡上，再后来是满岛。叫声一层一层往上走，走到山顶上，又从山顶上一层一层落回来。

火堆早灭了，灰还是温的。十一个人一个挨一个从石阶上爬起来。

奥斯卡手一翻，变出一把香肠，挨个分。分到薛峰手里那根，他想了想，又多变了一根塞过去。薛峰攥着两根香肠，看了看古月娜，见她没出声，低下头小口小口地吃。

``你这扽魂，''叶阳在旁边看着，``能变多少？''

``看饿不饿。''奥斯卡说，``饿了就多。''

叶阳没听懂，看了他一眼，也接了一根。

古月娜站在牌坊底下，把头发用一根素白的带子束了起来。蔡须坤走过去，在她边上站了一会儿。

``睡着了吗？''

``没怎么睡。''古月娜说，``你呢？''

``睡着了。''蔡须坤说，``十六年，头一回睡得这么沉。''

古月娜看了她一眼，没说什么，看向坡上那条往上去的路。

``走吧。''

\begin{center}\rule{0.5\linewidth}{0.5pt}\end{center}

牌坊后头的路是土的，踩实了，两边是田。

田埂上有鸡，田里也有鸡，屋顶上蹲着鸡，墙头上站着鸡。昨天从滩上往上看，只觉得鸡多；走进来了才知道，这座岛上的鸡，比人多出不知多少倍。

家家户户门口都有鸡舍。鸡舍是石头垒的，跟住人的屋子一样，屋顶也压着石块，有的比旁边的人屋还高。一户人家门口，鸡舍新刷了一层白灰，人屋的墙皮倒掉了一大块，没补。

霍雨浩走一路，记一路。

``头一户，鸡舍比屋大。''他嘴里念着，``第三户也是。第五户，屋里没人，鸡舍里点着灯。''

``点灯做什么？''小舞凑过去看。

``不知道。''霍雨浩说，``记下来，回头问。''

走过一个岔路口，路边立着一座小石龛，半人高，龛里供着一只石头刻的鸡，鸡跟前摆着一溜鸡蛋，一个挨一个，有的还沾着草屑，是今天早上刚放上去的。

一个七八岁的孩子捧着一只蛋跑过来，踮着脚把蛋放在最边上，冲石鸡拜了一拜，又跑走了。

``头一个蛋，先给坤神。''古月娜说。她看着那座小龛，``北边也有人这么讲究，可没这么讲究。''

奥斯卡凑到小舞边上，压着嗓子：``你说，岛上的人吃不吃鸡？''

话音没落，路边一户人家门口坐着的老头抬起眼皮，看了他一眼。那一眼看得奥斯卡后脖颈子一凉，他把嘴闭上了，一路再没提这个字。

人是有的。跟昨天不一样，今天一路上的门都开着，门口站着人。端着簸箕撒谷子的老太太，蹲在鸡舍边上捡蛋的孩子，扛着锄头往田里去的汉子。看见他们这一大队人过来，都停下手，看着，不说话，也不躲。

他们看得最久的，是蔡兴坤肩上那只黑子鸡。

黑子鸡今天没缩着。它蹲在蔡兴坤肩上，脖子伸得笔直，一路走一路看。走过一户人家门口，那家鸡舍里的鸡一齐把头抬了起来；再走过一户，又一齐抬了起来。抬起来了，就那么看着，一声不叫。

``它们在看什么？''诺当低声问。

蔡兴坤没答。她心口那股热，上了岛以后一直是静的，今天早上起来又有了一点，温温的，不烫。

蔡须坤走在她旁边。身子里那口井也是静的，底下那一团还睡着。这些天在海上，她跟那两股力量较了一路劲，眼前什么都看不清。上了这座岛，那股劲一松，她才觉出来，四周的东西一样一样地清楚了：风从哪边来，坡上哪棵树底下蹲着几只鸡，三户人家以外那个孩子手里的蛋还是热的。

清楚得她自己都有点不习惯。

\begin{center}\rule{0.5\linewidth}{0.5pt}\end{center}

问路是奥斯卡去问的。

他走到一户人家门口，那家老太太正端着簸箕撒谷子。奥斯卡变出一根香肠，双手递过去。

``大娘，跟您打听个人。''

老太太看了看香肠，没接。

``岛上有位大供奉，''奥斯卡说，``叫最菜萌新，您知道往哪儿找吗？''

老太太手里的簸箕停了。

她把奥斯卡从头到脚看了一遍，又越过他的肩膀，看了一眼后头站着的那一队人。眼睛在蔡兴坤肩上停了一下，停得比别处长。

``往上走。''她说，``走到没有田的地方，就是山脚。''

``山脚就能见着？''

``山脚见不着。''老太太又把簸箕端了起来，``山脚有人守着。''

``守着的是谁？''

老太太没再答，转身撒她的谷子去了。

奥斯卡站在那儿，把那根香肠又揣了回去。

他们往上走。一路上又问了三回。头一回问的是个扛锄头的汉子，汉子往上一指，说''山脚''；第二回问的是个捡蛋的孩子，孩子也往上一指，说''山脚''，说完就跑了；第三回问的是坐在门槛上编草鞋的老头，老头连头都没抬。

``山脚有人守着。''老头说，``你们过不去。''

``为什么过不去？''小舞问。

老头把手里那根草茎一折，搓了搓，接着编。

``那是坤神的地方。''

卢本韦在老头跟前蹲了下来。

``老人家，跟您打听个旧事。''她说，``一百多年前，有个大陆上来的人，姓卢，在岛上住过三年。您听说过吗？''

老头编草鞋的手停了。

``姓卢的。''他说，``我爷爷那辈的事了。''

``您爷爷提过他？''

``提过。''老头说，``说是个年轻后生，话少，天天一早起来帮人喂鸡，喂完了，就坐在山脚底下写东西。写了三年，走了。''

他抬手往上指了指。

``山脚边上有间石屋，空了一百多年，没人住。说是他住过的。''

卢本韦蹲在那儿，半天没站起来。奥斯卡在她身后站着，等她。

过了一会儿，她站起来，冲老头弯了弯腰，往上走了。走了几步，奥斯卡追上去，跟她并排走，没说话。

\begin{center}\rule{0.5\linewidth}{0.5pt}\end{center}

田一层一层往上，越往上越窄，到后来，最上头那一层田边上，只剩一溜菜畦，再往上，就没有田了。

前头是山。

山不算高，可是陡。一整面坡上长满了树，树底下是一溜石阶，从山脚一直往上铺，铺进了半山腰的雾里。雾是白的，不厚，跟海上那堵雾墙不一样，是飘着的，一缕一缕挂在树梢上。

石阶最底下那一级，立着两根石柱。没有门，没有匾，两根柱子中间空着，柱子顶上各蹲着一只石头鸡，比牌坊上那只刻得细，翅膀上的羽毛一根一根都看得出来。

石柱左边几丈远的地方，真有一间石屋。矮矮的，屋顶压着石块，门板歪了，门口长了半人高的草。窗台上落着一层鸡毛，几只鸡在门槛上蹲着，像是常来。

卢本韦看着那间石屋，脚底下慢了下来。

曾祖笔记里补的那几十页，写的都是魂灵。怎么合，怎么养，怎么不烧死人。关于这座岛，他几乎一个字都没写。可那几十页，是在这间屋子跟前写出来的。

两根石柱当中，站着一个女人。

谁也没看见她是从哪儿来的。

她穿一身深蓝的长衣，蓝得发黑，像入夜以后的天。衣裳上星星点点缀着些银白的小点，不知道是绣上去的还是生在布里的，日头底下看，一点一点地亮。头发也是黑的，挽着，插一根素银簪子。看年纪说不准，三十上下，也可能四十。脸很静，眼睛也静，像夏天夜里没有风的那种静。

她就那么站在两根石柱当中，两只手拢在袖子里。

蔡须坤的脚步停住了。

她身子里那口井，没有晃。可井水底下，有什么东西往下沉了一沉。

她见过的强者不算少。梅长老是九十七级，挡在她们前头的时候，背影像一堵墙。真霜长老宫变那一夜骑在冰熊王头上，一百五十米的熊瘸着一条腿，还把胡帅钉在宫门前钉了一整夜。phyrion 九十八级，宫变那一夜脚边跟着一头家猪那么大的白猪，那头猪往阵上一拱，一座城的大阵就碎了。

眼前这个女人，跟他们都不一样。

她身上什么都没放出来。没有魂力云，没有领域，没有一丝魂力往外漏。可蔡须坤站在离她二十来丈的地方，就觉得那两根石柱当中站着的，是一整座山。

封号斗罗没有这个分量。比那个深得多。

古月娜走在蔡须坤边上，也停了。蔡须坤用余光看见，她垂在身侧的手指头动了一下，又按住了。她身后有一点银光要往外冒，叫她硬生生压了回去。

后头的人也都觉出来了。诺当往前挪了半步，挡在蔡兴坤前头。唐三的手垂在身侧，五指张了张，又收拢。薛峰整个人缩到了叶阳背后，叶阳没回头，伸手把他往后拨了拨。

那女人开了口。声音不高，也不冷，一个字一个字，很清楚。

``这里生人勿进。请告诉我们，你们来这里干什么。''

``我们''两个字，她说得很自然。可两根石柱当中，只站着她一个。

山上的雾动了一下。

\begin{center}\rule{0.5\linewidth}{0.5pt}\end{center}

蔡须坤往前走了几步。

``天斗城来的。''她说，``扽魂殿大长老 test\_alpha0 门下，八个人。来找最菜萌新大供奉。''

那女人看着她，没有说话。

``大长老两百多年前来过这座岛。''蔡须坤说，``他说，大供奉认得他。''

卢本韦从队伍里走了出来，站到蔡须坤边上。

``我曾祖，卢承业。''她说，``一百三十六年前，在岛上住过三年。''

那女人的眼睛这才动了一下。

古月娜也往前走了一步。``我们三个，''她朝叶阳和薛峰偏了偏头，``是来岛上历练的。昨天在滩头碰上了他们，一块儿来的。''

那女人没看古月娜。她的眼睛从蔡须坤脸上挪开，落到了卢本韦身上。

看了一会儿，她从两根石柱当中走了出来，往下走了三级台阶。

离得近了，能看见她衣裳上那些银白的小点，真的在一闪一闪地亮，跟天上的星星一样，有的亮一点，有的暗一点。

她在卢本韦跟前站住了。

``你身上，''她说，``有一头母猪。''

卢本韦没动。她的扽魂没放出来，一丝都没放。

``你没放，我也闻得见。''那女人说，``这股味道，这座岛上上一回闻见，是一百三十多年前的事了。''

卢本韦的嘴唇抿紧了。

那女人没再问她。她转过头，看向卢本韦身后。

看蔡须坤，看蔡兴坤，看蔡兴坤肩上那只黑子鸡。

黑子鸡蹲在蔡兴坤肩上，脖子伸得笔直，跟她对看着。

她看得很久。看到后来，她那张一直很静的脸上，眉头一点一点皱了起来。

``你们两个，''她说，``把手伸出来。''

蔡须坤和蔡兴坤对看了一眼，把手伸了出去。

那女人没碰她们。她只是低下头，看着那两只手。蔡兴坤的袖口底下，那几道淡金的纹自己透了出来，一闪，一闪。蔡须坤手背上，也浮起了一道，很淡，跟昨天放本源时那种金纹不一样，颜色更暖一点，像早上刚出来的太阳。

蔡须坤自己都愣了一下。

那女人盯着那两道纹，看了很长时间。

``你们的血，''她说，声音比刚才低了，``是从哪儿来的？''

蔡兴坤张了张嘴，答不上来。

蔡须坤也答不上来。

那女人没等她们答。她退后一步，把拢在袖子里的手抽了出来。

她抬起右手，两根手指并着，往山上一指。

指尖上亮起一点光。银白的，很小，像天上掉下来的一颗星。那点光离了她的指尖，往山上飞，飞进半山腰的雾里，不见了。

过了一息。

山上有一只鸡叫了一声。

又一息，再往上，又一声。一声接一声，顺着那一溜石阶往上传，传过半山腰，传进雾里，一直传到看不见的山顶上。

``我叫夏夜星空。''那女人这才说，``坤神座前四大守卫，我是其中一个。''

她说''我们''的时候，原来说的是她们四个。另外三个，蔡须坤一个都没看见，可她这会儿忽然觉得，山上那片雾里，不止一双眼睛。

夏夜星空的目光从蔡须坤身上挪开，在古月娜脸上停了一下。

``你身上那条龙，''她说，``收好。''

古月娜的脸色变了变，没说话，把手攥紧了。

她回过头，看着蔡须坤。

``大供奉马上就到。''她说，``在他来之前，谁也别往上走一步，也别往下走一步。''

山顶上，最后那一声鸡叫落下来。

然后，山上的雾，从当中分开了一条缝。

\begin{center}\rule{0.5\linewidth}{0.5pt}\end{center}

石阶上有人下来了。

走得不快，一级一级的。先看见的是一双布鞋，鞋面上沾着山上的露水，后来是一件素色的长袍，洗得发白，可穿得一丝不乱，再后来，是一张老人的脸。

很老。头发全白了，在脑后挽得整整齐齐，胡子也是白的，不长，修得很短。脸上的褶子一道压一道，比大长老还要深。眼皮半垂着，看人的时候不怎么抬。

他每往下走一级，山上的雾就往两边让开一点。石阶两边树底下蹲着的鸡，一只一只把脑袋低了下去。

夏夜星空退到石柱边上，弯下腰。

``大供奉。''

老人嗯了一声，从她身边走过去，在最底下那级台阶上站住了。两只手背在身后，抬起眼皮，往这一大队人身上看了一遍。

被他这么一看，蔡须坤忽然想起大长老。两个人长得一点都不像，那股劲儿却像。站在那儿，什么也不放，可你就知道，这个人身后站着的是什么。

眼睛在蔡须坤和蔡兴坤身上停了一下，停得很短，像是记下了，没往下说。在古月娜身上也停了一下。

然后他看见了卢本韦。

他的眼皮抬起来了，抬得比刚才高。

他看着卢本韦，看了很久。看着看着，眉头皱了起来，像是在一大堆旧东西里翻找什么，翻了半天，没翻着。

``这孩子，''他说，声音沙沙的，``我在哪儿见过。''

``您没见过我。''卢本韦说。

``见过。''老人摇了摇头，``脸没见过。这双眼睛见过，还有站着的这个架势。站在那儿，跟一根桩子似的，谁推都不动。''

卢本韦往前走了一步。

``我的曾祖，是卢承业。''

老人的手在半空停了一下。

``卢承业。''他念了一遍这个名字，念得很慢。

``一百三十六年前，他在您这儿住了三年。''卢本韦说，``山脚那间石屋。''

老人没说话。他转过头，看了一眼石柱左边那间矮矮的石屋，门口的草长了半人高。

``天天一早起来喂鸡。''他说，``喂完了，坐在门口写东西。写一页，撕一页，撕了三年，最后带走的就那么几十页。''

``那几十页，在我手上。''卢本韦说。

老人回过头来，看着她。

``他还在吗？''

``不在了。''

老人点了点头，没再问。他背在身后的那只手，指头慢慢收拢了。

卢本韦从怀里摸出一封信，双手递过去。

``这是扽魂殿大长老 test\_alpha0 写给您的。他让我当面交给您。''

老人看了看那封信。信封上没写字，封口压着一方小小的朱印。

他伸手接了过来，看了一眼那方朱印。

``test\_alpha0。''他说，嘴角动了一下，``两百多年前他来的时候，刚坐上那个位子没几年，还是个毛头小子。站在我这山脚底下，问东问西，问了三天。''

他把封口揭开，抽出一张纸，展开来看。

纸不大，字不多，一笔一画写得很稳。

萌新兄：

一别两百余年，兄安否。当年兄在山脚下讲的那一万年，我推了一辈子，到今年才推明白一半。iblp 的良序化，今年刻完了。兄若还记得那个问了三天的毛头小子，他如今也三百一十七岁了。

持信者八人，是我修炼堂里带出来的孩子，自称果糕小队。闭关六个月，都过了 0 111 22。此番渡海，一为历练，二为魂灵那件事。卢家那孩子是承业的曾孙女，承业当年在兄处住过三年，兄当记得。

孩子们年轻，岛上的规矩不懂，冒犯之处，望兄担待。历练的事，兄看着安排，该吃的苦，让他们吃。他们若问起一万年前的事，兄看着答。

另有两个孩子身上的东西，兄见了自然明白，我不多写。

test\_alpha0 顿首

老人看完了，没说话，又从头看了一遍。

看到最后那一行，他抬起眼皮，往蔡须坤和蔡兴坤那边看了一眼，又把眼皮垂了下去。

他把信折好，揣进袖子里。

``毛头小子，''他说，``也老了。''

他这才把眼皮抬起来，看向蔡须坤和蔡兴坤。

这一回没有一扫而过。他看得很慢，从蔡须坤的脸看到她的手，又从蔡兴坤的手看到她肩上那只黑子鸡。黑子鸡蹲着，脖子伸得笔直，跟他对看。

``你们两个，叫什么？''

``蔡须坤。''

``蔡兴坤。''

老人点了点头，把这两个名字在嘴里过了一遍。

``蔡须坤，蔡兴坤。''他说，``你们身上，有一定的坤神血脉。''

山脚下一下子静了。小舞捂住了嘴。古月娜转过头，看着蔡须坤。

``不算多。''老人说，``可这些年，这座岛上没出过这样的血。坤神的继承人，你们两个，是眼下最好的人选。''

他停了一下。

``既然来了，那就准备好接受坤神的考验。''

``这期间，谁放弃了，就永远失去继承人的资格。没有第二回。''

他抬起手，往身后那座山指了指。山不高，可陡，一溜石阶从山脚一直铺进雾里。

``这座山，叫中分山。坤神殿在山上，我就住在那儿。''

``两年之内，你们两个登上中分山，找到坤神殿，算是第一考。过了这一考，你们才能真正接受坤神的传承。''

蔡兴坤肩上的黑子鸡，这时候叫了一声。很短，像是应了一下。

蔡须坤没有马上说话。她抬头看着那一溜石阶。石阶往上走，走进雾里，雾那头什么都看不见。从这儿看过去，这座山也不算高。可大供奉说的是两年。

老人看着她，又开了口。

``还有一件。''

``你们队里，有一个人，肉身已经没了。''

蔡须坤的背一下子绷直了。

背上那一处，温了一下，比这些天哪一回都温。

老人的眼睛落在她背上，像是隔着衣裳也看得见。

``他还在。''老人说，``魂还在，在你身上。可光有魂，回不来。''

``想让他回来，坤神的力量，是一个关键。''

队伍后头，奥斯卡的手伸进了怀里，攥住了那根拿油纸包了三层的香肠。攥得很紧，油纸都皱了。诺当站在他边上，没说话，把一只缠着布的手搭在了他肩上。

蔡须坤把一只手背到身后，按在那一处上。

那一处在她掌心底下，又温了一下。

她抬起头，看着那座山，看着山上那片雾。

\begin{center}\rule{0.5\linewidth}{0.5pt}\end{center}

大供奉没再往下说。他转过身，冲石柱边上的夏夜星空抬了抬下巴。

``山脚那间屋子，收拾出来，让他们住。''他说，``缺什么，你往下送。''

``是。''

``那三个也住下。''他朝古月娜那边偏了偏头，``上岛来历练的，就是岛上的客。''

古月娜愣了一下，弯腰行了一礼。

大供奉背着手，转身往石阶上走。走得跟下来的时候一样慢，一级一级的。雾在他前头又让开了一道缝。

人这才散开。奥斯卡头一个跑到石屋跟前，弯腰去拔门口那半人高的草。拔一把，就有一只鸡从草窝里扑棱着翅膀蹿出来，蹿出来四五只，他才拔出一条路。唐三去扶那扇歪了的门板，门轴锈死了，他托住门板往上一抬，吱的一声，门正了，门框上簌簌落下一层灰。小舞捂着鼻子，拿手往外扇。

卢本韦站在门口，没进去。屋里有一张矮桌，桌面上的灰积了一指厚。

蔡须坤没有跟过去。

她站在两根石柱当中，看着那个白头发的背影往上走。走到第七级，她开了口。

``大供奉。''

老人站住了，没回头。

``有件事，想跟您打听。''

``说。''

蔡须坤往身后看了一眼。石屋那边，奥斯卡还蹲在地上拔草，霍雨浩蹲在门口，拿炭条往本子上记着什么。古月娜站在石屋的墙根底下，背对着这边，头却微微偏着。

蔡须坤把眼睛收了回来。

``一百六十八年前，天州那一仗。''她说，``您知道多少？''

老人慢慢转过身来。他站在高处，低头看着她，眼皮还是半垂着。

``该知道的，都知道。''

``那一天，我在帝天身上。''

老人点了一下头，脸上没有一点意外。

``那一剑落下来，我都看见了。''蔡须坤说，``昨天在滩上，我跟古月娜说，换了谁站在那扇城门底下，都只能那样接。''

她停了一下。

``话是我说的。可我自己心里，一直没底。''

``他两百多年前来过您这儿，在山脚底下问了三天。问的是一万年前的事，问的是魂灵。''她仰着头，``问过三天魂灵的人，到了城门底下，怎么就只剩下那一剑了？''

山脚下静了一会儿。石阶两边的树底下蹲满了鸡，一只也没叫。

老人从那一级台阶上往下走了两级，站住了。

``那一战，''他说，``魂兽和他，都太冲动了。''

蔡须坤没出声。

``帝天以为，自己早就走过了极限斗罗那一层，天底下没有他接不住的人。''老人说，``可他碰上的，是 test\_alpha0。''

``那孩子的实力，我非常清楚。两百多年前他站在这儿，我就看出了七八分。这些年他走到了哪一层，我也大概有数。''

他顿了顿。

``除非有神级强者，否则谁也不可能真正击败他。''

蔡须坤想起那一天。帝天化作一道黑光往天州城扑过去，她在帝天身上，觉得出他那股劲。那股劲一直往前顶，一点往后收的意思都没有。他是真觉得自己能赢。

``帝天那天带着兽潮压到天州城下，没想过这一层。''老人说。

``那他呢？''

``他也是迫不得已。''老人说。

``他赶到的时候，天州前头那两座小城，已经成了废墟。急报上写着，总殿派去的两位长老死在了城外，连熊君跟前都没杀到。他落进城里，城墙上挂着血，满街的百姓拖家带口往东跑，有的跑着跑着倒下去，再没爬起来。''

``两位长老殉了职，一城百姓遭着这样的罪，他心里是有火的。''老人说，``带着这股火站到城门底下，那一剑下去，就收不住了。''

他背在身后的两只手松开了，垂下来，拢进了袖子里。

``魂灵的事，他也不是没考虑过。''

``两百多年前那三天，他在我这山脚底下问来问去，问的就是这个。回去以后，他一直搁在心上。''

``谁知道兽潮来得这么突然。他这边还没想明白，星斗大森林的兵，已经压到城下了。''

蔡须坤的嘴唇动了一下，没说出话来。

``他的身后，是百万民众。''老人说，``他坐的是大长老的位子。坐上那个位子，守护苍生的责任，就压在他一个人身上。''

蔡须坤低下了头。

她想起那天从帝天身上看出去的那一眼。城门底下站着一个人，手里提着剑，站得很直。他身后那扇城门裂开了一道缝，从缝里看进去，满街的人都在往东跑，有人背着人，有人抱着孩子，有人光着一只脚。

那时候她眼睛里只有门底下那个人。门缝里那些往东跑的人，她到今天才想起来。

``你在帝天身上待了一万年，又在他门下待了这些年。''老人说，``两头，你都站过。那一仗的账，往后你自己慢慢算。''

他顿了一下，往山下石屋那边看了一眼。卢本韦还站在门口，手扶着门框。

``魂灵的事情，我会教你们。''

蔡须坤抬起头来。

``但是坤神的考验，你们必须自己完成。''老人说，``谁也替不了，我也不行。''

``明天开始挑战吧。''

他转过身，接着往上走。走到雾跟前，雾往两边让开，他走了进去。雾在他身后又合上了，合得严严实实，像从来没开过。

\begin{center}\rule{0.5\linewidth}{0.5pt}\end{center}

蔡须坤在两根石柱当中站了好一会儿，才转身往石屋那边走。

古月娜还站在墙根底下。这回她没背着身子，脸朝着这边。

两个人隔着几步远，站住了。

``他说的那句，''古月娜先开了口，``魂兽太冲动。''

``嗯。''

古月娜低头看着脚底下的草，看了一会儿。

``帝天那会儿，是真以为自己打得过。''

她就说了这一句，转身进了石屋。

\begin{center}\rule{0.5\linewidth}{0.5pt}\end{center}

太阳偏西的时候，石屋收拾出来了。

屋子不大，一间半。外间一张矮桌，一口灶；里间一铺土炕，炕上铺的草烂成了灰，奥斯卡抱出去扔了，换上了新割的干草。夏夜星空下山送过一回东西，几床被褥，一袋米，一篮子鸡蛋，放下就走了，一句话也没说。

那篮鸡蛋，奥斯卡数了三遍，四十二个。他蹲在篮子跟前看了半天，到底没敢问能不能吃。

卢本韦在矮桌底下摸出了一方砚台。石头的，边角磕掉了一块，砚池里干透了，结着一层黑。她拿袖子擦了擦，擦不掉，就那么捧着，在门槛上坐了下来。

一百三十六年前，有个人天天坐在这道门槛上写东西。

天黑下来，屋里点了灯。

蔡兴坤坐在石屋外头一块石头上，黑子鸡缩在她怀里睡着了。蔡须坤走过去，在她边上坐下。

山上那片雾，到了夜里也没散。雾里很高很高的地方，亮着一点光，很小，说不清是星星，还是坤神殿里的灯。

``两年。''蔡兴坤说。

``两年。''蔡须坤说。

背上那一处，温了一下子

\chapter{第一百零四章：bms提升效应}\label{ux7b2cux4e00ux767eux96f6ux56dbux7ae0bmsux63d0ux5347ux6548ux5e94}

夏夜星空到石屋时，奥斯卡正举着一个鸡蛋。

``能吃。''她说。

奥斯卡立刻把蛋磕进锅里。小舞笑着去拿碗，诺当伸手接，碰到掌心的伤，又把手缩了回去。

十一人吃过早饭，夏夜星空在桌上摊开一卷山图。最高处写着二千五百，最底下的横线上写着五百。

``都是米，算往上升了多少。前五百米，要把这一层重新分析清楚。''

\NovelDisplay{M01159R01}{%
(0)(1,1,1)(2,2)%
}

``这一层，我们的魂力已经过了。''蔡须坤说。

``力量到了。''夏夜星空点头，``你们用投影序数走过了那道界。可这段里面，每一步为什么对应这个序数，哪一处是乘，哪一处是改参数，分析还没有走全。山上补这一课。''

蔡须坤看着那条横线。身子里的井静着，底下那团混沌还在睡。她把山图卷好，递给霍雨浩。

``从哪里开始？''

夏夜星空在两根石柱中间写下一行字。

\NovelDisplay{M01160R01}{%
(0)(1,1,1)=\psi(\Omega_\omega)%
}

``从你们早就用熟的地方。''

夏夜星空让十一人都带上纸笔。蔡须坤、霍雨浩、卢本韦和古月娜先到前排试阶，其余七人留在相邻横台复算，轮流将不同的父项标出来。前排看清受压的位置，再交给后排核对；换人登阶时，每个人仍须自己接完那段运转。唐三收绳、奥斯卡送水，都在各自的练习间隙进行。

霍雨浩拿炭条在石上分出三行。除第一行的数必须保留外，每列后面补齐用的零都省略；若中间的零之后还有非零坐标，仍保留它的位置。读父项和阶差时，须把省略的行按零补回。

\NovelDisplay{M01161R01}{%
a_1=\text{第一行},\quad a_2=\text{第二行},\quad a_3=\text{第三行};%
}
\NovelDisplay{M01161R02}{%
\quad \text{省略的末尾坐标按零参与计算。}%
}

霍雨浩在本子上写下记号约定。最外面的 \allowbreak{}\NovelInline{I05132}{\psi}\allowbreak{} 是可数层坍缩，函数下标按坍缩上界书写。

\NovelDisplay{M01162R01}{%
\psi=\psi_\Omega,\qquad\psi_{\Omega_2},\ \psi_{\Omega_3},\ \ldots\quad\text{均以上界作函数下标。}%
}

\NovelDisplay{M01163R01}{%
\psi(0)=1,\qquad\psi(\Omega)=\varepsilon_0,\qquad\psi_{\Omega_2}(0)=\Omega%
}

教材里见到旧编号 \allowbreak{}\NovelInline{I05133}{\psi_1,\psi_2,\ldots}\allowbreak{} 时，只按下面的对照换写，后续计算统一使用上界下标。

\NovelDisplay{M01164R01}{%
\psi_1=\psi_{\Omega_2},\qquad\psi_2=\psi_{\Omega_3},\qquad\ldots%
}

``下标换成这个写法，参数别跟着乱改。''古月娜说，``\allowbreak{}\NovelInline{I05134}{\Omega_2}\allowbreak{} 本身，仍是 \allowbreak{}\NovelInline{I05135}{\Omega_2}\allowbreak{}。''

遇到后继参数，要按所在的 BOCF 结构核对有限相加的基本列，不能一律沿用 MOCF 中``再起一层指数塔''的读法。霍雨浩把这一句也记在页首。

长式中的无限列和无限指数塔都是基本列极限的示意，不是一条真的写有无限多列的有限 BMS。具体到每次运转，只能使用有限项。

蔡须坤的右脚刚落上第一阶，肩背便往下一沉。她抬左脚，膝盖只抬到一半，胸口那口气就被压了回来。

卢本韦伸手托住她的胳膊，自己跨上同一阶，也弯了腰。霍雨浩扶着石沿上来，炭条在指间折成两截。古月娜最后落脚，银发贴在颈侧，呼吸顿了一顿。

他们都有超过这层的魂力。蔡须坤一口气提起自己修来的力量，脚底的纹却只亮了一半，几处原本连贯的运转在压力下互相牵扯。她越想一并带过去，胸口便越闷。

``慢一点。''卢本韦说，``先把最初那段留住，只动后面。''

蔡须坤依言收慢。沉重的阻力将运转拖得极缓，平时一掠而过的接缝，这时竟能一处处分辨。前段仍然稳着，后面的乘数却在另一处受阻。

霍雨浩蹲下，把受阻前后两种运转分开记。古月娜跟着试了一遍，抬手指向他圈出的右括号。

``两处的先后不同。她刚才把它们当成一层了。''

他们照基本列重新核对，再试。接缝处的牵扯松了，蔡须坤终于能把左脚移上来；肩上的重量仍在，汗从下颌滴到石阶上。

``压力把你们省过去的细节压出来了。''夏夜星空说，``看清一处，再分析一处。''

霍雨浩又在记录旁留了空白。身体察觉的差异，还要拿父项、阶差和对应式验证，不能只凭觉得顺畅就填等号。

唐三在后面系好蓝银草，奥斯卡把水和补给递过来。蔡须坤、霍雨浩、卢本韦、古月娜在第一处横台上摊开本子。一式核对，一次尝试；他们开始学着用这份压力辨认自己的力量。

\NovelDisplay{M01165R01}{%
(0)(1,1,1)(1)=\psi(\Omega_\omega)\times\omega%
}

末项 \allowbreak{}\NovelInline{I05136}{(1)}\allowbreak{} 把外部乘数推进到 \allowbreak{}\NovelInline{I05137}{\omega}\allowbreak{}，前面的 \allowbreak{}\NovelInline{I05138}{{\psi(\Omega_\omega)}}\allowbreak{} 保持原样。

\NovelDisplay{M01166R01}{%
(0)(1,1,1)(1)(2)=\psi(\Omega_\omega)\times\omega^\omega%
}

再接 \allowbreak{}\NovelInline{I05139}{(2)}\allowbreak{}，变化落在乘数的指数上：\allowbreak{}\NovelInline{I05140}{\omega\to\omega^\omega}\allowbreak{}。

\NovelDisplay{M01167R01}{%
(0)(1,1,1)(1)(2,1)=\psi(\Omega_{\omega})\times\psi(\Omega)%
}

\allowbreak{}\NovelInline{I05141}{(2,1)}\allowbreak{} 在此处生成乘数 \allowbreak{}\NovelInline{I05142}{\psi(\Omega)=\varepsilon_0}\allowbreak{}；它没有进入前面那个 \allowbreak{}\NovelInline{I05143}{\psi}\allowbreak{}。

\NovelDisplay{M01168R01}{%
(0)(1,1,1)(1)(2,1)(3,1)=\psi(\Omega_{\omega})\times\psi(\Omega^{2})%
}

连续的 \allowbreak{}\NovelInline{I05144}{(2,1)(3,1)}\allowbreak{} 把乘数内部参数由 \allowbreak{}\NovelInline{I05145}{\Omega}\allowbreak{} 推到 \allowbreak{}\NovelInline{I05146}{\Omega^2}\allowbreak{}。

``这一组只增加后面的乘数，前缀 \allowbreak{}\NovelInline{I05147}{(0)(1,1,1)}\allowbreak{} 对应的 \allowbreak{}\NovelInline{I05148}{{\psi(\Omega_\omega)}}\allowbreak{} 不变。''霍雨浩把乘号圈出来，继续往后核对。

\NovelDisplay{M01169R01}{%
(0)(1,1,1)(1)(2,1)(3,2)=\psi(\Omega_{\omega})\times\psi(\Omega_2)%
}

末列换成 \allowbreak{}\NovelInline{I05149}{(3,2)}\allowbreak{}，乘数的参数达到 \allowbreak{}\NovelInline{I05150}{\Omega_2}\allowbreak{}，外面仍是 \allowbreak{}\NovelInline{I05151}{{\psi(\Omega_\omega)}\times(\cdot)}\allowbreak{}。

\NovelDisplay{M01170R01}{%
(0)(1,1,1)(1)(2,1)(3,2)(4,3)=\psi(\Omega_{\omega})\times\psi(\Omega_3)%
}

再添 \allowbreak{}\NovelInline{I05152}{(4,3)}\allowbreak{}，参数进入 \allowbreak{}\NovelInline{I05153}{\Omega_3}\allowbreak{}。注意 \allowbreak{}\NovelInline{I05154}{\Omega_3}\allowbreak{} 是参数，不是函数下标。

\NovelDisplay{M01171R01}{%
(0)(1,1,1)(1)(2,1,1)= (0)(1,1,1)(1)(2,1)(3,2)(4,3)\cdots= \psi(\Omega_\omega)^2%
}

中间那条省略号表示基本列的极限；乘数逼近 \allowbreak{}\NovelInline{I05155}{\psi(\Omega_\omega)}\allowbreak{}，所以整式到 \allowbreak{}\NovelInline{I05156}{{\psi(\Omega_\omega)}^2}\allowbreak{}。

\NovelDisplay{M01172R01}{%
(0)(1,1,1)(1)(2,1,1)(1)(2,1,1)=\psi(\Omega_\omega)^3%
}

重复这里的 \allowbreak{}\NovelInline{I05157}{(1)(2,1,1)}\allowbreak{}，得到第三个 \allowbreak{}\NovelInline{I05158}{{\psi(\Omega_\omega)}}\allowbreak{} 因子，故为 \allowbreak{}\NovelInline{I05159}{{\psi(\Omega_\omega)}^3}\allowbreak{}。

\NovelDisplay{M01173R01}{%
(0)(1,1,1)(1)(2,1,1)(2)=\psi(\Omega_\omega)^\omega%
}

末尾 \allowbreak{}\NovelInline{I05160}{(2)}\allowbreak{} 把有限个 \allowbreak{}\NovelInline{I05161}{{\psi(\Omega_\omega)}}\allowbreak{} 因子的重复推进到 \allowbreak{}\NovelInline{I05162}{{\psi(\Omega_\omega)}^\omega}\allowbreak{}。

\NovelDisplay{M01174R01}{%
(0)(1,1,1)(1)(2,1,1)(2)(3,1,1)=\psi(\Omega_\omega)^{\psi(\Omega_\omega)}%
}

接上 \allowbreak{}\NovelInline{I05163}{(3,1,1)}\allowbreak{}，变化进入幂指数：\allowbreak{}\NovelInline{I05164}{\omega\to {\psi(\Omega_\omega)}}\allowbreak{}，底数仍为 \allowbreak{}\NovelInline{I05165}{{\psi(\Omega_\omega)}}\allowbreak{}。

幂继续向右嵌套。霍雨浩在最上层的指数外添了一对括号，古月娜依次指回每个括号所属的层次。

\NovelDisplay{M01175R01}{%
(0)(1,1,1)(1)(2,1,1)(2)(3,1,1)(3)(4,1,1)=\psi(\Omega_\omega)^{\psi(\Omega_\omega)^{\psi(\Omega_\omega)}}%
}

后段再复制一次这种结构，得到 \allowbreak{}\NovelInline{I05166}{{\psi(\Omega_\omega)}^{{\psi(\Omega_\omega)}^{\psi(\Omega_\omega)}}}\allowbreak{}；幂按右结合读成 \allowbreak{}\NovelInline{I05167}{{\psi(\Omega_\omega)}^{({\psi(\Omega_\omega)}^{\psi(\Omega_\omega)})}}\allowbreak{}。

第三日，二十米处，蔡须坤停住了。她继续增加外部乘数，胸口的牵扯却越来越重。霍雨浩圈出下一题的右括号，古月娜沿着它向内指。

``这次变的是参数。''

卢本韦陪她固定前段，只试这一处。压力使两种运转的差别显出来，他们再回到式子核对。

\NovelDisplay{M01176R01}{%
\psi(\Omega_\omega)\times\alpha,\qquad\psi(\Omega_\omega+\alpha)\quad\text{分别辨认乘法与参数中的加法。}%
}

\NovelDisplay{M01177R01}{%
(0)(1,1,1)(1,1)=\psi(\Omega_{\omega}+\Omega)%
}

此处 \allowbreak{}\NovelInline{I05168}{(1,1)}\allowbreak{} 已改变外层参数，增量为 \allowbreak{}\NovelInline{I05169}{\Omega}\allowbreak{}；不能把结果读成 \allowbreak{}\NovelInline{I05170}{{\psi(\Omega_\omega)}\times\Omega}\allowbreak{}。

这里的 \allowbreak{}\NovelInline{I05171}{(1,1)}\allowbreak{} 对应二行 BMS 的 \allowbreak{}\NovelInline{I05172}{(1,1)}\allowbreak{} 型结构。在本处，它让外层参数增加 \allowbreak{}\NovelInline{I05173}{\Omega}\allowbreak{}，其基本列仍逐层构造指数塔。不能沿用 MOCF 中``加一''的旧读法。

\NovelDisplay{M01178R01}{%
(0)(1,1,1)(1,1)(1,1)=\psi(\Omega_{\omega}+\Omega\times2)%
}

相同列再接一次，增加第二份 \allowbreak{}\NovelInline{I05174}{\Omega}\allowbreak{}，形成有序和 \allowbreak{}\NovelInline{I05175}{\Omega\times2}\allowbreak{}。

\NovelDisplay{M01179R01}{%
(0)(1,1,1)(1,1)(2)=\psi(\Omega_{\omega}+\Omega\times\omega)%
}

末尾的一行结构把这个有限重复推进到 \allowbreak{}\NovelInline{I05176}{\Omega\times\omega}\allowbreak{}。

\NovelDisplay{M01180R01}{%
(0)(1,1,1)(1,1)(2)(3)=\psi(\Omega_{\omega}+\Omega\times\omega^{\omega})%
}

再添一列一行结构，右乘部分由 \allowbreak{}\NovelInline{I05177}{\omega}\allowbreak{} 变为 \allowbreak{}\NovelInline{I05178}{\omega^\omega}\allowbreak{}。

\NovelDisplay{M01181R01}{%
(0)(1,1,1)(1,1)(2)(3,1)=\psi(\Omega_{\omega}+\Omega\times\psi(\Omega))%
}

固定的左因子 \allowbreak{}\NovelInline{I05179}{\Omega}\allowbreak{} 保留，后面的右因子改为 \allowbreak{}\NovelInline{I05180}{\psi(\Omega)}\allowbreak{}。

\NovelDisplay{M01182R01}{%
(0)(1,1,1)(1,1)(2)(3,1)(3,1)=\psi(\Omega_{\omega}+\Omega\times\psi(\Omega\times2))%
}

重复 \allowbreak{}\NovelInline{I05181}{(3,1)}\allowbreak{}，增加的是内层参数的一份 \allowbreak{}\NovelInline{I05182}{\Omega}\allowbreak{}，得到 \allowbreak{}\NovelInline{I05183}{\psi(\Omega\times2)}\allowbreak{}。

\NovelDisplay{M01183R01}{%
(0)(1,1,1)(1,1)(2)(3,1)(4,1)=\psi(\Omega_{\omega}+\Omega\times\psi(\Omega^{2}))%
}

第一坐标继续抬高，内层参数变成 \allowbreak{}\NovelInline{I05184}{\Omega^2}\allowbreak{}；这与上一式的 \allowbreak{}\NovelInline{I05185}{\Omega\times2}\allowbreak{} 不同。

\NovelDisplay{M01184R01}{%
(0)(1,1,1)(1,1)(2)(3,1)(4,2)=\psi(\Omega_{\omega}+\Omega\times\psi(\Omega_2))%
}

第二坐标升到二，内层参数进入 \allowbreak{}\NovelInline{I05186}{\Omega_2}\allowbreak{}，外层的 \allowbreak{}\NovelInline{I05187}{\Omega_\omega}\allowbreak{} 仍保留。

\NovelDisplay{M01185R01}{%
(0)(1,1,1)(1,1)(2)(3,1,1)=\psi(\Omega_{\omega}+\Omega\times\psi(\Omega_{\omega}))%
}

末列出现第三行的一，右因子达到 \allowbreak{}\NovelInline{I05188}{\psi(\Omega_\omega)}\allowbreak{}。

\allowbreak{}\NovelInline{I05189}{(2)}\allowbreak{} 在这一段隔出内层结构。其后的变化首先落在内层的 \allowbreak{}\NovelInline{I05190}{\psi}\allowbreak{}，外面那份 \allowbreak{}\NovelInline{I05191}{\Omega_\omega}\allowbreak{} 及相应左因子留在原处。

\NovelDisplay{M01186R01}{%
(0)(1,1,1)(1,1)(2)(3,1,1)(3,1)=\psi(\Omega_{\omega}+\Omega\times\psi(\Omega_{\omega}+\Omega))%
}

新接的 \allowbreak{}\NovelInline{I05192}{(3,1)}\allowbreak{} 在内层起作用，令内层参数 \allowbreak{}\NovelInline{I05193}{\Omega_\omega\to\Omega_\omega+\Omega}\allowbreak{}。

\NovelDisplay{M01187R01}{%
(0)(1,1,1)(1,1)(2)(3,1,1)(3,1)(4)(5,1,1)%
}
\NovelDisplay{M01187R02}{%
\quad =\psi(\Omega_{\omega}+\Omega\times\psi(\Omega_{\omega}+\Omega\times\psi(\Omega_{\omega})))%
}

\allowbreak{}\NovelInline{I05194}{(4)}\allowbreak{} 再隔出一层；最内层的 \allowbreak{}\NovelInline{I05195}{\psi(\Omega_\omega)}\allowbreak{} 先乘 \allowbreak{}\NovelInline{I05196}{\Omega}\allowbreak{}，再参与外层相加。

\NovelDisplay{M01188R01}{%
(0)(1,1,1)(1,1)(2,1)=\psi(\Omega_{\omega}+\Omega^{2})%
}

\allowbreak{}\NovelInline{I05197}{(1,1)(2,1)}\allowbreak{} 相接，在这一位置形成 \allowbreak{}\NovelInline{I05198}{\Omega^2}\allowbreak{}，作用范围要看整条 BMS。

\NovelDisplay{M01189R01}{%
(0)(1,1,1)(1,1)(2,1)(2)=\psi(\Omega_{\omega}+\Omega^{2}\times\omega)%
}

再接 \allowbreak{}\NovelInline{I05199}{(2)}\allowbreak{}，已有的 \allowbreak{}\NovelInline{I05200}{\Omega^2}\allowbreak{} 右乘 \allowbreak{}\NovelInline{I05201}{\omega}\allowbreak{}，幂次数暂不变化。

\NovelDisplay{M01190R01}{%
(0)(1,1,1)(1,1)(2,1)(2,1)=\psi(\Omega_{\omega}+\Omega^{3})%
}

换成 \allowbreak{}\NovelInline{I05202}{(2,1)}\allowbreak{}，幂次数增一，得到 \allowbreak{}\NovelInline{I05203}{\Omega^3}\allowbreak{}，不能误写成 \allowbreak{}\NovelInline{I05204}{\Omega^2\times2}\allowbreak{}。

\allowbreak{}\NovelInline{I05205}{(n,1,0)}\allowbreak{} 在当前结构中引入 \allowbreak{}\NovelInline{I05206}{\Omega}\allowbreak{} 的幂次；例如这里相邻的 \allowbreak{}\NovelInline{I05207}{(1,1)(2,1)}\allowbreak{} 给出 \allowbreak{}\NovelInline{I05208}{\Omega^2}\allowbreak{}。它的具体贡献依赖整条 BMS，不能脱离位置逐列换算。

\NovelDisplay{M01191R01}{%
(0)(1,1,1)(1,1)(2,1)(3)=\psi(\Omega_\omega+\Omega^\omega)%
}

后接 \allowbreak{}\NovelInline{I05209}{(3)}\allowbreak{}，把有限幂的推进收束为 \allowbreak{}\NovelInline{I05210}{\Omega^\omega}\allowbreak{}。

\NovelDisplay{M01192R01}{%
(0)(1,1,1)(1,1)(2,1)(3)(4,1,1)=\psi(\Omega_\omega+\Omega^{\psi(\Omega_\omega)})%
}

\allowbreak{}\NovelInline{I05211}{(4,1,1)}\allowbreak{} 让幂指数由 \allowbreak{}\NovelInline{I05212}{\omega}\allowbreak{} 进到 \allowbreak{}\NovelInline{I05213}{{\psi(\Omega_\omega)}}\allowbreak{}，底数 \allowbreak{}\NovelInline{I05214}{\Omega}\allowbreak{} 不变。

\NovelDisplay{M01193R01}{%
(0)(1,1,1)(1,1)(2,1)(3,1)=\psi(\Omega_\omega+\Omega^\Omega)%
}

末列改成 \allowbreak{}\NovelInline{I05215}{(3,1)}\allowbreak{}，在这个结构里生成 \allowbreak{}\NovelInline{I05216}{\Omega^\Omega}\allowbreak{}。

第八日，四十米处，古月娜把内层的一道银纹接错了。气息往胸口涌，她退回横台，等那口气平下去。

蔡须坤挪到她身边，指出自己昨日也受阻的位置。卢本韦照两种顺序各试一次，霍雨浩在纸上圈出内外两对括号。

``每次只认正在变的那一层。''

古月娜接过炭条，自己改了最后一笔。

\NovelDisplay{M01194R01}{%
(0)(1,1,1)(1,1)(2,2)=\psi(\Omega_{\omega}+\psi_{\Omega_2}(\Omega_2))%
}

这里新增上界为 \allowbreak{}\NovelInline{I05217}{\Omega_2}\allowbreak{} 的坍缩层，整体为 \allowbreak{}\NovelInline{I05218}{\psi(\Omega_\omega+{\psi_{\Omega_2}(\Omega_2)})}\allowbreak{}。

\NovelDisplay{M01195R01}{%
(0)(1,1,1)(1,1)(2,2)(2)=\psi(\Omega_{\omega}+\psi_{\Omega_2}(\Omega_2)\times\omega)%
}

固定 \allowbreak{}\NovelInline{I05219}{{\psi_{\Omega_2}(\Omega_2)}}\allowbreak{}，它在外层参数中右乘 \allowbreak{}\NovelInline{I05220}{\omega}\allowbreak{}；\allowbreak{}\NovelInline{I05221}{\psi_{\Omega_2}}\allowbreak{} 的参数仍为 \allowbreak{}\NovelInline{I05222}{\Omega_2}\allowbreak{}。

\NovelDisplay{M01196R01}{%
(0)(1,1,1)(1,1)(2,2)(2,1)=\psi(\Omega_{\omega}+\psi_{\Omega_2}(\Omega_2)\times\Omega)%
}

右乘部分由 \allowbreak{}\NovelInline{I05223}{\omega}\allowbreak{} 进到 \allowbreak{}\NovelInline{I05224}{\Omega}\allowbreak{}，仍未改变 \allowbreak{}\NovelInline{I05225}{{\psi_{\Omega_2}(\Omega_2)}}\allowbreak{} 的内部。

\NovelDisplay{M01197R01}{%
(0)(1,1,1)(1,1)(2,2)(2,1)(3,2)=\psi(\Omega_{\omega}+\psi_{\Omega_2}(\Omega_2)^{2})%
}

加上 \allowbreak{}\NovelInline{I05226}{(3,2)}\allowbreak{} 后出现第二份 \allowbreak{}\NovelInline{I05227}{{\psi_{\Omega_2}(\Omega_2)}}\allowbreak{} 因子，得到 \allowbreak{}\NovelInline{I05228}{{\psi_{\Omega_2}(\Omega_2)}^2}\allowbreak{}。

\NovelDisplay{M01198R01}{%
(0)(1,1,1)(1,1)(2,2)(2,1)(3,2)(3)=\psi(\Omega_{\omega}+\psi_{\Omega_2}(\Omega_2)^{\omega})%
}

末尾的一行结构把 \allowbreak{}\NovelInline{I05229}{{\psi_{\Omega_2}(\Omega_2)}}\allowbreak{} 的有限幂推进到 \allowbreak{}\NovelInline{I05230}{{\psi_{\Omega_2}(\Omega_2)}^\omega}\allowbreak{}。

\NovelDisplay{M01199R01}{%
(0)(1,1,1)(1,1)(2,2)(2,1)(3,2)(3,1)=\psi(\Omega_{\omega}+\psi_{\Omega_2}(\Omega_2)^{\Omega})%
}

末列第二坐标为一，指数层变成 \allowbreak{}\NovelInline{I05231}{\Omega}\allowbreak{}，得到 \allowbreak{}\NovelInline{I05232}{{\psi_{\Omega_2}(\Omega_2)}^\Omega}\allowbreak{}。

\NovelDisplay{M01200R01}{%
(0)(1,1,1)(1,1)(2,2)(2,1)(3,2)(3,1)(3,1)=\psi(\Omega_{\omega}+\psi_{\Omega_2}(\Omega_2)^{\Omega^{2}})%
}

重复同一列使幂指数中的 \allowbreak{}\NovelInline{I05233}{\Omega}\allowbreak{} 形成平方，结果为 \allowbreak{}\NovelInline{I05234}{{\psi_{\Omega_2}(\Omega_2)}^{\Omega^2}}\allowbreak{}。

\NovelDisplay{M01201R01}{%
(0)(1,1,1)(1,1)(2,2)(2,1)(3,2)(3,1)(4,1)=\psi(\Omega_{\omega}+\psi_{\Omega_2}(\Omega_2)^{\Omega^{\Omega}})%
}

提高第一坐标，幂指数由 \allowbreak{}\NovelInline{I05235}{\Omega^2}\allowbreak{} 进到 \allowbreak{}\NovelInline{I05236}{\Omega^\Omega}\allowbreak{}。

\NovelDisplay{M01202R01}{%
(0)(1,1,1)(1,1)(2,2)(2,1)(3,2)(3,1)(4,2)=\psi(\Omega_{\omega}+\psi_{\Omega_2}(\Omega_2)^{\psi_{\Omega_2}(\Omega_2)})%
}

末列第二坐标提高到二，幂指数本身到 \allowbreak{}\NovelInline{I05237}{{\psi_{\Omega_2}(\Omega_2)}}\allowbreak{}，结果是 \allowbreak{}\NovelInline{I05238}{{\psi_{\Omega_2}(\Omega_2)}^{\psi_{\Omega_2}(\Omega_2)}}\allowbreak{}。

\NovelDisplay{M01203R01}{%
(0)(1,1,1)(1,1)(2,2)(2,2)=\psi(\Omega_{\omega}+\psi_{\Omega_2}(\Omega_2\times2))%
}

重复 \allowbreak{}\NovelInline{I05239}{(2,2)}\allowbreak{} 改动的是 \allowbreak{}\NovelInline{I05240}{{\psi_{\Omega_2}(\Omega_2)}}\allowbreak{} 的内部参数：\allowbreak{}\NovelInline{I05241}{\Omega_2\to\Omega_2\times2}\allowbreak{}，并非外面乘二。

霍雨浩另写一条去掉外层背景的对照式，与卢本韦确认相同的参数推进仍在发生：

\NovelDisplay{M01204R01}{%
(0)(1,1)(2,2)(2,2)=\psi(\Omega_2\times2)%
}

回到上式，\allowbreak{}\NovelInline{I05242}{\psi_{\Omega_2}}\allowbreak{} 的参数由 \allowbreak{}\NovelInline{I05243}{\Omega_2}\allowbreak{} 进到 \allowbreak{}\NovelInline{I05244}{\Omega_2\times2}\allowbreak{}；外层背景没有把这一步变成普通乘二。

\NovelDisplay{M01205R01}{%
(0)(1,1,1)(1,1)(2,2)(3)=\psi(\Omega_{\omega}+\psi_{\Omega_2}(\Omega_2\times\omega))%
}

\allowbreak{}\NovelInline{I05245}{(3)}\allowbreak{} 使内层参数的右系数达到 \allowbreak{}\NovelInline{I05246}{\omega}\allowbreak{}，得到 \allowbreak{}\NovelInline{I05247}{\Omega_2\times\omega}\allowbreak{}。

\NovelDisplay{M01206R01}{%
(0)(1,1,1)(1,1)(2,2)(3,1)=\psi(\Omega_{\omega}+\psi_{\Omega_2}(\Omega_2\times\Omega))%
}

把末列改成 \allowbreak{}\NovelInline{I05248}{(3,1)}\allowbreak{}，同一处右系数达到 \allowbreak{}\NovelInline{I05249}{\Omega}\allowbreak{}。

\NovelDisplay{M01207R01}{%
(0)(1,1,1)(1,1)(2,2)(3,1)(3)=\psi(\Omega_{\omega}+\psi_{\Omega_2}(\Omega_2\times\Omega\times\omega))%
}

再接 \allowbreak{}\NovelInline{I05250}{(3)}\allowbreak{}，参数中形成 \allowbreak{}\NovelInline{I05251}{\Omega_2\times\Omega\times\omega}\allowbreak{}；乘法顺序不能对调。

\NovelDisplay{M01208R01}{%
(0)(1,1,1)(1,1)(2,2)(3,1)(3,1)=\psi(\Omega_{\omega}+\psi_{\Omega_2}(\Omega_2\times\Omega^{2}))%
}

重复 \allowbreak{}\NovelInline{I05252}{(3,1)}\allowbreak{}，在 \allowbreak{}\NovelInline{I05253}{\Omega_2}\allowbreak{} 右边形成 \allowbreak{}\NovelInline{I05254}{\Omega^2}\allowbreak{}。

\NovelDisplay{M01209R01}{%
(0)(1,1,1)(1,1)(2,2)(3,1)(4)=\psi(\Omega_{\omega}+\psi_{\Omega_2}(\Omega_2\times\Omega^{\omega}))%
}

提高第一坐标，内层参数右端成为 \allowbreak{}\NovelInline{I05255}{\Omega^\omega}\allowbreak{}。

\NovelDisplay{M01210R01}{%
(0)(1,1,1)(1,1)(2,2)(3,1)(4,1)=\psi(\Omega_{\omega}+\psi_{\Omega_2}(\Omega_2\times\Omega^{\Omega}))%
}

提高第二坐标到一，右端进一步成为 \allowbreak{}\NovelInline{I05256}{\Omega^\Omega}\allowbreak{}。

\NovelDisplay{M01211R01}{%
(0)(1,1,1)(1,1)(2,2)(3,1)(4,2)=\psi(\Omega_{\omega}+\psi_{\Omega_2}(\Omega_2\times\psi_{\Omega_2}(\Omega_2)))%
}

\allowbreak{}\NovelInline{I05257}{(3,1)(4,2)}\allowbreak{} 又引入一个 \allowbreak{}\NovelInline{I05258}{{\psi_{\Omega_2}(\Omega_2)}}\allowbreak{}，得到内层参数 \allowbreak{}\NovelInline{I05259}{\Omega_2\times {\psi_{\Omega_2}(\Omega_2)}}\allowbreak{}。

\NovelDisplay{M01212R01}{%
(0)(1,1,1)(1,1)(2,2)(3,1)(4,2)(2,2)%
}
\NovelDisplay{M01212R02}{%
\quad =\psi(\Omega_{\omega}+\psi_{\Omega_2}(\Omega_2\times\psi_{\Omega_2}(\Omega_2)+\Omega_2))%
}

再添 \allowbreak{}\NovelInline{I05260}{(2,2)}\allowbreak{}，在这整个参数后面加 \allowbreak{}\NovelInline{I05261}{\Omega_2}\allowbreak{}；左边的 \allowbreak{}\NovelInline{I05262}{\Omega_2\times {\psi_{\Omega_2}(\Omega_2)}}\allowbreak{} 保留。

\NovelDisplay{M01213R01}{%
(0)(1,1,1)(1,1)(2,2)(3,1)(4,2)(4,2)%
}
\NovelDisplay{M01213R02}{%
\quad =\psi(\Omega_{\omega}+\psi_{\Omega_2}(\Omega_2\times\psi_{\Omega_2}(\Omega_2\times2)))%
}

重复 \allowbreak{}\NovelInline{I05263}{(4,2)}\allowbreak{}，改动更内层的参数；外层的 \allowbreak{}\NovelInline{I05264}{\Omega_2\times(\cdot)}\allowbreak{} 不变。

像 \allowbreak{}\NovelInline{I05265}{(1,1)(2,2)}\allowbreak{}、\allowbreak{}\NovelInline{I05266}{(3,1)(4,2)}\allowbreak{} 这样具有对应间隔的结构，在本段会引入 \allowbreak{}\NovelInline{I05267}{{\psi_{\Omega_2}(\Omega_2)}}\allowbreak{}。续接相应的 \allowbreak{}\NovelInline{I05268}{(n,2)}\allowbreak{}，可在这层函数参数中加 \allowbreak{}\NovelInline{I05269}{\Omega_2}\allowbreak{}。更高层也有类似现象，须检查父项和所在层次，不能无条件套用。

\NovelDisplay{M01214R01}{%
(0)(1,1,1)(1,1)(2,2)(3,1)(4,2)(5)=\psi(\Omega_{\omega}+\psi_{\Omega_2}(\Omega_2\times\psi_{\Omega_2}(\Omega_2\times\omega)))%
}

末列 \allowbreak{}\NovelInline{I05270}{(5)}\allowbreak{} 将最内层的 \allowbreak{}\NovelInline{I05271}{\Omega_2\times2}\allowbreak{} 推进为 \allowbreak{}\NovelInline{I05272}{\Omega_2\times\omega}\allowbreak{}。

\NovelDisplay{M01215R01}{%
(0)(1,1,1)(1,1)(2,2)(3,1)(4,2)(5,1)%
}
\NovelDisplay{M01215R02}{%
\quad =\psi(\Omega_{\omega}+\psi_{\Omega_2}(\Omega_2\times\psi_{\Omega_2}(\Omega_2\times\Omega)))%
}

末列第二坐标变一，最内层的右系数到 \allowbreak{}\NovelInline{I05273}{\Omega}\allowbreak{}，其余各层保持原位置。

第十一日，卢本韦把两张纸并在一起。霍雨浩刚试完两种平方，肩背仍绷着；压力让其中一处先在内层受阻，另一处却牵动整个函数值。她按住纸角，请他先停在这处差别上。

\NovelDisplay{M01216R01}{%
\bigl(\psi_{\Omega_2}(\Omega_2)\bigr)^2,\qquad\psi_{\Omega_2}(\Omega_2^2)%
}

``一个平方在括号外，一个在里面。位置不同，要按各自那条 BMS 核对。''

霍雨浩重新核对第二坐标。蔡须坤与古月娜轮流再试，将感到受阻的先后和纸上的括号一一对照。

\NovelDisplay{M01217R01}{%
(0)(1,1,1)(1,1)(2,2)(3,2)=\psi(\Omega_{\omega}+\psi_{\Omega_2}(\Omega_2^{2}))%
}

这里改变的是函数内部的参数幂次，不能写成 \allowbreak{}\NovelInline{I05274}{{\psi_{\Omega_2}(\Omega_2)}^2}\allowbreak{}。

相隔 \allowbreak{}\NovelInline{I05275}{(1)}\allowbreak{} 的连续两项 \allowbreak{}\NovelInline{I05276}{(n,2)}\allowbreak{}，在此处引入参数的幂次结构：\allowbreak{}\NovelInline{I05277}{\Omega_2\to\Omega_2^2}\allowbreak{}。把这一步看成函数值平方，就会认错下面整组。

\NovelDisplay{M01218R01}{%
(0)(1,1,1)(1,1)(2,2)(3,2)(2)=\psi(\Omega_{\omega}+\psi_{\Omega_2}(\Omega_2^{2})\times\omega)%
}

现在取 \allowbreak{}\NovelInline{I05278}{{\psi_{\Omega_2}(\Omega_2^2)}}\allowbreak{} 右乘 \allowbreak{}\NovelInline{I05279}{\omega}\allowbreak{}。看清乘号位于 \allowbreak{}\NovelInline{I05280}{\psi_{\Omega_2}}\allowbreak{} 的右括号外。

\NovelDisplay{M01219R01}{%
(0)(1,1,1)(1,1)(2,2)(3,2)(2,1)=\psi(\Omega_{\omega}+\psi_{\Omega_2}(\Omega_2^{2})\times\Omega)%
}

右乘部分提高到 \allowbreak{}\NovelInline{I05281}{\Omega}\allowbreak{}，函数内部仍为 \allowbreak{}\NovelInline{I05282}{\Omega_2^2}\allowbreak{}。

\NovelDisplay{M01220R01}{%
(0)(1,1,1)(1,1)(2,2)(3,2)(2,1)(3,2)=\psi(\Omega_{\omega}+\psi_{\Omega_2}(\Omega_2^{2})\times\psi_{\Omega_2}(\Omega_2))%
}

末列引入右因子 \allowbreak{}\NovelInline{I05283}{{\psi_{\Omega_2}(\Omega_2)}}\allowbreak{}，故外层参数出现 \allowbreak{}\NovelInline{I05284}{{\psi_{\Omega_2}(\Omega_2^2)}\times {\psi_{\Omega_2}(\Omega_2)}}\allowbreak{}。

\NovelDisplay{M01221R01}{%
(0)(1,1,1)(1,1)(2,2)(3,2)(2,1)(3,2)(4,2)=\psi(\Omega_{\omega}+\psi_{\Omega_2}(\Omega_2^{2})^{2})%
}

再接一列后，右因子也达到 \allowbreak{}\NovelInline{I05285}{{\psi_{\Omega_2}(\Omega_2^2)}}\allowbreak{}，整段变为 \allowbreak{}\NovelInline{I05286}{{\psi_{\Omega_2}(\Omega_2^2)}^2}\allowbreak{}。

\NovelDisplay{M01222R01}{%
(0)(1,1,1)(1,1)(2,2)(3,2)(2,1)(3,2)(4,2)(3,1)%
}
\NovelDisplay{M01222R02}{%
\quad (4,2)(5,2)=\psi(\Omega_{\omega}+\psi_{\Omega_2}(\Omega_2^{2})^{\psi_{\Omega_2}(\Omega_2^{2})})%
}

后段继续嵌套，幂指数本身达到 \allowbreak{}\NovelInline{I05287}{{\psi_{\Omega_2}(\Omega_2^2)}}\allowbreak{}，整段为 \allowbreak{}\NovelInline{I05288}{{\psi_{\Omega_2}(\Omega_2^2)}^{\psi_{\Omega_2}(\Omega_2^2)}}\allowbreak{}。

\NovelDisplay{M01223R01}{%
(0)(1,1,1)(1,1)(2,2)(3,2)(2,2)=\psi(\Omega_{\omega}+\psi_{\Omega_2}(\Omega_2^{2}+\Omega_2))%
}

退回参数层看，新增 \allowbreak{}\NovelInline{I05289}{(2,2)}\allowbreak{} 产生的是 \allowbreak{}\NovelInline{I05290}{\Omega_2^2+\Omega_2}\allowbreak{}。

\allowbreak{}\NovelInline{I05291}{(2,2)}\allowbreak{} 在这组仍使相应参数增加 \allowbreak{}\NovelInline{I05292}{\Omega_2}\allowbreak{}。后面的每式都要检查增量落在哪一层，不能只看末列相同。

\NovelDisplay{M01224R01}{%
(0)(1,1,1)(1,1)(2,2)(3,2)(2,2)(3,1)=\psi(\Omega_{\omega}+\psi_{\Omega_2}(\Omega_2^{2}+\Omega_2\times\Omega))%
}

在新增的 \allowbreak{}\NovelInline{I05293}{\Omega_2}\allowbreak{} 后右乘 \allowbreak{}\NovelInline{I05294}{\Omega}\allowbreak{}，前面的 \allowbreak{}\NovelInline{I05295}{\Omega_2^2}\allowbreak{} 固定。

\NovelDisplay{M01225R01}{%
(0)(1,1,1)(1,1)(2,2)(3,2)(2,2)(3,1)(4)=\psi(\Omega_{\omega}+\psi_{\Omega_2}(\Omega_2^{2}+\Omega_2\times\Omega^{\omega}))%
}

\allowbreak{}\NovelInline{I05296}{(4)}\allowbreak{} 把后项提升为 \allowbreak{}\NovelInline{I05297}{\Omega_2\times\Omega^\omega}\allowbreak{}，不动前项。

\NovelDisplay{M01226R01}{%
(0)(1,1,1)(1,1)(2,2)(3,2)(2,2)(3,1)(4,1)%
}
\NovelDisplay{M01226R02}{%
\quad =\psi(\Omega_{\omega}+\psi_{\Omega_2}(\Omega_2^{2}+\Omega_2\times\Omega^{\Omega}))%
}

\allowbreak{}\NovelInline{I05298}{(4,1)}\allowbreak{} 则把同一幂指数推进到 \allowbreak{}\NovelInline{I05299}{\Omega}\allowbreak{}。

\NovelDisplay{M01227R01}{%
(0)(1,1,1)(1,1)(2,2)(3,2)(2,2)(3,1)(4,2)%
}
\NovelDisplay{M01227R02}{%
\quad =\psi(\Omega_{\omega}+\psi_{\Omega_2}(\Omega_2^{2}+\Omega_2\times\psi_{\Omega_2}(\Omega_2)))%
}

\allowbreak{}\NovelInline{I05300}{(4,2)}\allowbreak{} 在后项内引入 \allowbreak{}\NovelInline{I05301}{{\psi_{\Omega_2}(\Omega_2)}}\allowbreak{}，形成 \allowbreak{}\NovelInline{I05302}{\Omega_2^2+\Omega_2\times {\psi_{\Omega_2}(\Omega_2)}}\allowbreak{}。

\NovelDisplay{M01228R01}{%
(0)(1,1,1)(1,1)(2,2)(3,2)(2,2)(3,1)(4,2)(5,1)(6,2)=\psi(\Omega_{\omega}+\psi_{\Omega_2}(\Omega_2^{2}%
}
\NovelDisplay{M01228R02}{%
\quad +\Omega_2\times\psi_{\Omega_2}(\Omega_2\times\psi_{\Omega_2}(\Omega_2))))%
}

再加 \allowbreak{}\NovelInline{I05303}{(5,1)(6,2)}\allowbreak{}，\allowbreak{}\NovelInline{I05304}{{\psi_{\Omega_2}(\Omega_2)}}\allowbreak{} 的参数也含有 \allowbreak{}\NovelInline{I05305}{\Omega_2\times {\psi_{\Omega_2}(\Omega_2)}}\allowbreak{}，出现新的嵌套。

\NovelDisplay{M01229R01}{%
(0)(1,1,1)(1,1)(2,2)(3,2)(2,2)(3,1)(4,2)(5,2)%
}
\NovelDisplay{M01229R02}{%
\quad =\psi(\Omega_{\omega}+\psi_{\Omega_2}(\Omega_2^{2}+\Omega_2\times\psi_{\Omega_2}(\Omega_2^{2})))%
}

末项换成 \allowbreak{}\NovelInline{I05306}{(5,2)}\allowbreak{}，更内层参数成为 \allowbreak{}\NovelInline{I05307}{\Omega_2^2}\allowbreak{}，也就是引入 \allowbreak{}\NovelInline{I05308}{{\psi_{\Omega_2}(\Omega_2^2)}}\allowbreak{}。

\NovelDisplay{M01230R01}{%
(0)(1,1,1)(1,1)(2,2)(3,2)(2,2)(3,2)=\psi(\Omega_{\omega}+\psi_{\Omega_2}(\Omega_2^{2}\times2))%
}

两份 \allowbreak{}\NovelInline{I05309}{\Omega_2^2}\allowbreak{} 在参数里相加，得到 \allowbreak{}\NovelInline{I05310}{\Omega_2^2\times2}\allowbreak{}。

\NovelDisplay{M01231R01}{%
(0)(1,1,1)(1,1)(2,2)(3,2)(3)=\psi(\Omega_{\omega}+\psi_{\Omega_2}(\Omega_2^{2}\times\omega))%
}

\allowbreak{}\NovelInline{I05311}{(3)}\allowbreak{} 把这类有限重复推进到 \allowbreak{}\NovelInline{I05312}{\Omega_2^2\times\omega}\allowbreak{}。

\NovelDisplay{M01232R01}{%
(0)(1,1,1)(1,1)(2,2)(3,2)(3,1)=\psi(\Omega_{\omega}+\psi_{\Omega_2}(\Omega_2^{2}\times\Omega))%
}

末列第二坐标变一，相应右系数到 \allowbreak{}\NovelInline{I05313}{\Omega}\allowbreak{}。

\NovelDisplay{M01233R01}{%
(0)(1,1,1)(1,1)(2,2)(3,2)(3,1)(4,2)=\psi(\Omega_{\omega}+\psi_{\Omega_2}(\Omega_2^{2}\times\psi_{\Omega_2}(\Omega_2)))%
}

再接 \allowbreak{}\NovelInline{I05314}{(4,2)}\allowbreak{}，右因子到 \allowbreak{}\NovelInline{I05315}{{\psi_{\Omega_2}(\Omega_2)}}\allowbreak{}，其左边的 \allowbreak{}\NovelInline{I05316}{\Omega_2^2}\allowbreak{} 不动。

\NovelDisplay{M01234R01}{%
(0)(1,1,1)(1,1)(2,2)(3,2)(3,2)=\psi(\Omega_{\omega}+\psi_{\Omega_2}(\Omega_2^{3}))%
}

相邻两列的第二坐标都为二，参数的幂次由二升到三，得到 \allowbreak{}\NovelInline{I05317}{\Omega_2^3}\allowbreak{}。

\NovelDisplay{M01235R01}{%
(0)(1,1,1)(1,1)(2,2)(3,2)(4,1)=\psi(\Omega_\omega+\psi_{\Omega_2}(\Omega_2^{\Omega}))%
}

提高第一坐标，参数进入 \allowbreak{}\NovelInline{I05318}{\Omega_2^\Omega}\allowbreak{}。

\NovelDisplay{M01236R01}{%
(0)(1,1,1)(1,1)(2,2)(3,2)(4,2)=\psi(\Omega_\omega+\psi_{\Omega_2}(\Omega_2^{\Omega_2}))%
}

末列第二坐标变二，幂指数达到 \allowbreak{}\NovelInline{I05319}{\Omega_2}\allowbreak{}，结果为 \allowbreak{}\NovelInline{I05320}{\Omega_2^{\Omega_2}}\allowbreak{}。

第十五日，霍雨浩在台阶上蹲得两腿发颤，仍没把内外层分开。蔡须坤替他压住纸，卢本韦提醒他先保留共同前项。

他喘匀气，将下一组共同不变的部分描了一遍。

\NovelDisplay{M01237R01}{%
\psi\bigl(\Omega_\omega+\psi_{\Omega_2}(\boxed{\alpha})\bigr)%
}

``先盯方框。\allowbreak{}\NovelInline{I05321}{\alpha}\allowbreak{} 换了，不等于 \allowbreak{}\NovelInline{I05322}{\psi_{\Omega_2}}\allowbreak{} 也换了。''

\NovelDisplay{M01238R01}{%
(0)(1,1,1)(1,1)(2,2)(3,3)=\psi(\Omega_{\omega}+\psi_{\Omega_2}(\Omega_3))%
}

第二坐标进一步变三，参数跳到新的层次标记 \allowbreak{}\NovelInline{I05323}{\Omega_3}\allowbreak{}；函数下标仍为 \allowbreak{}\NovelInline{I05324}{\Omega_2}\allowbreak{}。

\NovelDisplay{M01239R01}{%
(0)(1,1,1)(1,1)(2,2)(3,3)(3,3)=\psi(\Omega_{\omega}+\psi_{\Omega_2}(\Omega_3\times2))%
}

重复 \allowbreak{}\NovelInline{I05325}{(3,3)}\allowbreak{}，在同一函数的参数中加第二份 \allowbreak{}\NovelInline{I05326}{\Omega_3}\allowbreak{}。

\NovelDisplay{M01240R01}{%
(0)(1,1,1)(1,1)(2,2)(3,3)(4)=\psi(\Omega_{\omega}+\psi_{\Omega_2}(\Omega_3\times\omega))%
}

末列 \allowbreak{}\NovelInline{I05327}{(4)}\allowbreak{} 给 \allowbreak{}\NovelInline{I05328}{\Omega_3}\allowbreak{} 添上右因子 \allowbreak{}\NovelInline{I05329}{\omega}\allowbreak{}。

\NovelDisplay{M01241R01}{%
(0)(1,1,1)(1,1)(2,2)(3,3)(4,1)=\psi(\Omega_{\omega}+\psi_{\Omega_2}(\Omega_3\times\Omega))%
}

把末列改为 \allowbreak{}\NovelInline{I05330}{(4,1)}\allowbreak{}，右因子推进到 \allowbreak{}\NovelInline{I05331}{\Omega}\allowbreak{}。

\NovelDisplay{M01242R01}{%
(0)(1,1,1)(1,1)(2,2)(3,3)(4,1)(5,2)=\psi(\Omega_{\omega}+\psi_{\Omega_2}(\Omega_3\times\psi_{\Omega_2}(\Omega_2)))%
}

再接 \allowbreak{}\NovelInline{I05332}{(5,2)}\allowbreak{}，右因子变为 \allowbreak{}\NovelInline{I05333}{{\psi_{\Omega_2}(\Omega_2)}}\allowbreak{}；它仍处于外面那层坍缩函数的参数中。

\NovelDisplay{M01243R01}{%
(0)(1,1,1)(1,1)(2,2)(3,3)(4,2)=\psi(\Omega_{\omega}+\psi_{\Omega_2}(\Omega_3\times\Omega_2))%
}

末列第二坐标到二，右因子达到 \allowbreak{}\NovelInline{I05334}{\Omega_2}\allowbreak{}。

\NovelDisplay{M01244R01}{%
(0)(1,1,1)(1,1)(2,2)(3,3)(4,2)(5,3)=\psi(\Omega_{\omega}+\psi_{\Omega_2}(\Omega_3\times\psi_{\Omega_3}(\Omega_3)))%
}

再接 \allowbreak{}\NovelInline{I05335}{(5,3)}\allowbreak{}，生成 \allowbreak{}\NovelInline{I05336}{\psi_{\Omega_3}(\Omega_3)}\allowbreak{}；这里才真的出现新的函数下标。

\NovelDisplay{M01245R01}{%
(0)(1,1,1)(1,1)(2,2)(3,3)(4,3)=\psi(\Omega_{\omega}+\psi_{\Omega_2}(\Omega_3^{2}))%
}

相邻列的第二坐标同为三，在参数中形成 \allowbreak{}\NovelInline{I05337}{\Omega_3^2}\allowbreak{}。

\NovelDisplay{M01246R01}{%
(0)(1,1,1)(1,1)(2,2)(3,3)(4,4)=\psi(\Omega_{\omega}+\psi_{\Omega_2}(\Omega_4))%
}

末列第二坐标进到四，参数到 \allowbreak{}\NovelInline{I05338}{\Omega_4}\allowbreak{}；外层函数仍没有换成 \allowbreak{}\NovelInline{I05339}{\psi_{\Omega_4}}\allowbreak{}。

第十九日，他们仍卡在六十米处。蔡须坤抬脚，又放下；熟练调动的力量到了这处接缝，在压力下分成两股互相牵扯。

古月娜指出她把内函数当成了第二份外层基数。霍雨浩写出两种结构，卢本韦扶她站稳，等她重新核对。

\NovelDisplay{M01247R01}{%
\Omega_\omega+{\psi_{\Omega_2}(\Omega_\omega)}\ \text{与}\ \Omega_\omega+\Omega_\omega\ \text{分开辨认}%
}

``先把这一道认清。我们都在这里。''

\NovelDisplay{M01248R01}{%
(0)(1,1,1)(1,1)(2,2,1)=\psi(\Omega_\omega+\psi_{\Omega_2}(\Omega_\omega))%
}

这是内层参数走到 \allowbreak{}\NovelInline{I05340}{\Omega_\omega}\allowbreak{}，整体尚非 \allowbreak{}\NovelInline{I05341}{\psi(\Omega_\omega\times2)}\allowbreak{}。

这是一个卡点。反复嵌套内函数，并不会立刻得到 \allowbreak{}\NovelInline{I05342}{\psi(\Omega_\omega\times2)}\allowbreak{}；此时先得到 \allowbreak{}\NovelInline{I05343}{\psi(\Omega_\omega+{\psi_{\Omega_2}(\Omega_\omega)})}\allowbreak{}，随后还有 \allowbreak{}\NovelInline{I05344}{\psi(\Omega_\omega+{\psi_{\Omega_2}(\Omega_\omega)}\times2)}\allowbreak{}、\allowbreak{}\NovelInline{I05345}{\psi(\Omega_\omega+{\psi_{\Omega_2}(\Omega_\omega)}\times\omega)}\allowbreak{} 等。外层后项到 \allowbreak{}\NovelInline{I05346}{\Omega_2}\allowbreak{} 仍是一道要单独跨过的界。蔡须坤把那处括号改正，接缝的牵扯才逐渐松开。

\NovelDisplay{M01249R01}{%
\psi(\Omega_\omega+\psi_{\Omega_2}(\Omega_\omega)\times2),\quad\psi(\Omega_\omega+\psi_{\Omega_2}(\Omega_\omega)\times\omega),\quad\dots%
}

两条都按完整形式核对。倍数与 \allowbreak{}\NovelInline{I05347}{\omega}\allowbreak{} 都在内函数外面，没有进入它的参数。

\NovelDisplay{M01250R01}{%
(0)(1,1,1)(1,1)(2,2,1)(2,2)=\psi(\Omega_{\omega}+\Omega_2)%
}

新增 \allowbreak{}\NovelInline{I05348}{(2,2)}\allowbreak{} 后，外层参数中的后项到 \allowbreak{}\NovelInline{I05349}{\Omega_2}\allowbreak{}；这一列的作用要结合前缀读。

夏夜星空站在横台下面，看着他们把这一列重新运转完。``这只是下一步。一步也得认清。''

\NovelDisplay{M01251R01}{%
(0)(1,1,1)(1,1)(2,2,1)(2,2)(3,1)=\psi(\Omega_{\omega}+\Omega_2\times\Omega)%
}

固定新出现的 \allowbreak{}\NovelInline{I05350}{\Omega_2}\allowbreak{}，让它右乘 \allowbreak{}\NovelInline{I05351}{\Omega}\allowbreak{}。

\NovelDisplay{M01252R01}{%
(0)(1,1,1)(1,1)(2,2,1)(2,2)(3,1)(4,2)=\psi(\Omega_{\omega}+\Omega_2\times\psi_{\Omega_2}(\Omega_2))%
}

再接 \allowbreak{}\NovelInline{I05352}{(4,2)}\allowbreak{}，右因子变为 \allowbreak{}\NovelInline{I05353}{{\psi_{\Omega_2}(\Omega_2)}}\allowbreak{}。

\NovelDisplay{M01253R01}{%
(0)(1,1,1)(1,1)(2,2,1)(2,2)(3,1)(4,2,1)=\psi(\Omega_{\omega}+\Omega_2\times\psi_{\Omega_2}(\Omega_{\omega}))%
}

末列的第三行变一，右因子改为 \allowbreak{}\NovelInline{I05354}{{\psi_{\Omega_2}(\Omega_\omega)}}\allowbreak{}，左因子仍是 \allowbreak{}\NovelInline{I05355}{\Omega_2}\allowbreak{}。

\NovelDisplay{M01254R01}{%
(0)(1,1,1)(1,1)(2,2,1)(2,2)(3,2)=\psi(\Omega_{\omega}+\Omega_2^{2})%
}

末列 \allowbreak{}\NovelInline{I05356}{(3,2)}\allowbreak{} 与前面的 \allowbreak{}\NovelInline{I05357}{(2,2)}\allowbreak{} 配合，直接形成 \allowbreak{}\NovelInline{I05358}{\Omega_2^2}\allowbreak{}。

这里 \allowbreak{}\NovelInline{I05359}{(2,2)(3,2)}\allowbreak{} 在外层参数中产生 \allowbreak{}\NovelInline{I05360}{\Omega_2^2}\allowbreak{}。与上面把平方放进 \allowbreak{}\NovelInline{I05361}{\psi_{\Omega_2}}\allowbreak{} 参数中的那组对照，作用层次已经不同。

\NovelDisplay{M01255R01}{%
(0)(1,1,1)(1,1)(2,2,1)(2,2)(3,3)=\psi(\Omega_{\omega}+\psi_{\Omega_3}(\Omega_3))%
}

第二坐标再抬高，后项成为 \allowbreak{}\NovelInline{I05362}{\psi_{\Omega_3}(\Omega_3)}\allowbreak{}，进入下一层坍缩。

\NovelDisplay{M01256R01}{%
(0)(1,1,1)(1,1)(2,2,1)(2,2)(3,3)(3,3)=\psi(\Omega_{\omega}+\psi_{\Omega_3}(\Omega_3\times2))%
}

重复 \allowbreak{}\NovelInline{I05363}{(3,3)}\allowbreak{}，在这层新函数的参数里增加一份 \allowbreak{}\NovelInline{I05364}{\Omega_3}\allowbreak{}。

\NovelDisplay{M01257R01}{%
(0)(1,1,1)(1,1)(2,2,1)(2,2)(3,3,1)=\psi(\Omega_\omega+\psi_{\Omega_3}(\Omega_\omega))%
}

末列出现第三行的一，将该函数的参数推进到 \allowbreak{}\NovelInline{I05365}{\Omega_\omega}\allowbreak{}。

\NovelDisplay{M01258R01}{%
(0)(1,1,1)(1,1)(2,2,1)(2,2)(3,3,1)(3,3)=\psi(\Omega_{\omega}+\Omega_3)%
}

追加 \allowbreak{}\NovelInline{I05366}{(3,3)}\allowbreak{}，外层参数的后项达到 \allowbreak{}\NovelInline{I05367}{\Omega_3}\allowbreak{}。

\NovelDisplay{M01259R01}{%
(0)(1,1,1)(1,1)(2,2,1)(2,2)(3,3,1)(3,3)%
}
\NovelDisplay{M01259R02}{%
\quad (4,2)(5,3)=\psi(\Omega_{\omega}+\Omega_3\times\psi_{\Omega_3}(\Omega_3))%
}

把新出现的 \allowbreak{}\NovelInline{I05368}{\Omega_3}\allowbreak{} 右乘 \allowbreak{}\NovelInline{I05369}{\psi_{\Omega_3}(\Omega_3)}\allowbreak{}，注意函数下标比 \allowbreak{}\NovelInline{I05370}{{\psi_{\Omega_2}(\Omega_2)}}\allowbreak{} 的下标高一层。

\NovelDisplay{M01260R01}{%
(0)(1,1,1)(1,1)(2,2,1)(2,2)(3,3,1)(3,3)%
}
\NovelDisplay{M01260R02}{%
\quad (4,2)(5,3,1)=\psi(\Omega_{\omega}+\Omega_3\times\psi_{\Omega_3}(\Omega_{\omega}))%
}

只改末列第三行，右因子变为 \allowbreak{}\NovelInline{I05371}{\psi_{\Omega_3}(\Omega_\omega)}\allowbreak{}。

\NovelDisplay{M01261R01}{%
(0)(1,1,1)(1,1)(2,2,1)(2,2)(3,3,1)(3,3)(4,3)=\psi(\Omega_{\omega}+\Omega_3^{2})%
}

这里在外层参数中形成 \allowbreak{}\NovelInline{I05372}{\Omega_3^2}\allowbreak{}，并未把最外面的 \allowbreak{}\NovelInline{I05373}{\psi}\allowbreak{} 平方。

\NovelDisplay{M01262R01}{%
(0)(1,1,1)(1,1)(2,2,1)(2,2)(3,3,1)(3,3)(4,4)=\psi(\Omega_{\omega}+\psi_{\Omega_4}(\Omega_4))%
}

第二坐标升到四，出现 \allowbreak{}\NovelInline{I05374}{\psi_{\Omega_4}(\Omega_4)}\allowbreak{}；这仍只是后项的继续增长。

他们把前面几组重新连起来，才走到这一阶段的结果。

第二十日午后，四人把此前的接缝又核对了一轮。蔡须坤踩实台阶，霍雨浩记下结果，古月娜从末列重新追根。卢本韦看着纸，轻轻敲了敲最初那个零。

``现在才轮到第二份。''

\NovelDisplay{M01263R01}{%
(0)(1,1,1)(1,1,1)=\psi(\Omega_\omega\times2)%
}

到这里，第二个 \allowbreak{}\NovelInline{I05375}{\Omega_\omega}\allowbreak{} 才真正进入最外层参数；前面整段都不能省掉。

末项 \allowbreak{}\NovelInline{I05376}{(1,1,1)}\allowbreak{} 在本段追根到 \allowbreak{}\NovelInline{I05377}{(0)}\allowbreak{}，阶差为 \allowbreak{}\NovelInline{I05378}{(1,1)}\allowbreak{}；各行要按父项关系判断是否随复制段抬升。教材先用带美元符号的占位符表示中间片段：

\NovelDisplay{M01264R01}{%
(0)(\$_0)(1,1)(\$_1)(2,2)(\$_2)\dots%
}

\NovelDisplay{M01265R01}{%
\$_i\longleftrightarrow\mathfrak b_i\qquad\text{都只代表一段待填的结构，不是一个数字列。}%
}

为免把占位符看成数值，中间的片段改用 \allowbreak{}\NovelInline{I05379}{\mathfrak b_i}\allowbreak{} 占位，记成下面的结构示意；\allowbreak{}\NovelInline{I05380}{\mathfrak b_i}\allowbreak{} 不是数字列。

\NovelDisplay{M01266R01}{%
(0)\,\mathfrak b_0\,(1,1)\,\mathfrak b_1\,(2,2)\,\mathfrak b_2\,\cdots%
}

其中 \allowbreak{}\NovelInline{I05381}{(1,1),(2,2),\ldots}\allowbreak{} 引出相应的 \allowbreak{}\NovelInline{I05382}{\psi_{\Omega_2},\psi_{\Omega_3},\ldots}\allowbreak{} 层次；由这条基本列的极限，外层再出现一份 \allowbreak{}\NovelInline{I05383}{\Omega_\omega}\allowbreak{}。``极限''不是做了某个很大的有限次数。

``同样的结构再走一次。''卢本韦说。

\NovelDisplay{M01267R01}{%
(0)(1,1,1)(1,1,1)(1,1,1)=\psi(\Omega_\omega\times3)%
}

再接一列 \allowbreak{}\NovelInline{I05384}{(1,1,1)}\allowbreak{}，多一份 \allowbreak{}\NovelInline{I05385}{\Omega_\omega}\allowbreak{}；此处作用范围尚未改变。

同样追根、复制、取对应极限，参数中再添一份 \allowbreak{}\NovelInline{I05386}{\Omega_\omega}\allowbreak{}。这一次，古月娜先指出根，霍雨浩核对后才落下标记。

\NovelDisplay{M01268R01}{%
(0)(1,1,1)(2)=\psi(\Omega_\omega\times\omega)%
}

末列 \allowbreak{}\NovelInline{I05387}{(2)}\allowbreak{} 把外层参数的有限系数推进到 \allowbreak{}\NovelInline{I05388}{\omega}\allowbreak{}。

\NovelDisplay{M01269R01}{%
(0)(1,1,1)(2)(1,1,1)=\psi(\Omega_\omega\times\omega+\Omega_\omega)%
}

在 \allowbreak{}\NovelInline{I05389}{\Omega_\omega\times\omega}\allowbreak{} 后再加 \allowbreak{}\NovelInline{I05390}{\Omega_\omega}\allowbreak{}；序数加法按书写顺序进行。

\NovelDisplay{M01270R01}{%
(0)(1,1,1)(2)(2)=\psi(\Omega_\omega\times\omega^2)%
}

再接 \allowbreak{}\NovelInline{I05391}{(2)}\allowbreak{}，右端系数到 \allowbreak{}\NovelInline{I05392}{\omega^2}\allowbreak{}，仍是同一个左因子。

\NovelDisplay{M01271R01}{%
(0)(1,1,1)(2)(3,1)=\psi(\Omega_{\omega}\times\psi(\Omega))%
}

\allowbreak{}\NovelInline{I05393}{(3,1)}\allowbreak{} 在这里把右系数推进到 \allowbreak{}\NovelInline{I05394}{\psi(\Omega)}\allowbreak{}。

\NovelDisplay{M01272R01}{%
(0)(1,1,1)(2)(3,1,1)=\psi(\Omega_\omega\times\psi(\Omega_\omega))%
}

第三行的一让右系数达到 \allowbreak{}\NovelInline{I05395}{{\psi(\Omega_\omega)}}\allowbreak{}，要看清这个 \allowbreak{}\NovelInline{I05396}{{\psi(\Omega_\omega)}}\allowbreak{} 仍位于最外层参数之内。

\NovelDisplay{M01273R01}{%
(0)(1,1,1)(2)(3,1,1)(3,1,1)=\psi(\Omega_\omega\times\psi(\Omega_\omega\times2))%
}

重复 \allowbreak{}\NovelInline{I05397}{(3,1,1)}\allowbreak{}，改变的是系数自身的坍缩参数，成为 \allowbreak{}\NovelInline{I05398}{\psi(\Omega_\omega\times2)}\allowbreak{}。

\NovelDisplay{M01274R01}{%
(0)(1,1,1)(2)(3,1,1)(4)(5,1,1)=\psi(\Omega_\omega\times\psi(\Omega_\omega\times\psi(\Omega_\omega)))%
}

再隔出一层，系数的参数包含 \allowbreak{}\NovelInline{I05399}{\Omega_\omega\times {\psi(\Omega_\omega)}}\allowbreak{}，形成更深的嵌套。

\NovelDisplay{M01275R01}{%
(0)(1,1,1)(2,1)=\psi(\Omega_\omega\times\Omega)%
}

最后，右系数走到 \allowbreak{}\NovelInline{I05400}{\Omega}\allowbreak{}，完成这一段分析；前面的 \allowbreak{}\NovelInline{I05401}{\Omega_\omega}\allowbreak{} 一直保留。

到此，\allowbreak{}\NovelInline{I05402}{(1,1,1)}\allowbreak{} 在所考察的位置上都表现为在 \allowbreak{}\NovelInline{I05403}{\psi}\allowbreak{} 参数内添一份 \allowbreak{}\NovelInline{I05404}{\Omega_\omega}\allowbreak{}。蔡须坤据此提出下一步的猜想：

\NovelDisplay{M01276R01}{%
(0)(1,1,1)(2,1)(1,1,1)\stackrel{?}{=}\psi(\Omega_\omega\times\Omega+\Omega_\omega)%
}

霍雨浩没有把问号擦掉。此前的局部规律，尚不能保证放在这个新位置仍然成立。

第二十日黄昏，七十米的石桩已经就在前方。

她照着猜想去接，后项刚亮，前面已经站稳的结构忽然向回牵。卢本韦托住她的胳膊，等她退回横台，才让她收住魂力。

``不是力量不够。''蔡须坤低声说。

她的魂力早已越过这层。如今压力拖慢运转，旧日熟练带过的接缝才逐处显露。霍雨浩把回扯发生的位置记在问号旁，请她明日仍在同一阶试。

古月娜蹲下，指尖停在末列的一上。她先追第一行，再追第二行，追到最后，目光落回前面那列 \allowbreak{}\NovelInline{I05405}{(2,1)}\allowbreak{}。

``后面这个一一一，牵动的还不只是它自己。''

夏夜星空把那张带问号的纸夹回本子。``明天从这里起，找中间的路。''

诺当将本子装进干布袋，奥斯卡递过水囊，一个人喝完，再递给下一个。蔡兴坤坐在石桩边，把黑子鸡抱回怀里。唐三等最后一人站稳，才收回蓝银草。

霍雨浩先将目标的前两项展开写在末页，留待下一日核对。

他们采用从零编号的基本列约定：零项只截去末列，一项接上一份完整复制段。展开朝较小的结构走；登山时做的是理解并重新组织已有力量，不能把展开后列数增加当成序数变大。

\NovelDisplay{M01277R01}{%
\bigl((0)(1,1,1)(2,1)(1,1,1)\bigr)[0]=(0)(1,1,1)(2,1),%
}
\NovelDisplay{M01277R02}{%
\bigl((0)(1,1,1)(2,1)(1,1,1)\bigr)[1]=(0)(1,1,1)(2,1)(1,1)(2,2,1)(3,2)%
}

登阶的前二十日，他们暂歇在七十米石桩下方。五百米的横线还在山图上头，谁也没有把猜想后的问号擦去。

\chapter{第一百零五章：替身挡住的那一步}\label{ux7b2cux4e00ux767eux96f6ux4e94ux7ae0ux66ffux8eabux6321ux4f4fux7684ux90a3ux4e00ux6b65}

第二十一日清早，蔡须坤把带问号的纸摊开。过去二十日，四人每天登阶受压，退回横台记录，再带着推导重试。一行式子，有时要试许多天。

\NovelDisplay{M01278R01}{%
(0)(1,1,1)(2,1)(1,1,1)\stackrel{?}{=}\psi(\Omega_\omega\times\Omega+\Omega_\omega)%
}

``今天不先填答案。''蔡须坤说，``找出这个一一一为什么接不上。''

七十米石桩下方，山道分出一层窄台。唐三先把蓝银草系在石桩上，另一端留给后面的人扶稳。夏夜星空站在窄台另一头，指尖向下，石面浮出三排细纹。

``护住身体可以，里面的次序，要各自走完。''

霍雨浩把前二十日的记录续在页首。以下的省略零都按三行补齐参与计算；式中没有下标的坍缩函数仍是可数层的那一个。

\NovelDisplay{M01279R01}{%
\psi=\psi_\Omega,\qquad\psi_{\Omega_2}(0)=\Omega%
}

``参数和乘数都写全。''霍雨浩对卢本韦说，``每层都能直接看见。''

卢本韦翻回前页，望着末尾的 \allowbreak{}\NovelInline{I05406}{(3,2)}\allowbreak{}，将旧记录中的 \allowbreak{}\NovelInline{I05407}{(3,1)}\allowbreak{} 圈出。

``缺的就是第二行这一点？''

``这一点怎么来的，得看完。''古月娜答。

\NovelDisplay{M01280R01}{%
(0)(1,1,1)(2,1)(1,1)=\psi(\Omega_\omega\times\Omega+\Omega)%
}

与上一章末式相比，新增的是最外层参数的后项 \allowbreak{}\NovelInline{I05408}{\Omega}\allowbreak{}，前项 \allowbreak{}\NovelInline{I05409}{\Omega_\omega\times\Omega}\allowbreak{} 固定；先不要把这个加号读成幂次。

\NovelDisplay{M01281R01}{%
(0)(1,1,1)(2,1)(1,1)(2)(3,1,1)=\psi(\Omega_{\omega}\times\Omega+\Omega\times\psi(\Omega_{\omega}))%
}

新增片段使后项成为 \allowbreak{}\NovelInline{I05410}{\Omega\times\psi(\Omega_\omega)}\allowbreak{}。乘法在最外层参数内，右边的坍缩函数必须先辨清作用范围。

\NovelDisplay{M01282R01}{%
(0)(1,1,1)(2,1)(1,1)(2,1)=\psi(\Omega_{\omega}\times\Omega+\Omega^{2})%
}

这里的后项改为 \allowbreak{}\NovelInline{I05411}{\Omega^2}\allowbreak{}。同样接在 \allowbreak{}\NovelInline{I05412}{(1,1)}\allowbreak{} 后面，末列的第二坐标不同，所得结构也不同。

\NovelDisplay{M01283R01}{%
(0)(1,1,1)(2,1)(1,1)(2,2)=\psi(\Omega_{\omega}\times\Omega+\psi_{\Omega_2}(\Omega_2))%
}

后项由普通的 \allowbreak{}\NovelInline{I05413}{\Omega}\allowbreak{} 幂进入 \allowbreak{}\NovelInline{I05414}{\psi_{\Omega_2}(\Omega_2)}\allowbreak{}；函数上界与函数参数各写一遍，不能合成一个下标。

\NovelDisplay{M01284R01}{%
(0)(1,1,1)(2,1)(1,1)(2,2,1)=\psi(\Omega_\omega\times\Omega+\psi_{\Omega_2}(\Omega_\omega))%
}

末列第三坐标从省略的零变成一，内层参数达到 \allowbreak{}\NovelInline{I05415}{\Omega_\omega}\allowbreak{}；内函数仍是 \allowbreak{}\NovelInline{I05416}{\psi_{\Omega_2}}\allowbreak{}，没有换上界。

到这里，几条仍可沿用此前方法。蔡须坤想把它们合成``再添一个 \allowbreak{}\NovelInline{I05417}{{\Omega_\omega}}\allowbreak{}''，卢本韦却感到两处受阻不同。霍雨浩请她们先比较各自的根。

\NovelDisplay{M01285R01}{%
(0)(1,1,1)(2,1)(1,1)(2,2,1)(3,1)=\psi(\Omega_\omega\times\Omega+\psi_{\Omega_2}(\Omega_\omega\times\Omega))%
}

\allowbreak{}\NovelInline{I05418}{(3,1)}\allowbreak{} 在内函数参数中生成右因子 \allowbreak{}\NovelInline{I05419}{\Omega}\allowbreak{}。外层也有一个 \allowbreak{}\NovelInline{I05420}{\Omega}\allowbreak{}，但两处与各自根列的阶差并不相同。

这就是第一处关键。外层代表 \allowbreak{}\NovelInline{I05421}{\Omega}\allowbreak{} 的列相对最初根的阶差，和内层代表 \allowbreak{}\NovelInline{I05422}{\Omega}\allowbreak{} 的列相对自己根的阶差，并不相同。末尾零虽不写，参与计算时仍在原处。

\NovelDisplay{M01286R01}{%
\Delta_{\rm outer}=(2,1),\qquad\Delta_{\rm inner}=(2),\qquad(2,1)\ne(2)%
}

\NovelDisplay{M01287R01}{%
\delta_1^{\rm outer}=2-0=2,\quad\delta_2^{\rm outer}=1-0=1,\quad%
}
\NovelDisplay{M01287R02}{%
\quad \delta_1^{\rm inner}=3-1=2,\quad\delta_2^{\rm inner}=1-1=0%
}

``第一行相等，第二行还没齐。''古月娜说。

\NovelDisplay{M01288R01}{%
\text{外层参照：}(0)\longrightarrow(2,1),\qquad\text{内层参照：}(1,1)\longrightarrow(3,1)%
}

\NovelDisplay{M01289R01}{%
\text{若内层对应列达到 }(3,2),\qquad3-1=2,\quad2-1=1%
}

这些箭头表示比较两列，不是把原有前缀真的改掉。要让新结构成立，必须经过能生成后一种阶差的中间基本列。额外的一轮结构，正是提升效应的来源。

\NovelDisplay{M01290R01}{%
(0)(1,1,1)(2,1)(1,1)(2,2,1)(3,1)(2,2)=\psi(\Omega_\omega\times\Omega+\Omega_2)%
}

加上 \allowbreak{}\NovelInline{I05423}{(2,2)}\allowbreak{} 后，最外层的后项达到 \allowbreak{}\NovelInline{I05424}{\Omega_2}\allowbreak{}。这里不是把前式内函数的右端直接换成一个字母，须沿完整基本列核对。

\NovelDisplay{M01291R01}{%
(0)(1,1,1)(2,1)(1,1)(2,2,1)(3,1)(2,2)(3,1)%
}
\NovelDisplay{M01291R02}{%
\quad (4,2)=\psi(\Omega_{\omega}\times\Omega+\Omega_2\times\psi_{\Omega_2}(\Omega_2))%
}

新的 \allowbreak{}\NovelInline{I05425}{(3,1)(4,2)}\allowbreak{} 作用于后项，给出 \allowbreak{}\NovelInline{I05426}{\Omega_2\times\psi_{\Omega_2}(\Omega_2)}\allowbreak{}；左边第一份 \allowbreak{}\NovelInline{I05427}{\Omega_2}\allowbreak{} 不在内函数中。

\NovelDisplay{M01292R01}{%
(0)(1,1,1)(2,1)(1,1)(2,2,1)(3,1)(2,2)(3,2)=\psi(\Omega_{\omega}\times\Omega+\Omega_2^{2})%
}

相应末列变为 \allowbreak{}\NovelInline{I05428}{(3,2)}\allowbreak{}，后项形成 \allowbreak{}\NovelInline{I05429}{\Omega_2^2}\allowbreak{}。这一步对应平方，不能写成 \allowbreak{}\NovelInline{I05430}{\Omega_2\times2}\allowbreak{}。

\NovelDisplay{M01293R01}{%
(0)(1,1,1)(2,1)(1,1)(2,2,1)(3,1)(2,2)(3,2)%
}
\NovelDisplay{M01293R02}{%
\quad (1,1)(2,2,1)=\psi(\Omega_{\omega}\times\Omega+\Omega_2^{2}+\psi_{\Omega_2}(\Omega_{\omega}))%
}

前面的平方保留，尾段重新开出 \allowbreak{}\NovelInline{I05431}{\psi_{\Omega_2}(\Omega_\omega)}\allowbreak{}。这是相加，不是把整个外层参数一起乘内函数。

\NovelDisplay{M01294R01}{%
(0)(1,1,1)(2,1)(1,1)(2,2,1)(3,1)(2,2)(3,2)(1,1)(2,2,1)(3,1)(2,2)(3,2)%
}
\NovelDisplay{M01294R02}{%
\quad =\psi(\Omega_{\omega}\times\Omega+\Omega_2^{2}+\psi_{\Omega_2}(\Omega_{\omega}\times\Omega+\Omega_2^{2}))%
}

内函数重新走过与外部相似的片段，参数中也出现 \allowbreak{}\NovelInline{I05432}{\Omega_2^2}\allowbreak{}；两个平方位于不同括号层次。

\NovelDisplay{M01295R01}{%
(0)(1,1,1)(2,1)(1,1)(2,2,1)(3,1)(2,2)%
}
\NovelDisplay{M01295R02}{%
\quad (3,2)(2,2)=\psi(\Omega_{\omega}\times\Omega+\Omega_2^{2}+\Omega_2)%
}

同一外部平方后面增加 \allowbreak{}\NovelInline{I05433}{\Omega_2}\allowbreak{}，有序和为 \allowbreak{}\NovelInline{I05434}{\Omega_2^2+\Omega_2}\allowbreak{}；这里不能交换两项的位置。

\NovelDisplay{M01296R01}{%
(0)(1,1,1)(2,1)(1,1)(2,2,1)(3,1)(2,2)(3,2)(2,2)(3,1)(4,2,1)=\psi(\Omega_{\omega}%
}
\NovelDisplay{M01296R02}{%
\quad \times\Omega+\Omega_2^{2}+\Omega_2\times\psi_{\Omega_2}(\Omega_{\omega}))%
}

只将第二项扩展成 \allowbreak{}\NovelInline{I05435}{\Omega_2\times\psi_{\Omega_2}(\Omega_\omega)}\allowbreak{}，第一项的平方保持不动。

\NovelDisplay{M01297R01}{%
(0)(1,1,1)(2,1)(1,1)(2,2,1)(3,1)(2,2)%
}
\NovelDisplay{M01297R02}{%
\quad (3,2)(2,2)(3,2)=\psi(\Omega_{\omega}\times\Omega+\Omega_2^{2}\times2)%
}

现在才有两份完整的 \allowbreak{}\NovelInline{I05436}{\Omega_2^2}\allowbreak{}，写作 \allowbreak{}\NovelInline{I05437}{\Omega_2^2\times2}\allowbreak{}；与上面的 \allowbreak{}\NovelInline{I05438}{\Omega_2^2+\Omega_2}\allowbreak{} 对照阅读。

\NovelDisplay{M01298R01}{%
(0)(1,1,1)(2,1)(1,1)(2,2,1)(3,1)(2,2)(3,2)(3)=\psi(\Omega_{\omega}\times\Omega+\Omega_2^{2}\times\omega)%
}

有限份平方的重复在这一基本列中走向右乘 \allowbreak{}\NovelInline{I05439}{\omega}\allowbreak{}；左边的平方没有变为 \allowbreak{}\NovelInline{I05440}{\Omega_2^\omega}\allowbreak{}。

\NovelDisplay{M01299R01}{%
(0)(1,1,1)(2,1)(1,1)(2,2,1)(3,1)(2,2)(3,2)(3,2)=\psi(\Omega_{\omega}\times\Omega+\Omega_2^{3})%
}

末列第二坐标升高后，多出的是同底数的一次幂，故为 \allowbreak{}\NovelInline{I05441}{\Omega_2^3}\allowbreak{}，不是三份平方。

\NovelDisplay{M01300R01}{%
(0)(1,1,1)(2,1)(1,1)(2,2,1)(3,1)(2,2)(3,3)=\psi(\Omega_{\omega}\times\Omega+\psi_{\Omega_3}(\Omega_3))%
}

后项进入新的坍缩上界 \allowbreak{}\NovelInline{I05442}{\Omega_3}\allowbreak{}，参数也为 \allowbreak{}\NovelInline{I05443}{\Omega_3}\allowbreak{}。下标升一层，与上一式多一次幂是两种变化。

\NovelDisplay{M01301R01}{%
(0)(1,1,1)(2,1)(1,1)(2,2,1)(3,1)(2,2)(3,3,1)=\psi(\Omega_\omega\times\Omega+\psi_{\Omega_3}(\Omega_\omega))%
}

函数上界固定为 \allowbreak{}\NovelInline{I05444}{\Omega_3}\allowbreak{}，参数从 \allowbreak{}\NovelInline{I05445}{\Omega_3}\allowbreak{} 进到 \allowbreak{}\NovelInline{I05446}{\Omega_\omega}\allowbreak{}；第三行的一要在完整列中判断。

\NovelDisplay{M01302R01}{%
(0)(1,1,1)(2,1)(1,1)(2,2,1)(3,1)(2,2)(3,3,1)%
}
\NovelDisplay{M01302R02}{%
\quad (4,1)= \psi(\Omega_\omega\times\Omega+\psi_{\Omega_3}(\Omega_\omega\times\Omega))%
}

新接 \allowbreak{}\NovelInline{I05447}{(4,1)}\allowbreak{} 把内函数参数改为 \allowbreak{}\NovelInline{I05448}{\Omega_\omega\times\Omega}\allowbreak{}。这一项将继续承担替身的作用。

\NovelDisplay{M01303R01}{%
(0)(1,1,1)(2,1)(1,1)(2,2,1)(3,1)(2,2)%
}
\NovelDisplay{M01303R02}{%
\quad (3,3,1)(4,1)(3,3)(4,3)=\psi(\Omega_{\omega}\times\Omega+\Omega_3^{2})%
}

内侧相应结构走完后，最外层后项为 \allowbreak{}\NovelInline{I05449}{\Omega_3^2}\allowbreak{}；此前最初的 \allowbreak{}\NovelInline{I05450}{\Omega_\omega\times\Omega}\allowbreak{} 仍未改变。

\NovelDisplay{M01304R01}{%
(0)(1,1,1)(2,1)(1,1)(2,2,1)(3,1)(2,2,1)= (0)(1,1,1)(2,1)(1,1)(2,2,1)%
}
\NovelDisplay{M01304R02}{%
\quad (3,1)(2,2)(3,3,1)(4,1)(3,3)(4,4,1)(5,1)\dots= \psi(\Omega_\omega\times(\Omega+1))%
}

此处写出的长链表示基本列的趋近，不是把无限多列当作一条有限 BMS。结果只增加一份 \allowbreak{}\NovelInline{I05451}{\Omega_\omega}\allowbreak{}，尚未得到目标末列 \allowbreak{}\NovelInline{I05452}{(1,1,1)}\allowbreak{}。

第三十日，蔡须坤仍在七十米石桩下方试着给整段加上第二行的一，膝盖立刻被压弯。卢本韦托住她；古月娜留在同一阶，只改变这一行，感到牵扯恰从那列传来。

``不是整块每个数都加。''古月娜用指尖点着 \allowbreak{}\NovelInline{I05453}{(3,1)}\allowbreak{} 的第二行。

这一行的一，父项追到前面的 \allowbreak{}\NovelInline{I05454}{(0)}\allowbreak{}，没有经过这次坏根 \allowbreak{}\NovelInline{I05455}{(1,1)}\allowbreak{}。所以复制时第一行可以抬升，第二行仍留在一。霍雨浩给这一列单独画出了祖先经过条件。

\NovelDisplay{M01305R01}{%
r=(1,1),\quad\Delta=(1,1),\quad\text{第二行祖先链：}\ (3,1)\longrightarrow(0)\quad\text{绕过 }r%
}

\NovelDisplay{M01306R01}{%
\text{该列第一份复制：}\ (3,1)\longmapsto(4,1),\qquad(4,2)\ \text{是此处误算}%
}

\NovelDisplay{M01307R01}{%
\text{同块中经过坏根的列：}\ (1,1)\longmapsto(2,2),\qquad(2,2,1)\longmapsto(3,3,1)%
}

较高行的父项沿前一行的父项链找更小的数，不能只沿整张表横着找最近的小数。第二行是否跟着加，是祖先关系决定的；末列的第三行只用于定根，不在这一阶差里加一。

教材把这个现象叫``替身''。\allowbreak{}\NovelInline{I05456}{(3,1)}\allowbreak{} 使眼前这个 \allowbreak{}\NovelInline{I05457}{(1,1)}\allowbreak{} 无法按原来的直觉直接接成 \allowbreak{}\NovelInline{I05458}{(1,1,1)}\allowbreak{}；它没有凭空创造新规则，而是让坏根和复制关系多走一层。

\NovelDisplay{M01308R01}{%
\text{被阻住的局部期待：}\quad(1,1)\longrightarrow(1,1,1)%
}

\NovelDisplay{M01309R01}{%
\text{本段实际先走：}\quad(2,2)\longrightarrow(2,2,1),\qquad\psi({\Omega_\omega}\times\Omega+{\Omega_\omega})%
}

卢本韦在 \allowbreak{}\NovelInline{I05459}{(3,1)}\allowbreak{} 下画了一面小盾。蔡须坤笑得有些喘：``原来卡在这里。''

四人轮流退到横台歇息，将刚走通的部分留在左边，右边从新背景重推。

\NovelDisplay{M01310R01}{%
(0)(1,1,1)(2,1)(1,1)(2,2,1)(3,1)(2,2,1)(1,1)%
}
\NovelDisplay{M01310R02}{%
\quad (2,2,1)= \psi(\Omega_\omega\times\Omega+\Omega_\omega+\psi_{\Omega_2}(\Omega_\omega))%
}

以已经得到的 \allowbreak{}\NovelInline{I05460}{\Omega_\omega\times\Omega+\Omega_\omega}\allowbreak{} 为背景，后面又开一层内函数；新开层不会自动消去前缀。

\NovelDisplay{M01311R01}{%
(0)(1,1,1)(2,1)(1,1)(2,2,1)(3,1)(2,2,1)(1,1)(2,2,1)(3,1)= \psi(\Omega_\omega\times\Omega%
}
\NovelDisplay{M01311R02}{%
\quad +\Omega_\omega+\psi_{\Omega_2}(\Omega_\omega\times\Omega))%
}

这一层内函数参数再次右乘 \allowbreak{}\NovelInline{I05461}{\Omega}\allowbreak{}，注意它在第二个加号之后，不能移到最外层前项上。

\NovelDisplay{M01312R01}{%
(0)(1,1,1)(2,1)(1,1)(2,2,1)(3,1)(2,2,1)(1,1)(2,2,1)(3,1)(2,2)(3,1)%
}
\NovelDisplay{M01312R02}{%
\quad (4,2,1)=\psi(\Omega_{\omega}\times\Omega+\Omega_{\omega}+\psi_{\Omega_2}(\Omega_{\omega}\times\Omega+\Omega_2\times\psi_{\Omega_2}(\Omega_{\omega})))%
}

新贡献落在最内侧的后项：\allowbreak{}\NovelInline{I05462}{\Omega_2\times\psi_{\Omega_2}(\Omega_\omega)}\allowbreak{}；由内到外读完两层坍缩，再回到外层和。

\NovelDisplay{M01313R01}{%
(0)(1,1,1)(2,1)(1,1)(2,2,1)(3,1)(2,2,1)(1,1)(2,2,1)(3,1)(2,2)(3,2)%
}
\NovelDisplay{M01313R02}{%
\quad =\psi(\Omega_{\omega}\times\Omega+\Omega_{\omega}+\psi_{\Omega_2}(\Omega_{\omega}\times\Omega+\Omega_2^{2}))%
}

最内侧的后项成为 \allowbreak{}\NovelInline{I05463}{\Omega_2^2}\allowbreak{}，外层第二份 \allowbreak{}\NovelInline{I05464}{\Omega_\omega}\allowbreak{} 仍与它隔着 \allowbreak{}\NovelInline{I05465}{\psi_{\Omega_2}}\allowbreak{}。

\NovelDisplay{M01314R01}{%
(0)(1,1,1)(2,1)(1,1)(2,2,1)(3,1)(2,2,1)(1,1)(2,2,1)(3,1)(2,2)(3,3)%
}
\NovelDisplay{M01314R02}{%
\quad =\psi(\Omega_{\omega}\times\Omega+\Omega_{\omega}+\psi_{\Omega_2}(\Omega_{\omega}\times\Omega+\psi_{\Omega_3}(\Omega_3)))%
}

在内层参数中引入 \allowbreak{}\NovelInline{I05466}{\psi_{\Omega_3}(\Omega_3)}\allowbreak{}，变化是上界更高的嵌套，外面一层上界仍为 \allowbreak{}\NovelInline{I05467}{\Omega_2}\allowbreak{}。

\NovelDisplay{M01315R01}{%
(0)(1,1,1)(2,1)(1,1)(2,2,1)(3,1)(2,2,1)(1,1)(2,2,1)(3,1)(2,2)(3,3,1)%
}
\NovelDisplay{M01315R02}{%
\quad = \psi(\Omega_\omega\times\Omega+\Omega_\omega+\psi_{\Omega_2}(\Omega_\omega\times\Omega+\psi_{\Omega_3}(\Omega_\omega)))%
}

最内函数参数进到 \allowbreak{}\NovelInline{I05468}{\Omega_\omega}\allowbreak{}；两个函数的下标分别是 \allowbreak{}\NovelInline{I05469}{\Omega_2}\allowbreak{} 与 \allowbreak{}\NovelInline{I05470}{\Omega_3}\allowbreak{}，不能因参数相同而合并。

\NovelDisplay{M01316R01}{%
(0)(1,1,1)(2,1)(1,1)(2,2,1)(3,1)(2,2,1)(1,1)(2,2,1)(3,1)(2,2,1)%
}
\NovelDisplay{M01316R02}{%
\quad = \psi(\Omega_\omega\times\Omega+\Omega_\omega+\psi_{\Omega_2}(\Omega_\omega\times\Omega+\Omega_\omega))%
}

内层参数也有了两项 \allowbreak{}\NovelInline{I05471}{\Omega_\omega\times\Omega+\Omega_\omega}\allowbreak{}，与外部背景同形；``同形''仍不等于在同一层。

\NovelDisplay{M01317R01}{%
(0)(1,1,1)(2,1)(1,1)(2,2,1)(3,1)(2,2,1)(2)=\psi(\Omega_\omega\times(\Omega+1)%
}
\NovelDisplay{M01317R02}{%
\quad +\psi_{\Omega_2}(\Omega_\omega\times(\Omega+1)+1))%
}

末尾一行结构在内层共同背景后加一。这里的 \allowbreak{}\NovelInline{I05472}{+1}\allowbreak{} 在 \allowbreak{}\NovelInline{I05473}{\psi_{\Omega_2}}\allowbreak{} 内，不能加到全式最后。

\NovelDisplay{M01318R01}{%
(0)(1,1,1)(2,1)(1,1)(2,2,1)(3,1)(2,2,1)(2,2)=\psi(\Omega_\omega\times(\Omega+1)+\Omega_2)%
}

后项达到外层可见的 \allowbreak{}\NovelInline{I05474}{\Omega_2}\allowbreak{}。括号外的背景也可写成 \allowbreak{}\NovelInline{I05475}{\Omega_\omega\times(\Omega+1)}\allowbreak{}，乘法顺序不变。

\NovelDisplay{M01319R01}{%
(0)(1,1,1)(2,1)(1,1)(2,2,1)(3,1)(2,2,1)(2,2)(3,1)(4,2,1)=\psi(\Omega_{\omega}\times\Omega%
}
\NovelDisplay{M01319R02}{%
\quad +\Omega_{\omega}+\Omega_2\times\psi_{\Omega_2}(\Omega_{\omega}))%
}

在外部 \allowbreak{}\NovelInline{I05476}{\Omega_2}\allowbreak{} 后接一份可数右因子，得到 \allowbreak{}\NovelInline{I05477}{\Omega_2\times\psi_{\Omega_2}(\Omega_\omega)}\allowbreak{}，乘号没有进入该函数参数。

\NovelDisplay{M01320R01}{%
(0)(1,1,1)(2,1)(1,1)(2,2,1)(3,1)(2,2,1)%
}
\NovelDisplay{M01320R02}{%
\quad (2,2)(3,2)=\psi(\Omega_{\omega}\times\Omega+\Omega_{\omega}+\Omega_2^{2})%
}

后项为 \allowbreak{}\NovelInline{I05478}{\Omega_2^2}\allowbreak{}，与第一轮同型；区别在于前面已经有第二份 \allowbreak{}\NovelInline{I05479}{\Omega_\omega}\allowbreak{}。

熟悉的 \allowbreak{}\NovelInline{I05480}{(2,2)(3,2)}\allowbreak{} 再次出现。卢本韦先指向最外层加号，请蔡须坤确认平方落在哪层。

\NovelDisplay{M01321R01}{%
(0)(1,1,1)(2,1)(1,1)(2,2,1)(3,1)(2,2,1)(2,2)%
}
\NovelDisplay{M01321R02}{%
\quad (3,2)(2,2)(3,2)=\psi(\Omega_{\omega}\times\Omega+\Omega_{\omega}+\Omega_2^{2}\times2)%
}

两份平方在最外层参数内相加，缩写为 \allowbreak{}\NovelInline{I05481}{\Omega_2^2\times2}\allowbreak{}；本式保留的两份 \allowbreak{}\NovelInline{I05482}{\Omega_\omega}\allowbreak{} 不计入这个二。

\NovelDisplay{M01322R01}{%
(0)(1,1,1)(2,1)(1,1)(2,2,1)(3,1)(2,2,1)%
}
\NovelDisplay{M01322R02}{%
\quad (2,2)(3,2)(3,2)=\psi(\Omega_{\omega}\times\Omega+\Omega_{\omega}+\Omega_2^{3})%
}

这一项将后部平方推进为立方。霍雨浩把幂指数上的三圈出，另把最外层两个加号连在同一条线上。

\NovelDisplay{M01323R01}{%
(0)(1,1,1)(2,1)(1,1)(2,2,1)(3,1)(2,2,1)%
}
\NovelDisplay{M01323R02}{%
\quad (2,2)(3,3)=\psi(\Omega_{\omega}\times\Omega+\Omega_{\omega}+\psi_{\Omega_3}(\Omega_3))%
}

后项出现 \allowbreak{}\NovelInline{I05483}{\psi_{\Omega_3}(\Omega_3)}\allowbreak{}。这轮也要重走更高上界，不能看到相似前缀就提前跳到结果。

\NovelDisplay{M01324R01}{%
(0)(1,1,1)(2,1)(1,1)(2,2,1)(3,1)(2,2,1)(2,2)%
}
\NovelDisplay{M01324R02}{%
\quad (3,3,1)(4,1)= \psi(\Omega_\omega\times(\Omega+1)+\psi_{\Omega_3}(\Omega_\omega\times\Omega))%
}

\allowbreak{}\NovelInline{I05484}{\psi_{\Omega_3}}\allowbreak{} 的参数达到 \allowbreak{}\NovelInline{I05485}{\Omega_\omega\times\Omega}\allowbreak{}；左边新出现的 \allowbreak{}\NovelInline{I05486}{(4,1)}\allowbreak{} 要与自己的根比较。

\NovelDisplay{M01325R01}{%
(0)(1,1,1)(2,1)(1,1)(2,2,1)(3,1)(2,2,1)(2,2)(3,3,1)(4,1)(3,3)=\psi(\Omega_{\omega}%
}
\NovelDisplay{M01325R02}{%
\quad \times\Omega+\Omega_{\omega}+\psi_{\Omega_3}(\Omega_{\omega}\times\Omega+\Omega_3))%
}

在这个内函数里增加 \allowbreak{}\NovelInline{I05487}{\Omega_3}\allowbreak{}，不是在全式最外层增加 \allowbreak{}\NovelInline{I05488}{\Omega_3}\allowbreak{}；右括号给出了作用范围。

\NovelDisplay{M01326R01}{%
(0)(1,1,1)(2,1)(1,1)(2,2,1)(3,1)(2,2,1)(2,2)(3,3,1)(4,1)(3,3)(4,1)%
}
\NovelDisplay{M01326R02}{%
\quad (5,2,1)=\psi(\Omega_{\omega}\times\Omega+\Omega_{\omega}+\psi_{\Omega_3}(\Omega_{\omega}\times\Omega+\Omega_3\times\psi_{\Omega_2}(\Omega_{\omega})))%
}

内层后项变为 \allowbreak{}\NovelInline{I05489}{\Omega_3\times\psi_{\Omega_2}(\Omega_\omega)}\allowbreak{}。这里更深处的上界回到 \allowbreak{}\NovelInline{I05490}{\Omega_2}\allowbreak{}，不要求层层严格增大。

\NovelDisplay{M01327R01}{%
(0)(1,1,1)(2,1)(1,1)(2,2,1)(3,1)(2,2,1)(2,2)(3,3,1)(4,1)(3,3)(4,2)%
}
\NovelDisplay{M01327R02}{%
\quad =\psi(\Omega_{\omega}\times\Omega+\Omega_{\omega}+\psi_{\Omega_3}(\Omega_{\omega}\times\Omega+\Omega_3\times\Omega_2))%
}

内层后项中的右因子为 \allowbreak{}\NovelInline{I05491}{\Omega_2}\allowbreak{}，形成 \allowbreak{}\NovelInline{I05492}{\Omega_3\times\Omega_2}\allowbreak{}；不可把两个不同基数乘积读成平方。

\NovelDisplay{M01328R01}{%
(0)(1,1,1)(2,1)(1,1)(2,2,1)(3,1)(2,2,1)(2,2)(3,3,1)(4,1)(3,3)(4,2)(5,3)%
}
\NovelDisplay{M01328R02}{%
\quad =\psi(\Omega_{\omega}\times\Omega+\Omega_{\omega}+\psi_{\Omega_3}(\Omega_{\omega}\times\Omega+\Omega_3\times\psi_{\Omega_3}(\Omega_3)))%
}

该右因子换成 \allowbreak{}\NovelInline{I05493}{\psi_{\Omega_3}(\Omega_3)}\allowbreak{}，其函数上界与外面一层相同，仍须保留两对括号。

\NovelDisplay{M01329R01}{%
(0)(1,1,1)(2,1)(1,1)(2,2,1)(3,1)(2,2,1)(2,2)(3,3,1)(4,1)(3,3)(4,2)%
}
\NovelDisplay{M01329R02}{%
\quad (5,3,1)=\psi(\Omega_{\omega}\times\Omega+\Omega_{\omega}+\psi_{\Omega_3}(\Omega_{\omega}\times\Omega+\Omega_3\times\psi_{\Omega_3}(\Omega_{\omega})))%
}

内侧函数的参数变为 \allowbreak{}\NovelInline{I05494}{\Omega_\omega}\allowbreak{}，外围 \allowbreak{}\NovelInline{I05495}{\psi_{\Omega_3}}\allowbreak{} 的参数仍以 \allowbreak{}\NovelInline{I05496}{\Omega_\omega\times\Omega}\allowbreak{} 起头。

\NovelDisplay{M01330R01}{%
(0)(1,1,1)(2,1)(1,1)(2,2,1)(3,1)(2,2,1)(2,2)(3,3,1)(4,1)(3,3)(4,3)%
}
\NovelDisplay{M01330R02}{%
\quad =\psi(\Omega_{\omega}\times\Omega+\Omega_{\omega}+\psi_{\Omega_3}(\Omega_{\omega}\times\Omega+\Omega_3^{2}))%
}

内层后项此时形成 \allowbreak{}\NovelInline{I05497}{\Omega_3^2}\allowbreak{}，两个 \allowbreak{}\NovelInline{I05498}{\Omega_3}\allowbreak{} 因子在同一参数层内；区别于前面的嵌套函数。

\NovelDisplay{M01331R01}{%
(0)(1,1,1)(2,1)(1,1)(2,2,1)(3,1)(2,2,1)(2,2)(3,3,1)(4,1)(3,3)(4,4,1)%
}
\NovelDisplay{M01331R02}{%
\quad =\psi(\Omega_{\omega}\times\Omega+\Omega_{\omega}+\psi_{\Omega_3}(\Omega_{\omega}\times\Omega+\psi_{\Omega_4}(\Omega_{\omega})))%
}

本式采用底稿当前保留的转换结果；它另有来源候选，未获指定表格确认。先记为待核的对应，不拿它单独推出新通则。

\NovelDisplay{M01332R01}{%
(0)(1,1,1)(2,1)(1,1)(2,2,1)(3,1)(2,2,1)(2,2)(3,3,1)(4,1)(3,3,1)%
}
\NovelDisplay{M01332R02}{%
\quad =\psi(\Omega_{\omega}\times\Omega+\Omega_{\omega}+\psi_{\Omega_3}(\Omega_{\omega}\times\Omega+\Omega_{\omega}))%
}

这一条同样保留待核标记。它与上一条的候选表达式相似，恰好说明不能只凭右端眼熟就交换两条 BMS。

\NovelDisplay{M01333R01}{%
(0)(1,1,1)(2,1)(1,1)(2,2,1)(3,1)(2,2,1)(2,2,1)=\psi(\Omega_\omega\times(\Omega+2))%
}

新接的 \allowbreak{}\NovelInline{I05499}{(2,2,1)}\allowbreak{} 又提供一份 \allowbreak{}\NovelInline{I05500}{\Omega_\omega}\allowbreak{}；两份合并为右系数二，得到 \allowbreak{}\NovelInline{I05501}{\Omega_\omega\times(\Omega+2)}\allowbreak{}。

第四十日，练习横台中段的石纹亮了。原来``添一份 \allowbreak{}\NovelInline{I05502}{{\Omega_\omega}}\allowbreak{}''的作用改由 \allowbreak{}\NovelInline{I05503}{(2,2,1)}\allowbreak{} 承担，替身 \allowbreak{}\NovelInline{I05504}{(3,1)}\allowbreak{} 仍在前面。霍雨浩对照压力下的记录，圈出变化的位置与背景。

\NovelDisplay{M01334R01}{%
(0)(1,1,1)(2,1)(1,1)(2,2,1)(3,1)(2,2,1)%
}
\NovelDisplay{M01334R02}{%
\quad (2,2,1)(2,2,1)= \psi(\Omega_\omega\times\Omega+\Omega_\omega\times3)%
}

同一位置再添一项，有限部分成为三；这里的三在乘法右侧括号内，尚未升到 \allowbreak{}\NovelInline{I05505}{\Omega\times3}\allowbreak{}。

\NovelDisplay{M01335R01}{%
(0)(1,1,1)(2,1)(1,1)(2,2,1)(3,1)(2,2,1)(3)=\psi(\Omega_\omega\times\Omega+\Omega_\omega\times\omega)%
}

有限个后加项的极限是 \allowbreak{}\NovelInline{I05506}{\Omega_\omega\times\omega}\allowbreak{}，前项 \allowbreak{}\NovelInline{I05507}{\Omega_\omega\times\Omega}\allowbreak{} 仍然存在。

\NovelDisplay{M01336R01}{%
(0)(1,1,1)(2,1)(1,1)(2,2,1)(3,1)(2,2,1)(3,1)=\psi(\Omega_\omega\times(\Omega\times2))%
}

已有一份 \allowbreak{}\NovelInline{I05508}{\Omega_\omega\times\Omega}\allowbreak{}，后面又形成同样一份，故为 \allowbreak{}\NovelInline{I05509}{\Omega_\omega\times(\Omega\times2)}\allowbreak{}。

\NovelDisplay{M01337R01}{%
(0)(1,1,1)(2,1)(1,1)(2,2,1)(3,1)(2,2,1)(3,1)%
}
\NovelDisplay{M01337R02}{%
\quad (1,1)(2,2,1)= \psi(\Omega_\omega\times\Omega\times2+\psi_{\Omega_2}(\Omega_\omega))%
}

固定两份大段之后，继续增加 \allowbreak{}\NovelInline{I05510}{\psi_{\Omega_2}(\Omega_\omega)}\allowbreak{}；这是更高背景上的同类尾部分析。

\NovelDisplay{M01338R01}{%
(0)(1,1,1)(2,1)(1,1)(2,2,1)(3,1)(2,2,1)(3,1)(1,1)(2,2,1)(3,1)(2,2,1)%
}
\NovelDisplay{M01338R02}{%
\quad (3,1)= \psi(\Omega_\omega\times\Omega\times2+\psi_{\Omega_2}(\Omega_\omega\times\Omega\times2))%
}

这次内函数参数已经包含 \allowbreak{}\NovelInline{I05511}{\Omega\times2}\allowbreak{}。原来的替身结构仍在，不能只把外部乘数改二、内部留一。

\NovelDisplay{M01339R01}{%
(0)(1,1,1)(2,1)(1,1)(2,2,1)(3,1)(2,2,1)%
}
\NovelDisplay{M01339R02}{%
\quad (3,1)(2,2)=\psi(\Omega_{\omega}\times\Omega\times2+\Omega_2)%
}

新尾部达到 \allowbreak{}\NovelInline{I05512}{\Omega_2}\allowbreak{}，但大段右因子仍是 \allowbreak{}\NovelInline{I05513}{\Omega\times2}\allowbreak{}，它没有因后项达到基数就自动升阶。

\NovelDisplay{M01340R01}{%
(0)(1,1,1)(2,1)(1,1)(2,2,1)(3,1)(2,2,1)%
}
\NovelDisplay{M01340R02}{%
\quad (3,1)(2,2)(3,2)=\psi(\Omega_{\omega}\times\Omega\times2+\Omega_2^{2})%
}

再得到外部后项 \allowbreak{}\NovelInline{I05514}{\Omega_2^2}\allowbreak{}；这一平方属于和的末项，与左边那一个有限乘数二各司其职。

\NovelDisplay{M01341R01}{%
(0)(1,1,1)(2,1)(1,1)(2,2,1)(3,1)(2,2,1)%
}
\NovelDisplay{M01341R02}{%
\quad (3,1)(2,2,1)= \psi(\Omega_\omega\times\Omega\times2+\Omega_\omega)%
}

外部尾项达到 \allowbreak{}\NovelInline{I05515}{\Omega_\omega}\allowbreak{}，于是两份大段之后又接一份小段，尚不能合并成第三份大段。

\NovelDisplay{M01342R01}{%
(0)(1,1,1)(2,1)(1,1)(2,2,1)(3,1)(3)=\psi(\Omega_\omega\times(\Omega\times\omega))%
}

右因子由有限份 \allowbreak{}\NovelInline{I05516}{\Omega}\allowbreak{} 进入 \allowbreak{}\NovelInline{I05517}{\Omega\times\omega}\allowbreak{}。括号保留乘法次序，不能写成 \allowbreak{}\NovelInline{I05518}{\omega\times\Omega}\allowbreak{}。

\NovelDisplay{M01343R01}{%
(0)(1,1,1)(2,1)(1,1)(2,2,1)(3,1)(3,1)=\psi(\Omega_\omega\times\Omega^{2})%
}

右因子达到 \allowbreak{}\NovelInline{I05519}{\Omega^2}\allowbreak{}，左因子仍为 \allowbreak{}\NovelInline{I05520}{\Omega_\omega}\allowbreak{}；不要误把结果看成 \allowbreak{}\NovelInline{I05521}{\Omega_\omega^2}\allowbreak{}。

\NovelDisplay{M01344R01}{%
(0)(1,1,1)(2,1)(1,1)(2,2,1)(3,1)(4)=\psi(\Omega_\omega\times\Omega^{\omega})%
}

本组推进的是右因子的指数，得到 \allowbreak{}\NovelInline{I05522}{\Omega^\omega}\allowbreak{}，不是以 \allowbreak{}\NovelInline{I05523}{\Omega_\omega}\allowbreak{} 为底的幂。

\NovelDisplay{M01345R01}{%
(0)(1,1,1)(2,1)(1,1)(2,2,1)(3,1)(4,1)=\psi(\Omega_\omega\times\Omega^{\Omega})%
}

右因子成为 \allowbreak{}\NovelInline{I05524}{\Omega^\Omega}\allowbreak{}。唐三在共同的左因子下画了一道蓝线，让每个人看清正在改变的位置。

\NovelDisplay{M01346R01}{%
(0)(1,1,1)(2,1)(1,1)(2,2,1)(3,1)(4,2)=\psi(\Omega_{\omega}\times\psi_{\Omega_2}(\Omega_2))%
}

右因子转为 \allowbreak{}\NovelInline{I05525}{\psi_{\Omega_2}(\Omega_2)}\allowbreak{}，从普通幂进入坍缩层；乘法仍在最外层参数内部。

\NovelDisplay{M01347R01}{%
(0)(1,1,1)(2,1)(1,1)(2,2,1)(3,1)(4,2)(4,2)=\psi(\Omega_{\omega}\times\psi_{\Omega_2}(\Omega_2\times2))%
}

该内函数的参数增加第二份 \allowbreak{}\NovelInline{I05526}{\Omega_2}\allowbreak{}。应读成函数参数乘二，不能读成函数值乘二。

\NovelDisplay{M01348R01}{%
(0)(1,1,1)(2,1)(1,1)(2,2,1)(3,1)(4,2)(5,1)=\psi(\Omega_{\omega}\times\psi_{\Omega_2}(\Omega_2\times\Omega))%
}

内函数参数的右因子达到 \allowbreak{}\NovelInline{I05527}{\Omega}\allowbreak{}，即 \allowbreak{}\NovelInline{I05528}{\Omega_2\times\Omega}\allowbreak{}；外部乘法与内部乘法分属两层。

\NovelDisplay{M01349R01}{%
(0)(1,1,1)(2,1)(1,1)(2,2,1)(3,1)(4,2)(5,1)%
}
\NovelDisplay{M01349R02}{%
\quad (6,2)=\psi(\Omega_{\omega}\times\psi_{\Omega_2}(\Omega_2\times\psi_{\Omega_2}(\Omega_2)))%
}

内函数参数里出现 \allowbreak{}\NovelInline{I05529}{\Omega_2\times\psi_{\Omega_2}(\Omega_2)}\allowbreak{}，同一上界的内函数又嵌入一次。

\NovelDisplay{M01350R01}{%
(0)(1,1,1)(2,1)(1,1)(2,2,1)(3,1)(4,2)(5,2)=\psi(\Omega_{\omega}\times\psi_{\Omega_2}(\Omega_2^{2}))%
}

这里得到的是 \allowbreak{}\NovelInline{I05530}{\psi_{\Omega_2}(\Omega_2^2)}\allowbreak{}，平方在函数内部；与 \allowbreak{}\NovelInline{I05531}{\psi_{\Omega_2}(\Omega_2)^2}\allowbreak{} 分开。

\NovelDisplay{M01351R01}{%
(0)(1,1,1)(2,1)(1,1)(2,2,1)(3,1)(4,2)(5,3)=\psi(\Omega_{\omega}\times\psi_{\Omega_2}(\Omega_3))%
}

右因子的函数参数进入 \allowbreak{}\NovelInline{I05532}{\Omega_3}\allowbreak{}，上界仍为 \allowbreak{}\NovelInline{I05533}{\Omega_2}\allowbreak{}。参数取值高于下标并不意味着把函数下标一同改掉。

\NovelDisplay{M01352R01}{%
(0)(1,1,1)(2,1)(1,1)(2,2,1)(3,1)(4,2,1)=\psi(\Omega_\omega\times\psi_{\Omega_2}(\Omega_\omega))%
}

参数达到 \allowbreak{}\NovelInline{I05534}{\Omega_\omega}\allowbreak{}，与最外层左因子用同一个符号；一处是函数参数，一处是乘数。

\NovelDisplay{M01353R01}{%
(0)(1,1,1)(2,1)(1,1)(2,2,1)(3,1)(4,2,1)(4,2,1)=\psi(\Omega_\omega\times\psi_{\Omega_2}(\Omega_\omega\times2))%
}

内部参数先形成两份 \allowbreak{}\NovelInline{I05535}{\Omega_\omega}\allowbreak{}。这不是把整个 \allowbreak{}\NovelInline{I05536}{\Omega_\omega\times\psi_{\Omega_2}(\Omega_\omega)}\allowbreak{} 复制两遍。

\NovelDisplay{M01354R01}{%
(0)(1,1,1)(2,1)(1,1)(2,2,1)(3,1)(4,2,1)(5,1)=\psi(\Omega_\omega\times\psi_{\Omega_2}(\Omega_\omega\times\Omega))%
}

内部参数成为 \allowbreak{}\NovelInline{I05537}{\Omega_\omega\times\Omega}\allowbreak{}，末尾 \allowbreak{}\NovelInline{I05538}{(5,1)}\allowbreak{} 在这层承担新的替身作用。

这次挡在内层的是 \allowbreak{}\NovelInline{I05539}{(5,1)}\allowbreak{}。它使对应的 \allowbreak{}\NovelInline{I05540}{(3,1)}\allowbreak{} 不能立刻完成所期待的进位。与第一轮一样，要继续补齐内侧参数，再回到外部系数。

\NovelDisplay{M01355R01}{%
\text{外部关注：}\quad\Omega_\omega\times\boxed{\psi_{\Omega_2}(\Omega_\omega\times\Omega)}%
}

\NovelDisplay{M01356R01}{%
\text{内部仍要比较：}\quad {\Omega_\omega}\times\Omega,\quad {\Omega_\omega}\times\Omega+{\Omega_\omega},\quad {\Omega_\omega}\times\Omega\times2%
}

第四十六日，霍雨浩在压力下辨出两层受阻的先后。古月娜沿父项链复核，蔡须坤与卢本韦各试一次，才将它记入本子。

\NovelDisplay{M01357R01}{%
(0)(1,1,1)(2,1)(1,1)(2,2,1)(3,1)(4,2,1)(5,1)(4,2)(5,1)(6,2,1)=\psi(\Omega_{\omega}%
}
\NovelDisplay{M01357R02}{%
\quad \times\psi_{\Omega_2}(\Omega_{\omega}\times\Omega+\Omega_2\times\psi_{\Omega_2}(\Omega_{\omega})))%
}

内层后项为 \allowbreak{}\NovelInline{I05541}{\Omega_2\times\psi_{\Omega_2}(\Omega_\omega)}\allowbreak{}；各层最初的 \allowbreak{}\NovelInline{I05542}{\Omega_\omega}\allowbreak{} 要沿右括号分别归属。

\NovelDisplay{M01358R01}{%
(0)(1,1,1)(2,1)(1,1)(2,2,1)(3,1)(4,2,1)(5,1)(4,2)(5,2)=\psi(\Omega_{\omega}%
}
\NovelDisplay{M01358R02}{%
\quad \times\psi_{\Omega_2}(\Omega_{\omega}\times\Omega+\Omega_2^{2}))%
}

同一后项形成 \allowbreak{}\NovelInline{I05543}{\Omega_2^2}\allowbreak{}，改变的仍是内函数参数的末项，而非外面共同的左因子。

\NovelDisplay{M01359R01}{%
(0)(1,1,1)(2,1)(1,1)(2,2,1)(3,1)(4,2,1)(5,1)(4,2)(5,3,1)= \psi(\Omega_\omega%
}
\NovelDisplay{M01359R02}{%
\quad \times\psi_{\Omega_2}(\Omega_\omega\times\Omega+\psi_{\Omega_3}(\Omega_\omega)))%
}

内层后项进入 \allowbreak{}\NovelInline{I05544}{\psi_{\Omega_3}(\Omega_\omega)}\allowbreak{}；要先求这一层结构，再组成上界为 \allowbreak{}\NovelInline{I05545}{\Omega_2}\allowbreak{} 的外一层。

\NovelDisplay{M01360R01}{%
(0)(1,1,1)(2,1)(1,1)(2,2,1)(3,1)(4,2,1)(5,1)%
}
\NovelDisplay{M01360R02}{%
\quad (4,2,1)= \psi(\Omega_\omega\times\psi_{\Omega_2}(\Omega_\omega\times\Omega+\Omega_\omega))%
}

内部后项达到 \allowbreak{}\NovelInline{I05546}{\Omega_\omega}\allowbreak{}，所以该参数是 \allowbreak{}\NovelInline{I05547}{\Omega_\omega\times\Omega+\Omega_\omega}\allowbreak{}，尚不等于平方。

\NovelDisplay{M01361R01}{%
(0)(1,1,1)(2,1)(1,1)(2,2,1)(3,1)(4,2,1)(5,1)%
}
\NovelDisplay{M01361R02}{%
\quad (4,2,1)(5,1)= \psi(\Omega_\omega\times\psi_{\Omega_2}(\Omega_\omega\times\Omega\times2))%
}

内部参数中形成两份 \allowbreak{}\NovelInline{I05548}{\Omega_\omega\times\Omega}\allowbreak{}，末尾 \allowbreak{}\NovelInline{I05549}{(5,1)}\allowbreak{} 仍在限制进一步进位。

\NovelDisplay{M01362R01}{%
(0)(1,1,1)(2,1)(1,1)(2,2,1)(3,1)(4,2,1)(5,1)%
}
\NovelDisplay{M01362R02}{%
\quad (6,2)=\psi(\Omega_{\omega}\times\psi_{\Omega_2}(\Omega_{\omega}\times\psi_{\Omega_2}(\Omega_2)))%
}

内部右因子进入 \allowbreak{}\NovelInline{I05550}{\psi_{\Omega_2}(\Omega_2)}\allowbreak{}，这一层也开始重复外面走过的坍缩结构。

\NovelDisplay{M01363R01}{%
(0)(1,1,1)(2,1)(1,1)(2,2,1)(3,1)(4,2,1)(5,1)%
}
\NovelDisplay{M01363R02}{%
\quad (6,2,1)= \psi(\Omega_\omega\times\psi_{\Omega_2}(\Omega_\omega\times\psi_{\Omega_2}(\Omega_\omega)))%
}

更内侧参数达到 \allowbreak{}\NovelInline{I05551}{\Omega_\omega}\allowbreak{}。相似片段再嵌一次，使基本列更长，不会因此变成简单的有限乘法。

\NovelDisplay{M01364R01}{%
(0)(1,1,1)(2,1)(1,1)(2,2,1)(3,2)=\psi(\Omega_\omega\times\Omega_2)%
}

此时外部右因子才达到 \allowbreak{}\NovelInline{I05552}{\Omega_2}\allowbreak{}，即 \allowbreak{}\NovelInline{I05553}{\psi(\Omega_\omega\times\Omega_2)}\allowbreak{}，完成第一轮系数提升。

\allowbreak{}\NovelInline{I05554}{(3,1)}\allowbreak{} 对应的位置终于变成 \allowbreak{}\NovelInline{I05555}{(3,2)}\allowbreak{}，本轮替身的作用结束。这里说的是前后两条不同 BMS 的相应位置，不是展开时把已经保留的前缀悄悄改了。

\NovelDisplay{M01365R01}{%
\psi({\Omega_\omega}\times\Omega)\quad\leadsto\quad\psi({\Omega_\omega}\times\Omega_2)%
}

第五十日，卢本韦终于在台阶上稳住这一轮。``提升的是后面的系数，最初左边的 \allowbreak{}\NovelInline{I05556}{{\Omega_\omega}}\allowbreak{} 还在。''

\NovelDisplay{M01366R01}{%
(0)(1,1,1)(2,1)(1,1)(2,2,1)(3,2)(1,1)(2,2,1)=\psi(\Omega_\omega\times\Omega_2+\psi_{\Omega_2}(\Omega_\omega))%
}

在新的共同前项 \allowbreak{}\NovelInline{I05557}{\Omega_\omega\times\Omega_2}\allowbreak{} 之后，从 \allowbreak{}\NovelInline{I05558}{\psi_{\Omega_2}(\Omega_\omega)}\allowbreak{} 开始核对尾部。

\NovelDisplay{M01367R01}{%
(0)(1,1,1)(2,1)(1,1)(2,2,1)(3,2)(1,1)(2,2,1)%
}
\NovelDisplay{M01367R02}{%
\quad (3,1)= \psi(\Omega_\omega\times\Omega_2+\psi_{\Omega_2}(\Omega_\omega\times\Omega))%
}

尾部内函数参数仍先形成 \allowbreak{}\NovelInline{I05559}{\Omega_\omega\times\Omega}\allowbreak{}，不因背景已到 \allowbreak{}\NovelInline{I05560}{\Omega_2}\allowbreak{} 就直接跳到同样背景。

\NovelDisplay{M01368R01}{%
(0)(1,1,1)(2,1)(1,1)(2,2,1)(3,2)(1,1)(2,2,1)(3,1)(2,2,1)= \psi(\Omega_\omega\times\Omega_2%
}
\NovelDisplay{M01368R02}{%
\quad +\psi_{\Omega_2}(\Omega_\omega\times\Omega+\Omega_\omega))%
}

内部参数增加一份 \allowbreak{}\NovelInline{I05561}{\Omega_\omega}\allowbreak{}，这仍是一道局部的小步，外部背景保持 \allowbreak{}\NovelInline{I05562}{\Omega_\omega\times\Omega_2}\allowbreak{}。

\NovelDisplay{M01369R01}{%
(0)(1,1,1)(2,1)(1,1)(2,2,1)(3,2)(1,1)(2,2,1)(3,1)(2,2,1)(3,1)= \psi(\Omega_\omega%
}
\NovelDisplay{M01369R02}{%
\quad \times\Omega_2+\psi_{\Omega_2}(\Omega_\omega\times\Omega\times2))%
}

内部参数的右因子形成 \allowbreak{}\NovelInline{I05563}{\Omega\times2}\allowbreak{}，继续保持乘法次序；二不属于最外层共同前项。

\NovelDisplay{M01370R01}{%
(0)(1,1,1)(2,1)(1,1)(2,2,1)(3,2)(1,1)(2,2,1)(3,1)(4,2)=\psi(\Omega_{\omega}\times\Omega_2%
}
\NovelDisplay{M01370R02}{%
\quad +\psi_{\Omega_2}(\Omega_{\omega}\times\psi_{\Omega_2}(\Omega_2)))%
}

内部右因子进入 \allowbreak{}\NovelInline{I05564}{\psi_{\Omega_2}(\Omega_2)}\allowbreak{}，意味着又有一层函数需要单独核对参数。

\NovelDisplay{M01371R01}{%
(0)(1,1,1)(2,1)(1,1)(2,2,1)(3,2)(1,1)(2,2,1)(3,1)(4,2,1)= \psi(\Omega_\omega\times\Omega_2%
}
\NovelDisplay{M01371R02}{%
\quad +\psi_{\Omega_2}(\Omega_\omega\times\psi_{\Omega_2}(\Omega_\omega)))%
}

更内层参数达到 \allowbreak{}\NovelInline{I05565}{\Omega_\omega}\allowbreak{}，尚未令最外面新尾项成为 \allowbreak{}\NovelInline{I05566}{\Omega_2}\allowbreak{}，括号层次不能省。

\NovelDisplay{M01372R01}{%
(0)(1,1,1)(2,1)(1,1)(2,2,1)(3,2)(1,1)(2,2,1)%
}
\NovelDisplay{M01372R02}{%
\quad (3,2)= \psi(\Omega_\omega\times\Omega_2+\psi_{\Omega_2}(\Omega_\omega\times\Omega_2))%
}

尾部函数参数终于追到共同背景 \allowbreak{}\NovelInline{I05567}{\Omega_\omega\times\Omega_2}\allowbreak{}。结构对齐后，后继变化才有了新的参照。

\NovelDisplay{M01373R01}{%
(0)(1,1,1)(2,1)(1,1)(2,2,1)(3,2)(2)=\psi(\Omega_\omega\times\Omega_2+\psi_{\Omega_2}(\Omega_\omega\times\Omega_2+1))%
}

一行末列把内函数参数改为共同背景加一。这个一只在内函数中，外部仍是原有背景加该函数。

\NovelDisplay{M01374R01}{%
(0)(1,1,1)(2,1)(1,1)(2,2,1)(3,2)(2,1)(3,2,1)%
}
\NovelDisplay{M01374R02}{%
\quad (4,2)= \psi(\Omega_\omega\times\Omega_2+\psi_{\Omega_2}(\Omega_\omega\times\Omega_2)^2)%
}

这里平方落在整个尾部函数值上，即 \allowbreak{}\NovelInline{I05568}{\bigl(\psi_{\Omega_2}(\Omega_\omega\times\Omega_2)\bigr)^2}\allowbreak{}，别把平方移入函数。

\NovelDisplay{M01375R01}{%
(0)(1,1,1)(2,1)(1,1)(2,2,1)(3,2)(2,2)=\psi(\Omega_\omega\times\Omega_2+\Omega_2)%
}

新的外部尾项达到 \allowbreak{}\NovelInline{I05569}{\Omega_2}\allowbreak{}，结束这一小段内函数的追赶。共同前项的右因子不变。

\NovelDisplay{M01376R01}{%
(0)(1,1,1)(2,1)(1,1)(2,2,1)(3,2)(2,2)(3,1)%
}
\NovelDisplay{M01376R02}{%
\quad (4,2,1)=\psi(\Omega_\omega\times\Omega_2+\Omega_2\times\psi_{\Omega_2}(\Omega_\omega))%
}

后项扩为 \allowbreak{}\NovelInline{I05570}{\Omega_2\times\psi_{\Omega_2}(\Omega_\omega)}\allowbreak{}，现在两处 \allowbreak{}\NovelInline{I05571}{\Omega_2}\allowbreak{} 各在不同层次。

\NovelDisplay{M01377R01}{%
(0)(1,1,1)(2,1)(1,1)(2,2,1)(3,2)(2,2)(3,1)(4,2,1)(5,2)=\psi(\Omega_{\omega}\times\Omega_2%
}
\NovelDisplay{M01377R02}{%
\quad +\Omega_2\times\psi_{\Omega_2}(\Omega_{\omega}\times\Omega_2))%
}

右因子函数的参数追到 \allowbreak{}\NovelInline{I05572}{\Omega_\omega\times\Omega_2}\allowbreak{}，不是把外部那个 \allowbreak{}\NovelInline{I05573}{\Omega_2}\allowbreak{} 也收进函数。

\NovelDisplay{M01378R01}{%
(0)(1,1,1)(2,1)(1,1)(2,2,1)(3,2)(2,2)(3,1)(4,2,1)(5,2)(4,2)=\psi(\Omega_{\omega}%
}
\NovelDisplay{M01378R02}{%
\quad \times\Omega_2+\Omega_2\times\psi_{\Omega_2}(\Omega_{\omega}\times\Omega_2+\Omega_2))%
}

该内函数的参数再加 \allowbreak{}\NovelInline{I05574}{\Omega_2}\allowbreak{}；先读内部和，再读外部乘积，最后读最外层的和。

\NovelDisplay{M01379R01}{%
(0)(1,1,1)(2,1)(1,1)(2,2,1)(3,2)(2,2)(3,2)=\psi(\Omega_\omega\times\Omega_2+\Omega_2^{2})%
}

外部尾项成为 \allowbreak{}\NovelInline{I05575}{\Omega_2^2}\allowbreak{}，它只是共同前项后面的平方，尚未改变共同前项的右因子。

\NovelDisplay{M01380R01}{%
(0)(1,1,1)(2,1)(1,1)(2,2,1)(3,2)(2,2)(3,3)=\psi(\Omega_{\omega}\times\Omega_2+\psi_{\Omega_3}(\Omega_3))%
}

尾项进入 \allowbreak{}\NovelInline{I05576}{\psi_{\Omega_3}(\Omega_3)}\allowbreak{}。这一轮对应更高的坍缩上界，需重新追踪父项。

\NovelDisplay{M01381R01}{%
(0)(1,1,1)(2,1)(1,1)(2,2,1)(3,2)(2,2)(3,3,1)=\psi(\Omega_\omega\times\Omega_2+\psi_{\Omega_3}(\Omega_\omega))%
}

尾部函数参数达到 \allowbreak{}\NovelInline{I05577}{\Omega_\omega}\allowbreak{}，函数下标仍为 \allowbreak{}\NovelInline{I05578}{\Omega_3}\allowbreak{}；读法与第一轮相同，背景不同。

\NovelDisplay{M01382R01}{%
(0)(1,1,1)(2,1)(1,1)(2,2,1)(3,2)(2,2)(3,3,1)%
}
\NovelDisplay{M01382R02}{%
\quad (4,2)= \psi(\Omega_\omega\times\Omega_2+\psi_{\Omega_3}(\Omega_\omega\times\Omega_2))%
}

内函数参数追到 \allowbreak{}\NovelInline{I05579}{\Omega_\omega\times\Omega_2}\allowbreak{}。新出现的 \allowbreak{}\NovelInline{I05580}{(4,2)}\allowbreak{} 仍有不随第二行阶差增加的部分。

再往上，新的替身是 \allowbreak{}\NovelInline{I05581}{(4,2)}\allowbreak{}。它相对自己根 \allowbreak{}\NovelInline{I05582}{(2,2)}\allowbreak{} 的第二行阶差仍为零，而本轮对应结构需要第二行差一。

\NovelDisplay{M01383R01}{%
4-2=2,\qquad2-2=0,\qquad3-2=1%
}

\NovelDisplay{M01384R01}{%
\text{现有内层阶差：}(2),\qquad\text{本轮须对齐的阶差：}(2,1)%
}

古月娜膝盖发颤，仍按父项链慢慢运转。压力让她辨出第二行的停顿，待推导核实，她才请霍雨浩记下。

\NovelDisplay{M01385R01}{%
(0)(1,1,1)(2,1)(1,1)(2,2,1)(3,2)(2,2)(3,3,1)(4,2)(3,3)(4,2)(5,3,1)(6,2)%
}
\NovelDisplay{M01385R02}{%
\quad (5,3)=\psi(\Omega_{\omega}\times\Omega_2+\Omega_3\times\psi_{\Omega_3}(\Omega_{\omega}\times\Omega_2+\Omega_3))%
}

外部尾项为 \allowbreak{}\NovelInline{I05583}{\Omega_3\times\psi_{\Omega_3}(\Omega_\omega\times\Omega_2+\Omega_3)}\allowbreak{}，乘法与内层加法别交换。

\NovelDisplay{M01386R01}{%
(0)(1,1,1)(2,1)(1,1)(2,2,1)(3,2)(2,2)(3,3,1)%
}
\NovelDisplay{M01386R02}{%
\quad (4,2)(3,3)(4,3)=\psi(\Omega_{\omega}\times\Omega_2+\Omega_3^{2})%
}

尾项形成 \allowbreak{}\NovelInline{I05584}{\Omega_3^2}\allowbreak{}，尚未令共同前项的右因子成为 \allowbreak{}\NovelInline{I05585}{\Omega_3}\allowbreak{}；完成尾部与完成整轮是两件事。

\NovelDisplay{M01387R01}{%
(0)(1,1,1)(2,1)(1,1)(2,2,1)(3,2)(2,2)(3,3,1)%
}
\NovelDisplay{M01387R02}{%
\quad (4,2)(3,3)(4,4,1)= \psi(\Omega_\omega\times\Omega_2+\psi_{\Omega_4}(\Omega_\omega))%
}

尾项继续进入 \allowbreak{}\NovelInline{I05586}{\psi_{\Omega_4}(\Omega_\omega)}\allowbreak{}，说明还要把更高上界的结构纳入这轮基本列。

\NovelDisplay{M01388R01}{%
(0)(1,1,1)(2,1)(1,1)(2,2,1)(3,2)(2,2)(3,3,1)(4,2)(3,3)(4,4,1)(5,2)(4,4)%
}
\NovelDisplay{M01388R02}{%
\quad (5,4)=\psi(\Omega_{\omega}\times\Omega_2+\Omega_4^{2})%
}

在指定结构中尾项达到 \allowbreak{}\NovelInline{I05587}{\Omega_4^2}\allowbreak{}。蔡须坤先描根，再写末列，逐行核对进位。

\NovelDisplay{M01389R01}{%
(0)(1,1,1)(2,1)(1,1)(2,2,1)(3,2)(2,2)%
}
\NovelDisplay{M01389R02}{%
\quad (3,3,1)(4,2)(3,3,1)= \psi(\Omega_\omega\times(\Omega_2+1))%
}

新末列 \allowbreak{}\NovelInline{I05588}{(3,3,1)}\allowbreak{} 提供一份 \allowbreak{}\NovelInline{I05589}{\Omega_\omega}\allowbreak{}，与共同前项合写为 \allowbreak{}\NovelInline{I05590}{\Omega_\omega\times(\Omega_2+1)}\allowbreak{}。

随后 \allowbreak{}\NovelInline{I05591}{(3,3)}\allowbreak{} 进到 \allowbreak{}\NovelInline{I05592}{(3,3,1)}\allowbreak{}，这一次由它承担增加一份 \allowbreak{}\NovelInline{I05593}{{\Omega_\omega}}\allowbreak{} 的作用。外观换了，判断办法没有换：先定共同背景，再看后项进入哪一层。

\NovelDisplay{M01390R01}{%
(0)(1,1,1)(2,1)(1,1)(2,2,1)(3,2)(2,2)(3,3,1)(4,2)(3,3,1)(3,3,1)%
}
\NovelDisplay{M01390R02}{%
\quad = \psi(\Omega_\omega\times\Omega_2+\Omega_\omega\times2)%
}

重复这份尾项，得到 \allowbreak{}\NovelInline{I05594}{\Omega_\omega\times\Omega_2+\Omega_\omega\times2}\allowbreak{}；右侧有限系数从一到二。

\NovelDisplay{M01391R01}{%
(0)(1,1,1)(2,1)(1,1)(2,2,1)(3,2)(2,2)(3,3,1)%
}
\NovelDisplay{M01391R02}{%
\quad (4,2)(3,3,1)(4)= \psi(\Omega_\omega\times(\Omega_2+\omega))%
}

有限尾项走到 \allowbreak{}\NovelInline{I05595}{\omega}\allowbreak{}，右侧括号内是 \allowbreak{}\NovelInline{I05596}{\Omega_2+\omega}\allowbreak{}，仍不等于 \allowbreak{}\NovelInline{I05597}{\Omega_2\times\omega}\allowbreak{}。

\NovelDisplay{M01392R01}{%
(0)(1,1,1)(2,1)(1,1)(2,2,1)(3,2)(2,2)(3,3,1)%
}
\NovelDisplay{M01392R02}{%
\quad (4,2)(3,3,1)(4,1)= \psi(\Omega_\omega\times\Omega_2+\Omega_\omega\times\Omega)%
}

后部右因子达到 \allowbreak{}\NovelInline{I05598}{\Omega}\allowbreak{}，成为 \allowbreak{}\NovelInline{I05599}{\Omega_\omega\times\Omega}\allowbreak{}；与共同前项仍以加法连接。

\NovelDisplay{M01393R01}{%
(0)(1,1,1)(2,1)(1,1)(2,2,1)(3,2)(2,2)(3,3,1)%
}
\NovelDisplay{M01393R02}{%
\quad (4,2)(3,3,1)(4,2)= \psi(\Omega_\omega\times(\Omega_2\times2))%
}

后部也形成一份 \allowbreak{}\NovelInline{I05600}{\Omega_\omega\times\Omega_2}\allowbreak{}，故共同右因子可写成 \allowbreak{}\NovelInline{I05601}{\Omega_2\times2}\allowbreak{}。

\NovelDisplay{M01394R01}{%
(0)(1,1,1)(2,1)(1,1)(2,2,1)(3,2)(2,2)(3,3,1)(4,2)(4)= \psi(\Omega_\omega\times\Omega_2\times\omega)%
}

有限份这类大段的极限为右乘 \allowbreak{}\NovelInline{I05602}{\omega}\allowbreak{}；这里的 \allowbreak{}\NovelInline{I05603}{\Omega_2\times\omega}\allowbreak{} 在最外层左因子之后。

\NovelDisplay{M01395R01}{%
(0)(1,1,1)(2,1)(1,1)(2,2,1)(3,2)(2,2)(3,3,1)(4,2)(4,2)= \psi(\Omega_\omega\times\Omega_2^{2})%
}

共同右因子进入 \allowbreak{}\NovelInline{I05604}{\Omega_2^2}\allowbreak{}，不应把它与整式 \allowbreak{}\NovelInline{I05605}{\Omega_\omega^2}\allowbreak{} 混淆，底数所属层次不同。

\NovelDisplay{M01396R01}{%
(0)(1,1,1)(2,1)(1,1)(2,2,1)(3,2)(2,2)(3,3,1)(4,2)(5,2)= \psi(\Omega_\omega\times\Omega_2^{\Omega_2})%
}

右因子的幂指数达到 \allowbreak{}\NovelInline{I05606}{\Omega_2}\allowbreak{}，即 \allowbreak{}\NovelInline{I05607}{\Omega_2^{\Omega_2}}\allowbreak{}，外面还保留左乘的 \allowbreak{}\NovelInline{I05608}{\Omega_\omega}\allowbreak{}。

\NovelDisplay{M01397R01}{%
(0)(1,1,1)(2,1)(1,1)(2,2,1)(3,2)(2,2)%
}
\NovelDisplay{M01397R02}{%
\quad (3,3,1)(4,2)(5,3)=\psi(\Omega_{\omega}\times\psi_{\Omega_3}(\Omega_3))%
}

共同右因子进入 \allowbreak{}\NovelInline{I05609}{\psi_{\Omega_3}(\Omega_3)}\allowbreak{}，第二轮也重复了``幂、坍缩、追赶''的过程。

\NovelDisplay{M01398R01}{%
(0)(1,1,1)(2,1)(1,1)(2,2,1)(3,2)(2,2)(3,3,1)%
}
\NovelDisplay{M01398R02}{%
\quad (4,2)(5,3)(5,3)=\psi(\Omega_{\omega}\times\psi_{\Omega_3}(\Omega_3\times2))%
}

内部参数增加第二份 \allowbreak{}\NovelInline{I05610}{\Omega_3}\allowbreak{}，并不是在函数值后面乘二，更不是复制整个最外层表达式。

\NovelDisplay{M01399R01}{%
(0)(1,1,1)(2,1)(1,1)(2,2,1)(3,2)(2,2)%
}
\NovelDisplay{M01399R02}{%
\quad (3,3,1)(4,2)(5,3,1)= \psi(\Omega_\omega\times\psi_{\Omega_3}(\Omega_\omega))%
}

内部参数达到 \allowbreak{}\NovelInline{I05611}{\Omega_\omega}\allowbreak{}。较高上界并不免去分析内部参数这一步。

\NovelDisplay{M01400R01}{%
(0)(1,1,1)(2,1)(1,1)(2,2,1)(3,2)(2,2)(3,3,1)%
}
\NovelDisplay{M01400R02}{%
\quad (4,2)(5,3,1)(6,1)= \psi(\Omega_\omega\times\psi_{\Omega_3}(\Omega_\omega\times\Omega))%
}

内部参数出现 \allowbreak{}\NovelInline{I05612}{\Omega_\omega\times\Omega}\allowbreak{}，还要继续走到本轮的 \allowbreak{}\NovelInline{I05613}{\Omega_2}\allowbreak{} 背景，再完成进位。

\NovelDisplay{M01401R01}{%
(0)(1,1,1)(2,1)(1,1)(2,2,1)(3,2)(2,2)(3,3,1)%
}
\NovelDisplay{M01401R02}{%
\quad (4,2)(5,3,1)(6,2)= \psi(\Omega_\omega\times\psi_{\Omega_3}(\Omega_\omega\times\Omega_2))%
}

内部参数追到 \allowbreak{}\NovelInline{I05614}{\Omega_\omega\times\Omega_2}\allowbreak{}，与本轮开始的共同背景对应上。

\NovelDisplay{M01402R01}{%
(0)(1,1,1)(2,1)(1,1)(2,2,1)(3,2)(2,2)(3,3,1)(4,3)= \psi(\Omega_\omega\times\Omega_3)%
}

第二轮系数提升完成，成为 \allowbreak{}\NovelInline{I05615}{\psi(\Omega_\omega\times\Omega_3)}\allowbreak{}。原来的 \allowbreak{}\NovelInline{I05616}{(4,2)}\allowbreak{} 此时才在相应位置进到 \allowbreak{}\NovelInline{I05617}{(4,3)}\allowbreak{}。

第五十七日，\allowbreak{}\NovelInline{I05618}{(4,2)}\allowbreak{} 的相应位置进到 \allowbreak{}\NovelInline{I05619}{(4,3)}\allowbreak{}，系数达到 \allowbreak{}\NovelInline{I05620}{\Omega_3}\allowbreak{}。霍雨浩与古月娜对照两轮：内部追赶、替身限制、阶差对齐，再显为外层系数提升。

\NovelDisplay{M01403R01}{%
(0)(1,1,1)(2,1)(1,1)(2,2,1)(3,2)(2,2)(3,3,1)%
}
\NovelDisplay{M01403R02}{%
\quad (4,3)(3,3)(4,4,1)(5,4)= \psi(\Omega_\omega\times\Omega_4)%
}

临近两个月时，把同样的父项分析再提升一层，得到右因子 \allowbreak{}\NovelInline{I05621}{\Omega_4}\allowbreak{}；长式显示每层都须补齐相应内部结构。

\NovelDisplay{M01404R01}{%
(0)(1,1,1)(2,1)(1,1,1)= (0)(1,1,1)(2,1)(1,1)(2,2,1)(3,2)(2,2)(3,3,1)%
}
\NovelDisplay{M01404R02}{%
\quad (4,3)(3,3)(4,4,1)(5,4)(4,4)(5,5,1)(6,5)\dots= \psi(\Omega_\omega^{2})%
}

任意有限层的系数提升组成基本列，才收束到 \allowbreak{}\NovelInline{I05622}{\psi(\Omega_\omega^2)}\allowbreak{}。省略号代表同一构造对任意有限层成立，绝不是任选一个大数便等于极限。

他们把替身一层层接下去。若用展开层次来示意，目标的内部可以排成下面的嵌套；这是一种说明基本列的非标准写法，不把无限字符串当成一条有限正规式。

\NovelDisplay{M01405R01}{%
\psi(\Omega_\omega\times\Omega+\psi_{\Omega_2}(\Omega_\omega\times\Omega_2+\psi_{\Omega_3}(\Omega_\omega\times\Omega_3+\dots)))%
}

每多一层，都要重新补齐那一层的内侧结构。随着有限层数没有上界，系数序列趋向 \allowbreak{}\NovelInline{I05623}{\Omega_\omega}\allowbreak{}；在这套对应的基本列解释下，整体收束为上面的 \allowbreak{}\NovelInline{I05624}{\psi(\Omega_\omega^2)}\allowbreak{}。

\NovelDisplay{M01406R01}{%
\sup_{n<\omega}\Omega_{n+1}=\Omega_\omega,\qquad\sup_{n<\omega}(\Omega_\omega\times\Omega_{n+1})=\Omega_\omega\times\Omega_\omega=\Omega_\omega^2%
}

这一句不是宣称任意函数都与上确界交换。这里核对的是前面逐层构造的基本列与教材给定对应，而非单凭 \allowbreak{}\NovelInline{I05625}{\psi}\allowbreak{} 这个字形推出来的连续性。

满两个月的那天，蔡须坤在七十米的石桩前站稳了。

肩上的压力仍沉。她把先前互相牵扯的几段分别接住，再按刚才确认的父项关系运转，末段终于平稳地续上。

卢本韦在旁边重新试了一遍；古月娜追完最后一条祖先链，朝霍雨浩点头。他翻回最初那张纸，才划去旧猜想。

\NovelDisplay{M01407R01}{%
\bigl((0)(1,1,1)(2,1)(1,1,1)\bigr)=\psi(\Omega_\omega^2),\bigl((0)%
}
\NovelDisplay{M01407R02}{%
\quad (1,1,1)(2,1)(1,1,1)\bigr)\ne\psi(\Omega_\omega\times\Omega+\Omega_\omega)%
}

从踏上第一阶，到分析清楚 \allowbreak{}\NovelInline{I05626}{(0)(1,1,1)(2,1)(1,1,1)}\allowbreak{}，他们用了整整两个月。霍雨浩的记录已换过一册，纸上每个回扯的位置，四人都曾在台阶上反复试过。

蔡须坤扶着石沿喘气。她把这次分析接回自己的力量，收好了笔记。

回营时，诺当掌心的旧伤已经收口。奥斯卡仍将三层油纸包好的那根香肠留在最里面。蔡须坤背上的温热贴着她；卢本韦接过本子，古月娜替霍雨浩压住被风翻起的纸。

石屋里，霍雨浩将本子放在旧砚台旁，写下这一段的三行。

\NovelDisplay{M01408R01}{%
\text{定根}\longrightarrow\text{逐行查祖先}\longrightarrow\text{只提升满足条件的复制坐标}%
}

\NovelDisplay{M01409R01}{%
\text{外层不变部分：}\ \Omega_\omega,\qquad\text{本段系数：}\ \Omega\rightsquigarrow\Omega_2\rightsquigarrow\Omega_3\rightsquigarrow\cdots%
}

\NovelDisplay{M01410R01}{%
\text{有限展开：}\ \bigl((0)(1,1,1)(2,1)(1,1,1)\bigr)[n], n<\omega;%
}
\NovelDisplay{M01410R02}{%
\text{目标本身：}\ \bigl((0)(1,1,1)(2,1)(1,1,1)\bigr)%
}
\NovelDisplay{M01410R03}{%
\text{不能用某个固定的 }\bigl((0)(1,1,1)(2,1)(1,1,1)\bigr)[n_0]\text{ 代替}%
}

两个月，他们终于分析清楚七十米这一层。五百米还在山上。

\chapter{第一百零六章：无处不在的提升}\label{ux7b2cux4e00ux767eux96f6ux516dux7ae0ux65e0ux5904ux4e0dux5728ux7684ux63d0ux5347}

登阶进入第三个月，蔡须坤从七十米处重新站稳。两个月的反复分析，使她习惯先在压力下放慢运转，辨出接缝，再与纸上的推导对照。

这次内层又有一处牵扯。卢本韦扶着石沿，与她各试一遍，古月娜沿旧记录追根。霍雨浩翻到新页，问：``再接一一一，会不会就是立方？''

霍雨浩把问题写在纸上，仍留问号。

\NovelDisplay{M01411R01}{%
(0)(1,1,1)(2,1)(1,1,1)=\psi(\Omega_\omega^2)%
}

\NovelDisplay{M01412R01}{%
\bigl((0)(1,1,1)(2,1)(1,1,1)\bigr)\,(1,1,1)\stackrel{?}{=}\psi({\Omega_\omega}^3)%
}

这里的并列表示在 BMS 右端续接数字列，不是把两个序数相乘。古月娜从最外层的和开始，银色的指尖停在第一对括号上。

``先看它加在什么位置。''

蔡须坤收慢魂力，将受阻的位置指给霍雨浩。四人仍按前两个月的办法，一处处分析已有力量的来路；其他同伴也在身后跟着核对。

\NovelDisplay{M01413R01}{%
(0)(1,1,1)(2,1)(1,1,1)(1,1)=\psi(\Omega_\omega^{2}+\Omega)%
}

以 \allowbreak{}\NovelInline{I05627}{{\Omega_\omega}^2}\allowbreak{} 为共同前项，新增的仍是外部参数中的 \allowbreak{}\NovelInline{I05628}{\Omega}\allowbreak{}；先核对这个小后项，再追到新末列。

\NovelDisplay{M01414R01}{%
(0)(1,1,1)(2,1)(1,1,1)(1,1)(2,1)=\psi(\Omega_{\omega}^{2}+\Omega^{2})%
}

后项形成 \allowbreak{}\NovelInline{I05629}{\Omega^2}\allowbreak{}，平方没有落到共同前项上。左边多一列，不等于把整个旧结构平方。

\NovelDisplay{M01415R01}{%
(0)(1,1,1)(2,1)(1,1,1)(1,1)(2,2)=\psi(\Omega_{\omega}^{2}+\psi_{\Omega_2}(\Omega_2))%
}

后项进入 \allowbreak{}\NovelInline{I05630}{\psi_{\Omega_2}(\Omega_2)}\allowbreak{}，是一个新内函数；外面共同的 \allowbreak{}\NovelInline{I05631}{{\Omega_\omega}^2}\allowbreak{} 保持不动。

\NovelDisplay{M01416R01}{%
(0)(1,1,1)(2,1)(1,1,1)(1,1)(2,2,1)=\psi(\Omega_\omega^{2}+\psi_{\Omega_2}(\Omega_\omega))%
}

内函数参数达到 \allowbreak{}\NovelInline{I05632}{{\Omega_\omega}}\allowbreak{}，函数上界仍为 \allowbreak{}\NovelInline{I05633}{\Omega_2}\allowbreak{}。与前章相似的尾部，要在新背景上重算。

\NovelDisplay{M01417R01}{%
(0)(1,1,1)(2,1)(1,1,1)(1,1)(2,2,1)(3,1)=\psi(\Omega_\omega^{2}+\psi_{\Omega_2}(\Omega_\omega\times\Omega))%
}

接上 \allowbreak{}\NovelInline{I05634}{(3,1)}\allowbreak{}，内函数参数成为 \allowbreak{}\NovelInline{I05635}{{\Omega_\omega}\times\Omega}\allowbreak{}；这一列再次限制下一步进位。

熟悉的 \allowbreak{}\NovelInline{I05636}{(3,1)}\allowbreak{} 又出现了。卢本韦先点出内层根，古月娜复核，霍雨浩才画上盾牌。蔡须坤记得此前的回扯，逐行核对阶差。

\NovelDisplay{M01418R01}{%
(0)(1,1,1)(2,1)(1,1,1)(1,1)(2,2,1)(3,1)%
}
\NovelDisplay{M01418R02}{%
\quad (2,2)=\psi(\Omega_{\omega}^{2}+\psi_{\Omega_2}(\Omega_{\omega}\times\Omega+\Omega_2))%
}

内函数参数的后项先增加 \allowbreak{}\NovelInline{I05637}{\Omega_2}\allowbreak{}，注意最外层参数并没有直接增加这一项。

\NovelDisplay{M01419R01}{%
(0)(1,1,1)(2,1)(1,1,1)(1,1)(2,2,1)(3,1)%
}
\NovelDisplay{M01419R02}{%
\quad (2,2)(3,2)=\psi(\Omega_{\omega}^{2}+\psi_{\Omega_2}(\Omega_{\omega}\times\Omega+\Omega_2^{2}))%
}

继续使内层后项形成 \allowbreak{}\NovelInline{I05638}{\Omega_2^2}\allowbreak{}，仍隔着 \allowbreak{}\NovelInline{I05639}{\psi_{\Omega_2}}\allowbreak{}；三个层次由括号清楚分开。

\NovelDisplay{M01420R01}{%
(0)(1,1,1)(2,1)(1,1,1)(1,1)(2,2,1)(3,1)(2,2)%
}
\NovelDisplay{M01420R02}{%
\quad (3,3,1)= \psi(\Omega^2_\omega+\psi_{\Omega_2}(\Omega_\omega\times\Omega+\psi_{\Omega_3}(\Omega_\omega)))%
}

最内侧出现 \allowbreak{}\NovelInline{I05640}{\psi_{\Omega_3}({\Omega_\omega})}\allowbreak{}，较高上界位于较深处，不能与外层函数合并。

\NovelDisplay{M01421R01}{%
(0)(1,1,1)(2,1)(1,1,1)(1,1)(2,2,1)(3,1)(2,2)(3,3,1)(4,1)= \psi(\Omega^2_\omega%
}
\NovelDisplay{M01421R02}{%
\quad +\psi_{\Omega_2}(\Omega_\omega\times\Omega+\psi_{\Omega_3}(\Omega_\omega\times\Omega)))%
}

这个最内函数参数又变为 \allowbreak{}\NovelInline{I05641}{{\Omega_\omega}\times\Omega}\allowbreak{}；两处相同的乘积，分别属于两层不同的参数。

\NovelDisplay{M01422R01}{%
(0)(1,1,1)(2,1)(1,1,1)(1,1)(2,2,1)(3,1)(2,2)(3,3,1)(4,1)(3,3)(4,3)%
}
\NovelDisplay{M01422R02}{%
\quad =\psi(\Omega_{\omega}^{2}+\psi_{\Omega_2}(\Omega_{\omega}\times\Omega+\Omega_3^{2}))%
}

内层后项达到 \allowbreak{}\NovelInline{I05642}{\Omega_3^2}\allowbreak{}，前面的 \allowbreak{}\NovelInline{I05643}{{\Omega_\omega}\times\Omega}\allowbreak{} 不动；这并未把最外层背景变成 \allowbreak{}\NovelInline{I05644}{{\Omega_\omega}^3}\allowbreak{}。

\NovelDisplay{M01423R01}{%
(0)(1,1,1)(2,1)(1,1,1)(1,1)(2,2,1)(3,1)%
}
\NovelDisplay{M01423R02}{%
\quad (2,2,1)=\psi(\Omega^2_\omega+\psi_{\Omega_2}(\Omega_\omega\times\Omega+\Omega_\omega))%
}

这一尾部形成 \allowbreak{}\NovelInline{I05645}{{\Omega_\omega}\times\Omega+{\Omega_\omega}}\allowbreak{}，仍全部包在 \allowbreak{}\NovelInline{I05646}{\psi_{\Omega_2}}\allowbreak{} 内，不能挪到最外层相加。

\NovelDisplay{M01424R01}{%
(0)(1,1,1)(2,1)(1,1,1)(1,1)(2,2,1)(3,1)%
}
\NovelDisplay{M01424R02}{%
\quad (2,2,1)(3,1)= \psi(\Omega^2_\omega+\psi_{\Omega_2}(\Omega_\omega\times\Omega\times2))%
}

内函数参数形成两份 \allowbreak{}\NovelInline{I05647}{{\Omega_\omega}\times\Omega}\allowbreak{}。右侧乘二作用于这一乘积，不属于共同背景 \allowbreak{}\NovelInline{I05648}{{\Omega_\omega}^2}\allowbreak{}。

\NovelDisplay{M01425R01}{%
(0)(1,1,1)(2,1)(1,1,1)(1,1)(2,2,1)(3,1)(3,1)=\psi(\Omega^2_\omega+\psi_{\Omega_2}(\Omega_\omega\times\Omega^2))%
}

内函数的参数是 \allowbreak{}\NovelInline{I05649}{{\Omega_\omega}\times\Omega^2}\allowbreak{}，底数 \allowbreak{}\NovelInline{I05650}{\Omega}\allowbreak{} 的平方位于右因子，不能看成 \allowbreak{}\NovelInline{I05651}{{\Omega_\omega}^2}\allowbreak{}。

\NovelDisplay{M01426R01}{%
(0)(1,1,1)(2,1)(1,1,1)(1,1)(2,2,1)(3,1)%
}
\NovelDisplay{M01426R02}{%
\quad (4,2,1)=\psi(\Omega_\omega^{2}+\psi_{\Omega_2}(\Omega_\omega\times\psi_{\Omega_2}(\Omega_\omega)))%
}

内部右因子进入 \allowbreak{}\NovelInline{I05652}{\psi_{\Omega_2}({\Omega_\omega})}\allowbreak{}，相同上界的内函数又嵌入一次；仍从最内层往外读。

\NovelDisplay{M01427R01}{%
(0)(1,1,1)(2,1)(1,1,1)(1,1)(2,2,1)(3,1)(4,2,1)(5,1)(6,2,1)= \psi(\Omega^2_\omega%
}
\NovelDisplay{M01427R02}{%
\quad +\psi_{\Omega_2}(\Omega_\omega\times\psi_{\Omega_2}(\Omega_\omega\times\psi_{\Omega_2}(\Omega_\omega))))%
}

这一式把上述内部结构再嵌一次，两个 \allowbreak{}\NovelInline{I05653}{{\Omega_\omega}}\allowbreak{} 的作用位置由各自右括号区分，不能删去中间函数。

\NovelDisplay{M01428R01}{%
(0)(1,1,1)(2,1)(1,1,1)(1,1)(2,2,1)(3,2)=\psi(\Omega_\omega^{2}+\psi_{\Omega_2}(\Omega_\omega\times\Omega_2))%
}

内函数参数的右因子达到 \allowbreak{}\NovelInline{I05654}{\Omega_2}\allowbreak{}，这轮 \allowbreak{}\NovelInline{I05655}{(3,1)}\allowbreak{} 的作用完成；最外层仍以 \allowbreak{}\NovelInline{I05656}{{\Omega_\omega}^2}\allowbreak{} 开头。

\allowbreak{}\NovelInline{I05657}{(3,1)}\allowbreak{} 对应的位置进到了 \allowbreak{}\NovelInline{I05658}{(3,2)}\allowbreak{}。内部乘积从 \allowbreak{}\NovelInline{I05659}{{\Omega_\omega}\times\Omega}\allowbreak{} 提升到 \allowbreak{}\NovelInline{I05660}{{\Omega_\omega}\times\Omega_2}\allowbreak{}，这轮确实发生了提升。

\NovelDisplay{M01429R01}{%
\text{内层：}\quad {\Omega_\omega}\times\Omega\rightsquigarrow {\Omega_\omega}\times\Omega_2,\qquad\text{外部共同前项：}\quad {\Omega_\omega}^2%
}

不过，内部的一轮提升并没有把最外层共同前项也改成立方。蔡须坤对照内外两处，气息才沿原路平下来；霍雨浩将这一处压力反应另记在推导旁。

\NovelDisplay{M01430R01}{%
(0)(1,1,1)(2,1)(1,1,1)(1,1)(2,2,1)(3,2)(2,2)%
}
\NovelDisplay{M01430R02}{%
\quad (3,2)=\psi(\Omega_{\omega}^{2}+\psi_{\Omega_2}(\Omega_{\omega}\times\Omega_2+\Omega_2^{2}))%
}

新内层背景之后再加 \allowbreak{}\NovelInline{I05661}{\Omega_2^2}\allowbreak{}；整体是 \allowbreak{}\NovelInline{I05662}{{\Omega_\omega}^2+\psi_{\Omega_2}(\cdots)}\allowbreak{}，不是两个外部平方的和。

\NovelDisplay{M01431R01}{%
(0)(1,1,1)(2,1)(1,1,1)(1,1)(2,2,1)(3,2)(2,2)(3,3,1)(4,2)= \psi(\Omega_\omega^{2}%
}
\NovelDisplay{M01431R02}{%
\quad +\psi_{\Omega_2}(\Omega_\omega\times\Omega_2+\psi_{\Omega_3}(\Omega_\omega\times\Omega_2)))%
}

内层后项进入 \allowbreak{}\NovelInline{I05663}{\psi_{\Omega_3}({\Omega_\omega}\times\Omega_2)}\allowbreak{}，对应列 \allowbreak{}\NovelInline{I05664}{(4,2)}\allowbreak{} 将承担下一轮替身。

新的替身是 \allowbreak{}\NovelInline{I05665}{(4,2)}\allowbreak{}。古月娜把它与前一章同一角色对照：同样要检查第二行祖先是否经过坏根，不能只因第一行加二便断定两行一起加。

\NovelDisplay{M01432R01}{%
(0)(1,1,1)(2,1)(1,1,1)(1,1)(2,2,1)(3,2)(2,2)(3,3,1)(4,2)(3,3)(4,2)%
}
\NovelDisplay{M01432R02}{%
\quad (5,3,1)(6,2)=\psi(\Omega_{\omega}^{2}+\psi_{\Omega_2}(\Omega_{\omega}\times\Omega_2+\Omega_3\times\psi_{\Omega_3}(\Omega_{\omega}\times\Omega_2)))%
}

内层后项是 \allowbreak{}\NovelInline{I05666}{\Omega_3}\allowbreak{} 乘一个坍缩函数，其参数也追到 \allowbreak{}\NovelInline{I05667}{{\Omega_\omega}\times\Omega_2}\allowbreak{}；两处相同背景分别归属。

\NovelDisplay{M01433R01}{%
(0)(1,1,1)(2,1)(1,1,1)(1,1)(2,2,1)(3,2)(2,2)(3,3,1)(4,2)(3,3)(4,3)%
}
\NovelDisplay{M01433R02}{%
\quad =\psi(\Omega_{\omega}^{2}+\psi_{\Omega_2}(\Omega_{\omega}\times\Omega_2+\Omega_3^{2}))%
}

在内函数参数中形成 \allowbreak{}\NovelInline{I05668}{\Omega_3^2}\allowbreak{}。平方仍是最内层和的后项，尚未替换它前面的右因子。

\NovelDisplay{M01434R01}{%
(0)(1,1,1)(2,1)(1,1,1)(1,1)(2,2,1)(3,2)(2,2)(3,3,1)(4,2)(3,3,1)%
}
\NovelDisplay{M01434R02}{%
\quad = \psi(\Omega_\omega^{2}+\psi_{\Omega_2}(\Omega_\omega\times(\Omega_2+1)))%
}

内函数参数可合写为 \allowbreak{}\NovelInline{I05669}{{\Omega_\omega}\times(\Omega_2+1)}\allowbreak{}，表示原有大段后再加一份 \allowbreak{}\NovelInline{I05670}{{\Omega_\omega}}\allowbreak{}。

\NovelDisplay{M01435R01}{%
(0)(1,1,1)(2,1)(1,1,1)(1,1)(2,2,1)(3,2)(2,2)(3,3,1)(4,2)(3,3,1)(3,3,1)%
}
\NovelDisplay{M01435R02}{%
\quad = \psi(\Omega^2_\omega+\psi_{\Omega_2}(\Omega_\omega\times\Omega_2+\Omega_\omega\times2))%
}

内部有限尾项由一到二，成为 \allowbreak{}\NovelInline{I05671}{{\Omega_\omega}\times\Omega_2+{\Omega_\omega}\times2}\allowbreak{}；二不能移进 \allowbreak{}\NovelInline{I05672}{\Omega_2}\allowbreak{} 的下标。

\NovelDisplay{M01436R01}{%
(0)(1,1,1)(2,1)(1,1,1)(1,1)(2,2,1)(3,2)(2,2)(3,3,1)(4,2)(3,3,1)(4,2)%
}
\NovelDisplay{M01436R02}{%
\quad = \psi(\Omega_\omega^{2}+\psi_{\Omega_2}(\Omega_\omega\times(\Omega_2\times2)))%
}

内部形成两份 \allowbreak{}\NovelInline{I05673}{{\Omega_\omega}\times\Omega_2}\allowbreak{}，右因子是 \allowbreak{}\NovelInline{I05674}{\Omega_2\times2}\allowbreak{}；最外层共同背景仍只出现一次。

\NovelDisplay{M01437R01}{%
(0)(1,1,1)(2,1)(1,1,1)(1,1)(2,2,1)(3,2)(2,2)%
}
\NovelDisplay{M01437R02}{%
\quad (3,3,1)(4,2)(4,2)= \psi(\Omega_\omega^{2}+\psi_{\Omega_2}(\Omega_\omega\times\Omega_2^{2}))%
}

内部右因子形成 \allowbreak{}\NovelInline{I05675}{\Omega_2^2}\allowbreak{}，与最外层 \allowbreak{}\NovelInline{I05676}{{\Omega_\omega}^2}\allowbreak{} 不同底数，必须分别读。

\NovelDisplay{M01438R01}{%
(0)(1,1,1)(2,1)(1,1,1)(1,1)(2,2,1)(3,2)(2,2)(3,3,1)(4,2)(5,3)=\psi(\Omega_{\omega}^{2}%
}
\NovelDisplay{M01438R02}{%
\quad +\psi_{\Omega_2}(\Omega_{\omega}\times\psi_{\Omega_3}(\Omega_3)))%
}

内部右因子进入 \allowbreak{}\NovelInline{I05677}{\psi_{\Omega_3}(\Omega_3)}\allowbreak{}，是第二轮追赶在更高上界上的展开。

\NovelDisplay{M01439R01}{%
(0)(1,1,1)(2,1)(1,1,1)(1,1)(2,2,1)(3,2)(2,2)(3,3,1)(4,2)(5,3,1)%
}
\NovelDisplay{M01439R02}{%
\quad = \psi(\Omega_\omega^{2}+\psi_{\Omega_2}(\Omega_\omega\times\psi_{\Omega_3}(\Omega_\omega)))%
}

该函数参数达到 \allowbreak{}\NovelInline{I05678}{{\Omega_\omega}}\allowbreak{}，外面的 \allowbreak{}\NovelInline{I05679}{\psi_{\Omega_2}}\allowbreak{} 仍完整保留；不能因为都是坍缩函数就压成一层。

\NovelDisplay{M01440R01}{%
(0)(1,1,1)(2,1)(1,1,1)(1,1)(2,2,1)(3,2)(2,2)%
}
\NovelDisplay{M01440R02}{%
\quad (3,3,1)(4,3)= \psi(\Omega_\omega^{2}+\psi_{\Omega_2}(\Omega_\omega\times\Omega_3))%
}

内函数参数达到 \allowbreak{}\NovelInline{I05680}{{\Omega_\omega}\times\Omega_3}\allowbreak{}，这一轮替身 \allowbreak{}\NovelInline{I05681}{(4,2)}\allowbreak{} 的作用完成，再往后按同样父项关系分析。

\allowbreak{}\NovelInline{I05682}{(4,2)}\allowbreak{} 这一轮结束后，内部乘积走到 \allowbreak{}\NovelInline{I05683}{{\Omega_\omega}\times\Omega_3}\allowbreak{}。继续对任意有限层作同样分析，内函数参数才会收束到 \allowbreak{}\NovelInline{I05684}{{\Omega_\omega}^2}\allowbreak{}。

\NovelDisplay{M01441R01}{%
(0)(1,1,1)(2,1)(1,1,1)(1,1)(2,2,1)(3,2)(2,2,1)=\psi(\Omega_\omega^{2}+\psi_{\Omega_2}(\Omega_\omega^{2}))%
}

内函数参数终于达到 \allowbreak{}\NovelInline{I05685}{{\Omega_\omega}^2}\allowbreak{}，与最外层共同背景同形；同形也仍隔着一层 \allowbreak{}\NovelInline{I05686}{\psi_{\Omega_2}}\allowbreak{}。

\NovelDisplay{M01442R01}{%
(0)(1,1,1)(2,1)(1,1,1)(1,1)(2,2,1)(3,2)(2,2,1)(2,2)(3,1)(4,2,1)(5,2)%
}
\NovelDisplay{M01442R02}{%
\quad (4,2,1)=\psi(\Omega_{\omega}^{2}+\Omega_2\times\psi_{\Omega_2}(\Omega_{\omega}^{2}))%
}

新外部后项为 \allowbreak{}\NovelInline{I05687}{\Omega_2\times\psi_{\Omega_2}({\Omega_\omega}^2)}\allowbreak{}。这次乘法已经在该内函数之外。

\NovelDisplay{M01443R01}{%
(0)(1,1,1)(2,1)(1,1,1)(1,1)(2,2,1)(3,2)(2,2,1)(2,2)(3,2)=\psi(\Omega_{\omega}^{2}+\Omega_2^{2})%
}

外部后项形成 \allowbreak{}\NovelInline{I05688}{\Omega_2^2}\allowbreak{}，对应的是前一步之后的新阶段；不能说最内层的平方自动跑到外面。

\NovelDisplay{M01444R01}{%
(0)(1,1,1)(2,1)(1,1,1)(1,1)(2,2,1)(3,2)%
}
\NovelDisplay{M01444R02}{%
\quad (2,2,1)(2,2)(3,3,1)= \psi(\Omega_\omega^{2}+\psi_{\Omega_3}(\Omega_\omega))%
}

外部后项进入 \allowbreak{}\NovelInline{I05689}{\psi_{\Omega_3}({\Omega_\omega})}\allowbreak{}，还要走过相应更高上界的尾部，才能收束为 \allowbreak{}\NovelInline{I05690}{{\Omega_\omega}}\allowbreak{}。

接下来的更高上界也按相应父项关系分析，直至最外层后项到达下面这一式。每个人都得把这一步重新运转出来，不能只抄邻人的答案。

\NovelDisplay{M01445R01}{%
(0)(1,1,1)(2,1)(1,1,1)(1,1,1)=\psi(\Omega_\omega^{2}+\Omega_\omega)%
}

新加的 \allowbreak{}\NovelInline{I05691}{(1,1,1)}\allowbreak{} 最终只在最外层参数后添一份 \allowbreak{}\NovelInline{I05692}{{\Omega_\omega}}\allowbreak{}。前面的 \allowbreak{}\NovelInline{I05693}{{\Omega_\omega}^2}\allowbreak{} 没有因此成为 \allowbreak{}\NovelInline{I05694}{{\Omega_\omega}^3}\allowbreak{}。

这次末列 \allowbreak{}\NovelInline{I05695}{(1,1,1)}\allowbreak{} 回到了``添一份 \allowbreak{}\NovelInline{I05696}{{\Omega_\omega}}\allowbreak{}''的作用；但不能说提升效应没有发生。

\NovelDisplay{M01446R01}{%
\text{内函数中完成：}\quad {\Omega_\omega}\times\Omega\rightsquigarrow {\Omega_\omega}\times\Omega_2\rightsquigarrow\cdots\rightsquigarrow {\Omega_\omega}^2%
}

\NovelDisplay{M01447R01}{%
\text{外层后项所经：}\quad\psi_{\Omega_2}({\Omega_\omega}^2),\quad\Omega_2^2,\quad\psi_{\Omega_3}({\Omega_\omega}),\quad\ldots,\quad {\Omega_\omega}%
}

提升发生在 \allowbreak{}\NovelInline{I05697}{\psi_{\Omega_2}}\allowbreak{} 及更深的结构里，先生成与外部背景同形的 \allowbreak{}\NovelInline{I05698}{{\Omega_\omega}^2}\allowbreak{}，再形成外层相应的 \allowbreak{}\NovelInline{I05699}{\Omega_2^2,\Omega_3^2}\allowbreak{} 等幂次结构。最终外部只表现为加一份 \allowbreak{}\NovelInline{I05700}{{\Omega_\omega}}\allowbreak{}，不能据此抹掉内部的工作。

\NovelDisplay{M01448R01}{%
\bigl((0)(1,1,1)(2,1)(1,1,1)\bigr)\,(1,1,1)=\psi({\Omega_\omega}^2+{\Omega_\omega}),\bigl((0)(1,1,1)(2,1)%
}
\NovelDisplay{M01448R02}{%
\quad (1,1,1)\bigr)\,(1,1,1)\ne\psi({\Omega_\omega}^3)%
}

霍雨浩划去旧答案。古月娜看过纸，又请卢本韦陪她在同一阶重试，两人确认受阻的先后与推导相符。

他们在阴影下休息了片刻，便接着分析外层后项。霍雨浩留出共同前项的位置，剩下的纸只用来标新变化。

\NovelDisplay{M01449R01}{%
(0)(1,1,1)(2,1)(1,1,1)(1,1,1)(1,1,1)=\psi(\Omega^2_\omega+\Omega_\omega\times2)%
}

再添一项，外部尾部成为 \allowbreak{}\NovelInline{I05701}{{\Omega_\omega}\times2}\allowbreak{}；共同前项为 \allowbreak{}\NovelInline{I05702}{{\Omega_\omega}^2}\allowbreak{}，两者的加法顺序不能倒过来。

\NovelDisplay{M01450R01}{%
(0)(1,1,1)(2,1)(1,1,1)(2)=\psi(\Omega_\omega^{2}+\Omega_\omega\times\omega)%
}

一行末列令有限尾项的重复走向 \allowbreak{}\NovelInline{I05703}{{\Omega_\omega}\times\omega}\allowbreak{}；这里不是把共同前项右乘 \allowbreak{}\NovelInline{I05704}{\omega}\allowbreak{}。

\NovelDisplay{M01451R01}{%
(0)(1,1,1)(2,1)(1,1,1)(2,1)=\psi(\Omega_\omega^{2}+\Omega_\omega\times\Omega)%
}

后项变为 \allowbreak{}\NovelInline{I05705}{{\Omega_\omega}\times\Omega}\allowbreak{}，于是重新出现``一个大段后还有一个较小大段''的背景。

右端又出现 \allowbreak{}\NovelInline{I05706}{(1,1,1)(2,1)}\allowbreak{} 的组合。前章已经见过它如何引出系数提升，现在仍要先让新内函数追到整个外部背景。

\NovelDisplay{M01452R01}{%
(0)(1,1,1)(2,1)(1,1,1)(2,1)(1,1)(2,2,1)(3,2)(2,2,1)(3,1)= \psi(\Omega_\omega^{2}+\Omega_\omega%
}
\NovelDisplay{M01452R02}{%
\quad \times\Omega+\psi_{\Omega_2}(\Omega_\omega^{2}+\Omega_\omega\times\Omega))%
}

新开的内函数参数追到整个外部背景 \allowbreak{}\NovelInline{I05707}{{\Omega_\omega}^2+{\Omega_\omega}\times\Omega}\allowbreak{}，对齐后才能分析下一轮提升。

\NovelDisplay{M01453R01}{%
(0)(1,1,1)(2,1)(1,1,1)(2,1)(1,1)(2,2,1)(3,2)(2,2,1)(3,1)(2,2,1)%
}
\NovelDisplay{M01453R02}{%
\quad = \psi(\Omega_\omega^{2}+\Omega_\omega\times(\Omega+1))%
}

在同一共同前项后又添一份 \allowbreak{}\NovelInline{I05708}{{\Omega_\omega}}\allowbreak{}，后部可合写为 \allowbreak{}\NovelInline{I05709}{{\Omega_\omega}\times(\Omega+1)}\allowbreak{}。

\NovelDisplay{M01454R01}{%
(0)(1,1,1)(2,1)(1,1,1)(2,1)(1,1)(2,2,1)%
}
\NovelDisplay{M01454R02}{%
\quad (3,2)(2,2,1)(3,2)= \psi(\Omega_\omega^{2}+\Omega_\omega\times\Omega_2)%
}

后部右因子达到 \allowbreak{}\NovelInline{I05710}{\Omega_2}\allowbreak{}，所以全式仍是 \allowbreak{}\NovelInline{I05711}{\psi({\Omega_\omega}^2+{\Omega_\omega}\times\Omega_2)}\allowbreak{}，不是三个 \allowbreak{}\NovelInline{I05712}{{\Omega_\omega}}\allowbreak{} 相乘。

\NovelDisplay{M01455R01}{%
(0)(1,1,1)(2,1)(1,1,1)(2,1)(1,1)(2,2,1)(3,2)(2,2,1)(3,2)(2,2)(3,3,1)%
}
\NovelDisplay{M01455R02}{%
\quad (4,3)(3,3,1)(4,3)= \psi(\Omega_\omega^{2}+\Omega_\omega\times\Omega_3)%
}

后部右因子进到 \allowbreak{}\NovelInline{I05713}{\Omega_3}\allowbreak{}，与前章的系数提升对应；把这一构造延伸到任意有限上界。

\NovelDisplay{M01456R01}{%
(0)(1,1,1)(2,1)(1,1,1)(2,1)(1,1,1)=\psi(\Omega_\omega^{2}\times2)%
}

此时才形成第二份完整的 \allowbreak{}\NovelInline{I05714}{{\Omega_\omega}^2}\allowbreak{}，写作 \allowbreak{}\NovelInline{I05715}{{\Omega_\omega}^2\times2}\allowbreak{}。两个平方相加与一个立方不同。

\NovelDisplay{M01457R01}{%
(0)(1,1,1)(2,1)(2)=\psi(\Omega_\omega^{2}\times\omega)%
}

这条基本列以有限份 \allowbreak{}\NovelInline{I05716}{{\Omega_\omega}^2}\allowbreak{} 为对应背景，结果为 \allowbreak{}\NovelInline{I05717}{{\Omega_\omega}^2\times\omega}\allowbreak{}；无需再接一个 \allowbreak{}\NovelInline{I05718}{(1,1,1)}\allowbreak{}。

这里不用为了凑一个整段，再额外续接 \allowbreak{}\NovelInline{I05719}{(1,1,1)}\allowbreak{}。教材给出的有限重复结构已有这条基本列的极限；应核对整段，而非按末列字形套上一章的答案。

\NovelDisplay{M01458R01}{%
(0)(1,1,1)(2,1),(0)(1,1,1)(2,1)(1,1,1)(2,1),\dots%
}

\NovelDisplay{M01459R01}{%
\text{对应基本列收束：}\quad\psi({\Omega_\omega}^2\times\omega)%
}

卢本韦圈住自己多抄的那列。``留着，下次也对照着试。''

顺着这一结构再向上，先辨认共同前项后的右乘部分。

\NovelDisplay{M01460R01}{%
(0)(1,1,1)(2,1)(2,1)=\psi(\Omega_\omega^{2}\times\Omega)%
}

共同前项右乘部分达到 \allowbreak{}\NovelInline{I05720}{\Omega}\allowbreak{}，为 \allowbreak{}\NovelInline{I05721}{{\Omega_\omega}^2\times\Omega}\allowbreak{}；这仍在立方之前的当前结构中。

右乘部分已经到达当前的参照层，这时再看新接第三行的一会引出什么。

\NovelDisplay{M01461R01}{%
(0)(1,1,1)(2,1)(2,1)(1,1,1)=\psi(\Omega_\omega^{3})%
}

这次末列 \allowbreak{}\NovelInline{I05722}{(1,1,1)}\allowbreak{} 令参数中的 \allowbreak{}\NovelInline{I05723}{{\Omega_\omega}^2\times\Omega}\allowbreak{} 经提升达到 \allowbreak{}\NovelInline{I05724}{{\Omega_\omega}^3}\allowbreak{}，作用取决于前面的背景。

这一次，\allowbreak{}\NovelInline{I05725}{(1,1,1)}\allowbreak{} 才使参数的幂次再增一。同样一个末列，在 \allowbreak{}\NovelInline{I05726}{{\Omega_\omega}^2}\allowbreak{} 之后只加 \allowbreak{}\NovelInline{I05727}{{\Omega_\omega}}\allowbreak{}，在 \allowbreak{}\NovelInline{I05728}{{\Omega_\omega}^2\times\Omega}\allowbreak{} 之后却能引起到 \allowbreak{}\NovelInline{I05729}{{\Omega_\omega}^3}\allowbreak{} 的提升，背景不能省。

\NovelDisplay{M01462R01}{%
(0)(1,1,1)(2,1)(2,1)(1,1,1)(2,1)(2,1)(1,1,1)=\psi(\Omega_\omega^{3}\times2)%
}

把对应的大段再接一份，形成 \allowbreak{}\NovelInline{I05730}{{\Omega_\omega}^3\times2}\allowbreak{}；与上一式比较，是立方后的有限重复。

\NovelDisplay{M01463R01}{%
(0)(1,1,1)(2,1)(2,1)(2,1)(1,1,1)=\psi(\Omega_\omega^{4})%
}

另一种续接把参数中的 \allowbreak{}\NovelInline{I05731}{{\Omega_\omega}}\allowbreak{} 因子增加到四个，故为 \allowbreak{}\NovelInline{I05732}{{\Omega_\omega}^4}\allowbreak{}。重复段与提高幂次要分别判定。

\NovelDisplay{M01464R01}{%
(0)(1,1,1)(2,1)(3)=\psi(\Omega_\omega^{\omega})%
}

有限幂在这条基本列中收束到 \allowbreak{}\NovelInline{I05733}{{\Omega_\omega}^\omega}\allowbreak{}；底数 \allowbreak{}\NovelInline{I05734}{{\Omega_\omega}}\allowbreak{} 固定，指数首先到可数极限。

\NovelDisplay{M01465R01}{%
(0)(1,1,1)(2,1)(3,1)=\psi(\Omega_\omega^{\Omega})%
}

指数进入 \allowbreak{}\NovelInline{I05735}{\Omega}\allowbreak{}，得到 \allowbreak{}\NovelInline{I05736}{{\Omega_\omega}^\Omega}\allowbreak{}；此后新的提升效应会进入指数层。

九十米处，蔡须坤忽然停住。肩背承受的重量还在，魂力里受阻的位置却变了；她试着加大外部右因子，胸口越发闷。古月娜在同阶复查，霍雨浩顺着两人的记录，把变化标进指数。

霍雨浩把纸举高：``下面提升的，开始是指数。''

\NovelDisplay{M01466R01}{%
\psi\bigl({\Omega_\omega}\times\boxed{\alpha}\bigr)\quad\text{与}\quad\psi\bigl({\Omega_\omega}^{\boxed{\alpha}}\bigr)%
}

卢本韦扶蔡须坤退回半步，古月娜沿括号往回指。一个是右因子，一个是幂指数；下面每一步，都要认清当前在哪一种结构里。

\NovelDisplay{M01467R01}{%
(0)(1,1,1)(2,1)(3,1)(1,1)(2,2,1)(3,1)=\psi(\Omega_\omega^{\Omega}+\psi_{\Omega_2}(\Omega_\omega\times\Omega))%
}

先在 \allowbreak{}\NovelInline{I05737}{{\Omega_\omega}^\Omega}\allowbreak{} 后开出 \allowbreak{}\NovelInline{I05738}{\psi_{\Omega_2}({\Omega_\omega}\times\Omega)}\allowbreak{}，共同前项还未发生指数提升。

\NovelDisplay{M01468R01}{%
(0)(1,1,1)(2,1)(3,1)(1,1)(2,2,1)(3,2)(2,2,1)=\psi(\Omega_\omega^{\Omega}+\psi_{\Omega_2}(\Omega_\omega^2))%
}

内函数参数达到 \allowbreak{}\NovelInline{I05739}{{\Omega_\omega}^2}\allowbreak{}，所以对应列必须已经是 \allowbreak{}\NovelInline{I05740}{(3,2)}\allowbreak{}；不能用尚未完成上一轮的 \allowbreak{}\NovelInline{I05741}{(3,1)}\allowbreak{} 替代。

这里倒数第二列是 \allowbreak{}\NovelInline{I05742}{(3,2)}\allowbreak{}，不能抄回 \allowbreak{}\NovelInline{I05743}{(3,1)}\allowbreak{}。内部参数达到 \allowbreak{}\NovelInline{I05744}{{\Omega_\omega}^2}\allowbreak{} 的过程，已经用完了前一轮替身。

\NovelDisplay{M01469R01}{%
\psi_{\Omega_2}({\Omega_\omega}\times\Omega)\quad\text{对应尚在追赶的内层，}%
}
\NovelDisplay{M01469R02}{%
\quad \psi_{\Omega_2}({\Omega_\omega}^2)\quad\text{对应完成该轮后的内层。}%
}

蔡须坤改正那一笔，在压力下重新运转；卢本韦等她站稳，才松手。霍雨浩将这次结果与前两个月的同类片段并排。

\NovelDisplay{M01470R01}{%
(0)(1,1,1)(2,1)(3,1)(1,1)(2,2,1)(3,2)(3,2)(2,2,1)= \psi(\Omega_\omega^{\Omega}+\psi_{\Omega_2}(\Omega_\omega^3))%
}

内函数参数的幂次再增一，达到 \allowbreak{}\NovelInline{I05745}{{\Omega_\omega}^3}\allowbreak{}，最外层共同背景仍是 \allowbreak{}\NovelInline{I05746}{{\Omega_\omega}^\Omega}\allowbreak{}。

\NovelDisplay{M01471R01}{%
(0)(1,1,1)(2,1)(3,1)(1,1)(2,2,1)(3,2)(4)=\psi(\Omega_\omega^{\Omega}+\psi_{\Omega_2}(\Omega_\omega^\omega))%
}

内函数参数达到 \allowbreak{}\NovelInline{I05747}{{\Omega_\omega}^\omega}\allowbreak{}，外面的 \allowbreak{}\NovelInline{I05748}{\psi_{\Omega_2}}\allowbreak{} 还没有撤去，括号不能省。

\NovelDisplay{M01472R01}{%
(0)(1,1,1)(2,1)(3,1)(1,1)(2,2,1)(3,2)(4,1)=\psi(\Omega_\omega^{\Omega}+\psi_{\Omega_2}(\Omega_\omega^{\Omega}))%
}

内函数参数追到 \allowbreak{}\NovelInline{I05749}{{\Omega_\omega}^\Omega}\allowbreak{}，这时新替身 \allowbreak{}\NovelInline{I05750}{(4,1)}\allowbreak{} 出现，接下来须处理指数所在层的阶差。

接着出现 \allowbreak{}\NovelInline{I05751}{(4,1)}\allowbreak{}。它限制的已是新一层关系：内侧结构必须再追赶，才能使主项的指数由 \allowbreak{}\NovelInline{I05752}{\Omega}\allowbreak{} 进入更高上界。

\NovelDisplay{M01473R01}{%
(0)(1,1,1)(2,1)(3,1)(1,1)(2,2,1)(3,2)(4,1)(2,2)(3,2)=\psi(\Omega_{\omega}^{\Omega}+\Omega_2^{2})%
}

外部后项为 \allowbreak{}\NovelInline{I05753}{\Omega_2^2}\allowbreak{}，这是新一轮的尾部，并非把前面的 \allowbreak{}\NovelInline{I05754}{{\Omega_\omega}^\Omega}\allowbreak{} 平方。

\NovelDisplay{M01474R01}{%
(0)(1,1,1)(2,1)(3,1)(1,1)(2,2,1)(3,2)(4,1)(2,2,1)= \psi(\Omega_\omega^{\Omega}+\Omega_\omega)%
}

外部后项到 \allowbreak{}\NovelInline{I05755}{{\Omega_\omega}}\allowbreak{}，最初替身仍在；这一步只是在主项后面添加一份 \allowbreak{}\NovelInline{I05756}{{\Omega_\omega}}\allowbreak{}。

\NovelDisplay{M01475R01}{%
(0)(1,1,1)(2,1)(3,1)(1,1)(2,2,1)(3,2)%
}
\NovelDisplay{M01475R02}{%
\quad (4,1)(2,2,1)(2,2,1)= \psi(\Omega_\omega^{\Omega}+\Omega_\omega\times2)%
}

尾部有两份 \allowbreak{}\NovelInline{I05757}{{\Omega_\omega}}\allowbreak{}，即 \allowbreak{}\NovelInline{I05758}{{\Omega_\omega}\times2}\allowbreak{}；不能把它并入主项的指数写成 \allowbreak{}\NovelInline{I05759}{{\Omega_\omega}^{\Omega+2}}\allowbreak{}。

\NovelDisplay{M01476R01}{%
(0)(1,1,1)(2,1)(3,1)(1,1)(2,2,1)(3,2)%
}
\NovelDisplay{M01476R02}{%
\quad (4,1)(2,2,1)(3,1)= \psi(\Omega_\omega^{\Omega}+\Omega_\omega\times\Omega)%
}

尾部达到 \allowbreak{}\NovelInline{I05760}{{\Omega_\omega}\times\Omega}\allowbreak{}，与共同前项 \allowbreak{}\NovelInline{I05761}{{\Omega_\omega}^\Omega}\allowbreak{} 之间仍是加法。

\NovelDisplay{M01477R01}{%
(0)(1,1,1)(2,1)(3,1)(1,1)(2,2,1)(3,2)(4,1)%
}
\NovelDisplay{M01477R02}{%
\quad (2,2,1)(3,2)(2,2,1)= \psi(\Omega_\omega^{\Omega}+\Omega_\omega^{2})%
}

尾部形成 \allowbreak{}\NovelInline{I05762}{{\Omega_\omega}^2}\allowbreak{}。这段先重走较低幂次，正是新提升不能凭外形一步跳过的原因。

\NovelDisplay{M01478R01}{%
(0)(1,1,1)(2,1)(3,1)(1,1)(2,2,1)(3,2)%
}
\NovelDisplay{M01478R02}{%
\quad (4,1)(2,2,1)(3,2)(4)= \psi(\Omega_\omega^{\Omega}+\Omega_\omega^{\omega})%
}

尾部达到 \allowbreak{}\NovelInline{I05763}{{\Omega_\omega}^\omega}\allowbreak{}，仍未形成第二份 \allowbreak{}\NovelInline{I05764}{{\Omega_\omega}^\Omega}\allowbreak{}；指数的差别要保留。

\NovelDisplay{M01479R01}{%
(0)(1,1,1)(2,1)(3,1)(1,1)(2,2,1)(3,2)%
}
\NovelDisplay{M01479R02}{%
\quad (4,1)(2,2,1)(3,2)(4,1)= \psi(\Omega_\omega^{\Omega}\times2)%
}

第二份 \allowbreak{}\NovelInline{I05765}{{\Omega_\omega}^\Omega}\allowbreak{} 形成，故可合成 \allowbreak{}\NovelInline{I05766}{{\Omega_\omega}^\Omega\times2}\allowbreak{}，这是系数的变化。

\NovelDisplay{M01480R01}{%
(0)(1,1,1)(2,1)(3,1)(1,1)(2,2,1)(3,2)(4,1)(3)= \psi(\Omega_\omega^{\Omega}\times\omega)%
}

有限份主项的重复走向右乘 \allowbreak{}\NovelInline{I05767}{\omega}\allowbreak{}，而非在指数里加 \allowbreak{}\NovelInline{I05768}{\omega}\allowbreak{}；后者是另一条结构。

\NovelDisplay{M01481R01}{%
(0)(1,1,1)(2,1)(3,1)(1,1)(2,2,1)(3,2)(4,1)(3,1)= \psi(\Omega_\omega^{\Omega}\times\Omega)%
}

主项右乘部分到 \allowbreak{}\NovelInline{I05769}{\Omega}\allowbreak{}，即 \allowbreak{}\NovelInline{I05770}{{\Omega_\omega}^\Omega\times\Omega}\allowbreak{}；此时尚未到 \allowbreak{}\NovelInline{I05771}{{\Omega_\omega}^{\Omega+1}}\allowbreak{}。

\NovelDisplay{M01482R01}{%
(0)(1,1,1)(2,1)(3,1)(1,1)(2,2,1)(3,2)(4,1)(3,2)(2,2,1)= \psi(\Omega_\omega^{\Omega+1})%
}

本式才在指数上加一，成为 \allowbreak{}\NovelInline{I05772}{{\Omega_\omega}^{\Omega+1}}\allowbreak{}；与上一式对照，它右乘的是底数 \allowbreak{}\NovelInline{I05773}{{\Omega_\omega}}\allowbreak{} 对应的整段。

\NovelDisplay{M01483R01}{%
(0)(1,1,1)(2,1)(3,1)(1,1)(2,2,1)(3,2)(4,1)(3,2)(4,1)= \psi(\Omega_\omega^{\Omega\times2})%
}

指数形成两份 \allowbreak{}\NovelInline{I05774}{\Omega}\allowbreak{}，即 \allowbreak{}\NovelInline{I05775}{{\Omega_\omega}^{\Omega\times2}}\allowbreak{}；不能读成 \allowbreak{}\NovelInline{I05776}{{\Omega_\omega}^\Omega\times2}\allowbreak{}。

\NovelDisplay{M01484R01}{%
(0)(1,1,1)(2,1)(3,1)(1,1)(2,2,1)(3,2)(4,1)(4,1)= \psi(\Omega_\omega^{\Omega^{2}})%
}

指数进入 \allowbreak{}\NovelInline{I05777}{\Omega^2}\allowbreak{}，指数里的幂与最外层底数的幂要逐层分开。

\NovelDisplay{M01485R01}{%
(0)(1,1,1)(2,1)(3,1)(1,1)(2,2,1)(3,2)(4,1)(5,1)= \psi(\Omega_\omega^{\Omega^\Omega})%
}

指数变成 \allowbreak{}\NovelInline{I05778}{\Omega^\Omega}\allowbreak{}，最外层仍是以 \allowbreak{}\NovelInline{I05779}{{\Omega_\omega}}\allowbreak{} 为底，幂按右结合读取。

\NovelDisplay{M01486R01}{%
(0)(1,1,1)(2,1)(3,1)(1,1)(2,2,1)(3,2)(4,1)(5,2)=\psi(\Omega_{\omega}^{\psi_{\Omega_2}(\Omega_2)})%
}

指数进入 \allowbreak{}\NovelInline{I05780}{\psi_{\Omega_2}(\Omega_2)}\allowbreak{}，坍缩函数这次位于幂指数里，不是外部相加的尾项。

\NovelDisplay{M01487R01}{%
(0)(1,1,1)(2,1)(3,1)(1,1)(2,2,1)(3,2)(4,1)(5,2,1)= \psi(\Omega_\omega^{\psi_{\Omega_2}(\Omega_\omega)})%
}

指数里的内函数参数达到 \allowbreak{}\NovelInline{I05781}{{\Omega_\omega}}\allowbreak{}，与前面右因子追赶的套路相似，但作用位置已经改变。

\NovelDisplay{M01488R01}{%
(0)(1,1,1)(2,1)(3,1)(1,1)(2,2,1)(3,2)(4,1)%
}
\NovelDisplay{M01488R02}{%
\quad (5,2,1)(6,2)(5,2,1)= \psi(\Omega_\omega^{\psi_{\Omega_2}(\Omega_\omega^2)})%
}

指数里的内函数参数达到 \allowbreak{}\NovelInline{I05782}{{\Omega_\omega}^2}\allowbreak{}；平方没有落在整个主项上，须保留最外层指数括号。

\NovelDisplay{M01489R01}{%
(0)(1,1,1)(2,1)(3,1)(1,1)(2,2,1)(3,2)(4,1)(5,2,1)(6,2)(7,1)(8,2,1)%
}
\NovelDisplay{M01489R02}{%
\quad = \psi(\Omega_\omega^{\psi_{\Omega_2}(\Omega_\omega^{\psi_{\Omega_2}(\Omega_\omega)})})%
}

指数函数参数又出现以 \allowbreak{}\NovelInline{I05783}{{\Omega_\omega}}\allowbreak{} 为底的幂，其指数仍是内函数。霍雨浩从最深处逐层添右括号。

\NovelDisplay{M01490R01}{%
(0)(1,1,1)(2,1)(3,1)(1,1)(2,2,1)(3,2)(4,2)=\psi(\Omega_\omega^{\Omega_2})%
}

新替身对应位置进到 \allowbreak{}\NovelInline{I05784}{(4,2)}\allowbreak{}，最外层主项的指数达到 \allowbreak{}\NovelInline{I05785}{\Omega_2}\allowbreak{}；不是前面那种系数变为 \allowbreak{}\NovelInline{I05786}{\Omega_2}\allowbreak{}。

\allowbreak{}\NovelInline{I05787}{(4,1)}\allowbreak{} 对应的位置进到 \allowbreak{}\NovelInline{I05788}{(4,2)}\allowbreak{}，这一轮结束。现在写在指数上的确实是 \allowbreak{}\NovelInline{I05789}{\Omega_2}\allowbreak{}，不是把外面的函数下标改成 \allowbreak{}\NovelInline{I05790}{\Omega_2}\allowbreak{}。

\NovelDisplay{M01491R01}{%
\psi({\Omega_\omega}^{\Omega_2})\quad\text{与}\quad\psi_{\Omega_2}({\Omega_\omega}^\Omega)\quad\text{分别辨认}%
}

\NovelDisplay{M01492R01}{%
(0)(1,1,1)(2,1)(3,1)(1,1)(2,2,1)(3,2)%
}
\NovelDisplay{M01492R02}{%
\quad (4,2)(2,2)(3,3,1)(4,3)(5,3)= \psi(\Omega_\omega^{\Omega_3})%
}

指数再进入 \allowbreak{}\NovelInline{I05791}{\Omega_3}\allowbreak{}，中间仍须按相应有限层的基本列补齐，长式展示了这一轮对应片段。

\NovelDisplay{M01493R01}{%
(0)(1,1,1)(2,1)(3,1)(1,1,1)=\psi(\Omega_\omega^{\Omega_\omega})%
}

指数提升到 \allowbreak{}\NovelInline{I05792}{{\Omega_\omega}}\allowbreak{}，主项成为 \allowbreak{}\NovelInline{I05793}{{\Omega_\omega}^{\Omega_\omega}}\allowbreak{}。它是双层同底数幂，尚非无穷高的塔。

这次是指数层的提升：\allowbreak{}\NovelInline{I05794}{\Omega\rightsquigarrow\Omega_2\rightsquigarrow\Omega_3\rightsquigarrow\cdots}\allowbreak{}。相应基本列收束后，主项指数达到 \allowbreak{}\NovelInline{I05795}{{\Omega_\omega}}\allowbreak{}，后面还会在更深指数里重复类似现象。

\NovelDisplay{M01494R01}{%
(0)(1,1,1)(2,1)(3,1)(1,1,1)(2,1)(1,1,1)=\psi(\Omega_\omega^{\Omega_\omega}+\Omega_\omega^2)%
}

在 \allowbreak{}\NovelInline{I05796}{{\Omega_\omega}^{\Omega_\omega}}\allowbreak{} 后增加 \allowbreak{}\NovelInline{I05797}{{\Omega_\omega}^2}\allowbreak{}，较小幂次位于和的后面，不能合写成主项指数加二。

\NovelDisplay{M01495R01}{%
(0)(1,1,1)(2,1)(3,1)(1,1,1)(2,1)(3,1)(1,1,1)=\psi(\Omega_\omega^{\Omega_\omega}\times2)%
}

两份完整的 \allowbreak{}\NovelInline{I05798}{{\Omega_\omega}^{\Omega_\omega}}\allowbreak{} 才合成 \allowbreak{}\NovelInline{I05799}{{\Omega_\omega}^{\Omega_\omega}\times2}\allowbreak{}。这依然是外部有限系数。

\NovelDisplay{M01496R01}{%
(0)(1,1,1)(2,1)(3,1)(2,1)=\psi(\Omega_\omega^{\Omega_\omega}\times\Omega)%
}

主项右乘 \allowbreak{}\NovelInline{I05800}{\Omega}\allowbreak{}，仍没有在指数中增加 \allowbreak{}\NovelInline{I05801}{\Omega}\allowbreak{}，要与下面两式分别比较。

\NovelDisplay{M01497R01}{%
(0)(1,1,1)(2,1)(3,1)(2,1)(1,1,1)=\psi(\Omega_\omega^{\Omega_\omega+1})%
}

主项指数增加一，为 \allowbreak{}\NovelInline{I05802}{{\Omega_\omega}^{{\Omega_\omega}+1}}\allowbreak{}；与 \allowbreak{}\NovelInline{I05803}{{\Omega_\omega}^{\Omega_\omega}\times\Omega}\allowbreak{} 是不同的结构。

\NovelDisplay{M01498R01}{%
(0)(1,1,1)(2,1)(3,1)(2,1)(3,1)=\psi(\Omega_\omega^{\Omega_\omega+\Omega})%
}

主项指数增加 \allowbreak{}\NovelInline{I05804}{\Omega}\allowbreak{}，为 \allowbreak{}\NovelInline{I05805}{{\Omega_\omega}^{{\Omega_\omega}+\Omega}}\allowbreak{}，这里的加号已在指数括号内部。

\NovelDisplay{M01499R01}{%
(0)(1,1,1)(2,1)(3,1)(2,1)(3,1)(1,1,1)=\psi(\Omega_\omega^{\Omega_\omega\times2})%
}

指数增加到第二份 \allowbreak{}\NovelInline{I05806}{{\Omega_\omega}}\allowbreak{}，成为 \allowbreak{}\NovelInline{I05807}{{\Omega_\omega}\times2}\allowbreak{}，所以主项是 \allowbreak{}\NovelInline{I05808}{{\Omega_\omega}^{{\Omega_\omega}\times2}}\allowbreak{}。

\NovelDisplay{M01500R01}{%
(0)(1,1,1)(2,1)(3,1)(3,1)=\psi(\Omega_\omega^{\Omega_\omega\times\Omega})%
}

指数中的右因子进到 \allowbreak{}\NovelInline{I05809}{\Omega}\allowbreak{}，成为 \allowbreak{}\NovelInline{I05810}{{\Omega_\omega}\times\Omega}\allowbreak{}；外面的底数仍为 \allowbreak{}\NovelInline{I05811}{{\Omega_\omega}}\allowbreak{}。

\NovelDisplay{M01501R01}{%
(0)(1,1,1)(2,1)(3,1)(3,1)(1,1,1)=\psi(\Omega_\omega^{\Omega_\omega^{2}})%
}

指数形成 \allowbreak{}\NovelInline{I05812}{{\Omega_\omega}^2}\allowbreak{}，整个主项是 \allowbreak{}\NovelInline{I05813}{{\Omega_\omega}^{{\Omega_\omega}^2}}\allowbreak{}；三处幂次要从最里侧向外辨认。

\NovelDisplay{M01502R01}{%
(0)(1,1,1)(2,1)(3,1)(4,1)=\psi(\Omega_\omega^{\Omega_\omega^{\Omega}})%
}

指数再形成 \allowbreak{}\NovelInline{I05814}{{\Omega_\omega}^\Omega}\allowbreak{}，整式中的幂按 \allowbreak{}\NovelInline{I05815}{{\Omega_\omega}^{({\Omega_\omega}^\Omega)}}\allowbreak{} 读取，不能按左结合算。

\NovelDisplay{M01503R01}{%
(0)(1,1,1)(2,1)(3,1)(4,1)(1,1,1)=\psi(\Omega_\omega^{\Omega_\omega^{\Omega_\omega}})%
}

指数成为 \allowbreak{}\NovelInline{I05816}{{\Omega_\omega}^{\Omega_\omega}}\allowbreak{}，得到三层塔 \allowbreak{}\NovelInline{I05817}{{\Omega_\omega}^{{\Omega_\omega}^{\Omega_\omega}}}\allowbreak{}；任意固定的三层、四层都还不是塔高无界的极限。

把每一个有限塔高都写完，不是把最后的省略号换成一个很大的整数。对于这套基本列，塔高没有有限上界；幂按右结合，层层送进更内的指数。

\NovelDisplay{M01504R01}{%
t_1={\Omega_\omega},\qquad t_2={\Omega_\omega}^{\Omega_\omega},\qquad t_3={\Omega_\omega}^{({\Omega_\omega}^{\Omega_\omega})},%
}
\NovelDisplay{M01504R02}{%
\quad t_{m+1}={\Omega_\omega}^{t_m}\quad(m\ge1)%
}

下面沿用教材对塔极限及 EBO 范围扩展记号的对应。等号中的无穷塔与无穷列是极限示意，不作为普通有限项参与逐列运算。

\NovelDisplay{M01505R01}{%
(0)(1,1,1)(2,1)(3,2)= (0)(1,1,1)(2,1)(3,1)(4,1)(5,1)\dots= \psi(\Omega_\omega^{\Omega_\omega^{\Omega_\omega^{\cdot^{\cdot^{\cdot}}}}})%
}
\NovelDisplay{M01505R02}{%
\quad = \psi(\varepsilon_{\Omega_\omega+1})= \psi(\Omega_{\omega+1})= \mathrm{TFBO}%
}

长链把有限塔高的共同构造收束为 TFBO；所写无穷塔与无穷列都是极限示意，最后两种参数写法沿用教材的扩展约定。

教材将 \allowbreak{}\NovelInline{I05818}{\mathrm{TFBO}}\allowbreak{} 与 \allowbreak{}\NovelInline{I05819}{\mathrm{BO}}\allowbreak{} 放在一起，对照 \allowbreak{}\NovelInline{I05820}{\mathrm{SCG}}\allowbreak{} 函数增长率研究中的上、下界；这里保留这一背景说明，其精确界值依赖相应函数与理论的约定，不在这道山路题里另作证明。

\NovelDisplay{M01506R01}{%
\psi(\Omega_{\omega+1})=\psi(\varepsilon_{\Omega_\omega+1})=\mathrm{TFBO}%
}

这一行与上式是同一教材约定的两种参数写法，沿用表格第703行的 EBO 扩展表达。

两者写法相似，不能由字形推断结构接近。眼前的塔已经把刚才那些系数提升送进一层又一层的指数；只看两个简称，便会漏掉全部来路。蔡须坤把最初的加号页与这一页并排，才看清这一段山道里，自己补全了多少旧时只会调动的力量。

在横台复盘时，众人先核对前章留下的夹页题与两处来源差异，再回到当前这一段。

霍雨浩又翻出教材夹页。第一条校核题采用原表第514行的数值，参数、乘数都按这条给定对应保留。

\NovelDisplay{M01507R01}{%
(0)(1,1,1)(2,1)(1,1)(2,2,1)(3,1)(4,2,1)(4,2)(4,1)(5,2,1)(5,2)= \psi(\Omega_\omega%
}
\NovelDisplay{M01507R02}{%
\quad \times\psi_{\Omega_2}(\Omega_\omega+\Omega_2+\psi_{\Omega_2}(\Omega_\omega+\Omega_2)))%
}

\NovelDisplay{M01508R01}{%
\psi\bigl(\Omega_\omega\times\psi_{\Omega_2}(\Omega_\omega+\Omega_2+\psi_{\Omega_2}(\Omega_\omega+\Omega_2))\bigr)%
}

薛峰从内向外读：先把 \allowbreak{}\NovelInline{I05821}{\Omega_\omega+\Omega_2}\allowbreak{} 放进 \allowbreak{}\NovelInline{I05822}{\psi_{\Omega_2}}\allowbreak{}，再将函数值加到 \allowbreak{}\NovelInline{I05823}{\Omega_\omega+\Omega_2}\allowbreak{} 后，把这个和放进外一层的 \allowbreak{}\NovelInline{I05824}{\psi_{\Omega_2}}\allowbreak{}，最后左乘 \allowbreak{}\NovelInline{I05825}{\Omega_\omega}\allowbreak{}，才进入最外层 \allowbreak{}\NovelInline{I05826}{\psi}\allowbreak{}。这条题的括号与前面几条不同，不能照眼熟的格式补出 \allowbreak{}\NovelInline{I05827}{\Omega_\omega\times\Omega}\allowbreak{}。

夹页另有两处尚待核定的来源差异。古月娜把它们单独圈了出来：它们在正文里暂取的表达式可以用来对照，但另一候选没有因此被证明相等。

\NovelDisplay{M01509R01}{%
\text{第一条原候选：}\psi(\Omega_\omega\times\Omega+\Omega_\omega+\psi_{\Omega_3}(\Omega_\omega\times\Omega+\Omega_\omega))%
}
\NovelDisplay{M01509R02}{%
\text{当前保留：}\psi(\Omega_\omega\times\Omega+\Omega_\omega+\psi_{\Omega_3}(\Omega_\omega\times\Omega+\psi_{\Omega_4}(\Omega_\omega)))%
}

\NovelDisplay{M01510R01}{%
\text{第二条原候选：}\psi(\Omega_\omega\times\Omega+\Omega_\omega+\Omega_3^2)%
}
\NovelDisplay{M01510R02}{%
\text{当前保留：}\psi(\Omega_\omega\times\Omega+\Omega_\omega+\psi_{\Omega_3}(\Omega_\omega\times\Omega+\Omega_\omega))%
}

``这两个差异，表格没有给出裁决。''霍雨浩写道，``保留问题，也是一种核对。''

夏夜星空没催他们继续走。霍雨浩从夹页取出另外三条表格校核题，放在横台中央。

``这三条按给定对应逐字核对。''他说，``不要凭刚才总结的局部规律改掉某个乘数。''

第一条，原表第618行：

\NovelDisplay{M01511R01}{%
(0)(1,1,1)(2,1)(2,1)(1,1)(2,2,1)(3,2)%
}
\NovelDisplay{M01511R02}{%
\quad (3,1)= \psi(\Omega_\omega^{2}\times\Omega+\psi_{\Omega_2}(\Omega_\omega\times\Omega))%
}

\NovelDisplay{M01512R01}{%
\psi\bigl(\Omega_\omega^2\times\Omega+\psi_{\Omega_2}(\Omega_\omega\times\Omega)\bigr)%
}

古月娜指出，外部前项 \allowbreak{}\NovelInline{I05828}{\Omega_\omega^2\times\Omega}\allowbreak{} 含平方，内函数参数 \allowbreak{}\NovelInline{I05829}{\Omega_\omega\times\Omega}\allowbreak{} 不含平方；外部共同背景与内函数参数要分别核对。这条题既没有把内函数升到上界 \allowbreak{}\NovelInline{I05830}{\Omega_3}\allowbreak{}，也没有把整个外部背景放进内函数。

第二条，原表第623行：

\NovelDisplay{M01513R01}{%
(0)(1,1,1)(2,1)(2,1)(1,1)(2,2,1)(3,2)(3,1)(2,2,1)(3,1)(2,2)(3,3,1)(4,3)%
}
\NovelDisplay{M01513R02}{%
\quad (4,1)= \psi(\Omega_\omega^{2}\times\Omega+\Omega_\omega\times\Omega+\psi_{\Omega_3}(\Omega_\omega\times\Omega))%
}

\NovelDisplay{M01514R01}{%
\psi\bigl(\Omega_\omega^2\times\Omega+\Omega_\omega\times\Omega+\psi_{\Omega_3}(\Omega_\omega\times\Omega)\bigr),\qquad\psi_{\Omega_3}\ne\psi_{\Omega_2}%
}

这一条在 \allowbreak{}\NovelInline{I05831}{\Omega_\omega^2\times\Omega}\allowbreak{} 后还有一份 \allowbreak{}\NovelInline{I05832}{\Omega_\omega\times\Omega}\allowbreak{}，最内侧函数上界也变了。霍雨浩逐对指过括号，卢本韦回到对应长列，古月娜再复核一次。

第三条，原表第632行：

\NovelDisplay{M01515R01}{%
(0)(1,1,1)(2,1)(2,1)(1,1)(2,2,1)(3,2)(3,1)(2,2,1)(3,2)(2,2,1)(2,2,1)%
}
\NovelDisplay{M01515R02}{%
\quad = \psi(\Omega_\omega\times(\Omega+1)+\Omega_\omega)%
}

\NovelDisplay{M01516R01}{%
\psi\bigl(\Omega_\omega\times(\Omega+1)+\Omega_\omega\bigr)=\psi\bigl(\Omega_\omega\times\Omega+\Omega_\omega\times2\bigr)%
}

这一条按原表保留，右端没有第一条那份 \allowbreak{}\NovelInline{I05833}{\Omega_\omega^2\times\Omega}\allowbreak{}。也不能为了与邻近公式整齐，擅自在这里补上平方。蔡须坤把原式重新抄了一遍，在等号下面画了两道线：从长列推测答案，与核对给定教材，是两种需要分别留下证据的工作。

三条校核结束，大家又把前面那条夹页题翻回来，确认四条都已收在本子里。古月娜发现霍雨浩漏了一个右括号。

``这个要留。''

霍雨浩补上，卢本韦将对应列也圈了出来。

接近一百二十米时，山风把本子的页角吹得哗哗响。夏夜星空在横台上写出下一阶段的预告：多重提升效应，以及提升链。

\NovelDisplay{M01517R01}{%
(0)(1,1,1)(2,1,1)(3,1)(2)%
}

这是教材后续讨论的一个边界位置。眼下不必抢先给它后面的所有结构填上闭式。蔡须坤抬头，一百二十米的石桩已经能够看见。再往上的结构，还要沿着后续分析逐段走清。

夏夜星空写回前五百米最终要重新理解的目标，留下一组纸面预演。

\NovelDisplay{M01518R01}{%
(0)(1,1,1)(2,2)%
}

``这一项以后还要走到。今天先在这里预演它的展开，脚下这段仍以一百二十米为止。''

他们已有超过这层的魂力。现在神压把熟练调动时掠过的缝隙一条条显出来：每个人都能放出那股力量，可谁若说不清最后两行怎么带动复制，脚下就停。

霍雨浩另取一页，只写数字列。可核对的前段 BOCF 对应仍照教材保留；超出这份对应表已说明的范围，就从 BMS 的父项与基本列继续分析，不凭相似字形猜一个坍缩表达式。

\NovelDisplay{M01519R01}{%
\bigl((0)(1,1,1)(2,2)\bigr)[0]=(0)(1,1,1)%
}

\NovelDisplay{M01520R01}{%
\bigl((0)(1,1,1)(2,2)\bigr)[1]=(0)(1,1,1)(2,1,1)%
}

\NovelDisplay{M01521R01}{%
\bigl((0)(1,1,1)(2,2)\bigr)[2]=(0)(1,1,1)(2,1,1)(3,1,1)%
}

\NovelDisplay{M01522R01}{%
\bigl((0)(1,1,1)(2,2)\bigr)[3]=(0)(1,1,1)(2,1,1)(3,1,1)(4,1,1)%
}

``不是只背这三项。''卢本韦说，``把下一项怎么来的也说清。''

最后列是 \allowbreak{}\NovelInline{I05834}{(2,2)}\allowbreak{}。它最高非零坐标在第二行，沿第一行父项链向左，第二行首先遇到更小值的父项是 \allowbreak{}\NovelInline{I05835}{(1,1,1)}\allowbreak{}。它就是本次坏根；复制段从它开始，截在末列之前。

\NovelDisplay{M01523R01}{%
\text{坏根}=(1,1,1),\qquad\delta_1=2-1=1,\qquad\delta_j=0\quad(j\ge2)%
}

\NovelDisplay{M01524R01}{%
\mathfrak b_0=(1,1,1),\qquad\mathfrak b_i=(i+1,1,1)\quad(i\ge0)%
}

增量只作用在最高非零行以下、并满足祖先经过条件的坐标。这里可动的是第一行，所以复制出来的第二、三行都保持一。保留下来的原前缀始终没有改写。

\NovelDisplay{M01525R01}{%
\bigl((0)(1,1,1)(2,2)\bigr)[n]=(0)(1,1,1)\prod_{i=1}^{n}(i+1,1,1)(n\ge0)%
}

这里的乘积符号专指数字列依次拼接，零项采用空拼接；不是序数乘法。霍雨浩在下面换一种写法，让拼接的位置更清楚。

\NovelDisplay{M01526R01}{%
\bigl((0)(1,1,1)(2,2)\bigr)[n]=(0)(1,1,1)\,\mathfrak b_1\mathfrak b_2\cdots\mathfrak b_n,\mathfrak b_i=(i+1,1,1)%
}

\NovelDisplay{M01527R01}{%
\mathfrak b_{i+1}=(i+2,1,1),\qquad\text{第二、三行始终各为一。}%
}

古月娜先区分前缀和复制段。压力使她无法将几处运转一起掠过，只好按第一行的增量逐处尝试，后两行保持原有秩序。

蔡须坤跟着试，第二行却也被惯性带动，胸口立刻绷紧。卢本韦让她停在此处，指向增量公式。

``最高非零的是第二行，它不加。第三行也不加。''

霍雨浩按父项条件重新核对。蔡须坤收回多带动的部分，再运转一次；接缝处的牵扯松开，肩上的压力依然沉重。古月娜抬头，问她第一项的根。

``那里末列三行都是一，展开时，还会是同一个根吗？''

霍雨浩将两条结构分开写。

\NovelDisplay{M01528R01}{%
(0)(1,1,1)(2,1,1),(0)(1,1,1)(2,1,1)(3,1,1)%
}

对 \allowbreak{}\NovelInline{I05836}{\bigl((0)(1,1,1)(2,1,1)\bigr)}\allowbreak{}，末列最高非零行已是第三行。在上一行的父项链上，前面的 \allowbreak{}\NovelInline{I05837}{(1,1,1)}\allowbreak{} 第三行并不小于一，须继续追到 \allowbreak{}\NovelInline{I05838}{(0)}\allowbreak{}，坏根因此改为最初那列。

\NovelDisplay{M01529R01}{%
\text{坏根}=(0),\qquad\delta_1=2-0=2,\qquad\delta_2=1-0=1,\qquad\delta_3=0%
}

\NovelDisplay{M01530R01}{%
\bigl((0)(1,1,1)(2,1,1)\bigr)[0]=(0)(1,1,1)%
}

\NovelDisplay{M01531R01}{%
\bigl((0)(1,1,1)(2,1,1)\bigr)[1]=(0)(1,1,1)(2,1)(3,2,1)%
}

\NovelDisplay{M01532R01}{%
\bigl((0)(1,1,1)(2,1,1)\bigr)[2]=(0)(1,1,1)(2,1)(3,2,1)(4,2)(5,3,1)%
}

\NovelDisplay{M01533R01}{%
\bigl((0)(1,1,1)(2,1,1)\bigr)[3]=(0)(1,1,1)(2,1)(3,2,1)(4,2)(5,3,1)(6,3)(7,4,1)%
}

第一份复制中，原根 \allowbreak{}\NovelInline{I05839}{(0)}\allowbreak{} 生成 \allowbreak{}\NovelInline{I05840}{(2,1)}\allowbreak{}，原来的 \allowbreak{}\NovelInline{I05841}{(1,1,1)}\allowbreak{} 生成 \allowbreak{}\NovelInline{I05842}{(3,2,1)}\allowbreak{}。第三行的一仍不动，原根省略的第三行零也仍为零；省略写法没有删掉计算所需的行。

\NovelDisplay{M01534R01}{%
\text{第 }i\text{ 份复制段：}\quad(2i,i)(2i+1,i+1,1)\qquad(i\ge1)%
}

\NovelDisplay{M01535R01}{%
\bigl((0)(1,1,1)(2,1,1)\bigr)[n]=(0)(1,1,1)\prod_{i=1}^{n}\bigl((2i,i)(2i+1,i+1,1)\bigr)%
}

乘积仍只表示拼接。与刚才 \allowbreak{}\NovelInline{I05843}{\bigl((0)(1,1,1)(2,2)\bigr)}\allowbreak{} 每次新增一列不同，\allowbreak{}\NovelInline{I05844}{\bigl((0)(1,1,1)(2,1,1)\bigr)}\allowbreak{} 每次复制一个两列段。行数看上去一样，坏根却不同，因此生成规律也不同。

前段可以核对的 BOCF 对应也显示了这一区别。注意下面的 \allowbreak{}\NovelInline{I05845}{\omega\times2}\allowbreak{}、\allowbreak{}\NovelInline{I05846}{\omega^2}\allowbreak{} 位于 \allowbreak{}\NovelInline{I05847}{\Omega}\allowbreak{} 的下标中。

\NovelDisplay{M01536R01}{%
\bigl((0)(1,1,1)(2,1,1)\bigr)[0]=\psi(\Omega_\omega),%
}
\NovelDisplay{M01536R02}{%
\bigl((0)(1,1,1)(2,1,1)\bigr)[1]=\psi(\Omega_{\omega\times2}),%
}
\NovelDisplay{M01536R03}{%
\bigl((0)(1,1,1)(2,1,1)\bigr)[2]=\psi(\Omega_{\omega\times3}),%
}
\NovelDisplay{M01536R04}{%
\bigl((0)(1,1,1)(2,1,1)\bigr)[3]=\psi(\Omega_{\omega\times4})%
}

\NovelDisplay{M01537R01}{%
\bigl((0)(1,1,1)(2,1,1)\bigr)=\psi(\Omega_{\omega^2}), \bigl((0)(1,1,1)(2,1)(1,1,1)\bigr)=\psi(\Omega_\omega^2)%
}

一个是 \allowbreak{}\NovelInline{I05848}{\Omega}\allowbreak{} 的下标变为 \allowbreak{}\NovelInline{I05849}{\omega^2}\allowbreak{}，另一个是整个 \allowbreak{}\NovelInline{I05850}{\Omega_\omega}\allowbreak{} 平方。古月娜把纸递给蔡须坤，她读了两遍，才把这两层放回各自的位置。

接着看 \allowbreak{}\NovelInline{I05851}{\bigl((0)(1,1,1)(2,1,1)(3,1,1)\bigr)}\allowbreak{}。末列第三行仍为一，追根仍到最初的 \allowbreak{}\NovelInline{I05852}{(0)}\allowbreak{}，但第一行的阶差已经变为三，复制段也包含了前面那列 \allowbreak{}\NovelInline{I05853}{(2,1,1)}\allowbreak{}。

\NovelDisplay{M01538R01}{%
\text{坏根}=(0),\qquad\delta_1=3,\qquad\delta_2=1,\qquad\delta_3=0%
}

\NovelDisplay{M01539R01}{%
\bigl((0)(1,1,1)(2,1,1)(3,1,1)\bigr)[0]=(0)(1,1,1)(2,1,1)%
}

\NovelDisplay{M01540R01}{%
\bigl((0)(1,1,1)(2,1,1)(3,1,1)\bigr)[1]=(0)(1,1,1)(2,1,1)(3,1)(4,2,1)(5,2,1)%
}

\NovelDisplay{M01541R01}{%
\bigl((0)(1,1,1)(2,1,1)(3,1,1)\bigr)[2]=(0)(1,1,1)%
}
\NovelDisplay{M01541R02}{%
\quad (2,1,1)(3,1)(4,2,1)(5,2,1)(6,2)(7,3,1)(8,3,1)%
}

\NovelDisplay{M01542R01}{%
\bigl((0)(1,1,1)(2,1,1)(3,1,1)\bigr)[3]=(0)(1,1,1)(2,1,1)(3,1)(4,2,1)(5,2,1)%
}
\NovelDisplay{M01542R02}{%
\quad (6,2)(7,3,1)(8,3,1)(9,3)(10,4,1)(11,4,1)%
}

蔡须坤指着展开式：``列数增加了，序数却在往下展开。要比较的是结构。''

夏夜星空点头。霍雨浩把 \allowbreak{}\NovelInline{I05854}{\bigl((0)(1,1,1)(2,1,1)\bigr)}\allowbreak{} 与 \allowbreak{}\NovelInline{I05855}{\bigl((0)(1,1,1)(2,1,1)(3,1,1)\bigr)}\allowbreak{} 纳入同一个有限族，让最后一段从三道例题变成一条可继续生成的规则。

\NovelDisplay{M01543R01}{%
(0)(1,1,1)(2,1,1)\cdots(m,1,1)(m\ge2)%
}

\NovelDisplay{M01544R01}{%
\bigl((0)(1,1,1)(2,2)\bigr)[n]=(0)(1,1,1)(2,1,1)\cdots(n+1,1,1)(n\ge1)%
}

对每个 \allowbreak{}\NovelInline{I05856}{\bigl((0)(1,1,1)\cdots(m,1,1)\bigr)}\allowbreak{}，最高非零行都是第三行；末列的父项链继续追到初始零，前面所有要复制的列组成同一块，第一行加 \allowbreak{}\NovelInline{I05857}{m}\allowbreak{}，第二行加一，第三行保持原值。

\NovelDisplay{M01545R01}{%
\text{坏根}=(0),\qquad\delta_1=m,\qquad\delta_2=1,\qquad\delta_3=0%
}

\NovelDisplay{M01546R01}{%
\mathfrak c_i^{(m)}=(im,i)\,(im+1,i+1,1)%
}
\NovelDisplay{M01546R02}{%
\quad \,(im+2,i+1,1)\cdots(im+m-1,i+1,1)\qquad(i\ge1)%
}

\NovelDisplay{M01547R01}{%
\bigl((0)(1,1,1)\cdots(m,1,1)\bigr)[k]=(0)(1,1,1)%
}
\NovelDisplay{M01547R02}{%
\quad \cdots(m-1,1,1)\,\mathfrak c_1^{(m)}\cdots\mathfrak c_k^{(m)}(k\ge0)%
}

第一份块的首列，来自原坏根；后面各列来自原复制段里的三行结构。每块恰有 \allowbreak{}\NovelInline{I05858}{m}\allowbreak{} 列。这样不但能写前几个样本，也能给任意指定的有限 \allowbreak{}\NovelInline{I05859}{m,k}\allowbreak{} 生成展开。

\NovelDisplay{M01548R01}{%
\#\mathfrak c_i^{(m)}=m,%
}
\NovelDisplay{M01548R02}{%
\#\bigl((0)(1,1,1)\cdots(m,1,1)\bigr)[k]=m+km,%
}
\NovelDisplay{M01548R03}{%
\#\bigl((0)(1,1,1)(2,2)\bigr)[n]=2+n%
}

\NovelDisplay{M01549R01}{%
\bigl((0)(1,1,1)\cdots(m,1,1)\bigr)[0]=(0)(1,1,1)\cdots(m-1,1,1)%
}

\NovelDisplay{M01550R01}{%
\mathfrak c_1^{(2)}=(2,1)(3,2,1),\qquad\mathfrak c_1^{(3)}=(3,1)(4,2,1)(5,2,1)%
}

霍雨浩照抄过两列块，到了三列块仍下意识只留两处承接。压力在第三列传来，他伏在石沿上喘了好一会儿，才将漏掉的位置记清。

蔡须坤在旁陪着分段重试。卢本韦在自己的运转中固定另外两处，与他对照第三处的差别；古月娜追完父项链，再让他独自走全。等三列各有了位置，霍雨浩终于站稳。

卢本韦递过水：``换四列。先写，再试。''

他们练了许多遍，却始终没有把某个很大的有限下标说成极限。山路要求的是能处理任意有限复制次数的构造，以及能辨认其所属的上层结构；这与把第一百遍的结果当成终点不同。

\NovelDisplay{M01551R01}{%
\forall m\ge2\ \forall k\ge0:\quad \bigl((0)(1,1,1)\cdots(m,1,1)\bigr)[k]\ \text{按同一父项与复制规则生成}%
}

\NovelDisplay{M01552R01}{%
\forall n<\omega: \bigl((0)(1,1,1)(2,2)\bigr)[n]\ \text{是有限展开项},%
}
\NovelDisplay{M01552R02}{%
\bigl((0)(1,1,1)(2,2)\bigr)[n]\ne \bigl((0)(1,1,1)(2,2)\bigr)%
}

纸面预演收好以后，夏夜星空让众人翻回一百二十米所对应的那页。更高处的展开暂留在后面；眼下要核的是从七十米接到这里，哪些结构已经能由他们各自说明。四人重新受压，后面的同伴也按自己的记录逐段尝试，没有拿预演的终点替脚下落笔。

日光移到山的另一边时，一百二十米石桩下方的最后一道横纹亮了。蔡须坤没有猛提魂力，只照着方才确认的前缀与复制段，把原本早已超过这一层的力量重新稳稳地接起来。

古月娜与卢本韦先后站稳，霍雨浩收好核对完的最后一页。四人让到横台侧面，等身后的同伴陆续上来。唐三收回蓝银草，叶阳与薛峰压住被风翻起的纸，蔡兴坤抱着黑子鸡坐到石边。

夏夜星空看过十一人的落脚处，才在山图最下方那道横线旁落下一点光。

``今天这一段，过了。''

蔡须坤仰头。云间还有二千三百八十米，坤神殿仍在上面。两年的第一考没有在这里结束，前五百米这一项目也还没有完成。

她背上那点温热轻轻贴过来。奥斯卡把三层油纸又按紧了一回，朝远处看了一眼，仍把它收进最里面。

``明天从这里接着。''卢本韦说。

众人把最后一页夹进本子，一起沿来路回营。那些曾经只凭熟练便掠过的接缝，如今在纸上有了根、有了阶差，也有了每个人改错时留下的笔迹。

\chapter{第一百零七章：TFBO之后道路}\label{ux7b2cux4e00ux767eux96f6ux4e03ux7ae0tfboux4e4bux540eux9053ux8def}

一百二十米的石桩上多了一条细细的炭痕。

蔡须坤蹲下去，用拇指将那条痕擦净。石桩没有变高，她的靴底也没有离开横台。昨日纸上预演到五百米的几组展开，仍夹在霍雨浩本子的后面；真正站稳的，是眼前这一段。

七十米那里，他们用了整整两个月。往上一段，每逢以为已经会了，肩上的压力便会从一处没有说清的接缝压下来。直到把前段分析接到这里，四人才在一百二十米的石纹上留下完整的落脚印。

卢本韦把水囊搁在桩后，伸手去拧自己湿透的袖口。

``本子能先到五百米，人不能跟着纸飞过去。''

霍雨浩正要翻最后那页，被她用两根手指按住。他笑了笑，将预演页折回去。古月娜把另一张白纸压在石面上，银发从肩头滑下来，发梢沾着清晨的水汽。

``昨天只看了它怎样向下展开。今天还要知道，它上面那些括号各自容着什么。''

蔡须坤将炭条掰成四截，一人一截。

``每个人都要写。看懂旁人的答案，到了台阶上未必接得起来。''

他们从刚刚站稳的位置重新写起。

\NovelDisplay{M01553R01}{%
(0)(1,1,1)(2,1)(3,2)=\psi(\Omega_{\omega+1})%
}

霍雨浩在纸边重申记号：\allowbreak{}\NovelInline{I05860}{\psi_{\Omega_2}}\allowbreak{}、\allowbreak{}\NovelInline{I05861}{\psi_{\Omega_{\omega+1}}}\allowbreak{}、\allowbreak{}\NovelInline{I05862}{\psi_{\Omega_{\omega\times2+1}}}\allowbreak{} 的下标均写坍缩上界。没有另标下标的 \allowbreak{}\NovelInline{I05863}{\psi}\allowbreak{} 仍指 \allowbreak{}\NovelInline{I05864}{\psi_{\Omega}}\allowbreak{}。数字列只省略末尾补齐用的零，第一行的零保留；读父项时，缺省坐标仍按零计算。

夏夜星空将通往上方的路让开。唐三先把蓝银草系在石柱上，绳尾沿外侧垂下；奥斯卡把温水送到横台边。夏夜星空只伸手摸过第一道石纹，压力便重新落到四人肩上。

蔡须坤脚跟向下沉了半寸。她没有急着提起更强的魂力，先把已经站稳的前段留住，让末尾新添的结构在缓慢运转中显出位置。

古月娜在她身边坐下，掌心贴住石面。

``先找这次增加的是哪一项。''

\NovelDisplay{M01554R01}{%
(0)(1,1,1)(2,1)(3,2)(1,1,1)=\psi(\Omega_{\omega+1}+\Omega_{\omega})%
}

卢本韦把最后那个一一一涂紫，问蔡须坤：``这次根会停在前面的一一一吗？''蔡须坤没有按字形回答，从末列的第三行往左追。霍雨浩将列号排在下面，古月娜逐个核对候选。

\NovelDisplay{M01555R01}{%
\textcolor[HTML]{B45309}{(0)}\textcolor[HTML]{B91C1C}{(1,1,1)}\textcolor[HTML]{B91C1C}{(2,1)}\textcolor[HTML]{B91C1C}{(3,2)}\textcolor[HTML]{7E22CE}{(1,1,1)}%
}

\NovelDisplay{M01556R01}{%
\text{原列数}=5, h=3, r=1,%
}
\NovelDisplay{M01556R02}{%
\text{坏根}=\textcolor[HTML]{B45309}{(0)},\text{坏部列数}=4,%
}
\NovelDisplay{M01556R03}{%
(\delta_1,\delta_2,\delta_3)=(1,1,0).%
}

\NovelDisplay{M01557R01}{%
\text{按坏部列序排列的提升阈值：}%
}
\NovelDisplay{M01557R02}{%
(2,2,2,2),%
}
\NovelDisplay{M01557R03}{%
\text{第 }j\text{ 列只允许前 }\tau_j\text{ 行应用阶差。}%
}

\NovelDisplay{M01558R01}{%
\text{原第 }1\text{ 列：}(0)%
}
\NovelDisplay{M01558R02}{%
\longmapsto(q,q), q\ge1.%
}

\NovelDisplay{M01559R01}{%
\text{原第 }4\text{ 列：}(3,2)%
}
\NovelDisplay{M01559R02}{%
\longmapsto(3+q,2+q), q\ge1.%
}

这里的增量按原坏部计算。最高非零行及其下方坐标保持原值；有非零阶差的更前行，还要逐列满足祖先条件。阈值中的列号按坏部次序重新编号，新增段的编号 q 为正整数。

\NovelDisplay{M01560R01}{%
\bigl((0)(1,1,1)(2,1)(3,2)(1,1,1)\bigr)[0]=\textcolor[HTML]{B45309}{(0)}\textcolor[HTML]{B91C1C}{(1,1,1)}\textcolor[HTML]{B91C1C}{(2,1)}\textcolor[HTML]{B91C1C}{(3,2)}%
}

\NovelDisplay{M01561R01}{%
\bigl((0)(1,1,1)(2,1)(3,2)(1,1,1)\bigr)[1]=\textcolor[HTML]{B45309}{(0)}%
}
\NovelDisplay{M01561R02}{%
\quad \textcolor[HTML]{B91C1C}{(1,1,1)}\textcolor[HTML]{B91C1C}{(2,1)}\textcolor[HTML]{B91C1C}{(3,2)}\textcolor[HTML]{1D4ED8}{(1,1)}\textcolor[HTML]{1D4ED8}{(2,2,1)}\textcolor[HTML]{1D4ED8}{(3,2)}\textcolor[HTML]{1D4ED8}{(4,3)}%
}

\NovelDisplay{M01562R01}{%
\bigl((0)(1,1,1)(2,1)(3,2)(1,1,1)\bigr)[2]=\textcolor[HTML]{B45309}{(0)}\textcolor[HTML]{B91C1C}{(1,1,1)}\textcolor[HTML]{B91C1C}{(2,1)}\textcolor[HTML]{B91C1C}{(3,2)}\textcolor[HTML]{1D4ED8}{(1,1)}\textcolor[HTML]{1D4ED8}{(2,2,1)}\textcolor[HTML]{1D4ED8}{(3,2)}%
}
\NovelDisplay{M01562R02}{%
\quad \textcolor[HTML]{1D4ED8}{(4,3)}\textcolor[HTML]{15803D}{(2,2)}\textcolor[HTML]{15803D}{(3,3,1)}\textcolor[HTML]{15803D}{(4,3)}\textcolor[HTML]{15803D}{(5,4)}%
}

\NovelDisplay{M01563R01}{%
\text{第零项列数}=4,%
}
\NovelDisplay{M01563R02}{%
\text{第一项列数}=4+4=8,%
}
\NovelDisplay{M01563R03}{%
\text{第二项列数}=4+2\times4=12.%
}

蔡须坤先盖住答案，按列数把第一项与第二项各写一次。霍雨浩核坐标，卢本韦看保留段有没有误留末列，古月娜再查比较方向。两项都从当前原式求得，不能把第一项再次展开的结果填到第二项的位置。

卢本韦将末列与前面的一一一连起来，炭条刚划到一半，蔡须坤便伸手按住了纸。

``我刚才就是这么看的。看见相同的列，就去找它们相同的作用。''

``所以这次你先别看右边。''古月娜说。

霍雨浩将写着序数的那半张纸折到背面。蔡须坤重新补齐每列的第三行，从末列最高的非零行出发。候选并非前面字形相同的一一一，而是最左端的零列。她顺着父项说明为什么能到那里，又把根后的四列逐个点过。末列移去，原来的四列留下；接下来的新增段由坏部生成。整个过程里，没有哪一步允许因为两列相同就直接把它们合并。

``那右边为什么还有原来的主项？''卢本韦问。

蔡须坤翻开折起的半页。这回她先读最外层的坍缩，再读参数里的加号。原来的主项在加号前面，后来出现的那一项接在后面。两者的位置写得清楚；所谓此处没有提升，说的是这个主项没有在当前的 BOCF 写法中被整体替换。局部展开里仍有阶差，不能因此声称所有坐标都不增加。

她说到这里，忽然停下来，将自己刚写的``没有增量''四字划去。

霍雨浩没有替她接。``把两件事分别说完。''

``BMS 的坐标怎样生成，是一件事。对应序数归一化之后，哪一段出现在什么位置，是另一件事。要连接起来，需要前缀和展开共同给出依据。''

古月娜点了点纸上的两种颜色。根的位置用橙色，原坏部其余列用红色，删去的末列用紫色；展开第一项中新接的段用蓝色，展开第二项再接的段用绿色。颜色标明每段的来历，同一颜色里的坐标仍须分别检查。遇到有阶差却不满足祖先条件的位置，才另用灰色标出。

卢本韦把自己的答案扣在膝上，请蔡须坤独自生成第一项。蔡须坤刚接完最后一列，肩后的压力便向左偏了一下。她下意识想用魂力将那一侧顶住，手抬起来又放下，回头检查自己漏掉的那行。

是她在简写之后忘了补零。

改过来以后，偏斜没有立刻消失。她又把刚才的说明从头说了一遍，确定那处坐标为什么不动，才重新尝试。压力仍然沉，左右两侧却渐渐落到同样的位置。

``可不是我把你按稳的。''卢本韦把手从她背后拿开，亮出空着的掌心。

蔡须坤笑了一声，笑过才发觉自己满手都是炭灰。霍雨浩递过湿布，没有将这次站稳写成实力又跨了一级。他只在页角记下：原先能调动的魂力，今天终于多了一处可以解释的结构。

夏夜星空等他们把颜色对应的来历都说清，才收回压在纸角的手。往后的连续几组，四人按相同读法自行核对，不再为每个等号另讲一遍开头。

此处没有提升。原理和 \allowbreak{}\NovelInline{I05865}{(0)(1,1,1)(2,1)(1,1,1)(1,1,1)}\allowbreak{} 处是有些类似的。局部提升仍发生在更内层的 BMS 展开结构中。归一化成 BOCF 后，当前式子的主参数 \allowbreak{}\NovelInline{I05866}{\Omega_{\omega+1}}\allowbreak{} 保持不变，这里的 \allowbreak{}\NovelInline{I05867}{(1,1,1)}\allowbreak{} 在其后添加一个 \allowbreak{}\NovelInline{I05868}{\Omega_\omega}\allowbreak{}。

\NovelDisplay{M01564R01}{%
(0)(1,1,1)(2,1)(3,2)(1,1,1)(2,1)=\psi(\Omega_{\omega+1}+\Omega_{\omega}\times\Omega)%
}

\NovelDisplay{M01565R01}{%
(0)(1,1,1)(2,1)(3,2)(1,1,1)(2,1)(1,1,1)=\psi(\Omega_{\omega+1}+\Omega_{\omega}^{2})%
}

这个还是可以提升的。

\NovelDisplay{M01566R01}{%
(0)(1,1,1)(2,1)(3,2)(1,1,1)(2,1)(3,1)=\psi(\Omega_{\omega+1}+\Omega_{\omega}^{\Omega})%
}

\NovelDisplay{M01567R01}{%
(0)(1,1,1)(2,1)(3,2)(1,1,1)(2,1)(3,1)(1,1,1)=\psi(\Omega_{\omega+1}+\Omega_{\omega}^{\Omega_{\omega}})%
}

\NovelDisplay{M01568R01}{%
(0)(1,1,1)(2,1)(3,2)(1,1,1)(2,1)(3,2)=\psi(\Omega_{\omega+1}+\psi_{\Omega_{\omega+1}}(\Omega_{\omega+1}))%
}

这里仍是相同 BMS 片段的重复。在这份 BOCF 笔记中，要按上式保留基数主项与后面的坍缩项，不能把它直接读成整个外层序数乘二。接下来继续看乘数怎样增长：

\NovelDisplay{M01569R01}{%
(0)(1,1,1)(2,1)(3,2)(2)=\psi(\Omega_{\omega+1}+\psi_{\Omega_{\omega+1}}(\Omega_{\omega+1})\times\omega)%
}

\NovelDisplay{M01570R01}{%
(0)(1,1,1)(2,1)(3,2)(2,1)=\psi(\Omega_{\omega+1}+\psi_{\Omega_{\omega+1}}(\Omega_{\omega+1})\times\Omega)%
}

\NovelDisplay{M01571R01}{%
(0)(1,1,1)(2,1)(3,2)(2,1)(1,1,1)=\psi(\Omega_{\omega+1}+\psi_{\Omega_{\omega+1}}(\Omega_{\omega+1})\times\Omega_{\omega})%
}

又是一个提升。过程和之前一模一样。

\NovelDisplay{M01572R01}{%
(0)(1,1,1)(2,1)(3,2)(2,1)(2,1)(1,1,1)%
}
\NovelDisplay{M01572R02}{%
\quad =\psi(\Omega_{\omega+1}+\psi_{\Omega_{\omega+1}}(\Omega_{\omega+1})\times\Omega_{\omega}^{2})%
}

\NovelDisplay{M01573R01}{%
(0)(1,1,1)(2,1)(3,2)(2,1)(3,1)(1,1,1)%
}
\NovelDisplay{M01573R02}{%
\quad =\psi(\Omega_{\omega+1}+\psi_{\Omega_{\omega+1}}(\Omega_{\omega+1})\times\Omega_{\omega}^{\Omega_{\omega}})%
}

\NovelDisplay{M01574R01}{%
(0)(1,1,1)(2,1)(3,2)(2,1)(3,2)=\psi(\Omega_{\omega+1}+\psi_{\Omega_{\omega+1}}(\Omega_{\omega+1})^{2})%
}

\NovelDisplay{M01575R01}{%
(0)(1,1,1)(2,1)(3,2)(3)=\psi(\Omega_{\omega+1}+\psi_{\Omega_{\omega+1}}(\Omega_{\omega+1})^{\omega})%
}

指数环节。

\NovelDisplay{M01576R01}{%
(0)(1,1,1)(2,1)(3,2)(3,1)=\psi(\Omega_{\omega+1}+\psi_{\Omega_{\omega+1}}(\Omega_{\omega+1})^{\Omega})%
}

\NovelDisplay{M01577R01}{%
(0)(1,1,1)(2,1)(3,2)(3,1)(1,1,1)=\psi(\Omega_{\omega+1}+\psi_{\Omega_{\omega+1}}(\Omega_{\omega+1})^{\Omega_{\omega}})%
}

\NovelDisplay{M01578R01}{%
(0)(1,1,1)(2,1)(3,2)(3,1)(4,1)(1,1,1)=\psi(\Omega_{\omega+1}%
}
\NovelDisplay{M01578R02}{%
\quad +\psi_{\Omega_{\omega+1}}(\Omega_{\omega+1})^{\Omega_{\omega}^{\Omega_{\omega}}})%
}

\NovelDisplay{M01579R01}{%
(0)(1,1,1)(2,1)(3,2)(3,1)(4,2)=\psi(\Omega_{\omega+1}%
}
\NovelDisplay{M01579R02}{%
\quad +\psi_{\Omega_{\omega+1}}(\Omega_{\omega+1})^{\psi_{\Omega_{\omega+1}}(\Omega_{\omega+1})})%
}

\allowbreak{}\NovelInline{I05869}{\psi_{\Omega_{\omega+1}}(\Omega_{\omega+1})}\allowbreak{} 的指数塔具形。

\NovelDisplay{M01580R01}{%
(0)(1,1,1)(2,1)(3,2)(3,1)(4,2)(4,1)(5,2)=\psi(\Omega_{\omega+1}%
}
\NovelDisplay{M01580R02}{%
\quad +\psi_{\Omega_{\omega+1}}(\Omega_{\omega+1})^{\psi_{\Omega_{\omega+1}}(\Omega_{\omega+1})^{\psi_{\Omega_{\omega+1}}(\Omega_{\omega+1})}})%
}

蔡须坤把相邻的两列重写了一遍，停笔时，手腕仍受着向下的牵扯。

``它们看着一样，所以我刚才想把它们当成两份相同的力量。''

卢本韦没有接答案，只把最外面的括号圈出来。古月娜从她圈出的地方向内指，霍雨浩另画一条线，标出读数时尚未进入的函数。

``重复的是左边的片段。''他说，``右边怎样组织，要把根和整个前缀一起看。''

蔡须坤重新试了两次。第一回仍按自己熟悉的合并方法，胸口顿时发闷；第二回先将基数主项留在原处，只动后面的坍缩项，牵扯才在纸上标记的位置停住。她靠在石沿上喘匀气，把这两回的差别都记了下来。

\NovelDisplay{M01581R01}{%
(0)(1,1,1)(2,1)(3,2)(3,2)=\psi(\Omega_{\omega+1}\times2)%
}

两列看着相同，古月娜却先让霍雨浩删去末列，再划复制段。阶差和新增列数分别写清，才能知道``展开第一项''究竟多出哪几列。

\NovelDisplay{M01582R01}{%
(0)(1,1,1)\textcolor[HTML]{B45309}{(2,1)}\textcolor[HTML]{B91C1C}{(3,2)}\textcolor[HTML]{7E22CE}{(3,2)}%
}

\NovelDisplay{M01583R01}{%
\text{原列数}=5, h=2, r=3,%
}
\NovelDisplay{M01583R02}{%
\text{坏根}=\textcolor[HTML]{B45309}{(2,1)},\text{坏部列数}=2,%
}
\NovelDisplay{M01583R03}{%
(\delta_1,\delta_2,\delta_3)=(1,0,0).%
}

\NovelDisplay{M01584R01}{%
\text{按坏部列序排列的提升阈值：}%
}
\NovelDisplay{M01584R02}{%
(1,1),%
}
\NovelDisplay{M01584R03}{%
\text{第 }j\text{ 列只允许前 }\tau_j\text{ 行应用阶差。}%
}

\NovelDisplay{M01585R01}{%
\text{原第 }3\text{ 列：}(2,1)%
}
\NovelDisplay{M01585R02}{%
\longmapsto(2+q,1), q\ge1.%
}

\NovelDisplay{M01586R01}{%
\text{原第 }4\text{ 列：}(3,2)%
}
\NovelDisplay{M01586R02}{%
\longmapsto(3+q,2), q\ge1.%
}

霍雨浩先检查根以前的保留段，蔡须坤再沿原坏部核阈值。每一新增段使用同一份原坏部与阶差，编号 q 为正整数；最高非零行及更后坐标不因段数增多而跟着改变。

\NovelDisplay{M01587R01}{%
\bigl((0)(1,1,1)(2,1)(3,2)(3,2)\bigr)[0]=(0)(1,1,1)\textcolor[HTML]{B45309}{(2,1)}\textcolor[HTML]{B91C1C}{(3,2)}%
}

\NovelDisplay{M01588R01}{%
\bigl((0)(1,1,1)(2,1)(3,2)(3,2)\bigr)[1]=(0)(1,1,1)\textcolor[HTML]{B45309}{(2,1)}\textcolor[HTML]{B91C1C}{(3,2)}\textcolor[HTML]{1D4ED8}{(3,1)}\textcolor[HTML]{1D4ED8}{(4,2)}%
}

\NovelDisplay{M01589R01}{%
\bigl((0)(1,1,1)(2,1)(3,2)(3,2)\bigr)[2]=(0)%
}
\NovelDisplay{M01589R02}{%
\quad (1,1,1)\textcolor[HTML]{B45309}{(2,1)}\textcolor[HTML]{B91C1C}{(3,2)}\textcolor[HTML]{1D4ED8}{(3,1)}\textcolor[HTML]{1D4ED8}{(4,2)}\textcolor[HTML]{15803D}{(4,1)}\textcolor[HTML]{15803D}{(5,2)}%
}

\NovelDisplay{M01590R01}{%
\text{第零项列数}=4,%
}
\NovelDisplay{M01590R02}{%
\text{第一项列数}=4+2=6,%
}
\NovelDisplay{M01590R03}{%
\text{第二项列数}=4+2\times2=8.%
}

古月娜将第二项的两段新增部分分别圈出，卢本韦从原坏部核它们各自的增量。蔡须坤专查那些本该不动的位置，霍雨浩最后逐列比较。各项虽然可以更长，仍位于原式之下。

接下来继续观察乘数与指数怎样随具体 BMS 结构增长。

\NovelDisplay{M01591R01}{%
(0)(1,1,1)(2,1)(3,2)(4)=\psi(\Omega_{\omega+1}\times\omega)%
}

\NovelDisplay{M01592R01}{%
(0)(1,1,1)(2,1)(3,2)(4)(5,1,1)(6,1)(7,2)=\psi(\Omega_{\omega+1}\times\psi(\Omega_{\omega+1}))%
}

\NovelDisplay{M01593R01}{%
(0)(1,1,1)(2,1)(3,2)(4,1)=\psi(\Omega_{\omega+1}\times\Omega)%
}

接下来算是一个关键点。眼前的 \allowbreak{}\NovelInline{I05870}{(1,1,1)}\allowbreak{} 能把乘数上的、指数上的 \allowbreak{}\NovelInline{I05871}{\Omega}\allowbreak{} 提升为 \allowbreak{}\NovelInline{I05872}{\Omega_\omega}\allowbreak{} ，那么对于 \allowbreak{}\NovelInline{I05873}{\psi}\allowbreak{} 里面的呢？答案是：

\NovelDisplay{M01594R01}{%
(0)(1,1,1)(2,1)(3,2)(4,1)(1,1)=\psi(\Omega_{\omega+1}\times\Omega+\Omega)%
}

霍雨浩先保留原前缀，等待下一项的变化落到可核对的位置。

\NovelDisplay{M01595R01}{%
(0)(1,1,1)(2,1)(3,2)(4,1)(1,1)(2,2,1)(3,2)%
}
\NovelDisplay{M01595R02}{%
\quad (4,3)=\psi(\Omega_{\omega+1}\times\Omega+\psi_{\Omega_2}(\Omega_{\omega+1}))%
}

\NovelDisplay{M01596R01}{%
(0)(1,1,1)(2,1)(3,2)(4,1)(1,1)(2,2,1)(3,2)(4,3)%
}
\NovelDisplay{M01596R02}{%
\quad (5,1)=\psi(\Omega_{\omega+1}\times\Omega+\psi_{\Omega_2}(\Omega_{\omega+1}\times\Omega))%
}

替身 \allowbreak{}\NovelInline{I05874}{(5,1)}\allowbreak{} 召唤。

\NovelDisplay{M01597R01}{%
(0)(1,1,1)(2,1)(3,2)(4,1)(1,1)(2,2,1)(3,2)(4,3)(5,1)(2,2)(3,2)%
}
\NovelDisplay{M01597R02}{%
\quad =\psi(\Omega_{\omega+1}\times\Omega+\Omega_2^{2})%
}

\NovelDisplay{M01598R01}{%
(0)(1,1,1)(2,1)(3,2)(4,1)(1,1)(2,2,1)(3,2)%
}
\NovelDisplay{M01598R02}{%
\quad (4,3)(5,1)(2,2,1)=\psi(\Omega_{\omega+1}\times\Omega+\Omega_{\omega})%
}

一个似曾相识的结构。蔡须坤在旧册里找到它的前例：

\NovelDisplay{I05875}{(0)(1,1,1)(2,1)(1,1)(2,2,1)(3,1)(2,2,1)=\psi(\Omega_{\omega}\times\Omega+\Omega_{\omega})}

在那之后，同样是 \allowbreak{}\NovelInline{I05876}{(2,2,1)}\allowbreak{} 替代了 \allowbreak{}\NovelInline{I05877}{(1,1,1)}\allowbreak{} 进行``添加 \allowbreak{}\NovelInline{I05878}{\Omega_\omega}\allowbreak{} ''的功能。

\NovelDisplay{M01599R01}{%
(0)(1,1,1)(2,1)(3,2)(4,1)(1,1)(2,2,1)(3,2)(4,3)(5,1)(2,2,1)(3,1)%
}
\NovelDisplay{M01599R02}{%
\quad =\psi(\Omega_{\omega+1}\times\Omega+\Omega_{\omega}\times\Omega)%
}

召唤\ldots\ldots 等等，这里的 \allowbreak{}\NovelInline{I05879}{(3,1)}\allowbreak{} 是替身吗？

\NovelDisplay{M01600R01}{%
(0)(1,1,1)(2,1)(3,2)(4,1)(1,1)(2,2,1)(3,2)(4,3)(5,1)(2,2,1)(3,1)(1,1)%
}
\NovelDisplay{M01600R02}{%
\quad (2,2,1)(3,2)(4,3)(5,1)(2,2,1)(3,1)=\psi(\Omega_{\omega+1}\times\Omega+\Omega_{\omega}\times\Omega%
}
\NovelDisplay{M01600R03}{%
\quad +\psi_{\Omega_2}(\Omega_{\omega+1}\times\Omega+\Omega_{\omega}\times\Omega))%
}

\NovelDisplay{M01601R01}{%
(0)(1,1,1)(2,1)(3,2)(4,1)(1,1)(2,2,1)(3,2)(4,3)(5,1)(2,2,1)(3,1)(2,2)%
}
\NovelDisplay{M01601R02}{%
\quad =\psi(\Omega_{\omega+1}\times\Omega+\Omega_{\omega}\times\Omega+\Omega_2)%
}

\NovelDisplay{M01602R01}{%
(0)(1,1,1)(2,1)(3,2)(4,1)(1,1)(2,2,1)(3,2)(4,3)(5,1)(2,2,1)(3,1)(2,2)%
}
\NovelDisplay{M01602R02}{%
\quad (3,2)=\psi(\Omega_{\omega+1}\times\Omega+\Omega_{\omega}\times\Omega+\Omega_2^{2})%
}

\NovelDisplay{M01603R01}{%
(0)(1,1,1)(2,1)(3,2)(4,1)(1,1)(2,2,1)(3,2)(4,3)(5,1)(2,2,1)(3,1)(2,2,1)%
}
\NovelDisplay{M01603R02}{%
\quad =\psi(\Omega_{\omega+1}\times\Omega+\Omega_{\omega}\times\Omega+\Omega_{\omega})%
}

纸上的紫色逐渐过渡成粉色。蔡须坤沿着那几列往下读，读到第二行时，忽然把笔拿了起来。

``这里一直还是一。''

``而旁边的列在增加。''霍雨浩说。

卢本韦把两条连线分别画到末端。第一条走过坏根，第二条绕过去，停在原来允许停下的父项上。古月娜请她将原式的列号也添上：长式一旦换行，相同的数字会出现在不同位置；只说`那个三一'，旁人未必知道她正在指哪一列。

蔡须坤按着列号重走。阶差的前两行都是非零的，这只是给出了可用的增量；具体到每一列，第二行是否能接受这个增量，还要看那一行的祖先关系。她原先把这两个条件揉在一起，算出阶差便将整段的前两行同时提高，结果恰好错在了粉色标出的那几个一上。

``我能记住这里不能动。''她说，``可下一次换了数字呢？''

``所以别背这个一。''古月娜将那一小块盖住，``从它的父项追。''

霍雨浩先收回精神共享，只留下十一人之间用于警戒的细线。蔡须坤没有了现成的路径，算到中段便慢下来。卢本韦等她算完，也用另一张纸追同一条链。两人的线在相邻位置分开，古月娜没有立即判谁正确，让她们分别说出每一步为什么能选到下一个候选。

卢本韦在第三步停住。她忽略了前一行施加的限制，单凭这一行的数值更小便往左跳了过去。

``那我刚才也不能算真的会了。''她将自己的线擦断，``我的答案对，走法错。''

蔡须坤将她擦掉的那一截另画到纸边。正确的结果不能替代选择父项的理由，碰巧绕过坏根，也不能保证换一道题仍会绕过。两人把这句留在错线旁，再去生成后面的新增段。

古月娜这时才解开盖住数字的布。粉色的一仍在原处，与两人重新求出的祖先链一致。

山风从横台下面卷上来，霍雨浩伸手压住纸，却没压住卢本韦被吹到脸上的一缕头发。她一边将头发别回耳后，一边指着纸笑：``就这一个一，害我给它讲了半天情。''

``它也没答应你。''蔡须坤说。

笑声让紧绷的气息松了一点。两人起身试步，蔡须坤负责指出压力落在哪个环节，卢本韦负责把那一环写回原式。霍雨浩观察两人的先后，却不把体内的感受直接命名为某个序数；古月娜再拿父项检查两人的说明。压力帮助两人找出含混处，最终能否解释清楚，仍要回到纸上。

这一次，两人在那处停顿前主动放慢。原先常常一并带起的第二行留在原位，前一行按自己的条件接上。脚下的石纹只亮了一小截，尚不到更高的里程桩。蔡须坤却看见了那一小截为什么能亮。

在 \allowbreak{}\NovelInline{I05880}{(3,1)}\allowbreak{} 之后的 \allowbreak{}\NovelInline{I05881}{(2,2,1)}\allowbreak{} 的作用仍然是添加 \allowbreak{}\NovelInline{I05882}{\Omega_\omega}\allowbreak{} 。这说明这里的 \allowbreak{}\NovelInline{I05883}{(3,1)}\allowbreak{} 并非替身。想要分析原因的话，可以试着写出它的展开式：

\NovelDisplay{M01604R01}{%
(0)(1,1,1)(2,1)(3,2)(4,1)\textcolor[HTML]{0000FF}{(1,1)}\textcolor[HTML]{0000FF}{(2,2,1)}\textcolor[HTML]{0000FF}{(3,2)}\textcolor[HTML]{0000FF}{(4,3)}\textcolor[HTML]{0000FF}{(5,1)}\textcolor[HTML]{0000FF}{(2,2,1)}\textcolor[HTML]{EE00FF}{(3,1)}\textcolor[HTML]{0066FF}{(2,2)}%
}
\NovelDisplay{M01604R02}{%
\quad \textcolor[HTML]{0066FF}{(3,3,1)}\textcolor[HTML]{0066FF}{(4,3)}\textcolor[HTML]{0066FF}{(5,4)}\textcolor[HTML]{0066FF}{(6,1)}\textcolor[HTML]{0066FF}{(3,3,1)}\textcolor[HTML]{EE66FF}{(4,1)}\textcolor[HTML]{00CCFF}{(3,3)}%
}
\NovelDisplay{M01604R03}{%
\quad \textcolor[HTML]{00CCFF}{(4,4,1)}\textcolor[HTML]{00CCFF}{(5,4)}\textcolor[HTML]{00CCFF}{(6,5)}\textcolor[HTML]{00CCFF}{(7,1)}\textcolor[HTML]{00CCFF}{(4,4,1)}\textcolor[HTML]{EECCFF}{(5,1)}\dots%
}

注意被标为紫色到粉色的这几项。可以看到它们的第二项由于祖先项不经过坏根 \allowbreak{}\NovelInline{I05884}{(1,1)}\allowbreak{} ，并没有得到提升。所以这里没有提升效应。这也是这个 \allowbreak{}\NovelInline{I05885}{(3,1)}\allowbreak{} 不算替身的原因。

\NovelDisplay{M01605R01}{%
(0)(1,1,1)(2,1)(3,2)(4,1)(1,1)(2,2,1)(3,2)(4,3)(5,1)(2,2,1)(3,1)(2,2,1)%
}
\NovelDisplay{M01605R02}{%
\quad (3,1)=\psi(\Omega_{\omega+1}\times\Omega+\Omega_{\omega}\times\Omega\times2)%
}

\NovelDisplay{M01606R01}{%
(0)(1,1,1)(2,1)(3,2)(4,1)(1,1)(2,2,1)(3,2)(4,3)(5,1)(2,2,1)(3,1)(3,1)%
}
\NovelDisplay{M01606R02}{%
\quad =\psi(\Omega_{\omega+1}\times\Omega+\Omega_{\omega}\times\Omega^{2})%
}

\NovelDisplay{M01607R01}{%
(0)(1,1,1)(2,1)(3,2)(4,1)(1,1)(2,2,1)(3,2)(4,3)(5,1)(2,2,1)(3,1)(4,2)%
}
\NovelDisplay{M01607R02}{%
\quad =\psi(\Omega_{\omega+1}\times\Omega+\Omega_{\omega}\times\psi_{\Omega_2}(\Omega_2))%
}

\NovelDisplay{M01608R01}{%
(0)(1,1,1)(2,1)(3,2)(4,1)(1,1)(2,2,1)(3,2)(4,3)(5,1)(2,2,1)(3,2)%
}
\NovelDisplay{M01608R02}{%
\quad =\psi(\Omega_{\omega+1}\times\Omega+\Omega_{\omega}\times\Omega_2)%
}

但现在到了 \allowbreak{}\NovelInline{I05886}{(3,2)}\allowbreak{} ，就没有那么好运了。 \allowbreak{}\NovelInline{I05887}{(3,2)}\allowbreak{} 的第二行的祖先项会经过坏根 \allowbreak{}\NovelInline{I05888}{(1,1)}\allowbreak{} ，因此：

\NovelDisplay{M01609R01}{%
(0)(1,1,1)(2,1)(3,2)(4,1)(1,1)(2,2,1)(3,2)(4,3)(5,1)(2,2,1)(3,2)(2,2,1)%
}
\NovelDisplay{M01609R02}{%
\quad =\psi(\Omega_{\omega+1}\times\Omega+\Omega_{\omega}^{2})%
}

提升效应虽迟但到。

\NovelDisplay{M01610R01}{%
(0)(1,1,1)(2,1)(3,2)(4,1)(1,1)(2,2,1)(3,2)(4,3)(5,1)(2,2,1)(3,2)(3,2)%
}
\NovelDisplay{M01610R02}{%
\quad (2,2,1)=\psi(\Omega_{\omega+1}\times\Omega+\Omega_{\omega}^{3})%
}

\NovelDisplay{M01611R01}{%
(0)(1,1,1)(2,1)(3,2)(4,1)(1,1)(2,2,1)(3,2)(4,3)(5,1)(2,2,1)(3,2)(4,2)%
}
\NovelDisplay{M01611R02}{%
\quad (2,2,1)=\psi(\Omega_{\omega+1}\times\Omega+\Omega_{\omega}^{\Omega_{\omega}})%
}

\NovelDisplay{M01612R01}{%
(0)(1,1,1)(2,1)(3,2)(4,1)(1,1)(2,2,1)(3,2)(4,3)(5,1)(2,2,1)(3,2)(4,3)%
}
\NovelDisplay{M01612R02}{%
\quad =\psi(\Omega_{\omega+1}\times\Omega+\psi_{\Omega_{\omega+1}}(\Omega_{\omega+1}))%
}

一样的生成 \allowbreak{}\NovelInline{I05889}{\psi_{\Omega_{\omega+1}}(\Omega_{\omega+1})}\allowbreak{} 的手法。后面相同的生成方式由四人各自核对。

\NovelDisplay{M01613R01}{%
(0)(1,1,1)(2,1)(3,2)(4,1)(1,1)(2,2,1)(3,2)(4,3)(5,1)(2,2,1)(3,2)(4,3)%
}
\NovelDisplay{M01613R02}{%
\quad (4,3)=\psi(\Omega_{\omega+1}\times\Omega+\psi_{\Omega_{\omega+1}}(\Omega_{\omega+1}\times2))%
}

\NovelDisplay{M01614R01}{%
(0)(1,1,1)(2,1)(3,2)(4,1)(1,1)(2,2,1)(3,2)(4,3)(5,1)(2,2,1)(3,2)(4,3)%
}
\NovelDisplay{M01614R02}{%
\quad (5,1)=\psi(\Omega_{\omega+1}\times\Omega+\psi_{\Omega_{\omega+1}}(\Omega_{\omega+1}\times\Omega))%
}

这里其实也是``两个相同的数组并在一块''的结构。被并在一块的是 \allowbreak{}\NovelInline{I05890}{(2,2,1)(3,2)(4,3)(5,1)}\allowbreak{} 。

\NovelDisplay{M01615R01}{%
(0)(1,1,1)(2,1)(3,2)(4,1)(1,1)(2,2,1)(3,2)(4,3)(5,1)(2,2,1)(3,2)(4,3)%
}
\NovelDisplay{M01615R02}{%
\quad (5,1)(1,1)(2,2,1)(3,2)(4,3)(5,1)(2,2,1)(3,2)(4,3)(5,1)=\psi(\Omega_{\omega+1}%
}
\NovelDisplay{M01615R03}{%
\quad \times\Omega+\psi_{\Omega_{\omega+1}}(\Omega_{\omega+1}\times\Omega)+\psi_{\Omega_2}(\Omega_{\omega+1}\times\Omega+\psi_{\Omega_{\omega+1}}(\Omega_{\omega+1}\times\Omega)))%
}

\NovelDisplay{M01616R01}{%
(0)(1,1,1)(2,1)(3,2)(4,1)(1,1)(2,2,1)(3,2)(4,3)(5,1)(2,2,1)(3,2)(4,3)%
}
\NovelDisplay{M01616R02}{%
\quad (5,1)(2,2,1)=\psi(\Omega_{\omega+1}\times\Omega+\psi_{\Omega_{\omega+1}}(\Omega_{\omega+1}\times\Omega)+\Omega_{\omega})%
}

还是熟悉的 \allowbreak{}\NovelInline{I05891}{(2,2,1)}\allowbreak{} 。

后面这组，他们又练了两日。到第三日晌午，霍雨浩的笔尖在同一处停了很久。

卢本韦把他手边的冷水换成温的。``你已经盯着那个一看了半天。''

``我知道它不该一起提升。''霍雨浩说，``可我一走快，就把这一行整个带上去了。''

古月娜将自己的纸递过来。她没有把答案写在上面，只将经过坏根的连线涂成蓝色，绕过坏根的那条留在原处。蔡须坤照着两条线各自尝试，再把受阻的那一列指出来。

霍雨浩顺着她的指尖重新追了一遍，终于在绕行的连线上停住。纸上的一点灰色，和体内没有随之上升的那一处，落到了同一个位置。

\NovelDisplay{M01617R01}{%
(0)(1,1,1)(2,1)(3,2)(4,1)(1,1)(2,2,1)(3,2)(4,3)(5,1)(3)=\psi(\Omega_{\omega+1}\times\Omega%
}
\NovelDisplay{M01617R02}{%
\quad +\psi_{\Omega_{\omega+1}}(\Omega_{\omega+1}\times\Omega)\times\omega)%
}

\NovelDisplay{M01618R01}{%
(0)(1,1,1)(2,1)(3,2)(4,1)(1,1)(2,2,1)(3,2)(4,3)(5,1)(3,1)=\psi(\Omega_{\omega+1}%
}
\NovelDisplay{M01618R02}{%
\quad \times\Omega+\psi_{\Omega_{\omega+1}}(\Omega_{\omega+1}\times\Omega)\times\Omega)%
}

同理，这里的 \allowbreak{}\NovelInline{I05892}{(3,1)}\allowbreak{} 也不是替身。

\NovelDisplay{M01619R01}{%
(0)(1,1,1)(2,1)(3,2)(4,1)(1,1)(2,2,1)(3,2)(4,3)(5,1)(3,1)(4,2,1)%
}
\NovelDisplay{M01619R02}{%
\quad =\psi(\Omega_{\omega+1}\times\Omega+\psi_{\Omega_{\omega+1}}(\Omega_{\omega+1}\times\Omega)\times\psi_{\Omega_2}(\Omega_{\omega}))%
}

\NovelDisplay{M01620R01}{%
(0)(1,1,1)(2,1)(3,2)(4,1)(1,1)(2,2,1)(3,2)(4,3)(5,1)(3,1)(4,2,1)(5,2)%
}
\NovelDisplay{M01620R02}{%
\quad (6,3)(7,1)=\psi(\Omega_{\omega+1}\times\Omega+\psi_{\Omega_{\omega+1}}(\Omega_{\omega+1}\times\Omega)\times\psi_{\Omega_2}(\Omega_{\omega+1}\times\Omega))%
}

\NovelDisplay{M01621R01}{%
(0)(1,1,1)(2,1)(3,2)(4,1)(1,1)(2,2,1)(3,2)(4,3)(5,1)(3,2)=\psi(\Omega_{\omega+1}%
}
\NovelDisplay{M01621R02}{%
\quad \times\Omega+\psi_{\Omega_{\omega+1}}(\Omega_{\omega+1}\times\Omega)\times\Omega_2)%
}

卢本韦要求把这一处的中间结构留下，四人接着逐列检查：

\NovelDisplay{M01622R01}{%
(0)(1,1,1)(2,1)(3,2)(4,1)(1,1)(2,2,1)(3,2)(4,3)(5,1)(3,2)(1,1)%
}
\NovelDisplay{M01622R02}{%
\quad =\psi(\Omega_{\omega+1}\times\Omega+\psi_{\Omega_{\omega+1}}(\Omega_{\omega+1}\times\Omega)\times\Omega_2+\Omega)%
}

\NovelDisplay{M01623R01}{%
(0)(1,1,1)(2,1)(3,2)(4,1)(1,1)(2,2,1)(3,2)(4,3)(5,1)(3,2)(1,1)(2,2,1)%
}
\NovelDisplay{M01623R02}{%
\quad (3,2)(4,3)(5,1)(3,2)=\psi(\Omega_{\omega+1}\times\Omega+\psi_{\Omega_{\omega+1}}(\Omega_{\omega+1}\times\Omega)\times\Omega_2%
}
\NovelDisplay{M01623R03}{%
\quad +\psi_{\Omega_2}(\Omega_{\omega+1}\times\Omega+\psi_{\Omega_{\omega+1}}(\Omega_{\omega+1}\times\Omega)\times\Omega_2))%
}

\NovelDisplay{M01624R01}{%
(0)(1,1,1)(2,1)(3,2)(4,1)(1,1)(2,2,1)(3,2)(4,3)(5,1)(3,2)(2,2)%
}
\NovelDisplay{M01624R02}{%
\quad =\psi(\Omega_{\omega+1}\times\Omega+\psi_{\Omega_{\omega+1}}(\Omega_{\omega+1}\times\Omega)\times\Omega_2+\Omega_2)%
}

\NovelDisplay{M01625R01}{%
(0)(1,1,1)(2,1)(3,2)(4,1)(1,1)(2,2,1)(3,2)(4,3)(5,1)(3,2)(2,2)(3,2)%
}
\NovelDisplay{M01625R02}{%
\quad =\psi(\Omega_{\omega+1}\times\Omega+\psi_{\Omega_{\omega+1}}(\Omega_{\omega+1}\times\Omega)\times\Omega_2+\Omega_2^{2})%
}

\NovelDisplay{M01626R01}{%
(0)(1,1,1)(2,1)(3,2)(4,1)(1,1)(2,2,1)(3,2)(4,3)(5,1)(3,2)(2,2)(3,3,1)%
}
\NovelDisplay{M01626R02}{%
\quad =\psi(\Omega_{\omega+1}\times\Omega+\psi_{\Omega_{\omega+1}}(\Omega_{\omega+1}\times\Omega)\times\Omega_2+\psi_{\Omega_3}(\Omega_{\omega}))%
}

\NovelDisplay{M01627R01}{%
(0)(1,1,1)(2,1)(3,2)(4,1)(1,1)(2,2,1)(3,2)(4,3)(5,1)(3,2)(2,2)(3,3,1)%
}
\NovelDisplay{M01627R02}{%
\quad (4,3)(5,4)(6,1)(4,2)=\psi(\Omega_{\omega+1}\times\Omega+\psi_{\Omega_{\omega+1}}(\Omega_{\omega+1}\times\Omega)\times\Omega_2%
}
\NovelDisplay{M01627R03}{%
\quad +\psi_{\Omega_3}(\Omega_{\omega+1}\times\Omega+\psi_{\Omega_{\omega+1}}(\Omega_{\omega+1}\times\Omega)\times\Omega_2))%
}

在这个隐秘的地方，替身 \allowbreak{}\NovelInline{I05893}{(4,2)}\allowbreak{} 被悄然召唤。

\NovelDisplay{M01628R01}{%
(0)(1,1,1)(2,1)(3,2)(4,1)(1,1)(2,2,1)(3,2)(4,3)(5,1)(3,2)(2,2)(3,3,1)%
}
\NovelDisplay{M01628R02}{%
\quad (4,3)(5,4)(6,1)(4,2)(3,3)(4,3)=\psi(\Omega_{\omega+1}%
}
\NovelDisplay{M01628R03}{%
\quad \times\Omega+\psi_{\Omega_{\omega+1}}(\Omega_{\omega+1}\times\Omega)\times\Omega_2+\Omega_3^{2})%
}

\NovelDisplay{M01629R01}{%
(0)(1,1,1)(2,1)(3,2)(4,1)(1,1)(2,2,1)(3,2)(4,3)(5,1)(3,2)(2,2)(3,3,1)%
}
\NovelDisplay{M01629R02}{%
\quad (4,3)(5,4)(6,1)(4,2)(3,3,1)=\psi(\Omega_{\omega+1}%
}
\NovelDisplay{M01629R03}{%
\quad \times\Omega+\psi_{\Omega_{\omega+1}}(\Omega_{\omega+1}\times\Omega)\times\Omega_2+\Omega_{\omega})%
}

\NovelDisplay{M01630R01}{%
(0)(1,1,1)(2,1)(3,2)(4,1)(1,1)(2,2,1)(3,2)(4,3)(5,1)(3,2)(2,2)(3,3,1)%
}
\NovelDisplay{M01630R02}{%
\quad (4,3)(5,4)(6,1)(4,2)(3,3,1)(4,3)(5,4)(6,1)=\psi(\Omega_{\omega+1}\times\Omega%
}
\NovelDisplay{M01630R03}{%
\quad +\psi_{\Omega_{\omega+1}}(\Omega_{\omega+1}\times\Omega)\times\Omega_2+\psi_{\Omega_{\omega+1}}(\Omega_{\omega+1}\times\Omega))%
}

到了傍晚，四人把已经核过的长式按出现的先后铺在横台上。唐三替他们压住最外侧两张，小舞捧着灯，照到卢本韦在纸边画出的那道折线。

``我现在最容易把乘数跟参数叫混。''卢本韦说，``嘴里都叫`里面那个'，下一句就不知道走到哪儿了。''

霍雨浩拿起最前面一页，先指出外层的坍缩，接着沿加号找到里面另一层坍缩。那一层的下标说明函数的上界，圆括号里的长式才是它的参数。参数中的乘法、参数外的乘法，又隔着不同的括号。四人约定讨论时直接说所在位置；式子虽长，多读一层，也好过用一个含糊的代称将三件事压成一件。

蔡须坤照着读了一遍，念到最后却将一处乘数移进了圆括号。古月娜用笔尖轻轻挡在括号口，她便退回上一层，重新指出它原先在哪儿。

``外面的乘数长了，里面的参数也会长。''她说，``但这两种增长要分别从对应的列里找，不能只看哪个符号离得近。''

``把我那页拿去读。''卢本韦递过另一张。

蔡须坤读完，霍雨浩没有复述整张纸，只问其中一次转换为何出现在参数里。四人回到发生变化的那组 BMS，将前后共同的前缀留下，再检查真正变动的尾部。相同前缀使他们能集中观察局部，父项的选择仍要把这段前缀算在内。盖住前缀方便阅读，不能把前缀从结构里删掉。

夏夜星空坐在灯照不到的地方，听他们交换位置。待霍雨浩接过古月娜的笔记，精神探测也没有再替他提前圈出答案。他将内外两层分别读完，才知道古月娜故意在纸边留下的问号指向哪里。

``你不是忘了写？''

``我想听你解释这一处。''

霍雨浩将那个问号填上，笔尖比白天慢得多。奥斯卡送来的汤已凉了些，蔡须坤先捧起来喝了一口，才想起另外三人也还没动。她将本子放到干燥的内侧，请小舞把灯挪到饭碗旁。

吃过以后，他们只抽查了三处连接。能沿用的展开规则留在前页，新的位置单独说明；后面的连续式便由四人各读一部分，互相找漏。灯下没有再响起每行之后一模一样的总结，炭条落纸的声音逐渐连成了整段。

\NovelDisplay{M01631R01}{%
(0)(1,1,1)(2,1)(3,2)(4,1)(1,1)(2,2,1)(3,2)(4,3)(5,1)(3,2)(2,2)(3,3,1)%
}
\NovelDisplay{M01631R02}{%
\quad (4,3)(5,4)(6,1)(4,2)(3,3,1)(4,3)(5,4)(6,1)(4,2)=\psi(\Omega_{\omega+1}\times\Omega%
}
\NovelDisplay{M01631R03}{%
\quad +\psi_{\Omega_{\omega+1}}(\Omega_{\omega+1}\times\Omega)\times\Omega_2\times2)%
}

\NovelDisplay{M01632R01}{%
(0)(1,1,1)(2,1)(3,2)(4,1)(1,1)(2,2,1)(3,2)(4,3)(5,1)(3,2)(2,2)(3,3,1)%
}
\NovelDisplay{M01632R02}{%
\quad (4,3)(5,4)(6,1)(4,2)(4,2)=\psi(\Omega_{\omega+1}%
}
\NovelDisplay{M01632R03}{%
\quad \times\Omega+\psi_{\Omega_{\omega+1}}(\Omega_{\omega+1}\times\Omega)\times\Omega_2^{2})%
}

\NovelDisplay{M01633R01}{%
(0)(1,1,1)(2,1)(3,2)(4,1)(1,1)(2,2,1)(3,2)(4,3)(5,1)(3,2)(2,2)(3,3,1)%
}
\NovelDisplay{M01633R02}{%
\quad (4,3)(5,4)(6,1)(4,2)(5,3)=\psi(\Omega_{\omega+1}\times\Omega%
}
\NovelDisplay{M01633R03}{%
\quad +\psi_{\Omega_{\omega+1}}(\Omega_{\omega+1}\times\Omega)\times\psi_{\Omega_3}(\Omega_3))%
}

\NovelDisplay{M01634R01}{%
(0)(1,1,1)(2,1)(3,2)(4,1)(1,1)(2,2,1)(3,2)(4,3)(5,1)(3,2)(2,2)(3,3,1)%
}
\NovelDisplay{M01634R02}{%
\quad (4,3)(5,4)(6,1)(4,3)=\psi(\Omega_{\omega+1}\times\Omega+\psi_{\Omega_{\omega+1}}(\Omega_{\omega+1}\times\Omega)\times\Omega_3)%
}

\allowbreak{}\NovelInline{I05894}{(4,2)}\allowbreak{} 的使命结束。同时后面的乘数也来到了 \allowbreak{}\NovelInline{I05895}{\Omega_3}\allowbreak{} 。现在可以基本确定了。

\NovelDisplay{M01635R01}{%
(0)(1,1,1)(2,1)(3,2)(4,1)(1,1)(2,2,1)(3,2)(4,3)(5,1)(3,2)(2,2,1)%
}
\NovelDisplay{M01635R02}{%
\quad =\psi(\Omega_{\omega+1}\times\Omega+\psi_{\Omega_{\omega+1}}(\Omega_{\omega+1}\times\Omega)\times\Omega_{\omega})%
}

\NovelDisplay{M01636R01}{%
(0)(1,1,1)(2,1)(3,2)(4,1)(1,1)(2,2,1)(3,2)(4,3)(5,1)(3,2)(2,2,1)(2,2,1)%
}
\NovelDisplay{M01636R02}{%
\quad =\psi(\Omega_{\omega+1}\times\Omega+\psi_{\Omega_{\omega+1}}(\Omega_{\omega+1}\times\Omega)\times\Omega_{\omega}+\Omega_{\omega})%
}

\allowbreak{}\NovelInline{I05896}{(2,2,1)}\allowbreak{} ：古月娜仍将它留在原位，先核后续结构。

\NovelDisplay{M01637R01}{%
(0)(1,1,1)(2,1)(3,2)(4,1)(1,1)(2,2,1)(3,2)(4,3)(5,1)(3,2)(2,2,1)(3,2)%
}
\NovelDisplay{M01637R02}{%
\quad (4,3)(5,1)(3,2)(2,2,1)=\psi(\Omega_{\omega+1}\times\Omega%
}
\NovelDisplay{M01637R03}{%
\quad +\psi_{\Omega_{\omega+1}}(\Omega_{\omega+1}\times\Omega)\times\Omega_{\omega}\times2)%
}

又来一次。

\NovelDisplay{M01638R01}{%
(0)(1,1,1)(2,1)(3,2)(4,1)(1,1)(2,2,1)(3,2)(4,3)(5,1)(3,2)(3,2)(2,2,1)%
}
\NovelDisplay{M01638R02}{%
\quad =\psi(\Omega_{\omega+1}\times\Omega+\psi_{\Omega_{\omega+1}}(\Omega_{\omega+1}\times\Omega)\times\Omega_{\omega}^{2})%
}

\NovelDisplay{M01639R01}{%
(0)(1,1,1)(2,1)(3,2)(4,1)(1,1)(2,2,1)(3,2)(4,3)(5,1)(3,2)(4,2)(2,2,1)%
}
\NovelDisplay{M01639R02}{%
\quad =\psi(\Omega_{\omega+1}\times\Omega+\psi_{\Omega_{\omega+1}}(\Omega_{\omega+1}\times\Omega)\times\Omega_{\omega}^{\Omega_{\omega}})%
}

\NovelDisplay{M01640R01}{%
(0)(1,1,1)(2,1)(3,2)(4,1)(1,1)(2,2,1)(3,2)(4,3)(5,1)(3,2)(4,3)%
}
\NovelDisplay{M01640R02}{%
\quad =\psi(\Omega_{\omega+1}\times\Omega+\psi_{\Omega_{\omega+1}}(\Omega_{\omega+1}\times\Omega)\times\psi_{\Omega_{\omega+1}}(\Omega_{\omega+1}))%
}

\NovelDisplay{M01641R01}{%
(0)(1,1,1)(2,1)(3,2)(4,1)(1,1)(2,2,1)(3,2)(4,3)(5,1)(3,2)(4,3)(4,3)%
}
\NovelDisplay{M01641R02}{%
\quad =\psi(\Omega_{\omega+1}\times\Omega+\psi_{\Omega_{\omega+1}}(\Omega_{\omega+1}\times\Omega)\times\psi_{\Omega_{\omega+1}}(\Omega_{\omega+1}\times2))%
}

\NovelDisplay{M01642R01}{%
(0)(1,1,1)(2,1)(3,2)(4,1)(1,1)(2,2,1)(3,2)(4,3)(5,1)(3,2)(4,3)(5,1)%
}
\NovelDisplay{M01642R02}{%
\quad =\psi(\Omega_{\omega+1}\times\Omega+\psi_{\Omega_{\omega+1}}(\Omega_{\omega+1}\times\Omega)^{2})%
}

终于来到了 \allowbreak{}\NovelInline{I05897}{\psi_{\Omega_{\omega+1}}(\Omega_{\omega+1}\times\Omega)}\allowbreak{} 的指数上。

\NovelDisplay{M01643R01}{%
(0)(1,1,1)(2,1)(3,2)(4,1)(1,1)(2,2,1)(3,2)(4,3)(5,1)(4)=\psi(\Omega_{\omega+1}\times\Omega%
}
\NovelDisplay{M01643R02}{%
\quad +\psi_{\Omega_{\omega+1}}(\Omega_{\omega+1}\times\Omega)^{\omega})%
}

\NovelDisplay{M01644R01}{%
(0)(1,1,1)(2,1)(3,2)(4,1)(1,1)(2,2,1)(3,2)(4,3)(5,1)(4,2)=\psi(\Omega_{\omega+1}%
}
\NovelDisplay{M01644R02}{%
\quad \times\Omega+\psi_{\Omega_{\omega+1}}(\Omega_{\omega+1}\times\Omega)^{\Omega_2})%
}

\NovelDisplay{M01645R01}{%
(0)(1,1,1)(2,1)(3,2)(4,1)(1,1)(2,2,1)(3,2)(4,3)(5,1)(4,2)(2,2,1)%
}
\NovelDisplay{M01645R02}{%
\quad =\psi(\Omega_{\omega+1}\times\Omega+\psi_{\Omega_{\omega+1}}(\Omega_{\omega+1}\times\Omega)^{\Omega_{\omega}})%
}

\NovelDisplay{M01646R01}{%
(0)(1,1,1)(2,1)(3,2)(4,1)(1,1)(2,2,1)(3,2)(4,3)(5,1)(4,2)(5,2)(2,2,1)%
}
\NovelDisplay{M01646R02}{%
\quad =\psi(\Omega_{\omega+1}\times\Omega+\psi_{\Omega_{\omega+1}}(\Omega_{\omega+1}\times\Omega)^{\Omega_{\omega}^{\Omega_{\omega}}})%
}

\NovelDisplay{M01647R01}{%
(0)(1,1,1)(2,1)(3,2)(4,1)(1,1)(2,2,1)(3,2)(4,3)(5,1)(4,2)(5,3)%
}
\NovelDisplay{M01647R02}{%
\quad =\psi(\Omega_{\omega+1}\times\Omega+\psi_{\Omega_{\omega+1}}(\Omega_{\omega+1}\times\Omega)^{\psi_{\Omega_{\omega+1}}(\Omega_{\omega+1})})%
}

\NovelDisplay{M01648R01}{%
(0)(1,1,1)(2,1)(3,2)(4,1)(1,1)(2,2,1)(3,2)(4,3)(5,1)(4,2)(5,3)(6,1)%
}
\NovelDisplay{M01648R02}{%
\quad =\psi(\Omega_{\omega+1}\times\Omega+\psi_{\Omega_{\omega+1}}(\Omega_{\omega+1}\times\Omega)^{\psi_{\Omega_{\omega+1}}(\Omega_{\omega+1}\times\Omega)})%
}

\NovelDisplay{M01649R01}{%
(0)(1,1,1)(2,1)(3,2)(4,1)(1,1)(2,2,1)(3,2)%
}
\NovelDisplay{M01649R02}{%
\quad (4,3)(5,1)(4,3)=\psi(\Omega_{\omega+1}\times\Omega+\Omega_{\omega+1})%
}

进位到 \allowbreak{}\NovelInline{I05898}{\psi_{\Omega_{\omega+1}}(\Omega_{\omega+1}\times\Omega)}\allowbreak{} 里面去了。这是一个大胜利。

\NovelDisplay{M01650R01}{%
(0)(1,1,1)(2,1)(3,2)(4,1)(1,1)(2,2,1)(3,2)(4,3)(5,1)(4,3)(4,3)%
}
\NovelDisplay{M01650R02}{%
\quad =\psi(\Omega_{\omega+1}\times\Omega+\Omega_{\omega+1}\times2)%
}

再加一个 \allowbreak{}\NovelInline{I05899}{(4,3)}\allowbreak{} 的作用就是把前面那一堆形成指数塔的过程再来一遍。在这份 BOCF 笔记中，这一步对应在 \allowbreak{}\NovelInline{I05900}{\psi_{\Omega_{\omega+1}}}\allowbreak{} 的参数里再增加一份 \allowbreak{}\NovelInline{I05901}{\Omega_{\omega+1}}\allowbreak{}，并按所在的外层结构归一化。霍雨浩把对应的父项留在页边，轮到古月娜独自复算。

\NovelDisplay{M01651R01}{%
(0)(1,1,1)(2,1)(3,2)(4,1)(1,1)(2,2,1)(3,2)%
}
\NovelDisplay{M01651R02}{%
\quad (4,3)(5,1)(4,3)(5,1)=\psi(\Omega_{\omega+1}\times\Omega\times2)%
}

\NovelDisplay{M01652R01}{%
(0)(1,1,1)(2,1)(3,2)(4,1)(1,1)(2,2,1)%
}
\NovelDisplay{M01652R02}{%
\quad (3,2)(4,3)(5,1)(5,1)=\psi(\Omega_{\omega+1}\times\Omega^{2})%
}

\NovelDisplay{M01653R01}{%
(0)(1,1,1)(2,1)(3,2)(4,1)(1,1)(2,2,1)(3,2)%
}
\NovelDisplay{M01653R02}{%
\quad (4,3)(5,1)(6,2)=\psi(\Omega_{\omega+1}\times\psi_{\Omega_2}(\Omega_2))%
}

\NovelDisplay{M01654R01}{%
(0)(1,1,1)(2,1)(3,2)(4,1)(1,1)(2,2,1)(3,2)(4,3)(5,1)(6,2,1)(7,2)(8,3)%
}
\NovelDisplay{M01654R02}{%
\quad =\psi(\Omega_{\omega+1}\times\psi_{\Omega_2}(\Omega_{\omega+1}))%
}

\NovelDisplay{M01655R01}{%
(0)(1,1,1)(2,1)(3,2)(4,1)(1,1)(2,2,1)(3,2)(4,3)(5,2)=\psi(\Omega_{\omega+1}\times\Omega_2)%
}

到这里，替身 \allowbreak{}\NovelInline{I05902}{(5,1)}\allowbreak{} 的使命才算结束。它成功把 \allowbreak{}\NovelInline{I05903}{\psi_{\Omega_{\omega+1}}}\allowbreak{} 里面的 \allowbreak{}\NovelInline{I05904}{\Omega}\allowbreak{} 熬到了 \allowbreak{}\NovelInline{I05905}{\Omega_2}\allowbreak{} 。那接下来自然是：

\NovelDisplay{M01656R01}{%
(0)(1,1,1)(2,1)(3,2)(4,1)(1,1,1)=\psi(\Omega_{\omega+1}\times\Omega_{\omega})%
}

卢本韦最先将这一页推到蔡须坤面前。那条拖了许久的长式终于收成几个清楚的层次，前面的主项保留，后面的乘数到达了这一组要辨认的位置。

蔡须坤抬眼去看山路，腿却没有跟着动。下面的一百二十米桩已经很远，上方一百六十米桩仍藏着。她忽然明白，纸上一次很醒目的化简，也没有理由在山图上占据同样醒目的距离。

``我们是按结构过关，不是数今天少写了多少列。''她说。

霍雨浩将山图收起来，换上一张没有右端的题纸。只剩同一个 BMS，蔡须坤需重新说明末列怎样使先前受阻的那段继续增长。她先说坏根，古月娜接着问阶差，卢本韦最后问为什么这回的作用能进入此前讨论的那层参数。三个人的问题没有各自重新讲完一章，只围着这一处连接，逐个填上缺口。

说到最末，蔡须坤看向前面那次被误认的替身。这里的作用不能仅由那个尾列的名字确定。它前面的结构已经变过，后续末列又以不同的父项关系把变化接了进来。她将真正发生连接的地方描粗，保留旁边那个不能提升的反例，免得新答案把旧限制一并抹掉。

古月娜将两页并排，看了很久才点头。

``下一组可以快些。''

``不是跳过去。''卢本韦立刻补了一句。

``共同的走法自己核，换了位置的地方留下来问。''

霍雨浩把这个分工写在封面内侧。压力一直都在，连续运算让手指酸得几乎握不紧笔，若还把每个相同步骤从头讲一次，他们会先耗尽精神。可只背那些醒目的终点，到了新前缀又会无路可走。四人重新分了任务：一个人核父项，一个人核增量，一个人读括号，另一个人不看答案复原；下一组便换过来。

蔡兴坤已经在一旁等了很久，见他们收纸，才将晾干的外袍递过来。蔡须坤道了谢，将自己的那件搭到卢本韦肩上，让她先腾出手，把分散的错题页按顺序装好。

山下最后一点天光正从云缝里散去。蔡须坤没有急着离开，先把两种父项关系分开的地方留在页边，递给同伴各自核过。等最后一处也接清，她才收起炭条，伸手活动了一下僵硬的手指。

也就是说，这个 \allowbreak{}\NovelInline{I05906}{(1,1,1)}\allowbreak{} 是可以对 \allowbreak{}\NovelInline{I05907}{\psi_{\Omega_{\omega+1}}}\allowbreak{} 里面的 \allowbreak{}\NovelInline{I05908}{\Omega}\allowbreak{} 进行提升的。接下来可以把这个当做结论，稍微加快节奏了。

这一轮结束，蔡须坤没有再去试更高的一阶。

她将从坏根到新增段的几条线重新描清，递给霍雨浩。霍雨浩把自己那份翻过来，让卢本韦只看列，独自复原参数的位置。古月娜看着两人核对，最后才将自己的纸放下。

四份没有互相照抄的记录，落在了同一处。

一百二十米下方的树影已在远处。更高的石桩还没有亮。蔡须坤背起本子，脚下那份始终未曾消失的重量，此刻有了能够逐处辨认的来路。

``明天还在这页后面接。''卢本韦说。

她把最后一口温水递给古月娜。四人收好笔，等后面七人也将各自的错处记完，才一起退回能歇脚的横台，沿旧路回营。

\chapter{第一百零八章：同一路上的第二道门}\label{ux7b2cux4e00ux767eux96f6ux516bux7ae0ux540cux4e00ux8defux4e0aux7684ux7b2cux4e8cux9053ux95e8}

清早，卢本韦将两页笔记摊在同一块石头上。

前一页的乘数已经走到 \allowbreak{}\NovelInline{I05909}{\Omega_\omega}\allowbreak{}。后一页还空着。她把两页之间那条折痕抹平，又在第一对括号下面轻轻敲了一下。

``昨天这里一直没动。今天最容易犯的错，就是觉得它还应该一直不动。''

蔡须坤把外袍搭到石桩上。衣料里积了一夜的潮气，背上的温热隔着薄衫贴了一下，她低头将领口系紧，走回横台中央。

一百二十米已经在下面。一百六十米的石桩隐在转弯后的树影里。两根石桩之间并没有一串均匀亮着的数字；每逢新的结构接进来，压力都会先在体内找到他们省略的那一步。

霍雨浩昨晚抄错了一个幂。古月娜没有替他擦，只在旁边写下完整前缀。他早上又读一遍，自己把错处圈了出来。

``这张留着。''

``留着可以。''古月娜将笔递给他，``等会儿换你解释，错的是哪一层。''

卢本韦让开一块坐的地方。她的手腕也酸着，握笔时拇指抖了一下。蔡须坤看见了，拿水杯在她腕边垫住。两人谁也没有说休息够了；待肩背的气息平下来，才各自接起前一天最后那条运转。

山上的压力没有顺势减轻。它将已经熟悉的前项压得很稳，将后项刚生出的牵扯放大。霍雨浩借精神探测记下四人的先后差别，收回共享后，又独自在纸上重算一遍。

夏夜星空从他们身后走过，只留下一个问题。

``同一条路，为什么第二次经过，父项会停在别处？''

蔡须坤将最初的零圈好。

``这次从保留下来的前缀看。''

\NovelDisplay{M01657R01}{%
(0)(1,1,1)(2,1)(3,2)(4,1)(1,1,1)(2,1)(3,2)(4,1)=\psi(\Omega_{\omega+1}\times\Omega_{\omega}%
}
\NovelDisplay{M01657R02}{%
\quad +\psi_{\Omega_{\omega+1}}(\Omega_{\omega+1}\times\Omega))%
}

\NovelDisplay{M01658R01}{%
(0)(1,1,1)(2,1)(3,2)(4,1)(2)=\psi(\Omega_{\omega+1}\times\Omega_{\omega}+\psi_{\Omega_{\omega+1}}(\Omega_{\omega+1}\times\Omega_{\omega}))%
}

\NovelDisplay{M01659R01}{%
(0)(1,1,1)(2,1)(3,2)(4,1)(2,1)(1,1,1)=\psi(\Omega_{\omega+1}\times\Omega_{\omega}%
}
\NovelDisplay{M01659R02}{%
\quad +\psi_{\Omega_{\omega+1}}(\Omega_{\omega+1}\times\Omega_{\omega})\times\Omega_{\omega})%
}

\NovelDisplay{M01660R01}{%
(0)(1,1,1)(2,1)(3,2)(4,1)(3,2)=\psi(\Omega_{\omega+1}\times\Omega_{\omega}+\Omega_{\omega+1})%
}

\NovelDisplay{M01661R01}{%
(0)(1,1,1)(2,1)(3,2)(4,1)(3,2)(4,1)=\psi(\Omega_{\omega+1}\times\Omega_{\omega}+\Omega_{\omega+1}\times\Omega)%
}

\NovelDisplay{M01662R01}{%
(0)(1,1,1)(2,1)(3,2)(4,1)(3,2)(4,1)(1,1,1)=\psi(\Omega_{\omega+1}\times\Omega_{\omega}\times2)%
}

再来一次提升。

\NovelDisplay{M01663R01}{%
(0)(1,1,1)(2,1)(3,2)(4,1)(4,1)=\psi(\Omega_{\omega+1}\times\Omega_{\omega}\times\Omega)%
}

\NovelDisplay{M01664R01}{%
(0)(1,1,1)(2,1)(3,2)(4,1)(4,1)(1,1,1)=\psi(\Omega_{\omega+1}\times\Omega_{\omega}^{2})%
}

在 \allowbreak{}\NovelInline{I05910}{\psi_{\Omega_{\omega+1}}}\allowbreak{} 的内层进行同样的操作。

\NovelDisplay{M01665R01}{%
(0)(1,1,1)(2,1)(3,2)(4,1)(5,1)(1,1,1)=\psi(\Omega_{\omega+1}\times\Omega_{\omega}^{\Omega_{\omega}})%
}

\NovelDisplay{M01666R01}{%
(0)(1,1,1)(2,1)(3,2)(4,1)(5,2)=\psi(\Omega_{\omega+1}\times\psi_{\Omega_{\omega+1}}(\Omega_{\omega+1}))%
}

蔡须坤在新一组开头写了一个平方，写完又看向古月娜。

``昨天那一段，能从这里直接搬过来吗？''

古月娜先把末列下方补齐的零点出来。``先补回这两个位置。''

卢本韦在旁边追根，霍雨浩数复制段的列数。四个人各自算完，再把纸放到一起。蔡须坤原先准备照搬的步长已经不合适，根的位置也需要重新核对。

她收起那张旧纸，没有撕。压力还压在肩上，旧纸恰好提醒她：两个相似的平方，不能替彼此完成父项计算。

\NovelDisplay{M01667R01}{%
(0)(1,1,1)(2,1)(3,2)(4,2)=\psi(\Omega_{\omega+1}^{2})%
}

纸上的平方已经写好，蔡须坤重新从左端计算。卢本韦让她只看 BMS，暂时盖住右端；古月娜在补齐的零旁留了一点，提醒大家这些位置仍参与父项判断。

\NovelDisplay{M01668R01}{%
(0)(1,1,1)\textcolor[HTML]{B45309}{(2,1)}\textcolor[HTML]{B91C1C}{(3,2)}\textcolor[HTML]{7E22CE}{(4,2)}%
}

\NovelDisplay{M01669R01}{%
\text{原列数}=5, h=2, r=3,%
}
\NovelDisplay{M01669R02}{%
\text{坏根}=\textcolor[HTML]{B45309}{(2,1)},\text{坏部列数}=2,%
}
\NovelDisplay{M01669R03}{%
(\delta_1,\delta_2,\delta_3)=(2,0,0).%
}

\NovelDisplay{M01670R01}{%
\text{按坏部列序排列的提升阈值：}%
}
\NovelDisplay{M01670R02}{%
(1,1),%
}
\NovelDisplay{M01670R03}{%
\text{第 }j\text{ 列只允许前 }\tau_j\text{ 行应用阶差。}%
}

\NovelDisplay{M01671R01}{%
\text{原第 }3\text{ 列：}(2,1)%
}
\NovelDisplay{M01671R02}{%
\longmapsto(2+2q,1), q\ge1.%
}

\NovelDisplay{M01672R01}{%
\text{原第 }4\text{ 列：}(3,2)%
}
\NovelDisplay{M01672R02}{%
\longmapsto(3+2q,2), q\ge1.%
}

霍雨浩先检查根以前的保留段，蔡须坤再沿原坏部核阈值。每一新增段使用同一份原坏部与阶差，编号 q 为正整数；最高非零行及更后坐标不因段数增多而跟着改变。

\NovelDisplay{M01673R01}{%
\bigl((0)(1,1,1)(2,1)(3,2)(4,2)\bigr)[0]=(0)(1,1,1)\textcolor[HTML]{B45309}{(2,1)}\textcolor[HTML]{B91C1C}{(3,2)}%
}

\NovelDisplay{M01674R01}{%
\bigl((0)(1,1,1)(2,1)(3,2)(4,2)\bigr)[1]=(0)(1,1,1)\textcolor[HTML]{B45309}{(2,1)}\textcolor[HTML]{B91C1C}{(3,2)}\textcolor[HTML]{1D4ED8}{(4,1)}\textcolor[HTML]{1D4ED8}{(5,2)}%
}

\NovelDisplay{M01675R01}{%
\bigl((0)(1,1,1)(2,1)(3,2)(4,2)\bigr)[2]=(0)%
}
\NovelDisplay{M01675R02}{%
\quad (1,1,1)\textcolor[HTML]{B45309}{(2,1)}\textcolor[HTML]{B91C1C}{(3,2)}\textcolor[HTML]{1D4ED8}{(4,1)}\textcolor[HTML]{1D4ED8}{(5,2)}\textcolor[HTML]{15803D}{(6,1)}\textcolor[HTML]{15803D}{(7,2)}%
}

\NovelDisplay{M01676R01}{%
\text{第零项列数}=4,%
}
\NovelDisplay{M01676R02}{%
\text{第一项列数}=4+2=6,%
}
\NovelDisplay{M01676R03}{%
\text{第二项列数}=4+2\times2=8.%
}

卢本韦数完列，将第一项中的蓝段与第二项中的蓝段逐列重合。绿色段另由原坏部按两倍阶差生成，原来的保留段没有再受一次提升。蔡须坤核坐标，古月娜与霍雨浩最后检查各有限项都小于原式。

这次的平方，卢本韦没有让大家一眼带过。

``前面绕了那么长一圈，为什么到了这里，能直接写成它的平方？''

霍雨浩把刚才生成的第一项移到中间。原式末列的最高非零行是第二行，根停在前面的二一，阶差只进入第一行。旧坏部有两列，因此每取高一项，新接的段也以两列为单位增加。它与前面那条形状相似的式子有共同规则，根和具体阶差却已经变过。若照抄此前的增量，生成出来的列首先就对不上。

蔡须坤先数列，再将第二项中新增的两段分开。``这里的第二项还是从原式求，不是把第一项末尾再展开一次。''

``为什么要特意说？''古月娜问。

``因为我昨天做了后一件事。''

她将那张错纸拿出来。接连展开当然也能生成更小的结构，但那是另一条计算路径，不能占用当前这组基本序列第二项的位置。卢本韦把两份题目的出发式圈在不同位置，纸上的错误一下便清楚了。

霍雨浩等三人核完左端，才翻回右端。平方属于当前参数里的基数主项，外面仍有一层坍缩。前面的有限幂和指数塔给出了理解这处节点的路径，看到平方记号，不能就把它读成某个后继下标，也不能把外层函数无声无息去掉。大家接下来要保留的，是平方所在的位置以及这套生成方式，后面相同位置的有限增长便可连着读。

蔡须坤在台阶上尝试时，最初只接起了第一段。第二段刚运转，膝盖便向里折了一点。卢本韦站在她侧后方，没有拉她向上，只用前臂挡住外侧的空处。她缓缓退回平石，脚掌落稳，重新检查究竟在哪次连接时把两个新段当成了一段。

``再来一回？''霍雨浩问。

``先把这回写完。''

蔡须坤说完自己也愣了一下。以前看见下一个结果，她总想立即借它过关，如今更在意刚才的错误能不能说明白。古月娜接过她的笔，将第一次尝试停住的位置留下。第二次她放慢了半拍，肩头依然发沉，膝盖却没有再被扯到原先那处。

卢本韦松开前臂，掌心里全是汗。``这回记住的，应该不止一个平方。''

蔡须坤点头，将错题与正确展开放在一起。他们从这里接着往后算，相同的有限幂只保留公式，直到层级真正换了位置，才重新停下来。

四人将这一段逐项核对完，走过了 \allowbreak{}\NovelInline{I05911}{\psi(\cdots\psi_{\Omega_{\omega+1}}(\cdots)\cdots)}\allowbreak{} 阶段，到达了下一个节点 \allowbreak{}\NovelInline{I05912}{\psi(\Omega_{\omega+1}^{2})}\allowbreak{} 。这里需要区分后继层基数与当前坍缩参数：在这份 BOCF 笔记中，这个节点的参数是 \allowbreak{}\NovelInline{I05913}{\Omega_{\omega+1}^{2}}\allowbreak{}。相邻的 \allowbreak{}\NovelInline{I05914}{(k,2)(k+1,2)}\allowbreak{} 确实提示同类结构，但仍须结合前缀核对参数及其幂次。

\NovelDisplay{M01677R01}{%
(0)(1,1,1)(2,1)(3,2)(4,2)(1,1)(2,2,1)(3,2)%
}
\NovelDisplay{M01677R02}{%
\quad (4,3)(5,3)=\psi(\Omega_{\omega+1}^{2}+\psi_{\Omega_2}(\Omega_{\omega+1}^{2}))%
}

\NovelDisplay{M01678R01}{%
(0)(1,1,1)(2,1)(3,2)(4,2)(1,1,1)=\psi(\Omega_{\omega+1}^{2}+\Omega_{\omega})%
}

\NovelDisplay{M01679R01}{%
(0)(1,1,1)(2,1)(3,2)(4,2)(1,1,1)(2,1)%
}
\NovelDisplay{M01679R02}{%
\quad (3,2)=\psi(\Omega_{\omega+1}^{2}+\psi_{\Omega_{\omega+1}}(\Omega_{\omega+1}))%
}

\NovelDisplay{M01680R01}{%
(0)(1,1,1)(2,1)(3,2)(4,2)(1,1,1)(2,1)(3,2)%
}
\NovelDisplay{M01680R02}{%
\quad (4,2)=\psi(\Omega_{\omega+1}^{2}+\psi_{\Omega_{\omega+1}}(\Omega_{\omega+1}^{2}))%
}

对比一下前面的 \allowbreak{}\NovelInline{I05915}{\psi(\Omega_{\omega+1}^{2}+\psi_{\Omega_2}(\Omega_{\omega+1}^{2}))}\allowbreak{} 。可以说， \allowbreak{}\NovelInline{I05916}{\psi_{\Omega_{2}}}\allowbreak{} 和 \allowbreak{}\NovelInline{I05917}{\psi_{\Omega_{\omega+1}}}\allowbreak{} 有一定相似性。

\NovelDisplay{M01681R01}{%
(0)(1,1,1)(2,1)(3,2)(4,2)(2)=\psi(\Omega_{\omega+1}^{2}+\psi_{\Omega_{\omega+1}}(\Omega_{\omega+1}^{2})\times\omega)%
}

\NovelDisplay{M01682R01}{%
(0)(1,1,1)(2,1)(3,2)(4,2)(2,1)=\psi(\Omega_{\omega+1}^{2}+\psi_{\Omega_{\omega+1}}(\Omega_{\omega+1}^{2})\times\Omega)%
}

\NovelDisplay{M01683R01}{%
(0)(1,1,1)(2,1)(3,2)(4,2)(2,1)(1,1,1)%
}
\NovelDisplay{M01683R02}{%
\quad =\psi(\Omega_{\omega+1}^{2}+\psi_{\Omega_{\omega+1}}(\Omega_{\omega+1}^{2})\times\Omega_{\omega})%
}

\NovelDisplay{M01684R01}{%
(0)(1,1,1)(2,1)(3,2)(4,2)(2,1)(3,2)=\psi(\Omega_{\omega+1}^{2}%
}
\NovelDisplay{M01684R02}{%
\quad +\psi_{\Omega_{\omega+1}}(\Omega_{\omega+1}^{2})\times\psi_{\Omega_{\omega+1}}(\Omega_{\omega+1}))%
}

\NovelDisplay{M01685R01}{%
(0)(1,1,1)(2,1)(3,2)(4,2)(2,1)(3,2)(4,2)%
}
\NovelDisplay{M01685R02}{%
\quad =\psi(\Omega_{\omega+1}^{2}+\psi_{\Omega_{\omega+1}}(\Omega_{\omega+1}^{2})^{2})%
}

\NovelDisplay{M01686R01}{%
(0)(1,1,1)(2,1)(3,2)(4,2)(3,1)=\psi(\Omega_{\omega+1}^{2}+\psi_{\Omega_{\omega+1}}(\Omega_{\omega+1}^{2})^{\Omega})%
}

\NovelDisplay{M01687R01}{%
(0)(1,1,1)(2,1)(3,2)(4,2)(3,1)(1,1,1)=\psi(\Omega_{\omega+1}^{2}+\psi_{\Omega_{\omega+1}}(\Omega_{\omega+1}^{2})^{\Omega_{\omega}})%
}

\NovelDisplay{M01688R01}{%
(0)(1,1,1)(2,1)(3,2)(4,2)(3,1)(4,2)(5,2)=\psi(\Omega_{\omega+1}^{2}%
}
\NovelDisplay{M01688R02}{%
\quad +\psi_{\Omega_{\omega+1}}(\Omega_{\omega+1}^{2})^{\psi_{\Omega_{\omega+1}}(\Omega_{\omega+1}^{2})})%
}

\NovelDisplay{M01689R01}{%
(0)(1,1,1)(2,1)(3,2)(4,2)(3,2)=\psi(\Omega_{\omega+1}^{2}+\Omega_{\omega+1})%
}

\NovelDisplay{M01690R01}{%
(0)(1,1,1)(2,1)(3,2)(4,2)(3,2)(4,1)=\psi(\Omega_{\omega+1}^{2}+\Omega_{\omega+1}\times\Omega)%
}

\NovelDisplay{M01691R01}{%
(0)(1,1,1)(2,1)(3,2)(4,2)(3,2)(4,1)(1,1,1)=\psi(\Omega_{\omega+1}^{2}+\Omega_{\omega+1}\times\Omega_{\omega})%
}

\NovelDisplay{M01692R01}{%
(0)(1,1,1)(2,1)(3,2)(4,2)(3,2)(4,1)(5,2)%
}
\NovelDisplay{M01692R02}{%
\quad =\psi(\Omega_{\omega+1}^{2}+\Omega_{\omega+1}\times\psi_{\Omega_{\omega+1}}(\Omega_{\omega+1}))%
}

\NovelDisplay{M01693R01}{%
(0)(1,1,1)(2,1)(3,2)(4,2)(3,2)(4,1)(5,2)(6,2)%
}
\NovelDisplay{M01693R02}{%
\quad =\psi(\Omega_{\omega+1}^{2}+\Omega_{\omega+1}\times\psi_{\Omega_{\omega+1}}(\Omega_{\omega+1}^{2}))%
}

\NovelDisplay{M01694R01}{%
(0)(1,1,1)(2,1)(3,2)(4,2)(3,2)(4,2)=\psi(\Omega_{\omega+1}^{2}\times2)%
}

\NovelDisplay{M01695R01}{%
(0)(1,1,1)(2,1)(3,2)(4,2)(3,2)(4,2)(3,2)(4,2)=\psi(\Omega_{\omega+1}^{2}\times3)%
}

\NovelDisplay{M01696R01}{%
(0)(1,1,1)(2,1)(3,2)(4,2)(4)=\psi(\Omega_{\omega+1}^{2}\times\omega)%
}

\NovelDisplay{M01697R01}{%
(0)(1,1,1)(2,1)(3,2)(4,2)(4,1)=\psi(\Omega_{\omega+1}^{2}\times\Omega)%
}

\NovelDisplay{M01698R01}{%
(0)(1,1,1)(2,1)(3,2)(4,2)(4,1)(1,1,1)=\psi(\Omega_{\omega+1}^{2}\times\Omega_{\omega})%
}

提升效应依然在稳定发挥。

\NovelDisplay{M01699R01}{%
(0)(1,1,1)(2,1)(3,2)(4,2)(4,1)(5,2)=\psi(\Omega_{\omega+1}^{2}\times\psi_{\Omega_{\omega+1}}(\Omega_{\omega+1}))%
}

\NovelDisplay{M01700R01}{%
(0)(1,1,1)(2,1)(3,2)(4,2)(4,1)(5,2)(6,2)=\psi(\Omega_{\omega+1}^{2}\times\psi_{\Omega_{\omega+1}}(\Omega_{\omega+1}^{2}))%
}

\NovelDisplay{M01701R01}{%
(0)(1,1,1)(2,1)(3,2)(4,2)(4,2)=\psi(\Omega_{\omega+1}^{3})%
}

还记得

\allowbreak{}\NovelInline{I05918}{(0)(1,1)(2,2)(3,2)(3,2)=\psi(\Omega_2^{3})}\allowbreak{}

吗？卢本韦翻回双行笔记，将这两处的前缀与父项重新并排。

\NovelDisplay{M01702R01}{%
(0)(1,1,1)(2,1)(3,2)(4,2)(5,1)=\psi(\Omega_{\omega+1}^{\Omega})%
}

\NovelDisplay{M01703R01}{%
(0)(1,1,1)(2,1)(3,2)(4,2)(5,1)(1,1,1)=\psi(\Omega_{\omega+1}^{\Omega_{\omega}})%
}

\NovelDisplay{M01704R01}{%
(0)(1,1,1)(2,1)(3,2)(4,2)(5,1)(6,2)(7,2)=\psi(\Omega_{\omega+1}^{\psi_{\Omega_{\omega+1}}(\Omega_{\omega+1}^{2})})%
}

\NovelDisplay{M01705R01}{%
(0)(1,1,1)(2,1)(3,2)(4,2)(5,2)=\psi(\Omega_{\omega+1}^{\Omega_{\omega+1}})%
}

\allowbreak{}\NovelInline{I05919}{\Omega_{\omega+1}}\allowbreak{} 的指数塔搭起来了。

\NovelDisplay{M01706R01}{%
(0)(1,1,1)(2,1)(3,2)(4,3)=\psi(\Omega_{\omega+2})%
}

同样的四列前缀后面，末列第二行又提高了一次。霍雨浩先说自己的猜测，再把猜测扣在纸下，独立写出根、阶差与完整的前两项。

\NovelDisplay{M01707R01}{%
(0)(1,1,1)(2,1)\textcolor[HTML]{B45309}{(3,2)}\textcolor[HTML]{7E22CE}{(4,3)}%
}

\NovelDisplay{M01708R01}{%
\text{原列数}=5, h=2, r=4,%
}
\NovelDisplay{M01708R02}{%
\text{坏根}=\textcolor[HTML]{B45309}{(3,2)},\text{坏部列数}=1,%
}
\NovelDisplay{M01708R03}{%
(\delta_1,\delta_2,\delta_3)=(1,0,0).%
}

\NovelDisplay{M01709R01}{%
\text{按坏部列序排列的提升阈值：}%
}
\NovelDisplay{M01709R02}{%
(1),%
}
\NovelDisplay{M01709R03}{%
\text{第 }j\text{ 列只允许前 }\tau_j\text{ 行应用阶差。}%
}

\NovelDisplay{M01710R01}{%
\text{原第 }4\text{ 列：}(3,2)%
}
\NovelDisplay{M01710R02}{%
\longmapsto(3+q,2), q\ge1.%
}

霍雨浩先检查根以前的保留段，蔡须坤再沿原坏部核阈值。每一新增段使用同一份原坏部与阶差，编号 q 为正整数；最高非零行及更后坐标不因段数增多而跟着改变。

\NovelDisplay{M01711R01}{%
\bigl((0)(1,1,1)(2,1)(3,2)(4,3)\bigr)[0]=(0)(1,1,1)(2,1)\textcolor[HTML]{B45309}{(3,2)}%
}

\NovelDisplay{M01712R01}{%
\bigl((0)(1,1,1)(2,1)(3,2)(4,3)\bigr)[1]=(0)(1,1,1)(2,1)\textcolor[HTML]{B45309}{(3,2)}\textcolor[HTML]{1D4ED8}{(4,2)}%
}

\NovelDisplay{M01713R01}{%
\bigl((0)(1,1,1)(2,1)(3,2)(4,3)\bigr)[2]=(0)(1,1,1)(2,1)\textcolor[HTML]{B45309}{(3,2)}\textcolor[HTML]{1D4ED8}{(4,2)}\textcolor[HTML]{15803D}{(5,2)}%
}

\NovelDisplay{M01714R01}{%
\text{第零项列数}=4,%
}
\NovelDisplay{M01714R02}{%
\text{第一项列数}=4+1=5,%
}
\NovelDisplay{M01714R03}{%
\text{第二项列数}=4+2\times1=6.%
}

蔡须坤先盖住答案，按列数把第一项与第二项各写一次。霍雨浩核坐标，卢本韦看保留段有没有误留末列，古月娜再查比较方向。两项都从当前原式求得，不能把第一项再次展开的结果填到第二项的位置。

古月娜在新的下标旁停笔时，山间响起了两声短促的鸟鸣。她听见声音，才发觉自己从清晨坐到这里，颈后早已僵得发酸。

卢本韦将水杯递给她，顺手要把上一组指数塔翻过去，却被霍雨浩留住。

``这一页还得放在旁边。''

他先让蔡须坤读上一组的幂与这里的下标。它们都只添了很短的一段，位置却不同。前面是以同一基数为底的指数结构，这里出现新的后继层基数；二者在给定的对应中有各自的 BMS 节点。若只认第二行的两个数字，忽略前缀里已经完成的结构，容易把局部的排列直接当作通用换名规则。

蔡须坤先指出最高非零行，再将坏根落在末列前的那一列。坏部只有一列，第一行每次按阶差生成下一段；他们熟悉的机械步骤并没有突然增加一条。真正需要重读的，是这一步在完整对应中怎样结束前面的指数过程，落到当前下标。

``所以展开算法会了，右边就都能靠形状猜出来？''卢本韦问。

``还不够。''古月娜说，``我们至少得核对它落在哪个上下文，前面的有限项怎样组织。不能拿几条相似的结果证明所有相似的列都有同样的名字。''

夏夜星空从石阶上走下来，听见这句，便将原先要收走的旧笔记又留在他们中间。四人各选一条前面的式子，指出共同的生成方式，也指出为什么这份说明只适用于当前已经核过的结构。熟悉的规则用来减少重算，不能替未检查的对应作保证。

霍雨浩将自己的那一页读得最快。读到最后，蔡须坤却问他，上界写在函数下标时究竟是哪一段。霍雨浩低头，才看见自己将两个位置相近的符号指成了同一个。

卢本韦忍住笑，将空杯转了半圈。``现在轮到你留下错页。''

他没有辩解，重新读完那两层。古月娜将大家的纸按前后摆好，留给晚上复核，免得疲倦时只因答案熟悉便匆匆放行。

这一天下山前，四人只在这里做了一次互问。接下来相同上界中的连续增长，便顺着已讲清的括号往下写。蔡须坤坐在最外侧，替同伴挡住风，直到最后一页墨迹干透，才将它们一并收好。

接下来的组织可以与此前到达 BO 之前的笔记逐项对照。

\NovelDisplay{M01715R01}{%
(0)(1,1,1)(2,1)(3,2)(4,3)(1,1,1)(2,1)(3,2)%
}
\NovelDisplay{M01715R02}{%
\quad (4,3)=\psi(\Omega_{\omega+2}+\psi_{\Omega_{\omega+1}}(\Omega_{\omega+2}))%
}

\NovelDisplay{M01716R01}{%
(0)(1,1,1)(2,1)(3,2)(4,3)(3,2)(4,3)=\psi(\Omega_{\omega+2}+\psi_{\Omega_{\omega+2}}(\Omega_{\omega+2}))%
}

\NovelDisplay{M01717R01}{%
(0)(1,1,1)(2,1)(3,2)(4,3)(4,3)=\psi(\Omega_{\omega+2}\times2)%
}

\NovelDisplay{M01718R01}{%
(0)(1,1,1)(2,1)(3,2)(4,3)(5,1)=\psi(\Omega_{\omega+2}\times\Omega)%
}

\NovelDisplay{M01719R01}{%
(0)(1,1,1)(2,1)(3,2)(4,3)(5,1)(1,1,1)=\psi(\Omega_{\omega+2}\times\Omega_{\omega})%
}

\NovelDisplay{M01720R01}{%
(0)(1,1,1)(2,1)(3,2)(4,3)(5,1)(6,2)=\psi(\Omega_{\omega+2}\times\psi_{\Omega_{\omega+1}}(\Omega_{\omega+1}))%
}

\NovelDisplay{M01721R01}{%
(0)(1,1,1)(2,1)(3,2)(4,3)(5,1)(6,2)(7,3)=\psi(\Omega_{\omega+2}\times\psi_{\Omega_{\omega+1}}(\Omega_{\omega+2}))%
}

\NovelDisplay{M01722R01}{%
(0)(1,1,1)(2,1)(3,2)(4,3)(5,2)=\psi(\Omega_{\omega+2}\times\Omega_{\omega+1})%
}

\NovelDisplay{M01723R01}{%
(0)(1,1,1)(2,1)(3,2)(4,3)(5,2)(6,3)=\psi(\Omega_{\omega+2}\times\psi_{\Omega_{\omega+2}}(\Omega_{\omega+2}))%
}

\NovelDisplay{M01724R01}{%
(0)(1,1,1)(2,1)(3,2)(4,3)(5,3)=\psi(\Omega_{\omega+2}^{2})%
}

\allowbreak{}\NovelInline{I05920}{\Omega_{\omega+2}}\allowbreak{} 。它的标志就是 \allowbreak{}\NovelInline{I05921}{(k,3)(k+1,3)}\allowbreak{} 了。

\NovelDisplay{M01725R01}{%
(0)(1,1,1)(2,1)(3,2)(4,3)(5,3)(6,3)=\psi(\Omega_{\omega+2}^{\Omega_{\omega+2}})%
}

\NovelDisplay{M01726R01}{%
(0)(1,1,1)(2,1)(3,2)(4,3)(5,4)=\psi(\Omega_{\omega+3})%
}

\NovelDisplay{M01727R01}{%
(0)(1,1,1)(2,1)(3,2)(4,3)(5,4)(6,4)=\psi(\Omega_{\omega+3}^{2})%
}

接下来的几日，他们把这一组逐项走完。古月娜将最末那列的第三行点亮时，蔡须坤脚下先传来一次短促的回牵。

``这一行终于参与了。''

霍雨浩把前后两式并排，卢本韦却让他先翻过纸，独自再找一遍根。古月娜靠着石壁等他，没有先去试下一阶。

这一次四人把答案对完，才轮流在压力下运转。第一层接缝稳住之后，还有更内的一处停顿。蔡须坤没有将那停顿硬推过去，退回横台，将它留给下一组分析。

一百六十米的石桩仍在更高处。

\NovelDisplay{M01728R01}{%
(0)(1,1,1)(2,1)(3,2,1)=\psi(\Omega_{\omega\times2})%
}

这次第三行不再是补齐的零。古月娜从第三行选根，蔡须坤核前一行的祖先限制；卢本韦把两种不同的追法并排放好，等展开结果自己给出区别。

\NovelDisplay{M01729R01}{%
(0)(1,1,1)\textcolor[HTML]{B45309}{(2,1)}\textcolor[HTML]{7E22CE}{(3,2,1)}%
}

\NovelDisplay{M01730R01}{%
\text{原列数}=4, h=3, r=3,%
}
\NovelDisplay{M01730R02}{%
\text{坏根}=\textcolor[HTML]{B45309}{(2,1)},\text{坏部列数}=1,%
}
\NovelDisplay{M01730R03}{%
(\delta_1,\delta_2,\delta_3)=(1,1,0).%
}

\NovelDisplay{M01731R01}{%
\text{按坏部列序排列的提升阈值：}%
}
\NovelDisplay{M01731R02}{%
(2),%
}
\NovelDisplay{M01731R03}{%
\text{第 }j\text{ 列只允许前 }\tau_j\text{ 行应用阶差。}%
}

\NovelDisplay{M01732R01}{%
\text{原第 }3\text{ 列：}(2,1)%
}
\NovelDisplay{M01732R02}{%
\longmapsto(2+q,1+q), q\ge1.%
}

霍雨浩先检查根以前的保留段，蔡须坤再沿原坏部核阈值。每一新增段使用同一份原坏部与阶差，编号 q 为正整数；最高非零行及更后坐标不因段数增多而跟着改变。

\NovelDisplay{M01733R01}{%
\bigl((0)(1,1,1)(2,1)(3,2,1)\bigr)[0]=(0)(1,1,1)\textcolor[HTML]{B45309}{(2,1)}%
}

\NovelDisplay{M01734R01}{%
\bigl((0)(1,1,1)(2,1)(3,2,1)\bigr)[1]=(0)(1,1,1)\textcolor[HTML]{B45309}{(2,1)}\textcolor[HTML]{1D4ED8}{(3,2)}%
}

\NovelDisplay{M01735R01}{%
\bigl((0)(1,1,1)(2,1)(3,2,1)\bigr)[2]=(0)(1,1,1)\textcolor[HTML]{B45309}{(2,1)}\textcolor[HTML]{1D4ED8}{(3,2)}\textcolor[HTML]{15803D}{(4,3)}%
}

\NovelDisplay{M01736R01}{%
\text{第零项列数}=3,%
}
\NovelDisplay{M01736R02}{%
\text{第一项列数}=3+1=4,%
}
\NovelDisplay{M01736R03}{%
\text{第二项列数}=3+2\times1=5.%
}

古月娜将第二项的两段新增部分分别圈出，卢本韦从原坏部核它们各自的增量。蔡须坤专查那些本该不动的位置，霍雨浩最后逐列比较。各项虽然可以更长，仍位于原式之下。

第三行的一出现以后，蔡须坤先把前面几页推远了些。

``我怕手还照着第二行的题写。''

霍雨浩将末列补齐，在它下方轻轻敲了三下。最高非零行换到第三行，选根便从第三行的父项出发。这里的根落在二一，根本身补齐的第三行是零；由根到末列前，坏部只有这一列。前两行可用的阶差分别求出，第三行不跟着加，新增段由这三件事共同确定。

卢本韦把手中的纸盖住。``我不用你的右边，也能写前两项。''

古月娜让她写。写完，蔡须坤检查列数，霍雨浩核每个坐标，古月娜最后问她为什么第三行仍保留原值。四人逐项对上以后，才重新看这里的对应：下标已经走到两个长度为小写欧米伽的段，原先第一个后继位置不能替它承担这个名字。

``我最容易把后继跟第二整段念成一个意思。''蔡须坤说。

她没有另起一串新例子，只把当前两处下标并排指给卢本韦。一处是在已有位置之后取后继，一处是按顺序接起两个同样长度的段。这里的加法和乘法保持次序，不能拿普通整数里熟悉的交换法则挪动它们。等会儿继续出现的加项，首先也要留在原来这一侧。

说完，她将纸翻回左端。若没有根与阶差的变化，只在右端把一个加号换成乘号，仍然解释不了自己究竟越过了什么。

古月娜起身时，肩背的银纹轻轻一晃。她也没有比别人轻松，先前追第三行时太专注，站直才觉出腿已经麻了。蔡须坤将石板放在一旁，扶着她等了一会儿。卢本韦没有催，反而接过笔，将古月娜刚才负责的末列再讲给霍雨浩听。

压力还在上方等着。他们重新落脚时，三行没有一齐被蛮力提起，而是按刚才核过的关系分别接住。脚下那种细小的牵扯并未全部消失，下一层的停顿已经先露了出来。

蔡须坤看向古月娜，两人都没有急着宣布已经完成两重提升。如今到达的，只是后面要继续分析的一个清楚起点。霍雨浩将那处尚未解释的停顿记在新页页首，让它接着问下去。

四人终于核清了这一组的末端： \allowbreak{}\NovelInline{I05922}{\psi(\Omega_{\omega\times2})}\allowbreak{} 。将它作为终点肯定是有原因的。因为在此之后，后面将出现``连续两段的提升效应''(这里还不是双重提升)。

\NovelDisplay{M01737R01}{%
(0)(1,1,1)(2,1)(3,2,1)(1,1,1)=\psi(\Omega_{\omega\times2}+\Omega_{\omega})%
}

\NovelDisplay{M01738R01}{%
(0)(1,1,1)(2,1)(3,2,1)(1,1,1)(2,1)(3,2)=\psi(\Omega_{\omega\times2}+\psi_{\Omega_{\omega+1}}(\Omega_{\omega+1}))%
}

\NovelDisplay{M01739R01}{%
(0)(1,1,1)(2,1)(3,2,1)(1,1,1)(2,1)(3,2)%
}
\NovelDisplay{M01739R02}{%
\quad (4,2)=\psi(\Omega_{\omega\times2}+\psi_{\Omega_{\omega+1}}(\Omega_{\omega+1}^{2}))%
}

\NovelDisplay{M01740R01}{%
(0)(1,1,1)(2,1)(3,2,1)(1,1,1)(2,1)(3,2)%
}
\NovelDisplay{M01740R02}{%
\quad (4,3)=\psi(\Omega_{\omega\times2}+\psi_{\Omega_{\omega+1}}(\Omega_{\omega+2}))%
}

\NovelDisplay{M01741R01}{%
(0)(1,1,1)(2,1)(3,2,1)(1,1,1)(2,1)(3,2,1)%
}
\NovelDisplay{M01741R02}{%
\quad =\psi(\Omega_{\omega\times2}+\psi_{\Omega_{\omega+1}}(\Omega_{\omega\times2}))%
}

被 \allowbreak{}\NovelInline{I05923}{\psi_{\Omega_{\omega+1}}}\allowbreak{} 稍微拦了一下。这是正常的。

\NovelDisplay{M01742R01}{%
(0)(1,1,1)(2,1)(3,2,1)(2,1)=\psi(\Omega_{\omega\times2}+\psi_{\Omega_{\omega+1}}(\Omega_{\omega\times2})\times\Omega)%
}

\NovelDisplay{M01743R01}{%
(0)(1,1,1)(2,1)(3,2,1)(2,1)(1,1,1)=\psi(\Omega_{\omega\times2}+\psi_{\Omega_{\omega+1}}(\Omega_{\omega\times2})\times\Omega_{\omega})%
}

\NovelDisplay{M01744R01}{%
(0)(1,1,1)(2,1)(3,2,1)(2,1)(3,2)=\psi(\Omega_{\omega\times2}%
}
\NovelDisplay{M01744R02}{%
\quad +\psi_{\Omega_{\omega+1}}(\Omega_{\omega\times2})\times\psi_{\Omega_{\omega+1}}(\Omega_{\omega+1}))%
}

\NovelDisplay{M01745R01}{%
(0)(1,1,1)(2,1)(3,2,1)(2,1)(3,2,1)=\psi(\Omega_{\omega\times2}+\psi_{\Omega_{\omega+1}}(\Omega_{\omega\times2})^{2})%
}

\NovelDisplay{M01746R01}{%
(0)(1,1,1)(2,1)(3,2,1)(3,1)=\psi(\Omega_{\omega\times2}+\psi_{\Omega_{\omega+1}}(\Omega_{\omega\times2})^{\Omega})%
}

\NovelDisplay{M01747R01}{%
(0)(1,1,1)(2,1)(3,2,1)(3,1)(4,2,1)=\psi(\Omega_{\omega\times2}%
}
\NovelDisplay{M01747R02}{%
\quad +\psi_{\Omega_{\omega+1}}(\Omega_{\omega\times2})^{\psi_{\Omega_{\omega+1}}(\Omega_{\omega\times2})})%
}

\NovelDisplay{M01748R01}{%
(0)(1,1,1)(2,1)(3,2,1)(3,2)=\psi(\Omega_{\omega\times2}+\Omega_{\omega+1})%
}

如果把后面三项单独拿出来并减去 \allowbreak{}\NovelInline{I05924}{(2,1)}\allowbreak{} 的话，它们就是 \allowbreak{}\NovelInline{I05925}{(0)(1,1,1)(1,1)}\allowbreak{} ，对应序数 \allowbreak{}\NovelInline{I05926}{\psi(\Omega_{\omega}+\Omega)}\allowbreak{} ，而这个是 \allowbreak{}\NovelInline{I05927}{\psi_{\Omega_{\omega+1}}(\Omega_{\omega\times2}+\Omega_{\omega+1})}\allowbreak{} ，相当于一个升级版。这种升级版的结构在BMS里面还会出现很多。

\NovelDisplay{M01749R01}{%
(0)(1,1,1)(2,1)(3,2,1)(3,2)(4,1)(5,2,1)(5,2)=\psi(\Omega_{\omega\times2}+\Omega_{\omega+1}%
}
\NovelDisplay{M01749R02}{%
\quad \times\psi_{\Omega_{\omega+1}}(\Omega_{\omega\times2}+\Omega_{\omega+1}))%
}

\NovelDisplay{M01750R01}{%
(0)(1,1,1)(2,1)(3,2,1)(3,2)(4,2)=\psi(\Omega_{\omega\times2}+\Omega_{\omega+1}^{2})%
}

\NovelDisplay{M01751R01}{%
(0)(1,1,1)(2,1)(3,2,1)(3,2)(4,2)(3,2)(4,2)=\psi(\Omega_{\omega\times2}+\Omega_{\omega+1}^{2}\times2)%
}

\NovelDisplay{M01752R01}{%
(0)(1,1,1)(2,1)(3,2,1)(3,2)(4,2)(4,1)(1,1,1)=\psi(\Omega_{\omega\times2}+\Omega_{\omega+1}^{2}\times\Omega_{\omega})%
}

\NovelDisplay{M01753R01}{%
(0)(1,1,1)(2,1)(3,2,1)(3,2)(4,2)(4,2)=\psi(\Omega_{\omega\times2}+\Omega_{\omega+1}^{3})%
}

\NovelDisplay{M01754R01}{%
(0)(1,1,1)(2,1)(3,2,1)(3,2)(4,2)(5,2)=\psi(\Omega_{\omega\times2}+\Omega_{\omega+1}^{\Omega_{\omega+1}})%
}

\NovelDisplay{M01755R01}{%
(0)(1,1,1)(2,1)(3,2,1)(3,2)(4,3)=\psi(\Omega_{\omega\times2}+\psi_{\Omega_{\omega+2}}(\Omega_{\omega+2}))%
}

\NovelDisplay{M01756R01}{%
(0)(1,1,1)(2,1)(3,2,1)(3,2)(4,3,1)=\psi(\Omega_{\omega\times2}+\psi_{\Omega_{\omega+2}}(\Omega_{\omega\times2}))%
}

\NovelDisplay{M01757R01}{%
(0)(1,1,1)(2,1)(3,2,1)(3,2)(4,3,1)(4,3)=\psi(\Omega_{\omega\times2}+\Omega_{\omega+2})%
}

\NovelDisplay{M01758R01}{%
(0)(1,1,1)(2,1)(3,2,1)(3,2)(4,3,1)(4,3)(5,3)=\psi(\Omega_{\omega\times2}+\Omega_{\omega+2}^{2})%
}

\NovelDisplay{M01759R01}{%
(0)(1,1,1)(2,1)(3,2,1)(3,2)(4,3,1)(4,3)%
}
\NovelDisplay{M01759R02}{%
\quad (5,4)=\psi(\Omega_{\omega\times2}+\psi_{\Omega_{\omega+3}}(\Omega_{\omega+3}))%
}

\NovelDisplay{M01760R01}{%
(0)(1,1,1)(2,1)(3,2,1)(3,2)(4,3,1)(4,3)%
}
\NovelDisplay{M01760R02}{%
\quad (5,4,1)=\psi(\Omega_{\omega\times2}+\psi_{\Omega_{\omega+3}}(\Omega_{\omega\times2}))%
}

\NovelDisplay{M01761R01}{%
(0)(1,1,1)(2,1)(3,2,1)(3,2,1)=\psi(\Omega_{\omega\times2}\times2)%
}

终于到了 \allowbreak{}\NovelInline{I05928}{\Omega_{\omega\times2}\times2}\allowbreak{} 。在这里， \allowbreak{}\NovelInline{I05929}{(3,2,1)}\allowbreak{} 的作用就类似于前面的 \allowbreak{}\NovelInline{I05930}{(1,1,1)}\allowbreak{} ，是添加 \allowbreak{}\NovelInline{I05931}{\Omega_{\omega\times2}}\allowbreak{} 。接下来，接下来的连接也需要重新追根了。

\NovelDisplay{M01762R01}{%
(0)(1,1,1)(2,1)(3,2,1)(4)=\psi(\Omega_{\omega\times2}\times\omega)%
}

卢本韦没有催，把下一行的空白留在纸上。

\NovelDisplay{M01763R01}{%
(0)(1,1,1)(2,1)(3,2,1)(4,1)=\psi(\Omega_{\omega\times2}\times\Omega)%
}

古月娜将末列重新补齐，提醒三人从这里接着追根。

\NovelDisplay{M01764R01}{%
(0)(1,1,1)(2,1)(3,2,1)(4,1)(1,1,1)=\psi(\Omega_{\omega\times2}\times\Omega_{\omega})%
}

日头移到对面山脊时，卢本韦把前后的两份记录对了许久。

``第一段到了这里。那第二段，是把刚才那套话再说一遍？''

``先别替它取名字。''蔡须坤将后面的空页压平，``看看根会不会仍在同一处。''

古月娜从刚刚到达的对应往回指，划开主项与乘数。第一段使乘数来到当前的位置，主项里原先那一段继续保留。接着还有中间坍缩结构，后面的末列要怎样作用于它，尚需新的前缀与父项关系。已知两个端点，不等于中间每一步都可以省掉。

霍雨浩将这一组保留下来的前缀抄到新页，特意不写右端。蔡须坤与卢本韦各追一遍根，两人的线在前面几列重合，到了下一处候选才出现需要讨论的分岔。古月娜只检查那处分岔，前面共同的部分便不再重新解释。

``这样才算接着学。''卢本韦说，``不然每天听见同样的开场，我都以为已经会了。''

她说着将自己早上那张错纸抽出来。错误不在完全陌生的地方，恰在她最有把握、准备一笔滑过去的那段。蔡须坤把纸接过去，按前缀、根、阶差依次指出发生变化的位置，再请她将与上一题相同的步骤自己接完。

两人讲到一半，蔡兴坤从下面端来热水。杯沿沾着一点泥，她先用袖角擦净，递过来的时候又瞧了一眼纸，问第一段是不是也还压在他们身上。

``还在。''蔡须坤将手掌放回膝上，``只是我们现在知道哪一处该怎么接。后面还有看不清的，才会把整段一起拉乱。''

霍雨浩听见这句，将原先写的``第一段压力消失''改成``第一段已能分辨''。山上的重量没有随着一个名词被说对便永久退去，疲惫时，他们仍可能在熟悉的接缝犯错。能复查、能重来，才使后面的分析有了落脚处。

天暗后，四人收起纸，沿旧路回到营地。蔡须坤在转弯处扶了卢本韦一把，卢本韦走稳，又回身接住古月娜的纸袋。霍雨浩最后确认无人落下，才将精神探测收窄到脚下。

第二天还要继续。那条没写完的父项链，仍清清楚楚留在新页上。

第一段提升效应开始显出。 \allowbreak{}\NovelInline{I05932}{(1,1,1)}\allowbreak{} 将 \allowbreak{}\NovelInline{I05933}{(4,1)}\allowbreak{} 代表的\allowbreak{}\NovelInline{I05934}{\Omega}\allowbreak{} 提升至 \allowbreak{}\NovelInline{I05935}{\Omega_\omega}\allowbreak{} 。

接着分析后面的结构。

回到石屋，四人将当天的记录摊到灯下，没有再翻新题。

白天被压得最紧的那处接缝，已经能在各自的演算中分开说清。霍雨浩将写错的幂留在页角，古月娜在旁边添上正确的父项；蔡须坤把两种运转的先后记完，卢本韦才将纸装进干布袋。

``后面还有一段。''

蔡须坤朝门外望了一眼：``知道。从第一段已经站稳的地方接。''

山图上，一百六十米的石桩就在转弯另一侧，却还没有落到他们脚下。霍雨浩没有提前画勾，只在当前页末留下一行空白，等明天把第二段的来路填进去。

\chapter{第一百零九章：bms的两重提升}\label{ux7b2cux4e00ux767eux96f6ux4e5dux7ae0bmsux7684ux4e24ux91cdux63d0ux5347}

翌日清早，风把上方石桩旁的叶子吹开了一瞬，``一百六十''的刻字露出来，又被摇回去的树影遮住。

卢本韦没有抬头。她正数面前两条祖先链，一条追到前面的三行列，一条继续往左。古月娜的指尖停在第二行，等她数完才往下移。

``你把第二段的根，拿去替第一段用了。''

卢本韦皱了皱眉，将纸转向蔡须坤。``你来找。我先不说我数到了哪里。''

蔡须坤从末列重新出发。压力把她的呼吸拖得很长，心里原先一闪而过的连线，此刻要一处处接住。她说出根的位置时，霍雨浩才把自己背面朝上的那张纸翻过来。

两人的标记相同。

``只对上这一处还不够。''霍雨浩把两张纸错开，``要把新段接进去，再看下一层为什么能跟着变。''

唐三替后面的人留好落脚处。小舞将一小包盐放在温水旁，奥斯卡把补给递上来，随即退到不妨碍四人演算的位置。蔡兴坤抱着黑子鸡蹲在边上，瞧见蔡须坤膝盖发颤，往她手边推来一块平整的石板。

蔡须坤接住石板，把下一个等号写在中间。她知道这一回将经过一百六十米，也知道脚下先要走完的，还包括第一段提升之后那些看着很像的中间式。

古月娜将已经画好的两条链放回各自的行。

``先让它们分别成立，再看怎样衔接。''

\NovelDisplay{M01765R01}{%
(0)(1,1,1)(2,1)(3,2,1)(4,1)(2,1)(3,2)=\psi(\Omega_{\omega\times2}\times\Omega_{\omega}%
}
\NovelDisplay{M01765R02}{%
\quad +\psi_{\Omega_{\omega+1}}(\Omega_{\omega\times2}\times\Omega_{\omega})\times\psi_{\Omega_{\omega+1}}(\Omega_{\omega+1}))%
}

\NovelDisplay{M01766R01}{%
(0)(1,1,1)(2,1)(3,2,1)(4,1)(3,2)=\psi(\Omega_{\omega\times2}\times\Omega_{\omega}+\Omega_{\omega+1})%
}

\NovelDisplay{M01767R01}{%
(0)(1,1,1)(2,1)(3,2,1)(4,1)(3,2)(4,1)(5,2,1)(6,1)(5,2)=\psi(\Omega_{\omega\times2}\times\Omega_{\omega}%
}
\NovelDisplay{M01767R02}{%
\quad +\Omega_{\omega+1}\times\psi_{\Omega_{\omega+1}}(\Omega_{\omega\times2}\times\Omega_{\omega}+\Omega_{\omega+1}))%
}

\NovelDisplay{M01768R01}{%
(0)(1,1,1)(2,1)(3,2,1)(4,1)(3,2)(4,2)=\psi(\Omega_{\omega\times2}\times\Omega_{\omega}+\Omega_{\omega+1}^{2})%
}

\NovelDisplay{M01769R01}{%
(0)(1,1,1)(2,1)(3,2,1)(4,1)(3,2)(4,3)%
}
\NovelDisplay{M01769R02}{%
\quad =\psi(\Omega_{\omega\times2}\times\Omega_{\omega}+\psi_{\Omega_{\omega+2}}(\Omega_{\omega+2}))%
}

\NovelDisplay{M01770R01}{%
(0)(1,1,1)(2,1)(3,2,1)(4,1)(3,2)(4,3,1)(5,1)(1,1,1)=\psi(\Omega_{\omega\times2}\times\Omega_{\omega}%
}
\NovelDisplay{M01770R02}{%
\quad +\psi_{\Omega_{\omega+2}}(\Omega_{\omega\times2}\times\Omega_{\omega}))%
}

\NovelDisplay{M01771R01}{%
(0)(1,1,1)(2,1)(3,2,1)(4,1)(3,2)(4,3,1)(5,1)%
}
\NovelDisplay{M01771R02}{%
\quad (4,3)(5,3)=\psi(\Omega_{\omega\times2}\times\Omega_{\omega}+\Omega_{\omega+2}^{2})%
}

\NovelDisplay{M01772R01}{%
(0)(1,1,1)(2,1)(3,2,1)(4,1)(3,2)(4,3,1)(5,1)(4,3)(5,4,1)(6,1)(5,4)(6,4)%
}
\NovelDisplay{M01772R02}{%
\quad =\psi(\Omega_{\omega\times2}\times\Omega_{\omega}+\Omega_{\omega+3}^{2})%
}

\NovelDisplay{M01773R01}{%
(0)(1,1,1)(2,1)(3,2,1)(4,1)(3,2,1)=\psi(\Omega_{\omega\times2}\times\Omega_{\omega}+\Omega_{\omega\times2})%
}

仍然是一个 \allowbreak{}\NovelInline{I05936}{(3,2,1)}\allowbreak{} 对应一个 \allowbreak{}\NovelInline{I05937}{\Omega_{\omega\times2}}\allowbreak{} 的样子。

\NovelDisplay{M01774R01}{%
(0)(1,1,1)(2,1)(3,2,1)(4,1)(3,2,1)(3,2,1)%
}
\NovelDisplay{M01774R02}{%
\quad =\psi(\Omega_{\omega\times2}\times\Omega_{\omega}+\Omega_{\omega\times2}\times2)%
}

\NovelDisplay{M01775R01}{%
(0)(1,1,1)(2,1)(3,2,1)(4,1)(3,2,1)(4,1)(1,1,1)=\psi(\Omega_{\omega\times2}\times\Omega_{\omega}\times2)%
}

\NovelDisplay{M01776R01}{%
(0)(1,1,1)(2,1)(3,2,1)(4,1)(4,1)(1,1,1)=\psi(\Omega_{\omega\times2}\times\Omega_{\omega}^{2})%
}

\NovelDisplay{M01777R01}{%
(0)(1,1,1)(2,1)(3,2,1)(4,1)(5,2)=\psi(\Omega_{\omega\times2}\times\psi_{\Omega_{\omega+1}}(\Omega_{\omega+1}))%
}

\NovelDisplay{M01778R01}{%
(0)(1,1,1)(2,1)(3,2,1)(4,1)(5,2,1)=\psi(\Omega_{\omega\times2}\times\psi_{\Omega_{\omega+1}}(\Omega_{\omega+1}))%
}

\NovelDisplay{M01779R01}{%
(0)(1,1,1)(2,1)(3,2,1)(4,1)(5,2,1)(6,1)(7,2,1)=\psi(\Omega_{\omega\times2}%
}
\NovelDisplay{M01779R02}{%
\quad \times\psi_{\Omega_{\omega+1}}(\Omega_{\omega\times2}\times\psi_{\Omega_{\omega+1}}(\Omega_{\omega\times2})))%
}

\NovelDisplay{M01780R01}{%
(0)(1,1,1)(2,1)(3,2,1)(4,2)=\psi(\Omega_{\omega\times2}\times\Omega_{\omega+1})%
}

但是现在情况就不一样了。还记得前面说的，``把后面三项拿出来减掉一个 \allowbreak{}\NovelInline{I05938}{(2,1)}\allowbreak{} ''吗？如果对这个式子这么做，四人会得到 \allowbreak{}\NovelInline{I05939}{(0)(1,1,1)(2,1)}\allowbreak{} 。而 \allowbreak{}\NovelInline{I05940}{(0)(1,1,1)(2,1)(1,1,1)}\allowbreak{} 有提升。所以接下来是：

\NovelDisplay{M01781R01}{%
(0)(1,1,1)(2,1)(3,2,1)(4,2)(3,2)(4,2)=\psi(\Omega_{\omega\times2}\times\Omega_{\omega+1}+\Omega_{\omega+1}^{2})%
}

\NovelDisplay{M01782R01}{%
(0)(1,1,1)(2,1)(3,2,1)(4,2)(3,2)(4,3,1)(5,2)=\psi(\Omega_{\omega\times2}\times\Omega_{\omega+1}%
}
\NovelDisplay{M01782R02}{%
\quad +\psi_{\Omega_{\omega+2}}(\Omega_{\omega\times2}\times\Omega_{\omega+1}))%
}

替身 \allowbreak{}\NovelInline{I05941}{(5,2)}\allowbreak{} 召唤。

\NovelDisplay{M01783R01}{%
(0)(1,1,1)(2,1)(3,2,1)(4,2)(3,2)(4,3,1)%
}
\NovelDisplay{M01783R02}{%
\quad (5,2)(4,3)=\psi(\Omega_{\omega\times2}\times\Omega_{\omega+1}+\Omega_{\omega+2})%
}

\NovelDisplay{M01784R01}{%
(0)(1,1,1)(2,1)(3,2,1)(4,2)(3,2)(4,3,1)(5,2)%
}
\NovelDisplay{M01784R02}{%
\quad (4,3)(5,3)=\psi(\Omega_{\omega\times2}\times\Omega_{\omega+1}+\Omega_{\omega+2}^{2})%
}

\NovelDisplay{M01785R01}{%
(0)(1,1,1)(2,1)(3,2,1)(4,2)(3,2)(4,3,1)(5,2)%
}
\NovelDisplay{M01785R02}{%
\quad (4,3,1)=\psi(\Omega_{\omega\times2}\times\Omega_{\omega+1}+\Omega_{\omega\times2})%
}

可以看到，现在添加 \allowbreak{}\NovelInline{I05942}{\Omega_{\omega\times2}}\allowbreak{} 的变成 \allowbreak{}\NovelInline{I05943}{(4,3,1)}\allowbreak{} 了。

\NovelDisplay{M01786R01}{%
(0)(1,1,1)(2,1)(3,2,1)(4,2)(3,2)(4,3,1)(5,2)%
}
\NovelDisplay{M01786R02}{%
\quad (4,3,1)(5,2)=\psi(\Omega_{\omega\times2}\times\Omega_{\omega+1}\times2)%
}

\NovelDisplay{M01787R01}{%
(0)(1,1,1)(2,1)(3,2,1)(4,2)(3,2)(4,3,1)(5,2)(5,2)=\psi(\Omega_{\omega\times2}\times\Omega_{\omega+1}^{2})%
}

\NovelDisplay{M01788R01}{%
(0)(1,1,1)(2,1)(3,2,1)(4,2)(3,2)(4,3,1)(5,2)%
}
\NovelDisplay{M01788R02}{%
\quad (6,3,1)=\psi(\Omega_{\omega\times2}\times\psi_{\Omega_{\omega+2}}(\Omega_{\omega\times2}))%
}

\NovelDisplay{M01789R01}{%
(0)(1,1,1)(2,1)(3,2,1)(4,2)(3,2)(4,3,1)(5,3)=\psi(\Omega_{\omega\times2}\times\Omega_{\omega+2})%
}

\NovelDisplay{M01790R01}{%
(0)(1,1,1)(2,1)(3,2,1)(4,2)(3,2)(4,3,1)(5,3)%
}
\NovelDisplay{M01790R02}{%
\quad (4,3)(5,4,1)(6,4)=\psi(\Omega_{\omega\times2}\times\Omega_{\omega+3})%
}

卢本韦在这一处抢先写出了两行，第三行却空着。

``我先猜不动。''

蔡须坤看着她的纸笑了，笑到一半，肩上的压力将那口气压回胸口。她扶住石沿，等气息平下来，才把第三行补齐。

``这回把猜字也写上。''

霍雨浩将末列的父项链写在她们旁边，古月娜负责逐列查哪些坐标真正允许提升。等第三行的零增量与前两行的非零增量分清，四人才把那枚问号擦去。

\NovelDisplay{M01791R01}{%
(0)(1,1,1)(2,1)(3,2,1)(4,2)(3,2,1)=\psi(\Omega_{\omega\times2}^{2})%
}

要检验第一段是否已经足够，卢本韦将最末那列单独圈出。霍雨浩没有简写前缀，让本次根以前的每一列都留在纸上。

\NovelDisplay{M01792R01}{%
(0)(1,1,1)\textcolor[HTML]{B45309}{(2,1)}\textcolor[HTML]{B91C1C}{(3,2,1)}\textcolor[HTML]{B91C1C}{(4,2)}\textcolor[HTML]{7E22CE}{(3,2,1)}%
}

\NovelDisplay{M01793R01}{%
\text{原列数}=6, h=3, r=3,%
}
\NovelDisplay{M01793R02}{%
\text{坏根}=\textcolor[HTML]{B45309}{(2,1)},\text{坏部列数}=3,%
}
\NovelDisplay{M01793R03}{%
(\delta_1,\delta_2,\delta_3)=(1,1,0).%
}

\NovelDisplay{M01794R01}{%
\text{按坏部列序排列的提升阈值：}%
}
\NovelDisplay{M01794R02}{%
(2,2,2),%
}
\NovelDisplay{M01794R03}{%
\text{第 }j\text{ 列只允许前 }\tau_j\text{ 行应用阶差。}%
}

\NovelDisplay{M01795R01}{%
\text{原第 }3\text{ 列：}(2,1)%
}
\NovelDisplay{M01795R02}{%
\longmapsto(2+q,1+q), q\ge1.%
}

\NovelDisplay{M01796R01}{%
\text{原第 }5\text{ 列：}(4,2)%
}
\NovelDisplay{M01796R02}{%
\longmapsto(4+q,2+q), q\ge1.%
}

霍雨浩先检查根以前的保留段，蔡须坤再沿原坏部核阈值。每一新增段使用同一份原坏部与阶差，编号 q 为正整数；最高非零行及更后坐标不因段数增多而跟着改变。

\NovelDisplay{M01797R01}{%
\bigl((0)(1,1,1)(2,1)(3,2,1)(4,2)(3,2,1)\bigr)[0]=(0)(1,1,1)\textcolor[HTML]{B45309}{(2,1)}\textcolor[HTML]{B91C1C}{(3,2,1)}\textcolor[HTML]{B91C1C}{(4,2)}%
}

\NovelDisplay{M01798R01}{%
\bigl((0)(1,1,1)(2,1)(3,2,1)(4,2)(3,2,1)\bigr)[1]=(0)%
}
\NovelDisplay{M01798R02}{%
\quad (1,1,1)\textcolor[HTML]{B45309}{(2,1)}\textcolor[HTML]{B91C1C}{(3,2,1)}\textcolor[HTML]{B91C1C}{(4,2)}\textcolor[HTML]{1D4ED8}{(3,2)}\textcolor[HTML]{1D4ED8}{(4,3,1)}\textcolor[HTML]{1D4ED8}{(5,3)}%
}

\NovelDisplay{M01799R01}{%
\bigl((0)(1,1,1)(2,1)(3,2,1)(4,2)(3,2,1)\bigr)[2]=(0)(1,1,1)\textcolor[HTML]{B45309}{(2,1)}\textcolor[HTML]{B91C1C}{(3,2,1)}\textcolor[HTML]{B91C1C}{(4,2)}%
}
\NovelDisplay{M01799R02}{%
\quad \textcolor[HTML]{1D4ED8}{(3,2)}\textcolor[HTML]{1D4ED8}{(4,3,1)}\textcolor[HTML]{1D4ED8}{(5,3)}\textcolor[HTML]{15803D}{(4,3)}\textcolor[HTML]{15803D}{(5,4,1)}\textcolor[HTML]{15803D}{(6,4)}%
}

\NovelDisplay{M01800R01}{%
\text{第零项列数}=5,%
}
\NovelDisplay{M01800R02}{%
\text{第一项列数}=5+3=8,%
}
\NovelDisplay{M01800R03}{%
\text{第二项列数}=5+2\times3=11.%
}

卢本韦数完列，将第一项中的蓝段与第二项中的蓝段逐列重合。绿色段另由原坏部按两倍阶差生成，原来的保留段没有再受一次提升。蔡须坤核坐标，古月娜与霍雨浩最后检查各有限项都小于原式。

``这才是第二段真正接上的地方。''霍雨浩说。

卢本韦却没有立刻放下笔。``你先说，为什么不是把前一段多做几次。''

霍雨浩将较早的一页移过来，让蔡须坤先指出第一段停在哪里。那一段使原先的乘数增长到已辨认的层级，后面仍有一层需要越过的中间结构。眼前的新末列接进来，父项与增量在当前前缀中重新组织，原先受阻的那一处才进一步到达第二个整段的位置。两个阶段前后衔接，各有来路；此时尚不能因为它们连续发生，便把它们统称为后面那种嵌套的双重提升。

蔡须坤将这句话拆成两次运转。第一次，她先稳住已经熟悉的连接，不去追更远处那一点牵扯。第二次，才从留下来的前缀接新末列。压力将两次停顿分得格外清楚，她记下先后，仍请古月娜回到列上检查，免得身体感觉近似，纸上的层级却读错。

古月娜看完，在第一段与第二段之间画了一道短线。

``一段完成，留下它的结果，再接后一段。后面讲到嵌套时，要拿这个区别回来比。''

卢本韦点头，随即把两段的颜色对调了一遍。她想看看四人是否真的按位置读，还是只记住某种颜色代表某种结果。蔡须坤果然在开头迟疑了一瞬，霍雨浩也抬起笔重新数了列号。古月娜没有出声，等他们把根的位置重新说全，才将那张换色纸收在反例页后面。

``颜色是帮我们追来历的。''蔡须坤说，``不能让颜色替我们选根。''

他们重新生成第一项，列数与坐标都对上以后，才将当前的平方放回完整对应。这里的平方落在基数主项上，既不删除外层坍缩，也不把前面两段提升压成普通的整数运算。卢本韦在纸边留下这一处位置说明，后面的同类有限增长便直接核公式。

四人起身时，横台外侧的雾散了一些。唐三将绳结往上挪，给他们留出更宽的转身处。蔡须坤走到一半，忽然又退回来，拿起那张换色的纸。

卢本韦看着她：``不是已经对过了？''

``晚上给我。我想不看旧页再算一遍。''

卢本韦将纸折好，塞进她的本子。古月娜已经站在下一处等着，霍雨浩关好水囊，四人这才重新接上那条向上的路。

完成了。这正是眼前的第二段提升效应，它将 \allowbreak{}\NovelInline{I05944}{\Omega_{\omega+1}}\allowbreak{} 提升到 \allowbreak{}\NovelInline{I05945}{\Omega_{\omega\times2}}\allowbreak{} 。

\NovelDisplay{M01801R01}{%
(0)(1,1,1)(2,1)(3,2,1)(4,2)(4,2)(3,2,1)=\psi(\Omega_{\omega\times2}^{3})%
}

\NovelDisplay{M01802R01}{%
(0)(1,1,1)(2,1)(3,2,1)(4,2)(5,1)(1,1,1)=\psi(\Omega_{\omega\times2}^{\Omega_{\omega}})%
}

指数上的第一段。

\NovelDisplay{M01803R01}{%
(0)(1,1,1)(2,1)(3,2,1)(4,2)(5,2)(3,2,1)=\psi(\Omega_{\omega\times2}^{\Omega_{\omega\times2}})%
}

指数上的第二段。

\NovelDisplay{M01804R01}{%
(0)(1,1,1)(2,1)(3,2,1)(4,2)(5,3)=\psi(\Omega_{\omega\times2+1})%
}

用和前面 \allowbreak{}\NovelInline{I05946}{\psi_{\Omega_{\omega+1}}(\Omega_{\omega+1})}\allowbreak{} 一样的方法构建指数塔，最后得到 \allowbreak{}\NovelInline{I05947}{\psi_{\Omega_{\omega\times2+1}}(\Omega_{\omega\times2+1})}\allowbreak{} 。

\NovelDisplay{M01805R01}{%
(0)(1,1,1)(2,1)(3,2,1)(4,2)(5,3)(5,3)=\psi(\Omega_{\omega\times2+1}\times2)%
}

\NovelDisplay{M01806R01}{%
(0)(1,1,1)(2,1)(3,2,1)(4,2)(5,3)(6,1)=\psi(\Omega_{\omega\times2+1}\times\Omega)%
}

\NovelDisplay{M01807R01}{%
(0)(1,1,1)(2,1)(3,2,1)(4,2)(5,3)(6,1)(1,1,1)=\psi(\Omega_{\omega\times2+1}\times\Omega_{\omega})%
}

\NovelDisplay{M01808R01}{%
(0)(1,1,1)(2,1)(3,2,1)(4,2)(5,3)(6,2)(3,2,1)=\psi(\Omega_{\omega\times2+1}\times\Omega_{\omega\times2})%
}

发生在 \allowbreak{}\NovelInline{I05948}{\psi_{\Omega_{\omega\times2+1}}}\allowbreak{} 内层的两段提升。原理都是一样的。

\NovelDisplay{M01809R01}{%
(0)(1,1,1)(2,1)(3,2,1)(4,2)(5,3)(6,2)(7,3)=\psi(\Omega_{\omega\times2+1}%
}
\NovelDisplay{M01809R02}{%
\quad \times\psi_{\Omega_{\omega\times2+1}}(\Omega_{\omega\times2+1}))%
}

\NovelDisplay{M01810R01}{%
(0)(1,1,1)(2,1)(3,2,1)(4,2)(5,3)(6,3)=\psi(\Omega_{\omega\times2+1}^{2})%
}

\NovelDisplay{M01811R01}{%
(0)(1,1,1)(2,1)(3,2,1)(4,2)(5,3)(6,4)=\psi(\Omega_{\omega\times2+2})%
}

\NovelDisplay{M01812R01}{%
(0)(1,1,1)(2,1)(3,2,1)(4,2)(5,3,1)=\psi(\Omega_{\omega\times3})%
}

第二段提升接通前，蔡须坤先确认旧段怎样保留。古月娜逐列指出可动的行，卢本韦照着第一项重新摆出新增段，再由霍雨浩独立生成第二项。

\NovelDisplay{M01813R01}{%
(0)(1,1,1)(2,1)(3,2,1)\textcolor[HTML]{B45309}{(4,2)}\textcolor[HTML]{7E22CE}{(5,3,1)}%
}

\NovelDisplay{M01814R01}{%
\text{原列数}=6, h=3, r=5,%
}
\NovelDisplay{M01814R02}{%
\text{坏根}=\textcolor[HTML]{B45309}{(4,2)},\text{坏部列数}=1,%
}
\NovelDisplay{M01814R03}{%
(\delta_1,\delta_2,\delta_3)=(1,1,0).%
}

\NovelDisplay{M01815R01}{%
\text{按坏部列序排列的提升阈值：}%
}
\NovelDisplay{M01815R02}{%
(2),%
}
\NovelDisplay{M01815R03}{%
\text{第 }j\text{ 列只允许前 }\tau_j\text{ 行应用阶差。}%
}

\NovelDisplay{M01816R01}{%
\text{原第 }5\text{ 列：}(4,2)%
}
\NovelDisplay{M01816R02}{%
\longmapsto(4+q,2+q), q\ge1.%
}

霍雨浩先检查根以前的保留段，蔡须坤再沿原坏部核阈值。每一新增段使用同一份原坏部与阶差，编号 q 为正整数；最高非零行及更后坐标不因段数增多而跟着改变。

\NovelDisplay{M01817R01}{%
\bigl((0)(1,1,1)(2,1)(3,2,1)(4,2)(5,3,1)\bigr)[0]=(0)(1,1,1)(2,1)(3,2,1)\textcolor[HTML]{B45309}{(4,2)}%
}

\NovelDisplay{M01818R01}{%
\bigl((0)(1,1,1)(2,1)(3,2,1)(4,2)(5,3,1)\bigr)[1]%
}
\NovelDisplay{M01818R02}{%
\quad =(0)(1,1,1)(2,1)(3,2,1)\textcolor[HTML]{B45309}{(4,2)}\textcolor[HTML]{1D4ED8}{(5,3)}%
}

\NovelDisplay{M01819R01}{%
\bigl((0)(1,1,1)(2,1)(3,2,1)(4,2)(5,3,1)\bigr)[2]%
}
\NovelDisplay{M01819R02}{%
\quad =(0)(1,1,1)(2,1)(3,2,1)\textcolor[HTML]{B45309}{(4,2)}\textcolor[HTML]{1D4ED8}{(5,3)}\textcolor[HTML]{15803D}{(6,4)}%
}

\NovelDisplay{M01820R01}{%
\text{第零项列数}=5,%
}
\NovelDisplay{M01820R02}{%
\text{第一项列数}=5+1=6,%
}
\NovelDisplay{M01820R03}{%
\text{第二项列数}=5+2\times1=7.%
}

蔡须坤先盖住答案，按列数把第一项与第二项各写一次。霍雨浩核坐标，卢本韦看保留段有没有误留末列，古月娜再查比较方向。两项都从当前原式求得，不能把第一项再次展开的结果填到第二项的位置。

后面连续三段的结构出现时，蔡须坤先伸出三根手指，随即又收了回来。

``光数三次还不够。''

卢本韦将三处对应的位置依次圈出。起点不同，所处的中间层级也不同；每一段将什么带到哪里，要分别读清。最后这些位置同属于下标按顺序延长的过程，才看得见整段规律。若把前面已经成立的一段复制三遍，根以前的前缀却没有随之保留，左端首先便不是当前这条 BMS。

霍雨浩让三人各负责一段，却不许只交自己的终点。蔡须坤要交给卢本韦一份能继续求父项的完整前缀，卢本韦再将结果交给古月娜；霍雨浩从最后一段倒着检查三次交接处。前一人少写一列，后一人就必须指出缺少的信息，不能凭记得右端而补答案。

第一次交接，卢本韦便将纸推了回来。

``根以前这列去哪儿了？''

蔡须坤低头，发现自己为了让纸更短，把一段当作已经用完的背景。那段虽然不直接接受这回增量，仍参与之后的祖先选择。她将它补回，重新接出后面的列，再递过去。

古月娜拿到第二份时也停了一次。卢本韦的坐标都对，括号却少了一层。霍雨浩用笔帽挡住右端，请两人只凭左端接下去。两人追到发生分歧的位置，才知道那层括号究竟保留了怎样的上下文。

第二轮交接慢了许多，四张纸终于从头到尾接上。

``这比各算一遍还累。''卢本韦将笔横在指间。

``各算各的，容易把同一个地方都想当然。''古月娜说。

蔡须坤把完整记录留在中央。压力下的运转也按交接分开，谁负责的那段含混，便由谁先停下来说明。霍雨浩的共享只传去受阻先后的观察，答案由每个人自己算完。他们能借彼此的眼睛发现问题，却不能借同伴已经理解的那部分代替自己的缺口。

黄昏时，十一人一起回营，薛峰将新写的几页与旧册编号接好。蔡须坤接过本子，看见此前两个月留下的那一叠错页仍在前面。它们厚得很，却没有白占地方。如今三次交接时反复需要保留的，正是那些曾让他们停在七十米之前的细节。

到 \allowbreak{}\NovelInline{I05949}{\Omega_{\omega\times3}}\allowbreak{} 了。那么接下来会是连续三段提升。

\NovelDisplay{M01821R01}{%
(0)(1,1,1)(2,1)(3,2,1)(4,2)(5,3,1)(5,3,1)=\psi(\Omega_{\omega\times3}\times2)%
}

\NovelDisplay{M01822R01}{%
(0)(1,1,1)(2,1)(3,2,1)(4,2)(5,3,1)(6,1)(1,1,1)=\psi(\Omega_{\omega\times3}\times\Omega_{\omega})%
}

第一段， \allowbreak{}\NovelInline{I05950}{\Omega}\allowbreak{} 到 \allowbreak{}\NovelInline{I05951}{\Omega_{\omega}}\allowbreak{} 。

\NovelDisplay{M01823R01}{%
(0)(1,1,1)(2,1)(3,2,1)(4,2)(5,3,1)(6,2)(3,2,1)=\psi(\Omega_{\omega\times3}\times\Omega_{\omega\times2})%
}

第二段， \allowbreak{}\NovelInline{I05952}{\Omega_{\omega+1}}\allowbreak{} 到 \allowbreak{}\NovelInline{I05953}{\Omega_{\omega\times2}}\allowbreak{} 。

\NovelDisplay{M01824R01}{%
(0)(1,1,1)(2,1)(3,2,1)(4,2)(5,3,1)(6,3)(5,3,1)=\psi(\Omega_{\omega\times3}^{2})%
}

第三段， \allowbreak{}\NovelInline{I05954}{\Omega_{\omega\times2+1}}\allowbreak{} 到 \allowbreak{}\NovelInline{I05955}{\Omega_{\omega\times3}}\allowbreak{} 。还是比较有规律的。

\NovelDisplay{M01825R01}{%
(0)(1,1,1)(2,1)(3,2,1)(4,2)(5,3,1)(6,3)(7,4)=\psi(\Omega_{\omega\times3+1})%
}

\NovelDisplay{M01826R01}{%
(0)(1,1,1)(2,1)(3,2,1)(4,2)(5,3,1)(6,3)(7,4,1)=\psi(\Omega_{\omega\times4})%
}

\NovelDisplay{M01827R01}{%
(0)(1,1,1)(2,1,1)=\psi(\Omega_{\omega^{2}})%
}

次日清早，四人来到一百六十米石桩下方的最后一级。霍雨浩先在纸上重新展开这三列。蔡须坤记得此前的预演，这次却要把每一列怎样参与增量都说全；古月娜请她先补回省略的零。

\NovelDisplay{M01828R01}{%
\textcolor[HTML]{B45309}{(0)}\textcolor[HTML]{B91C1C}{(1,1,1)}\textcolor[HTML]{7E22CE}{(2,1,1)}%
}

\NovelDisplay{M01829R01}{%
\text{原列数}=3, h=3, r=1,%
}
\NovelDisplay{M01829R02}{%
\text{坏根}=\textcolor[HTML]{B45309}{(0)},\text{坏部列数}=2,%
}
\NovelDisplay{M01829R03}{%
(\delta_1,\delta_2,\delta_3)=(2,1,0).%
}

\NovelDisplay{M01830R01}{%
\text{按坏部列序排列的提升阈值：}%
}
\NovelDisplay{M01830R02}{%
(2,2),%
}
\NovelDisplay{M01830R03}{%
\text{第 }j\text{ 列只允许前 }\tau_j\text{ 行应用阶差。}%
}

\NovelDisplay{M01831R01}{%
\text{原第 }1\text{ 列：}(0)%
}
\NovelDisplay{M01831R02}{%
\longmapsto(2q,q), q\ge1.%
}

\NovelDisplay{M01832R01}{%
\text{原第 }2\text{ 列：}(1,1,1)%
}
\NovelDisplay{M01832R02}{%
\longmapsto(1+2q,1+q,1), q\ge1.%
}

霍雨浩先检查根以前的保留段，蔡须坤再沿原坏部核阈值。每一新增段使用同一份原坏部与阶差，编号 q 为正整数；最高非零行及更后坐标不因段数增多而跟着改变。

\NovelDisplay{M01833R01}{%
\bigl((0)(1,1,1)(2,1,1)\bigr)[0]=\textcolor[HTML]{B45309}{(0)}\textcolor[HTML]{B91C1C}{(1,1,1)}%
}

\NovelDisplay{M01834R01}{%
\bigl((0)(1,1,1)(2,1,1)\bigr)[1]=\textcolor[HTML]{B45309}{(0)}\textcolor[HTML]{B91C1C}{(1,1,1)}\textcolor[HTML]{1D4ED8}{(2,1)}\textcolor[HTML]{1D4ED8}{(3,2,1)}%
}

\NovelDisplay{M01835R01}{%
\bigl((0)(1,1,1)(2,1,1)\bigr)[2]=\textcolor[HTML]{B45309}{(0)}\textcolor[HTML]{B91C1C}{(1,1,1)}\textcolor[HTML]{1D4ED8}{(2,1)}\textcolor[HTML]{1D4ED8}{(3,2,1)}\textcolor[HTML]{15803D}{(4,2)}\textcolor[HTML]{15803D}{(5,3,1)}%
}

\NovelDisplay{M01836R01}{%
\text{第零项列数}=2,%
}
\NovelDisplay{M01836R02}{%
\text{第一项列数}=2+2=4,%
}
\NovelDisplay{M01836R03}{%
\text{第二项列数}=2+2\times2=6.%
}

古月娜将第二项的两段新增部分分别圈出，卢本韦从原坏部核它们各自的增量。蔡须坤专查那些本该不动的位置，霍雨浩最后逐列比较。各项虽然可以更长，仍位于原式之下。

三列的结构落到纸上，卢本韦反而坐直了一点。

``越短越容易让我想直接念结果。''

蔡须坤将三列分别补齐。根回到最初的零列，坏部有两列；最高非零行是第三行，前两行的阶差并不相同。她先核第一行，再核第二行，没有把整列统一加上一个数。展开第二项仍以同一个原式为出发点，第二个新增段使用相应倍数的阶差，保留段不因此再加一次。

古月娜检查完，让她用没有颜色的纸重写。蔡须坤写到最末，霍雨浩问起第一行的零是否也要省。她摇头，把起始零列留在原处，只将每列末尾补齐用的零藏回简写。

``省掉的是显示的位置，规则里仍按补齐的列算。''

卢本韦逐个核列数，发现三列在展开以后可以变得很长。外形变长并没有使基本序列的项超过原式；相反，这些有限项都位于原式之下。此前按阅读顺序一路分析到这里，是在组织更高节点的对应，而对某个节点取基本序列的项，是朝它下面的有限结构看。两种方向同时出现在笔记里，不能混为一次不断变大的操作。

霍雨浩将上山的箭头与展开的箭头画在不同页边。这一次他没有写许多新符号，只把当前三列与它的前两项放在一起。蔡须坤看明白后，也拿自己的脚下作比较：他们站上更高的里程桩，仍可以回头分析较低的结构，回头本身没有让那较低结构变成更高。

``但会让我们知道脚下为什么站得住。''她说。

夏夜星空让四人先在最后一级下方轮流试接。蔡须坤的动作很慢；她能分辨原来两次纠缠的回牵，便把呼吸留给它们各自的位置。霍雨浩试完，将手按在一旁的石栏上，掌心的汗沿石缝积成一道细线。最上面那道横纹还没有亮，四人先将各自的结果带回纸上。

古月娜最后试过，没有先看他的记录，而是将自己的根与阶差说完。四份答案对上，卢本韦才靠着石沿坐下，等呼吸平稳，再准备最后一遍完整运转。

这一处对应下标里的平方，真正的多重提升还会接着把乘法送到更深的位置。他们先将眼前这页收好，留出一会儿喘息的工夫，再去看三列之后还能接些什么。

这一组的末端是 \allowbreak{}\NovelInline{I05956}{\psi(\Omega_{\omega^{2}})}\allowbreak{} ，对应了一个非常简洁的三项BMS。蔡须坤将下一个三列结构留在空页上，准备重新追根。

一百六十米的石纹亮起时，四人都没有立刻往上走。

蔡须坤用手背抹去下颌的汗，看了一遍眼前的三列，又回头看下面一百二十米的转弯。卢本韦站到她身侧，脚尖刚放稳便蹲了下来，将水囊放在两人之间。

霍雨浩在山图上记下这一处。

\NovelDisplay{M01837R01}{%
(0)(1,1,1)(2,1,1)\longleftrightarrow\quad160\text{ 米}%
}

``这里到了。''他说。

古月娜将刚才各自求出的展开夹到同一页后面，没有再让大家重算一遍。她指着末列第三行的一，确认列号与先前的根都已记全，才让霍雨浩合上这一页。

蔡须坤取出旧预演页，将今天独立求出的结果放在旁边。同样的三列，昨日还能凭熟悉直接念出对应，此时却已能逐项说清来路。她在页角添上日期，收住炭条。

夏夜星空确认四人的运转与记录对得上，才让后面的同伴按各自核过的次序陆续上台。唐三收回一段蓝银草，奥斯卡把还热着的补给递了上来。

\NovelDisplay{M01838R01}{%
(0)(1,1,1)(2,1,1)(1,1,1)=\psi(\Omega_{\omega^{2}}+\Omega_{\omega})%
}

标准开局。

\NovelDisplay{M01839R01}{%
(0)(1,1,1)(2,1,1)(1,1,1)(2,1)(3,2)=\psi(\Omega_{\omega^{2}}+\psi_{\Omega_{\omega+1}}(\Omega_{\omega+1}))%
}

\NovelDisplay{M01840R01}{%
(0)(1,1,1)(2,1,1)(1,1,1)(2,1)(3,2)(4,2)=\psi(\Omega_{\omega^{2}}+\psi_{\Omega_{\omega+1}}(\Omega_{\omega+1}^{2}))%
}

\NovelDisplay{M01841R01}{%
(0)(1,1,1)(2,1,1)(1,1,1)(2,1)(3,2,1)=\psi(\Omega_{\omega^{2}}+\psi_{\Omega_{\omega+1}}(\Omega_{\omega\times2}))%
}

\NovelDisplay{M01842R01}{%
(0)(1,1,1)(2,1,1)(1,1,1)(2,1)(3,2,1)(4,2)%
}
\NovelDisplay{M01842R02}{%
\quad (5,3,1)=\psi(\Omega_{\omega^{2}}+\psi_{\Omega_{\omega+1}}(\Omega_{\omega\times3}))%
}

\NovelDisplay{M01843R01}{%
(0)(1,1,1)(2,1,1)(1,1,1)(2,1)(3,2,1)(4,2,1)=\psi(\Omega_{\omega^{2}}+\psi_{\Omega_{\omega+1}}(\Omega_{\omega^{2}}))%
}

\NovelDisplay{M01844R01}{%
(0)(1,1,1)(2,1,1)(1,1,1)(2,1)(3,2,1)(4,2,1)(3,2)(4,1)(5,2,1)(6,2,1)%
}
\NovelDisplay{M01844R02}{%
\quad =\psi(\Omega_{\omega^{2}}+\Omega_{\omega+1}\times\psi_{\Omega_{\omega+1}}(\Omega_{\omega^{2}}))%
}

\NovelDisplay{M01845R01}{%
(0)(1,1,1)(2,1,1)(1,1,1)(2,1)(3,2,1)%
}
\NovelDisplay{M01845R02}{%
\quad (4,2,1)(3,2)(4,2)=\psi(\Omega_{\omega^{2}}+\Omega_{\omega+1}^{2})%
}

\NovelDisplay{M01846R01}{%
(0)(1,1,1)(2,1,1)(1,1,1)(2,1)(3,2,1)(4,2,1)(3,2,1)=\psi(\Omega_{\omega^{2}}+\Omega_{\omega\times2})%
}

\NovelDisplay{M01847R01}{%
(0)(1,1,1)(2,1,1)(1,1,1)(2,1)(3,2,1)(4,2,1)(3,2,1)(4,2)(5,3)=\psi(\Omega_{\omega^{2}}%
}
\NovelDisplay{M01847R02}{%
\quad +\psi_{\Omega_{\omega\times2+1}}(\Omega_{\omega\times2+1}))%
}

\NovelDisplay{M01848R01}{%
(0)(1,1,1)(2,1,1)(1,1,1)(2,1)(3,2,1)(4,2,1)(3,2,1)(4,2)(5,3,1)(6,3,1)%
}
\NovelDisplay{M01848R02}{%
\quad =\psi(\Omega_{\omega^{2}}+\psi_{\Omega_{\omega\times2+1}}(\Omega_{\omega^{2}}))%
}

\NovelDisplay{M01849R01}{%
(0)(1,1,1)(2,1,1)(1,1,1)(2,1)(3,2,1)(4,2,1)(3,2,1)(4,2)(5,3,1)(6,3,1)%
}
\NovelDisplay{M01849R02}{%
\quad (5,3,1)=\psi(\Omega_{\omega^{2}}+\Omega_{\omega\times3})%
}

\NovelDisplay{M01850R01}{%
(0)(1,1,1)(2,1,1)(1,1,1)(2,1,1)=\psi(\Omega_{\omega^{2}}\times2)%
}

从这一段开头到这里的过程，都只相当于是在做加法。是时候做乘法了。

\NovelDisplay{M01851R01}{%
(0)(1,1,1)(2,1,1)(2)=\psi(\Omega_{\omega^{2}}\times\omega)%
}

\NovelDisplay{M01852R01}{%
(0)(1,1,1)(2,1,1)(2,1)=\psi(\Omega_{\omega^{2}}\times\Omega)%
}

霍雨浩指向这个熟悉的结构，也即 \allowbreak{}\NovelInline{I05957}{\Omega_{\alpha}\times\Omega}\allowbreak{} 。就眼前已经核过的前缀，四人已经确认，在式子后面添加 \allowbreak{}\NovelInline{I05958}{(1,1,1)}\allowbreak{} 可以把 \allowbreak{}\NovelInline{I05959}{\Omega}\allowbreak{} 提升到 \allowbreak{}\NovelInline{I05960}{\Omega_{\omega}}\allowbreak{} ，而 \allowbreak{}\NovelInline{I05961}{(3,2,1)}\allowbreak{} , \allowbreak{}\NovelInline{I05962}{(5,3,1)}\allowbreak{} 等项也有提升效应，它们可以充当第二段、第三段的提升。但怎样让这个 \allowbreak{}\NovelInline{I05963}{\Omega}\allowbreak{} 提升到 \allowbreak{}\NovelInline{I05964}{\Omega_{\omega^2}}\allowbreak{} 呢？

\NovelDisplay{M01853R01}{%
(0)(1,1,1)(2,1,1)(2,1)(1,1,1)=\psi(\Omega_{\omega^{2}}\times\Omega_{\omega})%
}

\NovelDisplay{M01854R01}{%
(0)(1,1,1)(2,1,1)(2,1)(1,1,1)(1,1,1)=\psi(\Omega_{\omega^{2}}\times\Omega_{\omega}+\Omega_{\omega})%
}

\NovelDisplay{M01855R01}{%
(0)(1,1,1)(2,1,1)(2,1)(1,1,1)(2,1)=\psi(\Omega_{\omega^{2}}\times\Omega_{\omega}+\Omega_{\omega}\times\Omega)%
}

\NovelDisplay{M01856R01}{%
(0)(1,1,1)(2,1,1)(2,1)(1,1,1)(2,1)(1,1,1)=\psi(\Omega_{\omega^{2}}\times\Omega_{\omega}+\Omega_{\omega}^{2})%
}

\NovelDisplay{M01857R01}{%
(0)(1,1,1)(2,1,1)(2,1)(1,1,1)(2,1)(3,2)%
}
\NovelDisplay{M01857R02}{%
\quad =\psi(\Omega_{\omega^{2}}\times\Omega_{\omega}+\psi_{\Omega_{\omega+1}}(\Omega_{\omega+1}))%
}

\NovelDisplay{M01858R01}{%
(0)(1,1,1)(2,1,1)(2,1)(1,1,1)(2,1)(3,2,1)%
}
\NovelDisplay{M01858R02}{%
\quad (4,2,1)=\psi(\Omega_{\omega^{2}}\times\Omega_{\omega}+\psi_{\Omega_{\omega+1}}(\Omega_{\omega^{2}}))%
}

\NovelDisplay{M01859R01}{%
(0)(1,1,1)(2,1,1)(2,1)(1,1,1)(2,1)(3,2,1)(4,2,1)(4,1)(1,1,1)=\psi(\Omega_{\omega^{2}}%
}
\NovelDisplay{M01859R02}{%
\quad \times\Omega_{\omega}+\psi_{\Omega_{\omega+1}}(\Omega_{\omega^{2}}\times\Omega_{\omega}))%
}

注意这里的 \allowbreak{}\NovelInline{I05965}{(4,1)}\allowbreak{} 。它作为替身出现，但是本身代表的 \allowbreak{}\NovelInline{I05966}{\Omega}\allowbreak{} 已经由后面的 \allowbreak{}\NovelInline{I05967}{(1,1,1)}\allowbreak{} 提升到了 \allowbreak{}\NovelInline{I05968}{\Omega_\omega}\allowbreak{} 。

\NovelDisplay{M01860R01}{%
(0)(1,1,1)(2,1,1)(2,1)(1,1,1)(2,1)(3,2,1)(4,2,1)(4,1)(3,2)(4,1)(5,2,1)%
}
\NovelDisplay{M01860R02}{%
\quad (6,2,1)(6,1)(1,1,1)=\psi(\Omega_{\omega^{2}}\times\Omega_{\omega}+\Omega_{\omega+1}\times\psi_{\Omega_{\omega+1}}(\Omega_{\omega^{2}}\times\Omega_{\omega}))%
}

\NovelDisplay{M01861R01}{%
(0)(1,1,1)(2,1,1)(2,1)(1,1,1)(2,1)(3,2,1)(4,2,1)(4,1)(3,2)(4,2)%
}
\NovelDisplay{M01861R02}{%
\quad =\psi(\Omega_{\omega^{2}}\times\Omega_{\omega}+\Omega_{\omega+1}^{2})%
}

\NovelDisplay{M01862R01}{%
(0)(1,1,1)(2,1,1)(2,1)(1,1,1)(2,1)(3,2,1)(4,2,1)(4,1)(3,2)(4,3,1)%
}
\NovelDisplay{M01862R02}{%
\quad (5,3,1)(5,1)(4,3)(5,3)=\psi(\Omega_{\omega^{2}}\times\Omega_{\omega}+\Omega_{\omega+2}^{2})%
}

\NovelDisplay{M01863R01}{%
(0)(1,1,1)(2,1,1)(2,1)(1,1,1)(2,1)(3,2,1)%
}
\NovelDisplay{M01863R02}{%
\quad (4,2,1)(4,1)(3,2,1)=\psi(\Omega_{\omega^{2}}\times\Omega_{\omega}+\Omega_{\omega\times2})%
}

\NovelDisplay{M01864R01}{%
(0)(1,1,1)(2,1,1)(2,1)(1,1,1)(2,1)(3,2,1)(4,2,1)(4,1)(3,2,1)(4,2)(5,3)%
}
\NovelDisplay{M01864R02}{%
\quad =\psi(\Omega_{\omega^{2}}\times\Omega_{\omega}+\psi_{\Omega_{\omega\times2+1}}(\Omega_{\omega\times2+1}))%
}

\NovelDisplay{M01865R01}{%
(0)(1,1,1)(2,1,1)(2,1)(1,1,1)(2,1)(3,2,1)(4,2,1)(4,1)(3,2,1)(4,2)%
}
\NovelDisplay{M01865R02}{%
\quad (5,3,1)(6,3,1)(6,1)(1,1,1)=\psi(\Omega_{\omega^{2}}\times\Omega_{\omega}+\psi_{\Omega_{\omega\times2+1}}(\Omega_{\omega^{2}}\times\Omega_{\omega}))%
}

\NovelDisplay{M01866R01}{%
(0)(1,1,1)(2,1,1)(2,1)(1,1,1)(2,1)(3,2,1)(4,2,1)(4,1)(3,2,1)(4,2)%
}
\NovelDisplay{M01866R02}{%
\quad (5,3,1)(6,3,1)(6,1)(5,3)(6,3)=\psi(\Omega_{\omega^{2}}\times\Omega_{\omega}+\Omega_{\omega\times2+1}^{2})%
}

\NovelDisplay{M01867R01}{%
(0)(1,1,1)(2,1,1)(2,1)(1,1,1)(2,1)(3,2,1)(4,2,1)(4,1)(3,2,1)(4,2)%
}
\NovelDisplay{M01867R02}{%
\quad (5,3,1)(6,3,1)(6,1)(5,3,1)=\psi(\Omega_{\omega^{2}}\times\Omega_{\omega}+\Omega_{\omega\times3})%
}

\NovelDisplay{M01868R01}{%
(0)(1,1,1)(2,1,1)(2,1)(1,1,1)(2,1)(3,2,1)(4,2,1)(4,1)(3,2,1)(4,2,1)%
}
\NovelDisplay{M01868R02}{%
\quad =\psi(\Omega_{\omega^{2}}\times\Omega_{\omega}+\Omega_{\omega^{2}})%
}

可以看作 \allowbreak{}\NovelInline{I05969}{(3,2,1)(4,2,1)}\allowbreak{} 这两项作为一个整体共同完成了``添加 \allowbreak{}\NovelInline{I05970}{\Omega_{\omega^2}}\allowbreak{} ''的功能。本来这个功能应该由 \allowbreak{}\NovelInline{I05971}{(1,1,1)(2,1,1)}\allowbreak{} 完成，但被替身 \allowbreak{}\NovelInline{I05972}{(4,1)}\allowbreak{} 阻挡了。

\NovelDisplay{M01869R01}{%
(0)(1,1,1)(2,1,1)(2,1)(1,1,1)(2,1)(3,2,1)(4,2,1)(4,1)(3,2,1)(4,2,1)%
}
\NovelDisplay{M01869R02}{%
\quad (3,2,1)(4,2,1)=\psi(\Omega_{\omega^{2}}\times\Omega_{\omega}+\Omega_{\omega^{2}}\times2)%
}

\NovelDisplay{M01870R01}{%
(0)(1,1,1)(2,1,1)(2,1)(1,1,1)(2,1)(3,2,1)(4,2,1)(4,1)(3,2,1)(4,2,1)%
}
\NovelDisplay{M01870R02}{%
\quad (4,1)(1,1,1)=\psi(\Omega_{\omega^{2}}\times\Omega_{\omega}\times2)%
}

\NovelDisplay{M01871R01}{%
(0)(1,1,1)(2,1,1)(2,1)(1,1,1)(2,1)(3,2,1)%
}
\NovelDisplay{M01871R02}{%
\quad (4,2,1)(4,1)(4,1)(1,1,1)=\psi(\Omega_{\omega^{2}}\times\Omega_{\omega}^{2})%
}

\NovelDisplay{M01872R01}{%
(0)(1,1,1)(2,1,1)(2,1)(1,1,1)(2,1)(3,2,1)(4,2,1)%
}
\NovelDisplay{M01872R02}{%
\quad (4,1)(5,2)=\psi(\Omega_{\omega^{2}}\times\psi_{\Omega_{\omega+1}}(\Omega_{\omega+1}))%
}

\NovelDisplay{M01873R01}{%
(0)(1,1,1)(2,1,1)(2,1)(1,1,1)(2,1)(3,2,1)(4,2,1)(4,1)(5,2,1)(6,2,1)%
}
\NovelDisplay{M01873R02}{%
\quad =\psi(\Omega_{\omega^{2}}\times\psi_{\Omega_{\omega+1}}(\Omega_{\omega^{2}}))%
}

\NovelDisplay{M01874R01}{%
(0)(1,1,1)(2,1,1)(2,1)(1,1,1)(2,1)%
}
\NovelDisplay{M01874R02}{%
\quad (3,2,1)(4,2,1)(4,2)=\psi(\Omega_{\omega^{2}}\times\Omega_{\omega+1})%
}

过了一百六十米以后，蔡须坤又一次误判了自己已经会的部分。

她看见旧片段重新出现，便以为后面的提升只需逐段相加。古月娜没有立刻纠正，先让她指出当前那个起中介作用的结构，自己最后会走到哪里。

蔡须坤追到一半，停住了。

``它完成原来那次作用之后，还能继续被后面提升。''

霍雨浩将刚才几组中间式连起来。前面连续两段、三段的情形，是一段完成后接下一段；眼前逐渐显出的情形中，承担中间作用的结构本身，又接受了后续提升。被推动的对象没有在交接处退出，而是进入下一轮。于是要追的不只是每一段的结果，还有中间那个对象怎样改变，再把改变传给它所处的整层结构。

卢本韦将手伸到两张纸之间，本想将它们按前后分开，最后却将后一张压到前一张的边缘。

``这回不能只摆成一排了。''

他们没有因为想明白这个比喻便直接写结论。古月娜把末列的根重找一遍，蔡须坤确认好部哪些列保留，霍雨浩生成前两项，卢本韦再指出其中哪个结构接过了此前的作用。四人的说明都落在已列出的中间式上，才看得见嵌套到底从哪一步开始。

山上的压力也换了受阻的位置。蔡须坤将外面那层稳住以后，里面仍有一处把气息折回。她不能只把外面的力气再添大些，便请霍雨浩只记录两次停顿相隔的先后。古月娜检查她每次接起的前缀，卢本韦负责提醒她何时已经把同一个增量重复加到了保留段。

``停一下。''卢本韦说第三遍时，蔡须坤终于没有急着继续。

她退回半步，将手掌贴在冷石上。肩背并未受伤，只是连续尝试使她又想用熟练的运转盖住没懂的结构。她把第三次的错误写完，留到明天再核。霍雨浩也合上自己的本子，未曾因为还差一个明显的节点就劝她硬撑。

回营路上，卢本韦帮她背了纸袋，古月娜走在较低的一侧。蔡须坤空出的手扶着绳索，心里仍在想那个继续受提升的中间结构。等她终于能将两种情形分别说清，霍雨浩才在旁边答了一声，接着问它下一步怎样落到已写的公式里。

那一问，留到了新一页。

\allowbreak{}\NovelInline{I05973}{(4,1)}\allowbreak{} 的使命完成。但生成的 \allowbreak{}\NovelInline{I05974}{(4,2)}\allowbreak{} ，它又能被 \allowbreak{}\NovelInline{I05975}{(3,2,1)}\allowbreak{} 提升到 \allowbreak{}\NovelInline{I05976}{\Omega_{\omega\times2}}\allowbreak{} 。

一百六十米留在了这一册本子的中间。

回到营地，蔡须坤将那页角折起，后面仍有许多新写的列。纸面上，第二段的根与第一段的根各有位置；身体里，曾经揉成一团的两次牵扯，也已经能分别辨认。

卢本韦让她先收起本子，自己在炉边的空纸上接完白天留下的那一段。古月娜没有提醒，霍雨浩也把答案扣在掌心下。等她独自说完，四人才交换记录。

他们没有谁走出压力，只是越来越少用含混的名字遮住它。

收好最后一张错页，蔡须坤把湿纸放在炉边烘。奥斯卡端来热汤，诺当替大家挪开地上的杂物。她喝了半碗，才发觉握笔的手一直在抖。卢本韦看见了，将勺子递过来，顺手把最后一张还没干的纸翻了一面。

\chapter{第一百一十章：把整座山写进下标}\label{ux7b2cux4e00ux767eux4e00ux5341ux7ae0ux628aux6574ux5ea7ux5c71ux5199ux8fdbux4e0bux6807}

一百六十米的石桩留在了身后。

霍雨浩把最近几天的纸摆成一列，最初一张写着一个下标，最后一张的下标里已经有了乘法。卢本韦看了半晌，忽然将写得最满的那张抽了出来。

``要是只看纸有多长，这张该比前面的都高。''

``可它是展开出来的较小项。''蔡须坤将那张纸放回原来的位置，``先比列，再比你正在读的是哪一层。''

古月娜在空白处写下第一对括号。霍雨浩的精神力刚要延过去，她便伸手拦了一下。

``等我自己追完。''

她沿父项链走得很慢。肩头的银纹被压力压成一线，直到她将中途一列省掉的零补回，才重新分出两层。霍雨浩记下这一处，等她点头，再把自己的观察接到纸边。

蔡须坤坐在另一侧，试着把下面那层已经能熟练调动的魂力放慢。压力从指数的接缝传到下标的接缝，身体的感觉很清楚，名字却不能凭感觉填。她在纸上留出一片空白，先写末列，再写祖先。

卢本韦看了她一眼，将自己原先想写的答案也暂时收住。

``先给下标里的每一步找到来路。''

夏夜星空把被风吹散的纸压在石边，没有替他们命名新的层级。通往二百二十米的路还没有走完，远处五百米的横线仍在雾中。四人将山图折到只露出当前这一段，继续向上。

过了一百六十米，霍雨浩最先弄混的是下标里的加法。

他读得很熟，笔也没有停，写完才发现古月娜没有跟着往上接。她把纸抬起来，指尖停在一个小得几乎被他忽略的加号旁边。

``你把它移到外面去了。''

卢本韦沿着两对括号分别画线。蔡须坤照两种位置各试一次，清楚地感觉到回牵来自不同接缝。霍雨浩把共享收回，重新独立演算，直到纸上的括号也与两种反应一一对应。

他们在这块横台多留了一夜。翌日上山，谁也没有再把那两张纸合成一张。

\NovelDisplay{M01875R01}{%
(0)(1,1,1)(2,1,1)(2,1)(1,1,1)(2,1)(3,2,1)(4,2,1)(4,2)(3,2)(4,3,1)%
}
\NovelDisplay{M01875R02}{%
\quad (5,3,1)(5,2)(4,3,1)(5,3,1)=\psi(\Omega_{\omega^{2}}\times\Omega_{\omega+1}+\Omega_{\omega^{2}})%
}

下标开始变长，卢本韦特意挑出一条会产生不同提升阈值的式子。蔡须坤从坏根一路追过去，遇到绕行的位置便停下，由古月娜检查为什么不能把那一行整块抬高。

\NovelDisplay{M01876R01}{%
(0)(1,1,1)(2,1,1)(2,1)(1,1,1)(2,1)(3,2,1)(4,2,1)(4,2)\textcolor[HTML]{B45309}{(3,2)}\textcolor[HTML]{B91C1C}{(4,3,1)}%
}
\NovelDisplay{M01876R02}{%
\quad \textcolor[HTML]{B91C1C}{(5,3,1)}\textcolor[HTML]{B91C1C}{(5,\textcolor[HTML]{6B7280}{2})}\textcolor[HTML]{B91C1C}{(4,3,1)}\textcolor[HTML]{7E22CE}{(5,3,1)}%
}

\NovelDisplay{M01877R01}{%
\text{原列数}=15, h=3, r=10,%
}
\NovelDisplay{M01877R02}{%
\text{坏根}=\textcolor[HTML]{B45309}{(3,2)},\text{坏部列数}=5,%
}
\NovelDisplay{M01877R03}{%
(\delta_1,\delta_2,\delta_3)=(2,1,0).%
}

\NovelDisplay{M01878R01}{%
\text{按坏部列序排列的提升阈值：}%
}
\NovelDisplay{M01878R02}{%
(2,2,2,1,2),%
}
\NovelDisplay{M01878R03}{%
\text{第 }j\text{ 列只允许前 }\tau_j\text{ 行应用阶差。}%
}

\NovelDisplay{M01879R01}{%
\text{原第 }10\text{ 列：}(3,2)%
}
\NovelDisplay{M01879R02}{%
\longmapsto(3+2q,2+q), q\ge1.%
}

\NovelDisplay{M01880R01}{%
\text{原第 }13\text{ 列：}(5,2)%
}
\NovelDisplay{M01880R02}{%
\longmapsto(5+2q,2), q\ge1.%
}

\NovelDisplay{M01881R01}{%
\text{原第 }14\text{ 列：}(4,3,1)%
}
\NovelDisplay{M01881R02}{%
\longmapsto(4+2q,3+q,1), q\ge1.%
}

古月娜把阈值较低的那几列单独指出。阶差虽已求出，原坏部各列允许改变的行数仍不同；下方的灰色坐标保留原值，其余有资格的坐标按新增段的正整数编号 q 应用相应倍数的阶差。

\NovelDisplay{M01882R01}{%
\bigl((0)(1,1,1)(2,1,1)(2,1)(1,1,1)(2,1)(3,2,1)(4,2,1)(4,2)(3,2)(4,3,1)%
}
\NovelDisplay{M01882R02}{%
\quad (5,3,1)(5,2)(4,3,1)(5,3,1)\bigr)[0]=(0)(1,1,1)(2,1,1)(2,1)(1,1,1)(2,1)%
}
\NovelDisplay{M01882R03}{%
\quad (3,2,1)(4,2,1)(4,2)\textcolor[HTML]{B45309}{(3,2)}\textcolor[HTML]{B91C1C}{(4,3,1)}\textcolor[HTML]{B91C1C}{(5,3,1)}\textcolor[HTML]{B91C1C}{(5,\textcolor[HTML]{6B7280}{2})}\textcolor[HTML]{B91C1C}{(4,3,1)}%
}

\NovelDisplay{M01883R01}{%
\bigl((0)(1,1,1)(2,1,1)(2,1)(1,1,1)(2,1)(3,2,1)(4,2,1)(4,2)(3,2)(4,3,1)%
}
\NovelDisplay{M01883R02}{%
\quad (5,3,1)(5,2)(4,3,1)(5,3,1)\bigr)[1]=(0)(1,1,1)(2,1,1)(2,1)(1,1,1)(2,1)%
}
\NovelDisplay{M01883R03}{%
\quad (3,2,1)(4,2,1)(4,2)\textcolor[HTML]{B45309}{(3,2)}\textcolor[HTML]{B91C1C}{(4,3,1)}\textcolor[HTML]{B91C1C}{(5,3,1)}\textcolor[HTML]{B91C1C}{(5,\textcolor[HTML]{6B7280}{2})}%
}
\NovelDisplay{M01883R04}{%
\quad \textcolor[HTML]{B91C1C}{(4,3,1)}\textcolor[HTML]{1D4ED8}{(5,3)}\textcolor[HTML]{1D4ED8}{(6,4,1)}\textcolor[HTML]{1D4ED8}{(7,4,1)}\textcolor[HTML]{1D4ED8}{(7,2)}\textcolor[HTML]{1D4ED8}{(6,4,1)}%
}

\NovelDisplay{M01884R01}{%
\bigl((0)(1,1,1)(2,1,1)(2,1)(1,1,1)(2,1)(3,2,1)(4,2,1)(4,2)(3,2)(4,3,1)%
}
\NovelDisplay{M01884R02}{%
\quad (5,3,1)(5,2)(4,3,1)(5,3,1)\bigr)[2]=(0)(1,1,1)(2,1,1)(2,1)(1,1,1)(2,1)%
}
\NovelDisplay{M01884R03}{%
\quad (3,2,1)(4,2,1)(4,2)\textcolor[HTML]{B45309}{(3,2)}\textcolor[HTML]{B91C1C}{(4,3,1)}\textcolor[HTML]{B91C1C}{(5,3,1)}\textcolor[HTML]{B91C1C}{(5,\textcolor[HTML]{6B7280}{2})}\textcolor[HTML]{B91C1C}{(4,3,1)}\textcolor[HTML]{1D4ED8}{(5,3)}\textcolor[HTML]{1D4ED8}{(6,4,1)}%
}
\NovelDisplay{M01884R04}{%
\quad \textcolor[HTML]{1D4ED8}{(7,4,1)}\textcolor[HTML]{1D4ED8}{(7,2)}\textcolor[HTML]{1D4ED8}{(6,4,1)}\textcolor[HTML]{15803D}{(7,4)}\textcolor[HTML]{15803D}{(8,5,1)}\textcolor[HTML]{15803D}{(9,5,1)}\textcolor[HTML]{15803D}{(9,2)}\textcolor[HTML]{15803D}{(8,5,1)}%
}

\NovelDisplay{M01885R01}{%
\text{第零项列数}=14,%
}
\NovelDisplay{M01885R02}{%
\text{第一项列数}=14+5=19,%
}
\NovelDisplay{M01885R03}{%
\text{第二项列数}=14+2\times5=24.%
}

卢本韦数完列，将第一项中的蓝段与第二项中的蓝段逐列重合。绿色段另由原坏部按两倍阶差生成，原来的保留段没有再受一次提升。蔡须坤核坐标，古月娜与霍雨浩最后检查各有限项都小于原式。

\NovelDisplay{M01886R01}{%
(0)(1,1,1)(2,1,1)(2,1)(1,1,1)(2,1)(3,2,1)(4,2,1)(4,2)(3,2)(4,3,1)%
}
\NovelDisplay{M01886R02}{%
\quad (5,3,1)(5,3)=\psi(\Omega_{\omega^{2}}\times\Omega_{\omega+2})%
}

\NovelDisplay{M01887R01}{%
(0)(1,1,1)(2,1,1)(2,1)(1,1,1)(2,1)(3,2,1)%
}
\NovelDisplay{M01887R02}{%
\quad (4,2,1)(4,2)(3,2,1)=\psi(\Omega_{\omega^{2}}\times\Omega_{\omega\times2})%
}

就像这样。

\NovelDisplay{M01888R01}{%
(0)(1,1,1)(2,1,1)(2,1)(1,1,1)(2,1)(3,2,1)(4,2,1)(4,2)(3,2,1)(4,2)%
}
\NovelDisplay{M01888R02}{%
\quad (5,3,1)(6,3,1)(6,2)(5,3,1)(6,3,1)=\psi(\Omega_{\omega^{2}}\times\Omega_{\omega\times2}+\Omega_{\omega^{2}})%
}

现在又召唤了替身 \allowbreak{}\NovelInline{I05977}{(6,2)}\allowbreak{} 。

\NovelDisplay{M01889R01}{%
(0)(1,1,1)(2,1,1)(2,1)(1,1,1)(2,1)(3,2,1)(4,2,1)(4,2)(3,2,1)(4,2)%
}
\NovelDisplay{M01889R02}{%
\quad (5,3,1)(6,3,1)(6,3)=\psi(\Omega_{\omega^{2}}\times\Omega_{\omega\times2+1})%
}

\NovelDisplay{M01890R01}{%
(0)(1,1,1)(2,1,1)(2,1)(1,1,1)(2,1)(3,2,1)(4,2,1)(4,2)(3,2,1)(4,2)%
}
\NovelDisplay{M01890R02}{%
\quad (5,3,1)(6,3,1)(6,3)(5,3,1)=\psi(\Omega_{\omega^{2}}\times\Omega_{\omega\times3})%
}

\allowbreak{}\NovelInline{I05978}{(6,2)}\allowbreak{} 达成使命之后，生成的 \allowbreak{}\NovelInline{I05979}{(6,3)}\allowbreak{} 又可以被 \allowbreak{}\NovelInline{I05980}{(5,3,1)}\allowbreak{} 提升到 \allowbreak{}\NovelInline{I05981}{\Omega_{\omega\times3}}\allowbreak{} 。像这样一直持续下去，最后是------

\NovelDisplay{M01891R01}{%
(0)(1,1,1)(2,1,1)(2,1)(1,1,1)(2,1,1)=\psi(\Omega_{\omega^{2}}^{2})%
}

这便是``双重提升效应''的全貌。 \allowbreak{}\NovelInline{I05982}{(1,1,1)}\allowbreak{} 把 \allowbreak{}\NovelInline{I05983}{(2,1)}\allowbreak{} 提升到 \allowbreak{}\NovelInline{I05984}{\Omega_\omega}\allowbreak{} 之后，紧随其后的 \allowbreak{}\NovelInline{I05985}{(2,1,1)}\allowbreak{} 又把它进一步提升到 \allowbreak{}\NovelInline{I05986}{\Omega_{\omega^2}}\allowbreak{} 。造成这一现象的原因，根据四人刚才的分析，是每一步形成的替身自己，在完成使命之后，还会吃到一段提升，这相当于两个提升``套在了一起''，作用大增。这里的 \allowbreak{}\NovelInline{I05987}{(1,1,1)(2,1,1)}\allowbreak{} 可以看作构成了第一个``提升链''。在后面的分析当中，``提升链''可能会变得相当长。

\NovelDisplay{M01892R01}{%
(0)(1,1,1)(2,1,1)(2,1)(1,1,1)(2,1,1)(1,1,1)(2,1,1)=\psi(\Omega_{\omega^{2}}^{2}+\Omega_{\omega^{2}})%
}

回归主线。现在又是 \allowbreak{}\NovelInline{I05988}{(1,1,1)(2,1,1)}\allowbreak{} 共同完成``添加 \allowbreak{}\NovelInline{I05989}{\Omega_{\omega^2}}\allowbreak{} ''的作用了。

\NovelDisplay{M01893R01}{%
(0)(1,1,1)(2,1,1)(2,1)(1,1,1)(2,1,1)(2,1)(1,1,1)(2,1,1)=\psi(\Omega_{\omega^{2}}^{2}\times2)%
}

\NovelDisplay{M01894R01}{%
(0)(1,1,1)(2,1,1)(2,1)(2,1)(1,1,1)(2,1,1)=\psi(\Omega_{\omega^{2}}^{3})%
}

\NovelDisplay{M01895R01}{%
(0)(1,1,1)(2,1,1)(2,1)(3,1)=\psi(\Omega_{\omega^{2}}^{\Omega})%
}

\NovelDisplay{M01896R01}{%
(0)(1,1,1)(2,1,1)(2,1)(3,1)(1,1,1)=\psi(\Omega_{\omega^{2}}^{\Omega_{\omega}})%
}

\NovelDisplay{M01897R01}{%
(0)(1,1,1)(2,1,1)(2,1)(3,1)(1,1,1)(2,1,1)=\psi(\Omega_{\omega^{2}}^{\Omega_{\omega^{2}}})%
}

两层结构真正衔接起来以后，古月娜将前面留下的短线翻了出来。

``现在可以拿连续两段来比了。''

卢本韦读出此前的过程：第一段完成，保留结果，第二段再接着推动。蔡须坤读眼前这一组：中间起替身作用的结构完成原来的任务后，自己又接到一次提升，于是下一次交接所用的对象也变了。嵌套使这一组能够继续深入，不只是把前一段终点在旁边再写一遍。

霍雨浩请三人从最早一个可检查的中间式开始，把那个对象每次怎样换位指出来。他们没有为所有连续项再作口头总结，只在它交出作用、又接受新作用的地方停下。其余有限增长由各人按前页规则独立核过，连接处则交换着读。

``所以说双重，讲的是作用的组织。''蔡须坤说，``不是第三行多出来一个一，就一律叫双重。''

古月娜将左端的提升链圈出。链的识别来自这组完整对应与刚才追过的父项，不能凭列数另立一条通用口诀。她又请霍雨浩找回一个没有完成同样嵌套的旧式，放在旁边作区别。两页都有相近的片段，前缀与后续怎样接起来却不同。

卢本韦沿着那条区别画线，画到一半忽然笑了。

``我昨天还以为，只要再多抄几个相同的列，就能把路直接抄到山顶。''

蔡须坤看向她磨得很短的炭条。``那你得先问它够不够。''

两人笑过，将纸拿到台阶前。霍雨浩只共享周围的情况，先让古月娜独自尝试。她第一次在内层停住，没有向外再添魂力，而是退回自己原先用错前缀的位置；第二次把那个中间对象接对，外层已经稳住的部分便没有再被一起扯乱。

蔡须坤随后走了一遍，停顿与古月娜并不完全同拍。他们没有为了让感觉一致而强行调整，只把各自漏掉的分析说清。相同的公式要求相同的结构解释，身体在疲倦、呼吸上的差别仍然允许留下。压力是暴露问题的条件，无法替他们自动判断每个序数对应。

傍晚下山时，四人在旧石桩旁歇了一会儿。卢本韦将连续两段与嵌套两重的笔记放进同一册，分别折好页角。以后再听见相似的名字，他们便能翻回具体中间式，知道到底在比较哪一种组织。

\NovelDisplay{M01898R01}{%
(0)(1,1,1)(2,1,1)(2,1)(3,2)=\psi(\Omega_{\omega^{2}+1})%
}

\allowbreak{}\NovelInline{I05990}{\Omega_{\omega^2}}\allowbreak{} 的指数塔。又是 \allowbreak{}\NovelInline{I05991}{(2,1)(3,2)}\allowbreak{} ，这也算一个规律了。

\NovelDisplay{M01899R01}{%
(0)(1,1,1)(2,1,1)(2,1)(3,2)(3,2)=\psi(\Omega_{\omega^{2}+1}\times2)%
}

\NovelDisplay{M01900R01}{%
(0)(1,1,1)(2,1,1)(2,1)(3,2)(4,1)=\psi(\Omega_{\omega^{2}+1}\times\Omega)%
}

\NovelDisplay{M01901R01}{%
(0)(1,1,1)(2,1,1)(2,1)(3,2)(4,1)(1,1,1)(2,1,1)=\psi(\Omega_{\omega^{2}+1}\times\Omega_{\omega^{2}})%
}

\NovelDisplay{M01902R01}{%
(0)(1,1,1)(2,1,1)(2,1)(3,2)(4,1)(5,2)=\psi(\Omega_{\omega^{2}+1}\times\psi_{\Omega_{\omega^{2}+1}}(\Omega_{\omega^{2}+1}))%
}

\NovelDisplay{M01903R01}{%
(0)(1,1,1)(2,1,1)(2,1)(3,2)(4,2)=\psi(\Omega_{\omega^{2}+1}^{2})%
}

接下来，是 \allowbreak{}\NovelInline{I05992}{\Omega}\allowbreak{} 下标的猝不及防的增长。

\NovelDisplay{M01904R01}{%
(0)(1,1,1)(2,1,1)(2,1)(3,2)(4,3)=\psi(\Omega_{\omega^{2}+2})%
}

\NovelDisplay{M01905R01}{%
(0)(1,1,1)(2,1,1)(2,1)(3,2,1)=\psi(\Omega_{\omega^{2}+\omega})%
}

\NovelDisplay{M01906R01}{%
(0)(1,1,1)(2,1,1)(2,1)(3,2,1)(4,2)(5,3,1)=\psi(\Omega_{\omega^{2}+\omega\times2})%
}

\NovelDisplay{M01907R01}{%
(0)(1,1,1)(2,1,1)(2,1)(3,2,1)(4,2,1)=\psi(\Omega_{\omega^{2}\times2})%
}

同为三行，坏根未必回到最初的零。霍雨浩将根以前的列留在好部，先请卢本韦核查这段是否真的保持原样，再接新增段。

\NovelDisplay{M01908R01}{%
(0)(1,1,1)(2,1,1)\textcolor[HTML]{B45309}{(2,1)}\textcolor[HTML]{B91C1C}{(3,2,1)}\textcolor[HTML]{7E22CE}{(4,2,1)}%
}

\NovelDisplay{M01909R01}{%
\text{原列数}=6, h=3, r=4,%
}
\NovelDisplay{M01909R02}{%
\text{坏根}=\textcolor[HTML]{B45309}{(2,1)},\text{坏部列数}=2,%
}
\NovelDisplay{M01909R03}{%
(\delta_1,\delta_2,\delta_3)=(2,1,0).%
}

\NovelDisplay{M01910R01}{%
\text{按坏部列序排列的提升阈值：}%
}
\NovelDisplay{M01910R02}{%
(2,2),%
}
\NovelDisplay{M01910R03}{%
\text{第 }j\text{ 列只允许前 }\tau_j\text{ 行应用阶差。}%
}

\NovelDisplay{M01911R01}{%
\text{原第 }4\text{ 列：}(2,1)%
}
\NovelDisplay{M01911R02}{%
\longmapsto(2+2q,1+q), q\ge1.%
}

\NovelDisplay{M01912R01}{%
\text{原第 }5\text{ 列：}(3,2,1)%
}
\NovelDisplay{M01912R02}{%
\longmapsto(3+2q,2+q,1), q\ge1.%
}

霍雨浩先检查根以前的保留段，蔡须坤再沿原坏部核阈值。每一新增段使用同一份原坏部与阶差，编号 q 为正整数；最高非零行及更后坐标不因段数增多而跟着改变。

\NovelDisplay{M01913R01}{%
\bigl((0)(1,1,1)(2,1,1)(2,1)(3,2,1)(4,2,1)\bigr)[0]%
}
\NovelDisplay{M01913R02}{%
\quad =(0)(1,1,1)(2,1,1)\textcolor[HTML]{B45309}{(2,1)}\textcolor[HTML]{B91C1C}{(3,2,1)}%
}

\NovelDisplay{M01914R01}{%
\bigl((0)(1,1,1)(2,1,1)(2,1)(3,2,1)(4,2,1)\bigr)[1]%
}
\NovelDisplay{M01914R02}{%
\quad =(0)(1,1,1)(2,1,1)\textcolor[HTML]{B45309}{(2,1)}\textcolor[HTML]{B91C1C}{(3,2,1)}\textcolor[HTML]{1D4ED8}{(4,2)}\textcolor[HTML]{1D4ED8}{(5,3,1)}%
}

\NovelDisplay{M01915R01}{%
\bigl((0)(1,1,1)(2,1,1)(2,1)(3,2,1)(4,2,1)\bigr)[2]=(0)(1,1,1)(2,1,1)\textcolor[HTML]{B45309}{(2,1)}%
}
\NovelDisplay{M01915R02}{%
\quad \textcolor[HTML]{B91C1C}{(3,2,1)}\textcolor[HTML]{1D4ED8}{(4,2)}\textcolor[HTML]{1D4ED8}{(5,3,1)}\textcolor[HTML]{15803D}{(6,3)}\textcolor[HTML]{15803D}{(7,4,1)}%
}

\NovelDisplay{M01916R01}{%
\text{第零项列数}=5,%
}
\NovelDisplay{M01916R02}{%
\text{第一项列数}=5+2=7,%
}
\NovelDisplay{M01916R03}{%
\text{第二项列数}=5+2\times2=9.%
}

蔡须坤先盖住答案，按列数把第一项与第二项各写一次。霍雨浩核坐标，卢本韦看保留段有没有误留末列，古月娜再查比较方向。两项都从当前原式求得，不能把第一项再次展开的结果填到第二项的位置。

下标中的乘法出现以后，霍雨浩先把自己画在外层的一根线擦掉了。

``我把乘法读到了另一个地方。''

古月娜将那根线重新画在纸边，保留为错法，再请他指出新式的运算从哪一层开始。最外面仍是坍缩，里面是带下标的基数；乘法发生在这个下标的表达式中。下标中的平方与乘二有各自的位置，不能移成整座基数的平方或两倍，也不能与外层参数里原有的乘数互换。

霍雨浩一层层读完，蔡须坤才把纸转到自己面前。她以往常将下标当作一个已经取好的名字，念出它便准备往后走。如今里面还能继续做加法、乘法和幂，她要像对待参数那样，继续把它拆开，而不是用一句``下标变大''遮过去。

``先说它里面发生了什么，再说整条式子。''卢本韦提醒。

蔡须坤照着读了一遍。古月娜问她，根以前那些列为何还在。她回到左端，将好部与坏部重新划开。局部上需要讨论的运算越深，前面的上下文越容易被视为已经无关；恰恰因为下标里有自己的结构，保留段才更不能随意省去。

这一次四人做了一个没有新公式的小练习。霍雨浩只指位置，卢本韦说出那里属于哪层；古月娜遮住右端，蔡须坤指出左端要保留到哪里。随后轮换，谁把某层括号漏掉，就回到那层的起点重新读，不把整组从头讲起。

蔡须坤错了两次，第三次才不再把乘数说到外面去。压力压在肩上，原先那种名字念得很熟、下一步却无从落手的感觉逐渐有了区别。她知道自己还没有走到更高的石桩，但眼前这一层已经能在同伴突然换题时接住。

奥斯卡送水上来，见四人一人指纸、一人念括号，便将水囊轻轻放在旁边。他没打断，等霍雨浩抬起头才问要不要先歇。卢本韦看了看蔡须坤的手腕，替她答应下来。

休息时，纸仍摊着，没人急着将新下标缩成另一个方便的字母。那些长长的符号保留在原处，等四人再次拿笔，还能清楚看见每层运算的边界。

\NovelDisplay{M01917R01}{%
(0)(1,1,1)(2,1,1)(2,1)(3,2,1)(4,2,1)(4,2)(5,3,1)(6,3,1)=\psi(\Omega_{\omega^{2}\times3})%
}

\NovelDisplay{M01918R01}{%
(0)(1,1,1)(2,1,1)(2,1,1)=\psi(\Omega_{\omega^{3}})%
}

\NovelDisplay{M01919R01}{%
(0)(1,1,1)(2,1,1)(2,1,1)(1,1,1)(2,1,1)(2,1,1)=\psi(\Omega_{\omega^{3}}\times2)%
}

可以看作现在是 \allowbreak{}\NovelInline{I05993}{(1,1,1)(2,1,1)(2,1,1)}\allowbreak{} 三项在进行着``添加 \allowbreak{}\NovelInline{I05994}{\Omega_{\omega^3}}\allowbreak{} ''的功能。

\NovelDisplay{M01920R01}{%
(0)(1,1,1)(2,1,1)(2,1,1)(2,1)=\psi(\Omega_{\omega^{3}}\times\Omega)%
}

接下来，该是``三重提升效应''了。

\NovelDisplay{M01921R01}{%
(0)(1,1,1)(2,1,1)(2,1,1)(2,1)(1,1,1)=\psi(\Omega_{\omega^{3}}\times\Omega_{\omega})%
}

\NovelDisplay{M01922R01}{%
(0)(1,1,1)(2,1,1)(2,1,1)(2,1)(1,1,1)(2,1,1)=\psi(\Omega_{\omega^{3}}\times\Omega_{\omega^{2}})%
}

\NovelDisplay{M01923R01}{%
(0)(1,1,1)(2,1,1)(2,1,1)(2,1)(1,1,1)(2,1,1)(2,1,1)=\psi(\Omega_{\omega^{3}}^{2})%
}

这个过程也可以看作有一条 \allowbreak{}\NovelInline{I05995}{(1,1,1)(2,1,1)(2,1,1)}\allowbreak{} 的提升链，可以直接把 \allowbreak{}\NovelInline{I05996}{\Omega}\allowbreak{} 提升到 \allowbreak{}\NovelInline{I05997}{\Omega_{\omega^3}}\allowbreak{} 。

\NovelDisplay{M01924R01}{%
(0)(1,1,1)(2,1,1)(2,1,1)(2,1)(3,2)=\psi(\Omega_{\omega^{3}+1})%
}

\NovelDisplay{M01925R01}{%
(0)(1,1,1)(2,1,1)(2,1,1)(2,1)(3,2,1)(4,2,1)(4,2,1)=\psi(\Omega_{\omega^{3}\times2})%
}

\NovelDisplay{M01926R01}{%
(0)(1,1,1)(2,1,1)(2,1,1)(2,1,1)=\psi(\Omega_{\omega^{4}})%
}

一个 \allowbreak{}\NovelInline{I05998}{(2,1,1)}\allowbreak{} ，下标乘一个 \allowbreak{}\NovelInline{I05999}{\omega}\allowbreak{} 。

天阴下来，卢本韦在平石边点了一盏灯。

古月娜说光够亮，手却仍停在纸上。她刚才把有限重复的次数当成了下标的极限，最后一行迟迟接不稳。蔡须坤没有催，挪到她身边，将前几个有限项一条条写出来。

``这些都没有等于它。''

古月娜点头，将最后那个等号改成比较号。霍雨浩接着写任意有限次数的生成规则，卢本韦让两人轮流说清楚下一项，而不是只念已经抄好的几项。

灯烧去半截，古月娜终于放下笔，独自走完这一轮运转。她额角的汗冷了，抬手却将那两张有错的纸一并收进本子。

\NovelDisplay{M01927R01}{%
(0)(1,1,1)(2,1,1)(3)=\psi(\Omega_{\omega^{\omega}})%
}

灯下那张纸，古月娜一直留到了第二天。

``昨天我把有限次当成已经到达极限。''她说。

卢本韦没有翻那张错页，只让她指出当前末列的最高非零行。这里末列只在第一行非零，根停在紧靠前面的三行列；可提升的更前行已经没有，所以阶差全零。展开出的新增段并非每个坐标都要升高，而是依规则重复原坏部，有限项以这样的方式逐项生成。

蔡须坤先数第一项，又让霍雨浩不看她的纸生成第二项。两份答案接上以后，古月娜才解释为什么写出再多有限项也不能把最后一项当作极限。这里不存在一个``最大的有限次数''等着他们数到。给定任意一个有限位置，还可以按规则写出后面的项；当前更简洁的节点组织了整个基本序列，不等同于其中某个最终成员。

``我们昨天只列前三项，凭什么就说尽了？''卢本韦问。

``前三项帮助检查规则。''古月娜答，``任意有限项怎样生成，得由原来的规则说明。右边的对应也要放回完整结构，不能拿三个例子代替所有项的理由。''

霍雨浩将这句写到有限项旁边。四人已经有了可用的展开算法，这里要核的是算法在本题怎样选根、保留什么、重复什么；不另用几幅相像的图充当一份通用证明。对于当前给出的序数对应，他们同时保留原来的分析链，知道每个简写来自哪里。

蔡须坤想到脚下的山路。他们确实能一步步走完有限距离，眼前的序数极限却不能直接拿脚步数量来比。山的米数由考核设定，代表需要分析清楚的关口；它不是把序数当成普通长度，再按比例换算。否则短短几列与长长几页交替出现，早就会让地图失去意义。

古月娜听到这里，终于将昨天错用的等号划去。纸上留下一点淡墨，她没有反复涂黑，只在旁边写清为什么要改。

卢本韦把热水往她那边推。``现在能睡得着了？''

``还有下一层。''

``下一层也等吃完再说。''

蔡须坤先将本子合上，替古月娜留住那张尚未干透的纸。四人歇过以后再上台阶，没有抢着跨过这一处。压力仍沉，有限项与极限之间那条原先混着的界线，却已经能在他们的说明中清楚留下。

\NovelDisplay{M01928R01}{%
(0)(1,1,1)(2,1,1)(3)(1,1,1)(2,1,1)(3)=\psi(\Omega_{\omega^{\omega}}\times2)%
}

\NovelDisplay{M01929R01}{%
(0)(1,1,1)(2,1,1)(3)(2,1)(1,1,1)(2,1,1)(3)=\psi(\Omega_{\omega^{\omega}}^{2})%
}

\NovelDisplay{M01930R01}{%
(0)(1,1,1)(2,1,1)(3)(2,1)(3,2)=\psi(\Omega_{\omega^{\omega}+1})%
}

\NovelDisplay{M01931R01}{%
(0)(1,1,1)(2,1,1)(3)(2,1)(3,2)(4,2)=\psi(\Omega_{\omega^{\omega}+1}^{2})%
}

\NovelDisplay{M01932R01}{%
(0)(1,1,1)(2,1,1)(3)(2,1)(3,2,1)=\psi(\Omega_{\omega^{\omega}+\omega})%
}

\NovelDisplay{M01933R01}{%
(0)(1,1,1)(2,1,1)(3)(2,1)(3,2,1)(4,2,1)=\psi(\Omega_{\omega^{\omega}+\omega^{2}})%
}

\NovelDisplay{M01934R01}{%
(0)(1,1,1)(2,1,1)(3)(2,1)(3,2,1)(4,2,1)(5)=\psi(\Omega_{\omega^{\omega}\times2})%
}

\NovelDisplay{M01935R01}{%
(0)(1,1,1)(2,1,1)(3)(2,1,1)=\psi(\Omega_{\omega^{\omega+1}})%
}

有限重复与极限之间，蔡须坤留下这条作例题。古月娜先查尾部单独那列的祖先，霍雨浩按允许提升的行数生成新列，卢本韦则专门寻找不能同步增加的坐标。

\NovelDisplay{M01936R01}{%
\textcolor[HTML]{B45309}{(0)}\textcolor[HTML]{B91C1C}{(1,1,1)}\textcolor[HTML]{B91C1C}{(2,1,1)}\textcolor[HTML]{B91C1C}{(3)}\textcolor[HTML]{7E22CE}{(2,1,1)}%
}

\NovelDisplay{M01937R01}{%
\text{原列数}=5, h=3, r=1,%
}
\NovelDisplay{M01937R02}{%
\text{坏根}=\textcolor[HTML]{B45309}{(0)},\text{坏部列数}=4,%
}
\NovelDisplay{M01937R03}{%
(\delta_1,\delta_2,\delta_3)=(2,1,0).%
}

\NovelDisplay{M01938R01}{%
\text{按坏部列序排列的提升阈值：}%
}
\NovelDisplay{M01938R02}{%
(2,2,2,1),%
}
\NovelDisplay{M01938R03}{%
\text{第 }j\text{ 列只允许前 }\tau_j\text{ 行应用阶差。}%
}

\NovelDisplay{M01939R01}{%
\text{原第 }1\text{ 列：}(0)%
}
\NovelDisplay{M01939R02}{%
\longmapsto(2q,q), q\ge1.%
}

\NovelDisplay{M01940R01}{%
\text{原第 }4\text{ 列：}(3)%
}
\NovelDisplay{M01940R02}{%
\longmapsto(3+2q), q\ge1.%
}

古月娜把阈值较低的那几列单独指出。阶差虽已求出，原坏部各列允许改变的行数仍不同；下方的灰色坐标保留原值，其余有资格的坐标按新增段的正整数编号 q 应用相应倍数的阶差。

\NovelDisplay{M01941R01}{%
\bigl((0)(1,1,1)(2,1,1)(3)(2,1,1)\bigr)[0]=\textcolor[HTML]{B45309}{(0)}\textcolor[HTML]{B91C1C}{(1,1,1)}\textcolor[HTML]{B91C1C}{(2,1,1)}\textcolor[HTML]{B91C1C}{(3)}%
}

\NovelDisplay{M01942R01}{%
\bigl((0)(1,1,1)(2,1,1)(3)(2,1,1)\bigr)[1]=\textcolor[HTML]{B45309}{(0)}%
}
\NovelDisplay{M01942R02}{%
\quad \textcolor[HTML]{B91C1C}{(1,1,1)}\textcolor[HTML]{B91C1C}{(2,1,1)}\textcolor[HTML]{B91C1C}{(3)}\textcolor[HTML]{1D4ED8}{(2,1)}\textcolor[HTML]{1D4ED8}{(3,2,1)}\textcolor[HTML]{1D4ED8}{(4,2,1)}\textcolor[HTML]{1D4ED8}{(5)}%
}

\NovelDisplay{M01943R01}{%
\bigl((0)(1,1,1)(2,1,1)(3)(2,1,1)\bigr)[2]=\textcolor[HTML]{B45309}{(0)}\textcolor[HTML]{B91C1C}{(1,1,1)}\textcolor[HTML]{B91C1C}{(2,1,1)}\textcolor[HTML]{B91C1C}{(3)}\textcolor[HTML]{1D4ED8}{(2,1)}\textcolor[HTML]{1D4ED8}{(3,2,1)}%
}
\NovelDisplay{M01943R02}{%
\quad \textcolor[HTML]{1D4ED8}{(4,2,1)}\textcolor[HTML]{1D4ED8}{(5)}\textcolor[HTML]{15803D}{(4,2)}\textcolor[HTML]{15803D}{(5,3,1)}\textcolor[HTML]{15803D}{(6,3,1)}\textcolor[HTML]{15803D}{(7)}%
}

\NovelDisplay{M01944R01}{%
\text{第零项列数}=4,%
}
\NovelDisplay{M01944R02}{%
\text{第一项列数}=4+4=8,%
}
\NovelDisplay{M01944R03}{%
\text{第二项列数}=4+2\times4=12.%
}

古月娜将第二项的两段新增部分分别圈出，卢本韦从原坏部核它们各自的增量。蔡须坤专查那些本该不动的位置，霍雨浩最后逐列比较。各项虽然可以更长，仍位于原式之下。

\NovelDisplay{M01945R01}{%
(0)(1,1,1)(2,1,1)(3)(2,1,1)(3)=\psi(\Omega_{\omega^{\omega\times2}})%
}

\NovelDisplay{M01946R01}{%
(0)(1,1,1)(2,1,1)(3)(3)=\psi(\Omega_{\omega^{\omega^{2}}})%
}

\NovelDisplay{M01947R01}{%
(0)(1,1,1)(2,1,1)(3)(4)=\psi(\Omega_{\omega^{\omega^{\omega}}})%
}

\NovelDisplay{M01948R01}{%
(0)(1,1,1)(2,1,1)(3)(4,1)=\psi(\Omega_{\psi(\Omega)})%
}

下标第一次出现了 \allowbreak{}\NovelInline{I06000}{\psi}\allowbreak{} 。

\NovelDisplay{M01949R01}{%
(0)(1,1,1)(2,1,1)(3)(4,1)(2,1)(3,2)=\psi(\Omega_{\psi(\Omega)+1})%
}

\NovelDisplay{M01950R01}{%
(0)(1,1,1)(2,1,1)(3)(4,1)(5,1)=\psi(\Omega_{\psi(\Omega^{2})})%
}

\NovelDisplay{M01951R01}{%
(0)(1,1,1)(2,1,1)(3)(4,1,1)=\psi(\Omega_{\psi(\Omega_{\omega})})%
}

\NovelDisplay{M01952R01}{%
(0)(1,1,1)(2,1,1)(3)(4,1,1)(5,1,1)=\psi(\Omega_{\psi(\Omega_{\omega^{2}})})%
}

在下标的 \allowbreak{}\NovelInline{I06001}{\psi}\allowbreak{} 里面重复着一样的操作。

\NovelDisplay{M01953R01}{%
(0)(1,1,1)(2,1,1)(3)(4,1,1)(5,1,1)(6)(7,1)=\psi(\Omega_{\psi(\Omega_{\psi(\Omega)})})%
}

这么做的极限自然是------

\NovelDisplay{M01954R01}{%
(0)(1,1,1)(2,1,1)(3,1)=\psi(\Omega_{\Omega})%
}

下标里第一次出现完整的 Ω 时，四人没有跳过它的展开。蔡须坤在根的位置落笔，霍雨浩写阶差，卢本韦读出完整的第零项，古月娜接着核第一项和第二项。

\NovelDisplay{M01955R01}{%
\textcolor[HTML]{B45309}{(0)}\textcolor[HTML]{B91C1C}{(1,1,1)}\textcolor[HTML]{B91C1C}{(2,1,1)}\textcolor[HTML]{7E22CE}{(3,1)}%
}

\NovelDisplay{M01956R01}{%
\text{原列数}=4, h=2, r=1,%
}
\NovelDisplay{M01956R02}{%
\text{坏根}=\textcolor[HTML]{B45309}{(0)},\text{坏部列数}=3,%
}
\NovelDisplay{M01956R03}{%
(\delta_1,\delta_2,\delta_3)=(3,0,0).%
}

\NovelDisplay{M01957R01}{%
\text{按坏部列序排列的提升阈值：}%
}
\NovelDisplay{M01957R02}{%
(1,1,1),%
}
\NovelDisplay{M01957R03}{%
\text{第 }j\text{ 列只允许前 }\tau_j\text{ 行应用阶差。}%
}

\NovelDisplay{M01958R01}{%
\text{原第 }1\text{ 列：}(0)%
}
\NovelDisplay{M01958R02}{%
\longmapsto(3q), q\ge1.%
}

\NovelDisplay{M01959R01}{%
\text{原第 }3\text{ 列：}(2,1,1)%
}
\NovelDisplay{M01959R02}{%
\longmapsto(2+3q,1,1), q\ge1.%
}

霍雨浩先检查根以前的保留段，蔡须坤再沿原坏部核阈值。每一新增段使用同一份原坏部与阶差，编号 q 为正整数；最高非零行及更后坐标不因段数增多而跟着改变。

\NovelDisplay{M01960R01}{%
\bigl((0)(1,1,1)(2,1,1)(3,1)\bigr)[0]=\textcolor[HTML]{B45309}{(0)}\textcolor[HTML]{B91C1C}{(1,1,1)}\textcolor[HTML]{B91C1C}{(2,1,1)}%
}

\NovelDisplay{M01961R01}{%
\bigl((0)(1,1,1)(2,1,1)(3,1)\bigr)[1]=\textcolor[HTML]{B45309}{(0)}\textcolor[HTML]{B91C1C}{(1,1,1)}\textcolor[HTML]{B91C1C}{(2,1,1)}\textcolor[HTML]{1D4ED8}{(3)}\textcolor[HTML]{1D4ED8}{(4,1,1)}\textcolor[HTML]{1D4ED8}{(5,1,1)}%
}

\NovelDisplay{M01962R01}{%
\bigl((0)(1,1,1)(2,1,1)(3,1)\bigr)[2]=\textcolor[HTML]{B45309}{(0)}\textcolor[HTML]{B91C1C}{(1,1,1)}%
}
\NovelDisplay{M01962R02}{%
\quad \textcolor[HTML]{B91C1C}{(2,1,1)}\textcolor[HTML]{1D4ED8}{(3)}\textcolor[HTML]{1D4ED8}{(4,1,1)}\textcolor[HTML]{1D4ED8}{(5,1,1)}\textcolor[HTML]{15803D}{(6)}\textcolor[HTML]{15803D}{(7,1,1)}\textcolor[HTML]{15803D}{(8,1,1)}%
}

\NovelDisplay{M01963R01}{%
\text{第零项列数}=3,%
}
\NovelDisplay{M01963R02}{%
\text{第一项列数}=3+3=6,%
}
\NovelDisplay{M01963R03}{%
\text{第二项列数}=3+2\times3=9.%
}

卢本韦数完列，将第一项中的蓝段与第二项中的蓝段逐列重合。绿色段另由原坏部按两倍阶差生成，原来的保留段没有再受一次提升。蔡须坤核坐标，古月娜与霍雨浩最后检查各有限项都小于原式。

完整的 Ω 放进下标以后，蔡须坤先念了一遍，又回到最外面重读。

``我刚才把两个 Ω 当成同一个位置了。''

霍雨浩将纸向她靠近一点。外面这个基数由下标指定，里面的 Ω 则是下标表达式的一部分。字形相同，并不表示可以消去、相乘或者把它们沿纸面移到一起。此前那些加法、乘法与幂次已经教会他们区分层次，如今读得更深，边界仍然需要保留。

古月娜重新展开这条 BMS。末列的最高非零行回到第二行，根在最初的零列，三列坏部都留在原来的次序中。她先用第一项解释怎样接新段，再让卢本韦从同一个原式求第二项。两份展开都比原式更长，比较方向仍向原式之下。

蔡须坤把对应的层级放回这组有限结构旁边，没有把下标里的 Ω 解释成某个有限长式已经与它相等。前面各项给出向上理解这一关口的路径，极限节点则有自己的位置。说``规则与此前相同''，是说仍按父项、坏根和阶差求 BMS 的项；不等于任何内层运算都可凭表面相似任意替换。

夏夜星空这时请他们交换笔记。霍雨浩拿到卢本韦那页，先把自己熟悉的结论放到一边，只沿她写出的连接检查。蔡须坤读古月娜那页，问她中间的一段为何要保留；古月娜又拿蔡须坤的记录，追问她在哪里将同名符号分开。

没有谁从头讲到尾，问题只落在刚才这一层。

四人讲完，纸上的长式依旧很长，卢本韦却将原先为它预留的两个代号划去了。直接写出下标里的结构，虽然费纸，回头核查时便不会忘记自己代的是整项、参数还是其中一个基数。霍雨浩把桌边三张窄纸接成两张较宽的，仍将公式沿同一层括号分行，不让它挤成一条细线。

蔡须坤走出横台，回头看见古月娜将刚才的问题写到下一页。下标中的 Ω 还能不能接受后面的提升？他们此刻只把问题留住，还不凭它已经出现在里面，就替后面所有结构先下结论。

雾从山腰升上来，更高处的石桩只剩一道模糊的影子。卢本韦点好灯，把新页推到四人中央。他们从当前这一层继续，不再为前面的每个等号另作总结。

下标上被放上了一个 \allowbreak{}\NovelInline{I06002}{\Omega}\allowbreak{} 。可以看到，展开下标上的 \allowbreak{}\NovelInline{I06003}{\Omega}\allowbreak{} 的规则和之前是一样的，只需要把``下标''也当做运算的一种就可以了。不过，既然下标上出现了 \allowbreak{}\NovelInline{I06004}{\Omega}\allowbreak{} ，那么它可不可以被提升呢？

蔡须坤先提出猜想：下标里的 Ω 仍可能受到后续提升。霍雨浩保留猜想，再逐条写出可检查的中间结构。

霍雨浩将这组长式分成几页。蔡须坤先读外层，卢本韦追复制段，古月娜专查下标中的变化。每翻一页，四人都交换一次负责的位置。

接下来的几日，蔡须坤一直把那个下标中的 Ω 留在页首。

每做完一组，她便将已经核过的前缀与新页并排。先前提出的问题还在，后面的具体结构却要重新追根，不能将上一页的名字直接搬来。霍雨浩记下受阻的位置，卢本韦翻出对应的旧式，古月娜专查连接处。

``前面这段留住，''蔡须坤说，``从下一处真正换了位置的地方接。''

霍雨浩把新页翻开。卢本韦接过她手里的旧本子，替她压好页角；古月娜则在新页中央留出一条足够长的空白。

古月娜将笔尖落到末列：``先把这一题的根找清。''

\NovelDisplay{M01964R01}{%
(0)(1,1,1)(2,1,1)(3,1)(1,1)=\psi(\Omega_{\Omega}+\Omega)%
}

\NovelDisplay{M01965R01}{%
(0)(1,1,1)(2,1,1)(3,1)(1,1)(2,2,1)=\psi(\Omega_{\Omega}+\psi_{\Omega_2}(\Omega_{\omega}))%
}

\NovelDisplay{M01966R01}{%
(0)(1,1,1)(2,1,1)(3,1)(1,1)(2,2,1)(3,2,1)=\psi(\Omega_{\Omega}+\psi_{\Omega_2}(\Omega_{\omega^{2}}))%
}

\NovelDisplay{M01967R01}{%
(0)(1,1,1)(2,1,1)(3,1)(1,1)(2,2,1)(3,2,1)(4,1)=\psi(\Omega_{\Omega}+\psi_{\Omega_2}(\Omega_{\Omega}))%
}

\NovelDisplay{M01968R01}{%
(0)(1,1,1)(2,1,1)(3,1)(1,1)(2,2,1)(3,2,1)(4,1)(2,2)=\psi(\Omega_{\Omega}+\Omega_2)%
}

\NovelDisplay{M01969R01}{%
(0)(1,1,1)(2,1,1)(3,1)(1,1)(2,2,1)(3,2,1)(4,1)(2,2)(3,1)(4,2,1)(5,2,1)%
}
\NovelDisplay{M01969R02}{%
\quad (6,1)=\psi(\Omega_{\Omega}+\Omega_2\times\psi_{\Omega_2}(\Omega_{\Omega}))%
}

\NovelDisplay{M01970R01}{%
(0)(1,1,1)(2,1,1)(3,1)(1,1)(2,2,1)(3,2,1)(4,1)(2,2)(3,2)=\psi(\Omega_{\Omega}+\Omega_2^{2})%
}

\NovelDisplay{M01971R01}{%
(0)(1,1,1)(2,1,1)(3,1)(1,1)(2,2,1)(3,2,1)(4,1)%
}
\NovelDisplay{M01971R02}{%
\quad (2,2)(3,3,1)(4,3,1)(5,1)=\psi(\Omega_{\Omega}+\psi_{\Omega_3}(\Omega_{\Omega}))%
}

\NovelDisplay{M01972R01}{%
(0)(1,1,1)(2,1,1)(3,1)(1,1)(2,2,1)(3,2,1)(4,1)(2,2)(3,3,1)(4,3,1)(5,1)%
}
\NovelDisplay{M01972R02}{%
\quad (3,3)(4,3)=\psi(\Omega_{\Omega}+\Omega_3^{2})%
}

\NovelDisplay{M01973R01}{%
(0)(1,1,1)(2,1,1)(3,1)(1,1)(2,2,1)(3,2,1)(4,1)(2,2,1)=\psi(\Omega_{\Omega}+\Omega_{\omega})%
}

\NovelDisplay{M01974R01}{%
(0)(1,1,1)(2,1,1)(3,1)(1,1)(2,2,1)(3,2,1)%
}
\NovelDisplay{M01974R02}{%
\quad (4,1)(2,2,1)(2,2,1)=\psi(\Omega_{\Omega}+\Omega_{\omega}\times2)%
}

\NovelDisplay{M01975R01}{%
(0)(1,1,1)(2,1,1)(3,1)(1,1)(2,2,1)(3,2,1)%
}
\NovelDisplay{M01975R02}{%
\quad (4,1)(2,2,1)(3,2)(2,2,1)=\psi(\Omega_{\Omega}+\Omega_{\omega}^{2})%
}

\NovelDisplay{M01976R01}{%
(0)(1,1,1)(2,1,1)(3,1)(1,1)(2,2,1)(3,2,1)(4,1)(2,2,1)(3,2)(4,3)%
}
\NovelDisplay{M01976R02}{%
\quad =\psi(\Omega_{\Omega}+\psi_{\Omega_{\omega+1}}(\Omega_{\omega+1}))%
}

\NovelDisplay{M01977R01}{%
(0)(1,1,1)(2,1,1)(3,1)(1,1)(2,2,1)(3,2,1)(4,1)(2,2,1)(3,2)(4,3,1)%
}
\NovelDisplay{M01977R02}{%
\quad (5,3,1)(6,1)=\psi(\Omega_{\Omega}+\psi_{\Omega_{\omega+1}}(\Omega_{\Omega}))%
}

\NovelDisplay{M01978R01}{%
(0)(1,1,1)(2,1,1)(3,1)(1,1)(2,2,1)(3,2,1)(4,1)(2,2,1)(3,2)(4,3,1)%
}
\NovelDisplay{M01978R02}{%
\quad (5,3,1)(6,1)(4,3)(5,3)=\psi(\Omega_{\Omega}+\Omega_{\omega+1}^{2})%
}

\NovelDisplay{M01979R01}{%
(0)(1,1,1)(2,1,1)(3,1)(1,1)(2,2,1)(3,2,1)(4,1)(2,2,1)(3,2)(4,3,1)%
}
\NovelDisplay{M01979R02}{%
\quad (5,3,1)(6,1)(4,3,1)=\psi(\Omega_{\Omega}+\Omega_{\omega\times2})%
}

\NovelDisplay{M01980R01}{%
(0)(1,1,1)(2,1,1)(3,1)(1,1)(2,2,1)(3,2,1)%
}
\NovelDisplay{M01980R02}{%
\quad (4,1)(2,2,1)(3,2,1)=\psi(\Omega_{\Omega}+\Omega_{\omega^{2}})%
}

\NovelDisplay{M01981R01}{%
(0)(1,1,1)(2,1,1)(3,1)(1,1)(2,2,1)(3,2,1)(4,1)(2,2,1)(3,2,1)(2,2,1)%
}
\NovelDisplay{M01981R02}{%
\quad (3,2,1)=\psi(\Omega_{\Omega}+\Omega_{\omega^{2}}\times2)%
}

\NovelDisplay{M01982R01}{%
(0)(1,1,1)(2,1,1)(3,1)(1,1)(2,2,1)(3,2,1)(4,1)(2,2,1)(3,2,1)(3,2)(4,3)%
}
\NovelDisplay{M01982R02}{%
\quad =\psi(\Omega_{\Omega}+\psi_{\Omega_{\omega^{2}+1}}(\Omega_{\omega^{2}+1}))%
}

\NovelDisplay{M01983R01}{%
(0)(1,1,1)(2,1,1)(3,1)(1,1)(2,2,1)(3,2,1)(4,1)(2,2,1)(3,2,1)(3,2)%
}
\NovelDisplay{M01983R02}{%
\quad (4,3,1)(5,3,1)(6,1)(4,3)(5,3)=\psi(\Omega_{\Omega}+\Omega_{\omega^{2}+1}^{2})%
}

\NovelDisplay{M01984R01}{%
(0)(1,1,1)(2,1,1)(3,1)(1,1)(2,2,1)(3,2,1)(4,1)(2,2,1)(3,2,1)(3,2)%
}
\NovelDisplay{M01984R02}{%
\quad (4,3,1)(5,3,1)(6,1)(4,3,1)=\psi(\Omega_{\Omega}+\Omega_{\omega^{2}+\omega})%
}

\NovelDisplay{M01985R01}{%
(0)(1,1,1)(2,1,1)(3,1)(1,1)(2,2,1)(3,2,1)(4,1)(2,2,1)(3,2,1)(3,2)%
}
\NovelDisplay{M01985R02}{%
\quad (4,3,1)(5,3,1)(6,1)(4,3,1)(5,3,1)=\psi(\Omega_{\Omega}+\Omega_{\omega^{2}\times2})%
}

\NovelDisplay{M01986R01}{%
(0)(1,1,1)(2,1,1)(3,1)(1,1)(2,2,1)(3,2,1)%
}
\NovelDisplay{M01986R02}{%
\quad (4,1)(2,2,1)(3,2,1)(3,2,1)=\psi(\Omega_{\Omega}+\Omega_{\omega^{3}})%
}

\NovelDisplay{M01987R01}{%
(0)(1,1,1)(2,1,1)(3,1)(1,1)(2,2,1)(3,2,1)%
}
\NovelDisplay{M01987R02}{%
\quad (4,1)(2,2,1)(3,2,1)(4)=\psi(\Omega_{\Omega}+\Omega_{\omega^{\omega}})%
}

\NovelDisplay{M01988R01}{%
(0)(1,1,1)(2,1,1)(3,1)(1,1)(2,2,1)(3,2,1)(4,1)%
}
\NovelDisplay{M01988R02}{%
\quad (2,2,1)(3,2,1)(4)(5,1,1)=\psi(\Omega_{\Omega}+\Omega_{\psi(\Omega_{\omega})})%
}

\NovelDisplay{M01989R01}{%
(0)(1,1,1)(2,1,1)(3,1)(1,1)(2,2,1)(3,2,1)(4,1)(2,2,1)(3,2,1)(4)(5,1,1)%
}
\NovelDisplay{M01989R02}{%
\quad (6,1,1)(7,1)(5,1)(6,2,1)(7,2,1)(8,1)%
}
\NovelDisplay{M01989R03}{%
\quad (6,2)(7,2)=\psi(\Omega_{\Omega}+\Omega_{\psi(\Omega_{\Omega}+\Omega_2^{2})})%
}

\NovelDisplay{M01990R01}{%
(0)(1,1,1)(2,1,1)(3,1)(1,1)(2,2,1)(3,2,1)%
}
\NovelDisplay{M01990R02}{%
\quad (4,1)(2,2,1)(3,2,1)(4,1)=\psi(\Omega_{\Omega}\times2)%
}

\NovelDisplay{M01991R01}{%
(0)(1,1,1)(2,1,1)(3,1)(1,1)(2,2,1)(3,2,1)(4,1)(2,2,1)(3,2,1)(4,1)%
}
\NovelDisplay{M01991R02}{%
\quad (2,2,1)(3,2,1)(4,1)=\psi(\Omega_{\Omega}\times3)%
}

\NovelDisplay{M01992R01}{%
(0)(1,1,1)(2,1,1)(3,1)(1,1)(2,2,1)(3,2,1)(4,1)(3,1)=\psi(\Omega_{\Omega}\times\Omega)%
}

\NovelDisplay{M01993R01}{%
(0)(1,1,1)(2,1,1)(3,1)(1,1)(2,2,1)(3,2,1)(4,1)%
}
\NovelDisplay{M01993R02}{%
\quad (3,1)(4,2,1)(5,2,1)(6,1)=\psi(\Omega_{\Omega}\times\psi_{\Omega_2}(\Omega_{\Omega}))%
}

\NovelDisplay{M01994R01}{%
(0)(1,1,1)(2,1,1)(3,1)(1,1)(2,2,1)(3,2,1)(4,1)(3,2)=\psi(\Omega_{\Omega}\times\Omega_2)%
}

夜深，石屋里只剩灯芯轻轻炸开的一声。

霍雨浩将下标里的加号重新数完，古月娜又从末列追过一次祖先。蔡须坤把两人的记录摆在一起，确认差别只在刚才讨论的位置。卢本韦坐在门边，见三人终于放下笔，才把留着的热水递过去。

``明天那页，还得从根开始。''

蔡须坤接住杯子，掌心被烫得一暖。

``根找得清，长一点也能算。''

她说完又看了眼山图。二百二十米的石桩还在前面，五百米的横线更远，最高处仍是二千五百。她将山图卷好，放在十一人的行李之间，等第二天再一同带上去。

\chapter{第一百一十一章：链的尽头，二百二十米}\label{ux7b2cux4e00ux767eux4e00ux5341ux4e00ux7ae0ux94feux7684ux5c3dux5934ux4e8cux767eux4e8cux5341ux7c73}

蔡须坤第一次看清二百二十米的石桩，是在一个没有风的早晨。

它并不比前面那些石桩高，只是侧面刻着的列短得出奇。霍雨浩隔着几级台阶认出那几列，随即把目光收回来，落到眼前还没有分析完的长式上。

卢本韦笑了一下。

``到了跟前，倒不敢先念答案了？''

``我怕先念完，自己就当已经走到了。''

古月娜把纸向他那边推了一点。``那就先算现在这一条。''

他们从前一天留下的乘数接起。此时的压力不再只让手腕与膝盖发酸；一处下标没有安放清楚，整段运转便会在更内层折回来。蔡须坤将呼吸分得很缓，等那处回牵显出来，才让霍雨浩记下先后。

古月娜仍会停下来。卢本韦仍会写错。四人已不再急着把错误一笔涂黑：错列旁留下父项，错括号旁留下原先的读法，下一轮各自独立再算。到了这时，那些由压力显出的细节，才一点点成为能离开台阶、在白纸上也说清楚的规则。

唐三和小舞把后方的绳索重新系稳，叶阳将被雾气浸湿的纸晾在平石上，薛峰替大家数好页码。蔡兴坤蹲在石桩边，刚想让黑子鸡往上跳，被诺当用掌心轻轻拦住。

``这里让他们自己走。''

蔡须坤听见了，把最后一截炭条放到纸中央。

``先把整条链走清楚。''

\NovelDisplay{M01995R01}{%
(0)(1,1,1)(2,1,1)(3,1)(1,1)(2,2,1)(3,2,1)%
}
\NovelDisplay{M01995R02}{%
\quad (4,1)(3,2)(2,2,1)=\psi(\Omega_{\Omega}\times\Omega_{\omega})%
}

这一次，绕过坏根的坐标就在旧坏部中间。卢本韦说出所在的列号后，蔡须坤仍从父项链重走一遍，直到灰色标记能与祖先条件对应。

\NovelDisplay{M01996R01}{%
(0)(1,1,1)(2,1,1)(3,1)\textcolor[HTML]{B45309}{(1,1)}\textcolor[HTML]{B91C1C}{(2,2,1)}\textcolor[HTML]{B91C1C}{(3,2,1)}\textcolor[HTML]{B91C1C}{(4,\textcolor[HTML]{6B7280}{1})}\textcolor[HTML]{B91C1C}{(3,2)}\textcolor[HTML]{7E22CE}{(2,2,1)}%
}

\NovelDisplay{M01997R01}{%
\text{原列数}=10, h=3, r=5,%
}
\NovelDisplay{M01997R02}{%
\text{坏根}=\textcolor[HTML]{B45309}{(1,1)},\text{坏部列数}=5,%
}
\NovelDisplay{M01997R03}{%
(\delta_1,\delta_2,\delta_3)=(1,1,0).%
}

\NovelDisplay{M01998R01}{%
\text{按坏部列序排列的提升阈值：}%
}
\NovelDisplay{M01998R02}{%
(2,2,2,1,2),%
}
\NovelDisplay{M01998R03}{%
\text{第 }j\text{ 列只允许前 }\tau_j\text{ 行应用阶差。}%
}

\NovelDisplay{M01999R01}{%
\text{原第 }5\text{ 列：}(1,1)%
}
\NovelDisplay{M01999R02}{%
\longmapsto(1+q,1+q), q\ge1.%
}

\NovelDisplay{M02000R01}{%
\text{原第 }8\text{ 列：}(4,1)%
}
\NovelDisplay{M02000R02}{%
\longmapsto(4+q,1), q\ge1.%
}

\NovelDisplay{M02001R01}{%
\text{原第 }9\text{ 列：}(3,2)%
}
\NovelDisplay{M02001R02}{%
\longmapsto(3+q,2+q), q\ge1.%
}

古月娜把阈值较低的那几列单独指出。阶差虽已求出，原坏部各列允许改变的行数仍不同；下方的灰色坐标保留原值，其余有资格的坐标按新增段的正整数编号 q 应用相应倍数的阶差。

\NovelDisplay{M02002R01}{%
\bigl((0)(1,1,1)(2,1,1)(3,1)(1,1)(2,2,1)(3,2,1)(4,1)(3,2)(2,2,1)\bigr)[0]=(0)%
}
\NovelDisplay{M02002R02}{%
\quad (1,1,1)(2,1,1)(3,1)\textcolor[HTML]{B45309}{(1,1)}\textcolor[HTML]{B91C1C}{(2,2,1)}\textcolor[HTML]{B91C1C}{(3,2,1)}\textcolor[HTML]{B91C1C}{(4,\textcolor[HTML]{6B7280}{1})}\textcolor[HTML]{B91C1C}{(3,2)}%
}

\NovelDisplay{M02003R01}{%
\bigl((0)(1,1,1)(2,1,1)(3,1)(1,1)(2,2,1)(3,2,1)(4,1)(3,2)(2,2,1)\bigr)[1]=(0)%
}
\NovelDisplay{M02003R02}{%
\quad (1,1,1)(2,1,1)(3,1)\textcolor[HTML]{B45309}{(1,1)}\textcolor[HTML]{B91C1C}{(2,2,1)}\textcolor[HTML]{B91C1C}{(3,2,1)}%
}
\NovelDisplay{M02003R03}{%
\quad \textcolor[HTML]{B91C1C}{(4,\textcolor[HTML]{6B7280}{1})}\textcolor[HTML]{B91C1C}{(3,2)}\textcolor[HTML]{1D4ED8}{(2,2)}\textcolor[HTML]{1D4ED8}{(3,3,1)}\textcolor[HTML]{1D4ED8}{(4,3,1)}\textcolor[HTML]{1D4ED8}{(5,1)}\textcolor[HTML]{1D4ED8}{(4,3)}%
}

\NovelDisplay{M02004R01}{%
\bigl((0)(1,1,1)(2,1,1)(3,1)(1,1)(2,2,1)(3,2,1)(4,1)(3,2)(2,2,1)\bigr)[2]=(0)%
}
\NovelDisplay{M02004R02}{%
\quad (1,1,1)(2,1,1)(3,1)\textcolor[HTML]{B45309}{(1,1)}\textcolor[HTML]{B91C1C}{(2,2,1)}\textcolor[HTML]{B91C1C}{(3,2,1)}\textcolor[HTML]{B91C1C}{(4,\textcolor[HTML]{6B7280}{1})}\textcolor[HTML]{B91C1C}{(3,2)}\textcolor[HTML]{1D4ED8}{(2,2)}\textcolor[HTML]{1D4ED8}{(3,3,1)}\textcolor[HTML]{1D4ED8}{(4,3,1)}%
}
\NovelDisplay{M02004R03}{%
\quad \textcolor[HTML]{1D4ED8}{(5,1)}\textcolor[HTML]{1D4ED8}{(4,3)}\textcolor[HTML]{15803D}{(3,3)}\textcolor[HTML]{15803D}{(4,4,1)}\textcolor[HTML]{15803D}{(5,4,1)}\textcolor[HTML]{15803D}{(6,1)}\textcolor[HTML]{15803D}{(5,4)}%
}

\NovelDisplay{M02005R01}{%
\text{第零项列数}=9,%
}
\NovelDisplay{M02005R02}{%
\text{第一项列数}=9+5=14,%
}
\NovelDisplay{M02005R03}{%
\text{第二项列数}=9+2\times5=19.%
}

蔡须坤先盖住答案，按列数把第一项与第二项各写一次。霍雨浩核坐标，卢本韦看保留段有没有误留末列，古月娜再查比较方向。两项都从当前原式求得，不能把第一项再次展开的结果填到第二项的位置。

\NovelDisplay{M02006R01}{%
(0)(1,1,1)(2,1,1)(3,1)(1,1)(2,2,1)(3,2,1)%
}
\NovelDisplay{M02006R02}{%
\quad (4,1)(3,2)(2,2,1)(3,2,1)=\psi(\Omega_{\Omega}\times\Omega_{\omega^{2}})%
}

\NovelDisplay{M02007R01}{%
(0)(1,1,1)(2,1,1)(3,1)(1,1)(2,2,1)(3,2,1)%
}
\NovelDisplay{M02007R02}{%
\quad (4,1)(3,2)(2,2,1)(3,2,1)(4,1)=\psi(\Omega_{\Omega}^{2})%
}

\NovelDisplay{M02008R01}{%
(0)(1,1,1)(2,1,1)(3,1)(1,1)(2,2,1)(3,2,1)%
}
\NovelDisplay{M02008R02}{%
\quad (4,1)(3,2)(3,2)(2,2,1)(3,2,1)(4,1)=\psi(\Omega_{\Omega}^{3})%
}

\NovelDisplay{M02009R01}{%
(0)(1,1,1)(2,1,1)(3,1)(1,1)(2,2,1)(3,2,1)(4,1)(3,2)(4,1)=\psi(\Omega_{\Omega}^{\Omega})%
}

\NovelDisplay{M02010R01}{%
(0)(1,1,1)(2,1,1)(3,1)(1,1)(2,2,1)(3,2,1)(4,1)%
}
\NovelDisplay{M02010R02}{%
\quad (3,2)(4,2)(2,2,1)(3,2,1)(4,1)=\psi(\Omega_{\Omega}^{\Omega_{\Omega}})%
}

风停了，山上的雾却更浓。

霍雨浩将一根炭条握断时，卢本韦把剩下半根接过去，在两人中间续完最后一列。她写到第二坐标便停住，等蔡须坤追出父项，再让古月娜核第三行。

这回没有谁负责把整页讲完。一个人先说前缀，另一个人找会变的位置，第三个人指出为什么能变，最后那个人拿反例位置核查。轮到下一页，四人又将顺序换过来。

蔡须坤发现，自己在台阶上不再只等霍雨浩给答案。压力将问题压到她身上，她也能把问题指给旁人，陪着旁人从最初的零走一遍。

\NovelDisplay{M02011R01}{%
(0)(1,1,1)(2,1,1)(3,1)(1,1)(2,2,1)(3,2,1)(4,1)(3,2)(4,3)=\psi(\Omega_{\Omega+1})%
}

新的后继下标出现，霍雨浩先抬手挡了一下灯。蔡须坤沿他指的位置读完，却没有立刻往更高处试。

``这里变化的是下标里面的那一段，不是外面多加一份。''

古月娜让她指着括号重说。前一组主要在参数里接项、增长乘数，当前这一组则使下标中的结构抵达新的位置。相同的加号若处在不同层，不能共享一句不带位置的解释；外层仍保持当前坍缩，里面的基数由完整下标决定。

卢本韦将一张旧页叠在旁边。他们先前已经见过相似的后继过程，现在可以借熟悉的展开规则核查局部，仍需保留这回的全部前缀。她把两页共有的走法轻轻划过，只在连接不同的地方停下。这样既不为每项重讲一次，也不会因名字眼熟而漏过真正的新处。

霍雨浩先独自写父项，古月娜核末列最高非零行，蔡须坤生成前两项，卢本韦将它们放回长式读层次。分工完成，他们又换了一轮，确认每个人都能接过另一人的位置。

第二轮，蔡须坤在中段忘记了一个补零位置。她没有向同伴要答案，将纸放在膝上，自己从发生分歧的那列往前追。霍雨浩看见她停得很久，几次想共享过去，最后只把温水放到她身边。

``你先别传。''她说。

``没传。''

卢本韦笑了一下，背过身去给别的纸编号。古月娜也暂时合起答案，等蔡须坤把那条链重新追到坏根。过了好一阵，她才抬头，说出刚才漏掉的条件。霍雨浩与她核到同一处，才将本子重新摊开。

压力没有因这一声对上而立刻减轻。蔡须坤上阶时仍需要卢本韦在外侧照看，她也在古月娜停步时回身等了一会儿。四人真正多出来的，是可以各自说明哪里受阻，再共同找到同一个结构位置。

回到横台，蔡须坤把那张没有收到共享的纸收进本子。墨迹有些散，列号却清楚。她不再只将独立复算当作核对同伴的结果；自己能把断开的地方接回来，才使眼前这一层不至于在疲倦时重新变成一个模糊名字。

\NovelDisplay{M02012R01}{%
(0)(1,1,1)(2,1,1)(3,1)(1,1)(2,2,1)(3,2,1)%
}
\NovelDisplay{M02012R02}{%
\quad (4,1)(3,2)(4,3)(4,3)=\psi(\Omega_{\Omega+1}\times2)%
}

\NovelDisplay{M02013R01}{%
(0)(1,1,1)(2,1,1)(3,1)(1,1)(2,2,1)(3,2,1)%
}
\NovelDisplay{M02013R02}{%
\quad (4,1)(3,2)(4,3)(5,2)=\psi(\Omega_{\Omega+1}\times\Omega_2)%
}

\NovelDisplay{M02014R01}{%
(0)(1,1,1)(2,1,1)(3,1)(1,1)(2,2,1)(3,2,1)(4,1)(3,2)(4,3)(5,2)(2,2,1)%
}
\NovelDisplay{M02014R02}{%
\quad (3,2,1)(4,1)=\psi(\Omega_{\Omega+1}\times\Omega_{\Omega})%
}

\NovelDisplay{M02015R01}{%
(0)(1,1,1)(2,1,1)(3,1)(1,1)(2,2,1)(3,2,1)(4,1)(3,2)(4,3)(5,2)(6,3)%
}
\NovelDisplay{M02015R02}{%
\quad =\psi(\Omega_{\Omega+1}\times\psi_{\Omega_{\Omega+1}}(\Omega_{\Omega+1}))%
}

\NovelDisplay{M02016R01}{%
(0)(1,1,1)(2,1,1)(3,1)(1,1)(2,2,1)(3,2,1)%
}
\NovelDisplay{M02016R02}{%
\quad (4,1)(3,2)(4,3)(5,3)=\psi(\Omega_{\Omega+1}^{2})%
}

\NovelDisplay{M02017R01}{%
(0)(1,1,1)(2,1,1)(3,1)(1,1)(2,2,1)(3,2,1)(4,1)(3,2)(4,3,1)=\psi(\Omega_{\Omega+\omega})%
}

\NovelDisplay{M02018R01}{%
(0)(1,1,1)(2,1,1)(3,1)(1,1)(2,2,1)(3,2,1)(4,1)%
}
\NovelDisplay{M02018R02}{%
\quad (3,2)(4,3,1)(5,3)(6,4,1)=\psi(\Omega_{\Omega+\omega\times2})%
}

\NovelDisplay{M02019R01}{%
(0)(1,1,1)(2,1,1)(3,1)(1,1)(2,2,1)(3,2,1)%
}
\NovelDisplay{M02019R02}{%
\quad (4,1)(3,2)(4,3,1)(5,3,1)=\psi(\Omega_{\Omega+\omega^{2}})%
}

\NovelDisplay{M02020R01}{%
(0)(1,1,1)(2,1,1)(3,1)(1,1)(2,2,1)(3,2,1)(4,1)%
}
\NovelDisplay{M02020R02}{%
\quad (3,2)(4,3,1)(5,3,1)(6)(7,1)=\psi(\Omega_{\Omega+\psi(\Omega)})%
}

\NovelDisplay{M02021R01}{%
(0)(1,1,1)(2,1,1)(3,1)(1,1)(2,2,1)(3,2,1)%
}
\NovelDisplay{M02021R02}{%
\quad (4,1)(3,2)(4,3,1)(5,3,1)(6,1)=\psi(\Omega_{\Omega\times2})%
}

\NovelDisplay{M02022R01}{%
(0)(1,1,1)(2,1,1)(3,1)(1,1)(2,2,1)(3,2,1)(4,1)(3,2)(4,3,1)(5,3,1)(6,1)%
}
\NovelDisplay{M02022R02}{%
\quad (2,2)(3,2)=\psi(\Omega_{\Omega\times2}+\Omega_2^{2})%
}

\NovelDisplay{M02023R01}{%
(0)(1,1,1)(2,1,1)(3,1)(1,1)(2,2,1)(3,2,1)(4,1)(3,2)(4,3,1)(5,3,1)(6,1)%
}
\NovelDisplay{M02023R02}{%
\quad (2,2,1)(3,2,1)(4,1)=\psi(\Omega_{\Omega\times2}+\Omega_{\Omega})%
}

\NovelDisplay{M02024R01}{%
(0)(1,1,1)(2,1,1)(3,1)(1,1)(2,2,1)(3,2,1)(4,1)(3,2)(4,3,1)(5,3,1)(6,1)%
}
\NovelDisplay{M02024R02}{%
\quad (2,2,1)(3,2,1)(4,1)(3,2)(4,3)=\psi(\Omega_{\Omega\times2}+\psi_{\Omega_{\Omega+1}}(\Omega_{\Omega+1}))%
}

\NovelDisplay{M02025R01}{%
(0)(1,1,1)(2,1,1)(3,1)(1,1)(2,2,1)(3,2,1)(4,1)(3,2)(4,3,1)(5,3,1)(6,1)%
}
\NovelDisplay{M02025R02}{%
\quad (2,2,1)(3,2,1)(4,1)(3,2)(4,3)(5,2)(2,2,1)(3,2,1)(4,1)=\psi(\Omega_{\Omega\times2}%
}
\NovelDisplay{M02025R03}{%
\quad +\psi_{\Omega_{\Omega+1}}(\Omega_{\Omega+1}\times\Omega_{\Omega}))%
}

\NovelDisplay{M02026R01}{%
(0)(1,1,1)(2,1,1)(3,1)(1,1)(2,2,1)(3,2,1)(4,1)(3,2)(4,3,1)(5,3,1)(6,1)%
}
\NovelDisplay{M02026R02}{%
\quad (2,2,1)(3,2,1)(4,1)(3,2)(4,3,1)(5,3,1)%
}
\NovelDisplay{M02026R03}{%
\quad (6,1)=\psi(\Omega_{\Omega\times2}+\psi_{\Omega_{\Omega+1}}(\Omega_{\Omega\times2}))%
}

\NovelDisplay{M02027R01}{%
(0)(1,1,1)(2,1,1)(3,1)(1,1)(2,2,1)(3,2,1)(4,1)(3,2)(4,3,1)(5,3,1)(6,1)%
}
\NovelDisplay{M02027R02}{%
\quad (3,1)=\psi(\Omega_{\Omega\times2}+\psi_{\Omega_{\Omega+1}}(\Omega_{\Omega\times2})\times\Omega)%
}

\NovelDisplay{M02028R01}{%
(0)(1,1,1)(2,1,1)(3,1)(1,1)(2,2,1)(3,2,1)(4,1)(3,2)(4,3,1)(5,3,1)(6,1)%
}
\NovelDisplay{M02028R02}{%
\quad (3,2)(2,2,1)(3,2,1)(4,1)=\psi(\Omega_{\Omega\times2}+\psi_{\Omega_{\Omega+1}}(\Omega_{\Omega\times2})\times\Omega_{\Omega})%
}

\NovelDisplay{M02029R01}{%
(0)(1,1,1)(2,1,1)(3,1)(1,1)(2,2,1)(3,2,1)(4,1)(3,2)(4,3,1)(5,3,1)(6,1)%
}
\NovelDisplay{M02029R02}{%
\quad (3,2)(3,2)(2,2,1)(3,2,1)(4,1)=\psi(\Omega_{\Omega\times2}+\psi_{\Omega_{\Omega+1}}(\Omega_{\Omega\times2})\times\Omega_{\Omega}^{2})%
}

\NovelDisplay{M02030R01}{%
(0)(1,1,1)(2,1,1)(3,1)(1,1)(2,2,1)(3,2,1)(4,1)(3,2)(4,3,1)(5,3,1)(6,1)%
}
\NovelDisplay{M02030R02}{%
\quad (3,2)(4,3)=\psi(\Omega_{\Omega\times2}+\psi_{\Omega_{\Omega+1}}(\Omega_{\Omega\times2})\times\psi_{\Omega_{\Omega+1}}(\Omega_{\Omega+1}))%
}

\NovelDisplay{M02031R01}{%
(0)(1,1,1)(2,1,1)(3,1)(1,1)(2,2,1)(3,2,1)(4,1)(3,2)(4,3,1)(5,3,1)(6,1)%
}
\NovelDisplay{M02031R02}{%
\quad (3,2)(4,3,1)(5,3,1)(6,1)=\psi(\Omega_{\Omega\times2}+\psi_{\Omega_{\Omega+1}}(\Omega_{\Omega\times2})^{2})%
}

\NovelDisplay{M02032R01}{%
(0)(1,1,1)(2,1,1)(3,1)(1,1)(2,2,1)(3,2,1)(4,1)(3,2)(4,3,1)(5,3,1)(6,1)%
}
\NovelDisplay{M02032R02}{%
\quad (4,3)=\psi(\Omega_{\Omega\times2}+\Omega_{\Omega+1})%
}

\NovelDisplay{M02033R01}{%
(0)(1,1,1)(2,1,1)(3,1)(1,1)(2,2,1)(3,2,1)(4,1)(3,2)(4,3,1)(5,3,1)(6,1)%
}
\NovelDisplay{M02033R02}{%
\quad (4,3)(5,3)=\psi(\Omega_{\Omega\times2}+\Omega_{\Omega+1}^{2})%
}

\NovelDisplay{M02034R01}{%
(0)(1,1,1)(2,1,1)(3,1)(1,1)(2,2,1)(3,2,1)(4,1)(3,2)(4,3,1)(5,3,1)(6,1)%
}
\NovelDisplay{M02034R02}{%
\quad (4,3,1)=\psi(\Omega_{\Omega\times2}+\Omega_{\Omega+\omega})%
}

\NovelDisplay{M02035R01}{%
(0)(1,1,1)(2,1,1)(3,1)(1,1)(2,2,1)(3,2,1)(4,1)(3,2)(4,3,1)(5,3,1)(6,1)%
}
\NovelDisplay{M02035R02}{%
\quad (4,3,1)(5,3,1)(6,1)=\psi(\Omega_{\Omega\times2}\times2)%
}

\NovelDisplay{M02036R01}{%
(0)(1,1,1)(2,1,1)(3,1)(1,1)(2,2,1)(3,2,1)(4,1)(3,2)(4,3,1)(5,3,1)(6,1)%
}
\NovelDisplay{M02036R02}{%
\quad (5,2)=\psi(\Omega_{\Omega\times2}\times\Omega_2)%
}

\NovelDisplay{M02037R01}{%
(0)(1,1,1)(2,1,1)(3,1)(1,1)(2,2,1)(3,2,1)(4,1)(3,2)(4,3,1)(5,3,1)(6,1)%
}
\NovelDisplay{M02037R02}{%
\quad (5,2)(2,2,1)(3,2,1)(4,1)=\psi(\Omega_{\Omega\times2}\times\Omega_{\Omega})%
}

\NovelDisplay{M02038R01}{%
(0)(1,1,1)(2,1,1)(3,1)(1,1)(2,2,1)(3,2,1)(4,1)(3,2)(4,3,1)(5,3,1)(6,1)%
}
\NovelDisplay{M02038R02}{%
\quad (5,3)(4,3,1)(5,3,1)(6,1)=\psi(\Omega_{\Omega\times2}^{2})%
}

\NovelDisplay{M02039R01}{%
(0)(1,1,1)(2,1,1)(3,1)(1,1)(2,2,1)(3,2,1)(4,1)(3,2)(4,3,1)(5,3,1)(6,1)%
}
\NovelDisplay{M02039R02}{%
\quad (5,3)(6,3)(4,3,1)(5,3,1)(6,1)=\psi(\Omega_{\Omega\times2}^{\Omega_{\Omega\times2}})%
}

\NovelDisplay{M02040R01}{%
(0)(1,1,1)(2,1,1)(3,1)(1,1)(2,2,1)(3,2,1)(4,1)(3,2)(4,3,1)(5,3,1)(6,1)%
}
\NovelDisplay{M02040R02}{%
\quad (5,3)(6,4)=\psi(\Omega_{\Omega\times2+1})%
}

\NovelDisplay{M02041R01}{%
(0)(1,1,1)(2,1,1)(3,1)(1,1)(2,2,1)(3,2,1)(4,1)(3,2)(4,3,1)(5,3,1)(6,1)%
}
\NovelDisplay{M02041R02}{%
\quad (5,3)(6,4)(7,4)=\psi(\Omega_{\Omega\times2+1}^{2})%
}

\NovelDisplay{M02042R01}{%
(0)(1,1,1)(2,1,1)(3,1)(1,1)(2,2,1)(3,2,1)(4,1)(3,2)(4,3,1)(5,3,1)(6,1)%
}
\NovelDisplay{M02042R02}{%
\quad (5,3)(6,4,1)=\psi(\Omega_{\Omega\times2+\omega})%
}

\NovelDisplay{M02043R01}{%
(0)(1,1,1)(2,1,1)(3,1)(1,1)(2,2,1)(3,2,1)(4,1)(3,2)(4,3,1)(5,3,1)(6,1)%
}
\NovelDisplay{M02043R02}{%
\quad (5,3)(6,4,1)(7,4,1)=\psi(\Omega_{\Omega\times2+\omega^{2}})%
}

\NovelDisplay{M02044R01}{%
(0)(1,1,1)(2,1,1)(3,1)(1,1)(2,2,1)(3,2,1)(4,1)(3,2)(4,3,1)(5,3,1)(6,1)%
}
\NovelDisplay{M02044R02}{%
\quad (5,3)(6,4,1)(7,4,1)(8,1)=\psi(\Omega_{\Omega\times3})%
}

\NovelDisplay{M02045R01}{%
(0)(1,1,1)(2,1,1)(3,1)(1,1)(2,2,1)(3,2,1)(4,1)(3,2,1)=\psi(\Omega_{\Omega\times\omega})%
}

\NovelDisplay{M02046R01}{%
(0)(1,1,1)(2,1,1)(3,1)(1,1)(2,2,1)(3,2,1)%
}
\NovelDisplay{M02046R02}{%
\quad (4,1)(3,2,1)(3,2,1)=\psi(\Omega_{\Omega\times\omega^{2}})%
}

\NovelDisplay{M02047R01}{%
(0)(1,1,1)(2,1,1)(3,1)(1,1)(2,2,1)(3,2,1)(4,1)(3,2,1)(4,1)=\psi(\Omega_{\Omega^{2}})%
}

\NovelDisplay{M02048R01}{%
(0)(1,1,1)(2,1,1)(3,1)(1,1)(2,2,1)(3,2,1)(4,1)(4)=\psi(\Omega_{\Omega^{\omega}})%
}

\NovelDisplay{M02049R01}{%
(0)(1,1,1)(2,1,1)(3,1)(1,1)(2,2,1)(3,2,1)(4,1)(4,1)=\psi(\Omega_{\Omega^{\Omega}})%
}

\NovelDisplay{M02050R01}{%
(0)(1,1,1)(2,1,1)(3,1)(1,1)(2,2,1)(3,2,1)%
}
\NovelDisplay{M02050R02}{%
\quad (4,1)(4,1)(3,2,1)(4,1)=\psi(\Omega_{\Omega^{\Omega+1}})%
}

\NovelDisplay{M02051R01}{%
(0)(1,1,1)(2,1,1)(3,1)(1,1)(2,2,1)(3,2,1)%
}
\NovelDisplay{M02051R02}{%
\quad (4,1)(4,1)(3,2,1)(4,1)(4,1)=\psi(\Omega_{\Omega^{\Omega\times2}})%
}

\NovelDisplay{M02052R01}{%
(0)(1,1,1)(2,1,1)(3,1)(1,1)(2,2,1)(3,2,1)(4,1)(4,1)(4,1)=\psi(\Omega_{\Omega^{\Omega^{2}}})%
}

上方石桩已近，蔡须坤却在这一处退了半步。

卢本韦扶住她，没有往上拉。``先别踩第二脚。''

她刚才看见两段相同的列，便把后一段按前一段的名字接了进去。古月娜将原式放到她面前，霍雨浩只写父项，没有写右端。

蔡须坤沿那条链重走，发现相同的列落在不同前缀里，组织出的层级也不同。她松开抓紧石沿的手，重新摆正纸，再自己把最后那对括号写完。

走稳之后，她回头接过卢本韦的本子。

``下一道换我只给你写根。''

\NovelDisplay{M02053R01}{%
(0)(1,1,1)(2,1,1)(3,1)(1,1)(2,2,1)(3,2,1)(4,1)(5,1)=\psi(\Omega_{\Omega^{\Omega^{\Omega}}})%
}

下标里的指数塔写到第三层，卢本韦将笔横在纸上，先请其余三人从最外层往里读。

``别一口气念完。每进一层都指出括号。''

蔡须坤先认外层坍缩，再认基数下标；下标本身是幂，幂的指数里又有幂。她先前看见相似的列重复，想当然地给它们安上相同作用，正是没有保留它们所在的层级。如今同一个符号在不同层重新出现，仍然要带着完整上下文说明。

古月娜从里面往外核一遍，霍雨浩再读由左端进入右端时哪些结构保留。两种读法相互检查，谁若把里面的幂移到整座基数之外，下一句的括号便接不上。

卢本韦将三条读线分别停在不同深度。``我们不用每个式子都画。这里画清，后面仍按同一层走的就自己读；等层次又换了，再留一页。''

蔡须坤照着画了一遍，才看见自己刚才的错误并不在最大最深的那层，而在最熟悉的外侧。她为了赶快去看新的结构，把原来应该保留的那部分草草接过。山上的压力将那个草率处拖得很清楚，越想去追深处，越先被外层扯回来。

``那就先把外面接稳。''古月娜说。

霍雨浩没有再共享更多细节。他请蔡须坤闭眼复述自己保留的前缀，再睁眼核左端。卢本韦从末列开始追根，等她追完，才指出哪些新增段仍沿用此前的生成方式。四人各自重算以后，又把那个塔放回原来所在的下标，没有用一个字母遮住。

蔡须坤第二次试步慢了许多，第一处牵扯稳住后，深处那一层才逐渐显出自己的边界。她没有同时推起所有魂力，只让已经辨清的部分按先后接上。古月娜等在上一级，握住她伸来的手，让她自己把脚掌落稳。

卢本韦紧跟上来，站定后先喘了几口气，才想起纸还在霍雨浩手里。她回头伸手，霍雨浩却把干燥的一面朝内，折好才递给她。

雾水将石栏浸得发冷，纸上的塔已经不必靠一眼望到底才能读。他们知道怎样在中间停住，也知道重新接起时应从哪一层继续。

\NovelDisplay{M02054R01}{%
(0)(1,1,1)(2,1,1)(3,1)(1,1)(2,2,1)(3,2,1)(4,1)(5,2)=\psi(\Omega_{\psi_{\Omega_2}(\Omega_2)})%
}

\NovelDisplay{M02055R01}{%
(0)(1,1,1)(2,1,1)(3,1)(1,1)(2,2,1)(3,2,1)%
}
\NovelDisplay{M02055R02}{%
\quad (4,1)(5,2)(6,2)=\psi(\Omega_{\psi_{\Omega_2}(\Omega_2^{2})})%
}

\NovelDisplay{M02056R01}{%
(0)(1,1,1)(2,1,1)(3,1)(1,1)(2,2,1)(3,2,1)(4,1)(5,2,1)=\psi(\Omega_{\psi_{\Omega_2}(\Omega_{\omega})})%
}

\NovelDisplay{M02057R01}{%
(0)(1,1,1)(2,1,1)(3,1)(1,1)(2,2,1)(3,2,1)%
}
\NovelDisplay{M02057R02}{%
\quad (4,1)(5,2,1)(6,2,1)(7,1)=\psi(\Omega_{\psi_{\Omega_2}(\Omega_{\Omega})})%
}

\NovelDisplay{M02058R01}{%
(0)(1,1,1)(2,1,1)(3,1)(1,1)(2,2,1)(3,2,1)(4,1)(5,2,1)(6,2,1)(7,1)(8,2)%
}
\NovelDisplay{M02058R02}{%
\quad =\psi(\Omega_{\psi_{\Omega_2}(\Omega_{\psi_{\Omega_2}(\Omega_2)})})%
}

\NovelDisplay{M02059R01}{%
(0)(1,1,1)(2,1,1)(3,1)(1,1)(2,2,1)(3,2,1)(4,2)=\psi(\Omega_{\Omega_2})%
}

\NovelDisplay{M02060R01}{%
(0)(1,1,1)(2,1,1)(3,1)(1,1)(2,2,1)%
}
\NovelDisplay{M02060R02}{%
\quad (3,2,1)(4,2)(2,2)(3,2)=\psi(\Omega_{\Omega_2}+\Omega_2^{2})%
}

\NovelDisplay{M02061R01}{%
(0)(1,1,1)(2,1,1)(3,1)(1,1)(2,2,1)(3,2,1)%
}
\NovelDisplay{M02061R02}{%
\quad (4,2)(2,2)(3,3,1)=\psi(\Omega_{\Omega_2}+\psi_{\Omega_3}(\Omega_{\omega}))%
}

\NovelDisplay{M02062R01}{%
(0)(1,1,1)(2,1,1)(3,1)(1,1)(2,2,1)(3,2,1)(4,2)(2,2)(3,3,1)(4,3,1)(5,1)%
}
\NovelDisplay{M02062R02}{%
\quad =\psi(\Omega_{\Omega_2}+\psi_{\Omega_3}(\Omega_{\Omega}))%
}

\NovelDisplay{M02063R01}{%
(0)(1,1,1)(2,1,1)(3,1)(1,1)(2,2,1)(3,2,1)(4,2)(2,2)(3,3,1)(4,3,1)(5,2)%
}
\NovelDisplay{M02063R02}{%
\quad =\psi(\Omega_{\Omega_2}+\psi_{\Omega_3}(\Omega_{\Omega_2}))%
}

\NovelDisplay{M02064R01}{%
(0)(1,1,1)(2,1,1)(3,1)(1,1)(2,2,1)(3,2,1)(4,2)(2,2)(3,3,1)(4,3,1)(5,2)%
}
\NovelDisplay{M02064R02}{%
\quad (3,3)(4,3)=\psi(\Omega_{\Omega_2}+\Omega_3^{2})%
}

\NovelDisplay{M02065R01}{%
(0)(1,1,1)(2,1,1)(3,1)(1,1)(2,2,1)(3,2,1)(4,2)(2,2)(3,3,1)(4,3,1)(5,2)%
}
\NovelDisplay{M02065R02}{%
\quad (3,3,1)=\psi(\Omega_{\Omega_2}+\Omega_{\omega})%
}

\NovelDisplay{M02066R01}{%
(0)(1,1,1)(2,1,1)(3,1)(1,1)(2,2,1)(3,2,1)(4,2)(2,2)(3,3,1)(4,3,1)(5,2)%
}
\NovelDisplay{M02066R02}{%
\quad (3,3,1)(4,3,1)(5,2)=\psi(\Omega_{\Omega_2}\times2)%
}

\NovelDisplay{M02067R01}{%
(0)(1,1,1)(2,1,1)(3,1)(1,1)(2,2,1)(3,2,1)(4,2)(2,2)(3,3,1)(4,3,1)(5,2)%
}
\NovelDisplay{M02067R02}{%
\quad (4,3)(3,3,1)(4,3,1)(5,2)=\psi(\Omega_{\Omega_2}^{2})%
}

\NovelDisplay{M02068R01}{%
(0)(1,1,1)(2,1,1)(3,1)(1,1)(2,2,1)(3,2,1)(4,2)(2,2)(3,3,1)(4,3,1)(5,2)%
}
\NovelDisplay{M02068R02}{%
\quad (4,3)(5,4)=\psi(\Omega_{\Omega_2+1})%
}

\NovelDisplay{M02069R01}{%
(0)(1,1,1)(2,1,1)(3,1)(1,1)(2,2,1)(3,2,1)(4,2)(2,2)(3,3,1)(4,3,1)(5,2)%
}
\NovelDisplay{M02069R02}{%
\quad (4,3)(5,4)(6,4)=\psi(\Omega_{\Omega_2+1}^{2})%
}

\NovelDisplay{M02070R01}{%
(0)(1,1,1)(2,1,1)(3,1)(1,1)(2,2,1)(3,2,1)(4,2)(2,2)(3,3,1)(4,3,1)(5,2)%
}
\NovelDisplay{M02070R02}{%
\quad (4,3)(5,4,1)=\psi(\Omega_{\Omega_2+\omega})%
}

\NovelDisplay{M02071R01}{%
(0)(1,1,1)(2,1,1)(3,1)(1,1)(2,2,1)(3,2,1)(4,2)(2,2)(3,3,1)(4,3,1)(5,2)%
}
\NovelDisplay{M02071R02}{%
\quad (4,3)(5,4,1)(6,4,1)(7,2)=\psi(\Omega_{\Omega_2\times2})%
}

\NovelDisplay{M02072R01}{%
(0)(1,1,1)(2,1,1)(3,1)(1,1)(2,2,1)(3,2,1)(4,2)%
}
\NovelDisplay{M02072R02}{%
\quad (2,2)(3,3,1)(4,3,1)(5,2)(4,3,1)=\psi(\Omega_{\Omega_2\times\omega})%
}

\NovelDisplay{M02073R01}{%
(0)(1,1,1)(2,1,1)(3,1)(1,1)(2,2,1)(3,2,1)(4,2)(2,2)(3,3,1)(4,3,1)(5,2)%
}
\NovelDisplay{M02073R02}{%
\quad (4,3,1)(5,2)=\psi(\Omega_{\Omega_2^{2}})%
}

\NovelDisplay{M02074R01}{%
(0)(1,1,1)(2,1,1)(3,1)(1,1)(2,2,1)(3,2,1)(4,2)%
}
\NovelDisplay{M02074R02}{%
\quad (2,2)(3,3,1)(4,3,1)(5,2)(5,2)=\psi(\Omega_{\Omega_2^{\Omega_2}})%
}

\NovelDisplay{M02075R01}{%
(0)(1,1,1)(2,1,1)(3,1)(1,1)(2,2,1)(3,2,1)(4,2)(2,2)(3,3,1)(4,3,1)(5,2)%
}
\NovelDisplay{M02075R02}{%
\quad (6,3)=\psi(\Omega_{\psi_{\Omega_3}(\Omega_3)})%
}

\NovelDisplay{M02076R01}{%
(0)(1,1,1)(2,1,1)(3,1)(1,1)(2,2,1)(3,2,1)%
}
\NovelDisplay{M02076R02}{%
\quad (4,2)(2,2)(3,3,1)(4,3,1)(5,3)=\psi(\Omega_{\Omega_3})%
}

\NovelDisplay{M02077R01}{%
(0)(1,1,1)(2,1,1)(3,1)(1,1,1)=\psi(\Omega_{\Omega_{\omega}})%
}

提升链进入下标之后，霍雨浩把这条较短的式子单独拎出。短列没有替他们免去追根，古月娜要求先解释最初那列为什么保留，再写后续新增段。

\NovelDisplay{M02078R01}{%
\textcolor[HTML]{B45309}{(0)}\textcolor[HTML]{B91C1C}{(1,1,1)}\textcolor[HTML]{B91C1C}{(2,1,1)}\textcolor[HTML]{B91C1C}{(3,1)}\textcolor[HTML]{7E22CE}{(1,1,1)}%
}

\NovelDisplay{M02079R01}{%
\text{原列数}=5, h=3, r=1,%
}
\NovelDisplay{M02079R02}{%
\text{坏根}=\textcolor[HTML]{B45309}{(0)},\text{坏部列数}=4,%
}
\NovelDisplay{M02079R03}{%
(\delta_1,\delta_2,\delta_3)=(1,1,0).%
}

\NovelDisplay{M02080R01}{%
\text{按坏部列序排列的提升阈值：}%
}
\NovelDisplay{M02080R02}{%
(2,2,2,2),%
}
\NovelDisplay{M02080R03}{%
\text{第 }j\text{ 列只允许前 }\tau_j\text{ 行应用阶差。}%
}

\NovelDisplay{M02081R01}{%
\text{原第 }1\text{ 列：}(0)%
}
\NovelDisplay{M02081R02}{%
\longmapsto(q,q), q\ge1.%
}

\NovelDisplay{M02082R01}{%
\text{原第 }4\text{ 列：}(3,1)%
}
\NovelDisplay{M02082R02}{%
\longmapsto(3+q,1+q), q\ge1.%
}

霍雨浩先检查根以前的保留段，蔡须坤再沿原坏部核阈值。每一新增段使用同一份原坏部与阶差，编号 q 为正整数；最高非零行及更后坐标不因段数增多而跟着改变。

\NovelDisplay{M02083R01}{%
\bigl((0)(1,1,1)(2,1,1)(3,1)(1,1,1)\bigr)[0]=\textcolor[HTML]{B45309}{(0)}\textcolor[HTML]{B91C1C}{(1,1,1)}\textcolor[HTML]{B91C1C}{(2,1,1)}\textcolor[HTML]{B91C1C}{(3,1)}%
}

\NovelDisplay{M02084R01}{%
\bigl((0)(1,1,1)(2,1,1)(3,1)(1,1,1)\bigr)[1]=\textcolor[HTML]{B45309}{(0)}%
}
\NovelDisplay{M02084R02}{%
\quad \textcolor[HTML]{B91C1C}{(1,1,1)}\textcolor[HTML]{B91C1C}{(2,1,1)}\textcolor[HTML]{B91C1C}{(3,1)}\textcolor[HTML]{1D4ED8}{(1,1)}\textcolor[HTML]{1D4ED8}{(2,2,1)}\textcolor[HTML]{1D4ED8}{(3,2,1)}\textcolor[HTML]{1D4ED8}{(4,2)}%
}

\NovelDisplay{M02085R01}{%
\bigl((0)(1,1,1)(2,1,1)(3,1)(1,1,1)\bigr)[2]=\textcolor[HTML]{B45309}{(0)}\textcolor[HTML]{B91C1C}{(1,1,1)}\textcolor[HTML]{B91C1C}{(2,1,1)}\textcolor[HTML]{B91C1C}{(3,1)}\textcolor[HTML]{1D4ED8}{(1,1)}\textcolor[HTML]{1D4ED8}{(2,2,1)}%
}
\NovelDisplay{M02085R02}{%
\quad \textcolor[HTML]{1D4ED8}{(3,2,1)}\textcolor[HTML]{1D4ED8}{(4,2)}\textcolor[HTML]{15803D}{(2,2)}\textcolor[HTML]{15803D}{(3,3,1)}\textcolor[HTML]{15803D}{(4,3,1)}\textcolor[HTML]{15803D}{(5,3)}%
}

\NovelDisplay{M02086R01}{%
\text{第零项列数}=4,%
}
\NovelDisplay{M02086R02}{%
\text{第一项列数}=4+4=8,%
}
\NovelDisplay{M02086R03}{%
\text{第二项列数}=4+2\times4=12.%
}

古月娜将第二项的两段新增部分分别圈出，卢本韦从原坏部核它们各自的增量。蔡须坤专查那些本该不动的位置，霍雨浩最后逐列比较。各项虽然可以更长，仍位于原式之下。

真正发生在下标中的提升被认出来，蔡须坤反而问得很慢。

``这里还是要区分，坐标的增量和右边被改变的对象，对吗？''

霍雨浩把左端转过来。末列的最高非零行是第三行，根仍按该行的父项确定，前两行再按阈值应用阶差。规则本身没有新增一个``专门提升下标''的按钮。新处来自完整结构：由这些列组织的作用，落在原来下标里的 Ω 上，使当前对应的下标进一步改变。

古月娜沿右端指出，外面的基数仍有自己的下标，里面新出现的带下标基数又位于这个下标之中。若在这一处少写一层，便会把嵌套的对象移成外层参数。蔡须坤将每层边界各描一次，卢本韦再不看她的纸读回去。

``下标里面还可以有下标。''卢本韦读完，抬头看了眼更高处，``后面的路恐怕就要沿着它接了。''

``先把这一处收稳。''

蔡须坤把二百二十米的山图压到本子下面。石桩已经很近，她却不再像早先那样急着给下一页留一个答案。现在能说明的，是已列出的这一次提升；更长的链会怎样，需要继续检查对应的中间结构，不能只靠一句``既然这里可以，后面都可以''省去。

四人将这一处完整展开重做了一轮。卢本韦算第一项，古月娜算第二项，霍雨浩检查两项是否都来自原式，蔡须坤核保留段与新增段。随后他们把顺序调换，谁也不只守着自己最熟的角色。

最后一轮，卢本韦数错了列。她盯着自己的答案看了很久，终于发现不是坐标错，而是把原式的末列也留在了前缀中。古月娜把紫色末列轻轻圈住，删去以后，数目与坐标一起对上。

``我刚说短列不该漏查，自己就先漏了。''卢本韦摇头。

蔡须坤将那张纸收好：``留着。以后我漏的时候你再拿出来。''

霍雨浩将新的对应放回整段分析。前面的七十米、一百二十米、一百六十米没有被更深的名字淘汰，反而不断作为当前结构中的局部重新出现。他们能继续往里读，正因为那些地方仍保留着各自准确的来路。

卢本韦把磨短的炭条放回纸边。古月娜又核过连接处，霍雨浩才将这一组收好，交给蔡须坤从头复述。

既然下标 \allowbreak{}\NovelInline{I06005}{\Omega}\allowbreak{} 可以被 \allowbreak{}\NovelInline{I06006}{(1,1,1)}\allowbreak{} 提升为 \allowbreak{}\NovelInline{I06007}{\Omega_\omega}\allowbreak{} ，它自然也可以被更长一点的提升链提升到更大的序数。不过，在这一段要展示的，可不仅仅是这一点------

连续一整页的下标读完，卢本韦在纸边画了一条很长的线。

``一条链后面，还能接另一条。''

霍雨浩把她画的线分成几段，每段下面都留下已经核过的 BMS。古月娜沿各段查过去，确认连接处没有凭空消失的列。蔡须坤把最初一张有七十米记号的旧纸也翻出来，放在这页旁边。

两个月才看清的那一处，如今并没有失去作用。它仍在这些更长的结构里，承担着自己原先的位置。

她看了很久，才将纸合到一处。

``最后那几列，从第一项开始看。''

\NovelDisplay{M02087R01}{%
(0)(1,1,1)(2,1,1)(3,1)(1,1,1)(2,1)=\psi(\Omega_{\Omega_{\omega}}+\Omega_{\omega}\times\Omega)%
}

``这一处先按熟悉的结构走。''卢本韦将新页压平。

\NovelDisplay{M02088R01}{%
(0)(1,1,1)(2,1,1)(3,1)(1,1,1)(2,1,1)=\psi(\Omega_{\Omega_{\omega^{2}}})%
}

下标被 \allowbreak{}\NovelInline{I06008}{(1,1,1)(2,1,1)}\allowbreak{} 提升为 \allowbreak{}\NovelInline{I06009}{\Omega_{\omega^2}}\allowbreak{} 。中间的有限项也逐一列在页边。

\NovelDisplay{M02089R01}{%
(0)(1,1,1)(2,1,1)(3,1)(1,1,1)(2,1,1)(2,1,1)=\psi(\Omega_{\Omega_{\omega^{3}}})%
}

现在的提升链是 \allowbreak{}\NovelInline{I06010}{(1,1,1)(2,1,1)(2,1,1)}\allowbreak{} 。下标被提升为 \allowbreak{}\NovelInline{I06011}{\Omega_{\omega^3}}\allowbreak{} 。

\NovelDisplay{M02090R01}{%
(0)(1,1,1)(2,1,1)(3,1)(1,1,1)(2,1,1)(3)=\psi(\Omega_{\Omega_{\omega^{\omega}}})%
}

\NovelDisplay{M02091R01}{%
(0)(1,1,1)(2,1,1)(3,1)(1,1,1)(2,1,1)(3,1)=\psi(\Omega_{\Omega_{\Omega}})%
}

现在的提升链是 \allowbreak{}\NovelInline{I06012}{(1,1,1)(2,1,1)(3,1)}\allowbreak{} 。下标被提升为 \allowbreak{}\NovelInline{I06013}{\Omega_{\Omega}}\allowbreak{} 。在下标的下标上，又出现了一个 \allowbreak{}\NovelInline{I06014}{\Omega}\allowbreak{} ？

\NovelDisplay{M02092R01}{%
(0)(1,1,1)(2,1,1)(3,1)(1,1,1)(2,1,1)(3,1)(1,1,1)=\psi(\Omega_{\Omega_{\Omega_{\omega}}})%
}

没错，这个地方的 \allowbreak{}\NovelInline{I06015}{\Omega}\allowbreak{} 也会被提升。卢本韦低声笑了一下，将那个更深的下标重新描清。

\NovelDisplay{M02093R01}{%
(0)(1,1,1)(2,1,1)(3,1)(1,1,1)(2,1,1)%
}
\NovelDisplay{M02093R02}{%
\quad (3,1)(1,1,1)(2,1,1)(3,1)=\psi(\Omega_{\Omega_{\Omega_{\Omega}}})%
}

由于眼前的提升链以 \allowbreak{}\NovelInline{I06016}{(3,1)}\allowbreak{} 结尾，它代表着 \allowbreak{}\NovelInline{I06017}{\Omega}\allowbreak{} ，所以在这个提升链的后面还能再接上一段提升链。而这一段提升链的后面还能再接\ldots\ldots 它们一起，形成了一段无穷的``链式提升效应''。而在这无穷无尽的提升链的终点，是------

\NovelDisplay{M02094R01}{%
(0)(1,1,1)(2,1,1)(3,1)(2)=\psi(\Omega_{\Omega_{\Omega_{\ddots}}})=\psi(\Lambda)%
}

最后的末列只有第一行非零。蔡须坤将此前容易一并带动的后两行留在原处，卢本韦重新找到第一行的根；霍雨浩按零阶差生成有限项，古月娜再核它们各自低于边界。

\NovelDisplay{M02095R01}{%
(0)\textcolor[HTML]{B45309}{(1,1,1)}\textcolor[HTML]{B91C1C}{(2,1,1)}\textcolor[HTML]{B91C1C}{(3,1)}\textcolor[HTML]{7E22CE}{(2)}%
}

\NovelDisplay{M02096R01}{%
\text{原列数}=5, h=1, r=2,%
}
\NovelDisplay{M02096R02}{%
\text{坏根}=\textcolor[HTML]{B45309}{(1,1,1)},\text{坏部列数}=3,%
}
\NovelDisplay{M02096R03}{%
(\delta_1,\delta_2,\delta_3)=(0,0,0).%
}

\NovelDisplay{M02097R01}{%
\text{原第 }2\text{ 列：}(1,1,1)%
}
\NovelDisplay{M02097R02}{%
\longmapsto(1,1,1), q\ge1.%
}

\NovelDisplay{M02098R01}{%
\text{原第 }4\text{ 列：}(3,1)%
}
\NovelDisplay{M02098R02}{%
\longmapsto(3,1), q\ge1.%
}

霍雨浩先检查根以前的保留段，蔡须坤再沿原坏部核阈值。每一新增段使用同一份原坏部与阶差，编号 q 为正整数；最高非零行及更后坐标不因段数增多而跟着改变。

\NovelDisplay{M02099R01}{%
\bigl((0)(1,1,1)(2,1,1)(3,1)(2)\bigr)[0]=(0)\textcolor[HTML]{B45309}{(1,1,1)}\textcolor[HTML]{B91C1C}{(2,1,1)}\textcolor[HTML]{B91C1C}{(3,1)}%
}

\NovelDisplay{M02100R01}{%
\bigl((0)(1,1,1)(2,1,1)(3,1)(2)\bigr)[1]=(0)%
}
\NovelDisplay{M02100R02}{%
\quad \textcolor[HTML]{B45309}{(1,1,1)}\textcolor[HTML]{B91C1C}{(2,1,1)}\textcolor[HTML]{B91C1C}{(3,1)}\textcolor[HTML]{1D4ED8}{(1,1,1)}\textcolor[HTML]{1D4ED8}{(2,1,1)}\textcolor[HTML]{1D4ED8}{(3,1)}%
}

\NovelDisplay{M02101R01}{%
\bigl((0)(1,1,1)(2,1,1)(3,1)(2)\bigr)[2]=(0)\textcolor[HTML]{B45309}{(1,1,1)}\textcolor[HTML]{B91C1C}{(2,1,1)}\textcolor[HTML]{B91C1C}{(3,1)}\textcolor[HTML]{1D4ED8}{(1,1,1)}\textcolor[HTML]{1D4ED8}{(2,1,1)}%
}
\NovelDisplay{M02101R02}{%
\quad \textcolor[HTML]{1D4ED8}{(3,1)}\textcolor[HTML]{15803D}{(1,1,1)}\textcolor[HTML]{15803D}{(2,1,1)}\textcolor[HTML]{15803D}{(3,1)}%
}

\NovelDisplay{M02102R01}{%
\text{第零项列数}=4,%
}
\NovelDisplay{M02102R02}{%
\text{第一项列数}=4+3=7,%
}
\NovelDisplay{M02102R03}{%
\text{第二项列数}=4+2\times3=10.%
}

卢本韦数完列，将第一项中的蓝段与第二项中的蓝段逐列重合。绿色段另由原坏部按两倍阶差生成，原来的保留段没有再受一次提升。蔡须坤核坐标，古月娜与霍雨浩最后检查各有限项都小于原式。

最后一组有限项的 BOCF 对应也分别写出。它们显示下标链怎样一层层延伸，整条无限链仍由极限参数表示。

\NovelDisplay{M02103R01}{%
\text{展开第 }0\text{ 项所对应的序数：}\psi(\Omega_{\Omega})%
}
\NovelDisplay{M02103R02}{%
\text{展开第 }1\text{ 项所对应的序数：}\psi(\Omega_{\Omega_{\Omega}})%
}
\NovelDisplay{M02103R03}{%
\text{展开第 }2\text{ 项所对应的序数：}\psi(\Omega_{\Omega_{\Omega_{\Omega}}})%
}
\NovelDisplay{M02103R04}{%
\text{展开第 }3\text{ 项所对应的序数：}\psi(\Omega_{\Omega_{\Omega_{\Omega_{\Omega}}}})%
}

末列只剩一个二时，蔡须坤先在纸上将它单独圈出，随后又将圈擦去。

``不能只盯最后这一列。''

古月娜让她从第一行追根。根停在前面的一一一，最高非零行也是第一行，因此没有更前行可以施加阶差；后两行照原值保留。删去末列以后，坏部是从根到末列之前的那三列，新段依规则重复。展开第一项接一段，展开第二项接两段；这里每一项都由同一个原式求出，并非将已经求出的第一项再展开一次。

卢本韦将有限项数完，霍雨浩把它们各自的对应放在旁边。每增加一段，右端的下标链便继续深入；这些明确写出的有限项都低于原式，没有哪一个因为纸面已经写了许多层就变成无限链本身。

``那最后这条无限链，具体指什么？''蔡须坤问。

霍雨浩指出前文留下的序列：先由一个 Ω 开始，再将它放进下一个 Ω 的下标，重复形成嵌套，取整条序列的极限。那个极限是最小的 Ω 下标不动点，满足相应的固定点关系。无穷的省略号用于表示这条生成方式，不能只在任何长塔后面加省略号便声称得到了同一个对象。

古月娜随即点住最外层的坍缩。``这个仍要留下。''

蔡须坤原先正准备在山图旁写下不动点的名字，听见这句便停住。链的极限是坍缩所用的参数，EBO 则是外层坍缩后的可数序数；二者的位置不同。前面的分析曾多次要求保留括号，此时到了最醒目的终点，也不能将最后一层略去。

卢本韦把那张还没写完的山图翻到背面，先请蔡须坤完整说一遍。蔡须坤说根，霍雨浩追问任意有限项的生成；卢本韦说新增段，古月娜再问有限项与边界有何区别。每人都被问到最容易省略的一处，答案便没有只停在终点的名字上。

唐三已经将上方的绳结系好，没有催。十一人的行李放在较低的平石上，谁也没有因为石桩近在眼前便先越过去。蔡须坤看着那几份独立写完的展开，才发现手指一直攥着同一截炭条，指腹都被磨黑了。

她将炭条放下，先活动手指，再回到台阶前。四人没有同时发力，按各自已经说清的结构接起运转。压力仍然压着肩背，却不再把几处来路不同的牵扯揉成一团。蔡须坤向前伸手，接住卢本韦递来的本子。她先将一只脚落到最上一级的边缘，等霍雨浩核完末页，古月娜与卢本韦也分别收稳运转，才准备迈出最后一步。

这一段分析能够收住，考核却还在继续。更高的五百米目标保留在雾里，今天不必先给它写结论。他们先把眼前这一组从有限展开说到边界，再将最后那一点刻度留给真正落稳的脚掌。

\allowbreak{}\NovelInline{I06018}{\mathrm{EBO}}\allowbreak{}，全称为 \allowbreak{}\NovelInline{I06019}{\mathrm{Extended\ Buchholz\ Ordinal}}\allowbreak{}。这里 \allowbreak{}\NovelInline{I06020}{\Lambda}\allowbreak{} 表示最小的 \allowbreak{}\NovelInline{I06021}{\Omega}\allowbreak{} 下标不动点：它是序列 \allowbreak{}\NovelInline{I06022}{\{\Omega,\Omega_\Omega,\Omega_{\Omega_\Omega},\dots\}}\allowbreak{} 的极限，也是满足 \allowbreak{}\NovelInline{I06023}{\alpha=\Omega_\alpha}\allowbreak{} 的最小序数。BOCF 的坍缩参数可以写成这条无穷下标链，也可以记为 \allowbreak{}\NovelInline{I06024}{\Lambda}\allowbreak{}；EBO 则是坍缩后的可数序数 \allowbreak{}\NovelInline{I06025}{\psi(\Lambda)}\allowbreak{}。

二百二十米的横台，在最后一遍核对后亮了起来。

蔡须坤先收住气息，才把第二只脚踏上去。她肩上的重量仍在，脚底却不再有那种被未明接缝猛然扯回的感觉。卢本韦紧跟着站稳，扶住膝盖喘了几口气；古月娜收起指尖的银纹，霍雨浩把末页的墨迹吹干，才在山图上落下这一点。

\NovelDisplay{M02104R01}{%
(0)(1,1,1)(2,1,1)(3,1)(2)\longleftrightarrow\quad220\text{ 米}%
}

夏夜星空看过他们各自写出的规则，又请四人交换着解释。蔡须坤说清坏根，卢本韦说明新增段，霍雨浩接出任意有限项，古月娜指出有限项与所趋边界的分别。没有人用``后面都差不多''带过最后一句。

唐三将后面的同伴接上横台，小舞坐到石边，长长地呼出一口气。蔡兴坤把黑子鸡抱稳，诺当将装水的袋子放在众人中央。叶阳和薛峰一页页数过笔记，奥斯卡把刚热好的补给分开，先递给手仍在抖的几个人。

蔡须坤仰头。离前五百米还差二百八十米，离山顶还差二千二百八十米。两年的第一考依旧压在每个人心上，而这一段，终于能从最初的零说到最后的边界。

卢本韦将最后一页夹好。

``今天到这里。''

古月娜点了点头。霍雨浩把炭条的碎末收在纸角，没有急着给下一页起名字。蔡须坤背起本子，等十一人都歇过气，才沿着已经走稳的路一同回营。

\chapter{第一百一十二章：混沌暴走}\label{ux7b2cux4e00ux767eux4e00ux5341ux4e8cux7ae0ux6df7ux6c8cux66b4ux8d70}

回到山脚时，石屋里还没有点灯。

下山那一路，谁也没急着说下一段该怎么推。唐三接过了前面几人的行李，小舞将剩下的水分了一圈，叶阳和薛峰一路护着装笔记的包，到了屋前才放下来。

蔡须坤落在最后半步。横台上的压力已经离开肩背，胸口却还有一点闷，她只当是今天收力收得慢。卢本韦把水递过来时，她还笑着说，明天得先把磨破的鞋底补一补。

卢本韦把写着二百二十米的那页压进本子，正要去拿灯，身后忽然响了一声。

蔡须坤手里的水碗掉了。

她还站在门槛外，一只手按着胸口，另一只手扶住门框。指缝底下透出几道淡金的纹，亮到腕子，便被一层灰黑吞了回去。

霍雨浩先看见那层灰。

``都退开。''

话音刚落，门框就在蔡须坤掌下裂了。没有木屑飞出来，裂口却沿着纹理一路爬到门楣，像有什么东西从木头里抽走了。

她想松手，手指却不听使唤。

底下那团混沌醒了。

上岛以后，它一直沉在井底，没有动过。此刻却猛地翻上来，撞在本源下面。蔡须坤眼前一黑，连卢本韦扶住她的那只手，都隔了好一会儿才觉出来。

``别往屋里走。''古月娜说。

她伸手托住蔡须坤另一边肩膀。银色的光刚贴上去，灰气便沿着她的指尖缠过来。古月娜当即将光收窄，只留在衣料外侧，示意卢本韦向石柱旁的平石挪。

卢本韦掌下已经聚起魂力，见那点灰沿着银光攀动，便将力收回，改用肩膀抵住蔡须坤。她一手扶背，一手托住肘弯，让那只攥紧门框的手先松下来。

``别跟它抢。''古月娜低声说。

``我知道。''

卢本韦答得很快，脚下却没敢快。蔡须坤每挪一步，身上的气息都要断一下。短短几步路，她们停了两次，直到人坐稳，卢本韦才发现自己的手心也出了汗。

唐三将门前的行李推远。小舞叫叶阳、薛峰清开山道，诺当站在平石外侧，挡住想往前凑的蔡兴坤。

``先让雨浩看。''

蔡兴坤没再挤，抱着黑子鸡停在几步外。她的手也按着心口，袖口底下，一点暖金忽明忽暗。

夏夜星空从后面的石阶落下来，掌心按住两根石柱之间的地面。一层薄光绕过众人，罩住平石。逸出去的灰气撞在光上，发出很轻的一声裂响。

``今晚不上山了。''她说，``先稳住她。''

霍雨浩已经蹲下去，搭住蔡须坤的手腕。

``看着我。能听见吗？''

蔡须坤点了一下头，牙关咬得很紧。

``混沌。''她只说出两个字。

霍雨浩认得。码头上那一次，也是金纹与灰黑交替着从她腕下透出来。他将精神力放得极细，顺着她仍能自行运转的魂力往里探，想先隔开那两股互相冲撞的力量。

最外面的一层，他接住了。

蔡须坤发颤的手指停了一瞬。霍雨浩稍稍松了口气，将后面的精神力接过去。

那条细线忽然断了。

他甚至没有看清它在哪里断的。刚才还能辨认的脉络，在灰黑里翻了一下，前后的位置便全乱了。他留在外面的另一缕精神力，也跟着往里滑。

霍雨浩立即收手，喉咙里却先涌上来一股腥甜。

卢本韦看见他脸白了。

``怎么样？''

他摇头，抹去嘴角的血，又看向蔡须坤。

``别替我撑。把你自己的呼吸留下。''

她刚才分出了一点魂力，想将他陷进去的精神力送回来。就这一点，井底便又撞上来一次。平石从她身下裂开一道细缝，背后那处熟悉的温热也随之一紧。

蔡须坤立刻将手背到身后。

``你别出来。''

那点温热没有散，只贴在本源边缘，替她守住了最后一点清明。

她很清楚，那里面已经没有一具能替她挡在前面的身体。牢大留下的魂还在养着，每多用一分力，都要从那一点残存里支取。

她隔着衣料轻轻按了一下，将翻到嘴边的那声求援咽了回去。温热仍贴在那里，没有离开，她便靠着这一点，将霍雨浩的声音重新辨了出来。

霍雨浩听见她那句``你别出来''，手停了半息。

随后，他将一层极薄的寒气铺在蔡须坤手腕外侧。冰蝎的虚影在肩后亮起，寒气顺着她紊乱的魂力缓缓合拢，试着将正在外涌的那一截冻住。

冰结成了。

蔡须坤的指尖却先失去了血色。

冰层下面，那点灰黑仍在动，反而沿着停滞的经脉，向更细处分开。

霍雨浩猛地撤去寒气。

``不行。''

这一次，他没有再往里送力。

自己的手在抖，他便将手掌撑在石沿上，稳住以后，才重新握住蔡须坤没有被灰气缠住的指尖。

``我压不住它。''他说，``再冻下去，先伤的是你。''

古月娜将蔡须坤泛白的指尖护在掌间，等冰意退尽，才略略松手。

霍雨浩盯着那一点缓缓回来的血色，没有立刻换第三种办法。刚才极寒合拢时，他以为终于碰到了能封住的东西；可先停下来的，是蔡须坤本来就在挣扎的魂力。他将这处记住，强迫自己把视线从灰气上移开，先看她有没有在呼吸。

蔡须坤想说知道了，胸口却像被什么攥紧，只呼出半口气。

夏夜星空掌下的光忽然暗了一块。她另一只手抬起，在石柱上重重一按，一道细光顺着山阶窜进了雾里。

古月娜额角已经见汗。她将银光托在蔡须坤肩背外面，只护住身体，没有再去碰那团灰。

``雨浩，别让她断了回应。''

霍雨浩点头，一声一声叫她的名字。

蔡须坤听得见，却答不出来。井底正在往上翻，她自己那点熟悉的力量，被挤到了很远的地方。

山上的雾就在这时分开了。

最菜萌新走下最后几级石阶，素袍的下摆还沾着水。他只看了一眼，便伸手按住蔡须坤身前那片灰气。

夏夜星空掌下的薄光终于不再向内塌。

``别再往她身子里填。''老人说。

霍雨浩抬起头：``两次都没接住。精神力进去就断，极寒也封不住。''

最菜萌新点了点头，蹲到蔡须坤面前。

他没有将灰气一掌压回去。指尖在她腕前停住，一缕暗金从石柱底下升起，绕着她身外缓缓合拢。灰气向外一撞，暗金也跟着向外让了一寸。

石柱上，落下一点细细的石粉。

老人的眉头皱了起来。

``底下这一团，藏得真深。''

他将手移到蔡须坤眉前，仍隔着一点距离。

``听我说。外面的，我替你拦。你只找自己还能使唤的那一点。''

蔡须坤的嘴唇动了动。

``找\ldots\ldots 不着。''

``别往最深处找。刚才谁叫你，你听见了？''

她的目光艰难地转向霍雨浩。

``听见了，就从那里接。''

老人抬眼看他：``你别再探进去。只接她主动送出来的。''

霍雨浩将精神力留在她指尖之外，没有越过皮肤。

等了很久，一点微弱的回应终于碰到他。

``在。''他说，``我接到了。''

蔡须坤闭了一下眼。那一点是她自己的，没有被灰黑卷走，也没有被旁人的力量代替。她顺着这一点往回找，找到了指尖，找到了掌心，慢慢找回半只手。

混沌猛地翻了上来。

她的肩膀向前一折，最菜萌新托在身外的暗金同时沉下去。平石边沿无声地缺了一角，老人袖口一震，将那点外逸的灰拢进了掌间。

卢本韦伸手垫住蔡须坤的额头，没让她磕到石上。

``慢一点。我们都在。''

蔡须坤攥住了她的袖口。

最菜萌新等这一下攥稳，才将停在眉前的手往下移了半寸。

``就是这个。留住。''

她没有再试着把整团混沌按下去，只护着刚找回的那一点，将它与身下翻涌的灰黑分开。霍雨浩随她的呼吸一点点接稳，古月娜托住肩背，卢本韦的袖子一直攥在她手里。

老人掌间的暗金缓缓收拢。

灰气仍在冲，每撞一次，两根石柱就轻轻震一下。夏夜星空按着地面的手没有抬。

``只堵着，待我撤了，它还会撞。''最菜萌新看着那几缕灰，``先把这一截接进你自己的魂力运转里。慢慢收，别去动更深处。''

蔡须坤迟疑了一下。

``我接不住呢？''

``我拦着。慢慢来。''

她将刚找回的力量往回收，带着贴在边缘的那一缕混沌，一点点引回自己的魂力运转中。灰黑刚接进去，她的手指便猛地攥紧。霍雨浩没有添力，只将她主动送出的回应接稳，让她知道自己还握着哪一端。

那缕灰没有再冲出来。

它沿着刚刚稳住的经脉缓缓绕过，随她自己的魂力转了一圈，终于留下了。

蔡须坤额角的汗落到石上。她停了一口气，才去接第二缕。

这一缕刚经过腕间，便将熟悉的运转挤偏了些。她原想照着前一回接下去，指腹却忽然一麻，送向霍雨浩的回应也跟着淡了。

``先停。''他说。

她停在原处，最菜萌新也没有再放进下一缕。卢本韦将另一只袖子递到她手里，等她慢慢找回掌心的知觉。

蔡须坤重新接上时，比先前还慢。她不再想着一次走完，先等腕间那一处能顺畅接续，再让灰黑沿着自己的魂力往前挪。额上的汗落了一滴，又落一滴，霍雨浩始终没有替她往前送。

``接上了。''这一次，是她自己先说的。

最菜萌新把其余外涌的力量拦在掌下，等她收稳一段，再放开一段。起先每接一次，她都要停很久；后来，那些灰黑不再一味往外撞，顺着她已经接通的位置，逐渐并入她能够调动的魂力。

直到这回暴走涌出的最后一截也被吸收，蔡须坤才终于吸进一口完整的气。

她额头上全是冷汗。

过了几息，手腕上的灰也退了。

最菜萌新没有立即松手，又等了一阵，确认她能自己续上呼吸，才慢慢撤去掌间的力量。

``这一回涌出来的，收住了。''他说，``底下还有多少，先别去碰。''

霍雨浩垂下头，这才发现自己的衣背已经湿透。

夏夜星空仍守着石柱，直到最菜萌新向她点头，才将掌下的光松开少许。蔡兴坤抱着黑子鸡走近一步，张了张嘴，先把怀里的鸡交给诺当，才将水放到平石边。

``现在能喝吗？''

老人看了看蔡须坤：``她自己能咽，就让她抿一点。''

杯口碰到嘴唇时，蔡须坤才发现自己渴得厉害。她喝得很慢，蔡兴坤便一直托着杯底，没让那点水洒到她颤抖的手上。

奥斯卡将补给递到霍雨浩手里，没催他吃。蔡兴坤放下水杯，接回诺当怀里的黑子鸡。唐三走过来，与诺当一同将平石外的碎块清开，小舞从屋里抱来干净的被褥。

蔡须坤松开卢本韦的袖子，低低说了声抱歉。

卢本韦将皱透的衣料拉平，没理这一句，只问她：``冷不冷？''

她摇头。

吸收进体内的混沌之力，随她的魂力缓缓转过一周。

一层近乎透明的薄光，随之在她身周荡开。它沿着平石铺展，越过卢本韦垫在她背后的被褥，绕到霍雨浩脚边。外缘浮着几缕细碎的灰，刚要散开，又随她的呼吸缓缓合住。

霍雨浩忽然觉出，自己留在她指尖旁的精神细线稳了。

刚才他拼命接也接不住，此刻没有再添力，那点回应却清楚地留在了那里。古月娜掌下的银光也不再受牵扯，终于能慢慢收回。两人的力量各自运转，之间仍留着蔡须坤铺开的那层光。

蔡须坤低头看去。平石、石柱、围在身边的伙伴，都在那片光里。她只动了一个念头，最外缘那点想往诺当脚边扑的灰，便顺着边界绕开，安静地伏了下去。

她起先没敢再动，等了两息，才试着将最外侧那一点收近。薄光掠过一块碎石，停到诺当鞋尖之前，贴着她想留下的位置稳住了。

卢本韦没有放出母猪扽魂，只觉出掌间残留的魂力顺畅了些。古月娜将最后一缕银光撤开，霍雨浩也慢慢收细精神回应，三人各自的力量都没有混进那片薄光里。

``还在。''霍雨浩说。

蔡须坤知道他指的是自己。刚才要拼命才能留住的那一点，此刻顺着呼吸，清清楚楚地握在她手中。她没有让薄光继续铺开，只把几人的脚边留在里面。

最菜萌新的指尖仍停在蔡须坤腕前。正要撤开时，他忽然顿了一下。

她腕下浮起几道淡金的纹，只留在皮肤上。随着血脉重新平稳下来，一缕气息从中透出，被老人清楚地感受到了。

他抬起眼皮，重新看向她。

先前他辨出的是她与蔡兴坤身上共有的坤神血脉。眼下显露的，是另一股更深的血脉气息。混沌不再翻搅，将它遮住的力量退去，帝皇瑞兽特有的气息才完整透了出来。

老人的手缓缓收回，目光停在她脸上。

``帝皇瑞兽。''

声音很轻，霍雨浩却抬起了头。

古月娜看向蔡须坤，没有开口。卢本韦仍坐在她身侧，将被褥往她背后垫好，没有松开扶着她的手。

最菜萌新从几人的反应里看明白了，目光重新落回蔡须坤脸上。

``你从帝天身边出来的？''

``是。''

老人沉默了片刻，又看了一眼她掌心渐渐收敛的金纹。

``原来已经走到这里了。''

蔡须坤试着松开手指，石上那层薄光便随之动了一下。她低头望去。

``这是\ldots\ldots''

``你的领域。''最菜萌新说，``刚才吸收的混沌之力，经你自己运转，催生了这一层领域。不过刚觉醒，先收住。''

蔡须坤看见霍雨浩尚未擦净的嘴角，又看见卢本韦被攥皱的袖子。古月娜还蹲在身侧，其他人都围在几步以内，谁也没走远。

她将掌心轻轻收拢，光停在众人脚下，没有再往外延伸。

``就叫果糕领域。''

``要记到本子里吗？''霍雨浩问。

``记。''她声音还有点哑，``就这四个字，别给我添长。''

卢本韦嘴角终于动了一下：``刚才连话都说不出，名字倒定得快。''

蔡须坤想回她一句，力气没续上，只轻轻碰了碰她的袖口。

卢本韦没再逗她，将滑到膝边的被褥拉好。霍雨浩终于松开撑着石沿的手，坐稳了些。

领域只维持了片刻，边缘就开始发虚。蔡须坤依着老人的提醒，一点点将它收回，没再试着多留。最后一线薄光收回身周，她便疲惫地靠回被褥里。

山脚开始起风。唐三已经把屋里的床铺挪到背风处，诺当收好门前的碎片，又找来木条，将裂开的门框暂且支住。夏夜星空等蔡须坤身上的气息平复下来，才让众人把她扶进屋。

卢本韦与古月娜一左一右，走得比刚才还慢。霍雨浩想伸手，被小舞先按着肩膀坐下，奥斯卡将补给重新塞回他掌心。

``你也吃。''

他这才低头，发现自己还没咬过一口。

灯是叶阳点起来的，薛峰把桌边收出一块空地，蔡兴坤将水放在床头。最菜萌新接过唐三搬来的矮凳，坐在床沿外侧，夏夜星空则守在半开的门边。人都坐稳以后，屋里才渐渐有了说话声。

霍雨浩等她的呼吸重新平稳下来，才低声问：``大长老的良序之力，能不能帮她稳住底下的混沌？''

``或许能派上用场。''最菜萌新看了一眼她腕上退尽的灰，``说到他，你们知道大长老为什么实力这么强吗？''

几人望向老人。

``不只因为他有着世间唯一的同位共振三魂核，也因为他是良序之神的大供奉。''

古月娜替蔡须坤掖好被角，听到``三魂核''，手上的动作顿了一下。

霍雨浩抬起头：``第三颗也参与同位共振？三颗之间的牵引，要怎么稳住？''

``老夫和帝天，都是阴阳互补双魂核。''最菜萌新说，``两颗魂核彼此承接，既要互补，也得维持平衡。可一旦有了第三颗，每一颗的变化都会同时牵动另外两颗，另外两颗的变化又会传回来。三颗一同运转，结果就非常不可控。''

他抬起三根手指，随后慢慢收拢。

``大长老身上的良序之力，弥补了这一点。三核之间本会失去控制的牵引，在他体内能够有序地接续下去。他既能让三颗魂核持续运转，也能将它们的力量一起用出来。''

霍雨浩想起自己刚才探入那团混沌时，精神力还没来得及调转，就已被几股牵扯截断。直到此刻，他才稍稍明白，让三颗魂核始终安稳地运转，是怎样的本事。

``有了第三颗，还得让前两颗都重新适应它。''他低声说。

``正是。前两颗维持的平衡，也会因此改变。''老人答道。

霍雨浩慢慢松开握紧的手。登阶时，他们至少能停下来，把没看清的一处拆开重查；魂核在体内运转，却不能等他写完这一页，再去接下一步。大长老能随时出手，那份稳定便须一直维持着。

``所以，他的实力强于老夫。''

最菜萌新说得很平常，接着又道：``当年他来蓬莱岛，我们在这里切磋过七场。按他的胜负算，三胜四败。''

卢本韦抬眼：``那您赢了四场？''

``是。头五场，他三胜，老夫两胜。到最后两场，老夫借了这里的主场优势，才连胜两场。''老人朝山上看了一眼，``在蓬莱岛，坤神之力为我所用。最后那两场能击败他，靠的就是这一点。这七场的胜负，不能用来断定老夫的实力在他之上。''

山风从门缝里吹进来，灯芯晃了一下。蔡须坤望向老人，手指在被褥上轻轻动了动。

``那良序之力，和我身上的混沌\ldots\ldots''

``是同一个品级。''最菜萌新说，``良序之力来自良序神王Zermelo，与混沌之力同为宇宙本源力量。方才接住你涌出来的力量时，老夫已经感受到了。只是这股混沌之力的来历，还需要继续弄清。''

他停了一下，目光落到卢本韦身上。

``大长老的良序之力，甚至能压制母猪扽魂。''

卢本韦扶着蔡须坤肩头的手微微一顿。

她低头看了一眼自己的掌心，像是终于想通了一件搁置许久的事。

``古月娜的银龙扽魂，我能压制。''她说，``可一到良序圣殿，我就觉得自己的力量像是在被净化。我还以为，单纯是大长老的魂力太高了。''

古月娜也转过目光，等着老人的回答。

``不错。''最菜萌新说，``从力量来源上说，大长老能把帝天压制成那样，正有良序之力的缘故。你在圣殿里的感受，也不能只用魂力高低来解释。''

他略停了一下。

``不只你，连老夫也被这股力量压制过。''

霍雨浩抬起头，看了一眼老人刚才替蔡须坤挡住混沌的手。连这样的修为，也会受到良序之力的压制。

``并且，你们也知道，母猪扽魂是神级扽魂。''最菜萌新的目光重新落回卢本韦身上，``连它也会受到压制，你该明白，大长老身上的良序之力有多不寻常了。''

卢本韦缓缓合拢手指，想起在良序圣殿里，那股力气明明还在，却怎么也无法像平时一样尽情调动。当时没能说清的不适，此刻终于有了来由。

最菜萌新看着她，神色缓和了些。

``你甚至比你曾祖都优秀很多。''

卢本韦怔了怔，抬头看向老人。

``可刚才，我连她身上的灰气都不敢碰。''

``看见古月娜的银光受了牵扯，你就收住了手。这一步做得对。''最菜萌新说。

她听完，半晌才低低应了一声，将扶着蔡须坤的那只手放稳。

最菜萌新又道：``记住，你的扽魂也是魂灵。等你到达 \allowbreak{}\NovelInline{I06026}{(0)(1,1,1)(2,2,2)}\allowbreak{} 境界，自然会知道这一点。''

``魂灵\ldots\ldots''卢本韦轻轻重复了一遍。

老人没有继续解释。她等了片刻，便点了点头，将这句话连同那个境界一起记在心里。

古月娜仍望着最菜萌新。

``前辈，当年天州兽潮之后，大长老去找胡赞的事，您知道多少？''

``那年兽潮之后，他又到老夫这里来过一趟，留了一阵。''最菜萌新说，``这些事，他亲口和我说过。''

古月娜稍稍坐直：``既然良序之力连帝天都能压制，胡赞为什么还能牵制住他？''

``胡赞身上那股帅气力量，似乎超乎寻常。用了帅气教的禁术，抽取旁人的帅气填进自己身上，他就有了不弱于大长老的实力。''

老人缓缓收回搁在膝上的手。

``天州那一仗结束后，大长老往下追查，已经知道胡赞干的好事。他去找胡赞，本就是要追究这些事。可那时候，帅气教在外的行径，主要还是掠杀魂兽；藏在胡赞那股力量底下的东西，远没有那么容易查清。''

``而且，只要他们交起手来，胡赞就会抽取百姓的帅气。''

霍雨浩的眉头慢慢拧紧。

``一边打，一边从百姓身上补充？''

``是。''最菜萌新说，``战况越紧，他越会把主意打到百姓身上。两人本就没有谁能稳拿下谁，接着打，受害的百姓只会更多。所以，他们交过手，却没有真正放开到你死我活的地步。''

古月娜沉默了片刻，又问：``胡赞知道大长老有这些顾虑？''

``他当然知道。他甚至还提过，要与大长老联手，一致对付魂兽。''

古月娜的神色冷了下来。

``掠杀魂兽的是他，惹出兽潮的也是他。最后，他还想让大长老站到他那边？''

``正是这个意思。''最菜萌新说，``那是胡赞的打算，大长老没有答应。''

古月娜望着灯芯，过了一会儿才开口：``他躲在后面，让魂兽与守城的人打，自己还要拿城里的百姓补充力量。''

``大长老去找他时，也正是带着这笔账。''

屋里静了下来。胡赞没有在天州城外露过面，提到那些死伤时，却哪里都绕不开他。霍雨浩把手里的补给放下，记住了老人说的那句``超乎寻常''，没有急着拿刚听懂的三魂核去给胡赞排一个高低。

蔡须坤靠在被褥里，听到这里，轻轻攥住了身侧的衣料。卢本韦察觉到她的动作，将手覆了上去。古月娜看了两人一眼，没再追问。

最菜萌新重新看向蔡须坤。

``但力量的品级，和你现在能用出多少，是两回事。你刚才收进去的那些，先练到每一缕都能听你的话。大长老能稳住三颗魂核，你现在先把自己的气息稳住。''

蔡须坤试着调匀呼吸，没有再去碰井底。

``今晚留在山脚养着。''最菜萌新说，``底下那团混沌还没弄清。领域也要慢慢稳，别急着登阶。''

她点了点头，将手背到身后，轻轻贴住那处重新平稳下来的温热。

霍雨浩坐在旁边，没再放出精神探测。卢本韦拨亮了灯芯，古月娜将门掩上些，替她挡住门口的风。

最菜萌新又守了一阵，直到蔡须坤自己将魂力续过一周，腕上也没有再泛灰，才起身。夏夜星空送他到门外，回来时没有走远，在门边另放了一把椅子。

唐三把夜里的照应分好，没有再议明天能上多少级。叶阳和薛峰去整理散在外面的行李，小舞收拾床边的水，奥斯卡将热着的补给留了一份。

霍雨浩吃完，才拿笔在新纸上写下``果糕领域''。写到一半，他停住，望了一眼床上的人。

``刚才冷不冷？''

蔡须坤知道他问的是极寒合拢的那一下。

``冷。''她说，``你收得也快。''

他垂眼看了看自己已经不抖的手。蔡须坤伸出指尖，轻轻碰了一下他的手背，这一次没有夹着灰，也不用他探进本源里找。霍雨浩将纸压好，替她把手放回被里。

卢本韦坐回床沿，见她还睁着眼，便将那册本子挪到床边的桌角，让她躺着也能看见。

``二百二十米写着呢，睡一觉也丢不了。''

蔡须坤终于合上眼。古月娜坐在灯光照不到的另一侧，等那点呼吸渐渐拉长，才放松了一直绷着的肩背。

山上的石阶渐渐没进夜色。

写着二百二十米的那册本子，仍好好压在屋里的桌上。

\chapter{第一百一十三章：两套记号}\label{ux7b2cux4e00ux767eux4e00ux5341ux4e09ux7ae0ux4e24ux5957ux8bb0ux53f7}

蔡须坤重新走到山阶前，是在那场暴走过去几天以后。

这几天，最菜萌新没有再替她往深处探。她能做的事很少：把已经收下的混沌之力慢慢运转，觉得腕间发涩就停，歇过一阵，再从能够接续的地方开始。霍雨浩来过几次，每次都先问她能不能自己收住，没有再将精神力送进那口井。

卢本韦把磨破的鞋底补好了。针脚歪了一处，她原想拆，被蔡须坤接过去，连着另外一只一同收进包里。

``能踩稳就行。''

说这句话时，她的手没有再抖。古月娜仍等她将一周魂力自行运转完，才把门边的本子递过来。

十一人回到二百二十米横台，花的时间比下山时更长。旧路的压力没有消失，蔡须坤却没有抢着把果糕领域放出去。她先站稳，让自己确认，这一步是靠什么接住的。

夏夜星空在横台上等着，身边压着一叠新页。纸上两列记号隔得很开，一边反复出现I，另一边则密密地套着M与ψ。

卢本韦一眼看中了左边。蔡须坤翻到右边，读了半页，才发觉左手不知何时已经捏紧纸角。

``我试M这边。''她说。

``那我先理I。''卢本韦把自己的本子放到旁边，``你读不下去的时候，先告诉我停在哪层，别一个人闷着。''

霍雨浩把两本都翻开。薛峰没有挤到最前，只接过最后一页，拿炭条在页角标好次序。那时没人留意，他留出的空白，比旁人宽得多。

夏夜星空先让他们翻回旧页。早些时候经过的区间，在这份讲义里分别记作\allowbreak{}\NovelInline{I06027}{0 \sim \psi(\Omega_\omega)}\allowbreak{}与\allowbreak{}\NovelInline{I06028}{\psi(\Omega_\omega) \sim \psi(I)}\allowbreak{}。蔡须坤看见熟悉的左列，肩膀稍松了一点，又被右栏全新的写法拉住了目光。

页首给I一侧留的标识是\allowbreak{}\NovelInline{I06029}{I = 2\;1-2}\allowbreak{}，随后仍使用\allowbreak{}\NovelInline{I06030}{I}\allowbreak{}；另一套的标识是\allowbreak{}\NovelInline{I06031}{M = 2-2}\allowbreak{}，后面只保留\allowbreak{}\NovelInline{I06032}{M}\allowbreak{}这一层符号。霍雨浩用笔杆压住那两个标识，提醒大家，这是讲义规定的记法，不能把里面的数字、分号和横线拆成一道普通的减法题。

``也别看见两边最后相等，就把两边每一个字母都直接换掉。''他说。

蔡须坤点头，将M侧最先用到的规则单独抄了下来。

\NovelDisplay{M02105R01}{%
\psi_M(0) = 1%
}

\NovelDisplay{M02106R01}{%
\psi_M(x+1) = \psi_M(x) \times \omega%
}

她顺着第二条念，到了乘号后面停了一下。``右乘ω。顺序也得留住。''

\NovelDisplay{M02107R01}{%
\psi_M(\# \sim M) = 2\;1-\{\psi_M(\# \sim x)\}%
}

第三条后面的说明，是\allowbreak{}\NovelInline{I06033}{M}\allowbreak{}折叠容许点。讲义在这里没有给出更完整的构造细则。古月娜在旁边留了一道短线，叫他们先保留这条规则的原样，不要把位置标记当成一个已经算出的序数，往后面的每个括号里随便代。

卢本韦原来觉得，I侧短得多，这一次总能省些事。可她刚要把第一处M换成Ω，霍雨浩就让她看下一行：同一个更深的边界，在两套记号里占据的位置已经不同。她慢慢把笔收回来，与蔡须坤并排核完下面几组。

旧BMS的零补回脑中即可，纸上仍按此前的习惯省去每列末尾的零，只保留开头的零列。断行也只是续写，下一行不另起一个BMS。

\NovelDisplay{M02108R01}{%
(0)(1,1)%
}

\NovelDisplay{M02109R01}{%
= \psi(\Omega)%
}

\NovelDisplay{M02110R01}{%
= \psi(M)%
}

\NovelDisplay{M02111R01}{%
(0)(1,1,1)%
}

\NovelDisplay{M02112R01}{%
= \psi(\Omega_\omega)%
}

\NovelDisplay{M02113R01}{%
= \psi(M \times \omega)%
}

\NovelDisplay{M02114R01}{%
(0)(1,1,1)(2,1,1)%
}

\NovelDisplay{M02115R01}{%
= \psi(\Omega_{\omega^2})%
}

\NovelDisplay{M02116R01}{%
= \psi(M \times \omega^2)%
}

\NovelDisplay{M02117R01}{%
(0)(1,1,1)(2,1,1)(3,1)%
}

\NovelDisplay{M02118R01}{%
= \psi(\Omega_\Omega)%
}

\NovelDisplay{M02119R01}{%
= \psi(M \times \psi_M(M))%
}

\NovelDisplay{M02120R01}{%
(0)(1,1,1)(2,1,1)(3,1)(1,1,1)%
}

\NovelDisplay{M02121R01}{%
= \psi(\Omega_{\Omega_\omega})%
}

\NovelDisplay{M02122R01}{%
= \psi(M \times \psi_M(M \times \omega))%
}

\NovelDisplay{M02123R01}{%
(0)(1,1,1)(2,1,1)(3,1)(1,1,1)(2,1,1)(3,1)%
}

\NovelDisplay{M02124R01}{%
= \psi(\Omega_{\Omega_\Omega})%
}

\NovelDisplay{M02125R01}{%
= \psi(M \times \psi_M(M \times \psi_M(M)))%
}

\NovelDisplay{M02126R01}{%
(0)(1,1,1)(2,1,1)(3,1)(1,1,1)(2,1,1)(3,1)(1,1,1)(2,1,1)(3,1)%
}

\NovelDisplay{M02127R01}{%
= \psi(\Omega_{\Omega_{\Omega_\Omega}})%
}

\NovelDisplay{M02128R01}{%
= \psi(M \times \psi_M(M \times \psi_M(M \times \psi_M(M))))%
}

\NovelDisplay{M02129R01}{%
(0)(1,1,1)(2,1,1)(3,1)(2)%
}

\NovelDisplay{M02130R01}{%
= \psi(I)%
}

\NovelDisplay{M02131R01}{%
= \psi(M^2)%
}

两人核到旧横台对应的那一式时，都没有立即动脚。左列仍是她们几天前才真正说清的结构，右边却多出一套陌生的解释。

蔡须坤顺着M侧看完，低声说：``相等的是整条对应，不是随便挑两个内部符号，就能说它们同义。''

夏夜星空没有接着讲，只将下面的纸向前推了一寸。横台上面第一阶的压力，便在这一寸之后等着她们。

他们先用M侧的前两条规则核了几个有限的起点。由零出发，接一个后继，再接下一个，每一步右乘ω，先前认得的有限幂次便能够在这套记法里找到来路。

蔡须坤没有顺势将这条后继规则搬到所有极限上。她把第三条仍原样留在旁边：到了容许点的折叠，不能只把有限后继做过几回，就说已经给出一般的极限规则。

卢本韦原本打算只抄I侧，听到这里，也将M的规则留了一份。她用I讲得较短，仍要知道同伴为什么会在另一侧停住；否则她递出的帮助，可能只是一个蔡须坤早已知道、却仍接不回去的结果。

薛峰将最早的短式翻回正面，请两人分别指出最外层ψ在哪里。这个问题太短，谁都以为不必再答，真正遮住后面的长页以后，却替接下来的比较留住了一个可靠的起点。

最先露出来的，是在既有边界之后接有限的尾部。

卢本韦读I侧很快，读到加一时甚至没有停。蔡须坤却把M平方外面的那层ψ重新描了一遍。她容易跟住的是括号内的加一，容易漏掉的，是这个加一到底加在谁的参数中。

``你先说最后得到什么。''卢本韦提醒。

``一个最外层ψ的值。里面先有M平方，再接这一段。''

``那最外面的ψ还在？''

蔡须坤抬头看了她一眼，没嫌这个问题太简单。她刚才在心里读过一遍，确实险些把这一层省掉。

薛峰将两人的纸沿同一处等号对齐。他先不问深处为什么这样对应，只替她们查三件能当场看清的事：左列有没有抄少一列，外层括号有没有闭合，换到下一行时，前面的等号接的是哪一段。他找出一处漏写的右括号，蔡须坤补上以后，才发现自己先前那口气总接得不顺，正卡在这儿。

轮到Ω与Ω下标的几处时，卢本韦原想凭熟悉的名字替她收尾。霍雨浩却让两人分别把答案遮住。I侧给出了较短的整体结构，M侧还必须把函数的下标与它吃进的参数分开。下标里已有一个ψ，并不表示括号内的参数也跟着消失。

他们把这几行一起留下。每遇到一处看似能缩写的地方，先用左列认清是哪一个BMS，再读完整的右栏，免得只认出了一个结果的名字。

薛峰先将前面共同的几列描在纸背，再把每条真正多出来的尾部留在正面。这样翻过去看时，变化很醒目；翻回来时，完整左列仍然能被接上。他没有把共同前缀另取一个字母，让后来的人还要去别页猜它是什么。

蔡须坤照着试了一行，发现自己终于能把注意力留在当前参数中。她又将完整前缀补回正面，防止独立复算时忘记那段从哪里来。

``这一处看着像只加了一点。''她说，``可要是外面那层放错，我连加在哪里都不知道。''

霍雨浩便让她只说明眼前这一条，不提前猜以后所有加法的样子。先把当前对应读准确，再找下一条真正相同的位置，进度虽然慢些，前一页却不必因为后面一处变化而整本推翻。

\NovelDisplay{M02132R01}{%
(0)(1,1,1)(2,1,1)(3,1)(2)(1)%
}

\NovelDisplay{M02133R01}{%
= \psi(I+1)%
}

\NovelDisplay{M02134R01}{%
= \psi(M^2+1)%
}

\NovelDisplay{M02135R01}{%
(0)(1,1,1)(2,1,1)(3,1)(2)(1,1)%
}

\NovelDisplay{M02136R01}{%
= \psi(I+\Omega)%
}

\NovelDisplay{M02137R01}{%
= \psi(M^2+\psi_{\psi_M(M^2)}(1))%
}

\NovelDisplay{M02138R01}{%
\psi_{\psi_M(M^2)}(1) = \psi_I(1) = \Omega%
}

\NovelDisplay{M02139R01}{%
\psi_M(M)%
}

\NovelDisplay{M02140R01}{%
(0)(1,1,1)(2,1,1)(3,1)(2)(1,1)(2,2,1)(3,2,1)(4,2)(3)(2,2)%
}

\NovelDisplay{M02141R01}{%
= \psi(I+\Omega_2)%
}

\NovelDisplay{M02142R01}{%
= \psi(M^2+\psi_{\psi_M(M^2)}(2))%
}

\NovelDisplay{M02143R01}{%
(0)(1,1,1)(2,1,1)(3,1)(2)(1,1,1)%
}

\NovelDisplay{M02144R01}{%
= \psi(I+\Omega_\omega)%
}

\NovelDisplay{M02145R01}{%
= \psi(M^2+\psi_{\psi_M(M^2)}(\omega))%
}

\NovelDisplay{M02146R01}{%
(0)(1,1,1)(2,1,1)(3,1)(2)(1,1,1)(2,1,1)(3,1)%
}

\NovelDisplay{M02147R01}{%
= \psi(I+\Omega_\Omega)%
}

\NovelDisplay{M02148R01}{%
= \psi(M^2+\psi_{\psi_M(M^2)}(M))%
}

\NovelDisplay{M02149R01}{%
\psi_{\psi_M(M^2)}(M)=\psi_{\psi_M(M^2)}(\psi_M(M)) = \Omega_\Omega%
}

\NovelDisplay{M02150R01}{%
(0)(1,1,1)(2,1,1)(3,1)(2)(1,1,1)(2,1,1)(3,1)(1,1,1)%
}

\NovelDisplay{M02151R01}{%
= \psi(I+\Omega_\Omega)%
}

\NovelDisplay{M02152R01}{%
= \psi(M^2+\psi_{\psi_M(M^2)}(M \times \omega))%
}

\NovelDisplay{M02153R01}{%
\psi_{\psi_M(M^2)}(M \times \omega)=\Omega_{\Omega_\omega}%
}

\NovelDisplay{M02154R01}{%
(0)(1,1,1)(2,1,1)(3,1)(2)(1,1,1)(2,1,1)(3,1)(1,1,1)(2,1,1)(3,1)%
}

\NovelDisplay{M02155R01}{%
= \psi(I+\Omega_{\Omega_\omega})%
}

\NovelDisplay{M02156R01}{%
= \psi(M^2+\psi_{\psi_M(M^2)}(M \times \psi_M(M)))%
}

\NovelDisplay{M02157R01}{%
\psi_{\psi_M(M^2)}(M \times \psi_M(M))=\Omega_{\Omega_\Omega}%
}

\NovelDisplay{M02158R01}{%
(0)(1,1,1)(2,1,1)(3,1)(2)(1,1,1)(2,1,1)(3,1)(1,1,1)(2,1,1)(3,1)(1,1,1)%
}

\NovelDisplay{M02159R01}{%
= \psi(I+\Omega_{\Omega_\Omega})%
}

\NovelDisplay{M02160R01}{%
= \psi(M^2+\psi_{\psi_M(M^2)}(M \times \psi_M(M \times \psi_M(M))))%
}

\NovelDisplay{M02161R01}{%
(0)(1,1,1)(2,1,1)(3,1)(2)(1,1,1)(2,1,1)(3,1)(2)%
}

\NovelDisplay{M02162R01}{%
= \psi(I+\psi_I(I))%
}

\NovelDisplay{M02163R01}{%
= \psi(M^2+\psi_{\psi_M(M^2)}(M^2))%
}

\NovelDisplay{M02164R01}{%
\psi_{\psi_M(M^2)}(M^2) = \psi_I(I)%
}

到了讲义特别解释的那几处，蔡须坤终于问出了第一句真正难的问题：``既然有别的写法也能认出Ω，为什么这里偏要留着这一套？''

霍雨浩点着后面的几行。``眼下相同，不等于后面的接法就可以混写。这份表自己也在留标准型的位置。先按同一套接下去，等要比较别的写法时，再把条件一起带上。''

卢本韦把自己最初画的那根直接替换箭头擦掉了。她留下两边完整的式子，连那一点炭灰也没有急着吹走。

接下来换动的地方，藏进了函数的下标。

卢本韦第一次把这一段读完，蔡须坤没有应。她眼前M侧的ψ一个接一个，最里侧刚看清，外面包着它的那一层便从记忆里滑走。头顶的压力还在往肩上压，她不自觉地抬手揉了揉眼角，连纸上的加号都像在向两边挪。

霍雨浩没替她念。他用一张空纸遮住右栏，只让她读到一对能独立闭合的括号。等她说清这是谁的下标，才露出下一截。

``别一口气背完。先告诉我，外面这一层变了吗？''

蔡须坤看了许久，摇头。

``那就把没变的留着，找它里面真正动的地方。''

卢本韦也慢下来。I侧虽然短，下标内的加一仍然不是给整个式子加一。她一度将那一小段读到Ω外面，薛峰的炭条随即停在纸边，等她自己沿括号走回来。

山风卷起最上面的一张纸。唐三伸手按住，顺势将装本子的包往两人脚后挪了些，免得她们在原地调整时绊住。小舞把水递到卢本韦另一侧，没有让她空出正压着纸的手。那些没有直接写答案的人，也正在让这道题能继续写下去。

最后，蔡须坤将M侧的函数头读清，卢本韦才接着念I侧。几个相似的尾部被她们并排放在一起：有的只多一个有限项，有的把整个旧段重新放进参数里；不能看见相同的数字，就认作同一类增量。

\NovelDisplay{M02165R01}{%
(0)(1,1,1)(2,1,1)(3,1)(2)(2)%
}

\NovelDisplay{M02166R01}{%
= \psi(I+\psi_{\Omega_{\psi_I(I)+1}}(1))%
}

\NovelDisplay{M02167R01}{%
= \psi(M^2+\psi_{\psi_{\psi_M(M^2)}(M^2+1)}(1))%
}

\NovelDisplay{M02168R01}{%
\psi_{\psi_M(M^2)}(M^2+1) = \Omega_{\psi_I(I)+1}%
}

\NovelDisplay{M02169R01}{%
(0)(1,1,1)(2,1,1)(3,1)(2)(2)(2)%
}

\NovelDisplay{M02170R01}{%
= \psi(I+\psi_{\Omega_{\psi_I(I)+1}}(2))%
}

\NovelDisplay{M02171R01}{%
= \psi(M^2+\psi_{\psi_{\psi_M(M^2)}(M^2+1)}(2))%
}

\NovelDisplay{M02172R01}{%
(0)(1,1,1)(2,1,1)(3,1)(2,1)%
}

\NovelDisplay{M02173R01}{%
= \psi(I+\psi_{\Omega_{\psi_I(I)+1}}(\Omega))%
}

\NovelDisplay{M02174R01}{%
= \psi(M^2+\psi_{\psi_{\psi_M(M^2)}(M^2+1)}(M))%
}

\NovelDisplay{M02175R01}{%
(0)(1,1,1)(2,1,1)(3,1)(2,1)(1,1,1)%
}

\NovelDisplay{M02176R01}{%
= \psi(I+\psi_{\Omega_{\psi_I(I)+1}}(\Omega_\omega))%
}

\NovelDisplay{M02177R01}{%
= \psi(M^2+\psi_{\psi_{\psi_M(M^2)}(M^2+1)}(M \times \omega))%
}

\NovelDisplay{M02178R01}{%
(0)(1,1,1)(2,1,1)(3,1)(2,1)(1,1,1)(2,1,1)(3,1)(2)%
}

\NovelDisplay{M02179R01}{%
= \psi(I+\psi_{\Omega_{\psi_I(I)+1}}(I))%
}

\NovelDisplay{M02180R01}{%
= \psi(M^2+\psi_{\psi_{\psi_M(M^2)}(M^2+1)}(M^2))%
}

\NovelDisplay{M02181R01}{%
(0)(1,1,1)(2,1,1)(3,1)(2,1)(1,1,1)(2,1,1)(3,1)(2)(1,1,1)(2,1,1)(3,1)(2)%
}

\NovelDisplay{M02182R01}{%
= \psi(I+\psi_{\Omega_{\psi_I(I)+1}}(I)+\psi_I(I))%
}

\NovelDisplay{M02183R01}{%
= \psi(M^2+\psi_{\psi_{\psi_M(M^2)}(M^2+1)}(M^2)+\psi_{\psi_M(M^2)}(M^2))%
}

\NovelDisplay{M02184R01}{%
(0)(1,1,1)(2,1,1)(3,1)(2,1)(1,1,1)(2,1,1)(3,1)(2,1)%
}

\NovelDisplay{M02185R01}{%
= \psi(I+\psi_{\Omega_{\psi_I(I)+1}}(I)+\psi_{\Omega_{\psi_I(I)+1}}(\Omega))%
}

\NovelDisplay{M02186R01}{%
= \psi(M^2+\psi_{\psi_{\psi_M(M^2)}(M^2+1)}(M^2)+\psi_{\psi_{\psi_M(M^2)}(M^2+1)}(M))%
}

\NovelDisplay{M02187R01}{%
(0)(1,1,1)(2,1,1)(3,1)(2,1)(1,1,1)(2,1,1)(3,1)(2,1)(1,1,1)(2,1,1)%
}
\NovelDisplay{M02187R02}{%
(3,1)(2)%
}

\NovelDisplay{M02188R01}{%
= \psi(I+\psi_{\Omega_{\psi_I(I)+1}}(I)\times 2)%
}

\NovelDisplay{M02189R01}{%
= \psi(M^2+\psi_{\psi_{\psi_M(M^2)}(M^2+1)}(M^2)\times 2)%
}

\NovelDisplay{M02190R01}{%
(0)(1,1,1)(2,1,1)(3,1)(2,1)(2)%
}

\NovelDisplay{M02191R01}{%
= \psi(I+\psi_{\Omega_{\psi_I(I)+1}}(I+1))%
}

\NovelDisplay{M02192R01}{%
= \psi(M^2+\psi_{\psi_{\psi_M(M^2)}(M^2+1)}(M^2+1))%
}

\NovelDisplay{M02193R01}{%
(0)(1,1,1)(2,1,1)(3,1)(2,1)(2,1)%
}

\NovelDisplay{M02194R01}{%
= \psi(I+\psi_{\Omega_{\psi_I(I)+1}}(I+\Omega))%
}

\NovelDisplay{M02195R01}{%
= \psi(M^2+\psi_{\psi_{\psi_M(M^2)}(M^2+1)}(M^2+\psi_{\psi_M(M^2)}(1)))%
}

\NovelDisplay{M02196R01}{%
(0)(1,1,1)(2,1,1)(3,1)(2,1)(2,1)(1,1,1)(2,1,1)(3,1)(2)%
}

\NovelDisplay{M02197R01}{%
= \psi(I+\psi_{\Omega_{\psi_I(I)+1}}(I+\psi_I(I)))%
}

\NovelDisplay{M02198R01}{%
= \psi(M^2+\psi_{\psi_{\psi_M(M^2)}(M^2+1)}(M^2+\psi_{\psi_M(M^2)}(M^2)))%
}

\NovelDisplay{M02199R01}{%
(0)(1,1,1)(2,1,1)(3,1)(2,1)(2,1)(1,1,1)(2,1,1)(3,1)(2,1)%
}

\NovelDisplay{M02200R01}{%
= \psi(I+\psi_{\Omega_{\psi_I(I)+1}}(I+\psi_I(I))+\psi_{\Omega_{\psi_I(I)+1}}(\Omega))%
}

\NovelDisplay{M02201R01}{%
= \psi(M^2+\psi_{\psi_{\psi_M(M^2)}(M^2+1)}(M^2+\psi_{\psi_M(M^2)}(M^2))+\psi_{\psi_{\psi_M(M^2)}(M^2+1)}(M))%
}

\NovelDisplay{M02202R01}{%
(0)(1,1,1)(2,1,1)(3,1)(2,1)(2,1)(1,1,1)(2,1,1)(3,1)(2,1)(2,1)(1,1,1)%
}
\NovelDisplay{M02202R02}{%
(2,1,1)(3,1)(2)%
}

\NovelDisplay{M02203R01}{%
= \psi(I+\psi_{\Omega_{\psi_I(I)+1}}(I+\psi_I(I))\times 2)%
}

\NovelDisplay{M02204R01}{%
= \psi(M^2+\psi_{\psi_{\psi_M(M^2)}(M^2+1)}(M^2+\psi_{\psi_M(M^2)}(M^2))\times 2)%
}

她们在这一组里停了两个下午。第二天下山，卢本韦还在用手指比划那处下标，差点把水袋扣到本子上。薛峰把袋子接过去，翻到她白天改过的一行，让她只说这一个加号归哪层。

卢本韦说对了。他才把水还给她。

``你现在也会拿题拦路了。''她说。

薛峰笑了一下，却把那一行又默写了一遍。

M侧真正难以维持的位置，是一层函数的参数里，又装进同一类函数的时候。

蔡须坤最初还能凭声音记住顺序。等念到第三个相似的ψ，她忽然分不清自己是在重复前一层，还是已经走进新的参数。手底下写出来的括号全都闭合，式子的层次却已经错开了。

卢本韦从I侧找出相应的那几行，把外面共同的部分念给她听。蔡须坤跟到一半，胸口生出一点熟悉的闷，立刻停了下来。

她没有再碰井底。那天暴走时的感觉仍在，她怕自己因为着急，又去向更深处讨一分力。古月娜看见她按住手腕，便将纸接过来，先让她退到横台内侧。

``这页不会跑。''古月娜说。

蔡须坤歇了好一会儿，腕上的不适没有继续加重，才重新把本子接住。霍雨浩坐到她身边，没有共享一整篇答案，只把她刚才漏认的那对括号指出来。

``你念到这里的时候，以为已经回到外面了。其实还在参数里面。''

她沿着那一处走完，才看见真正的出口。

这一回她换了办法。每往内读一层，先说清这一层的下标，再说参数；读完以后，必须说自己回到了哪一层。卢本韦以I侧短式核整体，薛峰对照每次返回的位置，霍雨浩专听有没有哪一处跳过。古月娜在旁边守着她的呼吸，没有替她把剩下的层次念完。

几人的说话声变得很慢。纸面上那些看似相同的函数，也终于各自有了能被找到的位置。

\NovelDisplay{M02205R01}{%
(0)(1,1,1)(2,1,1)(3,1)(2,1)(2,1)(2)%
}

\NovelDisplay{M02206R01}{%
= \psi(I+\psi_{\Omega_{\psi_I(I)+1}}(I+\psi_I(I)+1))%
}

\NovelDisplay{M02207R01}{%
= \psi(M^2+\psi_{\psi_{\psi_M(M^2)}(M^2+1)}(M^2+\psi_{\psi_M(M^2)}(M^2)+1))%
}

\NovelDisplay{M02208R01}{%
(0)(1,1,1)(2,1,1)(3,1)(2,1)(2,1)(2,1)(2)%
}

\NovelDisplay{M02209R01}{%
= \psi(I+\psi_{\Omega_{\psi_I(I)+1}}(I+\psi_I(I)\times 2+1))%
}

\NovelDisplay{M02210R01}{%
= \psi(M^2+\psi_{\psi_{\psi_M(M^2)}(M^2+1)}(M^2+\psi_{\psi_M(M^2)}(M^2)\times 2+1))%
}

\NovelDisplay{M02211R01}{%
(0)(1,1,1)(2,1,1)(3,1)(2,1)(3)%
}

\NovelDisplay{M02212R01}{%
= \psi(I+\psi_{\Omega_{\psi_I(I)+1}}(I+\psi_{\Omega_{\psi_I(I)+1}}(1)))%
}

\NovelDisplay{M02213R01}{%
= \psi(M^2+\psi_{\psi_{\psi_M(M^2)}(M^2+1)}(M^2+\psi_{\psi_{\psi_M(M^2)}(M^2+1)}(1)))%
}

\NovelDisplay{M02214R01}{%
(0)(1,1,1)(2,1,1)(3,1)(2,1)(3)(2)%
}

\NovelDisplay{M02215R01}{%
= \psi(I+\psi_{\Omega_{\psi_I(I)+1}}(I+\psi_{\Omega_{\psi_I(I)+1}}(1)+1))%
}

\NovelDisplay{M02216R01}{%
= \psi(M^2+\psi_{\psi_{\psi_M(M^2)}(M^2+1)}(M^2+\psi_{\psi_{\psi_M(M^2)}(M^2+1)}(1)+1))%
}

\NovelDisplay{M02217R01}{%
(0)(1,1,1)(2,1,1)(3,1)(2,1)(3)(2,1)(3)%
}

\NovelDisplay{M02218R01}{%
= \psi(I+\psi_{\Omega_{\psi_I(I)+1}}(I+\psi_{\Omega_{\psi_I(I)+1}}(1)\times 2))%
}

\NovelDisplay{M02219R01}{%
= \psi(M^2+\psi_{\psi_{\psi_M(M^2)}(M^2+1)}(M^2+\psi_{\psi_{\psi_M(M^2)}(M^2+1)}(1)\times 2))%
}

\NovelDisplay{M02220R01}{%
(0)(1,1,1)(2,1,1)(3,1)(2,1)(3)(3)%
}

\NovelDisplay{M02221R01}{%
= \psi(I+\psi_{\Omega_{\psi_I(I)+1}}(I+\psi_{\Omega_{\psi_I(I)+1}}(2)))%
}

\NovelDisplay{M02222R01}{%
= \psi(M^2+\psi_{\psi_{\psi_M(M^2)}(M^2+1)}(M^2+\psi_{\psi_{\psi_M(M^2)}(M^2+1)}(2)))%
}

\NovelDisplay{M02223R01}{%
(0)(1,1,1)(2,1,1)(3,1)(2,1)(3,1)%
}

\NovelDisplay{M02224R01}{%
= \psi(I+\psi_{\Omega_{\psi_I(I)+1}}(I+\psi_{\Omega_{\psi_I(I)+1}}(M)))%
}

\NovelDisplay{M02225R01}{%
= \psi(M^2+\psi_{\psi_{\psi_M(M^2)}(M^2+1)}(M^2+\psi_{\psi_{\psi_M(M^2)}(M^2+1)}(M)))%
}

\NovelDisplay{M02226R01}{%
(0)(1,1,1)(2,1,1)(3,1)(2,1)(3,1)(1,1,1)(2,1,1)(3,1)(2)%
}

\NovelDisplay{M02227R01}{%
= \psi(I+\psi_{\Omega_{\psi_I(I)+1}}(I+\psi_{\Omega_{\psi_I(I)+1}}(I)))%
}

\NovelDisplay{M02228R01}{%
= \psi(M^2+\psi_{\psi_{\psi_M(M^2)}(M^2+1)}(M^2+\psi_{\psi_{\psi_M(M^2)}(M^2+1)}(M^2)))%
}

\NovelDisplay{M02229R01}{%
(0)(1,1,1)(2,1,1)(3,1)(2,1)(3,1)(2)%
}

\NovelDisplay{M02230R01}{%
= \psi(I+\psi_{\Omega_{\psi_I(I)+1}}(I+\psi_{\Omega_{\psi_I(I)+1}}(I)+1))%
}

\NovelDisplay{M02231R01}{%
= \psi(M^2+\psi_{\psi_{\psi_M(M^2)}(M^2+1)}(M^2+\psi_{\psi_{\psi_M(M^2)}(M^2+1)}(M^2)+1))%
}

\NovelDisplay{M02232R01}{%
(0)(1,1,1)(2,1,1)(3,1)(2,1)(3,1)(2,1)(3)%
}

\NovelDisplay{M02233R01}{%
= \psi(I+\psi_{\Omega_{\psi_I(I)+1}}(I+\psi_{\Omega_{\psi_I(I)+1}}(I)+\psi_{\Omega_{\psi_I(I)+1}}(1)))%
}

\NovelDisplay{M02234R01}{%
= \psi(M^2+\psi_{\psi_{\psi_M(M^2)}(M^2+1)}(M^2+\psi_{\psi_{\psi_M(M^2)}(M^2+1)}(M^2)+\psi_{\psi_{\psi_M(M^2)}(M^2+1)}(1)))%
}

\NovelDisplay{M02235R01}{%
(0)(1,1,1)(2,1,1)(3,1)(2,1)(3,1)(2,1)(3,1)(1,1,1)(2,1,1)(3,1)(2)%
}

\NovelDisplay{M02236R01}{%
= \psi(I+\psi_{\Omega_{\psi_I(I)+1}}(I+\psi_{\Omega_{\psi_I(I)+1}}(I)\times 2))%
}

\NovelDisplay{M02237R01}{%
= \psi(M^2+\psi_{\psi_{\psi_M(M^2)}(M^2+1)}(M^2+\psi_{\psi_{\psi_M(M^2)}(M^2+1)}(M^2)\times 2))%
}

\NovelDisplay{M02238R01}{%
(0)(1,1,1)(2,1,1)(3,1)(2,1)(3,1)(2,1)(3,1)(2,1)(3,1)(1,1,1)(2,1,1)%
}
\NovelDisplay{M02238R02}{%
(3,1)(2)%
}

\NovelDisplay{M02239R01}{%
= \psi(I+\psi_{\Omega_{\psi_I(I)+1}}(I+\psi_{\Omega_{\psi_I(I)+1}}(I)\times 3))%
}

\NovelDisplay{M02240R01}{%
= \psi(M^2+\psi_{\psi_{\psi_M(M^2)}(M^2+1)}(M^2+\psi_{\psi_{\psi_M(M^2)}(M^2+1)}(M^2)\times 3))%
}

\NovelDisplay{M02241R01}{%
(0)(1,1,1)(2,1,1)(3,1)(2,1)(3,1)(3)%
}

\NovelDisplay{M02242R01}{%
= \psi(I+\psi_{\Omega_{\psi_I(I)+1}}(I+\psi_{\Omega_{\psi_I(I)+1}}(I+1)))%
}

\NovelDisplay{M02243R01}{%
= \psi(M^2+\psi_{\psi_{\psi_M(M^2)}(M^2+1)}(M^2+\psi_{\psi_{\psi_M(M^2)}(M^2+1)}(M^2+1)))%
}

\NovelDisplay{M02244R01}{%
(0)(1,1,1)(2,1,1)(3,1)(2,1)(3,1)(3,1)(1,1,1)(2,1,1)(3,1)(2)%
}

\NovelDisplay{M02245R01}{%
= \psi(I+\psi_{\Omega_{\psi_I(I)+1}}(I+\psi_{\Omega_{\psi_I(I)+1}}(I+\psi_I(I))))%
}

\NovelDisplay{M02246R01}{%
= \psi(M^2+\psi_{\psi_{\psi_M(M^2)}(M^2+1)}(M^2+\psi_{\psi_{\psi_M(M^2)}(M^2+1)}(M^2+\psi_{\psi_M(M^2)}(M^2))))%
}

\NovelDisplay{M02247R01}{%
(0)(1,1,1)(2,1,1)(3,1)(2,1)(3,1)(4)%
}

\NovelDisplay{M02248R01}{%
= \psi(I+\psi_{\Omega_{\psi_I(I)+1}}(I+\psi_{\Omega_{\psi_I(I)+1}}(I+\psi_{\Omega_{\psi_I(I)+1}}(1))))%
}

\NovelDisplay{M02249R01}{%
= \psi(M^2+\psi_{\psi_{\psi_M(M^2)}(M^2+1)}(M^2+\psi_{\psi_{\psi_M(M^2)}(M^2+1)}(M^2+\psi_{\psi_{\psi_M(M^2)}(M^2+1)}(1))))%
}

唐三再来接应时，她们仍没把这组全写完。蔡须坤的肩背已经直不起来，卢本韦也抹着额角的汗，谁都没有说再撑一下就能过去。

她们把未完的一行夹好，留到第二天。那一晚，霍雨浩没有另出新题，只请薛峰把四人白天各自停住的位置记在空页上。第二天再看，最该补的是哪一段，终于不必从头猜。

蔡须坤在山脚收过一阵魂力，才重新把纸拿起来。

卢本韦问她，选M的时候是不是没想到会这样费力。她没有立即回答，翻回第一页，找出那条自己曾经连头尾都读不清的函数。

``那天灰气上来的时候，我只知道底下有东西。''她说，``现在遇到这些，我还是会慌。但至少能试着说清，哪一处是我的，哪一处是我刚才没有看明白的。''

她没有将混沌和M说成同一种力量。M只是纸上的表示，读清它不会替她平定井底，可一层层找回位置的过程，让她不再急着凭更大的力气挤过不懂的地方。

卢本韦听完，把自己的I册放下，陪她将今天最后一处返回位置重新接过。两人都困，念错了便停，确认以后再往下。直到古月娜推门进来，她们才发现灯芯已经短了一截，桌边的水也凉了。

等Ω重新出现在外面时，卢本韦险些以为自己终于走出了那片复杂的下标。

她把式子读到Ω便停住，脚下也顺势向上探了一寸。那一寸没有落稳，膝盖先弯了下去。蔡须坤从旁边抓住她的袖子，古月娜托住另一边，三人一齐退回原处。

纸没有错，是她没有读完。

``它带着谁的下标？''霍雨浩问。

卢本韦喘匀以后，低头把那一段重新说全。下标里旧有的函数值还在，后面的有限增量也还在。Ω露在外面，只是这一组对应的层次换了写法，不能因此把它底下托着的东西抹掉。

蔡须坤反过来拿M侧帮了她一次。这里的变化可以在同一个较深函数的参数中看见，她指着其中尚能清楚比较的两处，让卢本韦先看有限增量落在哪儿，再回I侧找那一处。

``我刚才只看见了熟悉的那个Ω。''卢本韦低声说。

``我刚才只看见了太多ψ。''蔡须坤应道。

两人对看了一眼，都没笑出来，倒是把各自最容易躲过去的那一栏，重新翻到朝着自己的方向。

这一段中也有把整项取两次的情况。霍雨浩让她们把乘二所在的层次读出来：它究竟留在最外层ψ的参数里，还是被包进了某一个下标，必须按眼前的括号判断。有限增量、整项重复与更换函数上界，写出来可能都只差很短一截，接续时却不能混成一种动作。

薛峰把这些不同的位置分在纸的两边，终于第一次不用等霍雨浩问，就指出卢本韦漏念的一层。

这一次，薛峰先核了左列最末的一处展开。末列是三、二，坏根在前面的二、一；最高非零行在第二行，所以不能将第二行也跟着加上阶差。按这一条原式，新增段所用的第一行增量是一，下面的坐标保持原值。

他把同一个原式的第零项、第一项、第二项分别写出。第一项与第二项各自从原式求来，第二项不靠将第一项再展开一次取得。

\NovelDisplay{M02250R01}{%
\text{展开第零项：}\quad(0)(1,1,1)(2,1,1)(3,1)(2,1)%
}

\NovelDisplay{M02251R01}{%
\text{展开第一项：}\quad(0)(1,1,1)(2,1,1)(3,1)(2,1)(3,1)%
}

\NovelDisplay{M02252R01}{%
\text{展开第二项：}\quad(0)(1,1,1)(2,1,1)(3,1)(2,1)(3,1)(4,1)%
}

卢本韦逐列看过，蔡须坤再看新增段的第二行有没有误升。更长的有限项仍低于原式，不能用列数多少取代BMS本身的比较。

这份复算能够确认左列的展开，M侧与I侧的完整对应仍要分别读清。薛峰将两类检查分开，才把纸交还给她们。

\NovelDisplay{M02253R01}{%
(0)(1,1,1)(2,1,1)(3,1)(2,1)(3,2)%
}

\NovelDisplay{M02254R01}{%
= \psi(I+\Omega_{\psi_I(I)+1})%
}

\NovelDisplay{M02255R01}{%
= \psi(M^2+\psi_{\psi_M(M^2)}(M^2+1))%
}

\NovelDisplay{M02256R01}{%
(0)(1,1,1)(2,1,1)(3,1)(2,1)(3,2)(1,1,1)(2,1,1)(3,1,1)(2)%
}

\NovelDisplay{M02257R01}{%
= \psi(I+\Omega_{\psi_I(I)+1}+\psi_I(I))%
}

\NovelDisplay{M02258R01}{%
= \psi(M^2+\psi_{\psi_M(M^2)}(M^2+1)+\psi_{\psi_M(M^2)}(M^2))%
}

\NovelDisplay{M02259R01}{%
(0)(1,1,1)(2,1,1)(3,1)(2,1)(3,2)(1,1,1)(2,1,1)(3,1,1)(2,1)(3,2)%
}

\NovelDisplay{M02260R01}{%
= \psi(I+\Omega_{\psi_I(I)+1}+\psi_{\Omega_{\psi_I(I)+1}}(I+\Omega_{\psi_I(I)+1}))%
}

\NovelDisplay{M02261R01}{%
= \psi(M^2+\psi_{\psi_M(M^2)}(M^2+1)+\psi_{\psi_{\psi_M(M^2)}(M^2+1)}(M^2+\psi_{\psi_M(M^2)}(M^2+1)))%
}

\NovelDisplay{M02262R01}{%
(0)(1,1,1)(2,1,1)(3,1)(2,1)(3,2)(2)%
}

\NovelDisplay{M02263R01}{%
= \psi(I+\Omega_{\psi_I(I)+1}+\psi_{\Omega_{\psi_I(I)+1}}(I+\Omega_{\psi_I(I)+1}+1))%
}

\NovelDisplay{M02264R01}{%
= \psi(M^2+\psi_{\psi_M(M^2)}(M^2+1)+\psi_{\psi_{\psi_M(M^2)}(M^2+1)}(M^2+\psi_{\psi_M(M^2)}(M^2+1)+1))%
}

\NovelDisplay{M02265R01}{%
(0)(1,1,1)(2,1,1)(3,1)(2,1)(3,2)(2,1)(1,1,1)(2,1,1)(3,1)(2)%
}

\NovelDisplay{M02266R01}{%
= \psi(I+\Omega_{\psi_I(I)+1}+\psi_{\Omega_{\psi_I(I)+1}}(I+\Omega_{\psi_I(I)+1}+\psi_I(I)))%
}

\NovelDisplay{M02267R01}{%
= \psi(M^2+\psi_{\psi_M(M^2)}(M^2+1)+\psi_{\psi_{\psi_M(M^2)}(M^2+1)}(M^2+\psi_{\psi_M(M^2)}(M^2+1)%
}
\NovelDisplay{M02267R02}{%
+\psi_{\psi_M(M^2)}(M^2)))%
}

\NovelDisplay{M02268R01}{%
(0)(1,1,1)(2,1,1)(3,1)(2,1)(3,2)(2,1)(3,2)%
}

\NovelDisplay{M02269R01}{%
= \psi(I+\Omega_{\psi_I(I)+1}+\psi_{\Omega_{\psi_I(I)+1}}(I+\Omega_{\psi_I(I)+1}+\psi_{\Omega_{\psi_I(I)+1}}(I+\Omega_{\psi_I(I)+1})))%
}

\NovelDisplay{M02270R01}{%
= \psi(M^2+\psi_{\psi_M(M^2)}(M^2+1)+\psi_{\psi_{\psi_M(M^2)}(M^2+1)}(M^2+\psi_{\psi_M(M^2)}(M^2+1)%
}
\NovelDisplay{M02270R02}{%
+\psi_{\psi_{\psi_M(M^2)}(M^2+1)}(M^2+\psi_{\psi_M(M^2)}(M^2+1))))%
}

\NovelDisplay{M02271R01}{%
(0)(1,1,1)(2,1,1)(3,1)(2,1)(3,2)(3)%
}

\NovelDisplay{M02272R01}{%
= \psi(I+\Omega_{\psi_I(I)+1}+\psi_{\Omega_{\psi_I(I)+1}}(I+\Omega_{\psi_I(I)+1}+\psi_{\Omega_{\psi_I(I)+1}}(I+\Omega_{\psi_I(I)+1}+1)))%
}

\NovelDisplay{M02273R01}{%
= \psi(M^2+\psi_{\psi_M(M^2)}(M^2+1)+\psi_{\psi_{\psi_M(M^2)}(M^2+1)}(M^2+\psi_{\psi_M(M^2)}(M^2+1)%
}
\NovelDisplay{M02273R02}{%
+\psi_{\psi_{\psi_M(M^2)}(M^2+1)}(M^2+\psi_{\psi_M(M^2)}(M^2+1)+1)))%
}

\NovelDisplay{M02274R01}{%
(0)(1,1,1)(2,1,1)(3,1)(2,1)(3,2)(3,1)(4,2)%
}

\NovelDisplay{M02275R01}{%
= \psi(I+\Omega_{\psi_I(I)+1}+\psi_{\Omega_{\psi_I(I)+1}}(I+\Omega_{\psi_I(I)+1}+\psi_{\Omega_{\psi_I(I)+1}}(I+\Omega_{\psi_I(I)+1}%
}
\NovelDisplay{M02275R02}{%
+\psi_{\Omega_{\psi_I(I)+1}}(I+\Omega_{\psi_I(I)+1}))))%
}

\NovelDisplay{M02276R01}{%
= \psi(M^2+\psi_{\psi_M(M^2)}(M^2+1)+\psi_{\psi_{\psi_M(M^2)}(M^2+1)}(M^2+\psi_{\psi_M(M^2)}(M^2+1)%
}
\NovelDisplay{M02276R02}{%
+\psi_{\psi_{\psi_M(M^2)}(M^2+1)}(M^2+\psi_{\psi_M(M^2)}(M^2+1)+\psi_{\psi_{\psi_M(M^2)}(M^2+1)}(M^2%
}
\NovelDisplay{M02276R03}{%
+\psi_{\psi_M(M^2)}(M^2+1)))))%
}

\NovelDisplay{M02277R01}{%
(0)(1,1,1)(2,1,1)(3,1)(2,1)(3,2)(3,2)%
}

\NovelDisplay{M02278R01}{%
= \psi(I+\Omega_{\psi_I(I)+1}\times 2)%
}

\NovelDisplay{M02279R01}{%
= \psi(M^2+\psi_{\psi_M(M^2)}(M^2+1)\times 2)%
}

\NovelDisplay{M02280R01}{%
(0)(1,1,1)(2,1,1)(3,1)(2,1)(3,2)(4,1)%
}

\NovelDisplay{M02281R01}{%
= \psi(I+\psi_{\Omega_{\psi_I(I)+2}}(\Omega) )%
}

\NovelDisplay{M02282R01}{%
= \psi(M^2+\psi_{\psi_{\psi_M(M^2)}(M^2+2)}(M))%
}

\NovelDisplay{M02283R01}{%
(0)(1,1,1)(2,1,1)(3,1)(2,1)(3,2)(4,1)(1,1,1)(2,1,1)(3,1)(2)%
}

\NovelDisplay{M02284R01}{%
= \psi(I+\psi_{\Omega_{\psi_I(I)+2}}(I) )%
}

\NovelDisplay{M02285R01}{%
= \psi(M^2+\psi_{\psi_{\psi_M(M^2)}(M^2+2)}(M^2))%
}

\NovelDisplay{M02286R01}{%
(0)(1,1,1)(2,1,1)(3,1)(2,1)(3,2)(4,1)(2)%
}

\NovelDisplay{M02287R01}{%
= \psi(I+\psi_{\Omega_{\psi_I(I)+2}}(I)+\psi_{\Omega_{\psi_I(I)+1}}(I+\psi_{\Omega_{\psi_I(I)+2}}(I)+1))%
}

\NovelDisplay{M02288R01}{%
= \psi(M^2+\psi_{\psi_{\psi_M(M^2)}(M^2+2)}(M^2)+\psi_{\psi_{\psi_M(M^2)}(M^2+1)}(M^2%
}
\NovelDisplay{M02288R02}{%
+\psi_{\psi_{\psi_M(M^2)}(M^2+2)}(M^2)+1))%
}

\NovelDisplay{M02289R01}{%
(0)(1,1,1)(2,1,1)(3,1)(2,1)(3,2)(4,1)(2,1)(3,2)%
}

\NovelDisplay{M02290R01}{%
= \psi(I+\psi_{\Omega_{\psi_I(I)+2}}(I)+\psi_{\Omega_{\psi_I(I)+1}}(I+\psi_{\Omega_{\psi_I(I)+2}}(I)+\psi_{\Omega_{\psi_I(I)+1}}(I+\Omega_{\psi_I(I)+1})))%
}

\NovelDisplay{M02291R01}{%
= \psi(M^2+\psi_{\psi_{\psi_M(M^2)}(M^2+2)}(M^2)+\psi_{\psi_{\psi_M(M^2)}(M^2+1)}(M^2+\psi_{\psi_{\psi_M(M^2)}(M^2+2)}(M^2)%
}
\NovelDisplay{M02291R02}{%
+\psi_{\psi_{\psi_M(M^2)}(M^2+1)}(M^2+\psi_{\psi_M(M^2)}(M^2+1))))%
}

\NovelDisplay{M02292R01}{%
(0)(1,1,1)(2,1,1)(3,1)(2,1)(3,2)(4,1)(2,1)(3,2)(4,1)%
}

\NovelDisplay{M02293R01}{%
= \psi(I+\psi_{\Omega_{\psi_I(I)+2}}(I)+\psi_{\Omega_{\psi_I(I)+1}}(I+\psi_{\Omega_{\psi_I(I)+2}}(I)+\psi_{\Omega_{\psi_I(I)+1}}(I+\psi_{\Omega_{\psi_I(I)+2}}(\Omega))))%
}

\NovelDisplay{M02294R01}{%
= \psi(M^2+\psi_{\psi_{\psi_M(M^2)}(M^2+2)}(M^2)+\psi_{\psi_{\psi_M(M^2)}(M^2+1)}(M^2+\psi_{\psi_{\psi_M(M^2)}(M^2+2)}(M^2)%
}
\NovelDisplay{M02294R02}{%
+\psi_{\psi_{\psi_M(M^2)}(M^2+1)}(M^2+\psi_{\psi_{\psi_M(M^2)}(M^2+2)}(M))))%
}

\NovelDisplay{M02295R01}{%
(0)(1,1,1)(2,1,1)(3,1)(2,1)(3,2)(4,1)(3,1)%
}

\NovelDisplay{M02296R01}{%
=\! \psi(I\!+\!\psi_{\Omega_{\psi_I(I)+2}}(I)\!+\!\psi_{\Omega_{\psi_I(I)+1}}(I\!+\!\psi_{\Omega_{\psi_I(I)+2}}(I)\!+\!\psi_{\Omega_{\psi_I(I)+1}}(I\!+\!\psi_{\Omega_{\psi_I(I)+2}}(I)\!+\!\Omega)))%
}

\NovelDisplay{M02297R01}{%
= \psi(M^2+\psi_{\psi_{\psi_M(M^2)}(M^2+2)}(M^2)+\psi_{\psi_{\psi_M(M^2)}(M^2+1)}(M^2+\psi_{\psi_{\psi_M(M^2)}(M^2+2)}(M^2)%
}
\NovelDisplay{M02297R02}{%
+\psi_{\psi_{\psi_M(M^2)}(M^2+1)}(M^2+\psi_{\psi_{\psi_M(M^2)}(M^2+2)}(M^2)+\psi_{\psi_M(M^2)}(1))))%
}

\NovelDisplay{M02298R01}{%
(0)(1,1,1)(2,1,1)(3,1)(2,1)(3,2)(4,1)(3,1)(4,2)(5,1)%
}

\NovelDisplay{M02299R01}{%
= \psi(I+\psi_{\Omega_{\psi_I(I)+2}}(I)+\psi_{\Omega_{\psi_I(I)+1}}(I+\psi_{\Omega_{\psi_I(I)+2}}(I)+\psi_{\Omega_{\psi_I(I)+1}}(I+\psi_{\Omega_{\psi_I(I)+2}}(I)%
}
\NovelDisplay{M02299R02}{%
+\psi_{\Omega_{\psi_I(I)+1}}(I+\psi_{\Omega_{\psi_I(I)+2}}(\Omega)))))%
}

\NovelDisplay{M02300R01}{%
= \psi(M^2+\psi_{\psi_{\psi_M(M^2)}(M^2+2)}(M^2)+\psi_{\psi_{\psi_M(M^2)}(M^2+1)}(M^2+\psi_{\psi_{\psi_M(M^2)}(M^2+2)}(M^2)%
}
\NovelDisplay{M02300R02}{%
+\psi_{\psi_{\psi_M(M^2)}(M^2+1)}(M^2+\psi_{\psi_{\psi_M(M^2)}(M^2+2)}(M^2)+\psi_{\psi_{\psi_M(M^2)}(M^2+1)}(M^2%
}
\NovelDisplay{M02300R03}{%
+\psi_{\psi_{\psi_M(M^2)}(M^2+2)}(M)))))%
}

\NovelDisplay{M02301R01}{%
(0)(1,1,1)(2,1,1)(3,1)(2,1)(3,2)(4,1)(3,2)%
}

\NovelDisplay{M02302R01}{%
= \psi(I+\psi_{\Omega_{\psi_I(I)+2}}(I)+\Omega_{\psi_I(I)+1})%
}

\NovelDisplay{M02303R01}{%
= \psi(M^2+\psi_{\psi_{\psi_M(M^2)}(M^2+2)}(M^2)+\psi_{\psi_M(M^2)}(M^2+1))%
}

傍晚，卢本韦将那只被自己扯皱的袖口抚平，还给蔡须坤。

``刚才拉得挺紧。''

``不然让你滚下去？''

卢本韦想起自己在山脚说过的那些话，没再嘴硬。她把第二天要核的纸接过来，先检查蔡须坤那一侧有没有被风打湿。

越往后，I侧越容易给人一种已经摸到规律的错觉。加一、加二，再到加三，几个字足以说完，M侧却仍要把对应的函数头带在前面。

蔡须坤这次没急着羡慕。她已经看见，自己若把同一个函数头漏掉，后面的有限增量就会像失去依托，仿佛谁都能拿来接在别处。

卢本韦先把每一处共同的下标完整念过，再单独说出新添的部分。薛峰核左列，霍雨浩核她有没有将前面学过的结果误装到新的位置，古月娜请唐三、小舞也从最短的一行开始复述。

诺当原本看着长页就发愁。他能记住前面几列，念到换行处却总想重新补一个零。薛峰把上下两行的接缝对齐，拿一根细枝沿整条左列慢慢划过去。那根细枝没有在换行处停，诺当也跟着把后半段接完了。

``只是纸折过去了。''薛峰说。

诺当点头，重新写过一次，终于没多出开头那一列。

蔡须坤听见这句话时，忽然也把自己右栏最深的一处接上了。她先前总在视觉上回到新行，意识便想跟着退到最外层。可断行只替纸省地方，没有替她结束任何一个括号。

她又从头读了一遍，这次没有靠猜。到最后一处有限增量时，卢本韦也接得很稳。两人将各自的纸交换，连着几天积在一处的困难，终于露出一条能逐段走完的路。

夏夜星空没有把下一组的答案递过来，只让十一人分别说明自己现在能站稳的这一阶。霍雨浩站在最外侧，听到谁停住，便留在那儿等谁重新说完。

\NovelDisplay{M02304R01}{%
(0)(1,1,1)(2,1,1)(3,1)(2,1)(3,2)(4,1)(3,2)(4,1)%
}

\NovelDisplay{M02305R01}{%
= \psi(I+\psi_{\Omega_{\psi_I(I)+2}}(I)+\psi_{\Omega_{\psi_I(I)+2}}(\Omega))%
}

\NovelDisplay{M02306R01}{%
= \psi(M^2+\psi_{\psi_{\psi_M(M^2)}(M^2+2)}(M^2)+\psi_{\psi_{\psi_M(M^2)}(M^2+2)}(M))%
}

\NovelDisplay{M02307R01}{%
(0)(1,1,1)(2,1,1)(3,1)(2,1)(3,2)(4,1)(4)%
}

\NovelDisplay{M02308R01}{%
= \psi(I+\psi_{\Omega_{\psi_I(I)+2}}(I+1))%
}

\NovelDisplay{M02309R01}{%
= \psi(M^2+\psi_{\psi_{\psi_M(M^2)}(M^2+2)}(M^2+1))%
}

\NovelDisplay{M02310R01}{%
(0)(1,1,1)(2,1,1)(3,1)(2,1)(3,2)(4,1)(4,1)%
}

\NovelDisplay{M02311R01}{%
= \psi(I+\psi_{\Omega_{\psi_I(I)+2}}(I+\Omega))%
}

\NovelDisplay{M02312R01}{%
= \psi(M^2+\psi_{\psi_{\psi_M(M^2)}(M^2+2)}(M^2+\psi_{\psi_M(M^2)}(1)))%
}

\NovelDisplay{M02313R01}{%
(0)(1,1,1)(2,1,1)(3,1)(2,1)(3,2)(4,1)(4,1)(4,1)%
}

\NovelDisplay{M02314R01}{%
= \psi(I+\psi_{\Omega_{\psi_I(I)+2}}(I+\psi_I(I)+\Omega))%
}

\NovelDisplay{M02315R01}{%
= \psi(M^2+\psi_{\psi_{\psi_M(M^2)}(M^2+2)}(M^2+\psi_{\psi_M(M^2)}(M^2)+\psi_{\psi_M(M^2)}(1)))%
}

\NovelDisplay{M02316R01}{%
(0)(1,1,1)(2,1,1)(3,1)(2,1)(3,2)(4,1)(5)%
}

\NovelDisplay{M02317R01}{%
= \psi(I+\psi_{\Omega_{\psi_I(I)+2}}(I+\psi_{\Omega_{\psi_I(I)+1}}(1)))%
}

\NovelDisplay{M02318R01}{%
= \psi(M^2+\psi_{\psi_{\psi_M(M^2)}(M^2+2)}(M^2+\psi_{\psi_{\psi_M(M^2)}(M^2+1)}(1)))%
}

\NovelDisplay{M02319R01}{%
(0)(1,1,1)(2,1,1)(3,1)(2,1)(3,2)(4,1)(5,2)%
}

\NovelDisplay{M02320R01}{%
= \psi(I+\psi_{\Omega_{\psi_I(I)+2}}(I+\psi_{\Omega_{\psi_I(I)+1}}(I+\Omega_{\psi_I(I)+1})))%
}

\NovelDisplay{M02321R01}{%
= \psi(M^2+\psi_{\psi_{\psi_M(M^2)}(M^2+2)}(M^2+\psi_{\psi_{\psi_M(M^2)}(M^2+1)}(M^2+\psi_{\psi_M(M^2)}(M^2+1))))%
}

\NovelDisplay{M02322R01}{%
(0)(1,1,1)(2,1,1)(3,1)(2,1)(3,2)(4,2)%
}

\NovelDisplay{M02323R01}{%
= \psi(I+\psi_{\Omega_{\psi_I(I)+2}}(I+\Omega_{\psi_I(I)+1}))%
}

\NovelDisplay{M02324R01}{%
= \psi(M^2+\psi_{\psi_{\psi_M(M^2)}(M^2+2)}(M^2+\psi_{\psi_M(M^2)}(M^2+1)))%
}

\NovelDisplay{M02325R01}{%
(0)(1,1,1)(2,1,1)(3,1)(2,1)(3,2)(4,2)(4,2)%
}

\NovelDisplay{M02326R01}{%
= \psi(I+\psi_{\Omega_{\psi_I(I)+2}}(I+\Omega_{\psi_I(I)+1}\times 2))%
}

\NovelDisplay{M02327R01}{%
= \psi(M^2+\psi_{\psi_{\psi_M(M^2)}(M^2+2)}(M^2+\psi_{\psi_M(M^2)}(M^2+1)\times 2))%
}

\NovelDisplay{M02328R01}{%
(0)(1,1,1)(2,1,1)(3,1)(2,1)(3,2)(4,2)(5,2)%
}

\NovelDisplay{M02329R01}{%
= \psi(I+\psi_{\Omega_{\psi_I(I)+2}}(I+\psi_{\Omega_{\psi_I(I)+2}}(I+\Omega_{\psi_I(I)+1})))%
}

\NovelDisplay{M02330R01}{%
= \psi(M^2+\psi_{\psi_{\psi_M(M^2)}(M^2+2)}(M^2+\psi_{\psi_{\psi_M(M^2)}(M^2+2)}(M^2+\psi_{\psi_M(M^2)}(M^2+1))))%
}

\NovelDisplay{M02331R01}{%
(0)(1,1,1)(2,1,1)(3,1)(2,1)(3,2)(4,3)%
}

\NovelDisplay{M02332R01}{%
= \psi(I+\Omega_{\psi_I(I)+2})%
}

\NovelDisplay{M02333R01}{%
= \psi(M^2+\psi_{\psi_M(M^2)}(M^2+2))%
}

\NovelDisplay{M02334R01}{%
(0)(1,1,1)(2,1,1)(3,1)(2,1)(3,2)(4,3)(2)%
}

\NovelDisplay{M02335R01}{%
= \psi(I+\Omega_{\psi_I(I)+2}+\psi_{\Omega_{\psi_I(I)+1}}(I+\Omega_{\psi_I(I)+2}+1))%
}

\NovelDisplay{M02336R01}{%
= \psi(M^2+\psi_{\psi_M(M^2)}(M^2+2)+\psi_{\psi_{\psi_M(M^2)}(M^2+1)}(M^2+\psi_{\psi_M(M^2)}(M^2+2)+1))%
}

\NovelDisplay{M02337R01}{%
(0)(1,1,1)(2,1,1)(3,1)(2,1)(3,2)(4,3)(3,2)%
}

\NovelDisplay{M02338R01}{%
= \psi(I+\Omega_{\psi_I(I)+2}+\Omega_{\psi_I(I)+1})%
}

\NovelDisplay{M02339R01}{%
= \psi(M^2+\psi_{\psi_M(M^2)}(M^2+2)+\psi_{\psi_M(M^2)}(M^2+1))%
}

\NovelDisplay{M02340R01}{%
(0)(1,1,1)(2,1,1)(3,1)(2,1)(3,2)(4,3)(3,2)(4,3)%
}

\NovelDisplay{M02341R01}{%
= \psi(I+\Omega_{\psi_I(I)+2}+\psi_{\Omega_{\psi_I(I)+2}}(I+\Omega_{\psi_I(I)+2}))%
}

\NovelDisplay{M02342R01}{%
= \psi(M^2+\psi_{\psi_M(M^2)}(M^2+2)+\psi_{\psi_{\psi_M(M^2)}(M^2+2)}(M^2+\psi_{\psi_M(M^2)}(M^2+2)))%
}

\NovelDisplay{M02343R01}{%
(0)(1,1,1)(2,1,1)(3,1)(2,1)(3,2)(4,3)(4,3)%
}

\NovelDisplay{M02344R01}{%
= \psi(I+\Omega_{\psi_I(I)+2}\times 2)%
}

\NovelDisplay{M02345R01}{%
= \psi(M^2+\psi_{\psi_M(M^2)}(M^2+2) \times 2)%
}

\NovelDisplay{M02346R01}{%
(0)(1,1,1)(2,1,1)(3,1)(2,1)(3,2)(4,3)(5)%
}

\NovelDisplay{M02347R01}{%
= \psi(I+\psi_{\Omega_{\psi_I(I)+3}}(1))%
}

\NovelDisplay{M02348R01}{%
= \psi(M^2+\psi_{\psi_{\psi_M(M^2)}(M^2+3)}(1))%
}

\NovelDisplay{M02349R01}{%
(0)(1,1,1)(2,1,1)(3,1)(2,1)(3,2)(4,3)(5,4)%
}

\NovelDisplay{M02350R01}{%
= \psi(I+\Omega_{\psi_I(I)+3})%
}

\NovelDisplay{M02351R01}{%
= \psi(M^2+\psi_{\psi_M(M^2)}(M^2+3))%
}

他们真正踏上二百四十米时，已不是第一次来试这几级石阶。

蔡须坤仍旧走得慢，却能在一处接续不畅之前，先认出是哪一层出了问题。卢本韦的I侧写得短，也终于不再把短当成容易。薛峰背着两人的旧页站到旁边，念出了最容易在断行处抄错的那一段，连他自己都怔了一下。

横台上能坐的位置不多。唐三先将水袋放下，等后面的人都到了，才去卸自己的包。蔡兴坤把黑子鸡搁在膝上，小舞让诺当将刚才那根细枝留下，夹进了本子里。

上方的石阶仍然没进雾里。蔡须坤翻到下一页，看见有限的加三之后，出现了一点更熟悉、也更难解释的ω。

她没有往前试，只把纸递给卢本韦。

``明天一起看这个。''

横台上的问答没有按谁走得快来结束。蔡须坤听见后面有人卡住，原本已收起的M册又被她展开。卢本韦则退到小舞身边，让她把当前这条最外面的组织先说清。

她们开始能够给出很小、却真正有用的帮助：指出一对漏认的括号，接住一处断行，或是等对方把原来的那口气重新喘匀。没人因此少学一条，后面的问话却终于不再只剩``这页我全看不懂''。

那晚回到山脚，蔡须坤把M侧摊在灯下，先将今天真能说清的部分勾好。卢本韦坐在另一边，将I侧逐行压平，看到自己险些跌下去的那处，伸手摸了一下膝盖。

霍雨浩没有催她们再读。薛峰却还留在桌前，把白天替同伴核过的那几行遮住，自己又写了一次。写错的地方很少，重写得最慢的，恰是他先前帮诺当接好的断行。

蔡须坤看见，便将灯往他那边挪了挪。

``明天你也坐前面。''她说。

薛峰抬起头，过了一会儿才答好。

二百四十米只比旧横台高二十米。这个数字写在册上很短，翻过去，纸却厚了不少。最后一行核完，古月娜推来的那杯热水已经凉了。蔡须坤喝了一口，才将本子收进包里。

\chapter{第一百一十四章：下标深处}\label{ux7b2cux4e00ux767eux4e00ux5341ux56dbux7ae0ux4e0bux6807ux6df1ux5904}

二百四十米以上的第一段，蔡须坤连续三天没有走过。

最难的一次，她已经把右侧那条M式读完，脚抬起来，才发觉自己读出的最后一层，其实仍在前一层的参数里面。那只脚在半空停了很久，最终原样收回。

卢本韦没有催。她自己的I侧也卡住了：有限的一、二、三都能说清，ω落进去以后，那条原先熟悉的下标，忽然显得比纸面长得多。

薛峰将她们用过的纸分成三摞。一摞已经能独立复算，一摞只在听别人解释时懂，最后一摞则是谁也不肯再凭印象抄下去的。纸摞分完，最厚的是中间那一摞。

``原来我会的是听你们说。''他低声道。

霍雨浩把炭条递给他。``那就从你能自己接的一行开始。''

上方的雾里没有露出新横台。风声贴着石壁转回来，十一人都在那几级台阶间找能歇脚的位置。唐三把绳子重新系过，留出够人退回的余地。古月娜站在蔡须坤身后，等她将呼吸稳住，才一起翻开新页。

ω进入的，是Ω下标内部的一处，并没有让最外面的ψ换一个位置。

卢本韦先看清了这一点，后面的念法却又快起来。她把熟悉的旧项略过，直接说加ω，霍雨浩便将后面两个字遮住，让她把前面也说一遍。

``先把它的来处接回来。''

她重新念完，终于听见自己省略的那一段。这里是在旧的函数值之后接ω；下一行再添有限的一，与把ω换成有限的一，是两种不同的写法。她把有限尾部留在原处，没有因为看见ω便将整条下标读成一个大概的极限。

蔡须坤在M侧看见的，是原有函数头仍在，参数里那段有限增长走到了ω。她原想用几个字说明白，读到第二行便停住：同一个参数里还有前面保留的M平方，它没有因为新的ω出现而退走。

两人开始交替读。卢本韦每读出I侧的一处，蔡须坤就指出M侧相应变化落在哪个括号里。霍雨浩只在两人说到不同层次时抬手，薛峰则在页边画短线，记下她们必须共同停顿的位置。

第一遍，短线画得很密。第二遍，卢本韦等蔡须坤接完函数头再往下说，线便少了一半。

她们仍然走不过那一级。膝盖在压力下发抖，背后的绳子却始终松着，谁都没有被人向上拖。夏夜星空等她们退回来，才让小舞也读一遍最短的那处。小舞停在和蔡须坤相同的位置，蔡须坤便坐下来，把自己刚刚找清的括号再指给她看。

\NovelDisplay{M02352R01}{%
(0)(1,1,1)(2,1,1)(3,1)(2,1)(3,2,1)%
}

\NovelDisplay{M02353R01}{%
= \psi(I+\Omega_{\psi_I(I)+\omega})%
}

\NovelDisplay{M02354R01}{%
= \psi(M^2+\psi_{\psi_M(M^2)}(M^2+\omega))%
}

\NovelDisplay{M02355R01}{%
(0)(1,1,1)(2,1,1)(3,1)(2,1)(3,2,1)(3,2)%
}

\NovelDisplay{M02356R01}{%
= \psi(I+\Omega_{\psi_I(I)+\omega}+\Omega_{\psi_I(I)+1})%
}

\NovelDisplay{M02357R01}{%
= \psi(M^2+\psi_{\psi_M(M^2)}(M^2+\omega)+\psi_{\psi_M(M^2)}(M^2+1))%
}

\NovelDisplay{M02358R01}{%
(0)(1,1,1)(2,1,1)(3,1)(2,1)(3,2,1)(3,2)(4,3,1)(4,3)%
}

\NovelDisplay{M02359R01}{%
= \psi(I+\Omega_{\psi_I(I)+\omega}+\Omega_{\psi_I(I)+2})%
}

\NovelDisplay{M02360R01}{%
= \psi(M^2+\psi_{\psi_M(M^2)}(M^2+\omega)+\psi_{\psi_M(M^2)}(M^2+2))%
}

\NovelDisplay{M02361R01}{%
(0)(1,1,1)(2,1,1)(3,1)(2,1)(3,2,1)(3,2,1)%
}

\NovelDisplay{M02362R01}{%
= \psi(I+\Omega_{\psi_I(I)+\omega}\times 2)%
}

\NovelDisplay{M02363R01}{%
= \psi(M^2+\psi_{\psi_M(M^2)}(M^2+\omega)\times 2)%
}

\NovelDisplay{M02364R01}{%
(0)(1,1,1)(2,1,1)(3,1)(2,1)(3,2,1)(4,1)%
}

\NovelDisplay{M02365R01}{%
= \psi(I+\psi_{\Omega_{\psi_I(I)+\omega+1}}(\Omega))%
}

\NovelDisplay{M02366R01}{%
= \psi(M^2+\psi_{\psi_{\psi_M(M^2)}(M^2+\omega+1)}(M))%
}

\NovelDisplay{M02367R01}{%
(0)(1,1,1)(2,1,1)(3,1)(2,1)(3,2,1)(4,1)(1,1,1)(2,1,1)(3,1)(2)%
}

\NovelDisplay{M02368R01}{%
= \psi(I+\psi_{\Omega_{\psi_I(I)+\omega+1}}(I))%
}

\NovelDisplay{M02369R01}{%
= \psi(M^2+\psi_{\psi_{\psi_M(M^2)}(M^2+\omega+1)}(M^2))%
}

\NovelDisplay{M02370R01}{%
(0)(1,1,1)(2,1,1)(3,1)(2,1)(3,2,1)(4,1)(3,2,1)%
}

\NovelDisplay{M02371R01}{%
= \psi(I+\psi_{\Omega_{\psi_I(I)+\omega+1}}(I)+\Omega_{\psi_I(I)+1})%
}

\NovelDisplay{M02372R01}{%
= \psi(M^2+\psi_{\psi_{\psi_M(M^2)}(M^2+\omega+1)}(M^2)+\psi_{\psi_M(M^2)}(M^2+\omega))%
}

\NovelDisplay{M02373R01}{%
(0)(1,1,1)(2,1,1)(3,1)(2,1)(3,2,1)(4,1)(3,2,1)(4,1)%
}

\NovelDisplay{M02374R01}{%
= \psi(I+\psi_{\Omega_{\psi_I(I)+\omega+1}}(I)+\psi_{\Omega_{\psi_I(I)+\omega+1}}(\Omega))%
}

\NovelDisplay{M02375R01}{%
= \psi(M^2+\psi_{\psi_{\psi_M(M^2)}(M^2+\omega+1)}(M^2)+\psi_{\psi_{\psi_M(M^2)}(M^2+\omega+1)}(M))%
}

\NovelDisplay{M02376R01}{%
(0)(1,1,1)(2,1,1)(3,1)(2,1)(3,2,1)(4,1)(4,1)%
}

\NovelDisplay{M02377R01}{%
= \psi(I+\psi_{\Omega_{\psi_I(I)+\omega+1}}(I+\Omega))%
}

\NovelDisplay{M02378R01}{%
= \psi(M^2+\psi_{\psi_{\psi_M(M^2)}(M^2+\omega+1)}(M^2+\psi_{\psi_M(M^2)}(1)))%
}

\NovelDisplay{M02379R01}{%
(0)(1,1,1)(2,1,1)(3,1)(2,1)(3,2,1)(4,1)(4,1)(3,2,1)(4,1)(4,1)%
}

\NovelDisplay{M02380R01}{%
= \psi(I+\psi_{\Omega_{\psi_I(I)+\omega+1}}(I+\psi_I(I))+\psi_{\Omega_{\psi_I(I)+\omega+1}}(I+\Omega))%
}

\NovelDisplay{M02381R01}{%
= \psi(M^2+\psi_{\psi_{\psi_M(M^2)}(M^2+\omega+1)}(M^2+\psi_{\psi_M(M^2)}(M^2))%
}
\NovelDisplay{M02381R02}{%
+\psi_{\psi_{\psi_M(M^2)}(M^2+\omega+1)}(M^2+\psi_{\psi_M(M^2)}(1)))%
}

\NovelDisplay{M02382R01}{%
(0)(1,1,1)(2,1,1)(3,1)(2,1)(3,2,1)(4,1)(4,1)(4)%
}

\NovelDisplay{M02383R01}{%
= \psi(I+\psi_{\Omega_{\psi_I(I)+\omega+1}}(I+\psi_I(I)+1))%
}

\NovelDisplay{M02384R01}{%
= \psi(M^2+\psi_{\psi_{\psi_M(M^2)}(M^2+\omega+1)}(M^2+\psi_{\psi_M(M^2)}(M^2)+1))%
}

\NovelDisplay{M02385R01}{%
(0)(1,1,1)(2,1,1)(3,1)(2,1)(3,2,1)(4,1)(5)%
}

\NovelDisplay{M02386R01}{%
= \psi(I+\psi_{\Omega_{\psi_I(I)+\omega+1}}(I+\psi_{\Omega_{\psi_I(I)+1}}(1)))%
}

\NovelDisplay{M02387R01}{%
= \psi(M^2+\psi_{\psi_{\psi_M(M^2)}(M^2+\omega+1)}(M^2+\psi_{\psi_{\psi_M(M^2)}(M^2+1)}(1)))%
}

\NovelDisplay{M02388R01}{%
(0)(1,1,1)(2,1,1)(3,1)(2,1)(3,2,1)(4,1)(5,2)%
}

\NovelDisplay{M02389R01}{%
= \psi(I+\psi_{\Omega_{\psi_I(I)+\omega+1}}(I+\psi_{\Omega_{\psi_I(I)+1}}(I+\Omega_{\psi_I(I)+1})))%
}

\NovelDisplay{M02390R01}{%
=\! \psi(M^2\!+\!\psi_{\psi_{\psi_M(M^2)}(M^2+\omega+1)}(M^2\!+\!\psi_{\psi_{\psi_M(M^2)}(M^2+1)}(M^2\!+\!\psi_{\psi_M(M^2)}(M^2\!+\!1))))%
}

\NovelDisplay{M02391R01}{%
(0)(1,1,1)(2,1,1)(3,1)(2,1)(3,2,1)(4,1)(5,2,1)%
}

\NovelDisplay{M02392R01}{%
= \psi(I+\psi_{\Omega_{\psi_I(I)+\omega+1}}(I+\psi_{\Omega_{\psi_I(I)+1}}(I+\Omega_{\psi_I(I)+\omega})))%
}

\NovelDisplay{M02393R01}{%
=\! \psi(M^2\!+\!\psi_{\psi_{\psi_M(M^2)}(M^2+\omega+1)}(M^2\!+\!\psi_{\psi_{\psi_M(M^2)}(M^2+1)}(M^2\!+\!\psi_{\psi_M(M^2)}(M^2\!+\!\omega))))%
}

\NovelDisplay{M02394R01}{%
(0)(1,1,1)(2,1,1)(3,1)(2,1)(3,2,1)(4,2)%
}

\NovelDisplay{M02395R01}{%
= \psi(I+\psi_{\Omega_{\psi_I(I)+\omega+1}}(I+\Omega_{\psi_I(I)+1}))%
}

\NovelDisplay{M02396R01}{%
= \psi(M^2+\psi_{\psi_{\psi_M(M^2)}(M^2+\omega+1)}(M^2+\psi_{\psi_M(M^2)}(M^2+1)))%
}

\NovelDisplay{M02397R01}{%
(0)(1,1,1)(2,1,1)(3,1)(2,1)(3,2,1)(4,2)(3,2,1)%
}

\NovelDisplay{M02398R01}{%
= \psi(I+\psi_{\Omega_{\psi_I(I)+\omega+1}}(I+\Omega_{\psi_I(I)+\omega}))%
}

\NovelDisplay{M02399R01}{%
= \psi(M^2+\psi_{\psi_{\psi_M(M^2)}(M^2+\omega+1)}(M^2+\psi_{\psi_M(M^2)}(M^2+\omega)))%
}

\NovelDisplay{M02400R01}{%
(0)(1,1,1)(2,1,1)(3,1)(2,1)(3,2,1)(4,2)(3,2,1)(4,2)(3,2,1)%
}

\NovelDisplay{M02401R01}{%
= \psi(I+\psi_{\Omega_{\psi_I(I)+\omega+1}}(I+\Omega_{\psi_I(I)+\omega})\times 2)%
}

\NovelDisplay{M02402R01}{%
= \psi(M^2+\psi_{\psi_{\psi_M(M^2)}(M^2+\omega+1)}(M^2+\psi_{\psi_M(M^2)}(M^2+\omega))\times 2)%
}

\NovelDisplay{M02403R01}{%
(0)(1,1,1)(2,1,1)(3,1)(2,1)(3,2,1)(4,2)(5,3)%
}

\NovelDisplay{M02404R01}{%
= \psi(I+\Omega_{\psi_I(I)+\omega+1})%
}

\NovelDisplay{M02405R01}{%
= \psi(M^2+\psi_{\psi_M(M^2)}(M^2+\omega+1))%
}

\NovelDisplay{M02406R01}{%
(0)(1,1,1)(2,1,1)(3,1)(2,1)(3,2,1)(4,2)(5,3)(6,4)%
}

\NovelDisplay{M02407R01}{%
= \psi(I+\Omega_{\psi_I(I)+\omega+2})%
}

\NovelDisplay{M02408R01}{%
= \psi(M^2+\psi_{\psi_M(M^2)}(M^2+\omega+2))%
}

那天她们没有越过新的刻度，却第一次把几条很相近的长式分别留在脑中。下山以后，薛峰把自己记过的停顿擦去一半，尝试不靠标记再读。

读到最深的一处，他还会断。可断开的位置，已经能由他自己找出来了。

夜里，卢本韦将白天的几种尾部分开，先不写新的对应，只请霍雨浩听自己把位置说清。她一直把有限的一念得很快，念到ω之后再添一，反倒需要停一下，确认这个一没有被自己吞进前面的部分。

蔡须坤在另一边核M侧，听见她停，便将同一个位置递过来。两人的难点并不相同，可只要共同的输入仍在，就能知道眼下应该一起看哪一个括号。

薛峰原本已准备回去睡，听见两人互相问，也留下来遮了一次答案。他第一次发现，自己在I侧能先看清的那一点，不必等到整套M全会了才算有用。把这一点说明白，蔡须坤便能从对应的位置继续往里找。

有限的尾部开始排成更长的一列。ω乘二，ω乘三，再到ω平方，纸面上用到的字都不陌生，陌生的是它们被放进的位置。

诺当把其中两行说反，蔡兴坤跟着也乱了。两人看着前面已经写好的几页，越想赶上，越容易跳过同样的括号。

卢本韦原本正等蔡须坤算M侧，听见后面的声音，便先退了两级。她没有把自己的整本I笔记塞过去，只挑出当前这条下标，把有限倍数与幂次落下的位置分别指出。

``先说乘的是谁。别急着说到最后。''

蔡兴坤照着说了一遍，才发现自己把乘三读成了第三层。诺当则在断行处多添了一个括号。薛峰把两人的纸叠在一起，错处恰好落在不同位置，他便让两人各自找出对方的一处。

上方，蔡须坤等到卢本韦回头，已经把M侧遇到的那个较短ψ项圈了出来。它在这份对应表里承担的是ω平方所在的位置，不能因为写起来比旁边短，就当成一个可以随处替换的普通常数。

她没有给它另起名字，只沿整条M式读到这里，再读出去。

``你刚才在教什么？''她问。

``教他们别把第几层和乘几次混着说。''

蔡须坤低头看了一眼自己的纸，悄悄将一个写在错误层次的三擦掉。卢本韦看见了，没有拆穿，只把两份纸重新对齐。

同一条长式在讲义上出现了两回。薛峰逐列核到末尾，确认两回连右栏都相同，便在自己的目录里只给它记一个推进位置，另一回留作遮住答案后的复算。重复印出的式子仍然保留，山路却没有因此凭空多出一级。

\NovelDisplay{M02409R01}{%
(0)(1,1,1)(2,1,1)(3,1)(2,1)(3,2,1)(4,2)(5,3,1)%
}

\NovelDisplay{M02410R01}{%
= \psi(I+\Omega_{\psi_I(I)+\omega\times 2})%
}

\NovelDisplay{M02411R01}{%
= \psi(M^2+\psi_{\psi_M(M^2)}(M^2+\omega\times 2))%
}

\NovelDisplay{M02412R01}{%
(0)(1,1,1)(2,1,1)(3,1)(2,1)(3,2,1)(4,2)(5,3,1)(6,3)(7,4,1)%
}

\NovelDisplay{M02413R01}{%
= \psi(I+\Omega_{\psi_I(I)+\omega\times 3})%
}

\NovelDisplay{M02414R01}{%
= \psi(M^2+\psi_{\psi_M(M^2)}(M^2+\omega\times 3))%
}

\NovelDisplay{M02415R01}{%
(0)(1,1,1)(2,1,1)(3,1)(2,1)(3,2,1)(4,2,1)%
}

\NovelDisplay{M02416R01}{%
= \psi(I+\Omega_{\psi_I(I)+\omega^2})%
}

\NovelDisplay{M02417R01}{%
= \psi(M^2+\psi_{\psi_M(M^2)}(M^2+\psi(2)))%
}

\NovelDisplay{M02418R01}{%
(0)(1,1,1)(2,1,1)(3,1)(2,1)(3,2,1)(4,2,1)(4,2)(5,3)%
}

\NovelDisplay{M02419R01}{%
= \psi(I+\Omega_{\psi_I(I)+\omega^2+1})%
}

\NovelDisplay{M02420R01}{%
= \psi(M^2+\psi_{\psi_M(M^2)}(M^2+\psi(2)+1))%
}

\NovelDisplay{M02421R01}{%
(0)(1,1,1)(2,1,1)(3,1)(2,1)(3,2,1)(4,2,1)(4,2)(5,3,1)%
}

\NovelDisplay{M02422R01}{%
= \psi(I+\Omega_{\psi_I(I)+\omega^2+\omega})%
}

\NovelDisplay{M02423R01}{%
= \psi(M^2+\psi_{\psi_M(M^2)}(M^2+\psi(2)+\omega))%
}

\NovelDisplay{M02424R01}{%
(0)(1,1,1)(2,1,1)(3,1)(2,1)(3,2,1)(4,2,1)(4,2)(5,3,1)(6,3)(7,4,1)%
}

\NovelDisplay{M02425R01}{%
= \psi(I+\Omega_{\psi_I(I)+\omega^2+\omega\times 2})%
}

\NovelDisplay{M02426R01}{%
= \psi(M^2+\psi_{\psi_M(M^2)}(M^2+\psi(2)+\omega\times 2))%
}

\NovelDisplay{M02427R01}{%
(0)(1,1,1)(2,1,1)(3,1)(2,1)(3,2,1)(4,2,1)(4,2)(5,3,1)(6,3,1)%
}

\NovelDisplay{M02428R01}{%
= \psi(I+\Omega_{\psi_I(I)+\omega^2 \times 2})%
}

\NovelDisplay{M02429R01}{%
= \psi(M^2+\psi_{\psi_M(M^2)}(M^2+\psi(2)\times 2))%
}

\NovelDisplay{M02430R01}{%
(0)(1,1,1)(2,1,1)(3,1)(2,1)(3,2,1)(4,2,1)(4,2)(5,3,1)(6,3,1)(6,3)%
}
\NovelDisplay{M02430R02}{%
(7,4,1)(8,4,1)%
}

\NovelDisplay{M02431R01}{%
= \psi(I+\Omega_{\psi_I(I)+\omega^2 \times 3})%
}

\NovelDisplay{M02432R01}{%
= \psi(M^2+\psi_{\psi_M(M^2)}(M^2+\psi(2)\times 3))%
}

\NovelDisplay{M02433R01}{%
(0)(1,1,1)(2,1,1)(3,1)(2,1)(3,2,1)(4,2,1)(4,2)(5,3,1)(6,3,1)(6,3)%
}
\NovelDisplay{M02433R02}{%
(7,4,1)(8,4,1)%
}

\NovelDisplay{M02434R01}{%
= \psi(I+\Omega_{\psi_I(I)+\omega^2 \times 3})%
}

\NovelDisplay{M02435R01}{%
= \psi(M^2+\psi_{\psi_M(M^2)}(M^2+\psi(2)\times 3))%
}

\NovelDisplay{M02436R01}{%
(0)(1,1,1)(2,1,1)(3,1)(2,1)(3,2,1)(4,2,1)(4,2,1)%
}

\NovelDisplay{M02437R01}{%
= \psi(I+\Omega_{\psi_I(I)+\omega^3})%
}

\NovelDisplay{M02438R01}{%
= \psi(M^2+\psi_{\psi_M(M^2)}(M^2+\psi(3)))%
}

一场雨落到纸边，最先模糊的是薛峰写的小字。他将那页护进衣襟，等雨歇才拿出来晾。

蔡须坤把自己较干的那一页借给他。她看着他将重复式两次独立写出，心里忽然松了一点：这几日总是旁人等她，她也终于能把自己已经核清的部分，交给需要的人。

后面的下标里，又装进了他们此前才费力认清的函数值。

卢本韦翻开I侧，第一眼只觉得长。等她认真读进去，才意识到，外面看似相似的Ω仍留着，里面的旧项却已经从加法走到更复杂的乘法与嵌套。每次看到熟悉的ψ，她都得先问自己：这一个现在是值，是某个函数的下标，还是整个Ω的下标中的一部分。

蔡须坤的M侧更密。她读到其中最长的一行时，手指顺着纸面挪到了尽头，意识却仍没有回到最外层。

``我丢了。''她忽然说。

卢本韦停住。``丢哪儿了？''

``不知道。''

她说完便垂下眼。最难堪的不是算不出下一项，而是连自己从哪里开始失去位置，都说不清。山上的压力还将她往下按，胸口的灰意跟着微微一动，她立刻将魂力收慢。

古月娜没有碰她腕下，只扶着她退到平处，留出能自己调息的距离。霍雨浩将最长的一行折起来，先露出最外面函数的头尾。

``先找回这一圈。里面可以等。''

蔡须坤看了很久，终于将完整的外层说出来。薛峰这时递来另一张纸，上面没有答案，只有她先前依次读过的几处入口。哪一处读完应当返回，他都按她说过的顺序留好了。

她顺着第一处往下找，第二处停住，第三处终于接上。卢本韦把I侧同一个位置递到旁边，让她看见，较短的表示也需要保留那一段，绝非因为换成M才多出了一道没有来由的难题。

这一次蔡须坤没有整行重抄。她只重写了丢失位置的那一小段，再亲手将它放回完整长式。纸上的接缝仍明显，脑中的路却终于没有断。

\NovelDisplay{M02439R01}{%
(0)(1,1,1)(2,1,1)(3,1)(2,1)(3,2,1)(4,2,1)(5,1)%
}

\NovelDisplay{M02440R01}{%
= \psi(I+\Omega_{\psi_I(I)+\Omega})%
}

\NovelDisplay{M02441R01}{%
= \psi(M^2+\psi_{\psi_M(M^2)}(M^2+\psi_{\psi_M(M^2)}(1)))%
}

\NovelDisplay{M02442R01}{%
(0)(1,1,1)(2,1,1)(3,1)(2,1)(3,2,1)(4,2,1)(5,1)(1,1,1)(2,1,1)(3,1)(2)%
}

\NovelDisplay{M02443R01}{%
= \psi(I+\Omega_{\psi_I(I)\times 2})%
}

\NovelDisplay{M02444R01}{%
= \psi(M^2+\psi_{\psi_M(M^2)}(M^2+\psi_{\psi_M(M^2)}(M^2)))%
}

\NovelDisplay{M02445R01}{%
(0)(1,1,1)(2,1,1)(3,1)(2,1)(3,2,1)(4,2,1)(5,1)(4,2)(5,3)%
}

\NovelDisplay{M02446R01}{%
= \psi(I+\Omega_{\psi_I(I)\times 2+1})%
}

\NovelDisplay{M02447R01}{%
= \psi(M^2+\psi_{\psi_M(M^2)}(M^2+\psi_{\psi_M(M^2)}(M^2)+1))%
}

\NovelDisplay{M02448R01}{%
(0)(1,1,1)(2,1,1)(3,1)(2,1)(3,2,1)(4,2,1)(5,1)(4,2)(5,3,1)(6,3,1)(7,1)%
}
\NovelDisplay{M02448R02}{%
(1,1,1)(2,1,1)(3,1)(2)%
}

\NovelDisplay{M02449R01}{%
= \psi(I+\Omega_{\psi_I(I)\times 3})%
}

\NovelDisplay{M02450R01}{%
= \psi(M^2+\psi_{\psi_M(M^2)}(M^2+\psi_{\psi_M(M^2)}(M^2)\times 2))%
}

\NovelDisplay{M02451R01}{%
(0)(1,1,1)(2,1,1)(3,1)(2,1)(3,2,1)(4,2,1)(5,1)(4,2,1)%
}

\NovelDisplay{M02452R01}{%
= \psi(I+\Omega_{\psi_I(I)\times \omega})%
}

\NovelDisplay{M02453R01}{%
= \psi(M^2+\psi_{\psi_M(M^2)}(M^2+\psi_{\psi_{\psi_M(M^2)}(M^2+1)}(1)))%
}

\NovelDisplay{M02454R01}{%
(0)(1,1,1)(2,1,1)(3,1)(2,1)(3,2,1)(4,2,1)(5,1)(4,2,1)(4,2,1)%
}

\NovelDisplay{M02455R01}{%
= \psi(I+\Omega_{\psi_I(I)\times \omega^2})%
}

\NovelDisplay{M02456R01}{%
= \psi(M^2+\psi_{\psi_M(M^2)}(M^2+\psi_{\psi_{\psi_M(M^2)}(M^2+1)}(2)))%
}

\NovelDisplay{M02457R01}{%
(0)(1,1,1)(2,1,1)(3,1)(2,1)(3,2,1)(4,2,1)(5,1)(4,2,1)(5,1)(4,2,1)%
}

\NovelDisplay{M02458R01}{%
= \psi(I+\Omega_{\psi_{\Omega_{\psi_I(I)+1}}(I+1)})%
}

\NovelDisplay{M02459R01}{%
= \psi(M^2+\psi_{\psi_M(M^2)}(M^2+\psi_{\psi_{\psi_M(M^2)}(M^2+1)}(M^2+1)))%
}

\NovelDisplay{M02460R01}{%
(0)(1,1,1)(2,1,1)(3,1)(2,1)(3,2,1)(4,2,1)(5,1)(5)%
}

\NovelDisplay{M02461R01}{%
= \psi(I+\Omega_{\psi_{\Omega_{\psi_I(I)+1}}(I+\psi_I(I)\times \omega)})%
}

\NovelDisplay{M02462R01}{%
= \psi(M^2+\psi_{\psi_M(M^2)}(M^2+\psi_{\psi_{\psi_M(M^2)}(M^2+1)}(M^2+\psi_{\psi_{\psi_M(M^2)}(M^2+1)}(1))))%
}

\NovelDisplay{M02463R01}{%
(0)(1,1,1)(2,1,1)(3,1)(2,1)(3,2,1)(4,2,1)(5,1)(5,1)(4,2,1)%
}

\NovelDisplay{M02464R01}{%
= \psi(I+\Omega_{\psi_{\Omega_{\psi_I(I)+1}}(I+\psi_{\Omega_{\psi_I(I)+1}}(I)+1)})%
}

\NovelDisplay{M02465R01}{%
= \psi(M^2+\psi_{\psi_M(M^2)}(M^2+\psi_{\psi_{\psi_M(M^2)}(M^2+1)}(M^2+\psi_{\psi_{\psi_M(M^2)}(M^2+1)}(M^2)+1)))%
}

\NovelDisplay{M02466R01}{%
(0)(1,1,1)(2,1,1)(3,1)(2,1)(3,2,1)(4,2,1)(5,1)(5,1)(5)%
}

\NovelDisplay{M02467R01}{%
= \psi(I+\Omega_{\psi_{\Omega_{\psi_I(I)+1}}(I+\psi_{\Omega_{\psi_I(I)+1}}(I+1))})%
}

\NovelDisplay{M02468R01}{%
= \psi(M^2+\psi_{\psi_M(M^2)}(M^2+\psi_{\psi_{\psi_M(M^2)}(M^2+1)}(M^2+\psi_{\psi_{\psi_M(M^2)}(M^2+1)}(M^2+1))))%
}

\NovelDisplay{M02469R01}{%
(0)(1,1,1)(2,1,1)(3,1)(2,1)(3,2,1)(4,2,1)(5,1)(6)%
}

\NovelDisplay{M02470R01}{%
= \psi(I+\Omega_{\psi_{\Omega_{\psi_I(I)+1}}(I+\psi_{\Omega_{\psi_I(I)+1}}(I+\psi_I(I)\times \omega))})%
}

\NovelDisplay{M02471R01}{%
= \psi(M^2+\psi_{\psi_M(M^2)}(M^2+\psi_{\psi_{\psi_M(M^2)}(M^2+1)}(M^2+\psi_{\psi_{\psi_M(M^2)}(M^2+1)}(M^2%
}
\NovelDisplay{M02471R02}{%
+\psi_{\psi_{\psi_M(M^2)}(M^2+1)}(1)))))%
}

\NovelDisplay{M02472R01}{%
(0)(1,1,1)(2,1,1)(3,1)(2,1)(3,2,1)(4,2,1)(5,1)(6,2)%
}

\NovelDisplay{M02473R01}{%
= \psi(I+\Omega_{\psi_{\Omega_{\psi_I(I)+1}}(I+\Omega_{\psi_I(I)+1})})%
}

\NovelDisplay{M02474R01}{%
= \psi(M^2+\psi_{\psi_M(M^2)}(M^2+\psi_{\psi_{\psi_M(M^2)}(M^2+1)}(M^2+\psi_{\psi_M(M^2)}(M^2+1))))%
}

\NovelDisplay{M02475R01}{%
(0)(1,1,1)(2,1,1)(3,1)(2,1)(3,2,1)(4,2,1)(5,1)(6,2,1)%
}

\NovelDisplay{M02476R01}{%
= \psi(I+\Omega_{\psi_{\Omega_{\psi_I(I)+1}}(I+\Omega_{\psi_I(I)+\omega})})%
}

\NovelDisplay{M02477R01}{%
= \psi(M^2+\psi_{\psi_M(M^2)}(M^2+\psi_{\psi_{\psi_M(M^2)}(M^2+1)}(M^2+\psi_{\psi_M(M^2)}(M^2+\omega))))%
}

\NovelDisplay{M02478R01}{%
(0)(1,1,1)(2,1,1)(3,1)(2,1)(3,2,1)(4,2,1)(5,1)(6,2,1)(7,2,1)(8,1)%
}

\NovelDisplay{M02479R01}{%
= \psi(I+\Omega_{\psi_{\Omega_{\psi_I(I)+1}}(I+\Omega_{\psi_{\Omega_{\psi_I(I)+1}}}(\Omega))})%
}

\NovelDisplay{M02480R01}{%
= \psi(M^2+\psi_{\psi_M(M^2)}(M^2+\psi_{\psi_{\psi_M(M^2)}(M^2+1)}(M^2+\psi_{\psi_M(M^2)}(M^2%
}
\NovelDisplay{M02480R02}{%
+\psi_{\psi_{\psi_{M}(M^2)}(M^2+1)}(M)))))%
}

\NovelDisplay{M02481R01}{%
(0)(1,1,1)(2,1,1)(3,1)(2,1)(3,2,1)(4,2,1)(5,1)(6,2,1)(7,2,1)(8,1)%
}

\NovelDisplay{M02482R01}{%
= \psi(I+\Omega_{\psi_{\Omega_{\psi_I(I)+1}}(I+\Omega_{\psi_{\Omega_{\psi_I(I)+1}}}(\Omega))})%
}

\NovelDisplay{M02483R01}{%
= \psi(M^2+\psi_{\psi_M(M^2)}(M^2+\psi_{\psi_{\psi_M(M^2)}(M^2+1)}(M^2+\psi_{\psi_M(M^2)}(M^2%
}
\NovelDisplay{M02483R02}{%
+\psi_{\psi_{\psi_{M}(M^2)}(M^2+1)}(M)))))%
}

\NovelDisplay{M02484R01}{%
(0)(1,1,1)(2,1,1)(3,1)(2,1)(3,2,1)(4,2,1)(5,1)(6,2,1)(7,2,1)(8,1)(9,2)%
}

\NovelDisplay{M02485R01}{%
= \psi(I+\Omega_{\psi_{\Omega_{\psi_I(I)+1}}(I+\Omega_{\psi_{\Omega_{\psi_I(I)+1}}}(I+\Omega_{\psi_I(I)+1}))})%
}

\NovelDisplay{M02486R01}{%
= \psi(M^2+\psi_{\psi_M(M^2)}(M^2+\psi_{\psi_{\psi_M(M^2)}(M^2+1)}(M^2+\psi_{\psi_M(M^2)}(M^2%
}
\NovelDisplay{M02486R02}{%
+\psi_{\psi_{\psi_{M}(M^2)}(M^2+1)}(M^2+\psi_{\psi_{M}(M^2)}(M^2+1))))))%
}

\NovelDisplay{M02487R01}{%
(0)(1,1,1)(2,1,1)(3,1)(2,1)(3,2,1)(4,2,1)(5,1)(6,2,1)(7,2,1)(8,1)(9,2,1)%
}
\NovelDisplay{M02487R02}{%
(10,2,1)(11,1)%
}

\NovelDisplay{M02488R01}{%
= \psi(I+\Omega_{\psi_{\Omega_{\psi_I(I)+1}}(I+\Omega_{\psi_{\Omega_{\psi_I(I)+1}}}(I+\Omega_{\psi_{\Omega_{\psi_I(I)+1}}(\Omega)}))})%
}

\NovelDisplay{M02489R01}{%
= \psi(M^2+\psi_{\psi_M(M^2)}(M^2+\psi_{\psi_{\psi_M(M^2)}(M^2+1)}(M^2+\psi_{\psi_M(M^2)}(M^2%
}
\NovelDisplay{M02489R02}{%
+\psi_{\psi_{\psi_{M}(M^2)}(M^2+1)}(M^2+\psi_{\psi_{M}(M^2)}(M^2+\psi_{\psi_{\psi_M(M^2)}(M^2+1)}(M)))))))%
}

当天返回营地，薛峰的目录里添上了``入口、返回、外层''几个字。他写得很用力，炭粉落在指节上。

霍雨浩看过，叫他试着替叶阳解释一行。起初薛峰还总回头等他补话，讲到后半段，却开始自己判断哪里应当先停。霍雨浩便没有再插进去。

蔡须坤后来才问，自己选M是不是拖慢了队伍。

卢本韦正晾着薛峰那张受潮的目录，听见便停住手。``我今天漏下标的时候，也有人等我。''

``可我这一边更长。''

``那就告诉我长在哪里。你不说，我只能看见你在发呆。''

蔡须坤握着纸，过了一会儿，才把自己最容易丢失的返回位置指出来。卢本韦照着看，很快也迷了一次，这才明白，刚才那种长久的沉默里，她其实已经试着走过了多少层。

两人没有换掉各自主攻的记号。I侧负责先认整体，M侧保留更细的函数组织，遇到明确的困难就把位置交出来。薛峰将她们新说清的地方补进目录，霍雨浩最后核过，几人的分工才从最初的一句约定，慢慢长成了能在难处真正接上的帮助。

更深的Ω下标终于一层一层叠了上去，随后在新的一行里，I侧出现了乘二。

卢本韦把两页放到一起，没敢立即说这就是把前面的过程再来一次。旧页中的每一层仍要按那一层自身的规则核对，不能因为最后写成了一个短式，就将中间都当成毫无区别的重复。

蔡须坤却在M侧看见了一处很清楚的变化。原先那些越来越深的下标，都被保留在更复杂的参数组织中；到了这一行，最外层ψ里面仍有M平方，接在后面的却换成一个明确的M侧函数值。

她将两行分别读出，没有把整个结果误读成两个最外层ψ的值相加。

``乘二在I里面这一层，外面的ψ还是同一个。''她说。

卢本韦点头，继而又问：``那M那边为什么不是把M平方也乘二？''

蔡须坤原想说因为I对应M平方，话到一半便停住。她们在第一天就提醒过自己，相同的坍缩值，不能拿来证明两套体系的内部符号可以随意互换。

霍雨浩将这个问题留在纸边，要求她只说眼前已经给出的对应。她便把M侧完整读完，明确它在这一式里如何写，没有顺势造出一个跨体系的通用替换规则。

古月娜看见她停住，反倒点了点头。她们此前最容易犯的错，是以为停下来就会退步，遇到缺少依据的地方也要赶着接。眼下这一步能够停住，后面的纸才不至于全压在一个未经说明的猜想上。

十一人分开复算这组末尾。薛峰第一次比霍雨浩更早说清一处外层位置，他自己愣了愣，又把左列从头核过，确认不是猜中才将纸递出去。

\NovelDisplay{M02490R01}{%
(0)(1,1,1)(2,1,1)(3,1)(2,1)(3,2,1)(4,2,1)(5,2)%
}

\NovelDisplay{M02491R01}{%
= \psi(I+\Omega_{\Omega_{\psi_I(I)+1}})%
}

\NovelDisplay{M02492R01}{%
= \psi(M^2+\psi_{\psi_M(M^2)}(M^2+\psi_{\psi_M(M^2)}(M^2+1)))%
}

\NovelDisplay{M02493R01}{%
(0)(1,1,1)(2,1,1)(3,1)(2,1)(3,2,1)(4,2,1)(5,2)(3,2,1)%
}

\NovelDisplay{M02494R01}{%
= \psi(I+\Omega_{\Omega_{\psi_I(I)+\omega}})%
}

\NovelDisplay{M02495R01}{%
= \psi(M^2+\psi_{\psi_M(M^2)}(M^2+\psi_{\psi_M(M^2)}(M^2+\omega)))%
}

\NovelDisplay{M02496R01}{%
(0)(1,1,1)(2,1,1)(3,1)(2,1)(3,2,1)(4,2,1)(5,2)(3,2,1)(4,2,1)(5,2)%
}

\NovelDisplay{M02497R01}{%
= \psi(I+\Omega_{\Omega_{\Omega_{\psi_I(I)+1}}})%
}

\NovelDisplay{M02498R01}{%
= \psi(M^2+\psi_{\psi_M(M^2)}(M^2+\psi_{\psi_M(M^2)}(M^2+\psi_{\psi_M(M^2)}(M^2+1))))%
}

\NovelDisplay{M02499R01}{%
(0)(1,1,1)(2,1,1)(3,1)(2,1)(3,2,1)(4,2,1)(5,2)(3,2,1)(4,2,1)(5,2)%
}
\NovelDisplay{M02499R02}{%
(3,2,1)(4,2,1)(5,2)%
}

\NovelDisplay{M02500R01}{%
= \psi(I+\Omega_{\Omega_{\Omega_{\Omega_{\psi_I(I)+1}}}})%
}

\NovelDisplay{M02501R01}{%
=\! \psi(M^2\!+\!\psi_{\psi_M(M^2)}(M^2\!+\!\psi_{\psi_M(M^2)}(M^2\!+\!\psi_{\psi_M(M^2)}(M^2\!+\!\psi_{\psi_M(M^2)}(M^2\!+\!1)))))%
}

\NovelDisplay{M02502R01}{%
(0)(1,1,1)(2,1,1)(3,1)(2,1)(3,2,1)(4,2,1)(5,2)(4)%
}

\NovelDisplay{M02503R01}{%
= \psi(I\times 2)%
}

\NovelDisplay{M02504R01}{%
= \psi(M^2+\psi_M(M^2))%
}

那晚，蔡须坤没有把M侧收起来。

卢本韦问她要不要换一天用I。她想了一阵，摇头：``我能借你的看路。但这边已经走到这里，总得知道自己是怎么走过来的。''

卢本韦没有再劝，只将灯移到两本中间。两种记号铺在同一张桌上，谁也不再想着只看自己较顺手的那一侧。

卢本韦又回头看了一遍M侧，将方才差点写出的直接替换式留在废纸上。

同一条BMS在两套记号里给出相等的完整值，并不自动说明两套内部表示有一个保留每种加法、乘法的逐符号替换。她若凭那一处熟悉的边界，就将I乘二写成M平方乘二，等于偷偷加上了讲义尚未给她的条件。

霍雨浩让她把这个条件说出来，她一时答不上，便没有继续用。

蔡须坤从当前的M式指出实际出现的那段函数值，把它重新放到最外层ψ的参数里。两人终于能够不靠一个想当然的替换，也把眼前这一条对照说清。那张废纸被卢本韦折起来，夹在最容易再犯同样错误的地方，留下的不是新公式，而是需要先问的一句话。

以I乘二为背景以后，旧时熟悉的尾部再一次出现。

这一次最先停笔的是薛峰。他把刚写的下标擦掉，翻回前章，找出当时含有旧函数值的那几行。

``这里的背景已经变了。''他说。

卢本韦顺着看过去。她在前面的参数里留住了乘二，下标里的旧函数值却还写着原来的参数。两个位置一个改、一个没改，整条式子读起来仍然顺口，真正比较时却已经不是这一页给出的那条。

她没有马上把错处全改掉，先将原来写下的版本保留在页边，再抄一遍当前讲义的写法。两份并排以后，蔡须坤才也看见自己M侧漏掉的一处：更深的参数中，新的背景需要被完整带回，不能只在最外面写过一次便算交代完。

霍雨浩将两人的位置对照，又叫薛峰说这处差异为什么值得停。

``只要后面还用到这个函数值，就得知道它是从哪一个参数求出来的。名字长得像，不能替它补回参数。''

薛峰说到最后，声音仍有点紧。他并没有突然变得什么都懂，但这一次，他能把自己发现的问题讲到旁人也看见为止。

唐三让出靠石壁的那一小块平处，叶阳将湿纸换成干纸，奥斯卡把补给分到几人手里。卢本韦咬了一口，才发现自己已经很久没有离开那几道下标。蔡须坤接过纸时看了薛峰一眼，主动把右栏朝他转过去。

``你再替我查一遍，别只查她的。''

薛峰坐下来，从她能够自己说明的最外层开始核。

\NovelDisplay{M02505R01}{%
(0)(1,1,1)(2,1,1)(3,1)(2,1)(3,2,1)(4,2,1)(5,2)(4)(1,1,1)(2,1,1)(3,1)(2)%
}

\NovelDisplay{M02506R01}{%
= \psi(I\times 2+\psi_I(I))%
}

\NovelDisplay{M02507R01}{%
= \psi(M^2+\psi_M(M^2)+\psi_{\psi_{M}(M^2)}(M^2))%
}

\NovelDisplay{M02508R01}{%
(0)(1,1,1)(2,1,1)(3,1)(2,1)(3,2,1)(4,2,1)(5,2)(4)(3,2)%
}

\NovelDisplay{M02509R01}{%
= \psi(I\times 2+\Omega_{\psi_I(I)+1})%
}

\NovelDisplay{M02510R01}{%
= \psi(M^2+\psi_M(M^2)+\psi_{\psi_{M}(M^2)}(M^2+1))%
}

\NovelDisplay{M02511R01}{%
(0)(1,1,1)(2,1,1)(3,1)(2,1)(3,2,1)(4,2,1)(5,2)(4)(3,2,1)%
}

\NovelDisplay{M02512R01}{%
= \psi(I\times 2+\Omega_{\psi_I(I)+\omega})%
}

\NovelDisplay{M02513R01}{%
= \psi(M^2+\psi_M(M^2)+\psi_{\psi_{M}(M^2)}(M^2+\omega))%
}

\NovelDisplay{M02514R01}{%
(0)(1,1,1)(2,1,1)(3,1)(2,1)(3,2,1)(4,2,1)(5,2)(4)(3,2,1)(4,2,1)(5,2)%
}

\NovelDisplay{M02515R01}{%
= \psi(I\times 2+\Omega_{\Omega_{\psi_I(I)+1}})%
}

\NovelDisplay{M02516R01}{%
= \psi(M^2+\psi_M(M^2)+\psi_{\psi_{M}(M^2)}(M^2+\psi_{\psi_{M}(M^2)}(M^2+1)))%
}

\NovelDisplay{M02517R01}{%
(0)(1,1,1)(2,1,1)(3,1)(2,1)(3,2,1)(4,2,1)(5,2)(4)(3,2,1)(4,2,1)(5,2)(4)%
}

\NovelDisplay{M02518R01}{%
= \psi(I\times 2+\psi_I(I\times 2))%
}

\NovelDisplay{M02519R01}{%
= \psi(M^2+\psi_M(M^2)+\psi_{\psi_{M}(M^2)}(M^2+\psi_{M}(M^2)))%
}

\NovelDisplay{M02520R01}{%
(0)(1,1,1)(2,1,1)(3,1)(2,1)(3,2,1)(4,2,1)(5,2)(4)(4)%
}

\NovelDisplay{M02521R01}{%
= \psi(I\times 2+\psi_{\Omega_{\psi_I(I\times 2)+1}}(1))%
}

\NovelDisplay{M02522R01}{%
= \psi(M^2+\psi_M(M^2)+\psi_{\psi_{\psi_{M}(M^2)}(M^2+\psi_{M}(M^2)+1)}(1))%
}

\NovelDisplay{M02523R01}{%
(0)(1,1,1)(2,1,1)(3,1)(2,1)(3,2,1)(4,2,1)(5,2)(4,1)(1,1,1)(2,1,1)%
}
\NovelDisplay{M02523R02}{%
(3,1)(2)%
}

\NovelDisplay{M02524R01}{%
= \psi(I\times 2+\psi_{\Omega_{\psi_I(I\times 2)+1}}(I))%
}

\NovelDisplay{M02525R01}{%
= \psi(M^2+\psi_M(M^2)+\psi_{\psi_{\psi_{M}(M^2)}(M^2+\psi_{M}(M^2)+1)}(M^2))%
}

\NovelDisplay{M02526R01}{%
(0)(1,1,1)(2,1,1)(3,1)(2,1)(3,2,1)(4,2,1)(5,2)(4,2)%
}

\NovelDisplay{M02527R01}{%
= \psi(I\times 2+\psi_{\Omega_{\psi_I(I\times 2)+1}}(I+\Omega_{\psi_I(I)+1}))%
}

\NovelDisplay{M02528R01}{%
= \psi(M^2+\psi_M(M^2)+\psi_{\psi_{\psi_{M}(M^2)}(M^2+\psi_{M}(M^2)+1)}(M^2+\psi_{\psi_M(M^2)}(M^2+1)))%
}

\NovelDisplay{M02529R01}{%
(0)(1,1,1)(2,1,1)(3,1)(2,1)(3,2,1)(4,2,1)(5,2)(4,2)(3,2,1)(4,2,1)%
}
\NovelDisplay{M02529R02}{%
(5,2)(4)%
}

\NovelDisplay{M02530R01}{%
= \psi(I\times 2+\psi_{\Omega_{\psi_I(I\times 2)+1}}(I\times 2))%
}

\NovelDisplay{M02531R01}{%
= \psi(M^2+\psi_M(M^2)+\psi_{\psi_{\psi_{M}(M^2)}(M^2+\psi_{M}(M^2)+1)}(M^2+\psi_M(M^2)))%
}

\NovelDisplay{M02532R01}{%
(0)(1,1,1)(2,1,1)(3,1)(2,1)(3,2,1)(4,2,1)(5,2)(4,2)(4)%
}

\NovelDisplay{M02533R01}{%
= \psi(I\times 2+\psi_{\Omega_{\psi_I(I\times 2)+1}}(I\times 2+1))%
}

\NovelDisplay{M02534R01}{%
= \psi(M^2+\psi_M(M^2)+\psi_{\psi_{\psi_{M}(M^2)}(M^2+\psi_{M}(M^2)+1)}(M^2+\psi_M(M^2)+1))%
}

\NovelDisplay{M02535R01}{%
(0)(1,1,1)(2,1,1)(3,1)(2,1)(3,2,1)(4,2,1)(5,2)(4,2)(5,3)%
}

\NovelDisplay{M02536R01}{%
= \psi(I\times 2+\Omega_{\psi_I(I\times 2)+1})%
}

\NovelDisplay{M02537R01}{%
= \psi(M^2+\psi_M(M^2)+\psi_{\psi_{M}(M^2)}(M^2+\psi_{M}(M^2)+1))%
}

\NovelDisplay{M02538R01}{%
(0)(1,1,1)(2,1,1)(3,1)(2,1)(3,2,1)(4,2,1)(5,2)(4,2)(5,3)(6,4)%
}

\NovelDisplay{M02539R01}{%
= \psi(I\times 2+\Omega_{\psi_I(I\times 2)+2})%
}

\NovelDisplay{M02540R01}{%
= \psi(M^2+\psi_M(M^2)+\psi_{\psi_{M}(M^2)}(M^2+\psi_{M}(M^2)+2))%
}

同伴之间的帮助，也从替人指一个错字，慢慢变成了守住一条完整的来路。谁要是说不清自己当前借用的是哪一页的结果，就先把那一处留空，等查回去再补。

夏夜星空看完他们互相留下的空处，没有催着擦净。那些还没有填上的位置，让她看见他们到底在哪里遇到了难题。

第二天一早，薛峰先用空白纸写下昨天改过的一条，随后把正面盖住，从共同的左列重新说明。

他先说自己已经核过哪一列，再说I侧改变的参数，最后才请蔡须坤指出M侧相应位置。蔡须坤听到一半，看见他没有沿着答案倒着猜，便将自己的纸完全收起来，等他讲完。

一遍结束，霍雨浩只补了一处需要更精确的话。薛峰低头将那一句改好，神色却与以往只等别人打勾时不同。他已经能先守住一段，让同伴来查它哪里不够，而不是等整份解释送到面前才开始跟着点头。

此后遇到更长的左列，卢本韦会先将纸交给他；M侧出现旧参数的时候，蔡须坤也会主动请他对照前页。队伍里靠前的位置，就这样一点点有了他的份。

再上山后，有限的三与四接连出现，接着，M侧有一处原先的函数头换了更深的参数组织。

卢本韦先读完I侧，发现这里走到的是I乘ω。她没有把有限的三、四划掉，也没有把最后一行读成``再多一次''。有限的每一次仍在自己的位置，ω这一行必须看完整BMS怎样变，不能将末尾一个数字当作万能的进位开关。

蔡须坤的难处，则落在一个很小的加号上。它藏在M侧函数的下标里，加号后面有另一个M；这个M一旦漏掉，后面接着读出的每一层，都会像是在旧页上重复。

她连续两次念漏。第三次，卢本韦将自己的I侧翻过去，没有让她继续照结果认路。

``先把你自己这一处说给我听。我的晚点再看。''

蔡须坤闭了闭眼，将那一层函数头重新读完。额角的汗慢慢滑到下巴，她没有擦，等自己明确听见那个M，才让卢本韦翻回左栏。

霍雨浩随即请她指出，为什么这一处变化不能塞进最外面那层ψ的参数里。她沿原来的括号，从函数下标读回主参数，终于将两个位置分开。

后面的人在这一段落下了些距离。蔡兴坤还要借助停顿标记，诺当要从原有接缝重新查，古月娜便留在他们旁边，一处处核完整的左列。唐三也没有急着向上赶，将能退回的绳结再摸过一次。

前面四人等到了黄昏。薛峰将最短的I式读给后面听，霍雨浩补足容易漏的前提，蔡须坤与卢本韦则交换着读各自的表示。十一人的答案并不同时出来，却终于都能从自己的纸接回脚下这一级。

\NovelDisplay{M02541R01}{%
(0)(1,1,1)(2,1,1)(3,1)(2,1)(3,2,1)(4,2,1)(5,2)(4,2)(5,3,1)%
}

\NovelDisplay{M02542R01}{%
= \psi(I\times 2+\Omega_{\psi_I(I\times 2)+\omega})%
}

\NovelDisplay{M02543R01}{%
= \psi(M^2+\psi_M(M^2)+\psi_{\psi_{M}(M^2)}(M^2+\psi_{M}(M^2)+\omega))%
}

\NovelDisplay{M02544R01}{%
(0)(1,1,1)(2,1,1)(3,1)(2,1)(3,2,1)(4,2,1)(5,2)(4,2)(5,3,1)(6,3,1)%
}

\NovelDisplay{M02545R01}{%
= \psi(I\times 2+\Omega_{\psi_I(I\times 2)+\omega^2})%
}

\NovelDisplay{M02546R01}{%
= \psi(M^2+\psi_M(M^2)+\psi_{\psi_{M}(M^2)}(M^2+\psi_{M}(M^2)+\psi(2)))%
}

\NovelDisplay{M02547R01}{%
(0)(1,1,1)(2,1,1)(3,1)(2,1)(3,2,1)(4,2,1)(5,2)(4,2)(5,3,1)(6,3,1)(7,3)%
}

\NovelDisplay{M02548R01}{%
= \psi(I\times 2+\Omega_{\Omega_{\psi_I(I\times 2)+1}})%
}

\NovelDisplay{M02549R01}{%
=\! \psi(M^2\!+\!\psi_M(M^2)\!+\!\psi_{\psi_{M}(M^2)}(M^2\!+\!\psi_{M}(M^2)\!+\!\psi_{\psi_{M}(M^2)}(M^2\!+\!\psi_{M}(M^2)\!+\!1)))%
}

\NovelDisplay{M02550R01}{%
(0)(1,1,1)(2,1,1)(3,1)(2,1)(3,2,1)(4,2,1)(5,2)(4,2)(5,3,1)(6,3,1)(7,3)%
}
\NovelDisplay{M02550R02}{%
(5,3,1)(6,3,1)(7,3)%
}

\NovelDisplay{M02551R01}{%
= \psi(I\times 2+\Omega_{\Omega_{\Omega_{\psi_I(I\times 2)+1}}})%
}

\NovelDisplay{M02552R01}{%
= \psi(M^2+\psi_M(M^2)+\psi_{\psi_{M}(M^2)}(M^2+\psi_{M}(M^2)+\psi_{\psi_{M}(M^2)}(M^2+\psi_{M}(M^2)%
}
\NovelDisplay{M02552R02}{%
+\psi_{\psi_{M}(M^2)}(M^2+\psi_{M}(M^2)+1))))%
}

\NovelDisplay{M02553R01}{%
(0)(1,1,1)(2,1,1)(3,1)(2,1)(3,2,1)(4,2,1)(5,2)(4,2)(5,3,1)(6,3,1)%
}
\NovelDisplay{M02553R02}{%
(7,3)(6)%
}

\NovelDisplay{M02554R01}{%
= \psi(I\times 3)%
}

\NovelDisplay{M02555R01}{%
= \psi(M^2+\psi_M(M^2)\times 2)%
}

\NovelDisplay{M02556R01}{%
(0)(1,1,1)(2,1,1)(3,1)(2,1)(3,2,1)(4,2,1)(5,2)(4,2)(5,3,1)(6,3,1)(7,3)%
}
\NovelDisplay{M02556R02}{%
(6,3)(7,4,1)(8,4,1)(9,4)(8)%
}

\NovelDisplay{M02557R01}{%
= \psi(I\times 4)%
}

\NovelDisplay{M02558R01}{%
= \psi(M^2+\psi_M(M^2)\times 3)%
}

\NovelDisplay{M02559R01}{%
(0)(1,1,1)(2,1,1)(3,1)(2,1,1)%
}

\NovelDisplay{M02560R01}{%
= \psi(I\times \omega)%
}

\NovelDisplay{M02561R01}{%
= \psi(M^2+\psi_{\psi_M(M^2+M)}(1))%
}

二百六十米的横台露出来时，卢本韦先靠着石壁坐了下去。

蔡须坤把水放在她手边，手腕依旧平稳。果糕领域在她方才最吃力的时候，只短暂铺到身周，等心神能重新留在纸上，就被她收了回去。它没有替她读懂那个藏在下标里的M。

薛峰落到横台上，向后看了一眼，才发现自己已经能与霍雨浩并排接住后面的问话。他把本子抱紧，没有说自己追上了，只主动去找诺当那一处仍旧不顺的换行。

回到石屋后，几人很久都没开口。

纸上从二百四十米到二百六十米的那一道短线，耗去了他们接连许多次上下山的力气。蔡须坤拆开鞋底补过的一处，看见线还牢牢连着，便又原样收好。

卢本韦把自己的I册递给她，M册则留到手边。她们没有交换负责的那一套记号，只交换今天最容易丢失位置的那几页。

``我也得知道你是从哪儿停下来的。''卢本韦说。

霍雨浩将两套笔记的共同左列核齐，再把炭条交给薛峰。薛峰接过去，没有先翻标准答案，开始从空白处重写当天最难的一段。

古月娜从门口进来，看见桌边的位置换了，便将水一杯杯放好。薛峰还没习惯坐在这样靠前的位置，伸手接杯时碰了一下纸角。蔡须坤把纸压住，等他继续往下写。

窗外的山仍旧很高。三百米尚在更上方，今晚他们先把二百六十米真正留在自己手中。

\chapter{第一百一十五章：谁也不能落下}\label{ux7b2cux4e00ux767eux4e00ux5341ux4e94ux7ae0ux8c01ux4e5fux4e0dux80fdux843dux4e0b}

二百六十米以上，石阶转过一段背阴的山壁。

雨水沿岩缝往下滴，落在脚边，声音很轻。人站得久了，鞋底却会一点点凉透。唐三将绳子绕过外侧的石柱，诺当蹲下去摸索结扣，小舞把能坐人的那一处先擦干。

蔡须坤在岩壁前展开M册，第一眼便认出了上一段末尾那截很难读顺的函数头。它没有离开，如今成了后面许多长式中都要保留的背景。

卢本韦把I册摆到旁边。背景在她这边只占很短的一段，可接在后面的东西越来越复杂，已经不能只凭看见旧名字，就迅速替蔡须坤指一条路。

薛峰走到两人身后，先核共同的左列。霍雨浩等他核完，才将另一份纸压上去。

``今天你先说哪里变化，我来听。''

薛峰拿着笔，没有立刻答。他从来没有在最开始就承担这样的问话。后面十一人各自的本子摊开，连蔡兴坤怀里的黑子鸡也安静下来，只有雨水仍从上方一滴滴落下。

最先接回来的，是原有背景之后那些已经学过的函数值。

薛峰从I侧说起，读完尾部，才请卢本韦补充。他漏念了一层外面的加法，卢本韦没有替他往后说，只在纸边敲了一下。他顺着重新找过，脸慢慢红起来，最终还是自己把漏掉的部分接全。

蔡须坤的M侧读得更慢。新的背景包含M平方与一个较深的函数值，后面又接更早学过的结构。她最容易在后半段认出熟悉的函数头，便以为已经返回整式开头。

这一天她换了位置，将整条式子的最外层指给薛峰看，让他在她每次读到中途时，先问一句现在有没有离开这一层。

薛峰起先只会照问。到第三行，他已经能在蔡须坤自己犹豫以前停笔，因为她说出的返回位置和眼前的闭括号并没有对上。

蔡须坤退回去重读，发现他指出的地方确实少接了一截。

``这回你拦得对。''她说。

薛峰慢慢松开捏紧的笔杆。霍雨浩看过那一处，没有再重复解释，只请他将原因讲给叶阳听。

叶阳听到中途便卡住。薛峰又从最外面开始，讲得很笨，每一句都要确认对方有没有跟到当前的位置。蔡须坤原本想继续写，听着两人的停顿，却也将自己还没真懂的一处找了出来。

原来自己能顺着别人的声音过去，与换一个人仍能独立讲清，差着很长一段。她把那一处重新圈回待核的页里，没因为已经走到二百六十米，就将它当成过去了的事。

\NovelDisplay{M02562R01}{%
(0)(1,1,1)(2,1,1)(3,1)(2,1,1)(1,1,1)(2,1,1)(3,1)(2)%
}

\NovelDisplay{M02563R01}{%
= \psi(I\times \omega+\psi_I(I))%
}

\NovelDisplay{M02564R01}{%
= \psi(M^2+\psi_{\psi_M(M^2+M)}(1)+\psi_{\psi_M(M^2)}(M^2))%
}

\NovelDisplay{M02565R01}{%
(0)(1,1,1)(2,1,1)(3,1)(2,1,1)(1,1,1)(2,1,1)(3,1)(2,1)(3,2)%
}

\NovelDisplay{M02566R01}{%
= \psi(I\times \omega+\psi_{\Omega_{\psi_I(I)+1}}(I+\Omega_{\psi_I(I)+1}))%
}

\NovelDisplay{M02567R01}{%
= \psi(M^2+\psi_{\psi_M(M^2+M)}(1)+\psi_{\psi_{\psi_M(M^2)}(M^2+1)}(M^2+\psi_{\psi_M(M^2)}(M^2+1)))%
}

\NovelDisplay{M02568R01}{%
(0)(1,1,1)(2,1,1)(3,1)(2,1,1)(1,1,1)(2,1,1)(3,1)(2,1)(3,2,1)(4,2,1)%
}
\NovelDisplay{M02568R02}{%
(5,2)(4)%
}

\NovelDisplay{M02569R01}{%
= \psi(I\times \omega+\psi_{\Omega_{\psi_I(I)+1}}(I\times 2))%
}

\NovelDisplay{M02570R01}{%
= \psi(M^2+\psi_{\psi_M(M^2+M)}(1)+\psi_{\psi_{\psi_M(M^2)}(M^2+1)}(M^2+\psi_M(M^2)))%
}

\NovelDisplay{M02571R01}{%
(0)(1,1,1)(2,1,1)(3,1)(2,1,1)(1,1,1)(2,1,1)(3,1)(2,1)(3,2,1)(4,2,1)(5,2)%
}
\NovelDisplay{M02571R02}{%
(4,2)(5,3,1)(6,3,1)(7,3)(6)%
}

\NovelDisplay{M02572R01}{%
= \psi(I\times \omega+\psi_{\Omega_{\psi_I(I)+1}}(I\times 3))%
}

\NovelDisplay{M02573R01}{%
= \psi(M^2+\psi_{\psi_M(M^2+M)}(1)+\psi_{\psi_{\psi_M(M^2)}(M^2+1)}(M^2+\psi_M(M^2)\times 2))%
}

\NovelDisplay{M02574R01}{%
(0)(1,1,1)(2,1,1)(3,1)(2,1,1)(1,1,1)(2,1,1)(3,1)(2,1)(3,2,1)(4,2,1)%
}
\NovelDisplay{M02574R02}{%
(5,2)(4,2,1)%
}

\NovelDisplay{M02575R01}{%
= \psi(I\times \omega+\psi_{\Omega_{\psi_I(I)+1}}(I\times \omega))%
}

\NovelDisplay{M02576R01}{%
= \psi(M^2+\psi_{\psi_M(M^2+M)}(1)+\psi_{\psi_{\psi_M(M^2)}(M^2+1)}(M^2+\psi_{\psi_M(M^2+M)}(1)))%
}

\NovelDisplay{M02577R01}{%
(0)(1,1,1)(2,1,1)(3,1)(2,1,1)(1,1,1)(2,1,1)(3,1)(2,1)(3,2,1)(4,2,1)%
}
\NovelDisplay{M02577R02}{%
(5,2)(4,2,1)(3,2)%
}

\NovelDisplay{M02578R01}{%
= \psi(I\times \omega+\Omega_{\psi_I(I)+1})%
}

\NovelDisplay{M02579R01}{%
= \psi(M^2+\psi_{\psi_M(M^2+M)}(1)+\psi_{\psi_M(M^2)}(M^2+1))%
}

\NovelDisplay{M02580R01}{%
(0)(1,1,1)(2,1,1)(3,1)(2,1,1)(1,1,1)(2,1,1)(3,1)(2,1)(3,2,1)(4,2,1)(5,2)%
}
\NovelDisplay{M02580R02}{%
(4,2,1)(3,2,1)(4,2,1)(5,2)(4)%
}

\NovelDisplay{M02581R01}{%
= \psi(I\times \omega+\psi_I(I\times 2))%
}

\NovelDisplay{M02582R01}{%
= \psi(M^2+\psi_{\psi_M(M^2+M)}(1)+\psi_{\psi_M(M^2)}(M^2+\psi_M(M^2)))%
}

\NovelDisplay{M02583R01}{%
(0)(1,1,1)(2,1,1)(3,1)(2,1,1)(1,1,1)(2,1,1)(3,1)(2,1)(3,2,1)(4,2,1)(5,2)%
}
\NovelDisplay{M02583R02}{%
(4,2,1)(3,2,1)(4,2,1)(5,2)(4,2)(5,3)%
}

\NovelDisplay{M02584R01}{%
= \psi(I\times \omega+\psi_{\Omega_{\psi_I(I\times 2)+1}}(I\times 2+\Omega_{\psi_I(I\times 2)+1}))%
}

\NovelDisplay{M02585R01}{%
= \psi(M^2+\psi_{\psi_M(M^2+M)}(1)+\psi_{\psi_{\psi_M(M^2)}(M^2+\psi_M(M^2)+1)}(M^2+\psi_M(M^2)%
}
\NovelDisplay{M02585R02}{%
+\psi_{\psi_M(M^2)}(M^2+\psi_M(M^2)+1)))%
}

\NovelDisplay{M02586R01}{%
(0)(1,1,1)(2,1,1)(3,1)(2,1,1)(1,1,1)(2,1,1)(3,1)(2,1)(3,2,1)(4,2,1)(5,2)%
}
\NovelDisplay{M02586R02}{%
(4,2,1)(3,2,1)(4,2,1)(5,2)(4,2)(5,3,1)(6,3,1)(7,3)(6)%
}

\NovelDisplay{M02587R01}{%
= \psi(I\times \omega+\psi_{\Omega_{\psi_I(I\times 2)+1}}(I\times 3))%
}

\NovelDisplay{M02588R01}{%
= \psi(M^2+\psi_{\psi_M(M^2+M)}(1)+\psi_{\psi_{\psi_M(M^2)}(M^2+\psi_M(M^2)+1)}(M^2+\psi_M(M^2)\times 2))%
}

\NovelDisplay{M02589R01}{%
(0)(1,1,1)(2,1,1)(3,1)(2,1,1)(1,1,1)(2,1,1)(3,1)(2,1)(3,2,1)(4,2,1)(5,2)%
}
\NovelDisplay{M02589R02}{%
(4,2,1)(3,2,1)(4,2,1)(5,2)(4,2)(5,3,1)(6,3,1)(7,3)(6,3,1)%
}

\NovelDisplay{M02590R01}{%
= \psi(I\times \omega+\psi_{\Omega_{\psi_I(I\times 2)+1}}(I\times \omega))%
}

\NovelDisplay{M02591R01}{%
= \psi(M^2+\psi_{\psi_M(M^2+M)}(1)+\psi_{\psi_{\psi_M(M^2)}(M^2+\psi_M(M^2)+1)}(M^2+\psi_{\psi_M(M^2+M)}(1)))%
}

这组纸没有在一天之内收完。雨停了一阵，又从岩壁上流下来。小舞换过两次擦地的布，古月娜则陪着后面的几人，将同一处外层位置反复认清。

薛峰下山前把自己的解释写在页背，发现字比正面还多。他没有再誊一遍，只把真正能替同伴指出来的一句话，留在最上面。

新的函数值接在背景之后，随后，I侧终于又出现了一个完整的I。

卢本韦读到这里，先没有动脚。她把前面的I乘ω与后面的I分开念过，再说明它们仍一起处在最外层ψ的参数里。

蔡须坤在M侧找到相应的那一段。它并没有跟前面的复杂背景合成一团难以辨认的长串，最后补入的那一部分仍能单独指出。她读完以后，请卢本韦再沿整条I式回看，确认自己找中的确是当前尾部，而不是前面已存在的另一处。

薛峰发现同一行在讲义里又印了一回。这一回不只左列相同，I侧与M侧也完全一致。他照先前的方法，将重复项完整保留，再遮住正面独立核过，没把它记成第二次越过新的高度。

这时，后面忽然传来一声急促的擦响。

诺当脚下一滑，伸手扶石壁，手指没能扣住。他正在解释的那一段也跟着断了，身子向外偏去。唐三立即收紧外侧的绳子，小舞扣住他的手腕，卢本韦已经从上方退回来，肩膀抵住他，将人送回干些的那块石面。

``先别接那句话。''她说，``站稳再说。''

诺当喘着气，仍想把刚才漏下的词补出来。薛峰蹲到他身边，将纸放低，先找回他原来指着的地方，没有从别人已经讲到的末尾开始催。

古月娜在外侧护着几人，等绳子重新松下来，才让唐三调整扣结。雨水从她发梢滴下，她没空去擦。前面四人原本可以继续往上试，此刻却都停住了自己的那一页。

等诺当终于完整说完，卢本韦才收回手，拧去袖口沾上的水。

\NovelDisplay{M02592R01}{%
(0)(1,1,1)(2,1,1)(3,1)(2,1,1)(1,1,1)(2,1,1)(3,1)(2,1)(3,2,1)(4,2,1)(5,2)%
}
\NovelDisplay{M02592R02}{%
(4,2,1)(3,2,1)(4,2,1)(5,2)(4,2)(5,3,1)(6,3,1)(7,3)(6,3,1)(5,3)%
}

\NovelDisplay{M02593R01}{%
= \psi(I\times \omega+\Omega_{\psi_I(I\times 2)+1})%
}

\NovelDisplay{M02594R01}{%
= \psi(M^2+\psi_{\psi_M(M^2+M)}(1)+\psi_{\psi_M(M^2)}(M^2+\psi_M(M^2)+1))%
}

\NovelDisplay{M02595R01}{%
(0)(1,1,1)(2,1,1)(3,1)(2,1,1)(1,1,1)(2,1,1)(3,1)(2,1)(3,2,1)(4,2,1)(5,2)%
}
\NovelDisplay{M02595R02}{%
(4,2,1)(3,2,1)(4,2,1)(5,2)(4,2)(5,3,1)(6,3,1)(7,3)(6,3,1)(5,3,1)%
}
\NovelDisplay{M02595R03}{%
(6,3,1)(7,3)(6)%
}

\NovelDisplay{M02596R01}{%
= \psi(I\times \omega+\psi_I(I\times 3))%
}

\NovelDisplay{M02597R01}{%
= \psi(M^2+\psi_{\psi_M(M^2+M)}(1)+\psi_{\psi_M(M^2)}(M^2+\psi_M(M^2)\times 2))%
}

\NovelDisplay{M02598R01}{%
(0)(1,1,1)(2,1,1)(3,1)(2,1,1)(1,1,1)(2,1,1)(3,1)(2,1)(3,2,1)(4,2,1)(5,2)%
}
\NovelDisplay{M02598R02}{%
(4,2,1)(3,2,1)(4,2,1)(5,2)(4,2)(5,3,1)(6,3,1)(7,3)(6,3,1)(5,3,1)(6,3,1)%
}
\NovelDisplay{M02598R03}{%
(7,3)(6,3)(7,4,1)(8,4,1)(9,4)(8,4,1)(7,4,1)(8,4,1)(9,4)(8)%
}

\NovelDisplay{M02599R01}{%
= \psi(I\times \omega+\psi_I(I\times 4))%
}

\NovelDisplay{M02600R01}{%
= \psi(M^2+\psi_{\psi_M(M^2+M)}(1)+\psi_{\psi_M(M^2)}(M^2+\psi_M(M^2)\times 3))%
}

\NovelDisplay{M02601R01}{%
(0)(1,1,1)(2,1,1)(3,1)(2,1,1)(1,1,1)(2,1,1)(3,1)(2,1,1)%
}

\NovelDisplay{M02602R01}{%
= \psi(I\times \omega+\psi_I(I\times \omega))%
}

\NovelDisplay{M02603R01}{%
= \psi(M^2+\psi_{\psi_M(M^2+M)}(1)+\psi_{\psi_M(M^2)}(M^2+\psi_{\psi_M(M^2+M)}(1)))%
}

\NovelDisplay{M02604R01}{%
(0)(1,1,1)(2,1,1)(3,1)(2,1,1)(1,1,1)(2,1,1)(3,1)(2,1,1)(1,1,1)%
}
\NovelDisplay{M02604R02}{%
(2,1,1)(3,1)(2,1,1)%
}

\NovelDisplay{M02605R01}{%
= \psi(I\times \omega+\psi_I(I\times \omega)\times 2)%
}

\NovelDisplay{M02606R01}{%
= \psi(M^2+\psi_{\psi_M(M^2+M)}(1)+\psi_{\psi_M(M^2)}(M^2+\psi_{\psi_M(M^2+M)}(1))\times 2)%
}

\NovelDisplay{M02607R01}{%
(0)(1,1,1)(2,1,1)(3,1)(2,1,1)(2)%
}

\NovelDisplay{M02608R01}{%
= \psi(I\times \omega+\psi_I(I\times \omega)\times \omega)%
}

\NovelDisplay{M02609R01}{%
= \psi(M^2+\psi_{\psi_M(M^2+M)}(1)+\psi_{\psi_{\psi_M(M^2)}(M^2+\psi_{\psi_M(M^2+M)}(1)+1)}(1))%
}

\NovelDisplay{M02610R01}{%
(0)(1,1,1)(2,1,1)(3,1)(2,1,1)(2,1)(1,1,1)(2,1,1)(3,1)(2,1,1)%
}

\NovelDisplay{M02611R01}{%
= \psi(I\times \omega+\psi_{\Omega_{\psi_I(I\times \omega)+1}}(I\times \omega))%
}

\NovelDisplay{M02612R01}{%
=\! \psi(M^2\!+\!\psi_{\psi_M(M^2+M)}(1)\!+\!\psi_{\psi_{\psi_M(M^2)}(M^2+\psi_{\psi_M(M^2+M)}(1)+1)}(M^2\!+\!\psi_{\psi_M(M^2+M)}(1)))%
}

\NovelDisplay{M02613R01}{%
(0)(1,1,1)(2,1,1)(3,1)(2,1,1)(2,1)(3,2)%
}

\NovelDisplay{M02614R01}{%
= \psi(I\times \omega+\Omega_{\psi_I(I\times \omega)+1})%
}

\NovelDisplay{M02615R01}{%
= \psi(M^2+\psi_{\psi_M(M^2+M)}(1)+\psi_{\psi_M(M^2)}(M^2+\psi_{\psi_M(M^2+M)}(1)+1))%
}

\NovelDisplay{M02616R01}{%
(0)(1,1,1)(2,1,1)(3,1)(2,1,1)(2,1)(3,2,1)%
}

\NovelDisplay{M02617R01}{%
= \psi(I\times \omega+\Omega_{\psi_I(I\times \omega)+\omega})%
}

\NovelDisplay{M02618R01}{%
= \psi(M^2+\psi_{\psi_M(M^2+M)}(1)+\psi_{\psi_M(M^2)}(M^2+\psi_{\psi_M(M^2+M)}(1)+\omega))%
}

\NovelDisplay{M02619R01}{%
(0)(1,1,1)(2,1,1)(3,1)(2,1,1)(2,1)(3,2,1)(4,2,1)(5,2)%
}

\NovelDisplay{M02620R01}{%
= \psi(I\times \omega+\Omega_{\Omega_{\psi_I(I\times \omega)+1}})%
}

\NovelDisplay{M02621R01}{%
= \psi(M^2+\psi_{\psi_M(M^2+M)}(1)+\psi_{\psi_M(M^2)}(M^2+\psi_{\psi_M(M^2+M)}(1)+\psi_{\psi_M(M^2)}(M^2%
}
\NovelDisplay{M02621R02}{%
+\psi_{\psi_M(M^2+M)}(1)+1)))%
}

\NovelDisplay{M02622R01}{%
(0)(1,1,1)(2,1,1)(3,1)(2,1,1)(2,1)(3,2,1)(4,2,1)(5,2)(4)%
}

\NovelDisplay{M02623R01}{%
= \psi(I\times \omega+I)%
}

\NovelDisplay{M02624R01}{%
= \psi(M^2+\psi_{\psi_M(M^2+M)}(1)+\psi_M(M^2))%
}

\NovelDisplay{M02625R01}{%
(0)(1,1,1)(2,1,1)(3,1)(2,1,1)(2,1)(3,2,1)(4,2,1)(5,2)(4)%
}

\NovelDisplay{M02626R01}{%
= \psi(I\times \omega+I)%
}

\NovelDisplay{M02627R01}{%
= \psi(M^2+\psi_{\psi_M(M^2+M)}(1)+\psi_M(M^2))%
}

蔡须坤把本子移到较干的一侧，薄薄的领域只在几人身周留了一会儿。各人的魂力不再因急乱而互相牵扯，肩上的压力却仍旧各自压着。

她没有让那层光陪着所有人向上走。诺当缓过来后，光便收回，她也重新从自己的那一式接起。刚才停下的地方依旧要核，没有谁因为帮过同伴，就能省掉这一遍。

诺当歇过以后，先去摸了摸自己滑倒时被绳子勒住的那一处，才坐回本子前。他原本打算不再问了，免得大家一直在这儿等，薛峰却已经将刚才未完的纸放到他手边。

``你停的是这儿，不是整本。''

诺当看了很久，终于将最外面的组织说出来。说到里面仍然卡，便把指尖停在卡住的位置，不再拿一句都差不多糊过去。

卢本韦听见，便放心地回到自己的那页。后面有了能够继续的入口，她也不必一边向上读，一边担心是不是谁已经不敢再开口。霍雨浩替蔡须坤核过一处，薛峰接着补另一处，四人之间的检查开始能够真正轮换。

I加在后面，接着走向有限倍数，再走向乘ω的更深层次。卢本韦这一次没有只盯着新出现的那个倍数，而是回到整个乘法的位置。

她将自己的停顿说给霍雨浩听：``这里前面的背景变了，后面能加的东西也得跟着看。不能抓住一个乘二，就认作前面同一行再抄两遍。''

霍雨浩让她接着解释右侧。她在M侧卡住，蔡须坤便指着其中同一个函数值，念到有限重复的地方，再念到下一行参数发生变化的地方。哪一处是已有整项的有限倍数，哪一处已经进入更深的参数组织，两人终于能够互相提醒。

薛峰没有等她们说完才抄。他先把共同的BMS核好，依每条式子实际的层次，给两个本子留出相同的对照位置。左列较长的一行跨到了下页，他把页角折出一个很小的记号，使后来的人不至于把下一页误当成新的起点。

蔡兴坤翻到那一处，果然先想补零，看到记号才停。她抬头向薛峰笑了一下，随即自己把那一列从头接完。

这一段没有谁摔倒，却比刚才还磨人。许多地方只错一点，整行念起来并不会立刻卡住。卢本韦念到最后，才被霍雨浩请回起点，肩背都像跟着沉了一遍。

她原想说自己刚才只是口误，看到纸上同样漏掉的一层，又将这句话咽下去。她擦去旧行，重新写全，再把解释交给薛峰来核。

薛峰核得很慢，却没有因为她是队里先前讲得快的人就放过那一处。卢本韦抬眼看他，最后倒先笑了：``行，就这么查。''

\NovelDisplay{M02628R01}{%
(0)(1,1,1)(2,1,1)(3,1)(2,1,1)(2,1)(3,2,1)(4,2,1)(5,2)(4,2)(5,3)%
}

\NovelDisplay{M02629R01}{%
= \psi(I\times \omega+I+\Omega_{\psi_{I}(I\times \omega + I)+1})%
}

\NovelDisplay{M02630R01}{%
= \psi(M^2+\psi_{\psi_M(M^2+M)}(1)+\psi_M(M^2)+\psi_{\psi_M(M^2)}(M^2+\psi_{\psi_M(M^2+M)}(1)%
}
\NovelDisplay{M02630R02}{%
+\psi_M(M^2)+1))%
}

\NovelDisplay{M02631R01}{%
(0)(1,1,1)(2,1,1)(3,1)(2,1,1)(2,1)(3,2,1)(4,2,1)(5,2)(4,2)(5,3,1)%
}
\NovelDisplay{M02631R02}{%
(6,3,1)(7,3)(6)%
}

\NovelDisplay{M02632R01}{%
= \psi(I\times \omega+I\times 2)%
}

\NovelDisplay{M02633R01}{%
= \psi(M^2+\psi_{\psi_M(M^2+M)}(1)+\psi_M(M^2)\times 2)%
}

\NovelDisplay{M02634R01}{%
(0)(1,1,1)(2,1,1)(3,1)(2,1,1)(2,1)(3,2,1)(4,2,1)(5,2)(4,2,1)%
}

\NovelDisplay{M02635R01}{%
= \psi(I\times \omega \times 2)%
}

\NovelDisplay{M02636R01}{%
= \psi(M^2+\psi_{\psi_M(M^2+M)}(1)\times 2)%
}

\NovelDisplay{M02637R01}{%
(0)(1,1,1)(2,1,1)(3,1)(2,1,1)(2,1)(3,2,1)(4,2,1)(5,2)(4,2,1)(4,2)(5,3,1)%
}
\NovelDisplay{M02637R02}{%
(6,3,1)(7,3)(6,3,1)%
}

\NovelDisplay{M02638R01}{%
= \psi(I\times \omega \times 3)%
}

\NovelDisplay{M02639R01}{%
= \psi(M^2+\psi_{\psi_M(M^2+M)}(1)\times 3)%
}

\NovelDisplay{M02640R01}{%
(0)(1,1,1)(2,1,1)(3,1)(2,1,1)(2,1,1)%
}

\NovelDisplay{M02641R01}{%
= \psi(I\times \omega^2)%
}

\NovelDisplay{M02642R01}{%
= \psi(M^2+\psi_{\psi_M(M^2+M)}(2))%
}

\NovelDisplay{M02643R01}{%
(0)(1,1,1)(2,1,1)(3,1)(2,1,1)(2,1,1)(2,1,1)%
}

\NovelDisplay{M02644R01}{%
= \psi(I\times \omega^3)%
}

\NovelDisplay{M02645R01}{%
= \psi(M^2+\psi_{\psi_M(M^2+M)}(3))%
}

\NovelDisplay{M02646R01}{%
(0)(1,1,1)(2,1,1)(3,1)(2,1,1)(3,1)%
}

\NovelDisplay{M02647R01}{%
= \psi(I\times \Omega)%
}

\NovelDisplay{M02648R01}{%
= \psi(M^2+\psi_{\psi_M(M^2+M)}(M))%
}

雨后的石阶亮了一层薄光。走到能站稳的地方，霍雨浩先回头，让后面的几人各自说出今天最容易误认的一处。

答案并不一样。有人怕漏列，有人会把有限倍数说进幂次，也有人只要看到重复的函数头，就记不起自己在哪个参数里。几人将这些差别留下，下次互核时，终于知道该问对方什么。

当乘号后面的因子本身又是一层ψ时，卢本韦也慢得和蔡须坤差不多了。

她能读清I侧的最外层，往里一进，却总会把因子内部的变化与外面添尾部的动作混在一起。她原先依靠较短写法维持的速度，到这里暂时停住。

蔡须坤的M侧没能替她立刻解围。新结构本就更长，前面的背景又必须完整保留，两人读到不同位置停下来，反倒一时不知道先帮谁。

``先把同一条左列找到。''霍雨浩说。

他将两人的纸对齐，让她们各自说出已经确定的最外层，再让薛峰从尾部核是否还是当前这一式。共同的左列确认以后，M侧的某一层才有了能够与I侧对照的位置。

她们每人只解释自己真能说明的那一截。另一人听懂，再往外接。霍雨浩听着，发现自己留在页边的一处旧记法也没有跟上当前参数，便当着几人的面停笔重写。

薛峰没有立即去接后半行，反倒指住了那一处：``前面那份参数，也得带回来。''

霍雨浩看过，点头。

这一下让薛峰自己也安静了片刻。他先前检查的，多半是旁人的抄写；现在连讲得最清楚的人，也会在疲惫时漏掉来路。核对不是谁永远给谁打勾，谁先看清，谁就该将这一处守住。

后面的唐三接过已核好的纸，却没直接照着向上走。他先让诺当把其中最外面的一层说出来，自己再解释尾部。古月娜坐到两人旁边，替他们守着当前所用的那一页，直到他们能不翻旧纸也说清为止。

\NovelDisplay{M02649R01}{%
(0)(1,1,1)(2,1,1)(3,1)(2,1,1)(3,1)(1,1,1)(2,1,1)(3,1)(2)%
}

\NovelDisplay{M02650R01}{%
= \psi(I\times \psi_I(I))%
}

\NovelDisplay{M02651R01}{%
= \psi(M^2+\psi_{\psi_M(M^2+M)}(M\times \psi_{\psi_M(M^2)}(M^2)))%
}

\NovelDisplay{M02652R01}{%
(0)(1,1,1)(2,1,1)(3,1)(2,1,1)(3,1)(1,1,1)(2,1,1)(3,1)(2,1)(3,2,1)(4,2,1)%
}
\NovelDisplay{M02652R02}{%
(5,2)(4,2,1)(5,1)(3,2)%
}

\NovelDisplay{M02653R01}{%
= \psi(I\times \psi_I(I)+\Omega_{\psi_I(I)+1})%
}

\NovelDisplay{M02654R01}{%
= \psi(M^2+\psi_{\psi_M(M^2+M)}(M\times \psi_{\psi_M(M^2)}(M^2))+\psi_{\psi_M(M^2)}(M^2+1))%
}

\NovelDisplay{M02655R01}{%
(0)(1,1,1)(2,1,1)(3,1)(2,1,1)(3,1)(1,1,1)(2,1,1)(3,1)(2,1)(3,2,1)(4,2,1)%
}
\NovelDisplay{M02655R02}{%
(5,2)(4,2,1)(5,1)(3,2,1)(4,2,1)(5,2)(4)%
}

\NovelDisplay{M02656R01}{%
= \psi(I\times \psi_I(I)+\psi_I(I\times 2))%
}

\NovelDisplay{M02657R01}{%
= \psi(M^2+\psi_{\psi_M(M^2+M)}(M\times \psi_{\psi_M(M^2)}(M^2))+\psi_{\psi_M(M^2)}(M^2+\psi_M(M^2)))%
}

\NovelDisplay{M02658R01}{%
(0)(1,1,1)(2,1,1)(3,1)(2,1,1)(3,1)(1,1,1)(2,1,1)(3,1)(2,1)(3,2,1)(4,2,1)%
}
\NovelDisplay{M02658R02}{%
(5,2)(4,2,1)(5,1)(3,2,1)(4,2,1)(5,2)(4,2,1)%
}

\NovelDisplay{M02659R01}{%
= \psi(I\times \psi_I(I)+\psi_I(I\times \omega))%
}

\NovelDisplay{M02660R01}{%
=\! \psi(M^2\!+\!\psi_{\psi_M(M^2+M)}(M\!\times\! \psi_{\psi_M(M^2)}(M^2))\!+\!\psi_{\psi_M(M^2)}(M^2\!+\!\psi_{\psi_M(M^2+M)}(1)))%
}

\NovelDisplay{M02661R01}{%
(0)(1,1,1)(2,1,1)(3,1)(2,1,1)(3,1)(1,1,1)(2,1,1)(3,1)(2,1)(3,2,1)(4,2,1)%
}
\NovelDisplay{M02661R02}{%
(5,2)(4,2,1)(5,1)(3,2,1)(4,2,1)(5,2)(4,2,1)(5,1)(1,1,1)(2,1,1)(3,1)(2)%
}

\NovelDisplay{M02662R01}{%
= \psi(I\times \psi_I(I)+\psi_I(I\times \psi_I(I)))%
}

\NovelDisplay{M02663R01}{%
= \psi(M^2+\psi_{\psi_M(M^2+M)}(M\times \psi_{\psi_M(M^2)}(M^2))+\psi_{\psi_M(M^2)}(M^2+\psi_{\psi_M(M^2+M)}(M%
}
\NovelDisplay{M02663R02}{%
\times \psi_{\psi_M(M^2)}(M^2))))%
}

\NovelDisplay{M02664R01}{%
(0)(1,1,1)(2,1,1)(3,1)(2,1,1)(3,1)(1,1,1)(2,1,1)(3,1)(2,1)(3,2,1)(4,2,1)%
}
\NovelDisplay{M02664R02}{%
(5,2)(4,2,1)(5,1)(4,2)(5,3)%
}

\NovelDisplay{M02665R01}{%
= \psi(I\times \psi_I(I)+\Omega_{\psi_I(I\times \psi_I(I))+1})%
}

\NovelDisplay{M02666R01}{%
= \psi(M^2+\psi_{\psi_M(M^2+M)}(M\times \psi_{\psi_M(M^2)}(M^2))+\psi_{\psi_M(M^2)}(M^2+\psi_{\psi_M(M^2+M)}(M%
}
\NovelDisplay{M02666R02}{%
\times \psi_{\psi_M(M^2)}(M^2))+1))%
}

\NovelDisplay{M02667R01}{%
(0)(1,1,1)(2,1,1)(3,1)(2,1,1)(3,1)(1,1,1)(2,1,1)(3,1)(2,1)(3,2,1)(4,2,1)%
}
\NovelDisplay{M02667R02}{%
(5,2)(4,2,1)(5,1)(4,2)(5,3,1)(6,3,1)(7,3)(6)%
}

\NovelDisplay{M02668R01}{%
= \psi(I\times \psi_I(I)+I)%
}

\NovelDisplay{M02669R01}{%
= \psi(M^2+\psi_{\psi_M(M^2+M)}(M\times \psi_{\psi_M(M^2)}(M^2))+\psi_M(M^2))%
}

\NovelDisplay{M02670R01}{%
(0)(1,1,1)(2,1,1)(3,1)(2,1,1)(3,1)(1,1,1)(2,1,1)(3,1)(2,1)(3,2,1)(4,2,1)%
}
\NovelDisplay{M02670R02}{%
(5,2)(4,2,1)(5,1)(4,2)(5,3,1)(6,3,1)(7,3)(6,3,1)%
}

\NovelDisplay{M02671R01}{%
= \psi(I\times \psi_I(I)+I\times \omega)%
}

\NovelDisplay{M02672R01}{%
= \psi(M^2+\psi_{\psi_M(M^2+M)}(M\times \psi_{\psi_M(M^2)}(M^2))+\psi_{\psi_M(M^2+M)}(1))%
}

\NovelDisplay{M02673R01}{%
(0)(1,1,1)(2,1,1)(3,1)(2,1,1)(3,1)(1,1,1)(2,1,1)(3,1)(2,1)(3,2,1)(4,2,1)%
}
\NovelDisplay{M02673R02}{%
(5,2)(4,2,1)(5,1)(4,2)(5,3,1)(6,3,1)(7,3)(6,3,1)(7,1)(1,1,1)(2,1,1)%
}
\NovelDisplay{M02673R03}{%
(3,1)(2)%
}

\NovelDisplay{M02674R01}{%
= \psi(I\times \psi_I(I)\times 2)%
}

\NovelDisplay{M02675R01}{%
= \psi(M^2+\psi_{\psi_M(M^2+M)}(M\times \psi_{\psi_M(M^2)}(M^2)) \times 2)%
}

\NovelDisplay{M02676R01}{%
(0)(1,1,1)(2,1,1)(3,1)(2,1,1)(3,1)(1,1,1)(2,1,1)(3,1)(2,1)(3,2,1)(4,2,1)%
}
\NovelDisplay{M02676R02}{%
(5,2)(4,2,1)(5,1)(4,2)(5,3,1)(6,3,1)(7,3)(6,3,1)(7,1)(6,3)(7,4,1)%
}
\NovelDisplay{M02676R03}{%
(8,4,1)(9,4)(8)%
}

\NovelDisplay{M02677R01}{%
= \psi(I\times \psi_I(I)\times 2+I)%
}

\NovelDisplay{M02678R01}{%
= \psi(M^2+\psi_{\psi_M(M^2+M)}(M\times \psi_{\psi_M(M^2)}(M^2)) \times 2 + \psi_M(M^2))%
}

\NovelDisplay{M02679R01}{%
(0)(1,1,1)(2,1,1)(3,1)(2,1,1)(3,1)(1,1,1)(2,1,1)(3,1)(2,1)(3,2,1)(4,2,1)%
}
\NovelDisplay{M02679R02}{%
(5,2)(4,2,1)(5,1)(4,2)(5,3,1)(6,3,1)(7,3)(6,3,1)(7,1)(6,3)(7,4,1)(8,4,1)%
}
\NovelDisplay{M02679R03}{%
(9,4)(8,4,1)(9)%
}

\NovelDisplay{M02680R01}{%
= \psi(I\times \psi_I(I)\times 3)%
}

\NovelDisplay{M02681R01}{%
= \psi(M^2+\psi_{\psi_M(M^2+M)}(M\times \psi_{\psi_M(M^2)}(M^2)) \times 3)%
}

当天收纸时，薛峰特意将这一组的左列多描了一遍。它们长得太像，后来一旦在别处再出现，就要先认清是不是同一条输入。

他还不知道这份准备很快会派上用场。霍雨浩看见他在页边补的目录，便将其余几份纸也交来，让他统一核过。

霍雨浩后来将自己白天的错处留给薛峰再查了一次。

薛峰原以为他会把旧行擦掉，见仍在，便把自己核过的来路也写在旁边。这里并没有因为讲解者说得顺，就免去对参数的追查；疲劳时最容易留下的，恰是那些过去许多次都没有读错的熟悉部分。

蔡须坤看过两人的纸，把自己最深的一处也留出来。卢本韦则要求他们反过来查她的I侧，哪怕她正读得很快，也照样拦。

奥斯卡送补给来时，看见四个人各拿着别人的一页，找了一会儿才弄清本子是谁的。他没催他们收，只把能留到明天的那份分开，让他们先将手里这口吃完。

更深的因子开始反复出现，整体仍留在乘号后面。

蔡须坤最吃力的那一行，几乎读到气息都乱了。她停下来，先将手从纸上撤开，等腕间的魂力重新顺畅接续，再回来找刚才读过的出口。

卢本韦没有让她重新从第一字开始。她把I侧整体读给她听，让她先确认当前这一行的外层究竟组织成什么样。随后，M侧才从同一个位置重新展开。

``如果先看完整个右栏，我会觉得哪里都要记。''蔡须坤说。

``那先把外面的留住，里面一次只补一处。''

她们用了很长时间才做到。添在尾部的一小段，与因子内部换入的新参数，纸面上可能用到相同的I，位置却必须一处处认。某些更短的右式，只能作为这一条已经给出的对应来读，还不能替代她们对完整左列的核查。

薛峰在这一段没有再只是站在旁边。他先独立写I侧，再请霍雨浩遮去答案，重新说明同一个BMS在哪一列与前式发生差别。答案不见以后，最外面的路竟仍能留在脑中，他这才敢把纸递给蔡须坤，试着去核M侧相同位置的作用。

他还不能像她那样读完所有深层函数，蔡须坤也没有让他装作能读。她将自己已说明的部分接给他，其余仍由自己从原式往下核。两人的帮助落在明确的位置，不必靠谁将整页都背下来维持。

天色暗后，夏夜星空请他们收纸。没有人争着再向上试一次。霍雨浩起身时，才发现自己按在膝上的那只手僵了很久，薛峰将本子接过去，让他先活动手指。

\NovelDisplay{M02682R01}{%
(0)(1,1,1)(2,1,1)(3,1)(2,1,1)(3,1)(1,1,1)(2,1,1)(3,1)(2,1)(3,2,1)(4,2,1)%
}
\NovelDisplay{M02682R02}{%
(5,2)(4,2,1)(5,1)(4,2,1)%
}

\NovelDisplay{M02683R01}{%
= \psi(I\times \psi_I(I)\times \omega)%
}

\NovelDisplay{M02684R01}{%
= \psi(M^2+\psi_{\psi_M(M^2+M)}(M\times \psi_{\psi_M(M^2)}(M^2)+1))%
}

\NovelDisplay{M02685R01}{%
(0)(1,1,1)(2,1,1)(3,1)(2,1,1)(3,1)(1,1,1)(2,1,1)(3,1)(2,1)(3,2,1)(4,2,1)%
}
\NovelDisplay{M02685R02}{%
(5,2)(4,2,1)(5,1)(4,2,1)(5,1)%
}

\NovelDisplay{M02686R01}{%
= \psi(I\times \psi_I(I)\times \Omega)%
}

\NovelDisplay{M02687R01}{%
= \psi(M^2+\psi_{\psi_M(M^2+M)}(M\times \psi_{\psi_M(M^2)}(M^2)+\psi_{\psi_M(M^2)}(1)))%
}

\NovelDisplay{M02688R01}{%
(0)(1,1,1)(2,1,1)(3,1)(2,1,1)(3,1)(1,1,1)(2,1,1)(3,1)(2,1)(3,2,1)(4,2,1)%
}
\NovelDisplay{M02688R02}{%
(5,2)(4,2,1)(5,1)(4,2,1)(5,1)(4,2,1)%
}

\NovelDisplay{M02689R01}{%
= \psi(I\times \psi_{\Omega_{\psi_I(I)+1}}(I+1))%
}

\NovelDisplay{M02690R01}{%
= \psi(M^2+\psi_{\psi_M(M^2+M)}(M\times \psi_{\psi_M(M^2)}(M^2)+\psi_{\psi_M(M^2)}(M^2)+1))%
}

\NovelDisplay{M02691R01}{%
(0)(1,1,1)(2,1,1)(3,1)(2,1,1)(3,1)(1,1,1)(2,1,1)(3,1)(2,1)(3,2,1)(4,2,1)%
}
\NovelDisplay{M02691R02}{%
(5,2)(4,2,1)(5,2)%
}

\NovelDisplay{M02692R01}{%
= \psi(I\times \Omega_{\psi_I(I)+1})%
}

\NovelDisplay{M02693R01}{%
= \psi(M^2+\psi_{\psi_M(M^2+M)}(M\times \psi_{\psi_M(M^2)}(M^2+1)))%
}

\NovelDisplay{M02694R01}{%
(0)(1,1,1)(2,1,1)(3,1)(2,1,1)(3,1)(1,1,1)(2,1,1)(3,1)(2,1)(3,2,1)(4,2,1)%
}
\NovelDisplay{M02694R02}{%
(5,2)(4,2,1)(5,2)(3,2,1)(4,2,1)(5,2)(4)%
}

\NovelDisplay{M02695R01}{%
= \psi(I\times \psi_I(I\times 2))%
}

\NovelDisplay{M02696R01}{%
= \psi(M^2+\psi_{\psi_M(M^2+M)}(M\times \psi_{\psi_M(M^2)}(M^2+\psi_M(M^2))))%
}

\NovelDisplay{M02697R01}{%
(0)(1,1,1)(2,1,1)(3,1)(2,1,1)(3,1)(1,1,1)(2,1,1)(3,1)(2,1)(3,2,1)(4,2,1)%
}
\NovelDisplay{M02697R02}{%
(5,2)(4,2,1)(5,2)(3,2,1)(4,2,1)(5,2)(4,2)(5,3,1)(6,3,1)(7,3)(6,3,1)%
}
\NovelDisplay{M02697R03}{%
(7,2)(6,3)%
}

\NovelDisplay{M02698R01}{%
= \psi(I\times \psi_I(I\times 2)+\Omega_{\psi_I(I\times 2) + 1})%
}

\NovelDisplay{M02699R01}{%
= \psi(M^2+\psi_{\psi_M(M^2+M)}(M\times \psi_{\psi_M(M^2)}(M^2+\psi_M(M^2)))+\psi_{\psi_{M}(M^2)}(M^2%
}
\NovelDisplay{M02699R02}{%
+\psi_M(M^2)+1))%
}

\NovelDisplay{M02700R01}{%
(0)(1,1,1)(2,1,1)(3,1)(2,1,1)(3,1)(1,1,1)(2,1,1)(3,1)(2,1)(3,2,1)(4,2,1)%
}
\NovelDisplay{M02700R02}{%
(5,2)(4,2,1)(5,2)(3,2,1)(4,2,1)(5,2)(4,2)(5,3,1)(6,3,1)(7,3)(6,3,1)(7,3)%
}

\NovelDisplay{M02701R01}{%
= \psi(I\times \Omega_{\psi_I(I\times 2) + 1})%
}

\NovelDisplay{M02702R01}{%
= \psi(M^2+\psi_{\psi_M(M^2+M)}(M\times \psi_{\psi_M(M^2)}(M^2+\psi_M(M^2)+1)))%
}

\NovelDisplay{M02703R01}{%
(0)(1,1,1)(2,1,1)(3,1)(2,1,1)(3,1)(1,1,1)(2,1,1)(3,1)(2,1)(3,2,1)(4,2,1)%
}
\NovelDisplay{M02703R02}{%
(5,2)(4,2,1)(5,2)(3,2,1)(4,2,1)(5,2)(4,2)(5,3,1)(6,3,1)(7,3)(6,3,1)(7,3)%
}
\NovelDisplay{M02703R03}{%
(5,3,1)(6,3,1)(7,3)(6,3)(7,4,1)(8,4,1)(9,4)(8,4,1)(9,3)(8,4)%
}

\NovelDisplay{M02704R01}{%
= \psi(I\times \Omega_{\psi_I(I\times 3) + 1})%
}

\NovelDisplay{M02705R01}{%
= \psi(M^2+\psi_{\psi_M(M^2+M)}(M\times \psi_{\psi_M(M^2)}(M^2+\psi_M(M^2)\times 2+1)))%
}

那晚，他们把新写的I侧与M侧分开存好，再用共同的BMS页号接在一起。

卢本韦不愿让两本越翻越乱，蔡须坤也不肯把M侧缩成只够自己看懂的几个暗号。薛峰替两人留了能真正对应起来的入口，霍雨浩再逐一核过。灯下的工作很慢，第二天上山时，四人却第一次没有先花半个时辰找昨晚丢失的位置。

最后几式把变化推进到更深的参数里，随后，I侧写出了I平方。

卢本韦看着这个短得近乎轻巧的结果，很久没有出声。她的膝盖仍在发抖，纸上却只有这么一点。

``别先只念结果。''薛峰说。

她回头看了他一眼，依言将前面的对应从左到右读过。到了M侧，蔡须坤接上自己的那一栏。那些几天来反复读不完的函数头仍在，最深的一段也仍要放回完整括号里，它们没有被这个短结果抹去。

霍雨浩随后请两人交换问法。蔡须坤要卢本韦说明I平方出现在哪一层，卢本韦则请她指出M侧当前函数的下标与参数分别到哪里结束。薛峰核左列，古月娜再请后面的几人轮流解释自己能看清的一处。

一轮结束，蔡须坤才将脚向上移。

她没有觉得肩上的重量骤然减轻。真正变得清楚的，是先前一受压就混作一团的几处接续。她能先守住自己已经认出的外层，再将内部那一段慢慢接回。卢本韦紧跟着站稳，霍雨浩将本子收起，薛峰也没有等人提醒便说出了自己的核查位置。

后面的人还需要时间。四人站在已经能够稳定停住的地方，分别接住不同的问话，没有用同一个答案把所有人的困难一并盖过。

古月娜最后陪着蔡兴坤把那一行接完，抬头时，二百八十米的横台已从雾中露了出来。她向前面四人点了一下头，才与其余人继续往上走。

\NovelDisplay{M02706R01}{%
(0)(1,1,1)(2,1,1)(3,1)(2,1,1)(3,1)(1,1,1)(2,1,1)(3,1)(2,1,1)%
}

\NovelDisplay{M02707R01}{%
= \psi(I\times \psi_I(I\times \omega))%
}

\NovelDisplay{M02708R01}{%
= \psi(M^2+\psi_{\psi_M(M^2+M)}(M\times \psi_{\psi_M(M^2)}(M^2+\psi_{\psi_M(M^2+M)}(1))))%
}

\NovelDisplay{M02709R01}{%
(0)(1,1,1)(2,1,1)(3,1)(2,1,1)(3,1)(1,1,1)(2,1,1)(3,1)(2,1,1)(2,1,1)%
}

\NovelDisplay{M02710R01}{%
= \psi(I\times \psi_I(I\times \omega^2))%
}

\NovelDisplay{M02711R01}{%
= \psi(M^2+\psi_{\psi_M(M^2+M)}(M\times \psi_{\psi_M(M^2)}(M^2+\psi_{\psi_M(M^2+M)}(2))))%
}

\NovelDisplay{M02712R01}{%
(0)(1,1,1)(2,1,1)(3,1)(2,1,1)(3,1)(1,1,1)(2,1,1)(3,1)(2,1,1)(3,1)%
}

\NovelDisplay{M02713R01}{%
= \psi(I\times\psi_I(I\times \Omega))%
}

\NovelDisplay{M02714R01}{%
= \psi(M^2+\psi_{\psi_M(M^2+M)}(M\times \psi_{\psi_M(M^2)}(M^2+\psi_{\psi_M(M^2+M)}(M))))%
}

\NovelDisplay{M02715R01}{%
(0)(1,1,1)(2,1,1)(3,1)(2,1,1)(3,1)(1,1,1)(2,1,1)(3,1)(2,1,1)(3,1)(1,1,1)%
}
\NovelDisplay{M02715R02}{%
(2,1,1)(3,1)(2)%
}

\NovelDisplay{M02716R01}{%
= \psi(I\times\psi_I(I\times \psi_I(I)))%
}

\NovelDisplay{M02717R01}{%
=\! \psi(M^2\!+\!\psi_{\psi_M(M^2+M)}(M\!\times\! \psi_{\psi_M(M^2)}(M^2\!+\!\psi_{\psi_M(M^2+M)}(M\!\times\! \psi_{\psi_M(M^2)}(M^2)))))%
}

\NovelDisplay{M02718R01}{%
(0)(1,1,1)(2,1,1)(3,1)(2,1,1)(3,1)(1,1,1)(2,1,1)(3,1)(2,1,1)(3,1)(1,1,1)%
}
\NovelDisplay{M02718R02}{%
(2,1,1)(3,1)(2,1,1)%
}

\NovelDisplay{M02719R01}{%
= \psi(I\times\psi_I(I\times \psi_I(I\times \omega)))%
}

\NovelDisplay{M02720R01}{%
= \psi(M^2+\psi_{\psi_M(M^2+M)}(M\times \psi_{\psi_M(M^2)}(M^2+\psi_{\psi_M(M^2+M)}(M\times \psi_{\psi_M(M^2)}(M^2%
}
\NovelDisplay{M02720R02}{%
+\psi_{\psi_{M}(M^2+M)}(1))))))%
}

\NovelDisplay{M02721R01}{%
(0)(1,1,1)(2,1,1)(3,1)(2,1,1)(3,1)(1,1,1)(2,1,1)(3,1)(2,1,1)(3,1)(1,1,1)%
}
\NovelDisplay{M02721R02}{%
(2,1,1)(3,1)(2,1,1)(3,1)%
}

\NovelDisplay{M02722R01}{%
= \psi(I\times\psi_I(I\times \psi_I(I\times \Omega)))%
}

\NovelDisplay{M02723R01}{%
= \psi(M^2+\psi_{\psi_M(M^2+M)}(M\times \psi_{\psi_M(M^2)}(M^2+\psi_{\psi_M(M^2+M)}(M\times \psi_{\psi_M(M^2)}(M^2%
}
\NovelDisplay{M02723R02}{%
+\psi_{\psi_{M}(M^2+M)}(M))))))%
}

\NovelDisplay{M02724R01}{%
(0)(1,1,1)(2,1,1)(3,1)(2,1,1)(3,1)(2)%
}

\NovelDisplay{M02725R01}{%
= \psi(I^2)%
}

\NovelDisplay{M02726R01}{%
= \psi(M^2+\psi_{\psi_M(M^2+M)}(M^2))%
}

十一人都到了横台上，绳子才被唐三收好。

卢本韦把自己坐的位置挪出一半，蔡须坤靠过来，连说话也懒得再分长句。薛峰仍看着刚写完的末行，嘴唇轻轻动着，像在确认它确实能不靠旁人的提示留在心里。

霍雨浩听了一阵，没有打断。他只在薛峰把完整的一遍说完以后，将水递过去。

这次，薛峰接水时没有撞到纸角。

从这里往下看，二百六十米的横台被山壁遮住，只剩绳子绕过的那一截石柱。

诺当说起自己滑倒的地方，仍有些不好意思。卢本韦摆手，叫他先把刚才真能自己讲清的那一处再留一份，回去替别的人核。

``下回我也会漏。''她说，``你看出来就拦。''

诺当点头，把那根细枝从本子里取出来。小舞将纸铺平，等他开始。

蔡须坤坐在另一边，看着M册从最早几条短式翻到现在，才发现自己仍有许多地方需要同伴提醒。可这些提醒已经能落到具体的某一层，不再只有一句笼统的``我没懂''。

她将那口气慢慢放下，手背贴了贴身后熟悉的温热。那点属于残魂的暖意安安静静贴着，等她的呼吸重新放平。

二百八十米写进山图时，四个人的纸先被收齐。随后，其他人的解释也一份份接上。山路没有因此变轻，横台上却再没有谁独自抱着一整页不知该向谁开口。

\chapter{第一百一十六章：三百米}\label{ux7b2cux4e00ux767eux4e00ux5341ux516dux7ae0ux4e09ux767eux7c73}

二百八十米之后，薛峰没有立刻翻过旧页。

新讲义的头几行放在他膝上，他却将前一天收好的两本册子重新打开，一页页往回找。卢本韦以为他又漏了某一处函数头，正要接话，便看见他已经把两条完整的左列并排放到了纸上。

``这条我们读过。''他说。

蔡须坤顺着看了一遍，确实一列也不差。她去核M侧，卢本韦看I侧，霍雨浩则将上一段的目录从头找来。

后面的人暂时坐下。唐三还没将绳结向上挪，古月娜也没有催前面的人先去试。几天以来，大家都盼着早一点看见三百米，此刻一处很小的疑问，却把十一人一同留在了原地。

薛峰将炭条放稳，手没有再像最初那样急着收回来。他知道自己还不能证明哪一种写法合适，但同一条左列出现过，就是出现过，不能因为新页摆在后面，就当它从未见过。

他们找出了三处需要另行核对的对应。

这一页的第一行，与旧册里I乘一个ψ值的那一行，有着相同的完整BMS；第二行也能在旧册的另一处乘法结构中找到。第四行的左列，更正是上一段末尾才用来认清I平方的原式。

可新页给它们写下的I侧表达，分别换成了带I平方的另一种结构。

卢本韦原来已经要照着新页读，听到这里，慢慢把手移开。``那不能把它们当成又往上走了三步。''

霍雨浩点头。右端写得不同，不能只凭字面便断言数值一定不同；坍缩表示中还会遇到规范化与同值写法。可讲义没有在这里给出足够说明，他们也没有当场完成相应证明，便不能悄悄把这种关系当成已经讲清的通用结论。

三行仍原样保留，旁边另记``同一输入的异文，待核对''。夹在其中的另一条，也按自己的左列单独核，没被这处问题一并抹掉。

蔡须坤最先有些不甘。她为了M侧那几层函数花了很久，眼下若留下一处未定，总像前面白走了一截。卢本韦把两本推到她面前，叫她看自己昨天已经能够独立解释的那些页。

``能说清的还在。没说清的，也得让它留着。''

薛峰将三处左列各自对应的旧页找齐，放在新页旁边。以前他帮大家核对换行，是怕抄错；现在他需要守住的是，同一个输入不能只因为讲义换了一页，就被当成另一道已经解决的新题。

夏夜星空看过他们留下的标记，没有要求他们为赶路将它擦掉。她让几人先把确实能够说明的结构整理好，待核的关系另留一处，不拿它来证明脚下的新台阶已经通过。

\NovelDisplay{M02727R01}{%
(0)(1,1,1)(2,1,1)(3,1)(2,1,1)(3,1)(1,1,1)(2,1,1)(3,1)(2)%
}

\NovelDisplay{M02728R01}{%
= \psi(I^2 + \psi_I(I))%
}

\NovelDisplay{M02729R01}{%
= \psi(M^2+\psi_{\psi_M(M^2+M)}(M^2) + \psi_{\psi_M(M^2)}(M^2))%
}

\NovelDisplay{M02730R01}{%
(0)(1,1,1)(2,1,1)(3,1)(2,1,1)(3,1)(1,1,1)(2,1,1)(3,1)(2,1,1)%
}

\NovelDisplay{M02731R01}{%
= \psi(I^2 + \psi_I(I\times \omega))%
}

\NovelDisplay{M02732R01}{%
= \psi(M^2+\psi_{\psi_M(M^2+M)}(M^2) + \psi_{\psi_M(M^2)}(M^2 + \psi_{\psi_M(M^2+M)}(1)))%
}

\NovelDisplay{M02733R01}{%
(0)(1,1,1)(2,1,1)(3,1)(2,1,1)(3,1)(1,1,1)(2,1,1)(3,1)(2,1,1)(3,1)(2)%
}

\NovelDisplay{M02734R01}{%
= \psi(I^2 + \psi_I(I^2))%
}

\NovelDisplay{M02735R01}{%
= \psi(M^2+\psi_{\psi_M(M^2+M)}(M^2) + \psi_{\psi_M(M^2)}(M^2 + \psi_{\psi_M(M^2+M)}(M^2)))%
}

\NovelDisplay{M02736R01}{%
(0)(1,1,1)(2,1,1)(3,1)(2,1,1)(3,1)(2)%
}

\NovelDisplay{M02737R01}{%
= \psi(I^2 + \psi_I(I^2)\times \omega)%
}

\NovelDisplay{M02738R01}{%
= \psi(M^2+\psi_{\psi_M(M^2+M)}(M^2) + \psi_{\psi_{\psi_M(M^2)}(M^2 + \psi_{\psi_M(M^2+M)}(M^2)+1)}(1))%
}

蔡须坤将笔从那三处移开，深吸一口气，翻到下一组。

薛峰没有得到一个现成的新答案，却终于能将疑问完整交给同伴。霍雨浩把他找出的旧页夹在最前，留待以后核查。那份目录第一次成了所有人都会主动来借的东西。

接下来的结构，继续在I平方之后补入更深的尾部。

卢本韦先把每一行的背景读出，再确认加在后面的那一段归哪个括号。她已经不急着把所有相近的结果排成一条严格上升的直线；有些左列不同，右侧却给出相同表达，仍要保留各自的结构，按当前这条BMS来说明。

蔡须坤则在M侧留住整段背景。那里面不只有M平方，还包括前面才辨清的较深函数值。她读到尾部时，常会本能地想将前面省成几个自己认得的暗号，随后又把笔停住。

卢本韦看见，便问她前面省掉的是什么。

她第一次只说旧背景，卢本韦没应。第二次，她将具体的函数头与参数重新念全，卢本韦才把后面的I侧接过来。

``你要是连自己都暂时说不全，我就没法靠这个帮你查。''

蔡须坤抬头看她，最后将几个暗号全部擦去。M侧难的地方，本就有很大一部分藏在被她想省掉的位置。纸上能省多少字，与脑中有没有留住结构，是两件必须分开看的事。

薛峰也没有将几处相同的右侧直接合并。他把完整左列分别核过，指出相同部分在哪里结束，哪一列仍有差别，再把它们交给后面的人复述。

古月娜陪着蔡兴坤、小舞看同一处。她没有替前面四人下结论，而是让后面的几人把自己需要的那层先讲清。听懂以后，再把当前这条左列独立接回去，谁也不靠一句``前面已经会了''略过。

\NovelDisplay{M02739R01}{%
(0)(1,1,1)(2,1,1)(3,1)(2,1,1)(3,1)(2,1)(3,2)%
}

\NovelDisplay{M02740R01}{%
= \psi(I^2 + \Omega_{\psi_I(I^2)+1})%
}

\NovelDisplay{M02741R01}{%
= \psi(M^2+\psi_{\psi_M(M^2+M)}(M^2) + \psi_{\psi_M(M^2)}(M^2 + \psi_{\psi_M(M^2+M)}(M^2)+1))%
}

\NovelDisplay{M02742R01}{%
(0)(1,1,1)(2,1,1)(3,1)(2,1,1)(3,1)(2,1)(3,2,1)(4,2,1)(5,2)(4)%
}

\NovelDisplay{M02743R01}{%
= \psi(I^2 + I)%
}

\NovelDisplay{M02744R01}{%
= \psi(M^2+\psi_{\psi_M(M^2+M)}(M^2) + \psi_M(M^2))%
}

\NovelDisplay{M02745R01}{%
(0)(1,1,1)(2,1,1)(3,1)(2,1,1)(3,1)(2,1)(3,2,1)(4,2,1)(5,2)(4,2)(5,3,1)%
}
\NovelDisplay{M02745R02}{%
(6,3,1)(7,3)(6)%
}

\NovelDisplay{M02746R01}{%
= \psi(I^2 + I\times 2)%
}

\NovelDisplay{M02747R01}{%
= \psi(M^2+\psi_{\psi_M(M^2+M)}(M^2) + \psi_M(M^2)\times 2)%
}

\NovelDisplay{M02748R01}{%
(0)(1,1,1)(2,1,1)(3,1)(2,1,1)(3,1)(2,1)(3,2,1)(4,2,1)(5,2)(4,2,1)%
}

\NovelDisplay{M02749R01}{%
= \psi(I^2 + I\times \omega)%
}

\NovelDisplay{M02750R01}{%
= \psi(M^2+\psi_{\psi_M(M^2+M)}(M^2) + \psi_{\psi_M(M^2+M)}(1))%
}

\NovelDisplay{M02751R01}{%
(0)(1,1,1)(2,1,1)(3,1)(2,1,1)(3,1)(2,1)(3,2,1)(4,2,1)(5,2)(4,2,1)(5,1)%
}
\NovelDisplay{M02751R02}{%
(1,1,1)(2,1,1)(3,1)(2)%
}

\NovelDisplay{M02752R01}{%
= \psi(I^2 + I\times \psi_I(I))%
}

\NovelDisplay{M02753R01}{%
= \psi(M^2+\psi_{\psi_M(M^2+M)}(M^2) + \psi_{\psi_M(M^2+M)}(M \times \psi_{\psi_M(M^2)}(M^2)))%
}

\NovelDisplay{M02754R01}{%
(0)(1,1,1)(2,1,1)(3,1)(2,1,1)(3,1)(2,1)(3,2,1)(4,2,1)(5,2)(4,2,1)(5,1)%
}
\NovelDisplay{M02754R02}{%
(1,1,1)(2,1,1)(3,1)(2,1,1)%
}

\NovelDisplay{M02755R01}{%
= \psi(I^2 + I\times \psi_I(I\times \omega))%
}

\NovelDisplay{M02756R01}{%
= \psi(M^2+\psi_{\psi_M(M^2+M)}(M^2) + \psi_{\psi_M(M^2+M)}(M \times \psi_{\psi_M(M^2)}(M^2%
}
\NovelDisplay{M02756R02}{%
+\psi_{\psi_M(M^2+M)}(1))))%
}

\NovelDisplay{M02757R01}{%
(0)(1,1,1)(2,1,1)(3,1)(2,1,1)(3,1)(2,1)(3,2,1)(4,2,1)(5,2)(4,2,1)(5,1)%
}
\NovelDisplay{M02757R02}{%
(1,1,1)(2,1,1)(3,1)(2,1,1)(3,1)(2)%
}

\NovelDisplay{M02758R01}{%
= \psi(I^2 + I\times \psi_I(I^2))%
}

\NovelDisplay{M02759R01}{%
= \psi(M^2+\psi_{\psi_M(M^2+M)}(M^2) + \psi_{\psi_M(M^2+M)}(M \times \psi_{\psi_M(M^2)}(M^2%
}
\NovelDisplay{M02759R02}{%
+\psi_{\psi_M(M^2+M)}(M^2))))%
}

\NovelDisplay{M02760R01}{%
(0)(1,1,1)(2,1,1)(3,1)(2,1,1)(3,1)(2,1)(3,2,1)(4,2,1)(5,2)(4,2,1)(5,1)%
}
\NovelDisplay{M02760R02}{%
(1,1,1)(2,1,1)(3,1)(2,1,1)(3,1)(2)%
}

\NovelDisplay{M02761R01}{%
= \psi(I^2 + I\times \psi_I(I^2))%
}

\NovelDisplay{M02762R01}{%
= \psi(M^2+\psi_{\psi_M(M^2+M)}(M^2) + \psi_{\psi_M(M^2+M)}(M \times \psi_{\psi_M(M^2)}(M^2%
}
\NovelDisplay{M02762R02}{%
+\psi_{\psi_M(M^2+M)}(M^2))))%
}

\NovelDisplay{M02763R01}{%
(0)(1,1,1)(2,1,1)(3,1)(2,1,1)(3,1)(2,1)(3,2,1)(4,2,1)(5,2)(4,2,1)(5,2)%
}

\NovelDisplay{M02764R01}{%
= \psi(I^2 + I\times \Omega_{\psi_I(I^2)+1})%
}

\NovelDisplay{M02765R01}{%
= \psi(M^2+\psi_{\psi_M(M^2+M)}(M^2) + \psi_{\psi_M(M^2+M)}(M \times \psi_{\psi_M(M^2)}(M^2%
}
\NovelDisplay{M02765R02}{%
+\psi_{\psi_M(M^2+M)}(M^2)+1)))%
}

\NovelDisplay{M02766R01}{%
(0)(1,1,1)(2,1,1)(3,1)(2,1,1)(3,1)(2,1)(3,2,1)(4,2,1)(5,2)(4,2,1)(5,2)%
}
\NovelDisplay{M02766R02}{%
(3,2,1)(4,2,1)(5,2)(4)%
}

\NovelDisplay{M02767R01}{%
= \psi(I^2 + I\times \psi_I(I^2+I))%
}

\NovelDisplay{M02768R01}{%
= \psi(M^2+\psi_{\psi_M(M^2+M)}(M^2) + \psi_{\psi_M(M^2+M)}(M \times \psi_{\psi_M(M^2)}(M^2%
}
\NovelDisplay{M02768R02}{%
+\psi_{\psi_M(M^2+M)}(M^2)+\psi_M(M^2))))%
}

\NovelDisplay{M02769R01}{%
(0)(1,1,1)(2,1,1)(3,1)(2,1,1)(3,1)(2,1)(3,2,1)(4,2,1)(5,2)(4,2,1)(5,2)%
}
\NovelDisplay{M02769R02}{%
(3,2,1)(4,2,1)(5,2)(4,2,1)%
}

\NovelDisplay{M02770R01}{%
= \psi(I^2 + I\times \psi_I(I^2+I\times \omega))%
}

\NovelDisplay{M02771R01}{%
= \psi(M^2+\psi_{\psi_M(M^2+M)}(M^2) + \psi_{\psi_M(M^2+M)}(M \times \psi_{\psi_M(M^2)}(M^2%
}
\NovelDisplay{M02771R02}{%
+\psi_{\psi_M(M^2+M)}(M^2)+\psi_{\psi_M(M^2+M)}(1))))%
}

\NovelDisplay{M02772R01}{%
(0)(1,1,1)(2,1,1)(3,1)(2,1,1)(3,1)(2,1)(3,2,1)(4,2,1)(5,2)(4,2,1)(5,2)%
}
\NovelDisplay{M02772R02}{%
(3,2,1)(4,2,1)(5,2)(4,2,1)(5,2)%
}

\NovelDisplay{M02773R01}{%
= \psi(I^2 + I\times \psi_I(I^2+I\times \omega))%
}

\NovelDisplay{M02774R01}{%
= \psi(M^2+\psi_{\psi_M(M^2+M)}(M^2) + \psi_{\psi_M(M^2+M)}(M \times \psi_{\psi_M(M^2)}(M^2%
}
\NovelDisplay{M02774R02}{%
+\psi_{\psi_M(M^2+M)}(M^2)+\psi_{\psi_M(M^2+M)}(M \times \psi_{\psi_M(M^2)}(M^2%
}
\NovelDisplay{M02774R03}{%
+\psi_{\psi_M(M^2+M)}(M^2)+1)))))%
}

一旬里，山脚的灯常常留到很晚。

蔡须坤写M，卢本韦核I，霍雨浩在两本之间认参数，薛峰则能先找出重复输入与重见的尾部。四人做事的次序并不固定，谁先看清，便将那一处留给其余人检查。起初还要几遍才能接上的讨论，慢慢能够在同一页里说完。

有限倍数之后，背景开始进入更明显的幂次结构。

卢本韦将I平方乘二与后面出现的I立方分别圈出，没有将它们当成一句``继续增大''带过。I侧短，恰好能让她先看清当前运算的位置；读到M侧，却仍要由蔡须坤指出相应变化在多深的参数之中。

蔡须坤最吃力的那一行，左列已经跨过好几行，M侧也跟着分开写。她的手指从第一行挪到第二行时，下意识地想换一口气，脑中的结构便断了一下。

薛峰立即将另一张纸推到下方，接住前一行留下的位置。上面的BMS没有重新开始，右侧也没有因为排版结束一个函数。她先把尚未闭合的那一层说明，再往下念，终于没有靠猜测把后半段拼进去。

霍雨浩随后换了一种问法，让她从最深处读回外面。开始时她几乎句句停，卢本韦便在同一张I页上等着，不提前报出最后的短结果。等蔡须坤真正回到外层，她才指出两个表示在这一条完整对应上如何相接。

外面的风渐渐小了，石壁上的水仍没有干。唐三将绳子重新收短，诺当替几人把能站稳的地方擦过。四人轮流起身试阶，另外三人就将纸留在对方够得着的位置。

这时，领先已经不必由谁说出来。蔡须坤能独立找回M侧的深层位置，卢本韦能在I侧先认清整体，霍雨浩能核两边接续的前提，薛峰则能在答案没摆好以前，先将共同的输入守住。

后面的同伴仍各有卡住的地方。他们把多出来的这一点余裕，分给了那些具体的停顿。

\NovelDisplay{M02775R01}{%
(0)(1,1,1)(2,1,1)(3,1)(2,1,1)(3,1)(2,1)(3,2,1)(4,2,1)(5,2)(4,2,1)%
}
\NovelDisplay{M02775R02}{%
(5,2)(4)%
}

\NovelDisplay{M02776R01}{%
= \psi(I^2 \times 2)%
}

\NovelDisplay{M02777R01}{%
= \psi(M^2+\psi_{\psi_M(M^2+M)}(M^2) \times 2)%
}

\NovelDisplay{M02778R01}{%
(0)(1,1,1)(2,1,1)(3,1)(2,1,1)(3,1)(2,1)(3,2,1)(4,2,1)(5,2)(4,2,1)(5,2)%
}
\NovelDisplay{M02778R02}{%
(4,2)(5,3,1)(6,3,1)(7,3)(6,3,1)(7,3)(6)%
}

\NovelDisplay{M02779R01}{%
= \psi(I^2 \times 3)%
}

\NovelDisplay{M02780R01}{%
= \psi(M^2+\psi_{\psi_M(M^2+M)}(M^2) \times 3)%
}

\NovelDisplay{M02781R01}{%
(0)(1,1,1)(2,1,1)(3,1)(2,1,1)(3,1)(2,1,1)%
}

\NovelDisplay{M02782R01}{%
= \psi(I^2 \times \omega)%
}

\NovelDisplay{M02783R01}{%
= \psi(M^2+\psi_{\psi_M(M^2+M)}(M^2 + 1))%
}

\NovelDisplay{M02784R01}{%
(0)(1,1,1)(2,1,1)(3,1)(2,1,1)(3,1)(2,1,1)(3,1)%
}

\NovelDisplay{M02785R01}{%
= \psi(I^2 \times \Omega)%
}

\NovelDisplay{M02786R01}{%
= \psi(M^2+\psi_{\psi_M(M^2+M)}(M^2 + \psi_{\psi_M(M^2)}(1)))%
}

\NovelDisplay{M02787R01}{%
(0)(1,1,1)(2,1,1)(3,1)(2,1,1)(3,1)(2,1,1)(3,1)(1,1,1)(2,1,1)(3,1)(2)%
}

\NovelDisplay{M02788R01}{%
= \psi(I^2 \times \psi_I(I))%
}

\NovelDisplay{M02789R01}{%
= \psi(M^2+\psi_{\psi_M(M^2+M)}(M^2 + \psi_{\psi_M(M^2)}(M^2)))%
}

\NovelDisplay{M02790R01}{%
(0)(1,1,1)(2,1,1)(3,1)(2,1,1)(3,1)(2,1,1)(3,1)(1,1,1)(2,1,1)(3,1)(2,1,1)%
}

\NovelDisplay{M02791R01}{%
= \psi(I^2 \times \psi_I(I\times \omega))%
}

\NovelDisplay{M02792R01}{%
= \psi(M^2+\psi_{\psi_M(M^2+M)}(M^2 + \psi_{\psi_M(M^2)}(M^2 + \psi_{\psi_M(M^2+M)}(1))))%
}

\NovelDisplay{M02793R01}{%
(0)(1,1,1)(2,1,1)(3,1)(2,1,1)(3,1)(2,1,1)(3,1)(1,1,1)(2,1,1)(3,1)%
}
\NovelDisplay{M02793R02}{%
(2,1,1)(3,1)(2)%
}

\NovelDisplay{M02794R01}{%
= \psi(I^2 \times \psi_I(I^2))%
}

\NovelDisplay{M02795R01}{%
= \psi(M^2+\psi_{\psi_M(M^2+M)}(M^2 + \psi_{\psi_M(M^2)}(M^2 + \psi_{\psi_M(M^2+M)}(M^2))))%
}

\NovelDisplay{M02796R01}{%
(0)(1,1,1)(2,1,1)(3,1)(2,1,1)(3,1)(2,1,1)(3,1)(1,1,1)(2,1,1)(3,1)%
}
\NovelDisplay{M02796R02}{%
(2,1,1)(3,1)(2,1,1)%
}

\NovelDisplay{M02797R01}{%
= \psi(I^2 \times \psi_I(I^2 \times \omega))%
}

\NovelDisplay{M02798R01}{%
= \psi(M^2+\psi_{\psi_M(M^2+M)}(M^2 + \psi_{\psi_M(M^2)}(M^2 + \psi_{\psi_M(M^2+M)}(M^2 + 1))))%
}

\NovelDisplay{M02799R01}{%
(0)(1,1,1)(2,1,1)(3,1)(2,1,1)(3,1)(2,1,1)(3,1)(1,1,1)(2,1,1)(3,1)(2,1,1)%
}
\NovelDisplay{M02799R02}{%
(3,1)(2,1,1)(3,1)%
}

\NovelDisplay{M02800R01}{%
= \psi(I^2 \times \psi_I(I^2 \times \Omega))%
}

\NovelDisplay{M02801R01}{%
=\! \psi(M^2\!+\!\psi_{\psi_M(M^2+M)}(M^2 \!+\! \psi_{\psi_M(M^2)}(M^2 \!+\! \psi_{\psi_M(M^2+M)}(M^2 \!+\! \psi_{\psi_M(M^2)}(1)))))%
}

\NovelDisplay{M02802R01}{%
(0)(1,1,1)(2,1,1)(3,1)(2,1,1)(3,1)(2,1,1)(3,1)(1,1,1)(2,1,1)(3,1)(2,1,1)%
}
\NovelDisplay{M02802R02}{%
(3,1)(2,1,1)(3,1)(1,1,1)(2,1,1)(3,1)(2,1,1)(3,1)(2,1,1)(3,1)%
}

\NovelDisplay{M02803R01}{%
= \psi(I^2 \times \psi_I(I^2 \times \psi_I(I^2 \times \Omega)))%
}

\NovelDisplay{M02804R01}{%
= \psi(M^2+\psi_{\psi_M(M^2+M)}(M^2 + \psi_{\psi_M(M^2)}(M^2 + \psi_{\psi_M(M^2+M)}(M^2 + \psi_{\psi_M(M^2)}(M^2%
}
\NovelDisplay{M02804R02}{%
+ \psi_{\psi_M(M^2+M)}(M^2 + \psi_{\psi_M(M^2)}(1)))))))%
}

\NovelDisplay{M02805R01}{%
(0)(1,1,1)(2,1,1)(3,1)(2,1,1)(3,1)(2,1,1)(3,1)(2)%
}

\NovelDisplay{M02806R01}{%
= \psi(I^3)%
}

\NovelDisplay{M02807R01}{%
= \psi(M^2+\psi_{\psi_M(M^2+M)}(M^2 + \psi_M(M^2)))%
}

\NovelDisplay{M02808R01}{%
(0)(1,1,1)(2,1,1)(3,1)(2,1,1)(3,1)(2,1,1)(3,1)(2,1,1)%
}

\NovelDisplay{M02809R01}{%
= \psi(I^3\times \omega)%
}

\NovelDisplay{M02810R01}{%
= \psi(M^2+\psi_{\psi_M(M^2+M)}(M^2 + \psi_M(M^2)+1))%
}

\NovelDisplay{M02811R01}{%
(0)(1,1,1)(2,1,1)(3,1)(2,1,1)(3,1)(2,1,1)(3,1)(2,1,1)(3,1)(2)%
}

\NovelDisplay{M02812R01}{%
= \psi(I^4)%
}

\NovelDisplay{M02813R01}{%
= \psi(M^2+\psi_{\psi_M(M^2+M)}(M^2 + \psi_M(M^2)\times 2))%
}

I立方与I四次方写过以后，卢本韦没有急着去猜更高的终点。她把当前已经核清的几种运算分开，又回头问了一遍：最外层ψ是不是还在？

这一次，蔡须坤、霍雨浩与薛峰几乎同时答了出来。几人听见彼此的声音，都笑了一下，又各自把完整理由补上。

幂次终于也不再只取有限的数。

卢本韦能够读出I的ω次方，接到后面仍会被相似的写法绕住。指数里加一，指数自身再乘一段，与整座幂之后接一个乘数，纸上往往只差很短的位置，不能靠声音里那一个``乘''字分辨。

霍雨浩把她两次停错的地方留在纸边，没有给每一种都另起一个新字母。她们已经被长式困住够久，眼下更需要把原来的完整结构读顺，而非再添一串只有自己记得住的代称。

蔡须坤在M侧也遇到同样的困难。更深的参数里，又放入同一类函数值，外侧的括号仍须一路保留。她先看最外层，再从最内层读回，每回返回都说清现在是在指数的对应中，还是已经回到整个坍缩参数。

薛峰听到中途，忽然指出一处：``你在这里读回外面了，手却还指着内层。''

蔡须坤停下，重新沿着他的指尖找过去。错的正是她刚才最想跳过的地方。

她把那一处独立写清，让薛峰再核，然后主动将下一条也交给他。最早那个只在页角留白的人，已经能在她最费力的那一侧，找出明确的帮助位置。

到I的I次方附近时，卢本韦第一次比蔡须坤晚了半拍。她没有抢着把话补完，先让对方说清M侧，再回自己的I侧查指数的括号。换了记号以后，谁都会有看得较早的时候，也会有需要别人等的时候。

古月娜听完两人的解释，将同一条递给唐三。她们在前面已经找到一条可走的路，后面的人仍需亲自将这一段接完。

\NovelDisplay{M02814R01}{%
(0)(1,1,1)(2,1,1)(3,1)(3)%
}

\NovelDisplay{M02815R01}{%
= \psi(I^\omega)%
}

\NovelDisplay{M02816R01}{%
= \psi(M^2+\psi_{\psi_M(M^2+M)}(M^2 + \psi_{\psi_M(M^2+M)}(1)))%
}

\NovelDisplay{M02817R01}{%
(0)(1,1,1)(2,1,1)(3,1)(3)(2,1,1)(3,1)(2)%
}

\NovelDisplay{M02818R01}{%
= \psi(I^{\omega+1})%
}

\NovelDisplay{M02819R01}{%
= \psi(M^2+\psi_{\psi_M(M^2+M)}(M^2 + \psi_{\psi_M(M^2+M)}(1)+\psi_M(M^2)))%
}

\NovelDisplay{M02820R01}{%
(0)(1,1,1)(2,1,1)(3,1)(3)(2,1,1)(3,1)(3)%
}

\NovelDisplay{M02821R01}{%
= \psi(I^{\omega\times 2})%
}

\NovelDisplay{M02822R01}{%
= \psi(M^2+\psi_{\psi_M(M^2+M)}(M^2 + \psi_{\psi_M(M^2+M)}(1)\times 2))%
}

\NovelDisplay{M02823R01}{%
(0)(1,1,1)(2,1,1)(3,1)(3)(3)%
}

\NovelDisplay{M02824R01}{%
= \psi(I^{\omega^2})%
}

\NovelDisplay{M02825R01}{%
= \psi(M^2+\psi_{\psi_M(M^2+M)}(M^2 + \psi_{\psi_M(M^2+M)}(2)))%
}

\NovelDisplay{M02826R01}{%
(0)(1,1,1)(2,1,1)(3,1)(3,1)%
}

\NovelDisplay{M02827R01}{%
= \psi(I^{\Omega})%
}

\NovelDisplay{M02828R01}{%
= \psi(M^2+\psi_{\psi_M(M^2+M)}(M^2 + \psi_{\psi_M(M^2+M)}(M)))%
}

\NovelDisplay{M02829R01}{%
(0)(1,1,1)(2,1,1)(3,1)(3,1)(1,1,1)(2,1,1)(3,1)(2)%
}

\NovelDisplay{M02830R01}{%
= \psi(I^{\psi_I(I)})%
}

\NovelDisplay{M02831R01}{%
= \psi(M^2+\psi_{\psi_M(M^2+M)}(M^2 + \psi_{\psi_M(M^2+M)}(M \times \psi_{\psi_M(M)}(M^2))))%
}

\NovelDisplay{M02832R01}{%
(0)(1,1,1)(2,1,1)(3,1)(3,1)(1,1,1)(2,1,1)(3,1)(2,1,1)(3,1)(2)%
}

\NovelDisplay{M02833R01}{%
= \psi(I^{\psi_I(I^2)})%
}

\NovelDisplay{M02834R01}{%
= \psi(M^2+\psi_{\psi_M(M^2+M)}(M^2 + \psi_{\psi_M(M^2+M)}(M \times \psi_{\psi_M(M)}(M^2%
}
\NovelDisplay{M02834R02}{%
+ \psi_{\psi_M(M^2+M)}(M^2)))))%
}

\NovelDisplay{M02835R01}{%
(0)(1,1,1)(2,1,1)(3,1)(3,1)(1,1,1)(2,1,1)(3,1)(3)%
}

\NovelDisplay{M02836R01}{%
= \psi(I^{\psi_I(I^\omega)})%
}

\NovelDisplay{M02837R01}{%
= \psi(M^2+\psi_{\psi_M(M^2+M)}(M^2 + \psi_{\psi_M(M^2+M)}(M \times \psi_{\psi_M(M)}(M^2%
}
\NovelDisplay{M02837R02}{%
+ \psi_{\psi_M(M^2+M)}(M^2 + \psi_{\psi_M(M^2+M)}(1))))))%
}

\NovelDisplay{M02838R01}{%
(0)(1,1,1)(2,1,1)(3,1)(3,1)(1,1,1)(2,1,1)(3,1)(3,1)%
}

\NovelDisplay{M02839R01}{%
= \psi(I^{\psi_I(I^\Omega)})%
}

\NovelDisplay{M02840R01}{%
= \psi(M^2+\psi_{\psi_M(M^2+M)}(M^2 + \psi_{\psi_M(M^2+M)}(M \times \psi_{\psi_M(M)}(M^2%
}
\NovelDisplay{M02840R02}{%
+ \psi_{\psi_M(M^2+M)}(M^2 + \psi_{\psi_M(M^2+M)}(M))))))%
}

\NovelDisplay{M02841R01}{%
(0)(1,1,1)(2,1,1)(3,1)(3,1)(2)%
}

\NovelDisplay{M02842R01}{%
= \psi(I^I)%
}

\NovelDisplay{M02843R01}{%
= \psi(M^2+\psi_{\psi_M(M^2+M)}(M^2 + \psi_{\psi_M(M^2+M)}(M^2)))%
}

\NovelDisplay{M02844R01}{%
(0)(1,1,1)(2,1,1)(3,1)(3,1)(2,1,1)%
}

\NovelDisplay{M02845R01}{%
= \psi(I^I\times \omega)%
}

\NovelDisplay{M02846R01}{%
= \psi(M^2+\psi_{\psi_M(M^2+M)}(M^2 + \psi_{\psi_M(M^2+M)}(M^2)+1))%
}

\NovelDisplay{M02847R01}{%
(0)(1,1,1)(2,1,1)(3,1)(3,1)(2,1,1)(3,1)(2)%
}

\NovelDisplay{M02848R01}{%
= \psi(I^{I + 1})%
}

\NovelDisplay{M02849R01}{%
= \psi(M^2+\psi_{\psi_M(M^2+M)}(M^2 + \psi_{\psi_M(M^2+M)}(M^2)+\psi_M(M^2)))%
}

\NovelDisplay{M02850R01}{%
(0)(1,1,1)(2,1,1)(3,1)(3,1)(2,1,1)(3,1)(3,1)(2)%
}

\NovelDisplay{M02851R01}{%
= \psi(I^{I \times 2})%
}

\NovelDisplay{M02852R01}{%
= \psi(M^2+\psi_{\psi_M(M^2+M)}(M^2 + \psi_{\psi_M(M^2+M)}(M^2) \times 2))%
}

\NovelDisplay{M02853R01}{%
(0)(1,1,1)(2,1,1)(3,1)(3,1)(2,1,1)(3,1)(3,1)(2,1,1)(3,1)(3,1)(2)%
}

\NovelDisplay{M02854R01}{%
= \psi(I^{I \times 3})%
}

\NovelDisplay{M02855R01}{%
= \psi(M^2+\psi_{\psi_M(M^2+M)}(M^2 + \psi_{\psi_M(M^2+M)}(M^2) \times 3))%
}

这一组收尾时，蔡须坤将手从纸上拿开，发现指腹已经被炭条磨黑。

卢本韦递来湿布，自己却还在低声念指数中的那一段。蔡须坤听出她没有漏，便没打断，只替她将快被风翻过去的页角压住。

幂塔继续往里延伸，纸面上却不能随着塔高无限拉长。

霍雨浩将最长的一式分成几行，先让每行在一对完整的括号或一段明确的运算之后续接。能留在同一行的就留着，真正超出纸宽的才换行。他没有靠缩小字迹，把那些最容易读错的下标挤成一团。

卢本韦从I侧检查幂次的括号。塔里的幂按原式各自的指数范围来读，不能因为视觉上换到左边，就把后半段当成新的底数。她把停顿说给薛峰听，薛峰则核左列有没有在相似的几列之间漏掉一次重复。

蔡须坤将M侧读到最深的一层。函数头重复得太多，她在其中一处停了很久，额角又沁出汗。

这一次她没有独自硬撑，直接把纸转向霍雨浩：``这里我回不去。''

霍雨浩沿她已经说明的来路追过去，找到应当返回的位置，却没有替她接完。卢本韦将I侧同一个层次留在旁边，薛峰用指尖守着下一处还未合上的括号。蔡须坤自己将最后那一段读回来，终于能指出为什么外面还剩着这一层。

她靠着石壁歇了一会儿，才重新起身。果糕领域短暂浮起，又被她收拢到身周，没有试图伸到上方替谁承接整段考核。她知道，真正需要留住的是自己清楚听见的那一步，不能把领域的安稳当成所有括号都已核完。

走到最后几条时，四人开始能够在不同的纸上同时复算，再交换着解释。相同之处核得更快，真正不同的位置反而停得更仔细。霍雨浩听完薛峰的整段说明，终于将本子交还给他，没有再让他从头补一遍。

\NovelDisplay{M02856R01}{%
(0)(1,1,1)(2,1,1)(3,1)(3,1)(3)%
}

\NovelDisplay{M02857R01}{%
= \psi(I^{I \times \omega})%
}

\NovelDisplay{M02858R01}{%
= \psi(M^2+\psi_{\psi_M(M^2+M)}(M^2 + \psi_{\psi_M(M^2+M)}(M^2+1)))%
}

\NovelDisplay{M02859R01}{%
(0)(1,1,1)(2,1,1)(3,1)(3,1)(3,1)%
}

\NovelDisplay{M02860R01}{%
= \psi(I^{I \times \Omega})%
}

\NovelDisplay{M02861R01}{%
= \psi(M^2+\psi_{\psi_M(M^2+M)}(M^2 + \psi_{\psi_M(M^2+M)}(M^2+\psi_{\psi_M(M^2)}(1))))%
}

\NovelDisplay{M02862R01}{%
(0)(1,1,1)(2,1,1)(3,1)(3,1)(3,1)(2)%
}

\NovelDisplay{M02863R01}{%
= \psi(I^{I^2})%
}

\NovelDisplay{M02864R01}{%
= \psi(M^2+\psi_{\psi_M(M^2+M)}(M^2 + \psi_{\psi_M(M^2+M)}(M^2+\psi_M(M^2))))%
}

\NovelDisplay{M02865R01}{%
(0)(1,1,1)(2,1,1)(3,1)(3,1)(3,1)(2,1,1)%
}

\NovelDisplay{M02866R01}{%
= \psi(I^{I^2}\times \omega)%
}

\NovelDisplay{M02867R01}{%
= \psi(M^2+\psi_{\psi_M(M^2+M)}(M^2 + \psi_{\psi_M(M^2+M)}(M^2+\psi_M(M^2))+1))%
}

\NovelDisplay{M02868R01}{%
(0)(1,1,1)(2,1,1)(3,1)(3,1)(3,1)(3)%
}

\NovelDisplay{M02869R01}{%
= \psi(I^{I^2\times \omega})%
}

\NovelDisplay{M02870R01}{%
= \psi(M^2+\psi_{\psi_M(M^2+M)}(M^2 + \psi_{\psi_M(M^2+M)}(M^2+\psi_M(M^2)+1)))%
}

\NovelDisplay{M02871R01}{%
(0)(1,1,1)(2,1,1)(3,1)(3,1)(3,1)(3,1)(2)%
}

\NovelDisplay{M02872R01}{%
= \psi(I^{I^3})%
}

\NovelDisplay{M02873R01}{%
= \psi(M^2+\psi_{\psi_M(M^2+M)}(M^2 + \psi_{\psi_M(M^2+M)}(M^2+\psi_M(M^2)\times 2)))%
}

\NovelDisplay{M02874R01}{%
(0)(1,1,1)(2,1,1)(3,1)(3,1)(3,1)(3,1)(3,1)(2)%
}

\NovelDisplay{M02875R01}{%
= \psi(I^{I^4})%
}

\NovelDisplay{M02876R01}{%
= \psi(M^2+\psi_{\psi_M(M^2+M)}(M^2 + \psi_{\psi_M(M^2+M)}(M^2+\psi_M(M^2)\times 3)))%
}

\NovelDisplay{M02877R01}{%
(0)(1,1,1)(2,1,1)(3,1)(4)%
}

\NovelDisplay{M02878R01}{%
= \psi(I^{I^\omega})%
}

\NovelDisplay{M02879R01}{%
= \psi(M^2+\psi_{\psi_M(M^2+M)}(M^2 + \psi_{\psi_M(M^2+M)}(M^2+\psi_{\psi_M(M^2+M)}(1))))%
}

\NovelDisplay{M02880R01}{%
(0)(1,1,1)(2,1,1)(3,1)(4,1)%
}

\NovelDisplay{M02881R01}{%
= \psi(I^{I^\Omega})%
}

\NovelDisplay{M02882R01}{%
= \psi(M^2+\psi_{\psi_M(M^2+M)}(M^2 + \psi_{\psi_M(M^2+M)}(M^2+\psi_{\psi_M(M^2+M)}(M))))%
}

\NovelDisplay{M02883R01}{%
(0)(1,1,1)(2,1,1)(3,1)(4,1)(2)%
}

\NovelDisplay{M02884R01}{%
= \psi(I^{I^I})%
}

\NovelDisplay{M02885R01}{%
= \psi(M^2+\psi_{\psi_M(M^2+M)}(M^2 + \psi_{\psi_M(M^2+M)}(M^2+\psi_{\psi_M(M^2+M)}(M^2))))%
}

\NovelDisplay{M02886R01}{%
(0)(1,1,1)(2,1,1)(3,1)(4,1)(2,1,1)%
}

\NovelDisplay{M02887R01}{%
= \psi(I^{I^I}\times \omega)%
}

\NovelDisplay{M02888R01}{%
= \psi(M^2+\psi_{\psi_M(M^2+M)}(M^2 + \psi_{\psi_M(M^2+M)}(M^2+\psi_{\psi_M(M^2+M)}(M^2))+1))%
}

\NovelDisplay{M02889R01}{%
(0)(1,1,1)(2,1,1)(3,1)(4,1)(3)%
}

\NovelDisplay{M02890R01}{%
= \psi(I^{I^I\times \omega})%
}

\NovelDisplay{M02891R01}{%
= \psi(M^2+\psi_{\psi_M(M^2+M)}(M^2 + \psi_{\psi_M(M^2+M)}(M^2+\psi_{\psi_M(M^2+M)}(M^2)+1)))%
}

\NovelDisplay{M02892R01}{%
(0)(1,1,1)(2,1,1)(3,1)(4,1)(3,1)(2)%
}

\NovelDisplay{M02893R01}{%
= \psi(I^{I^{I+1}})%
}

\NovelDisplay{M02894R01}{%
=\! \psi(M^2\!+\!\psi_{\psi_M(M^2+M)}(M^2 \!+\! \psi_{\psi_M(M^2+M)}(M^2\!+\!\psi_{\psi_M(M^2+M)}(M^2)\!+\!\psi_M(M^2))))%
}

\NovelDisplay{M02895R01}{%
(0)(1,1,1)(2,1,1)(3,1)(4,1)(3,1)(4)%
}

\NovelDisplay{M02896R01}{%
= \psi(I^{I^{I+\omega}})%
}

\NovelDisplay{M02897R01}{%
= \psi(M^2+\psi_{\psi_M(M^2+M)}(M^2 + \psi_{\psi_M(M^2+M)}(M^2+\psi_{\psi_M(M^2+M)}(M^2)%
}
\NovelDisplay{M02897R02}{%
+\psi_{\psi_M(M^2+M)}(1))))%
}

\NovelDisplay{M02898R01}{%
(0)(1,1,1)(2,1,1)(3,1)(4,1)(3,1)(4,1)(2)%
}

\NovelDisplay{M02899R01}{%
= \psi(I^{I^{I\times 2}})%
}

\NovelDisplay{M02900R01}{%
= \psi(M^2+\psi_{\psi_M(M^2+M)}(M^2 + \psi_{\psi_M(M^2+M)}(M^2+\psi_{\psi_M(M^2+M)}(M^2)\times 2)))%
}

\NovelDisplay{M02901R01}{%
(0)(1,1,1)(2,1,1)(3,1)(4,1)(4)%
}

\NovelDisplay{M02902R01}{%
= \psi(I^{I^{I\times \omega}})%
}

\NovelDisplay{M02903R01}{%
= \psi(M^2+\psi_{\psi_M(M^2+M)}(M^2 + \psi_{\psi_M(M^2+M)}(M^2+\psi_{\psi_M(M^2+M)}(M^2+1))))%
}

\NovelDisplay{M02904R01}{%
(0)(1,1,1)(2,1,1)(3,1)(4,1)(4,1)(2)%
}

\NovelDisplay{M02905R01}{%
= \psi(I^{I^{I^2}})%
}

\NovelDisplay{M02906R01}{%
=\! \psi(M^2\!+\!\psi_{\psi_M(M^2+M)}(M^2 \!+\! \psi_{\psi_M(M^2+M)}(M^2\!+\!\psi_{\psi_M(M^2+M)}(M^2\!+\!\psi_M(M^2)))))%
}

\NovelDisplay{M02907R01}{%
(0)(1,1,1)(2,1,1)(3,1)(4,1)(5)%
}

\NovelDisplay{M02908R01}{%
= \psi(I^{I^{I^\omega}})%
}

\NovelDisplay{M02909R01}{%
= \psi(M^2+\psi_{\psi_M(M^2+M)}(M^2 + \psi_{\psi_M(M^2+M)}(M^2+\psi_{\psi_M(M^2+M)}(M^2%
}
\NovelDisplay{M02909R02}{%
+\psi_{\psi_M(M^2+M)}(1)))))%
}

\NovelDisplay{M02910R01}{%
(0)(1,1,1)(2,1,1)(3,1)(4,1)(5,1)(2)%
}

\NovelDisplay{M02911R01}{%
= \psi(I^{I^{I^I}})%
}

\NovelDisplay{M02912R01}{%
= \psi(M^2+\psi_{\psi_M(M^2+M)}(M^2 + \psi_{\psi_M(M^2+M)}(M^2+\psi_{\psi_M(M^2+M)}(M^2%
}
\NovelDisplay{M02912R02}{%
+\psi_{\psi_M(M^2+M)}(M^2)))))%
}

三百米的石桩已在上方，刻痕只露出一半。卢本韦抬头看了很久，又低头将最后一条幂塔的对应完整看完。

她没有先走。蔡须坤也将M册留在手里，等待那一处被讲义用红色标出的末行。前面所有细长的幂塔，都还不是她们可以替最后一行预写答案的理由。

蔡须坤看见M平方加M，第一反应竟是自己早在许多页以前就读过它。

她将旧册翻回去，果然找到了。可那一次，它出现在更深的函数头里，是某一个内层函数所用的参数。现在这一行，M平方加M直接成为最外层ψ所吃进的整个参数，位置已经不同。

``先前在里面见过，和现在整条式子走到这里，不是一回事。''她说。

卢本韦顺着I侧找回同一处，Ω的下标也必须完整保留。她没有因最后的M式比前面短，就把那一截I侧的下标略去。

霍雨浩请薛峰从共同的BMS再核一遍。薛峰认清最后一列，指出它与此前那几条幂塔对应的不同位置，才把左列交给蔡须坤复述。

这一行写完，蔡须坤终于将手慢慢松开。卢本韦把两份纸压齐，霍雨浩等她们各自说完，再请其余人说明自己能够独立认出的部分。讲义在这一处以红色收束EBO到JO的这一段路，他们也将它留在完整的对应里，没有只摘走一个更短的结果。

薛峰又把左列的前三个展开项单独核了一回。末列四、二的坏根在三、一，新增段仍依当前坏部生成，第二行保留原值。每一项都从同一原式出发，先前经过的幂塔对应可以用来对照，不能替代这一次具体的根与阶差。

\NovelDisplay{M02913R01}{%
\text{展开第零项：}\quad(0)(1,1,1)(2,1,1)(3,1)%
}

\NovelDisplay{M02914R01}{%
\text{展开第一项：}\quad(0)(1,1,1)(2,1,1)(3,1)(4,1)%
}

\NovelDisplay{M02915R01}{%
\text{展开第二项：}\quad(0)(1,1,1)(2,1,1)(3,1)(4,1)(5,1)%
}

第一处发生差别的列，已经让这些有限项处在原式之下，后面写得再长，也不会靠增加列数追成原式本身。蔡须坤、卢本韦与霍雨浩分别核完，才把这份左列检查与两套完整对应放在一起。

\NovelDisplay{M02916R01}{%
\textcolor{red}{(0)(1,1,1)(2,1,1)(3,1)(4,2)}%
}

\NovelDisplay{M02917R01}{%
\textcolor{red}{= \psi(\Omega_{I+1})}%
}

\NovelDisplay{M02918R01}{%
\textcolor{red}{= \psi(M^2+M)}%
}

最先走到石桩前的，是蔡须坤、卢本韦、霍雨浩与薛峰。

蔡须坤先落稳脚，卢本韦等自己把I侧完整接续，才将另一只脚跟上。霍雨浩回头核过刚才有人问住的一处，薛峰则能不翻答案，将自己的整段说明说完。

三百米的横台亮起来。四人没有把本子收走，分别站到后面同伴够得着的位置。唐三、小舞、诺当、奥斯卡、蔡兴坤、叶阳与古月娜随后陆续站稳，最后一人到了，绳子才真正松下来。

卢本韦坐下以后，很久才伸手去接水。

蔡须坤的M册摊在膝上，右边那一页已经被她翻得发软。她将最后一条看过，又回到前面那三处待核的异文，确认没有因为看见终点就把疑问擦掉。能完整说明的路已经留下，尚需证明的关系也仍有清楚的位置。

薛峰把四人的纸收在一起，随后才去接其余人的本子。收过来时，他发现大家已经习惯先问他哪一页能找到共同的左列。他应着，一份份归好，手上的动作终于不再迟疑。

夏夜星空看过他们站稳的地方，目光在前面四人之间停了一会儿，才请大家先歇。

从二百二十米到这里，蔡须坤一直主攻更繁复的M表示，卢本韦沿I表示认清整体，霍雨浩守住两侧解释所依赖的条件，薛峰则从核页、接缝，走到了能够独立复算和指出问题。此刻四人的分析与接续确实更早稳住，其余同伴也各自走到了横台，没有谁被一道快答留在下方。

古月娜将卢本韦身后的包放平，又把蔡须坤的水接满。她仍坐在几人之间，听见薛峰替诺当接回一处换行，便放心地把自己的那页重新展开。

蔡须坤仰头，五百米仍在更高的雾里。离这一项的目标还有二百米，离中分山顶还有二千二百米，两年的考核也没有因为队伍在山脚休养就停下来。

她没有急着给下一页起名字，只将手背贴了贴身后那点平静的温热，接过卢本韦递来的本子。

``三百米，先写上。''

霍雨浩将笔递过去。薛峰压住页角，等新刻度写稳，十一人才把今天这一段真正收好。

\chapter{第一百一十七章：寻找心中的bms矩阵}\label{ux7b2cux4e00ux767eux4e00ux5341ux4e03ux7ae0ux5bfbux627eux5fc3ux4e2dux7684bmsux77e9ux9635}

三百米的刻度写下以后，蔡须坤原以为，第二天醒来时，自己会轻松一些。

她醒得很早。窗纸还蒙着一层灰白，屋外有人走动，水倒进盆里，响了一会儿便停了。她披好外衣出去，看见唐三坐在门前，小舞站在他身后，两人中间摊着昨天的笔记。

那一页她认得。到三百米之前，大家已经一起读过。

唐三用指节压着前面的几列，试着从中段往后接。他说到一个下标，停下来，翻回前页。小舞等他找到，自己接了一句，也停在了那里。

蔡须坤走近时，小舞先抬起头。

``吵醒你了？''

``没有。''

她在旁边坐下，看了看两人的纸。唐三已经将同一处重新算了两遍，第一遍的痕迹很浅，第二遍压得重，纸背都能摸出印子。

``昨天能接上的。''他说。

``昨天我们都在，你也听过前面的解释。''蔡须坤指了指被他压住的那部分，``今天单独拿出来，前面少了一截来路，就容易卡。''

唐三没有辩解。他将手移开，从最初那列零重新读。

蔡须坤听到一半，便知道他还要在这里花很久。小舞的问题稍有不同：她记得这一段怎样展开，却总在换一种接法时沿用上一段的判断。若只问昨天最后的结果，两人都能答；若将末尾改动一点，他们便得重新摸索。

门里传来翻身的声音。奥斯卡也醒了，披着衣服出来，看见三人，便默默去烧水。

蔡须坤望了一眼山。

三百米那根石桩藏在薄雾里。昨天十一人都站上去了，这一点没有变，可每个人真正能独立走稳的路，仍有长短。

她原本要说``我陪你们重做一遍''，话到嘴边，忽然想起更早以前的一张纸。

那时他们还在七十米以下。

几个短得不能再短的列，困住了他们许多天。后来，他们终于看清提升怎样沿着祖先链发生，纸上的一行旧答案，也被整个改掉。

早饭时，蔡须坤将那张纸找了出来。

页角已经起毛，七十米旁边的墨迹被水晕过一次。霍雨浩看见她翻到这里，放下了手里的碗。

``你想重查那一段？''

``我想问另一件事。''

蔡须坤将最初的前缀抄在空纸上。

\NovelDisplay{M02919R01}{%
(0)(1,1,1)(2,1)=\psi(\Omega_\omega\times\Omega)%
}

``这一段，我们已经知道它对应什么。可后面再接一个一一一，当时我们最希望得到什么？''

卢本韦看了一眼，不用翻本子便答了出来。

``再加一份\allowbreak{}\NovelInline{I06034}{\Omega_\omega}\allowbreak{}。''

蔡须坤将他们曾经猜过的式子写在右边，随即用笔尖点住，没有落下等号。

\NovelDisplay{M02920R01}{%
\psi(\Omega_\omega\times\Omega+\Omega_\omega)%
}

``结果，原版BMS里的提升没有让它停在这里。''

\NovelDisplay{M02921R01}{%
(0)(1,1,1)(2,1)(1,1,1)=\psi(\Omega_\omega^2)%
}

屋里的人都认得这一行。

为了弄清这个平方从哪里来，他们追过父项，找过替身，将有限层的嵌套一层层摊开。七十米的石桩，也是在那段漫长的推演之后才真正踩稳。

蔡须坤将笔停在两条式子之间。

``后面多接一段，前面原以为已经固定的结构，也会被它牵动。越到后面，这件事越深。我们能算下去，却总得回头确认，自己之前认出的东西还在不在原来的位置。''

古月娜靠在窗边，目光落在那个平方上。

``你想让提升消失？''

``我想知道，能不能另外定义一种矩阵。''

蔡须坤说完，自己先安静了片刻。

这个念头很轻，真正说出口时，却让她有些心虚。那几个月的辛苦都还在本子里，她知道自己不能因为难，便将已经求出的答案涂掉。

霍雨浩将旧纸挪开了一点，给她面前留出空处。

``原来的留着。你把新的想法写在另一边。''

蔡须坤点头，在新纸顶端写下了几个字。

理想不提升BMS。

她写得很慢，写完以后，又在下面补了一句：另行设想，展开规则未定。

``原版照原版的规则走。''她说，``我想做的是另一套。至少在我们希望保留的地方，一段接上去，就能看清它加了什么，不要总将前面的含义一起改掉。''

卢本韦将凳子往她旁边挪了挪。

``就是你心里的那个？''

``嗯。''

蔡须坤在新纸下方写出第一条期望对应。她特意用了箭头，将它和旧纸上的原版等式分开。

\NovelDisplay{M02922R01}{%
\text{理想不提升BMS的期望对应：}%
}

\NovelDisplay{M02923R01}{%
(0)(1,1,1)(2,1)(1,1,1)\longmapsto\psi(\Omega_\omega\times\Omega+\Omega_\omega)%
}

``在这套新设想里，后接的一一一，就只在参数末尾加一份\allowbreak{}\NovelInline{I06035}{\Omega_\omega}\allowbreak{}。再接一段，也是再加一份。''

\NovelDisplay{M02924R01}{%
(0)(1,1,1)(2,1)(1,1,1)(1,1,1)\longmapsto\psi(\Omega_\omega\times\Omega+\Omega_\omega\times2)%
}

小舞不自觉地坐直了些。

她沿着左边数过两段相同的列，又看向右边的两份\allowbreak{}\NovelInline{I06036}{\Omega_\omega}\allowbreak{}。这一次，她不用先将许多页旧讲义叠在一起，便能说出蔡须坤想让它怎样增长。

``要是都能这样\ldots\ldots''她说了一半，轻轻吐出一口气。

``我们的考验会容易很多。''卢本韦接道。

唐三没有接话，目光却在新纸上停了很久。

蔡须坤看见了。

她并没有想好，山阶是否会接受一套由他们自己定义的体系。可若能把这种直观的接续真正建立起来，至少同伴不必每一次都从复杂的提升里重新辨认来路。他们需要的那座桥，也许能短一些。

她将新纸向桌中央推过去。

``雨浩，本韦，古月娜，还有薛峰。你们陪我试一试。''

薛峰刚将碗放回去，闻言抬起头。

``我也一起定义？''

``你最先看出过同一条矩阵被写成两份不同的答案。''蔡须坤说，``这次我怕的，就是我们写着写着，自己先把自己骗过去。''

薛峰看了一眼那张空白远多于字迹的纸，坐了下来。

唐三将自己的旧本子收拢。

``我们先把昨天那段补稳。有需要复算的，留给我。''

蔡须坤看着他，想说什么，唐三却先将另一张干净纸递了过来。

``你试。我这里也继续。''

上午，五人第一次真正围着那张新纸坐好。

霍雨浩先问：``最初的零，还保留吗？''

``保留。''蔡须坤说。

``末尾省略的零呢？''薛峰问。

``也保留原来的写法。我们先不要将外形全改掉。''

卢本韦用指尖点了点最后的一一一。

``那就从它开始。接上来，只加后面那一段。前面的不动。''

这句话说出来很顺。

等到霍雨浩将笔递给她，请她写出展开第一项时，卢本韦却没有立刻接。

她先看蔡须坤，随后又望向古月娜。

``还是先找根？''

古月娜没有回答。她将原版的祖先关系在旁边画出，从末列往前，一道道短线越过了卢本韦想保留的分界。

``按原版找，关系会穿过去。''她说。

``那就在这里停住。''

蔡须坤在新接的那段前面画了一道竖线。笔尖压得很稳，仿佛这样便能将前后截开。

霍雨浩看了一会儿。

``凭什么在这里停？''

``因为后面是新接的。''

``我们看着它写上去，才知道哪一段是新接的。换个人，只拿到最后这条完整矩阵，他从哪里读出这道线？''

蔡须坤的笔没有抬起来。

同一条矩阵可以分几次写下，也可以一口气写完。若规则还要问它昨晚写到哪里、今天从哪里添起，矩阵本身便不再足够。

她将竖线擦淡。

``可以另外标出边界。''

``可以。''霍雨浩说，``但那就是新添的结构。边界怎么放，能不能移动，两种边界会不会给同一条矩阵不同的展开，都得写清楚。''

卢本韦终于将笔接了过去。她试着在每一段前面加一个小记号，写了三条，很快发现前两条还能分清，第三条里的嵌套却需要在记号里面再标记号。

她将笔搁下，揉了揉眉心。

``这样下去，又要认哪一层包着哪一层。''

新纸上最初那条干净的加法，仍然没有变。它下面却已经挤满了细线。

午后，薛峰提议换一个方向。

``先别想着整个新体系。只拿这条，规定它想对应的序数，再找能往它逼近的那些项。''

蔡须坤立刻将纸翻过来。

她已经习惯更复杂的M表示。若只从序数一边看，许多外表纠缠的矩阵确实能被整理到明确的层次里。卢本韦也取出I册，将自己较熟的那一段摊开。

有那么一会儿，桌边的气氛又活了起来。

蔡须坤说最外层参数，卢本韦指出相邻的变化，霍雨浩将几种可能的逼近写在旁边。古月娜没有打断，薛峰则另拿一页，等着把每个有限项换回新矩阵。

他等了很久。

``第一项写哪条？''

蔡须坤顺着自己的推导往回找。她能说明自己希望接近哪个值，也能在已有记号中讨论它附近的结构；可要将这些东西交给``理想不提升BMS''，纸上便只剩下最初那几个没有定义完的列。

``可以先挑一条\ldots\ldots''卢本韦说。

``怎么挑？''

薛峰问得很轻，却让她后面的话停住了。

``这次我们在这里挑了一条，下次换一个末列，还能按同一条规则挑出来吗？还是每遇到一条，都要重新查I册和M册？''

蔡须坤低头看见自己摊开的厚本子。

若新矩阵每一次展开，都得先请他们将旧体系里最难的部分全部算清，再替它选一条路，那么唐三和小舞只是得到了一份更短的抄本。真正要跨过的那些困难，仍然压在抄本背面。

她将M册合上了。

桌边忽然少了许多翻页声。

``先记着这条路。''霍雨浩说，``它能帮助我们提出要求。但这会儿，还没有给出矩阵自己的展开办法。''

蔡须坤应了一声，在纸角写下``未解决''。

到了傍晚，五人仍没有写出展开第一项。

他们下意识试过原版的做法。父项一旦照旧追下去，蔡须坤想避开的关系便又出现在纸上；若提前截断，便要解释为什么偏偏截在这里。若将坐标的变化全部停住，新的项怎样接续、怎样靠近期望的值，又成了另一个空处。

古月娜将一条这样的尝试从头读到尾。

``少掉原版的提升步骤，后面的构造不会自动补齐。''她说。

蔡须坤点头，没有替那张纸辩解。

她原以为最难的是想清自己究竟要什么。如今想要的东西就写在面前，短短一行，人人都看得懂；越过这一行，便没有一条能直接走通的路。

夜里，他们回到石屋。

唐三和小舞仍在桌边。两人将上午卡住的那处各自复述完，又找出一条接法相近的旧题对照。小舞看见他们进门，先往旁边挪了挪，给五人空出位置。

``怎么样？''

卢本韦摇头。

小舞将本子合起一半，没再问。

蔡须坤在她身边坐下，想接着帮她看刚才那道题，眼睛落到纸上，却过了好一会儿才辨清最外面的括号。她只得将水碗端起来，借着喝水缓了一阵。

那天晚上，大家都睡得很晚。

第二天，五人去了三百米横台。

山上的风比屋里凉，纸张用小石子压着，稍一松手便向上掀。霍雨浩试着让众人沿同一条设想各写一次，不许看彼此的纸。他想先知道，大家口中的``不提升''，究竟是不是同一件事。

结果，五份纸放到一起，只有最初的期望值相同。

卢本韦将一段看成整体，蔡须坤想保留的是它在最外层参数中的贡献；古月娜沿着层次判断哪里应当继续，薛峰则担心相同的末尾接在不同前缀后，是否还应当这样解释。

霍雨浩自己也留了两处空白。

蔡须坤一张张看过去，脸上渐渐失了血色。

她没有想到，昨天说过那么多次的``不动''，连他们五人都还没有说成同一种意思。

卢本韦将自己的纸抽回来，划掉了一条线。

``那就先只要求外层参数不被改。''

古月娜指向里面：``这里变了，外面还可能保持原样吗？''

``我说的是已经固定的那一段。''

``哪一段算已经固定？''

卢本韦攥着炭条，没有马上回答。

这句话，他们昨天已经问过一次。

风将石子旁边的一角吹起来，啪地打回石面。薛峰伸手按住，发现自己掌心也是凉的。

又一天过去，三百米上方的石阶没有被他们向前推进。

唐三和小舞也来试了几次。旧路上的一些接续比前一天稳了，新的地方却仍要退回来。诺当本想问他们新矩阵有没有眉目，见五人都低着头，将话吞了回去。

奥斯卡带来的水，直到凉透才被喝完。

蔡兴坤抱着黑子鸡坐在较低处。鸡想往她们压纸的小石子旁边啄，被她抱了回去。平时这一下总能引出几句笑，今天只听见纸页翻动。

第三天，雨落下来了。

五人挤回屋里。窗边放着晾湿纸的细绳，桌上却已经没有足够的空处。

薛峰将尝试分成几叠，又重新合在一起。他原来总能替大家找到共同的左列，指出漏掉的括号，将断开的说明接起来。现在，同一行下面没有一个人写出能让他复算的规则，他连从哪一处开始查都不知道。

霍雨浩察觉蔡须坤很久没动，抬头看她。

她正盯着``展开第一项''后面的那片空白。

``再歇一会儿？''

``我不累。''

她说完，便知道这句话没有什么用。手腕的确没有酸，混沌之力也安安静静地收在体内，可只要想到下一次还得从头解释同一个关口，她就不想抬笔。

卢本韦将她手边的杯子接满。

``你昨天说过一遍了。''

``说过什么？''

``只要把它当成整体，就\ldots\ldots''

卢本韦没有再往下复述。她看见蔡须坤的神情，自己也低下了头。

雨声贴着屋檐不断往下落。

古月娜忽然将桌角一张干净纸抽到面前。

``可我觉得，有些地方不必绕这么久。''

她的声音不高，却让另外四人都看了过来。

``按我们想要的那种直观对应，即使还没有整套理想定义，我也知道这两条该写成什么。''

她先写下``心中BMS的期望对应''，随后落笔。

\NovelDisplay{M02925R01}{%
(0)(1,1,1)(2,1)(3,1)\longmapsto\psi(\Omega_\omega\times\Omega^\Omega)%
}

``前面保留\allowbreak{}\NovelInline{I06037}{\Omega_\omega}\allowbreak{}，后面接出我们想辨认的那段幂次。这样看，就很清楚。''

她又写出第二条。

\NovelDisplay{M02926R01}{%
(0)(1,1,1)(2,1)(3,2,1)(4,2)\longmapsto\psi(\Omega_\omega\times\Omega_2)%
}

``若要表示\allowbreak{}\NovelInline{I06038}{\psi(\Omega_\omega\times\Omega_2)}\allowbreak{}，我心里对应的就是这一条。它们是新体系想保留的对应，原版的答案另放一边。有些结果，我们现在就能选出来，还需要先把所有定义都等齐吗？''

卢本韦望着纸，眉头松开了一点。

这两行比桌上那些画满边界的草稿清楚得多。尤其第二条，长长一段矩阵终于落到一个能够一眼认清的乘法结构上，仿佛阴了几天的天色忽然开了一道缝。

蔡须坤却没有接笔。

她将第一行读完，又读了一遍第二行，手指停在古月娜刚写好的箭头旁边。

``展开给我看。''

古月娜抬起眼。

``哪一项？''

``第一项就行。''

屋里的翻页声停了。

古月娜重新看向左边的矩阵。她能说明自己为什么喜欢这样的对应，能沿着熟悉的序数结构描述右边应当怎样接续；可若只允许从这条新矩阵出发，她的笔尖便停在了末列旁。

蔡须坤等了一会儿，将那支笔轻轻推近她。

古月娜没有拿。

``你已经知道我想表达什么。''

``我知道。''蔡须坤说，``可唐三明天拿到它，小舞明天拿到它，他们要从哪里开始？''

``我们可以教。''

``教哪条规则？''

古月娜的目光沉下来。

``可以慢慢补。你写第一条设想的时候，不也没有规则？''

蔡须坤被这句话问住了。

她看向自己最早写的那张纸。理想不提升BMS，几个字写得比下面的式子还大。那天早晨，她将它推到桌中央时，唐三和小舞也都看了很久。

她曾经多么希望，那两道目光里能轻松一些。

``是。''她低声说，``我也没有。''

古月娜没有继续争，却也没将纸收走。

蔡须坤将两人的纸并在一起，指尖在最初的加法与新添的两条乘法之间慢慢划过。

``所以我的也没做好。谁写的都一样。我们可以说，希望它对应这里；可要是就这么说自己已经知道它怎么运转，我不敢让他们照着练。''

卢本韦低头望向自己的掌心。

山阶上接错一处时，魂力会怎样倒卷，她们都尝过。眼前的箭头很干净，纸上没有哪一处疼痛，便更容易让人觉得后面的空白不妨晚一点填。

霍雨浩将古月娜写的纸压平。

``原版那两条，有原版的展开和对应。你现在写的，是新体系期望赋予同样外形的另一种含义。我们可以保留这两个目标，但目标本身还不能告诉我们第一项。''

他没有从旧册里抄来任何一个展开结果。

将原版的项直接搬过来，只会将他们想改的关系也一并带回来；若搬来以后再挑几处删改，仍然得回答为什么这样改。

古月娜看了他一会儿，终于垂下眼。

``你觉得这是空中楼阁。''她对蔡须坤说。

蔡须坤的喉咙动了一下。

``现在这样，是。''

这句话落下去，桌边许久没有人出声。

蔡须坤没觉得轻松。古月娜是被她叫来帮忙的，从第一天到现在，除了她自己的本子，古月娜还替他们查过许多已经被划掉的尝试。如今最难听的话，却由她亲口说了出来。

她伸手，将自己那张纸往古月娜面前推了一点。

``我的这一张也在里面。''

古月娜没有看她，将笔拿了起来。

蔡须坤以为她要将那两行擦掉，刚想开口，便看见她只在下方添了六个小字。

期望，尚无展开。

薛峰将那张纸接过去，也在自己最早那次尝试下面写上了相同的话。

字写完以后，他没有像平常那样将纸排齐，只将两手搁在膝上，慢慢低下头。

``我们以前总能找到一个具体错处。''他说，``一个坐标，一条祖先链，或者少掉的那一层。现在连错在哪儿都说不清。''

霍雨浩没有立刻安慰他。

他自己的草稿，正压在桌面最下方。几次设想都能开头，往后推不了多久，便要增加一种没有说明过的判断。他能找出同伴的问题，却没能拿出自己的答案。

卢本韦将杯子往里推，碰响了另一只碗。她平时总会顺口说几句什么，这一次只是将碗扶稳。

屋门轻轻开了。

唐三抱着一捆细柴进来，小舞跟在后面，袖口湿了半截。两人看见桌边的神情，将东西放下，谁也没有先问算到了哪一步。

蔡须坤这才发现，天已经快黑了。

唐三去看炉火。小舞将晾纸绳往里面挪了一点，免得窗缝里的雨再飘上去。她拿起古月娜写的那一张，读完上面的几行，目光在最下面停住。

``是不是为了我们，才停在这儿？''

蔡须坤立即摇头，却又没有将最初的心思全藏起来。

``我早上看你们复算，就想，能不能有一种更容易走清楚的写法。后来越想越多\ldots\ldots''

小舞将纸放回原处。

``我也盼着。''她说。

这三个字让蔡须坤的眼眶忽然有些发热。

她原本还准备了一些解释，想告诉小舞新体系为什么不能只靠对应建立，想说自己再试几天也许会有办法。小舞却已经在她旁边坐下，将自己今天独立做过的那道旧题摊开。

那一页也有涂改。最外面的等号下方，她很认真地留出了检查的位置。

``这处今天接上了。''

蔡须坤看了很久，点点头。

唐三添好柴，走到桌边。

``往后我们还会慢。''他说，``你想试的新东西，也不必急着拿来让我们用。没想好的就留着。''

他没有说自己马上能赶上。那样的话，蔡须坤听了也不会安心。

卢本韦伸手将小舞的纸接过来，想帮她再核一次，目光却落到了自己今天仍未解决的那条式子上。她垂着眼，好半晌才将两张纸分开放好。

晚饭比平时安静。

夏夜星空回来时，问起白天的进度。霍雨浩把几天的尝试摊给她看，从最初想冻结前段，到后来借已有序数反找矩阵，一条条说完。

她听过以后，又看了看桌角那个三百米的刻度。

``山阶上呢？''

霍雨浩摇头。

夏夜星空将湿外衣挂好，没有替那张空白写下第一项。

``明天的安排，你们今晚商量。旧的考验还在。''

蔡须坤低低应了一声。

两年不会因为他们想另建一套记号便停下来。五百米也仍在那里。她望着晾在绳上的纸，几天前刚到三百米时的喜悦，已经被一种更深的疲惫压了下去。

夜深以后，其他人陆续睡了。

五人仍围着那张桌子。卢本韦披着外衣，眼睛盯着灯芯；薛峰将几张草稿收进一个空封套，写标题时停了很久，最后只写了``心中BMS''。

霍雨浩将原版的笔记放在另一边。那些他们确实能够独立解释的页，没有被新设想覆盖。他伸手想将蔡须坤面前的纸也收起，她却轻轻按住了。

``再留一会儿。''

古月娜坐在她对面，许久才开口。

``我没觉得那两个目标不好。''

``我也没有。''

``但我确实写不出第一项。''

蔡须坤将目光移到她脸上。古月娜没有避开，眉眼间还带着白日争执之后未散尽的倦意。

蔡须坤点了一下头。

``我也是。''

古月娜伸手，把快要碰到灯座的纸拉远了些。

桌面上，一边是他们已经走过的原版BMS，一边是几条清楚、简洁、尚无展开规则的期望对应。雨已经小了，屋檐还在滴水，五个人谁也没有再提一个新的办法。

薛峰将封套放到蔡须坤手边。

她终于把那张纸折起，放了进去。折痕经过最下面的一行字时，她的手指停了一瞬，又慢慢压平。

那一行写着：

展开第一项：

后面还是空的。

\chapter{第一百一十八章：从能够展开的地方}\label{ux7b2cux4e00ux767eux4e00ux5341ux516bux7ae0ux4eceux80fdux591fux5c55ux5f00ux7684ux5730ux65b9}

雨停后的第二天，唐三将一页写满的纸放在蔡须坤床边，压纸的是一颗从鞋底剔下来的小石子。

上面没有新定义。他只将前一天独立做过的那一段重新算了一遍，又添了两个自己还不确定的问题。蔡须坤坐起来，看完第一页，伸手去翻后面，才发现只有这一张。

``这次没看答案。''唐三站在门口说。

她将纸捧近一些。末尾有一处涂改，原先接错的列被他轻轻划去，下面补了来路。那道她昨晚已经不愿再看的旧题，被一个比她走得慢的人，一步一步接完了。

蔡须坤找出自己的本子，回答他留下的第一个问题。答到一半，她的笔停了一下，又沿原式回看，最后将第二个问题圈了出来。

``这个要重新追。早饭以后一起。''

说完她才发现，自己已经将脚伸进了鞋里。

那只装着``心中BMS''的封套仍在枕边。她摸过封口，没有将里面的草稿抽出来，而是把它收进了行李最内层。她还愿意想那套理想不提升BMS，只是今天上山，要带能够实际展开的那一本。

卢本韦在门外等她，手里抱着I册。古月娜将干衣服搭在包上，替她拢好漏出一角的纸；两人望见彼此，都没有再提雨夜里的争执。

回到三百米横台时，十一人的呼吸已经重了。夏夜星空将新讲义放在石桩旁，第一页，恰是他们此前走完的那个JO节点。

霍雨浩没有让大家从下一页抢着开始。他请每个人先将这条旧路说一遍，自己站在最外侧听；薛峰把本子摊开，第一次没有预先留出一片等别人填满的大空白。

蔡须坤低下头，重新读起M侧。那些函数头仍然繁复，她却记得自己如何将最后一个括号接回最外层。

``从这里。''她说。

她将笔放到第一列旁边，没有越过那列零。下面每条完整对应，仍依次保留BMS、I表示和M表示；矩阵只省略末尾补齐的零，内部的零照原位置读回去。I与M沿用两册讲义的记法，换行也不另起一个矩阵。

\NovelDisplay{M02927R01}{%
(0)(1,1,1)(2,1,1)(3,1)(4,2)%
}

\NovelDisplay{M02928R01}{%
=\psi(\Omega_{I+1})%
}

\NovelDisplay{M02929R01}{%
= \psi(M^2+M)%
}

\NovelDisplay{M02930R01}{%
(0)(1,1,1)(2,1,1)(3,1)(4,2)(1,1,1)(2,1,1)(3,1)(2)%
}

\NovelDisplay{M02931R01}{%
=\psi(\Omega_{I+1}+\psi_{I}(I))%
}

\NovelDisplay{M02932R01}{%
= \psi(M^2+M+\psi_{\psi_{M}(M^2)}(M^2))%
}

\NovelDisplay{M02933R01}{%
(0)(1,1,1)(2,1,1)(3,1)(4,2)(1,1,1)(2,1,1)(3,1)(2,1,1)(3,1)(2)%
}

\NovelDisplay{M02934R01}{%
=\psi(\Omega_{I+1}+\psi_{I}(I^2))%
}

\NovelDisplay{M02935R01}{%
= \psi(M^2+M+\psi_{\psi_{M}(M^2)}(M^2+\psi_{\psi_{M}(M^2+M)}(M^2)))%
}

\NovelDisplay{M02936R01}{%
(0)(1,1,1)(2,1,1)(3,1)(4,2)(1,1,1)(2,1,1)(3,1)(2,1,1)(3,1)(2,1,1)%
}
\NovelDisplay{M02936R02}{%
(3,1)(2)%
}

\NovelDisplay{M02937R01}{%
=\psi(\Omega_{I+1}+\psi_{I}(I^3))%
}

\NovelDisplay{M02938R01}{%
= \psi(M^2+M+\psi_{\psi_{M}(M^2)}(M^2+\psi_{\psi_{M}(M^2+M)}(M^2+\psi_{M}(M^2))))%
}

卢本韦归纳道：``这里先固定JO背景，后面的I幂进入内函数参数，不能把它们读成整个外层式子的幂。''

\NovelDisplay{M02939R01}{%
(0)(1,1,1)(2,1,1)(3,1)(4,2)(1,1,1)(2,1,1)(3,1)(3)%
}

\NovelDisplay{M02940R01}{%
=\psi(\Omega_{I+1}+\psi_{I}(I^\omega))%
}

\NovelDisplay{M02941R01}{%
= \psi(M^2+M+\psi_{\psi_{M}(M^2)}(M^2+\psi_{\psi_{M}(M^2+M)}(M^2+\psi_{\psi_{M}(M^2+M)}(1))))%
}

\NovelDisplay{M02942R01}{%
(0)(1,1,1)(2,1,1)(3,1)(4,2)(1,1,1)(2,1,1)(3,1)(3,1)(2)%
}

\NovelDisplay{M02943R01}{%
=\psi(\Omega_{I+1}+\psi_{I}(I^I))%
}

\NovelDisplay{M02944R01}{%
= \psi(M^2+M+\psi_{\psi_{M}(M^2)}(M^2+\psi_{\psi_{M}(M^2+M)}(M^2+\psi_{\psi_{M}(M^2+M)}(M^2))))%
}

\NovelDisplay{M02945R01}{%
(0)(1,1,1)(2,1,1)(3,1)(4,2)(1,1,1)(2,1,1)(3,1)(4,2)%
}

\NovelDisplay{M02946R01}{%
=\psi(\Omega_{I+1}+\psi_{I}(\Omega_{I+1}))%
}

\NovelDisplay{M02947R01}{%
= \psi(M^2+M+\psi_{\psi_{M}(M^2)}(M^2+M))%
}

\NovelDisplay{M02948R01}{%
(0)(1,1,1)(2,1,1)(3,1)(4,2)(2)%
}

\NovelDisplay{M02949R01}{%
=\psi(\Omega_{I+1}+\psi_{\Omega_{\psi_{I}(\Omega_{I+1})+1}}(1))%
}

\NovelDisplay{M02950R01}{%
= \psi(M^2+M+\psi_{\psi_{\psi_{M}(M^2)}(M^2+M+1)}(1))%
}

这几条先让内参数越过有限幂与幂塔，再更换函数下标，M侧必须同时保留外面的\allowbreak{}\NovelInline{I06039}{M^2+M}\allowbreak{}背景。

霍雨浩先将一条已经读过的新式从答案里抽出，只留下共同的左列。唐三伸手想翻I册，被小舞按住了页角。

``先写矩阵。''她说。

薛峰将同一原式的第零项、第一项、第二项分别求出，唐三核保留段，小舞核每一段新接的列；三个取项都从原式出发，第二项没有拿第一项再展开来代替。

\NovelDisplay{M02951R01}{%
\text{展开第零项：}\quad(0)(1,1,1)(2,1,1)(3,1)(4,2)(1,1,1)(2,1,1)(3,1)%
}

\NovelDisplay{M02952R01}{%
\text{展开第一项：}\quad(0)(1,1,1)(2,1,1)(3,1)(4,2)(1,1,1)(2,1,1)(3,1)%
}
\NovelDisplay{M02952R02}{%
(1,1,1)(2,1,1)(3,1)%
}

\NovelDisplay{M02953R01}{%
\text{展开第二项：}\quad(0)(1,1,1)(2,1,1)(3,1)(4,2)(1,1,1)(2,1,1)(3,1)(1,1,1)%
}
\NovelDisplay{M02953R02}{%
(2,1,1)(3,1)(1,1,1)(2,1,1)(3,1)%
}

这三个有限项保留相同前段并增加明确的重复段，各自低于原式，列数增加不能代替原版BMS的比较。

小舞把第三张纸递过去时，蔡须坤才发觉，自己已经听着她们的复算往前挪了半步。她还没弄清下一页，可眼前这一步，确实有来路。

她将两只脚在石面上站稳，伸手把唐三的本子接了过来。

第三天，他们第一次在三百米以上留住了一处完整接续。

石阶没有因此变轻。蔡须坤落脚时仍要将上身略向前压，脚踝一点点找稳，才敢抬另一只脚。她说到内函数的下标时停了一次，卢本韦没有替她念完，只将I册翻到相同的位置。

``先告诉我，你现在在哪个括号里。''

蔡须坤闭了闭眼，将前后两层分别说出。霍雨浩等她说稳，才请小舞复述这一处的区别。小舞答得慢，站在后面的唐三却一直没有催。

风把页角吹到几人手背上，薛峰用掌心压住，直到所有人的声音都接上。

\NovelDisplay{M02954R01}{%
(0)(1,1,1)(2,1,1)(3,1)(4,2)(2,1)(1,1,1)(2,1,1)(3,1)(4,2)%
}

\NovelDisplay{M02955R01}{%
=\psi(\Omega_{I+1}+\psi_{\Omega_{\psi_{I}(\Omega_{I+1})+1}}(\Omega_{I+1}))%
}

\NovelDisplay{M02956R01}{%
= \psi(M^2+M+\psi_{\psi_{\psi_{M}(M^2)}(M^2+M+1)}(M^2+M))%
}

\NovelDisplay{M02957R01}{%
(0)(1,1,1)(2,1,1)(3,1)(4,2)(2,1)(3,2)%
}

\NovelDisplay{M02958R01}{%
=\psi(\Omega_{I+1}+\Omega_{\psi_{I}(\Omega_{I+1})+1})%
}

\NovelDisplay{M02959R01}{%
= \psi(M^2+M+\psi_{\psi_{M}(M^2)}(M^2+M+1))%
}

\NovelDisplay{M02960R01}{%
(0)(1,1,1)(2,1,1)(3,1)(4,2)(2,1)(3,2)(4,3)%
}

\NovelDisplay{M02961R01}{%
=\psi(\Omega_{I+1}+\Omega_{\psi_{I}(\Omega_{I+1})+2})%
}

\NovelDisplay{M02962R01}{%
= \psi(M^2+M+\psi_{\psi_{M}(M^2)}(M^2+M+2))%
}

\NovelDisplay{M02963R01}{%
(0)(1,1,1)(2,1,1)(3,1)(4,2)(2,1)(3,2,1)%
}

\NovelDisplay{M02964R01}{%
=\psi(\Omega_{I+1}+\Omega_{\psi_{I}(\Omega_{I+1})+\omega})%
}

\NovelDisplay{M02965R01}{%
= \psi(M^2+M+\psi_{\psi_{M}(M^2)}(M^2+M+\omega))%
}

霍雨浩指出，同样出现一个一或一个ω时，应先看它落在函数参数还是Ω的下标里，再决定怎样接续。

\NovelDisplay{M02966R01}{%
(0)(1,1,1)(2,1,1)(3,1)(4,2)(2,1)(3,2,1)(4,2,1)(5,2)(4)%
}

\NovelDisplay{M02967R01}{%
=\psi(\Omega_{I+1}+I)%
}

\NovelDisplay{M02968R01}{%
= \psi(M^2+M+\psi_{M}(M^2))%
}

\NovelDisplay{M02969R01}{%
(0)(1,1,1)(2,1,1)(3,1)(4,2)(2,1)(3,2,1)(4,2,1)(5,2)(4,2)(5,3,1)%
}
\NovelDisplay{M02969R02}{%
(6,3,1)(7,3)(6)%
}

\NovelDisplay{M02970R01}{%
=\psi(\Omega_{I+1}+I\times 2)%
}

\NovelDisplay{M02971R01}{%
= \psi(M^2+M+\psi_{M}(M^2)\times 2)%
}

\NovelDisplay{M02972R01}{%
(0)(1,1,1)(2,1,1)(3,1)(4,2)(2,1)(3,2,1)(4,2,1)(5,2)(4,2,1)%
}

\NovelDisplay{M02973R01}{%
=\psi(\Omega_{I+1}+I\times \omega)%
}

\NovelDisplay{M02974R01}{%
= \psi(M^2+M+\psi_{\psi_{M}(M^2+M)}(1))%
}

\NovelDisplay{M02975R01}{%
(0)(1,1,1)(2,1,1)(3,1)(4,2)(2,1)(3,2,1)(4,2,1)(5,2)(4,2,1)(5,2)(4)%
}

\NovelDisplay{M02976R01}{%
=\psi(\Omega_{I+1}+I^2)%
}

\NovelDisplay{M02977R01}{%
= \psi(M^2+M+\psi_{\psi_{M}(M^2+M)}(M^2))%
}

新增的尾部从两份I走向I的乘法与平方，前缀的JO值仍完整留在外层参数中。

\NovelDisplay{M02978R01}{%
(0)(1,1,1)(2,1,1)(3,1)(4,2)(2,1)(3,2,1)(4,2,1)(5,2)(4,2,1)(5,2)%
}
\NovelDisplay{M02978R02}{%
(4,2,1)(5,2)(4)%
}

\NovelDisplay{M02979R01}{%
=\psi(\Omega_{I+1}+I^3)%
}

\NovelDisplay{M02980R01}{%
= \psi(M^2+M+\psi_{\psi_{M}(M^2+M)}(M^2+\psi_{M}(M^2)))%
}

\NovelDisplay{M02981R01}{%
(0)(1,1,1)(2,1,1)(3,1)(4,2)(2,1)(3,2,1)(4,2,1)(5,2)(5)%
}

\NovelDisplay{M02982R01}{%
=\psi(\Omega_{I+1}+I^\omega)%
}

\NovelDisplay{M02983R01}{%
= \psi(M^2+M+\psi_{\psi_{M}(M^2+M)}(M^2+\psi_{\psi_{M}(M^2+M)}(1)))%
}

\NovelDisplay{M02984R01}{%
(0)(1,1,1)(2,1,1)(3,1)(4,2)(2,1)(3,2,1)(4,2,1)(5,2)(5,2)(4)%
}

\NovelDisplay{M02985R01}{%
=\psi(\Omega_{I+1}+I^I)%
}

\NovelDisplay{M02986R01}{%
= \psi(M^2+M+\psi_{\psi_{M}(M^2+M)}(M^2+\psi_{\psi_{M}(M^2+M)}(M^2)))%
}

\NovelDisplay{M02987R01}{%
(0)(1,1,1)(2,1,1)(3,1)(4,2)(2,1)(3,2,1)(4,2,1)(5,2)(6,3)%
}

\NovelDisplay{M02988R01}{%
=\psi(\Omega_{I+1}+\psi_{\Omega_{I+1}}(\Omega_{I+1}))%
}

\NovelDisplay{M02989R01}{%
= \psi(M^2+M+\psi_{\psi_{M}(M^2+M)}(M^2+M))%
}

蔡须坤将这一组的变化归在同一处：I幂的增长逐步进入更高内函数，而外层背景不因内层变长便消失。

五米的刻度并不起眼，藏在一处被磨得发亮的石棱边。卢本韦看见它时，先以为是旧划痕，伏低身才认出那个数。

``三百零五。''

蔡须坤没有马上笑。她回头数过身后的人，等小舞也在相邻的平处站稳，才将一直握紧的笔换到另一只手里。

奥斯卡把水递上来，诺当替他扶住背包。今天还没有人说这条路已经容易，他们只是终于能指出，刚才哪一处比昨天多懂了一点。

再往上，讲义将同一个内函数放回不同的位置，短短的加号开始让人反复受阻。

\NovelDisplay{M02990R01}{%
(0)(1,1,1)(2,1,1)(3,1)(4,2)(2,1)(3,2,1)(4,2,1)(5,2)(6,3)(4,2)(5,3,1)%
}
\NovelDisplay{M02990R02}{%
(6,3,1)(7,3)(6)%
}

\NovelDisplay{M02991R01}{%
=\psi(\Omega_{I+1}+\psi_{\Omega_{I+1}}(\Omega_{I+1})+I)%
}

\NovelDisplay{M02992R01}{%
= \psi(M^2+M+\psi_{\psi_{M}(M^2+M)}(M^2+M)+\psi_{M}(M^2))%
}

\NovelDisplay{M02993R01}{%
(0)(1,1,1)(2,1,1)(3,1)(4,2)(2,1)(3,2,1)(4,2,1)(5,2)(6,3)(4,2)(5,3,1)%
}
\NovelDisplay{M02993R02}{%
(6,3,1)(7,3)(8,4)%
}

\NovelDisplay{M02994R01}{%
=\psi(\Omega_{I+1}+\psi_{\Omega_{I+1}}(\Omega_{I+1})\times 2)%
}

\NovelDisplay{M02995R01}{%
= \psi(M^2+M+\psi_{\psi_{M}(M^2+M)}(M^2+M)\times 2)%
}

\NovelDisplay{M02996R01}{%
(0)(1,1,1)(2,1,1)(3,1)(4,2)(2,1,1)%
}

\NovelDisplay{M02997R01}{%
=\psi(\Omega_{I+1}+\psi_{\Omega_{I+1}}(\Omega_{I+1}+1))%
}

\NovelDisplay{M02998R01}{%
= \psi(M^2+M+\psi_{\psi_{M}(M^2+M)}(M^2+M+1))%
}

\NovelDisplay{M02999R01}{%
(0)(1,1,1)(2,1,1)(3,1)(4,2)(2,1,1)(2,1)(3,2,1)(4,2,1)(5,2)(6,3)%
}

\NovelDisplay{M03000R01}{%
=\psi(\Omega_{I+1}+\psi_{\Omega_{I+1}}(\Omega_{I+1}+1)+\psi_{\Omega_{I+1}}(\Omega_{I+1}))%
}

\NovelDisplay{M03001R01}{%
= \psi(M^2+M+\psi_{\psi_{M}(M^2+M)}(M^2+M+1)+\psi_{\psi_{M}(M^2+M)}(M^2+M))%
}

古月娜提醒他们，尾部另加I、重复整个内函数值，以及把内参数加一，是三种不同位置的改变。

小舞拿着本子过来时，蔡须坤刚将笔收起。

``还问吗？''她说。

小舞看了看她磨黑的手指，又把本子合了一点：``明天也行。''

蔡须坤本想点头，却看见夹在里面的那张纸已经添了许多小记号。最早的一个，被小舞在旁边画了三道短线，后面仍留着空处。

``这个问过几次了？''

``你讲过。我后来自己做，又不确定。''

蔡须坤将本子接过，记起自己上次说到这里，曾随口说这一处很简单。那时她惦记着后面的难题，没有看小舞究竟从哪里开始跟不上。

她把包重新放下，在石上坐好。

``先说你现在怎么想。''

小舞开头说得很轻，见她一直听，后面便快起来。唐三从旁边经过，听到熟悉的错处，也停下了脚。三人没有另取新题，只将这个问过多次的问题从头拆开。

等小舞自己找到先前误接的位置，天已经暗了。蔡须坤把那三道短线擦淡，留下她刚补的说明，才将本子还回去。

``以后卡住就问，别等我看你脸色才猜。''

小舞接过，嘴角动了一下：``那你也得说累。''

蔡须坤愣了愣，终于点头。

\NovelDisplay{M03002R01}{%
(0)(1,1,1)(2,1,1)(3,1)(4,2)(2,1,1)(2,1)(3,2,1)(4,2,1)(5,2)(6,3)(4,2,1)%
}

\NovelDisplay{M03003R01}{%
=\psi(\Omega_{I+1}+\psi_{\Omega_{I+1}}(\Omega_{I+1}+1)\times 2)%
}

\NovelDisplay{M03004R01}{%
= \psi(M^2+M+\psi_{\psi_{M}(M^2+M)}(M^2+M+1)\times 2)%
}

\NovelDisplay{M03005R01}{%
(0)(1,1,1)(2,1,1)(3,1)(4,2)(2,1,1)(2,1,1)%
}

\NovelDisplay{M03006R01}{%
=\psi(\Omega_{I+1}+\psi_{\Omega_{I+1}}(\Omega_{I+1}+2))%
}

\NovelDisplay{M03007R01}{%
= \psi(M^2+M+\psi_{\psi_{M}(M^2+M)}(M^2+M+2))%
}

\NovelDisplay{M03008R01}{%
(0)(1,1,1)(2,1,1)(3,1)(4,2)(2,1,1)(2,1,1)(2,1,1)%
}

\NovelDisplay{M03009R01}{%
=\psi(\Omega_{I+1}+\psi_{\Omega_{I+1}}(\Omega_{I+1}+3))%
}

\NovelDisplay{M03010R01}{%
= \psi(M^2+M+\psi_{\psi_{M}(M^2+M)}(M^2+M+3))%
}

\NovelDisplay{M03011R01}{%
(0)(1,1,1)(2,1,1)(3,1)(4,2)(2,1,1)(3)%
}

\NovelDisplay{M03012R01}{%
=\psi(\Omega_{I+1}+\psi_{\Omega_{I+1}}(\Omega_{I+1}+\omega))%
}

\NovelDisplay{M03013R01}{%
= \psi(M^2+M+\psi_{\psi_{M}(M^2+M)}(M^2+M+\omega))%
}

这四条在同一个内函数里推进有限后继与ω，读M式时要把这些增量留在同一对参数括号内。

\NovelDisplay{M03014R01}{%
(0)(1,1,1)(2,1,1)(3,1)(4,2)(2,1,1)(3,1)%
}

\NovelDisplay{M03015R01}{%
=\psi(\Omega_{I+1}+\psi_{\Omega_{I+1}}(\Omega_{I+1}+\Omega))%
}

\NovelDisplay{M03016R01}{%
= \psi(M^2+M+\psi_{\psi_{M}(M^2+M)}(M^2+M+\psi_{\psi_M(M^2)}(1)))%
}

\NovelDisplay{M03017R01}{%
(0)(1,1,1)(2,1,1)(3,1)(4,2)(2,1,1)(3,1)(2)%
}

\NovelDisplay{M03018R01}{%
=\psi(\Omega_{I+1}+\psi_{\Omega_{I+1}}(\Omega_{I+1}+I))%
}

\NovelDisplay{M03019R01}{%
= \psi(M^2+M+\psi_{\psi_{M}(M^2+M)}(M^2+M+\psi_M(M^2)))%
}

\NovelDisplay{M03020R01}{%
(0)(1,1,1)(2,1,1)(3,1)(4,2)(2,1,1)(3,1)(2,1,1)%
}

\NovelDisplay{M03021R01}{%
=\psi(\Omega_{I+1}+\psi_{\Omega_{I+1}}(\Omega_{I+1}+I+1))%
}

\NovelDisplay{M03022R01}{%
= \psi(M^2+M+\psi_{\psi_{M}(M^2+M)}(M^2+M+\psi_M(M^2)+1))%
}

\NovelDisplay{M03023R01}{%
(0)(1,1,1)(2,1,1)(3,1)(4,2)(2,1,1)(3,1)(2,1,1)%
}

\NovelDisplay{M03024R01}{%
=\psi(\Omega_{I+1}+\psi_{\Omega_{I+1}}(\Omega_{I+1}+I+1))%
}

\NovelDisplay{M03025R01}{%
= \psi(M^2+M+\psi_{\psi_{M}(M^2+M)}(M^2+M+\psi_M(M^2)+1))%
}

薛峰归纳道：``这一组把Ω与I放回具体参数，末尾两条是完全相同的复算，不能算作两个新节点。''

卢本韦原想把相同的一行跳过去，薛峰却将两页并在一起，逐列核到最后，确认那是原样重复。

``留作自己算的第二遍。''他说，``别算作又过了一阶。''

小舞便将前一次的右端遮住，重新读同一条。她少念了一处一，卢本韦刚要提醒，唐三已经在旁边轻轻敲了敲纸，她自己停下来，回去找到了。

蔡须坤望着这一幕，把原先准备圈作进度的记号擦去。重复的一页仍有用处，只是它不会替任何人增加高度。

\NovelDisplay{M03026R01}{%
(0)(1,1,1)(2,1,1)(3,1)(4,2)(2,1,1)(3,1)(2,1,1)(3,1)(2)%
}

\NovelDisplay{M03027R01}{%
=\psi(\Omega_{I+1}+\psi_{\Omega_{I+1}}(\Omega_{I+1}+I\times 2))%
}

\NovelDisplay{M03028R01}{%
= \psi(M^2+M+\psi_{\psi_{M}(M^2+M)}(M^2+M+\psi_M(M^2)\times 2))%
}

\NovelDisplay{M03029R01}{%
(0)(1,1,1)(2,1,1)(3,1)(4,2)(2,1,1)(3,1)(3)%
}

\NovelDisplay{M03030R01}{%
=\psi(\Omega_{I+1}+\psi_{\Omega_{I+1}}(\Omega_{I+1}+\psi_{\Omega_{I+1}}(1)))%
}

\NovelDisplay{M03031R01}{%
= \psi(M^2+M+\psi_{\psi_{M}(M^2+M)}(M^2+M+\psi_{\psi_M(M^2+M)}(1)))%
}

\NovelDisplay{M03032R01}{%
(0)(1,1,1)(2,1,1)(3,1)(4,2)(2,1,1)(3,1)(3)(3)%
}

\NovelDisplay{M03033R01}{%
=\psi(\Omega_{I+1}+\psi_{\Omega_{I+1}}(\Omega_{I+1}+\psi_{\Omega_{I+1}}(2)))%
}

\NovelDisplay{M03034R01}{%
= \psi(M^2+M+\psi_{\psi_{M}(M^2+M)}(M^2+M+\psi_{\psi_M(M^2+M)}(2)))%
}

\NovelDisplay{M03035R01}{%
(0)(1,1,1)(2,1,1)(3,1)(4,2)(2,1,1)(3,1)(3,1)(2)%
}

\NovelDisplay{M03036R01}{%
=\psi(\Omega_{I+1}+\psi_{\Omega_{I+1}}(\Omega_{I+1}+\psi_{\Omega_{I+1}}(I)))%
}

\NovelDisplay{M03037R01}{%
= \psi(M^2+M+\psi_{\psi_{M}(M^2+M)}(M^2+M+\psi_{\psi_M(M^2+M)}(M^2)))%
}

内参数里的加法进一步嵌入另一层ψ，卢本韦先分清每一个函数头，才把相应的M侧读回去。

\NovelDisplay{M03038R01}{%
(0)(1,1,1)(2,1,1)(3,1)(4,2)(2,1,1)(3,1)(3,1)(2,1,1)(3,1)(3,1)(2)%
}

\NovelDisplay{M03039R01}{%
=\psi(\Omega_{I+1}+\psi_{\Omega_{I+1}}(\Omega_{I+1}+\psi_{\Omega_{I+1}}(I)\times 2))%
}

\NovelDisplay{M03040R01}{%
= \psi(M^2+M+\psi_{\psi_{M}(M^2+M)}(M^2+M+\psi_{\psi_M(M^2+M)}(M^2)\times 2))%
}

\NovelDisplay{M03041R01}{%
(0)(1,1,1)(2,1,1)(3,1)(4,2)(2,1,1)(3,1)(4)%
}

\NovelDisplay{M03042R01}{%
=\psi(\Omega_{I+1}+\psi_{\Omega_{I+1}}(\Omega_{I+1}+\psi_{\Omega_{I+1}}(I\times \omega)))%
}

\NovelDisplay{M03043R01}{%
= \psi(M^2+M+\psi_{\psi_{M}(M^2+M)}(M^2+M+\psi_{\psi_M(M^2+M)}(M^2+\psi_{\psi_{M}(M^2+M)}(1))))%
}

\NovelDisplay{M03044R01}{%
(0)(1,1,1)(2,1,1)(3,1)(4,2)(2,1,1)(3,1)(4)(2,1,1)(3,1)(4)%
}

\NovelDisplay{M03045R01}{%
=\psi(\Omega_{I+1}+\psi_{\Omega_{I+1}}(\Omega_{I+1}+\psi_{\Omega_{I+1}}(I\times \omega)\times 2))%
}

\NovelDisplay{M03046R01}{%
=\! \psi(M^2\!+\!M\!+\!\psi_{\psi_{M}(M^2+M)}(M^2\!+\!M\!+\!\psi_{\psi_M(M^2+M)}(M^2\!+\!\psi_{\psi_{M}(M^2+M)}(1))\!\times\! 2))%
}

\NovelDisplay{M03047R01}{%
(0)(1,1,1)(2,1,1)(3,1)(4,2)(2,1,1)(3,1)(4)(3)%
}

\NovelDisplay{M03048R01}{%
=\psi(\Omega_{I+1}+\psi_{\Omega_{I+1}}(\Omega_{I+1}+\psi_{\Omega_{I+1}}(I\times \omega+1)))%
}

\NovelDisplay{M03049R01}{%
=\! \psi(M^2\!+\!M\!+\!\psi_{\psi_{M}(M^2+M)}(M^2\!+\!M\!+\!\psi_{\psi_M(M^2+M)}(M^2\!+\!\psi_{\psi_{M}(M^2+M)}(1)\!+\!1)))%
}

蔡须坤指出，某个内函数值乘二与那个内函数的参数加一，改变的位置不同，不能被相似的长前缀遮住。

\NovelDisplay{M03050R01}{%
(0)(1,1,1)(2,1,1)(3,1)(4,2)(2,1,1)(3,1)(4)(3,1)(2)%
}

\NovelDisplay{M03051R01}{%
=\psi(\Omega_{I+1}+\psi_{\Omega_{I+1}}(\Omega_{I+1}+\psi_{\Omega_{I+1}}(I\times \omega+I)))%
}

\NovelDisplay{M03052R01}{%
= \psi(M^2+M+\psi_{\psi_{M}(M^2+M)}(M^2+M+\psi_{\psi_M(M^2+M)}(M^2%
}
\NovelDisplay{M03052R02}{%
+\psi_{\psi_{M}(M^2+M)}(1)+\psi_{M}(M^2))))%
}

\NovelDisplay{M03053R01}{%
(0)(1,1,1)(2,1,1)(3,1)(4,2)(2,1,1)(3,1)(4)(3,1)(4)%
}

\NovelDisplay{M03054R01}{%
=\psi(\Omega_{I+1}+\psi_{\Omega_{I+1}}(\Omega_{I+1}+\psi_{\Omega_{I+1}}(I\times \omega \times 2)))%
}

\NovelDisplay{M03055R01}{%
=\! \psi(M^2\!+\!M\!+\!\psi_{\psi_{M}(M^2+M)}(M^2\!+\!M\!+\!\psi_{\psi_M(M^2+M)}(M^2\!+\!\psi_{\psi_{M}(M^2+M)}(1)\!\times\! 2)))%
}

\NovelDisplay{M03056R01}{%
(0)(1,1,1)(2,1,1)(3,1)(4,2)(2,1,1)(3,1)(4)(4)%
}

\NovelDisplay{M03057R01}{%
=\psi(\Omega_{I+1}+\psi_{\Omega_{I+1}}(\Omega_{I+1}+\psi_{\Omega_{I+1}}(\psi_{\Omega_{I+1}}(2))))%
}

\NovelDisplay{M03058R01}{%
= \psi(M^2+M+\psi_{\psi_{M}(M^2+M)}(M^2+M+\psi_{\psi_M(M^2+M)}(M^2+\psi_{\psi_{M}(M^2+M)}(2))))%
}

\NovelDisplay{M03059R01}{%
(0)(1,1,1)(2,1,1)(3,1)(4,2)(2,1,1)(3,1)(4,1)(2)%
}

\NovelDisplay{M03060R01}{%
=\psi(\Omega_{I+1}+\psi_{\Omega_{I+1}}(\Omega_{I+1}+\psi_{\Omega_{I+1}}(\psi_{\Omega_{I+1}}(I))))%
}

\NovelDisplay{M03061R01}{%
=\! \psi(M^2\!+\!M\!+\!\psi_{\psi_{M}(M^2+M)}(M^2\!+\!M\!+\!\psi_{\psi_M(M^2+M)}(M^2\!+\!\psi_{\psi_{M}(M^2+M)}(M^2))))%
}

这几条继续在更深参数中加入I、重复项或新ψ值，所有已有括号都必须在当前结构里逐层闭合。

到了这一处，蔡须坤又读不下去了。

她已经能说清最外层，也知道末尾那一小段该落在哪里，偏偏中间套住它的下标，总在她准备往下接时散开。她靠着石壁坐下，怕别人看见自己又停住，先说要喝水。

古月娜在她旁边坐下，没有拆穿。她拿起那张M式，只指出两个同样长的函数头，问其中一个吃进的参数到哪里结束。

蔡须坤找了很久，终于将指尖落准。古月娜便把另外一处交还给她，自己接过水袋。

``这两层我陪你核。后面的，你再带我。''

\NovelDisplay{M03062R01}{%
(0)(1,1,1)(2,1,1)(3,1)(4,2)(2,1,1)(3,1)(4,2)%
}

\NovelDisplay{M03063R01}{%
=\psi(\Omega_{I+1}+\psi_{\Omega_{I+1}}(\Omega_{I+1}+\psi_{\Omega_{I+1}}(\Omega_{I+1})))%
}

\NovelDisplay{M03064R01}{%
= \psi(M^2+M+\psi_{\psi_{M}(M^2+M)}(M^2+M+\psi_{\psi_M(M^2+M)}(M^2+M)))%
}

\NovelDisplay{M03065R01}{%
(0)(1,1,1)(2,1,1)(3,1)(4,2)(2,1,1)(3,1)(4,2)(2,1,1)(3,1)(4,2)%
}

\NovelDisplay{M03066R01}{%
=\psi(\Omega_{I+1}+\psi_{\Omega_{I+1}}(\Omega_{I+1}+\psi_{\Omega_{I+1}}(\Omega_{I+1})\times 2))%
}

\NovelDisplay{M03067R01}{%
= \psi(M^2+M+\psi_{\psi_{M}(M^2+M)}(M^2+M+\psi_{\psi_M(M^2+M)}(M^2+M)\times 2))%
}

\NovelDisplay{M03068R01}{%
(0)(1,1,1)(2,1,1)(3,1)(4,2)(3)%
}

\NovelDisplay{M03069R01}{%
=\psi(\Omega_{I+1}+\psi_{\Omega_{I+1}}(\Omega_{I+1}+\psi_{\Omega_{I+1}}(\Omega_{I+1}+1)))%
}

\NovelDisplay{M03070R01}{%
= \psi(M^2+M+\psi_{\psi_{M}(M^2+M)}(M^2+M+\psi_{\psi_M(M^2+M)}(M^2+M+1)))%
}

\NovelDisplay{M03071R01}{%
(0)(1,1,1)(2,1,1)(3,1)(4,2)(3,1)(2)%
}

\NovelDisplay{M03072R01}{%
=\psi(\Omega_{I+1}+\psi_{\Omega_{I+1}}(\Omega_{I+1}+\psi_{\Omega_{I+1}}(\Omega_{I+1}+I)))%
}

\NovelDisplay{M03073R01}{%
= \psi(M^2+M+\psi_{\psi_{M}(M^2+M)}(M^2+M+\psi_{\psi_M(M^2+M)}(M^2+M+\psi_{M}(M^2))))%
}

古月娜将这一组概括为：同一个内层背景仍在，变化依次落进重复的ψ值和它的后继参数。

\NovelDisplay{M03074R01}{%
(0)(1,1,1)(2,1,1)(3,1)(4,2)(3,1)(3)%
}

\NovelDisplay{M03075R01}{%
=\psi(\Omega_{I+1}+\psi_{\Omega_{I+1}}(\Omega_{I+1}+\psi_{\Omega_{I+1}}(\Omega_{I+1}+I+1)))%
}

\NovelDisplay{M03076R01}{%
=\! \psi(M^2\!+\!M\!+\!\psi_{\psi_{M}(M^2+M)}(M^2\!+\!M\!+\!\psi_{\psi_M(M^2+M)}(M^2\!+\!M\!+\!\psi_{M}(M^2)\!+\!1)))%
}

\NovelDisplay{M03077R01}{%
(0)(1,1,1)(2,1,1)(3,1)(4,2)(3,1)(3,1)(2)%
}

\NovelDisplay{M03078R01}{%
=\psi(\Omega_{I+1}+\psi_{\Omega_{I+1}}(\Omega_{I+1}+\psi_{\Omega_{I+1}}(\Omega_{I+1}+I\times 2)))%
}

\NovelDisplay{M03079R01}{%
=\! \psi(M^2\!+\!M\!+\!\psi_{\psi_{M}(M^2+M)}(M^2\!+\!M\!+\!\psi_{\psi_M(M^2+M)}(M^2\!+\!M\!+\!\psi_{M}(M^2)\!\times\! 2)))%
}

\NovelDisplay{M03080R01}{%
(0)(1,1,1)(2,1,1)(3,1)(4,2)(3,1)(4)%
}

\NovelDisplay{M03081R01}{%
=\psi(\Omega_{I+1}+\psi_{\Omega_{I+1}}(\Omega_{I+1}+\psi_{\Omega_{I+1}}(\Omega_{I+1}+I\times 2)))%
}

\NovelDisplay{M03082R01}{%
=\! \psi(M^2\!+\!M\!+\!\psi_{\psi_{M}(M^2+M)}(M^2\!+\!M\!+\!\psi_{\psi_M(M^2+M)}(M^2\!+\!M\!+\!\psi_{M}(M^2)\!\times\! 2)))%
}

\NovelDisplay{M03083R01}{%
(0)(1,1,1)(2,1,1)(3,1)(4,2)(3,1)(4,1)(2)%
}

\NovelDisplay{M03084R01}{%
=\psi(\Omega_{I+1}+\psi_{\Omega_{I+1}}(\Omega_{I+1}+\psi_{\Omega_{I+1}}(\Omega_{I+1}+I^2)))%
}

\NovelDisplay{M03085R01}{%
= \psi(M^2+M+\psi_{\psi_{M}(M^2+M)}(M^2+M+\psi_{\psi_M(M^2+M)}(M^2+M%
}
\NovelDisplay{M03085R02}{%
+\psi_{\psi_{M}(M^2+M)}(M^2))))%
}

回填继续加深时，每次都要说明新增项属于哪层参数，不能用一句``继续嵌套''省去实际位置。

\NovelDisplay{M03086R01}{%
(0)(1,1,1)(2,1,1)(3,1)(4,2)(3,1)(4,2)%
}

\NovelDisplay{M03087R01}{%
=\psi(\Omega_{I+1}+\psi_{\Omega_{I+1}}(\Omega_{I+1}+\psi_{\Omega_{I+1}}(\Omega_{I+1}+\psi_{\Omega_{I+1}}(\Omega_{I+1}))))%
}

\NovelDisplay{M03088R01}{%
= \psi(M^2+M+\psi_{\psi_{M}(M^2+M)}(M^2+M+\psi_{\psi_M(M^2+M)}(M^2+M%
}
\NovelDisplay{M03088R02}{%
+\psi_{\psi_{M}(M^2+M)}(M^2+M))))%
}

\NovelDisplay{M03089R01}{%
(0)(1,1,1)(2,1,1)(3,1)(4,2)(4)%
}

\NovelDisplay{M03090R01}{%
=\psi(\Omega_{I+1}+\psi_{\Omega_{I+1}}(\Omega_{I+1}+\psi_{\Omega_{I+1}}(\Omega_{I+1}+\psi_{\Omega_{I+1}}(\Omega_{I+1}+1))))%
}

\NovelDisplay{M03091R01}{%
= \psi(M^2+M+\psi_{\psi_{M}(M^2+M)}(M^2+M+\psi_{\psi_M(M^2+M)}(M^2+M%
}
\NovelDisplay{M03091R02}{%
+\psi_{\psi_{M}(M^2+M)}(M^2+M+1))))%
}

\NovelDisplay{M03092R01}{%
(0)(1,1,1)(2,1,1)(3,1)(4,2)(4,1)%
}

\NovelDisplay{M03093R01}{%
=\psi(\Omega_{I+1}+\psi_{\Omega_{I+1}}(\Omega_{I+1}+\psi_{\Omega_{I+1}}(\Omega_{I+1}+\psi_{\Omega_{I+1}}(\Omega_{I+1}+I))))%
}

\NovelDisplay{M03094R01}{%
= \psi(M^2+M+\psi_{\psi_{M}(M^2+M)}(M^2+M+\psi_{\psi_M(M^2+M)}(M^2+M%
}
\NovelDisplay{M03094R02}{%
+\psi_{\psi_{M}(M^2+M)}(M^2+M+\psi_{M}(M^2)))))%
}

\NovelDisplay{M03095R01}{%
(0)(1,1,1)(2,1,1)(3,1)(4,2)(4,1)(5,2)%
}

\NovelDisplay{M03096R01}{%
=\psi(\Omega_{I+1}+\psi_{\Omega_{I+1}}(\Omega_{I+1}+\psi_{\Omega_{I+1}}(\Omega_{I+1}+\psi_{\Omega_{I+1}}(\Omega_{I+1}+\psi_{\Omega_{I+1}}(\Omega_{I+1})))))%
}

\NovelDisplay{M03097R01}{%
= \psi(M^2+M+\psi_{\psi_{M}(M^2+M)}(M^2+M+\psi_{\psi_M(M^2+M)}(M^2+M%
}
\NovelDisplay{M03097R02}{%
+\psi_{\psi_{M}(M^2+M)}(M^2+M+\psi_{\psi_{M}(M^2+M)}(M^2+M)))))%
}

这一组的M式虽长，核对的重点仍是最深处新增的参数与外面逐层保留的同一背景。

蔡须坤在山脚重做了两晚，第三次上来时，仍在同一个接缝停了一下。

她心里猛地往下一沉，以为之前那些功夫全没留住。霍雨浩却将昨天那张纸取来，让她看自己今天已经独立接到哪里。

那处停顿，比昨天更靠里。

她把两张纸放在一起，才觉出掌心一直压着冷汗。原来困在长式里时，连自己确实往前走了一点，也会看不见。

卢本韦将她刚才说对的前半段念给唐三听，蔡须坤听完，自己接上了后面那层。三百一十米的刻线就在脚边，她这一次没有低头躲开。

\NovelDisplay{M03098R01}{%
(0)(1,1,1)(2,1,1)(3,1)(4,2)(4,2)%
}

\NovelDisplay{M03099R01}{%
=\psi(\Omega_{I+1}\times 2)%
}

\NovelDisplay{M03100R01}{%
= \psi(M^2+M+\psi_{M}(M^2+M))%
}

\NovelDisplay{M03101R01}{%
(0)(1,1,1)(2,1,1)(3,1)(4,2)(4,2)(2,1)(3,2,1)(4,2,1)(5,2)(4)%
}

\NovelDisplay{M03102R01}{%
=\psi(\Omega_{I+1}\times 2+I)%
}

\NovelDisplay{M03103R01}{%
= \psi(M^2+M+\psi_{M}(M^2+M)+\psi_{M}(M^2))%
}

\NovelDisplay{M03104R01}{%
(0)(1,1,1)(2,1,1)(3,1)(4,2)(4,2)(2,1)(3,2,1)(4,2,1)(5,2)(4)%
}

\NovelDisplay{M03105R01}{%
=\psi(\Omega_{I+1}\times 2+I)%
}

\NovelDisplay{M03106R01}{%
= \psi(M^2+M+\psi_{M}(M^2+M)+\psi_{M}(M^2))%
}

\NovelDisplay{M03107R01}{%
(0)(1,1,1)(2,1,1)(3,1)(4,2)(4,2)(2,1)(3,2,1)(4,2,1)(5,2)(6,3)%
}

\NovelDisplay{M03108R01}{%
=\psi(\Omega_{I+1}\times 2+\psi_{\Omega_{I+1}}(\Omega_{I+1}))%
}

\NovelDisplay{M03109R01}{%
= \psi(M^2+M+\psi_{M}(M^2+M)+\psi_{\psi_{M}(M^2+M)}(M^2+M))%
}

卢本韦说：``这里终于出现整份Ω下标对象的重复，但重复两条原式只帮助复算，不会再多添一次增长。''

\NovelDisplay{M03110R01}{%
(0)(1,1,1)(2,1,1)(3,1)(4,2)(4,2)(2,1)(3,2,1)(4,2,1)(5,2)(6,3)(6,3)%
}

\NovelDisplay{M03111R01}{%
=\psi(\Omega_{I+1}\times 2+\psi_{\Omega_{I+1}}(\Omega_{I+1}\times 2))%
}

\NovelDisplay{M03112R01}{%
= \psi(M^2+M+\psi_{M}(M^2+M)+\psi_{\psi_{M}(M^2+M)}(M^2+M+\psi_{M}(M^2+M)))%
}

\NovelDisplay{M03113R01}{%
(0)(1,1,1)(2,1,1)(3,1)(4,2)(4,2)(2,1,1)%
}

\NovelDisplay{M03114R01}{%
=\psi(\Omega_{I+1}\times 2+\psi_{\Omega_{I+1}}(\Omega_{I+1}\times 2+1))%
}

\NovelDisplay{M03115R01}{%
=\! \psi(M^2\!+\!M\!+\!\psi_{M}(M^2\!+\!M)\!+\!\psi_{\psi_{M}(M^2+M)}(M^2\!+\!M\!+\!\psi_{M}(M^2\!+\!M)\!+\!1))%
}

\NovelDisplay{M03116R01}{%
(0)(1,1,1)(2,1,1)(3,1)(4,2)(4,2)(3)%
}

\NovelDisplay{M03117R01}{%
=\psi(\Omega_{I+1}\times 2+\psi_{\Omega_{I+1}}(\Omega_{I+1}\times 2+\psi_{\Omega_{I+1}}(\Omega_{I+1}\times 2+1)))%
}

\NovelDisplay{M03118R01}{%
= \psi(M^2+M+\psi_{M}(M^2+M)+\psi_{\psi_{M}(M^2+M)}(M^2+M+\psi_{M}(M^2+M)%
}
\NovelDisplay{M03118R02}{%
+\psi_{\psi_{M}(M^2+M)}(M^2+M+\psi_{M}(M^2+M)+1)))%
}

\NovelDisplay{M03119R01}{%
(0)(1,1,1)(2,1,1)(3,1)(4,2)(4,2)(4)%
}

\NovelDisplay{M03120R01}{%
=\psi(\Omega_{I+1}\times 2+\psi_{\Omega_{I+1}}(\Omega_{I+1}\times 2+\psi_{\Omega_{I+1}}(\Omega_{I+1}\times 2+\psi_{\Omega_{I+1}}(\Omega_{I+1}\times 2+1))))%
}

\NovelDisplay{M03121R01}{%
= \psi(M^2+M+\psi_{M}(M^2+M)+\psi_{\psi_{M}(M^2+M)}(M^2+M+\psi_{M}(M^2+M)%
}
\NovelDisplay{M03121R02}{%
+\psi_{\psi_{M}(M^2+M)}(M^2+M+\psi_{M}(M^2+M)+\psi_{\psi_{M}(M^2+M)}(M^2+M%
}
\NovelDisplay{M03121R03}{%
+\psi_{M}(M^2+M)+1))))%
}

从两份背景继续向内回填时，外层的第二份背景仍需完整保留，不能被后面相似的函数头吞掉。

\NovelDisplay{M03122R01}{%
(0)(1,1,1)(2,1,1)(3,1)(4,2)(4,2)(4,2)%
}

\NovelDisplay{M03123R01}{%
=\psi(\Omega_{I+1}\times 3)%
}

\NovelDisplay{M03124R01}{%
= \psi(M^2+M+\psi_{M}(M^2+M)\times 2)%
}

\NovelDisplay{M03125R01}{%
(0)(1,1,1)(2,1,1)(3,1)(4,2)(5)%
}

\NovelDisplay{M03126R01}{%
=\psi(\Omega_{I+1}\times \omega)%
}

\NovelDisplay{M03127R01}{%
= \psi(M^2+M+\psi_{\psi_{M}(M^2+M\times 2)}(1))%
}

\NovelDisplay{M03128R01}{%
(0)(1,1,1)(2,1,1)(3,1)(4,2)(5,1)(2)%
}

\NovelDisplay{M03129R01}{%
=\psi(\Omega_{I+1}\times I)%
}

\NovelDisplay{M03130R01}{%
= \psi(M^2+M+\psi_{\psi_{M}(M^2+M\times 2)}(M^2))%
}

\NovelDisplay{M03131R01}{%
(0)(1,1,1)(2,1,1)(3,1)(4,2)(5,1)(6,2)%
}

\NovelDisplay{M03132R01}{%
=\psi(\Omega_{I+1}\times \psi_{\Omega_{I+1}}(\Omega_{I+1}))%
}

\NovelDisplay{M03133R01}{%
= \psi(M^2+M+\psi_{\psi_{M}(M^2+M\times 2)}(M^2+\psi_{\psi_{M}(M^2+M)}(M^2+M)))%
}

霍雨浩归纳道：``有限重复、右乘ω与右乘I在I侧分得很开，M侧也必须逐项说明它们各自落在哪个参数。''

一阵山风从石隙里穿过，叶阳手中的纸滑了出去。

唐三伏身抓住最上面一张，另一张却已经贴着石面向下掠。古月娜一手扶住蔡须坤，另一只手按在石沿，银光停在纸边，将它带回了平处。

谁也没有顺着那一下力道往上走。山阶里的分析还要各自接，古月娜只替他们保住了纸。

叶阳将湿了的一角晾开，拿出自己的抄页。那一段霍雨浩刚好没能抄全，便向他借过来补。平日安静留在后方的人，此刻也替前面的人接住了一个缺口。

\NovelDisplay{M03134R01}{%
(0)(1,1,1)(2,1,1)(3,1)(4,2)(5,1)(6,2)(7)%
}

\NovelDisplay{M03135R01}{%
=\psi(\Omega_{I+1}\times \psi_{\Omega_{I+1}}(\Omega_{I+1}\times \omega))%
}

\NovelDisplay{M03136R01}{%
=\! \psi(M^2\!+\!M\!+\!\psi_{\psi_{M}(M^2+M\times 2)}(M^2\!+\!\psi_{\psi_{M}(M^2+M)}(M^2\!+\!M\!+\!\psi_{\psi_{M}(M^2+M)}(1))))%
}

\NovelDisplay{M03137R01}{%
(0)(1,1,1)(2,1,1)(3,1)(4,2)(5,1)(6,2)(7,1)(8,2)%
}

\NovelDisplay{M03138R01}{%
=\psi(\Omega_{I+1}\times \psi_{\Omega_{I+1}}(\Omega_{I+1}\times \psi_{\Omega_{I+1}}(\Omega_{I+1})))%
}

\NovelDisplay{M03139R01}{%
= \psi(M^2+M+\psi_{\psi_{M}(M^2+M\times 2)}(M^2+\psi_{\psi_{M}(M^2+M)}(M^2+M%
}
\NovelDisplay{M03139R02}{%
+\psi_{\psi_{M}(M^2+M)}(M^2+\psi_{\psi_{M}(M^2+M)}(M^2+M)))))%
}

\NovelDisplay{M03140R01}{%
(0)(1,1,1)(2,1,1)(3,1)(4,2)(5,2)%
}

\NovelDisplay{M03141R01}{%
=\psi(\Omega_{I+1}^2)%
}

\NovelDisplay{M03142R01}{%
= \psi(M^2+M+\psi_{\psi_{M}(M^2+M\times 2)}(M^2+M))%
}

\NovelDisplay{M03143R01}{%
(0)(1,1,1)(2,1,1)(3,1)(4,2)(5,2)(5)%
}

\NovelDisplay{M03144R01}{%
=\psi(\Omega_{I+1}^2\times \omega)%
}

\NovelDisplay{M03145R01}{%
= \psi(M^2+M+\psi_{\psi_{M}(M^2+M\times 2)}(M^2+M+1))%
}

同一背景的乘数继续增长直到平方，读到平方时应回查乘数所属层次，不能仅凭公式变短就判断已经清楚。

夜里，卢本韦在灯下替蔡须坤誊清一张被水洇过的纸。

她认得大部分字，却在一个函数头那里停了很久，抄到后面，发现右括号合不上。蔡须坤洗过手回来，看她把原纸举到灯边，忍不住笑出了声。

``这里我自己就漏过。''

``那你还留着，让我再漏一遍？''

蔡须坤把纸放低，将自己后来补在背面的那一处翻给她。卢本韦看完，把笔往她手里一放，要她坐下说清楚再走。

说到第二遍，蔡须坤忽然停住，发现自己原来只会沿熟悉的顺序读，换成回答卢本韦的问题，仍有一截没讲完整。她脸上的笑慢慢收了，重新找来旧册。

卢本韦没有催，自己也把I册打开，留在同一张桌上。

两人读完，灯芯已短了许多。蔡须坤靠着椅背缓了片刻，先说今天就到这里。卢本韦替她收笔，指了指仍在门边等着的热水，两人这才发现，奥斯卡睡前还给她们各留了一碗。

水上还浮着一点热气。蔡须坤端起来，手指贴住碗壁，才觉出自己刚才一直握着凉笔。卢本韦将湿纸从灯边挪开，找了一处不会压到字的地方晾好。

``明早我先看这一页。''她说。

蔡须坤答应，又补了一句，要是自己仍没有醒，就先叫她吃饭。

等她合上本子，卢本韦又将笔放回盒里，让她明早醒来就能找见；两人此刻都记得，今天不必再把所有问题带进梦中。

卢本韦看了她一眼，笑着把碗碰过来。两人坐着将水喝完，才各自起身。门外的风已经小了，隔壁传来很轻的翻身声，蔡须坤收起最后一张纸时，也将明天还没接完的那一处好好留着，没有再想当夜把它赶完。

\NovelDisplay{M03146R01}{%
(0)(1,1,1)(2,1,1)(3,1)(4,2)(5,2)(5,2)%
}

\NovelDisplay{M03147R01}{%
=\psi(\Omega_{I+1}^3)%
}

\NovelDisplay{M03148R01}{%
= \psi(M^2+M+\psi_{\psi_{M}(M^2+M\times 2)}(M^2+M+\psi_{M}(M^2+M)))%
}

\NovelDisplay{M03149R01}{%
(0)(1,1,1)(2,1,1)(3,1)(4,2)(5,2)(6)%
}

\NovelDisplay{M03150R01}{%
=\psi(\Omega_{I+1}^\omega)%
}

\NovelDisplay{M03151R01}{%
= \psi(M^2+M+\psi_{\psi_{M}(M^2+M\times 2)}(M^2+M+\psi_{\psi_{M}(M^2+M\times 2)}(1)))%
}

\NovelDisplay{M03152R01}{%
(0)(1,1,1)(2,1,1)(3,1)(4,2)(5,2)(6,2)%
}

\NovelDisplay{M03153R01}{%
=\psi(\Omega_{I+1}^{\Omega_{I+1}})%
}

\NovelDisplay{M03154R01}{%
= \psi(M^2+M+\psi_{\psi_{M}(M^2+M\times 2)}(M^2+M+\psi_{\psi_{M}(M^2+M\times 2)}(M^2+M)))%
}

\NovelDisplay{M03155R01}{%
(0)(1,1,1)(2,1,1)(3,1)(4,2)(5,3)%
}

\NovelDisplay{M03156R01}{%
=\psi(\Omega_{I+2})%
}

\NovelDisplay{M03157R01}{%
= \psi(M^2+M\times 2)%
}

蔡须坤指出，这里从旧下标的高次幂转到\allowbreak{}\NovelInline{I06040}{\Omega_{I+2}}\allowbreak{}，新下标与旧对象的乘方需要分别辨认。

\NovelDisplay{M03158R01}{%
(0)(1,1,1)(2,1,1)(3,1)(4,2)(5,3)(4,2)%
}

\NovelDisplay{M03159R01}{%
=\psi(\Omega_{I+2}+\Omega_{I+1})%
}

\NovelDisplay{M03160R01}{%
= \psi(M^2+M\times 2+\psi_{M}(M^2+M))%
}

\NovelDisplay{M03161R01}{%
(0)(1,1,1)(2,1,1)(3,1)(4,2)(5,3)(4,2)(5,3)%
}

\NovelDisplay{M03162R01}{%
=\psi(\Omega_{I+2}+\psi_{\Omega_{I+2}}(\Omega_{I+2}))%
}

\NovelDisplay{M03163R01}{%
= \psi(M^2+M\times 2+\psi_{\psi_{M}(M^2+M\times 2)}(M^2+M\times 2))%
}

\NovelDisplay{M03164R01}{%
(0)(1,1,1)(2,1,1)(3,1)(4,2)(5,3)(5,3)%
}

\NovelDisplay{M03165R01}{%
=\psi(\Omega_{I+2}\times 2)%
}

\NovelDisplay{M03166R01}{%
= \psi(M^2+M\times 2+\psi_{M}(M^2+M\times 2))%
}

\NovelDisplay{M03167R01}{%
(0)(1,1,1)(2,1,1)(3,1)(4,2)(5,3)(6,4)%
}

\NovelDisplay{M03168R01}{%
=\psi(\Omega_{I+3})%
}

\NovelDisplay{M03169R01}{%
= \psi(M^2+M\times 3)%
}

新的Ω下标之后又出现尾加、内函数和整份重复，相邻的短结果仍由各自完整矩阵给出。

到三百二十米那天，雾没有散。

霍雨浩最先摸到石桩侧面的刻线，却没有立即报数。他回身看蔡须坤和卢本韦将最后一层各自接完，又等薛峰将原左列复述清楚，才用指节敲了敲那块石头。

``这里。''

蔡须坤伸手摸过去，指腹擦过两个浅浅的转折，才认出自己这十几天一直惦记的高度。卢本韦在她身边长出一口气，笑到一半，先被自己的喘息打断了。

唐三落在后面几阶，正在重新接一处下标。蔡须坤没有叫他快一点，她将本子放到他看得清的位置，等他讲完，再陪他核了一次。小舞在另一边收稳脚跟，古月娜将她的包接下，留出双手给她翻页。

十一人陆续站到横台上时，已经过了午后。

奥斯卡将剩下的水分完。蔡兴坤把黑子鸡放在背风处，诺当靠着石桩，半晌没有说话；叶阳替大家一张张收好被雾打湿的纸，薛峰将今天那处未定的背景项另夹了一页。

蔡须坤望着这些纸，忽然想起那张写不出第一项的理想矩阵。

它还没有答案。可她今天不再只会望着空白发愣。唐三早晨那道旧题、小舞独立求出的有限项、卢本韦让她重新辨认的一层括号，都已经留在这二十米里。

她把封套在包中压稳，将当日记录递给霍雨浩。

``下一段，还是一起。''

霍雨浩接过笔，在三百二十旁边落下一点。古月娜伸手扶她起身，卢本韦在另一侧背好本子；他们歇过这一阵，才沿旧路慢慢下山。

\chapter{第一百一十九章：把长路接在一起}\label{ux7b2cux4e00ux767eux4e00ux5341ux4e5dux7ae0ux628aux957fux8defux63a5ux5728ux4e00ux8d77}

到三百二十米以后，蔡须坤终于睡了一个没有被长式惊醒的觉。

第二天早上，她先把昨天的纸翻了一遍，再去叫卢本韦。对方已经醒了，正坐在床沿系鞋带，听见门响，便将I册递过来，让她看昨晚新添的两处说明。

蔡须坤站着读完，指出其中一处还没有说清的位置。卢本韦没有马上翻答案，拿着纸想了片刻，把那一页带到早饭桌上。

这一次，屋里没人急着问今天能走到哪里。唐三把自己的疑问放在桌边，小舞先替他找原式，薛峰边吃边核共同的前缀，直到奥斯卡敲了敲锅沿，几人才想起碗还拿在手里。

蔡须坤将那张纸收进包里，出了门，回头看见十一人的脚步陆续跟了上来。

三百二十米以上，卢本韦的脚跟磨破了。

她先没有说，直到小舞看见鞋边的颜色，才被大家叫回较宽的石阶。蔡须坤蹲下来替她解鞋带，针脚上一次补得歪，这次却牢牢撑着鞋底，没有再开。

``歇的时候也能问。''卢本韦将I册放在膝上。

``先让她包好。''古月娜把布条递过去。

大家便坐成半圈，小舞先读，唐三接问，薛峰将两人分歧的那处单独指出。卢本韦听了好一阵才插话，发现同伴已经把她原先要解释的部分自己说清了。

\NovelDisplay{M03170R01}{%
(0)(1,1,1)(2,1,1)(3,1)(4,2)(5,3)(6,4)(7,5)%
}

\NovelDisplay{M03171R01}{%
=\psi(\Omega_{I+4})%
}

\NovelDisplay{M03172R01}{%
= \psi(M^2+M\times 4)%
}

\NovelDisplay{M03173R01}{%
(0)(1,1,1)(2,1,1)(3,1)(4,2,1)%
}

\NovelDisplay{M03174R01}{%
=\psi(\Omega_{I+\omega})%
}

\NovelDisplay{M03175R01}{%
= \psi(M^2+M+\psi_{M}(M^2+M\times \omega))%
}

\NovelDisplay{M03176R01}{%
(0)(1,1,1)(2,1,1)(3,1)(4,2,1)(4,2,1)%
}

\NovelDisplay{M03177R01}{%
=\psi(\Omega_{I+\omega}\times 2)%
}

\NovelDisplay{M03178R01}{%
= \psi(M^2+M\times \omega+\psi_{M}(M^2+M\times \omega))%
}

\NovelDisplay{M03179R01}{%
(0)(1,1,1)(2,1,1)(3,1)(4,2,1)(5,2)%
}

\NovelDisplay{M03180R01}{%
=\psi(\Omega_{I+\omega}\times \Omega_{I+1})%
}

\NovelDisplay{M03181R01}{%
= \psi(M^2+M\times \omega+\psi_{\psi_{M}(M^2+M\times \omega+M)}(M^2+M))%
}

这组进入\allowbreak{}\NovelInline{I06041}{I+\omega}\allowbreak{}下标并继续乘法，其中第一条M侧的背景写法与相邻项有待另核，原式保留而不据此推出通则。

\NovelDisplay{M03182R01}{%
(0)(1,1,1)(2,1,1)(3,1)(4,2,1)(5,2)(4,2,1)%
}

\NovelDisplay{M03183R01}{%
=\psi(\Omega_{I+\omega}^2)%
}

\NovelDisplay{M03184R01}{%
= \psi(M^2+M\times \omega+\psi_{\psi_{M}(M^2+M\times \omega+M)}(M^2+M\times \omega))%
}

\NovelDisplay{M03185R01}{%
(0)(1,1,1)(2,1,1)(3,1)(4,2,1)(5,2)(6,3)%
}

\NovelDisplay{M03186R01}{%
=\psi(\Omega_{I+\omega+1})%
}

\NovelDisplay{M03187R01}{%
= \psi(M^2+M\times \omega+M)%
}

\NovelDisplay{M03188R01}{%
(0)(1,1,1)(2,1,1)(3,1)(4,2,1)(5,2)(6,3)(7,4)%
}

\NovelDisplay{M03189R01}{%
=\psi(\Omega_{I+\omega+2})%
}

\NovelDisplay{M03190R01}{%
= \psi(M^2+M\times \omega+M\times 2)%
}

\NovelDisplay{M03191R01}{%
(0)(1,1,1)(2,1,1)(3,1)(4,2,1)(5,2)(6,3,1)%
}

\NovelDisplay{M03192R01}{%
=\psi(\Omega_{I+\omega\times 2})%
}

\NovelDisplay{M03193R01}{%
= \psi(M^2+M\times \omega \times 2)%
}

有限后继与ω倍数进入下标后，M参数中的\allowbreak{}\NovelInline{I06042}{M}\allowbreak{}、\allowbreak{}\NovelInline{I06043}{M\times\omega}\allowbreak{}及其追加项必须按出现位置逐个读清。

\NovelDisplay{M03194R01}{%
(0)(1,1,1)(2,1,1)(3,1)(4,2,1)(5,2)(6,3,1)(7,3)(8,4,1)%
}

\NovelDisplay{M03195R01}{%
=\psi(\Omega_{I+\omega\times 3})%
}

\NovelDisplay{M03196R01}{%
= \psi(M^2+M\times \omega \times 3)%
}

\NovelDisplay{M03197R01}{%
(0)(1,1,1)(2,1,1)(3,1)(4,2,1)(5,2,1)%
}

\NovelDisplay{M03198R01}{%
=\psi(\Omega_{I+\omega^2})%
}

\NovelDisplay{M03199R01}{%
= \psi(M^2+M\times \psi(2))%
}

\NovelDisplay{M03200R01}{%
(0)(1,1,1)(2,1,1)(3,1)(4,2,1)(5,2,1)(5,2,1)%
}

\NovelDisplay{M03201R01}{%
=\psi(\Omega_{I+\omega^3})%
}

\NovelDisplay{M03202R01}{%
= \psi(M^2+M\times \psi(3))%
}

\NovelDisplay{M03203R01}{%
(0)(1,1,1)(2,1,1)(3,1)(4,2,1)(5,2,1)(6,1)%
}

\NovelDisplay{M03204R01}{%
=\psi(\Omega_{I+\Omega})%
}

\NovelDisplay{M03205R01}{%
= \psi(M^2+M\times \psi_{\psi_{M}(M^2)}(1))%
}

卢本韦归纳道：``ω的幂、Ω和I在下标中依次出现时，M侧对应的函数头也在变化，不能只换最后一个数字。''

走出三百二十米之后，蔡须坤第一次在别人提出问题之前，自己停在了式中一处背景项上。

讲义的相邻对应看起来十分接近，M侧却不能靠挪一个字便连成整套通则。她将那一条完整留在原处，旁边写下需要复核的地方，没有为赶路擅自改成更顺眼的答案。

薛峰接过来，将原左列也一并留下。霍雨浩请其余人各自读当前明确能说清的原式，未定的一处没有被拿去支撑整段新结论。

这一次，蔡须坤没有觉得留下疑问便等于后退。她已能分清，哪一页需要继续追，哪一段确实可以自己走。

\NovelDisplay{M03206R01}{%
(0)(1,1,1)(2,1,1)(3,1)(4,2,1)(5,2,1)(6,1)(2)%
}

\NovelDisplay{M03207R01}{%
=\psi(\Omega_{I\times 2})%
}

\NovelDisplay{M03208R01}{%
= \psi(M^2+M\times \psi_{M}(M^2))%
}

\NovelDisplay{M03209R01}{%
(0)(1,1,1)(2,1,1)(3,1)(4,2,1)(5,2,1)(6,1)(7,2)%
}

\NovelDisplay{M03210R01}{%
=\psi(\Omega_{\psi_{\Omega_{I+1}}(\Omega_{I+1})})%
}

\NovelDisplay{M03211R01}{%
= \psi(M^2+M\times\psi_{\psi_{M}(M^2+M)}(M^2+M))%
}

\NovelDisplay{M03212R01}{%
(0)(1,1,1)(2,1,1)(3,1)(4,2,1)(5,2,1)(6,2)%
}

\NovelDisplay{M03213R01}{%
= \psi(\Omega_{\Omega_{I+1}})%
}

\NovelDisplay{M03214R01}{%
= \psi(M^2+M\times\psi_{M}(M^2+M))%
}

\NovelDisplay{M03215R01}{%
(0)(1,1,1)(2,1,1)(3,1)(4,2,1)(5,2,1)(6,2)(4,2,1)(5,2,1)(6,2)%
}

\NovelDisplay{M03216R01}{%
= \psi(\Omega_{\Omega_{\Omega_{I+1}}})%
}

\NovelDisplay{M03217R01}{%
= \psi(M^2+M\times\psi_{M}(M^2+M\times\psi_{M}(M^2+M)))%
}

最后四条把下标继续放进ψ值和更深的Ω下标，已经核清的旧背景始终是辨认新层次的起点。

随后一次晨练时，卢本韦先认出了下一页那个很短的结果。

她将I册摆正，念出下标二，没有急着往上迈。经过上一段，她已经知道，读得出名字与在压力下完整接续，中间仍有一截要自己走。

蔡须坤在另一侧摊开M册。最外面是M平方的两份，写起来比前面许多嵌套短，她却将旧页也翻到旁边，确认这一次变化已经落到整个外层参数。

唐三随两人一起看。小舞蹲在较低一级，替薛峰按住卷起的纸角。山里前一夜又落过雨，石阶中间还存着一层凉意，谁的膝盖停得久了，都要慢慢伸直才能再弯下去。

霍雨浩用手指挡住结果，请唐三先指出末列要去哪里找父项。唐三向前数到一半，自己退回一列，重新开始；蔡须坤一直等他指准，才把右侧露出来。

``这次先看下标。''卢本韦说。

薛峰在页边记下两种容易混的读法，又将笔递给小舞。今天他不再只替前面四个人收纸，每一处自己能判断的关系，都请另外一人独立说过。

山阶上方仍被雾遮着。十一人的绳子松松留在身侧，能让人扶稳，里面的接续却一人也省不掉。

\NovelDisplay{M03218R01}{%
(0)(1,1,1)(2,1,1)(3,1)(4,2,1)(5,2,1)(6,2)(5)%
}

\NovelDisplay{M03219R01}{%
= \psi(I_2)%
}

\NovelDisplay{M03220R01}{%
= \psi(M^2\times 2)%
}

\NovelDisplay{M03221R01}{%
(0)(1,1,1)(2,1,1)(3,1)(4,2,1)(5,2,1)(6,2)(5)(2,1)(3,2,1)(4,2,1)(5,2)(4)%
}

\NovelDisplay{M03222R01}{%
= \psi(I_2+I)%
}

\NovelDisplay{M03223R01}{%
= \psi(M^2\times 2+\psi_{M}(M^2))%
}

\NovelDisplay{M03224R01}{%
(0)(1,1,1)(2,1,1)(3,1)(4,2,1)(5,2,1)(6,2)(5)(2,1)(3,2,1)(4,2,1)%
}
\NovelDisplay{M03224R02}{%
(5,2)(6,3)%
}

\NovelDisplay{M03225R01}{%
= \psi(I_2+\psi_{\Omega_{I+1}}(\Omega_{I+1}))%
}

\NovelDisplay{M03226R01}{%
= \psi(M^2\times 2+\psi_{\psi_{M}(M^2+M)}(M^2+M))%
}

\NovelDisplay{M03227R01}{%
(0)(1,1,1)(2,1,1)(3,1)(4,2,1)(5,2,1)(6,2)(5)(2,1)(3,2,1)(4,2,1)(5,2)%
}
\NovelDisplay{M03227R02}{%
(6,3,1)(7,3,1)(8,3)(7)%
}

\NovelDisplay{M03228R01}{%
= \psi(I_2+\psi_{\Omega_{I+1}}(I_2))%
}

\NovelDisplay{M03229R01}{%
= \psi(M^2\times 2+\psi_{\psi_{M}(M^2+M)}(M^2 \times 2))%
}

卢本韦归纳道：``\allowbreak{}\NovelInline{I06044}{I_2}\allowbreak{}是下标二的对象，随后出现的尾加与内函数都从这个背景接起，不能先把它读作\allowbreak{}\NovelInline{I06045}{I^2}\allowbreak{}。''

\NovelDisplay{M03230R01}{%
(0)(1,1,1)(2,1,1)(3,1)(4,2,1)(5,2,1)(6,2)(5)(2,1)(3,2,1)(4,2,1)(5,2)%
}
\NovelDisplay{M03230R02}{%
(6,3,1)(7,3,1)(8,3)(7)(4,2)(5,3,1)(6,3,1)(7,3)(6)%
}

\NovelDisplay{M03231R01}{%
= \psi(I_2+\psi_{\Omega_{I+1}}(I_2)+I)%
}

\NovelDisplay{M03232R01}{%
= \psi(M^2\times 2+\psi_{\psi_{M}(M^2+M)}(M^2 \times 2)+\psi_{M}(M^2))%
}

\NovelDisplay{M03233R01}{%
(0)(1,1,1)(2,1,1)(3,1)(4,2,1)(5,2,1)(6,2)(5)(2,1)(3,2,1)(4,2,1)(5,2)%
}
\NovelDisplay{M03233R02}{%
(6,3,1)(7,3,1)(8,3)(7)(4,2)(5,3,1)(6,3,1)(7,3)(8,4,1)(9,4,1)(10,4)(9)%
}

\NovelDisplay{M03234R01}{%
= \psi(I_2+\psi_{\Omega_{I+1}}(I_2)\times 2)%
}

\NovelDisplay{M03235R01}{%
= \psi(M^2\times 2+\psi_{\psi_{M}(M^2+M)}(M^2 \times 2)\times 2)%
}

\NovelDisplay{M03236R01}{%
(0)(1,1,1)(2,1,1)(3,1)(4,2,1)(5,2,1)(6,2)(5)(2,1,1)%
}

\NovelDisplay{M03237R01}{%
= \psi(I_2+\psi_{\Omega_{I+1}}(I_2+1))%
}

\NovelDisplay{M03238R01}{%
= \psi(M^2\times 2+\psi_{\psi_{M}(M^2+M)}(M^2 \times 2 + 1))%
}

\NovelDisplay{M03239R01}{%
(0)(1,1,1)(2,1,1)(3,1)(4,2,1)(5,2,1)(6,2)(5)(2,1,1)(3,1)(2)%
}

\NovelDisplay{M03240R01}{%
= \psi(I_2+\psi_{\Omega_{I+1}}(I_2+I))%
}

\NovelDisplay{M03241R01}{%
= \psi(M^2\times 2+\psi_{\psi_{M}(M^2+M)}(M^2 \times 2 + \psi_{M}(M^2)))%
}

同一背景之后，另加I、重复整个函数值、以及把I放进内参数，必须各自保留不同的括号位置。

\NovelDisplay{M03242R01}{%
(0)(1,1,1)(2,1,1)(3,1)(4,2,1)(5,2,1)(6,2)(5)(2,1,1)(3,1)(2)%
}

\NovelDisplay{M03243R01}{%
= \psi(I_2+\psi_{\Omega_{I+1}}(I_2+I))%
}

\NovelDisplay{M03244R01}{%
= \psi(M^2\times 2+\psi_{\psi_{M}(M^2+M)}(M^2 \times 2 + \psi_{M}(M^2)))%
}

\NovelDisplay{M03245R01}{%
(0)(1,1,1)(2,1,1)(3,1)(4,2,1)(5,2,1)(6,2)(5)(2,1,1)(3,1)(4,2,1)%
}
\NovelDisplay{M03245R02}{%
(5,2,1)(6,2)(5)%
}

\NovelDisplay{M03246R01}{%
= \psi(I_2+\psi_{\Omega_{I+1}}(I_2+\psi_{\Omega_{I+1}}(I_2)))%
}

\NovelDisplay{M03247R01}{%
= \psi(M^2\times 2+\psi_{\psi_{M}(M^2+M)}(M^2 \times 2 + \psi_{\psi_{M}(M^2+M)}(M^2\times 2)))%
}

\NovelDisplay{M03248R01}{%
(0)(1,1,1)(2,1,1)(3,1)(4,2,1)(5,2,1)(6,2)(5)(3)%
}

\NovelDisplay{M03249R01}{%
= \psi(I_2+\psi_{\Omega_{I+1}}(I_2+\psi_{\Omega_{I+1}}(I_2+1)))%
}

\NovelDisplay{M03250R01}{%
= \psi(M^2\times 2+\psi_{\psi_{M}(M^2+M)}(M^2 \times 2 + \psi_{\psi_{M}(M^2+M)}(M^2\times 2+1)))%
}

\NovelDisplay{M03251R01}{%
(0)(1,1,1)(2,1,1)(3,1)(4,2,1)(5,2,1)(6,2)(5)(3,1)(4,2,1)(5,2,1)(6,2)(5)%
}

\NovelDisplay{M03252R01}{%
= \psi(I_2+\psi_{\Omega_{I+1}}(I_2+\psi_{\Omega_{I+1}}(I_2+\psi_{\Omega_{I+1}}(I_2))))%
}

\NovelDisplay{M03253R01}{%
= \psi(M^2\times 2+\psi_{\psi_{M}(M^2+M)}(M^2 \times 2 + \psi_{\psi_{M}(M^2+M)}(M^2\times 2%
}
\NovelDisplay{M03253R02}{%
+\psi_{\psi_{M}(M^2+M)}(M^2\times 2))))%
}

薛峰指出，组首是上一条原式的重复复算，后面的真正变化才是新加的一层内参数回填。

蔡须坤在一处内参数上接快了。

魂力从刚认清的前段向后续时，胸口忽然像被细线扯住，后面几层一齐往回收。她立即停下，脚尖却已经探到上一阶边缘，重心一偏，卢本韦伸手托住了她的肘。

``退回来，先坐。''

她坐稳后，仍觉得刚才自己只错了一点。古月娜把那一对括号指出来，让她看清那一点究竟包着多大一段；霍雨浩等她将呼吸调匀，才请她从出错之前的位置重新说。

这一天只多走了几阶，蔡须坤却把那处停顿认真记了下来。

\NovelDisplay{M03254R01}{%
(0)(1,1,1)(2,1,1)(3,1)(4,2,1)(5,2,1)(6,2)(5)(4)%
}

\NovelDisplay{M03255R01}{%
= \psi(I_2+\psi_{\Omega_{I+1}}(I_2+\psi_{\Omega_{I+1}}(I_2+\psi_{\Omega_{I+1}}(I_2+1))))%
}

\NovelDisplay{M03256R01}{%
= \psi(M^2\times 2+\psi_{\psi_{M}(M^2+M)}(M^2 \times 2 + \psi_{\psi_{M}(M^2+M)}(M^2\times 2%
}
\NovelDisplay{M03256R02}{%
+\psi_{\psi_{M}(M^2+M)}(M^2\times 2+1))))%
}

\NovelDisplay{M03257R01}{%
(0)(1,1,1)(2,1,1)(3,1)(4,2,1)(5,2,1)(6,2)(5)(4,1)(5,2,1)(6,2,1)(7,2)(6)%
}

\NovelDisplay{M03258R01}{%
= \psi(I_2+\psi_{\Omega_{I+1}}(I_2+\psi_{\Omega_{I+1}}(I_2+\psi_{\Omega_{I+1}}(I_2+\psi_{\Omega_{I+1}}(I_2)))))%
}

\NovelDisplay{M03259R01}{%
= \psi(M^2\times 2+\psi_{\psi_{M}(M^2+M)}(M^2 \times 2 + \psi_{\psi_{M}(M^2+M)}(M^2\times 2%
}
\NovelDisplay{M03259R02}{%
+\psi_{\psi_{M}(M^2+M)}(M^2\times 2+\psi_{\psi_{M}(M^2+M)}(M^2\times 2)))))%
}

\NovelDisplay{M03260R01}{%
(0)(1,1,1)(2,1,1)(3,1)(4,2,1)(5,2,1)(6,2)(5)(4,2)%
}

\NovelDisplay{M03261R01}{%
= \psi(I_2+\Omega_{I+1})%
}

\NovelDisplay{M03262R01}{%
= \psi(M^2\times 2+\psi_{M}(M^2+M))%
}

\NovelDisplay{M03263R01}{%
(0)(1,1,1)(2,1,1)(3,1)(4,2,1)(5,2,1)(6,2)(5)(4,2,1)%
}

\NovelDisplay{M03264R01}{%
= \psi(I_2+\Omega_{I+\omega})%
}

\NovelDisplay{M03265R01}{%
= \psi(M^2\times 2+\psi_{M}(M^2+M\times \omega))%
}

内层回填继续加深后，这组又将尾部收束到Ω下标对象，前面的\allowbreak{}\NovelInline{I06046}{I_2}\allowbreak{}背景仍在原处。

\NovelDisplay{M03266R01}{%
(0)(1,1,1)(2,1,1)(3,1)(4,2,1)(5,2,1)(6,2)(5)(4,2,1)(5,2,1)(6,2)(5)%
}

\NovelDisplay{M03267R01}{%
= \psi(I_2+\psi_{I_2}(I_2))%
}

\NovelDisplay{M03268R01}{%
= \psi(M^2\times 2+\psi_{\psi_{M}(M^2\times 2)}(M^2\times 2))%
}

\NovelDisplay{M03269R01}{%
(0)(1,1,1)(2,1,1)(3,1)(4,2,1)(5,2,1)(6,2)(5,2)(4,2,1)(5,2,1)(6,2)(5)%
}

\NovelDisplay{M03270R01}{%
= \psi(I_2+\psi_{\Omega_{\psi_{I_2}(I_2)+1}}(I_2))%
}

\NovelDisplay{M03271R01}{%
= \psi(M^2\times 2+\psi_{\psi_{\psi_{M}(M^2\times 2)}(M^2\times 2+1)}(M^2\times 2))%
}

\NovelDisplay{M03272R01}{%
(0)(1,1,1)(2,1,1)(3,1)(4,2,1)(5,2,1)(6,2)(5,2)(6,3)%
}

\NovelDisplay{M03273R01}{%
= \psi(I_2+\Omega_{\psi_{I_2}(I_2)+1})%
}

\NovelDisplay{M03274R01}{%
= \psi(M^2\times 2+\psi_{\psi_{M}(M^2\times 2)}(M^2\times 2 + 1))%
}

\NovelDisplay{M03275R01}{%
(0)(1,1,1)(2,1,1)(3,1)(4,2,1)(5,2,1)(6,2)(5,2)(6,3,1)%
}

\NovelDisplay{M03276R01}{%
= \psi(I_2+\Omega_{\psi_{I_2}(I_2)+\omega})%
}

\NovelDisplay{M03277R01}{%
= \psi(M^2\times 2+\psi_{\psi_{M}(M^2\times 2)}(M^2\times 2 + \omega))%
}

霍雨浩提醒他们，函数下标开始包含先前求出的ψ值时，要先读清整个下标，再判断括号内喂入的参数。

\NovelDisplay{M03278R01}{%
(0)(1,1,1)(2,1,1)(3,1)(4,2,1)(5,2,1)(6,2)(5,2)(6,3,1)(7,3,1)(8,3)(7)%
}

\NovelDisplay{M03279R01}{%
= \psi(I_2\times 2)%
}

\NovelDisplay{M03280R01}{%
= \psi(M^2\times 2+\psi_{M}(M^2\times 2))%
}

\NovelDisplay{M03281R01}{%
(0)(1,1,1)(2,1,1)(3,1)(4,2,1)(5,2,1)(6,2)(5,2,1)%
}

\NovelDisplay{M03282R01}{%
= \psi(I_2\times \omega)%
}

\NovelDisplay{M03283R01}{%
= \psi(M^2\times 2+\psi_{\psi_{M}(M^2\times 2+M)}(1))%
}

\NovelDisplay{M03284R01}{%
(0)(1,1,1)(2,1,1)(3,1)(4,2,1)(5,2,1)(6,2)(5,2,1)(6,1)(2)%
}

\NovelDisplay{M03285R01}{%
= \psi(I_2\times I)%
}

\NovelDisplay{M03286R01}{%
= \psi(M^2\times 2+\psi_{\psi_{M}(M^2\times 2+M)}(M^2))%
}

\NovelDisplay{M03287R01}{%
(0)(1,1,1)(2,1,1)(3,1)(4,2,1)(5,2,1)(6,2)(5,2,1)(6,2)%
}

\NovelDisplay{M03288R01}{%
= \psi(I_2\times \Omega_{I+1})%
}

\NovelDisplay{M03289R01}{%
= \psi(M^2\times 2+\psi_{\psi_{M}(M^2\times 2+M)}(M^2+M))%
}

这组从\allowbreak{}\NovelInline{I06047}{I_2}\allowbreak{}的重复走向不同右乘因子，M侧须同时保留背景与为乘数设置的新函数头。

\NovelDisplay{M03290R01}{%
(0)(1,1,1)(2,1,1)(3,1)(4,2,1)(5,2,1)(6,2)(5,2,1)(6,2)(4,2,1)%
}

\NovelDisplay{M03291R01}{%
= \psi(I_2\times \Omega_{I+\omega})%
}

\NovelDisplay{M03292R01}{%
= \psi(M^2\times 2+\psi_{\psi_{M}(M^2\times 2+M)}(M^2+M\times \omega))%
}

\NovelDisplay{M03293R01}{%
(0)(1,1,1)(2,1,1)(3,1)(4,2,1)(5,2,1)(6,2)(5,2,1)(6,2)(4,2,1)(5,2,1)%
}
\NovelDisplay{M03293R02}{%
(6,2)(5)%
}

\NovelDisplay{M03294R01}{%
= \psi(I_2\times \psi_{I_2}(I_2))%
}

\NovelDisplay{M03295R01}{%
= \psi(M^2\times 2+\psi_{\psi_{M}(M^2\times 2+M)}(M^2+M\times \psi_{\psi_{M}(M^2\times 2)}(M^2\times 2)))%
}

\NovelDisplay{M03296R01}{%
(0)(1,1,1)(2,1,1)(3,1)(4,2,1)(5,2,1)(6,2)(5,2,1)(6,2)(4,2,1)(5,2,1)%
}
\NovelDisplay{M03296R02}{%
(6,2)(5,2,1)%
}

\NovelDisplay{M03297R01}{%
= \psi(I_2\times \psi_{I_2}(I_2\times \omega))%
}

\NovelDisplay{M03298R01}{%
= \psi(M^2\times 2+\psi_{\psi_{M}(M^2\times 2+M)}(M^2+M\times \psi_{\psi_{M}(M^2\times 2)}(M^2\times 2%
}
\NovelDisplay{M03298R02}{%
+\psi_{\psi_{M}(M^2\times 2+M)}(1))))%
}

\NovelDisplay{M03299R01}{%
(0)(1,1,1)(2,1,1)(3,1)(4,2,1)(5,2,1)(6,2)(5,2,1)(6,2)(4,2,1)(5,2,1)%
}
\NovelDisplay{M03299R02}{%
(6,2)(5,2,1)(6,2)%
}

\NovelDisplay{M03300R01}{%
= \psi(I_2\times \psi_{I_2}(I_2\times \Omega_{I+1}))%
}

\NovelDisplay{M03301R01}{%
= \psi(M^2\times 2+\psi_{\psi_{M}(M^2\times 2+M)}(M^2+M\times \psi_{\psi_{M}(M^2\times 2)}(M^2\times 2%
}
\NovelDisplay{M03301R02}{%
+\psi_{\psi_{M}(M^2\times 2+M)}(M^2+M))))%
}

蔡须坤归纳道：``这里增长的是右乘部分，乘数里的ψ参数怎样变化，要与外层固定的\allowbreak{}\NovelInline{I06048}{I_2}\allowbreak{}分开说明。''

卢本韦这次没有卡在很长的式子上。

她读到一个平方，忽然把它与先前那个下标二混在了一起。蔡须坤刚听出不对，卢本韦自己便停了，盯着纸看了一会儿，耳根慢慢发红。

``我还以为这个最不会错。''

唐三将两条共同前缀各抄了一次，叫她分别补完右侧。卢本韦写完，自己圈出改变的位置，长出一口气。

``往后再看见短式，也得认清它在哪里短。''

她把那张纸折进I册，没有因为已经走在前面，就将这次错处撕掉。

\NovelDisplay{M03302R01}{%
(0)(1,1,1)(2,1,1)(3,1)(4,2,1)(5,2,1)(6,2)(5,2,1)(6,2)(5)%
}

\NovelDisplay{M03303R01}{%
= \psi(I_2^2)%
}

\NovelDisplay{M03304R01}{%
= \psi(M^2\times 2+\psi_{\psi_{M}(M^2\times 2+M)}(M^2\times 2))%
}

\NovelDisplay{M03305R01}{%
(0)(1,1,1)(2,1,1)(3,1)(4,2,1)(5,2,1)(6,2)(5,2,1)(6,2)(5,2,1)(6,2)(5)%
}

\NovelDisplay{M03306R01}{%
= \psi(I_2^3)%
}

\NovelDisplay{M03307R01}{%
= \psi(M^2\times 2+\psi_{\psi_{M}(M^2\times 2+M)}(M^2\times 2 + \psi_{M}(M^2\times 2)))%
}

\NovelDisplay{M03308R01}{%
(0)(1,1,1)(2,1,1)(3,1)(4,2,1)(5,2,1)(6,2)(6,2)(5)%
}

\NovelDisplay{M03309R01}{%
= \psi(I_2^{I_2})%
}

\NovelDisplay{M03310R01}{%
= \psi(M^2\times 2+\psi_{\psi_{M}(M^2\times 2+M)}(M^2\times 2 + \psi_{\psi_{M}(M^2\times 2+M)}(M^2\times 2)))%
}

\NovelDisplay{M03311R01}{%
(0)(1,1,1)(2,1,1)(3,1)(4,2,1)(5,2,1)(6,2)(7,2)(5)%
}

\NovelDisplay{M03312R01}{%
= \psi(I_2^{I_2^{I_2}})%
}

\NovelDisplay{M03313R01}{%
= \psi(M^2\times 2+\psi_{\psi_{M}(M^2\times 2+M)}(M^2\times 2 + \psi_{\psi_{M}(M^2\times 2+M)}(M^2\times 2%
}
\NovelDisplay{M03313R02}{%
+ \psi_{\psi_{M}(M^2\times 2+M)}(M^2\times 2))))%
}

真正的\allowbreak{}\NovelInline{I06049}{I_2}\allowbreak{}平方与后续幂塔在这组出现，幂的结合次序与下标二的记号都必须分别保留。

\NovelDisplay{M03314R01}{%
(0)(1,1,1)(2,1,1)(3,1)(4,2,1)(5,2,1)(6,2)(7,3)%
}

\NovelDisplay{M03315R01}{%
= \psi(\Omega_{I_2+1})%
}

\NovelDisplay{M03316R01}{%
= \psi(M^2\times 2+M)%
}

\NovelDisplay{M03317R01}{%
(0)(1,1,1)(2,1,1)(3,1)(4,2,1)(5,2,1)(6,2)(7,3,1)%
}

\NovelDisplay{M03318R01}{%
= \psi(\Omega_{I_2+\omega})%
}

\NovelDisplay{M03319R01}{%
= \psi(M^2\times 2+M\times \omega)%
}

\NovelDisplay{M03320R01}{%
(0)(1,1,1)(2,1,1)(3,1)(4,2,1)(5,2,1)(6,2)(7,3,1)(8,3,1)(9,3)%
}

\NovelDisplay{M03321R01}{%
= \psi(\Omega_{\Omega_{I_2+1}})%
}

\NovelDisplay{M03322R01}{%
= \psi(M^2\times 2+M\times \psi_{M}(M^2\times 2 + M))%
}

\NovelDisplay{M03323R01}{%
(0)(1,1,1)(2,1,1)(3,1)(4,2,1)(5,2,1)(6,2)(7,3,1)(8,3,1)(9,3)(8)%
}

\NovelDisplay{M03324R01}{%
= \psi(I_3)%
}

\NovelDisplay{M03325R01}{%
= \psi(M^2\times 3)%
}

从幂塔之后转到更高Ω下标与\allowbreak{}\NovelInline{I06050}{I_3}\allowbreak{}时，M式更换了整个背景，不能仅在旧式末尾补一个数字。

\NovelDisplay{M03326R01}{%
(0)(1,1,1)(2,1,1)(3,1)(4,2,1)(5,2,1)(6,2)(7,3,1)(8,3,1)(9,3)(8,3)(9,4,1)%
}
\NovelDisplay{M03326R02}{%
(10,4,1)(11,4)(10)%
}

\NovelDisplay{M03327R01}{%
= \psi(I_3\times 2)%
}

\NovelDisplay{M03328R01}{%
= \psi(M^2\times 3 + \psi_{M}(M^2\times 3))%
}

\NovelDisplay{M03329R01}{%
(0)(1,1,1)(2,1,1)(3,1)(4,2,1)(5,2,1)(6,2)(7,3,1)(8,3,1)(9,3)(8,3,1)%
}

\NovelDisplay{M03330R01}{%
= \psi(I_3\times \omega)%
}

\NovelDisplay{M03331R01}{%
= \psi(M^2\times 3 + \psi_{\psi_{M}(M^2\times 3+M)}(1))%
}

\NovelDisplay{M03332R01}{%
(0)(1,1,1)(2,1,1)(3,1)(4,2,1)(5,2,1)(6,2)(7,3,1)(8,3,1)(9,3)(10,4)%
}

\NovelDisplay{M03333R01}{%
= \psi(\Omega_{I_3+1})%
}

\NovelDisplay{M03334R01}{%
= \psi(M^2\times 3 + M)%
}

\NovelDisplay{M03335R01}{%
(0)(1,1,1)(2,1,1)(3,1)(4,2,1)(5,2,1)(6,2)(7,3,1)(8,3,1)(9,3)(10,4,1)%
}
\NovelDisplay{M03335R02}{%
(11,4,1)(12,4)(11)%
}

\NovelDisplay{M03336R01}{%
= \psi(I_4)%
}

\NovelDisplay{M03337R01}{%
= \psi(M^2\times 4)%
}

这一组将\allowbreak{}\NovelInline{I06051}{I_3}\allowbreak{}的重复与后继Ω下标分别展开，再接到\allowbreak{}\NovelInline{I06052}{I_4}\allowbreak{}，下标推进仍需回看对应的完整矩阵。

新的左列只剩下几列。

小舞看见时松了口气，随即又收住笑，伸手先问薛峰要原式。他们从有限的下标一路走来，纸越铺越宽，眼前忽然短下来，恰好更需要弄清它收束了什么。

蔡须坤没有替她说完。小舞将前一组几个有限节点找回，沿着自己的记录看了很久，才抬起头。

``这里要重新按原式展开，不能拿前面最后一个有限数顶上去。''

霍雨浩点头，将炭条分给她和唐三。两人各拿一半，坐到了同一级。

\NovelDisplay{M03338R01}{%
(0)(1,1,1)(2,1,1)(3,1,1)%
}

\NovelDisplay{M03339R01}{%
= \psi(I_\omega)%
}

\NovelDisplay{M03340R01}{%
= \psi(M^2\times \omega)%
}

\NovelDisplay{M03341R01}{%
(0)(1,1,1)(2,1,1)(3,1,1)(1,1,1)%
}

\NovelDisplay{M03342R01}{%
= \psi(I_\omega+\Omega_{\omega})%
}

\NovelDisplay{M03343R01}{%
= \psi(M^2\times \omega+\psi_M(M\times \omega))%
}

\NovelDisplay{M03344R01}{%
(0)(1,1,1)(2,1,1)(3,1,1)(1,1,1)(2,1,1)(3,1)(2)%
}

\NovelDisplay{M03345R01}{%
= \psi(I_\omega+\psi_{I}(I))%
}

\NovelDisplay{M03346R01}{%
= \psi(M^2\times \omega+\psi_{\psi_{M}(M^2)}(M^2))%
}

\NovelDisplay{M03347R01}{%
(0)(1,1,1)(2,1,1)(3,1,1)(1,1,1)(2,1,1)(3,1)(2,1)(3,2)%
}

\NovelDisplay{M03348R01}{%
= \psi(I_\omega+\psi_{\Omega_{\psi_{I}(I)+1}}(I+\Omega_{\psi_{I}(I)+1}))%
}

\NovelDisplay{M03349R01}{%
= \psi(M^2\times \omega+\psi_{\psi_{\psi_{M}(M^2)}(M^2+1)}(M^2+\psi_{\psi_{M}(M^2)}(M^2+1)))%
}

有限下标到\allowbreak{}\NovelInline{I06053}{I_\omega}\allowbreak{}的转折改变了整体背景，后面较低层的Ω项和ψ项仍可以在新背景之后分别辨认。

薛峰让小舞将刚才那条很短的原矩阵独立展开。她先写第零项，写完以后停了停，确认第一行的零只在最初保留一次，再从原式重新求第一项。

唐三在另一张纸求第二项，没有接着展开她刚写的第一项。蔡须坤看两份结果，古月娜则按当前末列的最高非零行核对生成位置，四人各自重算过，才把三张纸放在一起。

\NovelDisplay{M03350R01}{%
\text{展开第零项：}\quad(0)(1,1,1)(2,1,1)%
}

\NovelDisplay{M03351R01}{%
\text{展开第一项：}\quad(0)(1,1,1)(2,1,1)(3,1)(4,2,1)(5,2,1)%
}

\NovelDisplay{M03352R01}{%
\text{展开第二项：}\quad(0)(1,1,1)(2,1,1)(3,1)(4,2,1)(5,2,1)(6,2)(7,3,1)(8,3,1)%
}

这三个项都来自同一条\allowbreak{}\NovelInline{I06054}{I_\omega}\allowbreak{}原矩阵，新增段随编号增长而改变，但任一有限项都仍在原式之下。

小舞按着自己的第一项，对上唐三的第二项，忽然笑了一声。

``原来这里，我现在真能自己算。''

她笑得很轻，很快又低头去看生成段。蔡须坤没有接一句更大的鼓励，只将笔递回去，让她把自己认清的规则也写下来。

\NovelDisplay{M03353R01}{%
(0)(1,1,1)(2,1,1)(3,1,1)(1,1,1)(2,1,1)(3,1)(2,1)(3,2,1)(4,2,1)(5,2)(4)%
}

\NovelDisplay{M03354R01}{%
= \psi(I_\omega+\psi_{\Omega_{\psi_{I}(I)+1}}(I\times 2))%
}

\NovelDisplay{M03355R01}{%
= \psi(M^2\times \omega+\psi_{\psi_{\psi_{M}(M^2)}(M^2+1)}(M^2+\psi_{M}(M^2)))%
}

\NovelDisplay{M03356R01}{%
(0)(1,1,1)(2,1,1)(3,1,1)(1,1,1)(2,1,1)(3,1)(2,1)(3,2,1)(4,2,1)(5,2)(4,2)%
}
\NovelDisplay{M03356R02}{%
(5,3,1)(6,3,1)(7,3)(6)%
}

\NovelDisplay{M03357R01}{%
= \psi(I_\omega+\psi_{\Omega_{\psi_{I}(I)+1}}(I\times 3))%
}

\NovelDisplay{M03358R01}{%
= \psi(M^2\times \omega+\psi_{\psi_{\psi_{M}(M^2)}(M^2+1)}(M^2+\psi_{M}(M^2)\times 2))%
}

\NovelDisplay{M03359R01}{%
(0)(1,1,1)(2,1,1)(3,1,1)(1,1,1)(2,1,1)(3,1)(2,1)(3,2,1)(4,2,1)%
}
\NovelDisplay{M03359R02}{%
(5,2)(4,2,1)%
}

\NovelDisplay{M03360R01}{%
= \psi(I_\omega+\psi_{\Omega_{\psi_{I}(I)+1}}(I\times \omega))%
}

\NovelDisplay{M03361R01}{%
= \psi(M^2\times \omega+\psi_{\psi_{\psi_{M}(M^2)}(M^2+1)}(M^2+\psi_{\psi_{M}(M^2+M)}(1)))%
}

\NovelDisplay{M03362R01}{%
(0)(1,1,1)(2,1,1)(3,1,1)(1,1,1)(2,1,1)(3,1)(2,1)(3,2,1)(4,2,1)(5,2)(6,3)%
}

\NovelDisplay{M03363R01}{%
= \psi(I_\omega+\psi_{\Omega_{\psi_{I}(I)+1}}(\Omega_{I+1}))%
}

\NovelDisplay{M03364R01}{%
= \psi(M^2\times \omega+\psi_{\psi_{\psi_{M}(M^2)}(M^2+1)}(M^2+M))%
}

卢本韦指出，当前函数头保持同一层次，变化在它吃进的I倍数和Ω下标对象中发生。

\NovelDisplay{M03365R01}{%
(0)(1,1,1)(2,1,1)(3,1,1)(1,1,1)(2,1,1)(3,1)(2,1)(3,2,1)(4,2,1)(5,2)%
}
\NovelDisplay{M03365R02}{%
(6,3,1)(7,3,1)(8,3)(7)%
}

\NovelDisplay{M03366R01}{%
= \psi(I_\omega+\psi_{\Omega_{\psi_{I}(I)+1}}(I_2))%
}

\NovelDisplay{M03367R01}{%
= \psi(M^2\times \omega+\psi_{\psi_{\psi_{M}(M^2)}(M^2+1)}(M^2\times 2))%
}

\NovelDisplay{M03368R01}{%
(0)(1,1,1)(2,1,1)(3,1,1)(1,1,1)(2,1,1)(3,1)(2,1)(3,2,1)(4,2,1)(5,2)%
}
\NovelDisplay{M03368R02}{%
(6,3,1)(7,3,1)(8,3)(9,4,1)(10,4,1)(11,4)(10)%
}

\NovelDisplay{M03369R01}{%
= \psi(I_\omega+\psi_{\Omega_{\psi_{I}(I)+1}}(I_3))%
}

\NovelDisplay{M03370R01}{%
= \psi(M^2\times \omega+\psi_{\psi_{\psi_{M}(M^2)}(M^2+1)}(M^2\times 3))%
}

\NovelDisplay{M03371R01}{%
(0)(1,1,1)(2,1,1)(3,1,1)(1,1,1)(2,1,1)(3,1)(2,1)(3,2,1)(4,2,1)(5,2,1)%
}

\NovelDisplay{M03372R01}{%
= \psi(I_\omega+\psi_{\Omega_{\psi_{I}(I)+1}}(I_\omega))%
}

\NovelDisplay{M03373R01}{%
= \psi(M^2\times \omega+\psi_{\psi_{\psi_{M}(M^2)}(M^2+1)}(M^2\times \omega))%
}

\NovelDisplay{M03374R01}{%
(0)(1,1,1)(2,1,1)(3,1,1)(1,1,1)(2,1,1)(3,1)(2,1)(3,2,1)(4,2,1)%
}
\NovelDisplay{M03374R02}{%
(5,2,1)(3,2)%
}

\NovelDisplay{M03375R01}{%
= \psi(I_\omega+\Omega_{\psi_{I}(I)+1})%
}

\NovelDisplay{M03376R01}{%
= \psi(M^2\times \omega+\psi_{\psi_{M}(M^2)}(M^2+1))%
}

这一组把内参数继续放大直到\allowbreak{}\NovelInline{I06055}{I_\omega}\allowbreak{}，随后出现相应Ω下标值，不能把这一收束当作任意ψ都可直接去掉。

三百三十五米的风很冷。

唐三替蔡须坤看着她已读完的背景，她将注意力留给里面最难的一层。念到一半，她忽然怀疑自己漏了一项，回头时，唐三已经把原来那段指给她看。

``这里还在。你刚才说过了。''

她没有立刻相信，又从头核了一遍，发现他指得正好。

唐三对最深的M下标仍讲不清楚，可自己已经能够负责的部分，他守得很稳。蔡须坤便将纸向他那边挪了些，两人的手指分别停在内外两层，终于把这一段接到了末尾。

\NovelDisplay{M03377R01}{%
(0)(1,1,1)(2,1,1)(3,1,1)(1,1,1)(2,1,1)(3,1)(2,1)(3,2,1)(4,2,1)(5,2,1)%
}
\NovelDisplay{M03377R02}{%
(3,2)(4,3,1)(5,3,1)(6,3,1)(4,3)%
}

\NovelDisplay{M03378R01}{%
= \psi(I_\omega+\Omega_{\psi_{I}(I)+2})%
}

\NovelDisplay{M03379R01}{%
= \psi(M^2\times \omega+\psi_{\psi_{M}(M^2)}(M^2+2))%
}

\NovelDisplay{M03380R01}{%
(0)(1,1,1)(2,1,1)(3,1,1)(1,1,1)(2,1,1)(3,1)(2,1)(3,2,1)(4,2,1)%
}
\NovelDisplay{M03380R02}{%
(5,2,1)(3,2,1)%
}

\NovelDisplay{M03381R01}{%
= \psi(I_\omega+\Omega_{\psi_{I}(I)+\omega})%
}

\NovelDisplay{M03382R01}{%
= \psi(M^2\times \omega+\psi_{\psi_{M}(M^2)}(M^2+\omega))%
}

\NovelDisplay{M03383R01}{%
(0)(1,1,1)(2,1,1)(3,1,1)(1,1,1)(2,1,1)(3,1)(2,1)(3,2,1)(4,2,1)(5,2,1)%
}
\NovelDisplay{M03383R02}{%
(3,2,1)(4,2,1)(5,2)(4)%
}

\NovelDisplay{M03384R01}{%
= \psi(I_\omega+\psi_{I}(I\times 2))%
}

\NovelDisplay{M03385R01}{%
= \psi(M^2\times \omega+\psi_{\psi_{M}(M^2)}(M^2+\psi_{M}(M^2)))%
}

\NovelDisplay{M03386R01}{%
(0)(1,1,1)(2,1,1)(3,1,1)(1,1,1)(2,1,1)(3,1)(2,1)(3,2,1)(4,2,1)(5,2,1)%
}
\NovelDisplay{M03386R02}{%
(3,2,1)(4,2,1)(5,2)(6,3)%
}

\NovelDisplay{M03387R01}{%
= \psi(I_\omega+\psi_{\Omega_{\psi_{I}(I\times 2)+1}}(I\times 2+\Omega_{\psi_{I}(I\times 2)+1}))%
}

\NovelDisplay{M03388R01}{%
= \psi(M^2\times \omega+\psi_{\psi_{\psi_{M}(M^2)}(M^2+\psi_{M}(M^2)+1)}(M^2+\psi_{M}(M^2)+\psi_{\psi_{M}(M^2)}(M^2%
}
\NovelDisplay{M03388R02}{%
+\psi_{M}(M^2)+1)))%
}

Ω下标中的有限后继与ω先改变相应的M参数，接着又由更深函数头描述另一段回填结构。

霍雨浩这一次没有接上蔡须坤的追问。

他低头核过原式，拿出前面的记录，仍没立即开口。卢本韦等了一会儿，才发现他刚才点头，只是在听问题，并没有已经想好答案。

``你也没弄清？''她问。

``这处还没有。''

霍雨浩说完，将自己的草稿放到中间。上面有两种接法，都只写到一半，后面各留着一个小小的问号。

蔡须坤原本攥紧的手松了一些。她将自己说得清的前段指给他看，请他先找这部分有没有问题，再一起辨后面的分歧。

薛峰接过纸，把两种尝试共有的左列重新核出，古月娜便从那一列往回问。霍雨浩答到自己不确定的地方，停下来听她说完，再接着想。

这场问答比平时慢。收纸时，蔡须坤看见他把两个问号都留着，其中一处添了很长的说明，另一处还没有答案。

她将自己的同页记录夹在旁边，没有再把等他开口当成这一段唯一的起点。

霍雨浩看见她写的追问，自己也提起笔想了一阵。两人将灯往中间移了些，等其余人的脚步靠近，才把纸重新展开。

\NovelDisplay{M03389R01}{%
(0)(1,1,1)(2,1,1)(3,1,1)(1,1,1)(2,1,1)(3,1)(2,1)(3,2,1)(4,2,1)(5,2,1)%
}
\NovelDisplay{M03389R02}{%
(3,2,1)(4,2,1)(5,2)(6,3,1)(7,3,1)(8,3)(7)%
}

\NovelDisplay{M03390R01}{%
= \psi(I_\omega+\psi_{\Omega_{\psi_{I}(I\times 2)+1}}(I\times 3))%
}

\NovelDisplay{M03391R01}{%
= \psi(M^2\times \omega+\psi_{\psi_{\psi_{M}(M^2)}(M^2+\psi_{M}(M^2)+1)}(M^2+\psi_{M}(M^2)\times 2))%
}

\NovelDisplay{M03392R01}{%
(0)(1,1,1)(2,1,1)(3,1,1)(1,1,1)(2,1,1)(3,1)(2,1)(3,2,1)(4,2,1)(5,2,1)%
}
\NovelDisplay{M03392R02}{%
(3,2,1)(4,2,1)(5,2)(6,3,1)(7,3,1)(8,3,1)%
}

\NovelDisplay{M03393R01}{%
= \psi(I_\omega+\psi_{\Omega_{\psi_{I}(I\times 2)+1}}(I_\omega))%
}

\NovelDisplay{M03394R01}{%
= \psi(M^2\times \omega+\psi_{\psi_{\psi_{M}(M^2)}(M^2+\psi_{M}(M^2)+1)}(M^2\times \omega))%
}

\NovelDisplay{M03395R01}{%
(0)(1,1,1)(2,1,1)(3,1,1)(1,1,1)(2,1,1)(3,1)(2,1)(3,2,1)(4,2,1)(5,2,1)%
}
\NovelDisplay{M03395R02}{%
(3,2,1)(4,2,1)(5,2)(6,3,1)(7,3,1)(8,3,1)(6,3)%
}

\NovelDisplay{M03396R01}{%
= \psi(I_\omega+\Omega_{\psi_{I}(I\times 2)+1})%
}

\NovelDisplay{M03397R01}{%
= \psi(M^2\times \omega+\psi_{\psi_{M}(M^2)}(M^2+\psi_{M}(M^2)+1))%
}

\NovelDisplay{M03398R01}{%
(0)(1,1,1)(2,1,1)(3,1,1)(1,1,1)(2,1,1)(3,1)(2,1)(3,2,1)(4,2,1)(5,2,1)%
}
\NovelDisplay{M03398R02}{%
(3,2,1)(4,2,1)(5,2)(6,3,1)(7,3,1)(8,3,1)(6,3,1)(7,3,1)(8,3)(7)%
}

\NovelDisplay{M03399R01}{%
= \psi(I_\omega+\psi_{I}(I\times 3))%
}

\NovelDisplay{M03400R01}{%
= \psi(M^2\times \omega+\psi_{\psi_{M}(M^2)}(M^2+\psi_{M}(M^2)\times 2))%
}

新的函数头围绕I的重复背景组织参数，M侧每一份旧项都要留在它原有的有序位置。

\NovelDisplay{M03401R01}{%
(0)(1,1,1)(2,1,1)(3,1,1)(1,1,1)(2,1,1)(3,1)(2,1,1)%
}

\NovelDisplay{M03402R01}{%
= \psi(I_\omega+\psi_{I}(I\times \omega))%
}

\NovelDisplay{M03403R01}{%
= \psi(M^2\times \omega+\psi_{\psi_{M}(M^2)}(M^2+\psi_{\psi_{M}(M^2+M)}(1)))%
}

\NovelDisplay{M03404R01}{%
(0)(1,1,1)(2,1,1)(3,1,1)(1,1,1)(2,1,1)(3,1)(4,2)%
}

\NovelDisplay{M03405R01}{%
= \psi(I_\omega+\psi_{I}(\Omega_{I+1}))%
}

\NovelDisplay{M03406R01}{%
= \psi(M^2\times \omega+\psi_{\psi_{M}(M^2)}(M^2+M))%
}

\NovelDisplay{M03407R01}{%
(0)(1,1,1)(2,1,1)(3,1,1)(1,1,1)(2,1,1)(3,1)(4,2,1)(5,2,1)(6,2)(5)%
}

\NovelDisplay{M03408R01}{%
= \psi(I_\omega+\psi_{I}(I_2))%
}

\NovelDisplay{M03409R01}{%
= \psi(M^2\times \omega+\psi_{\psi_{M}(M^2)}(M^2\times 2))%
}

\NovelDisplay{M03410R01}{%
(0)(1,1,1)(2,1,1)(3,1,1)(1,1,1)(2,1,1)(3,1)(4,2,1)(5,2,1)(6,2)(7,3)%
}

\NovelDisplay{M03411R01}{%
= \psi(I_\omega+\psi_{I}(\Omega_{I_2+1}))%
}

\NovelDisplay{M03412R01}{%
= \psi(M^2\times \omega+\psi_{\psi_{M}(M^2)}(M^2\times 2+M))%
}

这里继续比较I的ω倍数、Ω下标对象及更高I下标作为内参数时的区别，外层\allowbreak{}\NovelInline{I06056}{I_\omega}\allowbreak{}保持完整。

\NovelDisplay{M03413R01}{%
(0)(1,1,1)(2,1,1)(3,1,1)(1,1,1)(2,1,1)(3,1)(4,2,1)(5,2,1)(6,2)(7,3,1)%
}
\NovelDisplay{M03413R02}{%
(8,3,1)(9,3)(8)%
}

\NovelDisplay{M03414R01}{%
= \psi(I_\omega+\psi_{I}(I_3))%
}

\NovelDisplay{M03415R01}{%
= \psi(M^2\times \omega+\psi_{\psi_{M}(M^2)}(M^2\times 3))%
}

\NovelDisplay{M03416R01}{%
(0)(1,1,1)(2,1,1)(3,1,1)(1,1,1)(2,1,1)(3,1)(4,2,1)(5,2,1)(6,2,1)%
}

\NovelDisplay{M03417R01}{%
= \psi(I_\omega+\psi_{I}(I_\omega))%
}

\NovelDisplay{M03418R01}{%
= \psi(M^2\times \omega+\psi_{\psi_{M}(M^2)}(M^2\times \omega))%
}

\NovelDisplay{M03419R01}{%
(0)(1,1,1)(2,1,1)(3,1,1)(1,1,1)(2,1,1)(3,1)(4,2,1)(5,2,1)(6,2,1)(1,1,1)%
}
\NovelDisplay{M03419R02}{%
(2,1,1)(3,1)(4,2,1)(5,2,1)(6,2,1)%
}

\NovelDisplay{M03420R01}{%
= \psi(I_\omega+\psi_{I}(I_\omega)\times 2)%
}

\NovelDisplay{M03421R01}{%
= \psi(M^2\times \omega+\psi_{\psi_{M}(M^2)}(M^2\times \omega)\times 2)%
}

\NovelDisplay{M03422R01}{%
(0)(1,1,1)(2,1,1)(3,1,1)(1,1,1)(2,1,1)(3,1)(4,2,1)(5,2,1)(6,2,1)(2,1)%
}

\NovelDisplay{M03423R01}{%
= \psi(I_\omega+\psi_{\Omega_{\psi_{I}(I_\omega)+1}}(\Omega))%
}

\NovelDisplay{M03424R01}{%
= \psi(M^2\times \omega+\psi_{\psi_{\psi_{M}(M^2)}(M^2\times \omega+1)}(M))%
}

蔡须坤说：``先前的整个\allowbreak{}\NovelInline{I06057}{I_\omega}\allowbreak{}进入内参数以后，函数值重复与更换下标仍然是不同的变化。''

傍晚，小舞主动请卢本韦不要先给提示。

她将最容易混的那两条翻到一起，先独自读，再把同样的问题问给诺当。诺当答到一半忽然犹豫，小舞便将自己也曾误认过的位置指了出来。

两人说了很久。蔡须坤在旁边等着，原本以为需要自己接手，等小舞收笔时，却发现那一小段已经被她们核清。

卢本韦拿起水袋，笑了一下，眼睛里带着一整天没睡够的疲倦。她把水递给小舞，自己又翻开了还没读懂的下一条。

\NovelDisplay{M03425R01}{%
(0)(1,1,1)(2,1,1)(3,1,1)(1,1,1)(2,1,1)(3,1)(4,2,1)(5,2,1)(6,2,1)(2,1)%
}
\NovelDisplay{M03425R02}{%
(1,1,1)(2,1,1)(3,1)(4,2,1)(5,2,1)(6,2,1)%
}

\NovelDisplay{M03426R01}{%
= \psi(I_\omega+\psi_{\Omega_{\psi_{I}(I_\omega)+1}}(I_\omega))%
}

\NovelDisplay{M03427R01}{%
= \psi(M^2\times \omega+\psi_{\psi_{\psi_{M}(M^2)}(M^2\times \omega+1)}(M^2\times \omega))%
}

\NovelDisplay{M03428R01}{%
(0)(1,1,1)(2,1,1)(3,1,1)(1,1,1)(2,1,1)(3,1)(4,2,1)(5,2,1)(6,2,1)%
}
\NovelDisplay{M03428R02}{%
(2,1)(3,2)%
}

\NovelDisplay{M03429R01}{%
= \psi(I_\omega+\Omega_{\psi_{I}(I_\omega)+1})%
}

\NovelDisplay{M03430R01}{%
= \psi(M^2\times \omega+\psi_{\psi_{M}(M^2)}(M^2\times \omega+1))%
}

\NovelDisplay{M03431R01}{%
(0)(1,1,1)(2,1,1)(3,1,1)(1,1,1)(2,1,1)(3,1)(4,2,1)(5,2,1)(6,2,1)(2,1)%
}
\NovelDisplay{M03431R02}{%
(3,2,1)(4,2,1)(5,2)(4)%
}

\NovelDisplay{M03432R01}{%
= \psi(I_\omega+I)%
}

\NovelDisplay{M03433R01}{%
= \psi(M^2\times \omega+\psi_{M}(M^2))%
}

\NovelDisplay{M03434R01}{%
(0)(1,1,1)(2,1,1)(3,1,1)(1,1,1)(2,1,1)(3,1)(4,2,1)(5,2,1)(6,2,1)(2,1)%
}
\NovelDisplay{M03434R02}{%
(3,2,1)(4,2,1)(5,2)(4,2)(5,3,1)(6,3,1)(7,3)(6)%
}

\NovelDisplay{M03435R01}{%
= \psi(I_\omega+I\times 2)%
}

\NovelDisplay{M03436R01}{%
= \psi(M^2\times \omega+\psi_{M}(M^2)\times 2)%
}

这一组从更深函数及其Ω下标收束回I的有限重复，不能因为最后两条短些便省去前面的\allowbreak{}\NovelInline{I06058}{I_\omega}\allowbreak{}背景。

三百四十米的横台比下方宽一点，石缝里长着一簇极短的草。

蔡须坤走上去时，双腿已经不愿再弯。她将本子交给卢本韦，扶着石桩缓了好一会儿，才有余力回头。

薛峰站在后面两级，正在替自己刚才那句说明补一个条件。霍雨浩没有往前拉他，只将问题问得更具体；薛峰答完，脚下那处回牵终于松开，他自己抬起了脚。

唐三和小舞到得晚些。小舞坐下以后，先把上午独立算过的那一页重新翻出，确认页角没有被雨打坏，才仰起头喘气。唐三靠在另一侧，望着刚才走过的石阶，半晌才用手背擦了擦脸。

``这次知道自己卡在哪了。''他说。

卢本韦把水分给他们，手还在抖，拧了两次才拧开袋口。蔡须坤接过去替她拿稳，两人谁也没有笑话谁。

夏夜星空等十一人都缓过来，请他们分别说明刚才负责的一处。薛峰先答，小舞接着问，诺当将自己还没接顺的地方重新摊开，大家便又围着那簇小草坐成半圈。

蔡须坤在圈里听了很久。

她仍有读不下去的M式，古月娜也会在不熟的尾部停住，霍雨浩问过别人之后，自己还要回去重核。可此刻，只要说出具体的问题，身边便有人拿起本子，陪着从共同的那一列找来路。

往上还有三十米，才到这一阶段的目标。

蔡须坤将今天的纸收好，在封套上记下三百四十。古月娜替她拿住一端，等墨迹不再发亮，才与她一起把本子放回包里。

\chapter{第一百二十章：在风里守住彼此}\label{ux7b2cux4e00ux767eux4e8cux5341ux7ae0ux5728ux98ceux91ccux5b88ux4f4fux5f7cux6b64}

三百四十米的横台上，小舞第一个把本子重新打开。

她原想趁大家都在，将刚才不熟的那一处问完。问到一半，诺当也凑过来，指出自己在同一段里犯过的另一个错。两人的问题不同，翻回去找的却是同一条完整矩阵。

蔡须坤在旁边听了很久，直到她们都抬头，才将自己的M册放到中间。

``我这里也有一处。''

她把问题说出来，没有先告诉大家自己已经想过多少遍。卢本韦接住外层，古月娜从共同左列问起，薛峰找到了前面的一页，唐三则将刚才小舞辨清的位置留在一旁，免得他们又在同一个地方绕回去。

风从石桩后面吹来，将纸页一张张翻响。霍雨浩等大家把手按稳，才把这次要继续核的原式放到最上面。

\NovelDisplay{M03437R01}{%
(0)(1,1,1)(2,1,1)(3,1,1)(1,1,1)(2,1,1)(3,1)(4,2,1)(5,2,1)(6,2,1)(2,1)%
}
\NovelDisplay{M03437R02}{%
(3,2,1)(4,2,1)(5,2)(4,2,1)%
}

\NovelDisplay{M03438R01}{%
= \psi(I_\omega+I\times \omega)%
}

\NovelDisplay{M03439R01}{%
= \psi(M^2\times \omega+\psi_{\psi_{M}(M^2+M)}(1))%
}

\NovelDisplay{M03440R01}{%
(0)(1,1,1)(2,1,1)(3,1,1)(1,1,1)(2,1,1)(3,1)(4,2,1)(5,2,1)(6,2,1)(2,1)%
}
\NovelDisplay{M03440R02}{%
(3,2,1)(4,2,1)(5,2)(6,3)%
}

\NovelDisplay{M03441R01}{%
= \psi(I_\omega+\psi_{\Omega_{I+1}}(\Omega_{I+1}))%
}

\NovelDisplay{M03442R01}{%
= \psi(M^2\times \omega+\psi_{\psi_{M}(M^2+M)}(M^2+M))%
}

\NovelDisplay{M03443R01}{%
(0)(1,1,1)(2,1,1)(3,1,1)(1,1,1)(2,1,1)(3,1)(4,2,1)(5,2,1)(6,2,1)(2,1)%
}
\NovelDisplay{M03443R02}{%
(3,2,1)(4,2,1)(5,2,1)%
}

\NovelDisplay{M03444R01}{%
= \psi(I_\omega+\psi_{\Omega_{I+1}}(I_\omega))%
}

\NovelDisplay{M03445R01}{%
= \psi(M^2\times \omega+\psi_{\psi_{M}(M^2+M)}(M^2\times \omega))%
}

\NovelDisplay{M03446R01}{%
(0)(1,1,1)(2,1,1)(3,1,1)(1,1,1)(2,1,1)(3,1)(4,2,1)(5,2,1)(6,2,1)(2,1)%
}
\NovelDisplay{M03446R02}{%
(3,2,1)(4,2,1)(5,2,1)(3,2)%
}

\NovelDisplay{M03447R01}{%
= \psi(I_\omega+\psi_{\Omega_{I+1}}(I_\omega)+\Omega_{\psi_{I}(I_\omega)+1})%
}

\NovelDisplay{M03448R01}{%
= \psi(M^2\times \omega+\psi_{\psi_{M}(M^2+M)}(M^2\times \omega)+\psi_{\psi_{M}(M^2)}(M^2\times \omega+1))%
}

在同一个\allowbreak{}\NovelInline{I06059}{I_\omega}\allowbreak{}背景下，尾部的I倍数、内函数值与另加Ω下标值按各自层次相接。

第二天上山前，夏夜星空让大家交换了昨日的分工。

蔡须坤先听小舞解释共同前段，卢本韦把原先只让薛峰核的左列接过来，霍雨浩则等别人讲完，再回答其中一个具体问题。

小舞起初以为，这样会更慢。她站到蔡须坤平日的位置，面对后面的人，第一句话便忘了自己准备怎么开头。

唐三没有替她说。他将本子朝她转正，问了一个她昨晚刚独立答对的问题。

小舞从那里接起，声音渐渐稳下来。她讲到一处自己仍不确定的位置，抬头看蔡须坤，对方却先将她已经说清的部分复述了一遍，确认大家听到的是同一种意思。

等这一轮结束，山上的雾已经移过一片。多花的时间没有替他们增加高度，小舞却第一次知道，自己能把哪些东西真正交给别人。

蔡须坤背起包，看见唐三还在等她，也让他问了一个问题。她答得不快，答完以后，三人才一同去追前面的脚步。

\NovelDisplay{M03449R01}{%
(0)(1,1,1)(2,1,1)(3,1,1)(1,1,1)(2,1,1)(3,1)(4,2,1)(5,2,1)(6,2,1)(2,1)%
}
\NovelDisplay{M03449R02}{%
(3,2,1)(4,2,1)(5,2,1)(3,2,1)(4,2,1)(5,2)(4)%
}

\NovelDisplay{M03450R01}{%
= \psi(I_\omega+\psi_{\Omega_{I+1}}(I_\omega)+\psi_{I}(I_\omega+I))%
}

\NovelDisplay{M03451R01}{%
= \psi(M^2\times \omega+\psi_{\psi_{M}(M^2+M)}(M^2\times \omega)+\psi_{\psi_{M}(M^2)}(M^2\times \omega+\psi_{M}(M^2)))%
}

\NovelDisplay{M03452R01}{%
(0)(1,1,1)(2,1,1)(3,1,1)(1,1,1)(2,1,1)(3,1)(4,2,1)(5,2,1)(6,2,1)(2,1)%
}
\NovelDisplay{M03452R02}{%
(3,2,1)(4,2,1)(5,2,1)(3,2,1)(4,2,1)(5,2)(6,3)%
}

\NovelDisplay{M03453R01}{%
= \psi(I_\omega+\psi_{\Omega_{I+1}}(I_\omega)+\psi_{I}(I_\omega+\psi_{\Omega_{I+1}}(\Omega_{I+1})))%
}

\NovelDisplay{M03454R01}{%
= \psi(M^2\times \omega+\psi_{\psi_{M}(M^2+M)}(M^2\times \omega)+\psi_{\psi_{M}(M^2)}(M^2\times \omega%
}
\NovelDisplay{M03454R02}{%
+\psi_{\psi_{M}(M^2+M)}(M^2+M)))%
}

\NovelDisplay{M03455R01}{%
(0)(1,1,1)(2,1,1)(3,1,1)(1,1,1)(2,1,1)(3,1)(4,2,1)(5,2,1)(6,2,1)(2,1)%
}
\NovelDisplay{M03455R02}{%
(3,2,1)(4,2,1)(5,2,1)(3,2,1)(4,2,1)(5,2)(6,3,1)(7,3,1)(8,3,1)%
}

\NovelDisplay{M03456R01}{%
= \psi(I_\omega+\psi_{\Omega_{I+1}}(I_\omega)+\psi_{I}(I_\omega+\psi_{\Omega_{I+1}}(I_\omega)))%
}

\NovelDisplay{M03457R01}{%
=\! \psi(M^2\!\times\! \omega\!+\!\psi_{\psi_{M}(M^2+M)}(M^2\!\times\! \omega)\!+\!\psi_{\psi_{M}(M^2)}(M^2\!\times\! \omega\!+\!\psi_{\psi_{M}(M^2+M)}(M^2\!\times\! \omega)))%
}

\NovelDisplay{M03458R01}{%
(0)(1,1,1)(2,1,1)(3,1,1)(1,1,1)(2,1,1)(3,1)(4,2,1)(5,2,1)(6,2,1)(2,1)%
}
\NovelDisplay{M03458R02}{%
(3,2,1)(4,2,1)(5,2,1)(3,2,1)(4,2,1)(5,2)(6,3,1)(7,3,1)(8,3,1)(4,2)(5,3)%
}

\NovelDisplay{M03459R01}{%
= \psi(I_\omega+\psi_{\Omega_{I+1}}(I_\omega)+\Omega_{\psi_{I}(I_\omega+\psi_{\Omega_{I+1}}(I_\omega))+1})%
}

\NovelDisplay{M03460R01}{%
= \psi(M^2\times \omega+\psi_{\psi_{M}(M^2+M)}(M^2\times \omega)+\psi_{\psi_{M}(M^2)}(M^2\times \omega%
}
\NovelDisplay{M03460R02}{%
+\psi_{\psi_{M}(M^2+M)}(M^2\times \omega)+1))%
}

霍雨浩归纳道：``有多个尾项时，每项都从同一外层参数进入，后加的函数不能误装进前一个函数的括号。''

他们在三百四十米附近停了两天。

雨气留在石壁上，靠过去便冷，站久了又累。蔡兴坤将一件干外衣垫在平石上，诺当和叶阳轮流带补给上来，奥斯卡每次分完食物，都留在旁边跟着重新读过一遍。

霍雨浩的声音到午后渐渐哑了。薛峰让他停下来，自己接过前一段说明；说错一处，古月娜抬手示意，他便退到能说明的位置再来。

没有谁能一直讲下去，可每次有人歇下，旁边总还有人愿意将那张纸接住。

\NovelDisplay{M03461R01}{%
(0)(1,1,1)(2,1,1)(3,1,1)(1,1,1)(2,1,1)(3,1)(4,2,1)(5,2,1)(6,2,1)(2,1)%
}
\NovelDisplay{M03461R02}{%
(3,2,1)(4,2,1)(5,2,1)(3,2,1)(4,2,1)(5,2)(6,3,1)(7,3,1)(8,3,1)(4,2)%
}
\NovelDisplay{M03461R03}{%
(5,3,1)(6,3,1)(7,3)(6)%
}

\NovelDisplay{M03462R01}{%
= \psi(I_\omega+\psi_{\Omega_{I+1}}(I_\omega)+I)%
}

\NovelDisplay{M03463R01}{%
= \psi(M^2\times \omega+\psi_{\psi_{M}(M^2+M)}(M^2\times \omega)+\psi_{M}(M^2))%
}

\NovelDisplay{M03464R01}{%
(0)(1,1,1)(2,1,1)(3,1,1)(1,1,1)(2,1,1)(3,1)(4,2,1)(5,2,1)(6,2,1)(2,1,1)%
}

\NovelDisplay{M03465R01}{%
= \psi(I_\omega+\psi_{\Omega_{I+1}}(I_\omega+1))%
}

\NovelDisplay{M03466R01}{%
= \psi(M^2\times \omega+\psi_{\psi_{M}(M^2+M)}(M^2\times \omega+1))%
}

\NovelDisplay{M03467R01}{%
(0)(1,1,1)(2,1,1)(3,1,1)(1,1,1)(2,1,1)(3,1)(4,2,1)(5,2,1)(6,2,1)%
}
\NovelDisplay{M03467R02}{%
(2,1,1)(3,1)(2)%
}

\NovelDisplay{M03468R01}{%
= \psi(I_\omega+\psi_{\Omega_{I+1}}(I_\omega+I))%
}

\NovelDisplay{M03469R01}{%
= \psi(M^2\times \omega+\psi_{\psi_{M}(M^2+M)}(M^2\times \omega+\psi_{M}(M^2)))%
}

\NovelDisplay{M03470R01}{%
(0)(1,1,1)(2,1,1)(3,1,1)(1,1,1)(2,1,1)(3,1)(4,2,1)(5,2,1)(6,2,1)%
}
\NovelDisplay{M03470R02}{%
(2,1,1)(3,1)(4,2)%
}

\NovelDisplay{M03471R01}{%
= \psi(I_\omega+\psi_{\Omega_{I+1}}(I_\omega+\psi_{\Omega_{I+1}}(\Omega_{I+1})))%
}

\NovelDisplay{M03472R01}{%
= \psi(M^2\times \omega+\psi_{\psi_{M}(M^2+M)}(M^2\times \omega+\psi_{\psi_{M}(M^2+M)}(M^2+M)))%
}

这组将I尾项、内参数加一以及更深回填并排，先找到最早分歧的括号，才能说明它们为何不同。

\NovelDisplay{M03473R01}{%
(0)(1,1,1)(2,1,1)(3,1,1)(1,1,1)(2,1,1)(3,1)(4,2,1)(5,2,1)(6,2,1)(2,1,1)%
}
\NovelDisplay{M03473R02}{%
(3,1)(4,2,1)(5,2,1)(6,2,1)%
}

\NovelDisplay{M03474R01}{%
= \psi(I_\omega+\psi_{\Omega_{I+1}}(I_\omega+\psi_{\Omega_{I+1}}(I_\omega)))%
}

\NovelDisplay{M03475R01}{%
= \psi(M^2\times \omega+\psi_{\psi_{M}(M^2+M)}(M^2\times \omega+\psi_{\psi_{M}(M^2+M)}(M^2\times \omega)))%
}

\NovelDisplay{M03476R01}{%
(0)(1,1,1)(2,1,1)(3,1,1)(1,1,1)(2,1,1)(3,1)(4,2,1)(5,2,1)(6,2,1)(2,1,1)%
}
\NovelDisplay{M03476R02}{%
(3,1)(4,2,1)(5,2,1)(6,2,1)(2,1,1)(3,1)(4,2,1)(5,2,1)(6,2,1)%
}

\NovelDisplay{M03477R01}{%
= \psi(I_\omega+\psi_{\Omega_{I+1}}(I_\omega+\psi_{\Omega_{I+1}}(I_\omega)\times 2))%
}

\NovelDisplay{M03478R01}{%
= \psi(M^2\times \omega+\psi_{\psi_{M}(M^2+M)}(M^2\times \omega+\psi_{\psi_{M}(M^2+M)}(M^2\times \omega)\times 2))%
}

\NovelDisplay{M03479R01}{%
(0)(1,1,1)(2,1,1)(3,1,1)(1,1,1)(2,1,1)(3,1)(4,2,1)(5,2,1)(6,2,1)(3)%
}

\NovelDisplay{M03480R01}{%
= \psi(I_\omega+\psi_{\Omega_{I+1}}(I_\omega+\psi_{\Omega_{I+1}}(I_\omega+1)))%
}

\NovelDisplay{M03481R01}{%
= \psi(M^2\times \omega+\psi_{\psi_{M}(M^2+M)}(M^2\times \omega+\psi_{\psi_{M}(M^2+M)}(M^2\times \omega+1)))%
}

\NovelDisplay{M03482R01}{%
(0)(1,1,1)(2,1,1)(3,1,1)(1,1,1)(2,1,1)(3,1)(4,2,1)(5,2,1)(6,2,1)%
}
\NovelDisplay{M03482R02}{%
(3,1)(4,2)%
}

\NovelDisplay{M03483R01}{%
= \psi(I_\omega+\psi_{\Omega_{I+1}}(I_\omega+\psi_{\Omega_{I+1}}(I_\omega+\psi_{\Omega_{I+1}}(\Omega_{I+1}))))%
}

\NovelDisplay{M03484R01}{%
= \psi(M^2\times \omega+\psi_{\psi_{M}(M^2+M)}(M^2\times \omega+\psi_{\psi_{M}(M^2+M)}(M^2\times \omega%
}
\NovelDisplay{M03484R02}{%
+\psi_{\psi_{M}(M^2+M)}(M^2+M))))%
}

内函数的参数继续嵌套时，外面的共同背景和里面每次新增的整段背景都要同时保留。

这一段风尤其急。

霍雨浩原想让大家挤到一处背风的凹口，走近才看见里面仍在滴水，几个人站进去，连本子都难以展开。他试着将纸面上的位置共享给后方的人，自己却要一边看前面，一边照应声音传不到的地方，额角便沁出了汗。

古月娜发现以后，让他先将共享收窄。她站到风来的那一侧，将散乱的气流引开一小片，留出能听清问答的地方，肩背承受的山阶压力依旧在。

唐三把说话的人排得近一些，小舞接住中间那份本子。诺当不再隔着几个人喊，只等上一句传到自己这里，确认听清，才把自己的问题送回去。

蔡须坤听到一句被风截断的话，刚要凭印象补齐，叶阳抬手叫停，又从头念了一遍。

少掉的几个字，恰好决定了它指向哪一层。

她拿笔的手停了片刻，重新记下完整的问题。霍雨浩在旁边喘匀气息，见大家已经能自己传过这一段，才把余下的精神力收了回来。

那天他们只挪了很短一截。临走时，蔡须坤却特意将叶阳叫住，请他把刚才听出缺漏的那一句再问给自己。

\NovelDisplay{M03485R01}{%
(0)(1,1,1)(2,1,1)(3,1,1)(1,1,1)(2,1,1)(3,1)(4,2,1)(5,2,1)(6,2,1)(3,1)%
}
\NovelDisplay{M03485R02}{%
(4,2,1)(5,2,1)(6,2,1)%
}

\NovelDisplay{M03486R01}{%
= \psi(I_\omega+\psi_{\Omega_{I+1}}(I_\omega+\psi_{\Omega_{I+1}}(I_\omega+\psi_{\Omega_{I+1}}(I_\omega))))%
}

\NovelDisplay{M03487R01}{%
= \psi(M^2\times \omega+\psi_{\psi_{M}(M^2+M)}(M^2\times \omega+\psi_{\psi_{M}(M^2+M)}(M^2\times \omega%
}
\NovelDisplay{M03487R02}{%
+\psi_{\psi_{M}(M^2+M)}(M^2\times \omega))))%
}

\NovelDisplay{M03488R01}{%
(0)(1,1,1)(2,1,1)(3,1,1)(1,1,1)(2,1,1)(3,1)(4,2,1)(5,2,1)(6,2,1)(4)%
}

\NovelDisplay{M03489R01}{%
= \psi(I_\omega+\psi_{\Omega_{I+1}}(I_\omega+\psi_{\Omega_{I+1}}(I_\omega+\psi_{\Omega_{I+1}}(I_\omega+1))))%
}

\NovelDisplay{M03490R01}{%
= \psi(M^2\times \omega+\psi_{\psi_{M}(M^2+M)}(M^2\times \omega+\psi_{\psi_{M}(M^2+M)}(M^2\times \omega%
}
\NovelDisplay{M03490R02}{%
+\psi_{\psi_{M}(M^2+M)}(M^2\times \omega + 1))))%
}

\NovelDisplay{M03491R01}{%
(0)(1,1,1)(2,1,1)(3,1,1)(1,1,1)(2,1,1)(3,1)(4,2,1)(5,2,1)(6,2,1)(4,1)%
}
\NovelDisplay{M03491R02}{%
(5,2,1)(6,2,1)(7,2,1)%
}

\NovelDisplay{M03492R01}{%
= \psi(I_\omega+\psi_{\Omega_{I+1}}(I_\omega+\psi_{\Omega_{I+1}}(I_\omega+\psi_{\Omega_{I+1}}(I_\omega+\psi_{\Omega_{I+1}}(I_\omega)))))%
}

\NovelDisplay{M03493R01}{%
= \psi(M^2\times \omega+\psi_{\psi_{M}(M^2+M)}(M^2\times \omega+\psi_{\psi_{M}(M^2+M)}(M^2\times \omega%
}
\NovelDisplay{M03493R02}{%
+\psi_{\psi_{M}(M^2+M)}(M^2\times \omega +\psi_{\psi_{M}(M^2+M)}(M^2\times \omega)))))%
}

\NovelDisplay{M03494R01}{%
(0)(1,1,1)(2,1,1)(3,1,1)(1,1,1)(2,1,1)(3,1)(4,2,1)(5,2,1)(6,2,1)(4,2)%
}

\NovelDisplay{M03495R01}{%
= \psi(I_\omega+\Omega_{I+1})%
}

\NovelDisplay{M03496R01}{%
= \psi(M^2\times \omega+\psi_{M}(M^2+M))%
}

更深的回填最终又出现较短的Ω下标尾项，式子长度的缩短并不取消此前明确的层次关系。

\NovelDisplay{M03497R01}{%
(0)(1,1,1)(2,1,1)(3,1,1)(1,1,1)(2,1,1)(3,1)(4,2,1)(5,2,1)(6,2,1)(4,2)%
}
\NovelDisplay{M03497R02}{%
(5,3,1)(6,3,1)(7,3,1)(5,3)%
}

\NovelDisplay{M03498R01}{%
= \psi(I_\omega+\Omega_{I+2})%
}

\NovelDisplay{M03499R01}{%
= \psi(M^2\times \omega+\psi_{M}(M^2+M\times 2))%
}

\NovelDisplay{M03500R01}{%
(0)(1,1,1)(2,1,1)(3,1,1)(1,1,1)(2,1,1)(3,1)(4,2,1)(5,2,1)(6,2,1)(4,2,1)%
}

\NovelDisplay{M03501R01}{%
= \psi(I_\omega+\Omega_{I+\omega})%
}

\NovelDisplay{M03502R01}{%
= \psi(M^2\times \omega+\psi_{M}(M^2+M\times \omega))%
}

\NovelDisplay{M03503R01}{%
(0)(1,1,1)(2,1,1)(3,1,1)(1,1,1)(2,1,1)(3,1)(4,2,1)(5,2,1)(6,2,1)(4,2,1)%
}
\NovelDisplay{M03503R02}{%
(5,2,1)(6,2)(5)%
}

\NovelDisplay{M03504R01}{%
= \psi(I_\omega+\psi_{I_2}(I_2))%
}

\NovelDisplay{M03505R01}{%
= \psi(M^2\times \omega+\psi_{\psi_{M}(M^2\times 2)}(M^2\times 2))%
}

\NovelDisplay{M03506R01}{%
(0)(1,1,1)(2,1,1)(3,1,1)(1,1,1)(2,1,1)(3,1)(4,2,1)(5,2,1)(6,2,1)(4,2,1)%
}
\NovelDisplay{M03506R02}{%
(5,2,1)(6,2)(7,3)%
}

\NovelDisplay{M03507R01}{%
= \psi(I_\omega+\psi_{I_2}(\Omega_{I_2+1}))%
}

\NovelDisplay{M03508R01}{%
= \psi(M^2\times \omega+\psi_{\psi_{M}(M^2\times 2)}(M^2\times 2+M))%
}

从\allowbreak{}\NovelInline{I06060}{\Omega_{I+2}}\allowbreak{}与\allowbreak{}\NovelInline{I06061}{\Omega_{I+\omega}}\allowbreak{}转到\allowbreak{}\NovelInline{I06062}{I_2}\allowbreak{}这一层的ψ函数时，应重新核对函数头和内参数。

后面的矩阵越来越长。

薛峰将三张纸横着并起，在相接处标好列的次序，请唐三从一端读，小舞从另一端倒着核。两人的手指最终停在同一处，他才将中间那张纸交给蔡须坤。

蔡须坤问他是否已经能解释所有M式，他摇头，随即将自己核出的保留段说了一遍。

``这部分我能确定。后面要再问。''

她听完，将后半段留到自己面前。薛峰没有把不会的藏起来，也没有因为后面还不会，就放开自己已经能负责的这一截。

\NovelDisplay{M03509R01}{%
(0)(1,1,1)(2,1,1)(3,1,1)(1,1,1)(2,1,1)(3,1)(4,2,1)(5,2,1)(6,2,1)(4,2,1)%
}
\NovelDisplay{M03509R02}{%
(5,2,1)(6,2)(7,3,1)(8,3,1)(9,3,1)%
}

\NovelDisplay{M03510R01}{%
= \psi(I_\omega+\psi_{I_2}(I_\omega))%
}

\NovelDisplay{M03511R01}{%
= \psi(M^2\times \omega+\psi_{\psi_{M}(M^2\times 2)}(M^2\times \omega))%
}

\NovelDisplay{M03512R01}{%
(0)(1,1,1)(2,1,1)(3,1,1)(1,1,1)(2,1,1)(3,1)(4,2,1)(5,2,1)(6,2,1)(4,2,1)%
}
\NovelDisplay{M03512R02}{%
(5,2,1)(6,2)(7,3,1)(8,3,1)(9,3,1)(5,2)(6,3)%
}

\NovelDisplay{M03513R01}{%
= \psi(I_\omega+\Omega_{\psi_{I_2}(I_\omega)+1})%
}

\NovelDisplay{M03514R01}{%
= \psi(M^2\times \omega+\psi_{\psi_{M}(M^2\times 2)}(M^2\times \omega+1))%
}

\NovelDisplay{M03515R01}{%
(0)(1,1,1)(2,1,1)(3,1,1)(1,1,1)(2,1,1)(3,1)(4,2,1)(5,2,1)(6,2,1)(4,2,1)%
}
\NovelDisplay{M03515R02}{%
(5,2,1)(6,2)(7,3,1)(8,3,1)(9,3,1)(5,2)(6,3,1)(7,3,1)(8,3)(7)%
}

\NovelDisplay{M03516R01}{%
= \psi(I_\omega+I_2)%
}

\NovelDisplay{M03517R01}{%
= \psi(M^2\times \omega+\psi_{M}(M^2\times 2))%
}

\NovelDisplay{M03518R01}{%
(0)(1,1,1)(2,1,1)(3,1,1)(1,1,1)(2,1,1)(3,1)(4,2,1)(5,2,1)(6,2,1)(4,2,1)%
}
\NovelDisplay{M03518R02}{%
(5,2,1)(6,2)(7,3,1)(8,3,1)(9,3,1)(5,2)(6,3,1)(7,3,1)(8,3,1)%
}

\NovelDisplay{M03519R01}{%
= \psi(I_\omega+\psi_{\Omega_{I_2+1}}(I_\omega))%
}

\NovelDisplay{M03520R01}{%
= \psi(M^2\times \omega+\psi_{\psi_{M}(M^2\times 2 + M)}(M^2\times \omega))%
}

卢本韦归纳道：``这几条依次区分I侧的下标二、Ω后继下标和它们在尾部的位置，不能只追最后一列的外形。''

卢本韦的腿忽然抽了一下。

她正要抬脚，膝弯猛地收紧，整个人便停在原地。蔡须坤最先察觉，伸手让她扶住石壁，自己往下退了半级。

``先别动。''

卢本韦咬着牙，等那阵疼过去，还想把刚才那句话讲完。古月娜已经将她前臂上的本子解下，交给薛峰，叫她先到平处坐稳。

几人退回去时，后面的队伍也跟着停了。卢本韦低着头，连水都没有立刻接，直到唐三在她旁边坐下，伸手把自己绷得发紧的脚踝慢慢活动开。

``我也正想退。''他说。

小舞看了他一眼，知道他确实早就站得难受，只是见前面还在继续，便一直没有开口。她把水塞到两人中间，让他们各自拿稳。

蔡须坤蹲在卢本韦面前，问疼的地方有没有缓下。卢本韦点头，过了一会儿，又低低说了句抱歉。

``我早一点说，就不用让你们都跟着退。''

蔡须坤原想马上劝她别这样想，话到嘴边，却记起自己也不止一次把疲倦藏到下一阶。她只将同一袋水接过，喝了一口，坐到了她身边。

那天下山，他们重新分了背包。最厚的两册拆开由不同的人带，今天先读完的人也不再只站着等，等同伴收稳运转，就先把下一处落脚的空间让出来。

\NovelDisplay{M03521R01}{%
(0)(1,1,1)(2,1,1)(3,1,1)(1,1,1)(2,1,1)(3,1)(4,2,1)(5,2,1)(6,2,1)(4,2,1)%
}
\NovelDisplay{M03521R02}{%
(5,2,1)(6,2)(7,3,1)(8,3,1)(9,3,1)(5,2)(6,3,1)(7,3,1)(8,3,1)(6,3,1)%
}
\NovelDisplay{M03521R03}{%
(7,3,1)(8,3)(7)%
}

\NovelDisplay{M03522R01}{%
= \psi(I_\omega+\psi_{\Omega_{I_2+1}}(I_\omega)+\psi_{I_2}(I_\omega+I_2))%
}

\NovelDisplay{M03523R01}{%
=\! \psi(M^2\!\times\! \omega\!+\!\psi_{\psi_{M}(M^2\times 2 + M)}(M^2\!\times\! \omega)\!+\!\psi_{\psi_{M}(M^2\times 2)}(M^2\!\times\! \omega \!+\! \psi_{M}(M^2\!\times\! 2)))%
}

\NovelDisplay{M03524R01}{%
(0)(1,1,1)(2,1,1)(3,1,1)(1,1,1)(2,1,1)(3,1)(4,2,1)(5,2,1)(6,2,1)(4,2,1)%
}
\NovelDisplay{M03524R02}{%
(5,2,1)(6,2)(7,3,1)(8,3,1)(9,3,1)(5,2)(6,3,1)(7,3,1)(8,3,1)(6,3,1)%
}
\NovelDisplay{M03524R03}{%
(7,3,1)(8,3)(9,4,1)(10,4,1)(11,4,1)(7,3)(8,4,1)(9,4,1)(10,4)(9)%
}

\NovelDisplay{M03525R01}{%
= \psi(I_\omega+\psi_{\Omega_{I_2+1}}(I_\omega)+I_2)%
}

\NovelDisplay{M03526R01}{%
= \psi(M^2\times \omega+\psi_{\psi_{M}(M^2\times 2 + M)}(M^2\times \omega)+\psi_{M}(M^2\times 2))%
}

\NovelDisplay{M03527R01}{%
(0)(1,1,1)(2,1,1)(3,1,1)(1,1,1)(2,1,1)(3,1)(4,2,1)(5,2,1)(6,2,1)(4,2,1)%
}
\NovelDisplay{M03527R02}{%
(5,2,1)(6,2)(7,3,1)(8,3,1)(9,3,1)(5,2)(6,3,1)(7,3,1)(8,3,1)(6,3,1)%
}
\NovelDisplay{M03527R03}{%
(7,3,1)(8,3)(9,4,1)(10,4,1)(11,4,1)(7,3)(8,4,1)(9,4,1)(10,4)(11,5,1)%
}
\NovelDisplay{M03527R04}{%
(12,5,1)(13,5,1)%
}

\NovelDisplay{M03528R01}{%
= \psi(I_\omega+\psi_{\Omega_{I_2+1}}(I_\omega)\times 2)%
}

\NovelDisplay{M03529R01}{%
= \psi(M^2\times \omega+\psi_{\psi_{M}(M^2\times 2 + M)}(M^2\times \omega)\times 2)%
}

\NovelDisplay{M03530R01}{%
(0)(1,1,1)(2,1,1)(3,1,1)(1,1,1)(2,1,1)(3,1)(4,2,1)(5,2,1)(6,2,1)(4,2,1)%
}
\NovelDisplay{M03530R02}{%
(5,2,1)(6,2)(7,3,1)(8,3,1)(9,3,1)(5,2,1)%
}

\NovelDisplay{M03531R01}{%
= \psi(I_\omega+\psi_{\Omega_{I_2+1}}(I_\omega+1))%
}

\NovelDisplay{M03532R01}{%
= \psi(M^2\times \omega+\psi_{\psi_{M}(M^2\times 2 + M)}(M^2\times \omega+1))%
}

最长的一组仍围绕同一个\allowbreak{}\NovelInline{I06063}{I_\omega}\allowbreak{}背景组织尾项，重复整个内函数值与在内参数加一各有不同的完整矩阵。

最长的一条铺完时，三百五十米的刻线已经露在上面。

蔡须坤抬头看过两次，第二次才发现，自己又在用那根石桩催促眼前的推演。她将本子稍稍转向石壁，暂时挡住刻线，只留最后几组式子在眼前。

卢本韦没有劝她快一点。两人各自说完自己的表示，霍雨浩请薛峰从共同矩阵检查，再问唐三和小舞分别能指出哪一处变化。

答案回来得有快有慢，蔡须坤却都听完了，才将脚重新放到下一阶前。

\NovelDisplay{M03533R01}{%
(0)(1,1,1)(2,1,1)(3,1,1)(1,1,1)(2,1,1)(3,1)(4,2,1)(5,2,1)(6,2,1)(4,2,1)%
}
\NovelDisplay{M03533R02}{%
(5,2,1)(6,2)(7,3,1)(8,3,1)(9,3,1)(7,3)%
}

\NovelDisplay{M03534R01}{%
= \psi(I_\omega+\Omega_{I_2+1})%
}

\NovelDisplay{M03535R01}{%
= \psi(M^2\times \omega+\psi_{M}(M^2\times 2 + M))%
}

\NovelDisplay{M03536R01}{%
(0)(1,1,1)(2,1,1)(3,1,1)(1,1,1)(2,1,1)(3,1)(4,2,1)(5,2,1)(6,2,1)(4,2,1)%
}
\NovelDisplay{M03536R02}{%
(5,2,1)(6,2)(7,3,1)(8,3,1)(9,3,1)(7,3,1)(8,3,1)(9,3)(10,4,1)%
}
\NovelDisplay{M03536R03}{%
(11,4,1)(12,4,1)%
}

\NovelDisplay{M03537R01}{%
= \psi(I_\omega+\psi_{I_3}(I_\omega))%
}

\NovelDisplay{M03538R01}{%
= \psi(M^2\times \omega+\psi_{\psi_{M}(M^2\times 3)}(M^2\times \omega))%
}

\NovelDisplay{M03539R01}{%
(0)(1,1,1)(2,1,1)(3,1,1)(1,1,1)(2,1,1)(3,1)(4,2,1)(5,2,1)(6,2,1)(4,2,1)%
}
\NovelDisplay{M03539R02}{%
(5,2,1)(6,2)(7,3,1)(8,3,1)(9,3,1)(7,3,1)(8,3,1)(9,3)(10,4,1)(11,4,1)%
}
\NovelDisplay{M03539R03}{%
(12,4,1)(10,4)%
}

\NovelDisplay{M03540R01}{%
= \psi(I_\omega+\Omega_{I_3+1})%
}

\NovelDisplay{M03541R01}{%
= \psi(M^2\times \omega+\psi_{M}(M^2\times 3+M))%
}

最后三条依次保留\allowbreak{}\NovelInline{I06064}{\Omega_{I_2+1}}\allowbreak{}及\allowbreak{}\NovelInline{I06065}{I_3}\allowbreak{}层的函数和Ω下标尾项，为后面整份\allowbreak{}\NovelInline{I06066}{I_\omega}\allowbreak{}的重复留下清楚背景。

三百五十米以上，石阶向山壁里折了一次。

站在横台上时，还能看见下方几段绳索；转过去以后，风声忽然贴近耳边，脚下可落的位置却只剩窄窄的一条。蔡须坤抬头看了看，上面二十米，比她最初望过的二十米陡得多。

卢本韦将I册用带子固定在前臂，问她先站哪一边。蔡须坤把较宽的位置让给后面的人，自己靠近石壁，将M册抵在膝上。

``还是从昨天的背景接。''

她说得很稳，手指却在纸上停了一会儿。M式后面一长段开始重新包住前面的结构，稍不留意，便会把整份对象的重复，读成外层参数里那个熟悉因子的直接重复。

霍雨浩先请卢本韦读I侧，再让蔡须坤从M侧指出它实际保留的部分。两人说到不同层次时，他便抬手，等她们各自退回一层再接。

古月娜在前面试脚下的落点。薛峰回头向唐三报了一个可以放包的位置，小舞便将后方同伴的纸先收来，免得谁腾不出手时把整页落进雾里。

这一段没有人说只剩二十米。大家将昨天还能独立说清的那部分先留稳，才打开了新的第一页。

\NovelDisplay{M03542R01}{%
(0)(1,1,1)(2,1,1)(3,1,1)(1,1,1)(2,1,1)(3,1,1)%
}

\NovelDisplay{M03543R01}{%
= \psi(I_\omega\times 2)%
}

\NovelDisplay{M03544R01}{%
= \psi(M^2\times \omega+\psi_{M}(M^2\times \omega))%
}

\NovelDisplay{M03545R01}{%
(0)(1,1,1)(2,1,1)(3,1,1)(1,1,1)(2,1,1)(3,1,1)(1,1,1)(2,1,1)(3,1)(4,2,1)%
}
\NovelDisplay{M03545R02}{%
(5,2,1)(6,2,1)(4,2,1)(5,2,1)(6,2,1)%
}

\NovelDisplay{M03546R01}{%
= \psi(I_\omega\times 2+\psi_{I}(I_\omega \times 2))%
}

\NovelDisplay{M03547R01}{%
= \psi(M^2\times \omega+\psi_{M}(M^2\times \omega) + \psi_{\psi_{M}(M^2)}(M^2\times \omega + \psi_{M}(M^2\times \omega)))%
}

\NovelDisplay{M03548R01}{%
(0)(1,1,1)(2,1,1)(3,1,1)(1,1,1)(2,1,1)(3,1,1)(1,1,1)(2,1,1)(3,1)(4,2,1)%
}
\NovelDisplay{M03548R02}{%
(5,2,1)(6,2,1)(4,2,1)(5,2,1)(6,2,1)(4,2)%
}

\NovelDisplay{M03549R01}{%
= \psi(I_\omega\times 2+\Omega_{I+1})%
}

\NovelDisplay{M03550R01}{%
= \psi(M^2\times \omega+\psi_{M}(M^2\times \omega) + \psi_{M}(M^2+M))%
}

\NovelDisplay{M03551R01}{%
(0)(1,1,1)(2,1,1)(3,1,1)(1,1,1)(2,1,1)(3,1,1)(1,1,1)(2,1,1)(3,1)(4,2,1)%
}
\NovelDisplay{M03551R02}{%
(5,2,1)(6,2,1)(4,2,1)(5,2,1)(6,2,1)(4,2,1)(5,2,1)(6,2)(7,3,1)(8,3,1)%
}
\NovelDisplay{M03551R03}{%
(9,3,1)(7,3,1)(8,3,1)(9,3,1)%
}

\NovelDisplay{M03552R01}{%
= \psi(I_\omega\times 2+\psi_{I_2}(I_\omega\times 2))%
}

\NovelDisplay{M03553R01}{%
= \psi(M^2\times \omega+\psi_{M}(M^2\times \omega) + \psi_{\psi_{M}(M^2\times 2)}(M^2\times \omega+\psi_{M}(M^2\times \omega)))%
}

卢本韦归纳道：``整份\allowbreak{}\NovelInline{I06067}{I_\omega}\allowbreak{}重复后，M侧保留原背景并加入它的相应函数值，不能直接把它改写成背景因子的普通倍数。''

霍雨浩将刚才那条原矩阵再取出来，请大家各自求一次有限项。蔡须坤和卢本韦暂时合起右侧记号，唐三找末列的来路，小舞检查保留段，薛峰把不同编号的生成过程分别写在纸上。

古月娜等三份结果都放到面前，才将自己的复算翻过来。没有人拿前一项再展开来充作后一项，每一个编号都重新回到同一原式。

\NovelDisplay{M03554R01}{%
\text{展开第零项：}\quad(0)(1,1,1)(2,1,1)(3,1,1)(1,1,1)(2,1,1)%
}

\NovelDisplay{M03555R01}{%
\text{展开第一项：}\quad(0)(1,1,1)(2,1,1)(3,1,1)(1,1,1)(2,1,1)(3,1)(4,2,1)%
}
\NovelDisplay{M03555R02}{%
(5,2,1)(6,2,1)(4,2,1)(5,2,1)%
}

\NovelDisplay{M03556R01}{%
\text{展开第二项：}\quad(0)(1,1,1)(2,1,1)(3,1,1)(1,1,1)(2,1,1)(3,1)(4,2,1)%
}
\NovelDisplay{M03556R02}{%
(5,2,1)(6,2,1)(4,2,1)(5,2,1)(6,2)(7,3,1)(8,3,1)(9,3,1)(7,3,1)(8,3,1)%
}

这三个有限项各自保留当前原式的既有前段，并按原版规则生成新增段，不能因结果较长便把它当成原式本身。

蔡须坤将少写的一列补回，问薛峰为什么他能先看出。薛峰指了指自己分开的三张纸，让她看每一次实际从哪里开始。

她照着重做，终于在还没有写完之前，自己停在那处原先会漏的位置。

\NovelDisplay{M03557R01}{%
(0)(1,1,1)(2,1,1)(3,1,1)(1,1,1)(2,1,1)(3,1,1)(1,1,1)(2,1,1)(3,1)(4,2,1)%
}
\NovelDisplay{M03557R02}{%
(5,2,1)(6,2,1)(4,2,1)(5,2,1)(6,2,1)(4,2,1)(5,2,1)(6,2)(7,3,1)(8,3,1)%
}
\NovelDisplay{M03557R03}{%
(9,3,1)(7,3,1)(8,3,1)(9,3,1)(7,3)%
}

\NovelDisplay{M03558R01}{%
= \psi(I_\omega\times 2+\Omega_{I_2+1})%
}

\NovelDisplay{M03559R01}{%
= \psi(M^2\times \omega+\psi_{M}(M^2\times \omega) + \psi_{M}(M^2\times 2+M))%
}

\NovelDisplay{M03560R01}{%
(0)(1,1,1)(2,1,1)(3,1,1)(1,1,1)(2,1,1)(3,1,1)(1,1,1)(2,1,1)(3,1,1)%
}

\NovelDisplay{M03561R01}{%
= \psi(I_\omega\times 3)%
}

\NovelDisplay{M03562R01}{%
= \psi(M^2\times \omega+\psi_{M}(M^2\times \omega)\times 2)%
}

\NovelDisplay{M03563R01}{%
(0)(1,1,1)(2,1,1)(3,1,1)(2)%
}

\NovelDisplay{M03564R01}{%
= \psi(I_\omega\times \omega)%
}

\NovelDisplay{M03565R01}{%
= \psi(M^2\times \omega+\psi_{\psi_{M}(M^2\times \omega+M)}(1))%
}

\NovelDisplay{M03566R01}{%
(0)(1,1,1)(2,1,1)(3,1,1)(2,1)%
}

\NovelDisplay{M03567R01}{%
= \psi(I_\omega\times \Omega)%
}

\NovelDisplay{M03568R01}{%
= \psi(M^2\times \omega+\psi_{\psi_{M}(M^2\times \omega+M)}(M))%
}

三份背景、右乘ω与右乘Ω分别由各自矩阵给出，增长进入哪一层参数需要逐条确认。

晚上，蔡须坤把M册摊在灯下，许久没有翻页。

唐三整理好明天要带的纸，见她仍盯着同一处，便问起当初为什么选了更难的一边。

她起初只说想试，后来想了很久，才把手放在那一整段函数上。

``我以前总想先认出它有多大。认出来以后，就觉得这道题差不多结束了。可再换一个接法，里面哪里变了，我还是会怕。''

唐三听着，没有打断。

``从这边慢慢读，有时候确实更难。''她说，``但我能知道自己到底漏在哪一层。真要帮别人，也不能只把最后那个名字告诉他。''

卢本韦在桌子另一头合上I册，抬头问她，那自己选了较短的表示，是不是少走了一段。

蔡须坤赶紧摇头，让她看今天两人互相问住的位置。她没能凭M侧独自越过去，卢本韦也要靠共同矩阵才将I侧讲清。把各自看见的东西放在一起，才比一个人抱着最难的一页有用。

古月娜端着温水进来，听见最后几句，站在桌边停了一阵。

``那两个期望对应，我还留着。''她说。

蔡须坤伸手摸了摸包里的封套。

``我也留着。''

两人没有接着争哪一条应该怎样定义。古月娜将水分好，问她明天从哪里开始，蔡须坤便翻回自己今天真正停住的位置，让她看那一道很细的笔痕。

小舞从门边经过，看见四人又要继续，便把灯座往桌子中央挪了一点，叫她们先把水喝完。

\NovelDisplay{M03569R01}{%
(0)(1,1,1)(2,1,1)(3,1,1)(2,1)(1,1,1)(2,1,1)(3,1)(2)%
}

\NovelDisplay{M03570R01}{%
= \psi(I_\omega\times \psi_{I}(I))%
}

\NovelDisplay{M03571R01}{%
= \psi(M^2\times \omega+\psi_{\psi_{M}(M^2\times \omega+M)}(M\times \psi_{\psi_{M}(M^2)}(M^2)))%
}

\NovelDisplay{M03572R01}{%
(0)(1,1,1)(2,1,1)(3,1,1)(2,1)(1,1,1)(2,1,1)(3,1)(4,2,1)(5,2,1)(6,2,1)%
}

\NovelDisplay{M03573R01}{%
= \psi(I_\omega\times \psi_{I}(I_\omega))%
}

\NovelDisplay{M03574R01}{%
= \psi(M^2\times \omega+\psi_{\psi_{M}(M^2\times \omega+M)}(M\times \psi_{\psi_{M}(M^2)}(M^2\times \omega)))%
}

\NovelDisplay{M03575R01}{%
(0)(1,1,1)(2,1,1)(3,1,1)(2,1)(1,1,1)(2,1,1)(3,1)(4,2,1)(5,2,1)%
}
\NovelDisplay{M03575R02}{%
(6,2,1)(5,1)(2)%
}

\NovelDisplay{M03576R01}{%
= \psi(I_\omega\times I)%
}

\NovelDisplay{M03577R01}{%
= \psi(M^2\times \omega+\psi_{\psi_{M}(M^2\times \omega+M)}(M^2))%
}

\NovelDisplay{M03578R01}{%
(0)(1,1,1)(2,1,1)(3,1,1)(2,1)(1,1,1)(2,1,1)(3,1)(4,2,1)(5,2,1)(6,2,1)%
}
\NovelDisplay{M03578R02}{%
(5,1)(6,2,1)(7,2,1)(8,2,1)%
}

\NovelDisplay{M03579R01}{%
= \psi(I_\omega\times \psi_{\Omega_{I+1}}(I_\omega))%
}

\NovelDisplay{M03580R01}{%
= \psi(M^2\times \omega+\psi_{\psi_{M}(M^2\times \omega+M)}(M^2+\psi_{\psi_{M}(M^2+M)}(M^2\times \omega)))%
}

蔡须坤指出，乘数从较低ψ值走到I时，M侧在同一个新函数头内保留了各自不同的参数结构。

蔡须坤在内参数里看见同样的旧背景时，心口先紧了一下。

她想起刚登山那一阵：本以为接上一段便能保留前面的答案，结果整条祖先链都要重新追。如今更深的式子铺在面前，那种疲惫又沿着肩背慢慢爬上来。

卢本韦问她是否要停，她摇头，又在准备开口时点了点头。

两人在平处坐下。卢本韦只读最外面的乘法，把里面的空给她留着；古月娜拿出前一天核过的同层结构，薛峰找到共同的左列。蔡须坤看着三人的手各停在一处，终于把自己说不清的那层圈了出来。

``先只问这里。''她说。

\NovelDisplay{M03581R01}{%
(0)(1,1,1)(2,1,1)(3,1,1)(2,1)(1,1,1)(2,1,1)(3,1)(4,2,1)(5,2,1)(6,2,1)%
}
\NovelDisplay{M03581R02}{%
(5,1)(6,2,1)(7,2,1)(8,2,1)(7,1)(8,2,1)(9,2,1)(10,2,1)%
}

\NovelDisplay{M03582R01}{%
= \psi(I_\omega\times \psi_{\Omega_{I+1}}(I_\omega\times \psi_{\Omega_{I+1}}(I_\omega)))%
}

\NovelDisplay{M03583R01}{%
= \psi(M^2\times \omega+\psi_{\psi_{M}(M^2\times \omega+M)}(M^2+\psi_{\psi_{M}(M^2+M)}(M^2\times \omega%
}
\NovelDisplay{M03583R02}{%
+\psi_{\psi_{M}(M^2\times \omega+M)}(M^2+\psi_{\psi_{M}(M^2+M)}(M^2\times \omega)))))%
}

\NovelDisplay{M03584R01}{%
(0)(1,1,1)(2,1,1)(3,1,1)(2,1)(1,1,1)(2,1,1)(3,1)(4,2,1)(5,2,1)%
}
\NovelDisplay{M03584R02}{%
(6,2,1)(5,2)%
}

\NovelDisplay{M03585R01}{%
= \psi(I_\omega\times \Omega_{I+1})%
}

\NovelDisplay{M03586R01}{%
= \psi(M^2\times \omega+\psi_{\psi_{M}(M^2\times \omega+M)}(M^2+M))%
}

\NovelDisplay{M03587R01}{%
(0)(1,1,1)(2,1,1)(3,1,1)(2,1)(1,1,1)(2,1,1)(3,1)(4,2,1)(5,2,1)(6,2,1)%
}
\NovelDisplay{M03587R02}{%
(5,2)(4,2,1)(5,2,1)(6,2)(7,3,1)(8,3,1)(9,3,1)%
}

\NovelDisplay{M03588R01}{%
= \psi(I_\omega\times \psi_{I_2}(I_\omega))%
}

\NovelDisplay{M03589R01}{%
= \psi(M^2\times \omega+\psi_{\psi_{M}(M^2\times \omega+M)}(M^2+M\times \psi_{\psi_{M}(M^2\times 2)}(M^2\times \omega)))%
}

\NovelDisplay{M03590R01}{%
(0)(1,1,1)(2,1,1)(3,1,1)(2,1)(1,1,1)(2,1,1)(3,1)(4,2,1)(5,2,1)(6,2,1)%
}
\NovelDisplay{M03590R02}{%
(5,2)(4,2,1)(5,2,1)(6,2)(7,3,1)(8,3,1)(9,3,1)(8,2)(5)%
}

\NovelDisplay{M03591R01}{%
= \psi(I_\omega\times I_2)%
}

\NovelDisplay{M03592R01}{%
= \psi(M^2\times \omega+\psi_{\psi_{M}(M^2\times \omega+M)}(M^2\times 2))%
}

乘数里的嵌套再深，也要从完整外层背景开始读，不能把内层的一次回填误作整个式子的重复。

薛峰是在收纸时被蔡须坤叫住的。

他已经将每个人的笔记按顺序放好，准备跟叶阳去取明天的补给。蔡须坤却把自己那张漏列的草稿抽出来，让他留一会儿。

``下午为什么没往前站？''

薛峰看了一眼纸，知道她问的是什么。那时前面的几个人各自说完，他明明已经核清共同矩阵，却仍等着霍雨浩回头再问一次。

``怕我只是记住了你们讲的。''

蔡须坤没有说他想得太多。她将右侧遮住，换了一个问法，请他指出自己这张纸究竟在哪一列开始漏掉来路。

薛峰接过，沿保留段往后数，找到了那个位置。他说得很慢，又把相邻的一段核了一遍，确认漏的并不是自己最初猜的那一列。

蔡须坤听完，才将答案翻回来。

``刚才这个，你问谁了？''

``没有。''

她点点头，把笔递过去，让他将检查的先后写在页背。薛峰写到中间，自己改了一处，最后将原矩阵重新读全，才慢慢把纸放下。

放下笔后，他忽然抬头，问蔡须坤自己要是又核错了怎么办。蔡须坤正在收另外一页，闻言把纸放下，让他看自己刚补过的那一列。

``就把错的地方说出来，再核。''

``你不会觉得，我只会做这个还做不好？''

蔡须坤望着他，想起自己刚才怎样等别人先回答，过了一会儿才摇头。

``你昨天把我问住的时候，也没说我不该学M。''

薛峰愣了一下，低头看那张纸。他把手上沾着的炭末擦去。蔡须坤等他把最后一行看完，将自己的名字记在复核处，又给小舞留了一个位置，让她明天也来试。薛峰把本子收好，这次没再把自己写的说明藏到最下面。

门口的叶阳没有催，见他还在看自己写的那几行，便把他的包先背了起来。第二天，霍雨浩刚将本子摊开，薛峰就走到桌边，先说了自己今天准备核的那一段。

\NovelDisplay{M03593R01}{%
(0)(1,1,1)(2,1,1)(3,1,1)(2,1)(1,1,1)(2,1,1)(3,1)(4,2,1)(5,2,1)(6,2,1)%
}
\NovelDisplay{M03593R02}{%
(5,2)(4,2,1)(5,2,1)(6,2)(7,3,1)(8,3,1)(9,3,1)(8,3)%
}

\NovelDisplay{M03594R01}{%
= \psi(I_\omega\times \Omega_{I_2+1})%
}

\NovelDisplay{M03595R01}{%
= \psi(M^2\times \omega+\psi_{\psi_{M}(M^2\times \omega+M)}(M^2\times 2+M))%
}

\NovelDisplay{M03596R01}{%
(0)(1,1,1)(2,1,1)(3,1,1)(2,1)(1,1,1)(2,1,1)(3,1)(4,2,1)(5,2,1)(6,2,1)%
}
\NovelDisplay{M03596R02}{%
(5,2)(4,2,1)(5,2,1)(6,2)(7,3,1)(8,3,1)(9,3,1)(8,3)(7,3,1)(8,3,1)(9,3)%
}
\NovelDisplay{M03596R03}{%
(10,4,1)(11,4,1)(12,4,1)%
}

\NovelDisplay{M03597R01}{%
= \psi(I_\omega\times \psi_{I_3}(I_\omega))%
}

\NovelDisplay{M03598R01}{%
= \psi(M^2\times \omega+\psi_{\psi_{M}(M^2\times \omega+M)}(M^2\times 2+M\times \psi_{\psi_{M}(M^2\times 3)}(M^2\times \omega)))%
}

\NovelDisplay{M03599R01}{%
(0)(1,1,1)(2,1,1)(3,1,1)(2,1)(1,1,1)(2,1,1)(3,1,1)%
}

\NovelDisplay{M03600R01}{%
= \psi(I_\omega^2)%
}

\NovelDisplay{M03601R01}{%
= \psi(M^2\times \omega+\psi_{\psi_{M}(M^2\times \omega+M)}(M^2\times \omega))%
}

\NovelDisplay{M03602R01}{%
(0)(1,1,1)(2,1,1)(3,1,1)(2,1)(3,2)%
}

\NovelDisplay{M03603R01}{%
= \psi(\Omega_{I_\omega+1})%
}

\NovelDisplay{M03604R01}{%
= \psi(M^2\times \omega+M)%
}

这一组比较更高右乘因子并到达\allowbreak{}\NovelInline{I06068}{I_\omega}\allowbreak{}的平方，随后转到后继Ω下标，二者不能沿同一末尾动作混读。

\NovelDisplay{M03605R01}{%
(0)(1,1,1)(2,1,1)(3,1,1)(2,1)(3,2,1)(4,2,1)(5,2)(4)%
}

\NovelDisplay{M03606R01}{%
= \psi(I_{\omega+1})%
}

\NovelDisplay{M03607R01}{%
= \psi(M^2\times \omega+M^2)%
}

\NovelDisplay{M03608R01}{%
(0)(1,1,1)(2,1,1)(3,1,1)(2,1)(3,2,1)(4,2,1)(5,2)(4,2)(5,3,1)(6,3,1)%
}
\NovelDisplay{M03608R02}{%
(7,3)(6)%
}

\NovelDisplay{M03609R01}{%
= \psi(I_{\omega+1}\times 2)%
}

\NovelDisplay{M03610R01}{%
= \psi(M^2\times \omega+M^2+\psi_{M}(M^2\times \omega + M^2))%
}

\NovelDisplay{M03611R01}{%
(0)(1,1,1)(2,1,1)(3,1,1)(2,1)(3,2,1)(4,2,1)(5,2)(4,2,1)%
}

\NovelDisplay{M03612R01}{%
= \psi(I_{\omega+1}\times \omega)%
}

\NovelDisplay{M03613R01}{%
= \psi(M^2\times \omega+M^2+\psi_{\psi_{M}(M^2\times \omega + M^2+M)}(1))%
}

\NovelDisplay{M03614R01}{%
(0)(1,1,1)(2,1,1)(3,1,1)(2,1)(3,2,1)(4,2,1)(5,2)(6,3)%
}

\NovelDisplay{M03615R01}{%
= \psi(\Omega_{I_{\omega+1}})%
}

\NovelDisplay{M03616R01}{%
= \psi(M^2\times \omega+M^2+M)%
}

下标由ω转向\allowbreak{}\NovelInline{I06069}{\omega+1}\allowbreak{}后，应分别核对I对象的变化及它后面的Ω写法，组末的特殊下标对应保留待核。

站上三百五十五米时，蔡须坤先回头去看唐三。

最后那段，他比她更早解释清楚，却在落脚时被湿滑的石面绊了一下。此刻他已经站稳，把绳子留在身侧，正陪小舞重新指出自己刚才接过的那一处。

卢本韦坐下来，解开前臂上固定I册的带子。勒得太久，腕边留下两道浅印，古月娜将本子接走，她才慢慢活动手指。薛峰在较低处等叶阳，将今天没被雨打湿的一面朝内折好，再递给他。

蔡须坤摸过石桩，把这个高度念了一遍。

离三百七十米还有十五米。她想起一开始只盯着五百米时，眼前的每一次停顿都显得很长；如今她能说出今天自己走清了哪一段，也能指出身边是谁替她留住了不熟的地方。

大家没有在这块窄台上挤太久。奥斯卡分过水，夏夜星空听完他们今天各自说明的一处，便让十一人先回去歇。

蔡须坤把包背到肩上，走到转弯处，向下方的小舞伸出手。小舞扶着落稳，再把手伸给后面的人。

\chapter{第一百二十一章：IFP的到达}\label{ux7b2cux4e00ux767eux4e8cux5341ux4e00ux7ae0ifpux7684ux5230ux8fbe}

过了三百五十五米，古月娜将前一天留下的两处记号重新翻给大家看。

讲义里有一条Ω下标写法仍需核对，随后那条I下标中的后继，M侧也不能只凭相邻规律便直接换过去。原式都留在本子里，旁边写着具体的问题，今天没有人为了让后面的路看起来连贯，先将它们改成自己喜欢的答案。

霍雨浩请大家继续分清：矩阵的实际展开，可以照原版规则独立复算；两套表示之间某个特殊对应的来路，还需要另外说明。前者已经算清的部分仍在，后者也不能借此一并算作证明。

蔡须坤将疑问页夹好，抬头看向上方的雾。她已不再因为本子里还留着空处，就觉得刚走过的高度也全都失去了。

卢本韦把I册放到她旁边。唐三和小舞各自展开今天要核的那一页，等薛峰从共同左列接起，大家才继续向上。

\NovelDisplay{M03617R01}{%
(0)(1,1,1)(2,1,1)(3,1,1)(2,1)(3,2,1)(4,2,1)(5,2)(6,3,1)(7,3,1)(8,3)(7)%
}

\NovelDisplay{M03618R01}{%
= \psi(I_{\omega+2})%
}

\NovelDisplay{M03619R01}{%
= \psi(M^2\times \omega+M^2\times 2)%
}

\NovelDisplay{M03620R01}{%
(0)(1,1,1)(2,1,1)(3,1,1)(2,1)(3,2,1)(4,2,1)(5,2,1)%
}

\NovelDisplay{M03621R01}{%
= \psi(I_{\omega\times 2})%
}

\NovelDisplay{M03622R01}{%
= \psi(M^2\times \omega\times 2)%
}

\NovelDisplay{M03623R01}{%
(0)(1,1,1)(2,1,1)(3,1,1)(2,1)(3,2,1)(4,2,1)(5,2,1)(4,2)(5,3)%
}

\NovelDisplay{M03624R01}{%
= \psi(\Omega_{I_{\omega\times 2}+1})%
}

\NovelDisplay{M03625R01}{%
= \psi(M^2\times \omega\times 2+M)%
}

\NovelDisplay{M03626R01}{%
(0)(1,1,1)(2,1,1)(3,1,1)(2,1)(3,2,1)(4,2,1)(5,2,1)(4,2)(5,3,1)%
}
\NovelDisplay{M03626R02}{%
(6,3,1)(7,3)(6)%
}

\NovelDisplay{M03627R01}{%
= \psi(I_{\omega\times 2+1})%
}

\NovelDisplay{M03628R01}{%
= \psi(M^2\times \omega\times 2+M\times 2)%
}

卢本韦将下标后继与ω倍数的不同先标清，其中\allowbreak{}\NovelInline{I06070}{\omega\times2+1}\allowbreak{}的M侧尾项仍照讲义保留，暂不据此类推。

\NovelDisplay{M03629R01}{%
(0)(1,1,1)(2,1,1)(3,1,1)(2,1)(3,2,1)(4,2,1)(5,2,1)(4,2)(5,3,1)%
}
\NovelDisplay{M03629R02}{%
(6,3,1)(7,3,1)%
}

\NovelDisplay{M03630R01}{%
= \psi(I_{\omega\times 3})%
}

\NovelDisplay{M03631R01}{%
= \psi(M^2\times \omega\times 3)%
}

\NovelDisplay{M03632R01}{%
(0)(1,1,1)(2,1,1)(3,1,1)(2,1,1)%
}

\NovelDisplay{M03633R01}{%
= \psi(I_{\psi(2)})%
}

\NovelDisplay{M03634R01}{%
= \psi(M^2\times \psi(2))%
}

\NovelDisplay{M03635R01}{%
(0)(1,1,1)(2,1,1)(3,1,1)(2,1,1)(2,1,1)%
}

\NovelDisplay{M03636R01}{%
= \psi(I_{\psi(3)})%
}

\NovelDisplay{M03637R01}{%
= \psi(M^2\times \psi(3))%
}

\NovelDisplay{M03638R01}{%
(0)(1,1,1)(2,1,1)(3,1,1)(2,1,1)(3,1)%
}

\NovelDisplay{M03639R01}{%
= \psi(I_{\Omega})%
}

\NovelDisplay{M03640R01}{%
= \psi(M^2\times \psi_{M}(M))%
}

这几条从有限ω倍数进入ψ值及Ω下标，真正变化的是I的下标，不能先把整个\allowbreak{}\NovelInline{I06071}{I_\alpha}\allowbreak{}当成I的普通幂。

这一处，霍雨浩请大家暂时放下两个右端，只对共同的原版矩阵取项。

唐三先将末列最高非零行找出，小舞沿对应的父项往前核，薛峰据此生成三份有限项。古月娜检查保留段有没有被多动一次，蔡须坤与卢本韦则各自重算一遍，再将所有结果对到同一个原式。

\NovelDisplay{M03641R01}{%
\text{展开第零项：}\quad(0)(1,1,1)(2,1,1)(3,1,1)(2,1,1)%
}

\NovelDisplay{M03642R01}{%
\text{展开第一项：}\quad(0)(1,1,1)(2,1,1)(3,1,1)(2,1,1)(3)(4,1,1)(5,1,1)%
}
\NovelDisplay{M03642R02}{%
(6,1,1)(5,1,1)%
}

\NovelDisplay{M03643R01}{%
\text{展开第二项：}\quad(0)(1,1,1)(2,1,1)(3,1,1)(2,1,1)(3)(4,1,1)(5,1,1)(6,1,1)%
}
\NovelDisplay{M03643R02}{%
(5,1,1)(6)(7,1,1)(8,1,1)(9,1,1)(8,1,1)%
}

这三次取项分别依原矩阵生成，新增段的行坐标按当前规则变化，有限项再长也不能直接充作原式本身。

蔡须坤在第二项里少抄过一列，薛峰替她指出后，她当着大家的面补回去，又从最初那列零核了一次。

唐三问她为什么不只改错处，她将新增段连同保留段一起读完，才答：``我想确认，手里改的是自己刚才真正算出的那一项。''

\NovelDisplay{M03644R01}{%
(0)(1,1,1)(2,1,1)(3,1,1)(2,1,1)(3,1)(1,1)%
}

\NovelDisplay{M03645R01}{%
= \psi(I_{\Omega}+\Omega)%
}

\NovelDisplay{M03646R01}{%
= \psi(M^2\times \psi_{M}(M)+\psi_{M}(M))%
}

\NovelDisplay{M03647R01}{%
(0)(1,1,1)(2,1,1)(3,1,1)(2,1,1)(3,1)(1,1)(2,2,1)(3,2,1)(4,2,1)%
}
\NovelDisplay{M03647R02}{%
(3,2,1)(4,1)(2,2)%
}

\NovelDisplay{M03648R01}{%
= \psi(I_{\Omega}+\Omega_2)%
}

\NovelDisplay{M03649R01}{%
= \psi(M^2\times \psi_{M}(M)+\psi_{M}(M\times 2))%
}

\NovelDisplay{M03650R01}{%
(0)(1,1,1)(2,1,1)(3,1,1)(2,1,1)(3,1)(1,1)(2,2,1)(3,2,1)(4,2,1)%
}
\NovelDisplay{M03650R02}{%
(3,2,1)(4,1)(2,2,1)%
}

\NovelDisplay{M03651R01}{%
= \psi(I_{\Omega}+\Omega_\omega)%
}

\NovelDisplay{M03652R01}{%
= \psi(M^2\times \psi_{M}(M)+\psi_{M}(M\times \omega))%
}

\NovelDisplay{M03653R01}{%
(0)(1,1,1)(2,1,1)(3,1,1)(2,1,1)(3,1)(1,1)(2,2,1)(3,2,1)(4,2,1)(3,2,1)%
}
\NovelDisplay{M03653R02}{%
(4,1)(2,2,1)(3,2,1)(4,1)%
}

\NovelDisplay{M03654R01}{%
= \psi(I_{\Omega}+\Omega_\Omega)%
}

\NovelDisplay{M03655R01}{%
= \psi(M^2\times \psi_{M}(M)+\psi_{M}(M\times \psi_{M}(M)))%
}

新的\allowbreak{}\NovelInline{I06072}{I_\Omega}\allowbreak{}背景之后再接较低Ω项，M侧须完整保留前面的乘积及后加的函数值。

三百五十五米附近，唐三的脚下一滑。

他没有向上扑，扶住身侧的绳子，将重心慢慢退回原来的一级。小舞伸手按住他的肩，等他重新落稳，才将刚才没讲完的本子递回来。

``这里还是你的。''

唐三接过纸，看了一眼自己还在发抖的手，先把前面的背景说了一遍。蔡须坤在上面听着，没有替他补尾句，直到他将那一处函数头自己接清，才回答他后面的追问。

他再次抬脚时，脚尖先擦去石面上的一点水，站稳以后，才往旁边让出位置。

\NovelDisplay{M03656R01}{%
(0)(1,1,1)(2,1,1)(3,1,1)(2,1,1)(3,1)(1,1)(2,2,1)(3,2,1)(4,2,1)(3,2,1)%
}
\NovelDisplay{M03656R02}{%
(4,1)(2,2,1)(3,2,1)(4,2)%
}

\NovelDisplay{M03657R01}{%
= \psi(I_{\Omega}+\Omega_{\Omega_2})%
}

\NovelDisplay{M03658R01}{%
= \psi(M^2\times \psi_{M}(M)+\psi_{M}(M\times \psi_{M}(M\times 2)))%
}

\NovelDisplay{M03659R01}{%
(0)(1,1,1)(2,1,1)(3,1,1)(2,1,1)(3,1)(1,1)(2,2,1)(3,2,1)(4,2,1)(3,2,1)%
}
\NovelDisplay{M03659R02}{%
(4,1)(2,2,1)(3,2,1)(4,2)(3)%
}

\NovelDisplay{M03660R01}{%
= \psi(I_{\Omega}+\psi_{I}(I))%
}

\NovelDisplay{M03661R01}{%
= \psi(M^2\times \psi_{M}(M)+\psi_{\psi_{M}(M^2)}(M^2))%
}

\NovelDisplay{M03662R01}{%
(0)(1,1,1)(2,1,1)(3,1,1)(2,1,1)(3,1)(1,1)(2,2,1)(3,2,1)(4,2,1)(3,2,1)%
}
\NovelDisplay{M03662R02}{%
(4,1)(2,2,1)(3,2,1)(4,2)(5,3,1)(6,3,1)(7,3,1)%
}

\NovelDisplay{M03663R01}{%
= \psi(I_{\Omega}+\psi_{I}(I_\omega))%
}

\NovelDisplay{M03664R01}{%
= \psi(M^2\times \psi_{M}(M)+\psi_{\psi_{M}(M^2)}(M^2 \times \omega))%
}

\NovelDisplay{M03665R01}{%
(0)(1,1,1)(2,1,1)(3,1,1)(2,1,1)(3,1)(1,1)(2,2,1)(3,2,1)(4,2,1)(3,2,1)%
}
\NovelDisplay{M03665R02}{%
(4,1)(2,2,1)(3,2,1)(4,2)(5,3,1)(6,3,1)(7,3,1)(6,3,1)(7,1)%
}

\NovelDisplay{M03666R01}{%
= \psi(I_{\Omega}+\psi_{I}(I_\Omega))%
}

\NovelDisplay{M03667R01}{%
= \psi(M^2\times \psi_{M}(M)+\psi_{\psi_{M}(M^2)}(M^2 \times \psi_{M}(M)))%
}

尾部的Ω下标与I层ψ函数继续推进，新增函数各自拥有自己的参数，外层加法次序保持原式。

\NovelDisplay{M03668R01}{%
(0)(1,1,1)(2,1,1)(3,1,1)(2,1,1)(3,1)(1,1)(2,2,1)(3,2,1)(4,2,1)(3,2,1)%
}
\NovelDisplay{M03668R02}{%
(4,1)(2,2,1)(3,2,1)(4,2)(5,3,1)(6,3,1)(7,3,1)(6,3,1)(7,1)(3,2)(4,3,1)%
}
\NovelDisplay{M03668R03}{%
(5,3,1)(6,3)(5)%
}

\NovelDisplay{M03669R01}{%
= \psi(I_{\Omega}+I)%
}

\NovelDisplay{M03670R01}{%
= \psi(M^2\times \psi_{M}(M)+\psi_{M}(M^2))%
}

\NovelDisplay{M03671R01}{%
(0)(1,1,1)(2,1,1)(3,1,1)(2,1,1)(3,1)(1,1)(2,2,1)(3,2,1)(4,2,1)(3,2,1)%
}
\NovelDisplay{M03671R02}{%
(4,1)(2,2,1)(3,2,1)(4,2)(5,3,1)(6,3,1)(7,3,1)(6,3,1)(7,1)(3,2)(4,3,1)%
}
\NovelDisplay{M03671R03}{%
(5,3,1)(6,3)(7,4,1)(8,4,1)(9,4,1)(8,4,1)(9,1)%
}

\NovelDisplay{M03672R01}{%
= \psi(I_{\Omega}+\psi_{\Omega_{I+1}}(I_\Omega))%
}

\NovelDisplay{M03673R01}{%
= \psi(M^2\times \psi_{M}(M)+\psi_{\psi_{M}(M^2+M)}(M^2 \times \psi_{M}(M)))%
}

\NovelDisplay{M03674R01}{%
(0)(1,1,1)(2,1,1)(3,1,1)(2,1,1)(3,1)(1,1)(2,2,1)(3,2,1)(4,2,1)(3,2,1)%
}
\NovelDisplay{M03674R02}{%
(4,1)(2,2,1)(3,2,1)(4,2)(5,3,1)(6,3,1)(7,3,1)(6,3,1)(7,1)(3,2,1)%
}

\NovelDisplay{M03675R01}{%
= \psi(I_{\Omega}+\psi_{\Omega_{I+1}}(I_\Omega+1))%
}

\NovelDisplay{M03676R01}{%
= \psi(M^2\times \psi_{M}(M)+\psi_{\psi_{M}(M^2+M)}(M^2 \times \psi_{M}(M)+1))%
}

\NovelDisplay{M03677R01}{%
(0)(1,1,1)(2,1,1)(3,1,1)(2,1,1)(3,1)(1,1)(2,2,1)(3,2,1)(4,2,1)(3,2,1)%
}
\NovelDisplay{M03677R02}{%
(4,1)(2,2,1)(3,2,1)(4,2)(5,3,1)(6,3,1)(7,3,1)(6,3,1)(7,1)(5,3)%
}

\NovelDisplay{M03678R01}{%
= \psi(I_{\Omega}+\Omega_{I+1})%
}

\NovelDisplay{M03679R01}{%
= \psi(M^2\times \psi_{M}(M)+\psi_{M}(M^2+M))%
}

这组将尾加I、较高内函数及Ω下标并列，最早变化的函数头比最右边的短尾巴更值得先核。

\NovelDisplay{M03680R01}{%
(0)(1,1,1)(2,1,1)(3,1,1)(2,1,1)(3,1)(1,1)(2,2,1)(3,2,1)(4,2,1)(3,2,1)%
}
\NovelDisplay{M03680R02}{%
(4,1)(2,2,1)(3,2,1)(4,2)(5,3,1)(6,3,1)(7,3,1)(6,3,1)(7,1)(5,3,1)%
}
\NovelDisplay{M03680R03}{%
(6,3,1)(7,3)(6)%
}

\NovelDisplay{M03681R01}{%
= \psi(I_{\Omega}+\psi_{I_2}(I_2))%
}

\NovelDisplay{M03682R01}{%
= \psi(M^2\times \psi_{M}(M)+\psi_{\psi_{M}(M^2\times 2)}(M^2\times 2))%
}

\NovelDisplay{M03683R01}{%
(0)(1,1,1)(2,1,1)(3,1,1)(2,1,1)(3,1)(1,1)(2,2,1)(3,2,1)(4,2,1)(3,2,1)%
}
\NovelDisplay{M03683R02}{%
(4,1)(2,2,1)(3,2,1)(4,2)(5,3,1)(6,3,1)(7,3,1)(6,3,1)(7,1)(5,3,1)(6,3,1)%
}
\NovelDisplay{M03683R03}{%
(7,3)(8,4,1)(9,4,1)(10,4,1)(9,4,1)(10,1)(8,4)%
}

\NovelDisplay{M03684R01}{%
= \psi(I_{\Omega}+\Omega_{I_2+1})%
}

\NovelDisplay{M03685R01}{%
= \psi(M^2\times \psi_{M}(M)+\psi_{M}(M^2\times 2+M))%
}

\NovelDisplay{M03686R01}{%
(0)(1,1,1)(2,1,1)(3,1,1)(2,1,1)(3,1)(1,1)(2,2,1)(3,2,1)(4,2,1)(3,2,1)%
}
\NovelDisplay{M03686R02}{%
(4,1)(2,2,1)(3,2,1)(4,2,1)%
}

\NovelDisplay{M03687R01}{%
= \psi(I_{\Omega}+I_\omega)%
}

\NovelDisplay{M03688R01}{%
= \psi(M^2\times \psi_{M}(M)+\psi_{M}(M^2\times \omega))%
}

\NovelDisplay{M03689R01}{%
(0)(1,1,1)(2,1,1)(3,1,1)(2,1,1)(3,1)(1,1)(2,2,1)(3,2,1)(4,2,1)(3,2,1)%
}
\NovelDisplay{M03689R02}{%
(4,1)(2,2,1)(3,2,1)(4,2,1)(3,2)(4,3,1)(5,3,1)(6,3,1)(5,3,1)(6,1)(4,3)%
}

\NovelDisplay{M03690R01}{%
= \psi(I_{\Omega}+\Omega_{I_\omega+1})%
}

\NovelDisplay{M03691R01}{%
= \psi(M^2\times \psi_{M}(M)+\psi_{M}(M^2\times \omega+M))%
}

从\allowbreak{}\NovelInline{I06073}{I_2}\allowbreak{}层的尾部转到\allowbreak{}\NovelInline{I06074}{I_\omega}\allowbreak{}层时，每个下标对象都须使用当前明确的表示，不能只凭名字相似搬用旧项。

天色暗下来，石壁却仍反着一点冷光。

霍雨浩抬头看了几次，最后把笔合起。蔡须坤还想把这一组读完，刚念了半行，卢本韦便听出她将已经核过的外层也念乱了。

``今天到这里。''

她这次没有争。古月娜先替她背好包，薛峰将今天停下的位置写在最上面，免得明天大家只记得疲惫，忘了真正断在哪里。

回到山脚后，奥斯卡留着热水。十一人轮流将手浸暖，没有在饭桌上继续抢着证明自己还能往下读。蔡须坤睡前仍想起那条长式，却终于将本子放在了枕边，而没有重新点灯。

\NovelDisplay{M03692R01}{%
(0)(1,1,1)(2,1,1)(3,1,1)(2,1,1)(3,1)(1,1)(2,2,1)(3,2,1)(4,2,1)(3,2,1)%
}
\NovelDisplay{M03692R02}{%
(4,1)(2,2,1)(3,2,1)(4,2,1)(3,2)(4,3,1)(5,3,1)(6,3,1)(5,3,1)(6,1)(4,3,1)%
}
\NovelDisplay{M03692R03}{%
(5,3,1)(6,3)(5)%
}

\NovelDisplay{M03693R01}{%
= \psi(I_{\Omega}+I_{\omega+1})%
}

\NovelDisplay{M03694R01}{%
= \psi(M^2\times \psi_{M}(M)+\psi_{M}(M^2\times \omega+M^2))%
}

\NovelDisplay{M03695R01}{%
(0)(1,1,1)(2,1,1)(3,1,1)(2,1,1)(3,1)(1,1)(2,2,1)(3,2,1)(4,2,1)(3,2,1)%
}
\NovelDisplay{M03695R02}{%
(4,1)(2,2,1)(3,2,1)(4,2,1)(3,2)(4,3,1)(5,3,1)(6,3,1)(5,3,1)(6,1)(4,3,1)%
}
\NovelDisplay{M03695R03}{%
(5,3,1)(6,3,1)%
}

\NovelDisplay{M03696R01}{%
= \psi(I_{\Omega}+I_{\omega\times 2})%
}

\NovelDisplay{M03697R01}{%
= \psi(M^2\times \psi_{M}(M)+\psi_{M}(M^2\times \omega\times 2))%
}

\NovelDisplay{M03698R01}{%
(0)(1,1,1)(2,1,1)(3,1,1)(2,1,1)(3,1)(1,1)(2,2,1)(3,2,1)(4,2,1)(3,2,1)%
}
\NovelDisplay{M03698R02}{%
(4,1)(2,2,1)(3,2,1)(4,2,1)(3,2,1)%
}

\NovelDisplay{M03699R01}{%
= \psi(I_{\Omega}+I_{\psi(2)})%
}

\NovelDisplay{M03700R01}{%
= \psi(M^2\times \psi_{M}(M)+\psi_{M}(M^2\times \psi(2)))%
}

\NovelDisplay{M03701R01}{%
(0)(1,1,1)(2,1,1)(3,1,1)(2,1,1)(3,1)(1,1)(2,2,1)(3,2,1)(4,2,1)(3,2,1)%
}
\NovelDisplay{M03701R02}{%
(4,1)(2,2,1)(3,2,1)(4,2,1)(3,2,1)(4,1)%
}

\NovelDisplay{M03702R01}{%
= \psi(I_{\Omega} \times 2)%
}

\NovelDisplay{M03703R01}{%
= \psi(M^2\times \psi_{M}(M)+\psi_{M}(M^2\times \psi_{M}(M)))%
}

蔡须坤归纳道：``下标里的后继、ω倍数及ψ值有各自对应，末尾整份对象的重复又回到外层参数里。''

\NovelDisplay{M03704R01}{%
(0)(1,1,1)(2,1,1)(3,1,1)(2,1,1)(3,1)(1,1)(2,2,1)(3,2,1)(4,2,1)(3,2,1)%
}
\NovelDisplay{M03704R02}{%
(4,1)(3,2)(4,3)%
}

\NovelDisplay{M03705R01}{%
= \psi(\Omega_{I_{\Omega}+1})%
}

\NovelDisplay{M03706R01}{%
= \psi(M^2\times \psi_{M}(M)+M)%
}

\NovelDisplay{M03707R01}{%
(0)(1,1,1)(2,1,1)(3,1,1)(2,1,1)(3,1)(1,1)(2,2,1)(3,2,1)(4,2,1)(3,2,1)%
}
\NovelDisplay{M03707R02}{%
(4,1)(3,2)(4,3,1)(5,3,1)(6,3)(5)%
}

\NovelDisplay{M03708R01}{%
= \psi(I_{\Omega+1})%
}

\NovelDisplay{M03709R01}{%
= \psi(M^2\times \psi_{M}(M)+M^2)%
}

\NovelDisplay{M03710R01}{%
(0)(1,1,1)(2,1,1)(3,1,1)(2,1,1)(3,1)(1,1)(2,2,1)(3,2,1)(4,2,1)(3,2,1)%
}
\NovelDisplay{M03710R02}{%
(4,1)(3,2)(4,3,1)(5,3,1)(6,3,1)%
}

\NovelDisplay{M03711R01}{%
= \psi(I_{\Omega+\omega})%
}

\NovelDisplay{M03712R01}{%
= \psi(M^2\times \psi_{M}(M)+M^2\times \omega)%
}

\NovelDisplay{M03713R01}{%
(0)(1,1,1)(2,1,1)(3,1,1)(2,1,1)(3,1)(1,1)(2,2,1)(3,2,1)(4,2,1)(3,2,1)%
}
\NovelDisplay{M03713R02}{%
(4,1)(3,2)(4,3,1)(5,3,1)(6,3,1)(5,3,1)(6,1)%
}

\NovelDisplay{M03714R01}{%
= \psi(I_{\Omega\times 2})%
}

\NovelDisplay{M03715R01}{%
= \psi(M^2\times \psi_{M}(M)\times 2)%
}

这组把Ω后继下标、I的下标后继与ω增量分别放回M式，有限新增的M与新增的\allowbreak{}\NovelInline{I06075}{M^2}\allowbreak{}不可混认。

\NovelDisplay{M03716R01}{%
(0)(1,1,1)(2,1,1)(3,1,1)(2,1,1)(3,1)(1,1)(2,2,1)(3,2,1)(4,2,1)(3,2,1)%
}
\NovelDisplay{M03716R02}{%
(4,1)(3,2)(4,3,1)(5,3,1)(6,3,1)(5,3,1)(6,1)(5,3)(6,4,1)(7,4,1)%
}
\NovelDisplay{M03716R03}{%
(8,4,1)(7,4,1)(8,1)%
}

\NovelDisplay{M03717R01}{%
= \psi(I_{\Omega\times 3})%
}

\NovelDisplay{M03718R01}{%
= \psi(M^2\times \psi_{M}(M)\times 3)%
}

\NovelDisplay{M03719R01}{%
(0)(1,1,1)(2,1,1)(3,1,1)(2,1,1)(3,1)(1,1)(2,2,1)(3,2,1)(4,2,1)%
}
\NovelDisplay{M03719R02}{%
(3,2,1)(4,1)(3,2,1)%
}

\NovelDisplay{M03720R01}{%
= \psi(I_{\Omega\times \omega})%
}

\NovelDisplay{M03721R01}{%
= \psi(M^2\times \psi_{\psi_{M}(M\times 2)}(1))%
}

\NovelDisplay{M03722R01}{%
(0)(1,1,1)(2,1,1)(3,1,1)(2,1,1)(3,1)(1,1)(2,2,1)(3,2,1)(4,2,1)(3,2,1)%
}
\NovelDisplay{M03722R02}{%
(4,1)(3,2,1)(4,1)%
}

\NovelDisplay{M03723R01}{%
= \psi(I_{\Omega^2})%
}

\NovelDisplay{M03724R01}{%
= \psi(M^2\times \psi_{\psi_{M}(M\times 2)}(M))%
}

\NovelDisplay{M03725R01}{%
(0)(1,1,1)(2,1,1)(3,1,1)(2,1,1)(3,1)(1,1)(2,2,1)(3,2,1)(4,2,1)%
}
\NovelDisplay{M03725R02}{%
(3,2,1)(4,1)(5,2)%
}

\NovelDisplay{M03726R01}{%
= \psi(I_{\psi_{\Omega_{2}}(\Omega_{2})})%
}

\NovelDisplay{M03727R01}{%
= \psi(M^2\times \psi_{\psi_{M}(M\times 2)}(M\times 2))%
}

下标进入Ω倍数及其平方后，相应变化集中在M侧较低函数的参数中，外面的\allowbreak{}\NovelInline{I06076}{M^2}\allowbreak{}左因子仍保留。

第二次来到同一个位置，蔡须坤仍然受压，前一晚怎样也接不清的那对括号，却没有再混。

她说完，先不动脚，等霍雨浩问过一个相邻例子，才确认自己没有只背下昨晚的顺序。卢本韦把I侧接上，古月娜沿完整矩阵重新找过末列来路，薛峰最后核生成段。

三百六十米的标记露出来时，蔡须坤已经累得笑不出声。她低头按了按膝盖，才发现唐三也正望着同一个刻度。

两人对视了一眼，没有再说还差多少。她把笔递下去，让他先记下自己今天独立解释的一处。

\NovelDisplay{M03728R01}{%
(0)(1,1,1)(2,1,1)(3,1,1)(2,1,1)(3,1)(1,1)(2,2,1)(3,2,1)(4,2,1)%
}
\NovelDisplay{M03728R02}{%
(3,2,1)(4,2)%
}

\NovelDisplay{M03729R01}{%
= \psi(I_{\Omega_{2}})%
}

\NovelDisplay{M03730R01}{%
= \psi(M^2\times \psi_{M}(M\times 2))%
}

\NovelDisplay{M03731R01}{%
(0)(1,1,1)(2,1,1)(3,1,1)(2,1,1)(3,1)(1,1)(2,2,1)(3,2,1)(4,2,1)(3,2,1)%
}
\NovelDisplay{M03731R02}{%
(4,2)(2,2)(3,3,1)(4,3,1)(5,3,1)(4,3,1)(5,2)(3,3,1)(4,3,1)(5,3,1)%
}
\NovelDisplay{M03731R03}{%
(4,3,1)(5,2)%
}

\NovelDisplay{M03732R01}{%
= \psi(I_{\Omega_{2}}\times 2)%
}

\NovelDisplay{M03733R01}{%
= \psi(M^2\times \psi_{M}(M\times 2) + \psi_{M}(M^2\times \psi_{M}(M\times 2)))%
}

\NovelDisplay{M03734R01}{%
(0)(1,1,1)(2,1,1)(3,1,1)(2,1,1)(3,1)(1,1)(2,2,1)(3,2,1)(4,2,1)(3,2,1)%
}
\NovelDisplay{M03734R02}{%
(4,2)(2,2)(3,3,1)(4,3,1)(5,3,1)(4,3,1)(5,2)(4,3)(5,4,1)(6,4,1)(7,4,1)%
}

\NovelDisplay{M03735R01}{%
= \psi(I_{\Omega_{2}+\omega})%
}

\NovelDisplay{M03736R01}{%
= \psi(M^2\times \psi_{M}(M\times 2) + M^2\times \omega)%
}

\NovelDisplay{M03737R01}{%
(0)(1,1,1)(2,1,1)(3,1,1)(2,1,1)(3,1)(1,1)(2,2,1)(3,2,1)(4,2,1)(3,2,1)%
}
\NovelDisplay{M03737R02}{%
(4,2)(2,2)(3,3,1)(4,3,1)(5,3,1)(4,3,1)(5,2)(4,3)(5,4,1)(6,4,1)%
}
\NovelDisplay{M03737R03}{%
(7,4,1)(6,4,1)(7,2)%
}

\NovelDisplay{M03738R01}{%
= \psi(I_{\Omega_{2}\times 2})%
}

\NovelDisplay{M03739R01}{%
= \psi(M^2\times \psi_{M}(M\times 2) \times 2)%
}

I的Ω₂下标、整份对象重复和下标中的新增项，需要分别追溯到完整矩阵，不能把三种倍数放在同一层。

\NovelDisplay{M03740R01}{%
(0)(1,1,1)(2,1,1)(3,1,1)(2,1,1)(3,1)(1,1)(2,2,1)(3,2,1)(4,2,1)(3,2,1)%
}
\NovelDisplay{M03740R02}{%
(4,2)(2,2)(3,3,1)(4,3,1)(5,3,1)(4,3,1)(5,2)(4,3,1)%
}

\NovelDisplay{M03741R01}{%
= \psi(I_{\Omega_{2}\times \omega})%
}

\NovelDisplay{M03742R01}{%
= \psi(M^2\times \psi_{\psi_{M}(M\times 3)}(1))%
}

\NovelDisplay{M03743R01}{%
(0)(1,1,1)(2,1,1)(3,1,1)(2,1,1)(3,1)(1,1)(2,2,1)(3,2,1)(4,2,1)(3,2,1)%
}
\NovelDisplay{M03743R02}{%
(4,2)(2,2)(3,3,1)(4,3,1)(5,3,1)(4,3,1)(5,3)%
}

\NovelDisplay{M03744R01}{%
= \psi(I_{\Omega_{3}})%
}

\NovelDisplay{M03745R01}{%
= \psi(M^2\times \psi_{M}(M\times 3))%
}

\NovelDisplay{M03746R01}{%
(0)(1,1,1)(2,1,1)(3,1,1)(2,1,1)(3,1)(1,1,1)%
}

\NovelDisplay{M03747R01}{%
= \psi(I_{\Omega_{\omega}})%
}

\NovelDisplay{M03748R01}{%
= \psi(M^2\times \psi_{M}(M\times \omega))%
}

\NovelDisplay{M03749R01}{%
(0)(1,1,1)(2,1,1)(3,1,1)(2,1,1)(3,1)(1,1,1)(2,1)(3,2,1)(4,2,1)(5,2,1)%
}
\NovelDisplay{M03749R02}{%
(4,2,1)(5,1)(3,2)%
}

\NovelDisplay{M03750R01}{%
= \psi(I_{\Omega_{\omega}}+\Omega_{\omega+1})%
}

\NovelDisplay{M03751R01}{%
= \psi(M^2\times \psi_{M}(M\times \omega)+\psi_{M}(M\times \omega + M))%
}

古月娜指出，Ω下标继续走到Ωω时，M侧也保留明确的ω乘积写法，所有乘号与下标位置都必须读全。

\NovelDisplay{M03752R01}{%
(0)(1,1,1)(2,1,1)(3,1,1)(2,1,1)(3,1)(1,1,1)(2,1)(3,2,1)(4,2,1)(5,2,1)%
}
\NovelDisplay{M03752R02}{%
(4,2,1)(5,1)(3,2,1)(4,2,1)(5,2,1)(4,2,1)(5,1)(1,1,1)%
}

\NovelDisplay{M03753R01}{%
= \psi(I_{\Omega_{\omega}} \times 2)%
}

\NovelDisplay{M03754R01}{%
= \psi(M^2\times \psi_{M}(M\times \omega)+\psi_{M}(M^2\times \psi_{M}(M\times \omega)))%
}

\NovelDisplay{M03755R01}{%
(0)(1,1,1)(2,1,1)(3,1,1)(2,1,1)(3,1)(1,1,1)(2,1)(3,2,1)(4,2,1)(5,2,1)%
}
\NovelDisplay{M03755R02}{%
(4,2,1)(5,1)(4,2)(5,3,1)(6,3,1)(7,3,1)(6,3,1)(7,1)(1,1,1)%
}

\NovelDisplay{M03756R01}{%
= \psi(I_{\Omega_{\omega}\times 2})%
}

\NovelDisplay{M03757R01}{%
= \psi(M^2\times \psi_{M}(M\times \omega)\times 2)%
}

\NovelDisplay{M03758R01}{%
(0)(1,1,1)(2,1,1)(3,1,1)(2,1,1)(3,1)(1,1,1)(2,1)(3,2,1)(4,2,1)(5,2,1)%
}
\NovelDisplay{M03758R02}{%
(4,2,1)(5,1)(4,2,1)%
}

\NovelDisplay{M03759R01}{%
= \psi(I_{\Omega_{\omega}\times \omega})%
}

\NovelDisplay{M03760R01}{%
= \psi(M^2\times \psi_{\psi_{M}(M\times \omega+M)}(1))%
}

\NovelDisplay{M03761R01}{%
(0)(1,1,1)(2,1,1)(3,1,1)(2,1,1)(3,1)(1,1,1)(2,1)(3,2,1)(4,2,1)%
}
\NovelDisplay{M03761R02}{%
(5,2,1)(4,2,1)(5,2)%
}

\NovelDisplay{M03762R01}{%
= \psi(I_{\Omega_{\omega+1}})%
}

\NovelDisplay{M03763R01}{%
= \psi(M^2\times \psi_{M}(M\times \omega+M))%
}

对象本身乘二与I下标内的Ωω乘二在这组并排出现，M侧各自落在不同深度的函数结构。

越靠近新的边界，I侧的下标越像先前熟悉的整条表达式。

卢本韦有一次读错，自己先笑了，笑完又叹气。她把那一整段重新圈回下标里，抬头问蔡须坤，M侧相应的增长到底落在哪个函数。

蔡须坤答到一半，换她停住。两人并排坐着，谁也没能独自将这一页读完。

古月娜从较低处上来，听过她们各自说清的部分，将共同矩阵摆在中央。薛峰沿前缀核到分歧处，霍雨浩再从内层问起。那一页最终收起时，页边留下了五种轻重不同的笔迹。

\NovelDisplay{M03764R01}{%
(0)(1,1,1)(2,1,1)(3,1,1)(2,1,1)(3,1)(1,1,1)(2,1)(3,2,1)(4,2,1)(5,2,1)%
}
\NovelDisplay{M03764R02}{%
(4,2,1)(5,2)(3,2)(4,3,1)(5,3,1)(6,3,1)(5,3,1)(6,3)%
}

\NovelDisplay{M03765R01}{%
= \psi(I_{\Omega_{\omega+2}})%
}

\NovelDisplay{M03766R01}{%
= \psi(M^2\times \psi_{M}(M\times \omega+M\times 2))%
}

\NovelDisplay{M03767R01}{%
(0)(1,1,1)(2,1,1)(3,1,1)(2,1,1)(3,1)(1,1,1)(2,1)(3,2,1)(4,2,1)(5,2,1)%
}
\NovelDisplay{M03767R02}{%
(4,2,1)(5,2)(3,2,1)%
}

\NovelDisplay{M03768R01}{%
= \psi(I_{\Omega_{\omega\times 2}})%
}

\NovelDisplay{M03769R01}{%
= \psi(M^2\times \psi_{M}(M\times \omega \times 2))%
}

\NovelDisplay{M03770R01}{%
(0)(1,1,1)(2,1,1)(3,1,1)(2,1,1)(3,1)(1,1,1)(2,1)(3,2,1)(4,2,1)(5,2,1)%
}
\NovelDisplay{M03770R02}{%
(4,2,1)(5,2)(3,2,1)(4,2)(5,3,1)(6,3,1)(7,3,1)(6,3,1)(7,3)%
}

\NovelDisplay{M03771R01}{%
= \psi(I_{\Omega_{\omega\times 2+1}})%
}

\NovelDisplay{M03772R01}{%
= \psi(M^2\times \psi_{M}(M\times \omega \times 2+M))%
}

\NovelDisplay{M03773R01}{%
(0)(1,1,1)(2,1,1)(3,1,1)(2,1,1)(3,1)(1,1,1)(2,1)(3,2,1)(4,2,1)(5,2,1)%
}
\NovelDisplay{M03773R02}{%
(4,2,1)(5,2)(3,2,1)(4,2)(5,3,1)(6,3,1)(7,3,1)(6,3,1)(7,3)(5,3,1)%
}

\NovelDisplay{M03774R01}{%
= \psi(I_{\Omega_{\omega\times 3}})%
}

\NovelDisplay{M03775R01}{%
= \psi(M^2\times \psi_{M}(M\times \omega \times 3))%
}

同样的有限后继和ω倍数进入更深Ω下标后，先保留已定的下标背景，再认新增项的作用位置。

\NovelDisplay{M03776R01}{%
(0)(1,1,1)(2,1,1)(3,1,1)(2,1,1)(3,1)(1,1,1)(2,1,1)%
}

\NovelDisplay{M03777R01}{%
= \psi(I_{\Omega_{\psi(2)}})%
}

\NovelDisplay{M03778R01}{%
= \psi(M^2\times \psi_{M}(M\times \psi(2)))%
}

\NovelDisplay{M03779R01}{%
(0)(1,1,1)(2,1,1)(3,1,1)(2,1,1)(3,1)(1,1,1)(2,1,1)(3,1)(2)%
}

\NovelDisplay{M03780R01}{%
= \psi(I_{\psi_{I}(I)})%
}

\NovelDisplay{M03781R01}{%
= \psi(M^2\times \psi_{\psi_M(M^2)}(M^2))%
}

\NovelDisplay{M03782R01}{%
(0)(1,1,1)(2,1,1)(3,1,1)(2,1,1)(3,1)(1,1,1)(2,1,1)(3,1)(4,2,1)%
}
\NovelDisplay{M03782R02}{%
(5,2,1)(6,2,1)%
}

\NovelDisplay{M03783R01}{%
= \psi(I_{\psi_{I}(I_\omega)})%
}

\NovelDisplay{M03784R01}{%
= \psi(M^2\times \psi_{\psi_M(M^2)}(M^2\times \omega))%
}

\NovelDisplay{M03785R01}{%
(0)(1,1,1)(2,1,1)(3,1,1)(2,1,1)(3,1)(1,1,1)(2,1,1)(3,1)(4,2,1)(5,2,1)%
}
\NovelDisplay{M03785R02}{%
(6,2,1)(5,2,1)(6,1)%
}

\NovelDisplay{M03786R01}{%
= \psi(I_{\psi_{I}(I_\Omega)})%
}

\NovelDisplay{M03787R01}{%
= \psi(M^2\times \psi_{\psi_M(M^2)}(M^2\times \psi_{M}(M)))%
}

这里的I下标直接取先前ψ函数值，M侧的函数头随之变化，不能将整个下标缩成一个熟悉的单字母。

三百六十五米过后，蔡须坤又看见一处需要停笔的对应。

讲义在一条尾函数中写了\allowbreak{}\NovelInline{I06077}{I_2}\allowbreak{}，M侧相应位置却给出\allowbreak{}\NovelInline{I06078}{M^2+M}\allowbreak{}，与此前用于\allowbreak{}\NovelInline{I06079}{I_2}\allowbreak{}的\allowbreak{}\NovelInline{I06080}{M^2\times2}\allowbreak{}写法不能直接互相代换。她将下面这一小组原式照录，把有疑问的那一条单独注明待核，没有用它去证明后续节点。

古月娜接过纸，先看清她圈的是内参数，再点了点头。两人都没有提那晚关于空中楼阁的争执，手中的笔却在同一个位置停住了。

薛峰保留原左列，霍雨浩请其余人各自从能核清的原版展开继续。未解决的问题仍留在本子里，已经走稳的结构也没有因此被一并抹去。

\NovelDisplay{M03788R01}{%
(0)(1,1,1)(2,1,1)(3,1,1)(2,1,1)(3,1)(1,1,1)(2,1,1)(3,1)(4,2,1)(5,2,1)%
}
\NovelDisplay{M03788R02}{%
(6,2,1)(5,2,1)(6,1)(2)%
}

\NovelDisplay{M03789R01}{%
= \psi(I_I)%
}

\NovelDisplay{M03790R01}{%
= \psi(M^2\times \psi_M(M^2))%
}

\NovelDisplay{M03791R01}{%
(0)(1,1,1)(2,1,1)(3,1,1)(2,1,1)(3,1)(1,1,1)(2,1,1)(3,1)(4,2,1)(5,2,1)%
}
\NovelDisplay{M03791R02}{%
(6,2,1)(5,2,1)(6,1)(2,1)(3,2,1)(4,2,1)(5,2)(4)%
}

\NovelDisplay{M03792R01}{%
= \psi(I_I+I)%
}

\NovelDisplay{M03793R01}{%
= \psi(M^2\times \psi_M(M^2)+\psi_{M}(M^2))%
}

\NovelDisplay{M03794R01}{%
(0)(1,1,1)(2,1,1)(3,1,1)(2,1,1)(3,1)(1,1,1)(2,1,1)(3,1)(4,2,1)(5,2,1)%
}
\NovelDisplay{M03794R02}{%
(6,2,1)(5,2,1)(6,1)(2,1)(3,2,1)(4,2,1)(5,2)(6,3,1)(7,3,1)(8,3)(7)%
}

\NovelDisplay{M03795R01}{%
= \psi(I_I+\psi_{\Omega_{I+1}}(I_2))%
}

\NovelDisplay{M03796R01}{%
= \psi(M^2\times \psi_M(M^2)+\psi_{\psi_{M}(M^2+M)}(M^2+M))%
}

\NovelDisplay{M03797R01}{%
(0)(1,1,1)(2,1,1)(3,1,1)(2,1,1)(3,1)(1,1,1)(2,1,1)(3,1)(4,2,1)(5,2,1)%
}
\NovelDisplay{M03797R02}{%
(6,2,1)(5,2,1)(6,1)(2,1)(3,2,1)(4,2,1)(5,2,1)%
}

\NovelDisplay{M03798R01}{%
= \psi(I_I+\psi_{\Omega_{I+1}}(I_\omega))%
}

\NovelDisplay{M03799R01}{%
= \psi(M^2\times \psi_M(M^2)+\psi_{\psi_{M}(M^2+M)}(M^2\times \omega))%
}

这一组进入\allowbreak{}\NovelInline{I06081}{I_I}\allowbreak{}并继续加尾项，其中含\allowbreak{}\NovelInline{I06082}{I_2}\allowbreak{}的特殊内参数对应已另留待核，不能作为换写通则使用。

\NovelDisplay{M03800R01}{%
(0)(1,1,1)(2,1,1)(3,1,1)(2,1,1)(3,1)(1,1,1)(2,1,1)(3,1)(4,2,1)(5,2,1)%
}
\NovelDisplay{M03800R02}{%
(6,2,1)(5,2,1)(6,1)(2,1)(3,2,1)(4,2,1)(5,2,1)(4,2,1)(5,1)(2)%
}

\NovelDisplay{M03801R01}{%
= \psi(I_I+\psi_{\Omega_{I+1}}(I_I))%
}

\NovelDisplay{M03802R01}{%
= \psi(M^2\times \psi_M(M^2)+\psi_{\psi_{M}(M^2+M)}(M^2\times \psi_{M}(M^2)))%
}

\NovelDisplay{M03803R01}{%
(0)(1,1,1)(2,1,1)(3,1,1)(2,1,1)(3,1)(1,1,1)(2,1,1)(3,1)(4,2,1)(5,2,1)%
}
\NovelDisplay{M03803R02}{%
(6,2,1)(5,2,1)(6,1)(2,1,1)%
}

\NovelDisplay{M03804R01}{%
= \psi(I_I+\psi_{\Omega_{I+1}}(I_I+1))%
}

\NovelDisplay{M03805R01}{%
= \psi(M^2\times \psi_M(M^2)+\psi_{\psi_{M}(M^2+M)}(M^2\times \psi_{M}(M^2)+1))%
}

\NovelDisplay{M03806R01}{%
(0)(1,1,1)(2,1,1)(3,1,1)(2,1,1)(3,1)(1,1,1)(2,1,1)(3,1)(4,2,1)(5,2,1)%
}
\NovelDisplay{M03806R02}{%
(6,2,1)(5,2,1)(6,1)(4,2)%
}

\NovelDisplay{M03807R01}{%
= \psi(I_I+\Omega_{I+1})%
}

\NovelDisplay{M03808R01}{%
= \psi(M^2\times \psi_M(M^2)+\psi_{M}(M^2+M))%
}

\NovelDisplay{M03809R01}{%
(0)(1,1,1)(2,1,1)(3,1,1)(2,1,1)(3,1)(1,1,1)(2,1,1)(3,1)(4,2,1)(5,2,1)%
}
\NovelDisplay{M03809R02}{%
(6,2,1)(5,2,1)(6,1)(4,2,1)(5,2,1)(6,2,1)(5,2,1)(6,1)(2)%
}

\NovelDisplay{M03810R01}{%
= \psi(I_I\times 2)%
}

\NovelDisplay{M03811R01}{%
= \psi(M^2\times \psi_M(M^2)+\psi_{M}(M^2\times \psi_M(M^2)))%
}

尾部继续从内函数的自下标参数增长到整份\allowbreak{}\NovelInline{I06083}{I_I}\allowbreak{}的重复，每一个加一与乘二仍须留在自己的层次。

古月娜把一截松开的带子重新缝好，递给卢本韦时，指尖被针扎了一下。

伤口很小，她自己没在意，蔡须坤却看见了，顺手将干净布条递过去。古月娜接过，低头包住那一点红，忽然问她，昨天为什么有问题也没有立即开口。

``你那时正在帮小舞。''

``我也能听见你说。''

蔡须坤没接话。她确实已经习惯先等别人问，等身边的人都找到位置，才把自己卡住的地方重新拿出来；若等得太久，有时连先前的问题都在脑中揉成了一团。

古月娜将针收好，指了指她面前的本子。

``到你了。''

蔡须坤把那一页推过去，先说自己能确定的部分，再把最后一点犹疑指出来。古月娜问了两句，便也停住，转头去叫霍雨浩。

卢本韦从旁边伸过手，拿走一张没有写满的纸，给她们重排那段共同前缀。薛峰跟过来看，几个人慢慢围拢，原先一句没说完的问题，终于没有再被蔡须坤留到所有人走后。

等唐三也把自己的想法说完，蔡须坤伸手活动了一下发僵的指节，自己先笑了。古月娜把收好的针放进包里，坐回她旁边，等她将还没说完的后半句接上；窗外天色渐亮，桌边的人也都没有急着站起来，谁都还想听她再说几句。

\NovelDisplay{M03812R01}{%
(0)(1,1,1)(2,1,1)(3,1,1)(2,1,1)(3,1)(1,1,1)(2,1,1)(3,1)(4,2,1)(5,2,1)%
}
\NovelDisplay{M03812R02}{%
(6,2,1)(5,2,1)(6,1)(5,2)(6,3,1)(7,3,1)(8,3)(7)%
}

\NovelDisplay{M03813R01}{%
= \psi(I_{I+1})%
}

\NovelDisplay{M03814R01}{%
= \psi(M^2\times \psi_M(M^2)+M^2)%
}

\NovelDisplay{M03815R01}{%
(0)(1,1,1)(2,1,1)(3,1,1)(2,1,1)(3,1)(1,1,1)(2,1,1)(3,1)(4,2,1)(5,2,1)%
}
\NovelDisplay{M03815R02}{%
(6,2,1)(5,2,1)(6,1)(5,2)(6,3,1)(7,3,1)(8,3,1)(7,3,1)(8,1)(2)%
}

\NovelDisplay{M03816R01}{%
= \psi(I_{I\times 2})%
}

\NovelDisplay{M03817R01}{%
= \psi(M^2\times \psi_M(M^2)\times 2)%
}

\NovelDisplay{M03818R01}{%
(0)(1,1,1)(2,1,1)(3,1,1)(2,1,1)(3,1)(1,1,1)(2,1,1)(3,1)(4,2,1)(5,2,1)%
}
\NovelDisplay{M03818R02}{%
(6,2,1)(5,2,1)(6,1)(5,2,1)%
}

\NovelDisplay{M03819R01}{%
= \psi(I_{I\times \omega})%
}

\NovelDisplay{M03820R01}{%
= \psi(M^2\times \psi_M(M^2)\times \omega)%
}

\NovelDisplay{M03821R01}{%
(0)(1,1,1)(2,1,1)(3,1,1)(2,1,1)(3,1)(1,1,1)(2,1,1)(3,1)(4,2,1)(5,2,1)%
}
\NovelDisplay{M03821R02}{%
(6,2,1)(5,2,1)(6,1)(7,2)%
}

\NovelDisplay{M03822R01}{%
= \psi(I_{\psi_{\Omega_{I+1}}(\Omega_{I+1})})%
}

\NovelDisplay{M03823R01}{%
= \psi(M^2\times \psi_{\psi_M(M^2+M)}(M^2+M))%
}

卢本韦归纳道：``\allowbreak{}\NovelInline{I06084}{I_{I+1}}\allowbreak{}、\allowbreak{}\NovelInline{I06085}{I_{I\times2}}\allowbreak{}和\allowbreak{}\NovelInline{I06086}{I_{I\times\omega}}\allowbreak{}改变的是下标，这与前一组在外层重复整个\allowbreak{}\NovelInline{I06087}{I_I}\allowbreak{}不同。''

最后几页的纸已经被翻软。

小舞在蔡须坤背后换了一次位置，接过她抱得太久的包。蔡须坤说自己能拿，小舞便将两份本子都放到她面前，指着其中一处，让她看自己昨天补上的说明。

蔡须坤读过，抬头时，小舞已经将包背好。

她没再伸手去抢，只把那一处接法说给小舞听，再请她反问自己一个相邻问题。两人一问一答，到下一组开始时，手中的书页终于不再发颤。

\NovelDisplay{M03824R01}{%
(0)(1,1,1)(2,1,1)(3,1,1)(2,1,1)(3,1)(1,1,1)(2,1,1)(3,1)(4,2,1)(5,2,1)%
}
\NovelDisplay{M03824R02}{%
(6,2,1)(5,2,1)(6,2)%
}

\NovelDisplay{M03825R01}{%
= \psi(I_{\Omega_{I+1}})%
}

\NovelDisplay{M03826R01}{%
= \psi(M^2\times \psi_M(M^2+M))%
}

\NovelDisplay{M03827R01}{%
(0)(1,1,1)(2,1,1)(3,1,1)(2,1,1)(3,1)(1,1,1)(2,1,1)(3,1)(4,2,1)(5,2,1)%
}
\NovelDisplay{M03827R02}{%
(6,2,1)(5,2,1)(6,2)(4,2,1)(5,2,1)(6,2)(7,3,1)(8,3,1)(9,3,1)(8,3,1)%
}
\NovelDisplay{M03827R03}{%
(9,2)(5)%
}

\NovelDisplay{M03828R01}{%
= \psi(I_{I_2})%
}

\NovelDisplay{M03829R01}{%
= \psi(M^2\times \psi_M(M^2\times 2))%
}

\NovelDisplay{M03830R01}{%
(0)(1,1,1)(2,1,1)(3,1,1)(2,1,1)(3,1)(1,1,1)(2,1,1)(3,1)(4,2,1)(5,2,1)%
}
\NovelDisplay{M03830R02}{%
(6,2,1)(5,2,1)(6,2)(4,2,1)(5,2,1)(6,2)(7,3,1)(8,3,1)(9,3,1)(8,3,1)(9,3)%
}

\NovelDisplay{M03831R01}{%
= \psi(I_{\Omega_{I_2+1}})%
}

\NovelDisplay{M03832R01}{%
= \psi(M^2\times \psi_M(M^2\times 2+M))%
}

\NovelDisplay{M03833R01}{%
(0)(1,1,1)(2,1,1)(3,1,1)(2,1,1)(3,1)(1,1,1)(2,1,1)(3,1,1)%
}

\NovelDisplay{M03834R01}{%
= \psi(I_{I_\omega})%
}

\NovelDisplay{M03835R01}{%
= \psi(M^2\times \psi_M(M^2\times \omega))%
}

较高Ω对象与更大I下标继续进入I的下标位置时，M式中的乘积和相应函数值都按原有次序保留。

\NovelDisplay{M03836R01}{%
(0)(1,1,1)(2,1,1)(3,1,1)(2,1,1)(3,1)(1,1,1)(2,1,1)(3,1,1)(2,1,1)(3,1)%
}

\NovelDisplay{M03837R01}{%
= \psi(I_{I_\Omega})%
}

\NovelDisplay{M03838R01}{%
= \psi(M^2\times \psi_M(M^2\times \psi_{M}(M)))%
}

\NovelDisplay{M03839R01}{%
(0)(1,1,1)(2,1,1)(3,1,1)(2,1,1)(3,1)(1,1,1)(2,1,1)(3,1,1)(2,1,1)(3,1)%
}
\NovelDisplay{M03839R02}{%
(1,1,1)(2,1,1)(3,1)(4,2,1)(5,2,1)(6,2,1)(5,2,1)(6,2)(4,2,1)(5,2,1)%
}
\NovelDisplay{M03839R03}{%
(6,2,1)(5,2,1)(6,1)(2)%
}

\NovelDisplay{M03840R01}{%
= \psi(I_{I_I})%
}

\NovelDisplay{M03841R01}{%
= \psi(M^2\times \psi_M(M^2\times \psi_{M}(M^2)))%
}

\NovelDisplay{M03842R01}{%
(0)(1,1,1)(2,1,1)(3,1,1)(2,1,1)(3,1)(1,1,1)(2,1,1)(3,1,1)(2,1,1)(3,1)%
}
\NovelDisplay{M03842R02}{%
(1,1,1)(2,1,1)(3,1)(4,2,1)(5,2,1)(6,2,1)(5,2,1)(6,2)(4,2,1)(5,2,1)%
}
\NovelDisplay{M03842R03}{%
(6,2,1)(5,2,1)(6,2)%
}

\NovelDisplay{M03843R01}{%
= \psi(I_{I_{\Omega_{I+1}}})%
}

\NovelDisplay{M03844R01}{%
= \psi(M^2\times \psi_M(M^2\times \psi_{M}(M^2+M)))%
}

\NovelDisplay{M03845R01}{%
(0)(1,1,1)(2,1,1)(3,1,1)(2,1,1)(3,1)(1,1,1)(2,1,1)(3,1,1)(2,1,1)(3,1)%
}
\NovelDisplay{M03845R02}{%
(1,1,1)(2,1,1)(3,1,1)%
}

\NovelDisplay{M03846R01}{%
= \psi(I_{I_{I_\omega}})%
}

\NovelDisplay{M03847R01}{%
= \psi(M^2\times \psi_M(M^2\times \psi_{M}(M^2\times \omega)))%
}

当I下标本身又包含I下标时，先将最里面的对象认全，再逐层返回最外面的ψ，才能说明整条对应。

三百七十米的石桩离他们只有最后几级。

讲义的红色末行很短，最外面的M参数也只剩一个清楚的立方。蔡须坤没有像从前那样先去抓这个短结果，她先沿左列辨认变化，又请卢本韦将I侧那个双参数记号完整读出。

末行后面另附着一道很长的补充题。霍雨浩将它留作抵达前的回看：核清前面已经经过的自下标结构与整份重复，不再给它另加一段高度。

大家在最后那处平石上坐下，唐三接过长矩阵的一端，小舞看中段，薛峰逐列合回完整原式；十一人都说清自己负责的部分，才将红色末行与这一道补充一起收好。

\NovelDisplay{M03848R01}{%
(0)(1,1,1)(2,1,1)(3,1,1)(2,1,1)(3,1)(1,1,1)(2,1,1)(3,1,1)(2,1,1)(3,1)%
}
\NovelDisplay{M03848R02}{%
(1,1,1)(2,1,1)(3,1,1)(2,1,1)(3,1)%
}

\NovelDisplay{M03849R01}{%
= \psi(I_{I_{I_\Omega}})%
}

\NovelDisplay{M03850R01}{%
= \psi(M^2\times \psi_M(M^2\times \psi_{M}(M^2\times \psi_{M}(M))))%
}

\NovelDisplay{M03851R01}{%
\textcolor{red}{(0)(1,1,1)(2,1,1)(3,1,1)(2,1,1)(3,1)(2)}%
}

\NovelDisplay{M03852R01}{%
\textcolor{red}{= \psi(I(1,0))}%
}

\NovelDisplay{M03853R01}{%
\textcolor{red}{= \psi(M^3)}%
}

末组由更深自下标的回看接到IFP的双参数记号与\allowbreak{}\NovelInline{I06088}{M^3}\allowbreak{}，末尾长式另用于复核已经经过的整份重复。

最后起身时，蔡须坤的膝盖僵了。

卢本韦将手伸过来，她扶着站稳，没有借那只手向上跨。台阶里的压力仍在，肩背也仍然沉，可刚才那几处相互回扯的运转，已经能被她逐层分辨。

她先将一只脚落上去，停住，再将最后一处接续在心里从头走完。

霍雨浩站在旁边等。薛峰抱着收好的长页，低声回答唐三最后一个问题；古月娜将自己那一处重新说明过，才抬脚跟上。

三百七十米的横台很窄，站四个人便要错开肩膀。蔡须坤到了以后，先把本子放在石桩后面，再转身去看下方。

卢本韦第二次收脚时有些晃，古月娜扶住她，等她自己稳下，才把手松开。霍雨浩与薛峰随后站到平处，没有继续往上，分别向后面的人留出能够看清纸页的位置。

唐三用了更久。

他在最后几级将同一段重新说了两次，第二次才说清先前那处下标为什么不能与外层乘法混读。小舞听过他的解释，也将自己负责的生成段独立补完，才将脚跟从下面那一级慢慢收起。

诺当的手指冻得发白，奥斯卡把自己的纸先交给叶阳，腾出手替他解开缠在袖口的带子。蔡兴坤将黑子鸡护在包里，轮到自己时才重新拿笔；叶阳核过最后一处保留段，将薛峰的本子递回，终于也站上了横台。

绳子松下来的时候，十一人谁都没有马上说话。

小舞坐在石桩旁，把脸埋在膝上，过了一会儿才抬头。唐三将水递给她，自己仰着脸喘气；诺当终于笑出一声，笑完便扶着腰不愿再动。

蔡须坤伸手摸到石桩上的刻痕。

三百七十。

从三百米到这里，他们停过、退过，许多夜里将同一段算到不愿再看。她仍记得那张写不出第一项的纸，也记得自己曾以为空白会一直留在面前。

理想不提升BMS仍未定义出来，那些期望对应也没有被拿来替换原版答案。山阶上的七十米，是他们沿着原版BMS重新求项、重新辨认I与M的结构，靠自己一步步走出的。

夏夜星空到了，先让所有人歇过气，才问今天最后那条短式。

卢本韦读I侧，蔡须坤接M侧，霍雨浩补全两人说明所用的条件，薛峰将原矩阵与末列的区别指出。唐三和小舞也各说了一处自己如今能够独立接续的位置。

``这一段，走完了。''夏夜星空说。

蔡须坤听见，手指才从那道刻痕上松开。

她没有立刻望向更高的五百米，只低头接过古月娜递来的本子。里面还夹着几条待核的原文对应，封套里的理想矩阵也还没有新添一个取项；这些都被她好好收住，没有因为到了石桩便顺手划去。

卢本韦挨着她坐下，将自己的I册放到同一块干石上。她们的水已经凉了，喝下去仍觉得舒服。

``明天呢？''小舞问。

蔡须坤看向她，又看看身边一个个累得不想动的人。

``先把今天能自己算的，写给彼此看。''

小舞点头，低头把自己的页角抚平。霍雨浩将三百七十记上山图，薛峰压住纸，古月娜等墨干，才把它交回蔡须坤手中。

这一回，她合上本子时，知道明天还愿意把它打开。

\chapter{第一百二十二章：I(1,0)的路}\label{ux7b2cux4e00ux767eux4e8cux5341ux4e8cux7ae0i10ux7684ux8def}

三百七十米的石桩仍在身后，小舞已经向上试了两次。

每次都只多踩半级。她的鞋底贴住石面，腿却像被另一股力量留在原处；只要想用魂力硬抬，胸口便跟着发紧。

唐三把本子翻到上一夜，确认最后几条没有抄漏。蔡须坤却没有往后催。她把两本册子并排摊开，发现旧页末尾与新页开头之间，还隔着一整段没有独立走清的分析。

``这里。''她指给霍雨浩看。

夏夜星空取来另一叠纸，放在他们中间。最前面是熟悉的双参数I，另一边仍能接回M的立方；后面的括号一层套一层，远没有她最初望过去时那么简单。

到了第四个月，他们又在石桩下铺开了纸。薛峰先取出空页，卢本韦翻回I册，古月娜将几块压纸的石头往旁边挪，为后来的人留出地方。

蔡须坤望了一眼五百米那道横线，低头把新页摆正。

``这段先走清。''

先要认清的，是两套写法怎样描述同一个对象。

卢本韦读I侧，蔡须坤把M侧接在下面。她们没有另用一个大写字母藏住长式；读到下标便停，让霍雨浩指出那个下标自己的作用对象。

诺当问，既然同一个东西能换写，能不能挑短的一边就算过关。霍雨浩将两张纸转回原位，让他先说两边各把哪一层写在了哪里。

``换写以后，没看清的地方也得看清。''

诺当点头，在自己的纸上留出下标的位置。

\NovelDisplay{M03854R01}{%
M%
}

\NovelDisplay{M03855R01}{%
\Omega_{\psi_{I(1,0)}(I(1,0))+1}%
}

\NovelDisplay{M03856R01}{%
I(1,0) = \psi_{M}(M^3)%
}

\NovelDisplay{M03857R01}{%
\psi_{I(1,0)}(I(1,0)) = \psi_{\psi_{M}(M^3)}(M^3)%
}

卢本韦先给出I侧的作用对象，蔡须坤再将同一对象写成M侧，下标嵌套中的作用范围由两人分别指清。

小舞将连续出现的那条对应读了两遍，抬头问是不是抄重了。

夏夜星空让她照原页留着，又指出后面那个更复杂的对象。讲义在这里提醒，下标嵌套这一路写法也有尚不能照样处理的情况，眼前几个能写成，并不能替后面所有情况作保证。

小舞把准备划掉的一行重新描清。唐三原想往后替她接，见她仍在找下标，便先将手收了回去。

\NovelDisplay{M03858R01}{%
\psi_{I(1,0)}(I(1,0)) = \psi_{\psi_{M}(M^3)}(M^3)%
}

\NovelDisplay{M03859R01}{%
I_{\psi_{I(1,0)}(I(1,0))+1} = \psi_{\psi_{M}(M^3)}(M^3+1)%
}

\NovelDisplay{M03860R01}{%
\Omega_{\psi_{I(1,0)}(I(1,0))+1} = \psi_{\psi_{\psi_{M}(M^3)}(M^3+1)}(1)%
}

\NovelDisplay{M03861R01}{%
\Omega_{\psi_{I(\omega+1,0)}(I(\omega+1,0))+1}%
}

相同对应保留两次以便回看，而更深的下标也暴露出这种写法的局限，不能因为能写出短例便宣称所有情况都能处理。

接下来用到的，是讲义所说的欠层。

蔡须坤盯着那个不动点关系，最初只觉得纸面少了一层。霍雨浩请她重新读出条件：这里谈的是所列对象在给定映射下的不动点，写法的改变依赖这个关系，不能对任何ψ随手欠去一层。

她将手从最外面的括号移开，沿着里面的乘号重新指过，才把最后那个加项接回原位。

\NovelDisplay{M03862R01}{%
\psi_{I(1,0)}(I(1,0))%
}

\NovelDisplay{M03863R01}{%
\psi_{M}(M^2 \times x)%
}

\NovelDisplay{M03864R01}{%
\psi_{M}(M^2\times \psi_{\psi_{M}(M^3)}(M^3))%
}

\NovelDisplay{M03865R01}{%
\Omega_{\psi_{I(1,0)}(I(1,0))+1} = \psi_{M}(M^2\times \psi_{\psi_{M}(M^3)}(M^3)+M)%
}

这里的欠层依赖讲义说明的不动点关系，末尾的M仍在外层加项的位置，不能把它挪进最内侧的ψ。

``那一直这样写，不就行了？''蔡兴坤问。

夏夜星空让他先把下面几式读完。I侧的下标逐层嵌进自己，M侧也反复套入那段背景，越写越像只需照前一行多添一层。

蔡须坤的笔停住。她想起那张理想不提升BMS的纸，自己当时也是先看中了整齐，随后才发现没有展开的依据。

这一次，她没有把``看着整齐''写进结论里。

\NovelDisplay{M03866R01}{%
\psi_{I(1,0)}(I(1,0)\times 2)%
}

\NovelDisplay{M03867R01}{%
I_{\psi_{I(1,0)}(I(1,0))+1} = \psi_{M}(M^2\times \psi_{\psi_{M}(M^3)}(M^3) + M^2)%
}

\NovelDisplay{M03868R01}{%
I_{I_{\psi_{I(1,0)}(I(1,0))+1}} = \psi_{M}(M^2\times \psi_{M}(M^2\times \psi_{\psi_{M}(M^3)}(M^3) + M^2))%
}

\NovelDisplay{M03869R01}{%
I_{I_{I_{\psi_{I(1,0)}(I(1,0))+1}}} = \psi_{M}(M^2\times \psi_{M}(M^2\times \psi_{M}(M^2\times \psi_{\psi_{M}(M^3)}(M^3) + M^2)))%
}

连续嵌入I下标会让欠层式反复套入自身，小舞沿每一对括号指出新增的一层，没有把它看成普通乘三。

最后那组较低阶的例子，让唐三先开了口。

``这里是拿来说明写法的问题？''

霍雨浩点头。讲义把这样的连续套入同非标准形式相比，提出折中：欠层若带来所说的非标准形式，便至多欠一层，随后转回下标嵌套；并没有宣称这一办法已经解决了所有问题。

唐三在页边记下这句话，给后面真正的矩阵分析留出另一行。

\NovelDisplay{M03870R01}{%
\psi_{\psi_{M}(M^3)}(M^3+\psi_{M}(M^3))%
}

\NovelDisplay{M03871R01}{%
\varepsilon_1 = \psi(\Omega \times 2)%
}

\NovelDisplay{M03872R01}{%
\psi(\Omega+1),\psi(\psi(\Omega+1)),\psi(\psi(\psi(\Omega+1))) \cdots%
}

这组低阶对照用于说明欠层可能带来的非标准形式，因此遇到这种情形先至多欠一层，再回到下标嵌套。

回到石阶前，重压并没有因为读过说明就退开。

蔡须坤抬脚时，仍在一个下标里停住。卢本韦扶住她的手肘，让她先撤回来，再把两侧共同的背景指出。她低声读完，发现方才竟把一个新加项带进了背景本身。

薛峰没有替她补答案，只将原矩阵那张纸移到她膝前。她重新看末列，后背的冷汗才慢慢落下去。

\NovelDisplay{M03873R01}{%
(0)(1,1,1)(2,1,1)(3,1,1)(2,1,1)(3,1)(2)%
}

\NovelDisplay{M03874R01}{%
= \psi(I(1,0))%
}

\NovelDisplay{M03875R01}{%
= \psi(M^3)%
}

\NovelDisplay{M03876R01}{%
(0)(1,1,1)(2,1,1)(3,1,1)(2,1,1)(3,1)(2)(1,1,1)(2,1,1)(3,1,1)%
}

\NovelDisplay{M03877R01}{%
= \psi(I(1,0)+I_\omega)%
}

\NovelDisplay{M03878R01}{%
= \psi(M^3+\psi_{M}(M^2\times \omega))%
}

\NovelDisplay{M03879R01}{%
(0)(1,1,1)(2,1,1)(3,1,1)(2,1,1)(3,1)(2)(1,1,1)(2,1,1)(3,1,1)(2,1,1)%
}
\NovelDisplay{M03879R02}{%
(3,1)(2)%
}

\NovelDisplay{M03880R01}{%
= \psi(I(1,0)+\psi_{I(1,0)}(I(1,0)))%
}

\NovelDisplay{M03881R01}{%
= \psi(M^3+\psi_{M}(M^2\times \psi_{\psi_{M}(M^3)}(M^3)))%
}

\NovelDisplay{M03882R01}{%
(0)(1,1,1)(2,1,1)(3,1,1)(2,1,1)(3,1)(2,1)(3,2)%
}

\NovelDisplay{M03883R01}{%
= \psi(I(1,0)+\Omega_{\psi_{I(1,0)}(I(1,0))+1})%
}

\NovelDisplay{M03884R01}{%
= \psi(M^3+\psi_{M}(M^2\times \psi_{\psi_{M}(M^3)}(M^3)+M))%
}

回到原矩阵以后，先固定立方背景，再辨认新加项进入I侧和M侧的具体位置，不由字母相似猜展开。

有一天，叶阳将练习纸带到山脚，发现纸角全被汗浸湿了。

他原以为只有登阶的时候才会难受，坐回屋里，仍不自觉地攥着膝盖。白天那处接不上的地方，像也跟着他一道回来，只要旁边有人翻页，他便忍不住去看自己落下多少。

奥斯卡没有把别人的进度念给他听，先让他将纸晾开。两人坐在门边，看水痕慢慢退去，叶阳才说，自己其实已经找到了问题，却怕一开口，大家又得等。

``你不说，明天他们连该等哪一处都不知道。''

叶阳望着纸，过了一会儿，把原先压在手掌下的那一行露出来。薛峰回来时，他叫住对方，请他先听自己说完。

那次没有人立刻给出答案。叶阳却第一次能把问题说到末尾，心里那份只怕被看见的慌乱，也慢慢松了一点。

\NovelDisplay{M03885R01}{%
(0)(1,1,1)(2,1,1)(3,1,1)(2,1,1)(3,1)(2,1)(3,2)(4,3)%
}

\NovelDisplay{M03886R01}{%
= \psi(I(1,0)+\Omega_{\psi_{I(1,0)}(I(1,0))+2})%
}

\NovelDisplay{M03887R01}{%
= \psi(M^3+\psi_{M}(M^2\times \psi_{\psi_{M}(M^3)}(M^3)+M\times 2))%
}

\NovelDisplay{M03888R01}{%
(0)(1,1,1)(2,1,1)(3,1,1)(2,1,1)(3,1)(2,1)(3,2,1)%
}

\NovelDisplay{M03889R01}{%
= \psi(I(1,0)+\Omega_{\psi_{I(1,0)}(I(1,0))+\omega})%
}

\NovelDisplay{M03890R01}{%
= \psi(M^3+\psi_{M}(M^2\times \psi_{\psi_{M}(M^3)}(M^3)+M\times \omega))%
}

\NovelDisplay{M03891R01}{%
(0)(1,1,1)(2,1,1)(3,1,1)(2,1,1)(3,1)(2,1)(3,2,1)(4,2,1)(5,2)(4)%
}

\NovelDisplay{M03892R01}{%
= \psi(I(1,0)+I_{\psi_{I(1,0)}(I(1,0))+1})%
}

\NovelDisplay{M03893R01}{%
= \psi(M^3+\psi_{M}(M^2\times \psi_{\psi_{M}(M^3)}(M^3)+M^2))%
}

\NovelDisplay{M03894R01}{%
(0)(1,1,1)(2,1,1)(3,1,1)(2,1,1)(3,1)(2,1)(3,2,1)(4,2,1)(5,2,1)%
}

\NovelDisplay{M03895R01}{%
= \psi(I(1,0)+I_{\psi_{I(1,0)}(I(1,0))+\omega})%
}

\NovelDisplay{M03896R01}{%
= \psi(M^3+\psi_{M}(M^2\times \psi_{\psi_{M}(M^3)}(M^3)+M^2\times \omega))%
}

末尾的M倍数与平方处于不同层次，蔡须坤让唐三先指出哪个加号属于外框，才继续比较后继与ω。

午后，诺当将一张纸压在石沿，忽然发现手背被磨出了红痕。

他方才一直用力，想让笔尖在压力下也不抖，纸上的字反而越来越挤。叶阳把自己的垫板递来，让他先放松手，再从能看清的位置重写。

诺当说这样又要慢一会儿。叶阳低头看看自己尚未完成的那张，答说他也还在这里。

两人坐到石桩阴影里。诺当重新写出下标，终于能给每一对括号留下位置；叶阳则请他听一遍自己的问题，确认没有将外层加项念进里面。

蔡须坤下来取水时，看见他们正互相压着纸角，便把备用的笔放在旁边。她没有问今天比昨日多写几行，只等诺当抬头，告诉他收纸以后，大家仍在原来的地方一起吃饭。

诺当应了一声，将那一处写清。吃饭前，他自己把纸收好，没有再攥着发红的手背。

\NovelDisplay{M03897R01}{%
(0)(1,1,1)(2,1,1)(3,1,1)(2,1,1)(3,1)(2,1)(3,2,1)(4,2,1)(5,2,1)%
}
\NovelDisplay{M03897R02}{%
(4,2,1)(5,1)(2)%
}

\NovelDisplay{M03898R01}{%
= \psi(I(1,0)+I_{\psi_{I(1,0)}(I(1,0))\times 2})%
}

\NovelDisplay{M03899R01}{%
= \psi(M^3+\psi_{\psi_{M}(M^3)}(M^3+\psi_{\psi_{M}(M^3)}(M^3)))%
}

\NovelDisplay{M03900R01}{%
(0)(1,1,1)(2,1,1)(3,1,1)(2,1,1)(3,1)(2,1)(3,2,1)(4,2,1)(5,2,1)(4,2,1)%
}
\NovelDisplay{M03900R02}{%
(5,1)(4,2,1)(5,1)(2)%
}

\NovelDisplay{M03901R01}{%
= \psi(I(1,0)+I_{\psi_{I(1,0)}(I(1,0))\times 3})%
}

\NovelDisplay{M03902R01}{%
= \psi(M^3+\psi_{\psi_{M}(M^3)}(M^3+\psi_{\psi_{M}(M^3)}(M^3)\times 2))%
}

\NovelDisplay{M03903R01}{%
(0)(1,1,1)(2,1,1)(3,1,1)(2,1,1)(3,1)(2,1)(3,2,1)(4,2,1)(5,2,1)%
}
\NovelDisplay{M03903R02}{%
(4,2,1)(5,1)(5)%
}

\NovelDisplay{M03904R01}{%
= \psi(I(1,0)+I_{\psi_{I(1,0)}(I(1,0))\times \omega})%
}

\NovelDisplay{M03905R01}{%
= \psi(M^3+\psi_{\psi_{M}(M^3)}(M^3+\psi_{\psi_{M}(M^2\times \psi_{\psi_{M}(M^3)}(M^3)+M)}(1)))%
}

\NovelDisplay{M03906R01}{%
(0)(1,1,1)(2,1,1)(3,1,1)(2,1,1)(3,1)(2,1)(3,2,1)(4,2,1)(5,2,1)%
}
\NovelDisplay{M03906R02}{%
(4,2,1)(5,2)%
}

\NovelDisplay{M03907R01}{%
= \psi(I(1,0)+I_{\Omega_{\psi_{I(1,0)}(I(1,0))}+1})%
}

\NovelDisplay{M03908R01}{%
= \psi(M^3+\psi_{\psi_{M}(M^3)}(M^3+\psi_{M}(M^2\times \psi_{\psi_{M}(M^3)}(M^3)+M)))%
}

有限倍数与ω倍数改变的内部结构并不相同，不能把系数二、三、ω依次塞进同一个空格便算完成分析。

长式铺到这里，奥斯卡先把桌边的碗端走。

小舞说自己还不饿，他便答，先让纸放得下。叶阳将折着的那一截展开，才发现小舞读到的``最后一层''后面，原来还有一段完整的作用对象。

她脸红了，想把纸盖回去。蔡须坤按住另一头，将自己的M侧也摊开。

``我也从这里重新读。''

两人的笔一前一后移动。外面的加号留在外面，中间的ψ一层都没有被藏进省略号里。

\NovelDisplay{M03909R01}{%
(0)(1,1,1)(2,1,1)(3,1,1)(2,1,1)(3,1)(2,1)(3,2,1)(4,2,1)(5,2,1)(4,2,1)%
}
\NovelDisplay{M03909R02}{%
(5,2)(3,2,1)(4,2,1)(5,2)(4)%
}

\NovelDisplay{M03910R01}{%
= \psi(I(1,0)+I_{\psi_{I_{\psi_{I(1,0)}(I(1,0))+1}}(I(1,0)+I_{\psi_{I(1,0)}(I(1,0))+1})})%
}

\NovelDisplay{M03911R01}{%
= \psi(M^3+\psi_{\psi_{M}(M^3)}(M^3+\psi_{M}(M^2\times \psi_{\psi_{M}(M^3)}(M^3)+M%
}
\NovelDisplay{M03911R02}{%
\times \psi_{\psi_{M}(M^2\times \psi_{\psi_{M}(M^3)}(M^3)+M^2)}(M^3+\psi_{M}(M^2\times \psi_{\psi_{M}(M^3)}(M^3)+M^2)))))%
}

\NovelDisplay{M03912R01}{%
(0)(1,1,1)(2,1,1)(3,1,1)(2,1,1)(3,1)(2,1)(3,2,1)(4,2,1)(5,2,1)(4,2,1)%
}
\NovelDisplay{M03912R02}{%
(5,2)(3,2,1)(4,2,1)(5,2)(6,3)%
}

\NovelDisplay{M03913R01}{%
= \psi(I(1,0)+I_{\psi_{I_{\psi_{I(1,0)}(I(1,0))+1}}(I(1,0)+\Omega_{I_{\psi_{I(1,0)}(I(1,0))+1}+1})})%
}

\NovelDisplay{M03914R01}{%
= \psi(M^3+\psi_{\psi_{M}(M^3)}(M^3+\psi_{M}(M^2\times \psi_{\psi_{M}(M^3)}(M^3)+M%
}
\NovelDisplay{M03914R02}{%
\times \psi_{\psi_{M}(M^2\times \psi_{\psi_{M}(M^3)}(M^3)+M^2)}(M^3+\psi_{M}(M^2\times \psi_{\psi_{M}(M^3)}(M^3)+M^2+M)))))%
}

\NovelDisplay{M03915R01}{%
(0)(1,1,1)(2,1,1)(3,1,1)(2,1,1)(3,1)(2,1)(3,2,1)(4,2,1)(5,2,1)(4,2,1)%
}
\NovelDisplay{M03915R02}{%
(5,2)(3,2,1)(4,2,1)(5,2)(6,3,1)(7,3,1)(8,3,1)%
}

\NovelDisplay{M03916R01}{%
= \psi(I(1,0)+I_{\psi_{I_{\psi_{I(1,0)}(I(1,0))+1}}(I(1,0)+I_{\psi_{I(1,0)}(I(1,0))+\omega})})%
}

\NovelDisplay{M03917R01}{%
= \psi(M^3+\psi_{\psi_{M}(M^3)}(M^3+\psi_{M}(M^2\times \psi_{\psi_{M}(M^3)}(M^3)+M%
}
\NovelDisplay{M03917R02}{%
\times \psi_{\psi_{M}(M^2\times \psi_{\psi_{M}(M^3)}(M^3)+M^2)}(M^3+\psi_{M}(M^2\times \psi_{\psi_{M}(M^3)}(M^3)+M^2\times \omega)))))%
}

\NovelDisplay{M03918R01}{%
(0)(1,1,1)(2,1,1)(3,1,1)(2,1,1)(3,1)(2,1)(3,2,1)(4,2,1)(5,2,1)(4,2,1)%
}
\NovelDisplay{M03918R02}{%
(5,2)(3,2,1)(4,2,1)(5,2)(6,3,1)(7,3,1)(8,3,1)(7,3,1)(8,1)(1,1,1)(2,1,1)%
}
\NovelDisplay{M03918R03}{%
(3,1,1)(2,1,1)(3,1)(2)%
}

\NovelDisplay{M03919R01}{%
= \psi(I(1,0)+I_{\psi_{I_{\psi_{I(1,0)}(I(1,0))+1}}(I(1,0)+I_{\psi_{I(1,0)}(I(1,0))\times 2})})%
}

\NovelDisplay{M03920R01}{%
= \psi(M^3+\psi_{\psi_{M}(M^3)}(M^3+\psi_{M}(M^2\times \psi_{\psi_{M}(M^3)}(M^3)+M%
}
\NovelDisplay{M03920R02}{%
\times \psi_{\psi_{M}(M^2\times \psi_{\psi_{M}(M^3)}(M^3)+M^2)}(M^3+\psi_{\psi_{M}(M^3)}(M^3+\psi_{\psi_{M}(M^3)}(M^3))))))%
}

这里的长式保留全部中层ψ与下标，卢本韦先读外层I的下标，蔡须坤再沿M侧逐层找回它的作用对象。

第一回复核，他们把最长的一条暂时同右侧表示分开。

原矩阵有二十六列，末列只留下第一行的二，最高非零位置就在第一行。沿该行的父链回找，坏根落在第二十一列；第二十一至第二十五列合成五列的一段，末列删去后保留二十五列。

蔡须坤让唐三先说范围，自己没有因为曾在另一个例子里见过第一行提升，就给这里也添上一个差值。最高非零行已经是第一行，它前面没有更低编号的行需要那样提升。

三个人分别由同一原矩阵求出第零项、第一项和第二项。

\NovelDisplay{M03921R01}{%
\text{展开第零项：}\quad(0)(1,1,1)(2,1,1)(3,1,1)(2,1,1)(3,1)(2,1)(3,2,1)(4,2,1)%
}
\NovelDisplay{M03921R02}{%
(5,2,1)(4,2,1)(5,2)(3,2,1)(4,2,1)(5,2)(6,3,1)(7,3,1)(8,3,1)(7,3,1)(8,1)%
}
\NovelDisplay{M03921R03}{%
(1,1,1)(2,1,1)(3,1,1)(2,1,1)(3,1)%
}

\NovelDisplay{M03922R01}{%
\text{展开第一项：}\quad(0)(1,1,1)(2,1,1)(3,1,1)(2,1,1)(3,1)(2,1)(3,2,1)(4,2,1)%
}
\NovelDisplay{M03922R02}{%
(5,2,1)(4,2,1)(5,2)(3,2,1)(4,2,1)(5,2)(6,3,1)(7,3,1)(8,3,1)(7,3,1)(8,1)%
}
\NovelDisplay{M03922R03}{%
(1,1,1)(2,1,1)(3,1,1)(2,1,1)(3,1)(1,1,1)(2,1,1)(3,1,1)(2,1,1)(3,1)%
}

\NovelDisplay{M03923R01}{%
\text{展开第二项：}\quad(0)(1,1,1)(2,1,1)(3,1,1)(2,1,1)(3,1)(2,1)(3,2,1)(4,2,1)%
}
\NovelDisplay{M03923R02}{%
(5,2,1)(4,2,1)(5,2)(3,2,1)(4,2,1)(5,2)(6,3,1)(7,3,1)(8,3,1)(7,3,1)(8,1)%
}
\NovelDisplay{M03923R03}{%
(1,1,1)(2,1,1)(3,1,1)(2,1,1)(3,1)(1,1,1)(2,1,1)(3,1,1)(2,1,1)(3,1)%
}
\NovelDisplay{M03923R04}{%
(1,1,1)(2,1,1)(3,1,1)(2,1,1)(3,1)%
}

这三项分别保留二十五列、三十列和三十五列，取项编号控制生成份数，而不是把前一份结果再展开一次。

小舞数到第三十五列，抬头望了唐三一眼。唐三也刚数完，先前被长式逼出的慌乱，终于让位给一个能亲手检查的步骤。

\NovelDisplay{M03924R01}{%
(0)(1,1,1)(2,1,1)(3,1,1)(2,1,1)(3,1)(2,1)(3,2,1)(4,2,1)(5,2,1)(4,2,1)%
}
\NovelDisplay{M03924R02}{%
(5,2)(3,2,1)(4,2,1)(5,2)(6,3,1)(7,3,1)(8,3,1)(7,3,1)(8,2)%
}

\NovelDisplay{M03925R01}{%
= \psi(I(1,0)+I_{\psi_{I_{\psi_{I(1,0)}(I(1,0))+1}}(I(1,0)+I_{\Omega_{\psi_{I(1,0)}(I(1,0))+1}})})%
}

\NovelDisplay{M03926R01}{%
= \psi(M^3+\psi_{\psi_{M}(M^3)}(M^3+\psi_{M}(M^2\times \psi_{\psi_{M}(M^3)}(M^3)+M%
}
\NovelDisplay{M03926R02}{%
\times \psi_{\psi_{M}(M^2\times \psi_{\psi_{M}(M^3)}(M^3)+M^2)}(M^3+\psi_{\psi_{M}(M^3)}(M^3+\psi_{M}(M^2%
}
\NovelDisplay{M03926R03}{%
\times \psi_{\psi_{M}(M^3)}(M^3)+M))))))%
}

\NovelDisplay{M03927R01}{%
(0)(1,1,1)(2,1,1)(3,1,1)(2,1,1)(3,1)(2,1)(3,2,1)(4,2,1)(5,2,1)(4,2,1)%
}
\NovelDisplay{M03927R02}{%
(5,2)(3,2,1)(4,2,1)(5,2)(6,3,1)(7,3,1)(8,3,1)(7,3,1)(8,2)(4)%
}

\NovelDisplay{M03928R01}{%
= \psi(I(1,0)+I_{I_{\psi_{I(1,0)}(I(1,0))+1}})%
}

\NovelDisplay{M03929R01}{%
= \psi(M^3+\psi_{\psi_{M}(M^3)}(M^3+\psi_{M}(M^2\times \psi_{\psi_{M}(M^3)}(M^3)+M^2)))%
}

\NovelDisplay{M03930R01}{%
(0)(1,1,1)(2,1,1)(3,1,1)(2,1,1)(3,1)(2,1)(3,2,1)(4,2,1)(5,2,1)(4,2,1)%
}
\NovelDisplay{M03930R02}{%
(5,2)(3,2,1)(4,2,1)(5,2)(6,3,1)(7,3,1)(8,3,1)(7,3,1)(8,3)%
}

\NovelDisplay{M03931R01}{%
= \psi(I(1,0)+I_{\Omega_{I_{\psi_{I(1,0)}(I(1,0))+1}+1}})%
}

\NovelDisplay{M03932R01}{%
= \psi(M^3+\psi_{\psi_{M}(M^3)}(M^3+\psi_{M}(M^2\times \psi_{\psi_{M}(M^3)}(M^3)+M^2+M)))%
}

\NovelDisplay{M03933R01}{%
(0)(1,1,1)(2,1,1)(3,1,1)(2,1,1)(3,1)(2,1)(3,2,1)(4,2,1)(5,2,1)(4,2,1)%
}
\NovelDisplay{M03933R02}{%
(5,2)(3,2,1)(4,2,1)(5,2,1)%
}

\NovelDisplay{M03934R01}{%
= \psi(I(1,0)+I_{I_{\psi_{I(1,0)}(I(1,0))+\omega}})%
}

\NovelDisplay{M03935R01}{%
= \psi(M^3+\psi_{\psi_{M}(M^3)}(M^3+\psi_{M}(M^2\times \psi_{\psi_{M}(M^3)}(M^3)+M^2 \times \omega)))%
}

新增的M、平方与ω倍数都要放回各自所在的括号，越长的共同前段越不能凭熟悉跳读。

\NovelDisplay{M03936R01}{%
(0)(1,1,1)(2,1,1)(3,1,1)(2,1,1)(3,1)(2,1)(3,2,1)(4,2,1)(5,2,1)(4,2,1)%
}
\NovelDisplay{M03936R02}{%
(5,2)(3,2,1)(4,2,1)(5,2,1)(4,2,1)(5,2)%
}

\NovelDisplay{M03937R01}{%
= \psi(I(1,0)+I_{I_{\Omega_{\psi_{I(1,0)}(I(1,0))}+1}})%
}

\NovelDisplay{M03938R01}{%
= \psi(M^3+\psi_{\psi_{M}(M^3)}(M^3+\psi_{\psi_{M}(M^3)}(M^3+\psi_{M}(M^2\times \psi_{\psi_{M}(M^3)}(M^3)+M))))%
}

\NovelDisplay{M03939R01}{%
(0)(1,1,1)(2,1,1)(3,1,1)(2,1,1)(3,1)(2,1)(3,2,1)(4,2,1)(5,2,1)%
}
\NovelDisplay{M03939R02}{%
(4,2,1)(5,2)(4)%
}

\NovelDisplay{M03940R01}{%
= \psi(I(1,0)\times 2)%
}

\NovelDisplay{M03941R01}{%
= \psi(M^3+\psi_{M}(M^3))%
}

\NovelDisplay{M03942R01}{%
(0)(1,1,1)(2,1,1)(3,1,1)(2,1,1)(3,1)(2,1)(3,2,1)(4,2,1)(5,2,1)(4,2,1)%
}
\NovelDisplay{M03942R02}{%
(5,2)(4,2)(5,3)%
}

\NovelDisplay{M03943R01}{%
= \psi(I(1,0)\times 2+\Omega_{\psi_{I(1,0)}(I(1,0)\times 2)+1})%
}

\NovelDisplay{M03944R01}{%
= \psi(M^3+\psi_{M}(M^3)+\psi_{M}(M^2\times\psi_{\psi_{M}(M^3)}(M^3+\psi_{M}(M^3))+M))%
}

\NovelDisplay{M03945R01}{%
(0)(1,1,1)(2,1,1)(3,1,1)(2,1,1)(3,1)(2,1)(3,2,1)(4,2,1)(5,2,1)(4,2,1)%
}
\NovelDisplay{M03945R02}{%
(5,2)(4,2)(5,3,1)(6,3,1)(7,3)(6)%
}

\NovelDisplay{M03946R01}{%
= \psi(I(1,0)\times 2+I_{\psi_{I(1,0)}(I(1,0)\times 2)+1})%
}

\NovelDisplay{M03947R01}{%
= \psi(M^3+\psi_{M}(M^3)+\psi_{M}(M^2\times\psi_{\psi_{M}(M^3)}(M^3+\psi_{M}(M^3))+M^2))%
}

当I侧转入整个背景的倍数时，M侧留下原背景并另加折叠项，外面的倍数并未变成直接给立方乘二。

同一段用了几天，小舞的焦躁又冒了出来。

她问是不是只有自己总在这里停。古月娜没有指着前面几个人说他们也慢过，只问她这几天究竟是哪一层还接不清。

小舞先说全都乱。等唐三将右侧遮住，她指着原矩阵，从已经确认的前段说起，最后却准确停在了一个位置。

古月娜把空纸推近：``那就从这里问。''

小舞看了看，终于没有再将整页都揉掉。

\NovelDisplay{M03948R01}{%
(0)(1,1,1)(2,1,1)(3,1,1)(2,1,1)(3,1)(2,1)(3,2,1)(4,2,1)(5,2,1)(4,2,1)%
}
\NovelDisplay{M03948R02}{%
(5,2)(4,2)(5,3,1)(6,3,1)(7,3,1)%
}

\NovelDisplay{M03949R01}{%
= \psi(I(1,0)\times 2+I_{\psi_{I(1,0)}(I(1,0)\times 2)+\omega})%
}

\NovelDisplay{M03950R01}{%
= \psi(M^3+\psi_{M}(M^3)+\psi_{M}(M^2\times\psi_{\psi_{M}(M^3)}(M^3+\psi_{M}(M^3))+M^2\times \omega))%
}

\NovelDisplay{M03951R01}{%
(0)(1,1,1)(2,1,1)(3,1,1)(2,1,1)(3,1)(2,1)(3,2,1)(4,2,1)(5,2,1)(4,2,1)%
}
\NovelDisplay{M03951R02}{%
(5,2)(4,2)(5,3,1)(6,3,1)(7,3,1)(6,3,1)(7,3)%
}

\NovelDisplay{M03952R01}{%
= \psi(I(1,0)\times 2+I_{\Omega_{\psi_{I(1,0)}(I(1,0)\times 2)+1}})%
}

\NovelDisplay{M03953R01}{%
= \psi(M^3+\psi_{M}(M^3)+\psi_{\psi_{M}(M^3)}(M^3+\psi_{M}(M^3)+\psi_{M}(M^2\times\psi_{\psi_{M}(M^3)}(M^3%
}
\NovelDisplay{M03953R02}{%
+\psi_{M}(M^3))+M)))%
}

\NovelDisplay{M03954R01}{%
(0)(1,1,1)(2,1,1)(3,1,1)(2,1,1)(3,1)(2,1)(3,2,1)(4,2,1)(5,2,1)(4,2,1)%
}
\NovelDisplay{M03954R02}{%
(5,2)(4,2)(5,3,1)(6,3,1)(7,3,1)(6,3,1)(7,3)(6)%
}

\NovelDisplay{M03955R01}{%
= \psi(I(1,0)\times 3)%
}

\NovelDisplay{M03956R01}{%
= \psi(M^3+\psi_{M}(M^3)\times 2)%
}

\NovelDisplay{M03957R01}{%
(0)(1,1,1)(2,1,1)(3,1,1)(2,1,1)(3,1)(2,1)(3,2,1)(4,2,1)(5,2,1)(4,2,1)%
}
\NovelDisplay{M03957R02}{%
(5,2)(4,2)(5,3,1)(6,3,1)(7,3,1)(6,3,1)(7,3)(6,3)(7,4,1)(8,4,1)(9,4,1)%
}
\NovelDisplay{M03957R03}{%
(8,4,1)(9,4)(8)%
}

\NovelDisplay{M03958R01}{%
= \psi(I(1,0)\times 4)%
}

\NovelDisplay{M03959R01}{%
= \psi(M^3+\psi_{M}(M^3)\times 3)%
}

相邻两式的参数变化和整个I项的有限倍数分别处理，先辨认变化发生在哪里，再比较它对应的矩阵尾部。

蔡须坤后来将这几条留在同一张长页上。

她曾经喜欢把相同的前段一概盖住，觉得这样更容易看出变化；可真到抬脚时，正是那段被自己视作相同的背景，常因后面作用对象的改变，必须重新辨认。

卢本韦听她说完，拿I册同她核了一遍。两人都没有规定从此每次只能怎样读，只将这一次误读的起点留下，提醒自己相同的字面也要放在当前的作用范围里看。

霍雨浩回来，见她们仍在纸前，先问有没有吃饭。蔡须坤想说等会儿，卢本韦已经将笔放到桌上。

``留着明早再看。''

蔡须坤看她一眼，忽然也笑了。长页被压在石头底下，谁都没有为了赶饭将未完的那层胡乱接上。

第二天，她们各自从原式重新读起，头一天能说清的部分仍在，空着的地方也没有因此消失；这一次，她们能更早指出那个确切的空处。

\NovelDisplay{M03960R01}{%
(0)(1,1,1)(2,1,1)(3,1,1)(2,1,1)(3,1)(2,1,1)%
}

\NovelDisplay{M03961R01}{%
= \psi(I(1,0)\times \omega)%
}

\NovelDisplay{M03962R01}{%
= \psi(M^3+\psi_{\psi_{M}(M^3+M)}(1))%
}

\NovelDisplay{M03963R01}{%
(0)(1,1,1)(2,1,1)(3,1,1)(2,1,1)(3,1)(2,1,1)(2,1)(3,2,1)(4,2,1)(5,2,1)%
}
\NovelDisplay{M03963R02}{%
(4,2,1)(5,2)(4)%
}

\NovelDisplay{M03964R01}{%
= \psi(I(1,0)\times \omega+I(1,0))%
}

\NovelDisplay{M03965R01}{%
= \psi(M^3+\psi_{\psi_{M}(M^3+M)}(1)+\psi_{M}(M^3))%
}

\NovelDisplay{M03966R01}{%
(0)(1,1,1)(2,1,1)(3,1,1)(2,1,1)(3,1)(2,1,1)(2,1)(3,2,1)(4,2,1)(5,2,1)%
}
\NovelDisplay{M03966R02}{%
(4,2,1)(5,2)(4,2)(5,3,1)(6,3,1)(7,3,1)(6,3,1)(7,3)(6)%
}

\NovelDisplay{M03967R01}{%
= \psi(I(1,0)\times \omega+I(1,0)\times 2)%
}

\NovelDisplay{M03968R01}{%
= \psi(M^3+\psi_{\psi_{M}(M^3+M)}(1)+\psi_{M}(M^3)\times 2)%
}

\NovelDisplay{M03969R01}{%
(0)(1,1,1)(2,1,1)(3,1,1)(2,1,1)(3,1)(2,1,1)(2,1)(3,2,1)(4,2,1)(5,2,1)%
}
\NovelDisplay{M03969R02}{%
(4,2,1)(5,2)(4,2,1)%
}

\NovelDisplay{M03970R01}{%
= \psi(I(1,0)\times \omega \times 2)%
}

\NovelDisplay{M03971R01}{%
= \psi(M^3+\psi_{\psi_{M}(M^3+M)}(1) \times 2)%
}

ω倍数与ω倍数之后再加一份背景的位置不同，霍雨浩请每个人守住加法次序，避免把几条式子合成一种写法。

\NovelDisplay{M03972R01}{%
(0)(1,1,1)(2,1,1)(3,1,1)(2,1,1)(3,1)(2,1,1)(2,1,1)%
}

\NovelDisplay{M03973R01}{%
= \psi(I(1,0)\times \psi(2))%
}

\NovelDisplay{M03974R01}{%
= \psi(M^3+\psi_{\psi_{M}(M^3+M)}(2))%
}

\NovelDisplay{M03975R01}{%
(0)(1,1,1)(2,1,1)(3,1,1)(2,1,1)(3,1)(2,1,1)(3,1)%
}

\NovelDisplay{M03976R01}{%
= \psi(I(1,0)\times \Omega)%
}

\NovelDisplay{M03977R01}{%
= \psi(M^3+\psi_{\psi_{M}(M^3+M)}(M))%
}

\NovelDisplay{M03978R01}{%
(0)(1,1,1)(2,1,1)(3,1,1)(2,1,1)(3,1)(2,1,1)(3,1)(1,1,1)(2,1,1)(3,1,1)%
}

\NovelDisplay{M03979R01}{%
= \psi(I(1,0)\times I_\omega)%
}

\NovelDisplay{M03980R01}{%
= \psi(M^3+\psi_{\psi_{M}(M^3+M)}(M^2\times \omega))%
}

\NovelDisplay{M03981R01}{%
(0)(1,1,1)(2,1,1)(3,1,1)(2,1,1)(3,1)(2,1,1)(3,1)(1,1,1)(2,1,1)(3,1,1)%
}
\NovelDisplay{M03981R02}{%
(2,1,1)(3,1)(2)%
}

\NovelDisplay{M03982R01}{%
= \psi(I(1,0)\times \psi_{I(1,0)}(I(1,0)))%
}

\NovelDisplay{M03983R01}{%
= \psi(M^3+\psi_{\psi_{M}(M^3+M)}(M^2\times \psi_{\psi_{M}(M^3)}(M^3)))%
}

乘数逐渐变复杂时，负责折叠它的下标仍须完整保留，纸面上少写一层不能算成更好的简化。

那天从山上下来，蔡须坤仍把手放在胸口。

那里没有新的暴走，只是压力退开之后，疲倦才一点点显出来。卢本韦将M册从她手里接过，问她今天还想看哪一处。

蔡须坤指向明早要重做的原式，没有说还要再练几遍。霍雨浩将温水递来，坐到旁边，先听她把已经明白的部分讲完。

果糕领域只在她身边铺开很薄的一圈，让呼吸慢慢平下来。纸上的下标仍得由她自己认清。

\NovelDisplay{M03984R01}{%
(0)(1,1,1)(2,1,1)(3,1,1)(2,1,1)(3,1)(2,1,1)(3,1)(1,1,1)(2,1,1)(3,1,1)%
}
\NovelDisplay{M03984R02}{%
(2,1,1)(3,1)(2,1,1)%
}

\NovelDisplay{M03985R01}{%
= \psi(I(1,0)\times \psi_{I(1,0)}(I(1,0)\times \omega))%
}

\NovelDisplay{M03986R01}{%
= \psi(M^3+\psi_{\psi_{M}(M^3+M)}(M^2\times \psi_{\psi_{M}(M^3)}(M^3+\psi_{\psi_{M}(M^3+M)}(1))))%
}

\NovelDisplay{M03987R01}{%
(0)(1,1,1)(2,1,1)(3,1,1)(2,1,1)(3,1)(2,1,1)(3,1)(1,1,1)(2,1,1)(3,1,1)%
}
\NovelDisplay{M03987R02}{%
(2,1,1)(3,1)(2,1,1)(3,1)%
}

\NovelDisplay{M03988R01}{%
= \psi(I(1,0)\times \psi_{I(1,0)}(I(1,0)\times \Omega))%
}

\NovelDisplay{M03989R01}{%
= \psi(M^3+\psi_{\psi_{M}(M^3+M)}(M^2\times \psi_{\psi_{M}(M^3)}(M^3+\psi_{\psi_{M}(M^3+M)}(M))))%
}

\NovelDisplay{M03990R01}{%
(0)(1,1,1)(2,1,1)(3,1,1)(2,1,1)(3,1)(2,1,1)(3,1)(2)%
}

\NovelDisplay{M03991R01}{%
= \psi({I(1,0)}^2)%
}

\NovelDisplay{M03992R01}{%
= \psi(M^3+\psi_{\psi_{M}(M^3+M)}(M^3))%
}

\NovelDisplay{M03993R01}{%
(0)(1,1,1)(2,1,1)(3,1,1)(2,1,1)(3,1)(2,1,1)(3,1)(2,1)(3,2)%
}

\NovelDisplay{M03994R01}{%
= \psi({I(1,0)}^2+\Omega_{\psi_{I(1,0)}({I(1,0)}^2)+1})%
}

\NovelDisplay{M03995R01}{%
= \psi(M^3+\psi_{\psi_{M}(M^3+M)}(M^3)+\psi_{M}(M^2 \times \psi_{\psi_{M}(M^3)}(M^3%
}
\NovelDisplay{M03995R02}{%
+\psi_{\psi_{M}(M^3+M)}(M^3)) + M))%
}

平方本身与平方后继续添项属于相邻的不同步骤，薛峰把共同背景和后来变化分别指回原矩阵。

\NovelDisplay{M03996R01}{%
(0)(1,1,1)(2,1,1)(3,1,1)(2,1,1)(3,1)(2,1,1)(3,1)(2,1)(3,2,1)(4,2,1)%
}
\NovelDisplay{M03996R02}{%
(5,2)(4)%
}

\NovelDisplay{M03997R01}{%
= \psi({I(1,0)}^2+I_{\psi_{I(1,0)}({I(1,0)}^2)+1})%
}

\NovelDisplay{M03998R01}{%
= \psi(M^3+\psi_{\psi_{M}(M^3+M)}(M^3)+\psi_{M}(M^2 \times \psi_{\psi_{M}(M^3)}(M^3%
}
\NovelDisplay{M03998R02}{%
+\psi_{\psi_{M}(M^3+M)}(M^3)) + M^2))%
}

\NovelDisplay{M03999R01}{%
(0)(1,1,1)(2,1,1)(3,1,1)(2,1,1)(3,1)(2,1,1)(3,1)(2,1)(3,2,1)%
}
\NovelDisplay{M03999R02}{%
(4,2,1)(5,2,1)%
}

\NovelDisplay{M04000R01}{%
= \psi({I(1,0)}^2+I_{\psi_{I(1,0)}({I(1,0)}^2)+\omega})%
}

\NovelDisplay{M04001R01}{%
= \psi(M^3+\psi_{\psi_{M}(M^3+M)}(M^3)+\psi_{M}(M^2 \times \psi_{\psi_{M}(M^3)}(M^3%
}
\NovelDisplay{M04001R02}{%
+\psi_{\psi_{M}(M^3+M)}(M^3)) + M^2\times \omega))%
}

\NovelDisplay{M04002R01}{%
(0)(1,1,1)(2,1,1)(3,1,1)(2,1,1)(3,1)(2,1,1)(3,1)(2,1)(3,2,1)(4,2,1)%
}
\NovelDisplay{M04002R02}{%
(5,2,1)(4,2,1)(5,2)%
}

\NovelDisplay{M04003R01}{%
= \psi({I(1,0)}^2+I_{\Omega_{\psi_{I(1,0)}({I(1,0)}^2)+1}})%
}

\NovelDisplay{M04004R01}{%
= \psi(M^3+\psi_{\psi_{M}(M^3+M)}(M^3)+\psi_{\psi_{M}(M^3)}(M^3+\psi_{\psi_{M}(M^3+M)}(M^3)+\psi_{M}(M^2%
}
\NovelDisplay{M04004R02}{%
\times \psi_{\psi_{M}(M^3)}(M^3+\psi_{\psi_{M}(M^3+M)}(M^3)) + M)))%
}

\NovelDisplay{M04005R01}{%
(0)(1,1,1)(2,1,1)(3,1,1)(2,1,1)(3,1)(2,1,1)(3,1)(2,1)(3,2,1)(4,2,1)%
}
\NovelDisplay{M04005R02}{%
(5,2,1)(4,2,1)(5,2)(4)%
}

\NovelDisplay{M04006R01}{%
= \psi({I(1,0)}^2+I(1,0))%
}

\NovelDisplay{M04007R01}{%
= \psi(M^3+\psi_{\psi_{M}(M^3+M)}(M^3)+\psi_{M}(M^3))%
}

两侧嵌套的作用对象同时变长，先从最外层确认背景，才能看清末尾加项是在内层还是外层。

下一条的末列是\allowbreak{}\NovelInline{I06089}{(5,2)}\allowbreak{}，最高非零位置改在第二行。卢本韦指出第九列为坏根，霍雨浩请小舞从父链重新确认；第九至第十三列共五列，第一行的提升差为三，哪些位置实际受提升仍要回查祖先关系。

\NovelDisplay{M04008R01}{%
\text{展开第零项：}\quad(0)(1,1,1)(2,1,1)(3,1,1)(2,1,1)(3,1)(2,1,1)(3,1)(2,1)%
}
\NovelDisplay{M04008R02}{%
(3,2,1)(4,2,1)(5,2,1)(4,2,1)%
}

\NovelDisplay{M04009R01}{%
\text{展开第一项：}\quad(0)(1,1,1)(2,1,1)(3,1,1)(2,1,1)(3,1)(2,1,1)(3,1)(2,1)%
}
\NovelDisplay{M04009R02}{%
(3,2,1)(4,2,1)(5,2,1)(4,2,1)(5,1)(6,2,1)(7,2,1)(8,2,1)(7,2,1)%
}

\NovelDisplay{M04010R01}{%
\text{展开第二项：}\quad(0)(1,1,1)(2,1,1)(3,1,1)(2,1,1)(3,1)(2,1,1)(3,1)(2,1)%
}
\NovelDisplay{M04010R02}{%
(3,2,1)(4,2,1)(5,2,1)(4,2,1)(5,1)(6,2,1)(7,2,1)(8,2,1)(7,2,1)(8,1)%
}
\NovelDisplay{M04010R03}{%
(9,2,1)(10,2,1)(11,2,1)(10,2,1)%
}

三项分别有十三列、十八列和二十三列，坏部范围与提升条件都来自本条原矩阵，不能直接搬用上一条的处理。

蔡须坤将两条原式并排留好。纸上同样有一段五列，却已经不能只凭``五列''两个字说它们的生成过程相同。

\NovelDisplay{M04011R01}{%
(0)(1,1,1)(2,1,1)(3,1,1)(2,1,1)(3,1)(2,1,1)(3,1)(2,1)(3,2,1)(4,2,1)%
}
\NovelDisplay{M04011R02}{%
(5,2,1)(4,2,1)(5,2)(4,2,1)%
}

\NovelDisplay{M04012R01}{%
= \psi({I(1,0)}^2+I(1,0)\times \omega)%
}

\NovelDisplay{M04013R01}{%
= \psi(M^3+\psi_{\psi_{M}(M^3+M)}(M^3)+\psi_{\psi_{M}(M^3+M)}(1))%
}

\NovelDisplay{M04014R01}{%
(0)(1,1,1)(2,1,1)(3,1,1)(2,1,1)(3,1)(2,1,1)(3,1)(2,1)(3,2,1)(4,2,1)%
}
\NovelDisplay{M04014R02}{%
(5,2,1)(4,2,1)(5,2)(4,2,1)(5,2)%
}

\NovelDisplay{M04015R01}{%
= \psi({I(1,0)}^2+I(1,0)\times \Omega_{\psi_{I(1,0)}({I(1,0)}^2)+1})%
}

\NovelDisplay{M04016R01}{%
= \psi(M^3+\psi_{\psi_{M}(M^3+M)}(M^3)+\psi_{\psi_{M}(M^3+M)}(M^2\times \psi_{\psi_{M}(M^3)}(M^3%
}
\NovelDisplay{M04016R02}{%
+\psi_{\psi_{M}(M^3+M)}(M^3))+M))%
}

\NovelDisplay{M04017R01}{%
(0)(1,1,1)(2,1,1)(3,1,1)(2,1,1)(3,1)(2,1,1)(3,1)(2,1)(3,2,1)(4,2,1)%
}
\NovelDisplay{M04017R02}{%
(5,2,1)(4,2,1)(5,2)(4,2,1)(5,2)(3,2,1)(4,2,1)(5,2,1)%
}

\NovelDisplay{M04018R01}{%
= \psi({I(1,0)}^2+I(1,0)\times I_{\psi_{I(1,0)}({I(1,0)}^2)+\omega})%
}

\NovelDisplay{M04019R01}{%
= \psi(M^3+\psi_{\psi_{M}(M^3+M)}(M^3)+\psi_{\psi_{M}(M^3+M)}(M^2\times \psi_{\psi_{M}(M^3)}(M^3%
}
\NovelDisplay{M04019R02}{%
+\psi_{\psi_{M}(M^3+M)}(M^3))+M^2\times \omega))%
}

\NovelDisplay{M04020R01}{%
(0)(1,1,1)(2,1,1)(3,1,1)(2,1,1)(3,1)(2,1,1)(3,1)(2,1)(3,2,1)(4,2,1)%
}
\NovelDisplay{M04020R02}{%
(5,2,1)(4,2,1)(5,2)(4,2,1)(5,2)(3,2,1)(4,2,1)(5,2,1)(4,2,1)(5,2)%
}

\NovelDisplay{M04021R01}{%
= \psi({I(1,0)}^2+I(1,0)\times I_{\Omega_{\psi_{I(1,0)}({I(1,0)}^2)+1}})%
}

\NovelDisplay{M04022R01}{%
= \psi(M^3+\psi_{\psi_{M}(M^3+M)}(M^3)+\psi_{\psi_{M}(M^3+M)}(M^2\times \psi_{M}(M^2\times\psi_{\psi_{M}(M^3)}(M^3%
}
\NovelDisplay{M04022R02}{%
+\psi_{\psi_{M}(M^3+M)}(M^3))+M)))%
}

这里改变的是平方后那一项的乘数，不能把它误接进平方的指数，也不能交换左右加项的次序。

接近这次练习的末尾，诺当忽然叫住了薛峰。

他先指了乘数里的ψ，再指指数，问自己是不是又把两处认作同一层。薛峰刚要顺口回答，见他还没说完，便压住纸角等了一会儿。

诺当终于把完整的问题说出。卢本韦顺着I侧指出括号，蔡须坤在M侧找到对应位置，霍雨浩则将原矩阵交回诺当手里。

``下面这三项，你也试着自己做。''

\NovelDisplay{M04023R01}{%
(0)(1,1,1)(2,1,1)(3,1,1)(2,1,1)(3,1)(2,1,1)(3,1)(2,1)(3,2,1)(4,2,1)%
}
\NovelDisplay{M04023R02}{%
(5,2,1)(4,2,1)(5,2)(4,2,1)(5,2)(3,2,1)(4,2,1)(5,2,1)(4,2,1)(5,2)(4)%
}

\NovelDisplay{M04024R01}{%
= \psi({I(1,0)}^2+I(1,0)\times \psi_{I(1,0)}({I(1,0)}^2+I(1,0)))%
}

\NovelDisplay{M04025R01}{%
= \psi(M^3+\psi_{\psi_{M}(M^3+M)}(M^3)+\psi_{\psi_{M}(M^3+M)}(M^2\times \psi_{\psi_{M}(M^3)}(M^3%
}
\NovelDisplay{M04025R02}{%
+\psi_{\psi_{M}(M^3+M)}(M^3)+\psi_{M}(M^3))))%
}

\NovelDisplay{M04026R01}{%
(0)(1,1,1)(2,1,1)(3,1,1)(2,1,1)(3,1)(2,1,1)(3,1)(2,1)(3,2,1)(4,2,1)%
}
\NovelDisplay{M04026R02}{%
(5,2,1)(4,2,1)(5,2)(4,2,1)(5,2)(3,2,1)(4,2,1)(5,2,1)(4,2,1)(5,2)(4,2,1)%
}

\NovelDisplay{M04027R01}{%
= \psi({I(1,0)}^2+I(1,0)\times \psi_{I(1,0)}({I(1,0)}^2+I(1,0)\times \omega))%
}

\NovelDisplay{M04028R01}{%
= \psi(M^3+\psi_{\psi_{M}(M^3+M)}(M^3)+\psi_{\psi_{M}(M^3+M)}(M^2\times \psi_{\psi_{M}(M^3)}(M^3%
}
\NovelDisplay{M04028R02}{%
+\psi_{\psi_{M}(M^3+M)}(M^3)+\psi_{\psi_{M}(M^3+M)}(1))))%
}

\NovelDisplay{M04029R01}{%
(0)(1,1,1)(2,1,1)(3,1,1)(2,1,1)(3,1)(2,1,1)(3,1)(2,1)(3,2,1)(4,2,1)%
}
\NovelDisplay{M04029R02}{%
(5,2,1)(4,2,1)(5,2)(4,2,1)(5,2)(3,2,1)(4,2,1)(5,2,1)(4,2,1)(5,2)%
}
\NovelDisplay{M04029R03}{%
(4,2,1)(5,2)%
}

\NovelDisplay{M04030R01}{%
= \psi({I(1,0)}^2+I(1,0)\times \psi_{I(1,0)}({I(1,0)}^2+I(1,0)\times \Omega_{\psi_{I(1,0)}({I(1,0)}^2)+1}))%
}

\NovelDisplay{M04031R01}{%
= \psi(M^3+\psi_{\psi_{M}(M^3+M)}(M^3)+\psi_{\psi_{M}(M^3+M)}(M^2\times \psi_{\psi_{M}(M^3)}(M^3%
}
\NovelDisplay{M04031R02}{%
+\psi_{\psi_{M}(M^3+M)}(M^3)+\psi_{\psi_{M}(M^3+M)}(M^2\times\psi_{\psi_{M}(M^3)}(M^3%
}
\NovelDisplay{M04031R03}{%
+\psi_{\psi_{M}(M^3+M)}(M^3))+M))))%
}

乘数内部再出现ψ时，卢本韦逐层读清下标与参数，蔡须坤沿M侧指出相同位置，最后仍由原矩阵独立求项。

最后这一条仍以\allowbreak{}\NovelInline{I06090}{(5,2)}\allowbreak{}收尾，坏根仍在第九列，待生成的段却已经长到十四列；诺当先把第九至第二十二列圈全，再由同伴分别核三次独立取项。

\NovelDisplay{M04032R01}{%
\text{展开第零项：}\quad(0)(1,1,1)(2,1,1)(3,1,1)(2,1,1)(3,1)(2,1,1)(3,1)(2,1)%
}
\NovelDisplay{M04032R02}{%
(3,2,1)(4,2,1)(5,2,1)(4,2,1)(5,2)(4,2,1)(5,2)(3,2,1)(4,2,1)(5,2,1)%
}
\NovelDisplay{M04032R03}{%
(4,2,1)(5,2)(4,2,1)%
}

\NovelDisplay{M04033R01}{%
\text{展开第一项：}\quad(0)(1,1,1)(2,1,1)(3,1,1)(2,1,1)(3,1)(2,1,1)(3,1)(2,1)%
}
\NovelDisplay{M04033R02}{%
(3,2,1)(4,2,1)(5,2,1)(4,2,1)(5,2)(4,2,1)(5,2)(3,2,1)(4,2,1)(5,2,1)%
}
\NovelDisplay{M04033R03}{%
(4,2,1)(5,2)(4,2,1)(5,1)(6,2,1)(7,2,1)(8,2,1)(7,2,1)(8,2)(7,2,1)(8,2)%
}
\NovelDisplay{M04033R04}{%
(6,2,1)(7,2,1)(8,2,1)(7,2,1)(8,2)(7,2,1)%
}

\NovelDisplay{M04034R01}{%
\text{展开第二项：}\quad(0)(1,1,1)(2,1,1)(3,1,1)(2,1,1)(3,1)(2,1,1)(3,1)(2,1)%
}
\NovelDisplay{M04034R02}{%
(3,2,1)(4,2,1)(5,2,1)(4,2,1)(5,2)(4,2,1)(5,2)(3,2,1)(4,2,1)(5,2,1)%
}
\NovelDisplay{M04034R03}{%
(4,2,1)(5,2)(4,2,1)(5,1)(6,2,1)(7,2,1)(8,2,1)(7,2,1)(8,2)(7,2,1)(8,2)%
}
\NovelDisplay{M04034R04}{%
(6,2,1)(7,2,1)(8,2,1)(7,2,1)(8,2)(7,2,1)(8,1)(9,2,1)(10,2,1)(11,2,1)%
}
\NovelDisplay{M04034R05}{%
(10,2,1)(11,2)(10,2,1)(11,2)(9,2,1)(10,2,1)(11,2,1)(10,2,1)%
}
\NovelDisplay{M04034R06}{%
(11,2)(10,2,1)%
}

第零项二十二列、第一项三十六列、第二项五十列，长的是生成段，原式与展开项的比较仍要看最先不同的位置。

诺当的第二项写了很久。奥斯卡替他换过一次水，没有替他添列；等最后一列由他自己落下，薛峰才把原纸重新摆到旁边，一列列陪他数完。

第四个月将尽，他们站上了三百八十米。

只有十米。小舞摸着刻痕，第一句话已经到了嘴边，又慢慢咽回去。她想起那条终于能由自己求完的长矩阵，低头看了看手里的本子。

蔡须坤也在看。最前面几页留下了欠层写法的条件与限度，后面才是他们一条条核过的分析；两者没有被混成一套见什么都能套进去的捷径。

卢本韦把I册放在她旁边。她们的字迹仍很不同，同一处下标却终于能从两边相互指出。

霍雨浩记下日期，薛峰等最后的叶阳站稳，才将水壶交过去。五百米还在上面，第一考也仍未结束。

蔡须坤将纸收进封套，轻轻压平。

``明天接着来。''

\chapter{第一百二十三章：容许点进位}\label{ux7b2cux4e00ux767eux4e8cux5341ux4e09ux7ae0ux5bb9ux8bb8ux70b9ux8fdbux4f4d}

第五个月，三百八十米上方那几级，成了他们每天最熟悉的地方。

唐三能说出自己昨天停在哪里，小舞也知道今天要先问哪一个括号。可知道问题的位置，与亲自将它接清，中间仍隔着几次抬不起脚的清晨。

卢本韦在I册后面缝了新页。双参数的写法让她比从前更能把握结构，她说得越来越清楚，纸面上那些逗号却也成了蔡须坤容易忽略的地方。

``这里先别翻。''蔡须坤忽然道。

卢本韦停下，将本子转给她。蔡须坤指着一个参数，坦白自己方才只听懂了M侧，没跟上I侧究竟改了哪里。

薛峰坐近一些。霍雨浩把原矩阵摆在两本册子之间，十一人重新各取了一页纸。

从平方接向更高幂时，唐三又想快一点。

他认出共同前段，便准备直接报出后面的指数。小舞伸手指住幂外的ω，问他刚才那句究竟念的是乘ω，还是让ω进入指数。

唐三怔了一下，重新看纸。这回轮到妹妹等他，他低头将整个指数框画全，才把方才略过的位置重新念出。

\NovelDisplay{M04035R01}{%
(0)(1,1,1)(2,1,1)(3,1,1)(2,1,1)(3,1)(2,1,1)(3,1)(2,1)(3,2,1)(4,2,1)%
}
\NovelDisplay{M04035R02}{%
(5,2,1)(4,2,1)(5,2)(4,2,1)(5,2)(4)%
}

\NovelDisplay{M04036R01}{%
= \psi({I(1,0)}^2\times 2)%
}

\NovelDisplay{M04037R01}{%
= \psi(M^3+\psi_{\psi_{M}(M^3+M)}(M^3) \times 2)%
}

\NovelDisplay{M04038R01}{%
(0)(1,1,1)(2,1,1)(3,1,1)(2,1,1)(3,1)(2,1,1)(3,1)(2,1,1)%
}

\NovelDisplay{M04039R01}{%
= \psi({I(1,0)}^2\times \omega)%
}

\NovelDisplay{M04040R01}{%
= \psi(M^3+\psi_{\psi_{M}(M^3+M)}(M^3+1))%
}

\NovelDisplay{M04041R01}{%
(0)(1,1,1)(2,1,1)(3,1,1)(2,1,1)(3,1)(2,1,1)(3,1)(2,1,1)(3,1)(2)%
}

\NovelDisplay{M04042R01}{%
= \psi({I(1,0)}^3)%
}

\NovelDisplay{M04043R01}{%
= \psi(M^3+\psi_{\psi_{M}(M^3+M)}(M^3+\psi_{M}(M^3)))%
}

\NovelDisplay{M04044R01}{%
(0)(1,1,1)(2,1,1)(3,1,1)(2,1,1)(3,1)(2,1,1)(3,1)(2,1,1)(3,1)(2,1,1)(3,1)%
}

\NovelDisplay{M04045R01}{%
= \psi({I(1,0)}^4)%
}

\NovelDisplay{M04046R01}{%
= \psi(M^3+\psi_{\psi_{M}(M^3+M)}(M^3+\psi_{M}(M^3)\times 2))%
}

平方外乘ω、立方和四次幂在这里并列，读式时必须把指数与幂外乘数分开。

\NovelDisplay{M04047R01}{%
(0)(1,1,1)(2,1,1)(3,1,1)(2,1,1)(3,1)(3)%
}

\NovelDisplay{M04048R01}{%
= \psi({I(1,0)}^\omega)%
}

\NovelDisplay{M04049R01}{%
= \psi(M^3+\psi_{\psi_{M}(M^3+M)}(M^3+\psi_{\psi_{M}(M^3+M)}(1)))%
}

\NovelDisplay{M04050R01}{%
(0)(1,1,1)(2,1,1)(3,1,1)(2,1,1)(3,1)(3,1)%
}

\NovelDisplay{M04051R01}{%
= \psi({I(1,0)}^\Omega)%
}

\NovelDisplay{M04052R01}{%
= \psi(M^3+\psi_{\psi_{M}(M^3+M)}(M^3+\psi_{\psi_{M}(M^3+M)}(M)))%
}

\NovelDisplay{M04053R01}{%
(0)(1,1,1)(2,1,1)(3,1,1)(2,1,1)(3,1)(3,1)(1,1,1)(2,1,1)(3,1,1)%
}
\NovelDisplay{M04053R02}{%
(2,1,1)(3,1)(2)%
}

\NovelDisplay{M04054R01}{%
= \psi({I(1,0)}^{\psi_{I(1,0)}(I(1,0))})%
}

\NovelDisplay{M04055R01}{%
= \psi(M^3+\psi_{\psi_{M}(M^3+M)}(M^3+\psi_{\psi_{M}(M^3+M)}(M^2\times \psi_{\psi_{M}(M^3)}(M^3))))%
}

\NovelDisplay{M04056R01}{%
(0)(1,1,1)(2,1,1)(3,1,1)(2,1,1)(3,1)(3,1)(1,1,1)(2,1,1)(3,1,1)(2,1,1)%
}
\NovelDisplay{M04056R02}{%
(3,1)(2,1,1)(3,1)(2)%
}

\NovelDisplay{M04057R01}{%
= \psi({I(1,0)}^{\psi_{I(1,0)}({I(1,0)}^2)})%
}

\NovelDisplay{M04058R01}{%
= \psi(M^3+\psi_{\psi_{M}(M^3+M)}(M^3+\psi_{\psi_{M}(M^3+M)}(M^2\times \psi_{\psi_{M}(M^3)}(M^3+\psi_{M}(M^3)))))%
}

ω、Ω以及带ψ的对象依次进入指数，整条M侧的外框仍需保留，不能只留下最末那个字母。

\NovelDisplay{M04059R01}{%
(0)(1,1,1)(2,1,1)(3,1,1)(2,1,1)(3,1)(3,1)(1,1,1)(2,1,1)(3,1,1)%
}
\NovelDisplay{M04059R02}{%
(2,1,1)(3,1)(3)%
}

\NovelDisplay{M04060R01}{%
= \psi({I(1,0)}^{\psi_{I(1,0)}({I(1,0)}^\omega)})%
}

\NovelDisplay{M04061R01}{%
= \psi(M^3+\psi_{\psi_{M}(M^3+M)}(M^3+\psi_{\psi_{M}(M^3+M)}(M^2\times \psi_{\psi_{M}(M^3)}(M^3%
}
\NovelDisplay{M04061R02}{%
+\psi_{\psi_{M}(M^3+M)}(1)))))%
}

\NovelDisplay{M04062R01}{%
(0)(1,1,1)(2,1,1)(3,1,1)(2,1,1)(3,1)(3,1)(1,1,1)(2,1,1)(3,1,1)%
}
\NovelDisplay{M04062R02}{%
(2,1,1)(3,1)(3,1)%
}

\NovelDisplay{M04063R01}{%
= \psi({I(1,0)}^{\psi_{I(1,0)}({I(1,0)}^\Omega)})%
}

\NovelDisplay{M04064R01}{%
= \psi(M^3+\psi_{\psi_{M}(M^3+M)}(M^3+\psi_{\psi_{M}(M^3+M)}(M^2\times \psi_{\psi_{M}(M^3)}(M^3%
}
\NovelDisplay{M04064R02}{%
+\psi_{\psi_{M}(M^3+M)}(M)))))%
}

\NovelDisplay{M04065R01}{%
(0)(1,1,1)(2,1,1)(3,1,1)(2,1,1)(3,1)(3,1)(2)%
}

\NovelDisplay{M04066R01}{%
= \psi({I(1,0)}^{I(1,0)})%
}

\NovelDisplay{M04067R01}{%
= \psi(M^3+\psi_{\psi_{M}(M^3+M)}(M^3+\psi_{\psi_{M}(M^3+M)}(M^3)))%
}

\NovelDisplay{M04068R01}{%
(0)(1,1,1)(2,1,1)(3,1,1)(2,1,1)(3,1)(3,1)(2,1,1)(3,1)(2)%
}

\NovelDisplay{M04069R01}{%
= \psi({I(1,0)}^{I(1,0)+1})%
}

\NovelDisplay{M04070R01}{%
= \psi(M^3+\psi_{\psi_{M}(M^3+M)}(M^3+\psi_{\psi_{M}(M^3+M)}(M^3)+\psi_{M}(M^3)))%
}

指数里的ψ项与指数后继分属不同括号，古月娜让小舞指出最外面的幂以后，才向内读。

奥斯卡听到一半，发现蔡兴坤一直在旁边画圈。

圈越来越多，彼此叠到一起，蔡兴坤却说自己是在找规律。叶阳让他只留下同一层的圈，他看了许久，竟舍不得擦去任何一个。

卢本韦换了一张干净的纸，请他先读最外面的结构。第一层说清以后，再往里面添下一层；那片原先乱成一团的墨圈，终于变成了能够逐个指出的位置。

\NovelDisplay{M04071R01}{%
(0)(1,1,1)(2,1,1)(3,1,1)(2,1,1)(3,1)(3,1)(2,1,1)(3,1)(3,1)(2)%
}

\NovelDisplay{M04072R01}{%
= \psi({I(1,0)}^{I(1,0)\times 2})%
}

\NovelDisplay{M04073R01}{%
= \psi(M^3+\psi_{\psi_{M}(M^3+M)}(M^3+\psi_{\psi_{M}(M^3+M)}(M^3)\times 2))%
}

\NovelDisplay{M04074R01}{%
(0)(1,1,1)(2,1,1)(3,1,1)(2,1,1)(3,1)(3,1)(3)%
}

\NovelDisplay{M04075R01}{%
= \psi({I(1,0)}^{I(1,0)\times \omega})%
}

\NovelDisplay{M04076R01}{%
= \psi(M^3+\psi_{\psi_{M}(M^3+M)}(M^3+\psi_{\psi_{M}(M^3+M)}(M^3+1)))%
}

\NovelDisplay{M04077R01}{%
(0)(1,1,1)(2,1,1)(3,1,1)(2,1,1)(3,1)(3,1)(3,1)(2)%
}

\NovelDisplay{M04078R01}{%
= \psi({I(1,0)}^{{I(1,0)}^2})%
}

\NovelDisplay{M04079R01}{%
= \psi(M^3+\psi_{\psi_{M}(M^3+M)}(M^3+\psi_{\psi_{M}(M^3+M)}(M^3+\psi_{M}(M^3))))%
}

\NovelDisplay{M04080R01}{%
(0)(1,1,1)(2,1,1)(3,1,1)(2,1,1)(3,1)(4)%
}

\NovelDisplay{M04081R01}{%
= \psi({I(1,0)}^{{I(1,0)}^\omega})%
}

\NovelDisplay{M04082R01}{%
= \psi(M^3+\psi_{\psi_{M}(M^3+M)}(M^3+\psi_{\psi_{M}(M^3+M)}(M^3+\psi_{\psi_{M}(M^3+M)}(1))))%
}

同一底数的指数也有倍数、幂与嵌套，唐三把指数框画全，终于没有将幂内乘法念成幂外乘法。

古月娜有一次把纸转向蔡须坤，请她只看I侧。

蔡须坤先读了一句，便想翻回自己的M册。她能在另一边找到熟悉的结构，面对眼前这一串参数与下标，却仍会因为不习惯而跳过中间的停顿。

卢本韦等她把第一句话说完，才指出遗漏的位置，没有说换一套记号就会变简单。蔡须坤将那一处重新念过，又问了自己还不确定的一个下标。

轮到卢本韦读M侧，她同样停了下来。蔡须坤忽然明白，同伴从前把I册一直留在身边，并不是不愿看更难的东西，而是要先站在自己能说明的那一处，再往陌生的结构里走。

两人各将一张纸留给对方，没有在上面写满答案，只圈出需要重新核的边界。第二天再交换时，卢本韦指着已经由自己补清的一处，先请蔡须坤把问题重新问一遍。

蔡须坤照问。她听到回答，才将纸收回自己那边。

\NovelDisplay{M04083R01}{%
(0)(1,1,1)(2,1,1)(3,1,1)(2,1,1)(3,1)(4,1)(2)%
}

\NovelDisplay{M04084R01}{%
= \psi({I(1,0)}^{{I(1,0)}^{I(1,0)}})%
}

\NovelDisplay{M04085R01}{%
= \psi(M^3+\psi_{\psi_{M}(M^3+M)}(M^3+\psi_{\psi_{M}(M^3+M)}(M^3+\psi_{\psi_{M}(M^3+M)}(M^3))))%
}

\NovelDisplay{M04086R01}{%
(0)(1,1,1)(2,1,1)(3,1,1)(2,1,1)(3,1)(4,2)%
}

\NovelDisplay{M04087R01}{%
= \psi(\Omega_{I(1,0)+1})%
}

\NovelDisplay{M04088R01}{%
= \psi(M^3+M)%
}

\NovelDisplay{M04089R01}{%
(0)(1,1,1)(2,1,1)(3,1,1)(2,1,1)(3,1)(4,2)(1,1,1)(2,1,1)(3,1,1)%
}
\NovelDisplay{M04089R02}{%
(2,1,1)(3,1)%
}

\NovelDisplay{M04090R01}{%
= \psi(\Omega_{I(1,0)+1}+I_\Omega)%
}

\NovelDisplay{M04091R01}{%
= \psi(M^3+M+\psi_{M}(M^2\times \psi_{M}(M)))%
}

\NovelDisplay{M04092R01}{%
(0)(1,1,1)(2,1,1)(3,1,1)(2,1,1)(3,1)(4,2)(1,1,1)(2,1,1)(3,1,1)%
}
\NovelDisplay{M04092R02}{%
(2,1,1)(3,1)(2)%
}

\NovelDisplay{M04093R01}{%
= \psi(\Omega_{I(1,0)+1}+\psi_{I(1,0)}(I(1,0)))%
}

\NovelDisplay{M04094R01}{%
= \psi(M^3+M+\psi_{\psi_{M}(M^3)}(M^3))%
}

从更深的指数转到Ω下标以后，新的M加项属于给出的下标对应，不能沿用前一式的指数读法。

\NovelDisplay{M04095R01}{%
(0)(1,1,1)(2,1,1)(3,1,1)(2,1,1)(3,1)(4,2)(1,1,1)(2,1,1)(3,1,1)%
}
\NovelDisplay{M04095R02}{%
(2,1,1)(3,1)(2,1,1)%
}

\NovelDisplay{M04096R01}{%
= \psi(\Omega_{I(1,0)+1}+\psi_{I(1,0)}(I(1,0)\times \omega))%
}

\NovelDisplay{M04097R01}{%
= \psi(M^3+M+\psi_{\psi_{M}(M^3)}(M^3+\psi_{\psi_{M}(M^3+M)}(1)))%
}

\NovelDisplay{M04098R01}{%
(0)(1,1,1)(2,1,1)(3,1,1)(2,1,1)(3,1)(4,2)(1,1,1)(2,1,1)(3,1,1)%
}
\NovelDisplay{M04098R02}{%
(2,1,1)(3,1)(4,2)%
}

\NovelDisplay{M04099R01}{%
= \psi(\Omega_{I(1,0)+1}+\psi_{I(1,0)}(\Omega_{I(1,0)+1}))%
}

\NovelDisplay{M04100R01}{%
= \psi(M^3+M+\psi_{\psi_{M}(M^3)}(M^3+M))%
}

\NovelDisplay{M04101R01}{%
(0)(1,1,1)(2,1,1)(3,1,1)(2,1,1)(3,1)(4,2)(1,1,1)(2,1,1)(3,1,1)%
}
\NovelDisplay{M04101R02}{%
(2,1,1)(3,1)(4,2)%
}

\NovelDisplay{M04102R01}{%
= \psi(\Omega_{I(1,0)+1}+\psi_{I(1,0)}(\Omega_{I(1,0)+1}))%
}

\NovelDisplay{M04103R01}{%
= \psi(M^3+M+\psi_{\psi_{M}(M^3)}(M^3+M))%
}

\NovelDisplay{M04104R01}{%
(0)(1,1,1)(2,1,1)(3,1,1)(2,1,1)(3,1)(4,2)(2)%
}

\NovelDisplay{M04105R01}{%
= \psi(\Omega_{I(1,0)+1}+\psi_{\Omega_{\psi_{I(1,0)}(\Omega_{I(1,0)+1})+1}}(1))%
}

\NovelDisplay{M04106R01}{%
= \psi(M^3+M+\psi_{\psi_{M}(M^2 \times \psi_{\psi_{M}(M^3)}(M^3+M)+M)}(1))%
}

重复对应按原稿保留，而新折叠项出现时仍须另读它的完整下标，不能据重复项自行外推。

蔡须坤第一次在这一段真正退回，是因为一个熟悉的ψ下标。

她看见旧背景，便顺着先前的接法往下说，脚底的力道却突然向后拉。霍雨浩托住她的手臂，让她自己收回那只脚，再问下标虽然熟悉，它这回作用的对象是不是也相同。

她没有立刻答。等纸铺回平处，她沿着括号慢慢读，才发现自己只认出了名字，没有认出它这次真正接住的那一段。

\NovelDisplay{M04107R01}{%
(0)(1,1,1)(2,1,1)(3,1,1)(2,1,1)(3,1)(4,2)(2,1)(1,1,1)(2,1,1)(3,1,1)%
}
\NovelDisplay{M04107R02}{%
(2,1,1)(3,1)(4,2)%
}

\NovelDisplay{M04108R01}{%
= \psi(\Omega_{I(1,0)+1}+\psi_{\Omega_{\psi_{I(1,0)}(\Omega_{I(1,0)+1})+1}}(\Omega_{I(1,0)+1}))%
}

\NovelDisplay{M04109R01}{%
= \psi(M^3+M+\psi_{\psi_{M}(M^2 \times \psi_{\psi_{M}(M^3)}(M^3+M)+M)}(M^3+M))%
}

\NovelDisplay{M04110R01}{%
(0)(1,1,1)(2,1,1)(3,1,1)(2,1,1)(3,1)(4,2)(2,1)(3,2)%
}

\NovelDisplay{M04111R01}{%
= \psi(\Omega_{I(1,0)+1}+\Omega_{\psi_{I(1,0)}(\Omega_{I(1,0)+1})+1})%
}

\NovelDisplay{M04112R01}{%
= \psi(M^3+M+\psi_{M}(M^2 \times \psi_{\psi_{M}(M^3)}(M^3+M)+M))%
}

\NovelDisplay{M04113R01}{%
(0)(1,1,1)(2,1,1)(3,1,1)(2,1,1)(3,1)(4,2)(2,1)(3,2,1)(4,2,1)(5,2)(4)%
}

\NovelDisplay{M04114R01}{%
= \psi(\Omega_{I(1,0)+1}+I_{\psi_{I(1,0)}(\Omega_{I(1,0)+1})+1})%
}

\NovelDisplay{M04115R01}{%
= \psi(M^3+M+\psi_{M}(M^2 \times \psi_{\psi_{M}(M^3)}(M^3+M)+M^2))%
}

\NovelDisplay{M04116R01}{%
(0)(1,1,1)(2,1,1)(3,1,1)(2,1,1)(3,1)(4,2)(2,1)(3,2,1)(4,2,1)(5,2,1)%
}

\NovelDisplay{M04117R01}{%
= \psi(\Omega_{I(1,0)+1}+I_{\psi_{I(1,0)}(\Omega_{I(1,0)+1})+\omega})%
}

\NovelDisplay{M04118R01}{%
= \psi(M^3+M+\psi_{M}(M^2 \times \psi_{\psi_{M}(M^3)}(M^3+M)+M^2\times \omega))%
}

相同外框下的作用对象、直接加项和I项各自改变，蔡须坤逐一指出其中的M与平方没有落在同一位置。

\NovelDisplay{M04119R01}{%
(0)(1,1,1)(2,1,1)(3,1,1)(2,1,1)(3,1)(4,2)(2,1)(3,2,1)(4,2,1)%
}
\NovelDisplay{M04119R02}{%
(5,2,1)(4,2,1)(5,2)%
}

\NovelDisplay{M04120R01}{%
= \psi(\Omega_{I(1,0)+1}+I_{\Omega_{\psi_{I(1,0)}(\Omega_{I(1,0)+1})+1}})%
}

\NovelDisplay{M04121R01}{%
= \psi(M^3+M+\psi_{\psi_{M}(M^3)}(M^3+M+\psi_{M}(M^2 \times \psi_{\psi_{M}(M^3)}(M^3+M)+M)))%
}

\NovelDisplay{M04122R01}{%
(0)(1,1,1)(2,1,1)(3,1,1)(2,1,1)(3,1)(4,2)(2,1)(3,2,1)(4,2,1)(5,2,1)%
}
\NovelDisplay{M04122R02}{%
(4,2,1)(5,2)(4)%
}

\NovelDisplay{M04123R01}{%
= \psi(\Omega_{I(1,0)+1}+I(1,0))%
}

\NovelDisplay{M04124R01}{%
= \psi(M^3+M+\psi_{M}(M^3))%
}

\NovelDisplay{M04125R01}{%
(0)(1,1,1)(2,1,1)(3,1,1)(2,1,1)(3,1)(4,2)(2,1)(3,2,1)(4,2,1)(5,2,1)%
}
\NovelDisplay{M04125R02}{%
(4,2,1)(5,2)(4,2,1)%
}

\NovelDisplay{M04126R01}{%
= \psi(\Omega_{I(1,0)+1}+I(1,0)\times \omega)%
}

\NovelDisplay{M04127R01}{%
= \psi(M^3+M+\psi_{\psi_{M}(M^3+M)}(1))%
}

\NovelDisplay{M04128R01}{%
(0)(1,1,1)(2,1,1)(3,1,1)(2,1,1)(3,1)(4,2)(2,1)(3,2,1)(4,2,1)(5,2,1)%
}
\NovelDisplay{M04128R02}{%
(4,2,1)(5,2)(6,3)%
}

\NovelDisplay{M04129R01}{%
= \psi(\Omega_{I(1,0)+1}+\psi_{\Omega_{I(1,0)+1}}(\Omega_{I(1,0)+1}))%
}

\NovelDisplay{M04130R01}{%
= \psi(M^3+M+\psi_{\psi_{M}(M^3+M)}(M^3+M))%
}

原背景后的回接项有各自的ψ下标，外层加号不能因为内部太长就被吞进其中一层。

\NovelDisplay{M04131R01}{%
(0)(1,1,1)(2,1,1)(3,1,1)(2,1,1)(3,1)(4,2)(2,1,1)%
}

\NovelDisplay{M04132R01}{%
= \psi(\Omega_{I(1,0)+1}+\psi_{\Omega_{I(1,0)+1}}(\Omega_{I(1,0)+1}+1))%
}

\NovelDisplay{M04133R01}{%
= \psi(M^3+M+\psi_{\psi_{M}(M^3+M)}(M^3+M+1))%
}

\NovelDisplay{M04134R01}{%
(0)(1,1,1)(2,1,1)(3,1,1)(2,1,1)(3,1)(4,2)(4,2)%
}

\NovelDisplay{M04135R01}{%
= \psi(\Omega_{I(1,0)+1}\times 2)%
}

\NovelDisplay{M04136R01}{%
= \psi(M^3+M+\psi_{M}(M^3+M))%
}

\NovelDisplay{M04137R01}{%
(0)(1,1,1)(2,1,1)(3,1,1)(2,1,1)(3,1)(4,2)(5,3)%
}

\NovelDisplay{M04138R01}{%
= \psi(\Omega_{I(1,0)+2})%
}

\NovelDisplay{M04139R01}{%
= \psi(M^3+M\times 2)%
}

\NovelDisplay{M04140R01}{%
(0)(1,1,1)(2,1,1)(3,1,1)(2,1,1)(3,1)(4,2,1)%
}

\NovelDisplay{M04141R01}{%
= \psi(\Omega_{I(1,0)+\omega})%
}

\NovelDisplay{M04142R01}{%
= \psi(M^3+M\times \omega)%
}

下标后继与整个Ω项的倍数并不相同，这几条短式正好用来核对自己有没有读错作用范围。

有一晚，诺当把两条相同的右式圈在一起。

他想用它们节省一行，薛峰却请他把两条原矩阵也同时留下。讲义里的重复项可以照原样保留，却不能因此声称，所有相近的矩阵都能按这种方式互换。

诺当听完，又将原来那一行写回去。蔡须坤坐在旁边，没有把这当成浪费纸；她的册子里，也已经有太多外表相像却必须重新核对的地方。

\NovelDisplay{M04143R01}{%
(0)(1,1,1)(2,1,1)(3,1,1)(2,1,1)(3,1)(4,2,1)(5,2,1)(6,2)(5)%
}

\NovelDisplay{M04144R01}{%
= \psi(I_{I(1,0)+1})%
}

\NovelDisplay{M04145R01}{%
= \psi(M^3+M^2)%
}

\NovelDisplay{M04146R01}{%
(0)(1,1,1)(2,1,1)(3,1,1)(2,1,1)(3,1)(4,2,1)(5,2,1)(6,2)(7,3)%
}

\NovelDisplay{M04147R01}{%
= \psi(\Omega_{I_{I(1,0)+1}+1})%
}

\NovelDisplay{M04148R01}{%
= \psi(M^3+M^2+M)%
}

\NovelDisplay{M04149R01}{%
(0)(1,1,1)(2,1,1)(3,1,1)(2,1,1)(3,1)(4,2,1)(5,2,1)(6,2,1)%
}

\NovelDisplay{M04150R01}{%
= \psi(I_{I(1,0)+\omega})%
}

\NovelDisplay{M04151R01}{%
= \psi(M^3+M^2\times \omega)%
}

\NovelDisplay{M04152R01}{%
(0)(1,1,1)(2,1,1)(3,1,1)(2,1,1)(3,1)(4,2,1)(5,2,1)(6,2,1)%
}

\NovelDisplay{M04153R01}{%
= \psi(I_{I(1,0)+\omega})%
}

\NovelDisplay{M04154R01}{%
= \psi(M^3+M^2\times \omega)%
}

I下标的后继对应平方位置，重复的ω式照原稿并排保留，不将相似误当成另一条已证明的规律。

\NovelDisplay{M04155R01}{%
(0)(1,1,1)(2,1,1)(3,1,1)(2,1,1)(3,1)(4,2,1)(5,2,1)(6,2,1)(5,2,1)(6,1)(2)%
}

\NovelDisplay{M04156R01}{%
= \psi(I_{I(1,0)\times 2})%
}

\NovelDisplay{M04157R01}{%
= \psi(M^3+M^2\times \psi_{M}(M^3))%
}

\NovelDisplay{M04158R01}{%
(0)(1,1,1)(2,1,1)(3,1,1)(2,1,1)(3,1)(4,2,1)(5,2,1)(6,2,1)(5,2,1)(6,2)%
}

\NovelDisplay{M04159R01}{%
= \psi(I_{\Omega_{I(1,0)+1}})%
}

\NovelDisplay{M04160R01}{%
= \psi(M^3+M^2\times \psi_{M}(M^3+M))%
}

\NovelDisplay{M04161R01}{%
(0)(1,1,1)(2,1,1)(3,1,1)(2,1,1)(3,1)(4,2,1)(5,2,1)(6,2,1)(5,2,1)(6,2)(5)%
}

\NovelDisplay{M04162R01}{%
= \psi(I(1,1))%
}

\NovelDisplay{M04163R01}{%
= \psi(M^3\times 2)%
}

\NovelDisplay{M04164R01}{%
(0)(1,1,1)(2,1,1)(3,1,1)(2,1,1)(3,1)(4,2,1)(5,2,1)(6,2,1)(5,2,1)(6,2)(5)%
}
\NovelDisplay{M04164R02}{%
(1,1,1)(2,1,1)(3,1)(4,2,1)(5,2,1)(6,2,1)(5,2,1)(6,2)(7,3,1)(8,3,1)%
}
\NovelDisplay{M04164R03}{%
(9,3,1)(8,3,1)(9,3)(8)(2,1)(3,2,1)(4,2,1)(5,2)(4)%
}

\NovelDisplay{M04165R01}{%
= \psi(I(1,1)+I)%
}

\NovelDisplay{M04166R01}{%
= \psi(M^3\times 2+\psi_{M}(M^2))%
}

进入双参数I以后，第二参数的改变需要连同整条矩阵一起看，不能把逗号后的一个一当成外面的加一。

这一条有三十一列，末列最高非零位置在第一行，坏根位于第二十八列；去掉末列后保留三十列，第二十八至第三十列共三列作为生成段。

薛峰先独立圈出范围，请叶阳只核位置，不替自己写后面的列；唐三与小舞再分别由同一原式求另外两项。

\NovelDisplay{M04167R01}{%
\text{展开第零项：}\quad(0)(1,1,1)(2,1,1)(3,1,1)(2,1,1)(3,1)(4,2,1)(5,2,1)%
}
\NovelDisplay{M04167R02}{%
(6,2,1)(5,2,1)(6,2)(5)(1,1,1)(2,1,1)(3,1)(4,2,1)(5,2,1)(6,2,1)(5,2,1)%
}
\NovelDisplay{M04167R03}{%
(6,2)(7,3,1)(8,3,1)(9,3,1)(8,3,1)(9,3)(8)(2,1)(3,2,1)(4,2,1)(5,2)%
}

\NovelDisplay{M04168R01}{%
\text{展开第一项：}\quad(0)(1,1,1)(2,1,1)(3,1,1)(2,1,1)(3,1)(4,2,1)(5,2,1)%
}
\NovelDisplay{M04168R02}{%
(6,2,1)(5,2,1)(6,2)(5)(1,1,1)(2,1,1)(3,1)(4,2,1)(5,2,1)(6,2,1)(5,2,1)%
}
\NovelDisplay{M04168R03}{%
(6,2)(7,3,1)(8,3,1)(9,3,1)(8,3,1)(9,3)(8)(2,1)(3,2,1)(4,2,1)(5,2)%
}
\NovelDisplay{M04168R04}{%
(3,2,1)(4,2,1)(5,2)%
}

\NovelDisplay{M04169R01}{%
\text{展开第二项：}\quad(0)(1,1,1)(2,1,1)(3,1,1)(2,1,1)(3,1)(4,2,1)(5,2,1)%
}
\NovelDisplay{M04169R02}{%
(6,2,1)(5,2,1)(6,2)(5)(1,1,1)(2,1,1)(3,1)(4,2,1)(5,2,1)(6,2,1)(5,2,1)%
}
\NovelDisplay{M04169R03}{%
(6,2)(7,3,1)(8,3,1)(9,3,1)(8,3,1)(9,3)(8)(2,1)(3,2,1)(4,2,1)(5,2)(3,2,1)%
}
\NovelDisplay{M04169R04}{%
(4,2,1)(5,2)(3,2,1)(4,2,1)(5,2)%
}

三项分别有三十列、三十三列和三十六列，共同前段再长也不能遮住真正只有三列的坏部。

蔡须坤数过，发现自己最初因为矩阵太长，险些连前面一段也圈进去。她没有擦掉原来的记号，另用一道细线指出自己多圈的位置。

走到双参数I的新背景，卢本韦比先前快了一些。

她刚说完，抬头见小舞仍停在页首，才发现自己已经把参数里的那个变化一口带过。小舞说听不见她停在哪里，越听越像只要把M侧的乘二抄过去就行。

卢本韦将册子退回前一页，问她先前能确定哪一处。小舞指出以后，她才从那里接下去，声音也放慢了。

蔡须坤把这个停顿记在自己的M册里。

\NovelDisplay{M04170R01}{%
(0)(1,1,1)(2,1,1)(3,1,1)(2,1,1)(3,1)(4,2,1)(5,2,1)(6,2,1)(5,2,1)(6,2)(5)%
}
\NovelDisplay{M04170R02}{%
(1,1,1)(2,1,1)(3,1,1)%
}

\NovelDisplay{M04171R01}{%
= \psi(I(1,1)+I_\omega)%
}

\NovelDisplay{M04172R01}{%
= \psi(M^3\times 2+\psi_{M}(M^2\times \omega))%
}

\NovelDisplay{M04173R01}{%
(0)(1,1,1)(2,1,1)(3,1,1)(2,1,1)(3,1)(4,2,1)(5,2,1)(6,2,1)(5,2,1)(6,2)(5)%
}
\NovelDisplay{M04173R02}{%
(1,1,1)(2,1,1)(3,1,1)(2,1,1)(3,1)(4,2,1)(5,2,1)(6,2,1)(5,2,1)(6,2)(5)%
}

\NovelDisplay{M04174R01}{%
= \psi(I(1,1)+\psi_{I(1,0)}(I(1,1)))%
}

\NovelDisplay{M04175R01}{%
= \psi(M^3\times 2+\psi_{\psi_{M}(M^3)}(M^3\times 2))%
}

\NovelDisplay{M04176R01}{%
(0)(1,1,1)(2,1,1)(3,1,1)(2,1,1)(3,1)(4,2,1)(5,2,1)(6,2,1)(5,2,1)(6,2)(5)%
}
\NovelDisplay{M04176R02}{%
(2,1)(3,2,1)(4,2,1)(5,2,1)(4,2,1)(5,2)(4)%
}

\NovelDisplay{M04177R01}{%
= \psi(I(1,1)+I(1,0))%
}

\NovelDisplay{M04178R01}{%
= \psi(M^3\times 2+\psi_{M}(M^3))%
}

\NovelDisplay{M04179R01}{%
(0)(1,1,1)(2,1,1)(3,1,1)(2,1,1)(3,1)(4,2,1)(5,2,1)(6,2,1)(5,2,1)(6,2)(5)%
}
\NovelDisplay{M04179R02}{%
(2,1)(3,2,1)(4,2,1)(5,2,1)(4,2,1)(5,2)(6,3,1)(7,3,1)(8,3,1)(7,3,1)%
}
\NovelDisplay{M04179R03}{%
(8,3)(7)%
}

\NovelDisplay{M04180R01}{%
= \psi(I(1,1)+\psi_{\Omega_{I(1,0)+1}}(I(1,1)))%
}

\NovelDisplay{M04181R01}{%
= \psi(M^3\times 2+\psi_{\psi_{M}(M^2\times \psi_{M}(M^3)+M)}(M^3\times 2))%
}

新的双参数背景不改变每个ψ必须读清下标与作用对象的要求，薛峰先核前段，再问后面的回接。

这天复核以后，卢本韦把两张有改动的纸夹在最外面。

蔡须坤以为她明日还要先讲，伸手想取过来看。卢本韦却说，最外面的两处得先让自己独立问完，再拿给她核。

蔡须坤停住，将手收回。她习惯一看见同伴的疑问便坐到旁边，这一回才发觉，有时留下思考的位置，也是一种帮助。

到了第二天，卢本韦果然先来找她。两个问题有一处已经由她自己接清，另一处仍留着空行。蔡须坤听完，便也把自己的M册翻到昨夜留下的问号处。

霍雨浩上山时，两人正各自说明一半，再将问题交给对方。薛峰在旁边核原矩阵，等她们谈完，才指出纸上仍少读的一列；那一列由最初漏下它的人重新补好，谁也没把一回互助当成往后都不必独立确认的凭据。

两人看过那处补写，各自收回自己的本子，随后一同去取山上的水。

\NovelDisplay{M04182R01}{%
(0)(1,1,1)(2,1,1)(3,1,1)(2,1,1)(3,1)(4,2,1)(5,2,1)(6,2,1)(5,2,1)(6,2)%
}
\NovelDisplay{M04182R02}{%
(5)(2,1,1)%
}

\NovelDisplay{M04183R01}{%
= \psi(I(1,1)+\psi_{\Omega_{I(1,0)+1}}(I(1,1)+1))%
}

\NovelDisplay{M04184R01}{%
= \psi(M^3\times 2+\psi_{\psi_{M}(M^2\times \psi_{M}(M^3)+M)}(M^3\times 2+1))%
}

\NovelDisplay{M04185R01}{%
(0)(1,1,1)(2,1,1)(3,1,1)(2,1,1)(3,1)(4,2,1)(5,2,1)(6,2,1)(5,2,1)(6,2)%
}
\NovelDisplay{M04185R02}{%
(5)(4,2)%
}

\NovelDisplay{M04186R01}{%
= \psi(I(1,1)+\Omega_{I(1,0)+1})%
}

\NovelDisplay{M04187R01}{%
= \psi(M^3\times 2+\psi_{M}(M^2\times \psi_{M}(M^3)+M))%
}

\NovelDisplay{M04188R01}{%
(0)(1,1,1)(2,1,1)(3,1,1)(2,1,1)(3,1)(4,2,1)(5,2,1)(6,2,1)(5,2,1)(6,2)(5)%
}
\NovelDisplay{M04188R02}{%
(4,2,1)(5,2,1)(6,2,1)(5,2,1)(6,2)(5)%
}

\NovelDisplay{M04189R01}{%
= \psi(I(1,1)+\psi_{I(1,1)}(I(1,1)))%
}

\NovelDisplay{M04190R01}{%
= \psi(M^3\times 2+\psi_{\psi_{M}(M^3\times 2)}(M^3\times 2))%
}

\NovelDisplay{M04191R01}{%
(0)(1,1,1)(2,1,1)(3,1,1)(2,1,1)(3,1)(4,2,1)(5,2,1)(6,2,1)(5,2,1)%
}
\NovelDisplay{M04191R02}{%
(6,2)(5,2)(6,3)%
}

\NovelDisplay{M04192R01}{%
= \psi(I(1,1)+\Omega_{\psi_{I(1,1)}(I(1,1))+1})%
}

\NovelDisplay{M04193R01}{%
= \psi(M^3\times 2+\psi_{M}(M^2\times \psi_{\psi_{M}(M^3\times 2)}(M^3\times 2)+M))%
}

原来的背景换成双倍立方以后，相关下标也须按式中位置更新，不能只改最外面的乘二。

\NovelDisplay{M04194R01}{%
(0)(1,1,1)(2,1,1)(3,1,1)(2,1,1)(3,1)(4,2,1)(5,2,1)(6,2,1)(5,2,1)(6,2)%
}
\NovelDisplay{M04194R02}{%
(5,2)(6,3,1)(7,3,1)(8,3,1)(7,3,1)(8,3)(7)%
}

\NovelDisplay{M04195R01}{%
= \psi(I(1,1)\times 2)%
}

\NovelDisplay{M04196R01}{%
= \psi(M^3\times 2+\psi_{M}(M^3\times 2))%
}

\NovelDisplay{M04197R01}{%
(0)(1,1,1)(2,1,1)(3,1,1)(2,1,1)(3,1)(4,2,1)(5,2,1)(6,2,1)(5,2,1)%
}
\NovelDisplay{M04197R02}{%
(6,2)(5,2,1)%
}

\NovelDisplay{M04198R01}{%
= \psi(I(1,1)\times \omega)%
}

\NovelDisplay{M04199R01}{%
= \psi(M^3\times 2+\psi_{\psi_{M}(M^3\times 2+M)}(1))%
}

\NovelDisplay{M04200R01}{%
(0)(1,1,1)(2,1,1)(3,1,1)(2,1,1)(3,1)(4,2,1)(5,2,1)(6,2,1)(5,2,1)%
}
\NovelDisplay{M04200R02}{%
(6,2)(7,3)%
}

\NovelDisplay{M04201R01}{%
= \psi(\Omega_{I(1,1)+1})%
}

\NovelDisplay{M04202R01}{%
= \psi(M^3\times 2+M)%
}

\NovelDisplay{M04203R01}{%
(0)(1,1,1)(2,1,1)(3,1,1)(2,1,1)(3,1)(4,2,1)(5,2,1)(6,2,1)(5,2,1)(6,2)%
}
\NovelDisplay{M04203R02}{%
(7,3,1)(8,3,1)(9,3,1)(8,3,1)(9,3)(8)%
}

\NovelDisplay{M04204R01}{%
= \psi(I(1,2))%
}

\NovelDisplay{M04205R01}{%
= \psi(M^3\times 3)%
}

同一双参数对象的倍数、ω倍数与下标后继分开处理，才能看出它们各自与后面新背景的区别。

那天下午，薛峰第一次将自己的纸主动交给叶阳。

他本以为已经核完，叶阳却问同一背景后的两个加项能不能交换。薛峰差点说都在后面有什么区别，话刚到嘴边，便看见霍雨浩望过来。

他把原稿拉近，重新指出顺序，承认自己刚才答得太随意。叶阳没有追着笑，只拿回自己的纸，将这一问认真留下。

\NovelDisplay{M04206R01}{%
(0)(1,1,1)(2,1,1)(3,1,1)(2,1,1)(3,1)(4,2,1)(5,2,1)(6,2,1)(5,2,1)(6,2)%
}
\NovelDisplay{M04206R02}{%
(7,3,1)(8,3,1)(9,3,1)(8,3,1)(9,3)(8)(2,1)(3,2,1)(4,2,1)(5,2,1)%
}
\NovelDisplay{M04206R03}{%
(4,2,1)(5,2)(4)%
}

\NovelDisplay{M04207R01}{%
= \psi(I(1,2)+I(1,0))%
}

\NovelDisplay{M04208R01}{%
= \psi(M^3\times 3+\psi_{M}(M^3))%
}

\NovelDisplay{M04209R01}{%
(0)(1,1,1)(2,1,1)(3,1,1)(2,1,1)(3,1)(4,2,1)(5,2,1)(6,2,1)(5,2,1)(6,2)%
}
\NovelDisplay{M04209R02}{%
(7,3,1)(8,3,1)(9,3,1)(8,3,1)(9,3)(8)(5,2)(6,3,1)(7,3,1)(8,3,1)%
}
\NovelDisplay{M04209R03}{%
(7,3,1)(8,3)(7)%
}

\NovelDisplay{M04210R01}{%
= \psi(I(1,2)+I(1,1))%
}

\NovelDisplay{M04211R01}{%
= \psi(M^3\times 3+\psi_{M}(M^3\times 2))%
}

\NovelDisplay{M04212R01}{%
(0)(1,1,1)(2,1,1)(3,1,1)(2,1,1)(3,1)(4,2,1)(5,2,1)(6,2,1)(5,2,1)(6,2)%
}
\NovelDisplay{M04212R02}{%
(7,3,1)(8,3,1)(9,3,1)(8,3,1)(9,3)(8)(8,3)(9,4,1)(10,4,1)(11,4,1)%
}
\NovelDisplay{M04212R03}{%
(10,4,1)(11,4)(10)%
}

\NovelDisplay{M04213R01}{%
= \psi(I(1,2)\times 2)%
}

\NovelDisplay{M04214R01}{%
= \psi(M^3\times 3+\psi_{M}(M^3\times 3))%
}

\NovelDisplay{M04215R01}{%
(0)(1,1,1)(2,1,1)(3,1,1)(2,1,1)(3,1)(4,2,1)(5,2,1)(6,2,1)(5,2,1)(6,2)%
}
\NovelDisplay{M04215R02}{%
(7,3,1)(8,3,1)(9,3,1)(8,3,1)(9,3)(10,4,1)(11,4,1)(12,4,1)(11,4,1)%
}
\NovelDisplay{M04215R03}{%
(12,4)(11)%
}

\NovelDisplay{M04216R01}{%
= \psi(I(1,3))%
}

\NovelDisplay{M04217R01}{%
= \psi(M^3\times 4)%
}

I侧的第二参数在这里继续增加，M侧保留自己的折叠项，原序数加法的次序不能随口交换。

卢本韦把三次取项分开落在不同纸上，诺当先将三份保留段对齐，再请霍雨浩看有没有沿用上一条的坏根。

\NovelDisplay{M04218R01}{%
\text{展开第零项：}\quad(0)(1,1,1)(2,1,1)(3,1,1)(2,1,1)(3,1)(4,2,1)(5,2,1)%
}
\NovelDisplay{M04218R02}{%
(6,2,1)(5,2,1)(6,2)(7,3,1)(8,3,1)(9,3,1)(8,3,1)(9,3)(8)(8,3)(9,4,1)%
}
\NovelDisplay{M04218R03}{%
(10,4,1)(11,4,1)(10,4,1)(11,4)%
}

\NovelDisplay{M04219R01}{%
\text{展开第一项：}\quad(0)(1,1,1)(2,1,1)(3,1,1)(2,1,1)(3,1)(4,2,1)(5,2,1)%
}
\NovelDisplay{M04219R02}{%
(6,2,1)(5,2,1)(6,2)(7,3,1)(8,3,1)(9,3,1)(8,3,1)(9,3)(8)(8,3)(9,4,1)%
}
\NovelDisplay{M04219R03}{%
(10,4,1)(11,4,1)(10,4,1)(11,4)(9,4,1)(10,4,1)(11,4,1)(10,4,1)(11,4)%
}

\NovelDisplay{M04220R01}{%
\text{展开第二项：}\quad(0)(1,1,1)(2,1,1)(3,1,1)(2,1,1)(3,1)(4,2,1)(5,2,1)%
}
\NovelDisplay{M04220R02}{%
(6,2,1)(5,2,1)(6,2)(7,3,1)(8,3,1)(9,3,1)(8,3,1)(9,3)(8)(8,3)(9,4,1)%
}
\NovelDisplay{M04220R03}{%
(10,4,1)(11,4,1)(10,4,1)(11,4)(9,4,1)(10,4,1)(11,4,1)(10,4,1)(11,4)%
}
\NovelDisplay{M04220R04}{%
(9,4,1)(10,4,1)(11,4,1)(10,4,1)(11,4)%
}

三次生成的共同起点是眼前这条原矩阵，不是前一张纸上刚得到的结果。

诺当发现自己的第一项少了最初保留的一列，自己补回以后，又从头核了一次第二项。

当天傍晚，唐三没有急着下山。

他与小舞坐在石桩下，将各自准备问的问题写在不同纸上，再交换着看。小舞看到哥哥的问题，先想直接指答案，手伸到一半，却记起自己白天也是怎样被一句太快的回答弄得更乱。

她让唐三先说已经知道的那一段。唐三停过两次，才把完整的作用范围指出来；原先看着无从开口的疑问，也就落到了一个能够回答的位置。

轮到小舞时，他同样等着。海风从较低处吹来，两张纸压在膝上，谁也没有先去看上方的高度。

古月娜下来叫他们吃饭，见两人正轮着问，便在边上坐了一会儿。她核过其中一处，提醒唐三仍须回到原矩阵独立求项，才将带来的水留下。

后来收纸，唐三在两页角上各写了日期。那不是今天多走了多少的记录，而是他终于知道妹妹停在哪里，也终于肯让妹妹听自己还没弄懂的地方。

\NovelDisplay{M04221R01}{%
(0)(1,1,1)(2,1,1)(3,1,1)(2,1,1)(3,1,1)%
}

\NovelDisplay{M04222R01}{%
= \psi(I(1,\omega))%
}

\NovelDisplay{M04223R01}{%
= \psi(M^3\times \omega)%
}

\NovelDisplay{M04224R01}{%
(0)(1,1,1)(2,1,1)(3,1,1)(2,1,1)(3,1,1)(1,1,1)(2,1,1)(3,1,1)(2,1,1)(3,1)%
}
\NovelDisplay{M04224R02}{%
(4,2,1)(5,2,1)(6,2,1)(5,2,1)(6,2,1)%
}

\NovelDisplay{M04225R01}{%
= \psi(I(1,\omega)+\psi_{I(1,0)}(I(1,\omega)))%
}

\NovelDisplay{M04226R01}{%
= \psi(M^3\times \omega+\psi_{\psi_{M}(M^3)}(M^3\times \omega))%
}

\NovelDisplay{M04227R01}{%
(0)(1,1,1)(2,1,1)(3,1,1)(2,1,1)(3,1,1)(1,1,1)(2,1,1)(3,1,1)(2,1,1)(3,1)%
}
\NovelDisplay{M04227R02}{%
(4,2,1)(5,2,1)(6,2,1)(5,2,1)(6,2,1)(2,1)(3,2,1)(4,2,1)(5,2,1)%
}
\NovelDisplay{M04227R03}{%
(4,2,1)(5,2)(4)%
}

\NovelDisplay{M04228R01}{%
= \psi(I(1,\omega)+I(1,0))%
}

\NovelDisplay{M04229R01}{%
= \psi(M^3\times \omega+\psi_{M}(M^3))%
}

\NovelDisplay{M04230R01}{%
(0)(1,1,1)(2,1,1)(3,1,1)(2,1,1)(3,1,1)(1,1,1)(2,1,1)(3,1,1)(2,1,1)(3,1)%
}
\NovelDisplay{M04230R02}{%
(4,2,1)(5,2,1)(6,2,1)(5,2,1)(6,2,1)(4,2)%
}

\NovelDisplay{M04231R01}{%
= \psi(I(1,\omega)+\Omega_{I(1,0)+1})%
}

\NovelDisplay{M04232R01}{%
= \psi(M^3\times \omega+\psi_{M}(M^3+M))%
}

有限第二参数走向ω时，先确认新的背景，再辨认附加项，不能把整个极限段当成一张长表的最后一行。

唐三与小舞在这一段又落后了几日。

小舞先急，觉得其他人都已经把ω参数后的几条接上，自己却还在找下标。唐三陪她看，越看也越急，最后两个人谁都不再说话。

古月娜坐到旁边，将水放下，先让他们各自指出今天真正不会的一处。两张纸摊开，空白竟不在同一个地方。

他们对视一眼。唐三先听小舞问完，再让她反过来听自己那处。那天上山时，两个人都仍走得慢，却没有再把落后说成同一块无法拆开的石头。

\NovelDisplay{M04233R01}{%
(0)(1,1,1)(2,1,1)(3,1,1)(2,1,1)(3,1,1)(1,1,1)(2,1,1)(3,1,1)(2,1,1)(3,1)%
}
\NovelDisplay{M04233R02}{%
(4,2,1)(5,2,1)(6,2,1)(5,2,1)(6,2,1)(4,2,1)(5,2,1)(6,2,1)(5,2,1)(6,2)%
}
\NovelDisplay{M04233R03}{%
(7,3,1)(8,3,1)(9,3,1)(8,3,1)(9,3,1)%
}

\NovelDisplay{M04234R01}{%
= \psi(I(1,\omega)+\psi_{I(1,1)}(I(1,\omega)))%
}

\NovelDisplay{M04235R01}{%
= \psi(M^3\times \omega+\psi_{\psi_{M}(M^3\times 2)}(M^3\times \omega))%
}

\NovelDisplay{M04236R01}{%
(0)(1,1,1)(2,1,1)(3,1,1)(2,1,1)(3,1,1)(1,1,1)(2,1,1)(3,1,1)(2,1,1)(3,1)%
}
\NovelDisplay{M04236R02}{%
(4,2,1)(5,2,1)(6,2,1)(5,2,1)(6,2,1)(4,2,1)(5,2,1)(6,2,1)(5,2,1)(6,2)%
}
\NovelDisplay{M04236R03}{%
(7,3,1)(8,3,1)(9,3,1)(8,3,1)(9,3,1)(5,2)(6,3,1)(7,3,1)(8,3,1)(7,3,1)%
}
\NovelDisplay{M04236R04}{%
(8,3)(7)%
}

\NovelDisplay{M04237R01}{%
= \psi(I(1,\omega)+I(1,1))%
}

\NovelDisplay{M04238R01}{%
= \psi(M^3\times \omega+\psi_{M}(M^3\times 2))%
}

\NovelDisplay{M04239R01}{%
(0)(1,1,1)(2,1,1)(3,1,1)(2,1,1)(3,1,1)(1,1,1)(2,1,1)(3,1,1)%
}
\NovelDisplay{M04239R02}{%
(2,1,1)(3,1,1)%
}

\NovelDisplay{M04240R01}{%
= \psi(I(1,\omega)\times 2)%
}

\NovelDisplay{M04241R01}{%
= \psi(M^3\times \omega+\psi_{M}(M^3\times \omega))%
}

\NovelDisplay{M04242R01}{%
(0)(1,1,1)(2,1,1)(3,1,1)(2,1,1)(3,1,1)(2,1)(3,2)%
}

\NovelDisplay{M04243R01}{%
= \psi(\Omega_{I(1,\omega)+1})%
}

\NovelDisplay{M04244R01}{%
= \psi(M^3\times \omega+M)%
}

同一ω背景上作用的不同ψ拥有不同下标，卢本韦先讲清区别，霍雨浩再请后来的人从矩阵核对。

同一ω背景后的这条矩阵更长，蔡须坤让唐三先遮住右式，小舞与薛峰各求一项，等三人写完才交换纸页。

\NovelDisplay{M04245R01}{%
\text{展开第零项：}\quad(0)(1,1,1)(2,1,1)(3,1,1)(2,1,1)(3,1,1)(1,1,1)(2,1,1)%
}
\NovelDisplay{M04245R02}{%
(3,1,1)(2,1,1)(3,1)(4,2,1)(5,2,1)(6,2,1)(5,2,1)(6,2,1)(4,2,1)(5,2,1)%
}
\NovelDisplay{M04245R03}{%
(6,2,1)(5,2,1)(6,2)(7,3,1)(8,3,1)(9,3,1)(8,3,1)(9,3,1)(5,2)(6,3,1)%
}
\NovelDisplay{M04245R04}{%
(7,3,1)(8,3,1)(7,3,1)(8,3)%
}

\NovelDisplay{M04246R01}{%
\text{展开第一项：}\quad(0)(1,1,1)(2,1,1)(3,1,1)(2,1,1)(3,1,1)(1,1,1)(2,1,1)%
}
\NovelDisplay{M04246R02}{%
(3,1,1)(2,1,1)(3,1)(4,2,1)(5,2,1)(6,2,1)(5,2,1)(6,2,1)(4,2,1)(5,2,1)%
}
\NovelDisplay{M04246R03}{%
(6,2,1)(5,2,1)(6,2)(7,3,1)(8,3,1)(9,3,1)(8,3,1)(9,3,1)(5,2)(6,3,1)%
}
\NovelDisplay{M04246R04}{%
(7,3,1)(8,3,1)(7,3,1)(8,3)(6,3,1)(7,3,1)(8,3,1)(7,3,1)(8,3)%
}

\NovelDisplay{M04247R01}{%
\text{展开第二项：}\quad(0)(1,1,1)(2,1,1)(3,1,1)(2,1,1)(3,1,1)(1,1,1)(2,1,1)%
}
\NovelDisplay{M04247R02}{%
(3,1,1)(2,1,1)(3,1)(4,2,1)(5,2,1)(6,2,1)(5,2,1)(6,2,1)(4,2,1)(5,2,1)%
}
\NovelDisplay{M04247R03}{%
(6,2,1)(5,2,1)(6,2)(7,3,1)(8,3,1)(9,3,1)(8,3,1)(9,3,1)(5,2)(6,3,1)%
}
\NovelDisplay{M04247R04}{%
(7,3,1)(8,3,1)(7,3,1)(8,3)(6,3,1)(7,3,1)(8,3,1)(7,3,1)(8,3)(6,3,1)%
}
\NovelDisplay{M04247R05}{%
(7,3,1)(8,3,1)(7,3,1)(8,3)%
}

这三份结果分别保留三十二列、三十七列和四十二列，生成段的长度来自原列范围，不能由右式里的倍数猜出。

小舞的纸比从前整齐了些，仍有两处涂改。她没有拿唐三那份替换，只在自己的纸上把错接的边界重新划清。

\NovelDisplay{M04248R01}{%
(0)(1,1,1)(2,1,1)(3,1,1)(2,1,1)(3,1,1)(2,1)(3,2,1)(4,2,1)(5,2,1)%
}
\NovelDisplay{M04248R02}{%
(4,2,1)(5,2)(4)%
}

\NovelDisplay{M04249R01}{%
= \psi(I(1,\omega+1))%
}

\NovelDisplay{M04250R01}{%
= \psi(M^3\times \omega+M^3)%
}

\NovelDisplay{M04251R01}{%
(0)(1,1,1)(2,1,1)(3,1,1)(2,1,1)(3,1,1)(2,1)(3,2,1)(4,2,1)(5,2,1)(4,2,1)%
}
\NovelDisplay{M04251R02}{%
(5,2)(6,3,1)(7,3,1)(8,3,1)(7,3,1)(8,3)(7)%
}

\NovelDisplay{M04252R01}{%
= \psi(I(1,\omega+2))%
}

\NovelDisplay{M04253R01}{%
= \psi(M^3\times \omega+M^3\times 2)%
}

\NovelDisplay{M04254R01}{%
(0)(1,1,1)(2,1,1)(3,1,1)(2,1,1)(3,1,1)(2,1)(3,2,1)(4,2,1)(5,2,1)%
}
\NovelDisplay{M04254R02}{%
(4,2,1)(5,2,1)%
}

\NovelDisplay{M04255R01}{%
= \psi(I(1,\omega\times 2))%
}

\NovelDisplay{M04256R01}{%
= \psi(M^3\times \omega\times 2)%
}

\NovelDisplay{M04257R01}{%
(0)(1,1,1)(2,1,1)(3,1,1)(2,1,1)(3,1,1)(2,1,1)%
}

\NovelDisplay{M04258R01}{%
= \psi(I(1,\psi(2)))%
}

\NovelDisplay{M04259R01}{%
= \psi(M^3\times \psi(2))%
}

ω之后的有限尾部与ω倍数分属不同形式，写在参数里时仍要保留完整的括号。

最后几页，蔡须坤没有抢着先读M侧。

卢本韦指出第二参数，她便在对应里慢慢找到乘数的位置；遇到自己仍需回看的地方，先留出空行，再请霍雨浩核对作用范围。

她们各自选用熟悉的记号，却已能知道另一边的问题停在哪里。奥斯卡将补给分给后来的人，笑说这次总算不用端着碗站半天，等谁一句``马上就好''。

\NovelDisplay{M04260R01}{%
(0)(1,1,1)(2,1,1)(3,1,1)(2,1,1)(3,1,1)(2,1,1)(3,1)%
}

\NovelDisplay{M04261R01}{%
= \psi(I(1,\Omega))%
}

\NovelDisplay{M04262R01}{%
= \psi(M^3\times \psi_{M}(M))%
}

\NovelDisplay{M04263R01}{%
(0)(1,1,1)(2,1,1)(3,1,1)(2,1,1)(3,1,1)(2,1,1)(3,1)(1,1,1)(2,1,1)(3,1)%
}
\NovelDisplay{M04263R02}{%
(4,2,1)(5,2,1)(6,2,1)(5,2,1)(6,2,1)(5,2,1)(6,1)(2)%
}

\NovelDisplay{M04264R01}{%
= \psi(I(1,I))%
}

\NovelDisplay{M04265R01}{%
= \psi(M^3\times \psi_{M}(M^2))%
}

\NovelDisplay{M04266R01}{%
(0)(1,1,1)(2,1,1)(3,1,1)(2,1,1)(3,1,1)(2,1,1)(3,1)(1,1,1)(2,1,1)(3,1)%
}
\NovelDisplay{M04266R02}{%
(4,2,1)(5,2,1)(6,2,1)(5,2,1)(6,2,1)(5,2,1)(6,2)%
}

\NovelDisplay{M04267R01}{%
= \psi(I(1,\Omega_{I+1}))%
}

\NovelDisplay{M04268R01}{%
= \psi(M^3\times \psi_{M}(M^2+M))%
}

第二参数换成Ω及I侧对象时，蔡须坤保留相应折叠表示，不把它们当作可以直接替换的普通常数。

\NovelDisplay{M04269R01}{%
(0)(1,1,1)(2,1,1)(3,1,1)(2,1,1)(3,1,1)(2,1,1)(3,1)(1,1,1)(2,1,1)(3,1,1)%
}

\NovelDisplay{M04270R01}{%
= \psi(I(1,I_\omega))%
}

\NovelDisplay{M04271R01}{%
= \psi(M^3\times \psi_{M}(M^2\times \omega))%
}

\NovelDisplay{M04272R01}{%
(0)(1,1,1)(2,1,1)(3,1,1)(2,1,1)(3,1,1)(2,1,1)(3,1)(1,1,1)(2,1,1)(3,1,1)%
}
\NovelDisplay{M04272R02}{%
(2,1,1)(3,1)(4,2,1)(5,2,1)(6,2,1)(5,2,1)(6,2,1)(5,2,1)(6,1)(2)%
}

\NovelDisplay{M04273R01}{%
= \psi(I(1,I(1,0)))%
}

\NovelDisplay{M04274R01}{%
= \psi(M^3\times \psi_{M}(M^3))%
}

\NovelDisplay{M04275R01}{%
(0)(1,1,1)(2,1,1)(3,1,1)(2,1,1)(3,1,1)(2,1,1)(3,1)(1,1,1)(2,1,1)(3,1,1)%
}
\NovelDisplay{M04275R02}{%
(2,1,1)(3,1,1)%
}

\NovelDisplay{M04276R01}{%
= \psi(I(1,I(1,\omega)))%
}

\NovelDisplay{M04277R01}{%
= \psi(M^3\times \psi_{M}(M^3\times \omega))%
}

同一第一参数下的三个复杂第二参数各有对应，先看清参数边界，再回看M侧乘数里实际改变的一层。

第五个月末，十一人站在三百九十米。

小舞摸到刻痕以后，先看了一眼唐三。两人的本子里各有一页，留着这几日互相问过的问题；她的那页上，哥哥的字并没有替自己填满空白，只留下一个请她重新说明的位置。

蔡须坤在石桩旁坐下，望着I册里那个能一口读出的双参数，忽然又将M册翻开。

``你再从这里讲一次。''

卢本韦没有嫌她已经问过。她指回参数，蔡须坤循着对应，终于把原先仅凭熟悉念过的那一层说清。

山风掠过页角。薛峰把日期写好，霍雨浩等最后一人坐稳，才将水递出去。

离五百米还有一百一十米，他们已经知道，这个数不能靠多写几个等号缩短。

\chapter{第一百二十四章：二元I函数}\label{ux7b2cux4e00ux767eux4e8cux5341ux56dbux7ae0ux4e8cux5143iux51fdux6570}

第六个月，古月娜带回一叠更厚的纸。

不是哪条公式忽然变短了，而是旧页的折口已经磨破，再铺开时，有一个下标恰好断在纸缝里。小舞问能不能干脆只留短的写法，蔡须坤将那一段慢慢展开，沿着缝找到被挡住的括号。

``先把它看全。''

说完，她自己也怔了一下。几个月前，她最盼的还是能有一套不必走这些弯路的矩阵，现在最怕的，却是看似已经接通的地方，藏着没有真正说清的一层。

卢本韦摊开I册，第一参数由有限数走向ω，后面又将出现更多参数。M侧的幂也一层层深入，原来已经读熟的字母换了位置，便又成了新的难处。

霍雨浩把原矩阵放在桌中间，薛峰先请叶阳说今天准备核哪一列。

开头几式看着清楚，蔡兴坤答得很快。

古月娜问他，第一参数从一变成二，与第二参数从零变成一，是不是改了同一个位置。他先点头，随即又摇头，看了半天，才承认刚才只记住了右边出现四次幂。

没人催他继续猜。卢本韦沿I侧指出参数，蔡须坤在M侧找回相应位置；他重新读完，给两种改动画了不同的线。

\NovelDisplay{M04278R01}{%
(0)(1,1,1)(2,1,1)(3,1,1)(2,1,1)(3,1,1)(2,1,1)(3,1)(2)%
}

\NovelDisplay{M04279R01}{%
= \psi(I(2,0))%
}

\NovelDisplay{M04280R01}{%
= \psi(M^4)%
}

\NovelDisplay{M04281R01}{%
(0)(1,1,1)(2,1,1)(3,1,1)(2,1,1)(3,1,1)(2,1,1)(3,1)(2,1,1)%
}

\NovelDisplay{M04282R01}{%
= \psi(I(2,0)\times \omega)%
}

\NovelDisplay{M04283R01}{%
= \psi(M^4+\psi_{\psi_{M}(M^4+M)}(1))%
}

\NovelDisplay{M04284R01}{%
(0)(1,1,1)(2,1,1)(3,1,1)(2,1,1)(3,1,1)(2,1,1)(3,1)(4,2)%
}

\NovelDisplay{M04285R01}{%
= \psi(\Omega_{I(2,0)+1})%
}

\NovelDisplay{M04286R01}{%
= \psi(M^4+M)%
}

\NovelDisplay{M04287R01}{%
(0)(1,1,1)(2,1,1)(3,1,1)(2,1,1)(3,1,1)(2,1,1)(3,1)(4,2,1)(5,2,1)(6,2)(5)%
}

\NovelDisplay{M04288R01}{%
= \psi(I_{I(2,0)+1})%
}

\NovelDisplay{M04289R01}{%
= \psi(M^4+M^2)%
}

第一参数从一转到二以后，M侧出现四次幂背景，其上的加项仍需要逐条核对。

\NovelDisplay{M04290R01}{%
(0)(1,1,1)(2,1,1)(3,1,1)(2,1,1)(3,1,1)(2,1,1)(3,1)(4,2,1)(5,2,1)(6,2,1)%
}
\NovelDisplay{M04290R02}{%
(5,2,1)(6,2)(5)%
}

\NovelDisplay{M04291R01}{%
= \psi(I(1,I(2,0)+1))%
}

\NovelDisplay{M04292R01}{%
= \psi(M^4+M^3)%
}

\NovelDisplay{M04293R01}{%
(0)(1,1,1)(2,1,1)(3,1,1)(2,1,1)(3,1,1)(2,1,1)(3,1)(4,2,1)(5,2,1)(6,2,1)%
}
\NovelDisplay{M04293R02}{%
(5,2,1)(6,2,1)(5,2,1)(6,2)(5)%
}

\NovelDisplay{M04294R01}{%
= \psi(I(2,1))%
}

\NovelDisplay{M04295R01}{%
= \psi(M^4\times 2)%
}

\NovelDisplay{M04296R01}{%
(0)(1,1,1)(2,1,1)(3,1,1)(2,1,1)(3,1,1)(2,1,1)(3,1,1)%
}

\NovelDisplay{M04297R01}{%
= \psi(I(2,\omega))%
}

\NovelDisplay{M04298R01}{%
= \psi(M^4\times \omega)%
}

\NovelDisplay{M04299R01}{%
(0)(1,1,1)(2,1,1)(3,1,1)(2,1,1)(3,1,1)(2,1,1)(3,1,1)(2,1,1)(3,1)(2)%
}

\NovelDisplay{M04300R01}{%
= \psi(I(3,0))%
}

\NovelDisplay{M04301R01}{%
= \psi(M^5)%
}

第二参数的变化与第一参数的继续增加分别对应不同位置，不能把双参数I压成一个普通整数下标。

那晚，薛峰将今天的纸翻了三遍。

他不是还要把每一条记得更快，只是白天独立说明时，曾有一句话先从嘴里说出，随后才发现自己并未核过。他害怕这种熟练：同伴听见自己讲得稳，也许就会把没有说明的部分一并信了。

霍雨浩进门，见他在页边反复画线，便坐下来，让他把那一句原样说出。两人从原式回看，找到真正需要另核的位置，其他已确认的部分则仍好好留着。

``明天我先说这个。''薛峰道。

霍雨浩点头，没有让他把所有练过的东西全部推倒。薛峰将纸收好，第二天站到大家面前时，先指出昨日自己讲得太快的一处，再请后来的人把各自记录也拿来核对。

\NovelDisplay{M04302R01}{%
(0)(1,1,1)(2,1,1)(3,1,1)(2,1,1)(3,1,1)(2,1,1)(3,1,1)(2,1,1)(3,1,1)%
}
\NovelDisplay{M04302R02}{%
(2,1,1)(3,1)(2)%
}

\NovelDisplay{M04303R01}{%
= \psi(I(4,0))%
}

\NovelDisplay{M04304R01}{%
= \psi(M^6)%
}

\NovelDisplay{M04305R01}{%
(0)(1,1,1)(2,1,1)(3,1,1)(3)%
}

\NovelDisplay{M04306R01}{%
= \psi(I(\omega,0))%
}

\NovelDisplay{M04307R01}{%
= \psi(M^\omega)%
}

\NovelDisplay{M04308R01}{%
(0)(1,1,1)(2,1,1)(3,1,1)(3)(1,1,1)(2,1,1)(3,1,1)(3)%
}

\NovelDisplay{M04309R01}{%
= \psi(I(\omega,0)\times 2)%
}

\NovelDisplay{M04310R01}{%
= \psi(M^\omega+\psi_{M}(M^\omega))%
}

\NovelDisplay{M04311R01}{%
(0)(1,1,1)(2,1,1)(3,1,1)(3)(2,1)(3,2)%
}

\NovelDisplay{M04312R01}{%
= \psi(\Omega_{I(\omega,0)+1})%
}

\NovelDisplay{M04313R01}{%
= \psi(M^\omega+M)%
}

有限第一参数走向ω时，M侧的幂指数也进入ω层次，背景后的回接项仍不等于直接把背景翻倍。

霍雨浩发现，他们越是能熟练读出旧幂项，越容易在新背景后把它带错地方。

薛峰便让大家先说最外层的加号，再说后面一项。叶阳起初嫌这样慢，等自己一口读完后找不到括号的收处，才把纸退回去。

蔡须坤陪他重读，没有将自己的那份盖到他纸上。他指到能确认的位置，她才继续追问。

\NovelDisplay{M04314R01}{%
(0)(1,1,1)(2,1,1)(3,1,1)(3)(2,1)(3,2,1)(4,2,1)(5,2)(4)%
}

\NovelDisplay{M04315R01}{%
= \psi(I_{I(\omega,0)+1})%
}

\NovelDisplay{M04316R01}{%
= \psi(M^\omega+M^2)%
}

\NovelDisplay{M04317R01}{%
(0)(1,1,1)(2,1,1)(3,1,1)(3)(2,1)(3,2,1)(4,2,1)(5,2,1)(4,2,1)(5,2)(4)%
}

\NovelDisplay{M04318R01}{%
= \psi(I(1,I(\omega,0)+1))%
}

\NovelDisplay{M04319R01}{%
= \psi(M^\omega+M^3)%
}

\NovelDisplay{M04320R01}{%
(0)(1,1,1)(2,1,1)(3,1,1)(3)(2,1)(3,2,1)(4,2,1)(5,2,1)(4,2,1)(5,2,1)%
}
\NovelDisplay{M04320R02}{%
(4,2,1)(5,2)(4)%
}

\NovelDisplay{M04321R01}{%
= \psi(I(2,I(\omega,0)+1))%
}

\NovelDisplay{M04322R01}{%
= \psi(M^\omega+M^4)%
}

\NovelDisplay{M04323R01}{%
(0)(1,1,1)(2,1,1)(3,1,1)(3)(2,1)(3,2,1)(4,2,1)(5,2,1)(5)%
}

\NovelDisplay{M04324R01}{%
= \psi(I(\omega,1))%
}

\NovelDisplay{M04325R01}{%
= \psi(M^\omega\times 2)%
}

同一ω次幂背景上可再添不同低阶项，它们各自在M侧的位置要按完整式子保留。

\NovelDisplay{M04326R01}{%
(0)(1,1,1)(2,1,1)(3,1,1)(3)(2,1)(3,2,1)(4,2,1)(5,2,1)(5)(4,2)(5,3,1)%
}
\NovelDisplay{M04326R02}{%
(6,3,1)(7,3,1)(7)%
}

\NovelDisplay{M04327R01}{%
= \psi(I(\omega,2))%
}

\NovelDisplay{M04328R01}{%
= \psi(M^\omega\times 3)%
}

\NovelDisplay{M04329R01}{%
(0)(1,1,1)(2,1,1)(3,1,1)(3)(2,1,1)%
}

\NovelDisplay{M04330R01}{%
= \psi(I(\omega,\omega))%
}

\NovelDisplay{M04331R01}{%
= \psi(M^\omega\times \omega)%
}

\NovelDisplay{M04332R01}{%
(0)(1,1,1)(2,1,1)(3,1,1)(3)(2,1,1)(3,1)(2)%
}

\NovelDisplay{M04333R01}{%
= \psi(I(\omega+1,0))%
}

\NovelDisplay{M04334R01}{%
= \psi(M^{\omega+1})%
}

\NovelDisplay{M04335R01}{%
(0)(1,1,1)(2,1,1)(3,1,1)(3)(2,1,1)(3,1)(2,1,1)%
}

\NovelDisplay{M04336R01}{%
= \psi(I(\omega+1,0)\times \omega)%
}

\NovelDisplay{M04337R01}{%
= \psi(M^{\omega+1}+\psi_{M}(M^{\omega+1}))%
}

ω背景的倍数与指数后继在此相邻出现，读到看似相近的尾部，也必须先分清改变的是哪一层。

\NovelDisplay{M04338R01}{%
(0)(1,1,1)(2,1,1)(3,1,1)(3)(2,1,1)(3,1)(4,2)%
}

\NovelDisplay{M04339R01}{%
= \psi(\Omega_{I(\omega+1,0)+1})%
}

\NovelDisplay{M04340R01}{%
= \psi(M^{\omega+1}+M)%
}

\NovelDisplay{M04341R01}{%
(0)(1,1,1)(2,1,1)(3,1,1)(3)(2,1,1)(3,1)(4,2,1)(5,2,1)(6,2,1)(6)%
}

\NovelDisplay{M04342R01}{%
= \psi(I(\omega,I(\omega+1,0)+1))%
}

\NovelDisplay{M04343R01}{%
= \psi(M^{\omega+1}+M^\omega)%
}

\NovelDisplay{M04344R01}{%
(0)(1,1,1)(2,1,1)(3,1,1)(3)(2,1,1)(3,1)(4,2,1)(5,2,1)(6,2,1)(6)(5,2,1)%
}

\NovelDisplay{M04345R01}{%
= \psi(I(\omega,I(\omega+1,0)+\omega))%
}

\NovelDisplay{M04346R01}{%
= \psi(M^{\omega+1}+M^\omega\times \omega)%
}

\NovelDisplay{M04347R01}{%
(0)(1,1,1)(2,1,1)(3,1,1)(3)(2,1,1)(3,1)(4,2,1)(5,2,1)(6,2,1)(6)%
}
\NovelDisplay{M04347R02}{%
(5,2,1)(6,2)(5)%
}

\NovelDisplay{M04348R01}{%
= \psi(I(\omega+1,1))%
}

\NovelDisplay{M04349R01}{%
= \psi(M^{\omega+1}\times 2)%
}

新指数背景后再添ω次幂，与新指数背景本身的倍数不同，古月娜先指出加号所在的位置才请大家接读。

雨还没下，山上已经先闷了起来。

奥斯卡在后面连叫两声，蔡须坤都没有回头。她正想一口讲完一条长式，读到最内层，却突然忘了自己从哪一层开始。

卢本韦伸手将纸接住，让她先坐下。蔡须坤想说只差一点，抬头看见霍雨浩将水壶打开，最终还是把手从纸边松开。

那天下午，她重新回来，第一句便指出自己上午真正没接上的位置。

\NovelDisplay{M04350R01}{%
(0)(1,1,1)(2,1,1)(3,1,1)(3)(2,1,1)(3,1)(4,2,1)(5,2,1)(6,2,1)(6)(5,2,1)%
}
\NovelDisplay{M04350R02}{%
(6,2)(7,3,1)(8,3,1)(9,3,1)(9)(8,3,1)(9,3)(8)%
}

\NovelDisplay{M04351R01}{%
= \psi(I(\omega+1,2))%
}

\NovelDisplay{M04352R01}{%
= \psi(M^{\omega+1}\times 3)%
}

\NovelDisplay{M04353R01}{%
(0)(1,1,1)(2,1,1)(3,1,1)(3)(2,1,1)(3,1,1)%
}

\NovelDisplay{M04354R01}{%
= \psi(I(\omega+1,\omega))%
}

\NovelDisplay{M04355R01}{%
= \psi(M^{\omega+1}\times \omega)%
}

\NovelDisplay{M04356R01}{%
(0)(1,1,1)(2,1,1)(3,1,1)(3)(2,1,1)(3,1,1)(2,1,1)(3,1)(2)%
}

\NovelDisplay{M04357R01}{%
= \psi(I(\omega+2,0))%
}

\NovelDisplay{M04358R01}{%
= \psi(M^{\omega+2})%
}

\NovelDisplay{M04359R01}{%
(0)(1,1,1)(2,1,1)(3,1,1)(3)(2,1,1)(3,1,1)(2,1,1)(3,1,1)(2,1,1)(3,1)(2)%
}

\NovelDisplay{M04360R01}{%
= \psi(I(\omega+3,0))%
}

\NovelDisplay{M04361R01}{%
= \psi(M^{\omega+3})%
}

第二参数的倍数变化和第一参数的后继分别列出，不能因为几条矩阵很像便把同一种接法复制过去。

霍雨浩将这条原矩阵单独留下，让薛峰、叶阳与蔡须坤分别求一项，卢本韦只在三人写完后核对共同前段。

\NovelDisplay{M04362R01}{%
\text{展开第零项：}\quad(0)(1,1,1)(2,1,1)(3,1,1)(3)(2,1,1)(3,1)(4,2,1)(5,2,1)%
}
\NovelDisplay{M04362R02}{%
(6,2,1)(6)(5,2,1)(6,2)(7,3,1)(8,3,1)(9,3,1)(9)(8,3,1)(9,3)%
}

\NovelDisplay{M04363R01}{%
\text{展开第一项：}\quad(0)(1,1,1)(2,1,1)(3,1,1)(3)(2,1,1)(3,1)(4,2,1)(5,2,1)%
}
\NovelDisplay{M04363R02}{%
(6,2,1)(6)(5,2,1)(6,2)(7,3,1)(8,3,1)(9,3,1)(9)(8,3,1)(9,3)(7,3,1)(8,3,1)%
}
\NovelDisplay{M04363R03}{%
(9,3,1)(9)(8,3,1)(9,3)%
}

\NovelDisplay{M04364R01}{%
\text{展开第二项：}\quad(0)(1,1,1)(2,1,1)(3,1,1)(3)(2,1,1)(3,1)(4,2,1)(5,2,1)%
}
\NovelDisplay{M04364R02}{%
(6,2,1)(6)(5,2,1)(6,2)(7,3,1)(8,3,1)(9,3,1)(9)(8,3,1)(9,3)(7,3,1)(8,3,1)%
}
\NovelDisplay{M04364R03}{%
(9,3,1)(9)(8,3,1)(9,3)(7,3,1)(8,3,1)(9,3,1)(9)(8,3,1)(9,3)%
}

第零项、第一项和第二项都从同一原式出发，不能让刚读熟的幂指数替自己决定坏部范围。

薛峰最先指出自己多圈了哪一列，叶阳也把手里那份转过来，两人的错误没有被一张抄得更整齐的纸盖住。

有一日，夏夜星空让他们只带当天要核的几页上山，其余放在屋里。

小舞原以为包轻了，走起来便也能快。真正到了那一级，纸面虽然不再翻来翻去，心里的不确定却一点没少。她想回去取全本，古月娜问她，缺的是哪一处已经学过的说明。

她说不清，只觉得本子全在手边才安心。唐三看着自己同样收轻的包，忽然也把那句话低声说了一遍。

霍雨浩没有让两人空手硬猜。他请他们分别记下真正需要回查的内容，先将当前原式能独立确认的范围说清，回去以后再各自找那一页。

傍晚再看本子，唐三发现自己原以为必须带上整册才能处理的问题，其实只需回查其中一处条件；小舞则确实漏了另一段，要重新复核，纸少并没有替她把那段补好。

第二天上山，两人的包都比前一天厚了一点，各自装入需要的那几页。诺当看见，便问多带的是什么，小舞将自己的那处指出，第一次没有因为它比哥哥带的长就发脾气。

蔡须坤听过，也从封套里取出一张必须展开铺平的长页。她知道自己今天要查哪一处，放进包里时，心里反而比抱着整册四处翻找安稳。

山阶上的压力仍在，几个人却能更早分清，是身体需要休息，还是纸上还有条件没有接清。

后来再有谁问起这一处，小舞便先指自己带来的那张纸。她在重压下仍会停，也仍有弄不清的地方，却已不必先把整册翻乱，才能承认眼前这一处需要回查。

\NovelDisplay{M04365R01}{%
(0)(1,1,1)(2,1,1)(3,1,1)(3)(2,1,1)(3,1,1)(3)%
}

\NovelDisplay{M04366R01}{%
= \psi(I(\omega\times 2,0))%
}

\NovelDisplay{M04367R01}{%
= \psi(M^{\omega\times 2})%
}

\NovelDisplay{M04368R01}{%
(0)(1,1,1)(2,1,1)(3,1,1)(3)(2,1,1)(3,1,1)(3)(2,1,1)(3,1,1)(3)%
}

\NovelDisplay{M04369R01}{%
= \psi(I(\omega\times 3,0))%
}

\NovelDisplay{M04370R01}{%
= \psi(M^{\omega\times 3})%
}

\NovelDisplay{M04371R01}{%
(0)(1,1,1)(2,1,1)(3,1,1)(3)(3)%
}

\NovelDisplay{M04372R01}{%
= \psi(I(\psi(2),0))%
}

\NovelDisplay{M04373R01}{%
= \psi(M^{\psi(2)})%
}

\NovelDisplay{M04374R01}{%
(0)(1,1,1)(2,1,1)(3,1,1)(3,1)%
}

\NovelDisplay{M04375R01}{%
= \psi(I(\Omega,0))%
}

\NovelDisplay{M04376R01}{%
= \psi(M^{\psi_{M}(M)})%
}

第一参数的ω倍数、ψ值与Ω各有原稿给出的对应，跨度变大并不省去核对参数的步骤。

从双参数进入三参数时，小舞把那几个逗号数了好几次。

她曾经觉得参数越多，自己总能找到把短处绕过去的办法。真轮到面前，每个位置都要回答，才知道多出的格子并不会替她把原来的问题放走。

唐三将M侧放在她右边，指着幂指数，请她先说新背景。她说过一遍，又主动从第一参数重来，把新添的那个零也保留下来。

``这个零不能像矩阵末尾的补零那样省。''她说。

唐三点头，把自己刚才差点划去的位置重新描清。

\NovelDisplay{M04377R01}{%
(0)(1,1,1)(2,1,1)(3,1,1)(3,1)(1,1,1)(2,1,1)(3,1,1)(3,1)%
}

\NovelDisplay{M04378R01}{%
= \psi(I(I(\Omega,0),0))%
}

\NovelDisplay{M04379R01}{%
= \psi(M^{\psi_{M}(M^{\psi_{M}(M)})})%
}

\NovelDisplay{M04380R01}{%
(0)(1,1,1)(2,1,1)(3,1,1)(3,1)(2)%
}

\NovelDisplay{M04381R01}{%
= \psi(I(1,0,0))%
}

\NovelDisplay{M04382R01}{%
= \psi(M^M)%
}

\NovelDisplay{M04383R01}{%
(0)(1,1,1)(2,1,1)(3,1,1)(3,1)(2)(1,1,1)(2,1,1)(3,1,1)(3,1)(2)%
}

\NovelDisplay{M04384R01}{%
= \psi(I(1,0,0)+\psi_{I(1,0,0)}(I(1,0,0)))%
}

\NovelDisplay{M04385R01}{%
= \psi(M^M+\psi_{\psi_{M}(M^M)}(M^M))%
}

\NovelDisplay{M04386R01}{%
(0)(1,1,1)(2,1,1)(3,1,1)(3,1)(2,1)(3,2)%
}

\NovelDisplay{M04387R01}{%
= \psi(I(1,0,0)+\Omega_{\psi_{I(1,0,0)}(I(1,0,0))+1})%
}

\NovelDisplay{M04388R01}{%
= \psi(M^M+\psi_{M}(M^{\psi_{\psi_{M}(M^M)}(M^M)}+M))%
}

进入三参数I以后，M侧以M自身进入幂指数，后续折叠项也从这一新背景重新读起。

\NovelDisplay{M04389R01}{%
(0)(1,1,1)(2,1,1)(3,1,1)(3,1)(2,1)(3,2,1)(4,2,1)(5,2)(4)%
}

\NovelDisplay{M04390R01}{%
= \psi(I(1,0,0)+I_{\psi_{I(1,0,0)}(I(1,0,0))+1})%
}

\NovelDisplay{M04391R01}{%
= \psi(M^M+\psi_{M}(M^{\psi_{\psi_{M}(M^M)}(M^M)}+M^2))%
}

\NovelDisplay{M04392R01}{%
(0)(1,1,1)(2,1,1)(3,1,1)(3,1)(2,1)(3,2,1)(4,2,1)(5,2,1)(4,2,1)(5,2)(4)%
}

\NovelDisplay{M04393R01}{%
= \psi(I(1,0,0)+I(1,\psi_{I(1,0,0)}(I(1,0,0))+1))%
}

\NovelDisplay{M04394R01}{%
= \psi(M^M+\psi_{M}(M^{\psi_{\psi_{M}(M^M)}(M^M)}+M^3))%
}

\NovelDisplay{M04395R01}{%
(0)(1,1,1)(2,1,1)(3,1,1)(3,1)(2,1)(3,2,1)(4,2,1)(5,2,1)(4,2,1)(5,2,1)%
}
\NovelDisplay{M04395R02}{%
(4,2,1)(5,2)(4)%
}

\NovelDisplay{M04396R01}{%
= \psi(I(1,0,0)+I(2,\psi_{I(1,0,0)}(I(1,0,0))+1))%
}

\NovelDisplay{M04397R01}{%
= \psi(M^M+\psi_{M}(M^{\psi_{\psi_{M}(M^M)}(M^M)}+M^4))%
}

\NovelDisplay{M04398R01}{%
(0)(1,1,1)(2,1,1)(3,1,1)(3,1)(2,1)(3,2,1)(4,2,1)(5,2,1)(5)%
}

\NovelDisplay{M04399R01}{%
= \psi(I(1,0,0)+I(\omega,\psi_{I(1,0,0)}(I(1,0,0))+1))%
}

\NovelDisplay{M04400R01}{%
= \psi(M^M+\psi_{M}(M^{\psi_{\psi_{M}(M^M)}(M^M)}+M^\omega))%
}

相同作用对象后添加M的不同幂时，I侧改动的是不同参数位置，必须把整对括号和其中的逗号一起看清。

\NovelDisplay{M04401R01}{%
(0)(1,1,1)(2,1,1)(3,1,1)(3,1)(2,1)(3,2,1)(4,2,1)(5,2,1)(5,1)(1,1,1)%
}
\NovelDisplay{M04401R02}{%
(2,1,1)(3,1,1)(3,1)(2)%
}

\NovelDisplay{M04402R01}{%
= \psi(I(1,0,0)+I(\psi_{I(1,0,0)}(I(1,0,0)),1))%
}

\NovelDisplay{M04403R01}{%
= \psi(M^M+\psi_{M}(M^{\psi_{\psi_{M}(M^M)}(M^M)}\times 2))%
}

\NovelDisplay{M04404R01}{%
(0)(1,1,1)(2,1,1)(3,1,1)(3,1)(2,1)(3,2,1)(4,2,1)(5,2,1)(5,1)(3,2,1)%
}
\NovelDisplay{M04404R02}{%
(4,2,1)(5,2,1)(5,1)(1,1,1)(2,1,1)(3,1,1)(3,1)(2)%
}

\NovelDisplay{M04405R01}{%
= \psi(I(1,0,0)+I(\psi_{I(1,0,0)}(I(1,0,0)),1)\times 2)%
}

\NovelDisplay{M04406R01}{%
= \psi(M^M+\psi_{M}(M^{\psi_{\psi_{M}(M^M)}(M^M)}\times 2)\times 2)%
}

\NovelDisplay{M04407R01}{%
(0)(1,1,1)(2,1,1)(3,1,1)(3,1)(2,1)(3,2,1)(4,2,1)(5,2,1)(5,1)(4,2)(5,3)%
}

\NovelDisplay{M04408R01}{%
= \psi(I(1,0,0)+\Omega_{I(\psi_{I(1,0,0)}(I(1,0,0)),1)+1})%
}

\NovelDisplay{M04409R01}{%
= \psi(M^M+\psi_{M}(M^{\psi_{\psi_{M}(M^M)}(M^M)}\times 2+M))%
}

\NovelDisplay{M04410R01}{%
(0)(1,1,1)(2,1,1)(3,1,1)(3,1)(2,1)(3,2,1)(4,2,1)(5,2,1)(5,1)(4,2)(5,3,1)%
}
\NovelDisplay{M04410R02}{%
(6,3,1)(7,3,1)(7,1)(1,1,1)(2,1,1)(3,1,1)(3,1)(2)%
}

\NovelDisplay{M04411R01}{%
= \psi(I(1,0,0)+I(\psi_{I(1,0,0)}(I(1,0,0)),2))%
}

\NovelDisplay{M04412R01}{%
= \psi(M^M+\psi_{M}(M^{\psi_{\psi_{M}(M^M)}(M^M)}\times 3))%
}

同一个复杂参数的倍数与其下标后继分开表示，最外面的倍数不能移入内层的幂指数。

再遇到整个对象的倍数时，唐三先遮住右式，古月娜请小舞从末列回找生成位置；三个取项编号分别算完以后，才把结果并排核对。

\NovelDisplay{M04413R01}{%
\text{展开第零项：}\quad(0)(1,1,1)(2,1,1)(3,1,1)(3,1)(2,1)(3,2,1)(4,2,1)(5,2,1)%
}
\NovelDisplay{M04413R02}{%
(5,1)(3,2,1)(4,2,1)(5,2,1)(5,1)(1,1,1)(2,1,1)(3,1,1)(3,1)%
}

\NovelDisplay{M04414R01}{%
\text{展开第一项：}\quad(0)(1,1,1)(2,1,1)(3,1,1)(3,1)(2,1)(3,2,1)(4,2,1)(5,2,1)%
}
\NovelDisplay{M04414R02}{%
(5,1)(3,2,1)(4,2,1)(5,2,1)(5,1)(1,1,1)(2,1,1)(3,1,1)(3,1)(1,1,1)%
}
\NovelDisplay{M04414R03}{%
(2,1,1)(3,1,1)(3,1)%
}

\NovelDisplay{M04415R01}{%
\text{展开第二项：}\quad(0)(1,1,1)(2,1,1)(3,1,1)(3,1)(2,1)(3,2,1)(4,2,1)(5,2,1)%
}
\NovelDisplay{M04415R02}{%
(5,1)(3,2,1)(4,2,1)(5,2,1)(5,1)(1,1,1)(2,1,1)(3,1,1)(3,1)(1,1,1)(2,1,1)%
}
\NovelDisplay{M04415R03}{%
(3,1,1)(3,1)(1,1,1)(2,1,1)(3,1,1)(3,1)%
}

这里的三项各保留自己的生成份数，较长的一份也仍低于原矩阵，不能只数列数比较序数大小。

小舞从第一个不同位置比过，才把原来写在页角的``越长越大''划掉，另留下一句请自己回查的位置。

那几日，蔡须坤和卢本韦常常互换位置。

先由卢本韦从I侧提出问题，蔡须坤接M侧，再由霍雨浩检查她们是否在谈同一个作用范围。薛峰不时找到更早漏下的一列，便请两人暂时退回原矩阵。

蔡须坤起初有些不甘心，觉得长式已读到这么深，还要重新数列。等她自己沿父链追到真正的坏根，那点不甘才慢慢散了。

表示可以很长，落脚的位置却不能靠熟悉代替。

\NovelDisplay{M04416R01}{%
(0)(1,1,1)(2,1,1)(3,1,1)(3,1)(2,1)(3,2,1)(4,2,1)(5,2,1)(5,1)(4,2,1)%
}

\NovelDisplay{M04417R01}{%
= \psi(I(1,0,0)+I(\psi_{I(1,0,0)}(I(1,0,0)),\omega))%
}

\NovelDisplay{M04418R01}{%
= \psi(M^M+\psi_{M}(M^{\psi_{\psi_{M}(M^M)}(M^M)}\times \omega))%
}

\NovelDisplay{M04419R01}{%
(0)(1,1,1)(2,1,1)(3,1,1)(3,1)(2,1)(3,2,1)(4,2,1)(5,2,1)(5,1)(4,2,1)(5,2)%
}

\NovelDisplay{M04420R01}{%
= \psi(I(1,0,0)+I(\psi_{I(1,0,0)}(I(1,0,0)),\Omega_{\psi_{I(1,0,0)}(I(1,0,0))+1}))%
}

\NovelDisplay{M04421R01}{%
= \psi(M^M+\psi_{M}(M^{\psi_{\psi_{M}(M^M)}(M^M)}\times \psi_{M}(M^{\psi_{\psi_{M}(M^M)}(M^M)}+M)))%
}

\NovelDisplay{M04422R01}{%
(0)(1,1,1)(2,1,1)(3,1,1)(3,1)(2,1)(3,2,1)(4,2,1)(5,2,1)(5,1)(4,2,1)%
}
\NovelDisplay{M04422R02}{%
(5,2)(4)%
}

\NovelDisplay{M04423R01}{%
= \psi(I(1,0,0)+I(\psi_{I(1,0,0)}(I(1,0,0))+1,0))%
}

\NovelDisplay{M04424R01}{%
= \psi(M^M+\psi_{\psi_{M}(M^M)}(M^M+1))%
}

\NovelDisplay{M04425R01}{%
(0)(1,1,1)(2,1,1)(3,1,1)(3,1)(2,1)(3,2,1)(4,2,1)(5,2,1)(5,1)%
}
\NovelDisplay{M04425R02}{%
(4,2,1)(5,2,1)%
}

\NovelDisplay{M04426R01}{%
= \psi(I(1,0,0)+I(\psi_{I(1,0,0)}(I(1,0,0))+1,\omega))%
}

\NovelDisplay{M04427R01}{%
= \psi(M^M+\psi_{M}(M^{\psi_{\psi_{M}(M^M)}(M^M)+1}\times  \omega))%
}

参数的后继、ω与倍数在这里进入不同结构，霍雨浩让每个人先独立圈清作用范围，再与同伴核对。

连续几日的闷热过去，山下落了一场雨。

他们没有上阶。蔡兴坤抱着黑子鸡坐在门槛里，奥斯卡将漏水处的盆挪开，叶阳把受潮的长页一张张晾起。诺当想帮着搬，先问清页次，才敢将手伸过去。

蔡须坤坐在桌前，忽然发现自己听得见大家做事的声音，却看不进面前那个下标。她先以为只是困，闭上眼，脑子里仍是越来越深的括号，怎样也不愿停。

卢本韦没有让她马上换一个更短的例子。她只问，若今天不继续，最怕丢掉的究竟是什么。

蔡须坤望着纸，过了很久，说怕明天又得从头来。

霍雨浩将已经由她独立核过的几页放在一旁，让她先指出真正还不清楚的位置。那处不小，也不再像最初那样盖住所有东西。

她将这张夹在最外面，自己起身去帮叶阳压平晾好的纸。雨一直下到晚饭前，桌上的问题没有被谁偷着替她做完，她却终于能把那口一直压着的气缓缓呼出去。

\NovelDisplay{M04428R01}{%
(0)(1,1,1)(2,1,1)(3,1,1)(3,1)(2,1)(3,2,1)(4,2,1)(5,2,1)(5,1)(4,2,1)%
}
\NovelDisplay{M04428R02}{%
(5,2,1)(4,2,1)(5,2)(4)%
}

\NovelDisplay{M04429R01}{%
= \psi(I(1,0,0)+I(\psi_{I(1,0,0)}(I(1,0,0))+2,0))%
}

\NovelDisplay{M04430R01}{%
= \psi(M^M+\psi_{\psi_{M}(M^M)}(M^M+2))%
}

\NovelDisplay{M04431R01}{%
(0)(1,1,1)(2,1,1)(3,1,1)(3,1)(2,1)(3,2,1)(4,2,1)(5,2,1)(5,1)%
}
\NovelDisplay{M04431R02}{%
(4,2,1)(5,2,1)(5)%
}

\NovelDisplay{M04432R01}{%
= \psi(I(1,0,0)+I(\psi_{I(1,0,0)}(I(1,0,0))+\omega,0))%
}

\NovelDisplay{M04433R01}{%
= \psi(M^M+\psi_{\psi_{M}(M^M)}(M^M+\omega))%
}

\NovelDisplay{M04434R01}{%
(0)(1,1,1)(2,1,1)(3,1,1)(3,1)(2,1)(3,2,1)(4,2,1)(5,2,1)(5,1)(4,2,1)%
}
\NovelDisplay{M04434R02}{%
(5,2,1)(5,1)(1,1,1)(2,1,1)(3,1,1)(3,1)(2)%
}

\NovelDisplay{M04435R01}{%
= \psi(I(1,0,0)+I(\psi_{I(1,0,0)}(I(1,0,0))\times 2,0))%
}

\NovelDisplay{M04436R01}{%
= \psi(M^M+\psi_{\psi_{M}(M^M)}(M^M+\psi_{\psi_{M}(M^M)}(M^M)))%
}

\NovelDisplay{M04437R01}{%
(0)(1,1,1)(2,1,1)(3,1,1)(3,1)(2,1)(3,2,1)(4,2,1)(5,2,1)(5,2)%
}

\NovelDisplay{M04438R01}{%
= \psi(I(1,0,0)+I(\Omega_{\psi_{I(1,0,0)}(I(1,0,0))+1},0))%
}

\NovelDisplay{M04439R01}{%
= \psi(M^M+\psi_{\psi_{M}(M^M)}(M^M+\psi_{M}(M^{\psi_{\psi_{M}(M^M)}(M^M)}+M)))%
}

当第一参数本身已经含ψ时，追加的一、ω和折叠项仍有各自位置，不能拿同一串熟悉字母遮住区别。

\NovelDisplay{M04440R01}{%
(0)(1,1,1)(2,1,1)(3,1,1)(3,1)(2,1)(3,2,1)(4,2,1)(5,2,1)(5,2)(3,2,1)%
}
\NovelDisplay{M04440R02}{%
(4,2,1)(5,2,1)(5,2)%
}

\NovelDisplay{M04441R01}{%
= \psi(I(1,0,0)+I(I(\Omega_{\psi_{I(1,0,0)}(I(1,0,0))+1},0),0))%
}

\NovelDisplay{M04442R01}{%
= \psi(M^M+\psi_{\psi_{M}(M^M)}(M^M+\psi_{\psi_{M}(M^M)}(M^M+\psi_{M}(M^{\psi_{\psi_{M}(M^M)}(M^M)}+M))))%
}

\NovelDisplay{M04443R01}{%
(0)(1,1,1)(2,1,1)(3,1,1)(3,1)(2,1)(3,2,1)(4,2,1)(5,2,1)(5,2)(4)%
}

\NovelDisplay{M04444R01}{%
= \psi(I(1,0,0)\times 2)%
}

\NovelDisplay{M04445R01}{%
= \psi(M^M+\psi_{M}(M^M))%
}

\NovelDisplay{M04446R01}{%
(0)(1,1,1)(2,1,1)(3,1,1)(3,1)(2,1)(3,2,1)(4,2,1)(5,2,1)(5,2)(4,2)(5,3)%
}

\NovelDisplay{M04447R01}{%
= \psi(I(1,0,0)\times 2+\Omega_{\psi_{I(1,0,0)}(I(1,0,0)\times 2)+1})%
}

\NovelDisplay{M04448R01}{%
= \psi(M^M+\psi_{M}(M^M)+\psi_{M}(M^{\psi_{\psi_{M}(M^M)}(M^M+\psi_{M}(M^M))}+M))%
}

\NovelDisplay{M04449R01}{%
(0)(1,1,1)(2,1,1)(3,1,1)(3,1)(2,1)(3,2,1)(4,2,1)(5,2,1)(5,2)(4,2)(5,3,1)%
}
\NovelDisplay{M04449R02}{%
(6,3,1)(7,3,1)(7,3)(6)%
}

\NovelDisplay{M04450R01}{%
= \psi(I(1,0,0)\times 3)%
}

\NovelDisplay{M04451R01}{%
= \psi(M^M+\psi_{M}(M^M)\times 2)%
}

重复套入I参数与整个三参数I项的倍数分别处理，M侧新增的回接项也不能因此漏掉。

小舞有一回独自去过石阶前。

她想先把今天反复练过的几级走清，再叫唐三来看。等脚底的压力真正压来，纸上的一处停顿却变得怎么也接不上，她站在那里，越想越急。

后来赶到的是叶阳。他没有问她为什么不等大家，只先将纸移到她能看清的位置。小舞从头说过，才发现自己拿错了上一页的背景。

叶阳陪她退回石桩下。唐三来时，她把这件事先说出来，没等哥哥开口。

\NovelDisplay{M04452R01}{%
(0)(1,1,1)(2,1,1)(3,1,1)(3,1)(2,1,1)%
}

\NovelDisplay{M04453R01}{%
= \psi(I(1,0,0)\times \omega)%
}

\NovelDisplay{M04454R01}{%
= \psi(M^M+\psi_{\psi_{M}(M^M+M)}(1))%
}

\NovelDisplay{M04455R01}{%
(0)(1,1,1)(2,1,1)(3,1,1)(3,1)(2,1,1)(3,1)(4,2)%
}

\NovelDisplay{M04456R01}{%
= \psi(\Omega_{I(1,0,0)+1})%
}

\NovelDisplay{M04457R01}{%
= \psi(M^M+M)%
}

\NovelDisplay{M04458R01}{%
(0)(1,1,1)(2,1,1)(3,1,1)(3,1)(2,1,1)(3,1)(4,2,1)(5,2,1)(6,2)(5)%
}

\NovelDisplay{M04459R01}{%
= \psi(I_{I(1,0,0)+1})%
}

\NovelDisplay{M04460R01}{%
= \psi(M^M+M^2)%
}

\NovelDisplay{M04461R01}{%
(0)(1,1,1)(2,1,1)(3,1,1)(3,1)(2,1,1)(3,1)(4,2,1)(5,2,1)(6,2,1)%
}
\NovelDisplay{M04461R02}{%
(5,2,1)(6,2)(5)%
}

\NovelDisplay{M04462R01}{%
= \psi(I(1,I(1,0,0)+1))%
}

\NovelDisplay{M04463R01}{%
= \psi(M^M+M^3)%
}

在M自幂背景后添加M、平方或立方各有对应，短式能帮助回看位置，却不能替代原矩阵的计算。

\NovelDisplay{M04464R01}{%
(0)(1,1,1)(2,1,1)(3,1,1)(3,1)(2,1,1)(3,1)(4,2,1)(5,2,1)(6,2,1)(6)%
}

\NovelDisplay{M04465R01}{%
= \psi(I(\omega,I(1,0,0)+1))%
}

\NovelDisplay{M04466R01}{%
= \psi(M^M+M^\omega)%
}

\NovelDisplay{M04467R01}{%
(0)(1,1,1)(2,1,1)(3,1,1)(3,1)(2,1,1)(3,1)(4,2,1)(5,2,1)(6,2,1)(6,1)%
}

\NovelDisplay{M04468R01}{%
= \psi(I(\Omega,I(1,0,0)+1))%
}

\NovelDisplay{M04469R01}{%
= \psi(M^M+M^{\psi_{M}(M)})%
}

\NovelDisplay{M04470R01}{%
(0)(1,1,1)(2,1,1)(3,1,1)(3,1)(2,1,1)(3,1)(4,2,1)(5,2,1)(6,2,1)(6,1)%
}
\NovelDisplay{M04470R02}{%
(1,1,1)(2,1,1)(3,1,1)(3,1)(2)%
}

\NovelDisplay{M04471R01}{%
= \psi(I(\psi_{I(1,0,0)}(I(1,0,0)),I(1,0,0)+1))%
}

\NovelDisplay{M04472R01}{%
= \psi(M^M+M^{\psi_{\psi_{M}(M^M)}(M^M)})%
}

\NovelDisplay{M04473R01}{%
(0)(1,1,1)(2,1,1)(3,1,1)(3,1)(2,1,1)(3,1)(4,2,1)(5,2,1)(6,2,1)(6,1)%
}
\NovelDisplay{M04473R02}{%
(1,1,1)(2,1,1)(3,1,1)(3,1)(2,1,1)(3,1)(4,2)%
}

\NovelDisplay{M04474R01}{%
= \psi(I(\psi_{I(1,0,0)}(\Omega_{I(1,0,0)+1}),I(1,0,0)+1))%
}

\NovelDisplay{M04475R01}{%
= \psi(M^M+M^{\psi_{\psi_{M}(M^M)}(M^M+M)})%
}

第一参数继续变复杂时，M侧指数里的折叠也随之改变，先守住外层加号，再逐个读清指数。

\NovelDisplay{M04476R01}{%
(0)(1,1,1)(2,1,1)(3,1,1)(3,1)(2,1,1)(3,1)(4,2,1)(5,2,1)(6,2,1)(6,1)%
}
\NovelDisplay{M04476R02}{%
(1,1,1)(2,1,1)(3,1,1)(3,1)(2,1,1)(3,1)(4,2,1)(5,2,1)(6,2,1)(6,1)%
}

\NovelDisplay{M04477R01}{%
= \psi(I(\psi_{I(1,0,0)}(I(\Omega,I(1,0,0)+1)),I(1,0,0)+1))%
}

\NovelDisplay{M04478R01}{%
= \psi(M^M+M^{\psi_{\psi_{M}(M^M)}(M^M+M^{\psi_{M}(M)})})%
}

\NovelDisplay{M04479R01}{%
(0)(1,1,1)(2,1,1)(3,1,1)(3,1)(2,1,1)(3,1)(4,2,1)(5,2,1)(6,2,1)(6,1)(2)%
}

\NovelDisplay{M04480R01}{%
= \psi(I(I(1,0,0),1))%
}

\NovelDisplay{M04481R01}{%
= \psi(M^M+M^{\psi_{M}(M^M)})%
}

\NovelDisplay{M04482R01}{%
(0)(1,1,1)(2,1,1)(3,1,1)(3,1)(2,1,1)(3,1)(4,2,1)(5,2,1)(6,2,1)(6,1)(4,2)%
}

\NovelDisplay{M04483R01}{%
= \psi(I(I(1,0,0),1)+\Omega_{I(1,0,0)+1})%
}

\NovelDisplay{M04484R01}{%
= \psi(M^M+M^{\psi_{M}(M^M)}+\psi_{M}(M^M+M))%
}

\NovelDisplay{M04485R01}{%
(0)(1,1,1)(2,1,1)(3,1,1)(3,1)(2,1,1)(3,1)(4,2,1)(5,2,1)(6,2,1)(6,1)%
}
\NovelDisplay{M04485R02}{%
(4,2,1)(5,2,1)(6,2,1)(6,1)%
}

\NovelDisplay{M04486R01}{%
= \psi(I(I(1,0,0),1)+I(\Omega,I(1,0,0)+1))%
}

\NovelDisplay{M04487R01}{%
= \psi(M^M+M^{\psi_{M}(M^M)}+\psi_{M}(M^M+M^{\psi_{M}(M)}))%
}

新的嵌套参数与背景后的附加项并列，薛峰将它们各自指出，避免把两次作用合并成一次。

这一条只有二十一列，第二项却会长到六十列。唐三望着空纸，先问是不是哪里重复多了。

蔡须坤指向末列\allowbreak{}\NovelInline{I06091}{(6,1)}\allowbreak{}：最高非零位置在第二行，沿这一行的父链回找，坏根落到第一列；第一至第二十列整段都要纳入生成，第一行的差为六，实际受提升的位置仍依祖先条件判断。

她没有将那二十列换成一个临时代号。卢本韦将每一列铺全，请三个人各自从同一原式开始。

\NovelDisplay{M04488R01}{%
\text{展开第零项：}\quad(0)(1,1,1)(2,1,1)(3,1,1)(3,1)(2,1,1)(3,1)(4,2,1)(5,2,1)%
}
\NovelDisplay{M04488R02}{%
(6,2,1)(6,1)(1,1,1)(2,1,1)(3,1,1)(3,1)(2,1,1)(3,1)(4,2,1)(5,2,1)(6,2,1)%
}

\NovelDisplay{M04489R01}{%
\text{展开第一项：}\quad(0)(1,1,1)(2,1,1)(3,1,1)(3,1)(2,1,1)(3,1)(4,2,1)(5,2,1)%
}
\NovelDisplay{M04489R02}{%
(6,2,1)(6,1)(1,1,1)(2,1,1)(3,1,1)(3,1)(2,1,1)(3,1)(4,2,1)(5,2,1)(6,2,1)%
}
\NovelDisplay{M04489R03}{%
(6)(7,1,1)(8,1,1)(9,1,1)(9,1)(8,1,1)(9,1)(10,2,1)(11,2,1)(12,2,1)(12,1)%
}
\NovelDisplay{M04489R04}{%
(7,1,1)(8,1,1)(9,1,1)(9,1)(8,1,1)(9,1)(10,2,1)(11,2,1)(12,2,1)%
}

\NovelDisplay{M04490R01}{%
\text{展开第二项：}\quad(0)(1,1,1)(2,1,1)(3,1,1)(3,1)(2,1,1)(3,1)(4,2,1)(5,2,1)%
}
\NovelDisplay{M04490R02}{%
(6,2,1)(6,1)(1,1,1)(2,1,1)(3,1,1)(3,1)(2,1,1)(3,1)(4,2,1)(5,2,1)(6,2,1)%
}
\NovelDisplay{M04490R03}{%
(6)(7,1,1)(8,1,1)(9,1,1)(9,1)(8,1,1)(9,1)(10,2,1)(11,2,1)(12,2,1)(12,1)%
}
\NovelDisplay{M04490R04}{%
(7,1,1)(8,1,1)(9,1,1)(9,1)(8,1,1)(9,1)(10,2,1)(11,2,1)(12,2,1)(12)%
}
\NovelDisplay{M04490R05}{%
(13,1,1)(14,1,1)(15,1,1)(15,1)(14,1,1)(15,1)(16,2,1)(17,2,1)(18,2,1)%
}
\NovelDisplay{M04490R06}{%
(18,1)(13,1,1)(14,1,1)(15,1,1)(15,1)(14,1,1)(15,1)(16,2,1)%
}
\NovelDisplay{M04490R07}{%
(17,2,1)(18,2,1)%
}

第零项二十列、第一项四十列、第二项六十列，坏根较远便会生成很长的段，却不改变取项必须低于原式的比较。

写到最后，唐三的手指已经发酸。小舞替他压住纸，没有替他添末列；等他自己写完，才请他回头指出第一处不同。

诺当发现，自己已经能听出一些以前从不敢问的问题。

他问同一个对象的倍数为什么不能只在M侧照字面乘二。蔡须坤没有立即给出一条通用口诀，而是将本组几式并排，让他指出讲义实际保留的背景与后来加上的折叠项。

他说得很慢，最后却自己看见了不同作用范围。薛峰把纸交回他手里，请他仍从原矩阵求项，再核纸上的理解。

\NovelDisplay{M04491R01}{%
(0)(1,1,1)(2,1,1)(3,1,1)(3,1)(2,1,1)(3,1)(4,2,1)(5,2,1)(6,2,1)(6,1)%
}
\NovelDisplay{M04491R02}{%
(4,2,1)(5,2,1)(6,2,1)(6,1)(2)%
}

\NovelDisplay{M04492R01}{%
= \psi(I(I(1,0,0),1)\times 2)%
}

\NovelDisplay{M04493R01}{%
= \psi(M^M+M^{\psi_{M}(M^M)}+\psi_{M}(M^M+M^{\psi_{M}(M^M)}))%
}

\NovelDisplay{M04494R01}{%
(0)(1,1,1)(2,1,1)(3,1,1)(3,1)(2,1,1)(3,1)(4,2,1)(5,2,1)(6,2,1)(6,1)%
}
\NovelDisplay{M04494R02}{%
(5,2)(6,3)%
}

\NovelDisplay{M04495R01}{%
= \psi(\Omega_{I(I(1,0,0),1)+1})%
}

\NovelDisplay{M04496R01}{%
= \psi(M^M+M^{\psi_{M}(M^M)}+M)%
}

\NovelDisplay{M04497R01}{%
(0)(1,1,1)(2,1,1)(3,1,1)(3,1)(2,1,1)(3,1)(4,2,1)(5,2,1)(6,2,1)(6,1)(5,2)%
}
\NovelDisplay{M04497R02}{%
(6,3,1)(7,3,1)(8,3,1)(8,1)%
}

\NovelDisplay{M04498R01}{%
= \psi(I(\Omega,I(I(1,0,0),1)+1))%
}

\NovelDisplay{M04499R01}{%
= \psi(M^M+M^{\psi_{M}(M^M)}+M^{\psi_{M}(M)})%
}

\NovelDisplay{M04500R01}{%
(0)(1,1,1)(2,1,1)(3,1,1)(3,1)(2,1,1)(3,1)(4,2,1)(5,2,1)(6,2,1)(6,1)(5,2)%
}
\NovelDisplay{M04500R02}{%
(6,3,1)(7,3,1)(8,3,1)(8,1)(2)%
}

\NovelDisplay{M04501R01}{%
= \psi(I(I(1,0,0),2))%
}

\NovelDisplay{M04502R01}{%
= \psi(M^M+M^{\psi_{M}(M^M)}\times 2)%
}

整个对象的倍数、Ω下标和第二参数的增加在此分开，不能只凭M侧反复出现同一背景便认作重复。

\NovelDisplay{M04503R01}{%
(0)(1,1,1)(2,1,1)(3,1,1)(3,1)(2,1,1)(3,1)(4,2,1)(5,2,1)(6,2,1)(6,1)(5,2)%
}
\NovelDisplay{M04503R02}{%
(6,3,1)(7,3,1)(8,3,1)(8,1)(7,3)(8,4,1)(9,4,1)(10,4,1)(10,1)(2)%
}

\NovelDisplay{M04504R01}{%
= \psi(I(I(1,0,0),3))%
}

\NovelDisplay{M04505R01}{%
= \psi(M^M+M^{\psi_{M}(M^M)}\times 3)%
}

\NovelDisplay{M04506R01}{%
(0)(1,1,1)(2,1,1)(3,1,1)(3,1)(2,1,1)(3,1)(4,2,1)(5,2,1)(6,2,1)%
}
\NovelDisplay{M04506R02}{%
(6,1)(5,2,1)%
}

\NovelDisplay{M04507R01}{%
= \psi(I(I(1,0,0),\omega))%
}

\NovelDisplay{M04508R01}{%
= \psi(M^M+M^{\psi_{M}(M^M)}\times \omega)%
}

\NovelDisplay{M04509R01}{%
(0)(1,1,1)(2,1,1)(3,1,1)(3,1)(2,1,1)(3,1)(4,2,1)(5,2,1)(6,2,1)(6,1)%
}
\NovelDisplay{M04509R02}{%
(5,2,1)(6,1)(2)%
}

\NovelDisplay{M04510R01}{%
= \psi(I(I(1,0,0),I(1,0,0)))%
}

\NovelDisplay{M04511R01}{%
= \psi(M^M+M^{\psi_{M}(M^M)}\times \psi_{M}(M^M))%
}

\NovelDisplay{M04512R01}{%
(0)(1,1,1)(2,1,1)(3,1,1)(3,1)(2,1,1)(3,1)(4,2,1)(5,2,1)(6,2,1)%
}
\NovelDisplay{M04512R02}{%
(6,1)(5,2,1)(6,2)%
}

\NovelDisplay{M04513R01}{%
= \psi(I(I(1,0,0),\Omega_{I(1,0,0)+1}))%
}

\NovelDisplay{M04514R01}{%
= \psi(M^M+M^{\psi_{M}(M^M)}\times \psi_{M}(M^M+M))%
}

第二参数的有限倍数、ω和复杂对象保留各自对应，卢本韦没有把参数变化误写成整个式子的外乘。

\NovelDisplay{M04515R01}{%
(0)(1,1,1)(2,1,1)(3,1,1)(3,1)(2,1,1)(3,1)(4,2,1)(5,2,1)(6,2,1)(6,1)%
}
\NovelDisplay{M04515R02}{%
(5,2,1)(6,2)(5)%
}

\NovelDisplay{M04516R01}{%
= \psi(I(I(1,0,0)+1,0))%
}

\NovelDisplay{M04517R01}{%
= \psi(M^M+M^{\psi_{M}(M^M)+1})%
}

\NovelDisplay{M04518R01}{%
(0)(1,1,1)(2,1,1)(3,1,1)(3,1)(2,1,1)(3,1)(4,2,1)(5,2,1)(6,2,1)(6,1)%
}
\NovelDisplay{M04518R02}{%
(5,2,1)(6,2)(5,2,1)%
}

\NovelDisplay{M04519R01}{%
= \psi(I(I(1,0,0)+1,0)\times \omega)%
}

\NovelDisplay{M04520R01}{%
= \psi(M^M+M^{\psi_{M}(M^M)+1}\times \omega)%
}

\NovelDisplay{M04521R01}{%
(0)(1,1,1)(2,1,1)(3,1,1)(3,1)(2,1,1)(3,1)(4,2,1)(5,2,1)(6,2,1)(6,1)%
}
\NovelDisplay{M04521R02}{%
(5,2,1)(6,2)(7,3)%
}

\NovelDisplay{M04522R01}{%
= \psi(\Omega_{I(I(1,0,0)+1,0)+1})%
}

\NovelDisplay{M04523R01}{%
= \psi(M^M+M^{\psi_{M}(M^M)+1}+M)%
}

\NovelDisplay{M04524R01}{%
(0)(1,1,1)(2,1,1)(3,1,1)(3,1)(2,1,1)(3,1)(4,2,1)(5,2,1)(6,2,1)(6,1)%
}
\NovelDisplay{M04524R02}{%
(5,2,1)(6,2,1)%
}

\NovelDisplay{M04525R01}{%
= \psi(I(I(1,0,0)+1,\omega))%
}

\NovelDisplay{M04526R01}{%
= \psi(M^M+M^{\psi_{M}(M^M)+1}\times \omega)%
}

第一参数后继进入指数时，后面的ω倍数与Ω下标加项仍各有范围，不能在口述时抹去括号。

练到最后几天，大家开始将纸分成几种。

已经由自己核清的放在里层，仍有疑问的放在外面；书写太长、需要展开铺平的，另外夹一张纸条。古月娜不许谁因为嫌重，就把复杂那份从包里抽走。

奥斯卡听见，笑着说那自己就少带一件衣裳。卢本韦回头看他，他立刻补了一句，衣裳也带，吃饭也不能少。

笑过以后，蔡须坤把最末一张长页认真折好。这一次，折口避开了下标。

\NovelDisplay{M04527R01}{%
(0)(1,1,1)(2,1,1)(3,1,1)(3,1)(2,1,1)(3,1)(4,2,1)(5,2,1)(6,2,1)(6,1)%
}
\NovelDisplay{M04527R02}{%
(5,2,1)(6,2,1)(5,2,1)(6,2)(5)%
}

\NovelDisplay{M04528R01}{%
= \psi(I(I(1,0,0)+2,0))%
}

\NovelDisplay{M04529R01}{%
= \psi(M^M+M^{\psi_{M}(M^M)+2})%
}

\NovelDisplay{M04530R01}{%
(0)(1,1,1)(2,1,1)(3,1,1)(3,1)(2,1,1)(3,1)(4,2,1)(5,2,1)(6,2,1)(6,1)%
}
\NovelDisplay{M04530R02}{%
(5,2,1)(6,2,1)(6,1)%
}

\NovelDisplay{M04531R01}{%
= \psi(I(I(1,0,0)+\Omega,0))%
}

\NovelDisplay{M04532R01}{%
= \psi(M^M+M^{\psi_{M}(M^M)+\psi_{M}(M)})%
}

\NovelDisplay{M04533R01}{%
(0)(1,1,1)(2,1,1)(3,1,1)(3,1)(2,1,1)(3,1)(4,2,1)(5,2,1)(6,2,1)(6,1)%
}
\NovelDisplay{M04533R02}{%
(5,2,1)(6,2,1)(6,1)(2)%
}

\NovelDisplay{M04534R01}{%
= \psi(I(I(1,0,0)\times 2,0))%
}

\NovelDisplay{M04535R01}{%
= \psi(M^M+M^{\psi_{M}(M^M)\times 2})%
}

\NovelDisplay{M04536R01}{%
(0)(1,1,1)(2,1,1)(3,1,1)(3,1)(2,1,1)(3,1)(4,2,1)(5,2,1)(6,2,1)(6,1)(6)%
}

\NovelDisplay{M04537R01}{%
= \psi(I(I(1,0,0)\times \omega,0))%
}

\NovelDisplay{M04538R01}{%
= \psi(M^M+M^{\psi_{\psi_{M}(M^M+M)}(1)})%
}

这里并列第一参数的有限尾部、Ω尾部、倍数与ω倍数，蔡须坤依次指出M侧幂指数中的不同对象。

\NovelDisplay{M04539R01}{%
(0)(1,1,1)(2,1,1)(3,1,1)(3,1)(2,1,1)(3,1)(4,2,1)(5,2,1)(6,2,1)(6,1)(7,2)%
}

\NovelDisplay{M04540R01}{%
= \psi(I(\psi_{\Omega_{I(1,0,0)+1}}(\Omega_{I(1,0,0)+1}),0))%
}

\NovelDisplay{M04541R01}{%
= \psi(M^M+M^{\psi_{\psi_{M}(M^M+M)}(M^M+M)})%
}

\NovelDisplay{M04542R01}{%
(0)(1,1,1)(2,1,1)(3,1,1)(3,1)(2,1,1)(3,1)(4,2,1)(5,2,1)(6,2,1)(6,2)%
}

\NovelDisplay{M04543R01}{%
= \psi(I(\Omega_{I(1,0,0)+1},0))%
}

\NovelDisplay{M04544R01}{%
= \psi(M^M+M^{\psi_{M}(M^M+M)})%
}

\NovelDisplay{M04545R01}{%
(0)(1,1,1)(2,1,1)(3,1,1)(3,1)(2,1,1)(3,1)(4,2,1)(5,2,1)(6,2,1)(6,2)%
}
\NovelDisplay{M04545R02}{%
(4,2,1)(5,2,1)(6,2,1)(6,2)%
}

\NovelDisplay{M04546R01}{%
= \psi(I(I(\Omega_{I(1,0,0)+1},0),0))%
}

\NovelDisplay{M04547R01}{%
= \psi(M^M+M^{\psi_{M}(M^M+M^{\psi_{M}(M^M+M)})})%
}

复杂第一参数的进一步嵌套需要保留每层ψ，原矩阵的末列与坏根则依原版BMS另行确认。

上横台前一晚，古月娜又翻到那张理想矩阵。

纸仍在旧封套里。她看见那些曾经觉得应当成立的对应，心里仍会浮起一丝舍不得：如果真能按想象接下去，大家何必在这些台阶上一次次退回。

蔡须坤过来取纸，见她没合封套，便在旁边坐下。

古月娜没有再说自己已经知道它们该怎样对应，只问，以后还会不会继续想。

``会。''蔡须坤说。

她把正在核的原矩阵放在另一侧，指了指两张纸之间空出来的地方。

``但得知道是哪一套。这里还没定义的，不能拿到那边当已经走过。''

古月娜点头。她没有把理想那张撕掉，也没有让它压住今天真正要问的问题；两份纸各自留在原来的封套里。

第二天，唐三看见她先取出原版那份，才把自己的本子摊开。小舞接过水，在两人中间留出一个空位，等蔡须坤把昨日最后那处重新说明。

那一小段没有因为他们承认还有理想就变轻。可她们已能在同一处停下，不必先争哪张纸看起来更好。

四百米的横台在雾里露出来时，小舞先看见了石桩。

她却没有跑。脚底最后那一层压力还在，唐三也仍停在下面一阶。她把本子摊在能看清的地方，等自己说完那处参数，才慢慢收起脚跟。

霍雨浩已经站到平处，薛峰让后来的人留意湿滑的石沿。蔡须坤与卢本韦收好长页，古月娜等诺当上来，才将他的包递回。

六个月过去了。

蔡须坤看着石桩，没有立刻计算还差多少。她从那条六十列的第二项想回只有二十一列的原式，又想起自己曾经差点只凭纸的长短判断大小。

``今天那一条，明早还想再核。''她说。

卢本韦点头。两人将纸放在同一个封套里，里面的字迹与疑问都留得清楚。

四百米以上，山阶仍一层层没入云中。她们坐下，先等最后一个人把水喝完。

\chapter{第一百二十五章：多元I函数和M}\label{ux7b2cux4e00ux767eux4e8cux5341ux4e94ux7ae0ux591aux5143iux51fdux6570ux548cm}

第七个月，新页上的参数又多了起来。

蔡兴坤一开始想给每个位置画一种颜色，画到后来，纸上比包里的黑子鸡还热闹。叶阳在旁边看了半晌，问这些颜色究竟各指什么，他却一时说不清。

卢本韦取来另一张纸，让他先将参数从左往右说完整，再决定哪一处需要标记。蔡须坤把M侧放在旁边，霍雨浩请大家指出每一个幂指数的范围。

远处的五百米仍是第一考的目标，眼前这几级却足够让十一人用去许多天。

夏夜星空没有将更高处的讲义提前塞到他们手中。三百七十米以后这段从IFP走来的路，还有最后几层要由每个人亲自接清。

第三参数变化时，小舞下意识去看前面的幂。

她已习惯前两章那种读法，见右边仍有熟悉的自幂背景，便准备照原来的位置接。古月娜问她，参数改的是哪一格，她才把手收回I侧，从逗号重新数起。

蔡须坤没有替她将答案念完，而是等她先说清新背景，再接后面的附加项。

\NovelDisplay{M04548R01}{%
(0)(1,1,1)(2,1,1)(3,1,1)(3,1)(2,1,1)(3,1)(4,2,1)(5,2,1)(6,2,1)(6,2)(5)%
}

\NovelDisplay{M04549R01}{%
= \psi(I(1,0,1))%
}

\NovelDisplay{M04550R01}{%
= \psi(M^M\times 2)%
}

\NovelDisplay{M04551R01}{%
(0)(1,1,1)(2,1,1)(3,1,1)(3,1)(2,1,1)(3,1)(4,2,1)(5,2,1)(6,2,1)%
}
\NovelDisplay{M04551R02}{%
(6,2)(5,2,1)%
}

\NovelDisplay{M04552R01}{%
= \psi(I(1,0,1)\times \omega)%
}

\NovelDisplay{M04553R01}{%
= \psi(M^M\times 2 + \psi_{\psi_{M}(M^M\times 2 + M)}(1))%
}

\NovelDisplay{M04554R01}{%
(0)(1,1,1)(2,1,1)(3,1,1)(3,1)(2,1,1)(3,1)(4,2,1)(5,2,1)(6,2,1)(6,2)%
}
\NovelDisplay{M04554R02}{%
(5,2,1)(6,2)(7,3)%
}

\NovelDisplay{M04555R01}{%
= \psi(\Omega_{I(1,0,1)+1})%
}

\NovelDisplay{M04556R01}{%
= \psi(M^M\times 2 + M)%
}

\NovelDisplay{M04557R01}{%
(0)(1,1,1)(2,1,1)(3,1,1)(3,1)(2,1,1)(3,1)(4,2,1)(5,2,1)(6,2,1)(6,2)%
}
\NovelDisplay{M04557R02}{%
(5,2,1)(6,2)(7,3,1)(8,3,1)(9,3,1)(9,1)(2)%
}

\NovelDisplay{M04558R01}{%
= \psi(I(I(1,0,0),I(1,0,1)+1))%
}

\NovelDisplay{M04559R01}{%
= \psi(M^M\times 2 + M^{\psi_{M}(M^M)})%
}

第三参数增加后，M侧以整个自幂背景的倍数起步，追加项仍须按所在层次逐条读清。

\NovelDisplay{M04560R01}{%
(0)(1,1,1)(2,1,1)(3,1,1)(3,1)(2,1,1)(3,1)(4,2,1)(5,2,1)(6,2,1)(6,2)%
}
\NovelDisplay{M04560R02}{%
(5,2,1)(6,2)(7,3,1)(8,3,1)(9,3,1)(9,2)(5)%
}

\NovelDisplay{M04561R01}{%
= \psi(I(I(1,0,1),1))%
}

\NovelDisplay{M04562R01}{%
= \psi(M^M\times 2 + M^{\psi_{M}(M^M\times 2)})%
}

\NovelDisplay{M04563R01}{%
(0)(1,1,1)(2,1,1)(3,1,1)(3,1)(2,1,1)(3,1)(4,2,1)(5,2,1)(6,2,1)(6,2)%
}
\NovelDisplay{M04563R02}{%
(5,2,1)(6,2)(7,3,1)(8,3,1)(9,3,1)(9,3)%
}

\NovelDisplay{M04564R01}{%
= \psi(I(\Omega_{I(1,0,1)+1},0))%
}

\NovelDisplay{M04565R01}{%
= \psi(M^M\times 2 + M^{\psi_{M}(M^M\times 2+M)})%
}

\NovelDisplay{M04566R01}{%
(0)(1,1,1)(2,1,1)(3,1,1)(3,1)(2,1,1)(3,1)(4,2,1)(5,2,1)(6,2,1)(6,2)%
}
\NovelDisplay{M04566R02}{%
(5,2,1)(6,2)(7,3,1)(8,3,1)(9,3,1)(9,3)(8)%
}

\NovelDisplay{M04567R01}{%
= \psi(I(1,0,2))%
}

\NovelDisplay{M04568R01}{%
= \psi(M^M\times 3)%
}

\NovelDisplay{M04569R01}{%
(0)(1,1,1)(2,1,1)(3,1,1)(3,1)(2,1,1)(3,1)(4,2,1)(5,2,1)(6,2,1)(6,2)%
}
\NovelDisplay{M04569R02}{%
(5,2,1)(6,2)(7,3,1)(8,3,1)(9,3,1)(9,3)(8,3,1)(9,3)(10,4,1)(11,4,1)%
}
\NovelDisplay{M04569R03}{%
(12,4,1)(12,4)(11)%
}

\NovelDisplay{M04570R01}{%
= \psi(I(1,0,3))%
}

\NovelDisplay{M04571R01}{%
= \psi(M^M\times 4)%
}

同一背景下新的参数嵌套与第三参数增加分别出现，不能把它们都简写成在末尾加一。

先由薛峰、唐三与诺当各求一项，三次都把原矩阵留在身旁，蔡须坤等他们写完以后再看共同保留的前段。

\NovelDisplay{M04572R01}{%
\text{展开第零项：}\quad(0)(1,1,1)(2,1,1)(3,1,1)(3,1)(2,1,1)(3,1)(4,2,1)(5,2,1)%
}
\NovelDisplay{M04572R02}{%
(6,2,1)(6,2)(5,2,1)(6,2)(7,3,1)(8,3,1)(9,3,1)(9,3)(8,3,1)(9,3)(10,4,1)%
}
\NovelDisplay{M04572R03}{%
(11,4,1)(12,4,1)(12,4)%
}

\NovelDisplay{M04573R01}{%
\text{展开第一项：}\quad(0)(1,1,1)(2,1,1)(3,1,1)(3,1)(2,1,1)(3,1)(4,2,1)(5,2,1)%
}
\NovelDisplay{M04573R02}{%
(6,2,1)(6,2)(5,2,1)(6,2)(7,3,1)(8,3,1)(9,3,1)(9,3)(8,3,1)(9,3)(10,4,1)%
}
\NovelDisplay{M04573R03}{%
(11,4,1)(12,4,1)(12,4)(10,4,1)(11,4,1)(12,4,1)(12,4)%
}

\NovelDisplay{M04574R01}{%
\text{展开第二项：}\quad(0)(1,1,1)(2,1,1)(3,1,1)(3,1)(2,1,1)(3,1)(4,2,1)(5,2,1)%
}
\NovelDisplay{M04574R02}{%
(6,2,1)(6,2)(5,2,1)(6,2)(7,3,1)(8,3,1)(9,3,1)(9,3)(8,3,1)(9,3)(10,4,1)%
}
\NovelDisplay{M04574R03}{%
(11,4,1)(12,4,1)(12,4)(10,4,1)(11,4,1)(12,4,1)(12,4)(10,4,1)(11,4,1)%
}
\NovelDisplay{M04574R04}{%
(12,4,1)(12,4)%
}

第三参数的有限变化仍须落实到矩阵自身，三次独立取项不能借用另一人的中间结果。

诺当核到最后，自己指出少写的那一列。唐三等他补完，才将两份结果并排收好。

蔡兴坤有一天将自己落后的几页藏到包底。

他怕大家每回休息都要等自己，便只拿出已经能说清的一张。薛峰听过，又问今天真正停在哪一处，他笑着说都差不多，手却没有从包口松开。

黑子鸡恰在这时动了一下，几张纸从底下滑出来。奥斯卡先替他捡起，没有念上面的空白，只将原来的页次排好，放回他膝上。

蔡兴坤红着脸，过了一会儿，终于取出最下面那张。

薛峰坐回来，让他从开头说。说到真正不懂的地方，他这次没有接一句笑话，而是抬头问了出来。

后来同伴也仍要等他，可那份等待终于对上了明确的问题。他不必再把本子藏起来，奥斯卡也不必凭他笑没笑，猜今天究竟顺不顺。

\NovelDisplay{M04575R01}{%
(0)(1,1,1)(2,1,1)(3,1,1)(3,1)(2,1,1)(3,1,1)%
}

\NovelDisplay{M04576R01}{%
= \psi(I(1,0,\omega))%
}

\NovelDisplay{M04577R01}{%
= \psi(M^M\times \omega)%
}

\NovelDisplay{M04578R01}{%
(0)(1,1,1)(2,1,1)(3,1,1)(3,1)(2,1,1)(3,1,1)(2,1,1)(3,1)(2)%
}

\NovelDisplay{M04579R01}{%
= \psi(I(1,1,0))%
}

\NovelDisplay{M04580R01}{%
= \psi(M^{M+1})%
}

\NovelDisplay{M04581R01}{%
(0)(1,1,1)(2,1,1)(3,1,1)(3,1)(2,1,1)(3,1,1)(2,1,1)(3,1)(4,2)%
}

\NovelDisplay{M04582R01}{%
= \psi(\Omega_{I(1,1,0)+1})%
}

\NovelDisplay{M04583R01}{%
= \psi(M^{M+1}+M)%
}

\NovelDisplay{M04584R01}{%
(0)(1,1,1)(2,1,1)(3,1,1)(3,1)(2,1,1)(3,1,1)(2,1,1)(3,1)(4,2,1)(5,2,1)%
}
\NovelDisplay{M04584R02}{%
(6,2,1)(6,1)(2)%
}

\NovelDisplay{M04585R01}{%
= \psi(I(I(1,1,0),1))%
}

\NovelDisplay{M04586R01}{%
= \psi(M^{M+1}+M^{\psi_{M}(M^{M+1})})%
}

第三参数的ω与第二参数的后继在这里形成不同结构，M侧的幂外乘数和幂指数也就必须分开看。

蔡兴坤的颜色终于少了。

他只标自己容易混读的参数位置，其他地方留着原式。黑子鸡在桌边啄了一下，差点碰翻水，他急着去护，袖口又扫乱了几页。

奥斯卡先扶住水壶，叶阳把纸收回来，让他别急。等几个人重新坐好，蔡兴坤低声说，自己刚才又忘了读到哪里。

薛峰将原矩阵递过去：``从你能确定的地方说。''

蔡兴坤点头，没有拿方才的玩笑把问题岔过去。

\NovelDisplay{M04587R01}{%
(0)(1,1,1)(2,1,1)(3,1,1)(3,1)(2,1,1)(3,1,1)(2,1,1)(3,1)(4,2,1)%
}
\NovelDisplay{M04587R02}{%
(5,2,1)(6,2,1)(6,2)%
}

\NovelDisplay{M04588R01}{%
= \psi(I(\Omega_{I(1,1,0)+1},0))%
}

\NovelDisplay{M04589R01}{%
= \psi(M^{M+1}+M^{\psi_{M}(M^{M+1}+M)})%
}

\NovelDisplay{M04590R01}{%
(0)(1,1,1)(2,1,1)(3,1,1)(3,1)(2,1,1)(3,1,1)(2,1,1)(3,1)(4,2,1)(5,2,1)%
}
\NovelDisplay{M04590R02}{%
(6,2,1)(6,2)(5)%
}

\NovelDisplay{M04591R01}{%
= \psi(I(1,0,I(1,1,0)+1))%
}

\NovelDisplay{M04592R01}{%
= \psi(M^{M+1}+M^M)%
}

\NovelDisplay{M04593R01}{%
(0)(1,1,1)(2,1,1)(3,1,1)(3,1)(2,1,1)(3,1,1)(2,1,1)(3,1)(4,2,1)(5,2,1)%
}
\NovelDisplay{M04593R02}{%
(6,2,1)(6,2)(5,2,1)(6,2,1)(5,2,1)(6,2)%
}

\NovelDisplay{M04594R01}{%
= \psi(I(1,0,\Omega_{I(1,1,0)+1}))%
}

\NovelDisplay{M04595R01}{%
= \psi(M^{M+1}+M^M\times \psi_{M}(M^{M+1}+M))%
}

\NovelDisplay{M04596R01}{%
(0)(1,1,1)(2,1,1)(3,1,1)(3,1)(2,1,1)(3,1,1)(2,1,1)(3,1)(4,2,1)(5,2,1)%
}
\NovelDisplay{M04596R02}{%
(6,2,1)(6,2)(5,2,1)(6,2,1)(5,2,1)(6,2)(5)%
}

\NovelDisplay{M04597R01}{%
= \psi(I(1,1,1))%
}

\NovelDisplay{M04598R01}{%
= \psi(M^{M+1}\times 2)%
}

新的指数背景后还能再添旧背景，古月娜提醒小舞，那一处加号没有把两个幂指数合成一个。

\NovelDisplay{M04599R01}{%
(0)(1,1,1)(2,1,1)(3,1,1)(3,1)(2,1,1)(3,1,1)(2,1,1)(3,1)(4,2,1)(5,2,1)%
}
\NovelDisplay{M04599R02}{%
(6,2,1)(6,2)(5,2,1)(6,2,1)(5,2,1)(6,2)(7,3,1)(8,3,1)(9,3,1)(9,1)(2)%
}

\NovelDisplay{M04600R01}{%
= \psi(I(I(1,1,0),I(1,1,1)+1))%
}

\NovelDisplay{M04601R01}{%
= \psi(M^{M+1}\times 2+M^{\psi_{M}(M^{M+1})})%
}

\NovelDisplay{M04602R01}{%
(0)(1,1,1)(2,1,1)(3,1,1)(3,1)(2,1,1)(3,1,1)(2,1,1)(3,1)(4,2,1)(5,2,1)%
}
\NovelDisplay{M04602R02}{%
(6,2,1)(6,2)(5,2,1)(6,2,1)(5,2,1)(6,2)(7,3,1)(8,3,1)(9,3,1)(9,2)(5)%
}

\NovelDisplay{M04603R01}{%
= \psi(I(I(1,1,1),1))%
}

\NovelDisplay{M04604R01}{%
= \psi(M^{M+1}\times 2+M^{\psi_{M}(M^{M+1}\times 2)})%
}

\NovelDisplay{M04605R01}{%
(0)(1,1,1)(2,1,1)(3,1,1)(3,1)(2,1,1)(3,1,1)(2,1,1)(3,1)(4,2,1)(5,2,1)%
}
\NovelDisplay{M04605R02}{%
(6,2,1)(6,2)(5,2,1)(6,2,1)(5,2,1)(6,2)(7,3,1)(8,3,1)(9,3,1)(9,3)(8)%
}

\NovelDisplay{M04606R01}{%
= \psi(I(1,0,I(1,1,1)+1))%
}

\NovelDisplay{M04607R01}{%
= \psi(M^{M+1}\times 2+M^M)%
}

\NovelDisplay{M04608R01}{%
(0)(1,1,1)(2,1,1)(3,1,1)(3,1)(2,1,1)(3,1,1)(2,1,1)(3,1)(4,2,1)(5,2,1)%
}
\NovelDisplay{M04608R02}{%
(6,2,1)(6,2)(5,2,1)(6,2,1)(5,2,1)(6,2)(7,3,1)(8,3,1)(9,3,1)(9,3)(8,3,1)%
}
\NovelDisplay{M04608R03}{%
(9,3,1)(8,3,1)(9,3)(8)%
}

\NovelDisplay{M04609R01}{%
= \psi(I(1,1,2))%
}

\NovelDisplay{M04610R01}{%
= \psi(M^{M+1}\times 3)%
}

双倍新背景后的附加项仍各有其作用对象，熟悉的前段不构成省略后面下标的理由。

同一新背景的倍数变长以后，霍雨浩请小舞先圈原来的坏部，再由古月娜与叶阳分别求另外两项。

\NovelDisplay{M04611R01}{%
\text{展开第零项：}\quad(0)(1,1,1)(2,1,1)(3,1,1)(3,1)(2,1,1)(3,1,1)(2,1,1)(3,1)%
}
\NovelDisplay{M04611R02}{%
(4,2,1)(5,2,1)(6,2,1)(6,2)(5,2,1)(6,2,1)(5,2,1)(6,2)(7,3,1)(8,3,1)%
}
\NovelDisplay{M04611R03}{%
(9,3,1)(9,3)(8,3,1)(9,3,1)(8,3,1)(9,3)%
}

\NovelDisplay{M04612R01}{%
\text{展开第一项：}\quad(0)(1,1,1)(2,1,1)(3,1,1)(3,1)(2,1,1)(3,1,1)(2,1,1)(3,1)%
}
\NovelDisplay{M04612R02}{%
(4,2,1)(5,2,1)(6,2,1)(6,2)(5,2,1)(6,2,1)(5,2,1)(6,2)(7,3,1)(8,3,1)%
}
\NovelDisplay{M04612R03}{%
(9,3,1)(9,3)(8,3,1)(9,3,1)(8,3,1)(9,3)(7,3,1)(8,3,1)(9,3,1)(9,3)(8,3,1)%
}
\NovelDisplay{M04612R04}{%
(9,3,1)(8,3,1)(9,3)%
}

\NovelDisplay{M04613R01}{%
\text{展开第二项：}\quad(0)(1,1,1)(2,1,1)(3,1,1)(3,1)(2,1,1)(3,1,1)(2,1,1)(3,1)%
}
\NovelDisplay{M04613R02}{%
(4,2,1)(5,2,1)(6,2,1)(6,2)(5,2,1)(6,2,1)(5,2,1)(6,2)(7,3,1)(8,3,1)%
}
\NovelDisplay{M04613R03}{%
(9,3,1)(9,3)(8,3,1)(9,3,1)(8,3,1)(9,3)(7,3,1)(8,3,1)(9,3,1)(9,3)(8,3,1)%
}
\NovelDisplay{M04613R04}{%
(9,3,1)(8,3,1)(9,3)(7,3,1)(8,3,1)(9,3,1)(9,3)(8,3,1)(9,3,1)(8,3,1)(9,3)%
}

这三项各保留自己的列次序与生成份数，不因右侧反复出现相同背景就删除重复段。

小舞把第二项重新数过，终于指出自己第一次多添的一份究竟接在了哪里。

\NovelDisplay{M04614R01}{%
(0)(1,1,1)(2,1,1)(3,1,1)(3,1)(2,1,1)(3,1,1)(2,1,1)(3,1,1)%
}

\NovelDisplay{M04615R01}{%
= \psi(I(1,1,\omega))%
}

\NovelDisplay{M04616R01}{%
= \psi(M^{M+1}\times \omega)%
}

\NovelDisplay{M04617R01}{%
(0)(1,1,1)(2,1,1)(3,1,1)(3,1)(2,1,1)(3,1,1)(2,1,1)(3,1,1)(2,1,1)(3,1)(2)%
}

\NovelDisplay{M04618R01}{%
= \psi(I(1,2,0))%
}

\NovelDisplay{M04619R01}{%
= \psi(M^{M+2})%
}

\NovelDisplay{M04620R01}{%
(0)(1,1,1)(2,1,1)(3,1,1)(3,1)(2,1,1)(3,1,1)(2,1,1)(3,1,1)(2,1,1)(3,1)%
}
\NovelDisplay{M04620R02}{%
(4,2,1)(5,2,1)(6,2,1)(6,2)(5)%
}

\NovelDisplay{M04621R01}{%
= \psi(I(1,0,I(1,2,0)+1))%
}

\NovelDisplay{M04622R01}{%
= \psi(M^{M+2}+M^M)%
}

\NovelDisplay{M04623R01}{%
(0)(1,1,1)(2,1,1)(3,1,1)(3,1)(2,1,1)(3,1,1)(2,1,1)(3,1,1)(2,1,1)(3,1)%
}
\NovelDisplay{M04623R02}{%
(4,2,1)(5,2,1)(6,2,1)(6,2)(5,2,1)(6,2,1)(5,2,1)(6,2)(5)%
}

\NovelDisplay{M04624R01}{%
= \psi(I(1,1,I(1,2,0)+1))%
}

\NovelDisplay{M04625R01}{%
= \psi(M^{M+2}+M^{M+1})%
}

第二参数继续增加时，先辨认幂指数的变化，再看后来添入的旧幂项，不能把两者交换。

山阶上的一天，比纸上多出的几行更长。

蔡须坤已经能将M侧读清，抬脚时却仍被一个细小的停顿扯住。她看见上面那几级，心里一急，想用更大的魂力压过去，霍雨浩叫她的声音便从身后传来。

她自己收回脚，坐在平处，才发现刚才拿的是前一条原矩阵。两条的外框太像，她换过右式，手里的原式却没有跟着换。

卢本韦替她将两张纸分开，先请她重新指出末列。

\NovelDisplay{M04626R01}{%
(0)(1,1,1)(2,1,1)(3,1,1)(3,1)(2,1,1)(3,1,1)(2,1,1)(3,1,1)(2,1,1)(3,1)%
}
\NovelDisplay{M04626R02}{%
(4,2,1)(5,2,1)(6,2,1)(6,2)(5,2,1)(6,2,1)(5,2,1)(6,2,1)(5,2,1)(6,2)(5)%
}

\NovelDisplay{M04627R01}{%
= \psi(I(1,2,1))%
}

\NovelDisplay{M04628R01}{%
= \psi(M^{M+2}\times 2)%
}

\NovelDisplay{M04629R01}{%
(0)(1,1,1)(2,1,1)(3,1,1)(3,1)(2,1,1)(3,1,1)(2,1,1)(3,1,1)(2,1,1)(3,1,1)%
}

\NovelDisplay{M04630R01}{%
= \psi(I(1,2,\omega))%
}

\NovelDisplay{M04631R01}{%
= \psi(M^{M+2}\times \omega)%
}

\NovelDisplay{M04632R01}{%
(0)(1,1,1)(2,1,1)(3,1,1)(3,1)(2,1,1)(3,1,1)(2,1,1)(3,1,1)(2,1,1)(3,1,1)%
}
\NovelDisplay{M04632R02}{%
(2,1,1)(3,1)(2)%
}

\NovelDisplay{M04633R01}{%
= \psi(I(1,3,0))%
}

\NovelDisplay{M04634R01}{%
= \psi(M^{M+3})%
}

\NovelDisplay{M04635R01}{%
(0)(1,1,1)(2,1,1)(3,1,1)(3,1)(2,1,1)(3,1,1)(2,1,1)(3,1,1)(2,1,1)(3,1,1)%
}
\NovelDisplay{M04635R02}{%
(2,1,1)(3,1,1)(2,1,1)(3,1)(2)%
}

\NovelDisplay{M04636R01}{%
= \psi(I(1,4,0))%
}

\NovelDisplay{M04637R01}{%
= \psi(M^{M+4})%
}

指数里的有限后继与幂外倍数在此连续出现，每条矩阵仍要独立核对，不能只背右式的数目。

一天傍晚，薛峰先独自回到山脚。

他原本走在前面，那次却在收脚时看错石沿，膝盖重重碰了一下。伤并不深，大家还是让他先下去，蔡须坤听完后来的人说明，也将原来准备追上的几级留到次日。

薛峰坐在屋里，越等越觉得过意不去。叶阳回来取纸时，他想把大家的问题先核完，叶阳看见他总往门外望，便将纸放在桌上，让他先把伤处处理好。

``这几张明天还在。''

薛峰低头看了许久，没有再拿走所有人的本子。晚上他只核自己的那一份，将不清楚的位置留下，等同伴回来，先问今天谁还需要重说。

蔡须坤听见，便把白天停住的一处交给他。两人坐着讨论，卢本韦将M侧移到中间，霍雨浩指明原矩阵所在的页次；等话说清，薛峰才发现自己已经许久没有去想上面多走了几级。

第二日，他没有抢着走到最前。奥斯卡把水递来，蔡兴坤问他的腿，黑子鸡在包里动了两下，他便笑着让大家先把纸拿稳。

那几天的速度慢了一些，原先总等他先开口的人，也开始主动将自己的问题摆到前面。薛峰坐在旁边听，遇到能核清的一处才接话，伤处休息着，手里的路却没有因暂时站得低而丢掉。

回山的前一晚，薛峰又看见那张留在桌上的问题页。他没有将它替同伴写完，只取出对应的原矩阵，放在次日能一并找到的位置。叶阳回来，看过以后，将自己的纸收进包里，说这一处明早由自己先问。薛峰点头，低头仔细重新处理好膝上的布带。门外有风，他起身关好窗，让那几张纸安稳留在灯下。

\NovelDisplay{M04638R01}{%
(0)(1,1,1)(2,1,1)(3,1,1)(3,1)(2,1,1)(3,1,1)(3,1)%
}

\NovelDisplay{M04639R01}{%
= \psi(I(1,\omega,0))%
}

\NovelDisplay{M04640R01}{%
= \psi(M^{M+\omega})%
}

\NovelDisplay{M04641R01}{%
(0)(1,1,1)(2,1,1)(3,1,1)(3,1)(2,1,1)(3,1,1)(3,1)(2)%
}

\NovelDisplay{M04642R01}{%
= \psi(I(2,0,0))%
}

\NovelDisplay{M04643R01}{%
= \psi(M^{M\times 2})%
}

\NovelDisplay{M04644R01}{%
(0)(1,1,1)(2,1,1)(3,1,1)(3,1)(2,1,1)(3,1,1)(3,1)(2,1,1)%
}

\NovelDisplay{M04645R01}{%
= \psi(I(2,0,0)\times \omega)%
}

\NovelDisplay{M04646R01}{%
= \psi(M^{M\times 2}+\psi_{\psi_{M}(M^{M\times 2}+M)}(1))%
}

\NovelDisplay{M04647R01}{%
(0)(1,1,1)(2,1,1)(3,1,1)(3,1)(2,1,1)(3,1,1)(3,1)(2,1,1)(3,1)(4,2)%
}

\NovelDisplay{M04648R01}{%
= \psi(\Omega_{I(2,0,0)+1})%
}

\NovelDisplay{M04649R01}{%
= \psi(M^{M\times 2}+M)%
}

M加ω和M乘二都进入指数，却不是同一种增长写法，蔡须坤把指数框完整画出才让大家继续。

\NovelDisplay{M04650R01}{%
(0)(1,1,1)(2,1,1)(3,1,1)(3,1)(2,1,1)(3,1,1)(3,1)(2,1,1)(3,1)(4,2,1)%
}
\NovelDisplay{M04650R02}{%
(5,2,1)(6,2)(5)%
}

\NovelDisplay{M04651R01}{%
= \psi(I(1,0,I(2,0,0)+1))%
}

\NovelDisplay{M04652R01}{%
= \psi(M^{M\times 2}+M^M)%
}

\NovelDisplay{M04653R01}{%
(0)(1,1,1)(2,1,1)(3,1,1)(3,1)(2,1,1)(3,1,1)(3,1)(2,1,1)(3,1)(4,2,1)%
}
\NovelDisplay{M04653R02}{%
(5,2,1)(6,2)(5,2,1)(6,2,1)(5,2,1)(6,2)(5)%
}

\NovelDisplay{M04654R01}{%
= \psi(I(1,1,I(2,0,0)+1))%
}

\NovelDisplay{M04655R01}{%
= \psi(M^{M\times 2}+M^{M+1})%
}

\NovelDisplay{M04656R01}{%
(0)(1,1,1)(2,1,1)(3,1,1)(3,1)(2,1,1)(3,1,1)(3,1)(2,1,1)(3,1)(4,2,1)%
}
\NovelDisplay{M04656R02}{%
(5,2,1)(6,2)(5,2,1)(6,2,1)(5,2,1)(6,2,1)(5,2,1)(6,2)(5)%
}

\NovelDisplay{M04657R01}{%
= \psi(I(1,2,I(2,0,0)+1))%
}

\NovelDisplay{M04658R01}{%
= \psi(M^{M\times 2}+M^{M+2})%
}

\NovelDisplay{M04659R01}{%
(0)(1,1,1)(2,1,1)(3,1,1)(3,1)(2,1,1)(3,1,1)(3,1)(2,1,1)(3,1)(4,2,1)%
}
\NovelDisplay{M04659R02}{%
(5,2,1)(6,2)(5,2,1)(6,2,1)(6,2)%
}

\NovelDisplay{M04660R01}{%
= \psi(I(1,\Omega_{I(2,0,0)+1},0))%
}

\NovelDisplay{M04661R01}{%
= \psi(M^{M\times 2}+M^{M+\psi_{M}(M^{M\times 2}+M)})%
}

新指数背景后再添不同旧幂项时，加法次序照原式保留，不能为了对称就交换前后。

有一次，古月娜已经讲到末尾，唐三却还停在第一行。

他越怕落后，越不愿让她重说，最后只点头。小舞看了他一眼，先举起自己的纸，说她有一处没听明白。

古月娜停下来。唐三等妹妹问完，终于也将自己的问题指出。

那晚收纸时，他问小舞是不是早看出来了。她没答，只让他下次早一点问，别等自己也不知道从哪里开始说。

唐三点头，将白天那一处重新写在自己的页角。

\NovelDisplay{M04662R01}{%
(0)(1,1,1)(2,1,1)(3,1,1)(3,1)(2,1,1)(3,1,1)(3,1)(2,1,1)(3,1)(4,2,1)%
}
\NovelDisplay{M04662R02}{%
(5,2,1)(6,2)(5,2,1)(6,2,1)(6,2)(5)%
}

\NovelDisplay{M04663R01}{%
= \psi(I(2,0,1))%
}

\NovelDisplay{M04664R01}{%
= \psi(M^{M\times 2}\times 2)%
}

\NovelDisplay{M04665R01}{%
(0)(1,1,1)(2,1,1)(3,1,1)(3,1)(2,1,1)(3,1,1)(3,1)(2,1,1)(3,1)(4,2,1)%
}
\NovelDisplay{M04665R02}{%
(5,2,1)(6,2)(5,2,1)(6,2,1)(6,2)(5,2,1)(6,2)(7,3,1)(8,3,1)(9,3,1)(9,3)%
}
\NovelDisplay{M04665R03}{%
(8,3,1)(9,3,1)(9,3)(8)%
}

\NovelDisplay{M04666R01}{%
= \psi(I(2,0,2))%
}

\NovelDisplay{M04667R01}{%
= \psi(M^{M\times 2}\times 3)%
}

\NovelDisplay{M04668R01}{%
(0)(1,1,1)(2,1,1)(3,1,1)(3,1)(2,1,1)(3,1,1)(3,1)(2,1,1)(3,1,1)%
}

\NovelDisplay{M04669R01}{%
= \psi(I(2,0,\omega))%
}

\NovelDisplay{M04670R01}{%
= \psi(M^{M\times 2}\times \omega)%
}

\NovelDisplay{M04671R01}{%
(0)(1,1,1)(2,1,1)(3,1,1)(3,1)(2,1,1)(3,1,1)(3,1)(2,1,1)(3,1,1)%
}
\NovelDisplay{M04671R02}{%
(2,1,1)(3,1)(2)%
}

\NovelDisplay{M04672R01}{%
= \psi(I(2,1,0))%
}

\NovelDisplay{M04673R01}{%
= \psi(M^{M\times 2+1})%
}

第三参数增加和第二参数后继再次并列，卢本韦先读参数，霍雨浩再指出M侧哪一个乘数在幂外。

蔡须坤与卢本韦各自求项，薛峰另取一张纸独立完成第三份，再同两人核对末列之前的保留段。

\NovelDisplay{M04674R01}{%
\text{展开第零项：}\quad(0)(1,1,1)(2,1,1)(3,1,1)(3,1)(2,1,1)(3,1,1)(3,1)(2,1,1)%
}
\NovelDisplay{M04674R02}{%
(3,1)(4,2,1)(5,2,1)(6,2)(5,2,1)(6,2,1)(6,2)(5,2,1)(6,2)(7,3,1)(8,3,1)%
}
\NovelDisplay{M04674R03}{%
(9,3,1)(9,3)(8,3,1)(9,3,1)(9,3)%
}

\NovelDisplay{M04675R01}{%
\text{展开第一项：}\quad(0)(1,1,1)(2,1,1)(3,1,1)(3,1)(2,1,1)(3,1,1)(3,1)(2,1,1)%
}
\NovelDisplay{M04675R02}{%
(3,1)(4,2,1)(5,2,1)(6,2)(5,2,1)(6,2,1)(6,2)(5,2,1)(6,2)(7,3,1)(8,3,1)%
}
\NovelDisplay{M04675R03}{%
(9,3,1)(9,3)(8,3,1)(9,3,1)(9,3)(7,3,1)(8,3,1)(9,3,1)(9,3)(8,3,1)%
}
\NovelDisplay{M04675R04}{%
(9,3,1)(9,3)%
}

\NovelDisplay{M04676R01}{%
\text{展开第二项：}\quad(0)(1,1,1)(2,1,1)(3,1,1)(3,1)(2,1,1)(3,1,1)(3,1)(2,1,1)%
}
\NovelDisplay{M04676R02}{%
(3,1)(4,2,1)(5,2,1)(6,2)(5,2,1)(6,2,1)(6,2)(5,2,1)(6,2)(7,3,1)(8,3,1)%
}
\NovelDisplay{M04676R03}{%
(9,3,1)(9,3)(8,3,1)(9,3,1)(9,3)(7,3,1)(8,3,1)(9,3,1)(9,3)(8,3,1)(9,3,1)%
}
\NovelDisplay{M04676R04}{%
(9,3)(7,3,1)(8,3,1)(9,3,1)(9,3)(8,3,1)(9,3,1)(9,3)%
}

同一个原式的取项编号各自使用，不能把三份结果连起来当作更高编号的一次生成。

薛峰没有先去看她们的纸。等自己的末列落下，他才将原矩阵转回中间，逐列说明自己是怎样保留与生成的。

\NovelDisplay{M04677R01}{%
(0)(1,1,1)(2,1,1)(3,1,1)(3,1)(2,1,1)(3,1,1)(3,1)(2,1,1)(3,1,1)(2,1,1)%
}
\NovelDisplay{M04677R02}{%
(3,1,1)(2,1,1)(3,1)(2)%
}

\NovelDisplay{M04678R01}{%
= \psi(I(2,2,0))%
}

\NovelDisplay{M04679R01}{%
= \psi(M^{M\times 2+2})%
}

\NovelDisplay{M04680R01}{%
(0)(1,1,1)(2,1,1)(3,1,1)(3,1)(2,1,1)(3,1,1)(3,1)(2,1,1)(3,1,1)(3,1)%
}

\NovelDisplay{M04681R01}{%
= \psi(I(2,\omega,0))%
}

\NovelDisplay{M04682R01}{%
= \psi(M^{M\times 2+\omega})%
}

\NovelDisplay{M04683R01}{%
(0)(1,1,1)(2,1,1)(3,1,1)(3,1)(2,1,1)(3,1,1)(3,1)(2,1,1)(3,1,1)(3,1)(2)%
}

\NovelDisplay{M04684R01}{%
= \psi(I(3,0,0))%
}

\NovelDisplay{M04685R01}{%
= \psi(M^{M\times 3})%
}

\NovelDisplay{M04686R01}{%
(0)(1,1,1)(2,1,1)(3,1,1)(3,1)(2,1,1)(3,1,1)(3,1)(2,1,1)(3,1,1)(3,1)%
}
\NovelDisplay{M04686R02}{%
(2,1,1)(3,1,1)(3,1)(2)%
}

\NovelDisplay{M04687R01}{%
= \psi(I(4,0,0))%
}

\NovelDisplay{M04688R01}{%
= \psi(M^{M\times 4})%
}

有限第二参数、ω第二参数与第一参数增加对应不同指数，读式时要保留整条参数链。

四参数I出现在纸上，卢本韦先把整串读了两次。

她知道这类写法给了自己一条熟悉的路，也知道路并没有因此变得轻省。多出的位置仍要核，旧参数嵌进去以后，边界仍要重新指出。

蔡须坤听她说完，指着M侧的幂指数问了一句。卢本韦翻回前页，两人把旧背景放在新背景后面比较，薛峰在旁边等她们各自说清，再将原矩阵移到中间。

\NovelDisplay{M04689R01}{%
(0)(1,1,1)(2,1,1)(3,1,1)(3,1)(3,1)(2)%
}

\NovelDisplay{M04690R01}{%
= \psi(I(1,0,0,0))%
}

\NovelDisplay{M04691R01}{%
= \psi(M^{M^2})%
}

\NovelDisplay{M04692R01}{%
(0)(1,1,1)(2,1,1)(3,1,1)(3,1)(3,1)(2,1,1)(3,1)(4,2)%
}

\NovelDisplay{M04693R01}{%
= \psi(\Omega_{I(1,0,0,0)+1})%
}

\NovelDisplay{M04694R01}{%
= \psi(M^{M^2}+M)%
}

\NovelDisplay{M04695R01}{%
(0)(1,1,1)(2,1,1)(3,1,1)(3,1)(3,1)(2,1,1)(3,1)(4,2,1)(5,2,1)(6,2,1)%
}
\NovelDisplay{M04695R02}{%
(6,2)(5)%
}

\NovelDisplay{M04696R01}{%
= \psi(I(1,0,I(1,0,0,0)+1))%
}

\NovelDisplay{M04697R01}{%
= \psi(M^{M^2}+M^M)%
}

\NovelDisplay{M04698R01}{%
(0)(1,1,1)(2,1,1)(3,1,1)(3,1)(3,1)(2,1,1)(3,1)(4,2,1)(5,2,1)(6,2,1)(6,2)%
}
\NovelDisplay{M04698R02}{%
(5,2,1)(6,2,1)(5,2,1)(6,2)(5)%
}

\NovelDisplay{M04699R01}{%
= \psi(I(1,1,I(1,0,0,0)+1))%
}

\NovelDisplay{M04700R01}{%
= \psi(M^{M^2}+M^{M+1})%
}

四参数I接到M平方进入幂指数的阶段，随后添加的各个旧幂项仍保持原来的顺序。

\NovelDisplay{M04701R01}{%
(0)(1,1,1)(2,1,1)(3,1,1)(3,1)(3,1)(2,1,1)(3,1)(4,2,1)(5,2,1)(6,2,1)(6,2)%
}
\NovelDisplay{M04701R02}{%
(5,2,1)(6,2,1)(6,2)(5)%
}

\NovelDisplay{M04702R01}{%
= \psi(I(2,0,I(1,0,0,0)+1))%
}

\NovelDisplay{M04703R01}{%
= \psi(M^{M^2}+M^{M\times 2})%
}

\NovelDisplay{M04704R01}{%
(0)(1,1,1)(2,1,1)(3,1,1)(3,1)(3,1)(2,1,1)(3,1)(4,2,1)(5,2,1)(6,2,1)(6,2)%
}
\NovelDisplay{M04704R02}{%
(5,2,1)(6,2,1)(6,2)(5,2,1)(6,2,1)(6,2)(5)%
}

\NovelDisplay{M04705R01}{%
= \psi(I(3,0,I(1,0,0,0)+1))%
}

\NovelDisplay{M04706R01}{%
= \psi(M^{M^2}+M^{M\times 3})%
}

\NovelDisplay{M04707R01}{%
(0)(1,1,1)(2,1,1)(3,1,1)(3,1)(3,1)(2,1,1)(3,1)(4,2,1)(5,2,1)(6,2,1)%
}
\NovelDisplay{M04707R02}{%
(6,2)(6,2)(5)%
}

\NovelDisplay{M04708R01}{%
= \psi(I(1,0,0,1))%
}

\NovelDisplay{M04709R01}{%
= \psi(M^{M^2}\times 2)%
}

\NovelDisplay{M04710R01}{%
(0)(1,1,1)(2,1,1)(3,1,1)(3,1)(3,1)(2,1,1)(3,1,1)%
}

\NovelDisplay{M04711R01}{%
= \psi(I(1,0,0,\omega))%
}

\NovelDisplay{M04712R01}{%
= \psi(M^{M^2}\times \omega)%
}

同一新背景后添旧项与新背景本身的倍数分开求项，薛峰没有因为整段重复就省去其中的列。

再遇到高参数串中的旧幂项时，卢本韦先盖住M侧，让唐三、小舞与蔡兴坤从原矩阵分别求项。

\NovelDisplay{M04713R01}{%
\text{展开第零项：}\quad(0)(1,1,1)(2,1,1)(3,1,1)(3,1)(3,1)(2,1,1)(3,1)(4,2,1)%
}
\NovelDisplay{M04713R02}{%
(5,2,1)(6,2,1)(6,2)(5,2,1)(6,2,1)(6,2)(5,2,1)(6,2,1)(6,2)%
}

\NovelDisplay{M04714R01}{%
\text{展开第一项：}\quad(0)(1,1,1)(2,1,1)(3,1,1)(3,1)(3,1)(2,1,1)(3,1)(4,2,1)%
}
\NovelDisplay{M04714R02}{%
(5,2,1)(6,2,1)(6,2)(5,2,1)(6,2,1)(6,2)(5,2,1)(6,2,1)(6,2)(4,2,1)(5,2,1)%
}
\NovelDisplay{M04714R03}{%
(6,2,1)(6,2)(5,2,1)(6,2,1)(6,2)(5,2,1)(6,2,1)(6,2)%
}

\NovelDisplay{M04715R01}{%
\text{展开第二项：}\quad(0)(1,1,1)(2,1,1)(3,1,1)(3,1)(3,1)(2,1,1)(3,1)(4,2,1)%
}
\NovelDisplay{M04715R02}{%
(5,2,1)(6,2,1)(6,2)(5,2,1)(6,2,1)(6,2)(5,2,1)(6,2,1)(6,2)(4,2,1)(5,2,1)%
}
\NovelDisplay{M04715R03}{%
(6,2,1)(6,2)(5,2,1)(6,2,1)(6,2)(5,2,1)(6,2,1)(6,2)(4,2,1)(5,2,1)(6,2,1)%
}
\NovelDisplay{M04715R04}{%
(6,2)(5,2,1)(6,2,1)(6,2)(5,2,1)(6,2,1)(6,2)%
}

参数的长短不能代替坏根判断，三份结果的保留段与新增段都要回到原列位置核对。

蔡兴坤终于没有先问别人写到第几列，他指着自己圈定的那一段，等同伴核过，才把后面接全。

第一次把四参数串完整讲清以后，卢本韦却没有立刻起身。

她听见自己方才用了几个熟悉的停顿，想再问蔡须坤一遍，那些停顿究竟有没有对应到M侧真正的层次。蔡须坤将纸转来，两人从外向内重新指过，发现其中一处仍要更早收住括号。

霍雨浩坐到旁边，核过她们所指的作用范围，又把原矩阵递给薛峰。三人各自完成自己的那一份，才同后来的人重新对照。

卢本韦将原先那处口述改好，没有因为已经讲完一次，就让它沿着旧习惯留下去。蔡须坤看着她的补页，想起自己几个月前曾经盼着有一个足够整齐的右式，好让所有台阶一眼就能读完。

眼前的纸仍长，队伍也仍慢。她却能看见某一处错误被具体指出、由写下它的人亲手改正，心里那份只盯着更高处的急切，也就不必每回都压到同伴身上。

\NovelDisplay{M04716R01}{%
(0)(1,1,1)(2,1,1)(3,1,1)(3,1)(3,1)(2,1,1)(3,1,1)(2,1,1)(3,1)(2)%
}

\NovelDisplay{M04717R01}{%
= \psi(I(1,0,1,0))%
}

\NovelDisplay{M04718R01}{%
= \psi(M^{M^2+1})%
}

\NovelDisplay{M04719R01}{%
(0)(1,1,1)(2,1,1)(3,1,1)(3,1)(3,1)(2,1,1)(3,1,1)(2,1,1)(3,1,1)%
}

\NovelDisplay{M04720R01}{%
= \psi(I(1,0,1,\omega))%
}

\NovelDisplay{M04721R01}{%
= \psi(M^{M^2+1}\times \omega)%
}

\NovelDisplay{M04722R01}{%
(0)(1,1,1)(2,1,1)(3,1,1)(3,1)(3,1)(2,1,1)(3,1,1)(2,1,1)(3,1,1)%
}
\NovelDisplay{M04722R02}{%
(2,1,1)(3,1)(2)%
}

\NovelDisplay{M04723R01}{%
= \psi(I(1,0,2,0))%
}

\NovelDisplay{M04724R01}{%
= \psi(M^{M^2+2})%
}

\NovelDisplay{M04725R01}{%
(0)(1,1,1)(2,1,1)(3,1,1)(3,1)(3,1)(2,1,1)(3,1,1)(3,1)(2)%
}

\NovelDisplay{M04726R01}{%
= \psi(I(1,1,0,0))%
}

\NovelDisplay{M04727R01}{%
= \psi(M^{M^2+M})%
}

四参数中的不同位置对应指数里的不同加项，逗号分隔的参数因此不能随意互换。

练到更长的参数串，诺当先问了省写的问题。

夏夜星空让他保留讲义给出的写法，没有拿那个省写符号另编一套展开。原矩阵仍完整留着，能否独立取项，仍得由原版规则核对。

诺当听完，将自己为了省事画出的几个省略点擦掉。他这一次没有觉得是退回去了，反而终于知道，纸上哪些是已有的约定，哪些只是自己还没有说明的空处。

\NovelDisplay{M04728R01}{%
(0)(1,1,1)(2,1,1)(3,1,1)(3,1)(3,1)(2,1,1)(3,1,1)(3,1)(2,1,1)(3,1,1)%
}
\NovelDisplay{M04728R02}{%
(2,1,1)(3,1)(2)%
}

\NovelDisplay{M04729R01}{%
= \psi(I(1,1,1,0))%
}

\NovelDisplay{M04730R01}{%
= \psi(M^{M^2+M+1})%
}

\NovelDisplay{M04731R01}{%
(0)(1,1,1)(2,1,1)(3,1,1)(3,1)(3,1)(2,1,1)(3,1,1)(3,1)(2,1,1)(3,1,1)%
}
\NovelDisplay{M04731R02}{%
(3,1)(2)%
}

\NovelDisplay{M04732R01}{%
= \psi(I(1,2,0,0))%
}

\NovelDisplay{M04733R01}{%
= \psi(M^{M^2+M\times 2})%
}

\NovelDisplay{M04734R01}{%
(0)(1,1,1)(2,1,1)(3,1,1)(3,1)(3,1)(2,1,1)(3,1,1)(3,1)(3)%
}

\NovelDisplay{M04735R01}{%
= \psi(I(1,\omega,0,0))%
}

\NovelDisplay{M04736R01}{%
= \psi(M^{M^2+M\times \omega})%
}

\NovelDisplay{M04737R01}{%
(0)(1,1,1)(2,1,1)(3,1,1)(3,1)(3,1)(2,1,1)(3,1,1)(3,1)(3,1)(2)%
}

\NovelDisplay{M04738R01}{%
= \psi(I(2,0,0,0))%
}

\NovelDisplay{M04739R01}{%
= \psi(M^{M^2\times 2})%
}

指数里的后继、倍数和ω倍数依原式并列，蔡兴坤先指出自己所改的参数，才说明它在哪一层起作用。

\NovelDisplay{M04740R01}{%
(0)(1,1,1)(2,1,1)(3,1,1)(3,1)(3,1)(2,1,1)(3,1,1)(3,1)(3,1)(2,1,1)%
}
\NovelDisplay{M04740R02}{%
(3,1,1)(3,1)(3,1)(2)%
}

\NovelDisplay{M04741R01}{%
= \psi(I(3,0,0,0))%
}

\NovelDisplay{M04742R01}{%
= \psi(M^{M^2\times 3})%
}

\NovelDisplay{M04743R01}{%
(0)(1,1,1)(2,1,1)(3,1,1)(3,1)(3,1)(3)%
}

\NovelDisplay{M04744R01}{%
= \psi(I(\omega,0,0,0))%
}

\NovelDisplay{M04745R01}{%
= \psi(M^{M^2\times \omega})%
}

\NovelDisplay{M04746R01}{%
(0)(1,1,1)(2,1,1)(3,1,1)(3,1)(3,1)(3,1)(2)%
}

\NovelDisplay{M04747R01}{%
= \psi(I(1,0,0,0,0))%
}

\NovelDisplay{M04748R01}{%
= \psi(M^{M^3})%
}

\NovelDisplay{M04749R01}{%
(0)(1,1,1)(2,1,1)(3,1,1)(3,1)(4)%
}

\NovelDisplay{M04750R01}{%
= \psi(I(1 @ \omega))%
}

\NovelDisplay{M04751R01}{%
= \psi(M^{M^\omega})%
}

参数长度继续增加并进入原稿的省写后，M侧仍保留完整的幂层次，不由省写符号自行推造展开。

这一条只有十五列，末列最高非零位置在第一行，坏根落在第二列；第二至第十四列共十三列，删去末列后保留十四列。

蔡须坤指出，最高非零位置既在第一行，就不能给它硬套另一个例子里的第一行提升差；霍雨浩请大家由同一原矩阵独立求三项。

\NovelDisplay{M04752R01}{%
\text{展开第零项：}\quad(0)(1,1,1)(2,1,1)(3,1,1)(3,1)(3,1)(2,1,1)(3,1,1)(3,1)%
}
\NovelDisplay{M04752R02}{%
(3,1)(2,1,1)(3,1,1)(3,1)(3,1)%
}

\NovelDisplay{M04753R01}{%
\text{展开第一项：}\quad(0)(1,1,1)(2,1,1)(3,1,1)(3,1)(3,1)(2,1,1)(3,1,1)(3,1)%
}
\NovelDisplay{M04753R02}{%
(3,1)(2,1,1)(3,1,1)(3,1)(3,1)(1,1,1)(2,1,1)(3,1,1)(3,1)(3,1)(2,1,1)%
}
\NovelDisplay{M04753R03}{%
(3,1,1)(3,1)(3,1)(2,1,1)(3,1,1)(3,1)(3,1)%
}

\NovelDisplay{M04754R01}{%
\text{展开第二项：}\quad(0)(1,1,1)(2,1,1)(3,1,1)(3,1)(3,1)(2,1,1)(3,1,1)(3,1)%
}
\NovelDisplay{M04754R02}{%
(3,1)(2,1,1)(3,1,1)(3,1)(3,1)(1,1,1)(2,1,1)(3,1,1)(3,1)(3,1)(2,1,1)%
}
\NovelDisplay{M04754R03}{%
(3,1,1)(3,1)(3,1)(2,1,1)(3,1,1)(3,1)(3,1)(1,1,1)(2,1,1)(3,1,1)(3,1)(3,1)%
}
\NovelDisplay{M04754R04}{%
(2,1,1)(3,1,1)(3,1)(3,1)(2,1,1)(3,1,1)(3,1)(3,1)%
}

三项分别有十四列、二十七列和四十列，生成段虽长，取项规则仍由本条原矩阵决定。

卢本韦将最后那条四十列的结果收好，又请蔡须坤把原式留下。第二天上山，两份都要由自己重新看清。

I侧的对照在这里没有继续逐条列出，卢本韦便把I册暂时合到能确认的那一页。

她没有将M侧的长幂塔勉强译成自己想要的短串。蔡须坤逐层指给她看，霍雨浩等她先提出一个完整的问题，再将最外面的乘号与指数范围分开。

卢本韦认真听完，在补页上留下M侧原式。

``这一处先这样记。''她说。

蔡须坤点头，没有把她熟悉的那本册子从桌上拿走。

\NovelDisplay{M04755R01}{%
(0)(1,1,1)(2,1,1)(3,1,1)(3,1)(4)(2,1,1)%
}

\NovelDisplay{M04756R01}{%
= \psi(M^{M^\omega}\times \omega)%
}

\NovelDisplay{M04757R01}{%
(0)(1,1,1)(2,1,1)(3,1,1)(3,1)(4)(2,1,1)(3,1)(2)%
}

\NovelDisplay{M04758R01}{%
= \psi(M^{M^\omega+1})%
}

\NovelDisplay{M04759R01}{%
(0)(1,1,1)(2,1,1)(3,1,1)(3,1)(4)(2,1,1)(3,1,1)(3)%
}

\NovelDisplay{M04760R01}{%
= \psi(M^{M^\omega+\omega})%
}

\NovelDisplay{M04761R01}{%
(0)(1,1,1)(2,1,1)(3,1,1)(3,1)(4)(2,1,1)(3,1,1)(3,1)(2)%
}

\NovelDisplay{M04762R01}{%
= \psi(M^{M^\omega+M})%
}

这几条只列M侧的对应便按M侧读，幂外乘ω与指数内部的加项仍须明确区分。

\NovelDisplay{M04763R01}{%
(0)(1,1,1)(2,1,1)(3,1,1)(3,1)(4)(2,1,1)(3,1,1)(3,1)(2,1,1)(3,1,1)%
}

\NovelDisplay{M04764R01}{%
= \psi(M^{M^\omega+M}\times \omega)%
}

\NovelDisplay{M04765R01}{%
(0)(1,1,1)(2,1,1)(3,1,1)(3,1)(4)(2,1,1)(3,1,1)(3,1)(2,1,1)(3,1,1)%
}
\NovelDisplay{M04765R02}{%
(2,1,1)(3,1)(2)%
}

\NovelDisplay{M04766R01}{%
= \psi(M^{M^\omega+M+1})%
}

\NovelDisplay{M04767R01}{%
(0)(1,1,1)(2,1,1)(3,1,1)(3,1)(4)(2,1,1)(3,1,1)(3,1)(2,1,1)(3,1,1)%
}
\NovelDisplay{M04767R02}{%
(3,1)(2)%
}

\NovelDisplay{M04768R01}{%
= \psi(M^{M^\omega+M\times 2})%
}

\NovelDisplay{M04769R01}{%
(0)(1,1,1)(2,1,1)(3,1,1)(3,1)(4)(2,1,1)(3,1,1)(3,1)(3)%
}

\NovelDisplay{M04770R01}{%
= \psi(M^{M^\omega+M\times \omega})%
}

相同的嵌套幂背景后出现有限尾部与不同M幂项，先看最外层指数的括号，才能分清每次改动。

\NovelDisplay{M04771R01}{%
(0)(1,1,1)(2,1,1)(3,1,1)(3,1)(4)(2,1,1)(3,1,1)(3,1)(3,1)(2)%
}

\NovelDisplay{M04772R01}{%
= \psi(M^{M^\omega+M^2})%
}

\NovelDisplay{M04773R01}{%
(0)(1,1,1)(2,1,1)(3,1,1)(3,1)(4)(2,1,1)(3,1,1)(3,1)(3,1)(3,1)(2)%
}

\NovelDisplay{M04774R01}{%
= \psi(M^{M^\omega+M^3})%
}

\NovelDisplay{M04775R01}{%
(0)(1,1,1)(2,1,1)(3,1,1)(3,1)(4)(2,1,1)(3,1,1)(3,1)(4)%
}

\NovelDisplay{M04776R01}{%
= \psi(M^{M^\omega\times 2})%
}

\NovelDisplay{M04777R01}{%
(0)(1,1,1)(2,1,1)(3,1,1)(3,1)(4)(3)%
}

\NovelDisplay{M04778R01}{%
= \psi(M^{M^\omega\times \omega})%
}

更高旧幂进入指数、整个内层指数变成倍数以及再添一层指数分别列出，不能把它们口头缩成同一句。

最深的几层，曾让小舞连着两晚看不下去。

她觉得无论从哪一边读，最后都只剩一片重复的M。唐三陪着看，自己也渐渐分不清，纸页翻动的声音越来越急。

薛峰在门口听见，进来把灯移近，先让他们只指出最外一层。里面暂时盖住，外层说清以后，再往里露出一层。

两人慢慢往下读，第一次停住时，便知道停的是哪一个指数，而不是整座幂塔都塌了。

\NovelDisplay{M04779R01}{%
(0)(1,1,1)(2,1,1)(3,1,1)(3,1)(4)(3,1)(2)%
}

\NovelDisplay{M04780R01}{%
= \psi(M^{M^{\omega+1}})%
}

\NovelDisplay{M04781R01}{%
(0)(1,1,1)(2,1,1)(3,1,1)(3,1)(4)(3,1)(3,1)(2)%
}

\NovelDisplay{M04782R01}{%
= \psi(M^{M^{\omega+2}})%
}

\NovelDisplay{M04783R01}{%
(0)(1,1,1)(2,1,1)(3,1,1)(3,1)(4)(3,1)(4)%
}

\NovelDisplay{M04784R01}{%
= \psi(M^{M^{\omega\times 2}})%
}

\NovelDisplay{M04785R01}{%
(0)(1,1,1)(2,1,1)(3,1,1)(3,1)(4)(4)%
}

\NovelDisplay{M04786R01}{%
= \psi(M^{M^{\psi(2)}})%
}

幂塔中较内层的指数开始变化时，蔡须坤从外向内指出层次，再由卢本韦检查自己遗漏的括号。

最后一周，叶阳提出换一种互问的次序。

先不由蔡须坤或霍雨浩带着大家读，而是每个人从自己最不确定的一处开口；能回答的人也先说自己能确认到哪里，超过的地方便一起回到原式。

第一天做得很慢。诺当的问题还没问完整，蔡兴坤已经想起自己的另一处，奥斯卡提醒两人别同时说，自己又差点把原矩阵压到补给袋下面。

小舞忍不住笑，笑过却主动先听诺当把话说完。她指出其中一个括号，诺当从那里重读，随后也听出了她前一天尚未分清的指数范围。

蔡须坤坐在另一侧，没有因为自己已经看清就替他们把后面的式子一口讲过。薛峰将需要复核的位置记在空页，霍雨浩等讨论停下，再核同一问题有没有被混成两种对象。

那天谁也没有多登几米，收纸时却有几张原先只敢压在最底下的疑问，终于被它们的主人堂堂正正放到桌面中央。

\NovelDisplay{M04787R01}{%
(0)(1,1,1)(2,1,1)(3,1,1)(3,1)(4,1)%
}

\NovelDisplay{M04788R01}{%
= \psi(M^{M^{\psi_{M}(M)}})%
}

\NovelDisplay{M04789R01}{%
(0)(1,1,1)(2,1,1)(3,1,1)(3,1)(4,1)(2,1,1)(3,1,1)%
}

\NovelDisplay{M04790R01}{%
= \psi(M^{M^{M}}\times \omega)%
}

\NovelDisplay{M04791R01}{%
(0)(1,1,1)(2,1,1)(3,1,1)(3,1)(4,1)(2,1,1)(3,1,1)(2,1,1)(3,1)(2)%
}

\NovelDisplay{M04792R01}{%
= \psi(M^{M^{M}+1})%
}

\NovelDisplay{M04793R01}{%
(0)(1,1,1)(2,1,1)(3,1,1)(3,1)(4,1)(2,1,1)(3,1,1)(3,1)(2)%
}

\NovelDisplay{M04794R01}{%
= \psi(M^{M^{M}+M})%
}

同一幂塔的外乘与最外一层指数里的加项相邻出现，二者的位置差别必须在纸上完整保留。

\NovelDisplay{M04795R01}{%
(0)(1,1,1)(2,1,1)(3,1,1)(3,1)(4,1)(2,1,1)(3,1,1)(3,1)(3,1)(2)%
}

\NovelDisplay{M04796R01}{%
= \psi(M^{M^{M}+M^2})%
}

\NovelDisplay{M04797R01}{%
(0)(1,1,1)(2,1,1)(3,1,1)(3,1)(4,1)(2,1,1)(3,1,1)(3,1)(4,1)(2)%
}

\NovelDisplay{M04798R01}{%
= \psi(M^{M^{M}\times 2})%
}

\NovelDisplay{M04799R01}{%
(0)(1,1,1)(2,1,1)(3,1,1)(3,1)(4,1)(3,1)(2)%
}

\NovelDisplay{M04800R01}{%
= \psi(M^{M^{M+1}})%
}

\NovelDisplay{M04801R01}{%
(0)(1,1,1)(2,1,1)(3,1,1)(3,1)(4,1)(3,1)(4,1)(2)%
}

\NovelDisplay{M04802R01}{%
= \psi(M^{M^{M\times 2}})%
}

进入更高嵌套以后，每一次后继、倍数和平方都须找回所属指数，不能看见相同M便按重复略过。

最后几日，十一人把各自的原矩阵摆在石桩下。

蔡须坤、霍雨浩、卢本韦与薛峰仍较领先，等后来的人提问时，却不再从自己已读到的最深处开始说。诺当先说自己能确认的前段，叶阳指出问题，古月娜便从他们真正停住的位置重新接。

奥斯卡守着水和补给，谁要把纸放下歇一会儿，他便将能坐的地方让出来。没有谁的那一份，因为别人已经说清，就被算作一并完成。

\NovelDisplay{M04803R01}{%
(0)(1,1,1)(2,1,1)(3,1,1)(3,1)(4,1)(4,1)(2)%
}

\NovelDisplay{M04804R01}{%
= \psi(M^{M^{M^2}})%
}

\NovelDisplay{M04805R01}{%
(0)(1,1,1)(2,1,1)(3,1,1)(3,1)(4,1)(4,1)(4,1)(2)%
}

\NovelDisplay{M04806R01}{%
= \psi(M^{M^{M^3}})%
}

\NovelDisplay{M04807R01}{%
(0)(1,1,1)(2,1,1)(3,1,1)(3,1)(4,1)(5)%
}

\NovelDisplay{M04808R01}{%
= \psi(M^{M^{M^\omega}})%
}

最后几条逐步加深指数层次，队员仍从同一原矩阵独立取项，没有把幂塔向上再写一层当作通用证明。

\NovelDisplay{M04809R01}{%
(0)(1,1,1)(2,1,1)(3,1,1)(3,1)(4,1)(5,1)(2)%
}

\NovelDisplay{M04810R01}{%
= \psi(M^{M^{M^M}})%
}

\NovelDisplay{M04811R01}{%
(0)(1,1,1)(2,1,1)(3,1,1)(3,1)(4,1)(5,1)(6,1)(2)%
}

\NovelDisplay{M04812R01}{%
= \psi(M^{M^{M^{M^M}}})%
}

\NovelDisplay{M04813R01}{%
\textcolor{red}{(0)(1,1,1)(2,1,1)(3,1,1)(3,1)(4,2)}%
}

\NovelDisplay{M04814R01}{%
\textcolor{red}{= \psi(\varepsilon_{M+1})}%
}

\NovelDisplay{M04815R01}{%
\textcolor{red}{= \psi(\Omega_{M+1})}%
}

\NovelDisplay{M04816R01}{%
\textcolor{red}{= \psi(M_2)}%
}

末组照原稿接到三个右侧表示，SRO阶段的结束只表示这一段分析走完，第一考的五百米仍须继续攀登。

四百一十米的刻痕，是在第七个月末摸到的。

蔡须坤刚站稳，便回头看下方。卢本韦收起纸，霍雨浩让出位置，薛峰还在等唐三把最后一处独立说完。

小舞上来以后，坐了好一会儿，才有力气把水壶拧开。她望着几个月前还觉得不过一小段的石阶，忽然问，今天算不算终于把漏下的路接回来了。

蔡须坤将本子翻到开头。下标嵌套、欠层的条件、不同参数、越深越长的幂，都有它们真正被问清的地方；其中几条重复与仍需回看的对应，也照原样留着，没有被她顺手删去。

``这一段接回来了。''她说。

五百米仍在上方。

夏夜星空拿出下一叠纸，最前面的背景与这一页的末尾接上。蔡须坤先核了那一处，没有急着站起，卢本韦也将自己的补页收进册里。

奥斯卡把最后一份补给递给叶阳，蔡兴坤护着包里的黑子鸡，笑着向旁边挪了一点。

他们歇够这一口气，再走剩下的九十米。

\chapter{第一百二十六章：走过雨季}\label{ux7b2cux4e00ux767eux4e8cux5341ux516dux7ae0ux8d70ux8fc7ux96e8ux5b63}

三百七十米之后，他们没有立刻看见另一座横台。

天晴时能看见更高处的石阶，到了下午，雾又将它们一层层吞回去。小舞将五百米画在纸上，隔几天便拿出来看；唐三却发现，那道横线与他们之间，多的并不只是一百三十米。

他们又用去许多天，把前一阶段尚未接稳的部分补齐。夏夜星空没有把不同阶段的讲义首尾一贴，就说中间也已经走过。每个人都留下了自己的演算页，有些原先只靠同伴提醒的地方，必须合上笔记重新说清。

蔡须坤在封套里又摸到那张理想矩阵的纸。它仍然没有展开规则。她将纸收回去，取出真正用过的本子，在扉页添上这次登山的日期。

七十米之前，他们已经用了两个月。后来那些长短不同的停顿，也都留在山图里。她以前嫌这些日期刺眼，到了雨季才明白，不能只记住向上走的日子，把退回来的全当成没有发生。

新的一页，以原来已经核到的SRO阶段为起点。卢本韦将I册摊在左边，蔡须坤把M册放在右边，霍雨浩拿走了中间那块替谁都能压纸的小石头。

``今天先听问题。''他说。

薛峰将原矩阵转向大家。十一个人挤在雨棚下，先前每次都抢着往后翻的小舞，这一次先找出了第一列。

新背景中的字母并没有使压力轻一点。

蔡须坤第一次接起时，膝盖便往下沉。她原想像先前那样，将读熟的前段迅速带过去，脚底却在一个不起眼的停顿处黏住，后面几层力道同时往回扯。

霍雨浩扶住她的手肘，没有将魂力探进去替她压平。他等她自己撤回来，才问，刚才哪一处是她认清了，哪一处只是以为同前一页差不多。

她望着脚边的水，过了一会儿，把原先一口念过的那一段重新圈出来。

``这里。''

卢本韦将自己的本子收近一些，给她留出完整的位置。

\NovelDisplay{M04817R01}{%
(0)(1,1,1)(2,1,1)(3,1,1)(3,1)(4,2)%
}

\NovelDisplay{M04818R01}{%
= \psi(\Omega_{M+1})%
}

\NovelDisplay{M04819R01}{%
(0)(1,1,1)(2,1,1)(3,1,1)(3,1)(4,2)(1,1,1)(2,1,1)(3,1,1)(3,1)(4,2)%
}

\NovelDisplay{M04820R01}{%
= \psi(\Omega_{M+1}+\psi_{M}(\Omega_{M+1}))%
}

\NovelDisplay{M04821R01}{%
(0)(1,1,1)(2,1,1)(3,1,1)(3,1)(4,2)(2,1)(3,2)%
}

\NovelDisplay{M04822R01}{%
= \psi(\Omega_{M+1}+M)%
}

\NovelDisplay{M04823R01}{%
(0)(1,1,1)(2,1,1)(3,1,1)(3,1)(4,2)(2,1)(3,2,1)%
}

\NovelDisplay{M04824R01}{%
= \psi(\Omega_{M+1}+M\times \omega)%
}

霍雨浩先固定\allowbreak{}\NovelInline{I06092}{\Omega_{M+1}}\allowbreak{}这一背景，再让大家分清回接项、\allowbreak{}\NovelInline{I06093}{M}\allowbreak{}和\allowbreak{}\NovelInline{I06094}{M\times\omega}\allowbreak{}各自落在什么位置。

那几周，唐三曾以为自己已经找到了加快的办法。

他先把长式的共同部分抄成一页，后面每有变化，只留下不同的尾部。纸一下薄了许多，问起几条相邻的对应，他也能很快翻出来。

直到小舞让他合起本子，从一条完整矩阵重新说明，他才发现自己只记住了尾部的排列。问到它为什么仍接在原来的背景下，脑子里先浮起来的竟是那页纸放在哪里。

他拿笔补，越补越乱，最后将整叠推开，说也许自己根本不适合这样学。奥斯卡刚进门，听见这句，便坐到门槛上，将手里的补给搁在一旁。

``你刚才哪一句知道，哪一句不知道？''

唐三不想答。小舞把那条完整原式取回，遮住右边，只让他从第一列说。说了一阵，他自己停住，指出真正没接清的那个位置。

原来空着的只有这一处，先前独立核过的前段并没有一齐坏掉。

他坐回桌边，重新抄完整式子。奥斯卡把补给递过去，没说一听就懂才算有天赋，只问他晚饭要不要等小舞一块吃。唐三点头，小舞便将旁边那张空纸留给他，自己也从原式重做。

后来整理本子时，他仍留下了那份短记。只是每一条前面，都添回能够找全原矩阵的位置。

\NovelDisplay{M04825R01}{%
(0)(1,1,1)(2,1,1)(3,1,1)(3,1)(4,2)(2,1)(3,2,1)(4,2,1)(5,2)(4)%
}

\NovelDisplay{M04826R01}{%
= \psi(\Omega_{M+1}+M^2)%
}

\NovelDisplay{M04827R01}{%
(0)(1,1,1)(2,1,1)(3,1,1)(3,1)(4,2)(2,1)(3,2,1)(4,2,1)(5,2,1)%
}

\NovelDisplay{M04828R01}{%
= \psi(\Omega_{M+1}+M^2\times \omega)%
}

\NovelDisplay{M04829R01}{%
(0)(1,1,1)(2,1,1)(3,1,1)(3,1)(4,2)(2,1)(3,2,1)(4,2,1)(5,2,1)(4,2,1)%
}
\NovelDisplay{M04829R02}{%
(5,2)(4)%
}

\NovelDisplay{M04830R01}{%
= \psi(\Omega_{M+1}+M^3)%
}

\NovelDisplay{M04831R01}{%
(0)(1,1,1)(2,1,1)(3,1,1)(3,1)(4,2)(2,1)(3,2,1)(4,2,1)(5,2,1)(5)%
}

\NovelDisplay{M04832R01}{%
= \psi(\Omega_{M+1}+M^\omega)%
}

蔡须坤从平方读到有限幂与ω次幂，提醒唐三矩阵里多出的一段要按原式逐列核对，不能只凭幂次名称猜结果。

第一场大雨下来，棚顶漏了三处。

奥斯卡拿锅去接水，诺当在另一边搬本子，蔡兴坤则把黑子鸡往衣襟里塞。小舞护着自己的纸，嫌大家挪得慢，伸手一推，最上面那叠便滑到薛峰怀里。

薛峰接住了纸，却没接住已经开始说的下一句。他断在那里，怎么也找不回原来的次序，脸一下涨红。

``先别念。''古月娜说。

她用一根干树枝指回原矩阵，让他只说这一处。雨还在敲锅底，等到他把整句说完，奥斯卡才将满掉的那只锅换出去。

那天下午，他们没向上走。十一人的纸却终于不再挤成一摞，而是留出了能够轮流独立说明的位置。

\NovelDisplay{M04833R01}{%
(0)(1,1,1)(2,1,1)(3,1,1)(3,1)(4,2)(2,1)(3,2,1)(4,2,1)(5,2,1)(5,1)(1,1,1)%
}
\NovelDisplay{M04833R02}{%
(2,1,1)(3,1,1)(3,1)(4,2)%
}

\NovelDisplay{M04834R01}{%
= \psi(\Omega_{M+1}+M^{\psi_{M}(\Omega_{M+1})})%
}

\NovelDisplay{M04835R01}{%
(0)(1,1,1)(2,1,1)(3,1,1)(3,1)(4,2)(2,1)(3,2,1)(4,2,1)(5,2,1)(5,1)(1,1,1)%
}
\NovelDisplay{M04835R02}{%
(2,1,1)(3,1,1)(3,1)(4,2)(2,1)(3,2,1)(4,2,1)(5,2,1)(5,1)%
}

\NovelDisplay{M04836R01}{%
= \psi(\Omega_{M+1}+M^{\psi_{M}(\Omega_{M+1})}+\psi_{\psi_{M}(\Omega_{M+1}+M)}(\Omega_{M+1}+M^\Omega))%
}

\NovelDisplay{M04837R01}{%
(0)(1,1,1)(2,1,1)(3,1,1)(3,1)(4,2)(2,1)(3,2,1)(4,2,1)(5,2,1)(5,1)(2)%
}

\NovelDisplay{M04838R01}{%
= \psi(\Omega_{M+1}+M^{\psi_{M}(\Omega_{M+1})}+\psi_{\psi_{M}(\Omega_{M+1}+M)}(\Omega_{M+1}+M^{\psi_{M}(\Omega_{M+1})}+1))%
}

\NovelDisplay{M04839R01}{%
(0)(1,1,1)(2,1,1)(3,1,1)(3,1)(4,2)(2,1)(3,2,1)(4,2,1)(5,2,1)(5,1)(2,1)%
}
\NovelDisplay{M04839R02}{%
(1,1,1)(2,1,1)(3,1,1)(3,1)(4,2)%
}

\NovelDisplay{M04840R01}{%
= \psi(\Omega_{M+1}+M^{\psi_{M}(\Omega_{M+1})}+\psi_{\psi_{M}(\Omega_{M+1}+M)}(\Omega_{M+1}+M^{\psi_{M}(\Omega_{M+1})}+\psi_{M}(\Omega_{M+1})))%
}

进入嵌套以后，薛峰把幂指数与幂后面的加项分开，先认清括号属于哪一层，再追到最内侧的ψ。

蔡须坤后来发现，山阶的压力并非只让人难受。

平地上能凭熟悉一口带过的接续，到了这里，哪一层没有接稳，身体便常在相近的位置发涩。她起初只想用更大的魂力将涩处压过去，压得越重，后面的回扯也越重。

卢本韦陪她退回，等她将两个不同的接法分别说过，再请霍雨浩从原式核对。那处阻滞终于同纸上一个明确的问题对上，她才第一次觉得，逼自己停下的重压，也把从前模糊的一层分了出来。

身体的反馈不能替他们证明序数对应，却能让他们不再轻易忽略自己尚未说清的地方。她把这个发现告诉后来的人，让他们觉得停住时先找问题，不必每回都认作自己比同伴差。

\NovelDisplay{M04841R01}{%
(0)(1,1,1)(2,1,1)(3,1,1)(3,1)(4,2)(2,1)(3,2,1)(4,2,1)(5,2,1)(5,1)%
}
\NovelDisplay{M04841R02}{%
(2,1)(3,2)%
}

\NovelDisplay{M04842R01}{%
= \psi(\Omega_{M+1}+M^{\psi_{M}(\Omega_{M+1})}+\psi_{\psi_{M}(\Omega_{M+1}+M)}(\Omega_{M+1}+M^{\psi_{M}(\Omega_{M+1})}%
}
\NovelDisplay{M04842R02}{%
+\psi_{\psi_{M}(\Omega_{M+1}+M)}(\Omega_{M+1}+M)))%
}

\NovelDisplay{M04843R01}{%
(0)(1,1,1)(2,1,1)(3,1,1)(3,1)(4,2)(2,1)(3,2,1)(4,2,1)(5,2,1)(5,1)(2,1)%
}
\NovelDisplay{M04843R02}{%
(3,2,1)(4,2,1)(5,2,1)(5,1)%
}

\NovelDisplay{M04844R01}{%
= \psi(\Omega_{M+1}+M^{\psi_{M}(\Omega_{M+1})}+\psi_{\psi_{M}(\Omega_{M+1}+M)}(\Omega_{M+1}+M^{\psi_{M}(\Omega_{M+1})}%
}
\NovelDisplay{M04844R02}{%
+\psi_{\psi_{M}(\Omega_{M+1}+M)}(\Omega_{M+1}+M^\Omega)))%
}

\NovelDisplay{M04845R01}{%
(0)(1,1,1)(2,1,1)(3,1,1)(3,1)(4,2)(2,1)(3,2,1)(4,2,1)(5,2,1)(5,1)(3)%
}

\NovelDisplay{M04846R01}{%
= \psi(\Omega_{M+1}+M^{\psi_{M}(\Omega_{M+1})}+\psi_{\psi_{M}(\Omega_{M+1}+M)}(\Omega_{M+1}+M^{\psi_{M}(\Omega_{M+1})}%
}
\NovelDisplay{M04846R02}{%
+\psi_{\psi_{M}(\Omega_{M+1}+M)}(\Omega_{M+1}+M^{\psi_{M}(\Omega_{M+1})}+1)))%
}

\NovelDisplay{M04847R01}{%
(0)(1,1,1)(2,1,1)(3,1,1)(3,1)(4,2)(2,1)(3,2,1)(4,2,1)(5,2,1)(5,1)(3,2)%
}

\NovelDisplay{M04848R01}{%
= \psi(\Omega_{M+1}+M^{\psi_{M}(\Omega_{M+1})}+\psi_{M}(\Omega_{M+1}+M))%
}

这几式外框相近而内层作用对象不同，古月娜用笔尖逐处指出变化，没有把新添的ψ当成外面的普通系数。

\NovelDisplay{M04849R01}{%
(0)(1,1,1)(2,1,1)(3,1,1)(3,1)(4,2)(2,1)(3,2,1)(4,2,1)(5,2,1)(5,1)(3,2,1)%
}

\NovelDisplay{M04850R01}{%
= \psi(\Omega_{M+1}+M^{\psi_{M}(\Omega_{M+1})}+\psi_{M}(\Omega_{M+1}+M\times \omega))%
}

\NovelDisplay{M04851R01}{%
(0)(1,1,1)(2,1,1)(3,1,1)(3,1)(4,2)(2,1)(3,2,1)(4,2,1)(5,2,1)(5,1)(3,2,1)%
}
\NovelDisplay{M04851R02}{%
(4,2,1)(5,2,1)(5,1)%
}

\NovelDisplay{M04852R01}{%
= \psi(\Omega_{M+1}+M^{\psi_{M}(\Omega_{M+1})}+\psi_{M}(\Omega_{M+1}+M^\Omega))%
}

\NovelDisplay{M04853R01}{%
(0)(1,1,1)(2,1,1)(3,1,1)(3,1)(4,2)(2,1)(3,2,1)(4,2,1)(5,2,1)(5,1)(3,2,1)%
}
\NovelDisplay{M04853R02}{%
(4,2,1)(5,2,1)(5,1)(3,2,1)(4,2,1)(5,2,1)(5,1)%
}

\NovelDisplay{M04854R01}{%
= \psi(\Omega_{M+1}+M^{\psi_{M}(\Omega_{M+1})}+\psi_{M}(\Omega_{M+1}+M^{\psi_{M}(\Omega_{M+1})})+\psi_{M}(\Omega_{M+1}+M^\Omega))%
}

\NovelDisplay{M04855R01}{%
(0)(1,1,1)(2,1,1)(3,1,1)(3,1)(4,2)(2,1)(3,2,1)(4,2,1)(5,2,1)(5,1)(4)%
}

\NovelDisplay{M04856R01}{%
= \psi(\Omega_{M+1}+M^{\psi_{M}(\Omega_{M+1})}+\psi_{\psi_{M}(\Omega_{M+1}+M^{\psi_{M}(\Omega_{M+1})}+M)}(1))%
}

卢本韦先找回共同背景，确认ω、Ω及新增的折叠项都仍在各自的位置上，才让小舞继续向右读。

三百八十米的刻痕，摸到时比小舞想的浅。

她站到石桩旁，先笑了一声，转身要叫唐三，才发现他还停在下面两级。唐三没有望她，正用湿透的笔尖将一个外层加号描清，问霍雨浩是不是自己方才又把它接进了幂里。

小舞坐下等。她原来觉得，到了这里总该立刻知道自己强了多少，石阶却没有替他们把体内的力量排成一串清清楚楚的刻度。压力仍在，只比先前多出一小段能够从容呼吸的地方。

唐三终于上来，第一眼看见的是妹妹留在石面的水印。

``我先上来的。''小舞说。

``知道。''

``那你下次别光盯着本子。''

唐三点头，坐到她旁边，没有急着再往上走。

\NovelDisplay{M04857R01}{%
(0)(1,1,1)(2,1,1)(3,1,1)(3,1)(4,2)(2,1)(3,2,1)(4,2,1)(5,2,1)(5,1)(4,2)%
}

\NovelDisplay{M04858R01}{%
= \psi(\Omega_{M+1}+M^{\psi_{M}(\Omega_{M+1})}+\psi_{\psi_{M}(\Omega_{M+1}+M^{\psi_{M}(\Omega_{M+1})}+M)}(\Omega_{M+1}+M))%
}

\NovelDisplay{M04859R01}{%
(0)(1,1,1)(2,1,1)(3,1,1)(3,1)(4,2)(2,1)(3,2,1)(4,2,1)(5,2,1)(5,1)(4,2)%
}
\NovelDisplay{M04859R02}{%
(3,2,1)(4,2,1)(5,2,1)(5,1)(1,1,1)(2,1,1)(3,1,1)(3,1)(4,2)%
}

\NovelDisplay{M04860R01}{%
= \psi(\Omega_{M+1}+M^{\psi_{M}(\Omega_{M+1})}+\psi_{\psi_{M}(\Omega_{M+1}+M^{\psi_{M}(\Omega_{M+1})}+M)}(\Omega_{M+1}+M^{\psi_{M}(\Omega_{M+1})}))%
}

\NovelDisplay{M04861R01}{%
(0)(1,1,1)(2,1,1)(3,1,1)(3,1)(4,2)(2,1)(3,2,1)(4,2,1)(5,2,1)(5,1)%
}
\NovelDisplay{M04861R02}{%
(4,2)(5,3)%
}

\NovelDisplay{M04862R01}{%
= \psi(\Omega_{M+1}+M^{\psi_{M}(\Omega_{M+1})}+M)%
}

\NovelDisplay{M04863R01}{%
(0)(1,1,1)(2,1,1)(3,1,1)(3,1)(4,2)(2,1)(3,2,1)(4,2,1)(5,2,1)(5,1)%
}
\NovelDisplay{M04863R02}{%
(4,2)(5,3,1)%
}

\NovelDisplay{M04864R01}{%
= \psi(\Omega_{M+1}+M^{\psi_{M}(\Omega_{M+1})}+M\times \omega)%
}

尾部重新落回\allowbreak{}\NovelInline{I06095}{M}\allowbreak{}及其ω倍数时，前面已经构成的幂项仍须完整保留，不能为了短写而一并擦去。

\NovelDisplay{M04865R01}{%
(0)(1,1,1)(2,1,1)(3,1,1)(3,1)(4,2)(2,1)(3,2,1)(4,2,1)(5,2,1)(5,1)(4,2)%
}
\NovelDisplay{M04865R02}{%
(5,3,1)(6,3,1)(7,1)(1,1,1)(2,1,1)(3,1,1)(3,1)(4,2)%
}

\NovelDisplay{M04866R01}{%
= \psi(\Omega_{M+1}+M^{\psi_{M}(\Omega_{M+1})}\times 2)%
}

\NovelDisplay{M04867R01}{%
(0)(1,1,1)(2,1,1)(3,1,1)(3,1)(4,2)(2,1)(3,2,1)(4,2,1)(5,2,1)(5,1)(4,2,1)%
}

\NovelDisplay{M04868R01}{%
= \psi(\Omega_{M+1}+M^{\psi_{M}(\Omega_{M+1})}\times \omega)%
}

\NovelDisplay{M04869R01}{%
(0)(1,1,1)(2,1,1)(3,1,1)(3,1)(4,2)(2,1)(3,2,1)(4,2,1)(5,2,1)(5,1)(4,2,1)%
}
\NovelDisplay{M04869R02}{%
(5,1)(1,1,1)(2,1,1)(3,1,1)(3,1)(4,2)%
}

\NovelDisplay{M04870R01}{%
= \psi(\Omega_{M+1}+M^{\psi_{M}(\Omega_{M+1})}\times \psi_{M}(\Omega_{M+1}))%
}

\NovelDisplay{M04871R01}{%
(0)(1,1,1)(2,1,1)(3,1,1)(3,1)(4,2)(2,1)(3,2,1)(4,2,1)(5,2,1)(5,1)%
}
\NovelDisplay{M04871R02}{%
(4,2,1)(5,2)%
}

\NovelDisplay{M04872R01}{%
= \psi(\Omega_{M+1}+M^{\psi_{M}(\Omega_{M+1})}\times \psi_{M}(\Omega_{M+1}+M))%
}

这一组先改变幂项外面的乘数，唐三把乘号圈在指数框之外，终于没有再将两种增长写成同一处改动。

\NovelDisplay{M04873R01}{%
(0)(1,1,1)(2,1,1)(3,1,1)(3,1)(4,2)(2,1)(3,2,1)(4,2,1)(5,2,1)(5,1)(4,2,1)%
}
\NovelDisplay{M04873R02}{%
(5,2)(3,2,1)(4,2,1)(5,2,1)(5,1)(1,1,1)(2,1,1)(3,1,1)(3,1)(4,2)%
}

\NovelDisplay{M04874R01}{%
= \psi(\Omega_{M+1}+M^{\psi_{M}(\Omega_{M+1})}\times \psi_{M}(\Omega_{M+1}+M^{\psi_{M}(\Omega_{M+1})}))%
}

\NovelDisplay{M04875R01}{%
(0)(1,1,1)(2,1,1)(3,1,1)(3,1)(4,2)(2,1)(3,2,1)(4,2,1)(5,2,1)(5,1)(4,2,1)%
}
\NovelDisplay{M04875R02}{%
(5,2)(3,2,1)(4,2,1)(5,2,1)(5,1)(4,2,1)(5,2)%
}

\NovelDisplay{M04876R01}{%
= \psi(\Omega_{M+1}+M^{\psi_{M}(\Omega_{M+1})}\times \psi_{M}(\Omega_{M+1}+M^{\psi_{M}(\Omega_{M+1})}\times\psi_{M}(\Omega_{M+1}+M)))%
}

\NovelDisplay{M04877R01}{%
(0)(1,1,1)(2,1,1)(3,1,1)(3,1)(4,2)(2,1)(3,2,1)(4,2,1)(5,2,1)(5,1)%
}
\NovelDisplay{M04877R02}{%
(4,2,1)(5,2)(4)%
}

\NovelDisplay{M04878R01}{%
= \psi(\Omega_{M+1}+M^{\psi_{M}(\Omega_{M+1})+1})%
}

\NovelDisplay{M04879R01}{%
(0)(1,1,1)(2,1,1)(3,1,1)(3,1)(4,2)(2,1)(3,2,1)(4,2,1)(5,2,1)(5,1)(4,2,1)%
}
\NovelDisplay{M04879R02}{%
(5,2)(4,2)(5,3)%
}

\NovelDisplay{M04880R01}{%
= \psi(\Omega_{M+1}+M^{\psi_{M}(\Omega_{M+1})+1}+M)%
}

卢本韦与蔡须坤互换笔记，确认有的变化进入乘数的ψ内部，有的进入幂指数，外观接近也要分别回退。

霍雨浩没有让大家拿右侧表示猜矩阵的后继，而是重新从原矩阵独立求项；小舞先写第零项，唐三与薛峰再各自求出第一项、第二项。

蔡须坤把十九列从左往右标过。末列是\allowbreak{}\NovelInline{I06096}{(5,2)}\allowbreak{}，按省略补零的写法，最高非零位置在第二行；沿这一行的父链回找，生成所用的坏根落在第七列\allowbreak{}\NovelInline{I06097}{(2,1)}\allowbreak{}。

她先将末列盖住，留下前十八列，再把第七列至第十八列圈成完整的十二列。第一行的提升差为三，符合提升条件的位置随所加的那一份增加；第二、第三行不因此一起加三。

薛峰起初只圈了后半段，霍雨浩没有替他画好，让他从末列的第二行再追一次。他一直追回第七列，才将先前漏掉的那几列接回原来的顺序。

小舞将保留段、待生成的整段分别指给唐三看，自己没有照着他的纸下笔。

\NovelDisplay{M04881R01}{%
\text{展开第零项：}\quad(0)(1,1,1)(2,1,1)(3,1,1)(3,1)(4,2)(2,1)(3,2,1)(4,2,1)%
}
\NovelDisplay{M04881R02}{%
(5,2,1)(5,1)(4,2,1)(5,2)(3,2,1)(4,2,1)(5,2,1)(5,1)(4,2,1)%
}

\NovelDisplay{M04882R01}{%
\text{展开第一项：}\quad(0)(1,1,1)(2,1,1)(3,1,1)(3,1)(4,2)(2,1)(3,2,1)(4,2,1)%
}
\NovelDisplay{M04882R02}{%
(5,2,1)(5,1)(4,2,1)(5,2)(3,2,1)(4,2,1)(5,2,1)(5,1)(4,2,1)(5,1)(6,2,1)%
}
\NovelDisplay{M04882R03}{%
(7,2,1)(8,2,1)(8,1)(7,2,1)(8,2)(6,2,1)(7,2,1)(8,2,1)(8,1)(7,2,1)%
}

\NovelDisplay{M04883R01}{%
\text{展开第二项：}\quad(0)(1,1,1)(2,1,1)(3,1,1)(3,1)(4,2)(2,1)(3,2,1)(4,2,1)%
}
\NovelDisplay{M04883R02}{%
(5,2,1)(5,1)(4,2,1)(5,2)(3,2,1)(4,2,1)(5,2,1)(5,1)(4,2,1)(5,1)(6,2,1)%
}
\NovelDisplay{M04883R03}{%
(7,2,1)(8,2,1)(8,1)(7,2,1)(8,2)(6,2,1)(7,2,1)(8,2,1)(8,1)(7,2,1)(8,1)%
}
\NovelDisplay{M04883R04}{%
(9,2,1)(10,2,1)(11,2,1)(11,1)(10,2,1)(11,2)(9,2,1)(10,2,1)(11,2,1)%
}
\NovelDisplay{M04883R05}{%
(11,1)(10,2,1)%
}

三项都从同一个原矩阵出发，不能将求第二项偷换成再展开刚得到的第一项。

蔡须坤逐列看过，终于指出自己方才沿用了旧生成段的那一处。

第零项十八列，第一项三十列，第二项四十二列。后面生成的列会很长，可与原式比较时，先看的仍是第一次不同的位置；第一份新段在该处已经向较低的一侧接去，不能只数总列数便说展开后更大。

唐三将这句话说给诺当听。诺当没有照抄，而是指着原来末列的位置，自己找出新段第一次与它不同的那一行。

过了三百八十米，蔡须坤又犯过相同的错。

她已经能将右式完整说出，真要合上本子，回到矩阵末列，却仍用了前一天那个生成位置。脚底随即一松，人反而往后仰。卢本韦拉住她，另一只手将M册扣了起来。

``再看，是看今天的原式。''

蔡须坤喘了两口气，想把本子拿回去。卢本韦没松手，问她能不能先指出末列。

她停了一会儿，才明白对方不是在拦她。前一天求出的几项就在纸上，太熟了，熟到她还没认真看今天的矩阵，手已经准备照旧接下去。

她退回平处，自己将两条原式并排放好。霍雨浩没有催，陪她等山风把纸面吹干，才重新开始。

\NovelDisplay{M04884R01}{%
(0)(1,1,1)(2,1,1)(3,1,1)(3,1)(4,2)(2,1)(3,2,1)(4,2,1)(5,2,1)(5,1)%
}
\NovelDisplay{M04884R02}{%
(4,2,1)(5,2,1)%
}

\NovelDisplay{M04885R01}{%
= \psi(\Omega_{M+1}+M^{\psi_{M}(\Omega_{M+1})+1}\times \omega)%
}

\NovelDisplay{M04886R01}{%
(0)(1,1,1)(2,1,1)(3,1,1)(3,1)(4,2)(2,1)(3,2,1)(4,2,1)(5,2,1)(5,1)(4,2,1)%
}
\NovelDisplay{M04886R02}{%
(5,2,1)(4,2,1)(5,2)(4)%
}

\NovelDisplay{M04887R01}{%
= \psi(\Omega_{M+1}+M^{\psi_{M}(\Omega_{M+1})+2})%
}

\NovelDisplay{M04888R01}{%
(0)(1,1,1)(2,1,1)(3,1,1)(3,1)(4,2)(2,1)(3,2,1)(4,2,1)(5,2,1)(5,1)%
}
\NovelDisplay{M04888R02}{%
(4,2,1)(5,2,1)(5)%
}

\NovelDisplay{M04889R01}{%
= \psi(\Omega_{M+1}+M^{\psi_{M}(\Omega_{M+1})+\omega})%
}

\NovelDisplay{M04890R01}{%
(0)(1,1,1)(2,1,1)(3,1,1)(3,1)(4,2)(2,1)(3,2,1)(4,2,1)(5,2,1)(5,1)(4,2,1)%
}
\NovelDisplay{M04890R02}{%
(5,2,1)(5,1)(1,1,1)(2,1,1)(3,1,1)(3,1)(4,2)%
}

\NovelDisplay{M04891R01}{%
= \psi(\Omega_{M+1}+M^{\psi_{M}(\Omega_{M+1})\times 2})%
}

这里的后继、ω与倍数出现在指数内，霍雨浩让大家先指出上标的范围，再谈整个式子的外层结构。

\NovelDisplay{M04892R01}{%
(0)(1,1,1)(2,1,1)(3,1,1)(3,1)(4,2)(2,1)(3,2,1)(4,2,1)(5,2,1)(5,2)%
}

\NovelDisplay{M04893R01}{%
= \psi(\Omega_{M+1}+M^{\psi_{M}(\Omega_{M+1}+M)})%
}

\NovelDisplay{M04894R01}{%
(0)(1,1,1)(2,1,1)(3,1,1)(3,1)(4,2)(2,1)(3,2,1)(4,2,1)(5,2,1)(5,2)(4)%
}

\NovelDisplay{M04895R01}{%
= \psi(\Omega_{M+1}+M^{M})%
}

\NovelDisplay{M04896R01}{%
(0)(1,1,1)(2,1,1)(3,1,1)(3,1)(4,2)(2,1)(3,2,1)(4,2,1)(5,2,1)(5,2)%
}
\NovelDisplay{M04896R02}{%
(4,2,1)(5,2)(6,3)%
}

\NovelDisplay{M04897R01}{%
= \psi(\Omega_{M+1}+M^{M}+M)%
}

\NovelDisplay{M04898R01}{%
(0)(1,1,1)(2,1,1)(3,1,1)(3,1)(4,2)(2,1)(3,2,1)(4,2,1)(5,2,1)(5,2)(4,2,1)%
}
\NovelDisplay{M04898R02}{%
(5,2)(6,3,1)(7,3,1)(8,3,1)(8,3)(7)%
}

\NovelDisplay{M04899R01}{%
= \psi(\Omega_{M+1}+M^{M}\times 2)%
}

从折叠后的指数转到\allowbreak{}\NovelInline{I06098}{M}\allowbreak{}自身，仍须依讲义给出的对应辨认，不能将看见同一个字母当成已经走过全部中间步骤。

\NovelDisplay{M04900R01}{%
(0)(1,1,1)(2,1,1)(3,1,1)(3,1)(4,2)(2,1)(3,2,1)(4,2,1)(5,2,1)(5,2)%
}
\NovelDisplay{M04900R02}{%
(4,2,1)(5,2,1)%
}

\NovelDisplay{M04901R01}{%
= \psi(\Omega_{M+1}+M^{M}\times \omega)%
}

\NovelDisplay{M04902R01}{%
(0)(1,1,1)(2,1,1)(3,1,1)(3,1)(4,2)(2,1)(3,2,1)(4,2,1)(5,2,1)(5,2)(4,2,1)%
}
\NovelDisplay{M04902R02}{%
(5,2,1)(4,2,1)(5,2)(4)%
}

\NovelDisplay{M04903R01}{%
= \psi(\Omega_{M+1}+M^{M+1})%
}

\NovelDisplay{M04904R01}{%
(0)(1,1,1)(2,1,1)(3,1,1)(3,1)(4,2)(2,1)(3,2,1)(4,2,1)(5,2,1)(5,2)(4,2,1)%
}
\NovelDisplay{M04904R02}{%
(5,2,1)(5,2)(4)%
}

\NovelDisplay{M04905R01}{%
= \psi(\Omega_{M+1}+M^{M\times 2})%
}

\NovelDisplay{M04906R01}{%
(0)(1,1,1)(2,1,1)(3,1,1)(3,1)(4,2)(2,1)(3,2,1)(4,2,1)(5,2,1)(5,2)%
}
\NovelDisplay{M04906R02}{%
(5,2)(4)%
}

\NovelDisplay{M04907R01}{%
= \psi(\Omega_{M+1}+M^{M^2})%
}

蔡须坤将幂外乘ω与指数加一分开说明，薛峰则沿着原矩阵找出两者实际改变的列。

第一个没有人受伤的晴天，他们却吵了一回。

蔡须坤想趁天气好，多接两组。唐三和诺当还留着问题，便说可以在下面慢慢核，让前面的人先走。蔡须坤答应得很快，背上包才发现，卢本韦仍坐在原处。

``你们先去。''唐三说。

卢本韦看着蔡须坤：``你是真的想先去，还是怕留下来，今天就又算没有上山？''

蔡须坤被问得一怔。她想说上面的路也重要，话到嘴边，却看见诺当将那张没有说清的纸往身后收了一下。

她把包放了下来。

这一回，带头的几个人换到后面。薛峰不再先讲自己的结果，先请诺当指出卡住的第一处；霍雨浩则陪唐三，把同一个问题从另一边重新问出来。

蔡须坤坐到卢本韦旁边，过了一会儿才说：``我怕停。''

卢本韦没有说她不该怕，只将原矩阵移到了两人中间。

\NovelDisplay{M04908R01}{%
(0)(1,1,1)(2,1,1)(3,1,1)(3,1)(4,2)(2,1)(3,2,1)(4,2,1)(5,2,1)(5,2)%
}
\NovelDisplay{M04908R02}{%
(6,2)(4)%
}

\NovelDisplay{M04909R01}{%
= \psi(\Omega_{M+1}+M^{M^M})%
}

\NovelDisplay{M04910R01}{%
(0)(1,1,1)(2,1,1)(3,1,1)(3,1)(4,2)(2,1)(3,2,1)(4,2,1)(5,2,1)(5,2)(6,3)%
}

\NovelDisplay{M04911R01}{%
= \psi(\Omega_{M+1}+\psi_{\Omega_{M+1}}(\Omega_{M+1}))%
}

\NovelDisplay{M04912R01}{%
(0)(1,1,1)(2,1,1)(3,1,1)(3,1)(4,2)(2,1)(3,2,1)(4,2,1)(5,2,1)(5,2)(6,3)%
}
\NovelDisplay{M04912R02}{%
(1,1,1)(2,1,1)(3,1,1)(3,1)(4,2)%
}

\NovelDisplay{M04913R01}{%
= \psi(\Omega_{M+1}+\psi_{\Omega_{M+1}}(\Omega_{M+1})+\psi_{M}(\Omega_{M+1}))%
}

\NovelDisplay{M04914R01}{%
(0)(1,1,1)(2,1,1)(3,1,1)(3,1)(4,2)(2,1)(3,2,1)(4,2,1)(5,2,1)(5,2)(6,3)%
}
\NovelDisplay{M04914R02}{%
(1,1,1)(2,1,1)(3,1,1)(3,1)(4,2)(2,1)(3,2)%
}

\NovelDisplay{M04915R01}{%
= \psi(\Omega_{M+1}+\psi_{\Omega_{M+1}}(\Omega_{M+1})+\psi_{\psi_{M}(\Omega_{M+1}+M)}(\Omega_{M+1}+M))%
}

新的ψ作用于背景本身以后，后续加项仍各有自己的下标与作用对象，不是把外面的ψ再抄一遍便能代替。

\NovelDisplay{M04916R01}{%
(0)(1,1,1)(2,1,1)(3,1,1)(3,1)(4,2)(2,1)(3,2,1)(4,2,1)(5,2,1)(5,2)(6,3)%
}
\NovelDisplay{M04916R02}{%
(1,1,1)(2,1,1)(3,1,1)(3,1)(4,2)(2,1)(3,2,1)(4,2,1)(5,2,1)(5,2)(6,3)%
}

\NovelDisplay{M04917R01}{%
= \psi(\Omega_{M+1}+\psi_{\Omega_{M+1}}(\Omega_{M+1})+\psi_{\psi_{M}(\Omega_{M+1}+M)}(\Omega_{M+1}+\psi_{\Omega_{M+1}}(\Omega_{M+1})))%
}

\NovelDisplay{M04918R01}{%
(0)(1,1,1)(2,1,1)(3,1,1)(3,1)(4,2)(2,1)(3,2,1)(4,2,1)(5,2,1)(5,2)%
}
\NovelDisplay{M04918R02}{%
(6,3)(2)%
}

\NovelDisplay{M04919R01}{%
= \psi(\Omega_{M+1}+\psi_{\Omega_{M+1}}(\Omega_{M+1})+\psi_{\psi_{M}(\Omega_{M+1}+M)}(\Omega_{M+1}+\psi_{\Omega_{M+1}}(\Omega_{M+1})+1))%
}

\NovelDisplay{M04920R01}{%
(0)(1,1,1)(2,1,1)(3,1,1)(3,1)(4,2)(2,1)(3,2,1)(4,2,1)(5,2,1)(5,2)%
}
\NovelDisplay{M04920R02}{%
(6,3)(3,2)%
}

\NovelDisplay{M04921R01}{%
= \psi(\Omega_{M+1}+\psi_{\Omega_{M+1}}(\Omega_{M+1})+\psi_{M}(\Omega_{M+1}+M))%
}

\NovelDisplay{M04922R01}{%
(0)(1,1,1)(2,1,1)(3,1,1)(3,1)(4,2)(2,1)(3,2,1)(4,2,1)(5,2,1)(5,2)(6,3)%
}
\NovelDisplay{M04922R02}{%
(3,2,1)(4,2,1)(5,2,1)(5,2)(6,3)%
}

\NovelDisplay{M04923R01}{%
= \psi(\Omega_{M+1}+\psi_{\Omega_{M+1}}(\Omega_{M+1})+\psi_{M}(\Omega_{M+1}+\psi_{\Omega_{M+1}}(\Omega_{M+1})))%
}

小舞先读完整个外框，再由唐三指出内层哪一处换了作用对象，这几式才没有在口述中丢掉中间那层ψ。

卢本韦将右式盖住，叶阳只看矩阵重做三项，等他写完，古月娜再同原来的记录逐列相核。

\NovelDisplay{M04924R01}{%
\text{展开第零项：}\quad(0)(1,1,1)(2,1,1)(3,1,1)(3,1)(4,2)(2,1)(3,2,1)(4,2,1)%
}
\NovelDisplay{M04924R02}{%
(5,2,1)(5,2)(6,3)(1,1,1)(2,1,1)(3,1,1)(3,1)(4,2)(2,1)(3,2,1)%
}
\NovelDisplay{M04924R03}{%
(4,2,1)(5,2,1)(5,2)%
}

\NovelDisplay{M04925R01}{%
\text{展开第一项：}\quad(0)(1,1,1)(2,1,1)(3,1,1)(3,1)(4,2)(2,1)(3,2,1)(4,2,1)%
}
\NovelDisplay{M04925R02}{%
(5,2,1)(5,2)(6,3)(1,1,1)(2,1,1)(3,1,1)(3,1)(4,2)(2,1)(3,2,1)(4,2,1)%
}
\NovelDisplay{M04925R03}{%
(5,2,1)(5,2)(6,2)%
}

\NovelDisplay{M04926R01}{%
\text{展开第二项：}\quad(0)(1,1,1)(2,1,1)(3,1,1)(3,1)(4,2)(2,1)(3,2,1)(4,2,1)%
}
\NovelDisplay{M04926R02}{%
(5,2,1)(5,2)(6,3)(1,1,1)(2,1,1)(3,1,1)(3,1)(4,2)(2,1)(3,2,1)(4,2,1)%
}
\NovelDisplay{M04926R03}{%
(5,2,1)(5,2)(6,2)(7,2)%
}

嵌套式再长，矩阵的三次取项仍各有同一原式作为起点，不能靠省去最内层来得到更短的答案。

叶阳第一次没有问谁算得快，而是指着自己留下的保留段，请蔡须坤确认他没有漏列。

\NovelDisplay{M04927R01}{%
(0)(1,1,1)(2,1,1)(3,1,1)(3,1)(4,2)(2,1)(3,2,1)(4,2,1)(5,2,1)(5,2)%
}
\NovelDisplay{M04927R02}{%
(6,3)(4,2)(5,3)%
}

\NovelDisplay{M04928R01}{%
= \psi(\Omega_{M+1}+\psi_{\Omega_{M+1}}(\Omega_{M+1})+M)%
}

\NovelDisplay{M04929R01}{%
(0)(1,1,1)(2,1,1)(3,1,1)(3,1)(4,2)(2,1)(3,2,1)(4,2,1)(5,2,1)(5,2)(6,3)%
}
\NovelDisplay{M04929R02}{%
(4,2)(5,3,1)(6,3,1)(7,3)(7)%
}

\NovelDisplay{M04930R01}{%
= \psi(\Omega_{M+1}+\psi_{\Omega_{M+1}}(\Omega_{M+1})+M^M)%
}

\NovelDisplay{M04931R01}{%
(0)(1,1,1)(2,1,1)(3,1,1)(3,1)(4,2)(2,1)(3,2,1)(4,2,1)(5,2,1)(5,2)(6,3)%
}
\NovelDisplay{M04931R02}{%
(4,2)(5,3,1)(6,3,1)(7,3)(8,4)%
}

\NovelDisplay{M04932R01}{%
= \psi(\Omega_{M+1}+\psi_{\Omega_{M+1}}(\Omega_{M+1})\times 2)%
}

\NovelDisplay{M04933R01}{%
(0)(1,1,1)(2,1,1)(3,1,1)(3,1)(4,2)(2,1,1)%
}

\NovelDisplay{M04934R01}{%
= \psi(\Omega_{M+1}+\psi_{\Omega_{M+1}}(\Omega_{M+1}+1))%
}

外层的加法、乘二与内层参数的后继在此相邻出现，牢牢记住它们所在的括号，才能避免把不同改动混成一条规则。

有一天，山路上的旧水沟被雨冲开。

叶阳先看见石面上那道细细的裂口，让大家把包交过去，一个个从旁边走。小舞原想跳，抬腿时却发现脚踝在压力下发僵，只得重新扶住岩壁。

诺当走在最后，背着装纸的布袋。袋角被树根挂住，他下意识回身去扯，脚底踩空，半边身子一下落进沟里。

霍雨浩扑过去攥住他的手，薛峰趴下压住自己的重心。纸已经从袋口散出，奥斯卡只看了一眼，便将包丢下，双手抱住薛峰的腰。

古月娜与卢本韦从另一侧将人拉回，蔡须坤跪到地上，先去看诺当的脚。

他疼得说不出话，指着沟下漂走的纸，嘴唇动了几下。

``先看人。''蔡须坤说。

那句话说出口以后，她才发觉自己的手也在抖。

\NovelDisplay{M04935R01}{%
(0)(1,1,1)(2,1,1)(3,1,1)(3,1)(4,2)(2,1,1)(3,1)(2)%
}

\NovelDisplay{M04936R01}{%
= \psi(\Omega_{M+1}+\psi_{\Omega_{M+1}}(\Omega_{M+1}+M))%
}

\NovelDisplay{M04937R01}{%
(0)(1,1,1)(2,1,1)(3,1,1)(3,1)(4,2)(2,1,1)(3,1)(2)(1,1,1)(2,1,1)(3,1,1)%
}
\NovelDisplay{M04937R02}{%
(3,1)(4,2)(2,1,1)(3,1)%
}

\NovelDisplay{M04938R01}{%
= \psi(\Omega_{M+1}+\psi_{\Omega_{M+1}}(\Omega_{M+1}+M)+\psi_{M}(\Omega_{M+1}+\psi_{\Omega_{M+1}}(\Omega_{M+1}+\psi_{M}(M))))%
}

\NovelDisplay{M04939R01}{%
(0)(1,1,1)(2,1,1)(3,1,1)(3,1)(4,2)(2,1,1)(3,1)(2)(1,1,1)(2,1,1)(3,1,1)%
}
\NovelDisplay{M04939R02}{%
(3,1)(4,2)(2,1,1)(3,1)(2)%
}

\NovelDisplay{M04940R01}{%
= \psi(\Omega_{M+1}+\psi_{\Omega_{M+1}}(\Omega_{M+1}+M)+\psi_{\psi_{M}(\Omega_{M+1}+\psi_{\Omega_{M+1}}(\Omega_{M+1}+M))}(\Omega_{M+1}%
}
\NovelDisplay{M04940R02}{%
+\psi_{\Omega_{M+1}}(\Omega_{M+1}+M)))%
}

\NovelDisplay{M04941R01}{%
(0)(1,1,1)(2,1,1)(3,1,1)(3,1)(4,2)(2,1,1)(3,1)(2,1)(3,2,1)(4,2,1)(5,2,1)%
}
\NovelDisplay{M04941R02}{%
(5,2)(6,3)(4,2,1)(5,2)(4)%
}

\NovelDisplay{M04942R01}{%
= \psi(\Omega_{M+1}+\psi_{\Omega_{M+1}}(\Omega_{M+1}+M)+\psi_{M}(\Omega_{M+1}+\psi_{\Omega_{M+1}}(\Omega_{M+1}+M)))%
}

卢本韦把能用I册旧页对照的前段留在一旁，面对新的嵌套则照实际记号读，不替讲义编造一条更短的I等式。

回到雨棚以后，他们把刚才最容易混读的一条留下再做，霍雨浩要求三个人各自落笔，中途不交换纸页。

这一条共十七列，末列只留下第一行的二。沿第一行回找，坏根在第十列；于是末列之前的十六列保留，第十至第十六列这七列作为一段重新生成。

古月娜问这里还要不要照上一回给第一行加三。小舞已经伸手想写，听见以后才停住。最高非零位置已经落到第一行，它上面没有更低编号的行需要依那种差值提升，不能把前一个例子里的三当成每次展开都要加的常数。

三个人分别写完，才将各自的纸推到桌中间。

\NovelDisplay{M04943R01}{%
\text{展开第零项：}\quad(0)(1,1,1)(2,1,1)(3,1,1)(3,1)(4,2)(2,1,1)(3,1)(2)(1,1,1)%
}
\NovelDisplay{M04943R02}{%
(2,1,1)(3,1,1)(3,1)(4,2)(2,1,1)(3,1)%
}

\NovelDisplay{M04944R01}{%
\text{展开第一项：}\quad(0)(1,1,1)(2,1,1)(3,1,1)(3,1)(4,2)(2,1,1)(3,1)(2)(1,1,1)%
}
\NovelDisplay{M04944R02}{%
(2,1,1)(3,1,1)(3,1)(4,2)(2,1,1)(3,1)(1,1,1)(2,1,1)(3,1,1)(3,1)%
}
\NovelDisplay{M04944R03}{%
(4,2)(2,1,1)(3,1)%
}

\NovelDisplay{M04945R01}{%
\text{展开第二项：}\quad(0)(1,1,1)(2,1,1)(3,1,1)(3,1)(4,2)(2,1,1)(3,1)(2)(1,1,1)%
}
\NovelDisplay{M04945R02}{%
(2,1,1)(3,1,1)(3,1)(4,2)(2,1,1)(3,1)(1,1,1)(2,1,1)(3,1,1)(3,1)(4,2)%
}
\NovelDisplay{M04945R03}{%
(2,1,1)(3,1)(1,1,1)(2,1,1)(3,1,1)(3,1)(4,2)(2,1,1)(3,1)%
}

这里随取项编号增加的是按原规则生成的整段，三次展开的差别必须由矩阵自身说明。

诺当靠着床坐起来，将自己刚能说清的那一处也指给他们看，声音虽轻，却没有再含糊过去。

这一回三项分别是十六列、二十三列、三十列。蔡须坤将两个例子并排保留，一个提醒她看提升差，另一个提醒她，先前用过的提升不能凭熟悉直接搬来。

\NovelDisplay{M04946R01}{%
(0)(1,1,1)(2,1,1)(3,1,1)(3,1)(4,2)(2,1,1)(3,1)(4,2)%
}

\NovelDisplay{M04947R01}{%
= \psi(\Omega_{M+1}+\psi_{\Omega_{M+1}}(\Omega_{M+1}+M)+M)%
}

\NovelDisplay{M04948R01}{%
(0)(1,1,1)(2,1,1)(3,1,1)(3,1)(4,2)(2,1,1)(3,1,1)%
}

\NovelDisplay{M04949R01}{%
= \psi(\Omega_{M+1}+\psi_{\Omega_{M+1}}(\Omega_{M+1}+M+1))%
}

\NovelDisplay{M04950R01}{%
(0)(1,1,1)(2,1,1)(3,1,1)(3,1)(4,2)(2,1,1)(3,1,1)(2,1,1)(3,1)(2)%
}

\NovelDisplay{M04951R01}{%
= \psi(\Omega_{M+1}+\psi_{\Omega_{M+1}}(\Omega_{M+1}+M\times 2))%
}

\NovelDisplay{M04952R01}{%
(0)(1,1,1)(2,1,1)(3,1,1)(3,1)(4,2)(2,1,1)(3,1,1)(2,1,1)(3,1,1)%
}

\NovelDisplay{M04953R01}{%
= \psi(\Omega_{M+1}+\psi_{\Omega_{M+1}}(\Omega_{M+1}+M\times 2+1))%
}

这里连续比较内层参数的后继与倍数，诺当先说清变化在ψ里面，才独立接完自己的那一段。

诺当在山脚休息了几天，脚踝肿着，十一人的进度也就停在那里。

夏夜星空没有要求所有人坐在床边守他。能练的部分仍可以练，可谁也不能把他的那一份念过，就在山图上替他多画一道线。霍雨浩把次日的安排改了，前面几人轮流上山复核，傍晚再将当天各自的问题带回来。

丢掉的纸要重写。诺当最初只说对不起，听到第三回，小舞索性搬张凳子坐到床边，让他从第一列开始说。

``你说，唐三写，我看。''

诺当说得很慢。有些地方，他原先只是记得那页在袋里的位置，纸没了，才发现自己并没有真正记住。唐三没有替他补，薛峰在另一边留出空行，等他自己将问题指出来。

几天以后，重写的那叠比丢掉的薄。诺当抱着它，却头一回不怕被谁忽然问起中间的一段。

\NovelDisplay{M04954R01}{%
(0)(1,1,1)(2,1,1)(3,1,1)(3,1)(4,2)(2,1,1)(3,1,1)(3)%
}

\NovelDisplay{M04955R01}{%
= \psi(\Omega_{M+1}+\psi_{\Omega_{M+1}}(\Omega_{M+1}+M\times \omega))%
}

\NovelDisplay{M04956R01}{%
(0)(1,1,1)(2,1,1)(3,1,1)(3,1)(4,2)(2,1,1)(3,1,1)(3,1)(2)%
}

\NovelDisplay{M04957R01}{%
= \psi(\Omega_{M+1}+\psi_{\Omega_{M+1}}(\Omega_{M+1}+M^2))%
}

\NovelDisplay{M04958R01}{%
(0)(1,1,1)(2,1,1)(3,1,1)(3,1)(4,2)(2,1,1)(3,1,1)(3,1)(2,1,1)(3,1,1)%
}
\NovelDisplay{M04958R02}{%
(3,1)(2)%
}

\NovelDisplay{M04959R01}{%
= \psi(\Omega_{M+1}+\psi_{\Omega_{M+1}}(\Omega_{M+1}+M^2\times 2))%
}

\NovelDisplay{M04960R01}{%
(0)(1,1,1)(2,1,1)(3,1,1)(3,1)(4,2)(2,1,1)(3,1,1)(3,1)(3,1)(2)%
}

\NovelDisplay{M04961R01}{%
= \psi(\Omega_{M+1}+\psi_{\Omega_{M+1}}(\Omega_{M+1}+M^3))%
}

从ω倍数转到有限幂时，整个内层背景没有消失，蔡兴坤逐一找回共同前段，才把不同尾列标出来。

诺当重新上山前，蔡须坤将几条较短的练习挑给他，想着先让他找回站稳的感觉。

他看完，却把其中一条推了回来，指出她纸角漏记的那一处。蔡须坤本来还要解释自己留的是简记，回头找到原式，才发觉那一列确实不该省。

她将纸收近，认真补好。

``谢谢。''

诺当没想到她会这样说，手在布带上搓了一下，笑得有些不好意思。奥斯卡站在门边，看两人都低着头，便把原来要说的鼓励咽了回去。

那天下午，诺当又在较低处停了很久，蔡须坤照旧坐到旁边听。他并没有忽然赶上前面几人，却已不再觉得自己手里只有需要别人补好的空页。

晚上收纸时，诺当又请她先别翻下一页，自己将今日会的与还不会的分开说了一遍。蔡须坤听过，才发觉他问问题时终于没有先说对不起。她把那份纸留好，回到屋里，将大家第二天要核的原式逐一取出。卢本韦看着她，问是不是还要练很久。她摇头，笑了一下，把手里的笔放下，说今天能做的已经做过，剩下的，明天再一同来看。

\NovelDisplay{M04962R01}{%
(0)(1,1,1)(2,1,1)(3,1,1)(3,1)(4,2)(2,1,1)(3,1,1)(3,1)(4,1)(2)%
}

\NovelDisplay{M04963R01}{%
= \psi(\Omega_{M+1}+\psi_{\Omega_{M+1}}(\Omega_{M+1}+M^M))%
}

\NovelDisplay{M04964R01}{%
(0)(1,1,1)(2,1,1)(3,1,1)(3,1)(4,2)(2,1,1)(3,1,1)(3,1)(4,2)%
}

\NovelDisplay{M04965R01}{%
= \psi(\Omega_{M+1}+\psi_{\Omega_{M+1}}(\Omega_{M+1}+\psi_{\Omega_{M+1}}(\Omega_{M+1})))%
}

\NovelDisplay{M04966R01}{%
(0)(1,1,1)(2,1,1)(3,1,1)(3,1)(4,2)(3)%
}

\NovelDisplay{M04967R01}{%
= \psi(\Omega_{M+1}+\psi_{\Omega_{M+1}}(\Omega_{M+1}+\psi_{\Omega_{M+1}}(\Omega_{M+1}+1)))%
}

\NovelDisplay{M04968R01}{%
(0)(1,1,1)(2,1,1)(3,1,1)(3,1)(4,2)(3,1)(4,2)%
}

\NovelDisplay{M04969R01}{%
= \psi(\Omega_{M+1}+\psi_{\Omega_{M+1}}(\Omega_{M+1}+\psi_{\Omega_{M+1}}(\Omega_{M+1}+\psi_{\Omega_{M+1}}(\Omega_{M+1}))))%
}

一层ψ重新作用于先前的ψ项，读式顺序必须向内追清，不能把反复出现的同一背景误当作重复书写。

雨渐渐少了，奥斯卡将那只漏过水的布袋摊在门槛上。

蔡兴坤翻出针线，说自己能补。黑子鸡在旁边啄了两下，被他护开；叶阳按住袋角，诺当坐在另一侧，把新写的纸分成几份，不再全塞进同一个口袋。

蔡须坤没有过去替他们安排。她在屋里听唐三与小舞互问，听见小舞第三次纠正下标，便在自己的纸上也画了同样一道框。

卢本韦看见，问她是不是还在惦记今天没走完。

``惦记。''

蔡须坤抬头，望了一眼门口那几个认真补袋子的人。

``可明天还能走。''

\NovelDisplay{M04970R01}{%
(0)(1,1,1)(2,1,1)(3,1,1)(3,1)(4,2)(4)%
}

\NovelDisplay{M04971R01}{%
= \psi(\Omega_{M+1}+\psi_{\Omega_{M+1}}(\Omega_{M+1}+\psi_{\Omega_{M+1}}(\Omega_{M+1}+\psi_{\Omega_{M+1}}(\Omega_{M+1}+1))))%
}

\NovelDisplay{M04972R01}{%
(0)(1,1,1)(2,1,1)(3,1,1)(3,1)(4,2)(4,1)(5,2)%
}

\NovelDisplay{M04973R01}{%
= \psi(\Omega_{M+1}+\psi_{\Omega_{M+1}}(\Omega_{M+1}+\psi_{\Omega_{M+1}}(\Omega_{M+1}+\psi_{\Omega_{M+1}}(\Omega_{M+1}+\psi_{\Omega_{M+1}}(\Omega_{M+1})))))%
}

\NovelDisplay{M04974R01}{%
(0)(1,1,1)(2,1,1)(3,1,1)(3,1)(4,2)(4,2)%
}

\NovelDisplay{M04975R01}{%
= \psi(\Omega_{M+1}\times 2)%
}

\NovelDisplay{M04976R01}{%
(0)(1,1,1)(2,1,1)(3,1,1)(3,1)(4,2)(4,2)(2,1,1)%
}

\NovelDisplay{M04977R01}{%
= \psi(\Omega_{M+1}\times 2+\psi_{\Omega_{M+1}}(\Omega_{M+1}\times 2 + 1))%
}

嵌套转入背景的倍数以后，叶阳将括号与乘号重新分组，确认原式没有靠省略一段才变得简短。

\NovelDisplay{M04978R01}{%
(0)(1,1,1)(2,1,1)(3,1,1)(3,1)(4,2)(4,2)(4,2)%
}

\NovelDisplay{M04979R01}{%
= \psi(\Omega_{M+1}\times 3)%
}

\NovelDisplay{M04980R01}{%
(0)(1,1,1)(2,1,1)(3,1,1)(3,1)(4,2)(5,3)%
}

\NovelDisplay{M04981R01}{%
= \psi(\Omega_{M+2})%
}

\NovelDisplay{M04982R01}{%
(0)(1,1,1)(2,1,1)(3,1,1)(3,1)(4,2,1)%
}

\NovelDisplay{M04983R01}{%
= \psi(\Omega_{M+\omega})%
}

\NovelDisplay{M04984R01}{%
(0)(1,1,1)(2,1,1)(3,1,1)(3,1)(4,2,1)(5,2)(6,3)%
}

\NovelDisplay{M04985R01}{%
= \psi(\Omega_{M+\omega+1})%
}

末尾开始改变Ω的下标，小舞把下标框单独指给唐三看，不再将\allowbreak{}\NovelInline{I06099}{M+\omega}\allowbreak{}念成Ω外面的加项。

第五个月将尽时，十一人终于站上了四百米的横台。

诺当最后上来，脚上还缠着已经换薄的布带。他扶住石桩，没有马上坐，先把自己最后那次取项说完。小舞站在旁边，等他停下，才将水递过去。

远处已经能看见雨后的海。云从山腰绕过来，他们来时那段湿滑的沟，藏在更低处的雾里，一点也看不清了。

蔡须坤将日期写在山图上，笔尖在那个月字旁停了停。

她清楚地记得第一次踩上三百八十米的感觉，也记得诺当被拉回来的那一瞬。若只看图上多出的那几道横线，很难看见他们究竟用了多少个回头。

霍雨浩把原来分开的纸放回袋里，薛峰检查封口。卢本韦坐在蔡须坤身旁，将今天能独立说清的那几页折好，问她还要不要看上面。

蔡须坤抬头。更高处仍有石阶，也仍有他们没走懂的地方。

``看。''她说，``先歇一会儿再看。''

卢本韦把水壶递过来。蔡须坤接住，等最后一个人坐到横台上，才将本子合好。

\chapter{第一百二十七章：没有人替你走过的那一步}\label{ux7b2cux4e00ux767eux4e8cux5341ux4e03ux7ae0ux6ca1ux6709ux4ebaux66ffux4f60ux8d70ux8fc7ux7684ux90a3ux4e00ux6b65}

第六个月开始，山上连着晴了几天。

诺当已经能自己走到练习的位置，奥斯卡却仍习惯回头看一眼他的脚。看得多了，诺当忍不住问，自己是不是走得很怪。奥斯卡说是怕他又把本子全塞在同一个袋里，诺当便当着他的面，把两份纸分进不同的口袋。

小舞笑过，先向上试了一步。压力顺着腿骨压来，笑便立刻收了。她扶住岩壁，等呼吸平下来，再将刚才漏掉的下标重新指出。

这一阶段，最先走清的仍常是蔡须坤、霍雨浩、卢本韦与薛峰。几人却渐渐不再将早上走清多少，直接当成全队当天已经完成多少。唐三和小舞有自己的问题，叶阳、蔡兴坤、诺当与奥斯卡，也各自留着不同的空处。

卢本韦把I册放在膝上。翻到新的对照时，她看见一处可以接回自己熟悉的记法，指尖停了很久。

``这次从这里说。''她对蔡须坤道。

蔡须坤没有替她将M侧先念完，只将自己的纸转过来，等她把问题完整说出。

开始几日，卢本韦比谁都认真。

她将旧I页与新讲义放在一起，每次读到相同的字母便回头查。查得太久，等她抬头，蔡须坤已经将旁边那一段独立接了两遍，薛峰也有了新的问法。

第三次停住时，她忽然把本子合上。

``是不是只有我还在翻旧东西？''

蔡须坤望着她，一时没答。霍雨浩过来，先问她刚才翻回去，究竟找到了哪一处。卢本韦说完，他又问，若当时没有翻，那一处准备怎么解释。

她没有开口。

蔡须坤这才将自己的M册也合上，露出下面留着问号的一页。

``这里我也还在查。''她说。

\NovelDisplay{M04986R01}{%
(0)(1,1,1)(2,1,1)(3,1,1)(3,1)(4,2,1)(5,2,1)%
}

\NovelDisplay{M04987R01}{%
= \psi(\Omega_{M+\omega^2})%
}

\NovelDisplay{M04988R01}{%
(0)(1,1,1)(2,1,1)(3,1,1)(3,1)(4,2,1)(5,2,1)(6,1)(2)%
}

\NovelDisplay{M04989R01}{%
= \psi(\Omega_{M\times 2})%
}

\NovelDisplay{M04990R01}{%
(0)(1,1,1)(2,1,1)(3,1,1)(3,1)(4,2,1)(5,2,1)(6,2)%
}

\NovelDisplay{M04991R01}{%
= \psi(\Omega_{\Omega_{M+1}})%
}

\NovelDisplay{M04992R01}{%
(0)(1,1,1)(2,1,1)(3,1,1)(3,1)(4,2,1)(5,2,1)(6,2)(5)%
}

\NovelDisplay{M04993R01}{%
= \psi(I_{M+1})%
}

\NovelDisplay{M04994R01}{%
= \psi({M_2}^2)%
}

从Ω下标的幂与嵌套接到\allowbreak{}\NovelInline{I06100}{I_{M+1}}\allowbreak{}时，卢本韦终于找到I册可以直接接入的一页，同时保留了M侧给出的完整对照。

\NovelDisplay{M04995R01}{%
(0)(1,1,1)(2,1,1)(3,1,1)(3,1)(4,2,1)(5,2,1)(6,2)(5)(2,1,1)%
}

\NovelDisplay{M04996R01}{%
= \psi({M_2}^2+\psi_{\Omega_{M+1}}({M_2}^2+1))%
}

\NovelDisplay{M04997R01}{%
(0)(1,1,1)(2,1,1)(3,1,1)(3,1)(4,2,1)(5,2,1)(6,2)(5)(4,2)%
}

\NovelDisplay{M04998R01}{%
= \psi({M_2}^2+\Omega_{M+1})%
}

\NovelDisplay{M04999R01}{%
(0)(1,1,1)(2,1,1)(3,1,1)(3,1)(4,2,1)(5,2,1)(6,2)(5)(4,2,1)(5,2,1)%
}
\NovelDisplay{M04999R02}{%
(6,2)(5)%
}

\NovelDisplay{M05000R01}{%
= \psi({M_2}^2+\psi_{\psi_{M_2}({M_2}^2)}({M_2}^2))%
}

\NovelDisplay{M05001R01}{%
(0)(1,1,1)(2,1,1)(3,1,1)(3,1)(4,2,1)(5,2,1)(6,2)(5,2)(6,3,1)(7,3,1)%
}
\NovelDisplay{M05001R02}{%
(8,3)(7)%
}

\NovelDisplay{M05002R01}{%
= \psi({M_2}^2+\psi_{M_2}({M_2}^2))%
}

薛峰先固定\allowbreak{}\NovelInline{I06101}{M_2^2}\allowbreak{}的背景，再区分直接加项与带ψ的回接项，读到相同字母也没有省去下标。

\NovelDisplay{M05003R01}{%
(0)(1,1,1)(2,1,1)(3,1,1)(3,1)(4,2,1)(5,2,1)(6,2)(7,3)%
}

\NovelDisplay{M05004R01}{%
= \psi({M_2}^2+M_2)%
}

\NovelDisplay{M05005R01}{%
(0)(1,1,1)(2,1,1)(3,1,1)(3,1)(4,2,1)(5,2,1)(6,2)(7,3,1)(8,3,1)(9,3)(8)%
}

\NovelDisplay{M05006R01}{%
= \psi({M_2}^2 \times 2)%
}

\NovelDisplay{M05007R01}{%
(0)(1,1,1)(2,1,1)(3,1,1)(3,1)(4,2,1)(5,2,1)(6,2,1)%
}

\NovelDisplay{M05008R01}{%
= \psi({M_2}^2 \times \omega)%
}

\NovelDisplay{M05009R01}{%
(0)(1,1,1)(2,1,1)(3,1,1)(3,1)(4,2,1)(5,2,1)(6,2,1)(5,2,1)(6,2)(5)%
}

\NovelDisplay{M05010R01}{%
= \psi({M_2}^3)%
}

平方外的加法、倍数与立方依次出现，蔡须坤让每个人同时指出式子中的位置和矩阵中的那一段。

\NovelDisplay{M05011R01}{%
(0)(1,1,1)(2,1,1)(3,1,1)(3,1)(4,2,1)(5,2,1)(6,2,1)(6)%
}

\NovelDisplay{M05012R01}{%
= \psi({M_2}^\omega)%
}

\NovelDisplay{M05013R01}{%
(0)(1,1,1)(2,1,1)(3,1,1)(3,1)(4,2,1)(5,2,1)(6,2,1)(6,1)(2)%
}

\NovelDisplay{M05014R01}{%
= \psi({M_2}^{M})%
}

\NovelDisplay{M05015R01}{%
(0)(1,1,1)(2,1,1)(3,1,1)(3,1)(4,2,1)(5,2,1)(6,2,1)(6,2)(5)%
}

\NovelDisplay{M05016R01}{%
= \psi({M_2}^{M_2})%
}

\NovelDisplay{M05017R01}{%
(0)(1,1,1)(2,1,1)(3,1,1)(3,1)(4,2,1)(5,2,1)(6,2,1)(6,2)(5,2,1)(6,2)(7,3)%
}

\NovelDisplay{M05018R01}{%
= \psi({M_2}^{M_2}+M_2)%
}

底数仍是\allowbreak{}\NovelInline{I06102}{M_2}\allowbreak{}而指数逐步变化，不能把指数里的\allowbreak{}\NovelInline{I06103}{M}\allowbreak{}和底数的下标二混读成同一个符号。

太阳照在石面上，到正午便热得不能久坐。

古月娜在岩影下留好水，小舞喝过一口，又要起身。她今天已比唐三先读懂一处，便想再多走几级，等哥哥赶上时，自己也许已经不用停。

走到第二级，她才发觉眼前的字在晃。唐三叫了她一声，她没应，抬脚时鞋尖撞在石沿，整个人被逼得跪了下去。

霍雨浩与叶阳将她扶回阴影。她坐稳了，还想拿笔，古月娜伸手按住纸，让她先将刚才那句话说全。小舞只说出开头，便忘了自己原来要往哪一层接。

她低头坐了一会儿，将水喝完。

下午再练，她先问唐三，自己方才没说完的究竟是哪一处。唐三没有笑她，把留着空行的本子推过去，等她自己补上。

\NovelDisplay{M05019R01}{%
(0)(1,1,1)(2,1,1)(3,1,1)(3,1)(4,2,1)(5,2,1)(6,2,1)(6,2)(5,2,1)(6,2,1)%
}

\NovelDisplay{M05020R01}{%
= \psi({M_2}^{M_2}\times \omega)%
}

\NovelDisplay{M05021R01}{%
(0)(1,1,1)(2,1,1)(3,1,1)(3,1)(4,2,1)(5,2,1)(6,2,1)(6,2)(5,2,1)(6,2,1)%
}
\NovelDisplay{M05021R02}{%
(5,2,1)(6,2)(5)%
}

\NovelDisplay{M05022R01}{%
= \psi({M_2}^{M_2+1})%
}

\NovelDisplay{M05023R01}{%
(0)(1,1,1)(2,1,1)(3,1,1)(3,1)(4,2,1)(5,2,1)(6,2,1)(6,2)(5,2,1)%
}
\NovelDisplay{M05023R02}{%
(6,2,1)(6,2)(5)%
}

\NovelDisplay{M05024R01}{%
= \psi({M_2}^{M_2\times 2})%
}

\NovelDisplay{M05025R01}{%
(0)(1,1,1)(2,1,1)(3,1,1)(3,1)(4,2,1)(5,2,1)(6,2,1)(6,2)(6,2)(5)%
}

\NovelDisplay{M05026R01}{%
= \psi({M_2}^{{M_2}^2})%
}

这一组把幂外倍数、指数后继及指数内平方并列，古月娜先讲清三个位置，才让小舞重新复算。

\NovelDisplay{M05027R01}{%
(0)(1,1,1)(2,1,1)(3,1,1)(3,1)(4,2,1)(5,2,1)(6,2,1)(6,2)(7,3)%
}

\NovelDisplay{M05028R01}{%
= \psi(\Omega_{M_2+1})%
}

\NovelDisplay{M05029R01}{%
= \psi(M_3)%
}

\NovelDisplay{M05030R01}{%
(0)(1,1,1)(2,1,1)(3,1,1)(3,1)(4,2,1)(5,2,1)(6,2,1)(6,2)(7,3,1)%
}
\NovelDisplay{M05030R02}{%
(8,3,1)(9,3)(8)%
}

\NovelDisplay{M05031R01}{%
= \psi({M_3}^2)%
}

\NovelDisplay{M05032R01}{%
(0)(1,1,1)(2,1,1)(3,1,1)(3,1)(4,2,1)(5,2,1)(6,2,1)(6,2)(7,3,1)(8,3,1)%
}
\NovelDisplay{M05032R02}{%
(9,3,1)(9,3)(8)%
}

\NovelDisplay{M05033R01}{%
= \psi({M_3}^{M_3})%
}

\NovelDisplay{M05034R01}{%
(0)(1,1,1)(2,1,1)(3,1,1)(3,1)(4,2,1)(5,2,1)(6,2,1)(6,2)(7,3,1)(8,3,1)%
}
\NovelDisplay{M05034R02}{%
(9,3,1)(9,3)(10,4)%
}

\NovelDisplay{M05035R01}{%
= \psi(M_4)%
}

从\allowbreak{}\NovelInline{I06104}{M_3}\allowbreak{}到\allowbreak{}\NovelInline{I06105}{M_4}\allowbreak{}的写法沿用讲义，队员们核对的是对应的完整矩阵，没有把下标增加直接当作普通加一。

\NovelDisplay{M05036R01}{%
(0)(1,1,1)(2,1,1)(3,1,1)(3,1,1)%
}

\NovelDisplay{M05037R01}{%
= \psi(M_\omega)%
}

\NovelDisplay{M05038R01}{%
= \psi(N\times \omega)%
}

\NovelDisplay{M05039R01}{%
(0)(1,1,1)(2,1,1)(3,1,1)(3,1,1)(1,1,1)(2,1,1)(3,1,1)(3,1)(4,2)%
}

\NovelDisplay{M05040R01}{%
= \psi(M_\omega + \psi_{M}(M_2))%
}

\NovelDisplay{M05041R01}{%
(0)(1,1,1)(2,1,1)(3,1,1)(3,1,1)(1,1,1)(2,1,1)(3,1,1)(3,1)(4,2,1)(5,2,1)%
}
\NovelDisplay{M05041R02}{%
(6,2,1)(6,2)(7,3)%
}

\NovelDisplay{M05042R01}{%
= \psi(M_\omega + \psi_{M}(M_3))%
}

\NovelDisplay{M05043R01}{%
(0)(1,1,1)(2,1,1)(3,1,1)(3,1,1)(1,1,1)(2,1,1)(3,1,1)(3,1)(4,2,1)(5,2,1)%
}
\NovelDisplay{M05043R02}{%
(6,2,1)(6,2,1)%
}

\NovelDisplay{M05044R01}{%
= \psi(M_\omega + \psi_{M}(M_\omega))%
}

\allowbreak{}\NovelInline{I06106}{M_\omega}\allowbreak{}与\allowbreak{}\NovelInline{I06107}{N\times\omega}\allowbreak{}在这里给出对照，进入新的表示以后，后面接来的ψ仍要分别读清其下标。

薛峰有一晚主动留下来核纸。

他把几份记录按不同前段分好，想在第二天之前，把大家停过的地方都找回来。最初做得很顺，写到后来，他却开始将相似的共同背景误放到同一页，越核越觉得纸上的路已经接通。

天亮时，蔡须坤来敲门，见他还坐着，先问他吃过没有。薛峰答非所问，指着其中一页，说自己大概将那道难处走完了。

霍雨浩听见，坐下来请他从原矩阵开始说。不到半页，他便停在了自己抄错的前缀上。

薛峰的脸一点点白下来。他原想让大家少等一天，结果连自己手上的这一份也得重看。

蔡须坤将那页轻轻抽出，没有撕，只在纸角记下``昨夜''。她让他先吃早饭，今天去山上时，仍由他来指出这一处。

``错在哪，你已经看见了。''她说，``别连人也藏起来。''

\NovelDisplay{M05045R01}{%
(0)(1,1,1)(2,1,1)(3,1,1)(3,1,1)(1,1,1)(2,1,1)(3,1,1)(3,1)(4,2,1)(5,2,1)%
}
\NovelDisplay{M05045R02}{%
(6,2,1)(6,2,1)(2,1)(3,2)%
}

\NovelDisplay{M05046R01}{%
= \psi(M_\omega + M)%
}

\NovelDisplay{M05047R01}{%
(0)(1,1,1)(2,1,1)(3,1,1)(3,1,1)(1,1,1)(2,1,1)(3,1,1)(3,1)(4,2,1)(5,2,1)%
}
\NovelDisplay{M05047R02}{%
(6,2,1)(6,2,1)(2,1)(3,2,1)(4,2,1)(5,2,1)(5,2)(6,3)%
}

\NovelDisplay{M05048R01}{%
= \psi(M_\omega + \psi_{\psi_{M_2}(M_2)}(M_2))%
}

\NovelDisplay{M05049R01}{%
(0)(1,1,1)(2,1,1)(3,1,1)(3,1,1)(1,1,1)(2,1,1)(3,1,1)(3,1)(4,2,1)(5,2,1)%
}
\NovelDisplay{M05049R02}{%
(6,2,1)(6,2,1)(2,1)(3,2,1)(4,2,1)(5,2,1)(5,2)(6,3,1)(7,3,1)(8,3,1)%
}
\NovelDisplay{M05049R03}{%
(8,3)(9,4)%
}

\NovelDisplay{M05050R01}{%
= \psi(M_\omega + \psi_{\psi_{M_2}(M_2)}(M_3))%
}

\NovelDisplay{M05051R01}{%
(0)(1,1,1)(2,1,1)(3,1,1)(3,1,1)(1,1,1)(2,1,1)(3,1,1)(3,1)(4,2,1)(5,2,1)%
}
\NovelDisplay{M05051R02}{%
(6,2,1)(6,2,1)(2,1)(3,2,1)(4,2,1)(5,2,1)(5,2,1)%
}

\NovelDisplay{M05052R01}{%
= \psi(M_\omega + \psi_{\psi_{M_2}(M_2)}(M_\omega))%
}

相同的\allowbreak{}\NovelInline{I06108}{M_\omega}\allowbreak{}背景下出现多种嵌套，卢本韦守住最外面的加法，蔡须坤再把内层作用对象逐个接回。

\NovelDisplay{M05053R01}{%
(0)(1,1,1)(2,1,1)(3,1,1)(3,1,1)(1,1,1)(2,1,1)(3,1,1)(3,1)(4,2,1)(5,2,1)%
}
\NovelDisplay{M05053R02}{%
(6,2,1)(6,2,1)(2,1)(3,2,1)(4,2,1)(5,2,1)(5,2,1)(1,1,1)(2,1,1)(3,1,1)%
}
\NovelDisplay{M05053R03}{%
(3,1)(4,2,1)(5,2,1)(6,2,1)(6,2,1)%
}

\NovelDisplay{M05054R01}{%
= \psi(M_\omega + \psi_{\psi_{M_2}(M_2)}(M_\omega)+\psi_{M}(M_\omega))%
}

\NovelDisplay{M05055R01}{%
(0)(1,1,1)(2,1,1)(3,1,1)(3,1,1)(1,1,1)(2,1,1)(3,1,1)(3,1)(4,2,1)(5,2,1)%
}
\NovelDisplay{M05055R02}{%
(6,2,1)(6,2,1)(2,1)(3,2,1)(4,2,1)(5,2,1)(5,2,1)(1,1,1)(2,1,1)(3,1,1)%
}
\NovelDisplay{M05055R03}{%
(3,1)(4,2,1)(5,2,1)(6,2,1)(6,2,1)(2,1)(3,2)%
}

\NovelDisplay{M05056R01}{%
= \psi(M_\omega + \psi_{\psi_{M_2}(M_2)}(M_\omega)+\psi_{\psi_{M}(M_\omega+M)}(M_\omega+M))%
}

\NovelDisplay{M05057R01}{%
(0)(1,1,1)(2,1,1)(3,1,1)(3,1,1)(1,1,1)(2,1,1)(3,1,1)(3,1)(4,2,1)(5,2,1)%
}
\NovelDisplay{M05057R02}{%
(6,2,1)(6,2,1)(2,1)(3,2,1)(4,2,1)(5,2,1)(5,2,1)(1,1,1)(2,1,1)(3,1,1)%
}
\NovelDisplay{M05057R03}{%
(3,1)(4,2,1)(5,2,1)(6,2,1)(6,2,1)(2,1)(3,2,1)(4,2,1)(5,2,1)(5,2,1)%
}

\NovelDisplay{M05058R01}{%
= \psi(M_\omega + \psi_{\psi_{M_2}(M_2)}(M_\omega)+\psi_{\psi_{M}(M_\omega+M)}(M_\omega+ \psi_{\psi_{M_2}(M_2)}(M_\omega)))%
}

\NovelDisplay{M05059R01}{%
(0)(1,1,1)(2,1,1)(3,1,1)(3,1,1)(1,1,1)(2,1,1)(3,1,1)(3,1)(4,2,1)(5,2,1)%
}
\NovelDisplay{M05059R02}{%
(6,2,1)(6,2,1)(2,1)(3,2,1)(4,2,1)(5,2,1)(5,2,1)(3,2)%
}

\NovelDisplay{M05060R01}{%
= \psi(M_\omega + \psi_{\psi_{M_2}(M_2)}(M_\omega)+\psi_{M}(M_\omega+M))%
}

矩阵变长时，保留段与新增段都要留在原位，霍雨浩请薛峰先指出真正改变的部分，再解释它与右式的关系。

薛峰将昨夜那条重新取出，先把右式盖好，再从原矩阵独立求第零项、第一项和第二项，三次都请叶阳核对共同保留的前段。

原矩阵虽有二十八列，末列的第二行父项却落在紧前的第二十七列。需要重生成的只有那一列，第一行的差为一。薛峰先报出位置，让叶阳回到原式确认，再开始向后写。

他们由此得到的三项，只依次增加这一列所生成的一份；大段相同的前缀仍完整保留，不能因为原式很长，就凭感觉圈出一大段当作坏部。

\NovelDisplay{M05061R01}{%
\text{展开第零项：}\quad(0)(1,1,1)(2,1,1)(3,1,1)(3,1,1)(1,1,1)(2,1,1)(3,1,1)%
}
\NovelDisplay{M05061R02}{%
(3,1)(4,2,1)(5,2,1)(6,2,1)(6,2,1)(2,1)(3,2,1)(4,2,1)(5,2,1)(5,2,1)%
}
\NovelDisplay{M05061R03}{%
(1,1,1)(2,1,1)(3,1,1)(3,1)(4,2,1)(5,2,1)(6,2,1)(6,2,1)(2,1)%
}

\NovelDisplay{M05062R01}{%
\text{展开第一项：}\quad(0)(1,1,1)(2,1,1)(3,1,1)(3,1,1)(1,1,1)(2,1,1)(3,1,1)%
}
\NovelDisplay{M05062R02}{%
(3,1)(4,2,1)(5,2,1)(6,2,1)(6,2,1)(2,1)(3,2,1)(4,2,1)(5,2,1)(5,2,1)%
}
\NovelDisplay{M05062R03}{%
(1,1,1)(2,1,1)(3,1,1)(3,1)(4,2,1)(5,2,1)(6,2,1)(6,2,1)(2,1)(3,1)%
}

\NovelDisplay{M05063R01}{%
\text{展开第二项：}\quad(0)(1,1,1)(2,1,1)(3,1,1)(3,1,1)(1,1,1)(2,1,1)(3,1,1)%
}
\NovelDisplay{M05063R02}{%
(3,1)(4,2,1)(5,2,1)(6,2,1)(6,2,1)(2,1)(3,2,1)(4,2,1)(5,2,1)(5,2,1)%
}
\NovelDisplay{M05063R03}{%
(1,1,1)(2,1,1)(3,1,1)(3,1)(4,2,1)(5,2,1)(6,2,1)(6,2,1)(2,1)(3,1)(4,1)%
}

同一原式的三次独立取项要分别保留原矩阵的位置，不能把抄错的前缀带进后面才补救。

这一次，他把自己那处错误说在前面，等叶阳指出还有一列没念到，便退回去，完整重说。

\NovelDisplay{M05064R01}{%
(0)(1,1,1)(2,1,1)(3,1,1)(3,1,1)(1,1,1)(2,1,1)(3,1,1)(3,1)(4,2,1)(5,2,1)%
}
\NovelDisplay{M05064R02}{%
(6,2,1)(6,2,1)(2,1)(3,2,1)(4,2,1)(5,2,1)(5,2,1)(3,2,1)(4,2,1)%
}
\NovelDisplay{M05064R03}{%
(5,2,1)(5,2)(6,3)%
}

\NovelDisplay{M05065R01}{%
= \psi(M_\omega + \psi_{\psi_{M_2}(M_2)}(M_\omega)+\psi_{M}(M_\omega+\psi_{\psi_{M_2}(M_2)}(M_2)))%
}

\NovelDisplay{M05066R01}{%
(0)(1,1,1)(2,1,1)(3,1,1)(3,1,1)(1,1,1)(2,1,1)(3,1,1)(3,1)(4,2,1)(5,2,1)%
}
\NovelDisplay{M05066R02}{%
(6,2,1)(6,2,1)(2,1)(3,2,1)(4,2,1)(5,2,1)(5,2,1)(3,2,1)(4,2,1)(5,2,1)%
}
\NovelDisplay{M05066R03}{%
(5,2)(6,3,1)(7,3,1)(8,3,1)(8,3,1)%
}

\NovelDisplay{M05067R01}{%
= \psi(M_\omega + \psi_{\psi_{M_2}(M_2)}(M_\omega)+\psi_{M}(M_\omega+\psi_{\psi_{M_2}(M_2)}(M_\omega)))%
}

\NovelDisplay{M05068R01}{%
(0)(1,1,1)(2,1,1)(3,1,1)(3,1,1)(1,1,1)(2,1,1)(3,1,1)(3,1)(4,2,1)(5,2,1)%
}
\NovelDisplay{M05068R02}{%
(6,2,1)(6,2,1)(2,1)(3,2,1)(4,2,1)(5,2,1)(5,2,1)(3,2,1)(4,2,1)(5,2,1)%
}
\NovelDisplay{M05068R03}{%
(5,2)(6,3,1)(7,3,1)(8,3,1)(8,3,1)(4,2)(5,3)%
}

\NovelDisplay{M05069R01}{%
= \psi(M_\omega + \psi_{\psi_{M_2}(M_2)}(M_\omega)+M)%
}

\NovelDisplay{M05070R01}{%
(0)(1,1,1)(2,1,1)(3,1,1)(3,1,1)(1,1,1)(2,1,1)(3,1,1)(3,1)(4,2,1)(5,2,1)%
}
\NovelDisplay{M05070R02}{%
(6,2,1)(6,2,1)(2,1,1)%
}

\NovelDisplay{M05071R01}{%
= \psi(M_\omega + \psi_{\psi_{M_2}(M_2)}(M_\omega+1))%
}

外层背景不变而内层对象继续扩展，唐三终于学会先等小舞把一层读完，不在她开口时抢着补下一层。

\NovelDisplay{M05072R01}{%
(0)(1,1,1)(2,1,1)(3,1,1)(3,1,1)(1,1,1)(2,1,1)(3,1,1)(3,1)(4,2,1)(5,2,1)%
}
\NovelDisplay{M05072R02}{%
(6,2,1)(6,2,1)(2,1,1)(3,1)(2)%
}

\NovelDisplay{M05073R01}{%
= \psi(M_\omega + \psi_{\psi_{M_2}(M_2)}(M_\omega+M))%
}

\NovelDisplay{M05074R01}{%
(0)(1,1,1)(2,1,1)(3,1,1)(3,1,1)(1,1,1)(2,1,1)(3,1,1)(3,1)(4,2,1)(5,2,1)%
}
\NovelDisplay{M05074R02}{%
(6,2,1)(6,2,1)(2,1,1)(3,1)(2,1,1)(3,1)(4,2)%
}

\NovelDisplay{M05075R01}{%
= \psi(M_\omega + \psi_{\psi_{M_2}(M_2)}(M_\omega+M)+M)%
}

\NovelDisplay{M05076R01}{%
(0)(1,1,1)(2,1,1)(3,1,1)(3,1,1)(1,1,1)(2,1,1)(3,1,1)(3,1)(4,2,1)(5,2,1)%
}
\NovelDisplay{M05076R02}{%
(6,2,1)(6,2,1)(2,1,1)(3,1)(4,2)%
}

\NovelDisplay{M05077R01}{%
= \psi(M_\omega + \psi_{\psi_{M_2}(M_2)}(M_\omega+M)+M)%
}

\NovelDisplay{M05078R01}{%
(0)(1,1,1)(2,1,1)(3,1,1)(3,1,1)(1,1,1)(2,1,1)(3,1,1)(3,1)(4,2,1)(5,2,1)%
}
\NovelDisplay{M05078R02}{%
(6,2,1)(6,2,1)(2,1,1)(3,1,1)%
}

\NovelDisplay{M05079R01}{%
= \psi(M_\omega + \psi_{\psi_{M_2}(M_2)}(M_\omega+M+1))%
}

相邻式中既有相同右侧也有参数后继，大家照讲义保留每条矩阵，既不擅改对应，也不把相似当成展开证明。

那次以后，薛峰不再将所有人的纸包揽到自己面前。

他先与叶阳核同一个问题。叶阳说得慢，薛峰习惯在后面接一句，接到第三回，叶阳忽然将笔放下，问他能不能等自己说完。

薛峰一怔，收回手。

同一段重新开始，叶阳中间停了两回，终于自己找回前面的条件。等他说完，薛峰才指出刚才那处保留段应该再核。叶阳听懂了，反过来问出他纸上另一处没有写全的地方。

两个人都沉默片刻，随后各取回自己的纸。

第二天换诺当来，薛峰仍坐在原来那块石头旁。他这次先听，听到一个完整的问题，才把自己真正能确认的部分说给对方。

\NovelDisplay{M05080R01}{%
(0)(1,1,1)(2,1,1)(3,1,1)(3,1,1)(1,1,1)(2,1,1)(3,1,1)(3,1)(4,2,1)(5,2,1)%
}
\NovelDisplay{M05080R02}{%
(6,2,1)(6,2,1)(2,1,1)(3,1,1)(2,1,1)(3,1)(2)%
}

\NovelDisplay{M05081R01}{%
= \psi(M_\omega + \psi_{\psi_{M_2}(M_2)}(M_\omega+M\times 2))%
}

\NovelDisplay{M05082R01}{%
(0)(1,1,1)(2,1,1)(3,1,1)(3,1,1)(1,1,1)(2,1,1)(3,1,1)(3,1)(4,2,1)(5,2,1)%
}
\NovelDisplay{M05082R02}{%
(6,2,1)(6,2,1)(2,1,1)(3,1,1)(3,1)(4,2)%
}

\NovelDisplay{M05083R01}{%
= \psi(M_\omega + \psi_{\psi_{M_2}(M_2)}(M_\omega+\psi_{\psi_{M_2}(M_2)}(M_2)))%
}

\NovelDisplay{M05084R01}{%
(0)(1,1,1)(2,1,1)(3,1,1)(3,1,1)(1,1,1)(2,1,1)(3,1,1)(3,1)(4,2,1)(5,2,1)%
}
\NovelDisplay{M05084R02}{%
(6,2,1)(6,2,1)(2,1,1)(3,1,1)(3,1)(4,2,1)(5,2,1)(6,2,1)(6,2,1)%
}

\NovelDisplay{M05085R01}{%
= \psi(M_\omega + \psi_{\psi_{M_2}(M_2)}(M_\omega+\psi_{\psi_{M_2}(M_2)}(M_\omega)))%
}

\NovelDisplay{M05086R01}{%
(0)(1,1,1)(2,1,1)(3,1,1)(3,1,1)(1,1,1)(2,1,1)(3,1,1)(3,1)(4,2,1)(5,2,1)%
}
\NovelDisplay{M05086R02}{%
(6,2,1)(6,2,1)(4,2)%
}

\NovelDisplay{M05087R01}{%
= \psi(M_\omega + \psi_{M_2}(M_2))%
}

读到内层倍数与新的回接时，叶阳先为每一个加号找准所在的括号，避免把背景项接进错误的ψ。

\NovelDisplay{M05088R01}{%
(0)(1,1,1)(2,1,1)(3,1,1)(3,1,1)(1,1,1)(2,1,1)(3,1,1)(3,1)(4,2,1)(5,2,1)%
}
\NovelDisplay{M05088R02}{%
(6,2,1)(6,2,1)(4,2,1)(5,2,1)(6,2,1)(6,2)(7,3,1)(8,3,1)(9,3,1)(9,3,1)%
}

\NovelDisplay{M05089R01}{%
= \psi(M_\omega + \psi_{M_2}(M_\omega))%
}

\NovelDisplay{M05090R01}{%
(0)(1,1,1)(2,1,1)(3,1,1)(3,1,1)(1,1,1)(2,1,1)(3,1,1)(3,1)(4,2,1)(5,2,1)%
}
\NovelDisplay{M05090R02}{%
(6,2,1)(6,2,1)(4,2,1)(5,2,1)(6,2,1)(6,2)(7,3,1)(8,3,1)(9,3,1)%
}
\NovelDisplay{M05090R03}{%
(9,3,1)(5,2)(6,3)%
}

\NovelDisplay{M05091R01}{%
= \psi(M_\omega + M_2)%
}

\NovelDisplay{M05092R01}{%
(0)(1,1,1)(2,1,1)(3,1,1)(3,1,1)(1,1,1)(2,1,1)(3,1,1)(3,1)(4,2,1)(5,2,1)%
}
\NovelDisplay{M05092R02}{%
(6,2,1)(6,2,1)(4,2,1)(5,2,1)(6,2,1)(6,2)(7,3,1)(8,3,1)(9,3,1)%
}
\NovelDisplay{M05092R03}{%
(9,3,1)(7,3)%
}

\NovelDisplay{M05093R01}{%
= \psi(M_\omega + \psi_{M_3}(M_3))%
}

\NovelDisplay{M05094R01}{%
(0)(1,1,1)(2,1,1)(3,1,1)(3,1,1)(1,1,1)(2,1,1)(3,1,1)(3,1,1)%
}

\NovelDisplay{M05095R01}{%
= \psi(M_\omega \times 2)%
}

这里重新比较直接的\allowbreak{}\NovelInline{I06109}{M_2}\allowbreak{}、不同下标的ψ与整个背景的倍数，薛峰把几种变化分栏记下才没有相互覆盖。

在四百二十米附近，薛峰第一次主动走到最前面。

他原来总等霍雨浩先问清条件，再接着核。那天霍雨浩正陪诺当在下面重做，蔡须坤看过新页，仍有一处没有说清，薛峰却已经找到了能够先核的一小段。

他开口时声音很低，第一句便被风盖过去。卢本韦没听清，让他再说。他起先以为是自己讲错了，攥着纸想退回去，古月娜却指指旁边的岩影，让他到那里站好。

他重新说完，蔡须坤才低头，将自己原来那一处同他的记录对照。

``这里你接得比我清楚。''

薛峰抬头，像是不大相信。她没有再夸一遍，只请他从原矩阵开始，给后来的人完整说明。那一天，霍雨浩上来时，先坐到了听的人中间。

\NovelDisplay{M05096R01}{%
(0)(1,1,1)(2,1,1)(3,1,1)(3,1,1)(2,1)(3,2)%
}

\NovelDisplay{M05097R01}{%
= \psi(\Omega_{M_\omega+1})%
}

\NovelDisplay{M05098R01}{%
= \psi(M_{\omega+1})%
}

\NovelDisplay{M05099R01}{%
= \psi(N\times \omega + N)%
}

\NovelDisplay{M05100R01}{%
(0)(1,1,1)(2,1,1)(3,1,1)(3,1,1)(2,1)(3,2,1)(4,2,1)(5,2)(4)%
}

\NovelDisplay{M05101R01}{%
= \psi({M_{\omega+1}}^2)%
}

\NovelDisplay{M05102R01}{%
(0)(1,1,1)(2,1,1)(3,1,1)(3,1,1)(2,1)(3,2,1)(4,2,1)(5,2,1)(5,2)(6,3)%
}

\NovelDisplay{M05103R01}{%
= \psi(M_{\omega+2})%
}

\NovelDisplay{M05104R01}{%
(0)(1,1,1)(2,1,1)(3,1,1)(3,1,1)(2,1)(3,2,1)(4,2,1)(5,2,1)(5,2)(6,3,1)%
}
\NovelDisplay{M05104R02}{%
(7,3,1)(8,3,1)(8,3)(9,4)%
}

\NovelDisplay{M05105R01}{%
= \psi(M_{\omega+3})%
}

\allowbreak{}\NovelInline{I06110}{M_{\omega+1}}\allowbreak{}的对照含有\allowbreak{}\NovelInline{I06111}{N\times\omega+N}\allowbreak{}，下标的后继因此不能被随手写成外式末尾的一个自然数一。

\NovelDisplay{M05106R01}{%
(0)(1,1,1)(2,1,1)(3,1,1)(3,1,1)(2,1)(3,2,1)(4,2,1)(5,2,1)(5,2,1)%
}

\NovelDisplay{M05107R01}{%
= \psi(M_{\omega\times 2})%
}

\NovelDisplay{M05108R01}{%
= \psi(N\times \omega\times 2)%
}

\NovelDisplay{M05109R01}{%
(0)(1,1,1)(2,1,1)(3,1,1)(3,1,1)(2,1)(3,2,1)(4,2,1)(5,2,1)(5,2,1)(4,2)%
}
\NovelDisplay{M05109R02}{%
(5,3,1)(6,3,1)(7,3,1)(7,3,1)%
}

\NovelDisplay{M05110R01}{%
= \psi(M_{\omega\times 3})%
}

\NovelDisplay{M05111R01}{%
(0)(1,1,1)(2,1,1)(3,1,1)(3,1,1)(2,1,1)%
}

\NovelDisplay{M05112R01}{%
= \psi(M_{\psi(2)})%
}

\NovelDisplay{M05113R01}{%
= \psi(N\times \psi(2))%
}

\NovelDisplay{M05114R01}{%
(0)(1,1,1)(2,1,1)(3,1,1)(3,1,1)(2,1,1)(3,1)%
}

\NovelDisplay{M05115R01}{%
= \psi(M_{\psi_{M}(M)})%
}

\NovelDisplay{M05116R01}{%
= \psi(N\times \Omega)%
}

\NovelDisplay{M05117R01}{%
= \psi(N\times \psi_{\psi_{N}(N)}(N))%
}

从ω倍数的下标过渡到ψ给出的下标时，卢本韦保留各个折叠符号，逐条核对而没有套用普通乘法的交换。

蔡须坤的手又发过一次冷。

那天没有混沌之力暴走，也没有谁突然失控，只是接连几日用力以后，她本想一口带过的气息忽然涩住。她站在台阶上，半边肩膀都僵着，仍不愿先开口。

卢本韦看出她翻页变慢，将I册从她手里接过去。霍雨浩站近些，只问她，今天能不能到这里为止。

蔡须坤看着上面。那几级已经反复走过，放在平日，她原本不该在这里停。

``能。''她终于说。

回去以后，她自己慢慢运转已经压稳的那一段。果糕领域只在身边铺开很小一层，帮她将外界的纷乱与自己的呼吸分开；领域没有替她算出任何一项，也没有改变石阶认可的接续。

霍雨浩坐在门边，等她伸手找水时才把杯子递来。她握住，过了一会儿，指尖终于暖了。

\NovelDisplay{M05118R01}{%
(0)(1,1,1)(2,1,1)(3,1,1)(3,1,1)(2,1,1)(3,1)(1,1,1)(2,1,1)(3,1,1)(3,1)%
}
\NovelDisplay{M05118R02}{%
(4,2,1)(5,2,1)(6,2,1)(6,2,1)%
}

\NovelDisplay{M05119R01}{%
= \psi(M_{\psi_{M}(M_\omega)})%
}

\NovelDisplay{M05120R01}{%
= \psi(N\times \psi_{M}(M_\omega))%
}

\NovelDisplay{M05121R01}{%
(0)(1,1,1)(2,1,1)(3,1,1)(3,1,1)(2,1,1)(3,1)(1,1,1)(2,1,1)(3,1,1)(3,1)%
}
\NovelDisplay{M05121R02}{%
(4,2,1)(5,2,1)(6,2,1)(6,2,1)(5,2,1)(6,1)%
}

\NovelDisplay{M05122R01}{%
= \psi(M_{\psi_{M}(M_{\psi_{M}(M)})})%
}

\NovelDisplay{M05123R01}{%
= \psi(N\times \psi_{M}(M_\Omega))%
}

\NovelDisplay{M05124R01}{%
(0)(1,1,1)(2,1,1)(3,1,1)(3,1,1)(2,1,1)(3,1)(1,1,1)(2,1,1)(3,1,1)(3,1)%
}
\NovelDisplay{M05124R02}{%
(4,2,1)(5,2,1)(6,2,1)(6,2,1)(5,2,1)(6,1)(2)%
}

\NovelDisplay{M05125R01}{%
= \psi(M_{M})%
}

\NovelDisplay{M05126R01}{%
= \psi(N\times \psi_{N}(N))%
}

\NovelDisplay{M05127R01}{%
(0)(1,1,1)(2,1,1)(3,1,1)(3,1,1)(2,1,1)(3,1)(1,1,1)(2,1,1)(3,1,1)(3,1)%
}
\NovelDisplay{M05127R02}{%
(4,2,1)(5,2,1)(6,2,1)(6,2,1)(5,2,1)(6,2)%
}

\NovelDisplay{M05128R01}{%
= \psi(M_{\Omega_{M+1}})%
}

\NovelDisplay{M05129R01}{%
= \psi(N\times \Omega_{M+1})%
}

下标本身也含有嵌套时，先认清下标，再看整个\allowbreak{}\NovelInline{I06112}{M}\allowbreak{}项的对应，能避免将两次作用压成一次。

\NovelDisplay{M05130R01}{%
(0)(1,1,1)(2,1,1)(3,1,1)(3,1,1)(2,1,1)(3,1)(1,1,1)(2,1,1)(3,1,1)(3,1,1)%
}

\NovelDisplay{M05131R01}{%
= \psi(M_{M_\omega})%
}

\NovelDisplay{M05132R01}{%
= \psi(N\times \psi_{N}(N\times \omega))%
}

\NovelDisplay{M05133R01}{%
(0)(1,1,1)(2,1,1)(3,1,1)(3,1,1)(2,1,1)(3,1)(1,1,1)(2,1,1)(3,1,1)%
}
\NovelDisplay{M05133R02}{%
(3,1,1)(2,1,1)(3,1)%
}

\NovelDisplay{M05134R01}{%
= \psi(M_{M_{\psi_{M}(M)}})%
}

\NovelDisplay{M05135R01}{%
= \psi(N\times \psi_{N}(N\times \Omega))%
}

\NovelDisplay{M05136R01}{%
(0)(1,1,1)(2,1,1)(3,1,1)(3,1,1)(2,1,1)(3,1)(2)%
}

\NovelDisplay{M05137R01}{%
= \psi(M(1,0))%
}

\NovelDisplay{M05138R01}{%
= \psi(M(1;0))%
}

\NovelDisplay{M05139R01}{%
= \psi(N^2)%
}

\NovelDisplay{M05140R01}{%
(0)(1,1,1)(2,1,1)(3,1,1)(3,1,1)(2,1,1)(3,1)(2)(1,1,1)(2,1,1)(3,1,1)%
}
\NovelDisplay{M05140R02}{%
(3,1,1)(2,1,1)(3,1)(2)%
}

\NovelDisplay{M05141R01}{%
= \psi(M(1,0)+\psi_{M(1,0)}(M(1,0)))%
}

新的双参数M记号与\allowbreak{}\NovelInline{I06113}{N^2}\allowbreak{}在这里相接，逗号和分号都属于原有写法，不能像矩阵末尾的补零那样省去。

\NovelDisplay{M05142R01}{%
(0)(1,1,1)(2,1,1)(3,1,1)(3,1,1)(2,1,1)(3,1)(2,1)(3,2)%
}

\NovelDisplay{M05143R01}{%
= \psi(M(1,0)+\Omega_{\psi_{M(1,0)}(M(1,0))+1})%
}

\NovelDisplay{M05144R01}{%
(0)(1,1,1)(2,1,1)(3,1,1)(3,1,1)(2,1,1)(3,1)(2,1)(3,2,1)(4,2,1)%
}
\NovelDisplay{M05144R02}{%
(5,2,1)(5,2,1)%
}

\NovelDisplay{M05145R01}{%
= \psi(M(1,0)+M_{\psi_{M(1,0)}(M(1,0))+\omega})%
}

\NovelDisplay{M05146R01}{%
(0)(1,1,1)(2,1,1)(3,1,1)(3,1,1)(2,1,1)(3,1)(2,1)(3,2,1)(4,2,1)(5,2,1)%
}
\NovelDisplay{M05146R02}{%
(5,2,1)(4,2,1)(5,2)(4)%
}

\NovelDisplay{M05147R01}{%
= \psi(M(1,0)\times 2)%
}

\NovelDisplay{M05148R01}{%
(0)(1,1,1)(2,1,1)(3,1,1)(3,1,1)(2,1,1)(3,1)(4,2)%
}

\NovelDisplay{M05149R01}{%
= \psi(\Omega_{M(1,0)+1})%
}

唐三把双参数项外面的加法同参数内部的改动分开，第一次能在不翻答案时说清自己究竟改变了哪里。

\NovelDisplay{M05150R01}{%
(0)(1,1,1)(2,1,1)(3,1,1)(3,1,1)(2,1,1)(3,1)(4,2,1)(5,2,1)(6,2,1)(6,2,1)%
}

\NovelDisplay{M05151R01}{%
= \psi(M_{M(1,0)+1})%
}

\NovelDisplay{M05152R01}{%
= \psi(\psi_{M(1,1)}(1))%
}

\NovelDisplay{M05153R01}{%
(0)(1,1,1)(2,1,1)(3,1,1)(3,1,1)(2,1,1)(3,1)(4,2,1)(5,2,1)(6,2,1)%
}
\NovelDisplay{M05153R02}{%
(6,2,1)(5,2,1)(6,2)%
}

\NovelDisplay{M05154R01}{%
= \psi(M_{\Omega_{M(1,0)+1}})%
}

\NovelDisplay{M05155R01}{%
(0)(1,1,1)(2,1,1)(3,1,1)(3,1,1)(2,1,1)(3,1)(4,2,1)(5,2,1)(6,2,1)(6,2,1)%
}
\NovelDisplay{M05155R02}{%
(5,2,1)(6,2)(5)%
}

\NovelDisplay{M05156R01}{%
= \psi(M(1,1))%
}

\NovelDisplay{M05157R01}{%
= \psi(M(1;0)\times 2)%
}

\NovelDisplay{M05158R01}{%
= \psi(N^2+\psi_{N}(N^2))%
}

\NovelDisplay{M05159R01}{%
(0)(1,1,1)(2,1,1)(3,1,1)(3,1,1)(2,1,1)(3,1)(4,2,1)(5,2,1)(6,2,1)(6,2,1)%
}
\NovelDisplay{M05159R02}{%
(5,2,1)(6,2)(7,3)%
}

\NovelDisplay{M05160R01}{%
= \psi(\Omega_{M(1,1)+1})%
}

这一组同时用到新参数与折叠回接，蔡须坤让诺当先读最外框，再逐层把内层参数补齐。

唐三与小舞真正说起落后的事，是在一个傍晚。

小舞先问，唐三是不是觉得自己一直拖着他。唐三听见，愣了半天，说自己这几日反而常要她帮着看遗漏的那一列。

``你以前不是这么说的。''

唐三望着她，没有急着解释。他想起在铺子里吵钱、吵谁干的活多，也想起上山以后，自己常一边听妹妹问，一边已经准备替她回答。

``以前有些话，我说错了。''

小舞低头，把他们刚才一起算的纸铺平。她没有说一句就算了，只让他从明天起别在她没讲完时抢答案。

唐三答应。第二天她停得很久，他几次想开口，都先把手收了回去。等她自己说清，蔡须坤在旁边点头，小舞才悄悄松了攥紧的纸角。

\NovelDisplay{M05161R01}{%
(0)(1,1,1)(2,1,1)(3,1,1)(3,1,1)(2,1,1)(3,1)(4,2,1)(5,2,1)(6,2,1)(6,2,1)%
}
\NovelDisplay{M05161R02}{%
(5,2,1)(6,2)(7,3,1)(8,3,1)(9,3,1)(9,3,1)(8,3,1)(9,3)(8)%
}

\NovelDisplay{M05162R01}{%
= \psi(M(1,2))%
}

\NovelDisplay{M05163R01}{%
= \psi(M(1;0)\times 3)%
}

\NovelDisplay{M05164R01}{%
= \psi(N^2+\psi_{N}(N^2)\times 2)%
}

\NovelDisplay{M05165R01}{%
(0)(1,1,1)(2,1,1)(3,1,1)(3,1,1)(2,1,1)(3,1,1)%
}

\NovelDisplay{M05166R01}{%
= \psi(M(1,\omega))%
}

\NovelDisplay{M05167R01}{%
= \psi(M(1;0)\times \omega)%
}

\NovelDisplay{M05168R01}{%
= \psi(N^2+\psi_{N}(N^2)\times \omega)%
}

\NovelDisplay{M05169R01}{%
(0)(1,1,1)(2,1,1)(3,1,1)(3,1,1)(2,1,1)(3,1,1)(2,1,1)(3,1)%
}

\NovelDisplay{M05170R01}{%
= \psi(M(1,\Omega))%
}

\NovelDisplay{M05171R01}{%
= \psi(M(1;0)\times \Omega)%
}

\NovelDisplay{M05172R01}{%
= \psi(N^2+\psi_{N}(N^2)\times \Omega)%
}

\NovelDisplay{M05173R01}{%
(0)(1,1,1)(2,1,1)(3,1,1)(3,1,1)(2,1,1)(3,1,1)(2,1,1)(3,1)(1,1,1)(2,1,1)%
}
\NovelDisplay{M05173R02}{%
(3,1,1)(3,1,1)(2,1,1)(3,1,1)(2,1,1)(3,1)%
}

\NovelDisplay{M05174R01}{%
= \psi(M(1,M(1,\Omega)))%
}

\NovelDisplay{M05175R01}{%
= \psi(M(1;0)\times \psi_{M(1;0)}(M(1;0)\times \Omega))%
}

第二参数从有限项进入ω及Ω的情形，仍须保留分号另一侧的既定结构，不能只看最后一个字符判断层次。

\NovelDisplay{M05176R01}{%
(0)(1,1,1)(2,1,1)(3,1,1)(3,1,1)(2,1,1)(3,1,1)(2,1,1)(3,1)(2)%
}

\NovelDisplay{M05177R01}{%
= \psi(M(2,0))%
}

\NovelDisplay{M05178R01}{%
= \psi({M(1;0)}^2)%
}

\NovelDisplay{M05179R01}{%
(0)(1,1,1)(2,1,1)(3,1,1)(3,1,1)(2,1,1)(3,1,1)(2,1,1)(3,1)(2)(1,1,1)%
}
\NovelDisplay{M05179R02}{%
(2,1,1)(3,1,1)(3,1,1)(2,1,1)(3,1,1)(2,1,1)(3,1)(2)%
}

\NovelDisplay{M05180R01}{%
= \psi(M(2,0)+\psi_{M(2,0)}(M(2,0)))%
}

\NovelDisplay{M05181R01}{%
(0)(1,1,1)(2,1,1)(3,1,1)(3,1,1)(2,1,1)(3,1,1)(2,1,1)(3,1)(2,1)(3,2)%
}

\NovelDisplay{M05182R01}{%
= \psi(M(2,0)+\Omega_{\psi_{M(2,0)}(M(2,0))+1})%
}

\NovelDisplay{M05183R01}{%
(0)(1,1,1)(2,1,1)(3,1,1)(3,1,1)(2,1,1)(3,1,1)(2,1,1)(3,1)(2,1)(3,2,1)%
}
\NovelDisplay{M05183R02}{%
(4,2,1)(5,2,1)(5,2,1)%
}

\NovelDisplay{M05184R01}{%
= \psi(M(2,0)+M_{\psi_{M(2,0)}(M(2,0))+\omega})%
}

\allowbreak{}\NovelInline{I06114}{M(2,0)}\allowbreak{}与\allowbreak{}\NovelInline{I06115}{M(1;0)^2}\allowbreak{}的对照固定以后，随后的加项才能在这个共同背景下逐条比较。

\NovelDisplay{M05185R01}{%
(0)(1,1,1)(2,1,1)(3,1,1)(3,1,1)(2,1,1)(3,1,1)(2,1,1)(3,1)(2,1)(3,2,1)%
}
\NovelDisplay{M05185R02}{%
(4,2,1)(5,2,1)(5,2,1)(4,2,1)(5,2)(4)%
}

\NovelDisplay{M05186R01}{%
= \psi(M(2,0)+M(1,\psi_{M(2,0)}(M(2,0))+1))%
}

\NovelDisplay{M05187R01}{%
(0)(1,1,1)(2,1,1)(3,1,1)(3,1,1)(2,1,1)(3,1,1)(2,1,1)(3,1)(2,1)(3,2,1)%
}
\NovelDisplay{M05187R02}{%
(4,2,1)(5,2,1)(5,2,1)(4,2,1)(5,2,1)%
}

\NovelDisplay{M05188R01}{%
= \psi(M(2,0)+M(1,\psi_{M(2,0)}(M(2,0))+\omega))%
}

\NovelDisplay{M05189R01}{%
(0)(1,1,1)(2,1,1)(3,1,1)(3,1,1)(2,1,1)(3,1,1)(2,1,1)(3,1)(2,1)(3,2,1)%
}
\NovelDisplay{M05189R02}{%
(4,2,1)(5,2,1)(5,2,1)(4,2,1)(5,2,1)(4,2,1)(5,2)%
}

\NovelDisplay{M05190R01}{%
= \psi(M(2,0)+M(1,\Omega_{\psi_{M(2,0)}(M(2,0))+1}))%
}

\NovelDisplay{M05191R01}{%
(0)(1,1,1)(2,1,1)(3,1,1)(3,1,1)(2,1,1)(3,1,1)(2,1,1)(3,1)(2,1)(3,2,1)%
}
\NovelDisplay{M05191R02}{%
(4,2,1)(5,2,1)(5,2,1)(4,2,1)(5,2,1)(4,2,1)(5,2)(4)%
}

\NovelDisplay{M05192R01}{%
= \psi(M(2,0)\times 2)%
}

\NovelDisplay{M05193R01}{%
= \psi({M(1;0)}^2+\psi_{M(1;0)}({M(1;0)}^2))%
}

新添参数项和整个平方的倍数在这一组并列，卢本韦读到乘二时停了一下，先确认乘号确在外面。

\NovelDisplay{M05194R01}{%
(0)(1,1,1)(2,1,1)(3,1,1)(3,1,1)(2,1,1)(3,1,1)(2,1,1)(3,1)(2,1)(3,2,1)%
}
\NovelDisplay{M05194R02}{%
(4,2,1)(5,2,1)(5,2,1)(4,2,1)(5,2,1)(4,2,1)(5,2)(4)%
}

\NovelDisplay{M05195R01}{%
= \psi(M(2,0)\times 2)%
}

\NovelDisplay{M05196R01}{%
= \psi({M(1;0)}^2+\psi_{M(1;0)}({M(1;0)}^2))%
}

\NovelDisplay{M05197R01}{%
(0)(1,1,1)(2,1,1)(3,1,1)(3,1,1)(2,1,1)(3,1,1)(2,1,1)(3,1)(2,1,1)%
}

\NovelDisplay{M05198R01}{%
= \psi(M(2,0)\times \omega)%
}

\NovelDisplay{M05199R01}{%
(0)(1,1,1)(2,1,1)(3,1,1)(3,1,1)(2,1,1)(3,1,1)(2,1,1)(3,1)(4,2)%
}

\NovelDisplay{M05200R01}{%
= \psi(\Omega_{M(2,0)+1})%
}

\NovelDisplay{M05201R01}{%
(0)(1,1,1)(2,1,1)(3,1,1)(3,1,1)(2,1,1)(3,1,1)(2,1,1)(3,1)(4,2,1)(5,2,1)%
}
\NovelDisplay{M05201R02}{%
(6,2,1)(6,2,1)%
}

\NovelDisplay{M05202R01}{%
= \psi(M_{M(2,0)+\omega})%
}

讲义中的重复项仍保留作复核，随后再区分倍数、Ω下标和新的双参数尾部，不靠删去一行来掩盖差别。

到第七个月将尽，先前那张山图已经换了第二张外套。

蔡兴坤缝袋子时学的针脚派上了用场，他把旧布留在里面，又在外面补一层。黑子鸡不知从哪里衔来半根枯草，往布上一放，他看了半天，郑重地把草也收了起来。

``这算帮忙？''奥斯卡问。

蔡兴坤反问他，前几日只在旁边喊别忘了吃饭算不算。奥斯卡想了想，把最后一小块补给塞进嘴里，说自己下回多做一点。

笑声从屋里传出去。卢本韦正将那页新对照夹进I册后面的补页，听见以后，抬头看了蔡须坤一眼。蔡须坤也在笑，手边那处曾让她几日不肯停的问号，已经有了能独立复述的说明。

\NovelDisplay{M05203R01}{%
(0)(1,1,1)(2,1,1)(3,1,1)(3,1,1)(2,1,1)(3,1,1)(2,1,1)(3,1)(4,2,1)(5,2,1)%
}
\NovelDisplay{M05203R02}{%
(6,2,1)(6,2,1)(5,2,1)(6,2)(5)%
}

\NovelDisplay{M05204R01}{%
= \psi(M(1,M(2,0)+1))%
}

\NovelDisplay{M05205R01}{%
= \psi({M(1;0)}^2+{M(1;0)})%
}

\NovelDisplay{M05206R01}{%
(0)(1,1,1)(2,1,1)(3,1,1)(3,1,1)(2,1,1)(3,1,1)(2,1,1)(3,1)(4,2,1)(5,2,1)%
}
\NovelDisplay{M05206R02}{%
(6,2,1)(6,2,1)(5,2,1)(6,2,1)%
}

\NovelDisplay{M05207R01}{%
= \psi(M(1,M(2,0)+\omega))%
}

\NovelDisplay{M05208R01}{%
= \psi({M(1;0)}^2+{M(1;0)}\times \omega)%
}

\NovelDisplay{M05209R01}{%
(0)(1,1,1)(2,1,1)(3,1,1)(3,1,1)(2,1,1)(3,1,1)(2,1,1)(3,1)(4,2,1)(5,2,1)%
}
\NovelDisplay{M05209R02}{%
(6,2,1)(6,2,1)(5,2,1)(6,2,1)(5,2,1)(6,2)(5)%
}

\NovelDisplay{M05210R01}{%
= \psi(M(2,1))%
}

\NovelDisplay{M05211R01}{%
= \psi({M(1;0)}^2\times 2)%
}

\NovelDisplay{M05212R01}{%
(0)(1,1,1)(2,1,1)(3,1,1)(3,1,1)(2,1,1)(3,1,1)(2,1,1)(3,1)(4,2,1)(5,2,1)%
}
\NovelDisplay{M05212R02}{%
(6,2,1)(6,2,1)(5,2,1)(6,2,1)(5,2,1)(6,2)(7,3)%
}

\NovelDisplay{M05213R01}{%
= \psi(\Omega_{M(2,1)+1})%
}

\NovelDisplay{M05214R01}{%
(0)(1,1,1)(2,1,1)(3,1,1)(3,1,1)(2,1,1)(3,1,1)(2,1,1)(3,1)(4,2,1)(5,2,1)%
}
\NovelDisplay{M05214R02}{%
(6,2,1)(6,2,1)(5,2,1)(6,2,1)(5,2,1)(6,2)(7,3,1)(8,3,1)(9,3,1)(9,3,1)%
}
\NovelDisplay{M05214R03}{%
(8,3,1)(9,3,1)(8,3,1)(9,3)(8)%
}

\NovelDisplay{M05215R01}{%
= \psi(M(2,2))%
}

\NovelDisplay{M05216R01}{%
= \psi({M(1;0)}^2\times 3)%
}

这里先在平方背景后接新的参数项，再进入平方的倍数，蔡须坤让薛峰同时核对分号位置和共同保留段。

临上横台前，蔡须坤把本组里最长的矩阵交给卢本韦，自己求另一项，霍雨浩等两人分别写完，才让薛峰拿出他的第三份结果。

末列在这里又只剩第一行。卢本韦从二十六列的原式回找，指出第十八列为坏根，第十八至第二十五列合成八列的一段。她先把这个范围说全，才请蔡须坤核对。

第零项保留二十五列，第一项再生成八列，第二项再按自己的编号生成两份。她没有给这些列沿用上一个例子的提升差，读到末列最高非零位置时便重新依规则判断。

\NovelDisplay{M05217R01}{%
\text{展开第零项：}\quad(0)(1,1,1)(2,1,1)(3,1,1)(3,1,1)(2,1,1)(3,1,1)(2,1,1)%
}
\NovelDisplay{M05217R02}{%
(3,1)(4,2,1)(5,2,1)(6,2,1)(6,2,1)(5,2,1)(6,2,1)(5,2,1)(6,2)(7,3,1)%
}
\NovelDisplay{M05217R03}{%
(8,3,1)(9,3,1)(9,3,1)(8,3,1)(9,3,1)(8,3,1)(9,3)%
}

\NovelDisplay{M05218R01}{%
\text{展开第一项：}\quad(0)(1,1,1)(2,1,1)(3,1,1)(3,1,1)(2,1,1)(3,1,1)(2,1,1)%
}
\NovelDisplay{M05218R02}{%
(3,1)(4,2,1)(5,2,1)(6,2,1)(6,2,1)(5,2,1)(6,2,1)(5,2,1)(6,2)(7,3,1)%
}
\NovelDisplay{M05218R03}{%
(8,3,1)(9,3,1)(9,3,1)(8,3,1)(9,3,1)(8,3,1)(9,3)(7,3,1)(8,3,1)(9,3,1)%
}
\NovelDisplay{M05218R04}{%
(9,3,1)(8,3,1)(9,3,1)(8,3,1)(9,3)%
}

\NovelDisplay{M05219R01}{%
\text{展开第二项：}\quad(0)(1,1,1)(2,1,1)(3,1,1)(3,1,1)(2,1,1)(3,1,1)(2,1,1)%
}
\NovelDisplay{M05219R02}{%
(3,1)(4,2,1)(5,2,1)(6,2,1)(6,2,1)(5,2,1)(6,2,1)(5,2,1)(6,2)(7,3,1)%
}
\NovelDisplay{M05219R03}{%
(8,3,1)(9,3,1)(9,3,1)(8,3,1)(9,3,1)(8,3,1)(9,3)(7,3,1)(8,3,1)(9,3,1)%
}
\NovelDisplay{M05219R04}{%
(9,3,1)(8,3,1)(9,3,1)(8,3,1)(9,3)(7,3,1)(8,3,1)(9,3,1)(9,3,1)(8,3,1)%
}
\NovelDisplay{M05219R05}{%
(9,3,1)(8,3,1)(9,3)%
}

三份结果由同一个原式各自求得，重复段有多长必须逐列核清，不能由右式的乘三反推一个未经计算的拼接。

小舞看过以后，也取回自己的原矩阵，没有将同伴写出的结果直接抄作自己已经走过。

休息那晚，霍雨浩听见门外有人翻纸，出去看时，蔡须坤正想把一处刚记下的问题看完。

他没替她合本，只在旁边坐下，问她明天准备先从哪里开口。她指出那一处，又问他，今日等自己退回来时，是不是也着急。

``着急。''

她抬头。霍雨浩望着已经暗下去的山道，过了一会儿才接道：``可我知道你会回来。等你说清，总比我们各自把不知道的接过去好。''

蔡须坤将那页折好。这一次，她自己起身，把灯与本子一同带回了屋里。

屋里的灯亮着，卢本韦还在等她。蔡须坤把问题页留下，明日第一个要说的地方，已不用再用纸角藏着了。

四百三十米的横台，石面早已被长年的风磨得很平。

薛峰这一次最先看清了最后那处接续，却没有催后来的人。他在边上坐下，听叶阳将自己的问题独立说完，再转头看诺当。

唐三与小舞到了以后，先把本子分开，各自在页角留下日期。古月娜将空水壶收起，奥斯卡靠着石壁喘了半天，才想起袋里还有一份留给最后上来的蔡兴坤。

蔡须坤坐到卢本韦身旁。七个月过去，海边的天色似乎和第一天一样，变的却是她们手里那两本不断加厚的册子。

``I后面还接得下吗？''她轻声问。

卢本韦把补页摊给她看。能由旧表示直接说明的，留在原页；要用新记号的，就认真写在后面。她不再因为某一页不能包住眼前所有东西，就把先前学过的全算成白费。

``接得下。''她说，``先把这一页写清。''

霍雨浩把山图压好，等墨干。薛峰回头问明天谁先讲，小舞举起手，又看了唐三一眼。

``我先问。''她改口。

唐三点头。山风从横台上过去，十一人的纸角一齐动了动，没有谁急着先把别人的那份合上。

\chapter{第一百二十八章：风雪留下的空行}\label{ux7b2cux4e00ux767eux4e8cux5341ux516bux7ae0ux98ceux96eaux7559ux4e0bux7684ux7a7aux884c}

第八个月，山路旁的树叶开始一片片变色。

四百三十米以上，有一段石阶总迎着风。蔡须坤第一次走上去时，还能听见身后唐三在问；再走几级，声音便被风切断，只看见小舞的嘴在动。

她回身，等两人靠近些，才听清刚才那个问题。唐三问的是参数里的分组，小舞听成了外面的幂，已经在另一页上越接越远。

从那以后，他们问话之前先找能听清的地方。霍雨浩把几处避风的横台画出来，薛峰与叶阳轮流照看后来的人；古月娜带着水与干纸，等需要的人停下，再把它们递过去。

有时一天只能将一个问题从问到答走完。山图上空着的几格越来越多，蔡须坤却没有把日期划掉。那都是他们真正留在山里的日子。

\NovelDisplay{M05220R01}{%
(0)(1,1,1)(2,1,1)(3,1,1)(3,1,1)(2,1,1)(3,1,1)(2,1,1)(3,1,1)%
}

\NovelDisplay{M05221R01}{%
= \psi(M(2,\omega))%
}

\NovelDisplay{M05222R01}{%
= \psi({M(1;0)}^2\times \omega)%
}

\NovelDisplay{M05223R01}{%
(0)(1,1,1)(2,1,1)(3,1,1)(3,1,1)(2,1,1)(3,1,1)(2,1,1)(3,1,1)(2,1,1)%
}
\NovelDisplay{M05223R02}{%
(3,1)(2)%
}

\NovelDisplay{M05224R01}{%
= \psi(M(3,0))%
}

\NovelDisplay{M05225R01}{%
= \psi({M(1;0)}^3)%
}

\NovelDisplay{M05226R01}{%
(0)(1,1,1)(2,1,1)(3,1,1)(3,1,1)(2,1,1)(3,1,1)(2,1,1)(3,1,1)(2,1,1)%
}
\NovelDisplay{M05226R02}{%
(3,1,1)(2,1,1)(3,1)(2)%
}

\NovelDisplay{M05227R01}{%
= \psi(M(4,0))%
}

\NovelDisplay{M05228R01}{%
= \psi({M(1;0)}^4)%
}

\NovelDisplay{M05229R01}{%
(0)(1,1,1)(2,1,1)(3,1,1)(3,1,1)(2,1,1)(3,1,1)(3)%
}

\NovelDisplay{M05230R01}{%
= \psi(M(\omega,0))%
}

\NovelDisplay{M05231R01}{%
= \psi({M(1;0)}^\omega)%
}

从ω倍数到有限幂及ω次幂，双参数M的外层写法都要完整保留，不能只抄右侧最显眼的指数。

多出一个参数位置时，蔡须坤最初答得很快。

古月娜请她再指一次分隔符，她才发觉，自己口里说的是前一页，手上落的却已是新一页。她想重来，气息被压力压住，连原先稳住的起手也一并散了。

霍雨浩让她先退到平处。她没有硬撑，自己数着脚下的石阶，一步步退回去。

那天傍晚，她把那几级重新画到图上。卢本韦问是否要记得这样细，她点头，将自己答得太快的那句也留下来。

``下次看见它，我想先停一下。''

\NovelDisplay{M05232R01}{%
(0)(1,1,1)(2,1,1)(3,1,1)(3,1,1)(2,1,1)(3,1,1)(3,1)(2)%
}

\NovelDisplay{M05233R01}{%
= \psi(M(1,0,0))%
}

\NovelDisplay{M05234R01}{%
= \psi({M(1;0)}^{M(1;0)})%
}

\NovelDisplay{M05235R01}{%
(0)(1,1,1)(2,1,1)(3,1,1)(3,1,1)(2,1,1)(3,1,1)(3,1)(2,1,1)%
}

\NovelDisplay{M05236R01}{%
= \psi(M(1,0,\omega))%
}

\NovelDisplay{M05237R01}{%
= \psi({M(1;0)}^{M(1;0)}\times \omega)%
}

\NovelDisplay{M05238R01}{%
(0)(1,1,1)(2,1,1)(3,1,1)(3,1,1)(2,1,1)(3,1,1)(3,1)(2,1,1)(3,1)(2)%
}

\NovelDisplay{M05239R01}{%
= \psi(M(1,1,0))%
}

\NovelDisplay{M05240R01}{%
= \psi({M(1;0)}^{M(1;0)+1})%
}

\NovelDisplay{M05241R01}{%
(0)(1,1,1)(2,1,1)(3,1,1)(3,1,1)(2,1,1)(3,1,1)(3,1)(2,1,1)(3,1,1)(3,1)(2)%
}

\NovelDisplay{M05242R01}{%
= \psi(M(2,0,0))%
}

\NovelDisplay{M05243R01}{%
= \psi({M(1;0)}^{M(1;0)\times 2})%
}

再增加参数位置后，幂指数本身也升级为完整的M项，霍雨浩提醒小舞先读完上标再回到底数。

\NovelDisplay{M05244R01}{%
(0)(1,1,1)(2,1,1)(3,1,1)(3,1,1)(2,1,1)(3,1,1)(3,1)(3,1)(2)%
}

\NovelDisplay{M05245R01}{%
= \psi(M(1,0,0,0))%
}

\NovelDisplay{M05246R01}{%
= \psi({M(1;0)}^{{M(1;0)}^2})%
}

\NovelDisplay{M05247R01}{%
(0)(1,1,1)(2,1,1)(3,1,1)(3,1,1)(2,1,1)(3,1,1)(3,1)(4,2)%
}

\NovelDisplay{M05248R01}{%
= \psi(\Omega_{M(1;0)+1})%
}

\NovelDisplay{M05249R01}{%
(0)(1,1,1)(2,1,1)(3,1,1)(3,1,1)(2,1,1)(3,1,1)(3,1)(4,2,1)(5,2,1)%
}
\NovelDisplay{M05249R02}{%
(6,2,1)(6,2,1)%
}

\NovelDisplay{M05250R01}{%
= \psi(M_{M(1;0)+\omega})%
}

\NovelDisplay{M05251R01}{%
(0)(1,1,1)(2,1,1)(3,1,1)(3,1,1)(2,1,1)(3,1,1)(3,1)(4,2,1)(5,2,1)(6,2,1)%
}
\NovelDisplay{M05251R02}{%
(6,2,1)(5,2,1)(6,2)(5)%
}

\NovelDisplay{M05252R01}{%
= \psi(M(1,M(1;0)+1))%
}

\NovelDisplay{M05253R01}{%
= \psi(M(1;1))%
}

\NovelDisplay{M05254R01}{%
= \psi(N^2\times 2)%
}

逗号与分号的位置在这几式里承担不同分组，古月娜让大家逐字对照，没有凭外形相似交换参数。

\NovelDisplay{M05255R01}{%
(0)(1,1,1)(2,1,1)(3,1,1)(3,1,1)(2,1,1)(3,1,1)(3,1)(4,2,1)(5,2,1)(6,2,1)%
}
\NovelDisplay{M05255R02}{%
(6,2,1)(5,2,1)(6,2,1)(5,2,1)(6,2)(5)%
}

\NovelDisplay{M05256R01}{%
= \psi(M(2,M(1;0)+1))%
}

\NovelDisplay{M05257R01}{%
= \psi({M(1;1)}^2)%
}

\NovelDisplay{M05258R01}{%
(0)(1,1,1)(2,1,1)(3,1,1)(3,1,1)(2,1,1)(3,1,1)(3,1)(4,2,1)(5,2,1)(6,2,1)%
}
\NovelDisplay{M05258R02}{%
(6,2,1)(5,2,1)(6,2,1)(6,2)(5)%
}

\NovelDisplay{M05259R01}{%
= \psi(M(1,0,M(1;0)+1))%
}

\NovelDisplay{M05260R01}{%
= \psi({M(1;1)}^{M(1;1)})%
}

\NovelDisplay{M05261R01}{%
(0)(1,1,1)(2,1,1)(3,1,1)(3,1,1)(2,1,1)(3,1,1)(3,1)(4,2,1)(5,2,1)(6,2,1)%
}
\NovelDisplay{M05261R02}{%
(6,2,1)(5,2,1)(6,2,1)(6,2)(7,3)%
}

\NovelDisplay{M05262R01}{%
= \psi(\Omega_{M(1;1)+1})%
}

\NovelDisplay{M05263R01}{%
(0)(1,1,1)(2,1,1)(3,1,1)(3,1,1)(2,1,1)(3,1,1)(3,1)(4,2,1)(5,2,1)(6,2,1)%
}
\NovelDisplay{M05263R02}{%
(6,2,1)(5,2,1)(6,2,1)(6,2)(7,3,1)(8,3,1)(9,3,1)(9,3,1)(8,3,1)%
}
\NovelDisplay{M05263R03}{%
(9,3,1)(9,3)(8)%
}

\NovelDisplay{M05264R01}{%
= \psi(M(1;2))%
}

\NovelDisplay{M05265R01}{%
= \psi(N^2\times 3)%
}

背景换成\allowbreak{}\NovelInline{I06116}{M(1;1)}\allowbreak{}以后，外层平方与自幂仍须分别接续，不能把之前的背景名称照搬到新一行。

\NovelDisplay{M05266R01}{%
(0)(1,1,1)(2,1,1)(3,1,1)(3,1,1)(2,1,1)(3,1,1)(3,1,1)%
}

\NovelDisplay{M05267R01}{%
= \psi(M(1;\omega))%
}

\NovelDisplay{M05268R01}{%
= \psi(N^2\times \omega)%
}

\NovelDisplay{M05269R01}{%
(0)(1,1,1)(2,1,1)(3,1,1)(3,1,1)(2,1,1)(3,1,1)(3,1,1)(2,1,1)(3,1)(2)%
}

\NovelDisplay{M05270R01}{%
= \psi(M(1;1,0))%
}

\NovelDisplay{M05271R01}{%
= \psi(M(2;0))%
}

\NovelDisplay{M05272R01}{%
= \psi(N^3)%
}

\NovelDisplay{M05273R01}{%
(0)(1,1,1)(2,1,1)(3,1,1)(3,1,1)(2,1,1)(3,1,1)(3,1,1)(2,1,1)(3,1)(2,1,1)%
}

\NovelDisplay{M05274R01}{%
= \psi(M(1;1,0)\times \omega)%
}

\NovelDisplay{M05275R01}{%
(0)(1,1,1)(2,1,1)(3,1,1)(3,1,1)(2,1,1)(3,1,1)(3,1,1)(2,1,1)(3,1)(4,2)%
}

\NovelDisplay{M05276R01}{%
= \psi(\Omega_{M(1;1,0)+1})%
}

\allowbreak{}\NovelInline{I06117}{M(1;1,0)}\allowbreak{}与\allowbreak{}\NovelInline{I06118}{M(2;0)}\allowbreak{}给出同一处对照，后面的倍数与回接则继续沿各自完整表达核清。

真正的大风是在一个清晨来的。

十一人刚走到昨日的位置，便看见上方碎石不断落下。夏夜星空让他们退回避风处，叶阳先护住后面的蔡兴坤，奥斯卡伸手压住了几乎要被吹翻的包。

小舞仍望着上面。前一晚她刚能将一段独立说清，今日本来想趁记得牢，亲手走过去。

唐三没有劝她算了，只把那页拿出来，问能不能先在下面重新核一遍。

她看着纸角，过了一会儿，坐回去。那段没有因暂时不能抬脚就消失，她说过一遍，又听出自己昨晚还没留意的另一处停顿。

\NovelDisplay{M05277R01}{%
(0)(1,1,1)(2,1,1)(3,1,1)(3,1,1)(2,1,1)(3,1,1)(3,1,1)(2,1,1)(3,1)(4,2,1)%
}
\NovelDisplay{M05277R02}{%
(5,2,1)(6,2,1)(6,2,1)(5,2,1)(6,2)(6)%
}

\NovelDisplay{M05278R01}{%
= \psi(M(1,M(1;1,0)+1))%
}

\NovelDisplay{M05279R01}{%
= \psi(M(1;M(1;1,0)+1))%
}

\NovelDisplay{M05280R01}{%
(0)(1,1,1)(2,1,1)(3,1,1)(3,1,1)(2,1,1)(3,1,1)(3,1,1)(2,1,1)(3,1)(4,2,1)%
}
\NovelDisplay{M05280R02}{%
(5,2,1)(6,2,1)(6,2,1)(5,2,1)(6,2,1)(6,2)(6)%
}

\NovelDisplay{M05281R01}{%
= \psi(M(1,0,M(1;1,0)+1))%
}

\NovelDisplay{M05282R01}{%
= \psi({M(1;M(1;1,0)+1)}^{M(1;M(1;1,0)+1)})%
}

\NovelDisplay{M05283R01}{%
(0)(1,1,1)(2,1,1)(3,1,1)(3,1,1)(2,1,1)(3,1,1)(3,1,1)(2,1,1)(3,1)(4,2,1)%
}
\NovelDisplay{M05283R02}{%
(5,2,1)(6,2,1)(6,2,1)(5,2,1)(6,2,1)(6,2)(7,3)%
}

\NovelDisplay{M05284R01}{%
= \psi(\Omega_{M(1;M(1;1,0)+1)+1})%
}

\NovelDisplay{M05285R01}{%
(0)(1,1,1)(2,1,1)(3,1,1)(3,1,1)(2,1,1)(3,1,1)(3,1,1)(2,1,1)(3,1)(4,2,1)%
}
\NovelDisplay{M05285R02}{%
(5,2,1)(6,2,1)(6,2,1)(5,2,1)(6,2,1)(6,2)(7,3,1)(8,3,1)(9,3,1)(9,3,1)%
}
\NovelDisplay{M05285R03}{%
(8,3,1)(9,3)(9)%
}

\NovelDisplay{M05286R01}{%
= \psi(M(1,M(1;M(1;1,0)+1)+1))%
}

\NovelDisplay{M05287R01}{%
= \psi(M(1;M(1;1,0)+2))%
}

参数内部再次出现多元M时，薛峰按从外到内的次序逐层标记，才没有把内外两个后继错接。

\NovelDisplay{M05288R01}{%
(0)(1,1,1)(2,1,1)(3,1,1)(3,1,1)(2,1,1)(3,1,1)(3,1,1)(2,1,1)(3,1)(4,2,1)%
}
\NovelDisplay{M05288R02}{%
(5,2,1)(6,2,1)(6,2,1)(5,2,1)(6,2,1)(6,2,1)%
}

\NovelDisplay{M05289R01}{%
= \psi(M(1;M(1;1,0)+\omega))%
}

\NovelDisplay{M05290R01}{%
(0)(1,1,1)(2,1,1)(3,1,1)(3,1,1)(2,1,1)(3,1,1)(3,1,1)(2,1,1)(3,1)(4,2,1)%
}
\NovelDisplay{M05290R02}{%
(5,2,1)(6,2,1)(6,2,1)(5,2,1)(6,2,1)(6,2,1)(5,2,1)(6,2)%
}

\NovelDisplay{M05291R01}{%
= \psi(M(1;\Omega_{M(1;1,0)+1}))%
}

\NovelDisplay{M05292R01}{%
(0)(1,1,1)(2,1,1)(3,1,1)(3,1,1)(2,1,1)(3,1,1)(3,1,1)(2,1,1)(3,1)(4,2,1)%
}
\NovelDisplay{M05292R02}{%
(5,2,1)(6,2,1)(6,2,1)(5,2,1)(6,2,1)(6,2,1)(5,2,1)(6,2)(5)%
}

\NovelDisplay{M05293R01}{%
= \psi(M(1;1,1))%
}

\NovelDisplay{M05294R01}{%
= \psi(M(2;0)\times 2)%
}

\NovelDisplay{M05295R01}{%
(0)(1,1,1)(2,1,1)(3,1,1)(3,1,1)(2,1,1)(3,1,1)(3,1,1)(2,1,1)(3,1,1)%
}

\NovelDisplay{M05296R01}{%
= \psi(M(1;1,\omega))%
}

\NovelDisplay{M05297R01}{%
= \psi(M(2;0)\times \omega)%
}

从参数中的ω与Ω走到外层的倍数，读式仍要沿着分号两侧逐一确认，而不是把最后一段当成独立小式。

\NovelDisplay{M05298R01}{%
(0)(1,1,1)(2,1,1)(3,1,1)(3,1,1)(2,1,1)(3,1,1)(3,1,1)(2,1,1)(3,1,1)%
}
\NovelDisplay{M05298R02}{%
(2,1,1)(3,1)(2)%
}

\NovelDisplay{M05299R01}{%
= \psi(M(1;2,0))%
}

\NovelDisplay{M05300R01}{%
= \psi({M(2;0)}^2)%
}

\NovelDisplay{M05301R01}{%
(0)(1,1,1)(2,1,1)(3,1,1)(3,1,1)(2,1,1)(3,1,1)(3,1,1)(2,1,1)(3,1,1)%
}
\NovelDisplay{M05301R02}{%
(3,1)(2)%
}

\NovelDisplay{M05302R01}{%
= \psi(M(1;1,0,0))%
}

\NovelDisplay{M05303R01}{%
= \psi({M(2;0)}^{M(2;0)})%
}

\NovelDisplay{M05304R01}{%
(0)(1,1,1)(2,1,1)(3,1,1)(3,1,1)(2,1,1)(3,1,1)(3,1,1)(2,1,1)(3,1,1)%
}
\NovelDisplay{M05304R02}{%
(3,1)(4,2)%
}

\NovelDisplay{M05305R01}{%
= \psi(\Omega_{M(2;0)+1})%
}

\NovelDisplay{M05306R01}{%
(0)(1,1,1)(2,1,1)(3,1,1)(3,1,1)(2,1,1)(3,1,1)(3,1,1)(2,1,1)(3,1,1)(3,1)%
}
\NovelDisplay{M05306R02}{%
(4,2,1)(5,2,1)(6,2,1)(6,2,1)(5,2,1)(6,2)(5)%
}

\NovelDisplay{M05307R01}{%
= \psi(M(1,M(2;0)+1))%
}

\NovelDisplay{M05308R01}{%
= \psi(M(1;M(2;0)+1))%
}

\NovelDisplay{M05309R01}{%
= \psi(N^3+N^2)%
}

新的平方、自幂及加项共同围绕\allowbreak{}\NovelInline{I06119}{M(2;0)}\allowbreak{}展开，蔡须坤把共同背景留全，才让后面的\allowbreak{}\NovelInline{I06120}{N^3+N^2}\allowbreak{}有了对照位置。

\NovelDisplay{M05310R01}{%
(0)(1,1,1)(2,1,1)(3,1,1)(3,1,1)(2,1,1)(3,1,1)(3,1,1)(2,1,1)(3,1,1)(3,1)%
}
\NovelDisplay{M05310R02}{%
(4,2,1)(5,2,1)(6,2,1)(6,2,1)(5,2,1)(6,2,1)(6,2)(7,3)%
}

\NovelDisplay{M05311R01}{%
= \psi(\Omega_{M(1;M(2;0)+1)+1})%
}

\NovelDisplay{M05312R01}{%
(0)(1,1,1)(2,1,1)(3,1,1)(3,1,1)(2,1,1)(3,1,1)(3,1,1)(2,1,1)(3,1,1)(3,1)%
}
\NovelDisplay{M05312R02}{%
(4,2,1)(5,2,1)(6,2,1)(6,2,1)(5,2,1)(6,2,1)(6,2,1)%
}

\NovelDisplay{M05313R01}{%
= \psi(M(1;M(2;0)+\omega))%
}

\NovelDisplay{M05314R01}{%
= \psi(N^3+N^2\times \omega)%
}

\NovelDisplay{M05315R01}{%
(0)(1,1,1)(2,1,1)(3,1,1)(3,1,1)(2,1,1)(3,1,1)(3,1,1)(2,1,1)(3,1,1)(3,1)%
}
\NovelDisplay{M05315R02}{%
(4,2,1)(5,2,1)(6,2,1)(6,2,1)(5,2,1)(6,2,1)(6,2,1)(5,2,1)(6,2)(5)%
}

\NovelDisplay{M05316R01}{%
= \psi(M(1;1,M(2;0)+1))%
}

\NovelDisplay{M05317R01}{%
(0)(1,1,1)(2,1,1)(3,1,1)(3,1,1)(2,1,1)(3,1,1)(3,1,1)(2,1,1)(3,1,1)(3,1)%
}
\NovelDisplay{M05317R02}{%
(4,2,1)(5,2,1)(6,2,1)(6,2,1)(5,2,1)(6,2,1)(6,2,1)(5,2,1)(6,2,1)%
}
\NovelDisplay{M05317R03}{%
(6,2)(7,3)%
}

\NovelDisplay{M05318R01}{%
= \psi(M(2;1))%
}

\NovelDisplay{M05319R01}{%
= \psi(N^3\times 2)%
}

对照进入\allowbreak{}\NovelInline{I06121}{N^3}\allowbreak{}的下一组写法时，诺当先说明参数嵌套在哪里，再将外层倍数单独指出。

\NovelDisplay{M05320R01}{%
(0)(1,1,1)(2,1,1)(3,1,1)(3,1,1)(2,1,1)(3,1,1)(3,1,1)(2,1,1)%
}
\NovelDisplay{M05320R02}{%
(3,1,1)(3,1,1)%
}

\NovelDisplay{M05321R01}{%
= \psi(M(2;\omega))%
}

\NovelDisplay{M05322R01}{%
= \psi(N^3\times \omega)%
}

\NovelDisplay{M05323R01}{%
(0)(1,1,1)(2,1,1)(3,1,1)(3,1,1)(2,1,1)(3,1,1)(3,1,1)(2,1,1)(3,1,1)%
}
\NovelDisplay{M05323R02}{%
(3,1,1)(2,1,1)(3,1)(2)%
}

\NovelDisplay{M05324R01}{%
= \psi(M(2;1,0))%
}

\NovelDisplay{M05325R01}{%
= \psi(M(3;0))%
}

\NovelDisplay{M05326R01}{%
= \psi(N^4)%
}

\NovelDisplay{M05327R01}{%
(0)(1,1,1)(2,1,1)(3,1,1)(3,1,1)(2,1,1)(3,1,1)(3,1,1)(2,1,1)(3,1,1)%
}
\NovelDisplay{M05327R02}{%
(3,1,1)(2,1,1)(3,1)(4,2,1)(5,2,1)(6,2,1)(6,2,1)(5,2,1)(6,2)(5)%
}

\NovelDisplay{M05328R01}{%
= \psi(M(1,M(2;1,0)+1))%
}

\NovelDisplay{M05329R01}{%
= \psi(M(1;M(2;1,0)+1))%
}

\NovelDisplay{M05330R01}{%
(0)(1,1,1)(2,1,1)(3,1,1)(3,1,1)(2,1,1)(3,1,1)(3,1,1)(2,1,1)(3,1,1)%
}
\NovelDisplay{M05330R02}{%
(3,1,1)(2,1,1)(3,1)(4,2,1)(5,2,1)(6,2,1)(6,2,1)(5,2,1)(6,2,1)(6,2)(7,3)%
}

\NovelDisplay{M05331R01}{%
= \psi(\Omega_{M(1;M(2;1,0)+1)+1})%
}

从\allowbreak{}\NovelInline{I06122}{N^3}\allowbreak{}相关段接到\allowbreak{}\NovelInline{I06123}{N^4}\allowbreak{}，每条对应仍需回到原矩阵，不能把指数的四当成只需多写一个字。

风封住高处的路以后，他们在山脚留了数日。

屋里比先前挤。湿衣裳挂着，锅里的补给慢慢冒热气，黑子鸡在门边找出一点干燥的地方。蔡兴坤怕它啄纸，拿半张废稿围在旁边，才回到桌前。

薛峰请他先说今天这一组，蔡兴坤一开口便错了两个分组。若放在从前，他早要找一句玩笑岔开，这回却自己伸手，将原式重新摆正。

``别翻下一页。''他说，``我还没说清。''

奥斯卡端着碗走过来，便也没有催。他们那晚吃饭晚了些，最末一处问题却终于由蔡兴坤自己指出。

\NovelDisplay{M05332R01}{%
(0)(1,1,1)(2,1,1)(3,1,1)(3,1,1)(2,1,1)(3,1,1)(3,1,1)(2,1,1)(3,1,1)%
}
\NovelDisplay{M05332R02}{%
(3,1,1)(2,1,1)(3,1)(4,2,1)(5,2,1)(6,2,1)(6,2,1)(5,2,1)(6,2,1)(6,2,1)%
}

\NovelDisplay{M05333R01}{%
= \psi(M(1;M(2;1,0)+\omega))%
}

\NovelDisplay{M05334R01}{%
(0)(1,1,1)(2,1,1)(3,1,1)(3,1,1)(2,1,1)(3,1,1)(3,1,1)(2,1,1)(3,1,1)%
}
\NovelDisplay{M05334R02}{%
(3,1,1)(2,1,1)(3,1)(4,2,1)(5,2,1)(6,2,1)(6,2,1)(5,2,1)(6,2,1)(6,2,1)%
}
\NovelDisplay{M05334R03}{%
(5,2,1)(6,2)(5)%
}

\NovelDisplay{M05335R01}{%
= \psi(M(1;1,M(2;1,0)+1))%
}

\NovelDisplay{M05336R01}{%
= \psi(M(2;M(2;1,0)+1))%
}

\NovelDisplay{M05337R01}{%
(0)(1,1,1)(2,1,1)(3,1,1)(3,1,1)(2,1,1)(3,1,1)(3,1,1)(2,1,1)(3,1,1)%
}
\NovelDisplay{M05337R02}{%
(3,1,1)(2,1,1)(3,1)(4,2,1)(5,2,1)(6,2,1)(6,2,1)(5,2,1)(6,2,1)(6,2,1)%
}
\NovelDisplay{M05337R03}{%
(5,2,1)(6,2,1)(6,2)(7,3)%
}

\NovelDisplay{M05338R01}{%
= \psi(\Omega_{M(2;M(2;1,0)+1)+1})%
}

\NovelDisplay{M05339R01}{%
(0)(1,1,1)(2,1,1)(3,1,1)(3,1,1)(2,1,1)(3,1,1)(3,1,1)(2,1,1)(3,1,1)%
}
\NovelDisplay{M05339R02}{%
(3,1,1)(2,1,1)(3,1)(4,2,1)(5,2,1)(6,2,1)(6,2,1)(5,2,1)(6,2,1)(6,2,1)%
}
\NovelDisplay{M05339R03}{%
(5,2,1)(6,2,1)(6,2,1)%
}

\NovelDisplay{M05340R01}{%
= \psi(M(2;M(2;1,0)+\omega))%
}

多元参数把整个旧段包入新位置时，叶阳先读清参数边界，才接出新背景下的后继与倍数。

风没有停，霍雨浩便将这条矩阵留在桌上，让唐三、小舞与叶阳各求一项，再由薛峰核对他们是否从同一个原式出发。

\NovelDisplay{M05341R01}{%
\text{展开第零项：}\quad(0)(1,1,1)(2,1,1)(3,1,1)(3,1,1)(2,1,1)(3,1,1)(3,1,1)%
}
\NovelDisplay{M05341R02}{%
(2,1,1)(3,1,1)(3,1,1)(2,1,1)(3,1)(4,2,1)(5,2,1)(6,2,1)(6,2,1)(5,2,1)%
}
\NovelDisplay{M05341R03}{%
(6,2,1)(6,2,1)(5,2,1)(6,2)%
}

\NovelDisplay{M05342R01}{%
\text{展开第一项：}\quad(0)(1,1,1)(2,1,1)(3,1,1)(3,1,1)(2,1,1)(3,1,1)(3,1,1)%
}
\NovelDisplay{M05342R02}{%
(2,1,1)(3,1,1)(3,1,1)(2,1,1)(3,1)(4,2,1)(5,2,1)(6,2,1)(6,2,1)(5,2,1)%
}
\NovelDisplay{M05342R03}{%
(6,2,1)(6,2,1)(5,2,1)(6,2)(4,2,1)(5,2,1)(6,2,1)(6,2,1)(5,2,1)(6,2,1)%
}
\NovelDisplay{M05342R04}{%
(6,2,1)(5,2,1)(6,2)%
}

\NovelDisplay{M05343R01}{%
\text{展开第二项：}\quad(0)(1,1,1)(2,1,1)(3,1,1)(3,1,1)(2,1,1)(3,1,1)(3,1,1)%
}
\NovelDisplay{M05343R02}{%
(2,1,1)(3,1,1)(3,1,1)(2,1,1)(3,1)(4,2,1)(5,2,1)(6,2,1)(6,2,1)(5,2,1)%
}
\NovelDisplay{M05343R03}{%
(6,2,1)(6,2,1)(5,2,1)(6,2)(4,2,1)(5,2,1)(6,2,1)(6,2,1)(5,2,1)(6,2,1)%
}
\NovelDisplay{M05343R04}{%
(6,2,1)(5,2,1)(6,2)(4,2,1)(5,2,1)(6,2,1)(6,2,1)(5,2,1)(6,2,1)%
}
\NovelDisplay{M05343R05}{%
(6,2,1)(5,2,1)(6,2)%
}

外面能否登阶暂时受风阻拦，矩阵的独立取项仍须按原规则走完，不能借天气将尚未核清的列略过去。

小舞指着第二项最末的重复段，请叶阳再看一遍自己的记录，确认以后才伸手去端已经放凉的碗。

\NovelDisplay{M05344R01}{%
(0)(1,1,1)(2,1,1)(3,1,1)(3,1,1)(2,1,1)(3,1,1)(3,1,1)(2,1,1)(3,1,1)%
}
\NovelDisplay{M05344R02}{%
(3,1,1)(2,1,1)(3,1)(4,2,1)(5,2,1)(6,2,1)(6,2,1)(5,2,1)(6,2,1)(6,2,1)%
}
\NovelDisplay{M05344R03}{%
(5,2,1)(6,2,1)(6,2,1)(5,2,1)(6,2)(5)%
}

\NovelDisplay{M05345R01}{%
= \psi(M(2;1,1))%
}

\NovelDisplay{M05346R01}{%
= \psi(M(3;0)\times 2)%
}

\NovelDisplay{M05347R01}{%
= \psi(N^3+\psi_{N}(N^4))%
}

\NovelDisplay{M05348R01}{%
(0)(1,1,1)(2,1,1)(3,1,1)(3,1,1)(2,1,1)(3,1,1)(3,1,1)(2,1,1)(3,1,1)%
}
\NovelDisplay{M05348R02}{%
(3,1,1)(2,1,1)(3,1,1)%
}

\NovelDisplay{M05349R01}{%
= \psi(M(2;1,\omega))%
}

\NovelDisplay{M05350R01}{%
= \psi(M(3;0)\times \omega)%
}

\NovelDisplay{M05351R01}{%
(0)(1,1,1)(2,1,1)(3,1,1)(3,1,1)(2,1,1)(3,1,1)(3,1,1)(2,1,1)(3,1,1)%
}
\NovelDisplay{M05351R02}{%
(3,1,1)(2,1,1)(3,1,1)(3,1)(4,2)%
}

\NovelDisplay{M05352R01}{%
= \psi(\Omega_{M(3;0)+1})%
}

\NovelDisplay{M05353R01}{%
(0)(1,1,1)(2,1,1)(3,1,1)(3,1,1)(2,1,1)(3,1,1)(3,1,1)(2,1,1)(3,1,1)%
}
\NovelDisplay{M05353R02}{%
(3,1,1)(2,1,1)(3,1,1)(3,1,1)%
}

\NovelDisplay{M05354R01}{%
= \psi(M(3;\omega))%
}

\NovelDisplay{M05355R01}{%
= \psi(N^4\times \omega)%
}

到这里外层倍数与折叠项并列，大家保留讲义的完整右式，并用独立取项核对原矩阵自身的接续。

\NovelDisplay{M05356R01}{%
(0)(1,1,1)(2,1,1)(3,1,1)(3,1,1)(3)%
}

\NovelDisplay{M05357R01}{%
= \psi(M(\omega;0))%
}

\NovelDisplay{M05358R01}{%
= \psi(N^\omega)%
}

\NovelDisplay{M05359R01}{%
(0)(1,1,1)(2,1,1)(3,1,1)(3,1,1)(3)(2,1,1)%
}

\NovelDisplay{M05360R01}{%
= \psi(M(\omega;\omega))%
}

\NovelDisplay{M05361R01}{%
= \psi(N^\omega\times \omega)%
}

\NovelDisplay{M05362R01}{%
(0)(1,1,1)(2,1,1)(3,1,1)(3,1,1)(3)(2,1,1)(3,1)(2)%
}

\NovelDisplay{M05363R01}{%
= \psi(M(\omega;1,0))%
}

\NovelDisplay{M05364R01}{%
= \psi(M(\omega+1;0))%
}

\NovelDisplay{M05365R01}{%
= \psi(N^{\omega+1})%
}

\NovelDisplay{M05366R01}{%
(0)(1,1,1)(2,1,1)(3,1,1)(3,1,1)(3)(2,1,1)(3,1,1)%
}

\NovelDisplay{M05367R01}{%
= \psi(M(\omega;1,\omega))%
}

\NovelDisplay{M05368R01}{%
= \psi(M(\omega+1;0)\times \omega)%
}

\allowbreak{}\NovelInline{I06124}{M(\omega;0)}\allowbreak{}给出\allowbreak{}\NovelInline{I06125}{N^\omega}\allowbreak{}的对照，随后出现的分号后参数必须继续读全，不能只留下共同的幂名。

\NovelDisplay{M05369R01}{%
(0)(1,1,1)(2,1,1)(3,1,1)(3,1,1)(3)(2,1,1)(3,1,1)(3,1)(4,2)%
}

\NovelDisplay{M05370R01}{%
= \psi(\Omega_{M(\omega+1;0)})%
}

\NovelDisplay{M05371R01}{%
(0)(1,1,1)(2,1,1)(3,1,1)(3,1,1)(3)(2,1,1)(3,1,1)(3,1,1)%
}

\NovelDisplay{M05372R01}{%
= \psi(M(\omega+1;\omega))%
}

\NovelDisplay{M05373R01}{%
= \psi(N^{\omega+1}\times \omega)%
}

\NovelDisplay{M05374R01}{%
(0)(1,1,1)(2,1,1)(3,1,1)(3,1,1)(3)(2,1,1)(3,1,1)(3,1,1)(3)%
}

\NovelDisplay{M05375R01}{%
= \psi(M(\omega\times 2;0))%
}

\NovelDisplay{M05376R01}{%
= \psi(N^{\omega\times 2})%
}

\NovelDisplay{M05377R01}{%
(0)(1,1,1)(2,1,1)(3,1,1)(3,1,1)(3,1)%
}

\NovelDisplay{M05378R01}{%
= \psi(M(\Omega;0))%
}

\NovelDisplay{M05379R01}{%
= \psi(N^{\Omega})%
}

指数后继、指数倍数与Ω进入参数的情形在此相邻，卢本韦逐个指出改变的位置，避免用一条口令代替三种核对。

\NovelDisplay{M05380R01}{%
(0)(1,1,1)(2,1,1)(3,1,1)(3,1,1)(3,1)(2)%
}

\NovelDisplay{M05381R01}{%
= \psi(M(1,0;0))%
}

\NovelDisplay{M05382R01}{%
= \psi(M(1,1;0))%
}

\NovelDisplay{M05383R01}{%
= \psi(N^N)%
}

\NovelDisplay{M05384R01}{%
(0)(1,1,1)(2,1,1)(3,1,1)(3,1,1)(3,1)(2,1,1)%
}

\NovelDisplay{M05385R01}{%
= \psi(M(1,0;0)\times \omega)%
}

\NovelDisplay{M05386R01}{%
(0)(1,1,1)(2,1,1)(3,1,1)(3,1,1)(3,1)(2,1,1)(3,1)(4,2)%
}

\NovelDisplay{M05387R01}{%
= \psi(\Omega_{M(1,0;0)+1})%
}

\NovelDisplay{M05388R01}{%
(0)(1,1,1)(2,1,1)(3,1,1)(3,1,1)(3,1)(2,1,1)(3,1,1)%
}

\NovelDisplay{M05389R01}{%
= \psi(M(1,0;\omega))%
}

不同多元M写法在ψ下的对应照原材料保留，蔡须坤没有据此外推它们在任何上下文都可以任意互换。

古月娜有一回把旧封套翻了出来。

里面那张理想不提升BMS的纸，边角已经压得很平。她看了半天，叫蔡须坤过去，说自己仍记得当时那些看起来应当成立的对应。

蔡须坤坐下来，没有伸手把纸收走。

``我当时以为，既然能说出想要的右式，后面总能找到路。''古月娜道。

她将纸翻过来。背后原来准备写展开规则的位置，仍然没有完整的一行。

``现在知道，空着的就在这里。''

蔡须坤点头，把正在核的原版矩阵放到旁边。两张纸没有被混到一起。古月娜先将封套合好，再问她今天最末那一列准备怎样接。

\NovelDisplay{M05390R01}{%
(0)(1,1,1)(2,1,1)(3,1,1)(3,1,1)(3,1)(2,1,1)(3,1,1)(2,1,1)(3,1)(2)%
}

\NovelDisplay{M05391R01}{%
= \psi(M(1,0;1,0))%
}

\NovelDisplay{M05392R01}{%
= \psi({M(1,1;0)}^2)%
}

\NovelDisplay{M05393R01}{%
(0)(1,1,1)(2,1,1)(3,1,1)(3,1,1)(3,1)(2,1,1)(3,1,1)(3,1)(4,2)%
}

\NovelDisplay{M05394R01}{%
= \psi(\Omega_{M(1,1;0)+1})%
}

\NovelDisplay{M05395R01}{%
(0)(1,1,1)(2,1,1)(3,1,1)(3,1,1)(3,1)(2,1,1)(3,1,1)(3,1)(4,2,1)(5,2,1)%
}
\NovelDisplay{M05395R02}{%
(6,2,1)(6,2,1)%
}

\NovelDisplay{M05396R01}{%
= \psi(M_{M(1,1;0)+1})%
}

\NovelDisplay{M05397R01}{%
(0)(1,1,1)(2,1,1)(3,1,1)(3,1,1)(3,1)(2,1,1)(3,1,1)(3,1)(4,2,1)(5,2,1)%
}
\NovelDisplay{M05397R02}{%
(6,2,1)(6,2,1)(6)%
}

\NovelDisplay{M05398R01}{%
= \psi(M(\omega;M(1,1;0)))%
}

新的参数对象回接自身时，霍雨浩让唐三先说明每一个下标属于哪个ψ，再把对应完整念给小舞听。

\NovelDisplay{M05399R01}{%
(0)(1,1,1)(2,1,1)(3,1,1)(3,1,1)(3,1)(2,1,1)(3,1,1)(3,1)(4,2,1)(5,2,1)%
}
\NovelDisplay{M05399R02}{%
(6,2,1)(6,2,1)(6,2)(5)%
}

\NovelDisplay{M05400R01}{%
= \psi(M(1,0;M(1,1;0)+1))%
}

\NovelDisplay{M05401R01}{%
= \psi(M(1,1;1))%
}

\NovelDisplay{M05402R01}{%
= \psi(N^N\times 2)%
}

\NovelDisplay{M05403R01}{%
(0)(1,1,1)(2,1,1)(3,1,1)(3,1,1)(3,1)(2,1,1)(3,1,1)(3,1)(4,2,1)(5,2,1)%
}
\NovelDisplay{M05403R02}{%
(6,2,1)(6,2,1)(6,2)(5,2,1)(6,2,1)%
}

\NovelDisplay{M05404R01}{%
= \psi(M(1,0;M(1,1;0)+\omega))%
}

\NovelDisplay{M05405R01}{%
(0)(1,1,1)(2,1,1)(3,1,1)(3,1,1)(3,1)(2,1,1)(3,1,1)(3,1)(4,2,1)(5,2,1)%
}
\NovelDisplay{M05405R02}{%
(6,2,1)(6,2,1)(6,2)(5,2,1)(6,2,1)(6,2)(7,3)%
}

\NovelDisplay{M05406R01}{%
= \psi(\Omega_{M(1,1;1)+1})%
}

\NovelDisplay{M05407R01}{%
(0)(1,1,1)(2,1,1)(3,1,1)(3,1,1)(3,1)(2,1,1)(3,1,1)(3,1,1)%
}

\NovelDisplay{M05408R01}{%
= \psi(M(1,1;\omega))%
}

\NovelDisplay{M05409R01}{%
= \psi(N^N\times \omega)%
}

这里从\allowbreak{}\NovelInline{I06126}{N^N}\allowbreak{}的倍数继续向后，分号后面仍可能包含完整旧项，不能把它压成一个没有说明的省略号。

\NovelDisplay{M05410R01}{%
(0)(1,1,1)(2,1,1)(3,1,1)(3,1,1)(3,1)(2,1,1)(3,1,1)(3,1,1)(2,1,1)(3,1)(2)%
}

\NovelDisplay{M05411R01}{%
= \psi(M(1,1;1,0))%
}

\NovelDisplay{M05412R01}{%
= \psi(M(1,2;0))%
}

\NovelDisplay{M05413R01}{%
= \psi(N^{N+1})%
}

\NovelDisplay{M05414R01}{%
(0)(1,1,1)(2,1,1)(3,1,1)(3,1,1)(3,1)(2,1,1)(3,1,1)(3,1,1)(2,1,1)%
}
\NovelDisplay{M05414R02}{%
(3,1,1)(3,1)(4,2)%
}

\NovelDisplay{M05415R01}{%
= \psi(\Omega_{M(1,2;0)+1})%
}

\NovelDisplay{M05416R01}{%
(0)(1,1,1)(2,1,1)(3,1,1)(3,1,1)(3,1)(2,1,1)(3,1,1)(3,1,1)(2,1,1)%
}
\NovelDisplay{M05416R02}{%
(3,1,1)(3,1,1)%
}

\NovelDisplay{M05417R01}{%
= \psi(M(1,2;\omega))%
}

\NovelDisplay{M05418R01}{%
= \psi(N^{N+1}\times \omega)%
}

\NovelDisplay{M05419R01}{%
(0)(1,1,1)(2,1,1)(3,1,1)(3,1,1)(3,1)(2,1,1)(3,1,1)(3,1,1)(2,1,1)(3,1,1)%
}
\NovelDisplay{M05419R02}{%
(3,1,1)(2,1,1)(3,1,1)(3,1,1)%
}

\NovelDisplay{M05420R01}{%
= \psi(M(1,3;\omega))%
}

\NovelDisplay{M05421R01}{%
= \psi(N^{N+2}\times \omega)%
}

\allowbreak{}\NovelInline{I06127}{N^{N+1}}\allowbreak{}及其ω倍数的表示逐一落定，薛峰终于能在队友开口以前指出哪一处是指数、哪一处是乘数。

\NovelDisplay{M05422R01}{%
(0)(1,1,1)(2,1,1)(3,1,1)(3,1,1)(3,1)(2,1,1)(3,1,1)(3,1,1)(3,1)%
}

\NovelDisplay{M05423R01}{%
= \psi(M(1,\Omega;0))%
}

\NovelDisplay{M05424R01}{%
= \psi(N^{N+\Omega})%
}

\NovelDisplay{M05425R01}{%
(0)(1,1,1)(2,1,1)(3,1,1)(3,1,1)(3,1)(2,1,1)(3,1,1)(3,1,1)(3,1)(2)%
}

\NovelDisplay{M05426R01}{%
= \psi(M(2,0;0))%
}

\NovelDisplay{M05427R01}{%
= \psi(M(2,1;0))%
}

\NovelDisplay{M05428R01}{%
= \psi(N^{N\times 2})%
}

\NovelDisplay{M05429R01}{%
(0)(1,1,1)(2,1,1)(3,1,1)(3,1,1)(3,1)(2,1,1)(3,1,1)(3,1,1)(3,1)%
}
\NovelDisplay{M05429R02}{%
(2,1,1)(3,1,1)%
}

\NovelDisplay{M05430R01}{%
= \psi(M(2,0;\omega))%
}

\NovelDisplay{M05431R01}{%
(0)(1,1,1)(2,1,1)(3,1,1)(3,1,1)(3,1)(2,1,1)(3,1,1)(3,1,1)(3,1)(2,1,1)%
}
\NovelDisplay{M05431R02}{%
(3,1,1)(3,1,1)%
}

\NovelDisplay{M05432R01}{%
= \psi(M(2,1;\omega))%
}

\NovelDisplay{M05433R01}{%
= \psi(N^{N\times 2}\times \omega)%
}

参数首项转成倍数后，新的记号仍保留原分隔符，奥斯卡按原矩阵逐列复述，没再靠纸上的长度猜答案。

\NovelDisplay{M05434R01}{%
(0)(1,1,1)(2,1,1)(3,1,1)(3,1,1)(3,1)(2,1,1)(3,1,1)(3,1,1)(3,1)(2,1,1)%
}
\NovelDisplay{M05434R02}{%
(3,1,1)(3,1,1)(2,1,1)(3,1,1)(3,1,1)%
}

\NovelDisplay{M05435R01}{%
= \psi(M(2,2;\omega))%
}

\NovelDisplay{M05436R01}{%
= \psi(N^{N\times 2+1}\times \omega)%
}

\NovelDisplay{M05437R01}{%
(0)(1,1,1)(2,1,1)(3,1,1)(3,1,1)(3,1)(2,1,1)(3,1,1)(3,1,1)(3,1)(2,1,1)%
}
\NovelDisplay{M05437R02}{%
(3,1,1)(3,1,1)(3,1)(2)%
}

\NovelDisplay{M05438R01}{%
= \psi(M(3,0;0))%
}

\NovelDisplay{M05439R01}{%
= \psi(M(3,1;0))%
}

\NovelDisplay{M05440R01}{%
= \psi(N^{N\times 3})%
}

\NovelDisplay{M05441R01}{%
(0)(1,1,1)(2,1,1)(3,1,1)(3,1,1)(3,1)(3)%
}

\NovelDisplay{M05442R01}{%
= \psi(M(\omega,0;0))%
}

\NovelDisplay{M05443R01}{%
= \psi(N^{N\times \omega})%
}

\NovelDisplay{M05444R01}{%
(0)(1,1,1)(2,1,1)(3,1,1)(3,1,1)(3,1)(3)(2,1,1)%
}

\NovelDisplay{M05445R01}{%
= \psi(M(\omega,0;\omega))%
}

\NovelDisplay{M05446R01}{%
= \psi(N^{N\times \omega}\times \omega)%
}

指数进入\allowbreak{}\NovelInline{I06128}{N}\allowbreak{}的不同倍数时，蔡兴坤先核对ω的位置，确认它没有从指数里滑到整个幂外面。

第九个月的一天，蔡须坤独自醒得很早。

她昨夜写到这里，面对两条相近的右式，怎么也不肯停。此刻灯已经灭了，床边却留着一点淡淡的魂光。牢大的魂影停在那里，不曾像活着时那样去拿风箱、去搬凳子，只等着她慢慢坐起。

``又不睡？''

``睡过。''她轻声说，自己也觉得那话太轻。

牢大望着她摊开的纸，问今天哪一处能自己接清。蔡须坤指了前面，又指后面，将确定与尚待核对的分别说出。

魂影听完，没有给她另一个答案。

``先把能走的站稳。''

她记得他从前说过许多次，到了这一晚，才没有再将这句话听成催自己练得更多。天亮以后，她带着那张仍有空处的纸去找霍雨浩，第一句便说，这里她还不知道。

\NovelDisplay{M05447R01}{%
(0)(1,1,1)(2,1,1)(3,1,1)(3,1,1)(3,1)(3)(2,1,1)(3,1,1)%
}

\NovelDisplay{M05448R01}{%
= \psi(M(\omega,0;1,\omega))%
}

\NovelDisplay{M05449R01}{%
(0)(1,1,1)(2,1,1)(3,1,1)(3,1,1)(3,1)(3)(2,1,1)(3,1,1)(3,1,1)%
}

\NovelDisplay{M05450R01}{%
= \psi(M(\omega,1;\omega))%
}

\NovelDisplay{M05451R01}{%
= \psi(N^{N\times \omega+1}\times \omega)%
}

\NovelDisplay{M05452R01}{%
(0)(1,1,1)(2,1,1)(3,1,1)(3,1,1)(3,1)(3)(2,1,1)(3,1,1)(3,1,1)(3,1)(2)%
}

\NovelDisplay{M05453R01}{%
= \psi(M(\omega+1,1;0))%
}

\NovelDisplay{M05454R01}{%
= \psi(N^{N\times \omega + N+1})%
}

\NovelDisplay{M05455R01}{%
(0)(1,1,1)(2,1,1)(3,1,1)(3,1,1)(3,1)(3)(2,1,1)(3,1,1)(3,1,1)(3,1)(3)%
}

\NovelDisplay{M05456R01}{%
= \psi(M(\omega\times 2, 0;0))%
}

\NovelDisplay{M05457R01}{%
= \psi(N^{N\times \omega + N+1})%
}

这两条相邻右式照讲义原样留存，队员们不据此追加未经核对的等式，脚下的接续仍由矩阵独立展开确认。

\NovelDisplay{M05458R01}{%
(0)(1,1,1)(2,1,1)(3,1,1)(3,1,1)(3,1)(3)(3)%
}

\NovelDisplay{M05459R01}{%
= \psi(M(\psi(2),0;0))%
}

\NovelDisplay{M05460R01}{%
= \psi(N^{N\times \psi(2)})%
}

\NovelDisplay{M05461R01}{%
(0)(1,1,1)(2,1,1)(3,1,1)(3,1,1)(3,1)(3,1)(2)%
}

\NovelDisplay{M05462R01}{%
= \psi(M(1,0,0;0))%
}

\NovelDisplay{M05463R01}{%
= \psi(M(1,0,1;0))%
}

\NovelDisplay{M05464R01}{%
= \psi(N^{N^2})%
}

\NovelDisplay{M05465R01}{%
(0)(1,1,1)(2,1,1)(3,1,1)(3,1,1)(3,1)(3,1)(2,1,1)(3,1,1)%
}

\NovelDisplay{M05466R01}{%
= \psi(M(1,0,0;\omega))%
}

\NovelDisplay{M05467R01}{%
(0)(1,1,1)(2,1,1)(3,1,1)(3,1,1)(3,1)(3,1)(2,1,1)(3,1,1)(3,1,1)%
}

\NovelDisplay{M05468R01}{%
= \psi(M(1,0,1;\omega))%
}

\NovelDisplay{M05469R01}{%
= \psi(N^{N^2}\times \omega)%
}

参数再增加一层以后，\allowbreak{}\NovelInline{I06129}{N^{N^2}}\allowbreak{}出现为新的对照，蔡须坤与卢本韦分别说明外框与内层分组。

\NovelDisplay{M05470R01}{%
(0)(1,1,1)(2,1,1)(3,1,1)(3,1,1)(3,1)(3,1)(2,1,1)(3,1,1)(3,1,1)(2,1,1)%
}
\NovelDisplay{M05470R02}{%
(3,1,1)(3,1,1)%
}

\NovelDisplay{M05471R01}{%
= \psi(M(1,0,2;\omega))%
}

\NovelDisplay{M05472R01}{%
= \psi(N^{N^2+1}\times \omega)%
}

\NovelDisplay{M05473R01}{%
(0)(1,1,1)(2,1,1)(3,1,1)(3,1,1)(3,1)(3,1)(2,1,1)(3,1,1)(3,1,1)(3,1)(2)%
}

\NovelDisplay{M05474R01}{%
= \psi(M(1,1,1;0))%
}

\NovelDisplay{M05475R01}{%
= \psi(N^{N^2+N})%
}

\NovelDisplay{M05476R01}{%
(0)(1,1,1)(2,1,1)(3,1,1)(3,1,1)(3,1)(3,1)(2,1,1)(3,1,1)(3,1,1)(3,1)%
}
\NovelDisplay{M05476R02}{%
(2,1,1)(3,1,1)(3,1,1)%
}

\NovelDisplay{M05477R01}{%
= \psi(M(1,1,1;\omega))%
}

\NovelDisplay{M05478R01}{%
= \psi(N^{N^2+N}\times \omega)%
}

\NovelDisplay{M05479R01}{%
(0)(1,1,1)(2,1,1)(3,1,1)(3,1,1)(3,1)(3,1)(2,1,1)(3,1,1)(3,1,1)(3,1)%
}
\NovelDisplay{M05479R02}{%
(2,1,1)(3,1,1)(3,1,1)(3,1)(2)%
}

\NovelDisplay{M05480R01}{%
= \psi(M(1,2,1;0))%
}

\NovelDisplay{M05481R01}{%
= \psi(N^{N^2+N\times 2})%
}

\allowbreak{}\NovelInline{I06130}{N^2}\allowbreak{}上继续接后继与\allowbreak{}\NovelInline{I06131}{N}\allowbreak{}的加项时，必须先确定上标的范围，不能把内层的加法念成整个幂的加法。

蔡须坤将原式推到卢本韦面前，三次取项分别落在不同的纸上，古月娜等三人写全以后，再拿共同保留的前段相核。

\NovelDisplay{M05482R01}{%
\text{展开第零项：}\quad(0)(1,1,1)(2,1,1)(3,1,1)(3,1,1)(3,1)(3,1)(2,1,1)(3,1,1)%
}
\NovelDisplay{M05482R02}{%
(3,1,1)(3,1)(2,1,1)(3,1,1)(3,1,1)(3,1)%
}

\NovelDisplay{M05483R01}{%
\text{展开第一项：}\quad(0)(1,1,1)(2,1,1)(3,1,1)(3,1,1)(3,1)(3,1)(2,1,1)(3,1,1)%
}
\NovelDisplay{M05483R02}{%
(3,1,1)(3,1)(2,1,1)(3,1,1)(3,1,1)(3,1)(1,1,1)(2,1,1)(3,1,1)(3,1,1)(3,1)%
}
\NovelDisplay{M05483R03}{%
(3,1)(2,1,1)(3,1,1)(3,1,1)(3,1)(2,1,1)(3,1,1)(3,1,1)(3,1)%
}

\NovelDisplay{M05484R01}{%
\text{展开第二项：}\quad(0)(1,1,1)(2,1,1)(3,1,1)(3,1,1)(3,1)(3,1)(2,1,1)(3,1,1)%
}
\NovelDisplay{M05484R02}{%
(3,1,1)(3,1)(2,1,1)(3,1,1)(3,1,1)(3,1)(1,1,1)(2,1,1)(3,1,1)(3,1,1)(3,1)%
}
\NovelDisplay{M05484R03}{%
(3,1)(2,1,1)(3,1,1)(3,1,1)(3,1)(2,1,1)(3,1,1)(3,1,1)(3,1)(1,1,1)(2,1,1)%
}
\NovelDisplay{M05484R04}{%
(3,1,1)(3,1,1)(3,1)(3,1)(2,1,1)(3,1,1)(3,1,1)(3,1)(2,1,1)(3,1,1)%
}
\NovelDisplay{M05484R05}{%
(3,1,1)(3,1)%
}

参数的对应可以留着继续追问，当前原矩阵的三项却要先独立核实，不能让一种记号的疑问遮住另一处可确认的步骤。

她将仍需核对的对应另外夹好，想起清晨那道魂影说过的话，便先把已经能独立说清的那页交给了同伴。

\NovelDisplay{M05485R01}{%
(0)(1,1,1)(2,1,1)(3,1,1)(3,1,1)(3,1)(3,1)(2,1,1)(3,1,1)(3,1,1)(3,1)%
}
\NovelDisplay{M05485R02}{%
(3,1)(2)%
}

\NovelDisplay{M05486R01}{%
= \psi(M(2,0,1;0))%
}

\NovelDisplay{M05487R01}{%
= \psi(N^{N^2 \times 2})%
}

\NovelDisplay{M05488R01}{%
(0)(1,1,1)(2,1,1)(3,1,1)(3,1,1)(3,1)(3,1)(3,1)(2)%
}

\NovelDisplay{M05489R01}{%
= \psi(M(1,0,0,1;0))%
}

\NovelDisplay{M05490R01}{%
= \psi(N^{N^3})%
}

\NovelDisplay{M05491R01}{%
(0)(1,1,1)(2,1,1)(3,1,1)(3,1,1)(3,1)(4)%
}

\NovelDisplay{M05492R01}{%
= \psi(M(1 @ \omega;0))%
}

\NovelDisplay{M05493R01}{%
= \psi(N^{N^\omega})%
}

从指数倍数走到更高的有限幂，再到ω指数，霍雨浩让每个人都指出共同保留的段，而非只背末尾的字母。

\NovelDisplay{M05494R01}{%
(0)(1,1,1)(2,1,1)(3,1,1)(3,1,1)(3,1)(4,1)(2)%
}

\NovelDisplay{M05495R01}{%
= \psi(N^{N^N})%
}

\NovelDisplay{M05496R01}{%
(0)(1,1,1)(2,1,1)(3,1,1)(3,1,1)(3,1)(4,1)(5,1)(2)%
}

\NovelDisplay{M05497R01}{%
= \psi(N^{N^{N^N}})%
}

\NovelDisplay{M05498R01}{%
\textcolor[HTML]{FF0000}{(0)(1,1,1)(2,1,1)(3,1,1)(3,1,1)(3,1)(4,2)}%
}

\NovelDisplay{M05499R01}{%
\textcolor[HTML]{FF0000}{= \psi(\Omega_{N+1})}%
}

\NovelDisplay{M05500R01}{%
\textcolor[HTML]{FF0000}{= \psi(\Pi_2\text{ aft }\Pi_2-\Pi_2-\Pi_2)}%
}

长幂接到最后的反射模式表示时，大家照讲义把两种写法都留下，没有额外宣称这就给出了未列中段的全部等式。

第十个月，最菜萌新在山脚来看过他们一次。

他没把山上的压力抹平，也没有替谁读过一份尚未完成的考卷。他只看他们把不同的册子搬上桌，听蔡须坤说明多元M里哪些地方已经能独立取项，再听卢本韦指出自己旧I册能接回哪一段。

翻到附录，蔡兴坤忽然问，那些二和连字符，能不能先算成零。

最菜萌新看了他一眼。奥斯卡先把自己差点也问出口的话咽了回去。

``读的是什么记号，先分清。''最菜萌新道，``就像你们不能把矩阵里的逗号，当成账本上随手隔开的空白。''

薛峰取来新纸，将正文与附录分开放，给每处参数留出能够回查的位置。

第十个月结束时，他们站上了四百七十米。

风已经小了，石阶边却留着一层薄霜。唐三替小舞将袖口的绳结松开，奥斯卡蹲下，重新系好诺当鞋底磨松的那一段。蔡兴坤把黑子鸡护在包里，等叶阳站稳，才从后面跟上。

蔡须坤看见薛峰先在石桩旁坐下，将今天最难的一个问题留给后来的人看。他没有写一份人人照抄便能走的答案，只将自己核过的起点与仍需对方独立说明的位置分开。

卢本韦最后收好附录。经过这些日子，蔡须坤再看那册加厚的I页，已经很少先想到谁的记号更难。她知道卢本韦守住了什么，也知道自己哪一处仍要请她慢下来再问一遍。

古月娜把封套放回包底。空着的理想定义没有被当成已完成的功法，山图已记下四百七十米。

离第一考要求的五百米，还差三十米。

小舞仰头看了一会儿，忽然问，到了以后第一件事要做什么。

奥斯卡说先吃饭，蔡兴坤说黑子鸡也要吃。大家笑过，霍雨浩将最后那段山图展开，薛峰从旁压住纸角。

蔡须坤没有先答。她伸手摸过石桩，等自己的呼吸完全平下来，才望向更高的那片雾。

\chapter{第一百二十九章：一年之后，第一考完成}\label{ux7b2cux4e00ux767eux4e8cux5341ux4e5dux7ae0ux4e00ux5e74ux4e4bux540eux7b2cux4e00ux8003ux5b8cux6210}

第十一个月，他们再次从四百八十米向上。

从N记号那一段走到新的K背景，中间还有许多没有在这两份讲义里连续列出的路。他们用了许多天补课、求项、互问，在各自的本子上留下记录，才接到现在这一页。霍雨浩没有将两份材料的首尾直接画等号，薛峰也没有把中间那段当成只需换一个字母。

小舞头一回没有先问还剩多少米。她把早上独立核过的那张纸拿给唐三，请他听自己从原矩阵开始说。唐三抱着包，在岩影里坐下，等她说到真正不明白的那一处，才抬手指向纸面。

卢本韦在I册后面又添了补页。她知道自己熟悉的几种表示各能接到哪里，眼前的新模式便照新模式读，没有为了保住一个名字，替它编造并不存在的对应。

蔡须坤走在她旁边，手里的M册已经很旧。她翻过封面，先摸到的却是最初那张写不出展开规则的纸。

她把纸收好，拿起了今天真正要走的原式。

到了这里，压力不像一开始那样只往肩上压。

蔡须坤有时能轻松走过几级，下一刻却在相同的位置骤然停住。霍雨浩问她刚才是从哪一处起的，她才发现，自己又因为认出一个熟悉的外框，提前沿用前一条的接法。

她退回去，重新看原矩阵。薛峰先核共同前段，卢本韦再将新背景读完整，等她自己说清，三人才把纸交还。

五百米的石桩仍在前面。她等呼吸平稳，抬脚时先想起的，也不再只是还差多少。

\NovelDisplay{M05501R01}{%
(0)(1,1,1)(2,1,1)(3,1,1)(4,1)(5,2)%
}

\NovelDisplay{M05502R01}{%
= \psi(\Omega_{K+1})%
}

\NovelDisplay{M05503R01}{%
= \psi(2\;\mathrm{aft}\;3)%
}

\NovelDisplay{M05504R01}{%
(0)(1,1,1)(2,1,1)(3,1,1)(4,1)(5,2)(1,1,1)(2,1,1)(3,1,1)(4,1)(5,2)%
}

\NovelDisplay{M05505R01}{%
= \psi(\Omega_{K+1}+\psi_{K}(\Omega_{K+1}))%
}

\NovelDisplay{M05506R01}{%
(0)(1,1,1)(2,1,1)(3,1,1)(4,1)(5,2)(2,1)(3,2)%
}

\NovelDisplay{M05507R01}{%
= \psi(\Omega_{K+1}+K)%
}

\NovelDisplay{M05508R01}{%
(0)(1,1,1)(2,1,1)(3,1,1)(4,1)(5,2)(2,1)(3,2,1)(4,2,1)(5,2,1)(5,2,1)%
}

\NovelDisplay{M05509R01}{%
= \psi(\Omega_{K+1}+K^2\times \omega)%
}

新的背景明确换成\allowbreak{}\NovelInline{I06132}{\Omega_{K+1}}\allowbreak{}，卢本韦把I册旧页夹好，在补页上照新记号读，不以换字母代替先前的接续。

\NovelDisplay{M05510R01}{%
(0)(1,1,1)(2,1,1)(3,1,1)(4,1)(5,2)(2,1)(3,2,1)(4,2,1)(5,2,1)(6)%
}

\NovelDisplay{M05511R01}{%
= \psi(\Omega_{K+1}+K^\omega)%
}

\NovelDisplay{M05512R01}{%
(0)(1,1,1)(2,1,1)(3,1,1)(4,1)(5,2)(2,1)(3,2,1)(4,2,1)(5,2,1)(6)(4,2,1)%
}

\NovelDisplay{M05513R01}{%
= \psi(\Omega_{K+1}+K^\omega\times \omega)%
}

\NovelDisplay{M05514R01}{%
(0)(1,1,1)(2,1,1)(3,1,1)(4,1)(5,2)(2,1)(3,2,1)(4,2,1)(5,2,1)(6)(5,2,1)%
}

\NovelDisplay{M05515R01}{%
= \psi(\Omega_{K+1}+K^{\omega+1}\times \omega)%
}

\NovelDisplay{M05516R01}{%
(0)(1,1,1)(2,1,1)(3,1,1)(4,1)(5,2)(2,1)(3,2,1)(4,2,1)(5,2,1)(6)%
}
\NovelDisplay{M05516R02}{%
(5,2,1)(6)%
}

\NovelDisplay{M05517R01}{%
= \psi(\Omega_{K+1}+K^{\omega\times 2})%
}

这里依次改变\allowbreak{}\NovelInline{I06133}{K}\allowbreak{}的幂与幂外倍数，唐三先指明ω究竟处在指数里还是指数外，才让小舞接着读。

\NovelDisplay{M05518R01}{%
(0)(1,1,1)(2,1,1)(3,1,1)(4,1)(5,2)(2,1)(3,2,1)(4,2,1)(5,2,1)(6,2)%
}

\NovelDisplay{M05519R01}{%
= \psi(\Omega_{K+1}+K^{\psi_{K}(\Omega_{K+1}+K)})%
}

\NovelDisplay{M05520R01}{%
(0)(1,1,1)(2,1,1)(3,1,1)(4,1)(5,2)(2,1)(3,2,1)(4,2,1)(5,2,1)(6,2)(4)%
}

\NovelDisplay{M05521R01}{%
= \psi(\Omega_{K+1}+K^{K})%
}

\NovelDisplay{M05522R01}{%
(0)(1,1,1)(2,1,1)(3,1,1)(4,1)(5,2)(2,1)(3,2,1)(4,2,1)(5,2,1)(6,2)(4,2,1)%
}

\NovelDisplay{M05523R01}{%
= \psi(\Omega_{K+1}+K^{K}\times \omega)%
}

\NovelDisplay{M05524R01}{%
(0)(1,1,1)(2,1,1)(3,1,1)(4,1)(5,2)(2,1)(3,2,1)(4,2,1)(5,2,1)(6,2)%
}
\NovelDisplay{M05524R02}{%
(4,2,1)(5,2,1)%
}

\NovelDisplay{M05525R01}{%
= \psi(1-2\;1-\{\psi_K(\Omega_{K+1}+K^{K}\times x)\})%
}

幂指数包含ψ后，又转入反射模式的写法，古月娜提醒大家按讲义的语法分别辨认，不将模式串当作普通算式。

\NovelDisplay{M05526R01}{%
(0)(1,1,1)(2,1,1)(3,1,1)(4,1)(5,2)(2,1)(3,2,1)(4,2,1)(5,2,1)(6,2)%
}
\NovelDisplay{M05526R02}{%
(5,2)(6,3)%
}

\NovelDisplay{M05527R01}{%
= \psi(\Omega_{K+1}+K^{K+1}+K)%
}

\NovelDisplay{M05528R01}{%
(0)(1,1,1)(2,1,1)(3,1,1)(4,1)(5,2)(2,1)(3,2,1)(4,2,1)(5,2,1)(6,2)(5,2,1)%
}

\NovelDisplay{M05529R01}{%
= \psi(\Omega_{K+1}+K^{K+1}\times \omega)%
}

\NovelDisplay{M05530R01}{%
(0)(1,1,1)(2,1,1)(3,1,1)(4,1)(5,2)(2,1)(3,2,1)(4,2,1)(5,2,1)(6,2)%
}
\NovelDisplay{M05530R02}{%
(5,2,1)(6,2)(4)%
}

\NovelDisplay{M05531R01}{%
= \psi(\Omega_{K+1}+K^{K\times 2})%
}

\NovelDisplay{M05532R01}{%
(0)(1,1,1)(2,1,1)(3,1,1)(4,1)(5,2)(2,1)(3,2,1)(4,2,1)(5,2,1)(6,2)%
}
\NovelDisplay{M05532R02}{%
(6,2)(4)%
}

\NovelDisplay{M05533R01}{%
= \psi(\Omega_{K+1}+K^{K^2})%
}

指数后继、倍数与平方继续在\allowbreak{}\NovelInline{I06134}{K}\allowbreak{}背景下出现，薛峰先找出共同结构，再核每个不同的上标。

第十二个月初，奥斯卡将补给分得比从前更细。

蔡兴坤问是不是山下的粮快没了，他摇头，说是有些人一说今天能多走几级，便又准备不吃东西。小舞听出他指谁，拿了一份先塞到唐三手里。

唐三看她，她便自己也取了一份。

诺当的脚早已不疼，这时却伸手摸了摸旧伤的位置。他把自己的本子放好，坐下来，将方才急着想问的话先留在纸角。叶阳在旁边看过，等两人吃完，才把那个完整的问题重新拿出来。

蔡须坤与霍雨浩没有催。风从他们身边过去，纸上的墨慢慢干了。

\NovelDisplay{M05534R01}{%
(0)(1,1,1)(2,1,1)(3,1,1)(4,1)(5,2)(2,1)(3,2,1)(4,2,1)(5,2,1)(6,2)%
}
\NovelDisplay{M05534R02}{%
(7,2)(4)%
}

\NovelDisplay{M05535R01}{%
= \psi(\Omega_{K+1}+K^{K^K})%
}

\NovelDisplay{M05536R01}{%
(0)(1,1,1)(2,1,1)(3,1,1)(4,1)(5,2)(2,1)(3,2,1)(4,2,1)(5,2,1)(6,2)(7,3)%
}

\NovelDisplay{M05537R01}{%
= \psi(\Omega_{K+1}+\psi_{\Omega_{K+1}}(\Omega_{K+1}))%
}

\NovelDisplay{M05538R01}{%
(0)(1,1,1)(2,1,1)(3,1,1)(4,1)(5,2)(2,1)(3,2,1)(4,2,1)(5,2,1)(6,2)%
}
\NovelDisplay{M05538R02}{%
(7,3)(4,2)(5,3)%
}

\NovelDisplay{M05539R01}{%
= \psi(\Omega_{K+1}+\psi_{\Omega_{K+1}}(\Omega_{K+1})+K)%
}

\NovelDisplay{M05540R01}{%
(0)(1,1,1)(2,1,1)(3,1,1)(4,1)(5,2)(2,1)(3,2,1)(4,2,1)(5,2,1)(6,2)(7,3)%
}
\NovelDisplay{M05540R02}{%
(4,2)(5,3,1)(6,3,1)(7,3,1)(8,3)(6)%
}

\NovelDisplay{M05541R01}{%
= \psi(\Omega_{K+1}+\psi_{\Omega_{K+1}}(\Omega_{K+1})+K^K)%
}

新的ψ作用于背景本身以后，随后加上的\allowbreak{}\NovelInline{I06135}{K}\allowbreak{}及其幂仍在外层，不能误接进前一个ψ的下标。

\NovelDisplay{M05542R01}{%
(0)(1,1,1)(2,1,1)(3,1,1)(4,1)(5,2)(2,1)(3,2,1)(4,2,1)(5,2,1)(6,2)(7,3)%
}
\NovelDisplay{M05542R02}{%
(4,2)(5,3,1)(6,3,1)(7,3,1)(8,3)(9,3)(6)%
}

\NovelDisplay{M05543R01}{%
= \psi(\Omega_{K+1}+\psi_{\Omega_{K+1}}(\Omega_{K+1})+K^{K^K})%
}

\NovelDisplay{M05544R01}{%
(0)(1,1,1)(2,1,1)(3,1,1)(4,1)(5,2)(2,1)(3,2,1)(4,2,1)(5,2,1)(6,2)(7,3)%
}
\NovelDisplay{M05544R02}{%
(4,2)(5,3,1)(6,3,1)(7,3,1)(8,3)(9,4)%
}

\NovelDisplay{M05545R01}{%
= \psi(\Omega_{K+1}+\psi_{\Omega_{K+1}}(\Omega_{K+1})\times 2)%
}

\NovelDisplay{M05546R01}{%
(0)(1,1,1)(2,1,1)(3,1,1)(4,1)(5,2)(2,1,1)%
}

\NovelDisplay{M05547R01}{%
= \psi(\Omega_{K+1}+\psi_{\Omega_{K+1}}(\Omega_{K+1}+1))%
}

\NovelDisplay{M05548R01}{%
(0)(1,1,1)(2,1,1)(3,1,1)(4,1)(5,2)(2,1,1)(3,1,1)%
}

\NovelDisplay{M05549R01}{%
= \psi(1-2\;1-\{\Omega_{K+1}+\psi_{\Omega_{K+1}}(\Omega_{K+1}+x)\})%
}

这里的倍数与内层后继相邻，诺当将每个括号的末端都指完，才没有把两条不同的接法混成一处。

诺当将右式收在另一面，请叶阳与自己分别从同一矩阵求项，蔡须坤等三项都写全以后，才逐列看过保留段与新生成的部分。

\NovelDisplay{M05550R01}{%
\text{展开第零项：}\quad(0)(1,1,1)(2,1,1)(3,1,1)(4,1)(5,2)(2,1)(3,2,1)(4,2,1)%
}
\NovelDisplay{M05550R02}{%
(5,2,1)(6,2)(7,3)(4,2)(5,3,1)(6,3,1)(7,3,1)(8,3)(9,3)%
}

\NovelDisplay{M05551R01}{%
\text{展开第一项：}\quad(0)(1,1,1)(2,1,1)(3,1,1)(4,1)(5,2)(2,1)(3,2,1)(4,2,1)%
}
\NovelDisplay{M05551R02}{%
(5,2,1)(6,2)(7,3)(4,2)(5,3,1)(6,3,1)(7,3,1)(8,3)(9,3)(5,3,1)(6,3,1)%
}
\NovelDisplay{M05551R03}{%
(7,3,1)(8,3)(9,3)%
}

\NovelDisplay{M05552R01}{%
\text{展开第二项：}\quad(0)(1,1,1)(2,1,1)(3,1,1)(4,1)(5,2)(2,1)(3,2,1)(4,2,1)%
}
\NovelDisplay{M05552R02}{%
(5,2,1)(6,2)(7,3)(4,2)(5,3,1)(6,3,1)(7,3,1)(8,3)(9,3)(5,3,1)(6,3,1)%
}
\NovelDisplay{M05552R03}{%
(7,3,1)(8,3)(9,3)(5,3,1)(6,3,1)(7,3,1)(8,3)(9,3)%
}

同一条长背景上的取项必须各自由原式开始，不能把第三份结果当成前两份拼起来的总和。

诺当把自己的错误留在页角，等今天再遇到相同位置时，便先停下来重看原列。

\NovelDisplay{M05553R01}{%
(0)(1,1,1)(2,1,1)(3,1,1)(4,1)(5,2)(3,1)(4,2)%
}

\NovelDisplay{M05554R01}{%
= \psi(\Omega_{K+1}+\psi_{\Omega_{K+1}}(\Omega_{K+1}+K)+K)%
}

\NovelDisplay{M05555R01}{%
(0)(1,1,1)(2,1,1)(3,1,1)(4,1)(5,2)(3,1)(4,2,1)(5,2,1)(6,2,1)(7,2)(8,3)%
}

\NovelDisplay{M05556R01}{%
= \psi(\Omega_{K+1}+\psi_{\Omega_{K+1}}(\Omega_{K+1}+K)+\psi_{\Omega_{K+1}}(\Omega_{K+1}))%
}

\NovelDisplay{M05557R01}{%
(0)(1,1,1)(2,1,1)(3,1,1)(4,1)(5,2)(3,1)(4,2,1)(5,2,1)(6,2,1)(7,2)%
}
\NovelDisplay{M05557R02}{%
(8,3)(6,2)(7,3)%
}

\NovelDisplay{M05558R01}{%
= \psi(\Omega_{K+1}+\psi_{\Omega_{K+1}}(\Omega_{K+1}+K)\times 2+K)%
}

\NovelDisplay{M05559R01}{%
(0)(1,1,1)(2,1,1)(3,1,1)(4,1)(5,2)(3,1,1)%
}

\NovelDisplay{M05560R01}{%
= \psi(\Omega_{K+1}+\psi_{\Omega_{K+1}}(\Omega_{K+1}+K+1))%
}

相同背景下的外部加项、倍数与内层后继逐条保留，霍雨浩没有用一句看起来差不多跳过任何一条。

\NovelDisplay{M05561R01}{%
(0)(1,1,1)(2,1,1)(3,1,1)(4,1)(5,2)(3,1,1)(2,1,1)%
}

\NovelDisplay{M05562R01}{%
= \psi(\Omega_{K+1}+\psi_{\Omega_{K+1}}(\Omega_{K+1}+K+2))%
}

\NovelDisplay{M05563R01}{%
(0)(1,1,1)(2,1,1)(3,1,1)(4,1)(5,2)(3,1,1)(3,1)(4,2)%
}

\NovelDisplay{M05564R01}{%
= \psi(\Omega_{K+1}+\psi_{\Omega_{K+1}}(\Omega_{K+1}+K\times 2)+K)%
}

\NovelDisplay{M05565R01}{%
(0)(1,1,1)(2,1,1)(3,1,1)(4,1)(5,2)(3,1,1)(3,1,1)%
}

\NovelDisplay{M05566R01}{%
= \psi(\Omega_{K+1}+\psi_{\Omega_{K+1}}(\Omega_{K+1}+K\times 2+1))%
}

\NovelDisplay{M05567R01}{%
(0)(1,1,1)(2,1,1)(3,1,1)(4,1)(5,2)(3,1,1)(4,1)(2)%
}

\NovelDisplay{M05568R01}{%
= \psi(\Omega_{K+1}+\psi_{\Omega_{K+1}}(\Omega_{K+1}+K^2))%
}

读到内层参数的倍数与平方，蔡兴坤先指出完整的作用对象，再把尾部那一项交给叶阳复查。

\NovelDisplay{M05569R01}{%
(0)(1,1,1)(2,1,1)(3,1,1)(4,1)(5,2)(3,1,1)(4,1)(3,1)(4,2)%
}

\NovelDisplay{M05570R01}{%
= \psi(\Omega_{K+1}+\psi_{\Omega_{K+1}}(\Omega_{K+1}+K^2)+K)%
}

\NovelDisplay{M05571R01}{%
(0)(1,1,1)(2,1,1)(3,1,1)(4,1)(5,2)(3,1,1)(4,1)(3,1,1)%
}

\NovelDisplay{M05572R01}{%
= \psi(\Omega_{K+1}+\psi_{\Omega_{K+1}}(\Omega_{K+1}+K^2+1))%
}

\NovelDisplay{M05573R01}{%
(0)(1,1,1)(2,1,1)(3,1,1)(4,1)(5,2)(3,1,1)(4,1)(3,1,1)(4,1)(2)%
}

\NovelDisplay{M05574R01}{%
= \psi(\Omega_{K+1}+\psi_{\Omega_{K+1}}(\Omega_{K+1}+K^2\times 2))%
}

\NovelDisplay{M05575R01}{%
(0)(1,1,1)(2,1,1)(3,1,1)(4,1)(5,2)(3,1,1)(4,1)(4,1)(2)%
}

\NovelDisplay{M05576R01}{%
= \psi(\Omega_{K+1}+\psi_{\Omega_{K+1}}(\Omega_{K+1}+K^3))%
}

平方后的后继与平方本身的倍数在此不同，蔡须坤让卢本韦先划定上标，才将整个式子向后接。

\NovelDisplay{M05577R01}{%
(0)(1,1,1)(2,1,1)(3,1,1)(4,1)(5,2)(3,1,1)(4,1)(5,1)(2)%
}

\NovelDisplay{M05578R01}{%
= \psi(\Omega_{K+1}+\psi_{\Omega_{K+1}}(\Omega_{K+1}+K^K))%
}

\NovelDisplay{M05579R01}{%
(0)(1,1,1)(2,1,1)(3,1,1)(4,1)(5,2)(3,1,1)(4,1)(5,2)%
}

\NovelDisplay{M05580R01}{%
= \psi(\Omega_{K+1}+\psi_{\Omega_{K+1}}(\Omega_{K+1}+\psi_{\Omega_{K+1}}(\Omega_{K+1})))%
}

\NovelDisplay{M05581R01}{%
(0)(1,1,1)(2,1,1)(3,1,1)(4,1)(5,2)(4)%
}

\NovelDisplay{M05582R01}{%
= \psi(\Omega_{K+1}+\psi_{\Omega_{K+1}}(\Omega_{K+1}+\psi_{\Omega_{K+1}}(\Omega_{K+1}+1)))%
}

\NovelDisplay{M05583R01}{%
(0)(1,1,1)(2,1,1)(3,1,1)(4,1)(5,2)(4,1)(5,2)%
}

\NovelDisplay{M05584R01}{%
= \psi(\Omega_{K+1}+\psi_{\Omega_{K+1}}(\Omega_{K+1}+\psi_{\Omega_{K+1}}(\Omega_{K+1}+\psi_{\Omega_{K+1}}(\Omega_{K+1}))))%
}

嵌套加深时，共同的\allowbreak{}\NovelInline{I06136}{\Omega_{K+1}}\allowbreak{}仍逐层出现，不能因为写过一次就把后面同名的作用对象省掉。

反射模式一层层缩短，卢本韦却比先前读得慢。

唐三起初以为她累了，伸手想把本子接来。她摇头，指着前缀，让他看一个不起眼的分组位置。长式里容易看见的层次，在这种短串里反而挤得更近，嘴上少一个停顿，便可能念成另一种结构。

古月娜把几张旧页取出，放在这几条旁边。小舞先说自己以为相同的地方，又听薛峰指出还没读到的那一截，终于没有把简短当成可以略过。

那一天，最高的刻度只向上多了一点。晚上收纸时，卢本韦却把今天自己能独立说明的几种写法都留了下来。

\NovelDisplay{M05585R01}{%
(0)(1,1,1)(2,1,1)(3,1,1)(4,1)(5,2)(5,1)(6,2)%
}

\NovelDisplay{M05586R01}{%
= \psi(\Omega_{K+1}+\psi_{\Omega_{K+1}}(\Omega_{K+1}+\psi_{\Omega_{K+1}}(\Omega_{K+1}+\psi_{\Omega_{K+1}}(\Omega_{K+1}+\psi_{\Omega_{K+1}}(\Omega_{K+1})))))%
}

\NovelDisplay{M05587R01}{%
(0)(1,1,1)(2,1,1)(3,1,1)(4,1)(5,2)(5,2)%
}

\NovelDisplay{M05588R01}{%
= \psi(\Omega_{K+1}\times 2)%
}

\NovelDisplay{M05589R01}{%
(0)(1,1,1)(2,1,1)(3,1,1)(4,1)(5,2)(6,3)%
}

\NovelDisplay{M05590R01}{%
= \psi(\mathrm{2nd}\;2\;\mathrm{aft}\;3)%
}

\NovelDisplay{M05591R01}{%
(0)(1,1,1)(2,1,1)(3,1,1)(4,1)(5,2,1)%
}

\NovelDisplay{M05592R01}{%
= \psi(1-2\;\mathrm{aft}\;3)%
}

外层倍数接到二次及一重模式前缀后，大家按模式结构阅读，不把其中的一与二直接当成普通系数计算。

\NovelDisplay{M05593R01}{%
(0)(1,1,1)(2,1,1)(3,1,1)(4,1)(5,2,1)(6,2,1)(7,2,1)(7,2,1)%
}

\NovelDisplay{M05594R01}{%
= \psi(1-2-2\;\mathrm{aft}\;3)%
}

\NovelDisplay{M05595R01}{%
(0)(1,1,1)(2,1,1)(3,1,1)(4,1)(5,2,1)(6,2,1)(7,2,1)(8)%
}

\NovelDisplay{M05596R01}{%
= \psi((2-)^\omega\;\mathrm{aft}\;3)%
}

\NovelDisplay{M05597R01}{%
(0)(1,1,1)(2,1,1)(3,1,1)(4,1)(5,2,1)(6,2,1)(7,2,1)(8)(7,2,1)%
}

\NovelDisplay{M05598R01}{%
= \psi(1-(2-)^{\omega+1}\;\mathrm{aft}\;3)%
}

\NovelDisplay{M05599R01}{%
(0)(1,1,1)(2,1,1)(3,1,1)(4,1)(5,2,1)(6,2,1)(7,2,1)(8,2)(7,2,1)%
}

\NovelDisplay{M05600R01}{%
= \psi(1-(2-)^{(1,1)}\;\mathrm{aft}\;3)%
}

模式前缀的重复和带参数的重复在这一组分开，薛峰让小舞找出完整的被重复对象才继续往下。

\NovelDisplay{M05601R01}{%
(0)(1,1,1)(2,1,1)(3,1,1)(4,1)(5,2,1)(6,2,1)(7,2,1)(8,2)(7,2,1)%
}
\NovelDisplay{M05601R02}{%
(8,2)(7,2,1)%
}

\NovelDisplay{M05602R01}{%
= \psi(1-(2-)^{(2,1)}\;\mathrm{aft}\;3)%
}

\NovelDisplay{M05603R01}{%
(0)(1,1,1)(2,1,1)(3,1,1)(4,1)(5,2,1)(6,2,1)(7,2,1)(8,2)(9,3)%
}

\NovelDisplay{M05604R01}{%
= \psi(2\;\mathrm{aft}\;\text{2nd}\;3)%
}

\NovelDisplay{M05605R01}{%
(0)(1,1,1)(2,1,1)(3,1,1)(4,1)(5,2,1)(6,2,1)(7,2,1)(8,2)(9,3,1)(10,3,1)%
}
\NovelDisplay{M05605R02}{%
(11,3,1)(12,3)(13,4)%
}

\NovelDisplay{M05606R01}{%
= \psi(2\;\mathrm{aft}\;\text{3rd}\;3)%
}

\NovelDisplay{M05607R01}{%
(0)(1,1,1)(2,1,1)(3,1,1)(4,1,1)%
}

\NovelDisplay{M05608R01}{%
= \psi(1-3)%
}

从既有前缀过渡到新的\allowbreak{}\NovelInline{I06137}{1-3}\allowbreak{}模式时，卢本韦先辨明哪一段属于aft关系，没有将连字符自行改成加减号。

\NovelDisplay{M05609R01}{%
(0)(1,1,1)(2,1,1)(3,1,1)(4,1,1)(2,1)(3,2,1)(4,2,1)(5,2,1)(6,2)(7,3)%
}

\NovelDisplay{M05610R01}{%
= \psi(2\;\mathrm{aft}\;3\;\mathrm{aft}\;1-3)%
}

\NovelDisplay{M05611R01}{%
(0)(1,1,1)(2,1,1)(3,1,1)(4,1,1)(2,1)(3,2,1)(4,2,1)(5,2,1)(6,2,1)%
}

\NovelDisplay{M05612R01}{%
= \psi(\mathrm{2nd}\;1-3)%
}

\NovelDisplay{M05613R01}{%
(0)(1,1,1)(2,1,1)(3,1,1)(4,1,1)(2,1,1)%
}

\NovelDisplay{M05614R01}{%
= \psi(1-1-3)%
}

\NovelDisplay{M05615R01}{%
(0)(1,1,1)(2,1,1)(3,1,1)(4,1,1)(2,1,1)(3,1,1)%
}

\NovelDisplay{M05616R01}{%
= \psi(1-2\;1-3)%
}

相邻的新模式保留了不同层次的前缀，唐三依原矩阵寻找生成段，终于不再只按模式字串的长短判断。

\NovelDisplay{M05617R01}{%
(0)(1,1,1)(2,1,1)(3,1,1)(4,1,1)(2,1,1)(3,1,1)(3,1,1)%
}

\NovelDisplay{M05618R01}{%
= \psi(1-2-2\;1-3)%
}

\NovelDisplay{M05619R01}{%
(0)(1,1,1)(2,1,1)(3,1,1)(4,1,1)(2,1,1)(3,1,1)(4,1)(2)%
}

\NovelDisplay{M05620R01}{%
= \psi((2-)^{(1,0)}\;1-3)%
}

\NovelDisplay{M05621R01}{%
(0)(1,1,1)(2,1,1)(3,1,1)(4,1,1)(2,1,1)(3,1,1)(4,1)(2,1,1)%
}

\NovelDisplay{M05622R01}{%
= \psi(1-(2-)^{(1,0)}\;1-3)%
}

\NovelDisplay{M05623R01}{%
(0)(1,1,1)(2,1,1)(3,1,1)(4,1,1)(2,1,1)(3,1,1)(4,1)(3,1,1)%
}

\NovelDisplay{M05624R01}{%
= \psi(1-(2-)^{(1,1)}\;1-3)%
}

带参数的重复接在同一模式之前时，参数对与被重复段都要读全，古月娜让诺当重新指出它们的边界。

四百九十米的横台上，蔡兴坤忽然不肯再走。

她蹲在石边，先说自己只是想歇一会儿，等大家真停下，又将脸埋进手里。她见前面几个人已经能说得那么清楚，越靠近终点，越怕到了最后，自己仍是被别人一路带上来的。

薛峰没有给她列出一串已经通过的高度，只请她从今天最确定的一处开始说。蔡兴坤最初说得很轻，奥斯卡便把补给包抱近，坐在旁边听。

说到后面，她终于指着一处，坦白这里还得再问。

``这就能问了。''薛峰说。

蔡兴坤抬起头。黑子鸡从包边探出来，恰好啄了啄她垂下的袖口，她低头望了一眼，慢慢把本子从膝下拿回。

\NovelDisplay{M05625R01}{%
(0)(1,1,1)(2,1,1)(3,1,1)(4,1,1)(2,1,1)(3,1,1)(4,1)(3,1,1)(4,1)(2)%
}

\NovelDisplay{M05626R01}{%
= \psi((2-)^{(2,0)}\;1-3)%
}

\NovelDisplay{M05627R01}{%
(0)(1,1,1)(2,1,1)(3,1,1)(4,1,1)(2,1,1)(3,1,1)(4,1)(5,2)%
}

\NovelDisplay{M05628R01}{%
= \psi(2\;\mathrm{aft}\;3\;1-3)%
}

\NovelDisplay{M05629R01}{%
(0)(1,1,1)(2,1,1)(3,1,1)(4,1,1)(2,1,1)(3,1,1)(4,1,1)%
}

\NovelDisplay{M05630R01}{%
= \psi(1-3\;1-3)%
}

\NovelDisplay{M05631R01}{%
(0)(1,1,1)(2,1,1)(3,1,1)(4,1,1)(3)%
}

\NovelDisplay{M05632R01}{%
= \psi((3\;1-)^\omega)%
}

这里将整段模式接入新的重复结构，奥斯卡用三次独立取项的经验提醒自己，不能把第二项理解成再展开第一项。

\NovelDisplay{M05633R01}{%
(0)(1,1,1)(2,1,1)(3,1,1)(4,1,1)(3,1)(2,1,1)(3,1,1)(4,1,1)%
}

\NovelDisplay{M05634R01}{%
= \psi(1-(3\;1-)^{(1,1)})%
}

\NovelDisplay{M05635R01}{%
(0)(1,1,1)(2,1,1)(3,1,1)(4,1,1)(3,1)(4,2)%
}

\NovelDisplay{M05636R01}{%
= \psi(2\;\mathrm{aft}\;2-3)%
}

\NovelDisplay{M05637R01}{%
(0)(1,1,1)(2,1,1)(3,1,1)(4,1,1)(3,1,1)%
}

\NovelDisplay{M05638R01}{%
= \psi(1-2-3)%
}

\NovelDisplay{M05639R01}{%
(0)(1,1,1)(2,1,1)(3,1,1)(4,1,1)(3,1,1)(3,1,1)%
}

\NovelDisplay{M05640R01}{%
= \psi(1-2-2-3)%
}

模式的末段改变以后，蔡须坤先保留共同前缀，霍雨浩再让每个人独立说明原矩阵为何需要从当前末列重新求项。

\NovelDisplay{M05641R01}{%
(0)(1,1,1)(2,1,1)(3,1,1)(4,1,1)(3,1,1)(4,1)(2)%
}

\NovelDisplay{M05642R01}{%
= \psi((2-)^{(1,0)}3)%
}

\NovelDisplay{M05643R01}{%
(0)(1,1,1)(2,1,1)(3,1,1)(4,1,1)(3,1,1)(4,1)(3,1)(2)%
}

\NovelDisplay{M05644R01}{%
= \psi(((2-)^{(1,0)}3\;1-)^{(1,0)})%
}

\NovelDisplay{M05645R01}{%
(0)(1,1,1)(2,1,1)(3,1,1)(4,1,1)(3,1,1)(4,1)(3,1,1)%
}

\NovelDisplay{M05646R01}{%
= \psi(1-(2-)^{(1,1)}3)%
}

\NovelDisplay{M05647R01}{%
(0)(1,1,1)(2,1,1)(3,1,1)(4,1,1)(3,1,1)(4,1)(5,2)%
}

\NovelDisplay{M05648R01}{%
= \psi(2\;\mathrm{aft}\;3\;2-3)%
}

套入另一层带参数结构时，括号外的前缀仍要保留，叶阳照着原式指出了自己先前漏掉的那一段。

\NovelDisplay{M05649R01}{%
(0)(1,1,1)(2,1,1)(3,1,1)(4,1,1)(3,1,1)(4,1,1)%
}

\NovelDisplay{M05650R01}{%
= \psi(1-3\;2-3)%
}

\NovelDisplay{M05651R01}{%
(0)(1,1,1)(2,1,1)(3,1,1)(4,1,1)(3,1,1)(4,1,1)(3,1,1)%
}

\NovelDisplay{M05652R01}{%
= \psi(1-2-3\;2-3)%
}

\NovelDisplay{M05653R01}{%
(0)(1,1,1)(2,1,1)(3,1,1)(4,1,1)(3,1,1)(4,1,1)(3,1,1)(4,1,1)%
}

\NovelDisplay{M05654R01}{%
= \psi(1-3\;2-3\;2-3)%
}

\NovelDisplay{M05655R01}{%
(0)(1,1,1)(2,1,1)(3,1,1)(4,1,1)(4,1)(2)%
}

\NovelDisplay{M05656R01}{%
= \psi((3\;2-)^{(1,0)})%
}

整段模式的接续与重复在这里交替出现，卢本韦按原分组逐一复述，没有让同一个短前缀遮住结构差异。

蔡须坤那几日常走在前面，又常停下来。

她先将自己的问题弄清，再听后来的人问。霍雨浩在旁边核条件，卢本韦将不同写法的页签翻给需要的人，薛峰有时先找到了接法，却仍会退回较低的那级，让同伴不必仰头才能听见。

他们没有替别人落那最后一脚。

古月娜将这一点看得最清楚。她扶小舞起身时只托住手臂，等对方自己站稳便放开；诺当请她指出自己遗漏的位置，她也先让他把已经知道的部分说完。

到了晚上，各自本子里的字迹仍不相同，能独立走过的路却开始慢慢合到一处。

\NovelDisplay{M05657R01}{%
(0)(1,1,1)(2,1,1)(3,1,1)(4,1,1)(4,1)(2,1,1)(3,1,1)(4,1,1)%
}

\NovelDisplay{M05658R01}{%
= \psi(1-3\;1-(3\;2-)^{(1,0)})%
}

\NovelDisplay{M05659R01}{%
(0)(1,1,1)(2,1,1)(3,1,1)(4,1,1)(4,1)(2,1,1)(3,1,1)(4,1,1)(4,1)(2)%
}

\NovelDisplay{M05660R01}{%
= \psi((3\;2-)^{(1,0)}1-(3\;2-)^{(1,0)})%
}

\NovelDisplay{M05661R01}{%
(0)(1,1,1)(2,1,1)(3,1,1)(4,1,1)(4,1)(3,1)(2)%
}

\NovelDisplay{M05662R01}{%
= \psi(((3\;2-)^{(1,0)}\;1-)^{(1,0)})%
}

\NovelDisplay{M05663R01}{%
(0)(1,1,1)(2,1,1)(3,1,1)(4,1,1)(4,1)(3,1,1)%
}

\NovelDisplay{M05664R01}{%
= \psi(1-2-(3\;2-)^{(1,0)})%
}

外层族与内层族分开以后，薛峰才读清这一组的参数位置，随后又将它与下一条aft模式分别保存。

卢本韦让大家暂时放下短模式，从对应的原矩阵独立求三项，先由薛峰指明三次计算都保留了哪一段。

\NovelDisplay{M05665R01}{%
\text{展开第零项：}\quad(0)(1,1,1)(2,1,1)(3,1,1)(4,1,1)(4,1)(2,1,1)%
}
\NovelDisplay{M05665R02}{%
(3,1,1)(4,1,1)(4,1)%
}

\NovelDisplay{M05666R01}{%
\text{展开第一项：}\quad(0)(1,1,1)(2,1,1)(3,1,1)(4,1,1)(4,1)(2,1,1)(3,1,1)%
}
\NovelDisplay{M05666R02}{%
(4,1,1)(4,1)(1,1,1)(2,1,1)(3,1,1)(4,1,1)(4,1)(2,1,1)(3,1,1)(4,1,1)(4,1)%
}

\NovelDisplay{M05667R01}{%
\text{展开第二项：}\quad(0)(1,1,1)(2,1,1)(3,1,1)(4,1,1)(4,1)(2,1,1)(3,1,1)%
}
\NovelDisplay{M05667R02}{%
(4,1,1)(4,1)(1,1,1)(2,1,1)(3,1,1)(4,1,1)(4,1)(2,1,1)(3,1,1)(4,1,1)(4,1)%
}
\NovelDisplay{M05667R03}{%
(1,1,1)(2,1,1)(3,1,1)(4,1,1)(4,1)(2,1,1)(3,1,1)(4,1,1)(4,1)%
}

模式字串的简短不能代替矩阵计算，三次取项各自的生成段仍要按原版规则确认。

蔡兴坤沿自己的那份结果逐列读过，终于将方才只敢含糊过去的问题说完整。

\NovelDisplay{M05668R01}{%
(0)(1,1,1)(2,1,1)(3,1,1)(4,1,1)(4,1)(3,1,1)(4,1,1)%
}

\NovelDisplay{M05669R01}{%
= \psi(1-(3\;2-)^{(1,1)})%
}

\NovelDisplay{M05670R01}{%
(0)(1,1,1)(2,1,1)(3,1,1)(4,1,1)(4,1)(5,2)%
}

\NovelDisplay{M05671R01}{%
= \psi(2\;\mathrm{aft}\;3-3)%
}

\NovelDisplay{M05672R01}{%
(0)(1,1,1)(2,1,1)(3,1,1)(4,1,1)(4,1,1)%
}

\NovelDisplay{M05673R01}{%
= \psi(1-3-3)%
}

\NovelDisplay{M05674R01}{%
(0)(1,1,1)(2,1,1)(3,1,1)(4,1,1)(4,1,1)(2,1,1)%
}

\NovelDisplay{M05675R01}{%
= \psi(1-1-3-3)%
}

模式进入新的共同末段时，蔡兴坤将原矩阵中对应的列重新圈出，确认这不是只在纸面上多添一个三。

\NovelDisplay{M05676R01}{%
(0)(1,1,1)(2,1,1)(3,1,1)(4,1,1)(4,1,1)(2,1,1)(3,1,1)(4,1,1)(4,1,1)%
}

\NovelDisplay{M05677R01}{%
= \psi(3-3\;1-3-3)%
}

\NovelDisplay{M05678R01}{%
(0)(1,1,1)(2,1,1)(3,1,1)(4,1,1)(4,1,1)(3,1,1)%
}

\NovelDisplay{M05679R01}{%
= \psi(1-2-3-3)%
}

\NovelDisplay{M05680R01}{%
(0)(1,1,1)(2,1,1)(3,1,1)(4,1,1)(4,1,1)(3,1,1)(4,1,1)%
}

\NovelDisplay{M05681R01}{%
= \psi(1-3\;2-3-3)%
}

\NovelDisplay{M05682R01}{%
(0)(1,1,1)(2,1,1)(3,1,1)(4,1,1)(4,1,1)(3,1,1)(4,1,1)(4,1,1)%
}

\NovelDisplay{M05683R01}{%
= \psi(1-3-3\;2-3-3)%
}

连续的整段回接把矩阵拉长，唐三与小舞各自核一遍保留段，再交换完整的问题而不是只交换最后答案。

\NovelDisplay{M05684R01}{%
(0)(1,1,1)(2,1,1)(3,1,1)(4,1,1)(4,1,1)(3,1,1)(4,1,1)(4,1,1)(3,1,1)%
}
\NovelDisplay{M05684R02}{%
(4,1,1)(4,1,1)%
}

\NovelDisplay{M05685R01}{%
= \psi(1-3-3\;2-3-3\;2-3-3)%
}

\NovelDisplay{M05686R01}{%
(0)(1,1,1)(2,1,1)(3,1,1)(4,1,1)(4,1,1)(4,1)(2)%
}

\NovelDisplay{M05687R01}{%
= \psi((3-3\;2-)^{(1,0)})%
}

\NovelDisplay{M05688R01}{%
(0)(1,1,1)(2,1,1)(3,1,1)(4,1,1)(4,1,1)(4,1)(3,1,1)(4,1,1)(4,1,1)%
}

\NovelDisplay{M05689R01}{%
= \psi(1-(3-3\;2-)^{(1,1)})%
}

\NovelDisplay{M05690R01}{%
(0)(1,1,1)(2,1,1)(3,1,1)(4,1,1)(4,1,1)(4,1)(5,2)%
}

\NovelDisplay{M05691R01}{%
= \psi(2\;\mathrm{aft}\;3-3-3)%
}

重复族接到下一层aft模式时，大家先把当前几条取项核稳，才让眼前更短的新矩阵成为新的起点。

这回由唐三、小舞与奥斯卡分别求项，霍雨浩只在三人都写完以后核对，不在中途替任何一人接出末列。

\NovelDisplay{M05692R01}{%
\text{展开第零项：}\quad(0)(1,1,1)(2,1,1)(3,1,1)(4,1,1)(4,1,1)(3,1,1)(4,1,1)%
}
\NovelDisplay{M05692R02}{%
(4,1,1)(3,1,1)(4,1,1)%
}

\NovelDisplay{M05693R01}{%
\text{展开第一项：}\quad(0)(1,1,1)(2,1,1)(3,1,1)(4,1,1)(4,1,1)(3,1,1)(4,1,1)%
}
\NovelDisplay{M05693R02}{%
(4,1,1)(3,1,1)(4,1,1)(4,1)(5,2,1)(6,2,1)(7,2,1)(8,2,1)(8,2,1)(7,2,1)%
}
\NovelDisplay{M05693R03}{%
(8,2,1)(8,2,1)(7,2,1)(8,2,1)%
}

\NovelDisplay{M05694R01}{%
\text{展开第二项：}\quad(0)(1,1,1)(2,1,1)(3,1,1)(4,1,1)(4,1,1)(3,1,1)(4,1,1)%
}
\NovelDisplay{M05694R02}{%
(4,1,1)(3,1,1)(4,1,1)(4,1)(5,2,1)(6,2,1)(7,2,1)(8,2,1)(8,2,1)(7,2,1)%
}
\NovelDisplay{M05694R03}{%
(8,2,1)(8,2,1)(7,2,1)(8,2,1)(8,2)(9,3,1)(10,3,1)(11,3,1)(12,3,1)%
}
\NovelDisplay{M05694R04}{%
(12,3,1)(11,3,1)(12,3,1)(12,3,1)(11,3,1)(12,3,1)%
}

较长的重复段随编号改变，求项时既不能漏列，也不能把已经重复的整段再凭印象添一次。

小舞把第二项重新数过，向唐三指出自己前一天多添的那一段，随后自己划去错误，重新写全。

最后一段长坡，唐三与小舞又用了几天。

小舞有两次在同一个地方停住，第三次上来前，气得把笔搁到旁边。唐三见她这样，没有说自己从前也卡过，只拿起她的原式，陪她重新找最末一列。

她先指出一处，自己又摇头。唐三便等。

等得久了，小舞抬头，问他为什么今天不抢着说。

``答应过你。''

她看了他一会儿，忽然笑起来，把笔拿了回去。那次她仍没有立刻走过去，却在下山前，终于说清自己前两次是怎样把不同分组错接到了一起。

\NovelDisplay{M05695R01}{%
(0)(1,1,1)(2,1,1)(3,1,1)(4,1,1)(4,1,1)(4,1,1)%
}

\NovelDisplay{M05696R01}{%
= \psi(1-3-3-3)%
}

\NovelDisplay{M05697R01}{%
(0)(1,1,1)(2,1,1)(3,1,1)(4,1,1)(4,1,1)(4,1,1)(4,1,1)%
}

\NovelDisplay{M05698R01}{%
= \psi(1-3-3-3-3)%
}

\NovelDisplay{M05699R01}{%
(0)(1,1,1)(2,1,1)(3,1,1)(4,1,1)(5)%
}

\NovelDisplay{M05700R01}{%
= \psi((3-)^\omega)%
}

\NovelDisplay{M05701R01}{%
(0)(1,1,1)(2,1,1)(3,1,1)(4,1,1)(5,1)(2)%
}

\NovelDisplay{M05702R01}{%
= \psi((3-)^{(1,0)})%
}

同一短前缀在不同参数中反复出现，卢本韦仍逐个读清分组，没有因为临近五百米便把后面的记号一口念过。

\NovelDisplay{M05703R01}{%
(0)(1,1,1)(2,1,1)(3,1,1)(4,1,1)(5,1)(2,1,1)%
}

\NovelDisplay{M05704R01}{%
= \psi(1-(3-)^{(1,0)})%
}

\NovelDisplay{M05705R01}{%
(0)(1,1,1)(2,1,1)(3,1,1)(4,1,1)(5,1)(2,1,1)(3,1,1)(4,1,1)(5,1)(2)%
}

\NovelDisplay{M05706R01}{%
= \psi((3-)^{(1,0)}\;1-(3-)^{(1,0)})%
}

\NovelDisplay{M05707R01}{%
(0)(1,1,1)(2,1,1)(3,1,1)(4,1,1)(5,1)(3,1)(2)%
}

\NovelDisplay{M05708R01}{%
= \psi(((3-)^{(1,0)}\;1-)^{(1,0)})%
}

\NovelDisplay{M05709R01}{%
(0)(1,1,1)(2,1,1)(3,1,1)(4,1,1)(5,1)(3,1,1)%
}

\NovelDisplay{M05710R01}{%
= \psi(1-2-(3-)^{(1,0)})%
}

参数族之外再接整段模式时，霍雨浩提醒叶阳先分内外，不能把外部接续算成族参数的一部分。

\NovelDisplay{M05711R01}{%
(0)(1,1,1)(2,1,1)(3,1,1)(4,1,1)(5,1)(3,1,1)(4,1,1)(5,1)(2)%
}

\NovelDisplay{M05712R01}{%
= \psi((3-)^{(1,0)}\;2-(3-)^{(1,0)})%
}

\NovelDisplay{M05713R01}{%
(0)(1,1,1)(2,1,1)(3,1,1)(4,1,1)(5,1)(4,1)(2)%
}

\NovelDisplay{M05714R01}{%
= \psi(((3-)^{(1,0)}\;2-)^{(1,0)})%
}

\NovelDisplay{M05715R01}{%
(0)(1,1,1)(2,1,1)(3,1,1)(4,1,1)(5,1)(4,1,1)%
}

\NovelDisplay{M05716R01}{%
= \psi(1-(3-)^{(1,1)})%
}

\NovelDisplay{M05717R01}{%
(0)(1,1,1)(2,1,1)(3,1,1)(4,1,1)(5,1)(4,1,1)(4,1,1)%
}

\NovelDisplay{M05718R01}{%
= \psi(1-(3-)^{(1,2)})%
}

从两种带参数的结构转向更高aft模式，薛峰先指出完整旧段在哪里结束，再独立接上新的一处。

\NovelDisplay{M05719R01}{%
(0)(1,1,1)(2,1,1)(3,1,1)(4,1,1)(5,1)(4,1,1)(5,1)(2)%
}

\NovelDisplay{M05720R01}{%
= \psi((3-)^{(2,0)})%
}

\NovelDisplay{M05721R01}{%
(0)(1,1,1)(2,1,1)(3,1,1)(4,1,1)(5,1)(6,2)%
}

\NovelDisplay{M05722R01}{%
= \psi(2\;\mathrm{aft}\;4)%
}

\NovelDisplay{M05723R01}{%
(0)(1,1,1)(2,1,1)(3,1,1)(4,1,1)(5,1,1)%
}

\NovelDisplay{M05724R01}{%
= \psi(1-4)%
}

\NovelDisplay{M05725R01}{%
(0)(1,1,1)(2,1,1)(3,1,1)(4,1,1)(5,1,1)(2,1,1)%
}

\NovelDisplay{M05726R01}{%
= \psi(1-1-4)%
}

读到\allowbreak{}\NovelInline{I06138}{1-4}\allowbreak{}一组时，蔡须坤让每个人轮流给同伴指出原矩阵，确保队伍的进度没有遮住个人仍空着的一页。

终点只剩最后一小段时，夏夜星空没有往前站。

她在下面的横台留好水，等每个人都来过，又看着他们自己走回原来的位置。最菜萌新也到了，远远站在石阶外，既没有叫停，也没有给出哪一式的答案。

卢本韦看见他，手在自己的册子上轻轻按了一下。蔡须坤问怎么了，她只说，等上去以后，有一件事要亲自看看。

霍雨浩将最后几页按好，薛峰回头数过人。十一份本子都在，谁的那一份也没有被合到别人的记录里面。

\NovelDisplay{M05727R01}{%
(0)(1,1,1)(2,1,1)(3,1,1)(4,1,1)(5,1,1)(2,1,1)(3,1,1)(4,1,1)(5,1,1)%
}

\NovelDisplay{M05728R01}{%
= \psi(1-4\;1-4)%
}

\NovelDisplay{M05729R01}{%
(0)(1,1,1)(2,1,1)(3,1,1)(4,1,1)(5,1,1)(3,1)(4,2)%
}

\NovelDisplay{M05730R01}{%
= \psi(2\;\mathrm{aft}\;2-4)%
}

\NovelDisplay{M05731R01}{%
(0)(1,1,1)(2,1,1)(3,1,1)(4,1,1)(5,1,1)(3,1,1)%
}

\NovelDisplay{M05732R01}{%
= \psi(1-2-4)%
}

\NovelDisplay{M05733R01}{%
(0)(1,1,1)(2,1,1)(3,1,1)(4,1,1)(5,1,1)(3,1,1)(4,1,1)(5,1,1)%
}

\NovelDisplay{M05734R01}{%
= \psi(1-4\;2-4)%
}

相同的四结尾前缀与不同aft对象在此并列，古月娜以原式分组为准，没有让相近的字面替代实际核对。

\NovelDisplay{M05735R01}{%
(0)(1,1,1)(2,1,1)(3,1,1)(4,1,1)(5,1,1)(4,1)(5,2)%
}

\NovelDisplay{M05736R01}{%
= \psi(2\;\mathrm{aft}\;3-4)%
}

\NovelDisplay{M05737R01}{%
(0)(1,1,1)(2,1,1)(3,1,1)(4,1,1)(5,1,1)(4,1,1)%
}

\NovelDisplay{M05738R01}{%
= \psi(1-3-4)%
}

\NovelDisplay{M05739R01}{%
(0)(1,1,1)(2,1,1)(3,1,1)(4,1,1)(5,1,1)(4,1,1)(5,1,1)%
}

\NovelDisplay{M05740R01}{%
= \psi(1-4\;3-4)%
}

\NovelDisplay{M05741R01}{%
(0)(1,1,1)(2,1,1)(3,1,1)(4,1,1)(5,1,1)(5,1)(6,2)%
}

\NovelDisplay{M05742R01}{%
= \psi(2\;\mathrm{aft}\;4-4)%
}

这几条模式虽然简短，仍各有自己的前后结构，诺当逐条复述过，才将已经练稳的位置交给后来的人看。

\NovelDisplay{M05743R01}{%
(0)(1,1,1)(2,1,1)(3,1,1)(4,1,1)(5,1,1)(5,1,1)%
}

\NovelDisplay{M05744R01}{%
= \psi(1-4-4)%
}

\NovelDisplay{M05745R01}{%
(0)(1,1,1)(2,1,1)(3,1,1)(4,1,1)(5,1,1)(6)%
}

\NovelDisplay{M05746R01}{%
= \psi((4-)^\omega)%
}

\NovelDisplay{M05747R01}{%
(0)(1,1,1)(2,1,1)(3,1,1)(4,1,1)(5,1,1)(6,1)(2)%
}

\NovelDisplay{M05748R01}{%
= \psi((4-)^{(1,0)})%
}

\NovelDisplay{M05749R01}{%
(0)(1,1,1)(2,1,1)(3,1,1)(4,1,1)(5,1,1)(6,1)(2,1,1)%
}

\NovelDisplay{M05750R01}{%
= \psi(1-(4-)^{(1,0)})%
}

从末段重复进入带参数的模式族时，小舞先读清括号中的完整对象，没有再把上标只当成普通次数。

到最后那天，蔡须坤醒来时，窗外还很暗。

她没有像往常一样先把长式全念一遍，而是取出最初的山图，从第一次落脚的日期慢慢看到今天。七十米那一段用了两个月，后来的停顿、伤病、风雨与回头，都有各自的位置。

卢本韦起身，问她还看什么。

``正好一年。''

卢本韦走过来，望着图上的字，过了一会儿才将自己的I册拿起。两人出了门，霍雨浩与薛峰已经在外面等，唐三、小舞和后来的人也陆续跟上。

这一次，没有谁先说今天一定得过去。

\NovelDisplay{M05751R01}{%
(0)(1,1,1)(2,1,1)(3,1,1)(4,1,1)(5,1,1)(6,1)(3,1,1)%
}

\NovelDisplay{M05752R01}{%
= \psi(1-2-(4-)^{(1,0)})%
}

\NovelDisplay{M05753R01}{%
(0)(1,1,1)(2,1,1)(3,1,1)(4,1,1)(5,1,1)(6,1)(4,1,1)%
}

\NovelDisplay{M05754R01}{%
= \psi(1-3-(4-)^{(1,0)})%
}

\NovelDisplay{M05755R01}{%
(0)(1,1,1)(2,1,1)(3,1,1)(4,1,1)(5,1,1)(6,1)(5,1,1)%
}

\NovelDisplay{M05756R01}{%
= \psi(1-(4-)^{(1,1)})%
}

\NovelDisplay{M05757R01}{%
(0)(1,1,1)(2,1,1)(3,1,1)(4,1,1)(5,1,1)(6,1)(5,1,1)(6,1)(2)%
}

\NovelDisplay{M05758R01}{%
= \psi((4-)^{(2,0)})%
}

这一组改变的是模式族的参数与前缀，卢本韦守住对应关系，蔡须坤则回到每条原矩阵独立复算。

最后一次在营地复核时，古月娜请每个人从原矩阵重新求项，再任选三份独立结果，照取项编号并排放好。

\NovelDisplay{M05759R01}{%
\text{展开第零项：}\quad(0)(1,1,1)(2,1,1)(3,1,1)(4,1,1)(5,1,1)(6,1)(5,1,1)(6,1)%
}

\NovelDisplay{M05760R01}{%
\text{展开第一项：}\quad(0)(1,1,1)(2,1,1)(3,1,1)(4,1,1)(5,1,1)(6,1)(5,1,1)(6,1)%
}
\NovelDisplay{M05760R02}{%
(1,1,1)(2,1,1)(3,1,1)(4,1,1)(5,1,1)(6,1)(5,1,1)(6,1)%
}

\NovelDisplay{M05761R01}{%
\text{展开第二项：}\quad(0)(1,1,1)(2,1,1)(3,1,1)(4,1,1)(5,1,1)(6,1)(5,1,1)(6,1)%
}
\NovelDisplay{M05761R02}{%
(1,1,1)(2,1,1)(3,1,1)(4,1,1)(5,1,1)(6,1)(5,1,1)(6,1)(1,1,1)(2,1,1)%
}
\NovelDisplay{M05761R03}{%
(3,1,1)(4,1,1)(5,1,1)(6,1)(5,1,1)(6,1)%
}

独立取项的起点始终是同一原式，正确的编号与正确的保留段需要同时留下。

没有人再问能否只背最短的一份，大家各自收好纸，等第二天由自己走完。

临近五百米横台，最后一处压力忽然收紧。

蔡须坤停在那一级上，后面许多曾经见过的结构，一时都像要同时挤进眼前这条短矩阵。她没有挑一个最熟的答案硬接，先把那些页放开，只看原式本身。

霍雨浩与薛峰还在较低处，各自核对尚未走清的接续。卢本韦站在蔡须坤身侧，将笔收好，没有朝她念答案。蔡须坤等胸口那阵急促过去，才重新指出最后这几列。

身后的唐三也停下来，小舞在较低处喘气。诺当、奥斯卡、蔡兴坤与叶阳分别守着各自尚未完成的一处，古月娜给他们留出能够独立看清纸的位置。

上方的雾一点点散了。

\NovelDisplay{M05762R01}{%
(0)(1,1,1)(2,1,1)(3,1,1)(4,1,1)(5,1,1)(6,1)(7,2)%
}

\NovelDisplay{M05763R01}{%
= \psi(2\;\mathrm{aft}\;5)%
}

\NovelDisplay{M05764R01}{%
(0)(1,1,1)(2,1,1)(3,1,1)(4,1,1)(5,1,1)(6,1,1)%
}

\NovelDisplay{M05765R01}{%
= \psi(1-5)%
}

\NovelDisplay{M05766R01}{%
(0)(1,1,1)(2,1,1)(3,1,1)(4,1,1)(5,1,1)(6,1,1)(7,1)(8,2)%
}

\NovelDisplay{M05767R01}{%
= \psi(2\;\mathrm{aft}\;6)%
}

\NovelDisplay{M05768R01}{%
(0)(1,1,1)(2,1,1)(3,1,1)(4,1,1)(5,1,1)(6,1,1)(7,1,1)%
}

\NovelDisplay{M05769R01}{%
= \psi(1-6)%
}

\NovelDisplay{M05770R01}{%
\textcolor{red}{(0)(1,1,1)(2,2)}%
}

\NovelDisplay{M05771R01}{%
\textcolor{red}{= \psi(\mathrm{psd}.\Pi_\omega)}%
}

\NovelDisplay{M05772R01}{%
\textcolor{red}{= \psi(\lambda{\alpha}.{\alpha+1}-\Pi_{0})}%
}

最后三列矩阵的两种右侧表示照讲义保留，终点的接续仍回到原版BMS逐项展开，不以魂力强行压过未明白之处。

最终，蔡须坤将右侧表示翻到背面，从这条只有三列的原矩阵重新求项，卢本韦、霍雨浩与薛峰也各自独立重做一遍。

\NovelDisplay{M05773R01}{%
\text{展开第零项：}\quad(0)(1,1,1)%
}

\NovelDisplay{M05774R01}{%
\text{展开第一项：}\quad(0)(1,1,1)(2,1,1)%
}

\NovelDisplay{M05775R01}{%
\text{展开第二项：}\quad(0)(1,1,1)(2,1,1)(3,1,1)%
}

第零项、第一项和第二项都从这条三列原式求出，第三行的生成必须依原规则保留，不能把简短误读成无需展开。

蔡须坤看清那三份结果，手指一点点松开，先将自己的纸收好，随后才抬起脚。

五百米横台的石面很平。

蔡须坤落脚以后，先听见身侧传来一声很轻的喘息。卢本韦也收起了脚跟，手里的纸慢慢垂下来。两人对视片刻，谁都没有立刻出声。

霍雨浩与薛峰还在下面。古月娜站在另一侧的石阶上，唐三低头核纸，小舞坐在更低的平处，等呼吸缓下来才重新拿笔。奥斯卡、诺当、蔡兴坤与叶阳，也各有尚未走完的一段。

石阶间的风忽然停了。

一道声音从山间响起，清清楚楚落在众人心里。

``坤神第一考，登上中分山五百米。''

``蔡须坤、卢本韦，第一考通过。''

蔡须坤的手停在胸前。直到那两个名字真正响起，她才相信第一考要求的最后一级，已经落在脚下。

``自入考之日起，历时一年。''

五百米的刻痕，在石面上一点点亮起来。加在两人身上的考核压力缓缓撤去，下面的石阶却仍有各自的重量，更高处的路也仍隐在云里。

``看看你们这一年所得。''

蔡须坤闭上眼。考核留下的光，分成两列，清清楚楚落在她心中；身旁的卢本韦，也看清了同样的记录。

扽力：

\NovelDisplay{M05776R01}{%
(0)(1,1,1)(2,2)%
}

魂力：

\NovelDisplay{M05777R01}{%
(0)(1,1,1)(2,2,2)%
}

她们一直用哪一段能独立分析、哪一步能够站稳来辨认进度。压力底下，每一点增强都被接下来的艰难包住，直到此刻，才看清一年究竟积累了多少。

卢本韦的心里，另有一道极轻的回应。

那道向来随她意志出现的母猪扽魂，忽然像从她尚未真正看见的地方向她望来。那份回应有自己的停顿，带着她没有先行下达的意思。

``记住，你的扽魂也是魂灵。''

最菜萌新从前说过的话，这时重新浮起。她接住那份回应，睁眼去看蔡须坤，发现对方也正望着自己。

果糕领域只铺开很薄一层，留在两人脚边，没有伸进下方仍在考核的石阶。

夏夜星空走到横台边，将水递给她们，又向下面望去。

``每个人都要独立走到这里。你们先通过，其他人的那一份还在继续。''

霍雨浩抬头，向两人点了点头。薛峰也笑了一下，随即重新看自己的原式。唐三没有替小舞收纸，奥斯卡把水放到诺当身旁，等他自己将问题指出。

蔡兴坤望着五百米的刻痕，手在包带上握紧，又慢慢松开。她把黑子鸡护好，从自己尚未接清的位置重新开口。

蔡须坤看着他们，没有催，也没有从上面念答案。她把今日真正走清的几页收好，等需要问的人抬头。

这一年大家同走一条山路，停下的位置却从来不尽相同。此前每逢几人先到横台，她们便坐在旁边等；到了第一考最后一级，名字终于被单独报出，这份差别才显得格外清楚。

小舞低头望了一会儿自己的纸，伸手将被汗浸软的页角抚平。唐三问她要不要先下山歇，她摇头，先指给他看自己已经弄明白的那一段，再说剩下的那处准备怎样问。

横台外，最菜萌新向卢本韦轻轻点头。她想起那份回应，暂时没有追问，只把自己的册子好好收住。

``山顶还远着呢。''小舞忽然抬头道。

``整座中分山有两千五百米。''夏夜星空说，``第一考只要求五百米。等整个BMS矩阵分析完成，才是登顶。''

蔡须坤望向更高处，又看看自己的笔记。理想不提升BMS仍未定义，原版BMS也还有长路。

她与卢本韦坐到石桩旁，先将水喝完。

她们都没有急着往更高处再走。

从第一次站到中分山脚，到最先两人通过第一考，正好一年。下面的伙伴仍在努力，五百米刻痕上的光，安静地照着她们各自收好的纸。

\chapter{第一百三十章：双领域与第二考}\label{ux7b2cux4e00ux767eux4e09ux5341ux7ae0ux53ccux9886ux57dfux4e0eux7b2cux4e8cux8003}

蔡兴坤上到五百米时，天已经快黑了。

她最后一步迈得很慢。鞋尖落稳以后，肩膀还绷着，像怕自己只要一松劲，就会被石阶上那股力量重新压回去。

直到石桩上亮起光，她才抬起头。

山间的声音报出了她的名字。

蔡须坤在旁边等着，看见她将脚跟收上来，伸手接住了她怀里的黑子鸡。蔡兴坤空出手，先摸了摸石桩上的刻痕，又低头望向自己走来的那段路，半天没说话。

距蔡须坤和卢本韦最先通过，已经过去五天。

霍雨浩、薛峰先后上来了。古月娜比他们晚些，奥斯卡则是在当天上午走完最后一段，下山以后，第一件事就是将晚饭比平日多做一份。

唐三、叶阳、诺当与小舞还在下面。

尤其是诺当和小舞，越到最后，停得越久。有时同一处别人已经能站稳，他们还得回到较低的平处，重新辨认魂力与精神在压力里的接续。

蔡须坤坐在这里等过，却没有再朝下念答案。

``先喝水。''她说。

蔡兴坤接过水壶，喝到一半，忽然停住。

壶里的水晃了一下。

有一股热意，从胸口深处透出来。

她起先以为是方才登阶用力太久。可那股热并没有顺着喘息散开，反而沿体内几处熟悉的运转路线，缓缓走过一周。原先需要她费力续上的地方，被一缕极细的金光接住，随后又自行合到下一处。

蔡兴坤握住水壶的手紧了。

``你怎么了？''

蔡须坤刚问出口，自己也顿了一下。

她胸口同样亮起一点暖意。

那点光并非从蔡兴坤身上传来，而是在她体内回应。它早已留在那里，只是这几天她不断内视，也只觉出魂力在增厚，直到蔡兴坤的气息稳定下来，才终于认清了其中这一股。

两人的手背上，各有几道淡金的纹缓缓透出。

黑子鸡原本正往她袖口里钻，这时抬起脑袋，望了望她，又望了望蔡兴坤。

``先坐下。''

夏夜星空的声音从石阶旁传来。她走到两人面前，目光在金纹上停了片刻。

``坤神之力开始显现了。''

蔡兴坤抬头：``是第一考以后\ldots\ldots''

``你们已经能将它接住。''她道，``慢慢运转，别急着往外放。''

蔡兴坤依言坐好，水壶放在脚边，双手平平搭在膝上。那股金色的力量顺着她的魂力缓缓流动，走到尚不熟练的位置，便停得久些，没有因为她心急就替她冲过去。

蔡须坤闭上眼，却觉出另一种变化。

果糕领域先动了。

薄光沿着她的脚边铺开，边缘浮着极淡的灰。那是她当初压住混沌暴走、将其中一部分吸收以后形成的领域，经过这一年的修炼，收放已经比初醒时稳定许多。

如今，那点金光顺着魂力来到它的边缘。

两股力量一碰，蔡须坤肩膀微微一颤。

脚边的薄光忽然鼓起。原本安静伏在外缘的灰，像被什么从里面托起，往外散了半寸。她想先把它收回来，体内的金光却仍在向前，两处不同的节奏一时撞到了一起。

``蔡须坤。''

卢本韦已经走近。

她没有往那层光里添力，只将手伸到蔡须坤能够碰到的位置。蔡须坤握住她的手腕，听见夏夜星空让她先认清其中一股，再慢慢接另一股。

她将果糕领域的边缘收近。

灰色回到原来的位置，薄光贴着脚边伏稳。那点金色却没有随之消失，而是沿着另一个方向，自她身周缓缓展开。

金光很浅，起初只看得见一圈。

蔡须坤试着将意念放在那一圈上，它便越过脚边的碎石，向石桩延伸。再回看果糕领域，先前那一层仍然留在原处，随着她的呼吸轻轻起伏。

她把金光收回一些。

灰色的边缘没有跟着退。

蔡须坤睁开眼，怔了片刻。

卢本韦顺着她的目光看去，也看见地上两道并不重合的边界。

``两个？''

蔡须坤没有立即答，先将果糕领域收起。原先那片近乎透明的薄光缓缓隐去，金色仍在，石桩上几道细小的裂纹被照得清清楚楚。

等她重新放出果糕领域，两层光又各自留在了她想留下的位置。

夏夜星空蹲下，指尖停在金色边缘之外，没有伸进去。

``坤神领域。''她说。

蔡须坤掌心轻轻一动。

金光里传来一阵清晰的回应，与蔡兴坤体内正在运转的力量遥遥相应。她第一次发现，自己能够分辨这道回应来自哪里，又怎样从那条路线回到自己身上。

果糕领域则仍守着另一份节奏。它由吸收的混沌之力催生，灰色深处有她熟悉的牵扯，必须靠自己稳住。

两种感觉同时留在心里，谁也没有将另一种挤掉。

蔡须坤试着把它们都向外放。

脚边的细砂忽然向两侧滑开，石面上响起极轻的一声裂响。

卢本韦的手指收紧。

蔡须坤立即停住，先收回金色，再将果糕领域压到身周。光一层层退下去，她的额角已经见了汗。

``先到这里。''夏夜星空起身，``两套运转都在你手里了，怎样同时展开，还要慢慢练。''

``双领域\ldots\ldots''

霍雨浩从下面上来时，正听见蔡兴坤低声说出这三个字。

他看了看石上的两圈痕迹，又看向蔡须坤。她摊开手掌，金纹尚未完全隐去，目光却仍停在方才细砂滑开的地方。

``刚才再急一点，就收不稳。''她说。

霍雨浩点头，将她放到身侧的水递回来。

蔡兴坤坐在一旁，听见姐姐这样说，也慢慢松开了原先攥紧的手。她还没有形成领域，只能把新显现的坤神之力沿体内运转，走得不快，却已经能清楚分出每一次回应。

黑子鸡从蔡须坤怀里探出头来。蔡兴坤伸手接回去，低声道：``这回别乱动。''

鸡在她掌心啄了一下。

蔡须坤看见，终于笑出一声。

下山时，她没有再同时展开两层领域。金色收在体内，果糕领域只在需要调息时放出很薄的一线，陪她走完余下的台阶。

卢本韦一路没有说多少话。

她掌心里那份回应，却比前几日更清楚了。

有时她还没有主动召唤，母猪扽魂的气息便已经在她精神里轻轻动过；等她想仔细感受，它又安静下来，留着一种仿佛正在等她开口的停顿。

最初她以为，那只是突破以后多出来的细微感知。

可到了山脚，她将手按在门边，尚未转头，母猪扽魂却先向屋内的方向送来一点极轻的回应。

蔡须坤正在里面收起果糕领域。

卢本韦停在门口。

她没有下过那样的指令。

这几日，类似的感觉已经出现过几次。它能觉出同伴靠近，有时会先随她的心绪安静，有时又像另有自己的注意，不完全跟着她刚起的念头走。

她站了一会儿，轻轻将那道气息引出。

母猪的虚影在身后凝起。

这次，它没有一出现就随着她转向门内，而是先低下头，望了望自己踩着的地面。卢本韦回头，它才缓缓抬眼。

那一瞬，她忽然有一种极清楚的感觉。

它也在看她。

卢本韦的呼吸停了半拍。

``你一直都听得见？''

没有话语传来。

一股温厚的力量却自虚影里展开，沿着她的魂力，缓缓向外铺去。

院子里，古月娜正将手上的银光收回，忽然抬起头。

那缕银光凝住了。

并非她没有续力，而是银龙扽魂在那道气息靠近时，自身先有了一瞬收缩。她掌心一沉，抬手稳住，目光已经越过窗沿，落到门外。

奥斯卡刚在院中停下，手里凝起的补给光芒也跟着暗了一层。

``卢本韦？''

蔡须坤从屋里出来，先看见她身后的虚影，又看见一层极淡的光，正以她为中心向外扩。

卢本韦听见自己的名字，才骤然回神。

她想收，那股力量却已经沿着院子铺开。几名刚下山的队员同时停了脚，小舞本就疲惫，脸色立刻白了些，诺当扶住门框，体内尚未收尽的扽魂气息被压得一沉。

``先守住边界。''

夏夜星空落到院中，掌心一抬，将尚在外散的力量拦在院墙内。

卢本韦依言先停住扩张，再沿着那份熟悉又陌生的回应，一点点往回收。

母猪虚影望了她片刻，也随之缓缓低下头。

光终于停了。

院里几人的呼吸重新顺过来。古月娜松开手掌，银光散去，先看了一眼小舞，才走向门口。

卢本韦已经出了汗。

``刚才没收住。''她先对几人道。

小舞扶着膝盖缓了两口气，摇摇头：``先别再来一下就行。''

奥斯卡望着自己手掌，过了半天，才低声问：``领域？''

最菜萌新在院外应了一声。

老人不知何时已经到了，目光先落在卢本韦身后的虚影上，随后才看向她。

``现在，你该明白老夫当初说的那句话了。''

卢本韦慢慢点头。

``它有自己的回应。''

``继续听，别一有变化，就只想着让它照你的念头做。''

老人让她在院中坐下，等气息平稳，再将领域放出很小一圈。古月娜站在边缘，主动凝起一缕银龙扽魂的气息；奥斯卡也伸出手，只留下极弱的一点光。

这一次，卢本韦没有急着向外推。

她先接住母猪扽魂的回应，才让那圈光慢慢越过两人的脚边。

银光又沉了一分。

奥斯卡手里的光芒则明显收缩，连他自己都能清楚感觉到，武魂在那片范围内承受着一股来自品质的压制。换一口气再续，压制仍在，不能靠多添一点魂力就抹去。

古月娜望着自己的掌心，先将那一缕收回。

``从前你要专门对上来。''她说，``现在只要进了这一圈，就会受影响。''

最菜萌新点头。

``神级扽魂的品质，如今进一步体现在领域里。凡品质低于母猪的武魂，进了你的领域，都会受到压制。''

卢本韦抬眼：``所有低于它的？''

``是。''

老人看了看奥斯卡已经收起的手，又道：``但对方还有修为、经验，也还有别的手段。有了领域带来的品质压制，仗照样要自己打。''

卢本韦将光收近。

母猪虚影没有立即散去，而是安静留在她身后。她终于能在召出它以后，给那份不同于自身意志的回应留一点位置。

``猪神领域。''她轻声道。

最菜萌新没有替她换名字，只看着她又收放一回，等这一圈光完全回到掌中，才让夏夜星空撤去院墙边的力量。

蔡须坤走过来，坐到卢本韦身旁。

``你也多了一个。''

``你是两个。''

卢本韦说完，转头看了她一眼。两人都累得不轻，谁也没有再站起身试一次，只将各自尚未平静的力量慢慢收好。

晚饭比平日迟了些。

奥斯卡先给诺当和小舞盛了热的，又将卢本韦面前那只碗推近。蔡兴坤把黑子鸡护到屋角，见它不再往院子里钻，才端着自己的碗坐下。

夏夜星空等大家吃过，才将一块深色的石片放到桌上。

石片不大，表面有几道错开的细纹。她将手指按下去，屋里原本散乱的说话声便慢慢停了。

桌面上方，浮起一扇没有完全开启的门影。

门后看不见一条完整的路。

先露出来的是一段横着的石廊，廊尽头似乎还有台阶。可众人再看，那台阶竟向侧面伸去，上面挂着的断链缓缓摇摆，摇到一半，又像落进了另一处更低的空间。

远处一盏灯忽然翻过去。

小舞下意识偏头，以为自己看错了，门里的石廊却又露出另一面。方才像地面的地方，此刻悬着几只枯瘦的影子，关节紧贴石壁，无声地往前爬。

其中一只忽然停住。

它抬起头，眼睛在暗处亮了一下。

诺当手里的筷子顿住。

门影深处传来刮擦声，先是一处，随后越来越多。暗色的东西沿着裂缝涌过，有的像伏地的兽，有的只看得见细长的肢体，越往里看，越难分清它们究竟站在哪一层。

卢本韦放在桌下的手慢慢收紧。

蔡须坤望着那道门，体内刚稳定的坤神之力忽然轻轻一动。

一道声音随之落下。

``坤神第二考。''

``穿越错层地府。''

``限时一年。''

屋里很静。

直到门影缓缓收回石片，桌上的灯火重新显出原来的颜色，小舞才低声问：``里面那些，都是要打过去的？''

``沿途有大量怪物。''夏夜星空说，``它们会拦路，也会从你们未留意的层面扑出来。到时不止要有力气，还得辨清路，找到能够继续前进的位置。''

霍雨浩看着石片：``什么时候开始计时？''

``正式进入错层地府以后。一年是穿越的期限。''

她的目光从已通过的人身上掠过，又落到尚未完成第一考的四人面前。

``你们要一同进去，就得先等每个人通过第一考。还没有走到五百米的，先把自己的那一段完成。''

唐三点头，没有将手里的纸收起来。叶阳将石片让出的桌面重新抹平，给诺当留出放本子的地方。

小舞却看了看蔡须坤，又看卢本韦。

``你们早就到了。''

``早五天。''卢本韦道。

``以后还不知道要等多久。''

她说这句时声音很低，筷子也放下了。诺当低头望着碗，没有接话，方才门影里那双眼睛，仿佛还留在他面前。

蔡须坤伸手，把那块石片往桌子中央推了些。

``等会儿再想门里的事。今天你最后卡的是哪一处？''

小舞先摇头，又慢慢指向自己的本子。

``知道问题在哪，就从那里接。''蔡须坤说。

她没有说大家实力都一样。

这一年走下来，差别一直清清楚楚摆在那里。同样的矩阵层次，有人能更快认清，有人还得反复核对；同一份增长落在不同的武魂与身体里，能运转到什么地步，也各有不同。

蔡须坤与卢本韦最先到达，霍雨浩、薛峰紧跟在后，其他人却有自己的停顿。

把这份差别藏起来，第二天的石阶也不会轻一点。

``你走你的这一段。''她又道，``我们把各自该办的事办好，回来再一起走下一段。''

小舞抬起眼：``回来？''

蔡须坤还没答，最菜萌新却先看向了桌角那个旧封套。

``里面还是你们那套不提升的设想？''

古月娜将封套取来，推到桌子中间。

纸角已经磨软。最初写下的几条期望对应仍在，后来添过疑问，也划去过几处未经说明的接法，真正用于展开的规则却始终没有补全。

夏夜星空接过看了许久。

蔡须坤坐直了些。

``我们沿原版走到这里了。''她说，``可这张纸，还是不知道该怎么接下去。''

``你们先想好了哪些矩阵对应哪些序数，再希望展开能顺着这些对应走。''夏夜星空放下纸，``问题在于，希望成立的对应，并不能替你们把展开定义出来。''

古月娜指了指当年最想保住的那一处。

``若删去这里的提升，其他部分尽量不动呢？''

夏夜星空摇头。

``不能这样直接删。提升、祖先关系和生成的段，在原来的规则里彼此相关。删去一处，却照搬其他部分，有些本该降下去的位置就未必还会降。错误的改法，会让你们得到无穷降链。''

唐三抬头：``可我们展开的时候，本来就是往下。''

``往下的过程能很长，和沿同一条路永远下不完，是两回事。''

她取来空纸，写了一个较早就学过的例子。

任取正整数\allowbreak{}\NovelInline{I06139}{n}\allowbreak{}，先从\allowbreak{}\NovelInline{I06140}{\omega}\allowbreak{}落到它，再依次减到零：

\NovelDisplay{M05778R01}{%
\omega>n>n-1>\cdots>0%
}

``你能选很大的数，令这一趟走很久；但已经选定之后，这趟仍会结束。''

她又在下面写出另一串。

\NovelDisplay{M05779R01}{%
\alpha_0>\alpha_1>\alpha_2>\alpha_3>\cdots%
}

``若新规则容许同一条路这样一直走下去，每一步都宣称比前一步小，却始终没有尽头，那就不能照你们所宣称的大小关系，把这些全当成序数。''

蔡须坤望着那串向右延伸的符号，慢慢将原来的纸拉近。

``只看开头几项算得像，挡不住后面出现这种路。''

``对。你们得给出规则，还得处理规则产生的后果。''

夏夜星空将笔放下，声音仍很平。

``至于怎样定义出你们想要的这一套，我也无能为力。这里哪一步该保留、哪一步能够换，不能凭我说一句应当如此，就当它已经成立。''

古月娜望着纸，过了一会儿，轻轻点头。

蔡须坤没有再拿出一条更漂亮的期望对应。她将两张新写的说明夹在原纸后面，正要收回，最菜萌新却抬手压住了封套的一角。

``你们应当去问大长老。''

几个人同时抬头。

``他也会这些。''老人道，``你们想改的是规则，先把原版怎么展开、自己希望改成什么、还没有定义的地方，一样样拿给他看。别只带几条看着顺眼的等式过去。''

霍雨浩看向蔡须坤。

这几天，她已不止一次摸过这个旧封套。最早困在理想矩阵里的那阵沮丧虽然过去了，纸上的空白却仍在，谁都没有因为第一考通过，就将它算成一个已经解决的问题。

``大长老\ldots\ldots''

蔡须坤低声重复了一遍，原先将封套收起的手停在桌边。

最菜萌新却没有让她就此起身。

老人看向她手背上尚有余温的金纹，目光比方才沉了一些。

``不过，你们要快去快回。''

``蔡须坤，你的坤神之力已经出现了。''

屋内的风仿佛停了。

蔡须坤抬眼。

``你的第二考，只有复活牢大，才能真正通过。''

灯芯轻轻响了一声。

蔡须坤没有动。

那句话已经听清了，每个字都清楚，可落到一起时，她却像忽然找不到该怎样回应。握着封套的手慢慢攥紧，纸角在掌心里弯下去，直到卢本韦碰了碰她的手，她才骤然松开。

``复活？''

最菜萌新点头。

``穿越错层地府，是你们要走的路。对你而言，牢大也必须真正活过来，这一考才算走完。''

蔡兴坤手里的碗碰到桌面，发出很轻的一声。

奥斯卡站在灶边，原本准备收走空碗，此刻也停住了。霍雨浩望着老人，许久才问：``她现在体内的坤神之力，已经能做到这件事？''

``力量出现了，怎样用，你们还没有学会。''

老人将目光重新落回蔡须坤脸上。

``去见大长老，先把牢大如今的情形说清，再问该怎样走。魂还要养稳，别一听能复活，就急着把方才那点力量全灌进去。''

蔡须坤低下头。

她方才确实已经动了这样的念头。

胸口的金光随心绪一阵阵起伏，只要再往前引一点，似乎便能碰到那些仍温养着的、极轻的魂识。她想试，又忽然想起当年混沌翻上来时，牢大仅剩的魂怎样想替她挡在前面。

那时没有一具身体了。

每多动一点，都要从残存里支取。

蔡须坤将手按在心口，强迫自己把那股急着向前的力量慢慢收回。金光停住，果糕领域的薄光也随之平稳，身侧几人的呼吸重新清楚地落进耳中。

``我知道。''她声音很轻。

屋角这时浮起一点淡淡的魂光。

牢大的魂影没有凝得很实，轮廓在灯下显出来，仍像有一部分留在更远的暗处。他先望着蔡须坤，随后看了看桌边没有收完的纸。

蔡须坤的手从心口落下来。

她已经见过他这样出现，也听过他的声音。可最菜萌新那句话以后，再看这道薄薄的轮廓，她忽然注意到许多从前不肯一直看的地方。

桌边没有他的凳子。

奥斯卡盛饭时，仍会有一瞬想多盛一碗，随后停下。

上山缺一个人压纸的时候，她也仍会回头，回过头，才想起那只原来能伸来的手，早已不能替她按住纸角。

魂影看了她一会儿。

``先别急。''牢大说。

她眼眶忽然热了。

``这次能去问了。''

牢大没有伸手替她擦。那道魂影留在原处，声音依旧很轻。

``那就问清楚。''

蔡须坤点头，点过，又点了一次。她将封套里被自己攥弯的纸角慢慢抚平，像怕一松手，刚才听见的那句话便也会折回看不清的地方。

卢本韦将自己的I册合起，放到她旁边。

``我同你回去。''

蔡兴坤已经站了起来。

``我也去。''

她望了一眼姐姐手背上的金纹，又看向自己的掌心。这份新显现的力量，她还有太多没有弄清的地方；而牢大那道仍旧单薄的魂影，正安静留在两人面前。

奥斯卡把手里的空碗放下。

``补给我来备。路上你们别又只顾着看本子。''

古月娜最后开口。

``我同去。关于魂与身体的问题，我也要亲耳听大长老怎么说。''

蔡须坤望着几个人，一时没有接上话。

她先前想的是自己回去，越快越好。可等真要动身，才发现要带哪些纸、牢大的魂该怎样一路养稳、刚显现的两种领域又该如何收住，每一件都不能凭她此刻那一股热意去办。

最菜萌新看过五人，没有拦。

``都已经通过第一考，可以暂时离岛。记住，快去快回。''

霍雨浩一直坐在蔡须坤对面。

等众人的话停下，他才伸手，将桌边另一叠原矩阵取过来，按页次分好。

``我和薛峰留下。''

蔡须坤抬头：``你\ldots\ldots''

``这里还有四个人要继续。''霍雨浩道，``夏夜星空能护着他们，可有些问过许多遍的地方，我们知道他们各自停在哪里。你们去问大长老，我把这边能做的事做好。''

薛峰点头，将唐三与叶阳的纸分别放回两边。

``等你们回来，再一同进第二考。''

诺当的手停在膝上。

过了片刻，他将自己那张问题页拿出来，放在霍雨浩面前。

``明天，我先从这里重做。''

霍雨浩看过，将原矩阵转向他，让他自己压好纸。

小舞盯着桌面很久，才慢慢抬头。

``那你们回来时，我要是还没有走完\ldots\ldots''

``就接着走。''蔡须坤道。

``第二考要等我们全部过了第一考？''

夏夜星空点头：``你们说好一同进去，便先把人等齐。到时从正式入门起，才算那一年。''

小舞没有立刻再说话。唐三将水递来，等她喝过，自己坐到旁边，重新看她今晚指出的那处下标。

古月娜把未吃完的补给推到她面前。

``明早先吃饭，再上山。''

小舞望着她，终于点了点头。

当天夜里，屋里一直亮着灯。

蔡兴坤先去问了返程的船，又回来同奥斯卡分装路上的东西。她原想把这段时日带上山的整包笔记全塞进去，放到一半，自己停住，将仍需回查与要请教大长老的另外取出。

黑子鸡蹲在包边，几次想踩进去，被她轻轻托回干燥的地方。

古月娜检查水袋与干衣，又去同夏夜星空说了几句话。她回来以后，没有先催蔡须坤收拾，只将一只留好的空包放到床边。

卢本韦坐在另一头，母猪扽魂的虚影已经收起。那份回应却仍在，她每将心绪稳住一次，便能更清楚感到对方也安静下来。

她没有拿这份新感知去反复试自己的领域，只将猪神领域初醒时尚不稳定的边界记下，夹进准备带回去的册子里。

霍雨浩将最后几页交给蔡须坤。

最上面是原版BMS实际怎样展开，后面才夹着理想不提升BMS的期望对应。哪些尚未定义，哪些曾经只是猜想，他都留着原来的问号。

``这张也带着。''他又抽出一页。

是夏夜星空今晚写的两串下降关系。

蔡须坤接过来，看了一会儿，将它放在封套最后。

``我回来以后，把大长老说的从头讲给你们听。''

霍雨浩点头。

``别只记他说的结果。你真正没听懂的，也记下来。''

蔡须坤望着他，忽然笑了一下。

``知道。''

她又低头，去看另一张还没有落笔的纸。

上面要写的，是牢大如今的情况。她写下魂影仍能短暂出现，又写每次出现以后需要怎样温养，写到当初那具身体时，笔停了许久。

卢本韦没有从旁边替她接下去。

蔡须坤缓过那一阵，才把自己真正知道的与还需要询问的分开，慢慢写完。最后一笔落下，她将纸放在最外层，等墨干了，才收入封套。

牢大的魂影没有再凝出来。

那一点极轻的魂识安静养着，偶尔有一丝回应，足以让她知道仍在。她没有将新得的坤神之力灌过去，只把运转稳住，照往常那样留下温养的余地。

这一夜，她终于没有再试着多算一条矩阵。

第二天清晨，五人出门时，山阶上还留着雾。

诺当站在门边，将奥斯卡忘在桌角的一包东西递给他。奥斯卡接过，正要像平日那样说一句记得吃饭，却发现给留下的人备好的那份，诺当已经分开放在屋里。

他笑了笑，拍过包口。

``回来再换我做。''

叶阳帮蔡兴坤提到门口，等她把包背稳才松手。唐三将小舞尚未收好的长页压住，小舞却先走过来，替古月娜将一条垂下的系带重新扣好。

``别挂到路边。''她说。

古月娜低头看了一眼，点头。

霍雨浩与薛峰送到山路转弯处。夏夜星空留在石阶边，等五人回头，又向蔡须坤提醒了一遍：领域先稳住，路上若有变化便停下，不必为了赶路硬接。

蔡须坤应下，将封套往怀里收近。

最菜萌新站在更前面，直到他们走近，才将目光从海边收回。

``去吧。''

老人看向蔡须坤，没有再将复活那句话重复一遍，只让她把该问的都带到。

她点头。

五百米的石桩藏在雾里，已经看不见了。更高处仍有长路，而错层地府那扇尚未真正开启的门，还留在等他们归来的地方。

蔡兴坤先走下转弯处，回头等姐姐。奥斯卡背着补给，卢本韦将自己的册子收好，古月娜走在外侧，替几人看着湿滑的路沿。

蔡须坤最后回望了一眼。

霍雨浩抬起手，薛峰站在他身旁。再后面，小舞已经拿起今日要核的原式，唐三等她说完，诺当与叶阳也各自取回自己的那一页。

她将手放下，跟上前面的四人。

海风迎面吹来，怀里的封套被她压得很稳。最外面那张纸上，写着牢大的名字。

\chapter{第一百三十一章：归航}\label{ux7b2cux4e00ux767eux4e09ux5341ux4e00ux7ae0ux5f52ux822a}

五人走到海湾时，返程的船已经靠在岸边。

船比他们来时坐的舢板宽些，后半截支着低矮的篷，篷下堆着捆好的空筐。岛上操舟的人替他们挪出一块干燥地方，指了指船头，让行李先往那里放。

奥斯卡先上去试了试落脚的位置，回身接包。蔡兴坤将黑子鸡护在怀里，一只手搭住船舷，等船身晃过那一下，才迈了进去。

``坐稳了再松手。''奥斯卡说。

黑子鸡倒比她稳。它从她臂弯里探出头，朝岸上望了一眼，又缩了回去。

卢本韦最后核过一遍东西，将装着笔记的包塞到篷下。原版BMS的推演、理想不提升BMS的草稿，连同那些还没有写出答案的问题，都已分开收好。蔡须坤怀里还有一份封套，最外面夹着牢大如今的情形。

她上船后，先按了按封口。

纸没有受潮。

远处的山腰仍挂着雾。唐三、小舞、诺当与叶阳，今天还要继续走各自没有走完的第一考。霍雨浩与薛峰留下，陪他们核对推演，也照看上山后的进退。

蔡须坤朝山的方向看了一会儿。

``他们这时候该吃完了。''奥斯卡将水袋递给她，``小舞昨晚还说，今天要早点上去。''

``嗯。''

她接过水袋，没有再站在船边。

缆绳收起，船身缓缓离开岸石。最初几下还晃得厉害，到了稍深的水里，起伏反倒匀了。桨叶划开海面，一条窄窄的白痕留在船后，将岸上的鸡鸣声越拉越远。

蔡兴坤将背包往脚边挪了挪，给姐姐留出位置。

``等回去见了大长老，我先把纸摊好。''她说，``省得又像上次，一着急，拿出来的全混在一起。''

卢本韦在旁边应了一声。

``这次先把没定义的地方找给他看。''

蔡须坤点头。那张写着``展开第一项''的纸也在里面，后面仍是空的。她们沿原版BMS走过了一年的山阶，心中那套理想矩阵，却还缺少真正能用来展开的规则。

这一次回去，就是要把这些空白拿到大长老面前，请他指点。

船离岛岸更远了一些。

古月娜一直坐在靠外的位置，银发束在身后。风将几缕散发吹到脸侧，她抬手拢过，没有立刻回到篷下。

过了一会儿，她看向蔡须坤。

``来这边一下。''

蔡须坤将水袋交还奥斯卡，走到她身旁。

船边有一截系缆的木桩，两人隔着它站定。篷下，蔡兴坤正将黑子鸡从包口抱出来，奥斯卡低头整理补给；卢本韦看见她们走远，抬了一下头，又将包上的结重新系紧，没有跟来。

古月娜望着船后渐渐模糊的山影。

``族长当年，其实也后悔过。''

蔡须坤的手停在船舷上。

古月娜说起帝天时，声音比在岛上争论天州旧事时低了许多。

``那一战以后，星斗大森林万年来积攒的底蕴，消耗了大半。''

海浪撞过船侧，船身往上抬了一下。古月娜扶住木桩，等它重新落稳，才继续说。

``你也知道，那些魂兽不是几年、几十年就能长起来的。各族失去的强者，空下来的守地，还有跟着兽潮出去、再也没回来的同族\ldots\ldots 人回到森林，缺的那些却补不回来。''

蔡须坤没有接话。

她记得兽潮出林时的声音。脚步连着脚步，从森林深处一直压向外面，强者的气息在前方展开，后面的兽群便跟着往前走。那时看过去，仿佛整片林海都能一直这样推到远处。

后来回去的，少了许多。

古月娜垂下眼。

``他养伤的那些年，我听过他提起这件事。说攒下一万年的东西，不能再这样耗下去。''

她顿了顿。

``可他还是得坐在那里。''

蔡须坤转过头，看向她。

``龙族和其他族群，本来就有许多恩怨。''古月娜说，``争过领地，争过生息的地方，有的族群彼此杀过。那些事早在天州以前就有，不会因为大家都叫魂兽，就全忘了。''

``我知道。''蔡须坤轻声道。

她在帝天身边时，听过各族为林界与栖息地争执。有的事压下去几年，再来时，旧怨仍在，连告状的都已换成上一代留下的后辈。

``他能叫它们先停下来。''古月娜说，``有时候不服，也得先停。不然，一处打起来，另一处就趁机去抢，弱的族群往别处逃，又会撞进别人的领地。''

她握在木桩上的手指慢慢收紧。

``如果不是他维持着，整个魂兽族群，只会陷进没完没了的战争里。''

蔡须坤看着她，没有打断。

古月娜这才将一直压在心里的话说完。

``我知道他那次做了什么，也知道城里死了人。可若只说他不该出兵，不该去天州\ldots\ldots 森林里那些原本靠他守住的东西，也都还在。''

``他后悔，也不能就此把它们丢下。''

风从两人中间穿过去。

蔡须坤低头看了看扶在船舷上的手。掌背的金纹已经收敛，体内新显现的坤神之力沿熟悉的路线缓缓运行，牢大那一点魂息仍安静地受着温养。

她想了一会儿，才开口。

``我没忘记他守过我们。''

古月娜抬眼。

``那些争执，我也听过。每回到最后，都要等他来压住。''蔡须坤说，``他能让双方停手，可有些事停过以后，还会再来。''

``总不能不管。''

``当然要管。''

蔡须坤望向篷下。卢本韦正将一本册子压在膝上，腾出手替蔡兴坤接住快滑下去的外衣。母猪扽魂没有放出来，她偶尔停下动作时，却像仍能听见什么回应。

``所以，我们才要把魂灵这条路继续走下去。''

古月娜顺着她的目光看了一眼。

``大长老对这件事，怀疑了很多年。''

``他怕再死人。''蔡须坤说。

这句话出口，她自己也静了一下。卢守衡的那场事故，最初的几天同样平稳。直到后来，双方都没能停住那股失衡的力量。

她们如今带回去的，不能只有一句成功了。

``我们和魂灵怎样相处，力量不稳的时候怎样收住，还有哪些地方做不到，都要让他看见。''她继续道，``能活下来，也能一直一起活下去，才算把这条路往前走了一步。''

古月娜低声问：``你觉得，他会愿意帮我们？''

``会。''

蔡须坤答得很轻，却没有犹豫。

``大长老殿下已经看见我们的魂灵能成功了。他会看好我们的。''

船篷轻轻响了一声。她将被风吹到脸边的发丝拨开，又按了按怀里的封套。

``他当年也想过去森林里谈共存。现在我们真有了能拿出来给他看的东西，我想，他会愿意继续看我们怎样往下走。''

古月娜没有立即回答。

``人类会和魂兽和解的。''蔡须坤说，``有些地方已经开始了。狼王肯将自己的意愿告诉我们，雨浩身边的魂灵也一直陪着他。我们把能做成的留下，再把没有做好的一处处补上。''

她望着船侧起伏的海水，慢慢将后面的话说清。

``森林要养得住魂兽，人也要有能活下去的田地。不能每次缺了什么，就去抢另一边的；也不能总等到打起来，再让一个最强的人出来收场。''

``最后，让人类和魂兽一起维持生态平衡。''

古月娜看着她。

过了许久，她才问：``各族之间的旧怨呢？''

``也得慢慢谈。''蔡须坤说，``族长一直做的那些事，不会因为有了魂灵，就全不用做了。可若人类和魂兽能多一条共生的路，少一回猎杀、少一回报复，他也能少接下一桩由这些事结成的仇。''

古月娜握着木桩的手，渐渐松开了。

船后，蓬莱岛的轮廓只剩下很淡的一片。她又望了一眼，回过头时，目光落在蔡须坤怀里的封套上。

``理想BMS，连第一项还没写出来。''

``所以要回去请教他。''

``别到时候一着急，只顾着问你最想听的那几句。''

蔡须坤看向她，嘴角终于松了一点。

``你也在。你听见我漏了，就替我问。''

古月娜点头。

``我会问的。''

篷下，奥斯卡举起水袋，朝她们晃了晃。

``说完了先回来坐。站那儿一晃一晃的，看得我都想伸手扶。''

蔡须坤应了一声。古月娜先松开木桩，等她站稳，才与她一同往船篷走。

卢本韦将膝上的包往里挪，给两人空出位置。蔡须坤坐下，取出那个旧封套，古月娜就在她身旁伸手，替她按住了被海风掀起的纸角。

\begin{center}\rule{0.5\linewidth}{0.5pt}\end{center}

第四天清晨，船靠上了大陆的码头。

岸上有人扛着鱼筐来回走，木板踩得咚咚响。卖热汤的摊子支在坡下，白气从锅里往上冒，一个挑担的汉子嫌价钱贵，站在摊前与摊主争了好一会儿，最后还是掏出铜钱，端走一碗。

奥斯卡站在船边听着，笑了一下。

``可算回来了。''

蔡兴坤下船时，黑子鸡先从她怀里跳到岸上。它抖过翅膀，踩了两步，才回头等她。卢本韦接过最后一个包，重新背好，五人沿着通向内陆的路离开码头。

最初一段路上，往来的人还不少。越过海边的低坡，鱼腥气渐渐淡去，两旁换成田地。快到晌午时，他们看见一处村庄。

村子不大，几十户人家，屋顶上还飘着炊烟。一条引水的小渠从村边流过，渠旁的菜畦刚浇过水，泥土湿得发黑。

``先去问个路。''奥斯卡说，``水也该添了。''

蔡须坤点头，顺着田埂往村口走。

她先看见一只倒扣的竹筐。

筐边滚着几个青瓜，一个断了，断口朝上，里面的籽还亮着。旁边放着扁担，两只绳套摊在地上，像挑筐的人刚弯下腰，还没来得及将它扶起。

村口没有人。

奥斯卡朝最近的院门喊了一声。

``有人吗？我们过路，想讨些水。''

院里传来一点响动。

他刚要过去，古月娜便伸手拦了一下。

一条黄狗从门后出来，嘴边还沾着水。它没有冲他们叫，走到门槛旁便停住了，扭头望向屋内，喉咙里发出低低的呜声。

屋里没有人应。

卢本韦绕开地上的青瓜，走到能看见院内的位置。

``先别进去。''

院里搭着晾衣绳，一只木盆搁在绳下，半件湿衣还挂在盆沿。堂屋两扇门都开着，桌上摆了三副碗筷，其中一只碗里盛着小半碗饭，筷子横在碗口。

灶房那边，水正往外溢。

嗤的一声，落在烧红的柴上，又是一声。

蔡须坤等了片刻，再次开口。

``家里有人吗？''

只有那只狗又朝屋里呜了一声。

古月娜的目光沿院墙移向隔壁。那边的门也开着，靠墙摆着一张矮凳，凳面上铺着鞋底，针已经穿过去，麻线拖下来，一直落到地上。

做活的人不在。

``旁边呢？''蔡兴坤问。

没人立即回答。

五个人沿着村里的土路往前走，没有分开。奥斯卡喊过几次，后来自己也停了。

磨坊边的水轮仍在转，磨下的粉落进布袋，袋口已经漫出来了。路旁有一辆独轮车，车轮陷在浅坑里，车上两袋粮食用绳捆得结实。

有人走到这里，便将车留在了这里。

再往前，一家门口摆着孩子的泥玩具，泥还没干透，一只小手印清清楚楚按在上面。旁边一串晒着的红果，被风吹得碰在墙上，轻轻响了一下。

蔡兴坤抱紧黑子鸡。

``他们都去哪儿了？''

蔡须坤停在路口。她已经能看见村子另一端，那里没有聚集的人群，也没有通往别处的热闹声。

炊烟、流水、风吹过衣裳的声音，全都还在。

唯独没有一个人。

她走到一户院门外，将果糕领域缓缓放开，薄光只铺到五人身边，没有继续漫进屋内。与此同时，古月娜收拢心神，将一缕精神力沿近旁的院落探了过去。

几息后，她睁开眼。

``这些屋里，没有人的回应。''

卢本韦的手已经按在包带上。

``帅气教？''

这三个字让几个人同时静下来。

他们见过帅气教留下的东西，也听大长老讲过那些被抽干的尸身。一路上村庄还在，饭菜还在，人却全没有了，谁都无法轻易绕开这个念头。

古月娜没有立即点头。

``先看清。''

蔡须坤望向身旁几人。

``一起查。别单独进屋。''

他们先回到那座灶上仍烧着水的院子。黄狗伏在堂屋门边，见他们过来，只往旁边挪了挪，又将鼻子凑近门槛，好像还在等屋里的人出来。

奥斯卡停在院中，看着灶下的柴。

最外头一根只烧黑了一端，里面的火却旺，显然添进去没多久。水溢了许久，柴边湿成一片，火仍在往上舔。

他伸出脚，将那根快要滚出来的柴拨回灶内，免得火落到旁边的草篮上。

蔡须坤站在堂屋门外，往里面看。

桌下有只小布鞋，另一只放在门边，鞋口朝着屋里。柜门关着，靠墙的竹篮里压着钱袋。人若是收拾东西走了，至少不会连这一双鞋都没有穿齐。

卢本韦看过地面，又看窗下。

``没有打过。''

桌凳没有翻，门窗没有破，连墙边靠着的扫帚都好好立着。他们一路找来，也没有看见血迹或拖曳留下的痕迹。

``也可能被直接带走。''奥斯卡道。

``嗯。''蔡须坤应了一声，``现在还不能定。''

她没有碰桌上的饭碗。古月娜却忽然向前半步，目光落在那只碗旁边。

``这里。''

桌面上覆着一缕极淡的东西。

日光从窗里照进来，原先只看得见木纹。古月娜将掌心的一点银光放低，那缕痕迹才从木纹之间显出些许轮廓，薄得几乎没有颜色。

蔡兴坤也凑近了一点，却没跨进门。

``是什么？''

古月娜没有作答。她让银光靠近，尚未真正碰上，光的边缘便散开了一小片。

她立刻收手。

那道痕迹仍留在原处。

没有爆响，没有向外荡开的气浪，屋里的碗筷甚至没有动一下。古月娜盯着掌心，眉头慢慢蹙起。

``你受伤了？''蔡须坤问。

``没有。''

她重新凝起那一点银光，确认能顺着自己的心念收放，才将手垂下。

``可我没有分辨出它怎么起作用。''

卢本韦望向桌面，母猪扽魂的一点气息已在她身侧浮起。她没有让它向屋里压，只护着近旁几人，目光却比方才更沉。

``我没有感觉到被抽取。''

她等了一会儿，才将这句话补全。

``从进村到现在，都没有。''

``剩在这里的力量太淡，也可能已经做完了。''蔡须坤说。

她知道，不能只凭他们没有受影响，就断言这里安全。

古月娜却缓缓摇头。

``这道痕迹，和我在天州见过的不同。''

她退回院内，沿门框看了一遍，又停在灶房外。

同样的薄痕，不只留在饭桌上。

门框的木缝、盛水的缸沿，连院中石磨背着日光的一面，都有一点。它们太淡，刚进来时根本察觉不到；如今用极少的力量照过去，才能分辨出它们各自停留的地方。

蔡须坤也看见了。

那不像从某个人身上抽出力量，向一个地方聚去。眼前这些没有生命的东西上，也留着相同的残迹，散在整个院子里。

她将领域收近一点，先稳住自己的运转，没有试着覆盖那些痕迹。

``隔壁也看一眼。''

五人一同退到路上。

隔壁那张铺着鞋底的矮凳下，同样留着一点。

再远些，独轮车旁边的地面也有。

古月娜在第三处停下，终于抬起头。

``至少，不是我们见过的那种手段。''

卢本韦顺着土路望向村口。

``如果是另一种呢？帅气教也未必只会以前那些。''

``所以还要查。''古月娜说，``可这里确实有一股我们没见过的力量。它留下的东西，和抽取帅气、剥离魂灵，都对不上。''

蔡须坤没有立刻接话。

她想起灶房里添进去不久的柴，桌上的半碗饭，还有那只没有穿上的鞋。一路从村口过来，没有人来得及收拾，也看不出仓促奔逃的踪迹。

仿佛前一刻，这里还过着寻常的日子。

有人挑菜回来，有人坐在门边纳鞋底，孩子在路旁捏泥玩具。灶里的水开了，屋里的人端起碗，正要招呼家人吃饭。

随后，所有人的声音一起断了。

奥斯卡背上的补给还很沉。他伸手扶了扶包，忽然觉得刚才在码头听到的争价声，已经远得像隔了很久。

``先去井边。''蔡须坤说。

村中的井就在路尽头。井旁放着两只水桶，其中一只已经装满，另一只仍系在绳上，垂在井里。

打水的人也不在。

蔡兴坤走到姐姐身侧，黑子鸡缩在她臂弯里，这一次没有再探头去啄衣上的绳结。

``村里这些人，会不会\ldots\ldots''

她没有将后面的话说完。

蔡须坤望着井沿，没有立刻给她一个安心的答案。

``还没有找到。''她说，``先找。''

古月娜在井边蹲下，手没有落到石头上。

背着光的井沿下，又有那种极淡的痕迹。

比方才院子里的，更清楚一些。

她向另外四人抬了抬手，示意先停住。蔡须坤站到她身侧，卢本韦与奥斯卡也走近了些。

井里传来水滴落下的声音。

一声。

又一声。

五个人都站在井旁，却没有谁贸然向井底探去。

蔡须坤往井中看了一眼。

垂下的井绳旁，竟有一张女人的脸。

那张脸仰着，正对井口。两颊僵硬地绷紧，嘴角却向两侧咧开，露出一口森白的牙。湿发贴在额上，遮住半边眉眼，剩下的那只眼睛一动不动地盯着她。

蔡须坤的目光刚落下去，那只眼睛便转了转。

它看见她了。

蔡须坤立刻退开井沿。

``井里！''

话音未落，五人的魂力几乎同时腾起。

篮球在蔡须坤掌中凝实。她侧身让开井口正前方，果糕领域随之收紧，将另外四人护在身边。蔡兴坤松开臂弯，黑子鸡跃到她身前，双翅猛地张开，方才伏在怀里的温顺模样顿时消失。

卢本韦向前跨了半步，母猪扽魂在她身侧显现，四蹄重重踏住地面，鼻端喷出一口粗气。另一侧，银龙的虚影在古月娜身后舒展开来，银光掠过井沿，将石缝里尚未干透的水迹照得一清二楚。

奥斯卡退到四人中间，掌心已经多了几根刚刚凝成的香肠。他迅速分到同伴手里，自己也握住一根，目光始终没有离开井口。

没有人再凑过去查看。

井下却响起了笑声。

声音很低，贴着狭窄的井壁往上爬，听不清究竟离他们有多远。系着水桶的绳子忽然绷直，绞盘轻轻响了一声。

随后，那张脸缓缓升了上来。

先是湿发，再是额头和眼睛。女人的下巴越过井沿，脖颈、肩膀也一点点露出。她没有攀住井绳，双手垂在身侧，整个人就这么从井中浮起。

水珠顺着她的头发落下，滴在井沿上。

蔡兴坤下意识地攥紧了手里那根香肠，目光扫过女人垂下的双手，又立刻转回她脸上。

方才从上往下看，已觉得那张脸狰狞。如今离开井里的暗处，五人才看清，她脸上那副笑容几乎没有变过，眼珠却在他们之间逐个转动。

转到谁身上，便停一停。

最后，停在蔡须坤脸上。

``既然来了，就留下吧。''

女人的嘴唇动得很慢，声音却清楚地落进每个人耳中。

蔡须坤托住篮球，余光确认着同伴的位置。古月娜站在她斜后方，蔡兴坤守住另一侧，奥斯卡已经退到两人都能照应到的地方。

``村里的人在哪里？''她问。

女人仍旧看着她，嘴角一点点抬高。

井绳上的水桶撞了一下井壁。

空荡荡的村子里，那一声格外清楚。身后的屋门还开着，灶房的烟从瓦缝间慢慢升上去，却始终没有人应声。

卢本韦的脸色沉了下来。

她一路看着那些没有收拾的饭菜、落在地上的鞋，还有站在空屋门前的狗，此刻再听见这句轻飘飘的``留下''，胸口那股火便压不住了。

母猪扽魂随她向前挪了一步，挡住女人望向奥斯卡的视线。

``你算哪根葱，不服就来碰一下。''

这一次，井边那阵低低的笑声停了。

\chapter{第一百三十二章：破镜}\label{ux7b2cux4e00ux767eux4e09ux5341ux4e8cux7ae0ux7834ux955c}

女人脸上的笑容僵住了。

下一刻，她的身形骤然拔高，湿透的衣摆从井口抽出，在半空中拖下一串水珠。她抬起两只手，指尖相互一扣，井沿四周便响起了接连不断的轻响。

一面镜子从她身后浮起。

随后是第二面、第三面。

大大小小的镜面沿着井台向外展开，越过屋檐，悬到五人头顶。镜框边缘泛着暗光，有的只有巴掌大小，有的足有一扇门高。每一面都朝着不同的方向，偏偏都能照见他们。

古月娜侧身挡在奥斯卡旁边。

她身后的银龙才一摆尾，镜中的银龙也跟着动了，却比她更早将尾巴甩向同伴。

蔡兴坤眼前的路随之折过去。

明明只隔着两步，蔡须坤却出现在了另一座院子的门口。卢本韦的身影被几面镜子分开，一半站在井边，另一半竟已走到她身后。

四周响起脚步声。

那些敞着门的空屋里，像是同时有人往外走。她瞥见一截衣角，又看见有人端着碗，从镜面的边缘一闪而过。

蔡须坤手里的篮球微微一沉。

她握紧它，脚下仍踩着井边那块有缺口的石头。方才退开的距离没有变，奥斯卡递给她的香肠也还留在另一只手中。

``破阵。''

她没有去追逐镜中的人影。

由混沌之力催生的果糕领域猛地展开，贴着地面掠过井台；与此同时，一圈金光自她身周升起，坤神领域沿着灰色薄光的边缘向外铺开。

两股力量在她掌控下同时释放。

灰光先触到镜面，金光紧随其后。

头顶那片被折叠的天空骤然一颤。

一声脆响传来，所有错乱的倒影同时裂开。院门前的蔡须坤消失了，走到蔡兴坤身后的卢本韦也消失了。镜中那些端碗、提篮的人影还没有走出半步，就被金光从中穿过，散成一层薄薄的暗雾。

幻境瞬间破碎。

蔡兴坤重新看见姐姐就在自己身侧。古月娜的银龙仍守着原来的位置，尾尖根本没有扫向奥斯卡。

井台、瓦房、路边的水桶，一样样回到眼前。

远处灶房的烟仍在升。

空屋里没有人出来。

女人悬在井口上方，两条手臂猛地向内收拢。散到屋檐上的镜子接连崩碎，只剩几面较小的，仓促聚回她身边。

她第一次低头，看向了蔡须坤脚下那两层领域。

蔡须坤也在这一刻缓了口气。

同时展开双领域，魂力沿两条路线迅速流动，肩背已经有些发紧。她将范围收回井台附近，让灰光稳住边缘，金光留在自己和同伴之间。

卢本韦抬眼扫过剩下的镜子。

``现在就是母猪主场。''

母猪扽魂应声向前，鼻端喷出的气息吹得井旁尘土一滚。

女人十指一张，护在面前的镜子连成一片。随着力量重新汇聚，她颈侧与手臂上的深色纹路渐渐显现，几处原本藏在衣袖下的纹路也亮了起来，在身前映出一列清晰的矩阵。

\NovelDisplay{M05780R01}{%
(0)(1,1,1)(2,2,2)%
}

卢本韦看清以后，嘴角一挑。

``这点实力也敢班门弄斧。''

话说得轻，她的手却已经抬起来，指向女人身前那片镜面。

魂力沿着她的手臂涌出。母猪扽魂的背脊随之隆起，四肢骤然粗壮，原本只到她腰间的身躯迅速膨胀，转眼便比井台高出一大截。

蔡兴坤向侧方让开，黑子鸡振翅掠到半空。古月娜则将银光铺向两旁，护住紧邻井台的院墙和屋门。

卢本韦向前一压手掌。

母猪扽魂低下头，前蹄踏实，肩背间聚起一团厚重的光，迎着那片镜面撞了过去。

女人没有闪避。

她两只手掌抵住镜背，身上纹路陡然亮起。

第一轮力量碰撞，响声比五人预想的还要沉。

井绳猛地一甩，空水桶撞上内壁，绞盘上的木柄折断半截。镜面向内凹下，几乎贴到女人胸前。她的肩膀被压得一沉，脚下的水珠尽数震散。

卢本韦刚要再添一分力，眼神忽然变了。

凹下去的镜面深处，竟浮出同样的一团光。

连力量凝聚时的起伏，都与母猪扽魂方才发出的相近。

``退开！''

卢本韦猛地收臂，母猪扽魂四蹄向下一沉。

镜中那团光已经反冲出来。

它沿着方才的方向倒卷而回，重重撞在母猪扽魂胸前。巨大的身躯晃了一下，前蹄在地上拖出两条深痕，碎石一路翻到卢本韦脚边。

古月娜将护在两旁的银光向内一合，拦下散向院墙的余劲。蔡兴坤拉住奥斯卡退了半步，最后一块石头擦过黑子鸡的翅尖，砸在空地上。

卢本韦站稳，手指仍抵在母猪扽魂背后。

``还能把力量反射回来？''

女人身前的镜子重新挺直。

可镜面中央，已经有了一道细长的裂纹。

她藏在镜后的手也在发抖。指缝里淌下的水珠掺进一点红色，顺着镜框落了下来。颈侧刚刚亮起的纹路暗了一截，又被她强行催亮。

奥斯卡一直盯着她。

那道裂纹没有恢复。他看见女人想将身旁另一面小镜子拉过来，刚挪到一半，镜中的光便闪了两次。

``她的力量是有限的，我们一鼓作气！''

他说着将手里的香肠送进口中，迅速运转魂力，准备接应。

蔡须坤已经动了。

她从母猪扽魂侧旁穿过，脚下金光随步伐前移。果糕领域压住身边翻涌的余劲，坤神领域则贴着她的身体，越过井台旁的碎石，直逼女人面前。

卢本韦立即明白了她的意思。

母猪扽魂再度向前逼近，厚重的力量压住镜面正中，将女人的两只手牢牢牵在原处。

半空中的黑子鸡也扑了下来。

翅膀擦过一面试图转向蔡须坤的小镜子，镜框歪向一旁。蔡兴坤抬手一引，黑子鸡随即翻身，挡在女人想要退开的方向。

蔡须坤就在这一瞬欺近。

篮球砸向女人持镜的手腕。

女人猛地扭身，右手勉强从镜背抽出，五指直抓蔡须坤的肩头。她袖口里涌出浓重的黑气，带着一阵细碎的哭声，瞬间裹住了整条手臂。

几张模糊的脸从黑气里挤出来。

接着是手，是伸向蔡须坤脖颈的指爪。

阴冷的气息扑面压来，蔡须坤耳中骤然多出许多声音。有的在哭，有的在叫她回头，还有一道紧贴耳畔的声音，像是已经趴到了她背上。

她的后颈泛起一层寒意。

金光随即亮了。

坤神之力沿着肩背涌上，覆盖住她挥出的手臂。那些探进领域的黑色指爪刚碰到光，便剧烈蜷缩，附在她耳边的哭声也陡然尖了起来。

女人脸上的笑意彻底消失。

她想将那股冤魂之力收回，却已来不及。金光顺着相接的地方压过去，裹在袖口外的黑气一层层散开，几张模糊的脸跟着松脱，原本攥住蔡须坤衣领的一只黑手，也在她眼前失了形。

那股阴冷的力量撞上坤神之力，竟连片刻都支撑不住。

蔡须坤手中的篮球重重击中女人手腕。

五指顿时一松。

女人向后缩身，还想调动身侧的小镜子。古月娜早已守在那里，一道银光横穿过去，将才转到一半的镜面推偏。

蔡须坤踏上井沿，顺势抓住女人失力的右臂，身体向前一压，带着她越过井口，重重落到另一侧的空地上。

女人左手一翻，掌心再次涌出黑气。

这一次，蔡兴坤手背上的淡金纹路也亮了。

黑子鸡振翅落下，带起的金光擦着蔡须坤身侧扫过，将那股尚未聚拢的冤魂之力当场冲散。

女人猛地一颤。

蔡须坤抓紧这个空隙，膝盖抵住她的手臂，将她彻底压在地上。坤神领域随之收拢，金光贴住她的肩背，每有一点黑气浮起，便立刻将其压回。

身后传来一声脆响。

那面挡住母猪扽魂的镜子，终于从裂纹处崩开。

卢本韦及时偏转手掌，母猪扽魂将余下的力量压进井旁空地。土层向下陷了一块，碎掉的镜光从女人头顶落下，还没碰到地面就消失了。

古月娜上前，以银色魂力缠住女人的双腕，将两只手一并锁在背后。蔡须坤确认她无法再结出先前的手势，才稍稍减轻膝下的力道。

从蔡须坤冲出，到最后一面镜子破碎，不过短短片刻。

女人已经被生擒。

奥斯卡走到井台边，先看过几名同伴，确认都没有受伤，才将一根恢复香肠递给蔡须坤。

她接过去吃下，沿着领域运转了一周，肩背间紧绷的感觉终于缓了些。

卢本韦也收住了母猪扽魂继续膨胀的力量。庞大的身躯渐渐缩回原样，仍守在女人旁边，鼻尖离她垂下的湿发只有半尺。

``刚才不是还要留我们吗？''卢本韦低头看着她，``现在能好好说话了？''

女人贴着地面喘息，眼睛仍在几人脚边来回转动。

蔡兴坤向井里看了一眼。绳子还在晃，掉下去的木柄浮在水面上，除此以外，看不见别的东西。

院墙内的门，依旧开着。

远处那只黄狗不知什么时候走到了路口，停在那里，鼻子朝着风来的方向轻轻抽动，却不肯再靠近。

蔡须坤蹲下身。

她将女人额前的湿发拨到一旁，让那张脸彻底露在日光里。坤神之力仍压在对方身上，手腕处的银色束缚没有半点松动。

``看着我。''

女人的眼珠终于停住。

蔡须坤盯着她，一字一字地问：

``村里的人，被你弄到哪里去了？''

女人抿住了嘴。

她试着动了一下手腕，银色束缚随之收紧。肩背上的金光仍旧压着她，那点刚从指间凝起的黑气还未成形，便散在地上。

蔡须坤等了片刻。

``谁让你来的？''

女人终于抬起眼。

``我是宏舒集团的人。''

她说出这个名字时，声音比方才稳了一点，目光也重新扫过五人的脸，像是在等他们有什么反应。

卢本韦皱了皱眉。

``宏舒集团？''

``你们没见过我，很正常。''女人说，``我平日很少在外面露面。可宏舒集团的事，你们最好少管。''

母猪扽魂往前拱了一下。

女人的视线顿时被一只凑近的猪鼻子挡住，后面的话也停在嘴边。

``现在是我们在问你。''卢本韦说。

蔡须坤将她身上的金光稳住，示意卢本韦先让她继续说。

古月娜站到侧面，低头看着她。

``你们和帅气教，有什么关系？''

女人眼皮一跳。

这一下细小的变化，没有逃过几人的眼睛。

``来往而已。''她答得很快。

``什么来往？''

女人不再看古月娜，目光落到地上的一片水迹里。银色束缚没有松，井口上方也没有再出现镜光。

过了好一会儿，她才继续开口。

宏舒集团与帅气教的往来，远比她最初那句轻描淡写的回答频繁。双方长期保持联系，许多事情并不经过斗罗大陆上公开的渠道。

而宏舒集团的最高统治者，名叫王雨新。

``她是胡赞的妻子。''

蔡兴坤猛地抬起头。

``胡赞？''

女人没有改口。

井旁安静了一瞬。

先前检查村庄时，她们发现的力量残迹，与熟悉的帅气教手段并不相同。可谁也没想到，追查这股陌生的力量，竟会问出这样一层关系。

卢本韦转头与蔡须坤对视了一眼。

``难怪能来往得这么勤。''

蔡须坤没有接这句话。

她想起天州城外的兽潮，想起那些被抽走帅气的人，也想起她们此行原本要带回去的几页疑问。海上航行时，她还在反复琢磨该怎样向大长老请教，如今刚回大陆，手里便又多了一件必须告诉他的事。

``宏舒集团在哪里？''她问。

``星辰大陆。''

这个名字，让奥斯卡原本准备记下来的手停住了。

``另一个大陆？''

女人点了一下头。

``从斗罗大陆西岸出海，继续往西，大约三千里，就能到。''

奥斯卡看向古月娜，随后又看回女人。

他们刚从蓬莱岛渡海回来，对海上的距离已经有了切身感受。三千里的航程绝非抬脚便到，可也远没有遥不可及。

``这么近，为什么一直没有人探索？''蔡兴坤问。

``那里只有斗罗大陆的十分之一大。''

女人说，星辰大陆在海上长期不受重视。对于斗罗大陆上的两大帝国而言，那片土地一直处在他们目光之外，始终没有被真正深入探索。

宏舒集团便在那片大陆上。

而帅气教过去千年的根基，同样扎在那里。

说到这里，她顿了顿。

``你们以为他们是突然壮大的？''

卢本韦眼中的轻慢渐渐收了起来。

女人仍被压在地上，呼吸也没有完全平复，可她说出的这几句话，已经让井旁的几人重新想起了帅气教这些年来显露出的力量。

``千年？''古月娜问。

``千年。''

那段漫长的时间里，帅气教在星辰大陆上不断发展，积蓄力量。斗罗大陆上的人们没有察觉，更没有在它尚未壮大时，就找到它真正经营的地方。

等到它在这里显露出威胁，背后早已不是一朝一夕的积累。

蔡须坤的手指缓缓收紧。

她们在山上用了一年，才一点点走到五百米。每一次停步、每一次重来，她都记得清楚。

千年。

这个原本只在旧事里听过的词，此刻落到了一个仍在活动、仍在不断伤人的势力身上。

奥斯卡低下头，将刚才几个名字逐一记下。

宏舒集团，王雨新，星辰大陆。

在最后一个名字下面，他又添了一行：斗罗大陆西岸，再往西约三千里。

``她说的这些，得一起带给大长老。''他轻声道。

古月娜点了点头，手中的银光仍没有放松。

女人望着那本册子，忽然又开口：

``既然知道了，你们还要------''

``村里的人呢？''

蔡须坤打断了她。

女人抬起头，正撞上她的目光。

蔡须坤没有去看册子，也没有望向西边。她身旁就是那两只水桶，稍远的院子里，尚未收起的衣服被风吹得轻轻摆动。

这些人没有机会等她们先去追查另一片大陆。

``集团的事，我们会查。''蔡须坤说，``现在先把这里的人交代清楚。''

她将手掌按在女人肩上，稳稳压住对方尚未平息的力量。

``一个村子的人，不会因为你报出几个名字，就可以不见了。''

就在这时，村口传来一阵急促的脚步声。

古月娜先转过头，守在身侧的银龙虚影抬起了头。蔡兴坤将黑子鸡招到路边，看见一队人沿着土路快步赶来。

领头的两人，一个身着扽魂殿的服饰，另一个披着来不及整理的外袍。随行人员刚走进村口，便看见了地上的菜篮、敞开的院门，还有井台旁被撞碎的石块。

来人正是海潮市扽魂中殿的殿主与市长。

两人赶到井边，目光落在被压制的女人身上，脚步同时一顿。

``宏舒集团的人\ldots\ldots''

殿主看清她的面孔，又扫过井边尚未散尽的镜光，脸色沉了下来。

女人将脸偏向一旁。

市长蹲下看了一眼她被缚住的双手，随即起身，向小队拱手。

``近来市内作乱四起，查到的几条线索，都指向宏舒集团。没想到在这里抓到了人。''

他望向那些空着的屋舍，神情愈发凝重。

``让我们带回去，严加审问。''

蔡须坤没有立即撤去金光。

``这个村子的人都不见了。我们进来时，饭菜还放在桌上，灶里也还烧着火。''

她将方才见到的情形简要说了一遍，又指向女人。

``她藏在井里，用镜阵袭击我们，还能反射打向镜面的力量。她身上的冤魂之力，也要留意。''

殿主的目光沿着井沿转了一圈。

他抬手示意，随行人员立刻分出几人，沿着村中道路检查屋舍。市长又低声吩咐了几句，让人看守村口、清点各家情况，先把搜索失踪村民的事情接了过去。

``人命要紧。''市长对蔡须坤道，``我们会把这一条放在前面问。''

奥斯卡翻开册子，将女人刚才供述的几个名字，以及星辰大陆的方位，另抄了一份交给两人。

殿主看到王雨新与胡赞的关系时，眉头明显一紧。

他把那张纸折好，收入怀中。

``这些消息，我们会一并上报。''

随行人员取来了押解用的禁锢镣铐。古月娜配合着将女人的双腕锁住，又以银光封住几处尚在凝聚力量的纹路。蔡须坤缓缓收回金光，直到看见对方再也无法聚起那片镜面，才起身让开。

女人被押了起来。

方才还能悬在井口上方的人，此刻脚步虚浮，衣摆拖在泥地上。她经过母猪扽魂身边时，后者鼻子一动，她便下意识地避了一下。

卢本韦看着她被带到路旁，这才抬手，将母猪扽魂收回身侧。

殿主也在这时重新望向五人。

刚进村时，他的注意全在井边的异状与被擒的女人身上。如今站得近了，蔡须坤正在收起领域，古月娜的银光也逐渐退去，几人运转力量时显露出的气息，反而更加清楚。

他的目光从蔡须坤移到卢本韦，再看向另外三人。

``你们的修为\ldots\ldots 居然在我之上。''

话出口后，他自己也怔了一下。

五个人看上去都还年轻，站在一起，身上还带着赶路的风尘。若非亲眼见到，他很难把眼前的小队，与井边残留的这场交手联系起来。

卢本韦先笑了一声。

蔡兴坤低头摸了摸黑子鸡的背，也忍不住弯起嘴角。

蔡须坤与同伴交换了一下目光，五人一同向两人行礼。

``见过殿主，见过市长。''

蔡须坤直起身，笑着说：``我们刚从海外历练回来，路过这里，察觉不对，就进来看看。''

殿主回了一礼。

``多亏你们及时出手。''

市长看着路旁那只仍在徘徊的黄狗，向留在村里的人员招了招手，让他们先从附近几户查起，连井下和屋后的地方也不要漏过。

押送女人的人先行往城里去了。

村中终于重新有了走动与交谈的声音。有人检查灶房，有人逐户查看门窗，发现遗落在地上的东西，便留在原处做好记号。

殿主抬头望向通往村外的另一条路。

``消失了这么久的宏舒集团，居然又出现了。''

他握了握袖中的那张供述，转向小队。

``我们得快速巡查周边村镇，看看还有没有同样的情况。你们保重。''

``两位也保重。''蔡须坤说。

市长向几人点头，留下人员继续搜索，便与殿主带着其余人匆匆往村外赶去。

蔡须坤目送他们走出村口，又回头看了一眼井台。

井绳已不再晃动。那只装满水的木桶仍放在原处，桶沿挂着几滴将落未落的水珠。

奥斯卡合上册子，放回随身的包里。

这一次，准备带给大长老的东西，又多了几页。

离开村庄后，五人继续赶路。

宏舒集团的名字被记进册子里，连同星辰大陆的方位，一并留在奥斯卡随身的包中。一路经过城镇，他们也会留意沿途的消息，却没有再遇上井边那样的事。

半个月后，一座熟悉的城市出现在道路尽头。

蔡须坤先认出了入城的方向。

可真正走近时，她的脚步又慢了下来。

路比记忆中宽了许多，两旁铺着整齐的石板，过去容易积水的低处修了排水沟。新栽的树沿街排开，枝叶尚不浓密，树下却已放好了供人歇脚的长凳。

再往前，沿街的房屋也重新修整过。

有些旧铺面还在，门前却多了一段平坦的台阶，原本挤在路边的摊位被移到较宽敞的地方。推车与行人各走一侧，卖吃食的摊子升着热气，来往的人说着话，从他们身旁经过。

奥斯卡在一处路口停了一下。

``以前往学院去，是从这里转吧？''

蔡须坤看向他所指的方向。

那条路还在，只是路口外又接出一条新街。几排新修的楼房遮住了远处的景物，站在原来的位置，已经看不见记忆里的那一段院墙。

``应该是。''她说着，自己也笑了，``进去看看就知道了。''

五人顺着街道往里走。

越往前，穿着学院服饰的学生越多。有人抱着书匆匆经过，也有人三两成群，在路边争论着什么。前头一个年轻学生走得太急，险些撞上蔡兴坤，忙低头道歉，又抱紧怀里的东西跑远了。

蔡兴坤侧身让开，黑子鸡从她臂弯里探出头，跟着看了几眼。

``这里变化也太大了。''她说。

这一次，卢本韦没有马上接话。

她正望着街道尽头。

一片宽阔的院墙从那里延伸出去，越过路旁的树，又一直向远处铺开。围墙后面，几栋高楼的屋顶错落露出，楼与楼之间留着大片空处，还能看见新修的步道与成行的绿树。

五人不约而同地放慢了脚步。

史莱克学院到了。

可是眼前的学院，几乎让他们认不出来。

奥斯卡走到街边，沿着围墙望了好一会儿。

``那边也算学院的地方了？''

他抬手指向远处。如今围墙早已越过他记得的边界，将更大的一片土地圈了进去。几处新建筑上还搭着收尾用的木架，墙外则已经清理干净，种上了低矮的花木。

卢本韦又往另一头看。

``不止那边。你看这里，墙也往外移了。''

蔡须坤没有急着说话。

她试着将记忆里的学院放回眼前，一处处对照过去，最后发现，能够直接对得上的地方已经不多了。

最显眼的，还是前方的大门。

门前空出了一片平整宽阔的场地，两根高大的石柱立在左右，柱面雕着整齐的纹饰。上方的门额比从前高了许多，深色底面上，几个大字清清楚楚地映在日光里。

史莱克学院。

蔡须坤抬头看着那几个字，终于确信自己没有走错。

``连大门都换了。''奥斯卡轻声说。

卢本韦往前走了几步，又回头看看身后的街。

``刚进城时，我还以为只是城里修整了一番，没想到学院扩建成了这样。''

蔡兴坤将黑子鸡往怀里抱稳，踮脚朝门内看去。

从他们所在的位置，只能看见一段宽阔的路，路尽头有学生经过，再远处的楼舍便被树木挡住了。

古月娜站在几人身边，也顺着那条路望了一眼。

``进去看看？''

这句话让蔡须坤回过神来。

一路上，她惦记着要向大长老请教的事情，也惦记着还留在蓬莱岛的同伴。直到此刻，看见学院的名字，她才真切地想到，自己已经回到能够见到老师们的地方了。

她忽然很想听听他们的声音。

想知道这段时间学院发生了什么，也想把这一年里在山阶上遇到的难处，挑几件说给他们听。

卢本韦已经转回身。

``既然路过了，总得去看看老师们。''

``当然要看。''奥斯卡笑着应道。

他低头理了理赶路时压皱的衣襟，又将背包往肩上提了一点。蔡兴坤看见，也腾出一只手，拍去袖口沾着的灰。

蔡须坤忍不住笑了。

她将落到脸侧的一缕头发拨回耳后，和同伴一起向前走。

先前远远看见的石柱，渐渐露出更清楚的轮廓。门额下的阴影落在平整的地面上，几名学生从门内走出来，与他们擦肩而过。

五人一边走，一边指认着围墙后偶尔露出的景物，猜哪一处还留着旧日的模样。

不知不觉间，他们离学院大门，已经只剩下短短的一段路。

\chapter{第一百三十三章：重回史莱克}\label{ux7b2cux4e00ux767eux4e09ux5341ux4e09ux7ae0ux91cdux56deux53f2ux83b1ux514b}

走到校门前，蔡须坤才发现，门内还添了一间值守的小屋。

窗边摆着登记册，两名守卫轮流留意进出的学生。其中一人见五个人径直走来，迎上两步，抬手示意他们稍等。

``几位是来找人的？''

``来看看老师。''蔡须坤笑着答道，``我们以前在这里读书。''

守卫点了点头，将几人让到不挡路的一侧。

``请问是哪一届的？找哪位老师？''

奥斯卡报了名字。

守卫听到一半，翻册子的手便停住了。他重新抬头，从奥斯卡看到蔡须坤，又转向卢本韦，脸上渐渐露出一点迟疑。

``你们是\ldots\ldots''

他像是想起了什么，可看着几人的脸，一时又不敢确定。

这几年过去，蔡须坤的五官已经长开，身量也与当年不同。她将长发束在身后，站在那里，笑意仍旧温和，眼神却比记忆中沉静许多。卢本韦也不再是从前那副模样，连奥斯卡背着包站在旁边，都让人觉得既眼熟，又对不上。

守卫又看了一遍册子。

蔡兴坤低下头，忍不住笑了起来。

``看来不只是我们认不出学院了。''

卢本韦向旁边让出一点空地。

``脸认不出来，这个总不会认错吧。''

她抬起手，魂力在掌心轻轻一转。

一头母猪扽魂随即在她身侧凝实。

它四蹄落地，先甩了甩脑袋，随后抬起鼻子，朝街边飘来的食物香气闻了闻。卢本韦伸手按住它的背，它才不情不愿地把头转回来。

没有展开领域，也没有故意释放威压。

可那股属于神级扽魂的气息，仍让守卫握着册子的手紧了一下。

一旁正要进门的学生停住脚步，目光落在母猪扽魂身上，再也挪不开。

``母猪扽魂\ldots\ldots''

他身后一个年长些的学生忽然睁大了眼睛。

``是神级的母猪扽魂！''

这句话一出口，门前几个人同时回过头。

守卫终于将眼前的面孔，与记忆里的名字对上了。

``卢本韦？''

卢本韦笑着点头。

``是我。''

``那你们就是果糕小队！''

学生看向另外四人，声音不自觉地高了起来。

``学长学姐回来了！''

门内有人听见，转身跑进了学院。附近的学生也围近了些，有人低声向同伴解释，有人听到蔡须坤与奥斯卡的名字，立刻探头往这边看。

一时间，母猪扽魂比几个人还受瞩目。

它倒没有半点被围观的自觉，凑到卢本韦脚边，轻轻拱了拱她的小腿。

蔡须坤看了它一眼，笑着对卢本韦说：``果然还是它好认。''

``那当然。''卢本韦拍了拍母猪的背，``换了学院大门，也不能把我们忘了。''

守卫忙把登记册合起，笑着让开路。

``方才没认出来，几位别见怪。你们这几年变化实在太大了。''

``是我们回来得突然。''奥斯卡说。

几人正要往里走，门内便传来一道声音。

``人呢？''

一个学生在前面指路，话还没有说完，身后的人已经越过他，大步朝校门走来。

弗兰德到了。

他来得很快，走近以后却停了一下，隔着几步距离，将门前五人从头到脚仔细看了一遍。

蔡须坤望着那张熟悉的脸，心里忽然一热。

``院长。''

五人一同行礼。

弗兰德看着他们，先前赶路时紧绷的神情慢慢松了下来。

``还真是你们。''

他向前走了两步，抬手拍了拍奥斯卡的肩，又看向卢本韦身侧的母猪扽魂。

``到门口了，怎么也不先让人通报一声？''

``正要通报呢。''卢本韦笑道，``您就来了。''

弗兰德也笑了。

校门口聚的人越来越多，他抬头朝学生们挥了挥手。

``该上课的上课，别都堵在这里。人既然回来了，还怕见不着？''

学生们这才散开，却仍有人边走边回头。刚才跑去报信的那个学生喘着气站在旁边，脸上的兴奋还没消下去。

弗兰德将五人领进学院。

``走，先进去。你们这一趟，可让人等了好久。''

门后的路比从前宽了，路旁新添的树木投下斑驳的影子。蔡须坤才走了不远，便看见几道身影从楼舍间迎出来。

最前面那人肩背宽厚，脚步迈得很大。

``赵老师。''奥斯卡先喊了一声。

赵无极看见他们，脸上顿时露出笑容。

``舍得回来了？''

他走到跟前，刚想伸手拍一拍几人，目光却忽然停在蔡须坤身上。

随后，他又看向卢本韦。

这一下，他脸上的笑意变成了惊讶。

``等等，你们现在的修为\ldots\ldots''

跟来的几位老师也察觉到了不同。

五人一路将气息收在身内，方才门口也只让母猪扽魂出来露了个面。可站在近处，仍能感到他们运转力量时的深厚与稳定，与离开时已经大不相同。

弗兰德没有急着替他们回答，转头看着蔡须坤。

蔡须坤点了点头。

``这次历练，确实又进了一步。''

她稍稍放开收敛的气息，让老师们能够看清。卢本韦与奥斯卡也依次运转魂力，没有将力量向外推出，只在各自身周停留片刻，便重新收住。

老师们安静了一瞬。

赵无极看了几人好一会儿，终于摇头笑出声。

``你们真是怪物一样。''

旁边一位老师也忍不住笑了。

``出去历练几年，回来一次，就让我们吃惊一次。''

蔡兴坤抱着黑子鸡，听得有些不好意思。

``也没有那么容易。我们在同一个地方，卡过很久。''

``有时候一天走不了几步。''奥斯卡接道，``还得退回去，重新想明白才能再往上走。''

``肯退回来重想，也是本事。''赵无极说。

弗兰德看着几人沾灰的鞋履，将还要继续问话的老师们招呼到一旁。

``先让他们坐下喝口水。堵在路上，打算问到什么时候？''

这句话让众人都笑了起来。

五人被领进一间宽敞的会客室。

窗户开着，外面是新修的步道与一片草地。有人送来清水，又添了几样吃食。奥斯卡将背包卸下，坐进椅子里，肩膀终于松了下来。

蔡兴坤给黑子鸡挪了个位置，它便安静地伏在椅旁。卢本韦将母猪扽魂收起，捧着水杯，长长地呼出一口气。

半个月赶路时绷着的那股劲，到这里才缓下来。

弗兰德在对面坐下，看了看五人。

``其他人呢？没和你们一起回来？''

蔡须坤放下杯子，将蓬莱岛上的情况简要说了。

他们五人先行返回，唐三、小舞、诺当和叶阳还在继续完成第一考。霍雨浩与薛峰已经通过，留在那里照应同伴。等事情办妥，他们还要回去会合。

``这次回来，是有些修炼上的问题，想向大长老请教。''她说，``路过学院，就想先来看看老师们。''

弗兰德点了点头。

``人都知道互相照应，就好。''

他没有催他们将所有经历一口气讲完，只让几人先歇一歇。

蔡须坤喝了水，抬眼看向窗外。

不远处有一座很高的楼，楼前台阶分向两侧，几名学生抱着书走下来。再远一点，原本从这里就该能望到院墙，如今却还是学院里的路。

她忍不住问：``院长，现在学院到底扩到哪里了？''

弗兰德顺着她的目光望出去，眼中有了笑意。

``你们刚才看见的，还只是一部分。''

休息了一会儿，几人放好行李，跟着老师们走出了会客室。

直到真正沿着新修的路走过去，他们才明白，校门外所见的变化，远远算不上全部。

教学楼增加了好几处。楼与楼之间留着宽敞的通道，讲堂也比从前大了许多。经过一处开着门的教室时，蔡须坤朝里看了一眼，只见座位一排排向后延伸，后排略高，坐在那里也能清楚看见前方的演示台。

``以前一间教室不够坐，现在能分开授课。''一位老师介绍道，``不同进度的学生，也不用都挤在一起等。''

卢本韦向后看了好几眼。

``我们刚才从外面看，还以为那整栋都是一间大讲堂。''

``讲堂在另一边。''老师抬手指了指。

她顺着看去，果然又看见一座更宽的建筑，顿时笑着摇了摇头。

再往前，是扩建后的训练场。

空旷的场地被划分成不同区域，几队学生同时练习，彼此仍隔着很远。场边留有休息与观摩的位置，通往训练室的路也重新铺过，宽得足够几组学员并排行走。

赵无极在一旁说：``从前总得轮着用，现在不同课程可以同时安排，省下不少等候的工夫。''

奥斯卡望向训练场另一头。

``光这一片，就比我们记得的大了太多。''

``后面还有。''

这四个字，几人一路听了不止一回。

藏书的地方扩了，阅览的位置添了，训练后的休整区域也单独留了出来。宿舍楼往更远处排开，周围种着树，中间有足够学生往来与歇息的空地。

走过食堂时，蔡兴坤还朝门里多望了一眼。

``连吃饭的地方都这么大了。''

奥斯卡笑道：``这个我刚才就想问。''

卢本韦跟着往里看，视线落在几排整齐的长桌上。

``至少这回，不用担心几个人一坐，就把一张桌子占满了。''

老师们听得直笑。

弗兰德领着他们在一处树荫下停住，抬手指向远处的院墙。

``学生多了，老师也多了，原来那些地方自然不够用。扩建不能光把墙往外推，教室、训练场、宿舍，都得跟上。''

蔡须坤沿着他指的方向望去。

从前熟悉的学院，如今已经延展到她一眼望不完的地方。路上不断有学生经过，见到院长与老师们，便停下来行礼；有人认出他们，又兴奋地与同伴低声说话。

弗兰德提起城中的变化，几人才知道，这些年发展的也不止学院所在的这一座城。

周围几座城市都比从前兴盛，往来的道路修整了，商旅与求学的人多了，城内的住处、铺面和各项设施也陆续增加。几座城之间联系得更紧，过去需要绕行的地方，如今赶起路来已经方便许多。

``难怪进城的时候，总觉得哪里都变了。''蔡须坤说。

``以后还有得看。''弗兰德道，``既然回来了，就别急着只看这一条路。''

蔡须坤笑着应下。

他们又往前走了一段，在一处连接旧院区的转角停住。这里还留着几棵老树，树干粗了些，枝条伸到路上，投下熟悉的阴凉。

奥斯卡抬起头，看了好一会儿。

``这几棵，总算还认得。''

卢本韦伸手摸了一下树皮，刚要接话，身后忽然传来一道熟悉的声音。

``听说你们回来了，我还以为要去校门口找。''

蔡须坤的手停在身侧。

那个声音，她太熟悉了。

林中休整时，队伍遇险时，还有那些需要有人站在前面的时刻，她都曾听见过。

五人几乎同时转过身。

树荫外，一道熟悉的身影正沿着步道走来。来人看着他们，脚步稍稍缓了些，目光从五人的脸上一一掠过。

蔡兴坤先睁大了眼睛。

``梅长老？''

来人正是梅天狸。

她在几人面前停下，脸上慢慢露出笑意。

``让我好好看看你们。''

蔡须坤上前一步，和同伴一起向梅天狸行礼。

``梅长老。''

梅天狸点了点头，目光在五人身上停留了片刻。

方才还笑着说话的几个人，此刻站到她面前，都不由自主地挺直了些。卢本韦将手从树干上收回，奥斯卡也放下了原本搭在包带上的手。

梅天狸看在眼里，笑意更深了一点。

``这两年里，你们的成长，让我很惊喜。''

她能察觉到五人收在身内的力量。与从前相比，那份力量厚重了许多，运转时也更稳，连几人停步后调整呼吸的节奏，都透着长久磨炼留下的习惯。

``你们这次回来，应该是想去总殿，请教大长老吧？''

蔡须坤点头。

``是。有些问题，我们自己还想不明白，准备当面向他请教。''

``那便把问题整理好，带过去。''

梅天狸说着，朝步道另一头看了一眼。

几个学生正站在那里，原本还在低声交谈，见她望过来，立刻站直。可他们的目光仍忍不住在五人身上打转，显然还不愿意走。

弗兰德顺着看过去，笑着摇了摇头。

``消息传得倒快。''

一名学生鼓起勇气，向前走了两步。

``院长，我们就是想看看几位学长学姐。课上常听老师提起，今天才见到真人\ldots\ldots''

说到后面，他自己也有些不好意思。

卢本韦笑了，蔡兴坤则低头看了一眼怀里的黑子鸡。它懒洋洋地缩着脖子，对周围的热闹毫不关心。

梅天狸收回目光。

``不过，学弟学妹们的热情，确实很高。''

她转向小队，语气变得认真了些。

``你们刚赶了这么久的路，今天先好好休息。等明天把状态调整好，我在训练场布下结界，你们到时候全力出手，和我比上两招。''

奥斯卡微微一怔。

``和您切磋？''

``嗯。''梅天狸说，``也让学弟学妹们好好学一学，看看力量到了实战里，该怎样配合着用。''

那几个学生顿时睁大了眼睛。

弗兰德看了看小队，随后向梅天狸点头。

``观摩的位置，我来安排。''

梅天狸应了一声。

奥斯卡却有片刻没有说话。

眼前宽阔的步道与树荫渐渐淡下去，他忽然想起了第一次与梅天狸切磋时，那座悬在云间的试炼台。

迎面吹来的风很冷。

可真正交手以后，他几乎感觉不到风，只记得身上越来越沉，魂力被一层层抽空。所有人都在想办法，将各自还能拿出的那一点力量接到一起。

到了最后，他们几乎用尽了所有手段。

那道汇聚全队之力的蓝光冲出去时，他曾睁大眼睛，生怕错过了碰撞的那一刻。

随后便是强光、震动，还有贴着背脊传来的冰冷台面。

他连抬手都费劲，只能躺在那里，望着光芒慢慢退下去。

而梅天狸，依旧站在他们面前。

倾尽全力，最终只让她动了半寸。

那半寸，他一直记得。

奥斯卡回过神来，重新看向树荫下的梅天狸。

她说话的语气，与当初带他们上云台时一样平静。那时他还很难想象，几个人合在一起，拼到魂力耗尽，竟也只能做到那一步。

如今，又会怎样？

他下意识地看向蔡须坤。

蔡须坤也正望着他，眼中有一点笑意。

显然，记起那场切磋的，不止他一个。

卢本韦活动了一下手腕，先开口道：``好啊。''

她说得很快，随后又认真补了一句。

``这一年练成什么样了，也该好好试一试。''

蔡兴坤抬起头。

``我也想试。''

在中分山上，他们无数次回头，去看自己已经走过的台阶。可每过一处，又有新的压力压下来，很少能停在原地，完整地用一场交手检验自己的变化。

井边那场战斗结束得很快。女人的手段虽然诡异，却还没有让他们将这一年来积累的力量与配合尽数施展出来。

如今，梅天狸就站在眼前。

那位曾让他们清楚看见自身极限的长老，再次给了他们一次出手的机会。

蔡须坤望向古月娜。

古月娜轻轻点头。

奥斯卡也笑了起来。

``那今晚就把该准备的都准备好。明天我的补给，一定跟上。''

蔡须坤转回身，向梅天狸郑重应道：

``我们愿意。明天，一定全力以赴。''

梅天狸看着他们。

她记得当初几个人面对自己时，眼中那份强撑着的紧张，也记得他们一次次重新站起来，试着将配合衔接得更紧。

如今，这几个人仍然知道她的实力，却已能望着她，坦然地说想再试一场。

梅天狸缓缓点头。

``有这个心气，很好。''

她的目光从蔡须坤移向卢本韦，再落到另外三人身上。

``你们突破魂帝之后，武魂融合技完全可以提升到魂斗罗级别。''

卢本韦脸上的笑容稍稍收起，听得更加认真。

梅天狸继续道：``力量都长进了，彼此接续的方法也要跟着调整。这次来的是你们五个，今晚先把各自能接上哪一步、需要同伴在哪一刻配合，重新理清楚。''

奥斯卡点了点头，将这几句话记在心里。

``配合好，''梅天狸看着他们说，``别让我失望。''

``是！''

五人的回应几乎同时响起。

步道另一头，几个学生压低了声音，脸上的兴奋却怎么也藏不住。其中一人已经转过身，迫不及待地想将明日训练场上的约定告诉同伴。

蔡须坤望了一眼远处的新训练场。

那里此刻还有学生练习，偶尔传来一两声呼喝。宽阔的场地被午后的阳光照着，四周清清楚楚，再没有云台上翻涌的云雾。

她握了握手指，随后松开。

明天，便在这里，再向梅长老请教一次。

说定明日的切磋，弗兰德便领着众人折回了会客室。

一路绕过新修的楼舍，卢本韦又回头望了几眼。坐下以后，她终于忍不住问道：``院长，这么大的地方，建起来得花多少钱？''

弗兰德闻言，脸上露出笑容。

``多亏了卢先生的支持。学院目前资金无比充足，许多从前想做、却一直没能做的事情，如今都能慢慢办起来了。''

他说着，向卢本韦点了点头。

``回去见到你父亲，替我再道一声谢。''

卢本韦有些意外，随即笑着应下。

``好，我一定转告。''

弗兰德提起几处新添的设施，说到训练室与藏书的地方，语气中难掩满意。几位老师也接着说起新学年的安排，有了足够的场地与经费，他们能给学生做的事情，比从前多了许多。

梅天狸在旁听着，端起杯子，却没有立刻喝。

她望向窗外。

步道上仍有学生经过，远处的训练场偶尔传来呼喝声。这里的一切都在向前生长，看得见新楼，也看得见年轻人脸上的朝气。

可她的神情渐渐沉静下来。

``我有预感，这片大陆要发生大事情，就在不久的将来。''

会客室里的谈话慢慢停住。

弗兰德看向她。

``梅长老？''

梅天狸将杯子放下，还没有接着说，蔡须坤便转头与奥斯卡对视了一眼。

奥斯卡立刻拿过身旁的背包，抽出了那本一路带回来的册子。

``有件事，我们也正想告诉老师们。''蔡须坤说。

她从海潮市附近的空村讲起。

灶里未熄的火，桌上没吃完的饭，还有藏在井中、以镜阵袭击他们的女人。她将亲眼见到的情形说清，又讲到女人被生擒后供出的来历。

``她是宏舒集团的人。''

梅天狸的目光停住了。

蔡须坤继续道：``宏舒集团与帅气教往来频繁。按她的供述，集团的最高统治者王雨新，是胡赞的妻子。他们经营的地方，在斗罗大陆西面海上的星辰大陆。''

奥斯卡将记着方位的那一页摊开。

西岸，再往西约三千里。

那片大陆的面积，约为斗罗大陆的十分之一。帅气教过去千年，便在那里隐秘发展，积蓄起了如今的力量。

弗兰德原本搭在扶手上的手慢慢收紧。

赵无极也坐直了身体。

``千年的经营\ldots\ldots''

他低声重复了一遍，抬眼看向梅天狸。

梅天狸接过册子，将几条记载从头看了一遍。

``人现在在哪里？''

``交给海潮市扽魂中殿殿主与市长了。''蔡须坤答道，``他们会继续审问，也已经安排人搜寻失踪村民、巡查周边村镇。我们将供述另抄了一份，留给了他们。''

梅天狸点头，手指仍停在宏舒集团的名字旁。

``这些消息，也要带到总殿去。''

``我们会的。''

蔡须坤稍稍停了一下，随后抬起头。

``但我们还想再往前做一步。''

几位老师都望向她。

``帅气教和宏舒集团能彼此往来，积蓄这么久。我们这里出了事，却常常要等到人受了伤、村子空了，才开始追查。''

她看着桌上的册子，声音渐渐稳下来。

``我们想请总殿出面，号召两大帝国联合起来，共同对付这些邪恶势力。消息能互通，遇到袭击能及时驰援，查到线索，也能一起往下追。''

弗兰德缓缓点头，正要开口，蔡须坤又接着说：

``还要争取和魂兽们联合。''

这一次，赵无极直接抬起了头。

``和魂兽联合？''

旁边几位老师也愣住了。

他们原以为，小队接下来会谈怎样应敌、怎样将消息尽快送到各方手中。可听蔡须坤的意思，他们竟还准备化解人类与魂兽之间长久的敌意。

梅天狸看着她，神情也郑重起来。

``你们想做到哪一步？''

``一起对抗掠夺它们的势力，也一起找出今后能共同生活的办法。''蔡须坤说，``帅气教伤害人，也掠杀魂兽。那些被毁掉的村庄与栖息地，总不能一处处继续增加下去。''

古月娜接过话。

``魂兽之间，也有不同的族群和旧怨。要联合起来，需要有人愿意坐下来谈，需要能让各方相信的约定。''

她的目光落到册子上。

``人类遭到袭击，魂兽未必知道；魂兽的族群被掠杀，人类也未必听得到消息。如果能真正接上联系，至少有些事情，可以在失去更多生命之前一起处理。''

会客室内静了片刻。

窗外有人跑过，脚步声由近而远，很快消失在步道尽头。

弗兰德看着几人，脸上的惊讶还没有完全退去。

``你们这一趟出去，想的事情，已经这么大了。''

奥斯卡将册子翻到后面的空白处。

``所以，我们想建立一个传灵塔。''

他提笔，在纸上写下三个字。

传灵塔。

几人的视线一同落了过去。

``让人类与魂兽之间，真正有个能传递消息、沟通事情的地方。''奥斯卡说，``哪里遇到了危险，哪里愿意谈合作，消息可以送来，也有人负责传出去。''

蔡兴坤看着那三个字，轻声补充：``发生争端，也总得有个地方，让双方把事情说清楚。''

蔡须坤点头。

``还有魂灵。''

梅天狸望向她。

这两个字，她并不陌生。

在星斗大森林里，她亲眼见过狼王保下意识，也看见过它怎样回应蔡须坤。那场尝试中，每个人做过什么、等过多久，她都记得。

``你们想把那一次的共生，继续做下去。''她缓缓说道。

``嗯。''蔡须坤答道，``把其中真正想明白、能做稳的部分整理出来，让愿意尝试的人类和魂兽都能了解。''

她望着梅天狸，没有回避对方的目光。

``狼王当时愿意，是因为它自己作了决定。以后也应该这样。魂兽保有意识，人类得到伙伴，双方一起修炼，一起成长。''

卢本韦将手放在桌沿，接着说道：``传灵塔里也可以有人专门研究这些。成功的经过留下来，遇到的问题也留下来，后来的人才有东西可学。''

``霍雨浩在这方面帮过我们很多。''奥斯卡说，``等大家会合，也要和他一起商量。''

赵无极听到这里，终于慢慢呼出一口气。

他将册子拉近一些，重新看了一遍纸上的那三个字。

``老师，''蔡须坤轻声道，``魂兽也要有能生活、能繁衍的地方。人类修炼，需要资源，也需要安稳的家园。如果能通过魂灵一起发展，我们就能少一些彼此猎杀，多一些共同守护的东西。''

古月娜在她身侧点了点头。

弗兰德久久没有说话。

方才在院中，他看着扩建后的楼舍，满心都是学院往后能教出多少学生。此刻坐在他面前的学生，却已经将目光放到了院墙之外，放到了两大帝国与魂兽之间。

他伸手，将写着``传灵塔''的那页纸稍稍转正。

``这份设想，你们也整理好，带给大长老。''

梅天狸看着那三个字，缓缓点头。

``要让各方都听见，就从把你们知道的、见过的，还有准备去做的事情，说清楚开始。''

她抬眼望向五人。

``你们已经有了想法，也有过真正的共生经历。这些，都值得带到总殿去谈。''

蔡须坤应了一声，将桌上的几页纸拢到一起。

她在``传灵塔''三个字下面，落笔写道：

让人类与魂兽共同发展。

\chapter{第一百三十四章：再战梅长老}\label{ux7b2cux4e00ux767eux4e09ux5341ux56dbux7ae0ux518dux6218ux6885ux957fux8001}

第二天一早，五人便来到了训练场。

昨日远远望去，还只觉得这里宽阔。如今真正站在场边，蔡须坤才发现，从入口走到中央，竟还隔着很长一段路。

可比他们来得更早的，还有围观的学生。

训练场外已经聚起了许多人。前面的学生站在指定的位置，后面的人便往两侧挪，想找一处能看清场内的空隙。有人还拿着没吃完的早点，听见果糕小队到了，立刻转过身来。

``就是他们。''

``昨天门口那个学姐，她的扽魂真的是母猪？''

卢本韦听见了，朝那边笑了一下。

``一会儿就看到了。''

人群中响起几声笑，随即又被老师招呼着往后退开。

弗兰德与赵无极站在观摩区域前方。梅天狸则已经到了场内，掌心向下一按，一圈光从她脚边铺出去，沿着场地边缘升起，渐渐合成一层笼住训练场的结界。

光幕很淡，从外面仍能看清里面。

蔡兴坤抬头望了一眼，肩上的黑子鸡也跟着伸长了脖子。

``进来吧。''梅天狸说。

五人依次跨入结界。

奥斯卡将刚凝出的香肠分到同伴手中，自己也吃下一根。卢本韦活动着手指，走到蔡须坤侧后方。古月娜与蔡兴坤分向两边，几个人隔着合适的距离停住。

一夜休整后，赶路留下的疲乏已经退去。

蔡须坤握了握手，望向场地中央。

梅天狸站在那里，看见他们摆开的阵势，脸上露出一点笑意。

``好久没这样交手了。''

她抬起右手，朝五人招了招。

``上一次陪你们玩玩，还不够。这一次，尽管过来。''

话音刚落，蔡须坤已经冲了出去。

没有试探，也没有多余的停顿。

篮球在她掌中凝实，澄澈的处女之光沿着指尖覆上球面，紧接着，一缕灰色从她脚下涌起，随着果糕领域的展开，迅速缠入那层光里。

混沌之力。

两种力量相接，球面上的光骤然向内一缩。

下一瞬，狂暴的气息轰然荡开。

蔡须坤脚下的尘土向后翻起，整个人借着前冲的势头，将篮球压在掌下，直逼梅天狸身前。灰色与清光交错着旋转，球还没到，迎面冲来的风已经扯动了梅天狸的袖口。

梅天狸抬起拳头。

一层凝实的魂力裹住拳面，迎着那颗篮球抵了过去。

砰！

撞击声在结界里猛地炸开。

球与拳相接的地方，光芒向四周喷出。地面的尘土被压出一圈，又沿着场地向外疾冲，撞到结界上，激起一片清晰的波纹。

外面几个原本还在交谈的学生，同时闭上了嘴。

梅天狸的手臂微微一沉。

她眼中的笑意停了一瞬。

这股力量，比她预想的重得多。

最先压过来的是篮球与蔡须坤前冲的力道，可接住以后，球面深处竟还有一股翻涌不定的力量，沿着相接处接连撞来。她才稳住正面，混沌之力又将接触处的魂力扯偏了一线，澄澈的光随即从那一线中压了进去。

梅天狸屈了一下手肘，将原本平平抵出的拳头收近。

魂力重新凝聚，才把那股冲势稳稳接住。

蔡须坤的手臂也随之绷紧，掌心传来的反震让指节一阵发麻。她脚下用力，果糕领域贴着身侧一卷，将向外散去的混沌之力重新拢向球面。

力量又压上来一层。

梅天狸终于抬眼，认真看了她一眼。

``长进不小。''

蔡须坤没有答话。

她听见侧后方传来一声短促的呼喝。

``来了！''

卢本韦抬手向前，母猪扽魂已经落在场中。

四蹄刚触地，厚重的光便沿着它的肩背隆起。它的身躯迅速膨胀，背脊高高抬起，转眼间已大到足以将梅天狸整个人遮在身影里。

场外有人忍不住往前探了探。

``这么大！''

母猪扽魂却没有直着撞过去。

它前蹄一错，沉重的身躯竟贴着蔡须坤侧方转了半圈。地面被踩出一道弧形的痕迹，肩背间积蓄的力量却没有散，随着转身，斜斜撞向梅天狸的另一侧。

几乎同时，蔡须坤收球，借着反震向旁边退开。

两个动作接得极紧。

梅天狸才要将拳上的力道向前送，面前已经换成母猪扽魂压来的肩背。

她左手一抬，按在它额前，身体顺势侧转。

母猪扽魂的冲势被带偏，前蹄在地上犁出两道深痕。可还没等那股偏转的力量将它完全甩开，卢本韦的手指已经向内一勾。

它后腿一蹬，腰身随之一扭，竟在极短的距离里重新转过头。

梅天狸的掌心尚未离开，它便又压了回来。

结界外响起一片低低的惊呼。

方才有人见它体形巨大，还以为一旦冲过头，就很难再调转方向。可眼下这一转一压，快得让他们险些没有看清。

卢本韦站在不远处，目光始终盯着梅天狸的手。

母猪扽魂也仿佛明白了她的意思，随着她手腕的动作，四蹄接连换位，紧紧逼住梅天狸身侧。

蔡须坤已从另一边再次切入。

篮球贴着掌心划过一道短弧，处女之光与混沌之力同时涌起，向梅天狸换力时露出的空处压去。

梅天狸轻轻吸了一口气。

这一次，她的两只手都抬了起来。

左掌抵住母猪扽魂，右拳迎上篮球，两股魂力同时向外一震。

尘土猛地扬高。

蔡须坤肩头一晃，顺势退开。母猪扽魂也向后滑出一段，四蹄死死扣住地面，才重新站稳。

两人刚退，梅天狸已经向前踏了一步。

一股浑厚的拳风沿着地面压来。

``兴坤！''

蔡须坤的声音才响，黑子鸡已经从蔡兴坤肩头跃起，落向她背后。

漆黑的羽翼骤然张开。

羽根向两侧舒展，一根根羽毛泛着铁一般的暗光。随着蔡兴坤魂力涌入，双翼继续向外伸展，扇起的风将扑到脚边的尘土猛地卷开。

她脚下离地，向上拔起。

这一次，两翼之间的黑色魂力也随之展开，分成几道绷紧的光带，托住另外四人的腰背，将他们一同带离地面。

梅天狸的拳风从下方掠过。

卢本韦被带起的同时，手掌仍朝着地面。母猪扽魂低头伏身，将那股余劲扛过，随后抬头，重新追向梅天狸。

空中与地面，一下有了不同的进攻方向。

蔡兴坤向左侧振翼，四人随之移开。那几道光带始终牵在翼下，既托着他们的分量，也随着她调整方向，为同伴留出出手的角度。

``别离太远。''她提醒道。

``知道。''卢本韦应了一声。

她人在半空，母猪扽魂却已绕到梅天狸右侧，步伐比先前更快。上方看得清楚，她每一次指向，都让那庞大的身躯及时换一个位置。

古月娜也抬起了手。

银龙虚影在她身后展开，一层银光沿着黑翼下方铺开，将追上来的碎石与乱流挡住。她看准梅天狸转身的方向，手掌向外一引，一道银色魂力从侧方切入，逼向对方刚刚抬起的手腕。

梅天狸手势一换，银光便被她按偏。

可就在这短短的空隙里，蔡兴坤已将姐姐送到了另一侧。

蔡须坤顺着腰间的支撑向前一跃，借黑翼回扇的力道，再次俯冲。

梅天狸抬头。

球面的光已经压到眼前。

她右拳迎上，左手却没有跟着空出来。地面的母猪扽魂正转过身体，紧逼她的侧后方，古月娜的银光也随时可能再次切入。

梅天狸接住球，将蔡须坤向外带开。

蔡兴坤立刻振翼，翼下光带重新绷紧，把姐姐退回来的身形托住。蔡须坤顺着力道收臂，转瞬便回到队伍前侧。

这一来一回，竟没有在地上落脚。

结界外，先前还想喊几声助威的学生看得发怔。

赵无极站在人群前，低声道：``看他们换位。前一个人退回来，后面的攻击已经接上了。''

有人听见，连忙把目光从最亮的篮球上挪开，试着去看另外几人的动作。

这时，他们才发现，奥斯卡也一直没有停。

他在半空调整了一下位置，先将新凝出的补给送向同伴，随后两手缓缓分开。

几根带着红纹的香肠在掌间成形。

没有抛向队友，也没有送进嘴里。

他指间魂力一紧，那几根香肠竟向外延展，一截连着一截，转眼便拉成了柔韧的长条。

``那是什么？''

一个学生忍不住开口。

没人来得及回答他。

长条已经垂向地面。

最先落下的一根轻轻一弯，贴着地面滑开，绕过母猪扽魂踩出的深痕。第二根则从另一侧探过去，第三根仍留在空中，随着蔡兴坤移位，悄悄换了方向。

看起来，竟像几条正在生长的藤蔓。

可那上面流动的魂力，分明还是奥斯卡的香肠扽魂。

他将力量沿着长条一段段送出，手指随之轻轻屈伸。那些柔软的肠索便贴地游走，末端卷起，又慢慢放开，仿佛正试着寻找可以缠住的地方。

梅天狸的目光终于扫了过来。

奥斯卡没有躲开她的视线。

他记得，当初站在云台上，自己几乎将所有力气都用在了支撑同伴上。如今人在空中，俯看整个交手的范围，他也已经能伸手，去封住对手的一条退路。

``须坤，左边。''

他的声音很短。

蔡须坤立刻压低篮球，澄澈的光裹着灰色向梅天狸左侧逼近。

母猪扽魂也在卢本韦指引下换位，庞大的肩背横过半步，压住另一边。

古月娜掌中的银光停在恰好能够补上的位置。

而那几条红色的香肠藤蔓，正沿着两道攻击之间留下的空隙，无声地探向梅天狸脚侧。

一条贴着地面，绕向她的脚踝。

一条忽然扬起，末端弯成一道弧，直取她将要抬起的手腕。

还有一条从侧后方缓缓收紧，将她身旁剩下的空处，一点点圈了进去。

奥斯卡抬起双手。

指间的红纹，在这一刻同时亮了。

几条香肠藤蔓骤然合拢。

梅天狸脚下一错，侧身避过探向脚踝的红索，抬手一拂，将另一条索梢拨开。指间魂力顺着长条压了回去，奥斯卡双臂一震，掌心顿时有些发麻。

他立即松开那一截，余下的藤蔓却仍从侧面卷来。

梅天狸看准空隙，屈指弹出一道魂力，将几条红索一并震开。

就在这时，蔡须坤的声音从上方传来。

``不能保留了。''

她望着场中仍站得稳稳的梅天狸，掌心在篮球上收紧。

``兴坤，落地，合击！''

蔡兴坤立刻收拢半边羽翼，带着同伴斜斜落下。古月娜先踏住地面，银光向两侧一托，让卢本韦与奥斯卡先后站稳。

翼下的光带随即松开。

蔡须坤与蔡兴坤却没有停，姐妹二人并肩向前，气息在换步间迅速接上。

几人的领域几乎同时展开。

果糕领域沿着地面猛地扩张，灰色深处的混沌之力骤然翻涌。紧接着，坤神领域的金光从蔡须坤身周升起，贴着同伴铺出一片明亮的区域。

卢本韦也向前踏出一步。

猪神领域压向梅天狸所在的一侧，与母猪扽魂的气息相接。奥斯卡在她侧后方站定，将散开的香肠藤蔓迅速收回，掌心重新凝起两根粗大的香肠。

梅天狸的目光从地面掠过，重新落到蔡须坤脸上。

她的眼睛变了。

灰色气流沿着肩背向上攀，原本澄澈的眸子里，血红一层层漫了出来。混沌之力还在暴涨，连她脚边几块碎石，都被卷得悬了起来。

``姐。''蔡兴坤低低唤了一声。

``我知道。''

蔡须坤答得很快，声音却比方才低沉。

她将目光从梅天狸身上收回一瞬，看向妹妹。蔡兴坤看见那双血红的眼睛里仍有清楚的回应，便伸出手，与她相扣。

两人的魂力迅速交融。

处女之光裹住翻涌的灰色，黑子鸡收回的羽翼化成一道凝练的暗光，从另一侧接上。姐妹二人肩背同时绷紧，向前踏出的脚步落在了同一个节奏上。

``融合技------铁山靠！''

她们身前，一道叠合的肩背虚影骤然凝实。

那股冲势才向前推出，地面的尘土便被压向两边。篮球上的光纹一圈圈收紧，混沌之力在其中剧烈翻腾，沿着姐妹二人的合击向前奔涌。

另一侧，卢本韦已经抬起手。

母猪扽魂低下身，四蹄踏实，胸腹间的光猛地鼓起。奥斯卡将两根香肠向前送出，准确地落进它的鼻端。

红纹与母猪扽魂身上的光相接。

它鼻端喷出两股白气，紧接着，厚重的魂力尽数向前压缩。两根香肠表面的纹路接连亮起，绷得笔直。

``母猪炮！''

两道炮光同时轰出。

沉闷的响声让场外的学生齐齐一颤。炮弹从左右掠过，追上铁山靠凝聚的冲势，在卢本韦与奥斯卡的共同牵引下，向中央汇拢。

古月娜站在最后，手中已多了一杆长枪。

白银龙枪。

枪身银光流动，银龙虚影沿着她的手臂盘旋而下，最后尽数聚向枪尖。

她看准前方三股力量接合的位置，向前踏步，手臂猛地一送。

白银龙枪脱手飞出。

银光从姐妹二人上方掠过，准确地切入她们身前的合击之中。两侧炮光随之向内一合，围着枪身汇聚，狂暴的灰色则被裹在更深处，一层层向枪尖压去。

五人的力量，终于在这一刻接到了一起。

梅天狸的神情变了。

她昨日听见五人应战时，曾估量过他们如今能够做到哪一步。魂帝修为，配合神级扽魂与融合技，将合击推到魂斗罗级别，已经足够让她满意。

可眼前这股力量，还在继续攀升。

一道道密集的纹路在枪身四周翻起，向外展开，又在五人的合力下重新收拢。几路原本不同的力量相互接续，终于冲到了\allowbreak{}\NovelInline{I06141}{n}\allowbreak{}行BMS矩阵的极限。

\NovelDisplay{M05781R01}{%
Y(1,3)%
}

蔡须坤双目血红，肩背间的混沌之力几乎尽数灌入这一击。蔡兴坤紧紧接住她的力道，卢本韦与奥斯卡将炮光稳住，古月娜则盯着枪尖，保持那一线贯穿始终的银光。

梅天狸再没有只抬一只手。

她双脚踏开，身后的海狸虚影迅速显现。宽厚的尾影从身侧展开，沉重的魂力向掌间汇聚，将前方翻卷的尘土生生压低。

``海狸破！''

她双掌向前推出。

尾影随之横转，凝成一道厚重的冲击，迎向疾射而来的银色枪锋。灌入这一招的力量，同样稳稳停在了\allowbreak{}\NovelInline{I06142}{Y(1,3)}\allowbreak{}。

两股力量，在训练场中央正面相撞。

最先响起的，是一声极短的金铁震鸣。

随后，枪锋前方的光猛地向内塌了一点。

竟停住了。

银色、灰色与炮光汇成的冲势，抵着那道宽厚的海狸尾影，双方在碰撞处僵持。一边向前贯穿，另一边牢牢压住，谁也没能在接触的瞬间将对面击散。

梅天狸的瞳孔骤然一缩。

平分秋色。

这一击，真的接住了她注入\allowbreak{}\NovelInline{I06143}{Y(1,3)}\allowbreak{}力量的海狸破。

她抬眼望向光芒后方的五人，心中第一次生出如此清晰的惊愕。

这几个人，不过是魂帝。

五人合力，已经能打出封号斗罗级别的威力了吗？

下一瞬，僵持的光团轰然炸开。

强光将训练场中央尽数淹没，地面的碎石与尘土一同掀起。结界上密集的纹路骤然亮起，将横扫而来的冲击牢牢拦在场内。

外面的学生只觉眼前一白。

再看清时，五道身影已经被掀了出去。

姐妹二人的合击最先被反震分开。蔡兴坤下意识张开羽翼，想将姐姐托住，可才撑起一半，便被迎面压来的气浪震散了翼外的光。

卢本韦被震得双臂一麻，母猪扽魂连着滑退，身形迅速缩小。奥斯卡踉跄着抬起手，香肠长索才探出去，便随他整个人一同被掀向场边。

古月娜收回牵着龙枪的银光，转身护了一下离她最近的奥斯卡，也被后续的反震推了出去。

五人一路退到训练场边缘。

蔡须坤肩背触地，翻了半圈，才勉强用手掌撑住。蔡兴坤落在她身侧，手臂上擦出一道血痕。卢本韦单膝跪地，指间微微发抖；奥斯卡坐倒下去，胸口剧烈起伏，半天没能顺过那口气。

古月娜站稳后，脚下又滑了几步，才停在两人旁边。

场地中央，梅天狸也退了。

第一步落下，地面被踏出一道浅痕。

第二步，她衣袖猛地向后扬起。

第三步，第四步。

她将双掌缓缓压下，终于稳住身体。身后的海狸尾影渐渐散去，白银龙枪从褪去的光中落下，斜斜插进了前方的地面。

结界外，一时没有人说话。

方才所有人都在看五人被击飞，直到尘土落下，才有人注意到梅天狸脚边那几道连续后移的足痕。

赵无极看得分明，握在身侧的手不觉攥紧。

弗兰德也缓缓吸了一口气。

他们听见了海狸破的名字，看见了两股力量相持，也看见了梅天狸接招以后，实实在在退了好几步。

场内的切磋，已经停了。

梅天狸抬手稳住结界，随后快步走向五人。

``先别急着起身。''

她俯身伸出手，柔和的魂力随之铺开，落到几人肩背与手臂上。擦破的伤口一点点合拢，撞击后绷紧的筋肉也在温润的力量中渐渐舒缓。

奥斯卡总算顺过了气，撑着地面坐直。

他望了一眼场地中央。

当年在云台上，他们拼到魂力耗尽，只换来了那半寸的位移。如今，那几道清楚留在地上的足痕，就在他眼前。

他嘴角动了一下，想笑，胸口却仍有些发闷，只好先低头缓了两口气。

梅天狸检查过几人的伤势，目光随即落到蔡须坤身上。

她的眼睛，仍是血红的。

合击已经散了，身周那股灰色却还在翻腾。蔡须坤想将果糕领域收回，才聚起一点魂力，肩头便骤然一颤，撑在地上的手指跟着扣紧。

``姐！''

蔡兴坤立刻挪近，伸手扶住她。

梅天狸已经按住了蔡须坤的手腕。

``看着我。''

她没有直接用更强的力量撞向那股混沌，而是将柔和的魂力缓缓送入，先托住几处运转紊乱的脉络，让向外冲出的力量有了一处能够缓下来的地方。

``别再往外送。跟着我，先把这一段稳住。''

蔡须坤抬起头，呼吸仍有些急。

梅天狸的声音很近，一句句落进耳中。她攥着妹妹的手，勉强将注意从场地中央收回来，沿着长老引导的方向，重新接住体内尚未平息的坤神之力。

金光一点点回拢。

灰色在两股力量的引导下，终于不再向外翻涌。果糕领域的边缘缓缓退回，贴近身侧，又一层层收进体内。

过了好一会儿，那双眼睛里的血红才渐渐淡下去。

蔡须坤肩背一松，额角已经全是汗。

梅天狸没有立即撤手，等她的呼吸真正平稳下来，才将护在经脉外的魂力慢慢收回。

卢本韦与奥斯卡在一旁看着，脸上的笑意也收敛了。古月娜走近半步，确认蔡须坤能够自己稳住运转，才向她轻轻点头。

梅天狸望着五人，停了片刻。

``这一击，确实有了封号斗罗的威力。''

她看向场地中央，眼中的惊讶还没有完全退去。

``我原本估的，还是低了。''

随后，她重新看向蔡须坤，语气郑重下来。

``但这股混沌之力，你还没有完全掌握。刚才打出去以后，你已经很难自己把它收回来了。''

蔡须坤低头看了看尚在发颤的手，轻轻点头。

``非必要的时候，不要动用太多这股力量。''梅天狸说，``能放出去，也要能收得回来。''

蔡兴坤握着姐姐的手，又收紧了一点。

蔡须坤缓缓吐出一口气。

``我记住了。''

梅天狸等她缓过这一口气，神情才稍稍放松。

``好了，你们休息一下。从这里到总殿已经不远了，等状态稳住，就跟着我来吧。''

她看向五人。

``我带你们去见大长老。''

蔡须坤抬起头，眼中终于露出一点笑意。

``那就麻烦梅长老了。''

``先歇着。''

梅天狸抬手撤去结界，场外的声音这才清楚地传进来。学生们还在低声讨论方才的对撞，见老师们走近，纷纷向两旁让开。

弗兰德让人送来清水，将五人领到场边的休息处。

奥斯卡取出预先备好的补给，分给同伴，自己也靠着椅背坐下。卢本韦揉了揉仍有些发酸的手臂，没有再急着运转魂力。

古月娜将落在场中的白银龙枪收回，随后在蔡须坤身旁坐下。

蔡兴坤接过水，先递给姐姐。黑子鸡已经恢复了平日的大小，伏在她怀里，翅膀收得紧紧的。

几人安静地歇了一阵。

蔡须坤慢慢调匀呼吸，确认体内的混沌之力不再翻涌，才松开一直握着的手。她看了一眼同伴，见几人也都缓过劲来，便向梅天狸点了点头。

``可以了。''

梅天狸走近，又看过她的气息，才应道：``那就去取行李。''

奥斯卡起身时，特意检查了包里的册子。修炼时留下的问题、宏舒集团的供述，还有写着传灵塔设想的那几页纸，都收得好好的。

蔡须坤将自己的笔记放进去，替他扣紧了包带。

临走前，五人再次向老师们行礼。

``院长，赵老师，我们先走了。''

弗兰德点了点头。

``到了总殿，把该请教的事情问清楚。以后有空，再回来看看。''

``一定。''奥斯卡笑着应道。

赵无极拍了拍他的肩，又看了看另外几人。

``路上听梅长老安排。''

五人一一应下。

梅天狸在前面等着，见他们告别完，便领着几人沿步道往外走。

再经过校门时，守卫已经远远认出了他们，笑着让开了路。蔡须坤回头向身后的老师们挥了挥手，随后转过身，跟上梅天狸。

五人随着她离开学院，向总殿而去。

\chapter{第一百三十五章：遵循本心}\label{ux7b2cux4e00ux767eux4e09ux5341ux4e94ux7ae0ux9075ux5faaux672cux5fc3}

总殿的石阶出现在眼前时，蔡须坤下意识地扶了一下肩上的行囊。

梅天狸走在前面，没有催促。众人沿着熟悉的路往上，越过殿前的空地，终于再次看见那十二根白玉柱。

风从柱间穿过。

石台旁，test\_alpha0正低头看着一页纸，指尖停在几行细密的符号之间。听见脚步声，他将纸放下，抬起头来。

``回来了。''

蔡须坤与同伴一同行礼。

``大长老。''

梅天狸将五人送到殿内，简要说了在学院相遇与训练场切磋的经过。说到最后那次合击时，她停了一下。

``五人合力，已经能够打出\allowbreak{}\NovelInline{I06144}{Y(1,3)}\allowbreak{}的威力。''

test\_alpha0的目光在五人身上停留片刻，最后落到蔡须坤脸上。

``力量长进了，运转还得跟上。''

蔡须坤想起训练场上迟迟收不回的灰色，认真应了一声。

``梅长老已经提醒过我了。''

test\_alpha0点头，没有让他们一直站着。他将石台旁几张空着的椅子挪近，示意众人坐下。

奥斯卡卸下行囊，将几本册子放到桌边。路上的见闻、魂灵与传灵塔的设想，都分别夹着纸页；最上面，却仍压着一个旧封套。

蔡须坤把它拿到面前。

``我们这次回来，有几件事想请教您。这个问题，想了很久，也还是放不下。''

她将封套打开。

里面的纸有新有旧。最早那张，页角已经磨软，几条期望对应下面，仍留着当初没能写完的空白。

展开第一项。

后面没有结果。

test\_alpha0低头看了一会儿，又翻过霍雨浩与薛峰留下的几页说明。

``你们想另外定义一套不提升的矩阵。''

``是。''蔡须坤说，``原版BMS，我们仍照它自己的规则展开。可我们总在想，能不能有另一套，让某些接在后面的结构，不再把先前认出的含义一起牵动。''

卢本韦将另一张纸推过去。

``我们试过直接去掉一些提升。开始几项看着很顺，再往后，就说不清该怎么接了。''

``夏夜星空还让我们看过，删错了以后，怎样可能出现无穷降链。''古月娜说。

她看着纸上留下的问号，语气比当初在山下平和许多。

``只写几条希望得到的对应，定不出完整的展开。这个我们已经明白，可具体该留下哪些提升，仍然拿不准。''

``所以最菜萌新让你们来问我？''

蔡须坤点头。

test\_alpha0将几页纸排好，伸手从石台另一侧取来一叠自己的草稿。

``我完成了BMS的良序证明之后，也在想这个问题。''

五人的目光同时抬了起来。

``现在，我已经有了一套方案。''

他将最上面的一页放到他们面前。

纸上写着：

UPMS（unupgrading projection matrix system）。

蔡须坤看着那几个字母，原本按在旧纸上的手，慢慢松开了一点。

test\_alpha0拿起笔。

``先看文本定义。你们熟悉的基础，还用得上。''

他写下第一段。

UPMS的表达式，以及父项、祖先、后继、极限、根等概念，同BMS。对于展开，UPMS比BMS多一些``不提升''点，其他完全同BMS。

奥斯卡读完，先在自己的纸上圈住了``多一些''三个字。

``所以，要沿着原来那些关系找，再增加不提升的判定。''

``对。''

test\_alpha0继续往下写。

第一行都按BMS提升，末列的末行都按BMS不提升。因此，只需考虑第2行至末列行数减1的那些行。从第2行开始，逐行判定一项是不是不提升。BMS认为不提升的，都不提升；其他项，也有可能被判为不提升。

蔡兴坤沿着这段文字看了一遍。

``先把BMS本来就不提升的留下。剩下的，才可能再去掉一部分。''

``嗯。已经判为不提升的，也不能到后面又随手改回来。''

卢本韦抬头问：``末列只有两行的时候呢？''

test\_alpha0用笔尖点了点那个范围。

``这一轮新增判定的行就没有了。相关部分依原来的规则处理。''

他等五人记好，又取来一张空纸。

``接下来找区段。矩阵记作\allowbreak{}\NovelInline{I06145}{A}\allowbreak{}，\allowbreak{}\NovelInline{I06146}{A(i,k)}\allowbreak{}表示第\allowbreak{}\NovelInline{I06147}{i}\allowbreak{}列、第\allowbreak{}\NovelInline{I06148}{k}\allowbreak{}行的元素。''

从根列出发，往右寻找下一个\allowbreak{}\NovelInline{I06149}{k}\allowbreak{}-父项是根列的列。

对于第\allowbreak{}\NovelInline{I06150}{i}\allowbreak{}列、第\allowbreak{}\NovelInline{I06151}{k}\allowbreak{}行的元素，以第\allowbreak{}\NovelInline{I06152}{i}\allowbreak{}列为起点，向右逐列收集不小于下面这个门槛的列：

\NovelDisplay{M05782R01}{%
(A(i,1)+1,\ A(i,2)+1,\ \ldots,\ A(i,k)+1)%
}

遇到不符合的列，就停止，且不收集那一列。所得的矩阵，记作：

\NovelDisplay{M05783R01}{%
X(i,k+1)%
}

古月娜的目光停在门槛上。

``第\allowbreak{}\NovelInline{I06153}{i}\allowbreak{}列本身，怎么通过这个加一以后的门槛？''

test\_alpha0将起始列圈了出来。

``这里以它作为起点，先收入这一列。向右延长时，才拿后面的列逐一比较。`从第\allowbreak{}\NovelInline{I06154}{i}\allowbreak{}列开始'，说的是区段从哪里起。''

他将圈后的几列用一条短线连起来，到下一列前停住。

``符合，就接着收；第一次不符合，就停在它前面。不能越过它，再去收后面看着合适的列。比较大小，仍沿用这套矩阵原有的比较约定。''

古月娜点头，将起点旁补上一笔。

奥斯卡却又指向了下标。

``这里收集时写到第\allowbreak{}\NovelInline{I06155}{k}\allowbreak{}行，得到的是\allowbreak{}\NovelInline{I06156}{X(i,k+1)}\allowbreak{}。''

``记住这个层号。稍后用到\allowbreak{}\NovelInline{I06157}{X(i,k)}\allowbreak{}时，就把刚才那段定义中的\allowbreak{}\NovelInline{I06158}{k}\allowbreak{}换成\allowbreak{}\NovelInline{I06159}{k-1}\allowbreak{}。''

test\_alpha0另写出对应的门槛。

\NovelDisplay{M05784R01}{%
(A(i,1)+1,\ A(i,2)+1,\ \ldots,\ A(i,k-1)+1)%
}

``它对应的名字才是\allowbreak{}\NovelInline{I06160}{X(i,k)}\allowbreak{}。构造的层号与判定的行号，要接在一起读。''

蔡须坤将两行对照着看了一会儿，重新抄到同一张纸上。

这一次，她先写规则，再给每一个名字找到来处。

test\_alpha0接着写道：

本次处理的行号满足：

\NovelDisplay{M05785R01}{%
2\leq k<\text{末列行数}%
}

假设已经处理完前\allowbreak{}\NovelInline{I06161}{k-1}\allowbreak{}行的提升判定，现在处理第\allowbreak{}\NovelInline{I06162}{k}\allowbreak{}行。设第\allowbreak{}\NovelInline{I06163}{i}\allowbreak{}列、第\allowbreak{}\NovelInline{I06164}{k}\allowbreak{}行的父项是根列。它上方相邻的一项必须提升，否则这一项不提升。

卢本韦用指尖按住``已经处理完''几个字。

``先知道上面怎么判，才轮到下面。''

``对。用的是前面已经得到的提升标记。''

``如果它上面不提升，这一项就先失去提升的条件。''蔡须坤说。

test\_alpha0点头。

``然后，给比较准备一个统一的数。''

设\allowbreak{}\NovelInline{I06165}{A}\allowbreak{}所有元素中的最大值的两倍为\allowbreak{}\NovelInline{I06166}{x}\allowbreak{}。

\NovelDisplay{M05786R01}{%
x=2\max A%
}

``取的是整张\allowbreak{}\NovelInline{I06167}{A}\allowbreak{}里所有元素的最大值。''他补充道，``这一点也要记清。''

他将刚才圈出的区段移到纸的另一边。

设\allowbreak{}\NovelInline{I06168}{X'}\allowbreak{}由\allowbreak{}\NovelInline{I06169}{X(i,k)}\allowbreak{}这样得到：对于

\NovelDisplay{M05787R01}{%
1\leq s\leq k-1%
}

如果\allowbreak{}\NovelInline{I06170}{X(i,k)}\allowbreak{}中某个行标为\allowbreak{}\NovelInline{I06171}{s}\allowbreak{}的项，在原矩阵\allowbreak{}\NovelInline{I06172}{A}\allowbreak{}中判为提升，就给该项增加：

\NovelDisplay{M05788R01}{%
x-A(i,s)%
}

卢本韦低头读着，指尖在原矩阵与区段之间来回移动。

``每一行用这一行对应的\allowbreak{}\NovelInline{I06173}{A(i,s)}\allowbreak{}作参照。只有原来标记为提升的那些项，才加。''

``嗯。''

``没标记的，仍留原值。第\allowbreak{}\NovelInline{I06174}{k}\allowbreak{}行及以下，也不在这一步的改动范围。''古月娜补道。

``正是。''

test\_alpha0在原矩阵旁按住笔。

``这里造的是用于比较的\allowbreak{}\NovelInline{I06175}{X'}\allowbreak{}。原矩阵\allowbreak{}\NovelInline{I06176}{A}\allowbreak{}还留在原处，之后真正的展开，仍按前面说过的规则去做。''

奥斯卡长长地呼出一口气，才将这一段记完。

到这里，他们仍没能靠名字猜出所有结果。可每一步该查什么、用前面哪一处判定，已经有了能落笔的地方。

``接下来，找比较的另一边。''

test\_alpha0在末列旁点了一下。

在包含末列的那些\allowbreak{}\NovelInline{I06177}{X(i,k)}\allowbreak{}里，取最小的那个，记为\allowbreak{}\NovelInline{I06178}{Y(k)}\allowbreak{}；它相应的\allowbreak{}\NovelInline{I06179}{X'}\allowbreak{}，记为\allowbreak{}\NovelInline{I06180}{Y'}\allowbreak{}。

蔡兴坤将目光从末列移向前面的区段。

``先看哪些区段真的包含末列，再从其中选最小的。''

``对。选中的那个，也按刚才的规则得到它对应的变换结果。后面比较时，用的是\allowbreak{}\NovelInline{I06181}{X'}\allowbreak{}与\allowbreak{}\NovelInline{I06182}{Y'}\allowbreak{}。''

卢本韦翻回自己上一页的记录，将两个带撇的符号圈在一起。

test\_alpha0写下最后的分支。

如果比较结果为：

\NovelDisplay{M05789R01}{%
X'<Y'%
}

则\allowbreak{}\NovelInline{I06183}{X(i,k)}\allowbreak{}里，所有行号大于等于\allowbreak{}\NovelInline{I06184}{k}\allowbreak{}的元素，都不提升。随后去到这个\allowbreak{}\NovelInline{I06185}{X(i,k)}\allowbreak{}的右边一列，继续寻找\allowbreak{}\NovelInline{I06186}{k}\allowbreak{}-父项是根列的列。

若不满足这个严格小于的条件，则只去第\allowbreak{}\NovelInline{I06187}{i}\allowbreak{}列的右边，继续寻找\allowbreak{}\NovelInline{I06188}{k}\allowbreak{}-父项是根列的列。

奥斯卡刚要往下抄，忽然停笔。

``两次去的地方不一样。''

``哪里不一样？''

``第一种，要越过整个\allowbreak{}\NovelInline{I06189}{X(i,k)}\allowbreak{}。第二种，只从\allowbreak{}\NovelInline{I06190}{i}\allowbreak{}的下一列继续找。''

他在纸上画出一个短短的区段，将末端圈住，又在起点旁点了一下。

``如果都只挪一列，就会重新钻进已经判完的区段。''

test\_alpha0点头。

``还有判定的范围。''

古月娜接过笔，在那条区段下面划了一道线。

``小于时，整个区段里第\allowbreak{}\NovelInline{I06191}{k}\allowbreak{}行及以下的元素，都记为不提升。前\allowbreak{}\NovelInline{I06192}{k-1}\allowbreak{}行已经判过，仍保留它们的结果。''

``嗯。之后处理更低的行，也要带着这些已经确定的标记往下走。''

蔡须坤读着最后两句，抬起头。

``如果相等呢？''

``相等就没有满足严格小于。走后一个分支，从\allowbreak{}\NovelInline{I06193}{i}\allowbreak{}右边继续。''

她把那个小于号重新描清。

``规则写到这里，才又交回原来的展开。''test\_alpha0说，``父项、祖先、根，还有BMS本来的不提升判定，都还在。新加的，是你们刚才读过的这一套筛选。''

殿内安静了一会儿。

风翻起了桌边的一角，梅天狸伸手替他们压住。奥斯卡将整段流程又读了一遍，卢本韦指着自己画的两处扫描位置，默默对照。

蔡须坤望着面前写满的几张纸。

他们当初只想让某些项别被提升。真正要把这个愿望写成规则，竟还要说明怎样找区段、怎样比较、判定会落到哪里，以及下一次从哪一列继续。

每一处，都需要一个能够继续往下做的交代。

``大长老，''她轻声问，``这套就是您最后确定下来的理想矩阵吗？''

test\_alpha0没有立即回答。

他将自己的草稿翻过来，露出背面密密麻麻的涂改。有的条件被整段划去，有的在边上重新补了箭头，几处不同颜色的墨迹，显然不是同一天写的。

``我把UPMS发布在了公开的矩阵展开器里。最初的版本，以及后来两次修订的版本，都留在那里。''

他将地址抄在纸边。

\href{https://smilelee-lyx.github.io/ne-rewritten/}{矩阵展开器：UPMS、UPMS（edit）与UPMS（edit 2）}

``你们可以去看，也可以拿自己的矩阵继续试。''

卢本韦看着那几处涂改。

``您也改过很多地方。''

``最初的版本，并不理想。''test\_alpha0说，``甚至不一定是良序的。''

蔡须坤的目光停住了。

她原以为，那叠草稿拿出来的时候，自己一直找不到的答案，便已经安稳地躺在其中。

``后来我编辑过两次，有了edit 2版本。''

他将三种版本的名字依次写在纸上。

UPMS。

UPMS（edit）。

UPMS（edit 2）。

``可是到了edit 2，也仍然不能说，这就是你们心里最理想的IBMS。''

古月娜看着最下面的名字，过了一会儿才开口。

``少一些提升，也未必就恰好得到我们想要的那些对应。''

``对。''

``原版的良序证明，也不能直接挪给改过规则的体系。''奥斯卡想起封套后面那几页下降关系，慢慢接上，``改变了哪里，仍要重新查清它会产生什么后果。''

test\_alpha0点了点头。

``把程序做出来，能展开，能看到前面的若干项，是有用的。但你们要问它能走到哪里，所有地方是否都接得住，还得继续研究。''

卢本韦靠回椅背，揉了揉握笔的手。

``我们先前总以为，是自己修为不够，或者没学到某一条诀窍。只要见到您，听到那一句，所有地方就都能接上了。''

test\_alpha0看着她，嘴角动了一下。

``有些问题，我也还在想。''

这句话让五人都安静下来。

蔡须坤低头看着旧封套。

最初那一页，仍留着``期望，尚无展开''的字样。如今旁边多了完整写下的UPMS规则，却没有哪一笔，替他们把理想的对应一并证明了。

她伸手，将两叠纸分开放好。

``我会把这些带回去，和雨浩、薛峰一起重新看。''

``你自己呢？''test\_alpha0问。

蔡须坤抬头。

``你还想做那套矩阵吗？''

她没有马上回答。

当初，那些空白曾经让她很沮丧。她明明想帮同伴，却又让几个人陪着自己在一张没有展开的纸前坐到深夜。后来重新上山，终于走完第一考，这个念头仍藏在封套里，没有消失。

``想。''

她说出来以后，声音反而稳了。

``最早是看见唐三和小舞总在前面的结构里卡住。我想，如果能有更直观的写法，他们或许能更容易看清自己走到了哪里。''

``后来我自己也越来越想弄明白，为什么同样几列，接法一变，就会牵出那么多意想不到的东西。''

她停了一下，低头看向刚抄完的规则。

``现在，我还是想让这条路更清楚一点。只是不能再只写几条好看的对应，就说它已经有了。''

test\_alpha0听完，缓缓点头。

``那就记住你为什么开始。''

他将那张写满涂改的草稿，推到新定义旁边。

``你们不必纠结定义。''

卢本韦微微一怔，又看了一眼桌上那么多条条件。

test\_alpha0也看见了她的目光。

``该查的关系、该写清的规则，仍得一处处做好。我说的不必纠结，是不必把自己困在这里，觉得一定要先找出唯一、完美的那一套，才肯再往前走。''

``能说明白的，就留下。遇到不对的地方，回头改。还没有想明白，也照实记着。''

梅天狸坐在旁边，轻轻点了一下头。

``你们在山上，就是这样一次次重新站稳的。''

蔡须坤看向她，又望回桌上的纸。

原版BMS的笔记，理想矩阵的期望，UPMS这一套能够逐条追下去的尝试，都在眼前。它们各自走到了哪里，她终于能够分清，也不必急着将其中一张盖在另一张上。

她取来一页新纸，准备把仍需核对的地方另行记下。

test\_alpha0望着她落笔，声音缓下来。

``你想让同行的人看得更明白，想让自己真正弄懂，那就带着这个念头继续做。''

他看向五人。

``遵循本心，方才是最优解。''

蔡须坤应了一声，低头又看向纸上的式子。

``那我们真正拿到一个表达式时，该怎样去理解它？''

她将最早那张期望对应取出来，放在UPMS的定义旁边。

``逐条判定，我可以照着做。可算完以后，怎样看出接上的这一段，到底带来了什么？''

test\_alpha0将纸转向她。

``你们先记住，在这里讨论的这套不提升写法里，末尾加上\allowbreak{}\NovelInline{I06194}{(1,1,1)}\allowbreak{}，表达式内部就加一份\allowbreak{}\NovelInline{I06195}{\Omega_\omega}\allowbreak{}。''

他用笔尖在末尾那一列下面轻轻点了一下，随后指向右边的参数。

\NovelDisplay{M05790R01}{%
\psi(\Omega_\omega\times\Omega)%
}

接上一列\allowbreak{}\NovelInline{I06196}{(1,1,1)}\allowbreak{}，便读作：

\NovelDisplay{M05791R01}{%
\psi(\Omega_\omega\times\Omega+\Omega_\omega)%
}

``就是这么简单，不需要太复杂的技巧。''

卢本韦顺着他的笔尖看过去。

``加在里面。外面那一层\allowbreak{}\NovelInline{I06197}{\psi}\allowbreak{}仍留着。''

``对。再接一列，就再加一份。''

\NovelDisplay{M05792R01}{%
\psi(\Omega_\omega\times\Omega+\Omega_\omega\times2)%
}

test\_alpha0将两处新添的部分圈出来。

``你们要回归本源。先认清每一段在表达什么，再看它们怎样接起来。''

蔡须坤看了很久。

当初她拿出那张纸时，最希望的，也不过是让后面接上的一一一，清楚地带来这一份\allowbreak{}\NovelInline{I06198}{\Omega_\omega}\allowbreak{}。后来，他们试着解释的东西越来越多，最先想问的那件事，反而常常被压在许多页草稿下面。

古月娜将旧纸挪近一点。

``所以，读它的含义时，可以先抓住这个最直接的接法。具体怎么展开，还是回到前面那些判定去查。''

test\_alpha0点头。

``一个帮你理解，一个让你知道每一步怎样落下去。''

奥斯卡在旁边补记了几笔。蔡兴坤则将三行式子从上到下读了一遍，终于轻轻笑了。

``这样再拿给小舞看，她就能先知道我们想接上的是什么了。''

蔡须坤也笑着点头，将那张纸单独夹好。

殿内安静片刻，卢本韦却没有立即收笔。

她抬头看着test\_alpha0，像是还有一个问题，已经在心里放了很久。

``大长老殿下。''

``说。''

``您说过的双魂核，固然强大。可帝天也有阴阳互补双魂核，最菜前辈也是。为什么您当年能那么快击败帝天？''

她稍稍停了一下，又接着问道：

``还有，您能承接大长老这个位置，修为的根基，究竟与他们有什么不同？''

test\_alpha0将笔放下。

梅天狸也转过目光，静静听着。

``你们知道阴阳互补双魂核吗？''

五人点头。

``我们学过。''奥斯卡答道，``两颗魂核相互补充，也要维持平衡，才能一直运转。''

``整个大陆，能真正凝聚成功的，也只有几个人。''

test\_alpha0抬起手，两点微光在掌上显现，彼此牵动，缓缓流转。

``帝天有，最菜萌新也有。他们能够接续力量，持续出手，双魂核是其中很重要的根基。''

他停了一下。

``可我的魂核，是三个。''

第三点微光随之亮起。

卢本韦的视线立刻落了过去。

她想起最菜萌新在山脚说过的那句话------世间唯一的同位共振三魂核。

那时，她已经知道大长老有三颗。可听旁人说起，与此刻看见他亲手示意，仍是两种全然不同的感受。

test\_alpha0将三点微光稍稍展开。

``这是个投影，让你们看清它们之间的关系。''

三团光各有轮廓。

其中一团明显更大，另一团小一些，剩下一团则又小了一圈。五人很快看清，那三颗魂核的大小，竟没有任何两颗相同。

蔡须坤微微坐直。

``大小也不一样？''

``每一个，都不一样。''

test\_alpha0答得平静。

三点微光在掌上持续流转，彼此牵引。较大的那一点先动，另外两点便随之接上；力道再传回来，原先那一点也随之调整。

一轮过去，又是一轮。

谁也没有被另一团光吞进去，彼此之间的牵动，也没有在往复中越积越乱。

古月娜望着那片微光，指尖慢慢停在桌沿。

她记得最菜萌新说过，一旦添上第三颗，每一颗的变化都会牵动另外两颗，结果极难控制。如今再加上大小不等，维持这种运转的难处，便更难只用``两颗互补''的经验去想象。

``三颗，都在参与同位共振？''她问。

``都在。''

test\_alpha0抬眼看向她，又望向另外几人。

``它们大小不同，牵动彼此时的分量也不同。可我能让它们稳定、持久地运行，将三颗的力量接续起来。''

蔡须坤望着那三点光，忽然想起自己在训练场边，怎样都收不回那股混沌时的感觉。

她刚学会同时接住两种领域，就已经几次在换力时出了岔子。眼前的三颗魂核，却在一轮轮彼此牵动之后，仍安稳地留在各自的位置。

test\_alpha0的手掌轻轻一合，投影随之淡去。

``对我而言，这不是只能支撑片刻的一次爆发。持续运转，连续出手，它们也一样能稳住。''

卢本韦低头看向自己的笔记。

阴阳互补双魂核下面，她又添了一行。

同位共振三魂核。

写到这里，她的笔尖停了一下，随后将刚听到的新细节，也认真记了下去。

三颗大小各不相同，仍可长期稳定运行。

\chapter{第一百三十六章：三核之间}\label{ux7b2cux4e00ux767eux4e09ux5341ux516dux7ae0ux4e09ux6838ux4e4bux95f4}

卢本韦将最后几个字写完，笔却仍停在纸上。

三颗大小各不相同，仍可长期稳定运行。

她抬头看向test\_alpha0。

``大长老殿下，最菜前辈说过，第三颗魂核一旦加入，每一颗的变化，都会同时牵动另外两颗。它们又不一样大，究竟怎么才能一直稳住？''

``他当年看见的时候，也惊讶不已。''

test\_alpha0将手重新抬起，方才淡去的三核投影再次显现。

``三颗彼此牵动，偏一点，传回来的变化又会继续累积。过去许多魂师想到这里，便认为这是三体运动那样的无解难题，最后总要走向混沌。''

三点微光在掌上缓缓流转。

``可在我体内，我依靠良序神力，让这三个魂核稳稳运行。''

蔡须坤的目光停在了他掌心。

她记得最菜萌新说过，这股力量弥补了三核之间本来极难控制的牵引。如今亲眼看见那一轮轮接续，她才有了更清楚的感受。

梅天狸在一旁开口。

``这件事，亲自试过，就知道有多难。''

五人转过头。

梅天狸抬起手，将三团魂力凝在身前。她没有将它们引入体内，只让它们各自压紧，再用细细的魂力牵连起来。

``这是仿照魂核做的试运转。''她说，``我没有真正凝成三魂核。你们看它们怎样互相影响就好。''

最初，三团魂力都停在她控制的位置上。

一团向前移动，另外两团便随之偏转。梅天狸指尖轻轻一动，便将那点偏差接了回去。

``这样看，不像很难。''奥斯卡低声道。

``刚开始，是这样。''

梅天狸没有撤手。

她让其中一团稍稍增大，另外两团与它之间的牵引随之改变。一条魂力细线骤然绷紧，原本平缓的移动快了一点。

她及时压住，另一边却又偏开。

蔡兴坤看了一会儿，慢慢皱起眉。

``您刚稳住这里，那里又变了。''

``嗯。每一次调整，也会成为下一次变化的一部分。''

三团魂力仍在转动。

梅天狸的手始终很稳，送进去的力量却越来越密。她已经不能只顾着最初那点偏移，而要同时接住三处不断传回来的牵动。

几人说话之间，时间一点点过去。

一条细线忽然颤了一下。

紧接着，另一团魂力也向内一沉，原本留出的空隙迅速缩小。梅天狸将它拉回，最外侧那团又被带得偏离了先前的方向。

``几分钟而已。''她说。

她抬起另一只手，将三团魂力依次隔开，缓缓散去。

``以我的魂力去同时控制三个魂核，至多也只能维持一时。眼前只是仿照它们相互牵引来做的试运转，几分钟就会走向混沌和无序。我的力量越往里压，需要处理的变化反而越多。''

殿内静了片刻。

奥斯卡看看她刚才收回的手，又看向大长老掌上依旧平稳的三点光。

那边始终没有停。

古月娜低声问：``三颗如此，两颗之间的变化，也一样不能只顾眼前？''

``当然。''梅天狸说，``这也是许多98级斗罗，面对99级极限斗罗的门槛时，觉得无从着手的难题。''

卢本韦坐直了些。

``可突破99级，不是只要求两个魂核吗？''

``是。99级要求的是双魂核。''

梅天狸答得很清楚。

``我用三处试运转，是让你们看清相互牵动会怎样积累。到了两颗真正的魂核之间，也得解决怎样凝聚、怎样接续、怎样长久运行的问题。''

``单靠自己不断补上一股魂力，将它们暂时压住，和它们已经能够在体内持续循环，是两回事。''

蔡须坤低头看了一眼自己的手。

她在训练场边，曾拼命想把翻涌的灰色压回去。可越急，换力时便越容易失去节奏。梅天狸刚才掌上出现的变化，让她想起了那时的感觉。

``所以，看着只是多了一颗。''她缓缓说道，``可真正稳住以后，出手与恢复，都和原来不同了。''

梅天狸点头。

``两颗能够长期接续，战斗时能调动的力量、持续出手的能力，都会大不一样。98级与99级，看上去只隔一级，战斗力的差距却非常大。''

她看向test\_alpha0掌上的投影。

``双魂核已经让那么多人止步。他却能让三颗大小不等的魂核，始终处在稳定运转之中。''

test\_alpha0听到这里，将三点光稍稍移近。

``所以，最菜萌新当年才会惊讶。他自己走过阴阳互补双魂核那一关，知道添上第三颗以后，原先许多已经调好的关系，都要重新承接。''

``良序神力，就用在这里。''

一缕极淡的蓝光从三团微光之间显现。

它顺着彼此牵动的位置缓缓流过。较大那团的变化传向另外两处，另外两处又将各自的力量接续回来，细微的偏转在往复之间被稳稳调回。

``我让它参与三核之间的调控。每一次变化，都要能接上下一次；接回来以后，还得继续稳住。''

古月娜望着那缕蓝，终于轻轻点头。

``外面看见的是三颗一起转，里面却一直在调整。''

``对。''

test\_alpha0将其中最大的一团光单独照亮。

它的轮廓顿时清楚起来，另外两团显得小了许多。更深处有一层层细密的纹路展开，又在光中彼此衔接，远比方才看见的一个光点复杂。

卢本韦的视线落在那里。

``第三颗？''

``我的第三魂核，拥有来自Laver table的力量。''

test\_alpha0说着，将那片纹路稍稍展给他们看。

``也就是Laver表。第三魂核因此比另外两个核大很多，这也是我的力量来源。''

蔡须坤抬起头。

天州城外的蓝色漩涡、被大长老反复提起的那一剑，还有良序之神大供奉的身份，在这一刻，与面前的第三颗魂核接上了关系。

``和您当年在天州动用的，是同一个源头？''

``同一个源头。''

test\_alpha0点头。

``你们先前听我讲过，我从Laver表中裁削出iblp，去理解并承接其中一部分力量。第三魂核接住的，也与这个源头有关。''

他让三团微光重新一同运转。

``第三颗接入的力量很强，另外两颗也要能与它接续。它大得多，牵动另外两颗时，维持平衡就更困难。良序神力让这种运转能够一直持续下去。''

梅天狸望着那三团光，缓缓呼出一口气。

方才她只是短暂试运转，便不得不停下来。而面前这三颗，不仅数量更多，其中一颗承接的力量，竟还来自Laver表。

奥斯卡低头看着自己记下的``双魂核''，在旁边又留出了一片空处。

``所以，帝天即使有阴阳互补双魂核，也不能只拿魂核数量来和您比较。''

``要看核里承接的是什么，也要看能怎样把它用出来。''test\_alpha0说。

他看向五人，语气仍然平静。

``三核的接续，让我能持续调动力量；良序神力与第三魂核所承接的源头，又让这股力量有了你们在天州往事里听见的分量。''

``不过，体内的循环稳住，不等于Laver表的完整力量，我便能随意放出、再毫无遗漏地收回来。''

卢本韦的笔尖停了一下。

她想起大长老讲过的那把不知道有多长的刀。

``核里能承接的，仍是您已经能够容纳的那一部分。''她说。

``嗯。''

test\_alpha0收近掌上的光。

``你们现在看到的，是我体内能够稳定持续的运转。更深处的神力展开以后，会影响到多大的范围，那又是另一件需要弄清的事。''

蔡须坤看着最大那团光，又望向围绕它接续运行的另外两团。

她终于不再只盯着最亮、最大的那一处。

在那缕蓝光流过的地方，三颗魂核都在参与，每一处细小的调整，都接着下一次变化。

梅天狸没有再做试运转，只将双手收回膝上，安静看了一会儿。

卢本韦则低下头，在新记的那一行旁边，认真补上了第三魂核的来源。

Laver table。

写完以后，她抬起眼睛。大长老掌上的三点微光，仍在平稳地转着。

test\_alpha0将掌上的光收起，目光从卢本韦移向蔡须坤。

``卢本韦，蔡须坤，你们几个人的扽魂，品质都很高。''

他说着，向蔡须坤稍稍点头。

``尤其是你体内的混沌之力，一定不普通。你如今能够引出的这一部分，已经显出远超自身修为的分量。它真正的来头，还值得继续弄清。''

蔡须坤轻轻按住自己的手腕。

那股力量此刻安静着，可训练场边失去控制的感觉，她仍记得清楚。

test\_alpha0又望向卢本韦。

``帝天的血脉，论品级，甚至比不上你的神级母猪扽魂，又何谈与我的良序之力相比？''

卢本韦抬起头，眼中露出一点恍然。

从前，她听见帝天的名字，想到的总是那头连梅长老都无法挡住的凶兽。如今听大长老这样一层层说来，她才将血脉品质、魂核与实际修为，分辨得更清楚。

``不过，你们绝不要小瞧他。''

test\_alpha0的语气郑重下来。

``就算是当年在天州受了重伤，也远没有到一个普通99级极限斗罗，就能轻易将他打败的地步。''

古月娜抬眼看向他。

``漫长岁月积下来的修为，淬炼出的肉身，还有他对力量的掌握，都在那里。神级扽魂让你们有很高的起点，真正能用出来多少，还要看你们自己。''

卢本韦认真点头。

``我们明白。''

梅天狸在旁没有插话。她曾在森林里直面过帝天，也亲自接过五人今日的合击，知道这两者之间仍有多远。

蔡须坤安静了一会儿，伸手将桌边另一个封套取来。

``大长老殿下，说到人类和魂兽，我们还想和您谈一件事。''

她将封套里的纸抽出。

最上面那一页，写着三个字。

传灵塔。

test\_alpha0的目光落了下去。

``我们想到了一个完美的解决方案。''蔡须坤说，``创办传灵塔，用魂灵的方法，让人类和魂兽能够一起发展。''

这句话说出口时，她的声音比刚才更清晰。

在海上，她曾与古月娜谈过；到了学院，也已经把想法说给老师们听。如今真正坐在大长老面前，她终于将这几页纸，推到了他手边。

奥斯卡把放在下面的几张说明依次展开。

``愿意共生的魂兽与魂师，可以通过传灵塔彼此了解。魂兽保留自己的意识，成为共同修炼的伙伴。塔里也能有人继续研究，把能够稳定做到的办法整理下来，让后来的人有东西可学。''

``各地遇到的事情，也可以在那里传递消息。''蔡兴坤接着说，``哪里愿意合作，哪里发生了争端，总得有人听见，才能一起想办法。''

卢本韦将其中一页转向大长老。

``关于魂灵，我们已经有了真正做成的经历。可要让更多人和魂兽一起走，仍要把各自的意愿、能承受的力量，还有怎样共同生活，都谈清楚。''

她指尖停在纸面上，声音慢下来。

``这些，也可以成为传灵塔要做的事。''

test\_alpha0没有打断他们。

他拿起那几页纸，将关于魂灵的部分从头看了一遍。纸边留着不少补记，有的字写得急，有的显然是在赶路时添上去的。

梅天狸看着他翻页，轻声道：``他们在学院，已经和我们说过这份计划。''

蔡须坤点头，又望向大长老。

``我们想让人类修炼时，多一条能够与魂兽共同成长的路。魂兽也能继续活着，有自己的族群与栖息地。这样一来，许多因猎杀而结下的仇，才有机会慢慢少下去。''

古月娜将目光从纸上抬起。

``所以，我们还要去一趟星斗大森林。''

殿内的风声似乎清楚了一些。

test\_alpha0翻页的手停住了。

``见帝天？''

``也要见愿意听我们说话的魂兽。''蔡须坤答道，``人类与魂兽的事情，需要两边都坐下来。传灵塔想怎样办，魂灵如何相处，它们也应当能够提出自己的想法。''

古月娜轻轻点头。

``各族的顾虑未必一样。去到森林里，才能当面听它们说。''

卢本韦看了同伴一眼，将笔记合上。

``我们想把做成的带过去，也把还需要一起解决的地方带过去。''

test\_alpha0抬起头，看向这几个坐在自己面前的年轻人。

方才，他们还在追问他为何能够击败帝天。如今摆到桌上的，却是一份准备走进星斗大森林、与魂兽商谈共生的计划。

他将写着``传灵塔''的那页纸放在最上面，用手掌轻轻压平。

蔡须坤望着他的手，继续说道：

``这一次，我们想去那里，和它们好好谈一谈。''

test\_alpha0没有立刻回答。

他的手仍压在那页纸上，过了一会儿，才望向卢本韦。

``你们家，也有人这样准备过。''

卢本韦抬起眼。

``我曾祖？''

``比承业还要早。''

test\_alpha0轻轻摇头。

``卢守衡的父亲，和我是同年龄段的人。年轻的时候，我和他一起游历过山川河水。那时，我们还没有如今这些身份，走到哪里，便在哪里修炼，找人切磋，也见识各地不同的扽魂。''

提起那个人，他脸上浮出一点很淡的笑意。

``在封号斗罗之前，他的修炼速度，甚至比我还要快一点。''

蔡兴坤原本搭在桌沿的手顿住了。

几人方才才听过三颗魂核的来历，再看面前的大长老，很难想象当年竟有人能在修炼上走到他的前面。

卢本韦低声问：``也是母猪扽魂？''

``当然。''

test\_alpha0望着她。

``你要知道，母猪是猪神Opt的武魂。神级武魂带来的优势，从开始修炼时就已经有了。他的天资又好，那时候，我们都觉得，他往后一定能走得很远。''

卢本韦将手中的笔慢慢放下。

族谱上那些隔了许多辈的名字，她从小便知道。可有些人的一生，只剩下几行字。她从来没有听谁说过，守衡的父亲曾经与大长老并肩游历，甚至修炼得比他更快。

``可惜，成就封号斗罗的那一次，他太想再往前走了。''

test\_alpha0脸上的笑意淡下去。

``他想在那个时候，就提前凝聚阴阳互补双魂核。''

奥斯卡看向桌上还没有收起的三核笔记。

``提前去碰极限斗罗才需要迈过的那道关？''

``嗯。他看到了两个魂核相互接续的可能，也相信自己能够稳住。可真正开始凝聚以后，魂力之间的牵引，比预想的猛烈得多。''

test\_alpha0的目光停在笔记上的两个圆圈之间。

``他最终保住了命，也成为了封号斗罗。但那次魂力反噬，给他留下了病根。''

梅天狸没有说话，只将刚才用来演示的手缓缓收紧。

五人却都记得，那三团魂力失去平衡时，变化来得有多快。体外的试运转尚可及时散去，若那一切发生在经脉与魂核之间，付出的代价便完全不同。

``后来，他修炼到了98级，在这一关前停了一百来年。''test\_alpha0说，``以他的天资，本来有望成就极限斗罗。可那处暗疾一直跟着他，平日里能够维持，真要再次凝核，便成了绕不过去的阻碍。''

卢本韦看了眼自己的掌心。

母猪扽魂此刻没有显现，她却下意识地收拢了一下手指。

同样的武魂，同样让人期待的天赋。她刚听完神级品质的优势，紧接着便听见了这样一个结局。

``他后来还在研究魂灵？''古月娜轻声问。

``在研究。''

test\_alpha0点头。

``他曾秘密去见过凶兽。那一趟我没有同行，有些经过，是他回来后才告诉我的。他想知道，魂兽愿不愿意与人合作，也想知道，Opt留下来的魂灵构想，究竟能不能做成。''

古月娜的目光微微一动。

``所以，那头柔骨兔\ldots\ldots''

``正是他和凶兽谈判、合作的结果。''

卢本韦的呼吸停了一瞬。

她听大长老讲过那头兔子，八万年的修为，在卢家后山住了四十多年，自愿与卢守衡相合。可直到此刻，她才知道，那段相处的起点，还要追溯到更早的一辈人。

``它愿意来到人类身边，在你们家的后山住下。这件事，本身就已经很难得。''

test\_alpha0说得很慢。

``有了多年的相处，后来又得到了它的同意，他们才决定往下试。''

``动手融合的是守衡。守衡的父亲则参与了前面的研究，也把这些年谈下来的合作与信任，交到了儿子手里。''

蔡须坤望着纸上的``传灵塔''，没有出声。

那三个字下面，他们已经写下许多设想。愿意合作的魂兽，愿意尝试的魂师，可以让彼此坐下来说话的地方。

原来，一百多年前，已经有人一步步做到了其中的一部分。

``头四天，一切都好。''

test\_alpha0垂下眼睛。

``守衡能够听见它的想法，它也能回应守衡。那几天，他父亲来看他们的时候，脸上的神色，我到现在还记得。''

卢本韦没有追问第五天。

她知道第五天发生了什么。

八十七级的身体，终于承受不住那股八万年的魂力。人和魂兽都愿意，却仍没能从反噬里活下来。

``谁能想到，这场他亲自促成的合作，最后竟会要了他儿子的命。''

大长老的手从那页纸上移开。

``守衡死后，他一下子垮了。''

殿内静下来。窗外有风吹过，几页纸的一角轻轻掀起，又落回桌面。

``他去翻笔记，从最初的构想到最后一天的记录，一页一页地翻。找出一处可疑的地方，便以为也许就是那里。可看得越多，便越想起守衡已经不在了。''

test\_alpha0停了片刻。

``我去看他的时候，他问过我，若当年没有去见凶兽，没有把那头兔子带到人类身边，是不是就不会有这件事。''

古月娜的手指轻轻蜷起。

她想起了那个在后山住了四十多年的生灵。它也曾相信这场合作，愿意与一个人分享自己的生命。

``可兔子也死了。''她低声说。

``是。''

test\_alpha0看了她一眼。

``他知道。所以，他连自己早年那趟满怀希望的出行，都不敢再想。''

卢本韦的手移到行囊旁，摸到里面那本旧笔记的边缘。

她记得卢景川十三岁时，抱着它站在父亲的葬礼前。如今再往那幅情景后面看，那里还有一个刚失去儿子的老人。

``过了几个月，他也走了。''

test\_alpha0的声音很低。

``旧疾本就缠着他。守衡的死，又让他陷在悲痛里，再没能缓过来。''

卢本韦捏着行囊带子，很久没有动。

一场失败，带走了卢守衡与那头柔骨兔。几个月后，又带走了那个曾经走进森林，把合作的希望带回卢家的人。

而那时的卢景川，才十三岁。

蔡须坤望向同伴，伸手轻轻碰了碰她的手背。

卢本韦慢慢松开攥紧的带子，将旧笔记取出来，放在那几页新写的计划旁边。

两样纸的颜色相差很大。旧纸已经发黄，边缘磨得起了毛，新纸上的字迹还很清楚。

``所以，您听见我们要进星斗大森林，就想起了他。''

``嗯。''

test\_alpha0望着那本笔记。

``双魂核也好，魂灵也好，我都亲眼见过有人带着很好的天赋、很好的愿望走进去，最后没能回来。''

卢本韦低下头，将笔记翻开。

她没有急着翻到已经熟悉的结论，只扶住了那道被许多人翻过的书脊。

奥斯卡把桌上的纸向旁边移开一点，给它留出完整的位置。蔡兴坤也将原本压在上面的封套取走。

蔡须坤看着旧纸，过了一会儿才开口。

``我们会把这些事记住。''

她的手仍停在卢本韦旁边。

``见到它们的时候，我们也想听一听，那头柔骨兔的族人，是怎么记得这件事的。''

蔡须坤安静了一会儿，目光重新落到三核笔记上。

``大长老殿下，您的三颗魂核，都是什么时候凝聚的？''

test\_alpha0伸出一根手指。

``第一颗，79级。''

卢本韦重新拿起笔。

``第二颗，89级。第三颗，是在98级的时候凝聚的。''

三个等级依次落在纸上。蔡须坤看着最后那个数字，抬头望向他。

``也就是说，您在突破极限斗罗之前，就已经开始凝聚第三颗了？''

``是。不过，凝聚第三魂核的时候，我动用了良序之神的神力作为辅助。''

test\_alpha0的手指轻轻点在``98级''旁边。

``神力参与了凝聚，也帮助我稳住新魂核与另外两颗之间的牵引。你们刚才看见的长期运转，是从那一步走过来的。''

梅天狸看了一眼那三个数字。

她先前只演示了几分钟，便不得不散去三团魂力。真正让第三魂核在体内成形，还要让它从此持续运行，显然比那场试运转困难得多。

卢本韦在最后一行旁边，认真补上了神力辅助的说明。

``那守衡的父亲呢？''她问，``他尝试双魂核之前，第一颗是什么时候凝聚的？''

``69级。''

笔尖停住了。

test\_alpha0看着她，缓缓说道：``我们一起游历的那一年，他已经凝聚了第一颗魂核。那一年里，他的修炼速度确实比我还快一点。虽说论年纪，他还比我大三岁。''

卢本韦往前翻了一页，将这个等级添在那位先辈的经历旁。

69级，第一魂核。

后来成就封号斗罗时，尝试阴阳互补双魂核，遭到反噬。

前后两行靠得很近，她却看了许久。

``第一颗凝成以后，给了他很大的帮助。''test\_alpha0说，``也正因为有过那次成功，他后来才会觉得，第二颗也能够提前试一试。''

蔡须坤低头看向自己的手。

她能感觉到体内魂力缓缓流过，混沌之力此刻也很安静。可前一天在训练场上，那股力量骤然失控的感觉，仍然清楚得很。

``凝成第一颗，并不意味着第二颗也一定能顺利凝成。''她轻声道。

``嗯。多出一颗，要处理的牵引就会变。''

test\_alpha0望向五人，停了片刻。

``而你们的等级，实际上已经接近69级了。''

蔡兴坤抬起头，奥斯卡也下意识地坐直了一些。

69级。

方才听来，那还是一位先辈年轻时走到的地方。直到大长老亲口指出，他们才将这个数字与自己联系起来。

蓬莱岛上那一年的攀登，回到大陆后遇见的战斗，以及训练场上拼尽全力的合击，都已经留在了他们如今的魂力里。

卢本韦看向身旁的同伴，握笔的手慢慢收紧。

``那我们也快到能够考虑第一颗魂核的时候了。''

``可以开始了解，开始准备。''test\_alpha0说，``至于何时真正动手，要看你们各自能不能承受，能不能控制。''

蔡须坤点点头，将自己的笔记移近。

她先记下了大长老的三个等级，又在旁边写下69级。写到这里，她没有立刻往后添新的目标，而是翻回了前面关于魂核相互牵动的那一页。

这一次，她看得比方才更仔细。

\chapter{第一百三十七章：天州纪念馆}\label{ux7b2cux4e00ux767eux4e09ux5341ux4e03ux7ae0ux5929ux5ddeux7eaaux5ff5ux9986}

离开总殿时，蔡须坤将那本记着魂核的册子放进了行囊。

79级、89级、98级，三个数字下面，还留着她最后补记的几行字。她没有继续往后写，将册子收好，又检查了一遍夹在另一个封套里的传灵塔计划。

古月娜已经站在阶下。

她将地图摊在一块平整的石面上，指尖从总殿所在的位置向西划去，停在了一座城旁边。

``我们从天州过去。''

卢本韦俯下身，看清了她指着的两个字。

``天州。''

这个名字，他们听过太多遍了。

良序圣殿上，大长老讲过那里的一场兽潮。蔡须坤还记得，更早些时候，phyrion也提起过，年轻时送猪路过那座城，在城外看见了一个方圆两三里、几十丈深的坑。

坑底有一座石像。

她将那几句话转述给了同伴。那一夜，她还忙着照顾伤员，听见这些，也只能在心里勉强描出一个轮廓。

``过城以后，我来带路。''古月娜说，``先沿官道走，到林区再换方向。''

奥斯卡将最后一袋干粮扎好，抬手掂了掂。

``东西齐了。''

蔡兴坤把黑子鸡从石栏上抱下来。它原本正歪着脑袋看地图，被她抱起后，又探出头，朝地图上的天州啄了一下。

``你也要去？''卢本韦问。

黑子鸡叫了一声。

``它哪次没去。''蔡兴坤将地图往旁边挪了挪，免得真被它啄出个洞。

梅天狸送他们走下最后一段长阶。

``到了森林里，彼此照应着。''她看向五人，``有什么拿不准的，先停下来商量。''

几人认真应下。

蔡须坤最后回头望了一眼。总殿的高处仍浸在晨光里，白玉柱之间，有一线很淡的蓝。

她将行囊带子收紧，跟上了同伴。

这一次，前面带路的是古月娜。

头一天，他们走过了城郊与大片田地。古月娜不时停下来核对岔路，等后面几人跟上，再继续向西。

入夜以后，奥斯卡在避风处架起小锅。卢本韦取水回来，将水袋挨个放在石边；蔡兴坤则蹲在火旁，帮着将潮湿的柴翻到上面去烘。

蔡须坤借火光把地图看了一遍。

``明天过这条河？''

``嗯。''古月娜在她身边坐下，``桥在上游。走桥，少绕些路。''

第二天，他们跟着运货的车队过了桥。第三天，官道渐渐宽起来，沿途驿站也多了，往来的车上，有不少都挂着天州商行的旗。

到了第四天，远处的地平线上终于显出了一片城楼。

最初只有几道模糊的线。再往前走，高处的檐角、城门上的旗，还有两侧延伸出去的墙垛，才一处处清楚起来。

蔡兴坤停下脚步，仰头看了一会儿。

``那就是天州？''

古月娜点头。

官道两旁已经有了连片的房屋。车马还没有到城门前，叫卖声便先传了过来。茶棚里坐满了歇脚的商旅，几名搬货的伙计正合力将木箱抬上车，车夫拉着缰绳，给迎面而来的行人让路。

五人随着人流走近城门。

十丈高的城墙向两侧延伸，门楼还要更高一些。宽阔的门道足够几辆货车并行。两边守卫分开查验，登记、放行，身后的队伍仍在缓缓向前。

卢本韦抬头望向城门上方。

天州。

两个大字刻在石额上，底下是一道道车辙磨出来的浅痕。

等真正走过门洞，他们的脚步又慢了下来。

城里的街道沿着两侧铺开，石板平整，街边店铺一间挨着一间。绸缎铺将新到的布匹挂在檐下，药材行门口晒着一排竹匾，几家酒楼的招牌伸到街面上，伙计站在楼梯口招呼客人。

更远些的地方，还有新修的高楼。

楼与楼之间架着魂导灯，白日里没有亮，灯座却擦得很干净。运货的车从街中走过，行人熟练地向两边让开，待车过去，又重新汇成了一片热闹的人流。

``我原先总想着那堵挨打的城墙。''奥斯卡说。

他站在一间食铺外，望了望街口，又看向更远处的楼阁。

``没想过里面会是这样。''

蔡须坤也有同样的感觉。

大长老讲起天州时，门闩断了，炮管打红了，守城的人蹲在城墙后，等着下一次冲撞。她记得那些细节，却很少在那段话之外，去想这座城后来怎样过日子。

如今，蒸笼揭开，白气扑到了街上。

一个小孩子拉着母亲的衣角，仰头盯着刚出锅的包子。旁边的伙计将一碗热汤放上桌，碗底碰着木板，轻轻一响。

蔡兴坤看了一眼，怀里的黑子鸡也跟着看。

``先吃点东西？''她问。

奥斯卡已经向里望去。

``这家还有位置。''

五人在靠窗的桌边坐下。他们将行囊收在脚旁，店家很快端来热汤与面食，又指着墙上的木牌，问要不要添些小菜。

卢本韦点了两样。

蔡兴坤捧着碗喝了一口，肩膀慢慢松下来。

``这几天总算吃上一顿不用蹲着的了。''

``昨晚那块石头也还行。''奥斯卡说。

``你坐的那块平。''

卢本韦接过筷子，忍不住笑了一声。

古月娜坐在窗边，看着外面经过的人。

一辆卖花的板车停在街角，几个孩子围在旁边，挑来挑去，最后只买走了一小束。车上的老人也不催，弯下腰替他们将细绳系紧。

她看了片刻，才端起自己的碗。

近一百七十年过去了。

那场兽潮里的许多人，早已经不在。城中的屋舍、街道、门楼，却在一代又一代人的手里扩了出去，重新变成了眼前的模样。

吃过东西，五人没有急着离开。

他们顺着街道走了一段，补足水与干粮，又向路边的商户问清西门的方向。

``去旧战场？''那人听他们问起城外，抬手向西指了指，``出了门，沿路往前。到那里就看得见了。''

蔡须坤向他道了谢。

临到西门，街边仍旧有许多铺子。一个修车的师傅蹲在墙角敲着轮轴，面前摆了好几根新削的木条。旁边树下，有几个老人正围着棋盘争论。

直到跨过门洞，那些声音才一点点落到身后。

城外的路比官道窄些，两侧长着低矮的树。再往远处，仍有田地，田埂上有人弯腰干活。

蔡兴坤往前看了看。

``纪念的地方，就在这边？''

``应当是。''卢本韦说。

前方的地势略微抬起，路沿着一道缓坡向上。坡顶修着石栏，几个行人正在那里停步，原本并肩走着的两个人，也先后站住了。

古月娜走上坡顶。

她的脚步忽然停了。

蔡须坤随即跟上，刚越过最后一段树影，眼前便骤然空了下去。

那一瞬，她几乎忘了再往前走。

一个巨大的坑，铺在了五人面前。

他们站在近处的边沿，向前望去，目光先落到陡然下沉的坑壁上。黑亮的岩土一层层向下，到了深处，连光也暗了一截。两侧的弧线远远伸开，一边没入树后，一边绕向更远的坡地，一眼扫过去，竟看不全整个边缘。

蔡兴坤抱着黑子鸡的手收紧了。

``这么大\ldots\ldots''

没有人立刻接话。

他们在山上看过深谷，也曾在训练场上见过强大魂力留下的破坏。可眼前这道凹陷太整齐，近处的坑壁又太深，站在石栏边，连刚才还显得宽阔的道路，都忽然窄了下去。

风从下面吹上来，带着一点凉意。

卢本韦走近半步，扶住栏杆。

``方圆两三里。''

她想起出发时蔡须坤转述的旧话，声音放得很轻。

``几十丈深。''

那时只是听见两句话。如今真正站到了坑前，她才知道，自己原先想的还是小了。

黑色的坑壁上，有几条细裂。裂隙里生着草，风一吹，细叶便贴向岩面。更低处，一片光滑的黑仍然映着天，近一百七十年的风雨，也没有将它完全磨去。

``当年那一下，就在这里撞上的。''奥斯卡低声说。

神圣审判。

龙神爪。

蔡须坤记得大长老讲到这里时，坐在石台前，手只轻轻比了一下。

可眼前整片土地，都被那次相撞压了下去。

她抬头向后望。天州城的门楼还在路的另一头，旗子被风吹起，城中的热闹声已听不大清楚。

这段距离，就是当年被挡在城外的那一段。

她再转回身时，看见了坑底正中的石像。

方才整片下陷的地面占住了视野，他们竟一时没有顾得上往那里细看。

现在，那道青灰色的身影，清楚地立在了光里。

雕像比天州的城楼还高。长袍的下摆垂到基座上，几道衣褶沿着石面落下，深处留着岁月积下的暗色。双手稳稳拄在剑柄上，长剑向下，剑尖立在身前。

他面对着西方。

那是星斗大森林所在的方向。

``大长老\ldots\ldots''蔡兴坤喃喃道。

石像的眉骨、鼻梁、唇角，都刻得极细。远看时，五人先认出了那身形；等沿着石栏向旁边走，避开斜照的光，那张脸便完整地显了出来。

与他们前些天在良序圣殿里见到的test\_alpha0，几乎一模一样。

卢本韦抬起眼睛，停了很久。

殿上的大长老会翻他们的笔记，会停下来等他们问完，也会提起年轻时结伴游历的人。眼前这座石像，却将他的神情永远留在了那里。

平静，垂着眼，手压着剑柄。

那双石刻的手，比他们整个人都大。

奥斯卡沿着雕像往下看，忽然指了指基座。

``下面还有一座馆。''

厚重的基座最下方，开着一道门。门前有石阶，几盏灯已经亮起来，将门楣上的字照得很清楚。

天州兽潮纪念馆。

通往坑底的步道沿一侧缓缓下降，护栏一直延伸到馆前。有人正从下面往上走，手里拿着薄薄的册子，到了拐弯处，又回头望了一眼石像。

``那边还有旧纪念堂。''卢本韦说。

坑沿另一侧，几棵大树围着一座较矮的建筑。门前的石阶被踩得发白，堂外摆着新换的花。

蔡须坤还说过，那里刻着当年死去的人的名字。

他们没有急着往坑底走，先沿石栏站了一会儿。

一位老师领着几个孩子经过，孩子们起初还在说话，到了栏边，声音也渐渐小下去。最矮的那个踮起脚，想从两根石栏之间看见下面，老师便俯身牵住了他的手。

``就是这里？''孩子问。

``就是这里。''

老师指向坑底，又指向远处的城门。

孩子顺着她的手看过去，久久没有再问。

古月娜站在几人稍后的位置。

她看着那座面对森林的石像，银色的发被风吹到肩前。雕像脚下有人走过，身影很小。再往后，坑壁在暮色里渐渐沉下去，黑得比别处更深。

蔡须坤回头看向她。

``娜儿。''

古月娜收回目光。

``我在。''

蔡须坤顺着她方才看的方向望了一眼，又看向坑底那道亮着灯的门。

``进去看看？''

古月娜没有立刻回答。

她望向城门那边。几辆归城的车正慢慢经过，车上的人朝守卫挥手，门洞深处的灯一盏接着一盏亮了。

过了一会儿，她轻轻点头。

``走吧。''

五人转向通往坑底的步道。

石梯贴着坑壁向下，到了第一处转弯，五人不约而同地停住了。

上面的城墙还露在坑沿之外，下面的台阶却已绕出去很远。隔着栏杆望向对面，黑亮的坑壁高高立起，中间空出来的地方，大得让人一时找不到合适的东西来比。

蔡兴坤扶住栏杆，回头望了望天州。

``刚进城的时候，我还觉得，十丈高的城墙已经够震撼了。''

她再转回来，看向石梯下方。

``可这个坑\ldots\ldots 比史莱克整个训练场都要大。''

奥斯卡顺着她的目光看过去。

前些天，他们从训练场入口走到中央时，还感叹过学院将场地扩得如此宽阔。梅天狸站在中央，围观的学生沿结界排开，已经能拉出很长的一圈。

如今把那片场地放进眼前的巨坑里，四周竟还会剩下大片空处。而脚下层层向下的石梯，又在提醒他们，这里还有几十丈的深度。

``在学院，站在一边，还能看清另一边的人。''卢本韦轻声说。

她指了指坑底步道上的几个行人。

``你们看那里。''

几个人影正沿着护栏移动，若不是其中一人抬起手，远看几乎分辨不出动作。

蔡须坤的视线从那些人身上移开，落到斜贯坑壁的一道黑痕上。

她记得自己与同伴合力时，怎样用尽了全身的魂力，也记得梅长老接下那一击后，训练场中央翻起的碎石。

面对眼前的遗迹，她一时没有说话，只将搭在栏杆上的手慢慢收了回来。

``走慢些。''她提醒道，``前面还有转弯。''

古月娜已经侧过身，等着几人重新跟上。

路越往下，石像便越显得高。蔡兴坤几次仰起头，到后来，已经看不全那人的眉眼，只能看见双手与长剑压在更高处的光里。

奥斯卡将行囊往肩上提了提，等她走稳，才继续向前。卢本韦在转弯处让过一位上行的老人，又回身招呼同伴。

蔡须坤走到最后一段石阶时，再次抬起了头。

她想起几天前，大长老抬手示意三颗魂核的样子。

此刻，那张相同的面容，高高立在近一百七十年前留下的战场中央。

纪念馆门内有脚步声传来。

她扶了扶行囊，跟着同伴走到了门前。

馆里比外面安静。

五人刚跨过门槛，脚步声便沿石地传了出去。入口旁摆着几张参观图，一名头发花白的馆员正在收拾长桌，听见声音，抬头向他们点了点头。

``还来得及看。''他说，``沿这边进去，最后从左侧出来。''

蔡须坤道了谢，接过一张图。

正对入口的墙上，是一幅天州旧城图。

她先找到了西城门，再沿旁边标出的街巷往里看。旧城的轮廓比如今小了许多，有几条路的方向还认得出来，更多的房屋与街道，却已经与他们方才经过的景象对不上了。

``那家吃饭的铺子，应该在这里再往东。''蔡兴坤指了指图外。

她手指落下的地方，只有一片留白。

卢本韦站到她旁边，低头看了看角落里的说明。

``后来才扩出去的。''

再往前，是几扇厚玻璃。

其中一扇后面，横放着一截断裂的门闩。木头发黑，断口的纤维一根根翘着，外面箍着的铁已经变了形。

奥斯卡停下来，隔着玻璃看了好一会儿。

大长老讲过，兽潮攻城时，西城门的门闩断过。守军将能抬动的东西搬去堵门，城楼上的炮还在向外打。

如今，断木安静地躺在灯下。

旁边还有一段损坏的炮管，管口裂开，边缘朝外翻着。底下的牌子记录着它从何处收回，又在什么时候被送入馆中。

``这些，都是当年留下来的？''奥斯卡问。

馆员跟了过来。

``门闩和这段炮管是。再前面的阵盘图是后来照着记录复原的，下面写着。''

几人顺着他指的方向看去。

一张城防阵线图悬在墙上，西边几处标着破损的位置。旁边的文字中，蔡须坤很快找到了一个熟悉的名字。

Smile Lee。

那段说明记着他接管、修复城防阵线，并开启护城大阵的经过。

蔡兴坤将上面的字读了一遍，轻轻呼出一口气。

``城外那一下的余波，最后就是撞在这里。''

``嗯。''蔡须坤说。

她在图上看见了那一圈围着城池的阵线，又抬头望向墙的另一端。战场示意图里，城外三里处标出了巨坑所在的位置。

从两张图之间走过去，只需几步。

当年，那几里的距离里，却有人在城头守着，有人在城外迎上了帝天。

馆员见他们看得仔细，便将旁边一盏偏暗的灯拨亮些。

``第一次来？''

五人点头。

``先前听人讲过。''卢本韦说，``真正站到外面，才知道这个坑有多大。''

馆员笑了笑。

``常有人这么说。有人在图上量半天，出了馆，还要沿坑边再走一段。''

他没有催促，又回长桌那边去了。

几人继续向里。

展柜之间有几份军报的抄件，字迹急促，有些地方留着涂改。再往后，是守城人员的记录，以及后来寻访到的口述。

有一段，写的是守军将受伤的同伴从城楼上背下来。

另有一段，只记着一家人怎样挤进城门，后来又去哪里寻找失散的亲人。

蔡须坤逐页看过去，脚步越来越慢。

这些人的名字，她大多没有听过。

她记得的那些强者，在这些记录里，有时只出现了短短一行。更多的字，写着米粮、药物、伤员，写着哪条街腾出屋子，又从哪里找来了还能用的水缸。

蔡兴坤走到她身边，将参观图折小了一些。

``姐。''

蔡须坤转过头。

妹妹指着一页记录上的地名。

``就是刚才那条街。我们买水的时候，还在那里停过。''

蔡须坤看了看，认出了那几个字。

那条街上，如今有人卖布，有人修车，拐角处的花摊前还围着几个孩子。

她伸手按平了参观图的折角，没有马上说话。

古月娜一直站在旁边。

她看完最后一页，又回头望向入口。隔着两道门，外面还有一小片天光，石像的基座挡住了更远的视线。

蔡须坤走过去，轻轻叫了她一声。

``娜儿。''

``嗯？''

``你想先出去透透气吗？''

古月娜摇头。

``我还想看完。''

她的视线重新落回那些记录上。

``以前听见天州这个名字，我最先想到的，都是森林里留下的事。''她压低声音，``走到这里，才知道城里也有这么多人的名字。''

蔡须坤站在她身边，没有催她。

奥斯卡与卢本韦已经走到下一处展柜，见两人停着，也放慢了脚步。蔡兴坤将抱着的黑子鸡往臂弯里拢了拢，安静等在一旁。

过了一会儿，古月娜将最后一页看完，才向几人走去。

``走吧。''她说，``明天，我带你们进森林。''

五人离开馆时，坑沿只剩了一线晚霞。

他们沿石梯回到上面，又去了旧纪念堂。堂里刻满了名字，靠近门边的石壁前，放着几束新鲜的花。

几人没有久留，在门内站了一会儿，便向外退去。

回头时，蔡须坤看见古月娜最后走出门。她的步子很轻，到了台阶下，才重新抬头看向西边。

森林的方向，已经沉在夜色里。

他们当晚回城住下。

奥斯卡将明日要带的食物分好，卢本韦则把参观图压进了旧笔记旁边。蔡兴坤坐在床沿，用布擦去鞋上的灰，黑子鸡伏在她脚边，已经闭上了眼睛。

蔡须坤将写着传灵塔的几页纸重新取出，看了一遍，在空白处添了几行。

各地留下的伤亡记录。

愿意共生者的经历。

发生过什么，怎样解决，后来又怎样相处。

她写得不快。写完后，将纸收回封套，放到了明日取用方便的地方。

第二天清晨，古月娜敲响了房门。

五人吃过早饭，从西门出了城。

这一次，他们没有再沿纪念馆的步道下去。走到岔路口，古月娜停了片刻，朝着那座石像的方向望了一眼，随后转向通往林区的路。

官道渐渐窄了。

田地退到身后，草丛与灌木多起来。越往前走，树影便越密，风吹过叶片的声音，也渐渐盖过了远处车马的动静。

第一片高大的树冠出现在前方时，古月娜抬起了手。

几人随她停下。

``从这里进去。''她说，``跟着我。路上遇见魂兽，先让我来。''

蔡须坤点头，走到她身侧。

``娜儿，我们都在。''

古月娜望了她一眼，迈进了树荫。

奥斯卡走在队伍中间，卢本韦与蔡兴坤跟在后面。脚下的石路很快被湿土取代，草叶擦过衣角，树根在地面拱起，一道接着一道向林中延伸。

走过一片灌木时，古月娜忽然停住。

左前方有细小的响动。

卢本韦的手动了一下，随即又按住。蔡须坤已经看见，两棵树之间，有一双眼睛正望着他们。

灰色的身影伏在暗处，没有立刻扑出来。

古月娜向前半步，开口说了几句。

那声音很轻，节奏与人类的话语不同。灌木后的魂兽抬了抬头，目光从她身上移向后面的四人，喉间发出一声低低的回应。

古月娜又说了一句。

它站了起来。

蔡兴坤下意识地看向姐姐。蔡须坤没有动，只轻轻摇了一下头。

片刻后，那道灰影转身钻入林中。

枝叶晃了几下，很快停住。

``它去传话了。''古月娜说。

``往里面？''奥斯卡低声问。

``嗯。''

古月娜回头看向四人。

``我们先在这里等一会儿。''

蔡须坤将肩上的行囊卸下来，放到脚边。

封套里的纸没有被碰皱。她又检查了一次，将它重新收好，抬头望向那道灰影离去的方向。

这一次，他们带来的话，终于开始向森林深处传去了。

过了一会儿，林中更远的地方，传来一声短促的兽鸣。

古月娜侧耳听了听，随后向四人招手。

``可以往里走了。''

蔡须坤重新背起行囊，跟到她身边。古月娜拨开垂到路上的枝叶，待几人都走过，才松开手，继续向前。

树冠渐渐合拢。

脚下只剩一条隐约可辨的小径，古月娜却没有迟疑。她沿着树根之间的空隙走过去，不时回头看一眼，带着同伴一步步往森林深处行去。

他们的身影，很快隐进了重叠的树影里。

小队已经走进森林深处时，天州城外的巨坑上方，才出现了一道黑影。

纪念馆里已经有了零星的参观者，有人提着清扫的工具走上石阶。坑沿那座旧纪念堂前，几枝新换的花还带着水珠。

黑色的袍角掠过晨光，向石像落去。

来人自高处落下，身形在半空微微一缓，双足便落在了雕像的剑柄上。

宽大的黑袍随之垂落，将身形遮了大半。兜帽压得很低，从坑底望上去，只能看见一道立在石剑旁的轮廓。

那人站稳以后，抬起了头。

test\_alpha0的石像近在眼前。

从这里看去，雕像面上的每一道刻痕都清清楚楚。眉眼间积着一点尘，唇角的线条仍很平静，石刻的目光越过来人，望向远处的森林。

黑袍人原本还在轻轻摆动的衣袖，慢慢停了下来。

他仰着头，望着那张脸，怔住了。

下面有人走过，交谈声沿着基座传上来，很快又散在坑里的风中。

那人没有低头。

一只手从袖中抬起，停在半空，过了许久，又缓缓收回。

阳光沿着石像的肩头移下来，照亮了拄在剑柄上的双手，也照到了黑袍人的脚边。

他仍站在那里，望着雕像出了神。

过了许久，兜帽下才传出一声低语。

``不知不觉，居然这么多年都过去了。''

他的目光仍停在那张石刻的面容上，声音很轻，几乎混进了风里。

随后，他缓缓侧过头，望向星斗大森林所在的方向。

``人类与魂兽，要改变了。''

\chapter{第一百三十八章：重返星斗}\label{ux7b2cux4e00ux767eux4e09ux5341ux516bux7ae0ux91cdux8fd4ux661fux6597}

离开那片灌木以后，古月娜走在了更前面。

林间的小径很快断了，前方只有交错的树根与积着落叶的坡地。她停下来，抬头辨认了一下方向，又望向蔡须坤。

``到这里，就不用再藏着气息了。''

蔡须坤点头。

古月娜身上的银光先亮了起来。

那层光很淡，顺着发梢与肩头漫开，又缓缓收回衣袍附近。一股清晰的龙族气息却留在了林中，随着她向前迈步，沿着树影间的空隙传了出去。

蔡须坤也松开了对本源气息的遮掩。

淡金色的纹路从她手腕浮起，袖口遮住了大半，只在抬手时露出细细的一线。那股属于帝皇瑞兽的气息，随之融入了四周湿润的空气。

附近的枝叶忽然响了一阵。

一只伏在横枝上的魂兽抬起头，鼻翼轻轻动了动。它原本绷起的前爪慢慢放松，低下脑袋，沿着树枝向后退了几步，将几人头顶的位置让开。

更远处，也有脚步停住。

奥斯卡看向右边。

两头灰毛魂兽从灌木后露出身形。它们先望向古月娜，又将目光转到蔡须坤身上，嗅了嗅，随即伏低前身，尾巴也垂了下去。

古月娜没有停步，只向它们轻轻点了点头。

两头魂兽一直等到五人走过，才慢慢直起身。

蔡兴坤回头看了一眼。

``它们认出来了？''

``认得气息。''古月娜说。

她看向蔡须坤露在袖外的金纹。

``尤其是你的。''

风从两人身边穿过，淡金色的纹路在蔡须坤腕上停了片刻，又隐入袖口。

她没有立刻往前走。

从前寄在帝天身边时，她也曾感受过这样的注视。那时隔着一层深沉的龙息，许多魂兽远远停住，待帝天经过，再朝着金光垂下头。

如今，她自己站在林地上，能闻到潮湿的落叶，也能听清那两头魂兽的呼吸。

身后三个同伴的脚步停了下来。

奥斯卡原本正低头整理肩带，见蔡须坤没动，便也顺着她的目光望过去。蔡兴坤抱着黑子鸡，安静地站到了姐姐身边。

``它们看的是你。''古月娜轻声道。

蔡须坤点点头，将手从袖中伸出来。两头魂兽的耳朵动了动，却没有急着靠近。

她终于向它们笑了一下。

其中一头小心地抬起脑袋，鼻尖在风里嗅了嗅，随后退向一旁，将身后那条窄路也露了出来。

``走吧，娜儿。''

古月娜应了一声，重新转身。

蔡须坤跟上时，步子比方才慢了一点。她留意着两旁的枝叶，经过低垂的树枝时，也会伸手托一下，等后面的同伴过去，再轻轻放开。

一路上，这样的情景越来越多。

在溪边饮水的鹿形魂兽闻见气息，便抬起头，将盘曲的角向旁边偏开。林坡上的猿类魂兽不再拨弄树枝，蹲在高处安静看着他们。还有几只羽毛鲜亮的鸟，从枝梢落到矮一些的地方，收拢翅膀，等两人经过。

五人穿过一片高草时，前方本来还有几声低吼。

古月娜与蔡须坤走近以后，吼声很快停了。草丛两边各有一道身影退开，中间露出一条可以过去的空隙。

卢本韦走在后面，将这些变化都看在眼里。

她的手原先还搭在身侧，走过几处灌木以后，才渐渐放松下来。

``上次进星斗，我总盯着这些地方。''她指了指高草根部，``风一吹，就得看看里面藏着什么。''

奥斯卡低头避过一根横枝。

``我记得。你盯一边，我盯另一边。''

卢本韦轻轻应了一声。

从前赶路时，他们总得轮流留意身后。如今，魂兽从树后走出来，闻到前面两人的气息，便安静地让开。几次眼看着要撞上，最后都只剩一阵草叶轻响。

她仍留意着路旁，却终于能腾出目光，看看这些魂兽原本在做什么。

有的正在溪边喝水，有的守着巢穴，还有的带着几只幼兽，在树根下面觅食。

一只绒毛蓬松的小兽，便是在这时候跑出来的。

它从成年魂兽前爪旁钻过，鼻尖贴着地面嗅了几下，随后抬起脑袋，望向蔡须坤。

蔡须坤停住脚步。

小兽又向前挪了两步，尾尖抖了抖。

后面的成年魂兽站起来，走到它身侧，却没有将它叼回去，只将头压低，安静地望着面前几人。

``它想跟我们走？''蔡兴坤小声问。

古月娜看了一眼树根下的巢穴。

``它的家在这里。''

她顿了顿，眼中露出一点浅浅的笑意。

``只是想再靠近一些。''

蔡须坤蹲下来，将手伸到小兽面前。

它嗅过她的指尖，便将额头贴了上来。细软的绒毛擦过掌心，暖乎乎的。蔡须坤顺着摸了一下，它就眯起了眼睛。

蔡兴坤也蹲到一旁。

怀里的黑子鸡探出头，盯着那团绒毛看了片刻，忽然动了动喙。

``你别啄它。''蔡兴坤赶紧低声提醒。

黑子鸡转过脑袋，看了她一眼，又将头收了回去。

卢本韦忍着笑，先替几人让出了路。

蔡须坤摸了摸小兽的头，站起身。它还想往前跟，成年魂兽便走过来，轻轻碰了碰它，将它带回了树根旁。

五人又走出去一段，回头时，还能看见那只小兽从前爪之间探出的脑袋。

到了傍晚，古月娜选了一块临近溪流的空地。

空地背靠石坡，离水边有一段距离，脚下也干燥。几人放下行囊，奥斯卡绕着附近走了一圈，回来时，手里抱着些干枝。

``这里可以生火？''

古月娜看了看石地，又看向周围的树。

``用这边的石头围住，别让火星飘出去。''

卢本韦已经俯身搬起一块，将它放到干枝旁边。

蔡兴坤去取水，蔡须坤则与古月娜沿着营地外侧走了一小段。她们没有收起气息，附近的魂兽嗅过以后，便各自退到远一些的位置。

有几道身影停在林间，安静地伏了下来。

``它们不回去吗？''蔡须坤问。

``过一会儿就散。''古月娜说，``有些只是想在附近多待一阵。''

两人往回走时，一头鹿形魂兽正从另一侧经过。它朝她们低了低头，随后绕开营地，走向下游。

蔡须坤看着它消失在树后，脚步慢下来。

``娜儿。''

``嗯？''

``我以前知道这股气息意味着什么，可在人类的城里，很少有这样的感觉。''

她看向四周。

夜色正一点点落进森林。很近的地方有虫鸣，溪水在石缝里流，刚才跑到她身边的幼兽，此刻应当也已经回到巢中。

古月娜沿着她的目光望过去。

``你回来，它们自然会这样。''

蔡须坤安静了一会儿，才向火光的方向走去。

锅里的水已经开了。

奥斯卡分好食物，卢本韦拿树枝拨了拨火，将一根没有烧透的干柴推到里面。蔡兴坤坐在石头上，把自己的位置往旁边挪了一些，招呼姐姐过来。

五人围着火吃过东西，又将火堆收小。

林间偶尔有脚步经过，到了营地附近，便会停一下，再慢慢绕开。起初奥斯卡还会抬头看，听得多了，渐渐也能辨出那些声音是在往远处去。

卢本韦抱着膝盖，望向树后的暗处。

``我现在才发现，夜里原来有这么多动静。''

``你上回没听见？''蔡兴坤问。

``听见了。那时候总想着，会不会有东西扑过来。''

她说着，自己也笑了一下。

奥斯卡抬手指向溪流另一边。

``刚才那边是喝水的。来时走得重，回去就轻了。''

古月娜听了一会儿，点点头。

``嗯。它还带着一只小的。''

几人循声望去，只见树影晃动，什么也看不清。

夜深以后，几人轮流休息，火堆旁渐渐只剩干柴轻轻裂开的声音。

蔡须坤靠着石壁，忽然听见共生联系里传来狼王的声音。

``这里的水声，和我族里那道溪很像。''

她睁开眼，望向黑暗中的水面。

``想它们了？''

``有些想。以前到了这个时候，往往还有几只幼兽不肯回去。大一点的便绕着溪岸找，找到以后，故意叫得很重，吓它们一下。''

蔡须坤听着，嘴角慢慢动了动。

她原先总觉得，狼群在夜里发出的声音，是召集、警戒，或是发现了敌人。如今听狼王这样说，才知道那些低鸣里，也有寻回幼兽的催促。

``等到了湖边，你也将这些讲给它们听。''她轻声道。

狼王安静了一会儿。

``我想让族人亲眼看看我。''

``会有机会的。''

蔡须坤将肩上的外衣拢好，重新闭上眼睛。溪水还在向下游流，几道原本让她有些陌生的叫声，也渐渐能够分辨了。

第二天，赶路比前一天顺利了许多。

树木越来越高，许多树根从土里隆起，连成一道道低矮的坎。古月娜带着他们从根系间穿过，遇见湿软的洼地，便提前换向另一侧的坡。

途中经过一片长满藤蔓的地方，蔡须坤正想拨开挡路的枝条，上方忽然响起一阵沙沙声。

一头伏在高处的魂兽探出前肢，将垂下来的几根藤蔓拨向两旁。

古月娜抬头，向它说了一句。

那头魂兽低下脑袋，等五人都走过去，才轻轻松开。

奥斯卡回头望了一眼，忍不住道：``还替我们开路。''

``它听得懂谢谢吗？''蔡兴坤问。

古月娜望向她。

``你可以让它看见。''

蔡兴坤便停了一下，朝树上的魂兽挥了挥手。

对方望着她，过了片刻，也动了动垂在枝外的前爪。

她笑起来，转身追上了姐姐。

到了午后，一头巨熊出现在前方坡地上。

它的肩背高出两旁的灌木，坐着时便已占住了半片树影。卢本韦先看见它，随即停步，将身后的奥斯卡也拦了一下。

巨熊抬起头，鼻端喷出一口气。

银色与淡金色的气息顺着风送了过去。

它的动作顿住了。

那双原本直视过来的眼睛，缓缓向下。巨熊撑起身子，挪开挡路的前爪，又退到石坡旁，将几人面前的位置让了出来。

古月娜带着蔡须坤走过。

经过巨熊身侧时，蔡须坤微微放慢脚步，向它点了一下头。

巨熊俯低了脑袋。

身后三人跟上来，直到越过那段坡地，卢本韦才轻轻吐出一口气。

``它也认得。''

古月娜回头看了一眼。

``嗯。''

巨熊仍站在原地，待他们走远，才重新坐到石坡边。

走过巨熊占着的坡地以后，卢本韦又回头看了看。

石坡旁的树干上，留着几道很深的爪痕。有的已经被新长的树皮挤到两侧，有的还露着浅色木屑。沿着它们往外，能看见一条被踩实的小道，到了溪边便断了。

溪流另一侧，也有几道身影停在林间。

它们没有向这边走。巨熊坐回石旁以后，其中一头才缓缓转身，沿着自己的那片林地离开。

``这里也分地界？''卢本韦问。

古月娜看了一眼两岸。

``当然。水从哪里喝，幼兽在哪里长大，冬天还能找到多少食物，都和住的地方有关。''

``若是争起来呢？''

``有的族群自己谈，有的要找能让双方都停手的。谈不拢，就会打。''

蔡兴坤望了望身后的石坡。

``所以，族长也得管这些。''

古月娜点头，没有将话说得很轻松。

她从一处低矮的树根旁走过，指了指前方另一条绕开的路。

``我们走这边。那一片是它们休息的地方。''

几人跟着转向。

蔡须坤跨过溪岸上的石头时，又看了一眼两侧的林木。方才让路的魂兽，对她们有敬意，可各自要守住的地方仍在那里。

她将这一处方向记了下来，继续跟上古月娜。

再往里，魂兽身上的气息愈发厚重。

有些隔着一片树林，就能感觉到它们行走时带起的细微震动。可一旦嗅见两人经过，它们便会停下来，将身形从路旁移开，安静看着五人走过。

蔡须坤渐渐习惯了那些目光。

她也会在经过时放慢一点，留意对方的动作。有时只需轻轻颔首，便能看见一颗硕大的头颅低下来；有时则是一声短短的鸣叫，隔着树冠传向远处。

傍晚停下休息时，她问古月娜：``消息已经传到里面了吗？''

``应当传到了。''

古月娜抬起头，望着更深处的林影。

``我们的气息一直在。沿途又见到了这么多魂兽，里面很快就会知道。''

卢本韦将水袋递过来。

``那到了以后，就把带来的东西好好说一遍。''

蔡须坤接住，点了点头。

第三天，他们穿过了一片雨中的密林。

雨水沿宽大的叶片落下，打在肩头，脚下的泥土也渐渐湿滑。五人将行囊裹好，踩着树根和露出地面的石块前行。

遇上一道被雨水冲开的窄沟，古月娜先掠到了对岸。蔡须坤随即提气跃过，脚尖在露出的树根上轻轻一点，便稳稳落在她身侧。

蔡兴坤跟着过去，顺手托起一根垂下的湿枝，等后面两人落地，才将枝叶放回去。

奥斯卡抹去脸上的雨水，将恢复香肠分给同伴。

``补一点，路还长。''

卢本韦接过，先将一根递给前面的古月娜，又俯身将压住鞋面的背带收紧。几人略作停留，便继续往前。

雨停以后，林间腾起一层薄雾。

远处有魂兽从树后抬头，闻到气息，便向旁边退开。还有一头停在溪水中央，等五人来到岸边时，主动移到了下游。

他们在原地稍作休息，随后又赶了半日的路。

当晚落脚的地方，已经听不见林外那些熟悉的鸟鸣。夜间传来的声音更低，也更悠长。古月娜坐在石边听了一会儿，神情却很平静。

``明天再走一段，就到了。''她说。

从入林算起，已经赶了三天的路。第四天清晨，五人重新沿溪流向前。

晨雾还没有散，树冠间只有几束细光。脚下的水越来越清，贴着溪岸延伸的石头上，留下了几道很深的爪痕。

古月娜走到一株粗大的古树旁，停下脚步。

蔡须坤随她望向前方。

那里的林木比身后更密，树干间隔很远，却依然有枝叶在高处相接。越过低垂的藤叶，能看见一段向里延伸的缓坡，坡后还藏着更深的树影。

``我们已经接近核心区了。''古月娜说。

卢本韦收起水袋，奥斯卡也将肩上的行囊重新提好。蔡兴坤抱着黑子鸡，走到姐姐身边。

蔡须坤看向古月娜。

``娜儿，还是你带路。''

古月娜点头。

银色与淡金色的气息仍在林中缓缓散开。前方一处树影下，伏着的魂兽嗅了嗅，随即低下头，挪开了挡在石路边的身躯。

古月娜先走了过去。

四人紧跟着她，踏上了那段通向核心区的缓坡。

缓坡越往上，脚下的岩石便越多。

古月娜沿着山势向前，越过一片横生的树根，又从两块巨石之间穿了过去。蔡须坤紧跟在她身后，走到坡顶时，先听见了一阵比溪流更开阔的水声。

风从山的另一侧吹来。

前方的枝叶渐渐稀疏，几人翻过最后一段山脊，眼前忽然显出一片明亮的水色。

蔡兴坤的脚步停了。

山下是一座湖。

湖面铺在层层林木之间，晨雾还浮在水上，远岸的树影被遮住一半。阳光从云隙中落下来，将近处的水照得清透，细小的波纹一直荡到岸边，又在石头旁缓缓散开。

几道溪流从山间汇来，最后都没入了这片湖水。

奥斯卡走到两人身边，抬眼望去。

``生命之湖。''

这个名字从他口中说出来时，声音不自觉地放轻了些。

古月娜点头。

``到了。''

蔡须坤站在坡顶，望着那片熟悉的水。

上一次来到这里，他们带着伤，也带着牢大那缕脆弱的残魂。那时她只想着，能不能再找到一点办法，顾不上把四周看得很清楚。

如今重新站在高处，她才看见湖岸延伸得很远，林木之间还有一片片开阔的草地，水边的浅滩上留着许多不同的足迹。

其中一些，大得足以让几人并排站进去。

蔡兴坤顺着一串脚印望向旁边。

一头生着厚甲的魂兽伏在树下，背上的苔色与石头相近。方才站得远，她竟以为那只是一块隆起的岩石。直到它抬起头，一股深厚的气息才从树影里显了出来。

蔡兴坤抱着黑子鸡的手微微一紧。

``姐，那边树下那个\ldots\ldots''她压低声音，看向姐姐。

古月娜已经听见，回过头来。

``十万年。''

她说完，那头魂兽也嗅见了风中的气息。

它先望向古月娜，再看向蔡须坤，随即将头低了下去，撑起身体，慢慢挪到了树的另一边。

山坡下面，一片树冠也动了。

一只粗壮的前臂拨开枝叶，露出一张宽阔的兽面。它看向这边，目光停了片刻，同样垂下脑袋，没有挡住通往湖边的路。

卢本韦望了一眼，又看向更远处的林间。

那里还有。

先前几日，万年魂兽越来越多，他们已经有了准备。可真正翻过山，到了生命之湖附近，十万年魂兽出现的次数，仍让几人心中发紧。

有的在岸边休息，有的守在溪流汇入的地方，还有的只露出一截盘在古树后的身体。即使它们并没有释放威压，那种长久修炼积下的气息，也很难让人忽略。

``这里平日就有这么多？''卢本韦轻声问。

``湖的周围，是许多强大魂兽修行的地方。''古月娜答道，``再往外，各族也有自己的栖息地。''

她看向山下，原本平静的神情忽然变了一点。

蔡须坤顺着她的目光望过去。

湖边一块临水的石地上，站着一道高大的身影。

长发垂在身后，那人侧对着山坡，目光原本落在湖面上。似乎察觉到了他们的到来，他缓缓转过头，一双金色的眼瞳隔着尚未散去的薄雾，望了过来。

帝天。

蔡须坤一眼便认了出来。

可她的目光很快又停在了旁边。

碧姬站在离水不远的地方，一身翡翠色的衣袍，袖边映着细细的光。更后面，一道赤色的庞大兽影正伏在树下，身上火焰收得很低，只有呼吸间，才从鬃毛里透出几缕。

赤王。

湖侧另有一株古树，树冠极密。粗大的根须贴着地面延伸，到了水边才向下没入泥土，深沉的气息沿着整片树影缓缓流动。

那是万妖王所在的位置。

再往两侧看，还有几道气息各不相同的身影。

古月娜在坡顶停了片刻，才低声道：``居然都在。''

奥斯卡望向她。

``你说的是\ldots\ldots''

``如今留在森林里的这些凶兽。''

古月娜将视线从湖岸一一移过。

``平日里各有各的地方，很少会这样聚在一起。''

卢本韦没有再问。

她先前能辨出十万年魂兽与沿途那些魂兽的不同，如今再看湖边，这一圈存在的气息又更深。她望过去时，常常才觉出其中一层，后面便还有更多，让人无法轻易探到尽头。

前些天，他们五人合力接过梅长老的一招，心里还留着那份兴奋。此刻站在湖岸上方，她却不由自主地将肩背挺直了一些。

奥斯卡也收回了目光，将脚边一根枯枝轻轻拨开。

山风将雾吹开一些，临水那块石地显得更清楚了。

上一次来到这里，梅天狸与ddfg一前一后护着他们。两位长老停，他们便停；龙息中让出一条路，才有人回身示意他们跟上。

如今身边只有五个人。

奥斯卡将行囊从肩上换到手里，免得垂下的带子扫过石块。蔡兴坤伸手理了理被树枝勾乱的衣领，又替姐姐拂去肩头一片干叶。

蔡须坤转过头，见妹妹正看着自己。

``紧张？''

``有一点。''蔡兴坤老实道，``都在看这里。''

她说完，自己将呼吸放缓了些，抱着黑子鸡站直。

卢本韦低头检查了一下封套放的位置，随后向蔡须坤点了点头。

一路准备过的话，此刻都还记得。可真要走到这些强大的存在面前，亲口讲清楚，她才知道，先前在火边练习时少了多少分量。

``我们一起说。''奥斯卡低声道，``谁漏了，旁边的人补上。''

古月娜听见，回身望向几人。

``嗯。''

她稍稍等了一下，直到蔡须坤走到身边，才与她一同迈下山坡。身后三人也随即跟上，五道脚步声在石间接续着落下。

五人沿山石间的空处下行。靠近道路的魂兽闻到两人的气息，便相继移开，让出一段段通向湖岸的位置。

走到最后几棵树前，古月娜停了一下。

``族长。''

帝天已经转过身。

他的视线先落在她身上，又移向蔡须坤。两人都没有遮掩自己的气息，银色与淡金色在林间相接，清清楚楚地停在他面前。

``回来了。''

他的声音不高，湖边的几道身影却都望了过来。

``回来了。''古月娜答道。

蔡须坤也点了点头。

帝天看了她们一会儿，随后望向跟在后面的三人。

卢本韦、奥斯卡与蔡兴坤在两人身后站定，分别向湖边致意。

``这一次，没有长老同行？''帝天问。

``只有我们五个。''蔡须坤说，``其他同伴还在蓬莱岛，我们先回大陆办些事。''

帝天的目光从几人身上一一掠过。

``长进不少。''

这句话让蔡兴坤稍稍松了口气，抱在怀里的黑子鸡也探出头来。它看了一眼湖边，又安静地伏回她臂间。

碧姬向前走了几步。

她看着蔡须坤，语气比旁人柔和些。

``一路都没有收起气息，早就知道你们往这里来了。''

蔡须坤望向湖边聚集的身影。

``所以，诸位才会在这里？''

``你的气息回到森林，总会有人想来看看。''碧姬说。

赤王抬起了头，身上那几缕火焰轻轻晃动。

``这趟在人那边，待得够久。''

它看了看蔡须坤，又看向古月娜。语气里有询问，也有掩不住的在意。

古月娜走到蔡须坤旁边。

``去了蓬莱，也回过大陆。''她说，``遇见了许多事，有些想当面说清楚。''

帝天没有立刻接话。

他看着蔡须坤的手腕，那里有一线金纹，袖口稍稍移动，光便随之隐入衣下。

蔡须坤顺着他的目光低头，随后将手轻轻抬起。

一缕魂息从她身侧显现。

湖边的目光随之集中起来。

灭世狼王的魂影缓缓凝出，四爪落在石地上。它先看了一眼湖水，又望向帝天，随后向四周熟悉的气息低首致意。

``我也来了。''

它的声音清晰，轮廓没有像上次那样，在显形之后立刻晃动。

碧姬停住脚步，仔细看了看它。

帝天的目光也落在那道魂影上。

上一次，它带着重伤初成魂灵，连维持形体都很吃力。如今再次来到湖边，它已经能够自己站在那里，向旧识开口。

``看来，这段日子过得还好。''碧姬轻声说。

狼王点头。

``有些事情，正想同诸位说一说。''

蔡须坤将行囊卸下来，取出那个一路带来的封套。

纸张放在临近的一块平石上，卢本韦走近半步，替她压住被湖风掀起的一角。奥斯卡将其余几页依次展开，古月娜则站在她身侧，看向湖边众人。

最上面的那一页，写着三个字。

传灵塔。

帝天的目光落了下来。

``这是你们这次带回来的？''

``嗯。''

蔡须坤抬起头。

``我们想建立传灵塔，让愿意共生的魂兽和魂师，有地方彼此了解，也有办法把已经做成的经历留下来。''

湖边没有谁立刻出声。

她向身侧的狼王望了一眼，接着说道：

``在讲那些设想以前，我想先请诸位听一听，它这些日子过得怎么样。''

狼王走到了石地前方。

蔡须坤将手放下，没有跟过去。那道魂影离开她身侧，沿着湖岸缓缓走出几步，停在了一处平静的浅水前。

水中映着树影，也映着它略显透明的轮廓。

``刚开始的时候，我连站稳都难。''它说，``魂力一散，就只能回去休养。听得见外面的声音，想自己出来看一眼，却撑不了多久。''

碧姬走近了些。

``现在呢？''

狼王抬起一只前爪，又重新落下。湖边一片干叶被魂力拨动，沿着石面轻轻滑了一段。

``现在好多了。她修炼的时候，我也在梳理自己的魂息。有时能帮她稳住一点，有时是她帮我。到了蓬莱岛，跟着她一遍遍走山阶，这道联系也慢慢稳了下来。''

``起初，我很不习惯人类的日子。''

它说着，低头望了望自己的前爪。

``我从前休息，是找一块熟悉的石头，听族群的动静。成了魂灵以后，有时明明想醒着，魂力却不够；有时恢复过来，她又已经睡下了。''

蔡须坤轻轻点头。

那段日子，她自己也常常睡不安稳。偶尔半夜醒来，先辨认心口那缕残魂，再察看狼王的魂息，等两处都安静，才肯重新闭上眼。

``后来我学着留一点力量。''狼王说，``够用来听她身边的动静，也能在她需要时回应。她也渐渐知道，我什么时候疲倦，什么时候还想出来走一走。''

碧姬听得很认真。

``你们会商量？''

``会。有时她觉得我该休息，我却还想把一处没有听明白的东西问完。''

蔡须坤忍不住道：``有一回，都已经下山了，你还在问那一步。''

``你们第二天还要重新走。''

狼王抬起头，语气很认真。

``我若听懂，至少可以提醒一句。''

奥斯卡的嘴角慢慢扬起来，又在赤王望来时收了些。他记得那几天，蔡须坤时常走着走着便停下来，像在听什么；问过以后，才知道狼王仍惦记着白天的推演。

湖边的这道魂影，终于让那些只有他们知道的小事，也有了能被旁人听见的声音。

蔡须坤看着它，眼里也有了些笑意。

赤王却没有跟着笑。

``你现在说话，出来，收回魂力，都是自己在做？''

``是。''

狼王的身影缓缓淡下去，片刻之后，又在近处重新凝出。

蔡须坤始终站在原地，连手也没有抬。

碧姬伸出手，停在魂影前。

狼王向前走了一步，让一缕自己的魂息落在她指尖附近。翡翠色的微光轻轻流过，碧姬看了一会儿，收回了手。

``比上次凝实了许多。''

狼王点头。

``我的记忆还在。我的族群，过去住过的地方，也都记得。她不会替我回答这些。''

万妖王的根须在石地边缘动了一下。

``可你的身体没了。''

``没了。''

狼王答得很平静。

``这一点，没有变。我也不会说，换了这样的形态，就再没有遗憾。如今想回到族群身边，还得重新想，该怎样和它们一起生活。''

它望向林中，停了片刻。

``当年伤我的，是梅天狸。她后来折返救治，那一击也仍旧是她出的。这些，上一次来时，我已经向帝天说过。''

赤王看向帝天。帝天没有否认，只将目光留在狼王身上，等它说下去。

``但当时的伤，已经救不回原来的身体。我选了这条路，至少今日还能站在这里，同诸位说话。''

湖水轻轻拍着岸边的石头。

赤王盯着它看了许久，身上的火苗渐渐收低。

帝天终于开口。

``你想见族人？''

``想。''狼王抬起头，``也想让它们亲自看看，我如今是什么样子。''

帝天点了一下头。

``会传信给它们。''

狼王低首致意，没有再退回蔡须坤身后，而是在石地旁留了下来。

赤王将伏在地上的前爪慢慢收紧。

``它还能说话，我已经看见了。''

它望向蔡须坤手边那几页纸，原本低伏的火光又亮了一些。

``可若你们想凭这一个例子，就把更多族人送到人类建的塔里，我不赞成。''

蔡兴坤的笑意渐渐收住。

湖边安静下来。灭世狼王却没有退开，它抬起头，迎上赤王的目光，仍站在那块平石旁边。

\chapter{第一百三十九章：湖畔的商议}\label{ux7b2cux4e00ux767eux4e09ux5341ux4e5dux7ae0ux6e56ux7554ux7684ux5546ux8bae}

赤王的话没有被湖风吹散。

蔡须坤按住封套，抬头看着它。她方才看见狼王能够自己站稳，心里还松了一些，如今听见这句反对，又将那几页纸重新放平。

赤王撑起了身子。

火焰顺着肩背浮动，映得近旁几块石头微微发红。它并没有将力量推向五人，可这样一头庞然大物站在眼前，仍让奥斯卡下意识地停住了整理肩带的手。

``我愿意听你们说，是因为认得你们。''赤王说，``可一座塔立起来，往后来的人，我认得几个？''

它看向蔡须坤，又望向古月娜。

``你们两个会回森林，记得这里住着谁。那些只听说魂灵能让自己变强的人，也会这样想吗？''

古月娜站在原地，没有避开它的目光。

``所以才要有人一直在两边来往。''

``我听过许多这样的好话。''

赤王的声音沉了下去。

``有人进来要一块兽骨，说带走以后能保护更多的人。有人拿着药来，说能救一头受伤的魂兽，最后救下的东西却装进了他们的魂导器。''

``今日你们说，意识还在。明日换一个人来，他要力量，不愿听魂兽说话，又会怎样？''

蔡兴坤将抱着黑子鸡的手收紧了些，嘴唇动了动。蔡须坤轻轻碰了一下妹妹的手背，示意她先听完。

帝天望向赤王。

``有顾虑，就在这里说。今日也没有谁要你只凭一句话答应。''

赤王停了片刻，转头看向狼王。

``我信它还是它。可我不信每一个后来的人，都肯给自己的魂灵留着这份自由。''

灭世狼王缓缓抬起头。

``那你也该听我说完。''

赤王望着它，火焰稍稍低了一点。

狼王向平石旁又挪了半步，停在蔡须坤另一侧。五人没有挪动它的位置，卢本韦反而将压着纸角的手收近，为它让出了一块空处。

``我站在这里，是想同你们一起商量。''狼王说，``你可以不赞成她们的想法，但别只把我当作需要被解释的一道影子。''

湖边安静了一瞬。

碧姬看了看它，又望向赤王。

``它说得对。''

赤王没有反驳，重新伏低了前身。它仍没有答应，只将目光留在了狼王与那几页纸之间。

万妖王的声音再次响起。

``它的路走通了，你们便想让更多魂兽走？''

蔡须坤望向那片树影。

``我们想让它成为一种能被选择的办法。''

``健康的魂兽，也要舍掉身体？''

这次发问的是赤王。

``你们说共生，最后若只剩一个魂影，还要跟在人类身边，凭什么叫它们愿意？''

卢本韦下意识向前半步，蔡须坤已经先开了口。

``所以，我们不能只把成功的地方讲出来。''

她指向石地旁的狼王。

``它失去了什么，后来又怎样活下来，都得让想尝试的魂兽知道。能够继续留在森林里的，也应当有自己的选择。''

狼王看着赤王，接过了话。

``我也不会劝我的族群，都变得和我一样。''

``我们家也有人死在这条路上。''

卢本韦忽然开口。

湖边几道目光移向她。她没有立刻往下说，先从行囊里取出了那本旧笔记，放在新写的纸旁边。

纸边已经发黄，几处折痕被不同的人补过。赤王低头看了一眼，没有催促。

``卢守衡，是我曾祖的祖父。他那时八十七级，与一头八万年的柔骨兔尝试相合。那头兔子是自愿的，在我们家后山住了四十多年。''

``头四天，他们能够听见彼此的想法，都以为成功了。第五天失衡，人和魂兽都没能留下。''

碧姬的神色凝重起来。

``两边都死了？''

``是。''

卢本韦将书页翻到夹着纸签的地方。

``他父亲卢岳曾经去与凶兽谈过合作，把那头兔子带到人类身边。后来儿子死了，兔子也死了，他没熬过几个月，就带着旧疾与悲痛一起走了。''

她的手指停在书页边，没有将这场失败说得轻巧。

``我们来之前，又听大长老讲过这件事。那之后，家里一代一代继续研究，补的就是当年没有看清的地方。''

赤王望向那本旧笔记。

它先前想到的，大多是魂兽失去身体，被人类带走。可眼前这个年轻人翻开的纸页里，失败留下的空处，也同样落在了人类一家又一家的生活中。

``所以，你们现在知道怎样保证了？''它问。

``知道一些过去不知道的事，也做成过不同的共生。''卢本韦回答，``可不能因为前面有人成功，就不看后来这一对能不能相合。''

她将笔记中一处改过多次的位置指给碧姬看。

``双方能承受多少，本源怎样接上，出了问题怎样停，都得看清。将来真有人来问，不能只给他看最顺利的那一页。''

狼王低头看着旧纸。

``我最初也不知道，自己能撑到哪一步。''它说，``蔡须坤受着伤，还护着另一缕残魂。她会累，我也会。那时候，我们得互相回应，才能知道对方究竟还剩多少力量。''

它抬眼望向赤王。

``所以，我愿意讲。我经历过的地方，至少可以让后来者提前知道。至于走不走，要由他们自己决定。''

赤王没有立即回话。

它看着狼王停在纸旁的前爪，忽然发现，方才一直问问题的时候，自己竟很少想过，狼王也能在这座尚未建起的塔里做些什么。

赤王的目光从它身上移开，落回那几页纸。

蔡兴坤抱着黑子鸡，认真看了看湖边几道身影。

``我们也想把人类那边的事带来。''她说，``有人愿意与魂兽相处，可连从哪里开始说话都不知道。走进森林怕被当作来猎杀的，魂兽出了森林，又怕先被人围起来。''

``有个能先见面、先把话说清楚的地方，总会好些。''

万妖王沉默了一会儿。

``地方设在哪里？''

卢本韦抬起头。

``我们想过，在森林与人类城镇之间找一处。双方都能来，也能有人照看。可具体设在哪里，我们不能替这里的族群定下来。''

``建在人类的地界，最后掌事的还是人。''

``那就让魂兽也掌事。''卢本韦说。

湖边几道目光一同转向她。

她站稳了，继续道：``来谈的、留下修炼的、受伤需要照料的，事情不会只有一件。你们熟悉魂兽，我们熟悉人类，各自都得有人在。''

万妖王的一条根须从平石边缘伸过，停在纸上画出的几间屋舍旁。

``你们想的地方，都是给人住的。''

卢本韦俯下身，看向它指着的位置。

``这些只是最初随手画的。哪里不合适，可以改。''

``水里的族群，去了那里住哪里？栖在树上的，要落在哪一棵树上？人类习惯把门关起来，未必知道，别的族群走进去以后，还得能找到自己的路。''

它说着，将那道根须往旁边移了一点。

``这里不能全铺成石地。''

蔡须坤听完，在空白处记下了水岸与林带。

``我们一路走进来，也总在跟着娜儿辨路。到了人类的地方，许多魂兽恐怕也是一样。''

奥斯卡拿起另一页，重新摆到万妖王面前。

``那靠近溪水的地方先留出来。屋舍也不用把四面围死，来的族群不同，再问它们需要什么。''

他说得很具体，先前还有些僵硬的气氛，也随着这几处能够落在纸上的改动，稍稍松开。

碧姬伸手指了指两边之间的一块空地。

``受伤的也不能都挤在一起。有些需要水，有些怕湿。若要在那里照看，旁边得有人能分辨它们的习性。''

``人类那边也得有人负责。''卢本韦说，``建屋、食物与常用的东西，我家可以出一份。回去以后，我会同父亲将能拿出的部分定下来。''

赤王听到这里，终于又开口。

``你们出得多，往后谁说话？''

``事情是谁的，就得让谁开口。''卢本韦答道，``我们若连眼前几位的意思都不肯听，今日也不用站在这里。''

狼王向前挪了一步。

``我可以留下来帮忙。''

蔡须坤望向它。

``等那地方真正能接待来者，我可以同你一起过去。我能与人类说话，也知道许多魂兽害怕什么。到了陌生的屋舍前，若先见到一个能够听懂它们的同族，总会好些。''

``你愿意一直听它们问？''碧姬问。

狼王看了一眼自己身侧的年轻人。

``我问过她们那么多次，也该轮到我回答。''

蔡兴坤低头笑了一下，随后认真地将它方才的话写在了纸边。

赤王看见了那行新添的字。它原以为，狼王今日只是来证明人类的办法确实有效；到这时候，才听明白它是真的在为那座塔想事情。

奥斯卡从旁开口。

``哪怕开始时地方不大，只能接待少数人和魂兽，也可以先做起来。遇见不懂的，就当面问。我们这一路过来，古月娜替我们说了很多话。可总不能往后每一回，都只等她一个人来。''

古月娜望了他一眼，随后看向帝天。

``族长，我在天州进了纪念馆。''

帝天的目光微微一凝。

赤王伏在地上的前爪也停了。

古月娜没有避开这个名字。

``那里记着守城的人怎样搬药，怎样背走伤员，也记着许多人再没有回来。那些事，我们在森林里很少听见。''

她停了片刻。

``森林里发生过的事，城里的人也未必听见过。我想，若能有人一直在两边往来，至少下一次出了事情，能够先找到该问的人。''

``那里面，记着我们的名字吗？''赤王忽然问。

古月娜的手指停在纸边。

``记着兽潮，也有当年战场的图。''她回答，``可森林里的许多经历，没有写在那里面。''

赤王低低地应了一声，望向湖面。

``我还记得那一拳。''

它说得不响，蔡须坤却抬起了眼睛。

``万妖王的藤蔓碎了，我已经扑到他面前。他一拳打过来，我倒飞出去，撞进丘陵里。等我重新抬起头，前面的声音都变了。''

它抬了一下前爪，又慢慢落下。

``很多时候，想起天州，我先想到的就是自己从土里挣出来，看见那些已经不在的地方。''

``后来回到森林，有些族群的守地空着，许多熟悉的气息再也没回来。过了一年又一年，也还是空着。''

古月娜听完，没有劝它忘掉。

``可我这次走进那座城，看见了当年幸存者的后代。''她说，``他们在那里开铺子，带孩子买花，知道大长老救过自己的家。''

``我在馆里看见一截断掉的门闩，后来才明白，那一头也有人一直顶着门，等身后的亲人先走。''

赤王转头望向她。

她仍站在蔡须坤旁边，却没有像从前那样，只将自己听过的那一段反复说下去。

``森林里失去的，我记得。''古月娜说，``城里失去的，也是真的。若以后还只是从各自活下来的人嘴里听见仇，走出去时，眼前就又只剩敌人。''

帝天缓缓开口。

``那一战，我的判断，由我来担。''

几道目光随即落向他。

他的声音依旧沉稳，没有在这些年轻人面前回避自己当年的决定。

``死去的族人不能再回来，让我们重选一次。如今既然已经知道，猎杀与掠夺从何处来，目光就该先落到那些真正动手的人身上。''

赤王沉默了。

它仍记得那股将自己击飞的力量，胸前的旧痛却早已散去。真正留到今日的，是每逢想到人类，便会一同浮起的那些空地与失去的身影。

它看向狼王，又看向古月娜。

前者已经失去肉身，却仍在为族群想往后的日子；后者离开了森林，也终于带回了这里从前没有认真听过的话。

赤王低下头，望着纸面上还没有定下位置的那座塔，没有再说``人都一样''。

卢本韦听着，眼前浮起了纪念堂里那些刻在石壁上的名字。她将一角被风掀起的纸按好，没有再插话。

帝天的目光从古月娜移回纸面，又望向蔡须坤。

``你们去过总殿？''

``去过。''蔡须坤说。

``test\_alpha0也看过这些？''

``看过。我们把想法告诉了他，也听他说了许多以前的事。''

帝天的手指缓缓收拢，随后又松开。

蔡须坤没有借着大长老的名字继续往下压，只望着他，将话说完。

``他知道魂灵曾经怎样失败过，我们也带着那些记录。但今天到了这里，还想亲自听你们说。塔若要让两边都能走进去，就该让两边都知道，它将来要做什么。''

``知道，和相信，差得很远。''帝天说。

``我明白。''

蔡须坤垂眼看了看脚边的水。

``所以我们回来了。狼王也回来了。今天没能说清的，还可以继续说。''

碧姬弯身拿起其中一页。

纸上除了魂灵共生的设想，还有沿途添下的几行小字，记着受伤的生灵、可以留下的经历，以及怎样将消息传到另一边。

她翻到后面，目光忽然停住。

``宏舒集团？''

帝天抬起眼睛。

蔡须坤将那一页移近。

``我们回到大陆的时候，在海潮市附近遇见过他们的人。''

她把村庄里空无一人、井下藏着女人的经过讲了一遍。讲到那些镜子时，奥斯卡补上了力量被反射的情形，卢本韦则说了将其生擒、交给海潮市中殿与市长的经过。

赤王的火焰又亮了一些。

``也在掠夺魂兽？''

``那个人用的是冤魂之力。村子里的人去了哪里，还要等当地继续追查。''蔡须坤说，``但她供出的另一件事，和帅气教有关。''

``宏舒集团的最高统治者叫王雨新，是胡赞的妻子。''

湖边安静下来。

帝天的神色沉了下去。

蔡兴坤将另一张记着方位的纸递来，放到几人面前。

``据她所说，大陆西岸再往西约三千里，有一片星辰大陆，面积大约是斗罗大陆的十分之一。帅气教过去千年一直在那里发展，宏舒集团也在那里。''

``你们亲自去过那片大陆？''赤王问。

蔡须坤摇头。

``还没有。海上的方位与集团的关系，都是那个人被生擒后说的。''

她将村庄所在的位置指了出来，又把几人从码头一路走过去的方向讲清。奥斯卡在旁补上，村里当时还有温热的食物，许多人像是刚离开，连原本正在做的事都没收拾完。

``就是这样，整座村子空了。''

这一次，连赤王也没有立即接话。

蔡兴坤低头看着那张纸，想起倒在路旁的菜篮，以及被留在院子里的狗。刚看见天州城里那些忙碌的人时，她还想到过那个村子。

那里原本也该有人买东西、做饭，带着孩子回家。

``那女人的力量与帅气教不一样。''卢本韦说，``她用镜子将幻境铺开，也能反射一部分攻击。可她与宏舒集团、帅气教的来往，都是自己供出来的。''

``我们将她交给海潮市的中殿殿主与市长。他们说，城里也有宏舒集团作乱，正在巡查周边。''

狼王抬起头。

``这样的消息，若只留在人类那里，森林里便还是听不见。''

``是。''古月娜说，``这里发现的猎杀线索，外面也常常不知道。哪一伙人在做事，货物怎样运走，有时要两边合起来，才看得清楚。''

帝天望向图上那段空白。

``若有船，有货物，便总会留下经过的地方。''

他向万妖王看了一眼，后者的根须轻轻移过图边。

``过去那些通向西边的踪迹，也可以重新看。''万妖王说。

赤王听着，终于将注意从那座尚未建成的塔上移开。

它从前习惯看见猎杀便向外追，追到人类的城镇附近，又常常分不清哪一条路才通向真正的凶手。如今几个年轻人从海岸带回来的话，竟与森林里断断续续的线索接上了。

它仍不喜欢人类的那些势力，可若连谁在背后掠夺都弄不清，一腔怒火终究只能困在这里。

帝天的指尖仍停在那片标出的海域旁。

``这件事，我们会继续追查。''他说。

蔡须坤点头。

``我们也会把这里听见的消息带回去。无论建立传灵塔，还是对付这些势力，先让彼此能够传话，总比每一方都只知道自己遇见的那一点好。''

赤王抬头看着她。

``你们想让人类与魂兽联手？''

``想。''蔡须坤答道，``我们已经在学院说过。帅气教伤害魂兽，也在伤害人类。如今宏舒集团又出现了，有些事，需要一起面对。''

``人类那边，有多少人答应了？''

蔡须坤顿了一下。

``老师们听过，我们也来过总殿。但两大帝国与各地的人，还没有坐到一起。这件事，也得继续去做。''

``那你们也要回去问。''赤王说。

蔡须坤点头。

``会。''

``若有人不肯呢？''

``先找愿意的人。''她答道，``学院里已经有人愿意听，卢家也有做过的研究。我们先把能做到的拿出来，再让更多人看见。''

她停了片刻，望向石地旁的狼王。

``我先前说过，想到了一个完美的办法。可听你们真正问起来，才知道，有些地方我们只写了几行，落到不同族群身上，便是完全不同的一段日子。''

``那就继续补。''狼王说。

它说得很直接，蔡须坤也轻轻点了一下头。

``嗯。总得让这几页纸上的事，真的有人去做。''

赤王望着她，半晌没有开口。

它当然可以再说一句不信，将所有顾虑留在原处。几个年轻人离开以后，森林仍是森林，外面的人类仍在外面。遇见猎杀，再去追；遇见报复，再将族群聚起来。

这些年，它已经这样做过许多回。

可失去的那些，始终没有因此回来。

森林里新的幼兽慢慢长大，有的族群却还没有填上当年空出的位置。帝天能令争斗暂时停下，碧姬能替受伤者疗伤，万妖王也能察觉许多经过林地的踪迹。

他们都有各自的力量，仍没能让那些从外面伸进来的手彻底停住。

如今，海的另一边又有了宏舒集团。它们没有因为森林闭紧边界，便少觊觎这里的生灵。

狼王顺着赤王的目光，看见岸边一头带幼兽饮水的魂兽。那只小的还没有察觉这里正在谈什么，低下头喝过水，又忙着追一片浮到岸边的叶子。

``你若怕那座塔以后变了，就更该有人亲自去看。''狼王说。

赤王转过头。

``你倒愿意往人类身边去。''

``我已经在那里过了一段日子。''

狼王看向身旁的小队。

``我看见过他们害怕，也看见过他们疲倦。有些事，我原先只在森林里想，想出来的样子，与后来遇到的并不一样。''

``我想让族人听见，也想让自己继续看清楚。活下来了，总得为往后的事做一点打算。''

赤王听着，伏在石地上的前爪渐渐松开。

它还没有答应，但这一次，再看那几张纸，心里已经不只剩下反对。

狼王走到平石侧面，低头看向几人最初画下的那一圈位置。

``有一件事，也该添进去。''

蔡须坤拿起笔，等它说。

``别只让我这样的强大魂兽来讲。我的经历能让它们看见一条路，可年限不同、习性不同，能承受的也不同。''

它抬头望向湖岸外的林地。

``我的族人里，有的还不懂人类的话。若只听见族里最强的那个说可以，它们便跟着点头，这样的点头，未必真的知道自己要做什么。''

帝天看了它一会儿。

``你担心连同族的意思，也会被替它们说掉。''

``是。''

狼王转头看向蔡须坤。

``你曾等我自己回答。我希望后来走进去的，也能有这样的机会。会说话的就自己说，不会的，让熟悉它们的去听，别只问站在最前面的首领。''

古月娜点头，将笔记向蔡须坤面前移近了一些。

``这一点，我们一起记。''

卢本韦也开了口。

``父亲在养猪场做过的群体共鸣，和眼前这一种又有不同。雨浩身边也有魂灵，各自形成联系的经过，并不完全一样。将来要讲，就得将不同的条件分开讲。''

``不能看见蔡须坤与狼王相合成功，就让一个从没试过的人，照着同样的分量去接。''

碧姬轻轻颔首。

``身体、魂息与承受的力量，都要看。''

狼王听完，又望向赤王。

``我愿意参加最初的讲述，也愿意听别人把难处带来。可她还有自己的路，我也不可能永远替所有来者回答。''

``若有一天，我的经历可以被别的魂兽读懂，后来又有新的经历留下，那时它们便不用每次都等我们在场。''

蔡须坤将最后一句记完，抬眼望向它。

狼王的魂影仍有一部分映着湖光，可它已经不再只看自己的前爪。它一边想，一边将能做的事慢慢说出来，连万妖王也安静听着，没有打断。

赤王看了许久，终于将原本压在石地边缘的前爪收了回来。

``还有修行。''狼王忽然补充。

碧姬向它望来。

``成了这样以后，还能继续往前？''

``能。只是不再照着原来的身体去运转，每一步都要重新熟悉。''

狼王望向蔡须坤。

``她在人类那里学过许多方法。有些问题，我生前只靠年岁与力量慢慢磨，跟着她以后，才知道也能这样一层层拆开。''

``她自己也有看不清的地方。我熟悉的魂息，她未必熟悉；她擅长的分析，我又未必能想到。遇上了，我们就慢慢商量。''

蔡须坤轻轻点头。

``它有时候会先觉出一处变化，我再回去检查，才找到自己漏掉的地方。''

赤王听到这里，目光又在两者之间停了一会儿。

它最初看见的，是狼王失去身体以后，借着一个人类的魂力留存。如今听它们一件件说下来，才明白那道联系里，也有彼此补上的东西。

狼王抬起前爪，望着爪缘透过去的光。

``我知道自己付过什么代价。但往后还能学，还能把没有走完的地方继续走，我也确实盼着。''

它将前爪落回石地，站得很稳。

赤王望着它，终于没有再把那道魂影看作仅仅被救下的一点残存。

帝天看向她。

她没有移开目光。

旁边的古月娜轻轻抬起手，将被风吹到纸面上的一片叶子取走。碧姬仍拿着那页关于魂灵的记录，看了狼王一眼。

狼王已经在石地旁坐下，魂影随着呼吸轻轻起伏。

湖边诸多深厚的气息中，那道轮廓显得很小，却始终清楚地留在那里。

赤王忽然站了起来。

五人下意识地望过去，见它没有向前，只缓缓走到了临水的地方。

一层细小的波纹从它脚边散开。

它望着水中那道赤色的轮廓，想起自己方才说出的每一句反对。那些话没有哪一句是随意出口，可说完以后，若仍只是让这些年轻人离开，眼下的一切又能改变多少？

狼王回到森林，还能站在这里，已经让它亲眼看见了一种与从前不同的结果。

天州没有被抹去。纪念馆还在，巨坑还在，森林里失去的也都还在。可几个年轻人经过那里以后，带回来的没有另一场兽潮的号令。

他们在谈怎样让两边说上话，谈一只不懂人语的魂兽该由谁去听，谈狼王今后还能为族群做什么。

赤王回过头，望向帝天。

这些年，帝天一直站在诸多族群之间。强的要争，弱的要活，来往的旧怨有时压下去一阵，又换成另一副模样回来。它见过，也帮着挡过，却从未觉得这条路会自然走到尽头。

至于森林之外，帅气教没有停手，宏舒集团又在海上出现。若双方继续只隔着仇怨看彼此，藏在后面的那些人，反而能不断从空隙里伸手。

它的视线最后落到蔡须坤身上。

``我还是记得当年那些事。''

``我们知道。''蔡须坤说。

``以后听见有魂兽被人害死，我也不会因为今日谈过一座塔，便忍着不管。''

``该查的，要查。该救的，也要救。''古月娜接着道，``那座塔若能帮上忙，就该让你更早知道，究竟是谁做的。''

赤王看着她，沉默片刻，终于慢慢点了一下头。

``我方才一直在想，你们说的这条路，哪里还会出错。''

它重新走回平石旁。

``现在想过别的几条，才发现，照旧关着森林，照旧遇见谁便追着谁去报仇，也并不比它更好。''

``眼下这样的局面，让双方都能开口，也都有人在场照看，已经是我们能争取到的最好结果了。''

它说完，向那页写着``传灵塔''的纸低下头。

``我同意先做起来。''

蔡兴坤抬起眼，奥斯卡一直悬着的肩背也慢慢松了些。卢本韦没有急着去动笔，仍望着赤王，等它将后面的话说完。

``等要选地方，我也要看。该让族人知道的，我会说。若真有魂兽愿意出去，我也愿意出力照应。''

``可你们带来的人，也得像今日这样，肯听完我们的话。''

蔡须坤认真点头。

``会把愿意认真来谈的人带来。''

狼王抬起头，看了赤王一会儿。

``那到时，有些问题，我也问你。''

赤王低低哼了一声。

``你今日问得还少？''

狼王的耳朵轻轻动了一下。

湖边终于响起一声很轻的笑，随后又安静下来。没有谁将这番话当作所有仇怨都已消失，可当赤王重新伏下身时，它已主动留在了那几页纸旁边。

``可以让各族来听一听。''碧姬先开口，``有的事，听它亲自说，会比听我们转述更明白。''

万妖王没有立刻反对。

``也要让它们把问题问出来。''

狼王抬起头。

``我会回答。''

帝天站在湖边，望向远处的树林。几个呼吸以后，他转回身，将散在石面上的纸重新拢到一起。

``这几日，你们先留在湖边。''

蔡须坤望向他。

``我会让几个族群各自派代表过来，也会传信给狼群。''帝天说，``先将共生的事讲清楚。它们有什么顾虑，你们当面听。''

卢本韦按着纸角的手终于松了一些。

帝天又将那一页写着``传灵塔''的纸单独放到上面。

``至于第一座塔，地方要一起选。到了人类那边，你们也去找真正愿意来谈的人。''

蔡须坤接过纸，认真点头。

``我们会去。''

狼王没有离开平石，又向前挪了一点。

``我的名字，也可以写在愿意来讲的人里。''

蔡兴坤拿起笔，先看向它，见它点头，才认真落下几个字。

赤王看着那一笔一画，伸出前爪，将被风吹起的一角压住。

帝天扫过石地旁的几道身影，没有再催他们作出更多承诺。今日能说清的，已经由他们各自说了出来；需要带回去继续问的，也都留在纸上。

蔡须坤看向古月娜，轻声叫了一句``娜儿''。

古月娜走到她身边，与她一起将散开的几页重新排好。

最上面仍是``传灵塔''三个字，下面已添上许多从森林这一边说出的话。

赤王望着那些新字，慢慢收回前爪。

``先把今日商量的做好。''

``嗯。''蔡须坤抬起头，``一起做。''

\chapter{第一百四十章：故人现身}\label{ux7b2cux4e00ux767eux56dbux5341ux7ae0ux6545ux4ebaux73b0ux8eab}

赤王的前爪刚从纸边收回，远处忽然传来一个声音。

``就定在天州吧。''

声音不高，像是有人站在几步之外，听完了方才的话，顺口接上了一句。

可那片林子离湖岸还有很远。

蔡须坤抬头的时候，湖边的气息已经变了。

赤王猛地站起，收在鬃毛间的火焰一下蹿高。万妖王垂在石边的根须骤然缩回，近处几株古树的枝条同时绷紧，叶片擦过彼此，发出一片急促的声响。

碧姬转过身，袖边浮起了翡翠色的光。

湖岸外围，几头原本伏着休息的十万年魂兽几乎同时抬头。水边一声低吼才起，林中便有另一道气息应了上来。

五人方才好不容易放松的肩背，再次绷紧。

奥斯卡的手停在行囊旁，指间已有魂力凝聚。蔡兴坤弯腰将黑子鸡抱起，古月娜则向前挪了半步，站到了蔡须坤身侧。

灭世狼王转向来声处，魂影的边缘微微收拢。

直到这句话传来，他们才发现，那边已经有人。

万妖王的声音从树影里落下，比先前低了许多。

``他来了？''

它没有说名字。

赤王身上的火却在这一刻停了一停，随后又贴着脊背翻卷起来。

帝天始终没有离开原来的位置。

他的目光越过湖岸，落向远处两株古树之间。片刻以后，他抬起手，止住了外围将要传来的第二声兽吼。

``既然都来了，那就请出面共同商议。''

湖风吹过他的衣摆。他的声音很稳，听不出惊慌，也没有刻意拔高。

``我知道你对我们没有恶意。''

林中的叶声渐渐低下去。

一片黑色袍角从树后露了出来。

来人沿着林间的空隙走近，兜帽遮着眉眼，衣袖随着步子轻轻摆动。他没有凌空越过湖岸，也没有拨开挡路的枝条；经过一处低垂的树冠时，只稍稍低头，便从枝下走了过去。

在他身后，还有一个紫衣人。

那人身形清瘦，走得不紧不慢。紫色衣袍的下摆扫过草叶，露出一双沾着薄薄尘土的鞋。

卢本韦看见那道身影，原本握紧的手忽然松开了。

她先望了望紫衣人的脸，又看向前面的黑袍人，嘴唇动了一下，竟一时没有出声。

黑袍人走到林影之外，抬手将兜帽掀了下来。

蔡须坤愣住了。

几日前，在良序圣殿石台旁，耐心听他们讲那些矩阵设想的人，此刻就站在生命之湖前。

同样的眉眼，同样平静的神情。

大长老，test\_alpha0。

``大长老殿下？''

蔡兴坤先叫出了声，随即又看向他的身后。

卢本韦终于向前迈了一步。

``干爷爷！''

紫衣人抬眼望来，那双带着淡淡粉色的眼睛在她身上停了停，清瘦的脸上终于有了一点笑意。

``这回认得倒快。''

正是养猪场的第一股东，phyrion。

奥斯卡看看他，又看看大长老，手里刚凝起的魂力慢慢散了。

``您二位怎么一起来了？''

test\_alpha0走近一些，先向帝天微微颔首，才看向他们。

``在天州遇上了。''

他说得简单，蔡须坤却下意识地回头看了一眼古月娜。

他们穿过那座城的时候，谁也没有想到，大长老后来也会到那里。

phyrion的目光在五人身上转了一圈。

``听说你们回了大陆，去过学院，又往总殿走，我也出来走了一趟。''

``本想先见见人。到天州，正好碰上他。''

他说着，向test\_alpha0侧了侧头。

卢本韦已经走到他面前，仰起脸，又认真看了一遍。

``我父亲呢？''

``在家。''

phyrion伸手，将她肩上翻折的一角衣领理平。

``好着呢。你不在，他照样有做不完的事。''

卢本韦忍不住笑了一下。笑过以后，她才想起自己仍站在一群凶兽之间，赶紧侧开身，将湖边的位置让出来。

``我们正在说传灵塔的事。''

``听见了一些。''

phyrion望向平石上的旧笔记，脚步稍稍缓了缓。

旁边的赤王仍没有伏下去。

它的目光一直留在test\_alpha0身上，火焰沿着肩背缓缓流动。先前说起天州时，那一拳仍只存在于回忆里；此刻，出拳的人却已走到了湖边。

它胸前的伤早已好了，可看见那道身影，身体仍先于思绪绷紧。

万妖王的根须停在地面，没有重新探回纸边。

碧姬望了test\_alpha0一会儿，袖上的光缓缓淡去。

只有帝天的神情始终没有多少变化。

``你知道他们会来这里。''

``知道。''test\_alpha0说，``他们离开总殿前，提过这件事。''

``所以，你也来了。''

``魂灵的事，我早年便同人谈过。如今有了新的进展，我想来听听。''

test\_alpha0望向石地旁的灭世狼王，又将目光移回帝天。

``不过，你们方才谈得很好。我便在林外等了一会儿。''

帝天与他对视片刻。

湖面有风经过，一道细长的水纹从岸边推了出去。

蔡须坤站在两人之间稍远的位置，没有立刻插话。她在纪念馆里看过那场战斗留下的东西，也刚听赤王说过当年的感受，如今看着这两人平静相对，反而比先前更加留心。

那几页草图仍摆在平石上。

她伸手压住被风吹起的一角。

帝天的视线随之落下。

``我们刚谈到地方。还没有定下具体位置。''

``看见了。''

test\_alpha0低头，看了一眼纸上那个圈。

天州城西一带，待共同踏勘。

旁边添着许多新字，有些挤在最初的屋舍之间，有些写在页边。水岸、林带、道路，几处还用线牵在一起。

他看得很仔细，没有伸手将那些字抹掉。

``先把天州这一带列下来，具体地方，我们一起去看。''他说。

赤王终于开口。

``为何偏偏是天州？''

test\_alpha0抬起头。

赤王没有退，迎着他的目光，将话说完。

``那里发生过什么，你比他们更清楚。''

卢本韦的笑意慢慢收了回去。

蔡须坤望着赤王，没有插话。它方才已经答应帮忙，也愿意去看地方；可那场旧战横在中间，总得让它当面问上一句。

``我清楚。''test\_alpha0回答。

他没有避开天州的名字。

``我刚去过那里。''

古月娜抬起眼。

``纪念馆？''

``也看了。''

test\_alpha0停了片刻，目光越过平石，望向森林之外。

``你们进林以后，我在那座石像前停过一会儿。''

蔡须坤脑海里浮起那柄高大的石剑。石像的双手拄在剑柄上，目光朝着森林，坑底的人仰起头，也只能看见高处的一张脸。

她从未想过，那张脸的主人，也会站在那里，重新看一遍自己留下的东西。

``那个坑还在。''test\_alpha0说，``纪念堂里的名字，也还在。''

赤王身上的火焰轻轻伏下一层。

``森林里也有。''

``我知道。''

这一次，test\_alpha0的目光落回它身上。

``那些不会因为修起一座塔，就全都过去。''

赤王没有接话。

test\_alpha0看向帝天，继续道：``天州靠近森林，城外也有往来多年的路。你们担心族人进城以后受困，这几个人刚才已经把位置往林缘挪了。能顾及的，先放在一起看。''

``还有城里的人。''帝天说。

``他们的意思，也要问。''

test\_alpha0俯下身，指向草图上留出来的道路。

``受伤的魂兽要来，魂师也要来。若能在那里彼此见上一面，有些人再想起这片森林时，便能多记住一些活着的身影。''

他说完，将手收了回来。

``但地方选得合适，才做得下去。你们记下的水岸和退路，都留下。''

万妖王的根须终于动了一下。

``林带也要留。''

``留。''

``水不能断。''碧姬说。

test\_alpha0点头。

``踏勘时一起看。''

几句很短的回应以后，湖边那股始终绷着的气息，才慢慢松开一点。

赤王低头望着图，没有再将火焰向外推。它仍站着，爪尖却已从绷紧的石缝边移开。

帝天侧过身，让出了平石另一侧的位置。

``那就一起看。''

test\_alpha0走过去，停在与他相对的一边。

蔡须坤原本想将那张画得最清楚的纸递到大长老面前，手抬起一半，又停了停，转而将它摆在了两人中间。

test\_alpha0看见这个动作，目光在她脸上停了一瞬。

``这样就好。''

蔡须坤应了一声，重新将纸角压稳。

另一边，phyrion已经在旧笔记旁停下。

``带了这么远？''

卢本韦点头。

``大长老同我们说了卢岳先祖的事。刚才，前辈们也说，他们记得那头柔骨兔。''

phyrion的手指从发黄的纸边轻轻掠过，没有立刻将书拿起来。

``这几页，你曾祖翻得很勤。''

他望着夹在书里的纸签，声音低了一些。

``父亲晚年，添过，划过，到了最后也没有把它合上，说以后不看了。''

卢本韦安静下来。

她知道他说的父亲，是曾祖卢承业。许多从族谱里看来的往事，在这位长辈口中，却仍是亲手接过、亲眼看过的一叠纸。

phyrion终于抬头，看向湖边诸兽。

``失败的那一次，我们家记着。你们记着的，我也愿意听。''

碧姬轻轻点头。

``她刚才没有瞒着。''

``那就好。''

phyrion看了卢本韦一眼。

她本来已经准备了好几句解释，被这一眼看得停住，最后只认真点了点头。

他的目光又落到她身侧的母猪扽魂上。

母猪也望着他，耳朵微微动了动。

一团很淡的白光在紫袍下方浮起。奥斯卡还未看清，脚边已多出一头白玉似的小猪。

个头不大，通身雪白，眼睛带着柔润的粉色。

它站稳以后，先拱了拱phyrion的袍角，随后认出了卢本韦，慢悠悠地走到她腿边，拿额头蹭了一下。

卢本韦低头，忍不住伸手摸了摸它。

``你也来了。''

小白猪眯起眼，轻轻哼了一声。

奥斯卡盯着它，往昔那张暗金天幕被撞开的景象忽然又回到眼前。他下意识看了看平石，再看了看这头正在蹭人衣角的小猪，脚步往旁边让了一点。

phyrion瞥见了。

``今天不用它撞东西。''

奥斯卡咳了一声，挪开目光。

蔡兴坤却没忍住，低头笑了出来。黑子鸡从她臂弯里探头，看着那头小白猪，难得没有立刻叫。

近旁一头十万年魂兽凝视着那道小小的身影，没有靠近。

万妖王的根须也避开了它脚边的位置。

小白猪走过以后，phyrion才重新看向草图。

``本韦刚才说，家里可以出一份。我听见了。''

卢本韦抬头。

``我只说回去同父亲商量，还没\ldots\ldots''

``商量就是了。''

phyrion将她身边歪着的一页纸扶正。

``想做什么，要用多少东西，先列出来。能拿出多少，回去一起算。建到一半缺了，来回折腾，苦的是已经住进去的。''

卢本韦望着他，眼睛一点点亮起来。

``您愿意帮我们？''

``我都站在这里了。''

他又看了一眼旧笔记。

``总不能只让你带着这本书，来回跑。''

卢本韦嘴唇动了动，最后应得很轻。

``嗯。''

赤王看着这一老一少，过了一会儿才开口。

``人类出人、出东西。魂兽到了那里，仍要自己说话。''

``应该。''phyrion说。

他没有停顿，也没有因为面前是一头凶兽，便换一种哄劝的语气。

``不然，今日何必坐在这里。''

赤王望了他一会儿，缓缓伏低前身。

地面没有震动，火焰也随之收了回去。

灭世狼王站在原来的位置，看着平石旁添上的两道人影，直到此时才缓缓开口。

``那我的问题，也还算数？''

test\_alpha0转向它。

``算。''

狼王向前走近一点。

``我可以说自己的经历，可我不能替每一个族人答应。它们来了，有些话，你们还得重新听。''

``那就重新听。''

test\_alpha0看着它，没有越过它去问蔡须坤。

``你刚才提过，不同族群能承受的不同。以后整理经历，也请你把这一点一起留下。''

狼王望着他，过了片刻，点了一下头。

蔡须坤站在一旁，掌心压着纸，心里那点因大长老突然出现而生出的不安，终于缓缓落下。

她起初还怕，几位强者走到一起以后，方才那些一点点问出的顾虑，会被轻易盖过去。

如今狼王仍站在那里，万妖王还在看林带，赤王也没有收回自己的那句话。

她低下头，将最初写着``传灵塔''的那页纸，重新放到众人都能看清的位置。

古月娜走近，轻轻按住另一边。

``娜儿。''蔡须坤低声叫她。

古月娜抬眼，顺着她的目光看过去。

帝天站在平石一侧，test\_alpha0在另一侧。那两道身影之间，隔着一张还没有画完的草图。

古月娜的手指停了一会儿，随后将纸上的一小片草屑拈了起来。

``继续吧。''她轻声说。

帝天听见了，伸手点向林缘的标记。

``先说这里。附近哪些族群需要来，我来问。''

test\_alpha0看向天州的墨点。

``城里这边，我也去见见。''

phyrion将旧笔记往中间移近了一点。卢本韦拿起笔，奥斯卡也重新展开了记录物资的那页纸。

蔡兴坤蹲下身，将黑子鸡安置在脚边，抬手接住被风吹起的纸角。

湖岸外围，几头魂兽慢慢伏回原处。

万妖王的根须再次伸来，停在那条尚未画完的林带旁。

蔡须坤顺着它指的位置，将那条线继续画了下去。

test\_alpha0看着那条新添的线，过了一会儿，将手指移向天州城西的另一处。

蔡须坤认出那个位置，笔尖停住了。

那里是他们前几日去过的纪念馆。

``具体的地方，我有个提议。''

test\_alpha0抬起头，看向帝天。

``纪念馆那座我的雕像，拆了。原址改建，做成第一座传灵塔。''

卢本韦刚翻开的纸页停在半空。

奥斯卡也抬起头，以为自己听错了。

``您是说\ldots\ldots 那座雕像？''

``嗯。''

大长老的手仍落在草图上，神色没有变化。

蔡须坤眼前又浮起坑底那道高大的身影。石刻的长剑立在身前，双手拄着剑柄，比城楼还高的雕像面对森林，已经站了许多年。

大长老说要拆的时候，却只轻轻点了一下那处位置。

赤王缓缓抬起头。

``人类在那里纪念你。''

``我知道。''

test\_alpha0望了它一会儿，接着道：``纪念堂留下，馆里的记录也留下。死去的人叫什么，那一战怎样发生，都要继续让后来的人看见。''

他将手指从纸上收回来。

``至于我的石像，可以拆。那地方若能让受伤的魂兽得到照料，让愿意尝试共生的双方找到彼此，便值得这样用。''

碧姬望向草图，神色渐渐认真。

``可那是在坑底。''

``所以原来的石梯不够。''test\_alpha0说，``要重新看地势，修能够上下的坡道。水怎样引，林带怎样接，也都得一道安排。''

万妖王的根须沿着纸边移过，停在纪念馆与森林之间。

``来的族人，不能全靠飞进去。''

``那条路，要让它们自己走得通。''

test\_alpha0看向蔡须坤面前的纸。

``把你们方才画的通道接过来，一起看。到了现场，哪里不够，就接着改。''

蔡须坤慢慢点头，在纪念馆的标记外添了一道线。古月娜按着纸角，看着那道线穿过旧战场，一直连向林缘。

帝天沉默了一会儿。

``天州的人，会答应？''

``我去同他们说。''test\_alpha0回答，``建塔的事，也请他们一起来谈。纪念馆里的那些名字，不动。''

赤王望着他，没有立刻接话。

它曾经从那片土地里挣扎着站起，如今再想起那个方向，仍会先记起当年压下来的力量。它没有想到，最先提出拿掉那道巨大身影的人，竟是眼前的大长老。

湖边安静了片刻。

test\_alpha0重新看向帝天。

``还有掌事的人。塔建起来以后，总不能每出一件事，都只等我们几个赶来。''

帝天抬眼，示意他说下去。

``你们星斗方面，派出十位代表。我们扽魂殿、养猪场和帝国这边，共同派出十人。''

卢本韦将悬着的纸页放下，拿起笔。

``双方各十位？''

``各十位。''

test\_alpha0答得很清楚。

``每位代表一票，双方票数平等。以后人类与魂兽之间有了争端，你们也有权力调查、审议和处置。''

万妖王垂在纸边的根须停住了。

``若是人类伤了魂兽呢？''

``同样交到双方代表面前。''

test\_alpha0看向它。

``受害的魂兽可以来陈述，证据拿出来，由双方共同查明、共同裁处。不能因为犯事的是人，便把你们的代表挡在外面。''

``反过来，也一样。''帝天缓缓开口。

``一样。''

大长老点头。

``要管人兽争端，就得听到两边的声音。查清以后，该由谁承担，也由两边一起议定。''

赤王向前挪了一点。

``你说的处置，是我们真能参与定夺？''

``真能。''

test\_alpha0没有让它等。

``十个席位，都要有相应的权力。这一点，写进去。''

卢本韦握笔的手顿了顿，随即将那句话添在代表人数的下方。

她先前说过，事情是谁的，就该让谁开口。直到此刻，纸上才有了能让这句话往后继续算数的位置。

帝天低头看着她写下的字。

``若是十票对十票？''

``分歧没有说清，就继续查，继续议。''test\_alpha0说，``不能添一张只属于人类的裁决票。即使我在席，也只算一票。''

帝天看了他许久。

随后，他将手放到了草图旁边。

``星斗这十位，我会同各族商议。''

碧姬轻轻点头。灭世狼王也向前走近了些，低头望着那两行相对而列的人数。

人类，十。

星斗，十。

它看了一会儿，才抬头道：``有族人受了委屈，至少知道该到哪里去说。''

``也要有人听完。''碧姬补了一句。

``会写清楚。''蔡须坤说。

她另取了一页，将方才关于接待与照料的记录放在一旁。那几张纸上的事，终于能同代表们要做的事接到一起。

phyrion望着卢本韦笔下的字，开口道：``养猪场这边，我赞成。出了多少材料、多少工，也不多占一票。''

卢本韦抬头看他，随即认真将这句也记了下来。

test\_alpha0看向帝天。

``人类的十个名额，我会同养猪场和两大帝国商议。具体人选再定，双方平等这一条，先写进提议。''

帝天点了一下头。

``名单定下以后，双方见面，把该管的事一件件说清楚。''

赤王低头，鼻息掠过石面，将靠近它的纸角吹动了一下。它停了停，又伸出前爪，轻轻压住那一角。

``我的意见，也记着。''

卢本韦望向它。

``在记。''

蔡须坤将最初那张草图移到中间，在纪念馆的位置旁，写下了``第一座传灵塔''。

古月娜的目光随着笔尖落下。

那个位置，她曾站在石梯上仰望过。如今将要立起来的塔，会让森林里的族人也有一条能够走进去的路。

她伸手，压住纸的另一边。

``通往林缘的这一段，踏勘时，我带路。''

蔡须坤抬眼看她，轻轻应了一声。

湖风再次吹来时，草图的两端都有人按着。

新添的几页纸被分开摆好，奥斯卡与卢本韦低头核对记录。test\_alpha0仍站在平石旁，听蔡须坤说沿途见过的道路，phyrion则将压在旧笔记下面的一页纸抽出来，放到他们面前。

帝天走到临水的地方，停了下来。

万妖王的身影从树阴下移近，望了一眼平石旁的人，开口时将声音压得很低。

``兽神大人，我们真的答应他们吗？''

帝天没有回头。

万妖王停了停，声音更低了一些。

``今日说的是自愿。可怕是不久的将来，魂灵就变成`被迫自愿'了。''

它的目光落向那头站在蔡须坤身侧的狼王。

狼王能自己说话，能自己选择，可往后走进塔里的每一头魂兽，未必都有它这样的力量。若有人握着它们赖以活命的东西，再让它们点头，那张写着自愿的纸，又能留下多少真正的意愿？

帝天终于侧过头。

``那你想怎么样？''

他的声音同样很低，神情却仍平静。

万妖王张了张口。

帝天看向平石旁的黑袍人。

``他在这里这么久，我们没有一点发现。''

test\_alpha0正伸手扶住一页滑落的纸。蔡兴坤将旁边的石块递过去，他接在手里，压住纸角，动作没有半分急迫。

帝天的目光停在那里。

``他只要想，我们一个都活不下来。''

万妖王没有再开口。

先前那一句``就定在天州吧''，仍清清楚楚地留在它耳边。声音传来之前，根须探着地面，枝叶覆着林冠，湖边又有这么多强者，可谁也没有察觉林外已经站着人。

过了片刻，它缓缓低下头。

``只能看他们能不能做成这件事了。''

它的声音里仍有疑虑，说到后面，却多了一点很旧的迟滞。

``就像当年的那个人\ldots\ldots Opt。''

那个名字落下以后，帝天许久没有接话。

平石旁传来翻动书页的轻响。卢本韦正将新写的记录与卢家旧笔记放在一起，身侧的母猪轻轻动了动耳朵。

帝天望着那边，终于缓缓开口。

``已经一万年了。''

\chapter{第一百四十一章：传灵初试}\label{ux7b2cux4e00ux767eux56dbux5341ux4e00ux7ae0ux4f20ux7075ux521dux8bd5}

回到天州以后，蔡须坤才发现，湖边那几张纸，远远装不下建起一座传灵塔要做的事。

第一天，他们跟着城里的工匠重新走了一遍下到坑底的石阶。第二天，碧姬与几位族群代表沿着水道走出去，回来时，图上原本只画了一条细线的地方，被添上了好几处标记。

有的坡太陡，有的土层太薄。纪念馆到林缘之间那条看着宽敞的路，也并非所有地方都能让体型庞大的魂兽转过身。

万妖王指过一处老树生长的低地，那条准备穿过去的路，便在图上拐了一个弯。

蔡兴坤拿笔改过，抬头望向那边，又把弯画得缓了一些。

``这样，大一点的也能过去。''

奥斯卡扶着铺在石台上的图，等墨干了一点，才将它卷起。

古月娜走在林缘与城外之间，将两边没能说明白的话重新带过去。有时同一个问题，上午才说过，下午换了一个族群，又要从头问。

五人没有觉得烦。

他们曾以为，找到愿意点头的人，事情便能往前走上很远。真正站在这一片土地上，才知道那一声答应之后，还有许多脚下的路，要一段一段走出来。

消息从天州送出去，来的信也渐渐多了。

又过了几日，Smile Lee到了。

他下了车，便绕到纪念馆外，低头检查基座旁的几处旧阵纹。蔡须坤等人赶来时，他正蹲在一块石板旁，指间夹着一片薄薄的金属，将上面的灰一点点刮开。

``别踩这里。''

奥斯卡刚要落脚，听见声音，立刻把脚收了回去。

Smile Lee抬头看了他一眼。

``这道线还接着旧灯阵。拆到这里，得先断开。''

他把那片金属收回指间，站起身，才向五人点了点头。

``都回来了。''

卢本韦看着他手上的细灰，忽然觉得，这一句比方才满城的迎接声都要熟悉。

梅天狸是同他一起来的。

她从纪念馆门前走过来，先看五人的脸色，随后目光落到蔡须坤手腕上。

``这几天，混沌之力可还安稳？''

``安稳。''蔡须坤回答，``没有再像训练场那样用。''

``那就好。''

梅天狸抬头望向高处的雕像，许久没有说话。

脚手架已经搭到了基座上方，工匠将要拆取的部分一一标好，长剑与手臂之间系着粗索。真正动工前，Smile Lee还要将旧阵法的连接逐处理清。

雕像下面，仍有人进出纪念馆。

一位头发花白的老人停在门前，看着上面的绳索，忍不住问：``以后，是不是就没有这个了？''

蔡须坤转过头，还没有想好如何回答，test\_alpha0已经从馆里走了出来。

``石像会拆。''他说，``馆里的名字和记录，都留下。''

老人望向他。

起初，那双浑浊的眼睛只是停了一停，随后便慢慢睁大。他抬头看向雕像，又低下头，目光重新落在眼前人的脸上，嘴唇颤了几下，却没能说出话来。

``您\ldots\ldots 您是\ldots\ldots''

他伸手去扶栏杆，手指却一直在抖。

蔡须坤赶紧往前一步，test\_alpha0已先扶住老人的手臂。

``慢些。''

老人听见这两个字，眼泪忽然就落了下来。

``我小时候，听曾祖父讲过\ldots\ldots''

他想把话说清，声音却哽住了，只能抬起袖子擦了一把脸。

``当年兽潮来的时候，他还只是个普通的守城士兵。别人让他守哪里，他就守哪里，搬箭，抬木头，伤了人，就帮着往城楼底下送。''

老人紧紧攥着栏杆。

``他说，那一天，他以为自己回不了家了。''

馆门前几人渐渐静下来。

``他总同我说，后来城外来了一个人。天上都是黑的，又有蓝光亮起来，城墙在晃，人连站都站不稳\ldots\ldots 可最后，城还在。''

他说到这里，抬起头，已经满脸是泪。

``他活着回家了。以后才有我们这一家子。''

蔡须坤看着老人，忽然想起馆里那些急促的军报。纸上只是一行``守军退至城楼后''，落在眼前这一家人的日子里，却是一位曾祖父后来还能坐在家中，给孩子讲起从前。

老人还望着test\_alpha0。

``曾祖父说，他隔得远，没能当面谢您。每回带我到这里来，都要指着石像讲上一遍。''

他抬起手，似乎想将眼前的人再看清一些，手到了半空，又怕失礼，赶紧收住。

``我只见过这座像。我哪想过\ldots\ldots 有一日，您真的就站在我面前了。''

他再也说不下去，低着头，肩膀轻轻颤着。泪水顺着脸上的皱纹流下，他擦了一次，很快又落了一脸。

``谢谢您。大长老殿下\ldots\ldots 我替他，也谢谢您。''

test\_alpha0扶着他的手臂，待他站稳，才缓缓开口。

``你曾祖父守在城上，也护住了城里的人。他那一份，也该记着。''

老人张了张口，最后只是用力点头。

大长老陪他在门边停了一会儿，等那双手不再抖得那么厉害，才指向馆旁铺开的草图。

``这里往后要建传灵塔。纪念馆还留着，人和魂兽有了事，也能来这里找人商量。''

老人顺着他的手看过去，看见图上通向森林的那条路，又望了望仍立在高处的石像。

过了许久，他慢慢松开栏杆。

``名字还在，事情还记着，就好。''

他又拿袖子擦了擦脸，站直一些，向大长老认真行了一礼。蔡须坤扶着他走到台阶边，直到家人从上面赶来接住，才松开手。

梅天狸站在门边，看着他走远，将手里的图纸重新展开。

``那这一侧，得把老人和孩子走的路也留出来。''

她用指尖点了点。

``别让来看纪念馆的人，同运材料的车挤在一起。''

蔡兴坤俯身，将那一处也记了下来。

两大帝国的使团随后抵达。

天斗来的是雪星亲王。

车停下以后，他扶着车沿站稳，一眼看见卢本韦等人，脚步便停了停。

``原来是你们。''

奥斯卡与同伴一同行礼。亲王望着他们如今的模样，片刻以后，轻轻叹了一声。

``都长大了。''

他没有在门前说太久，问过几人的近况，便将目光转向馆外铺开的图纸。

星罗皇室派来的亲王比他晚到半日，带着几名随行官员。来人先入纪念馆看了一遍，出来时，手中已多了一份城西地界与水道的抄图。

双方的旗帜立在临时会场外，隔着一段空地。

更靠近森林的一边，也为星斗来的代表留出了位置。

最初，总有人在门前多站一会儿。

人类的护卫看见魂兽身影，会下意识握紧兵器；从林中出来的族群代表，闻见车马与金属混杂的气息，也不肯立刻沿着指定的路进去。

赤王停在场地边缘，没有催它们。碧姬过去说了几句，直到里面的人把挡路的几辆车移开，它们才慢慢向前。

蔡须坤在门边看着，将那些没有画在图上的迟疑，也记在了心里。

会场借用的是馆旁一间腾空的屋舍。

门开得很宽，临近空地的一侧撤去了几扇窗。桌上铺着重新誊过的图与条目，旁边仍压着小队最初带来的那几页纸。

有些族群不惯于坐进屋里，座位便一直排到门外，谁开口，里面也都听得见。

有些墨迹已经被手掌蹭淡，边上却添满了后来人的字。

商议没有一次便结束。

星斗方面要确认通道与水岸，也问出了受伤魂兽来塔以后，由谁照料。两大皇室则要听清塔的管辖范围，以及城中人遇到争端，该向谁陈述。

养猪场能出什么，扽魂殿又能派来哪些人，也都在桌上逐项说过。

有几处谈到傍晚，仍各有不同的意思，只得将纸留下，第二日接着议。

最后，选址终于定了下来。

第一座传灵塔，建在天州城西纪念馆雕像的原址。馆中的记录保留，周围另作修整，通向林缘的道路与适合不同族群停留的水岸、林带，一并纳入建造的安排。

蔡须坤在誊清的图上看见那个位置，停了一会儿。

她还记得，自己第一次沿石梯走下去时，脚下的巨坑大得让人忘记了说话。如今，这一片土地上将要多出来的东西，已经被许多人一同看过、问过，又一同放回了纸上。

接下来，才是传灵塔议事堂的权责。

``双方各十名代表，每人一票。''

雪星亲王将面前的条目读过，抬头看向另一侧。

帝天坐在那里，神情平静。

``这一条，照旧。''他说。

星罗亲王低头翻过一页。

``但若某项决议，已经让一方认为无法承受，仅凭其余人的票数，仍要强行施行呢？''

屋内安静了一会儿。

test\_alpha0看向帝天。

``那便将最后这一道权力，也对等留下。''

他将笔放在纸边，缓缓说道：``传灵塔议事堂中，我与帝天，各有一票否决权。''

蔡须坤抬起头。

``任何一人否决，那项决议就不能强行通过，须重新商议。''

大长老说完，帝天的目光也落到同一行字上。

``两人都是如此。''帝天说。

``都是如此。''

星罗亲王看了看两人，随后将那句话重新读了一遍。

普通表决，仍是人类十票、星斗十票。两人各自的否决权，只能阻止决议，不能将未获同意的事强行定下。

万妖王坐在星斗代表一侧，许久才慢慢开口。

``即使是已经开始讨论很久的事？''

``也一样。''test\_alpha0说。

帝天望向它，点了一下头。

万妖王没有再追问，手指从纸边收了回去。

它知道，写下这些字，还要看后来的人怎样做。可如今，至少那一道原本只存在于兽神手中的守护，也被摆到了人类代表都看得见的地方。

两位亲王又与随行人员商议片刻，终于将条目放回桌上。

``此项，一并列入。''

phyrion站在卢本韦身后，看着她将最后一句誊清。

墨还未干，帝天便先伸出手，将那张纸移近。等他留下名字，test\_alpha0接过笔，也在旁边落下了自己的名字。

纸被重新展平，送回诸位代表面前。

蔡须坤看着那两个相邻的名字，忽然想起湖边听见第一声``就定在天州吧''时，万妖王骤然收回的根须。

万妖王低头看过那两个名字，又将纸递向身旁的族群代表。桌边接着响起了翻页声。

会散以后，几人没有立刻离开。

奥斯卡将借来的笔一支支收好，蔡兴坤与卢本韦核过最后几页记录，古月娜则把踏勘图重新卷起，交给负责继续勘察的人。

蔡须坤留下了另一份旧封套。

里面是从养猪场带出的问情谷手记抄本。她将它取出来，看过页边的几处断缺，又慢慢放回桌上。

梅天狸走近，目光在封套上停了一下。

``还在想着造化丹？''

``嗯。''

蔡须坤轻轻按住心口。被她护着的那缕魂息仍然很弱，贴着送入的瑞气，安静地留在那里。

她没有将牢大唤出，只确认了一遍，便收住外层的力量。

卢本韦抱着旧笔记，站在她身旁。

``我们想先在塔里做一次魂灵共生试验。''

她说完，看了一眼身侧的灭世狼王。

狼王已经能够独自现身，能够讲出自己这一年的经历。可真正等着传灵塔开门的魂兽里，多数没有它那样的修为；来问的人，也未必拥有蔡须坤的本源与承受力。

``我可以再留些日子。''蔡须坤抬起头，``先把这次魂灵试验做好，之后再去寻找乾坤造化丹。''

``天地灵水的消息，也还要继续找。''古月娜说。

``我们都记着。''

蔡须坤将封套仔细合好，移到了案边。

那一缕魂息仍在心口。她没有将寻找牢大的事放下，只把眼前需要自己留下的这一段时间，认真留了出来。

``这一次，我来。''

奥斯卡将收好的笔放下，望向几位长辈。

卢本韦转头看他。

``你想好了？''

``想过了。''

奥斯卡的目光从她手里的笔记移开，停在院外。碧姬这几日带来过一头石甲獾，他陪着看过，也同古月娜一起，慢慢摸清了它的意愿。

``它肯，我也愿意。''

他说完，又看向蔡须坤。

``问情谷还要一起去。我也想再往前走些，到了该出力的时候，能多撑一阵。''

卢本韦没有立刻接话，过了片刻，才将旧笔记重新放平。

``那我们陪着你。''

试验定在馆旁一处清静的院子。

Smile Lee带人将施工那一侧的声音隔开，沿院边布下记录魂力变化的阵盘。几块薄晶片嵌在槽中，亮起来以后，细细的光纹便沿着刻度缓缓游走。

奥斯卡在旁边看他调过两次，又问清了几处纹路亮起时代表什么，才慢慢坐到院中。

平时这时候，他总在替别人准备恢复香肠。如今轮到自己坐下，手里空着，竟有一点不习惯。

蔡兴坤把水放到他身侧。

``要用的都备好了。你先顾着自己。''

奥斯卡点点头，看向院门。

碧姬陪着那头石甲獾走了进来。

它已有七万年修为。

厚重的背甲叠成数道灰黑色的脊，边缘磨出了旧痕，颈边的毛也白了许多。它走得很慢，每抬起一只前爪，都要等身后的重量稳住，才向前再挪一步。

可它的目光仍清醒。

到了门槛前，它抬头望了一眼奥斯卡，辨认出这几日总在水槽边陪它坐着的人，便又慢慢向前。

奥斯卡起身，挪开挡路的矮凳，将它常停着的那块平石放到近旁。

石甲獾伏下以后，才终于长长地呼出一口气。

它的生命，已经快走到尽头。

碧姬检查过许多回。温润的生机能让它舒服一些，让衰弱的血肉得到滋养，却留不住正在散去的寿元。这副身体还能陪它多久，已经很难再说。

``它的主意识仍完整。''碧姬说，``想过的事，记得的人，都还清楚。''

古月娜俯身，与它重新交流了一会儿。

``它记得昨晚你问的那句话。''她抬头，看向奥斯卡，``变成魂灵以后，还能不能回来看看。''

``能。''奥斯卡说，``等它愿意，我们一起回来。''

石甲獾望着他，背上的灰脊间透出一缕微弱的魂息。

古月娜安静听完，将它的意思传给众人。

``它愿意试。它知道原来的身体留不住了，可它还想看看，新的塔会是什么样子。''

奥斯卡将手放低，停在它能够碰到的位置。

``我也愿意。''

这一次，石甲獾缓缓抬起前爪，轻轻抵住了他的掌缘。

院中诸人各自站定。

两位亲王与其余代表留在外间，窗上的薄帘放下来，里面仍能看清，嘈杂的脚步声却渐渐远了。

卢本韦把旧笔记与新记录并排放好。翻到卢守衡那一页时，她的手停了一下，随后认真记下今日参与双方。

奥斯卡。

七万年石甲獾。

七万年，比当年那头柔骨兔少了一万年，可谁也没有因此便放松下来。年限之外，还有魂息怎样相接、奥斯卡能接住多少，以及老獾怎样保住自身意识，都要看清。

phyrion在近旁坐下，脚边的小白猪也安静伏着。

他先看过奥斯卡的香肠武魂，又同碧姬、古月娜一道，辨认石甲獾尚能自主收束的力量。

``它七万年的积累，还很厚。''他说，``身体弱下来了，不等于里面只剩一点魂力。''

奥斯卡慢慢点头。

``我知道。''

``先将联系接稳。它留着自己的那一份，你接得住的，再一点点相合。''

帝天守在石甲獾一侧，大长老站到阵盘边。几位强者都收敛了气息，院中重新安静下来。

蔡须坤坐在奥斯卡斜后方，收住外放的瑞兽本源，只在碧姬示意以后，送出一缕温润的瑞气，养住老獾将要散逸的魂息。

古月娜保持着与它的回应。

梅天狸照看奥斯卡的状态，蔡兴坤盯住记录阵盘。卢本韦坐在两份笔记之间，手指已经握紧了笔。

``开始吧。''phyrion说。

奥斯卡摊开手，香肠武魂在掌心凝出温和的光。

他平日里用惯了这份力量。给同伴恢复，稳住消耗，或将香肠引成长索去接住一处空隙，都知道魂力应当怎样送出去。

可此刻，先迎过来的，是另一道意识。

像爪尖踩进湿土，像沉重的背甲贴着石面慢慢放下，还有漫长岁月里，曾在林间见过的一小片晨光。

奥斯卡呼吸微微一顿。

石甲獾的魂息沿着掌缘接近，先停了停，随后才触到他送出的那一线力量。

那份联系很轻。

可等它真正接上，背后压着的分量，却让奥斯卡肩背猛地一沉。

七万年修为积下的魂力，顺着才接通的地方缓缓涌来。他原本熟悉的运转路线，一下变得拥挤。掌心温和的光急促地涨了两次，边缘竟露出几道撑不稳的裂隙。

奥斯卡下意识想将它们全部压回自己体内。

阵盘上的细光骤然抬高。

``这里变快了！''蔡兴坤立刻报出位置。

卢本韦抬头，看见奥斯卡额边已经渗出汗。

``别硬撑。哪里接不住，就说。''

奥斯卡张了张口，声音有些发紧。

``掌心接进来了，后面\ldots\ldots 走不动。''

phyrion俯下身，看过那道拥在一起的光。

``你还在照平时那样，先把一切收到自己这里，再往外分。''

他指向两者之间尚未闭合的联系。

``这里还没走通。给它留出接续的位置。''

古月娜也抬起头。

``它察觉你撑得吃力，想退。你别拉着。''

奥斯卡手指颤了一下，慢慢放松。

那头老獾的前爪始终抵在他的掌边。随着他将拥住的魂力收细，石甲獾也一点点收回过重的部分，原先几乎绷断的联系，终于缓了下来。

奥斯卡低低喘了一口气。

卢本韦已挪到他身旁，扶住了因用力太久而有些发僵的手腕。

``每回都让我们别逞强。到了你这里，也得算数。''

奥斯卡望了她一眼，额角仍有汗，嘴角却轻轻动了一下。

``记着了。''

他重新看向那头石甲獾，将手掌放得更稳。

这次，双方都慢下来。

先是一线回应，随后才是一点能够接住的魂力。石甲獾将力量收在自己的主意识周围，等奥斯卡运转过一周，再将下一段缓缓接上。

院中没有人催。

Smile Lee看着阵盘，适时将几处过于敏感的记录刻度分开。大长老的手停在阵纹外，维持外围波动平稳，却让那道共生联系仍由双方自己接续。

奥斯卡体内，第一处滞涩终于慢慢通了。

一道温厚的力量顺着手臂向内流转，经过胸腹，再重新回到两者之间。那头石甲獾伏在地上的身体轻轻一颤，背甲间的魂影，露出了模糊的轮廓。

奥斯卡心里猛地一紧。

那双眼睛还在望他。

古月娜的声音及时传来。

``它还在回应。''

奥斯卡接住那点回应，没有去攥住将要显出的魂影。他依着已经走通的那一段，一次又一次，将自己的魂力送稳。

灰黑色的背甲，微白的颈毛，还有一对厚实的前爪，渐渐从模糊的光里显出来。

碧姬守着原来的身体，直到最后那一点完整的主意识也稳稳接入，才慢慢收回手。

老獾伏在地上的身躯安静下去。

而那道魂影，仍留在奥斯卡面前。

它低头望了望自己透着光的前爪，停了一会儿，随后试着向前，轻轻碰了一下他的手。

奥斯卡掌心一热。

他清楚地感到了，那道与自己不同、却已经能够相互回应的意识。

``接上了。''

他说得很轻，眼睛却始终没有移开。

院外有人向前迈了一步，又立刻收住。卢本韦将那一刻记下，刚要抬头，阵盘上的光却再次亮了起来。

这一次，亮起的不只最初那一道。

厚重的魂力沿着已接稳的联系，一层层送入奥斯卡体内。他衣襟无风而动，原有的魂环接连亮起，掌心那道香肠武魂的光，变得前所未有地凝实。

梅天狸的目光微微一凝。

``魂力在突破。''

奥斯卡也察觉了。

那几处他在山阶上反复梳理、回到大陆后仍只能一点点磨开的关口，此刻被温厚而充沛的力量接连推开。过去积在体内、尚未完全续上的运转，也随着这一轮共生，逐渐连成了更深的一周。

一次。

又一次。

他本以为刚接稳的那一层已经足够，下一道增长却仍顺着前面的路，稳稳涌了上来。

石甲獾的魂影渐渐清楚，维系自身的力量仍留在它那里。它能感到奥斯卡每一处新的承接，便也一点点调整送来的分量。

奥斯卡没有再独自硬压，哪里过重，就缓一缓；哪一段已经走稳，便让那股力量继续向前。

许久，院中扬起的细尘终于缓缓落下。

最后一处转折接稳，奥斯卡体内原本翻涌的魂力，也终于停在新的高度。

他睁开眼，掌心浮起一列清楚的纹路。

魂力：

\NovelDisplay{M05793R01}{%
(0)(1,1,1)(2,2,1)(3,3)%
}

蔡兴坤望着那几列，手里的笔停住了。

卢本韦同样看清了。

她熟悉奥斯卡先前的气息，也知道他每回战斗以后，为了让整支队伍多撑一会儿，究竟要怎样将自己的消耗反复摊开。

如今，这份魂力仍是他的，却比先前厚实了许多。那道过去颇为明显的差距，已经随着接连的突破，迅速缩短。

``你快追上我们了。''

她说完，才发现自己仍扶着他的手腕，赶紧稍稍松开。

奥斯卡低头看过，又望向蔡须坤。

蔡须坤已经站了起来，眼中有很清楚的笑意。

``真的很近了。''

她与卢本韦最先通过第一考时，魂力已经走到更前面的那一段。奥斯卡此前仍落后一些，如今再感受他的气息，已经不能像从前那样，轻易分出一道明显的远近。

奥斯卡张了张口，似乎也想笑，却先长长呼出一口气。

``我刚才还以为，它只要肯留下，便已经很好了。''

他的目光落到脚边。

石甲獾的魂影正慢慢转身，像是在重新熟悉自己。它迈步时，已不用再等那副衰老的身体将重量一点点挪过来，前爪落下，魂息便随着联系轻轻回应。

它仍记得那块平石。

走了两步，便在石边伏下，抬头望着奥斯卡。

他缓缓俯身，将手放到它面前。

一段很长的记忆掠过心间。林中清晨的露水，石下阴凉的泥土，还有最后这些日子，身体越来越重，抬头却还能看见有人坐在旁边。

奥斯卡的手停了许久。

``以后还能一起走。''他说。

古月娜望着那头獾，片刻后，轻轻点头。

``它也是这个意思。''

碧姬终于将一直绷着的肩背松开些，示意同来的族群成员，妥善安置老獾留下的身体。

帝天则低头看向新成的魂灵。

``让它自己试一试。''

奥斯卡听明白了，没有催动它，只安静等着。

石甲獾先显出更多轮廓，随后又试着收近。第二次，它便将魂影缓缓隐去，那点属于它的回应，却仍清楚地留在奥斯卡心中。

过了一会儿，它又自己出来了。

它在平石旁绕了半圈，望向门外的光，最后重新伏下。

灭世狼王站在近旁，看着这一幕，耳朵轻轻动了一下。

卢本韦将结果认真写下。

奥斯卡，七万年石甲獾，魂灵融合成功。

下面是他新达到的魂力境界。

她写完，抬头看向phyrion，发现这位长辈也正望着纸上的那一行字。

``这几日，接着熟悉。''phyrion说，``刚走通的地方，要学会怎样让它一直转稳。''

奥斯卡点头。

``我们一起。''

他说这句时，掌心又接住了一点温厚的回应。

接下来的几日，小队轮流陪他调息。

第一夜，奥斯卡醒过两回。一次是刚增加的魂力转得太急，另一次，却是自己在睡梦中仍习惯性地将力量往外送，才把石甲獾惊醒。

它收了收魂息，他立刻跟着停住，半晌才反应过来，自己已经不必再像替整支队伍补给时那样，时刻把每一分力量都推到外面。

``它想让你也歇一会儿。''古月娜听过回应，低声转述。

奥斯卡靠着床沿，慢慢点头。

蔡须坤坐在桌边，帮他把那一段运转重新理过。卢本韦将温水递过去，见他还下意识想先问旁人要不要喝，便直接塞进了他手里。

``先喝你的。''

奥斯卡捧着杯子，终于笑了。

阵盘记下的细光，渐渐平稳。

新的魂力没有继续无止境地往上冲。随着几日调息，奥斯卡已能将它收放在自己想留的位置，连制作最寻常的恢复香肠，也比从前稳了许多。

他没有急着给新魂灵定下要做什么。石甲獾想出来，就让它到水槽边看看；想收回，他也等着那点魂息慢慢安静下去。

到了第五日，卢本韦比平时来得更早。

天还未亮，她已将旧笔记放在桌边，挨着记录晶片坐下。

第五天。

卢守衡与八万年柔骨兔，当年正是在这一天失衡。

她一页页看着新记录。最初魂力过重的那一刻，后来怎样缓下，哪一道接续先走通，奥斯卡又怎样将接连突破的力量稳住，都清楚留在上面。

phyrion进来时，见她仍坐着，便在旁边停下。

``还没放心？''

卢本韦轻轻点头。

门内，奥斯卡正慢慢调息。石甲獾伏在平石旁，片刻以后，自己抬起头，收回魂影，重新回到那道共生联系之中。

他的呼吸没有乱。

卢本韦等过天亮，又等过午后。直到那一日真正过去，才在日期旁补下那一行字。

第五日，魂力平稳，共生联系无异常。

她写完，手指仍停在纸边。奥斯卡走出来时，看见她的模样，也没有急着打趣，只在旁边坐下，陪她将最后一页看完。

第七日，奥斯卡带着新魂灵，重新走进临时会场。

两位亲王与星斗代表都在。那头石甲獾先伏在他身侧，望了一会儿，才慢慢向有光的地方挪了半步。

卢本韦把这七日的记录依次铺开，Smile Lee取下晶片，将最初过急的那一道波动，与后来趋于平稳的几段一并放给众人看。

有人留意到奥斯卡如今的魂力，目光久久停在他身上。

七万年魂灵带来的增长，真实得无需多作说明。可碧姬先看的，却是那头石甲獾的眼睛。

古月娜向它送去询问，很快便接到一段清楚的回应。

林间的旧巢，熟悉的族群，还有它自己想要保留的记忆，都没有随着转变消失。

奥斯卡也听见了，低头朝它笑了一下。

``它想等塔建好，再回来看看。''

帝天望着那道魂影，慢慢点头。

万妖王看过记录，又望向它自己选的位置，沉默许久，终于将纸递给旁边的代表。

test\_alpha0接过最后一页，认真看完。

``这次融合成功了。''他说，``你们留下的这些经过，往后还要接着记。''

奥斯卡点头，手掌轻轻落到石甲獾面前。

``我会的。''

卢本韦将另一份纸递给蔡兴坤。

``给雨浩也留一份。他回来以后，得看见我们怎样走过来的。''

蔡兴坤应下，将用过的纸按日期重新排好。

蔡须坤望着奥斯卡，再一次清楚感到那份已经逼近自己与卢本韦的魂力。她走过去，看见他正在替石甲獾挪动平石，便伸手帮他扶稳。

那头獾低下头，轻轻碰了碰石面。

奥斯卡掌心随之传来一点温厚的满足。

蔡须坤笑了笑，手却仍在石边停了一会儿。

她心口那缕魂息，还安静留在瑞气深处。魂灵的成功，不能替牢大补回失去的肉身与道基，乾坤造化丹的线索，她也仍记着。

可今日，在他们亲手留下的记录里，又多了一头魂兽继续往前走的经历。

她将平石放稳，抬头看向同伴。

``下一段路，也一起走。''

奥斯卡望着她，认真点头。

``一起。''

\chapter{第一百四十二章：临行之赠}\label{chapter:142}

第二天，奥斯卡起得比院中的工匠还早。

那块平石还留在水槽旁。石甲獾的魂影伏在上面，听见门响，缓缓抬起头。他走过去，没有立刻催动魂力，只将手停在它面前，等那一点熟悉的回应接过来。

昨夜没有再醒。

新增加的魂力沿着已经走熟的地方转过一周，回到掌心时，仍然平稳。他让石甲獾自己收回魂影，等了一会儿，又同它一起将那道联系放开。

这一次，也没有滞涩。

卢本韦从廊下走来，手里拿着最后一页观察记录。

``又试了一遍？''

``它醒得早。''奥斯卡笑了笑，``我也睡不住了。''

她走近，先看过他的脸色，才将纸递过去。最下面留了一处空白，等着他记下今晨的情况。

奥斯卡接过笔，俯在石桌上写完。石甲獾伏在他的脚边，看着那支笔来回移动，前爪轻轻蜷了一下。

``这一份留在这里。''卢本韦说，``还有一份给雨浩。路上的变化，我们接着记。''

奥斯卡点头，把纸角抚平。

他已经能在抬手之间察觉出自己增长的力量。可看着这张记满七日经过的纸，他最先想起的，仍是刚开始时，石甲獾那只迟迟没有收回去的前爪。

那时候，彼此都还不知道下一步会怎样。

如今他们要一起出远门了。

院子另一头传来蔡兴坤的声音。

``奥斯卡，你的那只袋子放哪儿？''

``外间，桌子底下。''

``那里面怎么还留着半卷擦阵盘的布？''

奥斯卡一怔，赶紧站起来。卢本韦已经笑了，替他将笔接了过去。

``去收拾吧。再迟些，兴坤要把整个院子替你带上了。''

外间的长桌上，几人的行囊已经摆开。

蔡兴坤把那半卷布放回原处，又将用过的记录晶片逐块归好。黑子鸡站在椅背上，低头盯着一条垂下来的肩带，她伸手挡了一下，它才收住将要探过去的喙。

古月娜拿来一只旧封套，放到蔡须坤面前。

里面是从养猪场抄出的问情谷手记。纸角被反复翻过，已有些发软，断缺的位置仍空着，没有人往上补字。

蔡须坤将它取出，与新找来的舆图并排铺好。

落日森林。

西北，八十余里。

她的手指停在那几行旧字旁。

``这里还是没写，从哪一处算起。''奥斯卡走过来，俯身看了一会儿。

``也可能是前面那页缺掉了。''蔡兴坤说。

梅天狸正从门外进来，听见这句，便在桌边停下。

``所以先到落日森林，再找能同手记对上的地方。''

她将舆图往近处移了移，指过沿途几座城镇，把几人查到的道路接在一起。

``这一路先按官道走，到了山路再看地形。旧记录里的水口、石壁，未必还留着原来的样子，进林以后，一处处认。''

蔡须坤点头，又看了一遍手记上的空白。

乾坤造化丹的线索，总算不再只是一句话。可从纸上这几行字，到真正找到它，中间仍隔着很长一段路。天地灵水也还没有可靠的下落，她将记着这件事的另一张纸夹好，没有因为即将动身，便将它放到脑后。

古月娜伸出手，替她压住被风掀起的一角。

``娜儿，等到了林边，你也别一个人先走。''

古月娜抬起头。

``我知道。''

蔡须坤望着她，神色仍很认真。

古月娜便又添了一句：``要探路，也让你看得见。''

梅天狸听着，将那份手记重新叠好。

``都一样。遇上不认识的东西，先叫人，别各自去试。''

她没有收走抄本，仍递回蔡须坤手里。

几人正要接着商量，门外忽然传来一声很轻的金属相碰声。

Smile Lee站在门边，袖口挽着，指间夹着一枚尚未收起的薄片。蔡兴坤赶紧腾出桌面，奥斯卡则把差点挡住门口的袋子提到一旁。

``先别装满。''Smile Lee说，``还有件东西，要带上。''

他走到桌前，将银灰色的储物魂导戒转过半圈。

一件掌心大小的魂导器落在软布上。

外壳呈暗银色，边沿收得很薄，正中嵌着一枚透亮的核心。几圈细密纹路藏在外壳里面，光从缝隙间透出，到了最深处，隐约有一点蓝。

紧接着，桌上又多了六枚小扣片。

它们比主件轻得多，背面留着细小的扣环，能够固定在衣物内侧。卢本韦原本以为，这是平常用来接续阵纹的部件，等那六枚扣片同时亮起，才察觉桌边的魂力已经轻轻沉了下去。

``你们此行凶险，务必把这件带上。''Smile Lee将主件翻过来，``防护用的，十级魂导器。''

奥斯卡伸到一半的手停住了。

``十级？''

蔡兴坤也低下头，重新看向那枚透亮的核心。

他们见过Smile Lee布阵，也亲眼看着他用几块薄晶片，记下奥斯卡与石甲獾之间最细微的变化。可眼前这件东西才刚从储物戒里取出，竟已让桌边原本散乱的几缕魂力，变得如此安静。

梅天狸走近了一些。

``你将它做成了？''

``这些年一直在改，最近才合上最后一层。''

Smile Lee将手指停在主件内侧，并没有催动最深处的力量。

``里面有两处接续，我请大长老帮着稳过。其余的阵列，一层也省不得。为了让它们接起来，我拆过的外壳，够摆满几张这样的桌子。''

卢本韦抬头望向他。

Smile Lee却已经拿起一枚扣片，递给离自己最近的蔡兴坤。

``先学会用。它能做什么，你们得自己摸清。''

六枚扣片被分到五人与梅天狸手中，主件则仍放在桌上。

``这些是引导扣。防护展开时，会顺着你们留在里面的魂力，把人一同接进去。别扣在容易掉的地方。''

蔡须坤仔细看过那几道细纹，才开口问：``以我们现在的魂力，能将它撑起来吗？''

``核心里的力量，我已经先存好。你们负责引动，不必临时把它灌满。''

Smile Lee将其中一枚翻过来，露出内侧稍稍凸起的一道纹。

``从这里进去，慢一点。先记住它怎样回应。''

一线极细的魂力沿纹路落下。

扣片边缘泛起银光，桌上的主件也随之亮了一下。六枚小扣之间，像有极细的光接了过去，片刻后，又被他收回。

蔡须坤照着试过一次。

魂力走到转折处时，她本能地想多送一点，Smile Lee的手指便停在了旁边。

``够了。''

她收住力量，等那一点银光自行亮稳，才察觉自己掌心仍留着很清楚的余地。

``触发用平常的魂力。''他又道，``你的混沌之力，先收好。''

蔡须坤认真点头。

其余几人也依次试了一遍。卢本韦习惯在魂力接通以后继续往前推，才多送出一点，扣片便轻轻热起来。她立刻停下，重新收细，直到那份温热缓缓退去。

奥斯卡看见，没忍住笑了一下。

``刚才你还笑我东西没收好。''

卢本韦将扣片扣稳，抬眼看他。

``轮到你了。''

奥斯卡立刻收住笑，老老实实接过主件，先问清背面的扣环该怎样打开。

等六个人都留下了自己的魂力，Smile Lee才让他们走到院中。

廊下的桌椅被挪开。蔡须坤将主件固定在衣内，另一只手按住引导扣，试着将那一线力量接过去。

银光先从她身侧浮起。

随后，离她最近的古月娜身上，也亮起了相同的光。光沿着几人的位置缓缓展开，越过奥斯卡与卢本韦之间的空隙，最后在上方合拢。

院中的风声忽然远了。

蔡兴坤抬起头，看见头顶是一层极淡的光。外面的树影仍然清楚，叶子被风吹过，却像隔着一段平静的水面。

奥斯卡低头望向脚边。

石甲獾方才显出的魂影，也一同留在了里面。它有些疑惑地抬起前爪，等奥斯卡将那点感受慢慢传过去，才重新伏好。

Smile Lee站在光幕外，将掌心按了上来。

他的手没有立刻用力。

一道压迫却已从触到的位置传开，光幕向内收了一点，蔡须坤肩头微微一紧，下意识看向了那只手。

光在她眼前缓缓流动。

她能看见受力的地方，却没有像自己撑起防护时那样，被那股分量直接压住手臂。主件里的力量沿内层阵纹接续，身旁几人的银光也仍然亮着。

片刻后，Smile Lee收回手。

光幕重新展平。

``试着往门边走。''他说。

蔡须坤先迈出一步。古月娜跟着她，卢本韦则回头照看奥斯卡与蔡兴坤。那层光随主件缓缓移动，到了门前，仍将几人稳稳兜在里面。

梅天狸走在最后，留意过每一处转折，才抬手示意他们停下。

``记住刚才的位置。真出了事，先靠过来，别站在外头问谁受了伤。''

蔡须坤收回引动的魂力。

光从上方退下，院中的风声重新清楚起来。她隔着衣料，摸到了主件微凉的外壳，心里那一点惊讶却没有因此退去。

Smile Lee看着五人，声音比方才低了一些。

``就算胡赞亲自来了，它也能撑上一会儿。''

院中静了下来。

卢本韦原本正要将衣扣重新系好，听见这个名字，手便停住了。

蔡须坤也没有马上说话。

她知道，自己如今的力量已经同从前很不一样。可那个名字后面，还有牢大留在身后的那一日，还有许多他们至今没能真正看清的东西。

Smile Lee在主件外沿轻轻点了一下。

``撑出来的时间，要拿来脱身。核心的光往下退时，就该准备走，别等它全暗。''

``若是同伴还在外面呢？''蔡须坤问。

``先把人接近。主件在你这里，你要看清他们的位置。''

他看了一眼六人身上的引导扣，又指向方才光幕合拢的地方。

``所以才让你们都试。临到那时候，别只剩戴着它的人知道怎么做。''

蔡须坤垂下目光，将方才每个人站在哪里，又在心里过了一遍。

梅天狸这时走到她身旁。

``还有我。''

她抬起手，衣袖下那枚引导扣，仍留着一线很淡的银光。

Smile Lee转头望着她，点了一下头。

``小狸会跟你们一块去。这一路，交给她。''

奥斯卡望向梅天狸，肩背终于松开了一些。

她却没有像在训练场上那样，笑着让他们再来试两招，只将衣袖放下，认真看过五人。

``我会护着你们。真遇上麻烦，要退的时候，听我说的方向，谁也别掉队。''

蔡兴坤点头，手掌从黑子鸡背上轻轻抚过。

``那探路的时候，我也先同您说一声。''

``嗯。能看见多远、什么时候回来，都说清。''

卢本韦将自己那枚引导扣又摸了一遍。

她记得第一次同梅天狸切磋时，几个人怎样将各自的力量凑在一起，也记得前些日子那次真正让长老退后的合击。如今，梅天狸仍然站在他们身边，面对要走的这段路，却已将他们当作必须彼此照看的同行人。

她把手放下，郑重应了一声。

``我们记住。''

回到屋里，蔡须坤取下主件，重新看过背面的扣环，才将它固定在最稳妥的位置。

Smile Lee等她系好外衣，又让六人将引导扣都走过一遍，确认回应还在，才收起桌上剩余的薄片。

``多谢前辈。''蔡须坤说。

另外四人也一同行礼。

他摆了摆手，目光在奥斯卡身上多停了一会儿。

``魂灵的记录，路上接着写。哪怕只是夜里多醒了一回，也别嫌小。''

``我会带回来。''奥斯卡答。

Smile Lee点头，转身出了门。院外很快又响起工匠喊他的声音，有一处旧阵纹，还等着他去断开。

夜里，几人将东西重新理好。

奥斯卡没有一口气把魂力全用在准备香肠上。他先备下赶路时便于取用的几份，剩下的等途中再做。卢本韦替他将其中一袋递给梅天狸，回来时，手里又多了一页已经核过的记录。

蔡兴坤坐在灯下，把给蓬莱岛的信写完。

魂灵试验成功、奥斯卡如今的情况，还有接下来前往落日森林的安排，都被她逐项写清。提到尚在岛上的同伴，她停了停，又在末尾添了一句，让他们安心完成第一考，有新的消息，便往总殿传。

五人看过，才将信封好，连同留给霍雨浩的那份记录，一并交给殿中传信的人。

蔡须坤回房后，屋里只剩下一盏小灯。

她没有急着躺下，而是将心神缓缓收近，先理顺自己的魂力，再将温润的瑞气送到那缕被护着的魂息旁。

牢大的回应很轻。

她等了片刻，才在心里同他说：``我们要出发了。''

那一点魂息没有立刻动。过了一会儿，才有极微弱的熟悉感，顺着她停稳的力量，轻轻触了过来。

蔡须坤的手在心口停住。

她没有继续将他唤出来，只把这一日学会怎样使用魂导器、梅天狸也会同行的事，慢慢说给他听。

``你再等一等。''

说完，她没有再催动更多力量。灯下的呼吸一点点平稳，那缕魂息也重新安静下去。

第二日清早，phyrion来了。

卢本韦正在门边束紧行囊，听见动静，立刻转过身。

``干爷爷。''

他仍穿着紫衣，小白猪跟在脚边，到了门槛前，慢悠悠地抬腿跨过。

``都好了？''

``好了。''

卢本韦答过，却又低头，将已经束紧的肩带理了一下。phyrion看见，走近替她把压在下面的一角衣料抽出来。

``别怕。你是养猪场的继承人，更是猪神Opt的继承者，猪神大人会保佑你们的。''

她慢慢抬起头。

这几日，在那张铺着旧笔记的桌边，她听过许多先前不知道的往事，也终于亲眼看见，一头魂兽能以另一种方式，将自己完整地留在同伴身旁。

她伸手探进衣内，将猪神令取出来。

令牌的分量仍同从前一样。她却看了很久，才重新合拢手指。

``我不会辜负养猪场的。''

phyrion望着她，轻轻点头。

``那就照顾好自己，也照顾好同伴。等你们回来，塔这边还有许多事要一块儿做。''

小白猪走近，拿额头拱了一下她的小腿。卢本韦蹲下摸过它的头，站起来时，已将猪神令重新放好。

门外，梅天狸正等着他们。

几人走过馆旁的院落，将最后一份记录交给留下的人。碧姬又看过石甲獾的状态，才将手收回。那道魂影伏在奥斯卡身侧，目光在熟悉的平石上停了一会儿，随后自己收进了共生联系之中。

古月娜站在廊下，向帝天点了点头。

``族长，我们走了。''

帝天看过五人，又望向梅天狸。

``路上小心。''

test\_alpha0站在门边，将蔡须坤取出来确认的旧封套看了一眼。

``到了地方，看见什么，先同手记对一对。对不上的地方，留着疑问，别急着往里补。''

蔡须坤将封套收稳，认真应下。

城西已经响起了第一阵锤声。

从石阶上望下去，旧雕像旁的脚手架又多了一层，运送石料的人沿新留出的通道慢慢走着。几人走到路口，回头望过一眼，才随着梅天狸，往城门方向去。

出了城，脚下渐渐换成了结实的土路。

起初两日，沿途车马还多。蔡须坤听着身边的脚步，总会不自觉地去确认心口那缕魂息。到后来，不用停步，也已能将它同自己的呼吸一并照看。

奥斯卡却需要更留心一些。

新增加的魂力用在步行上，起先比他想得顺畅。他一日走下来，甚至比从前轻松。可到了夜里，等人真正坐稳，那些被赶路带快的接续，才一点点显出不同来。

第三日傍晚，他刚要替几人准备恢复香肠，掌心的光忽然慢了一瞬。

卢本韦看见，伸手接过他背上的袋子。

``先坐。''

``只是刚才走得快了些。''

``那就把刚才这一段转稳。''

梅天狸在旁边停下，让其余人也在背风处歇脚。古月娜坐近，听过石甲獾送来的回应，才将一根细枝放到奥斯卡面前。

她在土上画出那一处滞住的位置。

``它跟着你收得太急，后面还留着一小段。先别往外送，等它接回来。''

奥斯卡照着慢慢调过，掌心的光才重新连上。他低头看了看自己的手，随后翻开随身带着的小册子，将这件事记了进去。

第二天再走，他便没有为了赶上前面的人，将每一步都催到同样快。

蔡兴坤发现以后，也将探路的距离收近了一些。黑子鸡飞出去，绕过前面的坡顶便回来，她把看见的路告诉众人，再跟着一同往前走。

没有人催。

梅天狸有时走在前面，有时落到队伍后方。遇上视线被遮住的弯路，她便先过去看一眼，等五人走近，才接着动身。

第五日下了一场雨。

雨来时，他们刚走进一段山路，路旁的沟很快积满了水。几人没有去挤附近那处已经松动的土坡，跟着梅天狸转到背风的石壁下，等最急的一阵先过去。

奥斯卡将包裹挪到干处，石甲獾的魂影也在脚边显了出来。

它伸出前爪，停在岩石边缘，望着雨水沿裂缝淌下去。那一点传回来的感觉，让奥斯卡忽然安静了一会儿。

``怎么了？''蔡兴坤问。

``它以前也待过这样的地方。''

他低头望着那道魂影。

这一路，它并不认识那些陌生的岔口，也不能替他们指出问情谷在哪里。可石壁里透出的凉意，雨落进泥土时的气息，却让它记起了从前在林间躲雨的山洞。那副衰老的身体，已经许久不能让它走到这样的地方。

蔡兴坤蹲下来，顺着它看的地方，望了一会儿。

``那这一路，能看不少。''

雨势稍缓，几人重新上路。

傍晚歇下时，梅天狸又让他们把引导扣各走过一次。魂力只接通到能认清回应的位置，便被收住，没有耗用核心中的防护力量。

蔡须坤在几人之间挪过两步，确认自己闭上眼，也能辨出同伴靠近的位置。卢本韦则把扣片换到更稳的一处，免得肩带来回蹭到它。

第八日，官道上的车辙渐渐淡了。

舆图被翻得越来越勤。山脚的旧路有一段已经被水冲坏，梅天狸带着几人绕过低处，在远一些的地方重新接上。每回停下认路，蔡须坤都会将手记取出来，核过之后，再原样放回。

那些仍对不上的字，她没有因为一路平安，就看得轻些。

古月娜走在她身旁，见她又将卷起的纸角抚平，便伸手接过舆图。

``我拿一会儿。''

蔡须坤松开手，将封套收好。

``娜儿，你累不累？''

古月娜摇头，看了一眼落在后面的奥斯卡与卢本韦。

``前面有平地，我们等一等。''

两人在路边停下。蔡兴坤也抱着黑子鸡走回来，没多久，后面的脚步声便接近了。

十级魂导器一直安稳地贴在蔡须坤衣内。

他们没有再将整层防护展开，可每逢整理衣物，那枚引导扣都能让她想起院中光幕合拢时，同伴站在自己身旁的距离。

到了第十三日，远处的树色终于连成了一片。

最初只是一道压在地平线上的深绿，走近以后，才渐渐看清高低交错的树冠。风从林中吹来，带着潮湿的叶气，与他们这些日子走过的开阔山道，很不一样。

古月娜在林边停下。

这里的气息同星斗大森林并不相同。她没有立刻往深处探，只将近旁的树木与地面看过，又回头望向身后的同伴。

蔡须坤也停住脚步。

从天州出发，整整十三天，他们终于走到了落日森林。

梅天狸上前，看过林缘一段被灌木遮住的旧径，才示意五人靠近。

``进去以后，先跟着我走。''

蔡兴坤收住黑子鸡，奥斯卡将记录册放回袋中。卢本韦确认同伴都已跟上，才转向那一片渐渐暗下来的林木。

蔡须坤摸了摸衣内的主件，又将心口那缕魂息轻轻护稳。

她与古月娜并肩跨过林缘。

脚下的土路，被柔软的落叶接住。前方，梅天狸拨开垂下的枝条，停在那里，等他们一个个走近。

六人的身影，终于一同进入了落日森林。

\backmatter
\end{document}

% VS Code XeLaTeX build verification.

% VS Code build check with ASCII compilation path.

% Trigger desktop VS Code build after project load.
